% Options for packages loaded elsewhere
\PassOptionsToPackage{unicode}{hyperref}
\PassOptionsToPackage{hyphens}{url}
\documentclass[
]{article}
\usepackage{xcolor}
\usepackage{amsmath,amssymb}
\setcounter{secnumdepth}{5}
\usepackage{iftex}
\ifPDFTeX
  \usepackage[T1]{fontenc}
  \usepackage[utf8]{inputenc}
  \usepackage{textcomp} % provide euro and other symbols
\else % if luatex or xetex
  \usepackage{unicode-math} % this also loads fontspec
  \defaultfontfeatures{Scale=MatchLowercase}
  \defaultfontfeatures[\rmfamily]{Ligatures=TeX,Scale=1}
\fi
% \usepackage{lmodern}
\ifPDFTeX\else
  % xetex/luatex font selection
\fi
% Use upquote if available, for straight quotes in verbatim environments
\IfFileExists{upquote.sty}{\usepackage{upquote}}{}
\IfFileExists{microtype.sty}{% use microtype if available
  \usepackage[]{microtype}
  \UseMicrotypeSet[protrusion]{basicmath} % disable protrusion for tt fonts
}{}
\makeatletter
\@ifundefined{KOMAClassName}{% if non-KOMA class
  \IfFileExists{parskip.sty}{%
    \usepackage{parskip}
  }{% else
    \setlength{\parindent}{0pt}
    \setlength{\parskip}{6pt plus 2pt minus 1pt}}
}{% if KOMA class
  \KOMAoptions{parskip=half}}
\makeatother
\usepackage{color}
\usepackage{fancyvrb}
\newcommand{\VerbBar}{|}
\newcommand{\VERB}{\Verb[commandchars=\\\{\}]}
\DefineVerbatimEnvironment{Highlighting}{Verbatim}{commandchars=\\\{\}}
% Add ',fontsize=\small' for more characters per line
\newenvironment{Shaded}{}{}
\newcommand{\AlertTok}[1]{\textcolor[rgb]{1.00,0.00,0.00}{\textbf{#1}}}
\newcommand{\AnnotationTok}[1]{\textcolor[rgb]{0.38,0.63,0.69}{\textbf{\textit{#1}}}}
\newcommand{\AttributeTok}[1]{\textcolor[rgb]{0.49,0.56,0.16}{#1}}
\newcommand{\BaseNTok}[1]{\textcolor[rgb]{0.25,0.63,0.44}{#1}}
\newcommand{\BuiltInTok}[1]{\textcolor[rgb]{0.00,0.50,0.00}{#1}}
\newcommand{\CharTok}[1]{\textcolor[rgb]{0.25,0.44,0.63}{#1}}
\newcommand{\CommentTok}[1]{\textcolor[rgb]{0.38,0.63,0.69}{\textit{#1}}}
\newcommand{\CommentVarTok}[1]{\textcolor[rgb]{0.38,0.63,0.69}{\textbf{\textit{#1}}}}
\newcommand{\ConstantTok}[1]{\textcolor[rgb]{0.53,0.00,0.00}{#1}}
\newcommand{\ControlFlowTok}[1]{\textcolor[rgb]{0.00,0.44,0.13}{\textbf{#1}}}
\newcommand{\DataTypeTok}[1]{\textcolor[rgb]{0.56,0.13,0.00}{#1}}
\newcommand{\DecValTok}[1]{\textcolor[rgb]{0.25,0.63,0.44}{#1}}
\newcommand{\DocumentationTok}[1]{\textcolor[rgb]{0.73,0.13,0.13}{\textit{#1}}}
\newcommand{\ErrorTok}[1]{\textcolor[rgb]{1.00,0.00,0.00}{\textbf{#1}}}
\newcommand{\ExtensionTok}[1]{#1}
\newcommand{\FloatTok}[1]{\textcolor[rgb]{0.25,0.63,0.44}{#1}}
\newcommand{\FunctionTok}[1]{\textcolor[rgb]{0.02,0.16,0.49}{#1}}
\newcommand{\ImportTok}[1]{\textcolor[rgb]{0.00,0.50,0.00}{\textbf{#1}}}
\newcommand{\InformationTok}[1]{\textcolor[rgb]{0.38,0.63,0.69}{\textbf{\textit{#1}}}}
\newcommand{\KeywordTok}[1]{\textcolor[rgb]{0.00,0.44,0.13}{\textbf{#1}}}
\newcommand{\NormalTok}[1]{#1}
\newcommand{\OperatorTok}[1]{\textcolor[rgb]{0.40,0.40,0.40}{#1}}
\newcommand{\OtherTok}[1]{\textcolor[rgb]{0.00,0.44,0.13}{#1}}
\newcommand{\PreprocessorTok}[1]{\textcolor[rgb]{0.74,0.48,0.00}{#1}}
\newcommand{\RegionMarkerTok}[1]{#1}
\newcommand{\SpecialCharTok}[1]{\textcolor[rgb]{0.25,0.44,0.63}{#1}}
\newcommand{\SpecialStringTok}[1]{\textcolor[rgb]{0.73,0.40,0.53}{#1}}
\newcommand{\StringTok}[1]{\textcolor[rgb]{0.25,0.44,0.63}{#1}}
\newcommand{\VariableTok}[1]{\textcolor[rgb]{0.10,0.09,0.49}{#1}}
\newcommand{\VerbatimStringTok}[1]{\textcolor[rgb]{0.25,0.44,0.63}{#1}}
\newcommand{\WarningTok}[1]{\textcolor[rgb]{0.38,0.63,0.69}{\textbf{\textit{#1}}}}
\usepackage{longtable,booktabs,array}
\newcounter{none} % for unnumbered tables
\usepackage{calc} % for calculating minipage widths
% Correct order of tables after \paragraph or \subparagraph
\usepackage{etoolbox}
\makeatletter
\patchcmd\longtable{\par}{\if@noskipsec\mbox{}\fi\par}{}{}
\makeatother
% Allow footnotes in longtable head/foot
\IfFileExists{footnotehyper.sty}{\usepackage{footnotehyper}}{\usepackage{footnote}}
\makesavenoteenv{longtable}
\setlength{\emergencystretch}{3em} % prevent overfull lines
\providecommand{\tightlist}{%
  \setlength{\itemsep}{0pt}\setlength{\parskip}{0pt}}
\usepackage[]{natbib}
\bibliographystyle{plainnat}
\makeatletter
\@ifpackageloaded{subcaption}{}{\usepackage{subcaption}}
\@ifpackageloaded{caption}{}{\usepackage{caption}}
\captionsetup[subfigure]{margin=0.5em}
\AtBeginDocument{%
\renewcommand*\figurename{Figure}
\renewcommand*\tablename{Table}
}
\AtBeginDocument{%
\renewcommand*\listfigurename{List of Figures}
\renewcommand*\listtablename{List of Tables}
}
\newcounter{pandoccrossref@subfigures@footnote@counter}
\newenvironment{pandoccrossrefsubfigures}{%
\setcounter{pandoccrossref@subfigures@footnote@counter}{0}
\begin{figure}\centering%
\gdef\global@pandoccrossref@subfigures@footnotes{}%
\DeclareRobustCommand{\footnote}[1]{\footnotemark%
\stepcounter{pandoccrossref@subfigures@footnote@counter}%
\ifx\global@pandoccrossref@subfigures@footnotes\empty%
\gdef\global@pandoccrossref@subfigures@footnotes{{##1}}%
\else%
\g@addto@macro\global@pandoccrossref@subfigures@footnotes{, {##1}}%
\fi}}%
{\end{figure}%
\addtocounter{footnote}{-\value{pandoccrossref@subfigures@footnote@counter}}
\@for\f:=\global@pandoccrossref@subfigures@footnotes\do{\stepcounter{footnote}\footnotetext{\f}}%
\gdef\global@pandoccrossref@subfigures@footnotes{}}
\@ifpackageloaded{float}{}{\usepackage{float}}
\floatstyle{ruled}
\@ifundefined{c@chapter}{\newfloat{codelisting}{h}{lop}}{\newfloat{codelisting}{h}{lop}[chapter]}
\floatname{codelisting}{Listing}
\newcommand*\listoflistings{\listof{codelisting}{List of Listings}}
\makeatother
\usepackage{bookmark}
\IfFileExists{xurl.sty}{\usepackage{xurl}}{} % add URL line breaks if available
\urlstyle{same}
\hypersetup{
  colorlinks=true,linkcolor=red,urlcolor=red,citecolor=red,anchorcolor=red,filecolor=red,
  pdfcreator={LaTeX via pandoc}}

\author{}
\date{}

\IfFileExists{seqsplit.sty}{\usepackage{seqsplit}}{\newcommand{\seqsplit}[1]{#1}}
\protected\def\breakseq#1{\seqsplit{#1}}
\protected\def\breaktt#1{\begingroup\ttfamily\seqsplit{#1}\endgroup}
\title{Introduction to Biology\\\normalsize A Generative Approach}
\author{Daniel Ari Friedman\\\footnotesize{Active Inference Institute}\\\footnotesize{\texttt{daniel@activeinference.institute}}\\\footnotesize{\href{https://orcid.org/0000-0001-6232-9096}{ORCID: 0000-0001-6232-9096}} \\ \footnotesize{\href{https://doi.org/10.5281/zenodo.20286478}{DOI: 10.5281/zenodo.20286478}} \\ \footnotesize{2026}}
\date{2026}

% -----------------------------------------------------------------------
% PAGE GEOMETRY  (config.yaml: layout.*)
% Near-zero margins: 2 mm all sides (matches config.yaml layout.margin_*_mm).
% -----------------------------------------------------------------------
\usepackage[
  top=2mm,
  bottom=2mm,
  left=2mm,
  right=2mm,
  headheight=12pt,
  includeheadfoot
]{geometry}

% -----------------------------------------------------------------------
% HYPERLINK COLOURS — red links regardless of pandoc hidelinks.
%
% Strategy: PassOptionsToPackage fires BEFORE hyperref is loaded by Pandoc,
% so colorlinks=true wins over hidelinks.  AtBeginDocument then re-applies
% the complete hypersetup as a belt-and-braces override.
% -----------------------------------------------------------------------
\PassOptionsToPackage{colorlinks=true,urlcolor=red,linkcolor=red,citecolor=red}{hyperref}

\AtBeginDocument{%
  \hypersetup{colorlinks=true,urlcolor=red,linkcolor=red,citecolor=red,
              anchorcolor=red,filecolor=red,menucolor=red}%
}

% -----------------------------------------------------------------------
% TYPOGRAPHY  (config.yaml: typography.*)
% Body: Linux Libertine O or STIX Two Text | Heading: Linux Biolinum O or Arial
% -----------------------------------------------------------------------
\usepackage{fontspec}
\IfFontExistsTF{Linux Libertine O}{\setmainfont{Linux Libertine O}}{%
  \IfFontExistsTF{STIX Two Text}{\setmainfont{STIX Two Text}}{}%
}
\IfFontExistsTF{Linux Biolinum O}{\setsansfont{Linux Biolinum O}}{%
  \IfFontExistsTF{Arial}{\setsansfont{Arial}}{}%
}
\IfFontExistsTF{DejaVu Sans Mono}{\setmonofont[Scale=0.88]{DejaVu Sans Mono}}{%
  \IfFontExistsTF{Menlo}{\setmonofont[Scale=0.82]{Menlo}}{}%
}
\usepackage{microtype} % Enhances typographic quality (kerning, tracking)
\emergencystretch=8em
\sloppy
\raggedbottom

% Text-macro fallbacks for symbols that Pandoc may emit from code or bookmarks.
% Literal Unicode cleanup is handled upstream by the PDF remapper so the preamble
% itself does not introduce missing-glyph warnings.
\newcommand{\TextSafePipe}{\rule[-0.15ex]{0.35pt}{1.2ex}}
\renewcommand{\textasciitilde}{about }
\renewcommand{\textbar}{\TextSafePipe}
\providecommand{\VerbBar}{|}
\renewcommand{\VerbBar}{\TextSafePipe}

% Body size 9/10.8 pt (matches config.yaml typography.base_font_size_pt)
\makeatletter
\renewcommand{\normalsize}{%
  \@setfontsize\normalsize{9}{10.8}%
}
\makeatother
\normalsize

% -----------------------------------------------------------------------
% MATHEMATICS  (equation numbering with \eqref support)
% -----------------------------------------------------------------------
\usepackage{amsmath}
\usepackage{amssymb}
\IfFontExistsTF{STIX Two Math}{\setmathfont{STIX Two Math}}{}

% -----------------------------------------------------------------------
% CROSS-REFERENCES  (cleveref — must load AFTER hyperref+amsmath)
% Enables \cref{sec:foo} / \Cref{fig:bar} to render as "section 3" / "Figure 4.2"
% without hard-coding the kind word. Works with every \label{sec:...},
% \label{fig:...}, \label{eq:...}, \label{tbl:...} in the manuscript.
% -----------------------------------------------------------------------
\usepackage[capitalise,noabbrev]{cleveref}
\crefname{section}{section}{sections}
\Crefname{section}{Section}{Sections}
\crefname{figure}{figure}{figures}
\Crefname{figure}{Figure}{Figures}
\crefname{table}{table}{tables}
\Crefname{table}{Table}{Tables}
\crefname{equation}{equation}{equations}
\Crefname{equation}{Equation}{Equations}

% -----------------------------------------------------------------------
% BIBLIOGRAPHY  (natbib for \citep / \citet author-year citations)
% Pandoc auto-injects \usepackage{natbib} and \bibliographystyle{plainnat}
% when citations are present, so we do NOT re-declare them here to avoid
% "Illegal, another \bibstyle command" from bibtex.
% -----------------------------------------------------------------------

% -----------------------------------------------------------------------
% GRAPHICS  (required for \includegraphics)
% -----------------------------------------------------------------------
\usepackage{graphicx}
\makeatletter
\@ifundefined{pandocbounded}{}{%
  \renewcommand*\pandocbounded[1]{%
    \sbox\pandoc@box{#1}%
    \Gscale@div\@tempa{0.72\textheight}{\dimexpr\ht\pandoc@box+\dp\pandoc@box\relax}%
    \Gscale@div\@tempb{\linewidth}{\wd\pandoc@box}%
    \ifdim\@tempb\p@<\@tempa\p@\let\@tempa\@tempb\fi%
    \ifdim\@tempa\p@<\p@\scalebox{\@tempa}{\usebox\pandoc@box}%
    \else\usebox{\pandoc@box}%
    \fi%
  }%
}
\renewcommand{\topfraction}{0.92}
\renewcommand{\bottomfraction}{0.82}
\renewcommand{\textfraction}{0.06}
\renewcommand{\floatpagefraction}{0.74}
\setcounter{topnumber}{4}
\setcounter{bottomnumber}{3}
\setcounter{totalnumber}{6}
\makeatother

% -----------------------------------------------------------------------
% SECTION HEADER STYLE — compact sans-serif headings
% -----------------------------------------------------------------------
\usepackage{titlesec}
\titleformat{\section}{\normalfont\large\bfseries\sffamily}{\thesection}{0.85em}{}
\titleformat{\subsection}{\normalfont\normalsize\bfseries\sffamily}{\thesubsection}{0.85em}{}
\titleformat{\subsubsection}{\normalfont\small\bfseries\sffamily}{\thesubsubsection}{0.85em}{}
\titlespacing*{\section}{0pt}{10pt plus 2pt minus 2pt}{4pt plus 1pt}
\titlespacing*{\subsection}{0pt}{8pt plus 2pt minus 2pt}{3pt plus 1pt}

% -----------------------------------------------------------------------
% TABLE OF CONTENTS
% -----------------------------------------------------------------------
\usepackage{tocloft}
\setcounter{tocdepth}{3}
\setcounter{secnumdepth}{3}
\setlength{\cftsecnumwidth}{3.6em}
\setlength{\cftsubsecnumwidth}{4.8em}
\setlength{\cftsubsubsecnumwidth}{6.2em}
\renewcommand{\cftsecleader}{\cftdotfill{\cftdotsep}}
\renewcommand{\cftsubsecleader}{\cftdotfill{\cftdotsep}}
\renewcommand{\cftsubsubsecleader}{\cftdotfill{\cftdotsep}}

% -----------------------------------------------------------------------
% FIGURE / TABLE CAPTIONS
% -----------------------------------------------------------------------
\usepackage[font=footnotesize,labelfont=bf,labelsep=period]{caption}

% -----------------------------------------------------------------------
% LINE SPACING  (config.yaml: layout.line_height = 1.28)
% -----------------------------------------------------------------------
\usepackage{setspace}
\setstretch{1.28}

% -----------------------------------------------------------------------
% HEADER / FOOTER
% -----------------------------------------------------------------------
\usepackage{fancyhdr}
\pagestyle{fancy}
\fancyhf{}
\fancyhead[L]{\footnotesize\sffamily\nouppercase{\leftmark}}
\fancyhead[R]{\footnotesize\sffamily Introduction to Biology}
\fancyfoot[C]{\footnotesize\thepage}
\renewcommand{\headrulewidth}{0.4pt}

\begin{document}

\begin{titlepage}
\thispagestyle{empty}
\setcounter{page}{1}
\centering
\vspace*{0.6cm}
{\Huge\sffamily\bfseries Introduction to Biology\par}
\vspace{0.35cm}
{\Large\sffamily A Generative Approach\par}
\vspace{0.8cm}
{\Large\bfseries Daniel Ari Friedman}\\[0.35em]
{\normalsize Active Inference Institute}\\[0.25em]
{\normalsize\texttt{daniel@activeinference.institute}}\\[0.25em]
{\normalsize\href{https://orcid.org/0000-0001-6232-9096}{ORCID: 0000-0001-6232-9096}}\\[0.5em]
\vfill
\includegraphics[width=0.98\textwidth,height=0.62\textheight,keepaspectratio]{\detokenize{/Users/4d/Documents/GitHub/template/projects_in_progress/biology_textbook/output/manuscript/assets/cover/biology_textbook_cover.png}}
\vfill
{\small Edition 1.0 -- 2026}
\end{titlepage}
\clearpage
\thispagestyle{empty}
\section*{Publishing Information}

\noindent{\Large\bfseries Introduction to Biology}\\
{\large A Generative Approach}\\[1.2em]
{\Large\bfseries Daniel Ari Friedman}\\[0.35em]
{\normalsize Active Inference Institute}\\[0.25em]
{\normalsize\texttt{daniel@activeinference.institute}}\\[0.25em]
{\normalsize\href{https://orcid.org/0000-0001-6232-9096}{ORCID: 0000-0001-6232-9096}}\\[0.5em]
\vspace{1.0em}
\noindent Edition 1.0 -- 2026\\
\noindent Text license: CC BY 4.0\\
\noindent Source-code license: Apache-2.0\\
\noindent DOI: \href{https://doi.org/10.5281/zenodo.20286478}{10.5281/zenodo.20286478}\\
\noindent Source repository: \href{https://github.com/docxology/biology_textbook}{https://github.com/docxology/biology\_textbook}\\
\vspace{0.8em}
\begin{center}
\begingroup
\setlength{\fboxsep}{8pt}
\fcolorbox{red!55!black}{red!3}{%
\begin{minipage}{0.88\textwidth}
\footnotesize\itshape “The tree which moves some to tears of joy is in the eyes of others only a green thing that stands in the way. Some see nature all ridicule and deformity... and some scarce see nature at all. But to the eyes of the man of imagination, nature is imagination itself.”
\par\vspace{0.35em}\raggedleft\upshape --- William Blake
\end{minipage}%
}
\endgroup
\end{center}
\vspace{0.6em}
\noindent\fcolorbox{black!35}{black!2}{%
\begin{minipage}{0.94\textwidth}
{\sffamily\bfseries Acknowledgements}\par\vspace{0.25em}
{\footnotesize This open textbook is built from a living chain of public scholarship, open-source tools, tested computational models, and classroom-facing biological explanation. The author thanks the students, educators, research communities, open-science maintainers, and standards bodies whose work makes a reproducible, adaptable biology text possible.}
\end{minipage}%
}
\par
\vspace{1.0em}
\noindent Suggested citation: Friedman, D. A. (2026). \textit{Introduction to Biology: A Generative Approach} (Edition 1.0). Active Inference Institute. \href{https://github.com/docxology/biology_textbook}{https://github.com/docxology/biology\_textbook}.

\vspace{1.0em}
\noindent This open textbook is generated from version-controlled Markdown, tested Python modules, programmatic figures, and rendered Mermaid diagrams. Corrections and improvements may be submitted via the source repository linked above.

\vfill
\noindent Accessibility note: the compact PDF is optimized for dense print. Reader-profile builds, HTML output, and source Markdown can be generated from the same manuscript materials.
\clearpage
\tableofcontents
\clearpage


\section*{Front Matter}\label{front-matter}
\addcontentsline{toc}{section}{Front Matter}

\subsection*{Dedication}\label{dedication}
\addcontentsline{toc}{subsection}{Dedication}

\emph{To most students who encounter biology not as a collection of
facts,\\
but as a living, computable, mathematically rich discipline ---\\
and to those who show them the way.}

\vspace{2em}

\emph{And to the open-science community whose shared code, data, and
scholarship\\
made a freely available, computationally grounded textbook possible.}

\newpage

\begin{center}\rule{0.5\linewidth}{0.5pt}\end{center}

\subsection*{Acknowledgements}\label{acknowledgements}
\addcontentsline{toc}{subsection}{Acknowledgements}

The author gratefully acknowledges the following contributions and
sources of inspiration:

\textbf{Scientific foundations.} This textbook builds on the
foundational scholarship of Alberts et al.~(\emph{Molecular Biology of
the Cell}), Campbell \& Reece (\emph{Biology}), Stryer
(\emph{Biochemistry}), Lodish et al.~(\emph{Molecular Cell Biology}),
the primary literature cited throughout each chapter, and a variety of
other sources. Specific experimental results are cited at point of use
with full author credit. The structure of this textbook is informed by
the open curriculum
\href{https://openstax.org/details/books/biology-2e}{OpenStax Biology
2e} and the \href{https://qubeshub.org/}{QUBES Hub} community of
quantitative biology educators.

\textbf{Computational tools and open source.} All figures were generated
with \href{https://matplotlib.org/}{Matplotlib} and
\href{https://numpy.org/}{NumPy}. Mermaid process diagrams were rendered
with the \href{https://github.com/mermaid-js/mermaid-cli}{Mermaid CLI}.
The textbook pipeline uses \href{https://pandoc.org/}{Pandoc} and
\href{https://www.latex-project.org/}{LaTeX} for PDF rendering.
Scholarship was enriched using
\href{https://www.perplexity.ai/}{Perplexity} AI-assisted literature
searches.

\textbf{Research project template.} The build pipeline, testing
infrastructure, and multi-project architecture are based on the Research
Project Template maintained in the
\href{https://github.com/docxology/template}{Research Project Template
repository}. This textbook's sources, figures, and tests live in the
dedicated book repository
\href{https://github.com/docxology/biology_textbook}{biology textbook
source repository}. The template is a comprehensive framework for
reproducible research, including version control, testing,
documentation, and publication. It is a model for this work, and the
template is used to build this textbook.

\newpage

\begin{center}\rule{0.5\linewidth}{0.5pt}\end{center}

\subsection*{About This Textbook}\label{about-this-textbook}
\addcontentsline{toc}{subsection}{About This Textbook}

\subsubsection*{Computational
Philosophy}\label{computational-philosophy}
\addcontentsline{toc}{subsubsection}{Computational Philosophy}

\emph{Introduction to Biology} integrates narrative explanation with
quantitative models. When a chapter introduces a model (enzyme kinetics,
predator--prey dynamics, action potentials), the corresponding
computation is implemented in the accompanying Python modules and used
to generate figures where appropriate.

This approach aims to help students move between mechanisms,
measurements, and models: what a system does, how we know, and what a
simple quantitative model can predict (and where it breaks).

\subsubsection*{What Makes This Textbook
Different}\label{what-makes-this-textbook-different}
\addcontentsline{toc}{subsubsection}{What Makes This Textbook Different}

{\def\LTcaptype{none} % do not increment counter
\begin{longtable}[]{@{}
  >{\raggedright\arraybackslash}p{(\linewidth - 4\tabcolsep) * \real{0.1750}}
  >{\raggedright\arraybackslash}p{(\linewidth - 4\tabcolsep) * \real{0.3250}}
  >{\raggedright\arraybackslash}p{(\linewidth - 4\tabcolsep) * \real{0.5000}}@{}}
\toprule\noalign{}
\begin{minipage}[b]{\linewidth}\raggedright
Feature
\end{minipage} & \begin{minipage}[b]{\linewidth}\raggedright
This Textbook
\end{minipage} & \begin{minipage}[b]{\linewidth}\raggedright
Traditional Textbook
\end{minipage} \\
\midrule\noalign{}
\endhead
\bottomrule\noalign{}
\endlastfoot
\textbf{Figures} & Generated from mathematical models; reproducible &
Static graphics \\
\textbf{Equations} & Derived and numerically validated & Stated without
derivation \\
\textbf{Code} & Python modules for every major model & Absent \\
\textbf{Diagrams} & Mermaid process diagrams; automatically rendered &
Hand-drawn \\
\textbf{Opening vignettes} & Landmark experiment narrative opening each
chapter & Absent \\
\textbf{Curriculum map} & Generated chapter/lab/question/standards
alignment & Usually external \\
\textbf{Master glossary} & 220 terms with etymology and chapter
cross-references & Chapter-end lists \\
\textbf{Open access} & CC BY 4.0; fully open source & Copyright
restricted \\
\textbf{Citations} & Inline, with year and journal & Often chapter-end
primarily \\
\end{longtable}
}

\subsubsection*{How to Navigate This
Book}\label{how-to-navigate-this-book}
\addcontentsline{toc}{subsubsection}{How to Navigate This Book}

The textbook is organized from systems-level orientation through
molecular, cellular, organismal, evolutionary, and ecological scales.
The entries below are generated from \breaktt{manuscript/config.yaml} so
navigation stays aligned with the rendered table of contents.

\begin{itemize}
\tightlist
\item
  \textbf{Unit 0 --- Systems Science and the Biology of Complexity:}
  Systems Science and the Logic of Emergence; Complex Adaptive Systems;
  Active Inference and the Free Energy Principle.
\item
  \textbf{Unit I --- Chemistry of Life:} Atoms, Molecules, and Chemical
  Bonds; Water --- The Molecule of Life; Biological Macromolecules;
  Enzymes and the Kinetics of Catalysis.
\item
  \textbf{Unit II --- The Cell:} Cell Theory and Cell Types; Cell
  Structure and Organelles; Membrane Structure and Transport; Cell
  Signalling and Communication.
\item
  \textbf{Unit III --- Energy and Metabolism:} Bioenergetics and
  Cellular Respiration; Photosynthesis; Metabolic Integration and
  Regulation.
\item
  \textbf{Unit IV --- Molecular Genetics:} DNA Replication and the Cell
  Cycle; Gene Expression --- Transcription and Translation; Mutations,
  CRISPR, and Genomics; Epigenetics and Gene Regulation.
\item
  \textbf{Unit V --- Classical Genetics and Heredity:} Mendelian
  Genetics and Heredity; Chromosomal Inheritance and Linkage; Population
  Genetics and Hardy-Weinberg Equilibrium.
\item
  \textbf{Unit VI --- Evolution:} Evolution --- Theory, Natural
  Selection, and Adaptation; Genetic Drift, Gene Flow, and Speciation;
  Phylogenetics and the Tree of Life.
\item
  \textbf{Unit VII --- Microbiology:} Bacteria, Archaea, and Viruses;
  Microbial Ecology and the Microbiome; Infectious Disease and Immunity.
\item
  \textbf{Unit VIII --- Botany --- Plant Biology:} Plant Structure,
  Water Relations, and Transport; Plant Reproduction and Development;
  Plant Responses to the Environment.
\item
  \textbf{Unit IX --- Zoology and Systems Physiology:} Circulation,
  Respiration, and Homeostasis; Nervous System and Neural Signalling;
  Action Potentials and Synaptic Transmission; Endocrine and Immune
  Systems.
\item
  \textbf{Unit X --- Ecology:} Population Ecology and Growth Models;
  Community Ecology and Species Interactions; Ecosystem Ecology and
  Biogeochemical Cycles; Biomes and Conservation Biology.
\item
  \textbf{Laboratory activities:} one companion lab follows each chapter
  in the same canonical order.
\item
  \textbf{Question banks:} one 30-item question bank follows each
  chapter in the same canonical order.
\item
  \textbf{Appendix A --- Curriculum Map:} reference material generated
  or ordered from the same manifest.
\item
  \textbf{Appendix B --- Instructor Orchestration Guide:} reference
  material generated or ordered from the same manifest.
\item
  \textbf{Appendix C --- Mathematical Review for Biology:} reference
  material generated or ordered from the same manifest.
\item
  \textbf{Appendix D --- Units, Physical Constants, and Biological
  Ranges:} reference material generated or ordered from the same
  manifest.
\item
  \textbf{Appendix E --- A Periodic Table for Biology:} reference
  material generated or ordered from the same manifest.
\item
  \textbf{Appendix F --- Master Glossary of Biological Terms:} reference
  material generated or ordered from the same manifest.
\item
  \textbf{Appendix G --- Index of Key Terms:} reference material
  generated or ordered from the same manifest.
\item
  \textbf{Source modules:}
  \texttt{src/biology/\textless{}domain\textgreater{}/} contains the
  tested Python implementations for the quantitative models used
  throughout the book.
\end{itemize}

\subsubsection*{Suggested reading paths}\label{suggested-reading-paths}
\addcontentsline{toc}{subsubsection}{Suggested reading paths}

{\def\LTcaptype{none} % do not increment counter
\begin{longtable}[]{@{}
  >{\raggedright\arraybackslash}p{(\linewidth - 4\tabcolsep) * \real{0.1600}}
  >{\raggedright\arraybackslash}p{(\linewidth - 4\tabcolsep) * \real{0.6400}}
  >{\raggedright\arraybackslash}p{(\linewidth - 4\tabcolsep) * \real{0.2000}}@{}}
\toprule\noalign{}
\begin{minipage}[b]{\linewidth}\raggedright
Path
\end{minipage} & \begin{minipage}[b]{\linewidth}\raggedright
Units / emphasis
\end{minipage} & \begin{minipage}[b]{\linewidth}\raggedright
Notes
\end{minipage} \\
\midrule\noalign{}
\endhead
\bottomrule\noalign{}
\endlastfoot
\textbf{AP / first-year survey} & I, II, III, IV (selections), V, VI
(selections), X & Skim Unit 0; prioritise metabolism and genetics core
narratives. \\
\textbf{Pre-health / majors} & Full I--IX + X; Unit 0 as orientation &
Add labs for quantitative skills; pair each physiology chapter with its
Python bridge. \\
\textbf{Ecology / environmental focus} & I--II (review), III
(photosynthesis), VI, VII, X & Emphasise population models,
biogeochemistry, conservation metrics in \texttt{ecology.py}. \\
\textbf{Computation-first} & Unit 0 + any unit & Read ``Bridge to
computation'' blocks first, then narrative; run
\breaktt{scripts/generate\_figures.py}. \\
\end{longtable}
}

\subsubsection*{Notation and
conventions}\label{notation-and-conventions}
\addcontentsline{toc}{subsubsection}{Notation and conventions}

\begin{itemize}
\tightlist
\item
  \textbf{Logarithms:} \(\ln\) = natural log; \(\log_{10}\) used where
  orders of magnitude matter (pH, viral titre, doubling time).
\item
  \textbf{Concentrations:} square brackets \([X]\) denote molarity
  unless a chapter states otherwise.
\item
  \textbf{Genetics:} italic gene symbols (\emph{lacZ}); protein products
  often Roman with capital initial (LacZ) where conventional.
\item
  \textbf{Units:} SI base units; physiology at 37 °C, pH 7.4, sea level
  unless noted.
\item
  \textbf{Chapter numbers} in the glossary and cross-links match the PDF
  table of contents (sequential numbering from \texttt{config.yaml},
  including Unit 0 when present).
\end{itemize}

\subsubsection*{Scholarship and source
practice}\label{scholarship-and-source-practice}
\addcontentsline{toc}{subsubsection}{Scholarship and source practice}

Biology changes because methods change: better microscopes, longer
reads, larger cohorts, improved models, and more inclusive sampling can
most revise a claim that once looked settled. When reading this book,
treat every important statement as part of a source chain:

{\def\LTcaptype{none} % do not increment counter
\begin{longtable}[]{@{}
  >{\raggedright\arraybackslash}p{(\linewidth - 4\tabcolsep) * \real{0.3333}}
  >{\raggedright\arraybackslash}p{(\linewidth - 4\tabcolsep) * \real{0.3333}}
  >{\raggedright\arraybackslash}p{(\linewidth - 4\tabcolsep) * \real{0.3333}}@{}}
\toprule\noalign{}
\begin{minipage}[b]{\linewidth}\raggedright
Source type
\end{minipage} & \begin{minipage}[b]{\linewidth}\raggedright
Best use
\end{minipage} & \begin{minipage}[b]{\linewidth}\raggedright
Question to ask
\end{minipage} \\
\midrule\noalign{}
\endhead
\bottomrule\noalign{}
\endlastfoot
Primary research article & Evidence for a specific method, dataset,
result, or mechanism & What exactly was measured, and in which organism,
cell type, population, or environment? \\
Review or textbook synthesis & Orientation across a field or controversy
& Which primary findings does the synthesis depend on, and are any newer
results likely to change it? \\
Public dataset or database & Reusable evidence for comparison,
reanalysis, and reproducibility & How were samples selected, processed,
filtered, and annotated? \\
Institutional report or guideline & Current public-health, clinical,
biodiversity, or policy status & What date, jurisdiction, and decision
context does the guidance assume? \\
Model or simulation & A testable simplification of a mechanism & Which
assumptions, parameters, or boundary conditions would change the
conclusion? \\
\end{longtable}
}

For recent or numeric claims, prefer the source closest to the
measurement and write one sentence naming what would change your
confidence. For example, a claim about antimicrobial resistance should
identify the organism-resistance pair, surveillance population, and
selection pressure; a claim about biodiversity loss should separate
population trends, extinction risk, land-use drivers, and value
judgments. This is why chapter answer keys ask for both a core response
and a scholarship check.

\subsubsection*{Textbook Concept Map}\label{textbook-concept-map}
\addcontentsline{toc}{subsubsection}{Textbook Concept Map}

The \textbf{eleven} instructional blocks of this textbook (Unit 0 plus
Units I--X) form an interdependent architecture. The diagram below shows
primary dependency paths and integrative threads.

\begin{figure}[htbp]
\centering
\includegraphics[width=0.82\linewidth,height=4.2in,keepaspectratio]{../figures/mermaid_inline/inline_mermaid_0001_de08675e7bd6.png}
\caption{Figure FM-1. Dependency map: Unit 0 (dashed links) is a conceptual scaffold for feedback, emergence, and inference - not a prerequisite for reading Unit I in detail, but useful orientation. Solid arrows show content dependencies among Units I-X. Units I-III provide the molecular foundation; IV-VI the hereditary and evolutionary framework; VII-IX apply these to major organismal groups; Unit X integrates at the population-to-biosphere scale.}
\end{figure}

\subsubsection*{Accessing source
materials}\label{accessing-source-materials}
\addcontentsline{toc}{subsubsection}{Accessing source materials}

{\def\LTcaptype{none} % do not increment counter
\begin{longtable}[]{@{}
  >{\raggedright\arraybackslash}p{(\linewidth - 2\tabcolsep) * \real{0.5000}}
  >{\raggedright\arraybackslash}p{(\linewidth - 2\tabcolsep) * \real{0.5000}}@{}}
\toprule\noalign{}
\begin{minipage}[b]{\linewidth}\raggedright
Resource
\end{minipage} & \begin{minipage}[b]{\linewidth}\raggedright
Location
\end{minipage} \\
\midrule\noalign{}
\endhead
\bottomrule\noalign{}
\endlastfoot
\textbf{Living source (Git)} &
\href{https://github.com/docxology/biology_textbook}{biology textbook
source repository} --- manuscript Markdown, \texttt{src/biology/}
modules, maintenance scripts, and tests \\
\textbf{Archival citation (DOI)} &
\href{https://doi.org/10.5281/zenodo.20286478}{Zenodo archival DOI
record} --- cite this record for a fixed edition snapshot \\
\textbf{Combined PDF / HTML} & Build from the repository with
\breaktt{uv\ run\ python\ scripts/03\_render\_pdf.py\ -\/-project\ biology\_textbook}
(from the template root) or the project's \breaktt{biology\_analysis.py}
workflow \\
\textbf{Figures and diagrams} & \breaktt{scripts/generate\_figures.py}
(matplotlib) and \breaktt{scripts/generate\_diagrams.py} (registered
Mermaid); inline Mermaid in chapters renders during PDF preprocessing \\
\textbf{Corrections and extensions} & Open an issue or pull request on
the book repository; substantive edits should keep chapter
\texttt{\textbackslash{}cref} labels and \texttt{config.yaml} ordering
in sync \\
\end{longtable}
}

Text is licensed \textbf{CC BY 4.0}; source code is \textbf{Apache-2.0}
(see \breaktt{manuscript/config.yaml} \texorpdfstring{\ensuremath{\to}}{->} \texttt{book}).

\subsection*{Course Planning Grid}\label{course-planning-grid}
\addcontentsline{toc}{subsection}{Course Planning Grid}

The table below provides a per-chapter difficulty rating (Level 1/3 to
Level 3/3), an estimated student reading time, and a suggested lecture
allotment. It is auto-generated by
\breaktt{scripts/insert\_chapter\_metadata.py} from the canonical table
of contents in \breaktt{manuscript/config.yaml} plus
\breaktt{src/biology/chapter\_metadata.py} --- edit those sources and
re-run the script to refresh this grid.

{\def\LTcaptype{none} % do not increment counter
\begin{longtable}[]{@{}
  >{\raggedright\arraybackslash}p{(\linewidth - 10\tabcolsep) * \real{0.0625}}
  >{\raggedright\arraybackslash}p{(\linewidth - 10\tabcolsep) * \real{0.1250}}
  >{\raggedright\arraybackslash}p{(\linewidth - 10\tabcolsep) * \real{0.1875}}
  >{\raggedright\arraybackslash}p{(\linewidth - 10\tabcolsep) * \real{0.2500}}
  >{\raggedright\arraybackslash}p{(\linewidth - 10\tabcolsep) * \real{0.1875}}
  >{\raggedright\arraybackslash}p{(\linewidth - 10\tabcolsep) * \real{0.1875}}@{}}
\toprule\noalign{}
\begin{minipage}[b]{\linewidth}\raggedright
\#
\end{minipage} & \begin{minipage}[b]{\linewidth}\raggedright
Unit
\end{minipage} & \begin{minipage}[b]{\linewidth}\raggedright
Chapter
\end{minipage} & \begin{minipage}[b]{\linewidth}\raggedright
Difficulty
\end{minipage} & \begin{minipage}[b]{\linewidth}\raggedright
Reading
\end{minipage} & \begin{minipage}[b]{\linewidth}\raggedright
Lecture
\end{minipage} \\
\midrule\noalign{}
\endhead
\bottomrule\noalign{}
\endlastfoot
0 & 0 & Systems Science and the Logic of Emergence & Level 2/3 & 35 min
& 50 min \\
0 & 0 & Complex Adaptive Systems & Level 2/3 & 35 min & 50 min \\
0 & 0 & Active Inference and the Free Energy Principle & Level 3/3 & 45
min & 75 min \\
1 & I & Atoms, Molecules, and Chemical Bonds & Level 1/3 & 40 min & 50
min \\
2 & I & Water --- The Molecule of Life & Level 1/3 & 40 min & 50 min \\
3 & I & Biological Macromolecules & Level 2/3 & 55 min & 75 min \\
4 & I & Enzymes and the Kinetics of Catalysis & Level 3/3 & 60 min & 75
min \\
5 & II & Cell Theory and Cell Types & Level 1/3 & 45 min & 50 min \\
6 & II & Cell Structure and Organelles & Level 2/3 & 50 min & 75 min \\
7 & II & Membrane Structure and Transport & Level 2/3 & 50 min & 75
min \\
8 & II & Cell Signalling and Communication & Level 3/3 & 55 min & 75
min \\
9 & III & Bioenergetics and Cellular Respiration & Level 3/3 & 60 min &
100 min \\
10 & III & Photosynthesis & Level 2/3 & 55 min & 75 min \\
11 & III & Metabolic Integration and Regulation & Level 3/3 & 60 min &
100 min \\
12 & IV & DNA Replication and the Cell Cycle & Level 2/3 & 55 min & 75
min \\
13 & IV & Gene Expression --- Transcription and Translation & Level 2/3
& 60 min & 100 min \\
14 & IV & Mutations, CRISPR, and Genomics & Level 2/3 & 55 min & 75
min \\
15 & IV & Epigenetics and Gene Regulation & Level 3/3 & 50 min & 75
min \\
16 & V & Mendelian Genetics and Heredity & Level 2/3 & 65 min & 100
min \\
17 & V & Chromosomal Inheritance and Linkage & Level 2/3 & 60 min & 75
min \\
18 & V & Population Genetics and Hardy-Weinberg Equilibrium & Level 3/3
& 75 min & 100 min \\
19 & VI & Evolution --- Theory, Natural Selection, and Adaptation &
Level 2/3 & 60 min & 75 min \\
20 & VI & Genetic Drift, Gene Flow, and Speciation & Level 3/3 & 60 min
& 75 min \\
21 & VI & Phylogenetics and the Tree of Life & Level 3/3 & 60 min & 100
min \\
22 & VII & Bacteria, Archaea, and Viruses & Level 2/3 & 65 min & 75
min \\
23 & VII & Microbial Ecology and the Microbiome & Level 2/3 & 60 min &
75 min \\
24 & VII & Infectious Disease and Immunity & Level 2/3 & 60 min & 75
min \\
25 & VIII & Plant Structure, Water Relations, and Transport & Level 2/3
& 55 min & 75 min \\
26 & VIII & Plant Reproduction and Development & Level 2/3 & 55 min & 75
min \\
27 & VIII & Plant Responses to the Environment & Level 2/3 & 55 min & 75
min \\
28 & IX & Circulation, Respiration, and Homeostasis & Level 3/3 & 60 min
& 100 min \\
29 & IX & Nervous System and Neural Signalling & Level 3/3 & 55 min & 75
min \\
30 & IX & Action Potentials and Synaptic Transmission & Level 3/3 & 55
min & 100 min \\
31 & IX & Endocrine and Immune Systems & Level 2/3 & 55 min & 75 min \\
32 & X & Population Ecology and Growth Models & Level 3/3 & 75 min & 100
min \\
33 & X & Community Ecology and Species Interactions & Level 2/3 & 80 min
& 100 min \\
34 & X & Ecosystem Ecology and Biogeochemical Cycles & Level 2/3 & 65
min & 75 min \\
35 & X & Biomes and Conservation Biology & Level 2/3 & 70 min & 75
min \\
& & \textbf{Totals} & & \textbf{2150 min (35 h)} & \textbf{2975 min (49
h)} \\
\end{longtable}
}

\newpage
\newpage

\newpage

\section*{Preface}\label{preface}
\addcontentsline{toc}{section}{Preface}

\subsection*{A Textbook Built With
Code}\label{a-textbook-built-with-code}
\addcontentsline{toc}{subsection}{A Textbook Built With Code}

Welcome to \emph{Introduction to Biology: A Generative Approach} --- an
open-source textbook covering introductory biology across \textbf{Unit
0} (systems science and complexity), \textbf{ten thematic units (I--X)},
and \textbf{38 core chapters}, plus optional \textbf{laboratories} and
\textbf{question banks} in the appendices. Where this text uses
quantitative models, the corresponding computations are implemented as
tested Python modules.

Whether we examine Michaelis--Menten enzyme kinetics, Lotka--Volterra
predator--prey dynamics, or Hodgkin--Huxley action potentials, the
underlying mathematical model exists as a working module in the
accompanying codebase. Many figures are generated programmatically.
Process diagrams are expressed using Mermaid where it improves clarity.
For instructors and independent learners, this supports inspection,
extension, and reproducible re-generation of results.

The philosophy is: \textbf{understand biology by computing biology}.

\textbf{How to read this book digitally:} The reference build is a
compact \textbf{PDF} with LaTeX typesetting. Not all screen readers
handle mathematical notation the same way in PDF; the same sources also
produce \textbf{HTML} for some workflows. The
\href{https://github.com/docxology/biology_textbook}{biology textbook
source repository} contains the manuscript, figures, tests, and code;
the \href{https://doi.org/10.5281/zenodo.20286478}{Zenodo archival DOI
record} identifies a fixed edition snapshot. If you need larger type or
margins, a \textbf{reader} typography profile is documented (paired
edits to \breaktt{manuscript/config.yaml} and
\breaktt{manuscript/preamble.md} --- see \breaktt{docs/accessibility.md}
in the project tree). The book does not re-express every formula in
natural language; where precision matters, work through the worked
examples and the linked \texttt{src/biology} modules.

\textbf{Accessible formats:} The compact PDF is optimized for print
density, not for every reader. For screen-reader review, HTML/MathJax
workflows usually preserve reading order and mathematical navigation
better than a dense PDF. Figures and Mermaid diagrams include alt text
plus captions, and the large-type reader profile increases margins, font
size, and line spacing without changing chapter order or cross-reference
IDs.

\begin{center}\rule{0.5\linewidth}{0.5pt}\end{center}

\subsection*{Five Big Ideas in Biology}\label{five-big-ideas-in-biology}
\addcontentsline{toc}{subsection}{Five Big Ideas in Biology}

Biology is a vast and diverse science, but five core concepts unify its
every subdiscipline. These ideas --- drawn from the AAAS \emph{Vision
and Change in Undergraduate Biology Education}
\citep{visionandchange2011} --- are introduced here as orienting themes
that recur across the units that follow.

\subsubsection*{Evolution: the unifying theory of
life}\label{evolution-the-unifying-theory-of-life}
\addcontentsline{toc}{subsubsection}{Evolution: the unifying theory of
life}

\emph{"Nothing in biology makes sense except in the light of
evolution."} --- Theodosius \citet{dobzhansky1973}

Evolution by natural selection explains the origin of species, the
genetic code being nearly comprehensive, the homology between a human
arm and a bat wing, and why influenza vaccines must be redesigned every
year. Units VI (Evolution), V (Genetics), and IV (Molecular Genetics)
develop this idea from population genetics algorithms to phylogenetic
tree reconstruction.

\subsubsection*{Structure and function are
inseparable}\label{structure-and-function-are-inseparable}
\addcontentsline{toc}{subsubsection}{Structure and function are
inseparable}

A phospholipid bilayer is 7 nm thick and amphipathic --- and those two
structural facts explain every property of biological membranes (Units I
and II). The \texorpdfstring{\ensuremath{\alpha}}{alpha}-helix of haemoglobin's subunits explains its
cooperativity and the Bohr effect (Units I and IX). The T-shaped
architecture of the mitochondrial ATP synthase explains rotary catalysis
(Unit III). In biology, whenever you ask \emph{how}, the answer is
typically embedded in \emph{shape}.

\subsubsection*{Information --- storage, transfer,
expression}\label{information-storage-transfer-expression}
\addcontentsline{toc}{subsubsection}{Information --- storage, transfer,
expression}

DNA is not merely a chemical; it is a \textbf{code}. An alphabet of four
nucleotides encodes a programme of 20 amino acids, creating an
effectively unlimited diversity of proteins. Unit IV (Molecular
Genetics) examines the molecular machinery that reads this code --- from
DNA helicases to ribosomes --- and how errors in the code drive disease
and evolution. Unit V (Classical Genetics) examines how the code is
transmitted between generations. Unit VII (Microbiology) shows how
viruses inject their own code into host cells.

\subsubsection*{Emergent properties: the whole exceeds the
parts}\label{emergent-properties-the-whole-exceeds-the-parts}
\addcontentsline{toc}{subsubsection}{Emergent properties: the whole
exceeds the parts}

A single neuron can fire or remain silent. A brain can think. A single
predator can collapse an entire intertidal community (Unit X). A slight
imbalance in NAD\texorpdfstring{\textsuperscript{+}}{+}/NADH ratio shunts metabolism from aerobic to anaerobic
(Unit III). Emergence --- complex behaviour arising from simple rules
--- is everywhere in biology, and understanding it requires systems
thinking: tracking flows, feedback loops, and nonlinear dynamics, not
just cataloguing parts.

\subsubsection*{Cells: the universal unit of
life}\label{cells-the-universal-unit-of-life}
\addcontentsline{toc}{subsubsection}{Cells: the universal unit of life}

Every living organism is composed of one or more cells, sharing a common
molecular toolkit: phospholipid membranes, DNA, ribosomes, ATP. Unit II
examines the cell as a physical and computational system. Units VII--IX
extend this to specialised cells: bacteria, plant cells, neurons, immune
cells. Understanding the cell is not a module to be completed --- it is
a perspective to carry through every subsequent chapter.

\begin{center}\rule{0.5\linewidth}{0.5pt}\end{center}

\subsection*{How Science Actually
Works}\label{how-science-actually-works}
\addcontentsline{toc}{subsection}{How Science Actually Works}

This textbook presents science not as a collection of established facts
but as a dynamic, self-correcting process of inquiry. Understanding how
scientific knowledge is generated, tested, and revised is as important
as knowing what is currently accepted.

\textbf{The hypothetico-deductive method.} A hypothesis is a testable,
falsifiable prediction. Barry Marshall and Robin Warren's hypothesis
that bacteria cause peptic ulcers (not stress) was so contrary to
medical consensus that Marshall famously drank a \emph{Helicobacter
pylori} culture to demonstrate the relationship \citep{marshall1984}.
Their Nobel Prize-winning work illustrates how a single decisive
experiment can overturn decades of received wisdom.

\textbf{Models and computation as scientific tools.} Modern biology is
inseparable from mathematical modelling. The Michaelis-Menten equation
(1913) is a model. The Hardy-Weinberg principle is a model. The
Hodgkin-Huxley equations (1952) were Nobel Prize-winning models. A model
is not a simplification that sacrifices truth; it is a commitment to
precision --- stating \emph{exactly} what you are and are not claiming.
The Python modules in this textbook implement these models so that you
can inspect their assumptions, vary their parameters, and see where they
break.

\textbf{Primary literature vs.~textbook synthesis.} This textbook cites
primary papers at point of use. When you see "(Author et al., Year,
\emph{Journal})" inline, that citation points to the experiment or
calculation that established the claim. Science journalism, textbooks,
and Wikipedia involve more layers of interpretation. Primary papers are
the ultimate authority --- and learning to read them is a core
competency this textbook aims to develop.

\textbf{Revision is not failure.} Every chapter in this book contains
examples of ideas revised or overturned: the flat earth was replaced by
a spherical one; Lamarckian inheritance was replaced by Darwinian
selection; the DNA double helix replaced protein as the genetic
material; the "lock-and-key" enzyme model was replaced by "induced fit".
A willingness to revise beliefs on evidence is not weakness --- it is
the engine of scientific progress.

\begin{center}\rule{0.5\linewidth}{0.5pt}\end{center}

\subsection*{How to Use the Unit
Introductions}\label{how-to-use-the-unit-introductions}
\addcontentsline{toc}{subsection}{How to Use the Unit Introductions}

Each unit begins with a \texttt{unit\_intro.md} page containing:

\begin{itemize}
\tightlist
\item
  \textbf{A narrative hook} explaining why this area of biology matters
  and what surprises it holds
\item
  \textbf{A landmark discoveries table} citing the experiments that
  established the field
\item
  \textbf{A concept map} (Mermaid diagram) showing how the chapters
  connect to each other and to other units
\item
  \textbf{A chapter roadmap} mapping each chapter to its core question
  and key equation
\end{itemize}

Use these pages as \textbf{advance organisers} --- read them before
starting the unit and return to them when you feel you have lost the
conceptual thread. The best moment to understand how a jigsaw piece fits
is before you lose sight of the box.

\begin{center}\rule{0.5\linewidth}{0.5pt}\end{center}

\subsection*{Scope and Organisation}\label{scope-and-organisation}
\addcontentsline{toc}{subsection}{Scope and Organisation}

The textbook proceeds from atoms to ecosystems, following the standard
introductory course arc:

{\def\LTcaptype{none} % do not increment counter
\begin{longtable}[]{@{}
  >{\raggedright\arraybackslash}p{(\linewidth - 4\tabcolsep) * \real{0.1818}}
  >{\raggedright\arraybackslash}p{(\linewidth - 4\tabcolsep) * \real{0.2273}}
  >{\raggedright\arraybackslash}p{(\linewidth - 4\tabcolsep) * \real{0.5909}}@{}}
\toprule\noalign{}
\begin{minipage}[b]{\linewidth}\raggedright
Unit
\end{minipage} & \begin{minipage}[b]{\linewidth}\raggedright
Title
\end{minipage} & \begin{minipage}[b]{\linewidth}\raggedright
Core Chapters
\end{minipage} \\
\midrule\noalign{}
\endhead
\bottomrule\noalign{}
\endlastfoot
\textbf{0} & Systems Science and the Biology of Complexity & Emergence,
complex adaptive systems, active inference \\
\textbf{I} & Chemistry of Life & Atoms \& bonds, water, macromolecules,
enzymes \\
\textbf{II} & The Cell & Cell theory, organelles, membranes,
signalling \\
\textbf{III} & Energy and Metabolism & Respiration, photosynthesis,
metabolic integration \\
\textbf{IV} & Molecular Genetics & DNA replication, gene expression,
CRISPR and genomics \\
\textbf{V} & Classical Genetics and Heredity & Mendelian genetics,
chromosomal inheritance, population genetics \\
\textbf{VI} & Evolution & Natural selection, genetic drift, speciation,
phylogenetics \\
\textbf{VII} & Microbiology & Bacteria, archaea, viruses, microbiome,
infectious disease \\
\textbf{VIII} & Botany --- Plant Biology & Plant structure,
reproduction, environmental responses \\
\textbf{IX} & Zoology and Systems Physiology & Circulation, respiration,
nervous system, endocrinology \\
\textbf{X} & Ecology & Population, community, ecosystem, biomes,
conservation \\
\end{longtable}
}

\begin{center}\rule{0.5\linewidth}{0.5pt}\end{center}

\subsection*{Reading paths by goal}\label{reading-paths-by-goal}
\addcontentsline{toc}{subsection}{Reading paths by goal}

{\def\LTcaptype{none} % do not increment counter
\begin{longtable}[]{@{}
  >{\raggedright\arraybackslash}p{(\linewidth - 4\tabcolsep) * \real{0.1081}}
  >{\raggedright\arraybackslash}p{(\linewidth - 4\tabcolsep) * \real{0.3784}}
  >{\raggedright\arraybackslash}p{(\linewidth - 4\tabcolsep) * \real{0.5135}}@{}}
\toprule\noalign{}
\begin{minipage}[b]{\linewidth}\raggedright
Goal
\end{minipage} & \begin{minipage}[b]{\linewidth}\raggedright
Where to start
\end{minipage} & \begin{minipage}[b]{\linewidth}\raggedright
How to use the code
\end{minipage} \\
\midrule\noalign{}
\endhead
\bottomrule\noalign{}
\endlastfoot
\textbf{Exam / course survey} & Unit intros + chapter summaries + the
companion question bank that follows each chapter & Answer odd-numbered
questions first; check module footers for \texttt{biology.*} imports. \\
\textbf{Wet-lab or clinical bridge} & Units II, IV, VII, IX & Read boxed
clinical / systems notes; pair with labs in the same unit. \\
\textbf{Modelling / CS} & Unit 0, then any unit's ``Bridge to
computation'' & Run examples with \texttt{uv\ run\ python} from the
project root; regenerate figures with
\breaktt{scripts/generate\_figures.py}. \\
\textbf{Ecology / field biology} & Units VI, X (+ Unit III
photosynthesis) & Focus on \texttt{ecology.py} functions cited in
chapter footers; work Lotka--Volterra and logistic examples by hand then
in Python. \\
\end{longtable}
}

\begin{center}\rule{0.5\linewidth}{0.5pt}\end{center}

\subsection*{How to Use This Book}\label{how-to-use-this-book}
\addcontentsline{toc}{subsection}{How to Use This Book}

Each chapter is designed for introductory biology learners and can be
read in two modes:

\begin{enumerate}
\def\labelenumi{\arabic{enumi}.}
\tightlist
\item
  \textbf{Concept-first}: Read the narrative, work the examples, and do
  the questions.
\item
  \textbf{Model-first}: Run the companion modules and treat the text as
  interpretation of model behavior.
\end{enumerate}

\textbf{For each concept, the integrated approach is:}

\begin{enumerate}
\def\labelenumi{\arabic{enumi}.}
\tightlist
\item
  \textbf{Conceptual narrative} --- the biological story and mechanism
  in plain language with precise vocabulary.
\item
  \textbf{Quantitative framework} --- equations, derivations, and worked
  numerical examples.
\item
  \textbf{Primary citations} --- key papers (author, year, journal)
  supporting mechanisms; citations are included where used.
\item
  \textbf{Python source} --- model code in the project's biology modules
  (importable and runnable).
\item
  \textbf{Figures} --- generated by \breaktt{src/visualization/plots.py}
  (\breaktt{ALL\_FIGURE\_GENERATORS}) and
  \breaktt{src/mermaid/biology\_diagrams.py}.
\item
  \textbf{Clinical connections} --- boxed examples connecting molecular
  mechanisms to human disease, therapy, and public health.
\item
  \textbf{Review and discussion questions} --- a 30-item companion
  question bank per chapter, plus end-of-chapter review questions, from
  calculation to synthesis (Unit 0 uses discussion prompts; other units
  mix numeric and conceptual items).
\end{enumerate}

\begin{center}\rule{0.5\linewidth}{0.5pt}\end{center}

\subsection*{Computational
Infrastructure}\label{computational-infrastructure}
\addcontentsline{toc}{subsection}{Computational Infrastructure}

The project's biology code is organized by domain:

\begin{Shaded}
\begin{Highlighting}[]
\NormalTok{src/biology/}
\NormalTok{├── biochemistry/    — Enzyme kinetics, macromolecule analysis}
\NormalTok{├── cell/            — Membrane biophysics, signalling, organelles}
\NormalTok{├── genetics/        — Mendelian ratios, Hardy{-}Weinberg, linkage mapping}
\NormalTok{├── physiology/      — Homeostasis, haemoglobin, cardiac models}
\NormalTok{├── ecology/         — Population dynamics, community models}
\NormalTok{├── evolution/       — Drift simulation, fitness landscapes}
\NormalTok{├── microbiology/    — Growth curves, MIC, viral cycles}
\NormalTok{├── botany/          — Water potential, transpiration, photosynthesis pathways}
\NormalTok{└── neuroscience/    — Action potentials, synapse models}
\end{Highlighting}
\end{Shaded}

\textbf{Diagram generation:}

\begin{Shaded}
\begin{Highlighting}[]
\NormalTok{src/mermaid/biology\_diagrams.py  — diagram factories registered in ALL\_BIOLOGY\_DIAGRAMS}
\NormalTok{src/visualization/plots.py    — matplotlib figure generators (ALL\_FIGURE\_GENERATORS)}
\end{Highlighting}
\end{Shaded}

Most code conforms to the project's \textbf{no-mocks policy}: tests use
real computations.

\begin{center}\rule{0.5\linewidth}{0.5pt}\end{center}

\subsection*{Pedagogical Standards}\label{pedagogical-standards}
\addcontentsline{toc}{subsection}{Pedagogical Standards}

\textbf{What this book assumes:} basic chemistry (atomic structure,
covalent bonds, pH), basic algebra, and optional access to Python. No
prior biology is required.

\textbf{Mathematical notation} is standard across most chapters:

\begin{itemize}
\tightlist
\item
  Concentrations in mol/L (M) or mmol/L (mM) as context requires
\item
  Reaction rates in M\texorpdfstring{\ensuremath{\cdot}}{.}s\texorpdfstring{\textsuperscript{-}}{-}\texorpdfstring{\textsuperscript{1}}{1} or \texorpdfstring{\ensuremath{\mu}}{mu}mol\texorpdfstring{\ensuremath{\cdot}}{.}min\texorpdfstring{\textsuperscript{-}}{-}\texorpdfstring{\textsuperscript{1}}{1}\texorpdfstring{\ensuremath{\cdot}}{.}mg protein\texorpdfstring{\textsuperscript{-}}{-}\texorpdfstring{\textsuperscript{1}}{1}
\item
  Genetic distances in centimorgans (cM)
\item
  Most physiological parameters indexed to standard conditions (37 °C,
  pH 7.4, sea level) unless stated
\end{itemize}

\textbf{Primary citations} follow the format Author (year),
\emph{Journal}, and are embedded inline. They are selective: landmark
results and a small number of current references where they materially
clarify the consensus.

\begin{center}\rule{0.5\linewidth}{0.5pt}\end{center}

\subsection*{Open Science Commitment}\label{open-science-commitment}
\addcontentsline{toc}{subsection}{Open Science Commitment}

This book is released under the Creative Commons Attribution 4.0
International Licence (CC BY 4.0): you may freely use, adapt, and
redistribute it with attribution. Most accompanying source code is
released under the Apache-2.0 Licence.

The entire textbook --- source markdown files, Python modules, figure
scripts, and the build pipeline --- is publicly available. The intent is
for this resource to be a living, improvable document rather than a
static product. Pull requests are welcome.

\begin{center}\rule{0.5\linewidth}{0.5pt}\end{center}

\subsection*{Acknowledgements}\label{acknowledgements-1}
\addcontentsline{toc}{subsection}{Acknowledgements}

This textbook was developed as part of the Research Project Template
infrastructure for open, reproducible science education. Gratitude to
the open-source communities behind \href{https://pandoc.org/}{Pandoc},
\href{https://www.latex-project.org/}{LaTeX},
\href{https://mermaid.js.org/}{Mermaid}, and
\href{https://matplotlib.org/}{matplotlib}.

\begin{center}\rule{0.5\linewidth}{0.5pt}\end{center}

\emph{Author:} \textbf{Daniel Ari Friedman}, Active Inference Institute
\emph{ORCID:}
\href{https://orcid.org/0000-0001-6232-9096}{0000-0001-6232-9096}
\emph{Edition 1.0, 2026} \emph{Licence: CC BY 4.0 (content) \texorpdfstring{\ensuremath{\cdot}}{.} Apache-2.0
(code)}

\newpage

\section*{Unit 0 --- Systems Science and the Biology of Complexity:
Introduction}\label{unit-0-systems-science-and-the-biology-of-complexity-introduction}
\addcontentsline{toc}{section}{Unit 0 --- Systems Science and the
Biology of Complexity: Introduction}

\textbf{Unit 0 \texorpdfstring{\ensuremath{\cdot}}{.} Introduction}

\begin{center}\rule{0.5\linewidth}{0.5pt}\end{center}

\subsection*{Why This Unit Matters}\label{why-this-unit-matters}
\addcontentsline{toc}{subsection}{Why This Unit Matters}

The modern study of biology began by reducing living phenomena to their
smallest parts --- cells, molecules, genes. This reductionist strategy
yielded extraordinary insight. Yet the more we learned about individual
components, the clearer it became that biology cannot be explained by
parts alone: a neuron firing is not a thought, a gene is not a
behaviour, a metabolic pathway is not a cell. What emerges from the
\emph{interaction} of parts --- properties no single component possesses
in isolation --- is the subject of \textbf{systems science}.

Unit 0 introduces three interlocking frameworks that recur in every
subsequent unit:

\begin{enumerate}
\def\labelenumi{\arabic{enumi}.}
\tightlist
\item
  \textbf{Systems science} --- the general theory of organised
  complexity: how hierarchical systems form, how feedback governs their
  behaviour, how new properties emerge.
\item
  \textbf{Complex adaptive systems} --- how populations of agents with
  local rules give rise to robust, evolvable, collective behaviour.
\item
  \textbf{Active inference and the free energy principle} --- a
  mathematically grounded account of how living agents maintain
  themselves by predicting and acting on their environments.
\end{enumerate}

Reading Unit 0 is optional but recommended: it supplies the vocabulary
of emergence, attractors, bifurcations, allostasis, and precision that
makes Units I--X fit together as a single coherent theory of living
organisation.

\begin{center}\rule{0.5\linewidth}{0.5pt}\end{center}

\subsection*{Landmark Discoveries}\label{landmark-discoveries}
\addcontentsline{toc}{subsection}{Landmark Discoveries}

{\def\LTcaptype{none} % do not increment counter
\begin{longtable}[]{@{}
  >{\raggedright\arraybackslash}p{(\linewidth - 4\tabcolsep) * \real{0.1429}}
  >{\raggedright\arraybackslash}p{(\linewidth - 4\tabcolsep) * \real{0.4286}}
  >{\raggedright\arraybackslash}p{(\linewidth - 4\tabcolsep) * \real{0.4286}}@{}}
\toprule\noalign{}
\begin{minipage}[b]{\linewidth}\raggedright
Year
\end{minipage} & \begin{minipage}[b]{\linewidth}\raggedright
Contribution
\end{minipage} & \begin{minipage}[b]{\linewidth}\raggedright
Significance
\end{minipage} \\
\midrule\noalign{}
\endhead
\bottomrule\noalign{}
\endlastfoot
1867 & Hermann von Helmholtz --- \emph{Handbook of Physiological Optics}
& First framing of perception as unconscious inference \\
1932 & Walter Cannon --- \emph{The Wisdom of the Body} & Coined
\emph{homeostasis}; proposed regulated internal milieu \\
1948 & Norbert Wiener --- \emph{Cybernetics} & Formalised feedback
control; bridged engineering and biology \\
1968 & Ludwig von Bertalanffy --- \emph{General System Theory} &
Articulated open systems and hierarchy across disciplines \\
1977 & Ilya Prigogine --- Nobel in Chemistry & Dissipative structures
and non-equilibrium self-organisation \\
1987 & Per Bak, Tang \& Wiesenfeld --- \emph{Phys. Rev.~Lett.} &
Self-organised criticality and power-law scaling in CAS \\
1992 & John Holland --- \emph{Adaptation in Natural and Artificial
Systems} & Genetic algorithms; formal CAS framework \\
2010 & Karl Friston --- \emph{Nature Rev.~Neurosci.} & Free energy
principle as a unified theory of brain function \\
2015 & Peter Sterling --- \emph{Allostasis} & Proactive predictive
regulation, contrasting homeostasis \\
\end{longtable}
}

\begin{center}\rule{0.5\linewidth}{0.5pt}\end{center}

\subsection*{Key Concepts and
Connections}\label{key-concepts-and-connections}
\addcontentsline{toc}{subsection}{Key Concepts and Connections}

\begin{itemize}
\tightlist
\item
  \textbf{System} --- a set of interacting components forming an
  integrated whole; open systems exchange matter and energy with their
  environment.
\item
  \textbf{Emergence} --- properties of a system that cannot be explained
  solely by properties of its components (e.g., consciousness emerges
  from neural circuits; life emerges from biochemistry).
\item
  \textbf{Feedback} --- output fed back as input, shaping future
  behaviour; negative feedback stabilises, positive feedback amplifies.
\item
  \textbf{Attractor} --- a region of phase space toward which
  trajectories converge; phase transitions move a system between
  attractors.
\item
  \textbf{Complex adaptive system (CAS)} --- a system of adaptive agents
  whose collective behaviour self-organises without central control.
\item
  \textbf{Free energy} --- a variational upper bound on surprise;
  minimising it unifies perception, action, and learning.
\item
  \textbf{Homeostasis vs.~allostasis} --- passive correction to a fixed
  set-point vs.~predictive adjustment of set-points based on context.
\end{itemize}

These ideas are deliberately general: every subsequent unit of the
textbook re-encounters them. Feedback appears in enzyme regulation (Unit
I), in signalling cascades (Unit II), in metabolic flux (Unit III), in
population dynamics (Unit X), and in neural control (Unit IX).
Attractors describe cell-fate decisions (Unit II), bistable genetic
switches (Unit IV), and ecosystem stable states (Unit X). Active
inference recurs as the unifying framework for physiology (Unit IX) and
behaviour.

\begin{center}\rule{0.5\linewidth}{0.5pt}\end{center}

\subsection*{Current Evidence Thread}\label{current-evidence-thread}
\addcontentsline{toc}{subsection}{Current Evidence Thread}

Systems-science and complexity claims are not evidenced the way a
single-gene knockout is; the evidence is whether a model's assumptions
survive contact with data. A claim earns confidence when its generative
model exposes the parameters and boundary that matter, when a
perturbation or time-series test could have falsified it but did not,
and when an explicit null or alternative model fails where it succeeds.
Across this unit --- emergence, adaptive agents, and active inference
--- read each idea as such a model and ask what observation would move
you. As you move through the chapters, keep a two-column note:
\textbf{claim} on the left, \textbf{evidence that would change my
confidence} on the right. By the end of the unit, each major idea should
be tied to a measurement, model, citation, or paper-based lab decision.

\subsection*{Chapter Roadmap}\label{chapter-roadmap}
\addcontentsline{toc}{subsection}{Chapter Roadmap}

\begin{figure}[htbp]
\centering
\includegraphics[width=0.82\linewidth,height=4.2in,keepaspectratio]{../figures/mermaid_inline/inline_mermaid_0002_be8d09ab0dd9.png}
\caption{Unit 0 roadmap: systems science introduces feedback and emergence, complex adaptive systems adds attractors and evolution, and active inference connects prediction and allostasis to later units.}
\end{figure}

{\def\LTcaptype{none} % do not increment counter
\begin{longtable}[]{@{}
  >{\raggedright\arraybackslash}p{(\linewidth - 4\tabcolsep) * \real{0.2414}}
  >{\raggedright\arraybackslash}p{(\linewidth - 4\tabcolsep) * \real{0.4483}}
  >{\raggedright\arraybackslash}p{(\linewidth - 4\tabcolsep) * \real{0.3103}}@{}}
\toprule\noalign{}
\begin{minipage}[b]{\linewidth}\raggedright
Chapter
\end{minipage} & \begin{minipage}[b]{\linewidth}\raggedright
Core question
\end{minipage} & \begin{minipage}[b]{\linewidth}\raggedright
Key tools
\end{minipage} \\
\midrule\noalign{}
\endhead
\bottomrule\noalign{}
\endlastfoot
0.1 Systems Science & \emph{How does organisation arise from
interaction?} & Feedback, hierarchy, Hill equation, delay oscillation \\
0.2 Complex Adaptive Systems & \emph{How do simple local rules produce
robust global behaviour?} & Phase space, bifurcation, fitness landscape,
power laws \\
0.3 Active Inference & \emph{How do living agents maintain themselves
against disorder?} & Bayesian inference, free energy, precision,
allostasis \\
\end{longtable}
}

\begin{center}\rule{0.5\linewidth}{0.5pt}\end{center}

\subsection*{Connections Across the
Textbook}\label{connections-across-the-textbook}
\addcontentsline{toc}{subsection}{Connections Across the Textbook}

\begin{itemize}
\tightlist
\item
  \textbf{Unit I --- Chemistry of Life}: thermodynamic gradients (Unit
  0.1) explain why proteins fold spontaneously; catalysis is a
  feedback-regulated modular system.
\item
  \textbf{Unit II --- The Cell}: organelles are modules; the
  fluid-mosaic membrane is a self-organised phase; signalling cascades
  are hierarchical feedback loops.
\item
  \textbf{Unit III --- Energy and Metabolism}: glycolysis and the TCA
  cycle are classic examples of allosteric feedback and feed-forward
  control.
\item
  \textbf{Unit IV--V --- Genetics}: regulatory networks, bistable
  switches, and epigenetic memory are CAS par excellence.
\item
  \textbf{Unit VI --- Evolution}: fitness landscapes from
  \cref{sec:unit_0_complex_adaptive_systems} are the formal substrate of
  selection.
\item
  \textbf{Unit IX --- Physiology and Neuroscience}: allostasis from
  \cref{sec:unit_0_active_inference} and predictive coding replace
  simple fixed-set-point homeostasis as the model of central control.
\item
  \textbf{Unit X --- Ecology}: phase transitions, alternative stable
  states, and tipping points apply the mathematics of
  \cref{sec:unit_0_complex_adaptive_systems} directly.
\end{itemize}

\begin{center}\rule{0.5\linewidth}{0.5pt}\end{center}

\subsection*{Computational Toolbox --- Unit
0}\label{computational-toolbox-unit-0}
\addcontentsline{toc}{subsection}{Computational Toolbox --- Unit 0}

The three Unit 0 chapters are conceptual rather than algorithmic, but
they motivate every piece of code used later:

\begin{itemize}
\tightlist
\item
  \textbf{Hill cooperative binding:} cell-biology helpers, especially
  \texttt{hill\_equation()}.
\item
  \textbf{Receptor occupancy:} cell-signalling helpers for ligand
  binding and response.
\item
  \textbf{Logistic dynamics:} ecology helpers for growth under carrying
  capacity.
\item
  \textbf{Lotka--Volterra oscillation:} ecology helpers for
  predator--prey cycles.
\item
  \textbf{Bayesian update:} neuroscience-style posterior inference,
  illustrated in \cref{sec:unit_0_active_inference}.
\end{itemize}

Each concrete chapter later in the textbook either uses one of these
helpers or gives the reader a chance to write the next one.

\begin{center}\rule{0.5\linewidth}{0.5pt}\end{center}

\emph{Source modules: \texttt{src/biology/} (general framework across
most domains).} \emph{Figures: \breaktt{src/mermaid/biology\_diagrams.py}
--- systems-level diagrams reused throughout the book.}

\newpage

\section{Systems Science and the Logic of
Emergence}\label{systems-science-and-the-logic-of-emergence}

\label{sec:unit_0_systems_science}

\begin{quote}
\textbf{Unit 0} \texorpdfstring{\ensuremath{\cdot}}{.} Level 2/3 \texorpdfstring{\ensuremath{\cdot}}{.} 35 min read \texorpdfstring{\ensuremath{\cdot}}{.} 50 min lecture \texorpdfstring{\ensuremath{\cdot}}{.}
Prerequisites: none
\end{quote}

\begin{center}\rule{0.5\linewidth}{0.5pt}\end{center}

\subsection{Learning Objectives}\label{learning-objectives}

By the end of this chapter, students will be able to:

\begin{enumerate}
\def\labelenumi{\arabic{enumi}.}
\tightlist
\item
  Define \emph{system} and identify boundary, components, environment,
  and relationships in biological examples.
\item
  Describe positive and negative feedback loops and give biological
  examples of each.
\item
  Explain what emergence is and why it cannot be predicted from
  component-level analysis alone.
\item
  Apply the concepts of hierarchy, modularity, and scale to cellular and
  organismal biology.
\item
  Distinguish between linear and nonlinear system behaviour.
\item
  Analyse delayed feedback loops and predict whether they will damp,
  oscillate, or destabilise.
\item
  Recognise scale-free network structure in biological data and explain
  its functional implications.
\item
  Apply the Hill equation, the chemostat equations, and simple
  oscillator models to quantitative problems.
\end{enumerate}

\subsubsection{Study Blueprint}\label{study-blueprint}

\begin{itemize}
\tightlist
\item
  \textbf{Big idea:} Biological explanation improves when parts,
  interactions, feedback, and scale are kept in view together.
\item
  \textbf{Core concepts:} systems boundaries, feedback, emergence, state
  variables.
\item
  \textbf{Framework alignment:} Vision \& Change: Systems, Structure and
  function; AP Biology: Systems Interactions; NGSS-style topics:
  Structure and Function, Interdependent Relationships in Ecosystems.
\item
  \textbf{Model or quantitative lens:} Box-and-arrow causal models with
  explicit inputs, outputs, and feedback signs.
\item
  \textbf{Data skill:} Translate a verbal biological system into
  variables, links, and testable predictions.
\item
  \textbf{Practice cadence:} Questions and Methods, Representing and
  Describing Data, Argumentation.
\item
  \textbf{Common misconception to repair:} A system is not just a list
  of parts; the interactions are part of the explanation.
\item
  \textbf{Primary lab:} \cref{sec:lab_unit_0_systems_science}.
\item
  \textbf{Question bank:} \cref{sec:q_unit_0_systems_science}.
\item
  \textbf{Transfer task:} Apply the same feedback map to a cell,
  organism, and ecosystem, then name what changes at each scale.
\item
  \textbf{Bridge to computation:}
  \breaktt{biology.crossref\_validator.validate}.
\end{itemize}

\begin{center}\rule{0.5\linewidth}{0.5pt}\end{center}

\subsection{Opening Vignette: The Birth of Systems
Thinking}\label{opening-vignette-the-birth-of-systems-thinking}

In 1948, Norbert Wiener published \emph{Cybernetics: or Control and
Communication in the Animal and the Machine}, arguing that the same
mathematical principles of feedback and self-regulation govern
thermostats, servomechanisms, and the human nervous system. Two decades
later, Ludwig von Bertalanffy's \emph{General System Theory} (1968)
\citep{bertalanffy1968} proposed that most complex phenomena --- from a
cell to a city --- share comprehensive organisational laws: openness,
hierarchy, equifinality, and emergent order.

Their insight was radical: biology cannot be understood by dissecting
organisms into parts and studying each in isolation. Instead, one must
study \emph{relationships} --- the flows of matter, energy, and
information that connect components and give rise to behaviour no
component possesses alone. When Walter Cannon coined
\textbf{\hyperref[gl:homeostasis]{homeostasis}} in 1932
\citep{cannon1932}, he was describing a systems-level property: the
blood glucose concentration of 5 mM is not a property of any single cell
but an emergent steady-state of millions of feedback interactions
between pancreatic \texorpdfstring{\ensuremath{\beta}}{beta}-cells, hepatocytes, skeletal myocytes, and
adipocytes.

This chapter introduces the core vocabulary and conceptual toolkit of
systems science --- the intellectual scaffolding that unifies every
subsequent chapter in this textbook. We will move from definitions (what
is a system?) through dynamics (feedback, delay, oscillation, chaos) to
architecture (hierarchy, modularity, scale-free networks), then close
with information, entropy, and the practical implications of a systems
perspective for biomedicine.

\begin{center}\rule{0.5\linewidth}{0.5pt}\end{center}

\subsection{What Is a System?}\label{what-is-a-system}

A \textbf{system} is a set of components that interact to produce
collective behaviour. Three elements define any system:

\begin{itemize}
\tightlist
\item
  \textbf{Components} --- the identifiable parts (atoms, molecules,
  cells, organisms, species).
\item
  \textbf{Relationships} --- the interactions and constraints between
  components (chemical bonds, metabolic pathways, predator--prey links).
\item
  \textbf{Boundary} --- the interface that separates the system from its
  environment, through which matter and energy may flow.
\end{itemize}

\subsubsection{Isolated, Closed, and Open
Systems}\label{isolated-closed-and-open-systems}

{\def\LTcaptype{none} % do not increment counter
\begin{longtable}[]{@{}llll@{}}
\toprule\noalign{}
Type & Energy exchange & Matter exchange & Biological occurrence \\
\midrule\noalign{}
\endhead
\bottomrule\noalign{}
\endlastfoot
Isolated & None & None & Idealised primarily --- not found in nature \\
Closed & Yes & No & Rare; some thermal models \\
\textbf{Open} & \textbf{Yes} & \textbf{Yes} & \textbf{Most living
systems} \\
\end{longtable}
}

Most living systems are open systems. They import free energy and matter
(food, sunlight), perform work, and export entropy (heat, excreta). This
continuous throughput is what allows them to maintain internal order ---
a thermodynamic feat that distinguishes the living from the non-living.

\subsubsection{Equifinality and
Multifinality}\label{equifinality-and-multifinality}

Bertalanffy emphasised that open systems often display
\textbf{equifinality}: the same final state can be reached from
different starting conditions and along different trajectories. A
vertebrate embryo robustly produces the species-typical body plan
despite considerable variation in initial cell positions, ploidy, or
even surgical perturbation. Conversely, \textbf{multifinality} describes
the same starting condition giving rise to different outcomes ---
identical twins acquiring distinct microbiomes, for example. Both
properties contradict naive linear cause-and-effect intuition and
motivate the systems-level vocabulary that follows.

\begin{quote}
\textbf{Concept Check 1:} A surgeon removes part of an early sea-urchin
embryo. The remaining cells reorganise to produce a smaller but
otherwise normal larva. Identify which open-system property
(equifinality, multifinality, or neither) this illustrates, and explain
why a watch --- a non-living complicated system --- would not behave
this way.
\end{quote}

\begin{center}\rule{0.5\linewidth}{0.5pt}\end{center}

\subsection{Feedback: The Grammar of
Self-Regulation}\label{feedback-the-grammar-of-self-regulation}

\textbf{Feedback} occurs when a fraction of a system's output is
returned as input, altering subsequent behaviour. Biological regulation
is almost entirely feedback-based.

\subsubsection{Negative (Stabilising)
Feedback}\label{negative-stabilising-feedback}

Negative feedback opposes deviations from a set point. It is the
mechanistic basis of \textbf{homeostasis}.

\textbf{Examples:}

\begin{itemize}
\tightlist
\item
  Body temperature regulation --- warming activates sweating and
  vasodilation, cooling activates shivering and vasoconstriction.
\item
  Blood glucose --- rising glucose triggers insulin secretion, promoting
  uptake; falling glucose triggers glucagon, promoting glycogenolysis.
\item
  \hyperref[gl:gene]{\textbf{Gene}} expression --- product inhibition of
  the first \hyperref[gl:enzyme]{\textbf{enzyme}} in a biosynthetic
  pathway (end-product inhibition).
\item
  Baroreceptor reflex --- a fall in arterial pressure unloads stretch
  receptors in the carotid sinus, increasing sympathetic tone within
  seconds.
\item
  Renin--angiotensin--aldosterone system --- falling renal perfusion
  pressure activates a hormonal cascade that retains sodium and water on
  a slower (minutes-to-hours) timescale.
\end{itemize}

\begin{quote}
\textbf{Connection (clinical) --- glucose, insulin, and type 2 diabetes}
Fasting plasma glucose near \textbf{5 mM} is maintained by a
negative-feedback loop: \texorpdfstring{\ensuremath{\beta}}{beta}-cells release \textbf{insulin} when glucose
rises; insulin increases \hyperref[gl:glut4]{\textbf{GLUT4}}-mediated
uptake in muscle and adipose and restrains hepatic glucose output;
\textbf{glucagon} opposes insulin when glucose falls. In \textbf{type 2
diabetes mellitus}, \textbf{insulin resistance} shifts the loop: the
same insulin concentration produces less effect, so higher glucose is
``needed'' to clear a meal; \texorpdfstring{\ensuremath{\beta}}{beta}-cells compensate until they fail. Drugs
target \emph{nodes} in this network --- \textbf{metformin} (hepatic
gluconeogenesis), \textbf{GLP-1 agonists} (incretin axis), \textbf{SGLT2
inhibitors} (renal glucose excretion) --- illustrating how systems
biology informs combination therapy rather than single-molecule ``magic
bullets.''
\end{quote}

\subsubsection{A Linear Negative-Feedback Worked
Example}\label{a-linear-negative-feedback-worked-example}

Consider a regulator with set point \(x^* = 5\) mM, current state \(x\),
and proportional response

\begin{equation}
\frac{dx}{dt} = -k\,(x - x^{*}) + d(t)
\label{eq:unit_0_negative_feedback}
\end{equation}

where \(k > 0\) is the gain and \(d(t)\) is a disturbance. With
\(k = 0.5\) min\(^{-1}\), an injection that lifts \(x\) to 8 mM at
\(t = 0\) produces an exponential return
\(x(t) - x^{*} = 3\,e^{-0.5\,t}\) mM. The deviation halves every
\(\ln 2 / k \approx 1.4\) minutes. Doubling the gain to \(k = 1\)
min\(^{-1}\) halves the half-time to 0.7 minutes. \textbf{Higher gain
returns the system faster}, but delayed feedback can turn excessive gain
into oscillation or instability. This trade-off --- speed versus
stability --- is the central design problem of every physiological
control system.

\begin{figure}[htbp]
\centering
\includegraphics[width=0.82\linewidth,height=4.2in,keepaspectratio]{../figures/mermaid_inline/inline_mermaid_0003_db74273bd95c.png}
\caption{Canonical negative-feedback control loop: set point to effector to output, with the detector returning an error signal that opposes deviation.}
\end{figure}

\subsubsection{Positive (Amplifying)
Feedback}\label{positive-amplifying-feedback}

Positive feedback amplifies deviations. It underlies rapid, switch-like
biological transitions.

\textbf{Examples:}

\begin{itemize}
\tightlist
\item
  \hyperref[gl:action-potential]{\textbf{Action potential}} initiation
  --- sodium influx depolarises membrane, opening more Na\texorpdfstring{\textsuperscript{+}}{+} channels
  (Hodgkin cycle).
\item
  Childbirth contraction --- oxytocin stimulates contractions, which
  stimulate more oxytocin release.
\item
  \hyperref[gl:apoptosis]{\textbf{Apoptosis}} --- executioner
  \hyperref[gl:caspase]{\textbf{caspase}}s activate upstream caspases,
  accelerating cell death.
\item
  Blood-clotting cascade --- thrombin generates more thrombin via factor
  V and factor VIII activation, producing exponential burst once a
  threshold is crossed.
\item
  LH surge --- rising oestradiol switches from negative to positive
  feedback on the hypothalamic-pituitary axis, triggering ovulation.
\end{itemize}

Positive feedback systems are inherently unstable unless ultimately
bounded by a negative feedback or a hard limit (refractory period,
resource exhaustion, anatomical constraint).

\begin{figure}[htbp]
\centering
\includegraphics[width=0.82\linewidth,height=4.2in,keepaspectratio]{../figures/mermaid_inline/inline_mermaid_0004_02999862bd7e.png}
\caption{Canonical positive-feedback loop driving rapid state transitions.}
\end{figure}

\begin{quote}
\textbf{Concept Check 2:} Childbirth and the action potential are both
positive-feedback processes, yet neither runs away forever. Identify the
\emph{bounding} mechanism in each case (anatomical, refractory, or
chemical) and explain why pure positive feedback without an explicit
bound would be lethal.
\end{quote}

\subsubsection{Feedback With Delay: Period and
Damping}\label{feedback-with-delay-period-and-damping}

Real biology is rarely instantaneous. Gene expression requires
\hyperref[gl:transcription]{\textbf{transcription}} and
\hyperref[gl:translation]{\textbf{translation}} (minutes to hours);
vascular responses require perfusion (seconds); neural reflexes require
axonal conduction (milliseconds). A delay \texorpdfstring{\ensuremath{\tau}}{tau} between detecting a
perturbation and responding to it fundamentally changes feedback
behaviour:

\begin{equation}
\frac{dx}{dt} = -k\,x(t - \tau)
\label{eq:unit_0_delay}
\end{equation}

\begin{itemize}
\tightlist
\item
  When \(k\tau \ll 1\): the system returns smoothly to equilibrium
  (overdamped).
\item
  When \(k\tau \approx 1/2\): damped oscillations appear --- the system
  overshoots, undershoots, and asymptotes over several cycles.
\item
  When \(k\tau > \pi/2\): the damping turns negative and the system
  breaks into \textbf{sustained oscillations} with angular frequency
  \(\omega \approx \sqrt{k/\tau}\).
\end{itemize}

This last regime is where biological \textbf{clocks} live. Circadian
oscillators (PER/CRY feedback, about 24 h), cardiac pacemakers
(HCN channel gating, about 1 s), and cell-cycle oscillators
(cyclin--CDK, about 24 h) most exploit delayed negative
feedback to generate rhythms. The period is approximately twice the
transcriptional and translational delay plus degradation lag.

\textbf{Worked Example --- Circadian period from delay.} Suppose
mammalian \emph{PER} transcription takes 45 min, mRNA export and
translation 30 min, and PER \hyperref[gl:protein]{\textbf{protein}}
degradation has half-life 6 h. The total negative-feedback delay is
\(\tau \approx 7.5\) h, and because oscillation period
\(\approx 2\pi\tau/\sqrt{k\tau}\) the characteristic rhythm lands near
24 h for plausible rate constants.
\hyperref[gl:mutation]{\textbf{Mutation}}s that shorten PER stability
(e.g.~familial advanced sleep-phase syndrome, \emph{PER2} S662G)
empirically shorten the period to about 20 h --- consistent
with reducing \texorpdfstring{\ensuremath{\tau}}{tau}. This is why ``biological clocks'' are as much about
delays as about reaction rates.

\subsubsection{Feed-Forward and Anticipatory
Control}\label{feed-forward-and-anticipatory-control}

Purely reactive feedback can primarily correct an error after it has
developed. Many biological systems add a \textbf{feed-forward} pathway
that uses an upstream signal to predict and pre-empt the disturbance.
Examples:

\begin{itemize}
\tightlist
\item
  The cephalic phase of digestion releases insulin before glucose enters
  the bloodstream, smoothing post-prandial spikes.
\item
  Sensory cortex anticipates expected stimuli and subtracts prediction
  from input, improving signal-to-noise.
\item
  Cerebellar internal models generate predicted motor outcomes before
  execution, allowing correction mid-movement.
\end{itemize}

Feed-forward without feedback is brittle --- any mis-prediction persists
uncorrected --- so real biology layers the two. This is the formal
bridge to \textbf{allostasis} (\cref{sec:unit_0_active_inference}):
predictive control of internal state with dynamic set-points,
anticipating rather than merely correcting.

\begin{quote}
\textbf{Connection (clinical) --- diabetic insulin pumps.} Hybrid
closed-loop pumps combine continuous glucose monitoring (feedback) with
announcements of meals or exercise (feed-forward). Pumps that respond
purely to measured glucose chronically lag behind meals because gut
absorption dynamics are faster than subcutaneous insulin uptake; adding
a feed-forward ``meal bolus'' reduces post-prandial excursions by
30--50\% without raising hypoglycaemia rates. The engineering principle
is identical to the cephalic phase the body evolved millions of years
ago.
\end{quote}

\begin{center}\rule{0.5\linewidth}{0.5pt}\end{center}

\subsection{Emergence}\label{emergence}

\textbf{Emergence} refers to properties of a system that cannot be
explained solely by properties of its components. Emergent properties
are \emph{relational} --- they arise from the pattern of interaction,
not from any single part.

\begin{quote}
``The whole is more than the sum of its parts.'' --- Aristotle,
\emph{Metaphysics} (paraphrase)
\end{quote}

\subsubsection{Levels of Biological
Emergence}\label{levels-of-biological-emergence}

{\def\LTcaptype{none} % do not increment counter
\begin{longtable}[]{@{}
  >{\raggedright\arraybackslash}p{(\linewidth - 4\tabcolsep) * \real{0.1471}}
  >{\raggedright\arraybackslash}p{(\linewidth - 4\tabcolsep) * \real{0.5588}}
  >{\raggedright\arraybackslash}p{(\linewidth - 4\tabcolsep) * \real{0.2941}}@{}}
\toprule\noalign{}
\begin{minipage}[b]{\linewidth}\raggedright
Level
\end{minipage} & \begin{minipage}[b]{\linewidth}\raggedright
Emergent properties
\end{minipage} & \begin{minipage}[b]{\linewidth}\raggedright
Components
\end{minipage} \\
\midrule\noalign{}
\endhead
\bottomrule\noalign{}
\endlastfoot
Molecular & Enzyme catalysis, membrane fluidity & Atoms and bonds \\
Cellular & Metabolism, replication, signalling &
\hyperref[gl:organelle]{\textbf{Organelle}}s and molecules \\
Tissue & Contractility, conductance, secretion & Cells \\
Organism & Consciousness, immunity, homeostasis & Tissues and organs \\
Population & Heredity, selection, drift & Individual organisms \\
Ecosystem & Nutrient cycling, energy flow & Populations and
\hyperref[gl:abiotic]{\textbf{abiotic}} factors \\
\end{longtable}
}

\subsubsection{Strong and Weak
Emergence}\label{strong-and-weak-emergence}

Philosophers distinguish:

\begin{itemize}
\tightlist
\item
  \textbf{Weak emergence} --- properties are novel relative to
  component-level description but \emph{in principle} derivable from it
  (e.g., liquidity from molecular interactions).
\item
  \textbf{Strong emergence} --- properties that cannot, even in
  principle, be reduced to component interaction (consciousness remains
  a contested candidate).
\end{itemize}

For practical biology, the key insight is that \emph{explaining
components is not the same as explaining the system}. A complete
molecular catalogue of a neuron does not explain perception; a complete
genome does not specify the body plan without the cytoplasmic and
developmental context that interprets it.

\begin{quote}
\textbf{Concept Check 3:} Liquidity is often given as a paradigm of weak
emergence: it is not a property of any single H\texorpdfstring{\textsubscript{2}}{2}O molecule but is in
principle derivable from the molecular interaction Hamiltonian. Identify
\emph{one} biological property that is plausibly weakly emergent and
\emph{one} that is contested as possibly strongly emergent. What would
it take to upgrade the contested property from ``we cannot yet derive
it'' to ``it is in principle underivable''?
\end{quote}

\begin{center}\rule{0.5\linewidth}{0.5pt}\end{center}

\subsection{Hierarchy and Scale}\label{hierarchy-and-scale}

Biological organisation is \textbf{hierarchical}. Each level exhibits
emergent properties relative to the level below, and is \emph{nested}
within levels above.

\begin{figure}[htbp]
\centering
\includegraphics[width=0.82\linewidth,height=4.2in,keepaspectratio]{../figures/mermaid_inline/inline_mermaid_0005_406d75f30ab9.png}
\caption{Nested hierarchy of biological organisation, from atoms and molecules up through cells and organisms to the biosphere, each level emergent from the one below.}
\end{figure}

\subsubsection{Cross-Scale Constraints}\label{cross-scale-constraints}

Higher-level organisation \emph{constrains} lower-level behaviour
(downward causation), while lower-level mechanisms \emph{generate}
higher-level properties (upward causation). For example:

\begin{itemize}
\tightlist
\item
  A cell's membrane potential (tissue-level context) constrains which
  genes are expressed.
\item
  A mutation (molecular level) can alter the organismal
  \hyperref[gl:phenotype]{\textbf{phenotype}} and hence population
  fitness.
\item
  Tissue architecture constrains diffusion gradients that, in turn,
  pattern gene expression (developmental biology's central feedback).
\end{itemize}

\subsubsection{Modularity}\label{modularity}

Living systems are modular --- composed of semi-independent subsystems
with defined interfaces. Modularity confers:

\begin{itemize}
\tightlist
\item
  \textbf{Robustness}: a failure in one module need not propagate to
  others.
\item
  \textbf{Evolvability}: modules can be rewired or duplicated with
  limited pleiotropic effects.
\item
  \textbf{Decomposability}: scientists can study modules in isolation
  (the logic underlying reductionism's power).
\end{itemize}

\begin{quote}
\textbf{Concept Check 4:} The MAPK cascade (Ras \texorpdfstring{\ensuremath{\to}}{->} Raf \texorpdfstring{\ensuremath{\to}}{->} MEK \texorpdfstring{\ensuremath{\to}}{->} ERK) is
often described as a module with a well-defined input (Ras activation)
and output (ERK phosphorylation). List two biological consequences of
this modularity --- one that is \emph{evolutionarily} useful (robust to
rewiring) and one that is \emph{clinically} worrisome (a single drug
target that blocks many cell types).
\end{quote}

\begin{center}\rule{0.5\linewidth}{0.5pt}\end{center}

\subsection{Nonlinearity and
Thresholds}\label{nonlinearity-and-thresholds}

Simple systems obey \emph{linear} relationships: double the input,
double the output. Most biological systems are \textbf{nonlinear} ---
small changes can produce large effects, or large changes may produce
negligible effects.

Key nonlinear phenomena in biology:

{\def\LTcaptype{none} % do not increment counter
\begin{longtable}[]{@{}
  >{\raggedright\arraybackslash}p{(\linewidth - 4\tabcolsep) * \real{0.3571}}
  >{\raggedright\arraybackslash}p{(\linewidth - 4\tabcolsep) * \real{0.3929}}
  >{\raggedright\arraybackslash}p{(\linewidth - 4\tabcolsep) * \real{0.2500}}@{}}
\toprule\noalign{}
\begin{minipage}[b]{\linewidth}\raggedright
Phenomenon
\end{minipage} & \begin{minipage}[b]{\linewidth}\raggedright
Description
\end{minipage} & \begin{minipage}[b]{\linewidth}\raggedright
Example
\end{minipage} \\
\midrule\noalign{}
\endhead
\bottomrule\noalign{}
\endlastfoot
Threshold effects & Response absent below a critical input, present
above & Action potential firing threshold (about \texorpdfstring{\ensuremath{-}}{-}55 mV) \\
Bistability & System flips between two stable states & Lac
\hyperref[gl:operon]{\textbf{operon}} switch; cell-cycle entry \\
Hysteresis & History-dependent state & Epigenetic gene silencing \\
Oscillations & Sustained rhythmic output & Circadian clocks; cardiac
pacemakers \\
Chaos & Aperiodic but deterministic dynamics & Some cardiac arrhythmias;
population cycles \\
\end{longtable}
}

\subsubsection{A Simple Nonlinear Equation: The Hill
Function}\label{a-simple-nonlinear-equation-the-hill-function}

Cooperative binding (e.g., haemoglobin--oxygen, transcription factors)
follows the \textbf{Hill equation}:

\begin{equation}
\theta = \frac{[L]^n}{K_d^n + [L]^n}
\label{eq:unit_0_hill}
\end{equation}

where \texorpdfstring{\ensuremath{\theta}}{theta} is fractional occupancy, \([L]\) is ligand concentration,
\(K_d\) is the dissociation constant, and \(n\) is the Hill coefficient.
For \(n > 1\), binding is cooperative, producing a sigmoidal
(switch-like) response. The slope at half-occupancy scales with \(n\),
so a Hill coefficient of \(n = 4\) produces a transition that is roughly
four times sharper than a \(n = 1\) saturable response.

\subsubsection{Worked Example: Calculating Cooperative
Binding}\label{worked-example-calculating-cooperative-binding}

\textbf{Problem:} A dimeric transcription factor binds to DNA
cooperatively with a Hill coefficient of \(n = 2\). The dissociation
constant \(K_d\) for the DNA binding site is \(5 \text{ nM}\). What
fraction of target \hyperref[gl:promoter]{\textbf{promoter}}s (\texorpdfstring{\ensuremath{\theta}}{theta}) will
be bound by the transcription factor when its intracellular
concentration is \(10 \text{ nM}\)? What would the occupancy be if
binding were non-cooperative (\(n = 1\))?

\textbf{Solution:}

\begin{enumerate}
\def\labelenumi{\arabic{enumi}.}
\item
  \textbf{Calculate occupancy for cooperative binding (\(n = 2\)):}

  Using the Hill equation:
  \[ \theta = \frac{[L]^n}{K_d^n + [L]^n}  \label{eq:unit_0_systems_science_item_1}\]

  Substitute \([L] = 10\) and \(K_d = 5\):
  \[ \theta = \frac{10^2}{5^2 + 10^2} = \frac{100}{25 + 100} = \frac{100}{125} = 0.80  \label{eq:unit_0_systems_science_item_2}\]

  The transcription factor occupies \textbf{80\%} of the promoters.
\item
  \textbf{Calculate occupancy for non-cooperative binding (\(n = 1\)):}

  \[ \theta = \frac{10^1}{5^1 + 10^1} = \frac{10}{5 + 10} = \frac{10}{15} \approx 0.67  \label{eq:unit_0_systems_science_item_3}\]

  Without cooperativity, primarily \textbf{67\%} of promoters are bound.
\end{enumerate}

This calculation demonstrates how positive cooperativity steepens the
dose-response curve, allowing a system to switch from off (unbound) to
on (bound) over a narrower range of ligand concentrations, underpinning
nonlinear threshold effects in biology.

\subsubsection{Bistability and the Toggle
Switch}\label{bistability-and-the-toggle-switch}

Combine two mutually inhibiting genes and you obtain a \textbf{toggle
switch} with two stable attractors and an unstable separator.
Mathematically the rate equations look like

\begin{equation}
\frac{du}{dt} = \frac{\alpha_1}{1 + v^{n}} - u, \qquad \frac{dv}{dt} = \frac{\alpha_2}{1 + u^{n}} - v
\label{eq:unit_0_toggle}
\end{equation}

For \(n > 1\) and balanced production rates
\(\alpha_1 \approx \alpha_2\), the system has three fixed points:
\((u_{\mathrm{high}}, v_{\mathrm{low}})\),
\((u_{\mathrm{low}}, v_{\mathrm{high}})\), and an unstable saddle in
between. Each cell commits to one of the two outcomes as a function of
initial conditions and noise --- the molecular embodiment of cell-fate
decision-making. Variations on this motif underlie phage \texorpdfstring{\ensuremath{\lambda}}{lambda}
lysis-vs-lysogeny, the lac operon's bistable response to lactose, and
embryonic stem-cell differentiation.

\begin{quote}
\textbf{Concept Check 5:} A biotechnology company engineers a synthetic
toggle switch into \emph{E. coli} and observes that, when both inducer
ligands are equal, roughly half the population sits in state A and half
in state B. Why does the population split rather than every cell
choosing the same state? Sketch the relevant phase portrait and identify
the role of stochastic gene expression noise.
\end{quote}

\begin{center}\rule{0.5\linewidth}{0.5pt}\end{center}

\subsection{Nonlinear Dynamics: Bifurcations, Limit Cycles, and
Chaos}\label{nonlinear-dynamics-bifurcations-limit-cycles-and-chaos}

Sections 2 and 5 hinted at oscillations and bistability. Nonlinear
dynamics \citep{strogatz2018} provides the formal vocabulary for these
and other behaviours.

\subsubsection{Phase Space and
Trajectories}\label{phase-space-and-trajectories}

The \textbf{phase space} of a system is the space of possible internal
states. A two-variable system (e.g.~predator--prey densities, or a
gene's mRNA and protein abundance) has a 2-D phase plane in which each
system trajectory is a curve. \textbf{Fixed points} are states where the
derivatives vanish; \textbf{limit cycles} are closed trajectories that
nearby trajectories spiral towards or away from.

\subsubsection{Bifurcations}\label{bifurcations}

A \textbf{bifurcation} is a qualitative change in the structure of phase
space as a control parameter varies. Three biologically important kinds:

\begin{itemize}
\tightlist
\item
  \textbf{Saddle-node (fold) bifurcation} --- two fixed points (one
  stable, one unstable) collide and annihilate. Underlies catastrophic
  transitions in lake eutrophication and ecological regime shifts.
\item
  \textbf{Hopf bifurcation} --- a stable fixed point loses stability and
  is replaced by a small-amplitude limit cycle. Underlies the
  \emph{birth} of biological oscillations (cardiac pacing, circadian
  rhythm onset in development).
\item
  \textbf{Pitchfork bifurcation} --- one fixed point splits into three,
  generating bistability. Underlies cell-fate commitment and
  morphogenetic symmetry breaking.
\end{itemize}

\subsubsection{Limit Cycles and Biological
Oscillators}\label{limit-cycles-and-biological-oscillators}

A \textbf{limit cycle} is a closed trajectory toward which neighbouring
trajectories converge. Biological oscillators that are robustly
limit-cycle:

\begin{itemize}
\tightlist
\item
  \textbf{Circadian clocks} (about 24 h period) --- TTFLs
  (transcription--translation feedback loops) involving
  \emph{CLOCK/BMAL1} and \emph{PER/CRY} in mammals; \emph{KaiABC} in
  cyanobacteria --- the latter operates without transcription,
  demonstrating that a phosphorylation cycle alone can implement a
  biological clock.
\item
  \textbf{Cardiac pacemaker} (about 1 s) --- sino-atrial node
  cells use HCN (``funny'') current and Ca\texorpdfstring{\textsuperscript{2}}{2}\texorpdfstring{\textsuperscript{+}}{+} cycling; their limit cycle
  is robust to \texorpdfstring{\ensuremath{\pm}}{+/-}20 \% perturbations of any single ionic conductance.
\item
  \textbf{Cell-cycle oscillator} (about 24 h) --- cyclin
  synthesis and APC-mediated cyclin destruction; phosphorylation of CDK
  targets reaches threshold, triggers mitosis, then resets the loop.
\item
  \textbf{Glycolytic oscillations} in yeast (about 30 s--10
  min) --- feedback between phosphofructokinase and downstream
  metabolites.
\end{itemize}

\subsubsection{Chaos and the Lorenz
Attractor}\label{chaos-and-the-lorenz-attractor}

When a deterministic nonlinear system has at least three dynamical
variables, sustained motion can be \textbf{chaotic} --- bounded,
aperiodic, and exquisitely sensitive to initial conditions. The Lorenz
system

\begin{equation}
\dot{x} = \sigma(y - x), \qquad \dot{y} = x(\rho - z) - y, \qquad \dot{z} = xy - \beta z
\label{eq:unit_0_lorenz}
\end{equation}

was originally derived as a toy model of atmospheric convection but
applies (with reinterpreted variables) to laser dynamics, glycolysis,
and population ecology. Trajectories are confined to a butterfly-shaped
\textbf{strange attractor} with non-integer (fractal) dimension. Two
trajectories that start within \(10^{-6}\) of each other diverge until
they separate by order unity within a finite time set by the system's
leading \textbf{Lyapunov exponent} \texorpdfstring{\ensuremath{\lambda}}{lambda}: the separation grows like
\(\delta(t) \approx \delta_0 e^{\lambda t}\).

Biological chaos is documented (with varying confidence) in:

\begin{itemize}
\tightlist
\item
  Polymorphic ventricular tachycardia (``torsades de pointes'' exhibits
  chaotic cardiac dynamics).
\item
  Some EEG dynamics in seizure transitions.
\item
  Population cycles of the Canadian lynx and Soay sheep, modulated by
  weather and density-dependent disease.
\item
  Single-neuron interspike intervals in some cortical recordings.
\end{itemize}

\begin{quote}
\textbf{Connection (clinical) --- defibrillation as state-space reset.}
Ventricular fibrillation can be modelled as the heart's pacemaker
leaving its limit cycle and entering a chaotic regime in which thousands
of small re-entrant circuits prevent coordinated contraction. A
defibrillator delivers a brief, high-energy shock that depolarises
essentially every myocyte simultaneously, ``resetting'' the system to
the basin of attraction of the normal sinus-rhythm limit cycle. The
intervention is a phase-space \emph{kick}, not a chemical correction ---
pure dynamical-systems medicine.
\end{quote}

\begin{quote}
\textbf{Concept Check 6:} A phase-plane portrait shows a stable spiral
that, as a parameter \texorpdfstring{\ensuremath{\mu}}{mu} is increased past a critical value \(\mu_c\),
becomes encircled by a small closed orbit. (a) Name the bifurcation. (b)
Predict whether the resulting oscillation amplitude depends linearly or
quadratically on \(\mu - \mu_c\) near the bifurcation. (c) Name one
cardiac or circadian disorder consistent with a Hopf bifurcation gone
wrong.
\end{quote}

\begin{center}\rule{0.5\linewidth}{0.5pt}\end{center}

\subsection{Biological Oscillators: From Genes to
Heartbeats}\label{biological-oscillators-from-genes-to-heartbeats}

Biological oscillators are textbook applications of delayed negative
feedback (\cref{eq:unit_0_delay}) and limit-cycle dynamics.

\subsubsection{The Circadian Clock}\label{the-circadian-clock}

The mammalian circadian clock is a \textbf{transcription--translation
feedback loop} (TTFL). CLOCK and BMAL1 heterodimerise and activate
transcription of \emph{Per1}, \emph{Per2}, \emph{Cry1}, and \emph{Cry2}.
PER and CRY proteins accumulate, translocate to the nucleus, and inhibit
CLOCK/BMAL1, closing a negative-feedback loop with a delay set by
transcription, translation, dimerisation, and post-translational
modification (notably casein-kinase-1\texorpdfstring{\ensuremath{\varepsilon}}{epsilon} phosphorylation of PER).
Mutations that alter CK1\texorpdfstring{\ensuremath{\varepsilon}}{epsilon} or PER2 phosphorylation sites shift the
clock's period from the wild-type about 24 h to as little as 20
h or as much as 28 h, providing molecular evidence for the delay-period
relationship of \cref{eq:unit_0_delay}.

\subsubsection{The Cardiac Pacemaker}\label{the-cardiac-pacemaker}

Sino-atrial node cells lack a true resting potential. Instead, their
membrane potential oscillates between roughly \texorpdfstring{\ensuremath{-}}{-}60 mV (maximum diastolic
potential) and 0 mV (peak of action potential). Three ionic currents ---
funny current \(I_f\), T-type Ca\texorpdfstring{\textsuperscript{2}}{2}\texorpdfstring{\textsuperscript{+}}{+} current \(I_{Ca,T}\), and L-type Ca\texorpdfstring{\textsuperscript{2}}{2}\texorpdfstring{\textsuperscript{+}}{+}
current \(I_{Ca,L}\) --- operate on different timescales to generate the
pacemaker limit cycle, which is robust to \texorpdfstring{\ensuremath{\pm}}{+/-}20 \% changes in any one
conductance.

\subsubsection{The Cell Cycle}\label{the-cell-cycle}

In rapidly dividing cells the cyclin--CDK system is the master
oscillator. M-cyclin accumulates linearly during interphase; once a
threshold is reached, M-CDK activity is amplified by a positive-feedback
bistable switch (Wee1 inactivation, Cdc25 activation), driving entry
into mitosis. APC/C-mediated cyclin destruction collapses M-CDK
activity, the bistable switch flips back, and the cycle resets. This is
a \textbf{relaxation oscillator}: a slow build-up followed by a fast
release, mathematically distinct from the harmonic oscillation of the
circadian clock.

\subsubsection{Common Architecture}\label{common-architecture}

{\def\LTcaptype{none} % do not increment counter
\begin{longtable}[]{@{}
  >{\raggedright\arraybackslash}p{(\linewidth - 4\tabcolsep) * \real{0.3226}}
  >{\raggedright\arraybackslash}p{(\linewidth - 4\tabcolsep) * \real{0.1935}}
  >{\raggedright\arraybackslash}p{(\linewidth - 4\tabcolsep) * \real{0.4839}}@{}}
\toprule\noalign{}
\begin{minipage}[b]{\linewidth}\raggedright
Oscillator
\end{minipage} & \begin{minipage}[b]{\linewidth}\raggedright
Period
\end{minipage} & \begin{minipage}[b]{\linewidth}\raggedright
Mechanism class
\end{minipage} \\
\midrule\noalign{}
\endhead
\bottomrule\noalign{}
\endlastfoot
Circadian (PER/CRY) & about 24 h & Delayed transcriptional
negative feedback \\
Cardiac (SAN) & about 1 s & Coupled ionic currents, limit
cycle \\
Cell cycle (cyclin--CDK) & hours--days & Relaxation oscillator with
bistable switch \\
Glycolysis (yeast PFK) & seconds--minutes & Allosteric metabolic
feedback \\
Calcium spikes & seconds & IP3R/RyR feedback with diffusion \\
\end{longtable}
}

The lesson: oscillators are not exotic; they are nature's default
solution to the problem of \emph{temporal coordination}, and the same
handful of feedback architectures appears across kingdoms and
timescales.

\begin{center}\rule{0.5\linewidth}{0.5pt}\end{center}

\subsection{Scale-Free Networks in
Biology}\label{scale-free-networks-in-biology}

Most biological systems are not random networks but \textbf{scale-free}:
their node-degree distribution follows a power law
\(P(k) \propto k^{-\gamma}\) with exponent \texorpdfstring{\ensuremath{\gamma}}{gamma} typically between 2 and 3
(\cref{sec:unit_0_complex_adaptive_systems} discusses power laws in
depth). A few \textbf{hub} nodes have many connections; most nodes have
few.

Biologically validated examples include:

\begin{itemize}
\tightlist
\item
  \textbf{Protein--protein interaction networks} of yeast, fly, and
  human (\texorpdfstring{\ensuremath{\gamma}}{gamma} \texorpdfstring{\ensuremath{\approx}}{~} 2.5).
\item
  \textbf{Metabolic networks} --- substrates such as ATP, NADH, and
  water are hubs connecting hundreds of reactions.
\item
  \textbf{Gene-regulatory networks} --- master regulators like
  \emph{p53} or \emph{MYC} control hundreds of downstream targets.
\item
  \textbf{Food webs} --- apex species with broad diets contrast with
  specialist links.
\item
  \textbf{Brain connectomes} --- cortical ``rich-club'' hubs
  disproportionately interconnect distant brain regions.
\end{itemize}

\subsubsection{Why Scale-Free?}\label{why-scale-free}

Two mechanisms generate scale-free architecture:

\begin{enumerate}
\def\labelenumi{\arabic{enumi}.}
\tightlist
\item
  \textbf{Preferential attachment} --- new nodes preferentially connect
  to already-popular nodes (Barabási-Albert model). In biology this maps
  onto \textbf{gene duplication}: a duplicated gene inherits its
  parent's interaction partners, so highly-connected proteins beget more
  highly-connected proteins.
\item
  \textbf{Optimisation under cost-and-benefit} --- connections have
  wiring cost; a few high-degree hubs minimise total path length while
  limiting cost.
\end{enumerate}

\subsubsection{Functional Consequences}\label{functional-consequences}

{\def\LTcaptype{none} % do not increment counter
\begin{longtable}[]{@{}
  >{\raggedright\arraybackslash}p{(\linewidth - 4\tabcolsep) * \real{0.2000}}
  >{\raggedright\arraybackslash}p{(\linewidth - 4\tabcolsep) * \real{0.3500}}
  >{\raggedright\arraybackslash}p{(\linewidth - 4\tabcolsep) * \real{0.4500}}@{}}
\toprule\noalign{}
\begin{minipage}[b]{\linewidth}\raggedright
Property
\end{minipage} & \begin{minipage}[b]{\linewidth}\raggedright
Random network
\end{minipage} & \begin{minipage}[b]{\linewidth}\raggedright
Scale-free network
\end{minipage} \\
\midrule\noalign{}
\endhead
\bottomrule\noalign{}
\endlastfoot
Robustness to random failure & Fragile & Robust (random hits usually
miss hubs) \\
Robustness to targeted attack on hubs & Robust & Fragile \\
Path length between random nodes & \(\propto \ln N\) &
\(\propto \ln \ln N\) (ultra-small) \\
Clustering coefficient & Low & High \\
\end{longtable}
}

This has direct biomedical implications. \emph{Random} genetic
perturbations (background mutation rates) usually miss hubs, so cells
tolerate them. But \emph{targeted} perturbation of hub proteins ---
exactly what cancer-driver mutations and many viral effectors do --- can
cripple the network. The same logic underlies the success of TP53, MYC,
and KRAS as cancer drivers: they sit at network hubs.

\begin{quote}
\textbf{Connection (medical) --- antibiotic targets and the metabolic
hub.} Many successful antibiotics target hub enzymes in bacterial
metabolism (e.g.~RNA polymerase, ribosome, dihydrofolate reductase).
Hubs are tempting because perturbing one node disables many downstream
pathways; they are also dangerous because human homologues, where they
exist, suffer comparable disruption --- explaining the side-effect
profiles of trimethoprim or rifampin.
\end{quote}

\begin{center}\rule{0.5\linewidth}{0.5pt}\end{center}

\subsection{Information, Entropy, and
Self-Organisation}\label{information-entropy-and-self-organisation}

Living systems process \textbf{information}: they detect signals
(genetic, chemical, mechanical, photonic) and use that information to
generate ordered responses. From a thermodynamic perspective:

\begin{equation}
\Delta G = \Delta H - T\Delta S
\label{eq:unit_0_gibbs}
\end{equation}

Life builds local order (decreases local entropy) by doing work, and in
so doing increases entropy of the surroundings. Self-organisation ---
spontaneous formation of ordered structures --- occurs when energy
dissipation enables pattern formation (e.g., Bénard cells, Turing
patterns in morphogenesis, the spiral waves of slime-mould aggregation).

\subsubsection{Shannon Information}\label{shannon-information}

Claude Shannon defined the information content of a discrete probability
distribution as

\begin{equation}
H(X) = -\sum_{i} p_i \log_2 p_i
\label{eq:unit_0_shannon}
\end{equation}

with units of bits when the logarithm is base 2. A signalling pathway
that distinguishes ``ligand absent'' from ``ligand present'' carries at
most one bit; a four-state developmental switch carries at most two; a
transcription factor that selects among 1000 target promoters carries up
to about 10 bits. Empirical studies of cytokine signalling
pathways suggest mammalian cells extract on the order of 1 bit per
pathway from ligand concentration --- far less than the upper bound,
because biological noise eats most of the channel capacity.

\subsubsection{Mutual Information and Signalling
Fidelity}\label{mutual-information-and-signalling-fidelity}

The \textbf{mutual information} \(I(X; Y) = H(Y) - H(Y \mid X)\)
measures how much knowing the input \(X\) reduces uncertainty about
output \(Y\). For a noisy biological channel,

\begin{equation}
I(X; Y) = \frac{1}{2} \log_2 \!\left( 1 + \frac{\sigma_{\mathrm{signal}}^{2}}{\sigma_{\mathrm{noise}}^{2}} \right)
\label{eq:unit_0_mutual_info}
\end{equation}

Doubling the signal-to-noise ratio adds about half a bit. The
diminishing returns of this expression explain why cells often build
\emph{parallel} signalling channels (multiple receptor tyrosine kinases,
multiple cytokine receptors) rather than refining one channel
indefinitely --- additional channels add capacity additively, not
logarithmically.

\begin{quote}
\textbf{Concept Check 7:} A yeast cell at steady state maintains an
intracellular K\texorpdfstring{\textsuperscript{+}}{+} concentration of 140 mM against an extracellular
concentration of 5 mM. From the perspective of the second law of
\hyperref[gl:thermodynamics]{\textbf{thermodynamics}}, is this local
concentration gradient a violation of entropy increase? Identify the
ultimate entropy source and the thermodynamic machinery (at least two
molecular components) that sustains the gradient.
\end{quote}

\begin{quote}
\textbf{Concept Check 8:} Empirical estimates suggest a single mammalian
cytokine pathway carries about 1 bit of mutual information about ligand
concentration. What does this imply about the distinguishability of
``low,'' ``medium,'' and ``high'' concentrations from a single pathway
alone, and why might inflammation cascades evolve to use multiple
cytokines rather than ever-finer measurements of one?
\end{quote}

\begin{center}\rule{0.5\linewidth}{0.5pt}\end{center}

\subsection{Systems Thinking in Biology: Practical
Implications}\label{systems-thinking-in-biology-practical-implications}

{\def\LTcaptype{none} % do not increment counter
\begin{longtable}[]{@{}
  >{\raggedright\arraybackslash}p{(\linewidth - 2\tabcolsep) * \real{0.2903}}
  >{\raggedright\arraybackslash}p{(\linewidth - 2\tabcolsep) * \real{0.7097}}@{}}
\toprule\noalign{}
\begin{minipage}[b]{\linewidth}\raggedright
Principle
\end{minipage} & \begin{minipage}[b]{\linewidth}\raggedright
Biomedical implication
\end{minipage} \\
\midrule\noalign{}
\endhead
\bottomrule\noalign{}
\endlastfoot
Emergence & Drug targets must be evaluated in system context, not in
isolation \\
Feedback & Blocking one regulatory loop often activates compensatory
loops \\
Hierarchy & A mutation's effect depends on its network context \\
Nonlinearity & Dose--response curves are rarely linear; threshold
effects matter \\
Modularity & Synthetic biology exploits modular parts to engineer novel
circuits \\
Scale-free & Hub targeting is potent but risks broad collateral
damage \\
\end{longtable}
}

\subsubsection{From Hill kinetics to a computational
switch}\label{from-hill-kinetics-to-a-computational-switch}

The earlier nonlinear-threshold discussion introduced the \textbf{Hill
equation} as a model of cooperative binding. The same mathematics
appears when transcription factors bind clustered sites, when oxygen
binds haemoglobin, and when receptors oligomerise. Running the project
code makes the \textbf{threshold} tangible:

\begin{Shaded}
\begin{Highlighting}[]
\ImportTok{from}\NormalTok{ biology.cell.cell\_biology }\ImportTok{import}\NormalTok{ hill\_equation}

\NormalTok{kd, n }\OperatorTok{=} \FloatTok{10.0}\NormalTok{, }\FloatTok{3.0}
\ControlFlowTok{for}\NormalTok{ L }\KeywordTok{in}\NormalTok{ (}\FloatTok{1.0}\NormalTok{, }\FloatTok{5.0}\NormalTok{, }\FloatTok{10.0}\NormalTok{, }\FloatTok{20.0}\NormalTok{):}
\NormalTok{    theta }\OperatorTok{=}\NormalTok{ hill\_equation(ligand\_concentration}\OperatorTok{=}\NormalTok{L, kd}\OperatorTok{=}\NormalTok{kd, hill\_coefficient}\OperatorTok{=}\NormalTok{n)}
    \BuiltInTok{print}\NormalTok{(}\SpecialStringTok{f"[L] = }\SpecialCharTok{\{}\NormalTok{L}\SpecialCharTok{:5.1f\}}\SpecialStringTok{ µM  →  fractional occupancy θ = }\SpecialCharTok{\{}\NormalTok{theta}\SpecialCharTok{:.3f\}}\SpecialStringTok{"}\NormalTok{)}
\end{Highlighting}
\end{Shaded}

For \textbf{n \textgreater{} 1}, a narrow concentration band separates
``mostly off'' from ``mostly on'' --- the molecular implementation of a
\textbf{bistable} or \textbf{switch-like} response that subsequent units
link to gene circuits, metabolism, and neural firing.

\subsubsection{Feed-forward and feedback motifs in gene
networks}\label{feed-forward-and-feedback-motifs-in-gene-networks}

Beyond simple loops, \textbf{network motifs} \citep{alon2019} recur in
transcriptional regulation: \textbf{coherent feed-forward loops} (two
parallel paths to the same target, both activating or both repressing)
can filter noise and enforce delays; \textbf{incoherent feed-forward
loops} can generate pulses or adaptation. These motifs appear thousands
of times in \emph{E. coli} and yeast regulatory graphs. Recognising them
helps you predict how a drug that blocks one edge may reroute flux
through another --- the \textbf{compensatory feedback} that often limits
single-target therapy. Tyson, Chen and Novak \citep{tyson2003} catalogue
the recurring ``sniffers, buzzers, toggles and blinkers'' of cell
biology, each of which corresponds to a small network motif with a
stereotyped temporal signature.

\subsubsection{Why systems biology is not ``fancy
reductionism''}\label{why-systems-biology-is-not-fancy-reductionism}

A naive reading of systems biology is that it builds bigger models of
more parts. The deeper claim is methodological: certain \emph{questions}
(will the loop oscillate? Is the basin deep enough to resist noise? Does
the network tolerate hub deletion?) cannot be answered by any amount of
better single-component data. They require \emph{system-level} metrics
--- gain, delay, basin depth, Lyapunov exponent, mutual information ---
which primarily acquire meaning at the level of the assembled system.
Reductionism delivers parts; systems biology delivers grammars.

\begin{center}\rule{0.5\linewidth}{0.5pt}\end{center}

\subsection{Current Evidence and Frontier
Biology}\label{current-evidence-and-frontier-biology}

For \textbf{Systems Science and the Logic of Emergence}, frontier
biology belongs inside the evidence logic of the chapter. Systems models
are useful when they expose assumptions, uncertainty, and failure modes
rather than merely producing elegant diagrams. The core reading question
is this: system boundary choice, feedback sign, and scale determine
whether a model explains or hides the biology.

\begin{itemize}
\tightlist
\item
  \textbf{What to verify:} identify the observation, model, assay, or
  dataset that would make the claim stronger or weaker.
\item
  \textbf{What to qualify:} state the scale, organism, cell type,
  environmental condition, or population where the claim is expected to
  hold.
\item
  \textbf{What to compare:} test at least one alternative explanation,
  baseline, or null model before treating the pattern as causal.
\item
  \textbf{What to cite:} distinguish primary evidence, review synthesis,
  public dataset, and institutional guidance; for recent or numeric
  claims, prefer the source closest to the measurement and state what
  has changed since it was published.
\end{itemize}

Treat every model as a claim about mechanism: define the system
boundary, identify the observable that would falsify the model, and
report the uncertainty that would change a decision.

\textbf{Source practice:} Use model-validation sources when available,
and state which observation would falsify the model rather than treating
a diagram as proof.

\subsection{Summary}\label{summary}

\begin{itemize}
\tightlist
\item
  A system is defined by components, relationships, and boundary;
  equifinality and multifinality distinguish living open systems from
  contrived isolated ones.
\item
  Most living systems are open systems that import free energy and
  export entropy.
\item
  Negative feedback achieves homeostasis; positive feedback generates
  rapid transitions; delay turns stabilising loops into oscillators.
\item
  Emergent properties arise from patterns of interaction and cannot be
  reduced to component properties alone.
\item
  Biological organisation is hierarchical, modular, predominantly
  nonlinear, and overwhelmingly scale-free in its network structure.
\item
  Bifurcations classify how qualitative dynamical behaviour changes ---
  saddle-node (catastrophes), Hopf (oscillation onset), pitchfork
  (symmetry-breaking commitment).
\item
  Biological oscillators (circadian, cardiac, cell-cycle) realise
  delayed-feedback limit cycles; chaos appears in some cardiac and
  population dynamics.
\item
  Hill-type cooperativity implements sharp molecular thresholds; the
  toggle switch implements bistable cell-fate decisions; network motifs
  (feed-forward loops) shape dynamics and compensation under
  perturbation.
\item
  Information theory bounds how much a noisy biological channel can
  convey, motivating parallel rather than refined signalling
  architecture.
\end{itemize}

\begin{center}\rule{0.5\linewidth}{0.5pt}\end{center}

\subsection{Key Terms}\label{key-terms}

\textbf{system} \texorpdfstring{\ensuremath{\cdot}}{.} \textbf{emergence} \texorpdfstring{\ensuremath{\cdot}}{.} \textbf{negative feedback} \texorpdfstring{\ensuremath{\cdot}}{.}
\textbf{positive feedback} \texorpdfstring{\ensuremath{\cdot}}{.} \textbf{homeostasis} \texorpdfstring{\ensuremath{\cdot}}{.} \textbf{allostasis}
\texorpdfstring{\ensuremath{\cdot}}{.} \textbf{hierarchy} \texorpdfstring{\ensuremath{\cdot}}{.} \textbf{modularity} \texorpdfstring{\ensuremath{\cdot}}{.} \textbf{nonlinearity} \texorpdfstring{\ensuremath{\cdot}}{.}
\textbf{bistability} \texorpdfstring{\ensuremath{\cdot}}{.} \textbf{bifurcation} \texorpdfstring{\ensuremath{\cdot}}{.} \textbf{limit cycle} \texorpdfstring{\ensuremath{\cdot}}{.}
\textbf{chaos} \texorpdfstring{\ensuremath{\cdot}}{.} \textbf{Lyapunov exponent} \texorpdfstring{\ensuremath{\cdot}}{.} \textbf{scale-free
network} \texorpdfstring{\ensuremath{\cdot}}{.} \textbf{hub} \texorpdfstring{\ensuremath{\cdot}}{.} \textbf{self-organisation} \texorpdfstring{\ensuremath{\cdot}}{.} \textbf{Hill
equation} \texorpdfstring{\ensuremath{\cdot}}{.} \textbf{toggle switch} \texorpdfstring{\ensuremath{\cdot}}{.} \textbf{open system} \texorpdfstring{\ensuremath{\cdot}}{.}
\textbf{network motif} \texorpdfstring{\ensuremath{\cdot}}{.} \textbf{feed-forward loop} \texorpdfstring{\ensuremath{\cdot}}{.} \textbf{mutual
information} \texorpdfstring{\ensuremath{\cdot}}{.} \textbf{equifinality}

\begin{center}\rule{0.5\linewidth}{0.5pt}\end{center}

\subsection{Discussion Questions}\label{discussion-questions}

\begin{enumerate}
\def\labelenumi{\arabic{enumi}.}
\item
  Name a biological process at each of three hierarchical levels
  (molecular, cellular, organismal). For each, identify what would be
  lost by studying primarily the level below it.
\item
  Describe a biological example of hysteresis --- a system where the
  current state depends on history, not just the present inputs. How
  might hysteresis contribute to cell identity or memory?
\item
  The bacterium \emph{E. coli} regulates the \emph{lac} operon with a
  bistable switch: cells either fully express or fully repress the
  operon, with very few in intermediate states. Why might bistability be
  advantageous compared to a graded (linear) response? Under what
  physiological contexts might a graded response be preferred?
\item
  A patient takes a drug that blocks a key kinase in an oncogenic
  signalling pathway. Initial tumour regression is observed, but within
  months the tumour re-grows. Propose at least two systems-level
  mechanisms (feedback compensation, alternative pathway activation, or
  selection) that could explain this relapse. How would systems science
  inform rational combination therapy?
\item
  Turing patterns --- periodic spatial self-organisation arising from
  reaction--diffusion kinetics --- have been proposed to explain animal
  coat markings, digit spacing, and hair follicle arrangement. What does
  the existence of Turing patterns in biology suggest about the
  relationship between physics, chemistry, and life?
\item
  Sketch the glucose--insulin--glucagon feedback graph. Annotate where
  metformin, GLP-1 receptor agonists, and SGLT2 inhibitors act. Why
  might adding a second drug class outperform doubling the dose of one
  class?
\item
  Using the Hill equation with \(n = 4\) and \(K_d = 8\,\mu\text{M}\),
  estimate \texorpdfstring{\ensuremath{\theta}}{theta} at \([L] = 4\,\mu\text{M}\) and at
  \([L] = 16\,\mu\text{M}\). How does cooperativity sharpen the
  transition compared with \(n = 1\)?
\item
  Explain how a delay \texorpdfstring{\ensuremath{\tau}}{tau} in a negative-feedback loop with gain \(k\) can
  convert a stable equilibrium into a sustained oscillation. Use the
  conditions \(k\tau \ll 1\), \(k\tau \approx 1/2\), and
  \(k\tau > \pi/2\) to predict three qualitatively different behaviours
  and give a biological example of each.
\item
  Protein--protein interaction networks are scale-free with a power-law
  exponent near 2.5. Why does this make cells robust to random gene loss
  yet vulnerable to viral effectors that target hub proteins (e.g.,
  adenovirus E1B targeting p53)? How would you design an antiviral that
  exploits this asymmetry?
\item
  The cardiac pacemaker is extraordinarily robust to \texorpdfstring{\ensuremath{\pm}}{+/-}20\% perturbation
  of any single ionic conductance. Why might this robustness be a
  general property of well-designed limit-cycle oscillators, and what
  does it predict about the genetic redundancy you would expect to find
  in clock genes?
\end{enumerate}

\begin{center}\rule{0.5\linewidth}{0.5pt}\end{center}

\subsection{Review Questions}\label{review-questions}

\begin{enumerate}
\def\labelenumi{\arabic{enumi}.}
\item
  State the three elements that define any system (components,
  relationships, boundary) and classify isolated, closed, and open
  systems by their matter and energy exchange. Why are most living
  systems necessarily open?
\item
  Define equifinality and multifinality and give a biological example of
  each. Explain why a surgically reduced sea-urchin embryo producing a
  normal small larva illustrates equifinality, whereas a damaged watch
  does not behave this way.
\item
  Contrast negative and positive feedback by sign and biological role.
  For one positive-feedback process (action potential, childbirth, or
  the clotting cascade), identify the explicit bounding mechanism and
  explain why unbounded positive feedback would be lethal.
\item
  Using the delayed-feedback equation \(dx/dt = -k\,x(t-\tau)\), predict
  the qualitative behaviour for \(k\tau \ll 1\), \(k\tau \approx 1/2\),
  and \(k\tau > \pi/2\). Explain why biological clocks (circadian,
  cardiac, cell-cycle) live in the last regime and how shortening PER
  stability shortens the circadian period.
\item
  Using the Hill equation with \(K_d = 5\) nM and a dimeric (\(n=2\))
  transcription factor at \([L]=10\) nM, compute fractional occupancy,
  then recompute for \(n=1\). Explain quantitatively how cooperativity
  sharpens a dose--response curve into a switch.
\item
  Distinguish weak from strong emergence. Give one biological property
  that is plausibly weakly emergent and one that is contested as
  possibly strongly emergent, and state what evidence would upgrade the
  contested case from ``not yet derivable'' to ``in principle
  underivable.''
\item
  Compare saddle-node, Hopf, and pitchfork bifurcations by what happens
  to fixed points and which biological transition each underlies
  (catastrophic regime shift, oscillation onset, cell-fate commitment).
  For a stable spiral that gains a small limit cycle as a parameter
  crosses \(\mu_c\), name the bifurcation and a disorder consistent with
  it failing.
\item
  Explain why scale-free protein-interaction networks
  (\(\gamma \approx 2.5\)) are robust to random gene loss but fragile to
  targeted hub attack. Use this asymmetry to explain why TP53, MYC, and
  KRAS are potent cancer drivers and to propose the logic of a
  hub-exploiting antiviral.
\item
  Evaluate, using the mutual-information expression
  \(I = \tfrac12\log_2(1 + \sigma^2_{\text{signal}}/\sigma^2_{\text{noise}})\),
  why a single cytokine pathway conveying about 1 bit cannot
  reliably distinguish ``low/medium/high'' concentrations, and why
  evolution favours parallel signalling channels over ever-finer
  measurement of one channel.
\item
  Synthesis: the chapter argues systems biology is ``not fancy
  reductionism.'' Defend or challenge this claim by identifying at least
  three system-level metrics (e.g.~gain, delay, basin depth, Lyapunov
  exponent, mutual information) that acquire meaning primarily at the
  assembled-system level, and give a concrete question that no amount of
  better single-component data could answer.
\end{enumerate}

\begin{center}\rule{0.5\linewidth}{0.5pt}\end{center}

\subsection{Further Reading and Source
Notes}\label{further-reading-and-source-notes}

\begin{itemize}
\tightlist
\item
  Bertalanffy, L. von (1968). \emph{General System Theory}
  \citep{bertalanffy1968}. Braziller.
\item
  Strogatz, S. H. (2018). \emph{Nonlinear Dynamics and Chaos} (2nd ed.)
  \citep{strogatz2018}. Westview Press.
\item
  Tyson, J. J., Chen, K. C., \& Novak, B. (2003). Sniffers, buzzers,
  toggles and blinkers: dynamics of regulatory and signalling pathways
  in the cell \citep{tyson2003}. \emph{Current Opinion in Cell Biology},
  15(2), 221--231.
\item
  Alon, U. (2019). \emph{An Introduction to Systems Biology: Design
  Principles of Biological Circuits} (2nd ed.) \citep{alon2019}. CRC
  Press.
\item
  Cannon, W. B. (1932). \emph{The Wisdom of the Body}
  \citep{cannon1932}. W. W. Norton.
\item
  Mitchell, P. (1961). Coupling of phosphorylation to electron and
  hydrogen transfer by a chemi-osmotic type of mechanism
  \citep{mitchell1961}. \emph{Nature}, 191, 144--148. \emph{(A worked
  example of a feedback motif powering metabolism.)}
\end{itemize}

\begin{center}\rule{0.5\linewidth}{0.5pt}\end{center}

\subsubsection{Companion Source Module}\label{companion-source-module}

\textbf{Systems Science and the Logic of Emergence} should leave a
reproducible trail from a biological claim to the code, figure, diagram,
or paper-based activity that can test it. Use the surfaces below to
inspect the chapter's assumptions, rerun the relevant model, or compare
the manuscript explanation with companion labs and figures.

{\def\LTcaptype{none} % do not increment counter
\begin{longtable}[]{@{}
  >{\raggedright\arraybackslash}p{(\linewidth - 2\tabcolsep) * \real{0.5000}}
  >{\raggedright\arraybackslash}p{(\linewidth - 2\tabcolsep) * \real{0.5000}}@{}}
\toprule\noalign{}
\begin{minipage}[b]{\linewidth}\raggedright
Surface
\end{minipage} & \begin{minipage}[b]{\linewidth}\raggedright
Use it for
\end{minipage} \\
\midrule\noalign{}
\endhead
\bottomrule\noalign{}
\endlastfoot
\breaktt{src/biology/cell/cell\_biology.py} (\texttt{hill\_equation},
\breaktt{receptor\_occupancy}, \breaktt{signal\_amplification}) & Turn
feedback, thresholds, and signalling gain into inspectable
calculations. \\
\breaktt{src/biology/ecology/ecology.py} (\texttt{logistic\_growth}) &
Compare linear intuition with bounded growth and carrying-capacity
dynamics. \\
\breaktt{src/biology/biochemistry/biochemistry.py}
(\breaktt{reaction\_free\_energy}) & Connect system directionality to
thermodynamic constraints. \\
\breaktt{src/visualization/plots.py} (\breaktt{plot\_logistic\_growth})
and \breaktt{src/mermaid/biology\_diagrams.py}
(\breaktt{population\_growth\_stages\_diagram}) & Check whether graphical
summaries preserve the same model assumptions. \\
\end{longtable}
}

\textbf{Reproducibility check:} change one parameter at a time, record
the sign of the response, and explain whether the result reflects
feedback, saturation, or an arbitrary boundary choice.
\textbf{Cross-reference:} pair this with
\cref{sec:unit_0_complex_adaptive_systems},
\cref{sec:unit_III_bioenergetics_and_respiration}, and
\cref{sec:unit_X_population_ecology}.

\newpage

\section{Complex Adaptive Systems}\label{complex-adaptive-systems}

\label{sec:unit_0_complex_adaptive_systems}

\begin{quote}
\textbf{Unit 0} \texorpdfstring{\ensuremath{\cdot}}{.} Level 2/3 \texorpdfstring{\ensuremath{\cdot}}{.} 35 min read \texorpdfstring{\ensuremath{\cdot}}{.} 50 min lecture \texorpdfstring{\ensuremath{\cdot}}{.}
Prerequisites: \cref{sec:unit_0_systems_science}
\end{quote}

\begin{center}\rule{0.5\linewidth}{0.5pt}\end{center}

\subsection{Learning Objectives}\label{learning-objectives-1}

By the end of this chapter, students will be able to:

\begin{enumerate}
\def\labelenumi{\arabic{enumi}.}
\tightlist
\item
  Define a complex adaptive system and identify its distinguishing
  characteristics.
\item
  Describe how local rules among agents give rise to global order
  without a central controller.
\item
  Explain the roles of selection, variation, and adaptation in
  biological CAS.
\item
  Identify phase transitions and attractors in biological systems.
\item
  Recognise CAS dynamics at cellular, organismal, population, and
  ecosystem levels.
\item
  Analyse fitness landscapes using the NK model and predict the
  consequences of ruggedness.
\item
  Apply the logistic map to demonstrate how a single nonlinear equation
  can generate fixed points, oscillations, and chaos.
\item
  Connect biological CAS to engineered analogues such as genetic
  algorithms, neural networks, and ant-colony optimisation.
\end{enumerate}

\subsubsection{Study Blueprint}\label{study-blueprint-1}

\begin{itemize}
\tightlist
\item
  \textbf{Big idea:} Local rules can generate population-level patterns
  without a central controller.
\item
  \textbf{Core concepts:} nonlinearity, self-organization, attractors,
  power laws.
\item
  \textbf{Framework alignment:} Vision \& Change: Systems, Structure and
  function; AP Biology: Systems Interactions; NGSS-style topics:
  Structure and Function, Interdependent Relationships in Ecosystems.
\item
  \textbf{Model or quantitative lens:} Agent-rule and scaling arguments
  for emergent biological patterns.
\item
  \textbf{Data skill:} Distinguish deterministic trends from stochastic
  variation in repeated simulations.
\item
  \textbf{Practice cadence:} Questions and Methods, Representing and
  Describing Data, Argumentation.
\item
  \textbf{Common misconception to repair:} Emergence is not mysterious;
  it is a reproducible consequence of interactions plus constraints.
\item
  \textbf{Primary lab:} \cref{sec:lab_unit_0_complex_adaptive_systems}.
\item
  \textbf{Question bank:} \cref{sec:q_unit_0_complex_adaptive_systems}.
\item
  \textbf{Transfer task:} Compare flocking, immune activation, and
  microbial biofilms as adaptive systems.
\item
  \textbf{Bridge to computation:}
  \breaktt{biology.ecology.ecology.logistic\_growth}.
\end{itemize}

\begin{center}\rule{0.5\linewidth}{0.5pt}\end{center}

\subsection{Opening Vignette: From Sandpiles to the Santa Fe
Institute}\label{opening-vignette-from-sandpiles-to-the-santa-fe-institute}

In 1987, physicist Per Bak placed grains of sand one-by-one onto a pile
and discovered a startling regularity: the pile spontaneously organised
itself to a ``critical'' state where avalanches across a wide size range
obeyed a power law --- many small, few large, with no characteristic
scale. He called this \textbf{self-organised criticality}
\citep{bak1987} and argued it was a signature of complex adaptive
systems throughout nature, from earthquakes to extinction cascades to
neural firing patterns (\emph{Physical Review Letters}, 1987).

That same year, the Santa Fe Institute was founded in New Mexico with an
audacious charter: to study complex systems across disciplinary
boundaries. Economists (W. Brian Arthur), biologists (Stuart Kauffman
\citep{kauffman1993}), computer scientists (John Holland
\citep{holland1992}), and physicists (Murray Gell-Mann) worked
shoulder-to-shoulder, discovering that the same mathematics ---
agent-based models, fitness landscapes, power-law distributions ---
described stock markets, immune systems, and ecosystems alike.

The central biological insight: the immune system does not have a
commander. The brain does not have a central processor. An ant colony
does not have a foreman. In each case, simple agents following local
rules generate intelligent collective behaviour --- \textbf{emergence}
without design. This chapter explores how that principle operates from
the molecular to the planetary scale, equipping you with tools (fitness
landscapes, attractor analysis, power-law diagnostics, agent-based
simulation, bifurcation theory) you will use throughout the rest of the
textbook.

\begin{center}\rule{0.5\linewidth}{0.5pt}\end{center}

\subsection{What Is a Complex Adaptive
System?}\label{what-is-a-complex-adaptive-system}

A \textbf{complex adaptive system (CAS)} is a network of many
interacting agents that:

\begin{enumerate}
\def\labelenumi{\arabic{enumi}.}
\tightlist
\item
  \textbf{Adapt} --- agents alter their behaviour based on experience or
  selection.
\item
  \textbf{Self-organise} --- ordered collective behaviour emerges from
  local interactions without central control.
\item
  \textbf{Co-evolve} --- the environment of each agent is partly
  constituted by the other agents; most adapt together.
\item
  \textbf{Exhibit nonlinear dynamics} --- small perturbations can
  cascade into large effects (or large inputs may have minimal impact).
\end{enumerate}

The central insight of CAS theory is that \textbf{complexity is
endogenous}: it is generated by the interaction rules themselves, not
imposed from outside.

CAS theory unifies vocabulary that previously belonged to separate
disciplines. The ``selection pressure'' of an evolutionary biologist,
the ``energy landscape'' of a chemist, the ``objective function'' of a
computer scientist, and the ``potential well'' of a physicist are
different names for the same mathematical object: a real-valued function
over agent state that rules drive the system to extremise.

\begin{center}\rule{0.5\linewidth}{0.5pt}\end{center}

\subsection{Agents and Agent Rules}\label{agents-and-agent-rules}

In a CAS, the fundamental unit is the \textbf{agent} --- an entity with
internal state, sensors, and action capabilities. Agents follow local
rules based on local information. No agent knows the global state.

\subsubsection{Boids and Flocking}\label{boids-and-flocking}

Computer scientist Craig \citet{reynolds1987} demonstrated that three
simple rules produce realistic flocking behaviour in simulated birds
(``boids''):

\begin{enumerate}
\def\labelenumi{\arabic{enumi}.}
\tightlist
\item
  \textbf{Cohesion} --- steer towards the average position of nearby
  boids.
\item
  \textbf{Alignment} --- steer toward the average heading of nearby
  boids.
\item
  \textbf{Separation} --- avoid crowding; steer away from nearby boids.
\end{enumerate}

No boid knows the global shape of the flock. The flock emerges from
purely local interaction.

This principle recurs throughout biology:

\begin{itemize}
\tightlist
\item
  Quorum sensing in bacteria produces
  \hyperref[gl:biofilm]{\textbf{biofilm}} formation.
\item
  Cytokine gradients direct immune cell recruitment.
\item
  \hyperref[gl:action-potential]{\textbf{Action potential}} propagation
  in cardiac muscle coordinates the heartbeat.
\item
  Schooling fish escape predators by aligning with neighbours faster
  than the predator can re-target.
\end{itemize}

\subsubsection{Ant-Colony Optimisation and
Stigmergy}\label{ant-colony-optimisation-and-stigmergy}

Ant foraging is a paradigmatic CAS. An ant deposits a pheromone trail as
it returns from a food source; subsequent ants probabilistically follow
stronger trails and reinforce them. Shorter trails are reinforced faster
(more round-trips per unit time), so the colony progressively converges
on near-optimal paths. This \textbf{stigmergic} mechanism ---
communication through environmental modification --- has been
re-engineered into ant-colony-optimisation algorithms used to route
packets in computer networks and schedule jobs in factories. A single
ant is nearly random; the colony is a problem-solver.

\subsubsection{Immune Surveillance}\label{immune-surveillance}

Each lymphocyte carries a single antigen receptor specificity. The
immune system as a whole performs continuous surveillance not because
any cell is ``smart'' but because billions of differently-specific cells
circulate, sample, and clonally expand on contact with cognate antigen.
Negative selection in the thymus removes self-reactive clones; affinity
maturation in the germinal centre runs micro-evolution within the lymph
node. The system has no commander; collective decision emerges from
local rules.

\begin{figure}[htbp]
\centering
\includegraphics[width=0.82\linewidth,height=4.2in,keepaspectratio]{../figures/mermaid_inline/inline_mermaid_0006_3065e0fb0d3a.png}
\caption{Local agent interactions and emergent collective behaviour.}
\end{figure}

\begin{quote}
\textbf{Concept Check 1:} A swarm-intelligence start-up wants to control
a fleet of drones using boid rules but discovers that during certain
wind conditions the flock fragments. Hypothesise which of the three
rules (cohesion, alignment, separation) most likely needs reweighting
and explain how a change in radius of ``nearby'' might restore
coherence.
\end{quote}

\begin{center}\rule{0.5\linewidth}{0.5pt}\end{center}

\subsection{Self-Organisation in
Biology}\label{self-organisation-in-biology}

\subsubsection{Slime Moulds: A Class
Example}\label{slime-moulds-a-class-example}

\emph{Dictyostelium discoideum} (slime mould) spends most of its life as
independent amoebae. When starved, individual cells begin secreting
cyclic AMP (cAMP), triggering a \textbf{collective aggregation} into a
multicellular slug and eventually a fruiting body. No single cell plans
the fruiting body; the entire developmental programme is orchestrated by
local cAMP gradients --- a spectacular example of biological
self-organisation.

\subsubsection{Termite Mounds}\label{termite-mounds}

Termites build elaborate air-conditioned mounds --- up to 9 metres tall
--- following primarily local construction rules (place soil pellets
where you detect a pheromone trail). The architecture emerges from
millions of local decisions, not from a blueprint read by any individual
termite.

\subsubsection{Embryonic Development}\label{embryonic-development}

The body plan of a vertebrate emerges from a single fertilised cell
through cell division, migration, and differentiation guided by local
chemical gradients (morphogens). The French flag model of positional
information --- cells adopt different fates depending on the local
concentration of a morphogen --- captures the CAS logic of development.

\subsubsection{Turing Patterns and
Reaction--Diffusion}\label{turing-patterns-and-reactiondiffusion}

Alan Turing's 1952 paper on morphogenesis showed that two interacting
chemicals --- a local activator and a long-range inhibitor --- can
spontaneously break spatial symmetry to produce stripes, spots, or
hexagonal patterns. The minimal model is

\begin{equation}
\frac{\partial u}{\partial t} = D_u \nabla^{2} u + f(u, v), \qquad \frac{\partial v}{\partial t} = D_v \nabla^{2} v + g(u, v)
\label{eq:unit_0_turing}
\end{equation}

with \(D_v \gg D_u\) (the inhibitor diffuses faster than the activator).
Empirical examples include zebrafish stripe patterning, mouse
hair-follicle spacing, and digit number in tetrapod limbs. Turing's
logic is the cleanest example of order-from-instability: the homogeneous
state is unstable to spatial perturbations precisely because the
inhibitor outruns the activator.

\begin{center}\rule{0.5\linewidth}{0.5pt}\end{center}

\subsection{Attractors and Phase
Space}\label{attractors-and-phase-space}

The concept of an \textbf{attractor} is central to understanding CAS
dynamics. An attractor is a set of states toward which a system tends to
evolve over time.

{\def\LTcaptype{none} % do not increment counter
\begin{longtable}[]{@{}
  >{\raggedright\arraybackslash}p{(\linewidth - 4\tabcolsep) * \real{0.3256}}
  >{\raggedright\arraybackslash}p{(\linewidth - 4\tabcolsep) * \real{0.2558}}
  >{\raggedright\arraybackslash}p{(\linewidth - 4\tabcolsep) * \real{0.4186}}@{}}
\toprule\noalign{}
\begin{minipage}[b]{\linewidth}\raggedright
Attractor type
\end{minipage} & \begin{minipage}[b]{\linewidth}\raggedright
Description
\end{minipage} & \begin{minipage}[b]{\linewidth}\raggedright
Biological example
\end{minipage} \\
\midrule\noalign{}
\endhead
\bottomrule\noalign{}
\endlastfoot
Fixed point & System converges to a single stable state & Biochemical
equilibrium; homeostatic set point \\
Limit cycle & System oscillates periodically & Circadian rhythm;
\hyperref[gl:cell-cycle]{\textbf{cell cycle}}; cardiac pacemaker \\
Torus & System oscillates quasi-periodically & Coupled biological
clocks \\
Strange attractor & Deterministic chaos; sensitive initial conditions &
Some ECG arrhythmias; neuronal population dynamics \\
\end{longtable}
}

\textbf{Phase space} is the abstract space in which most possible system
states are points. Attractors are embedded in phase space; trajectories
flow toward them. Understanding biological regulation means
understanding the attractor landscape of the relevant phase space.

\subsubsection{Basins of Attraction and Landscape
Geometry}\label{basins-of-attraction-and-landscape-geometry}

The \textbf{basin of attraction} of an attractor \(A\) is the set of
states from which trajectories flow to \(A\). Basins are separated by
\textbf{separatrices} --- codimension-one manifolds across which a small
perturbation flips the system between attractors. In Waddington's
classical metaphor for development, the embryo rolls down a
``landscape'' whose valleys (basins) correspond to differentiated cell
types and whose ridges (separatrices) correspond to bistable cell-fate
decisions. Modern single-cell RNA-seq data can be used to reconstruct an
empirical analogue of this landscape from observed cell-state densities.

The \textbf{depth} of a basin (the energy barrier to escape) determines
its noise robustness; the \textbf{width} determines the range of initial
conditions that converge to it. A deep, narrow basin is robust but
inflexible; a shallow, wide basin is flexible but easily perturbed.
Cancer cells often inhabit basins that are \emph{wider} than normal
(higher phenotypic plasticity) and \emph{shallower} (easier to perturb
stochastically), explaining their resistance to single-target therapy.

\begin{figure}[htbp]
\centering
\includegraphics[width=0.82\linewidth,height=4.2in,keepaspectratio]{../figures/mermaid_inline/inline_mermaid_0007_48da9fb8cdee.png}
\caption{Conceptual phase space with disjoint attractor basins.}
\end{figure}

\begin{quote}
\textbf{Concept Check 2:} Both cancer cells and healthy cells inhabit a
\hyperref[gl:gene]{\textbf{gene}}-expression phase space with multiple
attractors (proliferative, quiescent, senescent, apoptotic). Many
targeted therapies try to push malignant cells out of the proliferative
basin. Why does this often fail --- what does the landscape analogy
predict about recurrence when the basin barrier is lower than the
therapy's ``kick'' amplitude? Propose a CAS-aware rationale for drug
\emph{combinations} rather than single agents.
\end{quote}

\begin{center}\rule{0.5\linewidth}{0.5pt}\end{center}

\subsection{Phase Transitions and
Criticality}\label{phase-transitions-and-criticality}

A \textbf{phase transition} is a qualitative change in system behaviour
at a critical parameter value. In biology:

\begin{itemize}
\tightlist
\item
  \textbf{\hyperref[gl:protein]{Protein} folding} --- a polypeptide
  transitions from disordered to folded at a critical temperature.
\item
  \textbf{The lac \hyperref[gl:operon]{\textbf{operon}}} --- bistable
  switch with two stable attractors (ON and OFF) separated by an
  unstable threshold.
\item
  \textbf{The immune system} --- a subclinical infection (low pathogen
  load) may persist or be cleared; above a threshold, an acute
  inflammatory response is triggered.
\item
  \textbf{Ecosystems} --- beyond a critical level of nutrient loading, a
  clear lake can collapse into a eutrophic state --- an ecological phase
  transition.
\item
  \textbf{Cytoplasm condensation} --- at high macromolecular crowding
  the cytoplasm transitions from liquid-like to glass-like, drastically
  slowing diffusion. Stress granules and P-bodies are physiological
  liquid-liquid phase separations.
\end{itemize}

\subsubsection{The Ising Picture of Cell-Fate
Switching}\label{the-ising-picture-of-cell-fate-switching}

A toy mathematical bridge between physics and biology is the Ising
model, in which spins \(s_i \in \{+1, -1\}\) on a lattice interact
through

\begin{equation}
H = -J \sum_{\langle i, j \rangle} s_i s_j - h \sum_{i} s_i
\label{eq:unit_0_ising}
\end{equation}

with coupling \(J\) and external field \(h\). Below a critical
temperature \(T_c\) the system spontaneously magnetises (most spins
agree); above it the system is disordered. Mapping spin-up to ``gene
ON'' and spin-down to ``gene OFF'' recasts cell-fate commitment as a
magnetisation phase transition: at low effective ``temperature'' (low
gene-expression noise), tightly-coupled networks lock into one of a few
discrete cell types; at high temperature, expression patterns are
disordered. The same mathematics describes binary opinion dynamics in
social networks, illustrating how comprehensive these critical phenomena
are.

\subsubsection{Power Laws and Scale-Free
Behaviour}\label{power-laws-and-scale-free-behaviour}

Many CAS, especially those poised near a critical phase transition
(\textbf{critical systems}), display \textbf{power-law} distributions:

\begin{equation}
P(x) \propto x^{-\alpha}
\label{eq:unit_0_powerlaw}
\end{equation}

where \texorpdfstring{\ensuremath{\alpha}}{alpha} is the scaling exponent. Power-law distributions in biology
arise in:

\begin{itemize}
\tightlist
\item
  Protein interaction network degree distributions (\texorpdfstring{\ensuremath{\alpha}}{alpha} \texorpdfstring{\ensuremath{\approx}}{~} 2--3).
\item
  Species abundance distributions (Preston's lognormal-like, with
  power-law tail).
\item
  Neuronal avalanche sizes \citep{beggs2003} (\texorpdfstring{\ensuremath{\alpha}}{alpha} \texorpdfstring{\ensuremath{\approx}}{~} 1.5).
\item
  Forest fire size distributions and earthquake magnitudes
  (Gutenberg--Richter law).
\item
  Sizes of mass extinctions in the fossil record.
\item
  Citations of academic papers and population sizes of cities.
\end{itemize}

The hypothesis of \textbf{self-organised criticality} (Bak, Tang \&
Wiesenfeld, 1987 \citep{bak1987}) proposes that many living systems
automatically tune themselves to criticality, because critical systems
are maximally responsive and information-transmitting. Critical neural
networks have the longest correlation lengths, the highest dynamic range
of stimulus encoding, and the maximum mutual information between input
and output --- most desirable for a brain.

\subsubsection{Worked Example: Scaling in Power-Law
Distributions}\label{worked-example-scaling-in-power-law-distributions}

\textbf{Problem:} Neuronal avalanche sizes (the number of
\hyperref[gl:neuron]{\textbf{neuron}}s firing synchronously in a
cascade) in cortical slices have been empirically shown to follow a
power law probability distribution \(P(x) = C x^{-\alpha}\), where the
scaling exponent \(\alpha \approx 1.5\). If a small avalanche involving
\(x = 10\) neurons occurs with a probability of \(0.05\) (5\%), what is
the probability of a massive avalanche involving \(x = 100\) neurons?

\textbf{Solution:}

\textbf{Set up the probability ratio.} Because \(C\) is a constant, the
ratio of probabilities for any two event sizes depends primarily on the
ratio of those sizes to the negative power of \texorpdfstring{\ensuremath{\alpha}}{alpha} :

\[ \frac{P(x_2)}{P(x_1)} = \left(\frac{x_2}{x_1}\right)^{-\alpha}  \label{eq:unit_0_complex_adaptive_systems_item_1}\]

\textbf{Substitute the known values} (\(x_1 = 10\), \(x_2 = 100\),
\(\alpha = 1.5\)):

\[ \frac{P(100)}{P(10)} = \left(\frac{100}{10}\right)^{-1.5}  \label{eq:unit_0_complex_adaptive_systems_item_2}\]

\[ \frac{P(100)}{P(10)} = (10)^{-1.5}  \label{eq:unit_0_complex_adaptive_systems_item_3}\]

\[ (10)^{-1.5} = \frac{1}{\sqrt{10^3}} = \frac{1}{\sqrt{1000}} \approx \frac{1}{31.62} \approx 0.0316  \label{eq:unit_0_complex_adaptive_systems_item_4}\]

\textbf{Calculate the final probability.}

\[ P(100) = P(10) \times 0.0316 = 0.05 \times 0.0316 \approx 0.00158  \label{eq:unit_0_complex_adaptive_systems_item_5}\]

A neuronal avalanche involving 100 neurons occurs with a probability of
approximately \textbf{0.16\%}. Because power-law distributions possess a
``fat tail'' relative to Gaussian distributions, extreme outlier events
(massive avalanches) remain statistically probable, illustrating how a
CAS poised near criticality can periodically undergo massive,
system-wide state changes.

\begin{quote}
\textbf{Connection (clinical) --- epileptic seizures as critical
events.} Cortical recordings near seizure onset show increasing mutual
information between regions and rising statistical signatures of
approach to a critical bifurcation. This is consistent with the brain
operating just below a critical line, with seizures representing
excursions across it. Closed-loop responsive neurostimulation devices
(e.g.~NeuroPace) detect these pre-critical signatures and deliver
counter-stimulation to keep the system in the subcritical basin.
\end{quote}

\begin{center}\rule{0.5\linewidth}{0.5pt}\end{center}

\subsection{Adaptation and Evolution as CAS
Processes}\label{adaptation-and-evolution-as-cas-processes}

Evolution by \hyperref[gl:natural-selection]{\textbf{natural selection}}
is a CAS at the population level:

\begin{itemize}
\tightlist
\item
  \textbf{Agents}: individual organisms.
\item
  \textbf{Local rules}: reproduction, survival,
  \hyperref[gl:mutation]{\textbf{mutation}}, migration.
\item
  \textbf{Environment}: partly constituted by other organisms.
\item
  \textbf{Emergent outcome}: population-level adaptation,
  \hyperref[gl:speciation]{\textbf{speciation}}, co-evolution.
\end{itemize}

The evolutionary process is not guided toward any particular solution.
It is open-ended and path-dependent. Fitness landscapes capture the
attractor structure of adaptive evolution.

\subsubsection{Fitness Landscapes and Adaptive
Walks}\label{fitness-landscapes-and-adaptive-walks}

A \textbf{fitness landscape} \citep{kauffman1993} is a surface in which
each point represents a possible genotype and the height represents the
fitness that genotype confers. Selection drives populations ``uphill''
toward local peaks. The metaphor was first made concrete by Sewall
Wright (1932) and later formalised by Stuart Kauffman's \textbf{NK
models}, which vary the number of loci \(N\) and the epistatic coupling
\(K\) between them:

\begin{itemize}
\tightlist
\item
  \(K = 0\) (no epistasis) --- the landscape is a single smooth peak;
  evolution is a deterministic climb.
\item
  \(K = N - 1\) (maximal epistasis) --- the landscape is maximally
  rugged, with \(\sim 2^N / (N + 1)\) local peaks. Populations get
  trapped.
\item
  Intermediate \(K\) --- there are dozens to thousands of local peaks
  with valleys of mildly reduced fitness between them.
\end{itemize}

Mathematically, in an NK model with binary loci the fitness of a
genotype is

\begin{equation}
W(\mathbf{x}) = \frac{1}{N} \sum_{i=1}^{N} w_{i}\!\left(x_{i}, x_{i_{1}}, \ldots, x_{i_{K}}\right)
\label{eq:unit_0_nk}
\end{equation}

where each locus contribution \(w_i\) depends on its own state \(x_i\)
and on \(K\) partner loci. For \(K = 0\) each locus is independent, so
the global optimum can be found one locus at a time. For \(K > 0\) the
contributions become entangled, generating the rugged landscape.
Empirical fitness landscapes (for HIV-1 protease, the \emph{E. coli}
ribosomal protein S6, and \emph{Aspergillus} enzymes) consistently land
at intermediate \(K\) --- rugged enough to trap evolution on local
peaks, smooth enough that occasional valleys can be crossed.

Ruggedness predicts why evolution can take several different paths to
similar fitness and why identical starting populations (as in the Lenski
\emph{E. coli} long-term evolution experiment) reach \emph{different}
adaptive solutions. In cancer, the same principle explains why
single-agent targeted therapy traps tumours on a local peak; combination
therapy shifts the landscape so the old peak is no longer fit, forcing
movement into valleys where the tumour is more vulnerable.

\begin{quote}
\textbf{Connection (clinical) --- antimicrobial cycling.} Hospital
antibiotic stewardship sometimes deliberately \emph{cycles} between drug
classes, exploiting the fact that resistance mutations to one class may
carry fitness costs on another. Each cycle reshapes the fitness
landscape; resistance peaks that were locally fit under cefazolin become
valleys under ciprofloxacin, evicting the population. The same logic is
being explored for cancer (``evolutionary therapy'' with adaptive
dosing).
\end{quote}

\subsubsection{Genetic Algorithms: Evolution in a
Computer}\label{genetic-algorithms-evolution-in-a-computer}

John Holland's \textbf{genetic algorithm} \citep{holland1992} formalises
Darwinian evolution as an optimisation procedure. A population of
candidate solutions (each encoded as a ``chromosome'') is evaluated, the
fittest are selected for reproduction, and offspring are produced by
crossover (recombining parental chromosomes) and mutation. Genetic
algorithms have designed:

\begin{itemize}
\tightlist
\item
  Antenna geometries for NASA satellites that outperform human-designed
  alternatives.
\item
  Aerodynamic profiles for racing yachts.
\item
  Schedule optimisers for hospitals.
\item
  Trading strategies that survive market regime change.
\end{itemize}

The \emph{biological} lesson is the inverse: just as artificial
evolution solves design problems by sampling, recombining, and
selecting, biological evolution discovers the metabolic and
morphological inventions of life by exactly the same process --- with
the difference that biological ``fitness functions'' are themselves
changing because the agents alter their environment.

\subsubsection{Quorum sensing: a molecular vote without a
mayor}\label{quorum-sensing-a-molecular-vote-without-a-mayor}

Many bacteria secrete \textbf{autoinducer} molecules whose extracellular
concentration rises with local cell density. Below a threshold, the
signal is noise; above it, receptors synchronise gene expression ---
biofilm formation, bioluminescence (\emph{Vibrio}), or
\hyperref[gl:virulence]{\textbf{virulence}} factor deployment
(\emph{Staphylococcus}). No cell counts its neighbours; each cell
integrates a \textbf{local concentration} that \emph{implicitly} encodes
population size. Mathematically, positive feedback on autoinducer
synthesis once a threshold is crossed is a \textbf{bistable switch} ---
a CAS motif duplicated in synthetic gene circuits. Antibiotics that fail
to penetrate biofilms often select for \textbf{persister}
subpopulations, illustrating how multicellular bacterial collectives
occupy a different attractor than planktonic growth.

\subsubsection{Trophic cascades: top-down reshaping of food
webs}\label{trophic-cascades-top-down-reshaping-of-food-webs}

Removing an apex predator (wolves from Yellowstone pre-1995
\citep{ripple2012}; large sharks from coral reefs) can trigger
\textbf{cascading} changes in herbivores, plants, and even geomorphology
--- not because the predator ``controlled'' each species, but because
the \textbf{network of interactions} rewires. Reintroduction experiments
show that restoration is not the reverse movie of removal: history and
initial conditions matter (\textbf{hysteresis}), another CAS signature.
Levin's classic argument \citep{levin1998} that ecosystems and the
biosphere are themselves complex adaptive systems rests on exactly this
combination of nonlinear dynamics, history dependence, and multi-scale
interaction.

\subsubsection{The Immune System as CAS}\label{the-immune-system-as-cas}

The adaptive immune system is perhaps the most striking interorganismal
CAS:

\begin{itemize}
\tightlist
\item
  B-cell clonal expansion generates antigen-specific population growth.
\item
  Affinity maturation (somatic hypermutation + selection) is
  micro-evolution within the germinal centre.
\item
  Memory is a stable population-level attractor that persists after
  antigen clearance.
\end{itemize}

\begin{quote}
\textbf{Concept Check 3:} A bacterial population is grown for 1000
generations under cycling exposure to drugs A and B. Resistance
mutations to A typically reduce growth rate on B and vice versa. Sketch
how the population-mean fitness on each drug evolves over time, and
explain in CAS vocabulary why the drug-cycling regime keeps the
population on a ``treadmill'' rather than letting it climb a single
peak.
\end{quote}

\begin{center}\rule{0.5\linewidth}{0.5pt}\end{center}

\subsection{Noise, Stochasticity, and
Bifurcation}\label{noise-stochasticity-and-bifurcation}

Deterministic CAS models capture the \emph{average} trajectory but miss
two crucial features of real biology: \textbf{noise} (stochastic
fluctuations in component numbers) and \textbf{bifurcation} (qualitative
changes in system behaviour at parameter boundaries).

\subsubsection{Noise in Genetic
Switches}\label{noise-in-genetic-switches}

A bistable switch such as the \emph{lac} operon ideally sits in either a
fully-repressed or fully-induced state. In individual cells, however,
\hyperref[gl:transcription]{\textbf{transcription}} factor copy numbers
are low (tens to hundreds per cell), so thermal fluctuations produce
\textbf{stochastic transitions} between the two attractors. The rate of
spontaneous switching follows a Kramers-like relation:

\begin{equation}
k_{\mathrm{switch}} \approx \omega_{0} \, e^{-\Delta U / k_B T_{\mathrm{eff}}}
\label{eq:unit_0_kramers}
\end{equation}

where \(\Delta U\) is the height of the barrier between attractors and
\(T_{\mathrm{eff}}\) is an effective noise temperature set by molecule
copy number. This has direct clinical consequences: bacterial
\textbf{persister cells} escape antibiotics by stochastically entering a
dormant attractor; cancer cells sample a phenotypic landscape and
survive chemotherapy in rare sub-states --- a Darwinian process on a CAS
landscape rather than on genetic variation alone.

\subsubsection{Fold Bifurcations and Hysteresis: Lake
Eutrophication}\label{fold-bifurcations-and-hysteresis-lake-eutrophication}

\hyperref[gl:eutrophication]{\textbf{Eutrophication}} is the textbook
ecological example of a \textbf{fold bifurcation}. Let \(P\) =
phosphorus load and \(A\) = algal biomass. For low \(P\) the system has
a single clear-water equilibrium; as \(P\) exceeds a critical value
\(P_1\) the clear state loses stability and the lake flips to a turbid,
algae-dominated state. Crucially, \emph{reducing} \(P\) below \(P_1\)
does \textbf{not} restore the clear state --- recovery requires that
\(P\) falls below a lower threshold \(P_2 < P_1\). This
\textbf{hysteresis loop} is why eutrophication management is so
expensive: undoing the transition requires substantially more reduction
than preventing it in the first place. The same mathematics describes
coral-reef to algal-reef transitions and some epileptic-focus formation
in the brain.

\subsubsection{The Logistic Map: A Worked Bifurcation
Diagram}\label{the-logistic-map-a-worked-bifurcation-diagram}

The logistic map

\begin{equation}
x_{n+1} = r\,x_{n}\,(1 - x_{n})
\label{eq:unit_0_logistic_map}
\end{equation}

with \(r \in [0, 4]\) and \(x \in [0, 1]\) is one of the simplest
nonlinear systems and exhibits a complete cascade from order to chaos as
\(r\) increases:

{\def\LTcaptype{none} % do not increment counter
\begin{longtable}[]{@{}
  >{\raggedright\arraybackslash}p{(\linewidth - 2\tabcolsep) * \real{0.3871}}
  >{\raggedright\arraybackslash}p{(\linewidth - 2\tabcolsep) * \real{0.6129}}@{}}
\toprule\noalign{}
\begin{minipage}[b]{\linewidth}\raggedright
Range of \(r\)
\end{minipage} & \begin{minipage}[b]{\linewidth}\raggedright
Long-term behaviour
\end{minipage} \\
\midrule\noalign{}
\endhead
\bottomrule\noalign{}
\endlastfoot
\(0 \le r < 1\) & Population dies out: \(x_n \to 0\) \\
\(1 \le r < 3\) & Single stable fixed point \(x^{*} = 1 - 1/r\) \\
\(3 \le r < 3.449...\) & Period-2 limit cycle \\
\(3.449... \le r < 3.544...\) & Period-4 limit cycle \\
Period-doubling cascade & Periods 8, 16, 32, \ldots{} accumulating at
\(r_{\infty} \approx 3.5699\) \\
\(r \ge r_{\infty}\) & Chaos with periodic windows (notably period-3
near \(r = 3.83\)) \\
\(r = 4\) & Fully developed chaos on \([0, 1]\) \\
\end{longtable}
}

The successive period-doubling thresholds approach \(r_{\infty}\) at a
comprehensive ratio (the \textbf{Feigenbaum constant}
\(\delta \approx 4.6692\)) that recurs in any smooth unimodal
one-dimensional map and even in measurements of dripping faucets and
convection rolls. This \textbf{universality} is part of why the logistic
map matters: it predicts that systems as different as insect populations
and turbulent fluids will display the same quantitative route to chaos.

\textbf{Worked numerical iteration (\(r = 3.2\), \(x_0 = 0.5\)):}
\(x_1 = 3.2 \cdot 0.5 \cdot 0.5 = 0.800\).
\(x_2 = 3.2 \cdot 0.8 \cdot 0.2 = 0.512\).
\(x_3 = 3.2 \cdot 0.512 \cdot 0.488 = 0.799\).
\(x_4 = 3.2 \cdot 0.799 \cdot 0.201 = 0.514\). The iterates settle into
a period-2 oscillation between roughly 0.513 and 0.800 --- a
discrete-time biological oscillator from one line of code.

\begin{quote}
\textbf{Connection (ecology) --- boom-bust population cycles.} Some
insect populations (notably Australian budworm) show empirical
period-doubling and chaotic dynamics that are well-modelled by
logistic-like equations. Whether populations are ``really chaotic'' or
merely noisy depends on detailed model fitting; the pedagogical point is
that \emph{no exotic mechanism} is required to generate complex temporal
patterns --- a single nonlinear feedback suffices.
\end{quote}

\begin{quote}
\textbf{Concept Check 4:} Iterate the logistic map by hand for
\(r = 2.5\), \(x_0 = 0.1\), four steps. Predict the long-term value,
then iterate a few more steps to confirm. What does this tell you about
the basin of attraction of the fixed point?
\end{quote}

\begin{center}\rule{0.5\linewidth}{0.5pt}\end{center}

\subsection{Robustness and Fragility}\label{robustness-and-fragility}

CAS often exhibit \textbf{robustness} --- the ability to maintain
function in the face of perturbation --- paired with occasional
\textbf{catastrophic fragility} in specific directions.

Biological robustness is achieved through:

\begin{itemize}
\tightlist
\item
  \textbf{Redundancy} --- multiple components performing similar
  functions.
\item
  \textbf{Degeneracy} --- structurally distinct components can perform
  the same function in some contexts.
\item
  \textbf{Modularity} --- failure in one module is localised.
\item
  \textbf{Negative feedback} --- deviations are automatically opposed.
\item
  \textbf{Distributed encoding} --- important information is represented
  across many components, not localised in one.
\end{itemize}

However, the same architectural features that confer robustness to
typical perturbations make the system fragile to novel, coordinated, or
cross-module perturbations. Antibiotic resistance evolution,
autoimmunity, and cancer most reflect such fragility. This trade-off is
sometimes called the \textbf{robust-yet-fragile} signature of CAS:
highly optimised systems are \emph{more} robust to expected
perturbations and \emph{more} fragile to unexpected ones. Carlson and
Doyle named this ``highly optimised tolerance'' (HOT), and it explains
why mature ecosystems suffer catastrophically when novel invasive
species or climate regimes appear.

\begin{quote}
\textbf{Concept Check 5:} Consider a CAS with strong negative feedback,
modular organisation, and five-fold redundant sensors. Name one category
of perturbation it will handle easily and one category that will
nonetheless catastrophically break it. Why do the same architectural
features that create robustness also create specific fragilities?
\end{quote}

\begin{center}\rule{0.5\linewidth}{0.5pt}\end{center}

\subsection{CAS Across Biological
Scales}\label{cas-across-biological-scales}

{\def\LTcaptype{none} % do not increment counter
\begin{longtable}[]{@{}
  >{\raggedright\arraybackslash}p{(\linewidth - 6\tabcolsep) * \real{0.1190}}
  >{\raggedright\arraybackslash}p{(\linewidth - 6\tabcolsep) * \real{0.0714}}
  >{\raggedright\arraybackslash}p{(\linewidth - 6\tabcolsep) * \real{0.1429}}
  >{\raggedright\arraybackslash}p{(\linewidth - 6\tabcolsep) * \real{0.6667}}@{}}
\toprule\noalign{}
\begin{minipage}[b]{\linewidth}\raggedright
Scale
\end{minipage} & \begin{minipage}[b]{\linewidth}\raggedright
CAS
\end{minipage} & \begin{minipage}[b]{\linewidth}\raggedright
Agents
\end{minipage} & \begin{minipage}[b]{\linewidth}\raggedright
Emergent collective property
\end{minipage} \\
\midrule\noalign{}
\endhead
\bottomrule\noalign{}
\endlastfoot
Molecular & Gene regulatory network & Transcription factors, RNAs & Gene
expression pattern, cell fate \\
Cellular & Intracellular signalling & Signal proteins, second messengers
& Robust most-or-none responses \\
Tissue & Neural network & Neurons & Cognition, memory, perception \\
Organism & Immune system & B and T cells, innate cells & Adaptive
immunity \\
Population & Evolving population & Individual organisms & Adaptation,
speciation \\
Ecosystem & Community ecology & Species populations & Food web
structure, \hyperref[gl:succession]{\textbf{succession}} \\
Biosphere & Gaia / planetary metabolism & Biomes, ocean basins &
Atmospheric composition, climate \\
\end{longtable}
}

\begin{figure}[htbp]
\centering
\includegraphics[width=0.82\linewidth,height=4.2in,keepaspectratio]{../figures/mermaid_inline/inline_mermaid_0008_dc46e45ae847.png}
\caption{Cross-scale interactions: upward emergence and downward causation.}
\end{figure}

\subsubsection{A Gallery of Agent-Based
Models}\label{a-gallery-of-agent-based-models}

{\def\LTcaptype{none} % do not increment counter
\begin{longtable}[]{@{}
  >{\raggedright\arraybackslash}p{(\linewidth - 6\tabcolsep) * \real{0.2273}}
  >{\raggedright\arraybackslash}p{(\linewidth - 6\tabcolsep) * \real{0.1364}}
  >{\raggedright\arraybackslash}p{(\linewidth - 6\tabcolsep) * \real{0.2500}}
  >{\raggedright\arraybackslash}p{(\linewidth - 6\tabcolsep) * \real{0.3864}}@{}}
\toprule\noalign{}
\begin{minipage}[b]{\linewidth}\raggedright
Phenomenon
\end{minipage} & \begin{minipage}[b]{\linewidth}\raggedright
Agents
\end{minipage} & \begin{minipage}[b]{\linewidth}\raggedright
Local rules
\end{minipage} & \begin{minipage}[b]{\linewidth}\raggedright
Emergent property
\end{minipage} \\
\midrule\noalign{}
\endhead
\bottomrule\noalign{}
\endlastfoot
Bird flocks & Boids & Cohesion, alignment, separation & Coordinated
flight \\
Ant trails & Ants + pheromone & Deposit, follow, evaporate & Shortest
path \\
Slime mould & \emph{Dictyostelium} & cAMP secretion, chemotaxis &
Multicellular fruiting body \\
Schelling segregation & Households & Move if too few neighbours of
``type'' & Macroscopic segregation \\
Game of Life & Cells on grid & Birth/death rules from neighbourhood &
Gliders, oscillators, computers \\
SIR epidemic & S, I, R classes & Contact + recovery probability &
Outbreak peak, herd-immunity threshold \\
T-cell repertoire & Lymphocytes & Random receptor + clonal selection &
Adaptive specificity \\
\end{longtable}
}

The pattern is comprehensive: the rules are absurdly simple, the global
behaviour is rich. This is why agent-based modelling has become a
standard tool in epidemiology (COVID-19 contact networks), behavioural
ecology, and synthetic biology.

\begin{figure}[htbp]
\centering
\includegraphics[width=0.82\linewidth,height=4.2in,keepaspectratio]{../figures/mermaid_inline/inline_mermaid_0009_0666f051c2ee.png}
\caption{Agent-based modelling loop linking local rules to measured emergent behaviour.}
\end{figure}

\begin{center}\rule{0.5\linewidth}{0.5pt}\end{center}

\subsection{Synthesis and Practical
Implications}\label{synthesis-and-practical-implications}

CAS is not merely a metaphor; it is a \emph{toolkit} with specific
predictions:

\begin{itemize}
\tightlist
\item
  If you suspect a system sits at criticality, look for power-law
  statistics in size, duration, or wait-time distributions.
\item
  If you suspect bistability, look for hysteresis: probe the system in
  both forward and reverse directions and watch for trajectory mismatch.
\item
  If you suspect a rugged fitness landscape, look for path-dependent
  endpoints in replicate evolution experiments.
\item
  If you suspect agent-based emergence, ask which global property would
  \emph{not} survive removing the local rules --- that is the emergent
  component.
\item
  If you suspect chaos, estimate the leading Lyapunov exponent from
  time-series data and check for sensitive dependence on initial
  conditions.
\end{itemize}

These diagnostics turn CAS theory from rhetorical to actionable.

\begin{center}\rule{0.5\linewidth}{0.5pt}\end{center}

\subsection{Current Evidence and Frontier
Biology}\label{current-evidence-and-frontier-biology-1}

For \textbf{Complex Adaptive Systems}, frontier biology belongs inside
the evidence logic of the chapter. Systems models are useful when they
expose assumptions, uncertainty, and failure modes rather than merely
producing elegant diagrams. The core reading question is this: agent
rules, heterogeneity, stochasticity, and path dependence make biological
prediction conditional rather than absolute.

\begin{itemize}
\tightlist
\item
  \textbf{What to verify:} identify the observation, model, assay, or
  dataset that would make the claim stronger or weaker.
\item
  \textbf{What to qualify:} state the scale, organism, cell type,
  environmental condition, or population where the claim is expected to
  hold.
\item
  \textbf{What to compare:} test at least one alternative explanation,
  baseline, or null model before treating the pattern as causal.
\item
  \textbf{What to cite:} distinguish primary evidence, review synthesis,
  public dataset, and institutional guidance; for recent or numeric
  claims, prefer the source closest to the measurement and state what
  has changed since it was published.
\end{itemize}

Treat every model as a claim about mechanism: define the system
boundary, identify the observable that would falsify the model, and
report the uncertainty that would change a decision.

\textbf{Source practice:} Use model-validation sources when available,
and state which observation would falsify the model rather than treating
a diagram as proof.

\subsection{Summary}\label{summary-1}

\begin{itemize}
\tightlist
\item
  A complex adaptive system is a collection of interacting agents with
  adaptive rules that generate emergent collective behaviour.
\item
  Self-organisation arises from local interactions without central
  control; Turing patterns and Dictyostelium aggregation are
  paradigmatic examples.
\item
  Attractors define the landscape of stable system states; phase
  transitions separate distinct regimes.
\item
  Power laws and scale-free properties often signal systems poised near
  criticality; the Bak sandpile, neuronal avalanches, and earthquakes
  share statistical signatures.
\item
  Fitness landscapes (Wright's metaphor; Kauffman's NK formalism)
  capture the structure of adaptive evolution; ruggedness predicts
  replicate divergence and resistance evolution.
\item
  Genetic algorithms transplant biological evolution into computational
  design problems, illustrating that the
  search-by-variation-and-selection logic is substrate-independent.
\item
  The logistic map demonstrates that a single nonlinear feedback can
  produce fixed points, period doubling, and chaos at the Feigenbaum
  point.
\item
  Evolution, the immune system, neural networks, and ecosystems are most
  biological CAS --- and so is the engineered swarm of drones, the
  genetic algorithm, and the social network.
\end{itemize}

\begin{center}\rule{0.5\linewidth}{0.5pt}\end{center}

\subsection{Key Terms}\label{key-terms-1}

\textbf{complex adaptive system} \texorpdfstring{\ensuremath{\cdot}}{.} \textbf{agent} \texorpdfstring{\ensuremath{\cdot}}{.} \textbf{stigmergy} \texorpdfstring{\ensuremath{\cdot}}{.}
\textbf{self-organisation} \texorpdfstring{\ensuremath{\cdot}}{.} \textbf{attractor} \texorpdfstring{\ensuremath{\cdot}}{.} \textbf{basin of
attraction} \texorpdfstring{\ensuremath{\cdot}}{.} \textbf{phase space} \texorpdfstring{\ensuremath{\cdot}}{.} \textbf{phase transition} \texorpdfstring{\ensuremath{\cdot}}{.}
\textbf{critical system} \texorpdfstring{\ensuremath{\cdot}}{.} \textbf{power law} \texorpdfstring{\ensuremath{\cdot}}{.} \textbf{scale-free} \texorpdfstring{\ensuremath{\cdot}}{.}
\textbf{fitness landscape} \texorpdfstring{\ensuremath{\cdot}}{.} \textbf{NK model} \texorpdfstring{\ensuremath{\cdot}}{.} \textbf{genetic
algorithm} \texorpdfstring{\ensuremath{\cdot}}{.} \textbf{logistic map} \texorpdfstring{\ensuremath{\cdot}}{.} \textbf{Feigenbaum constant} \texorpdfstring{\ensuremath{\cdot}}{.}
\textbf{robustness} \texorpdfstring{\ensuremath{\cdot}}{.} \textbf{degeneracy} \texorpdfstring{\ensuremath{\cdot}}{.} \textbf{limit cycle} \texorpdfstring{\ensuremath{\cdot}}{.}
\textbf{strange attractor} \texorpdfstring{\ensuremath{\cdot}}{.} \textbf{Turing pattern} \texorpdfstring{\ensuremath{\cdot}}{.} \textbf{bistable
switch} \texorpdfstring{\ensuremath{\cdot}}{.} \textbf{hysteresis}

\begin{center}\rule{0.5\linewidth}{0.5pt}\end{center}

\subsection{Discussion Questions}\label{discussion-questions-1}

\begin{enumerate}
\def\labelenumi{\arabic{enumi}.}
\item
  The human gut \hyperref[gl:microbiome]{\textbf{microbiome}} contains
  about 10\texorpdfstring{\textsuperscript{1}}{1}\texorpdfstring{\textsuperscript{3}}{3} microbial cells from hundreds of species,
  interacting through metabolic cross-feeding, competition, and quorum
  sensing. Describe this as a CAS. What are the agents, local rules, and
  emergent properties? What could perturb the microbiome into a
  different attractor?
\item
  Explain why a cancer tumour is itself a CAS. What ``selection
  pressure'' drives cancer evolution? Why does targeting a single driver
  mutation often fail to permanently eliminate a tumour? How might CAS
  thinking guide combination therapies?
\item
  Neural networks in the brain and artificial neural networks in deep
  learning both exhibit emergent computation from simple local rules.
  Compare and contrast the biological and artificial cases in terms of
  (a) learning rules, (b) architecture, (c) energy consumption, and (d)
  generalisation vs.~overfitting.
\item
  Many ecosystems exhibit hysteresis: a lake may tolerate increasing
  nutrient loading up to a threshold, then abruptly collapse to a
  eutrophic state, which cannot be reversed simply by reducing nutrient
  input to pre-collapse levels. What does this asymmetry imply for
  conservation management? How does CAS theory explain the hysteresis?
\item
  The concept of ``self-organised criticality'' suggests that many
  biological networks operate near a critical point between order and
  disorder. What functional advantages might this confer on neural
  networks? What experimental evidence (e.g., neuronal avalanches
  \citep{beggs2003}) supports or challenges this hypothesis?
\item
  Explain quorum sensing in \emph{Pseudomonas aeruginosa} or
  \emph{Vibrio fischeri} using CAS vocabulary (agents, local rules,
  emergent behaviour). How might blocking autoinducer synthesis differ
  from traditional bactericidal antibiotics in terms of resistance
  evolution?
\item
  After wolf reintroduction in Yellowstone \citep{ripple2012}, willow
  recovery along streams was attributed partly to reduced elk browsing.
  What other biotic or \hyperref[gl:abiotic]{\textbf{abiotic}} variables
  could confound a simple ``predator \texorpdfstring{\ensuremath{\to}}{->} herbivore \texorpdfstring{\ensuremath{\to}}{->} plant'' chain? How
  would you design a study to separate correlation from causation?
\item
  The Lenski long-term evolution experiment in \emph{E. coli} has run
  for \textgreater75 000 generations across 12 replicate populations on
  a single fitness landscape. Roughly half the populations independently
  evolved hypermutator phenotypes, and one acquired the ability to use
  citrate aerobically --- but primarily after about 31 000
  generations of ``potentiating'' mutations. What does this pattern of
  divergent yet partially convergent evolution reveal about fitness
  landscape ruggedness, contingency, and the predictability of
  evolutionary outcomes?
\end{enumerate}

\begin{center}\rule{0.5\linewidth}{0.5pt}\end{center}

\subsection{Review Questions}\label{review-questions-1}

\begin{enumerate}
\def\labelenumi{\arabic{enumi}.}
\item
  List the four defining properties of a complex adaptive system
  (adaptation, self-organisation, co-evolution, nonlinear dynamics) and
  give one biological example that exhibits the complete set
  simultaneously.
\item
  Explain the chapter's claim that ``complexity is endogenous.'' Using
  the boids model, describe how three local rules (cohesion, alignment,
  separation) generate a global flock that no individual agent
  represents.
\item
  Distinguish a fixed-point attractor, a limit cycle, and a strange
  attractor, giving one biological example of each. For the limit cycle,
  name the feedback feature that sustains the oscillation.
\item
  Apply the basin-of-attraction landscape to cancer therapy. A targeted
  drug delivers a ``kick'' that displaces malignant cells from the
  proliferative basin. Explain, in terms of basin depth and barrier
  height, why recurrence is likely when the barrier is lower than the
  kick amplitude, and why combination therapy reshapes rather than
  merely perturbs the landscape.
\item
  Iterate the logistic map \(x_{n+1}=r x_n(1-x_n)\) by hand for
  \(r=3.2\), \(x_0=0.5\) for four steps, and classify the long-term
  behaviour using the chapter's \(r\)-regime table. Then explain what
  the Feigenbaum constant \(\delta\approx4.669\) implies about the
  \emph{universality} of the route to chaos across unrelated systems.
\item
  Analyse self-organised criticality: identify the statistical signature
  (power-law size distributions) and explain, using neuronal avalanches,
  why a system poised near criticality maximises dynamic range and
  mutual information. Why is this functionally advantageous for a brain?
\item
  Using Kauffman's NK model, explain how the epistatic coupling \(K\)
  controls landscape ruggedness. Predict the qualitative outcome of
  replicate evolution experiments (e.g.~the Lenski long-term \emph{E.
  coli} lines) at \(K=0\), intermediate \(K\), and \(K=N-1\), and
  connect ruggedness to the failure of single-agent targeted therapy.
\item
  Distinguish robustness, redundancy, and degeneracy. Then explain the
  ``robust-yet-fragile'' (highly-optimised-tolerance) signature: why do
  the same architectural features that confer robustness to expected
  perturbations create catastrophic fragility to novel, coordinated
  ones? Give one clinical or ecological example.
\item
  Evaluate the fold-bifurcation/hysteresis account of lake
  eutrophication. Why does reducing phosphorus below the original
  tipping point \(P_1\) fail to restore the clear state, and what does
  the lower threshold \(P_2\) imply for the cost asymmetry between
  prevention and remediation? Identify one other biological system
  governed by the same mathematics.
\item
  Synthesis: the chapter offers five diagnostic tests (power-law
  statistics, hysteresis probing, path-dependent replicate endpoints,
  the ``remove-the-local-rules'' emergence test, Lyapunov-exponent
  estimation). Choose a biological system of your choice, propose which
  two diagnostics you would run, state the falsifiable prediction each
  makes, and explain why a no-transition null model is essential before
  declaring an early-warning signal real.
\end{enumerate}

\begin{center}\rule{0.5\linewidth}{0.5pt}\end{center}

\subsection{Further Reading and Source
Notes}\label{further-reading-and-source-notes-1}

\begin{itemize}
\tightlist
\item
  Holland, J. H. (1992). \emph{Adaptation in Natural and Artificial
  Systems} (2nd ed.) \citep{holland1992}. MIT Press.
\item
  Kauffman, S. A. (1993). \emph{The Origins of Order: Self-Organisation
  and Selection in Evolution} \citep{kauffman1993}. Oxford University
  Press.
\item
  Bak, P., Tang, C., \& Wiesenfeld, K. (1987). Self-organized
  criticality \citep{bak1987}. \emph{Physical Review Letters}, 59(4),
  381--384.
\item
  Levin, S. A. (1998). Ecosystems and the biosphere as complex adaptive
  systems \citep{levin1998}. \emph{Ecosystems}, 1(5), 431--436.
\item
  Beggs, J. M., \& Plenz, D. (2003). Neuronal avalanches in neocortical
  circuits \citep{beggs2003}. \emph{Journal of Neuroscience}, 23(35),
  11167--11177.
\item
  Strogatz, S. H. (2018). \emph{Nonlinear Dynamics and Chaos} (2nd ed.)
  \citep{strogatz2018}. Westview Press.
\end{itemize}

\begin{center}\rule{0.5\linewidth}{0.5pt}\end{center}

\subsubsection{Companion Source Module}\label{companion-source-module-1}

\textbf{Complex Adaptive Systems} should leave a reproducible trail from
a biological claim to the code, figure, diagram, or paper-based activity
that can test it. Use the surfaces below to inspect the chapter's
assumptions, rerun the relevant model, or compare the manuscript
explanation with companion labs and figures.

{\def\LTcaptype{none} % do not increment counter
\begin{longtable}[]{@{}
  >{\raggedright\arraybackslash}p{(\linewidth - 2\tabcolsep) * \real{0.5000}}
  >{\raggedright\arraybackslash}p{(\linewidth - 2\tabcolsep) * \real{0.5000}}@{}}
\toprule\noalign{}
\begin{minipage}[b]{\linewidth}\raggedright
Surface
\end{minipage} & \begin{minipage}[b]{\linewidth}\raggedright
Use it for
\end{minipage} \\
\midrule\noalign{}
\endhead
\bottomrule\noalign{}
\endlastfoot
\breaktt{src/biology/ecology/ecology.py} (\texttt{lotka\_volterra},
\breaktt{biodiversity\_indices}, \texttt{logistic\_growth}) & Explore
feedback, interaction strength, and diversity metrics as emergent
summaries. \\
\breaktt{src/biology/evolution/evolution.py}
(\breaktt{simulate\_selection}, \breaktt{wright\_fisher\_drift},
\breaktt{fitness\_landscape\_1d}) & Compare selection, drift, and
landscape ruggedness as agent-level rules. \\
\breaktt{src/visualization/plots.py} (\breaktt{plot\_lotka\_volterra},
\breaktt{plot\_selection\_simulation}) & Inspect how oscillations and
allele-frequency trajectories depend on starting conditions. \\
\breaktt{src/mermaid/biology\_diagrams.py} (\texttt{food\_web\_diagram},
\breaktt{population\_growth\_stages\_diagram}) & Link network structure
to emergent population outcomes. \\
\end{longtable}
}

\textbf{Reproducibility check:} rerun or recalculate a scenario from two
initial states and ask whether convergence, hysteresis, or path
dependence is doing the explanatory work. \textbf{Cross-reference:} use
\cref{sec:unit_0_systems_science},
\cref{sec:unit_V_population_genetics}, and
\cref{sec:unit_X_community_ecology} as comparison cases.

\newpage

\section{Active Inference and the Free Energy
Principle}\label{active-inference-and-the-free-energy-principle}

\label{sec:unit_0_active_inference}

\begin{quote}
\textbf{Unit 0} \texorpdfstring{\ensuremath{\cdot}}{.} Level 3/3 \texorpdfstring{\ensuremath{\cdot}}{.} 45 min read \texorpdfstring{\ensuremath{\cdot}}{.} 75 min lecture \texorpdfstring{\ensuremath{\cdot}}{.}
Prerequisites: \cref{sec:unit_0_systems_science},
\cref{sec:unit_0_complex_adaptive_systems}
\end{quote}

\begin{center}\rule{0.5\linewidth}{0.5pt}\end{center}

\subsection{Learning Objectives}\label{learning-objectives-2}

By the end of this chapter, students will be able to:

\begin{enumerate}
\def\labelenumi{\arabic{enumi}.}
\tightlist
\item
  Explain what the free energy principle \citep{friston2010} states and
  why it is relevant to biology.
\item
  Describe the Bayesian brain hypothesis and how prediction-error
  minimisation drives perception.
\item
  Distinguish between perceptual inference (changing beliefs to match
  sensory input) and active inference \citep{friston2017} (acting to
  fulfil predictions).
\item
  Connect active inference to
  \hyperref[gl:homeostasis]{\textbf{homeostasis}}, allostasis
  \citep{sterling2015}, and interoception.
\item
  Define a Markov blanket and identify its four state-classes (internal,
  external, sensory, active) in a biological system.
\item
  Apply expected free energy to explain how an organism balances
  goal-directed reward and information-gathering.
\item
  Describe canonical predictive-processing microcircuits in the cortex
  and how they implement hierarchical inference.
\item
  Identify connections between active inference and other frameworks in
  the textbook (metabolism, neural signalling, evolution, immunology,
  computational psychiatry).
\end{enumerate}

\subsubsection{Study Blueprint}\label{study-blueprint-2}

\begin{itemize}
\tightlist
\item
  \textbf{Big idea:} Living systems reduce uncertainty by acting on the
  world as well as by updating internal models.
\item
  \textbf{Core concepts:} prediction error, Bayesian updating, policy
  selection, homeostasis.
\item
  \textbf{Framework alignment:} Vision \& Change: Systems, Structure and
  function; AP Biology: Systems Interactions; NGSS-style topics:
  Structure and Function, Interdependent Relationships in Ecosystems.
\item
  \textbf{Model or quantitative lens:} Bayesian belief updating and
  expected-free-energy-style policy comparison.
\item
  \textbf{Data skill:} Read a small probability table and update a
  prediction after new evidence.
\item
  \textbf{Practice cadence:} Questions and Methods, Representing and
  Describing Data, Argumentation.
\item
  \textbf{Common misconception to repair:} Active inference is not
  passive prediction; action changes the sensory data that arrive next.
\item
  \textbf{Primary lab:} \cref{sec:lab_unit_0_active_inference}.
\item
  \textbf{Question bank:} \cref{sec:q_unit_0_active_inference}.
\item
  \textbf{Transfer task:} Map prediction-error reasoning onto
  chemotaxis, thermoregulation, and attention.
\item
  \textbf{Bridge to computation:}
  \breaktt{biology.physiology.physiology.homeostasis\_response}.
\end{itemize}

\begin{center}\rule{0.5\linewidth}{0.5pt}\end{center}

\subsection{Opening Vignette: Helmholtz and the Unconscious
Inference}\label{opening-vignette-helmholtz-and-the-unconscious-inference}

In 1867, Hermann von Helmholtz proposed that perception is not passive
recording but \textbf{unconscious inference}: the brain interprets noisy
sensory signals by combining them with prior knowledge about the world.
When you perceive a flat photograph as three-dimensional, your brain is
making an inference --- filling in depth from shadow, texture, and
perspective cues that it has learned through a lifetime of experience.

For over a century, Helmholtz's insight remained a philosophical
curiosity. Then, in the early 2000s, neuroscientist Karl Friston
formulated the \textbf{free energy principle (FEP)} --- a mathematical
framework that makes Helmholtz's intuition precise. The FEP proposes
that most living systems, from single cells to entire organisms, can be
understood as systems that minimise \textbf{variational free energy}: a
tractable upper bound on the surprise (negative log-probability) of
their sensory observations \citep{friston2010}. The bacterium swimming
up a glucose gradient and the neuroscientist modelling the brain are,
under this framework, performing the same fundamental operation ---
inference under uncertainty.

The free energy principle does not merely explain perception. By
extending inference to \emph{action} --- organisms act to make their
predictions come true --- it provides a unified account of sensation,
movement, homeostasis, development, and even evolution
\citep{friston2017}. This chapter introduces this powerful framework and
shows how it connects to every major theme in this textbook.

\begin{center}\rule{0.5\linewidth}{0.5pt}\end{center}

\subsection{The Problem of Life}\label{the-problem-of-life}

Most living organisms face a fundamental challenge: they must maintain
their internal organisation in the face of a ceaselessly changing and
uncertain environment. To survive, an organism must:

\begin{itemize}
\tightlist
\item
  \textbf{Model} its environment --- build and update an internal
  representation of what is out there.
\item
  \textbf{Sense} --- gather information that reduces uncertainty about
  environmental states.
\item
  \textbf{Act} --- change the environment or its own internal state to
  remain within viable bounds.
\end{itemize}

This three-part loop --- \emph{model \texorpdfstring{\ensuremath{\to}}{->} sense \texorpdfstring{\ensuremath{\to}}{->} act} --- is the core of
what the \textbf{free energy principle (FEP)} formalises.

\begin{figure}[htbp]
\centering
\includegraphics[width=0.82\linewidth,height=4.2in,keepaspectratio]{../figures/mermaid_inline/inline_mermaid_0010_5abcffa1cee1.png}
\caption{Agent-environment loop under active inference.}
\end{figure}

\begin{center}\rule{0.5\linewidth}{0.5pt}\end{center}

\subsection{Markov Blankets: The Mathematical Skin of an
Agent}\label{markov-blankets-the-mathematical-skin-of-an-agent}

Before formalising free energy, we need a way to make rigorous the
boundary between an agent and its environment. The \textbf{Markov
blanket} provides exactly this.

\subsubsection{Definition}\label{definition}

For a set of random variables, the Markov blanket of a node is the
minimal set of nodes that, once conditioned on, render the node
statistically independent of every other variable in the network. For an
agent that persists in time, this generalises to a partition of states
into four disjoint classes:

\begin{itemize}
\tightlist
\item
  \textbf{External states} \texorpdfstring{\ensuremath{\eta}}{eta} --- the world outside the agent.
\item
  \textbf{Sensory states} \(s\) --- variables through which the world
  influences the agent (photoreceptor activations, mechanoreceptor
  strain, chemosensor occupancy).
\item
  \textbf{Active states} \(a\) --- variables through which the agent
  influences the world (motor neuron output, secretory activity, ciliary
  beating).
\item
  \textbf{Internal states} \texorpdfstring{\ensuremath{\mu}}{mu} --- everything inside the agent (gene
  expression, membrane potentials, synaptic weights).
\end{itemize}

The blanket itself is the union \(b = s \cup a\). The defining property
is that internal states are conditionally independent of external states
given the blanket:

\begin{equation}
P(\mu, \eta \mid s, a) = P(\mu \mid s, a)\,P(\eta \mid s, a)
\label{eq:unit_0_markov_blanket}
\end{equation}

The internal states ``see'' the world primarily through \(s\) and
``touch'' it primarily through \(a\).

\subsubsection{Why Markov Blankets
Matter}\label{why-markov-blankets-matter}

Markov blankets give a principled definition of an agent: anything that
maintains a recognisable boundary in state space \emph{has} a Markov
blanket, almost by definition. The same partition applies to a bacterium
(cell membrane), a neuron (axonal/dendritic membrane), an organism (skin
and sensorimotor apparatus), and a colony (collective foraging
interface). Because the partition is recursive --- the blanket of a
blanket is a blanket --- it naturally accommodates hierarchical
biological organisation. Here a \emph{generative model} is the
organism's internal probabilistic model that predicts its sensory inputs
from hypothesised hidden causes in the environment.

\begin{figure}[htbp]
\centering
\includegraphics[width=0.82\linewidth,height=4.2in,keepaspectratio]{../figures/mermaid_inline/inline_mermaid_0011_0a16fce1c099.png}
\caption{Markov-blanket partition: internal states couple to external states primarily through sensory and active states.}
\end{figure}

\begin{quote}
\textbf{Concept Check 1:} Identify sensory states, active states,
internal states, and external states for (a) an \emph{E. coli} cell
performing chemotaxis, (b) a single hepatocyte, and (c) a foraging
honeybee. Which states would be hardest to assign in each case, and why?
\end{quote}

\begin{center}\rule{0.5\linewidth}{0.5pt}\end{center}

\subsection{The Free Energy Principle}\label{the-free-energy-principle}

The free energy principle, developed principally by Karl Friston (UCL),
proposes that most living systems can be characterised as systems that
\textbf{minimise their free energy} --- a quantity that measures the
mismatch between an agent's internal model and the sensory data it
receives.

Formally, variational free energy \(F\) provides an upper bound on
\textbf{surprisal} (the log-probability of sensory observations under a
model):

\begin{equation}
F = \underbrace{D_{KL}\bigl[ Q(\mathbf{z}) \,\|\, P(\mathbf{z} \mid \mathbf{o}) \bigr]}_{\text{complexity (posterior mismatch)}} - \underbrace{\ln P(\mathbf{o})}_{\text{surprise (negative log-evidence)}}
\label{eq:unit_0_free_energy}
\end{equation}

where:

\begin{itemize}
\tightlist
\item
  \(Q(\mathbf{z})\) is the agent's approximate posterior belief over
  hidden causes \(\mathbf{z}\).
\item
  \(P(\mathbf{z}|\mathbf{o})\) is the true posterior.
\item
  \(D_{KL}\) is the Kullback--Leibler divergence.
\item
  \(P(\mathbf{o})\) is the probability of observations \(\mathbf{o}\)
  under the generative model.
\end{itemize}

Because the true posterior is generally intractable, agents minimise
\(F\) as a proxy for maximising their \textbf{evidence} (model fit to
sensory data).

Equivalently, free energy decomposes as:

\begin{equation}
F = \underbrace{\mathbb{E}_{Q}\bigl[ \ln Q(\mathbf{z}) - \ln P(\mathbf{z}) \bigr]}_{\text{complexity (prior mismatch)}} - \underbrace{\mathbb{E}_{Q}\bigl[ \ln P(\mathbf{o} \mid \mathbf{z}) \bigr]}_{\text{expected log-likelihood (accuracy)}}
\label{eq:unit_0_free_energy_decomp}
\end{equation}

The first term penalises beliefs that deviate from priors; the second
rewards beliefs that predict observations well. Minimising \(F\)
balances prior expectations against sensory evidence --- exactly the
logic of Bayesian inference.

\begin{quote}
\textbf{Concept Check 2:} An organism's ecological
\hyperref[gl:niche]{\textbf{niche}} can be described as the set of
states its generative model treats as probable. Why does a penguin
housed in a tropical aquarium experience persistent free energy, even
when fed and safe? Describe two distinct strategies the penguin might
use to minimise it and explain why behavioural change (return to cold
water) is usually faster than belief update.
\end{quote}

\begin{center}\rule{0.5\linewidth}{0.5pt}\end{center}

\subsection{The Bayesian Brain}\label{the-bayesian-brain}

The \textbf{Bayesian brain hypothesis} (Helmholtz 1867; Friston 2005;
Clark 2013) proposes that the brain is fundamentally a
\textbf{prediction machine}: it maintains a hierarchical generative
model of the world and continuously uses that model to predict incoming
sensory signals.

\subsubsection{Hierarchical Predictive
Coding}\label{hierarchical-predictive-coding}

In the predictive processing framework, the brain generates
\textbf{top-down predictions} from higher cortical areas toward lower
sensory areas. Lower areas compute the \textbf{prediction error} --- the
difference between what was predicted and what was received:

\begin{equation}
\varepsilon = \mathbf{o} - \hat{\mathbf{o}}
\label{eq:unit_0_prediction_error}
\end{equation}

Prediction errors propagate upward, driving updates to the generative
model. When predictions are accurate, prediction errors approach zero
and neural activity is minimised --- an energetically efficient coding
strategy.

\subsubsection{The Canonical
Microcircuit}\label{the-canonical-microcircuit}

Anatomical and physiological evidence suggests the cortex implements
predictive coding through a \textbf{canonical microcircuit} repeated
across cortical areas. The simplified picture:

\begin{itemize}
\tightlist
\item
  \textbf{Superficial pyramidal neurons} (layers II/III) send forward
  signals that empirically encode prediction errors, projecting to
  higher cortical areas.
\item
  \textbf{Deep pyramidal neurons} (layers V/VI) send backward signals
  that empirically encode predictions, projecting to lower areas.
\item
  \textbf{Inhibitory interneurons} modulate the \emph{gain} of error
  units, implementing the precision weighting described in the earlier
  predictive-processing discussion.
\end{itemize}

In the laminar reading of predictive coding, each cortical area is a
small inference engine: it receives ascending prediction errors from
below, compares them with descending predictions from above, and outputs
both an updated internal state (forward) and a refined prediction
(backward). The same architecture helps interpret visual cortical
hierarchies (V1\texorpdfstring{\ensuremath{\to}}{->}V2\texorpdfstring{\ensuremath{\to}}{->}V4\texorpdfstring{\ensuremath{\to}}{->}IT), auditory streams, and somatosensory pathways,
suggesting predictive coding is a broad candidate principle for cortical
computation rather than a settled one-size-fits-most algorithm.

\subsubsection{Precision and Attention}\label{precision-and-attention}

Each prediction error can be weighted by \textbf{precision} (inverse
variance): high precision means ``trust this channel.'' Mathematically,
a precision-weighted prediction error is

\begin{equation}
\xi = \pi \cdot (\mathbf{o} - \hat{\mathbf{o}})
\label{eq:unit_0_precision}
\end{equation}

where \(\pi = 1/\sigma^2\) is the precision. \textbf{Attention} --- from
spatial orienting to feature selection --- can be read as
\textbf{optimising precision} over sensory hierarchies: amplify
unexpected but reliable signals, attenuate predictable ones.
\textbf{Sensory attenuation} during self-generated movement (you tickle
yourself less than others tickle you) follows because the motor command
predicts somatic input, down-weighting proprioceptive precision so the
tickle does not register as surprising. Disorders that mis-estimate
precision (autism-spectrum hypotheses; schizophrenia-spectrum
\textbf{aberrant salience}) reinterpret noise as meaningful structure or
vice versa, producing false inference at the perceptual or interoceptive
level.

\subsubsection{Worked Example: Bayesian Updating and Prediction
Error}\label{worked-example-bayesian-updating-and-prediction-error}

\textbf{Problem:} A foraging bird maintains a prior belief about the
spatial location \(x\) of a food source, modeled as a Gaussian
distribution with mean \(\mu_{prior} = 10 \text{ m}\) and variance
\(\sigma^2_{prior} = 4\). The bird receives a noisy sensory observation
indicating food at \(o = 16 \text{ m}\), with a sensory variance
\(\sigma^2_{obs} = 12\). Calculate the prediction error and the new
updated (posterior) belief about the food's location.

\textbf{Solution:}

\begin{enumerate}
\def\labelenumi{\arabic{enumi}.}
\item
  \textbf{Calculate the prediction error:}

  The generative model's prediction is the mean of the prior
  (\(\hat{o} = 10\)).
  \[ \varepsilon = o - \hat{o} = 16 - 10 = 6 \text{ m}  \label{eq:unit_0_active_inference_item_1}\]
\item
  \textbf{Calculate the optimal updating weight (Kalman gain):}

  In Bayesian inference with Gaussians, the weight given to the sensory
  evidence depends on the relative precision (inverse variance) of the
  prior vs the observation.
  \[ K = \frac{\sigma^2_{prior}}{\sigma^2_{prior} + \sigma^2_{obs}} = \frac{4}{4 + 12} = \frac{4}{16} = 0.25  \label{eq:unit_0_active_inference_item_2}\]
\item
  \textbf{Calculate the posterior mean:}

  The new belief (\(\mu_{post}\)) is the old belief plus the prediction
  error scaled by the optimal weight:
  \[ \mu_{post} = \mu_{prior} + K \cdot \varepsilon  \label{eq:unit_0_active_inference_item_3}\]

  \[ \mu_{post} = 10 + 0.25 \cdot 6 = 10 + 1.5 = 11.5 \text{ m}  \label{eq:unit_0_active_inference_item_4}\]
\end{enumerate}

Due to the high uncertainty of the sensory observation compared to the
prior (\(\sigma^2_{obs} = 12\) vs \(\sigma^2_{prior} = 4\)), the bird
primarily slightly adjusts its internal model (moving beliefs from 10 to
11.5) despite a large prediction error of 6. This formalises how sensory
evidence shifts the generative model under the Free Energy Principle.

\begin{figure}[htbp]
\centering
\includegraphics[width=0.82\linewidth,height=4.2in,keepaspectratio]{../figures/mermaid_inline/inline_mermaid_0012_3c4d1d827422.png}
\caption{Predictive processing: top-down predictions, bottom-up errors, and action.}
\end{figure}

\begin{quote}
\textbf{Concept Check 3:} Repeat the worked example with
\(\sigma^2_{prior} = 12\) and \(\sigma^2_{obs} = 4\) (i.e.~a vague prior
and a precise observation). Compute the new Kalman gain and the
posterior mean for the same prior mean (10) and observation (16). What
does the comparison reveal about how priors and likelihoods compete?
\end{quote}

\begin{center}\rule{0.5\linewidth}{0.5pt}\end{center}

\subsection{Active Inference}\label{active-inference}

\textbf{Active inference} extends predictive coding from perception to
\textbf{action}. An organism can reduce prediction error in two ways:

\begin{enumerate}
\def\labelenumi{\arabic{enumi}.}
\tightlist
\item
  \textbf{Perceptual inference} --- update internal beliefs to match
  incoming sensory data (classic Bayesian inference).
\item
  \textbf{Active inference} --- act on the world to make sensory data
  match the predictions generated by the generative model.
\end{enumerate}

Under the FEP, behaviour is the fulfilment of prior expectations. An
organism with a prior that its body temperature will remain at 37 °C
will \emph{act} to make that prediction true when it is falsified by
cold exposure (seek warmth, shiver, vasoconstrict).

This reframes homeostasis: not as the passive correction of deviations,
but as the \textbf{active fulfilment of generative model predictions
about body state}.

\subsubsection{Homoeostasis
vs.~Allostasis}\label{homoeostasis-vs.-allostasis}

{\def\LTcaptype{none} % do not increment counter
\begin{longtable}[]{@{}
  >{\raggedright\arraybackslash}p{(\linewidth - 4\tabcolsep) * \real{0.1750}}
  >{\raggedright\arraybackslash}p{(\linewidth - 4\tabcolsep) * \real{0.2250}}
  >{\raggedright\arraybackslash}p{(\linewidth - 4\tabcolsep) * \real{0.6000}}@{}}
\toprule\noalign{}
\begin{minipage}[b]{\linewidth}\raggedright
Concept
\end{minipage} & \begin{minipage}[b]{\linewidth}\raggedright
Core idea
\end{minipage} & \begin{minipage}[b]{\linewidth}\raggedright
Active inference framing
\end{minipage} \\
\midrule\noalign{}
\endhead
\bottomrule\noalign{}
\endlastfoot
\textbf{Homeostasis} & Maintain a fixed set point & Precise prior over a
target internal state \\
\textbf{Allostasis} & Proactively adjust the set point in anticipation
of need & Updating prior expectations based on context \\
\end{longtable}
}

Allostasis --- popularised by \citet{sterling1988} --- is naturally
captured by active inference: the generative model changes its
predictions (set points) before the perturbation arrives, explaining why
the heart rate rises before exercise begins.

\begin{figure}[htbp]
\centering
\includegraphics[width=0.82\linewidth,height=4.2in,keepaspectratio]{../figures/mermaid_inline/inline_mermaid_0013_b335958206ee.png}
\caption{Reactive homeostasis versus predictive allostasis.}
\end{figure}

\subsubsection{Worked Example --- Free Energy Minimisation via
Action}\label{worked-example-free-energy-minimisation-via-action}

Consider an organism whose generative model encodes a prior on body
temperature \(T\):

\begin{equation}
P(T) = \mathcal{N}(37, 0.5^2) \qquad \text{(prior: body temperature should be 37 °C, tight precision)}
\label{eq:unit_0_temp_prior}
\end{equation}

Its interoceptive sensor is noisy,
\(P(o \mid T) = \mathcal{N}(T, 1.0^2)\). Suppose the sensor reports
\(o = 34\) °C (cold). The posterior belief about \(T\) is a Gaussian
with mean:

\begin{equation}
\mu_{\mathrm{post}} = \frac{\sigma^2_o \cdot 37 + \sigma^2_{\mathrm{prior}} \cdot 34}{\sigma^2_o + \sigma^2_{\mathrm{prior}}}
= \frac{1.0 \cdot 37 + 0.25 \cdot 34}{1.0 + 0.25}
= 36.4 \; \mathrm{°C}
\label{eq:unit_0_temp_posterior}
\end{equation}

The variational free energy associated with this observation is
approximately \(F_{\mathrm{before}} \approx 4.5\) nats (dominated by the
log-likelihood of \(o = 34\) under the prior). The agent has two ways to
reduce it:

\begin{enumerate}
\def\labelenumi{\arabic{enumi}.}
\tightlist
\item
  \textbf{Perceptual inference (passive)} --- accept
  \(\mu_{\mathrm{post}} = 36.4\) °C as its new belief. \(F\) drops to
  about 2.1 nats.
\item
  \textbf{Active inference} --- act to change the world so that future
  \(o\) matches the 37 °C prior. Shiver, vasoconstrict, seek warmth.
  After one minute of shivering the sensor reads 36.5 °C; after five
  minutes, 37.0 °C. \(F\) approaches 0.
\end{enumerate}

\textbf{Key insight:} acting reduces free energy \emph{more completely}
than updating beliefs, because the prior is tight --- the organism would
rather change the world than its expectations. This is the formal
expression of what Helmholtz intuited in 1867 and what the FEP
quantifies today: \emph{organisms do not merely represent the world;
they re-shape it to match their models.}

\subsubsection{Expected Free Energy and Policy
Selection}\label{expected-free-energy-and-policy-selection}

Selecting an action requires not only minimising current free energy but
also predicting the consequences of future actions. \textbf{Expected
free energy} \(G\) for a policy \texorpdfstring{\ensuremath{\pi}}{pi} (sequence of actions) is

\begin{equation}
G(\pi) = \underbrace{\mathbb{E}_{Q(o, z \mid \pi)}\!\bigl[\ln Q(z \mid \pi) - \ln Q(z \mid o, \pi)\bigr]}_{\text{epistemic value}} - \underbrace{\mathbb{E}_{Q(o \mid \pi)}\!\bigl[\ln P(o \mid C)\bigr]}_{\text{instrumental value}}
\label{eq:unit_0_expected_free_energy}
\end{equation}

where \(C\) encodes the agent's preferences (which observations are
valuable). The two terms decompose action selection into:

\begin{itemize}
\tightlist
\item
  \textbf{Epistemic (information) value} --- the expected reduction in
  uncertainty about latent causes; high for exploratory or ``curious''
  policies.
\item
  \textbf{Instrumental (pragmatic) value} --- the expected
  log-probability of preferred outcomes; high for goal-directed
  policies.
\end{itemize}

Policies are selected with probability \(Q(\pi) \propto e^{-G(\pi)}\):
the agent prefers policies with low expected free energy. Crucially,
when no information remains to be gained, the agent reduces to a
reward-maximiser; when no extrinsic reward differs across policies, the
agent reduces to a pure information-seeker (curiosity). This single
equation thereby unifies reinforcement learning, optimal experimental
design, and exploration--exploitation trade-offs that previously
occupied separate literatures.

\begin{quote}
\textbf{Concept Check 4:} A rat in a novel maze runs in apparently
random directions even when food is reliably available in one corner.
Decompose this behaviour into epistemic and instrumental value. What
would change if the rat had explored the maze for an hour first?
\end{quote}

\subsubsection{Epistemic vs.~Instrumental Value (worked numerical
example)}\label{epistemic-vs.-instrumental-value-worked-numerical-example}

Suppose an agent considers two policies, \(\pi_A\) and \(\pi_B\), with
epistemic values 1.0 and 3.0 nats and instrumental values 4.0 and 2.5
nats respectively. Then

\[G(\pi_A) = -(1.0 + 4.0) = -5.0, \qquad G(\pi_B) = -(3.0 + 2.5) = -5.5 \label{eq:unit_0_active_inference_item_5}\]

Both are roughly comparable, but \(\pi_B\) has slightly lower expected
free energy because its higher information gain compensates for somewhat
lower expected reward. Under softmax selection
\(Q(\pi_B)/Q(\pi_A) = e^{0.5} \approx 1.65\) --- the agent is about
1.65\texorpdfstring{\ensuremath{\times}}{x} more likely to choose the more informative policy. Reducing
temperature (sharpening the softmax) makes the choice nearly
deterministic; raising temperature flattens it. This single mechanism
can reproduce explore-exploit phenomena across species and contexts.

\begin{center}\rule{0.5\linewidth}{0.5pt}\end{center}

\subsection{Embodied Active Inference}\label{embodied-active-inference}

Active inference is inherently \textbf{embodied} and \textbf{enactive}
--- it requires a body that can act. The agent's generative model must
represent:

\begin{itemize}
\tightlist
\item
  \textbf{Exteroceptive states} --- sensory information about the
  external environment.
\item
  \textbf{Interoceptive states} --- sensory information about the
  internal body (viscera, blood chemistry, proprioception).
\item
  \textbf{Action affordances} --- possible actions the agent can take,
  and their predicted sensory consequences.
\end{itemize}

The body is not merely the carrier of the brain; the brain's generative
model is shaped by bodily form, evolutionary history, and developmental
experience. This \textbf{4E cognition} framework --- embodied, embedded,
enacted, extended --- is directly implied by active inference.

\begin{quote}
\textbf{Concept Check 5:} A patient with a pacemaker experiences
``phantom palpitations'' --- strong sensations of irregular heart rhythm
that the device rhythm itself rarely produces. Within the
active-inference framework, why might long-standing tight priors about
heart rhythm \emph{create} a prediction error when the rhythm is in fact
perfectly regular? What strategies --- pharmacological, cognitive, or
behavioural --- would widen those priors?
\end{quote}

\begin{center}\rule{0.5\linewidth}{0.5pt}\end{center}

\subsection{Active Inference and
Evolution}\label{active-inference-and-evolution}

The free energy principle provides a formal bridge between
individual-level adaptive behaviour and evolutionary dynamics:

\begin{itemize}
\tightlist
\item
  \hyperref[gl:natural-selection]{\textbf{Natural selection}} filters
  for organisms whose generative models minimise free energy (are
  accurate predictors of their ecological niche).
\item
  The \textbf{ecological niche} is itself the set of environmental
  states that the organism's generative model implicitly expects.
\item
  \textbf{Niche construction} --- organisms actively shape their
  environments to match their generative models --- is a direct
  prediction of active inference applied to evolutionary time scales.
\end{itemize}

The evolutionary emergence of increasingly sophisticated generative
models (from chemotaxis in bacteria to abstract planning in primates)
can be framed as the phylogenetic history of free energy minimisation.

\subsubsection{Free Energy Minimisation Beyond
Brains}\label{free-energy-minimisation-beyond-brains}

The FEP was developed in neuroscience but is not limited to it. Recent
work extends active inference to:

\begin{itemize}
\tightlist
\item
  \textbf{Single cells.} A bacterium running E. coli chemotaxis can be
  cast as an active-inference agent whose generative model encodes a
  prior over ``high glucose,'' whose sensors read receptor occupancy,
  and whose actions are tumble probability and run length.
  Berg-and-Purcell information bounds reappear as bounds on free energy
  minimisation.
\item
  \textbf{Immune system.} B-cell affinity maturation can be read as
  Bayesian inference over antigen ``causes,'' with somatic hypermutation
  as a proposal distribution and selection as a posterior update. Memory
  cells embody high-precision priors against re-infection.
\item
  \textbf{Development.} Morphogen gradients pattern tissues by reducing
  free energy of cell-fate ``predictions''; experimentally induced
  perturbations are absorbed when the developmental model can re-infer
  position from remaining cues.
\item
  \textbf{Evolution.} Natural selection itself can be framed as gradient
  descent on a long-timescale free energy functional, with the species'
  generative model encoded in its genome and the niche acting as the
  sensory environment. This is the most ambitious and the most contested
  of the FEP's extensions.
\end{itemize}

The unifying claim is that \textbf{any system that resists dispersion
--- that maintains its Markov blanket against thermodynamic noise ---
\emph{behaves as if} it minimises free energy}. The ``as-if'' is
important: the FEP is a description of what self-organising systems
necessarily do at large enough timescales, not a claim that any
particular molecule ``computes'' free energy.

\begin{center}\rule{0.5\linewidth}{0.5pt}\end{center}

\subsection{Computational Psychiatry: When Inference Goes
Wrong}\label{computational-psychiatry-when-inference-goes-wrong}

If the brain is a prediction machine, many psychiatric conditions admit
a generative-model interpretation. The framework is provisional but
increasingly testable.

\subsubsection{Anxiety and Interoceptive
Hyper-precision}\label{anxiety-and-interoceptive-hyper-precision}

If priors over ``safe body state'' are too tight (high precision),
ordinary interoceptive fluctuations (a normal heartbeat, gut
peristalsis) generate large prediction errors that the system
experiences as alarming. Therapeutic interventions that improve
interoceptive accuracy (mindfulness training, graded exposure,
beta-blocker pharmacology) can be re-interpreted as recalibrating prior
precisions and likelihood widths. Classical exposure therapy reads
naturally as deliberate prediction-error minimisation: the patient
repeatedly experiences a feared stimulus without consequence, and
Bayesian belief updating attenuates the conditioned association.

\subsubsection{Depression and the Dark-Room
Problem}\label{depression-and-the-dark-room-problem}

Depression has been modelled as an over-precise prior on negative
outcomes combined with under-weighted reward signals --- the agent's
generative model ``knows'' that nothing it does will help, so policies
with high epistemic or instrumental value are systematically discounted.
Predictive-processing accounts of anhedonia explain reduced exploratory
behaviour as a rational response to a prior of hopelessness. The
``dark-room problem'' --- why don't agents simply hide in a dark room to
minimise sensory surprise? --- is solved by noting that prolonged
darkness \emph{itself} generates surprise relative to the organism's
evolved generative model, which expects some level of stimulation.

\subsubsection{Schizophrenia and Aberrant
Salience}\label{schizophrenia-and-aberrant-salience}

Schizophrenia-spectrum hypotheses suggest a \emph{failure} of precision
regulation: prediction errors that ought to be down-weighted (because
they are noise) are instead amplified, while reliable signals are
attenuated. This produces ``aberrant salience,'' in which random noise
feels meaningful and structured signals feel hollow. Some accounts link
this directly to dopaminergic precision modulation in striatal circuits;
clozapine and other antipsychotics may, on this reading, restore
appropriate precision weighting. Hallucinations, on this view, are not
failures of perception but failures of inference --- top-down
predictions overruling weak sensory evidence.

\subsubsection{Autism-Spectrum Predictive
Differences}\label{autism-spectrum-predictive-differences}

A complementary account proposes that autism involves \emph{reduced}
precision on top-down priors (or \emph{increased} precision on bottom-up
sensory signals), so the world is experienced in raw sensory detail
without the smoothing influence of predictions. This may explain
hypersensitivity, attention to fine detail, and difficulty with rapidly
changing social cues --- most of which require high-prior precision to
integrate noisy data into coherent interpretations.

\begin{quote}
\textbf{Connection (clinical) --- limitations.} These accounts are
conceptually attractive but the empirical evidence for any specific
computational reading is still developing. They should be read as
\emph{hypotheses} that organise current research rather than diagnostic
frameworks. Patients deserve evidence-based treatments, and the FEP is a
complement to --- not a replacement for --- careful clinical psychiatry.
\end{quote}

\begin{quote}
\textbf{Concept Check 6:} A patient describes a constant feeling that
``something is about to go wrong'' despite no objective threat. Describe
how this could arise from (a) over-precise priors on negative outcomes,
(b) under-precise interoceptive likelihoods, or (c) hyperactive
dopaminergic precision. What experimental measurement might distinguish
the three?
\end{quote}

\begin{center}\rule{0.5\linewidth}{0.5pt}\end{center}

\subsection{Applications Beyond
Neuroscience}\label{applications-beyond-neuroscience}

The active-inference framework has been applied --- with varying degrees
of empirical support --- well beyond brain science.

\subsubsection{Immunology}\label{immunology}

The adaptive immune system can be framed as a multi-agent
active-inference system. Each B-cell embodies a hypothesis (its receptor
specificity); the population of B-cells maintains a distribution over
possible antigens; somatic hypermutation generates new hypotheses;
clonal selection performs Bayesian update; and memory cells are
high-precision priors that minimise future free energy of re-infection.
The thymic deletion of self-reactive T-cells corresponds to clipping the
prior to exclude hypotheses that would generate persistent self-reactive
prediction errors.

\subsubsection{Developmental Biology}\label{developmental-biology}

Embryonic development can be modelled as a hierarchical inference: each
cell receives morphogen ``observations,'' combines them with priors
encoded in its genome and current state, and performs an active update
(differentiation, migration, division). Robustness to perturbation is
then explained as \emph{Bayesian inference over a redundant generative
model}: removing one cue still leaves enough information to infer
position.

\subsubsection{Evolutionary Active
Inference}\label{evolutionary-active-inference}

If natural selection filters for free-energy-minimising organisms, the
long-time evolutionary dynamics can themselves be cast as inference over
a ``species-level'' generative model encoded in the genome and
implemented through development. Niche construction (beavers building
dams, plants oxygenating the atmosphere) becomes the evolutionary
analogue of action: organisms reshape the world to match their inherited
expectations.

\subsubsection{Ecology and Collective
Behaviour}\label{ecology-and-collective-behaviour}

Schooling fish, foraging ants, and coordinated bacteria can be analysed
as collective active-inference systems whose Markov blankets are larger
than any individual organism. The mathematical formalism predicts when
collective behaviour will outperform individual behaviour: when
individual sensory data are noisy and the environment is partially
observable, a collective Bayesian posterior is sharper than any
individual's.

\begin{center}\rule{0.5\linewidth}{0.5pt}\end{center}

\subsection{Connections to Other
Units}\label{connections-to-other-units}

{\def\LTcaptype{none} % do not increment counter
\begin{longtable}[]{@{}
  >{\raggedright\arraybackslash}p{(\linewidth - 2\tabcolsep) * \real{0.1333}}
  >{\raggedright\arraybackslash}p{(\linewidth - 2\tabcolsep) * \real{0.8667}}@{}}
\toprule\noalign{}
\begin{minipage}[b]{\linewidth}\raggedright
Unit
\end{minipage} & \begin{minipage}[b]{\linewidth}\raggedright
Active inference connection
\end{minipage} \\
\midrule\noalign{}
\endhead
\bottomrule\noalign{}
\endlastfoot
\textbf{Unit III --- Metabolism} & ATP production and allostasis
maintain the energetic constraints required for inference \\
\textbf{Unit VIII --- Botany} & Stomatal regulation and tropisms can be
modelled as active inference by plant systems \\
\textbf{Unit IX --- Neuroscience} & Predictive coding instantiates
active inference at the neural circuit level \\
\textbf{Unit X --- Ecology} & Niche construction and habitat selection
as ecosystem-level active inference \\
\end{longtable}
}

\subsubsection{Interoceptive active inference and psychiatry
(conceptual)}\label{interoceptive-active-inference-and-psychiatry-conceptual}

\textbf{Interoception} --- the sense of the visceral body --- supplies
observations about heart rate, gut distension, inflammation, and
temperature. If priors over ``safe body state'' are too tight, ordinary
fluctuations may generate persistent prediction-error signals
experienced as \textbf{anxiety}; if priors are too loose, allostatic
forecasting may fail. This is not a replacement for clinical diagnosis
but a \textbf{generative} framing that links behaviour (avoidance,
reassurance seeking, substance use) to error-minimisation strategies.
Therapies that improve interoceptive accuracy (certain mindfulness
protocols, graded exposure, interoceptive training) can be understood as
recalibrating precisions and priors --- though the neuroscience is still
being mapped experimentally.

\begin{center}\rule{0.5\linewidth}{0.5pt}\end{center}

\subsection{The Free Energy Principle in
Context}\label{the-free-energy-principle-in-context}

The FEP is a \textbf{theoretical framework}, not a claim about specific
neural mechanisms. It provides a unifying mathematical language, but
multiple neural implementations are consistent with it. It should be
understood alongside:

\begin{itemize}
\tightlist
\item
  \textbf{Reinforcement learning} (reward-maximising agents). Active
  inference subsumes RL in the limit where epistemic value is zero (no
  information-seeking).
\item
  \textbf{Optimal control theory} (minimum-cost action selection).
  Active inference subsumes OC in the limit where the generative model
  perfectly captures dynamics.
\item
  \textbf{Information-theoretic accounts} of neural coding. Predictive
  coding implements rate-distortion-optimal sensory compression.
\item
  \textbf{Cybernetics} (Wiener; Ashby). Cybernetic homeostasis is the
  special case of active inference with a fixed prior and reactive
  correction.
\end{itemize}

The FEP is most powerful as a \textbf{design principle}: living systems
\emph{behave as if} they are minimising free energy, just as the lens
\emph{behaves as if} it minimises the time of light travel (Fermat's
principle). Whether neurons ``really'' compute KL divergences or just
behave as if they do is a separate empirical question.

\begin{quote}
\textbf{Concept Check 7:} Reinforcement learning maximises expected
reward; optimal control minimises expected cost; active inference
minimises expected free energy. State a biological prediction that would
distinguish active inference from pure RL --- i.e.~a behaviour that an
RL agent would not produce but an active-inference agent would.
\end{quote}

\begin{center}\rule{0.5\linewidth}{0.5pt}\end{center}

\subsection{Current Evidence and Frontier
Biology}\label{current-evidence-and-frontier-biology-2}

For \textbf{Active Inference and the Free Energy Principle}, frontier
biology belongs inside the evidence logic of the chapter. Systems models
are useful when they expose assumptions, uncertainty, and failure modes
rather than merely producing elegant diagrams. The core reading question
is this: active-inference explanations must connect hidden states,
sensory evidence, action, and measurable prediction error.

\begin{itemize}
\tightlist
\item
  \textbf{What to verify:} identify the observation, model, assay, or
  dataset that would make the claim stronger or weaker.
\item
  \textbf{What to qualify:} state the scale, organism, cell type,
  environmental condition, or population where the claim is expected to
  hold.
\item
  \textbf{What to compare:} test at least one alternative explanation,
  baseline, or null model before treating the pattern as causal.
\item
  \textbf{What to cite:} distinguish primary evidence, review synthesis,
  public dataset, and institutional guidance; for recent or numeric
  claims, prefer the source closest to the measurement and state what
  has changed since it was published.
\end{itemize}

Treat every model as a claim about mechanism: define the system
boundary, identify the observable that would falsify the model, and
report the uncertainty that would change a decision.

\textbf{Source practice:} Use model-validation sources when available,
and state which observation would falsify the model rather than treating
a diagram as proof.

\subsection{Summary}\label{summary-2}

\begin{itemize}
\tightlist
\item
  The free energy principle proposes that living systems minimise
  variational free energy --- a measure of the gap between their
  generative model and sensory evidence.
\item
  Markov blankets formalise the boundary between an agent and its
  environment, partitioning most states into internal, sensory, active,
  and external classes.
\item
  The Bayesian brain maintains hierarchical generative models that
  generate top-down predictions; prediction errors drive belief
  updating; the canonical microcircuit implements this in cortical
  layers.
\item
  Active inference distinguishes perceptual inference (update beliefs)
  from active inference (act to fulfil predictions).
\item
  Homeostasis and allostasis are naturally captured as prior-driven
  active fulfilment of body-state predictions.
\item
  Expected free energy unifies exploration (epistemic value) and
  exploitation (instrumental value) in a single functional, recovering
  reinforcement learning and optimal experimental design as special
  cases.
\item
  Active inference connects to evolution (niche construction),
  metabolism (energetic constraints on inference), neuroscience
  (predictive coding circuits), immunology, development, and
  computational psychiatry.
\item
  \textbf{Precision} weights prediction errors; \textbf{attention} and
  \textbf{sensory attenuation} are interpretable as precision control;
  \textbf{interoception} links body state to psychiatric phenomenology
  at a systems level.
\item
  Computational psychiatry offers FEP-grounded provisional accounts of
  anxiety, depression, schizophrenia, and autism --- useful as research
  frameworks while clinical practice continues to rely on evidence-based
  treatment.
\end{itemize}

\begin{center}\rule{0.5\linewidth}{0.5pt}\end{center}

\subsection{Key Terms}\label{key-terms-2}

\textbf{free energy} \texorpdfstring{\ensuremath{\cdot}}{.} \textbf{variational free energy} \texorpdfstring{\ensuremath{\cdot}}{.}
\textbf{expected free energy} \texorpdfstring{\ensuremath{\cdot}}{.} \textbf{Markov blanket} \texorpdfstring{\ensuremath{\cdot}}{.}
\textbf{generative model} \texorpdfstring{\ensuremath{\cdot}}{.} \textbf{prior belief} \texorpdfstring{\ensuremath{\cdot}}{.} \textbf{prediction
error} \texorpdfstring{\ensuremath{\cdot}}{.} \textbf{predictive coding} \texorpdfstring{\ensuremath{\cdot}}{.} \textbf{canonical microcircuit} \texorpdfstring{\ensuremath{\cdot}}{.}
\textbf{Bayesian brain} \texorpdfstring{\ensuremath{\cdot}}{.} \textbf{active inference} \texorpdfstring{\ensuremath{\cdot}}{.}
\textbf{homeostasis} \texorpdfstring{\ensuremath{\cdot}}{.} \textbf{allostasis} \texorpdfstring{\ensuremath{\cdot}}{.} \textbf{interoception} \texorpdfstring{\ensuremath{\cdot}}{.}
\textbf{niche construction} \texorpdfstring{\ensuremath{\cdot}}{.} \textbf{4E cognition} \texorpdfstring{\ensuremath{\cdot}}{.} \textbf{precision}
\texorpdfstring{\ensuremath{\cdot}}{.} \textbf{sensory attenuation} \texorpdfstring{\ensuremath{\cdot}}{.} \textbf{aberrant salience} \texorpdfstring{\ensuremath{\cdot}}{.}
\textbf{epistemic value} \texorpdfstring{\ensuremath{\cdot}}{.} \textbf{instrumental value} \texorpdfstring{\ensuremath{\cdot}}{.} \textbf{policy
selection}

\begin{center}\rule{0.5\linewidth}{0.5pt}\end{center}

\subsection{Discussion Questions}\label{discussion-questions-2}

\begin{enumerate}
\def\labelenumi{\arabic{enumi}.}
\item
  Explain the difference between \textbf{perceptual inference} and
  \textbf{active inference} using a concrete physiological example
  (e.g., cardiovascular regulation during exercise, or immune response
  to infection). Why does the distinction matter clinically?
\item
  A patient with \textbf{interoceptive dysregulation} has a poorly
  calibrated internal model of body states --- they cannot accurately
  predict or infer their own visceral sensations. Using active
  inference, explain how this might manifest as anxiety, depression, or
  eating disorders. What therapeutic interventions might improve
  interoceptive precision?
\item
  The bacterium \emph{E. coli} performs chemotaxis --- swimming toward
  attractants, tumbling to reorient --- guided by a two-component
  receptor-kinase system. Can this behaviour be described in active
  inference terms? What would the ``generative model,'' ``prediction
  error,'' and ``action'' be for \emph{E. coli}? Identify its Markov
  blanket explicitly.
\item
  Consider the \textbf{allostatic} adjustment of blood pressure in
  anticipation of orthostatic challenge (standing up). Compare this to
  the homeostatic model where blood pressure is corrected \emph{after}
  it drops. Which better fits the empirical timeline of cardiovascular
  adjustment? How does this support an active inference account of
  autonomic regulation?
\item
  The free energy principle has been compared to Helmholtz's observation
  that perception is ``unconscious inference.'' Is this analogy
  illuminating or misleading? What does the FEP add beyond classical
  Bayesian perception theories \citep{friston2010}? What does it risk
  conflating?
\item
  In one paragraph, explain why down-regulating somatosensory precision
  during a self-generated movement might prevent you from tickling
  yourself. How could a failure of this mechanism feel subjectively?
\item
  Propose an experiment (behavioural or neuroimaging) that would test
  whether anxious individuals assign \textbf{too much} precision to
  interoceptive prediction errors compared with controls. What would
  falsify your hypothesis?
\item
  The expected free energy functional contains both an epistemic
  (information-gain) term and an instrumental (reward) term. Construct a
  hypothetical scenario in which a pure reinforcement-learning agent and
  a pure active-inference agent would behave differently, and explain
  which behaviour seems more biologically realistic.
\end{enumerate}

\begin{center}\rule{0.5\linewidth}{0.5pt}\end{center}

\subsection{Review Questions}\label{review-questions-2}

\begin{enumerate}
\def\labelenumi{\arabic{enumi}.}
\item
  Define the four state classes of a Markov blanket (\texorpdfstring{\ensuremath{\eta}}{eta}, \(s\), \(a\), \texorpdfstring{\ensuremath{\mu}}{mu})
  and state the conditional-independence property they enforce. Why does
  this partition make the boundary between an organism and its
  environment mathematically precise rather than merely intuitive?
\item
  Explain in your own words why minimising variational free energy \(F\)
  is a tractable proxy for maximising model evidence, given that the
  true posterior \(P(\mathbf{z}\mid\mathbf{o})\) is generally
  intractable. Which term of \(F\) penalises over-complex beliefs, and
  which rewards accurate ones?
\item
  Distinguish perceptual inference from active inference using the
  chapter's body-temperature worked example. The agent's free energy
  falls from about 4.5 nats to about 2.1 nats by
  belief update but approaches 0 by shivering. Explain why action
  reduces free energy \emph{more completely} when the temperature prior
  is tight.
\item
  Apply the Markov-blanket partition to \emph{E. coli} chemotaxis:
  assign concrete molecular identities to the sensory, active, internal,
  and external states, and identify which assignment is least clear-cut
  and why.
\item
  Using the expected-free-energy functional \(G(\pi)\), analyse the
  worked policy comparison (\(\pi_A\): epistemic 1.0, instrumental 4.0;
  \(\pi_B\): epistemic 3.0, instrumental 2.5). Recompute \(G\) for each,
  justify why \(\pi_B\) is selected slightly more often, and predict how
  lowering the softmax temperature changes the choice.
\item
  The chapter reframes homeostasis and allostasis in active-inference
  terms. Compare the two as claims about prior precision and set-point
  dynamics, and explain why allostasis predicts heart rate rising
  \emph{before} exercise begins whereas a purely homeostatic model does
  not.
\item
  Computational psychiatry interprets anxiety, depression,
  schizophrenia, and autism as different failures of precision
  regulation. Analyse how ``over-precise priors on negative outcomes''
  versus ``aberrant salience from amplified noise prediction errors''
  would produce different observable behaviours, and propose one
  measurement that could distinguish them.
\item
  Evaluate the claim that natural selection itself can be cast as
  gradient descent on a long-timescale free-energy functional. What does
  the ``as-if'' qualifier protect against, and what evidence would make
  this extension more (or less) defensible than the neuroscience
  applications?
\item
  Critically assess the FEP as a unifying framework that subsumes
  reinforcement learning and optimal control as limiting cases.
  Construct a biological prediction that an active-inference agent would
  make but a pure reward-maximiser would not, and explain why the
  difference is empirically meaningful.
\item
  Synthesis: the chapter argues that the current frontier is
  \emph{identifiability} --- very different generative models can fit
  the same behaviour. Design a research strategy (combining
  out-of-sample prediction, parameter recovery, and a falsifiable neural
  or behavioural observable) that would let a sceptic decide whether an
  active-inference model genuinely explains a given dataset rather than
  merely curve-fitting it.
\end{enumerate}

\begin{center}\rule{0.5\linewidth}{0.5pt}\end{center}

\subsection{Further Reading and Source
Notes}\label{further-reading-and-source-notes-2}

\begin{itemize}
\tightlist
\item
  Friston, K. (2010). The free-energy principle: A unified brain theory?
  \citep{friston2010} \emph{Nature Reviews Neuroscience}, 11(2),
  127--138.
\item
  Friston, K., FitzGerald, T., Rigoli, F., Schwartenbeck, P., \&
  Pezzulo, G. (2017). Active inference: A process theory
  \citep{friston2017}. \emph{Neural Computation}, 29(1), 1--49.
\item
  Sterling, P., \& Laughlin, S. (2015). \emph{Principles of Neural
  Design} \citep{sterling2015}. MIT Press.
\item
  Sterling, P., \& Eyer, J. (1988). Allostasis: A new paradigm to
  explain arousal pathology \citep{sterling1988}.
\item
  Clark, A. (2016). \emph{Surfing Uncertainty: Prediction, Action, and
  the Embodied Mind}. Oxford University Press.
\item
  Parr, T., Pezzulo, G., \& Friston, K. J. (2022). \emph{Active
  Inference: The Free Energy Principle in Mind, Brain, and Behavior}.
  MIT Press.
\end{itemize}

\begin{center}\rule{0.5\linewidth}{0.5pt}\end{center}

\subsubsection{Companion Source Module}\label{companion-source-module-2}

\textbf{Active Inference and the Free Energy Principle} should leave a
reproducible trail from a biological claim to the code, figure, diagram,
or paper-based activity that can test it. Use the surfaces below to
inspect the chapter's assumptions, rerun the relevant model, or compare
the manuscript explanation with companion labs and figures.

{\def\LTcaptype{none} % do not increment counter
\begin{longtable}[]{@{}
  >{\raggedright\arraybackslash}p{(\linewidth - 2\tabcolsep) * \real{0.5000}}
  >{\raggedright\arraybackslash}p{(\linewidth - 2\tabcolsep) * \real{0.5000}}@{}}
\toprule\noalign{}
\begin{minipage}[b]{\linewidth}\raggedright
Surface
\end{minipage} & \begin{minipage}[b]{\linewidth}\raggedright
Use it for
\end{minipage} \\
\midrule\noalign{}
\endhead
\bottomrule\noalign{}
\endlastfoot
\breaktt{src/biology/neuroscience/neuroscience.py}
(\breaktt{action\_potential\_hh}, \breaktt{hebbian\_weight\_update}) &
Connect prediction, update, and plasticity to measurable neural
variables. \\
\breaktt{src/biology/physiology/physiology.py}
(\breaktt{homeostasis\_response}) & Compare allostatic regulation with
error-correcting control. \\
\breaktt{src/biology/cell/cell\_biology.py}
(\breaktt{receptor\_occupancy}, \breaktt{signal\_amplification}) & Make
sensing and gain explicit rather than metaphorical. \\
\breaktt{src/mermaid/biology\_diagrams.py}
(\breaktt{nervous\_system\_reflex\_diagram},
\breaktt{hormone\_signaling\_diagram}) & Contrast reflex arcs, endocrine
loops, and inference-style control diagrams. \\
\end{longtable}
}

\textbf{Reproducibility check:} name the hidden state, observation,
action, and error term before treating a biological feedback loop as
active inference. \textbf{Cross-reference:} compare with
\cref{sec:unit_IX_nervous_system},
\cref{sec:unit_IX_circulation_respiration_homeostasis}, and
\cref{sec:unit_II_cell_signaling}.

\newpage

\section*{Unit I --- Chemistry of Life:
Introduction}\label{unit-i-chemistry-of-life-introduction}
\addcontentsline{toc}{section}{Unit I --- Chemistry of Life:
Introduction}

\subsection*{Why This Unit Matters}\label{why-this-unit-matters-1}
\addcontentsline{toc}{subsection}{Why This Unit Matters}

Every organism on Earth --- a bacterium dividing in a hydrothermal vent,
a redwood drawing water through 100 metres of trunk, a neuron firing in
your brain right now --- is made of the same about 25 elements
obeying the same quantum-mechanical rules identified over the last two
centuries. Chemistry is not merely a prerequisite for biology; it
\textbf{is} biology at its most fundamental level. Understanding the
electron orbitals of carbon explains why proteins fold. Understanding
electronegativity explains why DNA holds its shape. Understanding pH
buffers explains why a runner's muscles do not seize in their own acid.

This unit begins with atoms and ends with molecular machines. We trace
how the CHNOPS elements (carbon, hydrogen, nitrogen, oxygen, phosphorus,
sulfur) combine into functional groups, how functional groups assemble
into the four classes of biological macromolecules, and how enzymes ---
protein catalysts evolved over billions of years --- accelerate
reactions to the speeds that life demands. Along the way you will
encounter the Pauling electronegativity scale, the Henderson-Hasselbalch
equation, the Michaelis-Menten model, and the principle of chirality:
the mirror-image selectivity that means L-amino acids build your
proteins while D-sugars fuel your cells.

A thread connecting every chapter: \textbf{water}. The anomalous
properties of water --- its high boiling point, the cohesion-tension
mechanism in plants, its role as both solvent and reactant --- make
Earth habitable and make most aqueous biochemistry possible. As the
physical chemist Lawrence Henderson (1913) argued in \emph{The Fitness
of the Environment}, water is not an incidental solvent but the molecule
around which life has been shaped by evolution.

\begin{center}\rule{0.5\linewidth}{0.5pt}\end{center}

\subsection*{Landmark Discoveries}\label{landmark-discoveries-1}
\addcontentsline{toc}{subsection}{Landmark Discoveries}

{\def\LTcaptype{none} % do not increment counter
\begin{longtable}[]{@{}
  >{\raggedright\arraybackslash}p{(\linewidth - 8\tabcolsep) * \real{0.2407}}
  >{\raggedright\arraybackslash}p{(\linewidth - 8\tabcolsep) * \real{0.0741}}
  >{\raggedright\arraybackslash}p{(\linewidth - 8\tabcolsep) * \real{0.2963}}
  >{\raggedright\arraybackslash}p{(\linewidth - 8\tabcolsep) * \real{0.1667}}
  >{\raggedright\arraybackslash}p{(\linewidth - 8\tabcolsep) * \real{0.2222}}@{}}
\toprule\noalign{}
\begin{minipage}[b]{\linewidth}\raggedright
Discoverer(s)
\end{minipage} & \begin{minipage}[b]{\linewidth}\raggedright
Year
\end{minipage} & \begin{minipage}[b]{\linewidth}\raggedright
Journal / Source
\end{minipage} & \begin{minipage}[b]{\linewidth}\raggedright
Discovery
\end{minipage} & \begin{minipage}[b]{\linewidth}\raggedright
Significance
\end{minipage} \\
\midrule\noalign{}
\endhead
\bottomrule\noalign{}
\endlastfoot
Friedrich Wöhler & 1828 & \emph{Annalen der Physik} & Synthesis of urea
from inorganic ammonium cyanate & Disproved vitalism; showed organic
compounds obey the same chemistry as inorganic ones \\
Linus Pauling \& Robert Corey & 1951 & \emph{Proc. Natl. Acad. Sci.} &
\texorpdfstring{\ensuremath{\alpha}}{alpha}-helix and \texorpdfstring{\ensuremath{\beta}}{beta}-sheet structures of proteins & Established that protein
secondary structure is stabilised by H-bonds; first use of
electronegativity to predict biological structure \\
James Watson \& Francis Crick & 1953 & \emph{Nature} & Double-helix
structure of DNA & Anti-parallel, complementary base-pairing via
H-bonds; immediately suggested replication mechanism \\
Frederic Sanger & 1955 & \emph{Biochemical Journal} & Complete amino
acid sequence of bovine insulin & Proved proteins are defined covalent
sequences, not random polymers; introduced sequence concept \\
Leonor Michaelis \& Maud Menten & 1913 & \emph{Biochemische Zeitschrift}
& Enzyme-substrate saturation kinetics & Derived
\(v = V_{max}[S]/(K_m + [S])\); quantitative foundation of enzymology \\
Daniel Koshland & 1958 & \emph{Proc. Natl. Acad. Sci.} & Induced-fit
model of enzyme catalysis & Replaced lock-and-key; explained
conformational changes in enzyme-substrate interaction \\
Jacques Monod, Jeffries Wyman \& Jean-Pierre Changeux & 1965 & \emph{J.
Mol. Biol.} & MWC model of allosteric regulation & Showed enzymes have
distinct regulatory and catalytic sites; explains cooperative binding \\
\end{longtable}
}

\begin{center}\rule{0.5\linewidth}{0.5pt}\end{center}

\subsection*{Key Concepts and
Connections}\label{key-concepts-and-connections-1}
\addcontentsline{toc}{subsection}{Key Concepts and Connections}

\begin{figure}[htbp]
\centering
\includegraphics[width=0.82\linewidth,height=4.2in,keepaspectratio]{../figures/mermaid_inline/inline_mermaid_0014_d64907ce8bd6.png}
\caption{Unit I concept map - Chemistry of Life. Arrows show conceptual dependencies; colour groups: blue = atomic/molecular; orange = enzyme function.}
\end{figure}

\begin{center}\rule{0.5\linewidth}{0.5pt}\end{center}

\subsection*{Current Evidence Thread}\label{current-evidence-thread-1}
\addcontentsline{toc}{subsection}{Current Evidence Thread}

Read Unit I as the point where the chemistry of life becomes measurable.
Bond polarity, hydrogen bonding, and reaction thermodynamics are not
abstractions here -- they are quantities now read off cryo-EM maps,
ultrafast spectroscopy, and kinetic assays, and structure predictions
for macromolecules are accepted primarily once cryo-EM, NMR, or
mutagenesis confirms them. As you move through the chapters, keep a
two-column note: \textbf{claim} on the left, \textbf{evidence that would
change my confidence} on the right. By the end of the unit, each major
idea should be tied to a measurement, model, citation, or paper-based
lab decision.

\subsection*{Chapter Roadmap}\label{chapter-roadmap-1}
\addcontentsline{toc}{subsection}{Chapter Roadmap}

{\def\LTcaptype{none} % do not increment counter
\begin{longtable}[]{@{}
  >{\raggedright\arraybackslash}p{(\linewidth - 6\tabcolsep) * \real{0.1556}}
  >{\raggedright\arraybackslash}p{(\linewidth - 6\tabcolsep) * \real{0.1111}}
  >{\raggedright\arraybackslash}p{(\linewidth - 6\tabcolsep) * \real{0.2889}}
  >{\raggedright\arraybackslash}p{(\linewidth - 6\tabcolsep) * \real{0.4444}}@{}}
\toprule\noalign{}
\begin{minipage}[b]{\linewidth}\raggedright
Chapter
\end{minipage} & \begin{minipage}[b]{\linewidth}\raggedright
Title
\end{minipage} & \begin{minipage}[b]{\linewidth}\raggedright
Core Question
\end{minipage} & \begin{minipage}[b]{\linewidth}\raggedright
Key Equation / Model
\end{minipage} \\
\midrule\noalign{}
\endhead
\bottomrule\noalign{}
\endlastfoot
\textbf{1} & Atoms, Molecules, and Chemical Bonds & Why do atoms bond,
and how does bond type determine biological function? &
\(\Delta G^{\circ\prime} = -nF\Delta E^{\circ\prime}\) \\
\textbf{2} & Water --- The Molecule of Life & What makes water
anomalous, and why is that anomaly essential for life? & \(\pi = iCRT\)
(osmotic pressure) \\
\textbf{3} & Biological Macromolecules & How do four classes of polymer
encode and carry out life's functions? & Molecular weight, degree of
polymerisation \\
\textbf{4} & Enzymes and the Kinetics of Catalysis & How do enzymes
lower activation energy, and how is their activity regulated? &
\(v = V_{max}[S]/(K_m + [S])\) \\
\end{longtable}
}

\begin{center}\rule{0.5\linewidth}{0.5pt}\end{center}

\subsection*{Connections Across the
Textbook}\label{connections-across-the-textbook-1}
\addcontentsline{toc}{subsection}{Connections Across the Textbook}

Unit I establishes the chemical vocabulary for the entire textbook:

\begin{itemize}
\tightlist
\item
  The \textbf{bond energies and functional groups} introduced here
  reappear in Unit III (energy released by glucose oxidation), Unit IV
  (phosphodiester bonds of DNA), and Unit V (crossing-over in
  chromatin).
\item
  \textbf{Enzyme kinetics} (\cref{sec:unit_I_enzymes_and_kinetics})
  reappears in Unit III (PFK-1 allosteric regulation), Unit VII
  (antibiotic targets on bacterial enzymes), and Unit IX (haemoglobin
  co-operativity).
\item
  \textbf{pH and buffers} link directly to Unit II (lysosomal pH 4.5
  vs.~cytoplasmic pH 7.2), Unit III (mitochondrial proton gradient), and
  Unit IX (blood acid-base balance in clinical contexts).
\item
  \textbf{Chirality} of L-amino acids and D-sugars reappears in Unit VII
  (D-amino acids in bacterial cell walls) and Unit V (enzyme specificity
  in DNA repair).
\end{itemize}

\begin{quote}
\textbf{Before you begin:} You should be comfortable with scientific
notation, logarithms, and basic algebra. No prior chemistry or biology
is assumed.
\end{quote}

\subsection*{Computational Toolbox --- Unit
I}\label{computational-toolbox-unit-i}
\addcontentsline{toc}{subsection}{Computational Toolbox --- Unit I}

The following demonstrates the real \breaktt{src/biology/biochemistry}
functions used to generate figures and worked examples throughout Unit
I.

\begin{Shaded}
\begin{Highlighting}[]
\ImportTok{from}\NormalTok{ biology.biochemistry }\ImportTok{import}\NormalTok{ michaelis\_menten, glycolysis\_summary}

\CommentTok{\# Michaelis{-}Menten kinetics: lactate dehydrogenase (LDH)}
\CommentTok{\# Vmax = 12.0 mM/s, Km = 0.8 mM (pyruvate as substrate)}
\NormalTok{result }\OperatorTok{=}\NormalTok{ michaelis\_menten(substrate\_conc}\OperatorTok{=}\FloatTok{2.0}\NormalTok{, Vmax}\OperatorTok{=}\FloatTok{12.0}\NormalTok{, Km}\OperatorTok{=}\FloatTok{0.8}\NormalTok{)}
\BuiltInTok{print}\NormalTok{(}\SpecialStringTok{f"Reaction rate at [S]=2 mM: }\SpecialCharTok{\{}\NormalTok{result}\SpecialCharTok{.}\NormalTok{reaction\_rate}\SpecialCharTok{:.2f\}}\SpecialStringTok{ mM/s"}\NormalTok{)}
\CommentTok{\# Expected: Reaction rate at [S]=2 mM: 8.57 mM/s}
\CommentTok{\# (= 12.0 × 2.0 / (0.8 + 2.0); Vmax/2 occurs at [S]=Km=0.8 mM)}

\CommentTok{\# Glycolysis summary: net products per glucose}
\NormalTok{summary }\OperatorTok{=}\NormalTok{ glycolysis\_summary()}
\BuiltInTok{print}\NormalTok{(}\SpecialStringTok{f"Net ATP yield: }\SpecialCharTok{\{}\NormalTok{summary}\SpecialCharTok{.}\NormalTok{net\_atp}\SpecialCharTok{\}}\SpecialStringTok{ ATP per glucose"}\NormalTok{)}
\BuiltInTok{print}\NormalTok{(}\SpecialStringTok{f"NADH produced: }\SpecialCharTok{\{}\NormalTok{summary}\SpecialCharTok{.}\NormalTok{net\_nadh}\SpecialCharTok{\}}\SpecialStringTok{ NADH per glucose"}\NormalTok{)}
\BuiltInTok{print}\NormalTok{(}\SpecialStringTok{f"Steps represented: }\SpecialCharTok{\{}\BuiltInTok{len}\NormalTok{(summary.steps)}\SpecialCharTok{\}}\SpecialStringTok{"}\NormalTok{)}
\CommentTok{\# Expected:}
\CommentTok{\# Net ATP yield: 2 ATP per glucose}
\CommentTok{\# NADH produced: 2 NADH per glucose}
\CommentTok{\# Steps represented: 10}
\end{Highlighting}
\end{Shaded}

\begin{quote}
\textbf{Try it yourself:} Install the project environment
(\texttt{uv\ sync}) and run
\breaktt{python3\ -c\ "from\ biology.biochemistry\ import\ michaelis\_menten;\ print(michaelis\_menten(5,\ 10,\ 2))"}.
Modify substrate concentration to trace the full saturation curve.
\end{quote}

\begin{center}\rule{0.5\linewidth}{0.5pt}\end{center}

\emph{Source modules: \breaktt{src/biology/biochemistry/biochemistry.py}
--- enzyme kinetics models, macromolecule analysis.} \emph{Figures
generated by \breaktt{src/visualization/plots.py} (Michaelis-Menten
plots) and \breaktt{src/mermaid/biology\_diagrams.py}.}

\newpage

\section{Atoms, Molecules, and Chemical
Bonds}\label{atoms-molecules-and-chemical-bonds}

\label{sec:unit_I_atoms_molecules}

\begin{quote}
\textbf{Ch 1} \texorpdfstring{\ensuremath{\cdot}}{.} Level 1/3 \texorpdfstring{\ensuremath{\cdot}}{.} 40 min read \texorpdfstring{\ensuremath{\cdot}}{.} 50 min lecture \texorpdfstring{\ensuremath{\cdot}}{.}
Prerequisites: none
\end{quote}

\subsection{Learning Objectives}\label{learning-objectives-3}

By the end of this chapter, you should be able to:

\begin{enumerate}
\def\labelenumi{\arabic{enumi}.}
\tightlist
\item
  Describe the structure of an atom, including subatomic particles and
  their quantum mechanical behaviour.
\item
  Explain how
  \hyperref[gl:electronegativity]{\textbf{electronegativity}} governs
  covalent and ionic bond formation.
\item
  Distinguish between polar covalent, nonpolar covalent, and ionic bonds
  with biological examples.
\item
  Explain the roles of \hyperref[gl:hydrogen-bond]{\textbf{hydrogen
  bond}}s and \hyperref[gl:van-der-waals-forces]{\textbf{van der Waals
  forces}} in biological macromolecules.
\item
  Interpret the concept of \hyperref[gl:ph]{\textbf{pH}},
  \hyperref[gl:buffer]{\textbf{buffer}}s, and their physiological
  significance.
\item
  Describe \hyperref[gl:chirality]{\textbf{chirality}} and
  stereoisomerism in biological molecules.
\item
  Explain redox reactions and their central role in bioenergetics.
\item
  Apply isotope chemistry to biomedical imaging and radiotherapy.
\end{enumerate}

\subsubsection{Study Blueprint}\label{study-blueprint-3}

\begin{itemize}
\tightlist
\item
  \textbf{Big idea:} Atomic structure constrains the bonds, charges, and
  geometries that make biological molecules possible.
\item
  \textbf{Core concepts:} valence, electronegativity, formal charge,
  isotopes.
\item
  \textbf{Framework alignment:} Vision \& Change: Structure and
  function, Pathways and transformations of energy and matter; AP
  Biology: Energetics, Systems Interactions; NGSS-style topics: Matter
  and Energy in Organisms and Ecosystems, Structure and Function.
\item
  \textbf{Model or quantitative lens:} Formal charge and
  electronegativity-difference reasoning.
\item
  \textbf{Data skill:} Use tabular atomic data to predict polarity,
  solubility, and biological reactivity.
\item
  \textbf{Practice cadence:} Concept Explanation, Statistical Tests and
  Data Analysis, Argumentation.
\item
  \textbf{Common misconception to repair:} Weak interactions are not
  unimportant; many weak interactions together dominate structure and
  specificity.
\item
  \textbf{Primary lab:} \cref{sec:lab_unit_I_atoms_molecules}.
\item
  \textbf{Question bank:} \cref{sec:q_unit_I_atoms_molecules}.
\item
  \textbf{Transfer task:} Use atomic reasoning to explain a medical
  tracer, enzyme cofactor, or membrane-solubility problem.
\item
  \textbf{Bridge to computation:}
  \breaktt{biology.biochemistry.biochemistry.atp\_free\_energy}.
\end{itemize}

\begin{center}\rule{0.5\linewidth}{0.5pt}\end{center}

\begin{quote}
\textbf{Opening Vignette: The Radioactive Tracer That Exposed Cancer}

In 1973, chemist Alfred Wolf and colleagues at Brookhaven National
Laboratory synthesised a new molecule: fluorine-18-labelled deoxyglucose
(\(^{18}\)F-FDG). Because \(^{18}\)F decays with a half-life of just
109.8 minutes, emitting a positron that collides with a nearby electron
to produce a pair of detectable gamma rays, the molecule could be
tracked in real time inside a living body. Because FDG mimics glucose,
metabolically hungry cells take it up --- and cancer cells, with their
vastly elevated glucose consumption (the \textbf{Warburg effect}, first
described by Otto Warburg in 1924), accumulate it at concentrations
2--10\texorpdfstring{\ensuremath{\times}}{x} higher than surrounding normal tissue.

Today, PET (positron emission tomography) scanning using \(^{18}\)F-FDG
is performed over 2 million times per year in the United States alone,
guiding cancer diagnosis, staging, and treatment response assessment.
None of it would be possible without understanding the atomic and
subatomic physics of radioactive decay, the electronegativity that
distinguishes O--H from C--H bonds, and the biochemistry of glucose
transport --- precisely the subjects of this chapter.

\emph{Primary source: Warburg, O. (1924). Über den Stoffwechsel der
Carcinomzelle. Naturwissenschaften, 12(50), 1131--1137.}
\end{quote}

\begin{center}\rule{0.5\linewidth}{0.5pt}\end{center}

Life is fundamentally chemistry. Every organism, from a bacterium to a
blue whale, is composed of atoms obeying the same physical laws
discovered in the nineteenth and twentieth centuries. Understanding life
therefore begins with understanding atoms.

Of the 118 known elements, only about 25 are essential to life. Just six
elements --- carbon, hydrogen, nitrogen, oxygen, phosphorus, and sulfur
(CHNOPS) --- account for approximately 97\% of the mass of a human body.
Trace elements such as iron, zinc, copper, manganese, and selenium
perform critical catalytic roles despite being present in minuscule
quantities.

\subsubsection{Subatomic Particles}\label{subatomic-particles}

An atom consists of three primary subatomic particles:

{\def\LTcaptype{none} % do not increment counter
\begin{longtable}[]{@{}llll@{}}
\toprule\noalign{}
Particle & Charge & Mass (Da) & Location \\
\midrule\noalign{}
\endhead
\bottomrule\noalign{}
\endlastfoot
Proton & +1 & 1.007 & Nucleus \\
Neutron & 0 & 1.008 & Nucleus \\
Electron & \texorpdfstring{\ensuremath{-}}{-}1 & 0.000549 & Orbital shells \\
\end{longtable}
}

The \textbf{atomic number} (Z) equals the number of protons, defining
the element. The \textbf{mass number} (A) equals protons plus neutrons.
\textbf{Isotopes} share the same atomic number but differ in neutron
count.

\subsubsection{Isotopes and Their Biological
Applications}\label{isotopes-and-their-biological-applications}

Isotopes are classified as \textbf{stable} or \textbf{radioactive}
(radioisotopes). Radioisotopes undergo spontaneous nuclear decay,
emitting alpha particles, beta particles, or gamma rays.

{\def\LTcaptype{none} % do not increment counter
\begin{longtable}[]{@{}
  >{\raggedright\arraybackslash}p{(\linewidth - 6\tabcolsep) * \real{0.1250}}
  >{\raggedright\arraybackslash}p{(\linewidth - 6\tabcolsep) * \real{0.1607}}
  >{\raggedright\arraybackslash}p{(\linewidth - 6\tabcolsep) * \real{0.1786}}
  >{\raggedright\arraybackslash}p{(\linewidth - 6\tabcolsep) * \real{0.5357}}@{}}
\toprule\noalign{}
\begin{minipage}[b]{\linewidth}\raggedright
Isotope
\end{minipage} & \begin{minipage}[b]{\linewidth}\raggedright
Half-life
\end{minipage} & \begin{minipage}[b]{\linewidth}\raggedright
Decay Type
\end{minipage} & \begin{minipage}[b]{\linewidth}\raggedright
Biological/Medical Application
\end{minipage} \\
\midrule\noalign{}
\endhead
\bottomrule\noalign{}
\endlastfoot
\(^{14}\)C & 5,730 years & \(\beta^-\) & Radiocarbon dating of fossils
and archaeological specimens \\
\(^{3}\)H (tritium) & 12.3 years & \(\beta^-\) & Tracing metabolic
pathways in autoradiography \\
\(^{32}\)P & 14.3 days & \(\beta^-\) & Labelling DNA/RNA in molecular
biology \\
\(^{18}\)F & 109.8 min & \(\beta^+\) & PET (positron emission
tomography) imaging \\
\(^{131}\)I & 8.0 days & \(\beta^-\), \texorpdfstring{\ensuremath{\gamma}}{gamma} & Thyroid cancer radiotherapy \\
\(^{99m}\)Tc & 6.0 hours & \texorpdfstring{\ensuremath{\gamma}}{gamma} & SPECT imaging of bones, heart, kidneys \\
\end{longtable}
}

\begin{quote}
\textbf{Clinical Connection: PET Imaging and \(^{18}\)F-FDG}

Positron emission tomography (PET) exploits the short half-life of
\(^{18}\)F. The glucose analogue \(^{18}\)F-fluorodeoxyglucose
(\(^{18}\)F-FDG) is taken up by metabolically active cells. Because
cancer cells have elevated glucose uptake (the Warburg effect), PET
scans reveal tumours as ``hot spots.'' The emitted positron annihilates
with a nearby electron, producing two 511 keV gamma photons detected by
the scanner in coincidence.
\end{quote}

\begin{quote}
\textbf{Concept Check 1:} Carbon-14 has a half-life of 5,730 years. If a
fossilised bone contains 12.5\% of its original \(^{14}\)C,
approximately how old is it? (Hint: How many half-lives reduce 100\% to
12.5\%?)
\end{quote}

\subsubsection{Electron Shells, Orbitals, and
Valence}\label{electron-shells-orbitals-and-valence}

Electrons occupy discrete energy levels called \textbf{shells}
(principal quantum number n = 1, 2, 3, \ldots). Within each shell,
electrons occupy \textbf{orbitals} --- three-dimensional regions of
space where the probability of finding an electron exceeds 90\%.

{\def\LTcaptype{none} % do not increment counter
\begin{longtable}[]{@{}llll@{}}
\toprule\noalign{}
Shell (n) & Subshells & Orbitals & Max Electrons \\
\midrule\noalign{}
\endhead
\bottomrule\noalign{}
\endlastfoot
1 & 1s & 1 & 2 \\
2 & 2s, 2p & 4 & 8 \\
3 & 3s, 3p, 3d & 9 & 18 \\
\end{longtable}
}

The outermost occupied shell is the \textbf{valence shell} --- its
electron count determines chemical reactivity. Atoms are most stable
when their valence shell is complete (the \textbf{octet rule} for
second-period elements).

\textbf{Key biological elements and their valences:}

{\def\LTcaptype{none} % do not increment counter
\begin{longtable}[]{@{}lllll@{}}
\toprule\noalign{}
Element & Symbol & Atomic Number & Valence Electrons & Typical Bonds \\
\midrule\noalign{}
\endhead
\bottomrule\noalign{}
\endlastfoot
Hydrogen & H & 1 & 1 & 1 \\
Carbon & C & 6 & 4 & 4 \\
Nitrogen & N & 7 & 5 & 3 \\
Oxygen & O & 8 & 6 & 2 \\
Phosphorus & P & 15 & 5 & 5 \\
Sulfur & S & 16 & 6 & 2, 4, or 6 \\
\end{longtable}
}

Carbon's four valence electrons allow it to form four strong covalent
bonds in a tetrahedral geometry (\(sp^3\) hybridisation, bond angle
109.5 degrees), enabling virtually unlimited structural complexity ---
the molecular foundation of life.

\subsubsection{Orbital Hybridisation in Biological
Molecules}\label{orbital-hybridisation-in-biological-molecules}

Carbon's bonding versatility arises from orbital hybridisation --- the
linear combination of atomic \(s\) and \(p\) orbitals to form new hybrid
orbitals matched to molecular geometry.

\begin{itemize}
\tightlist
\item
  \textbf{\(sp^3\):} One \(s\) + three \(p\) orbitals -\textgreater{}
  four equivalent hybrid orbitals pointing to the corners of a
  tetrahedron (109.5 degrees). Each forms a single (\texorpdfstring{\ensuremath{\sigma}}{sigma}) bond. Example:
  methane, saturated carbons in \hyperref[gl:fatty-acid]{\textbf{fatty
  acid}} tails, the \texorpdfstring{\ensuremath{\alpha}}{alpha}-carbon of every amino acid.
\item
  \textbf{\(sp^2\):} One \(s\) + two \(p\) orbitals -\textgreater{}
  three coplanar hybrid orbitals (120 degrees) plus one unhybridised
  \(p\) orbital perpendicular to the plane that overlaps with another
  \(p\) orbital to form a \texorpdfstring{\ensuremath{\pi}}{pi}-bond. Example: C=C in unsaturated fatty
  acids, the carbonyl carbon of every
  \hyperref[gl:peptide-bond]{\textbf{peptide bond}}, aromatic rings.
\item
  \textbf{\(sp\):} One \(s\) + one \(p\) orbital -\textgreater{} two
  collinear hybrid orbitals (180 degrees) plus two \(p\) orbitals
  available for two \texorpdfstring{\ensuremath{\pi}}{pi}-bonds. Example: C\(\equiv\)N in hydrogen cyanide;
  the C\(\equiv\)C of acetylene; rare in biology except in some natural
  products.
\end{itemize}

Nitrogen exhibits analogous hybridisation: in primary amines
(--NH\(_2\)) and the \texorpdfstring{\ensuremath{\alpha}}{alpha}-amino group of amino acids, nitrogen is \(sp^3\)
with one lone pair (pyramidal geometry, bond angle about 107
degrees). In the planar peptide bond, however, the nitrogen lone pair is
delocalised into the carbonyl \texorpdfstring{\ensuremath{\pi}}{pi} system; the nitrogen behaves as
effectively \(sp^2\), locking the six atoms of the peptide unit
(C\(_\alpha\)--C(=O)--N(H)--C\(_\alpha\)) into a plane.

The \textbf{peptide bond} between amino acids has partial double-bond
character (about 40\%) because of resonance between the
\(sp^2\)-hybridised C and N atoms, enforcing planarity and the
\emph{trans} configuration that constrains
\hyperref[gl:protein]{\textbf{protein}} backbone geometry. The
\emph{cis} form is destabilised by about 8 kJ/mol relative to
\emph{trans} (except before proline, where the steric penalty is smaller
and about 5\% of X-Pro bonds are \emph{cis}).

\subsubsection{Lewis Structures and Formal
Charge}\label{lewis-structures-and-formal-charge}

A \textbf{Lewis structure} depicts every valence electron --- as bonding
pairs (lines) or lone pairs (dots) --- making explicit which atoms share
which electrons. The \textbf{formal charge} on an atom is what its
charge would be if most bonds were nonpolar (electrons shared equally):

\begin{equation}
\text{Formal charge} = (\text{valence electrons}) - (\text{lone-pair electrons}) - \tfrac{1}{2}(\text{bonding electrons})
\label{eq:unit_I_formal_charge}
\end{equation}

Formal charges sum to the molecular charge and help select among
possible Lewis structures: the best structure minimises formal charges
and places any negative formal charge on the most electronegative atom.

\textbf{Worked example --- formal charge on phosphate (PO\(_4^{3-}\)):}
Phosphate has central P bonded to four O atoms. In the most common
textbook structure, three P--O are single bonds (each O carries three
lone pairs) and one P=O is a double bond (the doubly bonded O carries
two lone pairs).

\begin{itemize}
\tightlist
\item
  \emph{Phosphorus (group 5, 5 valence electrons):} 0 lone-pair
  electrons; 5 bonds \(\times\) 2 electrons / 2 = 5 bonding electrons
  assigned. FC = 5 \texorpdfstring{\ensuremath{-}}{-} 0 \texorpdfstring{\ensuremath{-}}{-} 5 = \textbf{0}.
\item
  \emph{Each single-bonded O (with three lone pairs):} 6 lone-pair
  electrons; 1 bond \(\times\) 2 / 2 = 1 bonding electron assigned. FC =
  6 \texorpdfstring{\ensuremath{-}}{-} 6 \texorpdfstring{\ensuremath{-}}{-} 1 = \textbf{\texorpdfstring{\ensuremath{-}}{-}1} (\(\times\) 3 oxygens = \texorpdfstring{\ensuremath{-}}{-}3).
\item
  \emph{Double-bonded O (with two lone pairs):} 4 lone-pair electrons; 2
  bonds \(\times\) 2 / 2 = 2 bonding electrons. FC = 6 \texorpdfstring{\ensuremath{-}}{-} 4 \texorpdfstring{\ensuremath{-}}{-} 2 =
  \textbf{0}.
\end{itemize}

Total: 0 + 3(\texorpdfstring{\ensuremath{-}}{-}1) + 0 = \texorpdfstring{\ensuremath{-}}{-}3, matching the overall charge. By symmetry the
four oxygens are actually equivalent; resonance delocalises the three
negative charges and the double-bond character across the four P--O
bonds, giving an average bond order of 1.25 and bond length 154 pm
(intermediate between a P--O single bond at 163 pm and a P=O double bond
at 145 pm). This delocalisation explains why phosphate is such a stable
leaving group --- a key feature exploited by ATP, GTP, and the DNA/RNA
backbone.

\subsubsection{Electronegativity Drives Polarity in
Biology}\label{electronegativity-drives-polarity-in-biology}

The Pauling \textbf{electronegativity} (\texorpdfstring{\ensuremath{\chi}}{chi}) of an element predicts how
strongly it pulls bonding electrons. Differences in \texorpdfstring{\ensuremath{\chi}}{chi} across a bond
(\(\Delta\chi\)) determine the partial charges that drive
\hyperref[gl:hydrogen-bond]{\textbf{hydrogen bond}}ing, dipole--dipole
interactions, and acid-base behaviour throughout biology.

{\def\LTcaptype{none} % do not increment counter
\begin{longtable}[]{@{}llll@{}}
\toprule\noalign{}
Element & \texorpdfstring{\ensuremath{\chi}}{chi} (Pauling) & Period & Biological role \\
\midrule\noalign{}
\endhead
\bottomrule\noalign{}
\endlastfoot
F & 3.98 & 2 & Fluorinated drugs (5-FU, FDG) \\
O & 3.44 & 2 & Backbone of polar bonds in water, hydroxyl, carbonyl \\
N & 3.04 & 2 & Amino, amide, peptide, nucleobase H-bonds \\
Cl & 3.16 & 3 & Cl\(^-\) ion (electroneutrality, gastric HCl) \\
S & 2.58 & 3 & Disulfides, Fe--S clusters; weaker H-bonds than O \\
C & 2.55 & 2 & Backbone of organic chemistry \\
P & 2.19 & 3 & Phosphate (most polarity comes from O, not P) \\
H & 2.20 & 1 & H-bond donor when bonded to N/O/F \\
Na & 0.93 & 3 & Forms ionic Na\(^+\) in solution \\
K & 0.82 & 4 & Forms ionic K\(^+\) in solution \\
\end{longtable}
}

Three observations follow:

\begin{enumerate}
\def\labelenumi{\arabic{enumi}.}
\tightlist
\item
  \textbf{C--H bonds are nearly nonpolar} (\(\Delta\chi\) = 0.35), which
  is why the saturated tails of phospholipids and the methyl groups of
  valine/leucine/isoleucine drive the
  \hyperref[gl:hydrophobic-effect]{\textbf{hydrophobic effect}}.
\item
  \textbf{O--H and N--H bonds are strongly polar} (\(\Delta\chi\) = 1.24
  and 0.84), creating the partial charges that allow water, alcohols,
  and amides to form \hyperref[gl:hydrogen-bond]{\textbf{hydrogen
  bond}}s and to dissolve in aqueous environments.
\item
  \textbf{C=O is more polar than C--O} because in the carbonyl the \texorpdfstring{\ensuremath{\pi}}{pi}
  electrons pile up on oxygen; this polarity makes the carbonyl carbon
  electrophilic --- the site attacked by serine in trypsin, by water in
  \hyperref[gl:hydrolysis]{\textbf{hydrolysis}}, and by amine
  nucleophiles in transamination.
\end{enumerate}

The peptide bond's polarity (carbonyl O is \(\delta^-\), amide H is
\(\delta^+\)) is what enables the regular hydrogen bond pattern of
\texorpdfstring{\ensuremath{\alpha}}{alpha}-helices and \texorpdfstring{\ensuremath{\beta}}{beta}-sheets. Without the electronegativity difference between
O and C, there would be no protein secondary structure.

\begin{center}\rule{0.5\linewidth}{0.5pt}\end{center}

\subsection{Chemical Bonds}\label{chemical-bonds}

\begin{figure}[htbp]
\centering
\includegraphics[width=0.82\linewidth,height=4.2in,keepaspectratio]{../figures/mermaid_inline/inline_mermaid_0015_b5492bf82fbc.png}
\caption{Electronegativity and bond type. Classification of chemical bonds by electronegativity difference (DeltaEN). The boundaries (0.5 and 1.7) are approximate; real bonds exist on a continuum.}
\end{figure}

\subsubsection{Covalent Bonds}\label{covalent-bonds}

A \textbf{covalent bond} forms when two atoms share one or more pairs of
electrons. The strength of a covalent bond is approximately 150--950
kJ/mol, far greater than thermal energy at body temperature
(about 2.5 kJ/mol), making them the primary structural bonds of
organic molecules.

\textbf{Bond energies of biologically important bonds:}

{\def\LTcaptype{none} % do not increment counter
\begin{longtable}[]{@{}llll@{}}
\toprule\noalign{}
Bond & Energy (kJ/mol) & Bond Length (pm) & Biological Context \\
\midrule\noalign{}
\endhead
\bottomrule\noalign{}
\endlastfoot
C--C & 346 & 154 & Hydrocarbon chains \\
C=C & 614 & 134 & Unsaturated fatty acids \\
C--H & 413 & 109 & Ubiquitous in organics \\
C--O & 360 & 143 & Sugars, amino acids \\
C=O & 745 & 123 & Peptide bonds, carboxyl groups \\
C--N & 305 & 147 & Amino acids \\
O--H & 463 & 96 & Water, hydroxyl groups \\
N--H & 391 & 101 & Amino groups, peptide bonds \\
S--S & 266 & 205 & Disulfide bridges in proteins \\
P--O & 335 & 163 & DNA/RNA backbone, ATP \\
\end{longtable}
}

\textbf{Electronegativity} is the tendency of an atom to attract
electron pairs. The Pauling electronegativity scale (F = 3.98; O = 3.44;
N = 3.04; C = 2.55; H = 2.20) governs bond polarity:

\begin{itemize}
\tightlist
\item
  \textbf{Nonpolar covalent:} \(\Delta\chi\) \textless{} 0.5 (e.g.,
  C--C, C--H). Electrons shared equally. Found abundantly in lipid
  hydrocarbon tails.
\item
  \textbf{Polar covalent:} \(0.5 \le \Delta\chi < 1.7\) (e.g., O--H,
  N--H, C--O). Partial charges (\texorpdfstring{\ensuremath{\delta}}{delta}+ and \texorpdfstring{\ensuremath{\delta}}{delta}\texorpdfstring{\ensuremath{-}}{-}) arise. Critical for hydrogen
  bonding.
\item
  \textbf{Ionic:} \(\Delta\chi \ge 1.7\) (e.g., Na\(^+\)Cl\(^-\)). One
  atom effectively captures the electron(s). In solution, ionic
  compounds dissociate.
\end{itemize}

\textbf{Worked example:} Water (H\(_2\)O). Oxygen's electronegativity
(3.44) vs.~hydrogen's (2.20) gives \(\Delta\chi\) = 1.24 --- a polar
covalent bond. The two O--H bonds create a partial negative charge on
oxygen (\texorpdfstring{\ensuremath{\delta}}{delta}\texorpdfstring{\ensuremath{-}}{-} = \texorpdfstring{\ensuremath{-}}{-}0.834 in formal charge terms) and partial positive charges
on each hydrogen (\texorpdfstring{\ensuremath{\delta}}{delta}+ = +0.417 each). This polarity drives water's
exceptional biological properties (see
\cref{sec:unit_I_water_and_life}).

\subsubsection{Ionic Bonds}\label{ionic-bonds}

Ionic bonds form when one atom transfers electrons to another,
generating oppositely charged ions. In aqueous biological environments,
ionic bonds are typically disrupted by water molecules competing for the
ions. The \textbf{lattice energy} of NaCl is 787 kJ/mol, yet it
dissolves readily because of the high dielectric constant of water
(\(\varepsilon \approx 80\)).

In proteins, ionic interactions (\textbf{salt bridges}) between
oppositely charged residues stabilise higher-order structure. Aspartate
(Asp, D, --COO\(^-\)) and lysine (Lys, K, --NH\(_3^+\)) frequently form
salt bridges in \hyperref[gl:enzyme]{\textbf{enzyme}} active sites. A
typical salt bridge contributes 5--20 kJ/mol to protein stability, with
the strength depending on the local dielectric environment (buried salt
bridges are stronger because of lower effective dielectric constant).

\begin{quote}
\textbf{Concept Check 2:} Why do salt bridges buried in the hydrophobic
core of a protein contribute more to stability than those on the protein
surface? Consider the role of the dielectric constant.
\end{quote}

\subsubsection{Hydrogen Bonds}\label{hydrogen-bonds}

A \textbf{hydrogen bond} forms when a hydrogen atom covalently bonded to
an electronegative atom (N, O, or F) is attracted to another
electronegative atom. The bond energy is 2--40 kJ/mol --- weaker than
covalent bonds but stronger than van der Waals forces. Hydrogen bond
strengths and distances vary systematically with donor and acceptor
identity:

{\def\LTcaptype{none} % do not increment counter
\begin{longtable}[]{@{}
  >{\raggedright\arraybackslash}p{(\linewidth - 6\tabcolsep) * \real{0.1447}}
  >{\raggedright\arraybackslash}p{(\linewidth - 6\tabcolsep) * \real{0.3816}}
  >{\raggedright\arraybackslash}p{(\linewidth - 6\tabcolsep) * \real{0.1974}}
  >{\raggedright\arraybackslash}p{(\linewidth - 6\tabcolsep) * \real{0.2763}}@{}}
\toprule\noalign{}
\begin{minipage}[b]{\linewidth}\raggedright
H-bond type
\end{minipage} & \begin{minipage}[b]{\linewidth}\raggedright
Donor--Acceptor distance (pm)
\end{minipage} & \begin{minipage}[b]{\linewidth}\raggedright
Energy (kJ/mol)
\end{minipage} & \begin{minipage}[b]{\linewidth}\raggedright
Biological occurrence
\end{minipage} \\
\midrule\noalign{}
\endhead
\bottomrule\noalign{}
\endlastfoot
F--H\(\cdots\)F\(^-\) (low-barrier) & 240 & 50--170 & Rare in biology;
fluoride enzymes \\
O--H\(\cdots\)O\(^-\) (charged) & 250--260 & 25--40 & Salt bridges with
carboxylates \\
N\(^+\)--H\(\cdots\)O (charged) & 260--270 & 15--30 & Lys/Arg with
Asp/Glu \\
O--H\(\cdots\)O (neutral) & 270--280 & 15--25 & Water, hydroxyl
groups \\
N--H\(\cdots\)O=C (peptide) & 285--310 & 8--20 & \texorpdfstring{\ensuremath{\alpha}}{alpha}-helices, \texorpdfstring{\ensuremath{\beta}}{beta}-sheets \\
O--H\(\cdots\)N (neutral) & 280--295 & 10--20 & Ser--His in catalytic
triads \\
C--H\(\cdots\)O (weak) & 320--350 & 1--4 & \texorpdfstring{\ensuremath{\beta}}{beta}-sheet edges, base
stacking \\
S--H\(\cdots\)O & 320--360 & 4--10 & Cys hydrogen bonds (rare) \\
\end{longtable}
}

Despite their individual weakness, hydrogen bonds collectively confer
enormous structural stability:

\begin{itemize}
\tightlist
\item
  Water's liquid state at room temperature (vs.~the expected gas-phase
  at about \texorpdfstring{\ensuremath{-}}{-}100 degrees C as predicted from homologous
  compounds H\(_2\)S, H\(_2\)Se)
\item
  DNA double-helix stability: about 2 H-bonds per A--T base
  pair, about 3 H-bonds per G--C pair
\item
  Protein secondary structure: \texorpdfstring{\ensuremath{\alpha}}{alpha}-helices (intrachain H-bonds every 4
  residues), \texorpdfstring{\ensuremath{\beta}}{beta}-sheets (interchain H-bonds)
\item
  Cellulose microfibrils: interchain H-bonds between glucose polymers
  create immense tensile strength
\end{itemize}

\textbf{Geometry of hydrogen bonds:}

\begin{verbatim}
D–H ···· A
\end{verbatim}

Optimal: D--H\(\cdots\)A angle about 180 degrees; H\(\cdots\)A distance
about 1.8--2.0 angstroms; D\(\cdots\)A distance about 2.7--3.0
angstroms. Deviations from linearity weaken the bond approximately as
\(\cos^2\theta\), where \texorpdfstring{\ensuremath{\theta}}{theta} is the deviation from 180 degrees.

\textbf{Types of hydrogen bonds in biology:}

{\def\LTcaptype{none} % do not increment counter
\begin{longtable}[]{@{}llll@{}}
\toprule\noalign{}
Type & Donor & Acceptor & Example \\
\midrule\noalign{}
\endhead
\bottomrule\noalign{}
\endlastfoot
N--H\(\cdots\)O & Peptide N--H & Peptide C=O & \texorpdfstring{\ensuremath{\alpha}}{alpha}-helix backbone \\
O--H\(\cdots\)O & Water O--H & Carbonyl C=O & Hydration shells \\
N--H\(\cdots\)N & Adenine N--H & Thymine N & A--T base pair \\
O--H\(\cdots\)N & Serine O--H & Histidine N & Enzyme active sites \\
\end{longtable}
}

\subsubsection{Van der Waals Forces and Lipid
Packing}\label{van-der-waals-forces-and-lipid-packing}

Van der Waals forces arise from transient dipoles induced by quantum
fluctuations in electron clouds. They are extremely weak
(about 0.1--10 kJ/mol) but ubiquitous. Because they scale with
surface area, \textbf{nonpolar molecules with large surface areas} (such
as cholesterol or hydrophobic protein cores) benefit significantly from
van der Waals interactions.

There are three components of van der Waals forces:

\begin{enumerate}
\def\labelenumi{\arabic{enumi}.}
\tightlist
\item
  \textbf{Keesom forces} (permanent dipole--permanent dipole):
  orientation-dependent attraction between polar molecules.
\item
  \textbf{Debye forces} (permanent dipole--induced dipole): a polar
  molecule induces a dipole in a nearby nonpolar molecule.
\item
  \textbf{London dispersion forces} (induced dipole--induced dipole):
  the main attractive forces available between nonpolar molecules; arise
  from instantaneous fluctuations in electron density.
\end{enumerate}

The \textbf{van der Waals radius} defines the effective ``size'' of an
atom. When two non-bonded atoms approach closer than the sum of their
van der Waals radii, repulsive forces dominate (Pauli exclusion). This
steric effect constrains the conformational space of macromolecules, as
described by Ramachandran plots for proteins (see
\cref{sec:unit_I_macromolecules}).

\textbf{London dispersion forces and lipid packing.} The strength of
dispersion forces between two nonpolar surfaces grows roughly with the
polarisability and the contact area: longer alkyl chains have more
electrons whose instantaneous dipoles correlate more strongly. For two
parallel hydrocarbon chains, the dispersion energy of attraction is
approximately 2 kJ/mol per CH\(_2\) pair --- modest individually but
additive. A 16-carbon palmitic acid tail in close van der Waals contact
with a neighbouring tail in a phospholipid bilayer thus contributes
about 30 kJ/mol of attractive interaction per chain pair, which
(combined with the \hyperref[gl:hydrophobic-effect]{\textbf{hydrophobic
effect}}; see \cref{sec:unit_I_water_and_life}) holds the bilayer
together. Longer chains (24 carbons in sphingolipids) and saturated
chains (which can pack closely) give more rigid membranes; introducing a
single \emph{cis} double bond into oleic acid kinks the chain by
about 30 degrees, weakening dispersion contacts and raising
membrane fluidity --- the molecular reason vegetable oils are liquid
while butter is solid at room temperature.

These same forces explain why geckos can walk on walls: their toe pads
have millions of setae maximising van der Waals contact with smooth
surfaces, generating up to 10 N/cm\(^2\) of adhesive force.

\begin{center}\rule{0.5\linewidth}{0.5pt}\end{center}

\subsection{Functional Groups}\label{functional-groups}

Six major functional groups determine the chemical reactivity of organic
molecules:

\begin{figure}[htbp]
\centering
\includegraphics[width=0.82\linewidth,height=4.2in,keepaspectratio]{../figures/mermaid_inline/inline_mermaid_0016_ee38e756ca94.png}
\caption{Major functional groups. The six major functional groups of biological chemistry, their key properties, and representative biomolecules.}
\end{figure}

{\def\LTcaptype{none} % do not increment counter
\begin{longtable}[]{@{}
  >{\raggedright\arraybackslash}p{(\linewidth - 6\tabcolsep) * \real{0.3061}}
  >{\raggedright\arraybackslash}p{(\linewidth - 6\tabcolsep) * \real{0.1837}}
  >{\raggedright\arraybackslash}p{(\linewidth - 6\tabcolsep) * \real{0.2041}}
  >{\raggedright\arraybackslash}p{(\linewidth - 6\tabcolsep) * \real{0.3061}}@{}}
\toprule\noalign{}
\begin{minipage}[b]{\linewidth}\raggedright
Functional Group
\end{minipage} & \begin{minipage}[b]{\linewidth}\raggedright
Structure
\end{minipage} & \begin{minipage}[b]{\linewidth}\raggedright
Properties
\end{minipage} & \begin{minipage}[b]{\linewidth}\raggedright
Biological Role
\end{minipage} \\
\midrule\noalign{}
\endhead
\bottomrule\noalign{}
\endlastfoot
Hydroxyl & --OH & Polar, H-bond donor/acceptor & Alcohols, sugars,
serine \\
Carbonyl & C=O & Polar, H-bond acceptor & Ketones, aldehydes, amide
bonds \\
Carboxyl & --COOH & Polar, ionisable (pKaabout 4--5) & Amino
acids, fatty acids \\
Amino & --NH\(_2\) & Polar, weak base, H-bond donor & Amino acids,
\hyperref[gl:nucleotide]{\textbf{nucleotide}}s, ATP \\
Phosphate & --OPO\(_3^{2-}\) & Highly polar, charged at pH 7 & DNA/RNA
backbone, ATP, signalling \\
Sulfhydryl & --SH & Polar, oxidisable & Cysteine disulfide bridges \\
\end{longtable}
}

The \textbf{phosphate group} deserves special attention in biochemistry.
At physiological pH, phosphate groups carry two negative charges. In ATP
(adenosine triphosphate), the three sequential phosphate groups are
electrostatically repulsed, storing energy that is released upon
hydrolysis (about 30.5 kJ/mol under standard conditions;
about 50 kJ/mol under cellular conditions; see
\cref{sec:unit_III_bioenergetics_and_respiration}).

\textbf{Additional functional groups encountered in biochemistry:}

{\def\LTcaptype{none} % do not increment counter
\begin{longtable}[]{@{}
  >{\raggedright\arraybackslash}p{(\linewidth - 4\tabcolsep) * \real{0.4839}}
  >{\raggedright\arraybackslash}p{(\linewidth - 4\tabcolsep) * \real{0.2903}}
  >{\raggedright\arraybackslash}p{(\linewidth - 4\tabcolsep) * \real{0.2258}}@{}}
\toprule\noalign{}
\begin{minipage}[b]{\linewidth}\raggedright
Functional Group
\end{minipage} & \begin{minipage}[b]{\linewidth}\raggedright
Structure
\end{minipage} & \begin{minipage}[b]{\linewidth}\raggedright
Example
\end{minipage} \\
\midrule\noalign{}
\endhead
\bottomrule\noalign{}
\endlastfoot
Methyl & --CH\(_3\) & DNA methylation (epigenetic silencing) \\
Ester & --COO-- & Triacylglycerols, phospholipids \\
Ether & --O-- & Ether lipids (plasmalogens in myelin) \\
Amide & --CONH-- & Peptide bonds, asparagine side chain \\
Thioester & --COS-- & Acetyl-CoA (high-energy bond) \\
Imine (Schiff base) & C=N & Retinal in rhodopsin, PLP-dependent
enzymes \\
\end{longtable}
}

\begin{quote}
\textbf{Concept Check 3:} Acetyl-CoA contains a thioester bond (--COS--)
with a \(\Delta G^{\circ'}\) of hydrolysis of approximately --31.4
kJ/mol, comparable to ATP hydrolysis. Why might cells use thioester
bonds as ``activated'' carriers of acyl groups?
\end{quote}

\begin{center}\rule{0.5\linewidth}{0.5pt}\end{center}

\subsection{Resonance Structures and Delocalised
Electrons}\label{resonance-structures-and-delocalised-electrons}

Many biologically important molecules cannot be accurately described by
a single Lewis structure. \textbf{Resonance} occurs when electrons
(particularly \texorpdfstring{\ensuremath{\pi}}{pi} electrons and lone pairs) are delocalised across
multiple atoms.

\subsubsection{The Peptide Bond as a Resonance
Hybrid}\label{the-peptide-bond-as-a-resonance-hybrid}

The peptide bond (C--N) in proteins has significant double-bond
character because of resonance between two contributing structures:

\textbf{Structure I:} C=O with C--N single bond (nitrogen has a lone
pair)

\textbf{Structure II:} C--O\(^-\) with C=N\(^+\) double bond (lone pair
delocalised into \texorpdfstring{\ensuremath{\pi}}{pi} system)

The actual bond is a \textbf{resonance hybrid} --- the C--N bond length
(132 pm) is intermediate between a single C--N (147 pm) and a double C=N
(127 pm). This partial double-bond character enforces planarity and
restricts rotation around the peptide bond, which is fundamental to
protein backbone geometry.

\subsubsection{Resonance in Carboxylate
Ions}\label{resonance-in-carboxylate-ions}

When a carboxyl group (--COOH) loses its proton at physiological pH, the
resulting carboxylate (--COO\(^-\)) has two equivalent C--O bonds (bond
length 126 pm each) due to delocalisation of the negative charge over
both oxygen atoms. This resonance stabilisation explains why carboxylate
ions are weaker bases (pKa about 4--5) than alcohols (pKa
about 16).

\subsubsection{Aromatic Systems in
Biology}\label{aromatic-systems-in-biology}

The aromatic amino acids phenylalanine, tyrosine, and tryptophan contain
\textbf{aromatic rings} with fully delocalised \texorpdfstring{\ensuremath{\pi}}{pi} electrons. These
systems absorb UV light at characteristic wavelengths:

{\def\LTcaptype{none} % do not increment counter
\begin{longtable}[]{@{}
  >{\raggedright\arraybackslash}p{(\linewidth - 6\tabcolsep) * \real{0.1316}}
  >{\raggedright\arraybackslash}p{(\linewidth - 6\tabcolsep) * \real{0.2632}}
  >{\raggedright\arraybackslash}p{(\linewidth - 6\tabcolsep) * \real{0.4605}}
  >{\raggedright\arraybackslash}p{(\linewidth - 6\tabcolsep) * \real{0.1447}}@{}}
\toprule\noalign{}
\begin{minipage}[b]{\linewidth}\raggedright
Amino Acid
\end{minipage} & \begin{minipage}[b]{\linewidth}\raggedright
\(\lambda_{max}\) (nm)
\end{minipage} & \begin{minipage}[b]{\linewidth}\raggedright
\(\varepsilon\) (M\(^{-1}\) cm\(^{-1}\))
\end{minipage} & \begin{minipage}[b]{\linewidth}\raggedright
Application
\end{minipage} \\
\midrule\noalign{}
\endhead
\bottomrule\noalign{}
\endlastfoot
Tryptophan & 280 & 5,500 & Protein concentration (A280) \\
Tyrosine & 274 & 1,490 & Protein concentration \\
Phenylalanine & 257 & 200 & Minor contribution \\
\end{longtable}
}

The adenine and guanine bases in nucleic acids are also aromatic
heterocycles. Their \texorpdfstring{\ensuremath{\pi}}{pi}-stacking interactions contribute significantly to
DNA double-helix stability (see \cref{sec:unit_I_macromolecules}).

\begin{center}\rule{0.5\linewidth}{0.5pt}\end{center}

\subsection{Chirality and
Stereochemistry}\label{chirality-and-stereochemistry}

\subsubsection{Enantiomers and Biological
Specificity}\label{enantiomers-and-biological-specificity}

A \textbf{chiral} (asymmetric) carbon is bonded to four different
substituents. Such a carbon exists in two mirror-image forms called
\textbf{enantiomers} that cannot be superimposed. Life is remarkably
stereospecific:

\begin{itemize}
\tightlist
\item
  \textbf{L-amino acids} are used exclusively in proteins (D-amino acids
  are found primarily in bacterial cell walls and some antibiotics)
\item
  \textbf{D-sugars} predominate in metabolism (D-glucose, D-ribose)
\item
  Enzymes bind a single enantiomer --- the ``wrong'' one simply does not
  fit the active site
\end{itemize}

\subsubsection{R/S Configuration}\label{rs-configuration}

The Cahn-Ingold-Prelog (CIP) system assigns \textbf{R} (rectus) or
\textbf{S} (sinister) designations based on the priority of substituents
around the chiral centre. Priority rules:

\begin{enumerate}
\def\labelenumi{\arabic{enumi}.}
\tightlist
\item
  Higher atomic number = higher priority
\item
  If tied, move outward until a difference is found
\item
  View from opposite the lowest-priority group: clockwise = R,
  anticlockwise = S
\end{enumerate}

For \texorpdfstring{\ensuremath{\alpha}}{alpha}-amino acids, the L-configuration corresponds to S in most cases
(except cysteine, which is R due to sulfur's high atomic number).

\subsubsection{Geometric Isomerism in Fatty
Acids}\label{geometric-isomerism-in-fatty-acids}

Double bonds in unsaturated fatty acids adopt either \textbf{cis} (Z) or
\textbf{trans} (E) configurations:

\begin{itemize}
\tightlist
\item
  \textbf{cis} double bonds introduce a about 30-degree kink,
  disrupting membrane packing and increasing fluidity
\item
  \textbf{trans} double bonds (produced by industrial hydrogenation)
  behave like saturated chains, packing tightly
\end{itemize}

\begin{quote}
\textbf{Clinical Connection: Thalidomide and Chirality}

The drug thalidomide was prescribed as a racemic mixture (both R and S
forms) in the 1950s as a sedative for pregnant women. While the
R-enantiomer is an effective sedative, the S-enantiomer is teratogenic,
causing severe birth defects. This tragedy revolutionised pharmaceutical
regulation, leading to requirements for enantiomeric purity testing.
Tragically, even administering the pure R-form would not have prevented
harm, as thalidomide racemises \emph{in vivo}.
\end{quote}

\begin{quote}
\textbf{Concept Check 4:} Why can't enzymes catalyse reactions on both
enantiomers of a chiral substrate with equal efficiency? Consider the
three-dimensional geometry of the active site.
\end{quote}

\begin{center}\rule{0.5\linewidth}{0.5pt}\end{center}

\subsection{Redox Reactions in
Biology}\label{redox-reactions-in-biology}

\subsubsection{Oxidation and Reduction}\label{oxidation-and-reduction}

\textbf{Oxidation} is the loss of electrons; \textbf{reduction} is the
gain of electrons. The mnemonic ``OIL RIG'' (Oxidation Is Loss,
Reduction Is Gain) applies. In biology, redox reactions frequently
involve the transfer of hydrogen atoms (H\(^+\) + \(e^-\)) rather than
bare electrons.

\textbf{Oxidation states of carbon:}

\[\text{CH}_4 \xrightarrow{-2e^-} \text{CH}_3\text{OH} \xrightarrow{-2e^-} \text{HCHO} \xrightarrow{-2e^-} \text{HCOOH} \xrightarrow{-2e^-} \text{CO}_2 \tag{1.1} \label{eq:unit_I_atoms_molecules_item_1}\]

Each step represents a two-electron oxidation. Carbon in methane has
oxidation state --4; in CO\(_2\) it is +4. The complete oxidation of one
mole of glucose releases 2,870 kJ:

\[\text{C}_6\text{H}_{12}\text{O}_6 + 6\text{O}_2 \rightarrow 6\text{CO}_2 + 6\text{H}_2\text{O} \quad \Delta G^{\circ'} = -2,870 \; \text{kJ/mol} \tag{1.2} \label{eq:unit_I_atoms_molecules_item_2}\]

\subsubsection{Reduction Potentials}\label{reduction-potentials}

The \textbf{standard reduction potential} (\(E^{\circ'}\)) measures the
tendency of a half-reaction to gain electrons under standard biochemical
conditions (pH 7, 25 degrees C, 1 M concentrations):

{\def\LTcaptype{none} % do not increment counter
\begin{longtable}[]{@{}ll@{}}
\toprule\noalign{}
Half-Reaction & \(E^{\circ'}\) (V) \\
\midrule\noalign{}
\endhead
\bottomrule\noalign{}
\endlastfoot
O\(_2\) + 4H\(^+\) + 4\(e^-\) -\textgreater{} 2H\(_2\)O & +0.816 \\
Cytochrome c (Fe\(^{3+}\)) + \(e^-\) -\textgreater{} Cyt c (Fe\(^{2+}\))
& +0.254 \\
NAD\(^+\) + H\(^+\) + 2\(e^-\) -\textgreater{} NADH & --0.320 \\
2H\(^+\) + 2\(e^-\) -\textgreater{} H\(_2\) & --0.414 \\
\end{longtable}
}

The \textbf{free energy change} of a redox reaction is related to the
potential difference:

\[\Delta G^{\circ'} = -nF\Delta E^{\circ'} \tag{1.3} \label{eq:unit_I_atoms_molecules_item_3}\]

where \(n\) = number of electrons transferred and \(F\) = Faraday
constant (96,485 C/mol).

\subsubsection{Biological Electron
Carriers}\label{biological-electron-carriers}

The major electron carriers in metabolism are:

\begin{itemize}
\tightlist
\item
  \textbf{NAD\(^+\)/NADH:} carries two electrons as a hydride ion
  (H\(^-\)); central to catabolic pathways
\item
  \textbf{FAD/FADH\(_2\):} carries two electrons; tightly bound to
  enzymes (flavoproteins)
\item
  \textbf{NADP\(^+\)/NADPH:} carries two electrons; central to anabolic
  (biosynthetic) pathways
\item
  \textbf{Ubiquinone (CoQ)/Ubiquinol:} lipid-soluble electron carrier in
  the mitochondrial inner membrane
\item
  \textbf{Cytochromes:} iron-containing proteins that transfer single
  electrons via Fe\(^{2+}\)/Fe\(^{3+}\) cycling
\end{itemize}

\begin{quote}
\textbf{Concept Check 5:} In the electron transport chain, electrons
flow from NADH (\(E^{\circ'} = -0.320\) V) to O\(_2\)
(\(E^{\circ'} = +0.816\) V). Calculate \(\Delta G^{\circ'}\) for the
transfer of 2 electrons. Is this process thermodynamically favourable?
\end{quote}

\begin{center}\rule{0.5\linewidth}{0.5pt}\end{center}

\subsection{The Mole Concept and
Molarity}\label{the-mole-concept-and-molarity}

The \textbf{mole} (mol) is the SI unit for amount of substance, equal to
Avogadro's number (\(N_A\)):

\[N_A = 6.022 \times 10^{23} \; \text{mol}^{-1} \tag{1.4} \label{eq:unit_I_atoms_molecules_item_4}\]

\textbf{Molarity (M)} is defined as moles of solute per litre of
solution:

\[[C] = \frac{n}{V} \quad \text{(mol L}^{-1}\text{)} \tag{1.5} \label{eq:unit_I_atoms_molecules_item_5}\]

In biochemistry, cellular concentrations are typically in the millimolar
(mM) to micromolar (\texorpdfstring{\ensuremath{\mu}}{mu}M) range. ATP concentration in a liver cell is
approximately 3--5 mM. The dissociation constant \(K_d\) of a
high-affinity antibody for its antigen might be \(10^{-10}\) M (100 pM).

\textbf{Comparison of concentration scales used in biology:}

{\def\LTcaptype{none} % do not increment counter
\begin{longtable}[]{@{}
  >{\raggedright\arraybackslash}p{(\linewidth - 6\tabcolsep) * \real{0.1250}}
  >{\raggedright\arraybackslash}p{(\linewidth - 6\tabcolsep) * \real{0.3750}}
  >{\raggedright\arraybackslash}p{(\linewidth - 6\tabcolsep) * \real{0.1562}}
  >{\raggedright\arraybackslash}p{(\linewidth - 6\tabcolsep) * \real{0.3438}}@{}}
\toprule\noalign{}
\begin{minipage}[b]{\linewidth}\raggedright
Unit
\end{minipage} & \begin{minipage}[b]{\linewidth}\raggedright
Abbreviation
\end{minipage} & \begin{minipage}[b]{\linewidth}\raggedright
Value
\end{minipage} & \begin{minipage}[b]{\linewidth}\raggedright
Typical Use
\end{minipage} \\
\midrule\noalign{}
\endhead
\bottomrule\noalign{}
\endlastfoot
Molar & M & mol/L & Standard solutions \\
Millimolar & mM & \(10^{-3}\) M & Metabolite concentrations \\
Micromolar & \texorpdfstring{\ensuremath{\mu}}{mu}M & \(10^{-6}\) M & Enzyme concentrations \\
Nanomolar & nM & \(10^{-9}\) M & \hyperref[gl:hormone]{\textbf{Hormone}}
concentrations \\
Picomolar & pM & \(10^{-12}\) M & High-affinity ligand \(K_d\) \\
\end{longtable}
}

\begin{center}\rule{0.5\linewidth}{0.5pt}\end{center}

\subsection{pH and Buffers}\label{ph-and-buffers}

\subsubsection{The pH Scale}\label{the-ph-scale}

Water autoionises:

\[\text{H}_2\text{O} \rightleftharpoons \text{H}^+ + \text{OH}^- \tag{1.6} \label{eq:unit_I_atoms_molecules_item_6}\]

with the equilibrium constant at 25 degrees C:
\(K_w = [\text{H}^+][\text{OH}^-] = 10^{-14} \, \text{M}^2\).

The \textbf{pH} is defined as:

\[\text{pH} = -\log_{10}[\text{H}^+] \tag{1.7} \label{eq:unit_I_atoms_molecules_item_7}\]

At 37 degrees C (body temperature), \(K_w \approx 2.4 \times 10^{-14}\),
so neutral pH is about 6.81. The slight shift is physiologically
meaningful.

\textbf{Key physiological pH values:}

{\def\LTcaptype{none} % do not increment counter
\begin{longtable}[]{@{}
  >{\raggedright\arraybackslash}p{(\linewidth - 4\tabcolsep) * \real{0.4400}}
  >{\raggedright\arraybackslash}p{(\linewidth - 4\tabcolsep) * \real{0.0800}}
  >{\raggedright\arraybackslash}p{(\linewidth - 4\tabcolsep) * \real{0.4800}}@{}}
\toprule\noalign{}
\begin{minipage}[b]{\linewidth}\raggedright
Compartment
\end{minipage} & \begin{minipage}[b]{\linewidth}\raggedright
pH
\end{minipage} & \begin{minipage}[b]{\linewidth}\raggedright
Significance
\end{minipage} \\
\midrule\noalign{}
\endhead
\bottomrule\noalign{}
\endlastfoot
Blood plasma & 7.35--7.45 & Tight regulation; acidosis/alkalosis outside
this range \\
Lysosomes & 4.5--5.0 & Acid hydrolases active \\
Gastric juice & 1.5--3.5 & Pepsin active; kills pathogens \\
\hyperref[gl:cytoplasm]{\textbf{Cytoplasm}} (liver) & 7.2 & Most enzyme
optima \\
Mitochondrial matrix & 7.8--8.0 & pH gradient drives ATP synthesis \\
Pancreatic secretions & 8.0--8.5 & Neutralises gastric acid \\
\end{longtable}
}

\begin{quote}
\textbf{Clinical Connection: Diabetic Ketoacidosis}

In uncontrolled diabetes mellitus, insufficient insulin causes the body
to rely on fatty acid oxidation, producing excess acetoacetate and
\texorpdfstring{\ensuremath{\beta}}{beta}-hydroxybutyrate (ketone bodies). These are moderately strong acids
(pKa about 4.7), overwhelming the bicarbonate buffer system.
Blood pH can drop to 7.0--7.1, a life-threatening emergency called
\textbf{diabetic ketoacidosis (DKA)}. Treatment involves intravenous
insulin, fluid replacement, and careful potassium monitoring.
\end{quote}

\subsubsection{Henderson-Hasselbalch
Equation}\label{henderson-hasselbalch-equation}

For a weak acid HA with dissociation constant \(K_a\):

\begin{equation}
\text{pH} = \text{pK}_a + \log\frac{[\text{A}^-]}{[\text{HA}]}
\label{eq:unit_I_henderson_hasselbalch}
\end{equation}

A buffer resists pH change when {[}A\(^-\){]}/{[}HA{]} is near 1 (i.e.,
pH is about pK\(_a\)). The effective buffer range is typically
pK\(_a \pm\) 1. Outside this range, one component is depleted and buffer
capacity drops sharply.

\subsubsection{Buffer Capacity}\label{buffer-capacity}

The \textbf{buffer capacity} \texorpdfstring{\ensuremath{\beta}}{beta} measures how much strong acid or base a
buffer can absorb before pH shifts appreciably. Formally, \texorpdfstring{\ensuremath{\beta}}{beta} is the moles
of strong base added per litre per unit pH increase:

\begin{equation}
\beta = \frac{dC_b}{d\text{pH}} = 2.303\left([\text{H}^+] + [\text{OH}^-] + \frac{C_T \, K_a [\text{H}^+]}{(K_a + [\text{H}^+])^2}\right)
\label{eq:unit_I_buffer_capacity}
\end{equation}

where \(C_T = [\text{HA}] + [\text{A}^-]\) is the total buffer
concentration. The third term peaks when \(\mathrm{H}^+ = K_a\) (i.e.,
pH = pKa), where it reaches its maximum value of
\(\beta_{\max} = 2.303 \cdot C_T / 4 \approx 0.576\, C_T\). Two
practical consequences follow:

\begin{enumerate}
\def\labelenumi{\arabic{enumi}.}
\tightlist
\item
  \textbf{Buffer capacity scales linearly with total buffer
  concentration.} A 0.1 M phosphate buffer has 10\(\times\) the capacity
  of a 0.01 M phosphate buffer.
\item
  \textbf{Capacity is maximal at pH = pKa and falls off rapidly outside
  pKa \(\pm\) 1.} This is why the bicarbonate buffer (pKa = 6.1) is
  theoretically a poor blood buffer at pH 7.4 --- it is rescued
  primarily because the system is \emph{open} (CO\(_2\) is exhaled).
\end{enumerate}

\textbf{Worked buffer problem.} A biochemist needs to prepare 1.0 L of
0.1 M Tris buffer (pKa = 8.1) at pH 7.6 for an enzyme assay, starting
from solid Tris base (\(\text{A}^-\), MW 121) and 1 M HCl. How much of
each is required?

\emph{Solution:} From \cref{eq:unit_I_henderson_hasselbalch}:
\(7.6 = 8.1 + \log\frac{[\text{A}^-]}{[\text{HA}]}\), so
\(\log\frac{[\text{A}^-]}{[\text{HA}]} = -0.5\) and
\([\text{A}^-]/[\text{HA}] = 10^{-0.5} = 0.316\). With
\([\text{HA}] + [\text{A}^-] = 0.10\, \text{M}\):

\[[\text{A}^-] = 0.0240\,\text{M}, \quad [\text{HA}] = 0.0760\,\text{M}. \label{eq:unit_I_atoms_molecules_item_8}\]

Dissolve 0.10 mol \(\times\) 121 g/mol = \textbf{12.1 g of Tris base} in
about 800 mL water, then titrate with 1 M HCl until pH 7.6
(requiring about 76 mL HCl, since each mole of HCl converts one
mole of A\(^-\) to HA). Bring the volume to 1.0 L. Because pH is 0.5
units below pKa,
\(\beta \approx 0.45 \cdot 0.1 = 0.045\,\text{mol L}^{-1}\,\text{(pH unit)}^{-1}\)
--- adding 4.5 mmol of strong acid would lower the pH by
about 0.1 unit.

\begin{figure}[htbp]
\centering
\includegraphics[width=0.82\linewidth,height=4.2in,keepaspectratio]{../figures/mermaid_inline/inline_mermaid_0017_75546d74e4ec.png}
\caption{Bicarbonate buffer system. The bicarbonate buffer system. CO\$\_2\$ produced by metabolism combines with water (catalysed by carbonic anhydrase) to form carbonic acid, which dissociates to bicarbonate and H\$\textasciicircum{}+\$. The lungs and kidneys regulate the two ends of the equilibrium.}
\end{figure}

The \textbf{bicarbonate buffer system} in blood operates at pH 7.4 with:

\[\text{CO}_2 + \text{H}_2\text{O} \rightleftharpoons \text{H}_2\text{CO}_3 \rightleftharpoons \text{H}^+ + \text{HCO}_3^- \tag{1.9} \label{eq:unit_I_atoms_molecules_item_9}\]

where pKa\(_1\) = 6.1. Despite the pH/pKa offset of about 1.3
units, the bicarbonate system is effective because CO\(_2\) can be
rapidly exhaled to shift the equilibrium (an \textbf{open system}). The
normal ratio of {[}HCO\(_3^-\){]}/{[}CO\(_2\){]} is approximately 20:1.

The \textbf{phosphate buffer} (H\(_2\)PO\(_4^-\) / HPO\(_4^{2-}\); pKa =
6.8) buffers intracellular pH and is crucial in renal pH regulation.

\textbf{Protein buffers} --- particularly histidine residues (imidazole
pKa = 6.0) --- provide significant intracellular buffering. Haemoglobin
alone accounts for about 60\% of the buffering capacity of
blood because it contains 38 histidine residues per tetramer.

\subsubsection{Comparison of Biological pH Buffering
Systems}\label{comparison-of-biological-ph-buffering-systems}

{\def\LTcaptype{none} % do not increment counter
\begin{longtable}[]{@{}
  >{\raggedright\arraybackslash}p{(\linewidth - 10\tabcolsep) * \real{0.0769}}
  >{\raggedright\arraybackslash}p{(\linewidth - 10\tabcolsep) * \real{0.1795}}
  >{\raggedright\arraybackslash}p{(\linewidth - 10\tabcolsep) * \real{0.2308}}
  >{\raggedright\arraybackslash}p{(\linewidth - 10\tabcolsep) * \real{0.1667}}
  >{\raggedright\arraybackslash}p{(\linewidth - 10\tabcolsep) * \real{0.1410}}
  >{\raggedright\arraybackslash}p{(\linewidth - 10\tabcolsep) * \real{0.2051}}@{}}
\toprule\noalign{}
\begin{minipage}[b]{\linewidth}\raggedright
System
\end{minipage} & \begin{minipage}[b]{\linewidth}\raggedright
Conjugate pair
\end{minipage} & \begin{minipage}[b]{\linewidth}\raggedright
pKa (37 \(^\circ\)C)
\end{minipage} & \begin{minipage}[b]{\linewidth}\raggedright
Concentration
\end{minipage} & \begin{minipage}[b]{\linewidth}\raggedright
Compartment
\end{minipage} & \begin{minipage}[b]{\linewidth}\raggedright
Special features
\end{minipage} \\
\midrule\noalign{}
\endhead
\bottomrule\noalign{}
\endlastfoot
Bicarbonate & H\(_2\)CO\(_3\) / HCO\(_3^-\) & 6.10 (apparent) &
about 24 mM & Plasma, ECF & Open system: CO\(_2\) exhaled by
lungs \\
Phosphate & H\(_2\)PO\(_4^-\) / HPO\(_4^{2-}\) & 6.86 & 1--2 mM
(plasma); 100 mM (cells) & ICF, urine & Major intracellular buffer \\
Imidazole (His) & His-H\(^+\) / His & 6.00 & 38 sites/Hb tetramer & RBC
& Biggest blood buffer; Bohr effect \\
Ammonia & NH\(_4^+\) / NH\(_3\) & 9.25 & Variable & Renal tubule &
H\(^+\) excretion in acidosis \\
Carbonyl/amide & --COOH / --COO\(^-\) & 4--5 & Many sites & Proteins &
Side-chain buffering \\
Sulfonate & --SO\(_3\)H / --SO\(_3^-\) & \textless{} 1 &
Glycosaminoglycans & ECM & Typically ionised at cellular pH \\
\end{longtable}
}

The blood buffering hierarchy: \textbf{bicarbonate (about 75\%
of capacity) \textgreater{} haemoglobin (about 24\%)
\textgreater{} plasma proteins/phosphate (about 1\%)}. Despite
its inferior pKa, bicarbonate dominates because (i) it is at high
concentration and (ii) it is the main \emph{open} system, with CO\(_2\)
disposal by lungs (minutes timescale) and HCO\(_3^-\)
disposal/regeneration by kidneys (hours-to-days timescale). This
three-compartment design --- a chemistry buffer coupled to two
organ-system control loops --- is the body's solution to the problem
that any closed buffer will eventually be titrated to exhaustion.

\begin{quote}
\textbf{Concept Check 6:} A patient is hyperventilating after a panic
attack. Their pCO\(_2\) falls from 40 mmHg to 25 mmHg while
{[}HCO\(_3^-\){]} is briefly unchanged at 24 mM. Use
\cref{eq:unit_I_henderson_hasselbalch} to predict the resulting blood
pH. Will this be respiratory acidosis or alkalosis? Why does breathing
into a paper bag help?
\end{quote}

\begin{quote}
\textbf{Concept Check 7:} You are designing a buffer for an enzyme that
has its optimum at pH 6.5. Two candidates are available: 100 mM MES (pKa
6.15) and 100 mM HEPES (pKa 7.55). Using
\cref{eq:unit_I_buffer_capacity}, decide which provides greater buffer
capacity at pH 6.5 and explain your reasoning.
\end{quote}

\subsubsection{Titration Curves}\label{titration-curves}

A titration curve plots pH against the volume of added base (or acid).
For a monoprotic weak acid:

\begin{itemize}
\tightlist
\item
  At the \textbf{half-equivalence point}, pH = pKa (50\% ionised)
\item
  The \textbf{buffer region} extends from approximately pKa -- 1 to pKa
  + 1
\item
  At the \textbf{equivalence point}, most acid has been converted to
  conjugate base
\end{itemize}

Amino acids with ionisable side chains show \textbf{multiphasic
titration curves} with two or three pKa values:

{\def\LTcaptype{none} % do not increment counter
\begin{longtable}[]{@{}
  >{\raggedright\arraybackslash}p{(\linewidth - 8\tabcolsep) * \real{0.1250}}
  >{\raggedright\arraybackslash}p{(\linewidth - 8\tabcolsep) * \real{0.2750}}
  >{\raggedright\arraybackslash}p{(\linewidth - 8\tabcolsep) * \real{0.3250}}
  >{\raggedright\arraybackslash}p{(\linewidth - 8\tabcolsep) * \real{0.2500}}
  >{\raggedright\arraybackslash}p{(\linewidth - 8\tabcolsep) * \real{0.0250}}@{}}
\toprule\noalign{}
\begin{minipage}[b]{\linewidth}\raggedright
Amino Acid
\end{minipage} & \begin{minipage}[b]{\linewidth}\raggedright
pKa\(_1\) (\texorpdfstring{\ensuremath{\alpha}}{alpha}-COOH)
\end{minipage} & \begin{minipage}[b]{\linewidth}\raggedright
pKa\(_2\) (\texorpdfstring{\ensuremath{\alpha}}{alpha}-NH\(_3^+\))
\end{minipage} & \begin{minipage}[b]{\linewidth}\raggedright
pKa\(_R\) (side chain)
\end{minipage} & \begin{minipage}[b]{\linewidth}\raggedright
pI
\end{minipage} \\
\midrule\noalign{}
\endhead
\bottomrule\noalign{}
\endlastfoot
Glycine & 2.34 & 9.60 & --- & 5.97 \\
Aspartate & 2.09 & 9.82 & 3.86 (--COOH) & 2.77 \\
Histidine & 1.82 & 9.17 & 6.00 (imidazole) & 7.59 \\
Lysine & 2.18 & 8.95 & 10.53 (--NH\(_3^+\)) & 9.74 \\
\end{longtable}
}

The \textbf{isoelectric point} (pI) is the pH at which the amino acid
carries zero net charge. For amino acids without ionisable side chains,
pI = (pKa\(_1\) + pKa\(_2\))/2.

\begin{center}\rule{0.5\linewidth}{0.5pt}\end{center}

\subsection{Worked Examples}\label{worked-examples}

\textbf{Problem 1:} A solution contains 0.05 M acetic acid
(CH\(_3\)COOH, pKa = 4.76) and 0.15 M sodium acetate. What is the pH?

\emph{Solution:}
\[\text{pH} = 4.76 + \log\frac{0.15}{0.05} = 4.76 + \log 3 = 4.76 + 0.477 = 5.24 \tag{1.10} \label{eq:unit_I_atoms_molecules_item_10}\]

\textbf{Problem 2:} What fraction of histidine's imidazole group (pKa =
6.0) is protonated at pH 7.0?

\emph{Solution:}
\[\frac{[\text{HA}]}{[\text{HA}] + [\text{A}^-]} = \frac{1}{1 + 10^{\text{pH} - \text{pK}_a}} = \frac{1}{1 + 10^{7.0 - 6.0}} = \frac{1}{11} \approx 0.091 = 9.1\% \tag{1.11} \label{eq:unit_I_atoms_molecules_item_11}\]

\textbf{Problem 3:} Calculate the free energy released when 2 electrons
are transferred from NADH (\(E^{\circ'} = -0.320\) V) to O\(_2\)
(\(E^{\circ'} = +0.816\) V).

\emph{Solution:}
\[\Delta E^{\circ'} = E^{\circ'}_{\text{acceptor}} - E^{\circ'}_{\text{donor}} = +0.816 - (-0.320) = +1.136 \; \text{V} \tag{1.12} \label{eq:unit_I_atoms_molecules_item_12}\]

\[\Delta G^{\circ'} = -nF\Delta E^{\circ'} = -(2)(96{,}485)(1.136) = -219{,}213 \; \text{J/mol} \approx -219.2 \; \text{kJ/mol} \tag{1.13} \label{eq:unit_I_atoms_molecules_item_13}\]

This large negative \(\Delta G^{\circ'}\) drives
\hyperref[gl:oxidative-phosphorylation]{\textbf{oxidative
phosphorylation}}, ultimately producing about 2.5 ATP per NADH
(see \cref{sec:unit_III_bioenergetics_and_respiration}).

\textbf{Problem 4:} A patient's arterial blood shows {[}HCO\(_3^-\){]} =
12 mM and pCO\(_2\) = 40 mmHg. Using the Henderson-Hasselbalch equation
with {[}CO\(_2\){]} = 0.03 \(\times\) pCO\(_2\) mM and pKa = 6.1,
calculate the blood pH. Is this acidosis or alkalosis?

\emph{Solution:}
\[[\text{CO}_2] = 0.03 \times 40 = 1.2 \; \text{mM} \tag{1.14} \label{eq:unit_I_atoms_molecules_item_14}\]

\[\text{pH} = 6.1 + \log\frac{12}{1.2} = 6.1 + \log 10 = 6.1 + 1.0 = 7.1 \tag{1.15} \label{eq:unit_I_atoms_molecules_item_15}\]

The pH of 7.1 is below the normal range of 7.35--7.45, indicating
\textbf{metabolic acidosis} (low bicarbonate with normal pCO\(_2\)).
Normal {[}HCO\(_3^-\){]} is 24 mM; this patient's value is halved.

\textbf{Problem 5:} The Ka of lactic acid is \(1.38 \times 10^{-4}\).
Calculate the pKa. If a muscle cell produces lactic acid during intense
exercise, what percentage is ionised at intracellular pH 6.8?

\emph{Solution:}
\[\text{pK}_a = -\log(1.38 \times 10^{-4}) = 3.86 \tag{1.16} \label{eq:unit_I_atoms_molecules_item_16}\]

\[\frac{[\text{A}^-]}{[\text{HA}]} = 10^{\text{pH} - \text{pK}_a} = 10^{6.8 - 3.86} = 10^{2.94} = 871 \tag{1.17} \label{eq:unit_I_atoms_molecules_item_17}\]

\[\text{Fraction ionised} = \frac{871}{1 + 871} = 0.9989 = 99.9\% \tag{1.18} \label{eq:unit_I_atoms_molecules_item_18}\]

At physiological pH, the Henderson-Hasselbalch relationship predicts
that the ionised form (lactate\(^-\)) strongly dominates over protonated
lactic acid. This is why we typically refer to ``lactate'' rather than
``lactic acid'' in biological contexts.

\begin{center}\rule{0.5\linewidth}{0.5pt}\end{center}

\subsection{Key Comparison: Bond Types in
Biology}\label{key-comparison-bond-types-in-biology}

{\def\LTcaptype{none} % do not increment counter
\begin{longtable}[]{@{}
  >{\raggedright\arraybackslash}p{(\linewidth - 8\tabcolsep) * \real{0.1702}}
  >{\raggedright\arraybackslash}p{(\linewidth - 8\tabcolsep) * \real{0.1702}}
  >{\raggedright\arraybackslash}p{(\linewidth - 8\tabcolsep) * \real{0.1064}}
  >{\raggedright\arraybackslash}p{(\linewidth - 8\tabcolsep) * \real{0.2766}}
  >{\raggedright\arraybackslash}p{(\linewidth - 8\tabcolsep) * \real{0.2766}}@{}}
\toprule\noalign{}
\begin{minipage}[b]{\linewidth}\raggedright
Property
\end{minipage} & \begin{minipage}[b]{\linewidth}\raggedright
Covalent
\end{minipage} & \begin{minipage}[b]{\linewidth}\raggedright
Ionic
\end{minipage} & \begin{minipage}[b]{\linewidth}\raggedright
Hydrogen Bond
\end{minipage} & \begin{minipage}[b]{\linewidth}\raggedright
Van der Waals
\end{minipage} \\
\midrule\noalign{}
\endhead
\bottomrule\noalign{}
\endlastfoot
\textbf{Strength (kJ/mol)} & 150--950 & 20--200 (in solution) & 2--40 &
0.1--10 \\
\textbf{Distance} & 0.1--0.2 nm & 0.2--0.3 nm & 0.2--0.3 nm & 0.3--0.5
nm \\
\textbf{Directionality} & Highly directional & Non-directional &
Directional (about 180 degrees) & Non-directional \\
\textbf{In water} & Stable & Often disrupted & Dynamic & Very weak \\
\textbf{Biological role} & Molecular skeleton & Salt bridges,
electrolytes & Structure, recognition & Packing, complementarity \\
\textbf{Thermal disruption at 37 degrees C} & No & Sometimes &
Constantly breaking/reforming & Constantly fluctuating \\
\end{longtable}
}

\begin{center}\rule{0.5\linewidth}{0.5pt}\end{center}

\subsection{Computational Bridge}\label{computational-bridge}

Cellular \hyperref[gl:thermodynamics]{\textbf{thermodynamics}} uses the
same \(\Delta G = \Delta G^\circ + RT\ln Q\) form as introductory
chemistry. The project implements the concentration correction
explicitly:

\begin{Shaded}
\begin{Highlighting}[]
\ImportTok{from}\NormalTok{ biology.biochemistry }\ImportTok{import}\NormalTok{ reaction\_free\_energy}

\CommentTok{\# ATP hydrolysis: ΔG°\textquotesingle{} ≈ {-}30.5 kJ/mol; equal 1 mM pools → Q = 1, so ΔG = ΔG°\textquotesingle{}}
\NormalTok{dg }\OperatorTok{=}\NormalTok{ reaction\_free\_energy(}\OperatorTok{{-}}\FloatTok{30.5}\NormalTok{, product\_conc}\OperatorTok{=}\FloatTok{1e{-}3}\NormalTok{, reactant\_conc}\OperatorTok{=}\FloatTok{1e{-}3}\NormalTok{)}
\BuiltInTok{print}\NormalTok{(}\BuiltInTok{round}\NormalTok{(dg, }\DecValTok{2}\NormalTok{))  }\CommentTok{\# {-}30.5}
\end{Highlighting}
\end{Shaded}

\begin{quote}
\textbf{Clinical / systems note:} Arterial blood gas interpretation in
critical care applies the same logarithmic definition of pH and buffer
ratios you use with the Henderson--Hasselbalch equation; the
machine-reported bicarbonate and CO\(_2\) tension are coupled through an
open buffer system analogous to the models in this chapter.
\end{quote}

\begin{center}\rule{0.5\linewidth}{0.5pt}\end{center}

\subsection{Current Evidence and Frontier
Biology}\label{current-evidence-and-frontier-biology-3}

For \textbf{Atoms, Molecules, and Chemical Bonds}, frontier biology
belongs inside the evidence logic of the chapter. Chemistry-of-life
claims now connect classical bonding and thermodynamics with AI-guided
structure prediction and experimental validation. The core reading
question is this: molecular claims need charge, polarity, geometry,
concentration, and solvent context.

\begin{itemize}
\tightlist
\item
  \textbf{What to verify:} identify the observation, model, assay, or
  dataset that would make the claim stronger or weaker.
\item
  \textbf{What to qualify:} state the scale, organism, cell type,
  environmental condition, or population where the claim is expected to
  hold.
\item
  \textbf{What to compare:} test at least one alternative explanation,
  baseline, or null model before treating the pattern as causal.
\item
  \textbf{What to cite:} distinguish primary evidence, review synthesis,
  public dataset, and institutional guidance; for recent or numeric
  claims, prefer the source closest to the measurement and state what
  has changed since it was published.
\end{itemize}

Use AI biomolecular models as hypothesis generators: compare confidence,
conservation, solvent exposure, and assay evidence before turning a
predicted contact into a biological claim
\citep{abramson2024alphafold3}.

\textbf{Source practice:} For structure and interaction claims, cite
experimental structures when available and treat AlphaFold 3 or AFDB
complex predictions as hypotheses to validate with confidence metrics,
conservation, mutagenesis, binding, or cryo-EM/X-ray/NMR evidence
\citep{abramson2024alphafold3,velankar2026alphafolddb2025,emblebi2026alphafoldcomplexes}.

\subsection{Summary}\label{summary-3}

\begin{itemize}
\tightlist
\item
  Atoms consist of protons (nuclear charge), neutrons (nuclear mass),
  and electrons (chemical reactivity).
\item
  Electronegativity differences determine bond polarity: life exploits
  polar O--H and N--H bonds for hydrogen bonding.
\item
  Carbon's tetravalency and capacity for \(sp^3\)/\(sp^2\)/\(sp\)
  hybridisation is the structural basis of most organic molecules.
\item
  Resonance delocalisation stabilises peptide bonds, carboxylate ions,
  and aromatic rings.
\item
  Chirality (L-amino acids, D-sugars) is a hallmark of biological
  specificity.
\item
  Redox reactions (electron transfer) drive energy metabolism;
  \(\Delta G^{\circ'} = -nF\Delta E^{\circ'}\).
\item
  Isotopes serve as tracers (metabolic studies) and therapeutic agents
  (radiotherapy, PET imaging).
\item
  pH = \(-\log[\text{H}^+]\); buffers maintain physiological pH by
  resisting acid/base addition through weak acid/conjugate base pairs.
\item
  The bicarbonate buffer system is an open system regulated by lungs
  (CO\(_2\)) and kidneys (HCO\(_3^-\)).
\item
  \textbf{Connections:} See Unit III (bioenergetics) for how exergonic
  hydrolysis of ATP couples to cellular work; Unit IX (physiology) links
  blood gas chemistry to ventilation and renal compensation.
\end{itemize}

\subsection{Key Terms}\label{key-terms-3}

\begin{itemize}
\tightlist
\item
  \textbf{Atomic number (Z):} Number of protons; defines the element
\item
  \textbf{Isotope:} Atoms of the same element with different neutron
  counts
\item
  \textbf{Electronegativity:} Tendency of an atom to attract bonding
  electrons
\item
  \textbf{Polar covalent bond:} Unequal electron sharing producing
  partial charges
\item
  \textbf{Hydrogen bond:} Weak bond between H bonded to N/O/F and
  another electronegative atom
\item
  \textbf{Van der Waals force:} Weak attractions from transient or
  induced dipoles
\item
  \textbf{Functional group:} Characteristic chemical group determining
  reactivity
\item
  \textbf{Resonance:} Delocalisation of electrons across multiple atoms
\item
  \textbf{Chirality:} Property of a molecule with a non-superimposable
  mirror image
\item
  \textbf{Enantiomer:} One of two mirror-image forms of a chiral
  molecule
\item
  \textbf{Oxidation:} Loss of electrons
\item
  \textbf{Reduction:} Gain of electrons
\item
  \textbf{Reduction potential (\(E^{\circ'}\)):} Tendency of a
  half-reaction to gain electrons
\item
  \textbf{Molarity:} Concentration in moles per litre
\item
  \textbf{pH:} Negative log of hydrogen ion concentration
\item
  \textbf{Buffer:} Solution that resists pH change
\item
  \textbf{Henderson-Hasselbalch equation:} pH = pKa +
  log({[}A\(^-\){]}/{[}HA{]})
\item
  \textbf{pI (isoelectric point):} pH at which a molecule has zero net
  charge
\end{itemize}

\subsection{Review Questions}\label{review-questions-3}

\begin{enumerate}
\def\labelenumi{\arabic{enumi}.}
\tightlist
\item
  Draw electron-dot structures for CH\(_4\), NH\(_3\), and H\(_2\)O.
  Predict the bond angles using VSEPR theory and explain any deviations
  from ideal tetrahedral geometry.
\item
  Why does the boiling point of water (100 degrees C) far exceed that of
  H\(_2\)S (--60 degrees C), despite both having similar molecular
  masses?
\item
  A patient's blood pH is measured at 7.25. Is this acidosis or
  alkalosis? Using the Henderson-Hasselbalch equation, calculate the
  {[}HCO\(_3^-\){]}/{[}CO\(_2\){]} ratio this implies.
\item
  Explain why phosphate groups in DNA confer structural stability while
  also making DNA polyanionic.
\item
  Carbon-11 (\(^{11}\)C) has a half-life of 20.4 minutes. Why is this
  isotope useful for PET imaging of brain metabolism but impractical for
  studies lasting several hours?
\item
  Explain why the peptide bond is planar. Include the concept of
  resonance in your answer, and describe how this planarity constrains
  protein structure.
\item
  A pharmaceutical company synthesises a chiral drug as a racemic
  mixture. Explain why both enantiomers might not have the same
  therapeutic effect, and give a historical example.
\item
  Using the \(\Delta G^{\circ'} = -nF\Delta E^{\circ'}\) relationship,
  predict whether the following reaction is spontaneous: cytochrome c
  (reduced) donating an electron to O\(_2\). Refer to the table of
  standard reduction potentials.
\item
  Calculate the pI of aspartate (pKa\(_1\) = 2.09, pKa\(_R\) = 3.86,
  pKa\(_2\) = 9.82). At pH 7.4, what is the net charge on aspartate?
\item
  A buffer is prepared with 0.1 M NaH\(_2\)PO\(_4\) and 0.1 M
  Na\(_2\)HPO\(_4\) (pKa = 6.8). Calculate the pH. If 10 mL of 1 M HCl
  is added to 1 L of this buffer, what is the new pH?
\item
  Using \breaktt{reaction\_free\_energy} with \texorpdfstring{\ensuremath{\Delta}}{Delta}G°' = -30.5 kJ/mol,
  compute \texorpdfstring{\ensuremath{\Delta}}{Delta}G when {[}ADP{]}{[}Pi{]}/{[}ATP{]} = 0.1 vs.~2.0 (take
  concentrations in a consistent molar basis). Which condition better
  matches a working muscle fibre?
\item
  Compare how a \textbf{closed} phosphate buffer in a test tube and the
  \textbf{open} HCO\(_3^-\)/CO\(_2\) system in blood respond to a bolus
  of strong acid. Which system restores pH faster in vivo, and which
  organ systems do the work?
\item
  Using the formal-charge formula \cref{eq:unit_I_formal_charge}, draw a
  Lewis structure for the carboxylate ion --COO\(^-\) and assign the
  formal charge on each atom. Explain how resonance equalises the two
  C--O bond lengths.
\item
  Sort C--C, C=C, O--H, N--H, and C--H bonds by polarity using the
  Pauling electronegativity table. For each bond, predict whether the H
  atom (where present) can serve as a hydrogen bond donor.
\item
  Use \cref{eq:unit_I_buffer_capacity} to compare the buffer capacity at
  pH 7.4 of (a) 25 mM bicarbonate (pKa 6.1) and (b) 25 mM phosphate (pKa
  6.8). Why does the actual buffering importance of bicarbonate in blood
  far exceed what this calculation predicts?
\end{enumerate}

\begin{center}\rule{0.5\linewidth}{0.5pt}\end{center}

\subsection{Further Reading and Source
Notes}\label{further-reading-and-source-notes-3}

\begin{itemize}
\tightlist
\item
  Henderson (1913). \emph{The Fitness of the Environment}. Macmillan.
\item
  Pauling (1939). \emph{The Nature of the Chemical Bond}. Cornell
  University Press.
\item
  Lewis (1916). The atom and the molecule. \emph{Journal of the American
  Chemical Society}, 38.
\item
  Atkins \& de Paula (latest ed.). \emph{Atkins' Physical Chemistry}
  (chapters on chemical bonding and molecular structure). Oxford
  University Press.
\item
  Sanderson (1976). \emph{Chemical Bonds and Bond Energy}. Academic
  Press.
\item
  Schrödinger (1944). \emph{What Is Life?} Cambridge University Press.
\end{itemize}

\begin{center}\rule{0.5\linewidth}{0.5pt}\end{center}

\subsubsection{Companion Source Module}\label{companion-source-module-3}

\textbf{Atoms, Molecules, and Chemical Bonds} should leave a
reproducible trail from a biological claim to the code, figure, diagram,
or paper-based activity that can test it. Use the surfaces below to
inspect the chapter's assumptions, rerun the relevant model, or compare
the manuscript explanation with companion labs and figures.

{\def\LTcaptype{none} % do not increment counter
\begin{longtable}[]{@{}
  >{\raggedright\arraybackslash}p{(\linewidth - 2\tabcolsep) * \real{0.5000}}
  >{\raggedright\arraybackslash}p{(\linewidth - 2\tabcolsep) * \real{0.5000}}@{}}
\toprule\noalign{}
\begin{minipage}[b]{\linewidth}\raggedright
Surface
\end{minipage} & \begin{minipage}[b]{\linewidth}\raggedright
Use it for
\end{minipage} \\
\midrule\noalign{}
\endhead
\bottomrule\noalign{}
\endlastfoot
\breaktt{src/biology/biochemistry/biochemistry.py}
(\breaktt{reaction\_free\_energy}, \texttt{atp\_free\_energy}) & Tie
bonding and reaction direction to energy accounting. \\
\breaktt{src/biology/cell/cell\_biology.py} (\texttt{osmotic\_pressure},
\texttt{diffusion\_flux}) & Connect charge, solubility, and
concentration gradients to cell-scale outcomes. \\
\breaktt{src/mermaid/biology\_diagrams.py}
(\breaktt{macromolecule\_classification\_diagram}) & Place atomic and
bond-level concepts inside the larger biomolecule map. \\
\end{longtable}
}

\textbf{Reproducibility check:} for every molecular claim, write the
charge/polarity, solvent context, and unit-bearing quantity that would
make the claim testable. \textbf{Cross-reference:} extend the same logic
in \cref{sec:unit_I_water_and_life}, \cref{sec:unit_I_macromolecules},
and \cref{sec:unit_II_membrane_transport}.

\newpage

\section{Water --- The Molecule of
Life}\label{water-the-molecule-of-life}

\label{sec:unit_I_water_and_life}

\begin{quote}
\textbf{Ch 2} \texorpdfstring{\ensuremath{\cdot}}{.} Level 1/3 \texorpdfstring{\ensuremath{\cdot}}{.} 40 min read \texorpdfstring{\ensuremath{\cdot}}{.} 50 min lecture \texorpdfstring{\ensuremath{\cdot}}{.}
Prerequisites: \cref{sec:unit_I_atoms_molecules}
\end{quote}

\subsection{Learning Objectives}\label{learning-objectives-4}

\begin{enumerate}
\def\labelenumi{\arabic{enumi}.}
\tightlist
\item
  Explain how water's polarity generates its unusual physical
  properties.
\item
  Describe hydrophilic and hydrophobic interactions and their importance
  in cell architecture.
\item
  Quantify osmotic pressure using the van 't Hoff equation and derive it
  from first principles.
\item
  Apply Fick's First Law to describe diffusion across membranes with
  multiple biological applications.
\item
  Relate water's properties to
  \hyperref[gl:homeostasis]{\textbf{homeostasis}} and adaptation in
  living organisms.
\item
  Describe the structure and function of aquaporins and their clinical
  significance.
\item
  Explain colligative properties and their relevance to antifreeze
  organisms and cryobiology.
\item
  Define water activity and explain its importance for food science and
  microbiology.
\end{enumerate}

\subsubsection{Study Blueprint}\label{study-blueprint-4}

\begin{itemize}
\tightlist
\item
  \textbf{Big idea:} Water's polarity, hydrogen bonding, and ionization
  make cells physically possible.
\item
  \textbf{Core concepts:} hydrogen bonding, cohesion, pH, buffers.
\item
  \textbf{Framework alignment:} Vision \& Change: Structure and
  function, Pathways and transformations of energy and matter; AP
  Biology: Energetics, Systems Interactions; NGSS-style topics: Matter
  and Energy in Organisms and Ecosystems, Structure and Function.
\item
  \textbf{Model or quantitative lens:} Henderson-Hasselbalch and
  pH-scale calculations.
\item
  \textbf{Data skill:} Convert between pH, hydrogen ion concentration,
  and buffer ratios.
\item
  \textbf{Practice cadence:} Concept Explanation, Statistical Tests and
  Data Analysis, Argumentation.
\item
  \textbf{Common misconception to repair:} Water is not an inert
  background; it is an active participant in structure and reaction
  chemistry.
\item
  \textbf{Primary lab:} \cref{sec:lab_unit_I_water_and_life}.
\item
  \textbf{Question bank:} \cref{sec:q_unit_I_water_and_life}.
\item
  \textbf{Transfer task:} Transfer water-property reasoning to blood
  buffering, plant transport, or protein folding.
\item
  \textbf{Bridge to computation:}
  \breaktt{biology.biochemistry.biochemistry.atp\_free\_energy}.
\end{itemize}

\begin{center}\rule{0.5\linewidth}{0.5pt}\end{center}

\begin{quote}
\textbf{Opening Vignette: Life at the Boiling Point --- and Below Zero}

Water boils at 100 °C and freezes at 0 °C. By every prediction from
homologous compounds --- hydrogen sulfide (H\texorpdfstring{\textsubscript{2}}{2}S, boiling point \texorpdfstring{\ensuremath{-}}{-}60 °C),
hydrogen selenide (H\texorpdfstring{\textsubscript{2}}{2}Se, bp \texorpdfstring{\ensuremath{-}}{-}41 °C) --- water ``should'' be a gas at
room temperature. It is not, and that deviation is why you are alive.
The anomalous boiling point of water arises entirely from
\hyperref[gl:hydrogen-bond]{\textbf{hydrogen bond}}s: each water
molecule can donate two and accept two H-bonds, creating a dynamic,
three-dimensional network that requires exceptional energy to disrupt.
That high specific heat capacity (4.18 J/g\texorpdfstring{\ensuremath{\cdot}}{.}°C)
\hyperref[gl:buffer]{\textbf{buffer}}s organisms against temperature
swings. The high heat of vaporisation (2,260 J/g) makes sweating an
efficient coolant --- a marathon runner loses about 1.5 L/hour
but keeps their core temperature within 1 °C.

In 2012, scientists drilling into Lake Vostok --- buried 3.7 km beneath
Antarctic ice under 350 atmospheres of pressure --- recovered microbial
DNA from supercooled water isolated from the atmosphere for 15 million
years (Rogers et al., 2013, \emph{PLOS ONE}). Life persisted in
near-freezing darkness, sustained by water's extraordinary properties.
Everything in this chapter explains why.

\emph{Primary source: Rogers, S. O. et al.~(2013). Physiology and
phylogeny of microorganisms from Lake Vostok. PLOS ONE, 8(2), e56136.}
\end{quote}

\begin{center}\rule{0.5\linewidth}{0.5pt}\end{center}

Among known molecules on Earth, liquid water is unusually well suited as
the solvent of life. It is the most abundant molecule in living cells
(about 70\% of cell mass), and metabolism depends on water as
solvent, reactant, product, or thermal buffer. The following properties
--- each a direct consequence of its molecular structure --- make water
biologically hard to substitute.

Water is so central to biology that astrobiologists use the mantra
``follow the water'' when searching for extraterrestrial life. Every
known form of life requires liquid water, and the ``habitable zone''
around a star is defined primarily by the temperature range permitting
liquid water on a planet's surface.

\begin{center}\rule{0.5\linewidth}{0.5pt}\end{center}

\subsection{Molecular Structure and
Polarity}\label{molecular-structure-and-polarity}

The water molecule consists of one oxygen atom covalently bonded to two
hydrogen atoms. The O--H bond length is 0.096 nm; the H--O--H bond angle
is \textbf{104.5 degrees} (less than the tetrahedral
\citep{henderson1913} 109.5 degrees because the two lone pairs compress
the bonding pairs, as predicted by VSEPR theory).

Oxygen's high
\hyperref[gl:electronegativity]{\textbf{electronegativity}} (3.44 vs.~H
= 2.20) generates \textbf{partial charges}: \(\delta^-\) on oxygen and
\(\delta^+\) on each hydrogen. The resulting dipole moment (\texorpdfstring{\ensuremath{\mu}}{mu} = 1.85
Debye) is one of the largest for any small molecule.

Each water molecule can form \textbf{up to four hydrogen bonds}: two as
a hydrogen-bond donor (through its two O--H hydrogens) and two as a
hydrogen-bond acceptor (through its two lone pairs). This tetrahedral
hydrogen bonding network gives rise to many of water's exceptional
properties.

\begin{figure}[htbp]
\centering
\includegraphics[width=0.82\linewidth,height=4.2in,keepaspectratio]{../figures/mermaid_inline/inline_mermaid_0018_804368712927.png}
\caption{Water structure and hydrogen bonding. Water's molecular structure gives rise to a tetrahedral hydrogen-bonding network, which in turn produces five extraordinary physical properties. Each property has direct biological consequences.}
\end{figure}

\subsubsection{The Hydrogen Bond Network in
Detail}\label{the-hydrogen-bond-network-in-detail}

In liquid water at 25 degrees C, each molecule forms an average of
about 3.4 hydrogen bonds at any instant (out of a maximum 4).
These bonds are dynamic, with a lifetime of approximately 1--20
picoseconds. The entire hydrogen bond network rearranges on a timescale
of about 1 ps --- meaning water is simultaneously highly
structured and highly dynamic.

In ice, the hydrogen bond network is fully realised: every molecule
forms exactly 4 hydrogen bonds in a perfect tetrahedral geometry,
creating a hexagonal lattice. This explains both the lower density of
ice and the hexagonal symmetry of snowflakes.

\textbf{Quantitative dynamics of the H-bond network.} Modern femtosecond
infrared spectroscopy has measured several timescales that govern
water's behaviour:

{\def\LTcaptype{none} % do not increment counter
\begin{longtable}[]{@{}
  >{\raggedright\arraybackslash}p{(\linewidth - 4\tabcolsep) * \real{0.1944}}
  >{\raggedright\arraybackslash}p{(\linewidth - 4\tabcolsep) * \real{0.2500}}
  >{\raggedright\arraybackslash}p{(\linewidth - 4\tabcolsep) * \real{0.5556}}@{}}
\toprule\noalign{}
\begin{minipage}[b]{\linewidth}\raggedright
Process
\end{minipage} & \begin{minipage}[b]{\linewidth}\raggedright
Timescale
\end{minipage} & \begin{minipage}[b]{\linewidth}\raggedright
Biological relevance
\end{minipage} \\
\midrule\noalign{}
\endhead
\bottomrule\noalign{}
\endlastfoot
O--H stretch vibration & about 10 fs & Vibrational coupling of
H-bond network \\
Single H-bond lifetime (exchange) & 1--3 ps & Faster than enzyme
turnover; water ``lubricates'' catalysis \\
Translational diffusion (1 nm) & about 50 ps & Sets viscous
drag on macromolecules \\
Hydration shell residence time at protein surface & 10 ps -- 1 ns &
Couples to protein backbone fluctuations \\
Tightly bound waters (buried in protein) & \texorpdfstring{\ensuremath{\mu}}{mu}s -- ms & Visible in X-ray
structures; contribute to specificity \\
Bulk-water dielectric relaxation & about 8 ps & Determines
\(\varepsilon\) at MHz/GHz \\
\end{longtable}
}

Each water molecule undergoes about 10\(^{12}\) H-bond
rearrangements per second --- an extraordinary churn that nevertheless
preserves the average network structure. Biological molecules
effectively ``swim'' through this fluctuating cage; their conformational
dynamics are intimately coupled to the surrounding water (the
``slaving'' model of protein dynamics).

\begin{quote}
\textbf{Concept Check 1:} If each water molecule in liquid water forms
an average of 3.4 hydrogen bonds and each hydrogen bond is shared
between two molecules, how many hydrogen bonds exist per mole of liquid
water?
\end{quote}

\begin{center}\rule{0.5\linewidth}{0.5pt}\end{center}

\subsection{Physical Properties Explained by Hydrogen
Bonding}\label{physical-properties-explained-by-hydrogen-bonding}

\subsubsection{\texorpdfstring{High Heat Capacity (\(C_p\) = 4.18 J
g\(^{-1}\)
K\(^{-1}\))}{High Heat Capacity (C_p = 4.18 J g^(-1) K^(-1))}}\label{high-heat-capacity-c_p-4.18-j-g-1-k-1}

Heating water requires breaking hydrogen bonds before kinetic energy of
\hyperref[gl:translation]{\textbf{translation}}/rotation can increase.
This buffers temperature fluctuations in organisms --- an important
thermal homeostasis tool. Cells with high water content resist
temperature changes.

\emph{Comparison:} Ethanol has \(C_p\) = 2.44 J g\(^{-1}\) K\(^{-1}\)
(barely half of water's) despite similar molecular weight, because it
has fewer H-bonds per molecule.

The high heat capacity of water has global implications: coastal
climates are moderated by the ocean's thermal inertia. The Gulf Stream
carries about 1.4 petawatts (\(1.4 \times 10^{15}\) W) of heat
from the tropics to northern Europe, keeping Britain 5--10 degrees C
warmer than equivalent latitudes in Canada.

\textbf{Worked example:} How much heat is required to raise the
temperature of 70 kg of body water (a typical human) by 1 degree C?

\[Q = mC_p\Delta T = 70{,}000 \;\text{g} \times 4.18 \;\text{J g}^{-1}\text{K}^{-1} \times 1 \;\text{K} = 292{,}600 \;\text{J} \approx 293 \;\text{kJ} \tag{2.1} \label{eq:unit_I_water_and_life_item_1}\]

This is equivalent to the energy in approximately 70 kcal of food --- a
substantial metabolic investment. This explains why fever is
energetically costly: raising body temperature by even 2 degrees C
requires about 586 kJ, roughly 20\% of resting metabolic output
for one hour.

\subsubsection{\texorpdfstring{High Heat of Vaporisation
(\(\Delta_{\text{vap}}H\) = 44 kJ/mol at 37 degrees
C)}{High Heat of Vaporisation (\textbackslash Delta\_\{\textbackslash text\{vap\}\}H = 44 kJ/mol at 37 degrees C)}}\label{high-heat-of-vaporisation-delta_textvaph-44-kjmol-at-37-degrees-c}

Evaporating water requires breaking the hydrogen bond network. This
underlies sweating as a thermoregulatory strategy: evaporating 1 g of
sweat removes about 2.4 kJ of heat from the body. At metabolic
rates of about 80 W, a resting person must evaporate roughly
120 mL/h at 37 degrees C to maintain steady body temperature.

\textbf{Comparison with other solvents:}

{\def\LTcaptype{none} % do not increment counter
\begin{longtable}[]{@{}
  >{\raggedright\arraybackslash}p{(\linewidth - 6\tabcolsep) * \real{0.0897}}
  >{\raggedright\arraybackslash}p{(\linewidth - 6\tabcolsep) * \real{0.3974}}
  >{\raggedright\arraybackslash}p{(\linewidth - 6\tabcolsep) * \real{0.3077}}
  >{\raggedright\arraybackslash}p{(\linewidth - 6\tabcolsep) * \real{0.2051}}@{}}
\toprule\noalign{}
\begin{minipage}[b]{\linewidth}\raggedright
Solvent
\end{minipage} & \begin{minipage}[b]{\linewidth}\raggedright
\(\Delta_{\text{vap}}H\) (kJ/mol)
\end{minipage} & \begin{minipage}[b]{\linewidth}\raggedright
Boiling Point (degrees C)
\end{minipage} & \begin{minipage}[b]{\linewidth}\raggedright
Molecular Weight
\end{minipage} \\
\midrule\noalign{}
\endhead
\bottomrule\noalign{}
\endlastfoot
Water & 40.7 & 100 & 18 \\
Ethanol & 38.6 & 78 & 46 \\
Acetone & 31.3 & 56 & 58 \\
Diethyl ether & 26.5 & 35 & 74 \\
H\(_2\)S & 18.7 & --60 & 34 \\
\end{longtable}
}

Water has a disproportionately high heat of vaporisation relative to its
molecular weight --- a direct consequence of its extensive hydrogen
bonding.

\textbf{Thermodynamic decomposition.} The high \(\Delta_{\text{vap}}H\)
is dominated by the enthalpic cost of breaking H-bonds. On vaporisation
each molecule loses about 3.4 H-bonds; if each is worth about 10 kJ/mol
on rupture, the net enthalpy is about 34 kJ/mol from H-bonds alone after
correcting for double-counting of shared bonds, matched by about 7
kJ/mol from the work of expansion (\(P\Delta V\)) and other
contributions. The corresponding entropy gain on vaporisation is
\(\Delta_{\text{vap}}S \approx \Delta_{\text{vap}}H/T_b \approx 109\) J
mol\(^{-1}\) K\(^{-1}\) (Trouton's rule predicts about 85 J mol\(^{-1}\)
K\(^{-1}\) for ``normal'' liquids; water's higher value reflects the
unusually structured liquid state). The enormous
\(\Delta_{\text{vap}}H\) is what makes evaporative cooling so efficient:
each gram of sweat carries away the energy needed to disassemble its
share of the H-bond network.

\subsubsection{High Surface Tension (\texorpdfstring{\ensuremath{\gamma}}{gamma} = 72.8 mN/m at 20 degrees
C)}\label{high-surface-tension-ux3b3-72.8-mnm-at-20-degrees-c}

Surface molecules have fewer neighbours to form H-bonds with, creating a
net inward force. Water's surface tension is among the highest of common
liquids (except mercury). This allows:

\begin{itemize}
\tightlist
\item
  Insects (water striders) to walk on water --- their legs distribute
  body weight across about 2 cm\(^2\) of surface area,
  resulting in pressure below the surface tension threshold
\item
  Capillary action: water rises in narrow tubes against gravity,
  critical for tall plants (see
  \cref{sec:unit_VIII_plant_structure_and_water})
\item
  Pulmonary alveolar \hyperref[gl:surfactant]{\textbf{surfactant}}
  (dipalmitoylphosphatidylcholine, DPPC) reduces surface tension to
  about 25 mN/m to prevent lung collapse during expiration
\end{itemize}

\begin{quote}
\textbf{Clinical Connection: Neonatal Respiratory Distress Syndrome
(NRDS)}

Premature infants often lack sufficient pulmonary surfactant because
type II pneumocytes do not mature until about 35 weeks of
gestation. Without surfactant, alveolar surface tension is too high, and
alveoli collapse during expiration (atelectasis). The infant must expend
enormous energy to re-inflate collapsed alveoli with each breath.
Treatment involves administering exogenous surfactant (e.g., beractant)
directly into the trachea. The development of synthetic surfactant
therapy in the 1980s dramatically reduced neonatal mortality.
\end{quote}

\subsubsection{Density Anomaly --- Ice
Floats}\label{density-anomaly-ice-floats}

Liquid water is denser than ice because the hexagonal lattice of ice
(perfect tetrahedral H-bond geometry) is less compact than the
fluctuating, partially broken H-bond network of liquid water. This
causes:

\begin{itemize}
\tightlist
\item
  Ice floats, so aquatic ecosystems are insulated below, preventing
  freeze-solid
\item
  Maximum density at 4 degrees C, so lakes stratify with cold water at 4
  degrees C at the bottom
\item
  Floating sea ice reflects solar radiation (albedo effect),
  contributing to climate regulation
\end{itemize}

If ice were denser than water, as is the case for nearly most other
substances, lakes and oceans would freeze from the bottom up, making
most aquatic life unsustainable in temperate climates where water bodies
freeze seasonally.

\begin{quote}
\textbf{Concept Check 2:} If ice sank instead of floated, how would this
affect lake ecosystems during winter? Consider both the physical
arrangement and biological consequences.
\end{quote}

\begin{center}\rule{0.5\linewidth}{0.5pt}\end{center}

\subsection{Water as a Solvent}\label{water-as-a-solvent}

\subsubsection{Hydrophilic Interactions and the Hydration
Shell}\label{hydrophilic-interactions-and-the-hydration-shell}

\textbf{Hydrophilic} (``water-loving'') molecules carry polar groups or
charges that interact favourably with water molecules via hydrogen
bonding or electrostatic attraction. When ionic compounds dissolve in
water, each ion is surrounded by a \textbf{hydration shell} of oriented
water molecules:

\begin{itemize}
\tightlist
\item
  Na\(^+\): surrounded by about 6 water molecules (O pointing
  toward Na\(^+\))
\item
  Cl\(^-\): surrounded by about 8 water molecules (H pointing
  toward Cl\(^-\))
\item
  Mg\(^{2+}\): surrounded by about 6 water molecules held very
  tightly (residence time about 1 \texorpdfstring{\ensuremath{\mu}}{mu}s vs.~about 10 ps
  for Na\(^+\))
\end{itemize}

The energy cost of separating the ionic lattice is offset by the energy
released from forming the hydration shell. This balance is captured by
the \textbf{Born equation} for the free energy of ion solvation:

\[\Delta G_{\text{solv}} = -\frac{z^2 e^2 N_A}{8\pi\varepsilon_0 r}\left(1 - \frac{1}{\varepsilon_r}\right) \tag{2.2} \label{eq:unit_I_water_and_life_item_2}\]

where \(z\) is the ion charge, \(r\) is the ionic radius, and
\(\varepsilon_r\) is the dielectric constant of the solvent.

\begin{figure}[htbp]
\centering
\includegraphics[width=0.82\linewidth,height=4.2in,keepaspectratio]{../figures/mermaid_inline/inline_mermaid_0019_3d0d11815c85.png}
\caption{Sequence diagram for Hydrophilic Interactions and the Hydration Shell showing ordered interaction among Na+ Ion, Water Molecule 1, Water Molecule 2, and Water Molecule 3.}
\end{figure}

\subsubsection{The Dielectric Constant and Electrostatic
Screening}\label{the-dielectric-constant-and-electrostatic-screening}

Water's high dielectric constant (\(\varepsilon\) = 78.5 at 25 degrees
C) dramatically weakens electrostatic interactions between charges.
Coulomb's law in a medium:

\[F = \frac{q_1 q_2}{4\pi\varepsilon_0\varepsilon_r r^2} \tag{2.3} \label{eq:unit_I_water_and_life_item_3}\]

In water, the force between two charges is reduced by a factor of 78.5
compared to vacuum. This is why NaCl (lattice energy 787 kJ/mol)
dissolves readily in water but not in hexane (\(\varepsilon\) = 1.9).

\subsubsection{Hydrophobic Effect}\label{hydrophobic-effect}

\textbf{Hydrophobic} (``water-fearing'') molecules carry nonpolar groups
(C--H, C--C) that cannot H-bond with water. When forced into aqueous
solution, surrounding water molecules form a rigid, ordered cage
(\textbf{clathrate structure}) to maintain their H-bonding network. This
\textbf{decreases entropy} (\(\Delta S < 0\)).

The \textbf{hydrophobic effect} is therefore driven by entropy:
aggregating nonpolar molecules together minimises the clathrate-ordered
water, restoring freedom to bulk water molecules (\(\Delta S\)
increases). The free energy change is:

\[\Delta G = \Delta H - T\Delta S \tag{2.4} \label{eq:unit_I_water_and_life_item_4}\]

Because \(\Delta S\) is positive (entropy of the system increases) and
\(\Delta H \approx 0\) (van der Waals interactions in the nonpolar
aggregate are similar to water-water interactions), \(\Delta G < 0\) ---
the aggregation is spontaneous.

\textbf{Quantitative aspects:} The transfer free energy of a hydrocarbon
chain from water to a nonpolar environment is approximately --3.1 kJ/mol
per CH\(_2\) group. For a 16-carbon palmitic acid tail, the hydrophobic
driving force is roughly 50 kJ/mol per chain --- a substantial
thermodynamic force.

The hydrophobic effect explains: - \textbf{Lipid bilayer formation} ---
phospholipid tails aggregate, forming the barrier of cell membranes (see
\cref{sec:unit_II_membrane_transport}) -
\textbf{\hyperref[gl:protein]{Protein} folding} --- hydrophobic side
chains buried in the core; hydrophilic on surface - \textbf{Micelle
formation} --- soap molecules orient with polar heads facing water,
nonpolar tails inside - \textbf{Drug partitioning} --- the octanol-water
partition coefficient (\(\log P\)) predicts drug membrane permeability

\begin{quote}
\textbf{Concept Check 3:} Why does the hydrophobic effect become
stronger at higher temperatures? Consider the \(T\Delta S\) term in the
Gibbs free energy equation.
\end{quote}

\subsubsection{Amphipathic Molecules and
Self-Assembly}\label{amphipathic-molecules-and-self-assembly}

Molecules with both hydrophilic and hydrophobic regions are
\textbf{amphipathic} (also called amphiphilic). Their self-assembly
behaviour depends on geometry:

{\def\LTcaptype{none} % do not increment counter
\begin{longtable}[]{@{}
  >{\raggedright\arraybackslash}p{(\linewidth - 4\tabcolsep) * \real{0.1852}}
  >{\raggedright\arraybackslash}p{(\linewidth - 4\tabcolsep) * \real{0.5556}}
  >{\raggedright\arraybackslash}p{(\linewidth - 4\tabcolsep) * \real{0.2593}}@{}}
\toprule\noalign{}
\begin{minipage}[b]{\linewidth}\raggedright
Shape
\end{minipage} & \begin{minipage}[b]{\linewidth}\raggedright
Structure Formed
\end{minipage} & \begin{minipage}[b]{\linewidth}\raggedright
Example
\end{minipage} \\
\midrule\noalign{}
\endhead
\bottomrule\noalign{}
\endlastfoot
Cone (large head, small tail) & Micelles & Bile salts, SDS \\
Cylinder (head and tail cross-sections similar) & Bilayers / vesicles &
Phospholipids \\
Inverted cone (small head, large tail) & Inverted micelles & Some lipids
in hexagonal phase \\
\end{longtable}
}

The \textbf{critical micelle concentration (CMC)} is the concentration
above which amphipathic molecules spontaneously form micelles. For SDS,
CMC = 8.2 mM; for bile salts, CMC = 2--12 mM depending on the species.

\begin{center}\rule{0.5\linewidth}{0.5pt}\end{center}

\subsection{Osmotic Pressure --- The van 't Hoff
Equation}\label{osmotic-pressure-the-van-t-hoff-equation}

When a selectively permeable membrane separates solutions of different
solute concentrations, \textbf{water moves from low solute (high water
potential) to high solute (low water potential)} by
\hyperref[gl:osmosis]{\textbf{osmosis}}.

\subsubsection{Derivation of the van 't Hoff
Equation}\label{derivation-of-the-van-t-hoff-equation}

The osmotic pressure can be derived from the chemical potential of
water. The chemical potential of water in a solution is:

\[\mu_{\text{water}} = \mu^{\circ}_{\text{water}} + RT\ln a_w \tag{2.5} \label{eq:unit_I_water_and_life_item_5}\]

where \(a_w\) is the water activity. For dilute ideal solutions,
\(a_w \approx 1 - x_s\) where \(x_s\) is the mole fraction of solute.
Using the approximation \(\ln(1-x) \approx -x\) for small \(x\):

\[\mu_{\text{water}} \approx \mu^{\circ}_{\text{water}} - RTx_s \tag{2.6} \label{eq:unit_I_water_and_life_item_6}\]

The osmotic pressure is the external pressure needed to restore the
chemical potential to that of pure water:

\[\pi \bar{V} = RTx_s \tag{2.7} \label{eq:unit_I_water_and_life_item_7}\]

For dilute solutions, \(x_s \approx n_s/n_w\) and
\(n_w\bar{V} \approx V\), giving the \textbf{van 't Hoff equation}:

\[\pi = iCRT \tag{2.8} \label{eq:unit_I_water_and_life_item_8}\]

where: - \textbf{\(i\)} = van 't Hoff factor (number of particles per
formula unit; 2 for NaCl -\textgreater{} Na\(^+\) + Cl\(^-\)) -
\textbf{\(C\)} = molar solute concentration (mol L\(^{-1}\)) -
\textbf{\(R\)} = 8.314 J mol\(^{-1}\) K\(^{-1}\) - \textbf{\(T\)} =
temperature in Kelvin

\subsubsection{Worked Examples}\label{worked-examples-1}

\textbf{Example 1 --- Blood plasma osmotic pressure:} Blood plasma
osmolarity is about 0.3 Osm (made up of NaCl, glucose, proteins, etc.).
Using \(i \cdot C_{\text{total}}\) = 0.3 mol/L, T = 310 K:

\[\pi = 0.3 \times 8.314 \times 310 = 773 \; \text{kPa} \approx 7.6 \; \text{atm} \tag{2.9} \label{eq:unit_I_water_and_life_item_9}\]

This means a 7.6 atm pressure difference would be required to completely
prevent osmotic water movement --- illustrating why
\hyperref[gl:turgor-pressure]{\textbf{turgor pressure}} is critical for
plant cell rigidity (see
\cref{sec:unit_VIII_plant_structure_and_water}).

\textbf{Example 2 --- Sucrose solution:} What osmotic pressure is
generated by a 0.5 M sucrose solution at 25 degrees C? Sucrose does not
dissociate, so \(i = 1\):

\[\pi = 1 \times 0.5 \times 8.314 \times 298 = 1,239 \; \text{kPa} \approx 12.2 \; \text{atm} \tag{2.10} \label{eq:unit_I_water_and_life_item_10}\]

This is why plant cells can generate enormous turgor pressures ---
enough for roots to crack concrete.

\textbf{Example 3 --- Intravenous fluids:} A nurse prepares an isotonic
saline drip using NaCl. What concentration is needed? Plasma osmolarity
= 0.3 Osm; NaCl \(i = 2\):

\[C = \frac{0.3}{2} = 0.15 \; \text{M} = 0.15 \times 58.44 = 8.77 \; \text{g/L} \approx 0.9\% \tag{2.11} \label{eq:unit_I_water_and_life_item_11}\]

This is the basis for ``normal saline'' (0.9\% NaCl), almost universally
used in clinical settings.

\subsubsection{Tonicity and Cell
Behaviour}\label{tonicity-and-cell-behaviour}

\begin{figure}[htbp]
\centering
\includegraphics[width=0.82\linewidth,height=4.2in,keepaspectratio]{../figures/mermaid_inline/inline_mermaid_0020_cb442798a32d.png}
\caption{Tonicity and osmotic responses. Osmotic behaviour of animal and plant cells in solutions of different tonicity. Plant cells are protected from lysis by their rigid cell wall, which generates turgor pressure to counterbalance osmotic influx.}
\end{figure}

{\def\LTcaptype{none} % do not increment counter
\begin{longtable}[]{@{}
  >{\raggedright\arraybackslash}p{(\linewidth - 6\tabcolsep) * \real{0.1250}}
  >{\raggedright\arraybackslash}p{(\linewidth - 6\tabcolsep) * \real{0.3472}}
  >{\raggedright\arraybackslash}p{(\linewidth - 6\tabcolsep) * \real{0.2917}}
  >{\raggedright\arraybackslash}p{(\linewidth - 6\tabcolsep) * \real{0.2361}}@{}}
\toprule\noalign{}
\begin{minipage}[b]{\linewidth}\raggedright
Condition
\end{minipage} & \begin{minipage}[b]{\linewidth}\raggedright
Relative Solute {[}outside{]}
\end{minipage} & \begin{minipage}[b]{\linewidth}\raggedright
Red Blood Cell Effect
\end{minipage} & \begin{minipage}[b]{\linewidth}\raggedright
Plant Cell Effect
\end{minipage} \\
\midrule\noalign{}
\endhead
\bottomrule\noalign{}
\endlastfoot
Isotonic (0.9\% NaCl) & = inside & Normal biconcave disc & Flaccid \\
Hypotonic (distilled water) & \textless{} inside & Swells
-\textgreater{} lyses (haemolysis) & Turgid (ideal for growth) \\
Hypertonic (2.5\% NaCl) & \textgreater{} inside & Shrinks (crenation) &
Plasmolysis \\
\end{longtable}
}

\begin{quote}
\textbf{Clinical Connection: Osmotic Demyelination Syndrome}

When hyponatraemia (low blood sodium, \textless{} 135 mM) is corrected
too rapidly with hypertonic saline, brain cells --- which had adapted to
low osmolarity by losing organic osmolytes --- suddenly find themselves
in a hypertonic environment. Water rushes out of
\hyperref[gl:neuron]{\textbf{neuron}}s and oligodendrocytes, causing
demyelination of pontine neurons. This \textbf{osmotic demyelination
syndrome} (formerly called central pontine myelinolysis) causes
devastating neurological damage including ``locked-in syndrome.'' The
safe correction rate is less than 8 mM/24 h.
\end{quote}

\begin{quote}
\textbf{Concept Check 4:} A wilted lettuce leaf placed in cold water
becomes crisp again. Explain this observation using the van 't Hoff
equation and the concept of turgor pressure.
\end{quote}

\begin{center}\rule{0.5\linewidth}{0.5pt}\end{center}

\subsection{Diffusion --- Fick's First
Law}\label{diffusion-ficks-first-law}

Small molecules move passively down their \textbf{concentration
gradients} by Brownian motion --- diffusion. Fick's First Law describes
the net flux \(J\) (mol m\(^{-2}\) s\(^{-1}\)):

\[J = -D \frac{d[C]}{dx} \tag{2.12} \label{eq:unit_I_water_and_life_item_12}\]

where: - \textbf{\(D\)} = diffusion coefficient (m\(^2\) s\(^{-1}\)) ---
depends on molecule size, shape, and temperature - \textbf{\(d[C]/dx\)}
= concentration gradient (mol m\(^{-4}\))

The negative sign indicates flux is in the direction of decreasing
concentration.

\subsubsection{The Stokes-Einstein
Equation}\label{the-stokes-einstein-equation}

The diffusion coefficient is predicted by the Stokes-Einstein equation:

\[D = \frac{k_B T}{6\pi\eta r} \tag{2.13} \label{eq:unit_I_water_and_life_item_13}\]

where \(k_B\) is the Boltzmann constant, \(T\) is temperature, \texorpdfstring{\ensuremath{\eta}}{eta} is
solvent viscosity, and \(r\) is the hydrodynamic radius of the solute.
This equation predicts that:

\begin{itemize}
\tightlist
\item
  Larger molecules diffuse more slowly (\(D \propto 1/r\))
\item
  Higher temperature increases diffusion (\(D \propto T\))
\item
  More viscous media slow diffusion (\(D \propto 1/\eta\))
\end{itemize}

\textbf{Reference diffusion coefficients at 37 degrees C:}

{\def\LTcaptype{none} % do not increment counter
\begin{longtable}[]{@{}lll@{}}
\toprule\noalign{}
Molecule & D (m\(^2\) s\(^{-1}\)) & Hydrodynamic Radius (nm) \\
\midrule\noalign{}
\endhead
\bottomrule\noalign{}
\endlastfoot
O\(_2\) (in water) & \(2.1 \times 10^{-9}\) & 0.16 \\
H\(_2\)O (self-diffusion) & \(2.3 \times 10^{-9}\) & 0.14 \\
Glucose & \(6.7 \times 10^{-10}\) & 0.37 \\
ATP & \(3.7 \times 10^{-10}\) & 0.55 \\
Haemoglobin & \(6.9 \times 10^{-11}\) & 3.1 \\
IgG antibody & \(3.8 \times 10^{-11}\) & 5.4 \\
DNA (small \hyperref[gl:plasmid]{\textbf{plasmid}}) &
\(5.0 \times 10^{-13}\) & about 50 \\
\end{longtable}
}

\subsubsection{Mean Diffusion Distance and Biological
Design}\label{mean-diffusion-distance-and-biological-design}

\textbf{Mean diffusion distance:} \(\bar{x} = \sqrt{2Dt}\).

For O\(_2\) to diffuse 1 \texorpdfstring{\ensuremath{\mu}}{mu}m in water:
\(t = x^2/2D = (10^{-6})^2/(2 \times 2.1 \times 10^{-9}) \approx 0.24 \; \text{\texorpdfstring{\ensuremath{\mu}}{mu}s}\).
For 1 mm: \(t \approx 238 \; \text{s} \approx 4 \; \text{min}\).

This is why multicellular organisms exceeding about 1 mm
require \textbf{circulatory systems} --- diffusion alone is too slow. It
also explains why:

\begin{itemize}
\tightlist
\item
  Mitochondria are typically \textless{} 1 \texorpdfstring{\ensuremath{\mu}}{mu}m wide (O\(_2\) diffusion
  time \textless{} 1 \texorpdfstring{\ensuremath{\mu}}{mu}s)
\item
  Capillaries are spaced no more than about 100 \texorpdfstring{\ensuremath{\mu}}{mu}m apart in
  metabolically active tissues
\item
  Alveolar walls are about 0.2 \texorpdfstring{\ensuremath{\mu}}{mu}m thick (gas exchange must be nearly
  instantaneous)
\item
  Neurons require active transport along axons (diffusion over 1 m would
  take about 16 years!)
\end{itemize}

\subsubsection{Biological Applications of Fick's
Law}\label{biological-applications-of-ficks-law}

\textbf{Application 1 --- Oxygen delivery to muscle fibres:}

During exercise, the O\(_2\) concentration at the capillary wall is
approximately 0.13 mM, and at the centre of a muscle fibre (radius 25
\texorpdfstring{\ensuremath{\mu}}{mu}m) it drops to 0.01 mM. The flux across the fibre radius:

\[J = -D\frac{\Delta C}{\Delta x} = -(2.1 \times 10^{-9})\frac{(0.01 - 0.13) \times 10^3}{25 \times 10^{-6}} = 1.0 \times 10^{-2} \; \text{mol m}^{-2}\text{s}^{-1} \tag{2.14} \label{eq:unit_I_water_and_life_item_14}\]

\textbf{Application 2 --- CO\(_2\) removal from tissues:}

CO\(_2\) produced by \hyperref[gl:aerobic]{\textbf{aerobic}} metabolism
(\(D = 1.9 \times 10^{-9}\) m\(^2\)/s in water) must diffuse from
mitochondria to capillaries. Because CO\(_2\) diffuses about 20 times
faster in air than O\(_2\) in water, gas exchange in the lungs is
limited primarily by O\(_2\) diffusion, not CO\(_2\).

\textbf{Application 3 --- Neurotransmitter diffusion in the synaptic
cleft:}

The synaptic cleft is about 20 nm wide. For acetylcholine
(\(D \approx 5 \times 10^{-10}\) m\(^2\)/s):

\[t = \frac{x^2}{2D} = \frac{(20 \times 10^{-9})^2}{2 \times 5 \times 10^{-10}} = 4 \times 10^{-7} \; \text{s} = 0.4 \; \text{$\mu$s} \tag{2.15} \label{eq:unit_I_water_and_life_item_15}\]

This submicrosecond diffusion time ensures rapid neurotransmission, far
faster than the about 0.5 ms delay of synaptic transmission
(which is dominated by vesicle fusion and receptor activation).

\begin{quote}
\textbf{Concept Check 5:} A single-celled organism is spherical with a
radius of 500 \texorpdfstring{\ensuremath{\mu}}{mu}m. Using the diffusion time equation \(t = r^2/(6D)\) for
three-dimensional diffusion, calculate how long it would take O\(_2\) to
diffuse from the cell surface to the centre. Is this organism viable if
it relies solely on diffusion for O\(_2\) supply?
\end{quote}

\begin{center}\rule{0.5\linewidth}{0.5pt}\end{center}

\subsubsection{Solubility Rules and the Thermodynamics of
Dissolution}\label{solubility-rules-and-the-thermodynamics-of-dissolution}

Whether a solute dissolves in water is governed by the Gibbs free energy
of solution:

\begin{equation}
\Delta G_{\text{soln}} = \Delta H_{\text{soln}} - T\Delta S_{\text{soln}}
\label{eq:unit_I_dissolution_dG}
\end{equation}

The enthalpy term has three contributions: breaking solute-solute
interactions (\(\Delta H_1 > 0\)), breaking solvent-solvent interactions
(\(\Delta H_2 > 0\)), and forming solute-solvent interactions
(\(\Delta H_3 < 0\)). The sign and magnitude of
\(\Delta H_{\text{soln}} = \Delta H_1 + \Delta H_2 + \Delta H_3\)
depends on which set of interactions wins.

{\def\LTcaptype{none} % do not increment counter
\begin{longtable}[]{@{}
  >{\raggedright\arraybackslash}p{(\linewidth - 8\tabcolsep) * \real{0.0909}}
  >{\raggedright\arraybackslash}p{(\linewidth - 8\tabcolsep) * \real{0.1273}}
  >{\raggedright\arraybackslash}p{(\linewidth - 8\tabcolsep) * \real{0.3273}}
  >{\raggedright\arraybackslash}p{(\linewidth - 8\tabcolsep) * \real{0.3273}}
  >{\raggedright\arraybackslash}p{(\linewidth - 8\tabcolsep) * \real{0.1273}}@{}}
\toprule\noalign{}
\begin{minipage}[b]{\linewidth}\raggedright
Class
\end{minipage} & \begin{minipage}[b]{\linewidth}\raggedright
Example
\end{minipage} & \begin{minipage}[b]{\linewidth}\raggedright
Sign of \(\Delta H\)
\end{minipage} & \begin{minipage}[b]{\linewidth}\raggedright
Sign of \(\Delta S\)
\end{minipage} & \begin{minipage}[b]{\linewidth}\raggedright
Outcome
\end{minipage} \\
\midrule\noalign{}
\endhead
\bottomrule\noalign{}
\endlastfoot
Salt with high lattice E & NaCl & small (+) & (+) & Soluble, slightly
endothermic \\
Salt with strong hydration & CaCl\(_2\) & (\texorpdfstring{\ensuremath{-}}{-}) & (+) & Soluble,
exothermic (heat packs) \\
Salt with weak hydration & AgCl & (+) & (+) but small & Insoluble; H
wins \\
Polar covalent & sucrose, glucose & small & (+) & Soluble \\
Nonpolar in water & hexane, lipid tail & about 0 & \textbf{(\texorpdfstring{\ensuremath{-}}{-})}
& \textbf{Insoluble} (entropy! see hydrophobic effect) \\
Amphipathic & phospholipid, soap & mixed & mixed & Self-assembles into
bilayers/micelles \\
\end{longtable}
}

The textbook rule ``like dissolves like'' is a shorthand for: when
solute-solvent interactions are similar in strength to the bulk
interactions they replace, \(\Delta H_{\text{soln}} \approx 0\) and the
(usually positive) entropy of mixing drives dissolution. When
interactions differ greatly --- as for ionic salts in nonpolar solvents,
or alkanes in water --- \(\Delta H\) and/or \(\Delta S\) oppose
dissolution and the solute remains in its own phase.

\subsubsection{The Hydrophobic Effect
Quantified}\label{the-hydrophobic-effect-quantified}

The \textbf{hydrophobic effect} is the principal driving force for
protein folding, lipid bilayer assembly, and many ligand-binding
reactions. Despite its name, it is \textbf{not} a force; it is a
thermodynamic consequence of the difficulty water has accommodating
nonpolar surfaces.

When a nonpolar molecule is forced into water, the surrounding waters
cannot form their full complement of H-bonds with it. They preserve
their network by reorienting around the solute, sacrificing
translational and rotational entropy. The free energy cost is dominated
not by enthalpy but by \emph{entropy}:

{\def\LTcaptype{none} % do not increment counter
\begin{longtable}[]{@{}
  >{\raggedright\arraybackslash}p{(\linewidth - 6\tabcolsep) * \real{0.0864}}
  >{\raggedright\arraybackslash}p{(\linewidth - 6\tabcolsep) * \real{0.2346}}
  >{\raggedright\arraybackslash}p{(\linewidth - 6\tabcolsep) * \real{0.4444}}
  >{\raggedright\arraybackslash}p{(\linewidth - 6\tabcolsep) * \real{0.2346}}@{}}
\toprule\noalign{}
\begin{minipage}[b]{\linewidth}\raggedright
Process
\end{minipage} & \begin{minipage}[b]{\linewidth}\raggedright
\(\Delta H\) (kJ/mol)
\end{minipage} & \begin{minipage}[b]{\linewidth}\raggedright
\(T\Delta S\) at 25 \(^\circ\)C (kJ/mol)
\end{minipage} & \begin{minipage}[b]{\linewidth}\raggedright
\(\Delta G\) (kJ/mol)
\end{minipage} \\
\midrule\noalign{}
\endhead
\bottomrule\noalign{}
\endlastfoot
Transfer 1 mol benzene from oil to water & small (about 0) & \(-22\)
(entropy lost) & \(+22\) \\
Transfer 1 mol --CH\(_2\)-- from water to oil & small & \(+3.1\) &
\(-3.1\) (per CH\(_2\)) \\
Hydrocarbon aggregation in water (per CH\(_2\)) & small & \(+3.1\) &
\(-3.1\) \\
\end{longtable}
}

The entropic origin has a striking signature: the hydrophobic effect
\emph{strengthens} with temperature (up to about 70
\(^\circ\)C) because the \(T\Delta S\) term grows. This is why heat
denatures proteins partly by \emph{favouring} hydrophobic exposure (the
cold-denaturation paradox). A 16-carbon palmitate tail, contributing
about 3 kJ/mol per --CH\(_2\)--, generates about 50
kJ/mol of effective driving force into the membrane interior --- enough
to make membrane partitioning essentially irreversible at biological
concentrations.

\subsubsection{Colligative Property: Osmotic Pressure (van 't
Hoff)}\label{colligative-property-osmotic-pressure-van-t-hoff}

The \textbf{osmotic pressure} of a dilute solution is given by the
\textbf{van 't Hoff equation}:

\begin{equation}
\pi = i\,C\,R\,T
\label{eq:unit_I_vant_hoff}
\end{equation}

where \(i\) is the van 't Hoff factor, \(C\) is the molar solute
concentration, \(R\) is the gas constant, and \(T\) is the absolute
temperature. (We re-derive \cref{eq:unit_I_vant_hoff} from chemical
potential below.)

\textbf{Worked example: the osmotic pressure on a red blood cell.}
Plasma osmolarity is 290 mOsm at body temperature (310 K). What pressure
must the cell membrane withstand if placed in pure water? Using
\cref{eq:unit_I_vant_hoff} with \(iC\) = 0.290 osm L\(^{-1}\):

\[\pi = 0.290 \times 8.314\,\text{J mol}^{-1}\text{K}^{-1} \times 310\,\text{K} = 747\,\text{kPa} \approx 7.4\,\text{atm} \label{eq:unit_I_water_and_life_item_16}\]

This is the pressure of a vehicle tyre. The lipid bilayer cannot resist
such force --- placed in distilled water, a red blood cell swells and
lyses (haemolysis) within seconds. Plant cells, protected by their rigid
cellulose wall, can sustain similar pressures as turgor; this turgor is
what keeps lettuce crisp and what allows plant roots to crack pavement.

\subsubsection{Water Activity and Biological
Processes}\label{water-activity-and-biological-processes}

\textbf{Water activity} (\(a_w\)) is the effective ``thermodynamic
concentration'' of water and is governed by the same chemical potential
expression that gives osmotic pressure:

\begin{equation}
\mu_w = \mu_w^\circ + RT\ln a_w
\label{eq:unit_I_water_activity}
\end{equation}

Adding solutes lowers \(a_w\) below 1 (pure water), which (i) reduces
the rate of water-dependent reactions including hydrolysis, (ii)
inhibits microbial growth, and (iii) suppresses freezing. Most cellular
processes proceed at \(a_w \gtrsim 0.99\). When \(a_w\) falls below 0.6
(e.g., honey, raisins, dried meat), microbial growth halts entirely ---
the basis of most traditional preservation methods. We will quantify
this further in \S 8 below.

\subsubsection{Proton Hopping --- The Grotthuss
Mechanism}\label{proton-hopping-the-grotthuss-mechanism}

Protons (H\(^+\), actually the hydronium ion H\(_3\)O\(^+\)) diffuse
through water roughly 7\(\times\) faster than any other ion of similar
size. The reason is that protons do not move as discrete particles;
instead, the \emph{charge} hops from one water molecule to the next via
a chain of H-bond rearrangements --- the \textbf{Grotthuss mechanism}
(Theodor von Grotthuss, 1806).

Visualise a chain of H-bonded waters linking a proton donor to a proton
acceptor:

\[\text{H}_3\text{O}^+\cdots\text{H--O--H}\cdots\text{H--O--H}\cdots\text{O--H}\cdots\text{A}^- \label{eq:unit_I_water_and_life_item_17}\]

Concerted reorientation of the O--H bonds along the chain (each
requiring about 1 ps) transports a proton across many molecular
diameters in a few picoseconds, \emph{without} any individual water
moving more than fractions of an angstrom. The effective diffusion
coefficient of H\(^+\) in water (9.3 \(\times\) 10\(^{-9}\) m\(^2\)
s\(^{-1}\)) is dominated by Grotthuss hopping rather than vehicular
diffusion.

\textbf{Biological implications:} - \textbf{ATP synthase.} The F\(_o\)
subunit conducts protons across the mitochondrial inner membrane along a
chain of buried waters and protonatable residues (Asp, Glu, Lys);
rotational coupling drives ATP synthesis. The Grotthuss mechanism allows
fast H\(^+\) flux without bulk water flow. - \textbf{Cytochrome c
oxidase} and other proton pumps shuttle protons through ``proton wires''
--- ordered chains of waters and titratable residues. -
\textbf{Aquaporin proton exclusion.} Aquaporins must allow rapid
H\(_2\)O passage but \textbf{block} proton transport (otherwise membrane
potentials would collapse). They achieve this by reorienting a central
water so its O--H bonds break the H-bond chain, defeating the Grotthuss
mechanism while permitting ``vehicular'' water transport. -
\textbf{Photosystem II.} Splits water into O\(_2\) + 4H\(^+\) + 4e\(^-\)
on the lumenal side of the thylakoid; the four protons join the
proton-motive force via H-bond chains in the protein.

\subsubsection{Ice Nucleation and
Cryopreservation}\label{ice-nucleation-and-cryopreservation}

Pure liquid water can be \textbf{supercooled} to about --40 \(^\circ\)C
before homogeneous ice nucleation forces freezing. In real solutions,
\textbf{heterogeneous nucleation} (initiated by surfaces, dust, or
specialised proteins) usually freezes water far above --40 \(^\circ\)C.
Many bacteria (e.g., \emph{Pseudomonas syringae}) express ice-nucleating
proteins that template ice at --2 \(^\circ\)C, weaponising frost on
plant tissues so the bacteria can colonise the resulting wounds.

Cryopreservation of cells (sperm, oocytes, stem cells, organs) exploits
this thermodynamics. Two strategies:

\begin{enumerate}
\def\labelenumi{\arabic{enumi}.}
\tightlist
\item
  \textbf{Slow freezing with cryoprotectants.} Glycerol or DMSO at 1--3
  M lowers the freezing point colligatively (\cref{eq:unit_I_vant_hoff})
  and replaces water in hydrogen bonds at membrane surfaces, preventing
  the membrane damage caused by cellular dehydration. Slow cooling
  (0.5--1 \(^\circ\)C/min) lets water leave cells before intracellular
  ice forms.
\item
  \textbf{Vitrification.} Very high cryoprotectant concentrations
  (50\%+) plus rapid cooling (\textgreater10\(^4\) \(^\circ\)C/min)
  glasses water before any ice nucleates. Used for human oocytes and
  embryos with \textgreater90\% post-thaw survival.
\end{enumerate}

The wood frog (\emph{Rana sylvatica}) survives whole-body freezing for
weeks using a natural variant of strategy 1: liver glycogen is rapidly
converted to glucose at about 0 \(^\circ\)C, raising blood
glucose to about 250 mM (10\(\times\) hyperglycaemia by
mammalian standards). Glucose enters cells and lowers the freezing point
intracellularly, while extracellular ice forms harmlessly.

\begin{quote}
\textbf{Concept Check 6:} Using \cref{eq:unit_I_dissolution_dG}, explain
why ammonium nitrate (NH\(_4\)NO\(_3\)) dissolves \emph{spontaneously}
in water \emph{despite} an endothermic enthalpy of solution. (Hint:
instant cold packs use this reaction.)
\end{quote}

\subsection{Aquaporins --- Molecular Water
Channels}\label{aquaporins-molecular-water-channels}

\subsubsection{Discovery and Structure}\label{discovery-and-structure}

While water can slowly permeate lipid bilayers by diffusion, many cell
types require rapid water transport. In 1992, Peter Agre discovered
\textbf{aquaporins} (AQPs) --- integral membrane proteins that form
selective water channels. Agre received the Nobel Prize in Chemistry in
2003 for this discovery.

Aquaporins are homotetramers, with each subunit forming an independent
water pore. Key structural features:

\begin{itemize}
\tightlist
\item
  \textbf{Hourglass shape:} narrow constriction (2.8 angstroms diameter)
  allows primarily single-file water passage
\item
  \textbf{Two NPA motifs} (Asn-Pro-Ala): meet at the centre of the
  channel, creating a positive electrostatic barrier
\item
  \textbf{Aromatic/arginine (ar/R) selectivity filter:} four residues at
  the extracellular constriction determine selectivity
\item
  \textbf{Proton exclusion mechanism:} the NPA motifs reorient water
  molecules mid-channel, breaking the hydrogen-bond ``wire'' that would
  allow proton hopping
\end{itemize}

\subsubsection{Aquaporin Family and
Functions}\label{aquaporin-family-and-functions}

{\def\LTcaptype{none} % do not increment counter
\begin{longtable}[]{@{}
  >{\raggedright\arraybackslash}p{(\linewidth - 6\tabcolsep) * \real{0.2045}}
  >{\raggedright\arraybackslash}p{(\linewidth - 6\tabcolsep) * \real{0.3409}}
  >{\raggedright\arraybackslash}p{(\linewidth - 6\tabcolsep) * \real{0.1818}}
  >{\raggedright\arraybackslash}p{(\linewidth - 6\tabcolsep) * \real{0.2727}}@{}}
\toprule\noalign{}
\begin{minipage}[b]{\linewidth}\raggedright
Aquaporin
\end{minipage} & \begin{minipage}[b]{\linewidth}\raggedright
Primary Location
\end{minipage} & \begin{minipage}[b]{\linewidth}\raggedright
Function
\end{minipage} & \begin{minipage}[b]{\linewidth}\raggedright
Permeability
\end{minipage} \\
\midrule\noalign{}
\endhead
\bottomrule\noalign{}
\endlastfoot
AQP0 & Eye lens fibre cells & Lens transparency & Low water \\
AQP1 & Red blood cells, kidney proximal tubule & Rapid water
reabsorption & High water \\
AQP2 & Kidney collecting duct & ADH-regulated water reabsorption &
Regulated \\
AQP3 & Kidney, skin & Water and glycerol transport & Aquaglyceroporin \\
AQP4 & Brain astrocytes & Brain water homeostasis & High water \\
AQP5 & Salivary glands, lacrimal glands & Secretory water flow & High
water \\
AQP7 & Adipocytes & Glycerol release during lipolysis &
Aquaglyceroporin \\
\end{longtable}
}

A single AQP1 channel transports approximately \(3 \times 10^9\) water
molecules per second --- among the fastest known transport rates for any
channel protein.

\subsubsection{Aquaporins in Disease}\label{aquaporins-in-disease}

\begin{quote}
\textbf{Clinical Connection: Nephrogenic Diabetes Insipidus}

\hyperref[gl:mutation]{\textbf{Mutation}}s in the AQP2
\hyperref[gl:gene]{\textbf{gene}} cause \textbf{nephrogenic diabetes
insipidus (NDI)}, a condition where the kidneys cannot concentrate urine
despite adequate levels of antidiuretic
\hyperref[gl:hormone]{\textbf{hormone}} (ADH/vasopressin). Patients
produce up to 20 L of dilute urine per day and must drink enormous
quantities of water to avoid dehydration. The most common inherited form
involves autosomal recessive mutations that cause AQP2 to misfold and be
retained in the endoplasmic reticulum. Acquired NDI can result from
lithium therapy (used in bipolar disorder), which downregulates AQP2
expression via mechanisms involving glycogen synthase kinase-3\texorpdfstring{\ensuremath{\beta}}{beta}.
\end{quote}

\begin{quote}
\textbf{Clinical Connection: Cerebral Oedema and AQP4}

AQP4, the predominant water channel in brain astrocytes, plays a dual
role in cerebral oedema. In cytotoxic oedema (e.g., ischaemic stroke),
AQP4 facilitates water influx into swelling astrocytes, worsening
damage. In vasogenic oedema (e.g., brain tumours), AQP4 helps clear
excess extracellular water. AQP4 is also the target of autoantibodies in
neuromyelitis optica (NMO/Devic's disease), an autoimmune disorder that
mimics multiple sclerosis.
\end{quote}

\begin{center}\rule{0.5\linewidth}{0.5pt}\end{center}

\subsection{Colligative Properties}\label{colligative-properties}

\textbf{Colligative properties} depend on the number of dissolved
particles, not their identity. They are direct consequences of water
activity reduction by solutes.

\subsubsection{Boiling Point Elevation and Freezing Point
Depression}\label{boiling-point-elevation-and-freezing-point-depression}

\[\Delta T_b = iK_b m \quad \text{and} \quad \Delta T_f = iK_f m \tag{2.16} \label{eq:unit_I_water_and_life_item_18}\]

where \(K_b\) = 0.512 degrees C/m (ebullioscopic constant) and \(K_f\) =
1.86 degrees C/m (cryoscopic constant) for water, \(m\) is molality
(mol/kg solvent), and \(i\) is the van 't Hoff factor.

\textbf{Worked example:} Blood serum has an effective molality of
about 0.3 m (most solutes combined). Freezing point depression:

\[\Delta T_f = 1 \times 1.86 \times 0.3 = 0.558 \; \text{degrees C} \tag{2.17} \label{eq:unit_I_water_and_life_item_19}\]

Clinical osmometers measure this freezing point depression to determine
serum osmolality. Normal serum freezes at --0.56 degrees C,
corresponding to about 290 mOsm/kg.

\subsubsection{Antifreeze Organisms}\label{antifreeze-organisms}

Many organisms survive sub-zero temperatures using biological antifreeze
strategies:

{\def\LTcaptype{none} % do not increment counter
\begin{longtable}[]{@{}
  >{\raggedright\arraybackslash}p{(\linewidth - 4\tabcolsep) * \real{0.3200}}
  >{\raggedright\arraybackslash}p{(\linewidth - 4\tabcolsep) * \real{0.3200}}
  >{\raggedright\arraybackslash}p{(\linewidth - 4\tabcolsep) * \real{0.3600}}@{}}
\toprule\noalign{}
\begin{minipage}[b]{\linewidth}\raggedright
Organism
\end{minipage} & \begin{minipage}[b]{\linewidth}\raggedright
Strategy
\end{minipage} & \begin{minipage}[b]{\linewidth}\raggedright
Mechanism
\end{minipage} \\
\midrule\noalign{}
\endhead
\bottomrule\noalign{}
\endlastfoot
Arctic fish (\emph{Dissostichus}) & Antifreeze glycoproteins (AFGPs) &
Bind ice crystal surfaces, inhibit growth \\
Wood frog (\emph{Rana sylvatica}) & Freeze tolerance + glucose
cryoprotectant & Allows extracellular ice; glucose prevents
intracellular freezing \\
Antarctic midge (\emph{Belgica}) & Dehydration + trehalose & Removes
freezable water \\
Spruce bark beetle & Glycerol accumulation & Colligative freezing point
depression to --40 degrees C \\
Winter rye & Antifreeze proteins + supercooling & Survives to --30
degrees C \\
\end{longtable}
}

\textbf{Antifreeze proteins (AFPs)} do not work by colligative effects.
Instead, they bind to specific ice crystal faces via a flat, hydrophobic
ice-binding surface, preventing crystal growth. This creates a
difference between the melting point and the freezing point called
\textbf{thermal hysteresis} (typically 1--5 degrees C). Fish AFPs can
depress the freezing point of blood to --2.5 degrees C, matching the
temperature of polar seawater.

\begin{quote}
\textbf{Concept Check 7:} The spruce bark beetle accumulates glycerol to
about 2 M concentration in its haemolymph. Using the freezing
point depression equation with \(i = 1\) and \(K_f = 1.86\) degrees C/m,
estimate the freezing point depression. Why might the observed freezing
point be lower than this calculated value?
\end{quote}

\begin{center}\rule{0.5\linewidth}{0.5pt}\end{center}

\subsection{Water Activity and Microbial
Growth}\label{water-activity-and-microbial-growth}

\subsubsection{Definition of Water
Activity}\label{definition-of-water-activity}

\textbf{Water activity} (\(a_w\)) is the ratio of the vapour pressure of
water in a solution to that of pure water:

\[a_w = \frac{p}{p^{\circ}} \tag{2.18} \label{eq:unit_I_water_and_life_item_20}\]

Pure water has \(a_w = 1.0\). Adding solutes lowers \(a_w\). For an
ideal dilute solution, \(a_w \approx 1 - x_s\) where \(x_s\) is the mole
fraction of solute.

\subsubsection{Water Activity and Food
Preservation}\label{water-activity-and-food-preservation}

Microorganisms require a minimum \(a_w\) for growth:

{\def\LTcaptype{none} % do not increment counter
\begin{longtable}[]{@{}lll@{}}
\toprule\noalign{}
Organism Type & Minimum \(a_w\) & Example \\
\midrule\noalign{}
\endhead
\bottomrule\noalign{}
\endlastfoot
Most bacteria & 0.91 & \emph{Salmonella}, \emph{E. coli} \\
Most yeasts & 0.88 & \emph{Saccharomyces} \\
Most moulds & 0.80 & \emph{Aspergillus} \\
Halophilic bacteria & 0.75 & \emph{Halobacterium} \\
Xerophilic moulds & 0.65 & \emph{Xeromyces bisporus} \\
No microbial growth & \textless{} 0.60 & Honey, dried milk powder \\
\end{longtable}
}

Traditional food preservation methods most work by reducing \(a_w\):

\begin{itemize}
\tightlist
\item
  \textbf{Salting:} NaCl reduces \(a_w\) (salted fish, cured ham)
\item
  \textbf{Sugaring:} sucrose reduces \(a_w\) (jams, candied fruits,
  honey)
\item
  \textbf{Drying:} evaporation removes water (jerky, raisins, dried
  herbs)
\item
  \textbf{Smoking:} combines drying with antimicrobial phenols
\end{itemize}

\begin{quote}
\textbf{Concept Check 8:} Honey has a water activity of approximately
0.6 and an extremely high sugar content (about 80\% w/w).
Despite this, honey has been found in Egyptian tombs still preserved
after 3,000 years. Using the concept of water activity, explain why
honey is essentially immune to microbial spoilage. Why might diluted
honey be more susceptible to
\hyperref[gl:fermentation]{\textbf{fermentation}}?
\end{quote}

\begin{quote}
\textbf{Concept Check 9:} Aquaporins permit water passage at
\(\sim 3\times 10^9\) molecules/s yet block protons. Explain
qualitatively how the channel achieves this --- and why a hypothetical
``leaky'' aquaporin allowing Grotthuss proton hopping would be lethal to
the cell.
\end{quote}

\begin{center}\rule{0.5\linewidth}{0.5pt}\end{center}

\subsection{Water and Adaptation}\label{water-and-adaptation}

Because every biochemical reaction occurs in or at water, organisms in
water-scarce or water-hostile environments have evolved striking
adaptations to defend their hydration state. Three strategies recur
across the tree of life: tolerate water loss, prevent ice formation, and
economise water turnover.

\textbf{Desiccation tolerance (anhydrobiosis).} Tardigrades, the brine
shrimp \emph{Artemia}, many rotifers, and resurrection plants such as
\emph{Selaginella} can lose more than 95\% of their body water and
revive on rewetting. They replace the hydrogen bonds that normally hold
macromolecules in shape with the disaccharide \textbf{trehalose}
(together with late-embryogenesis-abundant, LEA, proteins), which
vitrifies the cytoplasm into a glass that arrests molecular motion and
prevents protein unfolding and membrane fusion. In this state metabolism
is undetectable and survival extends by orders of magnitude.

\textbf{Freeze avoidance and tolerance.} Because dissolved solute
depresses the freezing point primarily colligatively
(\(\Delta T_f = i K_f m\)), small molecules alone cannot protect a fish
in \(-1.9\ ^\circ\)C seawater. Antarctic notothenioid fish instead
secrete \textbf{antifreeze glycoproteins} that adsorb to nascent ice
crystals and arrest their growth, lowering the freezing point
non-colligatively (thermal hysteresis) by roughly \(1\ ^\circ\)C at
millimolar concentrations. Freeze-\emph{tolerant} wood frogs take the
opposite route: they nucleate ice in extracellular spaces and flood
cells with glucose and glycerol so that intracellular water rarely
freezes.

\textbf{Water economy.} A kangaroo rat can survive on metabolic water
alone without drinking: a nasal countercurrent exchanger recovers most
respiratory water vapour, and its kidney concentrates urine roughly
fivefold above plasma osmolarity. Such osmoregulatory extremes --- from
halophilic archaea balancing molar internal K\texorpdfstring{\textsuperscript{+}}{+} against a near-saturated
brine exterior, to euryhaline fish reversing the direction of their gill
ion pumps between fresh and salt water --- show that controlling water
potential and solute load (\cref{sec:unit_II_membrane_transport}) is as
much a target of natural selection as any catalytic site.

\subsection{Current Evidence and Frontier
Biology}\label{current-evidence-and-frontier-biology-4}

For \textbf{Water --- The Molecule of Life}, frontier biology belongs
inside the evidence logic of the chapter. Chemistry-of-life claims now
connect classical bonding and thermodynamics with AI-guided structure
prediction and experimental validation. The core reading question is
this: water's biological effects depend on hydrogen bonding, colligative
context, interfaces, temperature, and solute identity.

\begin{itemize}
\tightlist
\item
  \textbf{What to verify:} identify the observation, model, assay, or
  dataset that would make the claim stronger or weaker.
\item
  \textbf{What to qualify:} state the scale, organism, cell type,
  environmental condition, or population where the claim is expected to
  hold.
\item
  \textbf{What to compare:} test at least one alternative explanation,
  baseline, or null model before treating the pattern as causal.
\item
  \textbf{What to cite:} distinguish primary evidence, review synthesis,
  public dataset, and institutional guidance; for recent or numeric
  claims, prefer the source closest to the measurement and state what
  has changed since it was published.
\end{itemize}

Use AI biomolecular models as hypothesis generators: compare confidence,
conservation, solvent exposure, and assay evidence before turning a
predicted contact into a biological claim
\citep{abramson2024alphafold3}.

\textbf{Source practice:} For structure and interaction claims, cite
experimental structures when available and treat AlphaFold 3 or AFDB
complex predictions as hypotheses to validate with confidence metrics,
conservation, mutagenesis, binding, or cryo-EM/X-ray/NMR evidence
\citep{abramson2024alphafold3,velankar2026alphafolddb2025,emblebi2026alphafoldcomplexes}.

\subsection{Summary of Water's Exceptional
Properties}\label{summary-of-waters-exceptional-properties}

{\def\LTcaptype{none} % do not increment counter
\begin{longtable}[]{@{}
  >{\raggedright\arraybackslash}p{(\linewidth - 4\tabcolsep) * \real{0.1633}}
  >{\raggedright\arraybackslash}p{(\linewidth - 4\tabcolsep) * \real{0.4082}}
  >{\raggedright\arraybackslash}p{(\linewidth - 4\tabcolsep) * \real{0.4286}}@{}}
\toprule\noalign{}
\begin{minipage}[b]{\linewidth}\raggedright
Property
\end{minipage} & \begin{minipage}[b]{\linewidth}\raggedright
Value (25 degrees C)
\end{minipage} & \begin{minipage}[b]{\linewidth}\raggedright
Biological Consequence
\end{minipage} \\
\midrule\noalign{}
\endhead
\bottomrule\noalign{}
\endlastfoot
Specific heat & 4.18 J g\(^{-1}\) K\(^{-1}\) & Thermal buffer; clinical
hypothermia therapy \\
Heat of vaporisation & 2,442 J g\(^{-1}\) & Evaporative cooling
(sweating, \hyperref[gl:transpiration]{\textbf{transpiration}}) \\
Dielectric constant & 78.5 & Ion solvation; electrostatic screening \\
Surface tension & 72.8 mN m\(^{-1}\) & Capillary rise; alveolar
function \\
Density at 0 degrees C & 0.917 g cm\(^{-3}\) & Ice floats; aquatic
ecosystems insulated \\
Density at 4 degrees C & 1.000 g cm\(^{-3}\) & Maximum density drives
lake stratification \\
Viscosity & 0.89 mPa\(\cdot\)s & Blood flow; cytoplasmic diffusion
rates \\
Thermal conductivity & 0.606 W m\(^{-1}\) K\(^{-1}\) & Heat distribution
in tissues \\
\end{longtable}
}

\subsubsection{Extremophiles and Water}\label{extremophiles-and-water}

\textbf{Halophiles} (e.g., \emph{Halobacterium salinarum}) thrive in
salt-saturated brines where water activity \(a_w\) \textless{} 0.75.
Their protective strategy: accumulate compatible solutes (glycine
betaine, ectoine) to maintain osmotic balance without disrupting
\hyperref[gl:enzyme]{\textbf{enzyme}} function. Remarkably, halophilic
enzymes have evolved surfaces rich in acidic residues (Asp, Glu) that
create a hydration shell even in near-saturated salt solutions.

\textbf{Xerophiles} (e.g., \emph{Artemia} brine shrimp) survive complete
desiccation by replacing water with \textbf{trehalose}, which forms a
glassy matrix. Trehalose H-bonds to membrane phospholipid head groups in
place of water, preserving membrane architecture. Tardigrades (water
bears) can survive in a desiccated state (\textbf{tun}) for decades,
reviving upon rehydration.

\textbf{Thermophiles} face the challenge that hydrogen bonds weaken at
high temperatures. \emph{Thermus aquaticus} enzymes compensate with
increased salt bridges, more compact hydrophobic cores, and proline
substitutions that reduce backbone flexibility.

\textbf{Psychrophiles} in Antarctic ice maintain membrane fluidity by
incorporating polyunsaturated fatty acids and short-chain fatty acids.
Their enzymes have increased flexibility (more glycine residues, fewer
prolines) to maintain activity at sub-zero temperatures --- at the cost
of reduced thermal stability.

\begin{center}\rule{0.5\linewidth}{0.5pt}\end{center}

\subsection{Computational Bridge}\label{computational-bridge-1}

The van 't Hoff relation \(\pi = iCRT\) appears throughout this chapter.
The cell module implements the same expression in SI units (output in
pascals):

\begin{Shaded}
\begin{Highlighting}[]
\ImportTok{from}\NormalTok{ biology.cell }\ImportTok{import}\NormalTok{ osmotic\_pressure}

\CommentTok{\# about 0.15 M NaCl, i = 2, body temperature}
\NormalTok{pi\_pa }\OperatorTok{=}\NormalTok{ osmotic\_pressure(}\FloatTok{0.15}\NormalTok{, temperature\_K}\OperatorTok{=}\FloatTok{310.0}\NormalTok{, solute\_count}\OperatorTok{=}\DecValTok{2}\NormalTok{)}
\BuiltInTok{print}\NormalTok{(}\BuiltInTok{round}\NormalTok{(pi\_pa))  }\CommentTok{\# about 7.7e5 Pa for about isotonic NaCl (order of several atm)}
\end{Highlighting}
\end{Shaded}

\begin{quote}
\textbf{Clinical / systems note:} Nephrogenic diabetes insipidus from
\emph{AQP2} mutations illustrates how loss of selective water
permeability breaks the kidney's ability to concentrate urine despite
normal vasopressin signalling --- a channel-level failure of the osmotic
water flux you model with \texorpdfstring{\ensuremath{\pi}}{pi} and tonicity.
\end{quote}

\begin{center}\rule{0.5\linewidth}{0.5pt}\end{center}

\subsection{Summary}\label{summary-4}

\begin{itemize}
\tightlist
\item
  Water's polarity and H-bonding network underlie its high heat
  capacity, heat of vaporisation, surface tension, and solvent
  properties.
\item
  Hydrophilic/hydrophobic dichotomy drives membrane bilayer formation,
  protein folding, and micelle structure.
\item
  The hydration shell stabilises ions in solution; water's high
  dielectric constant screens electrostatic interactions.
\item
  Osmotic pressure (\(\pi = iCRT\)) governs water movement across
  semipermeable membranes with profound cellular consequences.
\item
  Fick's First Law (\(J = -D \cdot dC/dx\)) quantifies passive
  solute/gas diffusion; slow diffusion over mm distances necessitates
  circulatory systems.
\item
  Aquaporins are selective water channels critical for renal function,
  brain water homeostasis, and lens transparency; mutations cause
  nephrogenic diabetes insipidus.
\item
  Colligative properties (freezing point depression, boiling point
  elevation) are exploited by antifreeze organisms.
\item
  Water activity determines microbial growth limits and is the basis of
  traditional food preservation methods.
\item
  \textbf{Connections:} See Unit II (membrane transport and aquaporin
  structure), Unit IX (renal concentration gradients and ADH), and Unit
  VIII (transpiration and \hyperref[gl:xylem]{\textbf{xylem}} tension).
\end{itemize}

\subsection{Key Terms}\label{key-terms-4}

\begin{itemize}
\tightlist
\item
  \textbf{Hydrogen bond:} Weak electrostatic attraction between H bonded
  to N/O/F and another electronegative atom
\item
  \textbf{Hydrophilic:} Water-loving; molecules with polar or charged
  groups
\item
  \textbf{Hydrophobic effect:} Entropy-driven aggregation of nonpolar
  molecules in water
\item
  \textbf{Clathrate:} Ordered cage of water molecules around a nonpolar
  solute
\item
  \textbf{Amphipathic:} Molecules with both hydrophilic and hydrophobic
  regions
\item
  \textbf{Osmosis:} Net water movement across a semipermeable membrane
  down its concentration gradient
\item
  \textbf{Osmotic pressure (\texorpdfstring{\ensuremath{\pi}}{pi}):} Pressure required to prevent osmotic
  water flow
\item
  \textbf{Van 't Hoff equation:} \(\pi = iCRT\); relates osmotic
  pressure to solute concentration
\item
  \textbf{Tonicity:} Relative solute concentration affecting cell volume
\item
  \textbf{Diffusion:} Net movement of molecules down a concentration
  gradient
\item
  \textbf{Fick's First Law:} \(J = -D \cdot dC/dx\); relates flux to
  concentration gradient
\item
  \textbf{Diffusion coefficient (\(D\)):} Measure of molecular mobility
  in a medium
\item
  \textbf{Aquaporin:} Integral membrane protein forming a selective
  water channel
\item
  \textbf{Colligative property:} Solution property depending on particle
  number, not identity
\item
  \textbf{Water activity (\(a_w\)):} Effective concentration of water in
  a solution; determines microbial growth
\item
  \textbf{Antifreeze protein:} Protein that inhibits ice crystal growth
  via thermal hysteresis
\item
  \textbf{Dielectric constant (\(\varepsilon\)):} Measure of a medium's
  ability to screen electrostatic interactions
\end{itemize}

\subsection{Review Questions}\label{review-questions-4}

\begin{enumerate}
\def\labelenumi{\arabic{enumi}.}
\tightlist
\item
  Draw the hydrogen bonding network around a single water molecule. How
  many H-bonds can one water molecule form, and why?
\item
  Explain why water's heat of vaporisation is biologically important.
  Calculate the volume of sweat that must evaporate to remove 500 kJ of
  heat from the body.
\item
  A red blood cell is placed in a solution of 0.45\% NaCl (\(i = 2\), MW
  = 58.44). Calculate the osmolarity of this solution and predict what
  will happen to the cell.
\item
  Using Fick's First Law, explain why alveolar walls are about 0.2 \texorpdfstring{\ensuremath{\mu}}{mu}m
  thick. What would happen to gas exchange efficiency if they were 10 \texorpdfstring{\ensuremath{\mu}}{mu}m
  thick?
\item
  An aquaporin channel transports \(3 \times 10^9\) water molecules per
  second. If a red blood cell contains 200,000 AQP1 channels, how many
  moles of water pass through the cell membrane per second?
\item
  Explain the molecular mechanism by which aquaporins allow water
  passage but exclude protons (H\(^+\) / H\(_3\)O\(^+\)).
\item
  A food scientist wants to reduce the water activity of a fruit jam
  from 0.95 to 0.80 by adding sucrose. Qualitatively, how does this
  protect against microbial spoilage?
\item
  The wood frog (\emph{Rana sylvatica}) survives freezing by
  accumulating glucose in its cells. Explain, using colligative
  properties, how intracellular glucose prevents lethal ice crystal
  formation inside cells even while extracellular ice forms.
\item
  Compare the hydration shells of Na\(^+\) and Mg\(^{2+}\). Which ion
  has a more tightly bound hydration shell, and why? How does this
  affect the biological roles of each ion?
\item
  Derive the relationship between the mean diffusion distance and time
  (\(\bar{x} = \sqrt{2Dt}\)) starting from the concept of a random walk.
  Explain why this relationship makes circulatory systems necessary for
  organisms larger than about 1 mm.
\item
  Using \texttt{osmotic\_pressure}, compare the osmotic pressure (Pa) of
  0.30 M glucose (\(i=1\)) vs.~0.15 M NaCl (\(i=2\)) at 310 K. Which is
  closer to human plasma tonicity?
\item
  A child with nephrogenic diabetes insipidus excretes large volumes of
  dilute urine. Link this \hyperref[gl:phenotype]{\textbf{phenotype}} to
  aquaporin-2 trafficking in collecting ducts and the role of
  vasopressin (ADH) in increasing epithelial water permeability.
\item
  Using \cref{eq:unit_I_dissolution_dG}, predict whether n-octanol (a
  hydrocarbon-like alcohol) is more soluble in water at 25 \(^\circ\)C
  or 60 \(^\circ\)C. Assume \(\Delta H_{\text{soln}} \approx 0\) and the
  dissolution is dominated by the unfavourable entropy of cage
  formation. How does this relate to drug partition coefficients
  (\(\log P\))?
\item
  Explain how the Grotthuss mechanism allows ATP synthase to translocate
  about 10\(^4\) protons per second through its F\(_o\) ring
  without bulk water flow. What would happen if a mutation introduced a
  single Pro residue that broke the proton wire?
\item
  Cryopreservation protocols for human embryos use vitrification (rapid
  cooling with high cryoprotectant concentrations) rather than slow
  freezing. Using your understanding of ice nucleation and water
  activity, explain why vitrification is preferred for delicate cells.
\item
  Use \cref{eq:unit_I_water_activity} to quantify the difference in
  chemical potential of water in pure water (\(a_w = 1\)) versus a 1 M
  sucrose solution (\(a_w \approx 0.98\)) at 310 K. Compare your answer
  to the corresponding osmotic pressure from \cref{eq:unit_I_vant_hoff}.
  Are they consistent?
\end{enumerate}

\begin{center}\rule{0.5\linewidth}{0.5pt}\end{center}

\subsection{Further Reading and Source
Notes}\label{further-reading-and-source-notes-4}

\begin{itemize}
\tightlist
\item
  Henderson (1913). \emph{The Fitness of the Environment}. Macmillan.
\item
  Eisenberg \& Kauzmann (1969). \emph{The Structure and Properties of
  Water}. Oxford University Press.
\item
  Tanford (1980). \emph{The Hydrophobic Effect: Formation of Micelles
  and Biological Membranes}. Wiley.
\item
  Ball (2008). Water as an active constituent in cell biology.
  \emph{Chemical Reviews}, 108.
\item
  Pace (2001). The comprehensive nature of biochemistry.
  \emph{Proceedings of the National Academy of Sciences}, 98.
\item
  Franks (2000). \emph{Water: A Matrix of Life}. Royal Society of
  Chemistry.
\end{itemize}

\begin{center}\rule{0.5\linewidth}{0.5pt}\end{center}

\subsubsection{Companion Source Module}\label{companion-source-module-4}

\textbf{Water --- The Molecule of Life} should leave a reproducible
trail from a biological claim to the code, figure, diagram, or
paper-based activity that can test it. Use the surfaces below to inspect
the chapter's assumptions, rerun the relevant model, or compare the
manuscript explanation with companion labs and figures.

{\def\LTcaptype{none} % do not increment counter
\begin{longtable}[]{@{}
  >{\raggedright\arraybackslash}p{(\linewidth - 2\tabcolsep) * \real{0.5000}}
  >{\raggedright\arraybackslash}p{(\linewidth - 2\tabcolsep) * \real{0.5000}}@{}}
\toprule\noalign{}
\begin{minipage}[b]{\linewidth}\raggedright
Surface
\end{minipage} & \begin{minipage}[b]{\linewidth}\raggedright
Use it for
\end{minipage} \\
\midrule\noalign{}
\endhead
\bottomrule\noalign{}
\endlastfoot
\breaktt{src/biology/cell/cell\_biology.py} (\texttt{osmotic\_pressure},
\texttt{diffusion\_flux}) & Translate water potential, solute gradients,
and diffusion into quantitative predictions. \\
\breaktt{src/biology/botany/botany.py} (\texttt{water\_potential},
\breaktt{transpiration\_flux}) & Carry water chemistry into plant
transport and drought-response scenarios. \\
\breaktt{src/visualization/plots.py}
(\breaktt{plot\_light\_response\_curve}) & Use graph reading practice for
environmental-response curves with clear axes and units. \\
\end{longtable}
}

\textbf{Reproducibility check:} state temperature, solute identity,
concentration, and membrane permeability before generalising a
water-property claim. \textbf{Cross-reference:} connect molecular water
properties to \cref{sec:unit_VIII_plant_structure_and_water} and
\cref{sec:unit_II_membrane_transport}.

\newpage

\section{Biological Macromolecules}\label{biological-macromolecules}

\label{sec:unit_I_macromolecules}

\begin{quote}
\textbf{Ch 3} \texorpdfstring{\ensuremath{\cdot}}{.} Level 2/3 \texorpdfstring{\ensuremath{\cdot}}{.} 55 min read \texorpdfstring{\ensuremath{\cdot}}{.} 75 min lecture \texorpdfstring{\ensuremath{\cdot}}{.}
Prerequisites: \cref{sec:unit_I_atoms_molecules},
\cref{sec:unit_I_water_and_life}
\end{quote}

\subsection{Learning Objectives}\label{learning-objectives-5}

\begin{enumerate}
\def\labelenumi{\arabic{enumi}.}
\tightlist
\item
  Classify and compare the four classes of biological macromolecules.
\item
  Explain condensation (dehydration) and hydrolysis reactions as the
  basis of polymer chemistry.
\item
  Describe the structural levels of
  \hyperref[gl:protein]{\textbf{protein}}s and relate each level to
  function, including modern concepts such as intrinsically disordered
  proteins.
\item
  Explain nucleic acid structure, including A-form, B-form, and Z-form
  DNA, as well as RNA secondary structures.
\item
  Distinguish the structures and roles of carbohydrates and lipids,
  including glycoproteins and sphingolipids.
\item
  Describe ribozymes, their discovery, and their significance for the
  RNA World hypothesis.
\end{enumerate}

\subsubsection{Study Blueprint}\label{study-blueprint-5}

\begin{itemize}
\tightlist
\item
  \textbf{Big idea:} Macromolecule structure links monomer chemistry to
  biological function and information storage.
\item
  \textbf{Core concepts:} polymers, dehydration synthesis, protein
  structure, nucleic acids.
\item
  \textbf{Framework alignment:} Vision \& Change: Structure and
  function, Pathways and transformations of energy and matter; AP
  Biology: Energetics, Systems Interactions; NGSS-style topics: Matter
  and Energy in Organisms and Ecosystems, Structure and Function.
\item
  \textbf{Model or quantitative lens:} Polymerization energetics and
  sequence-to-structure reasoning.
\item
  \textbf{Data skill:} Classify molecules from structural evidence
  rather than names alone.
\item
  \textbf{Practice cadence:} Concept Explanation, Statistical Tests and
  Data Analysis, Argumentation.
\item
  \textbf{Common misconception to repair:} Structure is not decoration;
  small chemical changes can redirect function and recognition.
\item
  \textbf{Primary lab:} \cref{sec:lab_unit_I_macromolecules}.
\item
  \textbf{Question bank:} \cref{sec:q_unit_I_macromolecules}.
\item
  \textbf{Transfer task:} Explain how a mutation, lipid substitution, or
  glycosylation change propagates to phenotype.
\item
  \textbf{Bridge to computation:}
  \breaktt{biology.genetics.genetics.translate\_mrna}.
\end{itemize}

\begin{center}\rule{0.5\linewidth}{0.5pt}\end{center}

\begin{quote}
\textbf{Opening Vignette: The Protein That Won Two Nobel Prizes}

In 1955, Frederick Sanger published the complete amino acid sequence of
bovine insulin: 51 amino acids in two chains (A and B) linked by two
disulfide bonds. It had taken him 10 years and was considered one of the
most technically difficult achievements in biochemistry at the time. The
revelation was not just the sequence itself, but what it implied:
proteins are not random polymers --- they are precisely specified
sequences, encoded in \hyperref[gl:gene]{\textbf{gene}}s, determining
unique three-dimensional structures. This work earned Sanger his first
Nobel Prize in Chemistry (1958).

The same molecule --- insulin --- would later become the first human
protein produced by recombinant DNA technology (1982), when the human
insulin gene was inserted into \emph{E. coli} to manufacture
pharmaceutical insulin for diabetics. Before recombinant insulin, most
medical insulin came from pig and cow pancreas, with shortages,
immunological reactions, and ethical concerns. Today, nearly most 500
million diabetics worldwide who require insulin use a product whose
existence depends on understanding the four classes of biological
macromolecules this chapter explores.

\emph{Primary source: Sanger, F. \& Thompson, E. O. P. (1953). The
amino-acid sequence in the glycyl chain of insulin. Biochem. J., 53(3),
353--374.}
\end{quote}

\begin{center}\rule{0.5\linewidth}{0.5pt}\end{center}

Biological macromolecules fall into four major classes:
\textbf{carbohydrates}, \textbf{lipids}, \textbf{proteins}, and
\textbf{nucleic acids}. The first three are polymers of small repeating
units (monomers); lipids are an important exception --- they form
assemblies through noncovalent hydrophobic interactions rather than
covalent polymerisation.

\begin{figure}[htbp]
\centering
\includegraphics[width=0.82\linewidth,height=4.2in,keepaspectratio]{../figures/mermaid_inline/inline_mermaid_0021_9eb74518a708.png}
\caption{The four biological macromolecule families - carbohydrates, lipids, proteins, and nucleic acids - with their monomers, characteristic bonds, and cellular roles.}
\end{figure}

\subsubsection{Condensation and
Hydrolysis}\label{condensation-and-hydrolysis}

Most biological polymers are built by \textbf{condensation (dehydration)
reactions}, in which two monomers are joined with the release of one
water molecule per bond:

\[\text{Monomer}_1\text{--OH} + \text{H--Monomer}_2 \rightarrow \text{Monomer}_1\text{--Monomer}_2 + \text{H}_2\text{O} \tag{3.1} \label{eq:unit_I_macromolecules_item_1}\]

The reverse, \textbf{hydrolysis}, uses water to cleave the bond. These
reactions require \hyperref[gl:enzyme]{\textbf{enzyme}} catalysis (see
\cref{sec:unit_I_enzymes_and_kinetics}) and energy input for
condensation.

\textbf{\hyperref[gl:thermodynamics]{Thermodynamics} of bond formation:}

{\def\LTcaptype{none} % do not increment counter
\begin{longtable}[]{@{}
  >{\raggedright\arraybackslash}p{(\linewidth - 4\tabcolsep) * \real{0.1250}}
  >{\raggedright\arraybackslash}p{(\linewidth - 4\tabcolsep) * \real{0.5833}}
  >{\raggedright\arraybackslash}p{(\linewidth - 4\tabcolsep) * \real{0.2917}}@{}}
\toprule\noalign{}
\begin{minipage}[b]{\linewidth}\raggedright
Bond Type
\end{minipage} & \begin{minipage}[b]{\linewidth}\raggedright
\(\Delta G^{\circ'}\) of Hydrolysis (kJ/mol)
\end{minipage} & \begin{minipage}[b]{\linewidth}\raggedright
Required Energy Source
\end{minipage} \\
\midrule\noalign{}
\endhead
\bottomrule\noalign{}
\endlastfoot
Glycosidic & --15 to --20 & UDP-glucose (activated sugar) \\
Peptide & --8 to --16 & GTP (ribosomal
\hyperref[gl:translation]{\textbf{translation}}) \\
Phosphodiester & --25 & NTP (DNA/RNA polymerase) \\
Ester (lipid) & --20 & Acyl-CoA (activated fatty acid) \\
\end{longtable}
}

In each case, the monomers must first be ``activated'' by coupling to a
high-energy carrier molecule. This coupling makes the condensation
reaction thermodynamically favourable.

\subsubsection{Worked Example: The Thermodynamics of
Polymerisation}\label{worked-example-the-thermodynamics-of-polymerisation}

\textbf{Problem:} The hydrolysis of a
\hyperref[gl:peptide-bond]{\textbf{peptide bond}} releases approximately
\(\Delta G^{\circ'} = -10 \text{ kJ/mol}\). 1. What is the standard free
energy change (\(\Delta G^{\circ'}\)) for the condensation reaction to
form a single peptide bond? 2. Why does the cell require the hydrolysis
of high-energy bonds (like GTP, where
\(\Delta G^{\circ'} \approx -30 \text{ kJ/mol}\)) to drive this reaction
during protein synthesis in the
\hyperref[gl:ribosome]{\textbf{ribosome}}?

\textbf{Solution:}

\begin{enumerate}
\def\labelenumi{\arabic{enumi}.}
\item
  \textbf{Calculate the condensation \(\Delta G^{\circ'}\):} Because
  condensation is the exact reverse of hydrolysis, the standard free
  energy change has the same magnitude but the opposite sign:
  \[ \Delta G^{\circ'}_{\text{condensation}} = -(\Delta G^{\circ'}_{\text{hydrolysis}}) = -(-10 \text{ kJ/mol}) = +10 \text{ kJ/mol}  \label{eq:unit_I_macromolecules_item_2}\]

  As \(\Delta G^{\circ'} > 0\), the formation of a peptide bond is
  thermodynamically unfavourable (endergonic) and will not occur
  spontaneously.
\item
  \textbf{Calculate the coupled reaction thermodynamics:} To drive an
  endergonic reaction forward, the cell couples it to a highly exergonic
  reaction. During translation elongation, the ribosome couples the
  formation of a peptide bond to the hydrolysis of GTP down to GDP +
  P\(_i\).

  Summing the free energy changes:
  \[ \Delta G^{\circ'}_{\text{total}} = \Delta G^{\circ'}_{\text{condensation}} + \Delta G^{\circ'}_{\text{GTP\_hydrolysis}}  \label{eq:unit_I_macromolecules_item_3}\]

  \[ \Delta G^{\circ'}_{\text{total}} = (+10 \text{ kJ/mol}) + (-30 \text{ kJ/mol}) = -20 \text{ kJ/mol}  \label{eq:unit_I_macromolecules_item_4}\]
\end{enumerate}

By coupling the unfavourable peptide bond formation (+10 kJ/mol) to the
favourable hydrolysis of GTP (--30 kJ/mol), the overall coupled reaction
becomes highly favourable (--20 kJ/mol), driving protein synthesis
forward with near-perfect unidirectionality.

\begin{quote}
\textbf{Concept Check 1:} Hydrolysis reactions are thermodynamically
spontaneous (\(\Delta G < 0\)), yet polymers are kinetically stable in
cells. What prevents the spontaneous hydrolysis of your DNA, proteins,
and carbohydrates?
\end{quote}

\begin{center}\rule{0.5\linewidth}{0.5pt}\end{center}

\subsection{Carbohydrates}\label{carbohydrates}

\subsubsection{Monosaccharides}\label{monosaccharides}

The simplest carbohydrates are \textbf{monosaccharides} with the
empirical formula (CH\(_2\)O)\(_n\). They are classified by: -
\textbf{Number of carbons:} trioses (3C), tetroses (4C), pentoses (5C),
hexoses (6C) - \textbf{Functional group:} aldoses (aldehyde) or ketoses
(ketone)

Common examples:

{\def\LTcaptype{none} % do not increment counter
\begin{longtable}[]{@{}
  >{\raggedright\arraybackslash}p{(\linewidth - 8\tabcolsep) * \real{0.2381}}
  >{\raggedright\arraybackslash}p{(\linewidth - 8\tabcolsep) * \real{0.3333}}
  >{\raggedright\arraybackslash}p{(\linewidth - 8\tabcolsep) * \real{0.0476}}
  >{\raggedright\arraybackslash}p{(\linewidth - 8\tabcolsep) * \real{0.1905}}
  >{\raggedright\arraybackslash}p{(\linewidth - 8\tabcolsep) * \real{0.1905}}@{}}
\toprule\noalign{}
\begin{minipage}[b]{\linewidth}\raggedright
Sugar
\end{minipage} & \begin{minipage}[b]{\linewidth}\raggedright
Formula
\end{minipage} & \begin{minipage}[b]{\linewidth}\raggedright
n
\end{minipage} & \begin{minipage}[b]{\linewidth}\raggedright
Type
\end{minipage} & \begin{minipage}[b]{\linewidth}\raggedright
Role
\end{minipage} \\
\midrule\noalign{}
\endhead
\bottomrule\noalign{}
\endlastfoot
Glyceraldehyde & C\(_3\)H\(_6\)O\(_3\) & 3 & Aldotriose & Central
metabolite (\hyperref[gl:glycolysis]{\textbf{glycolysis}}) \\
Dihydroxyacetone & C\(_3\)H\(_6\)O\(_3\) & 3 & Ketotriose & Glycolysis
intermediate \\
Ribose & C\(_5\)H\(_{10}\)O\(_5\) & 5 & Aldopentose & RNA backbone, ATP,
NAD\(^+\) \\
Deoxyribose & C\(_5\)H\(_{10}\)O\(_4\) & 5 & Aldopentose & DNA
backbone \\
Glucose & C\(_6\)H\(_{12}\)O\(_6\) & 6 & Aldohexose & Primary fuel
molecule \\
Galactose & C\(_6\)H\(_{12}\)O\(_6\) & 6 & Aldohexose & Lactose
component \\
Fructose & C\(_6\)H\(_{12}\)O\(_6\) & 6 & Ketohexose & Fruit sugar,
sweetest natural sugar \\
Mannose & C\(_6\)H\(_{12}\)O\(_6\) & 6 & Aldohexose & Glycoprotein
glycans \\
\end{longtable}
}

Glucose, galactose, and fructose are \textbf{structural isomers} ---
identical molecular formulas, different atom arrangements. In solution,
glucose exists primarily in the \textbf{ring form} (pyranose,
about 99\%) rather than the open-chain aldehyde form. The two
ring anomers, \texorpdfstring{\ensuremath{\alpha}}{alpha}-D-glucose and \texorpdfstring{\ensuremath{\beta}}{beta}-D-glucose, differ primarily in the
orientation of the C1 hydroxyl group, but this distinction has enormous
biochemical significance --- enzymes discriminate exquisitely between
them.

\subsubsection{Disaccharides and
Polysaccharides}\label{disaccharides-and-polysaccharides}

Two monosaccharides join via a \textbf{glycosidic bond} to form a
disaccharide:

{\def\LTcaptype{none} % do not increment counter
\begin{longtable}[]{@{}
  >{\raggedright\arraybackslash}p{(\linewidth - 6\tabcolsep) * \real{0.2927}}
  >{\raggedright\arraybackslash}p{(\linewidth - 6\tabcolsep) * \real{0.1951}}
  >{\raggedright\arraybackslash}p{(\linewidth - 6\tabcolsep) * \real{0.0976}}
  >{\raggedright\arraybackslash}p{(\linewidth - 6\tabcolsep) * \real{0.4146}}@{}}
\toprule\noalign{}
\begin{minipage}[b]{\linewidth}\raggedright
Disaccharide
\end{minipage} & \begin{minipage}[b]{\linewidth}\raggedright
Monomers
\end{minipage} & \begin{minipage}[b]{\linewidth}\raggedright
Bond
\end{minipage} & \begin{minipage}[b]{\linewidth}\raggedright
Biological source
\end{minipage} \\
\midrule\noalign{}
\endhead
\bottomrule\noalign{}
\endlastfoot
Sucrose & Glucose + Fructose & \(\alpha(1\to2)\beta\) & Plants
(transport sugar) \\
Lactose & Glucose + Galactose & \(\beta(1\to4)\) & Milk \\
Maltose & Glucose + Glucose & \(\alpha(1\to4)\) & Starch digestion \\
Trehalose & Glucose + Glucose & \(\alpha(1\to1)\alpha\) & Insect
haemolymph, desiccation protectant \\
Cellobiose & Glucose + Glucose & \(\beta(1\to4)\) & Cellulose
degradation product \\
\end{longtable}
}

\textbf{Polysaccharides} are long chains of monosaccharides serving
structural or storage functions:

\begin{itemize}
\tightlist
\item
  \textbf{Starch (amylose + amylopectin):} \(\alpha(1\to4)\) linkages,
  branching via \(\alpha(1\to6)\) at intervals (about 25
  residues for amylopectin). Storage in plants.
\item
  \textbf{Glycogen:} Like amylopectin but more heavily branched (every
  about 10 residues). Glucose storage in animal liver
  (about 100 g) and muscle (about 400 g).
\item
  \textbf{Cellulose:} \(\beta(1\to4)\) glucose. Structural component of
  plant cell walls. Humans cannot digest it (no beta-glucosidase), but
  gut \hyperref[gl:microbiota]{\textbf{microbiota}} can; cellulose is
  dietary fibre.
\item
  \textbf{Chitin:} \(\beta(1\to4)\) N-acetylglucosamine. Structural
  component of fungal cell walls and arthropod exoskeletons --- the
  second most abundant biopolymer on Earth after cellulose.
\item
  \textbf{Hyaluronic acid:} alternating \(\beta(1\to4)\) glucuronic acid
  and \(\beta(1\to3)\) N-acetylglucosamine. Major component of
  extracellular matrix and synovial fluid.
\end{itemize}

The difference between starch and cellulose --- both glucose polymers
--- lies solely in the glycosidic bond geometry (\texorpdfstring{\ensuremath{\alpha}}{alpha} vs.~\texorpdfstring{\ensuremath{\beta}}{beta}). Yet this
single difference makes cellulose rigid and indigestible versus starch
flexible and readily hydrolysed.

\begin{figure}[htbp]
\centering
\includegraphics[width=0.82\linewidth,height=4.2in,keepaspectratio]{../figures/mermaid_inline/inline_mermaid_0022_817dc0a53fed.png}
\caption{Carbohydrate digestion and absorption. Carbohydrate digestion and absorption pathway. Note the key distinction: alpha-linked polysaccharides (starch) are digested by human enzymes, while beta-linked cellulose passes to the colon where microbial fermentation produces short-chain fatty acids.}
\end{figure}

\begin{quote}
\textbf{Clinical Connection: Lactose Intolerance}

Approximately 65--75\% of the global adult population has reduced
expression of lactase (the brush border enzyme that cleaves lactose)
after weaning --- the ancestral state called \textbf{lactase
non-persistence}. Undigested lactose reaches the colon, where bacterial
fermentation produces H\(_2\), CO\(_2\), and short-chain fatty acids,
causing bloating, flatulence, and diarrhoea. \textbf{Lactase
persistence} (the ability to digest lactose as an adult) evolved
independently in pastoral populations in Europe (about 7,500
years ago) and East Africa (about 3,000 years ago) due to
strong selection pressure from dairy farming. This is one of the
clearest examples of recent human evolution driven by cultural
practices.
\end{quote}

\begin{quote}
\textbf{Concept Check 2:} Glycogen is more heavily branched than
amylopectin (\texorpdfstring{\ensuremath{\alpha}}{alpha}1-\textgreater6 branch every about 10 residues
vs.~every about 25 residues). How does this increased branching
benefit an animal that needs rapid glucose mobilisation during a ``fight
or flight'' response?
\end{quote}

\subsubsection{The Glycosidic Bond: \texorpdfstring{\ensuremath{\alpha}}{alpha}
vs.~\texorpdfstring{\ensuremath{\beta}}{beta}}\label{the-glycosidic-bond-ux3b1-vs.-ux3b2}

The single greatest source of carbohydrate structural diversity is the
\textbf{glycosidic bond} geometry. When a hemiacetal hydroxyl (at the
anomeric carbon, C1 of the cyclic sugar) condenses with a hydroxyl of
another sugar, the resulting linkage can be either \textbf{\texorpdfstring{\ensuremath{\alpha}}{alpha}} (the OH/OR
on C1 points down in the standard Haworth projection, \emph{trans} to
the C6 substituent) or \textbf{\texorpdfstring{\ensuremath{\beta}}{beta}} (OH/OR points up, \emph{cis} to C6).

{\def\LTcaptype{none} % do not increment counter
\begin{longtable}[]{@{}
  >{\raggedright\arraybackslash}p{(\linewidth - 8\tabcolsep) * \real{0.1000}}
  >{\raggedright\arraybackslash}p{(\linewidth - 8\tabcolsep) * \real{0.1143}}
  >{\raggedright\arraybackslash}p{(\linewidth - 8\tabcolsep) * \real{0.3143}}
  >{\raggedright\arraybackslash}p{(\linewidth - 8\tabcolsep) * \real{0.1714}}
  >{\raggedright\arraybackslash}p{(\linewidth - 8\tabcolsep) * \real{0.3000}}@{}}
\toprule\noalign{}
\begin{minipage}[b]{\linewidth}\raggedright
Linkage
\end{minipage} & \begin{minipage}[b]{\linewidth}\raggedright
Geometry
\end{minipage} & \begin{minipage}[b]{\linewidth}\raggedright
Polysaccharide example
\end{minipage} & \begin{minipage}[b]{\linewidth}\raggedright
Conformation
\end{minipage} & \begin{minipage}[b]{\linewidth}\raggedright
Digestible by humans?
\end{minipage} \\
\midrule\noalign{}
\endhead
\bottomrule\noalign{}
\endlastfoot
\(\alpha(1\to4)\) (Glc--Glc) & bent & amylose, glycogen & helix & yes
(amylase) \\
\(\alpha(1\to6)\) & branchpoint & amylopectin, glycogen branches &
branched helix & yes (isomaltase) \\
\(\beta(1\to4)\) (Glc--Glc) & straight & cellulose & extended ribbons &
\textbf{no} (no \texorpdfstring{\ensuremath{\beta}}{beta}-glucosidase) \\
\(\beta(1\to3)\) & extended & callose, paramylon (algae) & linear &
no \\
\(\beta(1\to4)\) (GlcNAc--GlcNAc) & straight & chitin & extended & no \\
\(\alpha(1\to2)\) (Glc--Fru) & mixed & sucrose & non-reducing & yes
(sucrase) \\
\(\beta(1\to4)\) (Gal--Glc) & mixed & lactose & linear & yes (lactase)
until weaning \\
\end{longtable}
}

The biological consequences of this single stereochemical choice are
profound: \texorpdfstring{\ensuremath{\alpha}}{alpha}-linked glucose polymers (starch, glycogen) form helices that
pack as soft granules and are readily hydrolysed for energy storage.
\texorpdfstring{\ensuremath{\beta}}{beta}-linked glucose polymers (cellulose) form extended, hydrogen-bonded
sheets that pack into rigid microfibrils with tensile strength rivalling
steel by weight; nature uses them as structural material (plant cell
walls) precisely because they resist enzymatic attack.

\subsubsection{Branching, Linkages, and Polymer
Diversity}\label{branching-linkages-and-polymer-diversity}

Even within the same monosaccharide and same anomeric configuration,
\textbf{branching frequency} dramatically alters function. Glycogen
(animal storage) has \(\alpha(1\to6)\) branches every about 10
residues; amylopectin (plant storage) every about 25 residues.
Heavier branching gives glycogen many more non-reducing ends per
molecule --- and since both glycogen synthase and glycogen phosphorylase
work primarily at non-reducing ends, branching directly accelerates the
rate at which glucose can be added to or removed from the polymer.
During a ``fight or flight'' response, a hepatocyte mobilises glucose
from glycogen at about 10 mM/s; this rate would be unattainable
from a linear polymer with a single non-reducing end.

\subsubsection{Glycoproteins and
Proteoglycans}\label{glycoproteins-and-proteoglycans}

Many proteins are covalently modified with carbohydrate chains
(glycans):

\textbf{Glycoproteins} have relatively short, branched oligosaccharide
chains attached via: - \textbf{N-linked glycosylation:} to asparagine
(Asn) in the consensus sequence Asn-X-Ser/Thr (where X is any amino acid
except Pro). Built up co-translationally on dolichol phosphate in the
ER, then trimmed and modified in the Golgi. - \textbf{O-linked
glycosylation:} to serine (Ser) or threonine (Thr); built up
sequentially in the Golgi. - \textbf{Other linkages:} O-GlcNAc (a single
GlcNAc on Ser/Thr; common cytosolic/nuclear signal), C-mannose (on Trp),
GPI-anchor glycolipids (link some proteins to the membrane outer
leaflet).

{\def\LTcaptype{none} % do not increment counter
\begin{longtable}[]{@{}
  >{\raggedright\arraybackslash}p{(\linewidth - 4\tabcolsep) * \real{0.3043}}
  >{\raggedright\arraybackslash}p{(\linewidth - 4\tabcolsep) * \real{0.3478}}
  >{\raggedright\arraybackslash}p{(\linewidth - 4\tabcolsep) * \real{0.3478}}@{}}
\toprule\noalign{}
\begin{minipage}[b]{\linewidth}\raggedright
Feature
\end{minipage} & \begin{minipage}[b]{\linewidth}\raggedright
N-linked
\end{minipage} & \begin{minipage}[b]{\linewidth}\raggedright
O-linked
\end{minipage} \\
\midrule\noalign{}
\endhead
\bottomrule\noalign{}
\endlastfoot
Attachment residue & Asn (in Asn-X-Ser/Thr motif) & Ser, Thr \\
Linker sugar & GlcNAc & GalNAc (mucin-type), Gal, others \\
Site of attachment & ER (co-translational) & Golgi
(post-translational) \\
Initial structure & Dolichol-P-P-GlcNAc\(_2\)Man\(_9\)Glc\(_3\) & Single
sugar at a time \\
Examples & Antibodies, integrins, EPO, viral spike & Mucins, blood group
antigens, EGF repeats \\
Disease & Congenital disorders of glycosylation (CDG-I, II) & Tn
syndrome (cancer marker) \\
\end{longtable}
}

Glycoproteins are ubiquitous: antibodies (IgG), blood group antigens
(ABO), mucins (airway protection), and many secreted or surface proteins
carry glycans. Glycosylation affects protein folding, stability,
half-life, and cell-cell recognition. The SARS-CoV-2 spike protein is
heavily N-glycosylated, with at least 22 sites; these glycans shield the
protein from antibody neutralisation, and their removal is a strategy
under exploration for vaccine improvement.

\textbf{Proteoglycans} have very long, unbranched glycosaminoglycan
(GAG) chains:

{\def\LTcaptype{none} % do not increment counter
\begin{longtable}[]{@{}
  >{\raggedright\arraybackslash}p{(\linewidth - 6\tabcolsep) * \real{0.0909}}
  >{\raggedright\arraybackslash}p{(\linewidth - 6\tabcolsep) * \real{0.4242}}
  >{\raggedright\arraybackslash}p{(\linewidth - 6\tabcolsep) * \real{0.2424}}
  >{\raggedright\arraybackslash}p{(\linewidth - 6\tabcolsep) * \real{0.2424}}@{}}
\toprule\noalign{}
\begin{minipage}[b]{\linewidth}\raggedright
GAG
\end{minipage} & \begin{minipage}[b]{\linewidth}\raggedright
Repeating Unit
\end{minipage} & \begin{minipage}[b]{\linewidth}\raggedright
Location
\end{minipage} & \begin{minipage}[b]{\linewidth}\raggedright
Function
\end{minipage} \\
\midrule\noalign{}
\endhead
\bottomrule\noalign{}
\endlastfoot
Hyaluronic acid & GlcUA-GlcNAc & Synovial fluid, vitreous humour &
Lubrication, shock absorption \\
Chondroitin sulfate & GlcUA-GalNAc(SO\(_4\)) & Cartilage, bone &
Compression resistance \\
Heparan sulfate & GlcUA/IdoUA-GlcNAc(SO\(_4\)) & Cell surfaces, basement
membrane & Growth factor binding \\
Keratan sulfate & Gal-GlcNAc(SO\(_4\)) & Cornea, cartilage &
Transparency, hydration \\
\end{longtable}
}

The highly negatively charged GAGs attract enormous amounts of water
(via osmotic pressure), creating a hydrated gel that resists compression
--- this is why cartilage can support the body's weight.

\begin{center}\rule{0.5\linewidth}{0.5pt}\end{center}

\subsection{Lipids}\label{lipids}

Lipids are not polymers but are grouped together by their
\textbf{hydrophobicity} --- they do not form regular covalent polymers,
yet their self-assembly into bilayers, vesicles, and membranes is
central to cellular life.

\subsubsection{Fatty Acids}\label{fatty-acids}

Fatty acids are long-chain carboxylic acids. Food-relevant examples:

{\def\LTcaptype{none} % do not increment counter
\begin{longtable}[]{@{}
  >{\raggedright\arraybackslash}p{(\linewidth - 8\tabcolsep) * \real{0.2703}}
  >{\raggedright\arraybackslash}p{(\linewidth - 8\tabcolsep) * \real{0.1351}}
  >{\raggedright\arraybackslash}p{(\linewidth - 8\tabcolsep) * \real{0.3243}}
  >{\raggedright\arraybackslash}p{(\linewidth - 8\tabcolsep) * \real{0.1081}}
  >{\raggedright\arraybackslash}p{(\linewidth - 8\tabcolsep) * \real{0.1622}}@{}}
\toprule\noalign{}
\begin{minipage}[b]{\linewidth}\raggedright
Fatty Acid
\end{minipage} & \begin{minipage}[b]{\linewidth}\raggedright
Chain
\end{minipage} & \begin{minipage}[b]{\linewidth}\raggedright
Double Bonds
\end{minipage} & \begin{minipage}[b]{\linewidth}\raggedright
Name
\end{minipage} & \begin{minipage}[b]{\linewidth}\raggedright
Source
\end{minipage} \\
\midrule\noalign{}
\endhead
\bottomrule\noalign{}
\endlastfoot
CH\(_3\)(CH\(_2\))\(_{14}\)COOH & 16:0 & 0 & Palmitic acid (saturated) &
Palm oil, meat \\
CH\(_3\)(CH\(_2\))\(_{16}\)COOH & 18:0 & 0 & Stearic acid (saturated) &
Cocoa butter \\
C18:1 \(\Delta^9\) & 18:1 & 1 & Oleic acid (monounsaturated) & Olive
oil \\
C18:2 \(\Delta^{9,12}\) & 18:2 & 2 & Linoleic acid (\texorpdfstring{\ensuremath{\omega}}{omega}-6 PUFA) &
Vegetable oils \\
C18:3 \(\Delta^{9,12,15}\) & 18:3 & 3 & \texorpdfstring{\ensuremath{\alpha}}{alpha}-Linolenic acid (\texorpdfstring{\ensuremath{\omega}}{omega}-3 PUFA) &
Flaxseed, walnuts \\
C20:4 \(\Delta^{5,8,11,14}\) & 20:4 & 4 & Arachidonic acid (\texorpdfstring{\ensuremath{\omega}}{omega}-6) & Meat,
eggs \\
C20:5 \(\Delta^{5,8,11,14,17}\) & 20:5 & 5 & EPA (\texorpdfstring{\ensuremath{\omega}}{omega}-3) & Oily fish \\
C22:6 \(\Delta^{4,7,10,13,16,19}\) & 22:6 & 6 & DHA (\texorpdfstring{\ensuremath{\omega}}{omega}-3) & Fish oil,
brain \\
\end{longtable}
}

\textbf{Saturated} fatty acids have no double bonds and pack tightly ---
solid at room temperature (butter, lard). \textbf{Unsaturated} fatty
acids have \emph{cis} double bonds introducing kinks that prevent tight
packing --- liquid at room temperature (olive oil). \textbf{Trans} fats
have \emph{trans} double bonds that pack like saturated fats --- they
are artificial (from industrial hydrogenation) and have adverse
cardiovascular effects.

\textbf{Essential fatty acids:} Linoleic acid (\texorpdfstring{\ensuremath{\omega}}{omega}-6) and \texorpdfstring{\ensuremath{\alpha}}{alpha}-linolenic acid
(\texorpdfstring{\ensuremath{\omega}}{omega}-3) cannot be synthesised by humans (we lack \texorpdfstring{\ensuremath{\Delta}}{Delta}12 and \texorpdfstring{\ensuremath{\Delta}}{Delta}15 desaturases)
and must be obtained from the diet.

\begin{figure}[htbp]
\centering
\includegraphics[width=0.82\linewidth,height=4.2in,keepaspectratio]{../figures/mermaid_inline/inline_mermaid_0023_c4cbb747f919.png}
\caption{Lipid classes and roles. Classification of lipids. Unlike other macromolecules, lipids are classified by solubility (hydrophobicity) rather than by a common monomer structure. Note the diversity of roles: energy storage, membrane structure, signalling, and waterproofing.}
\end{figure}

\subsubsection{Glycerophospholipids}\label{glycerophospholipids}

Glycerophospholipids (also called phosphoglycerides) are the most
abundant membrane lipids. Their structure: glycerol backbone + 2 fatty
acid chains at positions sn-1 and sn-2 + phosphate group at sn-3 + polar
head group attached to the phosphate.

{\def\LTcaptype{none} % do not increment counter
\begin{longtable}[]{@{}
  >{\raggedright\arraybackslash}p{(\linewidth - 6\tabcolsep) * \real{0.1887}}
  >{\raggedright\arraybackslash}p{(\linewidth - 6\tabcolsep) * \real{0.2264}}
  >{\raggedright\arraybackslash}p{(\linewidth - 6\tabcolsep) * \real{0.2642}}
  >{\raggedright\arraybackslash}p{(\linewidth - 6\tabcolsep) * \real{0.3208}}@{}}
\toprule\noalign{}
\begin{minipage}[b]{\linewidth}\raggedright
Head Group
\end{minipage} & \begin{minipage}[b]{\linewidth}\raggedright
Phospholipid
\end{minipage} & \begin{minipage}[b]{\linewidth}\raggedright
Charge at \hyperref[gl:ph]{\textbf{pH}} 7
\end{minipage} & \begin{minipage}[b]{\linewidth}\raggedright
Location/Function
\end{minipage} \\
\midrule\noalign{}
\endhead
\bottomrule\noalign{}
\endlastfoot
Choline & Phosphatidylcholine (PC) & Zwitterionic & Outer leaflet; most
abundant \\
Ethanolamine & Phosphatidylethanolamine (PE) & Zwitterionic & Inner
leaflet; membrane curvature \\
Serine & Phosphatidylserine (PS) & Net negative & Inner leaflet; flipped
to outer in \hyperref[gl:apoptosis]{\textbf{apoptosis}} \\
Inositol & Phosphatidylinositol (PI) & Net negative & Inner leaflet;
PI3K signalling cascade \\
\end{longtable}
}

The \textbf{amphipathic} structure spontaneously forms bilayers in water
--- a thermodynamic consequence of the hydrophobic effect (see
\cref{sec:unit_I_water_and_life}).

\subsubsection{Quantitative Properties of the Lipid
Bilayer}\label{quantitative-properties-of-the-lipid-bilayer}

The ``fluid mosaic'' model of Singer and Nicolson (1972) proposed that
membranes are 2D fluids in which proteins and lipids diffuse laterally.
Quantitative measurements --- especially by fluorescence recovery after
photobleaching (FRAP) and single-particle tracking --- have refined this
picture into a richer, multi-domain landscape.

{\def\LTcaptype{none} % do not increment counter
\begin{longtable}[]{@{}
  >{\raggedright\arraybackslash}p{(\linewidth - 4\tabcolsep) * \real{0.3077}}
  >{\raggedright\arraybackslash}p{(\linewidth - 4\tabcolsep) * \real{0.5000}}
  >{\raggedright\arraybackslash}p{(\linewidth - 4\tabcolsep) * \real{0.1923}}@{}}
\toprule\noalign{}
\begin{minipage}[b]{\linewidth}\raggedright
Property
\end{minipage} & \begin{minipage}[b]{\linewidth}\raggedright
Typical value
\end{minipage} & \begin{minipage}[b]{\linewidth}\raggedright
Notes
\end{minipage} \\
\midrule\noalign{}
\endhead
\bottomrule\noalign{}
\endlastfoot
Bilayer thickness & 4--5 nm & Hydrophobic core about 3 nm;
head-group region about 1 nm each \\
Area per phospholipid & 60--70 \AA{}\(^2\) & Determined by NMR, X-ray on
monolayers \\
Lateral diffusion coefficient (\(D\), lipid) & 1--10 \texorpdfstring{\ensuremath{\mu}}{mu}m\(^2\)/s & A
lipid traverses 1 \texorpdfstring{\ensuremath{\mu}}{mu}m in about 1 s \\
Lateral diffusion coefficient (\(D\), integral protein) & 0.01--1
\texorpdfstring{\ensuremath{\mu}}{mu}m\(^2\)/s & 10--1000\(\times\) slower than lipids \\
Rotational correlation time (lipid) & about 100 ps & Acyl
chains ``wag'' on this timescale \\
Transbilayer flip-flop (phospholipid, spontaneous) & hours--days &
Energy barrier from polar head moving through hydrophobic core \\
Transbilayer flip-flop (catalysed by flippase) & seconds &
ATP-dependent; maintains lipid asymmetry \\
Bilayer surface tension & about 0 mN/m & Unlike air-water
interface; key for vesicle stability \\
Bending modulus (\texorpdfstring{\ensuremath{\kappa}}{kappa}) & 10--40 \(k_B T\) & Resistance to curvature;
depends on lipid shape, sterol content \\
Compressibility modulus (\(K_A\)) & 100--300 mN/m & Resistance to area
change \\
Permeability to water & 10--100 \texorpdfstring{\ensuremath{\mu}}{mu}m/s & Surprisingly high due to small
size; aquaporins boost about 10\(\times\) \\
Permeability to Na\(^+\) & \(10^{-12}\)--\(10^{-14}\) cm/s & Vanishingly
small without channels \\
\end{longtable}
}

\textbf{Lipid asymmetry.} The two leaflets of the plasma membrane have
very different compositions: the outer leaflet is enriched in PC,
sphingomyelin, and glycolipids; the inner leaflet in PE, PS, and PI.
This asymmetry is actively maintained by ATP-dependent flippases (move
lipids from outer to inner) and floppases (inner to outer). When PS
appears on the outer leaflet --- the result of activated scramblases
during apoptosis --- it serves as an ``eat-me'' signal for macrophages.
Loss of asymmetry is also a feature of activated platelets (where PS
exposure on the outer leaflet provides a surface for the coagulation
cascade).

\textbf{Lipid rafts.} Cholesterol and sphingolipids self-segregate into
liquid-ordered (\(L_o\)) microdomains within a sea of liquid-disordered
(\(L_d\)) phospholipids. These about 10--200 nm rafts
concentrate specific signalling proteins (GPI-anchored proteins,
Src-family kinases) and are now recognised as functional platforms for
receptor signalling, virus budding, and immune-cell synapse assembly.

\begin{quote}
\textbf{Concept Check 7:} A typical phospholipid diffuses
about 1 \texorpdfstring{\ensuremath{\mu}}{mu}m/s laterally in a fluid bilayer but flips between
leaflets primarily once every several hours unaided. Use the structure
of a phospholipid (charged head, hydrocarbon tails) and the energetics
of the hydrophobic effect to explain this 10\(^{12}\)-fold difference.
Why must cells expend ATP on flippases to break this asymmetry?
\end{quote}

\subsubsection{Sphingolipids}\label{sphingolipids}

Sphingolipids use \textbf{sphingosine} (rather than glycerol) as the
backbone. Ceramide (sphingosine + fatty acid) is the core structure:

\begin{itemize}
\tightlist
\item
  \textbf{Sphingomyelin:} ceramide + phosphocholine. Major component of
  the myelin sheath that insulates axons; critical for nerve impulse
  conduction velocity.
\item
  \textbf{Cerebrosides:} ceramide + single sugar (glucose or galactose).
  Abundant in brain.
\item
  \textbf{Gangliosides:} ceramide + complex oligosaccharide containing
  sialic acid (Neu5Ac). GM1 is the receptor for cholera toxin.
\end{itemize}

\begin{quote}
\textbf{Clinical Connection: Tay-Sachs Disease and Sphingolipid Storage}

Tay-Sachs disease results from deficiency of the lysosomal enzyme
hexosaminidase A, which normally degrades the ganglioside GM2. Without
this enzyme, GM2 accumulates in \hyperref[gl:neuron]{\textbf{neuron}}s,
causing progressive neurodegeneration. Symptoms appear at
about 6 months and include cherry-red macular spots, seizures,
and developmental regression. The disease is invariably fatal by age
3--5. Carrier frequency is about 1/30 in Ashkenazi Jewish
populations due to a \hyperref[gl:founder-effect]{\textbf{founder
effect}}. This is one of about 50 known \textbf{lysosomal
storage diseases}, most caused by deficient hydrolytic enzymes.
\end{quote}

\subsubsection{Sterols and Steroid
Hormones}\label{sterols-and-steroid-hormones}

\textbf{Cholesterol} has a four-ring steroid nucleus (three 6-carbon
rings + one 5-carbon ring) and a hydroxyl group at C3. In membranes,
cholesterol: - Inserts between phospholipids with its hydroxyl group
near the head groups and its rings among the acyl chains - At high
temperatures: reduces membrane fluidity by restricting phospholipid
movement - At low temperatures: prevents crystallisation by disrupting
ordered packing - Net effect: \textbf{homeoviscous regulation} that
maintains consistent membrane fluidity across a temperature range

Cholesterol is the precursor to most steroid
\hyperref[gl:hormone]{\textbf{hormone}}s: - \textbf{Glucocorticoids}
(\hyperref[gl:cortisol]{\textbf{cortisol}}): stress response,
anti-inflammatory - \textbf{Mineralocorticoids} (aldosterone):
sodium/potassium balance - \textbf{Androgens} (testosterone): male
sexual development - \textbf{Oestrogens} (oestradiol): female sexual
development - \textbf{Progestagens} (progesterone): pregnancy
maintenance - \textbf{Vitamin D\(_3\)} (cholecalciferol): calcium
absorption - \textbf{Bile acids} (cholic acid, chenodeoxycholic acid):
fat emulsification in the duodenum

\subsubsection{Waxes and Other Lipids}\label{waxes-and-other-lipids}

\begin{itemize}
\tightlist
\item
  \textbf{Waxes:} esters of long-chain fatty acids and long-chain
  alcohols; waterproofing (plant cuticle, bee honeycomb, whale head oil,
  ear wax)
\item
  \textbf{Eicosanoids:} signalling lipids derived from arachidonic acid
  (C20:4); include prostaglandins (inflammation, fever, pain),
  thromboxanes (platelet aggregation), and leukotrienes (immune
  response)
\item
  \textbf{Isoprenoids:} built from isoprene units (5C); include
  terpenes, carotenoids, ubiquinone (CoQ), and dolichol
\end{itemize}

\begin{quote}
\textbf{Concept Check 3:} Aspirin (acetylsalicylic acid) irreversibly
inhibits cyclooxygenase (COX), the enzyme that converts arachidonic acid
to prostaglandins and thromboxanes. Explain why aspirin reduces both
pain/inflammation and the risk of heart attacks.
\end{quote}

\begin{center}\rule{0.5\linewidth}{0.5pt}\end{center}

\subsection{Proteins}\label{proteins}

Proteins are the most structurally and functionally diverse
macromolecules. They serve as: - Structural elements (collagen, keratin,
\hyperref[gl:actin]{\textbf{actin}}, tubulin) - Enzymes (catalysts; see
\cref{sec:unit_I_enzymes_and_kinetics}) - Signalling molecules (growth
factors, receptors, G proteins) - Transport proteins (haemoglobin,
albumin, aquaporins) - Immune molecules (antibodies) - Motor proteins
(myosin, kinesin, dynein)

\subsubsection{Amino Acids}\label{amino-acids}

Most 20 canonical amino acids share a common backbone:

\begin{verbatim}
    NH₂
     |
H — C — COOH
     |
     R (side chain)
\end{verbatim}

The \textbf{\texorpdfstring{\ensuremath{\alpha}}{alpha}-amino group} (pKa about 9.0) and
\textbf{\texorpdfstring{\ensuremath{\alpha}}{alpha}-carboxyl group} (pKa about 2.0) flank the \texorpdfstring{\ensuremath{\alpha}}{alpha}-carbon,
to which the variable \textbf{R group} is attached. At physiological pH
7.4, the amino group is --NH\(_3^+\) and the carboxyl is --COO\(^-\),
forming a \textbf{zwitterion}.

\textbf{Classification by R group character:}

{\def\LTcaptype{none} % do not increment counter
\begin{longtable}[]{@{}
  >{\raggedright\arraybackslash}p{(\linewidth - 6\tabcolsep) * \real{0.1515}}
  >{\raggedright\arraybackslash}p{(\linewidth - 6\tabcolsep) * \real{0.2424}}
  >{\raggedright\arraybackslash}p{(\linewidth - 6\tabcolsep) * \real{0.3030}}
  >{\raggedright\arraybackslash}p{(\linewidth - 6\tabcolsep) * \real{0.3030}}@{}}
\toprule\noalign{}
\begin{minipage}[b]{\linewidth}\raggedright
Class
\end{minipage} & \begin{minipage}[b]{\linewidth}\raggedright
Examples
\end{minipage} & \begin{minipage}[b]{\linewidth}\raggedright
Side chain
\end{minipage} & \begin{minipage}[b]{\linewidth}\raggedright
Properties
\end{minipage} \\
\midrule\noalign{}
\endhead
\bottomrule\noalign{}
\endlastfoot
Nonpolar aliphatic & Ala, Val, Leu, Ile, Met, Pro & Hydrocarbon &
Hydrophobic core \\
Aromatic & Phe, Tyr, Trp & Benzene/indole rings & UV absorption,
stacking \\
Polar uncharged & Ser, Thr, Cys, Asn, Gln & --OH, --SH, --CONH\(_2\) &
H-bonding, disulfides \\
Positively charged & Lys, Arg, His & --NH\(_3^+\), guanidinium,
imidazole & Salt bridges, catalysis \\
Negatively charged & Asp, Glu & --COO\(^-\) & Salt bridges, metal
binding \\
\end{longtable}
}

\textbf{Special amino acids:} - \textbf{Glycine} (R = H) has no
stereogenic centre and maximal conformational flexibility --- often
found in tight turns - \textbf{Proline} has a cyclic side chain bonded
to the backbone nitrogen, restricting rotation and introducing kinks in
\texorpdfstring{\ensuremath{\alpha}}{alpha}-helices - \textbf{Cysteine} can form \textbf{disulfide bonds}
(--S--S--) with other cysteines under oxidising conditions (ER lumen,
extracellular space) - \textbf{Selenocysteine} (the ``21st amino acid'')
contains selenium instead of sulfur and is encoded by UGA
\hyperref[gl:codon]{\textbf{codon}}s with a special stem-loop (SECIS
element)

\subsubsection{The Peptide Bond}\label{the-peptide-bond}

A peptide bond forms by condensation between the carboxyl of one amino
acid and the amino group of the next:

\[\text{--NH--CHR}_1\text{--CO--NH--CHR}_2\text{--CO--} \ldots \tag{3.2} \label{eq:unit_I_macromolecules_item_5}\]

The peptide bond has \textbf{partial double-bond character}
(about 40\% double bond due to resonance), making it planar and
\textbf{trans} (almost typically; \emph{cis} occurs
about 0.03\% of the time, except before proline where
\emph{cis} occurs about 5\%). This rigidity limits rotation
around the C\(_\alpha\)--C bond to the backbone dihedral angles \texorpdfstring{\ensuremath{\varphi}}{phi} and \texorpdfstring{\ensuremath{\psi}}{psi}
(described on a Ramachandran plot), constraining protein conformational
space.

\subsubsection{Ramachandran Plots and Backbone Dihedral
Angles}\label{ramachandran-plots-and-backbone-dihedral-angles}

The polypeptide backbone has three torsional degrees of freedom per
residue:

\begin{itemize}
\tightlist
\item
  \textbf{\texorpdfstring{\ensuremath{\omega}}{omega}} (omega): rotation about the peptide C--N bond. Effectively
  fixed at 180\(^\circ\) (\emph{trans}) due to partial double-bond
  character; about 5\% of X-Pro bonds adopt \emph{cis} (\texorpdfstring{\ensuremath{\omega}}{omega} =
  0\(^\circ\)).
\item
  \textbf{\(\varphi\)} (phi): rotation about the N--C\(_\alpha\) bond.
\item
  \textbf{\texorpdfstring{\ensuremath{\psi}}{psi}} (psi): rotation about the C\(_\alpha\)--C bond.
\end{itemize}

Because each peptide unit is rigid and planar, primarily \(\varphi\) and
\texorpdfstring{\ensuremath{\psi}}{psi} can vary. The 1963 paper of G. N. Ramachandran plotted \(\varphi\)
vs.~\texorpdfstring{\ensuremath{\psi}}{psi} for most amino acids in known protein structures; the plot reveals
dramatic clustering caused by steric clashes (Pauli repulsion between
side chain, carbonyl, and amide atoms when the dihedrals approach
forbidden values).

{\def\LTcaptype{none} % do not increment counter
\begin{longtable}[]{@{}
  >{\raggedright\arraybackslash}p{(\linewidth - 4\tabcolsep) * \real{0.1053}}
  >{\raggedright\arraybackslash}p{(\linewidth - 4\tabcolsep) * \real{0.5439}}
  >{\raggedright\arraybackslash}p{(\linewidth - 4\tabcolsep) * \real{0.3509}}@{}}
\toprule\noalign{}
\begin{minipage}[b]{\linewidth}\raggedright
Region
\end{minipage} & \begin{minipage}[b]{\linewidth}\raggedright
Approximate (\(\varphi\), \texorpdfstring{\ensuremath{\psi}}{psi})
\end{minipage} & \begin{minipage}[b]{\linewidth}\raggedright
Secondary structure
\end{minipage} \\
\midrule\noalign{}
\endhead
\bottomrule\noalign{}
\endlastfoot
Right-handed \texorpdfstring{\ensuremath{\alpha}}{alpha}-helix & (--60\(^\circ\), --45\(^\circ\)) & \texorpdfstring{\ensuremath{\alpha}}{alpha}-helix \\
Left-handed \texorpdfstring{\ensuremath{\alpha}}{alpha}-helix & (+60\(^\circ\), +45\(^\circ\)) & Rare; primarily
Gly residues \\
\texorpdfstring{\ensuremath{\beta}}{beta}-sheet (parallel) & (--120\(^\circ\), +120\(^\circ\)) & Parallel
\texorpdfstring{\ensuremath{\beta}}{beta}-sheet \\
\texorpdfstring{\ensuremath{\beta}}{beta}-sheet (antiparallel) & (--140\(^\circ\), +135\(^\circ\)) &
Antiparallel \texorpdfstring{\ensuremath{\beta}}{beta}-sheet \\
Polyproline II & (--75\(^\circ\), +145\(^\circ\)) & Collagen, IDR
motifs \\
3\(_{10}\)-helix & (--49\(^\circ\), --26\(^\circ\)) & Tight helix at
N/C-termini \\
\end{longtable}
}

Two amino acids escape the standard restrictions:

\begin{itemize}
\tightlist
\item
  \textbf{Glycine} has primarily an H atom for its side chain --- no
  steric clash --- and so populates the entire \((\varphi, \psi)\)
  plane, including the disallowed regions used by left-handed helices.
\item
  \textbf{Proline's} side-chain ring fixes
  \(\varphi \approx -65^\circ\), severely restricting backbone
  flexibility. Proline cannot occur in a regular \texorpdfstring{\ensuremath{\alpha}}{alpha}-helix beyond position
  4 (no N--H proton available as H-bond donor) and is famous as a
  ``helix breaker.''
\end{itemize}

Real proteins occupy the allowed regions about 98\% of the
time; deviations into ``outlier'' regions are validated indicators of
either (a) genuine functional constraints (e.g., catalytic glycines in a
strained conformation) or (b) modelling errors in low-resolution X-ray
structures. Ramachandran statistics are now a standard quality metric
for crystallographic and AlphaFold-predicted structures.

\begin{quote}
\textbf{Concept Check 6:} Why does Pro almost rarely occur in the middle
of an \texorpdfstring{\ensuremath{\alpha}}{alpha}-helix, while Gly is over-represented in tight turns and loops?
Frame your answer in terms of \((\varphi, \psi)\) allowed regions on the
Ramachandran plot.
\end{quote}

\subsubsection{Levels of Protein
Structure}\label{levels-of-protein-structure}

\textbf{Primary structure:} The amino acid sequence --- the covalent
backbone. Written conventionally from N-terminus to C-terminus. Primary
structure is encoded by genes and determines most higher levels of
structure (Anfinsen's dogma).

\textbf{Secondary structure:} Local regular patterns of backbone
H-bonds: - \textbf{\texorpdfstring{\ensuremath{\alpha}}{alpha}-Helix:} 3.6 residues/turn, rise 0.15 nm/residue,
H-bond between residue \(i\) and \(i+4\). 100\% right-handed in natural
proteins. Helix dipole (N-terminus \texorpdfstring{\ensuremath{\delta}}{delta}+; C-terminus \texorpdfstring{\ensuremath{\delta}}{delta}-) affects function.
Helix-breaking residues: Pro (kinks), Gly (too flexible). -
\textbf{\texorpdfstring{\ensuremath{\beta}}{beta}-Sheet:} extended chains (strands) running parallel or
antiparallel, H-bonded between adjacent strands. Most \texorpdfstring{\ensuremath{\beta}}{beta}-sheets are
twisted. \texorpdfstring{\ensuremath{\beta}}{beta}-sheets are the structural basis of silk (fibroin) and amyloid
fibrils. - \textbf{\(3_{10}\)-helix:} tighter than \texorpdfstring{\ensuremath{\alpha}}{alpha}-helix (3
residues/turn), H-bond between \(i\) and \(i+3\). Less common; often
found at helix termini. - \textbf{Turns and loops:} allow the
polypeptide chain to reverse direction; typically on protein surface,
often functionally important (e.g., CDR loops in antibodies).

\textbf{Tertiary structure:} The complete 3D fold, stabilised by: -
Hydrophobic core packing (\hyperref[gl:dominant]{\textbf{dominant}}
driving force) - Disulfide bonds (Cys--S--S--Cys; formed in oxidising
compartments like ER lumen) - Salt bridges (Asp/Glu with Lys/Arg) -
H-bonds and van der Waals contacts

\textbf{Common protein folds:}

{\def\LTcaptype{none} % do not increment counter
\begin{longtable}[]{@{}
  >{\raggedright\arraybackslash}p{(\linewidth - 4\tabcolsep) * \real{0.1818}}
  >{\raggedright\arraybackslash}p{(\linewidth - 4\tabcolsep) * \real{0.5000}}
  >{\raggedright\arraybackslash}p{(\linewidth - 4\tabcolsep) * \real{0.3182}}@{}}
\toprule\noalign{}
\begin{minipage}[b]{\linewidth}\raggedright
Fold
\end{minipage} & \begin{minipage}[b]{\linewidth}\raggedright
Description
\end{minipage} & \begin{minipage}[b]{\linewidth}\raggedright
Example
\end{minipage} \\
\midrule\noalign{}
\endhead
\bottomrule\noalign{}
\endlastfoot
Globin fold & 8 \texorpdfstring{\ensuremath{\alpha}}{alpha}-helices surrounding a haem pocket & Myoglobin,
haemoglobin \\
TIM barrel & 8 alternating \(\beta\alpha\) units forming a barrel &
Triosephosphate isomerase \\
\texorpdfstring{\ensuremath{\beta}}{beta}-barrel & \texorpdfstring{\ensuremath{\beta}}{beta}-sheets curved into a closed barrel & Porins, GFP \\
Immunoglobulin fold & Two \texorpdfstring{\ensuremath{\beta}}{beta}-sheets packed face-to-face & Antibodies,
MHC \\
Leucine-rich repeat & Horseshoe of repeating units & Toll-like
receptors \\
\end{longtable}
}

\textbf{Quaternary structure:} Assembly of two or more polypeptide
chains (subunits). Examples: - \textbf{Haemoglobin:} \(\alpha_2\beta_2\)
tetramer; cooperative O\(_2\) binding - \textbf{Collagen:} triple helix
of three chains with Gly-X-Y repeats (X often Pro, Y often Hyp) -
\textbf{Immunoglobulin G:} two heavy chains + two light chains; Y-shaped
- \textbf{Ribosome:} RNA + \textgreater50 protein subunits

\textbf{Protein denaturation:} Disruption of
secondary/tertiary/quaternary structure by heat, pH extremes,
denaturants (urea, guanidinium), reducing agents (break disulfides).
Primary structure intact; activity lost. Often irreversible (cooked egg
white), but many small proteins refold spontaneously (Anfinsen's
experiment with ribonuclease A).

\subsubsection{Thermodynamics of Protein
Folding}\label{thermodynamics-of-protein-folding}

Protein folding is governed by the same Gibbs equation that governs most
spontaneous processes:

\begin{equation}
\Delta G_{\text{fold}} = \Delta H_{\text{fold}} - T\,\Delta S_{\text{fold}}
\label{eq:unit_I_protein_folding_dG}
\end{equation}

The total \(\Delta G_{\text{fold}}\) for a typical small globular
protein is primarily --20 to --60 kJ/mol --- the difference between two
enormous, partially cancelling terms.

\textbf{Decomposing the folding free energy:}

{\def\LTcaptype{none} % do not increment counter
\begin{longtable}[]{@{}
  >{\raggedright\arraybackslash}p{(\linewidth - 6\tabcolsep) * \real{0.2927}}
  >{\raggedright\arraybackslash}p{(\linewidth - 6\tabcolsep) * \real{0.0976}}
  >{\raggedright\arraybackslash}p{(\linewidth - 6\tabcolsep) * \real{0.4634}}
  >{\raggedright\arraybackslash}p{(\linewidth - 6\tabcolsep) * \real{0.1463}}@{}}
\toprule\noalign{}
\begin{minipage}[b]{\linewidth}\raggedright
Contribution
\end{minipage} & \begin{minipage}[b]{\linewidth}\raggedright
Sign
\end{minipage} & \begin{minipage}[b]{\linewidth}\raggedright
Magnitude (typical)
\end{minipage} & \begin{minipage}[b]{\linewidth}\raggedright
Origin
\end{minipage} \\
\midrule\noalign{}
\endhead
\bottomrule\noalign{}
\endlastfoot
Conformational entropy of backbone & \(-T\Delta S < 0\) unfavourable &
\(+500\) to \(+1500\) kJ/mol & Loss of \(\varphi\)/\texorpdfstring{\ensuremath{\psi}}{psi} freedom \\
Conformational entropy of side chains & unfavourable & \(+200\) to
\(+400\) kJ/mol & Side chains lock in place \\
Hydrophobic effect & favourable & \(-300\) to \(-1000\) kJ/mol & Burial
of nonpolar surface; entropy of water released \\
Backbone H-bonds & small & \(\pm 0\) to \(-100\) kJ/mol & Replace
H-bonds to water with intramolecular ones \\
Van der Waals packing & favourable & \(-200\) to \(-500\) kJ/mol & Tight
core packing \\
Disulfide bonds (if present) & favourable & \(-15\) to \(-30\) kJ/mol
per S--S & Topology constraint \\
\textbf{Net \(\Delta G_{\text{fold}}\)} & favourable, marginal &
\textbf{\(-20\) to \(-60\) kJ/mol} & Sum of large opposing terms \\
\end{longtable}
}

\textbf{The hydrophobic collapse model.} The dominant driving force is
the \hyperref[gl:hydrophobic-effect]{\textbf{hydrophobic effect}} (see
\cref{sec:unit_I_water_and_life}). Folding occurs in two stages:

\begin{enumerate}
\def\labelenumi{\arabic{enumi}.}
\tightlist
\item
  \textbf{Hydrophobic collapse} (microseconds): Water-exposed nonpolar
  side chains rapidly aggregate into a ``molten globule'' --- a compact
  but loosely packed intermediate with most secondary structure formed
  but tertiary contacts not yet locked in.
\item
  \textbf{Specific tertiary contact formation} (milliseconds--seconds):
  Slow rearrangement to the unique native fold; backbone H-bonds and van
  der Waals packing optimise.
\end{enumerate}

The marginal stability of folded proteins (\(|\Delta G| \approx\) a few
H-bonds' worth) has profound consequences: small mutations can
destabilise proteins enough to cause disease (e.g., a single
Glu6-\textgreater Val substitution causes sickle-cell haemoglobin to
polymerise), and many proteins exist in a fragile balance with
disordered or misfolded states.

\textbf{Levinthal's paradox and the folding funnel.} A 100-residue chain
with 3 backbone conformations per residue would have
\(3^{100} \approx 10^{47}\) possible conformations --- searching them
randomly at 10\(^{12}\) trials per second would take longer than the age
of the universe, yet proteins fold in milliseconds. The resolution is
the \textbf{folding funnel}: the energy landscape is biased so that
energy and conformational entropy \emph{both} decrease as the protein
nears the native state. The chain rolls down the funnel via many
parallel pathways rather than searching exhaustively. This conceptual
breakthrough is now supported by molecular dynamics simulations and by
structure-prediction systems. AlphaFold 3 extends the classroom question
from ``what fold does this protein adopt?'' to ``what complex does this
protein form with DNA, RNA, ligands, ions, or modified residues?'' It
should be read as a powerful hypothesis generator whose confidence
metrics and biochemical context still need experimental validation,
especially for dynamics, alternative conformations, and low-confidence
interfaces \citep{abramson2024alphafold3}.

\textbf{Current evidence: AI structure models and molecular interaction
claims.} The AlphaFold Protein Structure Database update reported for
the 2025 UniProt release improved the entry interface and structural
interpretation of hundreds of millions of predicted protein models,
while a 2026 EMBL-EBI/Google DeepMind/NVIDIA/SNU update began adding
predicted protein complexes, starting with high-confidence homodimers
\citep{velankar2026alphafolddb2025,emblebi2026alphafoldcomplexes}. For a
biology student, the key habit is not ``the computer solved the
protein'' but ``the model proposes a testable molecular hypothesis.'' A
predicted active-site pocket, ligand pose, or interface should be
checked against confidence metrics, conservation, mutagenesis, binding,
kinetics, cryo-EM/X-ray/NMR evidence, and whether the biological state
is static, dynamic, membrane-embedded, modified, or allosterically
regulated.

\begin{figure}[htbp]
\centering
\includegraphics[width=0.82\linewidth,height=4.2in,keepaspectratio]{../figures/mermaid_inline/inline_mermaid_0024_ea8285dc7e2b.png}
\caption{AI biomolecular modelling workflow. Predicted structures and complexes are evidence-generating tools; confidence metrics and independent assays determine whether a predicted contact supports a biological mechanism.}
\end{figure}

\subsubsection{Intrinsically Disordered Proteins
(IDPs)}\label{intrinsically-disordered-proteins-idps}

A revolutionary modern discovery challenged the classical ``structure =
function'' paradigm: approximately 30--50\% of eukaryotic proteins
contain \textbf{intrinsically disordered regions} (IDRs) --- segments
that do not adopt a fixed 3D structure under physiological conditions.

Properties of IDPs/IDRs: - Enriched in charged and polar residues (Lys,
Arg, Glu, Ser, Pro, Gly); depleted in hydrophobic residues - Exist as
dynamic ensembles of rapidly interconverting conformations - Often
undergo \textbf{coupled folding and binding} --- they fold upon
interacting with their binding partner

\textbf{Biological advantages of disorder:} - \textbf{Promiscuous
binding:} one IDR can bind multiple partners (hub proteins in signalling
networks) - \textbf{Large interaction surfaces:} extended conformations
provide large binding interfaces from a short sequence - \textbf{Rapid
binding kinetics:} ``fly-casting'' mechanism allows faster association
rates - \textbf{Regulation by post-translational modification:} many
phosphorylation sites are in disordered regions

\textbf{Examples:} - \textbf{p53 transactivation domain:} disordered;
folds into a helix upon binding MDM2 - \textbf{Tau protein:}
intrinsically disordered; pathological aggregation into neurofibrillary
tangles causes Alzheimer's disease - \textbf{Nucleoporin FG repeats:}
disordered regions form the permeability barrier of nuclear pores -
\textbf{\texorpdfstring{\ensuremath{\alpha}}{alpha}-Synuclein:} disordered; aggregation into Lewy bodies is the
hallmark of Parkinson's disease

\textbf{IDPs and biomolecular condensates --- liquid--liquid phase
separation (LLPS).} A revolutionary insight of the 2010s: many cellular
compartments have no membrane. The nucleolus, P-bodies, stress granules,
the centrosome, heterochromatin domains, and synaptic puncta are
\textbf{biomolecular condensates} formed by liquid--liquid phase
separation. The driving force is multivalent, weak interactions between
IDR-containing proteins and/or RNAs. Above a threshold concentration,
the polymer-rich solution demixes into a dense droplet phase coexisting
with a dilute phase, much like oil droplets in a vinaigrette.

Key features of LLPS-driven condensates: - \textbf{Multivalency}
(multiple weak interaction modules) is required; single-site
interactions give simple complexes, not phases. - \textbf{Aromatic
(\texorpdfstring{\ensuremath{\pi}}{pi}--\texorpdfstring{\ensuremath{\pi}}{pi}) and electrostatic (\(+\)/\(-\))} interactions between IDR
residues (especially Arg, Tyr, Phe, Glu, Asp) drive demixing. -
\textbf{Functional consequence:} condensates concentrate biochemical
reactants about 100-fold above bulk cytosol, dramatically
accelerating reactions. - \textbf{Disease link:} ALS-causing mutations
in FUS and TDP-43 transition liquid condensates into solid amyloid-like
aggregates (the ``liquid-to-solid'' transition). Understanding LLPS is
now central to neurodegeneration research.

This discovery transformed cell biology: the nucleus is not a uniform
compartment but a phase-organised mosaic of dozens of distinct liquid
bodies, each enriching specific molecules without membranes to maintain
identity.

\begin{quote}
\textbf{Clinical Connection: Amyloid Diseases and Protein Misfolding}

When normally soluble proteins or IDPs misfold and aggregate into
ordered, cross-\texorpdfstring{\ensuremath{\beta}}{beta} \textbf{amyloid fibrils}, the result can be devastating
disease. Alzheimer's disease involves amyloid-\texorpdfstring{\ensuremath{\beta}}{beta} plaques and tau tangles;
Parkinson's involves \texorpdfstring{\ensuremath{\alpha}}{alpha}-synuclein Lewy bodies; prion diseases (CJD, mad
cow disease) involve PrP\(^{Sc}\) amyloid. Type 2 diabetes involves
islet amyloid polypeptide (IAPP/amylin) deposits in pancreatic islets.
Understanding protein misfolding is one of the grand challenges of
molecular medicine.
\end{quote}

\begin{quote}
\textbf{Concept Check 4:} If a protein with an intrinsically disordered
region can bind multiple different partner proteins, what advantage does
this provide to \hyperref[gl:signal-transduction]{\textbf{signal
transduction}} networks? What are the risks?
\end{quote}

\subsubsection{Post-Translational Modifications
(PTMs)}\label{post-translational-modifications-ptms}

A typical eukaryotic protein is modified by a ``code'' of covalent
chemistry layered on top of the genetic sequence. Over 400 distinct PTMs
are now catalogued; the major ones are summarised below.

{\def\LTcaptype{none} % do not increment counter
\begin{longtable}[]{@{}
  >{\raggedright\arraybackslash}p{(\linewidth - 8\tabcolsep) * \real{0.0476}}
  >{\raggedright\arraybackslash}p{(\linewidth - 8\tabcolsep) * \real{0.2063}}
  >{\raggedright\arraybackslash}p{(\linewidth - 8\tabcolsep) * \real{0.2381}}
  >{\raggedright\arraybackslash}p{(\linewidth - 8\tabcolsep) * \real{0.2698}}
  >{\raggedright\arraybackslash}p{(\linewidth - 8\tabcolsep) * \real{0.2381}}@{}}
\toprule\noalign{}
\begin{minipage}[b]{\linewidth}\raggedright
PTM
\end{minipage} & \begin{minipage}[b]{\linewidth}\raggedright
Adding enzyme
\end{minipage} & \begin{minipage}[b]{\linewidth}\raggedright
Removing enzyme
\end{minipage} & \begin{minipage}[b]{\linewidth}\raggedright
Target residue(s)
\end{minipage} & \begin{minipage}[b]{\linewidth}\raggedright
Functional role
\end{minipage} \\
\midrule\noalign{}
\endhead
\bottomrule\noalign{}
\endlastfoot
Phosphorylation & Kinases (518 in human) & Phosphatases & Ser, Thr, Tyr
(rare on His) & Reversible ``switch''; signalling, enzyme activity \\
Acetylation & Acetyltransferases (HATs, GCN5) & Deacetylases (HDACs,
sirtuins) & Lys \(\varepsilon\)-NH\(_2\), N-terminus & Chromatin opening
(histones); metabolic regulation \\
Methylation & Methyltransferases (SET domain, PRMT) & Demethylases
(KDMs, LSD1) & Lys, Arg, His; DNA at C5 of cytosine & Gene
silencing/activation; epigenetic memory \\
Ubiquitination & E1-\textgreater E2-\textgreater E3 cascade & DUBs
(deubiquitinases) & Lys \(\varepsilon\)-NH\(_2\) & Protein turnover
(proteasome); signalling (K63 chains) \\
SUMOylation & SUMO E1, E2, E3 & SENPs & Lys (often in \texorpdfstring{\ensuremath{\Psi}}{Psi}KX(D/E) motif) &
Nuclear transport; transcription; DNA repair \\
Glycosylation & Glycosyltransferases & Glycosidases & Asn (N-linked);
Ser/Thr (O-linked); Trp (C-mannose) & Protein folding, half-life,
recognition \\
O-GlcNAc & OGT & OGA & Ser, Thr (cytosolic/nuclear) & Reciprocal
cross-talk with phosphorylation; nutrient sensing \\
Hydroxylation & Prolyl/lysyl hydroxylases & --- (irreversible) & Pro,
Lys (collagen, HIF-1\texorpdfstring{\ensuremath{\alpha}}{alpha}) & Collagen stability; oxygen sensing \\
Lipidation & NMTs, palmitoyltransferases & Acyl-protein thioesterases &
N-terminal Gly (myristate); Cys (palmitate, prenyl) & Membrane
targeting \\
ADP-ribosylation & PARPs & PARGs, ARH3 & Glu, Asp, Arg & DNA damage
signalling \\
Disulfide & PDI in ER & Reductants & Cys pairs & Structural; secreted
protein stability \\
Proteolytic cleavage & Proteases & (irreversible) & Specific peptide
bonds & Activation (zymogens); destruction; signal release \\
\end{longtable}
}

\textbf{Phosphorylation} is the most-studied PTM: about 30\% of
human proteins are phosphorylated on at least one residue at some point
in their life cycle, controlled by 518 kinases and about 150
phosphatases. A single phosphate group adds two negative charges and
about 30 kJ/mol of binding energy with phospho-recognition
domains (SH2, 14-3-3, FHA), enabling switch-like regulation of activity,
localisation, or partnerships. \textbf{Histone PTMs} (acetylation of
Lys, methylation of Lys/Arg, phosphorylation of Ser/Thr) write the
``histone code'' that organises chromatin into transcriptionally active
and silent states (see Unit IV).

\textbf{Ubiquitin} is itself a small (76-residue) protein that is
covalently attached to other proteins, usually through an isopeptide
bond between its C-terminal Gly76 and a Lys \(\varepsilon\)-NH\(_2\) on
the substrate. Polyubiquitin chains can grow on any of ubiquitin's seven
Lys residues (or its own N-terminus), creating different chain
topologies with different fates: K48 chains target proteins for
proteasomal destruction; K63 chains regulate signalling and DNA repair;
M1 (linear) chains regulate NF-\texorpdfstring{\ensuremath{\kappa}}{kappa}B inflammation. The 2004 Nobel Prize in
Chemistry recognised Aaron Ciechanover, Avram Hershko, and Irwin Rose
for discovering this regulated proteolysis system.

\begin{quote}
\textbf{Concept Check 8:} Phosphorylation of a Ser residue adds two
negative charges and about 80 Da of mass --- a tiny chemical
change. Yet kinase signalling cascades transmit information through
millions of phosphorylation events per second across thousands of
substrates. Identify two structural features of phosphate (charge,
geometry, H-bond capacity) that make it an ideal ``switch'' molecule for
biology, and contrast with why methylation (a smaller, neutral
modification) is used for slower, longer-lived regulation.
\end{quote}

\begin{center}\rule{0.5\linewidth}{0.5pt}\end{center}

\subsection{Nucleic Acids}\label{nucleic-acids}

\subsubsection{Nucleotide Monomers}\label{nucleotide-monomers}

Each \hyperref[gl:nucleotide]{\textbf{nucleotide}} consists of: 1. A
pentose sugar (ribose in RNA; 2'-deoxyribose in DNA) 2. A nitrogenous
base (purine or pyrimidine) 3. One to three phosphate groups

\textbf{Bases:} - Purines (two-ring): Adenine (A), Guanine (G) -
Pyrimidines (one-ring): Cytosine (C), Thymine (T; DNA primarily), Uracil
(U; RNA primarily)

The \textbf{nucleoside} is base + sugar. The \textbf{nucleotide} is
nucleoside + phosphate(s).

\textbf{Phosphodiester bond:} Nucleotides are linked 3'-\textgreater5'
by phosphodiester bonds, forming the sugar-phosphate backbone. This
backbone is highly uniform and carries a negative charge at every
phosphate (one negative charge per nucleotide at physiological pH).

\subsubsection{DNA Structure}\label{dna-structure}

DNA (deoxyribonucleic acid) is a double-stranded helix discovered by
Watson and Crick in 1953, guided by X-ray diffraction data from Franklin
and Wilkins and Chargaff's rules:

\begin{itemize}
\tightlist
\item
  \textbf{Chargaff's rules:} \hyperref[a]{A} = \hyperref[t]{T};
  \hyperref[g]{G} = \hyperref[c]{C} in any double-stranded DNA
\item
  \textbf{Complementary base pairs:} A--T (2
  \hyperref[gl:hydrogen-bond]{\textbf{hydrogen bond}}s) and G--C (3
  hydrogen bonds)
\item
  \textbf{Antiparallel strands:} one runs 5'-\textgreater3', its
  complement runs 3'-\textgreater5'
\end{itemize}

\textbf{Structural forms of DNA:}

{\def\LTcaptype{none} % do not increment counter
\begin{longtable}[]{@{}
  >{\raggedright\arraybackslash}p{(\linewidth - 6\tabcolsep) * \real{0.2800}}
  >{\raggedright\arraybackslash}p{(\linewidth - 6\tabcolsep) * \real{0.2400}}
  >{\raggedright\arraybackslash}p{(\linewidth - 6\tabcolsep) * \real{0.2400}}
  >{\raggedright\arraybackslash}p{(\linewidth - 6\tabcolsep) * \real{0.2400}}@{}}
\toprule\noalign{}
\begin{minipage}[b]{\linewidth}\raggedright
Feature
\end{minipage} & \begin{minipage}[b]{\linewidth}\raggedright
A-form
\end{minipage} & \begin{minipage}[b]{\linewidth}\raggedright
B-form
\end{minipage} & \begin{minipage}[b]{\linewidth}\raggedright
Z-form
\end{minipage} \\
\midrule\noalign{}
\endhead
\bottomrule\noalign{}
\endlastfoot
Helix direction & Right-handed & Right-handed & Left-handed \\
Base pairs/turn & 11 & 10.5 & 12 \\
Rise per bp & 0.26 nm & 0.34 nm & 0.37 nm \\
Helix diameter & 2.3 nm & 2.0 nm & 1.8 nm \\
Groove pattern & Deep narrow major, wide shallow minor & Wide major,
narrow minor & Flat major, narrow deep minor \\
Conditions & Dehydrated, RNA-DNA hybrids & Physiological & High salt,
alternating purine-pyrimidine \\
Biological role & dsRNA, RNA-DNA hybrids & Primary cellular form & Gene
regulation (\hyperref[gl:cpg-island]{\textbf{CpG island}}s?) \\
\end{longtable}
}

The \textbf{B-form} is the biologically dominant form at physiological
ionic strength: - Right-handed helix; 10.5 base pairs/turn; rise 0.34
nm/bp; pitch \texorpdfstring{\ensuremath{\approx}}{~}3.6 nm; diameter 2.0 nm - Major groove (wide,
about 1.2 nm) and minor groove (about 0.6 nm) ---
proteins bind these grooves

\textbf{Base stacking:} Beyond hydrogen bonding between complementary
bases, the flat aromatic rings of adjacent bases stack on top of each
other via \texorpdfstring{\ensuremath{\pi}}{pi}-\texorpdfstring{\ensuremath{\pi}}{pi} interactions and
\hyperref[gl:van-der-waals-forces]{\textbf{van der Waals forces}}. These
stacking interactions contribute more to duplex stability than hydrogen
bonding. This is why DNA melting temperature (\(T_m\)) increases with GC
content (three H-bonds + stronger stacking) and with ionic strength
(cations neutralise backbone charges).

\textbf{Quantitative base-stacking energies.} Nearest-neighbour stacking
energies (the contribution of each adjacent base-pair \emph{step}, not
each individual base) have been calibrated from melting curves and
computational chemistry:

{\def\LTcaptype{none} % do not increment counter
\begin{longtable}[]{@{}
  >{\raggedright\arraybackslash}p{(\linewidth - 6\tabcolsep) * \real{0.2015}}
  >{\raggedright\arraybackslash}p{(\linewidth - 6\tabcolsep) * \real{0.2090}}
  >{\raggedright\arraybackslash}p{(\linewidth - 6\tabcolsep) * \real{0.3284}}
  >{\raggedright\arraybackslash}p{(\linewidth - 6\tabcolsep) * \real{0.2612}}@{}}
\toprule\noalign{}
\begin{minipage}[b]{\linewidth}\raggedright
Step (5'-\textgreater3' top strand)
\end{minipage} & \begin{minipage}[b]{\linewidth}\raggedright
Stacking \(\Delta H\) (kJ/mol)
\end{minipage} & \begin{minipage}[b]{\linewidth}\raggedright
Stacking \(\Delta S\) (J mol\(^{-1}\) K\(^{-1}\))
\end{minipage} & \begin{minipage}[b]{\linewidth}\raggedright
\(\Delta G\) at 37 \(^\circ\)C (kJ/mol)
\end{minipage} \\
\midrule\noalign{}
\endhead
\bottomrule\noalign{}
\endlastfoot
AA / TT & --33.1 & --92.9 & --4.3 \\
AT / AT & --30.1 & --85.4 & --3.6 \\
TA / TA & --30.1 & --89.1 & --2.4 \\
CA / GT & --35.6 & --95.0 & --6.2 \\
GT / CA & --35.1 & --93.7 & --6.1 \\
CT / GA & --32.6 & --87.9 & --5.4 \\
GA / CT & --34.3 & --92.9 & --5.5 \\
CG / CG & --44.4 & --113.8 & --9.1 \\
GC / GC & --41.0 & --102.1 & --9.4 \\
GG / CC & --33.5 & --83.3 & --7.7 \\
\end{longtable}
}

Notice that \textbf{GC steps have the strongest stacking} (and therefore
highest melting temperature) --- this is why GC-rich regions (CpG
islands, gene promoters) require greater thermal energy to denature, and
why PCR primers are usually designed to be about 50\% GC. The
total free energy of a duplex is summed over most \(n - 1\) stacking
steps in an \(n\)-base-pair duplex, plus initiation/end corrections;
this nearest-neighbour model (Santa Lucia, 1998) predicts \(T_m\) values
within about 2 \(^\circ\)C of measurement and is the basis of
every PCR primer design tool.

\textbf{A-, B-, and Z-DNA in detail.} The three structural forms differ
in geometry and biological occurrence:

\begin{itemize}
\tightlist
\item
  \textbf{B-form DNA} (the canonical Watson-Crick form): right-handed,
  10.5 bp/turn, helical rise 0.34 nm, diameter 2.0 nm; major groove 1.2
  nm wide and 0.85 nm deep; minor groove 0.6 nm wide and 0.75 nm deep.
  Adopted under physiological hydration. The wide, deep major groove is
  the principal site of sequence-specific protein binding (most
  transcription factors read here).
\item
  \textbf{A-form} (right-handed, more compact): 11 bp/turn, rise 0.26
  nm, diameter 2.3 nm; deep, narrow major groove and shallow, broad
  minor groove. Adopted by RNA-DNA hybrids (e.g., during transcription)
  and by double-stranded RNA. Forced by the 2'-OH of ribose, which
  sterically excludes B-geometry. The A-form transition during
  transcription is essential for RNA polymerase to thread its product.
\item
  \textbf{Z-form} (left-handed): 12 bp/turn, rise 0.37 nm, diameter 1.8
  nm; zigzag backbone (hence ``Z''); flat major groove with no
  protein-binding pocket; narrow, deep minor groove. Adopted by
  sequences with alternating purine--pyrimidine patterns ((CG)\(_n\)
  being the canonical case) under high ionic strength or negative
  supercoiling. Z-DNA-binding proteins (ADAR1, ZBP1) play roles in
  innate immunity and CpG island regulation.
\end{itemize}

The conformational equilibrium B \(\rightleftharpoons\) Z is sensitive
to negative supercoiling generated behind a transcribing RNA polymerase;
transient Z-DNA tracts have been observed at active promoters and are
believed to participate in transcription regulation.

Human \hyperref[gl:genome]{\textbf{genome}}: \(3.2 \times 10^9\) base
pairs; if laid end-to-end, 2.2 m long but compressed into 6 \texorpdfstring{\ensuremath{\mu}}{mu}m nucleus
via \hyperref[gl:nucleosome]{\textbf{nucleosome}} packaging (DNA wound
around \hyperref[gl:histone]{\textbf{histone}} octamers in 147 bp
segments).

\subsubsection{Worked Example: DNA Melting Temperature from
\%GC}\label{worked-example-dna-melting-temperature-from-gc}

\textbf{Problem:} A laboratory amplifies a \(100\)-bp segment of a CpG
island whose G+C content is \(50\%\), in a buffer of
\([\text{Na}^+] = 50 \text{ mM}\). Using the length-corrected empirical
relation
\[ T_m = 81.5 + 0.41\,(\%GC) - \frac{675}{N}  \label{eq:unit_I_macromolecules_item_6}\]

estimate the melting temperature \(T_m\). Then repeat for a GC-rich
variant of the same length with \(\%GC = 70\), and interpret the
difference.

\textbf{Solution:}

\begin{enumerate}
\def\labelenumi{\arabic{enumi}.}
\item
  \textbf{Substitute the \(50\%\)-GC duplex} (\(\%GC = 50\),
  \(N = 100\), \([\text{Na}^+] = 50 \text{ mM}\)):
  \[ T_m = 81.5 + 0.41 \times 50 - \frac{675}{100} = 81.5 + 20.5 - 6.75  \label{eq:unit_I_macromolecules_item_7}\]

  \[ T_m = 95.25\ ^\circ\text{C}  \label{eq:unit_I_macromolecules_item_8}\]
\item
  \textbf{Substitute the \(70\%\)-GC duplex} (same \(N = 100\), same
  \([\text{Na}^+]\)):
  \[ T_m = 81.5 + 0.41 \times 70 - \frac{675}{100} = 81.5 + 28.7 - 6.75 = 103.45\ ^\circ\text{C}  \label{eq:unit_I_macromolecules_item_9}\]

  The \(20\)-percentage-point rise in G+C raises \(T_m\) by
  \(103.45 - 95.25 = 8.20\ ^\circ\text{C}\).
\end{enumerate}

Because every G--C pair contributes three hydrogen bonds and stronger
base-stacking than an A--T pair, a more GC-rich duplex of the same
length resists thermal denaturation more strongly --- which is why
GC-rich promoters and CpG islands melt last, and why PCR primers are
designed around predictable \(T_m\) values.

\subsubsection{RNA Structure and
Diversity}\label{rna-structure-and-diversity}

RNA is single-stranded and uses uracil instead of thymine. The 2'-OH on
ribose makes RNA more chemically labile than DNA (susceptible to
base-catalysed hydrolysis: the 2'-OH attacks the adjacent phosphate,
cleaving the backbone). This is why RNA's evolutionary niche is
short-lived information (mRNA, miRNA, sRNAs) while the longer-term
archive is DNA.

RNA folds on itself via intramolecular base pairing, forming complex
secondary and tertiary structures. Unlike DNA's relatively uniform
B-form duplex, RNA structure is wildly diverse --- ribozymes, tRNAs, and
ribosomal RNAs adopt 3D shapes as elaborate as folded proteins, often
with deep clefts and pockets that are functional active sites.

\textbf{RNA secondary structure elements:}

{\def\LTcaptype{none} % do not increment counter
\begin{longtable}[]{@{}
  >{\raggedright\arraybackslash}p{(\linewidth - 4\tabcolsep) * \real{0.2692}}
  >{\raggedright\arraybackslash}p{(\linewidth - 4\tabcolsep) * \real{0.4231}}
  >{\raggedright\arraybackslash}p{(\linewidth - 4\tabcolsep) * \real{0.3077}}@{}}
\toprule\noalign{}
\begin{minipage}[b]{\linewidth}\raggedright
Element
\end{minipage} & \begin{minipage}[b]{\linewidth}\raggedright
Description
\end{minipage} & \begin{minipage}[b]{\linewidth}\raggedright
Function
\end{minipage} \\
\midrule\noalign{}
\endhead
\bottomrule\noalign{}
\endlastfoot
Stem-loop (hairpin) & Complementary bases form a stem; unpaired bases
form a loop & Most common motif; regulatory elements \\
Internal loop & Mismatched bases within a stem & Protein binding
sites \\
Bulge & Unpaired nucleotides on one strand of a stem & Flexibility,
protein recognition \\
Pseudoknot & Loop base-pairs with downstream sequence & Ribosome
frameshifting, telomerase \\
G-quadruplex & Four guanines form a planar quartet via Hoogsteen H-bonds
& \hyperref[gl:telomere]{\textbf{Telomere}}s, oncogene regulation \\
Kissing loop & Two hairpin loops interact via complementary sequences &
Viral RNA dimerisation \\
\end{longtable}
}

\textbf{Pseudoknots} illustrate the topological complexity unique to
RNA. After a hairpin loop forms, \emph{unpaired} loop bases pair with a
complementary sequence located \emph{downstream}, creating a second
helix that threads through the first. The HIV-1 frameshift element is a
pseudoknot that pauses the ribosome, forcing a --1 frameshift to
translate the \emph{gag-pol} polyprotein. Telomerase contains a
pseudoknot in its RNA template that is essential for catalysis;
mutations in this pseudoknot cause dyskeratosis congenita (a hereditary
bone-marrow failure syndrome).

\textbf{RNA secondary structure prediction.} Whether a given RNA
sequence will adopt a particular fold can be predicted computationally
using thermodynamic models. The free energy of a fold is the sum of
nearest-neighbour stacking energies (similar to the DNA table above)
minus penalties for loops, bulges, and unpaired regions. The Zuker
algorithm finds the minimum-free-energy structure in \(O(n^3)\) time for
a sequence of length \(n\). Modern tools (RNAfold, ViennaRNA, NUPACK)
extend this to ensembles of suboptimal structures, base-pairing
probabilities, and pseudoknot prediction. Predicted folds are now
routinely used to design siRNAs, mRNA vaccines (e.g., the Pfizer/Moderna
COVID-19 mRNAs were sequence-optimised to favour stable secondary
structures and avoid immunogenic motifs), and guide RNAs for CRISPR
systems.

\textbf{Ribozymes revisited --- 3D structure shapes catalysis.} The
hammerhead ribozyme (about 40 nt, the smallest known) folds
into three coaxially stacked helices that position a 2'-OH for
nucleophilic attack on the adjacent phosphodiester, cleaving the
backbone via an in-line mechanism. The hepatitis delta ribozyme uses a
cytosine as a general acid catalyst with shifted pKa --- demonstrating
that RNA bases can perform acid-base chemistry analogous to enzyme side
chains.

\textbf{RNA types and functions:}

{\def\LTcaptype{none} % do not increment counter
\begin{longtable}[]{@{}
  >{\raggedright\arraybackslash}p{(\linewidth - 4\tabcolsep) * \real{0.2857}}
  >{\raggedright\arraybackslash}p{(\linewidth - 4\tabcolsep) * \real{0.4286}}
  >{\raggedright\arraybackslash}p{(\linewidth - 4\tabcolsep) * \real{0.2857}}@{}}
\toprule\noalign{}
\begin{minipage}[b]{\linewidth}\raggedright
RNA type
\end{minipage} & \begin{minipage}[b]{\linewidth}\raggedright
Abbreviation
\end{minipage} & \begin{minipage}[b]{\linewidth}\raggedright
Function
\end{minipage} \\
\midrule\noalign{}
\endhead
\bottomrule\noalign{}
\endlastfoot
Messenger RNA & mRNA & Carries protein-coding sequence \\
Transfer RNA & tRNA & Carrier of amino acids to ribosome \\
Ribosomal RNA & rRNA & Structural/catalytic core of ribosome \\
Small nuclear RNA & snRNA & Pre-mRNA splicing (spliceosome) \\
MicroRNA & miRNA & Post-transcriptional gene silencing \\
Long non-coding RNA & lncRNA &
\hyperref[gl:chromatin]{\textbf{Chromatin}} regulation, imprinting \\
Small interfering RNA & siRNA & RNA interference (gene silencing) \\
CRISPR RNA & crRNA & Guide for Cas nuclease (adaptive immunity) \\
Piwi-interacting RNA & piRNA & Transposon silencing in germ cells \\
\end{longtable}
}

\subsubsection{Ribozymes: Catalytic RNA}\label{ribozymes-catalytic-rna}

The \textbf{ribozymes} --- catalytically active RNAs --- demonstrate
that RNA was likely the primordial molecule of life in the \textbf{RNA
World} hypothesis.

\textbf{Discovery:} In 1982, Thomas Cech discovered that the Group I
\hyperref[gl:intron]{\textbf{intron}} of \emph{Tetrahymena} rRNA
catalyses its own excision (self-splicing) without any protein enzyme.
Shortly after, Sidney Altman showed that the RNA component of RNase P
catalyses tRNA maturation. Both received the \textbf{1989 Nobel Prize in
Chemistry} for demonstrating RNA catalysis.

\textbf{Known ribozymes:}

{\def\LTcaptype{none} % do not increment counter
\begin{longtable}[]{@{}
  >{\raggedright\arraybackslash}p{(\linewidth - 4\tabcolsep) * \real{0.2857}}
  >{\raggedright\arraybackslash}p{(\linewidth - 4\tabcolsep) * \real{0.2857}}
  >{\raggedright\arraybackslash}p{(\linewidth - 4\tabcolsep) * \real{0.4286}}@{}}
\toprule\noalign{}
\begin{minipage}[b]{\linewidth}\raggedright
Ribozyme
\end{minipage} & \begin{minipage}[b]{\linewidth}\raggedright
Reaction
\end{minipage} & \begin{minipage}[b]{\linewidth}\raggedright
Significance
\end{minipage} \\
\midrule\noalign{}
\endhead
\bottomrule\noalign{}
\endlastfoot
Group I intron & Phosphoester transfer (self-splicing) & First
discovered ribozyme \\
Group II intron & Self-splicing (lariat intermediate) & Ancestor of
spliceosome \\
RNase P & Endonuclease (tRNA 5'-end processing) & Essential in most
domains of life \\
Hammerhead ribozyme & Self-cleavage (phosphodiester) & Smallest known
ribozyme (about 40 nt) \\
Hepatitis delta ribozyme & Self-cleavage & Viral replication \\
\textbf{Ribosome (23S rRNA)} & \textbf{Peptidyl transfer} & \textbf{The
ribosome is a ribozyme!} \\
\end{longtable}
}

The discovery that the ribosome's peptidyl transferase activity resides
in rRNA (not ribosomal proteins) was arguably the most important
confirmation of the RNA World hypothesis. Thomas Steitz, Venkatraman
Ramakrishnan, and Ada Yonath shared the \textbf{2009 Nobel Prize in
Chemistry} for elucidating ribosome structure at atomic resolution.

\begin{quote}
\textbf{Clinical Connection: Ribozyme-Based Therapeutics}

Engineered ribozymes (e.g., hammerhead ribozymes targeting specific
mRNAs) have been explored as gene therapy tools. More practically,
\textbf{riboswitches} --- RNA elements that bind small molecules and
regulate gene expression --- are promising antibiotic targets. Because
riboswitches are found in bacteria but not humans, drugs that interfere
with riboswitch function could selectively kill pathogens. The
antibiotic roseoflavin targets the FMN riboswitch in Gram-positive
bacteria.
\end{quote}

\begin{quote}
\textbf{Concept Check 5:} The RNA World hypothesis proposes that RNA
preceded both DNA and proteins as the primordial macromolecule of life.
What three properties must RNA possess to fulfil this role, and does it
actually possess most three?
\end{quote}

\begin{center}\rule{0.5\linewidth}{0.5pt}\end{center}

\subsection{Key Comparisons}\label{key-comparisons}

\subsubsection{The Four Macromolecule Classes
Compared}\label{the-four-macromolecule-classes-compared}

{\def\LTcaptype{none} % do not increment counter
\begin{longtable}[]{@{}
  >{\raggedright\arraybackslash}p{(\linewidth - 8\tabcolsep) * \real{0.1489}}
  >{\raggedright\arraybackslash}p{(\linewidth - 8\tabcolsep) * \real{0.2766}}
  >{\raggedright\arraybackslash}p{(\linewidth - 8\tabcolsep) * \real{0.1277}}
  >{\raggedright\arraybackslash}p{(\linewidth - 8\tabcolsep) * \real{0.1702}}
  >{\raggedright\arraybackslash}p{(\linewidth - 8\tabcolsep) * \real{0.2766}}@{}}
\toprule\noalign{}
\begin{minipage}[b]{\linewidth}\raggedright
Feature
\end{minipage} & \begin{minipage}[b]{\linewidth}\raggedright
Carbohydrates
\end{minipage} & \begin{minipage}[b]{\linewidth}\raggedright
Lipids
\end{minipage} & \begin{minipage}[b]{\linewidth}\raggedright
Proteins
\end{minipage} & \begin{minipage}[b]{\linewidth}\raggedright
Nucleic Acids
\end{minipage} \\
\midrule\noalign{}
\endhead
\bottomrule\noalign{}
\endlastfoot
\textbf{Monomer} & Monosaccharide & Fatty acid / isoprene & Amino acid &
Nucleotide \\
\textbf{Bond} & Glycosidic & Ester / noncovalent & Peptide &
Phosphodiester \\
\textbf{Elements} & C, H, O & C, H, O (+ P, N) & C, H, O, N, S & C, H,
O, N, P \\
\textbf{Energy (kcal/g)} & 4 & 9 & 4 & --- \\
\textbf{Informational?} & Somewhat (glycan code) & No & Partially
(sequence) & Yes (genetic code) \\
\textbf{Typical size} & 1--10\(^6\) residues & 14--24C chains &
50--30,000 residues & 70--\(10^9\) nt \\
\textbf{Water solubility} & Most soluble & Insoluble & Variable &
Soluble (highly charged) \\
\end{longtable}
}

\subsubsection{Storage vs.~Structural
Polysaccharides}\label{storage-vs.-structural-polysaccharides}

{\def\LTcaptype{none} % do not increment counter
\begin{longtable}[]{@{}
  >{\raggedright\arraybackslash}p{(\linewidth - 4\tabcolsep) * \real{0.1148}}
  >{\raggedright\arraybackslash}p{(\linewidth - 4\tabcolsep) * \real{0.4098}}
  >{\raggedright\arraybackslash}p{(\linewidth - 4\tabcolsep) * \real{0.4754}}@{}}
\toprule\noalign{}
\begin{minipage}[b]{\linewidth}\raggedright
Feature
\end{minipage} & \begin{minipage}[b]{\linewidth}\raggedright
Starch/Glycogen (Storage)
\end{minipage} & \begin{minipage}[b]{\linewidth}\raggedright
Cellulose/Chitin (Structural)
\end{minipage} \\
\midrule\noalign{}
\endhead
\bottomrule\noalign{}
\endlastfoot
Bond type & \texorpdfstring{\ensuremath{\alpha}}{alpha}-glycosidic & \texorpdfstring{\ensuremath{\beta}}{beta}-glycosidic \\
Chain conformation & Helical (amylose) or branched & Extended, straight
chains \\
Inter-chain interactions & Weak & Strong H-bonds between chains \\
Digestibility (humans) & Yes (\texorpdfstring{\ensuremath{\alpha}}{alpha}-amylase) & No (no \texorpdfstring{\ensuremath{\beta}}{beta}-glucosidase) \\
Physical properties & Granular, amorphous & Fibrous, high tensile
strength \\
\end{longtable}
}

\begin{center}\rule{0.5\linewidth}{0.5pt}\end{center}

\subsection{Computational Bridge}\label{computational-bridge-2}

Base composition of nucleic acids feeds directly into melting
temperature and PCR design. The genetics module exposes a simple GC
fraction:

\begin{Shaded}
\begin{Highlighting}[]
\ImportTok{from}\NormalTok{ biology.genetics }\ImportTok{import}\NormalTok{ gc\_content}

\NormalTok{seq }\OperatorTok{=} \StringTok{"ATGCGCGCATAA"}
\BuiltInTok{print}\NormalTok{(}\BuiltInTok{round}\NormalTok{(gc\_content(seq), }\DecValTok{3}\NormalTok{))  }\CommentTok{\# 0.583}
\end{Highlighting}
\end{Shaded}

\begin{quote}
\textbf{Clinical / systems note:} Antisense oligonucleotide and siRNA
therapies rely on predictable hybridisation thermodynamics (GC content,
length, salt) to avoid off-target binding while silencing disease
\hyperref[gl:allele]{\textbf{allele}}s --- the same sequence statistics
you summarise with \(G{+}C\) fraction.
\end{quote}

\begin{center}\rule{0.5\linewidth}{0.5pt}\end{center}

\subsection{Current Evidence and Frontier
Biology}\label{current-evidence-and-frontier-biology-5}

For \textbf{Biological Macromolecules}, frontier biology belongs inside
the evidence logic of the chapter. Chemistry-of-life claims now connect
classical bonding and thermodynamics with AI-guided structure prediction
and experimental validation. The core reading question is this:
macromolecule explanations should connect sequence, structure, dynamics,
modification, interaction, and assay evidence.

\begin{itemize}
\tightlist
\item
  \textbf{What to verify:} identify the observation, model, assay, or
  dataset that would make the claim stronger or weaker.
\item
  \textbf{What to qualify:} state the scale, organism, cell type,
  environmental condition, or population where the claim is expected to
  hold.
\item
  \textbf{What to compare:} test at least one alternative explanation,
  baseline, or null model before treating the pattern as causal.
\item
  \textbf{What to cite:} distinguish primary evidence, review synthesis,
  public dataset, and institutional guidance; for recent or numeric
  claims, prefer the source closest to the measurement and state what
  has changed since it was published.
\end{itemize}

Use AI biomolecular models as hypothesis generators: compare confidence,
conservation, solvent exposure, and assay evidence before turning a
predicted contact into a biological claim
\citep{abramson2024alphafold3}.

\textbf{Source practice:} For structure and interaction claims, cite
experimental structures when available and treat AlphaFold 3 or AFDB
complex predictions as hypotheses to validate with confidence metrics,
conservation, mutagenesis, binding, or cryo-EM/X-ray/NMR evidence
\citep{abramson2024alphafold3,velankar2026alphafolddb2025,emblebi2026alphafoldcomplexes}.

\subsubsection{Current Evidence Map: AI Structure Claims Need
Validation}\label{current-evidence-map-ai-structure-claims-need-validation}

\begin{figure}[htbp]
\centering
\includegraphics[width=0.82\linewidth,height=4.2in,keepaspectratio]{../figures/mermaid_inline/inline_mermaid_0025_b99c1ee97911.png}
\caption{AI structure models are strongest when confidence, interface geometry, conservation, mutagenesis, and experimental assays converge rather than when a model is treated as final evidence \textbackslash{}citep\{abramson2024alphafold3,velankar2026alphafolddb2025,emblebi2026alphafoldcomplexes\}.}
\end{figure}

\subsection{Summary}\label{summary-5}

\begin{itemize}
\tightlist
\item
  The four biological macromolecule classes (carbohydrates, lipids,
  proteins, nucleic acids) perform most structural and functional roles
  in living systems.
\item
  Carbohydrates are linked by glycosidic bonds; proteins by peptide
  bonds; nucleic acids by phosphodiester bonds. Most built/broken by
  condensation/hydrolysis.
\item
  Carbohydrate digestion proceeds from mouth (salivary amylase) through
  intestine (pancreatic amylase + brush border enzymes); cellulose
  passes undigested to the colon.
\item
  Lipids are diverse: glycerophospholipids and sphingolipids form
  membranes; triacylglycerols store energy; sterols regulate membranes
  and serve as hormone precursors; eicosanoids mediate inflammation.
\item
  Protein structure has four hierarchical levels; function depends on 3D
  structure, which is lost upon denaturation. Intrinsically disordered
  proteins challenge the classical structure-function paradigm.
\item
  DNA exists in A, B, and Z forms; B-form dominates under physiological
  conditions. RNA forms complex secondary structures including
  pseudoknots and G-quadruplexes.
\item
  DNA's double helix uses complementary A--T and G--C pairing; this
  complementarity is the structural basis of heredity and information
  transfer.
\item
  Ribozymes (catalytic RNAs) demonstrate RNA's dual role as information
  carrier and catalyst, supporting the RNA World hypothesis. The
  ribosome itself is a ribozyme.
\item
  \textbf{Connections:} See
  \cref{sec:unit_IV_dna_replication_and_cell_cycle} and
  \cref{sec:unit_IV_gene_expression} for replication,
  \hyperref[gl:transcription]{\textbf{transcription}}, and translation;
  \cref{sec:unit_V_mendelian_genetics} for segregation of alleles; and
  \cref{sec:unit_I_enzymes_and_kinetics} for enzymes that polymerise and
  hydrolyse these macromolecules.
\end{itemize}

\subsection{Key Terms}\label{key-terms-5}

\begin{itemize}
\tightlist
\item
  \textbf{Condensation reaction:} Joining of monomers with loss of water
\item
  \textbf{Hydrolysis:} Cleavage of a bond by addition of water
\item
  \textbf{Glycosidic bond:} Covalent bond between monosaccharides
\item
  \textbf{Anomer:} Stereoisomers differing at the anomeric carbon (\texorpdfstring{\ensuremath{\alpha}}{alpha}
  vs.~\texorpdfstring{\ensuremath{\beta}}{beta})
\item
  \textbf{Glycoprotein:} Protein with covalently attached
  oligosaccharide chains
\item
  \textbf{Proteoglycan:} Protein with long glycosaminoglycan chains
\item
  \textbf{Amphipathic:} Having both hydrophilic and hydrophobic regions
\item
  \textbf{Sphingolipid:} Lipid built on a sphingosine backbone
\item
  \textbf{Eicosanoid:} Signalling lipid derived from arachidonic acid
\item
  \textbf{Peptide bond:} Amide bond between amino acids; partial
  double-bond character
\item
  \textbf{Ramachandran plot:} Map of allowed backbone dihedral angles \texorpdfstring{\ensuremath{\varphi}}{phi}
  and \texorpdfstring{\ensuremath{\psi}}{psi}
\item
  \textbf{Denaturing:} Loss of higher-order protein structure
\item
  \textbf{Intrinsically disordered protein (IDP):} Protein lacking fixed
  tertiary structure
\item
  \textbf{Amyloid fibril:} Ordered cross-\texorpdfstring{\ensuremath{\beta}}{beta} aggregate of misfolded
  protein
\item
  \textbf{Chargaff's rules:} \hyperref[a]{A} = \hyperref[t]{T} and
  \hyperref[g]{G} = \hyperref[c]{C} in double-stranded DNA
\item
  \textbf{Base stacking:} \texorpdfstring{\ensuremath{\pi}}{pi}-\texorpdfstring{\ensuremath{\pi}}{pi} interactions between adjacent bases in
  nucleic acids
\item
  \textbf{Ribozyme:} Catalytic RNA molecule
\item
  \textbf{RNA World:} Hypothesis that RNA preceded DNA and proteins in
  early life
\item
  \textbf{Pseudoknot:} RNA structure where loop bases pair with
  downstream sequence
\end{itemize}

\subsection{Review Questions}\label{review-questions-5}

\begin{enumerate}
\def\labelenumi{\arabic{enumi}.}
\tightlist
\item
  Draw the condensation reaction between two glucose molecules to form
  maltose. Label the glycosidic bond and identify whether it is \texorpdfstring{\ensuremath{\alpha}}{alpha} or \texorpdfstring{\ensuremath{\beta}}{beta}.
\item
  Explain why cellulose is indigestible by humans but starch is not,
  despite both being polymers of glucose.
\item
  A patient with suspected lactose intolerance undergoes a hydrogen
  breath test. Explain the biochemical basis of this test.
\item
  Compare the membrane structures formed by glycerophospholipids and
  sphingolipids. Why is sphingomyelin particularly abundant in the
  myelin sheath?
\item
  A protein biochemist denatures ribonuclease A with urea and
  \texorpdfstring{\ensuremath{\beta}}{beta}-mercaptoethanol, then slowly removes both denaturants. The enzyme
  regains full activity. What does this tell us about the information
  needed for protein folding?
\item
  Explain why G--C rich DNA has a higher melting temperature than A--T
  rich DNA. Include both hydrogen bonding and base stacking in your
  answer.
\item
  Compare the structural features of B-form and Z-form DNA. Under what
  conditions might Z-DNA form, and what biological significance might it
  have?
\item
  The ribosome is often called a ``ribozyme.'' What evidence supports
  this claim, and what are its implications for the RNA World
  hypothesis?
\item
  Intrinsically disordered proteins defy the classical
  structure-function paradigm. Describe two biological advantages of
  protein disorder and give a clinical example of disease caused by IDP
  aggregation.
\item
  A nutritionist claims that ``a calorie from fat is more energy-dense
  than a calorie from carbohydrate.'' Explain the biochemical basis for
  the statement that fats provide 9 kcal/g vs.~4 kcal/g for
  carbohydrates, considering the oxidation states of carbon.
\item
  For a 24 bp primer with \texttt{gc\_content} = 0.50, explain why
  raising annealing temperature in PCR reduces off-target amplification.
  How would raising GC to 0.75 change your Tm estimate qualitatively?
\item
  Compare \textbf{one} way lipids and proteins each contribute to
  membrane curvature generation during vesicular trafficking. Tie your
  answer to amphipathic molecular architecture.
\item
  Use \cref{eq:unit_I_protein_folding_dG} and the table of folding
  contributions to estimate the change in \(\Delta G_{\text{fold}}\)
  when a single buried Ile (about 140 \AA{}\(^2\) of nonpolar
  surface) is mutated to Ala (about 70 \AA{}\(^2\)). Take 0.10
  kJ mol\(^{-1}\) \AA{}\(^{-2}\) as the cost of exposing nonpolar
  surface. Predict whether the mutant will be folded at body
  temperature.
\item
  Sketch a Ramachandran plot showing the two main allowed regions (\texorpdfstring{\ensuremath{\alpha}}{alpha} and
  \texorpdfstring{\ensuremath{\beta}}{beta}) for a non-Gly, non-Pro residue. Mark where collagen's
  polyproline-II helix appears and explain why \texorpdfstring{\ensuremath{\alpha}}{alpha}-helix is forbidden
  after a proline.
\item
  Two sequences of equal length differ in GC content: 30\% vs.~70\%.
  Using the nearest-neighbour stacking table above, predict which has
  the higher melting temperature and estimate the difference in
  \(\Delta G\) (37 \(^\circ\)C) per 10 base-pair stretch.
\item
  Liquid--liquid phase separation forms the nucleolus, P-bodies, and
  stress granules. Explain why \textbf{multivalent} weak interactions
  are essential and why a protein with a single high-affinity binding
  site for its partner does not phase separate.
\item
  List three reasons why cells use post-translational modifications
  (PTMs) rather than relying solely on transcription/translation to
  regulate protein activity. Include time scales, energetic cost, and
  reversibility in your answer.
\end{enumerate}

\subsection{Further Reading and Source
Notes}\label{further-reading-and-source-notes-5}

\begin{itemize}
\tightlist
\item
  Pauling, Corey \& Branson (1951). The structure of proteins: Two
  hydrogen-bonded helical configurations of the polypeptide chain.
  \emph{Proceedings of the National Academy of Sciences}, 37.
\item
  Anfinsen (1973). Principles that govern the folding of protein chains.
  \emph{Science}, 181.
\item
  Watson \& Crick (1953). Molecular structure of nucleic acids: A
  structure for deoxyribose nucleic acid. \emph{Nature}, 171.
\item
  Kornberg (1974). \emph{Chromatin Structure: A Repeating Unit of
  Histones and DNA}. \emph{Science}, 184.
\item
  Fischer (1894). Einfluss der Configuration auf die Wirkung der Enzyme.
  \emph{Berichte der deutschen chemischen Gesellschaft}, 27.
\item
  Lehninger, Nelson \& Cox (latest ed.). \emph{Lehninger Principles of
  Biochemistry} (chapters on carbohydrates, lipids, proteins, and
  nucleotides). W. H. Freeman.
\item
  Banani, Lee, Hyman \& Rosen (2017). Biomolecular condensates:
  Organizers of cellular biochemistry. \emph{Nature Reviews Molecular
  Cell Biology}, 18.
\end{itemize}

\begin{center}\rule{0.5\linewidth}{0.5pt}\end{center}

\subsubsection{Companion Source Module}\label{companion-source-module-5}

\textbf{Biological Macromolecules} should leave a reproducible trail
from a biological claim to the code, figure, diagram, or paper-based
activity that can test it. Use the surfaces below to inspect the
chapter's assumptions, rerun the relevant model, or compare the
manuscript explanation with companion labs and figures.

{\def\LTcaptype{none} % do not increment counter
\begin{longtable}[]{@{}
  >{\raggedright\arraybackslash}p{(\linewidth - 2\tabcolsep) * \real{0.5000}}
  >{\raggedright\arraybackslash}p{(\linewidth - 2\tabcolsep) * \real{0.5000}}@{}}
\toprule\noalign{}
\begin{minipage}[b]{\linewidth}\raggedright
Surface
\end{minipage} & \begin{minipage}[b]{\linewidth}\raggedright
Use it for
\end{minipage} \\
\midrule\noalign{}
\endhead
\bottomrule\noalign{}
\endlastfoot
\breaktt{src/biology/biochemistry/biochemistry.py}
(\breaktt{reaction\_free\_energy}, \breaktt{glycolysis\_summary}) & Relate
polymer chemistry and hydrolysis to energy flow. \\
\breaktt{src/biology/genetics/genetics.py} (\texttt{dna\_complement},
\breaktt{transcribe\_dna\_to\_mrna}, \texttt{translate\_mrna}) & Connect
nucleic-acid structure to information transfer. \\
\breaktt{src/mermaid/biology\_diagrams.py}
(\breaktt{macromolecule\_classification\_diagram},
\breaktt{transcription\_translation\_diagram}) & Compare classification
diagrams with sequence-to-function pathways. \\
\end{longtable}
}

\textbf{Reproducibility check:} separate sequence, three-dimensional
structure, modification state, and assay evidence before claiming
function. \textbf{Cross-reference:} use this bridge with
\cref{sec:unit_I_enzymes_and_kinetics},
\cref{sec:unit_IV_gene_expression}, and
\cref{sec:unit_IV_mutations_and_genomics}.

\newpage

\section{Enzymes and the Kinetics of
Catalysis}\label{enzymes-and-the-kinetics-of-catalysis}

\label{sec:unit_I_enzymes_and_kinetics}

\begin{quote}
\textbf{Ch 4} \texorpdfstring{\ensuremath{\cdot}}{.} Level 3/3 \texorpdfstring{\ensuremath{\cdot}}{.} 60 min read \texorpdfstring{\ensuremath{\cdot}}{.} 75 min lecture \texorpdfstring{\ensuremath{\cdot}}{.}
Prerequisites: \cref{sec:unit_I_macromolecules}
\end{quote}

\subsection{Learning Objectives}\label{learning-objectives-6}

\begin{enumerate}
\def\labelenumi{\arabic{enumi}.}
\tightlist
\item
  Describe how \hyperref[gl:enzyme]{\textbf{enzyme}}s lower activation
  energy without changing the
  \hyperref[gl:thermodynamics]{\textbf{thermodynamics}} of a reaction.
\item
  Derive and interpret the Michaelis-Menten equation
  (\cref{eq:unit_I_enzymes_and_kinetics_item_7}) and its kinetic
  parameters using the steady-state assumption.
\item
  Explain competitive, uncompetitive, mixed, and irreversible inhibition
  with mathematical descriptions and clinical examples.
\item
  Describe \hyperref[gl:allosteric]{\textbf{allosteric}} regulation and
  cooperativity using the Hill equation, with ATCase as a model system.
\item
  Explain how \hyperref[gl:ph]{\textbf{pH}}, temperature, and cofactors
  modulate enzyme activity.
\item
  Classify enzymes using the EC numbering system, including the seventh
  class (translocases).
\item
  Describe modern applications of enzymes in medicine and industry.
\end{enumerate}

\subsubsection{Study Blueprint}\label{study-blueprint-6}

\begin{itemize}
\tightlist
\item
  \textbf{Big idea:} Enzymes accelerate reactions by stabilizing
  transition states without changing thermodynamic endpoints.
\item
  \textbf{Core concepts:} activation energy, Michaelis-Menten kinetics,
  inhibition, allostery.
\item
  \textbf{Framework alignment:} Vision \& Change: Structure and
  function, Pathways and transformations of energy and matter; AP
  Biology: Energetics, Systems Interactions; NGSS-style topics: Matter
  and Energy in Organisms and Ecosystems, Structure and Function.
\item
  \textbf{Model or quantitative lens:} Michaelis-Menten,
  Lineweaver-Burk, and inhibition-pattern calculations.
\item
  \textbf{Data skill:} Fit or interpret enzyme-rate data and identify
  which parameter changed.
\item
  \textbf{Practice cadence:} Concept Explanation, Statistical Tests and
  Data Analysis, Argumentation.
\item
  \textbf{Common misconception to repair:} A catalyst changes rate, not
  the equilibrium constant or the sign of delta G.
\item
  \textbf{Primary lab:} \cref{sec:lab_unit_I_enzymes_and_kinetics}.
\item
  \textbf{Question bank:} \cref{sec:q_unit_I_enzymes_and_kinetics}.
\item
  \textbf{Transfer task:} Apply enzyme-kinetic reasoning to drug dosing,
  metabolic control, or diagnostic assays.
\item
  \textbf{Bridge to computation:}
  \breaktt{biology.biochemistry.biochemistry.michaelis\_menten}.
\end{itemize}

\begin{center}\rule{0.5\linewidth}{0.5pt}\end{center}

\begin{quote}
\textbf{Opening Vignette: The Enzyme That Beats Geological Time}

Without the enzyme OMP decarboxylase, the spontaneous decarboxylation of
orotidine 5\texorpdfstring{'}{'}-monophosphate (OMP) would take approximately \textbf{78
million years} at body temperature --- longer than the entire duration
of dinosaur dominance of land vertebrates. OMP decarboxylase
accomplishes the same reaction in 25 milliseconds. That is a rate
enhancement of approximately \(10^{23}\)-fold --- the largest known for
any enzyme (Radzicka \& Wolfenden, 1995, \emph{Science}).

The mechanism is not conventional. OMP decarboxylase uses no metal ions,
no cofactors, and performs general acid-base catalysis primarily through
electrostatic destabilization of the substrate's ground state rather
than direct stabilisation of the transition state. Understanding how it
achieves such a feat challenged enzymologists for decades and continues
to inform the design of enzyme inhibitors as antifungal and
antiprotozoal drugs. OMP decarboxylase is the essential enzyme in the
pyrimidine biosynthesis pathway --- block it, and pathogens cannot make
the \hyperref[gl:nucleotide]{\textbf{nucleotide}}s they need to
replicate.

\emph{Primary source: Radzicka, A. \& Wolfenden, R. (1995). A proficient
enzyme. Science, 267(5194), 90--93.}
\end{quote}

\begin{center}\rule{0.5\linewidth}{0.5pt}\end{center}

A \textbf{catalyst} accelerates a chemical reaction without being
consumed and without altering the reaction equilibrium. \textbf{Enzymes}
are biological catalysts --- almost typically
\hyperref[gl:protein]{\textbf{protein}}s (a few are ribozymes; see
\cref{sec:unit_I_macromolecules}) --- that accelerate reaction rates by
factors of \(10^6\)--\(10^{23}\) compared with uncatalysed reaction
rates.

The free energy change of a reaction (\(\Delta G\)) dictates whether a
reaction is thermodynamically spontaneous (\(\Delta G < 0\)) or
non-spontaneous (\(\Delta G > 0\)). Enzymes do \textbf{not} alter
\(\Delta G\) --- they alter the \textbf{kinetics} by providing an
alternative mechanism with a \textbf{lower activation energy}
(\(\Delta G^{\ddagger}\)).

\textbf{Remarkable rate enhancements:}

{\def\LTcaptype{none} % do not increment counter
\begin{longtable}[]{@{}
  >{\raggedright\arraybackslash}p{(\linewidth - 6\tabcolsep) * \real{0.0984}}
  >{\raggedright\arraybackslash}p{(\linewidth - 6\tabcolsep) * \real{0.3279}}
  >{\raggedright\arraybackslash}p{(\linewidth - 6\tabcolsep) * \real{0.3115}}
  >{\raggedright\arraybackslash}p{(\linewidth - 6\tabcolsep) * \real{0.2623}}@{}}
\toprule\noalign{}
\begin{minipage}[b]{\linewidth}\raggedright
Enzyme
\end{minipage} & \begin{minipage}[b]{\linewidth}\raggedright
Uncatalysed \(t_{1/2}\)
\end{minipage} & \begin{minipage}[b]{\linewidth}\raggedright
Catalysed \(t_{1/2}\)
\end{minipage} & \begin{minipage}[b]{\linewidth}\raggedright
Rate Enhancement
\end{minipage} \\
\midrule\noalign{}
\endhead
\bottomrule\noalign{}
\endlastfoot
OMP decarboxylase & 78 million years & 25 ms & \(10^{23}\) \\
Staphylococcal nuclease & 130,000 years & about 1 ms &
\(10^{17}\) \\
Alkaline phosphatase & 7 years & about 1 ms & \(10^{11}\) \\
Carbonic anhydrase & 5 seconds & about 1 \texorpdfstring{\ensuremath{\mu}}{mu}s & \(10^7\) \\
Chorismate mutase & 7.4 hours & about 1 ms & \(10^7\) \\
\end{longtable}
}

OMP decarboxylase achieves the largest known rate enhancement: without
this enzyme, the spontaneous decarboxylation of orotidine
5'-monophosphate would take 78 million years --- longer than the age of
the dinosaurs.

\subsubsection{Transition State Theory}\label{transition-state-theory}

Every reaction proceeds through a \textbf{transition state} (TS) --- a
transient, high-energy configuration that cannot be isolated. The
\textbf{activation energy} (\(E_a\) or \(\Delta G^{\ddagger}\)) is the
energy required to reach this state from the ground state:

\[\text{Rate} = A \cdot e^{-\Delta G^{\ddagger}/RT} \tag{4.1} \label{eq:unit_I_enzymes_and_kinetics_item_1}\]

where \(A\) is the pre-exponential factor. Lowering
\(\Delta G^{\ddagger}\) by just 17 kJ/mol increases the rate by
\(e^7 \approx 1{,}000\)-fold. Lowering it by 34 kJ/mol gives a
million-fold increase.

\textbf{How enzymes lower \(\Delta G^{\ddagger}\):}

\begin{enumerate}
\def\labelenumi{\arabic{enumi}.}
\tightlist
\item
  \textbf{Proximity and orientation:} Substrates are brought together in
  optimal geometry at the active site, with the entropic cost already
  paid during binding. Effective concentration of reactants in the
  active site can exceed \(10^8\) M.
\item
  \textbf{Transition state stabilisation:} The enzyme active site binds
  the transition state \emph{more tightly} than either substrate or
  product (complementarity to TS, not substrate). This is the single
  most important mechanism.
\item
  \textbf{General acid-base catalysis:} Amino acid side chains (His,
  Asp, Glu, Lys) donate/accept protons simultaneously with bond
  breaking/forming, stabilising developing charges in the TS.
\item
  \textbf{Covalent catalysis:} Transient covalent enzyme-substrate
  intermediates (Ser proteases, Cys proteases, Lys in Schiff base
  enzymes). The covalent intermediate provides a lower-energy pathway.
\item
  \textbf{Metal ion catalysis:} Metal ions (Zn\(^{2+}\), Mg\(^{2+}\),
  Fe\(^{2+/3+}\)) stabilise negative charges, act as Lewis acids, or
  enable redox chemistry.
\item
  \textbf{Electrostatic catalysis:} The active site creates a
  microenvironment with a lower effective dielectric constant, enhancing
  electrostatic interactions.
\item
  \textbf{Desolvation:} Stripping water from substrate and active-site
  residues increases their reactivity (e.g., a ``naked'' carboxylate is
  a much stronger base than a hydrated one).
\end{enumerate}

\begin{quote}
\textbf{Concept Check 1:} Transition state analogues --- molecules that
mimic the transition state geometry --- are potent enzyme inhibitors.
Why are they typically much tighter binders than the substrate itself?
(Hint: Consider what the active site is optimised to bind.)
\end{quote}

\begin{figure}[htbp]
\centering
\includegraphics[width=0.82\linewidth,height=4.2in,keepaspectratio]{../figures/mermaid_inline/inline_mermaid_0026_ffd77fe37203.png}
\caption{Models of enzyme--substrate binding. Evolution of enzyme-substrate binding models. The lock-and-key model (1894) proposed rigid complementarity. The induced-fit model (1958) introduced conformational change upon binding. The modern conformational selection model recognises that enzymes sample multiple conformations, and substrates select the most complementary one.}
\end{figure}

\subsubsection{Active Site Architecture}\label{active-site-architecture}

The \textbf{active site} occupies roughly 1--10\% of the enzyme's total
surface but accounts for most catalytic power. Key features:

\begin{itemize}
\tightlist
\item
  \textbf{Specific geometry} complementary to the transition state (more
  so than to the substrate)
\item
  \textbf{Hydrophobic microenvironment} that enhances nucleophilicity
  and acid/base strength
\item
  \textbf{Flexibility:} conformational changes upon substrate binding
  (induced fit model)
\item
  \textbf{Conserved residues:} catalytic residues are highly conserved
  across species, even when surrounding sequences diverge
\end{itemize}

The \textbf{lock-and-key model} \citep{fischer1894} viewed the active
site as rigid. The \textbf{induced-fit model} \citep{koshland1958}
recognised that substrate binding induces conformational changes that
improve active-site complementarity --- validated by crystallographic
structures. The modern \textbf{conformational selection model} proposes
that the enzyme pre-exists in an ensemble of conformations, and the
substrate selects the optimal one, shifting the equilibrium.

\textbf{Case study --- Hexokinase:} X-ray crystallography reveals that
hexokinase undergoes a dramatic conformational change upon glucose
binding: two lobes of the enzyme close around the substrate like a jaw,
excluding water from the active site. This prevents the wasteful
hydrolysis of ATP (which would occur if water could access the
\texorpdfstring{\ensuremath{\gamma}}{gamma}-phosphate).

\begin{center}\rule{0.5\linewidth}{0.5pt}\end{center}

\subsection{Enzyme Classification (EC
Numbers)}\label{enzyme-classification-ec-numbers}

The International Union of Biochemistry and Molecular Biology (IUBMB)
classifies enzymes into \textbf{seven classes} by reaction type, each
with a four-digit EC (Enzyme Commission) number:

{\def\LTcaptype{none} % do not increment counter
\begin{longtable}[]{@{}
  >{\raggedright\arraybackslash}p{(\linewidth - 8\tabcolsep) * \real{0.1951}}
  >{\raggedright\arraybackslash}p{(\linewidth - 8\tabcolsep) * \real{0.0976}}
  >{\raggedright\arraybackslash}p{(\linewidth - 8\tabcolsep) * \real{0.3171}}
  >{\raggedright\arraybackslash}p{(\linewidth - 8\tabcolsep) * \real{0.1707}}
  >{\raggedright\arraybackslash}p{(\linewidth - 8\tabcolsep) * \real{0.2195}}@{}}
\toprule\noalign{}
\begin{minipage}[b]{\linewidth}\raggedright
EC Class
\end{minipage} & \begin{minipage}[b]{\linewidth}\raggedright
Name
\end{minipage} & \begin{minipage}[b]{\linewidth}\raggedright
Reaction Type
\end{minipage} & \begin{minipage}[b]{\linewidth}\raggedright
Example
\end{minipage} & \begin{minipage}[b]{\linewidth}\raggedright
EC Number
\end{minipage} \\
\midrule\noalign{}
\endhead
\bottomrule\noalign{}
\endlastfoot
1 & Oxidoreductases & Redox reactions & Lactate dehydrogenase & EC
1.1.1.27 \\
2 & Transferases & Group transfer & Hexokinase (phosphoryl) & EC
2.7.1.1 \\
3 & Hydrolases & Hydrolysis & Trypsin & EC 3.4.21.4 \\
4 & Lyases & Non-hydrolytic bond cleavage / addition to double bonds &
Fumarase & EC 4.2.1.2 \\
5 & Isomerases & Isomerisation & Phosphoglucose isomerase & EC
5.3.1.9 \\
6 & Ligases & Bond formation coupled to ATP/GTP hydrolysis &
\hyperref[gl:dna-ligase]{\textbf{DNA ligase}} & EC 6.5.1.1 \\
7 & \textbf{Translocases} & Movement of ions/molecules across membranes
& Na\(^+\)/K\(^+\)-ATPase & EC 7.2.2.6 \\
\end{longtable}
}

The \textbf{seventh class (translocases)} was added in 2018, recognising
the catalytic nature of active transport. Enzyme names follow the
pattern: \textbf{Substrate(s) + reaction type + ``-ase''} (e.g.,
pyruvate kinase, lactate dehydrogenase).

\begin{quote}
\textbf{Concept Check 2:} The Na\(^+\)/K\(^+\)-ATPase pumps 3 Na\(^+\)
out and 2 K\(^+\) in per ATP hydrolysed. It was traditionally classified
as an ATPase (EC 3.6). Why was reclassification to translocase (EC 7)
considered more appropriate?
\end{quote}

\begin{center}\rule{0.5\linewidth}{0.5pt}\end{center}

\subsection{The Michaelis-Menten
Equation}\label{the-michaelis-menten-equation}

\subsubsection{Detailed Derivation}\label{detailed-derivation}

In 1913, Leonor Michaelis and Maud Menten proposed a kinetic framework
for enzyme-catalysed reactions. The basic mechanism:

\[\text{E} + \text{S} \underset{k_{-1}}{\overset{k_1}{\rightleftharpoons}} \text{ES} \overset{k_2}{\rightarrow} \text{E} + \text{P} \tag{4.2} \label{eq:unit_I_enzymes_and_kinetics_item_2}\]

\textbf{Step 1: Write rate equations.}

The rate of ES formation:
\(\frac{d[\text{ES}]}{dt} = k_1[\text{E}][\text{S}] - k_{-1}[\text{ES}] - k_2[\text{ES}]\)

\textbf{Step 2: Apply the steady-state assumption}
(\(d[\text{ES}]/dt = 0\)).

This assumes that after an initial transient, {[}ES{]} reaches a
constant level because the rate of ES formation equals the rate of ES
breakdown:

\[k_1[\text{E}][\text{S}] = (k_{-1} + k_2)[\text{ES}] \tag{4.3} \label{eq:unit_I_enzymes_and_kinetics_item_3}\]

\textbf{Step 3: Define the Michaelis constant.}

\[K_m = \frac{k_{-1} + k_2}{k_1} \tag{4.4} \label{eq:unit_I_enzymes_and_kinetics_item_4}\]

\textbf{Step 4: Express \hyperref[e]{E} in terms of measurable
quantities.}

Total enzyme: \([\text{E}]_T = [\text{E}] + [\text{ES}]\), so
\([\text{E}] = [\text{E}]_T - [\text{ES}]\)

Substituting:
\(([\text{E}]_T - [\text{ES}])[\text{S}] = K_m[\text{ES}]\)

\[[\text{ES}] = \frac{[\text{E}]_T[\text{S}]}{K_m + [\text{S}]} \tag{4.5} \label{eq:unit_I_enzymes_and_kinetics_item_5}\]

\textbf{Step 5: Calculate the initial velocity.}

\[v_0 = k_2[\text{ES}] = \frac{k_2[\text{E}]_T[\text{S}]}{K_m + [\text{S}]} \tag{4.6} \label{eq:unit_I_enzymes_and_kinetics_item_6}\]

Since \(V_{max} = k_2[\text{E}]_T\) (maximum rate when most enzyme is
saturated):

\[\boxed{v_0 = \frac{V_{max}[\text{S}]}{K_m + [\text{S}]}} \tag{4.7} \label{eq:unit_I_enzymes_and_kinetics_item_7}\]

This is the \textbf{Michaelis-Menten equation} --- a rectangular
hyperbola (\cref{eq:unit_I_enzymes_and_kinetics_item_7}); plotting
\(v_0\) against \([\text{S}]\) traces the saturating curve shown in
\cref{fig:unit_I_michaelis_menten}, where the velocity rises steeply at
low substrate and asymptotically approaches \(V_{max}\).

\subsubsection{Interpretation of Kinetic
Parameters}\label{interpretation-of-kinetic-parameters}

\textbf{\(K_m\)} approximates the substrate affinity of the enzyme. When
\hyperref[s]{S} = \(K_m\), \(v = V_{max}/2\). A \textbf{low \(K_m\)}
means high affinity (half-saturation at low \hyperref[s]{S}).

Special cases: - When \(k_2 \ll k_{-1}\):
\(K_m \approx K_d = k_{-1}/k_1\) (true dissociation constant) - When
\(k_2 \gg k_{-1}\): \(K_m \approx k_2/k_1\) (the enzyme is so fast that
ES rarely dissociates back to E + S)

\textbf{\(V_{max}\) and \(k_{cat}\):}
\(V_{max} = k_{cat} \times [\text{E}]_T\). The \textbf{catalytic
constant} \(k_{cat}\) (also called the turnover number) is the number of
substrate molecules converted to product per enzyme molecule per second
when fully saturated.

\textbf{Catalytic efficiency:} \(k_{cat}/K_m\) (units: M\(^{-1}\)
s\(^{-1}\)). This is the gold standard for comparing enzyme power. The
diffusion limit sets the maximum possible
\(k_{cat}/K_m \approx 10^8\)--\(10^9\) M\(^{-1}\) s\(^{-1}\)
(kinetically ``perfect'' enzymes).

\textbf{Reference kinetic constants:}

{\def\LTcaptype{none} % do not increment counter
\begin{longtable}[]{@{}
  >{\raggedright\arraybackslash}p{(\linewidth - 10\tabcolsep) * \real{0.0723}}
  >{\raggedright\arraybackslash}p{(\linewidth - 10\tabcolsep) * \real{0.1084}}
  >{\raggedright\arraybackslash}p{(\linewidth - 10\tabcolsep) * \real{0.1205}}
  >{\raggedright\arraybackslash}p{(\linewidth - 10\tabcolsep) * \real{0.2410}}
  >{\raggedright\arraybackslash}p{(\linewidth - 10\tabcolsep) * \real{0.3855}}
  >{\raggedright\arraybackslash}p{(\linewidth - 10\tabcolsep) * \real{0.0723}}@{}}
\toprule\noalign{}
\begin{minipage}[b]{\linewidth}\raggedright
Enzyme
\end{minipage} & \begin{minipage}[b]{\linewidth}\raggedright
Substrate
\end{minipage} & \begin{minipage}[b]{\linewidth}\raggedright
\(K_m\) (mM)
\end{minipage} & \begin{minipage}[b]{\linewidth}\raggedright
\(k_{cat}\) (s\(^{-1}\))
\end{minipage} & \begin{minipage}[b]{\linewidth}\raggedright
\(k_{cat}/K_m\) (M\(^{-1}\)s\(^{-1}\))
\end{minipage} & \begin{minipage}[b]{\linewidth}\raggedright
Status
\end{minipage} \\
\midrule\noalign{}
\endhead
\bottomrule\noalign{}
\endlastfoot
Carbonic anhydrase & CO\(_2\) & 12 & \(10^6\) & \(8.3 \times 10^7\) &
Near-perfect \\
Fumarase & Fumarate & 0.005 & 800 & \(1.6 \times 10^8\) &
Near-perfect \\
Acetylcholinesterase & Acetylcholine & 0.095 & \(1.4 \times 10^4\) &
\(1.5 \times 10^8\) & Near-perfect \\
Catalase & H\(_2\)O\(_2\) & 25 & \(4 \times 10^7\) & \(1.6 \times 10^9\)
& Diffusion-limited \\
Lactate dehydrogenase & Pyruvate & 0.047 & 800 & \(1.7 \times 10^7\) &
Highly efficient \\
Hexokinase & Glucose & 0.15 & 650 & \(4.3 \times 10^6\) & Efficient \\
Alcohol dehydrogenase & Ethanol & 1.0 & 73 & \(7.3 \times 10^4\) &
Moderate \\
Lysozyme & Hexasaccharide & 0.006 & 0.5 & \(8.3 \times 10^4\) &
Moderate \\
\end{longtable}
}

\begin{quote}
\textbf{Concept Check 3:} Catalase converts H\(_2\)O\(_2\) to H\(_2\)O
and O\(_2\) with a \(k_{cat}\) of \(4 \times 10^7\) s\(^{-1}\). How many
molecules of H\(_2\)O\(_2\) does a single catalase molecule destroy per
millisecond? Why is such extreme speed necessary in the cell?
\end{quote}

\subsubsection{Lineweaver-Burk Analysis}\label{lineweaver-burk-analysis}

\begin{figure}[htbp]
\centering
\includegraphics[width=0.85\textwidth]{../figures/michaelis_menten.png}
\caption{Michaelis--Menten kinetics: initial velocity ($v_0$) versus substrate concentration with an uninhibited curve, a competitive-inhibitor curve, and the $V_{\max}$, $V_{\max}/2$, and $K_m$ reference lines annotated.}
\label{fig:unit_I_michaelis_menten}
\end{figure}

Inverting the Michaelis-Menten equation:

\[\frac{1}{v_0} = \frac{K_m}{V_{max}} \cdot \frac{1}{[\text{S}]} + \frac{1}{V_{max}} \tag{4.8} \label{eq:unit_I_enzymes_and_kinetics_item_8}\]

The \textbf{double-reciprocal (Lineweaver-Burk) plot} of \(1/v\)
vs.~\(1/[\text{S}]\) is linear: - Slope = \(K_m/V_{max}\) -
\(x\)-intercept = \(-1/K_m\) - \(y\)-intercept = \(1/V_{max}\)

\textbf{Worked example --- Determining \(K_m\) and \(V_{max}\):}

An enzyme is assayed at five substrate concentrations:

{\def\LTcaptype{none} % do not increment counter
\begin{longtable}[]{@{}llll@{}}
\toprule\noalign{}
\hyperref[s]{S} (mM) & \(v_0\) (\texorpdfstring{\ensuremath{\mu}}{mu}mol/min) & 1/\hyperref[s]{S}
(mM\(^{-1}\)) & \(1/v_0\) \\
\midrule\noalign{}
\endhead
\bottomrule\noalign{}
\endlastfoot
0.5 & 1.67 & 2.0 & 0.60 \\
1.0 & 2.50 & 1.0 & 0.40 \\
2.0 & 3.33 & 0.5 & 0.30 \\
4.0 & 4.00 & 0.25 & 0.25 \\
10.0 & 4.55 & 0.10 & 0.22 \\
\end{longtable}
}

From the Lineweaver-Burk plot: - \(y\)-intercept = 0.20, so
\(V_{max} = 1/0.20 = 5.0\) \texorpdfstring{\ensuremath{\mu}}{mu}mol/min - Slope =
\((0.60 - 0.20)/(2.0 - 0) = 0.20\), and slope \(= K_m/V_{max}\), so
\(K_m = 0.20 \times 5.0 = 1.0\) mM

\textbf{Limitations of the Lineweaver-Burk plot:} It distorts
experimental error at low \hyperref[s]{S} (high 1/\hyperref[s]{S}),
giving those points excessive weight. Modern software uses nonlinear
regression to fit the Michaelis-Menten equation directly, which is
statistically superior.

\begin{center}\rule{0.5\linewidth}{0.5pt}\end{center}

\subsection{Enzyme Inhibition}\label{enzyme-inhibition}

\begin{figure}[htbp]
\centering
\includegraphics[width=0.82\linewidth,height=4.2in,keepaspectratio]{../figures/mermaid_inline/inline_mermaid_0027_1177901e6d0b.png}
\caption{State diagram for Enzyme Inhibition showing transitions among E exists in solution, S binds (active site), S dissociates, and Catalysis (k\_cat).}
\end{figure}

\subsubsection{Competitive Inhibition}\label{competitive-inhibition}

A \textbf{competitive inhibitor} (I) resembles the substrate
structurally and competes for the same active site. Binding of I and S
are mutually exclusive:

\[E + I \rightleftharpoons EI \quad (K_i = [\text{E}][\text{I}]/[\text{EI}]) \tag{4.9} \label{eq:unit_I_enzymes_and_kinetics_item_9}\]

The apparent \(K_m\) increases (less substrate affinity at fixed
\hyperref[i]{I}) while \(V_{max}\) is unchanged (inhibition overcome by
high \hyperref[s]{S}):

\[v_0 = \frac{V_{max}[\text{S}]}{K_m\left(1 + \frac{[\text{I}]}{K_i}\right) + [\text{S}]} \tag{4.10} \label{eq:unit_I_enzymes_and_kinetics_item_10}\]

Defining \(\alpha = 1 + [\text{I}]/K_i\):

\[v_0 = \frac{V_{max}[\text{S}]}{\alpha K_m + [\text{S}]} \tag{4.11} \label{eq:unit_I_enzymes_and_kinetics_item_11}\]

On a Lineweaver-Burk plot: same \(y\)-intercept (\(1/V_{max}\)),
different slopes. Lines intersect on the \(y\)-axis.

\textbf{Clinical examples of
\hyperref[gl:competitive-inhibition]{\textbf{competitive inhibition}}:}

{\def\LTcaptype{none} % do not increment counter
\begin{longtable}[]{@{}
  >{\raggedright\arraybackslash}p{(\linewidth - 4\tabcolsep) * \real{0.2647}}
  >{\raggedright\arraybackslash}p{(\linewidth - 4\tabcolsep) * \real{0.3824}}
  >{\raggedright\arraybackslash}p{(\linewidth - 4\tabcolsep) * \real{0.3529}}@{}}
\toprule\noalign{}
\begin{minipage}[b]{\linewidth}\raggedright
Inhibitor
\end{minipage} & \begin{minipage}[b]{\linewidth}\raggedright
Target Enzyme
\end{minipage} & \begin{minipage}[b]{\linewidth}\raggedright
Clinical Use
\end{minipage} \\
\midrule\noalign{}
\endhead
\bottomrule\noalign{}
\endlastfoot
Methotrexate & Dihydrofolate reductase & Cancer chemotherapy \\
Statins (atorvastatin) & HMG-CoA reductase & Cholesterol-lowering \\
Sulfonamide antibiotics & Dihydropteroate synthase & Bacterial
infection \\
Allopurinol & Xanthine oxidase & Gout \\
Ethanol & Alcohol dehydrogenase & Methanol/ethylene glycol poisoning \\
\end{longtable}
}

\begin{quote}
\textbf{Clinical Connection: Statins and Cardiovascular Disease}

Statins (e.g., atorvastatin, rosuvastatin) are competitive inhibitors of
HMG-CoA reductase, the rate-limiting enzyme in cholesterol biosynthesis.
By mimicking the structure of HMG-CoA, statins bind the active site with
\(K_i\) values in the low nanomolar range (much lower than the \(K_m\)
for HMG-CoA, about 4 \texorpdfstring{\ensuremath{\mu}}{mu}M). This reduces hepatic cholesterol
synthesis, upregulates LDL receptors, and lowers circulating LDL
cholesterol by 30--50\%. Statins are among the most widely prescribed
drugs worldwide, with robust evidence for reducing cardiovascular events
and mortality.
\end{quote}

\subsubsection{Uncompetitive Inhibition}\label{uncompetitive-inhibition}

An \textbf{uncompetitive inhibitor} binds primarily to the ES complex,
not free enzyme:

\[v_0 = \frac{V_{max}[\text{S}]}{K_m + \alpha'[\text{S}]} \tag{4.12} \label{eq:unit_I_enzymes_and_kinetics_item_12}\]

where \(\alpha' = 1 + [\text{I}]/K_i'\). Both \(K_m\) and \(V_{max}\)
are decreased proportionally --- the ratio \(V_{max}/K_m\) (=
\(k_{cat}/K_m\)) is unchanged. On a Lineweaver-Burk plot:
\textbf{parallel lines} (both slope and intercepts change equally).

Uncompetitive inhibition is rare in single-substrate reactions but
important in two-substrate reactions and in multi-enzyme pathways.
\textbf{Lithium} is an uncompetitive inhibitor of inositol
monophosphatase, which may contribute to its therapeutic effect in
bipolar disorder.

\subsubsection{Mixed (Noncompetitive)
Inhibition}\label{mixed-noncompetitive-inhibition}

A \textbf{mixed inhibitor} binds to both free enzyme (E) and the ES
complex, at a site other than the active site:

\[v_0 = \frac{V_{max}[\text{S}]}{\alpha K_m + \alpha'[\text{S}]} \tag{4.13} \label{eq:unit_I_enzymes_and_kinetics_item_13}\]

where \(\alpha = 1 + [\text{I}]/K_i\) and
\(\alpha' = 1 + [\text{I}]/K_i'\).

\begin{itemize}
\tightlist
\item
  If \(\alpha = \alpha'\) (\(K_i = K_i'\)): \textbf{pure noncompetitive}
  inhibition. \(K_m\) unchanged, \(V_{max}\) decreased.
\item
  If \(\alpha \neq \alpha'\): \textbf{mixed} inhibition. Both \(K_m\)
  and \(V_{max}\) change.
\end{itemize}

On a Lineweaver-Burk plot: lines intersect to the left of the \(y\)-axis
(not on it).

\textbf{Pure noncompetitive inhibitor:} Heavy metal ions (Hg\(^{2+}\),
Pb\(^{2+}\)) that bind cysteine residues distant from the active site.

\subsubsection{Irreversible Inhibition}\label{irreversible-inhibition}

\textbf{Irreversible inhibitors} form covalent bonds with enzyme
residues, permanently inactivating the enzyme. The enzyme must be
replaced by new protein synthesis.

{\def\LTcaptype{none} % do not increment counter
\begin{longtable}[]{@{}
  >{\raggedright\arraybackslash}p{(\linewidth - 6\tabcolsep) * \real{0.2045}}
  >{\raggedright\arraybackslash}p{(\linewidth - 6\tabcolsep) * \real{0.1364}}
  >{\raggedright\arraybackslash}p{(\linewidth - 6\tabcolsep) * \real{0.2045}}
  >{\raggedright\arraybackslash}p{(\linewidth - 6\tabcolsep) * \real{0.4545}}@{}}
\toprule\noalign{}
\begin{minipage}[b]{\linewidth}\raggedright
Inhibitor
\end{minipage} & \begin{minipage}[b]{\linewidth}\raggedright
Target
\end{minipage} & \begin{minipage}[b]{\linewidth}\raggedright
Mechanism
\end{minipage} & \begin{minipage}[b]{\linewidth}\raggedright
Clinical Application
\end{minipage} \\
\midrule\noalign{}
\endhead
\bottomrule\noalign{}
\endlastfoot
Aspirin (acetylsalicylate) & COX-1 and COX-2 & Acetylates Ser530 &
Anti-inflammatory, antiplatelet \\
Penicillin / amoxicillin & Transpeptidase (DD-peptidase) & Acylates
active-site Ser & Antibiotic \\
DIPF / sarin & Acetylcholinesterase & Phosphorylates active-site Ser &
Nerve agent (sarin) \\
Omeprazole (Prilosec) & H\(^+\)/K\(^+\)-ATPase & Disulfide bond with
Cys813 & Proton pump inhibitor (acid reflux) \\
Clopidogrel (Plavix) & P2Y12 receptor & Disulfide bond & Antiplatelet \\
\end{longtable}
}

\begin{quote}
\textbf{Clinical Connection: Aspirin --- The Oldest Irreversible Enzyme
Inhibitor}

Aspirin (acetylsalicylic acid) irreversibly acetylates Ser530 in
cyclooxygenase (COX), preventing the conversion of arachidonic acid to
prostaglandins and thromboxane A\(_2\). Because platelets lack a nucleus
and cannot synthesise new protein, a single dose of aspirin inhibits
platelet COX-1 for the entire platelet lifespan (about 10
days). This is why low-dose aspirin (75--100 mg/day) is an effective
antiplatelet agent for preventing heart attacks and strokes, despite its
short plasma half-life of about 20 minutes.
\end{quote}

\begin{quote}
\textbf{Clinical Connection: Penicillin and Bacterial Cell Wall
Synthesis}

Penicillin and related \texorpdfstring{\ensuremath{\beta}}{beta}-lactam antibiotics are irreversible inhibitors
of bacterial transpeptidase (also called penicillin-binding protein,
PBP), the enzyme that cross-links peptidoglycan strands in the bacterial
cell wall. The \texorpdfstring{\ensuremath{\beta}}{beta}-lactam ring of penicillin mimics the D-Ala--D-Ala
terminus of the peptidoglycan substrate, forming a stable acyl-enzyme
intermediate. Without cross-linking, the cell wall loses integrity and
the bacterium lyses due to osmotic stress. Resistance arises when
bacteria produce \texorpdfstring{\ensuremath{\beta}}{beta}-lactamases that hydrolyse the \texorpdfstring{\ensuremath{\beta}}{beta}-lactam ring.
\end{quote}

\subsubsection{Lineweaver-Burk Patterns
Summary}\label{lineweaver-burk-patterns-summary}

{\def\LTcaptype{none} % do not increment counter
\begin{longtable}[]{@{}
  >{\raggedright\arraybackslash}p{(\linewidth - 6\tabcolsep) * \real{0.2419}}
  >{\raggedright\arraybackslash}p{(\linewidth - 6\tabcolsep) * \real{0.2097}}
  >{\raggedright\arraybackslash}p{(\linewidth - 6\tabcolsep) * \real{0.2903}}
  >{\raggedright\arraybackslash}p{(\linewidth - 6\tabcolsep) * \real{0.2581}}@{}}
\toprule\noalign{}
\begin{minipage}[b]{\linewidth}\raggedright
Inhibition Type
\end{minipage} & \begin{minipage}[b]{\linewidth}\raggedright
Apparent \(K_m\)
\end{minipage} & \begin{minipage}[b]{\linewidth}\raggedright
Apparent \(V_{max}\)
\end{minipage} & \begin{minipage}[b]{\linewidth}\raggedright
L-B Plot Pattern
\end{minipage} \\
\midrule\noalign{}
\endhead
\bottomrule\noalign{}
\endlastfoot
Competitive & Increases (\(\alpha K_m\)) & Unchanged & Same
\(y\)-intercept \\
Uncompetitive & Decreases (\(K_m/\alpha'\)) & Decreases
(\(V_{max}/\alpha'\)) & Parallel lines \\
Pure noncompetitive & Unchanged & Decreases (\(V_{max}/\alpha\)) & Same
\(x\)-intercept \\
Mixed & Changes & Decreases & Intersect left of \(y\)-axis \\
Irreversible & N/A & Decreases (less \([\text{E}]_T\)) & Same \(K_m\),
lower \(V_{max}\) \\
\end{longtable}
}

\begin{quote}
\textbf{Concept Check 4:} An enzyme is studied in the presence and
absence of an inhibitor. In the Lineweaver-Burk plot, the two lines are
parallel. What type of inhibition is this? Does the inhibitor bind the
free enzyme, the ES complex, or both?
\end{quote}

\begin{center}\rule{0.5\linewidth}{0.5pt}\end{center}

\subsection{Multi-substrate Kinetics: Bi-Bi
Mechanisms}\label{multi-substrate-kinetics-bi-bi-mechanisms}

The single-substrate Michaelis-Menten framework is a useful
approximation, but most enzymes use \emph{two} substrates (or one
substrate plus a cofactor). Reactions of the form

\[\text{A} + \text{B} \rightleftharpoons \text{P} + \text{Q} \label{eq:unit_I_enzymes_and_kinetics_item_14}\]

are known as \textbf{Bi-Bi reactions}. Cleland (1963) classified them
into three principal mechanisms based on the order in which substrates
bind and products release.

\subsubsection{Sequential (Ternary Complex)
Mechanisms}\label{sequential-ternary-complex-mechanisms}

Both substrates must be bound to the enzyme \emph{simultaneously} before
any product is released. The enzyme thus passes through a ternary EAB
complex.

\begin{itemize}
\tightlist
\item
  \textbf{Ordered Bi-Bi:} A binds before B; Q releases before P.
  Example: NAD\(^+\)-dependent dehydrogenases such as lactate
  dehydrogenase, where NAD\(^+\) binds first and NADH leaves last.
\item
  \textbf{Random Bi-Bi:} A and B can bind in either order; products can
  leave in either order. Example: creatine kinase.
\end{itemize}

The general rate equation for a sequential Bi-Bi reaction (assuming
rapid equilibrium and saturating second substrate) reduces to:

\begin{equation}
v_0 = \frac{V_{\max}\,[\text{A}]\,[\text{B}]}{K_{ia}K_b + K_b[\text{A}] + K_a[\text{B}] + [\text{A}][\text{B}]}
\label{eq:unit_I_bibi_sequential}
\end{equation}

where \(K_a\) and \(K_b\) are Michaelis constants for A and B, and
\(K_{ia}\) is the dissociation constant of A from the binary EA complex.
\textbf{Diagnostic Lineweaver-Burk pattern:} plotting \(1/v_0\)
vs.~\(1/[\text{A}]\) at several fixed \hyperref[b]{B} gives a family of
lines that \textbf{intersect to the left of the y-axis} (sequential
signature).

\subsubsection{Ping-Pong (Double-Displacement)
Mechanism}\label{ping-pong-double-displacement-mechanism}

The first substrate A binds, transfers a chemical group to the enzyme,
and the first product P is released, leaving the enzyme in a covalently
modified form (E\emph{). Then B binds to E}, accepts the group, and Q is
released. \textbf{At no time does the enzyme contain both substrates
simultaneously.} Examples: aminotransferases (PLP carries the amino
group between substrates), serine proteases (substrate A acylates the
active-site Ser; substrate B = water deacylates), and many kinases that
proceed via a phospho-enzyme intermediate.

The rate equation:

\begin{equation}
v_0 = \frac{V_{\max}\,[\text{A}]\,[\text{B}]}{K_b[\text{A}] + K_a[\text{B}] + [\text{A}][\text{B}]}
\label{eq:unit_I_bibi_pingpong}
\end{equation}

Notice the absent \(K_{ia}K_b\) term in the denominator.
\textbf{Diagnostic Lineweaver-Burk pattern:} the family of lines is
\textbf{parallel} --- a clean kinetic signature distinguishing ping-pong
from sequential mechanisms.

{\def\LTcaptype{none} % do not increment counter
\begin{longtable}[]{@{}
  >{\raggedright\arraybackslash}p{(\linewidth - 4\tabcolsep) * \real{0.1522}}
  >{\raggedright\arraybackslash}p{(\linewidth - 4\tabcolsep) * \real{0.6522}}
  >{\raggedright\arraybackslash}p{(\linewidth - 4\tabcolsep) * \real{0.1957}}@{}}
\toprule\noalign{}
\begin{minipage}[b]{\linewidth}\raggedright
Feature
\end{minipage} & \begin{minipage}[b]{\linewidth}\raggedright
Sequential (Random or Ordered)
\end{minipage} & \begin{minipage}[b]{\linewidth}\raggedright
Ping-Pong
\end{minipage} \\
\midrule\noalign{}
\endhead
\bottomrule\noalign{}
\endlastfoot
Ternary complex EAB? & Yes & No \\
Covalent enzyme intermediate E*? & No & Yes \\
Lineweaver-Burk pattern (vary A, fixed B) & Intersecting lines &
Parallel lines \\
Examples & LDH, hexokinase, CK & AspAT, chymotrypsin, GS \\
\end{longtable}
}

\begin{quote}
\textbf{Concept Check 5:} A biochemist measures the rate of an enzyme
reaction at varying \hyperref[a]{A} for several fixed \hyperref[b]{B}
values, and obtains a family of parallel lines on a Lineweaver-Burk
plot. Identify the mechanism and predict whether a covalent enzyme
intermediate exists.
\end{quote}

\begin{center}\rule{0.5\linewidth}{0.5pt}\end{center}

\subsection{Allosteric Regulation and
Cooperativity}\label{allosteric-regulation-and-cooperativity}

Allosteric enzymes are not governed by simple
\hyperref[gl:michaelis-menten-kinetics]{\textbf{Michaelis-Menten
kinetics}}. They have: - Multiple subunits, typically with identical
active sites - Regulatory sites distinct from active sites - Cooperative
(sigmoidal) substrate binding - Modulation by effector molecules
(activators or inhibitors)

\subsubsection{The Hill Equation}\label{the-hill-equation}

For an enzyme with \(n\) equivalent cooperative subunits, the fraction
of active sites occupied (\texorpdfstring{\ensuremath{\theta}}{theta}) follows the \textbf{Hill equation}:

\[\theta = \frac{[\text{S}]^n}{K_{0.5}^n + [\text{S}]^n} \tag{4.14} \label{eq:unit_I_enzymes_and_kinetics_item_15}\]

where \textbf{\(n\)} is the Hill coefficient: - \(n = 1\): no
cooperativity (hyperbolic, Michaelis-Menten) - \(n > 1\): positive
cooperativity (sigmoidal; binding one substrate increases affinity for
subsequent substrates) - \(n < 1\): negative cooperativity (binding
decreases affinity)

\textbf{Linearised Hill equation:}

\[\log\frac{\theta}{1 - \theta} = n\log[\text{S}] - n\log K_{0.5} \tag{4.15} \label{eq:unit_I_enzymes_and_kinetics_item_16}\]

A plot of \(\log[\theta/(1-\theta)]\) vs.~\(\log[\text{S}]\) (the
\textbf{Hill plot}) gives a straight line with slope \(n\) and
\(x\)-intercept at \(\log K_{0.5}\).

Haemoglobin (see \cref{sec:unit_IX_circulation_respiration_homeostasis})
has \(n \approx 2.8\) (not 4, because cooperativity is imperfect).
Sigmoidal kinetics give allosteric enzymes \textbf{switch-like}
behaviour --- small changes in \hyperref[s]{S} near \(K_{0.5}\) cause
large changes in activity.

\textbf{Worked example:} For an allosteric enzyme with \(K_{0.5}\) = 5
mM and \(n\) = 3:

At \hyperref[s]{S} = 4 mM:
\(\theta = \frac{4^3}{5^3 + 4^3} = \frac{64}{125 + 64} = \frac{64}{189} = 0.339\)
(33.9\%)

At \hyperref[s]{S} = 6 mM:
\(\theta = \frac{6^3}{5^3 + 6^3} = \frac{216}{125 + 216} = \frac{216}{341} = 0.634\)
(63.4\%)

A 50\% increase in \hyperref[s]{S} (from 4 to 6 mM) nearly doubles the
fractional saturation. This is the ``switch-like'' sensitivity of
cooperative systems.

\subsubsection{Concerted (MWC) and Sequential
Models}\label{concerted-mwc-and-sequential-models}

Two models describe the molecular mechanism of cooperativity:

\textbf{Monod-Wyman-Changeux (MWC) Concerted Model:} - Most subunits
exist in either the T-state (tense, low affinity) or R-state (relaxed,
high affinity) - Most subunits transition simultaneously (concerted) -
Substrate preferentially binds R-state, shifting
T\(\rightleftharpoons\)R equilibrium toward R - Activators stabilise R;
inhibitors stabilise T

\textbf{Koshland-Nemethy-Filmer (KNF) Sequential Model:} - Each subunit
can independently switch between T and R - Substrate binding to one
subunit induces conformational changes that propagate to adjacent
subunits - Explains negative cooperativity (which MWC cannot)

In reality, most allosteric systems show features of both models.

\subsubsection{Aspartate Transcarbamoylase (ATCase) --- A Model
Allosteric
Enzyme}\label{aspartate-transcarbamoylase-atcase-a-model-allosteric-enzyme}

\begin{figure}[htbp]
\centering
\includegraphics[width=0.82\linewidth,height=4.2in,keepaspectratio]{../figures/mermaid_inline/inline_mermaid_0028_642863369a4a.png}
\caption{ATCase allostery (CTP and ATP). Allosteric regulation of ATCase. CTP (the end product of the pyrimidine pathway) is a feedback inhibitor that stabilises the T-state. ATP (a purine nucleotide) is an activator that stabilises the R-state, coordinating purine and pyrimidine synthesis.}
\end{figure}

ATCase (catalyses the first committed step of pyrimidine biosynthesis)
is the textbook allosteric enzyme: - \textbf{Structure:} 6 catalytic
subunits (c\(_3\))\(_2\) + 6 regulatory subunits (r\(_2\))\(_3\) =
c\(_6\)r\(_6\) (about 310 kDa) - \textbf{Inhibited} by CTP
(end-product of the pathway --- \textbf{feedback inhibition}) -
\textbf{Activated} by ATP (signals adequate purine availability
-\textgreater{} coordinate pyrimidine synthesis) - Hill coefficient
\(n \approx 2.5\) for N-carbamoylaspartate synthesis - T-\textgreater R
transition involves a 12-angstrom expansion and 10-degree rotation of
catalytic trimers

This allows the cell to precisely match nucleotide biosynthesis to
demand. When CTP is abundant, ATCase is inhibited; when ATP is abundant
(indicating nucleotide precursors are available), ATCase is activated.

\subsubsection{Why Cooperativity Improves Pathway
Control}\label{why-cooperativity-improves-pathway-control}

A first-glance question: why has evolution invested in elaborate
allosteric machinery when simple Michaelis-Menten enzymes could in
principle perform every metabolic task? The answer lies in the
\emph{control properties} of cooperative enzymes.

Compare a hyperbolic enzyme (Hill coefficient \(n = 1\)) and a sigmoidal
enzyme (\(n = 4\)) in the small region around \(K_{0.5}\). For a 2-fold
increase in \hyperref[s]{S} near \(K_{0.5}\):

{\def\LTcaptype{none} % do not increment counter
\begin{longtable}[]{@{}
  >{\raggedright\arraybackslash}p{(\linewidth - 6\tabcolsep) * \real{0.0667}}
  >{\raggedright\arraybackslash}p{(\linewidth - 6\tabcolsep) * \real{0.3667}}
  >{\raggedright\arraybackslash}p{(\linewidth - 6\tabcolsep) * \real{0.3444}}
  >{\raggedright\arraybackslash}p{(\linewidth - 6\tabcolsep) * \real{0.2222}}@{}}
\toprule\noalign{}
\begin{minipage}[b]{\linewidth}\raggedright
Enzyme
\end{minipage} & \begin{minipage}[b]{\linewidth}\raggedright
\texorpdfstring{\ensuremath{\theta}}{theta} at \([S] = 0.5\,K_{0.5}\)
\end{minipage} & \begin{minipage}[b]{\linewidth}\raggedright
\texorpdfstring{\ensuremath{\theta}}{theta} at \([S] = 2\,K_{0.5}\)
\end{minipage} & \begin{minipage}[b]{\linewidth}\raggedright
Sensitivity (ratio)
\end{minipage} \\
\midrule\noalign{}
\endhead
\bottomrule\noalign{}
\endlastfoot
Michaelis-Menten (\(n=1\)) & 0.33 & 0.67 & 2.0\(\times\) \\
Cooperative (\(n=2\)) & 0.20 & 0.80 & 4.0\(\times\) \\
Cooperative (\(n=4\)) & 0.059 & 0.941 & 16\(\times\) \\
Cooperative (\(n=8\)) & 0.0039 & 0.9961 & 256\(\times\) \\
\end{longtable}
}

An 8-fold increase in Hill coefficient (from n=1 to n=8) gives a
128-fold increase in switching sensitivity around \(K_{0.5}\) --- each
doubling of n roughly squares the sensitivity ratio. This
\textbf{ultrasensitivity} confers three regulatory advantages:

\begin{enumerate}
\def\labelenumi{\arabic{enumi}.}
\tightlist
\item
  \textbf{Steeper response to small signal changes.} Glycolysis can be
  ramped up rapidly during exercise without requiring 10-fold changes in
  metabolite concentrations.
\item
  \textbf{Sharper response thresholds.} Below \(K_{0.5}\) the pathway is
  nearly off; above \(K_{0.5}\) nearly fully on. Enables clear
  ``decision-making'' at metabolic branch points.
\item
  \textbf{Insulation from noise.} Random fluctuations in \hyperref[s]{S}
  within the low-saturation regime cause much smaller activity changes
  than under MM kinetics. Stability and switch-like behaviour coexist.
\end{enumerate}

The same logic explains why haemoglobin's \(n \approx 2.8\) delivers
about 25\% of bound O\(_2\) between arterial (pO\(_2\) 100
mmHg) and venous (40 mmHg) blood, whereas hyperbolic myoglobin would
deliver about 10\% of the same change. Cooperativity is biology's
tunable amplifier.

\begin{quote}
\textbf{Concept Check 6:} Phosphofructokinase-1 (PFK-1) is an allosteric
enzyme in \hyperref[gl:glycolysis]{\textbf{glycolysis}} that is
activated by AMP and inhibited by ATP and citrate. Why does this
regulatory pattern make metabolic sense? (Consider the energy charge of
the cell.)
\end{quote}

\begin{center}\rule{0.5\linewidth}{0.5pt}\end{center}

\subsection{Environmental Factors Affecting Enzyme
Activity}\label{environmental-factors-affecting-enzyme-activity}

\subsubsection{Temperature}\label{temperature}

Enzyme activity generally increases with temperature (\(\uparrow\)
kinetic energy -\textgreater{} \(\uparrow\) collision frequency
-\textgreater{} \(\uparrow\) rate) according to the Arrhenius equation:

\[k = A \cdot e^{-E_a/RT} \tag{4.16} \label{eq:unit_I_enzymes_and_kinetics_item_17}\]

A 10 degrees C rise approximately doubles reaction rate
(\(Q_{10} \approx 2\)). However, \textbf{above the optimal temperature},
denaturation of the protein occurs (tertiary and quaternary structure
disrupts), sharply reducing activity. The optimal temperature for most
human enzymes: about 37 degrees C.

\textbf{\(Q_{10}\) calculation:}

\[Q_{10} = \left(\frac{k_2}{k_1}\right)^{10/(T_2 - T_1)} \tag{4.17} \label{eq:unit_I_enzymes_and_kinetics_item_18}\]

For most enzyme reactions, \(Q_{10}\) = 1.5--3.0. The \(Q_{10}\) rule
breaks down near the denaturation temperature.

\textbf{Extremophile enzymes (extremozymes):} - \emph{Thermus aquaticus}
(hot springs, 75 degrees C) Taq polymerase: thermostable DNA polymerase,
the workhorse of PCR (optimal 72 degrees C). Taq achieves
thermostability through increased salt bridges, compact hydrophobic
core, and proline substitutions. - \emph{Pyrococcus furiosus} (deep-sea
vent, 100 degrees C) Pfu polymerase: even more thermostable and has
3'-\textgreater5' proofreading activity (higher fidelity than Taq). -
\emph{Psychrobacter} (Antarctic, --20 degrees C) enzymes: more flexible
(more Gly, fewer Pro, fewer salt bridges), lower activation energy,
active at sub-zero temperatures but rapidly denature at 37 degrees C.

\subsubsection{pH}\label{ph}

Enzyme catalysis depends on the ionisation state of key active-site
residues. Activity-pH curves are typically bell-shaped:

{\def\LTcaptype{none} % do not increment counter
\begin{longtable}[]{@{}
  >{\raggedright\arraybackslash}p{(\linewidth - 6\tabcolsep) * \real{0.1364}}
  >{\raggedright\arraybackslash}p{(\linewidth - 6\tabcolsep) * \real{0.1818}}
  >{\raggedright\arraybackslash}p{(\linewidth - 6\tabcolsep) * \real{0.2273}}
  >{\raggedright\arraybackslash}p{(\linewidth - 6\tabcolsep) * \real{0.4545}}@{}}
\toprule\noalign{}
\begin{minipage}[b]{\linewidth}\raggedright
Enzyme
\end{minipage} & \begin{minipage}[b]{\linewidth}\raggedright
Location
\end{minipage} & \begin{minipage}[b]{\linewidth}\raggedright
Optimum pH
\end{minipage} & \begin{minipage}[b]{\linewidth}\raggedright
Key Catalytic Residue
\end{minipage} \\
\midrule\noalign{}
\endhead
\bottomrule\noalign{}
\endlastfoot
Pepsin & Stomach & 1.5--2.5 & Asp (must be protonated) \\
Salivary amylase & Mouth & 6.7--7.0 & Glu (acid/base catalyst) \\
Trypsin & Intestine & 7.5--8.5 & His57 (must be deprotonated) \\
Alkaline phosphatase & Intestine/bone & 8--10 & Ser (phosphorylated
intermediate) \\
Arginase & Liver & 9.5--10.0 & Mn\(^{2+}\) (binuclear metal centre) \\
\end{longtable}
}

The bell-shaped curve reflects the requirement for specific protonation
states: the ascending limb corresponds to deprotonation of one essential
residue; the descending limb to deprotonation of another.

\subsubsection{Cofactors and Coenzymes}\label{cofactors-and-coenzymes}

Many enzymes require non-protein cofactors for activity. The enzyme
without its cofactor is the \textbf{apoenzyme}; the complete, active
enzyme is the \textbf{holoenzyme}.

{\def\LTcaptype{none} % do not increment counter
\begin{longtable}[]{@{}
  >{\raggedright\arraybackslash}p{(\linewidth - 6\tabcolsep) * \real{0.1212}}
  >{\raggedright\arraybackslash}p{(\linewidth - 6\tabcolsep) * \real{0.2424}}
  >{\raggedright\arraybackslash}p{(\linewidth - 6\tabcolsep) * \real{0.2424}}
  >{\raggedright\arraybackslash}p{(\linewidth - 6\tabcolsep) * \real{0.3939}}@{}}
\toprule\noalign{}
\begin{minipage}[b]{\linewidth}\raggedright
Type
\end{minipage} & \begin{minipage}[b]{\linewidth}\raggedright
Examples
\end{minipage} & \begin{minipage}[b]{\linewidth}\raggedright
Function
\end{minipage} & \begin{minipage}[b]{\linewidth}\raggedright
Dietary Source
\end{minipage} \\
\midrule\noalign{}
\endhead
\bottomrule\noalign{}
\endlastfoot
\textbf{Metal ions} & Zn\(^{2+}\) (carboxypeptidase) & Lewis acid
catalysis & Red meat, shellfish \\
& Mg\(^{2+}\) (kinases) & Stabilise ATP phosphate & Green vegetables \\
& Fe\(^{2+/3+}\) (cytochromes) & Electron transfer & Red meat,
lentils \\
& Cu\(^{2+}\) (cytochrome c oxidase) & O\(_2\) reduction & Liver,
nuts \\
& Mn\(^{2+}\) (SOD, arginase) & Redox, hydrolysis & Whole grains \\
& Mo (xanthine oxidase) & Oxygen atom transfer & Legumes \\
& Se (glutathione peroxidase) & Peroxide reduction & Brazil nuts \\
\textbf{Coenzymes} & NAD\(^+\)/NADH (from niacin, B3) & Hydride transfer
& Meat, fish \\
& FAD/FADH\(_2\) (from riboflavin, B2) & Hydride transfer & Dairy,
eggs \\
& Coenzyme A (from pantothenate, B5) & Acyl group transfer &
Widespread \\
& Pyridoxal phosphate (from pyridoxine, B6) & Amino group transfer &
Poultry, fish \\
& Thiamine pyrophosphate (from thiamine, B1) & Decarboxylation & Whole
grains \\
& Biotin (B7) & Carboxylation & Eggs, nuts \\
& Tetrahydrofolate (from folate, B9) & One-carbon transfer & Leafy
greens \\
& 5'-deoxyadenosylcobalamin (from B12) & Rearrangements & Animal
products \\
\textbf{Prosthetic groups} & Haem (contains Fe) & O\(_2\) binding,
electron transfer & Synthesised endogenously \\
& Lipoic acid & Acyl group transfer & Synthesised endogenously \\
\end{longtable}
}

Most water-soluble vitamins (B vitamins) function as coenzyme
precursors. This is why vitamin deficiencies cause enzyme dysfunction
and disease:

{\def\LTcaptype{none} % do not increment counter
\begin{longtable}[]{@{}
  >{\raggedright\arraybackslash}p{(\linewidth - 4\tabcolsep) * \real{0.4419}}
  >{\raggedright\arraybackslash}p{(\linewidth - 4\tabcolsep) * \real{0.3953}}
  >{\raggedright\arraybackslash}p{(\linewidth - 4\tabcolsep) * \real{0.1628}}@{}}
\toprule\noalign{}
\begin{minipage}[b]{\linewidth}\raggedright
Vitamin Deficiency
\end{minipage} & \begin{minipage}[b]{\linewidth}\raggedright
Coenzyme Affected
\end{minipage} & \begin{minipage}[b]{\linewidth}\raggedright
Disease
\end{minipage} \\
\midrule\noalign{}
\endhead
\bottomrule\noalign{}
\endlastfoot
B1 (Thiamine) & TPP & Beriberi, Wernicke-Korsakoff syndrome \\
B2 (Riboflavin) & FAD & Ariboflavinosis (angular stomatitis) \\
B3 (Niacin) & NAD\(^+\) & Pellagra (dermatitis, diarrhoea, dementia) \\
B6 (Pyridoxine) & PLP & Peripheral neuropathy \\
B9 (Folate) & THF & Megaloblastic anaemia, neural tube defects \\
B12 (Cobalamin) & AdoCbl, MeCbl & Pernicious anaemia,
neurodegeneration \\
\end{longtable}
}

\begin{quote}
\textbf{Clinical Connection: ACE Inhibitors and Zinc-Dependent
Catalysis}

Angiotensin-converting enzyme (ACE) is a zinc metalloprotease that
converts angiotensin I to angiotensin II (a potent vasoconstrictor) and
degrades bradykinin (a vasodilator). \textbf{ACE inhibitors} (enalapril,
lisinopril, ramipril) chelate the active-site Zn\(^{2+}\), preventing
catalysis. By blocking angiotensin II production and bradykinin
degradation, ACE inhibitors reduce blood pressure and are first-line
therapy for hypertension and heart failure. The side effect of dry cough
(in about 10\% of patients) is caused by bradykinin
accumulation in the airways.
\end{quote}

\begin{quote}
\textbf{Concept Check 7:} Most water-soluble vitamins serve as coenzyme
precursors. Explain why these vitamins must be obtained from the diet
(unlike fat-soluble vitamins, which can be stored). What does this imply
about the recommended frequency of intake?
\end{quote}

\begin{center}\rule{0.5\linewidth}{0.5pt}\end{center}

\subsection{Industrial and Biotechnological
Enzymes}\label{industrial-and-biotechnological-enzymes}

Enzymes have transformed medicine and industry:

Modern enzyme biotechnology couples classical kinetics with
structure-guided design, directed evolution, high-throughput screening,
and increasingly AI-predicted interaction models. AlphaFold-class tools
can nominate an active-site geometry or protein-complex interface, but
an engineered-enzyme claim still needs measured \(k_{cat}\), \(K_M\),
specificity, stability, expression yield, and performance under
application conditions; a plausible structure is a hypothesis until
kinetics and controls support it
\citep{abramson2024alphafold3,velankar2026alphafolddb2025}.

\subsubsection{CRISPR-Cas9 as a Nuclease
Enzyme}\label{crispr-cas9-as-a-nuclease-enzyme}

The Cas9 protein from \emph{Streptococcus pyogenes} is an RNA-guided
endonuclease: - \textbf{Classification:} EC 3.1 (nuclease / hydrolase) -
\textbf{Mechanism:} The guide RNA (gRNA) forms a 20-nt complementary
complex with the target DNA. Cas9 unwinds the DNA and cleaves both
strands using two nuclease domains: RuvC cuts the non-target strand; HNH
cuts the target strand - \textbf{Specificity:} PAM (protospacer adjacent
motif, 5'-NGG-3') is required adjacent to the target sequence -
\textbf{Applications:} \hyperref[gl:gene]{\textbf{Gene}} editing, gene
therapy (sickle cell disease, \texorpdfstring{\ensuremath{\beta}}{beta}-thalassaemia), diagnostics (SHERLOCK,
DETECTR) - \textbf{Nobel Prize 2020:} Jennifer Doudna and Emmanuelle
Charpentier

\subsubsection{Industrial Applications}\label{industrial-applications}

{\def\LTcaptype{none} % do not increment counter
\begin{longtable}[]{@{}lll@{}}
\toprule\noalign{}
Enzyme & Industry & Application \\
\midrule\noalign{}
\endhead
\bottomrule\noalign{}
\endlastfoot
Lipases & Biodiesel, detergents & Transesterification, stain removal \\
Cellulases & Biofuel & Cellulosic ethanol production \\
Amylases & Food, brewing & Starch liquefaction, brewing \\
Proteases & Detergents, leather & Protein stain removal, dehairing \\
Lactase & Dairy & Lactose-free milk production \\
DNA polymerase (Taq) & Molecular biology & PCR amplification \\
Restriction enzymes & Biotechnology & Molecular cloning \\
Glucose isomerase & Food & High-fructose corn syrup \\
\end{longtable}
}

\subsubsection{Enzyme Engineering
Approaches}\label{enzyme-engineering-approaches}

\begin{itemize}
\tightlist
\item
  \textbf{Directed evolution:} Random mutagenesis + screening/selection
  for improved properties. Frances Arnold received the \textbf{2018
  Nobel Prize in Chemistry} for this approach.
\item
  \textbf{Rational design:} Structure-guided
  \hyperref[gl:mutation]{\textbf{mutation}} of specific residues based
  on mechanistic understanding.
\item
  \textbf{Immobilised enzymes:} Enzymes attached to solid supports for
  continuous-flow industrial processes. Advantages: reusability,
  stability, easy product separation.
\end{itemize}

\subsubsection{Quantitative Metrics for Engineered
Enzymes}\label{quantitative-metrics-for-engineered-enzymes}

When evaluating an engineered or evolved variant, several quantitative
figures of merit are routinely compared against the wild-type enzyme:

{\def\LTcaptype{none} % do not increment counter
\begin{longtable}[]{@{}
  >{\raggedright\arraybackslash}p{(\linewidth - 4\tabcolsep) * \real{0.1818}}
  >{\raggedright\arraybackslash}p{(\linewidth - 4\tabcolsep) * \real{0.3030}}
  >{\raggedright\arraybackslash}p{(\linewidth - 4\tabcolsep) * \real{0.5152}}@{}}
\toprule\noalign{}
\begin{minipage}[b]{\linewidth}\raggedright
Metric
\end{minipage} & \begin{minipage}[b]{\linewidth}\raggedright
Definition
\end{minipage} & \begin{minipage}[b]{\linewidth}\raggedright
What it tells you
\end{minipage} \\
\midrule\noalign{}
\endhead
\bottomrule\noalign{}
\endlastfoot
Fold-improvement in \(k_{cat}\) &
\(k_{cat}^{\text{var}} / k_{cat}^{\text{WT}}\) & Speed-up at saturating
\hyperref[s]{S} \\
Fold-improvement in \(k_{cat}/K_m\) &
\((k_{cat}/K_m)^{\text{var}} / (k_{cat}/K_m)^{\text{WT}}\) & Catalytic
efficiency at low \hyperref[s]{S} (closer to \emph{in vivo}) \\
Substrate scope & \# of new substrates accepted & Breadth of activity \\
Enantioselectivity (E-value) &
\(\ln[(1 - c)(1 - ee_S)] / \ln[(1 - c)(1 + ee_S)]\) & Discrimination
between enantiomers \\
Thermal stability (\(T_{50}\), \(T_m\)) & Temp at half-activity /
unfolding midpoint & Process robustness \\
Half-life under operating conditions & Time to lose 50\% activity &
Industrial economics \\
Specific activity & \texorpdfstring{\ensuremath{\mu}}{mu}mol/min/mg & Practical productivity \\
Total turnover number (TTN) & Total substrate molecules per enzyme
before death & Catalyst longevity \\
\end{longtable}
}

A successful directed evolution campaign will track several of these
metrics simultaneously across many rounds; the best-known successes
(e.g., engineered transaminases for sitagliptin manufacture) achieved
25,000-fold rate improvements and \textgreater10\(^6\) TTN over
about 10 evolutionary rounds.

\begin{center}\rule{0.5\linewidth}{0.5pt}\end{center}

\subsection{Worked Examples}\label{worked-examples-2}

\textbf{Problem 1:} An enzyme obeys Michaelis-Menten kinetics with
\(K_m\) = 2.0 mM and \(V_{max}\) = 100 \texorpdfstring{\ensuremath{\mu}}{mu}mol/min. Calculate \(v_0\) at
\hyperref[s]{S} = 0.5 mM, 2.0 mM, 10 mM, and 200 mM.

\emph{Solution:}

\[v_0 = \frac{100 \times [\text{S}]}{2.0 + [\text{S}]} \tag{4.18} \label{eq:unit_I_enzymes_and_kinetics_item_19}\]

{\def\LTcaptype{none} % do not increment counter
\begin{longtable}[]{@{}
  >{\raggedright\arraybackslash}p{(\linewidth - 4\tabcolsep) * \real{0.2093}}
  >{\raggedright\arraybackslash}p{(\linewidth - 4\tabcolsep) * \real{0.4651}}
  >{\raggedright\arraybackslash}p{(\linewidth - 4\tabcolsep) * \real{0.3256}}@{}}
\toprule\noalign{}
\begin{minipage}[b]{\linewidth}\raggedright
\hyperref[s]{S} (mM)
\end{minipage} & \begin{minipage}[b]{\linewidth}\raggedright
\(v_0\) (\texorpdfstring{\ensuremath{\mu}}{mu}mol/min)
\end{minipage} & \begin{minipage}[b]{\linewidth}\raggedright
\% of \(V_{max}\)
\end{minipage} \\
\midrule\noalign{}
\endhead
\bottomrule\noalign{}
\endlastfoot
0.5 & \(100 \times 0.5 / 2.5 = 20.0\) & 20\% \\
2.0 & \(100 \times 2.0 / 4.0 = 50.0\) & 50\% (by definition,
\hyperref[s]{S} = \(K_m\)) \\
10.0 & \(100 \times 10 / 12 = 83.3\) & 83.3\% \\
200.0 & \(100 \times 200 / 202 = 99.0\) & 99.0\% (near saturation) \\
\end{longtable}
}

\textbf{Problem 2:} In the presence of a competitive inhibitor at
\hyperref[i]{I} = 5 mM and \(K_i\) = 2.5 mM, what is the new apparent
\(K_m\)? At what \hyperref[s]{S} will \(v_0\) = \(V_{max}/2\)?

\emph{Solution:}
\[\alpha = 1 + \frac{[\text{I}]}{K_i} = 1 + \frac{5.0}{2.5} = 3.0 \tag{4.19} \label{eq:unit_I_enzymes_and_kinetics_item_20}\]

\[K_m^{app} = \alpha K_m = 3.0 \times 2.0 = 6.0 \; \text{mM} \tag{4.20} \label{eq:unit_I_enzymes_and_kinetics_item_21}\]

\(v_0 = V_{max}/2\) when \hyperref[s]{S} = \(K_m^{app}\) = 6.0 mM. The
inhibitor triples the substrate concentration needed for half-maximal
velocity.

\textbf{Problem 3:} The Hill coefficient for haemoglobin is 2.8, and
\(K_{0.5}\) for O\(_2\) is 26 mmHg. Calculate the fractional saturation
at: (a) pO\(_2\) = 40 mmHg (venous blood) (b) pO\(_2\) = 100 mmHg
(arterial blood)

\emph{Solution:}
\[\theta = \frac{[\text{S}]^n}{K_{0.5}^n + [\text{S}]^n} \tag{4.21} \label{eq:unit_I_enzymes_and_kinetics_item_22}\]

\begin{enumerate}
\def\labelenumi{(\alph{enumi})}
\tightlist
\item
  \(\theta = \frac{40^{2.8}}{26^{2.8} + 40^{2.8}} = \frac{40^{2.8}}{26^{2.8} + 40^{2.8}}\)
\end{enumerate}

\(40^{2.8} = e^{2.8 \ln 40} = e^{2.8 \times 3.689} = e^{10.33} = 30{,}600\)

\(26^{2.8} = e^{2.8 \ln 26} = e^{2.8 \times 3.258} = e^{9.12} = 9{,}160\)

\(\theta = \frac{30{,}600}{9{,}160 + 30{,}600} = \frac{30{,}600}{39{,}760} = 0.770\)
(77.0\%)

\begin{enumerate}
\def\labelenumi{(\alph{enumi})}
\setcounter{enumi}{1}
\tightlist
\item
  \(100^{2.8} = e^{2.8 \times 4.605} = e^{12.89} = 396{,}000\)
\end{enumerate}

\(\theta = \frac{396{,}000}{9{,}160 + 396{,}000} = 0.977\) (97.7\%)

The difference: haemoglobin loads O\(_2\) to 97.7\% in the lungs and
unloads to 77.0\% in tissues, delivering about 20\% of its
O\(_2\) per circuit. During exercise (venous pO\(_2\) drops to
about 20 mmHg), delivery increases substantially due to the
steep part of the sigmoidal curve.

\begin{center}\rule{0.5\linewidth}{0.5pt}\end{center}

\subsection{Computational Bridge}\label{computational-bridge-3}

The Michaelis--Menten hyperbola is implemented directly on
concentrations in \texorpdfstring{\ensuremath{\mu}}{mu}M:

\begin{Shaded}
\begin{Highlighting}[]
\ImportTok{from}\NormalTok{ biology.biochemistry }\ImportTok{import}\NormalTok{ michaelis\_menten}

\NormalTok{res }\OperatorTok{=}\NormalTok{ michaelis\_menten(}\FloatTok{5.0}\NormalTok{, Vmax}\OperatorTok{=}\FloatTok{120.0}\NormalTok{, Km}\OperatorTok{=}\FloatTok{10.0}\NormalTok{)}
\BuiltInTok{print}\NormalTok{(res.reaction\_rate, res.efficiency)  }\CommentTok{\# v and v/Vmax}
\end{Highlighting}
\end{Shaded}

\begin{quote}
\textbf{Clinical / systems note:} Many drugs are mechanism-based
inhibitors that shrink effective \(V_\mathrm{max}\) or raise apparent
\(K_m\); statins compete at HMG-CoA reductase's active site,
illustrating competitive themes at a clinically relevant scale.
\end{quote}

\begin{center}\rule{0.5\linewidth}{0.5pt}\end{center}

\subsection{Bridge to Systems: From Enzyme Kinetics to Metabolic
Flux}\label{bridge-to-systems-from-enzyme-kinetics-to-metabolic-flux}

A single enzyme's \(K_m\), \(V_{\max}\), and Hill coefficient become
inputs to a higher-level question: what is the \textbf{flux} (\(J\), in
mol L\(^{-1}\) s\(^{-1}\)) through a metabolic pathway, and which enzyme
limits it? This question is the bridge from this chapter to Unit III
(bioenergetics and respiration).

\subsubsection{Pathway Flux is Set by Multiple Enzymes, Not
One}\label{pathway-flux-is-set-by-multiple-enzymes-not-one}

The classical idea that a single ``rate-limiting step'' governs a
pathway has been replaced by \textbf{metabolic control analysis} (MCA),
developed in the 1970s by Heinrich, Rapoport, Kacser, and Burns. MCA
assigns each enzyme \(i\) a \textbf{flux control coefficient}:

\begin{equation}
C_i^J = \frac{\partial \ln J}{\partial \ln [E_i]}
\label{eq:unit_I_flux_control_coefficient}
\end{equation}

This is the fractional change in steady-state flux when the activity of
enzyme \(i\) is changed by a small fraction. The \textbf{summation
theorem} states that for any pathway:

\begin{equation}
\sum_i C_i^J = 1
\label{eq:unit_I_summation_theorem}
\end{equation}

Three lessons follow:

\begin{enumerate}
\def\labelenumi{\arabic{enumi}.}
\tightlist
\item
  \textbf{Control is distributed.} No single enzyme has \(C_i^J = 1\);
  instead, glycolytic flux in muscle is controlled by hexokinase
  (\(C^J \approx 0.3\)), PFK-1 (\(C^J \approx 0.3\)), pyruvate kinase
  (\(C^J \approx 0.2\)), and others.
\item
  \textbf{Targeting the ``rate-limiting'' enzyme often disappoints.}
  Doubling PFK-1 primarily doubles flux if \(C^J = 1\), which is rarely
  true. This is why metabolic engineering campaigns to improve product
  yields often increase one enzyme tenfold for a 1.5-fold flux gain.
\item
  \textbf{Allosteric regulation tunes the control distribution.} When
  ATP is high, AMP is low, and ATCase or PFK-1 sit in their T-state,
  \(C^J\) for those enzymes rises toward 1; the cell concentrates
  control where it can respond quickly to demand changes.
\end{enumerate}

\subsubsection{Enzyme Saturation Modes Determine Pathway
Behaviour}\label{enzyme-saturation-modes-determine-pathway-behaviour}

The local \([\text{S}] / K_m\) ratio for each enzyme determines whether
it operates near saturation or far below. Glycolytic enzymes typically
operate at \([\text{S}] \approx K_m\) (around half-saturation), where
the response to small \([\text{S}]\) changes is approximately linear and
control coefficients are large. Enzymes whose substrates are at
\(\gg K_m\) (saturated) have \(C^J \approx 0\) and exert no flux
control. Designing a pathway thus means choosing \(K_m\) values matched
to physiological metabolite concentrations --- a constraint visible in
the convergent evolution of glycolytic enzymes across most kingdoms of
life.

\subsubsection{Branching Points and
Connectivity}\label{branching-points-and-connectivity}

When a metabolite sits at a branch point (e.g., glucose-6-phosphate, fed
into glycolysis or the pentose phosphate pathway), the control
coefficient for one branch can exceed 1, balanced by a \emph{negative}
control coefficient on the other branch. Allosteric regulation at branch
points (NADPH-mediated inhibition of glucose-6-phosphate dehydrogenase;
citrate inhibition of PFK-1) allows the cell to apportion flux based on
demand.

Unit III revisits these concepts at the level of full pathways
(glycolysis, TCA cycle, electron transport) and shows how
phosphorylation cascades, \hyperref[gl:hormone]{\textbf{hormone}}-driven
changes in enzyme abundance, and substrate cycles allow flux to vary
100-fold between rest and exercise.

\begin{quote}
\textbf{Concept Check 8:} Insulin treatment of a hepatocyte reduces
gluconeogenic flux by about 10\(\times\) within minutes,
primarily by phosphorylating PFK-2/FBPase-2 to raise
{[}fructose-2,6-bisphosphate{]}. Without invoking transcription, explain
how this allosteric activator achieves a 10\(\times\) flux reduction
while individual enzyme abundances remain unchanged. Use the language of
\(C^J\) and \(K_m\).
\end{quote}

\begin{center}\rule{0.5\linewidth}{0.5pt}\end{center}

\subsection{Current Evidence and Frontier
Biology}\label{current-evidence-and-frontier-biology-6}

For \textbf{Enzymes and the Kinetics of Catalysis}, frontier biology
belongs inside the evidence logic of the chapter. Chemistry-of-life
claims now connect classical bonding and thermodynamics with AI-guided
structure prediction and experimental validation. The core reading
question is this: enzyme claims should separate binding, catalysis,
regulation, transport limits, and measurement conditions.

\begin{itemize}
\tightlist
\item
  \textbf{What to verify:} identify the observation, model, assay, or
  dataset that would make the claim stronger or weaker.
\item
  \textbf{What to qualify:} state the scale, organism, cell type,
  environmental condition, or population where the claim is expected to
  hold.
\item
  \textbf{What to compare:} test at least one alternative explanation,
  baseline, or null model before treating the pattern as causal.
\item
  \textbf{What to cite:} distinguish primary evidence, review synthesis,
  public dataset, and institutional guidance; for recent or numeric
  claims, prefer the source closest to the measurement and state what
  has changed since it was published.
\end{itemize}

Use AI biomolecular models as hypothesis generators: compare confidence,
conservation, solvent exposure, and assay evidence before turning a
predicted contact into a biological claim
\citep{abramson2024alphafold3}.

\textbf{Source practice:} For structure and interaction claims, cite
experimental structures when available and treat AlphaFold 3 or AFDB
complex predictions as hypotheses to validate with confidence metrics,
conservation, mutagenesis, binding, or cryo-EM/X-ray/NMR evidence
\citep{abramson2024alphafold3,velankar2026alphafolddb2025,emblebi2026alphafoldcomplexes}.

\subsubsection{Current Evidence Map: Enzyme Engineering Evidence
Chain}\label{current-evidence-map-enzyme-engineering-evidence-chain}

\begin{figure}[htbp]
\centering
\includegraphics[width=0.82\linewidth,height=4.2in,keepaspectratio]{../figures/mermaid_inline/inline_mermaid_0029_16fa34bca390.png}
\caption{Enzyme-engineering claims need rate, specificity, stability, and context; a better active-site story is not enough without quantitative kinetics.}
\end{figure}

\subsection{Summary}\label{summary-6}

\begin{itemize}
\tightlist
\item
  Enzymes are protein catalysts that lower the activation energy
  \(\Delta G^{\ddagger}\) without altering equilibrium.
\item
  Rate enhancement mechanisms include proximity/orientation, transition
  state stabilisation, acid-base catalysis, covalent catalysis, and
  metal ion catalysis.
\item
  The Michaelis-Menten equation:
  \(v_0 = V_{max}[\text{S}] / (K_m + [\text{S}])\); \(K_m\) approximates
  substrate affinity; \(k_{cat}\) = turnover number; \(k_{cat}/K_m\) =
  catalytic efficiency.
\item
  Enzyme inhibition: competitive (increases apparent \(K_m\)),
  uncompetitive (decreases both \(K_m\) and \(V_{max}\)),
  mixed/noncompetitive (\(V_{max}\) decreases), irreversible (covalent
  inactivation). Each has a characteristic Lineweaver-Burk pattern.
\item
  Allosteric enzymes exhibit sigmoidal kinetics (Hill equation); key for
  metabolic feedback control. ATCase is the paradigm for allosteric
  regulation.
\item
  pH, temperature, and cofactor availability most modulate enzyme
  activity. Vitamin deficiencies impair coenzyme-dependent enzymes.
\item
  Enzymes are classified into seven EC classes (the seventh,
  translocases, added in 2018).
\item
  Modern enzyme applications include CRISPR-Cas9 gene editing,
  industrial biocatalysis, and directed evolution.
\item
  \textbf{Connections:} See Unit III (pathway-level flux and regulated
  enzymes), Unit IV (polymerases and repair enzymes), and Unit VII
  (antibiotic targets in bacterial metabolism).
\end{itemize}

\subsection{Key Terms}\label{key-terms-6}

\begin{itemize}
\tightlist
\item
  \textbf{Activation energy (\(\Delta G^{\ddagger}\)):} Energy barrier
  between reactants and transition state
\item
  \textbf{Active site:} Region of enzyme that binds substrate and
  catalyses reaction
\item
  \textbf{Transition state:} Highest-energy intermediate along the
  reaction coordinate
\item
  \textbf{Michaelis constant (\(K_m\)):} Substrate concentration at
  half-maximal velocity
\item
  \textbf{Turnover number (\(k_{cat}\)):} Maximum substrate molecules
  converted per enzyme per second
\item
  \textbf{Catalytic efficiency (\(k_{cat}/K_m\)):} Overall measure of
  enzyme performance
\item
  \textbf{Competitive inhibition:} Inhibitor competes with substrate for
  active site; overcome by high \hyperref[s]{S}
\item
  \textbf{Uncompetitive inhibition:} Inhibitor binds primarily the ES
  complex
\item
  \textbf{Mixed inhibition:} Inhibitor binds both E and ES at an
  allosteric site
\item
  \textbf{Irreversible inhibition:} Covalent modification permanently
  inactivates the enzyme
\item
  \textbf{Allosteric regulation:} Modulation by effector binding at a
  site other than the active site
\item
  \textbf{Hill coefficient (\(n\)):} Measure of cooperativity; \(n > 1\)
  = positive cooperativity
\item
  \textbf{Feedback inhibition:} End product inhibits an early enzyme in
  its biosynthetic pathway
\item
  \textbf{Coenzyme:} Organic cofactor, often derived from a
  water-soluble vitamin
\item
  \textbf{Prosthetic group:} Tightly bound cofactor (e.g., haem, FAD)
\item
  \textbf{Holoenzyme:} Complete enzyme with most cofactors bound
\item
  \textbf{Apoenzyme:} Enzyme protein without its cofactor
\item
  \textbf{EC number:} Four-digit enzyme classification number assigned
  by IUBMB
\item
  \textbf{Transition state analogue:} Inhibitor that mimics the
  transition state geometry
\item
  \textbf{Directed evolution:} Laboratory technique for engineering
  improved enzymes
\end{itemize}

\subsection{Review Questions}\label{review-questions-6}

\begin{enumerate}
\def\labelenumi{\arabic{enumi}.}
\tightlist
\item
  Explain, using transition state theory, why an enzyme that binds the
  transition state \(10^6\)-fold more tightly than the substrate
  achieves a rate enhancement of approximately \(10^6\)-fold.
\item
  Derive the Michaelis-Menten equation starting from the steady-state
  assumption. State clearly most assumptions made.
\item
  An enzyme has \(K_m\) = 0.2 mM and \(V_{max}\) = 50 \texorpdfstring{\ensuremath{\mu}}{mu}mol/min. A
  competitive inhibitor (\(K_i\) = 0.1 mM) is added at 0.5 mM. Calculate
  the apparent \(K_m\) and the rate at \hyperref[s]{S} = 0.2 mM with and
  without inhibitor.
\item
  On a Lineweaver-Burk plot, two lines obtained with and without an
  inhibitor are parallel. Identify the type of inhibition and explain
  what molecular event is occurring.
\item
  The Hill coefficient for phosphofructokinase-1 is approximately 3.8.
  What does this tell you about the enzyme's response to changes in
  substrate concentration near \(K_{0.5}\)?
\item
  Explain how aspirin's mechanism of action (irreversible COX
  inhibition) accounts for both its anti-inflammatory properties and its
  use as an antiplatelet drug.
\item
  Why do statins lower blood cholesterol more effectively than simply
  reducing dietary cholesterol intake? Consider the relative
  contributions of endogenous synthesis and dietary absorption.
\item
  A food scientist wants to use an immobilised enzyme in a
  continuous-flow reactor. What are three advantages of enzyme
  immobilisation over using the free enzyme in solution?
\item
  CRISPR-Cas9 requires a PAM sequence (5'-NGG-3') adjacent to its
  target. Explain why this requirement exists and how it affects the
  enzyme's specificity. How does this relate to the concept of substrate
  recognition?
\item
  Pellagra (deficiency of niacin/vitamin B3) causes the ``three Ds'':
  dermatitis, diarrhoea, and dementia. Given that niacin is the
  precursor to NAD\(^+\), explain why this deficiency has such
  widespread systemic effects.
\item
  Using \texttt{michaelis\_menten}, compute \(v\) at \([S] = K_m/2\),
  \(K_m\), and \(2K_m\) for the same \(V_\mathrm{max}\). Verify
  numerically that \(v\) is \textbf{not} linear in \([S]\) even when
  \([S] \ll K_m\) is often approximated as linear.
\item
  Explain why \(k_\mathrm{cat}/K_m\) is a useful upper-bound metric for
  substrate discrimination under cellular subsaturating conditions, and
  give one evolutionary scenario where
  \hyperref[gl:natural-selection]{\textbf{natural selection}} would
  increase this ratio.
\item
  A two-substrate enzyme is studied at varying \hyperref[a]{A} for
  several fixed \hyperref[b]{B}. The Lineweaver-Burk plot reveals
  \textbf{intersecting} lines that meet to the \textbf{left of the
  y-axis}. Use \cref{eq:unit_I_bibi_sequential} and
  \cref{eq:unit_I_bibi_pingpong} to argue whether the mechanism is
  sequential or ping-pong.
\item
  Compare the switching sensitivity (fold-change in \texorpdfstring{\ensuremath{\theta}}{theta} for a 2\(\times\)
  change in \hyperref[s]{S} near \(K_{0.5}\)) of cooperative enzymes
  with \(n=1, 2, 4, 8\). Why do regulatory (rate-controlling) enzymes
  typically have \(n > 2\) while housekeeping enzymes are usually
  hyperbolic?
\item
  A directed-evolution campaign improves a transaminase's
  \(k_{cat}/K_m\) by 100\(\times\) over six rounds of
  mutation/selection. Discuss two metrics besides \(k_{cat}/K_m\) that
  the campaign should monitor in parallel, and explain why
  \(k_{cat}/K_m\) alone is not sufficient to declare success for an
  industrial process.
\item
  Using \cref{eq:unit_I_summation_theorem}, argue why doubling the
  abundance of a single ``rate-limiting'' glycolytic enzyme rarely
  doubles glycolytic flux. Suggest one cellular scenario where the flux
  control coefficient \(C^J\) for PFK-1 would approach 1.
\end{enumerate}

\begin{center}\rule{0.5\linewidth}{0.5pt}\end{center}

\subsection{Further Reading and Source
Notes}\label{further-reading-and-source-notes-6}

\begin{itemize}
\tightlist
\item
  Fischer (1894). Einfluss der Configuration auf die Wirkung der Enzyme.
  \emph{Berichte der deutschen chemischen Gesellschaft}, 27.
\item
  Koshland (1958). Application of a Theory of Enzyme Specificity to
  Protein Synthesis. \emph{Proceedings of the National Academy of
  Sciences}, 44.
\end{itemize}

\begin{center}\rule{0.5\linewidth}{0.5pt}\end{center}

\subsubsection{Companion Source Module}\label{companion-source-module-6}

\textbf{Enzymes and the Kinetics of Catalysis} should leave a
reproducible trail from a biological claim to the code, figure, diagram,
or paper-based activity that can test it. Use the surfaces below to
inspect the chapter's assumptions, rerun the relevant model, or compare
the manuscript explanation with companion labs and figures.

{\def\LTcaptype{none} % do not increment counter
\begin{longtable}[]{@{}
  >{\raggedright\arraybackslash}p{(\linewidth - 2\tabcolsep) * \real{0.5000}}
  >{\raggedright\arraybackslash}p{(\linewidth - 2\tabcolsep) * \real{0.5000}}@{}}
\toprule\noalign{}
\begin{minipage}[b]{\linewidth}\raggedright
Surface
\end{minipage} & \begin{minipage}[b]{\linewidth}\raggedright
Use it for
\end{minipage} \\
\midrule\noalign{}
\endhead
\bottomrule\noalign{}
\endlastfoot
\breaktt{src/biology/biochemistry/biochemistry.py}
(\texttt{michaelis\_menten}, \breaktt{competitive\_inhibition},
\texttt{enzyme\_rate\_curve}) & Reproduce saturation curves, inhibition
shifts, and parameter interpretation. \\
\breaktt{src/visualization/plots.py} (\breaktt{plot\_michaelis\_menten}) &
Check that graph shape, axes, and units match the kinetic equation. \\
\breaktt{src/mermaid/biology\_diagrams.py}
(\breaktt{enzyme\_kinetics\_diagram}) & Keep mechanism, substrate
binding, and regulation visually aligned. \\
\end{longtable}
}

\textbf{Reproducibility check:} report substrate range, enzyme amount,
temperature, pH, and rate units before comparing kinetic parameters.
\textbf{Cross-reference:} connect kinetics to
\cref{sec:unit_III_metabolic_integration} and
\cref{sec:unit_IV_gene_expression}.

\newpage

\section*{Unit II --- The Cell:
Introduction}\label{unit-ii-the-cell-introduction}
\addcontentsline{toc}{section}{Unit II --- The Cell: Introduction}

\subsection*{Why This Unit Matters}\label{why-this-unit-matters-2}
\addcontentsline{toc}{subsection}{Why This Unit Matters}

In 1665, Robert Hooke pressed a thin sliver of cork against the lens of
a compound microscope he had built himself and saw, for the first time
in human history, the architectural unit of life: the cell. He called
them \emph{cellulae} --- small rooms --- because they reminded him of
monks' chambers in a monastery. What Hooke saw in cork were dead walls.
Antoni van Leeuwenhoek, a decade later, saw living cells: bacteria
darting through pond water, human sperm, red blood cells coursing
through capillaries. The universe had grown smaller in every direction.

Today, three and a half centuries after those first observations, the
cell remains the central unit of biology --- and it turns out to be far
stranger and more sophisticated than Hooke could have imagined. A
typical human liver cell contains about 20,000 different
proteins, about 5 km of DNA, about 5,000 mitochondria,
and about 10 million ribosomes, most coordinated within a
volume of roughly 1,000 \texorpdfstring{\ensuremath{\mu}}{mu}m\texorpdfstring{\textsuperscript{3}}{3}. Every second, that cell performs thousands
of enzymatic reactions, processes about 20 receptor signals,
divides its mitochondria, degrades misfolded proteins in lysosomes at pH
4.5, and repairs about 10,000 DNA lesions in its nucleus ---
most simultaneously, without a central coordinator.

This unit examines the cell at three levels of resolution: the
\textbf{cellular} (organelle inventory and function), the
\textbf{molecular} (membrane biophysics, signal transduction), and the
\textbf{computational} (mathematical models of cell size constraints,
diffusion, and membrane potential). You will encounter the endosymbiotic
theory --- one of the most elegant and evidence-rich hypotheses in
biology --- and the minimal cell, which asks: what is the irreducible
instruction set for life?

\begin{center}\rule{0.5\linewidth}{0.5pt}\end{center}

\subsection*{Landmark Discoveries}\label{landmark-discoveries-2}
\addcontentsline{toc}{subsection}{Landmark Discoveries}

{\def\LTcaptype{none} % do not increment counter
\begin{longtable}[]{@{}
  >{\raggedright\arraybackslash}p{(\linewidth - 8\tabcolsep) * \real{0.2407}}
  >{\raggedright\arraybackslash}p{(\linewidth - 8\tabcolsep) * \real{0.0741}}
  >{\raggedright\arraybackslash}p{(\linewidth - 8\tabcolsep) * \real{0.2963}}
  >{\raggedright\arraybackslash}p{(\linewidth - 8\tabcolsep) * \real{0.1667}}
  >{\raggedright\arraybackslash}p{(\linewidth - 8\tabcolsep) * \real{0.2222}}@{}}
\toprule\noalign{}
\begin{minipage}[b]{\linewidth}\raggedright
Discoverer(s)
\end{minipage} & \begin{minipage}[b]{\linewidth}\raggedright
Year
\end{minipage} & \begin{minipage}[b]{\linewidth}\raggedright
Journal / Source
\end{minipage} & \begin{minipage}[b]{\linewidth}\raggedright
Discovery
\end{minipage} & \begin{minipage}[b]{\linewidth}\raggedright
Significance
\end{minipage} \\
\midrule\noalign{}
\endhead
\bottomrule\noalign{}
\endlastfoot
Robert Hooke & 1665 & \emph{Micrographia} & ``Cells'' in cork & Named
and conceptualised the cellular unit; launched microscopy as a
biological tool \\
Schleiden \& Schwann & 1838--39 & Multiple & Cell theory (plants +
animals) & First two postulates: most organisms are composed of cells \\
Rudolf Virchow & 1855 & \emph{Archiv f.~path. Anat.} & \emph{Omnis
cellula e cellula} & Third postulate: cells arise primarily from
pre-existing cells; foundational to cancer biology \\
George Palade & 1953 & \emph{J. Histochem. Cytochem.} & Ribosome
discovery via electron microscopy & Identified the protein synthesis
machinery; Nobel Prize, 1974 \\
Peter Mitchell & 1961 & \emph{Nature} & Chemiosmosis ---
proton-gradient-driven ATP synthesis & Explained mitochondrial ATP
production; later became central to Unit III \\
Lynn Margulis & 1967 & \emph{J. Theor. Biol.} & Formal endosymbiotic
theory & Proposed mitochondria and chloroplasts are descended from
engulfed bacteria \\
J. Craig Venter et al. & 2016 & \emph{Science} & JCVI-syn3.0: minimal
synthetic cell (473 genes) & Defined the lower boundary of
self-replicating cellular life \\
\end{longtable}
}

\begin{center}\rule{0.5\linewidth}{0.5pt}\end{center}

\subsection*{Key Concepts and
Connections}\label{key-concepts-and-connections-2}
\addcontentsline{toc}{subsection}{Key Concepts and Connections}

\begin{figure}[htbp]
\centering
\includegraphics[width=0.82\linewidth,height=4.2in,keepaspectratio]{../figures/mermaid_inline/inline_mermaid_0030_cf4b6cf2fa06.png}
\caption{Unit II concept map - The Cell. Purple = foundational theory; red = energy-related organelles (anticipating Unit III); blue = membrane; green = signalling.}
\end{figure}

\begin{center}\rule{0.5\linewidth}{0.5pt}\end{center}

\subsection*{Current Evidence Thread}\label{current-evidence-thread-2}
\addcontentsline{toc}{subsection}{Current Evidence Thread}

Read Unit II as a measurement story: cell biology is now captured as
live-cell biosensor and optogenetic recordings, cryo-electron tomograms,
spatial and single-cell atlases, and perturbational screens rather than
static diagrams alone. As you move through the chapters, keep a
two-column note: \textbf{claim} on the left, \textbf{evidence that would
change my confidence} on the right. By the end of the unit, each major
idea should be tied to a measurement, model, citation, or paper-based
lab decision.

\subsection*{Chapter Roadmap}\label{chapter-roadmap-2}
\addcontentsline{toc}{subsection}{Chapter Roadmap}

{\def\LTcaptype{none} % do not increment counter
\begin{longtable}[]{@{}
  >{\raggedright\arraybackslash}p{(\linewidth - 6\tabcolsep) * \real{0.1556}}
  >{\raggedright\arraybackslash}p{(\linewidth - 6\tabcolsep) * \real{0.1111}}
  >{\raggedright\arraybackslash}p{(\linewidth - 6\tabcolsep) * \real{0.2889}}
  >{\raggedright\arraybackslash}p{(\linewidth - 6\tabcolsep) * \real{0.4444}}@{}}
\toprule\noalign{}
\begin{minipage}[b]{\linewidth}\raggedright
Chapter
\end{minipage} & \begin{minipage}[b]{\linewidth}\raggedright
Title
\end{minipage} & \begin{minipage}[b]{\linewidth}\raggedright
Core Question
\end{minipage} & \begin{minipage}[b]{\linewidth}\raggedright
Key Equation / Model
\end{minipage} \\
\midrule\noalign{}
\endhead
\bottomrule\noalign{}
\endlastfoot
\textbf{5} & Cell Theory and Cell Types & What is a cell, and how did
cell theory transform biology? & \(A/V = 3/r\); \(t = x^2/2D\) \\
\textbf{6} & Cell Structure and Organelles & How do eukaryotic
organelles achieve compartmentalised function? & Nernst equation
(membrane potential) \\
\textbf{7} & Membrane Transport & How do cells control passage across
membranes? & Fick's First Law: \(J = -D(dC/dx)\) \\
\textbf{8} & Cell Signalling & How do extracellular signals change
intracellular behaviour? & Hill equation (co-operativity); cAMP second
messenger models \\
\end{longtable}
}

\begin{center}\rule{0.5\linewidth}{0.5pt}\end{center}

\subsection*{Connections Across the
Textbook}\label{connections-across-the-textbook-2}
\addcontentsline{toc}{subsection}{Connections Across the Textbook}

\begin{itemize}
\tightlist
\item
  \textbf{Membranes and transport}
  (\cref{sec:unit_II_membrane_transport}) are the physical substrate for
  the proton-motive force of mitochondria and chloroplasts (Unit III)
  and the action potential of neurons (Unit IX).
\item
  \textbf{Cell signalling} (\cref{sec:unit_II_cell_signaling}) reappears
  in Unit IV (nuclear signalling and gene regulation), Unit V
  (mitosis/meiosis checkpoints), Unit IX (hormone receptor cascades),
  and oncogene/tumour-suppressor biology.
\item
  \textbf{Endosymbiosis} directly informs Unit VI (mitochondrial DNA as
  molecular clock), Unit VII (organelle gene transfer to nucleus), and
  Unit VIII (chloroplast inheritance).
\item
  \textbf{Prokaryotic cell structure} (cell walls, flagella, pili)
  underpins Unit VII (Microbiology) --- antibiotic mechanisms, Gram
  staining, biofilm formation.
\end{itemize}

\begin{quote}
\textbf{Key vocabulary introduced here:} prokaryote, eukaryote,
organelle, nucleoid, plasma membrane, endosymbiosis, fluid-mosaic model,
ligand, receptor, second messenger, autophagy.
\end{quote}

\subsection*{Computational Toolbox --- Unit
II}\label{computational-toolbox-unit-ii}
\addcontentsline{toc}{subsection}{Computational Toolbox --- Unit II}

\begin{Shaded}
\begin{Highlighting}[]
\ImportTok{from}\NormalTok{ biology.cell }\ImportTok{import}\NormalTok{ nernst\_potential, goldman\_equation, IonConcentration}

\CommentTok{\# Nernst equation: equilibrium potential for K+ at 37°C}
\NormalTok{k\_ion }\OperatorTok{=}\NormalTok{ IonConcentration(name}\OperatorTok{=}\StringTok{"K+"}\NormalTok{, intracellular\_mM}\OperatorTok{=}\DecValTok{150}\NormalTok{, extracellular\_mM}\OperatorTok{=}\DecValTok{4}\NormalTok{, valence}\OperatorTok{=}\DecValTok{1}\NormalTok{)}
\NormalTok{e\_k }\OperatorTok{=}\NormalTok{ nernst\_potential(k\_ion)}
\BuiltInTok{print}\NormalTok{(}\SpecialStringTok{f"E\_K = }\SpecialCharTok{\{}\NormalTok{e\_k}\SpecialCharTok{:.1f\}}\SpecialStringTok{ mV"}\NormalTok{)}
\CommentTok{\# Expected: E\_K = {-}89.7 mV (close to resting potential of {-}70 mV; K+ is dominant)}

\CommentTok{\# Goldman{-}Hodgkin{-}Katz equation: resting membrane potential}
\CommentTok{\# Permeabilities at rest: P\_K = 1.0, P\_Na = 0.04, P\_Cl = 0.45}
\NormalTok{ions }\OperatorTok{=}\NormalTok{ [}
\NormalTok{    IonConcentration(}\StringTok{"K+"}\NormalTok{,  }\DecValTok{150}\NormalTok{, }\DecValTok{4}\NormalTok{,   valence}\OperatorTok{=}\DecValTok{1}\NormalTok{),}
\NormalTok{    IonConcentration(}\StringTok{"Na+"}\NormalTok{, }\DecValTok{15}\NormalTok{,  }\DecValTok{145}\NormalTok{, valence}\OperatorTok{=}\DecValTok{1}\NormalTok{),}
\NormalTok{    IonConcentration(}\StringTok{"Cl{-}"}\NormalTok{, }\DecValTok{10}\NormalTok{,  }\DecValTok{120}\NormalTok{, valence}\OperatorTok{={-}}\DecValTok{1}\NormalTok{),}
\NormalTok{]}
\NormalTok{permeabilities }\OperatorTok{=}\NormalTok{ \{}\StringTok{"K+"}\NormalTok{: }\FloatTok{1.0}\NormalTok{, }\StringTok{"Na+"}\NormalTok{: }\FloatTok{0.04}\NormalTok{, }\StringTok{"Cl{-}"}\NormalTok{: }\FloatTok{0.45}\NormalTok{\}}
\NormalTok{v\_rest }\OperatorTok{=}\NormalTok{ goldman\_equation(ions, permeabilities)}
\BuiltInTok{print}\NormalTok{(}\SpecialStringTok{f"V\_rest (GHK) = }\SpecialCharTok{\{}\NormalTok{v\_rest}\SpecialCharTok{:.1f\}}\SpecialStringTok{ mV"}\NormalTok{)}
\CommentTok{\# Expected: V\_rest (GHK) ≈ {-}70.0 mV}
\end{Highlighting}
\end{Shaded}

\begin{quote}
\textbf{Try it yourself:} Change \texttt{P\_Na} from 0.04 to 0.6
(simulating an action potential peak) and observe how \texttt{V\_rest}
shifts toward \texttt{E\_Na} (\texorpdfstring{\ensuremath{\approx}}{~} +60 mV).
\end{quote}

\begin{center}\rule{0.5\linewidth}{0.5pt}\end{center}

\emph{Source note: cell-biology helpers cover membrane biophysics,
signal transduction, and organelle inventory. Mermaid diagrams and
Nernst plots provide the visual companion.}

\newpage

\section{Cell Theory and Cell Types}\label{cell-theory-and-cell-types}

\label{sec:unit_II_cell_theory}

\begin{quote}
\textbf{Ch 5} \texorpdfstring{\ensuremath{\cdot}}{.} Level 1/3 \texorpdfstring{\ensuremath{\cdot}}{.} 45 min read \texorpdfstring{\ensuremath{\cdot}}{.} 50 min lecture \texorpdfstring{\ensuremath{\cdot}}{.}
Prerequisites: \cref{sec:unit_I_macromolecules}
\end{quote}

\subsection{Learning Objectives}\label{learning-objectives-7}

\begin{enumerate}
\def\labelenumi{\arabic{enumi}.}
\tightlist
\item
  State the three original tenets of cell theory (Schleiden, Schwann,
  Virchow) and the modern additions.
\item
  Compare the structural organisation of prokaryotic, eukaryotic, and
  archaeal cells at the molecular level.
\item
  Describe the endosymbiotic \citep{margulis1967} theory and evaluate
  the molecular evidence supporting it.
\item
  Classify cells by metabolic lifestyle (autotrophic vs.~heterotrophic)
  and oxygen relationship.
\item
  Quantify the surface-area-to-volume relationship and derive the
  mathematical constraints on cell size.
\item
  Explain microscopy techniques from bright-field to cryo-EM and
  super-resolution methods.
\item
  Describe the minimal cell concept (JCVI-syn3.0) and its implications
  for defining life.
\item
  Evaluate extremophile cells as evidence for the range of possible
  cellular adaptations and their relevance to astrobiology.
\end{enumerate}

\subsubsection{Study Blueprint}\label{study-blueprint-7}

\begin{itemize}
\tightlist
\item
  \textbf{Big idea:} Cells are bounded, evolving systems whose size and
  organization reflect physical constraints.
\item
  \textbf{Core concepts:} cell theory, surface area, microscopy,
  prokaryote/eukaryote comparison.
\item
  \textbf{Framework alignment:} Vision \& Change: Structure and
  function, Systems, Information flow, exchange, and storage; AP
  Biology: Systems Interactions, Information Storage and Transmission;
  NGSS-style topics: Structure and Function.
\item
  \textbf{Model or quantitative lens:} Surface-area-to-volume scaling.
\item
  \textbf{Data skill:} Infer cellular constraints from measurements,
  micrographs, and scale bars.
\item
  \textbf{Practice cadence:} Visual Representations, Questions and
  Methods, Argumentation.
\item
  \textbf{Common misconception to repair:} Cells are not just small bags
  of fluid; boundaries and internal organization create function.
\item
  \textbf{Primary lab:} \cref{sec:lab_unit_II_cell_theory}.
\item
  \textbf{Question bank:} \cref{sec:q_unit_II_cell_theory}.
\item
  \textbf{Transfer task:} Transfer scale reasoning to eggs, neurons,
  root hairs, and microbial colonies.
\item
  \textbf{Bridge to computation:}
  \breaktt{biology.cell.cell\_biology.diffusion\_flux}.
\end{itemize}

\begin{center}\rule{0.5\linewidth}{0.5pt}\end{center}

\begin{quote}
\textbf{Opening Vignette: The First Cell You Ever Saw}

In 1665, Robert Hooke pressed a thin sliver of cork against the lens of
a compound microscope he had constructed himself and published his
observations in \emph{Micrographia} --- a book so popular that Samuel
Pepys stayed up until 2 a.m. reading it. Hooke saw a regular array of
empty compartments --- the remnant cell walls of dead plant cells ---
and named them \textbf{cellulae} (from Latin, ``small rooms''). He had
no idea that cells were alive. That insight required a further 174 years
of microscopy, culminating in the cell theory formulations of
\citet{schleiden1838}, Schwann (1839), and \citet{virchow1855}.

Today, a single human body contains approximately 37.2 trillion cells
(Bianconi et al., 2013, \emph{Annals of Human Biology}) --- a number so
large that if laid end-to-end they would circle the Earth 200 times.
Every one of those cells descended, by an unbroken chain of cell
division, from a single fertilised egg roughly nine months before your
birth. And every cell alive on Earth today is the product of 3.7 billion
years of unbroken cell division since the origin of life --- not one
cell has ever been created from scratch by spontaneous generation since
Leuwenhoek looked through his microscope in 1674.

\emph{Primary source: Bianconi, E. et al.~(2013). An estimation of the
number of cells in the human body. Annals of Human Biology, 40(6),
463--471.}
\end{quote}

\begin{center}\rule{0.5\linewidth}{0.5pt}\end{center}

\subsection{The Cell as the Fundamental Unit of
Life}\label{the-cell-as-the-fundamental-unit-of-life}

\subsubsection{Historical Development of Cell
Theory}\label{historical-development-of-cell-theory}

The development of cell theory spans over two centuries and represents
one of the most important unifying principles in biology.

\textbf{Robert \citet{hooke1665}} first observed ``cells'' (actually
dead cell walls) in cork using a compound microscope he built himself.
He published his observations in \emph{Micrographia}, coining the term
``cell'' because the structures reminded him of monks' rooms (cellae) in
a monastery.

\textbf{Antonie van Leeuwenhoek (1674-1683)} constructed single-lens
microscopes achieving about 270x magnification and observed
living cells for the first time --- bacteria (``animalcules''),
protists, sperm cells, and red blood cells in capillary blood flow.

\textbf{Matthias \citet{schleiden1838}}, a botanist, concluded that most
plant tissues are composed of cells and that the cell is the basic unit
of plant structure.

\textbf{Theodor \citet{schwann1839}}, a zoologist and Schleiden's
colleague, extended this principle to animals, proposing that most
animal tissues are also composed of cells. Together, Schleiden and
Schwann formulated the first two tenets of cell theory.

\textbf{Rudolf \citet{virchow1855}} contributed the third tenet with his
famous dictum \emph{omnis cellula e cellula} (``every cell from a
cell''), establishing that cells arise primarily from pre-existing cells
through division. This principle demolished the lingering notion of
spontaneous generation and established the continuity of life.

\begin{figure}[htbp]
\centering
\includegraphics[width=0.82\linewidth,height=4.2in,keepaspectratio]{../figures/mermaid_inline/inline_mermaid_0031_9e71ef6de137.png}
\caption{Timeline of key milestones in cell theory development, from Hooke's first observations to the modern synthetic minimal cell.}
\end{figure}

\subsubsection{The Three Original
Postulates}\label{the-three-original-postulates}

The \textbf{cell theory}, formalised by \citet{schleiden1838},
\citet{schwann1839}, and \citet{virchow1855}, rests on three founding
principles:

\begin{enumerate}
\def\labelenumi{\arabic{enumi}.}
\tightlist
\item
  \textbf{Most living organisms are composed of one or more cells.}
\item
  \textbf{The cell is the basic structural and functional unit of
  cellular organisms.}
\item
  \textbf{Most cells arise from pre-existing cells} (\emph{omnis cellula
  e cellula} --- Virchow, 1855).
\end{enumerate}

A corollary of the third tenet is that life is a continuous lineage ---
every cell on Earth today descends without interruption from the first
cells, approximately 3.5--3.8 billion years ago. The oldest
microfossils, found in Pilbara (Western Australia) and Barberton (South
Africa), date to about 3.5 Ga and show morphologies consistent
with filamentous prokaryotes.

\subsubsection{Modern Additions to Cell
Theory}\label{modern-additions-to-cell-theory}

Modern cell biology has expanded the original three postulates:

\begin{enumerate}
\def\labelenumi{\arabic{enumi}.}
\setcounter{enumi}{3}
\item
  \textbf{The cell contains heritable information (DNA) that directs its
  activities and is passed to daughter cells.} This was established by
  the work of Avery, MacLeod, and McCarty (1944) and confirmed by
  \citet{hershey1952}.
\item
  \textbf{Most cells have the same basic chemical composition.} Most
  known cells use DNA as genetic material, RNA as intermediary,
  \hyperref[gl:protein]{\textbf{protein}}s as catalysts, and
  \hyperref[gl:phospholipid-bilayer]{\textbf{phospholipid bilayer}}s as
  membranes.
\item
  \textbf{Energy flow (metabolism and biochemistry) occurs within
  cells.} Most metabolic transformations ---
  \hyperref[gl:glycolysis]{\textbf{glycolysis}}, the TCA cycle,
  \hyperref[gl:oxidative-phosphorylation]{\textbf{oxidative
  phosphorylation}} --- take place within or across cell membranes.
\item
  \textbf{Cells contain the information necessary for their own
  reproduction.} The \hyperref[gl:genome]{\textbf{genome}} encodes not
  just structural components but the regulatory logic for cell division,
  differentiation, and programmed cell death.
\end{enumerate}

\begin{quote}
\textbf{Clinical Connection: Virchow and the Origin of Cancer Biology}
Virchow's principle \emph{omnis cellula e cellula} had a profound
clinical implication: if most cells come from pre-existing cells, then
cancer cells must also arise from normal cells through transformation.
Virchow himself applied this reasoning to pathology, founding the field
of cellular pathology. Today, cancer biology rests on the understanding
that \hyperref[gl:mutation]{\textbf{mutation}}s accumulate in somatic
cells, transforming them into malignant clones --- a direct intellectual
descendant of Virchow's third postulate. see
\cref{sec:unit_II_cell_signaling} (Cell Signalling) for oncogenes and
tumour suppressors.
\end{quote}

\begin{quote}
\textbf{Concept Check 1:} Viruses are not considered ``alive'' by cell
theory standards. List three properties of viruses that exclude them
from cell theory, and one property that challenges the boundary between
living and non-living.
\end{quote}

\begin{center}\rule{0.5\linewidth}{0.5pt}\end{center}

\subsection{Size, Scale, and the Surface-to-Volume
Constraint}\label{size-scale-and-the-surface-to-volume-constraint}

\subsubsection{Orders of Magnitude in
Biology}\label{orders-of-magnitude-in-biology}

Biology spans an enormous range of scales. Understanding this range is
essential for appreciating why cells occupy a specific size
\hyperref[gl:niche]{\textbf{niche}}.

{\def\LTcaptype{none} % do not increment counter
\begin{longtable}[]{@{}lll@{}}
\toprule\noalign{}
Structure & Approximate size & Scale \\
\midrule\noalign{}
\endhead
\bottomrule\noalign{}
\endlastfoot
Water molecule & 0.28 nm & Angstrom \\
Glucose molecule & 0.7 nm & Nanometre \\
Phospholipid bilayer thickness & 7--8 nm & Nanometre \\
\hyperref[gl:ribosome]{\textbf{Ribosome}} & 25 nm & Nanometre \\
HIV virus & 120 nm & Nanometre \\
Mitochondrion & 1--5 \texorpdfstring{\ensuremath{\mu}}{mu}m & Micrometre \\
\emph{E. coli} & 1--2 \texorpdfstring{\ensuremath{\mu}}{mu}m & Micrometre \\
Red blood cell & 7--8 \texorpdfstring{\ensuremath{\mu}}{mu}m & Micrometre \\
Typical animal cell & 10--30 \texorpdfstring{\ensuremath{\mu}}{mu}m & Micrometre \\
Typical plant cell & 30--100 \texorpdfstring{\ensuremath{\mu}}{mu}m & Micrometre \\
\emph{Thiomargarita namibiensis} & 750 \texorpdfstring{\ensuremath{\mu}}{mu}m & Sub-millimetre \\
Frog egg & 1 mm & Millimetre \\
Ostrich egg (cell) & 15 cm & Centimetre \\
\end{longtable}
}

\subsubsection{The Mathematics of Cell
Size}\label{the-mathematics-of-cell-size}

Cells are small --- almost without exception. Why?

The rate of metabolism (nutrient consumption, waste generation) scales
with \textbf{volume} (\(V \propto r^3\)). The rate of exchange with the
environment (nutrient uptake, waste removal) scales with \textbf{surface
area} (\(A \propto r^2\)). The ratio:

\[\frac{A}{V} = \frac{4\pi r^2}{\frac{4}{3}\pi r^3} = \frac{3}{r} \tag{5.1}\]

As \(r\) increases, \(A/V\) decreases. A sphere with \(r = 1\) \texorpdfstring{\ensuremath{\mu}}{mu}m has
\(A/V = 3\) \texorpdfstring{\ensuremath{\mu}}{mu}m\(^{-1}\). With \(r = 1\) mm (1,000x larger),
\(A/V = 0.003\) \texorpdfstring{\ensuremath{\mu}}{mu}m\(^{-1}\). Large cells simply cannot exchange
nutrients fast enough for their metabolic demand.

\subsection{Worked Example: Surface-to-Volume
Calculation}\label{worked-example-surface-to-volume-calculation}

\emph{Problem:} A spherical bacterium has radius \(r = 0.5\) \texorpdfstring{\ensuremath{\mu}}{mu}m.
Calculate (a) the surface area, (b) the volume, (c) the \(A/V\) ratio.
Then repeat for a spherical eukaryotic cell with \(r = 10\) \texorpdfstring{\ensuremath{\mu}}{mu}m and
compare.

\emph{Solution:}

For the bacterium (\(r = 0.5\) \texorpdfstring{\ensuremath{\mu}}{mu}m):

\[A = 4\pi r^2 = 4\pi(0.5)^2 = 3.14 \; \mu\text{m}^2 \tag{5.2}\]

\[V = \frac{4}{3}\pi r^3 = \frac{4}{3}\pi(0.5)^3 = 0.524 \; \mu\text{m}^3 \tag{5.3}\]

\[\frac{A}{V} = \frac{3.14}{0.524} = 6.0 \; \mu\text{m}^{-1} \tag{5.4}\]

For the eukaryotic cell (\(r = 10\) \texorpdfstring{\ensuremath{\mu}}{mu}m):

\[A = 4\pi(10)^2 = 1,257 \; \mu\text{m}^2 \tag{5.5}\]

\[V = \frac{4}{3}\pi(10)^3 = 4,189 \; \mu\text{m}^3 \tag{5.6}\]

\[\frac{A}{V} = \frac{1,257}{4,189} = 0.30 \; \mu\text{m}^{-1} \tag{5.7}\]

The bacterium has a 20-fold higher \(A/V\) ratio, enabling far more
efficient diffusion-based exchange per unit metabolic volume.

\subsubsection{Diffusion Time and Cell
Size}\label{diffusion-time-and-cell-size}

Diffusion time scales with the square of distance:

\[t = \frac{x^2}{2D} \tag{5.8}\]

where \(D\) is the diffusion coefficient. For a small molecule
(\(D \approx 10^{-9}\) m\(^2\)/s) to diffuse across a 1 \texorpdfstring{\ensuremath{\mu}}{mu}m cell takes
about 0.5 ms. Across a 1 mm cell, it takes about 500 s
(over 8 minutes). This places an absolute upper limit on cell size for
diffusion-dependent processes.

\subsubsection{Exceptions to Small Cell
Size}\label{exceptions-to-small-cell-size}

\textbf{Exceptions} to small cell size achieve large size by: -
\textbf{Increasing surface area:} intestinal cells have microvilli
(amplify \(A\) by about 600x); root hair cells extend long
projections - \textbf{Reducing metabolically active volume:} plant
vacuoles occupy \textgreater80\% of cell volume but are largely
metabolically inert - \textbf{Reducing effective diffusion distance:}
nerve axons can be \textgreater1 m long but about 1--20 \texorpdfstring{\ensuremath{\mu}}{mu}m in diameter;
active transport supplements diffusion - \textbf{Cytoplasmic streaming:}
in giant plant cells (e.g., \emph{Chara} internodal cells,
\textgreater10 cm), \hyperref[gl:actin]{\textbf{actin}}-myosin-driven
streaming circulates \hyperref[gl:cytoplasm]{\textbf{cytoplasm}} at
about 60 \texorpdfstring{\ensuremath{\mu}}{mu}m/s, overcoming diffusion limitations -
\textbf{Multinucleation:} skeletal muscle fibres (up to 30 cm long)
contain hundreds of nuclei, each governing a local cytoplasmic domain

\begin{quote}
\textbf{Concept Check 2:} \emph{Thiomargarita namibiensis} is a
bacterium with a diameter of about 750 \texorpdfstring{\ensuremath{\mu}}{mu}m --- visible to the
naked eye. How does it survive despite having an \(A/V\) ratio of about
0.008 \texorpdfstring{\ensuremath{\mu}}{mu}m\(^{-1}\)? (Hint: consider what fills \textgreater95\% of its
volume.)
\end{quote}

\subsubsection{The Cube Law and Why Giant Cells Are
Impossible}\label{the-cube-law-and-why-giant-cells-are-impossible}

The surface-to-volume problem can be sharpened into what physiologists
call the \textbf{cube law of metabolism}. If we model a cell as a sphere
of radius \(r\), mass \(m\), and uniform density \texorpdfstring{\ensuremath{\rho}}{rho}, then:

\begin{equation}
m = \rho \cdot \frac{4}{3}\pi r^3, \qquad A = 4\pi r^2
\label{eq:unit_II_cube_law_basis}
\end{equation}

Suppose every unit volume of cytoplasm consumes oxygen at a constant
rate \(q\) (mol O\(_2\) s\(^{-1}\) m\(^{-3}\)). The total metabolic
demand \(\dot{M}_\text{demand}\) scales as:

\[\dot{M}_\text{demand} = q \cdot V = q \cdot \tfrac{4}{3}\pi r^3 \tag{5.10}\]

Maximum oxygen flux through the membrane is set by Fick's law and is
bounded by surface area:

\[\dot{M}_\text{supply} = P_{O_2} \cdot \Delta[O_2] \cdot A = P_{O_2} \cdot \Delta[O_2] \cdot 4\pi r^2 \tag{5.11}\]

A viable cell requires
\(\dot{M}_\text{supply} \ge \dot{M}_\text{demand}\). Setting these equal
yields a \textbf{maximum admissible radius}:

\begin{equation}
r_\text{max} = \frac{3 \, P_{O_2} \, \Delta[O_2]}{q}
\label{eq:unit_II_rmax}
\end{equation}

For typical mammalian parameters (\(P_{O_2} \approx 1\times10^{-4}\)
m/s, \(\Delta[O_2] \approx 0.2\) mM, \(q \approx 1\) mol O\(_2\)
m\(^{-3}\) s\(^{-1}\)), \(r_\text{max} \approx 60\) \texorpdfstring{\ensuremath{\mu}}{mu}m --- close to the
upper bound for actively respiring animal cells. Larger cells must
reduce \(q\) (e.g., plant vacuoles), increase \(A\) (microvilli,
infoldings), or accept partial anoxia in their cores. This is the
\textbf{cube law of metabolism}: demand scales with \(r^3\), supply with
\(r^2\), so any uniform cell hits a hard ceiling on size.

\begin{quote}
\textbf{Concept Check 3a:} Using \cref{eq:unit_II_rmax} with the
parameters above, by what factor would \(q\) have to drop to allow a 1
mm diameter respiring cell to survive without surface amplification?
\end{quote}

\subsubsection{Quantitative Comparison: Prokaryote
vs.~Eukaryote}\label{quantitative-comparison-prokaryote-vs.-eukaryote}

Prokaryotes and eukaryotes occupy fundamentally different scales of
organisation. The differences in linear dimension translate to even
larger gaps in volume, surface area, and information content:

{\def\LTcaptype{none} % do not increment counter
\begin{longtable}[]{@{}
  >{\raggedright\arraybackslash}p{(\linewidth - 6\tabcolsep) * \real{0.1176}}
  >{\raggedright\arraybackslash}p{(\linewidth - 6\tabcolsep) * \real{0.3382}}
  >{\raggedright\arraybackslash}p{(\linewidth - 6\tabcolsep) * \real{0.3088}}
  >{\raggedright\arraybackslash}p{(\linewidth - 6\tabcolsep) * \real{0.2353}}@{}}
\toprule\noalign{}
\begin{minipage}[b]{\linewidth}\raggedright
Property
\end{minipage} & \begin{minipage}[b]{\linewidth}\raggedright
\emph{E. coli} (prokaryote)
\end{minipage} & \begin{minipage}[b]{\linewidth}\raggedright
HeLa cell (eukaryote)
\end{minipage} & \begin{minipage}[b]{\linewidth}\raggedright
Ratio (Euk/Prok)
\end{minipage} \\
\midrule\noalign{}
\endhead
\bottomrule\noalign{}
\endlastfoot
Length / diameter & 2 \texorpdfstring{\ensuremath{\mu}}{mu}m \texorpdfstring{\ensuremath{\times}}{x} 0.5 \texorpdfstring{\ensuremath{\mu}}{mu}m & about 20 \texorpdfstring{\ensuremath{\mu}}{mu}m &
about 10--40 \texorpdfstring{\ensuremath{\times}}{x} \\
Volume & about 1 \texorpdfstring{\ensuremath{\mu}}{mu}m\(^3\) & about 4,000 \texorpdfstring{\ensuremath{\mu}}{mu}m\(^3\) &
about 4,000 \texorpdfstring{\ensuremath{\times}}{x} \\
Surface area & about 6 \texorpdfstring{\ensuremath{\mu}}{mu}m\(^2\) & about 5,000 \texorpdfstring{\ensuremath{\mu}}{mu}m\(^2\)
& about 830 \texorpdfstring{\ensuremath{\times}}{x} \\
\(A/V\) & about 6 \texorpdfstring{\ensuremath{\mu}}{mu}m\(^{-1}\) & about 1.25 \texorpdfstring{\ensuremath{\mu}}{mu}m\(^{-1}\)
& 0.21 \texorpdfstring{\ensuremath{\times}}{x} (lower) \\
Genome size & 4.6 Mbp & 6.4 Gbp (diploid) & about 1,400 \texorpdfstring{\ensuremath{\times}}{x} \\
Protein number & about 3 \texorpdfstring{\ensuremath{\times}}{x} 10\(^6\) & about 2 \texorpdfstring{\ensuremath{\times}}{x}
10\(^{10}\) & about 7,000 \texorpdfstring{\ensuremath{\times}}{x} \\
Doubling time & 20 min (rich media) & about 24 h & 70 \texorpdfstring{\ensuremath{\times}}{x}
slower \\
ATP turnover (molecules/s) & about 10\(^9\) &
about 10\(^{10}\) & about 10 \texorpdfstring{\ensuremath{\times}}{x} \\
\end{longtable}
}

\textbf{Worked Example: Why eukaryotes need internal membranes.} A HeLa
cell with \(r = 10\) \texorpdfstring{\ensuremath{\mu}}{mu}m has about 1.25 \texorpdfstring{\ensuremath{\mu}}{mu}m\(^{-1}\) surface-to-volume ---
about 5-fold less than \emph{E. coli}. Yet the eukaryotic cell sustains
roughly 10\texorpdfstring{\ensuremath{\times}}{x} higher absolute metabolic flux. The arithmetic closes if the
eukaryote \emph{adds internal membrane surface area} equivalent to the
plasma membrane's deficit. Indeed, the inner mitochondrial membrane in a
single hepatocyte provides about \(3 \times 10^4\) \texorpdfstring{\ensuremath{\mu}}{mu}m\(^2\)
of additional respiratory surface --- roughly 6\texorpdfstring{\ensuremath{\times}}{x} the plasma membrane
area. Compartmentalisation is therefore not optional aesthetics: it is
the topological solution to the cube-law constraint on a body plan that
contains 1,000\texorpdfstring{\ensuremath{\times}}{x} more cytoplasm than a bacterium.

\begin{quote}
\textbf{Concept Check 3b:} A spherical eukaryotic cell with \(r = 10\)
\texorpdfstring{\ensuremath{\mu}}{mu}m has plasma-membrane area \(A_\text{PM} \approx 1{,}260\) \texorpdfstring{\ensuremath{\mu}}{mu}m\(^2\). If
respiratory demand scales with cytoplasmic volume and the cell needs 6\texorpdfstring{\ensuremath{\times}}{x}
the plasma-membrane area to support its metabolism, estimate the total
inner membrane area its mitochondria must provide.
\end{quote}

\begin{center}\rule{0.5\linewidth}{0.5pt}\end{center}

\subsection{Microscopy: Seeing the
Cell}\label{microscopy-seeing-the-cell}

Understanding cells requires seeing them. The history of cell biology is
inseparable from advances in microscopy.

\begin{figure}[htbp]
\centering
\includegraphics[width=0.82\linewidth,height=4.2in,keepaspectratio]{../figures/mermaid_inline/inline_mermaid_0032_ee5d89460c88.png}
\caption{Overview of microscopy techniques used in cell biology, categorised by resolution capability.}
\end{figure}

\subsubsection{Light Microscopy}\label{light-microscopy}

The \textbf{resolution limit} of light microscopy is governed by the
Abbe diffraction limit:

\[d = \frac{\lambda}{2 \cdot \text{NA}} \tag{5.9}\]

where \texorpdfstring{\ensuremath{\lambda}}{lambda} is the wavelength of light and NA is the numerical aperture of
the objective. For visible light (\(\lambda \approx 550\) nm) and a
high-quality oil-immersion objective (NA = 1.4): \(d \approx 200\) nm.

\textbf{Bright-field microscopy:} Transmitted light; specimens typically
stained (H\&E, \hyperref[gl:gram-stain]{\textbf{Gram stain}}). Simple
but limited contrast for living cells.

\textbf{Phase-contrast microscopy:} Converts phase differences (from
refractive index variations in transparent specimens) into amplitude
differences visible as contrast. Ideal for observing living, unstained
cells. Invented by Frits Zernike (Nobel Prize in Physics, 1953).

\textbf{Differential interference contrast (DIC/Nomarski):} Uses
polarised light split into two beams that pass through slightly
different optical paths in the specimen. Produces a pseudo-3D relief
image with excellent contrast. Eliminates the ``halo'' artifact of
phase-contrast.

\textbf{Fluorescence microscopy:} Specimens labelled with fluorophores
(antibodies conjugated to FITC, Cy3, Cy5; or genetically encoded GFP and
derivatives). Excitation light of specific wavelength causes fluorophore
emission at a longer wavelength. Enables specific protein localisation,
live-cell imaging, and multiplexing (multiple colours simultaneously).

\textbf{Confocal laser scanning microscopy:} Uses a pinhole to reject
out-of-focus light, producing optical sections about 0.5--1 \texorpdfstring{\ensuremath{\mu}}{mu}m
thick. Z-stacks can be reconstructed into 3D images. Essential for thick
specimens and co-localisation studies.

\textbf{Total internal reflection fluorescence (TIRF):} Evanescent wave
illuminates about 100 nm of the cell adjacent to the coverslip surface.
Used to study membrane dynamics, vesicle fusion events, and
single-molecule behaviour at the plasma membrane.

\subsubsection{Electron Microscopy}\label{electron-microscopy}

\textbf{Transmission electron microscopy (TEM):} Electrons pass through
ultrathin sections (about 50--70 nm) stained with heavy metals
(uranyl acetate, osmium tetroxide). Resolution about 0.1 nm.
Reveals ultrastructural details of
\hyperref[gl:organelle]{\textbf{organelle}}s, membranes, and
macromolecular complexes.

\textbf{Scanning electron microscopy (SEM):} Electrons scan the surface
of gold- or platinum-coated specimens. Produces stunning 3D topographic
images of cell surfaces, tissue architecture, and microorganisms.
Resolution about 1--5 nm.

\textbf{Cryo-electron microscopy (cryo-EM):} Specimens flash-frozen in
vitreous ice (no staining or fixation) and imaged at liquid nitrogen
temperature. Single-particle analysis reconstructs 3D structures of
proteins and complexes at near-atomic resolution (2--4 angstroms).
Awarded the Nobel Prize in Chemistry 2017 (Jacques Dubochet, Joachim
Frank, Richard Henderson). Cryo-EM has revolutionised structural
biology, enabling structures of membrane proteins, ribosomes, and viral
capsids that resist crystallisation.

\begin{quote}
\textbf{Clinical Connection: Cryo-EM and Drug Discovery} Cryo-EM has
become indispensable in pharmaceutical development. The structure of the
SARS-CoV-2 spike protein was solved by cryo-EM within weeks of the
pandemic onset, enabling rapid \hyperref[gl:vaccine]{\textbf{vaccine}}
design (Wrapp et al., 2020, \emph{Science}). Drug binding sites on ion
channels (e.g., TRPV1 pain receptor) and GPCRs have been mapped at
near-atomic resolution, accelerating rational drug design. see
\cref{sec:unit_II_membrane_transport} for membrane protein structures.
\end{quote}

\subsubsection{Super-Resolution
Microscopy}\label{super-resolution-microscopy}

Super-resolution methods bypass the Abbe diffraction limit, achieving
resolutions of 20--50 nm with light:

\textbf{STED (Stimulated Emission Depletion):} A donut-shaped depletion
beam suppresses fluorescence from the periphery of the excitation spot,
shrinking the effective point spread function. Resolution
about 30--50 nm. Developed by Stefan Hell (Nobel Prize in
Chemistry, 2014).

\textbf{STORM (Stochastic Optical Reconstruction Microscopy):}
Individual fluorophores are stochastically activated, imaged, and
precisely localised over thousands of frames. The composite image
achieves about 20 nm resolution. Co-developed by Xiaowei
Zhuang.

\textbf{PALM (Photoactivated Localisation Microscopy):} Similar to STORM
but uses photoactivatable fluorescent proteins (e.g., mEos, Dendra2).
Particularly suited for live-cell imaging of genetically encoded
constructs. Co-developed by Eric Betzig (Nobel Prize in Chemistry,
2014).

\begin{quote}
\textbf{Concept Check 3:} A researcher wants to study the real-time
dynamics of clathrin-coated pit formation at the plasma membrane of a
living cell. Which microscopy technique would be most appropriate and
why?
\end{quote}

\begin{center}\rule{0.5\linewidth}{0.5pt}\end{center}

\subsection{Prokaryotic Cells}\label{prokaryotic-cells}

Prokaryotes (\textbf{pro} = before, \textbf{karyon} = nucleus) lack
membrane-bound organelles. They include the domains \textbf{Bacteria}
and \textbf{\hyperref[gl:archaea]{Archaea}}.

\subsubsection{Key Features}\label{key-features}

{\def\LTcaptype{none} % do not increment counter
\begin{longtable}[]{@{}
  >{\raggedright\arraybackslash}p{(\linewidth - 4\tabcolsep) * \real{0.3182}}
  >{\raggedright\arraybackslash}p{(\linewidth - 4\tabcolsep) * \real{0.3636}}
  >{\raggedright\arraybackslash}p{(\linewidth - 4\tabcolsep) * \real{0.3182}}@{}}
\toprule\noalign{}
\begin{minipage}[b]{\linewidth}\raggedright
Feature
\end{minipage} & \begin{minipage}[b]{\linewidth}\raggedright
Bacteria
\end{minipage} & \begin{minipage}[b]{\linewidth}\raggedright
Archaea
\end{minipage} \\
\midrule\noalign{}
\endhead
\bottomrule\noalign{}
\endlastfoot
Diameter & 0.2--10 \texorpdfstring{\ensuremath{\mu}}{mu}m & 0.2--10 \texorpdfstring{\ensuremath{\mu}}{mu}m \\
Nucleus & Absent; nucleoid region & Absent; nucleoid \\
DNA & Circular; single \hyperref[gl:chromosome]{\textbf{chromosome}};
often \hyperref[gl:plasmid]{\textbf{plasmid}}s & Circular;
\hyperref[gl:histone]{\textbf{histone}}s present \\
Cell wall & Peptidoglycan (most) & Pseudopeptidoglycan or S-layer \\
Ribosomes & 70S (30S + 50S subunits) & 70S similar; unique rRNA \\
Membrane lipids & Ester-linked; straight-chain fatty acids &
Ether-linked; branched isoprenoids \\
\hyperref[gl:intron]{\textbf{Intron}}s & Rare & Present \\
\hyperref[gl:transcription]{\textbf{Transcription}} machinery &
Bacterial-type RNA polymerase (4 subunits) &
\hyperref[gl:eukaryote]{\textbf{Eukaryote}}-like RNA polymerase (12+
subunits) \\
DNA replication & DnaA origin recognition & Eukaryote-like ORC system \\
\hyperref[gl:crispr-cas]{\textbf{CRISPR-Cas}} & Common & Very common \\
\end{longtable}
}

\begin{figure}[htbp]
\centering
\includegraphics[width=0.82\linewidth,height=4.2in,keepaspectratio]{../figures/mermaid_inline/inline_mermaid_0033_2fa50940f014.png}
\caption{Classification diagram comparing Bacteria, Archaea, and Eukaryotes. Note that Archaea are phylogenetically closer to Eukaryotes in many molecular features (RNA polymerase, histones, replication machinery), while mitochondria and chloroplasts derive from bacterial endosymbionts.}
\end{figure}

\subsubsection{Prokaryotic Cell
Structure}\label{prokaryotic-cell-structure}

\begin{itemize}
\tightlist
\item
  \textbf{Plasma membrane:} phospholipid bilayer + integral proteins.
  Site of chemiosmotic ATP synthesis (no mitochondria). In
  photosynthetic bacteria (cyanobacteria, purple bacteria), the membrane
  invaginates to form thylakoid-like structures for the light reactions.
\item
  \textbf{Cell wall:} structurally rigid; resists osmotic lysis.

  \begin{itemize}
  \tightlist
  \item
    \textbf{Gram-positive} bacteria: thick peptidoglycan layer (20--80
    nm), teichoic acids, no outer membrane. Crystal violet stain
    retained (purple).
  \item
    \textbf{Gram-negative} bacteria: thin peptidoglycan
    (about 2--7 nm) + outer membrane with LPS
    (lipopolysaccharide). Crystal violet not retained (pink after
    safranin counter-stain). The periplasmic space between inner and
    outer membranes contains degradative
    \hyperref[gl:enzyme]{\textbf{enzyme}}s and transport proteins.
  \end{itemize}
\item
  \textbf{Capsule:} polysaccharide layer external to cell wall;
  anti-phagocytic (evades immune system);
  \hyperref[gl:biofilm]{\textbf{biofilm}} formation. \emph{Streptococcus
  pneumoniae} with capsule is virulent; without capsule is avirulent
  (Griffith's transformation experiment, 1928).
\item
  \textbf{Flagellum:} rotating protein filament powered by proton motive
  force (bacteria) or ATP (archaea); each rotation of the bacterial
  flagellar motor about 100 Hz in \emph{E. coli}; composed of
  flagellin protein subunits; the basal body contains a rotary motor
  with stator and rotor components.
\item
  \textbf{Pili (fimbriae):} protein tubes for attachment, conjugation,
  biofilm; Type IV pili can retract to generate motile force
  (``twitching motility''); F-pili mediate DNA transfer during
  conjugation.
\item
  \textbf{Nucleoid:} condensed region containing the circular chromosome
  (typically 1--6 Mbp); not bounded by a membrane; associated with
  nucleoid-associated proteins (NAPs: HU, H-NS, Fis, IHF) that organise
  chromosome topology.
\item
  \textbf{Plasmids:} small circular DNA molecules (1--200 kb) carrying
  accessory \hyperref[gl:gene]{\textbf{gene}}s (antibiotic resistance,
  \hyperref[gl:virulence]{\textbf{virulence}} factors, metabolic
  capabilities); replicate independently; horizontally transferable
  between cells.
\end{itemize}

\begin{quote}
\textbf{Clinical Connection: Gram Staining and Antibiotic Selection} The
Gram stain, developed by Hans Christian Gram in 1884, remains the most
important rapid diagnostic test in clinical microbiology. Gram-positive
infections (thick peptidoglycan) respond to beta-lactam antibiotics
(penicillins, cephalosporins) that inhibit peptidoglycan cross-linking
by transpeptidases (PBPs). Gram-negative infections are intrinsically
more resistant because the outer membrane excludes many antibiotics;
treatment often requires aminoglycosides or carbapenems that penetrate
through porins. The rise of multi-drug-resistant Gram-negative bacteria
(e.g., carbapenem-resistant \emph{Klebsiella pneumoniae}) is a major
public health crisis. see \cref{sec:unit_II_cell_structure} for membrane
permeability.
\end{quote}

\begin{quote}
\textbf{Concept Check 4:} Why are archaea resistant to antibiotics that
target peptidoglycan synthesis (e.g., penicillin, vancomycin)?
\end{quote}

\subsubsection{The Prokaryotic Cytoskeleton --- A Recent
Revolution}\label{the-prokaryotic-cytoskeleton-a-recent-revolution}

For nearly a century, textbooks asserted that prokaryotes lack a
\hyperref[gl:cytoskeleton]{\textbf{cytoskeleton}}. The discovery of
bacterial cytoskeletal homologues over the past 25 years has overturned
that view: prokaryotes contain dedicated structural and force-generating
filaments that are evolutionary precursors of every major eukaryotic
cytoskeletal class.

{\def\LTcaptype{none} % do not increment counter
\begin{longtable}[]{@{}
  >{\raggedright\arraybackslash}p{(\linewidth - 6\tabcolsep) * \real{0.2787}}
  >{\raggedright\arraybackslash}p{(\linewidth - 6\tabcolsep) * \real{0.3279}}
  >{\raggedright\arraybackslash}p{(\linewidth - 6\tabcolsep) * \real{0.1311}}
  >{\raggedright\arraybackslash}p{(\linewidth - 6\tabcolsep) * \real{0.2623}}@{}}
\toprule\noalign{}
\begin{minipage}[b]{\linewidth}\raggedright
Bacterial protein
\end{minipage} & \begin{minipage}[b]{\linewidth}\raggedright
Eukaryotic homologue
\end{minipage} & \begin{minipage}[b]{\linewidth}\raggedright
Function
\end{minipage} & \begin{minipage}[b]{\linewidth}\raggedright
Year established
\end{minipage} \\
\midrule\noalign{}
\endhead
\bottomrule\noalign{}
\endlastfoot
\textbf{FtsZ} & Tubulin (GTPase fold) & Z-ring assembly at midcell;
recruits divisome; constricts during binary fission & 1991 (Bi \&
Lutkenhaus) \\
\textbf{MreB} & Actin (ATP-binding fold) & Helical filaments under inner
membrane; defines rod-shaped cell geometry; coordinates peptidoglycan
synthesis & 2001 (Jones et al.) \\
\textbf{Crescentin (CreS)} & Intermediate filaments (lamins) & Curves
\emph{Caulobacter crescentus} into its vibrioid shape & 2003 (Ausmees et
al.) \\
\textbf{ParM} & Actin & Pushes plasmid copies apart during segregation &
2002 \\
\textbf{Bactofilins (BacA/B)} & None close & Polar localisation;
cell-wall remodelling & 2010 \\
\end{longtable}
}

\textbf{FtsZ and bacterial cytokinesis.} FtsZ polymerises
GTP-dependently into protofilaments that form a contractile
\textbf{Z-ring} at the future division site. The Z-ring recruits the
divisome (FtsA, FtsW, FtsI/PBP3) which synthesises septal peptidoglycan
as the membrane invaginates. FtsZ treadmills around the ring at
about 30 nm/s --- strikingly similar in mechanism to actin
treadmilling, yet evolutionarily a tubulin homologue. The GTP-bound
monomer fold superimposes on alpha/beta-tubulin with RMSD \textless{} 2
Å despite \textless{} 10\% sequence identity, providing direct
structural evidence that the eukaryotic mitotic machinery descends from
a bacterial cell-division protein.

\textbf{MreB and rod-shaped morphology.} MreB filaments rotate around
the bacterial circumference, dragging the \textbf{Rod complex} (RodA,
MreC, MreD, PBP2) that synthesises lateral peptidoglycan. Loss of MreB
converts rod-shaped \emph{E. coli} into spheres (cocci) within minutes;
expression in mutant cocci restores rod shape. MreB is therefore both
structural (filament) and a guidance system (positional cue for
cell-wall synthesis), exactly mirroring eukaryotic actin's dual roles.

\textbf{Crescentin and curvature.} \emph{Caulobacter crescentus} gets
its name from a comma-shaped curve produced by a single helical bundle
of CreS filaments running along the inner curve of the cell. Mutational
loss of \emph{creS} yields straight rods. CreS is biophysically a
coiled-coil of the same family as nuclear lamins and cytokeratins.

These discoveries closed a major textbook gap: \textbf{most three
eukaryotic cytoskeletal classes (microfilaments, microtubules,
intermediate filaments) have prokaryotic ancestors}. The cytoskeleton
predates compartmentalisation by at least 1.5 billion years.

\begin{quote}
\textbf{Concept Check 4b:} A pharmaceutical company is screening for
narrow-spectrum antibiotics that block bacterial cell division without
affecting human cells. Why is FtsZ a more promising target than MreB?
(Consider the structural and functional homology to eukaryotic
counterparts and what happens when each protein is inhibited.)
\end{quote}

\subsubsection{Detailed Comparison: Bacteria
vs.~Archaea}\label{detailed-comparison-bacteria-vs.-archaea}

The discovery of Archaea by \citet{woese1977} as a separate domain
transformed our view of life's deep structure. Archaea were initially
described as ``prokaryotes that live in extreme environments,'' but
molecular phylogenetics revealed that Archaea share many features with
eukaryotes that bacteria do not --- a fact now central to debates about
eukaryotic origin (the ``Asgard archaea'' hypothesis). The table below
extends \cref{sec:unit_II_cell_theory} earlier comparison with molecular
detail relevant to drug development, evolution, and astrobiology.

{\def\LTcaptype{none} % do not increment counter
\begin{longtable}[]{@{}
  >{\raggedright\arraybackslash}p{(\linewidth - 6\tabcolsep) * \real{0.1458}}
  >{\raggedright\arraybackslash}p{(\linewidth - 6\tabcolsep) * \real{0.1667}}
  >{\raggedright\arraybackslash}p{(\linewidth - 6\tabcolsep) * \real{0.1458}}
  >{\raggedright\arraybackslash}p{(\linewidth - 6\tabcolsep) * \real{0.5417}}@{}}
\toprule\noalign{}
\begin{minipage}[b]{\linewidth}\raggedright
Feature
\end{minipage} & \begin{minipage}[b]{\linewidth}\raggedright
Bacteria
\end{minipage} & \begin{minipage}[b]{\linewidth}\raggedright
Archaea
\end{minipage} & \begin{minipage}[b]{\linewidth}\raggedright
Functional / clinical note
\end{minipage} \\
\midrule\noalign{}
\endhead
\bottomrule\noalign{}
\endlastfoot
Cell wall & Peptidoglycan (N-acetylmuramic acid + N-acetylglucosamine,
beta-1,4 linked, peptide cross-bridges) & Pseudopeptidoglycan (NAG +
N-acetyltalosaminuronic acid, beta-1,3 linked) or S-layer or
polysaccharide & Penicillins/vancomycin target NAM and D-Ala-D-Ala ---
both absent in archaea \\
Membrane lipids & Glycerol-3-phosphate ester-linked to straight-chain
fatty acids & Glycerol-1-phosphate ether-linked to branched isoprenoids
(phytanyl); often tetraether monolayers & Stereochemistry of glycerol
backbone is opposite (G3P vs G1P) --- the ``lipid divide'' \\
Membrane topology & Bilayer & Bilayer or monolayer (tetraether) &
Monolayer membranes are heat-stable to 113 °C \\
RNA polymerase & 4--5 subunits (\(\alpha_2\beta\beta'\omega\));
rifampicin-sensitive & 12+ subunits, eukaryote-like;
rifampicin-resistant & Drug development implication \\
TATA-binding protein (TBP) & Absent & Present (homologous to eukaryotic
TBP) & Promoter recognition is eukaryote-like \\
Histones & Absent (nucleoid-associated proteins HU/H-NS) & H3/H4-like
archaeal histones, octameric tetramers & Asgard archaea have histones
structurally indistinguishable from eukaryotic \\
Replication initiator & DnaA (single ori per chromosome) & Cdc6/Orc1
family; multiple origins & Eukaryotic-style \\
DNA topoisomerase & Topo IV; Gyrase (negative supercoiling) & Reverse
gyrase (positive supercoiling) in hyperthermophiles & Reverse gyrase
stabilises DNA at \textgreater80 °C \\
tRNA introns & Rare & Common (BHB intron motif) & Used as taxonomic
marker \\
Translation initiator & fMet-tRNA\(^{fMet}\) (formylated) &
Met-tRNA\(_i^{Met}\) (not formylated) & Same as eukaryotes \\
Sensitivity: chloramphenicol, streptomycin & Sensitive & Resistant &
Antibiotic discrimination \\
Sensitivity: anisomycin, diphtheria toxin & Resistant & Sensitive & Both
target eukaryote-like ribosomal features \\
CRISPR--Cas & Present in about 50\% & Present in
about 90\% & Adaptive immune systems against viruses \\
Energy metabolism & Photosynthetic (oxygenic, anoxygenic),
chemoautotrophic, heterotrophic & Methanogens, halophiles, sulphate
reducers, ammonia oxidisers & Methanogenesis is exclusive to archaea \\
\end{longtable}
}

\textbf{Why this matters.} Archaea share information-processing
machinery (transcription, translation, replication) with eukaryotes but
housekeeping/metabolic machinery with bacteria. The leading
``two-domain'' tree (Embley \& Williams; Spang et al.~2015) places
eukaryotes \emph{within} the Asgard archaea, with the mitochondrial
endosymbiont contributing the bacterial-style metabolism. The
bacteria--archaea distinction is not a curiosity --- it is the
structural map for the deepest split in the tree of life.

\begin{quote}
\textbf{Concept Check 4c:} Why does penicillin kill many bacteria but no
archaea, even archaea with cell walls? List the molecular features
penicillin targets that distinguish bacterial from archaeal cell
envelopes.
\end{quote}

\begin{center}\rule{0.5\linewidth}{0.5pt}\end{center}

\subsection{Eukaryotic Cells}\label{eukaryotic-cells}

Eukaryotic cells (\textbf{eu} = true; \textbf{karyon} = nucleus) contain
membrane-bound organelles and a true nucleus. They include protists,
fungi, plants, and animals.

\subsubsection{Organelle Inventory}\label{organelle-inventory}

\textbf{Core eukaryotic organelles and functions:}

{\def\LTcaptype{none} % do not increment counter
\begin{longtable}[]{@{}
  >{\raggedright\arraybackslash}p{(\linewidth - 6\tabcolsep) * \real{0.2647}}
  >{\raggedright\arraybackslash}p{(\linewidth - 6\tabcolsep) * \real{0.2353}}
  >{\raggedright\arraybackslash}p{(\linewidth - 6\tabcolsep) * \real{0.2353}}
  >{\raggedright\arraybackslash}p{(\linewidth - 6\tabcolsep) * \real{0.2647}}@{}}
\toprule\noalign{}
\begin{minipage}[b]{\linewidth}\raggedright
Organelle
\end{minipage} & \begin{minipage}[b]{\linewidth}\raggedright
Membrane
\end{minipage} & \begin{minipage}[b]{\linewidth}\raggedright
Function
\end{minipage} & \begin{minipage}[b]{\linewidth}\raggedright
Unique to
\end{minipage} \\
\midrule\noalign{}
\endhead
\bottomrule\noalign{}
\endlastfoot
Nucleus & Double (NE) & DNA storage; transcription & Eukaryotes \\
Mitochondria & Double & ATP synthesis; TCA cycle;
\hyperref[gl:apoptosis]{\textbf{apoptosis}} & Eukaryotes \\
Chloroplast & Double &
\hyperref[gl:photosynthesis]{\textbf{Photosynthesis}} & Plants, algae \\
Rough ER & Single & Protein synthesis and folding & Eukaryotes \\
Smooth ER & Single & Lipid synthesis; Ca\(^{2+}\) storage &
Eukaryotes \\
Golgi apparatus & Single & Protein modification, sorting, export &
Eukaryotes \\
Lysosome & Single & Intracellular digestion
(\hyperref[gl:ph]{\textbf{pH}} 4.5) & Animals (mainly) \\
Vacuole & Single & Water regulation; storage & Plants (large) \\
Peroxisome & Single & Fatty acid beta-oxidation; H\(_2\)O\(_2\) detox &
Eukaryotes \\
Ribosome & None & \hyperref[gl:translation]{\textbf{Translation}} & Most
cells \\
\hyperref[gl:cytoskeleton]{\textbf{Cytoskeleton}} & None & Structural;
motility; intracellular transport & Eukaryotes \\
Centrosome & None (MTOC) & Microtubule organising centre & Animal
cells \\
\end{longtable}
}

This inventory should not be read as a set of isolated boxes. Eukaryotic
organelles exchange lipids, ions, metabolites, and stress signals
through membrane-contact sites; mitochondria divide and fuse, the ER
wraps around many organelles, lysosomes function as nutrient-sensing
hubs, and biomolecular condensates add non-membrane compartments for
RNA, protein, and signalling control. A modern cell-biological
explanation names the compartment, the traffic route, the time scale,
and the evidence type: microscopy, perturbation, biochemical
fractionation, or single-cell/spatial data.

\subsubsection{Compartmentalisation: The Eukaryotic
Advantage}\label{compartmentalisation-the-eukaryotic-advantage}

The defining feature of eukaryotic cells is
\textbf{compartmentalisation} --- the segregation of biochemical
processes into membrane-bound organelles. This provides several
advantages:

\begin{enumerate}
\def\labelenumi{\arabic{enumi}.}
\tightlist
\item
  \textbf{Incompatible reactions can occur simultaneously:} Protein
  synthesis (cytoplasm, pH 7.2) and protein degradation (lysosomes, pH
  4.5) are separated by lysosomal membranes.
\item
  \textbf{Concentration of substrates and enzymes:} Mitochondrial matrix
  concentrates TCA cycle enzymes, increasing reaction rates.
\item
  \textbf{Regulation of gene expression:} The nuclear envelope separates
  transcription (nucleus) from translation (cytoplasm), enabling
  post-transcriptional regulation (splicing, mRNA export, mRNA
  stability).
\item
  \textbf{Expanded membrane surface area:} Internal membranes (ER,
  Golgi, mitochondria) provide enormous surface area for
  membrane-associated reactions.
\end{enumerate}

\begin{center}\rule{0.5\linewidth}{0.5pt}\end{center}

\subsection{Endosymbiotic Theory}\label{endosymbiotic-theory}

Lynn Margulis formalised the \textbf{endosymbiotic theory} (1967):
mitochondria and chloroplasts are the descendants of free-living
bacteria engulfed by a proto-eukaryotic host about 1.5--2
billion years ago.

\subsubsection{Evidence for
Endosymbiosis}\label{evidence-for-endosymbiosis}

\begin{enumerate}
\def\labelenumi{\arabic{enumi}.}
\tightlist
\item
  \textbf{Double membranes:} Both organelles have inner (bacterial
  plasma membrane) and outer (phagocytic membrane) bilayers.
\item
  \textbf{Circular DNA:} Mitochondrial and chloroplast genomes are
  circular, like bacterial chromosomes.
\item
  \textbf{70S ribosomes:} Organelle ribosomes are structurally
  bacterial-type; sensitive to antibacterials (chloramphenicol,
  erythromycin) that inhibit 70S, not eukaryotic 80S ribosomes.
\item
  \textbf{Binary fission:} Organelles divide by binary fission,
  independently of nuclear cell division, using FtsZ-related proteins.
\item
  \textbf{Phylogenetics:} rRNA sequences of mitochondria are most
  similar to \textbf{alpha-proteobacteria} (e.g., \emph{Rickettsia});
  chloroplasts to \textbf{cyanobacteria}.
\item
  \textbf{Transfer of genes:} about 1,500 mitochondrial-origin
  genes now live in the nuclear genome (gene transfer to nucleus over
  evolutionary time).
\item
  \textbf{Cardiolipin:} The inner mitochondrial membrane contains
  cardiolipin, a lipid characteristic of bacterial membranes.
\item
  \textbf{Formylmethionine:} Mitochondrial protein synthesis initiates
  with N-formylmethionine, as in bacteria, not methionine as in
  eukaryotic cytoplasmic translation.
\end{enumerate}

\subsubsection{Primary and Secondary
Endosymbiosis}\label{primary-and-secondary-endosymbiosis}

\textbf{Primary \hyperref[gl:endosymbiosis]{\textbf{endosymbiosis}}}
(mitochondrion; then chloroplast) explains the origin of eukaryotes. The
mitochondrial endosymbiosis is believed to have occurred once in
evolutionary history, as most known eukaryotes either possess
mitochondria or have remnant mitochondria-derived organelles
(hydrogenosomes, mitosomes).

\textbf{Secondary endosymbiosis} (a eukaryote with chloroplast engulfed
by another eukaryote) explains secondary chloroplasts in
dinoflagellates, euglenids, and kelp (evidenced by 3--4 bounding
membranes). \textbf{Tertiary endosymbiosis} has also occurred in some
dinoflagellates.

\begin{quote}
\textbf{Concept Check 5:} If you treated a eukaryotic cell with
chloramphenicol (an antibiotic that inhibits 70S ribosomes), which
organelle's protein synthesis would be specifically affected? Would
cytoplasmic protein synthesis be affected? Explain.
\end{quote}

\subsubsection{LUCA: The Last Universal Common
Ancestor}\label{luca-the-last-universal-common-ancestor}

If cell theory is correct that most cells descend from pre-existing
cells \citep{virchow1855}, then the lineages of bacteria, archaea, and
eukaryotes converge backward in time on a single ancestral population:
the \textbf{Last Comprehensive Common Ancestor (LUCA)}. LUCA is not a
fossil --- it is a \emph{reconstruction} assembled from comparative
molecular data. Three classes of evidence let us infer its properties.

\begin{enumerate}
\def\labelenumi{\arabic{enumi}.}
\item
  \textbf{Almost universally conserved genes (the LUCA core).} Genes
  shared by Bacteria \emph{and} Archaea --- and absent from horizontally
  transferred outliers --- are candidates for LUCA inheritance. Modern
  reconstructions (Weiss et al.~2016, \emph{Nature Microbiology})
  identify about 355 such ancient gene families. LUCA almost
  certainly possessed: a fully fledged DNA\texorpdfstring{\ensuremath{\to}}{->}RNA\texorpdfstring{\ensuremath{\to}}{->}protein machinery (DNA
  polymerase, RNA polymerase, ribosome with rRNA, about 20
  aminoacyl-tRNA synthetases), the
  \hyperref[gl:genetic-code]{\textbf{genetic code}},
  Wood--Ljungdahl-pathway-like CO\(_2\) fixation, {[}Fe--S{]} and
  {[}Ni--Fe{]} metalloenzymes, ATP synthase (rotor/stator architecture),
  and a chemiosmotic membrane.
\item
  \textbf{Geochemistry-constrained habitat.} The metalloenzyme inventory
  (Ni, Fe, Mo, Co, W in unusual valences) suggests LUCA inhabited an
  \textbf{anaerobic, hydrothermal, alkaline environment} --- most
  parsimoniously interpreted as a serpentinising hydrothermal vent
  system on the early ocean floor (about 3.8 Ga). Dependence on
  natural proton gradients across thin inorganic membranes (FeS
  micropores in vent chimneys) preceded the modern membrane-spanning ATP
  synthase, suggesting LUCA was not yet a fully autonomous cell but a
  \emph{pre-cellular} metabolic unit (Lane \& Martin 2012).
\item
  \textbf{The DNA replication anomaly.} Bacterial and archaeal DNA
  polymerases (Family A vs.~Family B) share \emph{no} sequence
  similarity in their catalytic cores, despite both performing the same
  chemistry. The simplest explanation is that LUCA used
  \textbf{RNA-genome replication} and DNA replication evolved twice
  independently \emph{after} the bacterial--archaeal split. This argues
  that LUCA was a late RNA-world organism with a stable genetic code and
  ribosomes but was still transitioning from RNA to DNA storage.
\end{enumerate}

The reconstruction has profound implications. LUCA was not the first
cell --- it was the last common one. Earlier lineages may have existed
but left no descendants. The vast biochemical diversity of modern life
(oxygenic photosynthesis, methanogenesis, eukaryotic
compartmentalisation) most evolved \emph{after} LUCA, layered onto its
conserved core.

\begin{quote}
\textbf{Concept Check 5b:} If LUCA used naturally occurring proton
gradients (from inorganic vent chemistry) rather than self-generated
ones, which modern molecular machinery must have evolved \emph{before}
free-living cells became possible? (Consider what a free-living cell
must do that a vent-bound consortium need not.)
\end{quote}

\subsubsection{Viral Replication: Lytic vs.~Lysogenic
Cycles}\label{viral-replication-lytic-vs.-lysogenic-cycles}

Viruses are not cells (no membranes of their own metabolism, no
ribosomes, no autonomous reproduction) and so are excluded from cell
theory by definition. But viruses are cellular \emph{parasites}, and
their replication strategies illuminate what cells must do to maintain
integrity. Bacteriophages --- viruses that infect bacteria --- provide
the textbook example.

\textbf{Lytic cycle (e.g., T4 bacteriophage):} 1. \textbf{Attachment:}
tail fibres bind specific surface receptors (e.g., LPS, OmpC) ---
defines host range. 2. \textbf{Penetration:} the genome is injected into
the \hyperref[gl:cytoplasm]{\textbf{cytoplasm}} through a contractile
tail; capsid remains outside. 3. \textbf{Hijacking:} within
about 2 min, host RNA polymerase transcribes ``early'' phage
genes that degrade the host chromosome (T4 gene 46/47 nucleases) and
reprogramme cell metabolism toward phage production. 4.
\textbf{Replication and assembly:} phage DNA is replicated
(about 100 copies per cell); structural proteins self-assemble
into capsids; DNA is packaged via a portal motor (one of the strongest
molecular machines known: about 57 pN of force, comparable to
muscle myosin). 5. \textbf{Lysis:} about 25 min after
infection, phage lysozyme + holin perforate the cell envelope. The cell
bursts, releasing about 100--200 progeny phages. The host is
dead.

The lytic cycle is the textbook image of viral infection: rapid, fatal,
productive. The single-step \textbf{burst size} (200) and
\textbf{latency time} (25 min) determine the kinetics of phage
epidemics.

\textbf{Lysogenic cycle (e.g., lambda bacteriophage):} 1.
\textbf{Attachment, penetration:} same as lytic. 2.
\textbf{Integration:} the phage genome (linear at injection)
circularises and integrates into the host chromosome at a specific site
(\emph{attB}) via phage integrase. The integrated genome is now a
\textbf{prophage}. 3. \textbf{Quiescence:} the prophage is replicated
passively with the host chromosome at every cell division. Lambda CI
repressor blocks lytic gene expression. The bacterium grows and divides
normally with the prophage as a stowaway. This state can persist for
thousands of generations. 4. \textbf{Induction:} stressors (UV,
mitomycin C, low nutrients) activate the host SOS response \texorpdfstring{\ensuremath{\to}}{->} RecA
cleaves CI repressor \texorpdfstring{\ensuremath{\to}}{->} lytic genes are derepressed \texorpdfstring{\ensuremath{\to}}{->} prophage excises \texorpdfstring{\ensuremath{\to}}{->}
enters lytic cycle.

The lysogenic cycle has profound consequences. \textbf{Specialised
transduction} carries host genes between bacteria when imprecise
excision packages chromosomal DNA. Many bacterial \textbf{virulence
factors are phage-encoded}: diphtheria toxin (\texorpdfstring{\ensuremath{\beta}}{beta}-phage in
\emph{Corynebacterium diphtheriae}), cholera toxin (CTX\texorpdfstring{\ensuremath{\varphi}}{phi} in \emph{Vibrio
cholerae}), Shiga toxin (Stx phage in EHEC \emph{E. coli} O157:H7),
botulinum toxin types C and D. The lysogen-to-lytic switch ---
controlled by a single repressor --- is the prototype for many
developmental switches and inspired the operon model of
\citet{jacob1961}.

\begin{quote}
\textbf{Clinical Connection: Phage Therapy in the Antibiotic-Resistance
Era} Lytic bacteriophages can target antibiotic-resistant bacteria with
exquisite specificity (one phage typically infects one species or even
strain). Compassionate-use phage therapy has cured otherwise-fatal
infections (Steffanie Strathdee, \emph{A. baumannii} sepsis, 2017;
cystic fibrosis \emph{Mycobacterium abscessus}, 2019). Phages are now in
late-stage clinical trials for chronic urinary infections and
prosthetic-joint infections. Phage cocktails --- designed to evade
host-resistance evolution by targeting multiple receptors --- are the
closest realisation of the ``magic bullet'' concept since penicillin.
\end{quote}

\begin{quote}
\textbf{Concept Check 5c:} A \emph{V. cholerae} strain loses its CTX\texorpdfstring{\ensuremath{\varphi}}{phi}
prophage. The resulting strain can colonise the human gut but causes
mild diarrhoea. Why? What does this tell you about the relationship
between cell theory's exclusion of viruses and the practical impact of
viruses on cellular biology?
\end{quote}

\begin{center}\rule{0.5\linewidth}{0.5pt}\end{center}

\subsection{The Minimal Cell:
JCVI-syn3.0}\label{the-minimal-cell-jcvi-syn3.0}

In 2016, Craig Venter's team at the J. Craig Venter Institute created
\textbf{JCVI-syn3.0}, a synthetic organism with the smallest genome of
any self-replicating cell: \textbf{473 genes} (531 kb). This was
achieved by systematically deleting genes from \emph{Mycoplasma
mycoides} until primarily essential genes remained.

Key findings: - Of the 473 genes, \textbf{149 (31.5\%) have unknown
function} --- we do not understand why they are essential - The minimal
gene set includes genes for DNA replication, transcription, translation,
membrane synthesis, and central metabolism - No genes for cell wall
synthesis, stress responses, or secondary metabolism were required under
laboratory conditions - JCVI-syn3.0 grows more slowly than its parent
organism, with a doubling time of about 3 hours

This work defines the boundary between chemistry and life and raises
profound questions: What is the minimal instruction set for a living
system? Can we understand life by building it from scratch?

\begin{quote}
\textbf{Clinical Connection: Synthetic Biology and Minimal Cells} The
minimal cell concept has direct applications in biotechnology.
Engineered minimal cells could serve as ``chassis'' organisms for
producing pharmaceuticals, biofuels, or industrial chemicals with
predictable behaviour and reduced metabolic complexity. Understanding
essential gene sets also identifies new antibiotic targets --- genes
essential for bacterial survival that have no human homologues.
\end{quote}

\begin{center}\rule{0.5\linewidth}{0.5pt}\end{center}

\subsection{Extremophile Cells}\label{extremophile-cells}

Extremophiles are organisms (primarily archaea, but also some bacteria
and eukaryotes) that thrive in environments previously considered
incompatible with life. They demonstrate the remarkable adaptability of
cellular organisation.

\subsubsection{Types of Extremophiles}\label{types-of-extremophiles}

{\def\LTcaptype{none} % do not increment counter
\begin{longtable}[]{@{}
  >{\raggedright\arraybackslash}p{(\linewidth - 6\tabcolsep) * \real{0.0800}}
  >{\raggedright\arraybackslash}p{(\linewidth - 6\tabcolsep) * \real{0.2200}}
  >{\raggedright\arraybackslash}p{(\linewidth - 6\tabcolsep) * \real{0.3200}}
  >{\raggedright\arraybackslash}p{(\linewidth - 6\tabcolsep) * \real{0.3800}}@{}}
\toprule\noalign{}
\begin{minipage}[b]{\linewidth}\raggedright
Type
\end{minipage} & \begin{minipage}[b]{\linewidth}\raggedright
Environment
\end{minipage} & \begin{minipage}[b]{\linewidth}\raggedright
Example organism
\end{minipage} & \begin{minipage}[b]{\linewidth}\raggedright
Cellular adaptation
\end{minipage} \\
\midrule\noalign{}
\endhead
\bottomrule\noalign{}
\endlastfoot
Thermophile & \textgreater60 degrees C & \emph{Thermus aquaticus}
(source of Taq polymerase) & Reverse gyrase; saturated membrane
lipids \\
Hyperthermophile & \textgreater80 degrees C & \emph{Pyrolobus fumarii}
(113 degrees C) & Tetraether monolayer membranes; thermostable
proteins \\
Psychrophile & \textless15 degrees C & \emph{Psychrobacter} &
Unsaturated membrane lipids; cold-active enzymes \\
Halophile & \textgreater2 M NaCl & \emph{Halobacterium salinarum} & KCl
accumulation; acidic surface proteins \\
Acidophile & pH \textless3 & \emph{Ferroplasma acidarmanus} &
Impermeable membranes; proton pumps \\
Alkaliphile & pH \textgreater9 & \emph{Natronobacterium} &
Na\(^+\)/H\(^+\) antiporters \\
Barophile/piezophile & \textgreater100 atm & Deep-sea archaea &
Unsaturated lipids; pressure-stable proteins \\
Radioresistant & High radiation & \emph{Deinococcus radiodurans} &
Redundant DNA repair; Mn\(^{2+}\) antioxidant \\
\end{longtable}
}

\subsubsection{Implications for
Astrobiology}\label{implications-for-astrobiology}

Extremophiles expand our understanding of habitable environments in the
solar system:

\begin{itemize}
\tightlist
\item
  \textbf{Mars:} Psychrophilic and radiation-resistant organisms suggest
  that microbial life could persist in Martian subsurface brine
  environments
\item
  \textbf{Europa (Jupiter's moon):} Halophilic and barophilic organisms
  thrive under conditions similar to those predicted for Europa's
  subsurface ocean
\item
  \textbf{Enceladus (Saturn's moon):} Hydrothermal vent-associated
  chemolithoautotrophs provide models for life powered by water-rock
  chemistry
\item
  \textbf{Titan:} While no known organisms thrive in liquid methane, the
  existence of extremophiles encourages consideration of exotic
  biochemistries
\end{itemize}

\begin{quote}
\textbf{Concept Check 6:} \emph{Ferroplasma acidarmanus} lives at pH 0
(equivalent to battery acid), yet its internal pH is about 5.6.
Calculate the proton concentration gradient across its membrane. How
many orders of magnitude difference is this?
\end{quote}

\begin{center}\rule{0.5\linewidth}{0.5pt}\end{center}

\subsection{Cell Type Classification}\label{cell-type-classification}

\subsubsection{By Carbon/Energy Source}\label{by-carbonenergy-source}

{\def\LTcaptype{none} % do not increment counter
\begin{longtable}[]{@{}
  >{\raggedright\arraybackslash}p{(\linewidth - 6\tabcolsep) * \real{0.1053}}
  >{\raggedright\arraybackslash}p{(\linewidth - 6\tabcolsep) * \real{0.3421}}
  >{\raggedright\arraybackslash}p{(\linewidth - 6\tabcolsep) * \real{0.3421}}
  >{\raggedright\arraybackslash}p{(\linewidth - 6\tabcolsep) * \real{0.2105}}@{}}
\toprule\noalign{}
\begin{minipage}[b]{\linewidth}\raggedright
Type
\end{minipage} & \begin{minipage}[b]{\linewidth}\raggedright
Carbon source
\end{minipage} & \begin{minipage}[b]{\linewidth}\raggedright
Energy source
\end{minipage} & \begin{minipage}[b]{\linewidth}\raggedright
Examples
\end{minipage} \\
\midrule\noalign{}
\endhead
\bottomrule\noalign{}
\endlastfoot
Photoautotroph & CO\(_2\) & Light & Plants, cyanobacteria, algae \\
Photoheterotroph & Organic carbon & Light & Purple non-sulphur
bacteria \\
Chemoautotroph & CO\(_2\) & Inorganic chemicals & Nitrifiers, sulphur
oxidisers \\
Chemoheterotroph & Organic carbon & Organic chemicals & Animals, fungi,
most bacteria \\
\end{longtable}
}

\subsubsection{By Oxygen Relationship}\label{by-oxygen-relationship}

\begin{itemize}
\tightlist
\item
  \textbf{Obligate aerobe:} requires O\(_2\) (complex animals,
  \emph{Mycobacterium tuberculosis})
\item
  \textbf{Facultative anaerobe:} can grow with or without O\(_2\)
  (yeast, \emph{E. coli})
\item
  \textbf{Obligate anaerobe:} killed by O\(_2\) (\emph{Clostridium
  botulinum}, methanogens)
\item
  \textbf{Aerotolerant anaerobe:} not killed by O\(_2\) but does not use
  it (\emph{Lactobacillus})
\item
  \textbf{Microaerophile:} requires O\(_2\) but at lower concentrations
  than atmospheric (\emph{Campylobacter}, \emph{Helicobacter pylori})
\end{itemize}

\begin{quote}
\textbf{Concept Check 7:} A chemoautotrophic bacterium living near a
deep-sea hydrothermal vent uses H\(_2\)S as an electron donor and
CO\(_2\) as its carbon source. It rarely encounters sunlight. Is this
organism ultimately dependent on solar energy? Explain.
\end{quote}

\begin{center}\rule{0.5\linewidth}{0.5pt}\end{center}

\subsection{Computational Bridge}\label{computational-bridge-4}

The codebase tags canonical organelle inventories by cell type ---
useful when you compare compartmentalisation across the three domains:

\begin{Shaded}
\begin{Highlighting}[]
\ImportTok{from}\NormalTok{ biology.cell }\ImportTok{import}\NormalTok{ get\_organelles\_by\_cell\_type}

\NormalTok{animal }\OperatorTok{=}\NormalTok{ get\_organelles\_by\_cell\_type(}\StringTok{"animal"}\NormalTok{)}
\NormalTok{plant }\OperatorTok{=}\NormalTok{ get\_organelles\_by\_cell\_type(}\StringTok{"plant"}\NormalTok{)}
\BuiltInTok{print}\NormalTok{(}\BuiltInTok{len}\NormalTok{(animal), }\BuiltInTok{len}\NormalTok{(plant))}
\end{Highlighting}
\end{Shaded}

\begin{quote}
\textbf{Clinical / systems note:} Antibiotic discovery still exploits
differences between prokaryotic ribosomes, cell walls, and
topoisomerases and their human counterparts --- the same structural
dichotomy this unit formalises as cell theory.
\end{quote}

\begin{center}\rule{0.5\linewidth}{0.5pt}\end{center}

\subsection{Current Evidence and Frontier
Biology}\label{current-evidence-and-frontier-biology-7}

For \textbf{Cell Theory and Cell Types}, frontier biology belongs inside
the evidence logic of the chapter. Cell biology is increasingly measured
as live, spatial, single-cell, and perturbational data rather than
static diagrams alone. The core reading question is this: cell-theory
evidence now includes microscopy, lineage tracing, omics, and
synthetic-cell boundary tests.

\begin{itemize}
\tightlist
\item
  \textbf{What to verify:} identify the observation, model, assay, or
  dataset that would make the claim stronger or weaker.
\item
  \textbf{What to qualify:} state the scale, organism, cell type,
  environmental condition, or population where the claim is expected to
  hold.
\item
  \textbf{What to compare:} test at least one alternative explanation,
  baseline, or null model before treating the pattern as causal.
\item
  \textbf{What to cite:} distinguish primary evidence, review synthesis,
  public dataset, and institutional guidance; for recent or numeric
  claims, prefer the source closest to the measurement and state what
  has changed since it was published.
\end{itemize}

Ask what measurement scale is being claimed: nanometre structure,
single-cell transcript abundance, organelle dynamics, tissue context, or
organismal phenotype.

\textbf{Source practice:} For cell-state claims, distinguish microscopy,
live-cell perturbation, single-cell sequencing, spatial transcriptomics,
and biochemical assay evidence before making a causal statement.

\subsection{Summary}\label{summary-7}

\begin{itemize}
\tightlist
\item
  Cell theory: most organisms are cells; cells are the basic unit; cells
  come from cells. Modern additions include heritable DNA information
  and shared chemical composition.
\item
  The surface-area-to-volume ratio (\(A/V = 3/r\)) constrains cell size;
  large cells must increase \(A\) or reduce metabolic demand per unit
  volume. Diffusion time scales with \(x^2\), imposing additional size
  limits.
\item
  Microscopy techniques span light (bright-field, phase-contrast, DIC,
  fluorescence, confocal, TIRF), electron (TEM, SEM, cryo-EM), and
  super-resolution (STED, STORM, PALM) methods, with resolutions from
  about 200 nm to sub-nanometre.
\item
  Prokaryotes (Bacteria + Archaea) lack membrane-bound organelles; Gram
  staining discriminates based on cell wall structure. Archaea share
  molecular features with eukaryotes.
\item
  Eukaryotic organelles carry out compartmentalised functions;
  mitochondria and chloroplasts arose by endosymbiosis of bacterial
  ancestors.
\item
  The minimal cell (JCVI-syn3.0, 473 genes) defines the boundary of
  cellular life; extremophiles demonstrate life's adaptability to
  extreme environments.
\item
  \textbf{Connections:} See \cref{sec:unit_II_cell_structure} for
  organelle inventory, Unit VII for prokaryotic diversity, and Unit VI
  for common ancestry and LUCA reasoning.
\end{itemize}

\begin{center}\rule{0.5\linewidth}{0.5pt}\end{center}

\subsection{Review Questions}\label{review-questions-7}

\begin{enumerate}
\def\labelenumi{\arabic{enumi}.}
\item
  State the three original postulates of cell theory and the three
  modern additions. For each, identify the key experiment or discovery
  that established it.
\item
  A spherical cell has radius 5 \texorpdfstring{\ensuremath{\mu}}{mu}m. Calculate its surface area, volume,
  and \(A/V\) ratio. If the cell doubles its radius, by what factor does
  \(A/V\) change?
\item
  Compare and contrast the cell walls of Gram-positive bacteria,
  Gram-negative bacteria, and archaea. Why are archaea naturally
  resistant to penicillin?
\item
  List six lines of evidence supporting the endosymbiotic origin of
  mitochondria. Which single piece of evidence do you consider most
  compelling, and why?
\item
  Explain why TIRF microscopy would be preferred over confocal
  microscopy for studying single vesicle fusion events at the plasma
  membrane.
\item
  The diffusion coefficient of glucose in water is approximately
  \(6.7 \times 10^{-10}\) m\(^2\)/s. Calculate the time required for
  glucose to diffuse across (a) a 2 \texorpdfstring{\ensuremath{\mu}}{mu}m bacterium and (b) a 100 \texorpdfstring{\ensuremath{\mu}}{mu}m plant
  cell.
\item
  JCVI-syn3.0 has 473 genes, of which 149 have unknown function. Discuss
  the implications of this finding for our understanding of the minimal
  requirements for life.
\item
  Describe how \emph{Deinococcus radiodurans} survives radiation doses
  \textgreater1,000 times the lethal dose for humans. What cellular
  mechanisms enable this extreme resistance?
\item
  Compare the membrane lipids of bacteria, archaea, and eukaryotes. How
  do archaeal ether-linked isoprenoid lipids contribute to survival at
  extreme temperatures?
\item
  A pharmaceutical company is developing a new antibiotic. Using your
  knowledge of prokaryotic cell structure, suggest three molecular
  targets that are present in bacteria but absent in human cells.
\item
  Using \breaktt{get\_organelles\_by\_cell\_type}, contrast the organelle
  lists for \texttt{"animal"} and \texttt{"prokaryote"}. Which entries
  best illustrate endosymbiotic theory?
\item
  A spherical bacterium and a spherical eukaryote both rely on
  diffusion. If the eukaryote is 20\texorpdfstring{\ensuremath{\times}}{x} larger in radius, by what factor
  does the surface-to-volume ratio change relative to the bacterium?
\item
  Using \cref{eq:unit_II_cube_law_basis} and \cref{eq:unit_II_rmax},
  explain why every metabolically active animal cell is \textless{} 100
  \texorpdfstring{\ensuremath{\mu}}{mu}m in radius unless it has internal membranes, microvilli, or partial
  anoxia. Identify two cell types that violate this rule and explain
  \emph{how} they do so.
\item
  List three bacterial cytoskeletal proteins, their eukaryotic
  homologues, and the function of each. Why was the absence of ``the
  cytoskeleton'' in prokaryotes a textbook error rather than a fact?
\item
  Compare lytic and lysogenic phage cycles in terms of their (a) speed,
  (b) burst size, (c) consequences for the host cell, and (d)
  implications for evolution. Which strategy is more reminiscent of
  ``infectious disease'' and which of ``horizontal gene transfer''?
\item
  Outline the three classes of evidence used to reconstruct LUCA. What
  does the absence of conserved DNA polymerase tell us about LUCA's
  information system?
\end{enumerate}

\begin{center}\rule{0.5\linewidth}{0.5pt}\end{center}

\subsection{Discussion Questions}\label{discussion-questions-3}

These questions have no single correct answer --- they are designed to
provoke debate, integrate concepts across sections, and connect cell
biology to broader scientific and ethical questions.

\begin{enumerate}
\def\labelenumi{\arabic{enumi}.}
\item
  \textbf{Defining life.} Is JCVI-syn3.0 alive? It self-replicates,
  contains a genome, performs metabolism, and responds to its
  environment --- but every component was designed by humans. If the
  answer is ``yes,'' does that mean we have created life? If ``no,''
  what definition excludes it that does not also exclude obligate
  intracellular pathogens like \emph{Rickettsia} (which depend on host
  metabolism just as syn3.0 depends on enriched media)?
\item
  \textbf{Cell theory under stress.} Cell theory states that most cells
  come from pre-existing cells. But somatic-cell nuclear transfer
  (cloning), induced pluripotent stem cells (iPSCs from skin
  fibroblasts), and now in-vitro-generated organoids (mini-brains,
  mini-livers) push this principle. Are organoids ``organisms''? Do they
  have moral status if they begin to exhibit coordinated electrical
  activity?
\item
  \textbf{The 149-gene mystery.} JCVI-syn3.0 has 149 essential genes of
  unknown function. Choose one approach (CRISPR screening, AlphaFold
  structural prediction, evolutionary co-occurrence analysis, or
  chemical biology) and argue why it is most likely to crack this set
  first. What ethical or biosafety considerations follow from
  understanding the minimal genome?
\item
  \textbf{Astrobiology and the biosignature problem.} If life on Mars or
  Europa derives from a separate origin, it might lack DNA, ribosomes,
  and ATP --- making it invisible to instruments designed around
  terrestrial biochemistry. How might cell theory be revised if we
  discovered non-cellular life elsewhere? Conversely, what
  \emph{cellular} features might be comprehensive across any
  biochemistry (compartmentalisation? polymer-based information storage?
  membranes?) and why?
\item
  \textbf{Endosymbiosis and modern biotechnology.} Some researchers
  propose engineering new endosymbionts: introducing nitrogen-fixing
  bacteria into plant cells to eliminate fertiliser dependence. What
  lessons from the natural mitochondrial endosymbiosis (gene transfer to
  host nucleus, metabolite exchange, intimate cell-cycle coupling)
  suggest what would be needed for such a designed symbiosis to be
  stable across generations?
\item
  \textbf{The phage-resistance arms race.} Bacteria evolve CRISPR-Cas
  defences; phages evolve anti-CRISPR proteins (Acrs); bacteria evolve
  anti-anti-CRISPR systems. This Red Queen dynamic has continued for
  about 3 billion years. What does it predict about the
  long-term success of phage therapy for antibiotic-resistant infections
  in humans? How might phage-cocktail design exploit or sidestep this
  dynamic?
\end{enumerate}

\begin{center}\rule{0.5\linewidth}{0.5pt}\end{center}

\subsection{Further Reading and Source
Notes}\label{further-reading-and-source-notes-7}

\begin{itemize}
\tightlist
\item
  Sagan (1967). On the origin of mitosing cells. \emph{Journal of
  Theoretical Biology}, 14.
\item
  Schleiden (1838). Beitr\{"a\}ge zur Phytogenesis. \emph{M\{"u\}ller's
  Archiv f\{"u\}r Anatomie, Physiologie und wissenschaftliche Medicin}.
\item
  Virchow (1855). Die Cellularpathologie. \emph{Archiv f\{"u\}r
  pathologische Anatomie und Physiologie}, 8.
\item
  Hooke (1665). \emph{Micrographia}. Royal Society.
\item
  Schwann (1839). \emph{Mikroskopische Untersuchungen \{"u\}ber die
  \{"U\}bereinstimmung in der Struktur und dem Wachsthum der Thiere und
  Pflanzen}. Gebr\{"u\}der Borntraeger.
\item
  Hershey \& Chase (1952). Independent functions of viral protein and
  nucleic acid in growth of bacteriophage. \emph{Journal of General
  Physiology}, 36.
\end{itemize}

\begin{center}\rule{0.5\linewidth}{0.5pt}\end{center}

\subsection{Key Terms}\label{key-terms-7}

{\def\LTcaptype{none} % do not increment counter
\begin{longtable}[]{@{}
  >{\raggedright\arraybackslash}p{(\linewidth - 2\tabcolsep) * \real{0.2857}}
  >{\raggedright\arraybackslash}p{(\linewidth - 2\tabcolsep) * \real{0.7143}}@{}}
\toprule\noalign{}
\begin{minipage}[b]{\linewidth}\raggedright
Term
\end{minipage} & \begin{minipage}[b]{\linewidth}\raggedright
Definition
\end{minipage} \\
\midrule\noalign{}
\endhead
\bottomrule\noalign{}
\endlastfoot
\textbf{Cell theory} & Foundational principle: most organisms are
composed of cells, cells are the basic unit of life, most cells arise
from pre-existing cells \\
\textbf{Surface-area-to-volume ratio} & The ratio \(A/V = 3/r\) for a
sphere; constrains maximum cell size \\
\textbf{Prokaryote} & Cell lacking a membrane-bound nucleus; Bacteria
and Archaea \\
\textbf{Eukaryote} & Cell with a true membrane-bound nucleus and
organelles \\
\textbf{Gram staining} & Differential stain distinguishing
thick-peptidoglycan (Gram-positive) from thin-peptidoglycan + outer
membrane (Gram-negative) bacteria \\
\textbf{Endosymbiotic theory} & Mitochondria and chloroplasts originated
from engulfed bacteria (alpha-proteobacteria and cyanobacteria,
respectively) \\
\textbf{Photoautotroph} & Organism using light energy and CO\(_2\) as
carbon source \\
\textbf{Chemoheterotroph} & Organism using chemical energy from organic
compounds \\
\textbf{Obligate anaerobe} & Organism killed by oxygen; uses
\hyperref[gl:anaerobic]{\textbf{anaerobic}} metabolism exclusively \\
\textbf{Nucleoid} & Condensed DNA-containing region of prokaryotic
cells, not bounded by a membrane \\
\textbf{Peptidoglycan} & Polymer of NAG-NAM cross-linked by peptide
bridges; structural component of bacterial cell walls \\
\textbf{Cryo-EM} & Electron microscopy of vitrified (flash-frozen)
specimens; near-atomic resolution without staining \\
\textbf{JCVI-syn3.0} & Minimal synthetic cell with 473 genes; defines
the lower boundary of self-replicating life \\
\textbf{Extremophile} & Organism thriving in extreme environments
(temperature, pH, salinity, pressure, radiation) \\
\textbf{Diffusion limit} & Physical constraint on cell size imposed by
the relationship \(t = x^2/2D\) \\
\textbf{Compartmentalisation} & Segregation of biochemical processes
into membrane-bound organelles; a defining eukaryotic feature \\
\end{longtable}
}

\begin{center}\rule{0.5\linewidth}{0.5pt}\end{center}

\subsubsection{Companion Source Module}\label{companion-source-module-7}

\textbf{Cell Theory and Cell Types} should leave a reproducible trail
from a biological claim to the code, figure, diagram, or paper-based
activity that can test it. Use the surfaces below to inspect the
chapter's assumptions, rerun the relevant model, or compare the
manuscript explanation with companion labs and figures.

{\def\LTcaptype{none} % do not increment counter
\begin{longtable}[]{@{}
  >{\raggedright\arraybackslash}p{(\linewidth - 2\tabcolsep) * \real{0.5000}}
  >{\raggedright\arraybackslash}p{(\linewidth - 2\tabcolsep) * \real{0.5000}}@{}}
\toprule\noalign{}
\begin{minipage}[b]{\linewidth}\raggedright
Surface
\end{minipage} & \begin{minipage}[b]{\linewidth}\raggedright
Use it for
\end{minipage} \\
\midrule\noalign{}
\endhead
\bottomrule\noalign{}
\endlastfoot
\breaktt{src/biology/cell/cell\_biology.py}
(\breaktt{get\_organelles\_by\_cell\_type},
\breaktt{count\_membrane\_bound\_organelles}) & Turn cell-type
comparisons into explicit feature lists rather than memorised
diagrams. \\
\breaktt{src/mermaid/biology\_diagrams.py}
(\breaktt{organelle\_function\_diagram}) & Connect cell theory to
structure-function evidence. \\
\end{longtable}
}

\textbf{Reproducibility check:} identify the observation scale, specimen
state, and imaging limit before deciding what counts as evidence for a
cellular claim. \textbf{Cross-reference:} compare with
\cref{sec:unit_II_cell_structure} and
\cref{sec:unit_VII_bacteria_archaea_viruses}.

\newpage

\section{Cell Structure and
Organelles}\label{cell-structure-and-organelles}

\label{sec:unit_II_cell_structure}

\begin{quote}
\textbf{Ch 6} \texorpdfstring{\ensuremath{\cdot}}{.} Level 2/3 \texorpdfstring{\ensuremath{\cdot}}{.} 50 min read \texorpdfstring{\ensuremath{\cdot}}{.} 75 min lecture \texorpdfstring{\ensuremath{\cdot}}{.}
Prerequisites: \cref{sec:unit_II_cell_theory}
\end{quote}

\subsection{Learning Objectives}\label{learning-objectives-8}

\begin{enumerate}
\def\labelenumi{\arabic{enumi}.}
\tightlist
\item
  Describe the structure and function of each major eukaryotic
  \hyperref[gl:organelle]{\textbf{organelle}} in detail.
\item
  Explain the endomembrane system and the secretory/endocytic pathways,
  including coat \hyperref[gl:protein]{\textbf{protein}}s and sorting
  signals.
\item
  Describe the \hyperref[gl:cytoskeleton]{\textbf{cytoskeleton}}:
  microtubules, microfilaments, and intermediate filaments, including
  their associated motor proteins.
\item
  Explain how the nucleus controls \hyperref[gl:gene]{\textbf{gene}}
  expression through \hyperref[gl:chromatin]{\textbf{chromatin}}
  organisation and the \hyperref[gl:histone-code]{\textbf{histone
  code}}.
\item
  Relate organelle dysfunction to specific human diseases.
\item
  Describe mitochondrial dynamics (fusion, fission) and their role in
  cellular \hyperref[gl:homeostasis]{\textbf{homeostasis}}.
\item
  Explain the role of peroxisomes, centrosomes, and primary cilia in
  cell biology.
\item
  Describe cell adhesion molecules and their role in tissue
  architecture.
\end{enumerate}

\subsubsection{Study Blueprint}\label{study-blueprint-8}

\begin{itemize}
\tightlist
\item
  \textbf{Big idea:} Organelle structure partitions work, traffic,
  information, and energy inside eukaryotic cells.
\item
  \textbf{Core concepts:} organelles, endomembrane system, cytoskeleton,
  motor proteins.
\item
  \textbf{Framework alignment:} Vision \& Change: Structure and
  function, Systems, Information flow, exchange, and storage; AP
  Biology: Systems Interactions, Information Storage and Transmission;
  NGSS-style topics: Structure and Function.
\item
  \textbf{Model or quantitative lens:} Compartment-flow and
  motor-transport calculations.
\item
  \textbf{Data skill:} Trace a molecule through compartments using
  evidence from labels or perturbations.
\item
  \textbf{Practice cadence:} Visual Representations, Questions and
  Methods, Argumentation.
\item
  \textbf{Common misconception to repair:} Organelles are dynamic
  process hubs, not static textbook icons.
\item
  \textbf{Primary lab:} \cref{sec:lab_unit_II_cell_structure}.
\item
  \textbf{Question bank:} \cref{sec:q_unit_II_cell_structure}.
\item
  \textbf{Transfer task:} Apply compartment logic to secretion,
  apoptosis, infection, or cell division.
\item
  \textbf{Bridge to computation:}
  \breaktt{biology.cell.cell\_biology.ORGANELLES}.
\end{itemize}

\begin{center}\rule{0.5\linewidth}{0.5pt}\end{center}

\begin{quote}
\textbf{Opening Vignette: The World's Smallest Rotary Motor}

Buried in the envelope of a bacterium is a molecular machine that spins
at 100,000 rotations per minute, driven by a stream of protons flowing
down an electrochemical gradient. The \textbf{bacterial flagellar motor}
--- known since the 1970s from electron microscopy --- is roughly 45 nm
in diameter and converts ion-motive force directly into rotational
movement, propelling the cell at up to 30 body lengths per second. No
human-engineered motor approaches this combination of speed, size, and
efficiency.

The bacterial flagellar motor is assembled from over 40 different
proteins, each with a precise structural role: the MS ring anchors to
the inner membrane; the C ring acts as a cargo dock for proteins to be
secreted; the rod and hook transmit torque; the flagellar filament acts
as a propeller. The motor can switch rotation direction in under a
millisecond, enabling the tumbling runs that allow bacteria to chemotax
toward nutrients (Berg \& Brown, 1972, \emph{Nature}). Cell structure,
as this chapter shows, is not merely a container for biochemistry --- it
is the machinery of life itself.

\emph{Primary source: Berg, H. C. \& Brown, D. A. (1972). Chemotaxis in
Escherichia coli analysed by three-dimensional tracking. Nature,
239(5374), 500--504.}
\end{quote}

\begin{center}\rule{0.5\linewidth}{0.5pt}\end{center}

\subsection{Quantitative Organelle
Inventory}\label{quantitative-organelle-inventory}

Before examining each organelle individually, it is useful to fix the
dimensions and copy numbers in a single reference table. Sizes vary
modestly across cell types; the values below are typical for a cultured
mammalian cell (HeLa-like, about 20 \texorpdfstring{\ensuremath{\mu}}{mu}m diameter).

{\def\LTcaptype{none} % do not increment counter
\begin{longtable}[]{@{}
  >{\raggedright\arraybackslash}p{(\linewidth - 8\tabcolsep) * \real{0.1268}}
  >{\raggedright\arraybackslash}p{(\linewidth - 8\tabcolsep) * \real{0.2254}}
  >{\raggedright\arraybackslash}p{(\linewidth - 8\tabcolsep) * \real{0.2113}}
  >{\raggedright\arraybackslash}p{(\linewidth - 8\tabcolsep) * \real{0.2535}}
  >{\raggedright\arraybackslash}p{(\linewidth - 8\tabcolsep) * \real{0.1831}}@{}}
\toprule\noalign{}
\begin{minipage}[b]{\linewidth}\raggedright
Organelle
\end{minipage} & \begin{minipage}[b]{\linewidth}\raggedright
Linear dimension
\end{minipage} & \begin{minipage}[b]{\linewidth}\raggedright
Number per cell
\end{minipage} & \begin{minipage}[b]{\linewidth}\raggedright
Approximate volume
\end{minipage} & \begin{minipage}[b]{\linewidth}\raggedright
Membrane area
\end{minipage} \\
\midrule\noalign{}
\endhead
\bottomrule\noalign{}
\endlastfoot
Plasma membrane & 5 nm thick; 20 \texorpdfstring{\ensuremath{\mu}}{mu}m cell diameter & 1 & --- &
about 1,300 \texorpdfstring{\ensuremath{\mu}}{mu}m\(^2\) \\
Nucleus & 5--10 \texorpdfstring{\ensuremath{\mu}}{mu}m diameter & 1 & about 600 \texorpdfstring{\ensuremath{\mu}}{mu}m\(^3\) &
about 280 \texorpdfstring{\ensuremath{\mu}}{mu}m\(^2\) (NE) \\
Nuclear pore complex (NPC) & about 120 nm outer;
about 9 nm channel & about 3,000--5,000 & --- & --- \\
Mitochondrion & 0.5--10 \texorpdfstring{\ensuremath{\mu}}{mu}m long; 0.5 \texorpdfstring{\ensuremath{\mu}}{mu}m wide & 100--2,000 &
about 400 \texorpdfstring{\ensuremath{\mu}}{mu}m\(^3\) total & OMM about 5,000 \texorpdfstring{\ensuremath{\mu}}{mu}m\(^2\);
IMM about 30,000 \texorpdfstring{\ensuremath{\mu}}{mu}m\(^2\) (cristae 6\texorpdfstring{\ensuremath{\times}}{x} amplification) \\
Mitochondrial DNA (mtDNA) & 16,569 bp circular & about 10--100
copies per mitochondrion & --- & --- \\
Rough ER & sheets/tubules about 30--100 nm & one continuous
network & about 10\% cell volume & about 13,000
\texorpdfstring{\ensuremath{\mu}}{mu}m\(^2\) \\
Smooth ER & tubules about 30--60 nm & one network (continuous
with RER) & varies; large in liver/muscle & varies \\
Golgi stack & 4--8 cisternae; about 1 \texorpdfstring{\ensuremath{\mu}}{mu}m wide & typically 1
(cis--trans); fragmented during mitosis & about 1\% cell volume
& about 1,000 \texorpdfstring{\ensuremath{\mu}}{mu}m\(^2\) \\
Lysosome & 0.1--1 \texorpdfstring{\ensuremath{\mu}}{mu}m diameter & about 50--1,000 &
about 1\% cell volume & about 100 \texorpdfstring{\ensuremath{\mu}}{mu}m\(^2\) each \\
Peroxisome & 0.1--1 \texorpdfstring{\ensuremath{\mu}}{mu}m diameter & 50--1,000 & about 1\% cell
volume & --- \\
Centriole & 200 nm long \texorpdfstring{\ensuremath{\times}}{x} 250 nm wide & 2 (one centrosome) & --- &
--- \\
Primary cilium & 1--10 \texorpdfstring{\ensuremath{\mu}}{mu}m long \texorpdfstring{\ensuremath{\times}}{x} 250 nm wide & 1 & --- & --- \\
\hyperref[gl:ribosome]{\textbf{Ribosome}} (cytosolic, 80S) &
about 25 nm & about 10\(^7\) per cell & --- & --- \\
Ribosome (mitochondrial, 55S) & about 25 nm &
about 10\(^4\) per mitochondrion & --- & --- \\
Microtubule & 25 nm wide \texorpdfstring{\ensuremath{\times}}{x} variable length & thousands & --- & --- \\
\hyperref[gl:actin]{\textbf{Actin}} filament & 7 nm \texorpdfstring{\ensuremath{\times}}{x} variable length &
about 10\(^9\) G-actin monomers & --- & --- \\
Cytoplasmic vesicle & 50--200 nm & thousands transient & --- & --- \\
Proteasome (26S) & 45 nm long & about 30,000--500,000 per cell
& --- & --- \\
\end{longtable}
}

Three numerical patterns deserve attention. (1) The internal membrane
area exceeds the plasma membrane area by about 30-fold, so the
plasma membrane is a small fraction of cellular membrane surface. (2)
Mitochondria are the largest membrane reservoir (cristae alone are 20\texorpdfstring{\ensuremath{\times}}{x}
the plasma membrane), which is why their dysfunction propagates so
quickly through bioenergetics. (3) NPCs are far rarer than ribosomes
(10\(^3\) vs.~10\(^7\)) yet handle 100\texorpdfstring{\ensuremath{\times}}{x} more transport events per second
per channel --- they are arguably the most heavily used machinery in the
cell. The chapters that follow refer back to this table so that
quantitative arguments (e.g., ``if mitophagy disposes of 10\% of
mitochondria per day\ldots{}'') have grounded numbers.

\begin{center}\rule{0.5\linewidth}{0.5pt}\end{center}

\subsection{The Nucleus --- The Information
Centre}\label{the-nucleus-the-information-centre}

The \textbf{nucleus} (diameter typically 5--10 \texorpdfstring{\ensuremath{\mu}}{mu}m) is the defining
hallmark of the eukaryotic cell.

\subsubsection{Nuclear Envelope}\label{nuclear-envelope}

The nucleus is bounded by the \textbf{nuclear envelope} --- two
concentric \hyperref[gl:phospholipid-bilayer]{\textbf{phospholipid
bilayer}}s (inner and outer nuclear membrane). The outer nuclear
membrane is continuous with the rough ER. The \textbf{nuclear pore
complex (NPC)} penetrates both membranes at about 3,000 pores
per nucleus. Each NPC is a massive protein assembly (about 120
MDa; 34 distinct nucleoporin proteins, 8-fold rotational symmetry),
gating passage of macromolecules:

\begin{itemize}
\tightlist
\item
  Small molecules (\textless40 kDa) diffuse passively through the pore
  channel (diameter about 9 nm)
\item
  Larger molecules require \textbf{nuclear localisation sequences (NLS)}
  or \textbf{nuclear export sequences (NES)} recognised by
  importin/exportin proteins
\item
  The NPC can transport cargo up to about 39 nm in diameter
  (e.g., ribosomal subunits, mRNA-protein complexes)
\item
  Transport is powered by the RanGTP/RanGDP gradient: high RanGTP in the
  nucleus, high RanGDP in the
  \hyperref[gl:cytoplasm]{\textbf{cytoplasm}}
\end{itemize}

\textbf{Import mechanism:} Cargo with NLS binds importin-alpha/beta in
the cytoplasm, translocates through the NPC, and RanGTP in the nucleus
causes importin to release its cargo.

\textbf{Export mechanism:} Cargo with NES binds exportin (CRM1) + RanGTP
in the nucleus, translocates out, and RanGAP in the cytoplasm hydrolyses
RanGTP to RanGDP, releasing cargo.

\textbf{NPC architecture in numbers.} Each NPC is built of
about 30 distinct nucleoporin (Nup) species totalling
about 500--1,000 protein subunits arranged with 8-fold
rotational symmetry. The structural scaffold (the ``Y-complex'') forms
two outer rings; cytoplasmic and nuclear filaments project
about 50 nm from each face; FG-repeat nucleoporins (Nup98,
Nup62, Nsp1) line the central channel and form a hydrogel-like phase
that confers selectivity. A single human nucleus contains
about 3,000--5,000 NPCs (density about 10--20 per
\texorpdfstring{\ensuremath{\mu}}{mu}m\(^2\) of nuclear envelope), and each NPC handles roughly
\textbf{1,000 translocation events per second} at peak activity ---
meaning the entire nucleus exchanges about 10\(^7\)
macromolecules per second. The selective barrier is so efficient that
small molecules (\textless40 kDa) cross by simple diffusion in
\textless1 ms while large cargo without NLS/NES are excluded essentially
indefinitely. The energy budget is supplied by the \textbf{RanGTP/GDP
gradient} (10\(^4\):1 nuclear:cytoplasmic ratio for RanGTP), which is
itself maintained by RanGEF (RCC1, chromatin-bound) and RanGAP
(cytoplasmic): the result is a vectorial pump powered by GTP hydrolysis
but requiring no direct translocation ATPase.

\textbf{Pathological NPCs.} The Nup98 gene is fused to dozens of
partners in pediatric leukaemias (Nup98-HOXA9, Nup98-NSD1) --- the
FG-repeat domain provides a phase-separating prion-like domain that
drives oncogenic condensate formation. Nucleoporins also accumulate
damage with age; selective loss of NPC function is a hallmark of
post-mitotic neurons in Alzheimer's and ALS.

\subsubsection{Protein Targeting and
Sorting}\label{protein-targeting-and-sorting}

Every protein synthesised in the cell must reach a specific compartment,
but proteins are made (initially) in primarily two places: free
cytosolic ribosomes and ER-bound ribosomes. \textbf{Sorting signals}
within the polypeptide direct each protein to its destination. Most
signals are short, modular, and additive --- analogous to postal
addresses appended to the same envelope.

{\def\LTcaptype{none} % do not increment counter
\begin{longtable}[]{@{}
  >{\raggedright\arraybackslash}p{(\linewidth - 6\tabcolsep) * \real{0.1897}}
  >{\raggedright\arraybackslash}p{(\linewidth - 6\tabcolsep) * \real{0.2931}}
  >{\raggedright\arraybackslash}p{(\linewidth - 6\tabcolsep) * \real{0.3621}}
  >{\raggedright\arraybackslash}p{(\linewidth - 6\tabcolsep) * \real{0.1552}}@{}}
\toprule\noalign{}
\begin{minipage}[b]{\linewidth}\raggedright
Destination
\end{minipage} & \begin{minipage}[b]{\linewidth}\raggedright
Targeting signal
\end{minipage} & \begin{minipage}[b]{\linewidth}\raggedright
Recognition machinery
\end{minipage} & \begin{minipage}[b]{\linewidth}\raggedright
Mechanism
\end{minipage} \\
\midrule\noalign{}
\endhead
\bottomrule\noalign{}
\endlastfoot
Nucleus & Nuclear localisation sequence (NLS): clusters of basic
residues (e.g., PKKKRKV in SV40 T-antigen) & Importin-alpha (recognises
NLS) + Importin-beta (NPC binding) & Co- or post-translational; powered
by Ran-GTP gradient \\
Mitochondrion (matrix) & N-terminal amphipathic alpha-helix
(presequence; about 15--70 aa) & TOM complex (outer membrane);
TIM23 complex (inner membrane); Hsp70 motor & Post-translational;
unfolded protein threaded through; presequence cleaved by MPP \\
Mitochondrion (OMM, IMM) & Internal hydrophobic / cysteine motifs & TOM,
SAM, TIM22, MIA40 (IMS) & Multiple lateral and stop-transfer pathways \\
Chloroplast (stroma) & N-terminal transit peptide & TOC and TIC
complexes & Analogous to TOM/TIM \\
Peroxisome (matrix) & PTS1 (C-terminal SKL) or PTS2 (N-terminal RLx5HL)
& PEX5 (PTS1), PEX7 (PTS2) & Folded proteins imported (unique);
mono-ubiquitinated PEX5 recycles via PEX1/6 ATPases \\
ER (secretory pathway) & N-terminal signal peptide
(about 16--30 hydrophobic aa) & SRP (signal recognition
particle) \texorpdfstring{\ensuremath{\to}}{->} SRP receptor \texorpdfstring{\ensuremath{\to}}{->} Sec61 translocon & \textbf{Co-translational}
--- translation pauses until ribosome docks \\
ER \texorpdfstring{\ensuremath{\to}}{->} Golgi (anterograde) & Diacidic/dihydrophobic motifs (DXE, FF) &
COPII coat (Sar1, Sec23/24, Sec13/31) & Vesicle budding \\
Golgi \texorpdfstring{\ensuremath{\to}}{->} ER (retrieval) & C-terminal KDEL (lumenal) or KKXX (membrane) &
KDEL receptor (KDELR); COPI coat & Retrograde retrieval of escaped
resident proteins \\
Lysosome & Mannose-6-phosphate (M6P) added in cis-Golgi & M6P receptor
(MPR) & Sorting at trans-Golgi network \\
Plasma membrane & Default pathway (no retention signal) & Constitutive
secretion & Vesicles fuse continuously \\
\end{longtable}
}

\textbf{The SRP pathway in detail.} SRP is a ribonucleoprotein particle
(54 kDa SRP54 GTPase + 6 other proteins + 7SL RNA) that solves a
three-way coordination problem: it must (i) recognise the nascent signal
peptide as it emerges from the ribosome, (ii) pause translation to
prevent premature folding in the cytoplasm, and (iii) deliver the
ribosome--nascent chain to the ER membrane. SRP54 binds the hydrophobic
signal peptide; SRP9/14 contact the ribosome's elongation factor binding
site, halting translation. The complex docks at the SRP receptor (SR,
also a GTPase) on the ER membrane. Reciprocal GTP hydrolysis by SRP54
and SR releases SRP, hands the signal peptide to the \textbf{Sec61
translocon}, and translation resumes --- now co-translationally
threading the polypeptide into the ER lumen. The ``GTPase exchange''
mechanism is deeply conserved: SRP/SR systems occur in bacteria,
archaea, and eukaryotes, with lineage-specific architectural variations.

\textbf{TOM/TIM mitochondrial import.} Mitochondrial matrix proteins
face a unique problem: they must cross \emph{two} membranes while
remaining largely unfolded. The \textbf{TOM complex} (Translocase of
Outer Membrane) --- built around Tom40 (\texorpdfstring{\ensuremath{\beta}}{beta}-barrel channel), Tom20
(presequence receptor), Tom22, and small Tom proteins --- accepts the
presequence on the cytoplasmic side. \textbf{TIM23} (Translocase of
Inner Membrane) then takes over, threaded by the \textbf{PAM motor}
(Presequence translocase-Associated Motor) containing mtHsp70 that uses
ATP hydrolysis to ratchet the polypeptide into the matrix. The inner
membrane potential (\texorpdfstring{\ensuremath{\Delta}}{Delta}\texorpdfstring{\ensuremath{\Psi}}{Psi} \texorpdfstring{\ensuremath{\approx}}{~} \texorpdfstring{\ensuremath{-}}{-}180 mV) provides additional electrostatic
pulling force on the positively charged presequence. Once in the matrix,
MPP (Mitochondrial Processing Peptidase) cleaves the presequence and the
protein folds, often with mtHsp60/Hsp10 chaperonin assistance.

OMM \texorpdfstring{\ensuremath{\beta}}{beta}-barrel proteins (e.g., porin/VDAC, Tom40 itself) take a different
route: through TOM, into the IMS, then inserted laterally by the
\textbf{SAM complex} (Sorting and Assembly Machinery). Carrier proteins
(e.g., ADP/ATP translocase) of the inner membrane skip TIM23 entirely
and use \textbf{TIM22} for lateral release into the IMM. IMS proteins
use the \textbf{MIA40/Erv1 disulphide relay} for oxidative folding.

\textbf{Why this matters.} Failures of any single targeting pathway
produce specific, severe disease phenotypes --- testimony to how
hard-coded the cellular logistics are. Defects in PEX genes cause
peroxisome biogenesis disorders (Zellweger syndrome). Defects in TIM/TOM
components cause mitochondrial myopathies. ER signal peptide mutations
are now recognised in disorders ranging from preprovasopressin (familial
diabetes insipidus) to coagulation-factor deficiencies. The cell's
postal system has no general delivery option --- every misrouted package
is a disease.

\begin{quote}
\textbf{Concept Check 1b:} A researcher creates a fusion protein with
both an N-terminal mitochondrial presequence \emph{and} a C-terminal SKL
(peroxisomal PTS1). Where is the protein delivered? (Hint: consider
whether targeting is co-translational or post-translational, and which
signal acts first.)
\end{quote}

\subsubsection{Chromatin and
Nucleosomes}\label{chromatin-and-nucleosomes}

Human nuclear DNA (2 x 3.2 billion base pairs = 6.4 Gbp total) is
packaged into chromatin. The
\textbf{\hyperref[gl:nucleosome]{nucleosome}} is the basic packaging
unit: 147 bp of DNA wound about 1.75 times around an octamer of
\hyperref[gl:histone]{\textbf{histone}} proteins (2x H2A, H2B, H3, H4)
with linker H1 bridging adjacent nucleosomes.

\textbf{Levels of chromatin compaction:}

\begin{enumerate}
\def\labelenumi{\arabic{enumi}.}
\tightlist
\item
  \textbf{10 nm fibre} (beads on a string): nucleosome arrays;
  about 6-fold compaction of naked DNA
\item
  \textbf{30 nm fibre}: solenoid or zigzag arrangement (histone H1
  stabilises); about 40-fold compaction
\item
  \textbf{300 nm loops}: attached to a protein scaffold (condensin,
  cohesin); about 10,000-fold
\item
  \textbf{700 nm condensed chromatin}: mitotic
  \hyperref[gl:chromosome]{\textbf{chromosome}}
\item
  \textbf{1,400 nm} fully condensed: metaphase chromosome;
  about 10,000--20,000-fold compaction total
\end{enumerate}

\textbf{\hyperref[gl:euchromatin]{Euchromatin}} (transcriptionally
active; less compacted)
vs.~\textbf{\hyperref[gl:heterochromatin]{heterochromatin}}
(transcriptionally silent; highly compacted). Histone modifications
(acetylation = active; methylation context-dependent) regulate this
switch --- the \textbf{histone code}.

{\def\LTcaptype{none} % do not increment counter
\begin{longtable}[]{@{}
  >{\raggedright\arraybackslash}p{(\linewidth - 8\tabcolsep) * \real{0.2353}}
  >{\raggedright\arraybackslash}p{(\linewidth - 8\tabcolsep) * \real{0.1373}}
  >{\raggedright\arraybackslash}p{(\linewidth - 8\tabcolsep) * \real{0.1176}}
  >{\raggedright\arraybackslash}p{(\linewidth - 8\tabcolsep) * \real{0.2549}}
  >{\raggedright\arraybackslash}p{(\linewidth - 8\tabcolsep) * \real{0.2549}}@{}}
\toprule\noalign{}
\begin{minipage}[b]{\linewidth}\raggedright
Modification
\end{minipage} & \begin{minipage}[b]{\linewidth}\raggedright
Residue
\end{minipage} & \begin{minipage}[b]{\linewidth}\raggedright
Effect
\end{minipage} & \begin{minipage}[b]{\linewidth}\raggedright
Writer \hyperref[gl:enzyme]{\textbf{enzyme}}
\end{minipage} & \begin{minipage}[b]{\linewidth}\raggedright
Eraser enzyme
\end{minipage} \\
\midrule\noalign{}
\endhead
\bottomrule\noalign{}
\endlastfoot
Acetylation & H3K9ac, H3K27ac & Activation & HATs (p300/CBP) & HDACs \\
Methylation & H3K4me3 & Activation
(\hyperref[gl:promoter]{\textbf{promoter}}s) & SET1/MLL & KDM5/JARID \\
Methylation & H3K9me3 & Silencing (heterochromatin) & SUV39H1 &
KDM4/JMJD2 \\
Methylation & H3K27me3 & Silencing (Polycomb) & EZH2 (PRC2) &
KDM6/UTX \\
Phosphorylation & H3S10ph & Mitotic condensation & Aurora B & PP1 \\
Ubiquitination & H2BK120ub &
\hyperref[gl:transcription]{\textbf{Transcription}} elongation &
RNF20/40 & USP22 \\
\end{longtable}
}

\subsubsection{Nucleolus}\label{nucleolus}

The \textbf{nucleolus} (1--3 per nucleus) is not membrane-bound. It is
the site of ribosomal RNA (rRNA) transcription by RNA polymerase I and
\hyperref[gl:ribosome]{\textbf{ribosome}} assembly. It disassembles
during \hyperref[gl:mitosis]{\textbf{mitosis}} and re-forms around
\textbf{nucleolus organiser regions (NORs)} on chromosomes 13, 14, 15,
21, and 22 (in humans).

The nucleolus has three ultrastructural regions: - \textbf{Fibrillar
centre (FC):} rDNA gene clusters and RNA Pol I - \textbf{Dense fibrillar
component (DFC):} nascent rRNA processing; fibrillarin -
\textbf{Granular component (GC):} maturing pre-ribosomal particles

\begin{quote}
\textbf{Clinical Connection: The Nucleolus in Cancer} Nucleolar size and
number are often increased in cancer cells due to elevated ribosome
biogenesis driven by oncogenes (Myc, mTOR). Pathologists use nucleolar
prominence (AgNOR staining) as a prognostic marker in several cancers.
The RNA Pol I inhibitor CX-5461 is in clinical trials as an anti-cancer
agent that specifically targets ribosome biogenesis.
\end{quote}

\begin{quote}
\textbf{Concept Check 1:} A protein has both an NLS and an NES. Would it
be found in the nucleus or cytoplasm? How might the cell regulate its
localisation?
\end{quote}

\begin{center}\rule{0.5\linewidth}{0.5pt}\end{center}

\subsection{Mitochondria --- Power Plants of the
Cell}\label{mitochondria-power-plants-of-the-cell}

Mitochondria (singular: mitochondrion) are 1--5 \texorpdfstring{\ensuremath{\mu}}{mu}m long, often
branching, highly dynamic organelles. A liver hepatocyte contains
about 1,000--2,000 mitochondria, constituting
about 20\% of cell volume.

\subsubsection{Structural Organisation}\label{structural-organisation}

\begin{itemize}
\tightlist
\item
  \textbf{Outer mitochondrial membrane (OMM):} contains \textbf{porin}
  (VDAC) channels; freely permeable to molecules \textless5 kDa. Also
  contains MAM (mitochondria-associated ER membrane) contact sites for
  Ca\(^{2+}\) and lipid exchange.
\item
  \textbf{Intermembrane space (IMS):} equivalent to cytoplasm in ion
  composition; site of electron carrier cytochrome c; contains
  pro-apoptotic factors (Smac/DIABLO, AIF)
\item
  \textbf{Inner mitochondrial membrane (IMM):} highly folded into
  \textbf{cristae} (10x surface area amplification);
  \textbf{impermeable} to protons --- essential for
  \hyperref[gl:chemiosmosis]{\textbf{chemiosmosis}}; contains electron
  transport chain (ETC) complexes I--IV,
  \hyperref[gl:atp-synthase]{\textbf{ATP synthase}}. The cristae are
  shaped by the MICOS complex (mitochondrial contact site and cristae
  organising system) and OPA1.
\item
  \textbf{Matrix:} aqueous interior; contains mitochondrial DNA,
  ribosomes, TCA cycle enzymes, fatty acid oxidation enzymes, pyruvate
  dehydrogenase complex
\end{itemize}

\textbf{Crista density} correlates with energetic demand: sperm
mitochondria and cardiac muscle mitochondria have densely packed
cristae. Resting cells have fewer, wider cristae.

\subsubsection{Mitochondrial DNA (mtDNA)}\label{mitochondrial-dna-mtdna}

Human mtDNA: 16,569 bp, circular, encodes 37 genes: - 13 subunits of
respiratory chain complexes (e.g., ND1-6, COX1-3, ATP6, ATP8) - 22
transfer RNAs - 2 ribosomal RNAs

mtDNA is inherited \textbf{maternally} (sperm mitochondria are typically
ubiquitinated and destroyed in the fertilised egg). This enables
maternal-lineage \textbf{phylogenetic analysis} and forensic
identification. \textbf{Mitochondrial Eve} --- the most recent common
maternal ancestor of living humans --- lived
about 150,000--200,000 years ago in Africa.

\textbf{Mitochondrial diseases} arise from mtDNA
\hyperref[gl:mutation]{\textbf{mutation}}s (point mutations or
deletions): - \textbf{Leber hereditary optic neuropathy (LHON):}
mutations in ND genes leads to retinal ganglion cell degeneration and
blindness - \textbf{MELAS} (Mitochondrial Encephalomyopathy, Lactic
Acidosis, Stroke-like episodes): A3243G mutation in tRNA-Leu -
\textbf{Kearns-Sayre syndrome:} large mtDNA deletions; progressive
external ophthalmoplegia, cardiac conduction defects, retinitis
pigmentosa

\subsubsection{Mitochondrial Dynamics: Fusion and
Fission}\label{mitochondrial-dynamics-fusion-and-fission}

Mitochondria are not static organelles --- they constantly undergo
\textbf{fusion} (joining) and \textbf{fission} (splitting), forming a
dynamic network called the mitochondrial reticulum.

\begin{figure}[htbp]
\centering
\includegraphics[width=0.82\linewidth,height=4.2in,keepaspectratio]{../figures/mermaid_inline/inline_mermaid_0034_0d2889ec94eb.png}
\caption{Mitochondrial dynamics: the balance between fusion and fission controls mitochondrial morphology, quality control, and biogenesis. Damaged mitochondria are selectively removed by mitophagy via the PINK1/Parkin pathway.}
\end{figure}

\textbf{Fusion proteins:} - \textbf{MFN1/MFN2 (mitofusins):} GTPases on
the OMM; mediate outer membrane tethering and fusion - \textbf{OPA1:}
GTPase on the IMM; mediates inner membrane fusion and cristae
remodelling - Mutations in MFN2 cause \textbf{Charcot-Marie-Tooth type
2A} (peripheral neuropathy) - Mutations in OPA1 cause
\textbf{\hyperref[gl:dominant]{dominant} optic atrophy} (most common
inherited optic neuropathy)

\textbf{Fission proteins:} - \textbf{DRP1 (dynamin-related protein 1):}
cytosolic GTPase recruited to the OMM by adaptors (MFF, MiD49/51, FIS1);
assembles into a ring and constricts the mitochondrion -
\textbf{ER-mitochondria contact sites} mark fission sites; the ER wraps
around the mitochondrion before DRP1 recruitment - Fission is essential
for equal mitochondrial distribution during cell division and for
isolating damaged mitochondria for mitophagy

\textbf{Mitophagy (PINK1/Parkin pathway):} 1. In healthy mitochondria,
PINK1 kinase is imported and degraded by PARL protease 2. In damaged
mitochondria (depolarised), PINK1 accumulates on the OMM 3. PINK1
phosphorylates ubiquitin and recruits \textbf{Parkin} (E3 ubiquitin
ligase) 4. Parkin ubiquitinates OMM proteins, creating autophagy
receptor binding sites 5. Autophagosome engulfs the damaged
mitochondrion for lysosomal degradation

\begin{quote}
\textbf{Clinical Connection: Mitochondrial Dynamics and Parkinson's
Disease} Mutations in PINK1 and Parkin cause autosomal recessive
early-onset Parkinson's disease (PARK6 and PARK2). The failure of
mitophagy leads to accumulation of damaged mitochondria in dopaminergic
\hyperref[gl:neuron]{\textbf{neuron}}s of the substantia nigra, causing
oxidative stress, energy depletion, and neuronal death. This pathway is
a major therapeutic target. see \cref{sec:unit_II_cell_signaling} for
\hyperref[gl:apoptosis]{\textbf{apoptosis}} pathways.
\end{quote}

\begin{quote}
\textbf{Concept Check 2:} Predict the mitochondrial
\hyperref[gl:phenotype]{\textbf{phenotype}} in a cell with a
dominant-negative mutation in DRP1 that prevents GTP hydrolysis. Would
the mitochondria appear fragmented or hyperfused?
\end{quote}

\begin{center}\rule{0.5\linewidth}{0.5pt}\end{center}

\subsection{The Endomembrane System and Secretory
Pathway}\label{the-endomembrane-system-and-secretory-pathway}

Several organelles form a functionally integrated \textbf{endomembrane
system}:

\begin{figure}[htbp]
\centering
\includegraphics[width=0.82\linewidth,height=4.2in,keepaspectratio]{../figures/mermaid_inline/inline_mermaid_0035_42ef176149c8.png}
\caption{The endomembrane system and secretory/endocytic pathways. COPII vesicles carry cargo anterograde (ER to Golgi); COPI vesicles carry cargo retrograde (Golgi to ER). The trans-Golgi network sorts proteins to lysosomes (via M6P signal), plasma membrane, or secretory granules.}
\end{figure}

\subsubsection{Endoplasmic Reticulum}\label{endoplasmic-reticulum}

\textbf{Rough ER} (RER): studded with ribosomes; site of synthesis,
folding, and initial N-glycosylation of secreted and membrane proteins.
Molecular chaperones (BiP/GRP78, calnexin, calreticulin) ensure proper
folding; misfolded proteins are retrotranslocated for
\textbf{ER-associated degradation (ERAD)} by the proteasome.

The \textbf{signal recognition particle (SRP)} pathway directs secretory
proteins to the ER: 1. Ribosome begins translating mRNA; signal peptide
(about 16--30 hydrophobic amino acids) emerges 2. SRP binds
signal peptide and ribosome, pausing
\hyperref[gl:translation]{\textbf{translation}} 3. SRP-ribosome complex
docks at SRP receptor on ER membrane 4. Signal peptide inserts into the
\textbf{translocon} (Sec61 channel) 5. Translation resumes; polypeptide
is co-translationally threaded into the ER lumen 6. Signal peptidase
cleaves the signal peptide

\textbf{Smooth ER} (SER): lacks ribosomes; functions: - Lipid synthesis
(phospholipids, cholesterol, steroid
\hyperref[gl:hormone]{\textbf{hormone}}s) - Drug/toxin detoxification
(cytochrome P450 enzymes in liver SER) - Ca\(^{2+}\) storage and release
for signalling (sarcoplasmic reticulum in muscle = specialised SER) -
Glycogen metabolism

\textbf{Unfolded protein response (UPR):} Accumulation of misfolded
proteins triggers the UPR --- three ER stress sensor pathways (IRE1,
PERK, ATF6) that expand ER capacity, halt translation, and if
irresolvable, trigger apoptosis. Relevant in diabetes (beta-cells
overwhelmed by insulin demand), neurodegeneration, and cancer.

\textbf{ER quality control and ERAD (in detail).} Newly synthesised
glycoproteins entering the ER lumen are tagged with a
triple-glucosylated N-glycan (Glc\(_3\)Man\(_9\)GlcNAc\(_2\)).
Glucosidases I and II sequentially trim two glucoses, generating a
monoglucosylated form that is recognised by the lectin chaperones
\textbf{calnexin} (membrane-bound) and \textbf{calreticulin} (lumenal).
These chaperones hold the protein in a folding-competent state. If the
protein folds correctly, glucosidase II removes the last glucose and the
protein exits via COPII vesicles. If folding fails, \textbf{UGGT}
(UDP-glucose:glycoprotein glucosyltransferase) re-glucosylates it,
sending it back to calnexin/calreticulin --- the
\textbf{calnexin/calreticulin cycle}. Repeated failure marks the protein
for \textbf{ERAD (ER-associated degradation)}:

\begin{enumerate}
\def\labelenumi{\arabic{enumi}.}
\tightlist
\item
  \textbf{Recognition:} Mannose trimming by ER mannosidases (EDEM1/2/3)
  creates a ``give up'' signal --- no further folding attempts.
\item
  \textbf{Retrotranslocation:} The misfolded protein is recognised by
  ERAD components (HRD1/SEL1L, gp78, MARCH6) and threaded back to the
  cytoplasm through a retrotranslocation channel (likely Hrd1 itself).
\item
  \textbf{Polyubiquitination:} ER-membrane E3 ligases (Hrd1, gp78)
  attach K48-linked ubiquitin chains.
\item
  \textbf{Extraction:} The AAA+ ATPase \textbf{p97/VCP} (Cdc48 in yeast)
  uses ATP hydrolysis to mechanically extract the ubiquitinated
  polypeptide from the membrane.
\item
  \textbf{Degradation:} The proteasome digests the substrate to small
  peptides.
\end{enumerate}

\textbf{Three branches of the UPR.} When ERAD cannot keep up with the
load, three transmembrane sensors are activated:

{\def\LTcaptype{none} % do not increment counter
\begin{longtable}[]{@{}
  >{\raggedright\arraybackslash}p{(\linewidth - 6\tabcolsep) * \real{0.1935}}
  >{\raggedright\arraybackslash}p{(\linewidth - 6\tabcolsep) * \real{0.2903}}
  >{\raggedright\arraybackslash}p{(\linewidth - 6\tabcolsep) * \real{0.1935}}
  >{\raggedright\arraybackslash}p{(\linewidth - 6\tabcolsep) * \real{0.3226}}@{}}
\toprule\noalign{}
\begin{minipage}[b]{\linewidth}\raggedright
Sensor
\end{minipage} & \begin{minipage}[b]{\linewidth}\raggedright
Mechanism
\end{minipage} & \begin{minipage}[b]{\linewidth}\raggedright
Output
\end{minipage} & \begin{minipage}[b]{\linewidth}\raggedright
Time scale
\end{minipage} \\
\midrule\noalign{}
\endhead
\bottomrule\noalign{}
\endlastfoot
\textbf{IRE1} & Lumenal domain dimerises on misfolded proteins;
cytoplasmic domain has kinase + endoribonuclease activity & Splices XBP1
mRNA \texorpdfstring{\ensuremath{\to}}{->} active XBP1s transcription factor \texorpdfstring{\ensuremath{\to}}{->} genes for ER expansion, ERAD,
lipid synthesis & Minutes--hours (chronic adaptation) \\
\textbf{PERK} & Dimerises and trans-phosphorylates eIF2alpha & Global
translation attenuation; selective translation of ATF4 \texorpdfstring{\ensuremath{\to}}{->} CHOP & Minutes
(immediate response) \\
\textbf{ATF6} & Trafficks to Golgi; cleaved by S1P + S2P proteases &
Cytoplasmic fragment is a transcription factor for ER chaperones (BiP,
GRP94) & Hours \\
\end{longtable}
}

If the lesion is repaired, the UPR shuts down and homeostasis resumes.
If the stress is severe or prolonged, \textbf{CHOP} (downstream of
PERK/ATF4) and IRE1's secondary RIDD activity (Regulated IRE1-Dependent
Decay of mRNAs) bias the cell toward apoptosis. The same pathway is
therefore an adaptive thermostat \emph{and} a death timer --- the switch
is set by the integrated stress dose.

\textbf{Worked observation: insulin secretion and ER load.} A pancreatic
beta-cell synthesises about 10\(^6\) insulin molecules per
minute during glucose stimulation. Each must be folded with three
disulphide bonds, processed by PC1/3 and PC2 from proinsulin, and
packaged into secretory granules. ER folding capacity is therefore one
of the rate-limiting steps for insulin secretion. Chronic hyperglycaemia
(in obesity, type 2 diabetes) drives sustained UPR, eventual
CHOP-mediated apoptosis, and progressive beta-cell loss. Pharmacological
chemical chaperones (TUDCA, 4-PBA) and selective IRE1 inhibitors are in
trials to delay this collapse.

\begin{quote}
\textbf{Clinical Connection: ER Stress and Type 2 Diabetes} Pancreatic
beta-cells produce enormous quantities of insulin (about 1
million molecules per cell per minute during glucose stimulation). This
places extreme demands on the ER folding machinery. Chronic
hyperglycaemia and obesity increase insulin demand beyond ER capacity,
triggering the UPR. Prolonged UPR activation causes beta-cell apoptosis,
contributing to the progressive decline in insulin secretion seen in
type 2 diabetes. Pharmacological chaperones (e.g., TUDCA, 4-PBA) that
reduce ER stress are being investigated as diabetes therapeutics.
\end{quote}

\subsubsection{Golgi Apparatus}\label{golgi-apparatus}

The Golgi stack (cis-Golgi network, cis, medial, trans, trans-Golgi
network/TGN) receives cargo from the ER in \textbf{COPII-coated
vesicles} and processes proteins by: - O-glycosylation (adding sugars to
Ser/Thr) - Proteolytic processing (e.g., insulin prohormone to insulin)
- Phosphorylation of mannose-6-phosphate (M6P) residues, creating the
lysosomal targeting signal - Sorting to: secretory vesicles, plasma
membrane, lysosomes, or ER (retrograde COPI vesicles)

\textbf{Vesicle coat proteins and their roles:}

{\def\LTcaptype{none} % do not increment counter
\begin{longtable}[]{@{}
  >{\raggedright\arraybackslash}p{(\linewidth - 4\tabcolsep) * \real{0.4000}}
  >{\raggedright\arraybackslash}p{(\linewidth - 4\tabcolsep) * \real{0.4333}}
  >{\raggedright\arraybackslash}p{(\linewidth - 4\tabcolsep) * \real{0.1667}}@{}}
\toprule\noalign{}
\begin{minipage}[b]{\linewidth}\raggedright
Coat protein
\end{minipage} & \begin{minipage}[b]{\linewidth}\raggedright
Vesicle route
\end{minipage} & \begin{minipage}[b]{\linewidth}\raggedright
Cargo
\end{minipage} \\
\midrule\noalign{}
\endhead
\bottomrule\noalign{}
\endlastfoot
COPII (Sec13/31, Sec23/24) & ER to cis-Golgi & Newly synthesised
secretory proteins \\
COPI (coatomer) & Golgi to ER (retrograde) & ER-resident proteins (KDEL
retrieval) \\
Clathrin + adaptors (AP1, AP2) & TGN to endosome; PM to endosome &
Lysosomal enzymes (M6P); receptor-mediated
\hyperref[gl:endocytosis]{\textbf{endocytosis}} \\
\end{longtable}
}

\textbf{Brefeldin A} (BFA) inhibits GBF1 (a GEF for Arf1 GTPase needed
for COPI assembly) and dissolves the Golgi within minutes ---
demonstrating its dynamic membrane flux.

\subsubsection{Lysosomes}\label{lysosomes}

Lysosomes (\hyperref[gl:ph]{\textbf{pH}} 4.5--5.0, maintained by V-type
H\(^+\)-ATPase) contain about 60 acid hydrolases (proteases,
lipases, nucleases, glycosidases) that digest: - \textbf{Autophagy:}
cellular debris and aged organelles (macroautophagy, microautophagy,
chaperone-mediated autophagy) - \textbf{Phagocytosis:} bacteria and
debris (in macrophages, neutrophils) - \textbf{Receptor-mediated
endocytosis:} LDL-cholesterol uptake - \textbf{Extracellular digestion:}
osteoclasts secrete lysosomal contents to resorb bone

\textbf{Types of selective autophagy:} - \textbf{Mitophagy:} selective
removal of damaged mitochondria (PINK1/Parkin pathway) -
\textbf{Pexophagy:} selective removal of excess peroxisomes -
\textbf{Lipophagy:} selective degradation of lipid droplets -
\textbf{Ribophagy:} selective degradation of ribosomes during starvation
- \textbf{ER-phagy/reticulophagy:} selective removal of excess ER

\textbf{Lysosomal storage diseases} result from enzyme deficiencies: -
\textbf{Pompe disease:} alpha-1,4-glucosidase deficiency, glycogen
accumulates, cardiac/muscle failure - \textbf{Gaucher disease:}
glucocerebrosidase deficiency, sphingolipid accumulation in macrophages
- \textbf{Tay-Sachs disease:} hexosaminidase A deficiency, GM2
ganglioside accumulation in neurons, neurodegeneration -
\textbf{Niemann-Pick type C:} NPC1/NPC2 cholesterol transport
deficiency, cholesterol accumulates in lysosomes - Both Pompe and
Gaucher are treatable by \textbf{enzyme replacement therapy (ERT)},
where recombinant enzyme is infused intravenously and targeted to
lysosomes via M6P receptors

\begin{quote}
\textbf{Concept Check 3:} A patient has I-cell disease (mucolipidosis
II), in which the enzyme that adds M6P to lysosomal enzymes in the
cis-Golgi is defective. Predict the consequences for (a) lysosomal
enzyme targeting, (b) extracellular enzyme levels, and (c) intracellular
digestion.
\end{quote}

\begin{center}\rule{0.5\linewidth}{0.5pt}\end{center}

\subsection{The Cytoskeleton}\label{the-cytoskeleton}

The cytoskeleton is a dynamic protein network pervading the cytoplasm,
providing structural support, enabling cell movement, and directing
intracellular traffic.

\begin{figure}[htbp]
\centering
\includegraphics[width=0.82\linewidth,height=4.2in,keepaspectratio]{../figures/mermaid_inline/inline_mermaid_0036_9b04ee870582.png}
\caption{Cytoskeletal components and their associated motor proteins. Actin filaments are driven by myosin motors, microtubules by kinesin and dynein motors, and intermediate filaments have no associated motors.}
\end{figure}

\subsubsection{Actin Microfilaments}\label{actin-microfilaments}

Actin filaments (F-actin; diameter about 7 nm) are polar,
ATP-driven polymers of globular G-actin monomers. Properties: -
\textbf{Treadmilling:} net polymerisation at the (+) barbed end;
depolymerisation at the (-) pointed end, resulting in net movement of
the filament while maintaining constant length - \textbf{Branching:}
Arp2/3 complex creates about 70-degree branches (lamellipodia,
phagocytic cups); nucleated by WASP/WAVE family activators -
\textbf{Bundling:} Fimbrin (parallel tight bundles in microvilli),
alpha-actinin (looser bundles in stress fibres), villin (intestinal
brush border)

\textbf{Critical concentration and treadmilling --- the quantitative
basis.} A single actin monomer can add to (or dissociate from) either
end of a filament. At each end, the rate constants for assembly
(\(k_+\)) and disassembly (\(k_-\)) define an end-specific
\textbf{critical concentration}:

\begin{equation}
C_c^{\text{end}} = \frac{k_-}{k_+}
\label{eq:unit_II_critical_conc}
\end{equation}

When the free monomer concentration \([G\text{-actin}]\) exceeds
\(C_c\), the end grows; below \(C_c\), the end shrinks. Crucially, the
(+) barbed end and the (\texorpdfstring{\ensuremath{-}}{-}) pointed end have \emph{different} critical
concentrations because of their different geometries:

\[C_c^{(+)} \approx 0.1 \, \mu\text{M}, \qquad C_c^{(-)} \approx 0.6 \, \mu\text{M} \tag{6.2} \label{eq:unit_II_cell_structure_item_1}\]

In a steady state with \([G\text{-actin}]\) between these two values
(say 0.3 \texorpdfstring{\ensuremath{\mu}}{mu}M), the (+) end grows continuously while the (\texorpdfstring{\ensuremath{-}}{-}) end shrinks
continuously --- and the filament \emph{translates} through space at a
steady velocity even though its average length is constant. This is
\textbf{treadmilling}, and the velocity is:

\[v_\text{tread} = \delta \cdot k_+^{(+)} \cdot ([G] - C_c^{(+)}) \tag{6.3} \label{eq:unit_II_cell_structure_item_2}\]

where \(\delta = 2.7\) nm (axial rise per monomer; two monomers per
helical repeat = 5.4 nm). For lamellipodial actin in a migrating
fibroblast, treadmilling proceeds at about 0.1 \texorpdfstring{\ensuremath{\mu}}{mu}m/s at 37 °C
--- the molecular speed limit on which crawling cells move.

ATP hydrolysis on the filament shifts the (\texorpdfstring{\ensuremath{-}}{-})-end disassembly rate,
increasing the asymmetry between ends; cofilin and ADF
(actin-depolymerising factor) accelerate (\texorpdfstring{\ensuremath{-}}{-})-end disassembly; profilin
loads ATP onto G-actin to enrich the (+)-end-favoured pool. The net
result is that \textbf{treadmilling is energetically driven by ATP
hydrolysis} even though the polymerisation reaction itself does not
consume ATP.

\textbf{Microtubule dynamic instability --- the GTP cap.} Microtubules
treadmill very weakly; instead they exhibit \textbf{dynamic instability}
--- alternating phases of growth and rapid shrinkage at a single end.
The (+) end of a growing microtubule carries a ``GTP cap'' (a few
hundred GTP-tubulin subunits added before hydrolysis catches up). When
the cap is intact, polymerisation continues at about 1 \texorpdfstring{\ensuremath{\mu}}{mu}m/min.
Stochastic loss of the cap exposes GDP-tubulin, which has a strained
lateral lattice and undergoes rapid catastrophic depolymerisation at
about 10--30 \texorpdfstring{\ensuremath{\mu}}{mu}m/min (``\textbf{catastrophe}''). Catastrophic
shrinkage can be rescued by reformation of a GTP cap
(``\textbf{rescue}''). The four parameters --- growth rate \(v_g\),
shrinkage rate \(v_s\), catastrophe frequency \(f_c\), rescue frequency
\(f_r\) --- completely specify microtubule behaviour and are tuned by a
panoply of plus-end tracking proteins (+TIPs: EB1, CLIP-170) and
depolymerases (kinesin-13/MCAK).

\textbf{Motor proteins on actin:} - \textbf{Myosin II:} generates
contractile force (muscle sarcomere, cytokinesis, cell migration);
bipolar thick filaments; step size about 5--25 nm;
non-processive (releases after each step) - \textbf{Myosin V:}
processive motor for vesicle transport; walks ``hand-over-hand'' with 36
nm steps matching the F-actin helical repeat; transports melanosomes,
ER, secretory vesicles - \textbf{Myosin I:} single-headed; links
membrane to actin cortex; involved in endocytosis and membrane tension

Functions: cell shape, cytokinesis (contractile ring), muscle
contraction, phagocytosis, cell crawling (lamellipodia, filopodia),
intracellular organelle movement.

\subsubsection{Tubulin Microtubules}\label{tubulin-microtubules}

Microtubules (MTs; 25 nm diameter) are hollow tubes of
alpha/beta-tubulin heterodimers forming 13 protofilaments. Properties: -
\textbf{GTP-driven dynamic instability:} rapid switching between growth
(rescue) and catastrophic depolymerisation; GTP cap model --- growing
end has GTP-tubulin; loss of cap triggers depolymerisation -
\textbf{Polarity:} (+) end grows toward cell periphery; (-) end at
\textbf{MTOC/centrosome} - \textbf{Post-translational modifications:}
acetylation (stable MTs), tyrosination/detyrosination, polyglutamylation
--- the ``tubulin code'' directs motor protein binding

\textbf{Motor proteins on microtubules:} - \textbf{Kinesin-1
(conventional kinesin):} anterograde (+ end directed); dimeric,
processive; 8 nm step size matching tubulin repeat; transports vesicles,
mitochondria, mRNA toward cell periphery; speed about 0.8 \texorpdfstring{\ensuremath{\mu}}{mu}m/s
- \textbf{Cytoplasmic dynein:} retrograde (- end directed); large
complex (about 1.2 MDa) with dynactin cofactor; 8 nm steps;
transports ER, Golgi, endosomes, viruses toward cell centre; also drives
mitotic spindle positioning - \textbf{Kinesin-13 (MCAK):} depolymerises
MT ends; critical for spindle dynamics during mitosis

Functions: chromosome segregation (mitotic spindle), flagella/cilia
(axoneme: 9+2 MT arrangement), intracellular vesicle transport (axonal
transport in neurons), cell polarity, cell shape.

\subsubsection{Worked Example: Motor Protein Kinematics and
Energetics}\label{worked-example-motor-protein-kinematics-and-energetics}

\textbf{Problem:} A kinesin-1 motor protein is transporting a
neurotransmitter vesicle down a nerve axon along a microtubule. The
motor moves at a constant velocity of \(v = 0.80 \, \mu\text{m/s}\) and
takes fixed step sizes of \(d = 8 \, \text{nm}\) per step. 1. How many
steps does the kinesin take per second? 2. If kinesin hydrolyses exactly
one ATP molecule per 8 nm step, how many ATP molecules are consumed
during a \(2.4 \, \text{mm}\) transport journey along the axon?

\textbf{Solution:}

\begin{enumerate}
\def\labelenumi{\arabic{enumi}.}
\item
  \textbf{Calculate the stepping rate:} First, convert velocity to
  nanometres per second to match the step size units:
  \[ v = 0.80 \, \mu\text{m/s} = 800 \, \text{nm/s}  \label{eq:unit_II_cell_structure_item_3}\]

  Calculate the number of steps per second:
  \[ \text{Step rate} = \frac{v}{d} = \frac{800 \, \text{nm/s}}{8 \, \text{nm/step}} = 100 \text{ steps/s}  \label{eq:unit_II_cell_structure_item_4}\]

  The motor takes exactly \textbf{100 steps per second}.
\item
  \textbf{Calculate the energy consumption for the journey:} Convert the
  total distance from millimetres to nanometres:
  \[ D_{\text{total}} = 2.4 \, \text{mm} = 2.4 \times 10^6 \, \text{nm} = 2,400,000 \, \text{nm}  \label{eq:unit_II_cell_structure_item_5}\]

  Calculate the total number of steps:
  \[ \text{Total steps} = \frac{D_{\text{total}}}{d} = \frac{2,400,000 \, \text{nm}}{8 \, \text{nm/step}} = 300,000 \text{ steps}  \label{eq:unit_II_cell_structure_item_6}\]

  Since 1 step = 1 ATP molecule hydrolysed:
  \[ \text{Total ATP consumed} = 300,000 \text{ molecules of ATP}  \label{eq:unit_II_cell_structure_item_7}\]
\end{enumerate}

This calculation highlights the immense energy requirements of axonal
transport; maintaining trillions of synapses across the nervous system
commands a massive proportion of the body's total ATP production.

\begin{enumerate}
\def\labelenumi{\arabic{enumi}.}
\setcounter{enumi}{2}
\item
  \textbf{Calculate the transport time:} The total transport time at
  constant velocity is simply distance divided by velocity:
  \[ t = \frac{D_\text{total}}{v} = \frac{2.4 \times 10^{-3} \, \text{m}}{0.80 \times 10^{-6} \, \text{m/s}} = 3{,}000 \, \text{s} \approx 50 \, \text{min}  \label{eq:unit_II_cell_structure_item_8}\]

  So a kinesin-driven cargo crosses 2.4 mm of axon in just under an
  hour. For a giraffe motor neuron (axon length about 3 m), the
  same calculation yields about 43 days --- consistent with the
  historical observation by Weiss and Hiscoe (1948) that materials
  accumulate above an axonal ligature with a wave-front velocity of a
  few millimetres per day. Compare this with \textbf{passive diffusion}
  of the same vesicle: using \(t = x^2/2D\) with \(D \approx 10^{-12}\)
  m\(^2\)/s for a 100 nm vesicle, \(t \approx 1.4 \times 10^{12}\) s \texorpdfstring{\ensuremath{\approx}}{~}
  45,000 years. Active transport beats diffusion by \textbf{eight orders
  of magnitude} for cargoes of this size and distance. The cell's choice
  to spend about 10\(^5\) ATPs per vesicle is therefore not
  extravagance --- it is what makes the long axon viable at most.
\end{enumerate}

\textbf{Worked Example: Comparing kinesin and dynein transport.}

A dendritic spine 50 \texorpdfstring{\ensuremath{\mu}}{mu}m from the cell body needs both anterograde
(kinesin: 0.8 \texorpdfstring{\ensuremath{\mu}}{mu}m/s) and retrograde (dynein: 1.0 \texorpdfstring{\ensuremath{\mu}}{mu}m/s) cargoes. Calculate
the round-trip time and minimum ATP cost for a single cargo cycle.

\begin{itemize}
\tightlist
\item
  Anterograde transit:
  \(t_+ = 50 \, \mu\text{m} / 0.8 \, \mu\text{m/s} = 62.5\) s
\item
  Retrograde transit:
  \(t_- = 50 \, \mu\text{m} / 1.0 \, \mu\text{m/s} = 50\) s
\item
  Total transit time: \(\sim 112\) s (cargo handling and motor switching
  add another about 10 s)
\item
  Kinesin steps: \(50{,}000 \, \text{nm} / 8 \, \text{nm} = 6{,}250\)
  steps \texorpdfstring{\ensuremath{\to}}{->} 6,250 ATP
\item
  Dynein steps (more variable, average about 8 nm): another
  about 6,250 ATP
\item
  Total: about 12,500 ATP per round trip per cargo
\end{itemize}

A single dendrite that supports 100 active spines and turns over 10
cargoes per spine per minute therefore consumes about 10\(^7\)
ATPs/min on motor stepping alone --- roughly 1\% of a neuron's resting
ATP budget devoted to logistics.

\textbf{Taxol} (paclitaxel) stabilises MT and prevents depolymerisation,
arresting mitosis in metaphase. Used in chemotherapy for breast,
ovarian, and lung cancers. It is derived from the bark of the Pacific
yew tree (\emph{Taxus brevifolia}).

\textbf{Colchicine} destabilises MT by binding free tubulin and
preventing polymerisation. Used for gout (inhibits neutrophil migration)
and in cytogenetics (arrests cells in metaphase for karyotyping).

\textbf{Vincristine/vinblastine} (vinca alkaloids from
\emph{Catharanthus roseus}) also destabilise MT. Used in cancer
chemotherapy (lymphomas, leukaemias).

\subsubsection{Intermediate Filaments}\label{intermediate-filaments}

Intermediate filaments (IFs; about 10 nm diameter) are ropelike
cables of coiled-coil proteins. Unlike actin and MT, IFs are not polar
and have no associated motors --- they are \textbf{purely structural}:

{\def\LTcaptype{none} % do not increment counter
\begin{longtable}[]{@{}
  >{\raggedright\arraybackslash}p{(\linewidth - 6\tabcolsep) * \real{0.1750}}
  >{\raggedright\arraybackslash}p{(\linewidth - 6\tabcolsep) * \real{0.1750}}
  >{\raggedright\arraybackslash}p{(\linewidth - 6\tabcolsep) * \real{0.2000}}
  >{\raggedright\arraybackslash}p{(\linewidth - 6\tabcolsep) * \real{0.4500}}@{}}
\toprule\noalign{}
\begin{minipage}[b]{\linewidth}\raggedright
IF type
\end{minipage} & \begin{minipage}[b]{\linewidth}\raggedright
Protein
\end{minipage} & \begin{minipage}[b]{\linewidth}\raggedright
Location
\end{minipage} & \begin{minipage}[b]{\linewidth}\raggedright
Associated disease
\end{minipage} \\
\midrule\noalign{}
\endhead
\bottomrule\noalign{}
\endlastfoot
Keratin (types I/II) & Keratin 1--20 & Epithelial cells & Epidermolysis
bullosa simplex \\
Vimentin (type III) & Vimentin & Fibroblasts, endothelial & Diagnostic
marker for sarcomas \\
Desmin (type III) & Desmin & Muscle (Z-disc) & Desminopathy
(cardiomyopathy) \\
GFAP (type III) & GFAP & Astrocytes & Alexander disease
(leukodystrophy) \\
Neurofilament (type IV) & NF-L, NF-M, NF-H & Axons (determines axon
diameter) & ALS, Charcot-Marie-Tooth \\
Lamins (type V) & Lamin A/B/C & Nuclear lamina (inner NE) & Progeria,
Emery-Dreifuss muscular dystrophy \\
\end{longtable}
}

\textbf{Lamins} form a meshwork underlying the inner nuclear membrane,
giving the nucleus mechanical rigidity and organising chromatin.
Progerin (mutant lamin A with 50-aa deletion due to a point mutation
activating a cryptic splice site) causes \textbf{Hutchinson-Gilford
progeria syndrome} --- accelerated ageing with death typically by age 13
from cardiovascular disease.

\begin{quote}
\textbf{Clinical Connection: Cytoskeleton Diseases} \textbf{Alexander
disease} (GFAP mutations): Gain-of-function mutations in GFAP cause
Rosenthal fibre accumulation in astrocytes, leading to progressive
leukodystrophy with megalencephaly, seizures, and developmental
regression. \textbf{ALS (amyotrophic lateral sclerosis):} Some familial
ALS cases involve mutations in dynactin (DCTN1) or dynein, disrupting
retrograde axonal transport and causing motor neuron degeneration.
\textbf{Taxol mechanism in cancer therapy:} By hyperstabilising
microtubules, taxol prevents the dynamic instability required for
mitotic spindle function. Chromosomes cannot properly segregate,
activating the spindle assembly checkpoint and ultimately triggering
apoptosis in rapidly dividing tumour cells.
\end{quote}

\begin{quote}
\textbf{Concept Check 4:} Kinesin-1 moves toward the (+) end of
microtubules, while dynein moves toward the (-) end. In a neuron, which
direction does each motor move cargo along the axon? Which motor would
transport newly synthesised synaptic vesicle precursors from the cell
body to the synapse?
\end{quote}

\begin{center}\rule{0.5\linewidth}{0.5pt}\end{center}

\subsection{Centrosomes, Centrioles, and
Cilia}\label{centrosomes-centrioles-and-cilia}

\subsubsection{Centrosome/MTOC}\label{centrosomemtoc}

The \textbf{centrosome} is the primary microtubule organising centre
(MTOC) in animal cells. It consists of two \textbf{centrioles}
(barrel-shaped structures of 9 triplet microtubules arranged in a
pinwheel) surrounded by \textbf{pericentriolar material (PCM)}
containing gamma-tubulin ring complexes (gamma-TuRC) that nucleate new
microtubules.

During mitosis, centrosomes duplicate and migrate to opposite poles of
the cell, organising the mitotic spindle. Centrosome amplification
(\textgreater2 per cell) is common in cancer and contributes to
chromosomal instability.

\subsubsection{Cilia and Flagella}\label{cilia-and-flagella}

\textbf{Motile cilia} (9+2 axoneme: 9 outer doublet MTs + 2 central
singlet MTs, linked by dynein arms and nexin bridges) beat in
coordinated waves. Found on respiratory epithelium (about 200
per cell), fallopian tube epithelium, and ependymal cells lining brain
ventricles.

\textbf{Primary cilia} (9+0 axoneme: no central pair, non-motile) are
mechanosensory/chemosensory antennae present on nearly most mammalian
cells. Function as signalling hubs: - \textbf{Hedgehog signalling:}
Smoothened receptor relocates to the primary cilium upon Hedgehog ligand
binding - \textbf{Polycystin-1/2:} mechanosensitive Ca\(^{2+}\) channels
on kidney primary cilia; sense urine flow

\textbf{Ciliopathies} --- diseases caused by defective cilia: -
\textbf{Primary ciliary dyskinesia (Kartagener syndrome):} dynein arm
defects; immotile cilia; chronic respiratory infections, situs inversus
(50\%), male infertility - \textbf{Polycystic kidney disease (PKD):}
mutations in PKD1 (polycystin-1) or PKD2 (polycystin-2); defective
primary cilia mechanosensing; uncontrolled tubular cell proliferation
and cyst formation - \textbf{Bardet-Biedl syndrome:} defective ciliary
trafficking; obesity, retinitis pigmentosa, renal anomalies, polydactyly

\begin{quote}
\textbf{Concept Check 5:} Why does Kartagener syndrome sometimes cause
situs inversus (mirror-reversal of organ asymmetry)? Consider the role
of motile nodal cilia in establishing left-right body asymmetry during
embryonic development.
\end{quote}

\begin{center}\rule{0.5\linewidth}{0.5pt}\end{center}

\subsection{Peroxisomes}\label{peroxisomes}

Peroxisomes contain oxidases that generate hydrogen peroxide
(H\(_2\)O\(_2\)) as a byproduct of fatty acid beta-oxidation, and
\textbf{catalase} that destroys it:

\[2\text{H}_2\text{O}_2 \rightarrow 2\text{H}_2\text{O} + \text{O}_2 \quad (\text{catalase}) \tag{6.1} \label{eq:unit_II_cell_structure_item_9}\]

\subsubsection{Key Functions}\label{key-functions}

\begin{itemize}
\tightlist
\item
  \textbf{Very-long-chain fatty acid (VLCFA) beta-oxidation:} Shortens
  VLCFAs (\textgreater C22) to medium-chain products that are then
  transferred to mitochondria for complete oxidation
\item
  \textbf{Plasmalogen biosynthesis:} Ether-linked phospholipids
  essential for myelin sheaths; about 50\% of heart
  phospholipids are plasmalogens
\item
  \textbf{Bile acid synthesis:} Side chain oxidation of cholesterol
  intermediates
\item
  \textbf{Amino acid oxidation:} D-amino acid oxidase
\item
  \textbf{Reactive oxygen species management:} Both generation
  (H\(_2\)O\(_2\) from oxidases) and detoxification (catalase,
  peroxidase)
\item
  \textbf{Glyoxylate metabolism:} Alanine-glyoxylate aminotransferase
  (AGT) converts glyoxylate to glycine
\end{itemize}

\subsubsection{Peroxisome Biogenesis}\label{peroxisome-biogenesis}

Peroxisomes form by \textbf{growth and division} of existing peroxisomes
(similar to mitochondria) and also by \textbf{de novo formation} from
ER-derived vesicles. Peroxisomal matrix proteins are imported
post-translationally via \textbf{PTS1} (C-terminal SKL tripeptide) or
\textbf{PTS2} (N-terminal nonapeptide) signals, recognised by PEX5 and
PEX7 receptors.

\begin{quote}
\textbf{Clinical Connection: Peroxisome Biogenesis Disorders}
\textbf{Zellweger syndrome} (the most severe peroxisome biogenesis
disorder) results from mutations in PEX genes (most commonly PEX1).
Absent or non-functional peroxisomes lead to VLCFA accumulation,
plasmalogen deficiency, and bile acid synthesis failure. Affected
infants present with severe hypotonia, seizures, hepatomegaly, and
characteristic facial features, with death typically in the first year.
\textbf{X-linked adrenoleukodystrophy (X-ALD):} Mutation in ABCD1 (a
peroxisomal ABC transporter for VLCFA); VLCFA accumulation destroys
adrenal cortex and CNS myelin. Depicted in the film ``Lorenzo's Oil.''
\textbf{Primary hyperoxaluria type 1:} AGT mistargeted from peroxisomes
to mitochondria; oxalate accumulates; kidney stones and renal failure.
see \cref{sec:unit_II_membrane_transport} for ABC transporters.
\end{quote}

\begin{center}\rule{0.5\linewidth}{0.5pt}\end{center}

\subsection{Cell Adhesion and
Junctions}\label{cell-adhesion-and-junctions}

Multicellular organisms require cells to adhere to each other and to the
extracellular matrix (ECM). Four major families of cell adhesion
molecules (CAMs) and several junction types mediate this:

\subsubsection{Cell Adhesion Molecules}\label{cell-adhesion-molecules}

{\def\LTcaptype{none} % do not increment counter
\begin{longtable}[]{@{}
  >{\raggedright\arraybackslash}p{(\linewidth - 8\tabcolsep) * \real{0.1754}}
  >{\raggedright\arraybackslash}p{(\linewidth - 8\tabcolsep) * \real{0.1228}}
  >{\raggedright\arraybackslash}p{(\linewidth - 8\tabcolsep) * \real{0.3509}}
  >{\raggedright\arraybackslash}p{(\linewidth - 8\tabcolsep) * \real{0.2105}}
  >{\raggedright\arraybackslash}p{(\linewidth - 8\tabcolsep) * \real{0.1404}}@{}}
\toprule\noalign{}
\begin{minipage}[b]{\linewidth}\raggedright
CAM family
\end{minipage} & \begin{minipage}[b]{\linewidth}\raggedright
Binding
\end{minipage} & \begin{minipage}[b]{\linewidth}\raggedright
Ca\(^{2+}\) dependent?
\end{minipage} & \begin{minipage}[b]{\linewidth}\raggedright
Key examples
\end{minipage} & \begin{minipage}[b]{\linewidth}\raggedright
Function
\end{minipage} \\
\midrule\noalign{}
\endhead
\bottomrule\noalign{}
\endlastfoot
Cadherins & Homophilic & Yes & E-cadherin (epithelial), N-cadherin
(neural), VE-cadherin (endothelial) & Tissue-specific sorting; adherens
junctions \\
Integrins & Heterophilic (ECM ligands) & Yes (divalent cation) &
alpha5-beta1 (fibronectin), alphaV-beta3 (vitronectin), alpha2-beta1
(collagen) & Cell-ECM adhesion; focal adhesions; bidirectional
signalling \\
Selectins & Heterophilic (sugar ligands) & Yes & L-selectin
(leukocytes), P-selectin (platelets, endothelium), E-selectin
(endothelium) & Leukocyte rolling and homing \\
IgCAMs & Homo/heterophilic & No & NCAM, ICAM-1, VCAM-1 & Neural
development; immune cell adhesion \\
\end{longtable}
}

\subsubsection{Cell Junctions}\label{cell-junctions}

\begin{itemize}
\tightlist
\item
  \textbf{Tight junctions (zonula occludens):} Seal adjacent epithelial
  cells; prevent paracellular diffusion; claudins and occludin form the
  seal; ZO-1/2/3 link to actin cytoskeleton. Create apical-basal
  polarity.
\item
  \textbf{Adherens junctions (zonula adherens):} Cadherin-mediated;
  linked to actin cytoskeleton via catenins (alpha, beta, p120).
  E-cadherin loss is a hallmark of epithelial-mesenchymal transition
  (EMT) in cancer metastasis.
\item
  \textbf{Desmosomes (macula adherens):} Cadherin-mediated (desmogleins,
  desmocollins); linked to intermediate filaments (keratin) via
  plakoglobin and desmoplakin. Provide mechanical strength to skin and
  heart. Autoantibodies against desmoglein-3 cause \textbf{pemphigus
  vulgaris} (life-threatening blistering disease).
\item
  \textbf{Gap junctions:} Connexin hexamers (connexons) form channels
  between adjacent cells; allow passage of ions and small molecules
  (\textless1 kDa); enable electrical coupling (cardiac muscle
  synchrony), metabolic coupling, and intercellular signalling.
  Mutations in connexin-26 (GJB2) are the most common cause of
  hereditary deafness.
\item
  \textbf{Hemidesmosomes:} Integrin alpha6-beta4 links epithelial basal
  surface to basement membrane; connected to keratin intermediate
  filaments. Defects cause epidermolysis bullosa (skin blistering from
  minor trauma).
\end{itemize}

\begin{quote}
\textbf{Concept Check 6:} E-cadherin is frequently downregulated in
invasive carcinomas. How would loss of E-cadherin promote cancer
metastasis? Consider both adhesion and signalling functions.
\end{quote}

\begin{center}\rule{0.5\linewidth}{0.5pt}\end{center}

\subsection{The Proteasome --- Protein Quality
Control}\label{the-proteasome-protein-quality-control}

While lysosomes degrade proteins within membrane-bound compartments, the
\textbf{ubiquitin-proteasome system (UPS)} degrades misfolded, damaged,
or regulatory proteins in the cytoplasm and nucleus.

\subsubsection{The 26S Proteasome}\label{the-26s-proteasome}

The 26S proteasome is a about 2.5 MDa multi-subunit protease
complex consisting of: - \textbf{20S core particle:} barrel-shaped
structure of four stacked heptameric rings (alpha7-beta7-beta7-alpha7);
proteolytic activity (chymotrypsin-like, trypsin-like, and
\hyperref[gl:caspase]{\textbf{caspase}}-like) is sequestered within the
barrel interior - \textbf{19S regulatory particle (cap):} recognises
ubiquitinated substrates, deubiquitinates them, unfolds the polypeptide,
and threads it into the 20S core for degradation

\subsubsection{Ubiquitination --- The Degradation
Signal}\label{ubiquitination-the-degradation-signal}

Ubiquitin (76 amino acids, highly conserved) is conjugated to target
proteins via a three-enzyme cascade:

\begin{enumerate}
\def\labelenumi{\arabic{enumi}.}
\tightlist
\item
  \textbf{E1 (ubiquitin-activating enzyme):} Activates ubiquitin in an
  ATP-dependent reaction (2 E1 enzymes in humans)
\item
  \textbf{E2 (ubiquitin-conjugating enzyme):} Carries activated
  ubiquitin (about 40 E2s in humans)
\item
  \textbf{E3 (ubiquitin ligase):} Confers substrate specificity;
  transfers ubiquitin to the target protein (about 600 E3s in
  humans)
\end{enumerate}

A chain of at least 4 ubiquitin molecules linked via K48 serves as the
proteasomal degradation signal. Other ubiquitin chain types (K63,
M1/linear) serve non-degradative signalling functions (NF-kB activation,
DNA repair, endocytosis).

\textbf{Key E3 ligases and their substrates:}

{\def\LTcaptype{none} % do not increment counter
\begin{longtable}[]{@{}
  >{\raggedright\arraybackslash}p{(\linewidth - 4\tabcolsep) * \real{0.3462}}
  >{\raggedright\arraybackslash}p{(\linewidth - 4\tabcolsep) * \real{0.3462}}
  >{\raggedright\arraybackslash}p{(\linewidth - 4\tabcolsep) * \real{0.3077}}@{}}
\toprule\noalign{}
\begin{minipage}[b]{\linewidth}\raggedright
E3 Ligase
\end{minipage} & \begin{minipage}[b]{\linewidth}\raggedright
Substrate
\end{minipage} & \begin{minipage}[b]{\linewidth}\raggedright
Function
\end{minipage} \\
\midrule\noalign{}
\endhead
\bottomrule\noalign{}
\endlastfoot
MDM2 & p53 & Keeps p53 levels low in unstressed cells \\
APC/C (\hyperref[gl:anaphase]{\textbf{anaphase}}-promoting complex) &
Securin, cyclin B & Drives mitotic exit \\
SCF\(^{beta-TrCP}\) & Beta-catenin, IkB & Wnt signalling, NF-kB
regulation \\
VHL (von Hippel-Lindau) & HIF-1-alpha & Oxygen sensing; mutations cause
renal cell carcinoma \\
Parkin & OMM proteins & Mitophagy through the PINK1-Parkin
quality-control pathway \\
\end{longtable}
}

\begin{quote}
\textbf{Clinical Connection: Proteasome Inhibitors in Cancer}
\textbf{Bortezomib} (Velcade) inhibits the chymotrypsin-like activity of
the 20S proteasome. This causes accumulation of misfolded proteins, ER
stress, and apoptosis. Multiple myeloma cells are particularly sensitive
because they produce enormous quantities of immunoglobulins, generating
high proteasomal load. Bortezomib has transformed multiple myeloma
treatment. Carfilzomib and ixazomib are second-generation proteasome
inhibitors with improved profiles.
\end{quote}

\begin{center}\rule{0.5\linewidth}{0.5pt}\end{center}

\subsection{CRISPR Applications in Organelle
Biology}\label{crispr-applications-in-organelle-biology}

The CRISPR-Cas9 system has revolutionised the study of organelle biology
by enabling precise gene editing:

\begin{itemize}
\tightlist
\item
  \textbf{Fluorescent tagging of endogenous proteins:} CRISPR knock-in
  of GFP/mCherry at native loci enables live imaging of organelles at
  endogenous expression levels (avoiding overexpression artifacts)
\item
  \textbf{\hyperref[gl:genome]{Genome}-wide screens:} CRISPR knockout
  libraries have identified new genes required for mitophagy, autophagy,
  ER-phagy, and peroxisome biogenesis
\item
  \textbf{Base editing of mtDNA:} DddA-derived cytosine base editors
  (DdCBEs) can edit mtDNA without double-strand breaks, enabling
  creation of precise mitochondrial disease models
\item
  \textbf{Organelle-targeted degradation:} AID (auxin-inducible degron)
  tagged organelle proteins enable acute, reversible organelle
  disruption
\end{itemize}

\begin{center}\rule{0.5\linewidth}{0.5pt}\end{center}

\subsection{Computational Bridge}\label{computational-bridge-5}

Membrane-bound organelle counts are explicit in the cell model and can
be queried for teaching comparisons:

\begin{Shaded}
\begin{Highlighting}[]
\ImportTok{from}\NormalTok{ biology.cell }\ImportTok{import}\NormalTok{ get\_organelles\_by\_cell\_type, count\_membrane\_bound\_organelles}

\NormalTok{org }\OperatorTok{=}\NormalTok{ get\_organelles\_by\_cell\_type(}\StringTok{"animal"}\NormalTok{)}
\BuiltInTok{print}\NormalTok{(count\_membrane\_bound\_organelles(org))}
\end{Highlighting}
\end{Shaded}

\begin{quote}
\textbf{Clinical / systems note:} Mitochondrial disorders and
peroxisomal diseases (e.g.~Zellweger spectrum) show how loss of a single
organelle compartment produces multi-system phenotypes because
compartment-specific metabolites fail to reach the cytosol or nucleus.
\end{quote}

\begin{center}\rule{0.5\linewidth}{0.5pt}\end{center}

\subsubsection{DddA-Derived Base Editors: Editing the Mitochondrial
Genome Without a Double-Strand
Break}\label{ddda-derived-base-editors-editing-the-mitochondrial-genome-without-a-double-strand-break}

Mitochondrial DNA (mtDNA) encodes 13 proteins of the electron transport
chain and 22 tRNAs --- mutations here cause a spectrum of
\textbf{mitochondrial diseases} (MELAS, MERRF, Leber's hereditary optic
neuropathy) that affect roughly 1 in 5000 births. CRISPR--Cas9 could not
reach mtDNA because guide RNA import into mitochondria is inefficient,
and any DNA double-strand break in the polyploid (100--10 000 copies per
cell) mitochondrial genome triggers rapid linear-DNA degradation rather
than useful repair.

The breakthrough came from an unusual source: a \textbf{bacterial
interbacterial toxin (DddA)} from \emph{Burkholderia cenocepacia} that
deaminates cytosines in \emph{double-stranded} DNA (most cytidine
deaminases require single-stranded substrates). In 2020, the Liu lab
published \textbf{DddA-derived cytosine base editors (DdCBEs)} --- split
DddA halves fused to programmable TALE arrays that reconstitute
deaminase activity primarily when both halves bind adjacent
mitochondrial sites (\emph{Nature} 2020). Subsequent engineering
produced \textbf{zinc-finger--DdCBEs} (smaller, easier to import),
\textbf{mitoBEs} capable of A-to-G edits, and \textbf{TALED}
(TALE-linked deaminases) for A-to-I conversion. Efficiencies in human
cells now reach 30--50 \% heteroplasmy shift with bystander editing
\textless{} 5 \%.

Worked example: for the MELAS-causing mutation m.3243A\textgreater G in
tRNA\textsuperscript{Leu}, a DdCBE targeting the reverse strand near
position 3244 converts the pathogenic G back to A in about 30
\% of mtDNA copies --- above the threshold needed to cross the
\textbf{heteroplasmy clinical phenotype boundary} (about 80 \%
mutant copies for overt disease). A 2022 report corrected a murine
Ndufa10 mutation \emph{in vivo} via AAV-delivered DdCBE with functional
rescue of respiratory chain complex I. Cautions: off-target
mitochondrial editing rates of about 0.1--1 \% at 30--50
predicted sites; potential nuclear off-targets; the edit is heritable in
the female germline (with the ethical considerations of heritable genome
editing). The technology is a clean example of how \textbf{unusual
biochemistry in obscure organisms} (a bacterial toxin no textbook
mentioned in 2019) can become the keystone of a new therapeutic platform
within 4 years.

\begin{center}\rule{0.5\linewidth}{0.5pt}\end{center}

\subsection{Current Evidence and Frontier
Biology}\label{current-evidence-and-frontier-biology-8}

For \textbf{Cell Structure and Organelles}, frontier biology belongs
inside the evidence logic of the chapter. Cell biology is increasingly
measured as live, spatial, single-cell, and perturbational data rather
than static diagrams alone. The core reading question is this: organelle
function is dynamic, contact-mediated, and context-dependent rather than
a fixed list of compartments.

\begin{itemize}
\tightlist
\item
  \textbf{What to verify:} identify the observation, model, assay, or
  dataset that would make the claim stronger or weaker.
\item
  \textbf{What to qualify:} state the scale, organism, cell type,
  environmental condition, or population where the claim is expected to
  hold.
\item
  \textbf{What to compare:} test at least one alternative explanation,
  baseline, or null model before treating the pattern as causal.
\item
  \textbf{What to cite:} distinguish primary evidence, review synthesis,
  public dataset, and institutional guidance; for recent or numeric
  claims, prefer the source closest to the measurement and state what
  has changed since it was published.
\end{itemize}

Ask what measurement scale is being claimed: nanometre structure,
single-cell transcript abundance, organelle dynamics, tissue context, or
organismal phenotype.

\textbf{Source practice:} For cell-state claims, distinguish microscopy,
live-cell perturbation, single-cell sequencing, spatial transcriptomics,
and biochemical assay evidence before making a causal statement.

\subsubsection{Current Evidence Map: Membrane-Bound and Condensate
Organization}\label{current-evidence-map-membrane-bound-and-condensate-organization}

\begin{figure}[htbp]
\centering
\includegraphics[width=0.82\linewidth,height=4.2in,keepaspectratio]{../figures/mermaid_inline/inline_mermaid_0037_ce4977fe6a6c.png}
\caption{Condensates should be taught as regulatable cellular organization, not as a replacement for membrane-bound organelles or as proof of causality by appearance alone.}
\end{figure}

\subsection{Summary}\label{summary-8}

\begin{itemize}
\tightlist
\item
  The nucleus houses chromatin (DNA + histones); nuclear pores regulate
  macromolecule traffic via the Ran GTPase system; the nucleolus
  assembles ribosomes and is a cancer prognostic marker.
\item
  Mitochondria have their own circular DNA and ribosomes, consistent
  with endosymbiotic origin; cristae maximise IMM surface area for the
  ETC; fusion/fission dynamics maintain mitochondrial quality via
  PINK1/Parkin-mediated mitophagy.
\item
  The endomembrane system (RER, SER, Golgi, vesicles) coordinates
  protein secretion and glycosylation; COPII (anterograde), COPI
  (retrograde), and clathrin (endocytic) coat proteins direct vesicle
  traffic.
\item
  The cytoskeleton (actin, MTs, IFs) provides structure, enables
  motility, and directs intracellular transport via motor proteins
  (myosin, kinesin, dynein).
\item
  Peroxisomes perform VLCFA oxidation and ROS management;
  centrosomes/cilia function as MTOCs and sensory antennae; cell
  adhesion molecules and junctions maintain tissue architecture.
\item
  \textbf{Connections:} See \cref{sec:unit_II_cell_theory} for cell
  theory and \hyperref[gl:endosymbiosis]{\textbf{endosymbiosis}},
  \cref{sec:unit_II_membrane_transport} for membrane traffic, and Unit
  IX for cilia-linked physiology.
\end{itemize}

\begin{center}\rule{0.5\linewidth}{0.5pt}\end{center}

\subsection{Review Questions}\label{review-questions-8}

\begin{enumerate}
\def\labelenumi{\arabic{enumi}.}
\item
  Describe the mechanism by which a newly synthesised lysosomal enzyme
  is sorted from the trans-Golgi network to the lysosome. What would
  happen if the M6P phosphotransferase were defective?
\item
  Compare and contrast the three types of cytoskeletal filaments (actin,
  microtubules, intermediate filaments) in terms of diameter, polarity,
  \hyperref[gl:nucleotide]{\textbf{nucleotide}} dependence, associated
  motors, and dynamic behaviour.
\item
  Explain the PINK1/Parkin mitophagy pathway. Why do mutations in these
  genes specifically affect dopaminergic neurons in Parkinson's disease?
\item
  A patient presents with recurrent respiratory infections, male
  infertility, and situs inversus. What is the likely diagnosis, and
  what is the molecular basis?
\item
  Describe the unfolded protein response (UPR) and explain its three
  branches (IRE1, PERK, ATF6). Under what circumstances does the UPR
  switch from adaptive to apoptotic?
\item
  Explain how the nuclear pore complex achieves selective transport of
  macromolecules while allowing free diffusion of small molecules.
\item
  Compare the mechanisms by which taxol and colchicine affect
  microtubule dynamics. Both are used clinically --- for what different
  purposes?
\item
  Explain the role of primary cilia as signalling platforms. Why do
  PKD1/PKD2 mutations in primary cilia lead to polycystic kidney
  disease?
\item
  Describe the vesicle coat proteins COPII, COPI, and clathrin. For
  each, specify the transport route, the GTPase involved, and the type
  of cargo carried.
\item
  A cell biologist discovers a new organelle-associated protein.
  Describe a CRISPR-based experimental strategy to determine its
  function and subcellular localisation.
\item
  Using the organelle catalogue, explain why disrupting
  \textbf{mitochondrial} fusion/fission dynamics preferentially affects
  tissues with high ATP turnover.
\item
  Compare one structural difference between COPII-coated vesicles and
  clathrin-coated pits that explains their distinct cargo size ranges.
\end{enumerate}

\begin{center}\rule{0.5\linewidth}{0.5pt}\end{center}

\subsection{Further Reading and Source
Notes}\label{further-reading-and-source-notes-8}

\begin{itemize}
\tightlist
\item
  Sagan (1967). On the origin of mitosing cells. \emph{Journal of
  Theoretical Biology}, 14.
\item
  de Duve (1969). The lysosome in retrospect. \emph{Lysosomes in Biology
  and Pathology}, North-Holland.
\item
  Palade (1975). Intracellular aspects of the process of protein
  synthesis. \emph{Science}, 189.
\item
  Alberts et al.~(latest ed.). \emph{Molecular Biology of the Cell}
  (chapters on organelles and the endomembrane system). Garland Science.
\item
  Margulis (1981). \emph{Symbiosis in Cell Evolution}. W. H. Freeman.
\item
  Lane \& Martin (2010). The energetics of genome complexity.
  \emph{Nature}, 467.
\end{itemize}

\begin{center}\rule{0.5\linewidth}{0.5pt}\end{center}

\subsection{Key Terms}\label{key-terms-8}

{\def\LTcaptype{none} % do not increment counter
\begin{longtable}[]{@{}
  >{\raggedright\arraybackslash}p{(\linewidth - 2\tabcolsep) * \real{0.2857}}
  >{\raggedright\arraybackslash}p{(\linewidth - 2\tabcolsep) * \real{0.7143}}@{}}
\toprule\noalign{}
\begin{minipage}[b]{\linewidth}\raggedright
Term
\end{minipage} & \begin{minipage}[b]{\linewidth}\raggedright
Definition
\end{minipage} \\
\midrule\noalign{}
\endhead
\bottomrule\noalign{}
\endlastfoot
\textbf{Nuclear pore complex} & Massive protein assembly
(about 120 MDa) gating macromolecule transport between nucleus
and cytoplasm \\
\textbf{Nucleosome} & Basic chromatin packaging unit: 147 bp DNA wound
around histone octamer \\
\textbf{Histone code} & Combinatorial histone modifications that
regulate chromatin state and gene expression \\
\textbf{Cristae} & Folds of the inner mitochondrial membrane that
increase surface area for
\hyperref[gl:oxidative-phosphorylation]{\textbf{oxidative
phosphorylation}} \\
\textbf{ERAD} & ER-associated degradation; retrotranslocation of
misfolded ER proteins for proteasomal destruction \\
\textbf{Unfolded protein response} & Three-branch ER stress pathway
(IRE1, PERK, ATF6) that restores proteostasis or triggers apoptosis \\
\textbf{Treadmilling} & Actin filament behaviour: net polymerisation at
(+) end, depolymerisation at (-) end \\
\textbf{Dynamic instability} & Microtubule property: stochastic
switching between growth and rapid depolymerisation (GTP cap model) \\
\textbf{Kinesin} & Plus-end directed microtubule motor protein;
anterograde transport \\
\textbf{Dynein} & Minus-end directed microtubule motor protein;
retrograde transport \\
\textbf{Lamins} & Type V intermediate filament proteins forming the
nuclear lamina; mutations cause laminopathies \\
\textbf{Mitophagy} & Selective autophagy of damaged mitochondria;
mediated by PINK1/Parkin pathway \\
\textbf{COPII} & Coat protein complex mediating anterograde ER-to-Golgi
vesicle transport \\
\textbf{Primary cilium} & Non-motile 9+0 sensory cilium present on most
mammalian cells; signalling antenna \\
\textbf{Ciliopathy} & Disease caused by defective cilia structure or
function (e.g., PKD, Kartagener syndrome) \\
\textbf{Connexin} & Protein forming gap junction channels; allows
intercellular communication \\
\textbf{E-cadherin} & Calcium-dependent homophilic adhesion molecule;
loss promotes cancer metastasis \\
\end{longtable}
}

\begin{center}\rule{0.5\linewidth}{0.5pt}\end{center}

\subsubsection{Companion Source Module}\label{companion-source-module-8}

\textbf{Cell Structure and Organelles} should leave a reproducible trail
from a biological claim to the code, figure, diagram, or paper-based
activity that can test it. Use the surfaces below to inspect the
chapter's assumptions, rerun the relevant model, or compare the
manuscript explanation with companion labs and figures.

{\def\LTcaptype{none} % do not increment counter
\begin{longtable}[]{@{}
  >{\raggedright\arraybackslash}p{(\linewidth - 2\tabcolsep) * \real{0.5000}}
  >{\raggedright\arraybackslash}p{(\linewidth - 2\tabcolsep) * \real{0.5000}}@{}}
\toprule\noalign{}
\begin{minipage}[b]{\linewidth}\raggedright
Surface
\end{minipage} & \begin{minipage}[b]{\linewidth}\raggedright
Use it for
\end{minipage} \\
\midrule\noalign{}
\endhead
\bottomrule\noalign{}
\endlastfoot
\breaktt{src/biology/cell/cell\_biology.py} (\texttt{Organelle},
\breaktt{get\_organelles\_by\_cell\_type},
\breaktt{count\_membrane\_bound\_organelles}) & Connect organelle
inventories to cell type and function. \\
\breaktt{src/mermaid/biology\_diagrams.py}
(\breaktt{organelle\_function\_diagram},
\breaktt{membrane\_transport\_diagram}) & Keep compartment diagrams tied
to transport and interaction. \\
\end{longtable}
}

\textbf{Reproducibility check:} treat an organelle claim as conditional
on cell type, developmental state, and measurement method.
\textbf{Cross-reference:} use \cref{sec:unit_II_cell_theory},
\cref{sec:unit_II_membrane_transport}, and
\cref{sec:unit_III_bioenergetics_and_respiration}.

\newpage

\section{Membrane Structure and
Transport}\label{membrane-structure-and-transport}

\label{sec:unit_II_membrane_transport}

\begin{quote}
\textbf{Ch 7} \texorpdfstring{\ensuremath{\cdot}}{.} Level 2/3 \texorpdfstring{\ensuremath{\cdot}}{.} 50 min read \texorpdfstring{\ensuremath{\cdot}}{.} 75 min lecture \texorpdfstring{\ensuremath{\cdot}}{.}
Prerequisites: \cref{sec:unit_II_cell_structure},
\cref{sec:unit_I_water_and_life}
\end{quote}

\subsection{Learning Objectives}\label{learning-objectives-9}

\begin{enumerate}
\def\labelenumi{\arabic{enumi}.}
\tightlist
\item
  Describe the fluid mosaic model of the plasma membrane and recent
  updates from single-molecule studies.
\item
  Classify membrane \hyperref[gl:protein]{\textbf{protein}}s by their
  topology and function.
\item
  Distinguish passive from active transport and provide mechanistic
  explanations.
\item
  Derive the Nernst equation and apply it to membrane potential
  calculations.
\item
  Explain the Goldman-Hodgkin-Katz equation for multi-ion membrane
  potential.
\item
  Describe ion channel gating mechanisms and selectivity filters.
\item
  Explain the Na\(^+\)/K\(^+\)-ATPase cycle and its physiological
  significance.
\item
  Describe vesicular trafficking including SNARE-mediated membrane
  fusion.
\item
  Compare \hyperref[gl:endocytosis]{\textbf{endocytosis}} types:
  phagocytosis, macropinocytosis, clathrin-mediated, and caveolar.
\end{enumerate}

\subsubsection{Study Blueprint}\label{study-blueprint-9}

\begin{itemize}
\tightlist
\item
  \textbf{Big idea:} Membranes convert gradients, permeability, and
  selective transport into cellular physiology.
\item
  \textbf{Core concepts:} bilayers, diffusion, osmosis, electrochemical
  gradients.
\item
  \textbf{Framework alignment:} Vision \& Change: Structure and
  function, Systems, Information flow, exchange, and storage; AP
  Biology: Systems Interactions, Information Storage and Transmission;
  NGSS-style topics: Structure and Function.
\item
  \textbf{Model or quantitative lens:} Nernst, Goldman, osmotic, and
  facilitated-transport calculations.
\item
  \textbf{Data skill:} Interpret transport data from gradients, rates,
  and membrane potentials.
\item
  \textbf{Practice cadence:} Visual Representations, Questions and
  Methods, Argumentation.
\item
  \textbf{Common misconception to repair:} Equilibrium does not mean
  equal concentration; charge and permeability matter.
\item
  \textbf{Primary lab:} \cref{sec:lab_unit_II_membrane_transport}.
\item
  \textbf{Question bank:} \cref{sec:q_unit_II_membrane_transport}.
\item
  \textbf{Transfer task:} Transfer gradient logic to neurons, kidneys,
  roots, and mitochondrial membranes.
\item
  \textbf{Bridge to computation:}
  \breaktt{biology.cell.cell\_biology.nernst\_potential}.
\end{itemize}

\begin{figure}[htbp]
\centering
\includegraphics[width=0.85\textwidth]{../figures/ghk_permeability.png}
\caption{Goldman--Hodgkin--Katz permeability sensitivity. Membrane potential is plotted as relative sodium permeability increases while potassium permeability is held fixed and chloride permeability is varied.}
\label{fig:unit_II_ghk_permeability}
\end{figure}

\begin{center}\rule{0.5\linewidth}{0.5pt}\end{center}

\begin{quote}
\textbf{Opening Vignette: The Membrane Pump That Makes Cancer Cells
Immortal}

In 1976, July Ling and colleagues at MIT discovered that some tumour
cells were extraordinarily difficult to kill with chemotherapy drugs
because they expressed unusual amounts of a 170 kDa membrane protein ---
later named \textbf{P-glycoprotein} (P-gp) --- that used ATP hydrolysis
to actively export cytotoxic drugs from the cell interior (Juliano \&
Ling, 1976, \emph{Biochimica et Biophysica Acta}). A cancer cell
overexpressing P-gp can reduce the intracellular concentration of
doxorubicin, vincristine, paclitaxel, and dozens of other chemotherapy
agents to sub-lethal levels, rendering the entire drug arsenal
ineffective. This phenomenon --- \textbf{multidrug resistance (MDR)} ---
remains one of the greatest challenges in oncology.

P-gp is an ABC transporter: it has two
\hyperref[gl:nucleotide]{\textbf{nucleotide}}-binding domains (NBDs)
that hydrolyse ATP and two transmembrane domains (TMDs) that form the
drug-export channel. It exemplifies the principle that membranes are not
passive barriers --- they are dynamic, information-processing systems
that control exactly which molecules enter and exit cells, and at what
rate. Understanding that selectivity, and exploiting it therapeutically,
is the central theme of this chapter.

\emph{Primary source: Juliano, R. L. \& Ling, V. (1976). A surface
glycoprotein modulating drug permeability in Chinese hamster ovary cell
mutants. Biochimica et Biophysica Acta, 455(1), 152--162.}
\end{quote}

\begin{center}\rule{0.5\linewidth}{0.5pt}\end{center}

In 1972, Singer and Nicolson proposed the \textbf{fluid mosaic model}:
the plasma membrane is a lipid bilayer in which proteins float like
icebergs in a sea of lipids. Both the lipid and protein components can
diffuse laterally within the membrane plane (\textbf{fluid} component),
while the asymmetric distribution of lipids and proteins across the two
leaflets generates \textbf{mosaic} heterogeneity.

\subsubsection{Membrane Lipid
Composition}\label{membrane-lipid-composition}

The plasma membrane is not a homogeneous bilayer but contains distinct
regions of varying composition:

{\def\LTcaptype{none} % do not increment counter
\begin{longtable}[]{@{}
  >{\raggedright\arraybackslash}p{(\linewidth - 4\tabcolsep) * \real{0.2115}}
  >{\raggedright\arraybackslash}p{(\linewidth - 4\tabcolsep) * \real{0.4231}}
  >{\raggedright\arraybackslash}p{(\linewidth - 4\tabcolsep) * \real{0.3654}}@{}}
\toprule\noalign{}
\begin{minipage}[b]{\linewidth}\raggedright
Lipid Class
\end{minipage} & \begin{minipage}[b]{\linewidth}\raggedright
\% of membrane (animal)
\end{minipage} & \begin{minipage}[b]{\linewidth}\raggedright
Location preference
\end{minipage} \\
\midrule\noalign{}
\endhead
\bottomrule\noalign{}
\endlastfoot
Phosphatidylcholine (PC) & 25--30 & Outer leaflet \\
Phosphatidylethanolamine (PE) & 25 & Inner leaflet \\
Phosphatidylserine (PS) & 10 & Inner leaflet (negatively charged) \\
Sphingomyelin (SM) & 15 & Outer leaflet (rafts) \\
Cholesterol & 20--50 & Both leaflets (condenses rafts) \\
Phosphatidylinositol (PI) & 5 & Inner leaflet (signalling: PIP\(_2\),
PIP\(_3\)) \\
\end{longtable}
}

\textbf{Lipid asymmetry} is maintained by \textbf{flippases}
(ATP-dependent; P4-ATPases; move PS and PE to the inner leaflet) and
\textbf{floppases} (ATP-dependent; ABC transporters; move lipids to the
outer leaflet). \textbf{Scramblases} (Ca\(^{2+}\)-activated; TMEM16F)
randomise asymmetry during \hyperref[gl:apoptosis]{\textbf{apoptosis}}
--- externalised phosphatidylserine (PS) is the ``eat me'' signal
recognised by phagocyte receptors (TIM-4, BAI-1, Stabilin-2).

\textbf{Membrane fluidity} depends on lipid composition:
\[\text{Fluidity} \propto \frac{\text{degree of unsaturation} + \text{chain shortness}}{\text{cholesterol content (moderate)}} \tag{7.1} \label{eq:unit_II_membrane_transport_item_1}\]

Cholesterol has a biphasic effect: at low temperature, it disrupts
crystalline packing (fluidises); at high temperature, it restricts
excessive fluidity (condenses). This \textbf{buffering} effect maintains
membrane fluidity across a physiological temperature range.

\subsubsection{Lipid Rafts and Membrane
Microdomains}\label{lipid-rafts-and-membrane-microdomains}

\textbf{Lipid rafts} are dynamic, cholesterol- and sphingolipid-enriched
microdomains (10--200 nm) in the outer leaflet that preferentially
recruit certain proteins: - \textbf{GPI-anchored proteins} partition
into rafts - \textbf{Signalling receptors} (e.g., T cell receptor, B
cell receptor) cluster in rafts upon ligand binding - \textbf{Caveolins}
(caveolae) represent a specialised, stable form of lipid raft

\textbf{Controversy:} The existence and functional significance of lipid
rafts has been debated. Single-molecule tracking studies reveal that
raft-like domains are transient (about 10--20 ms lifetime),
small (about 10--20 nm), and form/dissolve dynamically. STED
super-resolution microscopy (Eggeling et al., 2009, \emph{Nature})
confirmed that sphingolipids and GPI-anchored proteins are transiently
confined in nanoscale membrane domains, supporting a dynamic raft model.

\subsubsection{Updates to the Fluid Mosaic
Model}\label{updates-to-the-fluid-mosaic-model}

Since 1972, several refinements have been made:

\begin{enumerate}
\def\labelenumi{\arabic{enumi}.}
\tightlist
\item
  \textbf{Membrane is more mosaic than fluid:} Up to 50\% of membrane
  area is occupied by proteins; crowding restricts lateral diffusion
\item
  \textbf{Cytoskeletal fences:} The \hyperref[gl:actin]{\textbf{actin}}
  cortex creates corrals (``picket fence'' model, Kusumi et al.) that
  compartmentalise membrane proteins into about 40--300 nm
  domains, restricting free diffusion
\item
  \textbf{Lipid asymmetry} is more extensive than originally appreciated
\item
  \textbf{Membrane curvature} is actively generated by BAR-domain
  proteins, ESCRT complexes, and coat proteins
\item
  \textbf{Transbilayer communication:} Inner and outer leaflet lipids
  can be coupled through interdigitating acyl chains
\end{enumerate}

\begin{quote}
\textbf{Concept Check 1:} During apoptosis, phosphatidylserine (PS)
moves from the inner to the outer leaflet. What
\hyperref[gl:enzyme]{\textbf{enzyme}} mediates this, and why is PS
externalisation a critical signal for phagocytic clearance?
\end{quote}

\begin{center}\rule{0.5\linewidth}{0.5pt}\end{center}

\subsection{Membrane Proteins}\label{membrane-proteins}

\textbf{Integral (transmembrane) proteins} span the bilayer via
alpha-helical segments (of about 20 hydrophobic amino acids
each, sufficient to span the about 3 nm hydrophobic core) or
beta-barrel structures (found primarily in outer membranes of
Gram-negative bacteria, mitochondria, and
\hyperref[gl:chloroplast]{\textbf{chloroplast}}s).

\textbf{Peripheral proteins} associate noncovalently with membrane
surfaces or integral proteins via electrostatic interactions,
\hyperref[gl:hydrogen-bond]{\textbf{hydrogen bond}}s, or hydrophobic
interactions with the lipid headgroup region.

\textbf{Lipid-anchored proteins} are covalently modified with: -
\textbf{GPI-anchors} (outer leaflet): e.g., alkaline phosphatase, CD59
(complement regulator) - \textbf{Myristoyl} groups (inner leaflet,
N-terminal): e.g., Src kinase - \textbf{Palmitoyl} groups (inner
leaflet, Cys residue): e.g., Ras (reversible; regulates membrane
association) - \textbf{Farnesyl/geranylgeranyl} groups (inner leaflet,
C-terminal CAAX): e.g., Ras, Rho GTPases

\subsubsection{Functions of Membrane
Proteins}\label{functions-of-membrane-proteins}

{\def\LTcaptype{none} % do not increment counter
\begin{longtable}[]{@{}
  >{\raggedright\arraybackslash}p{(\linewidth - 2\tabcolsep) * \real{0.5000}}
  >{\raggedright\arraybackslash}p{(\linewidth - 2\tabcolsep) * \real{0.5000}}@{}}
\toprule\noalign{}
\begin{minipage}[b]{\linewidth}\raggedright
Function
\end{minipage} & \begin{minipage}[b]{\linewidth}\raggedright
Examples
\end{minipage} \\
\midrule\noalign{}
\endhead
\bottomrule\noalign{}
\endlastfoot
Transport (channel) & Aquaporins, K\(^+\) channels, Cl\(^-\) channels,
mechanosensitive channels \\
Transport (carrier) & GLUT glucose transporters, amino acid carriers,
nucleoside transporters \\
Active transport (pump) & Na\(^+\)/K\(^+\)-ATPase, Ca\(^{2+}\)-ATPase
(SERCA), H\(^+\)/K\(^+\)-ATPase \\
ABC transporter & MDR1/P-glycoprotein, CFTR, ABCA1 (cholesterol
efflux) \\
Receptor & EGFR, insulin R, beta\(_2\)-adrenergic R, rhodopsin, TLR4 \\
Enzyme & Adenylyl cyclase, guanylyl cyclase, gamma-secretase \\
Cell adhesion & Integrins, cadherins, selectins, IgCAMs \\
Structural anchor & Ankyrin-spectrin (RBC), dystrophin-glycoprotein
complex \\
\end{longtable}
}

\begin{center}\rule{0.5\linewidth}{0.5pt}\end{center}

\subsection{Passive Transport}\label{passive-transport}

Passive transport requires no energy input --- molecules move down their
electrochemical gradient.

\begin{figure}[htbp]
\centering
\includegraphics[width=0.82\linewidth,height=4.2in,keepaspectratio]{../figures/mermaid_inline/inline_mermaid_0038_993b48293347.png}
\caption{Classification of membrane transport mechanisms, from simple diffusion through channels and carriers to primary and secondary active transport and bulk transport.}
\end{figure}

\subsubsection{Simple Diffusion}\label{simple-diffusion}

For uncharged molecules, flux obeys Fick's First Law:

\[J = -D \frac{d[C]}{dx} = P \cdot \Delta [C] \tag{7.2} \label{eq:unit_II_membrane_transport_item_2}\]

where \(P\) = permeability coefficient (m/s) =
\(D \cdot K_{\text{oil/water}} / d\) (membrane thickness).

Membrane permeability: - \textbf{High:} small nonpolar molecules
(O\(_2\), CO\(_2\), N\(_2\), ethanol, benzene) - \textbf{Moderate:}
small polar uncharged (water, glycerol, urea) - \textbf{Low:} large
polar (glucose, amino acids) - \textbf{Very low/none:} ions (Na\(^+\),
K\(^+\), Cl\(^-\), Ca\(^{2+}\)) --- require protein channels

The \textbf{partition coefficient} (\(K_{\text{oil/water}}\)) is the
strongest predictor of simple diffusion rate. Overton's rule (1899):
membrane permeability correlates with oil/water partition coefficient.

\subsubsection{Osmosis and the van't Hoff
Equation}\label{osmosis-and-the-vant-hoff-equation}

Although water is technically a small polar molecule, its transport
across cell membranes is fast and quantitative enough to deserve its own
treatment. \textbf{Osmosis} is the net diffusion of water across a
semi-permeable membrane in response to a difference in solute
concentrations. Because water mass is conserved, osmosis amounts to a
redistribution of \emph{volume} --- and uncontrolled osmotic swelling
will lyse a cell.

The thermodynamic driving force is the difference in \textbf{water
chemical potential} \(\mu_w\) between two compartments:

\[\mu_w = \mu_w^0 + RT \ln a_w + \bar{V}_w P \tag{7.24} \label{eq:unit_II_membrane_transport_item_3}\]

where \(a_w\) is water activity (\texorpdfstring{\ensuremath{\approx}}{~} mole fraction in dilute solution) and
\(\bar{V}_w\) is the partial molar volume of water (about 18
mL/mol). Each dissolved solute particle reduces \(a_w\) proportionally
--- this is the molecular origin of colligative properties
(boiling-point elevation, freezing-point depression, vapour-pressure
lowering). Setting \(\mu_w\) equal on both sides of a semi-permeable
membrane and solving for the pressure required to balance a solute
gradient yields \textbf{van't Hoff's law}:

\begin{equation}
\Pi = i M R T
\label{eq:unit_II_vant_hoff}
\end{equation}

where \texorpdfstring{\ensuremath{\Pi}}{Pi} is the \textbf{osmotic pressure} (in atm or Pa), \(M\) is the
solute molar concentration, \(R\) is the gas constant, \(T\) is absolute
temperature, and \(i\) is the \textbf{van't Hoff factor} --- the number
of effective particles produced per formula unit (1 for glucose; 2 for
NaCl; 3 for CaCl\(_2\) at full dissociation).

In SI units with \(R = 0.0821\) L\texorpdfstring{\ensuremath{\cdot}}{.}atm/mol/K and 37 °C (\(T = 310\) K), 1
mM of an ideal solute generates
\(\Pi = 1 \cdot 0.001 \cdot 0.0821 \cdot 310 = 0.0254\) atm \texorpdfstring{\ensuremath{\approx}}{~} 25 mmHg.
Plasma is about 290 mOsm/L, generating about 7,500
mmHg of total osmotic pressure --- but cells don't burst because they
are surrounded by isotonic interstitial fluid, and primarily the
\emph{difference} matters.

\textbf{Worked Example: Osmotic swelling of a red blood cell.}

A red blood cell is placed in a hypotonic solution (200 mOsm/L; plasma
is 290 mOsm/L). Calculate (a) the initial osmotic pressure difference
driving water inward, (b) the cell's predicted volume change to reach
equilibrium.

\begin{enumerate}
\def\labelenumi{(\alph{enumi})}
\tightlist
\item
  The osmotic pressure gradient is:
  \[ \Delta \Pi = (290 - 200) \times 10^{-3} \cdot 0.0821 \cdot 310 = 0.0900 \cdot 25.45 = 2.29 \text{ atm} \approx 1{,}740 \text{ mmHg}  \label{eq:unit_II_membrane_transport_item_4}\]
\end{enumerate}

This drives water into the cell at a flux of
\(J_w = p_f \cdot \Delta C\), where \(p_f \approx 0.02\) cm/s for an
AQP1-rich erythrocyte. For a typical cell
(\(A \approx 140 \, \mu\text{m}^2\), \(V \approx 90 \, \mu\text{m}^3\)),
volume doubling occurs in roughly 200 ms --- fast enough to be visible
in a microscope as the cell swells from biconcave disc to sphere.

\begin{enumerate}
\def\labelenumi{(\alph{enumi})}
\setcounter{enumi}{1}
\tightlist
\item
  At osmotic equilibrium, water content adjusts to make intracellular
  osmolarity = 200 mOsm/L. If initial intracellular osmolyte content is
  \(V_0 \cdot C_0 = 90 \cdot 290 = 26{,}100\) \texorpdfstring{\ensuremath{\mu}}{mu}mol\texorpdfstring{\ensuremath{\cdot}}{.}\texorpdfstring{\ensuremath{\mu}}{mu}m\(^3\)/L, the new
  equilibrium volume is:
  \[ V_\text{new} = V_0 \cdot \frac{C_0}{C_\text{new}} = 90 \cdot \frac{290}{200} = 130.5 \, \mu\text{m}^3  \label{eq:unit_II_membrane_transport_item_5}\]
\end{enumerate}

A 45\% volume increase. RBCs reach a critical volume of
about 150 \texorpdfstring{\ensuremath{\mu}}{mu}m\(^3\) before haemolysis (membrane area cannot
accommodate further sphere expansion), so the cell is on the brink of
lysis.

\textbf{Tonicity terms (clinical):} Hypotonic solutions cause cells to
swell (haemolysis if severe). Isotonic solutions (e.g., 0.9\% NaCl, 5\%
dextrose, lactated Ringer's) preserve volume. Hypertonic solutions cause
cells to shrink (crenation in RBCs). Aquaporin defects cause major fluid
balance disorders: \textbf{AQP2 mutations} cause nephrogenic diabetes
insipidus (water reabsorption fails despite vasopressin); \textbf{AQP4
autoantibodies} cause neuromyelitis optica (astrocyte water-handling
fails, edema and demyelination). The osmotic pressure of plasma proteins
(mostly albumin, about 25 mmHg) is the \textbf{oncotic
pressure} that retains fluid in the vasculature --- its loss in
nephrotic syndrome and liver failure causes generalised oedema.

\begin{quote}
\textbf{Concept Check 1b:} A patient with severe burns receives 4 L of
0.45\% NaCl (``half-normal saline'') rapidly. Calculate the tonicity
relative to plasma (290 mOsm/L) and predict the consequence for red
blood cells.
\end{quote}

\subsubsection{Facilitated Diffusion via
Channels}\label{facilitated-diffusion-via-channels}

\textbf{Ion channels} are integral membrane proteins forming aqueous
pores. They are: - \textbf{Selective:} ion selectivity filter determines
which ions pass - \textbf{Gated:} open/close in response to specific
stimuli - \textbf{Fast:} ion throughput rates of 10\(^7\)--10\(^8\)
ions/s (near diffusion limit)

\textbf{Gating mechanisms:}

\begin{itemize}
\tightlist
\item
  \textbf{Voltage-gated:} opened by membrane-potential change; examples
  include sodium channel Nav1.1--1.9, potassium channel Kv1--12, and
  calcium channel Cav1--3 families.
\item
  \textbf{Ligand-gated, extracellular:} opened by neurotransmitter
  binding; examples include nAChR, GABAA receptor, NMDAR, and AMPAR.
\item
  \textbf{Ligand-gated, intracellular:} opened by second messengers;
  examples include calcium-activated potassium channel BK and retinal
  cGMP-gated channels.
\item
  \textbf{Mechanosensitive:} opened by stretch or pressure; examples
  include Piezo1/2, TREK-1, and bacterial MscL.
\item
  \textbf{Temperature-gated:} opened by heat or cold; examples include
  TRPV1 (\textgreater43 degrees C, capsaicin) and TRPM8 (\textless26
  degrees C, menthol).
\item
  \textbf{Light-gated:} opened by photons; channelrhodopsin ChR2 is the
  standard optogenetic example.
\end{itemize}

\textbf{Selectivity filter of K\(^+\) channels} (MacKinnon, Nobel Prize
2003): The selectivity filter contains the signature sequence TVGYG.
Four carbonyl oxygens from each of the four subunits line the pore,
precisely mimicking the hydration shell of K\(^+\) (radius 1.33
angstroms). Na\(^+\) (radius 0.95 angstroms) is too small to be
coordinated effectively --- the energetic cost of dehydrating Na\(^+\)
without compensating coordination makes Na\(^+\) passage unfavourable.
This elegant size-exclusion mechanism achieves 10,000:1 selectivity for
K\(^+\) over Na\(^+\).

\textbf{Aquaporins (AQP):} Facilitate water transport (not solutes).
AQP1 (Peter Agre, Nobel Prize 2003) channels about 3 x 10\(^8\)
water molecules/s while excluding protons (the Grotthuss mechanism is
interrupted by a critical asparagine residue in the channel pore, and
the electrostatic field of the NPA motif reorients water molecules,
preventing H\(_3\)O\(^+\) passage).

\textbf{Aquaporin discovery and selectivity in detail.} Before 1992,
biologists could not explain how erythrocytes and renal tubule cells
achieve water permeability rates 100\texorpdfstring{\ensuremath{\times}}{x} faster than predicted from the
lipid bilayer alone. Peter Agre's group, while studying Rh-blood-group
proteins, noticed a contaminating 28 kDa membrane protein.
Reconstitution experiments (Preston et al., 1992, \emph{Science})
demonstrated that this protein --- christened \textbf{AQP1}
(aquaporin-1) --- conferred about 10\texorpdfstring{\ensuremath{\times}}{x} water permeability when
expressed in \emph{Xenopus} oocytes. Agre shared the 2003 Nobel Prize in
Chemistry. Thirteen aquaporin paralogues are now known in humans, with
tissue-specific expression: AQP2 (renal collecting duct,
vasopressin-regulated), AQP4 (astrocyte endfeet, target of NMO
autoantibodies), AQP5 (salivary/lacrimal glands), AQP7 (adipocytes,
glycerol channel).

The structural basis of selectivity is exquisite. Each AQP monomer is a
hexa-helical bundle with two short re-entrant helices (carrying the
conserved \textbf{NPA motifs} --- Asn-Pro-Ala) that meet in the middle
of the membrane to form the narrowest part of the pore
(about 2.8 Å, just wide enough for a single water molecule).
Two mechanisms exclude protons:

\begin{enumerate}
\def\labelenumi{\arabic{enumi}.}
\tightlist
\item
  \textbf{Electrostatic barrier:} The two NPA half-helices have positive
  macrodipoles meeting in the middle of the pore --- repelling
  H\(_3\)O\(^+\) but attracting electrically neutral water.
\item
  \textbf{Hydrogen-bond reorientation:} The central asparagines force
  passing water molecules to flip their hydrogen-bond donor/acceptor
  orientation. This breaks the proton-relay (Grotthuss) chain that would
  otherwise allow H\(^+\) to ``hop'' through the channel.
\end{enumerate}

The water permeability coefficient (\(p_f\), the \textbf{osmotic
permeability}) for a single AQP1 channel is
about \(3 \times 10^{-14}\) cm\(^3\)/s, corresponding to
about 3 \texorpdfstring{\ensuremath{\times}}{x} 10\(^9\) water molecules/s under typical osmotic
gradients. A red blood cell expressing about 200,000 AQP1
channels can therefore exchange its entire water volume in \textless100
ms --- the basis of erythrocyte swelling/shrinking responses.

\textbf{Clinical aquaporins.} Loss-of-function mutations in AQP2 cause
\textbf{nephrogenic diabetes insipidus} (kidneys cannot concentrate
urine despite normal vasopressin). Autoantibodies against AQP4 cause
\textbf{neuromyelitis optica (Devic's disease)}, a demyelinating
disorder distinct from MS. AQP1 is upregulated in many tumours and is
being explored as a target in glioma therapy.

\textbf{Clinical Connection: Channelopathies.} Ion channel
\hyperref[gl:mutation]{\textbf{mutation}}s cause a wide range of
diseases:

\begin{itemize}
\tightlist
\item
  \textbf{Long QT syndrome:} Mutations in cardiac potassium or sodium
  channels prolong repolarisation and can trigger fatal arrhythmias.
  Drug-development screens therefore test for hERG channel block.
\item
  \textbf{Cystic fibrosis:} Mutations in CFTR disrupt chloride
  transport. The common deltaF508 variant causes misfolding and ER
  retention; modulator combinations such as lumacaftor/ivacaftor and
  elexacaftor/tezacaftor/ivacaftor partly restore channel function.
\item
  \textbf{Malignant hyperthermia:} RyR1 mutations cause uncontrolled
  Ca\(^{2+}\) release from sarcoplasmic reticulum after exposure to
  volatile anaesthetics. See \cref{sec:unit_II_cell_signaling} for
  calcium signalling.
\end{itemize}

\begin{quote}
\textbf{Concept Check 2:} Tetrodotoxin (TTX, from pufferfish) blocks
voltage-gated Na\(^+\) channels by binding the selectivity filter from
the extracellular side. Predict the effects of TTX on (a)
\hyperref[gl:action-potential]{\textbf{action potential}} generation,
(b) nerve conduction, and (c) skeletal muscle contraction.
\end{quote}

\subsubsection{Ion Channel Gating
Kinetics}\label{ion-channel-gating-kinetics}

A single ion channel in a patch-clamp recording flickers between open
and closed states stochastically. The time-averaged current through a
population of \(N\) channels with single-channel conductance \texorpdfstring{\ensuremath{\gamma}}{gamma} (in pS =
pico-siemens) and open probability \(P_o\) is:

\begin{equation}
I = N \cdot P_o \cdot \gamma \cdot (V_m - E_\text{rev})
\label{eq:unit_II_macroscopic_current}
\end{equation}

where \(V_m\) is the membrane potential and \(E_\text{rev}\) is the
reversal potential (for a perfectly selective channel,
\(E_\text{rev} = E_\text{ion}\)). This decomposition is one of the most
powerful in biophysics: drugs can change channel function by altering
\(N\) (channel expression/internalisation), \(P_o\) (gating modulators,
allosteric drugs), or \texorpdfstring{\ensuremath{\gamma}}{gamma} (pore-block, modifications of selectivity
filter), and each can be measured separately.

\textbf{Single-channel conductances of representative channels (in pS):}

\begin{itemize}
\tightlist
\item
  \textbf{nAChR, muscle endplate:} conductance 30--50 pS; open
  probability about 0.85 with ACh; throughput about 3 x 10\^{}7 ions/s.
\item
  \textbf{BK, large-conductance calcium-activated potassium channel:}
  conductance 200--250 pS; open probability 0.1--0.9 depending on
  calcium; throughput up to 10\^{}8 ions/s.
\item
  \textbf{Kv1.x delayed rectifier:} conductance 10--20 pS; open
  probability about 0.7 at +20 mV; throughput about 10\^{}6 ions/s.
\item
  \textbf{Nav1.x:} conductance 15--20 pS; transient open probability
  below 0.5; throughput about 5 x 10\^{}6 ions/s during an action
  potential.
\item
  \textbf{L-type Cav1.2:} conductance about 25 pS; open probability
  about 0.3; throughput about 10\^{}5 ions/s.
\item
  \textbf{ClC-1 skeletal chloride channel:} conductance about 1 pS with
  gating-coupled pores; open probability about 0.5 at rest.
\item
  \textbf{CFTR chloride channel / ABC transporter:} conductance 8--10
  pS; PKA-dependent open probability 0.3--0.5.
\end{itemize}

\textbf{Open probability and gating models.} The open probability
\(P_o\) depends on the gating stimulus (voltage, ligand concentration,
mechanical force). For a voltage-gated channel with effective gating
charge \(z_g\) (typical about 4--6 for Na\(_V\),
about 6--10 for K\(_V\)), the equilibrium open probability
follows a Boltzmann relation:

\[P_o(V) = \frac{1}{1 + \exp\left(-\frac{z_g F (V - V_{1/2})}{RT}\right)} \tag{7.19} \label{eq:unit_II_membrane_transport_item_6}\]

where \(V_{1/2}\) is the half-activation voltage. This sigmoidal curve
is steep (about 5--10 mV per e-fold change) --- small voltage
perturbations cause large changes in \(P_o\), the basis of the action
potential's switch-like behaviour. For ligand-gated channels with \(n\)
binding sites and Hill coefficient \(h\):

\[P_o([L]) = \frac{[L]^h}{[L]^h + K_d^h} \tag{7.20} \label{eq:unit_II_membrane_transport_item_7}\]

For the muscle nAChR, \(h \approx 1.5\) and \(K_d \approx 30\) \texorpdfstring{\ensuremath{\mu}}{mu}M for
ACh: half-maximal activation at about 30 \texorpdfstring{\ensuremath{\mu}}{mu}M, full activation by
about 300 \texorpdfstring{\ensuremath{\mu}}{mu}M. The sub-millisecond opening of nAChRs after ACh
release at the neuromuscular junction (peak {[}ACh{]} about 1
mM at the postsynaptic membrane) ensures essentially complete channel
activation in every action potential.

\subsubsection{Electroneutral vs.~Electrogenic
Transport}\label{electroneutral-vs.-electrogenic-transport}

Whether a transporter contributes net charge to the membrane potential
depends on the \textbf{stoichiometry of the transport cycle}.

{\def\LTcaptype{none} % do not increment counter
\begin{longtable}[]{@{}
  >{\raggedright\arraybackslash}p{(\linewidth - 6\tabcolsep) * \real{0.1746}}
  >{\raggedright\arraybackslash}p{(\linewidth - 6\tabcolsep) * \real{0.2063}}
  >{\raggedright\arraybackslash}p{(\linewidth - 6\tabcolsep) * \real{0.4127}}
  >{\raggedright\arraybackslash}p{(\linewidth - 6\tabcolsep) * \real{0.2063}}@{}}
\toprule\noalign{}
\begin{minipage}[b]{\linewidth}\raggedright
Transporter
\end{minipage} & \begin{minipage}[b]{\linewidth}\raggedright
Stoichiometry
\end{minipage} & \begin{minipage}[b]{\linewidth}\raggedright
Net charge moved per cycle
\end{minipage} & \begin{minipage}[b]{\linewidth}\raggedright
Electrogenic?
\end{minipage} \\
\midrule\noalign{}
\endhead
\bottomrule\noalign{}
\endlastfoot
Na\(^+\)/K\(^+\)-ATPase & 3 Na\(^+\) out / 2 K\(^+\) in / 1 ATP & +1
outward & Yes \\
Ca\(^{2+}\)-ATPase (SERCA) & 2 Ca\(^{2+}\) in (to ER) / 2 H\(^+\) out /
1 ATP & 0 & No (electroneutral) \\
H\(^+\)/K\(^+\)-ATPase (gastric) & 1 H\(^+\) out / 1 K\(^+\) in / 1 ATP
& 0 & No \\
NCX (Na\(^+\)/Ca\(^{2+}\) exchanger) & 3 Na\(^+\) in / 1 Ca\(^{2+}\) out
& +1 inward & Yes \\
NHE1 (Na\(^+\)/H\(^+\) exchanger) & 1 Na\(^+\) in / 1 H\(^+\) out & 0 &
No \\
SGLT1 (Na\(^+\)/glucose symport) & 2 Na\(^+\) in / 1 glucose & +2 inward
& Yes \\
AE1 (Cl\(^-\)/HCO\(_3^-\) exchanger) & 1 Cl\(^-\) out / 1 HCO\(_3^-\) in
& 0 & No \\
\end{longtable}
}

\textbf{Energetic consequences.} Electroneutral transporters move
solutes ``for free'' with respect to the membrane potential --- their
thermodynamic feasibility depends primarily on chemical concentration
gradients. Electrogenic transporters, by contrast, are driven by
\emph{both} concentration and voltage. The Na\(^+\)/K\(^+\)-ATPase, for
example, would still hydrolyse ATP if Na\(^+\) and K\(^+\)
concentrations were equalised, because moving net positive charge
against the membrane potential (interior negative) costs additional
energy: the \textbf{electrochemical gradient} is the relevant
thermodynamic quantity:

\[\Delta G_\text{ion} = RT \ln \frac{[C]_\text{out}}{[C]_\text{in}} + z F V_m \tag{7.21} \label{eq:unit_II_membrane_transport_item_8}\]

For a 100\texorpdfstring{\ensuremath{\times}}{x} concentration gradient and 100 mV opposing voltage:
\(\Delta G \approx +12 + 9.6 = 21.6\) kJ/mol. ATP hydrolysis releases
about 50 kJ/mol under cellular conditions, so a single ATP can
drive about 2 ions of net charge against this combined gradient
--- exactly the stoichiometry observed in the Na\(^+\)/K\(^+\)-ATPase (3
Na\(^+\) out -- 2 K\(^+\) in = +1 net per ATP, with margin for
irreversibility).

\begin{quote}
\textbf{Clinical Connection: Channelopathy diagnostics.} When a patient
presents with arrhythmia, episodic weakness, or seizures, modern
diagnostics include \textbf{gene panels} (sequencing known channel
genes) followed by \textbf{functional reconstitution} of variants in
HEK293 or oocyte expression systems. Patch-clamp measurements decompose
the disease-causing change into \(N\), \(P_o\), and \texorpdfstring{\ensuremath{\gamma}}{gamma} effects: for
example, the cystic fibrosis \texorpdfstring{\ensuremath{\Delta}}{Delta}F508 mutation reduces \(N\) at the
membrane (trafficking defect, fixed by lumacaftor); G551D reduces
\(P_o\) (gating defect, fixed by ivacaftor). The combination drug
Orkambi addresses both --- a triumph of mechanism-guided pharmacology.
\end{quote}

\subsubsection{Facilitated Diffusion via Carriers
(Uniporters)}\label{facilitated-diffusion-via-carriers-uniporters}

\textbf{GLUT transporters} (SLC2A family; 14 isoforms in humans)
facilitate glucose diffusion:

{\def\LTcaptype{none} % do not increment counter
\begin{longtable}[]{@{}
  >{\raggedright\arraybackslash}p{(\linewidth - 6\tabcolsep) * \real{0.2973}}
  >{\raggedright\arraybackslash}p{(\linewidth - 6\tabcolsep) * \real{0.1622}}
  >{\raggedright\arraybackslash}p{(\linewidth - 6\tabcolsep) * \real{0.2703}}
  >{\raggedright\arraybackslash}p{(\linewidth - 6\tabcolsep) * \real{0.2703}}@{}}
\toprule\noalign{}
\begin{minipage}[b]{\linewidth}\raggedright
Transporter
\end{minipage} & \begin{minipage}[b]{\linewidth}\raggedright
Tissue
\end{minipage} & \begin{minipage}[b]{\linewidth}\raggedright
\(K_m\) (mM)
\end{minipage} & \begin{minipage}[b]{\linewidth}\raggedright
Regulation
\end{minipage} \\
\midrule\noalign{}
\endhead
\bottomrule\noalign{}
\endlastfoot
GLUT1 & Erythrocytes, brain endothelium, placenta & 1.5 &
Constitutive \\
GLUT2 & Liver, pancreatic beta-cells, small intestine & 17 & Low
affinity; glucose sensor \\
GLUT3 & \hyperref[gl:neuron]{\textbf{Neuron}}s & 1.4 & Constitutive;
high affinity for brain \\
\hyperref[gl:glut4]{\textbf{GLUT4}} & Muscle, adipose & 5 &
Insulin-stimulated PM insertion \\
GLUT5 & Small intestine & N/A & Fructose transporter \\
\end{longtable}
}

The \textbf{alternating access mechanism}: the carrier alternates
between outward-facing (substrate binds from extracellular side) and
inward-facing (substrate released to
\hyperref[gl:cytoplasm]{\textbf{cytoplasm}}) conformations. This is
slower than channel transport (about 10\(^2\)--10\(^4\)
molecules/s) but allows specificity and saturability.

\subsection{Worked Example: GLUT Transporter
Kinetics}\label{worked-example-glut-transporter-kinetics}

\emph{Problem:} GLUT1 has \(K_m = 1.5\) mM and \(V_{max} = 200\)
\texorpdfstring{\ensuremath{\mu}}{mu}mol/min per gram of membrane protein. Blood glucose is about 5
mM. At what fraction of \(V_{max}\) is GLUT1 operating?

\emph{Solution:}

Using the Michaelis-Menten equation:

\[v = \frac{V_{max} \cdot [S]}{K_m + [S]} = \frac{200 \times 5}{1.5 + 5} = \frac{1000}{6.5} = 154 \; \mu\text{mol/min/g} \tag{7.3} \label{eq:unit_II_membrane_transport_item_9}\]

\[\frac{v}{V_{max}} = \frac{154}{200} = 0.77 = 77\% \tag{7.4} \label{eq:unit_II_membrane_transport_item_10}\]

GLUT1 operates at 77\% of maximum capacity at normal blood glucose ---
providing a safety margin while ensuring high glucose flux to the brain.

\begin{center}\rule{0.5\linewidth}{0.5pt}\end{center}

\subsection{Active Transport}\label{active-transport}

Active transport moves solutes against their electrochemical gradient,
requiring energy (usually ATP hydrolysis or proton motive force).

\subsubsection{Primary Active Transport ---
Pumps}\label{primary-active-transport-pumps}

\textbf{Na\(^+\)/K\(^+\)-ATPase (sodium-potassium pump):} Exports 3
Na\(^+\) and imports 2 K\(^+\) per ATP consumed (overall electrogenic;
net +1 charge out). This pump consumes about 25\% of the body's
ATP (up to 70\% in neurons).

\begin{figure}[htbp]
\centering
\includegraphics[width=0.82\linewidth,height=4.2in,keepaspectratio]{../figures/mermaid_inline/inline_mermaid_0039_6d07c3fd2370.png}
\caption{The Post-Albers cycle of the Na\$\textasciicircum{}+\$/K\$\textasciicircum{}+\$-ATPase. The pump alternates between E1 (inward-facing, high Na\$\textasciicircum{}+\$ affinity) and E2 (outward-facing, high K\$\textasciicircum{}+\$ affinity) conformations. Phosphorylation by ATP and subsequent dephosphorylation drive the conformational changes.}
\end{figure}

This maintains: - Resting membrane potential (K\(^+\) gradient) - High
intracellular {[}K\(^+\){]} (about 140 mM) and low
{[}Na\(^+\){]} (about 12 mM) - Low intracellular Na\(^+\) that
drives secondary active transport - Cell volume regulation (preventing
osmotic swelling)

\textbf{Inhibition by ouabain/digitalis:} Cardiac glycosides block the
K\(^+\)-binding E2-P form. This elevates intracellular Na\(^+\), which
reduces Na\(^+\)/Ca\(^{2+}\) exchanger activity (NCX normally uses the
Na\(^+\) gradient to export Ca\(^{2+}\)), so {[}Ca\(^{2+}\){]}\(_i\)
rises, causing stronger cardiac contraction. Used for heart failure;
narrow therapeutic window (toxicity causes arrhythmias).

\textbf{Other primary active transport pumps:} -
\textbf{Ca\(^{2+}\)-ATPase (SERCA):} pumps Ca\(^{2+}\) from cytoplasm
into ER/SR lumen; maintains {[}Ca\(^{2+}\){]}\(_i\) at
about 100 nM (10,000-fold lower than extracellular); critical
for muscle relaxation - \textbf{H\(^+\)/K\(^+\)-ATPase:} gastric
parietal cells; pumps H\(^+\) into stomach lumen
(\hyperref[gl:ph]{\textbf{pH}} about 1); target of proton pump
inhibitors (omeprazole, lansoprazole) for acid reflux/ulcer treatment -
\textbf{V-type H\(^+\)-ATPase:} acidifies lysosomes, endosomes; does not
use a phosphorylated intermediate

\subsubsection{ABC Transporters}\label{abc-transporters}

\textbf{ATP-Binding Cassette (ABC) transporters} are a superfamily of
about 48 members in humans. They use ATP hydrolysis to
transport diverse substrates across membranes:

\begin{itemize}
\tightlist
\item
  \textbf{MDR1/P-glycoprotein (ABCB1):} Broad-specificity drug efflux
  pump; exports hydrophobic compounds from the cell. Overexpressed in
  many cancers, causing \textbf{multidrug resistance} --- tumour cells
  pump out chemotherapy drugs before they can act. Substrates include
  taxol, doxorubicin, vincristine.
\item
  \textbf{CFTR (ABCC7):} Unique ABC transporter that functions as a
  Cl\(^-\) channel. Mutations cause \textbf{cystic fibrosis} (see
  Clinical Connection above).
\item
  \textbf{ABCA1:} Cholesterol and phospholipid efflux to apoA-I;
  critical for HDL formation. Loss-of-function mutations cause
  \textbf{Tangier disease} (very low HDL, cholesterol deposition in
  tissues).
\item
  \textbf{TAP1/TAP2 (ABCB2/3):} Transport antigenic peptides from
  cytoplasm into ER lumen for MHC class I loading and immune
  presentation.
\end{itemize}

\subsubsection{Secondary Active
Transport}\label{secondary-active-transport}

Uses the Na\(^+\) electrochemical gradient (generated by the primary
pump) to drive uphill transport of other solutes:

\begin{itemize}
\tightlist
\item
  \textbf{Symport (co-transport, same direction):} SGLT1 (intestinal
  glucose): 2 Na\(^+\) + 1 glucose move in together; net uptake even if
  {[}glucose{]}\(_{\text{in}}\) \textgreater{}
  {[}glucose{]}\(_{\text{out}}\). SGLT2 (kidney proximal tubule): 1
  Na\(^+\) + 1 glucose; target of gliflozin drugs for type 2 diabetes.
\item
  \textbf{Antiport (exchange, opposite directions):} NCX
  (Na\(^+\)/Ca\(^{2+}\) exchanger): 3 Na\(^+\) in, 1 Ca\(^{2+}\) out;
  critical for cardiac Ca\(^{2+}\)
  \hyperref[gl:homeostasis]{\textbf{homeostasis}}. NHE (Na\(^+\)/H\(^+\)
  exchanger): Na\(^+\) in, H\(^+\) out; regulates intracellular pH.
\end{itemize}

\begin{quote}
\textbf{Clinical Connection: SGLT2 Inhibitors in Diabetes and Heart
Failure} SGLT2 inhibitors (empagliflozin, dapagliflozin, canagliflozin)
block glucose reabsorption in the kidney proximal tubule, causing
glucosuria (glucose loss in urine) and lowering blood glucose.
Remarkably, these drugs also reduce cardiovascular death and heart
failure hospitalisations, even in non-diabetic patients --- likely
through natriuresis, osmotic diuresis, and favourable metabolic effects.
They represent a rare case where a simple transport mechanism becomes a
blockbuster drug target. see \cref{sec:unit_III_metabolic_integration}
(Metabolic Integration) for insulin signalling.
\end{quote}

\begin{quote}
\textbf{Concept Check 3:} The antibiotic gramicidin forms a channel in
bacterial membranes that allows monovalent cations (Na\(^+\), K\(^+\))
to flow freely. Predict how gramicidin would affect (a) the bacterial
membrane potential, (b) the proton motive force, and (c) bacterial ATP
synthesis.
\end{quote}

\begin{center}\rule{0.5\linewidth}{0.5pt}\end{center}

\subsection{The Nernst Equation and Membrane
Potential}\label{the-nernst-equation-and-membrane-potential}

\subsubsection{Electrochemical
Potential}\label{electrochemical-potential}

The electrochemical potential of an ion is:

\[\tilde{\mu}_i = \mu_i^0 + RT\ln[C_i] + z_i F V \tag{7.5} \label{eq:unit_II_membrane_transport_item_11}\]

where \(z_i\) = ionic charge, \(F\) = Faraday constant (96,485 C/mol),
\(V\) = membrane potential.

At equilibrium (no net flux), the \textbf{Nernst equation} gives the
equilibrium potential \(E_i\):

\[E_i = \frac{RT}{z_i F} \ln \frac{[C_i]_{\text{out}}}{[C_i]_{\text{in}}} \tag{7.6} \label{eq:unit_II_membrane_transport_item_12}\]

At 37 degrees C (310 K), \(RT/F\) = 26.7 mV. Converting to log\(_{10}\):

\[E_i = \frac{61.5 \; \text{mV}}{z_i} \log_{10} \frac{[C_i]_{\text{out}}}{[C_i]_{\text{in}}} \tag{7.7} \label{eq:unit_II_membrane_transport_item_13}\]

\subsubsection{Derivation of the Nernst
Equation}\label{derivation-of-the-nernst-equation}

Starting from the condition of zero net electrochemical driving force at
equilibrium:

\[\Delta \tilde{\mu}_i = 0 \tag{7.8} \label{eq:unit_II_membrane_transport_item_14}\]

\[RT \ln \frac{[C_i]_{\text{in}}}{[C_i]_{\text{out}}} + z_i F (V_{\text{in}} - V_{\text{out}}) = 0 \tag{7.9} \label{eq:unit_II_membrane_transport_item_15}\]

\[z_i F \cdot E_i = -RT \ln \frac{[C_i]_{\text{in}}}{[C_i]_{\text{out}}} = RT \ln \frac{[C_i]_{\text{out}}}{[C_i]_{\text{in}}} \tag{7.10} \label{eq:unit_II_membrane_transport_item_16}\]

\[E_i = \frac{RT}{z_i F} \ln \frac{[C_i]_{\text{out}}}{[C_i]_{\text{in}}} \tag{7.11} \label{eq:unit_II_membrane_transport_item_17}\]

\subsection{Worked Example: Nernst
Potential}\label{worked-example-nernst-potential}

\emph{Problem:} Calculate the Nernst equilibrium potential for
Ca\(^{2+}\) at 37 degrees C, given {[}Ca\(^{2+}\){]}\(_{\text{out}}\) =
2.5 mM and {[}Ca\(^{2+}\){]}\(_{\text{in}}\) = 0.0001 mM (100 nM).

\emph{Solution:}

\[E_{Ca} = \frac{RT}{z_{Ca}F} \ln \frac{[Ca^{2+}]_{\text{out}}}{[Ca^{2+}]_{\text{in}}} = \frac{26.7 \; \text{mV}}{2} \ln \frac{2.5}{0.0001} \tag{7.12} \label{eq:unit_II_membrane_transport_item_18}\]

\[E_{Ca} = 13.35 \; \text{mV} \times \ln(25,000) = 13.35 \times 10.13 = +135 \; \text{mV} \tag{7.13} \label{eq:unit_II_membrane_transport_item_19}\]

The strongly positive \(E_{Ca}\) means that Ca\(^{2+}\) has a massive
driving force to enter cells. Brief Ca\(^{2+}\) channel openings can
therefore cause significant signalling events.

\textbf{Nernst potentials of key ions:}

{\def\LTcaptype{none} % do not increment counter
\begin{longtable}[]{@{}llll@{}}
\toprule\noalign{}
Ion & {[}inside{]} (mM) & {[}outside{]} (mM) & \(E_i\) (mV, 37 degrees
C) \\
\midrule\noalign{}
\endhead
\bottomrule\noalign{}
\endlastfoot
K\(^+\) & 140 & 5 & -89 \\
Na\(^+\) & 12 & 145 & +63 \\
Ca\(^{2+}\) & 0.0001 & 2.5 & +135 \\
Cl\(^-\) & 4 & 110 & -82 \\
\end{longtable}
}

\textbf{Worked Example: Nernst Potentials for K\(^+\), Na\(^+\), and
Cl\(^-\).}

Each row of the table above can be reproduced from
\cref{eq:unit_II_membrane_transport_nernst} below by careful
sign-handling. Using the simplified form at 37 °C:
\(E_i = (61.5 \text{ mV} / z_i) \log_{10}([\text{out}]/[\text{in}])\).

\begin{equation}
E_i = \frac{RT}{z_i F} \ln \frac{[C_i]_{\text{out}}}{[C_i]_{\text{in}}}
\label{eq:unit_II_membrane_transport_nernst}
\end{equation}

\emph{Potassium (z = +1):}
\(E_K = 61.5 \cdot \log_{10}(5/140) = 61.5 \cdot \log_{10}(0.0357) = 61.5 \cdot (-1.447) = -89.0 \text{ mV}\).
K\(^+\) is concentrated inside, so equilibrium drives K\(^+\) outward;
the equilibrium potential is negative (interior must be
about 89 mV more negative to halt K\(^+\) efflux). The resting
membrane potential (\(V_m \approx -70\) mV) is therefore \emph{less
negative} than \(E_K\) --- meaning K\(^+\) continues to leak slowly
outward at rest.

\emph{Sodium (z = +1):}
\(E_{Na} = 61.5 \cdot \log_{10}(145/12) = 61.5 \cdot \log_{10}(12.08) = 61.5 \cdot 1.082 = +66.6 \text{ mV}\).
(Tabulated values vary 63--67 mV depending on the assumed
concentrations.) Na\(^+\) is concentrated outside, so equilibrium drives
Na\(^+\) inward; the equilibrium potential is strongly positive. The
resting membrane potential is \emph{very far} from \(E_{Na}\) ---
meaning Na\(^+\) has a massive electrochemical force driving it inward,
restrained primarily by the low resting Na\(^+\) permeability.

\emph{Chloride (z = \texorpdfstring{\ensuremath{-}}{-}1):}
\(E_{Cl} = (61.5 / -1) \cdot \log_{10}(110/4) = -61.5 \cdot \log_{10}(27.5) = -61.5 \cdot 1.439 = -88.5 \text{ mV}\).
(Tabulated values vary 80--90 mV depending on cell type --- neurons have
higher {[}Cl\(^-\){]}\(_i\) than skeletal muscle.) Note the
\emph{negative} sign on \(z_\text{Cl}\) inverts the ratio's effect: a
higher {[}Cl\(^-\){]}\(_o\) than {[}Cl\(^-\){]}\(_i\) yields a
\emph{negative} equilibrium potential, the opposite of K\(^+\). In adult
neurons, \(E_\text{Cl}\) sits near \(V_m\), so Cl\(^-\) is essentially
at equilibrium and small Cl\(^-\) permeability changes (e.g.,
GABA\(_A\)-receptor opening) cause primarily modest hyperpolarisation by
stabilising \(V_m\) near \(E_\text{Cl}\). In neonatal neurons,
{[}Cl\(^-\){]}\(_i\) is higher (about 25 mM) due to immature
KCC2 expression, \(E_\text{Cl}\) becomes about \texorpdfstring{\ensuremath{-}}{-}40 mV, and GABA
is \emph{depolarising} --- a fact with major implications for early
brain development and neonatal seizures.

These three numbers --- \(E_K \approx -89\), \(E_{Na} \approx +63\),
\(E_{Cl} \approx -82\) --- together with the relative permeabilities
\(P_K : P_{Na} : P_{Cl}\) generate the entire repertoire of resting and
action potentials in excitable cells.

The resting membrane potential (-70 mV) lies close to \(E_K\) because
the resting cell membrane is about 25x more permeable to
K\(^+\) than Na\(^+\). K\(^+\) ``leaks'' out through K2P (two-pore
domain) leak channels; uncovered negative charges on impermeable
intracellular proteins (Donnan effect) and the electrogenic
Na\(^+\)/K\(^+\) pump (net export of +1 charge per cycle) contribute the
remainder.

\subsubsection{Goldman-Hodgkin-Katz
Equation}\label{goldman-hodgkin-katz-equation}

When multiple ions carry current, the resting membrane potential is
given by the \textbf{Goldman equation}:

\begin{equation}
V_m = \frac{RT}{F} \ln \frac{P_K[K^+]_o + P_{Na}[Na^+]_o + P_{Cl}[Cl^-]_i}{P_K[K^+]_i + P_{Na}[Na^+]_i + P_{Cl}[Cl^-]_o}
\label{eq:unit_II_ghk}
\end{equation}

Note: anions (Cl\(^-\)) appear with reversed subscripts because of their
negative charge.

\textbf{Derivation sketch.} The full Goldman-Hodgkin-Katz (GHK)
derivation begins from the Nernst-Planck equation for ionic flux \(J_i\)
across a membrane of thickness \(d\) in a constant electric field:

\[J_i = -D_i \left( \frac{dC_i}{dx} + z_i C_i \frac{F}{RT} \frac{dV}{dx} \right) \tag{7.22} \label{eq:unit_II_membrane_transport_item_20}\]

Integrating across the membrane assuming a constant field
(\(dV/dx = -V_m/d\)) and a single permeability \(P_i = D_i \beta_i / d\)
(where \(\beta_i\) is the partition coefficient between water and
membrane) gives the \textbf{GHK current equation}:

\[I_i = z_i F P_i \frac{z_i F V_m / RT}{1 - \exp(-z_i F V_m / RT)} \left( [C_i]_\text{in} - [C_i]_\text{out} \exp(-z_i F V_m / RT) \right) \tag{7.23} \label{eq:unit_II_membrane_transport_item_21}\]

At the resting membrane potential, the \emph{net} current must be zero:
\(\sum_i I_i = 0\). Solving this for K\(^+\), Na\(^+\), and Cl\(^-\)
(and treating Cl\(^-\) as monovalent anion) algebraically yields the GHK
voltage equation \cref{eq:unit_II_ghk}. The key conceptual takeaways
are: (i) the membrane potential is a \emph{weighted log average} of the
Nernst potentials, with weights set by permeability; (ii) the ion with
the highest permeability dominates; (iii) shifts in permeability ratios
(e.g., during the action potential) produce predictable shifts in
\(V_m\).

\subsection{Worked Example: Goldman
Equation}\label{worked-example-goldman-equation}

\emph{Problem:} Calculate \(V_m\) at 37 degrees C given relative
permeabilities \(P_K : P_{Na} : P_{Cl}\) = 1.0 : 0.04 : 0.45 and the ion
concentrations in the table above.

\emph{Solution:}

\[V_m = 26.7 \; \text{mV} \times \ln \frac{(1.0)(5) + (0.04)(145) + (0.45)(4)}{(1.0)(140) + (0.04)(12) + (0.45)(110)} \tag{7.15} \label{eq:unit_II_membrane_transport_item_22}\]

\[V_m = 26.7 \times \ln \frac{5.0 + 5.8 + 1.8}{140 + 0.48 + 49.5} = 26.7 \times \ln \frac{12.6}{190.0} \tag{7.16} \label{eq:unit_II_membrane_transport_item_23}\]

\[V_m = 26.7 \times \ln(0.0663) = 26.7 \times (-2.71) = -72.4 \; \text{mV} \tag{7.17} \label{eq:unit_II_membrane_transport_item_24}\]

This is close to the experimentally measured
\hyperref[gl:resting-potential]{\textbf{resting potential}} of
about -70 mV.

\begin{quote}
\textbf{Concept Check 4:} During an action potential, the permeability
to Na\(^+\) increases about 500-fold (from \(P_{Na}/P_K\) =
0.04 to about 20). Using the Goldman equation, predict what
happens to \(V_m\). Why does the membrane potential approach but not
quite reach \(E_{Na}\)?
\end{quote}

\begin{center}\rule{0.5\linewidth}{0.5pt}\end{center}

\subsection{Membrane Potential and the Action
Potential}\label{membrane-potential-and-the-action-potential}

While the Nernst and Goldman equations describe resting membrane
potential, the \textbf{action potential} is the defining electrical
event in excitable cells (neurons, muscle cells, some endocrine cells).

\subsubsection{Phases of the Action
Potential}\label{phases-of-the-action-potential}

\begin{enumerate}
\def\labelenumi{\arabic{enumi}.}
\item
  \textbf{Resting state} (\(V_m\) about  -70 mV):
  Voltage-gated Na\(^+\) and K\(^+\) channels are closed. K\(^+\) leak
  channels (K2P family) maintain the resting potential near \(E_K\).
\item
  \textbf{\hyperref[gl:depolarisation]{Depolarisation} to threshold}
  (about -55 mV): Graded potentials (e.g., from synaptic input)
  depolarise the membrane. If threshold is reached, voltage-gated
  Na\(^+\) channels (Na\$\_V\$1.x) open rapidly (activation gate).
\item
  \textbf{Rising phase} (depolarisation): Na\(^+\) influx drives \(V_m\)
  toward \(E_{Na}\) (+63 mV). Positive feedback: depolarisation opens
  more Na\(^+\) channels. This is the Hodgkin-Huxley regenerative cycle.
\item
  \textbf{Overshoot:} \(V_m\) briefly exceeds 0 mV (typically reaches
  +30 to +40 mV but does not reach \(E_{Na}\) because Na\(^+\) channel
  inactivation begins).
\item
  \textbf{Repolarisation:} Na\(^+\) channel inactivation (ball-and-chain
  mechanism, h-gate) closes channels within about 1 ms.
  Voltage-gated K\(^+\) channels (K\(_V\), delayed rectifier) open
  slowly, allowing K\(^+\) efflux, driving \(V_m\) back toward \(E_K\).
\item
  \textbf{Undershoot (afterhyperpolarisation):} K\(^+\) channels remain
  open transiently, overshooting below resting potential
  (about  -80 mV). K\(^+\) channels then close, and \(V_m\)
  returns to rest.
\item
  \textbf{Refractory periods:} Absolute refractory period (Na\(^+\)
  channels inactivated; no action potential possible). Relative
  refractory period (some Na\(^+\) channels recovered; stronger stimulus
  needed).
\end{enumerate}

\subsubsection{Hodgkin-Huxley Model}\label{hodgkin-huxley-model}

Hodgkin and Huxley (Nobel Prize 1963) described the action potential
mathematically using voltage-clamp experiments on the squid giant axon:

\[I_m = C_m \frac{dV}{dt} + g_K n^4 (V - E_K) + g_{Na} m^3 h (V - E_{Na}) + g_L (V - E_L) \tag{7.18} \label{eq:unit_II_membrane_transport_item_25}\]

where \(m\) = Na\(^+\) activation variable, \(h\) = Na\(^+\)
inactivation variable, \(n\) = K\(^+\) activation variable. This model
predicted the ionic basis of the action potential before the molecular
identity of ion channels was known.

\subsubsection{Saltatory Conduction}\label{saltatory-conduction}

In myelinated axons, myelin sheaths (formed by Schwann cells in PNS,
oligodendrocytes in CNS) insulate the axon, reducing membrane
capacitance. Action potentials ``jump'' between \textbf{nodes of
Ranvier} (gaps in myelin where Na\(^+\) channels are concentrated at
about 1,000/um\(^2\)). This increases conduction velocity from
about 1 m/s (unmyelinated C fibres) to about 120 m/s
(large myelinated A-alpha fibres).

\begin{quote}
\textbf{Clinical Connection: Multiple Sclerosis and Demyelination}
Multiple sclerosis (MS) is an autoimmune disease in which T cells and
antibodies attack CNS myelin. Demyelination exposes K\(^+\) channels
normally under the myelin sheath, causing K\(^+\) leakage and conduction
block. Symptoms include visual disturbances, motor weakness, and sensory
abnormalities. The drug 4-aminopyridine (dalfampridine) blocks exposed
K\(^+\) channels and partially restores conduction, improving walking
ability in MS patients.
\end{quote}

\begin{quote}
\textbf{Concept Check 6:} Local anaesthetics (e.g., lidocaine) block
voltage-gated Na\(^+\) channels by entering the channel pore from the
intracellular side in their uncharged form, then becoming protonated and
trapped. Why do local anaesthetics preferentially block small-diameter
pain fibres before large motor fibres?
\end{quote}

\begin{center}\rule{0.5\linewidth}{0.5pt}\end{center}

\subsection{Bulk Transport: Endocytosis and
Exocytosis}\label{bulk-transport-endocytosis-and-exocytosis}

\subsubsection{Exocytosis and SNARE-Mediated Membrane
Fusion}\label{exocytosis-and-snare-mediated-membrane-fusion}

\textbf{Exocytosis:} secretory vesicles fuse with the plasma membrane,
releasing contents extracellularly.

\begin{figure}[htbp]
\centering
\includegraphics[width=0.82\linewidth,height=4.2in,keepaspectratio]{../figures/mermaid_inline/inline_mermaid_0040_102315c3d59e.png}
\caption{SNARE-mediated vesicle fusion during regulated exocytosis. The v-SNARE on the vesicle (VAMP/synaptobrevin) and t-SNAREs on the target membrane (syntaxin + SNAP-25) form a tight four-helix bundle that drives membrane fusion. Synaptotagmin acts as the Ca\$\textasciicircum{}\{2+\}\$ sensor that triggers fusion in <1 ms.}
\end{figure}

\begin{itemize}
\tightlist
\item
  \textbf{Synaptic vesicle fusion:} Ca\(^{2+}\) entry through
  voltage-gated Ca\(^{2+}\) channels triggers synaptotagmin, fusion
  occurs in \textless0.2 ms (fastest biological membrane fusion)
\item
  \textbf{Insulin secretion:} glucose metabolism raises ATP/ADP ratio,
  closes K\(_{ATP}\) channels, depolarisation opens Ca\(^{2+}\)
  channels, Ca\(^{2+}\) triggers insulin granule exocytosis
\end{itemize}

\begin{quote}
\textbf{Clinical Connection: Botulinum Toxin and Tetanus Toxin} Both
botulinum toxin (Botox) and tetanus toxin are zinc metalloproteases that
cleave SNARE proteins: - \textbf{Botulinum toxin} (7 serotypes A--G)
cleaves VAMP, SNAP-25, or syntaxin at the neuromuscular junction,
preventing acetylcholine release, causing flaccid paralysis. Clinical
uses: dystonia, spasticity, cosmetic wrinkle reduction. -
\textbf{Tetanus toxin} cleaves VAMP in inhibitory interneurons of the
spinal cord, preventing GABA/glycine release, causing unopposed
excitatory activity and spastic paralysis (lockjaw).
\end{quote}

\subsubsection{SNARE Mechanism in
Detail}\label{snare-mechanism-in-detail}

The SNARE hypothesis (Rothman, Schekman, Südhof --- 2013 Nobel Prize in
Physiology or Medicine) explains how vesicles fuse with target membranes
with both spatial and temporal precision. Each fusion event requires a
quartet of helices contributed by a v-SNARE (on the vesicle) and
t-SNAREs (on the target membrane). The four-helix bundle
``\textbf{zippers}'' from the membrane-distal N-terminus toward the
membrane-proximal C-terminus, releasing about 35 \(k_B T\) of
free energy per SNARE complex --- enough to overcome the kinetic barrier
(about 25 \(k_B T\)) for membrane fusion.

\textbf{Compartment-specific SNARE pairs} ensure that each vesicle fuses
primarily with its correct target:

{\def\LTcaptype{none} % do not increment counter
\begin{longtable}[]{@{}
  >{\raggedright\arraybackslash}p{(\linewidth - 4\tabcolsep) * \real{0.5161}}
  >{\raggedright\arraybackslash}p{(\linewidth - 4\tabcolsep) * \real{0.2258}}
  >{\raggedright\arraybackslash}p{(\linewidth - 4\tabcolsep) * \real{0.2581}}@{}}
\toprule\noalign{}
\begin{minipage}[b]{\linewidth}\raggedright
Compartment pair
\end{minipage} & \begin{minipage}[b]{\linewidth}\raggedright
v-SNARE
\end{minipage} & \begin{minipage}[b]{\linewidth}\raggedright
t-SNAREs
\end{minipage} \\
\midrule\noalign{}
\endhead
\bottomrule\noalign{}
\endlastfoot
Synaptic vesicle \texorpdfstring{\ensuremath{\to}}{->} presynaptic membrane & VAMP2 (synaptobrevin) &
Syntaxin-1 + SNAP-25 \\
ER \texorpdfstring{\ensuremath{\to}}{->} cis-Golgi (COPII) & Sec22 / Bet1 & Syntaxin-5 + Membrin + GS27 \\
Endosome \texorpdfstring{\ensuremath{\to}}{->} late endosome & VAMP7 / VAMP8 & Syntaxin-7 + Vti1b +
Syntaxin-8 \\
Late endosome \texorpdfstring{\ensuremath{\to}}{->} lysosome & VAMP7 & Syntaxin-7 + Vti1b + Syntaxin-8 \\
Golgi \texorpdfstring{\ensuremath{\to}}{->} plasma membrane (constitutive) & VAMP2 / VAMP3 / VAMP4 &
Syntaxin-3/4 + SNAP-23 \\
\end{longtable}
}

\textbf{The five steps of SNARE-mediated fusion:}

\begin{enumerate}
\def\labelenumi{\arabic{enumi}.}
\tightlist
\item
  \textbf{Tethering.} Rab GTPases on the vesicle (e.g., Rab3 on synaptic
  vesicles, Rab1 on ER-derived vesicles) bind their effectors on the
  target membrane (long coiled-coil tethers like p115 or large
  multi-subunit complexes like the exocyst). Tethering is reversible and
  provides about 100 nm initial contact.
\item
  \textbf{Docking.} The four SNARE helices begin to pair through their
  N-terminal regions. SM proteins (Sec1/Munc18 family) chaperone the
  syntaxin partner, holding it in an open conformation. At synapses,
  Munc18 + Munc13 prepare syntaxin-1 for SNARE assembly.
\item
  \textbf{Priming.} Partial zippering brings the vesicle within
  about 5 nm of the target membrane. \textbf{Complexin} binds
  the half-zippered SNARE complex, clamping it in a metastable state
  ready to release on demand. Without complexin, vesicles fuse
  spontaneously; with complexin primarily, they cannot fuse --- both
  must be present for regulated fusion.
\item
  \textbf{Triggering.} At the synapse, action-potential-evoked
  Ca\(^{2+}\) entry through voltage-gated Ca\(^{2+}\) channels reaches
  about 10 \texorpdfstring{\ensuremath{\mu}}{mu}M at the active zone within microseconds.
  Synaptotagmin-1 (Ca\(^{2+}\) sensor with two C2 domains) binds
  Ca\(^{2+}\) and PIP\(_2\), displaces complexin, and releases the
  SNAREs to complete zippering. The whole event takes \textless0.2 ms
  --- the fastest known protein-mediated fusion.
\item
  \textbf{Disassembly.} After fusion, the cis-SNARE complex (the four
  helices now in the same membrane) is dismantled by \textbf{NSF}
  (N-ethylmaleimide-sensitive factor, an AAA+ ATPase) and its adaptor
  \textbf{alpha-SNAP}. This requires ATP hydrolysis (about 1
  ATP per SNARE complex) and recycles the SNAREs for the next round.
\end{enumerate}

\subsubsection{Vesicle Coat Proteins in
Detail}\label{vesicle-coat-proteins-in-detail}

Three major coat systems shape the secretory and endocytic pathways:
\textbf{COPII}, \textbf{COPI}, and \textbf{clathrin}. Each is built
around a small GTPase (Sar1 for COPII, Arf1 for COPI and clathrin) that
anchors to the membrane upon GDP\texorpdfstring{\ensuremath{\to}}{->}GTP exchange and recruits coat
subunits. Membrane curvature is generated by the geometry of the coat
itself.

{\def\LTcaptype{none} % do not increment counter
\begin{longtable}[]{@{}
  >{\raggedright\arraybackslash}p{(\linewidth - 10\tabcolsep) * \real{0.0678}}
  >{\raggedright\arraybackslash}p{(\linewidth - 10\tabcolsep) * \real{0.1017}}
  >{\raggedright\arraybackslash}p{(\linewidth - 10\tabcolsep) * \real{0.1864}}
  >{\raggedright\arraybackslash}p{(\linewidth - 10\tabcolsep) * \real{0.1864}}
  >{\raggedright\arraybackslash}p{(\linewidth - 10\tabcolsep) * \real{0.2542}}
  >{\raggedright\arraybackslash}p{(\linewidth - 10\tabcolsep) * \real{0.2034}}@{}}
\toprule\noalign{}
\begin{minipage}[b]{\linewidth}\raggedright
Coat
\end{minipage} & \begin{minipage}[b]{\linewidth}\raggedright
GTPase
\end{minipage} & \begin{minipage}[b]{\linewidth}\raggedright
Inner layer
\end{minipage} & \begin{minipage}[b]{\linewidth}\raggedright
Outer layer
\end{minipage} & \begin{minipage}[b]{\linewidth}\raggedright
Cargo selection
\end{minipage} & \begin{minipage}[b]{\linewidth}\raggedright
Vesicle size
\end{minipage} \\
\midrule\noalign{}
\endhead
\bottomrule\noalign{}
\endlastfoot
\textbf{COPII} & Sar1 & Sec23/24 (Sec24 is cargo selector via DXE/FF
motifs) & Sec13/31 (cuboctahedral cage) & Anterograde ER\texorpdfstring{\ensuremath{\to}}{->}Golgi & 60--80
nm \\
\textbf{COPI} & Arf1 & \texorpdfstring{\ensuremath{\beta}}{beta}-, \texorpdfstring{\ensuremath{\gamma}}{gamma}-, \texorpdfstring{\ensuremath{\delta}}{delta}-, \texorpdfstring{\ensuremath{\zeta}}{zeta}-COP (cargo selector via KKXX) & \texorpdfstring{\ensuremath{\alpha}}{alpha}-,
\texorpdfstring{\ensuremath{\beta}}{beta}'-, \texorpdfstring{\ensuremath{\varepsilon}}{epsilon}-COP & Retrograde Golgi\texorpdfstring{\ensuremath{\to}}{->}ER, intra-Golgi & 60--100 nm \\
\textbf{Clathrin} & Arf1 (TGN) or none (PM) & Adaptor proteins (AP1 at
TGN, AP2 at PM, AP3 at endosome, AP4) & Clathrin triskelions (3 heavy +
3 light chains) & Selective endocytosis; lysosomal sorting & 80--120
nm \\
\end{longtable}
}

\textbf{Clathrin triskelions} are 3-legged structures that self-assemble
into open polyhedral lattices (hexagons + pentagons, like a soccer
ball). Each clathrin-coated pit contains about 36 triskelions
arranged into a lattice that gradually curves the membrane into a deep
invagination. The vesicle is then released by \textbf{dynamin}, a GTPase
that wraps as a helical collar around the neck and uses GTP hydrolysis
to constrict and pinch off the bud (about 60--120 nm vesicle,
depending on local geometry). After release, the clathrin coat is
disassembled by Hsc70 + auxilin (J-domain co-chaperone) at the cost of
about 1 ATP per triskelion released.

\textbf{Vesicle traffic in numbers.} A typical mammalian cell sustains
about 10\(^4\) exocytic events and about 10\(^4\)
endocytic events per minute. The plasma membrane area equivalent of one
entire cell is internalised every about 30 minutes --- meaning
the membrane is in dynamic flux, with steady-state composition
maintained by precise SNARE-coupled bidirectional traffic. A typical
secretory neuron at full activity can fuse about 1000 synaptic
vesicles per second.

\subsubsection{Endocytosis}\label{endocytosis}

\textbf{Clathrin-mediated endocytosis (CME):} 1. Cargo receptors cluster
in clathrin-coated pits (adaptor AP2 links receptors to clathrin
triskelions) 2. Clathrin assembles into a polyhedral basket, deforming
the membrane into an invagination 3. Dynamin GTPase wraps around the
neck of the invagination and pinches it off (about 60--120 nm
vesicle) 4. Clathrin coat disassembles (uncoated by Hsc70 + auxilin) 5.
Vesicle fuses with early endosome - Examples: LDL receptor, transferrin
receptor, EGF receptor

\textbf{Caveolar endocytosis:} - 50--80 nm flask-shaped invaginations
rich in cholesterol and sphingolipids - Coated with caveolin-1 (integral
membrane protein with hairpin topology) - Functions: transcytosis across
endothelial cells, lipid homeostasis, signalling compartmentalisation

\textbf{Phagocytosis:} - Professional phagocytes (macrophages,
neutrophils, dendritic cells) - Actin-driven pseudopod extension engulfs
large particles (\textgreater0.5 \texorpdfstring{\ensuremath{\mu}}{mu}m): bacteria, dead cells, debris -
Receptors: Fc receptors (opsonised particles), complement receptors,
scavenger receptors, TLRs - Phagosome fuses with lysosomes to form
phagolysosome for degradation

\textbf{Macropinocytosis:} - Non-specific uptake of large volumes of
extracellular fluid - Actin-driven membrane ruffles collapse back onto
the cell surface, trapping fluid in large vesicles (0.2--5 \texorpdfstring{\ensuremath{\mu}}{mu}m) -
Important for antigen sampling by dendritic cells and for nutrient
acquisition by cancer cells (exploited by RAS-mutant tumours)

\begin{quote}
\textbf{Concept Check 5:} Familial hypercholesterolaemia (FH) can be
caused by mutations in the LDL receptor, the adaptor protein ARH, or the
PCSK9 protease. For each, explain how the mutation leads to elevated
blood LDL cholesterol. Which type responds to statin therapy?
\end{quote}

\begin{center}\rule{0.5\linewidth}{0.5pt}\end{center}

\subsection{Computational Bridge}\label{computational-bridge-6}

Nernst potentials for tabulated physiological ions are computed in
closed form:

\begin{Shaded}
\begin{Highlighting}[]
\ImportTok{from}\NormalTok{ biology.cell }\ImportTok{import}\NormalTok{ PHYSIOLOGICAL\_IONS, nernst\_potential}

\ControlFlowTok{for}\NormalTok{ ion }\KeywordTok{in}\NormalTok{ PHYSIOLOGICAL\_IONS:}
    \ControlFlowTok{if}\NormalTok{ ion.charge }\OperatorTok{==} \DecValTok{0}\NormalTok{:}
        \ControlFlowTok{continue}
    \ControlFlowTok{try}\NormalTok{:}
        \BuiltInTok{print}\NormalTok{(ion.ion, }\BuiltInTok{round}\NormalTok{(nernst\_potential(ion), }\DecValTok{2}\NormalTok{), }\StringTok{"mV"}\NormalTok{)}
    \ControlFlowTok{except} \PreprocessorTok{ValueError}\NormalTok{:}
        \ControlFlowTok{pass}
\end{Highlighting}
\end{Shaded}

\begin{quote}
\textbf{Clinical / systems note:} Long QT channelopathies and familial
hyperkalaemic paralysis are reminders that single-ion permeability or
gradient defects reshape the Goldman-style integrative potential you
approximate in silico.
\end{quote}

\begin{center}\rule{0.5\linewidth}{0.5pt}\end{center}

\subsubsection{AlphaFold-Predicted Transporter Structures: Computational
Biology Meets Membrane
Biophysics}\label{alphafold-predicted-transporter-structures-computational-biology-meets-membrane-biophysics}

Membrane transporter structures have historically been among the hardest
to solve experimentally: hydrophobic surfaces resist crystallisation,
cryo-EM requires stable purified protein, and conformational flexibility
(the very thing that makes transporters work) blurs reconstructions.
\textbf{AlphaFold2} (DeepMind, 2021), \textbf{AlphaFold-Multimer}
(2022), and AlphaFold 3 (2024) changed the workflow by making
high-quality structural hypotheses routine for many folded domains and
complexes, including many membrane proteins
\citep{abramson2024alphafold3,varadi2024alphafolddb}.

By 2024, the \textbf{AlphaFold Protein Structure Database} provided
structure coverage for more than 214 million protein sequences in
UniProt \citep{varadi2024alphafolddb}. This transforms downstream
biology: researchers can inspect predicted transmembrane helices and
cavities in an orphan transporter, nominate binding-site residues,
compare paralogues, and design mutagenesis or cryo-EM experiments before
a solved structure exists. The correct lesson is not that prediction
replaces structure determination; it is that prediction changes what
counts as a good first experiment.

Cautions are important for scientific literacy: AlphaFold predicts
likely static conformations, but transporters function by \emph{cycling}
between conformations (outward-open \texorpdfstring{\ensuremath{\leftrightarrow}}{<->} occluded \texorpdfstring{\ensuremath{\leftrightarrow}}{<->} inward-open); a single
model can miss gating, allostery, lipid dependence, and ligand-coupled
state changes. \textbf{Predicted Aligned Error (PAE)} plots flag
low-confidence inter-domain geometries, and \textbf{AlphaFold3}
integrates ligands, ions, nucleic acids, and other biomolecules into
interaction hypotheses \citep{abramson2024alphafold3}. The field now
often generates predicted structures \emph{before} requesting
experimental ones, but strong biological claims still require transport
assays, mutagenesis, binding data, or experimental structures.

\begin{center}\rule{0.5\linewidth}{0.5pt}\end{center}

\subsection{Current Evidence and Frontier
Biology}\label{current-evidence-and-frontier-biology-9}

For \textbf{Membrane Structure and Transport}, frontier biology belongs
inside the evidence logic of the chapter. Cell biology is increasingly
measured as live, spatial, single-cell, and perturbational data rather
than static diagrams alone. The core reading question is this: transport
claims require gradients, permeability, electrochemical driving force,
gating, and energy coupling.

\begin{itemize}
\tightlist
\item
  \textbf{What to verify:} identify the observation, model, assay, or
  dataset that would make the claim stronger or weaker.
\item
  \textbf{What to qualify:} state the scale, organism, cell type,
  environmental condition, or population where the claim is expected to
  hold.
\item
  \textbf{What to compare:} test at least one alternative explanation,
  baseline, or null model before treating the pattern as causal.
\item
  \textbf{What to cite:} distinguish primary evidence, review synthesis,
  public dataset, and institutional guidance; for recent or numeric
  claims, prefer the source closest to the measurement and state what
  has changed since it was published.
\end{itemize}

Ask what measurement scale is being claimed: nanometre structure,
single-cell transcript abundance, organelle dynamics, tissue context, or
organismal phenotype.

\textbf{Source practice:} For cell-state claims, distinguish microscopy,
live-cell perturbation, single-cell sequencing, spatial transcriptomics,
and biochemical assay evidence before making a causal statement.

\subsubsection{Current Evidence Map: Transporter Structure to
Function}\label{current-evidence-map-transporter-structure-to-function}

\begin{figure}[htbp]
\centering
\includegraphics[width=0.82\linewidth,height=4.2in,keepaspectratio]{../figures/mermaid_inline/inline_mermaid_0041_e85a4e71db9a.png}
\caption{Transporter structure claims need functional assays because transport depends on cycling among states, not a single static conformation \textbackslash{}citep\{abramson2024alphafold3,varadi2024alphafolddb\}.}
\end{figure}

\subsection{Summary}\label{summary-9}

\begin{itemize}
\tightlist
\item
  The plasma membrane is a fluid mosaic bilayer with asymmetric lipid
  distribution maintained by flippases. Lipid rafts are dynamic
  cholesterol-sphingolipid microdomains with signalling roles.
\item
  Ion channels achieve selectivity through precise selectivity filters
  (K\(^+\) channel: carbonyl oxygen coordination) and are gated by
  voltage, ligands, mechanical force, or temperature.
\item
  The Na\(^+\)/K\(^+\)-ATPase maintains ionic gradients essential for
  membrane potential, cell volume, and secondary active transport,
  consuming about 25\% of cellular ATP.
\item
  ABC transporters use ATP to transport diverse substrates; clinical
  significance includes multidrug resistance (MDR1) and cystic fibrosis
  (CFTR).
\item
  The Nernst equation gives the equilibrium potential of each ion; the
  Goldman equation gives the actual membrane potential for multiple ions
  simultaneously.
\item
  SNARE proteins mediate membrane fusion in exocytosis; synaptotagmin is
  the Ca\(^{2+}\) sensor for regulated secretion.
\item
  Endocytosis types (clathrin-mediated, caveolar, phagocytosis,
  macropinocytosis) internalise specific cargo or bulk fluid.
\item
  \textbf{Connections:} See \cref{sec:unit_I_water_and_life} for water
  and \hyperref[gl:osmosis]{\textbf{osmosis}}, Unit IX for excitable
  tissues, and \cref{sec:unit_II_cell_signaling} for receptor-mediated
  signalling.
\end{itemize}

\begin{center}\rule{0.5\linewidth}{0.5pt}\end{center}

\subsection{Review Questions}\label{review-questions-9}

\begin{enumerate}
\def\labelenumi{\arabic{enumi}.}
\item
  Derive the Nernst equation from the electrochemical potential
  equation. Calculate \(E_K\) at 20 degrees C (293 K) given
  {[}K\(^+\){]}\(_{\text{in}}\) = 140 mM, {[}K\(^+\){]}\(_{\text{out}}\)
  = 5 mM.
\item
  Explain how the Na\(^+\)/K\(^+\)-ATPase contributes to the resting
  membrane potential both directly (electrogenic pump) and indirectly
  (maintaining K\(^+\) gradient).
\item
  Compare the selectivity mechanisms of K\(^+\) channels and Na\(^+\)
  channels. How does each achieve \textgreater100:1 selectivity for its
  preferred ion?
\item
  A patient is treated with omeprazole for gastric ulcer disease. What
  is the molecular target of this drug, and how does inhibiting it
  reduce gastric acid secretion?
\item
  Explain why MDR1/P-glycoprotein overexpression causes multidrug
  resistance in cancer. Suggest two strategies to overcome this
  resistance.
\item
  Describe the complete SNARE fusion cycle for synaptic vesicle
  exocytosis, from vesicle docking to
  \hyperref[gl:snare-complex]{\textbf{SNARE complex}} disassembly. What
  role does Ca\(^{2+}\) play?
\item
  Using the Goldman equation, predict how the resting membrane potential
  would change if the membrane suddenly became equally permeable to
  Na\(^+\) and K\(^+\).
\item
  Compare clathrin-mediated endocytosis with phagocytosis in terms of
  vesicle size, coat proteins, cytoskeletal involvement, and cell types
  involved.
\item
  Explain the biphasic effect of cholesterol on membrane fluidity. Why
  is this important for cells that experience temperature fluctuations?
\item
  SGLT2 inhibitors cause glucose loss in urine. Calculate the
  approximate caloric loss per day if a diabetic patient loses 70 g of
  glucose per day in urine (glucose: 4 kcal/g).
\item
  Run the Nernst loop in the bridge code and identify which ion's
  equilibrium potential is closest to a typical neuronal resting
  potential; justify using permeability weighting qualitatively.
\item
  Explain why CFTR is called a \textbf{channel} yet sits in the ABC
  transporter superfamily, and how that relates to ATP usage vs.~passive
  flow.
\end{enumerate}

\begin{center}\rule{0.5\linewidth}{0.5pt}\end{center}

\subsection{Further Reading and Source
Notes}\label{further-reading-and-source-notes-9}

\begin{itemize}
\tightlist
\item
  Mitchell (1961). Coupling of phosphorylation to electron and hydrogen
  transfer by a chemi-osmotic type of mechanism. \emph{Nature}, 191.
\item
  Singer \& Nicolson (1972). The fluid mosaic model of the structure of
  cell membranes. \emph{Science}, 175.
\item
  Hodgkin \& Huxley (1952). A quantitative description of membrane
  current and its application to conduction and excitation in nerve.
  \emph{Journal of Physiology}, 117.
\item
  Skou (1957). The influence of some cations on an adenosine
  triphosphatase from peripheral nerves. \emph{Biochimica et Biophysica
  Acta}, 23.
\item
  Agre (2004). Aquaporin water channels (Nobel lecture).
  \emph{Bioscience Reports}, 24.
\item
  Simons \& Ikonen (1997). Functional rafts in cell membranes.
  \emph{Nature}, 387.
\end{itemize}

\begin{center}\rule{0.5\linewidth}{0.5pt}\end{center}

\subsection{Key Terms}\label{key-terms-9}

{\def\LTcaptype{none} % do not increment counter
\begin{longtable}[]{@{}
  >{\raggedright\arraybackslash}p{(\linewidth - 2\tabcolsep) * \real{0.2857}}
  >{\raggedright\arraybackslash}p{(\linewidth - 2\tabcolsep) * \real{0.7143}}@{}}
\toprule\noalign{}
\begin{minipage}[b]{\linewidth}\raggedright
Term
\end{minipage} & \begin{minipage}[b]{\linewidth}\raggedright
Definition
\end{minipage} \\
\midrule\noalign{}
\endhead
\bottomrule\noalign{}
\endlastfoot
\textbf{Fluid mosaic model} & Singer-Nicolson model of the membrane as a
lipid bilayer sea with floating protein icebergs \\
\textbf{Lipid raft} & Dynamic cholesterol/sphingolipid-enriched
microdomain that concentrates signalling proteins \\
\textbf{Flippase} & ATP-dependent enzyme maintaining lipid asymmetry by
moving PS/PE to the inner leaflet \\
\textbf{Aquaporin} & Water-selective channel; excludes protons via
electrostatic mechanism \\
\textbf{GLUT} & Glucose transporter family (SLC2A); 14 isoforms with
tissue-specific expression and \(K_m\) \\
\textbf{Na\(^+\)/K\(^+\)-ATPase} & Primary active pump; 3 Na\(^+\) out /
2 K\(^+\) in per ATP; electrogenic; 25\% of body ATP \\
\textbf{Symport} & Secondary active transport where two solutes move in
the same direction \\
\textbf{Antiport} & Secondary active transport where two solutes move in
opposite directions \\
\textbf{Nernst potential} & Equilibrium potential of a single ion
species: \(E_i = (RT/z_iF) \ln([C]_o/[C]_i)\) \\
\textbf{Goldman equation} & Membrane potential equation accounting for
permeabilities and concentrations of multiple ions \\
\textbf{Selectivity filter} & Narrow region of an ion channel that
determines ion selectivity (e.g., TVGYG in K\(^+\) channels) \\
\textbf{ABC transporter} & ATP-binding cassette transporter superfamily;
drug efflux (MDR1), Cl\(^-\) channel (CFTR) \\
\textbf{SNARE} & Soluble NSF attachment protein receptor; mediates
membrane fusion via four-helix bundle zippering \\
\textbf{Clathrin} & Triskelion coat protein forming polyhedral baskets
during receptor-mediated endocytosis \\
\textbf{Channelopathy} & Disease caused by ion channel mutations (e.g.,
Long QT syndrome, cystic fibrosis) \\
\textbf{Macropinocytosis} & Non-specific bulk uptake of extracellular
fluid via actin-driven membrane ruffling \\
\end{longtable}
}

\begin{center}\rule{0.5\linewidth}{0.5pt}\end{center}

\subsubsection{Companion Source Module}\label{companion-source-module-9}

\textbf{Membrane Structure and Transport} should leave a reproducible
trail from a biological claim to the code, figure, diagram, or
paper-based activity that can test it. Use the surfaces below to inspect
the chapter's assumptions, rerun the relevant model, or compare the
manuscript explanation with companion labs and figures.

{\def\LTcaptype{none} % do not increment counter
\begin{longtable}[]{@{}
  >{\raggedright\arraybackslash}p{(\linewidth - 2\tabcolsep) * \real{0.5000}}
  >{\raggedright\arraybackslash}p{(\linewidth - 2\tabcolsep) * \real{0.5000}}@{}}
\toprule\noalign{}
\begin{minipage}[b]{\linewidth}\raggedright
Surface
\end{minipage} & \begin{minipage}[b]{\linewidth}\raggedright
Use it for
\end{minipage} \\
\midrule\noalign{}
\endhead
\bottomrule\noalign{}
\endlastfoot
\breaktt{src/biology/cell/cell\_biology.py} (\texttt{nernst\_potential},
\texttt{goldman\_equation}, \texttt{osmotic\_pressure},
\texttt{diffusion\_flux}) & Reproduce electrochemical, osmotic, and
diffusive driving forces. \\
\breaktt{src/visualization/plots.py} (\breaktt{plot\_nernst\_potentials})
& Inspect ion-specific gradients and sign conventions. \\
\breaktt{src/mermaid/biology\_diagrams.py}
(\breaktt{membrane\_transport\_diagram}) & Separate channels, carriers,
pumps, and coupled transport. \\
\end{longtable}
}

\textbf{Reproducibility check:} list concentrations, permeability,
charge, temperature, and membrane orientation before interpreting
transport direction. \textbf{Cross-reference:} connect with
\cref{sec:unit_IX_action_potential_synapses} and
\cref{sec:unit_VIII_plant_structure_and_water}.

\newpage

\section{Cell Signalling and
Communication}\label{cell-signalling-and-communication}

\label{sec:unit_II_cell_signaling}

\begin{quote}
\textbf{Ch 8} \texorpdfstring{\ensuremath{\cdot}}{.} Level 3/3 \texorpdfstring{\ensuremath{\cdot}}{.} 55 min read \texorpdfstring{\ensuremath{\cdot}}{.} 75 min lecture \texorpdfstring{\ensuremath{\cdot}}{.}
Prerequisites: \cref{sec:unit_II_membrane_transport},
\cref{sec:unit_I_enzymes_and_kinetics}
\end{quote}

\subsection{Learning Objectives}\label{learning-objectives-10}

\begin{enumerate}
\def\labelenumi{\arabic{enumi}.}
\tightlist
\item
  Classify signalling molecules by their mode of delivery and explain
  signal amplification logic.
\item
  Describe the major families of cell-surface receptors (GPCRs, RTKs,
  ion channels) and their signalling mechanisms.
\item
  Explain the cAMP, PIP\(_2\)/DAG/IP\(_3\), and RTK-MAP kinase
  \hyperref[gl:signal-transduction]{\textbf{signal transduction}}
  cascades in detail.
\item
  Describe additional signalling pathways: Wnt/beta-catenin, JAK-STAT,
  and Notch.
\item
  Explain signal termination mechanisms: GTPase activity,
  phosphodiesterases, phosphatases, and receptor internalisation.
\item
  Describe the stages of \hyperref[gl:mitosis]{\textbf{mitosis}} and the
  role of CDK-cyclin complexes in \hyperref[gl:cell-cycle]{\textbf{cell
  cycle}} control.
\item
  Explain the mechanisms of \hyperref[gl:apoptosis]{\textbf{apoptosis}}
  (intrinsic and extrinsic pathways).
\item
  Explain cancer as dysregulated signalling and evaluate current
  targeted therapies.
\end{enumerate}

\subsubsection{Study Blueprint}\label{study-blueprint-10}

\begin{itemize}
\tightlist
\item
  \textbf{Big idea:} Cells communicate by converting external signals
  into regulated intracellular decisions.
\item
  \textbf{Core concepts:} receptors, second messengers, signal
  amplification, feedback.
\item
  \textbf{Framework alignment:} Vision \& Change: Structure and
  function, Systems, Information flow, exchange, and storage; AP
  Biology: Systems Interactions, Information Storage and Transmission;
  NGSS-style topics: Structure and Function.
\item
  \textbf{Model or quantitative lens:} Ligand-receptor occupancy and
  dose-response reasoning.
\item
  \textbf{Data skill:} Read pathway diagrams and infer the effect of
  agonists, antagonists, or mutations.
\item
  \textbf{Practice cadence:} Visual Representations, Questions and
  Methods, Argumentation.
\item
  \textbf{Common misconception to repair:} A pathway diagram is a causal
  model, not a memorization chart.
\item
  \textbf{Primary lab:} \cref{sec:lab_unit_II_cell_signaling}.
\item
  \textbf{Question bank:} \cref{sec:q_unit_II_cell_signaling}.
\item
  \textbf{Transfer task:} Apply signaling logic to hormones,
  neurotransmitters, immune receptors, or cancer mutations.
\item
  \textbf{Bridge to computation:}
  \breaktt{biology.physiology.physiology.homeostasis\_response}.
\end{itemize}

\begin{center}\rule{0.5\linewidth}{0.5pt}\end{center}

\begin{quote}
\textbf{Opening Vignette: When Signals Go Wrong --- HER2 Breast Cancer}

In 1987, Dennis Slamon's team at UCLA reported that about 25\% of breast
cancers showed amplification of the HER2
\hyperref[gl:gene]{\textbf{gene}}, which encodes a receptor tyrosine
kinase embedded in the plasma membrane (Slamon et al., 1987,
\emph{Science}). In normal cells, HER2 is transiently activated when it
binds its ligand --- triggering a cascade of phosphorylation events that
ultimately drive cell proliferation and survival. In tumours with HER2
amplification, the \hyperref[gl:protein]{\textbf{protein}} is expressed
at 40--100 times the normal level; even without ligand, HER2 molecules
are close enough together to activate each other constantly, delivering
an unrelenting ``grow and divide'' signal.

The consequences of understanding this single signalling pathway have
been transformative. In 1998, the FDA approved trastuzumab (Herceptin),
a monoclonal antibody that binds the extracellular domain of HER2 and
blocks its dimerisation and downstream signalling. In HER2-positive
breast cancer, trastuzumab therapy approximately doubled 5-year
survival. Every step in this therapeutic triumph required understanding
how cells receive, transduce, and amplify molecular signals ---
precisely the subject of this chapter.

\emph{Primary source: Slamon, D. J. et al.~(1987). Human breast cancer:
Correlation of relapse and survival with amplification of the HER-2/neu
oncogene. Science, 235(4785), 177--182.}
\end{quote}

\begin{center}\rule{0.5\linewidth}{0.5pt}\end{center}

Cells receive and interpret thousands of extracellular signals
simultaneously. Signalling systems share a general architecture:

\textbf{Signal (ligand) \texorpdfstring{\ensuremath{\to}}{->} Receptor (detection) \texorpdfstring{\ensuremath{\to}}{->} Transduction cascade
(amplification) \texorpdfstring{\ensuremath{\to}}{->} Effector \texorpdfstring{\ensuremath{\to}}{->} Response \texorpdfstring{\ensuremath{\to}}{->} Termination}

\subsubsection{Types of Cell Signalling by
Distance}\label{types-of-cell-signalling-by-distance}

{\def\LTcaptype{none} % do not increment counter
\begin{longtable}[]{@{}
  >{\raggedright\arraybackslash}p{(\linewidth - 6\tabcolsep) * \real{0.1333}}
  >{\raggedright\arraybackslash}p{(\linewidth - 6\tabcolsep) * \real{0.3667}}
  >{\raggedright\arraybackslash}p{(\linewidth - 6\tabcolsep) * \real{0.2333}}
  >{\raggedright\arraybackslash}p{(\linewidth - 6\tabcolsep) * \real{0.2667}}@{}}
\toprule\noalign{}
\begin{minipage}[b]{\linewidth}\raggedright
Mode
\end{minipage} & \begin{minipage}[b]{\linewidth}\raggedright
Description
\end{minipage} & \begin{minipage}[b]{\linewidth}\raggedright
Example
\end{minipage} & \begin{minipage}[b]{\linewidth}\raggedright
Distance
\end{minipage} \\
\midrule\noalign{}
\endhead
\bottomrule\noalign{}
\endlastfoot
\textbf{Endocrine} & \hyperref[gl:hormone]{\textbf{Hormone}}s in
bloodstream & Insulin, \hyperref[gl:cortisol]{\textbf{cortisol}} &
Metres \\
\textbf{Paracrine} & Local diffusion & Growth factors, prostaglandins &
\textless1 mm \\
\textbf{Autocrine} & Self-stimulation & Tumour self-growth signals, IL-2
in T cells & Same cell \\
\textbf{Juxtacrine} & Membrane-bound ligand/receptor & Notch-Delta,
ephrin-Eph & Cell contact \\
\textbf{Synaptic} & Neurotransmitter across synapse & Glutamate, GABA,
acetylcholine & about 20 nm \\
\end{longtable}
}

\subsubsection{Signal Transduction
Logic}\label{signal-transduction-logic}

Signal transduction cascades exhibit four fundamental properties:

\begin{enumerate}
\def\labelenumi{\arabic{enumi}.}
\item
  \textbf{Amplification:} One activated receptor activates multiple G
  proteins; each G protein activates multiple adenylyl cyclase
  molecules; each adenylyl cyclase produces many cAMP molecules. A
  single epinephrine molecule can trigger release of
  about 10\(^8\) glucose molecules from glycogen (10\(^6\)-fold
  amplification).
\item
  \textbf{Specificity:} Different cell types express different receptor
  subtypes, G proteins, and effectors. Epinephrine causes glycogen
  breakdown in liver (beta\(_2\)-adrenergic receptor, G\(_s\), cAMP) but
  smooth muscle relaxation in bronchi (same receptor, same second
  messenger, different downstream targets).
\item
  \textbf{Integration:} Multiple signals converge on shared effectors. A
  cell's response reflects the integrated sum of active signalling
  pathways. For example, cell proliferation requires simultaneous growth
  factor (RTK), integrin (adhesion), and survival factor signalling.
\item
  \textbf{Adaptation/desensitisation:} Prolonged stimulation leads to
  reduced response. Mechanisms include receptor phosphorylation (by
  GRKs), arrestin binding, receptor internalisation, and negative
  feedback loops \citep{alon2019}.
\end{enumerate}

\subsubsection{Signal Amplification Cascades --- A Quantitative
Treatment}\label{signal-amplification-cascades-a-quantitative-treatment}

Signal amplification is the architectural reason a single hormone
molecule can drive a macroscopic physiological response. Three concepts
make the amplification quantitative: gain per step, overall gain, and
dynamic range.

\textbf{Gain per step.} Each level of a cascade has a gain \(A_i\)
defined as the number of activated downstream molecules per active
upstream molecule (per unit time). For example, a single G\(_s\)\texorpdfstring{\ensuremath{\alpha}}{alpha}-GTP
activates about 10 adenylyl cyclase molecules during its
about 1 minute lifetime; each adenylyl cyclase produces
about 1,000 cAMP per second, so over 30 s the gain at this step
is roughly \(10 \times 30{,}000 = 3 \times 10^5\).

\textbf{Overall gain.} For an \(n\)-step cascade, the total
amplification is the product of step gains:

\begin{equation}
G_\text{total} = \prod_{i=1}^{n} A_i
\label{eq:unit_II_signaling_total_gain}
\end{equation}

This product structure means that \emph{removing one step} of a cascade
cuts gain by orders of magnitude, while \emph{adding one step}
multiplies it. The classic epinephrine\texorpdfstring{\ensuremath{\to}}{->}glycogenolysis cascade has 4
enzymatic amplification steps (receptor \texorpdfstring{\ensuremath{\to}}{->} G\(_s\) \texorpdfstring{\ensuremath{\to}}{->} adenylyl cyclase \texorpdfstring{\ensuremath{\to}}{->}
PKA \texorpdfstring{\ensuremath{\to}}{->} phosphorylase kinase \texorpdfstring{\ensuremath{\to}}{->} glycogen phosphorylase), each contributing
\(A_i \approx 10^{1.5}\), for \(G_\text{total} \approx 10^{6}\). A
single epinephrine molecule binding a \(\beta_2\) receptor releases
about \(10^8\) glucose molecules from glycogen --- exactly
the gain measured in liver perfusion experiments.

\textbf{Why so many steps?} A simpler cascade with the same gain ---
e.g., one step of \(A = 10^6\) --- would be biologically implausible:
few signalling enzymes have \(k_\text{cat} > 10^4\) s\(^{-1}\) (the
diffusion-limited outliers catalase and carbonic anhydrase are
exceptions), and no cell could sustain enzyme concentrations giving
\(10^6\)-fold amplification in a single reaction. The cascade
architecture solves this by \emph{multiplying} modest gains.

\textbf{Time delay and adaptation.} Each step also adds a time constant
\(\tau_i\) (the lifetime of the activated species). The total response
time is:

\[\tau_\text{total} \approx \sqrt{\sum_i \tau_i^2} \tag{8.0} \label{eq:unit_II_cell_signaling_item_1}\]

For visual transduction (rhodopsin \texorpdfstring{\ensuremath{\to}}{->} transducin \texorpdfstring{\ensuremath{\to}}{->} PDE6 \texorpdfstring{\ensuremath{\to}}{->} cGMP fall):
individual \(\tau_i\) are about 10 ms, total response
about 50 ms. For epinephrine: hundreds of milliseconds to
seconds. The cascade thus also serves as a \textbf{temporal filter} ---
fast inputs reach the effector quickly, slow inputs are integrated.

\textbf{Cooperativity sharpens the response.} A linear cascade has gain
but not steepness --- the dose--response is hyperbolic (Hill coefficient
1). To make a switch-like response, cells use: - \emph{Multiple binding
sites} (haemoglobin O\(_2\), \(n_H = 2.8\)). - \emph{Multi-site
phosphorylation requiring full occupancy} (ERK requires both Thr and Tyr
phosphorylation; effective \(n_H \approx 5\) in the MAP kinase cascade).
- \emph{Positive feedback loops} (ERK \texorpdfstring{\ensuremath{\to}}{->} SOS feedback creates
bistability).

The Huang--Ferrell analysis of the MAPK cascade showed mathematically
that \emph{three sequential switches} (each with \(n_H = 1.7\) from its
dual phosphorylation requirement) compose to give an overall Hill
coefficient near 5 --- converting a graded growth-factor input into an
essentially digital ERK output. This is the molecular substrate of the
cell's ``decision making'' between proliferation and quiescence.

\begin{quote}
\textbf{Concept Check 1a:} A six-step cascade has a gain of 10 per step.
What is the total gain? If a single step is removed (now five steps), by
what factor does the total gain change?
\end{quote}

\subsubsection{Quantitative Foundations of Signal
Transduction}\label{quantitative-foundations-of-signal-transduction}

\textbf{Receptor--ligand binding equilibrium.} The dissociation constant
\(K_d\) quantifies receptor affinity:

\[K_d = \frac{[\text{R}][\text{L}]}{[\text{RL}]} \tag{8.1} \label{eq:unit_II_cell_signaling_item_2}\]

where \([\text{R}]\) is free receptor concentration, \([\text{L}]\) is
free ligand concentration, and \([\text{RL}]\) is the receptor--ligand
complex. The fraction of receptors occupied at a given ligand
concentration is:

\[\theta = \frac{[\text{L}]}{[\text{L}] + K_d} \tag{8.2} \label{eq:unit_II_cell_signaling_item_3}\]

When \([\text{L}] = K_d\), exactly 50\% of receptors are occupied.
Typical \(K_d\) values: \hyperref[gl:insulin-receptor]{\textbf{insulin
receptor}} about 0.1 nM; epinephrine--beta\(_2\) receptor
about 1 \texorpdfstring{\ensuremath{\mu}}{mu}M.

\textbf{The Hill equation} models cooperative binding and switch-like
signalling responses:

\[\theta = \frac{[\text{L}]^{n_H}}{K_d^{n_H} + [\text{L}]^{n_H}} \tag{8.3} \label{eq:unit_II_cell_signaling_item_4}\]

where \(n_H\) is the Hill coefficient. \(n_H = 1\) gives a hyperbolic
(Michaelis--Menten-like) curve; \(n_H > 1\) produces a sigmoidal
(switch-like) response. The MAPK cascade achieves an effective
\(n_H \approx 5\), creating the ultrasensitive most-or-none ERK
activation observed experimentally (Huang \& Ferrell, 1996,
\emph{PNAS}).

\textbf{Signal amplification cascade.} If each step in a cascade has
amplification factor \(A_i\), the total amplification is:

\[G_{\text{total}} = \prod_{i=1}^{n} A_i = A_1 \times A_2 \times \cdots \times A_n \tag{8.4} \label{eq:unit_II_cell_signaling_item_5}\]

For the epinephrine--glycogen cascade (\(n = 4\) steps,
\(A_i \approx 10^{1.5}\) per step): \(G_{\text{total}} \approx 10^6\),
explaining how one hormone molecule triggers release of \(\sim 10^8\)
glucose molecules.

\subsubsection{Worked Example: Receptor Fractional
Occupancy}\label{worked-example-receptor-fractional-occupancy}

\textbf{Problem:} The epidermal growth factor receptor (EGFR) is an RTK
with a dissociation constant (\(K_d\)) for its ligand, EGF, of
approximately \(2.0 \text{ nM}\) (\(2.0 \times 10^{-9} \text{ M}\)).
Assuming standard non-cooperative binding behaviour (\(n_H = 1\)),
calculate the fractional occupancy (\texorpdfstring{\ensuremath{\theta}}{theta}) of the EGFR population on a cell
surface when the extracellular concentration of EGF is
\(6.0 \text{ nM}\). What concentration of EGF would be required to
achieve 90\% receptor occupancy?

\textbf{Solution:}

\begin{enumerate}
\def\labelenumi{\arabic{enumi}.}
\item
  \textbf{Calculate occupancy at 6.0 nM:} Using the receptor fractional
  occupancy equation:
  \[ \theta = \frac{[\text{L}]}{[\text{L}] + K_d}  \label{eq:unit_II_cell_signaling_item_6}\]

  Substitute the given values (\([\text{L}] = 6.0 \text{ nM}\),
  \(K_d = 2.0 \text{ nM}\)):
  \[ \theta = \frac{6.0}{6.0 + 2.0} = \frac{6.0}{8.0} = 0.75  \label{eq:unit_II_cell_signaling_item_7}\]

  Thus, \textbf{75\%} of the EGF receptors are bound to ligand.
\item
  \textbf{Calculate the ligand concentration required for 90\%
  occupancy:} Set \(\theta = 0.90\) and solve for \([\text{L}]\):
  \[ 0.90 = \frac{[\text{L}]}{[\text{L}] + 2.0}  \label{eq:unit_II_cell_signaling_item_8}\]

  \[ 0.90([\text{L}] + 2.0) = [\text{L}]  \label{eq:unit_II_cell_signaling_item_9}\]

  \[ 0.90[\text{L}] + 1.80 = [\text{L}]  \label{eq:unit_II_cell_signaling_item_10}\]

  \[ 1.80 = 0.10[\text{L}]  \label{eq:unit_II_cell_signaling_item_11}\]

  \[ [\text{L}] = \frac{1.80}{0.10} = 18.0 \text{ nM}  \label{eq:unit_II_cell_signaling_item_12}\]

  An EGF concentration of \textbf{18.0 nM} is required to saturate 90\%
  of the receptors. Note that achieving \(90\%\) occupancy requires a
  ligand concentration exactly \(9 \times K_d\).
\end{enumerate}

\begin{quote}
\textbf{Concept Check 1:} Epinephrine stimulates both
beta\(_2\)-adrenergic receptors (G\(_s\)-coupled, cAMP elevation) on
bronchial smooth muscle and alpha\(_1\)-adrenergic receptors
(G\(_q\)-coupled, IP\(_3\)/Ca\(^{2+}\)) on vascular smooth muscle. Why
does the same hormone cause bronchodilation in the lungs but
vasoconstriction in skin blood vessels?
\end{quote}

\subsubsection{Worked Example: Signal Amplification Through a Kinase
Cascade}\label{worked-example-signal-amplification-through-a-kinase-cascade}

\textbf{Problem:} The MAP kinase module is a three-tier cascade (Raf
\(\rightarrow\) MEK \(\rightarrow\) ERK) in which each activated kinase
phosphorylates and activates roughly \(10\) molecules of the next kinase
(\(A_1 = A_2 = A_3 = 10\)).

\begin{enumerate}
\def\labelenumi{\arabic{enumi}.}
\tightlist
\item
  Using the cascade gain relation, what is the total amplification of
  the three-tier module?
\item
  If \(50\) growth-factor receptors are activated, how many ERK
  molecules are activated?
\item
  What happens to the total gain if one tier is bypassed (a two-step
  cascade)?
\end{enumerate}

\textbf{Solution:}

\begin{enumerate}
\def\labelenumi{\arabic{enumi}.}
\item
  \textbf{Total gain is the product of the per-step gains:}
  \[ G_{\text{total}} = \prod_{i=1}^{3} A_i = A_1 \times A_2 \times A_3 = 10 \times 10 \times 10 = 10^3  \label{eq:unit_II_cell_signaling_item_13}\]

  The three-tier module amplifies the input \(1{,}000\)-fold.
\item
  \textbf{Scale by the number of activated receptors} (\(N_R = 50\)):
  \[ N_{\text{ERK}} = N_R \times G_{\text{total}} = 50 \times 10^3 = 5.0 \times 10^4  \label{eq:unit_II_cell_signaling_item_14}\]

  So \(50\) activated receptors yield about \(50{,}000\) activated ERK
  molecules.
\item
  \textbf{Drop one tier} (now a two-step cascade,
  \(G = A_1 \times A_2\)):
  \[ G_{\text{2-step}} = 10 \times 10 = 10^2  \label{eq:unit_II_cell_signaling_item_15}\]

  Removing a single tier cuts the gain by a factor of
  \(10^3 / 10^2 = 10\).
\end{enumerate}

Because the cascade gain is a product rather than a sum, each added tier
multiplies the response while a lost tier divides it --- a small number
of detected ligand molecules can therefore drive a large, switch-like
cellular decision, and dropping or inhibiting one kinase tier (a common
drug strategy) sharply attenuates the output.

\begin{center}\rule{0.5\linewidth}{0.5pt}\end{center}

\subsection{Cell-Surface Receptors}\label{cell-surface-receptors}

\subsubsection{G Protein-Coupled Receptors
(GPCRs)}\label{g-protein-coupled-receptors-gpcrs}

GPCRs are the largest receptor superfamily (about 800 in
humans; about 34\% of approved drug targets). They have 7
transmembrane alpha-helices, an extracellular N-terminus, and an
intracellular C-terminus. Binding of ligand activates the heterotrimeric
G protein (G-alpha, G-beta-gamma):

\begin{figure}[htbp]
\centering
\includegraphics[width=0.82\linewidth,height=4.2in,keepaspectratio]{../figures/mermaid_inline/inline_mermaid_0042_a36ac9dc7103.png}
\caption{The G\$\_s\$-cAMP-PKA signalling cascade and its termination mechanisms. Signal amplification occurs at multiple levels: one receptor activates multiple G proteins, each adenylyl cyclase produces about 1,000 cAMP/s, and PKA phosphorylates multiple substrates.}
\end{figure}

\textbf{G\(_s\) pathway (stimulatory):} Hormone (e.g., epinephrine) \texorpdfstring{\ensuremath{\to}}{->}
beta\(_2\)-adrenergic receptor \texorpdfstring{\ensuremath{\to}}{->} G\(_s\) \texorpdfstring{\ensuremath{\to}}{->} \textbf{adenylyl cyclase}
activated \texorpdfstring{\ensuremath{\to}}{->} ATP \texorpdfstring{\ensuremath{\to}}{->} \textbf{cAMP} \texorpdfstring{\ensuremath{\to}}{->} \textbf{Protein Kinase A (PKA)} \texorpdfstring{\ensuremath{\to}}{->}
phosphorylates multiple downstream targets (glycogen phosphorylase
kinase, CREB \hyperref[gl:transcription]{\textbf{transcription}} factor,
ion channels, etc.)

\textbf{G\(_i\) pathway (inhibitory):} Inhibits adenylyl cyclase,
reducing cAMP. Examples: opioid receptors (mu, delta, kappa), muscarinic
M2 receptors (cardiac parasympathetic), alpha\(_2\)-adrenergic
receptors.

\textbf{G\(_q\) pathway:} Activates \textbf{phospholipase C-beta} \texorpdfstring{\ensuremath{\to}}{->}
cleaves PIP\(_2\) \texorpdfstring{\ensuremath{\to}}{->} \textbf{IP\(_3\)} (Ca\(^{2+}\) release from ER via
IP\(_3\) receptors) + \textbf{DAG} (activates Protein Kinase C, PKC).
This dual second messenger system controls smooth muscle contraction,
platelet activation, and secretion.

\textbf{G\(_{12/13}\) pathway:} Activates Rho GEFs \texorpdfstring{\ensuremath{\to}}{->} Rho GTPase \texorpdfstring{\ensuremath{\to}}{->} ROCK
(Rho-associated kinase) \texorpdfstring{\ensuremath{\to}}{->} cytoskeletal remodelling, cell migration,
smooth muscle contraction.

\textbf{cAMP signalling components:} - \textbf{PDE (phosphodiesterase):}
hydrolyses cAMP to 5'-AMP, terminating the signal. Caffeine and
theophylline inhibit PDE \texorpdfstring{\ensuremath{\to}}{->} elevated cAMP \texorpdfstring{\ensuremath{\to}}{->} CNS stimulation,
bronchodilation. Sildenafil (Viagra) inhibits PDE5 (specific for cGMP in
penile smooth muscle). - \textbf{AKAP (A Kinase Anchoring Proteins):}
target PKA to specific subcellular locations, providing spatial
specificity to cAMP signalling. - \textbf{GTPase activity of G-alpha:}
G-alpha is a self-inactivating molecular switch (built-in GTPase;
k\(_{\text{cat}}\) about 1 min\(^{-1}\)). RGS (Regulators of
G-protein Signalling) proteins accelerate GTP hydrolysis, shortening
signal duration.

\subsubsection{The G-Protein Cycle in
Detail}\label{the-g-protein-cycle-in-detail}

Heterotrimeric G proteins (G\texorpdfstring{\ensuremath{\alpha}}{alpha} + G\texorpdfstring{\ensuremath{\beta}}{beta}\texorpdfstring{\ensuremath{\gamma}}{gamma}) are molecular switches that
integrate three biochemical events: GDP release, GTP binding, and GTP
hydrolysis. The cycle is canonical for most GTPases --- Ras, Rho, Rab,
Ran, and the heterotrimeric G proteins most follow the same logic, with
different specifics.

{\def\LTcaptype{none} % do not increment counter
\begin{longtable}[]{@{}
  >{\raggedright\arraybackslash}p{(\linewidth - 6\tabcolsep) * \real{0.1316}}
  >{\raggedright\arraybackslash}p{(\linewidth - 6\tabcolsep) * \real{0.2895}}
  >{\raggedright\arraybackslash}p{(\linewidth - 6\tabcolsep) * \real{0.3158}}
  >{\raggedright\arraybackslash}p{(\linewidth - 6\tabcolsep) * \real{0.2632}}@{}}
\toprule\noalign{}
\begin{minipage}[b]{\linewidth}\raggedright
State
\end{minipage} & \begin{minipage}[b]{\linewidth}\raggedright
Composition
\end{minipage} & \begin{minipage}[b]{\linewidth}\raggedright
Conformation
\end{minipage} & \begin{minipage}[b]{\linewidth}\raggedright
Signalling
\end{minipage} \\
\midrule\noalign{}
\endhead
\bottomrule\noalign{}
\endlastfoot
\textbf{Inactive (GDP)} & G\texorpdfstring{\ensuremath{\alpha}}{alpha}-GDP\texorpdfstring{\ensuremath{\cdot}}{.}G\texorpdfstring{\ensuremath{\beta}}{beta}\texorpdfstring{\ensuremath{\gamma}}{gamma} heterotrimer & ``Off'' & No
effector activation \\
\textbf{Receptor-bound} & G\texorpdfstring{\ensuremath{\alpha}}{alpha}-GDP\texorpdfstring{\ensuremath{\cdot}}{.}G\texorpdfstring{\ensuremath{\beta}}{beta}\texorpdfstring{\ensuremath{\gamma}}{gamma} + activated GPCR & GPCR catalyses
GDP release & Transient \\
\textbf{Empty} & G\texorpdfstring{\ensuremath{\alpha}}{alpha}\texorpdfstring{\ensuremath{\cdot}}{.}G\texorpdfstring{\ensuremath{\beta}}{beta}\texorpdfstring{\ensuremath{\gamma}}{gamma} + GPCR (very brief) & Nucleotide-free & Allows
GTP binding \\
\textbf{Active (GTP)} & G\texorpdfstring{\ensuremath{\alpha}}{alpha}-GTP + free G\texorpdfstring{\ensuremath{\beta}}{beta}\texorpdfstring{\ensuremath{\gamma}}{gamma} & ``On''; conformational
change in switch I/II/III & Activates effectors \\
\textbf{Hydrolysing} & G\texorpdfstring{\ensuremath{\alpha}}{alpha}-GTP intrinsic GTPase & Slow (\(k_\text{cat}\)
about 0.05 s\(^{-1}\), \(t_{1/2}\) about 15 s) &
Self-inactivation \\
\textbf{GAP-accelerated} & G\texorpdfstring{\ensuremath{\alpha}}{alpha}-GTP + RGS protein & Stabilises transition
state & \(t_{1/2}\) shortens to about 1 s \\
\textbf{Reassembly} & G\texorpdfstring{\ensuremath{\alpha}}{alpha}-GDP + G\texorpdfstring{\ensuremath{\beta}}{beta}\texorpdfstring{\ensuremath{\gamma}}{gamma} & Heterotrimer reforms & Cycle
closes \\
\end{longtable}
}

\textbf{Rate constants and cellular consequences.} For G\texorpdfstring{\ensuremath{\alpha}}{alpha}\(_s\), the
intrinsic GTP hydrolysis rate is about 0.05 s\(^{-1}\) (so the
active state lasts about 20 s without acceleration). RGS
proteins (RGS2, RGS4, RGS9) accelerate this 100--1,000\texorpdfstring{\ensuremath{\times}}{x}, shortening the
signal to about 50 ms. Disease example: RGS9 deficiency causes
\textbf{bradyopsia} --- a visual disorder where bright lights produce
after-images that linger because transducin (the rod photoreceptor G\texorpdfstring{\ensuremath{\alpha}}{alpha})
cannot be turned off rapidly.

\textbf{Why the cycle is reversible.} Each step is in principle
reversible at thermodynamic equilibrium. The cycle is driven
\emph{forward} by: - Receptor activation (provides the catalytic step
for GDP release); - High cellular {[}GTP{]}:{[}GDP{]} ratio
(about 10:1, maintained by metabolism); - Coupling of GTP
hydrolysis to effector dissociation (mass-action irreversibility).

Without metabolic GTP, the cycle would stall. ATP-depleted cells
therefore lose most GPCR-mediated signalling --- one reason why
ischaemic damage propagates so quickly in tissues.

\subsubsection{GPCR Desensitisation and
\texorpdfstring{\ensuremath{\beta}}{beta}-Arrestin}\label{gpcr-desensitisation-and-ux3b2-arrestin}

Continuous stimulation of a GPCR rapidly attenuates its response ---
within seconds to minutes, the system \textbf{desensitises}. This is a
textbook example of a homeostatic negative feedback loop and the
mechanistic basis for clinically important phenomena: tachyphylaxis
(rapid loss of drug effect), tolerance (e.g., to opioids, \texorpdfstring{\ensuremath{\beta}}{beta}-agonists),
and receptor downregulation.

\textbf{Three layers of desensitisation:}

\begin{enumerate}
\def\labelenumi{\arabic{enumi}.}
\tightlist
\item
  \textbf{Homologous desensitisation (rapid, agonist-specific;
  seconds).}

  \begin{itemize}
  \tightlist
  \item
    Active GPCR is phosphorylated by \textbf{GRKs (G-protein-coupled
    receptor kinases)}, especially GRK2 and GRK3. GRK2 specifically
    recognises \emph{active} receptor conformation; ligand-free or
    inactive receptors are not phosphorylated.
  \item
    Phosphorylated GPCR recruits \textbf{\texorpdfstring{\ensuremath{\beta}}{beta}-arrestin} (\texorpdfstring{\ensuremath{\beta}}{beta}-arrestin-1 or
    -2). \texorpdfstring{\ensuremath{\beta}}{beta}-arrestin binding sterically blocks G-protein coupling --- the
    receptor is now ``uncoupled'' from G\(_s\)/G\(_q\) but still on the
    cell surface.
  \end{itemize}
\item
  \textbf{Heterologous desensitisation (slower, cross-pathway; seconds
  to minutes).}

  \begin{itemize}
  \tightlist
  \item
    PKA and PKC phosphorylate GPCRs at distinct sites, attenuating
    responses to any agonist that activates that receptor (regardless of
    whether it's currently bound). This integrates signalling across
    pathways.
  \end{itemize}
\item
  \textbf{Receptor internalisation (slow, persistent; minutes to
  hours).}

  \begin{itemize}
  \tightlist
  \item
    \texorpdfstring{\ensuremath{\beta}}{beta}-arrestin recruits clathrin and AP2 \texorpdfstring{\ensuremath{\to}}{->} receptor is internalised in
    clathrin-coated vesicles \texorpdfstring{\ensuremath{\to}}{->} enters the early endosome.
  \item
    From the endosome: receptors can be \textbf{recycled} back to the
    plasma membrane (rapid, restores response within minutes --- the
    path of \texorpdfstring{\ensuremath{\beta}}{beta}\(_2\)-adrenergic receptor) or sorted to \textbf{lysosomes}
    for \textbf{downregulation} (slow, requires hours-to-days for new
    receptor synthesis to restore response --- the path of opioid
    receptors after chronic agonist exposure).
  \end{itemize}
\end{enumerate}

\textbf{Beyond G-protein signalling: \texorpdfstring{\ensuremath{\beta}}{beta}-arrestin as a signalling
scaffold.} Originally thought to ``merely'' terminate G-protein
signalling, \texorpdfstring{\ensuremath{\beta}}{beta}-arrestins are now recognised as scaffolds for their own
kinase cascades --- ERK, JNK, p38 MAPK can be activated by
\texorpdfstring{\ensuremath{\beta}}{beta}-arrestin-bound receptors \emph{after} internalisation, generating a
second wave of signalling distinct from the initial G-protein wave. Some
agonists (called \textbf{biased agonists}) preferentially activate one
pathway: - \textbf{Carvedilol} (\texorpdfstring{\ensuremath{\beta}}{beta}-blocker) inhibits G\(_s\) but
activates \texorpdfstring{\ensuremath{\beta}}{beta}-arrestin signalling --- this dual action may underlie its
superior efficacy in heart failure compared to other \texorpdfstring{\ensuremath{\beta}}{beta}-blockers. -
\textbf{TRV130 / oliceridine} (\texorpdfstring{\ensuremath{\mu}}{mu}-opioid agonist, FDA-approved 2020) is
biased toward G-protein signalling and away from \texorpdfstring{\ensuremath{\beta}}{beta}-arrestin --- this
reduces respiratory depression and constipation while preserving
analgesia.

\begin{quote}
\textbf{Concept Check 1b:} A patient with asthma uses an inhaled
\texorpdfstring{\ensuremath{\beta}}{beta}\(_2\)-agonist (albuterol) every 4 hours for several days. The drug
becomes progressively less effective (tachyphylaxis). Explain at the
molecular level --- which mechanism (homologous, heterologous,
internalisation/downregulation) is most likely responsible for each
time-scale of tolerance?
\end{quote}

\subsubsection{Second Messenger
Kinetics}\label{second-messenger-kinetics}

Second messengers are diffusible small molecules that connect receptors
to downstream effectors. Their effective signalling depends on three
quantitative parameters: production rate, diffusion (and spatial
spread), and degradation rate.

\textbf{cAMP kinetics.} Adenylyl cyclase produces cAMP at
about 1,000 molecules/s per active enzyme. cAMP degradation by
\textbf{phosphodiesterases (PDEs)} has half-lives ranging from
about 50 ms (cardiac muscle, near-membrane PDE3/PDE4
microdomains) to several seconds (cytoplasmic bulk). Diffusion
coefficient \(D \approx 300\) \texorpdfstring{\ensuremath{\mu}}{mu}m\(^2\)/s in cytoplasm; diffusion length
\(\ell = \sqrt{D \tau}\). With \(\tau_{1/2} = 100\) ms,
\(\ell \approx 5\) \texorpdfstring{\ensuremath{\mu}}{mu}m --- short enough to create spatial gradients
within a single cell. AKAP scaffolds and PDE microdomains create ``cAMP
nano-domains'' of just about 100 nm where local concentrations
can be 10--100\texorpdfstring{\ensuremath{\times}}{x} the cell average. This explains how a single cell can
use cAMP to signal \emph{locally} (e.g., individual ion channels)
without flooding the whole cell.

\textbf{The PDE family and pharmacology.} Eleven PDE families exist in
humans, each with tissue-specific expression:

{\def\LTcaptype{none} % do not increment counter
\begin{longtable}[]{@{}
  >{\raggedright\arraybackslash}p{(\linewidth - 8\tabcolsep) * \real{0.1562}}
  >{\raggedright\arraybackslash}p{(\linewidth - 8\tabcolsep) * \real{0.1406}}
  >{\raggedright\arraybackslash}p{(\linewidth - 8\tabcolsep) * \real{0.2656}}
  >{\raggedright\arraybackslash}p{(\linewidth - 8\tabcolsep) * \real{0.2500}}
  >{\raggedright\arraybackslash}p{(\linewidth - 8\tabcolsep) * \real{0.1875}}@{}}
\toprule\noalign{}
\begin{minipage}[b]{\linewidth}\raggedright
PDE family
\end{minipage} & \begin{minipage}[b]{\linewidth}\raggedright
Substrate
\end{minipage} & \begin{minipage}[b]{\linewidth}\raggedright
Tissue / function
\end{minipage} & \begin{minipage}[b]{\linewidth}\raggedright
Inhibitor (drug)
\end{minipage} & \begin{minipage}[b]{\linewidth}\raggedright
Clinical use
\end{minipage} \\
\midrule\noalign{}
\endhead
\bottomrule\noalign{}
\endlastfoot
PDE3 & cAMP \textgreater{} cGMP & Cardiac, vascular smooth muscle &
Milrinone, Cilostazol & Heart failure, claudication \\
PDE4 & cAMP & Inflammatory, CNS & Roflumilast, Apremilast & COPD,
psoriasis \\
\textbf{PDE5} & cGMP primarily & Vascular smooth muscle (penile,
pulmonary) & \textbf{Sildenafil (Viagra)}, Tadalafil & Erectile
dysfunction, pulmonary hypertension \\
PDE6 & cGMP & Photoreceptors & None therapeutic & (Mutated in retinitis
pigmentosa) \\
PDE9, PDE10 & cAMP, cGMP & Brain, smooth muscle & Various & Cognitive
disorders (trial) \\
\end{longtable}
}

\textbf{The Viagra story (a quintessential mechanism-to-medicine path).}
Pfizer scientists were originally testing sildenafil in 1989 as a
treatment for angina (the rationale: PDE5 inhibition \texorpdfstring{\ensuremath{\to}}{->} cGMP buildup \texorpdfstring{\ensuremath{\to}}{->}
vasodilation \texorpdfstring{\ensuremath{\to}}{->} coronary blood flow). Its anti-anginal effects were
modest, but male patients in trials reported a striking side effect.
Pfizer pivoted; sildenafil was approved for erectile dysfunction in 1998
and became a major pharmaceutical product. The mechanism is exquisitely
specific: nitric oxide \texorpdfstring{\ensuremath{\to}}{->} guanylyl cyclase \texorpdfstring{\ensuremath{\to}}{->} cGMP \texorpdfstring{\ensuremath{\to}}{->} PKG \texorpdfstring{\ensuremath{\to}}{->} relaxation of
penile smooth muscle. PDE5 normally degrades cGMP; sildenafil blocks
PDE5; cGMP accumulates \emph{primarily where it is being made} (i.e.,
where NO is released). This ``context-conditional'' pharmacology ---
drug effect primarily where the natural signal is active --- now informs
many other targeted therapies. Tadalafil (Cialis) has a longer half-life
(17.5 h vs.~4 h for sildenafil) due to its different binding kinetics in
the PDE5 active site.

\textbf{IP\(_3\) / Ca\(^{2+}\) kinetics.} IP\(_3\) has a half-life of
about 1 second (rapid dephosphorylation by
IP\(_3\)-5-phosphatase); Ca\(^{2+}\) signals are buffered by
about 99\% by cytoplasmic proteins (calbindin, calretinin) and
re-sequestered by SERCA in \textless1 second. Therefore Ca\(^{2+}\)
``spikes'' and ``waves'' propagate at controllable, intermediate speeds
(about 10--50 \texorpdfstring{\ensuremath{\mu}}{mu}m/s), enabling spatially organised signalling
such as the fertilisation Ca\(^{2+}\) wave that sweeps a sea urchin egg
in about 15 seconds.

\begin{figure}[htbp]
\centering
\includegraphics[width=0.82\linewidth,height=4.2in,keepaspectratio]{../figures/mermaid_inline/inline_mermaid_0043_2ee92fbbe7c5.png}
\caption{Major second messenger systems and their downstream effects. Each second messenger activates specific protein kinases or effectors that mediate distinct cellular responses.}
\end{figure}

\begin{quote}
\textbf{Clinical Connection: Cholera and Pertussis Toxins}
\textbf{Cholera toxin} ADP-ribosylates G\(_s\)-alpha, preventing GTP
hydrolysis. G\(_s\) is constitutively active \texorpdfstring{\ensuremath{\to}}{->} massive cAMP production
in intestinal epithelial cells \texorpdfstring{\ensuremath{\to}}{->} CFTR Cl\(^-\) channels open \texorpdfstring{\ensuremath{\to}}{->} Cl\(^-\)
and water secretion \texorpdfstring{\ensuremath{\to}}{->} severe watery diarrhoea (up to 20 L/day),
potentially fatal from dehydration. \textbf{Pertussis toxin} (whooping
cough) ADP-ribosylates G\(_i\)-alpha, preventing its activation. Without
G\(_i\) inhibition, adenylyl cyclase is overactive \texorpdfstring{\ensuremath{\to}}{->} elevated cAMP in
respiratory epithelial cells \texorpdfstring{\ensuremath{\to}}{->} impaired mucociliary clearance \texorpdfstring{\ensuremath{\to}}{->}
persistent cough. see \cref{sec:unit_II_membrane_transport} for CFTR.
\end{quote}

\subsubsection{Receptor Tyrosine Kinases
(RTKs)}\label{receptor-tyrosine-kinases-rtks}

RTKs are single-pass transmembrane receptors with a cytoplasmic tyrosine
kinase domain. Ligand binding (e.g., EGF, PDGF, insulin, FGF, VEGF) \texorpdfstring{\ensuremath{\to}}{->}
receptor \textbf{dimerisation} \texorpdfstring{\ensuremath{\to}}{->} \textbf{trans-autophosphorylation} of
tyrosines \texorpdfstring{\ensuremath{\to}}{->} phosphotyrosines recruit adaptor proteins via SH2 and PTB
domains.

\begin{figure}[htbp]
\centering
\includegraphics[width=0.82\linewidth,height=4.2in,keepaspectratio]{../figures/mermaid_inline/inline_mermaid_0044_a62d4ee8a5da.png}
\caption{The RTK/RAS/MAPK signalling cascade. Each step amplifies the signal: one activated EGFR can activate about 100 RAS molecules, producing about 10,000 activated ERK molecules.}
\end{figure}

Each step amplifies: one activated EGFR \texorpdfstring{\ensuremath{\to}}{->} about 100 RAS \texorpdfstring{\ensuremath{\to}}{->}
about 10,000 activated ERK molecules.

\subsubsection{Receptor Tyrosine Kinase Activation in
Detail}\label{receptor-tyrosine-kinase-activation-in-detail}

RTKs are activated by \textbf{ligand-induced dimerisation} --- and the
choreography is more elegant than the cartoon suggests. Different RTK
families use different dimerisation strategies:

\begin{enumerate}
\def\labelenumi{\arabic{enumi}.}
\tightlist
\item
  \textbf{Bivalent ligand} (e.g., PDGF, VEGF, SCF): the ligand is a
  dimer that bridges two receptor monomers. Binding stoichiometry is 2
  ligands : 2 receptors.
\item
  \textbf{Receptor-mediated dimerisation} (e.g., EGFR, HER2/3/4): the
  ligand is monomeric but binding induces a conformational change that
  exposes a ``dimerisation arm'' on the extracellular domain; two
  ligand-bound receptors then pair through these arms. Stoichiometry: 2
  ligands : 2 receptors.
\item
  \textbf{Pre-formed dimers} (e.g., insulin receptor, IGF-1R): exist as
  covalent (\texorpdfstring{\ensuremath{\alpha}}{alpha}\(_2\)\texorpdfstring{\ensuremath{\beta}}{beta}\(_2\)) tetramers held together by disulphide bonds.
  Ligand binding does not change quaternary structure but does induce a
  conformational change that brings the two intracellular kinase domains
  into productive juxtaposition.
\end{enumerate}

\textbf{Trans-autophosphorylation.} Upon dimerisation, the two
intracellular kinase domains phosphorylate each other on multiple
tyrosine residues in trans (kinase A phosphorylates kinase B and vice
versa). This serves three purposes: - Increases kinase activity
about 100-fold by stabilising the active conformation. -
Creates docking sites for downstream adaptor proteins. - Sets a
stoichiometric ``code'' --- different receptors phosphorylate different
tyrosine combinations, recruiting different adaptors.

\textbf{SH2 and PTB domains: phosphotyrosine readers.} Phosphorylated
tyrosines are recognised by two protein-domain types: - \textbf{SH2 (Src
homology 2)} domains, about 100 amino acids, recognise pTyr in
a sequence-specific context (the +1 to +3 residues C-terminal to pTyr
determine specificity). Each SH2 domain has a characteristic preference
(e.g., Grb2 SH2 prefers pYXNX; PI3K-p85 SH2 prefers pYXXM; Src family
SH2 prefers pYEEI). - \textbf{PTB (Phosphotyrosine Binding)} domains
recognise the residues \emph{N-terminal} to pTyr. Examples: Shc, IRS-1.

The combinatorial code of pTyr sites and SH2/PTB readers explains how
the same kinase domain (intrinsically not sequence-specific) can produce
highly specific signalling outputs.

\textbf{Specific RTK examples:}

{\def\LTcaptype{none} % do not increment counter
\begin{longtable}[]{@{}
  >{\raggedright\arraybackslash}p{(\linewidth - 6\tabcolsep) * \real{0.0625}}
  >{\raggedright\arraybackslash}p{(\linewidth - 6\tabcolsep) * \real{0.1250}}
  >{\raggedright\arraybackslash}p{(\linewidth - 6\tabcolsep) * \real{0.3542}}
  >{\raggedright\arraybackslash}p{(\linewidth - 6\tabcolsep) * \real{0.4583}}@{}}
\toprule\noalign{}
\begin{minipage}[b]{\linewidth}\raggedright
RTK
\end{minipage} & \begin{minipage}[b]{\linewidth}\raggedright
Ligand
\end{minipage} & \begin{minipage}[b]{\linewidth}\raggedright
Dimerisation type
\end{minipage} & \begin{minipage}[b]{\linewidth}\raggedright
Key adaptors recruited
\end{minipage} \\
\midrule\noalign{}
\endhead
\bottomrule\noalign{}
\endlastfoot
EGFR (HER1) & EGF, TGF-\texorpdfstring{\ensuremath{\alpha}}{alpha} & Receptor-mediated & Grb2 (\texorpdfstring{\ensuremath{\to}}{->} RAS); PLC-\texorpdfstring{\ensuremath{\gamma}}{gamma} (\texorpdfstring{\ensuremath{\to}}{->}
DAG/IP\(_3\)); STAT3 \\
HER2 & None known & Heterodimerises with HER1/3/4 & Same as above;
HER2-HER3 strongest mitogen \\
Insulin R & Insulin & Pre-formed & IRS-1/2 (\texorpdfstring{\ensuremath{\to}}{->} PI3K, Grb2); Shc \\
FGFR & FGF + heparan sulfate & Bivalent ligand + co-receptor & FRS2 (\texorpdfstring{\ensuremath{\to}}{->}
Grb2, Shp2); PLC-\texorpdfstring{\ensuremath{\gamma}}{gamma} \\
VEGFR2 & VEGF-A & Bivalent ligand & TSAd; PLC-\texorpdfstring{\ensuremath{\gamma}}{gamma}; Shc \\
Trk receptors & NGF, BDNF & Bivalent (TrkA: 1 NGF dimer / 2 receptors) &
Shc; PLC-\texorpdfstring{\ensuremath{\gamma}}{gamma}; FRS2 \\
\end{longtable}
}

\begin{quote}
\textbf{Clinical Connection: HER2 and Trastuzumab Reconsidered.}
Trastuzumab binds the membrane-proximal extracellular subdomain IV of
HER2 --- \emph{not} the dimerisation arm. So how does it work? Multiple
mechanisms: (1) it sterically interferes with HER2-HER2
self-dimerisation in amplified cancers; (2) it triggers
\textbf{antibody-dependent cellular cytotoxicity (ADCC)} by recruiting
NK cells via Fc receptors; (3) it accelerates HER2 internalisation and
degradation. Trastuzumab-resistant HER2\(^+\) tumours often retain ADCC
sensitivity, which is why \textbf{trastuzumab-emtansine (T-DM1)} --- an
antibody--drug conjugate carrying a microtubule poison --- works after
trastuzumab failure: it exploits residual ADCC and adds direct
cytotoxicity.
\end{quote}

\subsubsection{The MAP Kinase Cascade --- Three-Tier
Hierarchy}\label{the-map-kinase-cascade-three-tier-hierarchy}

The MAP kinase (MAPK) cascade is the canonical example of a multi-tier
signalling cascade. Eukaryotes have multiple parallel MAPK cascades; the
best-studied is the \textbf{classical (ERK) MAPK cascade}, but JNK and
p38 stress-activated cascades follow the same architecture:

\begin{enumerate}
\def\labelenumi{\arabic{enumi}.}
\tightlist
\item
  \textbf{MAPKKK (MAP kinase kinase kinase)} --- activated by upstream
  signals (e.g., Ras-GTP for ERK; MEKK1 for JNK; ASK1 for p38).
\item
  \textbf{MAPKK (MAP kinase kinase)} --- phosphorylated and activated by
  MAPKKK on two serines/threonines; in turn, dual-specificity kinase
  (Ser/Thr/Tyr) for the next tier.
\item
  \textbf{MAPK} --- phosphorylated by MAPKK on Thr-X-Tyr motif in the
  activation loop; both phosphorylations required for activity.
\end{enumerate}

{\def\LTcaptype{none} % do not increment counter
\begin{longtable}[]{@{}
  >{\raggedright\arraybackslash}p{(\linewidth - 10\tabcolsep) * \real{0.1750}}
  >{\raggedright\arraybackslash}p{(\linewidth - 10\tabcolsep) * \real{0.1500}}
  >{\raggedright\arraybackslash}p{(\linewidth - 10\tabcolsep) * \real{0.1250}}
  >{\raggedright\arraybackslash}p{(\linewidth - 10\tabcolsep) * \real{0.1000}}
  >{\raggedright\arraybackslash}p{(\linewidth - 10\tabcolsep) * \real{0.3000}}
  >{\raggedright\arraybackslash}p{(\linewidth - 10\tabcolsep) * \real{0.1500}}@{}}
\toprule\noalign{}
\begin{minipage}[b]{\linewidth}\raggedright
Cascade
\end{minipage} & \begin{minipage}[b]{\linewidth}\raggedright
MAPKKK
\end{minipage} & \begin{minipage}[b]{\linewidth}\raggedright
MAPKK
\end{minipage} & \begin{minipage}[b]{\linewidth}\raggedright
MAPK
\end{minipage} & \begin{minipage}[b]{\linewidth}\raggedright
Activated by
\end{minipage} & \begin{minipage}[b]{\linewidth}\raggedright
Output
\end{minipage} \\
\midrule\noalign{}
\endhead
\bottomrule\noalign{}
\endlastfoot
Classical (ERK) & RAF (A-, B-, C-RAF) & MEK1/2 & ERK1/2 & Growth factors
(RTKs); RAS-GTP & Proliferation, differentiation \\
JNK & MEKK1, MLK1-3, ASK1 & MKK4, MKK7 & JNK1/2/3 & Cytokines, UV,
stress & Stress response, apoptosis \\
p38 & MTK1, MEKK3, ASK1 & MKK3, MKK6 & p38 (\texorpdfstring{\ensuremath{\alpha}}{alpha}, \texorpdfstring{\ensuremath{\beta}}{beta}, \texorpdfstring{\ensuremath{\gamma}}{gamma}, \texorpdfstring{\ensuremath{\delta}}{delta}) & Inflammation,
osmotic stress & Cytokine production, cell-cycle arrest \\
ERK5 & MEKK2, MEKK3 & MEK5 & ERK5 & Growth factors, oxidative stress &
Proliferation, vascular development \\
\end{longtable}
}

\textbf{Why three tiers?} The architecture provides: -
\textbf{Amplification} --- three sequential gain stages (each
about 10\texorpdfstring{\ensuremath{\times}}{x}) \texorpdfstring{\ensuremath{\to}}{->} 1000\texorpdfstring{\ensuremath{\times}}{x} total amplification. -
\textbf{Ultrasensitivity} --- dual-phosphorylation requirement at the
MAPK tier creates an effective \(n_H \approx 5\) at the output,
converting graded input into switch-like ERK activity (Huang \& Ferrell
1996). - \textbf{Insulation by scaffolds} --- KSR1 binds RAF, MEK1, and
ERK1; this prevents undesired cross-activation between parallel MAPK
cascades that share components. Yeast pheromone-response cascade uses
\textbf{Ste5} scaffold; mammals use KSR1 (ERK), JIP1 (JNK), OSM (p38). -
\textbf{Negative feedback by specific phosphatases} --- DUSP1-6
(dual-specificity phosphatases / MKPs) selectively dephosphorylate
specific MAPKs; their expression is itself induced by ERK, creating a
delayed negative feedback that produces transient ERK activation.

The combination of cascade ultrasensitivity + scaffold insulation + DUSP
feedback is why activation of ``the same RAS-MAPK pathway'' can produce
\emph{opposite} outcomes depending on context: transient activation \texorpdfstring{\ensuremath{\to}}{->}
proliferation (ERK off before cell-cycle arrest); sustained activation \texorpdfstring{\ensuremath{\to}}{->}
differentiation (ERK on long enough to upregulate p21 and exit the
cycle); even higher activation \texorpdfstring{\ensuremath{\to}}{->} senescence (ERK-driven oncogene-induced
senescence). Cell fate is encoded in the \emph{temporal pattern} of ERK,
not just its peak amplitude.

\textbf{Key pathway components:}

{\def\LTcaptype{none} % do not increment counter
\begin{longtable}[]{@{}
  >{\raggedright\arraybackslash}p{(\linewidth - 4\tabcolsep) * \real{0.2143}}
  >{\raggedright\arraybackslash}p{(\linewidth - 4\tabcolsep) * \real{0.1905}}
  >{\raggedright\arraybackslash}p{(\linewidth - 4\tabcolsep) * \real{0.5952}}@{}}
\toprule\noalign{}
\begin{minipage}[b]{\linewidth}\raggedright
Component
\end{minipage} & \begin{minipage}[b]{\linewidth}\raggedright
Function
\end{minipage} & \begin{minipage}[b]{\linewidth}\raggedright
Human disease when mutated
\end{minipage} \\
\midrule\noalign{}
\endhead
\bottomrule\noalign{}
\endlastfoot
EGFR (ErbB1) & RTK; EGF binding & Non-small cell lung cancer \\
HER2 (ErbB2) & RTK; no known ligand; heterodimerises & Breast cancer
(amplified in about 25\%) \\
RAS (K-RAS, H-RAS, N-RAS) & Small GTPase; molecular switch & Pancreatic
(about 95\%), colorectal (about 40\%), lung
(about 30\%) \\
RAF (B-RAF) & Ser/Thr kinase (MAPKKK) & Melanoma (V600E,
about 50\%), thyroid cancer \\
NF1 (neurofibromin) & RAS-GAP (turns RAS off) & Neurofibromatosis type
1 \\
PTEN & PIP\(_3\) phosphatase (opposes PI3K) & Glioblastoma, prostate,
endometrial cancer \\
\end{longtable}
}

\subsubsection{Ligand-Gated Ion
Channels}\label{ligand-gated-ion-channels}

Direct coupling of ligand binding to ion flow, bypassing G proteins.
Fastest signalling (about  ms).

\begin{itemize}
\tightlist
\item
  \textbf{Nicotinic acetylcholine receptor (nAChR):} pentameric
  (alpha\(_2\)-beta-gamma-delta); 2 ACh bind \texorpdfstring{\ensuremath{\to}}{->} conformational change \texorpdfstring{\ensuremath{\to}}{->}
  Na\(^+\)/K\(^+\) channel opens \texorpdfstring{\ensuremath{\to}}{->}
  \hyperref[gl:depolarisation]{\textbf{depolarisation}} \texorpdfstring{\ensuremath{\to}}{->} muscle
  contraction
\item
  \textbf{GABA\(_A\) receptor:} Cl\(^-\) channel; GABA binding \texorpdfstring{\ensuremath{\to}}{->}
  Cl\(^-\) influx \texorpdfstring{\ensuremath{\to}}{->} hyperpolarisation \texorpdfstring{\ensuremath{\to}}{->} inhibition. Target of
  benzodiazepines (bind \hyperref[gl:allosteric]{\textbf{allosteric}}
  site \texorpdfstring{\ensuremath{\to}}{->} enhanced Cl\(^-\) flux \texorpdfstring{\ensuremath{\to}}{->} sedation/anxiolysis). Also target of
  barbiturates, ethanol, and general anaesthetics.
\item
  \textbf{NMDA glutamate receptor:} requires simultaneous glutamate +
  glycine binding AND membrane depolarisation to remove Mg\(^{2+}\)
  block; Ca\(^{2+}\) entry \texorpdfstring{\ensuremath{\to}}{->} CaMKII activation \texorpdfstring{\ensuremath{\to}}{->} LTP (Learning and
  memory). Coincidence detector for Hebbian learning.
\end{itemize}

\begin{quote}
\textbf{Concept Check 2:} Curare blocks nAChRs at the neuromuscular
junction. Predict the clinical effects. How does this differ from the
effect of an acetylcholinesterase inhibitor (e.g., neostigmine)?
\end{quote}

\begin{center}\rule{0.5\linewidth}{0.5pt}\end{center}

\subsection{Intracellular Receptors}\label{intracellular-receptors}

Lipophilic signalling molecules (steroid hormones, thyroid hormone,
retinoic acid, vitamin D) diffuse through the membrane and bind
\textbf{cytoplasmic or nuclear receptors} --- ligand-activated
transcription factors. The nuclear receptor superfamily includes 48
members in humans.

\textbf{Classic glucocorticoid pathway:} 1. Cortisol crosses membrane
(lipophilic) 2. Binds glucocorticoid receptor (GR) in
\hyperref[gl:cytoplasm]{\textbf{cytoplasm}} 3. Hsp90 dissociates; GR
dimerises 4. Translocates to nucleus via NLS 5. Binds glucocorticoid
response elements (GREs) in DNA (consensus: AGAACAnnnTGTTCT) 6. Recruits
coactivators (SRC-1, p300/CBP) \texorpdfstring{\ensuremath{\to}}{->} transcriptional activation

Target genes include anti-inflammatory (IkB-alpha, annexin A1),
gluconeogenic (PEPCK, G6Pase), and immunosuppressive genes.
Dexamethasone (synthetic glucocorticoid) is widely used for inflammation
and autoimmune disease.

\textbf{Other nuclear receptors:} - \textbf{Estrogen receptor (ER-alpha,
ER-beta):} breast development, bone density; tamoxifen (antagonist) for
breast cancer - \textbf{PPAR-gamma:} lipid metabolism, adipocyte
differentiation; thiazolidinediones (rosiglitazone) for type 2 diabetes
- \textbf{Thyroid hormone receptor (TR):} metabolic rate; T3 binding
activates genes for thermogenesis and development - \textbf{Vitamin D
receptor (VDR):} calcium
\hyperref[gl:homeostasis]{\textbf{homeostasis}}; calcitriol binding
activates genes for Ca\(^{2+}\) absorption

\begin{center}\rule{0.5\linewidth}{0.5pt}\end{center}

\subsection{Additional Signalling
Pathways}\label{additional-signalling-pathways}

\subsubsection{Wnt/Beta-Catenin Pathway}\label{wntbeta-catenin-pathway}

\textbf{Without Wnt signal:} Beta-catenin is phosphorylated by the
``destruction complex'' (APC + Axin + GSK-3-beta + CK1-alpha) \texorpdfstring{\ensuremath{\to}}{->}
ubiquitinated by beta-TrCP \texorpdfstring{\ensuremath{\to}}{->} proteasomal degradation. Target genes are
off.

\textbf{With Wnt signal:} Wnt binds Frizzled receptor + LRP5/6
co-receptor \texorpdfstring{\ensuremath{\to}}{->} Dishevelled recruited \texorpdfstring{\ensuremath{\to}}{->} destruction complex inhibited
(Axin sequestered to LRP6) \texorpdfstring{\ensuremath{\to}}{->} beta-catenin accumulates \texorpdfstring{\ensuremath{\to}}{->} enters nucleus \texorpdfstring{\ensuremath{\to}}{->}
binds TCF/LEF transcription factors \texorpdfstring{\ensuremath{\to}}{->} activates target genes (c-Myc,
cyclin D1, Axin2).

\textbf{Clinical significance:} APC
\hyperref[gl:mutation]{\textbf{mutation}}s (loss of destruction complex
function) are found in \textgreater80\% of colorectal cancers (familial
adenomatous polyposis, FAP). Beta-catenin is also critical for stem cell
maintenance in the intestinal crypt.

\subsubsection{JAK-STAT Pathway}\label{jak-stat-pathway}

Cytokine receptors (e.g., IL-6R, IFN receptors, erythropoietin receptor)
lack intrinsic kinase activity. Instead, they associate with
\textbf{Janus kinases (JAK1, JAK2, JAK3, TYK2)}:

\begin{enumerate}
\def\labelenumi{\arabic{enumi}.}
\tightlist
\item
  Cytokine binding \texorpdfstring{\ensuremath{\to}}{->} receptor dimerisation
\item
  JAKs trans-phosphorylate each other and the receptor
\item
  \textbf{STAT} proteins (Signal Transducers and Activators of
  Transcription) bind phosphotyrosines via SH2 domains
\item
  JAKs phosphorylate STATs
\item
  Phospho-STATs dimerise and translocate to the nucleus
\item
  STAT dimers bind DNA \texorpdfstring{\ensuremath{\to}}{->} activate target genes
\end{enumerate}

\textbf{Clinical significance:} - JAK2 V617F mutation: found in
\textgreater95\% of polycythaemia vera (constitutive JAK2 activation \texorpdfstring{\ensuremath{\to}}{->}
excess red blood cell production) - JAK inhibitors (ruxolitinib,
tofacitinib) are used for myeloproliferative disorders and rheumatoid
arthritis - STAT3 is constitutively active in many cancers \texorpdfstring{\ensuremath{\to}}{->} survival
and proliferation

\subsubsection{Notch Pathway}\label{notch-pathway}

Direct cell-cell signalling (juxtacrine): 1. Delta/Jagged ligand on one
cell binds Notch receptor on adjacent cell 2. Ligand
\hyperref[gl:endocytosis]{\textbf{endocytosis}} generates mechanical
force 3. ADAM protease cleaves Notch extracellular domain 4.
Gamma-secretase cleaves transmembrane domain, releasing \textbf{NICD}
(Notch intracellular domain) 5. NICD translocates to nucleus \texorpdfstring{\ensuremath{\to}}{->} binds CSL
(RBP-J) transcription factor \texorpdfstring{\ensuremath{\to}}{->} activates Hes/Hey target genes

\textbf{Clinical significance:} Notch1 activating mutations occur in
about 60\% of T-cell acute lymphoblastic leukaemia.
Gamma-secretase inhibitors are in clinical trials.

\subsubsection{Signal Termination
Mechanisms}\label{signal-termination-mechanisms}

{\def\LTcaptype{none} % do not increment counter
\begin{longtable}[]{@{}
  >{\raggedright\arraybackslash}p{(\linewidth - 4\tabcolsep) * \real{0.4091}}
  >{\raggedright\arraybackslash}p{(\linewidth - 4\tabcolsep) * \real{0.2727}}
  >{\raggedright\arraybackslash}p{(\linewidth - 4\tabcolsep) * \real{0.3182}}@{}}
\toprule\noalign{}
\begin{minipage}[b]{\linewidth}\raggedright
Mechanism
\end{minipage} & \begin{minipage}[b]{\linewidth}\raggedright
Target
\end{minipage} & \begin{minipage}[b]{\linewidth}\raggedright
Example
\end{minipage} \\
\midrule\noalign{}
\endhead
\bottomrule\noalign{}
\endlastfoot
GTPase activity & G proteins, RAS & Intrinsic GTPase + GAPs (RGS,
NF1) \\
Phosphodiesterases & cAMP, cGMP & PDE4 (lung), PDE5 (smooth muscle) \\
Protein phosphatases & Phosphoproteins & PP1, PP2A, calcineurin, PTEN,
MKPs \\
Receptor internalisation & Surface receptors & GRK/beta-arrestin \texorpdfstring{\ensuremath{\to}}{->}
clathrin-mediated endocytosis \\
Ubiquitin-proteasome & Signalling proteins & c-Cbl ubiquitinates EGFR;
SCF targets beta-catenin \\
Negative feedback & Pathway components & SOCS proteins inhibit JAK-STAT;
Sprouty inhibits RAS-MAPK \\
\end{longtable}
}

\begin{quote}
\textbf{Concept Check 3:} A mutation in NF1 (neurofibromin, a RAS-GAP)
causes neurofibromatosis type 1. Explain why loss of a GAP protein leads
to tumour formation, and predict the effect on RAS-GTP levels.
\end{quote}

\begin{center}\rule{0.5\linewidth}{0.5pt}\end{center}

\subsection{Crosstalk and Signal
Integration}\label{crosstalk-and-signal-integration}

Cells rarely respond to a single signal in isolation. Instead, multiple
pathways interact through \textbf{crosstalk} --- shared components,
convergent effectors, or mutual regulation. Signal integration is what
allows a cell to make complex decisions (proliferate, differentiate, or
die) based on the totality of its signalling environment.

\subsubsection{Convergence and
Divergence}\label{convergence-and-divergence}

\textbf{Convergence:} Multiple upstream signals activate the same
downstream effector. For example, both RTK/RAS/MAPK and Wnt/beta-catenin
pathways converge on cyclin D1 transcription, reinforcing the G1-to-S
transition. Both growth factor (PI3K/Akt) and integrin signalling
converge on mTORC1 to control cell growth.

\textbf{Divergence:} A single activated receptor triggers multiple
downstream cascades simultaneously. EGFR activation simultaneously
engages: - RAS-MAPK (proliferation) - PI3K-Akt (survival) -
PLC-gamma/IP\(_3\)/Ca\(^{2+}\) (immediate responses) - JAK-STAT (gene
expression) - c-Src (cytoskeletal remodelling)

\subsubsection{Scaffold Proteins and Signalling
Specificity}\label{scaffold-proteins-and-signalling-specificity}

Scaffold proteins physically assemble pathway components, ensuring
speed, specificity, and prevention of unwanted crosstalk:

{\def\LTcaptype{none} % do not increment counter
\begin{longtable}[]{@{}
  >{\raggedright\arraybackslash}p{(\linewidth - 4\tabcolsep) * \real{0.3478}}
  >{\raggedright\arraybackslash}p{(\linewidth - 4\tabcolsep) * \real{0.3043}}
  >{\raggedright\arraybackslash}p{(\linewidth - 4\tabcolsep) * \real{0.3478}}@{}}
\toprule\noalign{}
\begin{minipage}[b]{\linewidth}\raggedright
Scaffold
\end{minipage} & \begin{minipage}[b]{\linewidth}\raggedright
Pathway
\end{minipage} & \begin{minipage}[b]{\linewidth}\raggedright
Function
\end{minipage} \\
\midrule\noalign{}
\endhead
\bottomrule\noalign{}
\endlastfoot
KSR1 & RAS-MAPK & Brings RAF, MEK, and ERK together; promotes efficient
cascade activation \\
AKAP79/150 & cAMP-PKA & Anchors PKA, calcineurin (PP2B), and PKC near
substrates (e.g., ion channels) \\
Axin & Wnt & Core scaffold of the beta-catenin destruction complex \\
InaD & Drosophila phototransduction & PDZ-domain scaffold holding PLC,
TRP channel, and PKC \\
mTORC1/2 & PI3K-mTOR & Raptor (mTORC1) vs.~Rictor (mTORC2) determine
substrate access \\
\end{longtable}
}

Scaffold proteins explain how the same MAPK cascade can produce
different outcomes in different cell types: the scaffold determines
which substrates are phosphorylated.

\subsubsection{Positive and Negative Feedback
Loops}\label{positive-and-negative-feedback-loops}

\textbf{Positive feedback} creates switch-like (bistable
\citep{tyson2003}) responses: - ERK phosphorylates and activates SOS
(its own upstream GEF), creating an ultrasensitive on/off switch - p53
activates MDM2, which degrades p53 --- but also activates PUMA/NOXA,
creating a ``point of no return'' for apoptosis once the death threshold
is crossed - \hyperref[gl:caspase]{\textbf{Caspase}}-3 cleaves and
activates caspase-9 (amplification loop in apoptosis)

\textbf{Negative feedback} limits signal duration and creates
adaptation: - ERK phosphorylates SOS at inhibitory sites (slow negative
feedback counteracts fast positive feedback) - SOCS proteins induced by
JAK-STAT signalling bind JAKs and target them for degradation - Sprouty
(SPRY) proteins inhibit RAS-MAPK at the level of GRB2-SOS recruitment

\begin{quote}
\textbf{Concept Check 5:} The MAPK cascade exhibits both positive
feedback (ERK activates SOS) and negative feedback (ERK inhibits SOS at
different sites). How might the relative timing of these feedbacks
create a transient burst of ERK activity followed by adaptation? What
would happen if the negative feedback were eliminated by mutation?
\end{quote}

\begin{quote}
\textbf{Concept Check 6:} A patient with BRAF V600E melanoma is treated
with vemurafenib (BRAF inhibitor) alone. Initial response is dramatic,
but resistance develops in 6--12 months. Sequencing reveals upregulation
of CRAF, NRAS amplification, or MEK1 mutations in different patients.
Explain why each of these ``bypasses'' vemurafenib and why combination
therapy with trametinib (MEK inhibitor) achieves more durable responses.
\end{quote}

\begin{quote}
\textbf{Concept Check 7:} A neuron at rest contains
about 10\(^{\texorpdfstring{\ensuremath{-}}{-}7}\) M cytoplasmic {[}Ca\(^{2+}\){]}. Following an
action potential, {[}Ca\(^{2+}\){]} near a synaptic vesicle rises to
about 10\(^{\texorpdfstring{\ensuremath{-}}{-}5}\) M for about 1 ms before being
buffered and pumped back. Synaptotagmin has 2 C2 domains, each binding
2--3 Ca\(^{2+}\) ions cooperatively (Hill coefficient about 3).
Explain why this molecular architecture produces an essentially
``most-or-nothing'' exocytotic response to the brief calcium pulse, and
how this differs from the sustained calcium signals used in
CaMKII-mediated long-term potentiation.
\end{quote}

\begin{quote}
\textbf{Concept Check 8:} A scaffold protein binds both an upstream
kinase (RAF) and a downstream substrate (ERK) but does not bind the
intermediate kinase (MEK). Predict the qualitative effect on signal
transmission. What does this suggest about why most three components
must be present on the scaffold for productive signalling?
\end{quote}

\begin{figure}[htbp]
\centering
\includegraphics[width=0.82\linewidth,height=4.2in,keepaspectratio]{../figures/mermaid_inline/inline_mermaid_0045_a48c49019fcd.png}
\caption{Signal integration: multiple pathways converge on key decision nodes. Proliferative signals (RTK, Wnt, cytokines) converge on cyclin D1 transcription. Survival signals (PI3K/Akt) antagonise apoptosis by phosphorylating BAD. The cell's fate (proliferation, growth, or death) depends on the balance of active pathways.}
\end{figure}

\subsubsection{Computational Approaches to Signalling
Networks}\label{computational-approaches-to-signalling-networks}

Modern systems biology uses mathematical modelling to understand
signalling network behaviour:

\begin{itemize}
\tightlist
\item
  \textbf{Boolean networks:} model each node as ON/OFF; useful for large
  networks
\item
  \textbf{Ordinary differential equations (ODEs):} model concentration
  changes over time; quantitative predictions
\item
  \textbf{Stochastic modelling:} accounts for molecular noise in
  low-copy-number signalling (important in stem cell fate decisions)
\end{itemize}

The Huang-Ferrell model of the MAPK cascade demonstrated that the
cascade acts as an ultrasensitive switch: small changes in input (growth
factor concentration) produce most-or-none output (ERK activation). This
switch-like behaviour arises from the dual phosphorylation requirement
for ERK activation.

\begin{center}\rule{0.5\linewidth}{0.5pt}\end{center}

\subsection{The Cell Cycle --- Control Through CDK-Cyclin
Complexes}\label{the-cell-cycle-control-through-cdk-cyclin-complexes}

Cell division is controlled by cyclin-dependent kinases (CDKs) activated
by their partner cyclins (Nobel Prize 2001: Hartwell, Hunt, Nurse).

\subsubsection{Phases}\label{phases}

{\def\LTcaptype{none} % do not increment counter
\begin{longtable}[]{@{}llll@{}}
\toprule\noalign{}
Phase & Event & Duration (typical animal cell) & Key CDK-Cyclin \\
\midrule\noalign{}
\endhead
\bottomrule\noalign{}
\endlastfoot
G1 & Growth; protein synthesis & 6--12 h & CDK4/6-Cyclin D \\
S & DNA replication & 6--8 h & CDK2-Cyclin E/A \\
G2 & Growth; mitosis preparation & 3--5 h & CDK1-Cyclin A \\
M & Mitosis + cytokinesis & about 1 h & CDK1-Cyclin B (MPF) \\
\end{longtable}
}

\textbf{Checkpoints} ensure cell cycle quality: - \textbf{G1 checkpoint
(restriction point):} cell size adequate? DNA intact? Growth factors
present? CDK4/6-Cyclin D drives pRB phosphorylation \texorpdfstring{\ensuremath{\to}}{->} E2F release \texorpdfstring{\ensuremath{\to}}{->}
S-phase genes (cyclin E, DNA polymerase, thymidine kinase) -
\textbf{G2/M checkpoint:} DNA repair complete? CDK1-Cyclin B (MPF ---
maturation-promoting factor) triggers mitotic entry. Activated by CDC25
phosphatase; inhibited by Wee1 kinase. - \textbf{Spindle assembly
checkpoint (SAC):} most kinetochores attached to microtubules?
Mad2/BubR1 inhibit \hyperref[gl:anaphase]{\textbf{Anaphase}} Promoting
Complex (APC/C) until satisfied. APC/C then ubiquitinates securin \texorpdfstring{\ensuremath{\to}}{->}
separase released \texorpdfstring{\ensuremath{\to}}{->} cleaves cohesin \texorpdfstring{\ensuremath{\to}}{->} sister chromatid separation.

\subsubsection{Tumour Suppressors}\label{tumour-suppressors}

\begin{itemize}
\tightlist
\item
  \textbf{p53:} ``guardian of the
  \hyperref[gl:genome]{\textbf{genome}}''; activated by DNA damage (via
  ATM/ATR kinases, Chk1/Chk2) \texorpdfstring{\ensuremath{\to}}{->} arrests cell cycle (via p21/WAF1
  inhibiting CDK2) or triggers apoptosis (via PUMA, NOXA); mutated in
  about 50\% of human cancers. MDM2 is the E3 ubiquitin ligase
  that normally keeps p53 levels low; MDM2 inhibitors (nutlins) are in
  clinical trials.
\item
  \textbf{pRB:} binds and represses E2F in G1; CDK4/6 phosphorylation
  releases E2F \texorpdfstring{\ensuremath{\to}}{->} S phase entry; mutated/deleted in retinoblastoma and
  many cancers. CDK4/6 inhibitors (palbociclib, ribociclib) are used in
  HR\(^+\) breast cancer.
\end{itemize}

\begin{quote}
\textbf{Clinical Connection: CDK4/6 Inhibitors in Cancer} Palbociclib,
ribociclib, and abemaciclib specifically inhibit CDK4/6, preventing pRB
phosphorylation and maintaining E2F repression. This arrests cancer
cells in G1. These drugs have transformed treatment of hormone
receptor-positive, HER2-negative metastatic breast cancer, doubling
progression-free survival when combined with anti-estrogen therapy.
\end{quote}

\begin{center}\rule{0.5\linewidth}{0.5pt}\end{center}

\subsection{Apoptosis --- Programmed Cell
Death}\label{apoptosis-programmed-cell-death}

Apoptosis eliminates about 80 billion cells per day in the
adult human (equal to approximate cell proliferation). It differs
fundamentally from necrosis:

{\def\LTcaptype{none} % do not increment counter
\begin{longtable}[]{@{}
  >{\raggedright\arraybackslash}p{(\linewidth - 4\tabcolsep) * \real{0.2917}}
  >{\raggedright\arraybackslash}p{(\linewidth - 4\tabcolsep) * \real{0.3750}}
  >{\raggedright\arraybackslash}p{(\linewidth - 4\tabcolsep) * \real{0.3333}}@{}}
\toprule\noalign{}
\begin{minipage}[b]{\linewidth}\raggedright
Feature
\end{minipage} & \begin{minipage}[b]{\linewidth}\raggedright
Apoptosis
\end{minipage} & \begin{minipage}[b]{\linewidth}\raggedright
Necrosis
\end{minipage} \\
\midrule\noalign{}
\endhead
\bottomrule\noalign{}
\endlastfoot
Initiation & Genetically programmed & Accidental \\
Morphology & Cell shrinkage; DNA laddering; apoptotic bodies & Cell
rupture; inflammation \\
Inflammation & None (PS ``eat me'' signal) & Severe (DAMPs released) \\
ATP required & Yes & No \\
Caspases & Activated & Not involved \\
\end{longtable}
}

\subsubsection{Intrinsic (Mitochondrial)
Pathway}\label{intrinsic-mitochondrial-pathway}

Triggered by DNA damage, oxidative stress, growth factor withdrawal:

\begin{enumerate}
\def\labelenumi{\arabic{enumi}.}
\tightlist
\item
  \textbf{BH3-primarily proteins} (Bad, Bid, Puma, NOXA) activated by
  cellular stress
\item
  BH3-primarily proteins neutralise anti-apoptotic Bcl-2/Bcl-xL
\item
  \textbf{BAX/BAK} oligomerise in the OMM \texorpdfstring{\ensuremath{\to}}{->} form pores (MAC channel)
\item
  \textbf{Cytochrome c} released from intermembrane space into cytoplasm
\item
  Cytochrome c + Apaf-1 + ATP/dATP \texorpdfstring{\ensuremath{\to}}{->} \textbf{apoptosome} (7-subunit
  wheel, about 1 MDa)
\item
  Apoptosome activates \textbf{caspase-9} (initiator caspase)
\item
  Caspase-9 activates \textbf{caspase-3/-7} (executioner caspases)
\item
  Caspases cleave structural proteins (nuclear lamins, ICAD releasing
  CAD endonuclease \texorpdfstring{\ensuremath{\to}}{->} DNA fragmentation into about 180 bp
  nucleosomal ladders)
\end{enumerate}

\textbf{Bcl-2 family:} anti-apoptotic members (Bcl-2, Bcl-xL, Mcl-1)
bind and inhibit BAX/BAK. Bcl-2 overexpression in follicular lymphoma
(t(14;18) translocation) \texorpdfstring{\ensuremath{\to}}{->} cells cannot die. \textbf{Venetoclax} is a
BH3 mimetic drug that inhibits Bcl-2, restoring apoptosis in chronic
lymphocytic leukaemia (CLL).

\subsubsection{Extrinsic (Death Receptor)
Pathway}\label{extrinsic-death-receptor-pathway}

Triggered by extracellular death ligands (FasL, TNF-alpha, TRAIL):

\begin{enumerate}
\def\labelenumi{\arabic{enumi}.}
\tightlist
\item
  Ligand binds death receptor (FAS/CD95, TNFR1, DR4/DR5)
\item
  Death domain oligomers recruit FADD (Fas-associated death domain
  protein)
\item
  Procaspase-8 recruited \texorpdfstring{\ensuremath{\to}}{->} \textbf{DISC (Death-Inducing Signalling
  Complex)}
\item
  Caspase-8 auto-activated
\item
  In Type I cells: caspase-8 directly activates caspase-3 (sufficient
  DISC formation)
\item
  In Type II cells: caspase-8 cleaves Bid \texorpdfstring{\ensuremath{\to}}{->} tBid \texorpdfstring{\ensuremath{\to}}{->} activates
  mitochondrial pathway (amplification loop)
\end{enumerate}

\begin{center}\rule{0.5\linewidth}{0.5pt}\end{center}

\subsubsection{Other Forms of Regulated Cell
Death}\label{other-forms-of-regulated-cell-death}

Beyond classical apoptosis, several other regulated cell death pathways
have been identified:

\begin{itemize}
\tightlist
\item
  \textbf{Necroptosis:} Programmed necrosis triggered by TNF when
  caspase-8 is inhibited. RIPK1 and RIPK3 kinases phosphorylate MLKL,
  which oligomerises and permeabilises the plasma membrane. Important in
  host defence against viruses that express caspase inhibitors.
\item
  \textbf{Pyroptosis:} Inflammatory cell death mediated by gasdermin D
  (GSDMD). Inflammasome activation (NLRP3, NLRC4, AIM2) activates
  caspase-1, which cleaves GSDMD. The N-terminal fragment forms pores in
  the plasma membrane, releasing IL-1-beta and IL-18. Critical in
  \hyperref[gl:innate-immunity]{\textbf{innate immunity}} and sepsis.
\item
  \textbf{Ferroptosis:} Iron-dependent cell death characterised by lipid
  peroxidation. GPX4 (glutathione peroxidase 4) normally prevents lethal
  lipid ROS accumulation. GPX4 inhibition or glutathione depletion
  triggers ferroptosis. Relevant in neurodegeneration, kidney injury,
  and cancer therapy (some drug-resistant tumours are susceptible to
  ferroptosis inducers).
\end{itemize}

\begin{quote}
\textbf{Clinical Connection: Ferroptosis in Cancer Therapy}
Drug-resistant cancer cells that have undergone epithelial-mesenchymal
transition (EMT) become highly susceptible to ferroptosis. This
vulnerability arises because EMT upregulates polyunsaturated fatty acid
(PUFA) incorporation into membrane phospholipids, increasing
susceptibility to lipid peroxidation. Ferroptosis-inducing agents
(erastin, RSL3) are being investigated as therapies for
therapy-resistant cancers. Understanding cell death pathway diversity is
critical for developing next-generation anti-cancer strategies.
\end{quote}

\begin{center}\rule{0.5\linewidth}{0.5pt}\end{center}

\subsection{Cancer as Dysregulated
Signalling}\label{cancer-as-dysregulated-signalling}

Cancer results from accumulated mutations that activate oncogenes
(gain-of-function) and inactivate tumour suppressors (loss-of-function),
disrupting normal signalling control.

\subsubsection{Oncogenes and Their Normal
Counterparts}\label{oncogenes-and-their-normal-counterparts}

{\def\LTcaptype{none} % do not increment counter
\begin{longtable}[]{@{}
  >{\raggedright\arraybackslash}p{(\linewidth - 6\tabcolsep) * \real{0.2414}}
  >{\raggedright\arraybackslash}p{(\linewidth - 6\tabcolsep) * \real{0.2586}}
  >{\raggedright\arraybackslash}p{(\linewidth - 6\tabcolsep) * \real{0.3103}}
  >{\raggedright\arraybackslash}p{(\linewidth - 6\tabcolsep) * \real{0.1897}}@{}}
\toprule\noalign{}
\begin{minipage}[b]{\linewidth}\raggedright
Proto-oncogene
\end{minipage} & \begin{minipage}[b]{\linewidth}\raggedright
Normal function
\end{minipage} & \begin{minipage}[b]{\linewidth}\raggedright
Oncogenic mutation
\end{minipage} & \begin{minipage}[b]{\linewidth}\raggedright
Cancer type
\end{minipage} \\
\midrule\noalign{}
\endhead
\bottomrule\noalign{}
\endlastfoot
RAS (K-RAS, N-RAS) & Small GTPase (MAPK pathway) & Point mutations
(G12D, G12V, Q61L) lock in GTP state & Pancreatic, colorectal, lung \\
RAF (B-RAF) & Ser/Thr kinase (MAPK pathway) & V600E constitutive
activation & Melanoma, thyroid, colorectal \\
EGFR (ErbB1) & RTK (growth factor receptor) & Amplification; activating
mutations & Lung (NSCLC), glioblastoma \\
HER2 (ErbB2) & RTK (no ligand, heterodimerises) & Gene amplification
(about 25\% breast cancer) & Breast, gastric \\
MYC & Transcription factor (cell growth) & Amplification; translocation
(t(8;14)) & Burkitt lymphoma, many cancers \\
BCR-ABL & Constitutive tyrosine kinase & t(9;22) Philadelphia
\hyperref[gl:chromosome]{\textbf{chromosome}} & CML \\
PIK3CA & PI3K catalytic subunit & E545K, H1047R activating mutations &
Breast, endometrial, colorectal \\
\end{longtable}
}

\subsubsection{Targeted Cancer
Therapies}\label{targeted-cancer-therapies}

{\def\LTcaptype{none} % do not increment counter
\begin{longtable}[]{@{}
  >{\raggedright\arraybackslash}p{(\linewidth - 6\tabcolsep) * \real{0.1600}}
  >{\raggedright\arraybackslash}p{(\linewidth - 6\tabcolsep) * \real{0.2400}}
  >{\raggedright\arraybackslash}p{(\linewidth - 6\tabcolsep) * \real{0.3600}}
  >{\raggedright\arraybackslash}p{(\linewidth - 6\tabcolsep) * \real{0.2400}}@{}}
\toprule\noalign{}
\begin{minipage}[b]{\linewidth}\raggedright
Drug
\end{minipage} & \begin{minipage}[b]{\linewidth}\raggedright
Target
\end{minipage} & \begin{minipage}[b]{\linewidth}\raggedright
Mechanism
\end{minipage} & \begin{minipage}[b]{\linewidth}\raggedright
Cancer
\end{minipage} \\
\midrule\noalign{}
\endhead
\bottomrule\noalign{}
\endlastfoot
\textbf{Imatinib} (Gleevec) & BCR-ABL tyrosine kinase & ATP-competitive
inhibitor & CML \\
\textbf{Trastuzumab} (Herceptin) & HER2 extracellular domain &
Monoclonal antibody; ADCC + signalling block & HER2\(^+\) breast
cancer \\
\textbf{Vemurafenib} & B-RAF V600E & ATP-competitive inhibitor &
Melanoma \\
\textbf{Erlotinib/Gefitinib} & EGFR kinase domain & ATP-competitive
inhibitor & NSCLC \\
\textbf{Palbociclib} & CDK4/6 & Prevents pRB phosphorylation & HR\(^+\)
breast cancer \\
\textbf{Venetoclax} & Bcl-2 & BH3 mimetic; restores apoptosis & CLL \\
\textbf{Pembrolizumab} & PD-1 & Immune checkpoint inhibitor & Melanoma,
lung, many others \\
\textbf{Sotorasib} & K-RAS G12C & Covalent inhibitor (first RAS drug) &
NSCLC \\
\end{longtable}
}

\begin{quote}
\textbf{Clinical Connection: Imatinib --- The Paradigm of Targeted
Therapy} Chronic myeloid leukaemia (CML) is caused by the Philadelphia
chromosome --- a t(9;22) translocation creating the BCR-ABL fusion
protein (constitutively active tyrosine kinase). Before imatinib, 5-year
survival was about 30\%. Imatinib specifically inhibits BCR-ABL
by occupying the ATP-binding site, achieving complete cytogenetic
response in \textgreater80\% of patients and transforming CML into a
manageable chronic disease. This ``bench-to-bedside'' success story
(Brian Druker, Nicholas Lydon) demonstrated that understanding
signalling pathways at the molecular level can lead to revolutionary
therapies.
\end{quote}

\subsubsection{Clinical Pharmacology of Cell
Signalling}\label{clinical-pharmacology-of-cell-signalling}

Cell signalling is the most pharmacologically tractable layer of
cellular biology --- most major drug classes either activate, inhibit,
or compete with components of a signalling cascade. Understanding the
molecular logic of signalling therefore translates directly into
rational pharmacology. Three case studies illustrate the depth of this
connection.

\textbf{\texorpdfstring{\ensuremath{\beta}}{beta}-Blockers: from signal blockade to chronic-disease treatment.}
Propranolol (1964, James Black, Nobel Prize 1988) was the first
\texorpdfstring{\ensuremath{\beta}}{beta}-adrenergic receptor antagonist. It blocks \texorpdfstring{\ensuremath{\beta}}{beta}\(_1\) receptors on
cardiomyocytes (reducing heart rate and contractility, lowering oxygen
demand) and \texorpdfstring{\ensuremath{\beta}}{beta}\(_2\) receptors on bronchial smooth muscle (sometimes
causing wheeze in asthmatics). Modern \texorpdfstring{\ensuremath{\beta}}{beta}-blockers are graded by
selectivity: - \textbf{Non-selective:} propranolol, carvedilol --- block
\texorpdfstring{\ensuremath{\beta}}{beta}\(_1\), \texorpdfstring{\ensuremath{\beta}}{beta}\(_2\), sometimes \texorpdfstring{\ensuremath{\alpha}}{alpha}\(_1\). - \textbf{\texorpdfstring{\ensuremath{\beta}}{beta}\(_1\)-selective
(``cardioselective''):} metoprolol, atenolol, bisoprolol --- preferred
when bronchospasm is a concern. - \textbf{Inverse agonists:} carvedilol
--- actively reduces basal G-protein activity, not just blocks ligand
binding. - \textbf{Biased agonists/antagonists:} carvedilol, alprenolol
--- block G\(_s\) but activate \texorpdfstring{\ensuremath{\beta}}{beta}-arrestin signalling, possibly
explaining superior heart failure outcomes.

The mechanistic depth has grown over 60 years: from ``block adrenaline
binding'' \texorpdfstring{\ensuremath{\to}}{->} ``reduce cAMP'' \texorpdfstring{\ensuremath{\to}}{->} ``decouple from G-protein'' \texorpdfstring{\ensuremath{\to}}{->} ``biased
toward \texorpdfstring{\ensuremath{\beta}}{beta}-arrestin/ERK'' \texorpdfstring{\ensuremath{\to}}{->} ``remodel chronic gene expression patterns in
failing heart.'' A drug class that began as symptomatic therapy for
angina is now first-line in heart failure precisely \emph{because} it
modulates the chronic-stress signalling fingerprint.

\textbf{Statins: targeting the substrate of a regulatory cascade.}
HMG-CoA reductase catalyses the rate-limiting step of cholesterol
biosynthesis. Statins (lovastatin from \emph{Aspergillus} in 1976;
simvastatin, atorvastatin, rosuvastatin) are competitive inhibitors with
IC\(_{50}\) about 10 nM. Lower hepatic cholesterol synthesis \texorpdfstring{\ensuremath{\to}}{->}
upregulation of LDL receptor (via SREBP-2 transcription factor) \texorpdfstring{\ensuremath{\to}}{->}
increased LDL clearance from blood \texorpdfstring{\ensuremath{\to}}{->} about 30--50\% reduction
in LDL cholesterol. Beyond cholesterol, statins have \textbf{pleiotropic
effects} mediated by reduced isoprenoid (farnesyl, geranylgeranyl)
production: reduced membrane localisation of Rho/Rac GTPases
(anti-inflammatory), reduced platelet aggregation, improved endothelial
function. Many of these effects are arguably more important than
cholesterol lowering for cardiovascular outcomes --- supported by
clinical trials showing benefit even in patients with normal LDL.

\textbf{Imatinib in detail: lessons in resistance.} Imatinib (Gleevec,
2001) binds the inactive (``DFG-out'') conformation of the BCR-ABL
kinase domain at its ATP-binding site. The drug exploits a structural
quirk: the BCR-ABL kinase domain has a hydrophobic pocket that becomes
accessible primarily in the inactive conformation; this pocket is
\emph{not} present in most other kinases, providing remarkable
selectivity. Resistance develops through several mechanisms:

{\def\LTcaptype{none} % do not increment counter
\begin{longtable}[]{@{}
  >{\raggedright\arraybackslash}p{(\linewidth - 6\tabcolsep) * \real{0.2727}}
  >{\raggedright\arraybackslash}p{(\linewidth - 6\tabcolsep) * \real{0.2727}}
  >{\raggedright\arraybackslash}p{(\linewidth - 6\tabcolsep) * \real{0.2121}}
  >{\raggedright\arraybackslash}p{(\linewidth - 6\tabcolsep) * \real{0.2424}}@{}}
\toprule\noalign{}
\begin{minipage}[b]{\linewidth}\raggedright
Mechanism
\end{minipage} & \begin{minipage}[b]{\linewidth}\raggedright
Frequency
\end{minipage} & \begin{minipage}[b]{\linewidth}\raggedright
Example
\end{minipage} & \begin{minipage}[b]{\linewidth}\raggedright
Solution
\end{minipage} \\
\midrule\noalign{}
\endhead
\bottomrule\noalign{}
\endlastfoot
Point mutations in kinase domain & Most common & T315I (gatekeeper
residue) & Ponatinib (third-generation, designed against T315I) \\
BCR-ABL amplification & Less common & More fusion protein than drug can
inhibit & Higher dose or switch class \\
Activation of bypass pathways & Variable & SRC family kinases, PI3K &
Add SRC inhibitor (dasatinib) \\
Decreased intracellular drug & Variable & OCT1 transporter loss; MDR1
export & Drug-uptake-independent agents \\
\end{longtable}
}

This pattern --- initial dramatic response, eventual development of
resistance via mutation or bypass --- is now the norm in targeted
oncology. Modern strategies use \textbf{combination therapy} designed
against the resistance landscape: BRAF + MEK inhibitors in melanoma,
BCR-ABL + PI3K in CML resistance, EGFR + MEK in lung cancer.

\begin{quote}
\textbf{Concept Check 4b:} A patient with CML on imatinib for 3 years
develops disease progression. Sequencing reveals a T315I mutation in
BCR-ABL. Why does this single amino-acid change confer resistance, and
which next-generation drug would you choose? (Hint: T315 is the
``gatekeeper'' residue lining the imatinib binding pocket.)
\end{quote}

\subsubsection{Crosstalk and Signal Integration --- How Cells Process
Multiple
Signals}\label{crosstalk-and-signal-integration-how-cells-process-multiple-signals}

A real cell is not exposed to a single ligand at a time. A hepatocyte
receives glucagon, insulin, growth factors, fatty acids, cytokines, and
adhesion signals \emph{simultaneously}, and must integrate most of them
into a coherent metabolic and proliferative response. Three
architectural features make this integration possible:

\begin{enumerate}
\def\labelenumi{\arabic{enumi}.}
\tightlist
\item
  \textbf{Shared second messengers.} Many pathways feed into
  Ca\(^{2+}\), cAMP, or phosphoinositides. The cell's ``state'' is
  encoded in the time-integrated and spatially-resolved levels of these
  second messengers, not in any single signalling event.
\item
  \textbf{Coincidence detection.} Some effectors require \emph{two
  simultaneous} inputs: \textbf{PKC} is fully activated primarily when
  both DAG (membrane targeting) and Ca\(^{2+}\) (cytosolic increase) are
  present; \textbf{NMDA receptor} opens primarily when glutamate binds
  \emph{and} the membrane is sufficiently depolarised. Coincidence
  detection prevents false alarms.
\item
  \textbf{Conditional rules (logic gates).} Cells implement Boolean-like
  decision rules. Cell proliferation in normal epithelia requires growth
  factor (mitogen) AND adhesion (integrin) AND survival (PI3K/Akt). Loss
  of any one input causes anoikis (death by detachment) or quiescence.
  Cancer typically defeats these rules by activating each independently.
\end{enumerate}

\textbf{Concrete examples of crosstalk:}

\begin{itemize}
\tightlist
\item
  \textbf{Insulin--glucagon antagonism in liver.} Both signals reach the
  same hepatocyte. Insulin (via Akt) inhibits gluconeogenic gene
  expression; glucagon (via cAMP/PKA) activates it. The PEPCK and G6Pase
  promoters are bound by both insulin-responsive (FoxO1) and
  cAMP-responsive (CREB) transcription factors with opposing effects ---
  the integrated transcriptional output reflects the
  {[}insulin{]}/{[}glucagon{]} ratio.
\item
  \textbf{Growth factor + integrin convergence on mTORC1.} Growth factor
  \texorpdfstring{\ensuremath{\to}}{->} RTK \texorpdfstring{\ensuremath{\to}}{->} PI3K \texorpdfstring{\ensuremath{\to}}{->} Akt \texorpdfstring{\ensuremath{\to}}{->} TSC1/2 inhibition \texorpdfstring{\ensuremath{\to}}{->} Rheb-GTP \texorpdfstring{\ensuremath{\to}}{->} mTORC1. Integrin \texorpdfstring{\ensuremath{\to}}{->}
  FAK \texorpdfstring{\ensuremath{\to}}{->} PI3K \texorpdfstring{\ensuremath{\to}}{->} same Akt step. Both inputs are required for full mTORC1
  activation; loss of either prevents protein synthesis and cell growth,
  providing a checkpoint against unanchored proliferation.
\item
  \textbf{Stress \texorpdfstring{\ensuremath{\to}}{->} MAPK pathway crosstalk.} TNF-\texorpdfstring{\ensuremath{\alpha}}{alpha} activates IKK \texorpdfstring{\ensuremath{\to}}{->} NF-\texorpdfstring{\ensuremath{\kappa}}{kappa}B
  (survival) AND MEKK1 \texorpdfstring{\ensuremath{\to}}{->} JNK (apoptosis). The ratio determines outcome:
  brief TNF \texorpdfstring{\ensuremath{\to}}{->} NF-\texorpdfstring{\ensuremath{\kappa}}{kappa}B dominates \texorpdfstring{\ensuremath{\to}}{->} survival; prolonged TNF or with
  cyclohexamide (NF-\texorpdfstring{\ensuremath{\kappa}}{kappa}B-blocking) \texorpdfstring{\ensuremath{\to}}{->} JNK dominates \texorpdfstring{\ensuremath{\to}}{->} apoptosis.
\end{itemize}

\textbf{Computational modelling of integration.} Modern systems biology
builds these rules into ordinary differential equations (ODEs) or
Boolean networks. A 2009 paper by Janes \& Yaffe (\emph{Nature Reviews
MCB}) showed that 200+ measurements of signalling components in
TNF-treated cells could be reduced to a few \textbf{partial least
squares (PLS)} components that predicted apoptosis vs.~survival
outcomes. The lesson: cells appear to combine many inputs in a
low-dimensional decision space --- a finding that simplifies both basic
understanding and drug design (modulate the dominant principal
components, not 200 individual proteins).

\begin{quote}
\textbf{Concept Check 4:} Why would a combination of a BRAF inhibitor
(vemurafenib) and a MEK inhibitor (trametinib) be more effective than
either alone in BRAF V600E melanoma? Consider pathway reactivation
mechanisms.
\end{quote}

\begin{center}\rule{0.5\linewidth}{0.5pt}\end{center}

\subsection{Computational Bridge}\label{computational-bridge-7}

Cooperative ligand binding (Hill) captures RTK/RAS-MAPK dosage responses
in reduced form:

\begin{Shaded}
\begin{Highlighting}[]
\ImportTok{from}\NormalTok{ biology.cell }\ImportTok{import}\NormalTok{ hill\_equation}

\NormalTok{theta }\OperatorTok{=}\NormalTok{ hill\_equation(}\FloatTok{8.0}\NormalTok{, kd}\OperatorTok{=}\FloatTok{10.0}\NormalTok{, hill\_coefficient}\OperatorTok{=}\FloatTok{2.8}\NormalTok{)}
\BuiltInTok{print}\NormalTok{(}\BuiltInTok{round}\NormalTok{(theta, }\DecValTok{4}\NormalTok{))}
\end{Highlighting}
\end{Shaded}

\begin{quote}
\textbf{Clinical / systems note:} Dose-response steepness from
cooperativity is one reason combination targeted therapy can widen the
therapeutic window compared with single-agent saturation of one node.
\end{quote}

\begin{center}\rule{0.5\linewidth}{0.5pt}\end{center}

\subsubsection{Spatial Transcriptomics: Signalling in Tissue
Context}\label{spatial-transcriptomics-signalling-in-tissue-context}

Single-cell RNA-sequencing (scRNA-seq) dissociates tissues and so loses
the spatial information that is often the point --- you can measure what
each cell expresses but not \textbf{where it sits relative to its
neighbours}. \textbf{Spatial transcriptomics (ST)} recovers that
geometry by coupling \emph{in situ} RNA capture or detection to
high-throughput sequencing or imaging.

Two commercial platforms dominate. \textbf{10\texorpdfstring{\ensuremath{\times}}{x} Genomics Visium} prints a
glass slide with about 5000 barcoded spots (55 \texorpdfstring{\ensuremath{\mu}}{mu}m diameter,
about 100 \texorpdfstring{\ensuremath{\mu}}{mu}m spacing); a fresh-frozen tissue section is placed
on the slide, RNA diffuses into the spots, and each spot's transcriptome
is sequenced with its spot barcode. Resolution is sub-cellular to
small-cluster (1--10 cells per spot). \textbf{MERFISH (Multiplexed
Error-Robust FISH)} and \textbf{seqFISH+} instead use combinatorial
error-correcting fluorescent in situ hybridisation to detect thousands
of transcripts at \textbf{diffraction-limited (about 100 nm)
resolution} --- true single-molecule, single-cell spatial mapping.
\textbf{Slide-seq} and \textbf{Stereo-seq} push sequencing-based
approaches toward 10 \texorpdfstring{\ensuremath{\mu}}{mu}m and 500 nm spot sizes respectively.

Why this matters for signalling: paracrine and juxtacrine signalling ---
Notch--Delta, ephrin--Eph, WNT morphogen gradients --- depend on cells
\emph{being close to each other}. Spatial transcriptomics lets
investigators map ligand-expressing cells next to receptor-expressing
cells across a whole tissue section and infer \textbf{ligand-receptor
interaction networks in situ} (tools: CellChat, NicheNet,
SpatialCellChat, 2022-2024). In tumour biology, ST reveals
\hyperref[gl:niche]{\textbf{niches}} of exhausted T cells in
immunologically ``cold'' regions of the same tumour, motivating
spatially-guided immunotherapy combinations. In developmental biology,
ST in mouse embryos (Stereo-seq, \emph{Cell} 2022) produced the first
cell-resolved, tissue-resolved organogenesis atlas. The technology cost
has dropped about 10x every two years; a whole-slide Visium
experiment costs about \$2000 in 2024. For the student, ST is
the bridge between molecular cell biology (``what genes are on?'') and
anatomy (``where?'') --- and it is the platform on which the next decade
of in situ discovery will be built.

\begin{center}\rule{0.5\linewidth}{0.5pt}\end{center}

\subsection{Current Evidence and Frontier
Biology}\label{current-evidence-and-frontier-biology-10}

For \textbf{Cell Signalling and Communication}, frontier biology belongs
inside the evidence logic of the chapter. Cell biology is increasingly
measured as live, spatial, single-cell, and perturbational data rather
than static diagrams alone. The core reading question is this:
signalling explanations should include receptor context, dose, timing,
feedback, crosstalk, and cellular state.

\begin{itemize}
\tightlist
\item
  \textbf{What to verify:} identify the observation, model, assay, or
  dataset that would make the claim stronger or weaker.
\item
  \textbf{What to qualify:} state the scale, organism, cell type,
  environmental condition, or population where the claim is expected to
  hold.
\item
  \textbf{What to compare:} test at least one alternative explanation,
  baseline, or null model before treating the pattern as causal.
\item
  \textbf{What to cite:} distinguish primary evidence, review synthesis,
  public dataset, and institutional guidance; for recent or numeric
  claims, prefer the source closest to the measurement and state what
  has changed since it was published.
\end{itemize}

Ask what measurement scale is being claimed: nanometre structure,
single-cell transcript abundance, organelle dynamics, tissue context, or
organismal phenotype.

\textbf{Source practice:} For cell-state claims, distinguish microscopy,
live-cell perturbation, single-cell sequencing, spatial transcriptomics,
and biochemical assay evidence before making a causal statement.

\subsubsection{Current Evidence Map: Spatial Single-Cell Signalling
Evidence}\label{current-evidence-map-spatial-single-cell-signalling-evidence}

\begin{figure}[htbp]
\centering
\includegraphics[width=0.82\linewidth,height=4.2in,keepaspectratio]{../figures/mermaid_inline/inline_mermaid_0046_5910d9574b69.png}
\caption{A signalling claim is stronger when receptor state, ligand source, cell neighbourhood, and perturbation response point to the same mechanism.}
\end{figure}

\subsection{Summary}\label{summary-10}

\begin{itemize}
\tightlist
\item
  Cell signalling: ligand \texorpdfstring{\ensuremath{\to}}{->} receptor \texorpdfstring{\ensuremath{\to}}{->} cascade \texorpdfstring{\ensuremath{\to}}{->} response. GPCRs use
  cAMP or IP\(_3\)/DAG via heterotrimeric G proteins; RTKs use tyrosine
  phosphorylation and the RAS-MAPK cascade; ion channels open
  immediately upon ligand binding.
\item
  Second messengers (cAMP, cGMP, IP\(_3\), DAG, Ca\(^{2+}\), PIP\(_3\))
  amplify and diversify signals within the cell.
\item
  Signal termination involves GTPase activity, phosphodiesterases,
  protein phosphatases, receptor internalisation, and ubiquitin-mediated
  degradation.
\item
  The Wnt/beta-catenin, JAK-STAT, and Notch pathways control
  development, immune function, and stem cell maintenance.
\item
  The cell cycle is governed by CDK-cyclin complexes; checkpoints (G1,
  G2/M, SAC) ensure fidelity. p53 and pRB are the primary tumour
  suppressors.
\item
  Apoptosis (programmed cell death) proceeds via intrinsic
  (mitochondrial) or extrinsic (death receptor) pathways, both
  converging on executioner caspases.
\item
  Cancer results from dysregulated signalling (oncogenes + tumour
  suppressor loss); targeted therapies exploit specific pathway
  dependencies.
\item
  \textbf{Connections:} See Unit IV (growth-factor-driven
  transcription), Unit V (tumour suppressor
  \hyperref[gl:allele]{\textbf{allele}}s), Unit IX (endocrine first
  messengers).
\end{itemize}

\begin{center}\rule{0.5\linewidth}{0.5pt}\end{center}

\subsection{Review Questions}\label{review-questions-10}

\begin{enumerate}
\def\labelenumi{\arabic{enumi}.}
\item
  Draw the complete G\(_s\)-cAMP-PKA signalling cascade from epinephrine
  binding to glycogen breakdown. Identify most amplification steps and
  termination mechanisms.
\item
  Compare and contrast GPCRs and RTKs in terms of structure, mechanism
  of activation, downstream signalling, and examples of each.
\item
  Explain the RAS-MAPK pathway. Why is RAS one of the most commonly
  mutated oncogenes in human cancer?
\item
  Describe the Wnt/beta-catenin pathway. Why do APC mutations lead to
  colorectal cancer?
\item
  A patient has a gain-of-function mutation in JAK2 (V617F). Explain the
  molecular mechanism and predict the clinical
  \hyperref[gl:phenotype]{\textbf{phenotype}}.
\item
  Compare the intrinsic and extrinsic apoptosis pathways. Where do they
  converge, and what is the significance of BH3-primarily proteins?
\item
  Explain the mechanism of action of imatinib in CML. Why does
  resistance eventually develop, and how might it be overcome?
\item
  Describe four mechanisms by which cells terminate signalling
  responses. For each, give a specific molecular example.
\item
  Explain how CDK4/6 inhibitors (palbociclib) arrest the cell cycle in
  G1. Why are they used specifically in HR\(^+\) breast cancer?
\item
  Venetoclax is a BH3 mimetic that inhibits Bcl-2. Explain its mechanism
  of action and why it is effective in cancers that overexpress Bcl-2.
\item
  Using \texttt{hill\_equation}, compare fractional occupancy at
  \([L] = K_d\) for \(n = 1\) vs.~\(n = 4\). Which curve better matches
  oxygen binding to haemoglobin, and why?
\item
  Sketch how receptor internalisation provides \textbf{signal
  termination} for RTKs; link to one resistance mechanism in targeted
  cancer therapy.
\end{enumerate}

\begin{center}\rule{0.5\linewidth}{0.5pt}\end{center}

\subsection{Further Reading and Source
Notes}\label{further-reading-and-source-notes-10}

\begin{itemize}
\tightlist
\item
  Alon (2019). \emph{An Introduction to Systems Biology: Design
  Principles of Biological Circuits}. Chapman and Hall/CRC.
\item
  Tyson et al.~(2003). Sniffers, buzzers, toggles and blinkers: dynamics
  of regulatory and signaling pathways in the cell. \emph{Current
  Opinion in Cell Biology}, 15.
\end{itemize}

\begin{center}\rule{0.5\linewidth}{0.5pt}\end{center}

\subsection{Key Terms}\label{key-terms-10}

{\def\LTcaptype{none} % do not increment counter
\begin{longtable}[]{@{}
  >{\raggedright\arraybackslash}p{(\linewidth - 2\tabcolsep) * \real{0.2857}}
  >{\raggedright\arraybackslash}p{(\linewidth - 2\tabcolsep) * \real{0.7143}}@{}}
\toprule\noalign{}
\begin{minipage}[b]{\linewidth}\raggedright
Term
\end{minipage} & \begin{minipage}[b]{\linewidth}\raggedright
Definition
\end{minipage} \\
\midrule\noalign{}
\endhead
\bottomrule\noalign{}
\endlastfoot
\textbf{GPCR} & G protein-coupled receptor; 7-TM superfamily; largest
receptor family (about 800 in humans) \\
\textbf{Adenylyl cyclase} & \hyperref[gl:enzyme]{\textbf{Enzyme}}
converting ATP to cAMP; activated by G\(_s\), inhibited by G\(_i\) \\
\textbf{cAMP} & Cyclic AMP; second messenger activating PKA; degraded by
phosphodiesterases \\
\textbf{PKA} & Protein kinase A; cAMP-activated Ser/Thr kinase;
phosphorylates diverse substrates \\
\textbf{PIP\(_2\)} & Phosphatidylinositol 4,5-bisphosphate; PLC
substrate yielding IP\(_3\) + DAG \\
\textbf{IP\(_3\)} & Inositol 1,4,5-trisphosphate; releases Ca\(^{2+}\)
from ER via IP\(_3\) receptor channels \\
\textbf{RTK} & Receptor tyrosine kinase; dimerisation-activated;
autophosphorylation recruits SH2 proteins \\
\textbf{MAP kinase} & Mitogen-activated protein kinase (ERK); activated
by dual phosphorylation (Thr/Tyr) \\
\textbf{RAS} & Small GTPase; molecular switch; most commonly mutated
oncogene family \\
\textbf{CDK-cyclin} & Cyclin-dependent kinase paired with regulatory
cyclin; drives cell cycle progression \\
\textbf{Apoptosome} & Apaf-1/cytochrome c/caspase-9 complex; activates
executioner caspases \\
\textbf{Caspase} & Cysteine-aspartate protease; initiator (8, 9, 10) and
executioner (3, 6, 7) types \\
\textbf{Bcl-2} & Anti-apoptotic protein; inhibits BAX/BAK pore
formation; overexpressed in follicular lymphoma \\
\textbf{Wnt} & Secreted glycoprotein; activates beta-catenin signalling;
critical for development and stem cells \\
\textbf{JAK-STAT} & Janus kinase-Signal Transducer and Activator of
Transcription; cytokine signalling pathway \\
\textbf{Oncogene} & Mutated proto-oncogene with gain-of-function; drives
uncontrolled cell growth \\
\textbf{p53} & Tumour suppressor transcription factor; ``guardian of the
genome''; mutated in about 50\% of cancers \\
\end{longtable}
}

\begin{center}\rule{0.5\linewidth}{0.5pt}\end{center}

\subsubsection{Companion Source
Module}\label{companion-source-module-10}

\textbf{Cell Signalling and Communication} should leave a reproducible
trail from a biological claim to the code, figure, diagram, or
paper-based activity that can test it. Use the surfaces below to inspect
the chapter's assumptions, rerun the relevant model, or compare the
manuscript explanation with companion labs and figures.

{\def\LTcaptype{none} % do not increment counter
\begin{longtable}[]{@{}
  >{\raggedright\arraybackslash}p{(\linewidth - 2\tabcolsep) * \real{0.5000}}
  >{\raggedright\arraybackslash}p{(\linewidth - 2\tabcolsep) * \real{0.5000}}@{}}
\toprule\noalign{}
\begin{minipage}[b]{\linewidth}\raggedright
Surface
\end{minipage} & \begin{minipage}[b]{\linewidth}\raggedright
Use it for
\end{minipage} \\
\midrule\noalign{}
\endhead
\bottomrule\noalign{}
\endlastfoot
\breaktt{src/biology/cell/cell\_biology.py}
(\breaktt{receptor\_occupancy}, \texttt{hill\_equation},
\breaktt{signal\_amplification}) & Quantify ligand binding,
cooperativity, and cascade gain. \\
\breaktt{src/mermaid/biology\_diagrams.py}
(\breaktt{hormone\_signaling\_diagram},
\breaktt{immune\_response\_diagram}) & Compare receptor logic across
endocrine and immune examples. \\
\end{longtable}
}

\textbf{Reproducibility check:} specify dose, timing, receptor context,
feedback, and readout before inferring pathway causality.
\textbf{Cross-reference:} extend the same logic in
\cref{sec:unit_IX_endocrine_and_immune} and
\cref{sec:unit_IV_epigenetics_and_gene_regulation}.

\newpage

\section*{Unit III --- Energy and Metabolism:
Introduction}\label{unit-iii-energy-and-metabolism-introduction}
\addcontentsline{toc}{section}{Unit III --- Energy and Metabolism:
Introduction}

\subsection*{Why This Unit Matters}\label{why-this-unit-matters-3}
\addcontentsline{toc}{subsection}{Why This Unit Matters}

Life is, at its core, an ongoing argument with the second law of
thermodynamics. The second law demands that disorder increases in any
closed system; living organisms are ordered, and they maintain that
order at the price of consuming energy and exporting entropy to their
surroundings. A resting human adult dissipates approximately 80 watts
--- the same as an incandescent light bulb --- in the form of body heat,
CO\texorpdfstring{\textsubscript{2}}{2}, and metabolic waste. Every breath you take, every protein you
synthesise, every ion you pump against its gradient, uses energy
extracted from the oxidation of food.

The molecular currency of this energy economy is ATP (adenosine
triphosphate). At any moment, your cells contain about 50 g of
ATP, yet you consume roughly 40 kg of ATP every day at rest. This means
each ATP molecule is recycled approximately 800 times --- phosphorylated
from ADP in mitochondria, hydrolysed to ADP in the cytoplasm, and
regenerated in an endless chemical cycle. During vigorous exercise, ATP
turnover in muscle can exceed 0.5 kg per minute.

This unit traces energy from its entry point --- sunlight captured by
chlorophyll, or chemical bonds in food --- through the interconnected
pathways of glycolysis, the TCA cycle, oxidative phosphorylation, and
photosynthesis, to its exit as heat and CO\texorpdfstring{\textsubscript{2}}{2}. Each pathway is treated
quantitatively: you will derive ATP yields, calculate free energy
changes, and model the regulatory logic that switches metabolism between
fed and fasted states. The central thread is Peter Mitchell's
chemiosmosis (Nobel Prize 1978): ion gradients across membranes couple
electron flow to ATP synthesis, and the same principle operates in both
mitochondria and chloroplasts.

\begin{center}\rule{0.5\linewidth}{0.5pt}\end{center}

\subsection*{Landmark Discoveries}\label{landmark-discoveries-3}
\addcontentsline{toc}{subsection}{Landmark Discoveries}

{\def\LTcaptype{none} % do not increment counter
\begin{longtable}[]{@{}
  >{\raggedright\arraybackslash}p{(\linewidth - 8\tabcolsep) * \real{0.2407}}
  >{\raggedright\arraybackslash}p{(\linewidth - 8\tabcolsep) * \real{0.0741}}
  >{\raggedright\arraybackslash}p{(\linewidth - 8\tabcolsep) * \real{0.2963}}
  >{\raggedright\arraybackslash}p{(\linewidth - 8\tabcolsep) * \real{0.1667}}
  >{\raggedright\arraybackslash}p{(\linewidth - 8\tabcolsep) * \real{0.2222}}@{}}
\toprule\noalign{}
\begin{minipage}[b]{\linewidth}\raggedright
Discoverer(s)
\end{minipage} & \begin{minipage}[b]{\linewidth}\raggedright
Year
\end{minipage} & \begin{minipage}[b]{\linewidth}\raggedright
Journal / Source
\end{minipage} & \begin{minipage}[b]{\linewidth}\raggedright
Discovery
\end{minipage} & \begin{minipage}[b]{\linewidth}\raggedright
Significance
\end{minipage} \\
\midrule\noalign{}
\endhead
\bottomrule\noalign{}
\endlastfoot
Hans Krebs & 1937 & \emph{Enzymologia} & Citric acid (Krebs/TCA) cycle &
Identified the central hub of aerobic catabolism; Nobel Prize 1953 \\
Otto Warburg & 1923--1930 & \emph{Biochemische Zeitschrift} &
Respiratory enzyme and Warburg effect & First measured oxygen
consumption in cancer; discovered aerobic glycolysis \\
Peter Mitchell & 1961 & \emph{Nature} & Chemiosmotic hypothesis &
Proposed proton gradients drive ATP synthesis; revolutionary paradigm
shift; Nobel Prize 1978 \\
Paul Boyer \& John Walker & 1994--1997 & \emph{Biochemistry; J. Mol.
Biol.} & Binding-change mechanism of ATP synthase & Explained rotary
catalysis of F\texorpdfstring{\textsubscript{1}}{1}F\texorpdfstring{\textsubscript{0}}{0}-ATP synthase; Nobel Prize 1997 \\
Melvin Calvin, Andrew Benson \& James Bassham & 1950 & \emph{J. Am.
Chem. Soc.} & Calvin cycle (carbon fixation) & Traced \texorpdfstring{\textsuperscript{1}}{1}\texorpdfstring{\textsuperscript{4}}{4}C through
photosynthesis; established the biochemistry of C\texorpdfstring{\textsubscript{3}}{3} fixation; Nobel Prize
1961 \\
Racker \& Stoeckenius & 1974 & \emph{J. Biol. Chem.} & Bacteriorhodopsin
+ ATP synthase reconstitution & Direct experimental proof of Mitchell's
chemiosmosis using purified components \\
\end{longtable}
}

\begin{center}\rule{0.5\linewidth}{0.5pt}\end{center}

\subsection*{Key Concepts and
Connections}\label{key-concepts-and-connections-3}
\addcontentsline{toc}{subsection}{Key Concepts and Connections}

\begin{figure}[htbp]
\centering
\includegraphics[width=0.82\linewidth,height=4.2in,keepaspectratio]{../figures/mermaid_inline/inline_mermaid_0047_5f7536be6079.png}
\caption{Unit III concept map - Energy and Metabolism. Orange = ATP currency; red = oxidative phosphorylation; green = photosynthesis; purple = ATP synthase.}
\end{figure}

\begin{center}\rule{0.5\linewidth}{0.5pt}\end{center}

\subsection*{Current Evidence Thread}\label{current-evidence-thread-3}
\addcontentsline{toc}{subsection}{Current Evidence Thread}

Read this unit as one regulated network rather than three separate
pathways: respiration, the integrated fasting/fed economy, and
photosynthesis are most constrained by the same currencies of energy
charge, redox poise, and compartmentation, and a claim in one chapter is
primarily as strong as the flux, sensor, and condition it specifies.
Metabolism is now studied as a regulated network constrained by energy,
redox balance, compartmentation, and environment. As you move through
the chapters, keep a two-column note: \textbf{claim} on the left,
\textbf{evidence that would change my confidence} on the right. By the
end of the unit, each major idea should be tied to a measurement, model,
citation, or paper-based lab decision.

\subsection*{Chapter Roadmap}\label{chapter-roadmap-3}
\addcontentsline{toc}{subsection}{Chapter Roadmap}

{\def\LTcaptype{none} % do not increment counter
\begin{longtable}[]{@{}
  >{\raggedright\arraybackslash}p{(\linewidth - 6\tabcolsep) * \real{0.1556}}
  >{\raggedright\arraybackslash}p{(\linewidth - 6\tabcolsep) * \real{0.1111}}
  >{\raggedright\arraybackslash}p{(\linewidth - 6\tabcolsep) * \real{0.2889}}
  >{\raggedright\arraybackslash}p{(\linewidth - 6\tabcolsep) * \real{0.4444}}@{}}
\toprule\noalign{}
\begin{minipage}[b]{\linewidth}\raggedright
Chapter
\end{minipage} & \begin{minipage}[b]{\linewidth}\raggedright
Title
\end{minipage} & \begin{minipage}[b]{\linewidth}\raggedright
Core Question
\end{minipage} & \begin{minipage}[b]{\linewidth}\raggedright
Key Equation / Model
\end{minipage} \\
\midrule\noalign{}
\endhead
\bottomrule\noalign{}
\endlastfoot
\textbf{9} & Bioenergetics and Cellular Respiration & How do cells
extract and convert chemical energy from glucose? &
\(\Delta G = \Delta G^{\circ\prime} + RT\ln Q\); P/O ratios \\
\textbf{10} & Photosynthesis & How do plants capture solar energy and
fix CO\texorpdfstring{\textsubscript{2}}{2}? & Light-harvesting cross-section; Z-scheme; Calvin cycle
stoichiometry \\
\textbf{11} & Metabolic Integration & How do hormones and allosteric
signals coordinate anabolic and catabolic pathways? & AMPK activation;
insulin/glucagon signalling logic \\
\end{longtable}
}

\begin{center}\rule{0.5\linewidth}{0.5pt}\end{center}

\subsection*{Connections Across the
Textbook}\label{connections-across-the-textbook-3}
\addcontentsline{toc}{subsection}{Connections Across the Textbook}

\begin{itemize}
\tightlist
\item
  \textbf{ATP and free energy} introduced here underpin most driven
  processes in Unit IV (DNA replication, transcription, translation),
  Unit VIII (active transport in phloem), and Unit IX (Na\texorpdfstring{\textsuperscript{+}}{+}/K\texorpdfstring{\textsuperscript{+}}{+}-ATPase in
  neurons).
\item
  \textbf{The TCA cycle as metabolic hub} is integrated in
  \cref{sec:unit_III_metabolic_integration} and connects to amino acid
  catabolism (Unit IV), fatty acid oxidation, and nucleotide synthesis.
\item
  \textbf{Proton-motive force} is the conceptual bridge between
  mitochondrial ATP synthesis (this unit) and chloroplast ATP synthesis
  (\cref{sec:unit_III_photosynthesis}), and between bacterial physiology
  (Unit VII) and the evolution of eukaryotes (Unit II, endosymbiosis).
\item
  \textbf{Warburg effect} (cancer cells performing aerobic glycolysis)
  appears in Unit II (cell signalling oncogenes), Unit IV (genomic
  instability), and the clinical connection sections throughout.
\end{itemize}

\begin{quote}
\textbf{Key vocabulary introduced here:} free energy, enthalpy, entropy,
ATP hydrolysis, glycolysis, oxidative phosphorylation, proton-motive
force, chemiosmosis, electron transport chain, substrate-level
phosphorylation, fermentation, Calvin cycle, Warburg effect.
\end{quote}

\subsection*{Computational Toolbox --- Unit
III}\label{computational-toolbox-unit-iii}
\addcontentsline{toc}{subsection}{Computational Toolbox --- Unit III}

\begin{Shaded}
\begin{Highlighting}[]
\ImportTok{from}\NormalTok{ biology.biochemistry }\ImportTok{import}\NormalTok{ atp\_free\_energy, glycolysis\_summary}

\CommentTok{\# ATP hydrolysis free energy under physiological conditions}
\NormalTok{delta\_g }\OperatorTok{=}\NormalTok{ atp\_free\_energy(atp\_conc\_mM}\OperatorTok{=}\FloatTok{3.0}\NormalTok{, adp\_conc\_mM}\OperatorTok{=}\FloatTok{1.0}\NormalTok{, pi\_conc\_mM}\OperatorTok{=}\FloatTok{10.0}\NormalTok{)}
\BuiltInTok{print}\NormalTok{(}\SpecialStringTok{f"ΔG (ATP hydrolysis) = }\SpecialCharTok{\{}\NormalTok{delta\_g}\SpecialCharTok{:.1f\}}\SpecialStringTok{ kJ/mol"}\NormalTok{)}
\CommentTok{\# Expected: ΔG (ATP hydrolysis) ≈ {-}45.2 kJ/mol}
\CommentTok{\# (Standard ΔG°\textquotesingle{} = {-}30.5 kJ/mol; physiological conditions make it more negative}
\CommentTok{\#  due to low [ADP][Pi] and high [ATP] in active cells)}

\CommentTok{\# Comparing aerobic vs anaerobic ATP yield}
\NormalTok{aerobic }\OperatorTok{=}\NormalTok{ glycolysis\_summary()}
\BuiltInTok{print}\NormalTok{(}\SpecialStringTok{f"Glycolysis only: }\SpecialCharTok{\{}\NormalTok{aerobic}\SpecialCharTok{.}\NormalTok{net\_atp}\SpecialCharTok{\}}\SpecialStringTok{ ATP per glucose"}\NormalTok{)}
\CommentTok{\# 2 ATP (substrate{-}level only); aerobic adds about 32 more via oxidative phosphorylation}

\CommentTok{\# The complete aerobic yield}
\NormalTok{total\_aerobic\_atp }\OperatorTok{=}\NormalTok{ aerobic.net\_atp }\OperatorTok{+} \DecValTok{2} \OperatorTok{+} \DecValTok{10}\OperatorTok{*}\FloatTok{2.5} \OperatorTok{+} \DecValTok{2}\OperatorTok{*}\FloatTok{1.5}  \CommentTok{\# GTP + NADH + FADH2}
\BuiltInTok{print}\NormalTok{(}\SpecialStringTok{f"Aerobic total (estimate): about }\SpecialCharTok{\{}\NormalTok{total\_aerobic\_atp}\SpecialCharTok{:.0f\}}\SpecialStringTok{ ATP per glucose"}\NormalTok{)}
\CommentTok{\# Expected: Aerobic total (estimate): about 32 ATP per glucose}
\end{Highlighting}
\end{Shaded}

\begin{quote}
\textbf{Try it yourself:} The P/O ratio (ATP per oxygen atom) is 2.5 for
NADH and 1.5 for FADH\texorpdfstring{\textsubscript{2}}{2}. Calculate the theoretical maximum yield from one
glucose: 10 NADH \texorpdfstring{\ensuremath{\times}}{x} 2.5 + 2 FADH\texorpdfstring{\textsubscript{2}}{2} \texorpdfstring{\ensuremath{\times}}{x} 1.5 + 4 substrate-level = \_\_?
\end{quote}

\begin{center}\rule{0.5\linewidth}{0.5pt}\end{center}

\emph{Source modules: \breaktt{src/biology/biochemistry/biochemistry.py}
--- \breaktt{glycolysis\_summary()}, \texttt{atp\_free\_energy()},
\breaktt{michaelis\_menten()}.} \emph{Figures:
\breaktt{src/visualization/} (ETC diagrams, ATP yield tables);
\breaktt{src/mermaid/biology\_diagrams.py} (metabolic pathway diagrams).}

\newpage

\section{Bioenergetics and Cellular
Respiration}\label{bioenergetics-and-cellular-respiration}

\label{sec:unit_III_bioenergetics_and_respiration}

\begin{quote}
\textbf{Ch 9} \texorpdfstring{\ensuremath{\cdot}}{.} Level 3/3 \texorpdfstring{\ensuremath{\cdot}}{.} 60 min read \texorpdfstring{\ensuremath{\cdot}}{.} 100 min lecture \texorpdfstring{\ensuremath{\cdot}}{.}
Prerequisites: \cref{sec:unit_II_cell_structure},
\cref{sec:unit_I_enzymes_and_kinetics}
\end{quote}

\subsection{Learning Objectives}\label{learning-objectives-11}

\begin{enumerate}
\def\labelenumi{\arabic{enumi}.}
\tightlist
\item
  Define the laws of
  \hyperref[gl:thermodynamics]{\textbf{thermodynamics}} and apply them
  to biological reactions, distinguishing standard
  (\(\Delta G^{\circ\prime}\)) from physiological (\(\Delta G\)) free
  energies.
\item
  Explain the role of ATP as a central phosphoryl-group donor and
  short-term energy carrier, including the molecular basis of its high
  phosphoryl group transfer potential.
\item
  Describe most 10 steps of
  \hyperref[gl:glycolysis]{\textbf{glycolysis}} with
  \hyperref[gl:enzyme]{\textbf{enzyme}}s, intermediates, and regulatory
  mechanisms.
\item
  Explain the pyruvate dehydrogenase complex and its regulation.
\item
  Describe most 8 steps of the TCA cycle with carbon tracking, cofactor
  yields, and regulation.
\item
  Describe the electron transport chain complexes I--IV, including
  proton stoichiometry and the Q cycle.
\item
  Explain \hyperref[gl:atp-synthase]{\textbf{ATP synthase}} structure
  and the rotary catalysis mechanism, including the proton-motive force
  equation.
\item
  Calculate net ATP yield from glucose oxidation using modern P/O ratios
  (about 2.5 for NADH, about 1.5 for FADH\(_2\)) and
  explain how the ``30--32 ATP'' number is derived.
\item
  Compare substrate-level and
  \hyperref[gl:oxidative-phosphorylation]{\textbf{oxidative
  phosphorylation}} quantitatively across yield, speed, sustainability,
  and physiological niche.
\item
  Describe \hyperref[gl:fermentation]{\textbf{fermentation}} pathways,
  the Cori cycle, and the glucose-alanine cycle as inter-organ shuttles.
\item
  Explain gluconeogenesis (four bypass enzymes) and its reciprocal
  regulation with glycolysis.
\item
  Quantify reactive oxygen species (ROS) production at the ETC and the
  antioxidant defences that constrain it.
\item
  Distinguish anaplerotic and cataplerotic flux through the TCA cycle in
  the context of inter-organ metabolism.
\item
  Explain physiological (UCP1, brown adipose) and pharmacological (DNP)
  uncoupling at a molecular level.
\item
  Compute the thermodynamic efficiency of glucose oxidation from
  standard free energies of formation.
\end{enumerate}

\subsubsection{Study Blueprint}\label{study-blueprint-11}

\begin{itemize}
\tightlist
\item
  \textbf{Big idea:} Cells harvest free energy by coupling redox
  chemistry to phosphoryl transfer and ion gradients.
\item
  \textbf{Core concepts:} free energy, redox, glycolysis, oxidative
  phosphorylation.
\item
  \textbf{Framework alignment:} Vision \& Change: Pathways and
  transformations of energy and matter, Systems; AP Biology: Energetics,
  Systems Interactions; NGSS-style topics: Matter and Energy in
  Organisms and Ecosystems.
\item
  \textbf{Model or quantitative lens:} Delta G, ATP yield, and
  electron-carrier accounting.
\item
  \textbf{Data skill:} Track carbon, electrons, and ATP across a
  pathway.
\item
  \textbf{Practice cadence:} Representing and Describing Data,
  Statistical Tests and Data Analysis.
\item
  \textbf{Common misconception to repair:} ATP is not stored energy in a
  vague sense; it is a coupling currency with defined reaction
  chemistry.
\item
  \textbf{Primary lab:}
  \cref{sec:lab_unit_III_bioenergetics_and_respiration}.
\item
  \textbf{Question bank:}
  \cref{sec:q_unit_III_bioenergetics_and_respiration}.
\item
  \textbf{Transfer task:} Transfer energy accounting to exercise,
  fermentation, hypoxia, and mitochondrial disease.
\item
  \textbf{Bridge to computation:}
  \breaktt{biology.biochemistry.biochemistry.glycolysis\_summary}.
\end{itemize}

\begin{center}\rule{0.5\linewidth}{0.5pt}\end{center}

\begin{quote}
\textbf{Opening Vignette: The Heretic Who Was Right About ATP}

In 1961, British biochemist Peter Mitchell proposed something that
nearly the entire biochemistry community considered preposterous: that
the synthesis of ATP was not driven by a hypothetical high-energy
chemical intermediate but by a \textbf{proton gradient across a
membrane} --- a \textbf{proton-motive force} \citep{mitchell1961}.
Mitchell's ``chemiosmotic hypothesis'' met with fierce resistance.
Efraim Racker, one of biochemistry's most respected voices, called it
``an act of faith, not of reason.'' Mitchell spent years defending his
idea from private funds after leaving a university position.

In 1978, the Nobel Committee awarded Peter Mitchell the Nobel Prize in
Chemistry --- alone, without co-recipients --- specifically for the
chemiosmotic hypothesis. It was a stunning vindication. Today we know
that the about 28 molecules of ATP produced per glucose in
oxidative phosphorylation are most driven by the proton gradient that
Mitchell described. The mitochondrion --- and the
\hyperref[gl:chloroplast]{\textbf{chloroplast}} --- are essentially
biological batteries, storing energy as a proton electrochemical
gradient. This chapter explains exactly how they are charged and
discharged.

\emph{Primary source: Mitchell, P. (1961). Coupling of phosphorylation
to electron and hydrogen transfer by a chemi-osmotic type of mechanism.
Nature, 191(4784), 144--148.}
\end{quote}

\begin{center}\rule{0.5\linewidth}{0.5pt}\end{center}

\subsection{Thermodynamics of Life}\label{thermodynamics-of-life}

\subsubsection{First and Second Laws}\label{first-and-second-laws}

The \textbf{First Law of Thermodynamics:} Energy cannot be created or
destroyed --- primarily converted. Life obeys this law: every metabolic
reaction converts one form of energy to another, with total energy
conserved.

The \textbf{Second Law of Thermodynamics:} In any spontaneous process,
the total entropy (disorder) of the universe increases. This means no
energy conversion is 100\% efficient --- some energy is typically
dispersed as heat. Living organisms maintain internal order (low
entropy) at the expense of increasing entropy in their surroundings
(heat release, CO\(_2\) production, waste generation).

A useful schematic: a cell is a \textbf{dissipative structure} in the
sense of Prigogine --- its internal order is maintained by a continuous
through-flux of free energy (food, sunlight) and a continuous out-flux
of low-grade heat. Stop the through-flux, and the cell relaxes to
thermal equilibrium --- which is to say, dies.

\subsubsection{Gibbs Free Energy: A Derivation From First
Principles}\label{gibbs-free-energy-a-derivation-from-first-principles}

For a process at constant temperature \(T\) and constant pressure \(P\)
(the conditions inside virtually every cell), the relevant thermodynamic
potential is the \textbf{Gibbs free energy} \(G\). Starting from the
second law, we have for the universe (system + surroundings):

\[\Delta S_{\text{universe}} = \Delta S_{\text{system}} + \Delta S_{\text{surroundings}} \geq 0 \tag{9.0a}\]

Heat released by the system at constant pressure equals
\(-\Delta H_{\text{system}}\), and this heat increases the entropy of
the surroundings by \(-\Delta H_{\text{system}}/T\). Substituting:

\[\Delta S_{\text{universe}} = \Delta S_{\text{system}} - \frac{\Delta H_{\text{system}}}{T} \geq 0 \tag{9.0b}\]

Multiplying both sides by \(-T\) and dropping the ``system'' subscript
yields the \textbf{defining inequality of spontaneity at constant \(T\),
\(P\)}:

\begin{equation}
\Delta G = \Delta H - T\Delta S \leq 0
\label{eq:unit_III_gibbs_free_energy}
\end{equation}

\begin{itemize}
\tightlist
\item
  \textbf{\(\Delta G < 0\):} reaction is spontaneous (exergonic); energy
  released
\item
  \textbf{\(\Delta G > 0\):} reaction is non-spontaneous (endergonic);
  energy input required
\item
  \textbf{\(\Delta G = 0\):} system is at equilibrium (no net reaction)
\end{itemize}

This is profound: a quantity that depends primarily on the \emph{system}
(not the surroundings) tells us whether the \emph{universe's} entropy
will increase. It is the workhorse of biochemical thermodynamics.

\subsubsection{Standard Free Energy and Actual Free
Energy}\label{standard-free-energy-and-actual-free-energy}

The standard free energy change (\(\Delta G^{\circ\prime}\)) is measured
under biochemical standard conditions (1 M concentrations for reactants
except H\(^+\), \hyperref[gl:ph]{\textbf{pH}} 7.0, 25 degrees C, 1 atm,
and water at unit activity). Under cellular conditions (not standard
state), \(\Delta G\) differs from \(\Delta G^{\circ\prime}\):

\[\Delta G = \Delta G^{\circ\prime} + RT \ln Q \tag{9.2}\]

where \(Q\) = {[}products{]}/{[}reactants{]} (mass action ratio,
reaction quotient).

The relationship between \(\Delta G^{\circ\prime}\) and the equilibrium
constant:

\[\Delta G^{\circ\prime} = -RT \ln K_{eq} \tag{9.3}\]

\textbf{Standard versus physiological \(\Delta G^{\circ\prime}\) for key
hydrolysis reactions} (37 °C, pH 7.0, 1 mM Mg\(^{2+}\)):

{\def\LTcaptype{none} % do not increment counter
\begin{longtable}[]{@{}
  >{\raggedright\arraybackslash}p{(\linewidth - 6\tabcolsep) * \real{0.1081}}
  >{\raggedright\arraybackslash}p{(\linewidth - 6\tabcolsep) * \real{0.4459}}
  >{\raggedright\arraybackslash}p{(\linewidth - 6\tabcolsep) * \real{0.3784}}
  >{\raggedright\arraybackslash}p{(\linewidth - 6\tabcolsep) * \real{0.0676}}@{}}
\toprule\noalign{}
\begin{minipage}[b]{\linewidth}\raggedright
Reaction
\end{minipage} & \begin{minipage}[b]{\linewidth}\raggedright
\(\Delta G^{\circ\prime}\) (kJ/mol)
\end{minipage} & \begin{minipage}[b]{\linewidth}\raggedright
Cellular \(\Delta G\) (kJ/mol)
\end{minipage} & \begin{minipage}[b]{\linewidth}\raggedright
Notes
\end{minipage} \\
\midrule\noalign{}
\endhead
\bottomrule\noalign{}
\endlastfoot
Phosphoenolpyruvate (PEP) \texorpdfstring{\ensuremath{\to}}{->} Pyruvate + P\(_i\) & \(-61.9\) & about
\(-65\) & Highest-energy phosphate bond in cells \\
1,3-Bisphosphoglycerate \texorpdfstring{\ensuremath{\to}}{->} 3-PG + P\(_i\) & \(-49.3\) & about \(-52\) &
Drives ATP synthesis at glycolysis step 7 \\
Phosphocreatine \texorpdfstring{\ensuremath{\to}}{->} Creatine + P\(_i\) & \(-43.0\) & about \(-47\) &
Muscle ATP buffer (phosphagen) \\
\textbf{ATP \texorpdfstring{\ensuremath{\to}}{->} ADP + P\(_i\)} & \textbf{\(-30.5\)} & \textbf{\(-50\) to
\(-54\)} & The ``energy currency benchmark'' \\
ATP \texorpdfstring{\ensuremath{\to}}{->} AMP + PP\(_i\) & \(-45.6\) & about \(-50\) & PP\(_i\) further
hydrolysed (irreversibility) \\
Glucose-1-P \texorpdfstring{\ensuremath{\to}}{->} Glucose + P\(_i\) & \(-20.9\) & \(-21\) & Low-energy
ester \\
Glucose-6-P \texorpdfstring{\ensuremath{\to}}{->} Glucose + P\(_i\) & \(-13.8\) & \(-14\) & Low-energy
phosphate ester \\
\end{longtable}
}

\textbf{Key insight:} ATP sits in the \emph{middle} of this spectrum. It
is energetically capable of phosphorylating many low-energy compounds
(glucose, F6P) but can itself be re-synthesised from higher-energy
intermediates such as PEP and 1,3-BPG. This intermediate position is
what makes ATP an effective ``energy currency'' rather than the
highest-potential molecule in the cell.

\textbf{Why cells maintain far-from-equilibrium conditions:} At
equilibrium, \(\Delta G = 0\) and no net work can be done. Cells
maintain metabolite concentrations far from equilibrium
(\(Q \neq K_{eq}\)), keeping \(\Delta G\) strongly negative for
exergonic reactions. This is why cells must continuously consume
nutrients --- to maintain the far-from-equilibrium state that enables
life.

\subsubsection{Worked Example: Free Energy and
Equilibrium}\label{worked-example-free-energy-and-equilibrium}

\emph{Problem:} The \(\Delta G^{\circ\prime}\) for glucose-6-phosphate
isomerase (G6P to F6P) is +1.7 kJ/mol. Calculate \(K_{eq}\) at 37
degrees C. Is this reaction spontaneous under standard conditions?

\emph{Solution:}

\[\Delta G^{\circ\prime} = -RT \ln K_{eq} \tag{9.4}\]

\[1{,}700 = -(8.314)(310) \ln K_{eq} \tag{9.5}\]

\[\ln K_{eq} = \frac{-1{,}700}{2{,}577} = -0.660 \tag{9.6}\]

\[K_{eq} = e^{-0.660} = 0.517 \tag{9.7}\]

Under standard conditions, this reaction is slightly non-spontaneous
(\(\Delta G^{\circ\prime} > 0\)). However, in the cell, the mass action
ratio \(Q\) is kept well below \(K_{eq}\) because F6P is rapidly
consumed by PFK-1, making the actual \(\Delta G\) negative.

\subsubsection{Worked Example: Complete Glucose
Oxidation}\label{worked-example-complete-glucose-oxidation}

\emph{Problem:} Calculate the standard free energy of complete glucose
oxidation (glucose + 6 O\(_2\) \texorpdfstring{\ensuremath{\to}}{->} 6 CO\(_2\) + 6 H\(_2\)O) from tabulated
standard free energies of formation, and compare it with the energy
harvested as ATP.

\emph{Solution:} Using \(\Delta G^{\circ}_f\) values (kJ/mol) at 25 °C
--- glucose: \(-910.4\); CO\(_2\) (aq): \(-394.4\); H\(_2\)O (l):
\(-237.1\); O\(_2\): 0:

\[\Delta G^{\circ\prime}_{\text{rxn}} = \sum \Delta G^{\circ}_f(\text{products}) - \sum \Delta G^{\circ}_f(\text{reactants}) \tag{9.7a}\]

\[= [6(-394.4) + 6(-237.1)] - [(-910.4) + 6(0)] = -3{,}789 + 910.4 \approx -2{,}870 \; \text{kJ/mol} \tag{9.7b}\]

If 32 ATP are produced per glucose at a cellular hydrolysis \(\Delta G\)
of \(-50\) kJ/mol, the energy captured as ATP is
\(32 \times 50 = 1{,}600\) kJ/mol. The thermodynamic efficiency of
glucose oxidation in cells is therefore:

\[\eta = \frac{1{,}600}{2{,}870} \approx 56\% \tag{9.7c}\]

The remaining about 44\% is dissipated as heat --- useful for
mammals to maintain body temperature, but ultimately a payment to the
second law. For comparison: a modern combined-cycle gas turbine power
plant achieves about 60\% thermal efficiency at \textgreater800
K, while an automobile engine averages about 25\%. A
56\%-efficient isothermal energy converter operating at 310 K is, by any
engineering standard, remarkable.

\subsubsection{ATP as Energy Currency}\label{atp-as-energy-currency}

ATP (adenosine triphosphate) is the cell's central short-term energy
carrier and phosphoryl-group donor for many endergonic reactions. The
phrase ``energy currency'' is useful shorthand, but the cell also uses
ion gradients, redox cofactors such as NADH/NADPH, thioesters such as
acetyl-CoA, and activated carriers such as UDP-glucose.

The high negative \(\Delta G\) of ATP hydrolysis under cellular
conditions (\(-50\) to \(-54\) kJ/mol;
\(\Delta G^{\circ\prime} = -30.5\) kJ/mol) arises from: 1.
\textbf{Electrostatic repulsion} between closely spaced negative charges
on the gamma-phosphate (relieved on hydrolysis) 2. \textbf{Resonance
stabilisation} of inorganic phosphate (P\(_i\)) product --- P\(_i\) has
more resonance forms than the terminal phosphate of ATP 3.
\textbf{Hydration} (solvation enthalpy) of P\(_i\) and ADP --- products
are better hydrated than reactants 4. \textbf{Entropy increase} --- one
molecule becomes two

ATP couples exergonic reactions to endergonic ones via
\textbf{phosphoryl group transfer}. The cell maintains a very high
ATP/ADP ratio (about 10 in most cells; ATP about 3--5
mM, ADP about 0.5 mM), keeping the cellular \(\Delta G\) of
hydrolysis more negative than \(\Delta G^{\circ\prime}\).

\textbf{Daily ATP turnover:} A resting adult uses about 40 kg
of ATP per day (body contains about 50 g at any time, so each ATP
molecule is recycled about 800 times/day). During vigorous
exercise, ATP turnover can reach about 0.5 kg/min --- roughly
the body's mass in ATP equivalents flowing through the adenylate pool
every two hours.

\begin{quote}
\textbf{Concept Check 1:} The \(\Delta G^{\circ\prime}\) of ATP
hydrolysis is \(-30.5\) kJ/mol, but the actual cellular \(\Delta G\) is
approximately \(-54\) kJ/mol. Using the equation
\(\Delta G = \Delta G^{\circ\prime} + RT \ln Q\), calculate the mass
action ratio \(Q\) = {[}ADP{]}{[}P\(_i\){]}/{[}ATP{]} in the cell at 37
degrees C.
\end{quote}

\begin{quote}
\textbf{Concept Check 1b:} A cell experiences a sudden drop in ATP (from
4 mM to 1 mM) without changes in ADP or P\(_i\). Predict the qualitative
effect on (a) the cellular \(\Delta G\) of ATP hydrolysis, (b) the
ATP/ADP ratio, and (c) the activity of AMPK.
\end{quote}

\begin{center}\rule{0.5\linewidth}{0.5pt}\end{center}

\subsection{Glycolysis}\label{glycolysis}

Glycolysis (Greek: \emph{glykys} = sweet; \emph{lysis} = splitting)
converts one glucose (6C) to two pyruvate (3C), netting 2 ATP and 2
NADH. It occurs in the \textbf{\hyperref[gl:cytoplasm]{cytoplasm}} of
most cells and does not require oxygen.

\begin{figure}[htbp]
\centering
\includegraphics[width=0.82\linewidth,height=4.2in,keepaspectratio]{../figures/mermaid_inline/inline_mermaid_0048_843fd47ba6b4.png}
\caption{Flowchart for Glycolysis: Hexokinase ATPtoADP DeltaG = 16.7 kJ/mol IRREVERSIBLE"| G6P["Glucose-6-phosphate"], Glucose (6C, Glucose-6-phosphate, and Phosphoglucose isomerase DeltaG = +1.7"| F6P["Fructose-6-phosphate"] identify the main nodes and their directional relationships.}
\end{figure}

\emph{The 10 steps of glycolysis (Mermaid).} The energy investment phase
(steps 1--5) consumes 2 ATP. The energy payoff phase (steps 6--10)
produces 4 ATP and 2 NADH. The three irreversible steps (1, 3, 10) are
the regulatory control points.

\subsubsection{Ten Enzymatic Steps}\label{ten-enzymatic-steps}

{\def\LTcaptype{none} % do not increment counter
\begin{longtable}[]{@{}
  >{\raggedright\arraybackslash}p{(\linewidth - 8\tabcolsep) * \real{0.0952}}
  >{\raggedright\arraybackslash}p{(\linewidth - 8\tabcolsep) * \real{0.1429}}
  >{\raggedright\arraybackslash}p{(\linewidth - 8\tabcolsep) * \real{0.1905}}
  >{\raggedright\arraybackslash}p{(\linewidth - 8\tabcolsep) * \real{0.4524}}
  >{\raggedright\arraybackslash}p{(\linewidth - 8\tabcolsep) * \real{0.1190}}@{}}
\toprule\noalign{}
\begin{minipage}[b]{\linewidth}\raggedright
Step
\end{minipage} & \begin{minipage}[b]{\linewidth}\raggedright
Enzyme
\end{minipage} & \begin{minipage}[b]{\linewidth}\raggedright
Reaction
\end{minipage} & \begin{minipage}[b]{\linewidth}\raggedright
\(\Delta G\) (kJ/mol)
\end{minipage} & \begin{minipage}[b]{\linewidth}\raggedright
Notes
\end{minipage} \\
\midrule\noalign{}
\endhead
\bottomrule\noalign{}
\endlastfoot
1 & Hexokinase & Glucose + ATP \texorpdfstring{\ensuremath{\to}}{->} G6P + ADP & \(-16.7\) & Traps glucose
in cell \\
2 & Phosphoglucose isomerase & G6P \texorpdfstring{\ensuremath{\to}}{->} F6P & \(+1.7\) & Aldose to
ketose \\
3 & \textbf{PFK-1} & F6P + ATP \texorpdfstring{\ensuremath{\to}}{->} F1,6BP + ADP & \(-14.2\) &
\textbf{Committed step; rate-limiting} \\
4 & Aldolase & F1,6BP \texorpdfstring{\ensuremath{\to}}{->} DHAP + G3P & \(+23.8\) & Pulled forward by step
5 \\
5 & Triose phosphate isomerase & DHAP \texorpdfstring{\ensuremath{\to}}{->} G3P & \(+7.5\) & ``Perfect
enzyme'' (\(k_{cat}/K_m\) near diffusion limit) \\
6 & GAPDH & G3P + P\(_i\) + NAD\(^+\) \texorpdfstring{\ensuremath{\to}}{->} 1,3BPG + NADH & \(-6.3\) & First
NADH production \\
7 & Phosphoglycerate kinase & 1,3BPG + ADP \texorpdfstring{\ensuremath{\to}}{->} 3PG + ATP & \(-18.5\) &
First ATP production \\
8 & Phosphoglycerate mutase & 3PG \texorpdfstring{\ensuremath{\to}}{->} 2PG & \(+4.4\) & Phosphate shift \\
9 & Enolase & 2PG \texorpdfstring{\ensuremath{\to}}{->} PEP + H\(_2\)O & \(+1.8\) & Dehydration \\
10 & Pyruvate kinase & PEP + ADP \texorpdfstring{\ensuremath{\to}}{->} Pyruvate + ATP & \(-31.4\) & Largest
energy release \\
\end{longtable}
}

\subsubsection{Regulation of Glycolysis}\label{regulation-of-glycolysis}

\textbf{PFK-1 is the master regulator} of glycolysis. It integrates
energy status, biosynthetic needs, and hormonal signals
\citep{atkinson1968}:

{\def\LTcaptype{none} % do not increment counter
\begin{longtable}[]{@{}
  >{\raggedright\arraybackslash}p{(\linewidth - 2\tabcolsep) * \real{0.5000}}
  >{\raggedright\arraybackslash}p{(\linewidth - 2\tabcolsep) * \real{0.5000}}@{}}
\toprule\noalign{}
\begin{minipage}[b]{\linewidth}\raggedright
Activators
\end{minipage} & \begin{minipage}[b]{\linewidth}\raggedright
Inhibitors
\end{minipage} \\
\midrule\noalign{}
\endhead
\bottomrule\noalign{}
\endlastfoot
AMP (low energy) & ATP (high energy,
\hyperref[gl:allosteric]{\textbf{allosteric}} at regulatory site) \\
Fructose-2,6-bisphosphate (F2,6BP, most potent) & Citrate (TCA cycle
intermediate = adequate supply) \\
ADP & H\(^+\) (low pH; prevents excessive lactate acidosis) \\
P\(_i\) & \\
\end{longtable}
}

\textbf{Fructose-2,6-bisphosphate (F2,6BP)} is the most potent activator
of PFK-1 (lowers \(K_m\) for F6P by 50-fold). It is produced by
\textbf{PFK-2} (a bifunctional enzyme with both kinase and phosphatase
domains). Insulin activates PFK-2 kinase activity (producing F2,6BP,
stimulating glycolysis). Glucagon (via PKA) activates PFK-2 phosphatase
activity (destroying F2,6BP, inhibiting glycolysis in liver).

\subsubsection{Energy Accounting}\label{energy-accounting}

Net output from 1 glucose: - 2 ATP (from steps 7 and 10, minus 2 ATP
consumed at steps 1 and 3) - 2 NADH (from step 6; each NADH worth
about 2.5 ATP via oxidative phosphorylation in
\hyperref[gl:aerobic]{\textbf{aerobic}} conditions)

\begin{quote}
\textbf{Clinical Connection: The Warburg Effect} Cancer cells often
exhibit the \textbf{Warburg effect} --- high rates of glycolysis even in
the presence of oxygen (``aerobic glycolysis''). Rather than fully
oxidising glucose to CO\(_2\) via the TCA cycle and ETC, cancer cells
convert most glucose to lactate. This seems wasteful (2 ATP vs.~30--32
ATP per glucose), but it provides rapid ATP production and generates
biosynthetic precursors (G6P for pentose phosphate pathway, 3PG for
serine biosynthesis) needed for rapid cell division. PET scanning
exploits this by detecting tumour uptake of
\(^{18}\)F-fluorodeoxyglucose (FDG). see
\cref{sec:unit_III_metabolic_integration} for metabolic integration.
\end{quote}

\begin{quote}
\textbf{Concept Check 2:} Steps 4 and 5 of glycolysis have very positive
\(\Delta G^{\circ\prime}\) values (+23.8 and +7.5 kJ/mol). How do these
reactions proceed in vivo despite being thermodynamically unfavourable
under standard conditions?
\end{quote}

\begin{center}\rule{0.5\linewidth}{0.5pt}\end{center}

\subsection{Gluconeogenesis: Reverse But Not the
Reverse}\label{gluconeogenesis-reverse-but-not-the-reverse}

Gluconeogenesis is \emph{not} glycolysis run backwards. Three glycolytic
steps (hexokinase, PFK-1, pyruvate kinase) are \textbf{physiologically
irreversible} in the forward direction --- their cellular \(\Delta G\)
values are too negative. Gluconeogenesis bypasses each of these three
steps with dedicated enzymes that use different chemistry, often paying
additional ATP/GTP to make the reverse direction thermodynamically
favourable.

\subsubsection{The Four Bypass
Reactions}\label{the-four-bypass-reactions}

\begin{figure}[htbp]
\centering
\includegraphics[width=0.82\linewidth,height=4.2in,keepaspectratio]{../figures/mermaid_inline/inline_mermaid_0049_c71600fa6e4f.png}
\caption{Flowchart for The Four Bypass Reactions: Pyruvate (mitochondrion), Oxaloacetate, Malate, and Malate (cytoplasm) form the diagram's primary path or branches.}
\end{figure}

\emph{Gluconeogenesis from pyruvate to glucose (Mermaid).} Four bypass
enzymes (pyruvate carboxylase, PEPCK, FBPase-1, G6Pase) replace the
three irreversible glycolytic steps. The pyruvate-to-PEP conversion
requires two enzymes --- pyruvate carboxylase consumes ATP and PEPCK
consumes GTP --- so the bypass spans two energy-consuming steps.

\subsubsection{The Four Bypass Enzymes in
Detail}\label{the-four-bypass-enzymes-in-detail}

\textbf{1. Pyruvate carboxylase (PC):} A mitochondrial biotin-dependent
carboxylase that converts pyruvate + CO\(_2\) + ATP \texorpdfstring{\ensuremath{\to}}{->} oxaloacetate + ADP
+ P\(_i\). Two key features: (a) \textbf{allosterically activated by
acetyl-CoA}, providing a feed-forward signal that fat oxidation (which
produces acetyl-CoA) supplies the OAA needed for both gluconeogenesis
and TCA flux; (b) requires biotin as a covalently bound prosthetic group
that carries the activated CO\(_2\) (--N--COO\(^-\)). Biotin deficiency
(rare; most often from chronic raw-egg consumption due to avidin)
disrupts pyruvate carboxylase and three other carboxylases
simultaneously, producing lactic acidosis and ketogenic dysregulation.

\textbf{2. PEPCK (phosphoenolpyruvate carboxykinase):} A cytosolic (and
partially mitochondrial) GTP-dependent decarboxylase that converts OAA +
GTP \texorpdfstring{\ensuremath{\to}}{->} PEP + CO\(_2\) + GDP. This enzyme is unusual among gluconeogenic
enzymes in that it is regulated almost entirely at the level of
\emph{transcription} --- glucagon (via cAMP\texorpdfstring{\ensuremath{\to}}{->}PKA\texorpdfstring{\ensuremath{\to}}{->}CREB) and cortisol (via
the glucocorticoid receptor) potently induce PEPCK mRNA, while insulin
(via FoxO1 phosphorylation and exclusion from the nucleus) potently
represses it. PEPCK transcription has a half-time of about 1
hour, making it a slow but powerful long-term switch.

\textbf{3. Fructose-1,6-bisphosphatase (FBPase-1):} Cytosolic
phosphatase that converts F1,6BP + H\(_2\)O \texorpdfstring{\ensuremath{\to}}{->} F6P + P\(_i\)
(\(\Delta G^{\circ\prime} = -16.3\) kJ/mol). Critically regulated
\textbf{allosterically} by F2,6BP (strong inhibitor) and AMP (inhibitor)
--- both signals of glycolytic demand or low energy. The same F2,6BP
that \emph{activates} PFK-1 \emph{inhibits} FBPase-1, so a single second
messenger flips the entire futile-cycle node.

\textbf{4. Glucose-6-phosphatase (G6Pase):} A 9-transmembrane-helix
protein whose catalytic site faces the \textbf{lumen of the endoplasmic
reticulum}, accessed by a phosphate transporter (SLC37A4) and a glucose
efflux transporter (GLUT2). Critically, G6Pase is expressed almost
exclusively in \textbf{liver and renal cortex} --- muscle and brain lack
it, which is why these tissues cannot release free glucose to the blood
even though they produce G6P internally. Mutations in G6Pase cause
\textbf{von Gierke disease (GSD I)}: severe fasting hypoglycaemia with
hepatomegaly.

{\def\LTcaptype{none} % do not increment counter
\begin{longtable}[]{@{}
  >{\raggedright\arraybackslash}p{(\linewidth - 6\tabcolsep) * \real{0.4167}}
  >{\raggedright\arraybackslash}p{(\linewidth - 6\tabcolsep) * \real{0.2778}}
  >{\raggedright\arraybackslash}p{(\linewidth - 6\tabcolsep) * \real{0.1528}}
  >{\raggedright\arraybackslash}p{(\linewidth - 6\tabcolsep) * \real{0.1528}}@{}}
\toprule\noalign{}
\begin{minipage}[b]{\linewidth}\raggedright
Glycolytic step (irreversible)
\end{minipage} & \begin{minipage}[b]{\linewidth}\raggedright
Gluconeogenic bypass
\end{minipage} & \begin{minipage}[b]{\linewidth}\raggedright
Energy cost
\end{minipage} & \begin{minipage}[b]{\linewidth}\raggedright
Compartment
\end{minipage} \\
\midrule\noalign{}
\endhead
\bottomrule\noalign{}
\endlastfoot
Pyruvate kinase & \textbf{Pyruvate carboxylase} + \textbf{PEPCK} & +1
ATP + 1 GTP per pyruvate & Mitochondrion + cytoplasm \\
PFK-1 & \textbf{Fructose-1,6-bisphosphatase (FBPase-1)} & None
(hydrolysis is exergonic) & Cytoplasm \\
Hexokinase & \textbf{Glucose-6-phosphatase (G6Pase)} & None (hydrolysis
is exergonic) & ER lumen (liver, kidney primarily) \\
\end{longtable}
}

\textbf{Net cost:} Synthesising one glucose from two pyruvate consumes
\textbf{6 ATP equivalents} (2 ATP at pyruvate carboxylase \texorpdfstring{\ensuremath{\times}}{x} 2, 2 GTP at
PEPCK \texorpdfstring{\ensuremath{\times}}{x} 2, 2 ATP at PGK \texorpdfstring{\ensuremath{\times}}{x} 2 in the reverse direction). This is why
gluconeogenesis is restricted to tissues with substantial ATP-generating
capacity (liver, renal cortex).

\subsubsection{Reciprocal Regulation}\label{reciprocal-regulation}

The same allosteric and hormonal signals that activate glycolysis
inhibit gluconeogenesis, ensuring the two pathways do not run as a
futile cycle:

{\def\LTcaptype{none} % do not increment counter
\begin{longtable}[]{@{}
  >{\raggedright\arraybackslash}p{(\linewidth - 4\tabcolsep) * \real{0.1200}}
  >{\raggedright\arraybackslash}p{(\linewidth - 4\tabcolsep) * \real{0.3600}}
  >{\raggedright\arraybackslash}p{(\linewidth - 4\tabcolsep) * \real{0.5200}}@{}}
\toprule\noalign{}
\begin{minipage}[b]{\linewidth}\raggedright
Signal
\end{minipage} & \begin{minipage}[b]{\linewidth}\raggedright
Glycolysis (PFK-1)
\end{minipage} & \begin{minipage}[b]{\linewidth}\raggedright
Gluconeogenesis (FBPase-1)
\end{minipage} \\
\midrule\noalign{}
\endhead
\bottomrule\noalign{}
\endlastfoot
Fructose-2,6-bisphosphate & Strong activator & Strong inhibitor \\
AMP & Activator & Inhibitor \\
Citrate & Inhibitor & Activator \\
Acetyl-CoA & (no direct effect) & Activates pyruvate carboxylase \\
Insulin & \texorpdfstring{\ensuremath{\uparrow}}{up} PFK-2 kinase \texorpdfstring{\ensuremath{\to}}{->} \texorpdfstring{\ensuremath{\uparrow}}{up} F2,6BP \texorpdfstring{\ensuremath{\to}}{->} glycolysis ON & \texorpdfstring{\ensuremath{\downarrow}}{down} PEPCK / G6Pase
transcription \\
Glucagon & \texorpdfstring{\ensuremath{\uparrow}}{up} PFK-2 phosphatase \texorpdfstring{\ensuremath{\to}}{->} \texorpdfstring{\ensuremath{\downarrow}}{down} F2,6BP \texorpdfstring{\ensuremath{\to}}{->} glycolysis OFF & \texorpdfstring{\ensuremath{\uparrow}}{up} PEPCK /
G6Pase transcription \\
\end{longtable}
}

Glucagon-driven phosphorylation of PFK-2/FBPase-2 simultaneously turns
\emph{off} PFK-1 (glycolysis) and turns \emph{on} FBPase-1
(gluconeogenesis) by destroying F2,6BP. This is a textbook example of
how a single bifunctional enzyme can act as a metabolic switch.

\begin{quote}
\textbf{Concept Check 2b:} A patient with renal cortex damage but intact
liver presents with mild fasting hypoglycaemia after prolonged exercise.
Explain the contribution of renal gluconeogenesis to whole-body glucose
homeostasis, and why it is invisible in healthy individuals.
\end{quote}

\begin{center}\rule{0.5\linewidth}{0.5pt}\end{center}

\subsection{The Cori Cycle and the Glucose-Alanine
Cycle}\label{the-cori-cycle-and-the-glucose-alanine-cycle}

Inter-organ metabolism allows tissues to specialise: muscle generates
ATP rapidly via glycolysis, while liver shoulders the gluconeogenic
burden. Two complementary inter-organ shuttles dominate fasting and
exercise physiology.

\subsubsection{The Cori Cycle (Lactate
Shuttle)}\label{the-cori-cycle-lactate-shuttle}

\begin{figure}[htbp]
\centering
\includegraphics[width=0.82\linewidth,height=4.2in,keepaspectratio]{../figures/mermaid_inline/inline_mermaid_0050_bd9d9e663fc5.png}
\caption{Flowchart for The Cori Cycle (Lactate Shuttle): Skeletal muscle / RBCs, Glucose, Lactate, and Bloodstream form the diagram's primary path or branches.}
\end{figure}

\emph{Cori cycle: lactate from anaerobic muscle is regenerated to
glucose in the liver (Mermaid).} Each turn nets +2 ATP for muscle but \texorpdfstring{\ensuremath{-}}{-}6
ATP for liver, a 4-ATP net loss to the organism --- the metabolic cost
of supporting hypoxic tissue.

\subsubsection{The Glucose-Alanine
Cycle}\label{the-glucose-alanine-cycle}

When muscle proteins are catabolised, branched-chain amino acid carbon
skeletons donate amino groups to pyruvate, forming \textbf{alanine}.
Alanine carries both carbon and nitrogen safely to the liver:

\begin{enumerate}
\def\labelenumi{\arabic{enumi}.}
\tightlist
\item
  Muscle: Pyruvate + Glutamate \texorpdfstring{\ensuremath{\to}}{->} Alanine + \texorpdfstring{\ensuremath{\alpha}}{alpha}-Ketoglutarate (alanine
  aminotransferase)
\item
  Alanine travels via blood to liver
\item
  Liver: Alanine + \texorpdfstring{\ensuremath{\alpha}}{alpha}-KG \texorpdfstring{\ensuremath{\to}}{->} Pyruvate + Glutamate (alanine
  aminotransferase, reverse)
\item
  Glutamate releases NH\(_4^+\) into the urea cycle; pyruvate enters
  gluconeogenesis
\item
  Glucose returns to muscle via blood
\end{enumerate}

\begin{figure}[htbp]
\centering
\includegraphics[width=0.82\linewidth,height=4.2in,keepaspectratio]{../figures/mermaid_inline/inline_mermaid_0051_fd7bc334878a.png}
\caption{Flowchart for The Glucose-Alanine Cycle: Skeletal muscle, Branched-chain amino acid catabolism, -NH2 to pyruvate, and Pyruvate form the diagram's primary path or branches.}
\end{figure}

\emph{Glucose-alanine cycle: alanine simultaneously carries
gluconeogenic carbon and waste nitrogen from muscle to liver (Mermaid).}
The liver disposes of NH\(_4^+\) via the urea cycle and exports glucose
back to muscle.

\textbf{Why this is clever:} Alanine simultaneously transports
gluconeogenic carbon and waste nitrogen, replacing the need for two
separate transport systems. It also keeps free ammonia (toxic to the
brain) out of circulation --- nitrogen is delivered directly to the
urea-cycling liver.

\begin{quote}
\textbf{Concept Check 3:} During an acute sprint, muscle lactate output
rises from 1 mM/min to 8 mM/min. Estimate the additional hepatic ATP
demand (in ATP equivalents per minute) required to recycle this lactate
via the Cori cycle, assuming most lactate reaches the liver.
\end{quote}

\begin{center}\rule{0.5\linewidth}{0.5pt}\end{center}

\subsection{Pyruvate Oxidation and the Citric Acid
Cycle}\label{pyruvate-oxidation-and-the-citric-acid-cycle}

Under aerobic conditions, pyruvate enters the mitochondria (via the MPC,
mitochondrial pyruvate carrier) and is oxidised.

\subsubsection{Pyruvate Dehydrogenase Complex
(PDC)}\label{pyruvate-dehydrogenase-complex-pdc}

\[\text{Pyruvate (3C)} + \text{CoA} + \text{NAD}^+ \rightarrow \text{Acetyl-CoA (2C)} + \text{CO}_2 + \text{NADH} \tag{9.8}\]

Catalysed by PDC --- a 9.5 MDa multienzyme complex with three catalytic
components (E1: pyruvate dehydrogenase; E2: dihydrolipoyl
transacetylase; E3: dihydrolipoyl dehydrogenase) using \textbf{five
cofactors}: TPP (thiamine pyrophosphate, vitamin B\(_1\)), lipoic acid,
FAD (vitamin B\(_2\)), NAD\(^+\) (vitamin B\(_3\)), CoA (vitamin
B\(_5\)).

\textbf{PDC regulation:}

{\def\LTcaptype{none} % do not increment counter
\begin{longtable}[]{@{}
  >{\raggedright\arraybackslash}p{(\linewidth - 4\tabcolsep) * \real{0.3000}}
  >{\raggedright\arraybackslash}p{(\linewidth - 4\tabcolsep) * \real{0.4000}}
  >{\raggedright\arraybackslash}p{(\linewidth - 4\tabcolsep) * \real{0.3000}}@{}}
\toprule\noalign{}
\begin{minipage}[b]{\linewidth}\raggedright
Condition
\end{minipage} & \begin{minipage}[b]{\linewidth}\raggedright
PDC activity
\end{minipage} & \begin{minipage}[b]{\linewidth}\raggedright
Mechanism
\end{minipage} \\
\midrule\noalign{}
\endhead
\bottomrule\noalign{}
\endlastfoot
High acetyl-CoA, NADH, ATP & Inhibited & PDH kinase (PDK) phosphorylates
E1 \texorpdfstring{\ensuremath{\to}}{->} inactive \\
High CoA, NAD\(^+\), AMP, Ca\(^{2+}\) & Activated & PDH phosphatase
(PDP) dephosphorylates E1 \texorpdfstring{\ensuremath{\to}}{->} active \\
Insulin & Activated & Activates PDP in adipose tissue \\
Exercise (Ca\(^{2+}\)) & Activated & Ca\(^{2+}\) activates PDP \\
\end{longtable}
}

\begin{quote}
\textbf{Clinical Connection: PDC Deficiency} Pyruvate dehydrogenase
deficiency is the most common genetic cause of lactic acidosis. Without
PDC, pyruvate cannot enter the TCA cycle and is instead reduced to
lactate. Neurological damage results from the brain's dependence on
glucose oxidation. Treatment: \textbf{ketogenic diet} (high fat, very
low carbohydrate) provides acetyl-CoA directly from fatty acid
beta-oxidation, bypassing the PDC block. Dichloroacetate (DCA) inhibits
PDK, keeping residual PDC active.
\end{quote}

\subsubsection{TCA Cycle (Krebs Cycle / Citric Acid
Cycle)}\label{tca-cycle-krebs-cycle-citric-acid-cycle}

The TCA cycle occurs in the mitochondrial matrix. It oxidises acetyl-CoA
to CO\(_2\), generating reduced cofactors (NADH, FADH\(_2\)) that feed
the ETC.

\begin{figure}[htbp]
\centering
\includegraphics[width=0.82\linewidth,height=4.2in,keepaspectratio]{../figures/mermaid_inline/inline_mermaid_0052_4baeb11d3a60.png}
\caption{Flowchart for TCA Cycle (Krebs Cycle / Citric Acid Cycle): Oxaloacetate (4C), Citrate (6C), Isocitrate (6C), and alpha-Ketoglutarate (5C) form the diagram's primary path or branches.}
\end{figure}

\emph{The eight steps of the TCA cycle (Mermaid).} Each turn oxidises
one acetyl group (2C) to 2 CO\(_2\), producing 3 NADH, 1 FADH\(_2\), and
1 GTP. The cycle turns twice per glucose molecule.

\textbf{Eight steps with details:}

{\def\LTcaptype{none} % do not increment counter
\begin{longtable}[]{@{}
  >{\raggedright\arraybackslash}p{(\linewidth - 6\tabcolsep) * \real{0.1538}}
  >{\raggedright\arraybackslash}p{(\linewidth - 6\tabcolsep) * \real{0.2308}}
  >{\raggedright\arraybackslash}p{(\linewidth - 6\tabcolsep) * \real{0.3077}}
  >{\raggedright\arraybackslash}p{(\linewidth - 6\tabcolsep) * \real{0.3077}}@{}}
\toprule\noalign{}
\begin{minipage}[b]{\linewidth}\raggedright
Step
\end{minipage} & \begin{minipage}[b]{\linewidth}\raggedright
Enzyme
\end{minipage} & \begin{minipage}[b]{\linewidth}\raggedright
Reaction
\end{minipage} & \begin{minipage}[b]{\linewidth}\raggedright
Products
\end{minipage} \\
\midrule\noalign{}
\endhead
\bottomrule\noalign{}
\endlastfoot
1 & Citrate synthase & Acetyl-CoA + OAA \texorpdfstring{\ensuremath{\to}}{->} Citrate + CoA & -- \\
2 & Aconitase & Citrate \texorpdfstring{\ensuremath{\to}}{->} Isocitrate (via cis-aconitate) & -- \\
3 & \textbf{Isocitrate dehydrogenase} & Isocitrate \texorpdfstring{\ensuremath{\to}}{->} alpha-KG + CO\(_2\)
& \textbf{NADH} \\
4 & \textbf{alpha-KG dehydrogenase complex} & alpha-KG \texorpdfstring{\ensuremath{\to}}{->} Succinyl-CoA +
CO\(_2\) & \textbf{NADH} \\
5 & Succinyl-CoA synthetase & Succinyl-CoA \texorpdfstring{\ensuremath{\to}}{->} Succinate & \textbf{GTP} \\
6 & Succinate dehydrogenase (Complex II) & Succinate \texorpdfstring{\ensuremath{\to}}{->} Fumarate &
\textbf{FADH\(_2\)} \\
7 & Fumarase & Fumarate + H\(_2\)O \texorpdfstring{\ensuremath{\to}}{->} Malate & -- \\
8 & Malate dehydrogenase & Malate \texorpdfstring{\ensuremath{\to}}{->} OAA & \textbf{NADH} \\
\end{longtable}
}

\textbf{Carbon tracking:} The two carbons entering as acetyl-CoA are NOT
the same two carbons lost as CO\(_2\) in the first turn --- they are
released in subsequent turns. This can be demonstrated using
\(^{14}\)C-labelled acetyl-CoA.

\textbf{Net per turn:} 3 NADH + 1 FADH\(_2\) + 1 GTP + 2 CO\(_2\)

\textbf{Total from 2 turns (per glucose):} 6 NADH + 2 FADH\(_2\) + 2 GTP
+ 4 CO\(_2\)

\textbf{TCA cycle regulation:}

{\def\LTcaptype{none} % do not increment counter
\begin{longtable}[]{@{}
  >{\raggedright\arraybackslash}p{(\linewidth - 4\tabcolsep) * \real{0.2308}}
  >{\raggedright\arraybackslash}p{(\linewidth - 4\tabcolsep) * \real{0.3846}}
  >{\raggedright\arraybackslash}p{(\linewidth - 4\tabcolsep) * \real{0.3846}}@{}}
\toprule\noalign{}
\begin{minipage}[b]{\linewidth}\raggedright
Enzyme
\end{minipage} & \begin{minipage}[b]{\linewidth}\raggedright
Activators
\end{minipage} & \begin{minipage}[b]{\linewidth}\raggedright
Inhibitors
\end{minipage} \\
\midrule\noalign{}
\endhead
\bottomrule\noalign{}
\endlastfoot
Citrate synthase & OAA, acetyl-CoA (substrates) & ATP, NADH,
succinyl-CoA, citrate \\
Isocitrate DH & ADP, Ca\(^{2+}\) & ATP, NADH \\
alpha-KG DH & Ca\(^{2+}\), AMP & ATP, NADH, succinyl-CoA \\
\end{longtable}
}

\subsubsection{Anaplerotic and Cataplerotic
Reactions}\label{anaplerotic-and-cataplerotic-reactions}

The TCA cycle is not a closed loop in vivo --- intermediates flow in and
out for biosynthesis. \textbf{Anaplerotic} (``filling-up'') reactions
replenish drained intermediates; \textbf{cataplerotic} (``emptying'')
reactions remove intermediates for downstream pathways.

\textbf{Major anaplerotic inputs:}

{\def\LTcaptype{none} % do not increment counter
\begin{longtable}[]{@{}
  >{\raggedright\arraybackslash}p{(\linewidth - 8\tabcolsep) * \real{0.2222}}
  >{\raggedright\arraybackslash}p{(\linewidth - 8\tabcolsep) * \real{0.1667}}
  >{\raggedright\arraybackslash}p{(\linewidth - 8\tabcolsep) * \real{0.2222}}
  >{\raggedright\arraybackslash}p{(\linewidth - 8\tabcolsep) * \real{0.1667}}
  >{\raggedright\arraybackslash}p{(\linewidth - 8\tabcolsep) * \real{0.2222}}@{}}
\toprule\noalign{}
\begin{minipage}[b]{\linewidth}\raggedright
Reaction
\end{minipage} & \begin{minipage}[b]{\linewidth}\raggedright
Enzyme
\end{minipage} & \begin{minipage}[b]{\linewidth}\raggedright
Cofactor
\end{minipage} & \begin{minipage}[b]{\linewidth}\raggedright
Tissue
\end{minipage} & \begin{minipage}[b]{\linewidth}\raggedright
Function
\end{minipage} \\
\midrule\noalign{}
\endhead
\bottomrule\noalign{}
\endlastfoot
Pyruvate + CO\(_2\) + ATP \texorpdfstring{\ensuremath{\to}}{->} OAA & \textbf{Pyruvate carboxylase} & Biotin
(vitamin B\(_7\)) & Liver, kidney & Replenishes OAA; primary anaplerotic
source \\
Glutamate \texorpdfstring{\ensuremath{\to}}{->} \texorpdfstring{\ensuremath{\alpha}}{alpha}-KG + NH\(_4^+\) & Glutamate dehydrogenase & NAD(P)\(^+\) &
Liver & Amino-acid catabolism \\
Glutamine \texorpdfstring{\ensuremath{\to}}{->} Glutamate + NH\(_4^+\) \texorpdfstring{\ensuremath{\to}}{->} \texorpdfstring{\ensuremath{\alpha}}{alpha}-KG & Glutaminase + GDH & -- &
Most tissues & Amino-acid catabolism \\
Aspartate \texorpdfstring{\ensuremath{\to}}{->} OAA & Aspartate aminotransferase & PLP & Most tissues &
Direct anaplerosis \\
Propionyl-CoA \texorpdfstring{\ensuremath{\to}}{->} Succinyl-CoA & Propionyl-CoA carboxylase +
methylmalonyl-CoA mutase & Biotin, B\(_{12}\) & Liver & Odd-chain fatty
acid + Val/Ile/Met catabolism \\
\end{longtable}
}

Pyruvate carboxylase is \textbf{allosterically activated by acetyl-CoA}.
This is biologically elegant: when acetyl-CoA accumulates (high fat
oxidation, fed state) and OAA becomes limiting, the cell senses the
imbalance and replenishes OAA so that citrate synthase can continue.

\textbf{Glutamate dehydrogenase (GDH)} catalyses the reversible
deamination of glutamate to \texorpdfstring{\ensuremath{\alpha}}{alpha}-ketoglutarate, releasing NH\(_4^+\). It
uses \emph{either} NAD\(^+\) (catabolic direction) or NADP\(^+\)
(anabolic direction) and is allosterically regulated by ADP (activator)
and GTP (inhibitor) --- making it a sensor of mitochondrial energy
state. In hepatocytes, GDH-driven glutamate catabolism feeds \texorpdfstring{\ensuremath{\alpha}}{alpha}-KG into
the TCA cycle and NH\(_4^+\) into the urea cycle, integrating amino-acid
catabolism with both energy production and nitrogen disposal.

\textbf{Major cataplerotic outflows:}

{\def\LTcaptype{none} % do not increment counter
\begin{longtable}[]{@{}
  >{\raggedright\arraybackslash}p{(\linewidth - 4\tabcolsep) * \real{0.3929}}
  >{\raggedright\arraybackslash}p{(\linewidth - 4\tabcolsep) * \real{0.3393}}
  >{\raggedright\arraybackslash}p{(\linewidth - 4\tabcolsep) * \real{0.2679}}@{}}
\toprule\noalign{}
\begin{minipage}[b]{\linewidth}\raggedright
Intermediate withdrawn
\end{minipage} & \begin{minipage}[b]{\linewidth}\raggedright
Destination pathway
\end{minipage} & \begin{minipage}[b]{\linewidth}\raggedright
Example product
\end{minipage} \\
\midrule\noalign{}
\endhead
\bottomrule\noalign{}
\endlastfoot
\texorpdfstring{\ensuremath{\alpha}}{alpha}-Ketoglutarate & Glutamate / glutamine biosynthesis; collagen
hydroxylation & Glu, Gln, hydroxyproline \\
Succinyl-CoA & Heme biosynthesis (\texorpdfstring{\ensuremath{\delta}}{delta}-ALA synthase) & Heme, cytochromes \\
Oxaloacetate & Gluconeogenesis (via PEPCK); aspartate, asparagine
biosynthesis & Glucose, Asp, Asn \\
Citrate & Cytosolic acetyl-CoA for fatty acid synthesis; cholesterol &
Palmitate, cholesterol \\
\end{longtable}
}

\textbf{Clinical connection --- biotin deficiency:} Because pyruvate
carboxylase, propionyl-CoA carboxylase, and acetyl-CoA carboxylase most
depend on biotin, dietary biotin deficiency (rare; sometimes seen with
chronic raw-egg consumption due to avidin) causes lactic acidosis,
neurological symptoms, and ketogenic dysregulation simultaneously.

\begin{quote}
\textbf{Concept Check 4:} If fluoroacetate (a metabolic poison) inhibits
aconitase (step 2), predict the immediate effect on (a) citrate levels,
(b) isocitrate levels, (c) NADH production, and (d) ATP synthesis.
\end{quote}

\begin{center}\rule{0.5\linewidth}{0.5pt}\end{center}

\subsection{Oxidative Phosphorylation}\label{oxidative-phosphorylation}

Oxidative phosphorylation (OxPhos) accounts for about 90\% of
ATP production in aerobic organisms. It takes place at the inner
mitochondrial membrane.

\subsubsection{Electron Transport Chain
(ETC)}\label{electron-transport-chain-etc}

\begin{figure}[htbp]
\centering
\includegraphics[width=0.82\linewidth,height=4.2in,keepaspectratio]{../figures/mermaid_inline/inline_mermaid_0053_648c8b0187e7.png}
\caption{Sequence diagram for Electron Transport Chain (ETC) showing ordered interaction among NADH (matrix), Complex I (NADH:UQ oxidoreductase) 46 subunits, 1 MDa Pumps 4 H+, Ubiquinone (Q) Lipid-soluble mobile carrier, and Complex II (Succinate DH) 4 subunits No H+ pumping.}
\end{figure}

\emph{Electron flow through the ETC complexes I--IV and proton pumping
that generates the proton motive force (Mermaid).} ATP synthase (Complex
V) uses this PMF to synthesise ATP via rotary catalysis.

\textbf{Complex I (NADH:ubiquinone oxidoreductase):} - Largest ETC
complex (about 46 subunits in mammals, about 1 MDa) -
L-shaped structure: hydrophilic arm in matrix (electron transfer),
membrane arm (proton pumping) - Electron path: NADH \texorpdfstring{\ensuremath{\to}}{->} FMN \texorpdfstring{\ensuremath{\to}}{->} 7 Fe-S
clusters \texorpdfstring{\ensuremath{\to}}{->} ubiquinone (Q) - Pumps \textbf{4 H\(^+\)} per NADH oxidised -
Inhibited by \textbf{rotenone} (fish poison, pesticide) and
\textbf{piericidin A}

\textbf{Complex II (succinate:ubiquinone oxidoreductase):} - Also is
step 6 of the TCA cycle (succinate dehydrogenase) - 4 subunits; FAD
covalently bound - Electron path: succinate \texorpdfstring{\ensuremath{\to}}{->} FAD \texorpdfstring{\ensuremath{\to}}{->} Fe-S clusters \texorpdfstring{\ensuremath{\to}}{->}
ubiquinone (Q) - Pumps \textbf{0 H\(^+\)} (insufficient energy from
FADH\(_2\) oxidation to pump protons) - This is why FADH\(_2\) yields
less ATP than NADH

\textbf{Complex III (cytochrome bc\(_1\) complex):} - The \textbf{Q
cycle} mechanism (Mitchell, Nobel Prize 1978) \citep{mitchell1961}: 1.
QH\(_2\) binds Q\(_o\) site; one electron goes to Fe-S cluster \texorpdfstring{\ensuremath{\to}}{->}
cytochrome c\(_1\) \texorpdfstring{\ensuremath{\to}}{->} cytochrome c (high potential path) 2. Second
electron goes to cytochrome b\(_L\) \texorpdfstring{\ensuremath{\to}}{->} cytochrome b\(_H\) \texorpdfstring{\ensuremath{\to}}{->} Q at Q\(_i\)
site, forming Q\(^{\bullet-}\) (semiquinone) 3. Second QH\(_2\) repeats;
second electron to Q\(_i\) completes QH\(_2\) formation 4. Net: 2
QH\(_2\) oxidised, 1 Q reduced, 2 cytochrome c reduced, \textbf{4
H\(^+\) pumped per 2e\(^-\)} - Inhibited by \textbf{antimycin A}
(Q\(_i\) site) and \textbf{myxothiazol} (Q\(_o\) site)

\textbf{Complex IV (cytochrome c oxidase):} - Contains Cu\(_A\), heme a,
heme a\(_3\)-Cu\(_B\) centres - Transfers 4 electrons to O\(_2\):
\(\text{O}_2 + 4\text{e}^- + 4\text{H}^+_{\text{matrix}} \rightarrow 2\text{H}_2\text{O}\)
- Pumps \textbf{2 H\(^+\) per 2 electrons} (plus 2 H\(^+\) consumed in
water formation) - Inhibited by \textbf{cyanide} (CN\(^-\)),
\textbf{carbon monoxide} (CO), and \textbf{azide} (N\(_3^-\)) --- most
bind heme a\(_3\)-Cu\(_B\)

\begin{quote}
\textbf{Clinical Connection: Cyanide Poisoning} Cyanide irreversibly
inhibits Complex IV by binding to the Fe\(^{3+}\) of heme a\(_3\). Most
electron flow stops, the proton gradient collapses, and ATP synthesis
ceases. Cells switch to \hyperref[gl:anaerobic]{\textbf{anaerobic}}
glycolysis, producing lactic acid. Death results from inability to
produce sufficient ATP, especially in the brain and heart. Treatment:
hydroxocobalamin (vitamin B\(_{12a}\)) chelates cyanide, or sodium
thiosulphate converts it to less toxic thiocyanate.
\end{quote}

\subsubsection{Reactive Oxygen Species and Antioxidant
Defence}\label{reactive-oxygen-species-and-antioxidant-defence}

The ETC's job is to deliver four electrons safely onto O\(_2\) to
produce water --- but about 0.1--2\% of electrons ``leak''
prematurely, producing \textbf{reactive oxygen species (ROS)}.

\textbf{Sites of ROS production in the ETC:}

{\def\LTcaptype{none} % do not increment counter
\begin{longtable}[]{@{}
  >{\raggedright\arraybackslash}p{(\linewidth - 6\tabcolsep) * \real{0.0741}}
  >{\raggedright\arraybackslash}p{(\linewidth - 6\tabcolsep) * \real{0.2037}}
  >{\raggedright\arraybackslash}p{(\linewidth - 6\tabcolsep) * \real{0.1667}}
  >{\raggedright\arraybackslash}p{(\linewidth - 6\tabcolsep) * \real{0.5556}}@{}}
\toprule\noalign{}
\begin{minipage}[b]{\linewidth}\raggedright
Site
\end{minipage} & \begin{minipage}[b]{\linewidth}\raggedright
Primary ROS
\end{minipage} & \begin{minipage}[b]{\linewidth}\raggedright
Mechanism
\end{minipage} & \begin{minipage}[b]{\linewidth}\raggedright
Inhibitor that \emph{increases} ROS
\end{minipage} \\
\midrule\noalign{}
\endhead
\bottomrule\noalign{}
\endlastfoot
Complex I (FMN, Fe-S clusters) & Superoxide (O\(_2^{\bullet-}\)) into
matrix & Reverse electron transfer when ETC is reduced; FMNH\(_2\)
donates 1e\(^-\) to O\(_2\) & Rotenone \\
Complex III (Q\(_o\) site) & Superoxide into both matrix and IMS &
Q\(^{\bullet-}\) semiquinone donates 1e\(^-\) to O\(_2\) & Antimycin
A \\
Glycerol-3-P dehydrogenase, ETF-QO & Superoxide & Flavoprotein leakage &
-- \\
Monoamine oxidase (outer membrane) & H\(_2\)O\(_2\) & 2e\(^-\) donation
to O\(_2\) & -- \\
\end{longtable}
}

\textbf{The ROS cascade:}

\begin{equation}
\text{O}_2 \xrightarrow{1\text{e}^-} \text{O}_2^{\bullet-} \xrightarrow{\text{SOD}} \text{H}_2\text{O}_2 \xrightarrow{\text{Fe}^{2+}} \text{HO}^\bullet + \text{OH}^- + \text{Fe}^{3+}
\label{eq:unit_III_ros_cascade}
\end{equation}

The Fenton reaction (last step) is particularly dangerous: hydroxyl
radicals (HO\(^\bullet\)) are the most reactive biological oxidants and
damage essentially any nearby molecule (DNA, lipids, proteins).

\textbf{Antioxidant defence enzymes:}

{\def\LTcaptype{none} % do not increment counter
\begin{longtable}[]{@{}
  >{\raggedright\arraybackslash}p{(\linewidth - 6\tabcolsep) * \real{0.1765}}
  >{\raggedright\arraybackslash}p{(\linewidth - 6\tabcolsep) * \real{0.2353}}
  >{\raggedright\arraybackslash}p{(\linewidth - 6\tabcolsep) * \real{0.2353}}
  >{\raggedright\arraybackslash}p{(\linewidth - 6\tabcolsep) * \real{0.3529}}@{}}
\toprule\noalign{}
\begin{minipage}[b]{\linewidth}\raggedright
Enzyme
\end{minipage} & \begin{minipage}[b]{\linewidth}\raggedright
Reaction
\end{minipage} & \begin{minipage}[b]{\linewidth}\raggedright
Cofactor
\end{minipage} & \begin{minipage}[b]{\linewidth}\raggedright
Localisation
\end{minipage} \\
\midrule\noalign{}
\endhead
\bottomrule\noalign{}
\endlastfoot
Superoxide dismutase (SOD1, cytosolic) & 2 O\(_2^{\bullet-}\) + 2
H\(^+\) \texorpdfstring{\ensuremath{\to}}{->} H\(_2\)O\(_2\) + O\(_2\) & Cu, Zn & Cytoplasm, IMS \\
Superoxide dismutase (SOD2, mitochondrial) & 2 O\(_2^{\bullet-}\) + 2
H\(^+\) \texorpdfstring{\ensuremath{\to}}{->} H\(_2\)O\(_2\) + O\(_2\) & Mn & Mitochondrial matrix \\
Superoxide dismutase (SOD3, extracellular) & 2 O\(_2^{\bullet-}\) + 2
H\(^+\) \texorpdfstring{\ensuremath{\to}}{->} H\(_2\)O\(_2\) + O\(_2\) & Cu, Zn & Plasma, ECM \\
Catalase & 2 H\(_2\)O\(_2\) \texorpdfstring{\ensuremath{\to}}{->} 2 H\(_2\)O + O\(_2\) & Heme &
Peroxisomes \\
Glutathione peroxidase (GPx) & H\(_2\)O\(_2\) + 2 GSH \texorpdfstring{\ensuremath{\to}}{->} 2 H\(_2\)O +
GSSG & Selenocysteine & Cytoplasm, mitochondria \\
Peroxiredoxins (Prx1--6) & H\(_2\)O\(_2\) + 2 R-SH \texorpdfstring{\ensuremath{\to}}{->} 2 H\(_2\)O +
R-S-S-R & Cysteine & Multiple compartments \\
\end{longtable}
}

\textbf{Non-enzymatic antioxidants:} Glutathione (GSH; about 5
mM cytosolic), ascorbate (vitamin C), \texorpdfstring{\ensuremath{\alpha}}{alpha}-tocopherol (vitamin E,
lipid-soluble), ubiquinone (also acts as antioxidant in its reduced
QH\(_2\) form), and uric acid (about 50\% of plasma antioxidant
capacity in humans).

\textbf{Clinical relevance:}

\begin{itemize}
\tightlist
\item
  \textbf{Mitochondrial diseases:} Complex I and III deficiencies often
  present with elevated lactate AND elevated ROS damage markers
  (8-oxoguanine in mtDNA, malondialdehyde in lipids).
\item
  \textbf{Aging (mitochondrial free-radical theory):} Cumulative mtDNA
  damage from ROS may underlie age-related mitochondrial decline. Modern
  data suggests ROS are also signalling molecules, complicating the
  simple ``oxidative damage = aging'' model.
\item
  \textbf{Reperfusion injury:} During ischaemia, ETC carriers become
  highly reduced. Upon reoxygenation, a burst of ROS is released ---
  paradoxically, the \emph{return} of oxygen causes the worst tissue
  damage in stroke and heart attack.
\item
  \textbf{ALS:} Mutations in SOD1 cause about 20\% of familial
  amyotrophic lateral sclerosis. The pathology is gain-of-function
  (toxic SOD1 aggregates), not simple loss of antioxidant activity.
\end{itemize}

\begin{quote}
\textbf{Concept Check 4b:} Antimycin A blocks the Q\(_i\) site of
Complex III. Predict the effects on (a) electron flow upstream
(Complexes I, II), (b) downstream (Complex IV), (c) PMF magnitude, and
(d) ROS production. Explain why this combination of effects makes
antimycin A particularly cytotoxic.
\end{quote}

\subsubsection{ATP Synthase --- the Rotary
Nanomotor}\label{atp-synthase-the-rotary-nanomotor}

The \textbf{proton-motive force} (PMF) consists of two components ---
the electrical potential (\(\Delta\Psi\)) and the chemical gradient
(\texorpdfstring{\ensuremath{\Delta}}{Delta}pH):

\begin{equation}
\text{PMF} = \Delta\Psi - \frac{2.303RT}{F}\Delta\text{pH} \approx 180 \; \text{mV}
\label{eq:unit_III_proton_motive_force}
\end{equation}

In mitochondria, \(\Delta\Psi \approx 150\) mV contributes
about 80\% of the PMF, while \texorpdfstring{\ensuremath{\Delta}}{Delta}pH about 0.5 units
contributes about 20\%.

The free energy stored per mole of protons translocated against this PMF
is:

\[\Delta G_{\text{H}^+} = F \cdot \text{PMF} \approx 96{,}485 \times 0.180 \approx 17.4 \; \text{kJ/mol} \tag{9.7d}\]

So pumping about 3.67 H\(^+\) to make one ATP requires
about 64 kJ/mol of stored gradient energy --- the budget that
funds an ATP whose cellular hydrolysis \(\Delta G \approx -54\) kJ/mol.
The about 10 kJ/mol surplus is the thermodynamic driving force
that keeps net ATP synthesis flowing forward and sets the response time
of the system to ATP demand.

\textbf{ATP synthase (Complex V)} is a molecular rotary motor:

\textbf{F\(_1\) (catalytic head, matrix-facing):} -
alpha\(_3\)-beta\(_3\)-gamma-delta-epsilon - Three catalytic sites on
beta subunits (binding change mechanism, Boyer) - Each beta subunit
cycles through three states: \textbf{O} (open, empty), \textbf{L}
(loose, binds ADP + P\(_i\)), \textbf{T} (tight, catalyses ATP
synthesis) - 120-degree rotation of gamma converts O\texorpdfstring{\ensuremath{\to}}{->}L\texorpdfstring{\ensuremath{\to}}{->}T\texorpdfstring{\ensuremath{\to}}{->}O

\textbf{F\(_0\) (membrane-spanning proton channel):} - Subunit a (proton
channel) + c-ring (c\(_{10-15}\), varies by species) - Each c-subunit
binds one H\(^+\) from IMS, rotates, releases H\(^+\) to matrix - Full
c-ring rotation (360 degrees) produces 3 ATP

\textbf{Stoichiometry:} If the c-ring has 10 subunits, then 10 H\(^+\)
per revolution, producing 3 ATP, so about 3.3 H\(^+\) per ATP.
With 8 c-subunits (yeast), about 2.7 H\(^+\) per ATP.

Paul Boyer (binding change mechanism) and John Walker (crystal structure
of F\(_1\)) shared the Nobel Prize in Chemistry, 1997. The rotation of
the gamma subunit was directly visualised by Yoshida and colleagues
using fluorescent \hyperref[gl:actin]{\textbf{actin}} filaments attached
to gamma --- one of the most beautiful experiments in biochemistry.

\subsubsection{How the ``30--32 ATP'' Number Is Calculated: P/O
Ratios}\label{how-the-3032-atp-number-is-calculated-po-ratios}

The classical undergraduate value of 38 ATP/glucose came from
\emph{assumed} round-number stoichiometries (3 ATP/NADH, 2
ATP/FADH\(_2\), 1 ATP per substrate-level step). Modern biochemistry
uses \textbf{measured P/O ratios} based on actual c-ring stoichiometry
and proton transport:

\textbf{Proton pumping per substrate:}

{\def\LTcaptype{none} % do not increment counter
\begin{longtable}[]{@{}
  >{\raggedright\arraybackslash}p{(\linewidth - 8\tabcolsep) * \real{0.1579}}
  >{\raggedright\arraybackslash}p{(\linewidth - 8\tabcolsep) * \real{0.1579}}
  >{\raggedright\arraybackslash}p{(\linewidth - 8\tabcolsep) * \real{0.1930}}
  >{\raggedright\arraybackslash}p{(\linewidth - 8\tabcolsep) * \real{0.1754}}
  >{\raggedright\arraybackslash}p{(\linewidth - 8\tabcolsep) * \real{0.3158}}@{}}
\toprule\noalign{}
\begin{minipage}[b]{\linewidth}\raggedright
Substrate
\end{minipage} & \begin{minipage}[b]{\linewidth}\raggedright
Complex I
\end{minipage} & \begin{minipage}[b]{\linewidth}\raggedright
Complex III
\end{minipage} & \begin{minipage}[b]{\linewidth}\raggedright
Complex IV
\end{minipage} & \begin{minipage}[b]{\linewidth}\raggedright
Total H\(^+\) pumped
\end{minipage} \\
\midrule\noalign{}
\endhead
\bottomrule\noalign{}
\endlastfoot
1 NADH (matrix) & 4 & 4 & 2 & \textbf{10 H\(^+\)} \\
1 FADH\(_2\) (Complex II) & 0 & 4 & 2 & \textbf{6 H\(^+\)} \\
\end{longtable}
}

\textbf{Mammalian mitochondrial c-ring:} 8 subunits (recent cryo-EM
data). \textbf{ATP synthase stoichiometry:} 8 H\(^+\) per full
revolution \texorpdfstring{\ensuremath{\to}}{->} 3 ATP per revolution \texorpdfstring{\ensuremath{\to}}{->} \textbf{2.67 H\(^+\)/ATP}.
\textbf{Add P\(_i\) transport:} importing P\(_i\) via the phosphate
carrier costs \textbf{1 additional H\(^+\)} per ATP. So the
\emph{effective} H\(^+\) cost per ATP is approximately \textbf{3.67
H\(^+\)/ATP}.

\begin{equation}
\text{P/O ratio (NADH)} = \frac{10\;\text{H}^+}{3.67\;\text{H}^+/\text{ATP}} \approx 2.5\;\text{ATP per NADH}
\label{eq:unit_III_po_ratio_nadh}
\end{equation}

\begin{equation}
\text{P/O ratio (FADH}_2\text{)} = \frac{6\;\text{H}^+}{3.67\;\text{H}^+/\text{ATP}} \approx 1.5\;\text{ATP per FADH}_2
\label{eq:unit_III_po_ratio_fadh2}
\end{equation}

\textbf{Total ATP per glucose (aerobic):}

\begin{equation}
\text{ATP}_{\text{glucose}} = \underbrace{2}_{\text{glycolysis SLP}} + \underbrace{2}_{\text{TCA SLP (2 GTP)}} + \underbrace{10 \times 2.5}_{\text{NADH via OxPhos}} + \underbrace{2 \times 1.5}_{\text{FADH}_2\text{ via OxPhos}} \approx 32\;\text{ATP}
\label{eq:unit_III_atp_yield}
\end{equation}

(if cytoplasmic NADH uses the malate-aspartate shuttle;
about 30 ATP if the glycerol-3-P shuttle is used in muscle and
brain).

{\def\LTcaptype{none} % do not increment counter
\begin{longtable}[]{@{}llllll@{}}
\toprule\noalign{}
Stage & NADH & FADH\(_2\) & ATP direct & ATP via OxPhos & Subtotal \\
\midrule\noalign{}
\endhead
\bottomrule\noalign{}
\endlastfoot
Glycolysis & 2 & 0 & 2 & 2 \texorpdfstring{\ensuremath{\times}}{x} 2.5 = 5* & 7 \\
Pyruvate oxidation & 2 & 0 & 0 & 2 \texorpdfstring{\ensuremath{\times}}{x} 2.5 = 5 & 5 \\
TCA cycle & 6 & 2 & 2 & 6 \texorpdfstring{\ensuremath{\times}}{x} 2.5 + 2 \texorpdfstring{\ensuremath{\times}}{x} 1.5 = 18 & 20 \\
\textbf{Total} & \textbf{10} & \textbf{2} & \textbf{4} & \textbf{28} &
\textbf{about 30--32} \\
\end{longtable}
}

*Glycolytic NADH yields 2.5 ATP (malate-aspartate shuttle) or 1.5 ATP
(glycerol-3-phosphate shuttle), giving 30 or 32 ATP total.

\textbf{The glycerol-3-phosphate shuttle vs.~malate-aspartate shuttle:}

\begin{itemize}
\tightlist
\item
  Cytoplasmic NADH cannot cross the IMM directly.
\item
  \textbf{Malate-aspartate shuttle} (liver, heart, kidney): NADH
  electrons enter the matrix via OAA\texorpdfstring{\ensuremath{\to}}{->}malate\texorpdfstring{\ensuremath{\to}}{->}OAA at full reducing
  potential \texorpdfstring{\ensuremath{\to}}{->} \textbf{2.5 ATP each}.
\item
  \textbf{Glycerol-3-phosphate shuttle} (skeletal muscle, brain): NADH
  electrons enter directly into Q via FAD-linked mitochondrial G3PDH \texorpdfstring{\ensuremath{\to}}{->}
  \textbf{1.5 ATP each} (skips Complex I).
\end{itemize}

\subsubsection{Worked Example: Total ATP
Yield}\label{worked-example-total-atp-yield}

\emph{Problem:} Calculate the total ATP yield from one molecule of
glucose using the malate-aspartate shuttle for most glycolytic NADH.

\emph{Solution:}

{\def\LTcaptype{none} % do not increment counter
\begin{longtable}[]{@{}llll@{}}
\toprule\noalign{}
Source & Quantity & ATP per unit & ATP \\
\midrule\noalign{}
\endhead
\bottomrule\noalign{}
\endlastfoot
Glycolysis substrate-level & 2 ATP & 1 & 2 \\
TCA substrate-level (GTP) & 2 GTP & 1 & 2 \\
Glycolytic NADH (malate-aspartate) & 2 NADH & 2.5 & 5 \\
Pyruvate DH NADH & 2 NADH & 2.5 & 5 \\
TCA NADH & 6 NADH & 2.5 & 15 \\
TCA FADH\(_2\) & 2 FADH\(_2\) & 1.5 & 3 \\
\textbf{Total} & & & \textbf{32} \\
\end{longtable}
}

The thermodynamic efficiency:

\[\eta = \frac{32 \times 30.5}{2{,}870} \times 100\% = 34\% \tag{9.10}\]

Using the more realistic \emph{cellular} \(\Delta G\) of \(-50\) kJ/mol,
efficiency rises to about 56\%, comparable to the best internal
combustion engines and considerably better than most human-engineered
energy converters.

\begin{quote}
\textbf{Concept Check 5:} Why does FADH\(_2\) yield fewer ATP than NADH?
Trace the difference back to the ETC entry point and proton pumping
stoichiometry.
\end{quote}

\begin{quote}
\textbf{Concept Check 5b:} Consult \cref{eq:unit_III_atp_yield}. If a
mutation altered the c-ring stoichiometry from 8 to 12 subunits (raising
H\(^+\)/ATP from 2.67 to 4), recompute the P/O ratios for NADH and
FADH\(_2\) and the total ATP yield per glucose. Comment on what this
would mean for daily ATP turnover and whole-organism metabolism.
\end{quote}

\begin{center}\rule{0.5\linewidth}{0.5pt}\end{center}

\subsection{Substrate-Level vs.~Oxidative Phosphorylation: A
Quantitative
Comparison}\label{substrate-level-vs.-oxidative-phosphorylation-a-quantitative-comparison}

{\def\LTcaptype{none} % do not increment counter
\begin{longtable}[]{@{}
  >{\raggedright\arraybackslash}p{(\linewidth - 4\tabcolsep) * \real{0.1111}}
  >{\raggedright\arraybackslash}p{(\linewidth - 4\tabcolsep) * \real{0.4921}}
  >{\raggedright\arraybackslash}p{(\linewidth - 4\tabcolsep) * \real{0.3968}}@{}}
\toprule\noalign{}
\begin{minipage}[b]{\linewidth}\raggedright
Feature
\end{minipage} & \begin{minipage}[b]{\linewidth}\raggedright
Substrate-level phosphorylation
\end{minipage} & \begin{minipage}[b]{\linewidth}\raggedright
Oxidative phosphorylation
\end{minipage} \\
\midrule\noalign{}
\endhead
\bottomrule\noalign{}
\endlastfoot
Location & Cytoplasm (glycolysis), matrix (TCA) & Inner mitochondrial
membrane \\
Energy source & High-energy phosphorylated intermediate & Proton motive
force \\
O\(_2\) required? & No & Yes (terminal electron acceptor) \\
Energy ``currency'' intermediate & 1,3-BPG, PEP, succinyl-CoA &
\(\Delta\mu_{\text{H}^+}\) (electrochemical) \\
Maximum theoretical efficiency & about 60\% (per intermediate)
& about 70\% (PMF \texorpdfstring{\ensuremath{\to}}{->} ATP) \\
ATP per glucose (aerobic, 32 total) & 4 ATP (12.5\%) & 28 ATP
(87.5\%) \\
Speed of onset & Immediate (ms) & Seconds (requires PMF buildup) \\
Sustainability & Limited (substrate-bound) & Continuous (as long as fuel
+ O\(_2\)) \\
Examples & PGK, pyruvate kinase, succinyl-CoA synthetase & ATP
synthase \\
Inhibited by & Specific enzyme inhibitors & ETC inhibitors, uncouplers,
oligomycin \\
\end{longtable}
}

\textbf{Quantitative leverage:} Although substrate-level phosphorylation
contributes roughly 12.5\% of total ATP in this accounting, it is
critical in two scenarios:

\begin{enumerate}
\def\labelenumi{\arabic{enumi}.}
\tightlist
\item
  \textbf{Anaerobic conditions:} It is the \emph{primarily} source of
  ATP, so 4 ATP/glucose is everything.
\item
  \textbf{Sprint exercise:} It can produce ATP in milliseconds, before
  mitochondrial respiration ramps up; the phosphocreatine system extends
  this further.
\end{enumerate}

This illustrates why metabolic systems are layered: fast-but-low-yield
substrate-level phosphorylation handles transients while
slow-but-high-yield OxPhos handles steady state.

\textbf{Numerical leverage at the organism level.} A 70 kg adult at rest
uses about 1,400 kcal/day, equivalent to about 250 mol
ATP turned over. If primarily substrate-level phosphorylation operated,
with about 4 ATP/glucose, the body would need to oxidise
about 62 mol glucose/day (= 11 kg). Aerobic OxPhos at
about 32 ATP/glucose reduces the requirement to
about 8 mol = 1.4 kg glucose/day --- a more achievable dietary
intake. The 8-fold yield improvement of OxPhos is precisely what enables
active multicellular life on a planet with bulk-food-supply constraints.

\begin{center}\rule{0.5\linewidth}{0.5pt}\end{center}

\subsection{Fermentation}\label{fermentation}

Without O\(_2\) (or in cells that lack mitochondria), pyruvate is not
oxidised aerobically. Instead, \textbf{fermentation} regenerates
NAD\(^+\) so glycolysis can continue:

\textbf{Alcoholic fermentation} (yeast, some bacteria):
\[\text{Pyruvate} \xrightarrow[\text{TPP}]{\text{pyruvate decarboxylase}} \text{Acetaldehyde} + \text{CO}_2 \xrightarrow[\text{NADH} \rightarrow \text{NAD}^+]{\text{alcohol dehydrogenase}} \text{Ethanol} \tag{9.11}\]

\textbf{Lactic acid fermentation} (muscle, erythrocytes,
\emph{Lactobacillus}):
\[\text{Pyruvate} + \text{NADH} \xrightarrow{\text{lactate dehydrogenase}} \text{Lactate} + \text{NAD}^+ \tag{9.12}\]

Both pathways yield primarily \textbf{2 ATP per glucose} (from
glycolysis alone).

\textbf{The Cori Cycle:} In exercising muscle, lactate is exported to
the blood, taken up by the liver, and converted back to glucose via
\textbf{gluconeogenesis} at the cost of 6 ATP. This glucose returns to
the muscle via the blood. The Cori cycle shifts the metabolic burden of
gluconeogenesis from the oxygen-limited muscle to the well-oxygenated
liver.

\textbf{Pasteur Effect:} In the presence of oxygen, yeast switches from
fermentation to aerobic respiration, dramatically reducing glucose
consumption (because 32 ATP/glucose aerobically vs.~2 ATP/glucose by
fermentation --- 16x more efficient).

\begin{quote}
\textbf{Clinical Connection: Lactic Acidosis} Lactic acidosis occurs
when lactate production exceeds clearance. \textbf{Type A} (most
common): tissue hypoxia (shock, cardiac arrest, severe anaemia) forces
reliance on glycolysis. \textbf{Type B}: metabolic causes (liver failure
preventing lactate clearance, metformin in renal impairment, thiamine
deficiency preventing PDC function, mitochondrial diseases). Blood
lactate \textgreater4 mmol/L with pH \textless7.35 indicates lactic
acidosis. Treatment focuses on addressing the underlying cause and
restoring tissue oxygenation.
\end{quote}

\begin{quote}
\textbf{Clinical Connection: Mitochondrial Complex Disorders}
\hyperref[gl:mutation]{\textbf{Mutation}}s in nuclear or mitochondrial
\hyperref[gl:gene]{\textbf{gene}}s encoding ETC complex subunits cause a
spectrum of diseases: - \textbf{Complex I deficiency:} Most common ETC
disorder (about 30\% of paediatric mitochondrial disease);
Leigh syndrome (necrotising encephalopathy), optic atrophy,
cardiomyopathy - \textbf{Complex II deficiency:} Rare; associated with
paragangliomas and pheochromocytomas (SDH mutations are tumour
suppressors) - \textbf{Complex III deficiency:} Exercise intolerance,
lactic acidosis - \textbf{Complex IV (COX) deficiency:} Leigh syndrome,
cardiomyopathy; reversible COX deficiency in infants Most are
characterised by lactic acidosis, neurological dysfunction, and exercise
intolerance --- tissues with high energy demands (brain, heart, muscle)
are most affected.
\end{quote}

\begin{quote}
\textbf{Concept Check 6:} A patient with a mitochondrial Complex I
mutation has reduced NADH oxidation. Would you expect succinate-linked
respiration (via Complex II) to be normal or impaired in this patient?
Explain.
\end{quote}

\begin{center}\rule{0.5\linewidth}{0.5pt}\end{center}

\subsection{Uncoupling Proteins and Brown Adipose
Thermogenesis}\label{uncoupling-proteins-and-brown-adipose-thermogenesis}

In most cells, the primary ``exit'' for protons in the IMS back to the
matrix is through ATP synthase. \textbf{Uncoupling proteins (UCPs)}
provide an alternative path that dissipates the proton gradient as heat
--- uncoupling electron transport from ATP synthesis.

\subsubsection{UCP1: The Brown Adipose
Thermogenin}\label{ucp1-the-brown-adipose-thermogenin}

\textbf{Structure:} UCP1 is a 32-kDa, six-transmembrane-helix protein in
the IMM of brown adipocytes. It belongs to the SLC25 mitochondrial
carrier family.

\textbf{Mechanism (proton transport, not pure leak):}

\begin{equation}
\text{H}^+_{\text{IMS}} + \text{FA}^- \xrightarrow{\text{UCP1}} \text{H}^+_{\text{matrix}} + \text{FA}^-
\label{eq:unit_III_ucp1_mechanism}
\end{equation}

Long-chain fatty acid anions (FA\(^-\)) act as cofactors: they bind
H\(^+\) in the IMS, traverse the membrane via UCP1, and release H\(^+\)
in the matrix. UCP1 is \textbf{inhibited by purine nucleotides} (ATP,
ADP, GTP, GDP) and \textbf{activated by fatty acids} released from
intracellular triglyceride stores.

\textbf{Activation cascade in brown adipose tissue (BAT):}

\begin{figure}[htbp]
\centering
\includegraphics[width=0.82\linewidth,height=4.2in,keepaspectratio]{../figures/mermaid_inline/inline_mermaid_0054_6fd41e09a5eb.png}
\caption{Flowchart for UCP1: The Brown Adipose Thermogenin: Norepinephrine (cold, sympathetic activation), beta3-adrenergic receptor, Adenylyl cyclase, and increased cAMP form the diagram's primary path or branches.}
\end{figure}

\emph{Sympathetic activation of brown adipose thermogenesis (Mermaid).}
Norepinephrine binds \texorpdfstring{\ensuremath{\beta}}{beta}\(_3\)-adrenergic receptors \texorpdfstring{\ensuremath{\to}}{->} cAMP \texorpdfstring{\ensuremath{\to}}{->} PKA \texorpdfstring{\ensuremath{\to}}{->}
simultaneous lipolysis (acute fuel) and PGC-1\texorpdfstring{\ensuremath{\alpha}}{alpha} transcription (chronic
mitochondrial biogenesis). UCP1 short-circuits the proton gradient,
releasing energy as heat.

\subsubsection{Physiological Roles}\label{physiological-roles}

{\def\LTcaptype{none} % do not increment counter
\begin{longtable}[]{@{}
  >{\raggedright\arraybackslash}p{(\linewidth - 2\tabcolsep) * \real{0.2258}}
  >{\raggedright\arraybackslash}p{(\linewidth - 2\tabcolsep) * \real{0.7742}}@{}}
\toprule\noalign{}
\begin{minipage}[b]{\linewidth}\raggedright
Setting
\end{minipage} & \begin{minipage}[b]{\linewidth}\raggedright
Role of UCP1 / brown fat
\end{minipage} \\
\midrule\noalign{}
\endhead
\bottomrule\noalign{}
\endlastfoot
\textbf{Neonates} & Critical for non-shivering thermogenesis; lack
effective shivering due to small muscle mass; about 5\% body
weight is BAT \\
\textbf{Hibernating mammals} & Massive BAT depots maintain
about 5 °C body temp; explosive arousal phase warms the body in
hours \\
\textbf{Adult humans} & Active BAT discovered in 2009 in the
supraclavicular and paraspinal regions; activated by cold exposure;
declines with age and obesity \\
\textbf{Cold acclimation} & Chronic cold exposure increases BAT mass and
UCP1 expression --- a target for obesity therapy \\
\textbf{Beige (brite) adipocytes} & UCP1\(^+\) cells that emerge from
white adipose under cold/exercise; an inducible thermogenic
compartment \\
\end{longtable}
}

\textbf{Other UCPs:}

\begin{itemize}
\tightlist
\item
  \textbf{UCP2}: ubiquitously expressed; modest uncoupling activity;
  thought to attenuate ROS and modulate insulin secretion.
\item
  \textbf{UCP3}: skeletal muscle and BAT; possible role in lipid
  handling; debated thermogenic relevance.
\item
  \textbf{Plant UCPs} (PUMPs): present in mitochondria of
  e.g.~\emph{Arum} species, where they generate the heat
  (about 46 °C) that volatilises odours to attract pollinator
  insects.
\end{itemize}

\begin{quote}
\textbf{Clinical Connection: DNP --- A Cautionary Tale}
2,4-Dinitrophenol (DNP) is a synthetic chemical uncoupler that
dissipates PMF as heat, mimicking UCP1 across most tissues. In the
1930s, DNP was sold as a weight-loss drug --- patients lost
about 1 kg/week as their basal metabolic rate increased
about 50\%. But the therapeutic window is dangerously narrow:
at slightly higher doses, body temperature rises to \textgreater40 °C
(fatal hyperthermia), and unregulated tissue uncoupling causes
cataracts, neuropathy, and sudden death. DNP was banned in 1938 but
remains a recurrent source of bodybuilding deaths today. The takeaway:
physiological uncoupling is \emph{regulated} (UCP1 activated primarily
in BAT, primarily on demand) --- pharmacological uncoupling is
unregulated and potentially lethal.
\end{quote}

\begin{quote}
\textbf{Concept Check 6b:} A 3-month-old infant is brought to the ED in
mild hypothermia after exposure to cold. Imaging reveals normal BAT
volume and PET activity. Predict (a) whether shivering or non-shivering
thermogenesis dominates the infant's response, (b) the role of
\texorpdfstring{\ensuremath{\beta}}{beta}\(_3\)-AR signalling, and (c) the metabolic consequences of repeated
cold exposures.
\end{quote}

\begin{center}\rule{0.5\linewidth}{0.5pt}\end{center}

\subsection{Metabolic Poisons and
Inhibitors}\label{metabolic-poisons-and-inhibitors}

Understanding the ETC and oxidative phosphorylation has been greatly
advanced by studying specific inhibitors and metabolic poisons:

\subsubsection{ETC Inhibitors}\label{etc-inhibitors}

{\def\LTcaptype{none} % do not increment counter
\begin{longtable}[]{@{}
  >{\raggedright\arraybackslash}p{(\linewidth - 4\tabcolsep) * \real{0.4286}}
  >{\raggedright\arraybackslash}p{(\linewidth - 4\tabcolsep) * \real{0.2857}}
  >{\raggedright\arraybackslash}p{(\linewidth - 4\tabcolsep) * \real{0.2857}}@{}}
\toprule\noalign{}
\begin{minipage}[b]{\linewidth}\raggedright
Inhibitor
\end{minipage} & \begin{minipage}[b]{\linewidth}\raggedright
Target
\end{minipage} & \begin{minipage}[b]{\linewidth}\raggedright
Effect
\end{minipage} \\
\midrule\noalign{}
\endhead
\bottomrule\noalign{}
\endlastfoot
\textbf{Rotenone} & Complex I (NADH dehydrogenase) & Blocks electron
transfer from Fe-S clusters to ubiquinone; NADH-linked respiration
stops; used as fish poison and pesticide \\
\textbf{Antimycin A} & Complex III (Q\(_i\) site) & Blocks electron
transfer at the Q\(_i\) site of the Q cycle; generates superoxide \\
\textbf{Cyanide (CN\(^-\))} & Complex IV (CuB/a\(_3\)) & Binds
Fe\(^{3+}\) in cytochrome a\(_3\), blocking O\(_2\) reduction; rapidly
fatal \\
\textbf{Carbon monoxide (CO)} & Complex IV (haem a\(_3\)) & Competes
with O\(_2\) for binding; also binds haemoglobin (200x affinity vs
O\(_2\)) \\
\textbf{Hydrogen sulfide (H\(_2\)S)} & Complex IV & Similar to cyanide;
binds haem iron; toxic at \textgreater50 ppm in air \\
\end{longtable}
}

\subsubsection{Uncouplers}\label{uncouplers}

\textbf{Uncouplers} dissipate the proton gradient without inhibiting
electron transport, allowing respiration to continue at maximum rate
without ATP production:

\begin{itemize}
\tightlist
\item
  \textbf{2,4-dinitrophenol (DNP):} A lipid-soluble weak acid; carries
  H\(^+\) across the IMM. Was used as a weight-loss drug in the 1930s
  (burns fat by dissipating energy as heat) but was banned due to narrow
  therapeutic window and fatal hyperthermia.
\item
  \textbf{FCCP/CCCP:} Research uncouplers; commonly used in
  mitochondrial stress tests (Seahorse assays)
\item
  \textbf{Thermogenin (UCP1):} The physiological uncoupler in brown
  adipose tissue, developed in the thermogenesis discussion above.
\end{itemize}

\subsubsection{ATP Synthase Inhibitors}\label{atp-synthase-inhibitors}

\begin{itemize}
\tightlist
\item
  \textbf{Oligomycin:} Blocks the H\(^+\) channel in the c-ring of
  F\(_0\); prevents H\(^+\) flow and ATP synthesis; also halts electron
  transport (because PMF builds up, preventing further H\(^+\) pumping
  by the ETC)
\item
  \textbf{DCCD (dicyclohexylcarbodiimide):} Reacts with a specific
  Asp/Glu residue in the c-subunit, blocking H\(^+\) translocation
\end{itemize}

\begin{quote}
\textbf{Concept Check 7:} If oligomycin is added to mitochondria,
electron transport stops. If DNP (an uncoupler) is then added, electron
transport resumes. Explain these observations in terms of the
chemiosmotic model.
\end{quote}

\begin{quote}
\textbf{Concept Check 7b:} A researcher exposes isolated mitochondria
sequentially to: (1) succinate, (2) ADP + P\(_i\), (3) rotenone, (4)
DNP. For each addition, predict whether O\(_2\) consumption increases,
decreases, or stays the same, and explain why --- referring to
\cref{eq:unit_III_proton_motive_force} and \cref{eq:unit_III_atp_yield}.
\end{quote}

\begin{center}\rule{0.5\linewidth}{0.5pt}\end{center}

\subsection{Computational Bridge}\label{computational-bridge-8}

Cellular ATP hydrolysis free energy is adjusted by the mass-action ratio
\(Q = [\mathrm{ADP}][\mathrm{P_i}]/[\mathrm{ATP}]\):

\begin{Shaded}
\begin{Highlighting}[]
\ImportTok{from}\NormalTok{ biology.biochemistry }\ImportTok{import}\NormalTok{ atp\_free\_energy}

\BuiltInTok{print}\NormalTok{(}\BuiltInTok{round}\NormalTok{(atp\_free\_energy(atp\_conc\_mM}\OperatorTok{=}\FloatTok{3.0}\NormalTok{, adp\_conc\_mM}\OperatorTok{=}\FloatTok{1.0}\NormalTok{, pi\_conc\_mM}\OperatorTok{=}\FloatTok{10.0}\NormalTok{), }\DecValTok{2}\NormalTok{))}
\end{Highlighting}
\end{Shaded}

\begin{quote}
\textbf{Clinical / systems note:} Mitochondrial poisons (cyanide, carbon
monoxide) and uncouplers like DNP illustrate how collapsing PMF or
bypassing ATP synthase converts energy to heat --- the same gradient you
model when counting protons per ATP. UCP1 is a regulated example of the
same principle.
\end{quote}

\begin{center}\rule{0.5\linewidth}{0.5pt}\end{center}

\subsection{Current Evidence and Frontier
Biology}\label{current-evidence-and-frontier-biology-11}

For \textbf{Bioenergetics and Cellular Respiration}, frontier biology
belongs inside the evidence logic of the chapter. Metabolism is now
studied as a regulated network constrained by energy, redox balance,
compartmentation, and environment. The core reading question is this:
respiration claims should track electrons, protons, redox poise, ATP
yield, and uncoupling.

\begin{itemize}
\tightlist
\item
  \textbf{What to verify:} identify the observation, model, assay, or
  dataset that would make the claim stronger or weaker.
\item
  \textbf{What to qualify:} state the scale, organism, cell type,
  environmental condition, or population where the claim is expected to
  hold.
\item
  \textbf{What to compare:} test at least one alternative explanation,
  baseline, or null model before treating the pattern as causal.
\item
  \textbf{What to cite:} distinguish primary evidence, review synthesis,
  public dataset, and institutional guidance; for recent or numeric
  claims, prefer the source closest to the measurement and state what
  has changed since it was published.
\end{itemize}

A strong metabolic explanation names the flux, the limiting step, the
sensor, and the condition under which the pathway changes direction or
priority.

\textbf{Source practice:} For metabolic claims, keep the organism,
compartment, energetic state, and measurement method visible; a pathway
map is not enough without flux or concentration evidence.

\subsubsection{Current Evidence Map: Respiration Evidence
Accounting}\label{current-evidence-map-respiration-evidence-accounting}

\begin{figure}[htbp]
\centering
\includegraphics[width=0.82\linewidth,height=4.2in,keepaspectratio]{../figures/mermaid_inline/inline_mermaid_0055_05d0878efc1f.png}
\caption{ATP-yield claims are conditional on shuttle use, proton leak, coupling efficiency, tissue state, and measurement method rather than one fixed number.}
\end{figure}

\subsection{Summary}\label{summary-11}

\begin{itemize}
\tightlist
\item
  \(\Delta G = \Delta H - T\Delta S\)
  (\cref{eq:unit_III_gibbs_free_energy}); cellular reactions are coupled
  to ATP hydrolysis (\(\Delta G \approx -54\) kJ/mol) to drive
  endergonic processes. Cells maintain far-from-equilibrium conditions.
\item
  ATP turnover is about 40 kg/day; the high \(\Delta G\) of
  hydrolysis arises from charge repulsion, resonance stabilisation, and
  hydration of products. ATP sits in the \emph{middle} of the
  phosphoryl-transfer hierarchy, between very high (PEP, 1,3-BPG) and
  very low (glucose-6-P).
\item
  Glucose oxidation has a standard free energy of \(\approx -2{,}870\)
  kJ/mol; cells capture about 56\% of this as ATP, dissipating
  the rest as heat.
\item
  Glycolysis: 10 cytoplasmic steps; net 2 ATP + 2 NADH per glucose;
  regulated by PFK-1 (activated by AMP, F2,6BP; inhibited by ATP,
  citrate).
\item
  Gluconeogenesis: not the reverse of glycolysis --- uses four bypass
  enzymes (pyruvate carboxylase, PEPCK, FBPase-1, G6Pase) at a cost of 6
  ATP equivalents per glucose; reciprocally regulated with glycolysis
  via F2,6BP.
\item
  The Cori cycle and glucose-alanine cycle distribute metabolic burden
  between muscle and liver, supporting hypoxic tissues at a
  4-ATP-per-glucose net cost to the organism.
\item
  PDC links glycolysis to TCA; regulated by
  phosphorylation/dephosphorylation; deficiency causes lactic acidosis.
\item
  TCA cycle: 8 mitochondrial matrix steps; 6 NADH + 2 FADH\(_2\) + 2 GTP
  per glucose; regulated by isocitrate DH and \texorpdfstring{\ensuremath{\alpha}}{alpha}-KG DH; anaplerotic
  (pyruvate carboxylase, glutamate dehydrogenase) and cataplerotic flux
  balance the cycle.
\item
  ETC: Complex I (4 H\(^+\)), II (0 H\(^+\)), III (4 H\(^+\) via Q
  cycle), IV (2 H\(^+\)); final electron acceptor = O\(_2\). Modest
  electron leak produces ROS (\cref{eq:unit_III_ros_cascade}); SOD,
  catalase, GPx, and peroxiredoxins constrain damage.
\item
  ATP synthase: F\(_0\)F\(_1\) rotary motor; binding change mechanism;
  about 3.67 H\(^+\) per ATP (including phosphate import).
  Proton-motive force \texorpdfstring{\ensuremath{\approx}}{~} 180 mV (\cref{eq:unit_III_proton_motive_force}).
\item
  Modern ATP yield: about 30--32 ATP per glucose aerobically
  (\cref{eq:unit_III_atp_yield}, \cref{eq:unit_III_po_ratio_nadh},
  \cref{eq:unit_III_po_ratio_fadh2}); 2 ATP per glucose anaerobically.
\item
  Uncoupling proteins (UCP1) regulate physiological thermogenesis in
  brown adipose tissue; pharmacological uncouplers (DNP) dissipate the
  same PMF unregulated and are dangerous.
\item
  \textbf{Connections:} See \cref{sec:unit_III_photosynthesis} for
  light-driven phosphorylation,
  \cref{sec:unit_III_metabolic_integration} for hormonal integration,
  and Unit X for ecosystem productivity and reduced carbon.
\end{itemize}

\begin{center}\rule{0.5\linewidth}{0.5pt}\end{center}

\subsection{Review Questions}\label{review-questions-11}

\begin{enumerate}
\def\labelenumi{\arabic{enumi}.}
\item
  Explain why cells must maintain metabolite concentrations far from
  equilibrium. What would happen to ATP synthesis if the cell reached
  thermodynamic equilibrium?
\item
  Describe most 10 steps of glycolysis. Which three steps are
  irreversible and why are these the regulatory control points?
\item
  Compare the four bypass enzymes of gluconeogenesis with the three
  irreversible glycolytic steps they replace. Why is each bypass
  thermodynamically favourable in the gluconeogenic direction?
\item
  Explain the structure and regulation of the pyruvate dehydrogenase
  complex. Why does it require five cofactors?
\item
  Track the fate of both carbon atoms in acetyl-CoA through one turn of
  the TCA cycle. Are these the same carbons released as CO\(_2\)?
\item
  Distinguish between anaplerotic and cataplerotic reactions. Give two
  examples of each, and explain why pyruvate carboxylase is
  allosterically activated by acetyl-CoA.
\item
  Describe the Q cycle in Complex III. Why is this mechanism necessary,
  and how does it contribute to proton pumping?
\item
  Identify the major ETC sites of ROS production and the antioxidant
  enzymes that detoxify each downstream species. Why is the Fenton
  reaction particularly dangerous?
\item
  Explain the binding change mechanism of ATP synthase. How was the
  rotation of the gamma subunit experimentally demonstrated?
\item
  Calculate the ATP yield from complete oxidation of one molecule of
  glucose via glycolysis, PDC, TCA cycle, and oxidative phosphorylation
  (\cref{eq:unit_III_atp_yield}). Explain why the answer is 30 or 32
  depending on the NADH shuttle used, and trace each number back to
  specific P/O ratios.
\item
  Compare alcoholic and lactic acid fermentation. Why is fermentation
  essential even though it produces far less ATP than aerobic
  respiration?
\item
  Compare physiological (UCP1) and pharmacological (DNP) uncoupling. Why
  does the body tolerate UCP1 activity but not DNP, even though both
  dissipate the same PMF as heat?
\item
  A patient presents with exercise intolerance and lactic acidosis.
  Genetic testing reveals a mutation in the ND1 gene (Complex I
  subunit). Explain the biochemical basis of each symptom.
\item
  Explain the Pasteur effect and the Warburg effect. How do they
  represent opposite metabolic strategies?
\item
  Recompute \texttt{atp\_free\_energy} when cytosolic ATP drops to 1 mM
  while ADP rises to 4 mM (Pi fixed). Does hydrolysis become more or
  less exergonic, and why does that matter during exercise?
\item
  Using \breaktt{glycolysis\_summary()} from the same module, state the
  net ATP and NADH from one glucose without writing most ten steps.
\item
  Using \(\Delta G^{\circ\prime}\) values for ATP and PEP hydrolysis,
  compute the standard free energy change of pyruvate kinase (PEP + ADP
  \texorpdfstring{\ensuremath{\to}}{->} pyruvate + ATP). Compare with the cellular \(\Delta G\) and explain
  the irreversibility.
\item
  Compute the proton motive force from \(\Delta\Psi = -160\) mV and
  \(\Delta\text{pH} = 0.6\) at 37 °C using
  \cref{eq:unit_III_proton_motive_force}. Convert to free energy per
  mole of protons.
\end{enumerate}

\begin{center}\rule{0.5\linewidth}{0.5pt}\end{center}

\subsection{Further Reading and Source
Notes}\label{further-reading-and-source-notes-11}

\begin{itemize}
\tightlist
\item
  Mitchell (1961). Coupling of phosphorylation to electron and hydrogen
  transfer by a chemi-osmotic type of mechanism. \emph{Nature}, 191
  \citep{mitchell1961}.
\item
  Atkinson (1968). The regulation of enzyme activity in metabolism.
  \emph{Biochemical Journal}, 108 \citep{atkinson1968}.
\end{itemize}

\begin{center}\rule{0.5\linewidth}{0.5pt}\end{center}

\subsection{Key Terms}\label{key-terms-11}

{\def\LTcaptype{none} % do not increment counter
\begin{longtable}[]{@{}
  >{\raggedright\arraybackslash}p{(\linewidth - 2\tabcolsep) * \real{0.2857}}
  >{\raggedright\arraybackslash}p{(\linewidth - 2\tabcolsep) * \real{0.7143}}@{}}
\toprule\noalign{}
\begin{minipage}[b]{\linewidth}\raggedright
Term
\end{minipage} & \begin{minipage}[b]{\linewidth}\raggedright
Definition
\end{minipage} \\
\midrule\noalign{}
\endhead
\bottomrule\noalign{}
\endlastfoot
\textbf{Gibbs free energy (\(\Delta G\))} & Thermodynamic quantity
predicting reaction spontaneity; \(\Delta G < 0\) = spontaneous \\
\textbf{ATP} & Adenosine triphosphate; central phosphoryl-group donor
and short-term energy carrier; \(\Delta G\) hydrolysis \(\approx -54\)
kJ/mol (cellular) \\
\textbf{Glycolysis} & 10-step cytoplasmic pathway converting glucose to
2 pyruvate + 2 ATP + 2 NADH \\
\textbf{Gluconeogenesis} & Synthesis of glucose from non-carbohydrate
precursors using four bypass enzymes; 6 ATP equivalents per glucose \\
\textbf{PFK-1} & Phosphofructokinase-1; committed step of glycolysis;
allosterically regulated \\
\textbf{Pyruvate dehydrogenase complex} & Multienzyme complex linking
glycolysis to TCA; 5 cofactors; regulated by phosphorylation \\
\textbf{TCA cycle} & 8-step mitochondrial cycle oxidising acetyl-CoA;
produces 3 NADH + 1 FADH\(_2\) + 1 GTP per turn \\
\textbf{Anaplerotic reaction} & Reaction replenishing TCA cycle
intermediates (e.g., pyruvate carboxylase) \\
\textbf{Cataplerotic reaction} & Reaction removing TCA intermediates for
biosynthesis (e.g., citrate \texorpdfstring{\ensuremath{\to}}{->} acetyl-CoA for FA synthesis) \\
\textbf{Electron transport chain} & Complexes I--IV on IMM; transfers
electrons from NADH/FADH\(_2\) to O\(_2\), pumping H\(^+\) \\
\textbf{Q cycle} & Complex III mechanism doubling H\(^+\) pumping
efficiency per electron pair \\
\textbf{ATP synthase} & F\(_0\)F\(_1\) rotary motor using PMF to
synthesise ATP; about 3.67 H\(^+\)/ATP including P\(_i\)
import \\
\textbf{Proton motive force (PMF)} & Electrochemical gradient of H\(^+\)
across IMM; \(\approx 180\) mV; drives ATP synthesis \\
\textbf{P/O ratio} & ATP molecules produced per oxygen atom reduced;
about 2.5 for NADH, about 1.5 for FADH\(_2\) \\
\textbf{Reactive oxygen species (ROS)} & O\(_2^{\bullet-}\),
H\(_2\)O\(_2\), HO\(^\bullet\) produced by electron leak; constrained by
SOD, catalase, GPx \\
\textbf{Uncoupling protein (UCP1)} & IMM proton channel in brown adipose
tissue; dissipates PMF as heat \\
\textbf{Fermentation} & Anaerobic regeneration of NAD\(^+\); lactic acid
or ethanol pathways; 2 ATP per glucose \\
\textbf{Cori cycle} & Lactate (muscle) to glucose (liver) shuttle;
shifts gluconeogenic cost to liver \\
\textbf{Glucose-alanine cycle} & Inter-organ shuttle carrying both
carbon and waste nitrogen from muscle to liver \\
\textbf{Warburg effect} & Aerobic glycolysis in cancer cells; rapid ATP
+ biosynthetic precursors despite O\(_2\) \\
\textbf{Substrate-level phosphorylation} & Direct transfer of phosphoryl
group from high-energy intermediate to ADP \\
\end{longtable}
}

\begin{center}\rule{0.5\linewidth}{0.5pt}\end{center}

\subsubsection{Companion Source
Module}\label{companion-source-module-11}

\textbf{Bioenergetics and Cellular Respiration} should leave a
reproducible trail from a biological claim to the code, figure, diagram,
or paper-based activity that can test it. Use the surfaces below to
inspect the chapter's assumptions, rerun the relevant model, or compare
the manuscript explanation with companion labs and figures.

{\def\LTcaptype{none} % do not increment counter
\begin{longtable}[]{@{}
  >{\raggedright\arraybackslash}p{(\linewidth - 2\tabcolsep) * \real{0.5000}}
  >{\raggedright\arraybackslash}p{(\linewidth - 2\tabcolsep) * \real{0.5000}}@{}}
\toprule\noalign{}
\begin{minipage}[b]{\linewidth}\raggedright
Surface
\end{minipage} & \begin{minipage}[b]{\linewidth}\raggedright
Use it for
\end{minipage} \\
\midrule\noalign{}
\endhead
\bottomrule\noalign{}
\endlastfoot
\breaktt{src/biology/biochemistry/biochemistry.py}
(\breaktt{reaction\_free\_energy}, \texttt{atp\_free\_energy},
\breaktt{glycolysis\_summary}) & Track energy accounting across
glycolysis, respiration, and ATP coupling. \\
\breaktt{src/mermaid/biology\_diagrams.py}
(\breaktt{glycolysis\_pathway\_diagram},
\breaktt{atp\_synthesis\_diagram}) & Check pathway order and coupling
between electron flow and proton motive force. \\
\end{longtable}
}

\textbf{Reproducibility check:} name the electron donor, acceptor,
compartment, proton path, and ATP-yield assumption before comparing
respiratory claims. \textbf{Cross-reference:} connect with
\cref{sec:unit_III_photosynthesis},
\cref{sec:unit_III_metabolic_integration}, and
\cref{sec:unit_I_enzymes_and_kinetics}.

\newpage

\section{Photosynthesis}\label{photosynthesis}

\label{sec:unit_III_photosynthesis}

\begin{quote}
\textbf{Ch 10} \texorpdfstring{\ensuremath{\cdot}}{.} Level 2/3 \texorpdfstring{\ensuremath{\cdot}}{.} 55 min read \texorpdfstring{\ensuremath{\cdot}}{.} 75 min lecture \texorpdfstring{\ensuremath{\cdot}}{.}
Prerequisites: \cref{sec:unit_III_bioenergetics_and_respiration}
\end{quote}

\subsection{Learning Objectives}\label{learning-objectives-12}

\begin{enumerate}
\def\labelenumi{\arabic{enumi}.}
\tightlist
\item
  Describe \hyperref[gl:chloroplast]{\textbf{chloroplast}} structure and
  the organisation of the photosynthetic membrane.
\item
  Explain the light-harvesting antenna complexes, chlorophyll a/b ratio,
  and energy transfer mechanisms (FRET, exciton transfer).
\item
  Compute the energy of a photon and convert per-molecule values to
  per-mole values for biological calculations.
\item
  Describe Photosystem II, water splitting, and the Kok cycle in detail.
\item
  Explain the cytochrome b\(_6\)f complex and the Q cycle in
  chloroplasts.
\item
  Describe Photosystem I, linear and cyclic electron flow, and NADPH
  production.
\item
  Explain the Z-scheme of photosynthetic electron transport, including
  standard electrode potentials of each carrier.
\item
  Describe most three phases of the
  \hyperref[gl:calvin-cycle]{\textbf{Calvin cycle}}, including RuBisCO
  biochemistry, structure (L\(_8\)S\(_8\)), catalytic mechanism, and the
  CO\(_2\)/O\(_2\) specificity factor (\texorpdfstring{\ensuremath{\tau}}{tau}).
\item
  Explain the thioredoxin/ferredoxin redox regulation system and the
  activating role of stromal pH and Mg\(^{2+}\).
\item
  Explain photorespiration and its metabolic costs.
\item
  Compare C3, C4, and CAM
  \hyperref[gl:photosynthesis]{\textbf{photosynthesis}} strategies and
  their water-use efficiency, including quantitative trade-offs.
\item
  Evaluate photoprotection (NPQ, xanthophyll cycle) and artificial
  photosynthesis research as responses to excess light and energy
  demand.
\item
  Calculate ATP and NADPH requirements per CO\(_2\) fixed; compare
  biological and artificial photosynthesis efficiencies.
\end{enumerate}

\subsubsection{Study Blueprint}\label{study-blueprint-12}

\begin{itemize}
\tightlist
\item
  \textbf{Big idea:} Photosynthesis couples light-driven electron flow
  to carbon fixation and planetary productivity.
\item
  \textbf{Core concepts:} photosystems, electron transport, Calvin
  cycle, photorespiration.
\item
  \textbf{Framework alignment:} Vision \& Change: Pathways and
  transformations of energy and matter, Systems; AP Biology: Energetics,
  Systems Interactions; NGSS-style topics: Matter and Energy in
  Organisms and Ecosystems.
\item
  \textbf{Model or quantitative lens:} Photon, ATP/NADPH, and
  Calvin-cycle stoichiometry.
\item
  \textbf{Data skill:} Interpret light-response and carbon-fixation
  data.
\item
  \textbf{Practice cadence:} Representing and Describing Data,
  Statistical Tests and Data Analysis.
\item
  \textbf{Common misconception to repair:} Plants do not eat sunlight;
  they use light energy to reduce carbon using electrons and enzymes.
\item
  \textbf{Primary lab:} \cref{sec:lab_unit_III_photosynthesis}.
\item
  \textbf{Question bank:} \cref{sec:q_unit_III_photosynthesis}.
\item
  \textbf{Transfer task:} Transfer photosynthetic constraints to crops,
  algae, climate, or ecosystem productivity.
\item
  \textbf{Bridge to computation:}
  \breaktt{biology.botany.botany.photosynthesis\_rate}.
\end{itemize}

\begin{center}\rule{0.5\linewidth}{0.5pt}\end{center}

\begin{quote}
\textbf{Opening Vignette: Tracing Carbon From Air to Sugar}

In 1950, Melvin Calvin's team at the Lawrence Berkeley National
Laboratory began exposing algae (\emph{Chlorella}) to \texorpdfstring{\textsuperscript{1}}{1}\texorpdfstring{\textsuperscript{4}}{4}CO\texorpdfstring{\textsubscript{2}}{2} --- carbon
dioxide labelled with radioactive carbon-14 --- for very brief intervals
(as short as 5 seconds), then rapidly killing the cells and separating
their compounds by two-dimensional paper chromatography. By tracking
which molecules became radioactively labelled first, then in which
order, Calvin meticulously traced the path of carbon atoms from
atmospheric CO\texorpdfstring{\textsubscript{2}}{2} through a cycle of enzymatic reactions to eventually
emerge as glucose. The complete cycle --- the
\textbf{Calvin-Benson-Bassham cycle} --- was published in 1954, and
Calvin received the Nobel Prize in Chemistry alone in 1961.

The cycle's key \hyperref[gl:enzyme]{\textbf{enzyme}}, \textbf{RuBisCO}
(ribulose-1,5-bisphosphate carboxylase/oxygenase), is arguably the most
important enzyme on Earth: it is responsible for fixing approximately
10\texorpdfstring{\textsuperscript{1}}{1}\texorpdfstring{\textsuperscript{1}}{1} tonnes of atmospheric carbon per year, feeding essentially most
non-chemotrophic life. It is also probably the most abundant
\hyperref[gl:protein]{\textbf{protein}} on Earth (approximately 0.7 kg
per person on the planet's surface, an estimated 7 \texorpdfstring{\ensuremath{\times}}{x} 10\texorpdfstring{\textsuperscript{8}}{8} tonnes total),
and notably one of the slowest enzymes known, with a kcat of about 3--10
s\texorpdfstring{\textsuperscript{-}}{-}\texorpdfstring{\textsuperscript{1}}{1}. Improving RuBisCO efficiency is one of the holy grails of
agricultural biotechnology.

\emph{Primary source: Calvin, M. \& Benson, A. A. (1948). The path of
carbon in photosynthesis. Science, 107(2784), 476--480.}
\end{quote}

\begin{center}\rule{0.5\linewidth}{0.5pt}\end{center}

\subsection{Overview of
Photosynthesis}\label{overview-of-photosynthesis}

Photosynthesis is one of Earth's dominant chemical processes: it fixes
approximately \textbf{120 Gt of carbon} from CO\(_2\) per year, sustains
atmospheric oxygen, and supports nearly every heterotrophic food web.

\begin{equation}
6\text{CO}_2 + 6\text{H}_2\text{O} \xrightarrow{\text{light}} \text{C}_6\text{H}_{12}\text{O}_6 + 6\text{O}_2 \quad (\Delta G^{\circ\prime} = +2{,}870 \; \text{kJ/mol})
\label{eq:unit_III_photosynthesis_overall}
\end{equation}

The reaction is highly endergonic --- it requires light energy to drive
it. Photosynthesis occurs in the chloroplast (\textbf{plants and algae})
and at the plasma membrane (\textbf{cyanobacteria}). The process can be
divided into:

\begin{enumerate}
\def\labelenumi{\arabic{enumi}.}
\tightlist
\item
  \textbf{Light-dependent reactions} (thylakoid membrane): H\(_2\)O \texorpdfstring{\ensuremath{\to}}{->}
  O\(_2\) + NADPH + ATP
\item
  \textbf{Light-independent reactions / Calvin cycle} (stroma): CO\(_2\)
  + ATP + NADPH \texorpdfstring{\ensuremath{\to}}{->} G3P (sugar)
\end{enumerate}

\subsubsection{Chloroplast Structure}\label{chloroplast-structure}

\begin{itemize}
\tightlist
\item
  \textbf{Outer envelope:} freely permeable (porin-like channels)
\item
  \textbf{Inner envelope:} selective transporters (triose phosphate /
  phosphate antiporter, TPT)
\item
  \textbf{Thylakoid membrane:} site of light reactions; forms stacked
  \textbf{grana} (5--20 thylakoids per granum) + interconnecting
  \textbf{stroma lamellae}
\item
  \textbf{Lumen (thylakoid lumen):} acidified
  (\hyperref[gl:ph]{\textbf{pH}} about 5) by proton pumping;
  about 1,000-fold H\(^+\) gradient
\item
  \textbf{Stroma:} site of Calvin cycle; contains RuBisCO, enzymes,
  chloroplast DNA, 70S \hyperref[gl:ribosome]{\textbf{ribosome}}s
\end{itemize}

Thylakoid membranes are densely packed with protein complexes ---
about 70\% protein by mass, making them one of the most
protein-rich membranes in nature. Each reaction centre is associated
with about 250 antenna
\hyperref[gl:chlorophyll]{\textbf{chlorophyll}} molecules.

\textbf{Chloroplast \hyperref[gl:genome]{\textbf{genome}}:} 120--200 kb
circular DNA; encodes about 120
\hyperref[gl:gene]{\textbf{gene}}s including RuBisCO large subunit,
photosystem subunits, and ribosomal components. Like mitochondria,
chloroplasts originated from
\hyperref[gl:endosymbiosis]{\textbf{endosymbiosis}} (cyanobacteria; see
\cref{sec:unit_II_cell_theory}).

\begin{center}\rule{0.5\linewidth}{0.5pt}\end{center}

\subsection{Light Absorption and Energy
Transfer}\label{light-absorption-and-energy-transfer}

\subsubsection{Photosynthetic Pigments}\label{photosynthetic-pigments}

Chlorophylls absorb light primarily at:

\begin{itemize}
\tightlist
\item
  \textbf{Chlorophyll a:} 430 nm (blue-violet) and 680 nm (red)
\item
  \textbf{Chlorophyll b:} 450 nm and 640 nm (broader absorption;
  accessory pigment)
\item
  \textbf{Carotenoids:} 400--500 nm; \texorpdfstring{\ensuremath{\beta}}{beta}-carotene, lutein, zeaxanthin;
  also transfer energy to Chl a; function as photoprotectants
\end{itemize}

The \textbf{green wavelengths (about 500--600 nm)} are
reflected --- hence plant colour.

\subsubsection{Photon Energy and Quantum
Efficiency}\label{photon-energy-and-quantum-efficiency}

A \textbf{photon} of red light (\texorpdfstring{\ensuremath{\lambda}}{lambda} = 680 nm) carries:

\begin{equation}
E = \frac{hc}{\lambda} = \frac{(6.626 \times 10^{-34}\;\text{J\texorpdfstring{\ensuremath{\cdot}}{.}s})(3 \times 10^8\;\text{m/s})}{680 \times 10^{-9}\;\text{m}} = 2.93 \times 10^{-19}\;\text{J} = 1.82\;\text{eV}
\label{eq:unit_III_photon_energy}
\end{equation}

Per mole: \(E = N_A \times E_{\text{photon}} = 176\;\text{kJ/mol}\).

By comparison, blue light (440 nm) photons carry about 272
kJ/mol --- substantially more energy than red. The ``red drop'' in
photosynthetic action spectra reflects the fact that energy in excess of
the redox requirement is dissipated as heat \emph{whichever} photon is
absorbed; longer-wavelength light is therefore proportionally more
efficient on a per-photon basis.

\textbf{Quantum yield of photosynthesis:} The quantum yield (\texorpdfstring{\ensuremath{\Phi}}{Phi}) is the
molecules of product per photon absorbed. Two important quantum yields
are defined:

\begin{itemize}
\tightlist
\item
  \textbf{\texorpdfstring{\ensuremath{\Phi}}{Phi}\(_{\text{O}_2}\)} \texorpdfstring{\ensuremath{\approx}}{~} 0.10--0.12 mol O\(_2\) per mol photons
  absorbed (under optimal conditions)
\item
  \textbf{\texorpdfstring{\ensuremath{\Phi}}{Phi}\(_{\text{CO}_2}\)} \texorpdfstring{\ensuremath{\approx}}{~} 0.08--0.10 mol CO\(_2\) per mol photons
  (lower because of photorespiration in C3)
\end{itemize}

The \textbf{minimum quantum requirement} (1/\texorpdfstring{\ensuremath{\Phi}}{Phi}) is therefore
about 8--10 photons per O\(_2\) released --- close to the
theoretical minimum of 8 photons (4 photons through PS II + 4 through PS
I per O\(_2\) produced). The remaining inefficiency reflects energy
losses in antenna transfer, NPQ activation, and photorespiration.

\textbf{Quantum efficiency (energy basis):}

\begin{equation}
\eta_{\text{quantum}} = \frac{n_{\text{electrons}} \cdot \Delta E_{\text{redox}}}{n_{\text{photons}} \cdot E_{\text{photon}}}
\label{eq:unit_III_quantum_efficiency}
\end{equation}

For the Z-scheme: 4 electrons span 1.14 V from H\(_2\)O to NADPH,
requiring about 440 kJ/mol of electron flow. With 8 photons of
680 nm light supplying about 1,408 kJ/mol, the redox efficiency
is about 31\% --- the rest is consumed by overpotentials and
dissipated as heat.

\textbf{Comparison with engineered photovoltaics.} A commercial
single-junction silicon solar cell at 22\% efficiency converts solar
irradiance into electricity that must then power \emph{something} ---
and the chain from electricity \texorpdfstring{\ensuremath{\to}}{->} stored chemical fuel adds losses.
Plants, by contrast, deliver a stored fuel (sugar) directly. The
artificial-photosynthesis discussion returns to this comparison;
\cref{eq:unit_III_quantum_efficiency} sets the upper biological bound.

\subsubsection{Antenna Complexes and Light-Harvesting
Architecture}\label{antenna-complexes-and-light-harvesting-architecture}

Each reaction centre in plants is served by a peripheral antenna of
\textbf{light-harvesting complexes (LHC)}. LHCII alone is the most
abundant membrane protein on Earth (about 50\% of thylakoid
membrane protein).

\textbf{Chlorophyll a/b ratio:}

{\def\LTcaptype{none} % do not increment counter
\begin{longtable}[]{@{}
  >{\raggedright\arraybackslash}p{(\linewidth - 4\tabcolsep) * \real{0.2500}}
  >{\raggedright\arraybackslash}p{(\linewidth - 4\tabcolsep) * \real{0.4643}}
  >{\raggedright\arraybackslash}p{(\linewidth - 4\tabcolsep) * \real{0.2857}}@{}}
\toprule\noalign{}
\begin{minipage}[b]{\linewidth}\raggedright
Complex
\end{minipage} & \begin{minipage}[b]{\linewidth}\raggedright
Chl a:b ratio
\end{minipage} & \begin{minipage}[b]{\linewidth}\raggedright
Function
\end{minipage} \\
\midrule\noalign{}
\endhead
\bottomrule\noalign{}
\endlastfoot
LHCII (peripheral PS II antenna) & about 1.3 & Light
harvesting; mobile under state transitions \\
CP43, CP47 (core PS II antennae) & Most Chl a & Direct excitation
transfer to P680 \\
LHCI (PS I antenna) & about 3.5 & Light harvesting for P700 \\
Reaction centres (P680, P700) & Most Chl a & Charge separation
chromophores \\
\end{longtable}
}

The Chl a:b ratio of an entire leaf typically ranges from
about 2.5 (sun-adapted) to about 3.5 (shade-adapted
leaves invest more in Chl b--rich LHCII to broaden the absorbed
spectrum). Chl b's broader 450 nm/640 nm absorption is critical for
harvesting green-shifted light filtered through canopy chlorophyll
above.

\textbf{Carotenoid functions:}

\begin{enumerate}
\def\labelenumi{\arabic{enumi}.}
\tightlist
\item
  \textbf{Light harvesting} in the 450--500 nm range, where Chl absorbs
  poorly (about 30\% of the antenna's spectral coverage)
\item
  \textbf{Triplet quenching:} \(^3\)Chl* (long-lived triplet state,
  formed when singlet excitations are not used quickly) is rapidly
  quenched by carotenoids, preventing \(^1\)O\(_2\) generation
\item
  \textbf{Singlet oxygen scavenging:} any \(^1\)O\(_2\) that does form
  is quenched by carotenoids (the \emph{first} line of antioxidant
  defence)
\item
  \textbf{Structural stabilisation} of LHC protein folds
\end{enumerate}

\subsubsection{Energy Transfer
Mechanisms}\label{energy-transfer-mechanisms}

\textbf{Förster resonance energy transfer (FRET):} Dipole--dipole
coupling between donor (D*) and acceptor (A) chromophores transfers
excitation without electron transfer.

\begin{equation}
k_{\text{FRET}} = \frac{1}{\tau_D} \left( \frac{R_0}{r} \right)^6
\label{eq:unit_III_fret_rate}
\end{equation}

where \(r\) is donor--acceptor separation, \(\tau_D\) is the donor's
excited-state lifetime, and \(R_0\) is the \textbf{Förster radius}
(typically 5--10 nm in Chl networks). The \(r^6\) dependence makes FRET
extremely sensitive to distance --- doubling separation reduces transfer
rate 64-fold. \(R_0\) depends on:

\begin{itemize}
\tightlist
\item
  \textbf{Spectral overlap integral} \(J(\lambda)\): the donor emission
  must overlap the acceptor absorption (red-shifting toward the reaction
  centre)
\item
  \textbf{Orientation factor} \(\kappa^2\): relative orientation of
  donor emission and acceptor absorption transition dipoles (averages to
  2/3 in random orientations, but is precisely tuned in the LHCII
  crystal structure)
\end{itemize}

\textbf{Energy funnel architecture:} Antenna pigments are arranged with
higher-energy absorbing pigments at the periphery and lower-energy
pigments closer to the reaction centre. Excitation cascades ``downhill''
toward the reaction centre with each transfer step, achieving
about 95\% efficiency from initial absorption to charge
separation.

\textbf{Exciton transfer (Dexter mechanism):} When chromophores are
within about 1 nm, electronic wavefunctions overlap and
excitation transfers as a coherent exciton rather than incoherent FRET
hops. Recent 2D electronic spectroscopy (Fleming lab, 2007) revealed
that quantum coherence may persist for picoseconds at room temperature,
possibly contributing to the high transfer efficiency, though the
functional importance remains debated.

\begin{quote}
\textbf{Concept Check 1:} The antenna complex transfers excitation
energy to the reaction centre with about 95\% efficiency. If
100 photons are absorbed by the antenna, how many excitations reach the
reaction centre? What happens to the energy of the about 5
photons that do not reach the reaction centre?
\end{quote}

\begin{quote}
\textbf{Concept Check 1b:} \cref{eq:unit_III_fret_rate} shows that
\(k_{\text{FRET}} \propto 1/r^6\). If two chlorophylls separated by 5 nm
transfer at \(1 \times 10^{10}\) s\(^{-1}\), predict the rate at 7.5 nm
and at 10 nm. What does this mean for the spatial design of LHC
complexes?
\end{quote}

\begin{center}\rule{0.5\linewidth}{0.5pt}\end{center}

\subsection{Light-Dependent Reactions}\label{light-dependent-reactions}

\subsubsection{Photosystem II (PS II)}\label{photosystem-ii-ps-ii}

PS II is located primarily in appressed grana thylakoids. It is a large
complex (about 350 kDa per monomer; functions as a dimer). Key
components:

\begin{itemize}
\tightlist
\item
  \textbf{Reaction centre (P680):} Special pair of Chl a molecules that
  absorb at 680 nm; the primary electron donor
\item
  \textbf{D1 and D2 proteins:} Core reaction centre proteins (homologous
  to bacterial L and M subunits)
\item
  \textbf{Pheophytin (Pheo):} Primary electron acceptor (chlorophyll
  lacking Mg\(^{2+}\))
\item
  \textbf{Q\(_A\) and Q\(_B\):} Plastoquinone molecules; Q\(_A\) is
  tightly bound (one-electron acceptor), Q\(_B\) is the
  two-electron/two-proton mobile acceptor
\item
  \textbf{Water-splitting complex (OEC, oxygen-evolving complex):}
  Mn\(_4\)CaO\(_5\) cluster; performs the most thermodynamically
  demanding oxidation in biology
\end{itemize}

\textbf{Water oxidation:}

\begin{equation}
2\text{H}_2\text{O} \rightarrow \text{O}_2 + 4\text{H}^+ + 4\text{e}^- \quad (E^{\circ\prime} = +0.82 \text{ V})
\label{eq:unit_III_water_splitting}
\end{equation}

Each molecule of O\(_2\) requires 4 photons (one photon per electron
from water). The OEC's ability to extract electrons one at a time from
water --- accumulating four oxidising equivalents on the
Mn\(_4\)CaO\(_5\) cluster before releasing O\(_2\) in a single concerted
four-electron step --- was unprecedented in biology and is the
inspiration for synthetic water-splitting catalysts.

\textbf{The Kok cycle (S-state cycle):} The OEC cycles through 5
oxidation states (S\(_0\)--S\(_4\)), accumulating 4 oxidising
equivalents before releasing O\(_2\):

\[\text{S}_0 \xrightarrow{h\nu} \text{S}_1 \xrightarrow{h\nu} \text{S}_2 \xrightarrow{h\nu} \text{S}_3 \xrightarrow{h\nu} \text{S}_4 \xrightarrow{-\text{O}_2} \text{S}_0 \tag{10.4}\]

Each S-state transition involves one photon absorption and one electron
extraction from the Mn\(_4\)Ca cluster.

\textbf{Primary photochemistry:} P680 absorbs photon \texorpdfstring{\ensuremath{\to}}{->} excited state
P680* \texorpdfstring{\ensuremath{\to}}{->} donates electron to pheophytin (charge separation in
about 3 ps) \texorpdfstring{\ensuremath{\to}}{->} Q\(_A\) \texorpdfstring{\ensuremath{\to}}{->} Q\(_B\). After two electrons and two
protons, PQH\(_2\) (plastoquinol) leaves PS II and carries electrons to
Complex III (cytochrome b\(_6\)f).

\textbf{D1 protein damage:} P680\(^+\) is the strongest biological
oxidant (\(E^{\circ\prime} = +1.25\) V). The D1 protein is damaged by
this oxidative stress and must be replaced every 30--60 minutes --- the
most rapid protein turnover in the cell. This
\textbf{photoinhibition/repair cycle} requires a dedicated FtsH protease
and ribosomal machinery in the thylakoid.

\begin{quote}
\textbf{Clinical Connection: Herbicides Targeting PS II} \textbf{DCMU
(diuron)} and \textbf{atrazine} block the Q\(_B\) binding site on D1,
preventing electron flow from PS II. This kills plants by halting
photosynthesis. Atrazine is one of the most widely used herbicides
globally. Resistance has evolved in some weed species through a single
Ser264Gly \hyperref[gl:mutation]{\textbf{mutation}} in D1 that reduces
atrazine binding. Understanding PS II structure is essential for
developing new herbicides and for bioengineering more efficient
photosynthesis.
\end{quote}

\subsubsection{\texorpdfstring{Cytochrome b\(_6\)f
Complex}{Cytochrome b_6f Complex}}\label{cytochrome-b_6f-complex}

After PS II, PQH\(_2\) is oxidised by the \textbf{cytochrome b\(_6\)f
complex}, which pumps H\(^+\) into the lumen via the \textbf{Q-cycle},
analogous to Complex III in mitochondria:

\begin{enumerate}
\def\labelenumi{\arabic{enumi}.}
\tightlist
\item
  PQH\(_2\) binds at the luminal (Q\(_p\)) site
\item
  One electron passes to the Rieske Fe-S protein \texorpdfstring{\ensuremath{\to}}{->} cytochrome f \texorpdfstring{\ensuremath{\to}}{->}
  \textbf{plastocyanin} (PC, soluble Cu-protein in the lumen)
\item
  Second electron passes to cytochrome b\(_6\) (low potential) \texorpdfstring{\ensuremath{\to}}{->}
  cytochrome b\(_6\) (high potential) \texorpdfstring{\ensuremath{\to}}{->} PQ at the stromal (Q\(_n\)) site
\item
  Net: 2 H\(^+\) released to lumen per PQH\(_2\) oxidised, plus 2
  H\(^+\) consumed from stroma per PQ reduced
\end{enumerate}

The b\(_6\)f complex also produces superoxide as a byproduct, which can
trigger signalling cascades for acclimation to changing light
conditions.

\subsubsection{Photosystem I (PS I)}\label{photosystem-i-ps-i}

PS I reaction centre: \textbf{P700} (Chl a pair absorbing at 700 nm). PS
I is located primarily in non-appressed stroma lamellae and at the
margins of grana.

\textbf{Electron flow through PS I:} P700 + photon \texorpdfstring{\ensuremath{\to}}{->} P700* \texorpdfstring{\ensuremath{\to}}{->} A\(_0\)
(Chl a) \texorpdfstring{\ensuremath{\to}}{->} A\(_1\) (phylloquinone) \texorpdfstring{\ensuremath{\to}}{->} F\(_X\) (Fe-S) \texorpdfstring{\ensuremath{\to}}{->} F\(_A\)/F\(_B\)
(Fe-S) \texorpdfstring{\ensuremath{\to}}{->} \textbf{ferredoxin} (Fd, soluble Fe-S protein in stroma) \texorpdfstring{\ensuremath{\to}}{->}
\textbf{ferredoxin-NADP\(^+\) reductase (FNR)} \texorpdfstring{\ensuremath{\to}}{->} \textbf{NADPH}

\subsubsection{The Z-Scheme: Standard Electrode
Potentials}\label{the-z-scheme-standard-electrode-potentials}

\begin{figure}[htbp]
\centering
\includegraphics[width=0.82\linewidth,height=4.2in,keepaspectratio]{../figures/mermaid_inline/inline_mermaid_0056_eb22b5920e20.png}
\caption{Sequence diagram for The Z-Scheme: Standard Electrode Potentials showing ordered interaction among H2O (E degrees ' = +0.82 V), OEC (Mn4Ca), P680/PS II (E degrees ' = +1.25 V), and Plastoquinone Pool (PQH2).}
\end{figure}

\emph{The Z-scheme of photosynthetic electron transport (Mermaid).} Two
photons (absorbed by PS II and PS I) drive each electron from water
(\(E^{\circ\prime} = +0.82\) V) to NADPH (\(E^{\circ\prime} = -0.32\)
V), spanning a total potential of 1.14 V. The ``Z'' shape arises from
plotting electron carriers against their redox potential.

\textbf{Standard reduction potentials of key photosynthetic carriers:}

{\def\LTcaptype{none} % do not increment counter
\begin{longtable}[]{@{}
  >{\raggedright\arraybackslash}p{(\linewidth - 4\tabcolsep) * \real{0.1622}}
  >{\raggedright\arraybackslash}p{(\linewidth - 4\tabcolsep) * \real{0.7297}}
  >{\raggedright\arraybackslash}p{(\linewidth - 4\tabcolsep) * \real{0.1081}}@{}}
\toprule\noalign{}
\begin{minipage}[b]{\linewidth}\raggedright
Couple
\end{minipage} & \begin{minipage}[b]{\linewidth}\raggedright
\(E^{\circ\prime}\) (V, pH 7)
\end{minipage} & \begin{minipage}[b]{\linewidth}\raggedright
Role
\end{minipage} \\
\midrule\noalign{}
\endhead
\bottomrule\noalign{}
\endlastfoot
O\(_2\)/H\(_2\)O & \(+0.82\) & Electron donor (oxidised by S\(_4\)) \\
Tyr\(_Z^\bullet\)/Tyr\(_Z\) & \(+0.97\) & OEC-to-P680\(^+\) relay \\
\textbf{P680\(^+\)/P680 (ground)} & \textbf{\(+1.25\)} & The strongest
biological oxidant \\
**P680*/P680 (excited)** & \textbf{\(\sim -0.6\)} & After photon
absorption \\
Pheophytin (Pheo\(^-\)/Pheo) & \(-0.61\) & First acceptor in PS II \\
Q\(_A^-\)/Q\(_A\) & \(-0.13\) & Tightly bound plastoquinone \\
Q\(_B\)/QH\(_2\) & \(0.0\) to \(+0.10\) & Mobile plastoquinone pool \\
Cytochrome f & \(+0.36\) & b\(_6\)f exit point \\
Plastocyanin (Cu\(^{2+}\)/Cu\(^+\)) & \(+0.37\) & Lumen carrier to PS
I \\
\textbf{P700\(^+\)/P700 (ground)} & \textbf{\(+0.45\)} & PS I reaction
centre \\
**P700*/P700 (excited)** & \textbf{\(\sim -1.3\)} & Strongest biological
reductant \\
A\(_0\) (Chl a\(^-\)/Chl a) & \(-1.0\) & First PS I acceptor \\
Ferredoxin (Fe-S\(^{2+}\)/Fe-S\(^+\)) & \(-0.42\) & Stromal reductant \\
NADP\(^+\)/NADPH & \(-0.32\) & Final electron acceptor \\
\end{longtable}
}

The two large upward jumps (P680\texorpdfstring{\ensuremath{\to}}{->}P680\emph{, P700\texorpdfstring{\ensuremath{\to}}{->}P700}) require
absorption of red-light photons. The intervening downward drops drive
PMF buildup and NADPH formation. This gives the canonical ``Z'' shape.

The Gibbs free energy stored in moving 4 electrons from water to NADPH:

\begin{equation}
\Delta G^{\circ\prime} = -nF\Delta E = -(4)(96{,}485)(1.14) \approx -440\;\text{kJ/mol}_{\text{e}^-}
\label{eq:unit_III_z_scheme_energy}
\end{equation}

(N.B. negative sign indicates the \emph{reverse} reaction is
spontaneous; light energy is required to drive electrons \emph{uphill}
from water to NADPH.)

\begin{quote}
\textbf{Concept Check 2x:} Plastocyanin (\(E^{\circ\prime} = +0.37\) V)
reduces P700\(^+\) (\(E^{\circ\prime} = +0.45\) V). Calculate
\(\Delta G^{\circ\prime}\) for this electron transfer per mole of
electrons. Does the redox flow proceed downhill in the dark or primarily
after PS I activation?
\end{quote}

\subsubsection{Linear vs.~Cyclic Electron
Flow}\label{linear-vs.-cyclic-electron-flow}

\textbf{Linear electron flow (LEF):} H\(_2\)O \texorpdfstring{\ensuremath{\to}}{->} PS II \texorpdfstring{\ensuremath{\to}}{->} PQ \texorpdfstring{\ensuremath{\to}}{->} b\(_6\)f \texorpdfstring{\ensuremath{\to}}{->}
PC \texorpdfstring{\ensuremath{\to}}{->} PS I \texorpdfstring{\ensuremath{\to}}{->} Fd \texorpdfstring{\ensuremath{\to}}{->} NADPH. Produces both ATP and NADPH. ATP:NADPH ratio
about 1.3:1.

\textbf{Cyclic electron flow (CEF):} Ferredoxin returns electrons to the
b\(_6\)f complex (via PQ), generating additional PMF for ATP synthesis
without NADPH production. This adjusts the ATP:NADPH ratio to match
Calvin cycle demands (3:2 per CO\(_2\)).

CEF is mediated by two pathways: - \textbf{PGR5/PGRL1 pathway:} Major
route; ferredoxin \texorpdfstring{\ensuremath{\to}}{->} PQ via PGRL1 - \textbf{NDH-dependent pathway:}
Chloroplast NAD(P)H dehydrogenase complex; minor route

The functional necessity of CEF was first inferred from the gap between
the ATP demand of the Calvin cycle (3 ATP per CO\(_2\)) and what LEF
alone supplies (about 2.6 ATP per CO\(_2\) at 14:3 H\(^+\)/ATP
stoichiometry); CEF closes the gap. Hager's classic experiments on the
xanthophyll cycle \citep{hager1971} and subsequent work showed that the
magnitude of CEF varies dynamically with light and CO\(_2\) --- itself a
regulatory layer.

\subsubsection{Chemiosmotic ATP Synthesis
(Chloroplast)}\label{chemiosmotic-atp-synthesis-chloroplast}

The PMF generated by H\(^+\) accumulation in the lumen (from water
splitting + PQ reduction + b\(_6\)f Q-cycle) drives \textbf{chloroplast
\hyperref[gl:atp-synthase]{\textbf{ATP synthase}} (CF\(_1\)F\(_0\))}:

\begin{itemize}
\tightlist
\item
  Structure and mechanism analogous to mitochondrial ATP synthase
\item
  c-ring has 14 subunits in spinach chloroplasts (vs.~8--15 in
  mitochondria)
\item
  about 4.67 H\(^+\)/ATP (14 H\(^+\) per revolution / 3 ATP)
\end{itemize}

\textbf{Net output of light reactions (per 2 H\(_2\)O oxidised / 2 NADPH
produced):}

\begin{itemize}
\tightlist
\item
  1 O\(_2\)
\item
  2 NADPH
\item
  about 3 ATP (from linear flow; cyclic flow provides
  additional ATP)
\end{itemize}

\subsubsection{Worked Example: Energy Requirements for Carbon
Fixation}\label{worked-example-energy-requirements-for-carbon-fixation}

\emph{Problem:} Calculate the ATP and NADPH consumed to synthesise one
molecule of glucose, and determine how many photons must be absorbed
(assuming 8 photons per O\(_2\) for linear electron flow).

\emph{Solution:}

Net Calvin cycle stoichiometry per CO\(_2\) fixed: \textbf{3 ATP + 2
NADPH}, as derived in the Calvin-cycle stoichiometry discussion below.
Per glucose (6 CO\(_2\) fixed): \textbf{18 ATP + 12 NADPH}.

Linear electron flow produces NADPH at a fixed ratio of 2 NADPH per 4
e\(^-\) (4 photons per electron pair), so 12 NADPH require:

\[n_{\text{photons}} = 12 \times 4 = 48\;\text{photons (LEF primarily)} \tag{10.4a}\]

But linear flow produces about 12 ATP from this (12 \texorpdfstring{\ensuremath{\times}}{x} 1.0 ATP per NADPH
at about 1:1 stoichiometry --- far short of the 18 ATP
required). The shortfall of 6 ATP must be supplied by \textbf{cyclic
electron flow}, which is estimated to require about 6--9
additional photons. Total: \textbf{about 54--57 photons per
glucose}.

This is why the often-cited ``48 photons per glucose'' is a \emph{lower
bound} assuming perfect Z-scheme efficiency; real plants need somewhat
more to balance the ATP:NADPH ratio.

\subsubsection{Worked Example: Calvin Cycle
Stoichiometry}\label{worked-example-calvin-cycle-stoichiometry}

\emph{Problem:} Track the 18 ATP and 12 NADPH demand of the Calvin cycle
through the three phases for 6 CO\(_2\).

\emph{Solution:}

{\def\LTcaptype{none} % do not increment counter
\begin{longtable}[]{@{}
  >{\raggedright\arraybackslash}p{(\linewidth - 8\tabcolsep) * \real{0.2667}}
  >{\raggedright\arraybackslash}p{(\linewidth - 8\tabcolsep) * \real{0.2222}}
  >{\raggedright\arraybackslash}p{(\linewidth - 8\tabcolsep) * \real{0.2667}}
  >{\raggedright\arraybackslash}p{(\linewidth - 8\tabcolsep) * \real{0.1333}}
  >{\raggedright\arraybackslash}p{(\linewidth - 8\tabcolsep) * \real{0.1111}}@{}}
\toprule\noalign{}
\begin{minipage}[b]{\linewidth}\raggedright
Calvin phase
\end{minipage} & \begin{minipage}[b]{\linewidth}\raggedright
Per CO\(_2\)
\end{minipage} & \begin{minipage}[b]{\linewidth}\raggedright
Per 6 CO\(_2\)
\end{minipage} & \begin{minipage}[b]{\linewidth}\raggedright
Output
\end{minipage} & \begin{minipage}[b]{\linewidth}\raggedright
Notes
\end{minipage} \\
\midrule\noalign{}
\endhead
\bottomrule\noalign{}
\endlastfoot
Carboxylation (RuBisCO) & 0 ATP, 0 NADPH & 0, 0 & 12 \texorpdfstring{\ensuremath{\times}}{x} 3-PGA & 6 RuBP +
6 CO\(_2\) \texorpdfstring{\ensuremath{\to}}{->} 12 \texorpdfstring{\ensuremath{\times}}{x} 3-PGA \\
Reduction (PGK + GAPDH) & 2 ATP + 2 NADPH & \textbf{12 ATP + 12 NADPH} &
12 \texorpdfstring{\ensuremath{\times}}{x} G3P & Activation + reduction \\
Regeneration (5 G3P \texorpdfstring{\ensuremath{\to}}{->} 3 RuBP) & 1 ATP & \textbf{6 ATP} & 6 RuBP
regenerated & PRK consumes ATP \\
\textbf{Net inputs} & \textbf{3 ATP + 2 NADPH} & \textbf{18 ATP + 12
NADPH} & \textbf{2 G3P (= 1 hexose)} & Excess 2 G3P leaves cycle \\
\end{longtable}
}

This 18:12 = 3:2 ATP:NADPH demand is the \emph{target} the light
reactions must hit. Pure LEF delivers ATP:NADPH about 1.3, so plants
tune the LEF/CEF balance to close the gap; under light limitation, CEF
is upregulated.

\begin{quote}
\textbf{Concept Check 2:} The Calvin cycle requires 3 ATP and 2 NADPH
per CO\(_2\) fixed. Linear electron flow produces ATP and NADPH in
approximately a 1.3:1 ratio. How does cyclic electron flow help satisfy
the Calvin cycle's ATP:NADPH ratio of 3:2 (= 1.5:1)?
\end{quote}

\begin{center}\rule{0.5\linewidth}{0.5pt}\end{center}

\subsection{The Calvin Cycle (Light-Independent
Reactions)}\label{the-calvin-cycle-light-independent-reactions}

The Calvin cycle fixes CO\(_2\) into organic compounds in the stroma. It
has three phases:

\begin{figure}[htbp]
\centering
\includegraphics[width=0.82\linewidth,height=4.2in,keepaspectratio]{../figures/mermaid_inline/inline_mermaid_0057_4ad799402a8f.png}
\caption{Flowchart for The Calvin Cycle (Light-Independent Reactions): Phase 1: Carboxylation, RuBP (5C), 2 x 3-PGA (3C), and Phase 2: Reduction form the diagram's primary path or branches.}
\end{figure}

\emph{The three phases of the Calvin cycle (Mermaid).} For every 3
CO\(_2\) fixed, 9 ATP and 6 NADPH are consumed, and 1 net G3P (3C) is
produced. Six turns of the cycle produce one glucose molecule.

\subsubsection{RuBisCO: Structure and Catalytic
Mechanism}\label{rubisco-structure-and-catalytic-mechanism}

\textbf{Structure:} Plant RuBisCO is a hexadecamer (L\(_8\)S\(_8\)) ---
eight large subunits (about 55 kDa, plastid-encoded \emph{rbcL}
gene) and eight small subunits (about 15 kDa, nucleus-encoded
\emph{RbcS} gene family). The L subunits contain the catalytic site at
the L-L dimer interface (so 8 active sites per holoenzyme); the S
subunits play structural and modulatory roles, contributing to the cap
of the assembly. Total mass is about 520 kDa, making RuBisCO
one of the largest stromal soluble enzymes.

\textbf{Activation requires CO\(_2\) + Mg\(^{2+}\):} A separate,
\emph{non-substrate} CO\(_2\) molecule first reacts with Lys201 of the
active site to form a \textbf{carbamate}, which then chelates a
Mg\(^{2+}\) ion. Primarily the carbamylated, Mg\(^{2+}\)-bound enzyme
can bind RuBP and proceed with catalysis. The mechanism cleverly
\emph{senses} CO\(_2\) availability twice: once to activate the enzyme
and once during turnover.

\textbf{RuBisCO activase (an AAA+ ATPase)} removes inhibitory sugar
phosphates (CA1P, RuBP) from RuBisCO's active site, allowing
carbamylation. This step is \textbf{light-dependent} (requires ATP) ---
connecting Calvin cycle activation to the energy state of the
chloroplast.

\textbf{Catalytic mechanism (in seven steps):}

\begin{enumerate}
\def\labelenumi{\arabic{enumi}.}
\tightlist
\item
  RuBP binds and is deprotonated at C3 to form a 2,3-enediol
  intermediate (the rate-limiting enolisation)
\item
  The enediol attacks CO\(_2\) (or O\(_2\) --- see photorespiration)
\item
  C2-C3 bond cleavage produces a six-carbon \texorpdfstring{\ensuremath{\beta}}{beta}-keto acid intermediate
\item
  Hydration of the C3 carbon
\item
  C2-C3 cleavage releases the first 3-PGA
\item
  Stereospecific protonation produces the second 3-PGA
\item
  Product release
\end{enumerate}

The slow \(k_{\text{cat}}\) (about 3--10 s\(^{-1}\)) is dominated by the
enolisation step (step 1) and by the stringent steric constraints needed
to \emph{select} CO\(_2\) over O\(_2\) in step 2. Faster RuBisCOs leak
more O\(_2\) into the active site, producing photorespiration.

\textbf{The CO\(_2\)/O\(_2\) specificity factor (S\(_{c/o}\), also \texorpdfstring{\ensuremath{\tau}}{tau}):}

\begin{equation}
S_{c/o} = \tau = \frac{V_c K_o}{V_o K_c}
\label{eq:unit_III_specificity_factor}
\end{equation}

where \(V_c\), \(V_o\) are maximal carboxylase/oxygenase rates and
\(K_c\), \(K_o\) are Michaelis constants. Higher S\(_{c/o}\) = more
selective for CO\(_2\) over O\(_2\).

{\def\LTcaptype{none} % do not increment counter
\begin{longtable}[]{@{}
  >{\raggedright\arraybackslash}p{(\linewidth - 6\tabcolsep) * \real{0.1429}}
  >{\raggedright\arraybackslash}p{(\linewidth - 6\tabcolsep) * \real{0.1607}}
  >{\raggedright\arraybackslash}p{(\linewidth - 6\tabcolsep) * \real{0.6250}}
  >{\raggedright\arraybackslash}p{(\linewidth - 6\tabcolsep) * \real{0.0714}}@{}}
\toprule\noalign{}
\begin{minipage}[b]{\linewidth}\raggedright
Organism
\end{minipage} & \begin{minipage}[b]{\linewidth}\raggedright
S\(_{c/o}\) (\texorpdfstring{\ensuremath{\tau}}{tau})
\end{minipage} & \begin{minipage}[b]{\linewidth}\raggedright
\(k_{\text{cat}}\) (carb., s\(^{-1}\))
\end{minipage} & \begin{minipage}[b]{\linewidth}\raggedright
Note
\end{minipage} \\
\midrule\noalign{}
\endhead
\bottomrule\noalign{}
\endlastfoot
Higher plants (spinach) & 80--100 & 3.3 & C3 \\
Cyanobacteria (\emph{Synechococcus}) & 40--50 & 12 & Compensates with
carboxysomes \\
C4 plants (maize) & 70--80 & 4.5 & Compensates with C4 anatomy \\
Red algae (\emph{Galdieria}) & about 240 & 2.6 & Highest known;
very slow \\
Theoretical maximum & about 400 & -- & Limited by enediol
mechanism \\
\end{longtable}
}

\textbf{The fundamental trade-off:} Across the tree of life, \(S_{c/o}\)
and \(k_{\text{cat}}\) are \emph{negatively correlated}. Faster RuBisCOs
are less selective. This appears to be a thermodynamic constraint on the
enediol intermediate's reactivity --- the same electronic features that
make it react with CO\(_2\) also make it react with O\(_2\).

\subsubsection{Carboxylation (Carbon
Fixation)}\label{carboxylation-carbon-fixation}

\[\text{RuBP (5C)} + \text{CO}_2 \xrightarrow{\text{RuBisCO}} 2 \times \text{3-PGA (3C)} \tag{10.5}\]

\textbf{RuBisCO} is the most abundant protein on Earth
(about 0.7 billion tonnes globally; about 50\% of leaf
protein). Despite its central importance, RuBisCO has remarkable
limitations:

\begin{itemize}
\tightlist
\item
  \textbf{Slow:} \(k_{\text{cat}} = 3\)--\(10\) s\(^{-1}\) (typical
  enzymes: 10\(^2\)--10\(^3\) s\(^{-1}\))
\item
  \textbf{Indiscriminate:} also reacts with O\(_2\) (\textbf{oxygenase
  activity}), leading to \textbf{photorespiration} and wasting
  about 25\% of fixed carbon in C3 plants
\item
  Plants compensate for RuBisCO's slowness by producing enormous
  quantities of it
\end{itemize}

\subsubsection{Reduction}\label{reduction}

This is the primary reductive step of the Calvin cycle: the ATP and
NADPH generated by the light reactions are spent here to convert 3-PGA
into the triose phosphate G3P. For every 6 CO\(_2\) fixed, 12 G3P are
produced --- 2 leave as the cycle's net carbohydrate yield and 10 are
recycled to regenerate RuBP.

\[\text{3-PGA} + \text{ATP} \rightarrow \text{1,3-BPG} + \text{ADP} \quad (\text{phosphoglycerate kinase}) \tag{10.6}\]

\[\text{1,3-BPG} + \text{NADPH} \rightarrow \text{G3P} + \text{NADP}^+ + \text{P}_i \quad (\text{G3P dehydrogenase}) \tag{10.7}\]

\subsubsection{Regeneration of RuBP}\label{regeneration-of-rubp}

Five G3P molecules (15 carbons) are rearranged through a complex 10-step
pathway (involving transketolase, aldolase,
sedoheptulose-1,7-bisphosphatase, ribulose-5-phosphate epimerase, and
ribulose-5-phosphate isomerase) to regenerate 3 RuBP (15 carbons),
consuming 3 ATP.

\subsubsection{Stoichiometry}\label{stoichiometry}

\textbf{To fix 6 CO\(_2\) (net gain 1 hexose):}

\begin{itemize}
\tightlist
\item
  6 CO\(_2\) \texorpdfstring{\ensuremath{\times}}{x} (3 ATP + 2 NADPH per CO\(_2\)) = \textbf{18 ATP + 12
  NADPH}
\item
  Light reactions must supply: 18 ATP and 12 NADPH
\item
  Requires about 48 photons (8 per CO\(_2\); 4 photons for 2
  electrons through PS II + PS I \texorpdfstring{\ensuremath{\times}}{x} 2 electron pairs per NADPH) for
  NADPH; additional photons for cyclic electron flow to meet ATP demand
\end{itemize}

\subsubsection{Worked Example: Energy Efficiency of
Photosynthesis}\label{worked-example-energy-efficiency-of-photosynthesis}

\emph{Problem:} Calculate the energy efficiency of photosynthesis given
that 48 photons of red light (680 nm, 176 kJ/mol each) are required to
fix 6 CO\(_2\) into 1 glucose (\(\Delta G^{\circ\prime}\) of glucose
combustion = 2,870 kJ/mol).

\emph{Solution:}

Total light energy input: \(48 \times 176 = 8{,}448\) kJ/mol

Energy stored in glucose: \(2{,}870\) kJ/mol

\begin{equation}
\eta = \frac{2{,}870}{8{,}448} \times 100\% = 34\%
\label{eq:unit_III_photosynthesis_efficiency}
\end{equation}

This is remarkably efficient for an energy conversion process.
Whole-canopy photosynthetic efficiency in a real crop is typically
around 1--3\%, because most of the year's incoming photosynthetically
active radiation (PAR) is lost to (i) reflection, (ii) photorespiration,
(iii) leaf shading and saturation, (iv) plant respiration, and (v)
suboptimal water/nutrient supply.

\textbf{Starch vs.~sucrose export:}

\begin{itemize}
\tightlist
\item
  \textbf{Starch} (alpha-glucose polymer) stored in chloroplast stroma
  during the day \texorpdfstring{\ensuremath{\to}}{->} remobilised at night to fuel respiration
\item
  \textbf{Sucrose} (glucose-fructose disaccharide) exported to phloem
  via triose phosphate/phosphate antiporter (TPT) \texorpdfstring{\ensuremath{\to}}{->} transported to
  non-photosynthetic tissues (\emph{source-sink} relationship)
\end{itemize}

\begin{quote}
\textbf{Concept Check 3:} RuBisCO is the most abundant enzyme on Earth,
yet it is extremely slow (\(k_{cat}\) = 3--10 s\(^{-1}\)). Why hasn't
evolution produced a faster version? Consider the trade-off between
speed and specificity (CO\(_2\) vs.~O\(_2\) discrimination) in light of
the specificity-factor data above.
\end{quote}

\begin{quote}
\textbf{Concept Check 3b:} Spinach RuBisCO has \(\tau \approx 90\).
\emph{Galdieria} (red alga) RuBisCO has \(\tau \approx 240\) but
\(k_{\text{cat}} \approx 2.6\) s\(^{-1}\) (vs.~3.3 for spinach). At
ambient \([\text{CO}_2]/[\text{O}_2]\), which enzyme would maximize
CO\(_2\) flux per unit enzyme? When (in what environment) does the
trade-off favour the \emph{Galdieria} form?
\end{quote}

\begin{center}\rule{0.5\linewidth}{0.5pt}\end{center}

\subsection{Photorespiration}\label{photorespiration}

RuBisCO's oxygenase activity produces \textbf{2-phosphoglycolate} (2C),
which must be recycled through the photorespiratory pathway (C2 cycle):

\[\text{RuBP} + \text{O}_2 \xrightarrow{\text{RuBisCO oxygenase}} \text{3-PGA (3C)} + \text{2-phosphoglycolate (2C)} \tag{10.9}\]

The photorespiratory pathway spans three
\hyperref[gl:organelle]{\textbf{organelle}}s: 1. \textbf{Chloroplast:}
2-phosphoglycolate \texorpdfstring{\ensuremath{\to}}{->} glycolate (phosphatase) 2. \textbf{Peroxisome:}
glycolate \texorpdfstring{\ensuremath{\to}}{->} glyoxylate \texorpdfstring{\ensuremath{\to}}{->} glycine (generates H\(_2\)O\(_2\), destroyed by
catalase) 3. \textbf{Mitochondrion:} 2 glycine \texorpdfstring{\ensuremath{\to}}{->} serine + CO\(_2\) +
NH\(_3\) + NADH (glycine decarboxylase) 4. \textbf{Peroxisome:} serine \texorpdfstring{\ensuremath{\to}}{->}
hydroxypyruvate \texorpdfstring{\ensuremath{\to}}{->} glycerate 5. \textbf{Chloroplast:} glycerate \texorpdfstring{\ensuremath{\to}}{->} 3-PGA
(phosphorylation by ATP)

\textbf{Cost of photorespiration:}

\begin{itemize}
\tightlist
\item
  Loses 1 CO\(_2\) per 2 oxygenation events (25\% of fixed carbon)
\item
  Consumes ATP for recovery
\item
  Releases NH\(_3\) (must be re-assimilated)
\item
  At 25 degrees C and current atmospheric CO\(_2\), photorespiration
  reduces net photosynthesis by about 20--30\% in C3 plants
\end{itemize}

\begin{quote}
\textbf{Clinical Connection: Photorespiration and Global Food Security}
Reducing photorespiration in crop plants could increase yields by
20--40\%. Approaches include: - \textbf{Engineering a chloroplastic
glycolate bypass} (South et al., 2019, \emph{Science}): reduced
photorespiratory losses by 40\% and increased tobacco biomass -
\textbf{Introducing CO\(_2\)-concentrating mechanisms} (CCMs) from
cyanobacteria into C3 crops - \textbf{Engineering RuBisCO} with improved
CO\(_2\)/O\(_2\) specificity (challenging due to the speed-specificity
trade-off) These efforts matter because the UN World Population
Prospects 2024 projects roughly 9.7 billion people by 2050, while
climate stress, land limits, food loss, and inequitable access constrain
the food system \citep{un2024population}. Higher photosynthetic
efficiency can raise yield potential, but it is not a stand-alone
solution to nutrition, distribution, soil health, or water scarcity.
\end{quote}

\begin{center}\rule{0.5\linewidth}{0.5pt}\end{center}

\subsection{C3, C4, and CAM
Photosynthesis}\label{c3-c4-and-cam-photosynthesis}

\subsubsection{\texorpdfstring{Why CO\(_2\)-Concentrating Mechanisms
Evolved}{Why CO_2-Concentrating Mechanisms Evolved}}\label{why-co_2-concentrating-mechanisms-evolved}

The atmospheric CO\(_2\):O\(_2\) ratio (about 0.04\%/21\% \texorpdfstring{\ensuremath{\approx}}{~}
1:525) is hostile to RuBisCO. The fraction of carboxylation versus
oxygenation reactions is approximately:

\begin{equation}
\frac{v_c}{v_o} = S_{c/o} \cdot \frac{[\text{CO}_2]}{[\text{O}_2]}
\label{eq:unit_III_carb_oxy_ratio}
\end{equation}

For C3 plants at 25 °C, typical leaf-internal CO\(_2\) \texorpdfstring{\ensuremath{\approx}}{~} 7 \texorpdfstring{\ensuremath{\mu}}{mu}M and
O\(_2\) \texorpdfstring{\ensuremath{\approx}}{~} 250 \texorpdfstring{\ensuremath{\mu}}{mu}M, so \(v_c/v_o \texorpdfstring{\ensuremath{\approx}}{~} 80 \times (7/250) \approx 2.2\) ---
meaning about 30\% of RuBisCO turnovers are oxygenations. C4
and CAM plants have evolved spatial or temporal CO\(_2\)-concentrating
mechanisms (CCMs) that raise the local {[}CO\(_2\){]} around RuBisCO
about 10\texorpdfstring{\ensuremath{\times}}{x} above ambient, suppressing photorespiration.

\subsubsection{C3 Photosynthesis (First stable product:
3-PGA)}\label{c3-photosynthesis-first-stable-product-3-pga}

Standard Calvin cycle. Used by about 85\% of plant species
(wheat, rice, soybeans, most trees). Problem: \textbf{photorespiration}
is severe at high temperature and low CO\(_2\) (e.g., midday in summer,
when \hyperref[gl:stomata]{\textbf{stomata}} close to conserve water).

\subsubsection{C4 Photosynthesis (First stable product: oxaloacetate,
4C)}\label{c4-photosynthesis-first-stable-product-oxaloacetate-4c}

\begin{figure}[htbp]
\centering
\includegraphics[width=0.82\linewidth,height=4.2in,keepaspectratio]{../figures/mermaid_inline/inline_mermaid_0058_5c955f9bd359.png}
\caption{Flowchart for C4 Photosynthesis (First stable product: oxaloacetate, 4C): Mesophyll Cell, CO2 (atmospheric), Oxaloacetate (4C), and Malate (4C) form the diagram's primary path or branches.}
\end{figure}

\emph{C4 photosynthesis (Mermaid).} Spatial separation of initial
CO\(_2\) fixation (mesophyll cells, PEP carboxylase) and the Calvin
cycle (bundle sheath cells, RuBisCO). This CO\(_2\)-concentrating
mechanism virtually eliminates photorespiration.

\textbf{CO\(_2\)-concentrating mechanism:} CO\(_2\) is captured in
\textbf{mesophyll cells} by PEP carboxylase (no O\(_2\) reactivity;
\(K_m\) for CO\(_2\) about  2 \texorpdfstring{\ensuremath{\mu}}{mu}M, vs.~RuBisCO \(K_m\)
about 10 \texorpdfstring{\ensuremath{\mu}}{mu}M) as oxaloacetate \texorpdfstring{\ensuremath{\to}}{->} malate \texorpdfstring{\ensuremath{\to}}{->} transported to
\textbf{bundle sheath cells} via plasmodesmata \texorpdfstring{\ensuremath{\to}}{->} decarboxylated \texorpdfstring{\ensuremath{\to}}{->}
CO\(_2\) concentrated (10x above ambient) \texorpdfstring{\ensuremath{\to}}{->} RuBisCO operates at
saturating CO\(_2\) \texorpdfstring{\ensuremath{\to}}{->} negligible photorespiration.

C4 plants (maize, sugarcane, sorghum, millet) require \textbf{2
additional ATP per CO\(_2\)} fixed (total: 5 ATP + 2 NADPH per CO\(_2\))
but are more efficient at high temperature and high light. Spatial
separation of carboxylation (mesophyll) and Calvin cycle (bundle sheath)
requires specialised leaf anatomy = \textbf{Kranz anatomy} (German:
``wreath'').

\textbf{C4 subtypes} differ in the decarboxylation enzyme in bundle
sheath cells:

\begin{itemize}
\tightlist
\item
  \textbf{NADP-ME type:} NADP-malic enzyme (maize, sorghum, sugarcane)
\item
  \textbf{NAD-ME type:} NAD-malic enzyme (millet, Amaranthus)
\item
  \textbf{PEP-CK type:} PEP carboxykinase (guinea grass, Panicum)
\end{itemize}

\subsubsection{CAM Photosynthesis (Temporal
separation)}\label{cam-photosynthesis-temporal-separation}

\textbf{Crassulacean Acid Metabolism:} stomata open at \textbf{night}
(cool, high humidity \texorpdfstring{\ensuremath{\to}}{->} low
\hyperref[gl:transpiration]{\textbf{transpiration}}) \texorpdfstring{\ensuremath{\to}}{->} CO\(_2\) fixed by
PEP carboxylase \texorpdfstring{\ensuremath{\to}}{->} malate stored in vacuole (vacuolar pH drops from
about 7 to about 4 as malic acid accumulates). Stomata
\textbf{close during the day} (prevent water loss) \texorpdfstring{\ensuremath{\to}}{->} malate
decarboxylated \texorpdfstring{\ensuremath{\to}}{->} CO\(_2\) concentrated around RuBisCO \texorpdfstring{\ensuremath{\to}}{->} Calvin cycle
proceeds.

CAM plants (cacti, agaves, pineapple, orchids, jade plant) have
extremely high water-use efficiency (WUE: about 3--5x higher
than C3). The tradeoff: slow growth due to nighttime carbon storage
limitation and vacuolar capacity.

\subsubsection{Water-Use Efficiency}\label{water-use-efficiency}

Water-use efficiency (WUE) is defined as carbon gained per unit water
lost:

\begin{equation}
\text{WUE} = \frac{A}{E} = \frac{\text{net CO}_2 \text{ assimilation}}{\text{transpiration rate}}
\label{eq:unit_III_water_use_efficiency}
\end{equation}

The intrinsic WUE depends on the leaf-internal:atmospheric CO\(_2\)
gradient:

\begin{equation}
\text{WUE}_{\text{intrinsic}} = \frac{c_a - c_i}{1.6 \cdot (e_i - e_a)}
\label{eq:unit_III_intrinsic_wue}
\end{equation}

where \(c_a\) and \(c_i\) are atmospheric and intercellular CO\(_2\)
concentrations, \(e_a\) and \(e_i\) are ambient and intercellular water
vapour pressures, and 1.6 is the ratio of water-vapour to CO\(_2\)
diffusivity in air.

\textbf{Typical WUE values (mol CO\(_2\) fixed / mol H\(_2\)O
transpired):}

{\def\LTcaptype{none} % do not increment counter
\begin{longtable}[]{@{}
  >{\raggedright\arraybackslash}p{(\linewidth - 4\tabcolsep) * \real{0.2857}}
  >{\raggedright\arraybackslash}p{(\linewidth - 4\tabcolsep) * \real{0.4571}}
  >{\raggedright\arraybackslash}p{(\linewidth - 4\tabcolsep) * \real{0.2571}}@{}}
\toprule\noalign{}
\begin{minipage}[b]{\linewidth}\raggedright
Plant type
\end{minipage} & \begin{minipage}[b]{\linewidth}\raggedright
WUE (\texorpdfstring{\ensuremath{\times}}{x}10\(^{-3}\))
\end{minipage} & \begin{minipage}[b]{\linewidth}\raggedright
Mechanism
\end{minipage} \\
\midrule\noalign{}
\endhead
\bottomrule\noalign{}
\endlastfoot
C3 & 1--3 & Standard Calvin cycle; high transpiration \\
C4 & 3--6 & CCM allows lower \(c_i\) \texorpdfstring{\ensuremath{\to}}{->} smaller stomatal aperture \\
CAM & 10--40 & Nocturnal stomatal opening drastically reduces
\(e_i - e_a\) \\
\end{longtable}
}

CAM is the extreme: by opening stomata at night when air is cool and
humid (small \(e_i - e_a\)), CAM plants lose about 10\texorpdfstring{\ensuremath{\times}}{x} less
water per CO\(_2\) fixed than C3 plants do midday --- a huge advantage
in deserts but at the cost of slow growth.

\subsubsection{Comparison}\label{comparison}

{\def\LTcaptype{none} % do not increment counter
\begin{longtable}[]{@{}
  >{\raggedright\arraybackslash}p{(\linewidth - 6\tabcolsep) * \real{0.5000}}
  >{\raggedright\arraybackslash}p{(\linewidth - 6\tabcolsep) * \real{0.1429}}
  >{\raggedright\arraybackslash}p{(\linewidth - 6\tabcolsep) * \real{0.1429}}
  >{\raggedright\arraybackslash}p{(\linewidth - 6\tabcolsep) * \real{0.2143}}@{}}
\toprule\noalign{}
\begin{minipage}[b]{\linewidth}\raggedright
Feature
\end{minipage} & \begin{minipage}[b]{\linewidth}\raggedright
C3
\end{minipage} & \begin{minipage}[b]{\linewidth}\raggedright
C4
\end{minipage} & \begin{minipage}[b]{\linewidth}\raggedright
CAM
\end{minipage} \\
\midrule\noalign{}
\endhead
\bottomrule\noalign{}
\endlastfoot
First stable product & 3-PGA (3C) & Oxaloacetate (4C) & Oxaloacetate
(4C) \\
CO\(_2\) fixation enzyme & RuBisCO & PEP carboxylase (initial) & PEP
carboxylase (night) \\
Photorespiration & High (20--30\% loss) & Near zero & Near zero \\
Stomata pattern & Open day & Open day & Open night \\
Optimal temperature & 15--25 degrees C & 30--40 degrees C & Variable \\
Water use efficiency & Moderate & High & Very high \\
ATP cost per CO\(_2\) & 3 & 5 & 5.5--6.5 \\
Examples & Wheat, rice, trees & Maize, sugarcane, millet & Cactus,
agave, pineapple \\
\% of plant species & about 85\% & about 3\%
(about 7,500 species) & about 6--8\% \\
Leaf anatomy & Standard mesophyll & Kranz anatomy & Thick, succulent \\
\end{longtable}
}

\begin{quote}
\textbf{Concept Check 4:} Why is C4 photosynthesis advantageous at high
temperatures but not at low temperatures? Consider the temperature
dependence of RuBisCO oxygenase activity and the ATP cost of the C4
pump.
\end{quote}

\begin{quote}
\textbf{Concept Check 4b:} Bermuda grass (C4) and tall fescue (C3)
compete in temperate lawns. In a cool, wet spring, fescue dominates; in
a hot, dry summer, Bermuda dominates. Use
\cref{eq:unit_III_carb_oxy_ratio} and \cref{eq:unit_III_intrinsic_wue}
to explain this seasonal turnover quantitatively.
\end{quote}

\begin{center}\rule{0.5\linewidth}{0.5pt}\end{center}

\subsection{Regulation of the Calvin
Cycle}\label{regulation-of-the-calvin-cycle}

The Calvin cycle does not run in the dark --- several enzymes are
regulated by light-dependent mechanisms to prevent futile cycling with
\hyperref[gl:glycolysis]{\textbf{glycolysis}} and the oxidative pentose
phosphate pathway.

\subsubsection{The Thioredoxin/Ferredoxin Redox Regulation
System}\label{the-thioredoxinferredoxin-redox-regulation-system}

Light reduces the chloroplast thioredoxin pool through a precisely
engineered cascade:

\begin{equation}
\text{PS I} \to \text{Fd}_{\text{red}} \to \text{FTR} \to \text{Trx}_{\text{red}} \to \text{target enzyme (S–S} \to \text{2 SH)}
\label{eq:unit_III_thioredoxin_cascade}
\end{equation}

In detail:

\begin{enumerate}
\def\labelenumi{\arabic{enumi}.}
\tightlist
\item
  PS I reduces ferredoxin (Fd, \(E^{\circ\prime} = -0.42\) V).
\item
  Fd-thioredoxin reductase (FTR) --- a unique {[}4Fe-4S{]} enzyme ---
  transfers two electrons to thioredoxin, converting an internal
  disulfide to two free thiols.
\item
  Reduced thioredoxin (Trx-(SH)\(_2\)) reduces specific regulatory
  disulfides on target enzymes.
\item
  Reduction \emph{activates} the target enzyme (in most cases) by
  relieving an inhibitory conformational lock.
\end{enumerate}

\begin{figure}[htbp]
\centering
\includegraphics[width=0.82\linewidth,height=4.2in,keepaspectratio]{../figures/mermaid_inline/inline_mermaid_0059_86f89afd405c.png}
\caption{Flowchart for The Thioredoxin/Ferredoxin Redox Regulation System: Light, PS I, Ferredoxin (red), and Fd-Thioredoxin Reductase form the diagram's primary path or branches.}
\end{figure}

\emph{The thioredoxin/ferredoxin system as the master light-dependent
redox switch (Mermaid).} Calvin-cycle enzymes carry regulatory
disulfides that are reduced (activated) in the light and re-oxidised
(inactivated) in the dark. This prevents futile ATP/NADPH consumption
during darkness.

Thioredoxin-controlled Calvin-cycle targets include:

\begin{itemize}
\tightlist
\item
  \textbf{GAPDH (NADP-linked):} activated directly; catalyses the second
  reduction-phase step.
\item
  \textbf{Fructose-1,6-bisphosphatase (FBPase):} activated directly;
  supports carbon regeneration.
\item
  \textbf{Sedoheptulose-1,7-bisphosphatase (SBPase):} activated
  directly; a major regeneration-phase flux-control enzyme.
\item
  \textbf{PRK (phosphoribulokinase):} activated directly; regenerates
  RuBP.
\item
  \textbf{RuBisCO activase:} regulated indirectly; maintains RuBisCO in
  an active, carbamylated state.
\item
  \textbf{CP12:} reduction releases the GAPDH/PRK complex, reversing
  dark-state inactivation.
\end{itemize}

In the dark, the target disulfides re-oxidise (catalysed by 2-Cys
peroxiredoxins coupled to H\(_2\)O\(_2\)), inactivating the enzymes.
This prevents the Calvin cycle from consuming ATP and NADPH when they
are not being regenerated by the light reactions.

\subsubsection{\texorpdfstring{CO\(_2\) Concentration
Effects}{CO_2 Concentration Effects}}\label{co_2-concentration-effects}

The flux through the Calvin cycle is also acutely sensitive to CO\(_2\)
supply:

\begin{itemize}
\tightlist
\item
  At low {[}CO\(_2\){]} (e.g., closed stomata, drought), RuBisCO is the
  rate-limiting step and the cycle is ``carboxylation-limited.''
\item
  At intermediate {[}CO\(_2\){]}, the cycle becomes ``RuBP
  regeneration-limited'' --- flux limited by the supply of RuBP from the
  regeneration phase, which in turn depends on the rate of
  light-dependent ATP/NADPH supply.
\item
  At very high {[}CO\(_2\){]} (e.g., \textgreater{} ambient \texorpdfstring{\ensuremath{\times}}{x} 2), the
  cycle can become ``P\(_i\)-limited'' --- phosphate becomes scarce as
  triose phosphate is exported.
\end{itemize}

This three-zone behaviour underlies the classic Farquhar-von
Caemmerer-Berry (FvCB) photosynthesis model, the workhorse of crop and
ecosystem modelling.

\subsubsection{\texorpdfstring{Stromal pH and Mg\(^{2+}\) in the
Light}{Stromal pH and Mg^(2+) in the Light}}\label{stromal-ph-and-mg2-in-the-light}

In the light, H\(^+\) pumping from stroma to lumen has two activating
effects:

\begin{itemize}
\tightlist
\item
  Stromal pH rises from about 7.0 to about 8.0
  (RuBisCO and FBPase pH optimum is about 8.0--8.5)
\item
  Stromal {[}Mg\(^{2+}\){]} rises by about 3 mM as Mg\(^{2+}\)
  exits the lumen with the H\(^+\) flux. Mg\(^{2+}\) is required for
  RuBisCO carbamylation and FBPase activity.
\end{itemize}

This is a beautifully integrated example of physical chemistry: the
\emph{same} H\(^+\) transport that builds the PMF for ATP synthesis
simultaneously activates the carbon-fixation enzymes that consume that
ATP.

\subsubsection{CP12 Regulation}\label{cp12-regulation}

\textbf{CP12} is a small regulatory protein that forms a ternary complex
with GAPDH and PRK in the dark, inactivating both enzymes
simultaneously. Light-dependent reduction of CP12 by thioredoxin
releases GAPDH and PRK, activating the Calvin cycle.

\begin{quote}
\textbf{Concept Check 5:} If a mutation in ferredoxin-thioredoxin
reductase (FTR) prevented thioredoxin from being reduced, what would
happen to Calvin cycle activity in the light? Would the light reactions
be affected?
\end{quote}

\begin{center}\rule{0.5\linewidth}{0.5pt}\end{center}

\subsection{Photoprotection: Non-Photochemical
Quenching}\label{photoprotection-non-photochemical-quenching}

Excess light energy (when absorbed photons exceed the rate at which
electrons can be productively used) can damage the photosynthetic
apparatus through generation of reactive oxygen species, particularly
singlet oxygen (\(^1\)O\(_2\)) from triplet chlorophyll (\(^3\)Chl*).
Plants must dissipate excess excitation safely.

\subsubsection{The Three Components of
NPQ}\label{the-three-components-of-npq}

NPQ is operationally defined as fluorescence quenching that does not
arise from photochemistry:

\begin{equation}
\text{NPQ} = \frac{F_m - F_m'}{F_m'}
\label{eq:unit_III_npq_definition}
\end{equation}

where \(F_m\) is maximum fluorescence in the dark-adapted state and
\(F_m'\) is the maximum fluorescence in the light-adapted state. NPQ has
three kinetic components:

{\def\LTcaptype{none} % do not increment counter
\begin{longtable}[]{@{}
  >{\raggedright\arraybackslash}p{(\linewidth - 4\tabcolsep) * \real{0.3333}}
  >{\raggedright\arraybackslash}p{(\linewidth - 4\tabcolsep) * \real{0.3333}}
  >{\raggedright\arraybackslash}p{(\linewidth - 4\tabcolsep) * \real{0.3333}}@{}}
\toprule\noalign{}
\begin{minipage}[b]{\linewidth}\raggedright
Component
\end{minipage} & \begin{minipage}[b]{\linewidth}\raggedright
Timescale
\end{minipage} & \begin{minipage}[b]{\linewidth}\raggedright
Mechanism
\end{minipage} \\
\midrule\noalign{}
\endhead
\bottomrule\noalign{}
\endlastfoot
\textbf{qE (energy quenching)} & seconds & Low lumen pH; PsbS
protonation; xanthophyll cycle \\
\textbf{qT (state transitions)} & minutes & LHCII phosphorylation;
antenna redistribution between PS II and PS I \\
\textbf{qI (photoinhibition)} & hours & D1 damage and repair \\
\end{longtable}
}

\subsubsection{qE: The Xanthophyll Cycle in
Detail}\label{qe-the-xanthophyll-cycle-in-detail}

When light is in excess, the lumen pH falls below about 5.5,
triggering two parallel events:

\begin{enumerate}
\def\labelenumi{\arabic{enumi}.}
\item
  \textbf{PsbS protonation:} PsbS is a small (about 22 kDa)
  intrinsic membrane protein with two lumen-exposed glutamate residues
  that protonate at low pH. Protonation triggers a conformational change
  that brings antenna pigments into close contact, promoting energy
  transfer to a quenching site (a zeaxanthin-Chl heterodimer).
\item
  \textbf{Violaxanthin \texorpdfstring{\ensuremath{\to}}{->} zeaxanthin:} \textbf{Violaxanthin de-epoxidase
  (VDE)}, activated by low pH, catalyses sequential de-epoxidation
  \citep{hager1971}:
\end{enumerate}

\begin{equation}
\text{Violaxanthin (V)} \xrightarrow[+\text{ascorbate}]{\text{VDE, low pH}} \text{Antheraxanthin (A)} \xrightarrow[+\text{ascorbate}]{\text{VDE}} \text{Zeaxanthin (Z)}
\label{eq:unit_III_xanthophyll_cycle}
\end{equation}

Zeaxanthin (without epoxide groups) has a lower-energy first excited
state than chlorophyll, so it can act as an energy sink --- accepting
excitation from \(^1\)Chl* and dissipating it as heat. The reverse
reaction (Z \texorpdfstring{\ensuremath{\to}}{->} V) is catalysed by zeaxanthin epoxidase (ZE) in the dark
or low-light conditions.

\begin{figure}[htbp]
\centering
\includegraphics[width=0.82\linewidth,height=4.2in,keepaspectratio]{../figures/mermaid_inline/inline_mermaid_0060_987a5e8fe209.png}
\caption{Sequence diagram for qE: The Xanthophyll Cycle in Detail showing ordered interaction among Excess light, Lumen pH decreased (about 5), PsbS (2 Glu protonated), and Violaxanthin de-epoxidase.}
\end{figure}

\emph{The xanthophyll cycle and qE quenching cascade (Mermaid).} Lumen
acidification (signal of excess light) protonates PsbS and activates
VDE. Zeaxanthin appears at the LHCII quenching site within seconds and
dissipates excess excitation as heat. When light decreases, lumen pH
rises and the whole system reverses over minutes.

\subsubsection{Carotenoid Quenching of Triplet
States}\label{carotenoid-quenching-of-triplet-states}

Carotenoids (\texorpdfstring{\ensuremath{\beta}}{beta}-carotene, zeaxanthin) directly quench triplet chlorophyll
(\(^3\)Chl*) and singlet oxygen (\(^1\)O\(_2\)), converting the energy
to heat. Without carotenoids, photosynthesis is lethal --- mutants
lacking carotenoid biosynthesis (e.g., norflurazon-treated plants)
bleach and die in light.

The carotenoid triplet state (about 80 kJ/mol) is below both
\(^3\)Chl* (about 110 kJ/mol) and \(^1\)O\(_2\)
(about 94 kJ/mol), so energy transfer is downhill in both
cases.

\subsubsection{State Transitions (qT)}\label{state-transitions-qt}

When PS II is over-excited relative to PS I (e.g., under blue light,
which preferentially excites PS II), the plastoquinone pool becomes
reduced. Reduced PQ activates \textbf{STN7 kinase}, which phosphorylates
LHCII. Phosphorylated LHCII detaches from PS II in grana and migrates to
PS I in stroma lamellae, rebalancing excitation pressure. Under far-red
light (preferentially exciting PS I), PQ becomes oxidised, STN7 is
inactive, and a phosphatase (PPH1/TAP38) returns LHCII to PS II.

\subsubsection{Other Photoprotective
Mechanisms}\label{other-photoprotective-mechanisms}

\begin{enumerate}
\def\labelenumi{\arabic{enumi}.}
\item
  \textbf{Chloroplast movements:} In high light, chloroplasts align
  parallel to the light direction (low cross-section); in low light,
  they spread perpendicular (high cross-section). Mediated by
  phototropins.
\item
  \textbf{Reactive oxygen scavenging:} Superoxide dismutase (SOD),
  ascorbate peroxidase (APX), and the glutathione-ascorbate
  (Foyer-Halliwell) cycle detoxify ROS generated by PS I (Mehler
  reaction).
\item
  \textbf{Photorespiration as a safety valve:} Counter-intuitively,
  photorespiration consumes ATP and reductant, draining excess electron
  flow when CO\(_2\) is limiting --- a ``release valve'' that reduces
  ROS pressure on the chloroplast.
\end{enumerate}

\begin{quote}
\textbf{Concept Check 6:} Explain why a plant mutant lacking carotenoids
would die in the light but survive in the dark. What specific
photodamage mechanism would cause cell death?
\end{quote}

\begin{quote}
\textbf{Concept Check 6b:} A \emph{psbs} knockout \emph{Arabidopsis}
lacks the PsbS protein and shows almost no qE, although LHCII and the
xanthophyll cycle are intact. Predict the plant's growth phenotype under
(a) constant moderate light and (b) fluctuating light (sun \texorpdfstring{\ensuremath{\to}}{->} shade \texorpdfstring{\ensuremath{\to}}{->} sun
every minute). Why is the fluctuating-light phenotype more severe?
\end{quote}

\begin{center}\rule{0.5\linewidth}{0.5pt}\end{center}

\subsection{Evolutionary Origins of
Photosynthesis}\label{evolutionary-origins-of-photosynthesis}

Photosynthesis has a complex evolutionary history spanning
about 3.5 billion years:

\begin{enumerate}
\def\labelenumi{\arabic{enumi}.}
\item
  \textbf{Anoxygenic photosynthesis} (about 3.5 Ga): Early
  photosynthetic bacteria (e.g., purple bacteria, green sulfur bacteria)
  used a single reaction centre (Type I or Type II) and electron donors
  other than water (H\(_2\)S, Fe\(^{2+}\), H\(_2\)). No O\(_2\)
  produced.
\item
  \textbf{Origin of oxygenic photosynthesis} (about 2.7--3.0
  Ga): Cyanobacteria evolved the ability to link two photosystems (PS II
  and PS I) in series (Z-scheme), enabling water as an electron donor.
  This was the most transformative evolutionary innovation on Earth.
\item
  \textbf{Great Oxidation Event} (about 2.4 Ga): Cyanobacterial
  O\(_2\) production overwhelmed geological sinks (reduced iron,
  sulfur), causing atmospheric O\(_2\) to rise from \textless0.001\% to
  about 2\%. This was catastrophic for obligate anaerobes (the
  ``Oxygen Holocaust'') but enabled
  \hyperref[gl:aerobic]{\textbf{aerobic}} respiration
  (about 18x more efficient ATP production than
  \hyperref[gl:fermentation]{\textbf{fermentation}}).
\item
  \textbf{Primary endosymbiosis} (about 1.5 Ga): A
  cyanobacterium was engulfed by a eukaryotic ancestor, becoming the
  chloroplast. Evidence: double membrane, 70S ribosomes, circular DNA,
  sensitivity to chloramphenicol.
\item
  \textbf{Secondary and tertiary endosymbiosis:} Algae engulfed other
  algae, creating organelles with 3--4 membranes (e.g., chloroplasts of
  brown algae, diatoms, euglenoids). This accounts for the remarkable
  diversity of photosynthetic pigments and chloroplast structures across
  algal lineages.
\end{enumerate}

\begin{quote}
\textbf{Concept Check 7:} The reaction centres of PS II (D1/D2 proteins)
and purple bacteria (L/M subunits) share significant structural and
sequence homology, yet PS II can oxidise water and purple bacterial
reaction centres cannot. What structural innovation in PS II enabled
water oxidation?
\end{quote}

\begin{center}\rule{0.5\linewidth}{0.5pt}\end{center}

\subsection{Artificial Photosynthesis and Solar
Fuels}\label{artificial-photosynthesis-and-solar-fuels}

Inspired by natural photosynthesis, researchers are developing
artificial systems to convert solar energy into fuels. The diversity of
approaches reveals a tension between \emph{efficiency} (where
photovoltaics excel) and \emph{fuel storage} (where chemistry excels).

\subsubsection{Comparison: Natural vs.~Artificial Light
Conversion}\label{comparison-natural-vs.-artificial-light-conversion}

\textbf{Biological systems.}

\begin{itemize}
\tightlist
\item
  \textbf{C3 photosynthesis (whole plant):} 0.5--1.5\%
  solar-to-stored-carbon efficiency, producing glucose and starch. Its
  advantage is complete self-assembly and repair; its limit is low
  energy efficiency and dependence on water and mineral nutrients.
\item
  \textbf{C4 photosynthesis (for example, sugarcane):} 1.5--3.5\%
  solar-to-stored-carbon efficiency, producing glucose and sucrose. The
  CO\(_2\)-concentrating mechanism improves performance in hot, bright
  environments, but the anatomical and ATP costs restrict its ecological
  range.
\item
  \textbf{Theoretical biological maximum:} about 12\% for reduced carbon
  under idealised constraints. The main losses come from spectral
  mismatch, RuBisCO kinetics, photorespiration, respiration, and the
  ATP/NADPH cost of carbon fixation.
\end{itemize}

\textbf{Engineered energy systems.}

\begin{itemize}
\tightlist
\item
  \textbf{Silicon photovoltaic (single-junction):} about 22\% commercial
  electricity conversion. It is mature and inexpensive, but electricity
  is not a fuel and requires downstream storage.
\item
  \textbf{III--V tandem photovoltaics:} about 33\% in high-performance
  laboratory cells. They maximise single-cell photovoltaic efficiency
  but remain costly and materials-intensive.
\item
  \textbf{Photovoltaics plus electrolysis:} about 10--14\%
  solar-to-H\(_2\) efficiency. Hydrogen is a storable fuel, but
  compression, transport, leakage, and end-use infrastructure remain
  major constraints.
\end{itemize}

\textbf{Hybrid and chemical fuel systems.}

\begin{itemize}
\tightlist
\item
  \textbf{Photoelectrochemical cells:} about 4--10\% solar-to-H\(_2\)
  efficiency from water. Direct light-to-fuel integration is attractive,
  but catalyst durability and device stability are limiting.
\item
  \textbf{CO\(_2\) photoreduction with copper-family catalysts:} roughly
  1--5\% for mixtures such as CO, formate, and ethanol. The appeal is
  direct carbon capture into fuels; the challenge is selectivity.
\item
  \textbf{Bio-hybrids with engineered cyanobacteria:} about 0.5--4\% for
  products such as H\(_2\), ethanol, or alkanes. They are
  self-replicating and tunable but slow and fragile.
\item
  \textbf{Synthetic Mn\(_4\)Ca-inspired catalysts:} research-stage
  water-oxidation systems producing O\(_2\) from H\(_2\)O. They teach
  design principles from the oxygen-evolving complex but are not yet
  practical devices.
\end{itemize}

\subsubsection{Why Plants Lag PV (and Why It May Not
Matter)}\label{why-plants-lag-pv-and-why-it-may-not-matter}

A silicon solar cell at 22\% beats C3 photosynthesis at 1\% by an order
of magnitude --- but the comparison is misleading. PV converts photons
to electrons; plants convert photons to \emph{covalently stored carbon}.
To match plants, a PV system must add electrolysis, gas storage, and
CO\(_2\) reduction --- each step lossy. End-to-end PV \texorpdfstring{\ensuremath{\to}}{->} renewable
methanol systems currently achieve about 5--8\% solar-to-fuel
--- comparable to C4 photosynthesis. The biological system also
self-assembles, self-repairs, and runs at \textless1 sun without exotic
materials.

\textbf{Theoretical efficiency limits.} The ultimate biological limit is
set by the sum of (a) photon-to-electron efficiency in the antenna
(about 95\%), (b) thermalisation losses (longer-wavelength
photons supply about 176 kJ/mol; the enthalpy of a pair of e\(^-\)
across the Z-scheme is about 440 kJ --- so the lower bound is
about 1.14 V \texorpdfstring{\ensuremath{\times}}{x} 4 e\(^-\) / 8 photons \texorpdfstring{\ensuremath{\times}}{x} 176 kJ/mol =
about 31\% before downstream losses), and (c) Calvin cycle and
respiration costs (about 40\% loss). The product is a
theoretical about 12\% solar-to-biomass efficiency.

\subsubsection{Key Research Directions}\label{key-research-directions}

\begin{enumerate}
\def\labelenumi{\arabic{enumi}.}
\item
  \textbf{Artificial water splitting:} Metal oxide catalysts (e.g.,
  IrO\(_2\), Co-Pi) that oxidise water, coupled to photovoltaic cells.
  The goal: electrochemical production of H\(_2\) and O\(_2\) from water
  using sunlight.
\item
  \textbf{CO\(_2\) reduction:} Electrocatalytic or photocatalytic
  reduction of CO\(_2\) to formate, methanol, or hydrocarbons.
  Challenges include selectivity (CO\(_2\) vs.~H\(^+\) reduction) and
  catalyst durability.
\item
  \textbf{Bio-hybrid systems:} Living photosynthetic organisms
  (cyanobacteria, algae) engineered to produce hydrogen, ethanol, or
  alkanes directly from CO\(_2\) and sunlight.
\item
  \textbf{Synthetic minimal photosystem:} Recreating the core PS II
  Mn\(_4\)Ca water-splitting catalyst in synthetic molecules. The
  Mn\(_4\)CaO\(_5\) cluster structure (revealed by Shen et al., 2011,
  \emph{Nature}, at 1.9 Å resolution) guides biomimetic catalyst design.
\end{enumerate}

\begin{quote}
\textbf{Clinical Connection: Photodynamic Therapy (PDT)} Photodynamic
therapy exploits the same photochemistry that plants must protect
against. A photosensitiser (e.g., porphyrin derivative) is administered
to a patient and preferentially accumulates in tumour tissue. Upon
illumination with specific wavelengths, the photosensitiser generates
singlet oxygen (\(^1\)O\(_2\)), which destroys tumour cells.
Understanding the photophysics of chlorophyll-like molecules (energy
transfer, intersystem crossing, ROS generation) underpins PDT drug
design.
\end{quote}

\begin{center}\rule{0.5\linewidth}{0.5pt}\end{center}

\subsection{Computational Bridge}\label{computational-bridge-9}

Net leaf CO\(_2\) uptake is modelled as a saturating light response
minus dark respiration; \cref{fig:unit_III_light_response} shows how net
assimilation rises from the light compensation point, climbs steeply at
low irradiance, and levels off at light saturation, with the C4 curve
saturating at higher photon flux than C3:

\begin{figure}[htbp]
\centering
\includegraphics[width=0.85\textwidth]{../figures/light_response_curves.png}
\caption{Light-response curves for C3, C4, and CAM plants. Net CO2 assimilation is plotted against photosynthetic photon flux density, showing different saturation points and maximum assimilation rates.}
\label{fig:unit_III_light_response}
\end{figure}

\begin{Shaded}
\begin{Highlighting}[]
\ImportTok{from}\NormalTok{ biology.botany }\ImportTok{import}\NormalTok{ photosynthesis\_rate}

\NormalTok{anet }\OperatorTok{=}\NormalTok{ photosynthesis\_rate(}\FloatTok{500.0}\NormalTok{)  }\CommentTok{\# µmol PAR m⁻² s⁻¹}
\BuiltInTok{print}\NormalTok{(}\BuiltInTok{round}\NormalTok{(anet, }\DecValTok{2}\NormalTok{))}
\end{Highlighting}
\end{Shaded}

\begin{quote}
\textbf{Clinical / systems note:} Non-photochemical quenching and
carotenoid cycles in plants parallel clinical strategies that manage
photochemical ROS in tissues (e.g.~PDT dosimetry), where excess
excitation must be routed to harmless decay channels.
\end{quote}

\begin{center}\rule{0.5\linewidth}{0.5pt}\end{center}

\subsubsection{Engineering Synthetic Photosynthesis: From
Photorespiration Bypass to Directed Evolution of
Rubisco}\label{engineering-synthetic-photosynthesis-from-photorespiration-bypass-to-directed-evolution-of-rubisco}

Rubisco (ribulose-1,5-bisphosphate carboxylase/oxygenase) is the most
abundant protein on Earth, yet it is famously inefficient:
\(k_{\text{cat}}\) \texorpdfstring{\ensuremath{\approx}}{~} 3 s\texorpdfstring{\textsuperscript{-}}{-}\texorpdfstring{\textsuperscript{1}}{1} and an oxygenation side-reaction that drains
20--50 \% of carbon in warm climates through \textbf{photorespiration}.
Engineering around this rate-limiting step is the central problem of
crop productivity. Three complementary strategies have matured between
2019 and 2024.

\textbf{(1) Photorespiration bypass.} The Long lab (\emph{Science} 2019)
introduced a synthetic glycolate-metabolism pathway into tobacco that
shunts glycolate (the toxic product of Rubisco oxygenation) directly to
pyruvate in the chloroplast rather than sending it through the
peroxisome/mitochondrion loop. Field trials demonstrated a
\textbf{about 40 \% biomass yield increase} over 2 years,
without compromising photosynthetic rate. The same bypass is being
translated into cowpea, soybean, and rice via the Realising Increased
Photosynthetic Efficiency (RIPE) consortium. \textbf{(2)
Cyanobacterial-style carbon-concentrating mechanisms.} Cyanobacteria
compartmentalise Rubisco inside \textbf{carboxysomes} that concentrate
CO\texorpdfstring{\textsubscript{2}}{2} locally, suppressing the oxygenation reaction. Expressing the
bacterial \textbf{BicA} HCO\texorpdfstring{\textsubscript{3}}{3}\texorpdfstring{\textsuperscript{-}}{-} transporter plus carboxysome shell
proteins in C\texorpdfstring{\textsubscript{3}}{3} plants is a about 10-year synthetic-biology
goal; the Price/Long/Hanson labs have cleared the first milestones in
\emph{Nicotiana} and \emph{Chlamydomonas}. \textbf{(3) Directed
evolution of Rubisco itself.} Classical selection inside \emph{E. coli}
is limited because bacterial Rubisco is too different from plant
Rubisco. The 2024 work by the Savage / Kortemme labs used
\textbf{ribosome-display directed evolution} with MS-based screening to
isolate Rubisco variants with about 2\texorpdfstring{\ensuremath{\times}}{x} improved specificity
factor (Sc/o = 110 \texorpdfstring{\ensuremath{\to}}{->} 220) --- approaching the theoretical maximum set by
the enzyme's two-step mechanism.

Quantitative targets put the stakes in context: improving crop
photosynthetic efficiency by even 10--20 \% in major C3 crops would be
agronomically large, but the realised food-security gain would depend on
canopy architecture, sink strength, water and nitrogen supply, local
climate, post-harvest losses, and access. Molecular agriculture now has
the tools -- CRISPR-Cas9 in plants, synthetic biology standardisation,
and AI-assisted structural hypotheses for Rubisco and its chaperones --
to test ideas that were largely conceptual a decade ago
\citep{abramson2024alphafold3,varadi2024alphafolddb}.

\begin{center}\rule{0.5\linewidth}{0.5pt}\end{center}

\subsection{Current Evidence and Frontier
Biology}\label{current-evidence-and-frontier-biology-12}

For \textbf{Photosynthesis}, frontier biology belongs inside the
evidence logic of the chapter. Metabolism is now studied as a regulated
network constrained by energy, redox balance, compartmentation, and
environment. The core reading question is this: photosynthesis claims
should distinguish light capture, carbon fixation, photorespiration,
water stress, and canopy context.

\begin{itemize}
\tightlist
\item
  \textbf{What to verify:} identify the observation, model, assay, or
  dataset that would make the claim stronger or weaker.
\item
  \textbf{What to qualify:} state the scale, organism, cell type,
  environmental condition, or population where the claim is expected to
  hold.
\item
  \textbf{What to compare:} test at least one alternative explanation,
  baseline, or null model before treating the pattern as causal.
\item
  \textbf{What to cite:} distinguish primary evidence, review synthesis,
  public dataset, and institutional guidance; for recent or numeric
  claims, prefer the source closest to the measurement and state what
  has changed since it was published.
\end{itemize}

A strong metabolic explanation names the flux, the limiting step, the
sensor, and the condition under which the pathway changes direction or
priority.

\textbf{Source practice:} For metabolic claims, keep the organism,
compartment, energetic state, and measurement method visible; a pathway
map is not enough without flux or concentration evidence.

\subsubsection{Current Evidence Map: Photosynthesis Under
Stress}\label{current-evidence-map-photosynthesis-under-stress}

\begin{figure}[htbp]
\centering
\includegraphics[width=0.82\linewidth,height=4.2in,keepaspectratio]{../figures/mermaid_inline/inline_mermaid_0061_d4c7a9a49ee8.png}
\caption{Photosynthesis in a plant is a coupled carbon-water decision: high light cannot raise growth if CO2 entry, water status, or sink demand becomes limiting.}
\end{figure}

\subsection{Summary}\label{summary-12}

\begin{itemize}
\tightlist
\item
  Light reactions (thylakoid): light \texorpdfstring{\ensuremath{\to}}{->} O\(_2\) (water splitting) + NADPH
  + ATP (\cref{eq:unit_III_photosynthesis_overall},
  \cref{eq:unit_III_water_splitting}). PS II drives water oxidation
  (Mn\(_4\)CaO\(_5\) cluster, Kok cycle); cytochrome b\(_6\)f pumps
  protons via Q cycle; PS I drives NADP\(^+\) reduction.
\item
  Photon energy is given by \(E = hc/\lambda\)
  (\cref{eq:unit_III_photon_energy}); a 680 nm photon supplies 176
  kJ/mol. Quantum yield \texorpdfstring{\ensuremath{\Phi}}{Phi}\(_{\text{O}_2}\) \texorpdfstring{\ensuremath{\approx}}{~} 0.10 corresponds to
  about 10 photons per O\(_2\), close to the theoretical
  minimum of 8.
\item
  Antenna complexes (LHCII, LHCI) use FRET
  (\cref{eq:unit_III_fret_rate}) and exciton transfer to funnel
  excitation to reaction centres at about 95\% efficiency. Chl
  a/b ratios and carotenoid composition are tuned to spectral and
  protective demands.
\item
  The Z-scheme spans 1.14 V (\cref{eq:unit_III_z_scheme_energy}) from
  H\(_2\)O (\(+0.82\) V) through P680, P680\emph{, Pheo, PQ, b\(_6\)f,
  PC, P700, P700}, A\(_0\), Fd to NADPH (\(-0.32\) V), driven by two
  photons per electron.
\item
  RuBisCO is hexadecameric (L\(_8\)S\(_8\), 8 active sites), requires
  CO\(_2\) + Mg\(^{2+}\) for active-site carbamylation, and shows a
  fundamental specificity-factor / turnover trade-off
  (\cref{eq:unit_III_specificity_factor}). Plant
  \(S_{c/o} \approx 80\)--100; theoretical maximum about 400.
\item
  Calvin cycle (stroma): 18 ATP + 12 NADPH fix 6 CO\(_2\) \texorpdfstring{\ensuremath{\to}}{->} 1 glucose.
  Photorespiration costs C3 plants about 25\% of fixed carbon
  (the carboxylation/oxygenation ratio is set by
  \cref{eq:unit_III_carb_oxy_ratio}).
\item
  C4 uses Kranz anatomy (spatial separation); CAM uses temporal
  separation (night CO\(_2\) fixation). Both have near-zero
  photorespiration but higher ATP costs. Water-use efficiency
  (\cref{eq:unit_III_water_use_efficiency}) ranks CAM \textgreater{} C4
  \textgreater{} C3 by an order of magnitude.
\item
  The thioredoxin/ferredoxin system
  (\cref{eq:unit_III_thioredoxin_cascade}) couples PS I activity to
  Calvin-cycle enzyme activation via reversible disulfide reduction.
  Stromal pH 8.0 and Mg\(^{2+}\) release further activate RuBisCO and
  FBPase.
\item
  NPQ (\cref{eq:unit_III_npq_definition}) protects the photosynthetic
  apparatus through the xanthophyll cycle
  (\cref{eq:unit_III_xanthophyll_cycle}), state transitions, and
  photoinhibition repair. Carotenoids quench triplet Chl and singlet
  O\(_2\) \citep{hager1971}.
\item
  Artificial photosynthesis bridges biology (engineered cyanobacteria)
  and chemistry (PEC cells, Mn\(_4\)Ca biomimetics); silicon PV is more
  efficient but does not produce storable fuel directly.
\item
  \textbf{Connections:} See
  \cref{sec:unit_III_bioenergetics_and_respiration} for chemiosmotic
  coupling in mitochondria vs.~chloroplasts, Unit VIII for whole-plant
  carbon and water trade-offs, and Unit X for primary productivity.
\end{itemize}

\begin{center}\rule{0.5\linewidth}{0.5pt}\end{center}

\subsection{Review Questions}\label{review-questions-12}

\begin{enumerate}
\def\labelenumi{\arabic{enumi}.}
\item
  Describe the structure of a chloroplast and explain how the
  organisation of the thylakoid membrane is adapted for efficient
  photosynthesis.
\item
  Explain the Z-scheme of photosynthetic electron transport. Why are two
  photosystems needed to move electrons from water to NADP\(^+\)? Quote
  at least three intermediate redox potentials in your answer.
\item
  Calculate the energy of a 440 nm blue photon and a 680 nm red photon
  (per mole). Explain why both are equally effective at driving
  photochemistry despite different photon energies.
\item
  Describe the Kok cycle (S-state cycle) of water oxidation. Why must
  four photons be absorbed to produce one O\(_2\)?
\item
  Compare the Q cycle in Complex III of mitochondria with the Q cycle in
  cytochrome b\(_6\)f of chloroplasts. What is the common functional
  principle?
\item
  Describe RuBisCO's structure and catalytic mechanism. Why is the
  specificity factor \(S_{c/o}\) inversely correlated with
  \(k_{\text{cat}}\) across the tree of life?
\item
  Describe the three phases of the Calvin cycle. Calculate the total ATP
  and NADPH cost to fix 6 CO\(_2\) into one glucose molecule, and the
  minimum number of photons required.
\item
  Explain the thioredoxin/ferredoxin redox regulation system. Which
  Calvin cycle enzymes are activated, and what would happen in a
  \emph{trxR} (thioredoxin reductase) knockout in the light?
\item
  Compare C3, C4, and CAM photosynthesis in terms of leaf anatomy,
  biochemistry, water-use efficiency, and ecological
  \hyperref[gl:niche]{\textbf{niche}}. Use
  \cref{eq:unit_III_intrinsic_wue} to explain why CAM achieves the
  highest WUE.
\item
  Explain the xanthophyll cycle and qE non-photochemical quenching. Why
  is this mechanism essential for plant survival in fluctuating light?
\item
  Calculate the energy efficiency of photosynthesis if 48 photons of 680
  nm light are needed to fix 6 CO\(_2\). Compare this with the
  efficiency of a typical silicon solar cell (about 22\%).
\item
  Engineering reduced photorespiration has been proposed to increase
  crop yields. Describe two approaches and discuss potential risks.
\item
  Call \breaktt{photosynthesis\_rate} at 100 vs.~2000 \texorpdfstring{\ensuremath{\mu}}{mu}mol m\texorpdfstring{\textsuperscript{-}}{-}\texorpdfstring{\textsuperscript{2}}{2} s\texorpdfstring{\textsuperscript{-}}{-}\texorpdfstring{\textsuperscript{1}}{1}
  default parameters. At which irradiance are you nearer light
  saturation, and how does that relate to greenhouse shading trials?
\item
  Compare cyclic vs.~linear electron flow around PS I in terms of
  \textbf{ATP:NADPH} balance during the Calvin cycle.
\item
  Compute the Förster transfer rate at \(r = 7\) nm given \(R_0 = 5\) nm
  and \(\tau_D = 4\) ns using \cref{eq:unit_III_fret_rate}. Repeat for
  \(r = 10\) nm and comment on the spatial constraint imposed.
\item
  Spinach has \(\tau \approx 90\), atmospheric CO\(_2\) = 410 ppm.
  Compute \(v_c/v_o\) at chloroplastic CO\(_2\) = 7 \texorpdfstring{\ensuremath{\mu}}{mu}M and O\(_2\) = 250
  \texorpdfstring{\ensuremath{\mu}}{mu}M using \cref{eq:unit_III_carb_oxy_ratio}. What fraction of RuBisCO
  turnovers are oxygenations?
\end{enumerate}

\begin{center}\rule{0.5\linewidth}{0.5pt}\end{center}

\subsection{Further Reading and Source
Notes}\label{further-reading-and-source-notes-12}

\begin{itemize}
\tightlist
\item
  Mitchell (1961). Coupling of phosphorylation to electron and hydrogen
  transfer by a chemi-osmotic type of mechanism. \emph{Nature}, 191
  \citep{mitchell1961}.
\item
  Hager (1971). Die Reversiblen lichtabhängigen Xanthophyllumwandlungen
  im Chloroplasten. \emph{Berichte der Deutschen Botanischen
  Gesellschaft}, 84 \citep{hager1971}.
\end{itemize}

\begin{center}\rule{0.5\linewidth}{0.5pt}\end{center}

\subsection{Key Terms}\label{key-terms-12}

{\def\LTcaptype{none} % do not increment counter
\begin{longtable}[]{@{}
  >{\raggedright\arraybackslash}p{(\linewidth - 2\tabcolsep) * \real{0.2857}}
  >{\raggedright\arraybackslash}p{(\linewidth - 2\tabcolsep) * \real{0.7143}}@{}}
\toprule\noalign{}
\begin{minipage}[b]{\linewidth}\raggedright
Term
\end{minipage} & \begin{minipage}[b]{\linewidth}\raggedright
Definition
\end{minipage} \\
\midrule\noalign{}
\endhead
\bottomrule\noalign{}
\endlastfoot
\textbf{Photosystem II (PS II)} & Light-driven water oxidase; P680
reaction centre; O\(_2\) evolution via Mn\(_4\)CaO\(_5\) cluster \\
\textbf{Photosystem I (PS I)} & Light-driven NADP\(^+\) reductase; P700
reaction centre; produces NADPH via ferredoxin \\
\textbf{Z-scheme} & Complete electron transport path from H\(_2\)O
(\(+0.82\) V) to NADPH (\(-0.32\) V) via two photosystems \\
\textbf{RuBisCO} & Ribulose-1,5-bisphosphate carboxylase/oxygenase;
L\(_8\)S\(_8\); most abundant enzyme; CO\(_2\) fixation \\
\textbf{Specificity factor} (\(S_{c/o}\), \texorpdfstring{\ensuremath{\tau}}{tau}) & \(V_c K_o / V_o K_c\);
quantifies CO\(_2\) vs O\(_2\) preference of RuBisCO \\
\textbf{Calvin cycle} & Three-phase cycle fixing CO\(_2\) into G3P in
the stroma; 3 ATP + 2 NADPH per CO\(_2\) \\
\textbf{Thioredoxin} & Small redox protein that reduces regulatory
disulfides on Calvin-cycle enzymes in the light \\
\textbf{Photorespiration} & RuBisCO oxygenase activity; wastes
about 25\% of fixed carbon in C3 plants; releases CO\(_2\) \\
\textbf{C4 photosynthesis} & CO\(_2\) concentrating mechanism using PEP
carboxylase and Kranz anatomy \\
\textbf{CAM} & Crassulacean acid metabolism; temporal separation of
CO\(_2\) fixation (night) and Calvin cycle (day) \\
\textbf{Water-use efficiency (WUE)} & Carbon fixed per water lost; CAM
\textgreater{} C4 \textgreater{} C3 \\
\textbf{Kok cycle} & S-state cycle of the OEC; accumulates 4 oxidising
equivalents before O\(_2\) release \\
\textbf{Antenna complex} & Light-harvesting pigment-protein complex;
funnels energy to reaction centre \\
\textbf{FRET} & Förster Resonance Energy Transfer; dipole-dipole energy
transfer between pigments; rate \texorpdfstring{\ensuremath{\propto}}{propto} 1/r\(^6\) \\
\textbf{Quantum yield} & Molecules of product per photon absorbed;
\texorpdfstring{\ensuremath{\Phi}}{Phi}\(_{\text{O}_2}\) \texorpdfstring{\ensuremath{\approx}}{~} 0.10 \\
\textbf{Non-photochemical quenching (NPQ)} & Thermal dissipation of
excess light energy; comprises qE, qT, qI components \\
\textbf{Xanthophyll cycle} & Violaxanthin \texorpdfstring{\ensuremath{\leftrightarrow}}{<->} zeaxanthin interconversion;
regulates qE \\
\textbf{PsbS} & PS II protein with two lumen-exposed Glu residues;
protonated at low pH; triggers qE \\
\textbf{Plastocyanin} & Soluble copper protein carrying electrons from
b\(_6\)f to PS I in the lumen \\
\textbf{Kranz anatomy} & Wreath-like leaf anatomy of C4 plants; bundle
sheath surrounds vascular tissue \\
\textbf{Cyclic electron flow} & PS I-mediated; Fd \texorpdfstring{\ensuremath{\to}}{->} PQ \texorpdfstring{\ensuremath{\to}}{->} b\(_6\)f \texorpdfstring{\ensuremath{\to}}{->} PC \texorpdfstring{\ensuremath{\to}}{->}
PS I; produces ATP without NADPH \\
\end{longtable}
}

\begin{center}\rule{0.5\linewidth}{0.5pt}\end{center}

\subsubsection{Companion Source
Module}\label{companion-source-module-12}

\textbf{Photosynthesis} should leave a reproducible trail from a
biological claim to the code, figure, diagram, or paper-based activity
that can test it. Use the surfaces below to inspect the chapter's
assumptions, rerun the relevant model, or compare the manuscript
explanation with companion labs and figures.

{\def\LTcaptype{none} % do not increment counter
\begin{longtable}[]{@{}
  >{\raggedright\arraybackslash}p{(\linewidth - 2\tabcolsep) * \real{0.5000}}
  >{\raggedright\arraybackslash}p{(\linewidth - 2\tabcolsep) * \real{0.5000}}@{}}
\toprule\noalign{}
\begin{minipage}[b]{\linewidth}\raggedright
Surface
\end{minipage} & \begin{minipage}[b]{\linewidth}\raggedright
Use it for
\end{minipage} \\
\midrule\noalign{}
\endhead
\bottomrule\noalign{}
\endlastfoot
\breaktt{src/biology/botany/botany.py} (\breaktt{photosynthesis\_rate},
\breaktt{light\_response\_curve}) & Reproduce light-response and
environmental-limitation scenarios. \\
\breaktt{src/visualization/plots.py}
(\breaktt{plot\_light\_response\_curve}) & Inspect saturation,
compensation points, and axis labeling. \\
\breaktt{src/mermaid/biology\_diagrams.py}
(\breaktt{photosynthesis\_light\_dark\_diagram}) & Separate light
reactions, carbon fixation, and regulation. \\
\end{longtable}
}

\textbf{Reproducibility check:} report light intensity, CO2,
temperature, water status, and plant pathway before comparing
photosynthetic rates. \textbf{Cross-reference:} compare with
\cref{sec:unit_VIII_plant_responses} and
\cref{sec:unit_X_ecosystem_ecology}.

\newpage

\section{Metabolic Integration and
Regulation}\label{metabolic-integration-and-regulation}

\label{sec:unit_III_metabolic_integration}

\begin{quote}
\textbf{Ch 11} \texorpdfstring{\ensuremath{\cdot}}{.} Level 3/3 \texorpdfstring{\ensuremath{\cdot}}{.} 60 min read \texorpdfstring{\ensuremath{\cdot}}{.} 100 min lecture \texorpdfstring{\ensuremath{\cdot}}{.}
Prerequisites: \cref{sec:unit_III_bioenergetics_and_respiration},
\cref{sec:unit_III_photosynthesis}
\end{quote}

\subsection{Learning Objectives}\label{learning-objectives-13}

By the end of this chapter, you should be able to:

\begin{enumerate}
\def\labelenumi{\arabic{enumi}.}
\tightlist
\item
  Trace how the cell coordinates catabolic and anabolic pathways through
  shared cofactors and \hyperref[gl:allosteric]{\textbf{allosteric}}
  cross-regulation, and predict the metabolic state from cofactor and
  signal-molecule ratios.
\item
  Describe AMPK and mTORC1 as opposing energy-sensing hubs and explain
  how they toggle the metabolic state of the cell, including the
  molecular basis of mTOR's role as a nutrient sensor (Rag GTPases,
  Ragulator).
\item
  Describe the metabolic states of fed, post-absorptive, fasting, and
  starvation, including hormonal control and the timeline of fuel
  switching.
\item
  Explain insulin signalling in detail: the PI3K/Akt/mTOR pathway,
  \hyperref[gl:glut4]{\textbf{GLUT4}} translocation mechanism, and
  glycogen synthesis.
\item
  Explain glucagon signalling: cAMP/PKA activation of glycogenolysis,
  gluconeogenesis, and the phosphorylase kinase cascade.
\item
  Distinguish allosteric from transcriptional metabolic regulation given
  an \hyperref[gl:enzyme]{\textbf{enzyme}}-activity time-course, using
  gluconeogenesis bypass-enzyme control as the worked case.
\item
  Compare ATP yield and regulation of glucose versus fatty-acid
  oxidation, and calculate the net ATP yield from palmitoyl-CoA.
\item
  Explain ketone body metabolism and its role during starvation.
\item
  Describe AMPK structure (two catalytic + two regulatory subunits) and
  activation mechanisms (LKB1 vs CaMKK\texorpdfstring{\ensuremath{\beta}}{beta}).
\item
  Trace flux control through the glucose-fatty-acid (Randle) cycle under
  fed versus fasted conditions, and predict the resulting respiratory
  quotient.
\item
  Define substrate (futile) cycles and explain their control advantages.
\item
  Explain obesity, metabolic syndrome, and insulin resistance at the
  molecular level, including visceral fat as an endocrine organ.
\item
  Apply metabolic control analysis (MCA): flux control coefficients
  \(C_i^J\) and the summation theorem.
\item
  Calculate the adenylate energy charge from concentrations and predict
  the resulting AMPK activation state.
\item
  Describe modern metabolomics: NMR vs mass spectrometry approaches and
  identification challenges.
\end{enumerate}

\subsubsection{Study Blueprint}\label{study-blueprint-13}

\begin{itemize}
\tightlist
\item
  \textbf{Big idea:} Metabolism is a regulated network that reallocates
  flux across tissues, time, and nutrient states.
\item
  \textbf{Core concepts:} flux, energy charge, hormonal control,
  fed/fasted states.
\item
  \textbf{Framework alignment:} Vision \& Change: Pathways and
  transformations of energy and matter, Systems; AP Biology: Energetics,
  Systems Interactions; NGSS-style topics: Matter and Energy in
  Organisms and Ecosystems.
\item
  \textbf{Model or quantitative lens:} Energy charge, control points,
  and pathway-flux comparisons.
\item
  \textbf{Data skill:} Use pathway evidence to infer which metabolic
  state or tissue is active.
\item
  \textbf{Practice cadence:} Representing and Describing Data,
  Statistical Tests and Data Analysis.
\item
  \textbf{Common misconception to repair:} A pathway is not a one-way
  assembly line; reversibility and regulation define the real route.
\item
  \textbf{Primary lab:} \cref{sec:lab_unit_III_metabolic_integration}.
\item
  \textbf{Question bank:} \cref{sec:q_unit_III_metabolic_integration}.
\item
  \textbf{Transfer task:} Apply metabolic network reasoning to diabetes,
  fasting, exercise, or cancer metabolism.
\item
  \textbf{Bridge to computation:}
  \breaktt{biology.biochemistry.biochemistry.reaction\_free\_energy}.
\end{itemize}

\begin{center}\rule{0.5\linewidth}{0.5pt}\end{center}

\begin{quote}
\textbf{Opening Vignette: The Brain on Ketones --- Surviving Starvation}

The human brain is extraordinarily metabolically demanding: it consumes
roughly 20\% of the body's resting energy budget despite constituting
about 2\% of body mass. Under normal conditions, it runs almost
exclusively on glucose. But glucose stores --- maintained as liver
glycogen --- last only about 12--24 hours of fasting. So how does the
brain survive during prolonged starvation?

The liver activates a metabolic programme: free fatty acids released
from adipose tissue are not completely oxidised within the liver itself.
Instead, the liver packages their acetyl-CoA units into \textbf{ketone
bodies} (acetoacetate and \texorpdfstring{\ensuremath{\beta}}{beta}-hydroxybutyrate) and exports them to the
bloodstream. After 1--2 days of fasting, the brain begins importing
ketones; after 3--5 days, it derives up to 75\% of its energy from
ketone bodies. Circulating ketone concentrations rise from
about 0.1 mM (fed state) to 6--8 mM (prolonged starvation).
This elegant metabolic switch --- coordinated by insulin falling and
glucagon rising --- is one of the clearest examples of metabolic
integration studied in this chapter.

\emph{Primary source: Owen, O. E. et al.~(1967). Brain metabolism during
fasting. Journal of Clinical Investigation, 46(10), 1589--1595.}
\end{quote}

\begin{center}\rule{0.5\linewidth}{0.5pt}\end{center}

\subsection{Pathways Are Not Isolated}\label{pathways-are-not-isolated}

\hyperref[gl:glycolysis]{\textbf{Glycolysis}}, the TCA cycle,
beta-oxidation, and biosynthetic pathways do not operate independently.
They share intermediates, co-factors (NAD\(^+\)/NADH, NADP\(^+\)/NADPH,
CoA), and are subject to \textbf{allosteric cross-regulation}.
Integration of these pathways defines the \textbf{metabolic state} of
the cell.

The two fundamental modes: - \textbf{Catabolism} --- break down
nutrients \texorpdfstring{\ensuremath{\to}}{->} ATP + reduced cofactors (NADH, FADH\(_2\)) -
\textbf{Anabolism} --- use ATP + reduced cofactors (chiefly NADPH) \texorpdfstring{\ensuremath{\to}}{->}
biosynthesis of macromolecules

These modes are reciprocally regulated because of shared cofactors and
master regulatory kinases. The metabolic state is encoded in:

{\def\LTcaptype{none} % do not increment counter
\begin{longtable}[]{@{}
  >{\raggedright\arraybackslash}p{(\linewidth - 4\tabcolsep) * \real{0.2459}}
  >{\raggedright\arraybackslash}p{(\linewidth - 4\tabcolsep) * \real{0.3607}}
  >{\raggedright\arraybackslash}p{(\linewidth - 4\tabcolsep) * \real{0.3934}}@{}}
\toprule\noalign{}
\begin{minipage}[b]{\linewidth}\raggedright
Signal molecule
\end{minipage} & \begin{minipage}[b]{\linewidth}\raggedright
Fed state (high energy)
\end{minipage} & \begin{minipage}[b]{\linewidth}\raggedright
Fasted state (low energy)
\end{minipage} \\
\midrule\noalign{}
\endhead
\bottomrule\noalign{}
\endlastfoot
AMP/ATP ratio & Low & High \\
NADH/NAD\(^+\) ratio & High & Low (in
\hyperref[gl:cytoplasm]{\textbf{cytoplasm}}) \\
Acetyl-CoA & High & Low (cytoplasmic) \\
Citrate & High & Low \\
Malonyl-CoA & High & Low \\
Insulin/Glucagon ratio & High & Low \\
\end{longtable}
}

\begin{figure}[htbp]
\centering
\includegraphics[width=0.82\linewidth,height=4.2in,keepaspectratio]{../figures/mermaid_inline/inline_mermaid_0062_dafd60cfc7de.png}
\caption{Metabolic states from fed to starvation, showing the hormonal switches and metabolic pathway changes at each stage. Insulin dominates the fed state; glucagon dominates fasting; ketogenesis sustains the brain during starvation.}
\end{figure}

\subsubsection{Fasting Physiology: Fuel Switch
Timeline}\label{fasting-physiology-fuel-switch-timeline}

The body switches fuels in a programmed sequence as fasting progresses.
The transitions are coordinated by falling insulin and rising glucagon
(and later cortisol):

{\def\LTcaptype{none} % do not increment counter
\begin{longtable}[]{@{}
  >{\raggedright\arraybackslash}p{(\linewidth - 8\tabcolsep) * \real{0.2456}}
  >{\raggedright\arraybackslash}p{(\linewidth - 8\tabcolsep) * \real{0.2632}}
  >{\raggedright\arraybackslash}p{(\linewidth - 8\tabcolsep) * \real{0.2281}}
  >{\raggedright\arraybackslash}p{(\linewidth - 8\tabcolsep) * \real{0.1754}}
  >{\raggedright\arraybackslash}p{(\linewidth - 8\tabcolsep) * \real{0.0877}}@{}}
\toprule\noalign{}
\begin{minipage}[b]{\linewidth}\raggedright
Time post-meal
\end{minipage} & \begin{minipage}[b]{\linewidth}\raggedright
Predominant fuel
\end{minipage} & \begin{minipage}[b]{\linewidth}\raggedright
Liver process
\end{minipage} & \begin{minipage}[b]{\linewidth}\raggedright
Brain fuel
\end{minipage} & \begin{minipage}[b]{\linewidth}\raggedright
Notes
\end{minipage} \\
\midrule\noalign{}
\endhead
\bottomrule\noalign{}
\endlastfoot
0--4 h (absorptive) & Dietary glucose & Glycogenesis, lipogenesis &
Glucose & Insulin high, glucagon low \\
4--12 h (post-absorptive) & Liver glycogen & Glycogenolysis & Glucose &
Glucagon rises \\
12--24 h (early fasting) & Liver glycogen \texorpdfstring{\ensuremath{\to}}{->} muscle FA & Glycogen
depleting; gluconeogenesis ramping & Glucose (mostly from GNG) & Liver
glycogen mostly depleted by about 24 h \\
1--3 days & Muscle FA + early ketones & Gluconeogenesis (lactate,
alanine, glycerol); ketogenesis starting & Glucose 75\%, ketones 25\% &
Cortisol rises; protein catabolism increases \\
3--7 days & Ketones & Maximal ketogenesis & Glucose 25\%, ketones 75\% &
Brain adapts; muscle protein sparing \\
\textgreater{} 1--2 weeks & Ketones (sustained) & Ketogenesis at steady
state & Ketones 65--75\%, glucose 25--35\% & Resting metabolic rate
falls about 20\% \\
\textgreater{} 30+ days (starvation) & Body protein (terminal) &
Gluconeogenesis from muscle protein dominates & Glucose 50\%, ketones
50\% & Death imminent when about 50\% of protein lost \\
\end{longtable}
}

This fuel-switching timeline is one of the most elegant examples of
metabolic integration. It is also clinically relevant: bariatric
surgery, ketogenic diets, and prolonged fasting interventions most
exploit (or perturb) this same switch.

\begin{quote}
\textbf{Concept Check 1:} Why is it essential that glycolysis and
gluconeogenesis are reciprocally regulated and rarely fully active
simultaneously? What would happen if both ran at maximum rate?
\end{quote}

\begin{center}\rule{0.5\linewidth}{0.5pt}\end{center}

\subsection{Insulin Signalling: The PI3K/Akt/mTOR
Pathway}\label{insulin-signalling-the-pi3kaktmtor-pathway}

Insulin is the master anabolic \hyperref[gl:hormone]{\textbf{hormone}},
secreted by pancreatic beta-cells in response to elevated blood glucose.
Its signalling cascade is one of the most important in metabolic
regulation:

\[\text{Insulin} \to \text{IR (RTK)} \to \text{IRS-1/2} \to \text{PI3K} \to \text{PIP}_3 \to \text{PDK1 + mTORC2} \to \text{Akt (Thr308 + Ser473)} \tag{11.1}\]

\begin{figure}[htbp]
\centering
\includegraphics[width=0.82\linewidth,height=4.2in,keepaspectratio]{../figures/mermaid_inline/inline_mermaid_0063_09287921f06e.png}
\caption{Flowchart for Insulin Signalling: The PI3K/Akt/mTOR Pathway: Insulin, Insulin receptor (alpha2beta2 tyrosine kinase), IRS-1/2 (Tyr phosphorylated), and PI3K (p85/p110) form the diagram's primary path or branches.}
\end{figure}

\emph{Insulin signalling cascade from receptor to downstream metabolic
and growth effects (Mermaid).} Akt is the central node, integrating PDK1
(Thr308) and mTORC2 (Ser473) phosphorylations to drive glucose uptake,
glycogen synthesis, lipogenesis, and protein synthesis.

\textbf{Downstream Akt targets:}

{\def\LTcaptype{none} % do not increment counter
\begin{longtable}[]{@{}
  >{\raggedright\arraybackslash}p{(\linewidth - 4\tabcolsep) * \real{0.1818}}
  >{\raggedright\arraybackslash}p{(\linewidth - 4\tabcolsep) * \real{0.3030}}
  >{\raggedright\arraybackslash}p{(\linewidth - 4\tabcolsep) * \real{0.5152}}@{}}
\toprule\noalign{}
\begin{minipage}[b]{\linewidth}\raggedright
Target
\end{minipage} & \begin{minipage}[b]{\linewidth}\raggedright
Akt effect
\end{minipage} & \begin{minipage}[b]{\linewidth}\raggedright
Metabolic outcome
\end{minipage} \\
\midrule\noalign{}
\endhead
\bottomrule\noalign{}
\endlastfoot
AS160 (TBC1D4) & Phosphorylation \texorpdfstring{\ensuremath{\to}}{->} Rab-GAP inactivated & \textbf{GLUT4
vesicle exocytosis} \texorpdfstring{\ensuremath{\to}}{->} glucose uptake \\
GSK3 (glycogen synthase kinase 3) & Phosphorylation \texorpdfstring{\ensuremath{\to}}{->} inhibited &
\textbf{Glycogen synthase active} \texorpdfstring{\ensuremath{\to}}{->} glycogen synthesis \\
FoxO1 & Phosphorylation \texorpdfstring{\ensuremath{\to}}{->} nuclear exclusion & \textbf{PEPCK and G6Pase
\hyperref[gl:gene]{\textbf{gene}}s OFF} \texorpdfstring{\ensuremath{\to}}{->} gluconeogenesis inhibited \\
TSC2 & Phosphorylation \texorpdfstring{\ensuremath{\to}}{->} inhibited & \textbf{mTORC1 activated} \texorpdfstring{\ensuremath{\to}}{->}
\hyperref[gl:protein]{\textbf{protein}} synthesis, lipogenesis \\
BAD & Phosphorylation \texorpdfstring{\ensuremath{\to}}{->} binds 14-3-3 & \textbf{Anti-apoptotic} \texorpdfstring{\ensuremath{\to}}{->} cell
survival \\
PDE3B & Phosphorylation \texorpdfstring{\ensuremath{\to}}{->} activated & \textbf{cAMP degradation} \texorpdfstring{\ensuremath{\to}}{->}
opposes glucagon \\
\end{longtable}
}

\textbf{GLUT4 translocation mechanism (in detail):}

\begin{enumerate}
\def\labelenumi{\arabic{enumi}.}
\tightlist
\item
  \textbf{Basal state.} GLUT4 is sequestered in specialised
  intracellular vesicles (GSVs, GLUT4 storage vesicles)
  about 95\% of the time, with about 5\% at the plasma
  membrane. The vesicles are tethered by the inhibitory Rab-GAP
  \textbf{AS160 (TBC1D4)}, which keeps Rab10 in its GDP-bound (inactive)
  state.
\item
  \textbf{Insulin signal.} Akt phosphorylates AS160 on multiple sites
  (Thr642 most important), inhibiting its GAP activity.
\item
  \textbf{Rab10 activation.} With AS160 off, Rab10 accumulates in the
  GTP-bound active state, promoting vesicle docking at the plasma
  membrane.
\item
  \textbf{SNARE-mediated fusion.} Synaptobrevin-2 (VAMP2) on GSVs
  engages plasma-membrane syntaxin-4 + SNAP-23, driving membrane fusion.
\item
  \textbf{GLUT4 exposure.} GLUT4 appears at the cell surface, increasing
  glucose uptake 10--20-fold.
\item
  \textbf{Endocytosis on signal termination.} When insulin signal ends,
  GLUT4 is endocytosed back to GSVs via clathrin-coated pits and a slow
  constitutive recycling pathway.
\end{enumerate}

In type 2 diabetes, this cascade is disrupted at multiple levels (see
clinical box below) so that GLUT4 translocation is impaired even at high
circulating insulin.

\begin{quote}
\textbf{Clinical Connection: Insulin Resistance and Type 2 Diabetes} In
type 2 diabetes, insulin signalling is impaired at multiple levels: 1.
\textbf{IRS-1 serine phosphorylation} (by JNK, PKC-theta, IKK-beta)
prevents tyrosine phosphorylation by the
\hyperref[gl:insulin-receptor]{\textbf{insulin receptor}} 2.
\textbf{Ceramide accumulation} activates PP2A, which dephosphorylates
Akt 3. \textbf{DAG accumulation} activates PKC-epsilon (liver) and
PKC-theta (muscle), which phosphorylate IRS-1 on inhibitory serine
residues 4. \textbf{ER stress} (from lipid overload) activates JNK \texorpdfstring{\ensuremath{\to}}{->}
IRS-1 Ser307 phosphorylation The net result: GLUT4 translocation is
impaired, glycogen synthesis is reduced, and gluconeogenesis is not
suppressed. Metformin (first-line therapy) works partly through AMPK
activation and partly through inhibition of mitochondrial Complex I. see
\cref{sec:unit_II_cell_signaling} for RTK signalling.
\end{quote}

\begin{quote}
\textbf{Concept Check 2:} Thiazolidinediones (e.g., rosiglitazone)
activate PPAR-gamma in adipose tissue, increasing adipocyte
differentiation and lipid storage. Explain how increasing fat storage in
adipose tissue could paradoxically improve insulin sensitivity in muscle
and liver.
\end{quote}

\begin{center}\rule{0.5\linewidth}{0.5pt}\end{center}

\subsection{Glucagon Signalling: cAMP/PKA
Pathway}\label{glucagon-signalling-camppka-pathway}

Glucagon is secreted by pancreatic alpha-cells when blood glucose falls.
It acts primarily on the \textbf{liver} (hepatocytes express glucagon
receptors; muscle does not):

\textbf{Glucagon \texorpdfstring{\ensuremath{\to}}{->} Glucagon receptor (GPCR) \texorpdfstring{\ensuremath{\to}}{->} G\(_s\) \texorpdfstring{\ensuremath{\to}}{->} Adenylyl cyclase
\texorpdfstring{\ensuremath{\to}}{->} cAMP \texorpdfstring{\ensuremath{\to}}{->} PKA}

\begin{figure}[htbp]
\centering
\includegraphics[width=0.82\linewidth,height=4.2in,keepaspectratio]{../figures/mermaid_inline/inline_mermaid_0064_27dc0606c7c7.png}
\caption{Flowchart for Glucagon Signalling: cAMP/PKA Pathway: Glucagon (or epinephrine via beta-AR), Glucagon receptor / beta-AR (GPCR), Galphas, and Adenylyl cyclase form the diagram's primary path or branches.}
\end{figure}

\emph{Glucagon (or epinephrine) signalling cascade in the liver
(Mermaid).} The cAMP/PKA cascade simultaneously activates glycogenolysis
(via phosphorylase kinase) and gluconeogenesis (via CREB transcription)
while suppressing glycolysis (via PFK-2/pyruvate kinase
phosphorylation).

\subsubsection{The Phosphorylase Kinase Cascade --- A Signal
Amplifier}\label{the-phosphorylase-kinase-cascade-a-signal-amplifier}

The chain glucagon \texorpdfstring{\ensuremath{\to}}{->} GPCR \texorpdfstring{\ensuremath{\to}}{->} cAMP \texorpdfstring{\ensuremath{\to}}{->} PKA \texorpdfstring{\ensuremath{\to}}{->} \textbf{phosphorylase kinase}
(PhK) \texorpdfstring{\ensuremath{\to}}{->} glycogen phosphorylase is one of biology's most elegant
amplification cascades. Each step amplifies the signal:

{\def\LTcaptype{none} % do not increment counter
\begin{longtable}[]{@{}
  >{\raggedright\arraybackslash}p{(\linewidth - 2\tabcolsep) * \real{0.1481}}
  >{\raggedright\arraybackslash}p{(\linewidth - 2\tabcolsep) * \real{0.8519}}@{}}
\toprule\noalign{}
\begin{minipage}[b]{\linewidth}\raggedright
Step
\end{minipage} & \begin{minipage}[b]{\linewidth}\raggedright
Numerical amplification
\end{minipage} \\
\midrule\noalign{}
\endhead
\bottomrule\noalign{}
\endlastfoot
1 glucagon \texorpdfstring{\ensuremath{\to}}{->} 1 receptor (1:1) & 1\texorpdfstring{\ensuremath{\times}}{x} \\
1 receptor \texorpdfstring{\ensuremath{\to}}{->} about 100 G\texorpdfstring{\ensuremath{\alpha}}{alpha}\(_s\) activated (during signal) &
about 100\texorpdfstring{\ensuremath{\times}}{x} \\
1 adenylyl cyclase \texorpdfstring{\ensuremath{\to}}{->} about 100 cAMP/sec &
about 100\texorpdfstring{\ensuremath{\times}}{x} \\
4 cAMP \texorpdfstring{\ensuremath{\to}}{->} 1 PKA holoenzyme dissociation (releases 2 catalytic subunits) &
gating \\
1 PKA \texorpdfstring{\ensuremath{\to}}{->} about 100 phosphorylase kinase & about 100\texorpdfstring{\ensuremath{\times}}{x} \\
1 PhK \texorpdfstring{\ensuremath{\to}}{->} about 100 phosphorylase b \texorpdfstring{\ensuremath{\to}}{->} a & about 100\texorpdfstring{\ensuremath{\times}}{x} \\
1 phosphorylase a \texorpdfstring{\ensuremath{\to}}{->} about 100 glucose-1-P/sec &
about 100\texorpdfstring{\ensuremath{\times}}{x} \\
\textbf{End-to-end amplification} &
\textbf{about 10\(^7\)--10\(^8\)} \\
\end{longtable}
}

A single glucagon binding event releases a burst of
about 10\(^7\) glucose-1-phosphate molecules within seconds.
This is how a signal at femtomolar hormone concentrations produces a
millimolar metabolite response.

\textbf{PKA targets in the liver:}

{\def\LTcaptype{none} % do not increment counter
\begin{longtable}[]{@{}
  >{\raggedright\arraybackslash}p{(\linewidth - 4\tabcolsep) * \real{0.1818}}
  >{\raggedright\arraybackslash}p{(\linewidth - 4\tabcolsep) * \real{0.3030}}
  >{\raggedright\arraybackslash}p{(\linewidth - 4\tabcolsep) * \real{0.5152}}@{}}
\toprule\noalign{}
\begin{minipage}[b]{\linewidth}\raggedright
Target
\end{minipage} & \begin{minipage}[b]{\linewidth}\raggedright
PKA effect
\end{minipage} & \begin{minipage}[b]{\linewidth}\raggedright
Metabolic outcome
\end{minipage} \\
\midrule\noalign{}
\endhead
\bottomrule\noalign{}
\endlastfoot
Phosphorylase kinase & Activated & Glycogen phosphorylase activated \texorpdfstring{\ensuremath{\to}}{->}
\textbf{glycogenolysis} \\
Glycogen synthase & Inhibited (phosphorylated) & Glycogen synthesis
OFF \\
PFK-2/FBPase-2 & FBPase-2 activated \texorpdfstring{\ensuremath{\to}}{->} F2,6BP destroyed &
\textbf{Glycolysis inhibited} \\
Pyruvate kinase & Inhibited (phosphorylated) & Glycolysis slowed \\
CREB & Activated & \textbf{PEPCK and G6Pase gene
\hyperref[gl:transcription]{\textbf{transcription}}} \texorpdfstring{\ensuremath{\to}}{->}
gluconeogenesis \\
IP3R & Sensitised & Ca\(^{2+}\) release from ER \\
HSL (adipose, via \texorpdfstring{\ensuremath{\beta}}{beta}-AR) & Activated & Lipolysis (in adipose; not
liver) \\
\end{longtable}
}

\begin{center}\rule{0.5\linewidth}{0.5pt}\end{center}

\subsection{Gluconeogenesis}\label{gluconeogenesis}

Gluconeogenesis synthesises glucose from non-carbohydrate precursors
(lactate, glycerol, glucogenic amino acids). It occurs primarily in the
\textbf{liver} and to a lesser extent in the \textbf{kidney cortex}.

Gluconeogenesis is NOT simply the reverse of glycolysis. Three
irreversible glycolytic steps require \textbf{bypass enzymes}:

{\def\LTcaptype{none} % do not increment counter
\begin{longtable}[]{@{}
  >{\raggedright\arraybackslash}p{(\linewidth - 4\tabcolsep) * \real{0.5172}}
  >{\raggedright\arraybackslash}p{(\linewidth - 4\tabcolsep) * \real{0.3448}}
  >{\raggedright\arraybackslash}p{(\linewidth - 4\tabcolsep) * \real{0.1379}}@{}}
\toprule\noalign{}
\begin{minipage}[b]{\linewidth}\raggedright
Glycolytic step (irreversible)
\end{minipage} & \begin{minipage}[b]{\linewidth}\raggedright
Gluconeogenic bypass
\end{minipage} & \begin{minipage}[b]{\linewidth}\raggedright
Location
\end{minipage} \\
\midrule\noalign{}
\endhead
\bottomrule\noalign{}
\endlastfoot
Pyruvate kinase (PEP \texorpdfstring{\ensuremath{\to}}{->} Pyruvate) & \textbf{Pyruvate carboxylase}
(Pyruvate \texorpdfstring{\ensuremath{\to}}{->} OAA; biotin, ATP) + \textbf{PEPCK} (OAA \texorpdfstring{\ensuremath{\to}}{->} PEP; GTP) &
Mitochondria + cytoplasm \\
PFK-1 (F6P \texorpdfstring{\ensuremath{\to}}{->} F1,6BP) & \textbf{Fructose-1,6-bisphosphatase (FBPase-1)} &
Cytoplasm \\
Hexokinase (Glucose \texorpdfstring{\ensuremath{\to}}{->} G6P) & \textbf{Glucose-6-phosphatase (G6Pase)} &
ER membrane (liver, kidney primarily) \\
\end{longtable}
}

\textbf{Cost:} 6 ATP equivalents per glucose synthesised (4 ATP + 2
GTP).

\textbf{Regulation of gluconeogenesis:} - \textbf{Allosteric:} Pyruvate
carboxylase activated by acetyl-CoA; FBPase-1 inhibited by AMP and
F2,6BP - \textbf{Hormonal:} Glucagon (via CREB) induces PEPCK and G6Pase
genes; insulin (via FoxO1) represses them - \textbf{Reciprocal with
glycolysis:} F2,6BP activates PFK-1 (glycolysis) and inhibits FBPase-1
(gluconeogenesis), ensuring they do not run simultaneously

\begin{quote}
\textbf{Concept Check 3:} Ethanol metabolism in the liver generates
excess NADH (from alcohol dehydrogenase and aldehyde dehydrogenase).
Predict the effect of heavy alcohol consumption on gluconeogenesis and
the risk of hypoglycaemia.
\end{quote}

\begin{center}\rule{0.5\linewidth}{0.5pt}\end{center}

\subsection{Glycogen Metabolism}\label{glycogen-metabolism}

\subsubsection{Glycogen Structure}\label{glycogen-structure}

Glycogen is a branched polymer of glucose units: - \textbf{Alpha-1,4
linkages:} linear chains (about 8--12 glucose units between
branch points) - \textbf{Alpha-1,6 linkages:} branch points
(about  every 8--12 residues) - \textbf{Glycogenin:}
self-glucosylating primer protein at the core

Liver glycogen (about 100 g): blood glucose
\hyperref[gl:buffer]{\textbf{buffer}}. Muscle glycogen
(about 400 g): local fuel for contraction (rarely exported as
glucose because muscle lacks G6Pase).

\subsubsection{Key Enzymes}\label{key-enzymes}

\textbf{Glycogen phosphorylase:} Cleaves alpha-1,4 linkages from
non-reducing ends, releasing glucose-1-phosphate. Regulated by: -
\textbf{Allosteric activation:} AMP (muscle); glucose (liver,
inhibitory) - \textbf{Covalent modification:} Phosphorylase kinase
phosphorylates Ser14 \texorpdfstring{\ensuremath{\to}}{->} active \textbf{a} form - \textbf{Hormonal
cascade:} Glucagon/epinephrine \texorpdfstring{\ensuremath{\to}}{->} cAMP \texorpdfstring{\ensuremath{\to}}{->} PKA \texorpdfstring{\ensuremath{\to}}{->} phosphorylase kinase \texorpdfstring{\ensuremath{\to}}{->}
phosphorylase a

\textbf{Glycogen synthase:} Adds UDP-glucose to non-reducing ends
(alpha-1,4 linkages). Active when dephosphorylated (GSK3 phosphorylates
\texorpdfstring{\ensuremath{\to}}{->} inactive; insulin inhibits GSK3 \texorpdfstring{\ensuremath{\to}}{->} GS active).

\textbf{Branching enzyme:} Transfers about 7-residue segments
to create alpha-1,6 branches. Increases solubility and number of
non-reducing ends (more rapid mobilisation).

\textbf{Debranching enzyme:} Bifunctional: transferase activity moves 3
residues from branch, then alpha-1,6-glucosidase cleaves the branch
point, releasing free glucose.

\begin{quote}
\textbf{Clinical Connection: Glycogen Storage Diseases} |{} Type
|{} Enzyme deficiency |{} Features |{} |{}
---- |{} ----------------- |{} -------- |{}
|{} I (von Gierke) |{} Glucose-6-phosphatase |{}
Severe hypoglycaemia, hepatomegaly, lactic acidosis |{}
|{} II (Pompe) |{} Lysosomal alpha-glucosidase |{}
Cardiac/skeletal myopathy; treated by ERT |{} |{} V
(McArdle) |{} Muscle glycogen phosphorylase |{} Exercise
intolerance; ``second wind'' phenomenon |{} |{} VII
(Tarui) |{} Muscle PFK |{} Exercise intolerance,
haemolytic anaemia |{}
\end{quote}

\begin{center}\rule{0.5\linewidth}{0.5pt}\end{center}

\subsection{Fatty Acid Oxidation
(Beta-Oxidation)}\label{fatty-acid-oxidation-beta-oxidation}

Beta-oxidation occurs in the \textbf{mitochondrial matrix} and degrades
fatty acyl-CoA by removing 2-carbon units (acetyl-CoA) per cycle.

\subsubsection{Activation and Transport}\label{activation-and-transport}

\begin{enumerate}
\def\labelenumi{\arabic{enumi}.}
\tightlist
\item
  \textbf{Activation:} Fatty acid + CoA + ATP \texorpdfstring{\ensuremath{\to}}{->} Acyl-CoA + AMP +
  PP\(_i\) (acyl-CoA synthetase; outer mitochondrial membrane; costs 2
  ATP equivalents)
\item
  \textbf{Carnitine shuttle:} Acyl-CoA cannot cross the IMM. Instead:

  \begin{itemize}
  \tightlist
  \item
    \textbf{CPT-1} (outer face of IMM): acyl-CoA + carnitine \texorpdfstring{\ensuremath{\to}}{->}
    acylcarnitine + CoA
  \item
    \textbf{Translocase (CACT):} acylcarnitine crosses IMM
  \item
    \textbf{CPT-2} (inner face of IMM): acylcarnitine + CoA \texorpdfstring{\ensuremath{\to}}{->} acyl-CoA +
    carnitine
  \end{itemize}
\end{enumerate}

\textbf{CPT-1 is the rate-limiting step} of beta-oxidation.
\textbf{Malonyl-CoA} (the first committed intermediate of fatty acid
synthesis) \textbf{allosterically inhibits CPT-1}, ensuring that fatty
acid synthesis and oxidation do not occur simultaneously.

\subsubsection{Beta-Oxidation Steps (Each
Cycle)}\label{beta-oxidation-steps-each-cycle}

{\def\LTcaptype{none} % do not increment counter
\begin{longtable}[]{@{}
  >{\raggedright\arraybackslash}p{(\linewidth - 6\tabcolsep) * \real{0.1600}}
  >{\raggedright\arraybackslash}p{(\linewidth - 6\tabcolsep) * \real{0.2400}}
  >{\raggedright\arraybackslash}p{(\linewidth - 6\tabcolsep) * \real{0.3200}}
  >{\raggedright\arraybackslash}p{(\linewidth - 6\tabcolsep) * \real{0.2800}}@{}}
\toprule\noalign{}
\begin{minipage}[b]{\linewidth}\raggedright
Step
\end{minipage} & \begin{minipage}[b]{\linewidth}\raggedright
Enzyme
\end{minipage} & \begin{minipage}[b]{\linewidth}\raggedright
Reaction
\end{minipage} & \begin{minipage}[b]{\linewidth}\raggedright
Product
\end{minipage} \\
\midrule\noalign{}
\endhead
\bottomrule\noalign{}
\endlastfoot
1 & Acyl-CoA dehydrogenase & Acyl-CoA \texorpdfstring{\ensuremath{\to}}{->} trans-enoyl-CoA &
\textbf{FADH\(_2\)} \\
2 & Enoyl-CoA hydratase & trans-enoyl-CoA + H\(_2\)O \texorpdfstring{\ensuremath{\to}}{->}
L-3-hydroxyacyl-CoA & -- \\
3 & 3-Hydroxyacyl-CoA DH & L-3-hydroxyacyl-CoA \texorpdfstring{\ensuremath{\to}}{->} 3-ketoacyl-CoA &
\textbf{NADH} \\
4 & Thiolase & 3-ketoacyl-CoA + CoA \texorpdfstring{\ensuremath{\to}}{->} acyl-CoA(n-2) + acetyl-CoA &
\textbf{Acetyl-CoA} \\
\end{longtable}
}

Each cycle: removes 2 carbons as acetyl-CoA, produces 1 FADH\(_2\) + 1
NADH.

\subsubsection{ATP Yield from Palmitoyl-CoA
(C16:0)}\label{atp-yield-from-palmitoyl-coa-c160}

Summing the products of complete fatty-acid oxidation makes the ``fat is
energy-dense'' claim quantitative. Palmitoyl-CoA (16 carbons) undergoes
7 cycles of beta-oxidation:

{\def\LTcaptype{none} % do not increment counter
\begin{longtable}[]{@{}
  >{\raggedright\arraybackslash}p{(\linewidth - 6\tabcolsep) * \real{0.1944}}
  >{\raggedright\arraybackslash}p{(\linewidth - 6\tabcolsep) * \real{0.2222}}
  >{\raggedright\arraybackslash}p{(\linewidth - 6\tabcolsep) * \real{0.3333}}
  >{\raggedright\arraybackslash}p{(\linewidth - 6\tabcolsep) * \real{0.2500}}@{}}
\toprule\noalign{}
\begin{minipage}[b]{\linewidth}\raggedright
Product
\end{minipage} & \begin{minipage}[b]{\linewidth}\raggedright
Quantity
\end{minipage} & \begin{minipage}[b]{\linewidth}\raggedright
ATP per unit
\end{minipage} & \begin{minipage}[b]{\linewidth}\raggedright
Total ATP
\end{minipage} \\
\midrule\noalign{}
\endhead
\bottomrule\noalign{}
\endlastfoot
Acetyl-CoA (via TCA) & 8 & 10 (3 NADH x 2.5 + 1 FADH\(_2\) x 1.5 + 1
GTP) & 80 \\
FADH\(_2\) (beta-oxidation) & 7 & 1.5 & 10.5 \\
NADH (beta-oxidation) & 7 & 2.5 & 17.5 \\
Activation cost & & -2 & -2 \\
\textbf{Net total} & & & \textbf{106} \\
\end{longtable}
}

Some textbooks report 108 ATP (using older P/O ratios) or 129 ATP
(including palmitate rather than palmitoyl-CoA, and using maximal P/O
ratios). The exact number depends on assumed H\(^+\)/ATP stoichiometry
and shuttle costs. The key insight: fatty acids are energy-dense
(about 9 kcal/g vs.~4 kcal/g for carbohydrates) because they
are highly reduced.

\begin{quote}
\textbf{Clinical Connection: MCAD Deficiency} Medium-chain acyl-CoA
dehydrogenase (MCAD) deficiency is the most common inherited disorder of
fatty acid oxidation (about 1:10,000 births). During fasting,
patients cannot oxidise medium-chain fatty acids, leading to hypoketotic
hypoglycaemia (low ketones because acetyl-CoA production is impaired,
low glucose because gluconeogenesis lacks acetyl-CoA stimulation of
pyruvate carboxylase). Fatal Reye-like episodes can occur. Management:
avoid prolonged fasting; ensure regular carbohydrate intake.
\end{quote}

\begin{center}\rule{0.5\linewidth}{0.5pt}\end{center}

\subsection{Fatty Acid Synthesis}\label{fatty-acid-synthesis}

Fatty acid synthesis occurs in the \textbf{cytoplasm} and uses
fundamentally different chemistry from beta-oxidation:

{\def\LTcaptype{none} % do not increment counter
\begin{longtable}[]{@{}
  >{\raggedright\arraybackslash}p{(\linewidth - 4\tabcolsep) * \real{0.1707}}
  >{\raggedright\arraybackslash}p{(\linewidth - 4\tabcolsep) * \real{0.3415}}
  >{\raggedright\arraybackslash}p{(\linewidth - 4\tabcolsep) * \real{0.4878}}@{}}
\toprule\noalign{}
\begin{minipage}[b]{\linewidth}\raggedright
Feature
\end{minipage} & \begin{minipage}[b]{\linewidth}\raggedright
Beta-oxidation
\end{minipage} & \begin{minipage}[b]{\linewidth}\raggedright
Fatty acid synthesis
\end{minipage} \\
\midrule\noalign{}
\endhead
\bottomrule\noalign{}
\endlastfoot
Location & Mitochondrial matrix & Cytoplasm \\
Carrier & CoA & ACP (acyl carrier protein) \\
Reducing agent & FAD, NAD\(^+\) (oxidised) & NADPH (reduced) \\
2C unit & Acetyl-CoA (removed) & Malonyl-CoA (added) \\
Key enzyme & Multiple separate enzymes & \textbf{FAS} (fatty acid
synthase, one large multifunctional enzyme in animals) \\
Regulation & Inhibited by malonyl-CoA & Activated by citrate; inhibited
by palmitoyl-CoA \\
\end{longtable}
}

\textbf{Acetyl-CoA carboxylase (ACC)} catalyses the committed step:
acetyl-CoA + CO\(_2\) + ATP \texorpdfstring{\ensuremath{\to}}{->} malonyl-CoA. This is the key regulatory
enzyme: - \textbf{Activated by:} citrate (allosteric), insulin
(dephosphorylation) - \textbf{Inhibited by:} palmitoyl-CoA (product
inhibition), AMPK phosphorylation (energy depletion), glucagon (via PKA)

\begin{center}\rule{0.5\linewidth}{0.5pt}\end{center}

\subsection{Ketone Body Metabolism}\label{ketone-body-metabolism}

During prolonged fasting (\textgreater24--48 h), the liver converts
excess acetyl-CoA from fatty acid oxidation into \textbf{ketone bodies}:
acetoacetate, beta-hydroxybutyrate, and acetone.

\textbf{Synthesis (liver mitochondria):} 1. 2 Acetyl-CoA \texorpdfstring{\ensuremath{\to}}{->}
Acetoacetyl-CoA (thiolase) 2. Acetoacetyl-CoA + Acetyl-CoA \texorpdfstring{\ensuremath{\to}}{->} HMG-CoA
(HMG-CoA synthase) 3. HMG-CoA \texorpdfstring{\ensuremath{\to}}{->} Acetoacetate + Acetyl-CoA (HMG-CoA
lyase) 4. Acetoacetate \texorpdfstring{\ensuremath{\to}}{->} beta-Hydroxybutyrate (beta-hydroxybutyrate DH;
uses NADH) 5. Acetoacetate \texorpdfstring{\ensuremath{\to}}{->} Acetone (non-enzymatic decarboxylation;
exhaled)

\textbf{Utilisation (brain, heart, muscle):} Beta-hydroxybutyrate \texorpdfstring{\ensuremath{\to}}{->}
acetoacetate \texorpdfstring{\ensuremath{\to}}{->} acetoacetyl-CoA (via succinyl-CoA transferase, also
called thiophorase) \texorpdfstring{\ensuremath{\to}}{->} 2 acetyl-CoA \texorpdfstring{\ensuremath{\to}}{->} TCA cycle.

\textbf{Key points:} - The liver \textbf{cannot use} ketone bodies
(lacks thiophorase) - The brain normally uses primarily glucose but
during starvation adapts to derive \textbf{about 60--70\%} of
its energy from ketone bodies, sparing muscle protein - This
protein-sparing effect is essential for survival during prolonged
fasting - Diabetic ketoacidosis (DKA): uncontrolled type 1 diabetes \texorpdfstring{\ensuremath{\to}}{->}
absent insulin \texorpdfstring{\ensuremath{\to}}{->} massive lipolysis \texorpdfstring{\ensuremath{\to}}{->} excessive ketone body production \texorpdfstring{\ensuremath{\to}}{->}
metabolic acidosis (\hyperref[gl:ph]{\textbf{pH}} \textless7.3),
dehydration, potentially fatal

\begin{quote}
\textbf{Concept Check 4:} Why can the liver produce ketone bodies but
not use them? What is the physiological advantage of this arrangement?
\end{quote}

\begin{center}\rule{0.5\linewidth}{0.5pt}\end{center}

\subsection{AMPK --- The Cellular Energy
Gauge}\label{ampk-the-cellular-energy-gauge}

\textbf{AMPK} (AMP-activated protein kinase) is a heterotrimer with
\textbf{two catalytic \texorpdfstring{\ensuremath{\alpha}}{alpha}-subunits} (\texorpdfstring{\ensuremath{\alpha}}{alpha}1 and \texorpdfstring{\ensuremath{\alpha}}{alpha}2 isoforms), \textbf{two
regulatory \texorpdfstring{\ensuremath{\beta}}{beta}-subunits} (\texorpdfstring{\ensuremath{\beta}}{beta}1, \texorpdfstring{\ensuremath{\beta}}{beta}2), and the \textbf{\texorpdfstring{\ensuremath{\gamma}}{gamma}-subunit} that binds
adenine nucleotides --- though the assembled holoenzyme is
\texorpdfstring{\ensuremath{\alpha}}{alpha}\(_1\)\texorpdfstring{\ensuremath{\beta}}{beta}\(_1\)\texorpdfstring{\ensuremath{\gamma}}{gamma}\(_1\). (The classic textbook description is that AMPK is
``two catalytic + two regulatory subunits'' referring to the assembly's
binary modules: the catalytic \texorpdfstring{\ensuremath{\alpha}}{alpha}/\texorpdfstring{\ensuremath{\beta}}{beta} pair and the regulatory \texorpdfstring{\ensuremath{\beta}}{beta}/\texorpdfstring{\ensuremath{\gamma}}{gamma} pair.) The
\texorpdfstring{\ensuremath{\gamma}}{gamma}-subunit contains four CBS (cystathionine \texorpdfstring{\ensuremath{\beta}}{beta}-synthase) domains, of which
three bind adenine nucleotides --- AMP, ADP, and ATP --- competitively.
This competitive binding lets the \texorpdfstring{\ensuremath{\gamma}}{gamma}-subunit \emph{integrate} the entire
adenylate pool, not just one ratio.

\begin{figure}[htbp]
\centering
\includegraphics[width=0.82\linewidth,height=4.2in,keepaspectratio]{../figures/mermaid_inline/inline_mermaid_0065_0f9add5e4b9c.png}
\caption{AMPK as the cellular energy sensor. Low energy activates AMPK, which simultaneously promotes catabolic pathways and inhibits anabolic pathways. High energy inactivates AMPK, allowing mTORC1-driven anabolism.}
\end{figure}

\subsubsection{AMPK Activation: LKB1 vs
CaMKK\texorpdfstring{\ensuremath{\beta}}{beta}}\label{ampk-activation-lkb1-vs-camkkux3b2}

AMPK is activated when its catalytic \texorpdfstring{\ensuremath{\alpha}}{alpha}-subunit is phosphorylated at
\textbf{Thr172}. Two upstream kinases phosphorylate this site under
different physiological conditions:

{\def\LTcaptype{none} % do not increment counter
\begin{longtable}[]{@{}
  >{\raggedright\arraybackslash}p{(\linewidth - 6\tabcolsep) * \real{0.2381}}
  >{\raggedright\arraybackslash}p{(\linewidth - 6\tabcolsep) * \real{0.1111}}
  >{\raggedright\arraybackslash}p{(\linewidth - 6\tabcolsep) * \real{0.3016}}
  >{\raggedright\arraybackslash}p{(\linewidth - 6\tabcolsep) * \real{0.3492}}@{}}
\toprule\noalign{}
\begin{minipage}[b]{\linewidth}\raggedright
Upstream kinase
\end{minipage} & \begin{minipage}[b]{\linewidth}\raggedright
Trigger
\end{minipage} & \begin{minipage}[b]{\linewidth}\raggedright
Tissue distribution
\end{minipage} & \begin{minipage}[b]{\linewidth}\raggedright
Pathological connection
\end{minipage} \\
\midrule\noalign{}
\endhead
\bottomrule\noalign{}
\endlastfoot
\textbf{LKB1} (STK11) & Energy stress (high AMP); allosterically primed
by AMP binding to AMPK \texorpdfstring{\ensuremath{\gamma}}{gamma} & Ubiquitous; constitutively active &
Loss-of-function = Peutz-Jeghers syndrome, predisposes to
gastrointestinal cancers \\
\textbf{CaMKK\texorpdfstring{\ensuremath{\beta}}{beta}} & Elevated cytosolic Ca\(^{2+}\) (independent of AMP) &
Brain, muscle, T cells, endothelium & Important in T-cell activation,
osmotic stress \\
\textbf{TAK1} & Inflammatory cytokines (TGF-\texorpdfstring{\ensuremath{\beta}}{beta}, TNF-\texorpdfstring{\ensuremath{\alpha}}{alpha}) & Limited &
Crosstalk between inflammation and metabolism \\
\end{longtable}
}

\textbf{LKB1 is the dominant kinase under classical energy stress}
(exercise, hypoxia, starvation): AMP binding to AMPK \texorpdfstring{\ensuremath{\gamma}}{gamma} promotes LKB1
docking and Thr172 phosphorylation, while simultaneously protecting
Thr172 from dephosphorylation by PP2A and PP2C. AMP is therefore not
simply an ``activator'' but an \emph{allosteric primer} that shifts the
AMPK steady state from dephosphorylated to phosphorylated.

\textbf{CaMKK\texorpdfstring{\ensuremath{\beta}}{beta} provides a Ca\(^{2+}\)-coupled activation route} --- for
example, in muscle contraction (where Ca\(^{2+}\) release accompanies
ATP demand) or in T-cell activation following T-cell receptor
engagement. This route runs \textbf{in parallel} to the AMP/LKB1 pathway
and lets cells anticipate energy demand from Ca\(^{2+}\) signals before
AMP rises.

\subsubsection{Energy Charge}\label{energy-charge}

The \textbf{energy charge (EC)} of the adenylate system was introduced
by Atkinson \citep{atkinson1968} as a single scalar that captures the
fractional saturation of the adenylate pool with high-energy phosphate:

\begin{equation}
\text{E.C.} = \frac{[\text{ATP}] + \tfrac{1}{2}[\text{ADP}]}{[\text{ATP}] + [\text{ADP}] + [\text{AMP}]}
\label{eq:unit_III_energy_charge}
\end{equation}

The numerator counts each phosphoanhydride bond (ATP has two
phosphoanhydrides, ADP has one), while the denominator counts the total
adenylate pool. So E.C. = 1.0 means most adenylates are ATP; E.C. = 0
means most are AMP.

Most cells maintain EC approximately 0.85--0.95 \citep{atkinson1968}.
AMPK is highly sensitive in the range EC = 0.70--0.85 --- precisely the
window where rising AMP signals impending energy crisis.

\subsection{Worked Example: Energy
Charge}\label{worked-example-energy-charge}

\emph{Problem:} A resting skeletal-muscle biopsy gives adenylate pools
of {[}ATP{]} = 5.0 \texorpdfstring{\ensuremath{\mu}}{mu}mol/g, {[}ADP{]} = 0.6 \texorpdfstring{\ensuremath{\mu}}{mu}mol/g, and {[}AMP{]} = 0.05
\texorpdfstring{\ensuremath{\mu}}{mu}mol/g. Calculate the energy charge and predict whether AMPK is active.

\emph{Solution:}

\textbf{Step 1 --- State the relation and why ADP is half-weighted.} The
energy charge measures the fractional saturation of the adenylate pool
with transferable phosphoanhydride bonds:

\[EC = \frac{[\text{ATP}] + \tfrac{1}{2}[\text{ADP}]}{[\text{ATP}] + [\text{ADP}] + [\text{AMP}]}\]

ATP carries two high-energy phosphoanhydride bonds and ADP carries one,
so on a per-adenylate basis ADP holds exactly half the transferable
phosphate of ATP --- hence the \(\tfrac{1}{2}\) weighting. The
denominator is the conserved total adenylate pool
({[}ATP{]}+{[}ADP{]}+{[}AMP{]}), so \(EC = 1\) when every adenylate is
ATP and \(EC = 0\) when every adenylate is AMP.

\textbf{Step 2 --- Substitute the measured \texorpdfstring{\ensuremath{\mu}}{mu}mol/g values.}

\[EC = \frac{5.0 + \tfrac{1}{2}(0.6)}{5.0 + 0.6 + 0.05}\]

Numerator: \(5.0 + 0.5 \times 0.6 = 5.0 + 0.30 = 5.30\). Denominator:
\(5.0 + 0.6 + 0.05 = 5.65\). The shared \texorpdfstring{\ensuremath{\mu}}{mu}mol/g units cancel, so the
ratio is dimensionless.

\textbf{Step 3 --- Compute and interpret against the AMPK threshold.}

\[EC = \frac{5.30}{5.65} = 0.938 \tag{11.3}\]

At \(EC = 0.938\) the pool sits in the normal resting band of 0.85--0.95
and well above the AMPK-activation window of 0.70--0.85. Predict: AMPK
is largely inactive and the cell is in an energy-replete state favouring
anabolism. A fall toward \(EC \approx 0.80\) (rising AMP) would move the
pool into the AMPK-sensitive window and switch the cell toward
catabolism.

\begin{quote}
\textbf{Clinical Connection: Metformin and AMPK} Metformin, the most
widely prescribed diabetes drug, activates AMPK indirectly by inhibiting
mitochondrial Complex I, which raises the AMP/ATP ratio. AMPK
activation: - Increases glucose uptake in muscle - Inhibits hepatic
gluconeogenesis (via CREB coactivator CRTC2 phosphorylation) - Reduces
lipogenesis - May have anti-cancer effects (mTORC1 inhibition) Metformin
is also being investigated as an anti-ageing compound (TAME trial). see
\cref{sec:unit_III_bioenergetics_and_respiration} for Complex I.
\end{quote}

\begin{center}\rule{0.5\linewidth}{0.5pt}\end{center}

\subsection{Worked Example: Flux Control
Coefficient}\label{worked-example-flux-control-coefficient}

\emph{Problem:} In an isolated hepatocyte preparation, a small activator
raises PFK-1 activity by 20\% and the measured glycolytic flux rises by
5\%. A separate manipulation raises the pyruvate-transporter activity by
30\% and flux rises by 3\%. Calculate each flux control coefficient and
determine which step exerts more control over glycolytic flux.

\emph{Solution:}

\textbf{Step 1 --- State the operational form of the flux control
coefficient.} For a small fractional change in enzyme activity,
\(C_i^J = \dfrac{\partial \ln J}{\partial \ln e_i}\) is approximated by
the ratio of fractional changes:

\[C_i^J \approx \frac{\Delta J / J}{\Delta e_i / e_i} \tag{11.9}\]

\textbf{Step 2 --- Substitute for PFK-1.} Here \(\Delta e/e = 0.20\) and
\(\Delta J/J = 0.05\):

\[C_{\text{PFK-1}}^{J} = \frac{0.05}{0.20} = 0.25 \tag{11.10}\]

\textbf{Step 3 --- Substitute for the pyruvate transporter.} Here
\(\Delta e/e = 0.30\) and \(\Delta J/J = 0.03\):

\[C_{\text{transporter}}^{J} = \frac{0.03}{0.30} = 0.10\]

\textbf{Step 4 --- Interpret against the summation theorem.} Both
coefficients fall in the \(C_i^J \approx 0.1\)--\(0.3\) band reported
for glycolytic enzymes. PFK-1 (\(C^J = 0.25\)) exerts more than twice
the flux control of the pyruvate transporter (\(C^J = 0.10\)), so PFK-1
is the more rate-controlling step here --- yet because
\cref{eq:unit_III_mca_summation} requires \(\sum_i C_i^J = 1\), these
two steps together account for primarily \(0.25 + 0.10 = 0.35\) of total
control. The remaining \(0.65\) is distributed across other enzymes: no
single enzyme is \emph{the} rate-limiting step, which is why an
activator targeting PFK-1 alone produces primarily a fractional, not
proportional, increase in flux.

\begin{center}\rule{0.5\linewidth}{0.5pt}\end{center}

\subsection{mTORC1 --- The Growth Integrator and Nutrient
Sensor}\label{mtorc1-the-growth-integrator-and-nutrient-sensor}

\textbf{mTORC1} (mechanistic Target of Rapamycin Complex 1) integrates
amino acid, glucose, oxygen, and growth factor signals:

{\def\LTcaptype{none} % do not increment counter
\begin{longtable}[]{@{}
  >{\raggedright\arraybackslash}p{(\linewidth - 4\tabcolsep) * \real{0.1724}}
  >{\raggedright\arraybackslash}p{(\linewidth - 4\tabcolsep) * \real{0.5862}}
  >{\raggedright\arraybackslash}p{(\linewidth - 4\tabcolsep) * \real{0.2414}}@{}}
\toprule\noalign{}
\begin{minipage}[b]{\linewidth}\raggedright
Input
\end{minipage} & \begin{minipage}[b]{\linewidth}\raggedright
Sensing mechanism
\end{minipage} & \begin{minipage}[b]{\linewidth}\raggedright
Outcome
\end{minipage} \\
\midrule\noalign{}
\endhead
\bottomrule\noalign{}
\endlastfoot
Amino acids (Leu, Arg) & Rag GTPases (Sestrin2 senses Leu; CASTOR1
senses Arg) recruit mTORC1 to lysosomal surface & mTORC1 ON \\
Insulin/IGF-1 & PI3K \texorpdfstring{\ensuremath{\to}}{->} Akt \texorpdfstring{\ensuremath{\to}}{->} phosphorylates TSC2 \texorpdfstring{\ensuremath{\to}}{->} Rheb-GTP accumulates
& mTORC1 ON \\
Glucose & Low glucose \texorpdfstring{\ensuremath{\to}}{->} AMPK \texorpdfstring{\ensuremath{\to}}{->} TSC2/Raptor phosphorylation & mTORC1
OFF \\
Oxygen & Hypoxia \texorpdfstring{\ensuremath{\to}}{->} REDD1 \texorpdfstring{\ensuremath{\to}}{->} TSC1/2 activation & mTORC1 OFF \\
\end{longtable}
}

\subsubsection{Amino Acid Sensing: The Ragulator Complex and
Lysosome}\label{amino-acid-sensing-the-ragulator-complex-and-lysosome}

mTORC1 primarily senses amino acid sufficiency by being
\textbf{physically positioned} on the lysosomal surface. The structure
responsible is a multi-protein scaffold called the \textbf{Ragulator
complex} (LAMTOR1--5) that anchors mTORC1 to the cytoplasmic face of the
lysosomal membrane. The amino acid sensing cascade works as follows:

\begin{enumerate}
\def\labelenumi{\arabic{enumi}.}
\tightlist
\item
  \textbf{Specific amino acid sensors} detect intracellular
  concentrations:

  \begin{itemize}
  \tightlist
  \item
    \textbf{Sestrin2} binds and is inhibited by \textbf{leucine} ---
    when Leu is high, Sestrin2 is displaced and stops inhibiting GATOR2.
  \item
    \textbf{CASTOR1} binds and is inhibited by \textbf{arginine} ---
    same logic.
  \item
    \textbf{SAMTOR} senses S-adenosylmethionine (a proxy for
    methionine).
  \item
    \textbf{SLC38A9} senses lysosomal arginine and signals from the
    lysosomal membrane.
  \end{itemize}
\item
  \textbf{GATOR1/2 GAP/GEF complex.} GATOR1 is a GAP for the Rag GTPases
  (turning them off); GATOR2 inhibits GATOR1. When sensors release
  inhibition of GATOR2, GATOR1 is suppressed, allowing the Rag GTPases
  to enter their active conformation.
\item
  \textbf{Active Rag GTPases} (RagA/B-GTP + RagC/D-GDP) bind Raptor (a
  subunit of mTORC1) and recruit mTORC1 to the lysosomal surface.
\item
  \textbf{At the lysosome}, mTORC1 encounters Rheb-GTP --- provided that
  growth factor signalling has activated it via Akt \texorpdfstring{\ensuremath{\to}}{->} TSC2 inhibition \texorpdfstring{\ensuremath{\to}}{->}
  Rheb-GTP accumulation.
\item
  \textbf{Rheb-GTP fully activates mTORC1} kinase activity. Both signals
  (amino acids \texorpdfstring{\ensuremath{\to}}{->} Rag \texorpdfstring{\ensuremath{\to}}{->} lysosomal positioning \emph{and} growth factors \texorpdfstring{\ensuremath{\to}}{->}
  Rheb-GTP) must converge.
\end{enumerate}

This dual-input gating is why mTORC1 is sometimes described as a
\textbf{coincidence detector}: it primarily signals ``build'' when both
raw materials (amino acids) and external commands (insulin/IGF-1) are
present.

\subsubsection{Rapamycin: Mechanism of
Action}\label{rapamycin-mechanism-of-action}

\textbf{Rapamycin} is a macrolide originally isolated from
\emph{Streptomyces hygroscopicus} on Easter Island (Rapa Nui). It binds
to the immunophilin \textbf{FKBP12} in the cytoplasm; the
FKBP12-rapamycin complex then binds to the FRB domain of mTOR within
mTORC1, allosterically inhibiting kinase activity \emph{primarily in the
mTORC1 complex} (mTORC2 is largely unaffected on short timescales). This
inhibition: - Suppresses cap-dependent translation (4E-BP1
phosphorylation) - Suppresses ribosome biogenesis (S6K1 phosphorylation)
- Activates autophagy - Extends lifespan in yeast, worms, flies, and
mice (about 10--15\% extension at low doses)

Rapamycin and analogues (rapalogues: everolimus, temsirolimus) are
clinically approved as immunosuppressants (transplant rejection
prevention) and anticancer agents (renal cell carcinoma, certain breast
cancers). The TAME trial and large-scale rapamycin/metformin studies are
testing rapalogues as longevity interventions in humans.

\textbf{AMPK and mTORC1 reciprocal inhibition} creates a
\textbf{bistable metabolic switch:} energy-rich \texorpdfstring{\ensuremath{\to}}{->} mTORC1 on, AMPK off \texorpdfstring{\ensuremath{\to}}{->}
anabolism; energy-depleted \texorpdfstring{\ensuremath{\to}}{->} AMPK on, mTORC1 off \texorpdfstring{\ensuremath{\to}}{->} catabolism +
autophagy.

\begin{center}\rule{0.5\linewidth}{0.5pt}\end{center}

\subsection{The Randle Cycle (Glucose-Fatty Acid
Cycle)}\label{the-randle-cycle-glucose-fatty-acid-cycle}

Proposed by Randle et al.~(1963, \emph{The Lancet}), the Randle cycle
describes \textbf{reciprocal inhibition} between glucose and fatty acid
oxidation in heart and skeletal muscle:

\textbf{When fatty acids are high \texorpdfstring{\ensuremath{\to}}{->} glucose utilisation is suppressed:}
1. Beta-oxidation \texorpdfstring{\ensuremath{\to}}{->} elevated acetyl-CoA \texorpdfstring{\ensuremath{\to}}{->} \textbf{PDH kinase (PDK)
activated} \texorpdfstring{\ensuremath{\to}}{->} PDH phosphorylated/inactivated \texorpdfstring{\ensuremath{\to}}{->} less pyruvate oxidation 2.
Beta-oxidation \texorpdfstring{\ensuremath{\to}}{->} elevated citrate (exported to cytoplasm) \texorpdfstring{\ensuremath{\to}}{->}
\textbf{PFK-1 allosterically inhibited} \texorpdfstring{\ensuremath{\to}}{->} glycolysis slows 3. Elevated
intracellular acyl-CoA \texorpdfstring{\ensuremath{\to}}{->} PKC-theta (in muscle) \texorpdfstring{\ensuremath{\to}}{->} serine phosphorylation
of \textbf{IRS-1} \texorpdfstring{\ensuremath{\to}}{->} impaired insulin signalling \texorpdfstring{\ensuremath{\to}}{->} reduced GLUT4
translocation

\textbf{When glucose is high \texorpdfstring{\ensuremath{\to}}{->} fatty acid oxidation is suppressed:} 1.
Insulin \texorpdfstring{\ensuremath{\to}}{->} ACC activated \texorpdfstring{\ensuremath{\to}}{->} malonyl-CoA synthesis increases \texorpdfstring{\ensuremath{\to}}{->}
\textbf{CPT-1 inhibited} \texorpdfstring{\ensuremath{\to}}{->} beta-oxidation blocked 2. Elevated NADH from
glycolysis \texorpdfstring{\ensuremath{\to}}{->} inhibits beta-oxidation dehydrogenases

\subsubsection{Metabolic Flexibility}\label{metabolic-flexibility}

\textbf{Metabolic flexibility} = the ability to switch between glucose
and fat oxidation in response to fuel availability. Assessed by the
respiratory quotient (RQ):

\[RQ = \frac{\text{CO}_2 \text{ produced}}{\text{O}_2 \text{ consumed}} \tag{11.4}\]

\begin{itemize}
\tightlist
\item
  RQ = 1.0: pure glucose oxidation
\item
  RQ = 0.7: pure fat oxidation
\item
  Healthy individuals: RQ about 0.95 post-meal \texorpdfstring{\ensuremath{\to}}{->}
  about 0.75 fasting (rapid switching)
\item
  Obese/T2D individuals: RQ remains about 0.85 in fasting \texorpdfstring{\ensuremath{\to}}{->}
  ``metabolic inflexibility''
\end{itemize}

\begin{center}\rule{0.5\linewidth}{0.5pt}\end{center}

\subsection{Hormonal Coordination at the Whole-Organism
Level}\label{hormonal-coordination-at-the-whole-organism-level}

\begin{figure}[htbp]
\centering
\includegraphics[width=0.82\linewidth,height=4.2in,keepaspectratio]{../figures/mermaid_inline/inline_mermaid_0066_ccd169edacbf.png}
\caption{Hormonal coordination of metabolism across organs. Insulin (green) drives anabolic responses; glucagon (red) and cortisol (blue) drive catabolic responses. Each organ has distinct metabolic roles and receptor expression patterns.}
\end{figure}

{\def\LTcaptype{none} % do not increment counter
\begin{longtable}[]{@{}
  >{\raggedright\arraybackslash}p{(\linewidth - 8\tabcolsep) * \real{0.2121}}
  >{\raggedright\arraybackslash}p{(\linewidth - 8\tabcolsep) * \real{0.2424}}
  >{\raggedright\arraybackslash}p{(\linewidth - 8\tabcolsep) * \real{0.1515}}
  >{\raggedright\arraybackslash}p{(\linewidth - 8\tabcolsep) * \real{0.1818}}
  >{\raggedright\arraybackslash}p{(\linewidth - 8\tabcolsep) * \real{0.2121}}@{}}
\toprule\noalign{}
\begin{minipage}[b]{\linewidth}\raggedright
Hormone
\end{minipage} & \begin{minipage}[b]{\linewidth}\raggedright
Stimulus
\end{minipage} & \begin{minipage}[b]{\linewidth}\raggedright
Liver
\end{minipage} & \begin{minipage}[b]{\linewidth}\raggedright
Muscle
\end{minipage} & \begin{minipage}[b]{\linewidth}\raggedright
Adipose
\end{minipage} \\
\midrule\noalign{}
\endhead
\bottomrule\noalign{}
\endlastfoot
Insulin & High glucose & Glycogen synthesis; lipogenesis;
gluconeogenesis OFF & GLUT4; glycogen synthesis; protein synthesis &
GLUT4; lipogenesis; lipolysis OFF \\
Glucagon & Low glucose & Glycogenolysis; gluconeogenesis ON; ketogenesis
& Minimal (no glucagon receptor) & Lipolysis (via FFA release) \\
Epinephrine & Stress/exercise & Glycogenolysis; gluconeogenesis &
Glycogenolysis; glucose uptake & Lipolysis (beta\(_3\)-AR \texorpdfstring{\ensuremath{\to}}{->} cAMP \texorpdfstring{\ensuremath{\to}}{->} PKA \texorpdfstring{\ensuremath{\to}}{->}
HSL) \\
Cortisol & Chronic stress & Gluconeogenesis (PEPCK gene induction) &
Protein catabolism \texorpdfstring{\ensuremath{\to}}{->} AAs for GNG & Lipolysis (central redistribution) \\
Leptin & Adiposity signal & Central: inhibits NPY/AgRP \texorpdfstring{\ensuremath{\to}}{->} appetite
suppression & FA oxidation (via AMPK) & Lipogenesis suppressed \\
\end{longtable}
}

\begin{center}\rule{0.5\linewidth}{0.5pt}\end{center}

\subsection{Substrate (Futile) Cycles}\label{substrate-futile-cycles}

A \textbf{substrate cycle} involves two opposing metabolic reactions
operating simultaneously:

\[\text{Fructose-6-P} + \text{ATP} \xrightarrow{\text{PFK-1}} \text{Fructose-1,6-BP} + \text{ADP} \tag{11.5}\]

\[\text{Fructose-1,6-BP} + \text{H}_2\text{O} \xrightarrow{\text{FBPase-1}} \text{Fructose-6-P} + \text{P}_i \tag{11.6}\]

Net reaction: ATP \texorpdfstring{\ensuremath{\to}}{->} ADP + P\(_i\) (energy dissipated as heat)

\textbf{Why substrate cycles are useful (Newsholme hypothesis, 1984):}

\[\text{Amplification factor} = \frac{v_{\text{PFK}} + v_{\text{FBPase}}}{v_{\text{PFK}} - v_{\text{FBPase}}} \tag{11.7}\]

When both enzymes run at 50 U/min each and net flux is 1 U/min, a 10\%
change in PFK produces a 500\% change in net flux --- \textbf{metabolic
rheostat} behaviour. Substrate cycles provide ultrasensitive flux
control at the cost of ATP.

\textbf{Thermogenic substrate cycles:} - \textbf{Brown adipose tissue
(BAT):} UCP1 short-circuits the ETC proton gradient \texorpdfstring{\ensuremath{\to}}{->} heat
(non-shivering thermogenesis) - \textbf{Sarcolipin/SERCA cycle} in
skeletal muscle: sarcolipin uncouples SERCA Ca\(^{2+}\) pumping from
Ca\(^{2+}\) retention, dissipating ATP as heat - \textbf{Insect flight
muscle:} glycolytic futile cycles generate heat rapidly to reach flight
temperature (about 40 degrees C)

\begin{center}\rule{0.5\linewidth}{0.5pt}\end{center}

\subsection{Obesity and Metabolic
Syndrome}\label{obesity-and-metabolic-syndrome}

Metabolic syndrome is a cluster of conditions (central obesity, insulin
resistance, dyslipidaemia, hypertension, hyperglycaemia) that
dramatically increase cardiovascular disease risk.

\subsubsection{Visceral Fat as an Endocrine
Organ}\label{visceral-fat-as-an-endocrine-organ}

The shift in our understanding of adipose tissue --- from ``passive
triglyceride storage depot'' to ``active endocrine organ'' --- is one of
the most consequential paradigm changes in metabolism research of the
last 30 years. \textbf{Visceral adipose tissue (VAT)}, the fat depot
surrounding intra-abdominal organs, behaves very differently from
subcutaneous adipose tissue (SAT):

{\def\LTcaptype{none} % do not increment counter
\begin{longtable}[]{@{}
  >{\raggedright\arraybackslash}p{(\linewidth - 4\tabcolsep) * \real{0.1429}}
  >{\raggedright\arraybackslash}p{(\linewidth - 4\tabcolsep) * \real{0.4643}}
  >{\raggedright\arraybackslash}p{(\linewidth - 4\tabcolsep) * \real{0.3929}}@{}}
\toprule\noalign{}
\begin{minipage}[b]{\linewidth}\raggedright
Property
\end{minipage} & \begin{minipage}[b]{\linewidth}\raggedright
Subcutaneous adipose (SAT)
\end{minipage} & \begin{minipage}[b]{\linewidth}\raggedright
Visceral adipose (VAT)
\end{minipage} \\
\midrule\noalign{}
\endhead
\bottomrule\noalign{}
\endlastfoot
Adipocyte size & Smaller, more uniform & Hypertrophic, heterogeneous \\
Lipolysis sensitivity to catecholamines & Modest & High
(\texorpdfstring{\ensuremath{\beta}}{beta}\(_3\)-AR-driven) \\
Insulin sensitivity & High & Low \\
Drainage & Systemic circulation & \textbf{Hepatic portal vein} (direct
to liver) \\
Macrophage infiltration & Low (fed state) & High (M1-skewed) \\
Adipokine output & Adiponectin \textgreater{} leptin & Leptin
\textgreater{} adiponectin; \texorpdfstring{\ensuremath{\uparrow}}{up} TNF-\texorpdfstring{\ensuremath{\alpha}}{alpha}, IL-6, RBP4, MCP-1 \\
Cardiovascular disease risk & Modest & High \\
\end{longtable}
}

The portal-vein drainage means that VAT lipolytic products (FFAs,
glycerol) and pro-inflammatory cytokines (TNF-\texorpdfstring{\ensuremath{\alpha}}{alpha}, IL-6) hit the liver at
much higher concentrations than they do peripheral tissues --- driving
hepatic insulin resistance, NAFLD, and a chronic low-grade systemic
inflammatory state.

\textbf{Key adipokines:}

{\def\LTcaptype{none} % do not increment counter
\begin{longtable}[]{@{}
  >{\raggedright\arraybackslash}p{(\linewidth - 4\tabcolsep) * \real{0.2571}}
  >{\raggedright\arraybackslash}p{(\linewidth - 4\tabcolsep) * \real{0.5714}}
  >{\raggedright\arraybackslash}p{(\linewidth - 4\tabcolsep) * \real{0.1714}}@{}}
\toprule\noalign{}
\begin{minipage}[b]{\linewidth}\raggedright
Adipokine
\end{minipage} & \begin{minipage}[b]{\linewidth}\raggedright
Direction in obesity
\end{minipage} & \begin{minipage}[b]{\linewidth}\raggedright
Effect
\end{minipage} \\
\midrule\noalign{}
\endhead
\bottomrule\noalign{}
\endlastfoot
\textbf{Leptin} & \texorpdfstring{\ensuremath{\uparrow}}{up} (proportional to fat mass) & Hypothalamic POMC/AgRP
signaling --- appetite suppression; resistance develops with chronic
elevation \\
\textbf{Adiponectin} & \texorpdfstring{\ensuremath{\downarrow}}{down} in obesity & Insulin-sensitizing; AMPK
activator; anti-inflammatory \\
\textbf{TNF-\texorpdfstring{\ensuremath{\alpha}}{alpha}} & \texorpdfstring{\ensuremath{\uparrow}}{up} (from infiltrating macrophages) & IRS-1 Ser
phosphorylation via JNK/IKK-\texorpdfstring{\ensuremath{\beta}}{beta} \\
\textbf{IL-6} & \texorpdfstring{\ensuremath{\uparrow}}{up} & Hepatic CRP induction; insulin resistance in
liver \\
\textbf{MCP-1} & \texorpdfstring{\ensuremath{\uparrow}}{up} & Recruits more macrophages \texorpdfstring{\ensuremath{\to}}{->} vicious cycle \\
\textbf{RBP4} & \texorpdfstring{\ensuremath{\uparrow}}{up} & Promotes hepatic gluconeogenesis \\
\end{longtable}
}

\textbf{Molecular mechanisms of insulin resistance in obesity:} 1.
\textbf{Lipid overflow:} Adipose tissue capacity exceeded \texorpdfstring{\ensuremath{\to}}{->} ectopic
lipid deposition in liver (NAFLD) and muscle 2. \textbf{DAG
accumulation} \texorpdfstring{\ensuremath{\to}}{->} PKC-epsilon (liver) / PKC-theta (muscle) \texorpdfstring{\ensuremath{\to}}{->} IRS-1 Ser
phosphorylation 3. \textbf{Ceramide accumulation} \texorpdfstring{\ensuremath{\to}}{->} PP2A activation \texorpdfstring{\ensuremath{\to}}{->}
Akt dephosphorylation 4. \textbf{Adipose inflammation:} Hypertrophic
adipocytes secrete TNF-alpha, IL-6, MCP-1 \texorpdfstring{\ensuremath{\to}}{->} macrophage infiltration \texorpdfstring{\ensuremath{\to}}{->}
chronic low-grade inflammation \texorpdfstring{\ensuremath{\to}}{->} JNK/IKK-beta activation \texorpdfstring{\ensuremath{\to}}{->} IRS-1 serine
phosphorylation 5. \textbf{ER stress:} Lipid overload overwhelms ER
folding capacity \texorpdfstring{\ensuremath{\to}}{->} UPR \texorpdfstring{\ensuremath{\to}}{->} JNK \texorpdfstring{\ensuremath{\to}}{->} IRS-1 serine phosphorylation 6.
\textbf{Mitochondrial dysfunction:} Incomplete fatty acid oxidation \texorpdfstring{\ensuremath{\to}}{->}
ROS \texorpdfstring{\ensuremath{\to}}{->} oxidative stress \texorpdfstring{\ensuremath{\to}}{->} insulin resistance

\begin{quote}
\textbf{Clinical Connection: GLP-1 Receptor Agonists} Glucagon-like
peptide-1 (GLP-1) receptor agonists (semaglutide, liraglutide,
tirzepatide) have revolutionised treatment of both type 2 diabetes and
obesity. They: - Enhance glucose-dependent insulin secretion - Suppress
glucagon secretion - Slow gastric emptying (satiety) - Act on
hypothalamic appetite centres (weight loss of 15--20\%) - Reduce
cardiovascular events and NAFLD progression Tirzepatide is a dual
GIP/GLP-1 receptor agonist achieving even greater weight loss
(about 22.5\% in trials).
\end{quote}

\begin{quote}
\textbf{Concept Check 5:} Explain why
\hyperref[gl:aerobic]{\textbf{aerobic}} exercise improves insulin
sensitivity even without weight loss. Consider the effects of exercise
on AMPK, GLUT4, mitochondrial biogenesis, and intramyocellular lipid
turnover.
\end{quote}

\begin{center}\rule{0.5\linewidth}{0.5pt}\end{center}

\subsection{Metabolic Control
Analysis}\label{metabolic-control-analysis}

\textbf{Metabolic control analysis (MCA)} provides a quantitative
framework for understanding how control of metabolic flux is distributed
among pathway enzymes \citep{atkinson1968}.

\textbf{Flux control coefficient (\(C_i^J\)):} The fractional change in
pathway flux (\(J\)) caused by a fractional change in enzyme activity
(\(e_i\)), with most other variables held constant:

\[C_i^J = \frac{e_i}{J} \cdot \frac{\partial J}{\partial e_i} = \frac{\partial \ln J}{\partial \ln e_i} \tag{11.8}\]

Each \(C_i^J\) ranges from 0 (this enzyme has no control over flux) to
about 1 (this enzyme has near-total control). Negative values
are possible for enzymes that produce inhibitors of the pathway.

\textbf{Summation theorem:} The combined flux control coefficients in a
pathway sum to 1:

\begin{equation}
\sum_i C_i^J = 1
\label{eq:unit_III_mca_summation}
\end{equation}

This means control is \textbf{shared} among most enzymes, not
concentrated at a single ``rate-limiting step.'' In practice, one or a
few enzymes may have larger \(C_i^J\) values, but the concept of a
single rate-limiting step is an oversimplification.

For example, in glycolysis under physiological conditions, the flux
control coefficients for hexokinase, PFK-1, pyruvate kinase, and the
pyruvate transporter are most in the range
\(C_i^J \approx 0.1\)--\(0.3\) --- no single enzyme dominates, and the
precise distribution shifts with cellular state. This is why drug
therapies that target a single glycolytic enzyme often have surprisingly
modest effects on overall flux: the pathway redistributes control.

\subsubsection{Metabolomics: Measuring the Metabolic
State}\label{metabolomics-measuring-the-metabolic-state}

\textbf{Metabolomics} is the systematic measurement of small-molecule
metabolites --- the experimental complement to MCA, since you cannot
compute flux control coefficients without measuring concentrations and
fluxes.

{\def\LTcaptype{none} % do not increment counter
\begin{longtable}[]{@{}
  >{\raggedright\arraybackslash}p{(\linewidth - 8\tabcolsep) * \real{0.1667}}
  >{\raggedright\arraybackslash}p{(\linewidth - 8\tabcolsep) * \real{0.2292}}
  >{\raggedright\arraybackslash}p{(\linewidth - 8\tabcolsep) * \real{0.2083}}
  >{\raggedright\arraybackslash}p{(\linewidth - 8\tabcolsep) * \real{0.1875}}
  >{\raggedright\arraybackslash}p{(\linewidth - 8\tabcolsep) * \real{0.2083}}@{}}
\toprule\noalign{}
\begin{minipage}[b]{\linewidth}\raggedright
Approach
\end{minipage} & \begin{minipage}[b]{\linewidth}\raggedright
Sensitivity
\end{minipage} & \begin{minipage}[b]{\linewidth}\raggedright
Resolution
\end{minipage} & \begin{minipage}[b]{\linewidth}\raggedright
Strengths
\end{minipage} & \begin{minipage}[b]{\linewidth}\raggedright
Weaknesses
\end{minipage} \\
\midrule\noalign{}
\endhead
\bottomrule\noalign{}
\endlastfoot
\textbf{NMR} (\(^1\)H, \(^{13}\)C, \(^{31}\)P) & \texorpdfstring{\ensuremath{\mu}}{mu}M--mM
(about 10\(^4\) molecules) & Quantitative; structure-resolving
& Non-destructive; absolute quantification; detects isomers & Lower
sensitivity than MS; needs about  mg sample; chemical-shift
overlap in complex mixtures \\
\textbf{LC-MS} (HILIC + reversed phase) & nM--\texorpdfstring{\ensuremath{\mu}}{mu}M
(about 10\(^7\) molecules) & High sensitivity & Targeted and
untargeted; broad metabolite coverage & Ionization bias (different
molecules ionize differently); matrix effects; needs MS\(^2\) for
confident ID \\
\textbf{GC-MS} & nM (with derivatisation) & Very reproducible spectra &
Mature databases (NIST, GMD); quantitative & Requires volatile
derivatives; not most metabolites compatible \\
\textbf{CE-MS} & nM & Polar/charged metabolites & Excellent for
phosphometabolites, amines & Low robustness; less mature databases \\
\textbf{Imaging MS (DESI/MALDI-MSI)} & \texorpdfstring{\ensuremath{\mu}}{mu}M & Spatially resolved &
Tissue-level metabolic mapping & Lower mass accuracy than LC-MS; no
separation step \\
\end{longtable}
}

\textbf{Key challenges:}

\begin{enumerate}
\def\labelenumi{\arabic{enumi}.}
\tightlist
\item
  \textbf{Identification.} A typical untargeted LC-MS run detects
  5,000--30,000 features but about 10--20\% can be confidently
  identified. The remaining ``dark metabolome'' comprises mass spectra
  with no database match --- a major bottleneck.
\item
  \textbf{Compartmentation.} Metabolomics typically homogenises tissue,
  losing information about which metabolites were in mitochondria vs
  cytoplasm vs lysosome. Subcellular metabolomics (genetically encoded
  sensors, organelle isolation) is an emerging frontier.
\item
  \textbf{Flux \texorpdfstring{\ensuremath{\neq}}{!=} concentration.} A high concentration of an intermediate
  may reflect either high flux through that step \emph{or} a downstream
  block. Stable-isotope tracing (\(^{13}\)C-glucose,
  \(^{15}\)N-glutamine) is required to measure fluxes, not
  concentrations.
\end{enumerate}

These data feed directly into MCA models: combining flux measurements
(from \(^{13}\)C tracing) with elasticity coefficients (from in vitro
enzyme kinetics) gives an empirical \(C_i^J\) for each enzyme,
validating or revising textbook claims about which step is
``rate-limiting.''

\begin{quote}
\textbf{Concept Check 6:} A drug increases hexokinase activity by 50\%
in muscle. From fasting biopsies, the flux control coefficient
\(C_{HK}^J\) for hexokinase in glycolysis was estimated at 0.2. Use
\cref{eq:unit_III_mca_summation} to predict the maximum fractional
change in glycolytic flux. Why is the prediction an upper bound, not a
precise value?
\end{quote}

\begin{quote}
\textbf{Concept Check 7:} A patient with chronic alcohol consumption
shows elevated NADH/NAD\(^+\) ratio in cytoplasm and impaired
gluconeogenesis. Using \cref{eq:unit_III_energy_charge}, predict whether
the energy charge is preserved or depressed, and explain why fasting
hypoglycaemia is a clinical risk.
\end{quote}

\begin{center}\rule{0.5\linewidth}{0.5pt}\end{center}

\subsection{Computational Bridge}\label{computational-bridge-10}

Drug effects on pathway enzymes are often framed like competitive
inhibitors shifting apparent \(K_m\). The biochemistry module exposes
both curves numerically:

\begin{Shaded}
\begin{Highlighting}[]
\ImportTok{from}\NormalTok{ biology.biochemistry }\ImportTok{import}\NormalTok{ michaelis\_menten, competitive\_inhibition}

\NormalTok{base }\OperatorTok{=}\NormalTok{ michaelis\_menten(}\FloatTok{4.0}\NormalTok{, Vmax}\OperatorTok{=}\FloatTok{100.0}\NormalTok{, Km}\OperatorTok{=}\FloatTok{5.0}\NormalTok{)}
\NormalTok{blocked }\OperatorTok{=}\NormalTok{ competitive\_inhibition(}\FloatTok{4.0}\NormalTok{, }\FloatTok{100.0}\NormalTok{, }\FloatTok{5.0}\NormalTok{, inhibitor\_conc}\OperatorTok{=}\FloatTok{5.0}\NormalTok{, Ki}\OperatorTok{=}\FloatTok{2.5}\NormalTok{)}
\BuiltInTok{print}\NormalTok{(}\BuiltInTok{round}\NormalTok{(base.reaction\_rate, }\DecValTok{2}\NormalTok{), }\BuiltInTok{round}\NormalTok{(blocked.reaction\_rate, }\DecValTok{2}\NormalTok{))}
\end{Highlighting}
\end{Shaded}

\begin{quote}
\textbf{Clinical / systems note:} Metformin, SGLT2 inhibitors, and GLP-1
agonists reshape the same integrative network (hepatic glucose output,
peripheral uptake, renal glucose loss) that insulin and glucagon command
--- illustrating multi-node control rather than a single
``rate-limiting'' hormone.
\end{quote}

\begin{center}\rule{0.5\linewidth}{0.5pt}\end{center}

\subsection{Current Evidence and Frontier
Biology}\label{current-evidence-and-frontier-biology-13}

For \textbf{Metabolic Integration and Regulation}, frontier biology
belongs inside the evidence logic of the chapter. Metabolism is now
studied as a regulated network constrained by energy, redox balance,
compartmentation, and environment. The core reading question is this:
metabolic integration depends on compartmentation, hormone state,
nutrient availability, and time scale.

\begin{itemize}
\tightlist
\item
  \textbf{What to verify:} identify the observation, model, assay, or
  dataset that would make the claim stronger or weaker.
\item
  \textbf{What to qualify:} state the scale, organism, cell type,
  environmental condition, or population where the claim is expected to
  hold.
\item
  \textbf{What to compare:} test at least one alternative explanation,
  baseline, or null model before treating the pattern as causal.
\item
  \textbf{What to cite:} distinguish primary evidence, review synthesis,
  public dataset, and institutional guidance; for recent or numeric
  claims, prefer the source closest to the measurement and state what
  has changed since it was published.
\end{itemize}

A strong metabolic explanation names the flux, the limiting step, the
sensor, and the condition under which the pathway changes direction or
priority.

\textbf{Source practice:} For metabolic claims, keep the organism,
compartment, energetic state, and measurement method visible; a pathway
map is not enough without flux or concentration evidence.

\subsection{Summary}\label{summary-13}

\begin{itemize}
\tightlist
\item
  \textbf{AMPK} (heterotrimer of two catalytic \texorpdfstring{\ensuremath{\alpha}}{alpha} + two regulatory \texorpdfstring{\ensuremath{\beta}}{beta}/\texorpdfstring{\ensuremath{\gamma}}{gamma}
  modules) is activated by elevated AMP/ATP and ADP/ATP through Thr172
  phosphorylation by \textbf{LKB1} (energy stress) or \textbf{CaMKK\texorpdfstring{\ensuremath{\beta}}{beta}}
  (Ca\(^{2+}\)); it promotes catabolism (\texorpdfstring{\ensuremath{\beta}}{beta}-oxidation, glycolysis) and
  inhibits mTORC1 and anabolism; activates autophagy (via ULK1) and
  mitochondrial biogenesis (via PGC-1\texorpdfstring{\ensuremath{\alpha}}{alpha}).
\item
  \textbf{mTORC1} (activated by insulin + amino acids + oxygen) is
  positioned at the lysosome by Rag GTPases bound to the
  \textbf{Ragulator complex}; full activation requires both amino
  acid--driven lysosomal recruitment \emph{and} growth factor--driven
  Rheb-GTP. \textbf{Rapamycin} (FKBP12 complex) inhibits mTORC1.
\item
  \textbf{Energy charge} \cref{eq:unit_III_energy_charge} captures the
  fractional saturation of the adenylate pool with high-energy
  phosphate; cells maintain E.C. \texorpdfstring{\ensuremath{\approx}}{~} 0.85--0.95.
\item
  \textbf{Insulin} (PI3K/Akt/mTOR pathway) drives GLUT4 translocation,
  glycogen synthesis, lipogenesis, and suppresses gluconeogenesis. The
  phosphorylase kinase / GLUT4-vesicle / 14-3-3-BAD effectors are
  characteristic Akt outputs.
\item
  \textbf{Glucagon} (cAMP/PKA cascade) drives glycogenolysis (via
  phosphorylase kinase), gluconeogenesis (via CREB transcription), and
  ketogenesis. The hormonal cascade amplifies the signal
  about 10\(^7\)-fold from femtomolar hormone to millimolar
  metabolite.
\item
  \textbf{Gluconeogenesis} uses four bypass enzymes (pyruvate
  carboxylase, PEPCK, FBPase-1, G6Pase) and costs 6 ATP equivalents per
  glucose.
\item
  \textbf{Beta-oxidation} yields about 106 ATP per
  palmitoyl-CoA; CPT-1 is rate-limiting and inhibited by malonyl-CoA.
\item
  \textbf{Ketone bodies} (acetoacetate, \texorpdfstring{\ensuremath{\beta}}{beta}-hydroxybutyrate) sustain the
  brain during starvation; liver produces but cannot use them. Fasting
  fuel switching follows a programmed timeline from glycogen \texorpdfstring{\ensuremath{\to}}{->} fatty
  acids \texorpdfstring{\ensuremath{\to}}{->} ketones over hours-to-weeks.
\item
  \textbf{Randle cycle:} fatty acid--glucose reciprocal inhibition;
  disrupted in type 2 diabetes.
\item
  \textbf{Substrate cycles} enable ultrasensitive flux control and
  thermogenesis; amplification factor = (v\(_1\) + v\(_2\))/(v\(_1\) -
  v\(_2\)).
\item
  \textbf{Metabolic flexibility} (RQ switching) is impaired in
  obesity/T2D due to lipid intermediate accumulation and insulin
  resistance.
\item
  \textbf{Metabolic control analysis} (\cref{eq:unit_III_mca_summation})
  shows that flux control is distributed, not concentrated at a single
  rate-limiting step.
\item
  \textbf{Metabolomics} (NMR + LC-MS) provides the empirical basis for
  MCA, with \(^{13}\)C tracing distinguishing flux from concentration.
\item
  \textbf{Visceral adipose} is a metabolically active endocrine organ
  whose portal-vein drainage drives hepatic insulin resistance, NAFLD,
  and chronic systemic inflammation.
\item
  \textbf{Connections:} See \cref{sec:unit_II_cell_signaling} for
  \hyperref[gl:signal-transduction]{\textbf{signal transduction}} to
  mTOR, Unit IX for endocrine first messengers, and
  \cref{sec:unit_I_enzymes_and_kinetics} for enzyme-level kinetics.
\end{itemize}

\begin{center}\rule{0.5\linewidth}{0.5pt}\end{center}

\subsection{Review Questions}\label{review-questions-13}

\begin{enumerate}
\def\labelenumi{\arabic{enumi}.}
\item
  A patient presents with an inherited loss-of-function
  \hyperref[gl:mutation]{\textbf{mutation}} in LKB1 (the primary AMPK
  upstream kinase). Predict four metabolic consequences in skeletal
  muscle during fasting.
\item
  Derive the energy charge equation (\cref{eq:unit_III_energy_charge}).
  If {[}ATP{]} = 3.5 mM, {[}ADP{]} = 0.8 mM, {[}AMP{]} = 0.2 mM,
  calculate EC. Is AMPK likely to be active?
\item
  Describe the insulin signalling cascade from receptor binding to GLUT4
  translocation, including the role of AS160, Rab10, and the SNARE
  complex. At which steps does insulin resistance develop in type 2
  diabetes?
\item
  Compare gluconeogenesis with glycolysis. Identify the three bypass
  reactions and explain why they are necessary.
\item
  Calculate the ATP yield from complete oxidation of palmitoyl-CoA
  (C16:0) to CO\(_2\) and H\(_2\)O, showing most sources of ATP.
\item
  Explain ketone body synthesis and utilisation. Why can the brain use
  ketone bodies but the liver cannot? Sketch the fuel-switching timeline
  from fed to 30-day fast.
\item
  The Randle cycle operates in heart muscle. In a patient with type 2
  diabetes and chronically elevated plasma FFA: predict changes in
  cardiac PDH activity, cardiac efficiency, and susceptibility to
  ischaemia.
\item
  Compare the amplification advantages of the PFK-1/FBPase-1 substrate
  cycle in liver with UCP1-mediated thermogenesis in brown adipose
  tissue.
\item
  Explain how metformin improves glycaemic control. Describe its effects
  on Complex I, AMPK, and hepatic gluconeogenesis.
\item
  Describe three molecular mechanisms by which obesity causes insulin
  resistance. For each, identify the lipid species involved and the
  downstream kinase that phosphorylates IRS-1.
\item
  Using the \hyperref[gl:competitive-inhibition]{\textbf{competitive
  inhibition}} example in the bridge, explain why simply raising
  substrate concentration can overcome some drug interactions but not
  others (hint: \(K_m^{app}\) vs.~\(V_\mathrm{max}\)).
\item
  Estimate energy charge for {[}ATP{]}=2 mM, {[}ADP{]}=1.5 mM,
  {[}AMP{]}=0.5 mM. Is AMPK more likely active than at EC \texorpdfstring{\ensuremath{\approx}}{~} 0.9?
\item
  Describe how mTORC1 is recruited to the lysosomal surface by amino
  acid availability. What is the role of Sestrin2, GATOR1/2, and the Rag
  GTPases? Why is mTORC1 a ``coincidence detector'' rather than a
  single-input sensor?
\item
  Sketch the phosphorylase kinase amplification cascade and estimate the
  end-to-end signal gain from a single glucagon binding event to
  released glucose-1-phosphate. Why is this amplification
  physiologically necessary?
\item
  Compare the contributions of subcutaneous and visceral adipose tissue
  to whole-body insulin resistance in obesity. Why does visceral fat
  have a disproportionate effect on hepatic outcomes?
\item
  Compare NMR and LC-MS as metabolomics platforms. Which is preferable
  for (a) absolute quantification of glucose in serum, (b) discovery of
  unknown lipid species in muscle biopsies, (c) compartment-resolved
  metabolomics in subcellular fractions?
\end{enumerate}

\begin{center}\rule{0.5\linewidth}{0.5pt}\end{center}

\subsection{Further Reading and Source
Notes}\label{further-reading-and-source-notes-13}

\begin{itemize}
\tightlist
\item
  Atkinson (1968). The regulation of enzyme activity in metabolism.
  \emph{Biochemical Journal}, 108 \citep{atkinson1968}.
\item
  Mitchell (1961). Coupling of phosphorylation to electron and hydrogen
  transfer by a chemi-osmotic type of mechanism. \emph{Nature}, 191
  \citep{mitchell1961}.
\end{itemize}

\begin{center}\rule{0.5\linewidth}{0.5pt}\end{center}

\subsection{Key Terms}\label{key-terms-13}

{\def\LTcaptype{none} % do not increment counter
\begin{longtable}[]{@{}
  >{\raggedright\arraybackslash}p{(\linewidth - 2\tabcolsep) * \real{0.2857}}
  >{\raggedright\arraybackslash}p{(\linewidth - 2\tabcolsep) * \real{0.7143}}@{}}
\toprule\noalign{}
\begin{minipage}[b]{\linewidth}\raggedright
Term
\end{minipage} & \begin{minipage}[b]{\linewidth}\raggedright
Definition
\end{minipage} \\
\midrule\noalign{}
\endhead
\bottomrule\noalign{}
\endlastfoot
\textbf{AMPK} & AMP-activated protein kinase; principal energy-sensing
kinase; activated by elevated AMP/ATP via LKB1 or by Ca\(^{2+}\) via
CaMKK\texorpdfstring{\ensuremath{\beta}}{beta} \\
\textbf{mTORC1} & Mechanistic target of rapamycin complex 1;
growth-promoting kinase; activated by insulin, amino acids, and growth
factors at the lysosomal surface \\
\textbf{Ragulator} & Lysosomal scaffold (LAMTOR1--5) that anchors Rag
GTPases and recruits mTORC1 in response to amino acid availability \\
\textbf{Energy charge (EC)} & ({[}ATP{]} + 1/2{[}ADP{]}) /
({[}ATP{]}+{[}ADP{]}+{[}AMP{]}); normal EC approximately 0.85--0.95 \\
\textbf{Insulin} & Anabolic hormone from beta-cells; PI3K/Akt/mTOR
pathway; drives glucose uptake, glycogen synthesis, lipogenesis \\
\textbf{Glucagon} & Catabolic hormone from alpha-cells; cAMP/PKA; drives
glycogenolysis, gluconeogenesis, ketogenesis \\
\textbf{Gluconeogenesis} & Synthesis of glucose from non-carbohydrate
precursors; liver and kidney; costs 6 ATP \\
\textbf{CPT-1} & Carnitine palmitoyl transferase-1; rate-limiting for
mitochondrial fatty acid import; inhibited by malonyl-CoA \\
\textbf{Ketone bodies} & Acetoacetate, beta-hydroxybutyrate, acetone;
liver-produced; brain fuel during starvation \\
\textbf{Phosphorylase kinase} & Multi-subunit kinase activated by PKA;
phosphorylates glycogen phosphorylase to active ``a'' form \\
\textbf{Randle cycle} & Reciprocal inhibition of glucose and fatty acid
oxidation in heart and muscle \\
\textbf{Substrate cycle} & Two opposing reactions running
simultaneously; amplifies flux control; generates heat \\
\textbf{Metabolic flexibility} & Ability to switch between glucose and
fat oxidation; impaired in obesity and T2D \\
\textbf{SIRT1} & NAD\(^+\)-dependent deacetylase; caloric restriction
sensor; activates PGC-1alpha \\
\textbf{Metabolic syndrome} & Cluster of central obesity, insulin
resistance, dyslipidaemia, hypertension, hyperglycaemia \\
\textbf{Visceral adipose tissue} & Fat depot draining via portal vein to
liver; endocrine organ secreting TNF-\texorpdfstring{\ensuremath{\alpha}}{alpha}, IL-6, RBP4 \\
\textbf{Adipokines} & Hormones secreted by adipose tissue (leptin,
adiponectin, resistin, RBP4); regulate appetite and insulin
sensitivity \\
\textbf{Warburg effect} & Aerobic glycolysis in cancer cells; rapid ATP
+ biosynthetic precursors \\
\textbf{Metabolic control analysis} & Quantitative framework for flux
control distribution; summation theorem: sum of \(C_i^J\) = 1 \\
\textbf{Metabolomics} & Systematic measurement of small-molecule
metabolites; NMR and LC-MS as principal platforms \\
\end{longtable}
}

\begin{center}\rule{0.5\linewidth}{0.5pt}\end{center}

\subsubsection{Companion Source
Module}\label{companion-source-module-13}

\textbf{Metabolic Integration and Regulation} should leave a
reproducible trail from a biological claim to the code, figure, diagram,
or paper-based activity that can test it. Use the surfaces below to
inspect the chapter's assumptions, rerun the relevant model, or compare
the manuscript explanation with companion labs and figures.

{\def\LTcaptype{none} % do not increment counter
\begin{longtable}[]{@{}
  >{\raggedright\arraybackslash}p{(\linewidth - 2\tabcolsep) * \real{0.5000}}
  >{\raggedright\arraybackslash}p{(\linewidth - 2\tabcolsep) * \real{0.5000}}@{}}
\toprule\noalign{}
\begin{minipage}[b]{\linewidth}\raggedright
Surface
\end{minipage} & \begin{minipage}[b]{\linewidth}\raggedright
Use it for
\end{minipage} \\
\midrule\noalign{}
\endhead
\bottomrule\noalign{}
\endlastfoot
\breaktt{src/biology/biochemistry/biochemistry.py}
(\breaktt{glycolysis\_summary}, \texttt{atp\_free\_energy},
\breaktt{reaction\_free\_energy}) & Connect pathway summaries to energy
and redox constraints. \\
\breaktt{src/biology/physiology/physiology.py}
(\breaktt{homeostasis\_response}) & Compare cellular flux regulation with
organism-level homeostasis. \\
\breaktt{src/mermaid/biology\_diagrams.py}
(\breaktt{glycolysis\_pathway\_diagram},
\breaktt{hormone\_signaling\_diagram}) & Link metabolic pathways to
endocrine control. \\
\end{longtable}
}

\textbf{Reproducibility check:} state fed/fasted status, tissue,
compartment, and time scale before predicting pathway priority.
\textbf{Cross-reference:} use
\cref{sec:unit_III_bioenergetics_and_respiration},
\cref{sec:unit_I_enzymes_and_kinetics}, and
\cref{sec:unit_IX_endocrine_and_immune}.

\newpage

\section*{Unit IV --- Molecular Genetics:
Introduction}\label{unit-iv-molecular-genetics-introduction}
\addcontentsline{toc}{section}{Unit IV --- Molecular Genetics:
Introduction}

\subsection*{Why This Unit Matters}\label{why-this-unit-matters-4}
\addcontentsline{toc}{subsection}{Why This Unit Matters}

In 1944, Oswald Avery, Colin MacLeod, and Maclyn McCarty performed one
of the most transformative experiments in the history of science. They
showed that the "transforming principle" that could convert harmless
\emph{Streptococcus pneumoniae} into a lethal strain was not protein ---
as virtually everyone expected --- but deoxyribonucleic acid (DNA). The
announcement was met with scepticism. DNA seemed too simple: primarily
four nucleotide bases, arranged in an apparently monotonous double
helix. How could such a molecule encode the almost infinite variety of
life?

The answer, revealed over the next three decades, is linguistic in
structure: a four-letter alphabet (A, T, G, C) organised into
three-letter codons, read in a linear sequence, encoding twenty amino
acids, assembled by ribosomes into an effectively unlimited number of
proteins. This is the molecular logic of life: \textbf{information flows
from DNA to RNA to protein}, and control is exerted at every step. The
central dogma, articulated by Francis Crick in 1958, has been refined
but rarely overturned.

This unit examines the molecular machinery that reads, copies, and
expresses genetic information. You will study the multiprotein complexes
that replicate DNA with an error rate of fewer than one mistake per 10\texorpdfstring{\textsuperscript{9}}{9}
base pairs, the RNA polymerases that transcribe genes in a regulated,
stimulus-responsive manner, and the ribosome --- a 2.5 MDa
ribonucleoprotein machine that interprets the genetic code at
about 15 amino acids per second. You will also examine how
mutations drive evolution and disease, how the genome is structurally
organised, and how CRISPR-Cas9 has made precise genome editing a
clinical reality \citep{doudna2014}.

\begin{center}\rule{0.5\linewidth}{0.5pt}\end{center}

\subsection*{Landmark Discoveries}\label{landmark-discoveries-4}
\addcontentsline{toc}{subsection}{Landmark Discoveries}

{\def\LTcaptype{none} % do not increment counter
\begin{longtable}[]{@{}
  >{\raggedright\arraybackslash}p{(\linewidth - 8\tabcolsep) * \real{0.2407}}
  >{\raggedright\arraybackslash}p{(\linewidth - 8\tabcolsep) * \real{0.0741}}
  >{\raggedright\arraybackslash}p{(\linewidth - 8\tabcolsep) * \real{0.2963}}
  >{\raggedright\arraybackslash}p{(\linewidth - 8\tabcolsep) * \real{0.1667}}
  >{\raggedright\arraybackslash}p{(\linewidth - 8\tabcolsep) * \real{0.2222}}@{}}
\toprule\noalign{}
\begin{minipage}[b]{\linewidth}\raggedright
Discoverer(s)
\end{minipage} & \begin{minipage}[b]{\linewidth}\raggedright
Year
\end{minipage} & \begin{minipage}[b]{\linewidth}\raggedright
Journal / Source
\end{minipage} & \begin{minipage}[b]{\linewidth}\raggedright
Discovery
\end{minipage} & \begin{minipage}[b]{\linewidth}\raggedright
Significance
\end{minipage} \\
\midrule\noalign{}
\endhead
\bottomrule\noalign{}
\endlastfoot
Avery, MacLeod \& McCarty & 1944 & \emph{J. Exp. Med.} & DNA is the
genetic material & Overturned protein hypothesis; opened molecular
genetics \\
Watson \& Crick & 1953 & \emph{Nature} & Double-helix structure of DNA &
Anti-parallel complementary strands implied replication mechanism \\
Meselson \& Stahl & 1958 & \emph{Proc. Natl. Acad. Sci.} &
Semi-conservative DNA replication & Definitive proof of how DNA is
copied; used \texorpdfstring{\textsuperscript{1}}{1}\texorpdfstring{\textsuperscript{5}}{5}N density-gradient experiment \\
Nirenberg \& Matthaei & 1961 & \emph{Proc. Natl. Acad. Sci.} &
Deciphering the genetic code & Used cell-free translation of poly-U to
show UUU = phenylalanine; cracked the codon table \\
Jacob \& Monod & 1961 & \emph{J. Mol. Biol.} & \emph{lac} operon ---
gene regulatory model & First demonstration of gene circuit logic;
inducer-repressor-operator model \\
Alec Jeffreys & 1984 & \emph{Nature} & DNA fingerprinting (RFLP-based) &
Transformed forensic science, paternity testing, and population
genetics \\
Doudna \& Charpentier & 2012 & \emph{Science} & CRISPR-Cas9 programmable
genome editing & Created a precise molecular "scissors" for any target
sequence; Nobel Prize 2020 \\
\end{longtable}
}

\begin{center}\rule{0.5\linewidth}{0.5pt}\end{center}

\subsection*{Key Concepts and
Connections}\label{key-concepts-and-connections-4}
\addcontentsline{toc}{subsection}{Key Concepts and Connections}

\begin{figure}[htbp]
\centering
\includegraphics[width=0.82\linewidth,height=4.2in,keepaspectratio]{../figures/mermaid_inline/inline_mermaid_0067_802fdbae784c.png}
\caption{Unit IV concept map - Molecular Genetics. Dark = DNA; purple = regulation; orange = mutations; red = editing technology.}
\end{figure}

\begin{center}\rule{0.5\linewidth}{0.5pt}\end{center}

\subsection*{Current Evidence Thread}\label{current-evidence-thread-4}
\addcontentsline{toc}{subsection}{Current Evidence Thread}

Read this unit as molecular genetics built from stacked evidence layers:
DNA sequence and replication fidelity, transcription and translation
output, the chromatin and methylation state that gates that output, and
the variant calls and clinical interpretations that follow from most
three. Molecular genetics now spans single-reference sequences,
telomere-to-telomere assemblies, pangenome graphs, long-read sequencing,
CRISPR medicines, and ethical deployment. As you move through the
chapters, keep a two-column note: \textbf{claim} on the left,
\textbf{evidence that would change my confidence} on the right. By the
end of the unit, each major idea should be tied to a measurement, model,
citation, or paper-based lab decision.

\subsection*{Chapter Roadmap}\label{chapter-roadmap-4}
\addcontentsline{toc}{subsection}{Chapter Roadmap}

{\def\LTcaptype{none} % do not increment counter
\begin{longtable}[]{@{}
  >{\raggedright\arraybackslash}p{(\linewidth - 6\tabcolsep) * \real{0.1556}}
  >{\raggedright\arraybackslash}p{(\linewidth - 6\tabcolsep) * \real{0.1111}}
  >{\raggedright\arraybackslash}p{(\linewidth - 6\tabcolsep) * \real{0.2889}}
  >{\raggedright\arraybackslash}p{(\linewidth - 6\tabcolsep) * \real{0.4444}}@{}}
\toprule\noalign{}
\begin{minipage}[b]{\linewidth}\raggedright
Chapter
\end{minipage} & \begin{minipage}[b]{\linewidth}\raggedright
Title
\end{minipage} & \begin{minipage}[b]{\linewidth}\raggedright
Core Question
\end{minipage} & \begin{minipage}[b]{\linewidth}\raggedright
Key Equation / Model
\end{minipage} \\
\midrule\noalign{}
\endhead
\bottomrule\noalign{}
\endlastfoot
\textbf{12} & DNA Replication and the Cell Cycle & How is DNA copied
with such extraordinary accuracy, and how is replication coupled to cell
division? & Error rate: about 10\texorpdfstring{\textsuperscript{-}}{-}\texorpdfstring{\textsuperscript{9}}{9} per base; cell cycle
checkpoint models \\
\textbf{13} & Gene Expression & How is genetic information transcribed
and translated into functional proteins? & Rate equations for coupled
transcription-translation \\
\textbf{14} & Mutations, CRISPR, and Genomics & What causes mutations,
how are they repaired, and what does the human reference genome reveal?
& Jukes-Cantor distance; mutation rate \texorpdfstring{\ensuremath{\times}}{x} generation time \\
\textbf{15} & Epigenetics and Gene Regulation & How do chromatin,
methylation, and regulatory RNAs tune gene expression without changing
DNA sequence? & Methylation state; histone modification; regulatory
network logic \\
\end{longtable}
}

\begin{center}\rule{0.5\linewidth}{0.5pt}\end{center}

\subsection*{Connections Across the
Textbook}\label{connections-across-the-textbook-4}
\addcontentsline{toc}{subsection}{Connections Across the Textbook}

\begin{itemize}
\tightlist
\item
  \textbf{DNA replication} directly links to Unit V (meiosis and
  crossing-over in \cref{sec:unit_V_chromosomal_inheritance}) and Unit
  VI (molecular clock calculations using sequence divergence).
\item
  \textbf{Gene regulation} (operons, TFs, miRNA) reappears in Unit VIII
  (plant hormone signal transduction) and Unit IX (endocrine control of
  gene expression via steroid receptors).
\item
  \textbf{Mutations and DNA repair} are essential context for Unit VI
  (evolution as mutation + selection) and oncogene/tumour-suppressor
  biology in Unit II and Unit IX.
\item
  \textbf{CRISPR-Cas9} appears in clinical connection boxes throughout
  --- gene therapy (Unit V), antibiotic resistance (Unit VII), and plant
  engineering (Unit VIII).
\end{itemize}

\begin{quote}
\textbf{Key vocabulary introduced here:} nucleotide, base pair, template
strand, coding strand, codon, anticodon, spliceosome, intron, exon,
promoter, enhancer, transcription factor, epigenome, proto-oncogene,
restriction enzyme, CRISPR.
\end{quote}

\subsection*{Computational Toolbox --- Unit
IV}\label{computational-toolbox-unit-iv}
\addcontentsline{toc}{subsection}{Computational Toolbox --- Unit IV}

\begin{Shaded}
\begin{Highlighting}[]
\ImportTok{from}\NormalTok{ biology.genetics }\ImportTok{import}\NormalTok{ translate\_mrna, dna\_complement, punnett\_square}

\CommentTok{\# Translate an mRNA sequence into amino acids}
\NormalTok{mrna }\OperatorTok{=} \StringTok{"AUGUUUGAAGAACUUUAG"}
\NormalTok{protein }\OperatorTok{=}\NormalTok{ translate\_mrna(mrna)}
\BuiltInTok{print}\NormalTok{(}\SpecialStringTok{f"mRNA:    }\SpecialCharTok{\{}\NormalTok{mrna}\SpecialCharTok{\}}\SpecialStringTok{"}\NormalTok{)}
\BuiltInTok{print}\NormalTok{(}\SpecialStringTok{f"Protein: }\SpecialCharTok{\{}\StringTok{\textquotesingle{}{-}\textquotesingle{}}\SpecialCharTok{.}\NormalTok{join(protein)}\SpecialCharTok{\}}\SpecialStringTok{"}\NormalTok{)}
\CommentTok{\# Expected: Protein: Met{-}Phe{-}Glu{-}Glu{-}Leu (stop at UAG)}

\CommentTok{\# Generate the DNA complement of a coding sequence}
\NormalTok{dna\_template }\OperatorTok{=} \StringTok{"ATGTTCGAAATG"}
\NormalTok{complement }\OperatorTok{=}\NormalTok{ dna\_complement(dna\_template)}
\BuiltInTok{print}\NormalTok{(}\SpecialStringTok{f"5\textquotesingle{}→3\textquotesingle{} complement: }\SpecialCharTok{\{}\NormalTok{complement}\SpecialCharTok{\}}\SpecialStringTok{"}\NormalTok{)}
\CommentTok{\# Expected: 5\textquotesingle{}→3\textquotesingle{} complement: CATTTCGAACAT  (antiparallel)}

\CommentTok{\# Punnett square: monohybrid cross Aa × Aa}
\NormalTok{result }\OperatorTok{=}\NormalTok{ punnett\_square(}\StringTok{"Aa"}\NormalTok{, }\StringTok{"Aa"}\NormalTok{)}
\BuiltInTok{print}\NormalTok{(}\SpecialStringTok{f"Genotype ratios: }\SpecialCharTok{\{}\NormalTok{result}\SpecialCharTok{.}\NormalTok{genotype\_ratios}\SpecialCharTok{\}}\SpecialStringTok{"}\NormalTok{)}
\BuiltInTok{print}\NormalTok{(}\SpecialStringTok{f"Phenotype ratios: }\SpecialCharTok{\{}\NormalTok{result}\SpecialCharTok{.}\NormalTok{phenotype\_ratios}\SpecialCharTok{\}}\SpecialStringTok{"}\NormalTok{)}
\CommentTok{\# Expected:}
\CommentTok{\# Genotype ratios: \{\textquotesingle{}AA\textquotesingle{}: 0.25, \textquotesingle{}Aa\textquotesingle{}: 0.5, \textquotesingle{}aa\textquotesingle{}: 0.25\}}
\CommentTok{\# Phenotype ratios: \{\textquotesingle{}dominant\textquotesingle{}: 0.75, \textquotesingle{}recessive\textquotesingle{}: 0.25\}}
\end{Highlighting}
\end{Shaded}

\begin{quote}
\textbf{Try it yourself:} Change the mRNA to \texttt{AUGCAGUGA} --- note
that UGA is a stop codon. Compare proteins produced from synonymous
codon changes using the \texttt{GENETIC\_CODE} dictionary.
\end{quote}

\begin{center}\rule{0.5\linewidth}{0.5pt}\end{center}

\emph{Source modules: \breaktt{src/biology/genetics/} ---
\texttt{translate\_mrna()}, \texttt{dna\_complement()},
\texttt{punnett\_square()}, \texttt{hardy\_weinberg()}.} \emph{Figures:
\breaktt{src/mermaid/biology\_diagrams.py} (central dogma diagrams,
CRISPR mechanism).}

\newpage

\section{DNA Replication and the Cell
Cycle}\label{dna-replication-and-the-cell-cycle}

\label{sec:unit_IV_dna_replication_and_cell_cycle}

\begin{quote}
\textbf{Ch 12} \texorpdfstring{\ensuremath{\cdot}}{.} Level 2/3 \texorpdfstring{\ensuremath{\cdot}}{.} 55 min read \texorpdfstring{\ensuremath{\cdot}}{.} 75 min lecture \texorpdfstring{\ensuremath{\cdot}}{.}
Prerequisites: \cref{sec:unit_I_macromolecules},
\cref{sec:unit_II_cell_structure}
\end{quote}

\subsection{Learning Objectives}\label{learning-objectives-14}

\begin{enumerate}
\def\labelenumi{\arabic{enumi}.}
\tightlist
\item
  Describe the structure of DNA and explain why it is well-suited to
  information storage and replication.
\item
  Explain semiconservative replication and the key
  \hyperref[gl:enzyme]{\textbf{enzyme}}s at the replication fork.
\item
  Compare the replication machinery of prokaryotes and
  \hyperref[gl:eukaryote]{\textbf{eukaryote}}s.
\item
  Describe origin licensing, the role of MCM helicases, and how
  re-replication is prevented.
\item
  Explain the \hyperref[gl:telomere]{\textbf{telomere}} end-replication
  problem and the role of telomerase.
\item
  Describe the stages of the \hyperref[gl:cell-cycle]{\textbf{cell
  cycle}} and the function of cyclin-CDK complexes at each checkpoint.
\item
  Explain the DNA damage response pathway involving ATM, ATR, and p53.
\item
  Describe the stages of \hyperref[gl:mitosis]{\textbf{mitosis}} at the
  molecular level, including spindle assembly checkpoint mechanisms.
\item
  Explain the mechanism and biological importance of
  \hyperref[gl:meiosis]{\textbf{meiosis}}, including crossing over and
  \hyperref[gl:independent-assortment]{\textbf{independent assortment}}.
\item
  Describe DNA damage types and the key repair pathways, and connect
  repair defects to human disease.
\end{enumerate}

\subsubsection{Study Blueprint}\label{study-blueprint-14}

\begin{itemize}
\tightlist
\item
  \textbf{Big idea:} Genome copying is accurate because chemistry,
  enzyme proofreading, and checkpoints cooperate.
\item
  \textbf{Core concepts:} semiconservative replication, polymerase,
  proofreading, checkpoints.
\item
  \textbf{Framework alignment:} Vision \& Change: Information flow,
  exchange, and storage, Structure and function; AP Biology: Information
  Storage and Transmission, Systems Interactions; NGSS-style topics:
  Inheritance and Variation of Traits, Structure and Function.
\item
  \textbf{Model or quantitative lens:} Replication-fork speed,
  error-rate, and cell-cycle timing calculations.
\item
  \textbf{Data skill:} Interpret replication or cell-cycle data from
  timing, labeling, or checkpoint perturbations.
\item
  \textbf{Practice cadence:} Concept Explanation, Questions and Methods,
  Argumentation.
\item
  \textbf{Common misconception to repair:} High fidelity is not
  automatic; it is built from multiple partially redundant safeguards.
\item
  \textbf{Primary lab:}
  \cref{sec:lab_unit_IV_dna_replication_and_cell_cycle}.
\item
  \textbf{Question bank:}
  \cref{sec:q_unit_IV_dna_replication_and_cell_cycle}.
\item
  \textbf{Transfer task:} Transfer replication logic to cancer, aging,
  viral replication, or antibiotic targets.
\item
  \textbf{Bridge to computation:}
  \breaktt{biology.genetics.genetics.dna\_complement}.
\end{itemize}

\begin{center}\rule{0.5\linewidth}{0.5pt}\end{center}

\begin{quote}
\textbf{Opening Vignette --- The Double Helix's Debt to a Woman
Scientist}

In 1952, Rosalind Franklin produced the sharpest X-ray diffraction
photograph of DNA ever captured --- ``Photo 51.'' The image was obtained
without her permission by Maurice Wilkins and shown to James Watson.
Recognising in an instant that the helical periodicity and the 3.4 Å
spacing implied a double helix, Watson and Crick used this insight to
complete their landmark 1953 \emph{Nature} paper. Franklin, who died of
ovarian cancer in 1958 at age 37, rarely received the Nobel Prize
awarded to Watson, Crick, and Wilkins in 1962. Yet her replication
studies --- painstakingly mapping how DNA unwinds, is copied
semiconservatively, and repairs itself --- anticipated the entire field
of cell-cycle biology. Every chapter covering DNA replication and the
cell cycle is, in part, her story.
\end{quote}

\subsection{DNA as the Genetic
Material}\label{dna-as-the-genetic-material}

The Avery-MacLeod-McCarty experiment (1944) demonstrated that DNA (not
\hyperref[gl:protein]{\textbf{protein}}) carries genetic information by
showing that purified DNA from virulent \emph{Streptococcus pneumoniae}
could transform non-virulent strains. The Hershey-Chase blender
experiment (1952) confirmed this with radioactively labelled T4 phage:
\(^{32}\)P-labelled DNA (not \(^{35}\)S-labelled protein) entered
bacterial cells during infection.

\subsubsection{DNA Double Helix
Structure}\label{dna-double-helix-structure}

Watson and Crick's 1953 model of the B-form DNA helix, built on Rosalind
Franklin's X-ray diffraction data (Photo 51) and Chargaff's base-pairing
rules:

\begin{itemize}
\tightlist
\item
  \textbf{Antiparallel complementary strands} running 5' to 3' and 3' to
  5'
\item
  \textbf{Base pairing}: A-T (2
  \hyperref[gl:hydrogen-bond]{\textbf{hydrogen bond}}s), G-C (3 hydrogen
  bonds)
\item
  \textbf{Major groove}: wide (about 22 angstroms) and deep --
  the primary site for sequence-specific protein binding
  (\hyperref[gl:transcription]{\textbf{transcription}} factors,
  restriction enzymes). The pattern of hydrogen bond donors and
  acceptors in the major groove is unique for each base pair, allowing
  proteins to ``read'' the sequence without unwinding the helix.
\item
  \textbf{Minor groove}: narrow (about 12 angstroms) -- bound
  by certain drugs (e.g., netropsin, distamycin) and architectural
  proteins (e.g., HMG-box proteins, TBP)
\item
  \textbf{Right-handed helix}, 10.5 bp per turn, rise 0.34 nm/bp, pitch
  \texorpdfstring{\ensuremath{\approx}}{~}3.6 nm, diameter 2.0 nm
\end{itemize}

The phosphodiester backbone runs 5' to 3' and carries two negative
charges per \hyperref[gl:nucleotide]{\textbf{nucleotide}} at
physiological \hyperref[gl:ph]{\textbf{pH}} (\(pK_a < 1\)). This
negative charge means DNA repels other nucleic acids and requires
cationic counterions (Mg\(^{2+}\),
\hyperref[gl:histone]{\textbf{histone}}s, polyamines) for compaction and
stability.

\textbf{Alternative DNA conformations:}

{\def\LTcaptype{none} % do not increment counter
\begin{longtable}[]{@{}
  >{\raggedright\arraybackslash}p{(\linewidth - 8\tabcolsep) * \real{0.1364}}
  >{\raggedright\arraybackslash}p{(\linewidth - 8\tabcolsep) * \real{0.1591}}
  >{\raggedright\arraybackslash}p{(\linewidth - 8\tabcolsep) * \real{0.2045}}
  >{\raggedright\arraybackslash}p{(\linewidth - 8\tabcolsep) * \real{0.2273}}
  >{\raggedright\arraybackslash}p{(\linewidth - 8\tabcolsep) * \real{0.2727}}@{}}
\toprule\noalign{}
\begin{minipage}[b]{\linewidth}\raggedright
Form
\end{minipage} & \begin{minipage}[b]{\linewidth}\raggedright
Helix
\end{minipage} & \begin{minipage}[b]{\linewidth}\raggedright
bp/turn
\end{minipage} & \begin{minipage}[b]{\linewidth}\raggedright
Diameter
\end{minipage} & \begin{minipage}[b]{\linewidth}\raggedright
Conditions
\end{minipage} \\
\midrule\noalign{}
\endhead
\bottomrule\noalign{}
\endlastfoot
B-DNA & Right-handed & 10.5 & 2.0 nm & Physiological; most common \\
A-DNA & Right-handed & 11 & 2.6 nm & Dehydrated; RNA-DNA hybrids \\
Z-DNA & Left-handed & 12 & 1.8 nm & Alternating purine-pyrimidine; may
regulate transcription \\
\end{longtable}
}

\textbf{Information capacity:} The human haploid
\hyperref[gl:genome]{\textbf{genome}} contains approximately 3.1 billion
base pairs, equivalent to roughly 725 megabytes of digital data (2 bits
per base). If fully stretched, the DNA in a single cell would extend
approximately 2 meters, yet it is compacted into a nucleus approximately
6 micrometers in diameter. This compaction is achieved through a
hierarchy of packaging: 147 bp of DNA wraps around each histone octamer
to form a \hyperref[gl:nucleosome]{\textbf{nucleosome}} (the ``beads on
a string''), which further compacts into the 30-nm
\hyperref[gl:chromatin]{\textbf{chromatin}} fiber, loop domains
(anchored by CTCF and cohesin), and ultimately the mitotic
\hyperref[gl:chromosome]{\textbf{chromosome}}.

\begin{quote}
\textbf{Clinical Connection: Z-DNA and Autoimmunity} Antibodies against
Z-DNA are found in patients with systemic lupus erythematosus (SLE).
Z-DNA forms transiently behind moving RNA polymerase due to negative
supercoiling, and \hyperref[gl:mutation]{\textbf{mutation}}s in the
Z-DNA binding protein ADAR1 cause Aicardi-Goutieres syndrome, a severe
autoinflammatory disorder. Z-DNA-binding proteins such as ZBP1 (DAI)
also function as innate immune sensors, recognizing viral Z-form nucleic
acids.
\end{quote}

\begin{center}\rule{0.5\linewidth}{0.5pt}\end{center}

\subsection{DNA Replication}\label{dna-replication}

Replication is \textbf{semiconservative} -- each daughter double helix
retains one parental strand. This was demonstrated by
\citet{meselson1958} using \(^{15}\)N density gradient centrifugation in
\emph{E. coli}. After one generation in \(^{14}\)N medium, most DNA was
intermediate density (ruling out conservative replication). After two
generations, half was intermediate and half was light (ruling out
dispersive replication).

\subsubsection{Origins of Replication and
Licensing}\label{origins-of-replication-and-licensing}

Replication begins at specific genomic sites called \textbf{origins of
replication (ori)}:

\begin{itemize}
\tightlist
\item
  \textbf{Prokaryotes}: \emph{E. coli} has a single origin
  (\emph{oriC}), a 245-bp sequence containing five 9-bp DnaA-binding
  sites (DnaA boxes) and three AT-rich 13-mer repeats that melt first
  upon initiation
\item
  \textbf{Eukaryotes}: Human genome contains approximately 30,000-50,000
  origins; yeast origins are called \textbf{autonomously replicating
  sequences (ARS)}, each approximately 150 bp containing an 11-bp ARS
  consensus sequence (ACS)
\end{itemize}

\textbf{Origin licensing} ensures each origin fires primarily once per S
phase:

\begin{enumerate}
\def\labelenumi{\arabic{enumi}.}
\tightlist
\item
  In late M/early G1, the \textbf{Origin Recognition Complex (ORC)}
  binds origins constitutively (in yeast) or dynamically (in metazoans)
\item
  \textbf{Cdc6} and \textbf{Cdt1} load the \textbf{MCM2-7 helicase} (a
  hexameric ring) onto DNA as a double hexamer -- this is the
  ``licensed'' pre-replicative complex (pre-RC)
\item
  At the G1/S transition, \textbf{CDK} (S-CDK = cyclin A-CDK2) and
  \textbf{DDK} (Cdc7-Dbf4) phosphorylate MCM, activating the helicase
  and converting pre-RC to pre-initiation complex (pre-IC)
\item
  Re-licensing is prevented by CDK-mediated phosphorylation of Cdc6
  (targeting it for degradation) and by geminin binding to Cdt1
\end{enumerate}

\textbf{MCM licensing --- molecular detail.} The pre-RC formation
involves a precisely choreographed sequence:

\begin{enumerate}
\def\labelenumi{\arabic{enumi}.}
\tightlist
\item
  \textbf{ORC binding (G1):} ORC1--6 heterohexamer binds origins. ORC1
  contains an ATP-binding AAA+ ATPase fold; ATP hydrolysis triggers
  loading. In humans ORC1 levels oscillate (peaks late M/early G1;
  degraded in S/G2) --- ensuring temporal restriction of licensing.
\item
  \textbf{Cdc6 recruitment:} Cdc6 (also AAA+ ATPase) joins ORC, forming
  an ORC--Cdc6 complex that wraps DNA into a left-handed superhelix.
\item
  \textbf{Cdt1--MCM2-7 loading:} Cdt1 (the licensing factor) binds
  MCM2-7 and chaperones it to the ORC--Cdc6--origin complex. ATP
  hydrolysis by ORC and Cdc6 drives the MCM2-7 ring open at the
  MCM2/MCM5 interface, allowing dsDNA to thread through.
\item
  \textbf{Double hexamer assembly:} A second Cdt1--MCM2-7 is loaded in
  the opposite orientation, creating a head-to-head MCM2-7 double
  hexamer encircling dsDNA. The double hexamer is the licensed pre-RC.
\item
  \textbf{CDT1 destruction in S phase:} Once MCM is loaded, S-CDK
  phosphorylates Cdt1 \texorpdfstring{\ensuremath{\to}}{->} ubiquitin ligase SCF\textsuperscript{Skp2} \texorpdfstring{\ensuremath{\to}}{->}
  proteasomal degradation. \textbf{CRL4\textsuperscript{CDT2}} also
  targets Cdt1 for degradation when bound to chromatin (PCNA-coupled
  mechanism).
\item
  \textbf{Geminin sequesters Cdt1:} Geminin is an inhibitor that binds
  Cdt1, preventing it from re-licensing origins during S/G2/M. Geminin
  levels rise from S phase, peaking in M, and are degraded at anaphase
  by APC/C-Cdc20.
\item
  \textbf{Re-replication prevention is therefore three-fold:}

  \begin{itemize}
  \tightlist
  \item
    CDK-dependent Cdc6 export from nucleus and degradation
  \item
    CDK-dependent Cdt1 degradation
  \item
    Geminin-mediated Cdt1 sequestration
  \item
    Loss of any of these (e.g., siRNA against geminin) causes
    catastrophic re-replication, DNA damage, and apoptosis.
  \end{itemize}
\item
  \textbf{Origin firing in S phase:} A subset of licensed origins fires
  (about 10 \% per cell-cycle round in human cells); the
  remainder are dormant backups. Activation requires \textbf{DDK
  (Cdc7-Dbf4)} phosphorylation of MCM2-MCM4-MCM6 + \textbf{S-CDK}
  phosphorylation of Sld2/Sld3 \texorpdfstring{\ensuremath{\to}}{->} recruitment of \textbf{Cdc45} and
  \textbf{GINS} to form the \textbf{CMG helicase} (Cdc45-MCM-GINS) ---
  the active replicative helicase.
\end{enumerate}

\textbf{Replication timing domains.} Origin firing is not random in
time. Each about 1 Mb region has a characteristic
\textbf{replication timing}: - \textbf{Early-replicating domains} (S
phase first 1--3 hours): gene-rich, euchromatic, transcriptionally
active, A-compartment, marked by H3K27ac, H3K4me3, high accessibility -
\textbf{Mid-S replicating domains} (3--5 h): transition zones -
\textbf{Late-replicating domains} (5--8 h): heterochromatic, gene-poor,
B-compartment, marked by H3K9me3 and lamina-association (LADs)

Replication timing correlates with chromatin state and is reprogrammed
during cell-fate transitions. Late replication is associated with
elevated mutation rate (about 2\texorpdfstring{\ensuremath{\times}}{x} background) --- because the
late-S nucleotide pool is depleted of dNTPs (which were consumed
earlier), increasing polymerase mis-incorporation. This explains the
spatial mutation-rate variation observed in cancer genomes:
late-replicating regions accumulate about 2-fold more mutations
per Mb than early-replicating regions.

\textbf{Quantitative parameters of replication licensing:} - Number of
origins per genome: about 30,000--50,000 (about  one
per 50--100 kb) - MCM2-7 loaded per origin: 2 (head-to-head double
hexamer) - Total MCM2-7 in nucleus: about 10--20-fold excess
over origins \texorpdfstring{\ensuremath{\to}}{->} MCM ``paradox'' - Origins fired per cell cycle:
about 15,000 (about 30 \% of licensed origins) - Dormant
origins: provide backup if a fork stalls or fails --- mutation in ATR
(which suppresses dormant-origin firing under stress) causes Seckel
syndrome (microcephaly, growth failure)

\begin{quote}
\textbf{Clinical Connection: Meier-Gorlin Syndrome} Mutations in ORC1,
ORC4, ORC6, CDT1, or CDC6 cause Meier-Gorlin syndrome, a rare primordial
dwarfism characterized by short stature, absent or small patellae, and
microtia. This underscores the essential role of origin licensing in
normal growth and development.
\end{quote}

\subsubsection{The Replication Fork
Machinery}\label{the-replication-fork-machinery}

\begin{figure}[htbp]
\centering
\includegraphics[width=0.82\linewidth,height=4.2in,keepaspectratio]{../figures/mermaid_inline/inline_mermaid_0068_7fcdfa521831.png}
\caption{Eukaryotic DNA replication machinery at the replication fork. The leading strand is synthesized continuously, while the lagging strand requires repeated priming and ligation of Okazaki fragments.}
\end{figure}

\subsubsection{Key enzymes at the fork:}\label{key-enzymes-at-the-fork}

{\def\LTcaptype{none} % do not increment counter
\begin{longtable}[]{@{}
  >{\raggedright\arraybackslash}p{(\linewidth - 6\tabcolsep) * \real{0.2000}}
  >{\raggedright\arraybackslash}p{(\linewidth - 6\tabcolsep) * \real{0.2750}}
  >{\raggedright\arraybackslash}p{(\linewidth - 6\tabcolsep) * \real{0.2750}}
  >{\raggedright\arraybackslash}p{(\linewidth - 6\tabcolsep) * \real{0.2500}}@{}}
\toprule\noalign{}
\begin{minipage}[b]{\linewidth}\raggedright
Enzyme
\end{minipage} & \begin{minipage}[b]{\linewidth}\raggedright
Prokaryote
\end{minipage} & \begin{minipage}[b]{\linewidth}\raggedright
Eukaryote
\end{minipage} & \begin{minipage}[b]{\linewidth}\raggedright
Function
\end{minipage} \\
\midrule\noalign{}
\endhead
\bottomrule\noalign{}
\endlastfoot
Helicase & DnaB (5' to 3') & MCM2-7/CMG (3' to 5') & Unwinds double
helix; uses ATP hydrolysis (about 500-1000 bp/s in \emph{E.
coli}, about 50 bp/s in humans) \\
Single-strand binding & SSB & RPA (replication protein A) & Stabilizes
single-stranded regions, prevents re-annealing and nuclease attack \\
Primase & DnaG & DNA Pol alpha/primase complex & Synthesizes short RNA
primers (8-10 nt in prokaryotes; 7-12 nt RNA + 20 nt DNA in
eukaryotes) \\
Leading strand polymerase & DNA Pol III & DNA Pol epsilon & Synthesizes
DNA 5' to 3'; high processivity with clamp \\
Lagging strand polymerase & DNA Pol III & DNA Pol delta & Synthesizes
Okazaki fragments; also fills primer gaps \\
Sliding clamp & beta-clamp (homodimer) & PCNA (homotrimer) &
Processivity factor; tethers polymerase to DNA; increases processivity
from about 10 to \textgreater10,000 nt \\
Clamp loader & gamma complex & RFC (replication factor C) &
ATP-dependent loading of sliding clamp \\
Primer removal & DNA Pol I (5' to 3' exo) & RNase H1 + FEN1 (flap
endonuclease) & Removes RNA primers from Okazaki fragments \\
Ligase & \hyperref[gl:dna-ligase]{\textbf{DNA ligase}}
(NAD\(^+\)-dependent) & DNA ligase I (ATP-dependent) & Seals nicks
between Okazaki fragments \\
Topoisomerase & Gyrase (Topo II) & Topo I and Topo II & Relieves
positive supercoiling ahead of fork \\
\end{longtable}
}

\subsubsection{mRNA Vaccines: A Modern Application of Nucleic Acid
Technology}\label{mrna-vaccines-a-modern-application-of-nucleic-acid-technology}

\begin{figure}[htbp]
\centering
\includegraphics[width=0.82\linewidth,height=4.2in,keepaspectratio]{../figures/mermaid_inline/inline_mermaid_0069_4c0d48a0418b.png}
\caption{Mechanism of mRNA vaccines. Synthetic mRNA encoding a viral antigen is delivered to host cells, where it is translated into protein. The protein is processed and presented to both helper T cells (via MHC II) and cytotoxic T cells (via MHC I), generating both antibody and cellular immune responses.}
\end{figure}

\textbf{Molecular design}:

\begin{itemize}
\item
  \textbf{Sequence optimization}: The mRNA sequence is optimized for
  stability and translation efficiency. This includes:

  \begin{itemize}
  \tightlist
  \item
    \textbf{5' cap}: A modified guanosine (Cap-1: N7-methylated,
    2'-O-methylated) mimics natural eukaryotic mRNA and prevents
    recognition by innate immune sensors (RIG-I, MDA5)
  \item
    \textbf{5' UTR}: Contains Kozak consensus sequence
    (GCCACC{[}A/G{]}CC) for efficient translation initiation
  \item
    \textbf{Open reading frame (ORF)}: Encodes the target antigen (e.g.,
    SARS-CoV-2 spike protein). The sequence is codon-optimized to match
    human tRNA abundance, reducing rare codon clusters that could stall
    translation.
  \item
    \textbf{3' UTR}: Contains stability elements (e.g., from human
    alpha-globin or cytochrome P450 genes) and poly-A tail (100-150 nt)
    to protect from degradation and enhance translation.
  \item
    \textbf{Nucleoside modifications}: Pseudouridine (\texorpdfstring{\ensuremath{\Psi}}{Psi}) or
    N1-methylpseudouridine (m1\texorpdfstring{\ensuremath{\Psi}}{Psi}) replace uridine to reduce
    immunogenicity and increase translation efficiency.
  \end{itemize}
\item
  \textbf{In vitro transcription (IVT)}: The optimized DNA template
  (linearized plasmid or PCR product) is transcribed using T7 RNA
  polymerase in a cell-free system. Nucleoside triphosphates (NTPs)
  including modified UTP (e.g., m1\texorpdfstring{\ensuremath{\Psi}}{Psi}-TP) are added. The reaction produces
  a 5'-triphosphate-ended mRNA, which is then enzymatically capped.
\item
  \textbf{Purification and formulation}: The crude mRNA is treated with
  DNase to remove template DNA, then purified (often by tangential flow
  filtration or HPLC). The purified mRNA is formulated into
  \textbf{lipid nanoparticles (LNPs)} --- spherical vesicles
  about 80-100 nm in diameter composed of ionizable cationic
  lipids, cholesterol, phospholipids, and PEG-lipids. The ionizable
  lipid is positively charged at low pH (during formulation) but neutral
  at physiological pH, reducing toxicity. LNPs protect mRNA from
  degradation and facilitate cellular uptake.
\end{itemize}

\textbf{Cellular uptake and translation}:

\begin{enumerate}
\def\labelenumi{\arabic{enumi}.}
\tightlist
\item
  \textbf{Injection}: mRNA-LNP suspension is injected intramuscularly
  (typically deltoid). The injection site contains antigen-presenting
  cells (APCs) like dendritic cells and macrophages, as well as muscle
  cells.
\item
  \textbf{Uptake}: LNPs fuse with cell membranes or are endocytosed. In
  the endosome, the ionizable lipid becomes protonated, causing
  endosomal membrane disruption and release of mRNA into the cytoplasm.
\item
  \textbf{Translation}: Host ribosomes bind to the 5' cap and translate
  the mRNA into the antigen protein. The protein may be secreted (if
  signal peptide included) or remain intracellular.
\item
  \textbf{Processing and presentation}:

  \begin{itemize}
  \tightlist
  \item
    \textbf{MHC class II pathway}: Exogenous antigen taken up by APCs is
    degraded in lysosomes and presented on MHC II to CD4+ helper T
    cells.
  \item
    \textbf{MHC class I pathway}: Endogenous antigen (produced inside
    the cell) is degraded by the proteasome into peptides, transported
    to the ER via TAP, and presented on MHC I to CD8+ cytotoxic T cells.
  \end{itemize}
\end{enumerate}

\textbf{Immune activation}:

\begin{itemize}
\tightlist
\item
  \textbf{Humoral immunity}: B cells recognizing the antigen receive
  help from CD4+ T cells and differentiate into plasma cells producing
  neutralizing antibodies (IgG, IgM initially, then IgG). These
  antibodies block viral entry by binding to the spike protein
  receptor-binding domain.
\item
  \textbf{Cellular immunity}: CD8+ T cells recognize infected cells
  presenting viral peptides and kill them, preventing viral replication.
\item
  \textbf{Immune memory}: Long-lived plasma cells and memory B/T cells
  provide lasting protection.
\end{itemize}

\textbf{Advantages over traditional vaccines}:

\begin{itemize}
\tightlist
\item
  \textbf{No live virus}: Safer for immunocompromised individuals; no
  risk of reversion to virulence.
\item
  \textbf{Rapid design and manufacturing}: Once the viral sequence is
  known, mRNA can be designed in days and manufactured in weeks
  (cell-free synthesis). Traditional vaccines require growing virus in
  eggs or cell culture (months).
\item
  \textbf{Strong immune response}: Both antibody and T-cell responses
  are robust.
\item
  \textbf{Flexible platform}: The same LNP formulation can deliver mRNA
  for different antigens; the platform can be quickly adapted to new
  variants or emerging pathogens.
\end{itemize}

\textbf{Computational aspects}:

Modern mRNA vaccine design relies heavily on computational tools:

\begin{Shaded}
\begin{Highlighting}[]
\ImportTok{import}\NormalTok{ numpy }\ImportTok{as}\NormalTok{ np}
\ImportTok{from}\NormalTok{ Bio }\ImportTok{import}\NormalTok{ SeqIO}
\ImportTok{from}\NormalTok{ collections }\ImportTok{import}\NormalTok{ Counter}

\CommentTok{\# Example: Analyze GC content and codon usage of an mRNA sequence}
\KeywordTok{def}\NormalTok{ analyze\_mrna\_sequence(seq):}
    \CommentTok{"""Analyze key features of an mRNA vaccine candidate."""}
\NormalTok{    gc\_content }\OperatorTok{=}\NormalTok{ (seq.count(}\StringTok{\textquotesingle{}G\textquotesingle{}}\NormalTok{) }\OperatorTok{+}\NormalTok{ seq.count(}\StringTok{\textquotesingle{}C\textquotesingle{}}\NormalTok{)) }\OperatorTok{/} \BuiltInTok{len}\NormalTok{(seq) }\OperatorTok{*} \DecValTok{100}
    \CommentTok{\# Count codon frequencies}
\NormalTok{    codon\_table }\OperatorTok{=}\NormalTok{ \{\}}
    \ControlFlowTok{for}\NormalTok{ i }\KeywordTok{in} \BuiltInTok{range}\NormalTok{(}\DecValTok{0}\NormalTok{, }\BuiltInTok{len}\NormalTok{(seq)}\OperatorTok{{-}}\DecValTok{2}\NormalTok{, }\DecValTok{3}\NormalTok{):}
\NormalTok{        codon }\OperatorTok{=}\NormalTok{ seq[i:i}\OperatorTok{+}\DecValTok{3}\NormalTok{]}
        \ControlFlowTok{if}\NormalTok{ codon }\KeywordTok{in}\NormalTok{ codon\_table:}
\NormalTok{            codon\_table[codon] }\OperatorTok{+=} \DecValTok{1}
        \ControlFlowTok{else}\NormalTok{:}
\NormalTok{            codon\_table[codon] }\OperatorTok{=} \DecValTok{1}
    \ControlFlowTok{return}\NormalTok{ \{}
        \StringTok{\textquotesingle{}length\textquotesingle{}}\NormalTok{: }\BuiltInTok{len}\NormalTok{(seq),}
        \StringTok{\textquotesingle{}gc\_content\textquotesingle{}}\NormalTok{: }\BuiltInTok{round}\NormalTok{(gc\_content, }\DecValTok{2}\NormalTok{),}
        \StringTok{\textquotesingle{}codon\_counts\textquotesingle{}}\NormalTok{: codon\_table}
\NormalTok{    \}}

\CommentTok{\# Example: Codon optimization for human expression}
\KeywordTok{def}\NormalTok{ optimize\_codon\_usage(seq, codon\_usage\_table):}
    \CommentTok{"""Replace rare codons with optimal human{-}preferred codons."""}
\NormalTok{    optimized }\OperatorTok{=}\NormalTok{ []}
    \ControlFlowTok{for}\NormalTok{ i }\KeywordTok{in} \BuiltInTok{range}\NormalTok{(}\DecValTok{0}\NormalTok{, }\BuiltInTok{len}\NormalTok{(seq)}\OperatorTok{{-}}\DecValTok{2}\NormalTok{, }\DecValTok{3}\NormalTok{):}
\NormalTok{        codon }\OperatorTok{=}\NormalTok{ seq[i:i}\OperatorTok{+}\DecValTok{3}\NormalTok{]}
\NormalTok{        amino\_acid }\OperatorTok{=}\NormalTok{ translate\_codon(codon)  }\CommentTok{\# Assume a translation function}
        \CommentTok{\# Choose the most frequent codon for this amino acid in human cells}
\NormalTok{        optimal\_codon }\OperatorTok{=} \BuiltInTok{max}\NormalTok{(codon\_usage\_table[amino\_acid],}
\NormalTok{                           key}\OperatorTok{=}\NormalTok{codon\_usage\_table[amino\_acid].get)}
\NormalTok{        optimized.append(optimal\_codon)}
    \ControlFlowTok{return} \StringTok{\textquotesingle{}\textquotesingle{}}\NormalTok{.join(optimized)}

\CommentTok{\# Example: Predict secondary structure (simplified)}
\KeywordTok{def}\NormalTok{ predict\_secondary\_structure(seq):}
    \CommentTok{"""Calculate minimum free energy using a simplified nearest{-}neighbor model."""}
    \CommentTok{\# In practice, tools such as RNAfold (ViennaRNA) or NUPACK}
    \ControlFlowTok{return} \StringTok{"{-}10.5 kcal/mol"}  \CommentTok{\# illustrative fixed value for the textbook example}
\end{Highlighting}
\end{Shaded}

\textbf{Clinical impact}: The Pfizer-BioNTech BNT162b2 and Moderna
mRNA-1273 vaccines were developed in under a year from viral sequence
release to emergency use authorization, demonstrating the power of this
technology. Beyond COVID-19, mRNA vaccines are in development for
influenza, HIV, Zika, and personalized cancer vaccines.

\textbf{Concept Check 11.2}

\begin{quote}
\begin{enumerate}
\def\labelenumi{\arabic{enumi}.}
\tightlist
\item
  Why is the 5' cap important for mRNA vaccine stability and
  translation?
\item
  What is the role of lipid nanoparticles (LNPs) in mRNA vaccine
  delivery?
\item
  How does an mRNA vaccine activate both CD4+ helper T cells and CD8+
  cytotoxic T cells?
\item
  Explain why nucleoside-modified mRNA (e.g., pseudouridine) reduces
  immunogenicity compared to unmodified mRNA.
\item
  Design a simple Python function to calculate the GC content of an mRNA
  sequence and explain why GC content matters for vaccine stability.
\end{enumerate}
\end{quote}

\begin{center}\rule{0.5\linewidth}{0.5pt}\end{center}

\subsubsection{Leading and Lagging
Strands}\label{leading-and-lagging-strands}

DNA polymerase primarily adds nucleotides to the 3'-OH of an existing
primer, and can primarily synthesize in the 5' to 3' direction.
Therefore:

\begin{itemize}
\tightlist
\item
  \textbf{Leading strand}: synthesized continuously 5' to 3' in the same
  direction as fork movement. Requires primarily a single priming event.
  In eukaryotes, primarily synthesized by DNA Pol epsilon.
\item
  \textbf{Lagging strand}: synthesized discontinuously as
  \textbf{Okazaki fragments} (1,000-2,000 nt in prokaryotes; 100-200 nt
  in eukaryotes). Each fragment requires a new RNA primer laid down by
  primase. After extension by DNA Pol delta, the primers are removed by
  a multi-step process:

  \begin{enumerate}
  \def\labelenumi{\arabic{enumi}.}
  \tightlist
  \item
    \textbf{RNase H} degrades the RNA portion of the primer
  \item
    \textbf{FEN1} (flap endonuclease 1) removes the remaining 5' flap
    created when the upstream Okazaki fragment displaces the primer of
    the next fragment
  \item
    \textbf{DNA Pol delta} fills the gap
  \item
    \textbf{DNA Ligase I} seals the nick with a phosphodiester bond
    (consuming ATP)
  \end{enumerate}
\end{itemize}

\textbf{Replication fork speed}: \emph{E. coli} replicates its 4.6 Mb
genome from a single origin in approximately 40 minutes
(about 1,000 bp/s). Human cells replicate 3.2 Gb from
approximately 30,000 origins over 6-8 hours of S phase
(about 50 bp/s per fork, but thousands of forks operate
simultaneously).

\subsubsection{PCNA Sliding Clamp -- The Processivity
Engine}\label{pcna-sliding-clamp-the-processivity-engine}

PCNA is a homotrimeric ring (three identical subunits) that encircles
double-stranded DNA and slides freely. It acts as a \textbf{molecular
matchmaker}:

\begin{itemize}
\tightlist
\item
  \textbf{Face toward template junction}: binds DNA polymerase (Pol
  delta/epsilon), FEN1, and ligase I through a conserved PIP box motif
\item
  \textbf{Monoubiquitination at K164} (by RAD6/RAD18): recruits Y-family
  translesion synthesis (TLS) polymerases (Pol eta, Pol iota, Pol kappa,
  Rev1) when the fork stalls at DNA damage. TLS polymerases have relaxed
  active sites that can bypass lesions but at the cost of lower fidelity
\end{itemize}

\subsubsection{DNA Proofreading and
Fidelity}\label{dna-proofreading-and-fidelity}

The error rate of replication is remarkably low: approximately 1 error
per \(10^{9}\)-\(10^{10}\) base pairs copied. This fidelity results from
three layers:

\begin{enumerate}
\def\labelenumi{\arabic{enumi}.}
\tightlist
\item
  \textbf{Base selection by polymerase} (about \(10^{-5}\)
  error rate): The polymerase active site preferentially binds correctly
  paired nucleotides based on geometry and hydrogen bonding
\item
  \textbf{3' to 5' exonuclease proofreading}
  (about \(10^{-2}\) improvement): If a mismatch is
  incorporated, the polymerase pauses; the 3' to 5' exonuclease domain
  excises the incorrect nucleotide and the polymerase re-extends
\item
  \textbf{Mismatch repair (MMR)} (about \(10^{-3}\)
  improvement): Post-replicative scanning by MutS (MSH2/MSH6 in humans),
  which detects mismatches and small insertion/deletion loops
\end{enumerate}

\begin{equation}\text{Overall error rate} \approx 10^{-5} \times 10^{-2} \times 10^{-3} = 10^{-10} \text{ per bp per replication} \tag{12.1}\label{eq:error_rate}\end{equation}

\begin{quote}
\textbf{Clinical Connection: Lynch Syndrome (HNPCC)} Germline mutations
in MMR \hyperref[gl:gene]{\textbf{gene}}s (\emph{MLH1}, \emph{MSH2},
\emph{MSH6}, \emph{PMS2}) cause Lynch syndrome (hereditary non-polyposis
colorectal cancer). Affected individuals have a 50-80\% lifetime risk of
colorectal cancer and increased risk of endometrial, ovarian, and
gastric cancers. The hallmark is \textbf{microsatellite instability
(MSI)} -- expansion or contraction of short tandem repeats due to
uncorrected replication slippage. MSI-high tumors respond well to immune
checkpoint inhibitors (anti-PD-1), making Lynch syndrome diagnosis
therapeutically important.
\end{quote}

\subsubsection{Telomeres and Telomerase}\label{telomeres-and-telomerase}

\textbf{The end-replication problem}: Linear eukaryotic chromosomes
cannot replicate the very 5' end of the lagging strand (no upstream
sequence to place a primer). This means 50-200 bp of telomeric DNA are
lost with each cell division.

\textbf{Telomere structure}: Human telomeres consist of 5-15 kb of the
hexameric repeat TTAGGG, ending in a 150-200 nt single-stranded 3'
overhang. This overhang invades the upstream duplex to form a
\textbf{T-loop} (telomere loop), which is stabilized by the
\textbf{shelterin complex} (TRF1, TRF2, RAP1, TIN2, TPP1, POT1). The
T-loop hides the chromosome end from the DNA damage response, which
would otherwise recognize it as a double-strand break.

\textbf{The shelterin complex --- six-protein architecture and
function:}

{\def\LTcaptype{none} % do not increment counter
\begin{longtable}[]{@{}
  >{\raggedright\arraybackslash}p{(\linewidth - 4\tabcolsep) * \real{0.1842}}
  >{\raggedright\arraybackslash}p{(\linewidth - 4\tabcolsep) * \real{0.6053}}
  >{\raggedright\arraybackslash}p{(\linewidth - 4\tabcolsep) * \real{0.2105}}@{}}
\toprule\noalign{}
\begin{minipage}[b]{\linewidth}\raggedright
Subunit
\end{minipage} & \begin{minipage}[b]{\linewidth}\raggedright
DNA binding specificity
\end{minipage} & \begin{minipage}[b]{\linewidth}\raggedright
Function
\end{minipage} \\
\midrule\noalign{}
\endhead
\bottomrule\noalign{}
\endlastfoot
\textbf{TRF1} (telomeric repeat factor 1) & dsDNA TTAGGG (Myb domain) &
Negatively regulates telomerase access; controls telomere length \\
\textbf{TRF2} & dsDNA TTAGGG (Myb domain) & Stabilises T-loop;
suppresses ATM-mediated DNA damage signalling at telomeres \\
\textbf{TIN2} & TRF1, TRF2 partner (no DNA binding) & Bridge protein,
stabilises whole complex \\
\textbf{RAP1} & TRF2 partner & NF-\texorpdfstring{\ensuremath{\kappa}}{kappa}B pathway crosstalk; chromatin
contacts \\
\textbf{TPP1} & TIN2 partner; binds POT1 & Recruits telomerase via TEL
patch \\
\textbf{POT1} & ssDNA TTAGGG-overhang (OB-fold domain) & Suppresses ATR;
protects 3\texorpdfstring{'}{'} overhang \\
\end{longtable}
}

\textbf{T-loop formation in detail.} The 3\texorpdfstring{'}{'} G-rich overhang (150--200
nt) folds back and invades the upstream duplex, displacing one strand
and forming a D-loop (displacement loop) within the T-loop. TRF2 wraps
the DNA, inducing the strand exchange. The structure has been visualised
by electron microscopy and STORM super-resolution. Shelterin coats the
entire telomere --- about about 1 shelterin per 100 bp.

\textbf{Shelterin's six functions in protecting chromosome ends:} 1.
\textbf{Suppress ATM kinase} at telomeres (TRF2-dependent) 2.
\textbf{Suppress ATR kinase} at telomeres (POT1-dependent) 3.
\textbf{Inhibit NHEJ} at chromosome ends (preventing fusions; TRF2-RAP1)
4. \textbf{Inhibit HR} between telomeres (TRF2; preventing
telomere-telomere recombination) 5. \textbf{Recruit telomerase} in cells
where it is active (TPP1-TEL patch) 6. \textbf{Mediate semi-conservative
replication} through telomeres (CST complex with shelterin)

\textbf{Telomere replication problem in detail.} Two distinct issues: 1.
\textbf{Lagging-strand end-replication}: The very 5\texorpdfstring{'}{'} end of the lagging
strand cannot be replicated because there is no upstream template for
primer placement. After RNA primer removal, about 50--100 nt
are lost from the 5\texorpdfstring{'}{'} end of the lagging-strand-complement. 2.
\textbf{Leading-strand end-resection}: The leading strand is synthesised
continuously to the chromosome end, but specialised exonucleases then
resect the 5\texorpdfstring{'}{'} end to generate the 3\texorpdfstring{'}{'} overhang necessary for T-loop
formation.

Result: each replication cycle, about 50--200 bp of TTAGGG
repeats are lost from telomeres in the absence of telomerase.

\textbf{G-quadruplex regulation.} The G-rich telomere sequence can fold
into intramolecular G-quadruplexes (four G's stacked via Hoogsteen
bonds). G-quadruplexes are stabilised by K\texorpdfstring{\textsuperscript{+}}{+} but inhibit telomerase
access. Drugs that stabilise telomeric G-quadruplexes (e.g., RHPS4,
telomestatin, BRACO19) cause progressive telomere shortening and
senescence in cancer cells --- they are anti-cancer leads but limited by
toxicity in normal stem cells.

\textbf{G-quadruplexes}: The G-rich telomeric sequence can fold into
\textbf{G-quadruplex} structures -- stacked tetrads of four guanines
connected by Hoogsteen hydrogen bonds, stabilized by monovalent cations
(K\(^+\)). G-quadruplexes inhibit telomerase and are targets for
anti-cancer drug development.

\textbf{Telomerase RNA template --- TERC mechanism in detail.} TERC is a
451-nt RNA with several functional regions: (i) a \textbf{template
region} (5\texorpdfstring{'}{'}-CUAACCCUAAC-3\texorpdfstring{'}{'}, about 11 nt) that base-pairs with
the telomere 3\texorpdfstring{'}{'} overhang and dictates the TTAGGG sequence added; (ii) a
\textbf{pseudoknot} (essential structural motif); (iii) a
\textbf{CR4/CR5} stem-loop that binds TERT; (iv) an \textbf{H/ACA scaRNA
box} at the 3\texorpdfstring{'}{'} end (essential for stability and dyskerin binding). The
H/ACA box recruits dyskerin (DKC1) --- mutations in DKC1 cause X-linked
dyskeratosis congenita.

\textbf{Telomerase}: A specialized reverse transcriptase that extends
the 3' overhang:

\begin{equation}5'\text{-}(TTAGGG)_n\text{-OH-3'} \xrightarrow{\text{TERC RNA template: 3'-AAUCCC-5'}} 5'\text{-}(TTAGGG)_{n+1}\text{-OH-3'} \tag{12.2}\label{eq:telomerase}\end{equation}

\begin{itemize}
\tightlist
\item
  \textbf{TERT} (telomerase reverse transcriptase): the catalytic
  protein subunit
\item
  \textbf{TERC} (telomerase RNA component): 451-nt RNA containing the
  template sequence (in humans, 5'-CUAACCCUAAC-3') used to add TTAGGG
  repeats
\item
  After TERT extends the 3' overhang, conventional DNA Pol alpha/primase
  synthesizes the complementary C-rich strand
\end{itemize}

\textbf{Telomeres, aging, and cancer}: Somatic cells lack telomerase
activity and undergo progressive telomere shortening, eventually
reaching a critically short length that triggers replicative senescence
(the Hayflick limit, about 50-70 divisions for human
fibroblasts). This is a tumor-suppressive mechanism. However,
approximately 85\% of cancers reactivate telomerase (often through TERT
\hyperref[gl:promoter]{\textbf{promoter}} mutations, especially C228T
and C250T), and approximately 15\% use \textbf{ALT (alternative
lengthening of telomeres)} via
\hyperref[gl:recombination]{\textbf{recombination}} between telomeres
(associated with ATRX/DAXX mutations, with C-circles as biomarkers).

\begin{quote}
\textbf{Clinical Connection: Telomere Biology Disorders} Mutations in
TERT, TERC, DKC1, or other telomere maintenance genes cause
\textbf{dyskeratosis congenita} (DC), characterized by bone marrow
failure, pulmonary fibrosis, and mucosal leukoplakia. Patients have
critically short telomeres and premature aging features. Related
conditions include idiopathic pulmonary fibrosis and aplastic anemia,
forming a spectrum called \textbf{telomere biology disorders (TBDs)}.
\end{quote}

\textbf{Concept Check 11.1}

\begin{quote}
\begin{enumerate}
\def\labelenumi{\arabic{enumi}.}
\tightlist
\item
  Why is the lagging strand synthesized discontinuously?
\item
  What would happen to a eukaryotic cell if you knocked out FEN1?
\item
  A patient has MSI-high colorectal cancer. What is the most likely
  genetic defect? Why might this patient respond to immunotherapy?
\item
  Explain why telomerase is an attractive anti-cancer drug target, and
  what complications might arise from telomerase inhibition.
\end{enumerate}
\end{quote}

\begin{center}\rule{0.5\linewidth}{0.5pt}\end{center}

\subsection{The Cell Cycle}\label{the-cell-cycle-1}

The \textbf{cell cycle} is the ordered series of events by which a cell
duplicates its contents and divides. In a typical mammalian cell, the
cycle takes approximately 24 hours.

\begin{figure}[htbp]
\centering
\includegraphics[width=0.82\linewidth,height=4.2in,keepaspectratio]{../figures/mermaid_inline/inline_mermaid_0070_1ddf35a993ca.png}
\caption{The eukaryotic cell cycle. G1 contains the restriction point (R), after which the cell is committed to division independent of mitogenic signals. Each transition is guarded by checkpoints.}
\end{figure}

\subsubsection{Cell Cycle Phases}\label{cell-cycle-phases}

{\def\LTcaptype{none} % do not increment counter
\begin{longtable}[]{@{}
  >{\raggedright\arraybackslash}p{(\linewidth - 4\tabcolsep) * \real{0.1842}}
  >{\raggedright\arraybackslash}p{(\linewidth - 4\tabcolsep) * \real{0.5000}}
  >{\raggedright\arraybackslash}p{(\linewidth - 4\tabcolsep) * \real{0.3158}}@{}}
\toprule\noalign{}
\begin{minipage}[b]{\linewidth}\raggedright
Phase
\end{minipage} & \begin{minipage}[b]{\linewidth}\raggedright
Duration (typical)
\end{minipage} & \begin{minipage}[b]{\linewidth}\raggedright
Key Events
\end{minipage} \\
\midrule\noalign{}
\endhead
\bottomrule\noalign{}
\endlastfoot
\textbf{G1} (Gap 1) & 8-12 h & Cell growth;
\hyperref[gl:organelle]{\textbf{organelle}} duplication; assessment of
extracellular signals (mitogens, growth factors); passage through the
\textbf{restriction point (R)} commits the cell to divide \\
\textbf{S} (Synthesis) & 6-8 h & DNA replication; centrosome
duplication; histone synthesis;
\hyperref[gl:euchromatin]{\textbf{euchromatin}} replicates early,
\hyperref[gl:heterochromatin]{\textbf{heterochromatin}} late \\
\textbf{G2} (Gap 2) & 3-4 h & Continued growth; preparation for mitosis;
final DNA damage check \\
\textbf{M} (Mitosis) & about 1 h & Chromosome condensation,
alignment, separation, and cytokinesis \\
\textbf{G0} (Quiescence) & Variable & Reversible exit from the cell
cycle; \hyperref[gl:neuron]{\textbf{neuron}}s and muscle cells may
remain in G0 permanently \\
\end{longtable}
}

\subsubsection{Cyclin-CDK Complexes: The Cell Cycle
Engine}\label{cyclin-cdk-complexes-the-cell-cycle-engine}

The cell cycle is driven by the sequential activation and destruction of
\textbf{cyclin-CDK (cyclin-dependent kinase)} complexes:

{\def\LTcaptype{none} % do not increment counter
\begin{longtable}[]{@{}
  >{\raggedright\arraybackslash}p{(\linewidth - 6\tabcolsep) * \real{0.1346}}
  >{\raggedright\arraybackslash}p{(\linewidth - 6\tabcolsep) * \real{0.3654}}
  >{\raggedright\arraybackslash}p{(\linewidth - 6\tabcolsep) * \real{0.3077}}
  >{\raggedright\arraybackslash}p{(\linewidth - 6\tabcolsep) * \real{0.1923}}@{}}
\toprule\noalign{}
\begin{minipage}[b]{\linewidth}\raggedright
Phase
\end{minipage} & \begin{minipage}[b]{\linewidth}\raggedright
Cyclin-CDK Complex
\end{minipage} & \begin{minipage}[b]{\linewidth}\raggedright
Key Substrates
\end{minipage} & \begin{minipage}[b]{\linewidth}\raggedright
Function
\end{minipage} \\
\midrule\noalign{}
\endhead
\bottomrule\noalign{}
\endlastfoot
G1 (early) & Cyclin D - CDK4/6 & RB (retinoblastoma protein) & Partial
phosphorylation of RB; releases E2F to transcribe S-phase genes \\
G1/S & Cyclin E - CDK2 & RB (hyperphosphorylation), p27 & Commits cell
to S phase; triggers centrosome duplication \\
S & Cyclin A - CDK2 & Cdc6, Cdt1 (for degradation) & Fires replication
origins; prevents re-licensing \\
G2/M & Cyclin B - CDK1 (MPF) & Lamins, condensins, Golgi proteins &
Nuclear envelope breakdown; chromosome condensation; spindle
formation \\
\end{longtable}
}

\textbf{Regulation of CDK activity:}

\begin{itemize}
\tightlist
\item
  \textbf{Cyclin binding}: CDKs are inactive monomers without their
  cyclin partner
\item
  \textbf{Activating phosphorylation}: CAK (CDK-activating kinase)
  phosphorylates the T-loop
\item
  \textbf{Inhibitory phosphorylation}: Wee1 kinase phosphorylates CDK1
  at Tyr15 and Thr14 (inactive); Cdc25 phosphatase removes these
  (activates)
\item
  \textbf{CKI (CDK inhibitors)}: p21\(^{CIP1}\) (induced by p53),
  p27\(^{KIP1}\), p16\(^{INK4a}\) bind and inactivate specific
  cyclin-CDK complexes
\item
  \textbf{Cyclin destruction}: APC/C
  (\hyperref[gl:anaphase]{\textbf{anaphase}}-promoting
  complex/cyclosome) ubiquitinates cyclins A and B for proteasomal
  degradation, ending mitosis
\end{itemize}

\textbf{Quantitative CDK activity through the cell cycle.} Single-cell
biosensor measurements have provided a quantitative picture of CDK
activity dynamics:

{\def\LTcaptype{none} % do not increment counter
\begin{longtable}[]{@{}
  >{\raggedright\arraybackslash}p{(\linewidth - 6\tabcolsep) * \real{0.1739}}
  >{\raggedright\arraybackslash}p{(\linewidth - 6\tabcolsep) * \real{0.1957}}
  >{\raggedright\arraybackslash}p{(\linewidth - 6\tabcolsep) * \real{0.1739}}
  >{\raggedright\arraybackslash}p{(\linewidth - 6\tabcolsep) * \real{0.4565}}@{}}
\toprule\noalign{}
\begin{minipage}[b]{\linewidth}\raggedright
Cell-cycle phase
\end{minipage} & \begin{minipage}[b]{\linewidth}\raggedright
Total CDK activity
\end{minipage} & \begin{minipage}[b]{\linewidth}\raggedright
Active complexes
\end{minipage} & \begin{minipage}[b]{\linewidth}\raggedright
Approximate timing (mammalian, 24-h cycle)
\end{minipage} \\
\midrule\noalign{}
\endhead
\bottomrule\noalign{}
\endlastfoot
Early G1 (mitogen-stimulated) & about 5 \% of peak & Cyclin
D-CDK4/6 (low) & 0--4 h post-mitosis \\
Restriction point & about 20 \% & Cyclin D-CDK4/6 + early
Cyclin E-CDK2 & 6--8 h post-mitosis \\
G1/S transition & about 40 \% & Cyclin E-CDK2 high; Cyclin
A-CDK2 rising & 8--10 h \\
Mid-S & about 50 \% & Cyclin A-CDK2 dominant & 10--14 h \\
G2 & about 70 \% & Cyclin A-CDK2 + early Cyclin B-CDK1 & 14--18
h \\
G2/M transition & about 90 \% & Cyclin B-CDK1 (= MPF) rises
sharply & 18--20 h \\
Mitosis (M-phase) & \textbf{100 \% (peak)} & Cyclin B-CDK1 + Aurora A/B
kinases & 20--22 h \\
Anaphase / Telophase & Drops to \textless{} 10 \% & APC/C-Cdc20 destroys
cyclin B & 22--24 h \\
\end{longtable}
}

The peak-to-trough ratio is about 10-fold. This sharp
threshold-like activity profile is encoded by: - \textbf{Bistability} in
CDK1 activation: Wee1 inhibits CDK1; Cdc25 activates CDK1; both have
positive-feedback loops on themselves through CDK1 phosphorylation,
creating a bistable switch. - \textbf{Timing precision} from cyclin
synthesis and destruction: cyclin A, cyclin B, and securin most have
D-box and KEN-box motifs that direct APC/C-mediated ubiquitination at
precise times.

\subsubsection{Specific Checkpoint Failures in
Cancer}\label{specific-checkpoint-failures-in-cancer}

Cancer is essentially a disease of cell-cycle deregulation. Key
checkpoint mutations and their consequences:

\textbf{(1) RB1 loss (G1/S checkpoint failure)} - Normal:
hypophosphorylated RB binds E2F \texorpdfstring{\ensuremath{\to}}{->} E2F cannot transcribe S-phase genes
(DNA Pol \texorpdfstring{\ensuremath{\alpha}}{alpha}, PCNA, MCM, thymidine kinase) - Cancer: biallelic RB1
mutations or deletions (retinoblastoma; small-cell lung cancer; bladder;
osteosarcoma) - Consequence: constitutive E2F activity; cells enter S
phase regardless of mitogen status - Therapeutic relevance:
\textbf{CDK4/6 inhibitors (palbociclib, ribociclib, abemaciclib) are
ineffective in RB-null tumours} because the entire downstream pathway is
already deregulated.

\textbf{(2) p53 loss (G1/S and G2/M checkpoints failure)} - Normal: p53
senses DNA damage \texorpdfstring{\ensuremath{\to}}{->} induces p21 \texorpdfstring{\ensuremath{\to}}{->} arrests CDK2 (G1/S) and CDK1 (G2/M);
also induces apoptosis (BAX, PUMA, NOXA) - Cancer: about 50 \%
of cancers carry p53 mutations (most frequent mutated gene); often
dominant-negative (mutant p53 oligomerises with wild-type to inactivate
it) or gain-of-function (p53\textsuperscript{R175H},
p53\textsuperscript{R273H} stabilise oncogenic chromatin states) -
Consequence: cells with damaged DNA continue dividing; aneuploidy;
centrosome amplification; chromosomal instability - Therapeutic
relevance: PRIMA-1, APR-246 (eprenetapopt) re-fold mutant p53
conformations; HDM2 inhibitors (idasanutlin, milademetan) re-stabilise
wild-type p53 in tumours where p53 is wild-type but suppressed by HDM2
amplification.

\textbf{(3) WEE1/CDC25 deregulation (G2/M checkpoint failure)} - Normal:
Wee1 phosphorylates CDK1 (Tyr15) \texorpdfstring{\ensuremath{\to}}{->} inactive; Cdc25 dephosphorylates \texorpdfstring{\ensuremath{\to}}{->}
activates; balance controls mitotic entry timing. - Cancer-relevant
deregulation: - \textbf{CDC25 amplification} in many tumours (breast,
lung, pancreatic) \texorpdfstring{\ensuremath{\to}}{->} premature CDK1 activation \texorpdfstring{\ensuremath{\to}}{->} DNA-damaged cells enter
mitosis prematurely \texorpdfstring{\ensuremath{\to}}{->} mitotic catastrophe (paradoxically anti-tumour in
some contexts). - \textbf{WEE1 inhibition (adavosertib/MK-1775)} is
being trialled in p53-mutant tumours: without p53 (the G1/S guard), the
cell relies on Wee1-mediated G2/M arrest; inhibiting Wee1 forces entry
into mitosis with damaged DNA \texorpdfstring{\ensuremath{\to}}{->} mitotic catastrophe. This is
\textbf{synthetic lethality} with p53 loss.

\textbf{(4) APC/C-Cdc20 / SAC failures} - Normal: SAC (mitotic
checkpoint complex MCC: Mad2 + BubR1 + Bub3) prevents anaphase until
kinetochores are properly attached. - Cancer: \textbf{BUB1B (BubR1)}
loss-of-function causes mosaic variegated aneuploidy syndrome (childhood
cancers, microcephaly). - Tumour-targeting drugs: - \textbf{Taxanes
(paclitaxel, docetaxel)} stabilise microtubules \texorpdfstring{\ensuremath{\to}}{->} triggering SAC arrest
\texorpdfstring{\ensuremath{\to}}{->} mitotic catastrophe + apoptosis - \textbf{Vincristine, vinblastine,
vinorelbine} depolymerise microtubules \texorpdfstring{\ensuremath{\to}}{->} SAC arrest

\textbf{(5) MCM helicase / replication-stress tolerance} - Normal: ATR
activates Chk1 in response to stalled forks \texorpdfstring{\ensuremath{\to}}{->} suppresses dormant origin
firing under stress - Cancer: tumours often have constitutively elevated
replication stress (driven by oncogenic \emph{MYC} or \emph{RAS}); this
provides a vulnerability: - \textbf{ATR inhibitors (berzosertib,
ceralasertib)} block the cancer cell's ability to manage replication
stress \texorpdfstring{\ensuremath{\to}}{->} genome catastrophe - \textbf{CHK1 inhibitors (LY2603618,
prexasertib)} similarly target the replication-stress response

\textbf{Summary table --- checkpoint failures and therapeutic
targeting:}

{\def\LTcaptype{none} % do not increment counter
\begin{longtable}[]{@{}
  >{\raggedright\arraybackslash}p{(\linewidth - 4\tabcolsep) * \real{0.4091}}
  >{\raggedright\arraybackslash}p{(\linewidth - 4\tabcolsep) * \real{0.1364}}
  >{\raggedright\arraybackslash}p{(\linewidth - 4\tabcolsep) * \real{0.4545}}@{}}
\toprule\noalign{}
\begin{minipage}[b]{\linewidth}\raggedright
Checkpoint failure
\end{minipage} & \begin{minipage}[b]{\linewidth}\raggedright
Cancer
\end{minipage} & \begin{minipage}[b]{\linewidth}\raggedright
Therapeutic strategy
\end{minipage} \\
\midrule\noalign{}
\endhead
\bottomrule\noalign{}
\endlastfoot
RB1 loss (G1/S) & Retinoblastoma; SCLC; bladder & CDK4/6 inhibitors
INEFFECTIVE (downstream); requires alternative \\
p53 loss (G1/S, G2/M) & about 50 \% of cancers & Re-fold mutant
p53 (APR-246); WEE1 inhibitor (synthetic lethal); HDM2 inhibitor (if p53
wt) \\
BRCA1/2 loss (HR) & Breast, ovarian, prostate, pancreatic & PARP
inhibitors (synthetic lethal) \\
MMR loss (replication fidelity) & Lynch syndrome; MSI-high CRC,
endometrial & Anti-PD-1 (high neoantigen load) \\
BUB1B / SAC loss & Mosaic variegated aneuploidy syndrome & Microtubule
targeting (taxanes, vinca alkaloids) \\
MYC-driven replication stress & Lymphoma, neuroblastoma & ATR/CHK1
inhibitors \\
Cyclin D / CDK4-6 amplification & Mantle cell lymphoma; breast (HR+);
melanoma & CDK4/6 inhibitors (palbociclib, ribociclib, abemaciclib) ---
effective if RB intact \\
\end{longtable}
}

\subsubsection{Cell Cycle Checkpoints}\label{cell-cycle-checkpoints}

Checkpoints are surveillance mechanisms that halt the cell cycle if
conditions are not met:

\textbf{G1/S Checkpoint (Restriction Point)}: - The \textbf{RB-E2F
pathway} is the molecular switch. In early G1, hypophosphorylated RB
binds and inhibits E2F transcription factors. Mitogenic signals activate
cyclin D-CDK4/6, which partially phosphorylates RB. Cyclin E-CDK2
completes RB hyperphosphorylation, releasing E2F to transcribe genes for
DNA replication (DNA Pol alpha, PCNA, MCM, thymidine kinase). - This is
the ``point of no return'': once past R, the cell will complete S phase
even if mitogens are withdrawn.

\textbf{Intra-S Checkpoint}: - Monitors replication fork stalling and
DNA damage during S phase - Stalled forks activate \textbf{ATR kinase}
(via RPA-coated ssDNA), which phosphorylates \textbf{Chk1} kinase - Chk1
inhibits Cdc25A phosphatase, preventing activation of cyclin A-CDK2 and
cyclin E-CDK2 - Slows origin firing to prevent catastrophic fork
collapse

\textbf{G2/M Checkpoint}: - Prevents entry into mitosis with damaged or
incompletely replicated DNA - \textbf{ATM} (activated by DSBs) and
\textbf{ATR} (activated by ssDNA/stalled forks) phosphorylate
\textbf{Chk2} and \textbf{Chk1} respectively - Chk1/Chk2 phosphorylate
Cdc25C, creating a 14-3-3 binding site that sequesters Cdc25C in the
\hyperref[gl:cytoplasm]{\textbf{cytoplasm}}, preventing it from
activating cyclin B-CDK1 (MPF) - Simultaneously, p53 induces p21, which
directly inhibits cyclin B-CDK1

\textbf{Spindle Assembly Checkpoint (SAC)}: - Prevents anaphase until
most kinetochores are properly attached to spindle microtubules -
Unattached kinetochores catalyze formation of the \textbf{mitotic
checkpoint complex (MCC)}: Mad2, BubR1, Bub3, and Cdc20 - MCC binds and
inhibits \textbf{Cdc20}, preventing activation of APC/C - A single
unattached kinetochore is sufficient for complete arrest (Rieder et
al.~1995) - Once most kinetochores achieve bipolar attachment with
tension, SAC is silenced; Cdc20 activates APC/C, which ubiquitinates
securin and cyclin B - Securin degradation releases \textbf{separase},
which cleaves cohesin (Scc1/Rad21 subunit) -- triggering anaphase

\subsubsection{The DNA Damage Response: ATM/ATR and
p53}\label{the-dna-damage-response-atmatr-and-p53}

\begin{figure}[htbp]
\centering
\includegraphics[width=0.82\linewidth,height=4.2in,keepaspectratio]{../figures/mermaid_inline/inline_mermaid_0071_6d60c17cec8c.png}
\caption{The DNA damage response pathway. ATM responds to double-strand breaks; ATR responds to replication stress and single-stranded DNA. Both converge on p53, which decides between cell cycle arrest (allowing repair), apoptosis, or senescence depending on damage severity.}
\end{figure}

\textbf{p53: Guardian of the Genome}

p53 is the most frequently mutated gene in human cancers (approximately
50\% of cancers carry \emph{TP53} mutations). Normal p53 function:

\begin{itemize}
\tightlist
\item
  \textbf{Normally kept at low levels} by MDM2 (an E3 ubiquitin ligase)
  that targets p53 for proteasomal degradation (half-life approximately
  20 min)
\item
  \textbf{Stabilized by damage signals}: ATM/ATR-mediated
  phosphorylation of p53 at Ser15 and Chk2-mediated phosphorylation at
  Ser20 disrupt the p53-MDM2 interaction
\item
  \textbf{Acts as a transcription factor} (binds DNA as a tetramer) to
  induce:

  \begin{itemize}
  \tightlist
  \item
    \textbf{p21} (CDKN1A): CDK inhibitor that arrests the cell cycle at
    G1/S and G2/M
  \item
    \textbf{GADD45}: DNA repair and G2 arrest
  \item
    \textbf{14-3-3sigma}: G2 arrest
  \item
    \textbf{BAX, PUMA, NOXA}: pro-apoptotic Bcl-2 family members (if
    damage is irreparable)
  \item
    \textbf{MDM2} itself: negative feedback loop
  \end{itemize}
\end{itemize}

\begin{quote}
\textbf{Clinical Connection: Li-Fraumeni Syndrome} Germline
\hyperref[gl:heterozygous]{\textbf{heterozygous}} \emph{TP53} mutations
cause Li-Fraumeni syndrome (LFS). Patients have a greater than 90\%
lifetime cancer risk, developing sarcomas, breast cancer, brain tumors,
adrenocortical carcinoma, and leukemia -- often before age 30. The
remaining wild-type \hyperref[gl:allele]{\textbf{allele}} is lost
somatically (loss of heterozygosity, LOH), consistent with Knudson's
two-hit hypothesis. Many \emph{TP53} mutations are
\hyperref[gl:dominant]{\textbf{dominant}}-negative (mutant p53
tetramerizes with wild-type p53 and inactivates it) or gain-of-function.
\end{quote}

\subsubsection{Cancer and the Cell
Cycle}\label{cancer-and-the-cell-cycle}

Cancer is fundamentally a disease of uncontrolled cell cycle
progression. Key deregulated pathways:

{\def\LTcaptype{none} % do not increment counter
\begin{longtable}[]{@{}
  >{\raggedright\arraybackslash}p{(\linewidth - 6\tabcolsep) * \real{0.1940}}
  >{\raggedright\arraybackslash}p{(\linewidth - 6\tabcolsep) * \real{0.2388}}
  >{\raggedright\arraybackslash}p{(\linewidth - 6\tabcolsep) * \real{0.2836}}
  >{\raggedright\arraybackslash}p{(\linewidth - 6\tabcolsep) * \real{0.2836}}@{}}
\toprule\noalign{}
\begin{minipage}[b]{\linewidth}\raggedright
Gene/Pathway
\end{minipage} & \begin{minipage}[b]{\linewidth}\raggedright
Normal Function
\end{minipage} & \begin{minipage}[b]{\linewidth}\raggedright
Cancer Alteration
\end{minipage} & \begin{minipage}[b]{\linewidth}\raggedright
Associated Cancer
\end{minipage} \\
\midrule\noalign{}
\endhead
\bottomrule\noalign{}
\endlastfoot
RB1 & G1/S checkpoint & Loss-of-function mutations & Retinoblastoma
(childhood eye cancer); also lung, bladder \\
CDKN2A (p16) & Inhibits CDK4/6 & Deletion or methylation & Melanoma,
pancreatic, many others \\
CDK4 & G1 progression & Amplification or activating mutation & Melanoma,
glioblastoma \\
Cyclin D1 & G1 progression & Amplification (11q13) & Breast cancer,
mantle cell lymphoma \\
TP53 & G1/S and G2/M arrest; apoptosis & Loss-of-function or
dominant-negative mutation & Approximately 50\% of cancers \\
MDM2 & p53 degradation & Amplification & Sarcomas, other cancers \\
MYC & Transcription factor for cell cycle genes &
Amplification/translocation & Burkitt lymphoma (t(8;14)),
neuroblastoma \\
\end{longtable}
}

\textbf{CDK4/6 inhibitors} (palbociclib, ribociclib, abemaciclib) are
now standard of care in HR+/HER2- metastatic breast cancer. They block
cyclin D-CDK4/6, preventing RB phosphorylation and E2F release.

\textbf{Concept Check 11.2}

\begin{quote}
\begin{enumerate}
\def\labelenumi{\arabic{enumi}.}
\tightlist
\item
  What is the restriction point, and how does the RB-E2F pathway control
  it?
\item
  Why is p53 called the ``guardian of the genome''? Describe three
  outcomes of p53 activation.
\item
  Explain why CDK4/6 inhibitors would be ineffective in RB-null tumors.
\item
  A tumor cell has a homozygous deletion of CDKN2A (p16). Predict the
  effect on the cell cycle.
\end{enumerate}
\end{quote}

\begin{center}\rule{0.5\linewidth}{0.5pt}\end{center}

\subsection{Mitosis: Molecular Events at Each
Stage}\label{mitosis-molecular-events-at-each-stage}

Mitosis distributes one copy of each duplicated chromosome to each
daughter cell. In most mammalian cells it takes approximately 1 hour.

\textbf{Stages of mitosis -- molecular detail:}

{\def\LTcaptype{none} % do not increment counter
\begin{longtable}[]{@{}
  >{\raggedright\arraybackslash}p{(\linewidth - 4\tabcolsep) * \real{0.2059}}
  >{\raggedright\arraybackslash}p{(\linewidth - 4\tabcolsep) * \real{0.2941}}
  >{\raggedright\arraybackslash}p{(\linewidth - 4\tabcolsep) * \real{0.5000}}@{}}
\toprule\noalign{}
\begin{minipage}[b]{\linewidth}\raggedright
Phase
\end{minipage} & \begin{minipage}[b]{\linewidth}\raggedright
Duration
\end{minipage} & \begin{minipage}[b]{\linewidth}\raggedright
Molecular Events
\end{minipage} \\
\midrule\noalign{}
\endhead
\bottomrule\noalign{}
\endlastfoot
\textbf{Prophase} & about 25 min & Condensin I/II complexes
compact chromosomes (about 10,000-fold); centrosomes
(duplicated in S phase) begin migrating to opposite poles; Aurora A
kinase activates centrosome maturation \\
\textbf{Prometaphase} & about 15 min & Nuclear envelope
breakdown (NEB) by CDK1-mediated lamin phosphorylation (lamins A, B, C
depolymerize); kinetochores (CENP-A, CENP-C, KNL1, Ndc80 complex)
capture spindle microtubules; chromosomes undergo rapid
poleward/anti-poleward oscillations \\
\textbf{Metaphase} & about 5 min & Chromosomes aligned at the
metaphase plate (equidistant from both poles); bipolar attachment
generates tension across sister kinetochores; SAC monitors attachment \\
\textbf{Anaphase A} & about 5 min & Separase cleaves cohesin
Scc1 subunit; sister chromatids separate; kinetochore microtubules
shorten (depolymerization at kinetochore end via Pacman mechanism)
pulling chromosomes poleward at about 1 micrometer/min \\
\textbf{Anaphase B} & about 5 min & Interpolar microtubules
slide apart (kinesin-5/Eg5 pushes poles apart); astral microtubules pull
poles via cortical dynein \\
\textbf{Telophase} & about 10 min & Nuclear envelope re-forms
around each chromosome set (nuclear pore complexes reassemble);
chromosomes decondense; nuclear lamina reforms \\
\textbf{Cytokinesis} & about 5 min & \textbf{Animals}:
contractile ring of \hyperref[gl:actin]{\textbf{actin}} + myosin II
constricts at the cleavage furrow (position determined by central
spindle/RhoA); \textbf{Plants}: cell plate forms from vesicle fusion at
the phragmoplast \\
\end{longtable}
}

\textbf{Cohesin and separase}: Cohesin (Smc1/Smc3/Scc1/Scc3) holds
sister chromatids together from S phase. During prophase, most arm
cohesin is removed by the ``prophase pathway'' (Wapl + Polo-like kinase
1 phosphorylation of SA2 subunit). Centromeric cohesin is protected by
\textbf{shugoshin (Sgo1)} and protein phosphatase 2A (PP2A). At anaphase
onset, APC/C-Cdc20 ubiquitinates securin, releasing separase, which
cleaves the remaining centromeric Scc1 -- triggering sister chromatid
separation.

\begin{quote}
\textbf{Clinical Connection: Aneuploidy and Cancer Therapeutics} Defects
in the spindle assembly checkpoint cause chromosome instability (CIN), a
hallmark of cancer. Taxol (paclitaxel) and vincristine exploit
microtubule dynamics to kill cancer cells: taxol stabilizes microtubules
and vincristine destabilizes them, both activating the SAC and
triggering mitotic arrest followed by apoptosis. Mosaic variegated
aneuploidy (MVA) syndrome, caused by biallelic BubR1 mutations, leads to
growth retardation, microcephaly, and childhood cancers.
\end{quote}

\begin{center}\rule{0.5\linewidth}{0.5pt}\end{center}

\subsection{Meiosis: Generating Genetic
Diversity}\label{meiosis-generating-genetic-diversity}

Meiosis produces haploid (n) \hyperref[gl:gamete]{\textbf{gamete}}s from
\hyperref[gl:diploid]{\textbf{diploid (2n)}} precursors via two
sequential divisions without an intervening S phase:

\subsubsection{Meiosis I (Reductive
Division)}\label{meiosis-i-reductive-division}

{\def\LTcaptype{none} % do not increment counter
\begin{longtable}[]{@{}
  >{\raggedright\arraybackslash}p{(\linewidth - 2\tabcolsep) * \real{0.3889}}
  >{\raggedright\arraybackslash}p{(\linewidth - 2\tabcolsep) * \real{0.6111}}@{}}
\toprule\noalign{}
\begin{minipage}[b]{\linewidth}\raggedright
Stage
\end{minipage} & \begin{minipage}[b]{\linewidth}\raggedright
Key Events
\end{minipage} \\
\midrule\noalign{}
\endhead
\bottomrule\noalign{}
\endlastfoot
\textbf{Prophase I} & The longest and most complex stage; subdivided
into: leptotene (chromosomes condense; telomeres attach to nuclear
envelope forming ``bouquet''), zygotene (homologous chromosomes pair by
synapsis via the \textbf{synaptonemal complex}), pachytene (crossing
over occurs; \textbf{Spo11} makes programmed DSBs; \textbf{Dmc1} and
\textbf{Rad51} mediate strand invasion; about 1-3 crossovers
per chromosome), diplotene (synaptonemal complex disassembles; homologs
connected primarily at \textbf{chiasmata}), diakinesis (further
condensation; nuclear envelope breaks down) \\
\textbf{Metaphase I} & Bivalents (tetrads) align at metaphase plate;
homologs orient randomly (independent assortment) \\
\textbf{Anaphase I} & Homologs separate (reductive division); sister
chromatids remain joined at
\hyperref[gl:centromere]{\textbf{centromere}} (centromeric cohesin
protected by shugoshin/Rec8) \\
\textbf{Telophase I + Cytokinesis} & Two haploid cells, each with n
chromosomes (still replicated = sister chromatids joined) \\
\end{longtable}
}

\subsubsection{Meiosis II (Equational
Division)}\label{meiosis-ii-equational-division}

Essentially a mitotic division of each Meiosis I product. Sister
chromatids separate at anaphase II, producing 4 haploid cells total.

\subsubsection{Sources of Genetic
Diversity}\label{sources-of-genetic-diversity}

\begin{enumerate}
\def\labelenumi{\arabic{enumi}.}
\tightlist
\item
  \textbf{Independent assortment}: \(2^{23}\) = approximately 8.4
  million possible chromosome combinations for humans
\item
  \textbf{Crossing over}: generates novel recombinant haplotypes;
  average 1-3 crossovers per chromosome pair per meiosis = approximately
  30-50 total crossovers per meiosis
\item
  \textbf{Fertilization}: \(2 \times 8.4 \times 10^6\) combinations,
  further diversified by crossing over
\end{enumerate}

\textbf{Concept Check 11.3}

\begin{quote}
\begin{enumerate}
\def\labelenumi{\arabic{enumi}.}
\tightlist
\item
  Compare and contrast mitosis and meiosis I in terms of chromosome
  behavior at the metaphase plate.
\item
  Why must centromeric cohesin be protected during meiosis I but cleaved
  during meiosis II?
\item
  Calculate the number of possible gamete chromosome combinations for an
  organism with 2n = 8.
\end{enumerate}
\end{quote}

\begin{center}\rule{0.5\linewidth}{0.5pt}\end{center}

\subsection{DNA Damage and Repair}\label{dna-damage-and-repair}

DNA sustains approximately 10,000-100,000 lesions per cell per day from
endogenous (hydrolysis, oxidation, replication errors) and exogenous
(UV, chemicals, radiation) sources.

\subsubsection{Types of DNA Damage}\label{types-of-dna-damage}

{\def\LTcaptype{none} % do not increment counter
\begin{longtable}[]{@{}
  >{\raggedright\arraybackslash}p{(\linewidth - 4\tabcolsep) * \real{0.2549}}
  >{\raggedright\arraybackslash}p{(\linewidth - 4\tabcolsep) * \real{0.1569}}
  >{\raggedright\arraybackslash}p{(\linewidth - 4\tabcolsep) * \real{0.5882}}@{}}
\toprule\noalign{}
\begin{minipage}[b]{\linewidth}\raggedright
Damage Type
\end{minipage} & \begin{minipage}[b]{\linewidth}\raggedright
Source
\end{minipage} & \begin{minipage}[b]{\linewidth}\raggedright
Frequency (per cell per day)
\end{minipage} \\
\midrule\noalign{}
\endhead
\bottomrule\noalign{}
\endlastfoot
Depurination (AP sites) & Spontaneous hydrolysis &
about 5,000-10,000 \\
Deamination (C to U) & Spontaneous hydrolysis &
about 100-500 \\
Oxidized bases (8-oxoG) & Reactive oxygen species (ROS) &
about 10,000 \\
Single-strand breaks & ROS, topoisomerase failure &
about 20,000-40,000 \\
Double-strand breaks & Ionizing radiation, replication fork collapse &
about 10-50 \\
Pyrimidine dimers (CPD) & UV light (280-320 nm) & Variable (sun exposure
dependent) \\
Alkylated bases & Endogenous S-adenosylmethionine; alkylating agents &
about 3,000 \\
\end{longtable}
}

\subsubsection{Repair Pathways}\label{repair-pathways}

\textbf{Direct reversal}: - \textbf{Photolyase} (not in humans): uses
visible light energy to directly cleave cyclobutane pyrimidine dimers -
\textbf{MGMT} (O\(^6\)-methylguanine-DNA methyltransferase): transfers
the methyl group from O\(^6\)-methylguanine to an active site cysteine
(suicide enzyme -- each MGMT molecule works primarily once)

\textbf{Base excision repair (BER)}: 1. \textbf{DNA glycosylase} (11
types in humans, each recognizing specific lesions) cleaves the
N-glycosidic bond, creating an AP site 2. \textbf{APE1} (AP
endonuclease) cuts the phosphodiester backbone 5' to the AP site 3.
\textbf{DNA Pol beta} fills the single-nucleotide gap 4. \textbf{XRCC1 +
DNA Ligase III} seals the nick

\textbf{Nucleotide excision repair (NER)}: - Repairs bulky,
helix-distorting lesions (CPDs, 6-4 photoproducts, cisplatin adducts) -
\textbf{Global genome NER (GG-NER)}: XPC-RAD23B detects helix distortion
genome-wide - \textbf{Transcription-coupled NER (TC-NER)}: CSB detects
stalled RNA Pol II on the transcribed strand - Both pathways converge:
TFIIH (XPB + XPD helicases) unwinds about 30 bp; XPF-ERCC1 cuts
5' and XPG cuts 3'; the about 24-32 nt fragment is removed; DNA
Pol delta/epsilon fills the gap; ligase seals

\textbf{Mismatch repair (MMR)}: - MSH2/MSH6 (MutSalpha) detects single
base mismatches and 1-nt insertion/deletion loops - MSH2/MSH3 (MutSbeta)
detects 2-13 nt insertion/deletion loops - MLH1/PMS2 (MutLalpha) makes
an incision in the newly synthesized strand - Exonuclease 1 degrades the
error-containing strand; DNA Pol delta resynthesizes

\textbf{Double-strand break repair}: - \textbf{NHEJ} (active throughout
cell cycle): Ku70/Ku80 bind broken ends; DNA-PKcs recruited; Artemis
processes overhangs; XRCC4-Ligase IV seals. Fast but error-prone. -
\textbf{Homologous recombination (HR)} (active in S/G2): MRN complex
resects 5' ends; RPA coats 3' ssDNA; BRCA2 loads RAD51; strand invasion
into sister chromatid; error-free repair.

The choice between NHEJ and HR is regulated by \textbf{53BP1} (favors
NHEJ) versus \textbf{BRCA1} (promotes resection for HR). In S/G2, CDK
phosphorylation activates CtIP to promote resection.

\begin{quote}
\textbf{Clinical Connection: PARP Inhibitors and Synthetic Lethality}
\textbf{PARP inhibitors} (olaparib, niraparib, talazoparib) exploit
\textbf{synthetic lethality}: PARP1 is essential for single-strand break
repair. In BRCA1/2-deficient tumor cells (which cannot perform HR), PARP
inhibition forces reliance on error-prone NHEJ, leading to genomic
catastrophe and cell death. Normal cells with intact BRCA1/2 survive.
PARP inhibitors are standard of care for BRCA-mutant breast, ovarian,
prostate, and pancreatic cancers.
\end{quote}

\textbf{Concept Check 11.4}

\begin{quote}
\begin{enumerate}
\def\labelenumi{\arabic{enumi}.}
\tightlist
\item
  Why is NHEJ considered error-prone while HR is high-fidelity?
\item
  Explain the concept of synthetic lethality as it applies to PARP
  inhibitor therapy.
\item
  A patient with xeroderma pigmentosum develops melanoma at age 12. What
  repair pathway is deficient, and what type of DNA damage accumulates?
\item
  Compare BER and NER in terms of the type of damage repaired and the
  size of DNA removed.
\end{enumerate}
\end{quote}

\begin{center}\rule{0.5\linewidth}{0.5pt}\end{center}

\subsection{Worked Example: Telomere Shortening and the Hayflick
Limit}\label{worked-example-telomere-shortening-and-the-hayflick-limit}

\textbf{Problem}: A human fibroblast is born with telomeres of
\(10{,}000 \text{ bp}\) of TTAGGG repeats. In the absence of telomerase
it loses \(100 \text{ bp}\) of telomeric DNA per division (within the
stated \(50\)--\(200 \text{ bp}\) end-replication range). Replicative
senescence is triggered once telomeres reach a critically short length
of \(\approx 4{,}000 \text{ bp}\).

\begin{enumerate}
\def\labelenumi{(\alph{enumi})}
\tightlist
\item
  How many divisions can this cell undergo before senescence?
\end{enumerate}

\textbf{Solution}: The expendable telomere reserve is the starting
length minus the senescence threshold; the number of divisions is that
reserve divided by the loss per division:

\begin{equation}n = \frac{10{,}000 \text{ bp} - 4{,}000 \text{ bp}}{100 \text{ bp/division}} = \frac{6{,}000 \text{ bp}}{100 \text{ bp/division}} = 60 \text{ divisions} \tag{12.5}\label{eq:telomere_shortening_divisions}\end{equation}

\begin{enumerate}
\def\labelenumi{(\alph{enumi})}
\setcounter{enumi}{1}
\tightlist
\item
  How much total telomeric DNA is lost over that lifespan?
\end{enumerate}

\begin{equation}\Delta L = 60 \text{ divisions} \times 100 \text{ bp/division} = 6{,}000 \text{ bp} \tag{12.6}\label{eq:telomere_total_loss}\end{equation}

The predicted \(60\) divisions sits squarely within the measured
Hayflick limit of \(\approx 50\)--\(70\) divisions for human
fibroblasts, showing how the end-replication problem imposes a built-in
replicative clock that limits the proliferative capacity of somatic
cells and acts as a tumour-suppressive barrier.

\subsection{Worked Example: Replication Fork
Dynamics}\label{worked-example-replication-fork-dynamics}

\textbf{Problem}: An \emph{E. coli} cell with a 4.6 Mb genome replicates
bidirectionally from a single origin. The replication fork speed is
1,000 bp/s.

\begin{enumerate}
\def\labelenumi{(\alph{enumi})}
\tightlist
\item
  How long does it take to replicate the entire genome?
\end{enumerate}

\textbf{Solution}: Bidirectional replication means two forks travel in
opposite directions. Each fork must cover half the genome:

\begin{equation}t = \frac{4.6 \times 10^6 \text{ bp}}{2 \times 1000 \text{ bp/s}} = 2300 \text{ s} \approx 38.3 \text{ min} \tag{12.3}\label{eq:replication_time_bacteria}\end{equation}

\begin{enumerate}
\def\labelenumi{(\alph{enumi})}
\setcounter{enumi}{1}
\tightlist
\item
  \emph{E. coli} can divide every 20 minutes under optimal conditions.
  How is this possible?
\end{enumerate}

\textbf{Solution}: \emph{E. coli} uses \textbf{nested replication}
(multi-fork replication). A new round of replication initiates at
\emph{oriC} before the previous round completes. Thus, a cell dividing
every 20 minutes has multiple replication forks active simultaneously.

\begin{enumerate}
\def\labelenumi{(\alph{enumi})}
\setcounter{enumi}{2}
\tightlist
\item
  The eukaryotic replication fork speed is about 50 bp/s. If
  the human genome (3.2 Gb) had a single origin, how long would S phase
  take?
\end{enumerate}

\begin{equation}t = \frac{3.2 \times 10^9}{2 \times 50} = 3.2 \times 10^7 \text{ s} \approx 370 \text{ days} \tag{12.4}\label{eq:replication_time_euk}\end{equation}

This is why eukaryotes require approximately 30,000-50,000 origins to
complete replication in approximately 8 hours.

\begin{center}\rule{0.5\linewidth}{0.5pt}\end{center}

\subsection{Computational Bridge}\label{computational-bridge-11}

Complementary strand generation is the arithmetic behind
semiconservative replication:

\begin{Shaded}
\begin{Highlighting}[]
\ImportTok{from}\NormalTok{ biology.genetics }\ImportTok{import}\NormalTok{ dna\_complement}

\NormalTok{leading }\OperatorTok{=} \StringTok{"ATGCGATCGATCG"}
\NormalTok{lagging\_template }\OperatorTok{=}\NormalTok{ dna\_complement(leading)}
\BuiltInTok{print}\NormalTok{(lagging\_template)}
\end{Highlighting}
\end{Shaded}

\begin{quote}
\textbf{Clinical / systems note:} Defects in MMR and HR surveillance
convert replication infidelity into organ-specific cancer risks --- the
clinical mirror of error rates you estimate from polymerase fidelity.
\end{quote}

\begin{center}\rule{0.5\linewidth}{0.5pt}\end{center}

\subsection{Current Evidence and Frontier
Biology}\label{current-evidence-and-frontier-biology-14}

For \textbf{DNA Replication and the Cell Cycle}, frontier biology
belongs inside the evidence logic of the chapter. Molecular genetics now
spans single-reference sequences, telomere-to-telomere assemblies,
pangenome graphs, long-read sequencing, CRISPR medicines, and ethical
deployment. The core reading question is this: replication claims should
connect polymerase accuracy, checkpoint timing, damage response, and
cancer relevance.

\begin{itemize}
\tightlist
\item
  \textbf{What to verify:} identify the observation, model, assay, or
  dataset that would make the claim stronger or weaker.
\item
  \textbf{What to qualify:} state the scale, organism, cell type,
  environmental condition, or population where the claim is expected to
  hold.
\item
  \textbf{What to compare:} test at least one alternative explanation,
  baseline, or null model before treating the pattern as causal.
\item
  \textbf{What to cite:} distinguish primary evidence, review synthesis,
  public dataset, and institutional guidance; for recent or numeric
  claims, prefer the source closest to the measurement and state what
  has changed since it was published.
\end{itemize}

When a genomic claim depends on a reference, ask whether short reads,
structural variants, ancestry representation, phasing, or clinical
validation could change the interpretation
\citep{humanpangenome2023,fda2023casgevy,fda2024casgevythalassemia}.

\textbf{Source practice:} For genomics and editing claims, distinguish
discovery from clinical actionability, and cite reference resources,
regulatory records, or primary editing studies close to the claim
\citep{humanpangenome2023,fda2026casgevy,chalumeau2025primeediting}.

\subsection{Summary}\label{summary-14}

\begin{itemize}
\tightlist
\item
  \textbf{DNA structure}: B-form double helix with antiparallel strands;
  major groove allows sequence-specific protein binding; negative
  backbone charge requires cationic neutralization.
\item
  \textbf{Replication}: Semiconservative, bidirectional from licensed
  origins; leading strand continuous (Pol epsilon), lagging strand as
  Okazaki fragments (Pol delta); fidelity maintained by polymerase
  selectivity, 3' to 5' proofreading, and MMR (approximately
  \(10^{-10}\) error rate).
\item
  \textbf{Telomeres}: TTAGGG repeats protected by shelterin and T-loops;
  telomerase (TERT + TERC) extends 3' overhang in germ/stem cells;
  reactivated in most cancers; ALT pathway in approximately 15\%.
\item
  \textbf{Cell cycle}: G1-S-G2-M driven by sequential cyclin-CDK
  activation; checkpoints (G1/S restriction point, intra-S, G2/M, SAC)
  ensure fidelity; p53 integrates damage signals to decide between
  arrest, senescence, and apoptosis.
\item
  \textbf{Mitosis}: Prophase through cytokinesis; cohesin holds sister
  chromatids; separase cleaves at anaphase; SAC ensures bipolar
  attachment via MCC-mediated APC/C inhibition.
\item
  \textbf{Meiosis}: Meiosis I (reductive: homolog separation, crossing
  over at chiasmata) + Meiosis II (equational: sister chromatid
  separation) produce 4 haploid gametes; independent assortment and
  crossing over generate diversity.
\item
  \textbf{DNA repair}: BER (damaged bases), NER (bulky adducts), MMR
  (mismatches), NHEJ and HR (DSBs); defects cause cancer predisposition
  syndromes (Lynch, XP, BRCA-mutant cancers, Li-Fraumeni).
\item
  \textbf{Connections:} See \cref{sec:unit_IV_gene_expression} for gene
  expression downstream of replicated DNA, Unit V for segregation of
  sister chromatids, and Unit VII for bacterial replication timing.
\end{itemize}

\begin{center}\rule{0.5\linewidth}{0.5pt}\end{center}

\subsection{Review Questions}\label{review-questions-14}

\begin{enumerate}
\def\labelenumi{\arabic{enumi}.}
\tightlist
\item
  Compare the roles of DnaA/DnaB in \emph{E. coli} and ORC/MCM2-7 in
  eukaryotes during replication initiation.
\item
  Explain how the licensing system prevents re-replication. What happens
  if geminin is depleted?
\item
  Describe the three layers of replication fidelity and calculate the
  expected number of mutations per human cell division (3.2 Gb genome,
  error rate \(10^{-10}\) per bp).
\item
  A cell line has homozygous loss of RB1. Predict the effect on G1/S
  checkpoint function. Would CDK4/6 inhibitors be effective?
\item
  Compare the substrates and outcomes of APC/C-Cdc20 versus APC/C-Cdh1
  in mitotic exit.
\item
  Explain why crossing over in meiosis I is essential for proper
  chromosome segregation (beyond generating diversity).
\item
  A patient has a homozygous mutation in MSH2. What type of genomic
  instability would you observe? What cancer screening would you
  recommend?
\item
  Compare NHEJ and HR repair pathways. Why is 53BP1 vs.~BRCA1 balance
  critical for pathway choice?
\item
  Explain the mechanism of PARP inhibitor synthetic lethality in
  BRCA1/2-mutant tumors.
\item
  A researcher finds that a tumor has a TERT promoter C228T mutation.
  What is the functional consequence, and why does this confer a growth
  advantage?
\item
  Starting from a 50\% GC template, predict whether leading vs.~lagging
  synthesis faces more \textbf{secondary structure} risk in highly
  repetitive DNA. How does pol \texorpdfstring{\ensuremath{\delta}}{delta} processivity help?
\item
  Using the mutation-per-division estimate in Q3, scale to \textbf{1012}
  cells in a renewing epithelium. When does clonal selection become
  inevitable even if each division is nearly perfect?
\end{enumerate}

\begin{center}\rule{0.5\linewidth}{0.5pt}\end{center}

\subsection{Further Reading and Source
Notes}\label{further-reading-and-source-notes-14}

\begin{itemize}
\tightlist
\item
  Meselson \& Stahl (1958). The replication of \{DNA\} in \{Escherichia
  coli\}. \emph{Proceedings of the National Academy of Sciences}, 44.
\item
  Okazaki et al.~(1968). Mechanism of DNA chain growth: Discontinuous
  synthesis. \emph{Proceedings of the National Academy of Sciences}, 59.
\item
  Hartwell, Culotti, Pringle \& Reid (1974). Genetic control of the cell
  division cycle in yeast. \emph{Science}, 183.
\item
  Nurse (2000). A long twentieth century of the cell cycle and beyond.
  \emph{Cell}, 100.
\item
  Blackburn \& Gall (1978). A tandemly repeated sequence at the termini
  of the extrachromosomal rDNA in Tetrahymena. \emph{Journal of
  Molecular Biology}, 120.
\item
  Alberts et al.~(latest ed.). \emph{Molecular Biology of the Cell} (DNA
  replication and the cell cycle chapters). Garland Science.
\end{itemize}

\begin{center}\rule{0.5\linewidth}{0.5pt}\end{center}

\subsection{Key Terms}\label{key-terms-14}

\begin{enumerate}
\def\labelenumi{\arabic{enumi}.}
\tightlist
\item
  \textbf{Semiconservative replication} -- each daughter helix contains
  one parental and one newly synthesized strand
\item
  \textbf{Origin of replication} -- specific DNA sequence where
  replication initiates
\item
  \textbf{MCM helicase} -- hexameric helicase loaded during licensing;
  unwinds DNA at replication forks as part of the CMG complex
\item
  \textbf{PCNA} -- homotrimeric sliding clamp ring; tethers polymerases
  for processivity; K164-Ub recruits TLS polymerases
\item
  \textbf{Okazaki fragment} -- short DNA segment synthesized on the
  lagging strand (100-200 nt in eukaryotes)
\item
  \textbf{Telomerase} -- reverse transcriptase (TERT + TERC) that
  extends telomeric repeats using an RNA template
\item
  \textbf{Cyclin-CDK complex} -- heterodimer driving cell cycle
  transitions; cyclin provides specificity, CDK provides kinase activity
\item
  \textbf{Restriction point} -- point in late G1 after which the cell is
  committed to divide independent of mitogens
\item
  \textbf{Spindle assembly checkpoint (SAC)} -- monitors
  kinetochore-microtubule attachment; prevents premature anaphase via
  MCC
\item
  \textbf{Separase} -- protease that cleaves cohesin Scc1 subunit to
  trigger sister chromatid separation at anaphase
\item
  \textbf{p53} -- tumor suppressor transcription factor; induces cell
  cycle arrest, senescence, or apoptosis; mutated in about 50\%
  of cancers
\item
  \textbf{APC/C} -- E3 ubiquitin ligase that targets securin and mitotic
  cyclins for proteasomal degradation
\item
  \textbf{Synaptonemal complex} -- protein scaffold mediating homologous
  chromosome synapsis in meiosis I prophase
\item
  \textbf{Chiasma} -- physical connection between homologs at a
  crossover site during meiosis
\item
  \textbf{Mismatch repair (MMR)} -- post-replicative repair of base-base
  mismatches and small loops; defects cause Lynch syndrome
\item
  \textbf{Homologous recombination (HR)} -- error-free DSB repair using
  sister chromatid as template; requires BRCA1/2 and RAD51
\item
  \textbf{Synthetic lethality} -- cell death resulting from simultaneous
  loss of two compensatory pathways (basis for PARP inhibitor therapy)
\end{enumerate}

\begin{center}\rule{0.5\linewidth}{0.5pt}\end{center}

\subsubsection{Companion Source
Module}\label{companion-source-module-14}

\textbf{DNA Replication and the Cell Cycle} should leave a reproducible
trail from a biological claim to the code, figure, diagram, or
paper-based activity that can test it. Use the surfaces below to inspect
the chapter's assumptions, rerun the relevant model, or compare the
manuscript explanation with companion labs and figures.

{\def\LTcaptype{none} % do not increment counter
\begin{longtable}[]{@{}
  >{\raggedright\arraybackslash}p{(\linewidth - 2\tabcolsep) * \real{0.5000}}
  >{\raggedright\arraybackslash}p{(\linewidth - 2\tabcolsep) * \real{0.5000}}@{}}
\toprule\noalign{}
\begin{minipage}[b]{\linewidth}\raggedright
Surface
\end{minipage} & \begin{minipage}[b]{\linewidth}\raggedright
Use it for
\end{minipage} \\
\midrule\noalign{}
\endhead
\bottomrule\noalign{}
\endlastfoot
\breaktt{src/biology/genetics/genetics.py} (\texttt{dna\_complement},
\texttt{hamming\_distance}) & Test strand complementarity and
sequence-change reasoning. \\
\breaktt{src/mermaid/biology\_diagrams.py}
(\breaktt{dna\_replication\_diagram}, \breaktt{cell\_cycle\_diagram}) &
Keep replication forks, checkpoints, and cell-cycle stages aligned. \\
\end{longtable}
}

\textbf{Reproducibility check:} identify strand polarity, origin/fork
direction, repair pathway, and checkpoint readout before diagnosing
replication errors. \textbf{Cross-reference:} connect with
\cref{sec:unit_IV_mutations_and_genomics} and
\cref{sec:unit_V_chromosomal_inheritance}.

\newpage

\section{Gene Expression --- Transcription and
Translation}\label{gene-expression-transcription-and-translation}

\label{sec:unit_IV_gene_expression}

\begin{quote}
\textbf{Ch 13} \texorpdfstring{\ensuremath{\cdot}}{.} Level 2/3 \texorpdfstring{\ensuremath{\cdot}}{.} 60 min read \texorpdfstring{\ensuremath{\cdot}}{.} 100 min lecture \texorpdfstring{\ensuremath{\cdot}}{.}
Prerequisites: \cref{sec:unit_IV_dna_replication_and_cell_cycle}
\end{quote}

\subsection{Learning Objectives}\label{learning-objectives-15}

\begin{enumerate}
\def\labelenumi{\arabic{enumi}.}
\tightlist
\item
  Describe the mechanism of
  \hyperref[gl:transcription]{\textbf{transcription}} in prokaryotes and
  \hyperref[gl:eukaryote]{\textbf{eukaryote}}s, including
  \hyperref[gl:promoter]{\textbf{promoter}} recognition and elongation.
\item
  Compare prokaryotic and eukaryotic RNA polymerases in terms of subunit
  composition, promoter elements, and regulation.
\item
  Explain RNA processing in eukaryotes: 5' capping, splicing
  (spliceosome mechanism), and 3' polyadenylation.
\item
  Describe alternative splicing and its role in proteome diversity,
  including the mechanism of nonsense-mediated decay.
\item
  \textbf{Predict} the polypeptide produced from a given mRNA sequence,
  applying the genetic code and wobble pairing, including the effect of
  a single-base insertion on the reading frame.
\item
  Describe post-translational modifications and the ubiquitin-proteasome
  degradation pathway.
\item
  \textbf{Distinguish} transcriptional from post-transcriptional
  \hyperref[gl:gene]{\textbf{gene}} regulation (miRNA, siRNA, lncRNA,
  circRNA) given Northern-blot versus Western-blot data.
\item
  Describe epigenetic regulation: DNA methylation,
  \hyperref[gl:histone-code]{\textbf{histone code}}, and their roles in
  development and disease.
\item
  \textbf{Predict} lac- and trp-\hyperref[gl:operon]{\textbf{operon}}
  transcriptional output given specified inducer, repressor, and
  corepressor states.
\item
  Describe the role of Hox genes and homeodomain transcription factors
  in development.
\end{enumerate}

\subsubsection{Study Blueprint}\label{study-blueprint-15}

\begin{itemize}
\tightlist
\item
  \textbf{Big idea:} Gene expression translates sequence information
  into regulated RNA and protein products.
\item
  \textbf{Core concepts:} transcription, translation, codons, RNA
  processing.
\item
  \textbf{Framework alignment:} Vision \& Change: Information flow,
  exchange, and storage, Structure and function; AP Biology: Information
  Storage and Transmission, Systems Interactions; NGSS-style topics:
  Inheritance and Variation of Traits, Structure and Function.
\item
  \textbf{Model or quantitative lens:} Reading-frame, codon, and
  expression-output calculations.
\item
  \textbf{Data skill:} Convert DNA/RNA sequence data into predicted
  molecular products.
\item
  \textbf{Practice cadence:} Concept Explanation, Questions and Methods,
  Argumentation.
\item
  \textbf{Common misconception to repair:} A gene is not simply a
  protein recipe; context controls when, where, and how much product
  appears.
\item
  \textbf{Primary lab:} \cref{sec:lab_unit_IV_gene_expression}.
\item
  \textbf{Question bank:} \cref{sec:q_unit_IV_gene_expression}.
\item
  \textbf{Transfer task:} Apply expression logic to mutations,
  biotechnology, development, and disease diagnostics.
\item
  \textbf{Bridge to computation:}
  \breaktt{biology.genetics.genetics.translate\_mrna}.
\end{itemize}

\begin{center}\rule{0.5\linewidth}{0.5pt}\end{center}

\begin{quote}
\textbf{Opening Vignette --- Solving the Problem of the Messenger}

In 1961, Francis Crick and colleagues knew that information flowed from
DNA to \hyperref[gl:protein]{\textbf{protein}}, but the molecular
intermediary was a mystery. Working with Sydney Brenner and François
Jacob at Cambridge, Matthew Meselson performed a deceptively simple
experiment: he mixed radioactively labelled
\hyperref[gl:ribosome]{\textbf{ribosome}}s from bacteria infected with
T4 phage with unlabelled ribosomes, then centrifuged the extract through
a density gradient. The new viral proteins appeared on \emph{old}
unlabelled ribosomes --- proving that ribosomes are non-specific
``machines,'' and that the specificity must come from a short-lived
``messenger'' molecule. Brenner named it messenger RNA. Within a decade,
the entire transcription and translation apparatus --- RNA polymerase,
transfer RNA, \hyperref[gl:codon]{\textbf{codon}}s, the genetic code ---
had been decoded. This vignette illustrates how a single elegant
experiment can overturn received wisdom and open an entire research
programme.
\end{quote}

\subsection{Transcription}\label{transcription}

\subsubsection{Overview: The Central
Dogma}\label{overview-the-central-dogma}

Transcription produces RNA from a DNA template. The \textbf{central
dogma} \citep{crick1958} describes the flow of genetic information:

\begin{equation}\text{DNA} \xrightarrow{\text{transcription}} \text{RNA} \xrightarrow{\text{translation}} \text{Protein} \tag{13.1}\label{eq:central_dogma}\end{equation}

Additional information flows include: \textbf{reverse transcription}
(RNA to DNA; retroviruses, telomerase), \textbf{RNA replication}
(RNA-dependent RNA polymerases in RNA viruses), and \textbf{direct
translation regulation} (riboswitches, IREs).

\subsubsection{Prokaryotic
Transcription}\label{prokaryotic-transcription}

\textbf{RNA polymerase (RNAP)} in \emph{E. coli}: core
\hyperref[gl:enzyme]{\textbf{enzyme}} (\(\alpha_2\beta\beta'\omega\)) +
sigma (\texorpdfstring{\ensuremath{\sigma}}{sigma}) factor:

\begin{itemize}
\tightlist
\item
  The \texorpdfstring{\ensuremath{\sigma}}{sigma} factor recognizes the \textbf{promoter} -- specifically two
  conserved hexameric sequences:

  \begin{itemize}
  \tightlist
  \item
    \textbf{-10 box} (Pribnow box): consensus TATAAT (named after David
    Pribnow)
  \item
    \textbf{-35 box}: consensus TTGACA
  \end{itemize}
\item
  Promoter strength correlates with similarity to consensus; strong
  promoters (e.g., rRNA genes) match closely at both positions
\item
  The \(\sigma^{70}\) factor is the ``housekeeping'' sigma; alternative
  sigma factors direct RNAP to stress-response promoters (e.g.,
  \(\sigma^{32}\) for heat shock, \(\sigma^{54}\) for nitrogen
  starvation)
\end{itemize}

\textbf{Transcription cycle}: 1. \textbf{Closed complex}: RNAP + \texorpdfstring{\ensuremath{\sigma}}{sigma} binds
promoter DNA (double-stranded) 2. \textbf{Open complex}:
about 15 bp around the -10 region melt (AT-rich, easier to
denature); transcription bubble forms 3. \textbf{Initiation}: RNAP
synthesizes short abortive transcripts (2-12 nt); \texorpdfstring{\ensuremath{\sigma}}{sigma} factor released
after about 10 nt (``promoter escape'') 4. \textbf{Elongation}:
RNAP moves along template 3' to 5', synthesizing RNA 5' to 3' at
about 40-80 nt/s; transcription bubble is about 12-14
bp; about 8 bp RNA-DNA hybrid within the bubble 5.
\textbf{Termination}: two mechanisms

\textbf{Termination mechanisms}: - \textbf{Rho-independent (intrinsic)}:
GC-rich inverted repeat in nascent RNA forms a stem-loop hairpin
followed by a run of about 6-8 U residues (weak rU-dA base
pairs) -- hairpin destabilizes RNAP-RNA contact, U-tract promotes
dissociation - \textbf{Rho-dependent}: Rho protein (hexameric helicase)
loads onto rut (Rho utilization) sites on the nascent RNA, translocates
5' to 3' along the RNA (powered by ATP hydrolysis), catches up when RNAP
pauses, and unwinds the RNA-DNA hybrid to release the transcript

\subsubsection{Eukaryotic Transcription}\label{eukaryotic-transcription}

\textbf{Three nuclear RNA polymerases:}

{\def\LTcaptype{none} % do not increment counter
\begin{longtable}[]{@{}
  >{\raggedright\arraybackslash}p{(\linewidth - 6\tabcolsep) * \real{0.2143}}
  >{\raggedright\arraybackslash}p{(\linewidth - 6\tabcolsep) * \real{0.2381}}
  >{\raggedright\arraybackslash}p{(\linewidth - 6\tabcolsep) * \real{0.2619}}
  >{\raggedright\arraybackslash}p{(\linewidth - 6\tabcolsep) * \real{0.2857}}@{}}
\toprule\noalign{}
\begin{minipage}[b]{\linewidth}\raggedright
Feature
\end{minipage} & \begin{minipage}[b]{\linewidth}\raggedright
RNA Pol I
\end{minipage} & \begin{minipage}[b]{\linewidth}\raggedright
RNA Pol II
\end{minipage} & \begin{minipage}[b]{\linewidth}\raggedright
RNA Pol III
\end{minipage} \\
\midrule\noalign{}
\endhead
\bottomrule\noalign{}
\endlastfoot
\textbf{Products} & 45S pre-rRNA (18S, 5.8S, 28S rRNA) & Most mRNA; most
snRNA, miRNA & tRNA, 5S rRNA, U6 snRNA, 7SL RNA \\
\textbf{Location} & Nucleolus & Nucleoplasm & Nucleoplasm \\
\textbf{Number of subunits} & 14 & 12 & 17 \\
\textbf{Sensitivity to alpha-amanitin} & Resistant & Very sensitive (1
ug/mL) & Moderately sensitive (10 ug/mL) \\
\textbf{Promoter elements} & UCE + core (recognized by UBF + SL1) & TATA
box, Inr, DPE, BRE, \hyperref[gl:cpg-island]{\textbf{CpG island}}s &
Internal (type I, II) or upstream (type III) \\
\end{longtable}
}

The table is primarily the entry point. In eukaryotes, transcriptional
output depends on promoter architecture, enhancer-promoter looping,
chromatin accessibility, Mediator, pause release, and RNA processing.
RNA Pol II is especially regulated because its C-terminal domain is
phosphorylated in different patterns during initiation, elongation,
splicing, 3'-end formation, and termination. A strong gene-expression
claim therefore asks which polymerase is involved, which regulatory DNA
element is active, which chromatin state permits access, and whether
nascent-RNA, chromatin, or perturbation evidence supports the proposed
mechanism.

\subsubsection{Eukaryotic Promoter
Elements}\label{eukaryotic-promoter-elements}

The RNA Pol II \textbf{core promoter} can contain several elements (not
most are present in every promoter):

\begin{itemize}
\tightlist
\item
  \textbf{TATA box} (about 25-30 bp upstream of TSS; consensus
  TATAAAA): Bound by TBP (TATA-binding protein), a subunit of TFIID.
  Present in about 10-20\% of human gene promoters; associated
  with tissue-specific, highly regulated genes.
\item
  \textbf{Initiator element (Inr)}: Overlaps the TSS; consensus
  PyPyAN(T/A)PyPy; recognized by TAF1/TAF2 subunits of TFIID
\item
  \textbf{Downstream promoter element (DPE)}: Located
  about +28-33 relative to TSS; works with Inr in TATA-less
  promoters; recognized by TAF6/TAF9
\item
  \textbf{BRE (TFIIB recognition element)}: Immediately upstream or
  downstream of TATA box; directly contacts TFIIB
\item
  \textbf{CpG islands}: Regions of \textgreater200 bp with
  \textgreater50\% GC content and observed/expected CpG ratio
  \textgreater0.6; found at about 70\% of human gene promoters;
  typically unmethylated in active genes
\end{itemize}

\subsubsection{The Pre-Initiation Complex (PIC)
Assembly}\label{the-pre-initiation-complex-pic-assembly}

\begin{figure}[htbp]
\centering
\includegraphics[width=0.82\linewidth,height=4.2in,keepaspectratio]{../figures/mermaid_inline/inline_mermaid_0072_5707a07cccd9.png}
\caption{Assembly of the eukaryotic pre-initiation complex (PIC) at a TATA-containing promoter. TFIID binding initiates an ordered assembly cascade culminating in Pol II CTD phosphorylation and promoter escape.}
\end{figure}

{\def\LTcaptype{none} % do not increment counter
\begin{longtable}[]{@{}
  >{\raggedright\arraybackslash}p{(\linewidth - 4\tabcolsep) * \real{0.5152}}
  >{\raggedright\arraybackslash}p{(\linewidth - 4\tabcolsep) * \real{0.2424}}
  >{\raggedright\arraybackslash}p{(\linewidth - 4\tabcolsep) * \real{0.2424}}@{}}
\toprule\noalign{}
\begin{minipage}[b]{\linewidth}\raggedright
CTD Modification
\end{minipage} & \begin{minipage}[b]{\linewidth}\raggedright
Kinase
\end{minipage} & \begin{minipage}[b]{\linewidth}\raggedright
Effect
\end{minipage} \\
\midrule\noalign{}
\endhead
\bottomrule\noalign{}
\endlastfoot
Ser5-P & TFIIH (CDK7) & Recruits capping enzyme \\
Ser2-P & P-TEFb (CDK9) & Recruits splicing factors, polyadenylation
machinery \\
Ser7-P & CDK7 & snRNA gene-specific processing \\
\end{longtable}
}

\begin{quote}
\textbf{Clinical Connection: Transcription-Targeted Cancer Therapy}
Several cancers are ``transcription-addicted'' -- dependent on
super-enhancers driving oncogene expression. CDK7 inhibitors (THZ1) and
CDK9 inhibitors collapse super-enhancer-driven transcription programs.
BET bromodomain inhibitors (JQ1, OTX015) displace BRD4 from acetylated
\hyperref[gl:histone]{\textbf{histone}}s at enhancers, suppressing MYC
and other oncogene transcription. These represent a new class of
epigenetic cancer therapeutics.
\end{quote}

\begin{center}\rule{0.5\linewidth}{0.5pt}\end{center}

\newpage

\subsection{RNA Processing in
Eukaryotes}\label{rna-processing-in-eukaryotes}

\begin{figure}[htbp]
\centering
\includegraphics[width=0.82\linewidth,height=4.2in,keepaspectratio]{../figures/mermaid_inline/inline_mermaid_0073_8d677c52196e.png}
\caption{Eukaryotic mRNA processing pipeline. Pre-mRNA undergoes co-transcriptional capping, splicing, and polyadenylation before export. Transcripts with premature termination codons (PTCs) are eliminated by nonsense-mediated decay (NMD).}
\end{figure}

\subsubsection{5' Capping}\label{capping}

The 5' end of pre-mRNA receives a \textbf{7-methylguanosine cap}
(m\(^7\)G) added co-transcriptionally when the transcript is 20-30 nt
long. The capping reaction involves three enzymatic steps:

\begin{enumerate}
\def\labelenumi{\arabic{enumi}.}
\tightlist
\item
  \textbf{RNA triphosphatase} removes the terminal phosphate from the 5'
  end
\item
  \textbf{Guanylyltransferase} adds GMP in a 5'-5' triphosphate
  \hyperref[gl:linkage]{\textbf{linkage}} (unusual reverse orientation)
\item
  \textbf{Methyltransferase} adds a methyl group to the N-7 position of
  guanine (using SAM as donor)
\end{enumerate}

Cap functions: - Protects from 5' to 3' exonuclease degradation (Xrn1) -
Required for ribosome binding: eIF4E (eukaryotic initiation factor 4E)
specifically recognizes the m\(^7\)G cap - Signals nuclear export via
the cap-binding complex (CBC: CBP20/CBP80) - Enhances splicing of the
first \hyperref[gl:intron]{\textbf{intron}}

\subsubsection{3' Polyadenylation}\label{polyadenylation}

The 3' end is processed by cleavage and polyadenylation:

\begin{enumerate}
\def\labelenumi{\arabic{enumi}.}
\tightlist
\item
  \textbf{CPSF} (cleavage and polyadenylation specificity factor)
  recognizes the \textbf{AAUAAA} hexamer (the polyadenylation signal,
  present in about 90\% of human mRNAs)
\item
  \textbf{CstF} (cleavage stimulation factor) binds a downstream GU-rich
  or U-rich element
\item
  Cleavage occurs 10-30 nt downstream of AAUAAA
\item
  \textbf{Poly(A) polymerase (PAP)} adds about 200-250 adenine
  residues (not templated)
\end{enumerate}

Poly(A) tail functions: - Protects from 3' to 5' exonuclease degradation
(exosome) - Bound by \textbf{PABP} (poly(A)-binding protein), which
interacts with eIF4G to circularize the mRNA (closed-loop model),
enhancing translation efficiency - Progressive shortening
(\textbf{deadenylation} by CCR4-NOT complex) is the first and
rate-limiting step in most mRNA degradation pathways

\subsubsection{Pre-mRNA Splicing: The
Spliceosome}\label{pre-mrna-splicing-the-spliceosome}

Introns are removed from pre-mRNA by the \textbf{spliceosome}, one of
the cell's most complex molecular machines:

\begin{itemize}
\tightlist
\item
  \textbf{Composition}: 5 snRNPs (U1, U2, U4, U5, U6) +
  about 100-300 protein subunits; total mass 3-5 MDa
\item
  \textbf{snRNPs}: Each contains one snRNA (except U4/U6, which share) +
  about 7 Sm proteins + snRNP-specific proteins
\end{itemize}

\textbf{Consensus splice site sequences}: - 5' splice site (donor): GU
(almost invariant) - 3' splice site (acceptor): AG (almost invariant) -
Branch point: YNYURAY (Y = pyrimidine, R = purine, N = any); the
conserved \textbf{A} (adenosine) is the branch point nucleophile -
Polypyrimidine tract: about 15-20 pyrimidines between branch
point and 3' SS

\textbf{The splicing mechanism (two transesterification reactions)}:

\begin{enumerate}
\def\labelenumi{\arabic{enumi}.}
\tightlist
\item
  \textbf{First transesterification}: The 2'-OH of the branch point
  adenosine attacks the phosphodiester bond at the 5' splice site (GU).
  This produces:

  \begin{itemize}
  \tightlist
  \item
    A free 5' \hyperref[gl:exon]{\textbf{exon}} with a 3'-OH
  \item
    A lariat intermediate: the intron circularized via a 2'-5'
    phosphodiester bond at the branch point A
  \end{itemize}
\item
  \textbf{Second transesterification}: The 3'-OH of the free 5' exon
  attacks the phosphodiester bond at the 3' splice site (AG). This
  produces:

  \begin{itemize}
  \tightlist
  \item
    Ligated exons (5' exon joined to 3' exon)
  \item
    Released lariat intron (debranched and degraded)
  \end{itemize}
\end{enumerate}

\textbf{Spliceosome assembly cycle}: 1. \textbf{E complex}: U1 snRNP
base-pairs with 5' SS; U2AF65 binds polypyrimidine tract; U2AF35
contacts 3' SS AG; SF1/BBP binds branch point 2. \textbf{A complex}: U2
snRNP replaces SF1 at branch point (ATP-dependent); branch point A
bulges out 3. \textbf{B complex}: U4/U6.U5 tri-snRNP joins; massive
rearrangement 4. \textbf{B* complex}: U1 and U4 released; U6 replaces U1
at 5' SS; U6-U2 interaction positions catalytic metal ions 5. \textbf{C
complex}: First transesterification occurs 6. \textbf{C* complex}:
Second transesterification; exons ligated; lariat released 7.
\textbf{Post-spliceosomal complex}: snRNPs recycled; EJC (exon junction
complex) deposited about 20-24 nt upstream of each exon-exon
junction

\begin{quote}
\textbf{Clinical Connection: Spinal Muscular Atrophy and
Splice-Modifying Therapy} Spinal muscular atrophy (SMA) is caused by
loss of the \emph{SMN1} gene. The paralog \emph{SMN2} has a C-to-T
transition in exon 7 that weakens an exonic splicing enhancer, causing
about 90\% exon 7 skipping and production of a truncated,
unstable protein. \textbf{Nusinersen (Spinraza)} is an antisense
oligonucleotide (ASO) that binds an intronic splicing silencer (ISS-N1)
in \emph{SMN2} intron 7, promoting exon 7 inclusion and full-length SMN
protein production. Approved by FDA in 2016, it was one of the first
splice-modifying therapies. \textbf{Risdiplam (Evrysdi)}, an oral small
molecule that stabilizes U1 snRNP binding at the \emph{SMN2} exon 7 5'
SS, was approved in 2020.
\end{quote}

\subsubsection{Alternative Splicing}\label{alternative-splicing}

Alternative splicing generates multiple protein isoforms from a single
gene. An estimated \textgreater95\% of human multi-exon genes undergo
alternative splicing.

\textbf{Types of alternative splicing}:

{\def\LTcaptype{none} % do not increment counter
\begin{longtable}[]{@{}
  >{\raggedright\arraybackslash}p{(\linewidth - 4\tabcolsep) * \real{0.2222}}
  >{\raggedright\arraybackslash}p{(\linewidth - 4\tabcolsep) * \real{0.4444}}
  >{\raggedright\arraybackslash}p{(\linewidth - 4\tabcolsep) * \real{0.3333}}@{}}
\toprule\noalign{}
\begin{minipage}[b]{\linewidth}\raggedright
Type
\end{minipage} & \begin{minipage}[b]{\linewidth}\raggedright
Description
\end{minipage} & \begin{minipage}[b]{\linewidth}\raggedright
Example
\end{minipage} \\
\midrule\noalign{}
\endhead
\bottomrule\noalign{}
\endlastfoot
Cassette exon (exon skipping) & An exon is included or excluded &
Fibronectin: EDA and EDB exons (tissue-specific) \\
Alternative 5' splice site & Two different 5' SS compete & SV40 T/t
antigen \\
Alternative 3' splice site & Two different 3' SS compete & Adenovirus
E1A \\
Intron retention & An intron is retained in the mature mRNA & Common in
plants; increasingly recognized in mammals \\
Mutually exclusive exons & One of two (or more) exons is included,
rarely both & alpha-tropomyosin; DSCAM \\
\end{longtable}
}

\textbf{DSCAM in \emph{Drosophila}}: The Down syndrome cell adhesion
molecule gene contains 4 clusters of alternatively spliced exons: 12
alternatives for exon 4, 48 for exon 6, 33 for exon 9, and 2 for exon
17. Combinatorially: \(12 \times 48 \times 33 \times 2 = 38,016\)
possible mRNA isoforms -- more than twice the total number of genes in
the \emph{Drosophila} \hyperref[gl:genome]{\textbf{genome}}. DSCAM
diversity is essential for neuronal self-recognition and axon guidance.

\textbf{Quantitative scope in humans.} Genome-wide RNA-seq studies show
that \textbf{about 95 \% of human multi-exon genes undergo
alternative splicing}, with 92 \% producing more than one isoform per
gene at detectable levels (Wang et al.~2008; Pan et al.~2008). Average
number of detected isoforms per multi-exon gene: 4--7
(cell-type-dependent). Approximately 60 \% of alternative-splicing
events are tissue-specifically regulated. Estimates of the number of
distinct splicing events: - about 150,000 cassette-exon events
- about 60,000 alternative 5\texorpdfstring{'}{'}/3\texorpdfstring{'}{'} splice sites -
about 25,000 mutually exclusive exon events -
about 30,000 retained-intron events (especially in plants and
brain)

\textbf{Disease-causing splice site mutations --- clinical examples:} -
\textbf{\emph{BRCA1} IVS5+1G\textgreater A:} disrupts the consensus 5\texorpdfstring{'}{'}
splice site (GT \textgreater{} AT), causing exon 5 skipping and
frameshift; pathogenic for hereditary breast/ovarian cancer. -
\textbf{\emph{DMD} exon 51 deletion:} removes a frame-essential exon,
causing Duchenne muscular dystrophy. Antisense oligonucleotide therapies
(eteplirsen, casimersen) restore the reading frame by inducing exon 51
skipping. - \textbf{\emph{MAPT} IVS10+16C\textgreater T:} disrupts the
regulatory cis-element controlling tau exon 10 inclusion, causing
fronto-temporal dementia with parkinsonism (FTDP-17). Different intronic
mutations in \emph{MAPT} alter the 4R/3R tau ratio. -
\textbf{\emph{SMN2} C840T:} weakens an exonic splicing enhancer in
\emph{SMN2} exon 7, causing exon 7 skipping in 90 \% of \emph{SMN2}
transcripts and producing truncated, unstable SMN protein. This is the
basis of spinal muscular atrophy when \emph{SMN1} is also lost;
nusinersen (an antisense oligonucleotide) restores exon 7 inclusion. -
\textbf{\emph{LMNA} c.1824C\textgreater T (Hutchinson--Gilford
progeria):} activates a cryptic splice donor in exon 11, producing the
truncated lamin-A ``progerin'' that accumulates at the nuclear envelope,
driving accelerated ageing.

Approximately \textbf{15--20 \% of disease-causing point mutations}
disrupt splicing --- either at canonical splice sites (5\texorpdfstring{'}{'} GT, 3\texorpdfstring{'}{'} AG,
branch point), in regulatory elements (ESE, ESS, ISE, ISS), or by
creating cryptic splice sites. This is much higher than was estimated in
the 1990s, when primarily canonical splice-site mutations were
recognised.

\textbf{Regulatory elements controlling alternative splicing}: -
\textbf{Exonic splicing enhancers (ESEs)}: bound by SR proteins
(serine/arginine-rich); promote exon inclusion - \textbf{Exonic splicing
silencers (ESSs)}: bound by hnRNP proteins; promote exon skipping -
\textbf{Intronic splicing enhancers (ISEs)} and \textbf{intronic
splicing silencers (ISSs)}: analogous elements in introns

\subsubsection{Nonsense-Mediated Decay
(NMD)}\label{nonsense-mediated-decay-nmd}

NMD is an mRNA surveillance pathway that degrades transcripts containing
\textbf{premature termination codons (PTCs)}. This protects the cell
from producing truncated, potentially
\hyperref[gl:dominant]{\textbf{dominant}}-negative proteins.

\textbf{Mechanism}: During the first (``pioneer'') round of translation,
the ribosome encounters each exon junction complex (EJC). If translation
terminates \textgreater50-55 nt upstream of an EJC, the ribosome cannot
displace the EJC. The interaction of release factors (eRF1/eRF3) with
UPF1, followed by UPF1 interaction with the downstream EJC-bound
UPF2/UPF3, triggers UPF1 phosphorylation, mRNA decapping, deadenylation,
and degradation.

NMD degrades about 5-10\% of the human transcriptome, including
many alternatively spliced isoforms with PTCs (regulated unproductive
splicing and translation, or RUST).

\textbf{Concept Check 12.1}

\begin{quote}
\begin{enumerate}
\def\labelenumi{\arabic{enumi}.}
\tightlist
\item
  Why is the 5' cap structure a 5'-5' triphosphate linkage rather than
  the usual 3'-5'? What advantage does this provide?
\item
  A \hyperref[gl:mutation]{\textbf{mutation}} changes the branch point A
  to G in a critical intron. Predict the consequence for splicing.
\item
  Explain how nusinersen promotes exon 7 inclusion in SMN2 pre-mRNA.
\item
  Why does NMD require translation (and therefore nuclear export)?
\end{enumerate}
\end{quote}

\begin{center}\rule{0.5\linewidth}{0.5pt}\end{center}

\subsection{Translation}\label{translation}

\subsubsection{The Genetic Code}\label{the-genetic-code}

The genetic code maps 64 codons (\(4^3 = 64\)) to 20 amino acids plus
stop signals:

{\def\LTcaptype{none} % do not increment counter
\begin{longtable}[]{@{}
  >{\raggedright\arraybackslash}p{(\linewidth - 2\tabcolsep) * \real{0.4348}}
  >{\raggedright\arraybackslash}p{(\linewidth - 2\tabcolsep) * \real{0.5652}}@{}}
\toprule\noalign{}
\begin{minipage}[b]{\linewidth}\raggedright
Property
\end{minipage} & \begin{minipage}[b]{\linewidth}\raggedright
Description
\end{minipage} \\
\midrule\noalign{}
\endhead
\bottomrule\noalign{}
\endlastfoot
\textbf{Degenerate (redundant)} & Amino acids are encoded by 1-6
synonymous codons; Leu, Ser, Arg each have 6 codons; Met and Trp each
have 1 \\
\textbf{Unambiguous} & Each codon specifies exactly one amino acid (or
stop) \\
\textbf{Commaless and non-overlapping} & Read in consecutive,
non-overlapping triplets from a fixed start point \\
\textbf{Nearly comprehensive} & Exceptions in mitochondria
(\emph{Mycoplasma}: UGA = Trp; \emph{Tetrahymena}: UAA/UAG = Gln;
mammalian mitochondria: AGA/AGG = Stop) \\
\end{longtable}
}

\textbf{Start codon}: AUG (methionine) in eukaryotes; AUG
(formylmethionine) in prokaryotes. Rarely, GUG or UUG serve as
alternative start codons in prokaryotes.

\textbf{Stop codons}: UAA (ochre), UAG (amber), UGA (opal/umber).

\textbf{Wobble base pairing} \citep{crick1966}: The third codon position
(``wobble position'') pairs with the first
\hyperref[gl:anticodon]{\textbf{anticodon}} position with relaxed
stringency:

{\def\LTcaptype{none} % do not increment counter
\begin{longtable}[]{@{}ll@{}}
\toprule\noalign{}
Anticodon first position & Pairs with codon third position \\
\midrule\noalign{}
\endhead
\bottomrule\noalign{}
\endlastfoot
U & A, G \\
C & G primarily \\
A & U primarily \\
G & C, U \\
I (inosine) & A, C, U \\
\end{longtable}
}

Wobble explains how 45 tRNAs (in humans) can decode 61 sense codons.
\textbf{Suppressor tRNAs} are mutant tRNAs that can read stop codons as
amino acids, ``suppressing'' nonsense mutations.

\begin{figure}[htbp]
\centering
\includegraphics[width=0.82\linewidth,height=4.2in,keepaspectratio]{../figures/mermaid_inline/inline_mermaid_0074_1efbf9c5820e.png}
\caption{The ribosomal translation cycle in eukaryotes. Initiation involves cap-dependent scanning to the start codon. Elongation cycles through decoding, peptide bond formation, and translocation. Termination occurs when a stop codon is recognized by eRF1.}
\end{figure}

\subsubsection{Ribosome Structure}\label{ribosome-structure}

{\def\LTcaptype{none} % do not increment counter
\begin{longtable}[]{@{}
  >{\raggedright\arraybackslash}p{(\linewidth - 4\tabcolsep) * \real{0.1957}}
  >{\raggedright\arraybackslash}p{(\linewidth - 4\tabcolsep) * \real{0.4130}}
  >{\raggedright\arraybackslash}p{(\linewidth - 4\tabcolsep) * \real{0.3913}}@{}}
\toprule\noalign{}
\begin{minipage}[b]{\linewidth}\raggedright
Feature
\end{minipage} & \begin{minipage}[b]{\linewidth}\raggedright
Prokaryotic (70S)
\end{minipage} & \begin{minipage}[b]{\linewidth}\raggedright
Eukaryotic (80S)
\end{minipage} \\
\midrule\noalign{}
\endhead
\bottomrule\noalign{}
\endlastfoot
Small subunit & 30S (16S rRNA + about 21 proteins) & 40S (18S
rRNA + about 33 proteins) \\
Large subunit & 50S (23S + 5S rRNA + about 31 proteins) & 60S
(28S + 5.8S + 5S rRNA + about 49 proteins) \\
Sedimentation coefficient & 70S & 80S \\
Peptidyl transferase center & 23S rRNA & 28S rRNA \\
mRNA binding & Shine-Dalgarno sequence & 5' cap scanning \\
\end{longtable}
}

The ribosome has three tRNA-binding sites: - \textbf{A site
(aminoacyl)}: accepts incoming aminoacyl-tRNA - \textbf{P site
(peptidyl)}: holds the tRNA carrying the growing polypeptide chain -
\textbf{E site (exit)}: depleted (deacylated) tRNA exits here

\textbf{Critical insight}: The peptidyl transferase center (PTC) is
composed entirely of rRNA -- no protein comes within 18 A of the active
site. This demonstrates that the ribosome is a \textbf{ribozyme} (Nobel
Prize 2009: Ramakrishnan, Steitz, Yonath). This supports the \textbf{RNA
world hypothesis}: RNA catalyzed protein synthesis before proteins
existed.

\subsubsection{Translation Initiation, Elongation, and
Termination}\label{translation-initiation-elongation-and-termination}

\textbf{Prokaryotic initiation}: The \textbf{Shine-Dalgarno sequence}
(AGGAGG, about 5-10 nt upstream of AUG) base-pairs with the
anti-SD sequence at the 3' end of 16S rRNA, positioning the AUG in the P
site. IF1 blocks the A site; IF2-GTP escorts fMet-tRNA\(_f^{Met}\) to
the P site; IF3 prevents premature 50S joining. After start codon
recognition, IF2 hydrolyzes GTP, most IFs dissociate, and 50S joins.

\textbf{Eukaryotic initiation} (cap-dependent scanning model): 1.
\textbf{eIF4F complex} (eIF4E + eIF4G + eIF4A helicase) binds the
m\(^7\)G cap 2. eIF4G bridges to eIF3 on the 40S subunit and to PABP on
the poly(A) tail (circularization) 3. \textbf{43S pre-initiation
complex} (40S + eIF1 + eIF1A + eIF3 + eIF5 + ternary complex
{[}eIF2-GTP-Met-tRNA\(_i^{Met}\){]}) loads onto the mRNA near the cap 4.
The 43S complex \textbf{scans} 5' to 3' (powered by eIF4A helicase
activity) until it finds the first AUG in a favorable \textbf{Kozak
context} (RCCAUGG, R = purine; -3 and +4 positions most critical) 5.
Start codon recognition: eIF2 hydrolyzes GTP; eIF5B-GTP promotes 60S
joining; most eIFs released 6. \textbf{80S initiation complex} formed
with Met-tRNA\(_i^{Met}\) in P site

\textbf{Elongation cycle} (about 5-10 amino acids/s in
eukaryotes; about 15-20 aa/s in prokaryotes): 1.
\textbf{Decoding}: aminoacyl-tRNA delivered to A site as part of a
ternary complex with eEF1A (EF-Tu in prokaryotes) and GTP. Correct
codon-anticodon match triggers conformational change; GTP hydrolysis
releases eEF1A 2. \textbf{Peptide bond formation}: Peptidyl transferase
(28S rRNA catalytic activity) transfers the growing peptide from the
P-site tRNA to the A-site amino acid. The reaction is spontaneous
(\(\Delta G < 0\)) -- the aminoacyl ester bond is high-energy 3.
\textbf{Translocation}: eEF2-GTP (EF-G in prokaryotes) promotes
translocation; ribosome moves 3 nt (one codon) along mRNA; tRNAs shift:
A to P, P to E; E-site tRNA dissociates

\textbf{Termination}: When a stop codon enters the A site: -
\textbf{eRF1} (class I release factor) recognizes most three stop codons
(in eukaryotes) -- its shape mimics tRNA - \textbf{eRF3-GTP} (class II
release factor) stimulates eRF1 activity - Water molecule attacks the
peptidyl-tRNA ester bond, releasing the completed polypeptide -
\textbf{ABCE1} (ribosome recycling factor) splits the 80S ribosome into
subunits for re-use

\begin{quote}
\textbf{Clinical Connection: Antibiotics Targeting the Ribosome} The
structural differences between prokaryotic 70S and eukaryotic 80S
ribosomes make the bacterial ribosome an excellent drug target:

{\def\LTcaptype{none} % do not increment counter
\begin{longtable}[]{@{}lll@{}}
\toprule\noalign{}
Antibiotic & Target & Mechanism \\
\midrule\noalign{}
\endhead
\bottomrule\noalign{}
\endlastfoot
Tetracycline & 30S A site & Blocks aminoacyl-tRNA binding \\
Streptomycin & 30S decoding center & Causes misreading of mRNA \\
Chloramphenicol & 50S PTC & Inhibits peptide bond formation \\
Erythromycin & 50S exit tunnel & Blocks peptide exit \\
Linezolid & 50S A site & Prevents fMet-tRNA binding \\
\end{longtable}
}

These antibiotics are selectively toxic because they do not bind
eukaryotic ribosomes efficiently. However, mitochondrial ribosomes
resemble prokaryotic ribosomes, explaining side effects like ototoxicity
(aminoglycosides) in patients with mitochondrial rRNA mutations.
\end{quote}

\subsubsection{IRES Elements --- Cap-Independent Translation
Initiation}\label{ires-elements-cap-independent-translation-initiation}

While most eukaryotic translation initiates at the 5\texorpdfstring{'}{'} cap
(eIF4F-mediated scanning), some mRNAs can be translated even when
cap-dependent translation is shut down. They use \textbf{internal
ribosome entry sites (IRES)} --- RNA structural elements that directly
recruit the 40S ribosomal subunit without scanning from the cap.

\textbf{IRES classes (4 main types in eukaryotic and viral mRNAs):}

{\def\LTcaptype{none} % do not increment counter
\begin{longtable}[]{@{}
  >{\raggedright\arraybackslash}p{(\linewidth - 6\tabcolsep) * \real{0.1042}}
  >{\raggedright\arraybackslash}p{(\linewidth - 6\tabcolsep) * \real{0.1458}}
  >{\raggedright\arraybackslash}p{(\linewidth - 6\tabcolsep) * \real{0.5625}}
  >{\raggedright\arraybackslash}p{(\linewidth - 6\tabcolsep) * \real{0.1875}}@{}}
\toprule\noalign{}
\begin{minipage}[b]{\linewidth}\raggedright
Class
\end{minipage} & \begin{minipage}[b]{\linewidth}\raggedright
Example
\end{minipage} & \begin{minipage}[b]{\linewidth}\raggedright
Initiation factors required
\end{minipage} & \begin{minipage}[b]{\linewidth}\raggedright
Mechanism
\end{minipage} \\
\midrule\noalign{}
\endhead
\bottomrule\noalign{}
\endlastfoot
Type I & Picornaviruses (poliovirus, coxsackievirus) & Most canonical
eIFs except eIF4E & Internal AUG; PCBP2/3 binding \\
Type II & Picornaviruses (encephalomyocarditis virus, EMCV) & eIF1A,
eIF2, eIF3, eIF4A, eIF4G; \textbf{no eIF4E} & Direct 40S placement at
AUG; cleaved eIF4G is sufficient \\
Type III & Hepatitis C virus (HCV) & Primarily eIF2, eIF3 (and 40S) ---
minimal & Direct binding of 40S to a multi-domain RNA structure (domains
II, III, IV); IIIc binds 40S like a TBP-equivalent for ribosomes \\
Type IV & Cricket paralysis virus (CrPV); intergenic IRES & None --- no
initiator Met-tRNA needed & Pseudoknot mimics tRNA in P site;
translation starts at non-AUG codon \\
\end{longtable}
}

\textbf{Hepatitis C IRES --- structural mechanism in detail.} HCV IRES
is one of the most structurally characterised. The IRES adopts a defined
three-domain RNA fold: - \textbf{Domain II} binds the 40S subunit at the
head, positioning the mRNA channel - \textbf{Domain III} (with
sub-domains IIIa, IIIb, IIIc, IIId, IIIe, IIIf) directly contacts
ribosomal proteins; IIIabc binds eIF3 - \textbf{Domain IV} contains the
AUG start codon, positioned directly in the P site - This bypass of cap
recognition and scanning means HCV translation continues even when
cellular cap-dependent translation is suppressed.

\textbf{Cellular IRES elements} drive translation of stress-response
mRNAs when cap-dependent translation is inhibited (e.g., during
apoptosis when eIF4G is cleaved by caspase-3, or during heat shock when
eIF2\texorpdfstring{\ensuremath{\alpha}}{alpha} is phosphorylated). Examples: \textbf{c-Myc, p53, VEGF, BiP/HSPA5,
XIAP, DAP5, BCL2.} Stress-induced translation of pro-apoptotic
\emph{p53} via its IRES is one of the host's defences against viral
infection.

\textbf{Therapeutic relevance:} - \textbf{eIF4A inhibitors} (silvestrol,
zotatifin) preferentially affect cap-dependent translation of mRNAs with
structured 5\texorpdfstring{'}{'} UTRs --- bypassed by some IRESes. In trial for AML and
CTCL. - \textbf{Eltrombopag} (FDA-approved for thrombocytopenia)
destabilises HCV-IRES--40S contacts as an off-target mechanism.

\subsubsection{Ribosome Profiling --- A Window Into Translation
Dynamics}\label{ribosome-profiling-a-window-into-translation-dynamics}

\textbf{Ribosome profiling (Ribo-seq)} is a technique that captures and
sequences the about 28-nt mRNA fragments protected by
translating ribosomes. Each ``ribosome footprint'' identifies (a) which
mRNA is being translated, (b) which codon is in the A-site at the moment
of capture, and (c) the relative density of ribosomes along the message
--- quantifying \textbf{translation efficiency} at codon resolution.

\textbf{Key insights from ribosome profiling:}

\begin{enumerate}
\def\labelenumi{\arabic{enumi}.}
\tightlist
\item
  \textbf{Translation efficiency varies across mRNAs by \textgreater{}
  100-fold.} Some mRNAs have many more ribosomes loaded per unit of mRNA
  than others. This is regulated by 5\texorpdfstring{'}{'} UTR features (length, structure,
  upstream ORFs), codon usage in the ORF, and 3\texorpdfstring{'}{'} UTR sequences.
\item
  \textbf{Ribosome stalling on rare codons.} Codons recognised by
  low-abundance tRNAs (e.g., CGA-Arg, CGG-Arg in mammals) cause ribosome
  pausing detectable as elevated read density. Stalling can trigger:

  \begin{itemize}
  \tightlist
  \item
    \textbf{Co-translational protein folding signals} --- programmed
    pauses allow N-terminal protein domains to fold before C-terminal
    segments emerge.
  \item
    \textbf{Programmed -1 frameshifts} (e.g., HIV gag-pol, SARS-CoV-2
    ORF1ab/RDR-RP).
  \item
    \textbf{Ribosome quality control (RQC):} stalled ribosomes are split
    by Pelota--Hbs1L; the nascent peptide is ubiquitinated by NEMF--LTN1
    and degraded.
  \end{itemize}
\item
  \textbf{Upstream open reading frames (uORFs) regulate
  about 50 \% of mammalian mRNAs.} A short ORF in the 5\texorpdfstring{'}{'} UTR is
  translated first; the scanning ribosome must reinitiate downstream.
  uORF translation typically suppresses main-ORF translation.
  Stress-induced phosphorylation of eIF2\texorpdfstring{\ensuremath{\alpha}}{alpha} (the integrated stress
  response) increases scanning past inhibitory uORFs, paradoxically
  enhancing translation of stress-response mRNAs (e.g., \emph{ATF4,
  CHOP, GADD34}) that have inhibitory uORFs in their 5\texorpdfstring{'}{'} UTRs.
\item
  \textbf{Codon-level translation kinetics.} Each codon's average dwell
  time can be estimated from ribosome density. Codon-optimised mRNAs
  (e.g., synthetic mRNA vaccines, insulin) have shorter average dwell
  times and higher translation efficiency.
\item
  \textbf{Out-of-frame translation.} Ribo-seq has uncovered hidden ORFs
  (smORFs, 10--100 codons) in regions previously annotated as non-coding
  --- many of which encode functional micropeptides (e.g., MOTS-c from
  mitochondrial 12S rRNA; HOXB-AS3 micropeptide).
\end{enumerate}

\subsubsection{Nonsense-Mediated Decay (NMD) --- Mechanism in
Mechanistic
Detail}\label{nonsense-mediated-decay-nmd-mechanism-in-mechanistic-detail}

NMD destroys mRNAs containing premature termination codons (PTCs),
preventing the production of truncated proteins that would often act as
dominant-negative or gain-of-function pathogens. The key recognition
principle is the \textbf{EJC-distal PTC} rule: a stop codon is
``premature'' if the ribosome encounters it more than about 50
nt upstream of the last exon-exon junction.

\textbf{Key molecular components:}

{\def\LTcaptype{none} % do not increment counter
\begin{longtable}[]{@{}
  >{\raggedright\arraybackslash}p{(\linewidth - 2\tabcolsep) * \real{0.6923}}
  >{\raggedright\arraybackslash}p{(\linewidth - 2\tabcolsep) * \real{0.3077}}@{}}
\toprule\noalign{}
\begin{minipage}[b]{\linewidth}\raggedright
Component
\end{minipage} & \begin{minipage}[b]{\linewidth}\raggedright
Role
\end{minipage} \\
\midrule\noalign{}
\endhead
\bottomrule\noalign{}
\endlastfoot
\textbf{EJC (exon junction complex)} & Deposited 20--24 nt upstream of
every exon-exon junction during splicing; contains EIF4A3, MAGOH,
Y14/RBM8A, MLN51 (CASC3) \\
\textbf{UPF1 (RENT1)} & RNA helicase + ATPase; central NMD effector;
binds the 3\texorpdfstring{'}{'} UTR after stop codon \\
\textbf{UPF2 (RENT2)} & Bridges UPF1 to UPF3 \\
\textbf{UPF3A/UPF3B} & EJC component; bridges EJC to UPF1/2 \\
\textbf{SMG1} & PIKK-family kinase; phosphorylates UPF1 \\
\textbf{SMG5/6/7} & Recruit decapping (SMG7) and endonuclease cleavage
(SMG6) machinery \\
\textbf{eRF1, eRF3} & Translation termination factors; eRF3 interacts
with UPF1 at PTC \\
\end{longtable}
}

\textbf{Mechanism, step-by-step:}

\begin{enumerate}
\def\labelenumi{\arabic{enumi}.}
\tightlist
\item
  \textbf{mRNA enters first-round translation} (the ``pioneer round''),
  still bound by CBP80/20 cap-binding complex (not yet the cytoplasmic
  eIF4F).
\item
  \textbf{Ribosome translates from cap to first stop codon.} Each EJC
  encountered by the elongating ribosome is dislodged. EJCs deposited
  downstream of the actual stop codon (in the 3\texorpdfstring{'}{'} UTR) are NOT dislodged.
\item
  \textbf{Stop codon recognition.} eRF1 + eRF3-GTP enter the A site;
  eRF3 hydrolyses GTP; nascent peptide is released by hydrolysis of the
  peptidyl-tRNA ester bond.
\item
  \textbf{UPF1 recruitment.} eRF3 (after release-factor function)
  recruits UPF1 to the terminating ribosome.
\item
  \textbf{EJC-distal proximity check.} UPF1 scans the 3\texorpdfstring{'}{'} UTR; if it
  encounters an EJC within about 50 nt downstream, EJC-bound
  UPF2/UPF3B contact UPF1, activating NMD.
\item
  \textbf{SMG1 phosphorylation of UPF1} at numerous Ser/Thr-Q sites ---
  generates the ``phospho-UPF1'' signal.
\item
  \textbf{Decay execution:}

  \begin{itemize}
  \tightlist
  \item
    \textbf{SMG6 endonuclease} cleaves the mRNA near the PTC, exposing
    both fragments to exonucleases (XRN1 from 5\texorpdfstring{'}{'}; exosome from 3\texorpdfstring{'}{'}).
  \item
    \textbf{SMG5--SMG7} recruits CCR4--NOT (deadenylation) and DCP1/2
    (decapping) for canonical mRNA decay.
  \end{itemize}
\item
  \textbf{Ribosome recycling and UPF1 dephosphorylation} by PP2A
  complete the cycle.
\end{enumerate}

\textbf{Clinical NMD exploitation --- three therapeutic strategies:}

\begin{enumerate}
\def\labelenumi{\arabic{enumi}.}
\tightlist
\item
  \textbf{Read-through compounds} (e.g., \textbf{ataluren/Translarna})
  promote ribosomal mis-decoding at PTCs (especially UGA), allowing
  partial production of full-length protein. Approved in EU for
  nonsense-mutation Duchenne muscular dystrophy (cmDMD).
\item
  \textbf{NMD inhibition} (e.g., NMDI-1, SMG1 inhibitors in development)
  prolongs the half-life of PTC-containing mRNAs, increasing the
  substrate pool for read-through compounds.
\item
  \textbf{Antisense oligonucleotides for exon skipping} (e.g.,
  eteplirsen for DMD, nusinersen for SMA): exclude the PTC-containing
  exon, restoring the reading frame.
\end{enumerate}

NMD also has homeostatic functions beyond pathogen detection: it
degrades a large fraction of physiological alternatively-spliced
isoforms that contain PTCs (regulated unproductive splicing and
translation, RUST), particularly in neurons and during development. NMD
perturbation in flies and mice produces brain phenotypes consistent with
this regulatory role.

\subsubsection{Translation Fidelity --- Mechanism and Error
Rate}\label{translation-fidelity-mechanism-and-error-rate}

Translation fidelity is achieved through multiple kinetic checkpoints,
each contributing to the overall error rate of approximately \textbf{1
in 10\texorpdfstring{\textsuperscript{4}}{4} to 10\texorpdfstring{\textsuperscript{5}}{5} amino acids} misincorporated:

\begin{enumerate}
\def\labelenumi{\arabic{enumi}.}
\tightlist
\item
  \textbf{Aminoacyl-tRNA charging fidelity} by aminoacyl-tRNA
  synthetases (aaRS): each aaRS has high specificity for its cognate
  amino acid (10\texorpdfstring{\textsuperscript{-}}{-}\texorpdfstring{\textsuperscript{3}}{3}--10\texorpdfstring{\textsuperscript{-}}{-}\texorpdfstring{\textsuperscript{4}}{4} error in pre-transfer step). Many aaRSs (e.g.,
  IleRS, LeuRS, ValRS, AlaRS, PheRS, ProRS, ThrRS, MetRS) have
  \textbf{post-transfer editing domains} that hydrolyse mischarged tRNAs
  (e.g., Val-tRNA\^{}Ile is edited by IleRS's editing domain, achieving
  10\texorpdfstring{\textsuperscript{-}}{-}\texorpdfstring{\textsuperscript{5}}{5}--10\texorpdfstring{\textsuperscript{-}}{-}\texorpdfstring{\textsuperscript{6}}{6} overall error). The ``double-sieve'' mechanism: synthetic
  site selects amino acids that fit \emph{or are smaller than} the
  cognate; editing site rejects amino acids smaller than the cognate.
\item
  \textbf{Initial codon-anticodon recognition} (in the 30S/40S A site):
  Watson-Crick pairing at codon positions 1 and 2 (the wobble at
  position 3 is intentionally relaxed); A-form helix geometry sensing by
  16S rRNA (in bacteria: A1492, A1493, G530 --- the ``monitor bases'').
  Initial selection error rate about 10\texorpdfstring{\textsuperscript{-}}{-}\texorpdfstring{\textsuperscript{3}}{3}.
\item
  \textbf{EF-Tu / eEF1A conformational proofreading} (kinetic
  proofreading via GTP hydrolysis): EF-Tu (in bacteria) or eEF1A (in
  eukaryotes) brings aa-tRNA to the A site; correct codon-anticodon
  pairing accelerates GTP hydrolysis. Incorrect pairings have time to
  dissociate before GTP is hydrolysed (an irreversible step). After GTP
  hydrolysis, a second selection step (``accommodation'') allows further
  discrimination. This adds about 10\texorpdfstring{\textsuperscript{-}}{-}\texorpdfstring{\textsuperscript{3}}{3} to about 10\texorpdfstring{\textsuperscript{-}}{-}\texorpdfstring{\textsuperscript{4}}{4}
  improvement.
\item
  \textbf{Peptidyl transferase activity-coupled fidelity}: the rate of
  peptide bond formation depends on whether the tRNA is correctly
  accommodated; incorrect tRNAs have slower peptide-bond chemistry,
  providing another opportunity for dissociation.
\item
  \textbf{Quality control after translation}: misfolded proteins are
  degraded by the proteasome; ribosomes stalling on damaged mRNAs
  trigger ribosome-associated quality control (RQC) via Pelota-Hbs1L
  splitting and NEMF-LTN1-mediated nascent-peptide ubiquitination.
\end{enumerate}

\begin{equation}\text{Overall translation error} \approx 10^{-3} \times 10^{-3} \times \text{(post-transfer editing)} \approx 10^{-4}\text{ to }10^{-5} \tag{13.4}\label{eq:translation_fidelity}\end{equation}

For a typical about 500-amino acid protein, this error rate
predicts about 5 \% of nascent polypeptides contain at least
one misincorporation. Most are tolerated; a few drive cellular
dysfunction in aging or disease (statin-induced muscle damage from
low-fidelity translation; age-related aggregation diseases). Sub-cohort
of mistranslated proteins are degraded by the proteasome before causing
problems.

\textbf{Clinical translation:} Aminoglycosides (gentamicin, paromomycin)
bind 16S rRNA monitor bases, distorting the codon-recognition geometry
and increasing translation error rate (about 100-fold). This
explains both their bactericidal mechanism (catastrophic protein
synthesis errors) and their toxicities (cochlear hair-cell apoptosis
from mitochondrial mistranslation, especially in patients with the
A1555G 12S rRNA variant).

\textbf{Concept Check 12.2}

\begin{quote}
\begin{enumerate}
\def\labelenumi{\arabic{enumi}.}
\tightlist
\item
  Compare the roles of the Shine-Dalgarno sequence (prokaryotes) and
  Kozak sequence (eukaryotes) in translation initiation.
\item
  Why is the ribosome considered a ribozyme? What evidence supports
  this?
\item
  A patient has a mitochondrial 12S rRNA mutation (A1555G). Why might
  aminoglycoside antibiotics cause deafness in this patient?
\item
  Explain the wobble hypothesis and why fewer tRNAs than sense codons
  are needed.
\end{enumerate}
\end{quote}

\begin{center}\rule{0.5\linewidth}{0.5pt}\end{center}

\subsection{Post-Translational Modifications and Protein
Degradation}\label{post-translational-modifications-and-protein-degradation}

\subsubsection{Post-Translational Modifications
(PTMs)}\label{post-translational-modifications-ptms-1}

{\def\LTcaptype{none} % do not increment counter
\begin{longtable}[]{@{}
  >{\raggedright\arraybackslash}p{(\linewidth - 6\tabcolsep) * \real{0.2653}}
  >{\raggedright\arraybackslash}p{(\linewidth - 6\tabcolsep) * \real{0.1837}}
  >{\raggedright\arraybackslash}p{(\linewidth - 6\tabcolsep) * \real{0.3469}}
  >{\raggedright\arraybackslash}p{(\linewidth - 6\tabcolsep) * \real{0.2041}}@{}}
\toprule\noalign{}
\begin{minipage}[b]{\linewidth}\raggedright
Modification
\end{minipage} & \begin{minipage}[b]{\linewidth}\raggedright
Enzymes
\end{minipage} & \begin{minipage}[b]{\linewidth}\raggedright
Target Residues
\end{minipage} & \begin{minipage}[b]{\linewidth}\raggedright
Function
\end{minipage} \\
\midrule\noalign{}
\endhead
\bottomrule\noalign{}
\endlastfoot
\textbf{Phosphorylation} & Kinases / Phosphatases & Ser, Thr, Tyr &
\hyperref[gl:signal-transduction]{\textbf{Signal transduction}};
activation/inactivation; about 30\% of proteins are
phosphorylated \\
\textbf{Glycosylation} & Glycosyltransferases & Asn (N-linked), Ser/Thr
(O-linked) & Protein folding, stability, cell-cell recognition; starts
in ER, completed in Golgi \\
\textbf{Ubiquitination} & E1/E2/E3 ligase cascade & Lys (isopeptide
bond) & Mono-Ub: \hyperref[gl:endocytosis]{\textbf{endocytosis}},
histone regulation; Poly-Ub (K48-linked): proteasomal degradation;
K63-linked: signaling \\
\textbf{Acetylation} & HATs / HDACs & Lys (histones and non-histone
proteins) & \hyperref[gl:chromatin]{\textbf{Chromatin}} opening
(histones); p53 activation; metabolic regulation \\
\textbf{Methylation} & Methyltransferases / Demethylases & Lys, Arg
(histones) & Histone code: H3K4me3 = active; H3K27me3 = repressed;
H3K9me3 = \hyperref[gl:heterochromatin]{\textbf{heterochromatin}} \\
\textbf{SUMOylation} & E1/E2/E3 SUMO ligases & Lys & Nuclear transport,
transcription regulation, DNA repair; antagonizes ubiquitination at same
Lys \\
\textbf{Proteolytic cleavage} & Proteases & Specific peptide bonds &
Signal peptide removal; proinsulin to insulin;
\hyperref[gl:caspase]{\textbf{caspase}}-mediated
\hyperref[gl:apoptosis]{\textbf{apoptosis}} \\
\end{longtable}
}

\subsubsection{The Ubiquitin-Proteasome
System}\label{the-ubiquitin-proteasome-system}

The ubiquitin-proteasome system (UPS) is the primary mechanism for
targeted protein degradation:

\begin{enumerate}
\def\labelenumi{\arabic{enumi}.}
\tightlist
\item
  \textbf{E1 (ubiquitin-activating enzyme)}: Activates ubiquitin (76-aa
  protein) in an ATP-dependent reaction; forms E1about  Ub
  thioester bond; humans use 2 main E1 enzymes
\item
  \textbf{E2 (ubiquitin-conjugating enzyme)}: Accepts Ub from E1;
  about 40 E2 enzymes provide some specificity
\item
  \textbf{E3 (ubiquitin ligase)}: Provides substrate specificity;
  about 600 E3 ligases in humans. Two major families:

  \begin{itemize}
  \tightlist
  \item
    \textbf{RING E3s}: Transfer Ub directly from E2 to substrate (e.g.,
    APC/C, SCF complex, MDM2)
  \item
    \textbf{HECT E3s}: Accept Ub from E2, then transfer to substrate
  \end{itemize}
\item
  Poly-ubiquitination (K48-linked chains of at least 4 Ub) marks the
  substrate for degradation by the \textbf{26S proteasome}
\item
  \textbf{26S proteasome} (2.5 MDa): 20S catalytic core (barrel-shaped;
  contains chymotrypsin-like, trypsin-like, and caspase-like proteases)
  capped by 19S regulatory particles (recognize poly-Ub, unfold
  substrate, feed into 20S core)
\item
  Deubiquitinating enzymes (DUBs, about 100 in humans) recycle
  ubiquitin before substrate enters the 20S core
\end{enumerate}

\begin{quote}
\textbf{Clinical Connection: Proteasome Inhibitors in Cancer}
\textbf{Bortezomib (Velcade)} is a proteasome inhibitor approved for
multiple myeloma and mantle cell lymphoma. By blocking proteasomal
degradation, it causes accumulation of pro-apoptotic proteins and
misfolded proteins (ER stress), triggering cell death preferentially in
cancer cells with high protein synthesis rates.
\end{quote}

\begin{center}\rule{0.5\linewidth}{0.5pt}\end{center}

\subsection{Gene Regulation}\label{gene-regulation}

\subsubsection{Prokaryotic Gene Regulation: The Lac
Operon}\label{prokaryotic-gene-regulation-the-lac-operon}

The \emph{E. coli} lac operon \citep{jacob1961} (Jacob and Monod, 1961;
Nobel Prize 1965) was the founding model of gene regulation:

\textbf{Components}: \emph{lacZ} (beta-galactosidase), \emph{lacY}
(permease), \emph{lacA} (transacetylase), preceded by the promoter (P),
operator (O), and regulated by \emph{lacI} (repressor gene,
constitutively expressed from its own promoter).

\textbf{Negative control}: Lac repressor protein (tetramer, encoded by
\emph{lacI}) binds the operator sequences (O1 primary, O2 and O3
auxiliary -- DNA looping between O1 and O3 increases repression
50-fold). When lactose is present, its isomer \textbf{allolactose} (the
true inducer) binds the repressor, causing a conformational change that
reduces DNA-binding affinity by about 1000-fold. Repressor
released; RNAP can transcribe.

\textbf{Positive control (catabolite repression)}: CAP (catabolite
activator protein, also called CRP) + cAMP binds the CAP site upstream
of the lac promoter, bending DNA about 90 degrees and directly
contacting the alpha-CTD of RNAP to activate transcription
about 50-fold. When glucose is present, the phosphotransferase
system (PTS) keeps adenylyl cyclase inactive; {[}cAMP{]} is low; CAP
cannot bind; lac operon expression is reduced even if lactose is
present.

\textbf{Operon logic}: The lac operon is fully ON primarily when lactose
is present AND glucose is absent (high cAMP, CAP active, repressor
released). This makes biological sense: the cell uses glucose (more
efficient) before lactose.

\subsubsection{The Trp Operon (Repressible
System)}\label{the-trp-operon-repressible-system}

The \emph{E. coli} trp operon encodes five enzymes for tryptophan
biosynthesis:

\begin{itemize}
\tightlist
\item
  \textbf{Repression}: When tryptophan is abundant, Trp acts as a
  \textbf{co-repressor} -- it binds the TrpR aporepressor, activating it
  to bind the operator and block transcription
\item
  \textbf{Attenuation}: A leader peptide (trpL) containing two
  consecutive Trp codons (UGG UGG) is located upstream of the structural
  genes. When Trp-tRNA is abundant, the ribosome translates the leader
  rapidly, allowing formation of a terminator hairpin (3-4 stem loop) in
  the mRNA, causing premature termination. When Trp-tRNA is scarce, the
  ribosome stalls at the UGG codons, allowing an anti-terminator hairpin
  (2-3 stem loop) to form instead, permitting read-through transcription
\end{itemize}

This provides fine-tuning: repression gives about 70-fold
regulation; attenuation adds another about 8-fold; combined
about 600-fold.

\begin{quote}
\textbf{Concept Check:} Glucose is absent and lactose is present while a
\emph{lacI}\texorpdfstring{\textsuperscript{-}}{-}\texorpdfstring{\textsuperscript{-}}{-} strain also carries a non-functional \emph{crp} gene ---
predict whether the lac operon is transcribed and justify using both the
repressor and CAP--cAMP states.
\end{quote}

\subsubsection{Eukaryotic Gene Regulation: Multiple
Levels}\label{eukaryotic-gene-regulation-multiple-levels}

\textbf{Chromatin level}: - \textbf{SWI/SNF complexes} (BAF/PBAF in
mammals): ATP-dependent chromatin remodelers that slide, eject, or
restructure \hyperref[gl:nucleosome]{\textbf{nucleosome}}s to expose or
occlude regulatory sequences - \textbf{HATs} (histone
acetyltransferases, e.g., p300/CBP, PCAF): acetylate histone tails
(primarily H3K9, H3K14, H3K27, H4K16), neutralizing positive charge,
weakening histone-DNA interaction, and recruiting bromodomain-containing
factors -- associated with active transcription - \textbf{HDACs}
(histone deacetylases): remove acetyl groups -- associated with gene
silencing; HDAC inhibitors (vorinostat, romidepsin) are approved for
T-cell lymphoma

\textbf{Transcription factor combinatorics}: The human genome encodes
about 1,500-2,000 transcription factors. Gene expression is
controlled by the combinatorial binding of multiple TFs at enhancers,
which together produce richer logic than any single TF in isolation.
Three canonical principles:

\begin{enumerate}
\def\labelenumi{\arabic{enumi}.}
\tightlist
\item
  \textbf{Synergy:} Two TFs bound at adjacent sites produce more than
  the sum of their individual effects. Cooperative DNA binding (when one
  TF stabilises the binding of the next via direct protein--protein
  contact or by inducing DNA bending) and synergistic activation (when
  each TF independently recruits a different co-activator) both
  contribute. The interferon-\texorpdfstring{\ensuremath{\beta}}{beta} enhanceosome shows about 10\texorpdfstring{\ensuremath{\times}}{x}
  synergy: each individual TF drives about 5-fold transcription, but
  together they drive \textgreater{} 1000-fold induction.
\item
  \textbf{Quenching:} A repressor TF bound near an activator can inhibit
  the activator without displacing it (e.g., via histone deacetylation
  recruitment, or by sequestering coactivators in a non-productive
  complex). Examples: GAL80 quenches GAL4 in yeast; BCL6 quenches NF-\texorpdfstring{\ensuremath{\kappa}}{kappa}B
  in germinal-centre B cells.
\item
  \textbf{Enhanceosome:} A specific stereochemical assembly where the TF
  binding sites are positioned such that most eight (or more) factors
  must bind simultaneously and cooperatively to activate transcription.
  Spacing between binding sites is critical (loss of even one base-pair
  spacing destroys cooperative assembly). The IFN-\texorpdfstring{\ensuremath{\beta}}{beta} enhanceosome (3.7 kb
  upstream): NF-\texorpdfstring{\ensuremath{\kappa}}{kappa}B (p50/p65), IRF3, IRF7, ATF2/c-Jun, plus the
  architectural protein HMGA1 wraps the entire 55-bp regulatory region
  into a defined 3D conformation that primarily forms during viral
  infection.
\end{enumerate}

\textbf{Super-enhancers --- the high-output regulatory class.} First
operationally defined by Whyte and Young (2013), super-enhancers are
unusually large regulatory regions (median about 10 kb, often
\textgreater50 kb) characterised by: - \textbf{H3K27ac density:} at
least 10-fold higher than typical enhancers, often spanning multiple
H3K27ac peaks - \textbf{BRD4 occupancy:} the BET-family reader
concentrates here, as do Mediator and lineage-defining TFs -
\textbf{MED1/MED12 enrichment:} super-enhancers are constructed largely
from Mediator-bound elements - \textbf{Transcriptional output:} drives
high-output expression of cell-identity genes (e.g., \emph{MYC} in
cancer cell lines, lineage-defining TFs like \emph{SOX2} in ES cells,
\emph{PU.1} in haematopoietic cells) - \textbf{Phase separation:}
super-enhancer condensates concentrate Pol II, BRD4, Mediator, CDK7/9
--- see \cref{sec:unit_IV_epigenetics_and_gene_regulation} for detailed
mechanism - \textbf{Identification:} ROSE (Rank Ordering of
Super-Enhancers) algorithm --- rank H3K27ac signal across most
enhancers; the inflection point distinguishes super-enhancers (top
about 3 \% of enhancers, accounting for about 30 \% of
H3K27ac signal) from typical enhancers - \textbf{Therapeutic
implication:} BET inhibitors (JQ1, OTX015, molibresib, mivebresib)
preferentially collapse super-enhancer-driven transcription, providing
tumour selectivity. CDK7 (THZ1, SY-5609) and CDK9 (AZD4573) inhibitors
disrupt the condensate phosphorylation machinery.

The combination of synergy + enhanceosome organisation + super-enhancer
condensates explains why a few hundred lineage-defining genes can be
expressed at very high levels in each cell type, while the rest of the
genome is largely inactive.

\subsubsection{RNA-Based Regulation}\label{rna-based-regulation}

\textbf{\hyperref[gl:microrna]{microRNA (miRNA)}}: - about 22
nt single-stranded RNAs; about 2,000 in the human genome;
regulate about 60\% of protein-coding genes - Biogenesis:
pri-miRNA (transcribed by Pol II) \texorpdfstring{\ensuremath{\to}}{->} Drosha/DGCR8 cleave in nucleus to
pre-miRNA (hairpin) \texorpdfstring{\ensuremath{\to}}{->} Exportin-5 exports \texorpdfstring{\ensuremath{\to}}{->} \textbf{Dicer} cleaves to
about 22 nt duplex \texorpdfstring{\ensuremath{\to}}{->} one strand loaded into \textbf{RISC}
(RNA-induced silencing complex, containing Argonaute protein) -
Mechanism: miRNA guides RISC to complementary sites (typically in the 3'
UTR; seed sequence = positions 2-8); partial complementarity \texorpdfstring{\ensuremath{\to}}{->}
translational repression + mRNA deadenylation and decay; perfect
complementarity \texorpdfstring{\ensuremath{\to}}{->} mRNA cleavage (rare in animals)

\textbf{The miRNA pathway in mechanistic detail (mammalian):}

\begin{enumerate}
\def\labelenumi{\arabic{enumi}.}
\tightlist
\item
  \textbf{DICER processing:} DICER (PAZ + DUF + RNase IIIa + RNase IIIb
  domains) measures about 22 nt from the pre-miRNA terminus and
  cleaves to produce a about 22-nt duplex with 2-nt 3\texorpdfstring{'}{'}
  overhangs. DICER's PAZ domain binds the 3\texorpdfstring{'}{'} overhang of the pre-miRNA
  hairpin; the dsRBD anchors the substrate; the two RNase III domains
  cleave both strands in a single catalytic event.
\item
  \textbf{RISC loading complex (RLC):} TRBP + DICER + AGO2 form the RLC.
  The thermodynamic asymmetry rule determines which strand becomes the
  \textbf{guide} vs.~the \textbf{passenger}: the strand whose 5\texorpdfstring{'}{'} end is
  in a less-stable base-paired region is preferentially loaded as the
  guide. The passenger strand is degraded.
\item
  \textbf{AGO2 conformational dynamics:} AGO2 has four domains (N, PAZ,
  MID, PIWI). The MID domain anchors the guide's 5\texorpdfstring{'}{'} phosphate via a
  5\texorpdfstring{'}{'}-anchored binding pocket. The PAZ domain holds the 3\texorpdfstring{'}{'} end (when not
  engaged with target). The PIWI domain has the slicer (RNase H-like)
  active site that cleaves perfectly complementary targets.
\item
  \textbf{Target search and seed-pairing:} The guide miRNA's positions
  2--8 (the \textbf{seed sequence}) must base-pair with the mRNA target
  --- this is the primary recognition determinant. AGO2 conformationally
  remodels to ``extrude'' positions 2--8 in an A-form helix-ready
  geometry for rapid scanning.
\item
  \textbf{mRNA repression by deadenylation--decapping--decay:}

  \begin{itemize}
  \tightlist
  \item
    \textbf{CCR4--NOT complex} is recruited to the targeted mRNA via
    TNRC6/GW182 family proteins bound to AGO2.
  \item
    \textbf{CCR4 (CNOT6/CNOT6L) and POP2 (CNOT7/CNOT8)} deadenylases
    progressively remove the poly-A tail (deadenylation is the
    rate-limiting step).
  \item
    \textbf{DCP2 (decapping)} removes the 5\texorpdfstring{'}{'} cap, exposing the mRNA to
    \textbf{XRN1} 5\texorpdfstring{'}{'}-to-3\texorpdfstring{'}{'} exonuclease degradation.
  \item
    In parallel, the cytoplasmic \textbf{EXOSOME} complex degrades from
    the 3\texorpdfstring{'}{'} end.
  \end{itemize}
\item
  \textbf{Translational repression in parallel:} GW182 also disrupts
  eIF4F--PABP interaction, preventing closed-loop translation
  initiation. eIF4A helicase activity is reduced. Repression and decay
  are coupled but partially separable.
\end{enumerate}

\textbf{Seed sequence rules --- quantitative target prediction:}

{\def\LTcaptype{none} % do not increment counter
\begin{longtable}[]{@{}
  >{\raggedright\arraybackslash}p{(\linewidth - 4\tabcolsep) * \real{0.1957}}
  >{\raggedright\arraybackslash}p{(\linewidth - 4\tabcolsep) * \real{0.2174}}
  >{\raggedright\arraybackslash}p{(\linewidth - 4\tabcolsep) * \real{0.5870}}@{}}
\toprule\noalign{}
\begin{minipage}[b]{\linewidth}\raggedright
Seed type
\end{minipage} & \begin{minipage}[b]{\linewidth}\raggedright
Definition
\end{minipage} & \begin{minipage}[b]{\linewidth}\raggedright
Approximate target affinity
\end{minipage} \\
\midrule\noalign{}
\endhead
\bottomrule\noalign{}
\endlastfoot
6-mer & Perfect 6-bp seed (positions 2--7) & Weak (10--20 \%
repression) \\
7-mer-A1 & Seed + 3\texorpdfstring{'}{'} UTR adenosine at position 1 & Moderate
(about 30 \%) \\
7-mer-m8 & Seed + complementary base at position 8 & Moderate
(about 30 \%) \\
8-mer & 7-mer-A1 + 7-mer-m8 combined & Strong (about 50 \%) \\
3\texorpdfstring{'}{'} supplementary & Seed + 3\texorpdfstring{'}{'} end pairing (positions 13--16) & Modest
extra contribution \\
Centred site & 11-bp pairing centred on positions 4--14 & Moderate;
alternative class \\
\end{longtable}
}

Because each miRNA can have hundreds to thousands of targets
(TargetScan, miRanda, DIANA-microT predictions), miRNAs typically tune
target levels by 20--60 \% rather than completely silencing them ---
they are ``regulatory thermostats'' rather than on/off switches.

\textbf{siRNA (small interfering RNA)}: - 21 nt; derived from long dsRNA
(exogenous: viral, experimental; endogenous: transposons) - Biogenesis:
Dicer cleaves dsRNA; one strand loaded into RISC - Mechanism: perfect
complementarity to target \texorpdfstring{\ensuremath{\to}}{->} Argonaute-2 (Ago2) cleaves mRNA (``slicer''
activity) - \textbf{RNAi therapeutics}: Patisiran (Onpattro) -- first
FDA-approved siRNA drug (2018), targets TTR mRNA for transthyretin
amyloidosis; inclisiran (Leqvio) -- targets PCSK9 mRNA for
hypercholesterolemia (2021)

\textbf{Long non-coding RNA (lncRNA)}: - \textgreater200 nt;
\textgreater16,000 lncRNA genes in the human genome (likely more) -
Functions: \textbf{XIST} (X-inactivation), \textbf{HOTAIR} (recruits
PRC2 to silence HOXD cluster in trans), \textbf{MALAT1} (nuclear speckle
organization, splicing regulation), \textbf{NEAT1} (paraspeckle
formation)

\textbf{Circular RNA (circRNA)}: - Covalently closed RNA circles formed
by ``back-splicing'' (a downstream 5' SS joins to an upstream 3' SS) -
Functions: miRNA sponges (ciRS-7/CDR1as sponges miR-7, containing
\textgreater70 miR-7 binding sites), protein scaffolds, translated into
peptides in some cases

\begin{quote}
\textbf{Concept Check:} A circRNA bearing 70 miR-7 binding sites is
over-expressed in a cell --- predict the direction of change in the
protein output of a miR-7 target mRNA and explain the
post-transcriptional mechanism.
\end{quote}

\subsubsection{Epigenetic Regulation}\label{epigenetic-regulation}

\textbf{DNA methylation}: - Addition of a methyl group to the 5-position
of cytosine (5mC) in CpG dinucleotides - Catalyzed by \textbf{DNA
methyltransferases}: DNMT1 (maintenance methyltransferase; copies
methylation pattern after replication), DNMT3A/3B (de novo
methyltransferases) - \textbf{CpG islands} at gene promoters: usually
unmethylated in active genes; methylation recruits MeCP2 and MBD
proteins, which recruit HDACs and chromatin compactors, leading to
stable gene silencing - Aberrant methylation in cancer: global
hypomethylation (genome instability) + focal hypermethylation of tumor
suppressor promoters (e.g., RB1, p16/CDKN2A, BRCA1, MLH1)

\textbf{The histone code}:

{\def\LTcaptype{none} % do not increment counter
\begin{longtable}[]{@{}lll@{}}
\toprule\noalign{}
Modification & Location & Effect \\
\midrule\noalign{}
\endhead
\bottomrule\noalign{}
\endlastfoot
H3K4me3 & Promoters & Active transcription \\
H3K36me3 & Gene bodies & Active elongation \\
H3K27ac & Enhancers & Active enhancer \\
H3K27me3 & Promoters/enhancers & Polycomb-mediated silencing \\
H3K9me3 & Heterochromatin & Constitutive silencing (HP1 binding) \\
H4K20me3 & Heterochromatin & Repression \\
H3K4me1 & Enhancers & Poised enhancer (primed but not active) \\
\end{longtable}
}

\subsubsection{Gene Regulation in Development: Hox
Genes}\label{gene-regulation-in-development-hox-genes}

\textbf{Hox genes} encode homeodomain transcription factors that specify
positional identity along the anterior-posterior body axis:

\begin{itemize}
\tightlist
\item
  \textbf{Homeodomain}: about 60 amino acid DNA-binding domain
  (helix-turn-helix motif) recognizing AT-rich sequences
\item
  \textbf{Colinearity}: The order of Hox genes on the
  \hyperref[gl:chromosome]{\textbf{chromosome}} corresponds to their
  expression domains along the body axis (spatial colinearity) and the
  timing of their activation (temporal colinearity)
\item
  \textbf{Conservation}: Hox gene clusters are ancient and highly
  conserved from \emph{Drosophila} (2 clusters: ANT-C and BX-C, 8 genes
  total) to mammals (4 clusters: HOXA-HOXD, 39 genes total)
\item
  \textbf{Mutations}: Loss-of-function Hox mutations cause
  \textbf{homeotic transformations} (one body segment takes on the
  identity of another). In \emph{Drosophila}: \emph{Antennapedia}
  gain-of-function transforms antennae into legs; \emph{bithorax}
  loss-of-function transforms halteres into wings (creating a
  four-winged fly)
\end{itemize}

\begin{quote}
\textbf{Concept Check:} A posterior \emph{Hox} gene is ectopically
expressed in an anterior segment --- predict the homeotic transformation
that results and justify it from the spatial-colinearity rule.
\end{quote}

\textbf{Concept Check 12.3}

\begin{quote}
\begin{enumerate}
\def\labelenumi{\arabic{enumi}.}
\tightlist
\item
  Explain the dual regulation of the lac operon (negative + positive
  control). Under what conditions is the operon fully active?
\item
  How does the trp operon attenuation mechanism sense tryptophan levels?
\item
  Compare the functions of miRNA and siRNA. Why is RNAi a powerful
  research tool?
\item
  What is the histone code? Give two examples of activating marks and
  two of repressive marks.
\end{enumerate}
\end{quote}

\begin{center}\rule{0.5\linewidth}{0.5pt}\end{center}

\subsection{Worked Example: Reading Frame and Protein
Prediction}\label{worked-example-reading-frame-and-protein-prediction}

\textbf{Problem}: The following mRNA sequence is expressed in a
eukaryotic cell:

5'-AUGUUCAAGGACUAUUGCCCGUAGACUU-3'

\begin{enumerate}
\def\labelenumi{(\alph{enumi})}
\tightlist
\item
  Translate this sequence.
\end{enumerate}

\textbf{Solution}: Reading in triplets from AUG:

\begin{equation}\text{AUG-UUC-AAG-GAC-UAU-UGC-CCG-UAG-ACU-U} \tag{13.2}\label{eq:reading_frame}\end{equation}

{\def\LTcaptype{none} % do not increment counter
\begin{longtable}[]{@{}ll@{}}
\toprule\noalign{}
Codon & Amino Acid \\
\midrule\noalign{}
\endhead
\bottomrule\noalign{}
\endlastfoot
AUG & Met \\
UUC & Phe \\
AAG & Lys \\
GAC & Asp \\
UAU & Tyr \\
UGC & Cys \\
CCG & Pro \\
UAG & Stop (amber) \\
\end{longtable}
}

\textbf{Protein}: Met-Phe-Lys-Asp-Tyr-Cys-Pro (7 amino acids)

\begin{enumerate}
\def\labelenumi{(\alph{enumi})}
\setcounter{enumi}{1}
\tightlist
\item
  If a single G is inserted after the first codon
  (AUG\textbf{G}UUC\ldots), what happens?
\end{enumerate}

\textbf{Solution}: Frameshift mutation. New reading frame:

\begin{equation}\text{AUG-GUU-CAA-GGA-CUA-UUG-CCC-GUA-GAC-UU} \tag{13.3}\label{eq:frameshift}\end{equation}

Met-Val-Gln-Gly-Leu-Leu-Pro-Val-Asp\ldots{} (completely different
protein, no stop codon in this segment)

\begin{center}\rule{0.5\linewidth}{0.5pt}\end{center}

\subsection{Computational Bridge}\label{computational-bridge-12}

The same genetic code table used in the text is callable from
\texttt{biology.genetics}:

\begin{Shaded}
\begin{Highlighting}[]
\ImportTok{from}\NormalTok{ biology.genetics }\ImportTok{import}\NormalTok{ transcribe\_dna\_to\_mrna, translate\_mrna}

\NormalTok{mrna }\OperatorTok{=}\NormalTok{ transcribe\_dna\_to\_mrna(}\StringTok{"TACGGCTTGTTC"}\NormalTok{)}
\BuiltInTok{print}\NormalTok{(}\StringTok{" "}\NormalTok{.join(translate\_mrna(mrna)[:}\DecValTok{6}\NormalTok{]))}
\end{Highlighting}
\end{Shaded}

\begin{quote}
\textbf{Clinical / systems note:} NMD and cap-dependent scanning are
therapeutic pressure points: nonsense mutations in tumour suppressors
can be targeted with read-through compounds, while eIF4A inhibitors
attempt to collapse translation of highly structured oncogenic mRNAs.
\end{quote}

\begin{center}\rule{0.5\linewidth}{0.5pt}\end{center}

\subsection{Current Evidence and Frontier
Biology}\label{current-evidence-and-frontier-biology-15}

For \textbf{Gene Expression --- Transcription and Translation}, frontier
biology belongs inside the evidence logic of the chapter. Molecular
genetics now spans single-reference sequences, telomere-to-telomere
assemblies, pangenome graphs, long-read sequencing, CRISPR medicines,
and ethical deployment. The core reading question is this: expression
claims should separate transcription, RNA processing, translation,
localization, degradation, and feedback.

\begin{itemize}
\tightlist
\item
  \textbf{What to verify:} identify the observation, model, assay, or
  dataset that would make the claim stronger or weaker.
\item
  \textbf{What to qualify:} state the scale, organism, cell type,
  environmental condition, or population where the claim is expected to
  hold.
\item
  \textbf{What to compare:} test at least one alternative explanation,
  baseline, or null model before treating the pattern as causal.
\item
  \textbf{What to cite:} distinguish primary evidence, review synthesis,
  public dataset, and institutional guidance; for recent or numeric
  claims, prefer the source closest to the measurement and state what
  has changed since it was published.
\end{itemize}

When a genomic claim depends on a reference, ask whether short reads,
structural variants, ancestry representation, phasing, or clinical
validation could change the interpretation
\citep{humanpangenome2023,fda2023casgevy,fda2024casgevythalassemia}.

\textbf{Source practice:} For genomics and editing claims, distinguish
discovery from clinical actionability, and cite reference resources,
regulatory records, or primary editing studies close to the claim
\citep{humanpangenome2023,fda2026casgevy,chalumeau2025primeediting}.

\subsection{Summary}\label{summary-15}

\begin{itemize}
\tightlist
\item
  \textbf{Transcription}: RNAP reads the template strand 3' to 5' to
  produce RNA 5' to 3'. Prokaryotes use a single RNAP with sigma
  factors; eukaryotes use three RNA Polymerases plus general
  transcription factors and Mediator.
\item
  \textbf{Promoter elements}: TATA box (TBP binding), Inr, DPE, BRE, CpG
  islands; enhancers act over long distances via DNA looping.
\item
  \textbf{mRNA processing}: 5' m\(^7\)G cap (eIF4E recognition,
  stability), spliceosome-mediated intron removal (two
  transesterification reactions via lariat intermediate), 3' poly(A)
  tail (PABP, stability, translation).
\item
  \textbf{Alternative splicing}: \textgreater95\% of human multi-exon
  genes; generates proteome diversity; regulated by SR proteins and
  hnRNPs; NMD degrades transcripts with premature stop codons.
\item
  \textbf{Translation}: cap-dependent scanning (eukaryotes) or
  Shine-Dalgarno (prokaryotes); A/P/E sites; elongation cycle; genetic
  code (64 codons, wobble pairing, nearly comprehensive); ribosome is a
  ribozyme.
\item
  \textbf{Post-translational modifications}: phosphorylation,
  glycosylation, ubiquitination, acetylation, methylation, SUMOylation;
  ubiquitin-proteasome system for targeted degradation.
\item
  \textbf{Gene regulation}: prokaryotic operons (lac: inducible; trp:
  repressible with attenuation); eukaryotic: chromatin remodeling,
  histone code, DNA methylation, transcription factor combinatorics,
  miRNA/siRNA/lncRNA.
\item
  \textbf{\hyperref[gl:epigenetics]{Epigenetics}}: heritable changes in
  gene expression without DNA sequence change; DNA methylation (CpG
  islands), histone modifications; dysregulated in cancer.
\item
  \textbf{Development}: Hox genes specify positional identity;
  homeodomain TFs; colinearity of gene order and expression domain.
\item
  \textbf{Connections:} See
  \cref{sec:unit_IV_epigenetics_and_gene_regulation} for chromatin and
  RNAi, Unit VI for how expression variation fuels evolution, and
  \cref{sec:unit_II_cell_signaling} for signal-driven transcription
  factors.
\end{itemize}

\begin{center}\rule{0.5\linewidth}{0.5pt}\end{center}

\subsection{Review Questions}\label{review-questions-15}

\begin{enumerate}
\def\labelenumi{\arabic{enumi}.}
\tightlist
\item
  Compare the promoter recognition mechanisms in prokaryotes (sigma
  factor) and eukaryotes (GTF assembly). Why do eukaryotes require so
  many additional factors?
\item
  Describe the three enzymatic steps of 5' capping. Why is the cap
  essential for translation?
\item
  Explain the two transesterification reactions in splicing. What is the
  lariat intermediate and why does it form?
\item
  How does DSCAM alternative splicing generate 38,016 isoforms from a
  single gene? What biological function does this diversity serve?
\item
  Explain the NMD pathway. Why is it important to destroy mRNAs with
  premature stop codons?
\item
  Compare the mechanisms of eukaryotic and prokaryotic translation
  initiation. What is the Kozak sequence?
\item
  Describe the ubiquitin-proteasome pathway (E1/E2/E3 cascade). Why does
  the human genome encode about 600 E3 ligases but 2 main E1
  enzymes?
\item
  Compare the lac operon (inducible) and trp operon (repressible). How
  does attenuation fine-tune trp operon expression?
\item
  Describe three classes of non-coding RNAs and their roles in gene
  regulation.
\item
  What is the histone code? How do writers, readers, and erasers of
  histone marks coordinate gene regulation?
\item
  Insert a single-base frameshift in a toy 30 nt mRNA after the start
  codon and rerun \texttt{translate\_mrna} in Python. At what point do
  stop codons typically appear compared with in-frame controls?
\item
  Compare how a \textbf{uORF} upstream of a main AUG might tune
  translation initiation without changing DNA sequence.
\end{enumerate}

\begin{center}\rule{0.5\linewidth}{0.5pt}\end{center}

\subsection{Further Reading and Source
Notes}\label{further-reading-and-source-notes-15}

\begin{itemize}
\tightlist
\item
  Crick (1958). On protein synthesis. \emph{Symposia of the Society for
  Experimental Biology}, 12.
\item
  Crick (1966). The genetic code --- yesterday, today and tomorrow.
  \emph{Cold Spring Harbor Symposia on Quantitative Biology}, 31.
\item
  Jacob \& Monod (1961). Genetic regulatory mechanisms in the synthesis
  of proteins. \emph{Journal of Molecular Biology}, 3.
\end{itemize}

\begin{center}\rule{0.5\linewidth}{0.5pt}\end{center}

\subsection{Key Terms}\label{key-terms-15}

\begin{enumerate}
\def\labelenumi{\arabic{enumi}.}
\tightlist
\item
  \textbf{Central dogma} -- the flow of genetic information: DNA to RNA
  to protein
\item
  \textbf{Sigma factor} -- prokaryotic RNAP subunit that recognizes
  promoter sequences
\item
  \textbf{General transcription factors (GTFs)} -- proteins (TFIIA-H)
  required for basal Pol II transcription
\item
  \textbf{Mediator complex} -- coactivator bridging enhancer-bound TFs
  to basal machinery
\item
  \textbf{5' cap (m7G)} -- modified guanosine protecting mRNA 5' end;
  recognized by eIF4E
\item
  \textbf{Spliceosome} -- large RNP complex (U1, U2, U4, U5, U6 snRNPs)
  that removes introns
\item
  \textbf{Branch point} -- intronic adenosine whose 2'-OH attacks the 5'
  splice site in the first transesterification
\item
  \textbf{Alternative splicing} -- generation of multiple mRNA isoforms
  from one gene by differential exon inclusion
\item
  \textbf{Nonsense-mediated decay (NMD)} -- mRNA surveillance destroying
  transcripts with premature stop codons
\item
  \textbf{Wobble base pairing} -- relaxed pairing rules at codon
  position 3/anticodon position 1
\item
  \textbf{Kozak sequence} -- consensus context around eukaryotic AUG
  start codon (RCCAUGG)
\item
  \textbf{Peptidyl transferase} -- ribosomal RNA-based catalytic
  activity forming peptide bonds
\item
  \textbf{Ubiquitin-proteasome system} -- pathway for targeted protein
  degradation via polyubiquitination
\item
  \textbf{miRNA} -- about 22 nt non-coding RNA guiding RISC to
  target mRNAs for translational repression
\item
  \textbf{Epigenetics} -- heritable changes in gene expression without
  alterations to DNA sequence
\item
  \textbf{CpG island} -- GC-rich genomic region; methylation status
  controls gene activity
\item
  \textbf{Histone code} -- combinatorial histone modifications that
  regulate chromatin state and gene expression
\item
  \textbf{Homeodomain} -- 60-aa DNA-binding domain in Hox transcription
  factors specifying positional identity
\item
  \textbf{Operon} -- cluster of co-transcribed prokaryotic genes under
  shared regulatory control
\end{enumerate}

\begin{center}\rule{0.5\linewidth}{0.5pt}\end{center}

\subsubsection{Companion Source
Module}\label{companion-source-module-15}

\textbf{Gene Expression --- Transcription and Translation} should leave
a reproducible trail from a biological claim to the code, figure,
diagram, or paper-based activity that can test it. Use the surfaces
below to inspect the chapter's assumptions, rerun the relevant model, or
compare the manuscript explanation with companion labs and figures.

{\def\LTcaptype{none} % do not increment counter
\begin{longtable}[]{@{}
  >{\raggedright\arraybackslash}p{(\linewidth - 2\tabcolsep) * \real{0.5000}}
  >{\raggedright\arraybackslash}p{(\linewidth - 2\tabcolsep) * \real{0.5000}}@{}}
\toprule\noalign{}
\begin{minipage}[b]{\linewidth}\raggedright
Surface
\end{minipage} & \begin{minipage}[b]{\linewidth}\raggedright
Use it for
\end{minipage} \\
\midrule\noalign{}
\endhead
\bottomrule\noalign{}
\endlastfoot
\breaktt{src/biology/genetics/genetics.py}
(\breaktt{transcribe\_dna\_to\_mrna}, \texttt{translate\_mrna},
\texttt{gc\_content}) & Reproduce transcription, translation, codon
lookup, and sequence-composition checks. \\
\breaktt{src/mermaid/biology\_diagrams.py}
(\breaktt{transcription\_translation\_diagram},
\breaktt{mirna\_biogenesis\_diagram}) & Connect coding flow with RNA
regulation. \\
\end{longtable}
}

\textbf{Reproducibility check:} specify template strand, reading frame,
RNA-processing assumptions, and regulatory layer before interpreting
expression. \textbf{Cross-reference:} use
\cref{sec:unit_IV_epigenetics_and_gene_regulation} and
\cref{sec:unit_I_macromolecules}.

\newpage

\section{Mutations, CRISPR, and
Genomics}\label{mutations-crispr-and-genomics}

\label{sec:unit_IV_mutations_and_genomics}

\begin{quote}
\textbf{Ch 14} \texorpdfstring{\ensuremath{\cdot}}{.} Level 2/3 \texorpdfstring{\ensuremath{\cdot}}{.} 55 min read \texorpdfstring{\ensuremath{\cdot}}{.} 75 min lecture \texorpdfstring{\ensuremath{\cdot}}{.}
Prerequisites: \cref{sec:unit_IV_gene_expression}
\end{quote}

\subsection{Learning Objectives}\label{learning-objectives-16}

\begin{enumerate}
\def\labelenumi{\arabic{enumi}.}
\tightlist
\item
  Classify \hyperref[gl:mutation]{\textbf{mutation}}s by type (point,
  frameshift, chromosomal), molecular change, and functional
  consequence.
\item
  Describe mutagenic agents and assays for mutagenicity (Ames test).
\item
  Explain DNA repair mechanisms and their relationship to human disease.
\item
  Describe the CRISPR-Cas9 mechanism in detail, including PAM
  recognition, DSB creation, and repair outcomes.
\item
  Compare CRISPR-Cas9, base editing, and prime editing technologies.
\item
  Define the \hyperref[gl:genome]{\textbf{genome}} and describe key
  genomic features revealed by sequencing.
\item
  Explain next-generation sequencing technologies and their
  applications.
\item
  Describe GWAS methodology and its contributions to understanding
  complex disease.
\item
  Evaluate ethical issues in personal genomics and genetic testing.
\end{enumerate}

\subsubsection{Study Blueprint}\label{study-blueprint-16}

\begin{itemize}
\tightlist
\item
  \textbf{Big idea:} Genomic variation becomes biological consequence
  through sequence context, repair, and selection.
\item
  \textbf{Core concepts:} mutation classes, DNA repair, CRISPR, genome
  analysis.
\item
  \textbf{Framework alignment:} Vision \& Change: Information flow,
  exchange, and storage, Structure and function; AP Biology: Information
  Storage and Transmission, Systems Interactions; NGSS-style topics:
  Inheritance and Variation of Traits, Structure and Function.
\item
  \textbf{Model or quantitative lens:} Mutation-rate, edit-efficiency,
  and sequence-comparison calculations.
\item
  \textbf{Data skill:} Classify variants and predict likely molecular
  effect from sequence evidence.
\item
  \textbf{Practice cadence:} Concept Explanation, Questions and Methods,
  Argumentation.
\item
  \textbf{Common misconception to repair:} Not every mutation is
  harmful, and not every harmful mutation changes a protein sequence.
\item
  \textbf{Primary lab:} \cref{sec:lab_unit_IV_mutations_and_genomics}.
\item
  \textbf{Question bank:} \cref{sec:q_unit_IV_mutations_and_genomics}.
\item
  \textbf{Transfer task:} Transfer variant reasoning to cancer genomics,
  ancestry, gene therapy, or microbial evolution.
\item
  \textbf{Bridge to computation:}
  \breaktt{biology.genetics.genetics.hamming\_distance}.
\end{itemize}

\begin{center}\rule{0.5\linewidth}{0.5pt}\end{center}

\begin{quote}
\textbf{Opening Vignette --- One Letter Changes Everything}

In 1949, Linus Pauling published a paper in \emph{Science} titled
``Sickle Cell Anemia, a Molecular Disease'' --- the first time a human
illness was attributed to a chemical change in a specific
\hyperref[gl:protein]{\textbf{protein}}. Eight years later, Vernon
Ingram used the newly developed technique of protein fingerprinting to
identify the precise mutation: a single amino acid substitution,
glutamic acid to valine, at position 6 of the \texorpdfstring{\ensuremath{\beta}}{beta}-globin chain. Just one
\hyperref[gl:nucleotide]{\textbf{nucleotide}} change --- GAG to GTG ---
makes haemoglobin polymerise under low-oxygen conditions, distorting red
blood cells into rigid sickle shapes that clog capillaries. That one
point mutation causes pain crises, organ damage, and shortened lifespan
for millions worldwide. The sickle-cell story remains the most
compelling illustration of how a single mutational event can cascade
through molecular structure, protein function, cell physiology, and
whole-organism health.
\end{quote}

\subsection{Types of Mutations}\label{types-of-mutations}

A \textbf{mutation} is any heritable change in DNA sequence. Mutations
are the ultimate source of most genetic variation and therefore the raw
material for evolution.

\subsubsection{Mutation Rates: A Quantitative
Framework}\label{mutation-rates-a-quantitative-framework}

The \textbf{per-nucleotide per-generation mutation rate} in humans is
approximately \(\mu \approx 1.2 \times 10^{-8}\) per bp per generation
(Kong et al., 2012, \emph{Nature}). For a diploid genome of
\(2N \approx 6.4 \times 10^9\) bp, the expected number of de novo
mutations per offspring is:

\begin{equation}M = 2N \times \mu \approx 6.4 \times 10^9 \times 1.2 \times 10^{-8} \approx 70\text{--}80 \text{ mutations/generation} \tag{14.1}\label{eq:mutations_per_gen}\end{equation}

The probability of observing exactly \(k\) mutations in a region of
length \(L\) follows the \textbf{Poisson distribution}:

\begin{equation}P(k) = \frac{(\mu L)^k \, e^{-\mu L}}{k!} \tag{14.2}\label{eq:poisson_mut}\end{equation}

This is used to assess whether a \hyperref[gl:gene]{\textbf{gene}} or
genomic region has an excess of mutations relative to the background
rate --- a signature of positive selection or a mutational hotspot.

For genome-wide studies, the \textbf{Bonferroni-corrected significance
threshold} for \(m\) independent tests is:

\begin{equation}\alpha_{\text{corrected}} = \frac{\alpha}{m} = \frac{0.05}{10^6} = 5 \times 10^{-8} \tag{14.3}\label{eq:bonferroni_correction}\end{equation}

This is the standard genome-wide significance threshold for GWAS:
primarily associations with \(p < 5 \times 10^{-8}\) are considered
robust.

The \textbf{Ka/Ks ratio} (also called dN/dS or \texorpdfstring{\ensuremath{\omega}}{omega}) measures the ratio of
nonsynonymous to synonymous substitution rates, revealing the mode of
\hyperref[gl:natural-selection]{\textbf{natural selection}}:

\begin{equation}\omega = \frac{K_a}{K_s} = \frac{d_N}{d_S} \tag{14.4}\label{eq:omega_ratio}\end{equation}

where \(\omega < 1\) indicates purifying selection, \(\omega = 1\)
neutral evolution, and \(\omega > 1\) positive selection.

\subsubsection{Nucleotide-Level (Point)
Mutations}\label{nucleotide-level-point-mutations}

{\def\LTcaptype{none} % do not increment counter
\begin{longtable}[]{@{}
  >{\raggedright\arraybackslash}p{(\linewidth - 4\tabcolsep) * \real{0.1500}}
  >{\raggedright\arraybackslash}p{(\linewidth - 4\tabcolsep) * \real{0.3250}}
  >{\raggedright\arraybackslash}p{(\linewidth - 4\tabcolsep) * \real{0.5250}}@{}}
\toprule\noalign{}
\begin{minipage}[b]{\linewidth}\raggedright
Type
\end{minipage} & \begin{minipage}[b]{\linewidth}\raggedright
Description
\end{minipage} & \begin{minipage}[b]{\linewidth}\raggedright
Example Consequence
\end{minipage} \\
\midrule\noalign{}
\endhead
\bottomrule\noalign{}
\endlastfoot
\textbf{Transition} & Purine to purine (A to G) or pyrimidine to
pyrimidine (C to T) & Most common; often neutral or conservative amino
acid change \\
\textbf{Transversion} & Purine to pyrimidine or vice versa (A to C, A to
T, G to C, G to T) & Rarer (about 2:1 transition:transversion
ratio); more often disruptive \\
\textbf{Silent (synonymous)} & \hyperref[gl:codon]{\textbf{Codon}}
change producing the same amino acid (degeneracy) & Usually no
phenotypic effect; may affect splicing or mRNA stability \\
\textbf{Missense} & Codon change producing a different amino acid &
Conservative (e.g., Leu to Ile) or non-conservative (e.g., Glu to Val in
sickle cell) \\
\textbf{Nonsense} & Point mutation creating a premature stop codon &
Truncated protein; usually non-functional; often triggers NMD \\
\textbf{Splice site} & Mutation at
\hyperref[gl:intron]{\textbf{intron}}-\hyperref[gl:exon]{\textbf{exon}}
boundary (GT/AG dinucleotides) & Exon skipping, intron retention, or
cryptic splice site activation \\
\end{longtable}
}

\subsubsection{Insertions, Deletions, and
Frameshifts}\label{insertions-deletions-and-frameshifts}

\begin{itemize}
\tightlist
\item
  \textbf{Insertion}: One or more nucleotides added to the sequence
\item
  \textbf{Deletion}: One or more nucleotides removed from the sequence
\item
  \textbf{Frameshift}: An insertion or deletion that is NOT a multiple
  of 3 bp -- shifts the entire downstream reading frame, producing a
  completely altered amino acid sequence and usually encountering a
  premature stop codon
\item
  \textbf{In-frame insertion/deletion}: A multiple of 3 bp -- adds or
  removes amino acids without disrupting the reading frame (e.g., the
  deltaF508 mutation in CFTR is a 3-bp deletion removing Phe508, causing
  cystic fibrosis)
\end{itemize}

\subsubsection{Copy Number Variants and Structural
Mutations}\label{copy-number-variants-and-structural-mutations}

\begin{itemize}
\tightlist
\item
  \textbf{Copy number variants (CNVs)}: Duplications or deletions of
  segments from 1 kb to several Mb. The average human genome contains
  about 1,000 CNVs relative to the reference. CNVs contribute
  more base pairs of variation between individuals than SNPs.
\item
  \textbf{Inversions}: A chromosomal segment is reversed end-to-end.
  Pericentric inversions include the
  \hyperref[gl:centromere]{\textbf{centromere}}; paracentric do not.
  Inversions can disrupt genes at breakpoints or alter regulatory
  landscapes.
\item
  \textbf{Translocations}: A segment moves to a different
  \hyperref[gl:chromosome]{\textbf{chromosome}}. \textbf{Reciprocal
  translocations} exchange segments between two chromosomes.
  \textbf{Robertsonian translocations} fuse the long arms of two
  acrocentric chromosomes (13, 14, 15, 21, 22).
\end{itemize}

\begin{quote}
\textbf{Clinical Connection: The Philadelphia Chromosome} The t(9;22)
reciprocal translocation creates the \textbf{Philadelphia chromosome} --
a derivative chromosome 22 carrying the \textbf{BCR-ABL1} fusion gene.
This fusion produces a constitutively active tyrosine kinase that drives
chronic myelogenous leukemia (CML). The targeted tyrosine kinase
inhibitor \textbf{imatinib (Gleevec)} revolutionized CML therapy,
converting a fatal disease into a manageable chronic condition (5-year
survival increased from about 30\% to \textgreater90\%). This
is a paradigm for targeted cancer therapy.
\end{quote}

\subsubsection{Trinucleotide Repeat Expansion
Disorders}\label{trinucleotide-repeat-expansion-disorders}

A special class of mutations involving expansion of short tandem
repeats:

{\def\LTcaptype{none} % do not increment counter
\begin{longtable}[]{@{}
  >{\raggedright\arraybackslash}p{(\linewidth - 10\tabcolsep) * \real{0.1324}}
  >{\raggedright\arraybackslash}p{(\linewidth - 10\tabcolsep) * \real{0.1176}}
  >{\raggedright\arraybackslash}p{(\linewidth - 10\tabcolsep) * \real{0.1912}}
  >{\raggedright\arraybackslash}p{(\linewidth - 10\tabcolsep) * \real{0.2500}}
  >{\raggedright\arraybackslash}p{(\linewidth - 10\tabcolsep) * \real{0.1471}}
  >{\raggedright\arraybackslash}p{(\linewidth - 10\tabcolsep) * \real{0.1618}}@{}}
\toprule\noalign{}
\begin{minipage}[b]{\linewidth}\raggedright
Disease
\end{minipage} & \begin{minipage}[b]{\linewidth}\raggedright
Repeat
\end{minipage} & \begin{minipage}[b]{\linewidth}\raggedright
Normal Range
\end{minipage} & \begin{minipage}[b]{\linewidth}\raggedright
Pathogenic Range
\end{minipage} & \begin{minipage}[b]{\linewidth}\raggedright
Location
\end{minipage} & \begin{minipage}[b]{\linewidth}\raggedright
Mechanism
\end{minipage} \\
\midrule\noalign{}
\endhead
\bottomrule\noalign{}
\endlastfoot
Huntington disease & CAG (Gln) & 6-35 & \textgreater36 (full penetrance
\textgreater40) & HTT exon 1 & Toxic polyglutamine aggregation \\
Fragile X syndrome & CGG & 5-44 & \textgreater200 (full mutation) & FMR1
5' UTR & CpG methylation silences FMR1 \\
Myotonic dystrophy type 1 & CTG & 5-34 & \textgreater50 & DMPK 3' UTR &
Toxic RNA foci sequester MBNL1 \\
Friedreich ataxia & GAA & 5-33 & \textgreater66 & FXN intron 1 &
\hyperref[gl:heterochromatin]{\textbf{Heterochromatin}} formation
silences FXN \\
\end{longtable}
}

\textbf{Anticipation}: These diseases show \textbf{genetic anticipation}
-- earlier onset and increased severity in successive generations --
because the unstable repeats tend to expand further during DNA
replication (particularly during germ cell divisions).

\textbf{Concept Check 13.1}

\begin{quote}
\begin{enumerate}
\def\labelenumi{\arabic{enumi}.}
\tightlist
\item
  Classify the sickle cell mutation (GAG to GUG in beta-globin codon 6)
  by type.
\item
  Why are frameshift mutations generally more deleterious than missense
  mutations?
\item
  Explain why trinucleotide repeat diseases show anticipation.
\item
  The deltaF508 CFTR mutation is a 3-bp deletion. Why is this NOT a
  frameshift?
\end{enumerate}
\end{quote}

\begin{center}\rule{0.5\linewidth}{0.5pt}\end{center}

\subsection{Mutagenic Agents}\label{mutagenic-agents}

\subsubsection{Physical Mutagens}\label{physical-mutagens}

\textbf{UV radiation (UVB, 280-320 nm)}: - Creates \textbf{cyclobutane
pyrimidine dimers (CPDs)}: covalent bonds form between adjacent
pyrimidines (T-T most common) - Creates \textbf{6-4 photoproducts (6-4
PPs)}: covalent bond between positions 6 and 4 of adjacent pyrimidines -
Both distort the helix and block
replication/\hyperref[gl:transcription]{\textbf{transcription}} -
Repaired by nucleotide excision repair (NER); failure causes xeroderma
pigmentosum

\textbf{Ionizing radiation (X-rays, gamma rays)}: - Creates reactive
oxygen species (ROS) that attack bases and the sugar-phosphate backbone
- Causes single-strand breaks (SSBs) and double-strand breaks (DSBs) -
DSBs are the most dangerous lesion -- if unrepaired, lead to chromosomal
translocations, deletions, or cell death - Repaired by NHEJ and HR

\subsubsection{Chemical Mutagens}\label{chemical-mutagens}

{\def\LTcaptype{none} % do not increment counter
\begin{longtable}[]{@{}
  >{\raggedright\arraybackslash}p{(\linewidth - 6\tabcolsep) * \real{0.1750}}
  >{\raggedright\arraybackslash}p{(\linewidth - 6\tabcolsep) * \real{0.1500}}
  >{\raggedright\arraybackslash}p{(\linewidth - 6\tabcolsep) * \real{0.2750}}
  >{\raggedright\arraybackslash}p{(\linewidth - 6\tabcolsep) * \real{0.4000}}@{}}
\toprule\noalign{}
\begin{minipage}[b]{\linewidth}\raggedright
Agent
\end{minipage} & \begin{minipage}[b]{\linewidth}\raggedright
Type
\end{minipage} & \begin{minipage}[b]{\linewidth}\raggedright
Mechanism
\end{minipage} & \begin{minipage}[b]{\linewidth}\raggedright
Mutation Caused
\end{minipage} \\
\midrule\noalign{}
\endhead
\bottomrule\noalign{}
\endlastfoot
\textbf{EMS} (ethyl methanesulfonate) & Alkylating agent & Adds ethyl
group to O\(^6\) of guanine & G:C to A:T transitions \\
\textbf{Nitrogen mustard} & Bifunctional alkylating agent & Cross-links
DNA strands & Blocks replication; DSBs \\
\textbf{Nitrous acid (HNO\(_2\))} & Deaminating agent & Deaminates C to
U, A to hypoxanthine & C:G to T:A and A:T to G:C transitions \\
\textbf{Ethidium bromide} & Intercalating agent & Inserts between base
pairs, stretching the helix & Frameshift mutations
(insertions/deletions) \\
\textbf{5-Bromouracil (5-BU)} & Base analog & Incorporated in place of
T; tautomerizes to pair with G & A:T to G:C transitions \\
\textbf{Acridine orange} & Intercalating agent & Inserts between bases &
Frameshift mutations \\
\textbf{Benzo\hyperref[a]{a}pyrene} (cigarette smoke) & Polycyclic
aromatic hydrocarbon & Forms bulky adducts on guanine (after CYP1A1
activation) & G:C to T:A transversions; NER substrate \\
\end{longtable}
}

\textbf{Reactive oxygen species (ROS)}: Endogenous byproducts of
\hyperref[gl:aerobic]{\textbf{aerobic}} metabolism (superoxide
O\(_2^{-}\), hydroxyl radical OH\(\cdot\), hydrogen peroxide
H\(_2\)O\(_2\)). ROS oxidize guanine to \textbf{8-oxoguanine (8-oxoG)},
which mispairs with adenine, causing G:C to T:A transversions. Repaired
by BER (OGG1 glycosylase).

\subsubsection{The Ames Test}\label{the-ames-test}

The \textbf{Ames test} (Bruce Ames, 1973) is a bacterial mutagenicity
assay that uses \emph{Salmonella typhimurium} strains carrying various
\textbf{his\(^-\)} (histidine-requiring) mutations:

\textbf{Procedure}: 1. His\(^-\) bacteria are plated on minimal medium
(no histidine) 2. Test chemical applied (with and without liver
microsomal extract S9 to simulate mammalian metabolic activation) 3.
Count revertant colonies (his\(^+\)) that can grow without histidine 4.
Compare to spontaneous reversion rate (negative control) and known
mutagen (positive control)

\textbf{Interpretation}: A significant increase in revertant colonies
indicates the chemical is mutagenic. Different tester strains detect
different mutation types (e.g., TA98 detects frameshifts; TA100 detects
base substitutions).

\textbf{Correlation}: about 80-90\% of known carcinogens are
mutagenic in the Ames test. The Ames test is a standard regulatory tool
for chemical safety screening.

\begin{center}\rule{0.5\linewidth}{0.5pt}\end{center}

\subsection{DNA Damage and Repair}\label{dna-damage-and-repair-1}

The human genome sustains about 10,000-100,000 DNA lesions per
cell per day \citep{ward1988}. Multiple repair pathways maintain genomic
integrity.

\subsubsection{DNA Damage --- Quantitative
Background}\label{dna-damage-quantitative-background}

The endogenous and exogenous lesion load is enormous and well
characterised quantitatively:

{\def\LTcaptype{none} % do not increment counter
\begin{longtable}[]{@{}
  >{\raggedright\arraybackslash}p{(\linewidth - 4\tabcolsep) * \real{0.2444}}
  >{\raggedright\arraybackslash}p{(\linewidth - 4\tabcolsep) * \real{0.1333}}
  >{\raggedright\arraybackslash}p{(\linewidth - 4\tabcolsep) * \real{0.6222}}@{}}
\toprule\noalign{}
\begin{minipage}[b]{\linewidth}\raggedright
Lesion type
\end{minipage} & \begin{minipage}[b]{\linewidth}\raggedright
Source
\end{minipage} & \begin{minipage}[b]{\linewidth}\raggedright
Frequency (per cell per day)
\end{minipage} \\
\midrule\noalign{}
\endhead
\bottomrule\noalign{}
\endlastfoot
Depurinations / abasic (AP) sites & Spontaneous hydrolysis of
N-glycosidic bond & about 10,000 (purine more labile than
pyrimidine) \\
Depyrimidinations & Hydrolysis & about 500 \\
Cytosine deaminations (C \texorpdfstring{\ensuremath{\to}}{->} U) & Hydrolysis & 100--500 \\
5-methylcytosine deaminations (5mC \texorpdfstring{\ensuremath{\to}}{->} T) & Hydrolysis (faster than C
deamination) & about 50--100 (CpG-localized) \\
Oxidised guanines (8-oxoG) & ROS attack on guanine N7 &
about 10,000 \\
Other oxidised bases (FapyG, FapyA, 5-OH-cytosine) & ROS &
about 3,000 \\
Single-strand breaks (SSBs) & Direct breakage; sugar-radical chemistry &
20,000--50,000 \\
Inter-strand crosslinks (ICLs) & Endogenous aldehydes (acetaldehyde,
formaldehyde) & about 10--50 \\
\textbf{Replication-induced double-strand breaks} & Fork collapse at
SSBs, R-loops, repeats & \textbf{about 100--500 during
replication} (S-phase primarily) \\
Spontaneous DSBs (non-replicative) & Two coincident SSBs;
radical-induced & about 10--50 \\
UVB-induced CPDs & UV photons (sunlight, UVB 280--320 nm) & Highly
variable; up to 10\texorpdfstring{\textsuperscript{6}}{6} per skin cell per day with sun exposure \\
Alkylated bases (O\texorpdfstring{\textsuperscript{6}}{6}-methylguanine, N3-methylpurine) & Endogenous methyl
donors (S-adenosylmethionine) + exogenous alkylators &
about 3,000 \\
Topoisomerase-induced breaks & Stalled topoisomerase Ic-DNA complexes &
Variable; increased by camptothecin, etoposide \\
\end{longtable}
}

\textbf{Replication-associated DSBs} (about 100--500 per cell
during S phase) are particularly dangerous because they break a DNA
strand that is also undergoing semi-conservative replication, producing
daughter molecules with persistent breaks. These represent the bulk of
physiologically encountered DSBs and explain why HR, the high-fidelity
DSB-repair pathway, is preferentially active in S/G2 phase.

The cell maintains genomic integrity by directing each lesion to its
appropriate repair pathway. The remarkable per-base error rate of
replication (about 10\texorpdfstring{\textsuperscript{-}}{-}\texorpdfstring{\textsuperscript{1}}{1}\texorpdfstring{\textsuperscript{0}}{0}) is achieved through a combination of
polymerase fidelity, proofreading, mismatch repair (post-replicative),
and the various lesion-specific pathways below.

\begin{figure}[htbp]
\centering
\includegraphics[width=0.82\linewidth,height=4.2in,keepaspectratio]{../figures/mermaid_inline/inline_mermaid_0075_805c39e2ceba.png}
\caption{DNA damage types and their corresponding repair pathways. The choice of DSB repair pathway depends on cell cycle phase: NHEJ predominates in G1, while HR requires a sister chromatid template available in S/G2.}
\end{figure}

\subsubsection{Repair Pathway Details}\label{repair-pathway-details}

{\def\LTcaptype{none} % do not increment counter
\begin{longtable}[]{@{}
  >{\raggedright\arraybackslash}p{(\linewidth - 6\tabcolsep) * \real{0.2097}}
  >{\raggedright\arraybackslash}p{(\linewidth - 6\tabcolsep) * \real{0.2419}}
  >{\raggedright\arraybackslash}p{(\linewidth - 6\tabcolsep) * \real{0.2097}}
  >{\raggedright\arraybackslash}p{(\linewidth - 6\tabcolsep) * \real{0.3387}}@{}}
\toprule\noalign{}
\begin{minipage}[b]{\linewidth}\raggedright
Damage Type
\end{minipage} & \begin{minipage}[b]{\linewidth}\raggedright
Repair Pathway
\end{minipage} & \begin{minipage}[b]{\linewidth}\raggedright
Key Proteins
\end{minipage} & \begin{minipage}[b]{\linewidth}\raggedright
Disease if Defective
\end{minipage} \\
\midrule\noalign{}
\endhead
\bottomrule\noalign{}
\endlastfoot
UV dimers (CPD, 6-4PP) & Nucleotide excision repair (NER) & XPA, RPA,
XPD, XPB, XPF-ERCC1, XPG, PCNA, Pol delta & Xeroderma pigmentosum
(1,000x skin cancer risk) \\
Oxidized bases (8-oxoG) & Base excision repair (BER) & OGG1 glycosylase,
APE1, Pol beta, XRCC1, Ligase III & MUTYH-associated polyposis \\
Alkylated bases (O\(^6\)-meG) & Direct reversal & MGMT
(methyltransferase) & Glioblastoma (MGMT
\hyperref[gl:promoter]{\textbf{promoter}} methylation predicts
temozolomide response) \\
Replication mismatch & Mismatch repair (MMR) & MLH1, MSH2, MSH6, PMS2 &
Lynch syndrome (colorectal, endometrial cancer) \\
Double-strand break & NHEJ & Ku70/80, DNA-PKcs, Artemis, Ligase IV &
SCID (severe combined immunodeficiency) \\
Double-strand break & Homologous
\hyperref[gl:recombination]{\textbf{recombination}} (HR) & ATM, MRN,
BRCA1, BRCA2, RAD51, RPA & Hereditary breast/ovarian cancer \\
Interstrand crosslinks & Fanconi anemia pathway & FANC proteins (A-W),
then NER + HR & Fanconi anemia (bone marrow failure, cancer) \\
\end{longtable}
}

\begin{quote}
\textbf{Clinical Connection: MGMT Methylation and Glioblastoma}
Glioblastoma is the most aggressive brain tumor (median survival
about 15 months). The alkylating agent \textbf{temozolomide} is
a standard chemotherapy. Tumors with \textbf{MGMT promoter methylation}
(silencing MGMT expression) cannot repair temozolomide-induced
O\(^6\)-methylguanine lesions, making them more sensitive to the drug.
MGMT promoter methylation status is therefore a critical predictive
biomarker: methylated tumors have about 21-month median
survival vs about 12 months for unmethylated tumors when
treated with temozolomide (Hegi et al., 2005, \emph{NEJM}).
\end{quote}

\subsubsection{Base Excision Repair (BER) --- Mechanism in
Detail}\label{base-excision-repair-ber-mechanism-in-detail}

BER handles oxidised, deaminated, and alkylated single-base lesions. The
pathway has two main subdivisions:

\textbf{Short-patch BER (single-nucleotide replacement,
about 80 \% of events):} 1. \textbf{DNA glycosylase} recognises
the lesion and cleaves the N-glycosidic bond, producing an abasic (AP)
site. There are 11 DNA glycosylases in humans, each with substrate
specificity:

{\def\LTcaptype{none} % do not increment counter
\begin{longtable}[]{@{}
  >{\raggedright\arraybackslash}p{(\linewidth - 4\tabcolsep) * \real{0.3333}}
  >{\raggedright\arraybackslash}p{(\linewidth - 4\tabcolsep) * \real{0.5152}}
  >{\raggedright\arraybackslash}p{(\linewidth - 4\tabcolsep) * \real{0.1515}}@{}}
\toprule\noalign{}
\begin{minipage}[b]{\linewidth}\raggedright
Glycosylase
\end{minipage} & \begin{minipage}[b]{\linewidth}\raggedright
Lesion recognised
\end{minipage} & \begin{minipage}[b]{\linewidth}\raggedright
Notes
\end{minipage} \\
\midrule\noalign{}
\endhead
\bottomrule\noalign{}
\endlastfoot
\textbf{OGG1} & 8-oxoguanine paired with C & Most clinically important;
oxidative damage; bifunctional (lyase activity creates SSB) \\
\textbf{MUTYH} & Adenine misincorporated opposite 8-oxoG
(post-replication mismatch) & Loss-of-function \texorpdfstring{\ensuremath{\to}}{->} MUTYH-associated
polyposis (MAP) \\
\textbf{NTH1} & Oxidised pyrimidines (thymine glycol, dihydrouracil) &
Bifunctional \\
\textbf{NEIL1, NEIL2, NEIL3} & Ring-fragmented purines (FapyG, FapyA),
oxidised pyrimidines & Bifunctional; \texorpdfstring{\ensuremath{\beta}}{beta},\texorpdfstring{\ensuremath{\delta}}{delta}-elimination produces
3\texorpdfstring{'}{'}-phosphate end \\
\textbf{MPG (AAG)} & 3-methyladenine, 7-methylguanine, ethenoadenine &
Monofunctional (no lyase) \\
\textbf{UNG (uracil DNA glycosylase)} & Uracil (deamination product or U
mis-incorporation) & Most active; nuclear and mitochondrial isoforms \\
\textbf{TDG} & Thymine in T:G mismatch (5mC deamination); 5fC/5caC (TET
pathway) & Active in DNA demethylation \\
\textbf{SMUG1} & Uracil; backup to UNG & Monofunctional \\
\textbf{MBD4} & Thymine in T:G mismatch at CpG sites & Methylation-CpG
specific \\
\textbf{MYH (MUTYH)} & A from A:8-oxoG mispair & See above \\
\end{longtable}
}

\begin{enumerate}
\def\labelenumi{\arabic{enumi}.}
\setcounter{enumi}{1}
\tightlist
\item
  \textbf{AP endonuclease (APE1, also called APE)} cuts the DNA backbone
  5\texorpdfstring{'}{'} to the AP site, leaving a 3\texorpdfstring{'}{'}-OH and a 5\texorpdfstring{'}{'}-deoxyribose phosphate
  (dRP) terminus.
\item
  \textbf{DNA polymerase \texorpdfstring{\ensuremath{\beta}}{beta} (pol \texorpdfstring{\ensuremath{\beta}}{beta})} has both DNA polymerase activity
  (fills the 1-nt gap with the correct nucleotide) and \textbf{dRP lyase
  activity} (removes the 5\texorpdfstring{'}{'}-dRP).
\item
  \textbf{DNA ligase III + XRCC1} seals the resulting nick
  (ATP-dependent).
\end{enumerate}

\textbf{Long-patch BER (2--10 nucleotide replacement, about 20
\% of events):} - Used when the dRP cannot be removed (e.g., oxidised
dRP) or when the lesion creates a strand-displacement intermediate. -
DNA pol \texorpdfstring{\ensuremath{\delta}}{delta}/\texorpdfstring{\ensuremath{\varepsilon}}{epsilon} extends, displacing the 5\texorpdfstring{'}{'} flap; \textbf{FEN1} cleaves the
flap; \textbf{DNA ligase I} seals.

\textbf{OGG1 mechanism specifically} (the most studied glycosylase): -
OGG1 scans DNA for 8-oxoguanine. The 8-oxoG flips out of the helix into
OGG1's catalytic pocket. - Active-site lysine (K249) acts as
nucleophile, attacking the C1\texorpdfstring{'}{'} of the deoxyribose; the N-glycosidic bond
breaks via a Schiff-base intermediate. - A second active-site residue
(D268) hydrolyses the intermediate, releasing the 8-oxoG base and
leaving an AP site. - OGG1 is bifunctional: a \texorpdfstring{\ensuremath{\beta}}{beta}-elimination at the 3\texorpdfstring{'}{'}
phosphate produces a 3\texorpdfstring{'}{'}-\texorpdfstring{\ensuremath{\alpha}}{alpha},\texorpdfstring{\ensuremath{\beta}}{beta}-unsaturated aldehyde, which APE1 must
process before pol \texorpdfstring{\ensuremath{\beta}}{beta} can extend. - 8-oxoG is mutagenic because it pairs
with adenine (Hoogsteen face) instead of cytosine --- leading to G:C \texorpdfstring{\ensuremath{\to}}{->}
T:A transversions if not repaired.

\textbf{Clinical link --- MUTYH-associated polyposis (MAP):} Biallelic
loss-of-function mutations in MUTYH cause autosomal recessive colorectal
polyposis. The mechanism: OGG1 normally removes 8-oxoG, but if 8-oxoG
escapes repair and is replicated, it pairs with A. MUTYH normally
removes that A; without MUTYH, the A becomes fixed as a G:C \texorpdfstring{\ensuremath{\to}}{->} T:A
transversion. MAP tumours have a characteristic mutational signature
(transversions at GAA contexts).

\subsubsection{Comparing NER and MMR --- A Side-by-Side
Table}\label{comparing-ner-and-mmr-a-side-by-side-table}

The pathways have distinct substrates and proteins, often confused. Here
is a direct comparison:

{\def\LTcaptype{none} % do not increment counter
\begin{longtable}[]{@{}
  >{\raggedright\arraybackslash}p{(\linewidth - 4\tabcolsep) * \real{0.1148}}
  >{\raggedright\arraybackslash}p{(\linewidth - 4\tabcolsep) * \real{0.5410}}
  >{\raggedright\arraybackslash}p{(\linewidth - 4\tabcolsep) * \real{0.3443}}@{}}
\toprule\noalign{}
\begin{minipage}[b]{\linewidth}\raggedright
Feature
\end{minipage} & \begin{minipage}[b]{\linewidth}\raggedright
Nucleotide Excision Repair (NER)
\end{minipage} & \begin{minipage}[b]{\linewidth}\raggedright
Mismatch Repair (MMR)
\end{minipage} \\
\midrule\noalign{}
\endhead
\bottomrule\noalign{}
\endlastfoot
\textbf{Substrate} & Bulky helix-distorting lesions: CPDs, 6-4
photoproducts, cisplatin adducts, polycyclic aromatic hydrocarbons &
Mismatched bases (G-T, A-C); small (1--13 nt) insertion/deletion
loops \\
\textbf{Recognition} & Helix distortion or transcription stalling &
Mismatch geometry directly bound by MutS\texorpdfstring{\ensuremath{\alpha}}{alpha}/\texorpdfstring{\ensuremath{\beta}}{beta} \\
\textbf{Initiation timing} & Any cell-cycle phase & Post-replicative
(S/G2) \\
\textbf{Strand discrimination} & Both strands repaired & New strand
identified by transient nicks (5\texorpdfstring{'}{'} direction in eukaryotes; HMGB1) \\
\textbf{Primary recognition factor} & XPC-RAD23B (GG-NER); CSA-CSB
(TC-NER, transcription-coupled) & MSH2-MSH6 (MutS\texorpdfstring{\ensuremath{\alpha}}{alpha}; mismatches and 1-nt
loops); MSH2-MSH3 (MutS\texorpdfstring{\ensuremath{\beta}}{beta}; 2--13-nt loops) \\
\textbf{Verification factor} & TFIIH (XPB-helicase 3\texorpdfstring{'}{'}\texorpdfstring{\ensuremath{\to}}{->}5\texorpdfstring{'}{'}; XPD-helicase
5\texorpdfstring{'}{'}\texorpdfstring{\ensuremath{\to}}{->}3\texorpdfstring{'}{'}) --- unwinds about 25 bp & MLH1-PMS2 (MutL\texorpdfstring{\ensuremath{\alpha}}{alpha};
endonuclease) makes incision \\
\textbf{Excision} & XPF-ERCC1 (5\texorpdfstring{'}{'} cut), XPG (3\texorpdfstring{'}{'} cut); 24--32 nt fragment
removed & Exonuclease 1 (Exo1) degrades from MutL\texorpdfstring{\ensuremath{\alpha}}{alpha} cut to past mismatch;
about 150 nt fragment removed \\
\textbf{Synthesis} & Pol \texorpdfstring{\ensuremath{\delta}}{delta}/\texorpdfstring{\ensuremath{\varepsilon}}{epsilon} with PCNA, RFC & Pol \texorpdfstring{\ensuremath{\delta}}{delta} \\
\textbf{Ligation} & DNA ligase I & DNA ligase I \\
\textbf{Disease (loss-of-function)} & Xeroderma pigmentosum (XPA--XPG,
about 1,000\texorpdfstring{\ensuremath{\times}}{x} skin cancer); Cockayne syndrome (CSA, CSB;
premature aging) & Lynch syndrome / HNPCC (MLH1, MSH2, MSH6, PMS2;
50--80\% lifetime CRC risk; MSI-high tumours) \\
\textbf{Therapeutic relevance} & Cisplatin sensitivity correlates with
NER capacity & MSI-high tumours respond to anti-PD-1 (pembrolizumab) due
to high neoantigen load \\
\end{longtable}
}

\subsubsection{HR vs NHEJ --- Pathway Choice
Mechanism}\label{hr-vs-nhej-pathway-choice-mechanism}

The choice between homologous recombination (HR) and non-homologous end
joining (NHEJ) is regulated by cell-cycle phase and a competition at the
broken DNA ends.

\textbf{Pathway-determining factors:}

{\def\LTcaptype{none} % do not increment counter
\begin{longtable}[]{@{}
  >{\raggedright\arraybackslash}p{(\linewidth - 4\tabcolsep) * \real{0.1333}}
  >{\raggedright\arraybackslash}p{(\linewidth - 4\tabcolsep) * \real{0.4667}}
  >{\raggedright\arraybackslash}p{(\linewidth - 4\tabcolsep) * \real{0.4000}}@{}}
\toprule\noalign{}
\begin{minipage}[b]{\linewidth}\raggedright
Step
\end{minipage} & \begin{minipage}[b]{\linewidth}\raggedright
NHEJ-favouring
\end{minipage} & \begin{minipage}[b]{\linewidth}\raggedright
HR-favouring
\end{minipage} \\
\midrule\noalign{}
\endhead
\bottomrule\noalign{}
\endlastfoot
Cell cycle phase & G1, G0 & S, G2 (sister chromatid available as
template) \\
End sensor & \textbf{Ku70/Ku80} binds DSB ends rapidly & Slow
accommodation; replaces Ku \\
End processing & \textbf{DNA-PKcs} recruited; minimal end resection &
\textbf{MRN complex} (MRE11--RAD50--NBN) initiates short-range
resection \\
End-resection regulator & \textbf{53BP1} + \textbf{RIF1} + \textbf{PTIP}
+ \textbf{shieldin} complex (REV7-MAD2L2-SHLD1/2/3) --- protects ends
from resection & \textbf{CtIP} (CDK-phosphorylated in S/G2) +
\textbf{MRN} initiates resection; \textbf{EXO1} + \textbf{DNA2} +
\textbf{BLM} extend resection \\
Fidelity & Error-prone --- small deletions/insertions at junction &
High-fidelity (uses sister chromatid template) \\
Time scale & Fast (minutes) & Slow (hours) \\
Key proteins (after end-recognition) & DNA-PKcs, Artemis (nuclease),
XRCC4, Ligase IV, XLF/Cernunnos, PAXX & RPA \texorpdfstring{\ensuremath{\to}}{->} BRCA2 \texorpdfstring{\ensuremath{\to}}{->} RAD51 (filament) \texorpdfstring{\ensuremath{\to}}{->}
strand invasion \texorpdfstring{\ensuremath{\to}}{->} D-loop \texorpdfstring{\ensuremath{\to}}{->} resolution \\
Outcome & Direct ligation of ends (often with deletions \texorpdfstring{\ensuremath{\leq}}{<=} 50 bp) & Use
of sister chromatid as template; restores original sequence \\
\end{longtable}
}

\textbf{The CDK regulatory switch:} During G1, CDK activity is low; CtIP
is unphosphorylated; 53BP1-shieldin protects ends; Ku70/80 dominates \texorpdfstring{\ensuremath{\to}}{->}
NHEJ. During S/G2, CDK1/2 phosphorylates CtIP at Ser327 and Thr847,
activating its end-resection function; CtIP+MRN displaces Ku \texorpdfstring{\ensuremath{\to}}{->} HR is
initiated.

\textbf{The BRCA1 vs 53BP1 antagonism.} \textbf{BRCA1} promotes
resection (HR); \textbf{53BP1} blocks resection (NHEJ). The two compete
for binding at H4K20me-marked nucleosomes near the DSB. BRCA1's BRCT
domain reads \texorpdfstring{\ensuremath{\gamma}}{gamma}H2AX-bound MDC1; 53BP1's tandem Tudor domain reads
H4K20me2. Loss of 53BP1 in BRCA1-deficient cells \textbf{rescues HR}
(because resection is no longer blocked) --- a synthetic-lethal escape
mechanism in PARP-inhibitor resistance.

\textbf{The BRCA2 mechanism.} Once ssDNA is generated by resection, RPA
coats it. BRCA2 then loads RAD51 onto the RPA-coated ssDNA, displacing
RPA. This is the rate-limiting step of HR. Loss of BRCA2 prevents RAD51
loading \texorpdfstring{\ensuremath{\to}}{->} HR fails \texorpdfstring{\ensuremath{\to}}{->} cells become ``BRCAness''-positive \texorpdfstring{\ensuremath{\to}}{->} exquisitely
sensitive to PARP inhibitors (synthetic lethality, see
\cref{sec:unit_IV_dna_replication_and_cell_cycle}).

\textbf{Ku70/80 binding.} Ku70/80 forms a ring that threads onto the DSB
end with extreme avidity (Kd \texorpdfstring{\ensuremath{\approx}}{~} 1 nM). It positions the DNA-PKcs
catalytic subunit at the end and protects against resection. The Ku70/80
ring must be removed (by ubiquitination + proteasome) before HR can
proceed in S/G2. This explains why NHEJ has an immediate kinetic
advantage and HR requires time and active disassembly.

\textbf{Concept Check 13.2}

\begin{quote}
\begin{enumerate}
\def\labelenumi{\arabic{enumi}.}
\tightlist
\item
  Why is the Ames test performed with and without liver microsomal
  extract (S9)?
\item
  Explain why MGMT is called a ``suicide
  \hyperref[gl:enzyme]{\textbf{enzyme}}.''
\item
  A patient has mutations in both BRCA1
  \hyperref[gl:allele]{\textbf{allele}}s. Which repair pathway is
  defective? Why does this predispose to breast/ovarian cancer rather
  than most cancer types?
\item
  Compare the fidelity of NHEJ and HR. When would each be preferred?
\end{enumerate}
\end{quote}

\begin{center}\rule{0.5\linewidth}{0.5pt}\end{center}

\subsection{CRISPR-Cas9: Mechanism and
Applications}\label{crispr-cas9-mechanism-and-applications}

\textbf{Background}: CRISPR (Clustered Regularly Interspaced Short
Palindromic Repeats) is the adaptive immune system of bacteria and
\hyperref[gl:archaea]{\textbf{archaea}}. Jennifer Doudna and Emmanuelle
Charpentier received the 2020 Nobel Prize in Chemistry for demonstrating
that CRISPR-Cas9 could be programmed with a synthetic guide RNA for
precise genome editing (Jinek et al., 2012, \emph{Science}).

\subsubsection{Natural CRISPR Immunity}\label{natural-crispr-immunity}

In bacteria, CRISPR functions as immunological memory:

\begin{enumerate}
\def\labelenumi{\arabic{enumi}.}
\tightlist
\item
  \textbf{Adaptation}: When a bacterium survives phage infection, a
  short sequence from the phage genome (the \textbf{protospacer}) is
  integrated into the CRISPR array as a new \textbf{spacer} between
  repeat sequences
\item
  \textbf{Expression}: The CRISPR array is transcribed into pre-crRNA,
  which is processed (with tracrRNA in Type II systems) into mature
  crRNA
\item
  \textbf{Interference}: The crRNA guides Cas9 to complementary foreign
  DNA; Cas9 cleaves it, preventing subsequent infection
\end{enumerate}

\subsubsection{Engineered CRISPR-Cas9}\label{engineered-crispr-cas9}

\begin{figure}[htbp]
\centering
\includegraphics[width=0.82\linewidth,height=4.2in,keepaspectratio]{../figures/mermaid_inline/inline_mermaid_0076_2b34ba77d6a1.png}
\caption{The CRISPR-Cas9 genome editing mechanism. The sgRNA guides Cas9 to a PAM-adjacent target sequence where a DSB is created. Repair by NHEJ produces knockouts; HDR with a donor template allows precise editing. Newer variants (base editors, prime editors) avoid DSBs entirely.}
\end{figure}

\textbf{Key molecular details}: - \textbf{sgRNA (single guide RNA)}: A
fusion of crRNA (20-nt spacer complementary to the target) and tracrRNA
(scaffold that binds Cas9). The about 100-nt sgRNA is most that
is needed to redirect Cas9. - \textbf{PAM (protospacer adjacent motif)}:
For \emph{S. pyogenes} Cas9 (SpCas9), the PAM is \textbf{5'-NGG-3'} on
the non-target strand, immediately 3' of the target sequence. PAM
recognition is the first step; without PAM, Cas9 will not unwind or
cleave. - \textbf{Cleavage}: \textbf{HNH domain} cleaves the strand
complementary to the sgRNA (target strand); \textbf{RuvC domain} cleaves
the non-target strand. Both cuts occur about 3 bp upstream of
the PAM, creating a blunt-ended DSB.

\subsubsection{Beyond Cas9: Expanded
Toolkit}\label{beyond-cas9-expanded-toolkit}

{\def\LTcaptype{none} % do not increment counter
\begin{longtable}[]{@{}
  >{\raggedright\arraybackslash}p{(\linewidth - 4\tabcolsep) * \real{0.2000}}
  >{\raggedright\arraybackslash}p{(\linewidth - 4\tabcolsep) * \real{0.3667}}
  >{\raggedright\arraybackslash}p{(\linewidth - 4\tabcolsep) * \real{0.4333}}@{}}
\toprule\noalign{}
\begin{minipage}[b]{\linewidth}\raggedright
Tool
\end{minipage} & \begin{minipage}[b]{\linewidth}\raggedright
Mechanism
\end{minipage} & \begin{minipage}[b]{\linewidth}\raggedright
Application
\end{minipage} \\
\midrule\noalign{}
\endhead
\bottomrule\noalign{}
\endlastfoot
\textbf{Cas12a (Cpf1)} & Creates staggered DSB; T-rich PAM (TTTN);
self-processes crRNA array & Multiplexed editing; organisms with AT-rich
genomes \\
\textbf{Cas13 (RNA-targeting)} & Cleaves RNA (not DNA); collateral RNA
cleavage in bacteria & RNA knockdown; viral RNA detection (SHERLOCK
diagnostic) \\
\textbf{dCas9 (dead Cas9)} & Catalytically inactive (D10A + H840A
mutations); binds but does not cut & CRISPRi (transcriptional repression
with KRAB domain); CRISPRa (activation with VP64/p65/Rta) \\
\textbf{Base editors} & dCas9 or Cas9-nickase fused to deaminase & CBE
(cytidine base editor): C-to-T conversion; ABE (adenine base editor):
A-to-G conversion; no DSB needed \\
\textbf{Prime editors} & Cas9-H840A nickase fused to M-MLV reverse
transcriptase + pegRNA & ``Search-and-replace'': any point mutation,
small insertion, or small deletion; no DSB, no donor template needed \\
\end{longtable}
}

\textbf{CRISPR diagnostics:} SHERLOCK and DETECTR exploit collateral
nuclease activity after target recognition: Cas13 collateral RNA
cleavage and Cas12a collateral ssDNA cleavage convert a sequence match
into a fluorescent or lateral-flow signal
\citep{gootenberg2017sherlock,chen2018detectr}. That makes CRISPR a
diagnostic logic as well as an editing logic, but it does not remove the
usual validation burdens: sample preparation, contamination control,
limit of detection, clinical sensitivity/specificity, and variant
coverage still determine whether an assay is deployable.

\textbf{Current evidence: prime editing as promoter rewriting.} Prime
editing is especially useful when the desired change is not simply ``cut
here'' but ``install a small regulatory pattern.'' A 2025 \emph{Blood}
study used prime-editing guide RNAs to introduce multiple
hereditary-persistence-of-fetal-haemoglobin-like edits into the
\emph{HBG1/HBG2} gamma-globin promoters, aiming to reactivate fetal
haemoglobin without making a double-strand break
\citep{chalumeau2025primeediting}. Editing efficiency varied among donor
haematopoietic stem/progenitor cells, so this should be treated as early
translational evidence rather than an approved therapy. The concept is
powerful: regulatory DNA can be edited to change \emph{when and where} a
gene is expressed, not only whether its protein-coding sequence is
intact.

\subsubsection{Therapeutic Applications
(2024-2025)}\label{therapeutic-applications-2024-2025}

\textbf{Casgevy (exagamglogene autotemcel)}: First CRISPR-based therapy
approved by FDA for sickle cell disease (December 2023) and later for
transfusion-dependent beta-thalassaemia (January 2024), with EMA
authorisation following in 2024. FDA's 2026 product page lists exa-cel
for sickle cell disease and transfusion-dependent beta-thalassaemia in
patients 12 years of age and older, with updated approval documentation
\citep{fda2023casgevy,fda2024casgevythalassemia,fda2026casgevy}.

\begin{itemize}
\tightlist
\item
  \textbf{Indication}: Sickle cell disease (SCD) and
  transfusion-dependent beta-thalassemia
\item
  \textbf{Mechanism}: Ex vivo CRISPR editing of patient hematopoietic
  stem cells (HSCs) to disrupt the \emph{BCL11A} erythroid enhancer.
  BCL11A is a repressor of fetal hemoglobin (HbF, alpha2-gamma2).
  Disrupting its erythroid enhancer reactivates HbF production, which
  does not sickle and compensates for defective HbS.
\item
  \textbf{Results}: In SCD, the regulatory endpoint is freedom from
  severe vaso-occlusive crises after treatment; transfusion independence
  is the endpoint used for transfusion-dependent beta-thalassaemia.
  Early exa-cel studies also reported sustained fetal-haemoglobin
  induction, so the durable teaching point is endpoint specificity: the
  same edited cell product is judged against different disease-specific
  clinical outcomes
  \citep{fda2023casgevy,fda2024casgevythalassemia,fda2026casgevy}.
\end{itemize}

\textbf{NTLA-2001 (in vivo CRISPR)}: - \textbf{Indication}:
Transthyretin (ATTR) amyloidosis - \textbf{Mechanism}: Lipid
nanoparticle (LNP)-delivered Cas9 mRNA + sgRNA targeting the \emph{TTR}
gene in hepatocytes -- direct in vivo editing (no cell harvest needed) -
\textbf{Results}: Single infusion achieved 80-95\% serum TTR reduction
(Gillmore et al., 2021, \emph{NEJM})

\textbf{Safety considerations}: - \textbf{Off-target editing}: sgRNA may
bind sites with 1-3 mismatches. Mitigated by high-fidelity Cas9 variants
(SpCas9-HF1, eSpCas9, HiFi Cas9) and careful guide RNA design
(off-target prediction algorithms: Cas-OFFinder, CRISPOR) -
\textbf{Immunogenicity}: about 79\% of humans have pre-existing
antibodies to SpCas9; strategies include using Cas9 from other species
or transient delivery (mRNA/RNP vs.~viral vectors) - \textbf{NHEJ:HDR
ratio}: about 1000:1 in non-dividing cells (most
therapeutically relevant cells); S/G2 phase favors HDR; base editors and
prime editors circumvent this limitation

\begin{quote}
\textbf{Clinical Connection: Base Editing for Sickle Cell Disease} An
alternative approach uses adenine base editing to directly convert the
sickle mutation (GAG to GTG = Glu to Val) back to the wild-type or a
benign Makassar variant (GAG to GCG = Glu to Ala). This avoids DSBs
entirely and may be safer than Cas9-based approaches. Clinical trials
are underway (Beam Therapeutics, BEAM-101).
\end{quote}

\textbf{Mitochondrial DNA editing boundary:} Standard CRISPR-Cas systems
work poorly for mtDNA because guide RNAs are difficult to import into
mitochondria. DddA-derived cytosine base editors (DdCBEs) solve part of
the problem by using TALE-DddA protein fusions to install targeted
C\(\cdot\)G-to-T\(\cdot\)A edits in mtDNA without guide RNA
\citep{mok2020mitochondrialbaseediting}. This is a major tool advance,
not a general approved therapy: heteroplasmy, off-target editing,
delivery, selectable mutation classes, and tissue specificity remain the
current translational constraints
\citep{barrerapaez2023mitochondrialediting}.

\textbf{Concept Check 13.3}

\begin{quote}
\begin{enumerate}
\def\labelenumi{\arabic{enumi}.}
\tightlist
\item
  Why does Cas9 require a PAM sequence? What is the evolutionary
  rationale?
\item
  Compare NHEJ and HDR outcomes after Cas9-mediated DSB. Why is HDR
  difficult to achieve in post-mitotic cells?
\item
  What advantage does base editing offer over standard Cas9 editing for
  correcting point mutations?
\item
  Explain the mechanism of Casgevy: why does disrupting BCL11A treat
  sickle cell disease?
\end{enumerate}
\end{quote}

\begin{center}\rule{0.5\linewidth}{0.5pt}\end{center}

\subsection{Genomics: The Human Genome}\label{genomics-the-human-genome}

The \textbf{Human Genome Project} (1990-2003; Venter et al., 2001 and
International Human Genome Sequencing Consortium, 2001) produced the
first draft reference. The \textbf{T2T Consortium} (Nurk et al., 2022)
completed the first truly gapless reference genome (CHM13), filling
about 8\% of previously unresolvable sequence (centromeres,
\hyperref[gl:telomere]{\textbf{telomere}}s, segmental duplications, rDNA
arrays).

\subsubsection{Human Genome Composition}\label{human-genome-composition}

{\def\LTcaptype{none} % do not increment counter
\begin{longtable}[]{@{}
  >{\raggedright\arraybackslash}p{(\linewidth - 4\tabcolsep) * \real{0.3667}}
  >{\raggedright\arraybackslash}p{(\linewidth - 4\tabcolsep) * \real{0.3333}}
  >{\raggedright\arraybackslash}p{(\linewidth - 4\tabcolsep) * \real{0.3000}}@{}}
\toprule\noalign{}
\begin{minipage}[b]{\linewidth}\raggedright
Component
\end{minipage} & \begin{minipage}[b]{\linewidth}\raggedright
Fraction
\end{minipage} & \begin{minipage}[b]{\linewidth}\raggedright
Details
\end{minipage} \\
\midrule\noalign{}
\endhead
\bottomrule\noalign{}
\endlastfoot
Protein-coding exons & about 1.5\% &
about 19,000-20,000 genes; average gene about 27 kb;
average of 8.8 exons \\
Introns & about 26\% & Average intron about 3.4 kb;
some \textgreater100 kb \\
Transposable elements & about 46\% & LINEs (21\%, especially
L1), SINEs (13\%, especially Alu), DNA transposons (3\%), LTR
retrotransposons (9\%) \\
Intergenic sequence & about 15\% & Includes regulatory
elements, structural sequences \\
Non-coding RNA genes & about 5,000+ loci & lncRNAs, miRNAs,
snRNAs, snoRNAs, rRNAs, tRNAs \\
Simple sequence repeats & about 3\% & Microsatellites (1-6 bp
repeat units) and minisatellites (7-100 bp) \\
Segmental duplications & about 5\% & Blocks \textgreater1 kb
with \textgreater90\% sequence identity; hotspots for chromosomal
rearrangement \\
\end{longtable}
}

\subsubsection{Transposable Elements --- Mechanism and Genomic
Impact}\label{transposable-elements-mechanism-and-genomic-impact}

Transposable elements (TEs) make up about 46 \% of the human
genome --- far more than the about 1.5 \% occupied by
protein-coding exons. They are ancient mobile DNA elements that have
shaped genome architecture over evolutionary time. They fall into three
major classes:

\textbf{Class I --- RNA-mediated retrotransposons (copy-and-paste;
majority of human TEs):}

\begin{enumerate}
\def\labelenumi{\arabic{enumi}.}
\item
  \textbf{LINE-1 (Long Interspersed Nuclear Element 1)} ---
  about 21 \% of human genome; about 500,000 copies;
  about 6 kb full-length. Primarily about 80--100 are
  currently transposition-competent.

  \begin{itemize}
  \tightlist
  \item
    Encodes ORF1 (RNA-binding chaperone) and ORF2 (with
    reverse-transcriptase + endonuclease activities)
  \item
    \textbf{Retrotransposition mechanism (Target-Primed Reverse
    Transcription, TPRT):}

    \begin{enumerate}
    \def\labelenumii{\arabic{enumii}.}
    \tightlist
    \item
      L1 transcribed by Pol II \texorpdfstring{\ensuremath{\to}}{->} L1 mRNA (about 6 kb, capped,
      polyadenylated)
    \item
      L1 mRNA exported to cytoplasm, translated to ORF1p and ORF2p
    \item
      Cis-binding of ORF1p/ORF2p to their own mRNA forms a
      ribonucleoprotein (RNP)
    \item
      RNP enters nucleus
    \item
      ORF2p endonuclease nicks genomic DNA at TTAAAA-like sites,
      exposing 3\texorpdfstring{'}{'}-OH
    \item
      The 3\texorpdfstring{'}{'}-OH primes reverse transcription using L1 mRNA as template
    \item
      Second-strand DNA synthesis completes integration
    \item
      Target site duplication (TSD; 7--20 bp) flanks the new insertion,
      plus a poly-A tail at the 3\texorpdfstring{'}{'} end
    \end{enumerate}
  \item
    Most insertions are 5\texorpdfstring{'}{'}-truncated due to incomplete RT
  \item
    L1 endonuclease nicks about 100 sites/cell per day; most
    fail to complete TPRT (silenced by KRAB-ZFP-driven H3K9me3
    repression)
  \end{itemize}
\item
  \textbf{Alu elements (Short Interspersed Nuclear Element 1, SINE)} ---
  about 13 \% of human genome; about 1.1 million
  copies; about 300 bp.

  \begin{itemize}
  \tightlist
  \item
    Derived from 7SL RNA (signal recognition particle component);
    contain RNA Pol III internal promoter
  \item
    \textbf{Mobilisation requires L1 machinery} (parasitic on LINE-1's
    ORF2p RT/endonuclease)
  \item
    Transcription by Pol III; the Alu RNA recruits L1 ORF2p in
    \emph{trans} and inserts via the same TPRT mechanism
  \item
    Alu elements are densest in gene-rich regions; insertions in introns
    can disrupt splicing or create new alternative exons
  \end{itemize}
\item
  \textbf{SVA elements (SINE-VNTR-Alu)} --- about 3,000 copies;
  about 2 kb; younger than Alu (about 25 Mya). Hybrid
  composite.
\item
  \textbf{LTR retrotransposons (HERV --- human endogenous retroviruses)}
  --- about 9 \% of genome; relics of ancient retroviral
  insertions. Most are inactivated; some HERV-K elements remain intact
  and may be re-activated in cancer.
\end{enumerate}

\textbf{Class II --- DNA transposons (cut-and-paste; about 3 \%
of human genome):} - Most inactivated in modern human genome (no active
class II TE in humans) - Mechanism: a transposase enzyme (encoded by the
TE) recognises terminal inverted repeats (TIRs) flanking the element,
excises the element, and re-integrates it elsewhere - Mariner / Tc1,
hAT, P-element families most derive from class II - Hyperactive
synthetic DNA transposons (e.g., \textbf{Sleeping Beauty},
\textbf{PiggyBac}) are now used as gene-delivery tools in research and
CAR-T therapy

\textbf{Genomic and clinical impact of TE insertions:} -
about 100 \emph{de novo} L1 or Alu insertions per generation
(germline) - about 10--15 \% of disease-causing mutations are
TE insertions (e.g., Alu insertion in \emph{NF1} exon 6 causing
neurofibromatosis; L1 insertion in \emph{DMD} causing Duchenne) - TE
insertions in gene regulatory regions can create new enhancers or
promoters (about 25 \% of mammalian enhancers derive from TEs)
- Active L1 retrotransposition in the brain creates somatic mosaicism
(about 80 unique L1 insertions per neuron) potentially
contributing to neuronal diversity - Cancer cells often re-activate TE
expression (LINE-1 demethylation is a hallmark of global cancer
hypomethylation), contributing to chromosomal instability and
inflammation

\textbf{Defence mechanisms against TE activity:} - \textbf{DNA
methylation} of TE promoters (DNMT3A/B targets active TE families during
germline development) - \textbf{KRAB-ZFP transcriptional repression}:
\textgreater{} 350 KRAB-zinc finger proteins recognise specific TE
families and recruit KAP1/TRIM28 \texorpdfstring{\ensuremath{\to}}{->} SETDB1 \texorpdfstring{\ensuremath{\to}}{->} H3K9me3 silencing -
\textbf{piRNA pathway} in germline: piRNAs (26--32 nt) bind PIWI clade
Argonautes (PIWIL1/MILI, PIWIL4/MIWI2) and silence TEs via the
``ping-pong amplification'' cycle - \textbf{APOBEC3 deaminases}:
cytidine deaminases that mutate TE-derived ssDNA, restricting
retrotransposition and HIV integration

\subsubsection{Structural Variants --- Inversions, CNVs, and
Disease}\label{structural-variants-inversions-cnvs-and-disease}

\textbf{Inversions} are large segments of DNA flipped end-to-end. Unlike
SNPs, they are difficult to detect by short-read sequencing (the
inversion preserves sequence content, primarily orientation changes).
Long-read sequencing (PacBio HiFi, Oxford Nanopore) and optical mapping
(Bionano Saphyr) have revealed: - about 40,000 polymorphic
inversions per human genome (most \textless{} 50 kb) - \textbf{HbF
inversion regulators}: A rare inversion at the \emph{BCL11A} locus
disrupts the erythroid-specific enhancer of \emph{BCL11A}, derepressing
fetal haemoglobin (HbF) production. This explains the high HbF
persistence in some Sardinian families and provides the rationale for
therapeutic disruption of \emph{BCL11A} (Casgevy approved for sickle
cell disease). - \textbf{Pericentric inversion of chromosome 9}
(inv(9)(p11q12)): a common normal variant (about 1--3 \% of
population), classically considered benign, but recently associated with
subtle effects on fertility.

\textbf{Copy Number Variants (CNVs)} are duplications or deletions of
segments \textgreater{} 1 kb. Two major mechanisms: -
\textbf{Non-allelic homologous recombination (NAHR)} between segmental
duplications (low-copy repeats, LCRs) --- produces recurrent CNVs at
predictable breakpoints (e.g., 22q11.2 microdeletion, 16p11.2 micro-CNV)
- \textbf{Replication-based mechanisms (FoSTeS, MMBIR)} --- produce
non-recurrent CNVs at variable breakpoints

\textbf{CNVs in autism and schizophrenia:}

{\def\LTcaptype{none} % do not increment counter
\begin{longtable}[]{@{}
  >{\raggedright\arraybackslash}p{(\linewidth - 6\tabcolsep) * \real{0.1064}}
  >{\raggedright\arraybackslash}p{(\linewidth - 6\tabcolsep) * \real{0.0851}}
  >{\raggedright\arraybackslash}p{(\linewidth - 6\tabcolsep) * \real{0.1915}}
  >{\raggedright\arraybackslash}p{(\linewidth - 6\tabcolsep) * \real{0.6170}}@{}}
\toprule\noalign{}
\begin{minipage}[b]{\linewidth}\raggedright
Locus
\end{minipage} & \begin{minipage}[b]{\linewidth}\raggedright
Size
\end{minipage} & \begin{minipage}[b]{\linewidth}\raggedright
Phenotype
\end{minipage} & \begin{minipage}[b]{\linewidth}\raggedright
Frequency in cases vs controls
\end{minipage} \\
\midrule\noalign{}
\endhead
\bottomrule\noalign{}
\endlastfoot
\textbf{16p11.2} (deletion or duplication) & about 600 kb &
Autism, intellectual disability, obesity (deletion) & about 1
\% of autism cases vs 0.04 \% controls \\
\textbf{22q11.2} (deletion = DiGeorge syndrome / 22q11DS) &
about 3 Mb & Schizophrenia (about 25 \% of carriers),
congenital heart disease, palatal defects & about 25 \%
schizophrenia risk \\
\textbf{15q11--13} (duplication of maternal allele) & about 4
Mb & Autism, epilepsy & about 1 \% of autism \\
\textbf{NRXN1} (Neurexin-1) (heterozygous deletion) & Variable, exonic &
Autism, schizophrenia, intellectual disability & 5--10\texorpdfstring{\ensuremath{\times}}{x} enrichment \\
\textbf{CNTNAP2} (heterozygous loss) & Variable, exonic & Autism,
language delay, epilepsy & Modest enrichment \\
\textbf{3q29} (microdeletion) & about 1.6 Mb & Schizophrenia
(about 40\texorpdfstring{\ensuremath{\times}}{x} risk increase) & Rare but high penetrance \\
\textbf{22q11.2} (microduplication) & about 3 Mb & Increased
risk of schizophrenia & about 3\texorpdfstring{\ensuremath{\times}}{x} increase \\
\textbf{1q21.1} (deletion or duplication) & about 1.4 Mb &
Schizophrenia, autism, microcephaly (deletion) & about 10\texorpdfstring{\ensuremath{\times}}{x}
risk \\
\textbf{15q13.3} (deletion) & about 1.5 Mb & Schizophrenia,
epilepsy (deletion) & about 10\texorpdfstring{\ensuremath{\times}}{x} risk \\
\end{longtable}
}

\textbf{Why CNVs are so impactful:} Each CNV affects 1--5
dosage-sensitive genes simultaneously. The cumulative effect on neural
development is large because brain development depends on stoichiometric
protein interactions. Rare CNVs (occurring in \textless{} 1 \% of
population) collectively explain about 5 \% of schizophrenia
cases and about 10 \% of autism cases --- a substantial
fraction of risk despite each individual CNV being rare.

\subsubsection{Genomic Variation}\label{genomic-variation}

\begin{itemize}
\tightlist
\item
  \textbf{Single nucleotide polymorphisms (SNPs)}: about 1 per
  300 bp on average; the human genome contains about 4-5
  million SNPs per individual relative to the reference. Most are
  neutral; about 2\% in coding regions; about 0.5\%
  are missense.
\item
  \textbf{Indels}: about 500,000 per genome; usually 1-50 bp
\item
  \textbf{Copy number variants (CNVs)}: about 1,000 per genome;
  collectively encompass about 4 Mb
\item
  \textbf{Structural variants (SVs)}: inversions, translocations, large
  insertions; about 2,500 per genome
\item
  \textbf{Two individuals differ by about 0.1\% of their
  genomes (about 3 million SNPs)}, but most variation is within
  populations rather than between populations (Lewontin, 1972; Fst
  about 0.12).
\end{itemize}

\subsubsection{The ENCODE Project}\label{the-encode-project}

The \textbf{ENCODE} (Encyclopedia of DNA Elements) Project
(2003-present) systematically mapped functional elements in the human
genome:

\begin{itemize}
\tightlist
\item
  \textbf{Key finding}: \textgreater80\% of the genome has at least one
  biochemical function (\hyperref[gl:chromatin]{\textbf{chromatin}}
  accessibility, transcription factor binding,
  \hyperref[gl:histone]{\textbf{histone}} modification, RNA
  transcription)
\item
  \textbf{Revised ``junk DNA'' narrative}: Much non-coding DNA has
  regulatory function (enhancers, insulators, lncRNAs); however,
  ``biochemical activity'' does not necessarily mean ``biologically
  important'' (ongoing debate)
\item
  \textbf{about 400,000 enhancers} identified; most
  disease-associated SNPs (from GWAS) fall in non-coding regulatory
  regions, not protein-coding genes
\end{itemize}

\begin{center}\rule{0.5\linewidth}{0.5pt}\end{center}

\subsection{Next-Generation Sequencing
Technologies}\label{next-generation-sequencing-technologies}

\begin{figure}[htbp]
\centering
\includegraphics[width=0.82\linewidth,height=4.2in,keepaspectratio]{../figures/mermaid_inline/inline_mermaid_0077_380014b63899.png}
\caption{Genomic sequencing approaches. Short-read (Illumina) provides high throughput and accuracy for variant calling; long-read (Nanopore, PacBio HiFi) resolves structural variants, repeats, and enables de novo assembly.}
\end{figure}

\subsubsection{Pangenome Graph
Reasoning}\label{pangenome-graph-reasoning}

The old mental model of a genome reference was a single linear string.
That is still useful, but it hides ancestry-specific haplotypes,
structural variants, alternate gene copies, and repeat-resolved sequence
that do not align cleanly to one coordinate path. The Human Pangenome
Reference Consortium's draft reference introduced a graph built from 47
diverse individuals and 94 phased haplotypes \citep{humanpangenome2023}.
In a graph, a variant is not merely ``different from the reference''; it
is a path through an explicitly represented set of alternatives. This
improves reasoning about structural variation and representation, while
adding new validation obligations: graph alignment, annotation transfer,
clinical reporting, and population sampling must become part of the
claim.

\begin{figure}[htbp]
\centering
\includegraphics[width=0.82\linewidth,height=4.2in,keepaspectratio]{../figures/mermaid_inline/inline_mermaid_0078_74e762424472.png}
\caption{Pangenome graph reasoning. A graph reference turns population variation into explicit alternate paths, improving structural-variant interpretation while making validation and reporting more complex.}
\end{figure}

{\def\LTcaptype{none} % do not increment counter
\begin{longtable}[]{@{}
  >{\raggedright\arraybackslash}p{(\linewidth - 10\tabcolsep) * \real{0.1310}}
  >{\raggedright\arraybackslash}p{(\linewidth - 10\tabcolsep) * \real{0.1310}}
  >{\raggedright\arraybackslash}p{(\linewidth - 10\tabcolsep) * \real{0.1190}}
  >{\raggedright\arraybackslash}p{(\linewidth - 10\tabcolsep) * \real{0.1310}}
  >{\raggedright\arraybackslash}p{(\linewidth - 10\tabcolsep) * \real{0.2857}}
  >{\raggedright\arraybackslash}p{(\linewidth - 10\tabcolsep) * \real{0.2024}}@{}}
\toprule\noalign{}
\begin{minipage}[b]{\linewidth}\raggedright
Technology
\end{minipage} & \begin{minipage}[b]{\linewidth}\raggedright
Read Length
\end{minipage} & \begin{minipage}[b]{\linewidth}\raggedright
Accuracy
\end{minipage} & \begin{minipage}[b]{\linewidth}\raggedright
Throughput
\end{minipage} & \begin{minipage}[b]{\linewidth}\raggedright
Cost (30x Human Genome)
\end{minipage} & \begin{minipage}[b]{\linewidth}\raggedright
Key Application
\end{minipage} \\
\midrule\noalign{}
\endhead
\bottomrule\noalign{}
\endlastfoot
Illumina NovaSeq X & 150-300 bp & \textgreater99.9\% (Q30+) &
about 600 Gb/run & about \$200 & Variant calling,
RNA-seq, exome \\
PacBio Revio (HiFi) & 10-25 kb & \textgreater99.9\% (consensus) &
about 90 Gb/run & about \$1,000 & Structural variants,
de novo assembly, \hyperref[gl:epigenetics]{\textbf{epigenetics}} \\
Oxford Nanopore PromethION & \textgreater100 kb possible &
about 99\% (Q20+) & about 100 Gb/run &
about \$500 & Real-time sequencing, field deployment, direct
RNA sequencing \\
\end{longtable}
}

\begin{center}\rule{0.5\linewidth}{0.5pt}\end{center}

\subsection{Genome-Wide Association Studies
(GWAS)}\label{genome-wide-association-studies-gwas}

GWAS identify genetic variants (typically SNPs) associated with complex
traits or diseases by comparing allele frequencies between cases and
controls.

\textbf{Methodology}: 1. Genotype hundreds of thousands to millions of
SNPs across the genome in large cohorts (typically \textgreater10,000
individuals) 2. Test each SNP for association with the
\hyperref[gl:phenotype]{\textbf{phenotype}} using logistic regression
(binary traits) or linear regression (quantitative traits) 3. Apply a
stringent significance threshold (\(p < 5 \times 10^{-8}\)) to account
for about 1 million independent tests (Bonferroni-like
correction) 4. Display results on a \textbf{Manhattan plot}: genomic
position (x-axis) vs.~\(-\log_{10}(p)\) (y-axis); significant
associations appear as peaks above the genome-wide significance line

\textbf{Key concepts}: - \textbf{\hyperref[gl:linkage]{Linkage}
disequilibrium (LD)}: Non-random association of alleles at nearby loci.
A significant SNP may not be the causal variant but rather in LD with it
(a ``tag SNP''). - \textbf{Effect sizes}: Most GWAS-identified variants
have small effect sizes (odds ratio 1.1-1.3); complex traits are highly
polygenic - \textbf{Missing heritability}: GWAS variants collectively
explain primarily a fraction of heritability for most traits (e.g.,
about 25\% of height heritability, despite identifying
\textgreater700 loci)

\textbf{Examples of GWAS discoveries}: - Type 2 diabetes:
\textgreater400 loci, including TCF7L2 (strongest effect), PPARG, KCNJ11
- Age-related macular degeneration: CFH (complement factor H) -- led to
complement-targeted therapies - Autoimmune diseases: HLA region as the
strongest risk factor for type 1 diabetes, celiac disease, rheumatoid
arthritis

\begin{center}\rule{0.5\linewidth}{0.5pt}\end{center}

\subsection{Comparative Genomics and Molecular
Evolution}\label{comparative-genomics-and-molecular-evolution}

\textbf{Synteny}: Conservation of gene order between species. Humans and
mice share about 90\% of genes in syntenic blocks, despite
75-80 million years of divergence.

\textbf{Conserved non-coding elements (CNEs)}: about 5\% of the
human genome is under purifying selection (much more than the
about 1.5\% that encodes protein). These conserved regions
likely include enhancers, silencers, and structural elements essential
for gene regulation.

\textbf{Ultraconserved elements (UCEs)}: A particularly remarkable
subset --- about 481 segments of \texorpdfstring{\textgreater{}=}{>=} 200 bp that are perfectly
identical (100 \% sequence identity) across human, mouse, and rat
genomes (Bejerano et al., 2004 \emph{Science}). These elements have
remained unchanged for about 80 million years, implying extreme
purifying selection. UCEs cluster near developmental transcription
factors (e.g., \emph{DACH1}, \emph{POU3F2}) and are enriched in
enhancers controlling early embryonic patterning. Their function remains
incompletely understood --- paradoxically, knocking out individual UCEs
in mice produces primarily subtle phenotypes, raising questions about
why selection is so strong if redundancy is high.

\textbf{Comparative-genomics signatures of selection:}

{\def\LTcaptype{none} % do not increment counter
\begin{longtable}[]{@{}
  >{\raggedright\arraybackslash}p{(\linewidth - 4\tabcolsep) * \real{0.2250}}
  >{\raggedright\arraybackslash}p{(\linewidth - 4\tabcolsep) * \real{0.4250}}
  >{\raggedright\arraybackslash}p{(\linewidth - 4\tabcolsep) * \real{0.3500}}@{}}
\toprule\noalign{}
\begin{minipage}[b]{\linewidth}\raggedright
Signature
\end{minipage} & \begin{minipage}[b]{\linewidth}\raggedright
Detection method
\end{minipage} & \begin{minipage}[b]{\linewidth}\raggedright
Interpretation
\end{minipage} \\
\midrule\noalign{}
\endhead
\bottomrule\noalign{}
\endlastfoot
Conserved across species & PhastCons, PhyloP & Purifying selection
(function preserved) \\
Conserved + accelerated in lineage & HAR (Human Accelerated Regions) &
Adaptive change in human lineage (e.g., HAR1 \texorpdfstring{\ensuremath{\to}}{->} schizophrenia, autism
overlap) \\
Population-level fixation & Fst, XP-EHH & Recent positive selection in a
population (e.g., LCT lactase persistence in Europeans) \\
dN/dS \textgreater{} 1 & PAML, codeml & Positive selection at protein
level (e.g., immune genes, primate FOXP2) \\
TajimaD \textless{} 0 & Coalescent statistics & Recent selective
sweep \\
Long-range LD (haplotype block) & iHS, EHH & Recent adaptive
selection \\
\end{longtable}
}

\subsubsection{Mutation Rates --- Spontaneous vs
Induced}\label{mutation-rates-spontaneous-vs-induced}

Mutation rate varies enormously with cellular context, mutagen exposure,
and DNA repair capacity. The following table contrasts representative
spontaneous and induced mutation frequencies:

{\def\LTcaptype{none} % do not increment counter
\begin{longtable}[]{@{}
  >{\raggedright\arraybackslash}p{(\linewidth - 4\tabcolsep) * \real{0.1053}}
  >{\raggedright\arraybackslash}p{(\linewidth - 4\tabcolsep) * \real{0.8070}}
  >{\raggedright\arraybackslash}p{(\linewidth - 4\tabcolsep) * \real{0.0877}}@{}}
\toprule\noalign{}
\begin{minipage}[b]{\linewidth}\raggedright
Source
\end{minipage} & \begin{minipage}[b]{\linewidth}\raggedright
Mutation rate (per bp per generation/division)
\end{minipage} & \begin{minipage}[b]{\linewidth}\raggedright
Notes
\end{minipage} \\
\midrule\noalign{}
\endhead
\bottomrule\noalign{}
\endlastfoot
\textbf{Spontaneous mutations} & & \\
Human germline & 1.2 \texorpdfstring{\ensuremath{\times}}{x} 10\texorpdfstring{\textsuperscript{-}}{-}\texorpdfstring{\textsuperscript{8}}{8} per bp per generation & Trio sequencing
(Kong et al., 2012); paternal age \texorpdfstring{\ensuremath{\times}}{x} 2 from age 20 to 50 \\
Human somatic (most tissues) & 5 \texorpdfstring{\ensuremath{\times}}{x} 10\texorpdfstring{\textsuperscript{-}}{-}\texorpdfstring{\textsuperscript{9}}{9} to 10\texorpdfstring{\textsuperscript{-}}{-}\texorpdfstring{\textsuperscript{8}}{8} per cell division &
Stem-cell-population mutation rate \\
\emph{E. coli} (wild-type) & 5 \texorpdfstring{\ensuremath{\times}}{x} 10\texorpdfstring{\textsuperscript{-}}{-}\texorpdfstring{\textsuperscript{1}}{1}\texorpdfstring{\textsuperscript{0}}{0} per bp per generation & High
polymerase fidelity + MMR \\
\emph{E. coli} mutator (mutS-) & 5 \texorpdfstring{\ensuremath{\times}}{x} 10\texorpdfstring{\textsuperscript{-}}{-}\texorpdfstring{\textsuperscript{7}}{7} per bp per generation & 1000\texorpdfstring{\ensuremath{\times}}{x}
elevated due to MMR loss \\
\emph{S. cerevisiae} (yeast) & 2 \texorpdfstring{\ensuremath{\times}}{x} 10\texorpdfstring{\textsuperscript{-}}{-}\texorpdfstring{\textsuperscript{1}}{1}\texorpdfstring{\textsuperscript{0}}{0} per bp per generation & \\
\textbf{Induced mutations} & & \\
UV (sub-lethal, 10 J/m\texorpdfstring{\textsuperscript{2}}{2}) & 10\texorpdfstring{\textsuperscript{-}}{-}\texorpdfstring{\textsuperscript{5}}{5} to 10\texorpdfstring{\textsuperscript{-}}{-}\texorpdfstring{\textsuperscript{4}}{4} per bp at TT, TC sites & CPDs
and 6-4 photoproducts; localised to pyrimidine dimers \\
Ionizing radiation (1 Gy) & about 50--100 DSBs per cell & DSBs
cluster in damage tracks; many fatal if unrepaired \\
Cigarette smoke (chronic exposure) & about 150 missense
mutations/year in lung epithelium & Benzo\hyperref[a]{a}pyrene \texorpdfstring{\ensuremath{\to}}{->} G:C \texorpdfstring{\ensuremath{\to}}{->}
T:A transversions (signature 4) \\
EMS (ethyl methanesulfonate, 50 mM) & 5 \texorpdfstring{\ensuremath{\times}}{x} 10\texorpdfstring{\textsuperscript{-}}{-}\texorpdfstring{\textsuperscript{4}}{4} per bp per generation in
\emph{Drosophila} & Workhorse mutagen; G:C \texorpdfstring{\ensuremath{\to}}{->} A:T transitions \\
5-FU treatment (cancer chemo) & 100--1000-fold elevation in dividing
cells & Thymidylate synthase inhibition + mis-incorporation \\
\textbf{Disease-state elevation} & & \\
Lynch syndrome (MMR-deficient tumours) & 10\texorpdfstring{\textsuperscript{-}}{-}\texorpdfstring{\textsuperscript{5}}{5} per bp; microsatellite
instability (MSI) & Poll/Pold mutations also cause Lynch-like
phenotypes \\
MUTYH-associated polyposis & 10\texorpdfstring{\ensuremath{\times}}{x} G:C \texorpdfstring{\ensuremath{\to}}{->} T:A elevation &
OGG1+MUTYH-deficient \\
BRCA1/2-deficient cancer & High SV burden (BRCA1) or chromothripsis
(BRCA2) & HR-defective; relies on error-prone NHEJ/MMEJ \\
POLE proofreading-deficient cancer & 10\texorpdfstring{\textsuperscript{-}}{-}\texorpdfstring{\textsuperscript{4}}{4} to 10\texorpdfstring{\textsuperscript{-}}{-}\texorpdfstring{\textsuperscript{5}}{5} per bp &
``Ultramutator'' phenotype (about 100,000 mutations per
tumour); responds to anti-PD-1 \\
\end{longtable}
}

\textbf{Why the mutation rate is so finely tuned.} The germline mutation
rate of about 1.2 \texorpdfstring{\ensuremath{\times}}{x} 10\texorpdfstring{\textsuperscript{-}}{-}\texorpdfstring{\textsuperscript{8}}{8} produces about 70 \emph{de
novo} mutations per offspring, an upper limit set by cumulative effects
across generations: too high, and ``Muller's ratchet'' accumulates
deleterious mutations faster than selection can purge them; too low, and
there is insufficient variation for adaptation. Different organisms have
evolved different rates: high-fecundity microbes (10\texorpdfstring{\textsuperscript{-}}{-}\texorpdfstring{\textsuperscript{9}}{9} to 10\texorpdfstring{\textsuperscript{-}}{-}\texorpdfstring{\textsuperscript{1}}{1}\texorpdfstring{\textsuperscript{0}}{0}),
mammals (10\texorpdfstring{\textsuperscript{-}}{-}\texorpdfstring{\textsuperscript{8}}{8}), and viruses with their relaxed proofreading (HIV: 10\texorpdfstring{\textsuperscript{-}}{-}\texorpdfstring{\textsuperscript{5}}{5};
influenza: 10\texorpdfstring{\textsuperscript{-}}{-}\texorpdfstring{\textsuperscript{4}}{4}) reflect different optima between fidelity costs and
adaptive flexibility.

\textbf{Gene duplication and divergence} \citep{ohno1970}: Whole-genome
duplications (WGD) and segmental duplications create paralogous genes.
After duplication: - \textbf{Neofunctionalization}: One copy acquires a
new function - \textbf{Subfunctionalization}: Each copy retains a subset
of the original function - \textbf{Pseudogenization}: One copy
accumulates mutations and becomes a pseudogene

\textbf{Ka/Ks ratio} (dN/dS): The ratio of nonsynonymous (amino
acid-changing) to synonymous (silent) substitution rates: - Ka/Ks
\textless{} 1: purifying (negative) selection (most genes) - Ka/Ks = 1:
neutral evolution - Ka/Ks \textgreater{} 1: positive (adaptive)
selection (rare; examples: immune genes, reproductive proteins)

\begin{center}\rule{0.5\linewidth}{0.5pt}\end{center}

\subsection{Personal Genomics and Ethical
Issues}\label{personal-genomics-and-ethical-issues}

\subsubsection{Pharmacogenomics}\label{pharmacogenomics}

Genetic variation affects drug response:

Pharmacogenomics is strongest when it is treated as one evidence layer
in prescribing rather than a deterministic label. A genotype can alter
enzyme activity, transporter expression, immune recognition, or drug
target sensitivity, but clinical interpretation also depends on ancestry
representation in the evidence base, co-medications, age, liver and
kidney function, indication, and whether the variant is actually phased
or captured by the assay. Pangenome and long-read resources improve the
chance of detecting structural or ancestry-associated variants missed by
a single linear reference, but clinical use still requires validated
phenotype links and reporting standards \citep{humanpangenome2023}.

{\def\LTcaptype{none} % do not increment counter
\begin{longtable}[]{@{}
  >{\raggedright\arraybackslash}p{(\linewidth - 4\tabcolsep) * \real{0.1935}}
  >{\raggedright\arraybackslash}p{(\linewidth - 4\tabcolsep) * \real{0.1935}}
  >{\raggedright\arraybackslash}p{(\linewidth - 4\tabcolsep) * \real{0.6129}}@{}}
\toprule\noalign{}
\begin{minipage}[b]{\linewidth}\raggedright
Gene
\end{minipage} & \begin{minipage}[b]{\linewidth}\raggedright
Drug
\end{minipage} & \begin{minipage}[b]{\linewidth}\raggedright
Effect of Variant
\end{minipage} \\
\midrule\noalign{}
\endhead
\bottomrule\noalign{}
\endlastfoot
\textbf{CYP2D6} & Codeine, tamoxifen & Poor metabolizers: codeine
ineffective (cannot convert to morphine); ultra-rapid metabolizers:
toxicity risk \\
\textbf{CYP2C19} & Clopidogrel (Plavix) & Poor metabolizers: reduced
anti-platelet effect; increased cardiovascular risk \\
\textbf{VKORC1 + CYP2C9} & Warfarin & Variants explain
about 40\% of dosing variation; FDA-approved pharmacogenomic
labeling \\
**HLA-B*5701** & Abacavir (HIV) & Pre-treatment testing mandatory;
positive patients have about 50\% risk of hypersensitivity
reaction \\
\textbf{DPYD} & 5-fluorouracil & DPD deficiency: potentially fatal
toxicity; pre-treatment testing recommended in EU \\
\end{longtable}
}

\subsubsection{Ethical Considerations}\label{ethical-considerations}

\begin{itemize}
\tightlist
\item
  \textbf{Genetic privacy}: Genetic information reveals predisposition,
  not certainty. The \textbf{GINA Act} (Genetic Information
  Nondiscrimination Act, 2008, USA) prohibits genetic discrimination in
  employment and health insurance but does NOT cover life insurance,
  disability insurance, or long-term care insurance.
\item
  \textbf{Direct-to-consumer (DTC) testing}: Companies like 23andMe
  provide ancestry and health risk information. Concerns include:
  limited clinical validity for most variants, potential for anxiety,
  and data privacy (law enforcement access to databases).
\item
  \textbf{BRCA1/2 testing decisions}: Positive results require careful
  genetic counseling. Prophylactic mastectomy reduces breast cancer risk
  by \textgreater90\% but is psychologically significant. The 2013
  \emph{Myriad Genetics} Supreme Court decision (Association for
  Molecular Pathology v. Myriad Genetics) ruled that naturally occurring
  genes cannot be patented but \hyperref[gl:cdna]{\textbf{cDNA}} can be.
\item
  \textbf{Germline editing}: Editing human embryos raises profound
  ethical questions. The 2018 He Jiankui case (CCR5-edited babies in
  China) was widely condemned as premature, reckless, and ethically
  unacceptable. International consensus: germline editing for
  reproduction should not proceed until safety and efficacy are
  established and broad societal consensus is reached.
\end{itemize}

\begin{quote}
\textbf{Clinical Connection: Pharmacogenomic Testing in Practice} The
Clinical Pharmacogenetics Implementation Consortium (CPIC) publishes
evidence-based guidelines for 23+ gene-drug pairs. Pre-emptive
pharmacogenomic testing (genotyping a panel of pharmacogenes before any
drug is prescribed) is increasingly implemented in health systems. The
PREDICT study (Vanderbilt) and RIGHT study (Mayo Clinic) demonstrated
that pre-emptive testing is feasible and can prevent adverse drug
reactions.
\end{quote}

\begin{center}\rule{0.5\linewidth}{0.5pt}\end{center}

\subsection{Worked Example: Per-Gene De Novo Mutation
Probability}\label{worked-example-per-gene-de-novo-mutation-probability}

\textbf{Problem}: The human germline point-mutation rate is
\(\mu \approx 1.2 \times 10^{-8}\) per bp per generation. The dystrophin
gene (\emph{DMD}) has the largest coding sequence in the genome,
\(11{,}058 \text{ bp}\) (both alleles together giving a diploid coding
target \(L = 2 \times 11{,}058 \text{ bp}\)).

\begin{enumerate}
\def\labelenumi{(\alph{enumi})}
\tightlist
\item
  What is the expected number of de novo coding mutations in \emph{DMD}
  per offspring?
\end{enumerate}

\textbf{Solution}: The Poisson mean is the per-bp rate times the diploid
coding target length:

\begin{equation}\lambda = \mu L = 1.2 \times 10^{-8} \times 22{,}116 \text{ bp} = 2.65 \times 10^{-4} \text{ mutations/generation} \tag{14.5}\label{eq:dmd_cds_lambda}\end{equation}

\begin{enumerate}
\def\labelenumi{(\alph{enumi})}
\setcounter{enumi}{1}
\tightlist
\item
  Using the Poisson model, what is the probability that a given child
  carries no new coding mutation in \emph{DMD}?
\end{enumerate}

\begin{equation}P(k=0) = \frac{(\lambda)^0 e^{-\lambda}}{0!} = e^{-2.65 \times 10^{-4}} \approx 0.99973 \tag{14.6}\label{eq:dmd_cds_p0}\end{equation}

\begin{enumerate}
\def\labelenumi{(\alph{enumi})}
\setcounter{enumi}{2}
\tightlist
\item
  What is the probability of at least one new coding mutation in
  \emph{DMD}?
\end{enumerate}

\begin{equation}P(k \geq 1) = 1 - e^{-\lambda} \approx 2.65 \times 10^{-4} \approx \frac{1}{3{,}800} \text{ births} \tag{14.7}\label{eq:dmd_cds_pge1}\end{equation}

Even though any single child is overwhelmingly likely to inherit an
intact \emph{DMD} coding sequence, a rate of roughly \(1\) in
\(3{,}800\) births applied across a large population continually
regenerates new Duchenne muscular dystrophy alleles --- which is why
this X-linked disorder persists despite strong negative selection
against affected males.

\subsection{Worked Example: Mutation
Analysis}\label{worked-example-mutation-analysis}

\textbf{Problem}: A wild-type mRNA sequence reads:

5'-AUG-UUU-GGA-GAA-CUU-UAG-3'

\begin{enumerate}
\def\labelenumi{(\alph{enumi})}
\tightlist
\item
  Translate the wild-type sequence.
\end{enumerate}

{\def\LTcaptype{none} % do not increment counter
\begin{longtable}[]{@{}ll@{}}
\toprule\noalign{}
Codon & Amino Acid \\
\midrule\noalign{}
\endhead
\bottomrule\noalign{}
\endlastfoot
AUG & Met \\
UUU & Phe \\
GGA & Gly \\
GAA & Glu \\
CUU & Leu \\
UAG & Stop \\
\end{longtable}
}

\textbf{Protein}: Met-Phe-Gly-Glu-Leu (5 amino acids)

\begin{enumerate}
\def\labelenumi{(\alph{enumi})}
\setcounter{enumi}{1}
\tightlist
\item
  A transition mutation changes the 10th nucleotide (first position of
  codon 4) from G to A: GAA to AAA.
\end{enumerate}

{\def\LTcaptype{none} % do not increment counter
\begin{longtable}[]{@{}lll@{}}
\toprule\noalign{}
Codon & Wild-type AA & Mutant AA \\
\midrule\noalign{}
\endhead
\bottomrule\noalign{}
\endlastfoot
GAA & Glu & AAA = Lys \\
\end{longtable}
}

This is a \textbf{missense mutation} (Glu to Lys: non-conservative,
charge reversal).

\begin{enumerate}
\def\labelenumi{(\alph{enumi})}
\setcounter{enumi}{2}
\tightlist
\item
  A transition changes the 7th nucleotide (first position of codon 3)
  from G to A: GGA to AGA.
\end{enumerate}

GGA (Gly) to AGA (Arg): \textbf{missense} (non-conservative; small
nonpolar to large positive).

\begin{enumerate}
\def\labelenumi{(\alph{enumi})}
\setcounter{enumi}{3}
\tightlist
\item
  A deletion removes the 4th nucleotide (U): sequence becomes
  AUG-UUG-GAG-AAC-UUU-AG\ldots{}
\end{enumerate}

This is a \textbf{frameshift mutation}. New reading frame:
Met-Leu-Glu-Asn-Phe\ldots{} (completely different protein, no stop codon
in remaining sequence -- read-through into 3' UTR).

\begin{center}\rule{0.5\linewidth}{0.5pt}\end{center}

\subsection{Computational Bridge}\label{computational-bridge-13}

In silico \hyperref[gl:translation]{\textbf{translation}} checks
mutation consequences quickly:

\begin{Shaded}
\begin{Highlighting}[]
\ImportTok{from}\NormalTok{ biology.genetics }\ImportTok{import}\NormalTok{ translate\_mrna}

\NormalTok{wt }\OperatorTok{=}\NormalTok{ translate\_mrna(}\StringTok{"AUGGAAUAA"}\NormalTok{)}
\BuiltInTok{print}\NormalTok{(wt)}
\end{Highlighting}
\end{Shaded}

\begin{quote}
\textbf{Clinical / systems note:} Pharmacogenomic dosing (warfarin,
thiopurines, clopidogrel) is implemented as clinical algorithms over the
same variant catalogue GWAS discovers --- bridging population statistics
to individual prescriptions.
\end{quote}

\begin{center}\rule{0.5\linewidth}{0.5pt}\end{center}

\subsubsection{Post-Approval CRISPR Therapeutics: Casgevy and the Safety
Frontier}\label{post-approval-crispr-therapeutics-casgevy-and-the-safety-frontier}

The first CRISPR--Cas9 therapy approved by the FDA and EMA
(\textbf{Casgevy / exagamglogene autotemcel}, Vertex/CRISPR
Therapeutics) provides a quantitative benchmark for the entire field
\citep{fda2023casgevy,fda2024casgevythalassemia,fda2026casgevy}. The
therapy edits autologous CD34\texorpdfstring{\textsuperscript{+}}{+} haematopoietic stem cells \emph{ex vivo}
at the \textbf{BCL11A enhancer} to de-repress fetal haemoglobin (HbF),
then reinfuses them after conditioning chemotherapy. In the pivotal
trials, \textbf{28 of 29 sickle-cell patients became free of
vaso-occlusive crises for \textgreater{} 12 months}, and \textbf{39 of
42 \texorpdfstring{\ensuremath{\beta}}{beta}-thalassemia patients became transfusion-independent} --- outcomes
previously achievable primarily via allogeneic bone-marrow transplant
with its attendant graft-versus-host morbidity.

The safety discussion that accompanies this approval is instructive.
Off-target editing in CD34\texorpdfstring{\textsuperscript{+}}{+} cells is measured pre-infusion by
\textbf{GUIDE-seq} and \textbf{deep amplicon sequencing} at 100+
predicted sites --- with Casgevy, off-target editing fell below the
detection floor (about 0.1 \%) at most surveyed loci. Residual
concerns include \textbf{large structural variants} (kilobase-scale
deletions, inversions, chromothripsis) that amplicon-seq can miss;
long-read WGS of treated HSC clones is the emerging gold-standard check.
Longitudinal monitoring follows these patients for \textbf{15 years} for
malignancy, driven by the theoretical risk that any DNA double-strand
break event in a stem cell compartment could, over decades, seed a
clonal expansion. This regulatory framework --- efficacy by molecular
endpoint, safety by unbiased structural-variant surveillance,
surveillance by decade-scale cohort --- is now the template every
subsequent \emph{in vivo} CRISPR therapy (e.g., Verve's PCSK9 base
editor for hypercholesterolaemia) must meet.

\subsubsection{The Telomere-to-Telomere (T2T) Consortium: A Truly
Complete Human
Genome}\label{the-telomere-to-telomere-t2t-consortium-a-truly-complete-human-genome}

The human reference genome, from the Human Genome Project's 2003 draft
through GRCh38 (2013), contained \textbf{\textgreater{} 200 Mb of
unresolved gaps} --- chiefly heterochromatic satellite arrays at
centromeres, the short arms of acrocentric chromosomes (13, 14, 15, 21,
22), and segmental duplications. The \textbf{T2T-CHM13 assembly} (Nurk
et al., \emph{Science} 2022) closed these gaps using long-read
sequencing (PacBio HiFi, Oxford Nanopore Ultra-long) on a hydatidiform
mole (uniform haplotype, no
\hyperref[gl:heterozygous]{\textbf{heterozygous}} complications),
yielding the first \textbf{gapless 3.055 Gb human genome}.

What newly resolvable sequence revealed: \textbf{centromeric \texorpdfstring{\ensuremath{\alpha}}{alpha}-satellite
arrays} span 0.5--5 Mb per chromosome, organised into higher-order
repeat (HOR) units; the \textbf{rDNA arrays} on acrocentric p-arms vary
from 47 to 287 copies per haploid genome; roughly \textbf{2000
additional protein-coding genes and gene paralogs} were annotated (with
duplications in immunity, olfaction, and neurogenesis). Perhaps most
importantly, T2T demonstrated that prior ``dark'' regions harbour real
medical-genetic variation: \textbf{paracentric inversions at 16p11 and
22q11} --- known susceptibility loci for autism and schizophrenia ---
are present in about 10 \% of haplotypes but invisible to
GRCh38 mapping. The pangenome era extends T2T's logic: the reference is
now a graph, not a single string, and clinical genomic pipelines are
migrating cautiously because graph references improve representation
while adding alignment, annotation, and validation complexity
\citep{humanpangenome2023}.

\begin{center}\rule{0.5\linewidth}{0.5pt}\end{center}

\subsection{Current Evidence and Frontier
Biology}\label{current-evidence-and-frontier-biology-16}

For \textbf{Mutations, CRISPR, and Genomics}, frontier biology belongs
inside the evidence logic of the chapter. Molecular genetics now spans
single-reference sequences, telomere-to-telomere assemblies, pangenome
graphs, long-read sequencing, CRISPR medicines, and ethical deployment.
The core reading question is this: genomics claims should distinguish
variant discovery, pathogenic interpretation, ancestry representation,
and clinical actionability.

\begin{itemize}
\tightlist
\item
  \textbf{What to verify:} identify the observation, model, assay, or
  dataset that would make the claim stronger or weaker.
\item
  \textbf{What to qualify:} state the scale, organism, cell type,
  environmental condition, or population where the claim is expected to
  hold.
\item
  \textbf{What to compare:} test at least one alternative explanation,
  baseline, or null model before treating the pattern as causal.
\item
  \textbf{What to cite:} distinguish primary evidence, review synthesis,
  public dataset, and institutional guidance; for recent or numeric
  claims, prefer the source closest to the measurement and state what
  has changed since it was published.
\end{itemize}

When a genomic claim depends on a reference, ask whether short reads,
structural variants, ancestry representation, phasing, or clinical
validation could change the interpretation
\citep{humanpangenome2023,fda2023casgevy,fda2024casgevythalassemia}.

\textbf{Source practice:} For genomics and editing claims, distinguish
discovery from clinical actionability, and cite reference resources,
regulatory records, or primary editing studies close to the claim
\citep{humanpangenome2023,fda2026casgevy,chalumeau2025primeediting}.

\subsubsection{Current Evidence Map: Genome Editing From Variant to
Follow-Up}\label{current-evidence-map-genome-editing-from-variant-to-follow-up}

\begin{figure}[htbp]
\centering
\includegraphics[width=0.82\linewidth,height=4.2in,keepaspectratio]{../figures/mermaid_inline/inline_mermaid_0079_5a4c23b154f2.png}
\caption{Casgevy and prime-editing examples show why editing medicines require molecular endpoints, structural-variant surveillance, toxicity monitoring, and long follow-up \textbackslash{}citep\{fda2026casgevy,chalumeau2025primeediting\}.}
\end{figure}

\subsection{Summary}\label{summary-16}

\begin{itemize}
\tightlist
\item
  \textbf{Mutations} range from single nucleotide changes (transitions,
  transversions, silent, missense, nonsense) to large chromosomal
  rearrangements (deletions, duplications, inversions, translocations).
  Frameshifts and nonsense mutations are typically most disruptive.
\item
  \textbf{Mutagenic agents} include UV radiation (CPDs), alkylating
  agents (EMS), intercalators (ethidium bromide), and ROS (8-oxoG). The
  Ames test screens chemicals for mutagenicity.
\item
  \textbf{DNA repair}: BER (damaged bases), NER (bulky adducts), MMR
  (mismatches), NHEJ/HR (DSBs), Fanconi pathway (interstrand
  crosslinks). Defects cause cancer predisposition (XP, Lynch syndrome,
  BRCA-associated cancers, Fanconi anemia).
\item
  \textbf{CRISPR-Cas9}: sgRNA + Cas9 creates DSB at PAM-adjacent site;
  NHEJ for knockouts, HDR for precise edits. Base editors (C-to-T,
  A-to-G) and prime editors avoid DSBs. First approved therapy: Casgevy
  (sickle cell disease via BCL11A disruption).
\item
  \textbf{Human genome}: about 3.2 Gb; about 20,000
  protein-coding genes (1.5\%); 46\% transposable elements; ENCODE shows
  \textgreater80\% has biochemical activity. T2T consortium completed
  gapless reference (2022).
\item
  \textbf{NGS technologies}: Illumina (short-read, high throughput),
  PacBio HiFi and Nanopore (long-read, structural variant resolution).
\item
  \textbf{GWAS}: Identify SNP-trait associations at genome-wide
  significance (\(p < 5 \times 10^{-8}\)); most complex trait variants
  have small effect sizes; results displayed on Manhattan plots.
\item
  \textbf{Ethics}: GINA Act protections, pharmacogenomic testing
  (CYP2D6, warfarin dosing), BRCA testing decisions, germline editing
  moratorium.
\item
  \textbf{Connections:} See
  \cref{sec:unit_IV_dna_replication_and_cell_cycle} for repair pathways,
  Unit V for inheritance of variants, and Unit VII for CRISPR origins in
  bacterial immunity.
\end{itemize}

\begin{center}\rule{0.5\linewidth}{0.5pt}\end{center}

\subsection{Review Questions}\label{review-questions-16}

\begin{enumerate}
\def\labelenumi{\arabic{enumi}.}
\tightlist
\item
  Distinguish between transitions and transversions. Which is more
  common and why?
\item
  Explain why frameshift mutations are generally more deleterious than
  missense mutations. Give an exception where a missense mutation is
  highly pathogenic.
\item
  A patient has microsatellite instability in their colorectal tumor.
  What DNA repair pathway is likely defective? Name the genes involved.
\item
  Describe the CRISPR-Cas9 mechanism step by step, from guide design to
  DSB repair.
\item
  Compare base editing and prime editing. When would you choose each
  technology?
\item
  What is the ``end-replication problem'' of the Human Genome Project,
  and how did the T2T consortium solve it?
\item
  Explain GWAS methodology. Why does a Manhattan plot show peaks, and
  what does a significant peak mean?
\item
  What is linkage disequilibrium, and why is it important for
  interpreting GWAS results?
\item
  Describe the pharmacogenomic basis for warfarin dose variation. Which
  genes are involved?
\item
  Discuss the ethical issues raised by the He Jiankui CRISPR babies
  case. What international guidelines exist for germline editing?
\item
  Explain why a SNP in a \textbf{splice donor} can have larger effect
  sizes than a missense change in the same exon, using NMD concepts.
\item
  A GWAS hit lies in intronic sequence with no coding change. Propose
  two mechanistic classes of causal variants compatible with that
  pattern.
\end{enumerate}

\begin{center}\rule{0.5\linewidth}{0.5pt}\end{center}

\subsection{Further Reading and Source
Notes}\label{further-reading-and-source-notes-16}

\begin{itemize}
\tightlist
\item
  Ward (1988). In the Shadow of the Dinosaurs: Early Mesozoic Tetrapods.
  \emph{Science}, 240.
\item
  Ohno (1970). \emph{Evolution by Gene Duplication}. Springer-Verlag.
\item
  Doudna \& Charpentier (2014). The new frontier of genome engineering
  with \{CRISPR-Cas9\}. \emph{Science}, 346.
\end{itemize}

\begin{center}\rule{0.5\linewidth}{0.5pt}\end{center}

\subsection{Key Terms}\label{key-terms-16}

\begin{enumerate}
\def\labelenumi{\arabic{enumi}.}
\tightlist
\item
  \textbf{Transition} -- purine-to-purine or pyrimidine-to-pyrimidine
  point mutation
\item
  \textbf{Transversion} -- purine-to-pyrimidine (or reverse) point
  mutation
\item
  \textbf{Frameshift mutation} -- insertion/deletion not a multiple of 3
  bp, altering reading frame
\item
  \textbf{Nonsense-mediated decay (NMD)} -- mRNA surveillance pathway
  eliminating transcripts with premature stop codons
\item
  \textbf{Trinucleotide repeat expansion} -- unstable repeat growth
  causing anticipation in diseases like Huntington
\item
  \textbf{Ames test} -- bacterial mutagenicity assay using
  histidine-requiring Salmonella revertants
\item
  \textbf{CRISPR-Cas9} -- RNA-guided endonuclease for programmable
  genome editing
\item
  \textbf{PAM (protospacer adjacent motif)} -- short DNA sequence (NGG
  for SpCas9) required for Cas9 binding
\item
  \textbf{Base editor} -- dCas9/nickase fused to deaminase for C-to-T or
  A-to-G conversion without DSB
\item
  \textbf{Prime editor} -- Cas9-nickase + reverse transcriptase + pegRNA
  for search-and-replace editing
\item
  \textbf{GWAS (genome-wide association study)} -- unbiased scan for
  SNP-trait associations across the genome
\item
  \textbf{Linkage disequilibrium} -- non-random association of alleles
  at nearby loci
\item
  \textbf{Manhattan plot} -- genome-wide visualization of association
  p-values
\item
  \textbf{Pharmacogenomics} -- study of genetic variation affecting drug
  response
\item
  \textbf{ENCODE} -- Encyclopedia of DNA Elements; mapped functional
  elements across the human genome
\item
  \textbf{Copy number variant (CNV)} -- duplication or deletion of a
  genomic segment \textgreater1 kb
\item
  \textbf{Synthetic lethality} -- cell death from simultaneous loss of
  two compensatory pathways
\end{enumerate}

\begin{center}\rule{0.5\linewidth}{0.5pt}\end{center}

\subsubsection{Companion Source
Module}\label{companion-source-module-16}

\textbf{Mutations, CRISPR, and Genomics} should leave a reproducible
trail from a biological claim to the code, figure, diagram, or
paper-based activity that can test it. Use the surfaces below to inspect
the chapter's assumptions, rerun the relevant model, or compare the
manuscript explanation with companion labs and figures.

{\def\LTcaptype{none} % do not increment counter
\begin{longtable}[]{@{}
  >{\raggedright\arraybackslash}p{(\linewidth - 2\tabcolsep) * \real{0.5000}}
  >{\raggedright\arraybackslash}p{(\linewidth - 2\tabcolsep) * \real{0.5000}}@{}}
\toprule\noalign{}
\begin{minipage}[b]{\linewidth}\raggedright
Surface
\end{minipage} & \begin{minipage}[b]{\linewidth}\raggedright
Use it for
\end{minipage} \\
\midrule\noalign{}
\endhead
\bottomrule\noalign{}
\endlastfoot
\breaktt{src/biology/genetics/genetics.py} (\texttt{dna\_complement},
\texttt{translate\_mrna}, \texttt{hamming\_distance},
\breaktt{jukes\_cantor\_distance}) & Compare sequence variants, coding
effects, and corrected molecular distances. \\
\breaktt{src/mermaid/biology\_diagrams.py}
(\breaktt{transcription\_translation\_diagram},
\breaktt{dna\_replication\_diagram}) & Link mutation class to replication
and expression context. \\
\end{longtable}
}

\textbf{Reproducibility check:} distinguish discovery technology,
reference representation, variant class, evidence level, and clinical
actionability. \textbf{Cross-reference:} connect with
\cref{sec:unit_IV_dna_replication_and_cell_cycle},
\cref{sec:unit_IV_gene_expression}, and
\cref{sec:unit_VII_bacteria_archaea_viruses}.

\newpage

\section{Epigenetics and Gene
Regulation}\label{epigenetics-and-gene-regulation}

\label{sec:unit_IV_epigenetics_and_gene_regulation}

\begin{quote}
\textbf{Ch 15} \texorpdfstring{\ensuremath{\cdot}}{.} Level 3/3 \texorpdfstring{\ensuremath{\cdot}}{.} 50 min read \texorpdfstring{\ensuremath{\cdot}}{.} 75 min lecture \texorpdfstring{\ensuremath{\cdot}}{.}
Prerequisites: \cref{sec:unit_IV_gene_expression}
\end{quote}

\subsection{Learning Objectives}\label{learning-objectives-17}

\begin{enumerate}
\def\labelenumi{\arabic{enumi}.}
\tightlist
\item
  Define \hyperref[gl:epigenetics]{\textbf{epigenetics}} and distinguish
  epigenetic modifications from DNA sequence changes.
\item
  Describe \hyperref[gl:nucleosome]{\textbf{nucleosome}} structure and
  the levels of \hyperref[gl:chromatin]{\textbf{chromatin}} compaction
  from the 11-nm fibre to the metaphase
  \hyperref[gl:chromosome]{\textbf{chromosome}}.
\item
  Explain the major classes of \hyperref[gl:histone]{\textbf{histone}}
  modifications (acetylation, methylation, phosphorylation,
  ubiquitination, sumoylation) and how they are written, erased, and
  read.
\item
  Describe the mechanism and function of DNA methylation, including
  \hyperref[gl:cpg-island]{\textbf{CpG island}}s, the DNMT1/3A/3B
  \hyperref[gl:enzyme]{\textbf{enzyme}}s, TET-mediated demethylation,
  and the role of methylation in \hyperref[gl:gene]{\textbf{gene}}
  silencing.
\item
  Distinguish Polycomb (PRC1/PRC2) and Trithorax (MLL/COMPASS) systems
  and explain how they maintain repressive and activating states.
\item
  Compare ATP-dependent chromatin remodelling families (SWI/SNF, ISWI,
  CHD/NuRD, INO80) and their distinct mechanisms.
\item
  Explain X-chromosome inactivation (Lyonisation) and the role of the
  XIST lncRNA, escape genes, and skewing.
\item
  Describe genomic imprinting using IGF2/H19 and Prader--Willi/Angelman
  as paradigms.
\item
  Explain how 3D genome organisation --- TADs, loops, compartments, and
  biomolecular condensates --- shapes transcription.
\item
  Explain microRNA (miRNA \citep{fire1998}) biogenesis and the mechanism
  of RISC-mediated post-transcriptional silencing.
\item
  Model the maintenance of DNA methylation and Polycomb marks across
  cell divisions quantitatively, and predict the consequences of
  perturbing each layer.
\item
  Evaluate evidence for transgenerational epigenetic inheritance in
  humans, the mechanism of mitotic and meiotic transmission of
  epigenetic marks, and the clinical implications of epigenetic
  dysregulation in cancer and neurodevelopmental disease.
\end{enumerate}

\begin{figure}[htbp]
\centering
\includegraphics[width=0.85\textwidth]{../figures/methylation_heatmap.png}
\caption{Illustrative synthetic CpG methylation heatmap across indexed loci (rows) and indexed samples (columns). The colour scale reports beta methylation fraction from 0 to 1; the lower-methylation row band is a deterministic teaching pattern, not patient or cell-line data.}
\label{fig:unit_IV_methylation_heatmap}
\end{figure}

\subsubsection{Study Blueprint}\label{study-blueprint-17}

\begin{itemize}
\tightlist
\item
  \textbf{Big idea:} Cells create stable yet reversible expression
  states through chromatin, DNA marks, and regulatory circuits.
\item
  \textbf{Core concepts:} chromatin, methylation, histone modification,
  enhancers.
\item
  \textbf{Framework alignment:} Vision \& Change: Information flow,
  exchange, and storage, Structure and function; AP Biology: Information
  Storage and Transmission, Systems Interactions; NGSS-style topics:
  Inheritance and Variation of Traits, Structure and Function.
\item
  \textbf{Model or quantitative lens:} Regulatory-state and
  expression-ratio reasoning.
\item
  \textbf{Data skill:} Interpret chromatin or expression evidence from
  simple regulatory datasets.
\item
  \textbf{Practice cadence:} Concept Explanation, Questions and Methods,
  Argumentation.
\item
  \textbf{Common misconception to repair:} Epigenetic does not mean
  independent of DNA sequence or permanently inherited.
\item
  \textbf{Primary lab:}
  \cref{sec:lab_unit_IV_epigenetics_and_gene_regulation}.
\item
  \textbf{Question bank:}
  \cref{sec:q_unit_IV_epigenetics_and_gene_regulation}.
\item
  \textbf{Transfer task:} Apply regulation logic to differentiation,
  imprinting, cancer, or environmental responses.
\item
  \textbf{Bridge to computation:}
  \breaktt{biology.genetics.genetics.cpg\_methylation\_remaining}.
\end{itemize}

\begin{center}\rule{0.5\linewidth}{0.5pt}\end{center}

\begin{quote}
\textbf{Opening Vignette --- The Landscape That Changed Genetics}

In 1942, the developmental biologist Conrad Waddington drew a diagram
that would define a science. He sketched a marble rolling down a
hillscape of ridges and valleys --- each valley representing a stable
cell fate, each ridge a threshold that once crossed was difficult to
reverse. He called this the \textbf{epigenetic landscape}: the topology
of developmental possibilities. Waddington coined ``epigenetics'' from
the Greek \emph{epi-} (above) to describe the heritable changes in gene
expression that occurred \emph{above} or \emph{beyond} the DNA sequence
--- changes that could not be explained by Mendelian genetics alone. He
had no idea of the molecular mechanisms. That understanding would come
forty years later.

In 1961, Mary Lyon \citep{lyon1961} noticed something peculiar in female
mice \hyperref[gl:heterozygous]{\textbf{heterozygous}} for coat-colour
\hyperref[gl:mutation]{\textbf{mutation}}s: their fur was a mosaic, not
intermediate. She proposed that in every cell of a female mammal, one of
the two X chromosomes is randomly and permanently silenced --- a
hypothesis proven so thoroughly that it was renamed Lyonisation. The
silenced X does not have a different DNA sequence. It has a different
\emph{chemistry}: dense methylation, hypoacetylated histones, and a long
non-coding RNA called XIST that coats the entire chromosome. This was
the first clear demonstration that a whole chromosome could be heritably
silenced without changing a single base pair. Epigenetics had found its
molecular identity.

Three more revolutions followed. In 2000, Strahl and Allis
\citep{strahl2000} proposed that distinct combinations of histone
modifications --- a ``histone code'' --- encode regulatory information
beyond the DNA sequence itself, predicting an entire pharmacology of
writers, erasers, and readers that has now produced FDA-approved drugs.
In the 2000s, chromosome-conformation-capture (Hi-C) revealed that the
genome is folded into stereotyped contact domains (TADs) where enhancers
find their cognate promoters, and that disrupting these domains can
mis-wire developmental control. And in the 2010s, biomolecular
condensates and phase separation reframed the nucleus as a collection of
liquid-like assemblies that concentrate transcriptional machinery at
super-enhancers. These insights together turned epigenetics from a
metaphor into a quantitative, druggable discipline.
\end{quote}

\subsection{Chromatin Structure and Nucleosome
Organisation}\label{chromatin-structure-and-nucleosome-organisation}

\subsubsection{The Nucleosome: Fundamental
Unit}\label{the-nucleosome-fundamental-unit}

Eukaryotic DNA is not naked; it is complexed with
\hyperref[gl:protein]{\textbf{protein}}s to form \textbf{chromatin}. The
fundamental repeating unit is the \textbf{nucleosome}:

\begin{itemize}
\tightlist
\item
  \textbf{Histone octamer:} 2 copies each of H2A, H2B, H3, H4 ---
  forming a spool-like protein disk
\item
  \textbf{DNA wrapping:} about 147 bp of DNA wound 1.65 times
  around the histone octamer in a left-handed superhelix
\item
  \textbf{Linker DNA:} 10--80 bp connecting adjacent nucleosomes;
  associated with histone H1
\item
  \textbf{Bead-on-a-string:} nucleosomes connected by linker DNA form an
  11-nm fibre --- the first level of compaction
\end{itemize}

\textbf{Core histone structure:} Each core histone has: - A
\textbf{globular domain} forming the nucleosome disc surface through the
\textbf{histone fold} motif (three \texorpdfstring{\ensuremath{\alpha}}{alpha}-helices linked by two short loops;
pairs of histones form four-helix bundles: H3--H4 and H2A--H2B) - An
unstructured \textbf{N-terminal tail} extending beyond the disc (4--35
residues depending on histone) --- this is where most post-translational
modifications occur - A short \textbf{C-terminal tail} (especially
prominent on H2A and H2B) that also accepts modifications including
ubiquitination

\textbf{Histone variants:} In addition to canonical H2A/H2B/H3/H4,
mammalian cells express variant histones that confer specialised
properties on chromatin where they are deposited:

{\def\LTcaptype{none} % do not increment counter
\begin{longtable}[]{@{}
  >{\raggedright\arraybackslash}p{(\linewidth - 6\tabcolsep) * \real{0.1628}}
  >{\raggedright\arraybackslash}p{(\linewidth - 6\tabcolsep) * \real{0.1860}}
  >{\raggedright\arraybackslash}p{(\linewidth - 6\tabcolsep) * \real{0.4651}}
  >{\raggedright\arraybackslash}p{(\linewidth - 6\tabcolsep) * \real{0.1860}}@{}}
\toprule\noalign{}
\begin{minipage}[b]{\linewidth}\raggedright
Variant
\end{minipage} & \begin{minipage}[b]{\linewidth}\raggedright
Replaces
\end{minipage} & \begin{minipage}[b]{\linewidth}\raggedright
Deposition machinery
\end{minipage} & \begin{minipage}[b]{\linewidth}\raggedright
Function
\end{minipage} \\
\midrule\noalign{}
\endhead
\bottomrule\noalign{}
\endlastfoot
H3.3 & H3 & HIRA (gene bodies); ATRX/DAXX (heterochromatin) &
Replication-independent; marks transcribed/active regions \\
CENP-A (CenH3) & H3 & HJURP & Centromere identity; foundation for
kinetochore assembly \\
H2A.Z & H2A & SRCAP, INO80/SWR1 & Promoter-proximal; poises genes for
activation; insulates from heterochromatin \\
H2A.X & H2A & RAD51 / S139ph at DSBs & DNA-damage signalling (\texorpdfstring{\ensuremath{\gamma}}{gamma}H2AX);
foci visible by immunofluorescence \\
macroH2A & H2A & ATRX-dependent & Inactive X enrichment; gene
silencing \\
H2A.B (Bbd) & H2A & unclear & Active transcription; testis enriched \\
\end{longtable}
}

\subsubsection{Higher-Order Chromatin
Compaction}\label{higher-order-chromatin-compaction}

{\def\LTcaptype{none} % do not increment counter
\begin{longtable}[]{@{}
  >{\raggedright\arraybackslash}p{(\linewidth - 8\tabcolsep) * \real{0.1042}}
  >{\raggedright\arraybackslash}p{(\linewidth - 8\tabcolsep) * \real{0.1875}}
  >{\raggedright\arraybackslash}p{(\linewidth - 8\tabcolsep) * \real{0.1667}}
  >{\raggedright\arraybackslash}p{(\linewidth - 8\tabcolsep) * \real{0.3542}}
  >{\raggedright\arraybackslash}p{(\linewidth - 8\tabcolsep) * \real{0.1875}}@{}}
\toprule\noalign{}
\begin{minipage}[b]{\linewidth}\raggedright
Level
\end{minipage} & \begin{minipage}[b]{\linewidth}\raggedright
Structure
\end{minipage} & \begin{minipage}[b]{\linewidth}\raggedright
Diameter
\end{minipage} & \begin{minipage}[b]{\linewidth}\raggedright
Compaction factor
\end{minipage} & \begin{minipage}[b]{\linewidth}\raggedright
Mechanism
\end{minipage} \\
\midrule\noalign{}
\endhead
\bottomrule\noalign{}
\endlastfoot
0 & Naked DNA (B-form helix) & 2 nm & 1\texorpdfstring{\ensuremath{\times}}{x} & Watson--Crick base-pairing \\
1 & Nucleosome fibre (``beads on a string'') & 11 nm & about 6\texorpdfstring{\ensuremath{\times}}{x}
& Histone octamer wrapping \\
2 & 30-nm fibre (disputed \emph{in vivo}) & 30 nm & about 40\texorpdfstring{\ensuremath{\times}}{x} &
Nucleosome--nucleosome compaction \\
3 & Chromatin loops / TADs & 300 nm & about 1,000\texorpdfstring{\ensuremath{\times}}{x} & CTCF +
cohesin-defined loops; TADs \\
4 & A/B compartments & --- & variable & Active (A) vs.~repressive (B)
genomic neighbourhoods \\
5 & Chromatid (mitotic chromosome) & 700 nm & about 10,000\texorpdfstring{\ensuremath{\times}}{x} &
SMC condensin-mediated compaction \\
\end{longtable}
}

\subsubsection{Topologically Associating Domains
(TADs)}\label{topologically-associating-domains-tads}

Chromatin organises into DNA loops of 100 kb--1 Mb delineated by
\textbf{CTCF} (insulator protein) binding sites and \textbf{cohesin}
ring complexes. Within TADs, enhancers preferentially contact
\hyperref[gl:promoter]{\textbf{promoter}}s of the same TAD. TAD
boundaries are largely conserved across cell types and species. The
dominant model for TAD formation is \textbf{loop extrusion}: cohesin
loads on chromatin and reels DNA through its ring lumen ATP-dependently
until it stalls at convergently oriented CTCF binding sites, leaving
behind a chromatin loop. Quantitatively, mammalian genomes contain
about 3,000--10,000 TADs (about 1 Mb median), and
knockdown of CTCF or cohesin (RAD21, NIPBL) blurs or abolishes most
boundaries within hours.

\begin{equation}P_{\text{contact}}(s) \propto s^{-\alpha}, \quad \alpha \approx 1.0\text{–}1.2 \text{ for fractal globule (interphase)} \tag{15.1}\label{eq:hic_scaling}\end{equation}

\textbf{Hi-C methodology in detail.} Hi-C is a
chromosome-conformation-capture variant that measures contact frequency
between every pair of genomic loci genome-wide. The protocol:

\begin{enumerate}
\def\labelenumi{\arabic{enumi}.}
\tightlist
\item
  \textbf{Cross-link} chromatin with formaldehyde (1 \% for 10 min),
  preserving protein-mediated DNA contacts.
\item
  \textbf{Restriction digest} with a 4-cutter (DpnII, MboI) or 6-cutter
  (HindIII), creating sticky-ended fragments still tethered by
  cross-linked proteins.
\item
  \textbf{Fill-in with biotinylated nucleotides} to mark cleaved ends,
  then \textbf{proximity ligate} at low DNA concentration so that
  intermolecular ligations are improbable --- primarily fragments held
  together by cross-linked proteins (i.e., physically proximal \emph{in
  vivo}) ligate.
\item
  \textbf{Streptavidin pulldown} enriches biotinylated junction reads,
  paired-end Illumina sequencing identifies which genomic loci were
  joined.
\item
  \textbf{Heatmap normalisation} (ICE, KR-balanced) yields a contact
  matrix; diagonal-rich blocks are TADs; off-diagonal stripes identify
  loop anchors at convergent CTCF.
\end{enumerate}

Derivative methods refine the assay: \textbf{Micro-C} uses MNase
digestion to nucleosome resolution; \textbf{Capture-Hi-C} enriches for
promoter-anchored contacts (4DN consortium); \textbf{HiChIP / PLAC-seq}
combines ChIP for an active mark (H3K27ac) with proximity ligation;
\textbf{ChIA-PET} uses paired-end tagging on antibody-immunoprecipitated
chromatin; \textbf{Hi-C 3.0 / Pore-C} uses long-read multi-way contact
mapping.

Disruption of TAD boundaries (e.g., by large deletions, inversions, or
single-nucleotide CTCF site disruptions) can cause misregulation of
developmental genes --- a mechanism termed \textbf{enhancer hijacking}.
Lupiáñez et al.~(2015) showed that microdeletions removing TAD
boundaries near \emph{EPHA4} fuse adjacent regulatory landscapes,
mis-targeting limb enhancers to \emph{PAX3} or \emph{WNT6} and producing
congenital limb malformations (brachydactyly, syndactyly).

\textbf{Loop extrusion --- the dynamic model.} Single-molecule imaging
and live-cell Hi-C now support a kinetic loop-extrusion model with
quantitative parameters: - Cohesin (SMC1/SMC3 + RAD21/SCC1 +
STAG1/STAG2) is loaded onto chromatin by \textbf{NIPBL/MAU2} (the loader
complex). - It extrudes DNA into a growing loop at about 0.5--2
kb/s, consuming ATP at the SMC ATPase head. - \textbf{WAPL} unloads
cohesin (residence time about 20 min unloaded by WAPL,
vs.~\textgreater{} 1 h for the \textbf{cohesin-STAG2} ``stalled'' form
at CTCF anchors). - CTCF binds asymmetrically (a 19-bp consensus with a
defined orientation): convergent CTCF pairs trap extruding cohesin,
generating loops with measured loop sizes about 100 kb to 2 Mb.
- Acute auxin-induced degron depletion of RAD21 dissolves most loops
within about 15 min; restoration takes about 30--60
min --- confirming that TADs and loops are dynamic, not static
structures.

\subsubsection{A/B Compartments}\label{ab-compartments}

At the megabase scale, the genome partitions into two interaction
``compartments'':

\begin{itemize}
\tightlist
\item
  \textbf{A compartment} --- gene-rich, active (H3K4me3, H3K27ac,
  accessible chromatin), positioned toward the nuclear interior
\item
  \textbf{B compartment} --- gene-poor, repressed (H3K9me3,
  lamina-associated), positioned at the nuclear periphery
\end{itemize}

Compartment switching during cell-fate transitions correlates with
replication timing changes (early-replicating \texorpdfstring{\ensuremath{\leftrightarrow}}{<->} late-replicating) and
with large-scale gene expression remodelling. Hi-C (and its derivatives
Micro-C, Capture-Hi-C, ChIA-PET) measures these contacts genome-wide.

\begin{quote}
\textbf{Concept Check 1:} A patient harbours a balanced inversion that
breaks a CTCF site at the boundary of a TAD containing the \emph{SHH}
limb enhancer (ZRS, about 1 Mb upstream of \emph{SHH}). Predict
the developmental phenotype if the inversion places ZRS into a TAD
containing an unrelated proto-oncogene. What
chromosome-conformation-capture experiment would confirm enhancer
hijacking?
\end{quote}

\begin{center}\rule{0.5\linewidth}{0.5pt}\end{center}

\subsection{Histone Modifications --- The Histone
Code}\label{histone-modifications-the-histone-code}

Histone tails are subject to a large number of post-translational
modifications (PTMs). The \textbf{\hyperref[gl:histone-code]{histone
code} hypothesis} \citep{strahl2000} proposes that specific combinations
of marks specify distinct chromatin states and downstream outcomes.

\begin{figure}[htbp]
\centering
\includegraphics[width=0.82\linewidth,height=4.2in,keepaspectratio]{../figures/mermaid_inline/inline_mermaid_0080_2472a9838bd0.png}
\caption{Histone modification writers and the states they produce. Green = active marks; blue = active promoters; purple = Polycomb repression; grey = constitutive heterochromatin.}
\end{figure}

\subsubsection{The Histone Code: A Comprehensive Reference
Table}\label{the-histone-code-a-comprehensive-reference-table}

The histone-code reference table below systematically catalogs
\textgreater{} 25 modifications recurrently studied in mammalian
chromatin biology, paired with the \textbf{writer enzyme},
\textbf{eraser enzyme}, \textbf{reader domain} that recognises the mark,
the \textbf{histone position} (nucleosome face, exposed N-terminal tail,
C-terminal tail), the typical \textbf{genomic location}, and the
\textbf{functional effect} (activation, repression, or
context-dependent). This is the chromatin biologist's periodic table ---
the working vocabulary of writers, readers, and erasers that defines
every targeted therapy in clinical epigenetics.

\textbf{Active methylation marks}

{\def\LTcaptype{none} % do not increment counter
\begin{longtable}[]{@{}
  >{\raggedright\arraybackslash}p{(\linewidth - 6\tabcolsep) * \real{0.1053}}
  >{\raggedright\arraybackslash}p{(\linewidth - 6\tabcolsep) * \real{0.3947}}
  >{\raggedright\arraybackslash}p{(\linewidth - 6\tabcolsep) * \real{0.3158}}
  >{\raggedright\arraybackslash}p{(\linewidth - 6\tabcolsep) * \real{0.1842}}@{}}
\toprule\noalign{}
\begin{minipage}[b]{\linewidth}\raggedright
Mark
\end{minipage} & \begin{minipage}[b]{\linewidth}\raggedright
Writer / eraser
\end{minipage} & \begin{minipage}[b]{\linewidth}\raggedright
Main readers
\end{minipage} & \begin{minipage}[b]{\linewidth}\raggedright
Meaning
\end{minipage} \\
\midrule\noalign{}
\endhead
\bottomrule\noalign{}
\endlastfoot
H3K4me1 & MLL3/4 writes; LSD1 erases & BAF45c PHD, CHD7 & Poised or
active enhancers \\
H3K4me2 & MLL1/2 and SET1 write; LSD1 erases & CHD1, BPTF PHD & Active
promoters and 5-prime gene bodies \\
H3K4me3 & MLL/SET1 COMPASS writes; KDM5 removes & TAF3 PHD, ING1, BPTF &
Active promoters; recruits TFIID \\
H3K36me1/2 & NSD and ASH1L write; KDM2 removes & LEDGF PWWP, MRG15 &
Gene bodies and active enhancers \\
H3K36me3 & SETD2 writes; KDM2/KDM4 remove & DNMT3 PWWP, MRG15, LEDGF &
Transcribed gene bodies; suppresses cryptic initiation and recruits
MMR \\
H3K79me2/3 & DOT1L writes & 53BP1 Tudor, AF9, ENL & Active gene bodies;
important in MLL-fusion leukaemia \\
\end{longtable}
}

\textbf{Active acetylation and ubiquitination marks}

{\def\LTcaptype{none} % do not increment counter
\begin{longtable}[]{@{}
  >{\raggedright\arraybackslash}p{(\linewidth - 6\tabcolsep) * \real{0.1053}}
  >{\raggedright\arraybackslash}p{(\linewidth - 6\tabcolsep) * \real{0.3947}}
  >{\raggedright\arraybackslash}p{(\linewidth - 6\tabcolsep) * \real{0.3158}}
  >{\raggedright\arraybackslash}p{(\linewidth - 6\tabcolsep) * \real{0.1842}}@{}}
\toprule\noalign{}
\begin{minipage}[b]{\linewidth}\raggedright
Mark
\end{minipage} & \begin{minipage}[b]{\linewidth}\raggedright
Writer / eraser
\end{minipage} & \begin{minipage}[b]{\linewidth}\raggedright
Main readers
\end{minipage} & \begin{minipage}[b]{\linewidth}\raggedright
Meaning
\end{minipage} \\
\midrule\noalign{}
\endhead
\bottomrule\noalign{}
\endlastfoot
H3K9ac & GCN5/PCAF/p300 write; HDAC1-3 and SIRT1/6 erase & BRD4, TAF1 &
Active promoters; loosens histone-DNA contacts \\
H3K14ac & GCN5/PCAF write; HDAC1-3 erase & Bromodomains & Active
promoters; cooperates with H3K4me3 \\
H3K18ac and H3K23ac & p300/CBP write; HDACs erase & Bromodomains &
Promoters and enhancers; active transcription \\
H3K27ac & p300/CBP write; HDAC1-3 erase & BRD4, YEATS & Active enhancers
and super-enhancers \\
H3K56ac & p300/CBP write; HDAC1/SIRT1/SIRT6 erase & no dominant reader &
Newly synthesised H3; chromatin assembly \\
H4K5ac, H4K8ac, H4K12ac & HAT1/p300/CBP write; HDACs erase &
Bromodomains & Replication-coupled H4 deposition \\
H4K16ac & MOF writes; SIRT1/2 erase & context-dependent & Disrupts 30-nm
fibre compaction; active gene bodies \\
H2BK120ub & RNF20/40 writes; USP22 erases & crosstalk readers & Required
for H3K4me3 and H3K79me deposition \\
\end{longtable}
}

\textbf{Repressive methylation and Polycomb marks}

{\def\LTcaptype{none} % do not increment counter
\begin{longtable}[]{@{}
  >{\raggedright\arraybackslash}p{(\linewidth - 6\tabcolsep) * \real{0.1053}}
  >{\raggedright\arraybackslash}p{(\linewidth - 6\tabcolsep) * \real{0.3947}}
  >{\raggedright\arraybackslash}p{(\linewidth - 6\tabcolsep) * \real{0.3158}}
  >{\raggedright\arraybackslash}p{(\linewidth - 6\tabcolsep) * \real{0.1842}}@{}}
\toprule\noalign{}
\begin{minipage}[b]{\linewidth}\raggedright
Mark
\end{minipage} & \begin{minipage}[b]{\linewidth}\raggedright
Writer / eraser
\end{minipage} & \begin{minipage}[b]{\linewidth}\raggedright
Main readers
\end{minipage} & \begin{minipage}[b]{\linewidth}\raggedright
Meaning
\end{minipage} \\
\midrule\noalign{}
\endhead
\bottomrule\noalign{}
\endlastfoot
H3K9me1 & G9a/GLP/SETDB1 write; KDM3 removes & weak / context-dependent
& Active gene bodies; cooperates with elongation marks \\
H3K9me2 & G9a/GLP write; KDM3/LSD1 remove & HP1, CDYL & Facultative
heterochromatin and silenced euchromatin \\
H3K9me3 & SUV39H and SETDB1 write; KDM4 removes & HP1 alpha/beta/gamma,
CBX1/3/5 & Constitutive heterochromatin; centromeres; transposons \\
H3K27me1/2 & EZH2 writes; UTX/JMJD3 remove & EED for me2/3 contexts &
Active gene bodies at low level; me2 prevents inappropriate enhancer
activation \\
H3K27me3 & EZH2/PRC2 writes; UTX/JMJD3 remove & EED and CBX-PRC1 &
Polycomb domains; developmental-gene repression \\
H4K20me3 & SUV4-20H writes; PHF8 erases & 53BP1 Tudor & Pericentric
heterochromatin and DNA-damage response \\
H2AK119ub & RNF2/RING1B writes; BAP1 removes & JARID2-associated PRC2.2
& Stable Polycomb silencing layer \\
\end{longtable}
}

\textbf{Dynamic damage, cell-cycle, and metabolic marks}

{\def\LTcaptype{none} % do not increment counter
\begin{longtable}[]{@{}
  >{\raggedright\arraybackslash}p{(\linewidth - 6\tabcolsep) * \real{0.1053}}
  >{\raggedright\arraybackslash}p{(\linewidth - 6\tabcolsep) * \real{0.3947}}
  >{\raggedright\arraybackslash}p{(\linewidth - 6\tabcolsep) * \real{0.3158}}
  >{\raggedright\arraybackslash}p{(\linewidth - 6\tabcolsep) * \real{0.1842}}@{}}
\toprule\noalign{}
\begin{minipage}[b]{\linewidth}\raggedright
Mark
\end{minipage} & \begin{minipage}[b]{\linewidth}\raggedright
Writer / eraser
\end{minipage} & \begin{minipage}[b]{\linewidth}\raggedright
Main readers
\end{minipage} & \begin{minipage}[b]{\linewidth}\raggedright
Meaning
\end{minipage} \\
\midrule\noalign{}
\endhead
\bottomrule\noalign{}
\endlastfoot
H3S10ph & Aurora B/MSK write; PP1/PP2A erase & 14-3-3 proteins & Mitosis
and immediate-early gene activation \\
H3T3ph & Haspin writes; PP1 erases & Survivin/CPC & Inner-centromere
recruitment during mitosis \\
H2A.X-S139ph (gamma-H2AX) & ATM/ATR/DNA-PK write; PP2A/WIP1 erase & MDC1
BRCT & Double-strand-break signalling foci \\
H3T6ph & PRK1 writes; PP1 erases & context-dependent & Androgen-receptor
targets; H3K4 demethylation crosstalk \\
H4K12su and H2BK34su & PIAS/MMS21 write; SENP erases & SIM-containing
readers & Polycomb and DNA-damage coordination \\
H2A/H4 ADP-ribosylation & PARP1/2 write; PARG/ARH3 erase & Macrodomain
proteins & Chromatin relaxation for DNA repair \\
H3K18la & p300 writes; HDAC1-3 erase & YEATS domains & Lactate-linked
macrophage polarisation \\
H3K9bhb & p300 writes; HDAC1-3 erase & YEATS domains & Ketosis and
fasting response \\
H3K9cr, H3K9su, H3K4but & p300 or metabolic enzymes write; HDAC/sirtuin
systems erase & YEATS-family readers & Metabolism-linked activation in
enhancers, promoters, and gametogenesis \\
\end{longtable}
}

The bottom rows reflect a major recent expansion of the histone code:
short-chain acyl-CoAs derived from intermediary metabolism (lactate,
\texorpdfstring{\ensuremath{\beta}}{beta}-hydroxybutyrate, crotonyl-CoA, succinyl-CoA, butyryl-CoA) are
deposited by the same KAT/HAT enzymes (especially p300) that write
acetyl marks, directly coupling \textbf{metabolic state} to chromatin
state. The lactyl mark, in particular, has become the molecular
signature of the Warburg-effect--polarised tumour-associated macrophage.

\subsubsection{Histone Acetylation}\label{histone-acetylation}

\textbf{Writers:} Histone acetyltransferases (HATs) --- CBP/p300,
GCN5/PCAF, TIP60, MOF. \textbf{Erasers:} Histone deacetylases (HDACs)
--- 18 human HDACs in four classes (Class I/II/III/IV, where Class III =
NAD\texorpdfstring{\textsuperscript{+}}{+}-dependent sirtuins SIRT1--7). \textbf{Readers:} Bromodomains (e.g.,
BRD4, TFIID subunit TAF1, BAF180 of PBAF, BPTF). YEATS domains (AF9,
ENL) read both acetyl and longer acyl marks.

\textbf{Mechanism:} Acetylation of lysine \texorpdfstring{\ensuremath{\varepsilon}}{epsilon}-amino groups by transfer
from acetyl-CoA. Neutralises the positive charge on lysine,
\textbf{weakening histone--DNA electrostatic interactions}, loosening
chromatin and facilitating
\hyperref[gl:transcription]{\textbf{transcription}} factor binding.

\begin{equation}\text{Lys-NH}_3^+ + \text{acetyl-CoA} \xrightarrow{\text{HAT}} \text{Lys-NH-COCH}_3 + \text{CoA-SH} + \text{H}^+ \tag{15.2}\label{eq:hat}\end{equation}

\textbf{Key active marks:} - \textbf{H3K27ac:} Marks active enhancers
(distinguishes active from poised enhancers, which carry H3K27me3) -
\textbf{H3K9ac, H3K14ac, H4K16ac:} Found at active promoters and
transcribed gene bodies

\begin{quote}
\textbf{Clinical Connection --- HDAC inhibitors (full clinical
profile):} Histone deacetylase inhibitors (HDACi) are FDA-approved
anticancer drugs. - \textbf{Vorinostat (SAHA)} --- FDA approved 2006 for
cutaneous T-cell lymphoma (CTCL). Hydroxamate pan-HDACi. Median time to
response about 2 months; objective response about 30
\%. - \textbf{Romidepsin (Istodax)} --- FDA approved 2009 for CTCL and
2011 for peripheral T-cell lymphoma. Cyclic depsipeptide; class-I
selective. Distinct dose-limiting toxicities (GI, fatigue,
thrombocytopenia). - \textbf{Panobinostat (Farydak)} --- FDA approved
2015 for relapsed/refractory multiple myeloma in combination with
bortezomib + dexamethasone. Pan-HDACi; superior CD138+ plasma-cell
apoptosis when combined with proteasome inhibitor. - \textbf{Belinostat
(Beleodaq)} --- FDA 2014 for PTCL. - \textbf{Tucidinostat (Chidamide)}
--- China-approved 2014 for PTCL; first benzamide HDACi.

Mechanism in cancer: re-activation of silenced tumour-suppressor genes
(p21, gelsolin, RhoB, BCL6); accumulation of acetylation on non-histone
targets (HSP90 chaperone client release, p53 stabilisation, NF-\texorpdfstring{\ensuremath{\kappa}}{kappa}B
inhibition). Combination with \textbf{DNMT inhibitors} is being explored
to reverse layered silencing, a theme developed later in the
cancer-epigenetics discussion.
\end{quote}

\subsubsection{Histone Methylation}\label{histone-methylation}

\textbf{Writers:} Histone methyltransferases (HMTs) --- most contain SET
domains (except DOT1L, which has a 7\texorpdfstring{\ensuremath{\beta}}{beta}-strand methyltransferase fold
derived from class V methyltransferases). \textbf{Erasers:} Histone
demethylases --- KDM1/LSD1 (FAD-dependent; works on me1 and me2
primarily; cannot demethylate me3 because it requires a free lysine
\texorpdfstring{\ensuremath{\varepsilon}}{epsilon}-amino lone pair); KDM2--7 (Jumonji-domain, Fe(II)/\texorpdfstring{\ensuremath{\alpha}}{alpha}KG-dependent; can
remove me3 by hydroxymethyl-amine intermediate that decomposes to
formaldehyde and demethylated lysine). \textbf{Readers:} Chromodomains
(HP1 reads H3K9me3; CBX2/4/6/7/8 of PRC1 read H3K27me3); PHD fingers
(TAF3, BPTF, ING proteins read H3K4me3); Tudor domains (53BP1 reads
H4K20me2; SMN reads symmetrical R-methylation); MBT domains; WD40 (EED
reads H3K27me3); PWWP domains (read H3K36me3).

Unlike acetylation (typically activating), \textbf{methylation is
context-dependent}: the same modification on different residues has
opposite effects.

{\def\LTcaptype{none} % do not increment counter
\begin{longtable}[]{@{}
  >{\raggedright\arraybackslash}p{(\linewidth - 6\tabcolsep) * \real{0.1739}}
  >{\raggedright\arraybackslash}p{(\linewidth - 6\tabcolsep) * \real{0.2609}}
  >{\raggedright\arraybackslash}p{(\linewidth - 6\tabcolsep) * \real{0.2174}}
  >{\raggedright\arraybackslash}p{(\linewidth - 6\tabcolsep) * \real{0.3478}}@{}}
\toprule\noalign{}
\begin{minipage}[b]{\linewidth}\raggedright
Mark
\end{minipage} & \begin{minipage}[b]{\linewidth}\raggedright
Enzyme
\end{minipage} & \begin{minipage}[b]{\linewidth}\raggedright
State
\end{minipage} & \begin{minipage}[b]{\linewidth}\raggedright
Function
\end{minipage} \\
\midrule\noalign{}
\endhead
\bottomrule\noalign{}
\endlastfoot
H3K4me1 & MLL3/4 & Active & Enhancer elements (poised or active) \\
H3K4me3 & MLL1/SET1 & Active & Active promoters; read by TAF3 PHD \\
H3K36me3 & SETD2 & Active & Transcribed gene bodies; prevents cryptic
initiation \\
H3K27me3 & EZH2 (PRC2) & Repressive & Polycomb targets; poised/silent
developmental genes \\
H3K9me3 & SUV39H1/H2 & Repressive & Constitutive heterochromatin;
\hyperref[gl:centromere]{\textbf{centromere}}s; transposons \\
H4K20me3 & SUV4-20H & Repressive & Pericentric heterochromatin; DNA
damage response \\
\end{longtable}
}

\textbf{Bivalent domains:} In embryonic stem cells, many developmental
gene promoters carry \textbf{both H3K4me3 and H3K27me3} simultaneously
--- a ``bivalent'' state. These genes are poised for either rapid
activation (upon H3K27me3 removal by KDM6A/UTX) or repression (upon
H3K4me3 removal) during differentiation. About 2,500 promoters are
bivalent in human ES cells; most resolve to monovalent during lineage
commitment.

\subsubsection{Phosphorylation and
Ubiquitination}\label{phosphorylation-and-ubiquitination}

\textbf{H3S10ph (H3 Ser10 phosphorylation):} - Writer: Aurora B kinase
(\hyperref[gl:mitosis]{\textbf{mitosis}}), MSK1/2 (mitogenic signalling)
- Function: Chromosome condensation during mitosis; gene activation by
14-3-3 reader binding

\textbf{H2AX-S139ph (\texorpdfstring{\ensuremath{\gamma}}{gamma}H2AX):} - Writer: ATM/ATR kinases at double-strand
breaks - Function: Marks DSB sites; recruits MDC1 and DNA repair
machinery; visible as bright immunofluorescence foci at break sites

\textbf{H2AK119ub (H2A monoubiquitination):} - Writer: RNF2/RING1B (part
of PRC1 complex) - Function: Represses transcription elongation; part of
Polycomb silencing; read by Polycomb-like proteins

\begin{quote}
\textbf{Concept Check 2:} A cancer-associated EZH2 gain-of-function
mutation leads to \hyperref[gl:genome]{\textbf{genome}}-wide
hypermethylation of H3K27 and silencing of tumour-suppressor genes. To
what extent would an HDAC inhibitor (which increases histone
acetylation) be expected to re-activate these silenced loci? Explain by
reference to the dependencies between acetylation, methylation, and DNA
methylation described above --- and argue for a rational combination
therapy.
\end{quote}

\begin{quote}
\textbf{Concept Check 3:} Predict the genome-wide consequence of a
homozygous knockout of \textbf{SETD2} (the H3K36me3 writer) in an
embryonic stem cell. Consider effects on (i) cryptic intragenic
transcription, (ii) DNA mismatch repair recruitment, and (iii)
DNMT3B-mediated gene-body methylation. Which of these phenotypes most
directly explains why \emph{SETD2} loss-of-function is the most frequent
driver mutation in clear-cell renal cell carcinoma after VHL?
\end{quote}

\begin{center}\rule{0.5\linewidth}{0.5pt}\end{center}

\subsection{Polycomb and Trithorax --- Cellular Memory
Systems}\label{polycomb-and-trithorax-cellular-memory-systems}

The \textbf{Polycomb (PcG)} and \textbf{Trithorax (TrxG)} group proteins
were discovered in \emph{Drosophila} as gene families whose
loss-of-function alleles cause homeotic transformations of body
segments. They constitute the cell's principal \textbf{cellular memory
systems}: once a developmental gene is set ON or OFF in a progenitor,
PcG/TrxG complexes propagate that state through most subsequent mitotic
divisions of the lineage.

\subsubsection{Polycomb Repressive Complexes --- Structural
Detail}\label{polycomb-repressive-complexes-structural-detail}

\begin{figure}[htbp]
\centering
\includegraphics[width=0.82\linewidth,height=4.2in,keepaspectratio]{../figures/mermaid_inline/inline_mermaid_0081_44e01d6a348a.png}
\caption{Polycomb silencing operates as a feedback-amplified two-mark system: PRC2 writes H3K27me3, which recruits more PRC2 (positive feedback) and recruits canonical PRC1 (via CBX) which writes H2AK119ub. Variant PRC1 acts upstream at unmethylated CpG islands and seeds H2AK119ub which itself recruits PRC2. The mutual reinforcement explains how Polycomb domains span hundreds of kilobases and persist through cell divisions.}
\end{figure}

\textbf{PRC1 --- comprehensive subunit inventory:} - Catalytic core:
\textbf{RING1A} (RNF1) or \textbf{RING1B} (RNF2) E3 ubiquitin ligase + a
PCGF partner (PCGF1, PCGF2/MEL18, PCGF3, PCGF4/BMI1, PCGF5, PCGF6). -
\textbf{Canonical PRC1 (cPRC1):} RING1A/B + PCGF2/4 (MEL18/BMI1) + a
\textbf{CBX} (CBX2/4/6/7/8) chromobox subunit + PHC1/2/3 + SCMH1/L2. The
CBX chromodomain reads H3K27me3, anchoring cPRC1 onto Polycomb-marked
chromatin. PHC SAM-domain polymerisation drives \textbf{chromatin
compaction in cis} (head-to-tail oligomerisation creates a
phase-separated Polycomb body). - \textbf{Variant PRC1 (vPRC1):}
RING1A/B + PCGF1/3/5/6 + RYBP (or YAF2) instead of CBX. Recruited to
\textbf{unmethylated CpG islands} via KDM2B (PCGF1 complex,
``ncPRC1.1'') or via E2F6 (PCGF6). Deposits H2AK119ub \emph{upstream} of
H3K27me3 --- a key insight: vPRC1 acts first, then PRC2.2 reads
H2AK119ub via JARID2.

\textbf{PRC2 --- comprehensive subunit inventory:} - Catalytic core:
\textbf{EZH2} (or paralog \textbf{EZH1}; SET-domain methyltransferase).
EZH2 is faster but EZH1 is dominant in non-dividing cells. -
\textbf{EED} binds H3K27me3 product \texorpdfstring{\ensuremath{\to}}{->} allosteric activation of EZH2 (the
\textbf{read-write feedback} that generates broad Polycomb domains). -
\textbf{SUZ12} scaffold links the SET-domain catalytic core to the
DNA/RNA-recognition modules. - \textbf{RBBP4/7} histone chaperone
presents the H3 tail substrate. - \textbf{Accessory subunits define
context-dependent variants:} - \textbf{PRC2.1:} PCL1/2/3 (PHF1, MTF2,
PHF19) + EPOP or PALI1. The PHF Tudor domains read H3K36me3, restricting
PRC2.1 from active gene bodies. - \textbf{PRC2.2:} AEBP2 + JARID2.
JARID2 reads H2AK119ub deposited by variant PRC1 and stimulates PRC2
activity (closing the recruitment loop). - \textbf{Reaction kinetics:}
H3K27 \texorpdfstring{\ensuremath{\to}}{->} H3K27me1 (k\texorpdfstring{\textsubscript{1}}{1} \texorpdfstring{\ensuremath{\approx}}{~} 0.05 s\texorpdfstring{\textsuperscript{-}}{-}\texorpdfstring{\textsuperscript{1}}{1}) \texorpdfstring{\ensuremath{\to}}{->} H3K27me2 (k\texorpdfstring{\textsubscript{2}}{2} \texorpdfstring{\ensuremath{\approx}}{~} 0.02 s\texorpdfstring{\textsuperscript{-}}{-}\texorpdfstring{\textsuperscript{1}}{1}) \texorpdfstring{\ensuremath{\to}}{->} H3K27me3
(k\texorpdfstring{\textsubscript{3}}{3} \texorpdfstring{\ensuremath{\approx}}{~} 0.005 s\texorpdfstring{\textsuperscript{-}}{-}\texorpdfstring{\textsuperscript{1}}{1}), with each step about 3-fold slower;
EED-allosteric stimulation increases k\texorpdfstring{\textsubscript{3}}{3} by about 6-fold on
neighbouring nucleosomes carrying me3. - \textbf{Recruitment hierarchy:}
unmethylated CpG islands (via PRC2.2/JARID2 reading H2AK119ub from
variant PRC1, \emph{or} via PRC2.1/PCL with KDM2B); broad gene-body
recognition by EED reading existing me3; lncRNA recruitment (HOTAIR,
XIST repA, KCNQ1OT1, ANRIL).

\textbf{Mechanism of Polycomb--Trithorax antagonism:} PRC2.1 cannot act
on H3K36-methylated chromatin (PCL Tudor reads H3K36me3 antagonistically
--- the CPL paradox). Conversely, \textbf{ASH1L} (the Trithorax-group
H3K36 methyltransferase) deposits H3K36me2 at active gene bodies,
blocking PRC2 spreading. The cell uses this as a positive/negative
selection: H3K36me-active = Polycomb-restricted; H3K36me-absent =
Polycomb-permissive.

\subsubsection{Trithorax Group (TrxG) --- Antagonising
Polycomb}\label{trithorax-group-trxg-antagonising-polycomb}

{\def\LTcaptype{none} % do not increment counter
\begin{longtable}[]{@{}
  >{\raggedright\arraybackslash}p{(\linewidth - 4\tabcolsep) * \real{0.5000}}
  >{\raggedright\arraybackslash}p{(\linewidth - 4\tabcolsep) * \real{0.3333}}
  >{\raggedright\arraybackslash}p{(\linewidth - 4\tabcolsep) * \real{0.1667}}@{}}
\toprule\noalign{}
\begin{minipage}[b]{\linewidth}\raggedright
TrxG complex
\end{minipage} & \begin{minipage}[b]{\linewidth}\raggedright
Activity
\end{minipage} & \begin{minipage}[b]{\linewidth}\raggedright
Role
\end{minipage} \\
\midrule\noalign{}
\endhead
\bottomrule\noalign{}
\endlastfoot
MLL/COMPASS family (MLL1--4, SET1A/B) & H3K4 methyltransferase & Marks
active and poised promoters/enhancers \\
ASH1L & H3K36 methyltransferase & Antagonises Polycomb spread into
active genes \\
KDM6A/UTX, KDM6B/JMJD3 & H3K27me3 demethylases & Remove Polycomb marks
during differentiation \\
SWI/SNF (BAF) & ATP-dependent remodelling & Evicts PRC1/2; opens
chromatin \\
\end{longtable}
}

\textbf{The COMPASS family --- six MLL/SET1 complexes in mammals:} -
\textbf{MLL1/KMT2A} (mixed-lineage leukaemia 1) --- large promoter
H3K4me3 deposition; MLL1-AF4/AF9/ENL fusions cause infant acute
lymphoblastic or acute myeloid leukaemia. - \textbf{MLL2/KMT2B} ---
closely related to MLL1; redundant at most loci; selective at others. -
\textbf{MLL3/KMT2C} and \textbf{MLL4/KMT2D} --- write enhancer H3K4me1
(poised and active enhancers); KMT2D is among the most frequently
mutated chromatin-regulator genes in cancer (about 10 \% of
B-cell lymphoma and bladder cancer). - \textbf{SET1A/KMT2F} and
\textbf{SET1B/KMT2G} --- write the bulk of promoter H3K4me3 in adult
tissues; SET1A loss-of-function in MDS and CHIP.

Most six complexes share a \textbf{WRAD module}
(WDR5--RbBP5--ASH2L--DPY30) that allosterically activates the SET-domain
catalytic subunit by about 600-fold. This makes WRAD a
tractable allosteric drug target --- small-molecule WDR5--MLL interface
inhibitors (OICR-9429) are in early-phase trials for MLL-rearranged
leukaemia.

\textbf{The Polycomb--Trithorax balance} is maintained dynamically.
Differentiation cues that activate gene expression typically deploy
three concurrent steps: (i) UTX/JMJD3 removes H3K27me3, (ii) MLL/COMPASS
deposits H3K4me3, and (iii) BAF complexes evict residual PRC1.
Conversely, lineage repression deposits PRC1/PRC2 marks at the relevant
promoter.

\textbf{Clinical translation --- EZH2 inhibitors:} -
\textbf{Tazemetostat (Tazverik)} is a small-molecule SET-domain
inhibitor of EZH2, FDA-approved (2020) for \textbf{epithelioid sarcoma}
(which loses INI1/SMARCB1 --- a SWI/SNF subunit --- and becomes
hyper-dependent on EZH2 silencing) and \textbf{EZH2-mutant follicular
lymphoma}. Mechanism: SAM-competitive small molecule occupies the EZH2
SET-domain pocket. Approved dose 800 mg BID PO. Median PFS in INI1-loss
epithelioid sarcoma trial: 5.5 months vs 1.9 placebo; ORR 15 \%. -
\textbf{Valemetostat} (EZH1 + EZH2 dual inhibitor) was approved in 2022
in Japan for adult T-cell leukaemia/lymphoma (ATL); ORR 48 \%. -
\textbf{Ivosidenib + tazemetostat} combinations in IDH-mutant glioma are
in trial, exploiting the IDH-mutant tumour's dependence on
Polycomb-driven differentiation block.

The mechanism is paradigmatic of \textbf{synthetic lethality} in
chromatin: tumours with SWI/SNF loss become dependent on Polycomb
silencing of cell-cycle inhibitors, so EZH2 inhibition selectively
re-activates suppressors primarily in the cancer cells.

\begin{quote}
\textbf{Worked Example 1 --- Polycomb Spreading Model with Read-Write
Feedback}

\textbf{Setup:} A naive promoter has 2 nucleosomes carrying H3K27me3
(call this ``seed''), recruiting PRC2 with EED-driven allosteric
activation. PRC2 deposits H3K27me3 on the next nucleosome with rate
constant \emph{k} per generation. Each cell division dilutes existing
H3K27me3 by half (passive demethylation through nucleosome turnover,
since newly deposited histones at the replication fork are
unmethylated). The gene body comprises 50 nucleosomes.

\textbf{Question:} Derive the steady-state H3K27me3 density along the
gene body for \emph{k} = 1, 2, and 5 marks per generation. What is the
minimum \emph{k} needed for the silencing mark to spread fully (\texorpdfstring{\textgreater{}=}{>=} 90 \%
of nucleosomes methylated)?

\textbf{Solution:} Let \(m_n\) be the average number of H3K27me3-marked
nucleosomes at generation \(n\), with \(m_0 = 2\) (seed). Each
generation: - Half of existing marks are passively diluted:
\(m_{\text{post-replication}} = m_n / 2\). - PRC2 deposits \emph{k} new
marks at neighbouring nucleosomes (read-write spreading):
\(m_{n+1} = m_n / 2 + k\).

Steady-state when \(m_{\infty} = m_{\infty} / 2 + k\), so
\(m_{\infty} = 2k\) marked nucleosomes.

{\def\LTcaptype{none} % do not increment counter
\begin{longtable}[]{@{}lll@{}}
\toprule\noalign{}
\emph{k} (marks/gen) & Steady-state \(m_\infty\) & Gene-body coverage \\
\midrule\noalign{}
\endhead
\bottomrule\noalign{}
\endlastfoot
1 & 2 & 4 \% (silencing fails) \\
2 & 4 & 8 \% \\
5 & 10 & 20 \% \\
25 & 50 & 100 \% (saturated) \\
45 & 90 & 100 \% (rapid saturation) \\
\end{longtable}
}

\textbf{Insight:} For full coverage of a 50-nucleosome gene body (\texorpdfstring{\textgreater{}=}{>=} 45
marks at steady state), \emph{k} must exceed about 22 marks per
generation --- i.e., PRC2 must deposit \textgreater{} 11\texorpdfstring{\ensuremath{\times}}{x} the seed
quantity each cell cycle. This is achievable primarily with the
\textbf{EED--allosteric read-write loop}: each existing H3K27me3
catalytically recruits more PRC2 to neighbouring nucleosomes, creating
positive feedback. Without this loop (e.g., in EED-mutant cells),
Polycomb domains collapse over a few divisions. This is why
\textbf{EED-binding small molecules (EED226, MAK683)} are emerging as
alternative EZH2-pathway inhibitors that disrupt the feedback loop
selectively.
\end{quote}

\begin{center}\rule{0.5\linewidth}{0.5pt}\end{center}

\subsection{DNA Methylation}\label{dna-methylation}

\subsubsection{CpG Methylation
Mechanism}\label{cpg-methylation-mechanism}

In mammals, DNA methylation occurs almost exclusively on cytosine in the
context \textbf{5\texorpdfstring{'}{'}-CpG-3\texorpdfstring{'}{'}} dinucleotides (the ``p'' denotes the
phosphodiester bond between C and G). The methylated form is
\hyperref[gl:5-methylcytosine]{\textbf{5-methylcytosine (5mC)}}, often
called the ``fifth base.''

\begin{equation}\text{Cytosine} + \text{S-adenosylmethionine} \xrightarrow{\text{DNMT}} \text{5-methylcytosine} + \text{S-adenosylhomocysteine} \tag{15.3}\label{eq:DNMT}\end{equation}

\subsubsection{The DNA Methylation Machinery --- In
Depth}\label{the-dna-methylation-machinery-in-depth}

\textbf{Maintenance methyltransferase --- DNMT1:} the workhorse enzyme
that ensures heritability of CpG methylation across cell divisions. -
DNMT1 has a two-domain architecture: an N-terminal regulatory domain
(containing CXXC, BAH1/2, and replication-foci targeting RFTS
subdomains) and a C-terminal catalytic domain. - \textbf{Substrate
specificity:} preferentially methylates \textbf{hemi-methylated CpGs}
--- sites where the parental strand carries 5mC but the newly
synthesised daughter strand is unmodified. Affinity for hemi-methylated
CpGs is about 10--40-fold higher than for unmethylated CpGs. -
\textbf{Recruitment --- the UHRF1--DNMT1 axis:} the protein
\textbf{UHRF1} (``ubiquitin-like with PHD and RING finger domains 1'')
is the molecular bridge that targets DNMT1 to replicating
heterochromatin. UHRF1 has \textbf{five domains}: UBL (ubiquitin-like),
TTD (Tudor --- reads H3K9me2/3), PHD (reads unmethylated H3R2), SRA
(reads hemi-methylated CpG via base-flipping), and RING (E3 ubiquitin
ligase). The dual recognition of (a) hemi-methylated CpG (via SRA) AND
(b) H3K9me2/3 (via TTD) ensures that DNMT1 propagates methylation
primarily at chromatin states already marked for repression. -
\textbf{The mechanism cycle (per CpG):} 1. UHRF1 SRA domain flips the
5mC out of the helix, exposing the unmethylated daughter cytosine. 2.
UHRF1 RING ubiquitinates H3K18 (H3 tail) --- creating a docking site for
DNMT1's RFTS domain. 3. DNMT1 RFTS binds H3K18ub; DNMT1 catalytic domain
transfers a methyl group from SAM to the daughter cytosine. 4. SAH
(S-adenosyl-homocysteine) released; UHRF1 dissociates; cycle repeats. -
\textbf{Quantitative parameters:} maintenance efficiency per CpG \texorpdfstring{\ensuremath{\varepsilon}}{epsilon} \texorpdfstring{\ensuremath{\approx}}{~}
0.95 in healthy cells. Each replication fork carries about 50
CpGs/s of unmethylated daughter DNA into the lumen --- DNMT1 must fix
them within minutes before chromatin reassembly. -
\textbf{Loss-of-function:} mouse Dnmt1-knockout embryos die at E9.5 with
genome-wide demethylation, ectopic gene expression, and replication-fork
instability. Conditional knockouts in tissue-specific contexts produce
lineage-specific failures (haematopoietic stem-cell exhaustion;
intestinal epithelial collapse).

\textbf{De novo methyltransferases --- DNMT3A, DNMT3B:} - Both have an
N-terminal regulatory PWWP domain (reads H3K36me3) and an ADD domain
(reads unmethylated H3K4 = H3K4me0). PWWP recognition couples \emph{de
novo} methylation deposition to gene-body H3K36me3-marked regions; ADD
ensures methylation deposition is excluded from H3K4me3-marked active
promoters (a fundamental safeguard against silencing active genes). -
Establish \emph{new} methylation patterns during germline development,
gametogenesis, and early embryogenesis. - Both require \textbf{DNMT3L}
(catalytically inactive; missing the SET-like motif required for SAM
binding) as a regulatory partner that stimulates activity
about 15-fold and itself reads H3K4me0 --- coupling methylation
deposition to \emph{absence} of active marks. - DNMT3A: predominant in
oocyte, embryonic stem cells, haematopoietic stem cells;
loss-of-function mutations are frequent in \textbf{clonal
haematopoiesis} (about 5 \% of healthy individuals over age 70)
and AML (about 20 \% of cytogenetically normal AML); also
Tatton-Brown--Rahman syndrome (germline DNMT3A loss with overgrowth). -
DNMT3B: predominant in B cells, embryonic stem cells; loss-of-function
causes \textbf{ICF syndrome} (immunodeficiency, centromeric instability,
facial anomalies) due to demethylation of pericentromeric satellites.

\textbf{Demethylation --- TET family enzymes --- the active
demethylation pathway:} - TET1, TET2, TET3 are
Fe(II)/\texorpdfstring{\ensuremath{\alpha}}{alpha}-ketoglutarate-dependent dioxygenases (each about 2,000
amino acids) with conserved C-terminal catalytic domains. They
sequentially oxidise:

\begin{equation}5\text{mC} \xrightarrow{\text{TET}} 5\text{hmC} \xrightarrow{\text{TET}} 5\text{fC} \xrightarrow{\text{TET}} 5\text{caC} \xrightarrow{\text{TDG/BER}} \text{C} \tag{15.4}\label{eq:tet_oxidation}\end{equation}

\begin{itemize}
\tightlist
\item
  \textbf{Reaction details:} TET enzymes consume O\texorpdfstring{\textsubscript{2}}{2} and \texorpdfstring{\ensuremath{\alpha}}{alpha}KG, producing
  CO\texorpdfstring{\textsubscript{2}}{2} and succinate. Each oxidation step transfers an oxygen atom from
  O\texorpdfstring{\textsubscript{2}}{2} to the methyl carbon of cytosine.
\item
  \textbf{Active demethylation by base excision repair:}
  5-formylcytosine (5fC) and 5-carboxylcytosine (5caC) are recognised
  and excised by \textbf{thymine DNA glycosylase (TDG)}, which has a
  100-fold preference for 5fC/5caC over 5mC. TDG cleaves the
  N-glycosidic bond, leaving an abasic (AP) site. \textbf{APE1} then
  makes a single-nucleotide cut; \textbf{DNA Pol \texorpdfstring{\ensuremath{\beta}}{beta}} fills the gap with
  unmodified cytosine; \textbf{DNA ligase III + XRCC1} seal the nick.
  This is the canonical BER pathway adapted for methyl-cytosine erasure.
\item
  \textbf{5hmC is itself a stable mark with distinct biology, not merely
  an intermediate.} Different cell types have dramatically different
  5hmC/5mC ratios --- embryonic stem cells about 0.1 \%,
  neurons about 0.7 \% (Purkinje cells reach about 40
  \% of modified cytosines as 5hmC), cardiomyocytes about 0.5
  \%, tumour cells often \textless{} 0.05 \% (a \textbf{5hmC depletion
  signature} now used in liquid-biopsy diagnostics on platforms such as
  EpiCheck and GRAIL Galleri).
\item
  \textbf{Tissue-specific TET expression:} TET1 in stem cells (regulates
  pluripotency); TET2 in haematopoietic cells (clonal haematopoiesis);
  TET3 in zygote and neurons (post-fertilisation paternal demethylation;
  neural plasticity).
\item
  \textbf{Clinical relevance:} TET2 is the most frequently mutated
  epigenetic-writer gene in haematopoietic cancers --- loss-of-function
  TET2 mutations are present in about 15 \% of AML, 50 \% of
  CMML, and define \textbf{clonal haematopoiesis of indeterminate
  potential (CHIP)} (about 10 \% of people over age 70). The
  oncometabolite \textbf{2-hydroxyglutarate} (produced by mutant IDH1/2)
  competitively inhibits TET (and other \texorpdfstring{\ensuremath{\alpha}}{alpha}KG-dependent dioxygenases),
  causing genome-wide hypermethylation in IDH-mutant gliomas and AML
  (the \textbf{CIMP} phenotype).
\end{itemize}

\textbf{Methodological caveat:} Bisulphite sequencing (the traditional
5mC gold standard) \textbf{cannot distinguish 5mC from 5hmC} because
both resist bisulphite-induced deamination. Modern methods that resolve
oxidation states:

{\def\LTcaptype{none} % do not increment counter
\begin{longtable}[]{@{}
  >{\raggedright\arraybackslash}p{(\linewidth - 6\tabcolsep) * \real{0.1714}}
  >{\raggedright\arraybackslash}p{(\linewidth - 6\tabcolsep) * \real{0.1714}}
  >{\raggedright\arraybackslash}p{(\linewidth - 6\tabcolsep) * \real{0.2857}}
  >{\raggedright\arraybackslash}p{(\linewidth - 6\tabcolsep) * \real{0.3714}}@{}}
\toprule\noalign{}
\begin{minipage}[b]{\linewidth}\raggedright
Method
\end{minipage} & \begin{minipage}[b]{\linewidth}\raggedright
Signal
\end{minipage} & \begin{minipage}[b]{\linewidth}\raggedright
Resolution
\end{minipage} & \begin{minipage}[b]{\linewidth}\raggedright
Distinguishes
\end{minipage} \\
\midrule\noalign{}
\endhead
\bottomrule\noalign{}
\endlastfoot
Bisulphite-seq (BS-seq) & C\texorpdfstring{\ensuremath{\to}}{->}T primarily at unmethylated & Single-base &
5mC + 5hmC vs C \\
oxBS-seq & Chemical oxidation of 5hmC \texorpdfstring{\ensuremath{\to}}{->} 5fC, then BS & Single-base & 5mC
alone \\
TAB-seq & Tet-assisted bisulphite (\texorpdfstring{\ensuremath{\beta}}{beta}-glucosylates 5hmC, then
TET-oxidises 5mC) & Single-base & 5hmC alone \\
ACE-seq & APOBEC chemical sequencing & Single-base & 5hmC
enzymatically \\
CMS-seq & Cytosine 5-methylenesulfonate & Single-base & 5fC alone \\
caMAB-seq & Carboxymethylated-cytosine antibody & about 250-bp
& 5caC alone \\
EM-seq & Enzymatic methyl-seq (no bisulphite damage) & Single-base & 5mC
+ 5hmC vs C \\
\end{longtable}
}

\textbf{Clinical liquid-biopsy translation:} GRAIL's Galleri test
analyses cell-free DNA methylation patterns at \textgreater{} 100,000
CpG sites; multi-cancer early detection (MCED) sensitivity
about 67 \% for stage I--IV combined; tissue-of-origin
localisation accuracy about 93 \%. The technology exploits that
each tissue has a distinct methylation pattern, and tumours shed cfDNA
with their tissue-specific signature.

\subsubsection{CpG Islands and Gene
Regulation}\label{cpg-islands-and-gene-regulation}

\textbf{CpG islands (CGIs):} about 28,000 regions in the human
genome (about 500 bp--2 kb) with: - Observed/expected CpG ratio
\textgreater{} 0.6 - CG content \textgreater{} 55\% - Located near
about 70\% of mammalian gene promoters

\textbf{Normally unmethylated} in somatic cells. Methylation at CGI
promoters is strongly associated with \textbf{stable transcriptional
silencing}; the illustrative methylation heatmap in
\cref{fig:unit_IV_methylation_heatmap} should be read as a beta-value
pattern across loci and samples, not as patient or cell-line evidence.
In empirical data, silencing operates by: 1. Steric blockade of TF
binding (e.g., methylation at Sp1 sites) 2. Recruitment of methyl-CpG
binding proteins (MBD family: MBD1, MBD2, MBD3, MBD4; MeCP2; Kaiso) 3.
MBD-associated \textbf{NuRD complex} (MBD2/3-CHD4-HDAC1/2) \texorpdfstring{\ensuremath{\to}}{->}
deacetylation \texorpdfstring{\ensuremath{\to}}{->} chromatin compaction

\textbf{CpG depletion in the rest of the genome:} Methylated CpGs are
mutation hotspots (5mC \texorpdfstring{\ensuremath{\to}}{->} T deamination produces a C\texorpdfstring{\ensuremath{\to}}{->}T transition). Over
evolutionary time, most CpGs outside islands have been lost ---
explaining the genome-wide CpG suppression. CpG islands have escaped
depletion because methylation is normally absent and the deamination \texorpdfstring{\ensuremath{\to}}{->}
mutation pressure is removed.

\subsubsection{Quantitative Maintenance Through
Replication}\label{quantitative-maintenance-through-replication}

Let \(f_n\) be the fraction of CpGs methylated after \(n\) divisions in
the absence of \emph{de novo} methylation:

\begin{equation}f_n = f_0 \times \left(\frac{1-\epsilon}{2}\right)^n + f_0 \times \frac{1+\epsilon}{2} \tag{15.5}\label{eq:methyl_decay}\end{equation}

where \texorpdfstring{\ensuremath{\varepsilon}}{epsilon} is the maintenance efficiency of DNMT1 (typically
about 0.95 --- i.e.~each hemimethylated CpG is restored 95 \%
of the time at the replication fork). Even high maintenance efficiency
therefore produces measurable \textbf{stochastic demethylation} over
many divisions, which is why constitutively silenced genes also have
backup repressive marks (H3K27me3, H3K9me3, H2AK119ub).

\begin{quote}
\textbf{Worked Example 2 --- DNMT1 efficiency and the kinetics of clonal
demethylation}

\textbf{Setup:} A research group treats a cancer cell line with a DNMT1
small-molecule inhibitor that reduces \texorpdfstring{\ensuremath{\varepsilon}}{epsilon} from 0.95 to 0.50 (50 \%
maintenance per CpG per division). The cells divide every 24 h. The
cancer cell line carries an aberrantly hypermethylated \emph{MLH1}
promoter at \emph{f\texorpdfstring{\textsubscript{0}}{0}} = 0.95 (95 \% CpG methylation).

\textbf{Question:} How many days until methylation drops below the
\emph{f} = 0.30 threshold typically required for re-expression?

\textbf{Solution:} Using \cref{eq:methyl_decay},
\(f_n = f_0 \times \left(\frac{1-\epsilon}{2}\right)^n + f_0 \times \frac{1+\epsilon}{2}\)
--- but at \texorpdfstring{\ensuremath{\varepsilon}}{epsilon} = 0.5 the steady-state floor is
\(f_0 \cdot (1+0.5)/2 = 0.75 \cdot f_0 = 0.71\) for \emph{f\texorpdfstring{\textsubscript{0}}{0}} = 0.95.
The floor exceeds the 0.30 threshold, so methylation will not drop below
0.30 with \texorpdfstring{\ensuremath{\varepsilon}}{epsilon} = 0.5 alone. To re-express MLH1, you need either complete
DNMT1 inhibition (\texorpdfstring{\ensuremath{\varepsilon}}{epsilon} \texorpdfstring{\ensuremath{\to}}{->} 0) or active demethylation (TET upregulation).
With \texorpdfstring{\ensuremath{\varepsilon}}{epsilon} = 0:

\(f_n = f_0 / 2^n\)

{\def\LTcaptype{none} % do not increment counter
\begin{longtable}[]{@{}ll@{}}
\toprule\noalign{}
Generation \emph{n} & \(f_n\) \\
\midrule\noalign{}
\endhead
\bottomrule\noalign{}
\endlastfoot
0 & 0.95 \\
1 & 0.475 \\
2 & 0.238 \\
3 & 0.119 \\
\end{longtable}
}

Threshold crossed between generation 1 and 2 --- i.e., \textbf{2 cell
divisions \texorpdfstring{\ensuremath{\approx}}{~} 48 h} after starting full DNMT1 inhibition. This kinetic is
exactly why clinical responses to azacitidine/decitabine in MDS take
about 4--6 cycles (about 3--4 months) of continuous
suppression to manifest: maintaining \texorpdfstring{\ensuremath{\varepsilon}}{epsilon} near zero is necessary but
insufficient --- the cell must also divide several times to dilute the
existing mark.
\end{quote}

\begin{quote}
\textbf{Clinical Connection:} In cancer, two genome-wide patterns are
consistently found: 1. \textbf{Global hypomethylation}: Repetitive
elements (LINE-1, SINE) become unmethylated \texorpdfstring{\ensuremath{\to}}{->} chromosomal instability,
reactivation of transposons, mis-regulated lineage genes. 2.
\textbf{Focal CGI hypermethylation}: Tumour-suppressor gene promoters
(BRCA1, MLH1, VHL, CDKN2A/p16, MGMT, DAPK) become methylated \texorpdfstring{\ensuremath{\to}}{->}
silencing.

\textbf{DNMT inhibitors} (azacitidine, decitabine) are nucleoside
analogs incorporated into DNA; trapped DNMT1 forms a covalent adduct,
causing passive demethylation. Approved for myelodysplastic syndrome
(MDS) and AML. Combination with \textbf{HDAC inhibitors} is actively
trialled. \textbf{Liquid biopsies} (e.g., GRAIL Galleri, Guardant
Reveal) score hundreds of methylation features in cell-free DNA to
detect early-stage cancer with tissue-of-origin localisation.
\end{quote}

\begin{quote}
\textbf{Concept Check 4:} A patient is treated with decitabine for
high-risk MDS. After 4 cycles, the bone-marrow blast count is unchanged.
Bisulphite sequencing of blast DNA shows the targeted \emph{p15/CDKN2B}
CGI is now 30 \% methylated (from 90 \% at baseline), but transcript
levels remain undetectable. Propose three independent mechanisms that
could explain transcriptional silencing despite DNMT1 inhibition, and
design a bench experiment to distinguish them.
\end{quote}

\begin{center}\rule{0.5\linewidth}{0.5pt}\end{center}

\subsection{Genomic Imprinting --- Mechanisms and Reciprocal
Phenotypes}\label{genomic-imprinting-mechanisms-and-reciprocal-phenotypes}

\textbf{Genomic imprinting} is a form of epigenetic regulation in which
genes are expressed from a single parental
\hyperref[gl:allele]{\textbf{allele}} in a parent-of-origin-specific
manner. about 150 genes are imprinted in humans, organised into
about 30 imprinted clusters each controlled by an
\textbf{imprint control region (ICR)} that is differentially methylated
on the two parental chromosomes.

\subsubsection{Imprint Control Regions (ICRs) --- The Master
Switch}\label{imprint-control-regions-icrs-the-master-switch}

ICRs are \textbf{differentially methylated regions (DMRs)} that are: 1.
Established in the germline in a sex-specific pattern (paternal ICRs in
spermatogenesis; maternal ICRs in oogenesis). 2. Maintained throughout
somatic lineages --- they escape post-fertilisation reprogramming via
specialised protective mechanisms (e.g., \textbf{ZFP445} +
\textbf{TRIM28/KAP1} + \textbf{DPPA3} = the ICR-protection triad in
early embryos). 3. Read by sequence-specific factors (e.g.,
\textbf{CTCF} at unmethylated maternal IGF2/H19 ICR; methyl-binding
proteins at methylated paternal ICRs).

\textbf{The CTCF insulator model --- the canonical mechanism at
IGF2/H19:} CTCF binds CCCTC consensus motifs in the unmethylated ICR.
CTCF blocks the enhancer--promoter looping interaction by recruiting
cohesin to form a chromatin loop ``barrier.'' When the ICR is
methylated, CTCF cannot bind (its zinc fingers fail to recognise
methylated CpGs), so the enhancer is free to loop over and contact the
\emph{IGF2} promoter.

\subsubsection{The IGF2/H19 Locus --- Insulator-based Imprinting
(Detailed
Mechanism)}\label{the-igf2h19-locus-insulator-based-imprinting-detailed-mechanism}

The canonical example is the \textbf{IGF2/H19 locus} on chromosome
11p15:

\begin{figure}[htbp]
\centering
\includegraphics[width=0.82\linewidth,height=4.2in,keepaspectratio]{../figures/mermaid_inline/inline_mermaid_0082_dd4913419cc7.png}
\caption{Imprinting at the IGF2/H19 locus. Paternal allele (top): DMR is methylated, CTCF cannot bind; enhancer contacts IGF2 to IGF2 expressed, H19 silent. Maternal allele (bottom): DMR unmethylated, CTCF binds and insulates enhancer from IGF2; enhancer contacts H19 to H19 expressed, IGF2 silent.}
\end{figure}

\textbf{Functional consequences of the locus:} - \textbf{IGF2}
(insulin-like growth factor 2): paternally expressed; potent fetal
growth factor (signals through IGF1R / IGF2R). Drives placental and
fetal growth. - \textbf{H19}: maternally expressed; a 2.3-kb spliced
lncRNA; reservoir of miR-675 (which targets \emph{IGF1R}, creating a
feedback loop). In some tissues H19 itself acts as a tumour suppressor
by inhibiting IGF1R signalling.

\subsubsection{The 15q11--q13 Locus --- Prader--Willi / Angelman
(Comprehensive
Mechanism)}\label{the-15q11q13-locus-praderwilli-angelman-comprehensive-mechanism}

A different mechanism operates at the \textbf{15q11--q13} imprinted
cluster:

{\def\LTcaptype{none} % do not increment counter
\begin{longtable}[]{@{}
  >{\raggedright\arraybackslash}p{(\linewidth - 4\tabcolsep) * \real{0.1429}}
  >{\raggedright\arraybackslash}p{(\linewidth - 4\tabcolsep) * \real{0.5357}}
  >{\raggedright\arraybackslash}p{(\linewidth - 4\tabcolsep) * \real{0.3214}}@{}}
\toprule\noalign{}
\begin{minipage}[b]{\linewidth}\raggedright
Gene
\end{minipage} & \begin{minipage}[b]{\linewidth}\raggedright
Imprint pattern
\end{minipage} & \begin{minipage}[b]{\linewidth}\raggedright
Mechanism
\end{minipage} \\
\midrule\noalign{}
\endhead
\bottomrule\noalign{}
\endlastfoot
\emph{SNRPN, NDN, MAGEL2} & Paternally expressed & Maternal copies
silenced by methylation at the SNRPN ICR \\
\emph{UBE3A} & Maternally expressed (in neurons) & Paternal \emph{UBE3A}
silenced by an antisense lncRNA \emph{UBE3A-ATS} expressed primarily
from the paternal allele \\
\end{longtable}
}

\textbf{Mechanism details:} - The \textbf{SNURF-SNRPN ICR} sits
about 30 kb upstream of \emph{SNRPN} and is methylated on the
maternal allele (silencing maternal \emph{SNRPN, NDN, MAGEL2}) and
unmethylated on the paternal allele (where these genes are expressed). -
The unmethylated paternal ICR drives transcription of a long
polycistronic transcript that extends about 600 kb to produce
\emph{SNRPN, SNORD-cluster snoRNAs} (SNORD115, SNORD116), and
\emph{UBE3A-ATS} (the antisense to \emph{UBE3A}). - \emph{UBE3A-ATS}
transcription on the paternal allele suppresses paternal \emph{UBE3A}
expression --- \emph{primarily in neurons}. This is the most surprising
aspect of the locus: the imprint is \textbf{brain-specific}.

\textbf{Phenotypes --- the parental-conflict reciprocal:} -
\textbf{Prader--Willi syndrome}: Loss of \textbf{paternal} 15q11--q13
expression --- by paternal deletion (about 70\%), maternal
uniparental disomy (about 25\%), or imprinting centre defect
(about 3\%). Phenotype: neonatal hypotonia, hyperphagia and
obesity, intellectual disability, hypogonadism, growth-hormone
deficiency. The hypothalamic phenotype (especially the SNORD116 loss)
drives the food-intake dysregulation. - \textbf{Angelman syndrome}: Loss
of \textbf{maternal} \emph{UBE3A} expression in the same region --- by
maternal deletion (about 70\%), paternal uniparental disomy
(about 5\%), imprinting centre defect (about 5\%), or
\emph{UBE3A} point mutation (about 10\%). Phenotype: severe
intellectual disability, seizures, ataxic gait, ``happy puppet''
demeanour, absence of speech. UBE3A is an E3 ubiquitin ligase critical
for synaptic plasticity; loss in neurons disrupts excitatory/inhibitory
balance.

\textbf{Therapeutic angle --- antisense oligonucleotide reactivation:}
Activation of the silenced paternal \emph{UBE3A} by antisense
oligonucleotides targeting \emph{UBE3A-ATS} is in clinical trials for
Angelman syndrome (Ionis/Biogen ION-582/GTX-102, Phase 1/2 trials
underway 2024). The strategy: deplete the antisense lncRNA \texorpdfstring{\ensuremath{\to}}{->} de-repress
paternal \emph{UBE3A} \texorpdfstring{\ensuremath{\to}}{->} restore neuronal UBE3A protein. Intrathecal
delivery is required because UBE3A imprinting is brain-specific.

\subsubsection{Beckwith--Wiedemann and Silver--Russell --- Reciprocal
Disorders}\label{beckwithwiedemann-and-silverrussell-reciprocal-disorders}

The IGF2/H19 ICR also produces two reciprocal phenotypes --- a textbook
illustration of \textbf{dosage and parental conflict}:

\begin{itemize}
\tightlist
\item
  \textbf{Beckwith--Wiedemann syndrome (BWS):} biallelic IGF2 expression
  (loss of maternal imprint, \textbf{hypermethylation} of the maternal
  ICR at H19) \texorpdfstring{\ensuremath{\to}}{->} maternal allele now expresses IGF2 like the paternal \texorpdfstring{\ensuremath{\to}}{->}
  fetal overgrowth, macrosomia (birth weight \textgreater{} 4 kg),
  macroglossia, omphalocele, hypoglycaemia, hemihyperplasia, and a 7.5
  \% risk of childhood embryonal tumours (Wilms tumour, hepatoblastoma,
  neuroblastoma, rhabdomyosarcoma). Diagnosis: methylation analysis of
  the H19 DMR shows hypermethylation; some cases also have paternal UPD
  of 11p15 (about 20 \%).
\item
  \textbf{Silver--Russell syndrome (SRS):} biallelic H19 expression / no
  IGF2 (loss of paternal imprint, \textbf{hypomethylation} of the
  paternal ICR) \texorpdfstring{\ensuremath{\to}}{->} paternal allele now silent like the maternal \texorpdfstring{\ensuremath{\to}}{->} severe
  intrauterine growth restriction (IUGR), postnatal short stature, body
  asymmetry, characteristic triangular face, fifth-finger clinodactyly.
  Diagnosis: hypomethylation of the H19 DMR (about 50 \%) or
  maternal UPD of chromosome 7 (about 10 \%).
\end{itemize}

This ``mirror image'' pair illustrates the \textbf{parental-conflict
(kinship)} theory: paternally-expressed growth factors maximise
embryonic growth (paternal genes care less about maternal resources,
since the same father may not father subsequent siblings);
maternally-expressed antagonists restrain growth (the mother distributes
resources across multiple offspring).

\begin{quote}
\textbf{Concept Check 5:} A child is born with severe IUGR,
characteristic facial features, and asymmetric limbs. Methylation
analysis shows that both the paternal and maternal copies of the 11p15.5
ICR are \emph{hypomethylated}. Predict the likely diagnosis. Explain why
hypermethylation versus hypomethylation at the \emph{same} ICR produces
\emph{opposite} phenotypes (BWS vs SRS).
\end{quote}

\begin{center}\rule{0.5\linewidth}{0.5pt}\end{center}

\subsection{X-Chromosome Inactivation (Lyonisation) --- In
Depth}\label{x-chromosome-inactivation-lyonisation-in-depth}

In female placental mammals, dosage compensation for about 900
X-linked genes is achieved by transcriptionally silencing one of the two
X chromosomes in every somatic cell.

\textbf{Key features:} - \textbf{Random choice:} Either the maternal or
paternal X can be inactivated in each cell (choice is random in each
blastomere nucleus). Imprinted paternal X-inactivation occurs in mouse
extra-embryonic lineages but not in humans. - \textbf{Clonal
maintenance:} Once established (about  day 4.5 in human
embryo), the inactive state is mitotically heritable in most daughter
cells --- the same X is silenced in every descendant. - \textbf{Barr
body:} The inactive X (Xi) forms a dense heterochromatic body (Barr
body) visible by interphase cytology. - \textbf{Incomplete:}
about 15--25\% of X-linked genes \textbf{escape} inactivation.
Most escapees lie in the \textbf{pseudoautosomal regions (PAR1, PAR2)}
that recombine with the Y; others (e.g., \emph{KDM6A}, \emph{KDM5C},
\emph{DDX3X}, \emph{ZFX}) escape across the X. Escape genes contribute
to female-biased autoimmune disease and to Turner syndrome
haploinsufficiency phenotypes.

\textbf{Molecular mechanism --- a hierarchical cascade (timing and
milestones):}

{\def\LTcaptype{none} % do not increment counter
\begin{longtable}[]{@{}
  >{\raggedright\arraybackslash}p{(\linewidth - 4\tabcolsep) * \real{0.4884}}
  >{\raggedright\arraybackslash}p{(\linewidth - 4\tabcolsep) * \real{0.1163}}
  >{\raggedright\arraybackslash}p{(\linewidth - 4\tabcolsep) * \real{0.3953}}@{}}
\toprule\noalign{}
\begin{minipage}[b]{\linewidth}\raggedright
Time after initiation
\end{minipage} & \begin{minipage}[b]{\linewidth}\raggedright
Event
\end{minipage} & \begin{minipage}[b]{\linewidth}\raggedright
Molecular outcome
\end{minipage} \\
\midrule\noalign{}
\endhead
\bottomrule\noalign{}
\endlastfoot
0 min & Choice made; Xist transcription up-regulated on future Xi &
Allele-specific Xist transcription begins \\
0--60 min & Xist RNA spreads in cis & Coats 165 Mb Xi; entry sites at
LINE-1-rich domains \\
1--6 h & Repeat A recruits SPEN/SHARP & SPEN recruits HDAC3 \texorpdfstring{\ensuremath{\to}}{->}
genome-wide deacetylation \\
1--6 h & Repeat B/C/D recruits PRC2 & EZH2 deposits H3K27me3 on Xi \\
6--24 h & Repeat E recruits PTBP1, MATR3 & Nuclear matrix tethering of
Xi \\
6--24 h & PRC1 recruited downstream & RING1B writes H2AK119ub \\
1--7 days & DNMT3A/3B methylate CpG-island promoters & Stable
maintenance of silencing \\
2--14 days & Histone variants accumulate & macroH2A and H2A.Z reinforce
heterochromatin \\
7--21 days & LBR anchors Xi to nuclear lamina & Xi enters B compartment
(nuclear periphery) \\
\end{longtable}
}

\textbf{XIST regulation by TSIX:} The antisense lncRNA \emph{TSIX}
(transcribed across XIST in the opposite direction) keeps XIST repressed
on the future active X (Xa). The mutual exclusion of XIST/TSIX between
alleles enforces monoallelic XIST expression. \emph{TSIX} itself is
regulated by \emph{XITE} (X-inactivation intergenic transcription
element) and by chromatin-state asymmetries between the two alleles in
early embryogenesis.

\textbf{Reactivation in iPSC reprogramming:} When somatic female cells
are reprogrammed to induced pluripotent stem cells (iPSCs), the inactive
X reactivates --- XIST is silenced, H3K27me3 is lost, DNA methylation is
removed, and both alleles express X-linked genes again. This is a unique
window for studying XCI dynamics. iPSCs derived from female patients
with X-linked disease may have variable X-inactivation patterns after
differentiation, complicating disease modelling.

\textbf{Clinical implications --- X-linked diseases in heterozygous
females:}

\begin{equation}\text{P(affected)} \approx \int_0^1 g(x) \cdot \mathbf{1}[x \le x_{\text{threshold}}] \, dx \tag{15.6}\label{eq:xci_skew}\end{equation}

where \(g(x)\) is the distribution of X-inactivation skewing across
cells. \textbf{Skewed X-inactivation} (where one allele is silenced in
\textgreater75 \% of cells) explains why heterozygous female carriers of
X-linked diseases (Duchenne muscular dystrophy, Rett syndrome, X-linked
agammaglobulinaemia) show variable severity. Skewing arises when: -
Random inactivation by chance produces an extreme distribution -
Cell-autonomous selection eliminates clones with the active mutant
allele (e.g., immunodeficiencies) - A \emph{XIST} or \emph{XIC} mutation
biases choice - Skewing is age-dependent: fraction of women with
\textgreater{} 75 \% skewing rises from about 5 \% at birth to
about 25 \% by age 60 (clonal drift in haematopoiesis).

\textbf{Lyon's evidence} \citep{lyon1961}: Female mice heterozygous for
coat-colour mutations show patchy (mosaic) coat colour --- a direct
consequence of random X-inactivation producing patches of cells
expressing one allele or the other. The same mosaicism is seen in human
female carriers of X-linked albinism, X-linked anhidrotic ectodermal
dysplasia (Christ--Siemens--Touraine, with patchy sweat glands), and
X-linked retinoschisis (patchy retinal involvement).

\begin{quote}
\textbf{Worked Example 3 --- Predicting penetrance from XCI skewing}

\textbf{Setup:} A female patient is heterozygous for a Duchenne muscular
dystrophy (DMD) deletion. Muscle weakness is observed in cells where the
wild-type (active) allele is on the inactivated X. Assume that DMD
muscle fibres are syncytial (myofibres pool many nuclei); a fibre is
dystrophic if \textgreater{} 50 \% of its nuclei express the deleted
allele.

\textbf{Question:} If XCI is unbiased (mean skewing = 0.5, standard
deviation \texorpdfstring{\ensuremath{\sigma}}{sigma}\_skew = 0.05 across mononuclear cells), what fraction of
myofibres will be dystrophic? What if the patient has clonal-skewing
where \texorpdfstring{\ensuremath{\sigma}}{sigma}\_skew = 0.20?

\textbf{Solution:} A fibre is dystrophic if more than half its nuclei
have the wild-type allele inactivated. With \emph{N} = 20 nuclei per
fibre and a per-nucleus probability \emph{x} of expressing the deleted
allele (= probability wild-type X is inactivated), the fibre is
dystrophic when more than \emph{N}/2 nuclei out of \emph{N} express the
deleted allele.

For unbiased XCI, \emph{x} per nucleus \texorpdfstring{\ensuremath{\approx}}{~} 0.5, \texorpdfstring{\ensuremath{\sigma}}{sigma}\_x = 0.05. Per fibre,
the average fraction of mutant-expressing nuclei \texorpdfstring{\ensuremath{\approx}}{~} 0.5 \texorpdfstring{\ensuremath{\pm}}{+/-} 0.05/\texorpdfstring{\(\sqrt{\;}\)}{sqrt}20 \texorpdfstring{\ensuremath{\approx}}{~} 0.5
\texorpdfstring{\ensuremath{\pm}}{+/-} 0.011. Almost no fibres exceed 50 \% (by about \texorpdfstring{\ensuremath{\sigma}}{sigma}\_fibre \texorpdfstring{\ensuremath{\times}}{x} 1 \texorpdfstring{\ensuremath{\approx}}{~} 1 \%).

For clonal-skewing \texorpdfstring{\ensuremath{\sigma}}{sigma}\_skew = 0.20 across mononuclear precursors, mean
per-fibre fraction is 0.5 but with much wider variance: \texorpdfstring{\ensuremath{\sigma}}{sigma}\_fibre \texorpdfstring{\ensuremath{\approx}}{~}
0.20/\texorpdfstring{\(\sqrt{\;}\)}{sqrt}20 \texorpdfstring{\ensuremath{\approx}}{~} 0.045. Now about 25 \% of fibres exceed 50 \% mutant
expression and become dystrophic.

\textbf{Insight:} Manifesting Duchenne muscular dystrophy in
heterozygous female carriers (about 10 \% show muscle weakness)
is largely explained by \textbf{age-acquired XCI skewing} in
haematopoietic and myogenic precursors. Treatment strategies:
skewing-modulating ASOs to redirect XCI; or AAV-DMD gene therapy
(rebalancing dystrophin expression).
\end{quote}

\begin{quote}
\textbf{Concept Check 6:} A 4-year-old girl is diagnosed with severe
Rett syndrome. Sequencing reveals a heterozygous loss-of-function
mutation in \emph{MECP2}. XCI analysis shows 95 \% skewing toward
inactivation of the \emph{wild-type} allele. Explain (i) why this
skewing pattern produces severe disease, (ii) why some Rett patients
with the same mutation are mildly affected, and (iii) whether modulating
XCI skewing therapeutically could be a treatment strategy.
\end{quote}

\begin{center}\rule{0.5\linewidth}{0.5pt}\end{center}

\subsection{Chromatin Remodelling Complexes --- Mechanism in
Detail}\label{chromatin-remodelling-complexes-mechanism-in-detail}

In addition to covalent modifications, chromatin structure is actively
remodelled by ATP-dependent complexes that slide, eject, or restructure
nucleosomes. Most four families share a conserved \textbf{Snf2-family
ATPase} (SF2 helicase superfamily) but couple ATP hydrolysis to distinct
nucleosome operations.

\subsubsection{The Four Major Remodeller
Families}\label{the-four-major-remodeller-families}

\begin{figure}[htbp]
\centering
\includegraphics[width=0.82\linewidth,height=4.2in,keepaspectratio]{../figures/mermaid_inline/inline_mermaid_0083_0578d11a4c13.png}
\caption{Four ATP-dependent remodelling families, each with a distinct nucleosome operation. SWI/SNF: large bursts of sliding and ejection that expose regulatory DNA. ISWI: short, regular spacing for nucleosome arrays. CHD/NuRD: repressive sliding with HDAC coupling. INO80/SWR1: histone variant exchange (H2A.Z between H2A) at promoters and DNA damage sites.}
\end{figure}

{\def\LTcaptype{none} % do not increment counter
\begin{longtable}[]{@{}
  >{\raggedright\arraybackslash}p{(\linewidth - 10\tabcolsep) * \real{0.1132}}
  >{\raggedright\arraybackslash}p{(\linewidth - 10\tabcolsep) * \real{0.3208}}
  >{\raggedright\arraybackslash}p{(\linewidth - 10\tabcolsep) * \real{0.1132}}
  >{\raggedright\arraybackslash}p{(\linewidth - 10\tabcolsep) * \real{0.1698}}
  >{\raggedright\arraybackslash}p{(\linewidth - 10\tabcolsep) * \real{0.1509}}
  >{\raggedright\arraybackslash}p{(\linewidth - 10\tabcolsep) * \real{0.1321}}@{}}
\toprule\noalign{}
\begin{minipage}[b]{\linewidth}\raggedright
Family
\end{minipage} & \begin{minipage}[b]{\linewidth}\raggedright
Prototype complex
\end{minipage} & \begin{minipage}[b]{\linewidth}\raggedright
ATPase
\end{minipage} & \begin{minipage}[b]{\linewidth}\raggedright
Mechanism
\end{minipage} & \begin{minipage}[b]{\linewidth}\raggedright
Function
\end{minipage} & \begin{minipage}[b]{\linewidth}\raggedright
Disease
\end{minipage} \\
\midrule\noalign{}
\endhead
\bottomrule\noalign{}
\endlastfoot
SWI/SNF & BAF, PBAF (mammals); SWI/SNF (yeast); BAP, PBAP
(\emph{Drosophila}) & BRG1 (SMARCA4) or BRM (SMARCA2) & Large-scale
sliding and ejection & Activate transcription; prepare enhancers;
required for stem cell self-renewal & ARID1A mutated in
about 10\% cancers; SMARCB1/INI1 lost in malignant rhabdoid
tumour; SMARCA4 in lung cancer \\
ISWI & NURF, CHRAC, ACF, RSF, NoRC & SNF2H (SMARCA5), SNF2L (SMARCA1) &
Nucleosome spacing --- generates regular arrays & Chromatin assembly
post-replication; heterochromatin maintenance; transcriptional
repression & Williams syndrome (BAZ1B/WSTF deletion) \\
CHD / NuRD & NuRD (CHD3/4 + MBD3 + MTA1/2/3 + HDAC1/2 + RBBP4/7) & CHD1,
CHD3 (Mi-2\texorpdfstring{\ensuremath{\alpha}}{alpha}), CHD4 (Mi-2\texorpdfstring{\ensuremath{\beta}}{beta}) & Nucleosome sliding coupled to histone
deacetylation & Gene repression; lineage commitment; Polycomb-related
silencing & CHD7 in CHARGE syndrome; CHD8 in autism \\
INO80 / SWR1 & INO80, SWR1 (yeast); SRCAP (mammals) & INO80, SRCAP &
Histone variant exchange (H2A.Z \texorpdfstring{\ensuremath{\leftrightarrow}}{<->} canonical H2A; H2A.X \texorpdfstring{\ensuremath{\leftrightarrow}}{<->} H2A) &
Promoter-proximal H2A.Z deposition; DSB repair (\texorpdfstring{\ensuremath{\gamma}}{gamma}H2AX exchange);
centromere identity & Floating-Harbor syndrome (SRCAP mutations) \\
\end{longtable}
}

\textbf{Mechanism in detail --- the ``wave'' model of DNA
translocation:} Most remodellers translocate DNA through the nucleosome
by alternately gripping minor groove with two RecA-like ATPase lobes.
Each ATP cycle introduces a 1-bp DNA bulge that propagates around the
histone surface, effectively pulling DNA into the nucleosome from one
side and ejecting it from the other. The size of the ATP-driven movement
varies (1--10 bp per cycle) and explains the family-specific outcomes
(single-nucleosome repositioning vs.~processive slides vs.~ejection).
Substrate recognition differs: - \textbf{SWI/SNF} binds at SHL-2
(superhelical location -2 from the dyad), gripping nucleosomal DNA
about 20 bp from the histone--DNA boundary. - \textbf{ISWI}
binds the linker DNA flanking the nucleosome and uses an autoinhibitory
AutoN domain that is released by H4-tail acetylation. -
\textbf{CHD/NuRD} binds via a tandem chromodomain that reads H3K9me3
(CHD3) or H3K4me0 (CHD4), targeting the complex to silenced chromatin. -
\textbf{INO80} binds the H2A.Z-H2B dimer for selective exchange.

\subsubsection{SWI/SNF Specifically --- The Cancer-Most-Mutated
Remodeller}\label{swisnf-specifically-the-cancer-most-mutated-remodeller}

The mammalian BAF complex contains 12--15 subunits, including ATPase
BRG1 or BRM, ARID1A/B (DNA-binding), and SMARCB1/INI1. Its assembly is
highly cell-type-specific (npBAF in neural progenitors \texorpdfstring{\ensuremath{\to}}{->} nBAF in
post-mitotic neurons, with subunit swap at \emph{MIR9} induction; esBAF
in embryonic stem cells). BAF: - \textbf{Evicts PRC1/PRC2 from enhancers
within minutes} (this is the basis for the synthetic-lethal sensitivity
of SWI/SNF-mutant tumours to EZH2 inhibitors). - Generates
nucleosome-depleted regions (NDRs) at enhancers and promoters - Recruits
transcription factors (e.g., pioneer TFs FOXA1, GATA4) by exposing their
motifs - Specific complexes: - \textbf{canonical BAF (cBAF):} contains
ARID1A or ARID1B and DPF1/2/3 - \textbf{polybromo BAF (PBAF):} contains
BAF180/PBRM1, ARID2, and BRD7 --- required for IFN response and
differentiation - \textbf{non-canonical BAF (ncBAF/GBAF):} contains
BRD9, GLTSCR1/L, and BRG1 --- found in synovial sarcoma (where SS18-SSX
fusions hijack ncBAF)

\subsubsection{Targeted Epigenome Editing with dCas9
Fusions}\label{targeted-epigenome-editing-with-dcas9-fusions}

The CRISPR toolkit has expanded far beyond DNA cutting. A catalytically
dead \textbf{dCas9} (D10A + H840A) retains its ability to be guided to a
specific DNA sequence by a sgRNA but no longer cleaves --- turning it
into a programmable DNA-binding scaffold. Fusing dCas9 to an epigenetic
enzyme creates a targeted ``writer'' or ``eraser'' of chromatin state
\citep{doudna2014}:

{\def\LTcaptype{none} % do not increment counter
\begin{longtable}[]{@{}
  >{\raggedright\arraybackslash}p{(\linewidth - 4\tabcolsep) * \real{0.1463}}
  >{\raggedright\arraybackslash}p{(\linewidth - 4\tabcolsep) * \real{0.1951}}
  >{\raggedright\arraybackslash}p{(\linewidth - 4\tabcolsep) * \real{0.6585}}@{}}
\toprule\noalign{}
\begin{minipage}[b]{\linewidth}\raggedright
Fusion
\end{minipage} & \begin{minipage}[b]{\linewidth}\raggedright
Function
\end{minipage} & \begin{minipage}[b]{\linewidth}\raggedright
Experimental / clinical use
\end{minipage} \\
\midrule\noalign{}
\endhead
\bottomrule\noalign{}
\endlastfoot
dCas9--DNMT3A & Programmable DNA methylation & \textbf{CRISPRoff}
(heritable silencing without cutting DNA) \\
dCas9--TET1 & Programmable DNA demethylation & Rescue of fragile-X
\emph{FMR1} silencing; reactivation of tumour suppressors \\
dCas9--p300 & Programmable H3K27 acetylation & Programmable enhancer
activation \\
dCas9--KRAB & Recruits KAP1 \texorpdfstring{\ensuremath{\to}}{->} SETDB1 \texorpdfstring{\ensuremath{\to}}{->} H3K9me3 & \textbf{CRISPRi} ---
stable repression of any gene \\
dCas9--LSD1 & H3K4me1/2 demethylation & Enhancer decommissioning \\
dCas9--VPR (VP64-p65-Rta) & Transcriptional activation &
\textbf{CRISPRa} --- about 10-to-1000-fold induction \\
dCas9--PRC1 (RING1B--PCGF4) & H2AK119ub deposition & Polycomb-style
silencing \\
\end{longtable}
}

The clinical promise is substantial: a dCas9--DNMT3A targeting the
\emph{HBG1/2} promoter re-activates fetal haemoglobin in adult
\texorpdfstring{\ensuremath{\beta}}{beta}-thalassaemia patients without permanent genome modification.
CRISPRoff-induced silencing of \emph{PCSK9} in liver cells is being
pursued for hypercholesterolaemia. The key limitation is off-target
\hyperref[gl:epigenome]{\textbf{epigenome}} editing --- bystander
methylation at sequence-similar loci --- which the field now addresses
with paired-sgRNA and split-dCas9 designs.

\begin{quote}
{[}!NOTE{]} Unlike Cas9-mediated gene editing, epigenome editing is
\textbf{reversible}: the perturbation fades unless it is
self-reinforcing through mitotic heritability of DNA methylation. This
makes it attractive as a research tool (transient perturbation) and a
therapy (tunable, removable), but means durable clinical effects require
either continuous dCas9 expression or recruitment of self-propagating
chromatin states (e.g., H3K27me3 spreading by PRC2).
\end{quote}

\begin{quote}
\textbf{Clinical Connection:} \textbf{SWI/SNF subunits are mutated in
about 20\% of human cancers}, the highest mutation frequency of
any single chromatin regulator. ARID1A (about 10\%),
SMARCA4/BRG1 (about 5\%), SMARCB1/INI1 (childhood rhabdoid
tumours, near-almost universally biallelic loss), PBRM1
(about 40\% of clear-cell renal cell carcinoma). Many
SWI/SNF-mutant tumours become \textbf{synthetically lethal with EZH2}
(PRC2) inhibition because they over-rely on Polycomb repression of
cell-cycle inhibitors when SWI/SNF cannot evict PRC1/2.
\end{quote}

\begin{center}\rule{0.5\linewidth}{0.5pt}\end{center}

\subsection{Three-Dimensional Genome Organisation and Phase
Separation}\label{three-dimensional-genome-organisation-and-phase-separation}

\subsubsection{TADs, Loops, and Compartments --- Spatial
Layers}\label{tads-loops-and-compartments-spatial-layers}

The genome is folded across multiple length scales, and each scale
contributes to gene regulation:

{\def\LTcaptype{none} % do not increment counter
\begin{longtable}[]{@{}
  >{\raggedright\arraybackslash}p{(\linewidth - 6\tabcolsep) * \real{0.2000}}
  >{\raggedright\arraybackslash}p{(\linewidth - 6\tabcolsep) * \real{0.2500}}
  >{\raggedright\arraybackslash}p{(\linewidth - 6\tabcolsep) * \real{0.3000}}
  >{\raggedright\arraybackslash}p{(\linewidth - 6\tabcolsep) * \real{0.2500}}@{}}
\toprule\noalign{}
\begin{minipage}[b]{\linewidth}\raggedright
Length scale
\end{minipage} & \begin{minipage}[b]{\linewidth}\raggedright
Structural unit
\end{minipage} & \begin{minipage}[b]{\linewidth}\raggedright
Marker / detection
\end{minipage} & \begin{minipage}[b]{\linewidth}\raggedright
Functional role
\end{minipage} \\
\midrule\noalign{}
\endhead
\bottomrule\noalign{}
\endlastfoot
1 kb--100 kb & Promoter--enhancer loops & ChIA-PET, HiChIP, Capture-Hi-C
& Direct enhancer--TSS contact for activation \\
100 kb--1 Mb & TADs & Hi-C insulation score; CTCF/cohesin ChIP &
Constrains enhancer search to local genes \\
1 Mb--100 Mb & A/B compartments & Hi-C eigenvector & Active
vs.~repressed neighbourhoods \\
Chromosome & Chromosome territories & DNA FISH & Each chromosome
occupies a distinct nuclear domain \\
Nuclear scale & LADs, NADs, speckles & DamID, TSA-seq & Nuclear lamina,
nucleolus, nuclear speckle proximity \\
\end{longtable}
}

\textbf{LADs (Lamina-Associated Domains):} about 1,300 genomic
regions (median about 1 Mb) attached to the nuclear lamina via
lamin B receptor, marked by H3K9me2/me3, gene-poor and transcriptionally
repressive. Lamina detachment correlates with gene activation during
differentiation.

\subsubsection{Biomolecular Condensates and Phase
Separation}\label{biomolecular-condensates-and-phase-separation}

Chromatin is not a static polymer; many regulatory proteins undergo
\textbf{liquid--liquid phase separation} (LLPS) or form
\textbf{condensates} with elevated local concentration. Intrinsically
disordered regions (IDRs) on transcription factors (e.g., \textbf{BRD4},
\textbf{Mediator} subunits, \textbf{RNA Pol II} CTD-associated factors)
promote clustering at \textbf{super-enhancers} --- unusually large
clusters of enhancers densely occupied by Mediator, co-activators, and
active histone marks (\textbf{H3K27ac}). The resulting
\textbf{transcriptional condensate} concentrates the phosphorylation
machinery that releases promoter-proximal paused Pol II, explaining why
some loci fire at very high rates (oncogenes such as \emph{MYC} in
selected cancers).

\textbf{Quantitative criteria for LLPS in cells:} - IDR-rich proteins
above a saturation concentration \(c_{\text{sat}}\) - Multivalent
interactions (low-affinity, high-valency) between IDRs, with
\(K_d \sim 1\)--\(100\) \texorpdfstring{\ensuremath{\mu}}{mu}M - Round, dynamic droplets that fuse and
undergo FRAP recovery on the seconds-to-minutes timescale - Disassembly
upon 1,6-hexanediol (a hallmark, though imperfect, test)

\textbf{Examples in chromatin biology:}

{\def\LTcaptype{none} % do not increment counter
\begin{longtable}[]{@{}
  >{\raggedright\arraybackslash}p{(\linewidth - 4\tabcolsep) * \real{0.3333}}
  >{\raggedright\arraybackslash}p{(\linewidth - 4\tabcolsep) * \real{0.4000}}
  >{\raggedright\arraybackslash}p{(\linewidth - 4\tabcolsep) * \real{0.2667}}@{}}
\toprule\noalign{}
\begin{minipage}[b]{\linewidth}\raggedright
Condensate
\end{minipage} & \begin{minipage}[b]{\linewidth}\raggedright
Constituents
\end{minipage} & \begin{minipage}[b]{\linewidth}\raggedright
Function
\end{minipage} \\
\midrule\noalign{}
\endhead
\bottomrule\noalign{}
\endlastfoot
Heterochromatin foci & HP1\texorpdfstring{\ensuremath{\alpha}}{alpha}, H3K9me3 & Concentrates H3K9
methyltransferases; phase-separated repressive compartment \\
Nucleolus & NPM1, fibrillarin, Pol I, rRNA & Ribosome biogenesis;
multiphase (FC/DFC/GC sub-compartments) \\
Nuclear speckles & SRSF, SON, MALAT1 & Storage and assembly of splicing
factors \\
Cajal bodies & Coilin, snRNPs & snRNP and telomerase RNP biogenesis \\
PML bodies & PML, SUMO, p53 & DNA damage response; senescence; viral
defence \\
Super-enhancer condensates & BRD4, Mediator, Pol II CTD & Robust
transcription of cell-identity genes \\
Polycomb bodies & CBX2 (PRC1) phase-separated & Polycomb domain
compaction in \emph{cis} \\
\end{longtable}
}

\textbf{Conceptual link to TADs:} Condensates operate \emph{within} TADs
and at promoter--enhancer loops; disrupting CTCF boundaries can move an
oncogenic enhancer adjacent to a silent proto-oncogene (\textbf{enhancer
hijacking}), a structural-variant mechanism increasingly catalogued in
paediatric tumours.

\begin{figure}[htbp]
\centering
\includegraphics[width=0.82\linewidth,height=4.2in,keepaspectratio]{../figures/mermaid_inline/inline_mermaid_0084_989de778d28c.png}
\caption{Super-enhancer-associated condensate bridging enhancers to paused Pol II; BET inhibitors weaken acetyl-lysine engagement of BRD4 and dissolve the condensate; CDK8-mediated phosphorylation of Mediator disrupts condensate integrity.}
\end{figure}

\textbf{Clinical angle --- drugging condensates:} - \textbf{BET
bromodomain inhibitors} (e.g., JQ1, OTX015, birabresib, molibresib,
mivebresib) displace BRD4 from acetylated chromatin, collapsing
condensate-associated transcription at \emph{MYC} and other dependency
genes. - \textbf{CDK7 inhibitors} (THZ1, SY-5609) and \textbf{CDK9
inhibitors} target the kinases concentrated in transcription
condensates. - \textbf{EZH2 inhibitors} (tazemetostat) shrink
PRC2-mediated repressive domains. - \textbf{Tumour cells with
super-enhancer-driven oncogenes (MYCN-amplified neuroblastoma,
MLL-rearranged leukaemia, \emph{TAL1}-driven T-cell acute lymphoblastic
leukaemia) can be highly sensitive} to BET inhibition, but response
depends on enhancer wiring, compensatory transcription factors, and
therapeutic window rather than on super-enhancer status alone.

\begin{quote}
\textbf{Worked Example 4 --- Polymer Statistics in Hi-C:}

\textbf{Setup:} A human gene at chromosome 7p15 harbours an enhancer 700
kb upstream. We model intra-chromosomal contact probability with the
fractal-globule scaling \cref{eq:hic_scaling},
\(P(s) \propto s^{-\alpha}\), with \texorpdfstring{\ensuremath{\alpha}}{alpha} = 1.1.

\textbf{Question:} Estimate the probability of physical contact between
enhancer and promoter (s = 700,000 bp). Compare to (a) within-TAD
contact at s = 50 kb and (b) trans-chromosomal contact (genome-wide Hi-C
average \texorpdfstring{\ensuremath{\approx}}{~} 10\texorpdfstring{\textsuperscript{-}}{-}\texorpdfstring{\textsuperscript{6}}{6}).

\textbf{Solution:} Setting \emph{P}(50 kb) \texorpdfstring{\ensuremath{\approx}}{~} 1 (within-TAD baseline),
the relative \emph{P}(700 kb) = (700/50)\texorpdfstring{\textsuperscript{-}}{-}\texorpdfstring{\textsuperscript{1}}{1}\texorpdfstring{\ensuremath{\cdot}}{.}\texorpdfstring{\textsuperscript{1}}{1} = 14\texorpdfstring{\textsuperscript{-}}{-}\texorpdfstring{\textsuperscript{1}}{1}\texorpdfstring{\ensuremath{\cdot}}{.}\texorpdfstring{\textsuperscript{1}}{1} \texorpdfstring{\ensuremath{\approx}}{~} 0.063. The
enhancer--promoter contact probability is about 6 \% of the
within-TAD baseline.

Compare to trans-chromosomal: 6 \texorpdfstring{\ensuremath{\times}}{x} 10\texorpdfstring{\textsuperscript{-}}{-}\texorpdfstring{\textsuperscript{2}}{2} (intra-TAD-distant) versus 1 \texorpdfstring{\ensuremath{\times}}{x}
10\texorpdfstring{\textsuperscript{-}}{-}\texorpdfstring{\textsuperscript{6}}{6} (trans) = a \textbf{about 60,000-fold preference} for
intra-chromosomal contacts even at long distances.

\textbf{Insight:} CTCF/cohesin loop extrusion concentrates this
60,000-fold enrichment of intra-TAD contacts onto specific
enhancer--promoter pairs, raising effective contact probability another
10--100\texorpdfstring{\ensuremath{\times}}{x}. This is why disrupting a single CTCF anchor (about 20
bp deletion) can completely abolish enhancer--promoter looping and
silence the target gene --- even though the enhancer DNA is unchanged.
\end{quote}

\begin{center}\rule{0.5\linewidth}{0.5pt}\end{center}

\subsection{Non-Coding RNAs in Gene
Regulation}\label{non-coding-rnas-in-gene-regulation}

\subsubsection{MicroRNAs (miRNAs)}\label{micrornas-mirnas}

miRNAs are about 22 nt single-stranded RNAs that direct
post-transcriptional silencing. Over 2,000 annotated human miRNAs
collectively regulate about 60\% of protein-coding genes
\citep{fire1998}.

\begin{figure}[htbp]
\centering
\includegraphics[width=0.82\linewidth,height=4.2in,keepaspectratio]{../figures/mermaid_inline/inline_mermaid_0085_192a3bb982f8.png}
\caption{miRNA biogenesis and RISC-mediated silencing. Drosha in the nucleus produces pre-miRNA; DICER in the cytoplasm generates the duplex; AGO2 incorporates the guide strand into RISC; partial 3 UTR complementarity leads to translational repression and mRNA decay.}
\end{figure}

\textbf{Key oncogenic miRNA examples:} - \textbf{miR-21} (oncomiR):
Overexpressed in most cancers; targets PTEN, PDCD4, RECK (tumour
suppressors) - \textbf{miR-155} (oncomiR): Overexpressed in B-cell
lymphomas; targets SHIP1 (AKT suppressor) - \textbf{miR-34a} (tumour
suppressor miR): Downstream of p53; targets CDK6, BCL2, SNAIL;
methylated/silenced in many cancers

\subsubsection{Long Non-Coding RNAs
(lncRNAs)}\label{long-non-coding-rnas-lncrnas}

lncRNAs are \textgreater200 nt functional RNA transcripts with no
protein-coding potential. \textgreater100,000 annotated in the human
genome. Mechanisms of action are diverse:

\begin{itemize}
\tightlist
\item
  \textbf{XIST:} Coats the inactive X chromosome; recruits PRC2 to
  spread H3K27me3, as described in the X-inactivation section.
\item
  \textbf{HOTAIR:} Transcribed from HOXC; binds PRC2 to direct H3K27me3
  deposition at HOXD and other loci
\item
  \textbf{MALAT1:} Nuclear speckle-associated; regulates alternative
  splicing; highly expressed in cancer
\item
  \textbf{H19:} Reservoir for miR-675; tumour suppressor function;
  imprinted, as described in the genomic-imprinting section.
\item
  \textbf{NEAT1:} Paraspeckle scaffold; regulates gene expression by
  nuclear retention of specific mRNAs
\end{itemize}

\subsubsection{Small Interfering RNAs (siRNAs) and
piRNAs}\label{small-interfering-rnas-sirnas-and-pirnas}

\textbf{siRNAs:} 21--23 nt, perfect complementarity to target; processed
by DICER from long double-stranded RNA (dsRNA). In plants and nematodes,
siRNA pathways mediate transposon silencing and antiviral immunity. In
mammals, dsRNA triggers interferon responses rather than siRNA pathways
in somatic cells; siRNA silencing is more prominent in germ cells and
stem cells.

\textbf{piRNAs (PIWI-interacting RNAs):} 26--31 nt; DICER-independent
(processed by ``ping-pong'' amplification cycle involving PIWI clade
Argonautes: PIWIL1/MILI, PIWIL4/MIWI2). Essential for silencing
transposable elements in the germline. Loss of piRNA pathway in
\emph{Drosophila} or mice causes transposon derepression and
infertility.

\begin{center}\rule{0.5\linewidth}{0.5pt}\end{center}

\subsection{Epigenetic Reprogramming and
Inheritance}\label{epigenetic-reprogramming-and-inheritance}

\subsubsection{Mitotic Heritability of Epigenetic Marks --- Detailed
Mechanism}\label{mitotic-heritability-of-epigenetic-marks-detailed-mechanism}

For an epigenetic mark to be ``heritable,'' it must survive DNA
replication and cell division. Different marks have different
mechanisms:

{\def\LTcaptype{none} % do not increment counter
\begin{longtable}[]{@{}
  >{\raggedright\arraybackslash}p{(\linewidth - 4\tabcolsep) * \real{0.0645}}
  >{\raggedright\arraybackslash}p{(\linewidth - 4\tabcolsep) * \real{0.4839}}
  >{\raggedright\arraybackslash}p{(\linewidth - 4\tabcolsep) * \real{0.4516}}@{}}
\toprule\noalign{}
\begin{minipage}[b]{\linewidth}\raggedright
Mark
\end{minipage} & \begin{minipage}[b]{\linewidth}\raggedright
Mitotic heritability mechanism
\end{minipage} & \begin{minipage}[b]{\linewidth}\raggedright
Half-life through divisions
\end{minipage} \\
\midrule\noalign{}
\endhead
\bottomrule\noalign{}
\endlastfoot
5mC at CpGs & DNMT1/UHRF1 maintenance at replication fork & High
(\textgreater10 divisions, about 95\% efficiency) \\
H3K27me3 & Read-write feedback (PRC2 EED reads, EZH2 writes); inherited
on parental nucleosomes which are randomly distributed to both daughter
strands & Medium; restored within 1--2 divisions \\
H3K9me3 & Suv39H1/H2 reads HP1 \texorpdfstring{\ensuremath{\to}}{->} recruits more Suv39H; DNMT3A coupling &
High at constitutive heterochromatin \\
H3K4me3 & Inherited via parental nucleosomes; rapidly re-established by
transcription & Medium \\
H3K27ac & Diluted by half each division; restored by transcription
factor activity & Low; not stably heritable \\
H2AK119ub & Inherited via parental nucleosomes; restored by PRC1 &
Medium \\
\end{longtable}
}

\textbf{Key conceptual point --- parental nucleosome recycling:} During
DNA replication, parental nucleosomes are split between the two daughter
strands. The histone chaperones \textbf{MCM2} (binds H3-H4 dimers via
its H3-H4 binding domain), \textbf{ASF1} (general H3-H4 chaperone), and
\textbf{FACT} (Spt16-SSRP1, helps recycle H2A-H2B during transcription
and replication) mediate this recycling. The current model:

\begin{enumerate}
\def\labelenumi{\arabic{enumi}.}
\tightlist
\item
  \textbf{Parental nucleosome split:} As the replisome opens, parental
  histones are transiently released from DNA.
\item
  \textbf{Asymmetric recycling:} H3-H4 tetramers preferentially deposit
  on the \textbf{leading strand} (via Pol\texorpdfstring{\ensuremath{\varepsilon}}{epsilon}/MCM2 interaction) or
  \textbf{lagging strand} (via Pol\texorpdfstring{\ensuremath{\alpha}}{alpha}/MCM2-7 helicase-associated MCM2).
  Recent work shows this is biased toward the leading strand in
  proliferating cells (about 60:40), explaining
  strand-asymmetric mark inheritance.
\item
  \textbf{CAF-1 deposits new H3.1-H4:} The chromatin assembly factor
  CAF-1 (CHAF1A/B + RBBP4) is recruited by PCNA at the replication fork.
  CAF-1 deposits \textbf{new} H3.1-H4 tetramers (synthesized from
  S-phase histone gene bursts).
\item
  \textbf{Newly synthesised histones are deposited UNMODIFIED.} Daughter
  chromatin starts at half-density of any given mark.
\item
  \textbf{Restoration to full density depends on the read-write feedback
  loop} --- the existing parental marks recruit the writer enzyme, which
  copies the mark onto neighbouring new histones.
\end{enumerate}

\textbf{PRC2 propagation through replication.} EED reads H3K27me3 on a
parental nucleosome \texorpdfstring{\ensuremath{\to}}{->} allosterically activates EZH2 \texorpdfstring{\ensuremath{\to}}{->} EZH2 deposits
H3K27me3 on a neighbouring (newly assembled) nucleosome. Quantitative
imaging shows that PRC2 activity at recently replicated chromatin is
about 3-fold higher than at established chromatin ---
explaining how the mark ``fills in'' the daughter strand within
about 6 hours of fork passage.

\textbf{DNMT1/UHRF1 mechanism in detail at the fork.} UHRF1 is loaded
onto the replication fork by PCNA. UHRF1 SRA binds hemi-methylated CpG
(the parental strand carries 5mC; the daughter strand has unmodified C).
The UHRF1 RING domain ubiquitinates H3K18 --- creating a docking site
for DNMT1's RFTS domain. DNMT1 transfers a methyl group from SAM to the
daughter cytosine. This coupled mechanism ensures methylation is
restored within seconds of fork passage at active replication.

\subsubsection{Germline Reprogramming}\label{germline-reprogramming}

Somatic epigenetic marks must be erased and re-established in each
generation to prevent transmission of acquired somatic states:

\begin{enumerate}
\def\labelenumi{\arabic{enumi}.}
\tightlist
\item
  \textbf{Post-fertilisation reprogramming:} After fertilisation, the
  paternal genome undergoes rapid active demethylation (TET3-mediated
  5mC oxidation) within hours. The maternal genome is demethylated more
  slowly (replication-dependent passive demethylation). Both reach a
  methylation minimum at the blastocyst stage.
\item
  \textbf{Primordial germ cell (PGC) reprogramming:} PGCs migrate to the
  gonads (about  E7.5--E10.5 in mouse; about  weeks
  3--5 in human). They erase CpG methylation genome-wide (including
  imprint control regions) --- the most complete demethylation in the
  mammalian life cycle.
\item
  \textbf{Re-establishment:} DNMT3A/3B with DNMT3L re-methylate the
  genome in a sex-specific pattern during gametogenesis
  (prospermatogonia in males; oocyte growth in females). Imprinted loci
  are methylated in a sex-specific order: paternal imprints in
  spermatogonia (before meiosis); maternal imprints in growing oocytes
  (after meiosis I arrest, prior to ovulation).
\end{enumerate}

\subsubsection{Evidence for Transgenerational Epigenetic Inheritance in
Humans}\label{evidence-for-transgenerational-epigenetic-inheritance-in-humans}

The Dutch Hunger Winter (Hongerwinter, 1944--1945) provides the most
studied human evidence. Dutch civilians subjected to severe famine
(500--1,000 kcal/day) during WWII German occupation showed:

\begin{itemize}
\tightlist
\item
  \textbf{F1 offspring} (exposed in utero): Increased rates of obesity,
  diabetes, schizophrenia, and CVD in adult life --- consistent with
  developmental programming via altered methylation.
\item
  \textbf{F2 offspring} (children of F1): Increased rates of obesity and
  metabolic syndrome --- suggesting transmission across one germline
  generation.
\item
  \textbf{Mechanism:} Reduced methylation at IGF2 differentially
  methylated regions (DMRs) persists for decades in blood cells of F1
  individuals, detectable compared to unexposed siblings.
\end{itemize}

\textbf{Other lines of evidence:} - Överkalix cohort (Sweden): Paternal
grandfather's food supply during slow-growth period correlates with
grandsons' diabetes mortality. - \emph{Agouti} viable yellow
(\(A^{vy}\)) mouse model: Maternal methyl-donor diet (folate, B12,
methionine, choline) during pregnancy shifts coat-colour distribution by
altering methylation of an upstream IAP retrotransposon. -
piRNA-mediated transposon silencing: piRNAs in sperm carry information
about active transposons across generations.

\begin{quote}
{[}!NOTE{]} The evidence for true transgenerational epigenetic
inheritance in humans (affecting F2 and beyond without continued
environmental exposure) is suggestive but not yet definitive.
Confounding by shared post-natal environment and direct exposure of the
F1 germline (the F2 was a primordial germ cell in the F0 grandmother
during the F1 in utero exposure) remains difficult to exclude.
Mechanistically, piRNA-mediated transposon silencing and small-RNA
transmission in sperm provide plausible vehicles, but a clean causal
chain in humans has not been established.
\end{quote}

\begin{quote}
\textbf{Concept Check 7:} Imagine a dCas9 fusion protein is targeted to
the \emph{BRCA1} promoter in a breast-cancer cell line, carrying a
DNMT3A catalytic domain. Over several cell divisions, the promoter
becomes hypermethylated and transcription drops. Once dCas9 is
withdrawn, does the methylation persist, relax over mitoses, or become
heritable between cells? Explain in terms of DNMT1 maintenance
methylation, passive demethylation during S-phase, and the absence of
active demethylation by TET enzymes under these conditions.
\end{quote}

\begin{quote}
\textbf{Concept Check 8:} A scientist treats cells with the DNMT
inhibitor decitabine for 48 hours, then washes it out. Predicting that
DNMT1 will resume normal activity after washout, will the demethylation
persist, partially recover, or fully recover over the subsequent 10 cell
generations? Sketch a quantitative model assuming \texorpdfstring{\ensuremath{\varepsilon}}{epsilon} returns to 0.95
instantaneously after washout and starting \emph{f\texorpdfstring{\textsubscript{0}}{0}} = 0.10 (10 \%
methylated post-decitabine).
\end{quote}

\begin{center}\rule{0.5\linewidth}{0.5pt}\end{center}

\subsection{Cancer Epigenetics and Clinical Translation ---
Comprehensive}\label{cancer-epigenetics-and-clinical-translation-comprehensive}

Cancer cells display \textbf{systemic epigenetic dysregulation}:
typically global DNA hypomethylation (especially of repetitive
elements), focal hypermethylation of tumour-suppressor CGI promoters,
broad H3K27me3/H3K9me3 redistribution, and aberrant chromatin
remodelling complex composition. The therapeutic implication is that,
unlike genetic mutations, epigenetic states are \textbf{reversible} ---
and several drug classes now target each layer.

\subsubsection{The Cancer Epigenome --- Two Concurrent
Patterns}\label{the-cancer-epigenome-two-concurrent-patterns}

\textbf{Pattern 1: Global hypomethylation.} Repetitive elements (LINE-1,
SINE, IAP, satellite repeats), normally heavily methylated, lose
methylation in nearly most cancers. Consequences: - \textbf{Chromosomal
instability:} demethylation of pericentromeric satellite-2 sequences
disrupts heterochromatin, leading to chromosome bridges, lagging
chromosomes, and aneuploidy. - \textbf{LINE-1 reactivation:}
transposition of L1 elements can disrupt tumour-suppressor genes (e.g.,
L1 insertion in \emph{APC} contributes to colorectal cancer). -
\textbf{Mis-regulated lineage genes:} demethylation of normally
repressed germline genes (e.g., MAGE, NY-ESO-1) creates
\textbf{cancer-testis antigens} that can be exploited for immunotherapy
(e.g., NY-ESO-1 TCR-T cells). - \textbf{Mechanism:} likely loss of DNMT1
maintenance (UHRF1 dysregulation) and/or genome-wide TET2-mediated
demethylation in CHIP-derived clones.

\textbf{Pattern 2: Focal CGI hypermethylation.} Tumour-suppressor gene
promoters become hypermethylated, silencing them. Examples and clinical
relevance:

{\def\LTcaptype{none} % do not increment counter
\begin{longtable}[]{@{}
  >{\raggedright\arraybackslash}p{(\linewidth - 6\tabcolsep) * \real{0.0656}}
  >{\raggedright\arraybackslash}p{(\linewidth - 6\tabcolsep) * \real{0.2951}}
  >{\raggedright\arraybackslash}p{(\linewidth - 6\tabcolsep) * \real{0.3607}}
  >{\raggedright\arraybackslash}p{(\linewidth - 6\tabcolsep) * \real{0.2787}}@{}}
\toprule\noalign{}
\begin{minipage}[b]{\linewidth}\raggedright
Gene
\end{minipage} & \begin{minipage}[b]{\linewidth}\raggedright
Hypermethylated in
\end{minipage} & \begin{minipage}[b]{\linewidth}\raggedright
Functional consequence
\end{minipage} & \begin{minipage}[b]{\linewidth}\raggedright
Therapeutic angle
\end{minipage} \\
\midrule\noalign{}
\endhead
\bottomrule\noalign{}
\endlastfoot
\emph{MLH1} & CIMP-high colorectal cancer & Microsatellite instability &
Pembrolizumab in MSI-high tumours \\
\emph{BRCA1} & Triple-negative breast cancer & Homologous-recombination
deficiency & PARP inhibitors \\
\emph{VHL} & Renal cell carcinoma & HIF stabilisation & HIF-2\texorpdfstring{\ensuremath{\alpha}}{alpha}
inhibitors (belzutifan) \\
\emph{CDKN2A/p16} & Many cancers & RB1 pathway loss \texorpdfstring{\ensuremath{\to}}{->} unrestrained G1/S
& CDK4/6 inhibitors \\
\emph{MGMT} & Glioblastoma & DNA repair loss & Temozolomide
(sensitivity) \\
\emph{DAPK} & Lung, head/neck cancer & Apoptosis loss & Combination with
TRAIL agonists (clinical trials) \\
\end{longtable}
}

The two patterns can coexist --- a single tumour can have global
hypomethylation of repeats AND focal hypermethylation of tumour
suppressors. The mechanism appears to be partial: CGI-targeted DNMT3A/B
activity becomes mis-localised to actively transcribed regions, while
genome-wide DNMT1 maintenance fails at heterochromatin.

\subsubsection{Approved and Investigational Epigenetic Drugs ---
Comprehensive
Table}\label{approved-and-investigational-epigenetic-drugs-comprehensive-table}

{\def\LTcaptype{none} % do not increment counter
\begin{longtable}[]{@{}
  >{\raggedright\arraybackslash}p{(\linewidth - 8\tabcolsep) * \real{0.0769}}
  >{\raggedright\arraybackslash}p{(\linewidth - 8\tabcolsep) * \real{0.1231}}
  >{\raggedright\arraybackslash}p{(\linewidth - 8\tabcolsep) * \real{0.1077}}
  >{\raggedright\arraybackslash}p{(\linewidth - 8\tabcolsep) * \real{0.3077}}
  >{\raggedright\arraybackslash}p{(\linewidth - 8\tabcolsep) * \real{0.3846}}@{}}
\toprule\noalign{}
\begin{minipage}[b]{\linewidth}\raggedright
Class
\end{minipage} & \begin{minipage}[b]{\linewidth}\raggedright
Examples
\end{minipage} & \begin{minipage}[b]{\linewidth}\raggedright
Targets
\end{minipage} & \begin{minipage}[b]{\linewidth}\raggedright
Approved Indications
\end{minipage} & \begin{minipage}[b]{\linewidth}\raggedright
Clinical-trial highlights
\end{minipage} \\
\midrule\noalign{}
\endhead
\bottomrule\noalign{}
\endlastfoot
\textbf{DNMT inhibitors} & Azacitidine (Vidaza), decitabine (Dacogen);
oral cedazuridine/decitabine (Inqovi) & DNMT1 (covalent trap) & MDS, AML
(low-intensity), CMML & Aza + venetoclax: ORR 70 \% in elderly AML
(VIALE-A); decitabine-cedazuridine in CMML \\
\textbf{HDAC inhibitors} & Vorinostat (SAHA), romidepsin, panobinostat,
belinostat, tucidinostat & Pan-HDAC or class-selective & CTCL, PTCL,
multiple myeloma (panobinostat) & Romidepsin ORR 34 \% in PTCL;
panobinostat-bortezomib-dex in MM \\
\textbf{EZH2 / PRC2 inhibitors} & Tazemetostat, valemetostat & EZH2 SET
domain (and EZH1) & Epithelioid sarcoma (INI1-loss); EZH2-mut follicular
lymphoma; ATL (valemetostat) & Tazemetostat: median PFS 5.5 mo in
INI1-loss ES \\
\textbf{IDH1/IDH2 inhibitors} & Ivosidenib (IDH1), enasidenib (IDH2),
vorasidenib & Mutant IDH neomorphic activity (2-HG) & IDH-mut AML;
IDH-mut grade-2 glioma (vorasidenib, 2024) & Vorasidenib INDIGO trial:
median PFS 27 vs 11 mo \\
\textbf{BET inhibitors} & OTX015 (birabresib), molibresib, mivebresib,
BMS-986158 & BRD2/3/4 bromodomains & Investigational: NUT carcinoma,
MYC-driven tumours & NUT-carcinoma ORR about 35 \%; pancreatic
cancer combination trials \\
\textbf{LSD1 inhibitors} & Tranylcypromine, iadademstat, bomedemstat &
KDM1A/LSD1 demethylase & Investigational: AML, SCLC, MPN & Iadademstat +
azacitidine in AML; bomedemstat in essential thrombocythaemia \\
\textbf{Menin--MLL inhibitors} & Revumenib (FDA 2024) & Menin--MLL/KMT2A
interface & Approved KMT2A-rearranged AML & Revumenib AUGMENT-101 trial:
ORR 53 \%, CR 23 \% \\
\textbf{PROTAC bromodomain degraders} & ARV-771, dBET6, ARV-825 & BRD4
(degradation) & Investigational & Targets most three BET family members;
deeper BRD4 depletion than inhibition \\
\textbf{DOT1L inhibitors} & Pinometostat (EPZ-5676) & H3K79
methyltransferase & Investigational MLL-rearranged AML & Phase 2 ORR
about 10 \% monotherapy; combination strategies in trial \\
\textbf{CDK7/9 inhibitors} & THZ1, SY-5609 (CDK7); AZD4573 (CDK9) &
Transcription-cycle CDKs & Investigational & Disrupts
super-enhancer-driven oncogene transcription \\
\textbf{Combination DNMTi + HDACi or DNMTi + venetoclax} & Aza +
venetoclax (FDA 2018) & DNMT1 + BCL2 & Standard-of-care AML in unfit
patients & Aza-venetoclax: median OS 14.7 mo vs 9.6 mo aza alone \\
\end{longtable}
}

\subsubsection{Clinical Trial Vignettes --- Three Paradigm
Therapies}\label{clinical-trial-vignettes-three-paradigm-therapies}

\textbf{(1) Tazemetostat for INI1-loss epithelioid sarcoma.} In 2020,
tazemetostat became the first epigenetic-pathway-defined targeted
therapy. INI1/SMARCB1-loss epithelioid sarcoma has no chemotherapy or
surgical option in metastatic setting; the agent reactivates
SWI/SNF-suppressed targets via EZH2 inhibition. Pivotal trial
(NCT02601950): 62 patients enrolled, ORR 15 \%, median time to response
3.6 mo, median PFS 5.5 mo. The slow response kinetics reflect epigenetic
reprogramming requiring multiple cell divisions. FDA approval (Jan 2020)
was a landmark for synthetic-lethal epigenetic therapy
\citep{fda2020tazverik}.

\textbf{(2) Vorasidenib for IDH1-mutant grade-2 glioma.} In 2024,
vorasidenib became the first targeted therapy approved for grade-2
glioma. IDH1-R132H mutation produces oncometabolite 2-HG, which inhibits
TET demethylases \texorpdfstring{\ensuremath{\to}}{->} CIMP+ phenotype \texorpdfstring{\ensuremath{\to}}{->} genome-wide hypermethylation \texorpdfstring{\ensuremath{\to}}{->}
differentiation block. Vorasidenib is a brain-penetrant mutant-IDH1
inhibitor. INDIGO trial (NCT04164901): 331 patients with grade-2 IDH-mut
glioma post-surgery, median PFS 27.7 mo (vorasidenib) vs 11.1 mo
(placebo); 85 \% reduction in time-to-next-intervention. The trial
demonstrated that depleting 2-HG is sufficient to restore demethylation
and slow tumour growth without cytotoxicity.

\textbf{(3) Revumenib for KMT2A-rearranged AML.} In 2024, revumenib
became the first menin--MLL inhibitor approved by FDA. KMT2A-rearranged
(ex-MLL-rearranged) AML accounts for about 10 \% of paediatric
AML and about 5 \% of adult AML; menin is required for
MLL-fusion-driven HOXA9 expression. AUGMENT-101 trial (NCT04065399): ORR
53 \% in relapsed/refractory KMT2Ar AML, CR 23 \%, median DoR 6.4 mo.
The drug demonstrated that targeting protein--protein interaction
(menin-MLL) within an aberrant chromatin complex is therapeutically
viable.

\subsubsection{A Checklist for Interpreting Cancer
Epigenomes}\label{a-checklist-for-interpreting-cancer-epigenomes}

When you encounter a tumour \textbf{multi-omics} report, integrate:

\begin{enumerate}
\def\labelenumi{\arabic{enumi}.}
\tightlist
\item
  \textbf{DNA methylation} --- focal CGI hypermethylation (TSG
  silencing) vs.~global hypomethylation (genomic instability); CIMP (CpG
  island methylator phenotype) as a class.
\item
  \textbf{Histone marks} --- \textbf{H3K27me3} (PRC2 silencing),
  \textbf{H3K4me3} (promoter activity), \textbf{H3K27ac} (active
  enhancers), \textbf{H3K36me3} (transcribed gene bodies; loss in renal
  cell carcinoma with SETD2 mutation).
\item
  \textbf{Chromatin accessibility} (ATAC-seq) --- nucleosome-depleted
  regions at enhancers; super-enhancers identified by Mediator/H3K27ac
  density.
\item
  \textbf{3D genome} (Hi-C) --- TAD boundary loss, \textbf{enhancer
  hijacking}.
\item
  \textbf{Histone mutations (oncohistones)} --- H3K27M (paediatric
  DIPG), H3G34R/V (paediatric high-grade glioma), H3K36M
  (chondroblastoma): these are \emph{gain-of-function} mutants that act
  as PRC2 / SETD2 inhibitors in trans, producing genome-wide chromatin
  redistribution despite occurring in about 5--10 \% of nucleosomes.
\item
  \textbf{Chromatin-remodeller mutations} --- ARID1A, SMARCA4, SMARCB1,
  PBRM1 (SWI/SNF); KMT2C/D, EZH2 (PRC2).
\item
  \textbf{RNA} --- fusion oncoproteins that recruit aberrant chromatin
  complexes (\textbf{MLL} fusions to AF4/AF9/ENL, \textbf{SS18--SSX} in
  synovial sarcoma, \textbf{BRD4--NUT} in NUT carcinoma).
\end{enumerate}

This checklist mirrors how clinicians and trialists rationalise
\textbf{DNMT + HDAC} combinations, \textbf{EZH2} monotherapy in defined
molecular subtypes, and emerging \textbf{BET + BCL2} regimens in
haematological malignancies.

\begin{quote}
\textbf{Worked Example 5 --- CpG Island Methylation and Transcriptional
Silencing}

\textbf{Problem:} A molecular biologist treats cancer cells with
5-azacytidine (azacitidine), a DNMT inhibitor. They observe
re-expression of the tumour suppressor gene \emph{CDKN2B} (p15). Before
treatment, bisulfite sequencing showed 90\% CpG methylation at the
\emph{CDKN2B} promoter CGI. The drug is incorporated into DNA during
replication and forms a covalent adduct with DNMT1, preventing
maintenance methylation.

\textbf{(a)} After 3 cell divisions with the drug present, what
methylation level would you predict at the CGI if no de novo methylation
occurs?

\textbf{(b)} Why do HDAC inhibitors (e.g., vorinostat) often synergise
with DNMT inhibitors in cancer treatment?

\textbf{Solution:}

\textbf{(a)} Maintenance methylation requires DNMT1 to copy
hemi-methylated CpGs after replication. With DNMT1 blocked, after each
division sister strands receive no methylation:

\begin{equation}f_n = f_0 \times \left(\frac{1}{2}\right)^n \tag{15.7}\label{eq:f_n}\end{equation}

\begin{equation}f_3 = 0.90 \times \left(\frac{1}{2}\right)^3 = 0.90 \times 0.125 = 0.1125 \approx 11\% \tag{15.8}\label{eq:f_3}\end{equation}

After 3 divisions: about 11\% methylation (down from 90\%).
Since about 50\% methylation is the threshold for silencing at
most CGIs, \emph{CDKN2B} should be significantly re-expressed.

\textbf{(b)} Two distinct repressive layers maintain CGI-silenced genes:
1. \textbf{CpG methylation} \texorpdfstring{\ensuremath{\to}}{->} methyl-CpG binding proteins (MBD2) recruit
HDAC \texorpdfstring{\ensuremath{\to}}{->} deacetylated histones 2. \textbf{Deacetylated histones} \texorpdfstring{\ensuremath{\to}}{->} compact
chromatin, transcription factor exclusion

DNMT inhibition removes layer 1 but may not immediately open chromatin
if deacetylated histones persist. HDAC inhibitors (remove layer 2)
accumulate acetylation independently of methylation. \textbf{Together},
both layers are removed simultaneously \texorpdfstring{\ensuremath{\to}}{->} synergistic re-activation of
silenced tumour suppressors.
\end{quote}

\begin{quote}
\textbf{Worked Example 6 --- IDH1 Mutation and the Methylator Phenotype}

\textbf{Problem:} A glioma tumour carries the IDH1-R132H mutation. The
mutant enzyme converts \texorpdfstring{\ensuremath{\alpha}}{alpha}-ketoglutarate (\texorpdfstring{\ensuremath{\alpha}}{alpha}KG) to 2-hydroxyglutarate
(2-HG), which competitively inhibits \texorpdfstring{\ensuremath{\alpha}}{alpha}KG-dependent dioxygenases
including TET demethylases, KDM Jumonji-domain demethylases, and prolyl
hydroxylases. Predict three genome-wide chromatin consequences of
IDH1-R132H, and explain why vorasidenib is therapeutic.

\textbf{Solution:} 1. \textbf{Global hypermethylation of CpG islands}
(CIMP+): TET inhibition prevents demethylation \texorpdfstring{\ensuremath{\to}}{->} progressive CpG
accumulation. 2. \textbf{H3K27me3 and H3K9me3 accumulation}: Jumonji
KDM6A/B and KDM4 inhibition prevents demethylation of repressive marks \texorpdfstring{\ensuremath{\to}}{->}
genome-wide silencing of differentiation genes (the ``differentiation
block'' of IDH-mut AML and glioma). 3. \textbf{HIF-1\texorpdfstring{\ensuremath{\alpha}}{alpha} stabilisation}:
prolyl hydroxylase inhibition stabilises HIF-1\texorpdfstring{\ensuremath{\alpha}}{alpha} \texorpdfstring{\ensuremath{\to}}{->} pseudohypoxia \texorpdfstring{\ensuremath{\to}}{->}
angiogenesis, glycolysis upregulation.

Vorasidenib (FDA 2024 for IDH1-mut grade-2 glioma) crosses the
blood--brain barrier and inhibits mutant IDH1, depleting 2-HG and
partially reversing TET/KDM/PHD inhibition --- restoring demethylation
and re-engaging differentiation. Median progression-free survival
doubled (27 vs.~11 months) in the INDIGO trial.
\end{quote}

\begin{center}\rule{0.5\linewidth}{0.5pt}\end{center}

\subsection{Computational Bridge}\label{computational-bridge-14}

Promoter CpG islands are GC-rich; summarising composition links
chromatin biology to sequence statistics:

\begin{Shaded}
\begin{Highlighting}[]
\ImportTok{from}\NormalTok{ biology.genetics }\ImportTok{import}\NormalTok{ gc\_content}

\NormalTok{cgi\_like }\OperatorTok{=} \StringTok{"CGCGCGCGCGATATATAT"}
\BuiltInTok{print}\NormalTok{(}\BuiltInTok{round}\NormalTok{(gc\_content(cgi\_like), }\DecValTok{3}\NormalTok{))}
\end{Highlighting}
\end{Shaded}

Methylation decay through cell divisions is a one-line simulation:

\begin{Shaded}
\begin{Highlighting}[]
\KeywordTok{def}\NormalTok{ methylation\_decay(f0: }\BuiltInTok{float}\NormalTok{, generations: }\BuiltInTok{int}\NormalTok{, eps: }\BuiltInTok{float} \OperatorTok{=} \FloatTok{0.95}\NormalTok{) }\OperatorTok{{-}\textgreater{}} \BuiltInTok{list}\NormalTok{[}\BuiltInTok{float}\NormalTok{]:}
    \CommentTok{"""Predict CpG methylation fraction over \textasciigrave{}generations\textasciigrave{} divisions}
\CommentTok{    given maintenance efficiency \textasciigrave{}eps\textasciigrave{}.}

\CommentTok{    f\_\{n+1\} = (eps + (1 {-} eps) / 2) * f\_n  (per{-}CpG approximation)}
\CommentTok{    """}
\NormalTok{    factor }\OperatorTok{=}\NormalTok{ eps }\OperatorTok{+}\NormalTok{ (}\DecValTok{1} \OperatorTok{{-}}\NormalTok{ eps) }\OperatorTok{/} \DecValTok{2}
    \ControlFlowTok{return}\NormalTok{ [f0 }\OperatorTok{*}\NormalTok{ factor }\OperatorTok{**}\NormalTok{ n }\ControlFlowTok{for}\NormalTok{ n }\KeywordTok{in} \BuiltInTok{range}\NormalTok{(generations }\OperatorTok{+} \DecValTok{1}\NormalTok{)]}

\BuiltInTok{print}\NormalTok{([}\BuiltInTok{round}\NormalTok{(v, }\DecValTok{3}\NormalTok{) }\ControlFlowTok{for}\NormalTok{ v }\KeywordTok{in}\NormalTok{ methylation\_decay(}\FloatTok{0.90}\NormalTok{, }\DecValTok{5}\NormalTok{, eps}\OperatorTok{=}\FloatTok{0.0}\NormalTok{)])  }\CommentTok{\# DNMT1{-}null}
\CommentTok{\# [0.9, 0.45, 0.225, 0.113, 0.056, 0.028]}
\end{Highlighting}
\end{Shaded}

Setting \texorpdfstring{\ensuremath{\varepsilon}}{epsilon} = 0 simulates a DNMT1-null (or azacitidine-trapped) cell ---
the fraction halves each generation.

A Polycomb-spreading simulator with the EED--allosteric read-write
feedback:

\begin{Shaded}
\begin{Highlighting}[]
\KeywordTok{def}\NormalTok{ polycomb\_spread(seeds: }\BuiltInTok{int}\NormalTok{, generations: }\BuiltInTok{int}\NormalTok{, k: }\BuiltInTok{float}\NormalTok{, n\_nucleosomes: }\BuiltInTok{int} \OperatorTok{=} \DecValTok{50}\NormalTok{) }\OperatorTok{{-}\textgreater{}} \BuiltInTok{list}\NormalTok{[}\BuiltInTok{float}\NormalTok{]:}
    \CommentTok{"""Simulate H3K27me3 occupancy over generations.}

\CommentTok{    seeds: initial number of H3K27me3{-}marked nucleosomes}
\CommentTok{    k: marks deposited per generation per existing mark (read{-}write coupling)}
\CommentTok{    n\_nucleosomes: total nucleosomes available}
\CommentTok{    """}
\NormalTok{    history }\OperatorTok{=}\NormalTok{ [seeds]}
\NormalTok{    m }\OperatorTok{=}\NormalTok{ seeds}
    \ControlFlowTok{for}\NormalTok{ \_ }\KeywordTok{in} \BuiltInTok{range}\NormalTok{(generations):}
\NormalTok{        m }\OperatorTok{=}\NormalTok{ m }\OperatorTok{/} \DecValTok{2} \OperatorTok{+}\NormalTok{ k          }\CommentTok{\# passive halving + spreading}
\NormalTok{        m }\OperatorTok{=} \BuiltInTok{min}\NormalTok{(m, n\_nucleosomes)  }\CommentTok{\# saturation}
\NormalTok{        history.append(m)}
    \ControlFlowTok{return}\NormalTok{ history}

\BuiltInTok{print}\NormalTok{(polycomb\_spread(seeds}\OperatorTok{=}\DecValTok{2}\NormalTok{, generations}\OperatorTok{=}\DecValTok{5}\NormalTok{, k}\OperatorTok{=}\DecValTok{10}\NormalTok{))}
\CommentTok{\# [2, 11.0, 16.5, 18.25, 19.125, 19.5625]  → reaches about 20 marked nucleosomes}
\end{Highlighting}
\end{Shaded}

\begin{quote}
\textbf{Clinical / systems note:} Azacitidine/decitabine and HDAC
inhibitors are approved where malignant clones depend on silenced
differentiation programmes --- therapies that \textbf{reverse layers of
repression} rather than mutating oncogenes. Newer entries (tazemetostat
for EZH2; menin--MLL inhibitors for KMT2A-rearranged AML; vorasidenib
for IDH-mut glioma) extend the same logic to writers, readers, and
metabolic upstream regulators.
\end{quote}

\begin{center}\rule{0.5\linewidth}{0.5pt}\end{center}

\subsection{Current Evidence and Frontier
Biology}\label{current-evidence-and-frontier-biology-17}

For \textbf{Epigenetics and Gene Regulation}, frontier biology belongs
inside the evidence logic of the chapter. Molecular genetics now spans
single-reference sequences, telomere-to-telomere assemblies, pangenome
graphs, long-read sequencing, CRISPR medicines, and ethical deployment.
The core reading question is this: epigenetic claims require causal
perturbation, cell-type specificity, timing, and inheritance controls.

\textbf{Single-cell and spatial genomics boundary:} Single-cell RNA-seq,
single-cell ATAC-seq, immune profiling, and spatial transcriptomics now
let students ask which cell type and tissue neighbourhood carries a
regulatory claim rather than averaging the whole organ. The Human Cell
Atlas framed the goal as a reference map of human cell types, states,
locations, and lineages \citep{regev2017humancellatlas}; newer
multi-omic atlas work integrates single-cell modalities with spatial
transcriptomics across many tissues \citep{pan2024singlecellatlas}. The
evidence boundary is essential: a cluster marker or spatial
neighbourhood is a hypothesis about cell identity and interaction unless
validated by perturbation, orthogonal markers, sampling design, and
disease-relevant tissue context.

\begin{itemize}
\tightlist
\item
  \textbf{What to verify:} identify the observation, model, assay, or
  dataset that would make the claim stronger or weaker.
\item
  \textbf{What to qualify:} state the scale, organism, cell type,
  environmental condition, or population where the claim is expected to
  hold.
\item
  \textbf{What to compare:} test at least one alternative explanation,
  baseline, or null model before treating the pattern as causal.
\item
  \textbf{What to cite:} distinguish primary evidence, review synthesis,
  public dataset, and institutional guidance; for recent or numeric
  claims, prefer the source closest to the measurement and state what
  has changed since it was published.
\end{itemize}

When a genomic claim depends on a reference, ask whether short reads,
structural variants, ancestry representation, phasing, or clinical
validation could change the interpretation
\citep{humanpangenome2023,fda2023casgevy,fda2024casgevythalassemia}.

\textbf{Source practice:} For genomics and editing claims, distinguish
discovery from clinical actionability, and cite reference resources,
regulatory records, or primary editing studies close to the claim
\citep{humanpangenome2023,fda2026casgevy,chalumeau2025primeediting}.

\subsubsection{Current Evidence Map: Epigenetic Causality
Ladder}\label{current-evidence-map-epigenetic-causality-ladder}

\begin{figure}[htbp]
\centering
\includegraphics[width=0.82\linewidth,height=4.2in,keepaspectratio]{../figures/mermaid_inline/inline_mermaid_0086_4edc8987e9ae.png}
\caption{An epigenetic mark is not automatically a cause; causal claims need perturbation, timing, cell-type specificity, expression readout, and rescue or orthogonal evidence.}
\end{figure}

\subsection{Summary}\label{summary-17}

\begin{itemize}
\tightlist
\item
  \textbf{Chromatin} is organised around the nucleosome (147 bp DNA +
  histone octamer H2A/H2B/H3/H4). Histone variants (H3.3, CENP-A, H2A.Z,
  H2A.X, macroH2A) define specialised chromatin states. Higher-order
  compaction creates TADs (CTCF/cohesin) and A/B compartments.
\item
  \textbf{Histone acetylation} (HATs; activating, e.g., H3K27ac, H3K9ac,
  H4K16ac) and \textbf{deacetylation} (HDACs; repressive) regulate
  accessibility. Newer acyl marks (lactylation, \texorpdfstring{\ensuremath{\beta}}{beta}-hydroxybutyrylation,
  crotonylation, butyrylation) couple chromatin to metabolism.
\item
  \textbf{Histone methylation} has context-dependent effects. H3K4me3 =
  active promoters; H3K36me3 = gene bodies; H3K27me3 = Polycomb (PRC2)
  silencing; H3K9me3 = constitutive heterochromatin (HP1); H4K20me3 =
  pericentric heterochromatin. Bivalent domains (H3K4me3 + H3K27me3)
  mark poised developmental genes in ES cells.
\item
  \textbf{Polycomb / Trithorax cellular memory:} PRC2 (EZH2) writes
  H3K27me3 \texorpdfstring{\ensuremath{\to}}{->} recruits PRC1 (RING1B \texorpdfstring{\ensuremath{\to}}{->} H2AK119ub). MLL/COMPASS writes
  H3K4me3; UTX/JMJD3 erase H3K27me3; SWI/SNF evicts PRC1/2 --- together
  antagonising Polycomb at active genes.
\item
  \textbf{DNA methylation} at CpG islands (DNMT3A/3B \emph{de novo};
  DNMT1+UHRF1 maintenance) causes stable gene silencing. \textbf{TET
  enzymes} oxidise 5mC \texorpdfstring{\ensuremath{\to}}{->} 5hmC \texorpdfstring{\ensuremath{\to}}{->} 5fC \texorpdfstring{\ensuremath{\to}}{->} 5caC for active demethylation;
  5hmC is itself a stable mark. IDH1/2 mutations produce 2-HG, blocking
  TET and creating the CIMP phenotype.
\item
  \textbf{Genomic imprinting:} about 150 genes from
  about 30 ICR-controlled clusters expressed monoallelically by
  parent of origin. IGF2/H19 (CTCF-insulator), 15q11--q13
  (Prader--Willi/Angelman), 11p15 BWS/SRS exemplify mechanisms and
  reciprocal phenotypes.
\item
  \textbf{X-chromosome inactivation:} XIST lncRNA (cis-coating) \texorpdfstring{\ensuremath{\to}}{->}
  SPEN-HDAC3 + PRC2 \texorpdfstring{\ensuremath{\to}}{->} H3K27me3 + DNMT3 \texorpdfstring{\ensuremath{\to}}{->} CGI methylation;
  about 15--25 \% escape genes (mostly PAR). Skewed XCI
  explains variable X-linked disease in heterozygous females. iPSC
  reprogramming reactivates Xi.
\item
  \textbf{Chromatin remodellers:} SWI/SNF (BAF/PBAF --- slide+eject;
  mutated in about 20 \% cancers), ISWI (spacing), CHD/NuRD
  (repressive sliding + HDAC), INO80/SWR1 (H2A.Z exchange).
\item
  \textbf{3D genome:} TADs by loop extrusion (CTCF + cohesin); A/B
  compartments; LADs at the lamina; phase-separated condensates at
  super-enhancers (BRD4, Mediator, Pol II CTD); CDK8 dissolves
  condensates by phosphorylating Mediator.
\item
  \textbf{miRNA / lncRNA / piRNA:} DROSHA \texorpdfstring{\ensuremath{\to}}{->} DICER \texorpdfstring{\ensuremath{\to}}{->} RISC for miRNA;
  lncRNAs (XIST, HOTAIR, MALAT1, NEAT1) scaffold complexes; piRNAs
  silence transposons in germline.
\item
  \textbf{Mitotic / germline reprogramming:} Maintenance via DNMT1+UHRF1
  (H3K18ub-coupled) and read-write feedback for histone marks. CAF-1
  deposits new histones at the fork. Two waves of epigenetic erasure
  (post-fertilisation; PGC). Imprints protected from PGC erasure and
  re-set sex-specifically during gametogenesis.
\item
  \textbf{Cancer epigenetics \& clinical translation:} Approved drugs
  target each layer --- DNMTi, HDACi, EZH2i, IDH-mut inhibitors, BET
  inhibitors, menin--MLL inhibitors, oncohistone H3K27M biology,
  condensate-targeting agents.
\item
  \textbf{Connections:} See \cref{sec:unit_IV_gene_expression} for
  transcriptional output, \cref{sec:unit_IV_mutations_and_genomics} for
  CRISPR-based epigenome editing, Unit V for imprinting and pedigree
  patterns, and Unit IX for nuclear receptor chromatin targeting.
\end{itemize}

\begin{center}\rule{0.5\linewidth}{0.5pt}\end{center}

\subsection{Key Terms}\label{key-terms-17}

{\def\LTcaptype{none} % do not increment counter
\begin{longtable}[]{@{}
  >{\raggedright\arraybackslash}p{(\linewidth - 2\tabcolsep) * \real{0.2857}}
  >{\raggedright\arraybackslash}p{(\linewidth - 2\tabcolsep) * \real{0.7143}}@{}}
\toprule\noalign{}
\begin{minipage}[b]{\linewidth}\raggedright
Term
\end{minipage} & \begin{minipage}[b]{\linewidth}\raggedright
Definition
\end{minipage} \\
\midrule\noalign{}
\endhead
\bottomrule\noalign{}
\endlastfoot
\textbf{Epigenetics} & Heritable changes in gene expression that do not
alter the DNA sequence \\
\textbf{Nucleosome} & 147 bp DNA + histone octamer (H2A, H2B, H3, H4 \texorpdfstring{\ensuremath{\times}}{x}
2); fundamental chromatin unit \\
\textbf{Histone variant} & Non-canonical histone (H3.3, CENP-A, H2A.Z,
H2A.X, macroH2A) deposited at specialised loci \\
\textbf{Histone code} & Combinatorial histone modifications specify
chromatin states and downstream regulatory outcomes \\
\textbf{HAT / HDAC} & Histone acetyltransferase (writer) / deacetylase
(eraser); regulate H3K9/H3K14/H3K27/H4K16 acetylation \\
\textbf{H3K27me3} & PRC2/EZH2-deposited Polycomb repressive mark \\
\textbf{H3K4me3} & MLL/SET1-deposited active promoter mark \\
\textbf{H3K9me3} & SUV39H/SETDB1-deposited constitutive heterochromatin
mark; HP1 reader \\
\textbf{H3K36me3} & SETD2-deposited transcribed gene body mark \\
\textbf{PRC1 / PRC2} & Polycomb complexes; PRC2 writes H3K27me3 (EZH2),
PRC1 writes H2AK119ub (RING1B) \\
\textbf{Trithorax (TrxG)} & MLL/COMPASS, KDM6A/B, BAF --- antagonise
Polycomb \\
\textbf{CpG island} & about 500 bp--2 kb region of high CpG
density; methylation at promoter CGIs silences genes \\
\textbf{DNMT1} & Maintenance DNA methyltransferase; copies
hemimethylated CpGs at replication fork (UHRF1-recruited) \\
\textbf{DNMT3A/3B} & De novo methyltransferases; require DNMT3L
cofactor \\
\textbf{UHRF1} & Hemi-methyl-CpG reader (SRA); H3K9me reader (TTD);
recruits DNMT1 to fork \\
\textbf{TET enzymes} & Fe(II)/\texorpdfstring{\ensuremath{\alpha}}{alpha}KG-dependent dioxygenases; oxidise 5mC \texorpdfstring{\ensuremath{\to}}{->}
5hmC \texorpdfstring{\ensuremath{\to}}{->} 5fC \texorpdfstring{\ensuremath{\to}}{->} 5caC for active demethylation \\
\textbf{5-methylcytosine} & The ``fifth base''; methylated cytosine at
CpG dinucleotides \\
\textbf{5hmC} & 5-hydroxymethylcytosine; TET oxidation product; stable
mark in neurons \\
\textbf{Genomic imprinting} & Parent-of-origin-specific monoallelic gene
expression; about 150 human genes \\
\textbf{ICR / DMR} & Imprint control region / differentially methylated
region \\
\textbf{XIST} & X-inactive specific transcript; 17 kb lncRNA; coats
inactive X chromosome in cis \\
\textbf{Skewed XCI} & Non-50/50 X-inactivation distribution; explains
variable X-linked disease in females \\
\textbf{miRNA} & about 22 nt RNA; DROSHA \texorpdfstring{\ensuremath{\to}}{->} DICER pathway;
AGO2/RISC \texorpdfstring{\ensuremath{\to}}{->} translational repression/mRNA decay \\
\textbf{piRNA} & 26--31 nt germline RNA; PIWI pathway; silences
transposable elements \\
\textbf{TAD} & Topologically Associating Domain; chromatin loop defined
by CTCF/cohesin; enhancer-promoter contacts \\
\textbf{A/B compartment} & Megabase-scale active vs.~repressed chromatin
neighbourhoods \\
\textbf{LAD} & Lamina-Associated Domain; gene-poor regions tethered to
nuclear periphery \\
\textbf{LLPS / condensate} & Liquid-liquid phase separation;
concentrates IDR-containing regulators at hubs \\
\textbf{Super-enhancer} & Large enhancer cluster with dense
Mediator/BRD4/H3K27ac; drives high transcriptional output \\
\textbf{Bivalent domain} & Simultaneous H3K4me3 + H3K27me3 at
developmental gene promoters in ES cells \\
\textbf{SWI/SNF (BAF)} & ATP-dependent remodeller; mutated in
about 20 \% cancers (ARID1A, SMARCB1, PBRM1) \\
\textbf{ISWI / CHD/NuRD / INO80} & Other remodeller families: spacing /
repressive sliding+HDAC / H2A.Z exchange \\
\textbf{CAF-1} & Chromatin assembly factor 1; deposits new H3.1-H4 at
replication fork via PCNA \\
\textbf{Azacitidine / decitabine} & DNMT inhibitors; trap DNMT1;
approved for MDS/AML \\
\textbf{Vorinostat (SAHA)} & HDAC inhibitor; approved CTCL; synergises
with DNMT-i \\
\textbf{Tazemetostat} & EZH2 inhibitor; FDA-approved for SMARCB1-loss
epithelioid sarcoma and EZH2-mut FL \\
\textbf{BET inhibitor} & Blocks bromodomain--acetyl-lysine engagement
(e.g., BRD4); collapses oncogenic transcriptional condensates \\
\textbf{Vorasidenib} & Mutant-IDH1 inhibitor; FDA-approved 2024 for
grade-2 glioma; depletes 2-HG, restores TET activity \\
\textbf{Revumenib} & Menin--MLL inhibitor; FDA-approved 2024 for
KMT2A-rearranged AML \\
\textbf{Oncohistone} & Recurrent gain-of-function histone mutation
(H3K27M, H3G34R/V, H3K36M) acting in trans \\
\textbf{CIMP} & CpG-island methylator phenotype; cluster of
hypermethylated CGIs (IDH-mut glioma, MLH1-silent CRC) \\
\end{longtable}
}

\begin{center}\rule{0.5\linewidth}{0.5pt}\end{center}

\subsection{Discussion Questions}\label{discussion-questions-4}

These open-ended questions require integration across multiple sections
and connection to clinical or experimental practice. Suitable for
graduate-style discussion or take-home essay.

\begin{enumerate}
\def\labelenumi{\arabic{enumi}.}
\item
  \textbf{Cellular memory under perturbation.} PRC2 (H3K27me3) is
  regularly described as a ``cellular memory'' system. Experimentally,
  how would you distinguish whether Polycomb memory at a developmental
  gene is maintained by (a) the read-write feedback loop alone, (b) the
  parental nucleosome inheritance mechanism, or (c) DNA methylation
  acting downstream? Propose two CRISPR-based perturbations that would
  discriminate between these mechanisms.
\item
  \textbf{Aging and the epigenetic clock.} Steve Horvath's epigenetic
  clock predicts chronological age from CpG methylation patterns at
  about 353 sites with median error of 3.6 years. Why should
  DNA methylation drift with age? Propose three biological mechanisms
  (mitotic-error, oxidative-stress, lineage-replacement). What
  therapeutic implications would follow if epigenetic age could be
  reversed (e.g., via Yamanaka factors, partial reprogramming)?
\item
  \textbf{The ``junk DNA'' reckoning and 80 \% functional ENCODE claim.}
  ENCODE (2012) claimed \texorpdfstring{\textgreater{}=}{>=} 80 \% of the human genome has biochemical
  activity (chromatin accessibility, TF binding, histone marks, RNA
  transcription). Critics argue most of these activities are ``noise''
  --- neutral binding to GC-rich sequences. Take a position: design an
  experiment that would distinguish biologically functional enhancer
  activity from ``spurious'' biochemical activity at a candidate
  non-coding region.
\item
  \textbf{Synthetic biology with phase separation.} Engineering
  programmable transcription factors with IDR-mediated condensate
  formation could provide programmable control over transcription rate.
  Propose three design principles --- and three failure modes --- for an
  engineered IDR--dCas9 fusion that would create a tunable
  super-enhancer-like condensate.
\item
  \textbf{Trans-generational epigenetic inheritance: real or artefact?}
  Critique the human evidence (Dutch Hunger Winter, Överkalix, Holocaust
  trauma cohorts). What four experimental designs would conclusively
  demonstrate transgenerational inheritance through sperm/oocyte (rather
  than direct exposure of F1 germline)? Why is this question
  particularly difficult in mammals compared to \emph{C. elegans} or
  \emph{Arabidopsis}?
\item
  \textbf{Clinical sequencing and the epigenome.} A 70-year-old woman
  has DNMT3A R882H clonal-haematopoiesis (CHIP) detected on routine
  whole-exome sequencing. She has no AML and no cytopenia. As her
  physician, what counselling would you provide regarding: (i) AML risk,
  (ii) cardiovascular risk (CHIP elevates CHD risk independent of AML),
  (iii) lifestyle modifications, (iv) whether prophylactic treatment
  should be considered, and (v) what monitoring frequency you would
  recommend?
\end{enumerate}

\begin{center}\rule{0.5\linewidth}{0.5pt}\end{center}

\subsection{Review Questions}\label{review-questions-17}

\begin{enumerate}
\def\labelenumi{\arabic{enumi}.}
\tightlist
\item
  A gene has 95\% CpG methylation at its promoter CGI. A researcher
  applies 5-azacytidine for 5 cell divisions with no de novo
  methylation. Using the formula \(f_n = f_0 / 2^n\), what methylation
  level remains? If the transcription threshold is 30\%, at which
  division does the gene first become potentially active?
\item
  Explain why H3K27me3 is a \textbf{repressive} mark while H3K4me3 is an
  \textbf{activating} mark, even though both involve lysine methylation
  on histone H3. Include the relevant enzymes (HMT and KDM) and reader
  proteins for each.
\item
  Draw and annotate the IGF2/H19 imprinting mechanism. Explain why a
  maternal deletion at the ICR/DMR causes Beckwith-Wiedemann syndrome
  (biallelic IGF2 expression), while a paternal deletion causes
  Silver-Russell syndrome (biallelic IGF2 silencing).
\item
  A female patient presents with Rett syndrome. Genetic testing shows a
  missense mutation in \emph{MECP2} (methyl-CpG binding protein 2) on
  one X chromosome. Explain why Rett syndrome shows variable severity
  between affected females, using the concept of X-chromosome
  inactivation skewing.
\item
  A cancer biopsy shows the following ChIP-seq results at a
  tumour-suppressor gene promoter: H3K27me3 high, H3K4me3 absent,
  H3K27ac absent, and 85\% CpG methylation. A second tumour from the
  same patient shows H3K27me3 low, H3K27ac present, H3K4me3 high, and
  5\% methylation --- with the tumour-suppressor gene transcribed.
  Propose two mechanisms by which the epigenetic state could have
  switched, and describe the specific enzymes involved.
\item
  Compare piRNAs and miRNAs in terms of: (a) size, (b) biogenesis
  pathway (DICER-dependent vs.~independent), (c) primary biological
  function in mammals, and (d) what happens when the pathway is
  disrupted genetically.
\item
  A developmental biologist shows that a \emph{Drosophila} gene is in a
  ``bivalent domain'' in neural stem cells (H3K4me3 + H3K27me3). After
  neural differentiation, H3K27me3 is removed by KDM6A. What molecular
  events would you expect at this gene promoter? What histone
  remodelling complex would facilitate nucleosome repositioning for full
  activation?
\item
  The Dutch Hunger Winter cohort shows that caloric restriction in
  pregnant women led to measurable methylation differences at the IGF2
  DMR in their children, persisting for 60 years. Critically evaluate
  this as evidence for transgenerational epigenetic inheritance, noting
  the key confounders and what additional data would conclusively
  demonstrate germline transmission.
\item
  A clinical trial combines azacitidine (DNMT inhibitor) with vorinostat
  (HDAC inhibitor) for treating a myeloid cancer. Using your knowledge
  of the two repressive layers at silenced CGI promoters, predict which
  order of drug administration would be more effective. Support your
  reasoning with the molecular mechanism.
\item
  CRISPR-dCas9 (catalytically dead Cas9) fused to DNMT3A can write DNA
  methylation at specific genomic loci. Design an experiment using this
  tool to test whether the methylation at the \emph{BRCA1} promoter in a
  cancer cell line is the \textbf{cause} of silencing (not merely a
  consequence). Describe the expected result if methylation is causal
  vs.~if there is an upstream histone modification driving silencing.
\item
  Explain how a \textbf{BET bromodomain inhibitor} might reduce
  \emph{MYC} transcription without mutating the \emph{MYC} locus. How
  would you test whether the drug acts via loss of BRD4 at
  super-enhancers rather than a global transcription shutdown?
\item
  A paediatric tumour harbours a \textbf{CTCF} site deletion that fuses
  two \textbf{TADs}. Propose how this structural variant could activate
  an oncogene --- and one experiment (Hi-C, CRISPR excision rescue, or
  luciferase reporter) to support enhancer hijacking.
\item
  A patient with myelodysplastic syndrome has a TET2 loss-of-function
  mutation. Predict the genome-wide methylation phenotype, and explain
  why this patient might respond to azacitidine \emph{better} than a
  TET2-wild-type patient.
\item
  A diffuse intrinsic pontine glioma harbours an H3K27M oncohistone.
  Explain why about 10 \% of nucleosomes carrying this mutation
  can produce a genome-wide H3K27me3 reduction (gain-of-function effect
  on PRC2 in trans). What is the rationale for EZH2 inhibitors in this
  setting?
\item
  Quantitatively compare the half-lives of H3K27ac and 5mC marks across
  cell divisions. Why does this difference imply that long-term cellular
  memory is encoded primarily in DNA methylation and Polycomb marks, not
  acetylation?
\item
  A glioma is diagnosed with IDH1-R132H mutation. Plasma 2-HG is 200 \texorpdfstring{\ensuremath{\mu}}{mu}M
  (vs.~\textless{} 10 \texorpdfstring{\ensuremath{\mu}}{mu}M in healthy individuals). The patient is treated
  with vorasidenib (mutant-IDH1 inhibitor). After 3 months, plasma 2-HG
  drops to 12 \texorpdfstring{\ensuremath{\mu}}{mu}M. Predict the time-course of DNA methylation reversal at
  CIMP+ loci. Why does clinical response (radiographic improvement)
  typically lag 6--12 months behind 2-HG normalisation?
\item
  A laboratory generates a mouse with a homozygous deletion of the
  \emph{Cbx7} gene (encoding a PRC1 chromobox subunit). Predict the
  haematopoietic phenotype. Would you expect the phenotype to be more or
  less severe than a \emph{Bmi1} (PCGF4) homozygous deletion? Justify
  your answer in terms of redundancy among canonical-PRC1 paralogs.
\item
  A 35-year-old woman is found to have CHIP with a \emph{DNMT3A} R882H
  mutation at 15 \% VAF. Three years later she develops AML. Sequencing
  of the AML shows the same DNMT3A R882H plus a new \emph{FLT3-ITD}
  mutation. Discuss the implications: (i) why is DNMT3A loss often a
  CHIP-precursor mutation rather than the AML-driver itself? (ii) How
  does R882H differ from a complete loss-of-function? (iii) What
  treatment would target both the epigenetic and signalling layers?
\end{enumerate}

\begin{center}\rule{0.5\linewidth}{0.5pt}\end{center}

\subsection{Further Reading and Source
Notes}\label{further-reading-and-source-notes-17}

\begin{itemize}
\tightlist
\item
  Strahl \& Allis (2000). The language of covalent histone
  modifications. \emph{Nature}, 403.
\item
  Lyon (1961). Gene action in the X-chromosome of the mouse (\emph{Mus
  musculus L.}). \emph{Nature}, 190.
\item
  Fire et al.~(1998). Potent and specific genetic interference by
  double-stranded RNA in \emph{Caenorhabditis elegans}. \emph{Nature},
  391.
\item
  Doudna \& Charpentier (2014). The new frontier of genome engineering
  with CRISPR-Cas9. \emph{Science}, 346.\emph{Figures:
  \breaktt{src/visualization/plots.py} ---
  \breaktt{plot\_methylation\_heatmap()}}
\end{itemize}

\begin{center}\rule{0.5\linewidth}{0.5pt}\end{center}

\subsubsection{Companion Source
Module}\label{companion-source-module-17}

\textbf{Epigenetics and Gene Regulation} should leave a reproducible
trail from a biological claim to the code, figure, diagram, or
paper-based activity that can test it. Use the surfaces below to inspect
the chapter's assumptions, rerun the relevant model, or compare the
manuscript explanation with companion labs and figures.

{\def\LTcaptype{none} % do not increment counter
\begin{longtable}[]{@{}
  >{\raggedright\arraybackslash}p{(\linewidth - 2\tabcolsep) * \real{0.5000}}
  >{\raggedright\arraybackslash}p{(\linewidth - 2\tabcolsep) * \real{0.5000}}@{}}
\toprule\noalign{}
\begin{minipage}[b]{\linewidth}\raggedright
Surface
\end{minipage} & \begin{minipage}[b]{\linewidth}\raggedright
Use it for
\end{minipage} \\
\midrule\noalign{}
\endhead
\bottomrule\noalign{}
\endlastfoot
\breaktt{src/biology/genetics/genetics.py}
(\breaktt{cpg\_methylation\_remaining},
\breaktt{histone\_modification\_state}) & Convert methylation maintenance
and histone-mark claims into explicit state checks. \\
\breaktt{src/visualization/plots.py}
(\breaktt{plot\_methylation\_heatmap}) & Inspect whether heatmap
interpretation depends on color alone or includes labels. \\
\breaktt{src/mermaid/biology\_diagrams.py}
(\breaktt{mirna\_biogenesis\_diagram}, \breaktt{x\_inactivation\_diagram})
& Compare RNA-mediated and chromatin-mediated regulation. \\
\end{longtable}
}

\textbf{Reproducibility check:} require cell type, developmental time,
perturbation evidence, and inheritance control before calling a mark
causal. \textbf{Cross-reference:} use
\cref{sec:unit_IV_gene_expression},
\cref{sec:unit_IV_mutations_and_genomics}, and
\cref{sec:unit_V_chromosomal_inheritance}.

\newpage

\section*{Unit V --- Classical Genetics and Heredity:
Introduction}\label{unit-v-classical-genetics-and-heredity-introduction}
\addcontentsline{toc}{section}{Unit V --- Classical Genetics and
Heredity: Introduction}

\subsection*{Why This Unit Matters}\label{why-this-unit-matters-5}
\addcontentsline{toc}{subsection}{Why This Unit Matters}

In 1866, an Augustinian friar in Brno published a mathematical analysis
of 29,000 pea plants he had grown in the monastery garden. Gregor Mendel
showed that heritable traits are passed from parent to offspring in
discrete, predictable ratios --- 3:1 for monohybrid crosses, 9:3:3:1 for
dihybrid crosses. His paper, "Versuche über Pflanzenhybriden," went
unread for 34 years. When it was rediscovered simultaneously by de
Vries, Correns, and Tschermak in 1900, it ignited a revolution: the laws
of segregation and independent assortment described the behaviour of
something that Mendel had rarely seen --- chromosomes, structures not
visualised until the 1880s, and connected to genes primarily by Thomas
Hunt Morgan's fruit fly experiments in 1910.

The synthesis of Mendel's mathematical laws with chromosomal theory, and
later with population genetics, constitutes one of the great
intellectual achievements of twentieth-century science. It unified
patterns of inheritance visible to naked-eye observation with the
molecular machinery of DNA replication and meiosis. It transformed
agriculture, medicine, forensic science, and evolutionary theory
simultaneously. Today, genome-wide association studies (GWAS) identify
thousands of loci contributing to human traits; pedigree analysis
diagnoses rare recessive disorders; and population genetics models
predict drug-resistance evolution in pathogens.

This unit bridges the gap between the molecular genetics of Unit IV and
the evolutionary dynamics of Unit VI. You will derive Mendelian ratios
from first principles, work through the mathematics of incomplete
dominance, codominance, and epistasis, map gene positions using
recombination frequencies, and apply the Hardy-Weinberg equations to
natural populations.

\begin{center}\rule{0.5\linewidth}{0.5pt}\end{center}

\subsection*{Landmark Discoveries}\label{landmark-discoveries-5}
\addcontentsline{toc}{subsection}{Landmark Discoveries}

{\def\LTcaptype{none} % do not increment counter
\begin{longtable}[]{@{}
  >{\raggedright\arraybackslash}p{(\linewidth - 8\tabcolsep) * \real{0.2407}}
  >{\raggedright\arraybackslash}p{(\linewidth - 8\tabcolsep) * \real{0.0741}}
  >{\raggedright\arraybackslash}p{(\linewidth - 8\tabcolsep) * \real{0.2963}}
  >{\raggedright\arraybackslash}p{(\linewidth - 8\tabcolsep) * \real{0.1667}}
  >{\raggedright\arraybackslash}p{(\linewidth - 8\tabcolsep) * \real{0.2222}}@{}}
\toprule\noalign{}
\begin{minipage}[b]{\linewidth}\raggedright
Discoverer(s)
\end{minipage} & \begin{minipage}[b]{\linewidth}\raggedright
Year
\end{minipage} & \begin{minipage}[b]{\linewidth}\raggedright
Journal / Source
\end{minipage} & \begin{minipage}[b]{\linewidth}\raggedright
Discovery
\end{minipage} & \begin{minipage}[b]{\linewidth}\raggedright
Significance
\end{minipage} \\
\midrule\noalign{}
\endhead
\bottomrule\noalign{}
\endlastfoot
Gregor Mendel & 1866 & \emph{Verh. naturf. Vereines Brünn} & Laws of
segregation and independent assortment & Mathematical foundation of
genetics; rediscovered 1900 \\
Walther Flemming & 1882 & \emph{Zellsubstanz, Kern und Zelltheilung} &
Chromosomes during mitosis ("chromatin") & First visualisation of
chromosomal segregation \\
Thomas Hunt Morgan & 1910 & \emph{Science} & Sex-linked inheritance in
\emph{Drosophila} & Proved genes reside on chromosomes; Nobel Prize
1933 \\
Alfred Sturtevant & 1913 & \emph{J. Exp. Zool.} & First genetic map (6
X-linked \emph{Drosophila} genes) & Showed recombination frequency \texorpdfstring{\ensuremath{\propto}}{propto} map
distance \\
Barbara McClintock & 1950 & \emph{Proc. Natl. Acad. Sci.} & Transposable
elements ("jumping genes") & Showed genomes are dynamic; Nobel Prize
1983 \\
Hardy \& Weinberg & 1908 & \emph{Science; Jahreshefte d.~Ver.
f.~Vaterländ. Naturk.} & Population genetic equilibrium & Allele
frequencies stable without evolution; quantitative tool for
deviations \\
Human Genome Project Consortium & 2001 & \emph{Nature} & Draft human
genome sequence & about 3.2 Gb; about 20,000
protein-coding genes; foundation for GWAS \\
\end{longtable}
}

\begin{center}\rule{0.5\linewidth}{0.5pt}\end{center}

\subsection*{Key Concepts and
Connections}\label{key-concepts-and-connections-5}
\addcontentsline{toc}{subsection}{Key Concepts and Connections}

\begin{figure}[htbp]
\centering
\includegraphics[width=0.82\linewidth,height=4.2in,keepaspectratio]{../figures/mermaid_inline/inline_mermaid_0087_02dd2b2d5b31.png}
\caption{Flowchart for Key Concepts and Connections: 1 phenotypic ratio; Punnett square)"], pea genetics Mendel's Laws (Segregation + Independent Assortment), data Monohybrid Cross (3:1 phenotypic ratio; Punnett square), and 3:3:1; two unlinked genes)"] form the diagram's primary path or branches.}
\end{figure}

\textbf{Unit V concept map --- Classical Genetics and Heredity.}

\begin{center}\rule{0.5\linewidth}{0.5pt}\end{center}

\subsection*{Current Evidence Thread}\label{current-evidence-thread-5}
\addcontentsline{toc}{subsection}{Current Evidence Thread}

Treat inheritance in this unit as something that is \emph{evidenced},
not just asserted: pedigrees and controlled crosses test transmission
rules, allele-frequency time series test which evolutionary force is
acting, and biobank-scale association data test whether a `Mendelian'
trait is really single-gene. Classical genetics remains essential, but
modern interpretation adds penetrance, polygenicity, structural
variation, ancestry-aware inference, and uncertainty in risk prediction.
As you move through the chapters, keep a two-column note: \textbf{claim}
on the left, \textbf{evidence that would change my confidence} on the
right. By the end of the unit, each major idea should be tied to a
measurement, model, citation, or paper-based lab decision.

\subsection*{Chapter Roadmap}\label{chapter-roadmap-5}
\addcontentsline{toc}{subsection}{Chapter Roadmap}

{\def\LTcaptype{none} % do not increment counter
\begin{longtable}[]{@{}
  >{\raggedright\arraybackslash}p{(\linewidth - 6\tabcolsep) * \real{0.1556}}
  >{\raggedright\arraybackslash}p{(\linewidth - 6\tabcolsep) * \real{0.1111}}
  >{\raggedright\arraybackslash}p{(\linewidth - 6\tabcolsep) * \real{0.2889}}
  >{\raggedright\arraybackslash}p{(\linewidth - 6\tabcolsep) * \real{0.4444}}@{}}
\toprule\noalign{}
\begin{minipage}[b]{\linewidth}\raggedright
Chapter
\end{minipage} & \begin{minipage}[b]{\linewidth}\raggedright
Title
\end{minipage} & \begin{minipage}[b]{\linewidth}\raggedright
Core Question
\end{minipage} & \begin{minipage}[b]{\linewidth}\raggedright
Key Equation / Model
\end{minipage} \\
\midrule\noalign{}
\endhead
\bottomrule\noalign{}
\endlastfoot
\textbf{16} & Mendelian Genetics & How do Mendel's laws explain the
inheritance of discrete traits? & Binomial expansion; chi-square test
for goodness of fit \\
\textbf{17} & Chromosomal Inheritance & How do chromosomes carry genes,
and what happens when chromosomal segregation errors occur? &
Recombination frequency \texorpdfstring{\ensuremath{\theta}}{theta}; map function; trisomy probability \\
\textbf{18} & Population Genetics & How do allele frequencies change ---
or not change --- in populations? & Hardy-Weinberg:
\(p^2 + 2pq + q^2 = 1\); \(\Delta p\) under selection; \(F_{ST}\) \\
\end{longtable}
}

\begin{center}\rule{0.5\linewidth}{0.5pt}\end{center}

\subsection*{Connections Across the
Textbook}\label{connections-across-the-textbook-5}
\addcontentsline{toc}{subsection}{Connections Across the Textbook}

\begin{itemize}
\tightlist
\item
  \textbf{Meiosis and crossing-over}
  (\cref{sec:unit_V_chromosomal_inheritance}) build directly on DNA
  replication and chromosome structure from Unit IV, and provide the
  mechanistic basis for recombination in Unit VI (evolution).
\item
  \textbf{Hardy-Weinberg equilibrium}
  (\cref{sec:unit_V_population_genetics}) is the null model for both
  natural selection (\cref{sec:unit_VI_evolution_and_selection}) and
  genetic drift (\cref{sec:unit_VI_genetic_drift_and_speciation}).
\item
  \textbf{Pedigree analysis} and \textbf{inheritance patterns} reappear
  throughout Unit IX (genetic basis of cardiovascular disease, endocrine
  disorders, immune deficiencies).
\item
  \textbf{Epistasis and polygenic traits} link to Unit VI (quantitative
  trait loci, QTL mapping) and Unit X (population-level phenotypic
  variation in ecology).
\end{itemize}

\begin{quote}
\textbf{Key vocabulary introduced here:} allele, locus, genotype,
phenotype, dominant, recessive, homozygous, heterozygous, test cross,
chi-square test, centimorgan (cM), Hardy-Weinberg equilibrium,
inbreeding coefficient, linkage disequilibrium, aneuploidy,
nondisjunction.
\end{quote}

\subsection*{Computational Toolbox --- Unit
V}\label{computational-toolbox-unit-v}
\addcontentsline{toc}{subsection}{Computational Toolbox --- Unit V}

\begin{Shaded}
\begin{Highlighting}[]
\ImportTok{from}\NormalTok{ biology.genetics }\ImportTok{import}\NormalTok{ hardy\_weinberg, hamming\_distance, chi\_squared\_test}

\CommentTok{\# Hardy{-}Weinberg equilibrium: cystic fibrosis (autosomal recessive)}
\CommentTok{\# Carrier frequency q ≈ 1/22 in Northern Europeans → q ≈ 0.045, p ≈ 0.955}
\NormalTok{hw }\OperatorTok{=}\NormalTok{ hardy\_weinberg(p}\OperatorTok{=}\FloatTok{0.955}\NormalTok{, q}\OperatorTok{=}\FloatTok{0.045}\NormalTok{)}
\BuiltInTok{print}\NormalTok{(}\SpecialStringTok{f"AA (unaffected non{-}carrier): }\SpecialCharTok{\{}\NormalTok{hw}\SpecialCharTok{.}\NormalTok{p\_squared}\SpecialCharTok{:.4f\}}\SpecialStringTok{ = }\SpecialCharTok{\{}\NormalTok{hw}\SpecialCharTok{.}\NormalTok{p\_squared}\OperatorTok{*}\DecValTok{100}\SpecialCharTok{:.1f\}}\SpecialStringTok{\%"}\NormalTok{)}
\BuiltInTok{print}\NormalTok{(}\SpecialStringTok{f"Aa (carrier):               }\SpecialCharTok{\{}\NormalTok{hw}\SpecialCharTok{.}\NormalTok{two\_pq}\SpecialCharTok{:.4f\}}\SpecialStringTok{ = }\SpecialCharTok{\{}\NormalTok{hw}\SpecialCharTok{.}\NormalTok{two\_pq}\OperatorTok{*}\DecValTok{100}\SpecialCharTok{:.1f\}}\SpecialStringTok{\%"}\NormalTok{)}
\BuiltInTok{print}\NormalTok{(}\SpecialStringTok{f"aa (affected):              }\SpecialCharTok{\{}\NormalTok{hw}\SpecialCharTok{.}\NormalTok{q\_squared}\SpecialCharTok{:.5f\}}\SpecialStringTok{ = 1 in }\SpecialCharTok{\{}\DecValTok{1}\OperatorTok{/}\NormalTok{hw}\SpecialCharTok{.}\NormalTok{q\_squared}\SpecialCharTok{:.0f\}}\SpecialStringTok{"}\NormalTok{)}
\CommentTok{\# Expected:}
\CommentTok{\# AA (unaffected non{-}carrier): 0.9120 = 91.2\%}
\CommentTok{\# Aa (carrier):               0.0860 = 8.6\%  (≈ 1 in 23 — matches observed carrier rate)}
\CommentTok{\# aa (affected):              0.00203 = 1 in 493}

\CommentTok{\# Sequence distance between two short alleles}
\NormalTok{dist }\OperatorTok{=}\NormalTok{ hamming\_distance(}\StringTok{"ATGCTAGC"}\NormalTok{, }\StringTok{"ATGATAGT"}\NormalTok{)}
\BuiltInTok{print}\NormalTok{(}\SpecialStringTok{f"Sequence differences: }\SpecialCharTok{\{}\NormalTok{dist}\SpecialCharTok{\}}\SpecialStringTok{"}\NormalTok{)}
\CommentTok{\# Expected: Sequence differences: 2}

\CommentTok{\# Chi{-}square test: 3:1 vs observed 290:110 (n=400 plants)}
\NormalTok{chi }\OperatorTok{=}\NormalTok{ chi\_squared\_test(observed}\OperatorTok{=}\NormalTok{[}\DecValTok{290}\NormalTok{, }\DecValTok{110}\NormalTok{], expected}\OperatorTok{=}\NormalTok{[}\DecValTok{300}\NormalTok{, }\DecValTok{100}\NormalTok{])}
\BuiltInTok{print}\NormalTok{(}\SpecialStringTok{f"χ² = }\SpecialCharTok{\{}\NormalTok{chi}\SpecialCharTok{.}\NormalTok{chi\_squared}\SpecialCharTok{:.2f\}}\SpecialStringTok{, p ≈ }\SpecialCharTok{\{}\NormalTok{chi}\SpecialCharTok{.}\NormalTok{p\_value\_approx}\SpecialCharTok{:.3f\}}\SpecialStringTok{"}\NormalTok{)}
\CommentTok{\# Expected: χ² ≈ 1.33, p ≈ 0.25 (not significant; consistent with 3:1)}
\end{Highlighting}
\end{Shaded}

\begin{quote}
\textbf{Try it yourself:} Change carrier frequency to
\texttt{q\ =\ 0.02} (rare disease). What fraction of affected
individuals are born to two carrier parents?
\end{quote}

\begin{center}\rule{0.5\linewidth}{0.5pt}\end{center}

\emph{Source note: genetics helpers support Punnett squares,
Hardy-Weinberg checks, sequence distances, and chi-squared tests.}
\emph{Figures: \breaktt{src/visualization/} (Mendelian ratio histograms,
Hardy-Weinberg surface plots).}

\newpage

\section{Mendelian Genetics and
Heredity}\label{mendelian-genetics-and-heredity}

\label{sec:unit_V_mendelian_genetics}

\begin{quote}
\textbf{Ch 16} \texorpdfstring{\ensuremath{\cdot}}{.} Level 2/3 \texorpdfstring{\ensuremath{\cdot}}{.} 65 min read \texorpdfstring{\ensuremath{\cdot}}{.} 100 min lecture \texorpdfstring{\ensuremath{\cdot}}{.}
Prerequisites: \cref{sec:unit_IV_dna_replication_and_cell_cycle}
\end{quote}

\subsection{Learning Objectives}\label{learning-objectives-18}

\begin{enumerate}
\def\labelenumi{\arabic{enumi}.}
\tightlist
\item
  Describe Mendel \citep{mendel1866}'s experimental system and explain
  why \emph{Pisum sativum} was ideal for genetic analysis.
\item
  State Mendel's laws of segregation and
  \hyperref[gl:independent-assortment]{\textbf{independent assortment}}
  and explain their mechanistic basis in
  \hyperref[gl:meiosis]{\textbf{meiosis}}.
\item
  Solve monohybrid, dihybrid, and trihybrid cross problems using Punnett
  squares and probability.
\item
  Describe extensions to Mendelian genetics: incomplete dominance,
  codominance, multiple \hyperref[gl:allele]{\textbf{allele}}s,
  pleiotropy, epistasis, polygenic inheritance, genomic imprinting, and
  maternal effect.
\item
  Explain penetrance, expressivity, and phenocopy with disease examples.
\item
  Apply the chi-squared test for goodness of fit to evaluate Mendelian
  ratios.
\item
  Construct and interpret pedigrees for autosomal
  \hyperref[gl:dominant]{\textbf{dominant}}, autosomal recessive,
  X-linked recessive, X-linked dominant, and mitochondrial inheritance
  patterns.
\item
  Solve multi-step genetics problems involving probability rules
  (product rule, sum rule, conditional probability).
\item
  Connect Mendelian segregation to chromosome behavior at meiosis and
  explain the molecular basis of human autosomal-dominant and
  autosomal-recessive disorders.
\item
  Outline how genome-wide association studies (GWAS) extend Mendelian
  thinking to polygenic traits in human populations.
\end{enumerate}

\begin{figure}[htbp]
\centering
\includegraphics[width=0.85\textwidth]{../figures/punnett_AaxAa.png}
\caption{Punnett square for a monohybrid cross Aa \texorpdfstring{\ensuremath{\times}}{x} Aa. Each cell shows the zygote genotype; the 3:1 phenotype ratio follows directly from the 1:2:1 genotype ratio.}
\label{fig:unit_V_punnett_square}
\end{figure}

\subsubsection{Study Blueprint}\label{study-blueprint-18}

\begin{itemize}
\tightlist
\item
  \textbf{Big idea:} Inheritance patterns arise from chromosome
  behaviour, probability, and genotype-to-phenotype mapping.
\item
  \textbf{Core concepts:} segregation, independent assortment,
  pedigrees, epistasis.
\item
  \textbf{Framework alignment:} Vision \& Change: Information flow,
  exchange, and storage, Evolution; AP Biology: Information Storage and
  Transmission, Evolution; NGSS-style topics: Inheritance and Variation
  of Traits, Natural Selection and Evolution.
\item
  \textbf{Model or quantitative lens:} Punnett, binomial, pedigree, and
  chi-square calculations.
\item
  \textbf{Data skill:} Use family or cross data to infer genotype
  probabilities.
\item
  \textbf{Practice cadence:} Statistical Tests and Data Analysis,
  Representing and Describing Data.
\item
  \textbf{Common misconception to repair:} Dominant does not mean
  common, stronger, or better.
\item
  \textbf{Primary lab:} \cref{sec:lab_unit_V_mendelian_genetics}.
\item
  \textbf{Question bank:} \cref{sec:q_unit_V_mendelian_genetics}.
\item
  \textbf{Transfer task:} Transfer inheritance reasoning to counseling,
  breeding, model organisms, and risk prediction.
\item
  \textbf{Bridge to computation:}
  \breaktt{biology.genetics.genetics.chi\_squared\_test}.
\end{itemize}

\begin{center}\rule{0.5\linewidth}{0.5pt}\end{center}

\begin{quote}
\textbf{Opening Vignette --- The Forgotten Friar and the Laws of
Inheritance}

Gregor Mendel spent eight years (1856--1863) growing and
cross-pollinating nearly 28,000 pea plants in the garden of St.~Thomas's
Abbey in Brno, recording data with statistical rigour unprecedented in
biology at the time. He presented his findings to the Brno Natural
History Society in 1865, and published them in the society's journal in
1866. The audience of 40 scientists was politely unenthusiastic. For 35
years, his paper languished in obscurity --- unread, uncited,
misunderstood. Then in 1900, three botanists working independently ---
Hugo de Vries, Carl Correns, and Erich von Tschermak --- each
rediscovered the same ratios Mendel had documented and, searching the
literature, found his paper. Overnight, an unknown friar's meticulous
garden statistics became the foundation of modern genetics. Mendel
rarely knew he had founded a science.
\end{quote}

\subsection{Mendel's Experiments: Historical
Context}\label{mendels-experiments-historical-context}

Gregor Johann Mendel (1822-1884), an Augustinian friar and
physics-trained scientist, conducted breeding experiments on garden peas
(\emph{Pisum sativum}) from 1856 to 1863 in the monastery garden at Brno
(now Czech Republic). His 1866 paper, ``Versuche uber
Pflanzen-Hybriden'' (Experiments on Plant Hybrids), was largely ignored
until 1900 when it was independently ``rediscovered'' by de Vries,
Correns, and von Tschermak.

\textbf{Why peas were ideal:}

{\def\LTcaptype{none} % do not increment counter
\begin{longtable}[]{@{}
  >{\raggedright\arraybackslash}p{(\linewidth - 2\tabcolsep) * \real{0.4583}}
  >{\raggedright\arraybackslash}p{(\linewidth - 2\tabcolsep) * \real{0.5417}}@{}}
\toprule\noalign{}
\begin{minipage}[b]{\linewidth}\raggedright
Advantage
\end{minipage} & \begin{minipage}[b]{\linewidth}\raggedright
Significance
\end{minipage} \\
\midrule\noalign{}
\endhead
\bottomrule\noalign{}
\endlastfoot
Distinct, discontinuous traits & 7 characters with clear
dominant/recessive forms (no blending) \\
True-breeding lines available & Commercially available pure varieties;
years of inbreeding \\
Short generation time & about 3 months seed to seed \\
Self-compatible flowers & Could self-fertilize (controlled crosses
possible by emasculation) \\
Large number of offspring & Each pod contains about 7 seeds;
massive sample sizes \\
Easy to cross-pollinate & Large flowers with accessible reproductive
parts \\
\end{longtable}
}

\textbf{Mendel's seven characters:}

{\def\LTcaptype{none} % do not increment counter
\begin{longtable}[]{@{}lll@{}}
\toprule\noalign{}
Character & Dominant & Recessive \\
\midrule\noalign{}
\endhead
\bottomrule\noalign{}
\endlastfoot
Seed shape & Round & Wrinkled \\
Seed color & Yellow & Green \\
Flower color & Purple & White \\
Flower position & Axial & Terminal \\
Pod shape & Inflated & Constricted \\
Pod color & Green & Yellow \\
Stem length & Tall & Short \\
\end{longtable}
}

Mendel's genius was \textbf{quantitative}: he counted offspring in each
class and applied \textbf{binomial probability} reasoning -- the first
application of statistics to biology. He analyzed about 28,000
pea plants over 8 years.

\begin{center}\rule{0.5\linewidth}{0.5pt}\end{center}

\subsection{Mendel's Laws}\label{mendels-laws}

\subsubsection{Law of Segregation (Mendel's First
Law)}\label{law-of-segregation-mendels-first-law}

\textbf{Statement}: Each organism possesses two alleles for each
\hyperref[gl:gene]{\textbf{gene}}. These alleles \textbf{segregate}
(separate) during \hyperref[gl:gamete]{\textbf{gamete}} formation, such
that each gamete receives exactly one allele. Upon fertilization, the
diploid state is restored.

\textbf{Molecular basis}: Separation of homologous
\hyperref[gl:chromosome]{\textbf{chromosomes}} during
\textbf{\hyperref[gl:anaphase]{anaphase I of meiosis}}. The two alleles
of a gene reside on homologous chromosomes; when homologs are pulled to
opposite poles, each gamete receives one allele.

\begin{figure}[htbp]
\centering
\includegraphics[width=0.82\linewidth,height=4.2in,keepaspectratio]{../figures/mermaid_inline/inline_mermaid_0088_6f1f5e3cab2e.png}
\caption{Molecular basis of the Law of Segregation. A heterozygous (Yy) individual produces equal proportions of Y and y gametes because homologous chromosomes separate at anaphase I of meiosis.}
\end{figure}

\subsubsection{The Molecular Basis of Mendel's Laws: From Bivalents to
Binomials}\label{the-molecular-basis-of-mendels-laws-from-bivalents-to-binomials}

Mendel formulated his laws in 1866 without any knowledge of chromosomes,
meiosis, or DNA. He treated alleles as abstract ``factors'' --- discrete
particles that combined and segregated according to mathematical rules.
The chromosomal interpretation came nearly four decades later, when
Walter \citet{sutton1902} and Theodor Boveri independently noticed that
chromosomes in meiosis behaved exactly the way Mendel's factors had to
behave for his laws to hold:

{\def\LTcaptype{none} % do not increment counter
\begin{longtable}[]{@{}
  >{\raggedright\arraybackslash}p{(\linewidth - 2\tabcolsep) * \real{0.4894}}
  >{\raggedright\arraybackslash}p{(\linewidth - 2\tabcolsep) * \real{0.5106}}@{}}
\toprule\noalign{}
\begin{minipage}[b]{\linewidth}\raggedright
Mendel's abstract law
\end{minipage} & \begin{minipage}[b]{\linewidth}\raggedright
Chromosomal mechanism
\end{minipage} \\
\midrule\noalign{}
\endhead
\bottomrule\noalign{}
\endlastfoot
Each individual has two factors per trait & Diploid cells have two
homologs per chromosome \\
Factors separate equally into gametes & Homologs disjoin at anaphase I;
each gamete inherits one \\
Factors at different loci assort independently & Non-homologous
bivalents orient independently at metaphase I \\
Factors are particulate, not blended & Chromosomes maintain integrity
across generations (apart from crossing over) \\
\end{longtable}
}

The synthesis of Mendel's laws with the chromosome theory of inheritance
--- completed by Thomas Hunt Morgan's \emph{Drosophila} work between
1910 and 1915 --- is one of the central convergences in biology. Every
Punnett square is, at heart, a model of homologous-chromosome behavior
in meiosis (\cref{sec:unit_V_chromosomal_inheritance}). The 3:1 ratio is
not a numerical coincidence; it is the inevitable algebraic shadow of
two homologs segregating independently into haploid gametes that then
unite at fertilization. When meiosis fails --- non-disjunction,
translocation, or imprinting errors --- Mendelian ratios fail
predictably with it.

\subsubsection{Law of Independent Assortment (Mendel's Second
Law)}\label{law-of-independent-assortment-mendels-second-law}

\textbf{Statement}: Alleles at different gene loci are distributed to
gametes \textbf{independently} of one another -- the inheritance of one
gene does not affect the inheritance of another.

\textbf{Molecular basis}: Random orientation of non-homologous bivalents
at \textbf{metaphase I}. Each bivalent can orient with either homolog
facing either pole, and this orientation is independent of the other
bivalents.

\textbf{Important exception}: This law applies only to genes on
\textbf{different chromosomes} (or very far apart on the same
chromosome, \textgreater50 cM). Genes on the same chromosome that are
close together exhibit \textbf{\hyperref[gl:linkage]{linkage}} and
violate independent assortment (see \emph{Chromosomal Inheritance and
Linkage}).

\textbf{Concept Check 16.1}

\begin{quote}
\begin{enumerate}
\def\labelenumi{\arabic{enumi}.}
\tightlist
\item
  If the law of segregation is based on meiosis I, what specific
  cellular event during meiosis II is relevant?
\item
  A student claims that Mendel's laws prove genes are on chromosomes. Is
  this historically accurate? When was this connection established?
\item
  If two genes are 5 cM apart on the same chromosome, do they assort
  independently? Explain.
\end{enumerate}
\end{quote}

\begin{center}\rule{0.5\linewidth}{0.5pt}\end{center}

\subsection{Monohybrid Crosses}\label{monohybrid-crosses}

\subsubsection{The Standard Monohybrid
Cross}\label{the-standard-monohybrid-cross}

\textbf{Worked example -- seed color:}

\begin{equation}
P: \quad YY \text{ (yellow)} \times yy \text{ (green)}
\label{eq:unit_V_mendelian_p_cross}
\end{equation}

\begin{equation}
F_1: \quad Yy \text{ (most yellow -- Y is dominant)}
\label{eq:unit_V_mendelian_f1}
\end{equation}

\begin{equation}
F_1 \times F_1: \quad Yy \times Yy
\label{eq:unit_V_mendelian_f1_self}
\end{equation}

\textbf{Punnett square:}

{\def\LTcaptype{none} % do not increment counter
\begin{longtable}[]{@{}lll@{}}
\toprule\noalign{}
& \textbf{Y} & \textbf{y} \\
\midrule\noalign{}
\endhead
\bottomrule\noalign{}
\endlastfoot
\textbf{Y} & YY & Yy \\
\textbf{y} & Yy & yy \\
\end{longtable}
}

\textbf{F\(_2\) genotype ratio}: \(\frac{1}{4}\) YY : \(\frac{2}{4}\) Yy
: \(\frac{1}{4}\) yy = \textbf{1:2:1}

\textbf{F\(_2\) \hyperref[gl:phenotype]{\textbf{phenotype}} ratio}:
\(\frac{3}{4}\) yellow : \(\frac{1}{4}\) green = \textbf{3:1}, the
canonical heterozygote-cross outcome visualised in
\cref{fig:unit_V_punnett_square}, where the 3:1 phenotype ratio falls
out directly from the underlying 1:2:1 genotype ratio.

\subsubsection{The Testcross}\label{the-testcross}

To determine whether an individual showing the dominant phenotype is
homozygous (YY) or heterozygous (Yy), cross it with a \textbf{homozygous
recessive} (yy):

\begin{itemize}
\tightlist
\item
  If \textbf{YY x yy}: the offspring are Yy (yellow) -- no recessive
  offspring
\item
  If \textbf{Yy x yy}: \(\frac{1}{2}\) Yy (yellow) : \(\frac{1}{2}\) yy
  (green) -- 1:1 ratio reveals heterozygosity
\end{itemize}

\subsubsection{Probability Rules}\label{probability-rules}

\textbf{Product rule (AND)}: The probability of two independent events
both occurring = product of individual probabilities.

\begin{equation}
P(A \text{ and } B) = P(A) \times P(B)
\label{eq:unit_V_mendelian_product_rule}
\end{equation}

\textbf{Sum rule (OR)}: The probability of either of two mutually
exclusive events occurring = sum of individual probabilities.

\begin{equation}
P(A \text{ or } B) = P(A) + P(B)
\label{eq:unit_V_mendelian_sum_rule}
\end{equation}

\subsection{Worked Example: 16.1: Two heterozygous parents (Yy x Yy).
What is the probability
that:}\label{worked-example-16.1-two-heterozygous-parents-yy-x-yy.-what-is-the-probability-that}

\begin{enumerate}
\def\labelenumi{(\alph{enumi})}
\tightlist
\item
  The first child is yy?
\end{enumerate}

\[P(yy) = \frac{1}{4} \tag{16.1} \label{eq:unit_V_mendelian_genetics_item_1}\]

\begin{enumerate}
\def\labelenumi{(\alph{enumi})}
\setcounter{enumi}{1}
\tightlist
\item
  The first two children are both yy?
\end{enumerate}

\[P = \frac{1}{4} \times \frac{1}{4} = \frac{1}{16} \quad \text{(product rule; independent events)} \tag{16.2} \label{eq:unit_V_mendelian_genetics_item_2}\]

\begin{enumerate}
\def\labelenumi{(\alph{enumi})}
\setcounter{enumi}{2}
\tightlist
\item
  Among three children, exactly one is yy?
\end{enumerate}

\[P = \binom{3}{1}\left(\frac{1}{4}\right)^1\left(\frac{3}{4}\right)^2 = 3 \times \frac{1}{4} \times \frac{9}{16} = \frac{27}{64} \approx 0.422 \tag{16.3} \label{eq:unit_V_mendelian_genetics_item_3}\]

\begin{center}\rule{0.5\linewidth}{0.5pt}\end{center}

\subsection{Dihybrid and Trihybrid
Crosses}\label{dihybrid-and-trihybrid-crosses}

\subsubsection{The Dihybrid Cross}\label{the-dihybrid-cross}

Two independently assorting genes: seed color (Y/y) and seed shape (R =
round, r = wrinkled):

\[P: \quad YYRR \times yyrr \tag{16.4} \label{eq:unit_V_mendelian_genetics_item_4}\]

\[F_1: \quad YyRr \text{ (yellow, round)} \tag{16.5} \label{eq:unit_V_mendelian_genetics_item_5}\]

\[F_1 \times F_1: \quad YyRr \times YyRr \tag{16.6} \label{eq:unit_V_mendelian_genetics_item_6}\]

Each parent produces 4 gamete types in equal frequency: YR, Yr, yR, yr.

\textbf{F\(_2\) phenotype ratio}: 9 Y\_R\_ : 3 Y\_rr : 3 yyR\_ : 1 yyrr
= \textbf{9:3:3:1}

This ratio is the product of two independent 3:1 ratios:

\begin{equation}
\left(\frac{3}{4}Y\_ + \frac{1}{4}yy\right) \times \left(\frac{3}{4}R\_ + \frac{1}{4}rr\right) = \frac{9}{16} + \frac{3}{16} + \frac{3}{16} + \frac{1}{16}
\label{eq:unit_V_mendelian_9331}
\end{equation}

\subsubsection{The Forked-Line (Branch Diagram)
Method}\label{the-forked-line-branch-diagram-method}

For a trihybrid cross (AaBbCc x AaBbCc), the Punnett square would
require 64 cells. Instead, use the forked-line method:

\[\text{A locus:} \quad \tfrac{3}{4}A\_ \text{ and } \tfrac{1}{4}aa \tag{16.7} \label{eq:unit_V_mendelian_genetics_item_7}\]
\[\text{B locus:} \quad \tfrac{3}{4}B\_ \text{ and } \tfrac{1}{4}bb \tag{16.8} \label{eq:unit_V_mendelian_genetics_item_8}\]
\[\text{C locus:} \quad \tfrac{3}{4}C\_ \text{ and } \tfrac{1}{4}cc \tag{16.9} \label{eq:unit_V_mendelian_genetics_item_9}\]

Probability of A\_B\_C\_ =
\(\frac{3}{4} \times \frac{3}{4} \times \frac{3}{4} = \frac{27}{64}\)

Total number of phenotypic classes = \(2^3 = 8\)

Number of genotypic classes = \(3^3 = 27\)

\subsection{Worked Example: Trihybrid Cross
Probabilities}\label{worked-example-trihybrid-cross-probabilities}

In a trihybrid cross AaBbCc x AaBbCc, what fraction of offspring will
be:

\begin{enumerate}
\def\labelenumi{(\alph{enumi})}
\tightlist
\item
  AaBbCc (triple heterozygous)?
\end{enumerate}

\[P = \frac{2}{4} \times \frac{2}{4} \times \frac{2}{4} = \frac{8}{64} = \frac{1}{8} \tag{16.10} \label{eq:unit_V_mendelian_genetics_item_10}\]

\begin{enumerate}
\def\labelenumi{(\alph{enumi})}
\setcounter{enumi}{1}
\tightlist
\item
  Homozygous for most three loci (AA BB CC or aa bb cc etc.)?
\end{enumerate}

Each locus has probability \(\frac{1}{4}\) + \(\frac{1}{4}\) =
\(\frac{1}{2}\) of being homozygous (either AA or aa).

\[P(\text{homozygous at most 3}) = \frac{1}{2} \times \frac{1}{2} \times \frac{1}{2} = \frac{1}{8} \tag{16.11} \label{eq:unit_V_mendelian_genetics_item_11}\]

\begin{center}\rule{0.5\linewidth}{0.5pt}\end{center}

\subsection{Extensions to Mendelian
Genetics}\label{extensions-to-mendelian-genetics}

\subsubsection{Incomplete Dominance}\label{incomplete-dominance}

Neither allele is completely dominant; the heterozygote displays an
\textbf{intermediate} phenotype:

\begin{itemize}
\tightlist
\item
  \textbf{Snapdragon flower color}: \(C^R C^R\) (red) x \(C^W C^W\)
  (white) \texorpdfstring{\ensuremath{\to}}{->} \(C^R C^W\) (pink)
\item
  F\(_2\): 1 red : 2 pink : 1 white (\textbf{1:2:1 phenotype ratio =
  genotype ratio})
\end{itemize}

Molecular basis: In many cases, gene dosage explains the intermediate
phenotype. One copy of \(C^R\) produces half the amount of red pigment
compared to two copies.

\subsubsection{Codominance}\label{codominance}

Both alleles are \textbf{fully and simultaneously expressed} in the
heterozygote:

\begin{itemize}
\tightlist
\item
  \textbf{ABO blood groups}: \(I^A I^A\) or \(I^A i\) (type A)
  vs.~\(I^B I^B\) or \(I^B i\) (type B) vs.~\(I^A I^B\) (type AB -- both
  A and B antigens present on red blood cell surface)
\item
  Molecular basis: \(I^A\) encodes N-acetylgalactosaminyltransferase
  (adds GalNAc to H antigen); \(I^B\) encodes galactosyltransferase
  (adds Gal); both \hyperref[gl:enzyme]{\textbf{enzyme}}s are functional
  in heterozygotes; \(i\) encodes a nonfunctional enzyme
\end{itemize}

\subsubsection{Multiple Alleles}\label{multiple-alleles}

A gene can have more than two alleles in a population (though any
individual still carries at most two):

\textbf{ABO blood group system} (three alleles: \(I^A\), \(I^B\),
\(i\)):

\begin{table}[htbp]
\label{tbl:unit_V_mendelian_genetics_1}
\centering
\small
\begin{tabular}{@{}p{2.8cm}p{0.8cm}p{2.0cm}p{2.6cm}p{2.6cm}p{3.0cm}@{}}
\toprule
\textbf{Genotype} & \textbf{Phen.} & \textbf{Antigens} & \textbf{Antibodies} & \textbf{Can Donate To} & \textbf{Can Receive From} \\
\midrule
$I^A I^A$ or $I^A i$ & A & A & Anti-B & A, AB & A, O \\
$I^B I^B$ or $I^B i$ & B & B & Anti-A & B, AB & B, O \\
$I^A I^B$ & AB & A and B & Neither & AB only & Universal recipient \\
$ii$ & O & Neither (H antigen only) & Anti-A and Anti-B & Universal donor & O only \\
\bottomrule
\end{tabular}
\end{table}

\textbf{Rabbit coat color} (4 alleles at the \emph{C} locus): \(c^+\)
(full color) \textgreater{} \(c^{ch}\) (chinchilla) \textgreater{}
\(c^h\) (Himalayan) \textgreater{} \(c\) (albino). Dominance hierarchy
gives 10 possible genotypes and 5 phenotypic classes.

\begin{quote}
\textbf{Clinical Connection: Rh Disease of the Newborn} The Rh blood
group system involves the RhD \hyperref[gl:protein]{\textbf{protein}}.
Rh-negative mothers (RhD-/RhD-) carrying an Rh-positive fetus can become
sensitized during delivery when fetal RhD+ red blood cells enter
maternal circulation. Maternal IgG anti-RhD antibodies cross the
placenta in subsequent pregnancies, destroying fetal RhD+ erythrocytes
(hemolytic disease of the fetus and newborn, HDFN). Prevention:
\textbf{anti-RhD immunoglobulin (RhoGAM)} administered at 28 weeks and
within 72 hours of delivery clears fetal RhD+ cells before maternal
sensitization occurs.
\end{quote}

\subsubsection{Pleiotropy}\label{pleiotropy}

A single gene affects \textbf{multiple phenotypic traits}:

\textbf{Sickle cell disease (HbS)}: A single
\hyperref[gl:nucleotide]{\textbf{nucleotide}} change (GAG to GTG,
Glu6Val) in the beta-globin gene (\emph{HBB}) produces sickle
hemoglobin. Under low oxygen tension, HbS polymerizes, deforming red
blood cells into rigid sickle shapes. This single
\hyperref[gl:mutation]{\textbf{mutation}} causes:

\begin{itemize}
\tightlist
\item
  Vaso-occlusive crises (pain)
\item
  Hemolytic anemia (shortened RBC lifespan: about 17 days
  vs.~about 120 days)
\item
  Splenic sequestration and autosplenectomy
\item
  Stroke (large vessel occlusion)
\item
  Acute chest syndrome
\item
  Avascular necrosis of bones
\item
  Renal papillary necrosis
\item
  Retinopathy
\end{itemize}

\textbf{Phenylketonuria (PKU)}: Phenylalanine hydroxylase (PAH)
deficiency leads to phenylalanine accumulation, causing intellectual
disability, fair skin/hair (reduced melanin from phenylalanine
competition with tyrosine), musty body odor, and eczema. Newborn
screening and phenylalanine-restricted diet prevent neurological damage.

\textbf{Marfan syndrome}: Mutations in \emph{FBN1} (fibrillin-1) affect
connective tissue throughout the body: tall stature, arachnodactyly,
lens subluxation, mitral valve prolapse, and aortic root dilation (risk
of aortic dissection).

\subsubsection{Epistasis: When One Gene Masks
Another}\label{epistasis-when-one-gene-masks-another}

\textbf{Epistasis} describes the interaction between two or more genes
in which the phenotype produced by one gene depends on the genotype at a
second locus. The 9:3:3:1 dihybrid ratio applies primarily when two
independently assorting genes act independently on the phenotype. When
the genes interact --- when one gene's protein product feeds the
substrate that another gene's product modifies, or when one gene
controls whether another gene's product is even visible --- the dihybrid
F\(_2\) ratio is \textbf{modified} in characteristic ways. Each modified
ratio is a genetic signature of a specific kind of biochemical
interaction.

\begin{figure}[htbp]
\centering
\includegraphics[width=0.82\linewidth,height=4.2in,keepaspectratio]{../figures/mermaid_inline/inline_mermaid_0089_ed7f69a05270.png}
\caption{Modified dihybrid ratios produced by different types of epistasis. Each ratio is a modification of the standard 9:3:3:1 ratio.}
\end{figure}

\paragraph{Recessive Epistasis (9:3:4) --- Labrador coat
color}\label{recessive-epistasis-934-labrador-coat-color}

Recessive epistasis occurs when \textbf{homozygosity at one locus masks
the phenotype produced by a second locus}. The Labrador retriever coat
color illustrates the principle through a two-step pigment pathway:

\begin{itemize}
\tightlist
\item
  \textbf{B locus} (Brown): determines whether eumelanin is black
  (\(B\)) or brown/chocolate (\(b\)).
\item
  \textbf{E locus} (Extension): determines whether eumelanin is
  \textbf{deposited at most} in the hair shaft. \(E\) alleles deposit
  pigment; \(ee\) homozygotes block deposition entirely, producing
  yellow coats regardless of B genotype.
\end{itemize}

{\def\LTcaptype{none} % do not increment counter
\begin{longtable}[]{@{}lll@{}}
\toprule\noalign{}
Genotype & Pigment chemistry & Coat \\
\midrule\noalign{}
\endhead
\bottomrule\noalign{}
\endlastfoot
B\_E\_ & Black eumelanin deposited & Black \\
bbE\_ & Brown eumelanin deposited & Chocolate \\
B\_ee & Black eumelanin made but \textbf{not deposited} & Yellow \\
bbee & Brown eumelanin made but \textbf{not deposited} & Yellow \\
\end{longtable}
}

Because \(ee\) is \textbf{epistatic} to the B locus, the standard
9:3:3:1 collapses into \textbf{9 black : 3 chocolate : 4 yellow} in an
F\(_2\) from BbEe \texorpdfstring{\ensuremath{\times}}{x} BbEe.

\paragraph{Worked Example 16.3a: Recessive
epistasis}\label{worked-example-16.3a-recessive-epistasis}

In a cross BbEe \texorpdfstring{\ensuremath{\times}}{x} BbEe, what is the probability of a chocolate Labrador
puppy?

\[P(\text{chocolate}) = P(bb) \times P(E\_) = \tfrac{1}{4} \times \tfrac{3}{4} = \tfrac{3}{16} \label{eq:unit_V_mendelian_genetics_item_12}\]

\paragraph{Dominant Epistasis (12:3:1) --- Squash fruit
color}\label{dominant-epistasis-1231-squash-fruit-color}

Dominant epistasis occurs when \textbf{the dominant allele at one locus
masks expression at a second locus}. White-fruited summer squash
(\emph{Cucurbita pepo}) is a classic case. Allele \(W\) produces white
fruit and is \textbf{dominant} to \(w\) (colored fruit). At a second
locus, \(Y\) produces yellow and \(y\) produces green pigment --- but
primarily when \(ww\) allows pigment to develop.

{\def\LTcaptype{none} % do not increment counter
\begin{longtable}[]{@{}ll@{}}
\toprule\noalign{}
Genotype & Phenotype \\
\midrule\noalign{}
\endhead
\bottomrule\noalign{}
\endlastfoot
W\_Y\_ & White \\
W\_yy & White \\
wwY\_ & Yellow \\
wwyy & Green \\
\end{longtable}
}

The F\(_2\) from WwYy \texorpdfstring{\ensuremath{\times}}{x} WwYy yields \textbf{12 white : 3 yellow : 1
green}, written 12:3:1. Dominant \(W\) shuts down the entire pathway
downstream of the Y locus.

\paragraph{Duplicate Recessive Epistasis (9:7) --- Sweet pea flower
color}\label{duplicate-recessive-epistasis-97-sweet-pea-flower-color}

Sweet pea (\emph{Lathyrus odoratus}) flower color requires the
\textbf{complementary action} of two unlinked genes, both encoding
biosynthetic enzymes in the anthocyanin pathway. Bateson and Punnett
(1905) crossed two pure-breeding white strains and obtained purple
F\(_1\) offspring --- an unexpected result that reflected
complementation between the recessive defects in the two parental lines.
The F\(_2\) ratio is \textbf{9 purple : 7 white}.

{\def\LTcaptype{none} % do not increment counter
\begin{longtable}[]{@{}lll@{}}
\toprule\noalign{}
Genotype & Both enzymes functional? & Color \\
\midrule\noalign{}
\endhead
\bottomrule\noalign{}
\endlastfoot
C\_P\_ & Yes --- full pathway active & Purple \\
C\_pp & No --- pathway blocked at second step & White \\
ccP\_ & No --- pathway blocked at first step & White \\
ccpp & No --- pathway blocked at both steps & White \\
\end{longtable}
}

Either single-locus block (cc or pp) produces white; primarily
doubly-dominant individuals build pigment. This \textbf{9:7 ratio} is a
hallmark of \textbf{complementary gene action} in linear biosynthetic
pathways.

\paragraph{Duplicate Dominant Epistasis (15:1) --- Shepherd's purse seed
shape}\label{duplicate-dominant-epistasis-151-shepherds-purse-seed-shape}

When \textbf{either dominant allele alone is sufficient} to produce the
phenotype, primarily the double-recessive shows the alternative trait.
F\(_2\) ratio: \textbf{15 phenotype A : 1 phenotype B}. Shepherd's purse
seed-capsule shape is the classic example: triangular (\(A\_\) or
\(B\_\)) versus narrow (\(aabb\)).

\paragraph{Combinatorics summary}\label{combinatorics-summary}

{\def\LTcaptype{none} % do not increment counter
\begin{longtable}[]{@{}
  >{\raggedright\arraybackslash}p{(\linewidth - 4\tabcolsep) * \real{0.5122}}
  >{\raggedright\arraybackslash}p{(\linewidth - 4\tabcolsep) * \real{0.2683}}
  >{\raggedright\arraybackslash}p{(\linewidth - 4\tabcolsep) * \real{0.2195}}@{}}
\toprule\noalign{}
\begin{minipage}[b]{\linewidth}\raggedright
Modified F\(_2\) ratio
\end{minipage} & \begin{minipage}[b]{\linewidth}\raggedright
Mechanism
\end{minipage} & \begin{minipage}[b]{\linewidth}\raggedright
Example
\end{minipage} \\
\midrule\noalign{}
\endhead
\bottomrule\noalign{}
\endlastfoot
9:3:4 & Recessive epistasis & Labrador coat color (B/E) \\
12:3:1 & Dominant epistasis & White squash fruit (W/Y) \\
9:7 & Duplicate recessive (complementation) & Sweet pea flower color \\
15:1 & Duplicate dominant & Shepherd's purse capsules \\
13:3 & Dominant + recessive interaction & White Leghorn \texorpdfstring{\ensuremath{\times}}{x} white
Wyandotte chickens \\
\end{longtable}
}

In every case the ratio sums to 16, betraying its origin in 9:3:3:1 ---
epistasis is \textbf{typically a regrouping of the same dihybrid
classes}, rarely the introduction of new ones.

\paragraph{Epistasis in mouse coat color: agouti \texorpdfstring{\ensuremath{\times}}{x}
albino}\label{epistasis-in-mouse-coat-color-agouti-albino}

A second worked example consolidates the recessive-epistasis logic. In
the laboratory mouse, the \textbf{A locus} (agouti) determines whether
hairs have a yellow band (A\_) or are uniformly black (aa). At the
\textbf{C locus}, the C allele is required to deposit any pigment --- cc
homozygotes are albino regardless of the A genotype. A cross of AaCc \texorpdfstring{\ensuremath{\times}}{x}
AaCc produces F\(_2\) offspring in a \textbf{9 agouti : 3 black : 4
albino} ratio, with the albino class combining both A\_cc and aacc
genotypes (because pigment deposition fails in either case).

The \textbf{9:3:4 signature ratio} therefore appears whenever:

\begin{enumerate}
\def\labelenumi{\arabic{enumi}.}
\tightlist
\item
  A gene at one locus produces a substrate or precursor.
\item
  A gene at a second locus deposits, modifies, or expresses the product
  of the first.
\item
  Homozygosity for the loss-of-function allele at the second locus
  completely blocks phenotypic expression.
\end{enumerate}

This logic underlies the genetic basis of many recessive epistatic
systems: \textbf{Bombay phenotype} at the H locus (cc-equivalent for ABO
blood groups), \textbf{albinism in mammals} masking coat-pattern loci,
and \textbf{white-flower mutants} in many ornamental plants.

\paragraph{Worked Example 16.3b: Duplicate dominant ratio in shepherd's
purse}\label{worked-example-16.3b-duplicate-dominant-ratio-in-shepherds-purse}

In a cross of two heterozygous shepherd's purse plants (AaBb \texorpdfstring{\ensuremath{\times}}{x} AaBb),
what fraction of F\(_2\) offspring are predicted to have \textbf{narrow}
seed capsules (the alternative phenotype)?

Narrow capsules require the double-recessive genotype aabb. Each locus
contributes \(\tfrac{1}{4}\) probability of double-recessive at that
locus:

\[P(\text{narrow}) = P(aa) \times P(bb) = \tfrac{1}{4} \times \tfrac{1}{4} = \tfrac{1}{16} \label{eq:unit_V_mendelian_genetics_item_13}\]

The remaining \(\tfrac{15}{16}\) of offspring have triangular capsules.
The 15:1 ratio is the genetic signature of \textbf{functional
redundancy} --- when either gene alone is sufficient, primarily the
simultaneous loss of both produces the alternative phenotype. The same
logic explains why \textbf{paralogous gene families} often show
15:1-like behavior in knockout studies: knocking out either paralog
alone produces no phenotype, but the double knockout reveals the shared
function.

\subsubsection{Polygenic Inheritance}\label{polygenic-inheritance}

\textbf{Multiple genes} contribute additively to a single continuous
(quantitative) trait:

\begin{itemize}
\tightlist
\item
  \textbf{Human skin color}: At least 7 genes contribute (SLC24A5,
  SLC45A2, TYR, TYRP1, OCA2, KITLG, MC1R and others). Each gene has
  alleles that contribute incrementally to melanin production. The
  result is a \textbf{continuous distribution} (approximately
  normal/Gaussian) rather than discrete classes.
\item
  \textbf{Human height}: about 700 GWAS loci identified; each
  has a small effect (1-2 mm per allele). Environmental factors
  (nutrition, health) also contribute significantly. Heritability
  about 0.8 in well-nourished populations.
\end{itemize}

The number of phenotypic classes for n contributing loci (each with two
alleles, additive effects) = \(2n + 1\).

\subsubsection{Penetrance and
Expressivity}\label{penetrance-and-expressivity}

Mendelian ratios assume that every individual carrying a given genotype
displays the expected phenotype, and displays it identically. In
reality, the genotype--phenotype map is leaky on \textbf{two distinct
axes} --- \emph{whether} the phenotype appears at most (penetrance) and
\emph{how strongly} it appears when it does (expressivity).

\begin{itemize}
\tightlist
\item
  \textbf{Penetrance} is the \textbf{proportion of individuals with a
  given genotype who actually display the expected phenotype}. A trait
  with 100\% penetrance is \textbf{completely penetrant}; a trait with
  less than 100\% is \textbf{incompletely penetrant}.
\item
  \textbf{Expressivity} is the \textbf{degree or severity of phenotype
  expression} among individuals who \emph{are} penetrant. Variable
  expressivity means the same genotype produces a range of phenotypic
  intensities.
\item
  \textbf{Phenocopy} is an \textbf{environmentally induced phenotype
  that mimics a genetic condition} without the underlying genotype.
  Phenocopies can confound pedigree analysis if not recognized.
\end{itemize}

\paragraph{Disease examples}\label{disease-examples}

{\def\LTcaptype{none} % do not increment counter
\begin{longtable}[]{@{}
  >{\raggedright\arraybackslash}p{(\linewidth - 4\tabcolsep) * \real{0.3793}}
  >{\raggedright\arraybackslash}p{(\linewidth - 4\tabcolsep) * \real{0.3448}}
  >{\raggedright\arraybackslash}p{(\linewidth - 4\tabcolsep) * \real{0.2759}}@{}}
\toprule\noalign{}
\begin{minipage}[b]{\linewidth}\raggedright
Phenomenon
\end{minipage} & \begin{minipage}[b]{\linewidth}\raggedright
Disorder
\end{minipage} & \begin{minipage}[b]{\linewidth}\raggedright
Detail
\end{minipage} \\
\midrule\noalign{}
\endhead
\bottomrule\noalign{}
\endlastfoot
Reduced penetrance & \textbf{BRCA1/BRCA2} breast/ovarian cancer &
about 70--80\% lifetime breast cancer risk in BRCA1 carriers;
some carriers rarely develop disease \\
Reduced penetrance & \textbf{Familial hypercholesterolemia} (LDLR
mutations) & Variable age of onset for cardiovascular disease; modified
by diet and other genes \\
Reduced penetrance & \textbf{Retinoblastoma (RB1)} & about 90\%
penetrance; some heterozygous carriers escape because the second hit
(loss of the wild-type allele) does not occur in retinal precursor
cells \\
Variable expressivity & \textbf{Neurofibromatosis type 1 (NF1)} &
\emph{NF1} tumor suppressor (encodes neurofibromin, a Ras-GAP); same
mutation produces café-au-lait macules in some individuals, severe
plexiform neurofibromas, optic gliomas, and skeletal dysplasia in
others. Loss-of-heterozygosity (second hit) at \emph{NF1} in Schwann
cell lineages produces the variable tumor burden \\
Variable expressivity & \textbf{Polydactyly} (e.g., \emph{GLI3}
mutations, Greig cephalopolysyndactyly) & Affected family members
carrying the identical mutation may show one extra digit on one hand,
two extra digits on both hands and feet, or post-axial duplication
varying side-to-side. Limb-bud stochastic patterning during weeks 4--7
produces this variability \\
Variable expressivity & \textbf{Marfan syndrome (FBN1)} & Aortic root
size varies dramatically among affected family members carrying the same
mutation \\
Variable expressivity & \textbf{Huntington's disease} & Age of onset
varies inversely with CAG repeat length but also among individuals with
identical repeat counts (modifier loci) \\
Phenocopy & \textbf{Thalidomide-induced phocomelia} & Limb shortening
from drug exposure (1958--1962) mimicked the rare genetic condition \\
Phenocopy & \textbf{Congenital rubella syndrome} & Maternal rubella
infection causes deafness and cataracts that resemble genetic
syndromes \\
\end{longtable}
}

The biological causes of incomplete penetrance and variable expressivity
include \textbf{modifier genes} (other loci that influence how the
primary mutation manifests), \textbf{stochastic developmental events}
(random fluctuations during embryogenesis), \textbf{environmental
triggers} (diet, infection, smoking), \textbf{epigenetic state}
(methylation patterns inherited or established during development), and
--- for X-linked traits in females --- \textbf{random X-inactivation
patterns} (\cref{sec:unit_V_chromosomal_inheritance}).

For genetic counseling, incomplete penetrance complicates risk
prediction: a ``negative'' family member is not certainly mutation-free,
and a ``positive'' family history is not a assurance of disease.
Penetrance estimates are typically reported as age-specific cumulative
risks (e.g., ``80\% by age 70 for BRCA1'') rather than single numbers.

\subsubsection{Genomic Imprinting}\label{genomic-imprinting}

Some genes are expressed in a \textbf{parent-of-origin-specific} manner
due to differential DNA methylation established during gametogenesis:

\textbf{Prader-Willi syndrome vs.~Angelman syndrome} (chromosome
15q11-13): - \textbf{Prader-Willi syndrome}: Deletion/loss of the
\textbf{paternal} 15q11-13 region (or maternal uniparental disomy).
Genes in this region are normally expressed primarily from the paternal
chromosome. Features: hypotonia in infancy, hyperphagia and obesity,
intellectual disability, short stature. - \textbf{Angelman syndrome}:
Deletion/loss of the \textbf{maternal} 15q11-13 region (specifically the
\emph{UBE3A} gene, which is maternally expressed in
\hyperref[gl:neuron]{\textbf{neuron}}s). Features: severe intellectual
disability, seizures, ataxic gait, happy demeanor with frequent
laughter.

Same chromosomal region, but different disease depending on which
parent's copy is affected -- a clear demonstration that the
\hyperref[gl:genome]{\textbf{genome}} ``remembers'' parental origin.

\subsubsection{Maternal Effect Genes}\label{maternal-effect-genes}

The \textbf{mother's genotype} (not the offspring's) determines the
offspring's phenotype:

\begin{itemize}
\tightlist
\item
  \textbf{Bicoid in \emph{Drosophila}}: The bicoid mRNA is deposited
  into the oocyte by the mother. It localizes to the anterior end and is
  translated after fertilization, producing a
  \hyperref[gl:transcription]{\textbf{transcription}} factor gradient
  that specifies anterior structures (head, thorax). A mother homozygous
  for a \emph{bicoid} mutation produces embryos lacking anterior
  structures -- regardless of the embryo's own genotype.
\end{itemize}

\textbf{Concept Check 16.2}

\begin{quote}
\begin{enumerate}
\def\labelenumi{\arabic{enumi}.}
\tightlist
\item
  Distinguish between codominance and incomplete dominance at the
  molecular level.
\item
  Two yellow Labradors are mated. Can they produce chocolate puppies?
  Explain.
\item
  If a trait is 70\% penetrant, what fraction of heterozygous carriers
  will NOT show the phenotype?
\item
  A mother with Prader-Willi syndrome (paternal deletion) has a child.
  Can the child inherit Prader-Willi? Angelman? Explain.
\item
  Distinguish reduced penetrance from variable expressivity. Give one
  disease example of each.
\end{enumerate}
\end{quote}

\begin{center}\rule{0.5\linewidth}{0.5pt}\end{center}

\subsection{Human Genetic Disorders: Mechanism Meets
Mendel}\label{human-genetic-disorders-mechanism-meets-mendel}

Human genetics has historically been the inverse of Mendel's pea garden
--- instead of designing crosses, clinicians infer genotypes from
pedigrees and clinical features. The core inheritance patterns
(autosomal dominant, autosomal recessive, X-linked, mitochondrial)
follow directly from Mendelian principles applied to specific
chromosomal architecture. The molecular mechanisms behind individual
disorders, however, illustrate the rich biology that makes Mendel's
abstract ``factors'' into proteins, enzymes, and regulatory RNAs.

\subsubsection{Autosomal Dominant Disorders: Huntington's Disease and
CAG Repeat
Expansion}\label{autosomal-dominant-disorders-huntingtons-disease-and-cag-repeat-expansion}

\textbf{Huntington's disease (HD)} is an adult-onset neurodegenerative
disorder caused by mutations in the \emph{HTT} gene on chromosome
4p16.3. It is \textbf{autosomal dominant} with nearly complete
penetrance by age 65 --- a rare example of a fully penetrant late-onset
disorder. The clinical features include progressive chorea (involuntary
writhing movements), psychiatric symptoms (depression, irritability,
psychosis), and cognitive decline (executive dysfunction, dementia).
Median survival from symptom onset is approximately 15--20 years.

\textbf{Molecular mechanism}: The \emph{HTT} gene contains a polymorphic
CAG trinucleotide repeat in its first exon, encoding a polyglutamine
(polyQ) tract in the huntingtin protein. The number of CAG repeats
determines disease status:

{\def\LTcaptype{none} % do not increment counter
\begin{longtable}[]{@{}
  >{\raggedright\arraybackslash}p{(\linewidth - 2\tabcolsep) * \real{0.4483}}
  >{\raggedright\arraybackslash}p{(\linewidth - 2\tabcolsep) * \real{0.5517}}@{}}
\toprule\noalign{}
\begin{minipage}[b]{\linewidth}\raggedright
CAG repeats
\end{minipage} & \begin{minipage}[b]{\linewidth}\raggedright
Disease status
\end{minipage} \\
\midrule\noalign{}
\endhead
\bottomrule\noalign{}
\endlastfoot
\texorpdfstring{\ensuremath{\leq}}{<=} 26 & Normal --- no risk; stable inheritance \\
27--35 & Intermediate --- usually unaffected, but unstable in
transmission and may expand into the affected range in offspring
(especially through paternal transmission) \\
36--39 & Reduced penetrance --- disease may or may not develop in a
normal lifespan \\
\texorpdfstring{\textgreater{}=}{>=} 40 & Fully penetrant --- disease will develop if individual lives long
enough \\
\end{longtable}
}

The expanded polyQ tract causes the huntingtin protein to misfold and
aggregate, forming intranuclear inclusions in neurons (especially in the
striatum). The precise pathogenic mechanism remains debated --- gain of
toxic function, loss of normal function, and disruption of axonal
transport most contribute.

\textbf{Anticipation}: Disease severity worsens and age of onset
decreases across generations because the CAG repeat is \textbf{unstable
during meiosis}, especially during spermatogenesis. A father with 42
repeats may transmit 50 repeats to his son; that son's children may
inherit 60 or more. \textbf{Juvenile-onset HD} (onset before age 20)
almost typically reflects paternal transmission of severely expanded
alleles (\textgreater60 repeats). This phenomenon --- anticipation ---
is also observed in fragile X syndrome, myotonic dystrophy, and several
spinocerebellar ataxias, most of which are \textbf{trinucleotide repeat
expansion disorders}.

\textbf{Inverse correlation}: Within the affected range, age of onset
correlates inversely with repeat number, but with substantial variance
attributable to modifier genes (e.g., \emph{FAN1}, \emph{MSH3}) that
influence somatic repeat instability. This explains why two siblings
with the same CAG count can have onset ages differing by a decade or
more.

\textbf{Genetic counseling}: Because HD is autosomal dominant, each
child of an affected parent has a 50\% risk. Predictive testing in
unaffected adults raises difficult ethical questions --- there is no
preventive therapy, and testing carries psychological consequences
(depression, suicidality) that mandate pre-test and post-test genetic
counseling. Many at-risk individuals choose not to be tested.

\subsubsection{Autosomal Recessive Disorders: PKU and the Logic of
Newborn
Screening}\label{autosomal-recessive-disorders-pku-and-the-logic-of-newborn-screening}

\textbf{Phenylketonuria (PKU)} is an autosomal recessive disorder caused
by mutations in the \emph{PAH} gene on chromosome 12q23.2 encoding
\textbf{phenylalanine hydroxylase}, the liver enzyme that converts
phenylalanine to tyrosine. PKU illustrates how a single biochemical
defect cascades into multiple phenotypes (pleiotropy) and how early
detection plus dietary intervention can prevent the most devastating
manifestations.

\textbf{Biochemistry}: PAH is a tetrameric enzyme requiring
tetrahydrobiopterin (BH4) as cofactor. In its absence, phenylalanine
accumulates to toxic concentrations (\textgreater20 mg/dL versus normal
\textless2 mg/dL) while tyrosine --- the precursor of dopamine, melanin,
and thyroid hormone --- becomes deficient. The metabolic consequences
include:

\begin{itemize}
\tightlist
\item
  \textbf{Phenylalanine accumulation} in blood and brain, where excess
  phenylalanine competes with other large neutral amino acids (tyrosine,
  tryptophan, leucine) for the L-type amino acid transporter (LAT1) at
  the blood-brain barrier. Brain tyrosine and tryptophan deficiency
  disrupts catecholamine and serotonin synthesis.
\item
  \textbf{Phenylketones}: when phenylalanine is in excess, alternative
  metabolic pathways produce phenylpyruvate, phenyllactate, and
  phenylacetate --- the ``ketones'' detectable in urine that gave the
  disease its name.
\item
  \textbf{Tyrosine deficiency} reduces melanin synthesis (causing fair
  skin, blond hair, and blue eyes --- distinctive features in untreated
  patients) and may contribute to neurological symptoms.
\end{itemize}

\textbf{Clinical features (untreated)}: severe intellectual disability
(IQ typically \textless50 by age 1), seizures (in about 25\% of
untreated patients), microcephaly, motor abnormalities, eczema,
characteristic musty body odor (from phenylacetate excretion), and
reduced melanin pigmentation. Untreated PKU was a leading cause of
intellectual disability before comprehensive newborn screening.

\textbf{Newborn screening}: Robert Guthrie developed the \textbf{Guthrie
test} in 1962 --- a bacterial inhibition assay using \emph{Bacillus
subtilis} requiring phenylalanine. A heel-prick blood spot from a
24--72-hour-old newborn is dried on filter paper and screened en masse.
Modern programs use tandem mass spectrometry to measure dozens of
metabolites simultaneously. Newborn screening for PKU was the first
widespread genetic screening program and remains the prototype for
population screening programs worldwide.

\textbf{Treatment}: A \textbf{phenylalanine-restricted diet} initiated
within the first weeks of life prevents intellectual disability almost
entirely. Patients consume specially formulated medical foods low in
phenylalanine but supplemented with tyrosine. Treatment is lifelong ---
relaxation of dietary restrictions in adolescence or adulthood causes
neuropsychiatric symptoms (depression, anxiety, executive dysfunction)
and, in pregnant women with PKU, \textbf{maternal PKU syndrome}
(intellectual disability, microcephaly, congenital heart disease in the
fetus, even when the fetus is heterozygous). New therapies include
sapropterin (synthetic BH4, useful for about 30\% of patients
with residual PAH activity) and pegvaliase (PEGylated phenylalanine
ammonia lyase enzyme replacement, approved 2018).

\textbf{Population genetics}: PKU disease frequency is
about 1/10,000--15,000 in northern European populations, giving
a carrier frequency of about 1/55--60. PKU is more common in
Ireland (about 1/4,500), Turkey
(about 1/3,000--6,500), and parts of China
(about 1/15,000--17,000), with population-specific allele
spectra that allow targeted carrier screening.

\subsubsection{Other key autosomal
disorders}\label{other-key-autosomal-disorders}

{\def\LTcaptype{none} % do not increment counter
\begin{longtable}[]{@{}
  >{\raggedright\arraybackslash}p{(\linewidth - 6\tabcolsep) * \real{0.2500}}
  >{\raggedright\arraybackslash}p{(\linewidth - 6\tabcolsep) * \real{0.3250}}
  >{\raggedright\arraybackslash}p{(\linewidth - 6\tabcolsep) * \real{0.1500}}
  >{\raggedright\arraybackslash}p{(\linewidth - 6\tabcolsep) * \real{0.2750}}@{}}
\toprule\noalign{}
\begin{minipage}[b]{\linewidth}\raggedright
Disorder
\end{minipage} & \begin{minipage}[b]{\linewidth}\raggedright
Inheritance
\end{minipage} & \begin{minipage}[b]{\linewidth}\raggedright
Gene
\end{minipage} & \begin{minipage}[b]{\linewidth}\raggedright
Mechanism
\end{minipage} \\
\midrule\noalign{}
\endhead
\bottomrule\noalign{}
\endlastfoot
\textbf{Cystic fibrosis} & AR & \emph{CFTR} & Chloride channel defect;
thick mucus; \texorpdfstring{\ensuremath{\Delta}}{Delta}F508 most common \\
\textbf{Sickle cell disease} & AR & \emph{HBB} & Glu6Val in \texorpdfstring{\ensuremath{\beta}}{beta}-globin;
sickle hemoglobin polymerization \\
\textbf{Tay-Sachs disease} & AR & \emph{HEXA} & \texorpdfstring{\ensuremath{\beta}}{beta}-hexosaminidase A
deficiency; GM2 ganglioside accumulation; founder effect in Ashkenazi
Jews \\
\textbf{Achondroplasia} & AD & \emph{FGFR3} & Constitutively active
FGFR3; chondrocyte proliferation defect; about 80\% de novo
mutations \\
\textbf{Marfan syndrome} & AD & \emph{FBN1} & Fibrillin-1 connective
tissue defect \\
\textbf{Familial hypercholesterolemia} & AD & \emph{LDLR} & Defective
LDL receptor; extreme hypercholesterolemia \\
\textbf{Hereditary nonpolyposis colon cancer (Lynch)} & AD & \emph{MLH1,
MSH2, MSH6, PMS2} & DNA mismatch repair defects \\
\end{longtable}
}

\paragraph{CFTR mutation classes --- a paradigm for genotype--phenotype
maps}\label{cftr-mutation-classes-a-paradigm-for-genotypephenotype-maps}

The 2,000+ identified \emph{CFTR} variants are grouped into \textbf{six
classes} by their effect on protein production and function --- a
classification that now drives \textbf{mutation-specific therapy}:

{\def\LTcaptype{none} % do not increment counter
\begin{longtable}[]{@{}
  >{\raggedright\arraybackslash}p{(\linewidth - 6\tabcolsep) * \real{0.2059}}
  >{\raggedright\arraybackslash}p{(\linewidth - 6\tabcolsep) * \real{0.2353}}
  >{\raggedright\arraybackslash}p{(\linewidth - 6\tabcolsep) * \real{0.2941}}
  >{\raggedright\arraybackslash}p{(\linewidth - 6\tabcolsep) * \real{0.2647}}@{}}
\toprule\noalign{}
\begin{minipage}[b]{\linewidth}\raggedright
Class
\end{minipage} & \begin{minipage}[b]{\linewidth}\raggedright
Defect
\end{minipage} & \begin{minipage}[b]{\linewidth}\raggedright
Examples
\end{minipage} & \begin{minipage}[b]{\linewidth}\raggedright
Therapy
\end{minipage} \\
\midrule\noalign{}
\endhead
\bottomrule\noalign{}
\endlastfoot
I & No protein synthesis (premature stop codons) & W1282X, G542X &
Read-through agents (ataluren --- limited efficacy) \\
II & Misfolding / proteasomal degradation & \textbf{\texorpdfstring{\ensuremath{\Delta}}{Delta}F508}
(about 70\% of CF alleles in Northern Europeans) & Correctors
(lumacaftor, tezacaftor, elexacaftor) \\
III & Channel does not open (``gating defect'') & G551D & Potentiators
(ivacaftor) \\
IV & Reduced channel conductance & R117H & Potentiators \\
V & Reduced protein levels (splicing defects) & 3849+10kb C\texorpdfstring{\ensuremath{\to}}{->}T & Mixed \\
VI & Reduced channel stability at cell surface & rescued \texorpdfstring{\ensuremath{\Delta}}{Delta}F508 &
Stabilizers \\
\end{longtable}
}

The combination therapy \textbf{elexacaftor/tezacaftor/ivacaftor
(Trikafta)}, approved 2019, is effective for about 90\% of CF
patients (those carrying at least one \texorpdfstring{\ensuremath{\Delta}}{Delta}F508) and has dramatically
improved life expectancy --- a paradigm for translating mendelian
genetics into precision medicine.

\subsubsection{Achondroplasia and Marfan as Autosomal Dominant
Models}\label{achondroplasia-and-marfan-as-autosomal-dominant-models}

\textbf{Achondroplasia} is the most common form of dwarfism
(about 1/25,000 births), caused by gain-of-function mutations
in \emph{FGFR3} (fibroblast growth factor receptor 3). The vast majority
(about 98\%) of cases involve a single recurrent mutation:
\textbf{G380R} in the transmembrane domain. The mutation produces a
constitutively active receptor that inhibits chondrocyte proliferation
and endochondral ossification, shortening long bones while sparing
membranous bones (skull vault, vertebrae). Approximately \textbf{80\% of
cases are de novo mutations} arising in the paternal germline, and risk
increases with paternal age --- illustrating that autosomal dominant
disease can persist in populations even when affected individuals have
reduced reproductive fitness, because new mutations replenish the
disease allele each generation. Homozygous achondroplasia (two affected
parents --- about 1 in 4 pregnancies if both are heterozygous)
is lethal in infancy due to severe rib cage abnormalities.

\textbf{Marfan syndrome} illustrates how a single connective-tissue
gene, \emph{FBN1} (fibrillin-1), produces pleiotropic effects across
organ systems. Fibrillin-1 forms microfibrils that scaffold elastic
fibers and sequester latent TGF-\texorpdfstring{\ensuremath{\beta}}{beta}. Mutations cause both
\textbf{structural failure} (loose connective tissue \texorpdfstring{\ensuremath{\to}}{->} tall stature,
long limbs, lens dislocation, joint hypermobility) and
\textbf{dysregulated TGF-\texorpdfstring{\ensuremath{\beta}}{beta} signaling} (aortic root dilation \texorpdfstring{\ensuremath{\to}}{->} dissection
risk). Variable expressivity is dramatic: aortic root sizes among
affected family members carrying the same mutation can vary by
2--3-fold, depending on modifier genes and environment. Modern
management includes annual echocardiography, \texorpdfstring{\ensuremath{\beta}}{beta}-blockers or losartan to
slow aortic dilation, and elective aortic root replacement when diameter
exceeds about 5 cm --- extending life expectancy from a
historical median of about 32 years to nearly normal.

\begin{quote}
\textbf{Clinical Connection: Why Are Some AR Diseases Common?} Cystic
fibrosis carrier frequency of about 1/25 in Northern Europeans
far exceeds what mutation-selection balance alone predicts. Hypothesized
heterozygote advantage against historical pathogens --- possibly
\textbf{typhoid fever} (CFTR is a receptor for \emph{Salmonella typhi})
or \textbf{cholera} (heterozygotes may have reduced fluid secretion in
cholera infection) --- could explain the elevated frequency. Sickle cell
heterozygosity confers strong protection against severe \emph{Plasmodium
falciparum} malaria, and Tay-Sachs heterozygotes may have had protection
against tuberculosis in pre-industrial European cities. These hypotheses
remain debated, but they illustrate how Mendelian disease alleles can be
maintained by ancient selective pressures we no longer experience.
\end{quote}

\subsubsection{X-linked disorders: the molecular basis of male
susceptibility}\label{x-linked-disorders-the-molecular-basis-of-male-susceptibility}

Males are hemizygous for X-linked genes (one X, one Y), so a single
recessive allele on the X chromosome produces the full phenotype in a
male. Females, with two X chromosomes, must be homozygous (or display
unfavorable random X-inactivation patterns) to be fully affected. This
asymmetry --- recognized first by Morgan's \emph{white}-eyed
\emph{Drosophila} --- produces the diagnostic pedigree pattern of
``skipping generations through carrier mothers'' that distinguishes
X-linked recessive inheritance.

\paragraph{Hemophilia A --- F8 and Factor
VIII}\label{hemophilia-a-f8-and-factor-viii}

\textbf{Hemophilia A} (about 1 in 5,000 male births) is caused
by mutations in the \emph{F8} gene at Xq28 encoding \textbf{coagulation
factor VIII}, a 2,332-amino-acid plasma glycoprotein that forms a tenase
complex with activated factor IXa to activate factor X in the intrinsic
coagulation pathway. The clinical phenotype is dose-dependent on
residual factor VIII activity:

{\def\LTcaptype{none} % do not increment counter
\begin{longtable}[]{@{}
  >{\raggedright\arraybackslash}p{(\linewidth - 4\tabcolsep) * \real{0.2000}}
  >{\raggedright\arraybackslash}p{(\linewidth - 4\tabcolsep) * \real{0.4400}}
  >{\raggedright\arraybackslash}p{(\linewidth - 4\tabcolsep) * \real{0.3600}}@{}}
\toprule\noalign{}
\begin{minipage}[b]{\linewidth}\raggedright
Severity
\end{minipage} & \begin{minipage}[b]{\linewidth}\raggedright
Factor VIII activity
\end{minipage} & \begin{minipage}[b]{\linewidth}\raggedright
Bleeding pattern
\end{minipage} \\
\midrule\noalign{}
\endhead
\bottomrule\noalign{}
\endlastfoot
Severe & \textless{} 1\% & Spontaneous joint and muscle hemorrhage;
treated prophylactically \\
Moderate & 1--5\% & Bleeding after minor trauma \\
Mild & 6--40\% & Bleeding primarily after surgery or major trauma \\
\end{longtable}
}

The most common mutation in severe hemophilia A is the \textbf{intron 22
inversion}, which disrupts the \emph{F8} gene and accounts for
about 45\% of severe cases. The gene is unusually large (186
kb, 26 exons), and its 32-kb intron 22 contains the \emph{F8A1} gene
with two homologous repeats elsewhere on Xq28; non-allelic homologous
recombination during male meiosis produces the inversion. Modern
treatment uses \textbf{recombinant factor VIII concentrates} (Kogenate,
Advate) for replacement therapy, \textbf{emicizumab} (a bispecific
antibody mimicking factor VIIIa), and increasingly \textbf{gene
therapy}: in 2022 the FDA approved Hemgenix (etranacogene dezaparvovec)
for hemophilia B, with similar AAV-based gene therapies for hemophilia A
in advanced trials.

The Romanov royal family tragedy is the most famous historical case:
Queen Victoria was a carrier (likely de novo mutation), and through her
descendants the allele entered the Russian, Spanish, and German royal
families. The hemophilia of Tsarevich Alexei contributed to the
political collapse of the Romanov dynasty.

\paragraph{Color blindness --- OPN1LW and
OPN1MW}\label{color-blindness-opn1lw-and-opn1mw}

\textbf{Red-green color blindness} affects about 8\% of males
and about 0.5\% of females of European ancestry. The disorder
reflects defects in two tandem opsin genes on Xq28:

\begin{itemize}
\tightlist
\item
  \textbf{\emph{OPN1LW}} (long-wavelength opsin) --- sensitive to red
  light (about 560 nm peak)
\item
  \textbf{\emph{OPN1MW}} (medium-wavelength opsin) --- sensitive to
  green light (about 530 nm peak)
\end{itemize}

The two genes are \textgreater96\% identical and lie head-to-tail in a
tandem array (1--6 copies of \emph{OPN1MW} per X chromosome) at Xq28.
The high sequence similarity makes them prone to \textbf{unequal
crossing over} during female meiosis, producing chromosomes with
deleted, duplicated, or hybrid (chimeric) opsin genes:

\begin{itemize}
\tightlist
\item
  \textbf{Protan defects} (red-blindness): loss or alteration of
  \emph{OPN1LW}
\item
  \textbf{Deutan defects} (green-blindness): loss or alteration of
  \emph{OPN1MW}
\end{itemize}

Hybrid genes --- formed when the upstream portion of one opsin fuses
with the downstream portion of the other --- produce shifted spectral
sensitivities (anomalous trichromacy). The unique vulnerability of this
gene cluster to misalignment during recombination explains why color
blindness is far more common than other X-linked traits despite no
obvious selective advantage.

\paragraph{Duchenne muscular dystrophy --- DMD and
dystrophin}\label{duchenne-muscular-dystrophy-dmd-and-dystrophin}

\textbf{Duchenne muscular dystrophy (DMD)} affects about 1 in
3,500 male births and is caused by mutations in the \textbf{\emph{DMD}}
gene at Xp21.2 encoding \textbf{dystrophin} --- at 2.4 megabases the
largest known human gene (79 exons), accounting for about 0.1\%
of the entire genome. Dystrophin links the actin cytoskeleton inside
muscle fibers to the dystrophin-glycoprotein complex spanning the
sarcolemma, transmitting force during contraction and protecting the
membrane from mechanical stress.

The genotype--phenotype rule for \emph{DMD} mutations is the
\textbf{reading-frame rule}: deletions that disrupt the reading frame
produce no functional dystrophin and cause \textbf{Duchenne} (severe,
wheelchair by age 12, cardiac/respiratory failure by about 25
years), whereas in-frame deletions produce truncated but partially
functional dystrophin and cause \textbf{Becker muscular dystrophy}
(milder, much later onset, normal lifespan possible). This rule has
therapeutic consequences: \textbf{exon-skipping antisense
oligonucleotides} mask specific exons during splicing to convert an
out-of-frame Duchenne mutation into an in-frame Becker-like protein.
\textbf{Eteplirsen} (FDA-approved 2016) skips exon 51 and is applicable
to about 13\% of DMD patients; \textbf{golodirsen} (exon 53)
and \textbf{viltolarsen} (exon 53) treat additional subgroups.
CRISPR-mediated exon excision is in clinical trials.

The very large size of \emph{DMD} makes it the most mutation-prone gene
in the human genome, with about 1/3 of cases representing de
novo mutations rather than inheritance from a carrier mother --- a fact
that affects recurrence-risk counseling for affected families.

\paragraph{Female carriers --- random X-inactivation and ``manifesting
carriers''}\label{female-carriers-random-x-inactivation-and-manifesting-carriers}

Female carriers of X-linked recessive disorders are typically clinically
normal, but about 10--20\% of \emph{DMD} female carriers show
muscle weakness, cramping, or even cardiomyopathy. The reason is
\textbf{skewed X-inactivation}: by chance, the X chromosome carrying the
wild-type allele is preferentially inactivated in many of her muscle
cells, leaving the mutant X expressed. The clinical severity depends on
the proportion of cells with skewed X-inactivation patterns. Similar
manifesting-carrier phenotypes occur in hemophilia A
(about 10\% of carriers have prolonged bleeding) and ornithine
transcarbamylase deficiency (urea cycle disorder; manifesting carriers
may experience hyperammonemic crises).

\paragraph{Carrier frequency vs.~affected-male frequency: a
population-genetic
asymmetry}\label{carrier-frequency-vs.-affected-male-frequency-a-population-genetic-asymmetry}

A diagnostic feature of X-linked recessive inheritance is that
\textbf{carrier frequency in females greatly exceeds affected-female
frequency}, even though both sexes draw from the same X-linked allele
frequency \(q\). The arithmetic follows directly from hemizygosity in
males and Hardy-Weinberg expectations in females
(\cref{sec:unit_V_population_genetics}):

{\def\LTcaptype{none} % do not increment counter
\begin{longtable}[]{@{}
  >{\raggedright\arraybackslash}p{(\linewidth - 4\tabcolsep) * \real{0.3704}}
  >{\raggedright\arraybackslash}p{(\linewidth - 4\tabcolsep) * \real{0.3333}}
  >{\raggedright\arraybackslash}p{(\linewidth - 4\tabcolsep) * \real{0.2963}}@{}}
\toprule\noalign{}
\begin{minipage}[b]{\linewidth}\raggedright
Quantity
\end{minipage} & \begin{minipage}[b]{\linewidth}\raggedright
Formula
\end{minipage} & \begin{minipage}[b]{\linewidth}\raggedright
Reason
\end{minipage} \\
\midrule\noalign{}
\endhead
\bottomrule\noalign{}
\endlastfoot
Affected males & \(q\) & Hemizygous: one mutant allele suffices \\
Affected females & \(q^2\) & Homozygous: both X alleles must be
mutant \\
Carrier females (heterozygous) & \(2pq \approx 2q\) for small \(q\) &
Heterozygous and asymptomatic with random X-inactivation \\
\end{longtable}
}

Worked numbers for the three X-linked recessive disorders above
(Northern European frequencies):

\begin{itemize}
\tightlist
\item
  \textbf{Hemophilia A}: affected-male frequency \(\approx 1/5{,}000\),
  so \(q \approx 2 \times 10^{-4}\). Carrier-female frequency
  \(2pq \approx 4 \times 10^{-4}\) (about 1 in 2,500 women).
  Affected-female frequency \(q^2 \approx 4 \times 10^{-8}\)
  (about 1 in 25 million) --- meaning carriers outnumber
  affected females by about 10,000-fold.
\item
  \textbf{Duchenne muscular dystrophy}: affected-male frequency
  \(\approx 1/3{,}500\), \(q \approx 2.9 \times 10^{-4}\).
  Carrier-female frequency \(\approx 1/1{,}750\) women.
\item
  \textbf{Red-green color blindness} (high allele frequency):
  affected-male frequency \(\approx 0.08\), so \(q \approx 0.08\).
  Carrier-female frequency \(2pq \approx 0.147\), and affected-female
  frequency \(q^2 \approx 0.0064\) --- recovering the empirical
  about 0.5--1\% color-blind female frequency.
\end{itemize}

Two consequences follow. First, \textbf{most pathogenic X-linked alleles
in the population reside in unaffected female carriers}, who are
invisible to phenotypic screening. Second, when a new disease-causing
X-linked mutation arises (frequently de novo, especially for
\emph{DMD}), it persists for many generations through carrier mothers
before reaching another affected male --- producing the characteristic
``skipping generations through unaffected females'' pattern in
pedigrees. Genetic counseling for X-linked disorders therefore
prioritizes \textbf{carrier testing of female relatives of affected
males}, even when those women appear clinically normal.

\textbf{Concept Check 16.5}

\begin{quote}
\begin{enumerate}
\def\labelenumi{\arabic{enumi}.}
\tightlist
\item
  NF1 shows highly variable expressivity within a single family carrying
  the same mutation. Identify three biological reasons why identical
  genotypes can produce dramatically different phenotypes.
\item
  In a population where 1 in 4,000 males is affected with hemophilia A,
  what fraction of women are carriers? What is the probability that a
  randomly chosen woman is a carrier?
\item
  A man's brother has Duchenne muscular dystrophy (X-linked recessive).
  The brother's mother is confirmed by molecular testing to be a
  non-carrier (the affected brother arose from a de novo mutation).
  Should the man worry that his own future sons could inherit the
  disease through him? Explain.
\item
  Compare the dihybrid F\(_2\) ratio expected when (a) two genes act
  independently in separate biosynthetic pathways and (b) two genes
  encode sequential enzymes in the same pathway where both are required
  for product formation. Which produces 9:3:3:1, which produces 9:7, and
  why?
\item
  A researcher studying Huntington's disease finds that two siblings
  carry identical \emph{HTT} CAG repeat lengths (45 each) but have onset
  ages differing by 12 years. What is this discrepancy a manifestation
  of, and what general lesson does it teach about the genotype-phenotype
  map?
\end{enumerate}
\end{quote}

\begin{center}\rule{0.5\linewidth}{0.5pt}\end{center}

\subsection{GWAS: Mendel Meets the
Genome}\label{gwas-mendel-meets-the-genome}

Mendel studied traits with simple, dichotomous inheritance --- round
vs.~wrinkled, yellow vs.~green. Most human traits and diseases, however,
are \textbf{polygenic} (influenced by many loci) and
\textbf{multifactorial} (influenced by both genes and environment).
Height, body mass index, blood pressure, type 2 diabetes, schizophrenia,
and most common diseases do not segregate as Mendelian traits in
pedigrees --- they are \textbf{complex traits}. The modern tool for
dissecting their genetic architecture is the \textbf{genome-wide
association study (GWAS)}.

\subsubsection{How GWAS works}\label{how-gwas-works}

GWAS scans the genome for \textbf{statistical association} between
genetic variants --- typically common single-nucleotide polymorphisms
(\hyperref[gl:snp]{\textbf{SNP}}s with minor allele frequency \texorpdfstring{\textgreater{}=}{>=} 1\%) ---
and a phenotype, comparing thousands of cases to thousands of controls.
The basic design is a \textbf{case-control association test}: at each
SNP, compare the allele frequency in cases versus controls and ask
whether the difference exceeds chance expectation.

\begin{figure}[htbp]
\centering
\includegraphics[width=0.82\linewidth,height=4.2in,keepaspectratio]{../figures/mermaid_inline/inline_mermaid_0090_bab9b1c04d20.png}
\caption{The GWAS workflow, from cohort recruitment through genome-wide significance testing to functional follow-up. The genome-wide significance threshold of p < 5x10- corresponds to a Bonferroni correction for about 10 independent common variants tested across the genome.}
\end{figure}

\subsubsection{What GWAS has revealed}\label{what-gwas-has-revealed}

Since the first successful GWAS (age-related macular degeneration,
2005), tens of thousands of genome-wide significant associations have
been catalogued. Several lessons have emerged:

\begin{enumerate}
\def\labelenumi{\arabic{enumi}.}
\tightlist
\item
  \textbf{Most common diseases are highly polygenic.} Schizophrenia:
  \textgreater270 loci. Type 2 diabetes: \textgreater400 loci. Height:
  \textgreater12,000 contributing variants. The ``one-gene,
  one-disease'' Mendelian model rarely applies to common conditions.
\item
  \textbf{Effect sizes are small.} Individual common variants typically
  increase disease risk by 1.05--1.3-fold (odds ratios of 1.05--1.3).
  The largest common-variant effects are typically in immune disorders
  (HLA region, OR up to 5--10 for specific autoimmune conditions).
\item
  \textbf{Most GWAS hits are non-coding.} Approximately 90\% of
  GWAS-associated variants lie outside protein-coding regions, in
  regulatory elements (enhancers, promoters, 3\texorpdfstring{'}{'} UTRs). They alter gene
  \textbf{expression} rather than protein sequence --- connecting GWAS
  to the regulatory biology of
  \cref{sec:unit_IV_epigenetics_and_gene_regulation}.
\item
  \textbf{Pleiotropy is the rule, not the exception.} Many GWAS loci
  affect multiple traits. A locus near \emph{FTO} influences both BMI
  and diabetes risk. A locus near \emph{TCF7L2} affects type 2 diabetes,
  fasting glucose, and bone mineral density.
\item
  \textbf{``Missing heritability''} persists --- common GWAS variants
  typically explain 5--25\% of trait heritability, far less than twin
  studies suggest is genetic. The gap reflects rare variants
  (insufficient power), structural variation (poorly captured by SNP
  arrays), gene-environment interactions, and non-additive effects
  (epistasis, dominance).
\end{enumerate}

\subsubsection{Polygenic risk scores
(PRS)}\label{polygenic-risk-scores-prs}

GWAS results can be combined into a \textbf{polygenic risk score} --- a
weighted sum of risk alleles across many SNPs:

\[\text{PRS}_i = \sum_{j=1}^{M} \beta_j \cdot g_{ij} \label{eq:unit_V_mendelian_genetics_item_14}\]

where \(g_{ij}\) is the count of risk alleles (0, 1, or 2) at SNP \(j\)
in individual \(i\), and \(\beta_j\) is the GWAS effect estimate. PRS
provides individual-level risk predictions that, for some diseases,
equal or exceed the predictive power of family history. PRS is becoming
clinically useful in cardiovascular disease, breast cancer, and prostate
cancer screening, though equity concerns remain because most GWAS were
conducted in European-ancestry populations and PRS performance degrades
when transferred to other ancestries.

\subsubsection{GWAS as the modern face of Mendel's
program}\label{gwas-as-the-modern-face-of-mendels-program}

Mendel's deepest contribution was not the laws themselves but the
\textbf{method}: count phenotypes, treat heredity as a probabilistic
system, and test genetic hypotheses against quantitative ratios. GWAS is
Mendel's program rescaled to millions of loci and billions of genotype
calls --- every Manhattan plot is a vast, statistically corrected
dihybrid cross conducted across the entire genome.

\begin{quote}
\textbf{Clinical Connection: When GWAS Becomes Drug Discovery} GWAS hits
sometimes identify proteins that become drug targets. The discovery that
\emph{PCSK9} loss-of-function variants confer protection against
cardiovascular disease (low LDL cholesterol, no apparent harm) led
directly to the development of monoclonal antibodies (evolocumab,
alirocumab) and small interfering RNAs (inclisiran) inhibiting PCSK9 ---
now major cholesterol-lowering therapies. Similarly, \emph{HMGCR}
variants supported statins as a drug class. Genetics-informed drugs are
roughly twice as likely to succeed in clinical trials as drugs developed
without genetic support.
\end{quote}

\subsubsection{The Manhattan Plot and Genome-Wide
Significance}\label{the-manhattan-plot-and-genome-wide-significance}

The standard visual display of GWAS results is the \textbf{Manhattan
plot} --- a scatter of \(-\log_{10}(p)\) values plotted against genomic
position. The horizontal axis traverses chromosomes 1 through 22 (and X)
in order; each point is a single SNP. Strong associations form vertical
``skyscrapers'' of co-associated SNPs in the same locus (because nearby
SNPs are in linkage disequilibrium), giving the plot its skyline-like
appearance.

The genome-wide significance threshold is conventionally set at
\(p < 5 \times 10^{-8}\) --- the line drawn horizontally across most
Manhattan plots. This threshold derives from a \textbf{Bonferroni
correction} for multiple testing:

\[p_{\text{threshold}} = \frac{\alpha}{\text{number of independent tests}} = \frac{0.05}{10^6} = 5 \times 10^{-8} \label{eq:unit_V_mendelian_genetics_item_15}\]

The denominator of \(\sim 10^6\) reflects the number of approximately
independent common variants in the human genome (after accounting for
linkage disequilibrium between adjacent SNPs). A ``suggestive''
threshold of \(p < 10^{-5}\) is sometimes reported for hits that warrant
follow-up but do not meet genome-wide significance. Replication in an
independent cohort is the gold standard for confirming a true
association.

\subsubsection{Polygenic Risk Scores: Detailed
Example}\label{polygenic-risk-scores-detailed-example}

Consider type 2 diabetes (T2D), one of the most extensively GWAS-studied
complex diseases. Over 400 independent loci have been identified, each
with small individual effects (odds ratios typically 1.05--1.20). A
polygenic risk score combines these signals into an individual-level
prediction:

\begin{equation}
\text{PRS}_i = \sum_{j=1}^{M} \beta_j \cdot g_{ij}
\label{eq:unit_V_mendelian_prs}
\end{equation}

where \(g_{ij}\) is the dosage of the risk allele (0, 1, or 2) at SNP
\(j\) for individual \(i\), \(\beta_j\) is the GWAS-derived effect size
for that SNP (usually log-odds ratio for case-control studies), and
\(M\) is the number of SNPs included (often hundreds of thousands when
modern PRS methods like LDpred or PRS-CS are used).

\textbf{Clinical performance}: Modern PRS for T2D can identify
individuals in the top percentile of genetic risk who have
about 3--5\texorpdfstring{\ensuremath{\times}}{x} the population baseline lifetime risk of developing
the disease --- comparable in predictive power to having a single
first-degree relative with the disease, but available decades before
clinical onset and therefore actionable through lifestyle modification.

\textbf{Equity caveat}: Most large GWAS to date have been conducted in
European-ancestry populations. PRS performance degrades substantially
when applied to other ancestries, with predictive accuracy 50--75\%
lower in African-ancestry individuals. Closing this ``PRS performance
gap'' requires diverse GWAS cohorts and cross-ancestry methods --- an
active area of research and a matter of healthcare equity.

\subsubsection{Clinical Interpretation and Governance
Checkpoint}\label{clinical-interpretation-and-governance-checkpoint}

A PRS should be treated as a calibrated, ancestry- and context-dependent
risk model, not as a diagnosis. A responsible report states the
discovery cohort, target population, calibration method, confidence
interval, and action threshold; checks whether the risk strata were
validated in the patient's ancestry and care setting; and interprets the
score alongside family history, monogenic testing when indicated, and
modifiable clinical risk factors. Current PRS accuracy is several-fold
higher in European-ancestry cohorts than in other cohorts, so premature
deployment can widen health gaps if diverse validation does not keep
pace \citep{martin2019prsdisparities}.

Governance is part of the biological interpretation. Consent should
specify secondary use, data retention, reidentification risk, and which
relatives might be indirectly implicated. In the United States, GINA
addresses health-insurance and employment discrimination, but it does
not extend to life, disability, or long-term-care insurance
\citep{hhs2009gina}. A genetic-counseling workflow therefore separates
what the score estimates from what clinicians, insurers, employers,
families, or law-enforcement agencies are permitted to do with it.

\textbf{Concept Check 16.3}

\begin{quote}
\begin{enumerate}
\def\labelenumi{\arabic{enumi}.}
\tightlist
\item
  Why is the genome-wide significance threshold \(p < 5 \times 10^{-8}\)
  and not the conventional \(p < 0.05\)?
\item
  A father with 42 CAG repeats has Huntington's disease. His daughter
  inherits the affected allele. What can you predict about her CAG count
  and age of onset?
\item
  PKU is treatable with diet, yet newborn screening remains
  comprehensive in many countries. Why is comprehensive screening
  preferred over targeted screening of high-risk families?
\item
  Most GWAS hits lie outside protein-coding regions. What does this tell
  us about the genetic architecture of common diseases?
\item
  A patient's PRS for coronary artery disease places her in the top 1\%
  of the population. Her standard cholesterol panel is normal at age 30.
  How should this affect her clinical management, and what are the
  limitations of acting on PRS alone?
\item
  Why does PRS performance degrade when applied across ancestries? What
  practical and ethical issues does this raise for clinical deployment?
\end{enumerate}
\end{quote}

\begin{center}\rule{0.5\linewidth}{0.5pt}\end{center}

\subsection{\texorpdfstring{The Chi-Squared (\(\chi^2\))
Test}{The Chi-Squared (chi^2) Test}}\label{the-chi-squared-chi2-test}

The \(\chi^2\) test for goodness of fit evaluates whether observed
offspring ratios deviate significantly from expected Mendelian ratios.

\subsubsection{Formula and Procedure}\label{formula-and-procedure}

\begin{equation}
\chi^2 = \sum_{i=1}^{k} \frac{(O_i - E_i)^2}{E_i}
\label{eq:unit_V_mendelian_chi_squared}
\end{equation}

Where: - \(O_i\) = observed count in category \(i\) - \(E_i\) = expected
count in category \(i\) (based on the hypothesis) - \(k\) = number of
categories

\textbf{Degrees of freedom (df)} = \(k - 1\)

\textbf{Decision rule}: Compare calculated \(\chi^2\) to critical value
from the chi-squared distribution at significance level
\(\alpha = 0.05\). If \(\chi^2_{calc} > \chi^2_{crit}\), \textbf{reject}
the null hypothesis (the data do not fit the expected ratio).

\textbf{Critical values:}

{\def\LTcaptype{none} % do not increment counter
\begin{longtable}[]{@{}lll@{}}
\toprule\noalign{}
df & \(\alpha = 0.05\) & \(\alpha = 0.01\) \\
\midrule\noalign{}
\endhead
\bottomrule\noalign{}
\endlastfoot
1 & 3.841 & 6.635 \\
2 & 5.991 & 9.210 \\
3 & 7.815 & 11.345 \\
4 & 9.488 & 13.277 \\
\end{longtable}
}

\subsubsection{Worked Example}\label{worked-example}

\textbf{Problem}: F\(_2\) from a monohybrid cross, 100 offspring
observed: 68 dominant, 32 recessive. Expected ratio: 3:1 (H\(_0\): the
trait follows simple Mendelian inheritance).

\textbf{Step 1}: Calculate expected values. - Expected dominant =
\(100 \times \frac{3}{4} = 75\) - Expected recessive =
\(100 \times \frac{1}{4} = 25\)

\textbf{Step 2}: Calculate \(\chi^2\).

\[\chi^2 = \frac{(68-75)^2}{75} + \frac{(32-25)^2}{25} = \frac{49}{75} + \frac{49}{25} = 0.653 + 1.960 = 2.613 \tag{16.12} \label{eq:unit_V_mendelian_genetics_item_16}\]

\textbf{Step 3}: Determine df = 2 - 1 = 1.

\textbf{Step 4}: Compare to critical value.
\(\chi^2_{crit}(df=1, \alpha=0.05) = 3.841\).

Since \(2.613 < 3.841\), we \textbf{fail to reject} H\(_0\). The
observed data are consistent with a 3:1 Mendelian ratio
(\(p \approx 0.106\)).

\subsubsection{Worked Example: Dihybrid Chi-Squared
Test}\label{worked-example-dihybrid-chi-squared-test}

\textbf{Problem}: A dihybrid cross produces: 315 round yellow, 108 round
green, 101 wrinkled yellow, 32 wrinkled green (total = 556). Test the
9:3:3:1 hypothesis.

\textbf{Expected values}: \(\frac{9}{16} \times 556 = 312.75\);
\(\frac{3}{16} \times 556 = 104.25\);
\(\frac{3}{16} \times 556 = 104.25\);
\(\frac{1}{16} \times 556 = 34.75\)

\[\chi^2 = \frac{(315-312.75)^2}{312.75} + \frac{(108-104.25)^2}{104.25} + \frac{(101-104.25)^2}{104.25} + \frac{(32-34.75)^2}{34.75} \tag{16.13} \label{eq:unit_V_mendelian_genetics_item_17}\]

\[= 0.016 + 0.135 + 0.101 + 0.218 = 0.470 \tag{16.14} \label{eq:unit_V_mendelian_genetics_item_18}\]

df = 4 - 1 = 3. \(\chi^2_{crit}(df=3, \alpha=0.05) = 7.815\).

Since \(0.470 < 7.815\), we fail to reject H\(_0\). The data are
consistent with a 9:3:3:1 ratio (\(p > 0.90\)). These are, in fact,
Mendel's actual data for seed shape and color.

\begin{quote}
\textbf{Clinical Connection: Using Chi-Squared in Genetic Counseling} In
clinical genetics, the chi-squared test (or exact binomial tests for
small families) helps evaluate whether a family's disease pattern is
consistent with a specific inheritance model. For example, if a rare
disease appears in 5 of 20 siblings (25\%), a chi-squared test could
evaluate whether this is consistent with autosomal recessive inheritance
(expected 25\% affected). However, for small family sizes, exact
probability calculations are preferred.
\end{quote}

\begin{center}\rule{0.5\linewidth}{0.5pt}\end{center}

\subsection{Pedigree Analysis}\label{pedigree-analysis}

\subsubsection{Pedigree Symbols and
Conventions}\label{pedigree-symbols-and-conventions}

\begin{itemize}
\tightlist
\item
  \textbf{Squares}: males; \textbf{Circles}: females
\item
  \textbf{Filled symbols}: affected individuals
\item
  \textbf{Half-filled}: carrier (when known)
\item
  \textbf{Horizontal line}: mating pair
\item
  \textbf{Vertical line}: parent-offspring connection
\item
  \textbf{Double horizontal line}: consanguineous mating (mating between
  relatives)
\item
  \textbf{Roman numerals}: generations (I, II, III\ldots)
\item
  \textbf{Arabic numerals}: individuals within a generation (1, 2,
  3\ldots)
\end{itemize}

\subsubsection{Inheritance Pattern
Recognition}\label{inheritance-pattern-recognition}

\begin{figure}[htbp]
\centering
\includegraphics[width=0.82\linewidth,height=4.2in,keepaspectratio]{../figures/mermaid_inline/inline_mermaid_0091_e26ca442a73b.png}
\caption{Decision flowchart for pedigree analysis. Systematic questions about affected individuals across generations distinguish inheritance patterns.}
\end{figure}

{\def\LTcaptype{none} % do not increment counter
\begin{longtable}[]{@{}
  >{\raggedright\arraybackslash}p{(\linewidth - 2\tabcolsep) * \real{0.3000}}
  >{\raggedright\arraybackslash}p{(\linewidth - 2\tabcolsep) * \real{0.7000}}@{}}
\toprule\noalign{}
\begin{minipage}[b]{\linewidth}\raggedright
Pattern
\end{minipage} & \begin{minipage}[b]{\linewidth}\raggedright
Key Diagnostic Clues
\end{minipage} \\
\midrule\noalign{}
\endhead
\bottomrule\noalign{}
\endlastfoot
\textbf{Autosomal dominant} & Affected in every generation; unaffected
parents produce primarily unaffected children; males and females equally
affected; affected individuals have at least one affected parent \\
\textbf{Autosomal recessive} & Often skips generations; affected
children from unaffected (carrier) parents; consanguinity increases
risk; males and females equally affected \\
\textbf{X-linked recessive} & Males affected far more than females;
carrier mothers pass to about 50\% of sons; affected fathers
cannot pass to sons (but most daughters are carriers); no male-to-male
transmission \\
\textbf{X-linked dominant} & Both sexes affected but more females than
males; affected father passes the allele to daughters but not sons;
affected heterozygous mother passes to about 50\% of sons and
about 50\% of daughters \\
\textbf{Mitochondrial} & Most children of an affected mother are
affected; affected father does NOT transmit; maternal inheritance
primarily \\
\textbf{Y-linked} & Primarily males affected; most sons of affected
father are affected \\
\end{longtable}
}

\subsubsection{Worked Pedigree Problem}\label{worked-pedigree-problem}

\textbf{Problem}: A pedigree shows: - Generation I: Unaffected father,
unaffected mother - Generation II: 3 unaffected daughters, 2 affected
sons, 1 unaffected son - One affected son (Generation II) marries an
unaffected woman - Generation III: 2 unaffected sons, 1 unaffected
daughter

\textbf{Analysis}: 1. Trait skips Generation I (unaffected parents have
affected children) -- suggests \textbf{recessive} 2. Primarily males
affected -- suggests \textbf{X-linked} 3. Affected father (Gen II) has
unaffected sons (Gen III) -- consistent with X-linked (father gives Y to
sons) 4. Most daughters of affected father are unaffected -- but they
are obligate carriers (\(X^AX^a\)) if the mother is \(X^AX^A\) or
\(X^AY\) -- wait, the mother is female so \(X^AX^A\) or \(X^AX^a\)

\textbf{Conclusion}: \textbf{X-linked recessive} inheritance. Mother in
Gen I is a carrier (\(X^HX^h\)); Father in Gen I is unaffected
(\(X^HY\)).

Probability that Gen III daughter (from affected father \(X^hY\) x
unaffected mother \(X^HX^H\)) is a carrier:

Most daughters of an affected father receive his \(X^h\): probability =
\textbf{1 (100\% carriers)}.

\subsection{Worked Example: 16.4: Conditional probability in
pedigrees.}\label{worked-example-16.4-conditional-probability-in-pedigrees.}

A couple seeks genetic counseling. The woman's brother has cystic
fibrosis (CF, autosomal recessive, \(q = 0.02\) in the general
population). The man has no family history.

\begin{enumerate}
\def\labelenumi{(\alph{enumi})}
\tightlist
\item
  What is the probability the woman is a carrier?
\end{enumerate}

Her parents must both be carriers (Cc x Cc). Among unaffected siblings,
the probability of being a carrier:

\[P(\text{Cc} | \text{unaffected}) = \frac{P(\text{Cc})}{P(\text{unaffected})} = \frac{2/4}{3/4} = \frac{2}{3} \tag{16.15} \label{eq:unit_V_mendelian_genetics_item_19}\]

\begin{enumerate}
\def\labelenumi{(\alph{enumi})}
\setcounter{enumi}{1}
\tightlist
\item
  What is the probability the man is a carrier?
\end{enumerate}

Using Hardy-Weinberg: carrier frequency =
\(2pq \approx 2(0.98)(0.02) \approx \frac{1}{25}\)

\begin{enumerate}
\def\labelenumi{(\alph{enumi})}
\setcounter{enumi}{2}
\tightlist
\item
  What is the probability their first child will have CF?
\end{enumerate}

\[P = P(\text{mom carrier}) \times P(\text{dad carrier}) \times P(\text{child cc} | \text{both carriers}) = \frac{2}{3} \times \frac{1}{25} \times \frac{1}{4} = \frac{2}{300} = \frac{1}{150} \tag{16.16} \label{eq:unit_V_mendelian_genetics_item_20}\]

\textbf{Concept Check 16.4}

\begin{quote}
\begin{enumerate}
\def\labelenumi{\arabic{enumi}.}
\tightlist
\item
  An affected father and carrier mother have children. For an X-linked
  recessive trait, what fraction of daughters will be affected? What
  fraction of sons?
\item
  Why does consanguinity increase the risk of autosomal recessive
  disorders?
\item
  Can mitochondrial inheritance be distinguished from autosomal dominant
  maternal inheritance in a single pedigree? What additional information
  would help?
\end{enumerate}
\end{quote}

\begin{center}\rule{0.5\linewidth}{0.5pt}\end{center}

\subsection{Worked Example: Complex Multi-Locus
Genetics}\label{worked-example-complex-multi-locus-genetics}

\textbf{Problem}: In mice, coat color is controlled by two independently
assorting loci. At the A locus: A (agouti) is dominant to a
(non-agouti/black). At the C locus: C (colored) is dominant to c
(albino, epistatic -- no pigment regardless of A genotype). This is
recessive epistasis.

\begin{enumerate}
\def\labelenumi{(\alph{enumi})}
\tightlist
\item
  Cross AaCc x AaCc. Predict the F\(_2\) phenotype ratio.
\end{enumerate}

Using branching: - A locus: \(\frac{3}{4}\) A\_ (agouti pattern) :
\(\frac{1}{4}\) aa (black pattern) - C locus: \(\frac{3}{4}\) C\_
(colored) : \(\frac{1}{4}\) cc (albino)

{\def\LTcaptype{none} % do not increment counter
\begin{longtable}[]{@{}lll@{}}
\toprule\noalign{}
Genotype & Phenotype & Fraction \\
\midrule\noalign{}
\endhead
\bottomrule\noalign{}
\endlastfoot
A\_C\_ & Agouti & \(\frac{9}{16}\) \\
aaC\_ & Black & \(\frac{3}{16}\) \\
A\_cc & Albino & \(\frac{3}{16}\) \\
aacc & Albino & \(\frac{1}{16}\) \\
\end{longtable}
}

\textbf{Phenotype ratio: 9 agouti : 3 black : 4 albino (9:3:4)}

\begin{enumerate}
\def\labelenumi{(\alph{enumi})}
\setcounter{enumi}{1}
\tightlist
\item
  Among the albino offspring, what fraction are heterozygous at both
  loci (AaCc)?
\end{enumerate}

Wait -- albino mice are cc, so they must be \_\_cc. Among A\_cc
offspring (fraction \(\frac{3}{16}\)), the Aa fraction = \(\frac{2}{3}\)
of A\_. And most are cc (no Cc among albinos).

The question asks for AaCc, but cc individuals cannot be Cc. Therefore
the answer is \textbf{0} -- no albino offspring can be AaCc.

If the question intended Aacc: Among most albino mice
(\(\frac{4}{16}\)), the fraction that is Aacc =
\(\frac{2/16}{4/16} = \frac{1}{2}\).

\begin{center}\rule{0.5\linewidth}{0.5pt}\end{center}

\subsection{Computational Bridge}\label{computational-bridge-15}

Mendelian ratios are enumerated exactly by \texttt{punnett\_square}:

\begin{Shaded}
\begin{Highlighting}[]
\ImportTok{from}\NormalTok{ biology.genetics }\ImportTok{import}\NormalTok{ punnett\_square}

\NormalTok{r }\OperatorTok{=}\NormalTok{ punnett\_square(}\StringTok{"Aa"}\NormalTok{, }\StringTok{"Aa"}\NormalTok{)}
\BuiltInTok{print}\NormalTok{(r.genotype\_ratios, r.phenotype\_ratios)}
\end{Highlighting}
\end{Shaded}

\begin{quote}
\textbf{Clinical / systems note:} Pedigree risk counselling translates
identical probability trees into prior odds for carrier screening (CF,
SMA, Ashkenazi panels) before any DNA test is ordered.
\end{quote}

\begin{center}\rule{0.5\linewidth}{0.5pt}\end{center}

\subsection{Current Evidence and Frontier
Biology}\label{current-evidence-and-frontier-biology-18}

For \textbf{Mendelian Genetics and Heredity}, frontier biology belongs
inside the evidence logic of the chapter. Classical genetics remains
essential, but modern interpretation adds penetrance, polygenicity,
structural variation, ancestry-aware inference, and uncertainty in risk
prediction. The core reading question is this: Mendelian patterns are
starting models that must be qualified by penetrance, linkage,
environment, and sampling.

\begin{itemize}
\tightlist
\item
  \textbf{What to verify:} identify the observation, model, assay, or
  dataset that would make the claim stronger or weaker.
\item
  \textbf{What to qualify:} state the scale, organism, cell type,
  environmental condition, or population where the claim is expected to
  hold.
\item
  \textbf{What to compare:} test at least one alternative explanation,
  baseline, or null model before treating the pattern as causal.
\item
  \textbf{What to cite:} distinguish primary evidence, review synthesis,
  public dataset, and institutional guidance; for recent or numeric
  claims, prefer the source closest to the measurement and state what
  has changed since it was published.
\end{itemize}

A good genetics answer separates the Mendelian transmission model from
the evidence needed to use it in a population, family, or clinical
setting.

\textbf{Source practice:} For inheritance and population claims,
separate the model assumptions from sampling, ancestry representation,
penetrance, linkage, and environment.

\subsection{Summary}\label{summary-18}

\begin{itemize}
\tightlist
\item
  \textbf{Mendel's experiments}: Garden peas, quantitative analysis of 7
  characters, statistical reasoning -- the first application of
  statistics to biology.
\item
  \textbf{Law of Segregation}: Two alleles separate at meiosis I
  (anaphase I); each gamete gets one allele.
\item
  \textbf{Law of Independent Assortment}: Genes on different chromosomes
  orient independently at metaphase I; exceptions for linked genes.
\item
  \textbf{Molecular basis of Mendel's laws}: Sutton-Boveri chromosome
  theory unifies abstract ``factors'' with homolog behavior at meiosis;
  every Punnett square is a model of meiotic chromosome dynamics.
\item
  \textbf{Monohybrid F\(_2\)}: 3:1 phenotype ratio; 1:2:1 genotype
  ratio. Testcross confirms genotype.
\item
  \textbf{Dihybrid F\(_2\)}: 9:3:3:1 phenotype ratio (product of two
  independent 3:1 ratios).
\item
  \textbf{Non-Mendelian inheritance}: Incomplete dominance (1:2:1
  phenotype), codominance (both alleles expressed), multiple alleles
  (ABO), pleiotropy (sickle cell), polygenic traits (continuous
  distribution), epistasis (modified dihybrid ratios --- 9:3:4, 12:3:1,
  9:7, 15:1, 13:3), genomic imprinting (parent-of-origin expression),
  maternal effect (mother's genotype determines offspring phenotype).
\item
  \textbf{Penetrance and expressivity}: Reduced penetrance (BRCA1,
  retinoblastoma) and variable expressivity (NF1, Marfan) explain why
  identical genotypes produce divergent clinical outcomes.
\item
  \textbf{Human genetic disorders}: Huntington's disease (CAG repeat
  expansion, autosomal dominant, anticipation); PKU (autosomal
  recessive, screened almost universally at birth, dietary therapy); CF,
  sickle cell, Tay-Sachs as paradigm autosomal recessive disorders
  maintained at elevated frequencies by likely heterozygote advantage.
\item
  \textbf{GWAS}: Genome-wide association studies extend Mendel's logic
  to millions of common variants, revealing that complex traits are
  highly polygenic, that most causal variants are non-coding, and that
  polygenic risk scores can predict disease risk at the individual
  level.
\item
  \textbf{Chi-squared test}: \(\chi^2 = \sum \frac{(O-E)^2}{E}\); df =
  k-1; compare to critical value at \(\alpha = 0.05\).
\item
  \textbf{Pedigree analysis}: Autosomal dominant (every generation,
  equal sex ratio), autosomal recessive (skips generations,
  consanguinity), X-linked recessive (males \textgreater\textgreater{}
  females, no male-to-male), mitochondrial (maternal primarily).
\item
  \textbf{Connections:} See \cref{sec:unit_V_chromosomal_inheritance}
  for linkage and sex chromosomes, \cref{sec:unit_V_population_genetics}
  for allele frequencies in populations, and Unit IV for the molecular
  basis of alleles.
\end{itemize}

\begin{center}\rule{0.5\linewidth}{0.5pt}\end{center}

\subsection{Review Questions}\label{review-questions-18}

\begin{enumerate}
\def\labelenumi{\arabic{enumi}.}
\tightlist
\item
  Explain why Mendel's laws were not widely recognized for 34 years
  after their publication. What developments in cell biology (1900) made
  them comprehensible?
\item
  Perform a trihybrid cross AaBbCc x AaBbCc. What fraction of offspring
  will show most three dominant phenotypes? What fraction will be
  heterozygous at most three loci?
\item
  A snapdragon cross between two pink flowers (\(C^RC^W\) x \(C^RC^W\))
  produces 240 offspring. How many of each color would you expect?
  Perform a chi-squared test if the observed counts are 55 red, 128
  pink, 57 white.
\item
  In the ABO blood group system, a type A mother and type B father have
  a type O child. What are the parents' genotypes? What is the
  probability their next child will be type AB?
\item
  Describe three examples of pleiotropy and explain how one gene can
  affect multiple organ systems.
\item
  A Labrador cross produces 48 black, 16 chocolate, and 20 yellow
  puppies. Perform a chi-squared test for a 9:3:4 ratio.
\item
  Explain the molecular basis of Prader-Willi and Angelman syndromes.
  Why do deletions of the same chromosomal region cause different
  diseases?
\item
  A woman whose father had hemophilia A marries an unaffected man. What
  is the probability their first son will be affected? Their first
  daughter?
\item
  Using the Hardy-Weinberg equation, calculate the carrier frequency for
  phenylketonuria (PKU) if the disease frequency is 1/10,000.
\item
  Explain why the chi-squared test can rarely prove that a hypothesis is
  correct -- primarily fail to reject it.
\item
  For a dominant disorder with 80\% penetrance, how does risk to
  offspring of an affected \texorpdfstring{\ensuremath{\times}}{x} unaffected cross change compared with
  complete penetrance?
\item
  Explain why mitochondrial pedigrees show \textbf{most-maternal}
  transmission yet may still include unaffected mothers.
\item
  A man with 38 CAG repeats in \emph{HTT} fathers two children. Predict
  the inheritance, repeat-count distribution, and clinical phenotypes of
  his offspring. What does this illustrate about anticipation?
\item
  Explain why comprehensive newborn screening for PKU has been more
  cost-effective than family-history-based screening.
\item
  A GWAS for type 2 diabetes identifies a SNP with odds ratio 1.15
  (\(p = 2 \times 10^{-12}\)). Why is this a ``successful'' finding even
  though the odds ratio is small?
\end{enumerate}

\begin{center}\rule{0.5\linewidth}{0.5pt}\end{center}

\subsection{Further Reading and Source
Notes}\label{further-reading-and-source-notes-18}

\begin{itemize}
\tightlist
\item
  Mendel (1866). Versuche \{"u\}ber Pflanzenhybriden.
  \emph{Verhandlungen des naturforschenden Vereines in Br\{"u\}nn}, 4.
\item
  Bateson (1909). \emph{Mendel's Principles of Heredity}. Cambridge
  University Press.
\item
  Morgan (1910). Sex limited inheritance in Drosophila. \emph{Science},
  32.
\item
  Sturtevant (1913). The linear arrangement of six sex-linked factors in
  Drosophila, as shown by their mode of association. \emph{Journal of
  Experimental Zoology}, 14.
\item
  Punnett (1907). \emph{Mendelism}. Macmillan.
\item
  Griffiths et al.~(latest ed.). \emph{Introduction to Genetic Analysis}
  (Mendelian inheritance chapters). W. H. Freeman.
\end{itemize}

\begin{center}\rule{0.5\linewidth}{0.5pt}\end{center}

\subsection{Key Terms}\label{key-terms-18}

\begin{enumerate}
\def\labelenumi{\arabic{enumi}.}
\tightlist
\item
  \textbf{Allele} -- alternative form of a gene at a given locus
\item
  \textbf{Dominant} -- allele whose phenotype is expressed in
  heterozygotes
\item
  \textbf{Recessive} -- allele whose phenotype is masked in
  heterozygotes
\item
  \textbf{Homozygous} -- two identical alleles at a locus (AA or aa)
\item
  \textbf{Heterozygous} -- two different alleles at a locus (Aa)
\item
  \textbf{Genotype} -- the genetic constitution of an organism at
  specified loci
\item
  \textbf{Phenotype} -- the observable characteristics resulting from
  genotype and environment
\item
  \textbf{Testcross} -- cross of an individual with dominant phenotype
  to a homozygous recessive
\item
  \textbf{Incomplete dominance} -- heterozygote phenotype intermediate
  between two homozygotes
\item
  \textbf{Codominance} -- both alleles fully expressed in the
  heterozygote
\item
  \textbf{Epistasis} -- interaction where one gene masks the expression
  of another gene
\item
  \textbf{Pleiotropy} -- one gene influencing multiple phenotypic traits
\item
  \textbf{Polygenic inheritance} -- multiple genes contributing
  additively to a continuous trait
\item
  \textbf{Penetrance} -- proportion of individuals with a genotype who
  express the expected phenotype
\item
  \textbf{Expressivity} -- degree of phenotype expression among
  penetrant individuals
\item
  \textbf{Phenocopy} -- environmentally induced phenotype mimicking a
  genetic condition
\item
  \textbf{Anticipation} -- earlier onset and increased severity across
  generations, characteristic of trinucleotide repeat expansion
  disorders
\item
  \textbf{Genomic imprinting} -- parent-of-origin-specific gene
  expression via epigenetic marks
\item
  \textbf{Chi-squared test} -- statistical test comparing observed to
  expected frequencies
\item
  \textbf{Pedigree} -- diagram of family phenotypes across generations
  used to determine inheritance patterns
\item
  \textbf{Consanguinity} -- mating between related individuals;
  increases homozygosity
\item
  \textbf{GWAS} -- genome-wide association study; statistical scan of
  common variants for trait associations across the genome
\item
  \textbf{Polygenic risk score (PRS)} -- weighted sum of risk alleles
  producing an individual-level genetic risk estimate
\item
  \textbf{Trinucleotide repeat expansion} -- mutation class in which an
  unstable triplet repeat expands across generations, causing diseases
  such as Huntington's, fragile X, and myotonic dystrophy
\end{enumerate}

\begin{center}\rule{0.5\linewidth}{0.5pt}\end{center}

\subsubsection{Companion Source
Module}\label{companion-source-module-18}

\textbf{Mendelian Genetics and Heredity} should leave a reproducible
trail from a biological claim to the code, figure, diagram, or
paper-based activity that can test it. Use the surfaces below to inspect
the chapter's assumptions, rerun the relevant model, or compare the
manuscript explanation with companion labs and figures.

{\def\LTcaptype{none} % do not increment counter
\begin{longtable}[]{@{}
  >{\raggedright\arraybackslash}p{(\linewidth - 2\tabcolsep) * \real{0.5000}}
  >{\raggedright\arraybackslash}p{(\linewidth - 2\tabcolsep) * \real{0.5000}}@{}}
\toprule\noalign{}
\begin{minipage}[b]{\linewidth}\raggedright
Surface
\end{minipage} & \begin{minipage}[b]{\linewidth}\raggedright
Use it for
\end{minipage} \\
\midrule\noalign{}
\endhead
\bottomrule\noalign{}
\endlastfoot
\breaktt{src/biology/genetics/genetics.py} (\texttt{punnett\_square},
\texttt{hardy\_weinberg}, \texttt{chi\_squared\_test}) & Reproduce
inheritance ratios, equilibrium expectations, and goodness-of-fit
tests. \\
\breaktt{src/visualization/plots.py} (\breaktt{plot\_punnett\_square}) &
Check genotype and phenotype tables visually. \\
\breaktt{src/mermaid/biology\_diagrams.py}
(\breaktt{mendelian\_cross\_diagram}) & Link segregation logic to
diagrammed crosses. \\
\end{longtable}
}

\textbf{Reproducibility check:} state genotype notation, dominance
model, sample size, and statistical expectation before interpreting a
ratio. \textbf{Cross-reference:} compare with
\cref{sec:unit_V_chromosomal_inheritance} and
\cref{sec:unit_V_population_genetics}.

\newpage

\section{Chromosomal Inheritance and
Linkage}\label{chromosomal-inheritance-and-linkage}

\label{sec:unit_V_chromosomal_inheritance}

\begin{quote}
\textbf{Ch 17} \texorpdfstring{\ensuremath{\cdot}}{.} Level 2/3 \texorpdfstring{\ensuremath{\cdot}}{.} 60 min read \texorpdfstring{\ensuremath{\cdot}}{.} 75 min lecture \texorpdfstring{\ensuremath{\cdot}}{.}
Prerequisites: \cref{sec:unit_V_mendelian_genetics}
\end{quote}

\subsection{Learning Objectives}\label{learning-objectives-19}

\begin{enumerate}
\def\labelenumi{\arabic{enumi}.}
\tightlist
\item
  Explain the \hyperref[gl:chromosome]{\textbf{chromosome}} theory of
  heredity and its experimental support from Sutton, Boveri, and Morgan.
\item
  Describe chromosome architecture (centromere, telomere,
  heterochromatin/euchromatin) and connect it to chromosome behavior at
  meiosis.
\item
  Compare meiosis and mitosis at the level of cellular events, ploidy,
  recombination, and outcomes.
\item
  Describe the sex determination systems used across different
  organisms.
\item
  Explain the molecular mechanism of X-inactivation, including the role
  of XIST RNA, and outline broader dosage-compensation strategies.
\item
  Describe X-linked inheritance patterns and calculate expected
  offspring ratios for X-linked traits.
\item
  Explain meiotic non-disjunction (meiosis I vs.~meiosis II), aneuploidy
  outcomes, and the maternal age effect.
\item
  Define genetic \hyperref[gl:linkage]{\textbf{linkage}} and
  \hyperref[gl:recombination]{\textbf{recombination}}; calculate map
  distance from cross data.
\item
  Perform three-point test cross analysis to determine
  \hyperref[gl:gene]{\textbf{gene}} order and map distances.
\item
  Classify chromosomal rearrangements (translocations, inversions,
  deletions, duplications) and their clinical consequences, including
  the Philadelphia chromosome, DiGeorge, and Cri-du-chat syndromes.
\item
  Connect chromosome behavior to genomic imprinting and uniparental
  disomy (cross-link to
  \cref{sec:unit_IV_epigenetics_and_gene_regulation}).
\item
  Calculate inbreeding coefficients and predict their effect on
  offspring homozygosity and fitness.
\end{enumerate}

\subsubsection{Study Blueprint}\label{study-blueprint-19}

\begin{itemize}
\tightlist
\item
  \textbf{Big idea:} Genes travel on chromosomes, so linkage,
  recombination, and chromosome structure shape inheritance.
\item
  \textbf{Core concepts:} linkage, recombination, sex linkage,
  chromosomal rearrangements.
\item
  \textbf{Framework alignment:} Vision \& Change: Information flow,
  exchange, and storage, Evolution; AP Biology: Information Storage and
  Transmission, Evolution; NGSS-style topics: Inheritance and Variation
  of Traits, Natural Selection and Evolution.
\item
  \textbf{Model or quantitative lens:} Recombination frequency and
  three-point mapping.
\item
  \textbf{Data skill:} Infer gene order or chromosomal mechanism from
  offspring counts.
\item
  \textbf{Practice cadence:} Statistical Tests and Data Analysis,
  Representing and Describing Data.
\item
  \textbf{Common misconception to repair:} Independent assortment
  applies to unlinked loci, not to every pair of genes.
\item
  \textbf{Primary lab:} \cref{sec:lab_unit_V_chromosomal_inheritance}.
\item
  \textbf{Question bank:} \cref{sec:q_unit_V_chromosomal_inheritance}.
\item
  \textbf{Transfer task:} Transfer linkage reasoning to disease mapping,
  breeding, and genome assemblies.
\item
  \textbf{Bridge to computation:}
  \breaktt{biology.genetics.genetics.infer\_three\_point\_order}.
\end{itemize}

\begin{center}\rule{0.5\linewidth}{0.5pt}\end{center}

\begin{quote}
\textbf{Opening Vignette --- A Fly's White Eye Opens the Chromosome Era}

In 1910, Thomas Hunt Morgan was cultivating thousands of
\emph{Drosophila melanogaster} in a ``fly room'' at Columbia University
when he noticed a single male with white eyes among thousands of
red-eyed flies. When he crossed this mutant with normal red-eyed females
and counted the offspring across two generations, a striking pattern
emerged: white eyes appeared almost exclusively in male offspring.
Morgan immediately recognized that the white-eye trait was linked to the
X chromosome --- the first gene mapped to a specific chromosome in any
organism. The discovery of sex-linkage validated the chromosome theory
of heredity (which Mendel's work alone could not prove), won Morgan the
1933 Nobel Prize, and launched the entire field of chromosome mapping.
The humble fruit fly, with its four chromosome pairs, short generation
time, and thousands of offspring, remains one of the most productive
research organisms in genetics to this day.
\end{quote}

\subsection{The Chromosome Theory of
Heredity}\label{the-chromosome-theory-of-heredity}

In 1902, Walter Sutton (studying grasshopper \emph{Brachystola magna})
and Theodor Boveri (studying sea urchin embryos) independently proposed
that chromosomes are the physical carriers of Mendel's ``factors''
(genes):

{\def\LTcaptype{none} % do not increment counter
\begin{longtable}[]{@{}
  >{\raggedright\arraybackslash}p{(\linewidth - 2\tabcolsep) * \real{0.5000}}
  >{\raggedright\arraybackslash}p{(\linewidth - 2\tabcolsep) * \real{0.5000}}@{}}
\toprule\noalign{}
\begin{minipage}[b]{\linewidth}\raggedright
Mendel's Observation
\end{minipage} & \begin{minipage}[b]{\linewidth}\raggedright
Chromosomal Parallel
\end{minipage} \\
\midrule\noalign{}
\endhead
\bottomrule\noalign{}
\endlastfoot
Genes exist in pairs in somatic cells & Chromosomes exist in homologous
pairs (diploid) \\
\hyperref[gl:allele]{\textbf{Allele}}s segregate during
\hyperref[gl:gamete]{\textbf{gamete}} formation & Homologous chromosomes
separate at \hyperref[gl:meiosis]{\textbf{meiosis}} I \\
Different genes assort independently & Non-homologous chromosomes orient
independently at metaphase I \\
Gametes carry one allele per gene & Gametes are haploid (one chromosome
from each pair) \\
\end{longtable}
}

\textbf{Morgan's confirmation (1910-1915)}: Thomas Hunt Morgan's
experiments with \emph{Drosophila melanogaster} provided the definitive
proof. His discovery of \textbf{white-eyed males} (X-linked recessive)
demonstrated that a specific gene (white, \emph{w}) segregated with the
X chromosome, proving that genes reside on chromosomes. Morgan received
the Nobel Prize in Physiology or Medicine in 1933.

\begin{center}\rule{0.5\linewidth}{0.5pt}\end{center}

\subsection{Chromosome Architecture: The Hardware of
Heredity}\label{chromosome-architecture-the-hardware-of-heredity}

A eukaryotic chromosome is far more than a passive package of DNA --- it
is a structurally organized object whose every architectural feature has
a functional role in transmission, expression, and protection of the
genetic material. Understanding chromosome behavior requires
understanding chromosome \textbf{anatomy}: the metaphase chromosome in
\cref{fig:unit_V_chromosome_structure} shows how the centromere,
telomeres, and heterochromatic and euchromatic domains discussed below
are arranged along the sister chromatids.

\subsubsection{Centromeres}\label{centromeres}

The \textbf{centromere} is the chromosomal region where sister
chromatids are held together after DNA replication and where the
\textbf{kinetochore} assembles to attach to spindle microtubules during
cell division. Centromeres position chromosomes correctly at metaphase
and are essential for the orderly disjunction of chromatids at anaphase.

\begin{itemize}
\tightlist
\item
  \textbf{Position-based classification}: Metacentric (centromere
  central --- chromosome 1), submetacentric (slightly off-center ---
  chromosome 6), acrocentric (near one end --- chromosome 21),
  telocentric (terminal --- present in some species but not in normal
  humans).
\item
  \textbf{Sequence composition}: Human centromeres contain
  \textbf{alpha-satellite DNA} --- tandem repeats of a 171-bp monomer
  organized into higher-order arrays spanning 0.3--5 Mb. The DNA
  sequence is rapidly evolving, but centromere identity is
  \textbf{epigenetically defined} by the histone variant
  \textbf{CENP-A}, which replaces canonical H3 in centromeric
  nucleosomes and recruits the kinetochore.
\item
  \textbf{Robertsonian translocations} involve fusion of two acrocentric
  chromosomes through their centromeres --- a phenomenon that depends on
  the unique architecture of acrocentric short arms and is revisited in
  the structural-rearrangements discussion.
\end{itemize}

\subsubsection{Telomeres}\label{telomeres}

The ends of linear eukaryotic chromosomes pose a structural problem: the
cell must distinguish a true chromosome end from a double-strand break
(which would otherwise trigger DNA damage signaling). \textbf{Telomeres}
solve this with specialized terminal structures.

\begin{itemize}
\tightlist
\item
  \textbf{Sequence}: Human telomeres consist of tandem repeats of
  \textbf{TTAGGG} (5--15 kb at birth), bound by the \textbf{shelterin
  complex} (TRF1, TRF2, POT1, TIN2, TPP1, RAP1) that masks the
  chromosome end and inhibits inappropriate repair responses.
\item
  \textbf{The end-replication problem}: DNA polymerase cannot fully
  replicate the lagging strand, so chromosomes shorten by 50--200 bp per
  division. Without compensation, cells become senescent after
  about 50 divisions (the \textbf{Hayflick limit}).
\item
  \textbf{Telomerase}: A specialized reverse transcriptase (catalytic
  subunit TERT, RNA template TERC) extends telomeres in stem cells, germ
  cells, and about 85\% of cancers --- making telomerase a
  cancer therapeutic target.
\item
  \textbf{Clinical relevance}: \textbf{Dyskeratosis congenita} results
  from telomerase mutations and causes premature aging, bone marrow
  failure, and cancer. \textbf{Werner syndrome} (helicase mutations)
  accelerates telomere shortening.
\end{itemize}

\subsubsection{Euchromatin and
Heterochromatin}\label{euchromatin-and-heterochromatin}

Chromatin is not uniform --- it exists in two functionally and
structurally distinct states.

{\def\LTcaptype{none} % do not increment counter
\begin{longtable}[]{@{}
  >{\raggedright\arraybackslash}p{(\linewidth - 4\tabcolsep) * \real{0.2250}}
  >{\raggedright\arraybackslash}p{(\linewidth - 4\tabcolsep) * \real{0.3250}}
  >{\raggedright\arraybackslash}p{(\linewidth - 4\tabcolsep) * \real{0.4500}}@{}}
\toprule\noalign{}
\begin{minipage}[b]{\linewidth}\raggedright
Feature
\end{minipage} & \begin{minipage}[b]{\linewidth}\raggedright
Euchromatin
\end{minipage} & \begin{minipage}[b]{\linewidth}\raggedright
Heterochromatin
\end{minipage} \\
\midrule\noalign{}
\endhead
\bottomrule\noalign{}
\endlastfoot
Compaction & Open, accessible & Densely packed \\
Histone marks & H3K4me3, H3K27ac, H3K36me3 & H3K9me3, H3K27me3,
H4K20me3 \\
Replication timing & Early S phase & Late S phase \\
GC content & Higher & Lower \\
Recombination & Frequent & Suppressed \\
Transcription & Active genes & Silenced/repetitive \\
Examples & Most gene-rich regions & Centromeres, telomeres, inactive X
(Barr body) \\
\end{longtable}
}

\begin{itemize}
\tightlist
\item
  \textbf{Constitutive heterochromatin} is permanent and structural ---
  centromeric and pericentric regions, telomeres, the Y chromosome
  long-arm heterochromatin, and tandemly repeated satellite DNAs.
  Constitutive heterochromatin maintains chromosome architecture and
  silences transposable elements.
\item
  \textbf{Facultative heterochromatin} is \textbf{inducible} --- formed
  in specific cell types or developmental states. The classic example is
  the \textbf{inactive X chromosome (Barr body)} in female mammals,
  formed \emph{de novo} in early embryos through XIST-mediated silencing
  as developed in the X-inactivation section.
\item
  \textbf{Position-effect variegation (PEV)}: When a gene is relocated
  by translocation or transposition into a heterochromatic region, it
  can be silenced --- but with stochastic, mosaic patterns that depend
  on heterochromatin spreading. Classic \emph{Drosophila} PEV
  experiments (the \emph{white\^{}m4} allele) revealed how
  heterochromatin propagates and provided the first hints of
  histone-based silencing.
\end{itemize}

\begin{figure}[htbp]
\centering
\includegraphics[width=0.85\textwidth]{../figures/chromosome_structure.png}
\caption{Anatomy of a metaphase chromosome: two sister chromatids joined at the centromere (containing CENP-A nucleosomes that recruit the kinetochore), capped by telomeres of TTAGGG repeats bound by the shelterin complex, with constitutive heterochromatin flanking the centromere and facultative heterochromatin at silenced loci such as the inactive X.}
\label{fig:unit_V_chromosome_structure}
\end{figure}

\begin{center}\rule{0.5\linewidth}{0.5pt}\end{center}

\subsection{Meiosis vs.~Mitosis: Two Cell-Division
Programs}\label{meiosis-vs.-mitosis-two-cell-division-programs}

The chromosomal theory of inheritance rests on a single comparative
observation --- that \textbf{meiosis differs systematically from
mitosis} in ways that exactly produce Mendelian segregation. The two
programs share machinery (spindle, cohesin, kinetochores) but diverge in
chromosome behavior.

{\def\LTcaptype{none} % do not increment counter
\begin{longtable}[]{@{}
  >{\raggedright\arraybackslash}p{(\linewidth - 4\tabcolsep) * \real{0.3333}}
  >{\raggedright\arraybackslash}p{(\linewidth - 4\tabcolsep) * \real{0.3333}}
  >{\raggedright\arraybackslash}p{(\linewidth - 4\tabcolsep) * \real{0.3333}}@{}}
\toprule\noalign{}
\begin{minipage}[b]{\linewidth}\raggedright
Feature
\end{minipage} & \begin{minipage}[b]{\linewidth}\raggedright
Mitosis
\end{minipage} & \begin{minipage}[b]{\linewidth}\raggedright
Meiosis
\end{minipage} \\
\midrule\noalign{}
\endhead
\bottomrule\noalign{}
\endlastfoot
\textbf{Number of divisions} & One (M phase) & Two (Meiosis I, Meiosis
II) \\
\textbf{DNA replication} & Once before division & Once before Meiosis I;
none between MI and MII \\
\textbf{Daughter cells per parent} & 2 & 4 \\
\textbf{Ploidy change} & Diploid \texorpdfstring{\ensuremath{\to}}{->} diploid (2n \texorpdfstring{\ensuremath{\to}}{->} 2n) & Diploid \texorpdfstring{\ensuremath{\to}}{->} haploid
(2n \texorpdfstring{\ensuremath{\to}}{->} n) \\
\textbf{Homolog pairing} & None & Synapsis with synaptonemal complex at
prophase I \\
\textbf{Crossing over (recombination)} & None (rare exceptions) &
Required at prophase I; about 1.5--3 crossovers per bivalent in
humans \\
\textbf{Anaphase I event} & N/A & Homologs disjoin (sisters remain
joined) \\
\textbf{Anaphase II event} & Sisters disjoin & Sisters disjoin \\
\textbf{Genetic outcome} & Identical genetic copies & Genetically unique
haploid gametes \\
\textbf{Cohesin protection} & None at centromere; cleaved at anaphase &
\textbf{Shugoshin} protects centromeric cohesin until anaphase II \\
\textbf{Cellular function} & Growth, repair, asexual reproduction &
Sexual reproduction; gamete formation \\
\textbf{Failure consequences} & Tumor formation, aneuploid somatic cells
& Aneuploid gametes (Down syndrome, miscarriage) \\
\end{longtable}
}

\begin{figure}[htbp]
\centering
\includegraphics[width=0.82\linewidth,height=4.2in,keepaspectratio]{../figures/mermaid_inline/inline_mermaid_0092_a8a19795ee09.png}
\caption{Mitosis vs. meiosis: a single division producing two diploid copies versus a two-division program producing four haploid, recombined gametes. The defining event of meiosis is anaphase I, when homologs (not sisters) separate.}
\end{figure}

The combination of (1) homolog pairing in prophase I, (2) crossing over,
(3) random orientation of bivalents at metaphase I, and (4) reductional
division at anaphase I produces gametes that are \textbf{genetically
unique} and that satisfy Mendel's laws. Each of these four events can
fail --- and each failure has a recognizable disease signature.

\subsubsection{Molecular details of meiotic
recombination}\label{molecular-details-of-meiotic-recombination}

The orderly behavior of bivalents at metaphase I depends on
\textbf{physical connections} (chiasmata) created by reciprocal
recombination during prophase I. The molecular events that generate
these chiasmata are conserved across eukaryotes and are essential ---
failure of recombination is a leading cause of meiotic non-disjunction
in human oogenesis.

\begin{figure}[htbp]
\centering
\includegraphics[width=0.82\linewidth,height=4.2in,keepaspectratio]{../figures/mermaid_inline/inline_mermaid_0093_10bbf3532472.png}
\caption{Molecular events of meiotic recombination from SPO11 DSB induction through RAD51/DMC1-mediated strand invasion to chiasma formation. Approximately 10\% of programmed double-strand breaks resolve as crossovers; the remainder become non-crossover gene conversions.}
\end{figure}

\textbf{Key molecular players:}

\begin{itemize}
\tightlist
\item
  \textbf{SPO11} is a topoisomerase-VI-like enzyme that creates
  \textbf{programmed double-strand breaks} by covalent attachment to 5\texorpdfstring{'}{'}
  ends of DNA. In humans, about 250 DSBs are generated per
  meiosis (about 10\texorpdfstring{\ensuremath{\times}}{x} the number of crossovers --- most DSBs are
  repaired without crossover). SPO11 binding sites are determined by
  \textbf{PRDM9}, which trimethylates H3K4 at specific DNA motifs to
  mark recombination hotspots. \emph{Spo11} knockout mice are sterile
  because no chiasmata form, leading to bivalent disjunction failure and
  aneuploid gametes.
\item
  \textbf{MRE11--RAD50--NBS1 (MRN complex)} binds the DSB and resects
  the 5\texorpdfstring{'}{'} ends, exposing 3\texorpdfstring{'}{'} single-stranded overhangs about 1--2
  kb long.
\item
  \textbf{RAD51 and DMC1} are RecA-family recombinases that polymerize
  on the single-stranded DNA, forming a nucleoprotein filament that
  searches for homology in the homologous chromosome and catalyzes
  strand invasion. DMC1 is meiosis-specific and is essential for
  inter-homolog (rather than inter-sister) recombination.
\item
  \textbf{MSH4--MSH5} and \textbf{MLH1--MLH3} are required for the major
  crossover pathway (class I); their mutations reduce crossover
  frequency and increase non-disjunction. \textbf{MUS81--EME1} mediates
  a minor (about 10\%) class II crossover pathway.
\item
  \textbf{Synaptonemal complex}: A protein scaffold (SYCP1, SYCP2,
  SYCP3) that physically holds homologs together along their entire
  length during pachytene. The complex provides the platform on which
  crossovers form and assures proper pairing.
\end{itemize}

The molecular specifics matter because \textbf{mutations in many of
these genes cause human infertility, recurrent pregnancy loss, and
elevated aneuploidy risk}: \emph{SYCP3} mutations are associated with
azoospermia and pregnancy loss; \emph{MSH4} mutations with primary
ovarian insufficiency; \emph{DMC1} mutations with sterility; and reduced
\emph{PRDM9} function correlates with diminished ovarian reserve.
Meiotic recombination is far from a passive structural process --- it is
an actively regulated DNA-damage event without which sexual reproduction
fails.

\begin{center}\rule{0.5\linewidth}{0.5pt}\end{center}

\subsection{Sex Chromosomes and Sex
Determination}\label{sex-chromosomes-and-sex-determination}

\begin{figure}[htbp]
\centering
\includegraphics[width=0.82\linewidth,height=4.2in,keepaspectratio]{../figures/mermaid_inline/inline_mermaid_0094_80df69de47bb.png}
\caption{Sex determination systems across organisms. The XX/XY system is used by mammals, with the SRY gene on the Y chromosome being the master switch. Other organisms use ZW, XO, haplodiploidy, or temperature-dependent systems.}
\end{figure}

\subsubsection{Mammalian Sex Determination
(XX/XY)}\label{mammalian-sex-determination-xxxy}

\textbf{The Y chromosome and SRY}: - Human Y chromosome:
about 57 Mb; carries about 78
\hyperref[gl:protein]{\textbf{protein}}-coding genes (compared to
about 900 on the X) - \textbf{SRY} (sex-determining region Y):
Located at Yp11.31; encodes \textbf{TDF} (testis-determining factor), a
\hyperref[gl:transcription]{\textbf{transcription}} factor with an
HMG-box DNA-binding domain - SRY activates \textbf{SOX9}, which drives
Sertoli cell differentiation in the bipotential gonad - Sertoli cells
produce \textbf{anti-Mullerian \hyperref[gl:hormone]{\textbf{hormone}}
(AMH)}, causing regression of Mullerian ducts - Leydig cells produce
\textbf{testosterone}, driving Wolffian duct development (epididymis,
vas deferens, seminal vesicles) and external male genitalia

\textbf{Pseudoautosomal regions (PARs)}: - \textbf{PAR1} (2.6 Mb at
Xp/Yp tips): contains about 24 genes; mandatory crossover
during male meiosis ensures proper X-Y segregation - \textbf{PAR2} (320
kb at Xq/Yq tips): smaller; crossover not required - Genes in PARs
escape X-inactivation (expressed from both X chromosomes in females) -
\textbf{SHOX gene} (in PAR1): Short stature homeobox; haploinsufficiency
causes short stature in Turner syndrome; extra copies cause tall stature
in sex chromosome polysomies

\textbf{Evidence for SRY as the master switch}: - XX males (de la
Chapelle syndrome): about 1/20,000; caused by translocation of
SRY to the X chromosome - XY females (Swyer syndrome): SRY
\hyperref[gl:mutation]{\textbf{mutation}}s lead to female
\hyperref[gl:phenotype]{\textbf{phenotype}} with streak gonads despite
46,XY karyotype

\subsubsection{Non-Mammalian Sex
Determination}\label{non-mammalian-sex-determination}

{\def\LTcaptype{none} % do not increment counter
\begin{longtable}[]{@{}
  >{\raggedright\arraybackslash}p{(\linewidth - 8\tabcolsep) * \real{0.1860}}
  >{\raggedright\arraybackslash}p{(\linewidth - 8\tabcolsep) * \real{0.2558}}
  >{\raggedright\arraybackslash}p{(\linewidth - 8\tabcolsep) * \real{0.1395}}
  >{\raggedright\arraybackslash}p{(\linewidth - 8\tabcolsep) * \real{0.1860}}
  >{\raggedright\arraybackslash}p{(\linewidth - 8\tabcolsep) * \real{0.2326}}@{}}
\toprule\noalign{}
\begin{minipage}[b]{\linewidth}\raggedright
System
\end{minipage} & \begin{minipage}[b]{\linewidth}\raggedright
Organisms
\end{minipage} & \begin{minipage}[b]{\linewidth}\raggedright
Male
\end{minipage} & \begin{minipage}[b]{\linewidth}\raggedright
Female
\end{minipage} & \begin{minipage}[b]{\linewidth}\raggedright
Key Gene
\end{minipage} \\
\midrule\noalign{}
\endhead
\bottomrule\noalign{}
\endlastfoot
ZW/ZZ & Birds, snakes, butterflies & ZZ & ZW & DMRT1 (Z-linked; dosage
determines sex) \\
XX/XO & Grasshoppers, some nematodes & XO & XX & X:autosome ratio (C.
elegans) \\
Haplodiploidy & Hymenoptera (bees, ants, wasps) & Haploid (n) &
\hyperref[gl:diploid]{\textbf{Diploid (2n)}} & CSD locus (complementary
sex determiner) \\
Temperature-dependent & Crocodilians, some turtles &
Temperature-dependent & Temperature-dependent & Aromatase (converts
testosterone to estradiol) \\
\end{longtable}
}

\begin{center}\rule{0.5\linewidth}{0.5pt}\end{center}

\subsection{X-Inactivation and Dosage
Compensation}\label{x-inactivation-and-dosage-compensation}

\subsubsection{Lyon Hypothesis and
X-Inactivation}\label{lyon-hypothesis-and-x-inactivation}

Mary \citet{lyon1961} proposed that in each female somatic cell, one of
the two X chromosomes is randomly inactivated early in embryonic
development (about  day 5.5 in mouse, about  day 12-16
in human). This inactivation is:

\begin{itemize}
\tightlist
\item
  \textbf{Random}: Either the maternal or paternal X can be inactivated
  in any given cell
\item
  \textbf{Clonal}: Once established, most daughter cells maintain the
  same inactive X
\item
  \textbf{Result}: Adult females are \textbf{genetic mosaics} -- a
  patchwork of cells expressing the maternal X and cells expressing the
  paternal X
\end{itemize}

\begin{figure}[htbp]
\centering
\includegraphics[width=0.82\linewidth,height=4.2in,keepaspectratio]{../figures/mermaid_inline/inline_mermaid_0095_88ba3c6367df.png}
\caption{The molecular mechanism of X-inactivation. XIST RNA coats the future inactive X chromosome, recruiting Polycomb complexes that establish repressive histone marks and DNA methylation, ultimately condensing the chromosome into a Barr body.}
\end{figure}

\subsubsection{Molecular Mechanism}\label{molecular-mechanism}

\begin{enumerate}
\def\labelenumi{\arabic{enumi}.}
\tightlist
\item
  \textbf{XIST} (X-inactive specific transcript): A 17 kb lncRNA
  transcribed primarily from the X that will be inactivated
\item
  XIST RNA \textbf{coats} the Xi chromosome \emph{in cis} -- spreading
  outward from the X-inactivation center (XIC) at Xq13
\item
  XIST recruits \textbf{PRC2} (Polycomb repressive complex 2), which
  catalyzes histone H3K27 trimethylation
\item
  Additional silencing factors: \textbf{SMCHD1} (structural maintenance
  of chromosomes flexible hinge domain-containing 1) further compacts
  \hyperref[gl:chromatin]{\textbf{chromatin}}
\item
  \textbf{DNA methylation} at \hyperref[gl:cpg-island]{\textbf{CpG
  island}}s provides stable, long-term silencing
\item
  The silenced X is visible as a \textbf{Barr body} at the nuclear
  periphery (sex chromatin test)
\end{enumerate}

\textbf{Escape from X-inactivation}: about 15\% of X-linked
genes ``escape'' inactivation and are expressed from both the active and
inactive X. These include genes in the pseudoautosomal regions and some
genes with Y-chromosome homologs. Escape genes explain some phenotypic
features of X-chromosome aneuploidy.

\textbf{Barr body count} = number of X chromosomes minus 1: - 46,XX: 1
Barr body - 47,XXX: 2 Barr bodies - 45,X (Turner): 0 Barr bodies -
47,XXY (Klinefelter): 1 Barr body

\subsubsection{Dosage Compensation Across Species and Beyond Sex
Chromosomes}\label{dosage-compensation-across-species-and-beyond-sex-chromosomes}

{\def\LTcaptype{none} % do not increment counter
\begin{longtable}[]{@{}
  >{\raggedright\arraybackslash}p{(\linewidth - 4\tabcolsep) * \real{0.3448}}
  >{\raggedright\arraybackslash}p{(\linewidth - 4\tabcolsep) * \real{0.3793}}
  >{\raggedright\arraybackslash}p{(\linewidth - 4\tabcolsep) * \real{0.2759}}@{}}
\toprule\noalign{}
\begin{minipage}[b]{\linewidth}\raggedright
Organism
\end{minipage} & \begin{minipage}[b]{\linewidth}\raggedright
Mechanism
\end{minipage} & \begin{minipage}[b]{\linewidth}\raggedright
Result
\end{minipage} \\
\midrule\noalign{}
\endhead
\bottomrule\noalign{}
\endlastfoot
\textbf{Mammals} & X-inactivation (XIST) & One X silenced in XX females;
equal dosage with XY males \\
\textbf{Drosophila} & X upregulation in males (MSL/roX RNA complex) &
Male single X transcribed about 2x; equals female XX \\
\textbf{C. elegans} & X downregulation in XX hermaphrodites (DCC
complex) & Each X transcribed about 0.5x; equals male XO single
X \\
\end{longtable}
}

\textbf{Dosage compensation beyond X-inactivation}: The principle that
gene dosage must be balanced applies far more broadly than the sex
chromosomes. Several other mechanisms address dosage imbalances:

\begin{itemize}
\tightlist
\item
  \textbf{Autosomal dosage sensitivity}: Most genes tolerate small
  dosage changes (1.5x or 2x) without phenotype, but
  \textbf{dosage-sensitive genes} --- typically transcription factors,
  signaling components, and members of multi-protein complexes ---
  produce phenotypes when copy number changes. The \textbf{Williams
  syndrome} deletion (7q11.23) and its reciprocal duplication produce
  opposite cognitive-behavioral phenotypes precisely because of dosage
  sensitivity in the \emph{GTF2I} gene region.
\item
  \textbf{Gene balance hypothesis}: Genes whose products participate in
  stoichiometric complexes (e.g., ribosomal proteins, transcription
  complexes, RNA polymerases) cannot tolerate copy-number changes
  because subunit imbalance disrupts complex assembly. This explains why
  most autosomal monosomies are lethal even when the underlying gene
  mutations would not be.
\item
  \textbf{Compensatory feedback at promoters}: In \emph{Drosophila} and
  yeast, deletion of one allele of a dosage-sensitive gene often
  triggers transcriptional upregulation of the remaining allele,
  partially restoring protein levels (cis-compensation).
\item
  \textbf{microRNA buffering}: Individual miRNAs can buffer the
  expression of dozens of target genes against perturbation. The
  miR-17about 92 cluster, for example, buffers cell-cycle gene
  dosage during proliferation.
\item
  \textbf{Polycomb-mediated maintenance of dosage}: Active and inactive
  states established at gene clusters (Hox genes, the imprinted regions
  of chromosome 15) are maintained through cell division by Polycomb
  (silencing) and Trithorax (activation) complexes, ensuring stable gene
  dosage in each tissue.
\end{itemize}

The core insight is that \textbf{gene dosage is itself a regulated
variable}, defended by multiple parallel mechanisms. X-inactivation is
the most spectacular example because an entire chromosome is silenced,
but it is one strategy among many.

\begin{quote}
\textbf{Clinical Connection: Calico Cats and X-Inactivation} The calico
coat pattern in cats is a visible demonstration of X-inactivation. The
\emph{O} gene (orange) is X-linked: \(X^O\) produces orange pigment;
\(X^o\) produces black pigment.
\hyperref[gl:heterozygous]{\textbf{Heterozygous}} females (\(X^OX^o\))
are mosaics -- patches of orange (cells with \(X^o\) inactivated) and
black (cells with \(X^O\) inactivated). White patches come from a
separate autosomal spotting gene. Nearly most calico cats are female;
rare calico males are usually 47,XXY (Klinefelter).
\end{quote}

\begin{center}\rule{0.5\linewidth}{0.5pt}\end{center}

\subsection{X-Linked Inheritance}\label{x-linked-inheritance}

\subsubsection{X-Linked Recessive
Inheritance}\label{x-linked-recessive-inheritance}

For X-linked recessive traits, males are affected more frequently
because they are hemizygous (a single X):

\begin{itemize}
\tightlist
\item
  \textbf{Affected males}: \(X^aY\) -- one copy of the recessive allele
  is sufficient
\item
  \textbf{Affected females}: \(X^aX^a\) -- must be homozygous (very
  rare; requires carrier mother AND affected father)
\item
  \textbf{Carrier females}: \(X^AX^a\) -- phenotypically normal
  (usually); can show mosaic expression due to random X-inactivation
\end{itemize}

\textbf{Key X-linked recessive conditions:}

{\def\LTcaptype{none} % do not increment counter
\begin{longtable}[]{@{}
  >{\raggedright\arraybackslash}p{(\linewidth - 8\tabcolsep) * \real{0.2157}}
  >{\raggedright\arraybackslash}p{(\linewidth - 8\tabcolsep) * \real{0.1176}}
  >{\raggedright\arraybackslash}p{(\linewidth - 8\tabcolsep) * \real{0.1961}}
  >{\raggedright\arraybackslash}p{(\linewidth - 8\tabcolsep) * \real{0.2157}}
  >{\raggedright\arraybackslash}p{(\linewidth - 8\tabcolsep) * \real{0.2549}}@{}}
\toprule\noalign{}
\begin{minipage}[b]{\linewidth}\raggedright
Condition
\end{minipage} & \begin{minipage}[b]{\linewidth}\raggedright
Gene
\end{minipage} & \begin{minipage}[b]{\linewidth}\raggedright
Location
\end{minipage} & \begin{minipage}[b]{\linewidth}\raggedright
Frequency
\end{minipage} & \begin{minipage}[b]{\linewidth}\raggedright
Key Features
\end{minipage} \\
\midrule\noalign{}
\endhead
\bottomrule\noalign{}
\endlastfoot
\textbf{Hemophilia A} & \emph{F8} (Factor VIII) & Xq28 & 1/5,000 males &
Prolonged bleeding; joint hemorrhages; treated with recombinant Factor
VIII \\
\textbf{Hemophilia B} & \emph{F9} (Factor IX) & Xq27.1 & 1/30,000 males
& Similar to hemophilia A; ``Christmas disease'' \\
\textbf{Duchenne muscular dystrophy} & \emph{DMD} (dystrophin) & Xp21.2
& 1/3,500 males & Largest human gene (2.4 Mb); progressive muscle
wasting; frameshift mutations cause complete dystrophin loss; onset age
2-5; wheelchair by age 12; cardiac/respiratory failure \\
\textbf{Becker muscular dystrophy} & \emph{DMD} (dystrophin) & Xp21.2 &
1/18,000 males & In-frame deletions; partially functional dystrophin;
milder course \\
\textbf{Red-green color blindness} & \emph{OPN1LW/OPN1MW} & Xq28 & 8\%
of males, 0.5\% of females & Deutan (green) or protan (red) deficiency;
unequal crossing over between tandem opsin genes \\
\textbf{G6PD deficiency} & \emph{G6PD} & Xq28 & about 400
million affected worldwide & Hemolytic anemia triggered by oxidative
stress (fava beans, certain drugs); heterozygote advantage against
malaria \\
\end{longtable}
}

\subsection{Worked Example: X-Linked Recessive
Inheritance}\label{worked-example-x-linked-recessive-inheritance}

A carrier woman (\(X^HX^h\) for hemophilia) marries an unaffected man
(\(X^HY\)):

{\def\LTcaptype{none} % do not increment counter
\begin{longtable}[]{@{}lll@{}}
\toprule\noalign{}
& \(X^H\) & \(Y\) \\
\midrule\noalign{}
\endhead
\bottomrule\noalign{}
\endlastfoot
\(X^H\) & \(X^HX^H\) (normal female) & \(X^HY\) (normal male) \\
\(X^h\) & \(X^HX^h\) (carrier female) & \(X^hY\) (affected male) \\
\end{longtable}
}

Results: \(\frac{1}{4}\) normal daughters, \(\frac{1}{4}\) carrier
daughters, \(\frac{1}{4}\) normal sons, \(\frac{1}{4}\) affected sons.

Among sons primarily: \(\frac{1}{2}\) affected. Among daughters
primarily: \(\frac{1}{2}\) carriers.

\begin{quote}
\textbf{Clinical Connection: Duchenne Muscular Dystrophy and
\hyperref[gl:exon]{\textbf{Exon}} Skipping} The \emph{DMD} gene (2.4 Mb,
79 exons) is a frequent target of deletions. Out-of-frame deletions
cause Duchenne (no dystrophin); in-frame deletions cause Becker (partial
dystrophin). \textbf{Exon-skipping therapy} uses antisense
oligonucleotides (ASOs) to mask specific exons during splicing,
converting an out-of-frame deletion to an in-frame deletion (Duchenne to
Becker phenotype). \textbf{Eteplirsen} (FDA-approved 2016) targets exon
51; applicable to about 13\% of DMD patients. Newer ASOs target
other exons.
\end{quote}

\textbf{Concept Check 17.5}

\begin{quote}
\begin{enumerate}
\def\labelenumi{\arabic{enumi}.}
\tightlist
\item
  A pedigree shows an affected trait that appears primarily in males, is
  rarely transmitted father-to-son, and reaches grandsons through
  phenotypically normal daughters of affected men: explain why these
  three features together identify X-linked recessive inheritance, and
  state which single observation would instead force an
  autosomal-recessive interpretation.
\end{enumerate}
\end{quote}

\begin{center}\rule{0.5\linewidth}{0.5pt}\end{center}

\subsection{Linkage and Recombination}\label{linkage-and-recombination}

\subsubsection{Morgan's Experiments}\label{morgans-experiments}

Thomas Hunt Morgan (1910-1915) discovered that some genes do NOT assort
independently -- they are \textbf{linked} on the same chromosome.

In a testcross for body color (b: black recessive) and wing type (vg:
vestigial recessive) in \emph{Drosophila}:

\textbf{Expected (if unlinked)}: 1:1:1:1 ratio of four phenotype classes
\textbf{Observed}: Excess \textbf{parental types} (gray-normal,
black-vestigial) and deficiency of \textbf{recombinant types}
(gray-vestigial, black-normal)

This means b and vg are on the same chromosome, and \textbf{crossing
over} during meiosis generates recombinants at a frequency proportional
to the physical distance between the genes.

\subsubsection{Map Distance and Recombination
Frequency}\label{map-distance-and-recombination-frequency}

\textbf{Recombination frequency (RF)}:

\begin{equation}
RF = \frac{\text{number of recombinant offspring}}{\text{total offspring}} \times 100\%
\label{eq:unit_V_chromosomal_rf_definition}
\end{equation}

\textbf{Map distance} in centimorgans (cM): 1 cM = 1\% recombination
frequency.

\begin{itemize}
\tightlist
\item
  RF \textless{} 50\%: genes are \textbf{linked} (on the same
  chromosome)
\item
  RF = 50\%: genes are \textbf{unlinked} (on different chromosomes OR so
  far apart on the same chromosome that at least one crossover typically
  occurs between them)
\item
  RF can rarely exceed 50\% because a single crossover between two loci
  produces 50\% recombinant and 50\% parental gametes (two chromatids
  involved out of four)
\end{itemize}

\subsubsection{Three-Point Test Cross}\label{three-point-test-cross}

The three-point test cross determines the order of three linked genes
and the distances between them simultaneously. It is more efficient than
three separate two-point crosses and reveals \textbf{interference}.

\subsection{\texorpdfstring{Worked Example: 17.2: In \emph{Drosophila},
three linked genes on chromosome
2:}{Worked Example: 17.2: In Drosophila, three linked genes on chromosome 2:}}\label{worked-example-17.2-in-drosophila-three-linked-genes-on-chromosome-2}

\begin{itemize}
\tightlist
\item
  dp (dumpy wings): recessive
\item
  b (black body): recessive
\item
  cn (cinnabar eyes): recessive
\end{itemize}

A triply heterozygous female (dp+ b+ cn+ / dp b cn) is testcrossed to a
triply homozygous recessive male (dp b cn / dp b cn). The offspring are:

{\def\LTcaptype{none} % do not increment counter
\begin{longtable}[]{@{}lll@{}}
\toprule\noalign{}
Phenotype class & Count & Classification \\
\midrule\noalign{}
\endhead
\bottomrule\noalign{}
\endlastfoot
dp+ b+ cn+ & 350 & Parental \\
dp b cn & 345 & Parental \\
dp+ b cn & 62 & Single crossover (region II) \\
dp b+ cn+ & 58 & Single crossover (region II) \\
dp+ b+ cn & 40 & Single crossover (region I) \\
dp b cn+ & 38 & Single crossover (region I) \\
dp+ b cn+ & 5 & Double crossover \\
dp b+ cn & 2 & Double crossover \\
\textbf{Total} & \textbf{900} & \\
\end{longtable}
}

\textbf{Step 1: Determine gene order.} Compare the parental classes
(most frequent) to the double crossover classes (least frequent). The
gene that switches position in the double crossover class is the
\textbf{middle gene}.

\begin{itemize}
\tightlist
\item
  Parental: dp+ b+ cn+; dp b cn
\item
  Double crossover: dp+ b cn+; dp b+ cn
\end{itemize}

The b gene has switched relative to dp and cn. Therefore, the gene order
is: \textbf{dp -- b -- cn} (b is in the middle).

\textbf{Step 2: Calculate map distances.}

\[\text{Distance dp-b} = \frac{40 + 38 + 5 + 2}{900} \times 100 = \frac{85}{900} \times 100 = 9.4 \text{ cM} \tag{17.1}\]

\[\text{Distance b-cn} = \frac{62 + 58 + 5 + 2}{900} \times 100 = \frac{127}{900} \times 100 = 14.1 \text{ cM} \tag{17.2}\]

\[\text{Total map distance dp-cn} = 9.4 + 14.1 = 23.5 \text{ cM} \tag{17.3}\]

\textbf{Step 3: Calculate interference.}

\[\text{Expected double crossovers} = 0.094 \times 0.141 \times 900 = 11.9 \tag{17.4}\]

\[\text{Observed double crossovers} = 5 + 2 = 7 \tag{17.5}\]

\[\text{Coefficient of coincidence (CoC)} = \frac{\text{observed}}{\text{expected}} = \frac{7}{11.9} = 0.59 \tag{17.6}\]

\[\text{Interference (I)} = 1 - \text{CoC} = 1 - 0.59 = 0.41 \tag{17.7}\]

An interference of 0.41 means that 41\% of expected double crossovers
did not occur -- a crossover in one region partially suppresses
crossovers in the adjacent region (\textbf{positive interference}).

\subsubsection{Mapping Functions}\label{mapping-functions}

Simple RF underestimates true genetic distance because multiple
crossovers between distant loci are undetectable (double crossovers
restore parental configuration).

\textbf{Haldane mapping function} (1919, assumes no chromatid
interference):

\begin{equation}
m = -\frac{1}{2}\ln(1 - 2\theta) \quad \text{(morgans)}
\label{eq:unit_V_chromosomal_haldane}
\end{equation}

where \texorpdfstring{\ensuremath{\theta}}{theta} = observed RF. As \(\theta \to 0.5\), \(m \to \infty\).

\textbf{Kosambi mapping function} (1944, accounts for positive chromatid
interference):

\begin{equation}
m = \frac{1}{4}\ln\left(\frac{1 + 2\theta}{1 - 2\theta}\right)
\label{eq:unit_V_chromosomal_kosambi}
\end{equation}

Kosambi's function better fits empirical data for moderate distances.

\subsubsection{LOD Score Analysis}\label{lod-score-analysis}

\textbf{LOD (logarithm of odds) score}: Statistical method for testing
linkage in human pedigrees (where controlled crosses are ethically
unfeasible).

\begin{equation}
LOD = \log_{10}\frac{L(\theta)}{L(\theta = 0.5)}
\label{eq:unit_V_chromosomal_lod}
\end{equation}

Where \(L(\theta)\) is the likelihood of the data given recombination
fraction \texorpdfstring{\ensuremath{\theta}}{theta}, and \(L(0.5)\) is the likelihood under no linkage.

\begin{itemize}
\tightlist
\item
  \textbf{LOD \textgreater= 3}: Evidence for linkage (\texorpdfstring{\ensuremath{\theta}}{theta} value estimates
  map distance)
\item
  \textbf{LOD \textless= -2}: Evidence against linkage
\item
  \textbf{-2 \textless{} LOD \textless{} 3}: Inconclusive
\end{itemize}

\subsubsection{Recombination Hotspots and Sex Differences in Map
Length}\label{recombination-hotspots-and-sex-differences-in-map-length}

Recombination is \textbf{not distributed uniformly} across the genome
--- it concentrates in narrow \textbf{hotspots} of about 1--2
kb separated by larger cold regions, and the total map length differs
sharply between male and female meiosis.

\begin{itemize}
\tightlist
\item
  \textbf{PRDM9 and hotspots}: In humans and mice, the zinc-finger
  histone methyltransferase \textbf{PRDM9} binds specific DNA motifs and
  trimethylates H3K4 and H3K36 at the binding site, recruiting the
  meiotic double-strand break machinery (SPO11). Crossovers
  preferentially form at these PRDM9-marked sites. PRDM9 is itself
  rapidly evolving --- its zinc-finger array changes between populations
  and species, and \emph{Prdm9} knockout mice are sterile due to meiotic
  defects. \textbf{PRDM9 is one of very few genes implicated in
  mammalian speciation through hybrid sterility.}
\item
  \textbf{Hotspot rates}: Within hotspots, recombination rates can reach
  10--100\texorpdfstring{\ensuremath{\times}}{x} the genome-wide average. The \textbf{HLA region} on
  chromosome 6 contains some of the most active hotspots, contributing
  to the immense allelic diversity of MHC genes.
\item
  \textbf{Cold regions}: Centromeric and pericentric heterochromatin is
  recombination-suppressed, both because of structural constraints and
  because crossovers near centromeres can disrupt segregation.
  Recombination is also reduced near telomeres (in some species) and in
  regions of strong purifying selection.
\item
  \textbf{Sex differences in map length}: The female human genetic map
  is roughly \textbf{1.5--1.8\texorpdfstring{\ensuremath{\times}}{x} longer} than the male map
  (about 4,400 cM in females vs.~about 2,700 cM in
  males). This reflects more crossovers per meiosis in oogenesis
  (about 70 per cell) than spermatogenesis (about 50
  per cell), and dramatic regional differences --- recombination is
  enriched at telomeres in male meiosis (where obligate crossovers occur
  to ensure disjunction) and more uniformly distributed in female
  meiosis.
\item
  \textbf{Consequences}: Sex-specific maps must be used in linkage
  analysis when transmissions go through one sex predominantly. Large
  family genetic studies routinely report male, female, and sex-averaged
  maps separately.
\end{itemize}

\paragraph{Why do females have more
crossovers?}\label{why-do-females-have-more-crossovers}

The pronounced sex difference in recombination rate has several proposed
explanations, and most empirical studies suggest more than one is
operating simultaneously:

\begin{enumerate}
\def\labelenumi{\arabic{enumi}.}
\tightlist
\item
  \textbf{Oocyte time available}: Female meiosis includes a prolonged
  pachytene during which homolog pairing and crossover designation
  occur. Spermatogenesis is faster and may impose stricter checkpoints
  that limit crossover-pathway entry.
\item
  \textbf{Telomeric vs.~interstitial recombination}: Male meiosis
  enriches crossovers at telomeres because the \textbf{obligate
  crossover} rule (every bivalent must have at least one chiasma to
  segregate properly) is satisfied near chromosome ends, where pairing
  can complete in less time. Female meiosis has the time to distribute
  crossovers more uniformly, including in interstitial regions distant
  from telomeres.
\item
  \textbf{Selection for accuracy in oogenesis}: Because errors during
  the long oocyte arrest are catastrophic, female meiosis may have
  evolved higher recombination rates to ensure \textbf{at least one
  chiasma per bivalent}. Achiasmate bivalents segregate randomly,
  producing aneuploid gametes --- the female reproductive system
  tolerates this poorly because each oocyte represents irreplaceable
  investment.
\item
  \textbf{Sexually antagonistic selection on chromosome conformation}:
  Some authors propose that male and female germlines have different
  optima for genome stability, with selection acting in opposite
  directions on regulators of recombination.
\end{enumerate}

\paragraph{Haldane's rule and
recombination}\label{haldanes-rule-and-recombination}

Haldane's classic 1922 observation extends beyond sex-determination
biology to recombination patterns: in many species the
\textbf{homogametic sex} (XX females in mammals; ZZ males in birds)
shows higher recombination rates than the heterogametic sex. In
\emph{Drosophila}, the asymmetry is extreme --- males have \textbf{no
meiotic recombination at most}, while females have normal
crossover-mediated recombination. In humans, the pattern is muted but
present, with females having about 1.6\texorpdfstring{\ensuremath{\times}}{x} the male map length.
The \textbf{achiasmate male meiosis} of \emph{Drosophila} and several
other dipterans is supported by alternative bivalent-tethering
mechanisms (the \emph{Mnm} complex) that maintain homolog pairing
without crossovers. The mammalian intermediate pattern reflects partial
achiasmatism --- males still recombine, but with reduced rate and biased
localization to telomeres. Why these sex differences evolved remains an
open question and is one of the most robust empirical generalizations in
the genetics of meiosis.

\textbf{Concept Check 17.1}

\begin{quote}
\begin{enumerate}
\def\labelenumi{\arabic{enumi}.}
\tightlist
\item
  In a three-point cross, why are double crossover classes typically the
  least frequent?
\item
  If genes A and B have RF = 45\%, are they linked? Explain.
\item
  What does positive interference tell you about the biology of crossing
  over?
\item
  PRDM9 binds specific DNA motifs to direct recombination hotspots. What
  does its rapid evolution suggest about hotspot evolution between human
  populations?
\end{enumerate}
\end{quote}

\begin{center}\rule{0.5\linewidth}{0.5pt}\end{center}

\subsection{Chromosomal Abnormalities}\label{chromosomal-abnormalities}

\begin{figure}[htbp]
\centering
\includegraphics[width=0.82\linewidth,height=4.2in,keepaspectratio]{../figures/mermaid_inline/inline_mermaid_0096_2ae4e4cd7ccc.png}
\caption{Classification of chromosomal abnormalities. Numerical abnormalities include aneuploidy and polyploidy; structural abnormalities include deletions, duplications, inversions, and translocations.}
\end{figure}

\subsubsection{Aneuploidy: Non-Disjunction in Meiosis I vs.~Meiosis
II}\label{aneuploidy-non-disjunction-in-meiosis-i-vs.-meiosis-ii}

\textbf{Non-disjunction} is the failure of chromosomes to separate
properly during cell division and is the principal source of human
aneuploidy \citep{hassold2001}. It can occur in either meiosis I or
meiosis II, and the two failure modes have \textbf{distinguishable
genetic signatures}.

\begin{figure}[htbp]
\centering
\includegraphics[width=0.82\linewidth,height=4.2in,keepaspectratio]{../figures/mermaid_inline/inline_mermaid_0097_8b4bc56a0e9b.png}
\caption{Two distinct molecular failure modes of meiotic non-disjunction. MI errors produce trisomic gametes carrying both parental homologs; MII errors produce trisomic gametes carrying two sister copies of a single homolog. Centromere-region marker analysis distinguishes the two - clinically important because they have different mechanisms and risk factors.}
\end{figure}

This MI/MII distinction is clinically important because: -
\textbf{about 75\% of trisomy 21 cases originate in maternal
meiosis I} (the cohesin-decay mechanism described below). - Paternal
trisomy 21 (about 5--10\% of cases) is typically meiosis II in
origin. - Trisomy 18 has a higher proportion of MII errors than trisomy
21, suggesting different mechanistic pressures.

\textbf{Risk factors for non-disjunction}:

{\def\LTcaptype{none} % do not increment counter
\begin{longtable}[]{@{}
  >{\raggedright\arraybackslash}p{(\linewidth - 4\tabcolsep) * \real{0.2963}}
  >{\raggedright\arraybackslash}p{(\linewidth - 4\tabcolsep) * \real{0.2963}}
  >{\raggedright\arraybackslash}p{(\linewidth - 4\tabcolsep) * \real{0.4074}}@{}}
\toprule\noalign{}
\begin{minipage}[b]{\linewidth}\raggedright
Factor
\end{minipage} & \begin{minipage}[b]{\linewidth}\raggedright
Effect
\end{minipage} & \begin{minipage}[b]{\linewidth}\raggedright
Mechanism
\end{minipage} \\
\midrule\noalign{}
\endhead
\bottomrule\noalign{}
\endlastfoot
\textbf{Maternal age (\textgreater35)} & Strong, exponential increase in
non-disjunction risk & Cohesin decay during prolonged dictyate arrest \\
\textbf{Achiasmate bivalents} & Strong increase & A bivalent without a
chiasma cannot stay paired through MI; mis-segregates \\
\textbf{Telomeric crossovers primarily} & Increased & Distal-biased
chiasmata are unstable; inadequate to retain bivalent integrity \\
\textbf{Smoking} & Modest increase & Oxidative damage to oocyte DNA and
proteins \\
\textbf{Folate deficiency} & Modest increase & Impaired methylation may
affect centromeric heterochromatin \\
\textbf{Ovarian endometriosis} & Modest increase & Inflammatory effects
on oocyte quality \\
\end{longtable}
}

\textbf{Autosomal trisomies:}

{\def\LTcaptype{none} % do not increment counter
\begin{longtable}[]{@{}
  >{\raggedright\arraybackslash}p{(\linewidth - 8\tabcolsep) * \real{0.1667}}
  >{\raggedright\arraybackslash}p{(\linewidth - 8\tabcolsep) * \real{0.1667}}
  >{\raggedright\arraybackslash}p{(\linewidth - 8\tabcolsep) * \real{0.1970}}
  >{\raggedright\arraybackslash}p{(\linewidth - 8\tabcolsep) * \real{0.3182}}
  >{\raggedright\arraybackslash}p{(\linewidth - 8\tabcolsep) * \real{0.1515}}@{}}
\toprule\noalign{}
\begin{minipage}[b]{\linewidth}\raggedright
Condition
\end{minipage} & \begin{minipage}[b]{\linewidth}\raggedright
Karyotype
\end{minipage} & \begin{minipage}[b]{\linewidth}\raggedright
Key Features
\end{minipage} & \begin{minipage}[b]{\linewidth}\raggedright
Live Birth Frequency
\end{minipage} & \begin{minipage}[b]{\linewidth}\raggedright
Survival
\end{minipage} \\
\midrule\noalign{}
\endhead
\bottomrule\noalign{}
\endlastfoot
\textbf{Down syndrome} & 47,+21 & Intellectual disability (IQ 25-75);
congenital heart defects (40-50\%); increased Alzheimer risk (APP gene
on chr 21); characteristic facies & 1/700 & Most survive to adulthood
(median about 60 years) \\
\textbf{Edwards syndrome} & 47,+18 & Severe; clenched fists with
overlapping fingers; rocker-bottom feet; cardiac defects & 1/5,000 &
about 5-10\% survive to 1 year \\
\textbf{Patau syndrome} & 47,+13 & Holoprosencephaly; polydactyly; cleft
lip/palate; cardiac defects & 1/10,000 & about 5-10\% survive
to 1 year \\
\end{longtable}
}

\textbf{Sex chromosome aneuploidies:}

{\def\LTcaptype{none} % do not increment counter
\begin{longtable}[]{@{}
  >{\raggedright\arraybackslash}p{(\linewidth - 8\tabcolsep) * \real{0.1864}}
  >{\raggedright\arraybackslash}p{(\linewidth - 8\tabcolsep) * \real{0.1864}}
  >{\raggedright\arraybackslash}p{(\linewidth - 8\tabcolsep) * \real{0.2203}}
  >{\raggedright\arraybackslash}p{(\linewidth - 8\tabcolsep) * \real{0.2203}}
  >{\raggedright\arraybackslash}p{(\linewidth - 8\tabcolsep) * \real{0.1864}}@{}}
\toprule\noalign{}
\begin{minipage}[b]{\linewidth}\raggedright
Condition
\end{minipage} & \begin{minipage}[b]{\linewidth}\raggedright
Karyotype
\end{minipage} & \begin{minipage}[b]{\linewidth}\raggedright
Barr Bodies
\end{minipage} & \begin{minipage}[b]{\linewidth}\raggedright
Key Features
\end{minipage} & \begin{minipage}[b]{\linewidth}\raggedright
Frequency
\end{minipage} \\
\midrule\noalign{}
\endhead
\bottomrule\noalign{}
\endlastfoot
\textbf{Turner syndrome} & 45,X & 0 & Female; short stature (SHOX
haploinsufficiency); webbed neck; coarctation of aorta; streak gonads;
infertility & 1/2,500 females \\
\textbf{Klinefelter syndrome} & 47,XXY & 1 & Male; tall; gynecomastia;
small testes; reduced fertility; learning difficulties & 1/600 males \\
\textbf{Triple-X} & 47,XXX & 2 & Female; usually tall; mild learning
difficulties; fertile & 1/1,000 females \\
\textbf{Jacobs syndrome} & 47,XYY & 0 & Male; tall; usually normal
phenotype and fertility & 1/1,000 males \\
\end{longtable}
}

\textbf{Maternal age effect}: Non-disjunction risk for chromosome 21
increases exponentially with maternal age:

{\def\LTcaptype{none} % do not increment counter
\begin{longtable}[]{@{}ll@{}}
\toprule\noalign{}
Maternal Age & Risk of Trisomy 21 \\
\midrule\noalign{}
\endhead
\bottomrule\noalign{}
\endlastfoot
20 & about 1/1,500 \\
25 & about 1/1,250 \\
30 & about 1/900 \\
35 & about 1/270 \\
40 & about 1/100 \\
45 & about 1/50 \\
49 & about 1/12 \\
\end{longtable}
}

\begin{figure}[htbp]
\centering
\includegraphics[width=0.82\linewidth,height=4.2in,keepaspectratio]{../figures/mermaid_inline/inline_mermaid_0098_208a090d4b71.png}
\caption{The cohesin-decay model of the maternal age effect. Female oocytes enter MI prophase before birth and remain arrested for decades. Loss of cohesin during this prolonged arrest weakens chromosome cohesion, producing exponentially rising non-disjunction risk with maternal age.}
\end{figure}

\paragraph{Quantitative model of the maternal age
curve}\label{quantitative-model-of-the-maternal-age-curve}

The empirical risk of trisomy 21 with maternal age is well fit by an
exponential function:

\begin{equation}
P(\text{T21} \mid \text{age } a) \approx P_0 \, e^{k(a - a_0)}
\label{eq:unit_V_maternal_age_risk}
\end{equation}

with \(P_0 \approx 6.7 \times 10^{-4}\) at age \(a_0 = 20\) and
\(k \approx 0.18\) per year for ages 30--45. This corresponds to a
\textbf{doubling of risk every about 4 years} in the late 30s
and early 40s. Between ages 45 and 49, the curve rises even more steeply
as the oocyte pool nears exhaustion. Because the underlying mechanism
--- cohesin decay during dictyate arrest --- accumulates throughout
reproductive life, the exponential form reflects compounding error rates
per oocyte rather than a sudden threshold. Reproductive endocrinology
uses this curve to set screening recommendations; many countries trigger
non-invasive prenatal testing offers automatically at maternal ages
above 35.

Approximately 90\% of trisomy 21 conceptions and 99\% of trisomy 18
conceptions arise in maternal meiosis; paternal contributions are minor.
The asymmetry reflects the fundamental difference between oogenesis
(limited oocyte pool, decades-long arrest, gradual cohesin decay) and
spermatogenesis (continuous stem-cell division, no arrest, robust
quality control through germ-cell apoptosis).

\subsubsection{Structural Chromosomal Rearrangements: Detailed
Survey}\label{structural-chromosomal-rearrangements-detailed-survey}

Beyond losses and gains of whole chromosomes, \textbf{structural
rearrangements} cause a wide range of human disease. Each rearrangement
type has characteristic mechanisms and clinical signatures.

\paragraph{Translocations}\label{translocations}

A \textbf{translocation} is an exchange of chromosomal material between
non-homologous chromosomes. Two major classes exist.

\textbf{Reciprocal translocations} exchange segments between two
non-homologous chromosomes. Carriers of \emph{balanced} reciprocal
translocations usually retain full genetic dosage --- just rearranged
--- and are typically phenotypically normal. However, their gametes face
severe segregation problems at meiosis I: the four chromosomes involved
must form a quadrivalent at metaphase I, and \textbf{alternate
segregation} produces balanced gametes (about 50\% of meioses).
\textbf{Adjacent-1 and adjacent-2 segregation} produces unbalanced
gametes, leading to recurrent miscarriage (10--25\% per pregnancy) and a
1--10\% risk of an unbalanced live-born child.

\textbf{The Philadelphia Chromosome --- A Cancer Paradigm}:

The \textbf{t(9;22)(q34;q11.2)} reciprocal translocation fuses the
\textbf{BCR} gene (chromosome 22) to the \textbf{ABL1} gene (chromosome
9), creating the \textbf{BCR-ABL1} fusion gene on the derivative
chromosome 22 (the Philadelphia chromosome,
\hyperref[gl:ph]{\textbf{pH}}).

\begin{itemize}
\tightlist
\item
  \textbf{BCR-ABL1} encodes a constitutively active tyrosine kinase (210
  kDa in CML; 190 kDa in Philadelphia-positive acute lymphoblastic
  leukaemia)
\item
  Activates RAS/MAPK, PI3K/AKT, and JAK/STAT signaling -- driving
  uncontrolled proliferation
\item
  \textbf{Imatinib (Gleevec)}: ATP-competitive inhibitor of BCR-ABL1
  tyrosine kinase domain

  \begin{itemize}
  \tightlist
  \item
    CML 5-year survival: about 30\% (pre-imatinib) to
    \textgreater90\% (post-imatinib)
  \item
    Resistance mutations (e.g., T315I ``gatekeeper'' mutation)
    necessitated second-generation (dasatinib, nilotinib) and
    third-generation (ponatinib) inhibitors
  \item
    \textbf{Asciminib}: First
    \hyperref[gl:allosteric]{\textbf{allosteric}} BCR-ABL1 inhibitor
    (binds myristoyl pocket, not ATP site); approved 2021
  \end{itemize}
\end{itemize}

The Philadelphia chromosome story is one of the great triumphs of
molecular targeted therapy: a structural rearrangement \texorpdfstring{\ensuremath{\to}}{->} a fusion
oncogene \texorpdfstring{\ensuremath{\to}}{->} a kinase domain \texorpdfstring{\ensuremath{\to}}{->} a small-molecule drug \texorpdfstring{\ensuremath{\to}}{->} a chronic disease
that is no longer fatal for most patients.

\textbf{Robertsonian translocations} fuse two acrocentric chromosomes
(chromosomes 13, 14, 15, 21, 22 in humans) at their centromeres,
producing one large derivative chromosome and losing the small short
arms. Carriers have 45 chromosomes but a normal genome content (the lost
short arms contain primarily ribosomal RNA gene clusters, present in
multiple copies elsewhere). The clinical consequence appears in
offspring: a Robertsonian carrier can produce gametes with two copies of
one chromosome, leading to trisomic offspring. \textbf{Robertsonian
translocation t(14;21) is the cause of approximately 4\% of Down
syndrome cases}, with a recurrence risk of 10--15\% (much higher than
the population rate of about 1/700) --- making karyotype
analysis essential for genetic counseling of recurrent Down syndrome
families.

\paragraph{Worked Example 17.3: Robertsonian translocation
segregation}\label{worked-example-17.3-robertsonian-translocation-segregation}

A phenotypically normal woman has a Robertsonian translocation
rob(14;21), giving her a karyotype 45,XX,rob(14;21). At meiosis, the
trivalent formed by chromosome 14, chromosome 21, and the rob(14;21)
derivative can segregate in \textbf{six} different ways, producing six
possible gamete types --- three balanced and three unbalanced:

{\def\LTcaptype{none} % do not increment counter
\begin{longtable}[]{@{}
  >{\raggedright\arraybackslash}p{(\linewidth - 6\tabcolsep) * \real{0.1039}}
  >{\raggedright\arraybackslash}p{(\linewidth - 6\tabcolsep) * \real{0.2727}}
  >{\raggedright\arraybackslash}p{(\linewidth - 6\tabcolsep) * \real{0.5065}}
  >{\raggedright\arraybackslash}p{(\linewidth - 6\tabcolsep) * \real{0.1169}}@{}}
\toprule\noalign{}
\begin{minipage}[b]{\linewidth}\raggedright
Gamete
\end{minipage} & \begin{minipage}[b]{\linewidth}\raggedright
Chromosomes carried
\end{minipage} & \begin{minipage}[b]{\linewidth}\raggedright
After fertilization with normal sperm
\end{minipage} & \begin{minipage}[b]{\linewidth}\raggedright
Outcome
\end{minipage} \\
\midrule\noalign{}
\endhead
\bottomrule\noalign{}
\endlastfoot
Normal alternate & normal 14 + normal 21 & 46,XX or 46,XY &
Karyotypically normal \\
Translocation alternate & rob(14;21) primarily & 45,XX,rob(14;21) &
Phenotypically normal carrier \\
Adjacent --- extra 14 & rob(14;21) + normal 14 & 46,+14,rob --- trisomy
14 equivalent & Embryonic lethal (early miscarriage) \\
Adjacent --- missing 14 & normal 21 primarily & 45,--14 --- monosomy 14
& Embryonic lethal \\
Adjacent --- extra 21 & rob(14;21) + normal 21 & 46,+21,rob ---
\textbf{translocation Down syndrome} & Live-born child with Down
syndrome \\
Adjacent --- missing 21 & normal 14 primarily & 45,--21 --- monosomy 21
& Embryonic lethal \\
\end{longtable}
}

The empirical recurrence risk of Down syndrome from a maternal
rob(14;21) carrier is \textbf{10--15\%} --- far below the theoretical
1/3 viable expectation because most monosomic and trisomy-14 conceptuses
miscarry early, and because adjacent segregation patterns are not
equally likely. \textbf{For a paternal carrier} the recurrence risk is
much lower (about 1--2\%), because sperm with unbalanced
chromosome content are largely incapable of fertilization (sperm
selection effects). This sex difference in recurrence --- same
translocation, 10\texorpdfstring{\ensuremath{\times}}{x} different risk depending on the parent of origin ---
is one of the most clinically important consequences of male versus
female meiotic biology and dramatically affects genetic counseling.

\paragraph{Inversions}\label{inversions}

An \textbf{inversion} is a chromosomal segment that has been reversed
end-to-end. Two types exist based on whether the centromere is included.

\begin{itemize}
\tightlist
\item
  \textbf{Pericentric inversions} include the centromere; the inverted
  segment spans both arms.
\item
  \textbf{Paracentric inversions} are confined to one arm; the
  centromere is outside the inverted region.
\end{itemize}

Like balanced reciprocal translocations, inversion carriers are
typically phenotypically normal. The clinical problem arises in meiosis:
to pair correctly with the non-inverted homolog, the inverted segment
must form an \textbf{inversion loop}. Crossovers within the loop produce
unbalanced gametes:

\begin{itemize}
\tightlist
\item
  \textbf{Pericentric inversions}: Crossovers inside the loop generate
  gametes with duplications and deficiencies (and corresponding
  live-born offspring with congenital abnormalities). Recurrence risk
  depends on inversion size and crossover frequency, ranging from 1\% to
  15\%.
\item
  \textbf{Paracentric inversions}: Crossovers inside the loop generate
  dicentric and acentric chromosomes that are typically lost during
  meiosis or early embryonic development; consequently, paracentric
  inversion carriers have higher rates of miscarriage but lower rates of
  liveborn unbalanced offspring than pericentric inversion carriers.
\item
  \textbf{Inversion 9} is a common pericentric inversion
  (about 1\% of Europeans) that is generally considered a
  normal variant --- though some studies suggest mildly elevated risk
  for recurrent pregnancy loss.
\end{itemize}

\paragraph{Deletion Syndromes}\label{deletion-syndromes}

\textbf{Deletions} remove chromosomal material and produce
\textbf{haploinsufficiency} --- one functional copy is not enough for
normal function. Some classic deletion syndromes:

{\def\LTcaptype{none} % do not increment counter
\begin{longtable}[]{@{}
  >{\raggedright\arraybackslash}p{(\linewidth - 6\tabcolsep) * \real{0.2222}}
  >{\raggedright\arraybackslash}p{(\linewidth - 6\tabcolsep) * \real{0.2222}}
  >{\raggedright\arraybackslash}p{(\linewidth - 6\tabcolsep) * \real{0.2444}}
  >{\raggedright\arraybackslash}p{(\linewidth - 6\tabcolsep) * \real{0.3111}}@{}}
\toprule\noalign{}
\begin{minipage}[b]{\linewidth}\raggedright
Syndrome
\end{minipage} & \begin{minipage}[b]{\linewidth}\raggedright
Deletion
\end{minipage} & \begin{minipage}[b]{\linewidth}\raggedright
Frequency
\end{minipage} & \begin{minipage}[b]{\linewidth}\raggedright
Key Features
\end{minipage} \\
\midrule\noalign{}
\endhead
\bottomrule\noalign{}
\endlastfoot
\textbf{Cri-du-chat} & 5p15.2 & 1/20,000--50,000 & Distinctive
high-pitched mewing cry in infancy (laryngeal abnormalities);
microcephaly; severe intellectual disability; characteristic facies;
cardiac defects \\
\textbf{DiGeorge / 22q11.2 deletion} & 22q11.2 (3 Mb) & 1/4,000 &
Conotruncal cardiac defects (tetralogy of Fallot, interrupted aortic
arch); thymic hypoplasia (T-cell immunodeficiency); hypocalcemia
(parathyroid hypoplasia); cleft palate; characteristic facies; learning
disabilities; schizophrenia risk increased about 25\texorpdfstring{\ensuremath{\times}}{x} \\
\textbf{Williams syndrome} & 7q11.23 (1.55 Mb, 26 genes including
\emph{ELN}) & 1/7,500--20,000 & ``Elfin'' facies; supravalvular aortic
stenosis (elastin haploinsufficiency); hypercalcemia; intellectual
disability with \textbf{paradoxically preserved language} and
hypersociability \\
\textbf{Prader-Willi syndrome} & 15q11-q13 paternal deletion & 1/15,000
& Hypotonia; hyperphagia; obesity; intellectual disability; short
stature (see imprinting below) \\
\textbf{Angelman syndrome} & 15q11-q13 maternal deletion & 1/15,000 &
Severe intellectual disability; seizures; ataxic gait; characteristic
happy demeanor with frequent laughter \\
\textbf{Wolf-Hirschhorn} & 4p16.3 & 1/50,000 & Severe growth
retardation; ``Greek warrior helmet'' facies; intellectual disability;
seizures \\
\textbf{WAGR syndrome} & 11p13 & rare & Wilms tumor + Aniridia +
Genitourinary abnormalities + intellectual disability (Range) \\
\textbf{Smith-Magenis} & 17p11.2 & 1/25,000 & Intellectual disability;
sleep disturbances (inverted melatonin cycle); behavioral
abnormalities \\
\end{longtable}
}

\textbf{DiGeorge / 22q11.2 deletion syndrome} deserves special attention
as one of the most common microdeletion syndromes and a paradigm for
\textbf{contiguous gene syndromes} --- phenotypes resulting from
haploinsufficiency of multiple adjacent genes simultaneously. The
deletion is typically caused by \textbf{non-allelic homologous
recombination (NAHR)} between low-copy repeats (LCRs) flanking the
region. Within the deleted region, \textbf{TBX1} is the key driver of
cardiac and craniofacial features; other deleted genes contribute to
immune, neurological, and behavioral phenotypes. The 22q11.2 deletion
confers an approximately 25-fold increased risk of schizophrenia, making
it one of the strongest known genetic risk factors for psychotic
illness.

\textbf{Cri-du-chat syndrome} (5p deletion) was first described in 1963
by Jérôme Lejeune (who had identified trisomy 21 four years earlier).
The distinctive cat-like cry in infancy results from laryngeal cartilage
abnormalities. Survival to adulthood is common but with severe
intellectual disability requiring lifelong care.

\paragraph{Duplications}\label{duplications}

Duplications (extra copies of chromosomal segments) typically produce
\textbf{gain-of-function} phenotypes --- opposite to deletions. Many
duplication syndromes are reciprocal to known deletion syndromes:

\begin{itemize}
\tightlist
\item
  \textbf{22q11.2 duplication syndrome}: Reciprocal to DiGeorge. Highly
  variable, often mild; some carriers are phenotypically normal.
\item
  \textbf{Charcot-Marie-Tooth disease type 1A}: Duplication of 17p11.2
  containing \emph{PMP22} causes peripheral neuropathy through PMP22
  overexpression.
\item
  \textbf{MECP2 duplication syndrome} (Xq28): Severe intellectual
  disability and recurrent infections; reciprocal to Rett syndrome
  (caused by \emph{MECP2} loss-of-function mutations in girls).
\end{itemize}

The reciprocal relationship between deletion and duplication phenotypes
confirms gene dosage as the proximate cause and reinforces the
dosage-sensitivity principle developed in the chromosomal-rearrangements
section.

\begin{quote}
\textbf{Clinical Connection: Prenatal Screening for Chromosomal
Abnormalities} \textbf{Cell-free fetal DNA (cfDNA) testing} (NIPT --
non-invasive prenatal testing): Analyzes fetal DNA fragments circulating
in maternal blood (from about 10 weeks gestation). Detects
trisomies 21, 18, 13 and sex chromosome aneuploidies with
\textgreater99\% sensitivity and \textgreater99.5\% specificity for
trisomy 21. Has largely replaced first-trimester combined screening
(nuchal translucency + biochemical markers) as the primary screening
tool. Diagnostic confirmation still requires amniocentesis or chorionic
villus sampling (CVS) for karyotyping or chromosomal microarray.
\textbf{Chromosomal microarray (CMA)} is now the first-line test for
fetal abnormalities or developmental delay, detecting copy-number
changes at sub-megabase resolution.
\end{quote}

\begin{center}\rule{0.5\linewidth}{0.5pt}\end{center}

\subsection{Genomic Imprinting and Uniparental
Disomy}\label{genomic-imprinting-and-uniparental-disomy}

Mendel's laws assume that the phenotypic effect of an allele is
\textbf{independent of which parent transmitted it}. For most loci this
is true. But about 100--200 human genes show \textbf{genomic
imprinting} --- they are expressed from a single parental allele, with
the other allele silenced by parent-of-origin-specific epigenetic marks
(methylation, chromatin) established during gametogenesis
(\cref{sec:unit_IV_epigenetics_and_gene_regulation}).

\subsubsection{The principle of
imprinting}\label{the-principle-of-imprinting}

For an imprinted gene: - A maternally-imprinted (paternally-expressed)
gene contributes primarily the paternal allele to the offspring's
expression. - A paternally-imprinted (maternally-expressed) gene
contributes primarily the maternal allele. - The ``silent'' allele is
not deleted --- it carries methylation marks acquired during egg or
sperm formation that are preserved through embryogenesis.

This violates Mendel's assumption because the phenotype now depends on
\textbf{which parent contributed which allele}, even at otherwise
equivalent loci.

\subsubsection{Canonical imprinted
regions}\label{canonical-imprinted-regions}

{\def\LTcaptype{none} % do not increment counter
\begin{longtable}[]{@{}
  >{\raggedright\arraybackslash}p{(\linewidth - 4\tabcolsep) * \real{0.1250}}
  >{\raggedright\arraybackslash}p{(\linewidth - 4\tabcolsep) * \real{0.2969}}
  >{\raggedright\arraybackslash}p{(\linewidth - 4\tabcolsep) * \real{0.5781}}@{}}
\toprule\noalign{}
\begin{minipage}[b]{\linewidth}\raggedright
Region
\end{minipage} & \begin{minipage}[b]{\linewidth}\raggedright
Imprinting status
\end{minipage} & \begin{minipage}[b]{\linewidth}\raggedright
Disease (loss of expressed allele)
\end{minipage} \\
\midrule\noalign{}
\endhead
\bottomrule\noalign{}
\endlastfoot
\textbf{15q11--q13} & Paternal expression of \emph{SNRPN}, \emph{NDN},
\emph{MAGEL2}, others; maternal expression of \emph{UBE3A} in neurons &
Loss of paternal expression \texorpdfstring{\ensuremath{\to}}{->} \textbf{Prader-Willi syndrome}; loss of
maternal \emph{UBE3A} \texorpdfstring{\ensuremath{\to}}{->} \textbf{Angelman syndrome} \\
\textbf{11p15.5} & Paternal expression of \emph{IGF2}; maternal
expression of \emph{H19} and \emph{CDKN1C} & Disrupted imprinting \texorpdfstring{\ensuremath{\to}}{->}
\textbf{Beckwith-Wiedemann syndrome} (overgrowth, Wilms tumor risk);
reciprocal disruption \texorpdfstring{\ensuremath{\to}}{->} \textbf{Russell-Silver syndrome} (growth
restriction) \\
\textbf{14q32} & \emph{DLK1, MEG3} imprinted & Maternal UPD14 \texorpdfstring{\ensuremath{\to}}{->} Temple
syndrome; paternal UPD14 \texorpdfstring{\ensuremath{\to}}{->} Kagami-Ogata syndrome \\
\end{longtable}
}

\subsubsection{Mechanisms causing imprinting
disorders}\label{mechanisms-causing-imprinting-disorders}

Imprinting disorders arise through several distinct mechanisms --- and a
key clinical insight is that \textbf{the underlying mutation type is
independent of the disease phenotype}, which depends primarily on
whether the active parental allele is present:

\begin{enumerate}
\def\labelenumi{\arabic{enumi}.}
\item
  \textbf{Deletion of the active allele} (about 70--75\% of
  Prader-Willi and Angelman cases). A microdeletion of 15q11-q13 on the
  paternally inherited chromosome causes Prader-Willi; the same
  microdeletion on the maternally inherited chromosome causes Angelman.
  The deletion is identical at the DNA level --- the disease depends on
  which parent transmitted it.
\item
  \textbf{Uniparental disomy (UPD)} (about 25\% of
  Prader-Willi, about 5\% of Angelman cases). Both copies of
  chromosome 15 are inherited from one parent. Maternal UPD15 (no
  paternal contribution) causes Prader-Willi; paternal UPD15 (no
  maternal contribution) causes Angelman. UPD usually arises through
  \textbf{trisomy rescue} --- a trisomic conceptus loses one chromosome
  to revert to disomy; if the lost chromosome was from the parent with a
  single copy, the result is UPD.

  \textbf{Isodisomy vs.~heterodisomy.} UPD comes in two genetically
  distinct flavors that depend on \emph{when} the meiotic error
  occurred:

  \begin{itemize}
  \tightlist
  \item
    \textbf{Heterodisomy} --- both inherited copies derive from
    \textbf{different homologs} of the parent (the two copies are
    heterozygous for parental polymorphisms). Heterodisomy arises from
    \textbf{meiosis I non-disjunction} followed by trisomy rescue. The
    two homologs of one parent are inherited together, and the offspring
    is heterozygous at centromeric markers from that parent.
  \item
    \textbf{Isodisomy} --- both inherited copies are \textbf{identical}
    sister-chromatid copies of a single parental homolog. Isodisomy
    arises from \textbf{meiosis II non-disjunction} plus trisomy rescue,
    or from a postzygotic mitotic error. The offspring is
    \textbf{homozygous along the entire chromosome} for the
    parent-of-origin's allele.
  \end{itemize}

  Isodisomy is clinically important because it can \textbf{unmask
  recessive alleles}: if the contributing parent is a carrier for a
  recessive disease on the affected chromosome, isodisomy produces
  homozygosity for the recessive allele and the offspring is affected
  --- even though a single parent was a carrier. Documented examples
  include cystic fibrosis from maternal isodisomy 7 and rod monochromacy
  from paternal isodisomy 14. Mixed UPD (heterodisomy at the centromere
  with isodisomic distal regions) reflects a meiosis I error followed by
  a crossover; the relative proportions of iso- and heterodisomic
  regions trace the meiotic crossover that occurred in the parent's
  germline.
\item
  \textbf{Imprinting center mutations} (about 1--2\% of cases).
  Mutations in the imprinting control regions disrupt the establishment
  or maintenance of methylation marks, producing functional UPD without
  sequence changes elsewhere.
\item
  \textbf{\emph{UBE3A} point mutations} (about 10\% of Angelman
  cases). Mutations in the \emph{UBE3A} gene itself, which is expressed
  primarily from the maternal allele in neurons, cause Angelman
  syndrome.
\end{enumerate}

\subsubsection{Connection to chromosome
behavior}\label{connection-to-chromosome-behavior}

The imprinting story directly connects molecular biology (DNA
methylation), gametogenesis (where parent-of-origin marks are
established), meiosis (where non-disjunction can produce UPD via trisomy
rescue), and clinical genetics (where the same chromosomal lesion
produces different diseases depending on parental origin). For details
on the establishment and erasure of imprinting marks during germline
development, see \cref{sec:unit_IV_epigenetics_and_gene_regulation}.

\begin{quote}
\textbf{Clinical Connection: Why Both Parents Are Needed} Imprinting
explains why mammalian parthenogenesis (development from an unfertilized
egg) and androgenesis (development from sperm primarily) fail. Mouse
experiments combining two pronuclei from one parent --- gynogenetic
embryos with two maternal pronuclei or androgenetic embryos with two
paternal pronuclei --- fail at distinct embryonic stages, demonstrating
that \textbf{maternal and paternal genomes are not interchangeable}.
Imprinted genes have evolved to enforce biparental contribution in
mammalian reproduction.
\end{quote}

\textbf{Concept Check 17.3}

\begin{quote}
\begin{enumerate}
\def\labelenumi{\arabic{enumi}.}
\tightlist
\item
  A trisomy 21 fetus inherits two centromeric markers from the mother
  --- one identical to grandmother's and one identical to grandfather's
  --- and one paternal centromere. Was the non-disjunction event in
  maternal meiosis I or maternal meiosis II? Explain.
\item
  Using \cref{eq:unit_V_maternal_age_risk} with
  \(P_0 = 6.7 \times 10^{-4}\), \(a_0 = 20\), and \(k = 0.18\), estimate
  the trisomy 21 risk at maternal age 38. How does it compare to age 28?
\item
  A child with cystic fibrosis is born to a mother known to be a CFTR
  carrier and a father whose CFTR sequencing shows two wild-type alleles
  (paternity confirmed). What chromosomal mechanism could explain the
  diagnosis, and what kind of UPD would it require?
\item
  Why does \emph{Drosophila} hyper-transcribe the single male X
  (MOF/MSL) while \emph{C. elegans} halves transcription of both XX
  hermaphrodite chromosomes (DCC), yet mammals silence one X (XIST)?
  What general biological problem do most three solve?
\end{enumerate}
\end{quote}

\begin{center}\rule{0.5\linewidth}{0.5pt}\end{center}

\subsection{Recombination Mapping: Physical vs.~Genetic
Distance}\label{recombination-mapping-physical-vs.-genetic-distance}

\textbf{Genetic distance} (cM) does not always correspond linearly to
\textbf{physical distance} (Mb):

\begin{equation}
1 \text{ cM} \approx 1 \text{ Mb (genome-wide average in humans)}
\label{eq:unit_V_chromosomal_cm_to_mb}
\end{equation}

However, this ratio varies dramatically: - \textbf{Hotspots}:
Recombination rates can be 10-100x the average in small regions
(about 1-2 kb); determined by PRDM9 protein (zinc finger domain
recognizes specific DNA motifs) - \textbf{Coldspots}: Centromeric and
heterochromatic regions have very low recombination - \textbf{Sex
differences}: Female genetic maps are about 1.5-1.8x longer
than male maps (more crossovers in female meiosis)

\textbf{Concept Check 17.4}

\begin{quote}
\begin{enumerate}
\def\labelenumi{\arabic{enumi}.}
\tightlist
\item
  From a three-point testcross of 1\{,\}000 progeny you score 440 + 435
  parentals, 48 + 45 single crossovers in one region, 14 + 12 single
  crossovers in the other, and 4 + 2 double crossovers: use the
  double-crossover class to fix the gene order, then compute both map
  distances and explain why the directly summed parental-to-parental
  distance underestimates the true physical separation.
\end{enumerate}
\end{quote}

\begin{center}\rule{0.5\linewidth}{0.5pt}\end{center}

\subsection{Inbreeding and the Inbreeding
Coefficient}\label{inbreeding-and-the-inbreeding-coefficient}

\subsubsection{The Inbreeding Coefficient
(F)}\label{the-inbreeding-coefficient-f}

The \textbf{inbreeding coefficient (F)} is the probability that an
individual has two alleles at a locus that are identical by descent
(IBD) -- inherited from the same ancestral allele through both parents.

\textbf{Calculation for offspring of first-cousin mating}:

First cousins share one pair of grandparents. For a specific rare
recessive allele (frequency q) in one grandparent:

\begin{equation}
F = \frac{1}{16} = 0.0625 \quad \text{(for first-cousin offspring)}
\label{eq:unit_V_chromosomal_F_first_cousin}
\end{equation}

\textbf{General formula using path coefficients} \citep{wright1922}:

\begin{equation}
F = \sum\left(\frac{1}{2}\right)^{n_1 + n_2 + 1}
\label{eq:unit_V_chromosomal_F_path}
\end{equation}

where \(n_1\) and \(n_2\) are the number of links from one parent to the
common ancestor and from the common ancestor to the other parent, summed
over most paths through most common ancestors.

\subsubsection{Consequences of
Inbreeding}\label{consequences-of-inbreeding}

\begin{itemize}
\tightlist
\item
  \textbf{Increased homozygosity}: \(F\) proportion of loci become
  homozygous beyond random expectation
\item
  \textbf{Exposure of recessive alleles}: Rare deleterious recessive
  alleles are more likely to become homozygous
\item
  \textbf{Inbreeding depression}: Reduced fitness (fertility, viability)
  in inbred populations
\item
  \textbf{Examples}: Amish (Ellis-van Creveld syndrome, 1/5,000
  vs.~1/200,000 in general population due to
  \hyperref[gl:founder-effect]{\textbf{founder effect}} +
  consanguinity); Finnish disease heritage (36 rare diseases enriched in
  Finland due to founder effects)
\end{itemize}

\begin{quote}
\textbf{Clinical Connection: Consanguinity and Genetic Disease} In
populations where consanguineous marriage is common (parts of Middle
East, South Asia, North Africa: 20-50\% of marriages are
consanguineous), the incidence of autosomal recessive disorders is
significantly elevated. First-cousin offspring have about 6-8\%
risk of congenital abnormality (vs.~about 2-3\% in outbred
populations). Genetic counseling and carrier screening programs are
important in these communities.
\end{quote}

\textbf{Concept Check 17.2}

\begin{quote}
\begin{enumerate}
\def\labelenumi{\arabic{enumi}.}
\tightlist
\item
  Why are nearly most calico cats female? What would the karyotype of a
  rare calico male be?
\item
  Explain why maternal age increases the risk of trisomy 21 but paternal
  age does not have the same effect.
\item
  Calculate the inbreeding coefficient for offspring of a half-sibling
  mating (sharing one parent).
\item
  Why does the Philadelphia chromosome produce a constitutively active
  kinase? What is the normal function of ABL1?
\item
  A child has Prader-Willi syndrome with a normal-appearing chromosome
  15 by FISH. What molecular tests would you order, and what mechanisms
  might be responsible?
\item
  Distinguish reciprocal from Robertsonian translocation. Why are
  Robertsonian carriers especially relevant to Down syndrome counseling?
\end{enumerate}
\end{quote}

\begin{center}\rule{0.5\linewidth}{0.5pt}\end{center}

\subsection{Worked Example: Linkage
Analysis}\label{worked-example-linkage-analysis}

\textbf{Problem}: In a testcross, the following offspring are observed
from a female heterozygous for three linked genes (A, B, C) crossed with
a homozygous recessive male:

{\def\LTcaptype{none} % do not increment counter
\begin{longtable}[]{@{}lll@{}}
\toprule\noalign{}
Class & Phenotype & Count \\
\midrule\noalign{}
\endhead
\bottomrule\noalign{}
\endlastfoot
1 & A B C & 412 \\
2 & a b c & 405 \\
3 & A b C & 82 \\
4 & a B c & 78 \\
5 & A B c & 32 \\
6 & a b C & 28 \\
7 & A b c & 3 \\
8 & a B C & 2 \\
Total & & 1,042 \\
\end{longtable}
}

\begin{enumerate}
\def\labelenumi{(\alph{enumi})}
\tightlist
\item
  Determine gene order.
\end{enumerate}

Parental classes: 1 and 2 (most frequent); DCO classes: 7 and 8 (least
frequent).

Compare: Parental = A B C / a b c; DCO = A b c / a B C.

The \textbf{B} gene has switched relative to A and C. Gene order:
\textbf{A -- B -- C}.

\begin{enumerate}
\def\labelenumi{(\alph{enumi})}
\setcounter{enumi}{1}
\tightlist
\item
  Calculate map distances.
\end{enumerate}

\[d_{A-B} = \frac{82 + 78 + 3 + 2}{1042} \times 100 = \frac{165}{1042} \times 100 = 15.8 \text{ cM} \tag{17.8}\]

\[d_{B-C} = \frac{32 + 28 + 3 + 2}{1042} \times 100 = \frac{65}{1042} \times 100 = 6.2 \text{ cM} \tag{17.9}\]

\begin{enumerate}
\def\labelenumi{(\alph{enumi})}
\setcounter{enumi}{2}
\tightlist
\item
  Calculate interference.
\end{enumerate}

\[\text{Expected DCO} = 0.158 \times 0.062 \times 1042 = 10.2 \tag{17.10}\]

\[\text{CoC} = \frac{5}{10.2} = 0.49; \quad I = 1 - 0.49 = 0.51 \tag{17.11}\]

\begin{center}\rule{0.5\linewidth}{0.5pt}\end{center}

\subsection{Computational Bridge}\label{computational-bridge-16}

Recombination and mutation models often begin from pairwise sequence
disparity:

\begin{Shaded}
\begin{Highlighting}[]
\ImportTok{from}\NormalTok{ biology.genetics }\ImportTok{import}\NormalTok{ hamming\_distance}

\BuiltInTok{print}\NormalTok{(hamming\_distance(}\StringTok{"ATGCATGC"}\NormalTok{, }\StringTok{"ATACATGC"}\NormalTok{))}
\end{Highlighting}
\end{Shaded}

\begin{quote}
\textbf{Clinical / systems note:} Non-invasive prenatal testing counts
fetal DNA fragments mapped to chromosomes --- high-throughput
karyotyping that detects trisomies without invasive sampling.
\end{quote}

\begin{center}\rule{0.5\linewidth}{0.5pt}\end{center}

\subsection{Current Evidence and Frontier
Biology}\label{current-evidence-and-frontier-biology-19}

For \textbf{Chromosomal Inheritance and Linkage}, frontier biology
belongs inside the evidence logic of the chapter. Classical genetics
remains essential, but modern interpretation adds penetrance,
polygenicity, structural variation, ancestry-aware inference, and
uncertainty in risk prediction. The core reading question is this:
chromosome-scale inheritance depends on recombination, segregation,
structural variation, and dosage compensation.

\begin{itemize}
\tightlist
\item
  \textbf{What to verify:} identify the observation, model, assay, or
  dataset that would make the claim stronger or weaker.
\item
  \textbf{What to qualify:} state the scale, organism, cell type,
  environmental condition, or population where the claim is expected to
  hold.
\item
  \textbf{What to compare:} test at least one alternative explanation,
  baseline, or null model before treating the pattern as causal.
\item
  \textbf{What to cite:} distinguish primary evidence, review synthesis,
  public dataset, and institutional guidance; for recent or numeric
  claims, prefer the source closest to the measurement and state what
  has changed since it was published.
\end{itemize}

A good genetics answer separates the Mendelian transmission model from
the evidence needed to use it in a population, family, or clinical
setting.

\textbf{Source practice:} For inheritance and population claims,
separate the model assumptions from sampling, ancestry representation,
penetrance, linkage, and environment.

\subsubsection{Structural Variation, Consent, and Chromosome-Scale
Counseling}\label{structural-variation-consent-and-chromosome-scale-counseling}

Chromosome-scale inheritance is clinically useful because it links a
visible or sequenced rearrangement to a reproductive risk. It is also
easy to overstate. Karyotypes, chromosomal microarrays, short-read
sequencing, long-read sequencing, and graph-reference alignment see
different classes of variation: balanced translocations, copy-number
variants, repeat expansions, inversions, and complex rearrangements can
move between visible cytogenetics and sequence-level evidence. The Human
Pangenome Reference highlights why a single linear reference can miss or
misrepresent variants carried by under-sampled ancestries
\citep{humanpangenome2023}.

For prenatal, carrier, or infertility testing, the counseling note
should name the assay and its blind spots, distinguish screening from
diagnosis, state whether a variant is pathogenic or uncertain, and
explain family implications without treating relatives as automatic
research subjects. Consent and data-retention language matter because
chromosomal findings can expose non-paternity, consanguinity, incidental
adult-onset risks, or information relevant to relatives who did not
choose testing.

\subsection{Summary}\label{summary-19}

\begin{itemize}
\tightlist
\item
  \textbf{Chromosome theory}: \citet{sutton1902} proposed chromosomes
  carry genes; Morgan proved this with X-linked \emph{white} in
  \emph{Drosophila}.
\item
  \textbf{Chromosome architecture}: Centromeres (CENP-A,
  alpha-satellite, kinetochore docking), telomeres (TTAGGG, shelterin,
  telomerase), and the heterochromatin/euchromatin distinction shape
  both mechanics and expression.
\item
  \textbf{Meiosis vs.~mitosis}: A two-division reductional/equational
  program with synapsis, crossing over, and reductional segregation at
  MI is the cellular basis of Mendel's laws.
\item
  \textbf{Sex determination}: XX/XY (mammals, SRY gene), ZW/ZZ (birds),
  XO (grasshoppers), haplodiploidy (hymenoptera), temperature-dependent
  (reptiles).
\item
  \textbf{X-inactivation}: Lyon hypothesis (1961); XIST RNA coats Xi;
  Polycomb-mediated H3K27me3; DNA methylation; Barr body.
  about 15\% of genes escape.
\item
  \textbf{Dosage compensation}: X-inactivation in mammals, MSL
  upregulation in \emph{Drosophila}, DCC downregulation in \emph{C.
  elegans}; broader principles include autosomal dosage sensitivity,
  gene-balance constraints, miRNA buffering, and Polycomb maintenance.
\item
  \textbf{X-linked recessive inheritance}: Males affected
  \textgreater\textgreater{} females; no male-to-male transmission;
  carrier mothers pass to 50\% of sons.
\item
  \textbf{Linkage and recombination}: RF \textless{} 50\% indicates
  linkage; 1 cM \texorpdfstring{\ensuremath{\approx}}{~} 1\% recombination; three-point cross determines gene
  order and detects interference; Haldane and Kosambi mapping functions
  correct for undetected double crossovers; PRDM9 directs hotspots;
  female maps are about 1.5--1.8\texorpdfstring{\ensuremath{\times}}{x} longer than male maps.
\item
  \textbf{Non-disjunction}: MI errors produce trisomies inheriting both
  parental homologs; MII errors produce trisomies inheriting two sister
  copies of one homolog. The cohesin-decay model explains the maternal
  age effect.
\item
  \textbf{Chromosomal abnormalities}: Aneuploidy from non-disjunction
  (trisomies 21, 18, 13; Turner 45,X; Klinefelter 47,XXY); structural
  rearrangements include reciprocal and Robertsonian translocations
  (Philadelphia chromosome \texorpdfstring{\ensuremath{\to}}{->} BCR-ABL1 \texorpdfstring{\ensuremath{\to}}{->} CML), pericentric and
  paracentric inversions, deletion syndromes (DiGeorge 22q11.2,
  Cri-du-chat 5p, Williams 7q11.23, Wolf-Hirschhorn 4p), and reciprocal
  duplication syndromes.
\item
  \textbf{Imprinting and uniparental disomy}: about 100--200
  imprinted human genes; deletion of an expressed allele, UPD via
  trisomy rescue, imprinting-center mutations, or \emph{UBE3A} point
  mutations most cause Prader-Willi or Angelman depending on parental
  origin.
\item
  \textbf{Inbreeding}: F = probability of identity by descent; increases
  homozygosity; exposes recessive alleles; inbreeding depression.
\item
  \textbf{Connections:} See
  \cref{sec:unit_IV_dna_replication_and_cell_cycle} for meiosis errors,
  \cref{sec:unit_V_population_genetics} for \(F_{ST}\) and drift, and
  \cref{sec:unit_IV_epigenetics_and_gene_regulation} for imprinting
  disorders.
\end{itemize}

\begin{center}\rule{0.5\linewidth}{0.5pt}\end{center}

\subsection{Review Questions}\label{review-questions-19}

\begin{enumerate}
\def\labelenumi{\arabic{enumi}.}
\tightlist
\item
  Explain Morgan's evidence that the \emph{white} gene is on the X
  chromosome. What cross results would differ if \emph{white} were
  autosomal recessive?
\item
  Describe the role of XIST RNA in X-inactivation. Why do
  about 15\% of genes escape inactivation, and what are the
  phenotypic consequences?
\item
  A woman is a carrier for both hemophilia A (F8, Xq28) and red-green
  color blindness (OPN1LW, Xq28). These genes are about 6 cM
  apart. What proportion of her sons will have both conditions? A
  single? Neither?
\item
  Perform a three-point test cross analysis given the following data:
  {[}provide data if assigned{]}
\item
  Explain the maternal age effect on non-disjunction. What molecular
  mechanism has been proposed?
\item
  Why is trisomy 21 the most common viable autosomal trisomy? (Hint:
  consider chromosome 21 gene content.)
\item
  Describe the BCR-ABL1 fusion and explain why imatinib was a
  breakthrough in cancer therapy. What is the T315I resistance mutation?
\item
  Compare NHEJ and HR as mechanisms for joining chromosome translocation
  breakpoints. Which is more likely to produce the Philadelphia
  chromosome?
\item
  Calculate the inbreeding coefficient for offspring of a double
  first-cousin mating (where both parents are first cousins through
  independent lineages).
\item
  Explain why dosage compensation is necessary and compare the three
  known mechanisms (mammals, \emph{Drosophila}, \emph{C. elegans}).
\item
  A child carries a de novo balanced reciprocal translocation with no
  gene interruption. When is genetic counselling still indicated?
\item
  Contrast \textbf{Robertsonian} vs.~\textbf{reciprocal} translocations
  regarding segregation products at meiosis I.
\item
  Distinguish meiosis I from meiosis II non-disjunction. How would
  centromere-region marker analysis identify the failure mode in a
  trisomic conception?
\item
  A 22q11.2 deletion is found in a child with congenital heart disease.
  List four other phenotypic systems likely affected, and explain why a
  single contiguous deletion produces such diverse features.
\item
  Why does maternal UPD15 cause Prader-Willi syndrome rather than
  Angelman syndrome?
\end{enumerate}

\begin{center}\rule{0.5\linewidth}{0.5pt}\end{center}

\subsection{Further Reading and Source
Notes}\label{further-reading-and-source-notes-19}

\begin{itemize}
\tightlist
\item
  Lyon (1961). Gene action in the \{X\}-chromosome of the mouse (\{Mus
  musculus L.\}). \emph{Nature}, 190.
\item
  Hassold \& Hunt (2001). To err (meiotically) is human: The genesis of
  human aneuploidy. \emph{Nature Reviews Genetics}, 2.
\item
  Wright (1922). Coefficients of inbreeding and relationship.
  \emph{American Naturalist}, 56.
\item
  Sutton (1902). On the morphology of the chromosome group in
  \{Brachystola magna\}. \emph{Biological Bulletin}, 4.
\end{itemize}

\begin{center}\rule{0.5\linewidth}{0.5pt}\end{center}

\subsection{Key Terms}\label{key-terms-19}

\begin{enumerate}
\def\labelenumi{\arabic{enumi}.}
\tightlist
\item
  \textbf{Chromosome theory of heredity} -- genes are located on
  chromosomes; chromosomal behavior in meiosis explains Mendelian laws
\item
  \textbf{Centromere} -- chromosomal region holding sister chromatids;
  site of kinetochore assembly; CENP-A defines identity epigenetically
\item
  \textbf{Telomere} -- terminal TTAGGG repeats bound by shelterin;
  protected by telomerase in stem and germ cells
\item
  \textbf{Heterochromatin} -- densely packed, transcriptionally silenced
  chromatin; constitutive vs.~facultative
\item
  \textbf{SRY} -- sex-determining region Y; master switch for male
  development in mammals
\item
  \textbf{X-inactivation} -- silencing of one X chromosome in female
  somatic cells for dosage compensation
\item
  \textbf{XIST} -- long non-coding RNA that coats and silences the
  inactive X chromosome
\item
  \textbf{Barr body} -- condensed, transcriptionally inactive X
  chromosome visible at nuclear periphery
\item
  \textbf{Hemizygous} -- having a single copy of a gene (e.g., X-linked
  genes in males)
\item
  \textbf{Linkage} -- tendency of genes on the same chromosome to be
  inherited together
\item
  \textbf{Recombination frequency (RF)} -- proportion of recombinant
  offspring; measure of genetic distance
\item
  \textbf{Centimorgan (cM)} -- unit of genetic distance; 1 cM = 1\%
  recombination frequency
\item
  \textbf{Three-point test cross} -- cross using three linked markers to
  determine gene order and distances simultaneously
\item
  \textbf{Interference} -- suppression of nearby crossovers; I = 1 -
  coefficient of coincidence
\item
  \textbf{PRDM9} -- zinc-finger histone methyltransferase that directs
  meiotic recombination hotspots
\item
  \textbf{Non-disjunction} -- failure of chromosomes to separate during
  meiosis, leading to aneuploidy
\item
  \textbf{Aneuploidy} -- abnormal chromosome number (monosomy, trisomy)
\item
  \textbf{Reciprocal translocation} -- exchange of segments between two
  non-homologous chromosomes
\item
  \textbf{Robertsonian translocation} -- centromeric fusion of two
  acrocentric chromosomes
\item
  \textbf{Pericentric inversion} -- inversion that includes the
  centromere
\item
  \textbf{Paracentric inversion} -- inversion confined to one chromosome
  arm
\item
  \textbf{Philadelphia chromosome} -- t(9;22) translocation producing
  BCR-ABL1 fusion; hallmark of CML
\item
  \textbf{Inbreeding coefficient (F)} -- probability that two alleles at
  a locus are identical by descent
\item
  \textbf{LOD score} -- statistical measure of evidence for genetic
  linkage in pedigrees
\item
  \textbf{Pseudoautosomal region} -- segments at X and Y chromosome tips
  that recombine during male meiosis
\item
  \textbf{Dosage compensation} -- mechanism equalizing X-linked gene
  expression between XX and XY individuals
\item
  \textbf{Uniparental disomy (UPD)} -- both copies of a chromosome
  inherited from one parent; cause of imprinting disorders via trisomy
  rescue
\item
  \textbf{Genomic imprinting} -- parent-of-origin-specific expression
  resulting from gametic methylation marks
\end{enumerate}

\begin{center}\rule{0.5\linewidth}{0.5pt}\end{center}

\subsubsection{Companion Source
Module}\label{companion-source-module-19}

\textbf{Chromosomal Inheritance and Linkage} should leave a reproducible
trail from a biological claim to the code, figure, diagram, or
paper-based activity that can test it. Use the surfaces below to inspect
the chapter's assumptions, rerun the relevant model, or compare the
manuscript explanation with companion labs and figures.

{\def\LTcaptype{none} % do not increment counter
\begin{longtable}[]{@{}
  >{\raggedright\arraybackslash}p{(\linewidth - 2\tabcolsep) * \real{0.5000}}
  >{\raggedright\arraybackslash}p{(\linewidth - 2\tabcolsep) * \real{0.5000}}@{}}
\toprule\noalign{}
\begin{minipage}[b]{\linewidth}\raggedright
Surface
\end{minipage} & \begin{minipage}[b]{\linewidth}\raggedright
Use it for
\end{minipage} \\
\midrule\noalign{}
\endhead
\bottomrule\noalign{}
\endlastfoot
\breaktt{src/biology/genetics/genetics.py}
(\breaktt{recombination\_frequency}, \texttt{genetic\_distance},
\breaktt{infer\_three\_point\_order}) & Convert offspring counts into
linkage maps and gene order. \\
\breaktt{src/visualization/plots.py}
(\breaktt{plot\_chromosome\_structure}) & Connect cytogenetic structure
to inheritance patterns. \\
\breaktt{src/mermaid/biology\_diagrams.py}
(\breaktt{chromosome\_inheritance\_diagram},
\breaktt{x\_inactivation\_diagram}) & Compare segregation, linkage, and
dosage compensation. \\
\end{longtable}
}

\textbf{Reproducibility check:} specify phase, recombinant classes,
crossover assumptions, and mapping limits before inferring chromosome
structure. \textbf{Cross-reference:} use
\cref{sec:unit_V_mendelian_genetics},
\cref{sec:unit_IV_epigenetics_and_gene_regulation}, and
\cref{sec:unit_V_population_genetics}.

\newpage

\section{Population Genetics and Hardy-Weinberg
Equilibrium}\label{population-genetics-and-hardy-weinberg-equilibrium}

\label{sec:unit_V_population_genetics}

\begin{quote}
\textbf{Ch 18} \texorpdfstring{\ensuremath{\cdot}}{.} Level 3/3 \texorpdfstring{\ensuremath{\cdot}}{.} 75 min read \texorpdfstring{\ensuremath{\cdot}}{.} 100 min lecture \texorpdfstring{\ensuremath{\cdot}}{.}
Prerequisites: \cref{sec:unit_V_mendelian_genetics},
\cref{sec:unit_V_chromosomal_inheritance}
\end{quote}

\subsection{Learning Objectives}\label{learning-objectives-20}

\begin{enumerate}
\def\labelenumi{\arabic{enumi}.}
\tightlist
\item
  Define \hyperref[gl:allele]{\textbf{allele}} frequency, genotype
  frequency, and \hyperref[gl:gene]{\textbf{gene}} pool.
\item
  State the Hardy-Weinberg \citep{weinberg1908} principle, derive the
  equilibrium equation, and list most five assumptions.
\item
  Apply \hyperref[gl:hardy-weinberg-equilibrium]{\textbf{Hardy-Weinberg
  equilibrium}}
  (Equationabout \eqref{eq:population_genetics_1}) to
  calculate carrier frequencies and estimate disease prevalence.
\item
  Test a population for HWE using the chi-squared test.
\item
  Describe \hyperref[gl:natural-selection]{\textbf{natural selection}}:
  fitness, selection coefficient, dominance coefficient, and types of
  selection.
\item
  Calculate allele frequency change under selection and predict
  equilibrium conditions for balancing selection.
\item
  Explain \hyperref[gl:genetic-drift]{\textbf{genetic drift}}, effective
  population size, bottleneck and
  \hyperref[gl:founder-effect]{\textbf{founder effect}}s.
\item
  Describe gene flow and its effects on population differentiation
  (Fst).
\item
  Apply \hyperref[gl:mutation]{\textbf{mutation}}-selection balance to
  estimate equilibrium frequencies.
\item
  Explain the \hyperref[gl:neutral-theory]{\textbf{neutral theory}} of
  molecular evolution, Ka/Ks ratios, and coalescent theory.
\item
  Use \hyperref[gl:molecular-clock]{\textbf{molecular clock}} methods to
  estimate divergence times.
\end{enumerate}

\subsubsection{Study Blueprint}\label{study-blueprint-20}

\begin{itemize}
\tightlist
\item
  \textbf{Big idea:} Allele frequencies change when assumptions about
  random mating, population size, and fitness are violated.
\item
  \textbf{Core concepts:} Hardy-Weinberg, selection, mutation, genetic
  drift.
\item
  \textbf{Framework alignment:} Vision \& Change: Information flow,
  exchange, and storage, Evolution; AP Biology: Information Storage and
  Transmission, Evolution; NGSS-style topics: Inheritance and Variation
  of Traits, Natural Selection and Evolution.
\item
  \textbf{Model or quantitative lens:} Hardy-Weinberg and
  allele-frequency recurrence calculations.
\item
  \textbf{Data skill:} Estimate genotype or allele frequencies from
  population data.
\item
  \textbf{Practice cadence:} Statistical Tests and Data Analysis,
  Representing and Describing Data.
\item
  \textbf{Common misconception to repair:} Hardy-Weinberg is a null
  model, not a claim that populations do not evolve.
\item
  \textbf{Primary lab:} \cref{sec:lab_unit_V_population_genetics}.
\item
  \textbf{Question bank:} \cref{sec:q_unit_V_population_genetics}.
\item
  \textbf{Transfer task:} Apply population-genetic reasoning to
  screening, conservation, and pathogen evolution.
\item
  \textbf{Bridge to computation:}
  \breaktt{biology.genetics.genetics.hardy\_weinberg}.
\end{itemize}

\begin{center}\rule{0.5\linewidth}{0.5pt}\end{center}

\begin{quote}
\textbf{Opening Vignette --- Algebra, Moths, and the Mathematics of
Populations}

In 1908, a cricket enthusiast named G.H. Hardy --- then England's most
distinguished pure mathematician --- was irritated by a biologist's
claim that \hyperref[gl:dominant]{\textbf{dominant}} alleles
automatically increase in frequency over time. Hardy spent an afternoon
at the cricket match and, in a brief letter to \emph{Science}, produced
the algebraic proof that in the absence of disturbing forces, allele
frequencies remain constant generation after generation. A German
physician, Wilhelm Weinberg, had independently derived the same
principle. The Hardy-Weinberg equilibrium became the null hypothesis of
\hyperref[gl:population-genetics]{\textbf{population genetics}} --- and
its deviations became the heartbeat of evolutionary biology. Nowhere is
this better illustrated than the peppered moth (\emph{Biston betularia})
during England's Industrial Revolution: pollution blackened tree bark,
shifting the population from mostly pale moths to mostly dark (melanic)
ones within decades --- textbook natural selection altering allele
frequencies in real time.
\end{quote}

\subsubsection{Chapter Roadmap}\label{chapter-roadmap-6}

The chapter is long because it bridges three levels of time and scale
--- read it as a narrowing lens:

\begin{itemize}
\tightlist
\item
  \textbf{The Hardy-Weinberg null model.} Alleles, genotype frequencies,
  the equilibrium theorem and its assumptions; \texorpdfstring{\ensuremath{\chi}}{chi}\texorpdfstring{\textsuperscript{2}}{2} tests of fit.
\item
  \textbf{Four forces that perturb it.} Selection, drift, gene flow, and
  mutation. Each is a named violation of an HWE assumption.
\item
  \textbf{Consequences at the \hyperref[gl:genome]{\textbf{genome}}
  scale.} Neutral theory and molecular evolution, F-statistics and
  population structure, coalescent theory and molecular clocks.
\end{itemize}

If you are reading for the core Mendelian-genetics course, prioritise
the Hardy-Weinberg model and the four evolutionary forces. The
genome-scale extensions connect population genetics to molecular
evolution (Unit VI) and are helpful but optional.

\begin{figure}[htbp]
\centering
\includegraphics[width=0.85\textwidth]{../figures/genetic_drift_trajectories.png}
\caption{Wright--Fisher genetic-drift trajectories for the same starting allele frequency under three effective population sizes. Smaller populations show larger stochastic swings and faster approach to fixation or loss.}
\label{fig:unit_V_genetic_drift_trajectories}
\end{figure}

\subsection{The Gene Pool}\label{the-gene-pool}

A \textbf{population} is a group of individuals of the same species
living in the same area at the same time and that interbreed. The
\textbf{gene pool} is the population's allele set across loci.

\subsubsection{Allele and Genotype
Frequencies}\label{allele-and-genotype-frequencies}

For a biallelic locus with alleles A (frequency \(p\)) and a (frequency
\(q\)):

\begin{equation}
p + q = 1
\label{eq:population_genetics_1}
\end{equation}

\textbf{Calculating allele frequencies from genotype counts:}

Given \(N\) total diploid individuals with \(N_{AA}\) homozygous
dominant, \(N_{Aa}\) \hyperref[gl:heterozygous]{\textbf{heterozygous}},
and \(N_{aa}\) homozygous recessive:

\begin{equation}
p = \frac{2N_{AA} + N_{Aa}}{2N}
\label{eq:population_genetics_2}
\end{equation}

\begin{equation}
q = \frac{2N_{aa} + N_{Aa}}{2N}
\label{eq:population_genetics_3}
\end{equation}

\textbf{Genotype frequencies:}

\begin{equation}
F_{AA} + F_{Aa} + F_{aa} = 1
\label{eq:population_genetics_4}
\end{equation}

Note that allele frequencies can be calculated from genotype frequencies
but not vice versa (without additional assumptions), because different
combinations of genotype frequencies can yield the same allele
frequencies.

\subsection{Worked Example: 18.1: In a population of 1,000 individuals:
360 AA, 480 Aa, 160
aa.}\label{worked-example-18.1-in-a-population-of-1000-individuals-360-aa-480-aa-160-aa.}

\begin{equation}
p = \frac{2(360) + 480}{2(1000)} = \frac{1200}{2000} = 0.60
\label{eq:population_genetics_5}
\end{equation}

\begin{equation}
q = \frac{2(160) + 480}{2(1000)} = \frac{800}{2000} = 0.40
\label{eq:population_genetics_6}
\end{equation}

Check: \(p + q = 0.60 + 0.40 = 1.00\) (correct).

\textbf{Concept Check 18.5}

\begin{quote}
\begin{enumerate}
\def\labelenumi{\arabic{enumi}.}
\tightlist
\item
  A population of 1\{,\}000 has 250 AA, 500 Aa, and 250 aa: show that
  although \(p = q = 0.5\), the genotype counts are not those
  Hardy-Weinberg predicts, and decide from the direction of the
  heterozygote discrepancy whether inbreeding or heterozygote advantage
  is the more likely cause.
\end{enumerate}
\end{quote}

\begin{center}\rule{0.5\linewidth}{0.5pt}\end{center}

\subsection{Hardy-Weinberg
Equilibrium}\label{hardy-weinberg-equilibrium}

The \textbf{Hardy-Weinberg principle} (G.H. Hardy and Wilhelm Weinberg,
1908) states: in a large, randomly mating population with no
evolutionary forces acting, allele frequencies and genotype frequencies
remain constant from generation to generation.

\subsubsection{Derivation}\label{derivation}

Under random mating, \hyperref[gl:gamete]{\textbf{gamete}}s combine
randomly. The probability of drawing an A allele from the gene pool is
\(p\), and of drawing an a allele is \(q\):

\begin{equation}
P(AA) = p \times p = p^2
\label{eq:population_genetics_7}
\end{equation} \begin{equation}
P(Aa) = p \times q + q \times p = 2pq
\label{eq:population_genetics_8}
\end{equation} \begin{equation}
P(aa) = q \times q = q^2
\label{eq:population_genetics_9}
\end{equation}

Therefore, at equilibrium:

\begin{equation}
\boxed{p^2 + 2pq + q^2 = 1}
\label{eq:population_genetics_10}
\end{equation}

This can also be written as \((p + q)^2 = 1\), which is simply the
binomial expansion.

\textbf{Key insight}: HW equilibrium is reached in a \textbf{single
generation} of random mating (for autosomal loci), and allele
frequencies do not change across generations.

\begin{figure}[htbp]
\centering
\includegraphics[width=0.82\linewidth,height=4.2in,keepaspectratio]{../figures/mermaid_inline/inline_mermaid_0099_e2230bd94153.png}
\caption{The five assumptions of Hardy-Weinberg equilibrium and the evolutionary forces that result from their violation. HWE serves as a null model -- deviations indicate evolutionary forces are acting.}
\end{figure}

\subsubsection{Testing for Hardy-Weinberg
Equilibrium}\label{testing-for-hardy-weinberg-equilibrium}

\textbf{Chi-squared test for HWE:} A real population is tested against
HWE as the null model. Estimate the allele frequencies \(p\) and \(q\)
from the observed genotype counts, compute the genotype counts
\emph{expected} under HWE (\(Np^2\), \(2Npq\), \(Nq^2\) for a sample of
\(N\) individuals), and form \(\chi^2 = \sum (O-E)^2 / E\) over the
genotype classes. The degrees of freedom equal the number of genotype
classes minus the number of alleles --- for a biallelic locus,
\(3 - 2 = 1\) (one degree is spent estimating \(p\) from the same data).
A statistic exceeding the critical value \(\chi^2_{0.05,\,1} = 3.84\)
rejects HWE and indicates that at least one assumption --- random
mating, no selection, no drift, no migration, no mutation --- is
violated.

\paragraph{Worked Example: MN Blood Group
Equilibrium}\label{worked-example-mn-blood-group-equilibrium}

MN blood group data for codominant alleles M and N in a population of
1,000:

{\def\LTcaptype{none} % do not increment counter
\begin{longtable}[]{@{}lll@{}}
\toprule\noalign{}
Genotype & Observed & Allele Frequency \\
\midrule\noalign{}
\endhead
\bottomrule\noalign{}
\endlastfoot
MM & 298 & \\
MN & 489 & \\
NN & 213 & \\
\end{longtable}
}

Step 1: Calculate allele frequencies.

\begin{equation}
p(M) = \frac{2(298) + 489}{2000} = \frac{1085}{2000} = 0.5425
\label{eq:population_genetics_11}
\end{equation}

\begin{equation}
q(N) = \frac{2(213) + 489}{2000} = \frac{915}{2000} = 0.4575
\label{eq:population_genetics_12}
\end{equation}

Step 2: Calculate expected genotype frequencies under HWE.

\begin{equation}
E(MM) = p^2 \times 1000 = (0.5425)^2 \times 1000 = 294.3
\label{eq:population_genetics_13}
\end{equation}

\begin{equation}
E(MN) = 2pq \times 1000 = 2(0.5425)(0.4575) \times 1000 = 496.4
\label{eq:population_genetics_14}
\end{equation}

\begin{equation}
E(NN) = q^2 \times 1000 = (0.4575)^2 \times 1000 = 209.3
\label{eq:population_genetics_15}
\end{equation}

Step 3: Chi-squared test.

\begin{equation}
\chi^2 = \frac{(298-294.3)^2}{294.3} + \frac{(489-496.4)^2}{496.4} + \frac{(213-209.3)^2}{209.3}
\label{eq:population_genetics_16}
\end{equation}

\begin{equation}
= 0.047 + 0.110 + 0.065 = 0.222
\label{eq:population_genetics_17}
\end{equation}

Degrees of freedom = 3 genotypes - 1 constraint (total N) - 1 estimated
parameter (p) = \textbf{1 df}.

\(\chi^2_{crit}(1 \text{ df}, \alpha = 0.05) = 3.841\)

Since \(0.222 < 3.841\), we \textbf{fail to reject} HWE. This population
is in Hardy-Weinberg equilibrium for the MN locus.

\subsubsection{HWE as a Clinical Tool}\label{hwe-as-a-clinical-tool}

\textbf{Estimating carrier frequencies:}

Hardy-Weinberg estimates are useful clinical first approximations, not
diagnoses. Carrier-frequency calculations assume random mating, correct
disease prevalence, complete ascertainment, and a well-defined
population; those assumptions can fail when founder effects,
consanguinity, population stratification, penetrance, or reference bias
shape the sample. Modern screening therefore combines the HWE null model
with ancestry-aware variant panels, family history, and increasingly
graph-aware or long-read evidence for loci where structural variants
affect interpretation \citep{humanpangenome2023}.

\subsection{Worked Example: Cystic Fibrosis Carrier
Frequency}\label{worked-example-cystic-fibrosis-carrier-frequency}

Cystic fibrosis (CF) prevalence in Northern Europeans has disease (aa)
frequency \(\frac{1}{2,500}\).

\begin{equation}
q^2 = \frac{1}{2500} \implies q = \frac{1}{50} = 0.02
\label{eq:population_genetics_18}
\end{equation}

\begin{equation}
p = 1 - q = 0.98
\label{eq:population_genetics_19}
\end{equation}

\begin{equation}
\text{Carrier frequency} = 2pq = 2(0.98)(0.02) = 0.0392 \approx \frac{1}{25}
\label{eq:population_genetics_20}
\end{equation}

Approximately 1 in 25 Northern Europeans is a CF carrier -- critical
information for genetic counseling.

\textbf{ABO blood group HWE analysis} (three alleles: \(I^A\), \(I^B\),
\(i\)):

For a triallelic system: \(p(I^A) + q(I^B) + r(i) = 1\)

Genotype frequencies under HWE:

\begin{equation}
\text{Type A} = p^2 + 2pr, \quad \text{Type B} = q^2 + 2qr, \quad \text{Type AB} = 2pq, \quad \text{Type O} = r^2
\label{eq:population_genetics_21}
\end{equation}

\begin{quote}
\textbf{Clinical Connection: Newborn Screening and Hardy-Weinberg}
Hardy-Weinberg predictions guide newborn screening programs. For PKU (q
approximately 0.01, disease frequency 1/10,000), carrier frequency
approximately 2\%. For sickle cell disease in African Americans (q
approximately 0.05), carrier frequency approximately 9.5\%. For
Tay-Sachs in Ashkenazi Jewish populations (q approximately 0.018),
carrier frequency approximately 1/28. These calculations determine
cost-effectiveness of carrier screening programs.
\end{quote}

\textbf{Concept Check 18.1}

\begin{quote}
\begin{enumerate}
\def\labelenumi{\arabic{enumi}.}
\tightlist
\item
  Can genotype frequencies be in HWE if allele frequencies are changing?
  Explain.
\item
  Inbreeding does not change allele frequencies but does change genotype
  frequencies. Explain how.
\item
  Why does it take a single generation of random mating to restore HWE
  for an autosomal locus?
\item
  For a disease with q = 0.001, what is the expected carrier frequency?
  What is the ratio of carriers to affected individuals?
\end{enumerate}
\end{quote}

\begin{center}\rule{0.5\linewidth}{0.5pt}\end{center}

\subsection{Natural Selection}\label{natural-selection}

\subsubsection{Fitness and Selection}\label{fitness-and-selection}

\textbf{Absolute \hyperref[gl:fitness]{\textbf{fitness (w)}}}: Expected
number of offspring produced by an individual with a given genotype.

\textbf{Relative fitness (w)}: Fitness of a genotype relative to the
most fit genotype (which is set to 1).

\textbf{\hyperref[gl:selection-coefficient]{Selection coefficient (s)}}:
The reduction in fitness: \(s = 1 - w\).

\textbf{Dominance coefficient (h)}: Describes the fitness of the
heterozygote relative to the two homozygotes:

{\def\LTcaptype{none} % do not increment counter
\begin{longtable}[]{@{}ll@{}}
\toprule\noalign{}
Genotype & Fitness \\
\midrule\noalign{}
\endhead
\bottomrule\noalign{}
\endlastfoot
AA & \(w_{AA} = 1\) \\
Aa & \(w_{Aa} = 1 - hs\) \\
aa & \(w_{aa} = 1 - s\) \\
\end{longtable}
}

\begin{itemize}
\tightlist
\item
  \(h = 0\): A completely dominant (Aa has same fitness as AA)
\item
  \(h = 0.5\): Additive (codominant fitness effect)
\item
  \(h = 1\): A completely recessive (Aa has same fitness as aa)
\item
  \(h < 0\): Overdominance (heterozygote advantage)
\end{itemize}

\subsubsection{Change in Allele Frequency Under
Selection}\label{change-in-allele-frequency-under-selection}

For selection against a recessive homozygote (h = 0):

\begin{equation}
\Delta q = \frac{-spq^2}{\bar{w}}
\label{eq:population_genetics_22}
\end{equation}

where \(\bar{w} = 1 - sq^2\) is the mean fitness of the population.

\textbf{Important result}: Selection against a recessive allele becomes
very slow at low frequencies because the allele is ``hidden'' in
heterozygotes. The time to reduce \(q\) from 0.01 to 0.001 is much
longer than from 0.5 to 0.01.

For selection against a dominant allele: much faster because every
carrier is exposed to selection.

\subsection{Worked Example: Recessive Deleterious
Allele}\label{worked-example-recessive-deleterious-allele}

In a population with \(q = 0.3\) (frequency of a recessive deleterious
allele) and \(s = 0.2\) (20\% fitness reduction for aa homozygotes):

\begin{equation}
\bar{w} = 1 - s \cdot q^2 = 1 - 0.2(0.09) = 1 - 0.018 = 0.982
\label{eq:population_genetics_23}
\end{equation}

\begin{equation}
\Delta q = \frac{-0.2 \times 0.7 \times 0.09}{0.982} = \frac{-0.0126}{0.982} = -0.0128
\label{eq:population_genetics_24}
\end{equation}

After one generation: \(q' = 0.3 - 0.0128 = 0.287\)

\subsubsection{Types of Selection}\label{types-of-selection}

\begin{figure}[htbp]
\centering
\includegraphics[width=0.82\linewidth,height=4.2in,keepaspectratio]{../figures/mermaid_inline/inline_mermaid_0100_0b2cc58eea47.png}
\caption{Types of natural selection and their effects on allele frequency distributions. Directional selection shifts the mean; stabilizing selection reduces variance; disruptive selection increases variance; balancing selection maintains polymorphism.}
\end{figure}

\subsubsection{Heterozygote Advantage (Balancing
Selection)}\label{heterozygote-advantage-balancing-selection}

When the heterozygote has higher fitness than either homozygote
(\(w_{Aa} > w_{AA}\) and \(w_{Aa} > w_{aa}\)), both alleles are
maintained at a \textbf{stable equilibrium}:

\begin{equation}
\hat{q} = \frac{s_1}{s_1 + s_2}
\label{eq:population_genetics_25}
\end{equation}

where \(s_1\) = selection against AA and \(s_2\) = selection against aa.

\textbf{The sickle cell-malaria paradigm:}

In malaria-endemic regions of sub-Saharan Africa: - \textbf{HbA/HbA}
(normal): susceptible to \emph{Plasmodium falciparum} malaria;
\(w_{AA} = 1 - s_1\), where \(s_1 \approx 0.1-0.15\) - \textbf{HbA/HbS}
(sickle trait): about 10-15x reduced malaria mortality; highest
fitness; \(w_{Aa} = 1\) - \textbf{HbS/HbS} (sickle cell disease): severe
disease; \(w_{aa} = 1 - s_2\), where \(s_2 \approx 0.8-1.0\)

Equilibrium frequency of HbS allele:

\begin{equation}
\hat{q} = \frac{s_1}{s_1 + s_2} = \frac{0.12}{0.12 + 0.90} = 0.118
\label{eq:population_genetics_26}
\end{equation}

This predicts HbS frequency of about 12\%, remarkably close to
the observed about 10-20\% in malaria-endemic West Africa.

\textbf{Other examples of heterozygote advantage:} - \textbf{HbC}
(beta-globin E6K): heterozygotes also resistant to malaria; found in
West Africa - \textbf{G6PD deficiency}: heterozygous females have mosaic
RBCs; malaria resistance - \textbf{CF heterozygotes}: possibly resistant
to cholera or typhoid (debated); explains high carrier frequency
(about 4\% in Europeans) - \textbf{MHC/HLA diversity}:
heterozygotes present more pathogen antigens; frequency-dependent
selection also contributes

\textbf{Concept Check 18.2}

\begin{quote}
\begin{enumerate}
\def\labelenumi{\arabic{enumi}.}
\tightlist
\item
  Why does selection against a recessive allele become slower as the
  allele frequency decreases?
\item
  If heterozygote advantage maintains the sickle cell allele at
  about 12\% in malaria-endemic regions, predict what happens
  to HbS frequency in a population that migrates to a malaria-free
  region.
\item
  Calculate the equilibrium frequency of a deleterious recessive allele
  if \(s_1 = 0.05\) and \(s_2 = 0.50\).
\end{enumerate}
\end{quote}

\begin{center}\rule{0.5\linewidth}{0.5pt}\end{center}

\subsection{Genetic Drift}\label{genetic-drift}

In finite populations, allele frequencies change from generation to
generation by \textbf{random sampling error} -- this is \textbf{genetic
drift}.

\subsubsection{Wright-Fisher Model}\label{wright-fisher-model}

The Wright-Fisher model is the simplest mathematical framework for
drift: - Non-overlapping generations - Constant population size \(N\) -
Random mating - No selection, mutation, or migration

Each generation, \(2N\) allele copies are drawn randomly (with
replacement) from the parental gene pool.

\textbf{Variance in allele frequency change per generation:}

\begin{equation}
\sigma^2_{\Delta p} = \frac{pq}{2N_e}
\label{eq:population_genetics_27}
\end{equation}

where \(N_e\) = \textbf{effective population size} (the size of an ideal
Wright-Fisher population that would experience the same rate of drift).

\textbf{Key results of drift:} - Over time, alleles are either
\textbf{fixed} (\(p = 1\)) or \textbf{lost} (\(p = 0\)) - The
probability of fixation for a new neutral allele = its initial frequency
= \(\frac{1}{2N}\) - Time to fixation of a new neutral mutation =
\(4N_e\) generations - Heterozygosity decreases at rate
\(\frac{1}{2N_e}\) per generation

\begin{figure}[htbp]
\centering
\includegraphics[width=0.82\linewidth,height=4.2in,keepaspectratio]{../figures/mermaid_inline/inline_mermaid_0101_62257ac9f73f.png}
\caption{Genetic drift in populations of different sizes. Small populations experience rapid allele frequency changes and loss of diversity. Bottleneck events dramatically reduce effective population size, causing lasting genetic effects.}
\end{figure}

\subsubsection{\texorpdfstring{Effective Population Size
(\(N_e\))}{Effective Population Size (N_e)}}\label{effective-population-size-n_e}

\(N_e\) is usually \textbf{much smaller} than census size \(N\) because:

{\def\LTcaptype{none} % do not increment counter
\begin{longtable}[]{@{}
  >{\raggedright\arraybackslash}p{(\linewidth - 4\tabcolsep) * \real{0.2353}}
  >{\raggedright\arraybackslash}p{(\linewidth - 4\tabcolsep) * \real{0.5000}}
  >{\raggedright\arraybackslash}p{(\linewidth - 4\tabcolsep) * \real{0.2647}}@{}}
\toprule\noalign{}
\begin{minipage}[b]{\linewidth}\raggedright
Factor
\end{minipage} & \begin{minipage}[b]{\linewidth}\raggedright
Effect on \(N_e\)
\end{minipage} & \begin{minipage}[b]{\linewidth}\raggedright
Formula
\end{minipage} \\
\midrule\noalign{}
\endhead
\bottomrule\noalign{}
\endlastfoot
Unequal sex ratio & Reduces \(N_e\) &
\(N_e = \frac{4N_mN_f}{N_m + N_f}\) \\
Fluctuating population size & Reduces \(N_e\) & \$N\_e = \$ harmonic
mean of \(N\) over generations \\
Variance in reproductive success & Reduces \(N_e\) & High variance (few
individuals produce most offspring) reduces \(N_e\) \\
\end{longtable}
}

The three workhorse formulas for \(N_e\) are presented below --- each
corresponds to a distinct biological scenario, and conservation
geneticists routinely combine them to estimate \(N_e\) from real
demographic data \citep{frankham1995, frankham2014}.

\paragraph{\texorpdfstring{Sex-ratio \(N_e\)
formula}{Sex-ratio N_e formula}}\label{sex-ratio-n_e-formula}

When the breeding population has unequal numbers of males (\(N_m\)) and
females (\(N_f\)), the effective population size is governed by the
rarer sex because each offspring receives one allele from each:

\begin{equation}
N_e = \frac{4 N_m N_f}{N_m + N_f}
\label{eq:unit_V_ne_sex}
\end{equation}

\textbf{Worked example --- sex-biased breeding:} Consider a herd of 100
elk consisting of 10 breeding males and 90 breeding females (a strongly
female-biased operational sex ratio common in harem-forming ungulates):

\[N_e = \frac{4 \times 10 \times 90}{10 + 90} = \frac{3{,}600}{100} = 36 \label{eq:unit_V_population_genetics_item_1}\]

Despite a census size of 100, the effective population size is primarily
\textbf{36} --- barely a third of the count. The intuition: every
offspring's paternal allele must come from one of just 10 males, so 90\%
of females are functionally redundant for tracking allelic diversity
through paternal lineages. This asymmetry is why captive-breeding
programs aggressively manage sex ratios and why studbook software
optimizes founder representation rather than census numbers. If the same
population had 50 males and 50 females, \(N_e\) would equal 100 --- the
full census size.

\paragraph{\texorpdfstring{Bottleneck \(N_e\)
formula}{Bottleneck N_e formula}}\label{bottleneck-n_e-formula}

For a population that experiences a single-generation bottleneck of size
\(N_b\) embedded within \(t\) generations of size \(N_0\), the effective
population size over the entire interval is the \textbf{harmonic mean}
of generation sizes:

\begin{equation}
\frac{1}{N_e} = \frac{1}{t} \sum_{i=1}^{t} \frac{1}{N_i}
\label{eq:unit_V_ne_bottle}
\end{equation}

The reciprocal-mean form is the same harmonic structure given
equivalently as:

\begin{equation}
N_e = \frac{t}{\sum_{i=1}^{t} \frac{1}{N_i}}
\label{eq:unit_V_ne_harmonic}
\end{equation}

Because the harmonic mean is dominated by the smallest term, \textbf{a
single bottleneck generation can drag \(N_e\) far below the long-term
arithmetic mean of population size}. As an example: 99 generations of
\(N_0 = 10{,}000\) plus one generation of \(N_b = 10\) yields:

\[N_e = \frac{100}{99/10{,}000 + 1/10} = \frac{100}{0.0099 + 0.1} = \frac{100}{0.1099} \approx 910 \label{eq:unit_V_population_genetics_item_2}\]

--- not the arithmetic mean of about 9,901. \textbf{One bad
generation undoes a century of large population size}, and this
asymmetry is precisely why bottleneck events leave such durable genetic
signatures.

\paragraph{\texorpdfstring{\(N_e\) for fluctuating
populations}{N_e for fluctuating populations}}\label{n_e-for-fluctuating-populations}

For populations that vary cyclically (boom-bust dynamics in microbes,
seasonally migrating species, populations with cyclic predator-prey
dynamics):

\begin{equation}
N_e \approx \frac{1}{\frac{1}{T}\sum_{i=1}^{T}\frac{1}{N_i}}
\label{eq:unit_V_population_Ne_fluctuating}
\end{equation}

Drosophila in seasonal climates oscillate between \(N_b \approx 10^3\)
in winter and \(N_0 \approx 10^7\) in summer --- and the harmonic mean
yields \(N_e\) close to \(4 \times N_b\), not anywhere near the summer
maximum.

\paragraph{\texorpdfstring{\(N_e\) with variance in reproductive
success}{N_e with variance in reproductive success}}\label{n_e-with-variance-in-reproductive-success}

When the variance in offspring number (\(V_k\)) exceeds the mean
(\(\bar{k} = 2\) for a stable population), \(N_e\) is reduced:

\begin{equation}
N_e \approx \frac{4N - 2}{V_k + 2}
\label{eq:unit_V_population_Ne_variance}
\end{equation}

In species with extreme reproductive skew --- elephant seals (one
beachmaster male sires nearly the offspring), salmon (a few females
spawn the majority of fertilized eggs), or hierarchical primates ---
\(N_e\) may be 10--100\texorpdfstring{\ensuremath{\times}}{x} smaller than census size. \textbf{Conservation
programs that calculate \(N_e\) from stud-book pedigree data routinely
find \(N_e \ll N\)}, with significant implications for genetic-rescue
planning.

\textbf{Human \(N_e\)}: Despite about 8 billion people alive
today, human \(N_e\) estimated from genetic diversity is about
10,000-15,000, reflecting an ancient bottleneck about 70,000
years ago (possibly associated with the Toba supervolcanic eruption,
though this is debated). The harmonic-mean nature of \(N_e\) explains
why such a long-ago bottleneck still dominates the genetic diversity
calculation despite billions of intervening generations of population
growth.

\subsubsection{Bottleneck and Founder
Effects}\label{bottleneck-and-founder-effects}

\textbf{Bottleneck effect}: Severe reduction in population size causes
random loss of alleles: - \textbf{Cheetah}: Extreme genetic homogeneity;
can accept skin grafts from unrelated individuals; \(N_e\) probably
dropped to hundreds about 10,000 years ago - \textbf{Northern
elephant seal}: Hunted to about 20 individuals by 1890s; now
\textgreater200,000 but almost no genetic variation at many loci

\textbf{Founder effect}: A small group colonizes a new area, carrying a
non-representative sample of the original population's alleles: -
\textbf{Amish/Old Order Mennonites}: Founded by about 200
individuals in 18th century; high frequency of Ellis-van Creveld
syndrome, maple syrup urine disease, glutaric aciduria type I -
\textbf{Finnish disease heritage}: 36 rare genetic diseases at elevated
frequencies due to sequential founder events - \textbf{Afrikaners}: High
frequency of porphyria variegata (traced to a single Dutch couple who
arrived in 1688)

\textbf{Concept Check 18.4}

\begin{quote}
\begin{enumerate}
\def\labelenumi{\arabic{enumi}.}
\tightlist
\item
  An allele tracked for ten generations rises smoothly and monotonically
  from \(q = 0.10\) to \(q = 0.45\) in a census population of millions,
  whereas in a population of 30 it instead jumps erratically up and down
  between 0.05 and 0.60: explain which trajectory is the signature of
  natural selection and which of genetic drift, and why population size
  --- not the size of the change --- is the deciding feature.
\end{enumerate}
\end{quote}

\begin{center}\rule{0.5\linewidth}{0.5pt}\end{center}

\subsection{Gene Flow (Migration)}\label{gene-flow-migration}

Movement of individuals (and their alleles) between populations.

\subsubsection{One-Island Model}\label{one-island-model}

If a fraction \(m\) of population 1 is replaced by migrants from
population 2 each generation:

\begin{equation}
p_1' = (1-m)p_1 + mp_2
\label{eq:population_genetics_28}
\end{equation}

\begin{equation}
\Delta p_1 = m(p_2 - p_1)
\label{eq:population_genetics_29}
\end{equation}

Gene flow \textbf{homogenizes} populations -- reducing differentiation.
Even a small amount of gene flow (\(Nm > 1\) migrant per generation) is
sufficient to prevent populations from diverging by drift alone.

\subsubsection{\texorpdfstring{Measuring Population Differentiation:
\(F_{ST}\)}{Measuring Population Differentiation: F_(ST)}}\label{measuring-population-differentiation-f_st}

\textbf{Wright's F-statistics} quantify the distribution of genetic
variation within and among populations:

\begin{equation}
F_{ST} = \frac{H_T - H_S}{H_T}
\label{eq:population_genetics_30}
\end{equation}

where \(H_T\) = expected heterozygosity in the total population and
\(H_S\) = average expected heterozygosity within subpopulations.

{\def\LTcaptype{none} % do not increment counter
\begin{longtable}[]{@{}ll@{}}
\toprule\noalign{}
\(F_{ST}\) Value & Interpretation \\
\midrule\noalign{}
\endhead
\bottomrule\noalign{}
\endlastfoot
0 & No differentiation (panmixia) \\
0-0.05 & Little differentiation \\
0.05-0.15 & Moderate differentiation \\
0.15-0.25 & Great differentiation \\
\textgreater0.25 & Very great differentiation \\
\end{longtable}
}

\subsubsection{\texorpdfstring{Interpreting \(F_{ST}\) in Human
Populations}{Interpreting F_(ST) in Human Populations}}\label{interpreting-f_st-in-human-populations}

\textbf{Quick interpretation reference for \(F_{ST}\)}:
continental-scale comparisons between human populations consistently
fall in the \textbf{0.10--0.15 range} --- moderate differentiation by
Wright's classification. This empirical magnitude provides a benchmark
when interpreting \(F_{ST}\) from any other species: humans are
\emph{less} differentiated than most large mammals, \emph{comparably}
differentiated to some highly mobile species (large carnivores,
migratory birds), and \emph{more} differentiated than panmictic marine
species with planktonic larvae.

\textbf{Human \(F_{ST} \approx 0.10\)--\(0.15\)}: Human populations are
about 10--15\% differentiated -- meaning about 85--90\% of
human genetic variation exists within populations, not between them
(Lewontin, 1972; Rosenberg et al., 2002). This is lower than most large
mammals, reflecting recent common ancestry (the species is about 300,000
years old) and substantial gene flow throughout human history.

\paragraph{\texorpdfstring{Pairwise \(F_{ST}\) between human
populations}{Pairwise F_(ST) between human populations}}\label{pairwise-f_st-between-human-populations}

Continental-scale \(F_{ST}\) values from the 1000 Genomes and HGDP
projects illustrate the typical magnitudes:

{\def\LTcaptype{none} % do not increment counter
\begin{longtable}[]{@{}ll@{}}
\toprule\noalign{}
Population pair & Approximate \(F_{ST}\) \\
\midrule\noalign{}
\endhead
\bottomrule\noalign{}
\endlastfoot
West African vs.~East Asian & 0.10--0.15 \\
West African vs.~European & 0.10--0.13 \\
European vs.~East Asian & 0.07--0.11 \\
Within Europe (e.g., French vs.~Russian) & 0.005--0.02 \\
Within East Asia (Han Chinese vs.~Japanese) & 0.005--0.015 \\
Within sub-Saharan Africa (Yoruba vs.~Maasai) & 0.02--0.05 \\
Globally averaged across most human populations & about 0.12 \\
\end{longtable}
}

The largest within-Africa \(F_{ST}\) values (sub-Saharan vs.~Khoisan)
reach about 0.10, comparable to between-continental comparisons
elsewhere --- a genetic signature of Africa's status as the origin and
longest-inhabited continent, with the deepest population-genetic
divergences globally.

\paragraph{\texorpdfstring{What \(F_{ST} = 0.12\) does and does not
mean}{What F_(ST) = 0.12 does and does not mean}}\label{what-f_st-0.12-does-and-does-not-mean}

\textbf{It DOES mean}: A small but real fraction of human genetic
variation is partitioned between populations. Allele frequencies do
differ --- for some loci substantially. Genes under local selection
(skin pigmentation: SLC24A5, MC1R; lactase persistence: LCT/MCM6;
high-altitude adaptation: EPAS1) show \textbf{far higher locus-specific
\(F_{ST}\) values} (0.4--0.9) than the genome-wide average, reflecting
the differential selective pressures of latitude, diet, altitude, and
pathogens.

\textbf{It does NOT mean}: Human populations are sharply discrete,
biologically distinct ``races.'' The reverse: 88\% of variation is
shared \textbf{within} any given population, meaning two random
individuals from the same continental population are nearly as
genetically different as two random individuals from different
continents. The continuous-variation pattern dominates: human variation
is graded geographically, with no sharp boundaries between any defined
groups (the \textbf{clinal} pattern). Lewontin's 1972 observation has
been confirmed by every subsequent dataset, including whole-genome
sequencing of millions of individuals.

\paragraph{\texorpdfstring{\(F_{ST}\) as a tool for detecting
selection}{F_(ST) as a tool for detecting selection}}\label{f_st-as-a-tool-for-detecting-selection}

Because population-genetic models predict the \textbf{genome-wide}
\(F_{ST}\) distribution under neutral evolution, \textbf{outlier loci}
with unusually high \(F_{ST}\) can flag positive selection at specific
genes. The lactase persistence allele (rs4988235) shows
\(F_{ST} \approx 0.7\) between Northern Europeans (allele frequency
about 0.7) and East Asians (frequency about 0.01) ---
a textbook example of recent strong directional selection driving
population differentiation. Similarly, EDAR (East Asian hair texture,
sweat glands), DARC (West African malaria resistance), and SLC24A5
(light skin in Europeans) show \(F_{ST}\) outliers consistent with
population-specific selective sweeps.

\textbf{Isolation by distance}: \(F_{ST}\) increases with geographic
distance. \textbf{Ring species} (e.g., Ensatina salamanders around
California's Central Valley) demonstrate that populations at the
endpoints of a geographic ring may be reproductively isolated despite
continuous gene flow along the ring. Within human populations,
\(F_{ST}\) correlates approximately linearly with geographic distance up
to 5,000 km, a signature of the stepping-stone migration pattern out of
Africa.

\begin{center}\rule{0.5\linewidth}{0.5pt}\end{center}

\subsection{Mutation}\label{mutation}

\subsubsection{Mutation-Selection
Balance}\label{mutation-selection-balance}

For a deleterious recessive allele maintained by mutation-selection
balance:

Forward mutation rate \texorpdfstring{\ensuremath{\mu}}{mu} (A -\textgreater{} a) introduces new copies;
selection removes them:

\begin{equation}
\hat{q} = \sqrt{\frac{\mu}{s}} \quad \text{(for completely recessive, s = 1 for lethals)}
\label{eq:population_genetics_31}
\end{equation}

\subsection{Worked Example: Mutation-Selection
Balance}\label{worked-example-mutation-selection-balance}

For a recessive lethal allele with \(\mu = 10^{-5}\):

\begin{equation}
\hat{q} = \sqrt{\frac{10^{-5}}{1}} = 0.00316
\label{eq:population_genetics_32}
\end{equation}

\begin{equation}
\text{Disease frequency} = q^2 = 10^{-5} = 1/100,000
\label{eq:population_genetics_33}
\end{equation}

\begin{equation}
\text{Carrier frequency} = 2pq \approx 2(0.00316) = 0.00632 \approx 1/158
\label{eq:population_genetics_34}
\end{equation}

For \textbf{CF} (\(s \approx 1\), \(q \approx 0.02\)): Implied
\(\mu \approx q^2 s = (0.02)^2(1) = 4 \times 10^{-4}\). This is much
higher than the typical point mutation rate (\(\sim 10^{-8}\) per
\hyperref[gl:nucleotide]{\textbf{nucleotide}} per generation),
suggesting that \textbf{heterozygote advantage} or other selective
forces maintain CF carrier frequency above what mutation-selection
balance alone would predict.

For a deleterious \textbf{dominant} allele:

\begin{equation}
\hat{q} \approx \frac{\mu}{hs}
\label{eq:population_genetics_35}
\end{equation}

Since selection acts on heterozygotes (who are much more common than
rare homozygotes), dominant alleles reach much lower equilibrium
frequencies.

\begin{center}\rule{0.5\linewidth}{0.5pt}\end{center}

\subsection{Molecular Evolution}\label{molecular-evolution}

\subsubsection{\texorpdfstring{Neutral Theory
\citep{kimura1968}}{Neutral Theory}}\label{neutral-theory}

Motoo Kimura's \textbf{neutral theory of molecular evolution} proposes
that most molecular variation within and between species is selectively
\textbf{neutral} -- not affected by natural selection, but evolving by
genetic drift.

\textbf{Key predictions:} 1. The rate of neutral substitution = the
neutral mutation rate \(\mu_n\), \textbf{independent of population size}
- Proof: Rate of new neutral mutations = \(2N \times \mu_n\);
probability each fixes = \(\frac{1}{2N}\); substitution rate =
\(2N\mu_n \times \frac{1}{2N} = \mu_n\) 2. Heterozygosity at equilibrium
= \(\frac{4N_e\mu}{1 + 4N_e\mu}\) (for diploid, biallelic loci) 3.
Synonymous substitution rate (\(K_S\)) should be approximately constant
across genes (molecular clock)

\subsubsection{The Nearly Neutral
Theory}\label{the-nearly-neutral-theory}

Tomoko \citet{ohta1973} extended neutral theory: mutations with
selection coefficients \(|s| < \frac{1}{2N_e}\) behave as if neutral
(drift overwhelms selection). In large populations, weakly deleterious
mutations are efficiently removed; in small populations, they can drift
to fixation. This predicts slightly higher substitution rates in small
populations -- observed in island species and endangered populations.

\subsubsection{Ka/Ks (dN/dS) Ratio}\label{kaks-dnds-ratio}

The ratio of nonsynonymous (amino acid-changing, Ka or dN) to synonymous
(silent, Ks or dS) substitution rates:

{\def\LTcaptype{none} % do not increment counter
\begin{longtable}[]{@{}
  >{\raggedright\arraybackslash}p{(\linewidth - 4\tabcolsep) * \real{0.2258}}
  >{\raggedright\arraybackslash}p{(\linewidth - 4\tabcolsep) * \real{0.4839}}
  >{\raggedright\arraybackslash}p{(\linewidth - 4\tabcolsep) * \real{0.2903}}@{}}
\toprule\noalign{}
\begin{minipage}[b]{\linewidth}\raggedright
Ka/Ks
\end{minipage} & \begin{minipage}[b]{\linewidth}\raggedright
Interpretation
\end{minipage} & \begin{minipage}[b]{\linewidth}\raggedright
Example
\end{minipage} \\
\midrule\noalign{}
\endhead
\bottomrule\noalign{}
\endlastfoot
\textless\textless{} 1 & \textbf{Purifying (negative) selection} & Most
genes; \hyperref[gl:histone]{\textbf{histone}}s (Ka/Ks about 
0.005) \\
= 1 & \textbf{Neutral evolution} & Pseudogenes \\
\textgreater\textgreater{} 1 & \textbf{Positive (adaptive) selection} &
MHC antigen-binding sites; reproductive
\hyperref[gl:protein]{\textbf{protein}}s (e.g., bindin in sea
urchins) \\
\end{longtable}
}

\textbf{Detecting positive selection}: Methods like the
McDonald-Kreitman test compare within-species polymorphism to
between-species divergence at synonymous and nonsynonymous sites. An
excess of nonsynonymous divergence relative to polymorphism indicates
positive selection.

\begin{center}\rule{0.5\linewidth}{0.5pt}\end{center}

\subsection{Population Structure}\label{population-structure}

\subsubsection{Wright's F-Statistics}\label{wrights-f-statistics}

Wright defined three hierarchical F-statistics:

\begin{itemize}
\tightlist
\item
  \textbf{\(F_{IS}\)} (inbreeding coefficient of individuals relative to
  subpopulation): Measures departure from random mating within a
  subpopulation (excess or deficiency of heterozygotes)
\item
  \textbf{\(F_{ST}\)} (fixation index of subpopulations relative to
  total): Measures genetic differentiation among subpopulations
\item
  \textbf{\(F_{IT}\)} (inbreeding coefficient of individuals relative to
  total): Overall departure from HWE in the total population
\end{itemize}

Relationship: \((1 - F_{IT}) = (1 - F_{IS})(1 - F_{ST})\)

\subsubsection{STRUCTURE Analysis}\label{structure-analysis}

The program STRUCTURE (Pritchard et al., 2000) uses Bayesian clustering
to assign individuals to \(K\) populations based on multilocus genotype
data. It has been widely used in human population genetics to identify
population structure (though the number of clusters \(K\) is a modeling
choice, not a biological fact -- an important caveat against
over-interpretation).

\begin{center}\rule{0.5\linewidth}{0.5pt}\end{center}

\subsection{Coalescent Theory}\label{coalescent-theory}

\subsubsection{The Coalescent Framework}\label{the-coalescent-framework}

Rather than tracing alleles forward in time (classical population
genetics), \textbf{coalescent theory} \citep{kingman1982} traces allele
lineages \textbf{backward in time} to their \textbf{most recent common
ancestor (MRCA)}.

\textbf{Key results:} - For 2 alleles in a diploid population of size
\(N\): Expected coalescence time = \(2N_e\) generations - For \(k\)
alleles: Expected total time to MRCA = \(4N_e(1 - 1/k)\) generations -
As \(k \to \infty\): TMRCA \(\to 4N_e\) generations

\textbf{Mitochondrial Eve}: The TMRCA for most human mitochondrial DNA
lineages is estimated at about 150,000-200,000 years ago. This
does NOT mean there was a single woman alive at that time -- it means
most other maternal lineages have gone extinct by chance (drift). Many
women were alive, but a single lineage survived to the present through
unbroken maternal descent.

\textbf{Y-chromosomal Adam}: The TMRCA for most human Y
\hyperref[gl:chromosome]{\textbf{chromosome}}s is estimated at
about 200,000-300,000 years ago. Same caveat applies.

\subsubsection{Molecular Clocks}\label{molecular-clocks}

If neutral substitutions accumulate at a roughly constant rate
(\(\mu_n\)), sequence divergence between species provides a ``clock''
for estimating divergence times.

\textbf{Jukes-Cantor correction} for observed divergence \(d\):

\begin{equation}
d_{JC} = -\frac{3}{4}\ln\left(1 - \frac{4d}{3}\right)
\label{eq:population_genetics_36}
\end{equation}

This corrects for multiple substitutions at the same site (back
mutations, parallel mutations, convergent mutations).

\textbf{Divergence time}:

\begin{equation}
t = \frac{d_{JC}}{2\mu}
\label{eq:population_genetics_37}
\end{equation}

The factor of 2 accounts for divergence along both lineages since the
split.

\subsection{Worked Example: Molecular Clock
Estimate}\label{worked-example-molecular-clock-estimate}

Human-chimpanzee cytochrome c has roughly 1\% sequence divergence; use a
neutral mutation rate for this gene of approximately
\(2 \times 10^{-9}\) substitutions/site/year.

\begin{equation}
d_{JC} = -\frac{3}{4}\ln\left(1 - \frac{4(0.01)}{3}\right) = -0.75 \times \ln(0.9867) = -0.75 \times (-0.01339) = 0.01004
\label{eq:population_genetics_38}
\end{equation}

\begin{equation}
t = \frac{0.01004}{2 \times 2 \times 10^{-9}} = \frac{0.01004}{4 \times 10^{-9}} = 2.51 \times 10^6 \text{ years} \approx 2.5 \text{ Ma}
\label{eq:population_genetics_39}
\end{equation}

(Note: The actual human-chimp divergence is estimated at 6-7 Ma using
whole-genome data and calibrated clocks with fossil constraints.
Cytochrome c evolves slowly and gives an underestimate.)

\textbf{Calibrated molecular clocks using fossil data:}

{\def\LTcaptype{none} % do not increment counter
\begin{longtable}[]{@{}ll@{}}
\toprule\noalign{}
Comparison & Estimated Divergence \\
\midrule\noalign{}
\endhead
\bottomrule\noalign{}
\endlastfoot
Human-Chimpanzee & 6-7 Ma \\
Human-Gorilla & 8-10 Ma \\
Human-Orangutan & 14-16 Ma \\
Human-Old World Monkeys & 25-30 Ma \\
Human-Mouse & 75-80 Ma \\
Human-Chicken & about 310 Ma \\
\end{longtable}
}

\begin{quote}
\textbf{Clinical Connection: Molecular Clock and HIV Forensics} The
molecular clock has been used in HIV forensic cases. In the 1998
Louisiana case (State of Louisiana v. Richard Schmidt), phylogenetic
analysis of HIV sequences showed that the viral strains in the victim
and the defendant's HIV-positive patient were more closely related to
each other than to any other strains, consistent with direct
transmission. Coalescent analysis estimated the transmission date,
providing critical forensic evidence.
\end{quote}

\textbf{Concept Check 18.3}

\begin{quote}
\begin{enumerate}
\def\labelenumi{\arabic{enumi}.}
\tightlist
\item
  What is the probability that a new neutral mutation in a diploid
  population of \(N_e = 5,000\) will eventually fix? How many
  generations would this take on average?
\item
  Explain why synonymous substitution rates are approximately constant
  across genes (molecular clock) while nonsynonymous rates vary widely.
\item
  ``Mitochondrial Eve'' lived about 200,000 years ago. Does
  this mean she was the sole woman alive at that time? Explain.
\item
  Why does the Jukes-Cantor correction give a larger value than the raw
  observed divergence?
\end{enumerate}
\end{quote}

\begin{center}\rule{0.5\linewidth}{0.5pt}\end{center}

\subsection{Worked Example: Comprehensive Population
Genetics}\label{worked-example-comprehensive-population-genetics}

\textbf{Problem}: In a population of 10,000 grasshoppers, you observe
the following genotype counts at a color locus with codominant alleles:

{\def\LTcaptype{none} % do not increment counter
\begin{longtable}[]{@{}ll@{}}
\toprule\noalign{}
Genotype & Count \\
\midrule\noalign{}
\endhead
\bottomrule\noalign{}
\endlastfoot
RR (red) & 4,200 \\
Rr (orange) & 4,800 \\
rr (yellow) & 1,000 \\
\end{longtable}
}

\begin{enumerate}
\def\labelenumi{(\alph{enumi})}
\tightlist
\item
  Calculate allele frequencies.
\end{enumerate}

\begin{equation}
p(R) = \frac{2(4200) + 4800}{20000} = \frac{13200}{20000} = 0.66
\label{eq:population_genetics_40}
\end{equation}

\begin{equation}
q(r) = \frac{2(1000) + 4800}{20000} = \frac{6800}{20000} = 0.34
\label{eq:population_genetics_41}
\end{equation}

\begin{enumerate}
\def\labelenumi{(\alph{enumi})}
\setcounter{enumi}{1}
\tightlist
\item
  Is this population in HWE?
\end{enumerate}

Expected under HWE: - \(E(RR) = (0.66)^2 \times 10000 = 4356\) -
\(E(Rr) = 2(0.66)(0.34) \times 10000 = 4488\) -
\(E(rr) = (0.34)^2 \times 10000 = 1156\)

\begin{equation}
\chi^2 = \frac{(4200-4356)^2}{4356} + \frac{(4800-4488)^2}{4488} + \frac{(1000-1156)^2}{1156}
\label{eq:population_genetics_42}
\end{equation}

\begin{equation}
= \frac{24336}{4356} + \frac{97344}{4488} + \frac{24336}{1156} = 5.59 + 21.69 + 21.05 = 48.33
\label{eq:population_genetics_43}
\end{equation}

\(\chi^2_{crit}(1 \text{ df}, \alpha = 0.05) = 3.841\).

Since \(48.33 >> 3.841\), we \textbf{reject HWE}. This population is NOT
in Hardy-Weinberg equilibrium. The excess of heterozygotes (observed
4,800 vs.~expected 4,488) could indicate \textbf{heterozygote advantage}
(balancing selection) or \textbf{negative assortative mating}
(like-unlike pairing).

\begin{center}\rule{0.5\linewidth}{0.5pt}\end{center}

\subsection{Computational Bridge}\label{computational-bridge-17}

Selection recursions used in the chapter are iterated in
\breaktt{simulate\_selection}:

\begin{Shaded}
\begin{Highlighting}[]
\ImportTok{from}\NormalTok{ biology.evolution }\ImportTok{import}\NormalTok{ Population, simulate\_selection}

\NormalTok{pop }\OperatorTok{=}\NormalTok{ Population(}\StringTok{"demo"}\NormalTok{, p}\OperatorTok{=}\FloatTok{0.2}\NormalTok{, q}\OperatorTok{=}\FloatTok{0.8}\NormalTok{, fitness\_AA}\OperatorTok{=}\FloatTok{1.0}\NormalTok{, fitness\_Aa}\OperatorTok{=}\FloatTok{0.95}\NormalTok{, fitness\_aa}\OperatorTok{=}\FloatTok{0.5}\NormalTok{)}
\NormalTok{traj }\OperatorTok{=}\NormalTok{ simulate\_selection(pop, }\DecValTok{30}\NormalTok{)}
\BuiltInTok{print}\NormalTok{(}\BuiltInTok{round}\NormalTok{(traj[}\OperatorTok{{-}}\DecValTok{1}\NormalTok{].p, }\DecValTok{4}\NormalTok{))}
\end{Highlighting}
\end{Shaded}

\begin{quote}
\textbf{Clinical / systems note:} Pathogen resistance evolution (TB,
HIV, malaria) is tracked with the same allele-frequency dynamics;
combination therapy aims to prevent peaks of \(p\) for any single
resistance mutation.
\end{quote}

\begin{center}\rule{0.5\linewidth}{0.5pt}\end{center}

\subsection{Current Evidence and Frontier
Biology}\label{current-evidence-and-frontier-biology-20}

For \textbf{Population Genetics and Hardy-Weinberg Equilibrium},
frontier biology belongs inside the evidence logic of the chapter.
Classical genetics remains essential, but modern interpretation adds
penetrance, polygenicity, structural variation, ancestry-aware
inference, and uncertainty in risk prediction. The core reading question
is this: allele-frequency explanations should name the force, parameter
values, assumptions, and data needed to distinguish forces.

\begin{itemize}
\tightlist
\item
  \textbf{What to verify:} identify the observation, model, assay, or
  dataset that would make the claim stronger or weaker.
\item
  \textbf{What to qualify:} state the scale, organism, cell type,
  environmental condition, or population where the claim is expected to
  hold.
\item
  \textbf{What to compare:} test at least one alternative explanation,
  baseline, or null model before treating the pattern as causal.
\item
  \textbf{What to cite:} distinguish primary evidence, review synthesis,
  public dataset, and institutional guidance; for recent or numeric
  claims, prefer the source closest to the measurement and state what
  has changed since it was published.
\end{itemize}

A good genetics answer separates the Mendelian transmission model from
the evidence needed to use it in a population, family, or clinical
setting.

\textbf{Source practice:} For inheritance and population claims,
separate the model assumptions from sampling, ancestry representation,
penetrance, linkage, and environment.

\subsubsection{Current Evidence Map: Ancestry-Aware Variant
Interpretation}\label{current-evidence-map-ancestry-aware-variant-interpretation}

\begin{figure}[htbp]
\centering
\includegraphics[width=0.82\linewidth,height=4.2in,keepaspectratio]{../figures/mermaid_inline/inline_mermaid_0102_a237f5df08ad.png}
\caption{Graph references can reduce reference bias, but population-genetic interpretation still depends on sampling design, assumptions, and validation \textbackslash{}citep\{humanpangenome2023\}.}
\end{figure}

\subsubsection{Forensic and Population-Inference
Boundaries}\label{forensic-and-population-inference-boundaries}

Population-genetic statistics are summaries of samples, not identities.
A PCA cluster, ancestry proportion, or \(F_{ST}\) estimate depends on
the reference panel, marker set, filtering, and demographic model;
changing any of these can shift the interpretation. In forensic or
outbreak settings, phylogenetic closeness can support a transmission
hypothesis, but direction and responsibility require timing, sampling
density, chain-of-custody, and plausible exposure pathways. Good
practice reports the null model, the comparison database, the
uncertainty interval, and the alternative histories still compatible
with the data.

\subsection{Summary}\label{summary-20}

\begin{itemize}
\tightlist
\item
  \textbf{Hardy-Weinberg equilibrium}: \(p^2 + 2pq + q^2 = 1\) in large,
  randomly mating populations with no mutation, migration, or selection.
  Reached in one generation. Serves as a null model for detecting
  evolutionary forces.
\item
  \textbf{Natural selection}: Fitness (\(w\)), selection coefficient
  (\(s\)), dominance coefficient (\(h\)). Directional selection shifts
  allele frequencies; balancing selection (heterozygote advantage)
  maintains polymorphism at stable equilibrium
  (\(\hat{q} = s_1/(s_1 + s_2)\)).
\item
  \textbf{Genetic drift}: Random allele frequency changes in finite
  populations; variance \(\propto 1/2N_e\); fixation or loss inevitable.
  Bottleneck and founder effects cause non-representative gene pools.
\item
  \textbf{Gene flow}: Homogenizes populations;
  \(\Delta p_1 = m(p_2 - p_1)\). \(F_{ST}\) measures population
  differentiation (about 0.12 for humans).
\item
  \textbf{Mutation-selection balance}: \(\hat{q} = \sqrt{\mu/s}\) for
  deleterious recessives. CF carrier frequency suggests heterozygote
  advantage beyond mutation alone.
\item
  \textbf{Neutral theory} (Kimura): Most molecular variation is neutral;
  substitution rate = neutral mutation rate (independent of \(N\)).
  Ka/Ks ratio distinguishes purifying, neutral, and positive selection.
\item
  \textbf{Coalescent theory}: Traces lineages backward; TMRCA for
  mitochondrial DNA about 200,000 years; for Y chromosome
  about 200,000-300,000 years.
\item
  \textbf{Molecular clock}: \(t = d_{JC}/2\mu\); calibrated with
  fossils; widely used for estimating divergence times. Human-chimp
  about 6-7 Ma.
\item
  \textbf{Connections:} See Unit VI (tree thinking and substitution
  models), Unit X (metapopulation structure), Unit VII (microbial
  effective population sizes).
\end{itemize}

\begin{center}\rule{0.5\linewidth}{0.5pt}\end{center}

\subsection{Review Questions}\label{review-questions-20}

\begin{enumerate}
\def\labelenumi{\arabic{enumi}.}
\tightlist
\item
  Derive the Hardy-Weinberg equation from first principles, starting
  with random mating.
\item
  A population has the following genotype counts: AA = 200, Aa = 500, aa
  = 300. Calculate allele frequencies and test for HWE.
\item
  Sickle cell trait (HbAS) provides about 10-fold protection
  against severe \emph{P. falciparum} malaria. If \(s_1 = 0.10\) (cost
  of malaria susceptibility in HbAA) and \(s_2 = 0.85\) (cost of sickle
  cell disease in HbSS), calculate the equilibrium frequency of HbS.
\item
  A population of 100 individuals has \(p = 0.5\). Calculate the
  expected variance in \(\Delta p\) per generation due to drift. What is
  the expected heterozygosity loss per generation?
\item
  Explain why \(N_e\) is usually much less than census size \(N\). Give
  three factors that reduce \(N_e\).
\item
  A population on a remote island was founded by 10 individuals 200
  years ago. Predict the genetic consequences using the founder effect
  concept.
\item
  Compare directional, stabilizing, disruptive, and balancing selection.
  Give an example of each.
\item
  Using the Jukes-Cantor correction, calculate the evolutionary distance
  for two sequences with 20\% observed divergence.
\item
  Explain why the neutral theory predicts that substitution rate equals
  mutation rate regardless of population size. Show the mathematical
  proof.
\item
  What is \(F_{ST}\)? Calculate \(F_{ST}\) for two populations with
  allele frequencies \(p_1 = 0.8\) and \(p_2 = 0.2\).
\item
  If \breaktt{simulate\_selection} is run for 200 generations with
  \(s=0.01\) against the \(a\) allele, when does \(\Delta p\) per
  generation approximately linearise vs.~saturate?
\item
  Explain why \textbf{soft sweeps} complicate signatures of selection
  relative to classic hard sweeps.
\end{enumerate}

\begin{center}\rule{0.5\linewidth}{0.5pt}\end{center}

\subsection{Further Reading and Source
Notes}\label{further-reading-and-source-notes-20}

\begin{itemize}
\tightlist
\item
  Weinberg (1908). \{"U\}ber den Nachweis der Vererbung beim Menschen.
  \emph{Jahreshefte des Vereins f\{"u\}r vaterl\{"a\}ndische Naturkunde
  in W\{"u\}rttemberg}, 64.
\item
  Kimura (1968). Evolutionary Rate at the Molecular Level.
  \emph{Nature}, 217.
\item
  Ohta (1973). Slightly deleterious and nearly neutral mutations in
  molecular evolution. \emph{Proceedings of the National Academy of
  Sciences}, 70.
\item
  Kingman (1982). The coalescent. \emph{Stochastic Processes and their
  Applications}, 13.
\end{itemize}

\begin{center}\rule{0.5\linewidth}{0.5pt}\end{center}

\subsection{Key Terms}\label{key-terms-20}

\begin{enumerate}
\def\labelenumi{\arabic{enumi}.}
\tightlist
\item
  \textbf{Gene pool} -- the allele set carried by a population across
  its loci
\item
  \textbf{Allele frequency} -- the proportion of sampled gene copies
  carrying a specified allele
\item
  \textbf{Hardy-Weinberg equilibrium} -- state in which allele and
  genotype frequencies remain constant across generations
\item
  \textbf{Fitness (w)} -- relative reproductive success of a genotype
\item
  \textbf{Selection coefficient (s)} -- magnitude of selection against a
  genotype (\(s = 1 - w\))
\item
  \textbf{Heterozygote advantage (overdominance)} -- heterozygote has
  higher fitness than either homozygote
\item
  \textbf{Genetic drift} -- random changes in allele frequency due to
  finite population size
\item
  \textbf{Effective population size (\(N_e\))} -- size of an ideal
  population experiencing the same drift as the actual population
\item
  \textbf{Bottleneck effect} -- drastic reduction in population size
  causing loss of genetic diversity
\item
  \textbf{Founder effect} -- reduced genetic diversity when a small
  group establishes a new population
\item
  \textbf{Gene flow} -- movement of alleles between populations via
  migration
\item
  \textbf{\(F_{ST}\) (fixation index)} -- measure of genetic
  differentiation among populations
\item
  \textbf{Mutation-selection balance} -- equilibrium allele frequency
  determined by mutation introducing and selection removing deleterious
  alleles
\item
  \textbf{Neutral theory} -- most molecular variation is selectively
  neutral and evolves by drift
\item
  \textbf{Ka/Ks ratio (dN/dS)} -- ratio of nonsynonymous to synonymous
  substitution rates; indicates selection regime
\item
  \textbf{Coalescent theory} -- framework tracing allele lineages
  backward to their most recent common ancestor
\item
  \textbf{Molecular clock} -- assumption that neutral substitutions
  accumulate at a roughly constant rate
\item
  \textbf{Jukes-Cantor distance} -- correction for multiple
  substitutions at the same site:
  \(d_{JC} = -\frac{3}{4}\ln(1-\frac{4d}{3})\)
\item
  \textbf{TMRCA} -- time to the most recent common ancestor of a set of
  allele lineages
\end{enumerate}

\begin{center}\rule{0.5\linewidth}{0.5pt}\end{center}

\subsubsection{Companion Source
Module}\label{companion-source-module-20}

\textbf{Population Genetics and Hardy-Weinberg Equilibrium} should leave
a reproducible trail from a biological claim to the code, figure,
diagram, or paper-based activity that can test it. Use the surfaces
below to inspect the chapter's assumptions, rerun the relevant model, or
compare the manuscript explanation with companion labs and figures.

{\def\LTcaptype{none} % do not increment counter
\begin{longtable}[]{@{}
  >{\raggedright\arraybackslash}p{(\linewidth - 2\tabcolsep) * \real{0.5000}}
  >{\raggedright\arraybackslash}p{(\linewidth - 2\tabcolsep) * \real{0.5000}}@{}}
\toprule\noalign{}
\begin{minipage}[b]{\linewidth}\raggedright
Surface
\end{minipage} & \begin{minipage}[b]{\linewidth}\raggedright
Use it for
\end{minipage} \\
\midrule\noalign{}
\endhead
\bottomrule\noalign{}
\endlastfoot
\breaktt{src/biology/genetics/genetics.py} (\texttt{hardy\_weinberg},
\texttt{chi\_squared\_test}, \breaktt{jukes\_cantor\_distance}) & Test
equilibrium, goodness of fit, and molecular-distance assumptions. \\
\breaktt{src/biology/evolution/evolution.py}
(\breaktt{simulate\_selection}, \breaktt{wright\_fisher\_drift},
\breaktt{molecular\_clock\_divergence\_time}) & Compare deterministic and
stochastic allele-frequency change. \\
\breaktt{src/visualization/plots.py}
(\breaktt{plot\_selection\_simulation}) & Inspect trajectories and
sampling effects. \\
\end{longtable}
}

\textbf{Reproducibility check:} state population size, mating model,
selection coefficient, migration, mutation, and sampling uncertainty
before attributing allele-frequency change. \textbf{Cross-reference:}
connect with \cref{sec:unit_V_mendelian_genetics},
\cref{sec:unit_VI_evolution_and_selection}, and
\cref{sec:unit_VI_genetic_drift_and_speciation}.

\newpage

\section*{Unit VI --- Evolution:
Introduction}\label{unit-vi-evolution-introduction}
\addcontentsline{toc}{section}{Unit VI --- Evolution: Introduction}

\subsection*{Why This Unit Matters}\label{why-this-unit-matters-6}
\addcontentsline{toc}{subsection}{Why This Unit Matters}

On 24 November 1859, Charles Darwin published \emph{On the Origin of
Species by Means of Natural Selection}. The book sold out on its first
day. Within a decade, the concept of evolution by natural selection had
overturned centuries of natural theology and provided biology with its
first unifying theory. Theodosius Dobzhansky's 1973 essay captured this
with a sentence that has become a biologist's creed: \emph{"Nothing in
biology makes sense except in the light of evolution."}

Darwin's central insight --- that heritable variation + differential
reproductive success produces cumulative directional change in
populations over generations --- was elegant but incomplete. He had no
mechanism for inheritance. The Modern Synthesis (1930s--1940s), forged
by Fisher, Haldane, Wright, Dobzhansky, and Mayr, married Darwinian
selection to Mendelian genetics and population genetics, giving
evolution a precise mathematical foundation. A further deepening came
with the neutral theory (Kimura, 1968): not most genetic change is
driven by selection; much molecular evolution is stochastic, driven by
genetic drift. Understanding evolution now requires integrating
selection, drift, mutation, gene flow, and non-random mating --- the
five forces that perturb Hardy-Weinberg equilibrium.

This unit treats evolution quantitatively. You will simulate allele
frequency trajectories under selection and drift, derive the molecular
clock equation and apply it to evolutionary divergence times, model
speciation as a bifurcation in population structure, and reconstruct
phylogenetic trees using maximum-likelihood and Bayesian methods. The
tools here are the same used in clinical epidemiology (tracking
antibiotic resistance alleles), crop improvement (genomic selection),
and forensic genetics (ancestral inference from SNP profiles).

\begin{center}\rule{0.5\linewidth}{0.5pt}\end{center}

\subsection*{Landmark Discoveries}\label{landmark-discoveries-6}
\addcontentsline{toc}{subsection}{Landmark Discoveries}

{\def\LTcaptype{none} % do not increment counter
\begin{longtable}[]{@{}
  >{\raggedright\arraybackslash}p{(\linewidth - 8\tabcolsep) * \real{0.2407}}
  >{\raggedright\arraybackslash}p{(\linewidth - 8\tabcolsep) * \real{0.0741}}
  >{\raggedright\arraybackslash}p{(\linewidth - 8\tabcolsep) * \real{0.2963}}
  >{\raggedright\arraybackslash}p{(\linewidth - 8\tabcolsep) * \real{0.1667}}
  >{\raggedright\arraybackslash}p{(\linewidth - 8\tabcolsep) * \real{0.2222}}@{}}
\toprule\noalign{}
\begin{minipage}[b]{\linewidth}\raggedright
Discoverer(s)
\end{minipage} & \begin{minipage}[b]{\linewidth}\raggedright
Year
\end{minipage} & \begin{minipage}[b]{\linewidth}\raggedright
Journal / Source
\end{minipage} & \begin{minipage}[b]{\linewidth}\raggedright
Discovery
\end{minipage} & \begin{minipage}[b]{\linewidth}\raggedright
Significance
\end{minipage} \\
\midrule\noalign{}
\endhead
\bottomrule\noalign{}
\endlastfoot
Darwin \& Wallace & 1858--59 & \emph{Proc. Linn. Soc.}; \emph{On the
Origin of Species} & Natural selection as mechanism of evolution &
Provided the first mechanistic explanation for biological diversity \\
Gregor Mendel (rediscovered) & 1900 & de Vries, Correns, Tschermak &
Mendelian inheritance = mechanism for variation & Fused genetics with
evolution; launched Modern Synthesis \\
Ronald Fisher & 1930 & \emph{The Genetical Theory of Natural Selection}
& Mathematical synthesis of Darwinism + Mendelism & Fundamental theorem
of natural selection; \(\Delta q = spq^2/\bar{W}\) \\
Sewall Wright & 1931 & \emph{Genetics} & Genetic drift and adaptive
landscapes & Showed small populations diverge by chance;
\(\sigma^2(\Delta p) = pq/2N_e\) \\
Motoo Kimura & 1968 & \emph{Nature} & Neutral theory of molecular
evolution & Most amino acid substitutions are selectively neutral;
molecular clock \\
Woese \& Fox & 1977 & \emph{Proc. Natl. Acad. Sci.} & Three domains of
life (Archaea identified by rRNA) & Completely restructured the
comprehensive tree of life \\
Svante Pääbo et al. & 2010 & \emph{Science} & Neanderthal genome
sequencing & Showed interbreeding between modern humans and
Neanderthals; Nobel Prize 2022 \\
\end{longtable}
}

\begin{center}\rule{0.5\linewidth}{0.5pt}\end{center}

\subsection*{Key Concepts and
Connections}\label{key-concepts-and-connections-6}
\addcontentsline{toc}{subsection}{Key Concepts and Connections}

\begin{figure}[htbp]
\centering
\includegraphics[width=0.82\linewidth,height=4.2in,keepaspectratio]{../figures/mermaid_inline/inline_mermaid_0103_2f6f54b51fdb.png}
\caption{Unit VI concept map - Evolution. Red = selection; blue = drift; green = phylogenetics; orange = molecular clock.}
\end{figure}

\begin{center}\rule{0.5\linewidth}{0.5pt}\end{center}

\subsection*{Current Evidence Thread}\label{current-evidence-thread-6}
\addcontentsline{toc}{subsection}{Current Evidence Thread}

Treat this unit as a converging body of evidence for evolution: the
fossil record documenting transitional forms and deep time, comparative
genomics revealing shared ancestry and the molecular footprints of
selection, real-time selection measured in experimental and wild
populations, and phylogenies that reconstruct the branching history
those data imply. Evolutionary claims are strongest when they combine
mechanism, comparative evidence, population process, and explicit
uncertainty. As you move through the chapters, keep a two-column note:
\textbf{claim} on the left, \textbf{evidence that would change my
confidence} on the right. By the end of the unit, each major idea should
be tied to a measurement, model, citation, or paper-based lab decision.

\subsection*{Chapter Roadmap}\label{chapter-roadmap-7}
\addcontentsline{toc}{subsection}{Chapter Roadmap}

{\def\LTcaptype{none} % do not increment counter
\begin{longtable}[]{@{}
  >{\raggedright\arraybackslash}p{(\linewidth - 6\tabcolsep) * \real{0.1556}}
  >{\raggedright\arraybackslash}p{(\linewidth - 6\tabcolsep) * \real{0.1111}}
  >{\raggedright\arraybackslash}p{(\linewidth - 6\tabcolsep) * \real{0.2889}}
  >{\raggedright\arraybackslash}p{(\linewidth - 6\tabcolsep) * \real{0.4444}}@{}}
\toprule\noalign{}
\begin{minipage}[b]{\linewidth}\raggedright
Chapter
\end{minipage} & \begin{minipage}[b]{\linewidth}\raggedright
Title
\end{minipage} & \begin{minipage}[b]{\linewidth}\raggedright
Core Question
\end{minipage} & \begin{minipage}[b]{\linewidth}\raggedright
Key Equation / Model
\end{minipage} \\
\midrule\noalign{}
\endhead
\bottomrule\noalign{}
\endlastfoot
\textbf{19} & Evolution and Natural Selection & How does natural
selection change allele frequencies, and how fast? &
\(\Delta q = -spq^2/\bar{W}\); fitness surface; selection coefficient \\
\textbf{20} & Genetic Drift and Speciation & How do chance events in
small populations lead to divergence and new species? &
\(\sigma^2(\Delta p) = pq/2N_e\); isolation index; speciation models \\
\textbf{21} & Phylogenetics and the Tree of Life & How do we reconstruct
evolutionary history from molecular data? & Jukes-Cantor correction;
molecular clock: \(t = d_{JC}/2\mu\) \\
\end{longtable}
}

\begin{center}\rule{0.5\linewidth}{0.5pt}\end{center}

\subsection*{Connections Across the
Textbook}\label{connections-across-the-textbook-6}
\addcontentsline{toc}{subsection}{Connections Across the Textbook}

\begin{itemize}
\tightlist
\item
  \textbf{Hardy-Weinberg equilibrium} in
  \cref{sec:unit_V_population_genetics} is the null model disrupted by
  most five evolutionary forces analysed here.
\item
  \textbf{Molecular clock} and \textbf{Jukes-Cantor distances} connect
  to Unit IV (mutation rates in DNA replication) and Unit VII (phylogeny
  of pathogens and antibiotic-resistance evolution).
\item
  \textbf{Speciation} links to Unit X (biogeography, island species-area
  relationship, MacArthur-Wilson model).
\item
  \textbf{Adaptive evolution} of immune genes connects to Unit IX (MHC
  diversity and pathogen immunity), and antibiotic resistance evolution
  motivates the clinical sections of Unit VII.
\end{itemize}

\begin{quote}
\textbf{Key vocabulary introduced here:} fitness, selection coefficient,
genetic drift, effective population size (\(N_e\)), founder effect,
bottleneck, molecular clock, phylogeny, clade, synapomorphy, allopatric
speciation, sympatric speciation, neutral theory, molecular systematics,
maximum likelihood, bootstrap support.
\end{quote}

\subsection*{Computational Toolbox --- Unit
VI}\label{computational-toolbox-unit-vi}
\addcontentsline{toc}{subsection}{Computational Toolbox --- Unit VI}

\begin{Shaded}
\begin{Highlighting}[]
\ImportTok{from}\NormalTok{ biology.evolution }\ImportTok{import}\NormalTok{ simulate\_selection, simulate\_drift}

\CommentTok{\# Natural selection on a beneficial mutation (s = 0.02, starting frequency p0 = 0.01)}
\ImportTok{from}\NormalTok{ biology.evolution }\ImportTok{import}\NormalTok{ Population}

\NormalTok{initial }\OperatorTok{=}\NormalTok{ Population(}\StringTok{"beneficial allele"}\NormalTok{, p}\OperatorTok{=}\FloatTok{0.01}\NormalTok{, q}\OperatorTok{=}\FloatTok{0.99}\NormalTok{, fitness\_AA}\OperatorTok{=}\FloatTok{1.02}\NormalTok{, fitness\_Aa}\OperatorTok{=}\FloatTok{1.01}\NormalTok{, fitness\_aa}\OperatorTok{=}\FloatTok{1.0}\NormalTok{)}
\NormalTok{trajectory }\OperatorTok{=}\NormalTok{ simulate\_selection(initial, generations}\OperatorTok{=}\DecValTok{200}\NormalTok{)}
\BuiltInTok{print}\NormalTok{(}\SpecialStringTok{f"Generation 1:   p = }\SpecialCharTok{\{}\NormalTok{trajectory[}\DecValTok{0}\NormalTok{]}\SpecialCharTok{.}\NormalTok{p}\SpecialCharTok{:.4f\}}\SpecialStringTok{"}\NormalTok{)}
\BuiltInTok{print}\NormalTok{(}\SpecialStringTok{f"Generation 100: p = }\SpecialCharTok{\{}\NormalTok{trajectory[}\DecValTok{99}\NormalTok{]}\SpecialCharTok{.}\NormalTok{p}\SpecialCharTok{:.4f\}}\SpecialStringTok{"}\NormalTok{)}
\BuiltInTok{print}\NormalTok{(}\SpecialStringTok{f"Generation 200: p = }\SpecialCharTok{\{}\NormalTok{trajectory[}\DecValTok{199}\NormalTok{]}\SpecialCharTok{.}\NormalTok{p}\SpecialCharTok{:.4f\}}\SpecialStringTok{"}\NormalTok{)}

\CommentTok{\# Genetic drift: small population bottleneck (Ne=50) vs. large (Ne=1000)}
\ImportTok{import}\NormalTok{ random}\OperatorTok{;}\NormalTok{ random.seed(}\DecValTok{42}\NormalTok{)}
\NormalTok{small\_pop }\OperatorTok{=}\NormalTok{ simulate\_drift(p}\OperatorTok{=}\FloatTok{0.5}\NormalTok{, N}\OperatorTok{=}\DecValTok{50}\NormalTok{, generations}\OperatorTok{=}\DecValTok{200}\NormalTok{)}
\NormalTok{large\_pop }\OperatorTok{=}\NormalTok{ simulate\_drift(p}\OperatorTok{=}\FloatTok{0.5}\NormalTok{, N}\OperatorTok{=}\DecValTok{1000}\NormalTok{, generations}\OperatorTok{=}\DecValTok{200}\NormalTok{)}
\BuiltInTok{print}\NormalTok{(}\SpecialStringTok{f"Ne=50  terminal p:  }\SpecialCharTok{\{}\NormalTok{small\_pop[}\OperatorTok{{-}}\DecValTok{1}\NormalTok{]}\SpecialCharTok{:.3f\}}\SpecialStringTok{  (fixed=}\SpecialCharTok{\{}\StringTok{\textquotesingle{}yes\textquotesingle{}} \ControlFlowTok{if}\NormalTok{ small\_pop[}\OperatorTok{{-}}\DecValTok{1}\NormalTok{] }\KeywordTok{in}\NormalTok{ (}\DecValTok{0}\NormalTok{,}\DecValTok{1}\NormalTok{) }\ControlFlowTok{else} \StringTok{\textquotesingle{}no\textquotesingle{}}\SpecialCharTok{\}}\SpecialStringTok{)"}\NormalTok{)}
\BuiltInTok{print}\NormalTok{(}\SpecialStringTok{f"Ne=1000 terminal p: }\SpecialCharTok{\{}\NormalTok{large\_pop[}\OperatorTok{{-}}\DecValTok{1}\NormalTok{]}\SpecialCharTok{:.3f\}}\SpecialStringTok{"}\NormalTok{)}
\CommentTok{\# Small populations show large random fluctuations; large populations stay near 0.5.}
\end{Highlighting}
\end{Shaded}

\begin{quote}
\textbf{Try it yourself:} Run \texttt{simulate\_drift} 100 times with
\texttt{N=20} and count how many runs fix the allele (p=1.0) versus lose
it (p=0.0). Theory predicts fixation = p\texorpdfstring{\textsubscript{0}}{0} = 0.5.
\end{quote}

\begin{center}\rule{0.5\linewidth}{0.5pt}\end{center}

\emph{Source note: the evolution module supports selection, drift,
fitness-landscape, molecular-clock, sequence-distance, and
isolation-index examples. Figures and Mermaid diagrams provide the
selection and phylogenetic visuals.}

\newpage

\section{Evolution --- Theory, Natural Selection, and
Adaptation}\label{evolution-theory-natural-selection-and-adaptation}

\label{sec:unit_VI_evolution_and_selection}

\begin{quote}
\textbf{Ch 19} \texorpdfstring{\ensuremath{\cdot}}{.} Level 2/3 \texorpdfstring{\ensuremath{\cdot}}{.} 60 min read \texorpdfstring{\ensuremath{\cdot}}{.} 75 min lecture \texorpdfstring{\ensuremath{\cdot}}{.}
Prerequisites: \cref{sec:unit_V_population_genetics}
\end{quote}

\subsection{Learning Objectives}\label{learning-objectives-21}

By the end of this chapter, you should be able to:

\begin{enumerate}
\def\labelenumi{\arabic{enumi}.}
\tightlist
\item
  Trace the historical development of evolutionary thought from Lamarck
  through Darwin \citep{darwin1858}, Wallace, and the Modern Synthesis.
\item
  Evaluate five independent lines of evidence supporting evolution and
  explain why they converge on the same conclusion.
\item
  Define fitness mathematically and distinguish directional,
  stabilizing, disruptive, sexual, kin, and frequency-dependent
  selection.
\item
  Explain adaptation, exaptation, and the constraints that limit what
  \hyperref[gl:natural-selection]{\textbf{natural selection}}
  \citep{williams1966} can achieve.
\item
  Distinguish microevolution from macroevolution and evaluate punctuated
  equilibrium versus phyletic gradualism.
\item
  Describe coevolutionary dynamics including arms races and the Red
  Queen hypothesis.
\item
  Calculate a selection coefficient from observed allele-frequency
  change and predict the phenotypic response to selection using the
  breeder's equation \(R = h^2 S\).
\item
  Analyse direct observations of evolution in real time --- the \emph{E.
  coli} long-term evolution experiment, Galápagos finch beak shifts, and
  cane-toad spatial sorting --- and explain how each demonstrates
  measurable change within human timescales.
\item
  Evaluate the ethical responsibilities raised when evolutionary theory
  becomes a design framework, including gain-of-function research, gene
  drives, and heritable genome editing.
\end{enumerate}

\subsubsection{Study Blueprint}\label{study-blueprint-21}

\begin{itemize}
\tightlist
\item
  \textbf{Big idea:} Natural selection is differential reproductive
  success acting on heritable variation in context.
\item
  \textbf{Core concepts:} variation, fitness, adaptation, selection
  coefficient.
\item
  \textbf{Framework alignment:} Vision \& Change: Evolution, Systems; AP
  Biology: Evolution, Systems Interactions; NGSS-style topics: Natural
  Selection and Evolution, Interdependent Relationships in Ecosystems.
\item
  \textbf{Model or quantitative lens:} Selection coefficient and
  allele-frequency trajectory calculations.
\item
  \textbf{Data skill:} Interpret fitness data and distinguish selection
  from other forces.
\item
  \textbf{Practice cadence:} Visual Representations, Statistical Tests
  and Data Analysis, Argumentation.
\item
  \textbf{Common misconception to repair:} Evolution is not
  goal-directed progress; it is local change in populations under
  constraints.
\item
  \textbf{Primary lab:} \cref{sec:lab_unit_VI_evolution_and_selection}.
\item
  \textbf{Question bank:} \cref{sec:q_unit_VI_evolution_and_selection}.
\item
  \textbf{Transfer task:} Transfer selection reasoning to antibiotics,
  pesticide resistance, cancer, or climate adaptation.
\item
  \textbf{Bridge to computation:}
  \breaktt{biology.evolution.evolution.simulate\_selection}.
\end{itemize}

\begin{center}\rule{0.5\linewidth}{0.5pt}\end{center}

\begin{quote}
\textbf{Opening Vignette --- Thirteen Finches, One Revelation}

In 1835, H.M.S. Beagle anchored in the Galápagos Islands for five weeks.
Charles Darwin explored several islands, collecting birds and noting
that their beaks differed island to island --- but he initially assumed
they were different species entirely. It was the ornithologist John
Gould who, examining Darwin's specimens back in London in 1837, declared
most 13 birds finches, each adapted to a different food source.
Primarily then did Darwin understand what he had seen: a single
ancestral population, blown across 900 km of ocean, had diversified into
13 species to exploit different ecological
\hyperref[gl:niche]{\textbf{niche}}s. The insight sparked his theory of
natural selection. A century and a half later, Peter and Rosemary Grant
spent 40 years documenting beak evolution in \emph{Geospiza fortis}
during drought years --- measuring beaks with calipers on every
individual --- providing the most comprehensive real-time record of
natural selection in a wild population ever assembled.
\end{quote}

\subsection{Historical Context of Evolutionary
Thought}\label{historical-context-of-evolutionary-thought}

\subsubsection{Pre-Darwinian Ideas}\label{pre-darwinian-ideas}

Long before Darwin, naturalists wrestled with the observation that
organisms seem exquisitely matched to their environments. Three key
thinkers set the stage for evolutionary biology.

\textbf{Jean-Baptiste Lamarck (1744--1829)} proposed two mechanisms of
evolutionary change in his \emph{Philosophie Zoologique} (1809). First,
the \textbf{law of use and disuse} held that organs used extensively
become stronger and larger, while unused organs deteriorate. Second, the
\textbf{inheritance of acquired characteristics} proposed that
modifications gained during an organism's lifetime pass to offspring. A
blacksmith's children, by this logic, would inherit strong arms. While
the inheritance mechanism was wrong, Lamarck deserves credit for
proposing that species change over time -- a radical departure from the
fixity of species doctrine.

\textbf{Georges Cuvier (1769--1832)} established comparative anatomy and
paleontology as rigorous disciplines. Examining fossils in the Paris
Basin, he documented that species in deeper (older) strata differed from
living forms. His explanation was \textbf{catastrophism}: periodic
catastrophic events (floods, volcanism) destroyed regional faunas, which
were then replaced by immigration from unaffected areas. Cuvier
vehemently opposed transmutation of species but ironically provided the
fossil evidence that later supported it.

\textbf{Charles Lyell (1797--1875)} countered catastrophism with
\textbf{uniformitarianism} in his \emph{Principles of Geology}
(1830--1833). He argued that the same geological processes operating
today -- erosion, sedimentation, volcanism -- have typically operated,
and at roughly the same rates. The Earth, therefore, must be immensely
old, providing the deep time necessary for gradual biological change.
Darwin carried the first volume of Lyell's \emph{Principles} aboard HMS
Beagle, and it profoundly shaped his thinking.

\subsubsection{Darwin's Voyage on HMS Beagle
(1831--1836)}\label{darwins-voyage-on-hms-beagle-18311836}

Charles Darwin boarded HMS Beagle as a gentleman naturalist at age 22.
The five-year circumnavigation exposed him to patterns that would
crystallize into the theory of natural selection:

\begin{itemize}
\tightlist
\item
  \textbf{Galapagos Islands}: Darwin collected finches, mockingbirds,
  and tortoises from different islands. He initially failed to label
  specimens by island -- an oversight he later called his greatest
  regret. Ornithologist John Gould identified the finches as distinct
  but related species, each with bill morphology matching its food
  source. The mockingbirds showed similar island-specific variation.
  Giant tortoises differed in carapace shape between islands (saddleback
  on arid islands with tall vegetation, dome-shaped on humid islands
  with low vegetation).
\item
  \textbf{South American fauna}: Fossil glyptodonts (giant armadillos)
  in Argentina resembled living armadillos -- a clear example of descent
  with modification rather than separate creation. The rhea (South
  American ostrich) filled an ecological niche analogous to the African
  ostrich, suggesting geographic replacement driven by similar selective
  pressures rather than independent design.
\item
  \textbf{Coral atolls and geological uplift}: Darwin observed raised
  marine terraces in Chile after the 1835 earthquake, confirming Lyell's
  gradualism.
\end{itemize}

\subsubsection{Wallace's Independent
Discovery}\label{wallaces-independent-discovery}

\textbf{Alfred Russel Wallace (1823--1913)}, working as a specimen
collector in the Malay Archipelago, independently conceived natural
selection during a bout of malaria on the island of Ternate in February
1858. His letter to Darwin -- the famous \textbf{Ternate letter} --
outlined a mechanism virtually identical to Darwin's unpublished theory.
Lyell and Joseph Hooker arranged a joint presentation to the Linnean
Society of London on July 1, 1858, reading Wallace's paper alongside
extracts from Darwin's 1844 essay and an 1857 letter to Asa Gray. Darwin
then rushed to complete \emph{On the Origin of Species}, published
November 24, 1859.

\subsubsection{The Three Conditions for Natural
Selection}\label{the-three-conditions-for-natural-selection}

Natural selection operates whenever three conditions are met
simultaneously:

\begin{enumerate}
\def\labelenumi{\arabic{enumi}.}
\tightlist
\item
  \textbf{Heritable variation}: Individuals in a population differ in
  traits, and at least some of that variation is genetically transmitted
  to offspring.
\item
  \textbf{Differential survival and reproduction}: Some variants survive
  longer, reproduce more, or both -- contributing disproportionately to
  the next generation.
\item
  \textbf{Environmental pressure}: Resources are limited; not most
  individuals can survive and reproduce equally. The environment
  determines which traits are advantageous.
\end{enumerate}

When most three conditions hold, \hyperref[gl:allele]{\textbf{allele}}
frequencies shift across generations -- populations evolve.

\subsubsection{The Modern Synthesis and
Beyond}\label{the-modern-synthesis-and-beyond}

The \textbf{Modern Synthesis} (also called the Neo-Darwinian Synthesis)
emerged in the 1930s--1950s through the work of Theodosius Dobzhansky
(\emph{Genetics and the Origin of Species}, 1937), Ernst Mayr
(\emph{Systematics and the Origin of Species}, 1942), George Gaylord
Simpson (\emph{Tempo and Mode in Evolution}, 1944), and Julian Huxley
(\emph{Evolution: The Modern Synthesis}, 1942). This intellectual
movement unified Darwin's mechanism of natural selection with Mendelian
genetics, mathematical
\hyperref[gl:population-genetics]{\textbf{population genetics}} (Fisher,
Wright, Haldane), paleontology, and systematics into a coherent
framework. The Modern Synthesis established that:

\begin{itemize}
\tightlist
\item
  Evolution is the change in allele frequencies within populations over
  time
\item
  Natural selection is the primary mechanism of adaptive evolution
\item
  Gradual accumulation of small genetic changes produces large-scale
  evolutionary patterns
\item
  \hyperref[gl:speciation]{\textbf{Speciation}} typically involves the
  evolution of \hyperref[gl:reproductive-isolation]{\textbf{reproductive
  isolation}} between geographically separated populations
\end{itemize}

The Modern Synthesis remains the foundation of evolutionary biology,
though later work on evo-devo, plasticity, niche construction,
epigenetic inheritance, symbiosis, and cultural evolution has expanded
the synthesis.

\begin{quote}
\textbf{Concept Check 1:} Lamarck proposed inheritance of acquired
characteristics. Why does modern genetics reject this mechanism? Can you
think of any biological phenomenon that superficially resembles
Lamarckian inheritance? (Hint: consider epigenetic inheritance.)
\end{quote}

\begin{center}\rule{0.5\linewidth}{0.5pt}\end{center}

\subsection{Evidence for Evolution}\label{evidence-for-evolution}

Five independent lines of evidence converge to support evolution. Each
alone would be suggestive; together they constitute one of the most
strongly supported theories in science.

\begin{figure}[htbp]
\centering
\includegraphics[width=0.82\linewidth,height=4.2in,keepaspectratio]{../figures/mermaid_inline/inline_mermaid_0104_ca339ead0138.png}
\caption{Flowchart for Evidence for Evolution: Evidence for Evolution, Fossil Record, Comparative Anatomy, and Embryology form the diagram's primary path or branches.}
\end{figure}

\subsubsection{Fossil Record}\label{fossil-record}

The fossil record documents the history of life through mineralized
remains preserved in sedimentary rock. Fossils form through several
processes: permineralization (minerals replace organic material),
compression (organisms flattened in sediment), amber preservation
(organisms entombed in tree resin), and trace fossils (footprints,
burrows, coprolites). The record is inherently incomplete --
fossilization requires specific conditions (rapid burial, absence of
scavengers, appropriate mineral chemistry), and soft-bodied organisms
are rarely preserved. Despite this incompleteness, key transitional
fossils bridge major evolutionary gaps:

\begin{itemize}
\tightlist
\item
  \textbf{Tiktaalik roseae} (375 Mya, Late Devonian): Discovered in 2004
  on Ellesmere Island, Arctic Canada, by Neil Shubin's team. This
  ``fishapod'' possesses fish features (scales, fins, gills) alongside
  tetrapod features (flat head, neck, wrist bones capable of
  weight-bearing, ribs for lung ventilation). It bridges the
  fish-to-tetrapod transition.
\item
  \textbf{Archaeopteryx lithographica} (150 Mya, Late Jurassic): Found
  in Solnhofen limestone, Germany. Combines dinosaur features (teeth,
  clawed fingers, bony tail) with bird features (asymmetric flight
  feathers, wishbone). Now understood as one member of a diverse
  radiation of feathered dinosaurs.
\item
  \textbf{Whale evolution}: One of the best-documented evolutionary
  transitions. \emph{Pakicetus} (50 Mya) was a terrestrial carnivore
  with cetacean ear bones. \emph{Ambulocetus} (49 Mya) was semi-aquatic
  with large hind limbs for swimming. \emph{Rodhocetus} (47 Mya) had
  reduced hind limbs and a more streamlined body. \emph{Basilosaurus}
  (37 Mya) was fully aquatic with vestigial hind limbs. Modern cetaceans
  retain vestigial pelvic bones embedded in muscle.
\end{itemize}

\subsubsection{Comparative Anatomy}\label{comparative-anatomy}

\textbf{Homologous structures} share a common developmental and
evolutionary origin but may serve different functions. The vertebrate
forelimb illustrates this principle: the human arm (manipulation), bat
wing (flight), whale flipper (swimming), and horse leg (running) most
contain the same bones -- humerus, radius, ulna, carpals, metacarpals,
phalanges -- arranged in the same relative positions. The functional
differences reflect modification of a shared ancestral blueprint.

\textbf{Analogous structures} arise from convergent evolution: the
dolphin fin and shark fin are superficially similar but have completely
different internal anatomy and developmental origins. Similarly, bird
wings and insect wings evolved independently.

\textbf{Vestigial structures} are reduced remnants of organs that were
functional in ancestors: the human coccyx (remnant of a tail), arrector
pili muscles (produce goose bumps -- useful for fur-covered ancestors,
functionless in largely hairless humans), wisdom teeth (third molars
suited to ancestral diets of tough plant material), and the whale pelvis
(tiny bones embedded in body wall, remnants of hind limbs).

\subsubsection{Embryology}\label{embryology}

\textbf{Pharyngeal arches} appear in the embryos of most vertebrates --
fish, amphibians, reptiles, birds, and mammals. In fish, these arches
develop into gills and their supporting structures. In terrestrial
vertebrates, the same embryonic structures are repurposed for entirely
different functions: the first pharyngeal arch forms the jaw and middle
ear ossicles (malleus and incus); the second arch forms the stapes,
styloid process, and much of the hyoid bone; the third and fourth arches
contribute to laryngeal and tracheal cartilages. This shared
developmental program, using homologous embryonic structures for
radically different adult functions, is compelling evidence for common
ancestry.

\textbf{Haeckel's biogenetic law} (``ontogeny recapitulates phylogeny'')
proposed that embryonic development replays evolutionary history in
sequence. While this formulation is an overstatement -- development does
not literally replay adult ancestral stages -- the underlying
observation remains valid: embryos of related species share
developmental stages that diverge progressively later in development,
reflecting their shared genetic toolkit (particularly the conserved Hox
\hyperref[gl:gene]{\textbf{gene}} clusters that pattern the
anterior-posterior body axis across most bilaterians).

Human embryos possess a tail at approximately week 5 of development,
containing 10--12 developing caudal vertebrae. By week 8, programmed
cell death (\hyperref[gl:apoptosis]{\textbf{apoptosis}}) reduces this to
the 3--5 fused vertebrae of the adult coccyx. In rare developmental
anomalies, the apoptotic program fails partially, and a ``vestigial
tail'' is present at birth -- a striking reminder of our evolutionary
heritage. Similarly, human embryos develop lanugo (fine body hair) at
approximately week 20, a remnant of the fur coat that covered our
mammalian ancestors.

\subsubsection{Biogeography}\label{biogeography}

The geographic distribution of organisms reflects evolutionary history,
geological events, and patterns of dispersal and vicariance. Darwin and
Wallace both recognized that biogeography provides powerful evidence for
evolution -- the distribution of organisms makes sense primarily in the
context of evolutionary history and plate tectonics:

\begin{itemize}
\tightlist
\item
  \textbf{Marsupial distribution}: Marsupials dominate Australia
  (kangaroos, koalas, wombats) but are largely absent from
  placental-dominated continents. This reflects the breakup of Gondwana:
  marsupials diversified in isolation on the Australian continent after
  its separation from Antarctica approximately 45 Mya.
\item
  \textbf{\hyperref[gl:island-biogeography]{Island biogeography}}:
  Oceanic islands harbor endemic species derived from mainland
  colonizers. Hawaiian honeycreepers, Galapagos finches, and Madagascar
  lemurs most represent adaptive radiations from single founding
  lineages.
\item
  \textbf{Wallace Line}: A biogeographic boundary running between Borneo
  and Sulawesi, and between Bali and Lombok, separating the Asian and
  Australian faunal regions. Deep ocean trenches (reaching depths of
  1,500 meters or more) prevented land bridges even during Pleistocene
  glacial periods when sea levels dropped by approximately 120 meters,
  maintaining distinct faunas on either side. West of the line:
  placental mammals (tigers, rhinoceroses, primates). East of the line:
  marsupials and monotremes (kangaroos, possums, platypuses). The line
  was identified by Wallace in 1859 based on his extensive collecting in
  the Malay Archipelago, and its correspondence with deep-water tectonic
  boundaries was confirmed much later by plate tectonic theory.
\end{itemize}

\begin{quote}
\textbf{Real-World Connection: Biogeography and Conservation Hotspots}

Understanding biogeographic patterns is critical for conservation.
Oceanic islands and other biogeographic isolates harbor disproportionate
numbers of endemic species -- organisms found nowhere else. Madagascar
covers 0.4\% of Earth's land surface but harbors approximately 5\% of
known species, nearly 90\% of which are endemic. The island has been
isolated from Africa for approximately 88 million years, allowing a
unique biota to evolve in isolation. Conservation of island biotas
requires understanding their evolutionary history: small, isolated
populations are inherently vulnerable to extinction from habitat loss,
introduced predators, and genetic erosion. The current extinction crisis
disproportionately affects island species -- approximately 75\% of
recorded bird and mammal extinctions since 1500 have occurred on
islands. Because ``recorded extinctions'' are biased toward vertebrates
and well-surveyed islands, use the index as a directional warning
signal, not as a complete inventory of biodiversity loss.
\end{quote}

\subsubsection{Molecular Evidence}\label{molecular-evidence}

Molecular data provide the most quantitative and testable evidence for
evolution. The molecular evidence is particularly powerful because it
generates precise, quantitative predictions that can be verified
independently:

\begin{itemize}
\tightlist
\item
  \textbf{Cytochrome c similarity}: This mitochondrial
  \hyperref[gl:protein]{\textbf{protein}} is found in most
  \hyperref[gl:aerobic]{\textbf{aerobic}} organisms. Human and
  chimpanzee cytochrome c sequences are 100\% identical. Human and
  rhesus monkey differ by 1 amino acid. Human and yeast differ by
  approximately 40\% -- yet retain the same basic function. The pattern
  of similarity matches phylogenetic predictions exactly.
\item
  \textbf{Pseudogenes}: Non-functional DNA sequences that were once
  functional genes. Humans and chimpanzees share the same inactivating
  \hyperref[gl:mutation]{\textbf{mutation}} in the \(\psi\eta\)-globin
  pseudogene -- an event so specific that independent occurrence is
  highly improbable. Shared pseudogenes are powerful evidence for common
  descent.
\item
  \textbf{Endogenous retroviruses (ERVs)}: Approximately 8\% of the
  human \hyperref[gl:genome]{\textbf{genome}} consists of ancient
  retroviral sequences inserted into the germline of our ancestors.
  Humans and chimpanzees share many ERV insertion sites at identical
  genomic locations -- each representing an independent ancestral
  infection event, providing a molecular ``fossil record'' within the
  genome.
\item
  \textbf{C-value paradox}: Genome size does not correlate with
  organismal complexity. The onion genome (about 16 Gb) is five times
  larger than the human genome (about 3.2 Gb). This reflects the
  accumulation of transposable elements and other non-coding sequences
  through evolutionary time, not design.
\item
  \textbf{Molecular phylogenies match morphological phylogenies}: When
  phylogenetic trees are constructed independently from DNA sequences
  and from morphological characters, they produce the same branching
  patterns. This convergence of independent evidence is powerful
  confirmation that the trees reflect real historical relationships, not
  artifacts of either data type.
\end{itemize}

\begin{quote}
\textbf{Real-World Connection: Forensic Phylogenetics}

Molecular phylogenetics has been used in criminal investigations. In the
landmark 1998 case of a Louisiana gastroenterologist accused of
infecting a former lover with HIV, phylogenetic analysis of viral
sequences demonstrated that the victim's HIV strain was nested within
the phylogenetic cluster from one of the physician's patients --
consistent with deliberate injection of contaminated blood. The viral
phylogeny, combined with epidemiological evidence, contributed to the
physician's conviction. This case established molecular phylogenetics as
admissible forensic evidence and demonstrated that evolutionary biology
has applications far beyond the academic laboratory.
\end{quote}

\begin{quote}
\textbf{Concept Check 2:} A creationist argues that similarities between
species reflect common design, not common descent. How would you use
shared pseudogenes and endogenous retroviruses to counter this argument?
\end{quote}

\begin{center}\rule{0.5\linewidth}{0.5pt}\end{center}

\subsection{Mechanisms of Natural
Selection}\label{mechanisms-of-natural-selection}

\begin{figure}[htbp]
\centering
\includegraphics[width=0.85\textwidth]{../figures/selection_simulation.png}
\caption{Natural selection simulation: allele-frequency change over generations under directional selection, balancing selection via heterozygote advantage, and disruptive selection via heterozygote disadvantage.}
\label{fig:unit_VI_selection_simulation}
\end{figure}

\subsubsection{Fitness}\label{fitness}

\textbf{Darwinian fitness} (\(W\)) measures relative reproductive
success -- the contribution of a genotype to the next generation
compared to other genotypes:

\begin{equation}
W = \frac{\text{reproductive output of genotype}}{\text{mean reproductive output of population}}
\label{eq:evolution_and_selection_1}
\end{equation}

Fitness is typically relative and context-dependent. An allele
conferring antibiotic resistance has \(W > 1\) in the presence of
antibiotics but may have \(W < 1\) in their absence (due to the
metabolic cost of resistance).

For a single locus with two alleles, if we assign the favored homozygote
fitness \(W_{AA} = 1\), then:

\begin{equation}
W_{Aa} = 1 - hs \quad \text{and} \quad W_{aa} = 1 - s
\label{eq:evolution_and_selection_2}
\end{equation}

where \(s\) is the selection coefficient and \(h\) is the dominance
coefficient.

\begin{figure}[htbp]
\centering
\includegraphics[width=0.82\linewidth,height=4.2in,keepaspectratio]{../figures/mermaid_inline/inline_mermaid_0105_07a35c33686d.png}
\caption{Fisher's Fundamental Theorem of Natural Selection. (1930): The rate of increase in mean population fitness equals the additive genetic variance in fitness:.}
\end{figure}

\begin{equation}
\frac{d\bar{W}}{dt} = V_A(W)
\label{eq:evolution_and_selection_3}
\end{equation}

This theorem implies that natural selection typically increases mean
fitness (though the environment may change, altering what ``fit''
means). It also implies that populations under strong selection rapidly
deplete additive genetic variance -- the more efficiently selection
operates, the faster it erodes its own fuel.

\subsubsection{Directional Selection}\label{directional-selection}

Directional selection shifts the mean
\hyperref[gl:phenotype]{\textbf{phenotype}} toward one extreme. The
population's phenotype distribution moves in one direction across
generations, in contrast to the stabilising and disruptive regimes whose
divergent allele-frequency trajectories are simulated in
\cref{fig:unit_VI_selection_simulation}.

\textbf{Industrial melanism in \emph{Biston betularia}}: Before the
Industrial Revolution, the typical peppered moth was light-colored,
camouflaged against lichen-covered tree bark. The melanic (dark) form,
controlled by a \hyperref[gl:dominant]{\textbf{dominant}} allele at the
\emph{cortex} locus, increased from approximately 1\% in 1848 to 98\% in
Manchester by 1898 as industrial soot darkened tree trunks. Light moths
were conspicuous to bird predators on dark bark. Following the Clean Air
Act (1956), lichen recovered, and the light form returned to dominance
-- a textbook example of directional selection reversing direction.

\textbf{Antibiotic resistance}: Random mutations conferring resistance
arise at low frequency in bacterial populations. In the presence of
antibiotics, resistant bacteria have enormous selective advantage
(\(s \approx 0.1\)--\(1.0\)). With generation times of 20--30 minutes,
resistance alleles can reach fixation within days -- evolution
observable within a single clinical episode. Multiple mechanisms confer
resistance: target modification (altered penicillin-binding proteins in
MRSA), enzymatic degradation (beta-lactamases that hydrolyze
penicillin), efflux pumps (that actively export antibiotics from the
cell), and membrane permeability changes (reduced porin expression in
Gram-negative bacteria).

The selection coefficient can be estimated from the rate of allele
frequency change:

\begin{equation}
s \approx \frac{\Delta p}{p(1-p)} \cdot \frac{1}{\Delta t}
\label{eq:evolution_and_selection_4}
\end{equation}

\subsubsection{Worked Example: Estimating the Selection
Coefficient}\label{worked-example-estimating-the-selection-coefficient}

\textbf{Problem:} A novel antibiotic resistance allele in a hospital
population of \emph{Staphylococcus aureus} is initially present at a
frequency of \(p = 0.01\). After \(\Delta t = 50\) generations of
continuous antibiotic exposure, the resistance allele reaches a
frequency of \(p' = 0.06\). Assuming the change is driven entirely by
selection, estimate the selection coefficient (\(s\)) favoring this
resistance allele.

\textbf{Solution:}

\begin{enumerate}
\def\labelenumi{\arabic{enumi}.}
\item
  \textbf{Identify the change in allele frequency (\(\Delta p\)):}
  \[ \Delta p = p' - p = 0.06 - 0.01 = 0.05  \label{eq:unit_VI_evolution_and_selection_item_1}\]
\item
  \textbf{Calculate the selection coefficient using the approximation
  formula:} Substitute the initial frequency \(p = 0.01\), the change
  \(\Delta p = 0.05\), and the time \(\Delta t = 50\) generations:
  \[ s \approx \frac{0.05}{0.01(1 - 0.01)} \cdot \frac{1}{50}  \label{eq:unit_VI_evolution_and_selection_item_2}\]

  \[ s \approx \frac{0.05}{0.0099} \cdot 0.02  \label{eq:unit_VI_evolution_and_selection_item_3}\]

  \[ s \approx 5.0505 \cdot 0.02 \approx 0.101  \label{eq:unit_VI_evolution_and_selection_item_4}\]
\end{enumerate}

The selection coefficient is approximately \textbf{0.10}. This means the
resistant bacteria have a roughly 10\% fitness advantage over the
sensitive bacteria in the presence of the antibiotic, allowing the
resistance allele to rapidly sweep through the population.

In the absence of antibiotics, resistance alleles often carry a
\textbf{fitness cost} (the resistant genotype grows more slowly than the
sensitive genotype). This predicts that resistance frequencies should
decline when antibiotic use is reduced. However, compensatory mutations
can eliminate the fitness cost of resistance without eliminating
resistance itself -- making resistance effectively irreversible in many
clinical settings.

\subsubsection{Stabilizing Selection}\label{stabilizing-selection}

Stabilizing selection favors intermediate phenotypes, reducing variance
without shifting the mean. It is the most common form of selection in
nature.

\textbf{Human birth weight}: \citet{karn1951} documented that neonatal
mortality is lowest at birth weights of approximately 3.5 kg. Both very
low-birth-weight
(\$\textless{}\(2 kg) and very high-birth-weight (\)\textgreater\$4.5
kg) infants suffer increased mortality from complications including
respiratory distress and obstructed labor, respectively. Modern medical
intervention has relaxed this selection somewhat but has not eliminated
it.

\textbf{Clutch size in birds}: David Lack proposed that the most common
clutch size maximizes the number of surviving offspring. Too few eggs
waste reproductive potential; too many eggs lead to inadequate
provisioning and higher chick mortality. Experimental studies in great
tits (\emph{Parus major}) confirmed this: broods artificially enlarged
beyond the modal clutch size fledged more chicks in the short term but
produced lighter, lower-quality offspring with reduced survival -- the
intermediate clutch size maximizes lifetime reproductive success.

\textbf{The stabilizing selection paradox}: If stabilizing selection
constantly removes extreme phenotypes and their underlying genetic
variation, why does heritable variation persist? Several mechanisms
maintain variation under stabilizing selection:

\begin{itemize}
\tightlist
\item
  Recurrent mutation continuously introduces new alleles
\item
  Balancing selection at underlying loci (heterozygote advantage,
  frequency-dependent selection)
\item
  Genotype-by-environment interactions: different genotypes may be
  ``intermediate'' in different environments
\item
  Pleiotropic constraints: alleles affecting the stabilized trait may
  also affect other traits under different selective regimes
\item
  Epistatic interactions among loci can maintain genetic variation even
  when the phenotypic distribution appears stable
\end{itemize}

\subsubsection{Disruptive (Diversifying)
Selection}\label{disruptive-diversifying-selection}

Disruptive selection favors both phenotypic extremes at the expense of
intermediate forms, potentially producing a bimodal distribution.

\textbf{Black-bellied seedcracker (\emph{Pyrenestes ostrinus})}: In
Cameroon, this finch exhibits a striking bimodal distribution of bill
sizes. Small-billed birds efficiently crack soft sedge seeds;
large-billed birds crack hard sedge seeds. Intermediate-billed birds
perform poorly on both seed types and have lower fitness. Remarkably,
this polymorphism is controlled by a single locus with large effect on
bill width, demonstrating that major adaptive variation can have a
simple genetic basis. Disruptive selection in this system maintains the
polymorphism because heterozygotes (with intermediate bills) have lower
fitness than either homozygote -- the opposite of heterozygote
advantage.

Disruptive selection is of particular evolutionary interest because, if
strong enough and accompanied by assortative mating (large-billed birds
preferring to mate with other large-billed birds, and vice versa), it
can drive \textbf{\hyperref[gl:sympatric-speciation]{sympatric
speciation}} -- the formation of new species without geographic
isolation (see \cref{sec:unit_VI_genetic_drift_and_speciation}).

\subsubsection{Sexual Selection}\label{sexual-selection}

Darwin recognized that many traits -- peacock tails, elk antlers,
birdsong complexity -- reduce survival but enhance mating success.
\textbf{Sexual selection} operates through two mechanisms:
\textbf{intrasexual selection} (typically male-male competition for
access to mates or to resources mates require) and \textbf{intersexual
selection} (mate choice, typically by females among displaying males).

\paragraph{Intrasexual selection: competition within the same
sex}\label{intrasexual-selection-competition-within-the-same-sex}

Intrasexual selection acts when members of one sex (most commonly males)
compete directly with one another for reproductive access. Three modes
of competition predominate:

\begin{itemize}
\tightlist
\item
  \textbf{Direct contest competition} (combat, ritualized fighting):
  Males fight to displace rivals. Red deer (\emph{Cervus elaphus}) stags
  engage in roaring contests followed by antler-locked pushing matches;
  mass and antler size predict victory and harem retention. Bighorn
  sheep (\emph{Ovis canadensis}) rams meet head-on at velocities of
  about 36 km/h; horn size and the bony reinforcement of the
  skull are direct outcomes of intrasexual selection. Northern elephant
  seal (\emph{Mirounga angustirostris}) bulls grow to 3-4\texorpdfstring{\ensuremath{\times}}{x} the female
  mass --- the most extreme sexual size dimorphism in mammals ---
  because beachmaster males monopolize harems of 30--100 females through
  direct combat.
\item
  \textbf{Resource defense (territoriality)}: Males defend territories
  containing critical resources (food, nest sites, oviposition
  substrate); females mate with the holder of the best territory.
  Red-winged blackbirds (\emph{Agelaius phoeniceus}) defend marsh
  territories whose vegetation density predicts female settlement.
  Cichlid males defend pebble-bottom spawning sites whose orientation to
  currents affects egg survival. The trait under selection is the
  \textbf{capacity to acquire and defend} the territory, not necessarily
  a morphological ornament.
\item
  \textbf{Scramble competition}: When females are widely distributed and
  rapidly receptive, the male who locates and reaches receptive females
  fastest wins. \emph{Drosophila} species in the wild engage in scramble
  competition for emerging virgin females; selection favors visual
  acuity, locomotor performance, and pheromone detection. In
  thirteen-lined ground squirrels, males that emerge earliest from
  hibernation locate the most receptive females. Scramble competition
  selects for \textbf{endurance, sensory acuity, and search efficiency}
  rather than weaponry.
\end{itemize}

A fourth mode, \textbf{sperm competition}, operates after copulation
when females mate with multiple males within a single fertile period.
Selection acts on sperm number, sperm motility, and copulatory plug
deposition. Testis-to-body-mass ratio is a quantitative correlate of
sperm competition intensity: chimpanzees (highly polyandrous) have
testes about 3\texorpdfstring{\ensuremath{\times}}{x} the size of gorillas (harem-guarding) on a
body-mass-corrected basis.

\paragraph{Intersexual selection: mate
choice}\label{intersexual-selection-mate-choice}

In intersexual selection, members of one sex (typically females) choose
among displaying members of the other. The peacock's elaborate tail is
the canonical example. Three complementary hypotheses explain female
preference for costly ornaments --- they are not mutually exclusive, and
empirical examples often combine elements of the three mechanisms.

\textbf{Sensory exploitation} provides a fourth, mechanistically
distinct hypothesis: female preferences may arise initially from
pre-existing sensory biases unrelated to mate choice (e.g., sensitivity
to a frequency band useful for prey detection), and males evolve traits
that exploit those biases. The túngara frog (\emph{Engystomops
pustulosus}) ``chuck'' call exploits a sensory bias inherited from a
common ancestor: females of related species that have rarely produced
``chucks'' themselves still prefer males that add the syllable. Sensory
exploitation grounds preference in \textbf{pre-existing neural
architecture} rather than in genetic correlation or signal honesty.

\paragraph{Honest signaling (Zahavi handicap
principle)}\label{honest-signaling-zahavi-handicap-principle}

Amotz Zahavi (1975) proposed that \textbf{costly ornaments are honest
signals of male quality} precisely because they are costly. A male who
can survive the metabolic, predation, and parasite-related costs of an
elaborate tail must have superior underlying genetic quality (immunity,
foraging ability, parasite resistance). Females that select males with
the most extravagant ornaments preferentially choose mates who carry
``good genes'' --- the genes their offspring will inherit.

The handicap principle predicts that \textbf{ornaments should be
condition-dependent}: primarily individuals in good condition can
produce the most extreme display. \citet{hamilton1982} extended this
with the \textbf{parasite-mediated sexual selection hypothesis}: bright
male coloration in birds, bright tail color in fish, and similar
ornaments depend on physiological condition that itself depends on
parasite resistance. Female choice based on ornament brilliance is
therefore choice for parasite-resistance alleles.

\textbf{Empirical support}: In the satin bowerbird
(\emph{Ptilonorhynchus violaceus}), males build elaborate bowers and
decorate them with blue objects to attract females. Males with brighter
plumage and more extensive bowers also have lower parasite loads ---
consistent with the handicap principle. In barn swallows (\emph{Hirundo
rustica}), males with longer outer tail feathers (a sexually selected
trait) carry fewer feather mites.

\paragraph{Fisher-Lande runaway
selection}\label{fisher-lande-runaway-selection}

R.A. Fisher (1930) and Russell Lande (1981) developed the
\textbf{runaway selection model}, which can produce arbitrarily extreme
ornaments through a positive-feedback loop between female preference and
male trait --- without requiring that the ornament signal anything about
quality.

The model proceeds in three stages:

\begin{enumerate}
\def\labelenumi{\arabic{enumi}.}
\tightlist
\item
  \textbf{Initial preference}: A small genetic predisposition for some
  male trait variant exists in females (e.g., females prefer males with
  slightly longer tails because long tails happen to indicate flight
  performance).
\item
  \textbf{Genetic correlation}: When females with the preference allele
  mate with males possessing the preferred trait, their
  \textbf{offspring inherit both genetic linkage groups} --- daughters
  receive the preference allele \emph{and} sons receive the trait
  allele. The two loci become genetically correlated.
\item
  \textbf{Runaway}: Now selection on the trait drives correlated change
  in the preference, and vice versa. The two loci can spiral together to
  extreme values, \textbf{even when the trait becomes maladaptive} by
  other survival measures, because the trait is now selected purely by
  its effect on mating success.
\end{enumerate}

Runaway is balanced eventually by natural selection against the ornament
becoming dangerous (predation, energetic costs). The equilibrium can be
at moderate ornament size or, depending on parameters, can produce the
bizarre extremes seen in birds of paradise and peacocks.

\paragraph{Good genes model}\label{good-genes-model}

The \textbf{good genes model} is closely related to but distinguishable
from the handicap principle. The handicap requires that the ornament
itself be costly and condition-dependent. The good genes model primarily
requires that some heritable component of male quality (immunity, growth
rate, viability, foraging skill) be correlated with the trait females
prefer. The female who chooses a higher-quality male produces
higher-quality offspring through both \textbf{autosomal good genes}
(offspring inherit advantageous viability alleles) and \textbf{viability
genes}.

Empirical tests are difficult because they require demonstrating
heritable variance in fitness correlated with the trait --- a very high
data bar.

\paragraph{Worked scenario: peacock tail
evolution}\label{worked-scenario-peacock-tail-evolution}

Imagine an ancestral peafowl population in which a small subset of
females have a slight genetic preference for males with longer tails
(``preference allele'' frequency = 0.05). Long-tailed males are 5\% more
likely to be mated than short-tailed males, but they pay a 10\% survival
cost from increased predation. The genetic correlation between
preference and trait builds over generations --- initially weakly, but
accelerating once the correlation becomes substantial. After
about 100 generations of runaway, both preference frequency and
tail length have escalated dramatically. The system reaches equilibrium
when the survival cost (10\% per unit tail length) balances the mating
advantage (5\% \texorpdfstring{\ensuremath{\times}}{x} tail length), producing tail lengths far beyond what
natural selection alone would favor.

\textbf{The peacock's tail thus is} simultaneously: (a) an honest signal
of genetic quality (Zahavi handicap), (b) the equilibrium of a Fisherian
runaway between preference and trait, and (c) a marker for ``good
genes'' (parasite resistance) \citep{hamilton1982}. Most three
mechanisms can operate together. Modern empirical work emphasizes that
real peacock tail evolution --- and most sexual ornaments --- combines
these forces in proportions that vary across species and contexts.

\subsubsection{Kin Selection}\label{kin-selection}

\textbf{Hamilton's rule} (1964) explains altruistic behavior: an
altruistic act is favored when:

\begin{equation}
rB > C
\label{eq:unit_VI_hamilton_rule}
\end{equation}

where \(r\) is the coefficient of relatedness between actor and
recipient, \(B\) is the reproductive benefit to the recipient, and \(C\)
is the reproductive cost to the actor.

Coefficients of relatedness:

\begin{itemize}
\tightlist
\item
  Full siblings: \(r = 0.5\)
\item
  Half-siblings: \(r = 0.25\)
\item
  First cousins: \(r = 0.125\)
\item
  Parent--offspring: \(r = 0.5\)
\end{itemize}

\textbf{Eusociality in Hymenoptera}: In haplodiploid species (ants,
bees, wasps), females are diploid and males are haploid. Full sisters
share \(r = 0.75\) (they share most of their father's genes plus half of
their mother's). This unusually high relatedness may predispose
haplodiploidy systems toward eusociality. Modern understanding
emphasizes that eusociality evolves when ecological conditions
(defensible nests, progressive provisioning, high predation risk for
solitary individuals) combine with kin structure to make helping more
profitable than independent reproduction.

\subsubsection{Evolutionary Game Theory: Hawk--Dove and the
ESS}\label{evolutionary-game-theory-hawkdove-and-the-ess}

Evolutionary game theory analyzes situations in which an individual's
payoff depends on the \textbf{strategies adopted by others} in the
population. The framework was introduced by John Maynard Smith and
George Price (1973) to explain ritualized animal contests. The
cornerstone concept is the \textbf{evolutionarily stable strategy
(ESS)}: a strategy that, once common in a population, cannot be invaded
by any rare alternative strategy.

The classic example is the \textbf{Hawk--Dove game}. Two individuals
contest a resource of value \(V\). A ``Hawk'' escalates the contest,
risking injury at cost \(C\) if it loses. A ``Dove'' displays but
retreats from escalation. The expected payoffs:

{\def\LTcaptype{none} % do not increment counter
\begin{longtable}[]{@{}rcc@{}}
\toprule\noalign{}
Opponent \texorpdfstring{\ensuremath{\to}}{->} & Hawk & Dove \\
\midrule\noalign{}
\endhead
\bottomrule\noalign{}
\endlastfoot
\textbf{Hawk plays} & \(\frac{V-C}{2}\) & \(V\) \\
\textbf{Dove plays} & \(0\) & \(\frac{V}{2}\) \\
\end{longtable}
}

When two Hawks meet, each wins half the time and pays the cost half the
time, giving expected payoff \((V-C)/2\). A Hawk against a Dove takes
the resource (\(V\)). Two Doves share peacefully (\(V/2\) each).

If \(V > C\) (resource exceeds the cost of injury), pure Hawk is the ESS
--- escalation typically pays. If \(V < C\) (more interesting biological
case), neither pure strategy is stable: pure Dove can be invaded by
Hawks (who exploit the peaceful population), but pure Hawk can be
invaded by Doves (who avoid the cost of constant fighting). The stable
outcome is a \textbf{mixed ESS} with frequencies determined by the
cost--benefit structure:

\begin{equation}
p^*_{\text{Hawk}} = \frac{V}{C}
\label{eq:unit_VI_ess_hawkdove}
\end{equation}

\paragraph{Numerical illustration of the Hawk--Dove
ESS}\label{numerical-illustration-of-the-hawkdove-ess}

Suppose the contested resource is worth \(V = 10\) (e.g., a breeding
territory worth 10 offspring) and the cost of injury in a fight is
\(C = 40\) (a serious wound costing 40 future offspring). Because
\(V < C\), the mixed ESS predicts \(p^*_{\text{Hawk}} = 10/40 = 0.25\)
--- twenty-five percent of contests are escalated to fighting. At
equilibrium, the \textbf{expected payoff is identical} for both
strategies:

\begin{itemize}
\tightlist
\item
  \(E[\text{Hawk}] = 0.25 \cdot \frac{10-40}{2} + 0.75 \cdot 10 = -3.75 + 7.5 = 3.75\)
\item
  \(E[\text{Dove}] = 0.25 \cdot 0 + 0.75 \cdot 5 = 3.75\)
\end{itemize}

This frequency-dependence is the engine that maintains the polymorphism:
any deviation from \(p = 0.25\) creates a fitness gradient that pulls
the population back. A population of 30\% Hawks rewards Doves more than
Hawks; a population of 20\% Hawks rewards Hawks more than Doves. The ESS
is \textbf{negatively frequency-dependent at the strategy level} --- the
rarer strategy typically has the higher payoff.

When everyone plays Hawk with probability \(V/C\), no individual gains
by deviating. The same equilibrium can be realized in two ways: each
individual mixes Hawk and Dove with probability \(V/C\), or the
population is a polymorphism with fraction \(V/C\) committed Hawks and
\(1 - V/C\) committed Doves. \textbf{Behavioral dimorphisms in spider
mating tactics, frog calling strategies, and male-morph polymorphisms in
side-blotched lizards} match Hawk--Dove or related ESS predictions
quantitatively.

Game theory generalizes to many other biological problems: parental
investment (Trivers 1972), public-goods cooperation (Tragedy of the
Commons), the evolution of virulence in pathogens, and host--parasite
coevolution. The ESS concept formalizes Darwin's insight that ``the
rules of the game'' can themselves be evolutionary outcomes ---
selection acts on strategies, not just on traits.

\subsubsection{The Breeder's Equation}\label{the-breeders-equation}

While population-genetic models track allele frequencies under
selection, \textbf{quantitative genetics} describes the response of
continuously varying phenotypes (height, milk yield, beak depth) to
selection. The central result is the \textbf{breeder's equation}:

\begin{equation}
R = h^2 S
\label{eq:unit_VI_breeders_equation}
\end{equation}

where \(R\) is the \textbf{response to selection} (the change in mean
phenotype across one generation), \(S\) is the \textbf{selection
differential} (the difference between the mean of selected parents and
the mean of the entire parental population), and \(h^2\) is the
\textbf{narrow-sense heritability} --- the proportion of phenotypic
variance attributable to additive genetic variance:

\begin{equation}
h^2 = \frac{V_A}{V_P} = \frac{V_A}{V_A + V_D + V_E}
\label{eq:unit_VI_heritability}
\end{equation}

with \(V_A\) = additive genetic variance, \(V_D\) = dominance variance,
\(V_E\) = environmental variance.

\paragraph{Worked example: Galapagos
finches}\label{worked-example-galapagos-finches}

Peter and Rosemary Grant measured beak depth in \emph{Geospiza fortis}
during the 1977 Galapagos drought. The selection differential (mean beak
depth of survivors minus mean of original population) was \(S = +0.6\)
mm. Heritability of beak depth, estimated from parent--offspring
regression, was \(h^2 = 0.74\). The predicted response is:

\[R = h^2 \times S = 0.74 \times 0.6 \text{ mm} = 0.44 \text{ mm} \label{eq:unit_VI_evolution_and_selection_item_5}\]

The observed response in the next generation (offspring of the 1977
survivors) was +0.43 mm --- a strikingly precise prediction confirming
both the heritability estimate and the breeder's equation. The drought
selected for deeper beaks (better at cracking large, hard seeds that
remained when small seeds were exhausted), and this trait responded to
selection within a single generation.

\paragraph{Worked example: artificial selection on plant
height}\label{worked-example-artificial-selection-on-plant-height}

Suppose a maize breeder wishes to increase mean plant height. The base
population has mean height \(\bar{z}_0 = 200\) cm with phenotypic
standard deviation \(\sigma_P = 15\) cm. The breeder selects the tallest
10\% of plants as parents of the next generation, whose mean height is
\(\bar{z}_{\text{selected}} = 205\) cm. Heritability has been estimated
by parent--offspring regression at \(h^2 = 0.8\) (a high but realistic
value for plant height under controlled conditions).

\begin{enumerate}
\def\labelenumi{\arabic{enumi}.}
\tightlist
\item
  \textbf{Selection differential}:
  \(S = \bar{z}_{\text{selected}} - \bar{z}_0 = 205 - 200 = 5\) cm.
\item
  \textbf{Predicted response}: \(R = h^2 \cdot S = 0.8 \times 5 = 4\)
  cm. The next generation's mean is predicted at
  \(\bar{z}_1 = 200 + 4 = 204\) cm.
\item
  \textbf{After multiple generations}: If \(h^2\) stays roughly constant
  and the selection differential is reapplied each generation, the
  population gains about 4 cm per generation in expectation.
  After 10 generations under unrelaxed selection, mean height could rise
  to about 240 cm --- although in practice \(h^2\) erodes as
  additive variance is depleted, and physiological constraints (lodging,
  mechanical failure under wind load) eventually limit the response.
\end{enumerate}

This calculation illustrates two practical points. \textbf{First,
\(h^2\) matters more than the apparent strength of selection}: doubling
the selection differential to \(S = 10\) would produce \(R = 8\) cm, but
halving \(h^2\) to 0.4 (more realistic for height in a wild population)
produces \(R = 2\) cm even at \(S = 5\). \textbf{Second, response
decelerates over time} as additive variance is consumed by directional
selection --- long-term selection experiments (the Illinois corn
protein/oil lines, now over 100 generations) eventually plateau as
genetic variance is exhausted unless mutation refills the supply.

\paragraph{Why the breeder's equation
works}\label{why-the-breeders-equation-works}

The equation reflects a simple algebraic truth: primarily the
\textbf{additive} genetic component of phenotypic variation is reliably
transmitted from parents to offspring. Dominance variance (\(V_D\))
reflects allele-pairing combinations that are scrambled by meiosis.
Epistatic interactions (\(V_I\)) similarly do not transmit faithfully.
Environmental variance (\(V_E\)) is non-heritable. Thus the
\textbf{realized} response equals the additive heritability times the
selection that was actually applied. The breeder's equation is the
workhorse of agricultural selection (dairy cattle, corn, tomatoes),
animal breeding (racehorses, dogs), and the analysis of natural
selection in wild populations.

\paragraph{Limitations and the multivariate breeder's
equation}\label{limitations-and-the-multivariate-breeders-equation}

The simple form assumes a single trait under direct selection, with
stable heritability across generations and a stable environment. In
practice, traits are correlated through pleiotropy and linkage;
selection on one trait drags correlated traits along. The
\textbf{multivariate breeder's equation} generalizes:

\[\Delta\bar{z} = G P^{-1} S \label{eq:unit_VI_evolution_and_selection_item_6}\]

where \(\Delta\bar{z}\) is the vector of phenotypic responses, \(G\) is
the genetic covariance matrix, \(P\) is the phenotypic covariance
matrix, and \(S\) is the vector of selection differentials. This
Lande--Arnold framework (1983) underlies modern multivariate
quantitative genetics and explains why trait evolution can be
\textbf{constrained by genetic correlations} even when selection on each
trait individually is strong.

\subsubsection{Evolvability, Modularity, and Developmental
Constraints}\label{evolvability-modularity-and-developmental-constraints}

The capacity of a population to respond to selection is itself a
heritable, evolving property called \textbf{evolvability}. Two
organizational features make some lineages dramatically more evolvable
than others.

\paragraph{Modularity}\label{modularity-1}

A \textbf{modular} genotype--phenotype map partitions traits into
semi-independent units, each controlled by largely separate gene
networks. Selection on one module --- say, fin shape in cichlids, or
beak depth in finches --- leaves other modules undisturbed. Modular
architectures evolve faster because beneficial modifications in one
module do not derail traits in other modules. The \textbf{butterfly wing
eyespot system} is a paradigm of modularity: each eyespot is patterned
by a conserved gene-regulatory module (Distal-less, Engrailed, Spalt)
that can be redeployed independently across wing positions, allowing
rapid diversification of eyespot patterns across species without
disturbing wing structure or other traits.

In contrast, \textbf{highly integrated} (non-modular) systems show
\textbf{pleiotropic constraint}: every mutation affects many traits
simultaneously, most beneficial mutations also produce deleterious
side-effects, and adaptive evolution is slow. The \textbf{vertebrate
skull} is intermediate --- highly integrated for biomechanical reasons
(the bones must articulate properly), but with enough modularity that
snakes, birds, and toothed whales evolved radically divergent skull
shapes from the same ancestral plan.

\paragraph{Developmental constraints, canalization, and cryptic
variation}\label{developmental-constraints-canalization-and-cryptic-variation}

\textbf{Canalization} is the developmental buffering that produces the
same phenotype across a range of genetic and environmental
perturbations. Conrad Waddington (1942) introduced the metaphor of the
\textbf{epigenetic landscape}, in which developmental trajectories
follow valleys of low resistance toward stable phenotypic outcomes.
Canalization implies that \textbf{genetic variation is hidden} beneath
robust development --- many segregating alleles do not affect phenotype
because development buffers their effects.

When buffering is overwhelmed --- by extreme environmental stress, by a
chaperone disruption (HSP90 inhibition reveals abundant cryptic
variation in \emph{Drosophila}), or by a major mutation ---
\textbf{cryptic genetic variation} is exposed. The newly visible
phenotypic variation can then respond to selection. \textbf{Genetic
assimilation} is the process by which a phenotype originally produced
primarily under stress becomes fixed and is produced under most
conditions, having recruited cryptic variation along the way.

\textbf{Cryptic variation as evolutionary fuel}: Recent studies in
zebrafish, threespine stickleback, and \emph{Drosophila} show that
populations with deeper reservoirs of cryptic variation respond to novel
selection pressures faster than populations with shallow reservoirs.
Modular development plus cryptic variation thus form a paired pair of
features that increase evolvability --- modularity by allowing
independent change, cryptic variation by storing adaptive material.

\textbf{Inclusive fitness}: Hamilton's insight extended the concept of
fitness beyond direct reproduction. An individual's \textbf{inclusive
fitness} includes both its direct fitness (own offspring) and the
fitness gained indirectly by helping relatives reproduce --- weighted by
the coefficient of relatedness. A sterile worker bee with zero direct
fitness can still have high inclusive fitness if she helps her queen
mother produce hundreds of full sisters (\(r = 0.75\) each).

\subsubsection{Frequency-Dependent
Selection}\label{frequency-dependent-selection}

\textbf{Negative frequency-dependent selection}: Rare phenotypes have a
fitness advantage. As a phenotype becomes common, its fitness decreases.
This maintains polymorphism.

\textbf{Side-blotched lizards (\emph{Uta stansburiana})}: Males have
three throat color morphs engaged in rock-paper-scissors dynamics.
Orange-throated males are aggressive and defeat blue-throated males.
Blue-throated males form pair bonds that resist sneaking by
yellow-throated males. Yellow-throated males (female mimics)
successfully sneak copulations from orange-throated males. No morph can
dominate permanently; frequencies cycle over approximately 6-year
periods --- a real-world realization of a multi-strategy ESS.

\textbf{Positive frequency-dependent selection}: The common phenotype
has an advantage (e.g., warning coloration in Mullerian mimicry ---
predators learn to avoid the most common pattern).

\begin{quote}
\textbf{Concept Check 3:} A population of side-blotched lizards has
roughly equal frequencies of orange, blue, and yellow male morphs. Each
year, biologists notice that the frequency of orange males rises sharply
for two years, then yellow males rise as orange decline, then blue males
rise as yellow decline --- and the cycle repeats. Why does negative
frequency-dependent selection produce \textbf{cycles} rather than a
stable equilibrium with most three morphs at constant frequencies?
(Hint: think about what happens when one morph's fitness depends on the
\emph{current} frequencies of the other two, not a one-step-ahead
average.)
\end{quote}

\begin{quote}
\textbf{Concept Check on quantitative genetics:} A dairy farmer
practices artificial selection on milk yield, with selection
differential \(S = 1{,}000\) kg/lactation and observed response
\(R = 200\) kg in the next generation. (a) What is the implied
narrow-sense heritability? (b) The same selection regime applied for 20
generations produces about 3,000 kg total response (not 4,000 kg as
naive multiplication would predict). Give two reasons the response
decelerates.
\end{quote}

\subsubsection{Balancing Selection}\label{balancing-selection}

Several mechanisms maintain genetic polymorphism within populations,
collectively termed \textbf{balancing selection}:

\begin{itemize}
\tightlist
\item
  \textbf{Heterozygote advantage (overdominance)}: The heterozygote has
  higher fitness than either homozygote. The classic example is sickle
  cell anemia: \(HbA/HbS\) heterozygotes have resistance to
  \emph{Plasmodium falciparum} malaria without the severe anemia of
  \(HbS/HbS\) homozygotes. Equilibrium frequency of \(HbS\) in
  malaria-endemic regions is approximately \(q = h_1/(h_1 + h_2)\) where
  \(h_1\) and \(h_2\) are the selection coefficients against the two
  homozygotes.
\item
  \textbf{Negative frequency-dependent selection}: As described above,
  rare phenotypes have higher fitness, preventing fixation.
\item
  \textbf{Spatially varying selection}: Different alleles are favored in
  different environments within the species' range. Migration connects
  subpopulations, maintaining polymorphism across the metapopulation.
\item
  \textbf{Temporally varying selection}: The fitness ranking of
  genotypes fluctuates across seasons or years, preventing any single
  allele from reaching fixation.
\end{itemize}

\begin{quote}
\textbf{Real-World Connection: Antibiotic Resistance as Evolution in
Real Time}

Methicillin-resistant \emph{Staphylococcus aureus} (MRSA) illustrates
evolution by natural selection operating on a clinically devastating
timescale. Penicillin was introduced in 1943; penicillin-resistant
\emph{S. aureus} appeared by 1947. Methicillin was introduced in 1961;
MRSA was reported by 1962. The ESKAPE pathogens (\emph{Enterococcus
faecium}, \emph{Staphylococcus aureus}, \emph{Klebsiella pneumoniae},
\emph{Acinetobacter baumannii}, \emph{Pseudomonas aeruginosa},
\emph{Enterobacter} species) represent the most urgent antibiotic
resistance threats. Each has evolved resistance through the standard
Darwinian mechanism: heritable variation (resistance mutations or
horizontally acquired genes), differential reproduction (resistant
bacteria survive antibiotic exposure), and environmental pressure
(antibiotic use in hospitals and agriculture). The evolution of
resistance is predictable, repeatable, and observable in real time --
providing one of the clearest demonstrations that natural selection is
not merely a historical phenomenon but an ongoing process with immediate
consequences for human health.
\end{quote}

\begin{center}\rule{0.5\linewidth}{0.5pt}\end{center}

\subsection{Adaptation and Its Limits}\label{adaptation-and-its-limits}

\subsubsection{Adaptation and
Exaptation}\label{adaptation-and-exaptation}

An \textbf{adaptation} is a heritable feature shaped by natural
selection that increases fitness in the organism's current environment.
Adaptations are not perfect solutions but rather the best available
modifications of pre-existing structures.

An \textbf{exaptation} is a feature that was originally shaped by
natural selection for one function but has been co-opted for a different
function:

\begin{itemize}
\tightlist
\item
  \textbf{Feathers} evolved initially for
  \hyperref[gl:thermoregulation]{\textbf{thermoregulation}} in small
  theropod dinosaurs. Flight capability was a later co-option of an
  existing structure.
\item
  \textbf{Swim bladders} in bony fish are homologous with lungs in
  terrestrial vertebrates. The ancestral structure was likely a
  vascularized outpocketing of the pharynx used for gas exchange in
  oxygen-poor water, later modified for buoyancy control in derived fish
  lineages and for terrestrial respiration in the lineage leading to
  tetrapods.
\end{itemize}

\subsubsection{Phylogenetic Constraints}\label{phylogenetic-constraints}

Evolution cannot start from scratch. Natural selection can primarily
modify existing structures, constrained by the organism's
\textbf{bauplan} (body plan):

\begin{itemize}
\tightlist
\item
  \textbf{Panda's ``thumb''}: The giant panda's thumb is not a true
  digit but an enlarged radial sesamoid bone. The true thumb is
  committed to the bear-like paw. Natural selection co-opted a wrist
  bone to create a functional but imperfect food-handling structure -- a
  testament to evolutionary tinkering rather than design.
\item
  \textbf{Recurrent laryngeal nerve}: In mammals, this branch of the
  vagus nerve loops from the brainstem down around the aortic arch and
  back up to the larynx. In giraffes, this detour adds approximately 4.5
  meters of unnecessary nerve length. The path makes developmental sense
  given the ancestral fish anatomy (where the nerve innervated the sixth
  gill arch) but is a suboptimal arrangement in a long-necked mammal.
\end{itemize}

\subsubsection{Developmental
Constraints}\label{developmental-constraints}

Body plans established during embryonic development are deeply conserved
and difficult to alter fundamentally. The basic organization of body
segments, limb positions, and organ systems is established by deeply
conserved developmental gene networks (Hox genes, signaling pathways
such as Hedgehog, Wnt, and BMP). Tetrapods are ancestrally four-limbed
-- not because four is optimal for every ecological situation but
because the tetrapod body plan was established approximately 375 Mya and
is embedded so deeply in the developmental genetic toolkit that altering
limb number would require simultaneously reorganizing multiple
interconnected developmental pathways. No tetrapod has evolved six
limbs, despite the potential utility of additional appendages. Snakes
have lost their limbs entirely (through modification of Hox gene
expression domains), demonstrating that reduction is possible even when
addition is not.

Similarly, the vertebrate eye is built ``backwards'' -- with
photoreceptors behind the retinal \hyperref[gl:neuron]{\textbf{neuron}}s
and blood vessels, creating a blind spot where the optic nerve exits.
Cephalopod eyes evolved independently with photoreceptors facing
forward, lacking a blind spot. The vertebrate arrangement persists
because the ancestral developmental program cannot easily be reorganized
without disrupting the entire eye structure.

\subsubsection{Trade-offs}\label{trade-offs}

Organisms cannot maximize most fitness components simultaneously:

\begin{itemize}
\tightlist
\item
  \textbf{Immune function versus reproduction}: Mounting an immune
  response is energetically costly. In many species, individuals
  investing heavily in reproduction show suppressed immune function and
  increased susceptibility to parasites.
\item
  \textbf{Speed versus endurance}: Cheetahs achieve 110 km/h but sustain
  peak sprints for roughly 200--300 meters. Wolves run at 55 km/h but
  can maintain pursuit for kilometers.
\item
  \textbf{Crypsis versus mate attraction}: Brightly colored males
  attract both mates and predators. The optimal coloration balances
  sexual selection for conspicuousness against natural selection for
  camouflage. In guppies (\emph{Poecilia reticulata}), males in
  low-predation streams are more colorful than those in high-predation
  streams -- demonstrating the trade-off between sexual selection
  (favoring bright colors) and natural selection (favoring drab colors
  that avoid predator detection).
\item
  \textbf{Fecundity versus offspring quality}: Organisms face a
  fundamental trade-off between producing many small offspring
  (r-strategy) or few large offspring (K-strategy). Each egg or seed
  requires maternal investment; allocating more to each offspring
  reduces the total number that can be produced.
\item
  \textbf{Current reproduction versus future survival}: Investing
  heavily in current reproduction often reduces an organism's
  probability of surviving to reproduce again. Semelparous organisms
  (salmon, annual plants) invest everything in a single reproductive
  event and then die. Iteroparous organisms (most mammals, perennial
  plants) reproduce multiple times but invest less per reproductive
  event.
\end{itemize}

\subsubsection{Coevolution}\label{coevolution}

\textbf{Coevolution} occurs when two or more species exert reciprocal
selective pressures on each other, driving evolutionary change in both
lineages.

\textbf{Predator-prey arms race}: Rough-skinned newts (\emph{Taricha
granulosa}) produce tetrodotoxin (TTX), one of the most potent
neurotoxins known. Common garter snakes (\emph{Thamnophis sirtalis}) in
sympatric populations have evolved resistance through amino acid
substitutions in their sodium channel genes. The level of newt toxicity
and snake resistance are geographically correlated -- populations with
more toxic newts have more resistant snakes, and vice versa. This
represents a classic coevolutionary arms race.

\textbf{Plant-herbivore coevolution}: Brassicaceae (mustard family)
plants produce glucosinolates as herbivore deterrents. Specialist
herbivores (Pieridae butterflies) have evolved glucosinolate-specific
detoxification \hyperref[gl:enzyme]{\textbf{enzyme}}s. This escalation
of defense and counter-defense has been linked to diversification in
both plant and insect lineages.

\textbf{Red Queen hypothesis} (Van Valen, 1973): Species must
continually evolve simply to maintain their current fitness relative to
coevolving partners, particularly parasites. Like the Red Queen in
\emph{Through the Looking-Glass}, organisms must keep running just to
stay in place. This hypothesis explains the maintenance of sexual
reproduction: the genetic variation generated by
\hyperref[gl:recombination]{\textbf{recombination}} provides a moving
target for parasites adapted to common host genotypes.

\begin{quote}
\textbf{Real-World Connection: Coevolution in Agriculture}

The coevolutionary arms race between crops and pathogens has enormous
agricultural significance. When a new resistance gene is deployed in a
crop variety, pathogen populations rapidly evolve
\hyperref[gl:virulence]{\textbf{virulence}} alleles that overcome it --
often within 3--5 years. This ``boom-and-bust cycle'' has driven the
development of gene pyramiding (stacking multiple resistance genes),
multiline varieties (mixtures of genetically different cultivars), and
integrated pest management. The Irish Potato Famine (1845--1849)
resulted partly from the genetic uniformity of Irish potato crops -- a
single clone susceptible to \emph{Phytophthora infestans}. Modern
agricultural genetics applies evolutionary principles to stay ahead of
rapidly evolving pathogens.
\end{quote}

\begin{quote}
\textbf{Concept Check 4:} The panda's ``thumb'' is often cited as
evidence against intelligent design and for evolution. Explain why an
imperfect structure supports the evolutionary explanation better than a
designed explanation.
\end{quote}

\begin{center}\rule{0.5\linewidth}{0.5pt}\end{center}

\subsection{Macroevolution and Evolutionary
Rates}\label{macroevolution-and-evolutionary-rates}

\subsubsection{Microevolution and
Macroevolution}\label{microevolution-and-macroevolution}

\textbf{Microevolution} refers to changes in allele frequencies within
populations over relatively short timescales -- the processes described
across
\cref{sec:unit_V_population_genetics,sec:unit_VI_evolution_and_selection}.
These include mutation, natural selection,
\hyperref[gl:genetic-drift]{\textbf{genetic drift}}, and gene flow
operating within species. \textbf{Macroevolution} encompasses
large-scale patterns observed over geological time: the origin of major
body plans, mass extinctions and subsequent radiations, long-term trends
in biodiversity, and the emergence of evolutionary novelties such as the
vertebrate eye or the flower.

A central question in evolutionary biology is whether macroevolution is
simply the accumulation of microevolutionary changes over vast
timescales, or whether additional processes operate at higher levels of
biological organization. Most evolutionary biologists accept that the
same fundamental mechanisms (mutation, selection, drift, gene flow)
underlie both scales, but that emergent properties (species selection,
developmental constraints, contingency) shape macroevolutionary
patterns.

\subsubsection{Punctuated Equilibrium}\label{punctuated-equilibrium}

\textbf{Punctuated equilibrium} \citep{gould1972} proposes that most
species experience long periods of morphological stasis (millions of
years) interrupted by brief episodes of rapid evolutionary change
associated with speciation events. This pattern contrasts with
\textbf{phyletic gradualism}, which predicts slow, continuous
morphological change.

Evidence supporting punctuated equilibrium: - The fossil record
frequently shows species appearing abruptly (in geological terms),
persisting with little change, and disappearing suddenly. - Stasis is
the dominant pattern for many well-documented fossil lineages
(bryozoans, trilobites, mollusks). - Rapid change during speciation is
consistent with peripatric speciation in small, isolated populations
where fossil preservation is unlikely.

\paragraph{Stasis as the dominant pattern in the fossil
record}\label{stasis-as-the-dominant-pattern-in-the-fossil-record}

The most surprising aspect of \citet{gould1972}'s argument was not that
change can be rapid but that \textbf{stasis is so prevalent}.
Quantitative analyses of well-sampled fossil sequences confirm: most
species, once established, exhibit little net morphological change for
millions of years. Cheetham's analysis of Caribbean bryozoan
(\emph{Metrarabdotos}) species across the past 15 million years measured
46 morphological traits in 1,200 specimens; lineages persisted for 2--10
million years with mean morphological drift below the standard deviation
of within-population variation. Williamson's snail (\emph{Bulinus})
sequences from Lake Turkana show prolonged stasis broken by rapid bursts
of change correlated with regional climate shifts. The \textbf{stasis
itself requires explanation} --- even with shifting environments and
ongoing mutation, species somehow remain morphologically constrained.
Proposed mechanisms include strong stabilizing selection, gene-flow
homogenization across the species range, developmental canalization
(\cref{sec:unit_VI_evolution_and_selection}), and frequency-dependent
selection at species boundaries.

\paragraph{Correlated evolutionary
change}\label{correlated-evolutionary-change}

A complementary observation is that morphological change, when it
occurs, is \textbf{correlated across traits}. The Galapagos finch beak
responds to drought as an integrated unit (depth, width, and length most
change together), reflecting shared developmental and genetic
architecture. In the Cheetham bryozoan data, the burst of change at
speciation events involves multiple coordinated character shifts ---
rarely a single trait changing alone. This pattern has two
interpretations: (i) genetic correlations among traits cause selection
on one trait to produce correlated responses in others, as predicted by
the multivariate breeder's equation; (ii) speciation events reorganize
developmental modules in coordinated ways. Both processes likely
contribute, and disentangling them is an active research program in
evolutionary developmental biology.

The debate between punctuated equilibrium and gradualism is largely
resolved: both patterns occur, and the dominant mode varies among
lineages and environments. The significance of punctuated equilibrium
lies not in negating Darwinian mechanisms but in recognizing that the
tempo of evolution is uneven -- periods of rapid change (often
associated with speciation events in peripheral populations) alternate
with long intervals of morphological conservatism (stabilizing selection
maintaining adaptation to a stable environment).

\subsubsection{Evolutionary Novelties and Key
Innovations}\label{evolutionary-novelties-and-key-innovations}

\textbf{Key innovations} are traits that open new adaptive zones,
enabling diversification. Examples include:

\begin{itemize}
\tightlist
\item
  The evolution of \textbf{jaws} in gnathostome vertebrates (about 450
  Mya), enabling active predation and driving the diversification of
  jawed vertebrates
\item
  The evolution of the \textbf{amniotic egg} (about 340 Mya), freeing
  tetrapods from dependence on water for reproduction
\item
  The evolution of \textbf{flowers} in
  \hyperref[gl:angiosperm]{\textbf{angiosperm}}s (about 130 Mya),
  enabling coevolution with pollinators and explosive diversification
\item
  The evolution of \textbf{flight} in insects (about 350 Mya), birds
  (about 150 Mya), and bats (about 55 Mya), each opening vast ecological
  opportunities
\end{itemize}

Key innovations often trigger \textbf{adaptive radiations} (see
\cref{sec:unit_VI_genetic_drift_and_speciation}) by allowing lineages to
exploit previously inaccessible resources.

\subsubsection{Mass Extinctions}\label{mass-extinctions}

\begin{figure}[htbp]
\centering
\includegraphics[width=0.82\linewidth,height=4.2in,keepaspectratio]{../figures/mermaid_inline/inline_mermaid_0106_6bf79bf45270.png}
\caption{Geological timeline of life's major milestones and the five mass extinctions, with the proportion of species lost at each event.}
\end{figure}

Five major mass extinctions punctuate the history of life:

\begin{enumerate}
\def\labelenumi{\arabic{enumi}.}
\tightlist
\item
  \textbf{End-Ordovician} (about 445 Mya): Glaciation of Gondwana caused
  sea-level drop; approximately 85\% of marine species lost. Two pulses
  of extinction separated by approximately 1 million years.
\item
  \textbf{Late Devonian} (about 375 Mya): Prolonged extinction over
  approximately 20 million years; 75\% of species lost. Causes debated
  -- possible anoxia, climate change, or bolide impact.
\item
  \textbf{End-Permian} (about 252 Mya): The ``Great Dying'' --
  approximately 96\% of marine species and 70\% of terrestrial
  vertebrate species lost. Linked to Siberian Traps volcanism (massive
  flood basalts releasing CO\(_2\), SO\(_2\), and methane, causing ocean
  acidification, anoxia, and extreme warming).
\item
  \textbf{End-Triassic} (about 201 Mya): Approximately 80\% of species
  lost. Associated with Central Atlantic Magmatic Province volcanism.
  Cleared ecological space for dinosaur diversification.
\item
  \textbf{End-Cretaceous (KPg)} (about 66 Mya): Approximately 75\% of
  species lost, including most non-avian dinosaurs. Caused by the
  Chicxulub asteroid impact (Yucatan Peninsula, Mexico) producing a 180
  km diameter crater. Impact winter, acid rain, and global fires
  devastated ecosystems. Deccan Traps volcanism may have contributed to
  pre-impact stress.
\end{enumerate}

\textbf{Current Holocene/Anthropocene extinction}: Background extinction
rate is approximately 0.1--1 species per million species-years. Current
rates are estimated at 100--1,000 times higher due to habitat
destruction, overexploitation, pollution, invasive species, and climate
change. The International Union for Conservation of Nature (IUCN)
estimates that approximately 28\% of assessed species are threatened
with extinction.

\subsubsection{The Modern Extended Evolutionary
Synthesis}\label{the-modern-extended-evolutionary-synthesis}

The \textbf{Modern Synthesis} (1930s--1950s) unified Darwinian natural
selection with Mendelian genetics and population genetics. The
\textbf{Extended Evolutionary Synthesis} (EES) proposes additional
mechanisms that may play significant evolutionary roles:

\begin{itemize}
\tightlist
\item
  \textbf{Developmental plasticity}: A single genotype can produce
  multiple phenotypes in response to environmental conditions.
  Plasticity may facilitate (or hinder) evolutionary change by exposing
  hidden genetic variation to selection.
\item
  \textbf{Epigenetic inheritance}: Heritable changes in gene expression
  not encoded in DNA sequence (DNA methylation,
  \hyperref[gl:histone]{\textbf{histone}} modifications) can persist
  across generations in some organisms.
\item
  \textbf{Niche construction}: Organisms modify their own environments,
  altering the selective pressures they and other species experience.
  Beaver dams, earthworm soil modification, and human agriculture are
  examples.
\item
  \textbf{Cultural evolution}: In species with social learning,
  culturally transmitted behaviors can evolve and influence genetic
  evolution. Tool use in chimpanzees, song dialects in birds, and the
  entirety of human cultural evolution exemplify this process.
\end{itemize}

The status of the EES is actively debated. Laland \emph{et al.} (2015,
\emph{Nature}) argued that the EES has distinctive empirical content
because plasticity, inclusive inheritance, and niche construction cause
biased variation, motivate organisms as agents, and make evolution more
predictable than random mutation and selection alone. Wray \emph{et al.}
(2014, \emph{Evolution}) and others contested that the original Modern
Synthesis already has sufficient conceptual flexibility to accommodate
these phenomena within its mathematical framework. Pigliucci \& Müller
(2010, \emph{Cambridge University Press}) provided the most
comprehensive EES framework.

{\def\LTcaptype{none} % do not increment counter
\begin{longtable}[]{@{}
  >{\raggedright\arraybackslash}p{(\linewidth - 4\tabcolsep) * \real{0.1184}}
  >{\raggedright\arraybackslash}p{(\linewidth - 4\tabcolsep) * \real{0.3816}}
  >{\raggedright\arraybackslash}p{(\linewidth - 4\tabcolsep) * \real{0.5000}}@{}}
\toprule\noalign{}
\begin{minipage}[b]{\linewidth}\raggedright
Dimension
\end{minipage} & \begin{minipage}[b]{\linewidth}\raggedright
Classic Modern Synthesis (MS)
\end{minipage} & \begin{minipage}[b]{\linewidth}\raggedright
Extended Evolutionary Synthesis (EES)
\end{minipage} \\
\midrule\noalign{}
\endhead
\bottomrule\noalign{}
\endlastfoot
\textbf{Mechanism focus} & Selection on random genetic mutations & Plus:
developmental bias, niche construction,
\hyperref[gl:epigenetics]{\textbf{epigenetics}}, learning \\
\textbf{Inheritance} & Genetic (DNA sequence) & Plus: epigenetic,
cultural, ecological \\
\textbf{Phenotypic variation} & Genotype \texorpdfstring{\ensuremath{\to}}{->} phenotype (near 1-to-1
mapping) & Many genotypes \texorpdfstring{\ensuremath{\to}}{->} many phenotypes (reaction norms,
polyphenism) \\
\textbf{Evolutionary tempo} & Gradual (Fisher's infinitesimal model) &
Plus: rapid phenotypic evolution via plasticity \texorpdfstring{\ensuremath{\to}}{->} genetic
assimilation \\
\textbf{Organism role} & Passive target of selection & Active
evolvability agent (niche construction, behavior-driven evolution) \\
\textbf{Example} & Antibiotic resistance via mutation & Sticklebacks
rapidly evolving freshwater morphs via plastic bone density \\
\end{longtable}
}

\paragraph{Niche construction in
detail}\label{niche-construction-in-detail}

\textbf{Niche construction} --- organisms modifying their environment in
ways that change selective pressures on themselves and other species ---
is one of the most empirically grounded EES claims. Rather than viewing
the environment as a fixed selective filter, niche-construction theory
(Odling-Smee, Laland, Feldman 2003) treats the
\textbf{organism--environment relationship as bidirectional}: organisms
inherit not only genes but also a modified environment from their
ancestors (an ``ecological inheritance''). Examples:

\begin{itemize}
\tightlist
\item
  \textbf{Beaver dams} create wetland environments that persist for
  centuries and modify selection on the beavers themselves (favoring
  traits useful in lodge-based aquatic life), on aquatic plants and fish
  (which evolve in response to slowed water and altered hydrology), and
  on terrestrial herbivores in surrounding meadows.
\item
  \textbf{Earthworms} modify soil pH, structure, drainage, and nutrient
  cycling. Plants in earthworm-modified soils experience radically
  different selective regimes than they would in unmodified soils.
  Earthworm presence has measurably changed plant community composition
  and microbial diversity over millennia.
\item
  \textbf{Cyanobacteria} transformed Earth's atmosphere from anoxic to
  oxygenic over the Great Oxidation Event (about 2.4 Ga),
  simultaneously enabling the evolution of aerobic respiration and
  constraining (or extinguishing) the anaerobic life forms that had
  previously dominated. The largest niche-construction event in Earth
  history.
\item
  \textbf{Human agriculture} is the most extreme recent example: humans
  modified plant and animal genomes through artificial selection while
  simultaneously modifying their own selective environment (creating new
  disease pressures from livestock, dietary shifts, sedentary
  lifestyles).
\end{itemize}

\paragraph{Phenotypic plasticity and developmental
induction}\label{phenotypic-plasticity-and-developmental-induction}

\textbf{Phenotypic plasticity} is the production of different phenotypes
by the same genotype in response to different environments. The same
caterpillar genotype can produce a green form in summer and a brown form
in autumn (\emph{Nemoria arizonaria}). The same \emph{Daphnia} genotype
produces a helmeted, defended form in the presence of fish predator cues
and an unhelmeted form in their absence.

Plasticity matters for evolution because \textbf{plastic responses can
later be assimilated by genetic change} (genetic assimilation,
Waddington 1953). When a plastic response is consistently advantageous,
selection will favor genetic backgrounds that produce that response
constitutively --- even when the original environmental cue is absent.
This pathway, sometimes called ``\textbf{plasticity-led evolution},''
allows populations to respond to novel environments faster than
mutation-led evolution alone would predict. Documented cases include
rapid evolution of freshwater stickleback morphs from anadromous
ancestors, and the freshwater colonization of marine sticklebacks where
ancestral plasticity in body armor preceded eventual genetic fixation of
low-armor freshwater forms.

\paragraph{Epigenetic inheritance across
generations}\label{epigenetic-inheritance-across-generations}

Most epigenetic marks are reset during gametogenesis and embryogenesis,
but documented examples of \textbf{transgenerational epigenetic
inheritance} (TEI) exist:

\begin{itemize}
\tightlist
\item
  \textbf{\emph{A. thaliana FWA}} epi-allele: an unmethylated state at
  the \emph{FWA} locus produces flowering-time variation that is
  heritable across many generations without DNA sequence change.
\item
  \textbf{Dutch Hunger Winter cohort}: Individuals exposed prenatally to
  severe famine (1944--1945) showed altered DNA methylation patterns at
  \emph{IGF2} and other loci, and increased risk of metabolic disease in
  adulthood. Limited evidence suggests these patterns may persist into
  the next generation.
\item
  \textbf{\emph{C. elegans}} small-RNA inheritance: Specific small RNAs
  can be transmitted across generations in \emph{C. elegans}, where they
  have been documented to silence target genes for over 80 generations
  under some conditions.
\item
  \textbf{Mammalian paramutation}: At a few loci (\emph{Kit} in mice,
  \emph{agouti viable yellow}), parental phenotypes influence offspring
  phenotypes through epigenetic mechanisms despite identical DNA
  sequences.
\item
  \textbf{Rodent transgenerational stress and odor-conditioning}: In the
  Dias \& Ressler (2014) experiments, male mice conditioned to fear
  acetophenone odor produced F1 and F2 offspring (rarely themselves
  exposed) with elevated startle to that specific odor and detectable
  changes in the olfactory bulb glomerular map. The effect attenuates by
  F3, consistent with \textbf{2--3 generations of detectable
  transgenerational transmission} in rodents --- a pattern recurring
  across stress, dietary-restriction, and toxin-exposure paradigms in
  the rodent transgenerational literature.
\end{itemize}

The evolutionary significance of TEI in mammals (especially humans)
remains debated --- the pervasive epigenetic reprogramming in early
embryos limits its scope, and most documented mammalian effects fade
within 2--3 generations. Nonetheless, the existence of even rare TEI
events forces a more nuanced view of inheritance than the strictly
genetic view of the Modern Synthesis.

\begin{quote}
\textbf{Real-World Connection: The Sixth Mass Extinction}

Current extinction rates vastly exceed background levels. The Living
Planet Report 2024 documents a 73\% average decline in monitored
vertebrate population indices since 1970; this is a trend in monitored
populations, not a census of wild animals \citep{wwf2024livingplanet}.
Unlike previous mass extinctions driven by geological or astronomical
events, the current crisis is driven largely by human land use,
exploitation, pollution, invasive species, and climate change.
Conservation biology applies evolutionary principles to preserve
biodiversity: maintaining genetic diversity within populations,
preserving evolutionary potential, and protecting phylogenetically
distinct lineages.
\end{quote}

\begin{quote}
\textbf{Concept Check 5:} Punctuated equilibrium predicts long periods
of stasis interrupted by rapid change. How would you distinguish this
pattern from gradual change that simply was not preserved in the fossil
record?
\end{quote}

\begin{center}\rule{0.5\linewidth}{0.5pt}\end{center}

\subsection{Evolution Observed in Real
Time}\label{evolution-observed-in-real-time}

One of the most powerful confirmations of evolutionary theory is the
direct observation of evolution occurring within human timescales.

\subsubsection{Laboratory Evolution}\label{laboratory-evolution}

\textbf{Richard Lenski's Long-Term Evolution Experiment (LTEE)}: Since
1988, 12 populations of \emph{Escherichia coli} have been propagated
daily in minimal glucose medium at Michigan State University. By 2024,
these populations have undergone over 80,000 generations. Key results
include:

\begin{itemize}
\tightlist
\item
  Fitness increased rapidly in the first 2,000 generations, then more
  slowly -- consistent with Fisher's geometric model of adaptation.
\item
  Around generation 31,500, one population evolved the ability to
  metabolize citrate under aerobic conditions -- a trait previously
  considered a defining characteristic distinguishing \emph{E. coli}
  from other enteric bacteria. This required multiple mutations
  occurring in a specific order, demonstrating the role of historical
  contingency in evolution.
\item
  Mutations in mutator genes (DNA repair) increased the mutation rate in
  some lineages, accelerating adaptation but also accumulating slightly
  deleterious mutations.
\end{itemize}

\subsubsection{Natural Experiments}\label{natural-experiments}

\textbf{Darwin's finch beak evolution}: Peter and Rosemary Grant's
40-year field study on Daphne Major island documented directional
selection on beak depth in medium ground finches (\emph{Geospiza
fortis}) during the 1977 drought (beak depth increased 4\%) and the 1983
El Nino (beak depth decreased as small seeds became abundant again).
This provided real-time quantitative evidence that selection can produce
measurable morphological change within a single generation.

\textbf{Anolis lizard evolution}: Jonathan Losos and colleagues
demonstrated rapid evolutionary divergence in Caribbean anole lizards.
When populations of \emph{Anolis sagrei} were experimentally introduced
to small islands, they evolved shorter hindlimbs within approximately 10
generations -- an adaptation to the narrower perching surfaces available
on small islands versus the mainland.

\textbf{Cane toad invasion of Australia}: Since their introduction in
1935, cane toads (\emph{Rhinella marina}) have expanded across northern
Australia at accelerating rates -- the invasion front advanced
approximately 10 km per year in the 1940s but now advances over 55 km
per year. Toads at the invasion front have evolved longer legs and
greater endurance compared to toads in established populations -- they
are literally evolving to be better invaders. This is \textbf{spatial
sorting}: individuals at the front encounter and mate with other
front-runners, concentrating dispersal-enhancing alleles at the leading
edge.

Meanwhile, native predators in newly colonized areas have evolved
rapidly in response. Red-bellied black snakes (\emph{Pseudechis
porphyriacus}) and green tree snakes (\emph{Dendrelaphis punctulatus})
have evolved smaller heads relative to body size (reducing their ability
to consume the large, toxic toads -- a behavioral and morphological
adaptation that reduces mortality). Some snake populations have also
evolved increased physiological resistance to toad toxins
(bufadienolides). This represents rapid, reciprocal coevolutionary
response to an invasive species.

\begin{quote}
\textbf{Concept Check 6:} The LTEE citrate-utilization mutation required
multiple genetic changes occurring in a specific sequence. How does this
observation relate to the concept of historical contingency in
evolution? Would you expect the same mutation to arise in most 12
replicate populations?
\end{quote}

\subsubsection{Punctuated Equilibrium versus Phyletic Gradualism:
Empirical
Signatures}\label{punctuated-equilibrium-versus-phyletic-gradualism-empirical-signatures}

The choice between \textbf{punctuated equilibrium} and \textbf{phyletic
gradualism} is not strictly either/or --- both modes are observed in
different lineages, and the empirical signatures of each differ:

{\def\LTcaptype{none} % do not increment counter
\begin{longtable}[]{@{}
  >{\raggedright\arraybackslash}p{(\linewidth - 4\tabcolsep) * \real{0.1636}}
  >{\raggedright\arraybackslash}p{(\linewidth - 4\tabcolsep) * \real{0.3636}}
  >{\raggedright\arraybackslash}p{(\linewidth - 4\tabcolsep) * \real{0.4727}}@{}}
\toprule\noalign{}
\begin{minipage}[b]{\linewidth}\raggedright
Feature
\end{minipage} & \begin{minipage}[b]{\linewidth}\raggedright
Phyletic gradualism
\end{minipage} & \begin{minipage}[b]{\linewidth}\raggedright
Punctuated equilibrium
\end{minipage} \\
\midrule\noalign{}
\endhead
\bottomrule\noalign{}
\endlastfoot
\textbf{Tempo} & Slow, continuous change & Long stasis interrupted by
rapid bursts \\
\textbf{Locus of change} & Throughout lineage history & Concentrated at
speciation events \\
\textbf{Population structure} & Large, panmictic populations & Small
peripheral isolates (peripatric) \\
\textbf{Predicted fossil pattern} & Smooth morphological intermediates &
Abrupt transitions between successive forms \\
\textbf{Stasis explanation} & Apparent stasis = preservation gaps &
Genuine evolutionary stasis from stabilizing selection \\
\textbf{Speciation linkage} & Speciation not associated with
morphological change & Speciation drives morphological change \\
\end{longtable}
}

\citet{gould1972} introduced punctuated equilibrium based on
observations of marine invertebrate fossil sequences (bryozoans,
trilobites) where species persisted essentially unchanged for millions
of years, then were replaced by morphologically distinct relatives
without obvious gradual transitional sequences. Critics initially
countered that the gaps were artifacts of incomplete preservation.
Subsequent quantitative work has shown that \textbf{stasis is real} in
many groups: dense fossil sequences from sites with high preservation
(Caribbean coral reefs, Plio-Pleistocene foraminifera) confirm that
species often remain morphologically constant for geologic intervals.
\textbf{Phyletic gradualism} is also observed --- particularly in
foraminifera and some mammalian lineages --- but punctuated equilibrium
with extended stasis is the dominant pattern across well-sampled
lineages.

The biological mechanism for stasis remains debated: stabilizing
selection on a well-adapted phenotype, \textbf{canalization} of
development that buffers genetic variation, \textbf{genetic constraints}
that limit morphological evolution, or simply persistence in widespread
species that average across local selective pressures. The mechanism of
\textbf{punctuation} is more clearly understood --- peripatric
speciation in small isolated populations allows rapid morphological
change through founder effects and altered selective regimes, after
which the new species expands and stabilizes.

\subsubsection{Species Selection: Macroevolution as a Higher-Level
Process}\label{species-selection-macroevolution-as-a-higher-level-process}

If species themselves vary in heritable properties that affect their own
propagation and persistence, then \textbf{species selection} can occur
--- analogous to natural selection, but operating on species rather than
on individuals. Species selection requires that:

\begin{enumerate}
\def\labelenumi{\arabic{enumi}.}
\tightlist
\item
  \textbf{Variation among species} in traits affecting either speciation
  rate or extinction rate.
\item
  \textbf{Heritability}: daughter species inherit those traits from
  parent species.
\item
  \textbf{Differential reproduction/persistence}: species with favored
  traits diversify more rapidly or persist longer.
\end{enumerate}

Examples that may reflect species selection (though contested in the
literature):

\begin{itemize}
\tightlist
\item
  \textbf{Dispersal ability}: Marine invertebrates with
  \textbf{planktotrophic larvae} (broadly dispersing) speciate slowly
  but persist long; species with \textbf{direct development} (poorly
  dispersing) speciate rapidly but go extinct often. This produces a
  macroevolutionary trade-off in lineage diversification.
\item
  \textbf{Asexual versus sexual reproduction}: Asexual lineages tend to
  be evolutionarily ephemeral despite outcompeting sexual relatives in
  the short term. Species-rich and ancient clades are overwhelmingly
  sexual --- suggesting species selection has favored sexuality despite
  its twofold cost.
\item
  \textbf{Body size}: In many lineages, body size shows a long-term
  increase (``Cope's rule'') that cannot easily be explained by
  individual-level selection alone. Species selection on extinction risk
  may contribute.
\end{itemize}

Whether species selection is a genuinely separate force from cumulative
individual-level selection remains controversial. The strongest case is
for traits like \textbf{dispersal ability} that affect
speciation/extinction rates intrinsically rather than individual fitness
--- the property is irreducibly species-level. Under this view,
microevolutionary mechanisms still produce most evolutionary change, but
emergent macroevolutionary patterns reflect higher-order selection on
species-level properties.

These examples of real-time evolution are important not only as
scientific confirmations but also as pedagogical tools: they demonstrate
that evolution is not merely a historical narrative about the distant
past but a present, ongoing, and experimentally testable process.

\begin{quote}
\textbf{Clinical Connection --- Antibiotic resistance as natural
selection in hospital time.} \emph{Staphylococcus aureus} isolates from
UK hospitals have been genome-sequenced longitudinally since 1960. Each
new \texorpdfstring{\ensuremath{\beta}}{beta}-lactam antibiotic introduced (penicillin \texorpdfstring{\ensuremath{\to}}{->} methicillin \texorpdfstring{\ensuremath{\to}}{->}
vancomycin \texorpdfstring{\ensuremath{\to}}{->} linezolid \texorpdfstring{\ensuremath{\to}}{->} daptomycin) triggered the spread of resistance
alleles within 2--15 years, driven by
\hyperref[gl:plasmid]{\textbf{plasmid}}-borne \emph{mecA}, \emph{vanA},
or chromosomal mutations. This is directional selection with an imposed
fitness landscape whose peaks move with every new drug. Combination
therapy (two antibiotics targeting different molecular machinery) slows
resistance because a bacterium now requires two simultaneous mutations
--- the product of two low probabilities. The same Fisherian selection
mathematics that describes pepper-moth evolution quantitatively predicts
multi-drug regimen design in infectious disease and cancer.
\end{quote}

\begin{quote}
\textbf{Clinical Connection --- Darwinian oncology.} A tumour is a
population of cells with heritable genetic variation and differential
reproductive success under treatment pressure --- a Darwinian system.
Classic cytotoxic chemotherapy selects rapidly-dividing cells; the
surviving slow-cyclers (the \emph{persister} population) become the seed
for relapse. Modern ``adaptive therapy'' exploits frequency-dependent
selection: by intermittently \emph{withholding} drug once the tumour
shrinks, a sensitive subpopulation is kept alive to competitively
suppress resistant clones. Trials in metastatic prostate cancer
(abiraterone cycling) show median progression-free survival gains of
10--18 months versus continuous dosing, directly derived from
evolutionary game theory.
\end{quote}

\begin{center}\rule{0.5\linewidth}{0.5pt}\end{center}

\subsection{Computational Bridge}\label{computational-bridge-18}

Mean fitness \(\bar W\) as a function of \(p\) is tabulated by
\breaktt{fitness\_landscape\_1d}:

\begin{Shaded}
\begin{Highlighting}[]
\ImportTok{from}\NormalTok{ biology.evolution }\ImportTok{import}\NormalTok{ fitness\_landscape\_1d}

\NormalTok{p\_vals }\OperatorTok{=}\NormalTok{ [}\FloatTok{0.0}\NormalTok{, }\FloatTok{0.25}\NormalTok{, }\FloatTok{0.5}\NormalTok{, }\FloatTok{0.75}\NormalTok{, }\FloatTok{1.0}\NormalTok{]}
\NormalTok{w\_bar }\OperatorTok{=}\NormalTok{ fitness\_landscape\_1d(p\_vals, fitness\_AA}\OperatorTok{=}\FloatTok{1.0}\NormalTok{, fitness\_Aa}\OperatorTok{=}\FloatTok{0.95}\NormalTok{, fitness\_aa}\OperatorTok{=}\FloatTok{0.7}\NormalTok{)}
\BuiltInTok{print}\NormalTok{([}\BuiltInTok{round}\NormalTok{(w, }\DecValTok{4}\NormalTok{) }\ControlFlowTok{for}\NormalTok{ w }\KeywordTok{in}\NormalTok{ w\_bar])}
\end{Highlighting}
\end{Shaded}

\begin{quote}
\textbf{Clinical / systems note:} Antimicrobial stewardship explicitly
manipulates \(s\) and \(N_e\) (host population colonisation load) to
keep resistance alleles from fixing --- the same selection algebra
taught for moths and finches.
\end{quote}

\begin{center}\rule{0.5\linewidth}{0.5pt}\end{center}

\subsubsection{Evolution in the Synthetic Age: CRISPR Gain-of-Function
Research and Evolutionary
Ethics}\label{evolution-in-the-synthetic-age-crispr-gain-of-function-research-and-evolutionary-ethics}

The same experimental tools that make evolution a tractable laboratory
science --- CRISPR-Cas9 genome engineering, directed evolution via
phage-assisted continuous evolution (PACE), and \emph{in silico}
fitness-landscape modelling --- have forced evolutionary biologists into
ethics that Darwin rarely anticipated. The 2011--2012 \textbf{H5N1
gain-of-function experiments} (Fouchier lab, Kawaoka lab) demonstrated
that primarily a handful of spike-protein mutations were needed to make
avian influenza airborne-transmissible in ferrets --- effectively
\emph{running a natural experiment in the laboratory} on a pathogen that
kills about 60 \% of infected humans. The dual-use research of
concern (DURC) framework that emerged from the ensuing moratorium
(2012--2014, 2017--present) now governs any research whose foreseeable
output is ``enhanced potential pandemic pathogens.''

Beyond viruses, the ethical landscape expands rapidly: \textbf{gene
drives} (Crisanti lab, \emph{Nat. Biotechnol.} 2018) use CRISPR-Cas9 to
bias inheritance toward a designed allele at rates approaching 99 \% ---
demonstrated in caged \emph{Anopheles gambiae} populations eradicated in
7--11 generations. Gene drives could plausibly eliminate malaria, but
also represent an irrevocable alteration of an evolutionary trajectory
involving an entire species. \textbf{CRISPR germline editing} in humans
(He Jiankui, 2018) targeting CCR5 to confer HIV resistance produced two
living children (Lulu and Nana) with off-target effects in the same gene
of unknown phenotypic consequence, and generated a unanimous
international moratorium on heritable genome editing.
\textbf{Resurrection biology} (Colossal Biosciences woolly-mammoth
project, 2024) uses Asian-elephant fibroblast editing to re-create
approximations of extinct phenotypes, provoking debate about whether a
mammoth is ``brought back'' or ``created anew.''

The common thread: the tools of evolutionary biology are now potent
enough that the theory is no longer merely descriptive. Students
entering this field must internalise that evolutionary theory is
actively being used as a design framework (for directed evolution of
enzymes, for gene-drive deployment, for pandemic preparedness) ---
evolutionary literacy now carries an ethical dimension inseparable from
the scientific one. The Asilomar Conference (1975), the DURC frameworks
(2012 onward), and the heritable genome-editing summits (2015, 2018,
2023) represent the biology community confronting responsibilities that
basic evolutionary training must prepare us for.

\begin{center}\rule{0.5\linewidth}{0.5pt}\end{center}

\subsection{Current Evidence and Frontier
Biology}\label{current-evidence-and-frontier-biology-21}

For \textbf{Evolution --- Theory, Natural Selection, and Adaptation},
frontier biology belongs inside the evidence logic of the chapter.
Evolutionary claims are strongest when they combine mechanism,
comparative evidence, population process, and explicit uncertainty. The
core reading question is this: selection claims need fitness components,
ecological context, genetic variation, and alternative hypotheses.

\begin{itemize}
\tightlist
\item
  \textbf{What to verify:} identify the observation, model, assay, or
  dataset that would make the claim stronger or weaker.
\item
  \textbf{What to qualify:} state the scale, organism, cell type,
  environmental condition, or population where the claim is expected to
  hold.
\item
  \textbf{What to compare:} test at least one alternative explanation,
  baseline, or null model before treating the pattern as causal.
\item
  \textbf{What to cite:} distinguish primary evidence, review synthesis,
  public dataset, and institutional guidance; for recent or numeric
  claims, prefer the source closest to the measurement and state what
  has changed since it was published.
\end{itemize}

Distinguish adaptation from drift, phylogenetic signal from convergence,
and historical explanation from a testable prediction about present-day
data.

\textbf{Source practice:} For evolutionary claims, prefer evidence that
compares alternatives such as selection, drift, gene flow, constraint,
convergence, and shared ancestry.

\subsection{Summary}\label{summary-21}

\begin{itemize}
\tightlist
\item
  Pre-Darwinian thinkers (Lamarck, Cuvier, Lyell) established concepts
  of species change, extinction, and deep geological time that set the
  stage for Darwin and Wallace.
\item
  Five independent lines of evidence support evolution: fossil record
  (transitional forms), comparative anatomy (homology, vestigiality),
  embryology (shared developmental stages), biogeography (geographic
  distributions reflecting history), and molecular biology (sequence
  similarity, shared pseudogenes, ERVs).
\item
  Natural selection requires heritable variation, differential
  reproduction, and environmental pressure. It operates through multiple
  modes: directional, stabilizing, disruptive, sexual, kin, and
  frequency-dependent selection.
\item
  Fitness (\(W\)) is relative reproductive success. Hamilton's rule
  (\(rB > C\)) extends selection theory to explain altruism among
  relatives.
\item
  Adaptation is constrained by phylogenetic history, developmental
  programs, and trade-offs. Exaptations demonstrate co-option of
  existing structures for new functions.
\item
  Coevolution produces arms races (newt-snake toxin resistance) and
  maintains the Red Queen dynamic.
\item
  Macroevolutionary patterns include punctuated equilibrium, mass
  extinctions (5 major events plus the ongoing 6th), and the Extended
  Evolutionary Synthesis incorporating plasticity, epigenetics, niche
  construction, and cultural evolution.
\item
  \textbf{Connections:} See \cref{sec:unit_V_population_genetics} for
  formal \(p\) recursions,
  \cref{sec:unit_VI_genetic_drift_and_speciation} for drift
  vs.~selection, and Unit X for spatial sorting and range expansion.
\end{itemize}

\begin{center}\rule{0.5\linewidth}{0.5pt}\end{center}

\subsection{Key Terms}\label{key-terms-21}

{\def\LTcaptype{none} % do not increment counter
\begin{longtable}[]{@{}
  >{\raggedright\arraybackslash}p{(\linewidth - 2\tabcolsep) * \real{0.2857}}
  >{\raggedright\arraybackslash}p{(\linewidth - 2\tabcolsep) * \real{0.7143}}@{}}
\toprule\noalign{}
\begin{minipage}[b]{\linewidth}\raggedright
Term
\end{minipage} & \begin{minipage}[b]{\linewidth}\raggedright
Definition
\end{minipage} \\
\midrule\noalign{}
\endhead
\bottomrule\noalign{}
\endlastfoot
\textbf{Natural selection} & Differential survival and reproduction of
individuals based on heritable variation in fitness-related traits \\
\textbf{\hyperref[gl:fitness]{Fitness (w)}} & Relative reproductive
success of a genotype compared to the population mean \\
\textbf{Directional selection} & Selection favoring one phenotypic
extreme, shifting the population mean \\
\textbf{Stabilizing selection} & Selection favoring intermediate
phenotypes, reducing phenotypic variance \\
\textbf{Disruptive selection} & Selection favoring both phenotypic
extremes at the expense of intermediates \\
\textbf{Sexual selection} & Selection arising from variation in mating
success; includes intersexual (mate choice) and intrasexual
(competition) components \\
\textbf{Kin selection} & Natural selection favoring altruistic behaviors
toward genetic relatives, governed by Hamilton's rule (\(rB > C\)) \\
\textbf{Frequency-dependent selection} & Fitness of a phenotype depends
on its frequency in the population; negative form maintains
polymorphism \\
\textbf{Adaptation} & A heritable trait shaped by natural selection that
increases fitness in the current environment \\
\textbf{Exaptation} & A trait originally shaped by selection for one
function that has been co-opted for a different function \\
\textbf{Coevolution} & Reciprocal evolutionary change between
interacting species driven by mutual selective pressures \\
\textbf{Red Queen hypothesis} & Organisms must continually evolve to
maintain fitness relative to coevolving species \\
\textbf{Punctuated equilibrium} & Macroevolutionary pattern of long
stasis interrupted by rapid change during speciation \\
\textbf{Mass extinction} & Geologically rapid loss of a large percentage
of species; five major events documented \\
\textbf{Convergent evolution} & Independent evolution of similar
features in unrelated lineages subject to similar selective pressures \\
\textbf{Vestigial structure} & Reduced remnant of an organ that was
functional in an ancestor \\
\textbf{Homologous structures} & Features in different species sharing
common developmental and evolutionary origin \\
\textbf{Analogous structures} & Features in different species similar in
function but with different developmental and evolutionary origins \\
\textbf{Modern Synthesis} & Unification of Darwinian selection with
Mendelian genetics, population genetics, paleontology, and systematics
(1930s--1950s) \\
\textbf{Extended Evolutionary Synthesis} & Expansion of the Modern
Synthesis to include developmental plasticity, epigenetic inheritance,
niche construction, and cultural evolution \\
\textbf{Inclusive fitness} & Total fitness including direct reproduction
plus indirect fitness gained by helping relatives reproduce \\
\textbf{Phyletic gradualism} & Model proposing slow, continuous
morphological change throughout a lineage's history \\
\textbf{Uniformitarianism} & Principle that the same geological
processes operating today have typically operated at similar rates \\
\textbf{Catastrophism} & Doctrine that Earth's features and fossil
\hyperref[gl:succession]{\textbf{succession}} result from sudden,
violent events \\
\textbf{\hyperref[gl:selection-coefficient]{Selection coefficient (s)}}
& A measure of the fitness difference between genotypes; \(s = 0\) means
no selection, \(s = 1\) means the disfavored genotype is lethal \\
\textbf{Overdominance} & Heterozygote advantage; the
\hyperref[gl:heterozygous]{\textbf{heterozygous}} genotype has higher
fitness than either homozygote \\
\end{longtable}
}

\begin{center}\rule{0.5\linewidth}{0.5pt}\end{center}

\subsection{Review Questions}\label{review-questions-21}

\begin{enumerate}
\def\labelenumi{\arabic{enumi}.}
\item
  Darwin and Wallace independently proposed natural selection. Describe
  the three conditions necessary for natural selection to operate, and
  explain why Mendelian genetics was essential for completing the
  theory.
\item
  Explain how \emph{Tiktaalik roseae} serves as a transitional fossil.
  What specific features bridge the gap between fish and tetrapods? Why
  was its discovery predicted by paleontologists based on the age and
  type of rock formations?
\item
  Distinguish between homologous and analogous structures. For each,
  provide an example not discussed in this chapter, and explain how you
  would determine which category a given structure belongs to.
\item
  A population of moths is subject to directional selection with
  selection coefficient \(s = 0.05\) favoring a dominant allele
  currently at frequency \(p = 0.01\). Using
  \(\Delta p \approx sp(1-p)\), estimate the change in allele frequency
  after one generation. How many generations would be required for the
  allele to reach \(p = 0.50\) (qualitative answer)?
\item
  Explain why the recurrent laryngeal nerve in giraffes is considered
  evidence for evolution rather than design. What developmental and
  phylogenetic constraints explain its circuitous path?
\item
  Using Hamilton's rule, determine whether a sterile worker bee (\(C\) =
  lifetime reproduction about 50 offspring) should help her queen mother
  produce additional sisters (\(B\) = 100 additional sisters) given
  \(r = 0.75\) for full sisters in haplodiploid species. Show your
  calculation.
\item
  Compare and contrast punctuated equilibrium with phyletic gradualism.
  What types of fossil evidence would distinguish between these two
  models? Can both patterns coexist in the same clade?
\item
  The End-Permian extinction eliminated approximately 96\% of marine
  species. Describe two proposed causes and explain how mass extinctions
  can paradoxically promote long-term diversification by opening
  ecological opportunities for surviving lineages.
\item
  Define coevolution and provide an example of a coevolutionary arms
  race. How does the Red Queen hypothesis explain the maintenance of
  sexual reproduction in the face of its twofold cost?
\item
  The Extended Evolutionary Synthesis proposes mechanisms beyond those
  in the Modern Synthesis. Choose one (developmental plasticity,
  epigenetic inheritance, niche construction, or cultural evolution) and
  explain how it might influence evolutionary trajectories in ways not
  captured by standard population genetics models.
\item
  At which \(p\) is \(\bar W\) maximised for the fitness set in the
  bridge code? Relate to equilibrium under one-locus selection.
\item
  Give one example where \textbf{kin selection} predicts a behaviour
  that violates individual fitness maximisation yet spreads in
  haplodiploid hymenoptera. \#\# Further Reading and Source Notes
\end{enumerate}

\begin{itemize}
\tightlist
\item
  Darwin \& Wallace (1858). On the Tendency of Species to form
  Varieties. \emph{Journal of the Proceedings of the Linnean Society:
  Zoology}, 3.
\item
  Williams (1966). \emph{Adaptation and Natural Selection}. Princeton
  University Press.
\item
  Karn \& Penrose (1951). Birth weight and gestation time in relation to
  maternal age, parity and infant survival. \emph{Annals of Eugenics},
  16.
\item
  Hamilton \& Zuk (1982). Heritable true fitness and bright birds: A
  role for parasites?. \emph{Science}, 218.
\item
  Gould \& Eldredge (1972). Punctuated equilibria: an alternative to
  phyletic gradualism. \emph{Models in Paleobiology}.
\item
  Dobzhansky (1973). Nothing in biology makes sense except in the light
  of evolution. \emph{The American Biology Teacher}, 35.
\end{itemize}

\begin{center}\rule{0.5\linewidth}{0.5pt}\end{center}

\subsubsection{Companion Source
Module}\label{companion-source-module-21}

\textbf{Evolution --- Theory, Natural Selection, and Adaptation} should
leave a reproducible trail from a biological claim to the code, figure,
diagram, or paper-based activity that can test it. Use the surfaces
below to inspect the chapter's assumptions, rerun the relevant model, or
compare the manuscript explanation with companion labs and figures.

{\def\LTcaptype{none} % do not increment counter
\begin{longtable}[]{@{}
  >{\raggedright\arraybackslash}p{(\linewidth - 2\tabcolsep) * \real{0.5000}}
  >{\raggedright\arraybackslash}p{(\linewidth - 2\tabcolsep) * \real{0.5000}}@{}}
\toprule\noalign{}
\begin{minipage}[b]{\linewidth}\raggedright
Surface
\end{minipage} & \begin{minipage}[b]{\linewidth}\raggedright
Use it for
\end{minipage} \\
\midrule\noalign{}
\endhead
\bottomrule\noalign{}
\endlastfoot
\breaktt{src/biology/evolution/evolution.py}
(\breaktt{selection\_one\_generation}, \breaktt{simulate\_selection},
\breaktt{fitness\_landscape\_1d}) & Reproduce selection trajectories and
landscape reasoning. \\
\breaktt{src/visualization/plots.py}
(\breaktt{plot\_selection\_simulation}) & Compare fitness assumptions
with plotted allele-frequency change. \\
\breaktt{src/mermaid/biology\_diagrams.py}
(\breaktt{natural\_selection\_diagram}) & Keep variation, inheritance,
differential survival, and adaptation distinct. \\
\end{longtable}
}

\textbf{Reproducibility check:} define fitness component, environment,
heritable variation, and alternative explanation before calling a trait
adaptive. \textbf{Cross-reference:} use
\cref{sec:unit_V_population_genetics} and
\cref{sec:unit_VI_genetic_drift_and_speciation}.

\newpage

\section{Genetic Drift, Gene Flow, and
Speciation}\label{genetic-drift-gene-flow-and-speciation}

\label{sec:unit_VI_genetic_drift_and_speciation}

\begin{quote}
\textbf{Ch 20} \texorpdfstring{\ensuremath{\cdot}}{.} Level 3/3 \texorpdfstring{\ensuremath{\cdot}}{.} 60 min read \texorpdfstring{\ensuremath{\cdot}}{.} 75 min lecture \texorpdfstring{\ensuremath{\cdot}}{.}
Prerequisites: \cref{sec:unit_VI_evolution_and_selection}
\end{quote}

\subsection{Learning Objectives}\label{learning-objectives-22}

By the end of this chapter, you should be able to:

\begin{enumerate}
\def\labelenumi{\arabic{enumi}.}
\tightlist
\item
  Define \hyperref[gl:genetic-drift]{\textbf{genetic drift}} and explain
  why its effects are strongest in small populations, using the
  Wright-Fisher model mathematically.
\item
  Distinguish bottleneck and
  \hyperref[gl:founder-effect]{\textbf{founder effect}}s and describe
  their genetic consequences with real-world examples.
\item
  Calculate effective population size (\(N_e\)) and explain why it is
  typically much smaller than census population size.
\item
  Describe \hyperref[gl:gene]{\textbf{gene}} flow, its effects on
  population differentiation, and how \(F_{ST}\) quantifies population
  structure.
\item
  Compare and contrast species concepts (biological, morphological,
  phylogenetic, ecological, cohesion) and evaluate their strengths and
  limitations.
\item
  Explain allopatric, peripatric, parapatric, and
  \hyperref[gl:sympatric-speciation]{\textbf{sympatric speciation}}
  mechanisms, including the role of polyploidy and adaptive radiation.
\item
  Differentiate prezygotic from postzygotic reproductive barriers and
  explain Haldane's rule and reinforcement as outcomes of selection on
  isolation.
\item
  Calculate the fixation probability and mean fixation time of a neutral
  mutation under Kimura's neutral theory and explain why the neutral
  substitution rate is independent of population size.
\item
  Evaluate evidence for archaic introgression in modern humans,
  including Neanderthal and Denisovan ancestry, and explain how adaptive
  introgression can transfer beneficial alleles between lineages.
\end{enumerate}

\subsubsection{Study Blueprint}\label{study-blueprint-22}

\begin{itemize}
\tightlist
\item
  \textbf{Big idea:} Chance, population structure, and barriers to gene
  flow can generate divergence even without adaptive change.
\item
  \textbf{Core concepts:} drift, effective population size, gene flow,
  speciation.
\item
  \textbf{Framework alignment:} Vision \& Change: Evolution, Systems; AP
  Biology: Evolution, Systems Interactions; NGSS-style topics: Natural
  Selection and Evolution, Interdependent Relationships in Ecosystems.
\item
  \textbf{Model or quantitative lens:} Fixation probability, effective
  population size, and migration-selection balance.
\item
  \textbf{Data skill:} Distinguish stochastic from directional change in
  allele-frequency data.
\item
  \textbf{Practice cadence:} Visual Representations, Statistical Tests
  and Data Analysis, Argumentation.
\item
  \textbf{Common misconception to repair:} Random does not mean
  patternless; stochastic processes have predictable distributions.
\item
  \textbf{Primary lab:}
  \cref{sec:lab_unit_VI_genetic_drift_and_speciation}.
\item
  \textbf{Question bank:}
  \cref{sec:q_unit_VI_genetic_drift_and_speciation}.
\item
  \textbf{Transfer task:} Apply drift reasoning to endangered
  populations, founder effects, and island radiations.
\item
  \textbf{Bridge to computation:}
  \breaktt{biology.evolution.evolution.simulate\_drift}.
\end{itemize}

\begin{center}\rule{0.5\linewidth}{0.5pt}\end{center}

\begin{quote}
\textbf{Opening Vignette ---
\hyperref[gl:speciation]{\textbf{Speciation}} on the Underground}

Beneath London's streets, in the tunnels of the Underground railway
system, a mosquito silently diverged from its surface relatives.
\emph{Culex pipiens pipiens} lives above ground, hibernates through
winter, and feeds on birds. Its descendant, \emph{Culex pipiens
molestus}, rarely hibernates, bites mammals, and can breed without a
blood meal --- perfectly adapted to the warm, mammal-rich tunnels of the
Tube. Genetic analysis published in 1999 confirmed that the two are
reproductively isolated, despite sharing the same city. The London
Underground mosquito has become a textbook example of parapatric
speciation without geographic separation --- showing that
\hyperref[gl:reproductive-isolation]{\textbf{reproductive isolation}}
can evolve from a founder event and novel selection pressures even
within a few kilometers and a few hundred years. Tube passengers helped
document evolution in progress every time they got bitten.
\end{quote}

\subsection{Genetic Drift}\label{genetic-drift-1}

\subsubsection{Definition and
Significance}\label{definition-and-significance}

\textbf{Genetic drift} is the random change in
\hyperref[gl:allele]{\textbf{allele}} frequencies across generations
that results from the sampling error inherent in finite populations.
Unlike \hyperref[gl:natural-selection]{\textbf{natural selection}},
drift is \textbf{undirected} -- it is equally likely to increase or
decrease the frequency of any allele, regardless of its effect on
fitness. Drift is a consequence of probability, not adaptation.

The power of drift is inversely proportional to population size. In
large populations, the law of large numbers ensures that actual allele
frequencies closely match expected frequencies. In small populations,
random sampling can cause dramatic, unpredictable fluctuations -- and
ultimately, the fixation or loss of alleles.

\subsubsection{Mathematical Treatment}\label{mathematical-treatment}

Under the \textbf{Wright-Fisher model}, a population of \(N\) diploid
individuals contains \(2N\) gene copies at each locus. Each generation,
the new population is formed by randomly sampling \(2N\) alleles (with
replacement) from the current generation's gene pool.

If the current frequency of allele \(A\) is \(p\), the number of \(A\)
alleles in the next generation follows a binomial distribution:

\begin{equation}
P(k \mid p) = \binom{2N}{k} p^k (1-p)^{2N-k}
\label{eq:genetic_drift_and_speciation_1}
\end{equation}

Key properties:

\begin{itemize}
\tightlist
\item
  \textbf{Expected frequency}: \(E[p'] = p\) (drift has no directional
  bias)
\item
  \textbf{Variance per generation}:
  \(\text{Var}[\Delta p] = \frac{p(1-p)}{2N_e}\)
\end{itemize}

The variance expression reveals that doubling population size halves the
variance in allele frequency change per generation.

\textbf{Expected time to fixation}: For a neutral allele currently at
frequency \(p\), the probability of eventual fixation is simply \(p\)
(its current frequency). The expected time to fixation, conditional on
fixation occurring, is:

\begin{equation}
\bar{t}_{\text{fix}} \approx -4N_e \cdot \frac{(1-p)}{p} \cdot \ln(1-p)
\label{eq:genetic_drift_and_speciation_2}
\end{equation}

For a newly arisen neutral \hyperref[gl:mutation]{\textbf{mutation}}
(\(p = 1/2N\)), the expected time to fixation is approximately \(4N_e\)
generations.

\textbf{Heterozygosity decline}: Expected heterozygosity declines
geometrically:

\begin{equation}
H_t = H_0 \left(1 - \frac{1}{2N_e}\right)^t
\label{eq:genetic_drift_and_speciation_3}
\end{equation}

The half-life of heterozygosity is
\(t_{1/2} = N_e \ln 2 \approx 0.693 N_e\) generations.

\begin{figure}[htbp]
\centering
\includegraphics[width=0.82\linewidth,height=4.2in,keepaspectratio]{../figures/mermaid_inline/inline_mermaid_0107_b74cc0f94112.png}
\caption{State diagram for Mathematical Treatment showing transitions among p = 0.5, High variance, Allele likely fixed or lost, and p = 0.0 or 1.0.}
\end{figure}

\subsubsection{Effective Population
Size}\label{effective-population-size}

The \textbf{effective population size} (\(N_e\)) is the size of an ideal
Wright-Fisher population that would experience the same magnitude of
genetic drift as the actual population. In nearly most real populations,
\(N_e\) is substantially smaller than the census size \(N\) because of:

\textbf{Unequal sex ratio}: If the number of breeding females (\(N_f\))
differs from the number of breeding males (\(N_m\)):

\begin{equation}
\frac{1}{N_e} = \frac{1}{4N_f} + \frac{1}{4N_m}
\label{eq:genetic_drift_and_speciation_4}
\end{equation}

\subsubsection{Worked Example: Effective Population
Size}\label{worked-example-effective-population-size}

\textbf{Problem:} In a breeding colony of northern elephant seals,
extreme male-male competition limits mating opportunities to primarily
the most \hyperref[gl:dominant]{\textbf{dominant}} males. Suppose a
specific breeding group contains \(N_f = 40\) sexually mature females,
but primarily \(N_m = 1\) dominant ``beachmaster'' male successfully
mates with most of them. 1. What is the census size (\(N\)) of the
breeding adults? 2. What is the effective population size (\(N_e\))?

\textbf{Solution:}

\begin{enumerate}
\def\labelenumi{\arabic{enumi}.}
\item
  \textbf{Calculate the census size:}
  \[ N = N_f + N_m = 40 + 1 = 41 \text{ adults}  \label{eq:unit_VI_genetic_drift_and_speciation_item_1}\]
\item
  \textbf{Calculate the effective population size:} Using the rearranged
  formula \(N_e = \frac{4N_m N_f}{N_m + N_f}\):
  \[ N_e = \frac{4(1)(40)}{1 + 40}  \label{eq:unit_VI_genetic_drift_and_speciation_item_2}\]

  \[ N_e = \frac{160}{41} \approx 3.9  \label{eq:unit_VI_genetic_drift_and_speciation_item_3}\]
\end{enumerate}

Even though there are 41 breeding animals, the genetic diversity
transmitted to the next generation is equivalent to an ideal
Wright-Fisher population of \textbf{fewer than 4 individuals}. This
extreme reproductive skew severely limits the effective population size
and subjects the colony to powerful genetic drift.

\textbf{Variance in reproductive success}: If the variance in offspring
number (\(V_k\)) exceeds the mean (\(\bar{k} = 2\) for a stable
population):

\begin{equation}
N_e = \frac{4N - 2}{V_k + 2}
\label{eq:genetic_drift_and_speciation_5}
\end{equation}

Species with high reproductive skew (most individuals leave zero
offspring while a few leave many) have dramatically reduced \(N_e\).

\textbf{Fluctuating population size}: \(N_e\) is the harmonic mean of
population sizes across generations:

\begin{equation}
\frac{1}{N_e} = \frac{1}{t} \sum_{i=1}^{t} \frac{1}{N_i}
\label{eq:genetic_drift_and_speciation_6}
\end{equation}

The harmonic mean is dominated by the smallest values. A single
generation bottleneck of \(N = 10\) followed by 99 generations of
\(N = 10{,}000\) yields \(N_e \approx 909\) -- not the arithmetic mean
of about 9,900.

\textbf{Human effective population size}: Despite a current census size
of about 8 billion, human \(N_e \approx 10{,}000\)--\(15{,}000\)
(estimated from genomic diversity). This reflects severe ancestral
bottlenecks, including the Out-of-Africa event approximately 70,000
years ago.

\subsubsection{The Neutral Theory and Nearly Neutral
Theory}\label{the-neutral-theory-and-nearly-neutral-theory}

Motoo Kimura's \textbf{\hyperref[gl:neutral-theory]{neutral theory} of
molecular evolution} (1968) proposed that the majority of substitutions
observed at the molecular level are selectively neutral -- neither
advantageous nor deleterious. Under neutrality, the rate of substitution
equals the mutation rate (\(k = \mu\)), independent of population size.
This seemingly paradoxical result arises because, while drift fixes
neutral alleles more slowly in large populations
(\(\bar{t}_{\text{fix}} = 4N_e\)), more neutral mutations arise per
generation in large populations (\(2N_e \mu\)). These effects cancel
exactly.

Tomoko Ohta extended this to the \textbf{nearly neutral theory} (1973),
which recognizes that most mutations are slightly deleterious. The fate
of a slightly deleterious mutation depends on population size: if
\(|s| < 1/(2N_e)\), selection is too weak to overcome drift, and the
mutation behaves effectively as neutral. In small populations, more
mutations fall into this ``nearly neutral'' category, leading to faster
accumulation of slightly deleterious substitutions. This predicts that
species with small \(N_e\) (e.g., large-bodied vertebrates) should
accumulate more slightly deleterious mutations than species with large
\(N_e\) (e.g., bacteria) -- a prediction confirmed by comparative
genomics.

\subsubsection{Bottleneck Effect}\label{bottleneck-effect}

A \textbf{population bottleneck} occurs when a population undergoes a
dramatic, temporary reduction in size. The genetic consequences are
severe and often irreversible:

\begin{itemize}
\tightlist
\item
  Loss of rare alleles (alleles at low frequency are most likely to be
  lost by chance)
\item
  Reduction in heterozygosity
\item
  Random fixation of alleles unrelated to fitness
\item
  Increased inbreeding and expression of deleterious recessives
\end{itemize}

\textbf{Cheetah (\emph{Acinonyx jubatus})}: Cheetahs experienced at
least two severe bottlenecks -- one approximately 100,000 years ago and
another during the Late Pleistocene (about 10,000--12,000 years ago).
The genetic consequences are dramatic:

\begin{itemize}
\tightlist
\item
  \hyperref[gl:nucleotide]{\textbf{Nucleotide}} diversity
  \(\pi \approx 0.0001\) -- comparable to inbred laboratory mice
\item
  MHC near-monomorphism: skin grafts between unrelated cheetahs are not
  rejected, indicating almost no immune system genetic variation
\item
  High frequency of sperm abnormalities (about 70\% morphologically
  abnormal)
\item
  Extreme susceptibility to feline coronavirus and other pathogens
\item
  Conservation implication: even wild cheetah populations suffer
  inbreeding depression
\end{itemize}

\textbf{Northern elephant seal (\emph{Mirounga angustirostris})}: Hunted
to approximately 20 individuals by the 1890s, the population has since
recovered to over 200,000. However, genetic diversity at allozyme loci
is essentially zero -- most individuals are genetically nearly identical
at many loci. The population is a genetic time capsule of a tiny
founding group.

\begin{quote}
\textbf{Real-World Connection: Conservation Genetics -- The Cheetah
Crisis}

The cheetah illustrates why genetic diversity matters for species
survival. With minimal MHC variation, a single pathogen could
potentially devastate the entire species -- every individual lacks the
immune diversity that protects genetically variable populations.
Conservation geneticists use molecular markers to assess genetic health
of endangered populations, guide captive breeding programs to maximize
\(N_e\), and identify genetically distinct populations that should be
managed separately. The Florida panther recovery program successfully
increased genetic diversity by introducing Texas pumas -- a
controversial but effective genetic rescue. Similar approaches are being
considered for cheetahs, using the more genetically diverse East African
populations to supplement Southern African populations.
\end{quote}

\subsubsection{Founder Effect}\label{founder-effect}

The \textbf{founder effect} occurs when a small number of individuals
establish a new population, carrying primarily a non-representative
sample of the source population's genetic variation. Unlike bottlenecks,
founder effects involve colonization of new territory, and the resulting
population may remain small for many generations.

{\def\LTcaptype{none} % do not increment counter
\begin{longtable}[]{@{}
  >{\raggedright\arraybackslash}p{(\linewidth - 4\tabcolsep) * \real{0.1081}}
  >{\raggedright\arraybackslash}p{(\linewidth - 4\tabcolsep) * \real{0.3784}}
  >{\raggedright\arraybackslash}p{(\linewidth - 4\tabcolsep) * \real{0.5135}}@{}}
\toprule\noalign{}
\begin{minipage}[b]{\linewidth}\raggedright
Case
\end{minipage} & \begin{minipage}[b]{\linewidth}\raggedright
Founding event
\end{minipage} & \begin{minipage}[b]{\linewidth}\raggedright
Genetic consequence
\end{minipage} \\
\midrule\noalign{}
\endhead
\bottomrule\noalign{}
\endlastfoot
\textbf{Ellis-van Creveld syndrome in Amish} & One couple (Samuel King
and wife) emigrated from Europe to Lancaster County, Pennsylvania, about
1744 & Frequency of EVS allele: 1 in 8
(\hyperref[gl:heterozygous]{\textbf{heterozygous}} carriers) vs.~1 in
60,000 in general population \\
\textbf{Finnish disease heritage} & Multiple founding events about 4,000
years ago; population remained small and isolated & Over 30 genetic
diseases enriched in Finns (e.g., congenital nephrotic syndrome,
aspartylglucosaminuria) \\
\textbf{Tristan da Cunha} & 15 original settlers on remote South
Atlantic island (1816) & High frequency of asthma and retinitis
pigmentosa traced to specific founders \\
\textbf{Pingelap atoll achromatopsia} & Typhoon in about 1775 reduced
population to about 20 survivors; chief carried achromatopsia allele &
about 10\% of population has complete color blindness (vs.~1 in 30,000
globally) \\
\end{longtable}
}

\paragraph{Documented founder cases in
detail}\label{documented-founder-cases-in-detail}

\textbf{Old Order Amish, Lancaster County (Pennsylvania).} The Lancaster
County Amish trace their ancestry to about 200 founding
families (about 700 individuals) who emigrated from the German
Palatinate and Switzerland between 1720 and 1770; the population has
subsequently grown to about 45,000 with extensive endogamy. The
genetic consequences are dramatic. \textbf{Ellis--van Creveld syndrome
(EvC)} --- an autosomal-recessive chondrodysplasia (short-limbed
dwarfism, polydactyly, dental anomalies, cardiac septal defects) --- has
a Lancaster Amish prevalence of approximately \textbf{1 in 200 live
births}, compared with about \textbf{1 in 60,000--70,000} in
the general European-descent population. Coalescent analysis traces
every Amish EvC allele to a single founding couple --- Samuel King and
his wife --- who emigrated in about 1744. This is a textbook
case of how a \textbf{single founder chromosome}, with its specific
\emph{EVC} mutation, rose to high frequency through genetic drift in a
small endogamous population. Other recessive diseases enriched among the
Lancaster Amish include glutaric aciduria type 1, maple syrup urine
disease, and a specific \emph{KCNQ1} long QT-syndrome mutation.

\textbf{Pingelap atoll, Federated States of Micronesia.} Pingelap is a
small (1.8 km\texorpdfstring{\textsuperscript{2}}{2}) coral atoll in the Caroline Islands. In approximately
\textbf{1775}, \textbf{Typhoon Lengkieki} struck the island, killing
most residents and reducing the surviving population to roughly
\textbf{20 individuals} --- a severe single-generation bottleneck. One
survivor, the paramount chief Nahnmwarki Mwanenihsed, was a heterozygous
carrier of an autosomal recessive \emph{CNGB3} mutation causing
\textbf{complete achromatopsia} (total color blindness, photophobia,
severe visual impairment, with cone photoreceptors entirely
non-functional from birth). Genealogical and molecular analyses,
popularised in Oliver Sacks's \emph{The Island of the Colorblind}
(1996), trace nearly every modern Pingelap achromatopsia case to this
single founder chromosome. The condition, which afflicts roughly
\textbf{1 in 30,000--50,000 people globally}, occurs in approximately
\textbf{5--10 \% of the modern Pingelap population} (with carrier
frequency around 30 \%) --- a \textgreater1,000-fold elevation produced
by a single chance founder allele surviving the typhoon.

\textbf{French Canadians of Quebec.} Quebec was colonized between 1608
and 1759 by approximately \textbf{8,500 French settlers} --- a small
founding population that subsequently expanded to about 6
million descendants while remaining geographically and culturally
isolated from broader European-descent populations through mid-twentieth
century. This combination of a small founder population, large
subsequent expansion, and limited gene flow produced one of the world's
best-studied founder populations. Specific recessive disorders enriched
among French Canadians include:

\begin{itemize}
\tightlist
\item
  \textbf{Tay-Sachs disease}: A different \emph{HEXA} mutation than the
  Ashkenazi founder allele --- the French Canadian variant (a 7.6-kb
  deletion) reaches a carrier frequency of about 1/14 in some
  Quebec subpopulations, comparable to the Ashkenazi rate.
\item
  \textbf{Tyrosinemia type I}: Caused by a \emph{FAH}
  (fumarylacetoacetate hydrolase) splice-site mutation, with a carrier
  frequency of about 1/14 in the Saguenay-Lac-Saint-Jean region
  (vs.~about 1/100,000 worldwide). The single Quebec founder
  was traced to a Norman settler from the 1600s.
\item
  \textbf{Autosomal recessive spastic ataxia of Charlevoix-Saguenay
  (ARSACS)}: A neurodegenerative disorder reaching
  about 1/1,500 in the Charlevoix region, caused by a specific
  \emph{SACS} mutation absent elsewhere.
\item
  \textbf{Hereditary fructose intolerance, congenital lactic acidosis,
  oculopharyngeal muscular dystrophy}: Most show elevated frequencies in
  specific French Canadian subpopulations traceable to identifiable
  founder events.
\end{itemize}

These three cases illustrate that the founder effect is \textbf{not
selection} but \textbf{probability}: a recessive allele survives the
bottleneck because the founder happens to carry it, and subsequent
endogamous mating compounds homozygosity over generations. Modern
genetic medicine in such populations leverages founder-effect
concentration: a single targeted assay at one chromosomal site can
screen for the population-specific high-prevalence allele, which would
be cost-prohibitive in genetically heterogeneous outbred populations.

\begin{quote}
\textbf{Concept Check 1:} A population of 100 individuals experiences a
bottleneck to 5 individuals for one generation, then immediately
recovers to 100. Using the harmonic mean formula, calculate \(N_e\) over
these two generations. How does this compare to the arithmetic mean?
\end{quote}

\begin{center}\rule{0.5\linewidth}{0.5pt}\end{center}

\subsection{Gene Flow and Mutation}\label{gene-flow-and-mutation}

\subsubsection{Gene Flow}\label{gene-flow}

\textbf{Gene flow} (migration) is the movement of alleles between
populations through the dispersal and reproduction of individuals (or
their \hyperref[gl:gamete]{\textbf{gamete}}s, such as pollen). Gene flow
has two primary evolutionary effects:

\begin{enumerate}
\def\labelenumi{\arabic{enumi}.}
\tightlist
\item
  \textbf{Homogenization}: Gene flow tends to equalize allele
  frequencies between populations, counteracting divergence due to drift
  or local selection.
\item
  \textbf{Introduction of variation}: Migrants introduce novel alleles
  that may not have arisen locally, increasing genetic diversity within
  recipient populations.
\end{enumerate}

\subsubsection{\texorpdfstring{Measuring Population Differentiation:
\(F_{ST}\)}{Measuring Population Differentiation: F_(ST)}}\label{measuring-population-differentiation-f_st-1}

Wright's \textbf{fixation index} (\(F_{ST}\)) quantifies the proportion
of total genetic variation attributable to differences between
populations:

\begin{equation}
F_{ST} = \frac{H_T - H_S}{H_T} = \frac{\sigma^2_p}{\bar{p}(1-\bar{p})}
\label{eq:genetic_drift_and_speciation_7}
\end{equation}

where \(H_T\) is the expected heterozygosity of the total (pooled)
population and \(H_S\) is the mean expected heterozygosity within
subpopulations.

\begin{itemize}
\tightlist
\item
  \(F_{ST} = 0\): No differentiation (panmixia -- gene flow fully
  homogenizes populations)
\item
  \(F_{ST} = 1\): Complete differentiation (populations fixed for
  different alleles)
\end{itemize}

\textbf{Wright's island model}: At migration-drift equilibrium:

\begin{equation}
F_{ST} \approx \frac{1}{1 + 4N_e m}
\label{eq:unit_VI_fst_nm}
\end{equation}

where \(m\) is the proportion of each subpopulation composed of migrants
per generation. This compact relationship --- equating differentiation
to the inverse of the \textbf{scaled migration rate} \(4N_em\) --- is
one of the most-used results in molecular ecology. It reveals a striking
insight: even very low migration rates (\(m\)) can prevent substantial
differentiation. \(N_em = 1\) (just one migrant per generation) yields
\(F_{ST} = 0.20\); \(N_em = 4\) yields \(F_{ST} = 0.06\). \textbf{One
migrant per generation is the rule of thumb for keeping populations from
drifting apart}, regardless of population size --- because the
homogenizing effect of migration scales with \(m \cdot N_e\) while the
diverging effect of drift scales with \(1/N_e\), leaving the
dimensionless product \(N_em\) as the controlling parameter.

The relationship can be inverted for empirical use: given a measured
\(F_{ST}\), the implied number of migrants per generation is

\[N_em \approx \frac{1 - F_{ST}}{4 F_{ST}}. \label{eq:unit_VI_genetic_drift_and_speciation_item_4}\]

For \(F_{ST} = 0.05\) (moderate human population structure),
\(N_em \approx 4.75\) --- about five effective migrants between
subpopulations per generation, sufficient to keep differentiation
modest.

\subsubsection{Isolation by Distance}\label{isolation-by-distance}

\textbf{Isolation by distance (IBD)}: In continuously distributed
populations, gene flow decreases with geographic distance. This produces
a positive correlation between pairwise \(F_{ST}\) (or genetic distance)
and geographic distance, detectable by a \textbf{Mantel test} (matrix
correlation between genetic and geographic distance matrices).

IBD patterns confirm that most natural populations show continuous
genetic gradients (clines) rather than sharp boundaries. Human genetic
variation, for example, shows strong IBD patterns corresponding to the
out-of-Africa expansion routes.

\subsubsection{Secondary Contact and Hybrid
Zones}\label{secondary-contact-and-hybrid-zones}

When previously allopatric populations come back into contact, the
outcome depends on the degree of reproductive isolation that evolved
during separation:

\begin{itemize}
\tightlist
\item
  \textbf{Fusion}: If reproductive isolation is weak, populations merge
  and homogenize.
\item
  \textbf{Reinforcement}: If hybrids have reduced fitness, selection
  strengthens prezygotic barriers, as developed in the
  reproductive-isolation discussion.
\item
  \textbf{Stable hybrid zone}: A narrow geographic zone of hybridization
  maintained by a balance between dispersal into the zone and selection
  against hybrids --- called a \textbf{tension zone}.
\end{itemize}

\paragraph{Tension zone width and the dispersal-selection
balance}\label{tension-zone-width-and-the-dispersal-selection-balance}

Following Barton \& Hewitt (1985), the equilibrium \textbf{width} \(w\)
of a tension zone follows:

\begin{equation}
w \approx \frac{\sigma}{\sqrt{s}}
\label{eq:unit_VI_tension_zone}
\end{equation}

where \texorpdfstring{\ensuremath{\sigma}}{sigma} is the per-generation dispersal distance and \(s\) is the
selection coefficient against hybrids. Tension zones are
\textbf{independent of geography} in a key sense: their width is a
function primarily of dispersal and selection, not of the underlying
environmental gradient. The same hybrid zone can move across the
landscape (as one parental population expands at the expense of the
other) while maintaining its width. Tension zones become trapped at
landscape features (mountain passes, river crossings) where dispersal
rates dip locally --- explaining why hybrid zones often coincide with
geographic obstacles even when the species' ranges are broadly
continuous.

\textbf{Empirical examples}:

\begin{itemize}
\tightlist
\item
  \textbf{European crow hybrid zone} (carrion crow \emph{Corvus corone}
  \texorpdfstring{\ensuremath{\times}}{x} hooded crow \emph{C. cornix}): The hybrid zone runs through Europe
  from Italy and Denmark up through Russia, with a width of
  about 20--30 km. \emph{Corone} and \emph{cornix} differ at a
  small number of pigmentation loci. Genome-wide \(F_{ST}\) across the
  zone is dramatic at color-loci but minimal elsewhere --- a ``porous''
  species barrier consistent with reinforcement primarily at loci
  affected by mate choice.
\item
  \textbf{Northern flicker} (\emph{Colaptes auratus}) hybrid zone:
  yellow-shafted (eastern) \texorpdfstring{\ensuremath{\times}}{x} red-shafted (western) flickers form a
  narrow hybrid zone running about 3,000 km along the central
  Great Plains, with width about 100--200 km. Despite extensive
  interbreeding in the zone, the parental forms persist on either side
  --- assortative mating and reduced hybrid fitness preserve the
  polymorphism.
\item
  \textbf{\emph{Bombina bombina} \texorpdfstring{\ensuremath{\times}}{x} \emph{B. variegata}} toads: A 5--8 km
  tension zone runs through Central Europe. Genome-wide \(F_{ST}\)
  patterns reveal both barrier loci (under strong selection against
  hybrids) and freely-introgressing neutral loci. The genome behaves
  like a mosaic of differently permeable regions.
\end{itemize}

\paragraph{Reinforcement and Dobzhansky-Muller
incompatibilities}\label{reinforcement-and-dobzhansky-muller-incompatibilities}

When hybrids have reduced fitness, \textbf{reinforcement} strengthens
prezygotic barriers --- natural selection favors mate-discrimination
alleles that prevent hybridization. The genetic basis of postzygotic
isolation typically involves \textbf{Dobzhansky-Muller
incompatibilities} (BDMI): pairs of alleles at different loci that
function fine in their respective parental backgrounds but produce
sterile or inviable hybrids when combined.

The classic two-locus BDMI: ancestral genotype \(A_1B_1\). Population 1
evolves to \(A_2B_1\) (the new \(A_2\) is tested against \(B_1\) ---
works fine). Population 2 evolves to \(A_1B_2\) (the new \(B_2\) is
tested against \(A_1\) --- works fine). The hybrid genotype \(A_2B_2\)
has rarely been tested by selection; if \(A_2\) and \(B_2\) interact
deleteriously, hybrids fail. Because each parental population
independently accumulates substitutions, the number of potential
incompatibilities grows as \(k^2\) (Orr's snowball) --- explaining why
isolation often appears suddenly after a long period of gradual
divergence.

Sympatric divergence (without geographic separation) is rare because
gene flow continuously breaks up co-adapted gene complexes.
Reinforcement can complete divergence after an initial allopatric phase:
secondary contact selects for prezygotic isolation alleles that prevent
the formation of low-fitness hybrid combinations. \textbf{The genomic
signature of reinforcement} --- sharper differentiation at mate-choice
and behavioral-isolation loci than at genome-wide background --- has
been documented in \emph{Drosophila} species pairs and in cichlid
radiations.

\subsubsection{Coalescent Theory and the
TMRCA}\label{coalescent-theory-and-the-tmrca}

The \textbf{coalescent} is a backward-in-time view of population
genealogies that has become the dominant statistical framework in modern
population genetics \citep{kingman1982}. Rather than tracking allele
frequencies forward through generations, the coalescent traces ancestral
lineages of currently sampled alleles backward until they merge
(``coalesce'') at common ancestors.

\paragraph{Expected coalescence time}\label{expected-coalescence-time}

For a sample of \(k\) alleles drawn from a Wright-Fisher population of
effective size \(N_e\) diploid individuals (so \(2N_e\) allele copies),
the expected time for any two specific lineages to coalesce is \(2N_e\)
generations. Generalizing to a sample of \(k\) lineages, the expected
total time to the most recent common ancestor (TMRCA) is:

\begin{equation}
E[T_{\text{MRCA}}] = 4N_e \left(1 - \frac{1}{k}\right)
\label{eq:unit_VI_tmrca}
\end{equation}

For large samples (\(k \to \infty\)), \(E[T_{\text{MRCA}}] \to 4N_e\)
generations. The coalescent has several remarkable properties:

\begin{itemize}
\tightlist
\item
  \textbf{Most coalescences happen recently}: the expected time for the
  \textbf{last} two lineages to coalesce (\(2N_e\) generations) is half
  of the total tree depth. The genealogy is ``top-heavy'': coalescent
  events are most rapid at the tips of the tree, slowest near the root.
\item
  \textbf{Independent of forward dynamics}: Under neutrality, the
  topology of the coalescent tree is independent of mutations. Mutations
  are sprinkled onto branches at rate \texorpdfstring{\ensuremath{\mu}}{mu}; the expected number of
  segregating sites in a sample of \(k\) is
  \(\theta \sum_{i=1}^{k-1} 1/i\), where \(\theta = 4N_e\mu\).
\item
  \textbf{Variance}: TMRCA has a wide distribution. The variance of
  TMRCA grows with \(k\) but the \textbf{most likely TMRCA} values are
  surprisingly broad --- different loci sampled from the same population
  yield TMRCAs spanning an order of magnitude.
\end{itemize}

\paragraph{Why ``Mitochondrial Eve'' was not the only woman
alive}\label{why-mitochondrial-eve-was-not-the-only-woman-alive}

Coalescent logic explains why the TMRCA of human mitochondrial DNA
(about 150,000 years ago) does \textbf{not} mean a single
ancestress. Many women were alive 150 kya --- the coalescent estimate
refers primarily to the lineage of mtDNA that survived to the present.
Most other maternal lineages were lost by drift along the way; their
loss is statistically inevitable in finite populations. Indeed, of \(k\)
ancestral lineages alive in the founding population, most but one will
eventually be lost to drift --- with the surviving lineage tracing back
through ``Eve'' by definition. The surprising thing is not that an MRCA
exists, but that the coalescent pinpoints when she lived.

\paragraph{Coalescent applications}\label{coalescent-applications}

\begin{itemize}
\tightlist
\item
  \textbf{Effective population size estimation} from segregating-site
  counts uses the relationship \(\theta = 4N_e\mu\).
\item
  \textbf{Detecting selective sweeps} identifies regions with
  anomalously short TMRCAs --- a recent advantageous mutation drags
  linked variation toward fixation, collapsing the genealogy.
\item
  \textbf{Demographic inference} from extended coalescent models
  reconstructs population growth, bottlenecks, and migration history.
  PSMC (pairwise sequentially Markovian coalescent) and SMC++ are
  standard tools.
\end{itemize}

\paragraph{TMRCA for two alleles: derivation and
intuition}\label{tmrca-for-two-alleles-derivation-and-intuition}

The single most-cited coalescent result is the expected coalescence time
of two randomly sampled alleles in a Wright--Fisher diploid population
of effective size \(N_e\). In each generation, two specific allele
copies have probability \(1/(2N_e)\) of sharing the same parent (because
a parent is chosen uniformly from \(2N_e\) possibilities). The
coalescence time is therefore geometrically distributed with success
probability \(1/(2N_e)\), giving:

\begin{equation}
E[T_{\text{MRCA, 2-alleles}}] = 2N_e \text{ generations}
\label{eq:unit_VI_tmrca_pairwise}
\end{equation}

For \(N_e = 10{,}000\) (a typical estimate for ancestral humans), the
expected pairwise coalescence is \textbf{20,000 generations} \texorpdfstring{\ensuremath{\approx}}{~} 500,000
years (assuming 25-year generation time). This figure is the
\textbf{statistical baseline} against which observed pairwise sequence
divergences are interpreted; deeper coalescences indicate population
structure or balancing selection, shallower coalescences suggest
bottlenecks or selective sweeps.

For a sample of \(k\) alleles, total tree depth (root to most recent
tip) is \(4N_e(1 - 1/k)\) --- the value introduced in
\cref{eq:unit_VI_tmrca}. For two alleles (\(k = 2\)), this reduces to
\(2N_e\), recovering the pairwise result of
\cref{eq:unit_VI_tmrca_pairwise}.

\paragraph{Structured coalescent}\label{structured-coalescent}

When a population is subdivided into subpopulations with limited
migration, the coalescent becomes \textbf{structured}: lineages can
primarily coalesce when they are in the same subpopulation, and
migration events transfer lineages between subpopulations. Two lineages
drawn from the same subpopulation typically coalesce on the timescale
\(2N_e\) (within-deme coalescence). Two lineages drawn from
\textbf{different subpopulations} must first meet in the same deme
before they can coalesce --- adding a waiting time inversely
proportional to the migration rate \(m\). The result is that
\textbf{between-population pairwise TMRCA exceeds within-population
pairwise TMRCA}, with the difference scaling as \(\sim 1/(2m)\)
generations.

This produces an empirical signature: \(F_{ST}\) between subpopulations
grows as the fraction of the genealogy spent in the ``between'' rather
than ``within'' coalescent regime. Conversely, \textbf{migration reduces
between-population TMRCA} because lineages can quickly transfer between
subpopulations; high-migration models behave like a single panmictic
population with effective size equal to the sum of deme sizes. Modern
coalescent inference tools (msprime, IM, PhyloGeoTools) explicitly model
structured-coalescent dynamics to estimate ancient migration rates and
population-divergence times --- a key application in human population
history (e.g., African \texorpdfstring{\ensuremath{\leftrightarrow}}{<->} Eurasian migration after Out-of-Africa) and in
conservation genetics of fragmented populations.

\subsubsection{Mutation as the Ultimate Source of
Variation}\label{mutation-as-the-ultimate-source-of-variation}

\textbf{Mutation} is the ultimate source of most genetic variation.
Without mutation, evolution would eventually cease as selection and
drift exhausted existing variation.

\begin{itemize}
\tightlist
\item
  \textbf{De novo mutation rate in humans}: approximately
  \(1.2 \times 10^{-8}\) per nucleotide per generation, corresponding to
  about 70 new mutations per individual per generation.
\item
  \textbf{Paternal age effect}: Male germline cells undergo many more
  mitotic divisions than female germline cells. Each year of paternal
  age adds approximately 1--2 additional mutations to offspring. This
  contributes to the increased incidence of some genetic conditions
  (e.g., achondroplasia, autism spectrum disorder) with advanced
  paternal age.
\end{itemize}

\subsubsection{Mutation-Selection
Balance}\label{mutation-selection-balance-1}

For a deleterious recessive allele maintained by recurrent mutation
against purifying selection, the equilibrium frequency is:

\begin{equation}
\hat{q} \approx \sqrt{\frac{\mu}{s}}
\label{eq:genetic_drift_and_speciation_9}
\end{equation}

where \texorpdfstring{\ensuremath{\mu}}{mu} is the mutation rate and \(s\) is the selection coefficient
against the homozygous genotype. This explains why many genetic diseases
persist at low frequencies -- new mutations continually reintroduce them
even as selection removes them.

\textbf{Example}: Cystic fibrosis (CF) is caused by homozygosity for
loss-of-function alleles in the \emph{CFTR} gene. With
\(\mu \approx 6.7 \times 10^{-7}\) and \(s \approx 1\) (CF was
historically lethal before modern medicine), mutation-selection balance
predicts \(\hat{q} \approx \sqrt{6.7 \times 10^{-7}} \approx 0.0008\).
However, the actual carrier frequency in European populations is
approximately 1 in 25 (\(q \approx 0.02\)), far exceeding the
mutation-selection prediction. This suggests an additional force --
possibly heterozygote advantage against cholera, typhoid fever, or
tuberculosis -- maintains the CF allele at elevated frequency, similar
to the sickle cell / malaria relationship.

\begin{quote}
\textbf{Real-World Connection: Mutation Rates and Cancer}

The same mutational processes that drive evolution also drive cancer.
Somatic mutations accumulate in cells throughout an organism's lifetime
at rates of approximately \(0.5\)--\(5 \times 10^{-9}\) per nucleotide
per cell division. Larger and longer-lived organisms should, by simple
probability, experience more cancer -- yet elephants (with 100 times
more cells than humans) have lower cancer rates than humans. This is
\textbf{Peto's paradox}. The resolution involves the evolution of
enhanced cancer suppression mechanisms: elephants carry approximately 20
copies of the \emph{TP53} tumor suppressor gene (humans have one), and
their cells are hypersensitive to DNA damage, triggering
\hyperref[gl:apoptosis]{\textbf{apoptosis}} rather than attempting
repair. This represents an evolutionary adaptation to large body size.
\end{quote}

\begin{quote}
\textbf{Concept Check 2:} Using Wright's island model, calculate the
equilibrium \(F_{ST}\) for two populations with \(N_e = 500\) and
\(m = 0.005\). How many migrants per generation does this represent?
Would you consider these populations significantly differentiated?
\end{quote}

\begin{center}\rule{0.5\linewidth}{0.5pt}\end{center}

\subsection{Species Concepts}\label{species-concepts}

Defining what constitutes a ``species'' is one of the most contentious
issues in biology. Multiple species concepts exist, each emphasizing
different aspects of biological reality.

\subsubsection{Biological Species Concept
(BSC)}\label{biological-species-concept-bsc}

Ernst \citet{mayr1942} defined species as ``groups of actually or
potentially interbreeding natural populations that are reproductively
isolated from other such groups.'' This concept emphasizes
\textbf{reproductive isolation} as the defining criterion.

\textbf{Strengths}: Biologically meaningful; focuses on the process
(gene flow) that maintains species cohesion.

\textbf{Limitations}: Cannot be applied to asexual organisms, fossils
(where reproductive behavior is unobservable), or ring species;
hybridization between ``good'' species is common (e.g., about 10\% of
bird species hybridize).

\subsubsection{Morphological/Typological Species
Concept}\label{morphologicaltypological-species-concept}

Species are distinguished by consistent morphological differences. This
is the oldest concept, rooted in Linnaeus's original classifications.

\textbf{Strengths}: Practical and widely applicable; the basis of most
field identification.

\textbf{Limitations}: Subjective -- how much difference constitutes a
species? Fails for cryptic species (morphologically identical but
reproductively isolated) and polymorphic species (morphologically
variable within a single species).

\subsubsection{Phylogenetic Species Concept
(PSC)}\label{phylogenetic-species-concept-psc}

A species is the smallest diagnosable monophyletic group -- the smallest
cluster of organisms that share a common ancestor and can be
distinguished from other such clusters by unique character combinations
(synapomorphies).

\textbf{Strengths}: Applicable to any organism (sexual, asexual,
fossil); consistent with modern systematics.

\textbf{Limitations}: Tends to ``oversplit'' -- recognizing more species
than the BSC from the same organisms. With sufficient genetic data,
nearly every population is diagnosable.

\subsubsection{Ecological Species
Concept}\label{ecological-species-concept}

A species is a set of organisms exploiting a single ecological
\hyperref[gl:niche]{\textbf{niche}} (Hutchinson's hypervolume). This
concept focuses on the ecological role rather than reproductive
boundaries.

\textbf{Strengths}: Applicable to asexual organisms; emphasizes the role
of selection in maintaining species identity.

\textbf{Limitations}: Niche boundaries are often fuzzy; different
populations of the same species may occupy different niches.

\subsubsection{Cohesion Species Concept}\label{cohesion-species-concept}

Alan \citet{templeton1989} proposed the most comprehensive definition: a
species is the most inclusive group of organisms having the potential
for \textbf{genetic exchangeability} (gene flow) and/or
\textbf{demographic exchangeability} (ecological interchangeability).
This concept subsumes elements of the BSC and ecological concepts.

\subsubsection{Ring Species}\label{ring-species}

\textbf{Ring species} demonstrate speciation as a gradual process. A
series of populations forms a ring around a geographic barrier, with
neighboring populations interbreeding freely. However, at the point
where the ring closes, the terminal populations are reproductively
isolated despite being connected by a chain of interbreeding
intermediates.

\begin{itemize}
\tightlist
\item
  \textbf{Larus gulls}: Around the Arctic, a chain of herring gull and
  lesser black-backed gull populations form a ring. Adjacent populations
  interbreed, but the end members (in Europe) coexist as distinct,
  reproductively isolated species.
\item
  \textbf{Ensatina salamanders}: In California, populations encircle the
  Central Valley. Adjacent populations hybridize, but where the ring
  closes in southern California, the terminal forms (\emph{E. e.
  eschscholtzii} and \emph{E. e. klauberi}) rarely interbreed.
\end{itemize}

Ring species challenge the BSC because reproductive isolation is not
absolute but graded.

\begin{quote}
\textbf{Concept Check 3:} A researcher discovers two populations of
beetles that are morphologically identical but reproductively isolated
in the laboratory. Would the BSC and morphological species concept agree
on whether these are one species or two? Which concept would you favor,
and why?
\end{quote}

\begin{center}\rule{0.5\linewidth}{0.5pt}\end{center}

\subsection{Speciation Mechanisms}\label{speciation-mechanisms}

\begin{figure}[htbp]
\centering
\includegraphics[width=0.82\linewidth,height=4.2in,keepaspectratio]{../figures/mermaid_inline/inline_mermaid_0108_80611ce8a896.png}
\caption{Flowchart for Speciation Mechanisms: Speciation, Allopatric, Peripatric, and Parapatric form the diagram's primary path or branches.}
\end{figure}

\subsubsection{Allopatric Speciation}\label{allopatric-speciation}

\textbf{Allopatric speciation} -- speciation resulting from geographic
isolation -- is the most common and best-documented mode. A geographic
barrier (mountain range, ocean, ice sheet, river) divides a population,
preventing gene flow. The isolated populations then diverge through a
combination of natural selection (adapting to different local
environments) and genetic drift, eventually accumulating enough
differences that reproductive isolation is a byproduct of divergence.

\textbf{Isthmus of Panama}: The formation of the Central American land
bridge approximately 3 Mya divided many marine populations into
Caribbean and Pacific components. The resulting \textbf{geminate species
pairs} (closely related species on either side of the isthmus) provide
natural experiments in allopatric divergence. Snapping shrimp
(\emph{Alpheus} spp.) separated by the isthmus show increasing
reproductive isolation correlated with divergence time.

\textbf{Dobzhansky-Muller incompatibilities}: The genetic basis of
reproductive isolation often involves epistatic interactions between
independently evolved alleles. Consider two populations with ancestral
genotype \(A_1B_1\). Population 1 evolves to \(A_2B_1\) and Population 2
evolves to \(A_1B_2\). Each derived genotype is functional because the
new allele was tested against the existing genetic background. But
hybrids (\(A_2B_2\)) carry a combination that has rarely been tested by
natural selection -- and this untested epistatic combination may cause
developmental failure, sterility, or inviability. The number of
potential incompatibilities grows as the square of the number of
substitutions (\textbf{Orr's snowball effect}):

\begin{equation}
\text{Incompatibilities} \propto k^2
\label{eq:genetic_drift_and_speciation_10}
\end{equation}

This accelerating accumulation explains why reproductive isolation often
appears suddenly after a period of gradual divergence.

\subsubsection{Peripatric Speciation}\label{peripatric-speciation}

\textbf{Peripatric speciation} occurs when a small population at the
periphery of a species' range becomes isolated. The small founding
population experiences strong genetic drift (founder effect), which,
combined with potentially different selective pressures at the range
margin, can drive rapid divergence.

\begin{itemize}
\tightlist
\item
  \textbf{Hawaiian honeycreepers}: Approximately 55 species evolved from
  one or two founding finch-like ancestors that colonized the Hawaiian
  Islands approximately 5 Mya. The radiation produced an extraordinary
  diversity of bill shapes -- from long, curved bills for nectar feeding
  (\emph{Drepanis}) to massive, parrot-like bills for seed cracking
  (\emph{Pseudonestor}) to thin, warbler-like bills for insect gleaning
  (\emph{Oreomystis}). Many species are now critically endangered or
  extinct.
\item
  \textbf{Darwin's finches}: Fifteen species on the Galapagos Islands,
  diversified from a single South American ancestor approximately 1--2
  Mya. Beak depth and width vary with seed size and hardness.
  \textbf{Character displacement} -- divergence in beak morphology is
  greater where two species co-occur (sympatry) than where each occurs
  alone (allopatry) -- provides evidence that competition drives
  morphological divergence.
\end{itemize}

\subsubsection{Parapatric Speciation}\label{parapatric-speciation}

\textbf{Parapatric speciation} occurs along an environmental gradient
without a clear geographic barrier. Adjacent populations experience
different selective pressures, and if selection against intermediates is
strong enough to overcome gene flow, divergence can proceed.

\textbf{Sweet vernal grass (\emph{Anthoxanthum odoratum})}: On
heavy-metal-contaminated mine tailings in Wales, grass populations have
evolved heavy metal tolerance within meters of the mine boundary.
Tolerant and non-tolerant populations grow adjacent to each other.
Flowering time has shifted between the two ecotypes, providing a
prezygotic barrier -- tolerant plants flower approximately one week
earlier than adjacent non-tolerant plants. Wind-mediated gene flow
between the two populations is substantial (pollen travels freely), yet
the ecotypes maintain their distinctness because of strong selection
against non-tolerant genotypes on contaminated soil and against tolerant
genotypes (which pay a fitness cost) on clean soil. This is parapatric
speciation in progress -- divergence without geographic isolation,
driven by strong disruptive selection.

\textbf{Hybrid zones as natural laboratories}: The parapatric boundary
between ecotypes often produces a narrow \textbf{cline} -- a geographic
gradient in allele frequency. The width of the cline reflects the
balance between selection (which narrows it) and gene flow (which
broadens it). Cline width \(w\) can be estimated as:

\begin{equation}
w \approx \frac{\sigma}{\sqrt{s}}
\label{eq:genetic_drift_and_speciation_11}
\end{equation}

where \texorpdfstring{\ensuremath{\sigma}}{sigma} is the dispersal distance per generation and \(s\) is the
selection coefficient against immigrants. Narrow clines indicate strong
selection relative to dispersal.

\subsubsection{Sympatric Speciation}\label{sympatric-speciation}

\textbf{Sympatric speciation} -- speciation within a single, freely
mixing population -- has long been controversial because gene flow
should homogenize populations. However, several mechanisms can overcome
gene flow:

\textbf{Cichlid fish in Lake Victoria}: Approximately 500 species
evolved within a single lake basin in approximately 15,000 years -- one
of the most spectacular adaptive radiations known. Sexual selection on
male coloration, driven by female preference under different light
environments (clear versus turbid water), appears to be a primary
driver. When experimental turbidity eliminates the ability to
distinguish colors, species boundaries break down and hybridization
occurs -- demonstrating that assortative mating by color maintains
species isolation.

\textbf{Apple maggot fly (\emph{Rhagoletis pomonella})}: The ancestral
host plant is hawthorn. Following the introduction of domestic apples to
North America in the 1600s, a population shifted to apple as a host.
Apple and hawthorn races now differ in emergence timing (matching fruit
ripening), host preference, and multiple genomic regions under divergent
selection. This represents host race formation -- an early stage of
sympatric speciation.

\textbf{Polyploidy in plants}:
Whole-\hyperref[gl:genome]{\textbf{genome}} duplication creates
\textbf{instantaneous reproductive isolation} because polyploid
offspring cannot successfully mate with diploid parents (triploid
offspring are sterile or inviable due to meiotic irregularities --- odd
chromosome number prevents proper homolog pairing at meiosis I,
producing unbalanced gametes that fail to develop). The triploid block
is the mechanistic basis for why polyploidy can speciate in a
\textbf{single generation}: a tetraploid individual that arises by
chromosome non-disjunction is reproductively isolated from its diploid
parents from the moment of formation.

\begin{itemize}
\tightlist
\item
  \textbf{Autopolyploidy}: Genome doubling within a single species
  without hybridization. Potato (\emph{Solanum tuberosum}) is
  autotetraploid (4n = 48). Alfalfa (\emph{Medicago sativa}) is
  autotetraploid. Some wild \emph{Tolmiea} (piggyback plant) populations
  exhibit autotetraploid forms sympatric with diploid forms.
  Autopolyploids face less severe meiotic challenges than allopolyploids
  because their four homologs share recent ancestry, but
  \textbf{multivalent formation at meiosis I} (where four near-identical
  homologs may pair as quadrivalents rather than two bivalents) can
  produce aneuploid gametes at higher rates than in diploids.
  Reproductive isolation can arise rapidly because an autotetraploid
  (4n) crossed with a diploid (2n) parent often produces low-fertility
  triploid (3n) offspring.
\item
  \textbf{Allopolyploidy}: Hybridization between two distinct species
  followed by \hyperref[gl:chromosome]{\textbf{chromosome}} doubling.
  The chromosome doubling restores fertility because each chromosome now
  has a homolog to pair with at meiosis I (the parental homolog),
  eliminating the meiotic irregularities that would otherwise sterilize
  the F1 hybrid. Major allopolyploid crops illustrate the pattern across
  the angiosperms:

  \begin{itemize}
  \tightlist
  \item
    \textbf{Bread wheat (\emph{Triticum aestivum})}: an allohexaploid
    (2n = 42 = AABBDD) derived from three diploid ancestors through two
    rounds of hybridization and genome doubling, completed within the
    last about 8,000--10,000 years coinciding with the origin
    of agriculture.
  \item
    \textbf{Cultivated tobacco (\emph{Nicotiana tabacum})}: an
    allotetraploid (2n = 48 = SSTT), formed about 200,000 years
    ago by hybridization between \emph{N. sylvestris} (S genome) and an
    \emph{N. tomentosiformis}-like progenitor (T genome).
  \item
    \textbf{Cultivated cotton (\emph{Gossypium hirsutum}, upland
    cotton)}: an allotetraploid (2n = 52 = AADD) formed
    about 1--2 Mya by hybridization between an Old World
    A-genome species and a New World D-genome species --- followed by
    trans-oceanic dispersal of A-genome ancestors to the Americas.
    Modern upland cotton accounts for about 95\% of global
    cotton production.
  \end{itemize}
\end{itemize}

\paragraph{How widespread is
polyploidy?}\label{how-widespread-is-polyploidy}

Polyploidy is a major mode of sympatric speciation in plants and is
\textbf{far more common than once thought}. Recent genomic analyses
indicate that \textbf{about 35\% of extant angiosperm species}
are polyploid or have polyploid ancestors detectable in the recent past,
and many flowering plant lineages show evidence of paleopolyploid events
(whole-genome duplications followed by gene loss and rediploidisation)
during their evolutionary history. The seed plants share at least one
whole-genome duplication about 320 Mya, and the flowering
plants share a more recent duplication about 150 Mya. Major
crop genomes (rice, maize, soybean, \emph{Brassica}, banana, sugarcane)
bear paleopolyploid signatures. \textbf{Polyploidy is not a peculiarity
of a few weedy lineages} --- it is a recurring engine of plant evolution
that contributes substantially to angiosperm diversification, novel gene
functions arising from duplicated copies (sub- and
neofunctionalisation), and rapid speciation. In animals, polyploidy is
rarer because dosage-sensitive development and sex-determination systems
often make genome doubling costly, but it has occurred in salmonid fish,
\emph{Xenopus} frogs, and some invertebrates.

\paragraph{Wheat: a paradigm of agricultural
allopolyploidy}\label{wheat-a-paradigm-of-agricultural-allopolyploidy}

The polyploid origin of bread wheat is one of the best-documented
allopolyploid speciation events in agriculture. Approximately 500,000
years ago, \emph{Triticum urartu} (genome AA) hybridized with
\emph{Aegilops speltoides} (genome BB), producing the tetraploid emmer
wheat (\emph{Triticum turgidum}, AABB). Some 8,000 years ago, emmer
crossed with \emph{Aegilops tauschii} (genome DD), producing modern
hexaploid bread wheat (AABBDD). Each genome retains about 85\%
of its ancestral gene content; the \textbf{homoeologous} chromosome sets
pair specifically with their partners thanks to the \emph{Ph1} locus on
chromosome 5B, which suppresses pairing between non-identical
chromosomes and prevents chaos at meiosis. Modern wheat breeding takes
advantage of the genetic diversity in most three subgenomes ---
beneficial alleles can be sourced from any of the diploid ancestor
lineages and introduced into bread wheat by chromosomal manipulations.

\paragraph{The Brassica triangle (Triangle of
U)}\label{the-brassica-triangle-triangle-of-u}

The cultivated \emph{Brassica} species form one of the most elegant
illustrations of allopolyploid speciation, formalized by Korean botanist
Woo Jang-choon (U Nagaharu) in 1935:

\begin{figure}[htbp]
\centering
\includegraphics[width=0.82\linewidth,height=4.2in,keepaspectratio]{../figures/mermaid_inline/inline_mermaid_0109_5d5b4983bdd4.png}
\caption{The Triangle of U: three diploid Brassica species (corners) and the three allotetraploid species (sides) formed by their pairwise hybridizations and chromosome doublings. Each allopolyploid carries the chromosome complements of its two diploid ancestors.}
\end{figure}

Each diploid ancestor contributes its full chromosome set to the
corresponding allotetraploid; the resulting hexaploid is reproductively
isolated from most parents. The Triangle of U accounts for cabbage
(\emph{B. oleracea}, CC), turnip (\emph{B. rapa}, AA), and mustard
(\emph{B. nigra}, BB) plus their three allotetraploid descendants ---
canola (\emph{B. napus}, AACC), brown mustard (\emph{B. juncea}, AABB),
and Ethiopian mustard (\emph{B. carinata}, BBCC). The triangle forms one
of the most important agricultural-genetic systems in the world;
understanding the polyploid relationships drives breeding programs and
resistance-gene transfer between species.

An estimated 15\% of \hyperref[gl:angiosperm]{\textbf{angiosperm}}
speciation events involve polyploidy, and up to 70\% of grass genomes
show evidence of ancient polyploid ancestry --- including the deep
polyploidy events that gave rise to the cereals (rice, maize, wheat,
barley) that feed the human species.

\begin{center}\rule{0.5\linewidth}{0.5pt}\end{center}

\subsection{Reproductive Isolation}\label{reproductive-isolation}

Reproductive isolation is the hallmark of speciation under the BSC.
Barriers to reproduction are classified as \textbf{prezygotic}
(preventing formation of a hybrid \hyperref[gl:zygote]{\textbf{zygote}})
or \textbf{postzygotic} (reducing fitness of hybrid offspring).

\subsubsection{Haldane's Rule}\label{haldanes-rule}

\textbf{Haldane's rule} (1922): In interspecific hybrids, when a single
sex is inviable or sterile, it is the \textbf{heterogametic sex} (XY
males in mammals and flies; ZW females in birds and butterflies). Three
complementary mechanisms explain this pattern:

\begin{enumerate}
\def\labelenumi{\arabic{enumi}.}
\tightlist
\item
  \textbf{Dominance theory}: Recessive incompatibility alleles on the X
  (or Z) chromosome are fully exposed in the heterogametic sex (which
  has a single copy), but masked in the homogametic sex (which may carry
  a compatible allele on the other X/Z chromosome).
\item
  \textbf{Faster-X evolution}: X-linked genes evolve faster than
  autosomal genes because beneficial recessive mutations are immediately
  exposed to selection in hemizygous males. This accelerated divergence
  produces more X-linked incompatibilities.
\item
  \textbf{Meiotic drive and sex chromosome conflict}: Segregation
  distortion elements on sex chromosomes can cause hybrid sterility by
  disrupting normal chromosome segregation during
  \hyperref[gl:meiosis]{\textbf{meiosis}}.
\end{enumerate}

Haldane's rule has been confirmed across a broad range of animal taxa
examined, making it one of the most robust generalizations in speciation
biology.

\subsubsection{Prezygotic Barriers}\label{prezygotic-barriers}

{\def\LTcaptype{none} % do not increment counter
\begin{longtable}[]{@{}
  >{\raggedright\arraybackslash}p{(\linewidth - 4\tabcolsep) * \real{0.4286}}
  >{\raggedright\arraybackslash}p{(\linewidth - 4\tabcolsep) * \real{0.3214}}
  >{\raggedright\arraybackslash}p{(\linewidth - 4\tabcolsep) * \real{0.2500}}@{}}
\toprule\noalign{}
\begin{minipage}[b]{\linewidth}\raggedright
Barrier type
\end{minipage} & \begin{minipage}[b]{\linewidth}\raggedright
Mechanism
\end{minipage} & \begin{minipage}[b]{\linewidth}\raggedright
Example
\end{minipage} \\
\midrule\noalign{}
\endhead
\bottomrule\noalign{}
\endlastfoot
\textbf{Temporal} & Different breeding seasons or times of day & Two
species of field cricket: \emph{Gryllus veletis} breeds in spring;
\emph{G. pennsylvanicus} breeds in fall \\
\textbf{Habitat} & Different microhabitats reduce encounter rate &
\emph{Rhagoletis} apple vs.~hawthorn races live on different host
trees \\
\textbf{Behavioral} & Species-specific courtship signals (song, dance,
pheromones, color) & Firefly \emph{Photinus} species: each has a unique
flash pattern recognized primarily by conspecific females \\
\textbf{Mechanical} & Incompatible reproductive structures &
\emph{Salvia} flower species: bee-pollinated species have different
petal structures than hummingbird-pollinated species, preventing pollen
transfer \\
\textbf{Gametic} & Sperm-egg recognition molecules are species-specific
& Sea urchin bindin \hyperref[gl:protein]{\textbf{protein}}:
species-specific sperm-egg binding prevents cross-fertilization \\
\end{longtable}
}

\subsubsection{Postzygotic Barriers}\label{postzygotic-barriers}

\textbf{Hybrid inviability}: Hybrids fail to develop properly due to
Dobzhansky-Muller incompatibilities. Gene regulatory networks that
function within each parental species produce dysfunctional combinations
in hybrids.

\textbf{Hybrid sterility}: Hybrids develop normally but cannot
reproduce. The mule (horse \(\times\) donkey) is the classic example.
Horses have \(2n = 64\) chromosomes; donkeys have \(2n = 62\). The mule
has \(2n = 63\) -- odd chromosome number prevents proper pairing at
meiosis I, resulting in sterility.

\textbf{Hybrid breakdown}: F1 hybrids may be viable and fertile, but F2
or backcross offspring show reduced fitness. This occurs when
\hyperref[gl:recombination]{\textbf{recombination}} in the F1 generation
breaks up co-adapted gene complexes, producing incompatible allele
combinations in the F2.

\subsubsection{Reinforcement}\label{reinforcement}

\textbf{Reinforcement} (also called the Wallace effect) occurs when
natural selection strengthens prezygotic barriers after secondary
contact. If hybrids have low fitness (due to postzygotic barriers),
individuals who mate with members of their own species leave more viable
offspring than those who hybridize. This selects for enhanced mate
discrimination.

Evidence: In \emph{Drosophila} species pairs, female mate discrimination
is stronger in \textbf{sympatric} populations (where the two species
co-occur and could hybridize) than in \textbf{allopatric} populations
(where a single species is present). This pattern -- called
\textbf{reproductive character displacement} -- is a predicted signature
of reinforcement.

\subsubsection{Hybridization and
Introgression}\label{hybridization-and-introgression}

Hybridization does not inevitably lead to reproductive isolation
breakdown. \textbf{Introgression} -- the incorporation of alleles from
one species into the gene pool of another through backcrossing -- can be
an important source of adaptive variation.

\textbf{Neanderthal introgression into modern humans}: Non-African
modern humans carry approximately 1--4\% Neanderthal-derived DNA,
resulting from interbreeding that occurred approximately 50,000--60,000
years ago in the Middle East. Specific Neanderthal alleles have been
maintained by selection because they provided adaptive advantages:

\begin{itemize}
\tightlist
\item
  \textbf{HLA immune genes}: Neanderthal-derived HLA alleles increased
  immune diversity in migrating modern human populations.
\item
  \textbf{Keratin genes}: Neanderthal variants affecting skin and hair
  may have helped modern humans adapt to cold Eurasian climates.
\end{itemize}

\textbf{Denisovan introgression}: Melanesian populations carry
approximately 4--6\% Denisovan-derived DNA. The \textbf{EPAS1} gene
variant that helps Tibetans adapt to high altitude was introgressed from
Denisovans -- a striking example of \textbf{adaptive introgression}
where a borrowed allele provides significant fitness benefits in a new
environment.

Recent evidence also suggests introgression from an unknown ``ghost
lineage'' into West African populations, indicating that archaic
admixture was more widespread than initially recognized.

\begin{quote}
\textbf{Real-World Connection: Archaic Human Introgression}

The discovery that modern humans carry DNA from extinct hominin species
has revolutionized our understanding of human evolution. Far from a
simple out-of-Africa replacement, human evolution involved repeated
episodes of admixture. Neanderthal-derived alleles influence modern
human traits ranging from immune function to depression risk to skin
pigmentation. Some Neanderthal alleles are beneficial (immune genes),
while others are deleterious (associated with increased risk for certain
diseases). Natural selection is gradually purging Neanderthal DNA from
functionally important regions of the genome, while neutral or
beneficial segments persist. This ongoing process of post-introgression
selection provides a window into the fitness consequences of
hybridization over evolutionary time.
\end{quote}

\begin{quote}
\textbf{Concept Check 4:} A horse (\(2n = 64\)) crossed with a donkey
(\(2n = 62\)) produces a mule (\(2n = 63\)). Explain, in terms of
meiosis, why the mule is sterile. What type of reproductive barrier does
this represent?
\end{quote}

\begin{center}\rule{0.5\linewidth}{0.5pt}\end{center}

\subsection{Adaptive Radiation}\label{adaptive-radiation}

\textbf{Adaptive radiation} is the rapid diversification of a single
ancestral lineage into multiple species occupying diverse ecological
niches. Radiations are typically triggered by one of three
circumstances: mass extinction (opening ecological opportunities),
colonization of a new, underexploited territory, or the evolution of a
key innovation that opens new adaptive zones.

\begin{figure}[htbp]
\centering
\includegraphics[width=0.82\linewidth,height=4.2in,keepaspectratio]{../figures/mermaid_inline/inline_mermaid_0110_ee90e97d08fa.png}
\caption{Sequence diagram for Adaptive Radiation showing ordered interaction among Ancestral Population, Geographic Barrier Forms, Population 1, and Population 2.}
\end{figure}

\subsubsection{Hawaiian Honeycreepers}\label{hawaiian-honeycreepers}

The Hawaiian honeycreeper radiation (family Fringillidae, tribe
Drepanidini) produced approximately 55 species from one or two founding
finch-like ancestors that arrived approximately 5 Mya. Bill morphology
diversified to exploit a remarkable range of food sources:

\begin{itemize}
\tightlist
\item
  \textbf{Nectarivores} (long, curved bills): \emph{Drepanis},
  \emph{Hemignathus} -- probe tubular flowers for nectar
\item
  \textbf{Seed crackers} (massive bills): \emph{Pseudonestor},
  \emph{Chloridops} -- crush hard seeds and nuts
\item
  \textbf{Insectivores} (thin, warbler-like bills): \emph{Oreomystis},
  \emph{Paroreomyza} -- glean insects from bark and foliage
\item
  \textbf{Generalists}: \emph{Telespiza} -- omnivorous diet
\end{itemize}

Tragically, over one-third of honeycreeper species are now extinct, and
most survivors are critically endangered, primarily due to avian malaria
(transmitted by introduced mosquitoes), habitat loss, and introduced
predators.

\subsubsection{Cichlid Fish in African Great
Lakes}\label{cichlid-fish-in-african-great-lakes}

The cichlid radiations in Lakes Victoria (about 500 species, about
15,000 years), Malawi (about 800 species, about 2 Mya), and Tanganyika
(about 250 species, about 10 Mya) represent the most species-rich
vertebrate adaptive radiations known.

Key features driving cichlid diversification:

\begin{itemize}
\tightlist
\item
  \textbf{Trophic specialization}: Jaw morphology (particularly the
  pharyngeal jaw apparatus -- a second set of jaws in the throat) has
  diversified to process algae, plankton, snails, fish, scales, and even
  eyes of other fish.
\item
  \textbf{Sexual selection on color}: Female mate choice based on male
  nuptial coloration drives assortative mating. In Lake Victoria,
  species that are distinguished solely by color can hybridize when
  visual cues are disrupted by turbidity.
\item
  \textbf{Rapid speciation}: Genetic divergence between species is often
  minimal, suggesting that speciation was driven primarily by sexual
  selection and ecological specialization rather than accumulated
  genetic distance.
\end{itemize}

\subsubsection{Galapagos Finches}\label{galapagos-finches}

Darwin's finches (15 species across the Galapagos Islands and Cocos
Island) diversified from a single South American ancestor approximately
1--2 Mya. The Grants' long-term research on Daphne Major island
documented natural selection in real time:

\begin{itemize}
\tightlist
\item
  During the 1977 drought, large seeds predominated. Finches with
  deeper, stronger beaks survived preferentially -- directional
  selection increased mean beak depth by 4\% in a single generation.
\item
  \textbf{Character displacement}: On islands where medium ground
  finches (\emph{Geospiza fortis}) co-occur with large ground finches
  (\emph{G. magnirostris}), \emph{G. fortis} beaks are smaller than on
  islands where \emph{G. fortis} occurs alone. This divergence reduces
  competition for seeds.
\end{itemize}

\subsubsection{Mammalian Radiation After K-Pg
Extinction}\label{mammalian-radiation-after-k-pg-extinction}

The extinction of non-avian dinosaurs 66 Mya released enormous
ecological space. Mammals, which had been small, nocturnal insectivores
for over 100 million years, rapidly diversified into most modern orders
within approximately 10 million years. This radiation produced bats
(flight), whales (marine adaptation), elephants (megaherbivory),
primates (arboreality and eventually bipedalism), and the full range of
modern mammalian diversity.

\begin{quote}
\textbf{Concept Check 5:} What conditions are necessary for an adaptive
radiation to occur? Why do island archipelagos and post-extinction
periods seem particularly conducive to adaptive radiation?
\end{quote}

\begin{quote}
\textbf{Concept Check 6:} Apply the structured-coalescent intuition. Two
human populations (effective size \(N_e \approx 10{,}000\) each)
exchange migrants at rate \(m = 0.001\) per generation. (a) What is the
expected pairwise coalescence time within each subpopulation? (b) Why is
the pairwise coalescence time \textbf{longer} for two alleles drawn from
different subpopulations than for two alleles drawn from the same
subpopulation? (c) How does this asymmetry generate \(F_{ST}\)?
\end{quote}

\begin{quote}
\textbf{Concept Check 7:} Bread wheat (\emph{Triticum aestivum}) is an
allohexaploid (AABBDD). Cultivated tobacco is an allotetraploid (SSTT).
Cultivated upland cotton is an allotetraploid (AADD). For each, identify
(i) how many parental species contributed, (ii) why the immediate F1
hybrid was likely sterile before chromosome doubling, and (iii) why the
polyploid is reproductively isolated from each diploid parent. Why is
the analogous story for \textbf{autopolyploidy} less common as a route
to speciation in animals than in plants?
\end{quote}

\begin{quote}
\textbf{Clinical Connection --- Founder effects and Mendelian disease in
isolated populations.} Ashkenazi Jews (about 500-year
bottleneck, \(N_e \approx 250-400\)), Finns (isolation since the
Neolithic), Amish, and French Canadians carry elevated frequencies of
specific recessive alleles (Tay-Sachs \emph{HEXA}, BRCA1/2 185delAG,
Canavan disease \emph{ASPA}, phenylketonuria \emph{PAH}). In many cases
these alleles rose from a single founder chromosome, detectable as a
shared haplotype. Modern targeted screens (e.g.~the Ashkenazi panel) are
possible precisely \emph{because} drift + bottleneck has concentrated a
handful of distinctive alleles; in outbred populations, the allelic
heterogeneity for the same gene is often too great for panel-based
screening to be useful.
\end{quote}

\begin{quote}
\textbf{Clinical Connection --- Speciation of SARS-CoV-2 variants.}
Allopatric speciation analogues appear in real time during pandemics. A
single index case seeding a geographic region generates a local founder
population; drift plus selection for transmissibility produce variants
(Alpha, Delta, Omicron) that then contact and partially displace each
other. Phylogenetic analysis of \textgreater15 million SARS-CoV-2
genomes has documented exactly the processes this chapter formalises ---
drift at low incidence, bottleneck at importation, selection under
\hyperref[gl:vaccine]{\textbf{vaccine}} pressure --- on timescales of
weeks rather than millennia, making clinical virology the fastest-moving
evolutionary laboratory in human history.
\end{quote}

\begin{center}\rule{0.5\linewidth}{0.5pt}\end{center}

\subsection{Computational Bridge}\label{computational-bridge-19}

Wright--Fisher sampling noise is simulated explicitly:

\begin{Shaded}
\begin{Highlighting}[]
\ImportTok{from}\NormalTok{ biology.evolution }\ImportTok{import}\NormalTok{ simulate\_drift}

\NormalTok{traj }\OperatorTok{=}\NormalTok{ simulate\_drift(}\FloatTok{0.5}\NormalTok{, N}\OperatorTok{=}\DecValTok{30}\NormalTok{, generations}\OperatorTok{=}\DecValTok{50}\NormalTok{, rng\_seed}\OperatorTok{=}\DecValTok{7}\NormalTok{)}
\BuiltInTok{print}\NormalTok{(}\BuiltInTok{round}\NormalTok{(traj[}\OperatorTok{{-}}\DecValTok{1}\NormalTok{], }\DecValTok{4}\NormalTok{))}
\end{Highlighting}
\end{Shaded}

\begin{quote}
\textbf{Clinical / systems note:} Founder effects in human isolates
(Amish, Finns) mirror conservation breeding programs where small \(N\)
amplifies drift --- informing genetic screening panels.
\end{quote}

\begin{center}\rule{0.5\linewidth}{0.5pt}\end{center}

\subsubsection{Worked Example: Kimura's Neutral Fixation
Probability}\label{worked-example-kimuras-neutral-fixation-probability}

Motoo Kimura's \textbf{neutral theory of molecular evolution} (1968)
predicts the fate of neutral mutations (selection coefficient \(s = 0\))
under drift alone. The probability that a new neutral mutation (initial
frequency \(1/(2N)\) in a diploid population) \emph{eventually reaches
fixation} is equal to its initial frequency:

\[u_{\text{neutral}} = \dfrac{1}{2N_e} \label{eq:unit_VI_genetic_drift_and_speciation_item_5}\]

\textbf{Worked example.} In a diploid population of \(N_e = 10^4\), a
new neutral mutation has a fixation probability of
\(u = 1/(2 \times 10^4) = 5 \times 10^{-5}\) --- one in 20 000. Of the
mutations that do fix, the mean time to fixation is
\(\bar{t}_{\text{fix}} = 4N_e = 40\,000\) generations; \emph{most}
neutral mutations in fact go extinct within about \(2N_e\)
generations.

At the molecular level, the \textbf{neutral substitution rate} equals
the per-individual-per-generation mutation rate \texorpdfstring{\ensuremath{\mu}}{mu} (remarkable:
substitution rate is \emph{independent of population size} for neutral
alleles, because
\({\rm rate} = 2N\mu \times u = 2N\mu \times 1/(2N) = \mu\)). This is
the molecular clock's theoretical basis: if a site evolves neutrally,
the rate of divergence between lineages is constant regardless of their
\(N_e\) histories. For \textbf{slightly deleterious} mutations, Ohta's
near-neutral theory modifies this to
\(u \approx \dfrac{1 - e^{-2s}}{1 - e^{-4N_e s}}\), which for
\(|4N_e s| < 1\) reduces to the neutral result but for \(|4N_e s| > 1\)
predicts sharp selection-dominated fixation or loss. \textbf{Biological
significance}: the 1 \% sequence divergence between human and chimpanzee
at synonymous sites (about 6 million years of separation,
about \(6 \times 10^5\) generations) gives
\(\mu \approx 1.6 \times 10^{-8}\) per site per generation ---
consistent with pedigree-based mutation rate estimates from trio
whole-genome sequencing (deCODE 2012, 2023). Kimura's simple formula
thus \emph{quantitatively anchors} the molecular clock to the fossil
record.

\begin{center}\rule{0.5\linewidth}{0.5pt}\end{center}

\subsection{Current Evidence and Frontier
Biology}\label{current-evidence-and-frontier-biology-22}

For \textbf{Genetic Drift, Gene Flow, and Speciation}, frontier biology
belongs inside the evidence logic of the chapter. Evolutionary claims
are strongest when they combine mechanism, comparative evidence,
population process, and explicit uncertainty. The core reading question
is this: speciation claims should separate gene flow, reproductive
isolation, demographic history, and genomic architecture.

\begin{itemize}
\tightlist
\item
  \textbf{What to verify:} identify the observation, model, assay, or
  dataset that would make the claim stronger or weaker.
\item
  \textbf{What to qualify:} state the scale, organism, cell type,
  environmental condition, or population where the claim is expected to
  hold.
\item
  \textbf{What to compare:} test at least one alternative explanation,
  baseline, or null model before treating the pattern as causal.
\item
  \textbf{What to cite:} distinguish primary evidence, review synthesis,
  public dataset, and institutional guidance; for recent or numeric
  claims, prefer the source closest to the measurement and state what
  has changed since it was published.
\end{itemize}

Distinguish adaptation from drift, phylogenetic signal from convergence,
and historical explanation from a testable prediction about present-day
data.

\textbf{Source practice:} For evolutionary claims, prefer evidence that
compares alternatives such as selection, drift, gene flow, constraint,
convergence, and shared ancestry.

\subsection{Summary}\label{summary-22}

\begin{itemize}
\tightlist
\item
  \textbf{Genetic drift} is the random change in allele frequencies due
  to finite population size. Its effects are strongest in small
  populations, where \(\text{Var}[\Delta p] = p(1-p)/(2N_e)\).
\item
  \textbf{Effective population size} (\(N_e\)) is typically much smaller
  than census size due to unequal sex ratios, variance in reproductive
  success, and population fluctuations.
\item
  \textbf{Bottlenecks} (cheetah, northern elephant seal) and
  \textbf{founder effects} (Amish, Finnish disease heritage)
  dramatically reduce genetic diversity, with consequences persisting
  for thousands of generations.
\item
  \textbf{Gene flow} homogenizes populations; \(F_{ST}\) measures
  differentiation; Wright's island model predicts
  \(F_{ST} \approx 1/(1 + 4N_em)\).
\item
  \textbf{Species concepts} differ in emphasis: BSC (reproductive
  isolation), morphological (appearance), PSC (monophyly), ecological
  (niche), cohesion (exchangeability). Ring species challenge discrete
  species boundaries.
\item
  \textbf{Allopatric speciation} (geographic isolation) is most common;
  \textbf{peripatric} involves small peripheral populations;
  \textbf{parapatric} occurs along environmental gradients;
  \textbf{sympatric} requires strong disruptive selection or polyploidy.
\item
  \textbf{Reproductive isolation} includes prezygotic (temporal,
  habitat, behavioral, mechanical, gametic) and postzygotic
  (inviability, sterility, breakdown) barriers. Reinforcement
  strengthens prezygotic barriers after secondary contact.
\item
  \textbf{Hybridization and introgression} contribute to adaptation
  (Neanderthal DNA in modern humans; Denisovan EPAS1 in Tibetans).
\item
  \textbf{Adaptive radiations} (honeycreepers, cichlids, finches,
  post-K-Pg mammals) demonstrate rapid diversification following
  ecological opportunity.
\item
  \textbf{Connections:} See \cref{sec:unit_V_population_genetics} for
  \(F_{ST}\) and migration, \cref{sec:unit_VI_phylogenetics} for gene
  trees vs.~species trees, and \cref{sec:unit_X_biomes_and_conservation}
  for conservation genetics.
\end{itemize}

\begin{center}\rule{0.5\linewidth}{0.5pt}\end{center}

\subsection{Key Terms}\label{key-terms-22}

{\def\LTcaptype{none} % do not increment counter
\begin{longtable}[]{@{}
  >{\raggedright\arraybackslash}p{(\linewidth - 2\tabcolsep) * \real{0.2857}}
  >{\raggedright\arraybackslash}p{(\linewidth - 2\tabcolsep) * \real{0.7143}}@{}}
\toprule\noalign{}
\begin{minipage}[b]{\linewidth}\raggedright
Term
\end{minipage} & \begin{minipage}[b]{\linewidth}\raggedright
Definition
\end{minipage} \\
\midrule\noalign{}
\endhead
\bottomrule\noalign{}
\endlastfoot
\textbf{Genetic drift} & Random change in allele frequencies due to
sampling error in finite populations \\
\textbf{Effective population size (\(N_e\))} & Size of an ideal
Wright-Fisher population experiencing equivalent drift to the real
population \\
\textbf{Bottleneck effect} & Severe temporary reduction in population
size causing loss of genetic diversity \\
\textbf{Founder effect} & Establishment of a new population by a small,
genetically non-representative group of colonizers \\
\textbf{Gene flow} & Movement of alleles between populations through
migration and reproduction \\
\textbf{\(F_{ST}\)} & Fixation index measuring the proportion of genetic
variation between (versus within) populations \\
\textbf{Isolation by distance} & Positive correlation between genetic
differentiation and geographic distance \\
\textbf{Biological species concept} & Species defined as reproductively
isolated groups of interbreeding populations \citep{mayr1942} \\
\textbf{Phylogenetic species concept} & Species defined as the smallest
diagnosable monophyletic group \\
\textbf{Allopatric speciation} & Speciation resulting from geographic
isolation of populations \\
\textbf{Sympatric speciation} & Speciation within a single geographic
area without physical barriers \\
\textbf{Peripatric speciation} & Speciation involving a small peripheral
population experiencing founder effects \\
\textbf{Dobzhansky-Muller incompatibility} & Epistatic incompatibility
in hybrids from independently evolved alleles at interacting loci \\
\textbf{Reinforcement} & Natural selection strengthening prezygotic
barriers when hybrids have reduced fitness \\
\textbf{Adaptive radiation} & Rapid diversification of a lineage into
multiple species exploiting diverse ecological niches \\
\textbf{Allopolyploidy} & Speciation via hybridization between species
followed by whole-genome duplication \\
\textbf{Introgression} & Incorporation of alleles from one species into
another through hybridization and backcrossing \\
\textbf{Ring species} & Chain of interbreeding populations forming a
geographic ring, with terminal populations reproductively isolated \\
\end{longtable}
}

\begin{center}\rule{0.5\linewidth}{0.5pt}\end{center}

\subsection{Review Questions}\label{review-questions-22}

\begin{enumerate}
\def\labelenumi{\arabic{enumi}.}
\item
  A population of 200 individuals has initial heterozygosity
  \(H_0 = 0.80\). Calculate the expected heterozygosity after 50
  generations of pure drift (no mutation, no selection). How many
  generations until heterozygosity drops below 0.40?
\item
  Explain why the effective population size of northern elephant seals
  is far smaller than their current census size of 200,000. What
  specific aspects of their biology and history contribute to low
  \(N_e\)?
\item
  Two island populations have \(N_e = 300\) each. Calculate the
  equilibrium \(F_{ST}\) if (a) \(m = 0.001\), (b) \(m = 0.01\), and (c)
  \(m = 0.1\) using Wright's island model. What biological conclusion
  can you draw about the relationship between migration rate and
  population differentiation?
\item
  The Amish founder effect produced a frequency of the Ellis-van Creveld
  allele approximately 8,750 times higher than in the general
  population. Explain, using the concepts of founder effect and genetic
  drift, how a rare allele can reach such high frequency. Would natural
  selection alone predict this outcome?
\item
  Compare the biological and phylogenetic species concepts as applied to
  a ring species such as \emph{Ensatina} salamanders. Which concept
  handles this situation more satisfactorily, and why?
\item
  Describe the formation of bread wheat (\emph{Triticum aestivum})
  through allopolyploidy. Why does polyploidy create instantaneous
  reproductive isolation? What is the evolutionary significance of the
  observation that about 15\% of angiosperm speciation events involve
  polyploidy?
\item
  Neanderthal DNA constitutes approximately 1--4\% of non-African modern
  human genomes. Explain how adaptive introgression of the EPAS1 allele
  from Denisovans illustrates that hybridization can be evolutionarily
  beneficial rather than merely a breakdown of species boundaries.
\item
  Compare and contrast the cichlid radiation in Lake Victoria with
  Darwin's finch radiation on the Galapagos. What factors (sexual
  selection, ecological opportunity, geographic isolation) played
  different roles in each case? Why has the cichlid radiation produced
  far more species?
\item
  A conservation biologist is managing a population of 25 endangered
  wolves. The population has \(H_0 = 0.65\). (a) How rapidly will
  heterozygosity decline per generation? (b) If genetic rescue
  introduces 5 individuals from a genetically distinct population per
  generation, qualitatively describe how this affects the trajectory.
  (c) What are the risks and benefits of genetic rescue?
\item
  Explain why Dobzhansky-Muller incompatibilities accumulate as the
  square of the number of substitutions (Orr's snowball effect). What
  does this predict about the relationship between genetic divergence
  time and the strength of reproductive isolation?
\item
  Re-run \texttt{simulate\_drift} with \(N=200\) vs.~\(N=20\) for the
  same number of generations. Which trajectory shows larger
  \(|p_{50}-p_0|\) on average?
\item
  Contrast \textbf{reinforcement} with \textbf{fusion} at secondary
  contact; give one genomic signature expected under each. \#\# Further
  Reading and Source Notes
\end{enumerate}

\begin{itemize}
\tightlist
\item
  Mayr (1942). \emph{Systematics and the Origin of Species}. Columbia
  University Press.
\item
  Templeton (1989). The meaning of species and speciation: A genetic
  perspective. Sinauer.
\end{itemize}

\begin{center}\rule{0.5\linewidth}{0.5pt}\end{center}

\subsubsection{Companion Source
Module}\label{companion-source-module-22}

\textbf{Genetic Drift, Gene Flow, and Speciation} should leave a
reproducible trail from a biological claim to the code, figure, diagram,
or paper-based activity that can test it. Use the surfaces below to
inspect the chapter's assumptions, rerun the relevant model, or compare
the manuscript explanation with companion labs and figures.

{\def\LTcaptype{none} % do not increment counter
\begin{longtable}[]{@{}
  >{\raggedright\arraybackslash}p{(\linewidth - 2\tabcolsep) * \real{0.5000}}
  >{\raggedright\arraybackslash}p{(\linewidth - 2\tabcolsep) * \real{0.5000}}@{}}
\toprule\noalign{}
\begin{minipage}[b]{\linewidth}\raggedright
Surface
\end{minipage} & \begin{minipage}[b]{\linewidth}\raggedright
Use it for
\end{minipage} \\
\midrule\noalign{}
\endhead
\bottomrule\noalign{}
\endlastfoot
\breaktt{src/biology/evolution/evolution.py}
(\breaktt{wright\_fisher\_drift}, \texttt{simulate\_drift},
\texttt{isolation\_index}) & Compare stochastic drift, bottlenecks, and
isolation measures. \\
\breaktt{src/mermaid/biology\_diagrams.py} (\breaktt{speciation\_diagram},
\breaktt{phylogenetic\_tree\_diagram}) & Connect reproductive isolation
to lineage divergence. \\
\end{longtable}
}

\textbf{Reproducibility check:} report effective population size,
migration, selection possibility, and time scale before assigning
divergence to drift or speciation. \textbf{Cross-reference:} compare
with \cref{sec:unit_V_population_genetics},
\cref{sec:unit_VI_evolution_and_selection}, and
\cref{sec:unit_VI_phylogenetics}.

\newpage

\section{Phylogenetics and the Tree of
Life}\label{phylogenetics-and-the-tree-of-life}

\label{sec:unit_VI_phylogenetics}

\begin{quote}
\textbf{Ch 21} \texorpdfstring{\ensuremath{\cdot}}{.} Level 3/3 \texorpdfstring{\ensuremath{\cdot}}{.} 60 min read \texorpdfstring{\ensuremath{\cdot}}{.} 100 min lecture \texorpdfstring{\ensuremath{\cdot}}{.}
Prerequisites: \cref{sec:unit_VI_genetic_drift_and_speciation}
\end{quote}

\subsection{Learning Objectives}\label{learning-objectives-23}

By the end of this chapter, you should be able to:

\begin{enumerate}
\def\labelenumi{\arabic{enumi}.}
\tightlist
\item
  Explain what phylogenetics reveals about evolutionary relationships,
  ancestral states, divergence timing, and biogeographic origins.
\item
  Read and interpret phylogenetic trees, distinguishing between clades,
  paraphyletic groups, polyphyletic groups, and key character types
  (synapomorphy, symplesiomorphy, homoplasy).
\item
  Compare phylogenetic reconstruction methods: distance-based (UPGMA,
  Neighbor-Joining), maximum parsimony, maximum likelihood, and Bayesian
  inference.
\item
  Apply \hyperref[gl:molecular-clock]{\textbf{molecular clock}}
  calculations to estimate divergence times and evaluate the assumptions
  and limitations of clock models.
\item
  Describe the three-domain tree of life, including endosymbiotic
  origins of \hyperref[gl:organelle]{\textbf{organelle}}s and the
  significance of Asgard \hyperref[gl:archaea]{\textbf{archaea}}.
\item
  Outline human phylogeny, including fossil hominin diversity and
  archaic introgression events.
\item
  Calculate a Jukes-Cantor corrected distance from observed sequence
  divergence and estimate a divergence time, explaining why the
  correction is needed and why mitochondrial estimates differ from
  nuclear calibrations.
\item
  Explain how long-read sequencing (PacBio HiFi, Oxford Nanopore)
  resolves deep-divergence and polyploid phylogenies that short reads
  cannot, using the animal tree of life and ancient-DNA examples.
\item
  Evaluate a phylogenetic claim by identifying the dataset, model, and
  null hypothesis needed to distinguish shared ancestry from convergence
  and to assess sampling and calibration uncertainty.
\end{enumerate}

\subsubsection{Study Blueprint}\label{study-blueprint-23}

\begin{itemize}
\tightlist
\item
  \textbf{Big idea:} Phylogenies are evidence-based hypotheses about
  ancestry, not ladders of progress.
\item
  \textbf{Core concepts:} homology, tree topology, parsimony, molecular
  clocks.
\item
  \textbf{Framework alignment:} Vision \& Change: Evolution, Systems; AP
  Biology: Evolution, Systems Interactions; NGSS-style topics: Natural
  Selection and Evolution, Interdependent Relationships in Ecosystems.
\item
  \textbf{Model or quantitative lens:} Tree-distance, parsimony, and
  molecular-clock calculations.
\item
  \textbf{Data skill:} Read trees correctly and map traits or sequences
  onto branches.
\item
  \textbf{Practice cadence:} Visual Representations, Statistical Tests
  and Data Analysis, Argumentation.
\item
  \textbf{Common misconception to repair:} Living species are cousins,
  not ancestors of one another.
\item
  \textbf{Primary lab:} \cref{sec:lab_unit_VI_phylogenetics}.
\item
  \textbf{Question bank:} \cref{sec:q_unit_VI_phylogenetics}.
\item
  \textbf{Transfer task:} Transfer tree thinking to pathogens,
  conservation units, gene families, and development.
\item
  \textbf{Bridge to computation:}
  \breaktt{biology.genetics.genetics.jukes\_cantor\_distance}.
\end{itemize}

\begin{center}\rule{0.5\linewidth}{0.5pt}\end{center}

\begin{quote}
\textbf{Opening Vignette --- The Tree That Changed Everything}

In the 1970s, Carl Woese was alone in his University of Illinois
laboratory, painstakingly sequencing ribosomal RNA by subjecting it to
nuclease digestion and then separating fragments on two-dimensional
electrophoresis gels --- an agonisingly laborious technique that took
months per organism. When he sequenced methanogens (thought to be
unusual bacteria), the rRNA pattern was unlike any bacterium. It was
unlike any \hyperref[gl:eukaryote]{\textbf{eukaryote}} too. Woese
concluded they formed a third domain of life --- the Archaea. His 1977
\emph{PNAS} paper was widely dismissed for nearly a decade: a
microbiologist claiming to overturn a century of classification with RNA
fingerprints? Yet the phylogenetic evidence was incontrovertible, and by
the 1990s the three-domain tree had rearranged most of biology. The
lesson: molecular sequence data can reveal evolutionary relationships
invisible to morphology, and great paradigm shifts are often dismissed
before they are accepted.
\end{quote}

\subsection{What Phylogenetics Tells
Us}\label{what-phylogenetics-tells-us}

\textbf{Phylogenetics} is the study of evolutionary relationships among
organisms, inferred from heritable characters (DNA sequences,
\hyperref[gl:protein]{\textbf{protein}} sequences, morphological traits,
behavioral traits). A phylogenetic analysis produces a tree (or network)
that represents the historical pattern of descent from common ancestors.

Phylogenetics addresses four fundamental questions:

\begin{enumerate}
\def\labelenumi{\arabic{enumi}.}
\tightlist
\item
  \textbf{Evolutionary relationships}: Which organisms are most closely
  related? This is the foundation of modern biological classification --
  taxonomy based on shared ancestry rather than superficial similarity.
\item
  \textbf{Ancestral character states}: What did the ancestors look like?
  By mapping traits onto phylogenies, we can reconstruct the sequence in
  which characters evolved (e.g., when did flight originate in
  insects?).
\item
  \textbf{Timing of diversification}: When did lineages diverge?
  Molecular clocks calibrated with fossil data allow estimation of
  divergence dates even when the fossil record is incomplete.
\item
  \textbf{Biogeographic origins}: Where did lineages originate and how
  did they disperse? Phylogeography combines phylogenetics with
  geographic data to trace the spatial history of lineages.
\end{enumerate}

\subsubsection{Applications Beyond
Systematics}\label{applications-beyond-systematics}

Phylogenetics has become indispensable across biology and medicine:

\begin{itemize}
\tightlist
\item
  \textbf{Drug discovery}: Phylogenetic analysis identifies which
  organisms are likely to produce useful bioactive compounds. If a
  species in one clade produces a valuable compound, closely related
  species in the same clade are promising candidates for novel variants.
\item
  \textbf{Epidemiology}: Genomic phylogenetics tracks pathogen spread in
  real time (SARS-CoV-2, HIV, Ebola). Phylogenetic clustering reveals
  transmission chains, identifies superspreader events, and dates
  outbreak origins.
\item
  \textbf{Conservation}: \textbf{Phylogenetic diversity (PD)} metrics
  identify lineages that represent the greatest amount of unique
  evolutionary history. Losing a phylogenetically isolated species
  (e.g., the tuatara, sole survivor of order Rhynchocephalia) eliminates
  more evolutionary history than losing one species from a large,
  diverse clade.
\end{itemize}

\begin{quote}
\textbf{Concept Check 1:} Why might a conservation strategy based on
phylogenetic diversity differ from one based simply on species counts?
Give an example of a species with high phylogenetic distinctiveness.
\end{quote}

\subsubsection{Phylogenetic Diversity and the EDGE
Framework}\label{phylogenetic-diversity-and-the-edge-framework}

\textbf{Phylogenetic diversity (PD)}, formalized by Daniel Faith (1992),
is the \textbf{total branch length} of the phylogeny subtending a given
set of species. A clade of 100 species that diverged 5 Ma has lower PD
than a clade of 30 species spanning 200 Ma, because PD measures
evolutionary history rather than species count alone. The \textbf{EDGE
(Evolutionarily Distinct and Globally Endangered)} initiative of the
Zoological Society of London uses PD-based metrics to prioritize
conservation:

\begin{equation}
\text{EDGE score} = \log(1 + \text{ED}) + \text{GE} \cdot \log(2)
\label{eq:unit_VI_edge_score}
\end{equation}

where \textbf{ED} (evolutionary distinctness) is each species' share of
the phylogeny it sits in, weighted by branch length, and \textbf{GE}
(global endangerment) is the IUCN Red List category as a numerical score
(Least Concern = 0, Critically Endangered = 4). Top EDGE species are
both phylogenetically isolated and at high extinction risk --- their
loss would erase irreplaceable amounts of evolutionary history.

\textbf{EDGE flagship examples}:

\begin{itemize}
\tightlist
\item
  \textbf{Tuatara} (\emph{Sphenodon punctatus}) --- sole extant member
  of order Rhynchocephalia, sister group to most squamates (lizards and
  snakes). The tuatara's lineage diverged about 250 Ma, making
  it one of the most phylogenetically distinct vertebrates alive.
  Restricted to about 30 small islands off New Zealand.
\item
  \textbf{Pygmy hippopotamus} (\emph{Choeropsis liberiensis}) --- sole
  survivor of one of two extant hippo lineages, with deep evolutionary
  history.
\item
  \textbf{Aye-aye} (\emph{Daubentonia madagascariensis}) --- sole member
  of the family Daubentoniidae, an evolutionarily distinct lemur lineage
  on Madagascar.
\item
  \textbf{Yangtze giant softshell turtle} (\emph{Rafetus swinhoei}) ---
  fewer than four known living individuals; an entire evolutionary line
  nearly extinct.
\end{itemize}

The EDGE approach contrasts sharply with species-count-based
conservation. A reserve protecting 100 closely related warbler species
preserves much less evolutionary history than a reserve protecting 30
phylogenetically diverse tetrapods --- a counterintuitive but
biologically meaningful prioritization rule.

\begin{center}\rule{0.5\linewidth}{0.5pt}\end{center}

\subsection{Tree Terminology and Reading
Phylogenies}\label{tree-terminology-and-reading-phylogenies}

\begin{figure}[htbp]
\centering
\includegraphics[width=0.82\linewidth,height=4.2in,keepaspectratio]{../figures/mermaid_inline/inline_mermaid_0111_66222deabd2a.png}
\caption{Flowchart for tree terminology and reading phylogenies: root (MRCA of sampled taxa), internal node (divergence event), outgroup, and branch labels define how to read relatedness.}
\end{figure}

\textbf{Key: A and B are sister taxa. The clade (A, B, C) is
monophyletic. The outgroup roots the tree.}

\subsubsection{Fundamental Tree
Components}\label{fundamental-tree-components}

\textbf{Root}: The node representing the most recent common ancestor
(MRCA) of the sampled taxa in the tree. An \textbf{outgroup} -- a taxon
known to be outside the group of interest (the ingroup) -- is used to
determine the position of the root on an unrooted tree.

\textbf{Internal nodes}: Each internal node represents a divergence
event -- a \hyperref[gl:speciation]{\textbf{speciation}} event where an
ancestral lineage split into two daughter lineages. Nodes are
hypothetical ancestors; they are not observed directly but inferred from
the data.

\textbf{Branches}: The lines connecting nodes. In a \textbf{cladogram},
branch lengths have no meaning (primarily topology matters). In a
\textbf{phylogram}, branch lengths represent the amount of evolutionary
change (e.g., substitutions per site). In a \textbf{chronogram}
(ultrametric tree), branch lengths represent time.

\textbf{Tips (leaves)}: The terminal nodes, representing observed taxa
-- either extant species or fossil specimens.

\textbf{Clade (monophyletic group)}: An ancestor and its descendant
lineage. Clades are the preferred units in modern phylogenetic
classification because they preserve common ancestry. The clade Mammalia
includes the descendants of the MRCA of mammals -- bats, whales, humans,
and monotremes alike.

\textbf{Sister groups}: Two clades that share an immediate common
ancestor. In the diagram above, Taxon A and Taxon B are sister taxa;
their clade is sister to Taxon C.

\textbf{Polytomy}: An unresolved node where three or more lineages
diverge simultaneously. A ``hard'' polytomy reflects genuine
simultaneous speciation (rare); a ``soft'' polytomy reflects
insufficient data to resolve the branching order (common).

\subsubsection{Paraphyletic and Polyphyletic
Groups}\label{paraphyletic-and-polyphyletic-groups}

\textbf{Paraphyletic group}: An ancestor and SOME (but not most) of its
descendants. The traditional ``Reptilia'' (turtles, lizards, snakes,
crocodilians) is paraphyletic because it excludes birds, which are
nested within the reptile clade. Modern systematics recognizes that
crocodilians are more closely related to birds than to lizards -- a fact
that the traditional grouping obscures.

\textbf{Polyphyletic group}: A group whose members do NOT share a most
recent common ancestor exclusive to that group. ``Warm-blooded animals''
(birds + mammals) is polyphyletic because endothermy evolved
independently in each lineage. Polyphyletic groups reflect convergent
evolution, not shared ancestry.

\subsubsection{Character Types}\label{character-types}

\textbf{Synapomorphy}: A shared derived character that defines a clade.
The amniotic egg is a synapomorphy of Amniota (reptiles, birds, and
mammals). Synapomorphies are the characters most informative for
identifying clades.

\textbf{Symplesiomorphy}: A shared ancestral character inherited from a
deeper ancestor. The vertebral column is shared by most vertebrates but
does not diagnose any particular clade within vertebrates -- it is a
symplesiomorphy for any subgroup of vertebrates.

\textbf{Autapomorphy}: A unique derived character found in a single
taxon. Feathers are an autapomorphy of Aves (or more precisely, of
Maniraptora within the dinosaur phylogeny, though some non-avian
maniraptorans also had feathers).

\textbf{Homoplasy}: A character state that appears similar in two or
more taxa but was NOT inherited from their common ancestor. Homoplasy
arises through \textbf{convergent evolution} (independent evolution of
similar features -- e.g., wings in bats and birds) or \textbf{reversal}
(return to an ancestral state). Homoplasy confounds phylogenetic
analysis because it falsely suggests shared ancestry.

\subsubsection{Rooting Trees}\label{rooting-trees}

An \textbf{unrooted tree} shows the relationships among taxa but does
not indicate the direction of evolutionary time. To root a tree (and
thereby determine which character states are ancestral versus derived),
one of two approaches is used:

\begin{itemize}
\tightlist
\item
  \textbf{Outgroup rooting}: Include a taxon known to be outside the
  group of interest. The root is placed on the branch connecting the
  outgroup to the ingroup. This is the most common method. For example,
  when constructing a phylogeny of mammals, a reptile (e.g., a
  crocodilian) serves as the outgroup.
\item
  \textbf{Midpoint rooting}: Place the root at the midpoint of the
  longest path between any two taxa. This assumes approximately equal
  rates of evolution across lineages -- an assumption that is often
  violated.
\end{itemize}

\subsubsection{Reading Trees Correctly}\label{reading-trees-correctly}

A critical skill is reading phylogenetic trees by \textbf{topology}
(branching pattern) rather than by the visual arrangement of tips. Trees
can be rotated around any internal node without changing the
evolutionary relationships they depict. The common error of reading
relationships from left to right along the tips (treating one end as
``primitive'' and the other as ``advanced'') is incorrect. Evolution
does not proceed along a ladder from simple to complex; every living
species is equally evolved --- each has been evolving for the same
amount of time since the last comprehensive common ancestor.

\paragraph{Common tree-thinking
misconceptions}\label{common-tree-thinking-misconceptions}

Decades of cognitive-science research on student tree-reading reveal a
recurring set of errors. Recognizing these misconceptions is the
foundation of phylogenetic literacy:

{\def\LTcaptype{none} % do not increment counter
\begin{longtable}[]{@{}
  >{\raggedright\arraybackslash}p{(\linewidth - 2\tabcolsep) * \real{0.5769}}
  >{\raggedright\arraybackslash}p{(\linewidth - 2\tabcolsep) * \real{0.4231}}@{}}
\toprule\noalign{}
\begin{minipage}[b]{\linewidth}\raggedright
Misconception
\end{minipage} & \begin{minipage}[b]{\linewidth}\raggedright
Correction
\end{minipage} \\
\midrule\noalign{}
\endhead
\bottomrule\noalign{}
\endlastfoot
\textbf{Ladder thinking}: Trees show a progression from ``primitive''
tips on one side to ``advanced'' tips on the other. & Trees can be
rotated around any node; tip order has no biological meaning. Every tip
is equally distant from the root in time. \\
\textbf{Reading along the tips}: Adjacent tips are necessarily more
closely related than non-adjacent tips. & Adjacency on the page does not
imply evolutionary closeness. Look at the \textbf{branching topology}
--- count the nodes shared on the path back to the MRCA. \\
\textbf{``More evolved''}: Some tips (``higher'' organisms) are more
evolved than others. & Most living species have been evolving for the
same time. There is no ``scala naturae''. A bacterium's lineage is as
long as a human's. \\
\textbf{The root is the oldest extant species}: The root represents an
existing ancestor whose features can be observed today. & The root is a
hypothetical reconstructed ancestor, not a living species. Living
species are at the tips, not the root. \\
\textbf{A polytomy means simultaneous origin}: Three branches emerging
from a single node mean three lineages arose at the same instant. & Most
polytomies are ``soft'' --- they reflect insufficient data to resolve
close branching events, not actual simultaneous origin. Hard polytomies
(true simultaneous divergence) are rare. \\
\textbf{Branch length typically means time}: Most trees show time on a
horizontal axis. & Primarily \textbf{chronograms} (ultrametric trees)
show time. \textbf{Cladograms} show primarily topology.
\textbf{Phylograms} show evolutionary change (substitutions). Typically
check the scale bar. \\
\textbf{Convergent traits = relatedness}: Similar traits (wings in bats
and birds, dorsal fins in dolphins and sharks) imply close relationship.
& Convergent traits are \textbf{homoplasy} --- they reflect similar
selection pressures, not shared ancestry. Use synapomorphies, not
overall similarity, to define clades. \\
\textbf{Cousins versus ancestors}: Living species can be ancestors of
other living species. & Living species are sister taxa to each other,
rarely ancestor--descendant. The MRCA of two living species is a
hypothetical extinct ancestor, even if one species is ``older'' or
``less changed.'' \\
\end{longtable}
}

\paragraph{Tips for reading phylogenetic
trees}\label{tips-for-reading-phylogenetic-trees}

\begin{enumerate}
\def\labelenumi{\arabic{enumi}.}
\tightlist
\item
  \textbf{Identify the root} --- it represents the most recent common
  ancestor of the sampled taxa shown. Check the scale bar and any axis
  labels.
\item
  \textbf{Trace lineages back to the MRCA} --- don't read across tips
  horizontally. Find where two taxa share their nearest common ancestor
  by going \textbf{down} to the node where their lineages converge.
\item
  \textbf{Look at the branching pattern, not the layout} --- the same
  topology can be displayed in many visually distinct ways (rectangular,
  slanted, radial, circular).
\item
  \textbf{Distinguish character types} --- mapped traits should be
  inferred from synapomorphies (shared derived characters), not from
  overall similarity.
\item
  \textbf{Note the scale} --- a chronogram with branch lengths in
  millions of years tells a different story from a phylogram with branch
  lengths in substitutions/site.
\end{enumerate}

\begin{quote}
\textbf{Concept Check 2:} Consider the traditional classification
``Reptilia'' that includes turtles, lizards, snakes, and crocodilians
but excludes birds. Is this group monophyletic, paraphyletic, or
polyphyletic? Explain your reasoning based on the phylogenetic position
of birds relative to crocodilians.
\end{quote}

\begin{center}\rule{0.5\linewidth}{0.5pt}\end{center}

\subsection{Phylogenetic Reconstruction
Methods}\label{phylogenetic-reconstruction-methods}

\begin{figure}[htbp]
\centering
\includegraphics[width=0.82\linewidth,height=4.2in,keepaspectratio]{../figures/mermaid_inline/inline_mermaid_0112_065c5c6c1e41.png}
\caption{Flowchart for Phylogenetic Reconstruction Methods: Phylogenetic Methods, Distance-Based, Character-Based, and UPGMA form the diagram's primary path or branches.}
\end{figure}

\subsubsection{Distance Methods}\label{distance-methods}

\textbf{UPGMA (Unweighted Pair Group Method with Arithmetic Mean)}: The
simplest clustering algorithm. It assumes a strict molecular clock
(equal rates of evolution across most lineages) and produces an
ultrametric tree where most tips are equidistant from the root. UPGMA
works by iteratively joining the two most similar sequences and
recalculating distances.

\textbf{Limitation}: If the molecular clock assumption is violated (and
it usually is), UPGMA produces incorrect topologies. It is not
recommended for most phylogenetic analyses but remains useful for
constructing guide trees in sequence alignment.

\textbf{Neighbor-Joining (NJ)}: Developed by \citet{saitou1987}, NJ does
NOT assume a molecular clock. It works by star decomposition: starting
with an unresolved star tree, it iteratively identifies the pair of taxa
whose joining minimizes the total branch length. NJ is fast (\(O(n^3)\)
time complexity), produces additive trees with unequal branch lengths,
and provides a good starting tree for more sophisticated analyses. It is
appropriate for large datasets where ML or Bayesian methods would be
computationally prohibitive.

\subsubsection{Worked Example: One UPGMA Clustering Step from a Distance
Matrix}\label{worked-example-one-upgma-clustering-step-from-a-distance-matrix}

\textbf{Problem:} Four taxa (A, B, C, D) have the symmetric pairwise
distance matrix below (substitutions/site). Perform the first UPGMA
clustering step: identify the closest pair, place their common node, and
recompute the reduced distance matrix.

{\def\LTcaptype{none} % do not increment counter
\begin{longtable}[]{@{}lllll@{}}
\toprule\noalign{}
& A & B & C & D \\
\midrule\noalign{}
\endhead
\bottomrule\noalign{}
\endlastfoot
\textbf{A} & --- & 0.10 & 0.40 & 0.42 \\
\textbf{B} & 0.10 & --- & 0.38 & 0.40 \\
\textbf{C} & 0.40 & 0.38 & --- & 0.12 \\
\textbf{D} & 0.42 & 0.40 & 0.12 & --- \\
\end{longtable}
}

\textbf{Solution:}

\begin{enumerate}
\def\labelenumi{\arabic{enumi}.}
\item
  \textbf{Identify the smallest off-diagonal distance.} Scanning the
  matrix, the minimum is \(d_{AB} = 0.10\) substitutions/site. Join A
  and B into cluster \((AB)\).
\item
  \textbf{Place the internal node.} UPGMA builds an ultrametric tree, so
  the node \(U\) joining A and B sits at half the pair distance; the
  branch length from each of A and B up to \(U\) is

  \begin{equation}
  h_U = \frac{d_{AB}}{2} = \frac{0.10}{2} = 0.05 \text{ substitutions/site}
  \label{eq:phylogenetics_8}
  \end{equation}
\item
  \textbf{Recompute distances from the new cluster} as the size-weighted
  arithmetic mean over member pairs (here \(|A| = |B| = 1\)):

  \begin{equation}
  d_{(AB),C} = \frac{d_{AC} + d_{BC}}{2} = \frac{0.40 + 0.38}{2} = 0.39, \qquad
  d_{(AB),D} = \frac{d_{AD} + d_{BD}}{2} = \frac{0.42 + 0.40}{2} = 0.41
  \label{eq:phylogenetics_9}
  \end{equation}

  The distance \(d_{CD} = 0.12\) is unchanged.
\end{enumerate}

\textbf{Interpretation:} A and B cluster first at node height 0.05
substitutions/site; in the reduced matrix the smallest remaining
distance is \(d_{CD} = 0.12\), so the next step joins C and D, yielding
the tree \(((A,B),(C,D))\) --- exactly what an equal-rate (clock-like)
dataset should produce.

\subsubsection{Maximum Parsimony}\label{maximum-parsimony}

\textbf{Maximum parsimony} selects the tree (or trees) requiring the
fewest total character-state changes. This is an application of Occam's
razor -- the simplest explanation is preferred.

\textbf{Procedure}: For each possible tree topology, map character
changes onto branches to find the minimum number of steps. The tree
requiring the fewest total steps across most characters is the most
parsimonious.

\textbf{Strengths}: Intuitive; does not require an explicit model of
evolution; works well when evolutionary rates are low and homoplasy is
rare.

\textbf{Limitations}: - \textbf{Computationally NP-hard}: The number of
possible tree topologies grows super-exponentially with the number of
taxa. For \(n\) taxa, the number of unrooted bifurcating trees is
\((2n-5)!! = (2n-5)(2n-7)(2n-9)\cdots(3)(1)\). For 10 taxa: 2,027,025
trees. For 20 taxa: \(> 10^{21}\) trees. Heuristic search algorithms are
required. - \textbf{Long-branch attraction (LBA)}: When two distantly
related lineages evolve rapidly, they accumulate homoplasies (convergent
substitutions) that parsimony interprets as synapomorphies, incorrectly
grouping them together. LBA is a systematic error -- adding more data
makes it worse, not better.

\subsubsection{Worked Example: Computing a Parsimony Score for a Small
Character
Matrix}\label{worked-example-computing-a-parsimony-score-for-a-small-character-matrix}

\textbf{Problem:} Four taxa (A, B, C, D) are scored for four binary
morphological characters (state 0 or 1). Using the fixed rooted tree
\(((A,B),(C,D))\), compute the total parsimony score (minimum number of
character-state changes).

{\def\LTcaptype{none} % do not increment counter
\begin{longtable}[]{@{}lllll@{}}
\toprule\noalign{}
Taxon & Char 1 & Char 2 & Char 3 & Char 4 \\
\midrule\noalign{}
\endhead
\bottomrule\noalign{}
\endlastfoot
\textbf{A} & 0 & 0 & 1 & 0 \\
\textbf{B} & 0 & 1 & 1 & 0 \\
\textbf{C} & 1 & 1 & 0 & 1 \\
\textbf{D} & 1 & 1 & 0 & 1 \\
\end{longtable}
}

\textbf{Solution:}

Apply the Fitch bottom-up algorithm to each character independently. At
each internal node take the intersection of its two children's state
sets; if the intersection is empty, take the union and add one step.
Internal nodes: \(X\) = ancestor of \((A,B)\), \(Y\) = ancestor of
\((C,D)\), \(R\) = root.

\begin{enumerate}
\def\labelenumi{\arabic{enumi}.}
\item
  \textbf{Char 1} (\(A{=}0, B{=}0, C{=}1, D{=}1\)):
  \(X = \{0\}\cap\{0\}=\{0\}\) (0 steps); \(Y = \{1\}\cap\{1\}=\{1\}\)
  (0 steps); \(R = \{0\}\cap\{1\}=\varnothing \Rightarrow \{0,1\}\)
  (\textbf{1 step}). Subtotal = 1.
\item
  \textbf{Char 2} (\(A{=}0, B{=}1, C{=}1, D{=}1\)):
  \(X = \{0\}\cap\{1\}=\varnothing \Rightarrow \{0,1\}\) (\textbf{1
  step}); \(Y = \{1\}\cap\{1\}=\{1\}\) (0 steps);
  \(R = \{0,1\}\cap\{1\}=\{1\}\) (0 steps). Subtotal = 1.
\item
  \textbf{Char 3} (\(A{=}1, B{=}1, C{=}0, D{=}0\)): \(X=\{1\}\) (0
  steps); \(Y=\{0\}\) (0 steps);
  \(R=\{1\}\cap\{0\}=\varnothing\Rightarrow\{0,1\}\) (\textbf{1 step}).
  Subtotal = 1.
\item
  \textbf{Char 4} (\(A{=}0, B{=}1, C{=}1, D{=}1\)): same pattern as Char
  2 \texorpdfstring{\ensuremath{\to}}{->} \textbf{1 step}. Subtotal = 1.
\item
  \textbf{Total parsimony score:}

  \begin{equation}
  S = 1 + 1 + 1 + 1 = 4 \text{ steps}
  \label{eq:phylogenetics_10}
  \end{equation}
\end{enumerate}

\textbf{Interpretation:} The tree \(((A,B),(C,D))\) requires 4 changes;
the two alternative unrooted topologies \(((A,C),(B,D))\) and
\(((A,D),(B,C))\) each require 7, so \(((A,B),(C,D))\) is the
most-parsimonious tree --- Characters 1 and 3 are the
parsimony-informative synapomorphies that group \((A,B)\) apart from
\((C,D)\).

\subsubsection{Maximum Likelihood (ML)}\label{maximum-likelihood-ml}

Maximum likelihood selects the tree that maximizes the probability of
the observed data given the tree topology, branch lengths, and an
explicit model of sequence evolution:

\begin{equation}
L(\text{tree}, \theta) = \prod_{i=1}^{L} P(\text{site}_i \mid \text{tree}, \theta)
\label{eq:phylogenetics_1}
\end{equation}

where \(L\) is the alignment length and \texorpdfstring{\ensuremath{\theta}}{theta} represents the parameters of
the substitution model.

\textbf{Substitution models} describe the probability of one
\hyperref[gl:nucleotide]{\textbf{nucleotide}} changing to another over a
given evolutionary time:

{\def\LTcaptype{none} % do not increment counter
\begin{longtable}[]{@{}
  >{\raggedright\arraybackslash}p{(\linewidth - 4\tabcolsep) * \real{0.1923}}
  >{\raggedright\arraybackslash}p{(\linewidth - 4\tabcolsep) * \real{0.4231}}
  >{\raggedright\arraybackslash}p{(\linewidth - 4\tabcolsep) * \real{0.3846}}@{}}
\toprule\noalign{}
\begin{minipage}[b]{\linewidth}\raggedright
Model
\end{minipage} & \begin{minipage}[b]{\linewidth}\raggedright
Description
\end{minipage} & \begin{minipage}[b]{\linewidth}\raggedright
Parameters
\end{minipage} \\
\midrule\noalign{}
\endhead
\bottomrule\noalign{}
\endlastfoot
\textbf{JC69} (Jukes-Cantor, 1969) & Most substitutions equally likely;
equal base frequencies & 0 free parameters \\
\textbf{K2P} (Kimura 2-parameter, 1980) & Transitions \(\neq\)
transversions; equal base frequencies & 1 (Ti/Tv ratio) \\
\textbf{HKY85} (Hasegawa-Kishino-Yano) & Ti \(\neq\) Tv + unequal base
frequencies & 4 \\
\textbf{GTR} (General Time Reversible) & Most 6 substitution types have
independent rates + unequal base frequencies & 9 \\
\textbf{GTR+\texorpdfstring{\ensuremath{\Gamma}}{Gamma}+I} & GTR + gamma-distributed rate variation across sites +
proportion of invariant sites & 11 \\
\end{longtable}
}

\textbf{Model selection} uses information criteria: \textbf{AIC}
\(= 2k - 2\ln L\) (where \(k\) = number of parameters) or \textbf{BIC}
\(= k\ln n - 2\ln L\) (where \(n\) = sample size). The model with the
lowest AIC or BIC provides the best balance of fit and complexity.
Tools: ModelFinder (implemented in IQ-TREE), jModelTest2.

\textbf{Software}: RAxML (Randomized Axelerated Maximum Likelihood),
IQ-TREE (fast ML with ultrafast bootstrap UFBoot2), FastTree
(approximate ML for very large alignments).

\subsubsection{Bayesian Inference}\label{bayesian-inference}

Bayesian phylogenetics uses Bayes' theorem to estimate the posterior
probability of trees:

\begin{equation}
P(\text{tree} \mid \text{data}) = \frac{P(\text{data} \mid \text{tree}) \cdot P(\text{tree})}{P(\text{data})}
\label{eq:phylogenetics_2}
\end{equation}

Since \(P(\text{data})\) -- the marginal likelihood -- is
computationally intractable (it requires summing over most possible
trees), Bayesian inference uses \textbf{Markov Chain Monte Carlo (MCMC)}
to sample trees from the posterior distribution. The Metropolis-Hastings
algorithm proposes modifications to the current tree, accepts changes
that improve the likelihood, and occasionally accepts changes that
decrease it (allowing escape from local optima).

\textbf{Output}: A set of posterior trees summarized as a majority-rule
consensus tree. Each node receives a \textbf{posterior probability (PP)}
-- the proportion of sampled trees containing that node. PP
\(\geq 0.95\) is considered statistically significant support.

\textbf{Software}: MrBayes (standard Bayesian phylogenetics),
BEAST/BEAST2 (specialized for time-calibrated chronograms, molecular
clock models, and phylogeography).

\paragraph{Bayesian phylogenetics in
practice}\label{bayesian-phylogenetics-in-practice}

A typical Bayesian phylogenetic workflow consists of the following
steps. \textbf{(1) Specify priors}: a tree prior (often a Yule
pure-birth or birth--death process for species trees, or a coalescent
prior for population-genetic data), a clock-rate prior (e.g., a
lognormal distribution centered at a calibrated rate), and
substitution-model parameter priors. \textbf{(2) Run two or more
independent MCMC chains} for \(10^7\)--\(10^9\) generations, sampling
the tree topology, branch lengths, and model parameters at intervals
(typically every 1,000 generations). \textbf{(3) Diagnose convergence}:
discard the first 10--25\% of samples as ``burn-in'', then verify that
independent chains have reached the same stationary distribution by
inspecting Effective Sample Size (ESS \textgreater{} 200 for most
parameters), Potential Scale Reduction Factor (PSRF \texorpdfstring{\ensuremath{\approx}}{~} 1.0), and visual
trace plots in Tracer. \textbf{(4) Summarize}: produce a maximum clade
credibility (MCC) tree, with each node labeled by its posterior
probability and (for chronograms) its 95\% highest posterior density
(HPD) age interval.

\textbf{Comparison with bootstrap support.} A widely cited rule of thumb
is that \textbf{posterior probability \texorpdfstring{\textgreater{}=}{>=} 0.95 is approximately equivalent
to bootstrap support \texorpdfstring{\textgreater{}=}{>=} 70\%}, but the two measures are not
interchangeable. Posterior probabilities tend to be higher than
bootstrap values for the same node --- sometimes considerably so when
the substitution model is misspecified. The interpretation differs as
well: a 95\% posterior probability means ``given the model and data,
there is a 95\% probability that this clade is correct''; a 70\%
bootstrap value means ``70\% of pseudoreplicate datasets recover this
clade''. Modern best practice is to report both measures, and to treat
single nodes with PP \textless{} 0.95 or bootstrap \textless{} 70 as
\textbf{unresolved}.

\paragraph{Strict clock vs.~relaxed clock: when to use
each}\label{strict-clock-vs.-relaxed-clock-when-to-use-each}

Bayesian phylogenetic dating uses one of two molecular-clock frameworks.
A \textbf{strict clock} assumes a single substitution rate across the
entire tree; this is appropriate for \textbf{closely related taxa}
(within-species or recent congeners), where the assumption is testable
and often defensible --- e.g., dating SARS-CoV-2 lineages within a
single pandemic, or dating Native American mtDNA lineages relative to
the Bering Strait crossing. A \textbf{relaxed clock} allows substitution
rates to vary across branches, drawn from a prior distribution
(uncorrelated lognormal, exponential, or autocorrelated random walk).
Relaxed clocks are appropriate for \textbf{divergent taxa with rate
heterogeneity} --- for example, mammalian phylogenies where rodents
evolve 5--10\texorpdfstring{\ensuremath{\times}}{x} faster than cetaceans, or insect phylogenies where
holometabolous lineages have generation-time-correlated rate shifts. The
choice of clock is \textbf{testable}: a likelihood-ratio test or
Bayes-factor comparison between strict and relaxed clock models,
evaluated under the same data and tree prior, indicates which is
supported. In practice, modern phylogeneticists routinely default to
relaxed clocks and primarily revert to strict clocks when relaxed-clock
parameter estimates indicate near-uniform rates.

\subsubsection{Bootstrap Analysis}\label{bootstrap-analysis}

\textbf{Non-parametric bootstrap} \citep{felsenstein1985} assesses
branch support by resampling:

\begin{enumerate}
\def\labelenumi{\arabic{enumi}.}
\tightlist
\item
  Create pseudoreplicates by randomly sampling alignment columns with
  replacement (same length as original alignment).
\item
  Reconstruct a tree from each pseudoreplicate (typically 100--1,000
  replicates).
\item
  The \textbf{bootstrap value} for each branch is the percentage of
  pseudoreplicate trees containing that branch.
\end{enumerate}

\textbf{Interpretation}: Bootstrap \(\geq 70\)\% is generally considered
moderate to strong support. Bootstrap values are conservative estimates
-- a true clade may receive bootstrap support below 70\% if the
supporting signal is concentrated in a few alignment positions that are
not always sampled.

\textbf{Bootstrap vs.~Bayesian PP}: Bayesian posterior probabilities
tend to be higher than bootstrap values for the same node. A common rule
of thumb treats bootstrap support of at least 70\% and posterior
probability of at least 0.95 as broadly comparable thresholds, not
interchangeable measures. However, PP can be inflated by model
misspecification --- when the substitution model fails to capture
features of the data, posterior probabilities become falsely confident.
\textbf{Both metrics should be reported} for each node, and discordance
(high PP, low bootstrap, or vice versa) is itself diagnostic of either
signal heterogeneity or model violation. Modern best practice typically
reports nodes as well-supported primarily when \textbf{both} bootstrap
\(\geq 70\%\) \textbf{and} PP \(\geq 0.95\) are achieved.

\textbf{MCMC convergence diagnostics}: A Bayesian phylogeny depends
strongly on the convergence of the MCMC chains that produced it; high
posterior probabilities are misleading when chains have not mixed across
tree space. Standard diagnostics include:

\begin{enumerate}
\def\labelenumi{\arabic{enumi}.}
\tightlist
\item
  \textbf{Effective sample size (ESS)}: For each parameter (tree length,
  individual branch lengths, model parameters), the number of
  independent samples after autocorrelation correction. ESS values
  \(\geq 200\) are typically required for confident inference.
\item
  \textbf{Multiple independent runs}: The same analysis run from
  different random starting points should converge to indistinguishable
  posterior distributions. Discrepancies signal that the chains are
  stuck in local optima.
\item
  \textbf{Trace plots}: Visual inspection of parameter values across
  MCMC iterations should show a stationary ``fuzzy caterpillar''
  pattern, not directional trends or sticking at single values.
\item
  \textbf{Burn-in removal}: The first 10--25\% of MCMC iterations are
  typically discarded as the chain finds the high-posterior region of
  parameter space.
\end{enumerate}

Tools like \textbf{Tracer} (BEAST suite) automate these checks, and
submitting trees without convergence verification is now considered
methodologically inadequate in published phylogenomics.

\subsubsection{Gene Tree versus Species Tree
Discordance}\label{gene-tree-versus-species-tree-discordance}

A central insight of modern phylogenomics is that \textbf{gene trees
frequently disagree with species trees}, even when the underlying
biological process is straightforward speciation. The two main
biological causes:

\paragraph{Incomplete lineage sorting
(ILS)}\label{incomplete-lineage-sorting-ils}

When two speciation events occur in rapid succession, ancestral
polymorphism may not have time to sort completely between the daughter
lineages. Some loci will retain the ancestral relationship while others
will track the species tree --- producing a mosaic of gene trees that
disagree with each other and with the true species history. ILS is most
severe when:

\begin{itemize}
\tightlist
\item
  \textbf{Effective population size of the ancestor is large} (\(4N_e\)
  about  time between speciations)
\item
  \textbf{Speciation events are temporally close} (the ``anomaly zone''
  of phylogenetics)
\item
  \textbf{Recombination is high} so that adjacent loci have independent
  histories
\end{itemize}

The classic case is the human--chimpanzee--gorilla trichotomy:
about 30\% of human nuclear loci show \textbf{gene trees
grouping human with gorilla, or chimp with gorilla}, rather than the
species-tree topology of (human, chimp) sister to gorilla. The
discordance reflects ILS during the rapid succession of
human--chimp--gorilla speciation events about 6--8 Ma.

\textbf{Coalescent-based methods} (ASTRAL, MP-EST, *BEAST) explicitly
model gene-tree discordance under the multi-species coalescent and infer
the \textbf{species tree} as the topology that maximizes likelihood
across many gene trees. Concatenation methods (which simply combine most
loci into one alignment) can be \textbf{statistically inconsistent}
under high ILS --- they sometimes converge to the wrong tree as more
data are added.

\paragraph{Horizontal gene transfer
(HGT)}\label{horizontal-gene-transfer-hgt}

In prokaryotes, gene-tree discordance often reflects \textbf{horizontal
acquisition} rather than incomplete sorting. A gene tree that places a
thermophilic bacterium within the archaea may indicate that the
bacterium acquired the gene from an archaeal source through HGT, rather
than the species tree placing it incorrectly. Detecting HGT requires:

\begin{enumerate}
\def\labelenumi{\arabic{enumi}.}
\tightlist
\item
  \textbf{Phylogenetic incongruence}: gene tree topology contradicts the
  consensus species tree.
\item
  \textbf{Atypical sequence features}: GC content, codon usage bias, or
  dinucleotide signature differing from the host genome.
\item
  \textbf{Patchy phylogenetic distribution}: gene present in distantly
  related organisms while absent from close relatives.
\end{enumerate}

Tools like \textbf{DistAL} and reconciliation methods (Notung,
Ranger-DTL) systematically identify HGT events by comparing gene trees
with reference species trees.

\begin{quote}
\textbf{Real-World Connection: COVID-19 Phylogenetics}

The SARS-CoV-2 pandemic demonstrated the power of real-time
phylogenetics. Within weeks of the first
\hyperref[gl:genome]{\textbf{genome}} sequence being shared (January
2020), phylogenetic analysis confirmed the virus's origin in the
betacoronavirus clade related to bat coronaviruses. Platforms like
Nextstrain (nextstrain.org) provided continuously updated phylogenetic
trees from millions of viral genomes deposited in GISAID. Bayesian
time-calibrated analyses (BEAST) dated the most recent common ancestor
to approximately November 2019. Phylogenetic tracking identified the
emergence and global spread of Variants of Concern (Alpha, Delta,
Omicron), revealing that each arose from a single geographic origin and
spread globally through human travel. Genomic epidemiology became a
standard public health tool, guiding variant surveillance,
\hyperref[gl:vaccine]{\textbf{vaccine}} updating, and travel policy. The
pandemic validated decades of investment in phylogenetic methodology and
demonstrated that evolutionary biology has direct, immediate public
health applications.
\end{quote}

\begin{quote}
\textbf{Concept Check 3:} Why is long-branch attraction a problem
specifically for maximum parsimony but less so for maximum likelihood?
(Hint: consider whether each method uses an explicit model of
substitution rates.)
\end{quote}

\begin{center}\rule{0.5\linewidth}{0.5pt}\end{center}

\subsection{Molecular Clocks}\label{molecular-clocks-1}

\subsubsection{The Molecular Clock
Hypothesis}\label{the-molecular-clock-hypothesis}

\citet{zuckerkandl1965} observed that the number of amino acid
differences between homologous proteins from different species is
roughly proportional to the time since those species diverged. This
\textbf{molecular clock} hypothesis -- that neutral substitutions
accumulate at an approximately constant rate -- allows divergence times
to be estimated from sequence data.

The theoretical basis comes from Kimura's
\hyperref[gl:neutral-theory]{\textbf{neutral theory}} (1968): the rate
of neutral substitution equals the neutral
\hyperref[gl:mutation]{\textbf{mutation}} rate, \texorpdfstring{\ensuremath{\mu}}{mu}, regardless of
population size. For strictly neutral mutations:

\begin{equation}
k = \mu
\label{eq:phylogenetics_3}
\end{equation}

where \(k\) is the substitution rate per site per generation.

\subsubsection{Clock Calculations}\label{clock-calculations}

\textbf{Jukes-Cantor correction}: Observed sequence divergence
underestimates true divergence because of multiple substitutions at the
same site (saturation). The JC69 correction accounts for this:

\begin{equation}
d_{JC} = -\frac{3}{4} \ln\left(1 - \frac{4}{3}p\right)
\label{eq:phylogenetics_4}
\end{equation}

where \(p\) is the observed proportion of differing sites.

\textbf{Divergence time estimation}:

\begin{equation}
t = \frac{d_{JC}}{2\mu}
\label{eq:phylogenetics_5}
\end{equation}

The factor of 2 accounts for substitutions accumulating independently in
both lineages since their divergence from a common ancestor.

\subsubsection{Worked Example: Estimating Divergence Time using
Jukes-Cantor}\label{worked-example-estimating-divergence-time-using-jukes-cantor}

\textbf{Problem:} Human and chimpanzee mitochondrial \emph{cytochrome b}
sequences differ by approximately 1.5\% (\(p = 0.015\)). Assume a
mitochondrial substitution rate of \(\mu \approx 2 \times 10^{-8}\)
substitutions per site per year. 1. Calculate the Jukes-Cantor distance
(\(d_{JC}\)). 2. Estimate the divergence time (\(t\)) between the two
lineages.

\textbf{Solution:}

\begin{enumerate}
\def\labelenumi{\arabic{enumi}.}
\item
  \textbf{Calculate the Jukes-Cantor distance:} Using the Jukes-Cantor
  correction formula to account for multiple substitutions at the same
  site: \begin{equation}
  d_{JC} = -\frac{3}{4} \ln\left(1 - \frac{4}{3}(0.015)\right) \approx 0.01511 \text{ substitutions/site}
  \label{eq:phylogenetics_6}
  \end{equation}
\item
  \textbf{Estimate the divergence time:} Apply the divergence time
  estimation formula, remembering the factor of 2 (since mutations
  accumulate along both lineages branching from the MRCA):
  \begin{equation}
  t = \frac{d_{JC}}{2\mu} = \frac{0.01511}{2 \times 2 \times 10^{-8}} \approx 3.78 \times 10^5 \text{ years}
  \label{eq:phylogenetics_7}
  \end{equation}
\end{enumerate}

\emph{Note:} This estimate of about 378,000 years for mtDNA is much
lower than the accepted species divergence time of about 6--7 Mya. This
occurs because mitochondrial \hyperref[gl:gene]{\textbf{gene}}s evolve
significantly faster than the genome-wide average, and substituting a
long-term nuclear calibration rate yields an artificially recent
divergence for rapidly evolving loci.

\subsubsection{Rate Calibration}\label{rate-calibration}

Molecular clocks must be calibrated against independent time estimates:

\begin{itemize}
\tightlist
\item
  \textbf{Fossil record}: Fossils provide \textbf{minimum age
  constraints} -- a fossil establishes that a lineage existed by that
  time, but the actual divergence must be older. The oldest known fossil
  of a clade provides a minimum bound on the clade's age.
\item
  \textbf{Biogeographic events}: The formation of the Isthmus of Panama
  (about 3 Mya), the separation of Australia from Antarctica (about 45
  Mya), or the formation of volcanic islands (known radiometric ages)
  provide calibration points.
\end{itemize}

\subsubsection{Strict vs.~Relaxed
Clocks}\label{strict-vs.-relaxed-clocks}

\textbf{Strict molecular clock}: Assumes a constant substitution rate
across most lineages. This assumption is often violated because
substitution rates vary with generation time (shorter-generation
organisms evolve faster per unit time), metabolic rate, DNA repair
efficiency, and population size (through the interplay of drift and
selection).

\textbf{Relaxed molecular clock}: Allows the rate to vary among lineages
while still estimating divergence times. Two main approaches: -
\textbf{Autocorrelated rates}: Closely related lineages have similar
rates (rates evolve gradually along the tree). Implemented in BEAST. -
\textbf{Uncorrelated rates}: Each branch draws its rate independently
from a statistical distribution (e.g., lognormal). Implemented in
BEAST2.

\subsubsection{Limitations}\label{limitations}

\begin{itemize}
\tightlist
\item
  \textbf{Rate heterogeneity}: Substitution rates can vary dramatically
  among lineages. Rodents evolve approximately 5--10 times faster than
  cetaceans at synonymous sites.
\item
  \textbf{Generation time effect}: Species with shorter generation times
  accumulate more mutations per unit time because there are more DNA
  replication events per year.
\item
  \textbf{Saturation}: At very large divergence times, multiple
  substitutions at the same site erase the phylogenetic signal. Slowly
  evolving genes (18S rRNA, about 1--2 \(\times 10^{-9}\) substitutions
  per site per year) are needed for deep divergences.
\end{itemize}

\subsubsection{Tip Dating and Bayesian Skyline
Plots}\label{tip-dating-and-bayesian-skyline-plots}

Modern molecular clock analyses go beyond simple point estimates of
divergence times:

\begin{itemize}
\tightlist
\item
  \textbf{Tip dating}: When sequences are sampled at different calendar
  dates (e.g., viral genomes collected over the course of an epidemic),
  the sampling dates themselves provide calibration points. BEAST
  software uses tip dates to estimate substitution rates directly from
  serially sampled data --- no fossil calibration needed. This approach
  was critical for dating the SARS-CoV-2 origin, HIV emergence, and
  influenza evolution.
\item
  \textbf{Bayesian Skyline Plots}: These reconstruct changes in
  effective population size (\(N_e\)) through time from a single
  population sample. By analyzing the distribution of coalescence times
  in a gene genealogy, skyline plots reveal historical population
  expansions, bottlenecks, and declines. Applied to human mtDNA, skyline
  analyses confirm the Out-of-Africa expansion approximately 60--70 kya.
\end{itemize}

\subsubsection{A timeline of major events in life's
phylogeny}\label{a-timeline-of-major-events-in-lifes-phylogeny}

Phylogenetic methods, when calibrated against the geological record,
reconstruct a deep-time timeline of life on Earth. The major events span
four billion years:

{\def\LTcaptype{none} % do not increment counter
\begin{longtable}[]{@{}
  >{\raggedright\arraybackslash}p{(\linewidth - 4\tabcolsep) * \real{0.1667}}
  >{\raggedright\arraybackslash}p{(\linewidth - 4\tabcolsep) * \real{0.6000}}
  >{\raggedright\arraybackslash}p{(\linewidth - 4\tabcolsep) * \real{0.2333}}@{}}
\toprule\noalign{}
\begin{minipage}[b]{\linewidth}\raggedright
Era
\end{minipage} & \begin{minipage}[b]{\linewidth}\raggedright
Approximate date
\end{minipage} & \begin{minipage}[b]{\linewidth}\raggedright
Event
\end{minipage} \\
\midrule\noalign{}
\endhead
\bottomrule\noalign{}
\endlastfoot
Hadean--Eoarchean & about 4.5--4.0 Ga & Earth forms
(about 4.54 Ga); oceans condense (about 4.4 Ga);
abiogenesis events (debated dates) \\
Eoarchean & about 3.8--3.5 Ga & Earliest evidence of life
(graphite isotope signatures from Greenland; stromatolites in
Australia) \\
LUCA & about 3.8 Ga & Last Comprehensive Common Ancestor ---
root of the tree of life. Inferred from comprehensive molecular
machinery (rRNA, ribosomal proteins, ATP synthase, key tRNAs). LUCA was
likely an anaerobic, hyperthermophilic prokaryote dependent on H\texorpdfstring{\textsubscript{2}}{2}/CO\texorpdfstring{\textsubscript{2}}{2}
chemistry near hydrothermal vents. \\
Mesoarchean--Paleoproterozoic & about 3.0--2.4 Ga &
Cyanobacteria evolve oxygenic photosynthesis; \textbf{Great Oxidation
Event} (about 2.4 Ga) transforms atmosphere from anoxic to
oxic \\
Paleoproterozoic & about 2.0 Ga & LECA (Last Eukaryotic Common
Ancestor) --- already possessed a nucleus, mitochondria (from
\texorpdfstring{\ensuremath{\alpha}}{alpha}-proteobacterial endosymbiont), endomembrane system, cytoskeleton, and
meiotic sexual reproduction. The mitochondrial endosymbiosis is dated to
about 2.0--1.5 Ga. \\
Mesoproterozoic & about 1.6--1.0 Ga & Plastid endosymbiosis
(cyanobacterium \texorpdfstring{\ensuremath{\to}}{->} primary plastid in archaeplastid ancestor);
diversification of eukaryotic supergroups \\
Neoproterozoic & about 1.0 Ga -- 540 Ma & Snowball Earth
events; first multicellular eukaryotes; \textbf{Ediacaran biota}
(about 575--540 Ma) --- soft-bodied multicellular organisms,
mostly enigmatic in affinity \\
Cambrian & about 540--485 Ma & \textbf{Cambrian explosion}:
most modern animal phyla appear in the fossil record within
about 25 million years. Evolution of mineralized skeletons,
predator-prey ecosystems, complex sensory systems \\
Ordovician & about 485--445 Ma & Diversification of marine
invertebrates; first land plants (bryophytes); \textbf{end-Ordovician
extinction} (about 445 Ma): about 85\% of marine
species lost \\
Silurian--Devonian & about 445--360 Ma & Vascular plants
colonize land; jawed fish (Gnathostomata) diversify; \textbf{first
tetrapods} (about 375 Ma; \emph{Tiktaalik}); \textbf{late
Devonian extinction} \\
Carboniferous & about 360--299 Ma & Coal-forest ecosystems;
first amniotes (reptilian-line ancestors); insects undergo first major
radiation \\
Permian & about 299--252 Ma & Synapsids (mammal ancestors)
dominate; \textbf{end-Permian extinction (Great Dying,
about 252 Ma)}: about 96\% of marine species lost ---
the largest extinction event in life's history \\
Triassic--Jurassic & about 252--145 Ma & Dinosaurs diversify;
first mammals (about 225 Ma) and first birds (Late Jurassic);
\textbf{end-Triassic extinction} (about 201 Ma) \\
Cretaceous & about 145--66 Ma & Flowering plants (angiosperms)
emerge and diversify; mammals remain small but radiate quietly;
\textbf{K-Pg extinction} (about 66 Ma): asteroid impact +
Deccan Traps volcanism eliminate non-avian dinosaurs \\
Cenozoic & about 66 Ma -- present & Mammalian and bird
radiations fill ecological niches; primates evolve; \emph{Homo sapiens}
appears (about 300 ka); \textbf{Holocene/Anthropocene
extinction} ongoing \\
\end{longtable}
}

The central insight from this timeline: life persists across cataclysms
but is reshuffled. The five mass extinctions did not erase the tree of
life but pruned it dramatically --- and \textbf{post-extinction
radiations} repeatedly produce explosive diversification as surviving
lineages occupy emptied ecological space (mammals after the K-Pg, modern
teleost fish after the Permian, etc.).

\subsubsection{Ancestral Sequence Reconstruction
(ASR)}\label{ancestral-sequence-reconstruction-asr}

One of the most striking applications of phylogenetic methods is
\textbf{ancestral sequence reconstruction} --- inferring the DNA or
protein sequences of extinct ancestors from their living descendants.
Given a tree and a substitution model, ASR uses \textbf{maximum
likelihood} or \textbf{Bayesian} methods to estimate the most probable
ancestral state at each internal node and at each site.

\paragraph{How ASR works}\label{how-asr-works}

For each site in a multiple sequence alignment, ASR computes:

\[P(\text{ancestor} = X \mid \text{tip data, tree, model}) \label{eq:unit_VI_phylogenetics_item_1}\]

using the \textbf{Felsenstein pruning algorithm} (which efficiently sums
over most possible ancestral states at internal nodes). Marginal
reconstructions estimate the most probable state at each node
independently; joint reconstructions estimate the most probable combined
ancestral states across the whole tree. Confidence is reported as the
posterior probability for each reconstructed residue --- high-confidence
sites have a single dominant residue (\textgreater95\% probability),
while low-confidence sites have several plausible alternatives.

\paragraph{``Lazarus proteins'': resurrecting extinct
enzymes}\label{lazarus-proteins-resurrecting-extinct-enzymes}

Once an ancestral sequence is reconstructed, it can be
\textbf{synthesized} as a real protein in the laboratory --- a ``Lazarus
protein'' brought back from extinction. This approach has revolutionized
our understanding of how protein function evolves:

\begin{itemize}
\tightlist
\item
  \textbf{Ancestral steroid receptors} (Joe Thornton's lab, since the
  early 2000s): The mineralocorticoid and glucocorticoid receptors split
  about 450 Mya. Resurrection of the ancestral receptor showed
  it was a glucocorticoid-like receptor; specificity for cortisol versus
  aldosterone evolved later through a small number of historical
  substitutions. The reconstruction made testable predictions about
  which substitutions enabled the functional shift --- confirmed
  experimentally by mutagenesis on the modern proteins.
\item
  \textbf{Ancestral elongation factors} (Akanuma et al.~2013, Gaucher et
  al.~2008): Ancestral bacterial EF-Tu proteins reconstructed from
  extant sequences proved highly thermostable, with melting temperatures
  of 60--73°C. The thermostability profile across nodes traces planetary
  cooling --- providing \textbf{biological thermometers} that measure
  the temperature history of Earth.
\item
  \textbf{Ancestral \texorpdfstring{\ensuremath{\beta}}{beta}-lactamases}: Reconstructed enzymes from before the
  antibiotic era reveal ``promiscuous'' ancestral activity --- the
  ancestral enzymes catalyzed multiple reactions efficiently. Modern
  descendants are more specialized but less catalytically robust. This
  pattern of ancestral promiscuity \texorpdfstring{\ensuremath{\to}}{->} descendant specialization is
  widespread in enzyme evolution.
\item
  \textbf{Ancestral rhodopsins} illuminate the spectral history of
  vision: by reconstructing visual pigments at deep nodes in the
  vertebrate tree, researchers have inferred the wavelengths of light
  that ancestral fish, amphibians, and reptiles could detect --- and
  found that color vision was repeatedly lost and re-evolved across many
  lineages.
\end{itemize}

\paragraph{Practical applications of
ASR}\label{practical-applications-of-asr}

ASR is more than a historical exercise --- it is now a
\textbf{protein-engineering platform}. Resurrected ancestral proteins
are routinely \textbf{more thermostable} and \textbf{more catalytically
robust} than their modern descendants --- a pattern repeatedly observed
across ancestral \texorpdfstring{\ensuremath{\beta}}{beta}-lactamases, EF-Tu elongation factors, alcohol
dehydrogenases, and rhodopsins. The hypothesised explanation is that
ancestral proteins evolved under broader environmental tolerances (early
Earth was warmer; ancestral metabolisms more flexible), and that
\textbf{modern descendants are specialised but brittle} --- selection
has pruned away robustness in favor of niche-specific performance. The
applied consequence: industrial enzyme engineering increasingly uses ASR
to \textbf{prospect for thermostable scaffolds}, then optimises the
resurrected protein for specific substrate specificity. Ancestral
\emph{Bacillus} \texorpdfstring{\ensuremath{\alpha}}{alpha}-amylases have been engineered into industrial
laundry-detergent enzymes that retain function at 60--80 °C; ancestral
\emph{Pseudomonas} esterases have been used as starting points for
plastic-degrading enzymes (PETase variants); and pharmaceutical
biotechnology uses ASR to design more thermostable vaccine antigens
(e.g.,

The 16S rRNA gene is ideal for deep phylogenetics because it is present
in all cellular organisms, contains both highly conserved regions (for
universal primer design) and variable regions (for discriminating taxa),
and is long enough (about 1,550 bp) for robust phylogenetic inference.

\subsubsection{Domain Bacteria}\label{domain-bacteria}

Bacteria are the most metabolically diverse domain of life. Major phyla
include:

\begin{itemize}
\tightlist
\item
  \textbf{Proteobacteria}: The largest and most diverse bacterial
  phylum.

  \begin{itemize}
  \tightlist
  \item
    Alpha-proteobacteria: includes \emph{Rickettsia} (obligate
    intracellular pathogen) and the SAR11 clade (most abundant organisms
    in the ocean). The ancestor of \textbf{mitochondria} was an
    alpha-proteobacterium.
  \item
    Gamma-proteobacteria: includes \emph{Escherichia coli},
    \emph{Pseudomonas}, \emph{Salmonella}, \emph{Vibrio cholerae}.
  \item
    Epsilon-proteobacteria: includes \emph{Helicobacter pylori} (causes
    gastric ulcers and gastric cancer).
  \end{itemize}
\item
  \textbf{Firmicutes}: Gram-positive, low-GC content. Includes
  \emph{Bacillus} (anthrax), \emph{Clostridium} (tetanus, botulism),
  \emph{Staphylococcus}, \emph{Lactobacillus} (probiotic fermenters),
  \emph{Streptococcus}.
\item
  \textbf{Bacteroidetes}: \hyperref[gl:dominant]{\textbf{Dominant}} gut
  anaerobes; critical for polysaccharide digestion. \emph{Bacteroides}
  species compose 30--40\% of human fecal bacteria.
\item
  \textbf{Actinobacteria}: Gram-positive, high-GC content.
  \emph{Streptomyces} produces over two-thirds of clinically used
  antibiotics. \emph{Mycobacterium tuberculosis} causes tuberculosis.
\item
  \textbf{Cyanobacteria}: Oxygenic photosynthesizers that produced
  Earth's oxygen atmosphere (Great Oxidation Event, about 2.4 Ga). The
  ancestor of \textbf{\hyperref[gl:chloroplast]{chloroplast}s} was a
  cyanobacterium.
\item
  \textbf{Spirochetes}: Helical morphology. \emph{Treponema pallidum}
  (syphilis), \emph{Borrelia burgdorferi} (Lyme disease).
\item
  \textbf{Chlamydiae}: Obligate intracellular
  \hyperref[gl:parasite]{\textbf{parasite}}s with a unique biphasic
  developmental cycle.
\end{itemize}

\subsubsection{Domain Archaea}\label{domain-archaea}

Archaea were originally found primarily in extreme environments, but
culture-independent methods (metagenomics) have revealed that they are
ubiquitous -- in soils, oceans, and the human gut.

\begin{itemize}
\tightlist
\item
  \textbf{Euryarchaeota}: The most diverse archaeal phylum.

  \begin{itemize}
  \tightlist
  \item
    Methanogens (\emph{Methanobacterium}, \emph{Methanosarcina}):
    produce methane as a metabolic byproduct; dominant in
    \hyperref[gl:anaerobic]{\textbf{anaerobic}} environments (wetlands,
    ruminant guts, landfills).
  \item
    Extreme halophiles (\emph{Halobacterium}): thrive in salt-saturated
    environments; use bacteriorhodopsin for light-driven proton pumping.
  \item
    Thermoacidophiles (\emph{Thermoplasma}): grow at
    \hyperref[gl:ph]{\textbf{pH}} 1--2 and temperatures up to 60 degrees
    C.
  \end{itemize}
\item
  \textbf{Crenarchaeota}: Many are hyperthermophiles. \emph{Sulfolobus}
  grows at 80 degrees C and pH 2--3. \emph{Thermoproteus} is an
  anaerobic sulfur-reducing hyperthermophile.
\item
  \textbf{TACK superphylum}: Thaumarchaeota (ammonia-oxidizing archaea
  -- major players in the nitrogen cycle), Aigarchaeota, Crenarchaeota,
  and Korarchaeota.
\item
  \textbf{Asgard archaea}: The most significant discovery in archaeal
  biology in decades. Named after Norse mythology:

  \begin{itemize}
  \tightlist
  \item
    \textbf{Lokiarchaeota}: Discovered via metagenome-assembled genomes
    from Loki's Castle hydrothermal vent, Arctic Mid-Ocean Ridge.
    Contains \textbf{eukaryotic signature proteins} including homologs
    of \hyperref[gl:actin]{\textbf{actin}}, tubulin-like GTPases, ESCRT
    membrane remodeling complex, and Arp2/3 complex regulators.
  \item
    \textbf{Thorarchaeota}, \textbf{Odinarchaeota},
    \textbf{Heimdallarchaeota}: Additional Asgard lineages with
    progressively more eukaryotic-like features.
  \item
    The Asgard archaea are the \textbf{closest known living relatives of
    eukaryotes}, supporting the ``two-domain'' tree (Eocyte hypothesis)
    in which eukaryotes arose from within Archaea rather than as a
    separate domain.
  \end{itemize}
\end{itemize}

\subsubsection{Domain Eukarya}\label{domain-eukarya}

The \textbf{Last Eukaryotic Common Ancestor (LECA)} already possessed a
nucleus, mitochondria, endomembrane system,
\hyperref[gl:cytoskeleton]{\textbf{cytoskeleton}}, and sexual
reproduction. Modern eukaryotic diversity is organized into several
supergroups:

\begin{itemize}
\tightlist
\item
  \textbf{Amorphea}:

  \begin{itemize}
  \tightlist
  \item
    \textbf{Amoebozoa}: Amoebas, slime molds (\emph{Dictyostelium}),
    \emph{Entamoeba histolytica} (amoebic dysentery).
  \item
    \textbf{Opisthokonta}: Fungi + Metazoa (animals) + choanoflagellates
    (closest unicellular relatives of animals). This grouping -- placing
    fungi as the sister group of animals rather than plants -- was one
    of the great surprises of molecular phylogenetics.
  \end{itemize}
\item
  \textbf{Diaphoretickes}:

  \begin{itemize}
  \tightlist
  \item
    \textbf{SAR clade}: Stramenopiles (diatoms, brown algae, oomycetes)
    + Alveolata (\emph{Plasmodium}, dinoflagellates, ciliates) +
    Rhizaria (foraminifera, radiolarians).
  \item
    \textbf{Archaeplastida}: Glaucophyta + Rhodophyta (red algae) +
    Chloroplastida (green algae + land plants). This clade acquired
    plastids through \textbf{primary
    \hyperref[gl:endosymbiosis]{\textbf{endosymbiosis}}} with a
    cyanobacterium.
  \end{itemize}
\item
  \textbf{Excavata} (debated monophyly): Deep-branching protists
  including \emph{Giardia} (intestinal parasite), \emph{Trichomonas}
  (STI pathogen), \emph{Trypanosoma} (sleeping sickness, Chagas
  disease), and \emph{Euglena}.
\end{itemize}

\subsubsection{Endosymbiotic Theory}\label{endosymbiotic-theory-1}

Lynn Margulis (1967, then writing as Lynn Sagan) proposed that
mitochondria and chloroplasts originated as free-living bacteria
engulfed by ancestral eukaryotic cells. Decades of molecular evidence
have confirmed this theory:

\textbf{Primary endosymbiosis}: - \textbf{Mitochondria}: Descended from
an alpha-proteobacterium. Evidence: mitochondria have their own circular
DNA, replicate by binary fission, have double membranes (inner membrane
= bacterial membrane), use bacterial-type
\hyperref[gl:ribosome]{\textbf{ribosome}}s (70S), and their gene
sequences cluster with alpha-proteobacteria in phylogenetic analyses. -
\textbf{Chloroplasts}: Descended from a cyanobacterium. Evidence:
chloroplasts have circular DNA, 70S ribosomes, double membranes, and
gene sequences grouping with cyanobacteria. Thylakoid membranes resemble
cyanobacterial internal membranes.

\textbf{Secondary endosymbiosis}: A eukaryotic cell engulfs a
photosynthetic eukaryote (which already has a primary plastid). This
explains the three or four membranes surrounding plastids in groups such
as euglenids (two extra membranes from the engulfed cell), diatoms,
dinoflagellates, and cryptophytes. Some secondary plastids even retain a
vestigial nucleus (nucleomorph) from the engulfed alga.

\subsubsection{Horizontal Gene Transfer and the Web of
Life}\label{horizontal-gene-transfer-and-the-web-of-life}

\textbf{Horizontal gene transfer (HGT)} -- the transfer of genetic
material between organisms outside of parent-to-offspring transmission
-- is widespread in prokaryotes and fundamentally challenges the tree
metaphor for representing life's history.

Mechanisms of HGT: - \textbf{Transformation}: Uptake of free
environmental DNA through competence proteins. - \textbf{Transduction}:
Bacteriophage-mediated DNA transfer. - \textbf{Conjugation}: Direct
cell-to-cell transfer via pili; the primary route for antibiotic
resistance \hyperref[gl:plasmid]{\textbf{plasmid}} transfer.

An estimated 20--30\% of genes in a typical prokaryotic genome have been
acquired by HGT. In \emph{E. coli} MG1655, \citet{lawrence1998}
identified approximately 755 genes (about 18\% of the genome) acquired
by HGT, identifiable by atypical GC content,
\hyperref[gl:codon]{\textbf{codon}} usage, and phylogenetic
incongruence.

The prevalence of HGT means that prokaryotic evolution is better
represented by a \textbf{network} or \textbf{web} of life rather than a
strictly bifurcating tree. For eukaryotes, vertical inheritance
dominates the nuclear genome, but organellar genomes and some nuclear
genes (acquired by endosymbiotic gene transfer, EGT) record the history
of endosymbiosis.

\begin{figure}[htbp]
\centering
\includegraphics[width=0.82\linewidth,height=4.2in,keepaspectratio]{../figures/mermaid_inline/inline_mermaid_0113_dbe0ca0e02a1.png}
\caption{The three domains of life with endosymbiotic organelle origins: mitochondria from alpha-Proteobacteria, chloroplasts from Cyanobacteria, and the Asgard-archaea root of Eukarya.}
\end{figure}

\begin{quote}
\textbf{Real-World Connection: Antibiotic Resistance and HGT}

The spread of antibiotic resistance among pathogenic bacteria is
fundamentally a problem of horizontal gene transfer. Resistance genes --
encoding \hyperref[gl:enzyme]{\textbf{enzyme}}s that degrade antibiotics
(beta-lactamases), modify antibiotics (aminoglycoside
acetyltransferases), or pump them out of the cell (efflux pumps) -- are
frequently carried on mobile genetic elements: plasmids, transposons,
and integrons. A single conjugation event can transfer multi-drug
resistance from a harmless gut bacterium to a pathogenic strain. The
NDM-1 (New Delhi metallo-beta-lactamase) gene, first identified in 2008,
has spread globally within years via plasmid-mediated HGT, conferring
resistance to nearly most beta-lactam antibiotics including carbapenems
-- the ``last resort'' drugs. Understanding HGT is essential for
combating the antibiotic resistance crisis.
\end{quote}

\begin{quote}
\textbf{Concept Check 5:} Explain why the discovery of Asgard archaea
(particularly Lokiarchaeota) supports a two-domain rather than
three-domain tree of life. What eukaryotic signature proteins were found
in Lokiarchaeota genomes?
\end{quote}

\begin{center}\rule{0.5\linewidth}{0.5pt}\end{center}

\subsection{Human Phylogeny}\label{human-phylogeny}

\subsubsection{Great Apes and Molecular
Phylogeny}\label{great-apes-and-molecular-phylogeny}

Humans belong to the superfamily \textbf{Hominoidea} (great apes and
lesser apes). Molecular phylogenetics has resolved the relationships
within this group with high confidence:

\begin{itemize}
\tightlist
\item
  \textbf{Gibbons} (family Hylobatidae): lesser apes; diverged from the
  great ape lineage approximately 18--20 Mya.
\item
  \textbf{Orangutans} (\emph{Pongo}): diverged approximately 14 Mya. Two
  species: Bornean (\emph{P. pygmaeus}) and Sumatran (\emph{P. abelii}),
  plus the recently described Tapanuli orangutan (\emph{P.
  tapanuliensis}).
\item
  \textbf{Gorillas} (\emph{Gorilla}): diverged approximately 9--10 Mya.
  Two species: western (\emph{G. gorilla}) and eastern (\emph{G.
  beringei}).
\item
  \textbf{Chimpanzees and bonobos} (\emph{Pan}): Our closest living
  relatives. Humans and chimpanzees share approximately 98.7\% DNA
  sequence identity across alignable regions. Divergence approximately
  6--7 Mya.
\item
  \textbf{Humans} (\emph{Homo sapiens}): The primarily surviving species
  of genus \emph{Homo}.
\end{itemize}

The molecular phylogeny contradicts earlier morphology-based
classifications that grouped orangutans, gorillas, and chimpanzees as
``great apes'' separate from humans. DNA data unambiguously place humans
within the great apes, as the sister taxon of \emph{Pan}.

\subsubsection{Fossil Hominins}\label{fossil-hominins}

The human fossil record is one of the richest and most intensively
studied:

{\def\LTcaptype{none} % do not increment counter
\begin{longtable}[]{@{}
  >{\raggedright\arraybackslash}p{(\linewidth - 4\tabcolsep) * \real{0.2000}}
  >{\raggedright\arraybackslash}p{(\linewidth - 4\tabcolsep) * \real{0.4571}}
  >{\raggedright\arraybackslash}p{(\linewidth - 4\tabcolsep) * \real{0.3429}}@{}}
\toprule\noalign{}
\begin{minipage}[b]{\linewidth}\raggedright
Species
\end{minipage} & \begin{minipage}[b]{\linewidth}\raggedright
Approximate dates
\end{minipage} & \begin{minipage}[b]{\linewidth}\raggedright
Key features
\end{minipage} \\
\midrule\noalign{}
\endhead
\bottomrule\noalign{}
\endlastfoot
\emph{Sahelanthropus tchadensis} & about 7 Mya & Oldest possible
hominin; anterior foramen magnum suggests bipedality \\
\emph{Ardipithecus ramidus} & about 4.4 Mya & Partial skeleton
(``Ardi''); woodland bipedality with grasping feet \\
\emph{Australopithecus afarensis} & about 3.9--2.9 Mya & ``Lucy'' (AL
288-1); bipedal with small brain (about 430 cc); Laetoli footprints \\
\emph{Australopithecus africanus} & about 3.3--2.1 Mya & South African
``gracile'' australopith; Taung child \\
\emph{Homo habilis} & about 2.4--1.4 Mya & First stone tools (Oldowan
industry); brain about 600--700 cc \\
\emph{Homo ergaster/erectus} & about 1.8 Mya--110 kya & First hominin to
leave Africa; Acheulean hand axes; brain about 900--1100 cc; controlled
use of fire \\
\emph{Homo heidelbergensis} & about 700--200 kya & Common ancestor of
Neanderthals and modern humans; brain about 1200 cc \\
\emph{Homo neanderthalensis} & about 400--40 kya & Adapted to cold
European/Western Asian environments; brain about 1500 cc; burial of
dead, use of pigments \\
\emph{Homo sapiens} & about 300 kya--present & Anatomically modern;
fully symbolic culture; originated in Africa \\
\end{longtable}
}

\subsubsection{Out of Africa}\label{out-of-africa}

The \textbf{Out of Africa} hypothesis, supported by both genetic and
fossil evidence, proposes that anatomically modern \emph{Homo sapiens}
evolved in Africa and subsequently dispersed to populate the rest of the
world, largely replacing (but partially admixing with) archaic hominin
populations:

\begin{itemize}
\tightlist
\item
  \textbf{Mitochondrial Eve}: Coalescence analysis of mtDNA places the
  most recent common ancestor of living human mitochondrial lineages in
  Africa approximately 150--200 kya.
\item
  \textbf{Y-chromosomal Adam}: The most recent common ancestor of living
  human Y \hyperref[gl:chromosome]{\textbf{chromosome}}s dates to
  approximately 200--340 kya.
\item
  \textbf{Multiple dispersal events}: Genetic evidence suggests at least
  two major dispersal waves -- an early dispersal to Australia/Oceania
  (about 65 kya) and a later dispersal populating Europe and Asia (about
  45--50 kya).
\end{itemize}

\subsubsection{Archaic Introgression}\label{archaic-introgression}

One of the most remarkable discoveries of the genomic era is that modern
humans carry DNA from extinct hominin species, acquired through
interbreeding:

\begin{itemize}
\tightlist
\item
  \textbf{Neanderthal admixture}: Non-African modern human populations
  carry approximately 1--4\% Neanderthal-derived DNA. Interbreeding
  occurred approximately 50,000--60,000 years ago in the Middle East.
  Neanderthal \hyperref[gl:allele]{\textbf{allele}}s contribute to
  traits including skin pigmentation, hair texture, immune function, and
  susceptibility to certain diseases.
\item
  \textbf{Denisovan admixture}: Melanesian and Australian Aboriginal
  populations carry approximately 4--6\% Denisovan-derived DNA. The
  Denisovans are known primarily from a few fragmentary fossils from
  Denisova Cave, Siberia, and a mandible from Tibet, but their genomic
  legacy is substantial.
\item
  \textbf{Adaptive introgression}: Some introgressed alleles have been
  maintained by \hyperref[gl:natural-selection]{\textbf{natural
  selection}} because they conferred adaptive advantages:

  \begin{itemize}
  \tightlist
  \item
    \textbf{EPAS1}: A Denisovan-derived allele of this
    \hyperref[gl:transcription]{\textbf{transcription}} factor helps
    Tibetans adapt to high-altitude hypoxia by blunting the
    erythropoietic response to low oxygen, preventing polycythemia.
  \item
    \textbf{HLA alleles}: Neanderthal and Denisovan-derived HLA variants
    increased immune diversity in modern human populations expanding
    into new pathogen environments.
  \item
    \textbf{BNC2}: A Neanderthal-derived skin pigmentation gene at high
    frequency in European populations.
  \end{itemize}
\item
  \textbf{Ghost lineage introgression}: Analysis of West African genomes
  (Yoruba, Mende) reveals approximately 2--19\% of their ancestry from
  an unknown archaic hominin that diverged from the modern human lineage
  approximately 625 kya -- a ``ghost'' population for which no fossils
  have been identified.
\end{itemize}

\begin{quote}
\textbf{Concept Check 6:} Explain why ``Mitochondrial Eve'' does not
mean there was a single woman alive at that time. What does the
coalescence of mitochondrial lineages actually tell us?
\end{quote}

\begin{quote}
\textbf{Concept Check 7:} A phylogenomic analysis of 1,000 mammalian
loci finds that 70\% of gene trees support topology A, 20\% support
topology B, and 10\% support topology C, where A places mouse with rat
(sister taxa), B groups mouse with primates, and C groups rat with
primates. Is this pattern consistent with incomplete lineage sorting,
horizontal gene transfer, or model misspecification? Which
species-tree-inference method would be most appropriate here?
\end{quote}

\begin{quote}
\textbf{Concept Check 8:} A researcher reconstructs an ancestral steroid
receptor by maximum likelihood from 50 modern sequences. The
reconstructed ancestral residue at one critical position has posterior
probability 0.55 (versus 0.40 for the second-most-likely residue). What
does this low confidence imply about the resurrected protein, and how
should it influence the experimental design? (Consider: should both
candidate ancestors be synthesized?)
\end{quote}

\begin{quote}
\textbf{Concept Check 9:} A SARS-CoV-2 lineage emerges in late 2021 with
a novel constellation of mutations not directly descended from any of
the previously detected variants. The Bayesian time-calibrated tree
dates the most recent common ancestor of this lineage to early 2020.
What scenarios could explain this temporal mismatch --- sustained
cryptic transmission, evolution in an immunocompromised host, or
zoonotic re-emergence? What additional data would discriminate among
these hypotheses?
\end{quote}

\begin{quote}
\textbf{Clinical Connection --- Genomic epidemiology.} Phylogenetic
trees built from pathogen genomes are now routine public-health
infrastructure. For HIV, each new clinical isolate is placed on a
national phylogeny to identify transmission clusters and guide focused
intervention. For hospital-acquired infections (MRSA, \emph{C.
difficile}), whole-genome sequencing plus phylogenetic analysis
distinguishes ward-level outbreaks from sporadic community acquisitions
--- a 2013 Cambridge study traced 15 distinct \emph{C. difficile}
lineages in a single hospital, revealing that about 35 \% of cases were
nosocomial. For SARS-CoV-2, the real-time GISAID phylogeny guided
variant-specific vaccine updates.
\end{quote}

\begin{quote}
\textbf{Clinical Connection --- Identifying zoonotic spillover sources.}
When Ebola, SARS-CoV, MERS, or H5N1 emerge, the first question is: where
did it come from? Phylogenetic analysis places the human isolate within
the tree of known animal reservoirs, identifying the ancestor branch.
SARS-CoV-2's closest known relatives are horseshoe-bat
(\emph{Rhinolophus}) coronaviruses from Yunnan, China, sharing
about 96 \% sequence identity --- a divergence time estimated
by molecular clock at 40--70 years. These analyses direct surveillance
(which wild populations to screen) and inform policy on wildlife-market
regulation.
\end{quote}

\begin{center}\rule{0.5\linewidth}{0.5pt}\end{center}

\subsection{Computational Bridge}\label{computational-bridge-20}

Strict-clock time from fractional divergence (per site) is:

\begin{Shaded}
\begin{Highlighting}[]
\ImportTok{from}\NormalTok{ biology.evolution }\ImportTok{import}\NormalTok{ molecular\_clock\_divergence\_time}

\NormalTok{years }\OperatorTok{=}\NormalTok{ molecular\_clock\_divergence\_time(}\FloatTok{1e{-}9}\NormalTok{, }\FloatTok{0.02}\NormalTok{)}
\BuiltInTok{print}\NormalTok{(years)}
\end{Highlighting}
\end{Shaded}

\begin{quote}
\textbf{Clinical / systems note:} Pathogen phylogenomics during
outbreaks uses the same clock logic with externally calibrated rates
(serial sampling) rather than fossils.
\end{quote}

\begin{center}\rule{0.5\linewidth}{0.5pt}\end{center}

\subsubsection{Long-Read Sequencing and Deep-Divergence
Phylogenetics}\label{long-read-sequencing-and-deep-divergence-phylogenetics}

The accuracy of any phylogenetic tree is bounded by the fidelity of its
underlying sequence data --- and for billion-year divergences the
short-read Illumina platform (150--300 bp fragments) has been the
bottleneck. Short reads mis-assemble through repetitive regions,
collapse paralogs, and cannot span structural variants.
\textbf{Long-read platforms} --- \textbf{PacBio HiFi} (HIFI: 15--25 kb
reads at 99.9 \% accuracy via circular consensus sequencing) and
\textbf{Oxford Nanopore Ultra-long reads} (100 kb -- 2 Mb reads at
97--99 \% raw accuracy, 99.9 \% after polishing) --- have transformed
the field since 2020.

Three phylogenetic problems that long reads solved: \textbf{(1) The
animal tree of life.} Sponges vs.~ctenophores as the earliest-branching
animal clade --- a 10-year controversy --- was re-adjudicated with 1178
long-read-assembled orthologs (Redmond \& McLysaght, \emph{Nat. Commun.}
2021); the analysis supported \textbf{ctenophores-first}, implying
convergent evolution of the nervous system (or its loss in sponges).
\textbf{(2) Polyploid plant genomes.} The hexaploid wheat genome (5 \texorpdfstring{\ensuremath{\times}}{x}
the human genome size) was resolvable primarily with HiFi + Hi-C
scaffolding (\emph{Nature} 2023), enabling the first accurate
subgenome-level phylogeny of grass evolution. \textbf{(3) Ancient DNA
phylogenetics.} PacBio HiFi on Pleistocene-age samples now reaches
\textless{} 50 000-year mammoths and permafrost wolves; Nanopore
ultra-long sequencing on sediment-trapped eDNA (eDNA-seq) reconstructs
environments without individual organisms --- extending phylogenetics to
ecosystems rather than single species.

\textbf{Practical workflow for a 2025 phylogenomic study}: HiFi
sequencing to 30\texorpdfstring{\ensuremath{\times}}{x} coverage (about \$3000 per vertebrate-scale
genome, 2024 prices) \texorpdfstring{\ensuremath{\to}}{->} assembly with hifiasm \texorpdfstring{\ensuremath{\to}}{->} BUSCO completeness check
\texorpdfstring{\ensuremath{\to}}{->} orthologue identification with OrthoFinder \texorpdfstring{\ensuremath{\to}}{->} multiple sequence
alignment with MAFFT-einsi \texorpdfstring{\ensuremath{\to}}{->} species-tree inference with ASTRAL under
the multi-species coalescent \texorpdfstring{\ensuremath{\to}}{->} dating with MCMCTree against
molecular-clock calibrations from the fossil record. The bottleneck has
shifted from sequence generation to \textbf{taxonomic sampling} --- we
have the tools, we need the specimens.

\begin{center}\rule{0.5\linewidth}{0.5pt}\end{center}

\subsection{Current Evidence and Frontier
Biology}\label{current-evidence-and-frontier-biology-23}

For \textbf{Phylogenetics and the Tree of Life}, frontier biology
belongs inside the evidence logic of the chapter. Evolutionary claims
are strongest when they combine mechanism, comparative evidence,
population process, and explicit uncertainty. The core reading question
is this: phylogenetic confidence depends on sampling, model choice,
homology, conflict among loci, and calibration.

\begin{itemize}
\tightlist
\item
  \textbf{What to verify:} identify the observation, model, assay, or
  dataset that would make the claim stronger or weaker.
\item
  \textbf{What to qualify:} state the scale, organism, cell type,
  environmental condition, or population where the claim is expected to
  hold.
\item
  \textbf{What to compare:} test at least one alternative explanation,
  baseline, or null model before treating the pattern as causal.
\item
  \textbf{What to cite:} distinguish primary evidence, review synthesis,
  public dataset, and institutional guidance; for recent or numeric
  claims, prefer the source closest to the measurement and state what
  has changed since it was published.
\end{itemize}

Distinguish adaptation from drift, phylogenetic signal from convergence,
and historical explanation from a testable prediction about present-day
data.

\textbf{Source practice:} For evolutionary claims, prefer evidence that
compares alternatives such as selection, drift, gene flow, constraint,
convergence, and shared ancestry.

\subsection{Summary}\label{summary-23}

\begin{itemize}
\tightlist
\item
  \textbf{Phylogenetics} reconstructs evolutionary relationships from
  molecular and morphological data, enabling classification, divergence
  dating, biogeographic reconstruction, and applications in medicine,
  conservation, and drug discovery.
\item
  \textbf{Tree terminology}: Root (MRCA), nodes (divergence events),
  branches (lineages), tips (observed taxa), clades (monophyletic
  groups). Paraphyletic and polyphyletic groups are invalid in modern
  systematics. Synapomorphies define clades; symplesiomorphies and
  homoplasies do not.
\item
  \textbf{Reconstruction methods}: UPGMA (assumes clock; not
  recommended), Neighbor-Joining (fast, no clock), Maximum Parsimony
  (fewest changes; vulnerable to long-branch attraction), Maximum
  Likelihood (model-based; GTR+\texorpdfstring{\ensuremath{\Gamma}}{Gamma}+I; RAxML, IQ-TREE), Bayesian Inference
  (MCMC sampling; posterior probabilities; MrBayes, BEAST).
\item
  \textbf{Molecular clocks}:
  \(d_{JC} = -\frac{3}{4}\ln(1 - \frac{4}{3}p)\); \(t = d/(2\mu)\);
  calibrated from fossils and biogeographic events; strict vs.~relaxed
  clocks accommodate rate variation.
\item
  \textbf{Three domains}: Bacteria (metabolically diverse; source of
  mitochondria and chloroplasts), Archaea (Asgard archaea closest to
  eukaryotes), Eukarya (Amorphea, Diaphoretickes, Excavata).
  Endosymbiosis produced mitochondria (alpha-proteobacterium) and
  chloroplasts (cyanobacterium).
\item
  \textbf{HGT} is pervasive in prokaryotes (20--30\% of genes); the tree
  of life is more accurately a web for prokaryotes. HGT drives
  antibiotic resistance spread.
\item
  \textbf{Human phylogeny}: Diverged from chimpanzees about 6--7 Mya;
  fossil record spans \emph{Sahelanthropus} to \emph{Homo sapiens};
  Out-of-Africa dispersal; archaic introgression from Neanderthals
  (1--4\% in non-Africans) and Denisovans (4--6\% in Melanesians);
  adaptive introgression (EPAS1, HLA alleles).
\item
  \textbf{Connections:} See \cref{sec:unit_VII_bacteria_archaea_viruses}
  for HGT vs.~vertical signal, \cref{sec:unit_IV_mutations_and_genomics}
  for variant interpretation, and
  \cref{sec:unit_X_biomes_and_conservation} for phylogenetic diversity
  in triage.
\end{itemize}

\begin{center}\rule{0.5\linewidth}{0.5pt}\end{center}

\subsection{Key Terms}\label{key-terms-23}

{\def\LTcaptype{none} % do not increment counter
\begin{longtable}[]{@{}
  >{\raggedright\arraybackslash}p{(\linewidth - 2\tabcolsep) * \real{0.2857}}
  >{\raggedright\arraybackslash}p{(\linewidth - 2\tabcolsep) * \real{0.7143}}@{}}
\toprule\noalign{}
\begin{minipage}[b]{\linewidth}\raggedright
Term
\end{minipage} & \begin{minipage}[b]{\linewidth}\raggedright
Definition
\end{minipage} \\
\midrule\noalign{}
\endhead
\bottomrule\noalign{}
\endlastfoot
\textbf{Phylogenetics} & The study of evolutionary relationships among
organisms, inferred from heritable characters \\
\textbf{Clade} & A monophyletic group: an ancestor and its complete
descendant lineage \\
\textbf{Synapomorphy} & A shared derived character that defines a
clade \\
\textbf{Symplesiomorphy} & A shared ancestral character inherited from a
more distant ancestor; uninformative for grouping \\
\textbf{Homoplasy} & Similarity not due to shared ancestry; arises from
convergent evolution or reversal \\
\textbf{Outgroup} & A taxon known to fall outside the ingroup; used to
root phylogenetic trees \\
\textbf{Maximum parsimony} & Phylogenetic method selecting the tree
requiring the fewest total character-state changes \\
\textbf{Maximum likelihood} & Method selecting the tree maximizing
\(P(\text{data} \mid \text{tree, model})\) \\
\textbf{Bayesian inference} & Method estimating the posterior
probability of trees using MCMC sampling \\
\textbf{Posterior probability} & \(P(\text{tree} \mid \text{data})\);
Bayesian branch support; \(\geq 0.95\) is significant \\
\textbf{Bootstrap} & Resampling technique for assessing branch support;
\(\geq 70\)\% indicates moderate-strong support \\
\textbf{Molecular clock} & Hypothesis that neutral substitutions
accumulate at approximately constant rate, enabling divergence time
estimation \\
\textbf{Substitution model} & Mathematical description of nucleotide or
amino acid replacement probabilities (e.g., JC69, GTR+\texorpdfstring{\ensuremath{\Gamma}}{Gamma}+I) \\
\textbf{Horizontal gene transfer} & Transfer of genetic material between
organisms outside parent-offspring inheritance \\
\textbf{Endosymbiosis} & Origin of organelles (mitochondria,
chloroplasts) from free-living bacteria engulfed by ancestral
eukaryotes \\
\textbf{Asgard archaea} & Archaeal superphylum (Lokiarchaeota,
Thorarchaeota, etc.) that are the closest known relatives of
eukaryotes \\
\textbf{Phylogenetic diversity (PD)} & Total branch length of a
phylogeny containing a set of species; used in conservation
prioritization \\
\textbf{LECA} & Last Eukaryotic Common Ancestor; possessed nucleus,
mitochondria, endomembrane system \\
\textbf{Introgression} & Incorporation of alleles from one species into
another through hybridization and backcrossing \\
\end{longtable}
}

\begin{center}\rule{0.5\linewidth}{0.5pt}\end{center}

\subsection{Review Questions}\label{review-questions-23}

\begin{enumerate}
\def\labelenumi{\arabic{enumi}.}
\item
  Define monophyletic, paraphyletic, and polyphyletic groups. For each,
  provide an example from vertebrate taxonomy and explain why primarily
  monophyletic groups are valid in modern systematics.
\item
  Reconstruct a parsimony analysis for four taxa (A, B, C, D) given the
  following binary characters: Character 1: A=1, B=1, C=0, D=0;
  Character 2: A=1, B=0, C=1, D=0; Character 3: A=1, B=1, C=1, D=0.
  Which tree requires the fewest steps? Is there homoplasy?
\item
  A researcher obtains a 600 bp alignment of the COI gene from two
  beetle species, finding 42 differences. (a) Calculate the observed
  divergence (\(p\)). (b) Apply the Jukes-Cantor correction. (c) Using a
  COI substitution rate of \(2.3 \times 10^{-2}\) per site per million
  years, estimate the divergence time.
\item
  Compare maximum likelihood and Bayesian inference as phylogenetic
  methods. What are the advantages and disadvantages of each? When would
  you choose one over the other?
\item
  Explain why the GTR+\texorpdfstring{\ensuremath{\Gamma}}{Gamma}+I substitution model is preferred over simpler
  models (JC69, K2P) for most real phylogenetic datasets. What
  biological realities does each additional parameter capture?
\item
  Describe the evidence supporting the endosymbiotic origin of
  mitochondria from alpha-proteobacteria. List at least four independent
  lines of evidence.
\item
  Explain how Asgard archaea challenge the three-domain tree of life.
  What specific eukaryotic features have been found in Asgard genomes,
  and what does this imply about the origin of eukaryotes?
\item
  The SARS-CoV-2 pandemic relied heavily on phylogenetic analysis.
  Describe how Bayesian time-calibrated phylogenetics (BEAST) was used
  to (a) date the origin of the pandemic, (b) track the emergence of
  Variants of Concern, and (c) reconstruct transmission chains. How is a
  molecular clock calibrated for a virus with no fossil record?
\item
  Modern non-African humans carry approximately 1--4\% Neanderthal DNA.
  Explain the process by which this DNA was acquired and describe one
  example of adaptive introgression. Why is Neanderthal DNA being
  gradually purged from functionally important genomic regions?
\item
  A conservation biologist must choose between protecting Species X (one
  of 50 species in a diverse, recently radiated clade) and Species Y
  (the sole surviving member of an ancient lineage). Using the concept
  of phylogenetic diversity, argue for which species should receive
  priority and explain the logic behind the EDGE conservation framework.
\item
  Why does uncorrected \(p\)-distance \textbf{underestimate} deep
  divergences, and when is Jukes--Cantor inadequate compared with GTR+\texorpdfstring{\ensuremath{\Gamma}}{Gamma}?
\item
  Name one scenario where a \textbf{network} rather than a bifurcating
  tree better represents genome evolution. \#\# Further Reading and
  Source Notes
\end{enumerate}

\begin{itemize}
\tightlist
\item
  Saitou \& Nei (1987). The neighbor-joining method: a new method for
  reconstructing phylogenetic trees. \emph{Molecular Biology and
  Evolution}, 4.
\item
  Felsenstein (1985). Confidence limits on phylogenies: An approach
  using the bootstrap. \emph{Evolution}, 39.
\item
  Zuckerkandl \& Pauling (1965). Evolutionary divergence and convergence
  in proteins. \emph{Evolving Genes and Proteins}.
\item
  Lawrence \& Ochman (1998). Molecular archaeology of the \{Escherichia
  coli\} genome. \emph{Proceedings of the National Academy of Sciences},
  95.
\item
  Woese \& Fox (1977). Phylogenetic structure of the prokaryotic domain:
  The primary kingdoms. \emph{Proceedings of the National Academy of
  Sciences}, 74.
\end{itemize}

\begin{center}\rule{0.5\linewidth}{0.5pt}\end{center}

\subsubsection{Companion Source
Module}\label{companion-source-module-23}

\textbf{Phylogenetics and the Tree of Life} should leave a reproducible
trail from a biological claim to the code, figure, diagram, or
paper-based activity that can test it. Use the surfaces below to inspect
the chapter's assumptions, rerun the relevant model, or compare the
manuscript explanation with companion labs and figures.

{\def\LTcaptype{none} % do not increment counter
\begin{longtable}[]{@{}
  >{\raggedright\arraybackslash}p{(\linewidth - 2\tabcolsep) * \real{0.5000}}
  >{\raggedright\arraybackslash}p{(\linewidth - 2\tabcolsep) * \real{0.5000}}@{}}
\toprule\noalign{}
\begin{minipage}[b]{\linewidth}\raggedright
Surface
\end{minipage} & \begin{minipage}[b]{\linewidth}\raggedright
Use it for
\end{minipage} \\
\midrule\noalign{}
\endhead
\bottomrule\noalign{}
\endlastfoot
\breaktt{src/biology/evolution/evolution.py}
(\breaktt{molecular\_clock\_divergence\_time}) & Translate genetic
distance and rate assumptions into divergence-time estimates. \\
\breaktt{src/biology/genetics/genetics.py} (\texttt{hamming\_distance},
\breaktt{jukes\_cantor\_distance}) & Compare raw and corrected sequence
distances. \\
\breaktt{src/mermaid/biology\_diagrams.py}
(\breaktt{phylogenetic\_tree\_diagram}) & Keep topology, branch length,
and interpretation visually distinct. \\
\end{longtable}
}

\textbf{Reproducibility check:} state alignment quality, homology
assumption, substitution model, sampling, and calibration before
treating a tree as history. \textbf{Cross-reference:} use
\cref{sec:unit_VI_genetic_drift_and_speciation} and
\cref{sec:unit_VII_bacteria_archaea_viruses}.

\newpage

\section*{Unit VII --- Microbiology:
Introduction}\label{unit-vii-microbiology-introduction}
\addcontentsline{toc}{section}{Unit VII --- Microbiology: Introduction}

\subsection*{Why This Unit Matters}\label{why-this-unit-matters-7}
\addcontentsline{toc}{subsection}{Why This Unit Matters}

The living world is microbial. In mass, in metabolic diversity, and in
evolutionary history, microorganisms dominate life on Earth. The total
number of prokaryotic cells on Earth (about 10\texorpdfstring{\textsuperscript{3}}{3}\texorpdfstring{\textsuperscript{0}}{0}) exceeds the
number of stars in the observable universe by a factor of 100,000. The
deepest branches of the tree of life are most microbial. Photosynthesis
--- the ultimate source of energy for most ecosystems --- was invented
by cyanobacteria more than 2.7 billion years ago. Nitrogen fixation,
without which terrestrial agriculture would collapse, is performed
exclusively by bacteria and archaea. The atmosphere itself was
transformed by microbial metabolism long before any animal breathed it.

Medical microbiology became possible primarily after Robert Koch's
postulates (1877--1884) established the germ theory of disease --- the
revolutionary idea that specific microorganisms cause specific
infections. Alexander Fleming's observation in 1928 that
\emph{Penicillium} mould killed staphylococcal colonies initiated the
antibiotic era, saving an estimated 200 million lives over the following
century. But pathogens evolve. The emergence of multidrug-resistant
organisms (MRSA, XDR-TB, carbapenem-resistant \emph{Enterobacteriaceae})
represents one of the most urgent public health crises of the 21st
century: an estimated 1.27 million deaths were attributable to
antimicrobial resistance in 2019 (Murray et al., 2022, \emph{The
Lancet}).

This unit combines cell biology (structure of bacteria, archaea, and
viruses), quantitative microbiology (growth curves, MIC calculations,
viral replication kinetics), and ecology (the microbiome, biogeochemical
cycles, infectious disease dynamics). You will apply the S-I-R
epidemiological model and calculate R\texorpdfstring{\textsubscript{0}}{0}; model bacterial growth and
antibiotic kill curves; and examine the microbiome as a complex
ecological community whose disruption underlies diseases from \emph{C.
difficile} colitis to metabolic syndrome.

\begin{center}\rule{0.5\linewidth}{0.5pt}\end{center}

\subsection*{Landmark Discoveries}\label{landmark-discoveries-7}
\addcontentsline{toc}{subsection}{Landmark Discoveries}

{\def\LTcaptype{none} % do not increment counter
\begin{longtable}[]{@{}
  >{\raggedright\arraybackslash}p{(\linewidth - 8\tabcolsep) * \real{0.2407}}
  >{\raggedright\arraybackslash}p{(\linewidth - 8\tabcolsep) * \real{0.0741}}
  >{\raggedright\arraybackslash}p{(\linewidth - 8\tabcolsep) * \real{0.2963}}
  >{\raggedright\arraybackslash}p{(\linewidth - 8\tabcolsep) * \real{0.1667}}
  >{\raggedright\arraybackslash}p{(\linewidth - 8\tabcolsep) * \real{0.2222}}@{}}
\toprule\noalign{}
\begin{minipage}[b]{\linewidth}\raggedright
Discoverer(s)
\end{minipage} & \begin{minipage}[b]{\linewidth}\raggedright
Year
\end{minipage} & \begin{minipage}[b]{\linewidth}\raggedright
Journal / Source
\end{minipage} & \begin{minipage}[b]{\linewidth}\raggedright
Discovery
\end{minipage} & \begin{minipage}[b]{\linewidth}\raggedright
Significance
\end{minipage} \\
\midrule\noalign{}
\endhead
\bottomrule\noalign{}
\endlastfoot
Antoni van Leeuwenhoek & 1676 & \emph{Philos. Trans. R. Soc.} & First
observation of bacteria ("animalcules") & Founded microbiology; showed
microorganisms are everywhere \\
Louis Pasteur & 1859--1861 & \emph{Ann. Sci. Nat.} & Disproof of
spontaneous generation & Established germ theory; revolutionised
medicine and food safety \\
Robert Koch & 1876--1884 & \emph{Investigations into the Etiology of
Traumatic Infective Diseases} & Koch's postulates; anthrax, TB, cholera
& Rigorous criteria for linking pathogen to disease \\
Alexander Fleming & 1928 & \emph{Brit. J. Exp. Pathol.} & Penicillin
discovery & Launched the antibiotic era; Nobel Prize 1945 \\
Carl Woese \& George Fox & 1977 & \emph{Proc. Natl. Acad. Sci.} &
Archaea as third domain (16S rRNA) & Restructured the comprehensive tree
of life; archaea are closer to eukaryotes than to bacteria \\
David Relman et al. & 1999--2006 & \emph{Proc. Natl. Acad. Sci.;
Science} & Human microbiome characterisation by 16S rRNA & Revealed
about 10\texorpdfstring{\textsuperscript{1}}{1}\texorpdfstring{\textsuperscript{3}}{3} microbial cells as an integral part of human
biology \\
Murray et al. & 2022 & \emph{The Lancet} & Global mortality from
antimicrobial resistance (1.27 million deaths, 2019) & Quantified the
scale of the AMR crisis as a call to action \\
\end{longtable}
}

\begin{center}\rule{0.5\linewidth}{0.5pt}\end{center}

\subsection*{Key Concepts and
Connections}\label{key-concepts-and-connections-7}
\addcontentsline{toc}{subsection}{Key Concepts and Connections}

\begin{figure}[htbp]
\centering
\includegraphics[width=0.82\linewidth,height=4.2in,keepaspectratio]{../figures/mermaid_inline/inline_mermaid_0114_1cfe4327f5fb.png}
\caption{Unit VII concept map - Microbiology. Red = bacteria; purple = viruses; green = microbiome; orange = epidemiological models.}
\end{figure}

\begin{center}\rule{0.5\linewidth}{0.5pt}\end{center}

\subsection*{Current Evidence Thread}\label{current-evidence-thread-7}
\addcontentsline{toc}{subsection}{Current Evidence Thread}

In microbiology, a claim is primarily as strong as the way the organism
was observed: read this unit as an evidence trail in which microbial
life is established through culture, sequencing, and surveillance rather
than as a list of named taxa and diseases. Microbiology and infectious
disease now require One Health reasoning across people, animals,
environments, genomics, and antimicrobial stewardship. As you move
through the chapters, keep a two-column note: \textbf{claim} on the
left, \textbf{evidence that would change my confidence} on the right. By
the end of the unit, each major idea should be tied to a measurement,
model, citation, or paper-based lab decision.

\subsection*{Chapter Roadmap}\label{chapter-roadmap-8}
\addcontentsline{toc}{subsection}{Chapter Roadmap}

{\def\LTcaptype{none} % do not increment counter
\begin{longtable}[]{@{}
  >{\raggedright\arraybackslash}p{(\linewidth - 6\tabcolsep) * \real{0.1556}}
  >{\raggedright\arraybackslash}p{(\linewidth - 6\tabcolsep) * \real{0.1111}}
  >{\raggedright\arraybackslash}p{(\linewidth - 6\tabcolsep) * \real{0.2889}}
  >{\raggedright\arraybackslash}p{(\linewidth - 6\tabcolsep) * \real{0.4444}}@{}}
\toprule\noalign{}
\begin{minipage}[b]{\linewidth}\raggedright
Chapter
\end{minipage} & \begin{minipage}[b]{\linewidth}\raggedright
Title
\end{minipage} & \begin{minipage}[b]{\linewidth}\raggedright
Core Question
\end{minipage} & \begin{minipage}[b]{\linewidth}\raggedright
Key Equation / Model
\end{minipage} \\
\midrule\noalign{}
\endhead
\bottomrule\noalign{}
\endlastfoot
\textbf{22} & Bacteria, Archaea, and Viruses & What distinguishes the
three major types of microorganism, and how do they replicate? &
Doubling time: \(t_d = \ln 2 / \mu\); viral burst size \\
\textbf{23} & Microbial Ecology & What roles do microorganisms play in
ecosystems and the human body? & MIC, MBC; diversity indices (Shannon
\(H'\)) \\
\textbf{24} & Infectious Disease & How do pathogens cause disease, and
how do epidemics spread? & SIR model: \(R_0 = \beta / \gamma\); basic
reproductive number \\
\end{longtable}
}

\begin{center}\rule{0.5\linewidth}{0.5pt}\end{center}

\subsection*{Connections Across the
Textbook}\label{connections-across-the-textbook-7}
\addcontentsline{toc}{subsection}{Connections Across the Textbook}

\begin{itemize}
\tightlist
\item
  \textbf{Prokaryotic cell structure} in \cref{sec:unit_II_cell_theory}
  and \cref{sec:unit_II_cell_structure} provides the foundation for
  understanding antibiotic mechanisms (cell wall, ribosome 70S targets).
\item
  \textbf{Viral replication} links to Unit IV (CRISPR discovered as a
  bacterially-encoded anti-phage defence system; RNA viruses and reverse
  transcriptase).
\item
  \textbf{Antibiotic resistance evolution} is a direct application of
  Unit VI (natural selection, fitness, molecular clocks on resistance
  genes).
\item
  \textbf{Infectious disease epidemiology} (SIR model) connects to Unit
  X (disease ecology, parasite-host dynamics in community ecology).
\item
  \textbf{Microbiome} connects to Unit IX (gut-brain axis, immune
  regulation) and Unit X (decomposer role in nutrient cycling).
\end{itemize}

\begin{quote}
\textbf{Key vocabulary introduced here:} peptidoglycan, Gram stain,
endotoxin (LPS), exotoxin, minimum inhibitory concentration (MIC), basic
reproductive number (R\texorpdfstring{\textsubscript{0}}{0}), prophage, bacteriophage, quorum sensing,
biofilm, prion, viroid, microbiome, dysbiosis, Koch's postulates.
\end{quote}

\subsection*{Computational Toolbox --- Unit
VII}\label{computational-toolbox-unit-vii}
\addcontentsline{toc}{subsection}{Computational Toolbox --- Unit VII}

\begin{Shaded}
\begin{Highlighting}[]
\ImportTok{from}\NormalTok{ biology.microbiology }\ImportTok{import}\NormalTok{ bacterial\_growth\_curve, doubling\_time}

\CommentTok{\# Bacterial growth: E. coli at 37°C, doubling time about 20 min (0.333 h)}
\NormalTok{growth }\OperatorTok{=}\NormalTok{ bacterial\_growth\_curve(N0}\OperatorTok{=}\FloatTok{1e3}\NormalTok{, doubling\_time\_hr}\OperatorTok{=}\DecValTok{1}\OperatorTok{/}\DecValTok{3}\NormalTok{, t\_end\_hr}\OperatorTok{=}\DecValTok{4}\NormalTok{, lag\_phase\_hr}\OperatorTok{=}\DecValTok{0}\NormalTok{)}
\BuiltInTok{print}\NormalTok{(}\SpecialStringTok{f"Start: }\SpecialCharTok{\{}\NormalTok{growth}\SpecialCharTok{.}\NormalTok{populations[}\DecValTok{0}\NormalTok{]}\SpecialCharTok{:.0f\}}\SpecialStringTok{ cells"}\NormalTok{)}
\BuiltInTok{print}\NormalTok{(}\SpecialStringTok{f"After 4 h: }\SpecialCharTok{\{}\NormalTok{growth}\SpecialCharTok{.}\NormalTok{populations[}\OperatorTok{{-}}\DecValTok{1}\NormalTok{]}\SpecialCharTok{:.2e\}}\SpecialStringTok{ cells"}\NormalTok{)}
\BuiltInTok{print}\NormalTok{(}\SpecialStringTok{f"Growth rate: }\SpecialCharTok{\{}\NormalTok{growth}\SpecialCharTok{.}\NormalTok{growth\_rate\_per\_hr}\SpecialCharTok{:.2f\}}\SpecialStringTok{ h⁻¹"}\NormalTok{)}
\CommentTok{\# Expected:}
\CommentTok{\# Start: 1000 cells}
\CommentTok{\# After 4 h: 4.10e+06 cells}
\CommentTok{\# Growth rate: 2.08 h⁻¹}

\CommentTok{\# Doubling time from two measurements}
\NormalTok{td }\OperatorTok{=}\NormalTok{ doubling\_time(N0}\OperatorTok{=}\FloatTok{1e3}\NormalTok{, Nt}\OperatorTok{=}\FloatTok{4.096e6}\NormalTok{, elapsed\_time\_hr}\OperatorTok{=}\DecValTok{4}\NormalTok{)}
\BuiltInTok{print}\NormalTok{(}\SpecialStringTok{f"Estimated doubling time: }\SpecialCharTok{\{}\NormalTok{td}\OperatorTok{*}\DecValTok{60}\SpecialCharTok{:.1f\}}\SpecialStringTok{ min"}\NormalTok{)}
\CommentTok{\# Expected: Estimated doubling time: 20.0 min}
\end{Highlighting}
\end{Shaded}

\begin{quote}
\textbf{Try it yourself:} Change \texttt{doubling\_time\_hr} from
\texttt{1/3} to \texttt{1.0} and compare final population size after 4
hours.
\end{quote}

\begin{center}\rule{0.5\linewidth}{0.5pt}\end{center}

\emph{Source modules: \breaktt{src/biology/microbiology/} ---
\breaktt{bacterial\_growth\_curve()}, \texttt{doubling\_time()},
\breaktt{mic\_fold\_dilution()}, \breaktt{VIRAL\_REPLICATION\_CYCLES}.}
\emph{Figures: \breaktt{src/visualization/} (growth curves, SIR
dynamics); \breaktt{src/mermaid/biology\_diagrams.py} (viral replication
cycle diagrams).}

\newpage

\section{Bacteria, Archaea, and
Viruses}\label{bacteria-archaea-and-viruses}

\label{sec:unit_VII_bacteria_archaea_viruses}

\begin{quote}
\textbf{Ch 22} \texorpdfstring{\ensuremath{\cdot}}{.} Level 2/3 \texorpdfstring{\ensuremath{\cdot}}{.} 65 min read \texorpdfstring{\ensuremath{\cdot}}{.} 75 min lecture \texorpdfstring{\ensuremath{\cdot}}{.}
Prerequisites: \cref{sec:unit_II_cell_structure}
\end{quote}

\subsection{Learning Objectives}\label{learning-objectives-24}

By the end of this chapter, you should be able to:

\begin{enumerate}
\def\labelenumi{\arabic{enumi}.}
\tightlist
\item
  Compare the three domains \citep{woese1977} of life and explain how
  rRNA phylogeny established the
  Bacteria-\hyperref[gl:archaea]{\textbf{Archaea}}-Eukarya
  classification.
\item
  Describe bacterial cell structure including peptidoglycan
  architecture, Gram staining outcomes, and the clinical significance of
  LPS.
\item
  Trace the \textbf{peptidoglycan biosynthesis pathway} through
  MurA-MurF, lipid I, and lipid II, identifying the antibiotic targeting
  each step and the resistance mechanism that defeats it.
\item
  Classify bacteria by morphology, arrangement, and metabolic type, and
  identify major pathogenic groups.
\item
  Explain horizontal \hyperref[gl:gene]{\textbf{gene}} transfer
  mechanisms (transformation, transduction, conjugation) and their role
  in bacterial evolution and antibiotic resistance \citep{lawrence1998}.
\item
  Describe \textbf{two-component signal transduction}, \textbf{biofilm}
  formation, and \textbf{quorum sensing} as the molecular basis of
  bacterial sociality, and link them to chronic infections and
  device-associated disease.
\item
  Distinguish Archaea from Bacteria at the molecular level, including
  ether-linked isoprenoid lipids and tetraether monolayers, and describe
  the significance of Asgard archaea for eukaryotic origins.
\item
  Compare the Baltimore classification of viruses, contrast lytic and
  lysogenic cycles, and describe the HIV life cycle and its drug
  targets.
\item
  Explain CRISPR-Cas as a bacterial adaptive immune system, compare
  Types I/II/III, and assess \textbf{phage therapy} as a clinical
  alternative to antibiotics.
\item
  Apply the bacterial growth equation \(N(t) = N_0 \cdot 2^{t/t_d}\) to
  predict cell densities from generation time and infer doubling time
  from optical-density data.
\end{enumerate}

\subsubsection{Study Blueprint}\label{study-blueprint-24}

\begin{itemize}
\tightlist
\item
  \textbf{Big idea:} Microbial diversity reflects different cell
  architectures, genomes, metabolisms, and evolutionary histories.
\item
  \textbf{Core concepts:} prokaryotic structure, archaea, viruses,
  horizontal gene transfer.
\item
  \textbf{Framework alignment:} Vision \& Change: Evolution, Systems,
  Structure and function; AP Biology: Evolution, Systems Interactions;
  NGSS-style topics: Structure and Function, Interdependent
  Relationships in Ecosystems.
\item
  \textbf{Model or quantitative lens:} Growth-rate, genome-size, and
  dilution calculations.
\item
  \textbf{Data skill:} Interpret microbial observations from growth,
  sequence, or structural evidence.
\item
  \textbf{Practice cadence:} Questions and Methods, Representing and
  Describing Data, Argumentation.
\item
  \textbf{Common misconception to repair:} Viruses are not simply tiny
  bacteria; they use fundamentally different replication logic.
\item
  \textbf{Primary lab:}
  \cref{sec:lab_unit_VII_bacteria_archaea_viruses}.
\item
  \textbf{Question bank:}
  \cref{sec:q_unit_VII_bacteria_archaea_viruses}.
\item
  \textbf{Transfer task:} Apply microbial diversity reasoning to
  antibiotics, biotechnology, ecology, and outbreaks.
\item
  \textbf{Bridge to computation:}
  \breaktt{biology.microbiology.microbiology.bacterial\_growth\_curve}.
\end{itemize}

\begin{center}\rule{0.5\linewidth}{0.5pt}\end{center}

\begin{quote}
\textbf{Opening Vignette --- The Rebel Biologist Who Turned Bacteria
into \hyperref[gl:organelle]{\textbf{Organelle}}s}

In 1966, Lynn Margulis --- an unknown 28-year-old junior faculty member
at Boston University --- submitted a paper proposing that mitochondria
and \hyperref[gl:chloroplast]{\textbf{chloroplast}}s were once
free-living bacteria engulfed by a host cell \citep{margulis1967}. The
serial \hyperref[gl:endosymbiosis]{\textbf{endosymbiosis}} hypothesis
was rejected by more than fifteen journals before \emph{Journal of
Theoretical Biology} finally published it. Her evidence: the matching
size and \hyperref[gl:ribosome]{\textbf{ribosome}} type of mitochondria
and bacteria, their independent division by binary fission, and the
circular DNA they share with prokaryotes. When molecular sequencing
arrived in the 1970s, it confirmed everything --- mitochondria are
phylogenetically embedded within the \texorpdfstring{\ensuremath{\alpha}}{alpha}-Proteobacteria, and chloroplasts
within the Cyanobacteria. Margulis had been completely right. The
bacterial origin of eukaryotic organelles is now foundational to cell
biology, and stands as one of the most important corrections to the
naive view that bacteria are too simple to matter for complex life.
\end{quote}

\subsection{Bacteria -- Domain Overview and Cell
Structure}\label{bacteria-domain-overview-and-cell-structure}

\begin{figure}[htbp]
\centering
\includegraphics[width=0.85\textwidth]{../figures/bacterial_growth.png}
\caption{Bacterial growth curve: four phases (lag, exponential, stationary, death) plotted as $\log(\text{cell number})$ versus time.}
\label{fig:unit_VII_bacterial_growth}
\end{figure}

\subsubsection{The Three Domains of
Life}\label{the-three-domains-of-life}

In 1977, Carl Woese and George Fox revolutionized our understanding of
life's diversity by analyzing ribosomal RNA (rRNA) sequences across
organisms \citep{woese1977}. Their phylogenetic analysis revealed that
life comprises three fundamental domains: \textbf{Bacteria},
\textbf{Archaea}, and \textbf{Eukarya}. This replaced the older
five-kingdom system that had grouped bacteria and archaea together as
``Monera.''

The 16S rRNA gene (in prokaryotes) and 18S rRNA gene (in
\hyperref[gl:eukaryote]{\textbf{eukaryote}}s) serve as molecular
chronometers because they are:

\begin{itemize}
\tightlist
\item
  Almost universally present in most cellular life
\item
  Functionally constrained (slow evolutionary rate in conserved regions)
\item
  Sufficiently variable in hypervariable regions (V1-V9) for taxonomic
  resolution
\item
  Rarely subject to horizontal gene transfer
\end{itemize}

\subsubsection{Prokaryotic Cell
Features}\label{prokaryotic-cell-features}

Bacteria and Archaea share the prokaryotic cell plan, which differs
fundamentally from the eukaryotic organization:

{\def\LTcaptype{none} % do not increment counter
\begin{longtable}[]{@{}
  >{\raggedright\arraybackslash}p{(\linewidth - 4\tabcolsep) * \real{0.2647}}
  >{\raggedright\arraybackslash}p{(\linewidth - 4\tabcolsep) * \real{0.3824}}
  >{\raggedright\arraybackslash}p{(\linewidth - 4\tabcolsep) * \real{0.3529}}@{}}
\toprule\noalign{}
\begin{minipage}[b]{\linewidth}\raggedright
Feature
\end{minipage} & \begin{minipage}[b]{\linewidth}\raggedright
Prokaryotes
\end{minipage} & \begin{minipage}[b]{\linewidth}\raggedright
Eukaryotes
\end{minipage} \\
\midrule\noalign{}
\endhead
\bottomrule\noalign{}
\endlastfoot
Nucleus & No membrane-bound nucleus; nucleoid region & Membrane-bound
nucleus with nuclear pores \\
Ribosomes & 70S (30S + 50S subunits) & 80S (40S + 60S subunits) \\
\hyperref[gl:chromosome]{\textbf{Chromosome}} & Typically circular;
nucleoid-associated \hyperref[gl:protein]{\textbf{protein}}s (HU, H-NS,
IHF) & Linear chromosomes; \hyperref[gl:histone]{\textbf{histone}}
octamer packaging \\
Cell division & Binary fission (FtsZ ring) &
\hyperref[gl:mitosis]{\textbf{Mitosis}} with spindle apparatus \\
Internal membranes & Rare (thylakoids in cyanobacteria) & Extensive
endomembrane system \\
Flagella & Proton motive force-driven rotation (not ATP) &
Dynein-powered 9+2 microtubule arrangement \\
\end{longtable}
}

The bacterial flagellum is a remarkable molecular machine. Unlike
eukaryotic flagella, which bend in a wave-like motion powered by
ATP-dependent dynein motors, the bacterial flagellum rotates like a
propeller. The basal body contains a rotary motor driven by the proton
motive force (or Na\(^+\) gradient in some marine species), spinning at
up to 1,700 revolutions per second in some species. The flagellar motor
can reverse direction, enabling the ``run and tumble'' chemotaxis
behavior in \emph{E. coli}.

\subsubsection{Bacterial Cell Wall:
Peptidoglycan}\label{bacterial-cell-wall-peptidoglycan}

The bacterial cell wall is composed of \textbf{peptidoglycan} (also
called murein), a mesh-like polymer unique to bacteria. Its structure
consists of:

\begin{itemize}
\tightlist
\item
  \textbf{Glycan strands}: Alternating units of N-acetylglucosamine
  (NAG) and N-acetylmuramic acid (NAM) linked by \texorpdfstring{\ensuremath{\beta}}{beta}-1,4 glycosidic bonds
\item
  \textbf{Tetrapeptide side chains}: Attached to NAM; typically L-Ala --
  D-Glu -- meso-DAP (or L-Lys) -- D-Ala
\item
  \textbf{Cross-links}: Peptide bridges connecting adjacent glycan
  strands, catalyzed by transpeptidases (penicillin-binding proteins,
  PBPs)
\end{itemize}

The presence of \textbf{D-amino acids} (D-glutamate, D-alanine) is
unusual in biology -- most proteins use exclusively L-amino acids.
D-amino acids confer resistance to host proteases, which evolved to
cleave L-amino acid \hyperref[gl:peptide-bond]{\textbf{peptide bond}}s.

\subsubsection{Peptidoglycan Biosynthesis: A Step-by-Step Antibiotic
Atlas}\label{peptidoglycan-biosynthesis-a-step-by-step-antibiotic-atlas}

Peptidoglycan biosynthesis proceeds in three spatial compartments ---
cytoplasm, inner-membrane interface, and periplasm --- and nearly every
step has a clinically deployed antibiotic that targets it. Understanding
the pathway is therefore equivalent to understanding the chemotherapy of
bacterial infection. The pathway begins with the housekeeping sugar
\textbf{UDP-N-acetylglucosamine (UDP-GlcNAc)} and ends with a covalently
cross-linked sacculus surrounding the entire cell.

\begin{figure}[htbp]
\centering
\includegraphics[width=0.82\linewidth,height=4.2in,keepaspectratio]{../figures/mermaid_inline/inline_mermaid_0115_ce65da087915.png}
\caption{Cytoplasm-to-periplasm peptidoglycan biosynthesis with five major antibiotic targets - fosfomycin (MurA), D-cycloserine (D-Ala-D-Ala ligase / Alr), bacitracin (C55 carrier recycling), vancomycin (D-Ala-D-Ala terminus of Lipid II), and beta-lactams (PBP transpeptidation).}
\end{figure}

{\def\LTcaptype{none} % do not increment counter
\begin{longtable}[]{@{}
  >{\raggedright\arraybackslash}p{(\linewidth - 8\tabcolsep) * \real{0.0923}}
  >{\raggedright\arraybackslash}p{(\linewidth - 8\tabcolsep) * \real{0.2769}}
  >{\raggedright\arraybackslash}p{(\linewidth - 8\tabcolsep) * \real{0.0923}}
  >{\raggedright\arraybackslash}p{(\linewidth - 8\tabcolsep) * \real{0.3846}}
  >{\raggedright\arraybackslash}p{(\linewidth - 8\tabcolsep) * \real{0.1538}}@{}}
\toprule\noalign{}
\begin{minipage}[b]{\linewidth}\raggedright
Step
\end{minipage} & \begin{minipage}[b]{\linewidth}\raggedright
Enzyme / process
\end{minipage} & \begin{minipage}[b]{\linewidth}\raggedright
Drug
\end{minipage} & \begin{minipage}[b]{\linewidth}\raggedright
Mechanism of inhibition
\end{minipage} & \begin{minipage}[b]{\linewidth}\raggedright
Spectrum
\end{minipage} \\
\midrule\noalign{}
\endhead
\bottomrule\noalign{}
\endlastfoot
1 & MurA (UDP-GlcNAc enolpyruvyl transferase) & \textbf{Fosfomycin} &
Phosphoenolpyruvate analog; covalently inhibits MurA active-site Cys via
epoxide ring opening & Broad (UTIs, \emph{E. coli},
\emph{Enterococcus}) \\
2 & MurB (UDP-GlcNAc-EP reductase) & (No clinical drug; preclinical
leads) & Reduces enolpyruvyl ether to D-lactate & --- \\
3 & MurC, MurD, MurE, MurF (amino-acid ligases) & (Preclinical
pyrazolopyrimidines, hydrazides) & Block ATP-dependent ligation of
L-Ala, D-Glu, meso-DAP/L-Lys, and D-Ala--D-Ala & Investigational \\
4 & Alr / Ddl (D-Ala racemase, D-Ala-D-Ala ligase) &
\textbf{D-cycloserine} & D-Ala structural analog; covalently traps both
racemase and ligase active-site residues & Anti-TB (second-line) \\
5 & MraY (phospho-MurNAc-pentapeptide translocase) &
\textbf{Tunicamycin} (research primarily); muraymycins (preclinical) &
Blocks attachment of pentapeptide to C55 lipid carrier (Lipid I
formation) & Toxic to humans (also blocks N-glycosylation) \\
6 & MurG (UDP-GlcNAc:Lipid I GlcNAc transferase) & (Preclinical) &
Blocks Lipid I \texorpdfstring{\ensuremath{\to}}{->} Lipid II conversion & Investigational \\
7 & C55-PP recycling (BacA/UppP phosphatase) & \textbf{Bacitracin} &
Sequesters undecaprenyl-pyrophosphate, blocking dephosphorylation;
bactoprenol pool exhausted & Topical (Gram-positive) \\
8 & Lipid II flipping (MurJ flippase) & \textbf{Lysobactin} family
(preclinical) & Trap Lipid II in cytoplasm & Investigational \\
9 & Lipid II (D-Ala-D-Ala terminus) & \textbf{Vancomycin, teicoplanin,
dalbavancin, oritavancin} & Glycopeptide H-bonds to D-Ala-D-Ala;
sterically blocks PBP access and transglycosylation & MRSA, \emph{C.
difficile}, VRE-susceptible enterococci \\
10 & Transglycosylation + transpeptidation by PBPs &
\textbf{Penicillins, cephalosporins, carbapenems, monobactams} &
\texorpdfstring{\ensuremath{\beta}}{beta}-lactam ring acylates PBP active-site serine (suicide substrate; mimics
D-Ala--D-Ala) & Broad \\
11 & Lipid II / membrane & \textbf{Daptomycin} (membrane),
\textbf{lipoglycopeptides}, \textbf{nisin} (lantibiotic) & Bind Lipid II
/ depolarize membrane & MRSA, VRE \\
\end{longtable}
}

\textbf{Resistance recapitulates the pathway.} Each drug elicits a
corresponding resistance trick: VanA-type vancomycin resistance in
enterococci (VRE) replaces D-Ala-D-Ala with D-Ala-D-Lac, abolishing one
of the five hydrogen bonds that vancomycin exploits (\texorpdfstring{\textgreater{}=}{>=} 1000-fold MIC
increase; the binding free-energy loss is about 12 kJ/mol per
H-bond); MRSA expresses \textbf{PBP2a} (encoded by \emph{mecA}) with low
\texorpdfstring{\ensuremath{\beta}}{beta}-lactam affinity (K\(_i\) rises about  1000\texorpdfstring{\ensuremath{\times}}{x}); \texorpdfstring{\ensuremath{\beta}}{beta}-lactamases
hydrolyse the \texorpdfstring{\ensuremath{\beta}}{beta}-lactam ring before it reaches its PBP target, with
extended-spectrum \texorpdfstring{\ensuremath{\beta}}{beta}-lactamases (ESBLs) and carbapenemases (KPC, NDM,
OXA-48) progressively defeating each new generation of \texorpdfstring{\ensuremath{\beta}}{beta}-lactams. The
pathway thus serves as a unifying framework --- biosynthetic step, drug,
and resistance mechanism are three views of the same chemistry.

\begin{quote}
\textbf{Clinical Connection: Why Vancomycin Stopped Working --- and How
Daptomycin Took Over} Vancomycin was introduced in 1958 and remained
dependable for almost 30 years. The first VRE isolates appeared in
1986--1988 in European and US hospitals, driven by agricultural use of
avoparcin (a glycopeptide growth promoter in livestock); by the early
2000s VRE had spread globally. The molecular trick is elegant: a single
ester linkage (D-Ala-D-Lac instead of the D-Ala-D-Ala amide) costs the
bacterium nothing but eliminates one of the five H-bonds anchoring
vancomycin. Daptomycin (FDA-approved 2003) bypasses the resistance
entirely --- it inserts into the cell membrane in a
Ca\(^{2+}\)-dependent manner, depolarising it without requiring a
peptidoglycan binding site. Daptomycin resistance is now emerging too
(mprF and yycG mutations alter membrane charge), continuing the arms
race.
\end{quote}

\subsubsection{Gram Staining: The Foundational Diagnostic
Test}\label{gram-staining-the-foundational-diagnostic-test}

The \hyperref[gl:gram-stain]{\textbf{Gram stain}}, developed by Hans
Christian Gram in 1884, remains the single most important initial test
in clinical microbiology. The procedure exploits differences in cell
wall architecture:

\begin{figure}[htbp]
\centering
\includegraphics[width=0.82\linewidth,height=4.2in,keepaspectratio]{../figures/mermaid_inline/inline_mermaid_0116_fdf180a01853.png}
\caption{Gram-stain decision tree producing purple (Gram-positive) versus pink (Gram-negative) results from differences in peptidoglycan thickness and outer-membrane architecture.}
\end{figure}

\begin{itemize}
\tightlist
\item
  \textbf{Lipid A}: The toxic moiety; embedded in the outer membrane;
  activates TLR4 on innate immune cells
\item
  \textbf{Core oligosaccharide}: Short sugar chain linking Lipid A to
  the O-antigen
\item
  \textbf{O-antigen}: Variable polysaccharide extending outward;
  serotype-specific; target of antibodies
\end{itemize}

When Gram-negative bacteria lyse during infection, released LPS triggers
massive innate immune activation through the TLR4-MD2 receptor complex.
This can lead to a cytokine cascade (TNF-\texorpdfstring{\ensuremath{\alpha}}{alpha}, IL-1\texorpdfstring{\ensuremath{\beta}}{beta}, IL-6), systemic
inflammation, disseminated intravascular coagulation (DIC), and septic
shock -- a life-threatening condition with mortality rates of 30-50\%.

\begin{quote}
\textbf{Clinical Connection: Gram Staining in Diagnosis} The Gram stain
is typically the first test performed on clinical specimens (blood
cultures, cerebrospinal fluid, sputum). Results within 15 minutes guide
empirical antibiotic therapy before culture results are available (24-72
hours). For bacterial meningitis, Gram-positive diplococci suggest
\emph{Streptococcus pneumoniae} (treat with ceftriaxone + vancomycin),
while Gram-negative diplococci suggest \emph{Neisseria meningitidis}
(treat with ceftriaxone alone). This rapid distinction can be
life-saving in an emergency where hours matter.
\end{quote}

\begin{quote}
\textbf{Concept Check 1:} A novel bacterium is isolated from a deep-sea
hydrothermal vent. Gram staining produces a pink result. Predict the
cell wall architecture and explain why the alcohol decolorization step
removed the crystal violet-iodine complex from this organism.
\end{quote}

\begin{quote}
\textbf{Concept Check 1b:} A clinical microbiology lab tests a bacterial
isolate against six wall-targeted antibiotics. The isolate is resistant
to vancomycin (MIC \textgreater{} 256 \texorpdfstring{\ensuremath{\mu}}{mu}g/mL) but susceptible to
ampicillin, fosfomycin, and bacitracin. Sequencing reveals a \emph{vanA}
operon. Using the peptidoglycan biosynthesis pathway, explain (a) which
step is altered, (b) why vancomycin fails, and (c) why ampicillin still
works.
\end{quote}

\begin{quote}
\textbf{Concept Check 1c:} Fosfomycin is a phosphoenolpyruvate (PEP)
analog that covalently inactivates MurA. \emph{Pseudomonas aeruginosa}
is intrinsically resistant to fosfomycin despite having a functional
MurA. Hypothesise two non-target-based resistance mechanisms (think
about how the drug enters the cell) and propose an experiment using a
\emph{P. aeruginosa} MurA-knockout complemented with the \emph{E. coli}
allele to test whether MurA is responsible.
\end{quote}

\begin{center}\rule{0.5\linewidth}{0.5pt}\end{center}

\subsection{Bacterial Diversity and Metabolic
Types}\label{bacterial-diversity-and-metabolic-types}

\subsubsection{Morphology and
Arrangement}\label{morphology-and-arrangement}

Bacterial cells exhibit characteristic shapes that aid in
identification:

{\def\LTcaptype{none} % do not increment counter
\begin{longtable}[]{@{}
  >{\raggedright\arraybackslash}p{(\linewidth - 4\tabcolsep) * \real{0.2333}}
  >{\raggedright\arraybackslash}p{(\linewidth - 4\tabcolsep) * \real{0.4333}}
  >{\raggedright\arraybackslash}p{(\linewidth - 4\tabcolsep) * \real{0.3333}}@{}}
\toprule\noalign{}
\begin{minipage}[b]{\linewidth}\raggedright
Shape
\end{minipage} & \begin{minipage}[b]{\linewidth}\raggedright
Description
\end{minipage} & \begin{minipage}[b]{\linewidth}\raggedright
Examples
\end{minipage} \\
\midrule\noalign{}
\endhead
\bottomrule\noalign{}
\endlastfoot
Coccus & Spherical & \emph{Staphylococcus aureus}, \emph{Streptococcus
pyogenes} \\
Bacillus & Rod-shaped & \emph{Escherichia coli}, \emph{Bacillus
subtilis} \\
Spirillum & Rigid spiral & \emph{Campylobacter jejuni},
\emph{Helicobacter pylori} \\
Spirochete & Flexible helix with axial filaments & \emph{Treponema
pallidum}, \emph{Borrelia burgdorferi} \\
Vibrio & Comma-shaped & \emph{Vibrio cholerae}, \emph{Vibrio
parahaemolyticus} \\
Pleomorphic & Variable shape & \emph{Mycoplasma pneumoniae} (no cell
wall) \\
\end{longtable}
}

Cell arrangements reflect division plane and post-division adhesion
patterns:

\begin{itemize}
\tightlist
\item
  \textbf{Diplo-} (pairs): \emph{Diplococcus}, \emph{Neisseria}
\item
  \textbf{Strepto-} (chains): \emph{Streptococcus} -- divides in one
  plane, cells remain attached
\item
  \textbf{Staphylo-} (grape-like clusters): \emph{Staphylococcus} --
  divides in multiple planes
\item
  \textbf{Tetrad} (groups of 4): \emph{Micrococcus} -- divides in two
  perpendicular planes
\item
  \textbf{Sarcina} (cuboidal packets of 8): \emph{Sarcina ventriculi} --
  divides in three planes
\end{itemize}

\subsubsection{Metabolic Diversity}\label{metabolic-diversity}

Bacteria exhibit unparalleled metabolic diversity -- far exceeding that
of eukaryotes combined. This diversity is classified by energy source
(photo- vs chemo-) and carbon source (auto- vs hetero-):

\textbf{Photoautotrophs} use light energy to fix CO\(_2\):

\begin{itemize}
\tightlist
\item
  \emph{Cyanobacteria} (oxygenic
  \hyperref[gl:photosynthesis]{\textbf{photosynthesis}}): Use water as
  electron donor; produce O\(_2\); responsible for the Great Oxidation
  Event (about 2.4 Ga); ancestors of chloroplasts
\item
  Purple and green sulfur bacteria (anoxygenic): Use H\(_2\)S or S\(^0\)
  as electron donor; no O\(_2\) production; bacteriochlorophyll absorbs
  at longer wavelengths (800-1050 nm)
\end{itemize}

\textbf{Chemoautotrophs (lithotrophs)} derive energy from inorganic
chemical oxidation:

\begin{itemize}
\tightlist
\item
  \emph{Nitrosomonas}: NH\(_3\) -\textgreater{} NO\(_2^-\) (ammonia
  oxidation; \hyperref[gl:nitrification]{\textbf{nitrification}} step 1)
\item
  \emph{Nitrobacter}: NO\(_2^-\) -\textgreater{} NO\(_3^-\) (nitrite
  oxidation; nitrification step 2)
\item
  \emph{Thiobacillus}: S\(^0\) -\textgreater{} SO\(_4^{2-}\) (sulfur
  oxidation)
\item
  \emph{Acidithiobacillus ferrooxidans}: Fe\(^{2+}\) -\textgreater{}
  Fe\(^{3+}\) (iron oxidation; used industrially in bioleaching of
  copper and gold ores)
\end{itemize}

\textbf{Heterotrophs} (the vast majority) obtain carbon from organic
compounds. They are further classified by oxygen relationship:

{\def\LTcaptype{none} % do not increment counter
\begin{longtable}[]{@{}
  >{\raggedright\arraybackslash}p{(\linewidth - 4\tabcolsep) * \real{0.1765}}
  >{\raggedright\arraybackslash}p{(\linewidth - 4\tabcolsep) * \real{0.5588}}
  >{\raggedright\arraybackslash}p{(\linewidth - 4\tabcolsep) * \real{0.2647}}@{}}
\toprule\noalign{}
\begin{minipage}[b]{\linewidth}\raggedright
Type
\end{minipage} & \begin{minipage}[b]{\linewidth}\raggedright
O\(_2\) Requirement
\end{minipage} & \begin{minipage}[b]{\linewidth}\raggedright
Example
\end{minipage} \\
\midrule\noalign{}
\endhead
\bottomrule\noalign{}
\endlastfoot
Obligate aerobe & Requires O\(_2\) & \emph{Mycobacterium
tuberculosis} \\
Obligate anaerobe & Killed by O\(_2\) & \emph{Clostridium botulinum} \\
Facultative anaerobe & Grows with or without O\(_2\) & \emph{Escherichia
coli} \\
Aerotolerant anaerobe & Does not use O\(_2\) but tolerates it &
\emph{Lactobacillus} \\
Microaerophilic & Requires low O\(_2\) (2-10\%) & \emph{Campylobacter
jejuni}, \emph{Helicobacter pylori} \\
\end{longtable}
}

\subsubsection{Major Pathogenic Groups}\label{major-pathogenic-groups}

{\def\LTcaptype{none} % do not increment counter
\begin{longtable}[]{@{}
  >{\raggedright\arraybackslash}p{(\linewidth - 6\tabcolsep) * \real{0.1404}}
  >{\raggedright\arraybackslash}p{(\linewidth - 6\tabcolsep) * \real{0.2632}}
  >{\raggedright\arraybackslash}p{(\linewidth - 6\tabcolsep) * \real{0.2105}}
  >{\raggedright\arraybackslash}p{(\linewidth - 6\tabcolsep) * \real{0.3860}}@{}}
\toprule\noalign{}
\begin{minipage}[b]{\linewidth}\raggedright
Phylum
\end{minipage} & \begin{minipage}[b]{\linewidth}\raggedright
Key Pathogens
\end{minipage} & \begin{minipage}[b]{\linewidth}\raggedright
Gram Stain
\end{minipage} & \begin{minipage}[b]{\linewidth}\raggedright
Clinical Significance
\end{minipage} \\
\midrule\noalign{}
\endhead
\bottomrule\noalign{}
\endlastfoot
Proteobacteria & \emph{E. coli}, \emph{Salmonella}, \emph{Helicobacter},
\emph{Vibrio}, \emph{Neisseria}, \emph{Pseudomonas} & Gram-negative &
UTIs, enteric infections, peptic ulcers, cholera, meningitis, nosocomial
infections \\
Firmicutes & \emph{Streptococcus}, \emph{Staphylococcus},
\emph{Clostridium}, \emph{Enterococcus} & Gram-positive & Pharyngitis,
skin infections, tetanus, botulism, hospital-acquired infections \\
Actinobacteria & \emph{Mycobacterium tuberculosis}, \emph{M. leprae},
\emph{Corynebacterium} & Gram-positive (acid-fast) & Tuberculosis,
leprosy, diphtheria \\
Spirochaetes & \emph{Treponema pallidum}, \emph{Borrelia burgdorferi},
\emph{Leptospira} & Poorly staining & Syphilis, Lyme disease,
leptospirosis \\
\end{longtable}
}

\emph{Mycobacterium} species deserve special mention: their cell walls
contain \textbf{mycolic acids} -- long-chain (\(C_{60}\)-\(C_{90}\))
branched fatty acids that form a waxy, hydrophobic barrier. This barrier
makes them resistant to Gram staining (they are classified as
\textbf{acid-fast} bacteria, stained with the Ziehl-Neelsen method using
carbol fuchsin), resistant to many antibiotics, and resistant to
desiccation -- contributing to \emph{M. tuberculosis} survival in
aerosolized droplet nuclei for hours.

\begin{quote}
\textbf{Concept Check 2:} \emph{Helicobacter pylori}
\citep{marshall1984} is a microaerophilic, spiral-shaped, Gram-negative
bacterium that colonizes the human stomach. Explain how its oxygen
requirement and morphology are adaptive for its ecological
\hyperref[gl:niche]{\textbf{niche}}, and predict what would happen if
you tried to culture it under standard
\hyperref[gl:aerobic]{\textbf{aerobic}} conditions on a typical agar
plate.
\end{quote}

\begin{center}\rule{0.5\linewidth}{0.5pt}\end{center}

\subsection{Bacterial Genetics, Signalling and
Sociality}\label{bacterial-genetics-signalling-and-sociality}

\subsubsection{Binary Fission}\label{binary-fission}

Bacteria reproduce asexually by binary fission: the circular chromosome
is replicated bidirectionally from a single origin of replication
(\emph{oriC}), the cell elongates, a septum forms at mid-cell
(coordinated by the FtsZ ring -- a tubulin homolog), and two daughter
cells separate. Generation times vary enormously:

\begin{itemize}
\tightlist
\item
  \emph{Escherichia coli}: about 20 minutes under optimal
  conditions (37 degrees C, rich media)
\item
  \emph{Mycobacterium tuberculosis}: 12-24 hours (contributing to slow
  disease progression and prolonged treatment)
\item
  \emph{Treponema pallidum}: about 33 hours (cannot be cultured
  in vitro)
\end{itemize}

Population growth during exponential phase follows the \textbf{bacterial
growth equation}:

\begin{equation}
N(t) = N_0 \cdot 2^{t/t_d}
\label{eq:unit_VII_bacterial_growth}
\end{equation}

where \(N_0\) is the initial population, \(t\) is time, and \(t_d\) is
the doubling (generation) time. Equivalently, the natural-log form
\(\ln N(t) = \ln N_0 + \mu t\) with \(\mu = (\ln 2)/t_d\) (the
\textbf{specific growth rate}, units of \(\text{time}^{-1}\)) gives a
linear plot of \(\ln N\) versus \(t\) during exponential phase, whose
slope is the experimentally measured \texorpdfstring{\ensuremath{\mu}}{mu}.

\emph{(Note: the exponent \(t/t_d\) is often represented as \(n\), the
number of generations).}

\subsubsection{Worked Example: Calculating Bacterial
Growth}\label{worked-example-calculating-bacterial-growth}

\textbf{Problem:} A microbiologist inoculates a broth with \(N_0 = 100\)
cells of \emph{Vibrio natriegens}, which has a generation time of
\(t_d = 15\) minutes under optimal conditions. Assuming the culture
immediately enters the exponential growth phase and maintains it without
nutrient depletion, how many cells will be present after 2 hours?

\textbf{Solution:}

\begin{enumerate}
\def\labelenumi{\arabic{enumi}.}
\item
  \textbf{Calculate the total time (\(t\)) in consistent units:}
  \[ t = 2 \text{ hours} = 120 \text{ minutes}  \label{eq:unit_VII_bacteria_archaea_viruses_item_1}\]
\item
  \textbf{Calculate the number of generations (\(n\)):}
  \[ n = \frac{t}{t_d} = \frac{120}{15} = 8 \text{ generations}  \label{eq:unit_VII_bacteria_archaea_viruses_item_2}\]
\item
  \textbf{Calculate the final population size \(N(t)\) from
  \cref{eq:unit_VII_bacterial_growth}:}
  \[ N(t) = 100 \cdot 2^8  \label{eq:unit_VII_bacteria_archaea_viruses_item_3}\]

  \[ N(t) = 100 \cdot 256 = 25,600 \text{ cells}  \label{eq:unit_VII_bacteria_archaea_viruses_item_4}\]
\end{enumerate}

After 2 hours, the population will reach \textbf{25,600 cells}. The
rapid doubling of bacteria means that even a small initial inoculum can
quickly reach high densities in favorable environments. After 4 hours
(16 generations), the same culture would reach
\(\approx 6.5 \times 10^6\) cells; after 6 hours (24 generations),
\(\approx 1.7 \times 10^9\) cells --- at which point nutrient limitation
forces the culture into stationary phase.

\subsubsection{Worked Example: Inferring Generation Time from an OD
Growth
Curve}\label{worked-example-inferring-generation-time-from-an-od-growth-curve}

\textbf{Problem:} A microbiology student grows \emph{E. coli} in LB
broth at 37 °C and reads optical density at 600 nm
(\(\mathrm{OD}_{600}\)) every 30 minutes:

{\def\LTcaptype{none} % do not increment counter
\begin{longtable}[]{@{}lllllll@{}}
\toprule\noalign{}
Time (min) & 0 & 30 & 60 & 90 & 120 & 150 \\
\midrule\noalign{}
\endhead
\bottomrule\noalign{}
\endlastfoot
\(\mathrm{OD}_{600}\) & 0.05 & 0.10 & 0.20 & 0.40 & 0.80 & 1.60 \\
\end{longtable}
}

Estimate the generation time \(t_d\) during exponential growth.

\textbf{Solution:}

During exponential phase, \(\mathrm{OD}_{600}\) is proportional to cell
density, so \(\mathrm{OD}(t) = \mathrm{OD}_0 \cdot 2^{t/t_d}\). Take the
natural log of both sides:

\[ \ln \mathrm{OD}(t) = \ln \mathrm{OD}_0 + \frac{t}{t_d}\ln 2  \label{eq:unit_VII_bacteria_archaea_viruses_item_5}\]

A plot of \(\ln \mathrm{OD}\) versus \(t\) has slope
\(\mu = (\ln 2)/t_d\). Using the first and last points,

\[ \mu = \frac{\ln(1.60) - \ln(0.05)}{150 - 0} = \frac{0.470 - (-2.996)}{150} = \frac{3.466}{150} \approx 0.0231\ \mathrm{min}^{-1}.  \label{eq:unit_VII_bacteria_archaea_viruses_item_6}\]

Therefore,

\[ t_d = \frac{\ln 2}{\mu} = \frac{0.693}{0.0231} \approx 30\ \mathrm{minutes}.  \label{eq:unit_VII_bacteria_archaea_viruses_item_7}\]

The doubling time is \textbf{about 30 minutes}, consistent with
rich-media growth at 37 °C. Inspection of the table confirms the answer:
\(\mathrm{OD}_{600}\) doubles every 30 min from \(t = 30\) onwards. Real
growth curves deviate at the extremes (lag at the start, stationary
plateau at the end), as the full four-phase curve in
\cref{fig:unit_VII_bacterial_growth} shows; the linear-on-log-plot
region is the main window where \(t_d\) is well defined.

\subsubsection{Horizontal Gene Transfer}\label{horizontal-gene-transfer}

While binary fission produces clonal offspring, \textbf{horizontal gene
transfer (HGT)} is the major driver of bacterial evolution and genetic
diversity \citep{lawrence1998}. Three mechanisms operate:

\textbf{Transformation} involves the uptake of naked DNA from the
environment. Some species are naturally competent -- they express
surface proteins that bind, import, and integrate exogenous DNA.
\emph{Streptococcus pneumoniae} was the organism in Griffith's famous
1928 experiment demonstrating transformation (the ``transforming
principle,'' later identified as DNA by Avery, MacLeod, and McCarty in
1944). In molecular biology, artificial competence is induced by
treating \emph{E. coli} with CaCl\(_2\) and heat shock (42 degrees C for
45 seconds), which transiently increases membrane permeability.

\textbf{Transduction} is phage-mediated DNA transfer between bacteria:

\begin{itemize}
\tightlist
\item
  \emph{Generalized transduction}: During lytic replication, bacterial
  DNA fragments are accidentally packaged into phage capsids instead of
  phage DNA. These defective phage particles inject bacterial DNA into
  new host cells, where it can recombine into the chromosome.
\item
  \emph{Specialized transduction}: A lysogenic phage (e.g.,
  bacteriophage \texorpdfstring{\ensuremath{\lambda}}{lambda}) excises imprecisely from the host chromosome,
  carrying flanking bacterial genes (\emph{gal} or \emph{bio} genes
  flanking the \texorpdfstring{\ensuremath{\lambda}}{lambda} \emph{attB} site). These hybrid phage particles
  transduce specific genes at high frequency.
\end{itemize}

\textbf{Conjugation} requires direct cell-to-cell contact via a sex
pilus encoded by the \textbf{F factor} (fertility
\hyperref[gl:plasmid]{\textbf{plasmid}}). The F\(^+\) donor cell extends
an F pilus that contacts an F\(^-\) recipient, retracting to bring cells
together. A single strand of the F plasmid is transferred and replicated
in the recipient, converting it to F\(^+\). When the F factor integrates
into the chromosome, the cell becomes \textbf{Hfr} (high frequency
\hyperref[gl:recombination]{\textbf{recombination}}), and conjugation
can transfer chromosomal genes to the recipient at high frequency.

\textbf{R-factor plasmids} (resistance plasmids) carry multiple
antibiotic resistance genes and are transferred by conjugation -- this
is the primary mechanism by which antibiotic resistance spreads among
bacterial populations, including between different species.

\subsubsection{Plasmids and Mobile Genetic
Elements}\label{plasmids-and-mobile-genetic-elements}

{\def\LTcaptype{none} % do not increment counter
\begin{longtable}[]{@{}
  >{\raggedright\arraybackslash}p{(\linewidth - 4\tabcolsep) * \real{0.3784}}
  >{\raggedright\arraybackslash}p{(\linewidth - 4\tabcolsep) * \real{0.2432}}
  >{\raggedright\arraybackslash}p{(\linewidth - 4\tabcolsep) * \real{0.3784}}@{}}
\toprule\noalign{}
\begin{minipage}[b]{\linewidth}\raggedright
Plasmid Type
\end{minipage} & \begin{minipage}[b]{\linewidth}\raggedright
Example
\end{minipage} & \begin{minipage}[b]{\linewidth}\raggedright
Significance
\end{minipage} \\
\midrule\noalign{}
\endhead
\bottomrule\noalign{}
\endlastfoot
R-plasmids & NDM-1 (New Delhi metallo-\texorpdfstring{\ensuremath{\beta}}{beta}-lactamase) & Carbapenem
resistance; global spread \\
\hyperref[gl:virulence]{\textbf{Virulence}} plasmids & pINV
(\emph{Shigella}) & Encode invasion machinery \\
Ti plasmid & \emph{Agrobacterium tumefaciens} & Crown gall disease; used
in plant biotechnology for gene transfer \\
Col plasmids & ColE1 & Encode bacteriocins (antimicrobial peptides) \\
Metabolic plasmids & TOL (\emph{Pseudomonas putida}) & Toluene
degradation; bioremediation \\
\end{longtable}
}

\textbf{Transposons} (mobile genetic elements) contribute to
\hyperref[gl:genome]{\textbf{genome}} plasticity. Insertion sequences
(IS elements) are the simplest, encoding primarily transposase flanked
by inverted repeats. Composite transposons (e.g., Tn3 carrying
ampicillin resistance) consist of a resistance gene flanked by IS
elements. Transposons can ``jump'' between chromosome and plasmid,
facilitating the assembly of multi-drug resistance cassettes.

\subsubsection{Two-Component Signal
Transduction}\label{two-component-signal-transduction}

Bacteria sense their environment through ubiquitous
\textbf{two-component systems (TCS)} --- the prokaryotic equivalent of
receptor-tyrosine-kinase + transcription-factor cascades in eukaryotes.
Each system is built from just two proteins:

\begin{enumerate}
\def\labelenumi{\arabic{enumi}.}
\tightlist
\item
  A membrane-anchored \textbf{sensor histidine kinase} (HK) with an
  extracellular sensing domain and a cytoplasmic kinase domain. Stimulus
  binding induces ATP-dependent \textbf{autophosphorylation} of a
  conserved histidine residue in the HK dimer (the phosphate is
  transferred from the \texorpdfstring{\ensuremath{\gamma}}{gamma}-phosphate of ATP to the imidazole nitrogen of
  the histidine sidechain --- a high-energy phospho-amide bond,
  \(\Delta G \approx -50\) kJ/mol on hydrolysis).
\item
  A cytoplasmic \textbf{response regulator} (RR) with a receiver domain
  and an output (usually DNA-binding) domain. The HK transfers its
  phosphate to a conserved aspartate on the RR (a phospho-anhydride;
  faster hydrolysis kinetics, \(t_{1/2} \sim\) seconds to minutes),
  activating it as a transcription factor (or, less commonly, as an
  enzyme or motor regulator).
\end{enumerate}

\begin{figure}[htbp]
\centering
\includegraphics[width=0.82\linewidth,height=4.2in,keepaspectratio]{../figures/mermaid_inline/inline_mermaid_0117_917b675ef1e3.png}
\caption{Two-component signal transduction. Stimulus to HK autophosphorylation on histidine to phosphotransfer to RR aspartate to DNA binding and gene regulation; intrinsic phosphatase activity resets the system.}
\end{figure}

The two-component circuit is therefore a \textbf{molecular switch} with
three states (off / phosphorylated / phosphatase-reset), tunable by the
relative kinase versus phosphatase activities of the HK. Some HKs
(``classical'' HKs like EnvZ and PhoQ) act as kinase under stimulus and
phosphatase otherwise --- a single protein computing both ON and OFF,
with the stimulus ratio shifting the equilibrium. This bifunctionality
makes the steady-state level of RRabout  P a near-linear
readout of the input --- exactly what is wanted for analog signal
processing inside a cell.

{\def\LTcaptype{none} % do not increment counter
\begin{longtable}[]{@{}
  >{\raggedright\arraybackslash}p{(\linewidth - 6\tabcolsep) * \real{0.2973}}
  >{\raggedright\arraybackslash}p{(\linewidth - 6\tabcolsep) * \real{0.1351}}
  >{\raggedright\arraybackslash}p{(\linewidth - 6\tabcolsep) * \real{0.1081}}
  >{\raggedright\arraybackslash}p{(\linewidth - 6\tabcolsep) * \real{0.4595}}@{}}
\toprule\noalign{}
\begin{minipage}[b]{\linewidth}\raggedright
Two-component system
\end{minipage} & \begin{minipage}[b]{\linewidth}\raggedright
Stimulus
\end{minipage} & \begin{minipage}[b]{\linewidth}\raggedright
Output
\end{minipage} & \begin{minipage}[b]{\linewidth}\raggedright
Clinical/biological significance
\end{minipage} \\
\midrule\noalign{}
\endhead
\bottomrule\noalign{}
\endlastfoot
\textbf{EnvZ--OmpR} (\emph{E. coli}) & Osmolarity & Switches porin
expression OmpC \texorpdfstring{\ensuremath{\leftrightarrow}}{<->} OmpF & Outer-membrane permeability; \texorpdfstring{\ensuremath{\beta}}{beta}-lactam entry \\
\textbf{PhoQ--PhoP} (\emph{Salmonella}) & Low Mg\texorpdfstring{\textsuperscript{2}}{2}\texorpdfstring{\textsuperscript{+}}{+}, antimicrobial
peptides, low pH (phagosome) & Lipid A modification; \emph{pmrHFIJKLM} &
Resistance to host defensins and polymyxins \\
\textbf{AgrC--AgrA} (\emph{S. aureus}) & Autoinducing peptide AIP
(quorum) & RNAIII; toxin production & Switches biofilm OFF, virulence
ON \\
\textbf{VraS--VraR} (\emph{S. aureus}) & Cell-wall stress (\texorpdfstring{\ensuremath{\beta}}{beta}-lactams,
vancomycin) & Cell-wall stimulon & Vancomycin tolerance \\
\textbf{CheA--CheY} (chemotaxis) & Attractants/repellents & Flagellar
motor switching (CW \texorpdfstring{\ensuremath{\leftrightarrow}}{<->} CCW) & Run-and-tumble swimming \\
\textbf{WalK--WalR} (\emph{B. subtilis}, \emph{S. aureus}) & Cell-wall
integrity & Autolysin and PG-synthesis genes & Essential --- drug target
lead \\
\end{longtable}
}

Key features that make TCSs central to bacterial biology: they are
\textbf{modular} (HKs and RRs from one organism often work with
components from another), they are \textbf{fast} (sub-second response),
and there are \textbf{typically 30--80 systems per genome}, allowing
combinatorial integration of many environmental signals. The minimal
\emph{Mycoplasma genitalium} genome has 0 TCSs (it lives in a stable
host environment); free-living soil bacteria like \emph{Myxococcus
xanthus} encode \textgreater{} 100. Several recent antibacterial leads
(e.g., walkmycin against the WalK kinase, an essential cell-wall TCS)
target HKs precisely because they are absent from human cells.

\subsubsection{Biofilms and Quorum
Sensing}\label{biofilms-and-quorum-sensing}

Most bacteria in nature do not live as solitary planktonic cells; they
live in \textbf{biofilms} --- surface-attached, matrix-encased
communities that behave more like a tissue than a free-swimming
population. Biofilm formation proceeds through a stereotyped life cycle,
coordinated by \textbf{quorum sensing (QS)} --- cell-density-dependent
communication via diffusible \textbf{autoinducers} (AHLs in
Gram-negatives, AIPs in Gram-positives, AI-2 across species).

\begin{figure}[htbp]
\centering
\includegraphics[width=0.82\linewidth,height=4.2in,keepaspectratio]{../figures/mermaid_inline/inline_mermaid_0118_30f4d049f95f.png}
\caption{Biofilm life cycle from planktonic cell through reversible/irreversible attachment, microcolony, mature biofilm, and quorum-triggered dispersal - coordinated by intracellular c-di-GMP and intercellular autoinducers.}
\end{figure}

\begin{itemize}
\tightlist
\item
  \textbf{Exopolysaccharides} (PIA in \emph{S. aureus}; alginate, Pel,
  Psl in \emph{P. aeruginosa}; cellulose and PNAG in enterics) ---
  provide structural cohesion and protection.
\item
  \textbf{Extracellular DNA (eDNA)} --- released by autolysis; a major
  load-bearing component and a substrate for horizontal gene transfer at
  frequencies up to 1000\texorpdfstring{\ensuremath{\times}}{x} higher than in planktonic cells.
\item
  \textbf{Matrix proteins} --- amyloid-like curli fibres (\emph{E.
  coli}), Bap proteins, type IV pili.
\item
  \textbf{Lipids and outer-membrane vesicles} --- concentrate hydrolytic
  enzymes and signalling molecules.
\end{itemize}

\textbf{Quorum sensing molecules.}

{\def\LTcaptype{none} % do not increment counter
\begin{longtable}[]{@{}
  >{\raggedright\arraybackslash}p{(\linewidth - 8\tabcolsep) * \real{0.1786}}
  >{\raggedright\arraybackslash}p{(\linewidth - 8\tabcolsep) * \real{0.3036}}
  >{\raggedright\arraybackslash}p{(\linewidth - 8\tabcolsep) * \real{0.1964}}
  >{\raggedright\arraybackslash}p{(\linewidth - 8\tabcolsep) * \real{0.1786}}
  >{\raggedright\arraybackslash}p{(\linewidth - 8\tabcolsep) * \real{0.1429}}@{}}
\toprule\noalign{}
\begin{minipage}[b]{\linewidth}\raggedright
QS class
\end{minipage} & \begin{minipage}[b]{\linewidth}\raggedright
Signal molecule
\end{minipage} & \begin{minipage}[b]{\linewidth}\raggedright
Producers
\end{minipage} & \begin{minipage}[b]{\linewidth}\raggedright
Receptor
\end{minipage} & \begin{minipage}[b]{\linewidth}\raggedright
Output
\end{minipage} \\
\midrule\noalign{}
\endhead
\bottomrule\noalign{}
\endlastfoot
\textbf{AHLs} & N-acyl-homoserine-lactones (e.g., 3-oxo-C\texorpdfstring{\textsubscript{1}}{1}\texorpdfstring{\textsubscript{2}}{2}-HSL) &
Gram-negatives (\emph{P. aeruginosa} LasI/LasR, \emph{V. fischeri}
LuxI/LuxR) & LuxR-family TF & Bioluminescence, virulence, biofilm \\
\textbf{AIPs} & Cyclic autoinducing peptides (5--10 aa, thiolactone
bridge) & Gram-positives (\emph{S. aureus} Agr, \emph{S. pneumoniae}
ComD) & Membrane HK (TCS) & Virulence/biofilm switch \\
\textbf{AI-2} & Furanosyl-borate diester (LuxS) & Cross-species &
LuxP/LsrB & Inter-species coordination \\
\textbf{PQS} & Pseudomonas quinolone signal
(2-heptyl-3-hydroxy-4-quinolone) & \emph{P. aeruginosa} & PqsR & Iron
acquisition, virulence \\
\textbf{DSF} & cis-2-unsaturated fatty acids & \emph{Xanthomonas},
\emph{Burkholderia} & RpfC HK & Biofilm dispersal, virulence \\
\end{longtable}
}

\textbf{Why biofilms matter clinically.} The CDC estimates that
\textbf{about 65--80 \% of chronic and device-associated human
infections involve biofilms}. They are 10--1000\texorpdfstring{\ensuremath{\times}}{x} more antibiotic
tolerant than planktonic cells because of (1) reduced antibiotic
penetration through EPS --- vancomycin diffusion can be retarded by
100\texorpdfstring{\ensuremath{\times}}{x}; (2) \textbf{persister cells} --- metabolically dormant
subpopulations not killed by growth-dependent antibiotics, comprising
0.001--1 \% of biofilm cells; (3) localised enzyme accumulation (e.g.,
\texorpdfstring{\ensuremath{\beta}}{beta}-lactamase concentrated 100--1000\texorpdfstring{\ensuremath{\times}}{x} in matrix); and (4) anaerobic
interior microenvironments where aminoglycosides (which require an
electrochemical gradient for uptake) are inactive. Canonical clinical
biofilms include cystic fibrosis lung infections (\emph{P. aeruginosa}
alginate biofilms), prosthetic joint and heart-valve infections
(\emph{S. aureus}, \emph{S. epidermidis}), catheter-associated UTIs
(\emph{E. coli}, \emph{P. mirabilis}), chronic wounds, and dental plaque
(\emph{Streptococcus}, \emph{Porphyromonas}, \emph{Fusobacterium}).

\textbf{Quorum quenching as therapy.} Because most virulence factors are
quorum-regulated, blocking QS \emph{attenuates} virulence without
killing bacteria --- sidestepping the strong selection for resistance
that lytic antibiotics impose. Strategies include: \textbf{(1)}
lactonases (AiiA from \emph{Bacillus}, AhlD) and acylases that hydrolyse
AHLs; \textbf{(2)} synthetic AIP analogs that competitively inhibit
AgrC; \textbf{(3)} halogenated furanones (originally from the seaweed
\emph{Delisea pulchra}) that destabilise LuxR-family receptors;
\textbf{(4)} PqsR antagonists (e.g., M64); \textbf{(5)} c-di-GMP
signalling disruptors. Several are in early clinical trials (notably for
\emph{P. aeruginosa} in cystic fibrosis), though ``anti-virulence''
compounds face the challenge that they do not directly clear infection
--- they must work synergistically with antibiotics or host immunity.

\subsubsection{Endospores}\label{endospores}

Certain Firmicutes -- primarily \emph{Bacillus} (central or subterminal
spores) and \emph{Clostridium} (terminal spores) -- produce
\textbf{endospores} in response to nutrient depletion. The endospore is
the most resistant biological structure known:

\begin{itemize}
\tightlist
\item
  \textbf{Core}: Dehydrated \hyperref[gl:cytoplasm]{\textbf{cytoplasm}};
  calcium-dipicolinic acid complex stabilizes DNA; small acid-soluble
  spore proteins (SASPs) protect DNA from UV and desiccation
\item
  \textbf{Resistance}: Survives boiling (100 degrees C), desiccation, UV
  radiation, and chemical disinfectants
\item
  \textbf{Sterilization}: Requires autoclaving at 121 degrees C, 15 psi
  pressure, for 15-20 minutes
\item
  \textbf{Germination}: Triggered by favorable conditions (nutrients,
  water); rapid return to vegetative growth
\end{itemize}

\begin{quote}
\textbf{Clinical Connection: Endospores and Hospital Infection Control}
\emph{Clostridioides difficile} endospores persist on hospital surfaces
for months, resisting alcohol-based hand sanitizers (which are effective
against vegetative bacteria). Healthcare workers must use soap and water
(physical removal) and bleach-based surface disinfectants to prevent
\emph{C. difficile} transmission. \emph{Bacillus anthracis} spores were
used as a bioterrorism agent in the 2001 U.S. anthrax letter attacks,
contaminating postal facilities and requiring years of decontamination.
\end{quote}

\begin{quote}
\textbf{Concept Check 2b:} A research lab knocks out the \emph{agrA}
gene (response regulator of the Agr quorum-sensing TCS) in \emph{S.
aureus}. The mutant forms thicker, more persistent biofilms in vitro but
is dramatically less virulent in a mouse skin-abscess model. Reconcile
these two observations using the AgrA on/off switch between biofilm and
toxin production.
\end{quote}

\begin{quote}
\textbf{Concept Check 2c:} A patient develops a \emph{Pseudomonas
aeruginosa} bloodstream infection following indwelling-catheter
placement. Blood cultures show vancomycin-susceptible cells with MIC = 1
\texorpdfstring{\ensuremath{\mu}}{mu}g/mL, yet the patient fails 4 weeks of intravenous vancomycin while the
catheter remains in place. Explain --- using biofilm penetration
kinetics, persister biology, and oxygen gradients --- why the planktonic
MIC is a poor predictor of clinical efficacy in catheter-associated
infection. What therapeutic action almost typically succeeds where
antibiotic monotherapy fails?
\end{quote}

\begin{center}\rule{0.5\linewidth}{0.5pt}\end{center}

\subsection{Archaea}\label{archaea}

\subsubsection{Distinguishing Features}\label{distinguishing-features}

Archaea are prokaryotes that were long classified with bacteria but
differ fundamentally at the molecular level:

{\def\LTcaptype{none} % do not increment counter
\begin{longtable}[]{@{}
  >{\raggedright\arraybackslash}p{(\linewidth - 6\tabcolsep) * \real{0.2432}}
  >{\raggedright\arraybackslash}p{(\linewidth - 6\tabcolsep) * \real{0.2703}}
  >{\raggedright\arraybackslash}p{(\linewidth - 6\tabcolsep) * \real{0.2432}}
  >{\raggedright\arraybackslash}p{(\linewidth - 6\tabcolsep) * \real{0.2432}}@{}}
\toprule\noalign{}
\begin{minipage}[b]{\linewidth}\raggedright
Feature
\end{minipage} & \begin{minipage}[b]{\linewidth}\raggedright
Bacteria
\end{minipage} & \begin{minipage}[b]{\linewidth}\raggedright
Archaea
\end{minipage} & \begin{minipage}[b]{\linewidth}\raggedright
Eukarya
\end{minipage} \\
\midrule\noalign{}
\endhead
\bottomrule\noalign{}
\endlastfoot
Membrane lipids & Ester-linked fatty acids & \textbf{Ether-linked
isoprenoid (phytanyl) chains}; some form tetraether monolayers &
Ester-linked fatty acids \\
Cell wall & Peptidoglycan (murein) & \textbf{No peptidoglycan};
pseudomurein (methanogens), S-layer protein, or none & No peptidoglycan;
cellulose/chitin in plants/fungi \\
RNA polymerase & Single, simple (4-5 subunits) & \textbf{Multiple
subunits (12+), resembles eukaryotic Pol II} & Three RNA polymerases
(Pol I, II, III) \\
Initiator tRNA & Formyl-methionine & \textbf{Methionine} (like
eukaryotes) & Methionine \\
Histones & HU, H-NS (not true histones) & \textbf{True histone homologs}
& Histone octamers (H2A, H2B, H3, H4) \\
\hyperref[gl:intron]{\textbf{Intron}}s & Rare & Present in some tRNA
genes & Abundant \\
\end{longtable}
}

\subsubsection{Archaeal Lipid Biochemistry: Why Hot Vents Don't Melt
Membranes}\label{archaeal-lipid-biochemistry-why-hot-vents-dont-melt-membranes}

The single most diagnostic biochemical feature of Archaea is the
architecture of their membrane lipids. Three differences from
bacterial/eukaryotic phospholipids are simultaneously present, and most
three contribute to extreme stability:

\begin{enumerate}
\def\labelenumi{\arabic{enumi}.}
\tightlist
\item
  \textbf{Ether linkages} instead of ester linkages between glycerol and
  the hydrocarbon chains. Ether bonds (C--O--C) are far more chemically
  inert than ester bonds (C--O--C(=O)) --- they resist acid hydrolysis,
  base hydrolysis, oxidation, and high-temperature cleavage.
\item
  \textbf{Isoprenoid (phytanyl) chains} (built from C5 isoprene units,
  with branched methyl groups every 4 carbons) instead of straight-chain
  fatty acids. Methyl branches cause kinks that prevent tight chain
  crystallization at low temperatures (preventing membrane gelling)
  while van-der-Waals interactions between branches stiffen the membrane
  against fluidization at high temperatures.
\item
  \textbf{Glycerol stereochemistry}: archaeal glycerol is
  \emph{sn}-glycerol-1-phosphate (G-1-P); bacterial/eukaryotic glycerol
  is \emph{sn}-glycerol-3-phosphate (G-3-P). The two are enantiomers ---
  a deep, ``lipid divide'' between archaea and the rest of cellular life
  that probably reflects independent invention of phospholipid
  biosynthesis after the LUCA.
\end{enumerate}

Hyperthermophilic archaea (\emph{Sulfolobus}, \emph{Thermococcus},
\emph{Pyrodictium}) take this further by linking two phytanyl chains
end-to-end, forming \textbf{glycerol-dialkyl-glycerol-tetraether (GDGT)}
lipids. Two GDGT molecules then span the entire membrane as a
\textbf{monolayer} (rather than a bilayer), with two glycerol headgroups
(one on each surface) and one continuous hydrocarbon block in the
middle. The covalent backbone makes the membrane behave mechanically
like a single rigid sheet --- it cannot delaminate, leak, or peel apart
even at 113 °C.

Quantitatively: a typical bacterial membrane fluidizes (loses lipid
order) above about 60 °C and disintegrates above
about 100 °C; \emph{Pyrodictium} GDGT monolayers retain
integrity at 113 °C and pH 1, conditions that hydrolyse most
ester-linked lipids in seconds. \textbf{Cyclopentane rings} within the
GDGT chains can be added or removed on the fly to fine-tune membrane
fluidity to growth temperature --- a remarkable example of biochemical
homeoviscous adaptation. Because eukaryotic and bacterial membranes lack
these adaptations, penicillin-class and polymyxin-class antibiotics that
act on bacterial cell envelopes are generally ineffective against
archaea.

\subsubsection{Major Archaeal Groups}\label{major-archaeal-groups}

\textbf{Euryarchaeota} encompasses the broadest metabolic diversity:

\begin{itemize}
\item
  \textbf{Methanogens}: Strict anaerobes that produce methane as a
  metabolic end product. The key reaction is:
  \(\text{CO}_2 + 4\text{H}_2 \rightarrow \text{CH}_4 + 2\text{H}_2\text{O}\).
  Methanogens inhabit \hyperref[gl:anaerobic]{\textbf{anaerobic}}
  environments including wetlands, rice paddies, landfills, and the
  rumen of cattle. \emph{Methanobrevibacter smithii} is the
  \hyperref[gl:dominant]{\textbf{dominant}} archaeon in the human gut.
  Globally, methanogenesis produces approximately 1 billion tonnes of
  CH\(_4\) annually, contributing approximately 10\% of anthropogenic
  greenhouse gas emissions.
\item
  \textbf{Extreme halophiles}: \emph{Halobacterium salinarum} thrives in
  saturated salt solutions (about 5 M NaCl). It maintains
  osmotic balance by accumulating intracellular KCl to equimolar
  concentrations. Its purple membrane contains
  \textbf{bacteriorhodopsin}, a light-driven proton pump that generates
  ATP via \hyperref[gl:chemiosmosis]{\textbf{chemiosmosis}} -- a form of
  phototrophy independent of
  \hyperref[gl:chlorophyll]{\textbf{chlorophyll}}. The salt flats
  visible from space (pink/red coloration in San Francisco Bay salt
  ponds) owe their color to halophilic archaea.
\item
  \textbf{Thermoacidophiles}: \emph{Thermoplasma acidophilum} grows
  optimally at 59 degrees C and \hyperref[gl:ph]{\textbf{pH}} 2.
  Notably, it lacks a cell wall entirely, surviving in hot acid through
  a reinforced lipid membrane.
\end{itemize}

\textbf{Crenarchaeota} includes the most extreme hyperthermophiles:

\begin{itemize}
\tightlist
\item
  \emph{Sulfolobus solfataricus}: 80 degrees C, pH 2-3; model organism
  for archaeal molecular biology
\item
  \emph{Pyrodictium occultum}: Grows at 113 degrees C -- among the
  highest known temperatures for life
\item
  \emph{Thermoproteus tenax}: Autotrophic sulfur metabolism at 85
  degrees C
\item
  These organisms possess \textbf{reverse gyrase}, an
  \hyperref[gl:enzyme]{\textbf{enzyme}} that introduces positive
  supercoils into DNA, stabilizing the double helix at extreme
  temperatures
\end{itemize}

\subsubsection{Asgard Archaea: The Closest Living Relatives of
Eukaryotes}\label{asgard-archaea-the-closest-living-relatives-of-eukaryotes}

\textbf{Asgard Archaea} represent perhaps the most important discovery
in evolutionary biology of the past decade. In 2015, Ettema and
colleagues identified novel archaeal lineages in metagenomic data from
deep-sea sediments near Loki's Castle hydrothermal vent field in the
Arctic Mid-Ocean Ridge. The clade has since expanded into a large
superphylum named for Norse mythology:

\begin{itemize}
\tightlist
\item
  \textbf{Lokiarchaeota} (Loki's Castle, 2015) --- the founding lineage;
  \textgreater{} 100 published MAGs.
\item
  \textbf{Thorarchaeota} (named for Thor) --- abundant in marine
  sediments.
\item
  \textbf{Odinarchaeota}, \textbf{Heimdallarchaeota},
  \textbf{Helarchaeota}, \textbf{Hermodarchaeota},
  \textbf{Wukongarchaeota} --- discovered in successive metagenomic
  surveys 2017--2024.
\end{itemize}

\textbf{Eukaryotic signature proteins (ESPs).} What makes Asgard archaea
revolutionary is that they encode an unprecedented cluster of proteins
previously thought to be eukaryote-specific:

{\def\LTcaptype{none} % do not increment counter
\begin{longtable}[]{@{}
  >{\raggedright\arraybackslash}p{(\linewidth - 4\tabcolsep) * \real{0.2500}}
  >{\raggedright\arraybackslash}p{(\linewidth - 4\tabcolsep) * \real{0.3636}}
  >{\raggedright\arraybackslash}p{(\linewidth - 4\tabcolsep) * \real{0.3864}}@{}}
\toprule\noalign{}
\begin{minipage}[b]{\linewidth}\raggedright
ESP class
\end{minipage} & \begin{minipage}[b]{\linewidth}\raggedright
Asgard homolog
\end{minipage} & \begin{minipage}[b]{\linewidth}\raggedright
Eukaryotic role
\end{minipage} \\
\midrule\noalign{}
\endhead
\bottomrule\noalign{}
\endlastfoot
\textbf{Actin} & Lokiactin (about  60 \% identity to
eukaryotic actin) & Cytoskeleton, cell shape \\
\textbf{Profilin / gelsolin} & Loki-profilin & Actin polymerization
control \\
\textbf{Rab / Arf GTPases} & Lokirhabs & Membrane trafficking, vesicle
coats \\
\textbf{ESCRT-I/II/III} & Lokiarchaeal ESCRT homologs & Membrane
scission, multivesicular bodies \\
\textbf{Tubulin (sometimes)} & Heimdall artubulins & Microtubule
cytoskeleton (eukaryotes primarily) \\
\textbf{Ubiquitin-like modifiers} & Loki SAMPs/Ub-like & Protein
degradation tagging \\
\textbf{Eukaryotic-style ribosomal proteins} & Several rps homologs &
Translation \\
\end{longtable}
}

\textbf{Phylogenetic placement.} Maximum-likelihood and Bayesian
analyses of conserved markers consistently place Asgard archaea
(especially Heimdallarchaeota) as the \textbf{sister group to most
eukaryotes}, supporting the \textbf{``two-domain tree''} of life
(Bacteria as one domain; Archaea + Eukarya as a single,
rooted-from-within-Archaea clade). This is the \textbf{eocyte
hypothesis} validated at high resolution. The evolutionary inference:
eukaryotes arose by \textbf{endosymbiosis between an Asgard-like
archaeal host} (which contributed the cytoplasm, cytoskeleton, and
information-processing genes) \textbf{and an \texorpdfstring{\ensuremath{\alpha}}{alpha}-proteobacterial
endosymbiont} (which became the mitochondrion).

\textbf{Cultivation breakthrough.} In 2020, \emph{Candidatus
Prometheoarchaeum syntrophicum} --- the first cultured Asgard archaeon
--- was reported by Imachi \emph{et al.} (\emph{Nature}) after 12 years
of effort using extremely slow-growth bioreactor cultures fed with
peptides and hydrogen. The organism grows obligately syntrophically with
\emph{Halodesulfovibrio} (sulfate-reducing) and a methanogenic archaeon
(consuming H\(_2\) produced by \emph{Prometheoarchaeum}). Its doubling
time is about  14--25 days; cells are tiny (about 
0.5 \texorpdfstring{\ensuremath{\mu}}{mu}m) and \textbf{extend long branching tentacle-like protrusions}
thought to deliver hydrogen to syntrophic partners. The remarkable
observation is that those tentacles morphologically resemble what one
might predict for an archaeal cell preparing to engulf a bacterial
endosymbiont --- the \textbf{``E\(^3\) model'' of eukaryogenesis}
(Entangle, Engulf, Endogenise) directly suggested by
\emph{Prometheoarchaeum} morphology.

\textbf{Implications.} The discovery of Asgard archaea has narrowed the
``evolutionary gap'' between prokaryotes and eukaryotes from a vast
unbridgeable distance to a graded continuum. Many of the molecular
machines once thought to be uniquely eukaryotic (membrane-trafficking
ESCRTs, actin cytoskeleton, ubiquitin signalling) were already present
in the archaeal common ancestor of eukaryotes. Eukaryogenesis was a
\emph{combinatorial} event: pre-existing archaeal ESPs + acquisition of
a mitochondrion + nucleus + later innovations. Two outstanding questions
remain: (1) how exactly did the first eukaryotic cell engulf the
proto-mitochondrion (phagocytosis-first, syntrophy-first, or
virus-mediated?); (2) which Asgard subgroup is the closest living
relative --- Heimdallarchaeota (current consensus) or a yet-undiscovered
lineage?

\subsubsection{Archaeal Biotechnology}\label{archaeal-biotechnology}

Archaeal enzymes, evolved for extreme conditions, have enormous
biotechnological value:

\begin{itemize}
\tightlist
\item
  \textbf{Taq polymerase} from \emph{Thermus aquaticus} (technically a
  bacterium, but often discussed alongside thermophilic archaea): 70
  degrees C optimum; enabled the
  \hyperref[gl:polymerase-chain-reaction]{\textbf{polymerase chain
  reaction (PCR)}} -- arguably the single most important technique in
  molecular biology; Kary Mullis, Nobel Prize 1993
\item
  \textbf{Pfu polymerase} from \emph{Pyrococcus furiosus}: 3'-to-5'
  proofreading exonuclease activity; higher fidelity than Taq; used when
  accuracy is critical (cloning, mutagenesis)
\item
  \textbf{Vent polymerase} from \emph{Thermococcus litoralis}:
  thermostable with proofreading
\item
  Industrial applications: archaeal proteases, lipases, and amylases for
  high-temperature industrial processes; extremophilic enzymes for
  detergents, food processing, and biofuel production
\end{itemize}

\begin{quote}
\textbf{Concept Check 3:} Explain why antibiotics that target
peptidoglycan synthesis (e.g., penicillin, vancomycin) are ineffective
against archaea. What does this tell us about the evolutionary
relationship between bacterial and archaeal cell walls?
\end{quote}

\begin{quote}
\textbf{Concept Check 3b:} Hyperthermophiles use GDGT tetraether
monolayers; mesophilic archaea use diether bilayers. Predict,
qualitatively, how membrane permeability to small molecules (water,
protons, ATP) should differ between the two architectures and what this
implies for chemiosmotic ATP synthesis at 100 °C.
\end{quote}

\begin{quote}
\textbf{Concept Check 3c:} Asgard archaea encode actin homologs that
function in vitro indistinguishably from eukaryotic actin. If the
proto-eukaryote inherited actin from an archaeal ancestor, what testable
prediction does this make about the timing of cytoskeletal evolution
relative to the mitochondrial endosymbiosis? Design a
comparative-genomic experiment that could distinguish between the
alternatives.
\end{quote}

\begin{center}\rule{0.5\linewidth}{0.5pt}\end{center}

\subsection{Viruses}\label{viruses}

\subsubsection{General Features}\label{general-features}

Viruses are obligate intracellular
\hyperref[gl:parasite]{\textbf{parasite}}s that cannot carry out
metabolism independently. They lack ribosomes, cannot generate ATP, and
require host cell machinery for replication. Whether viruses are
``alive'' remains a philosophical and definitional question -- they
exhibit heredity and evolution but not autonomous metabolism or cellular
organization.

A complete virus particle (virion) consists of:

\begin{itemize}
\tightlist
\item
  \textbf{Genome}: DNA or RNA (rarely both); single-stranded or
  double-stranded; linear or circular; segmented or non-segmented
\item
  \textbf{Capsid}: Protein shell assembled from capsomeres; icosahedral
  (20 triangular faces -- \emph{adenovirus}), helical (\emph{tobacco
  mosaic virus}), or complex (\emph{bacteriophage T4})
\item
  \textbf{Envelope} (some viruses): Lipid bilayer derived from host cell
  membrane during budding; contains viral glycoproteins essential for
  host cell attachment and entry
\end{itemize}

\subsubsection{Baltimore Classification}\label{baltimore-classification}

David Baltimore (Nobel Prize 1975) classified viruses by genome type and
replication strategy into seven classes:

{\def\LTcaptype{none} % do not increment counter
\begin{longtable}[]{@{}
  >{\raggedright\arraybackslash}p{(\linewidth - 6\tabcolsep) * \real{0.1522}}
  >{\raggedright\arraybackslash}p{(\linewidth - 6\tabcolsep) * \real{0.1739}}
  >{\raggedright\arraybackslash}p{(\linewidth - 6\tabcolsep) * \real{0.4565}}
  >{\raggedright\arraybackslash}p{(\linewidth - 6\tabcolsep) * \real{0.2174}}@{}}
\toprule\noalign{}
\begin{minipage}[b]{\linewidth}\raggedright
Class
\end{minipage} & \begin{minipage}[b]{\linewidth}\raggedright
Genome
\end{minipage} & \begin{minipage}[b]{\linewidth}\raggedright
Replication Strategy
\end{minipage} & \begin{minipage}[b]{\linewidth}\raggedright
Examples
\end{minipage} \\
\midrule\noalign{}
\endhead
\bottomrule\noalign{}
\endlastfoot
I & dsDNA & DNA -\textgreater{} mRNA (host RNA Pol) & Herpesviruses,
adenoviruses, poxviruses, bacteriophage T4 \\
II & ssDNA & ssDNA -\textgreater{} dsDNA -\textgreater{} mRNA &
Parvoviruses, bacteriophage \texorpdfstring{\ensuremath{\varphi}}{phi}X174 \\
III & dsRNA & dsRNA -\textgreater{} mRNA (viral RdRp) & Reoviruses,
rotaviruses \\
IV & (+)ssRNA & RNA serves directly as mRNA & Poliovirus, rhinovirus,
SARS-CoV-2, hepatitis C \\
V & (-)ssRNA & RNA -\textgreater{} mRNA (viral RdRp) & Influenza,
rabies, Ebola, measles \\
VI & (+)ssRNA-RT & RNA -\textgreater{} DNA (reverse transcriptase)
-\textgreater{} mRNA & HIV, HTLV \\
VII & dsDNA-RT & dsDNA -\textgreater{} RNA -\textgreater{} dsDNA
(reverse transcriptase) & Hepatitis B \\
\end{longtable}
}

\subsubsection{Bacteriophage Life
Cycles}\label{bacteriophage-life-cycles}

\begin{figure}[htbp]
\centering
\includegraphics[width=0.82\linewidth,height=4.2in,keepaspectratio]{../figures/mermaid_inline/inline_mermaid_0119_d56546b6ad20.png}
\caption{State diagram of the phage lambda lytic / lysogenic decision and the SOS-induced switch from lysogeny to lytic growth.}
\end{figure}

\begin{enumerate}
\def\labelenumi{\arabic{enumi}.}
\tightlist
\item
  \textbf{Adsorption}: Tail fibers (long tail fibers recognize LPS;
  short tail fibers make irreversible contact) bind the surface of
  \emph{E. coli}
\item
  \textbf{DNA injection}: The tail sheath contracts like a syringe,
  driving the tail tube through the outer membrane and injecting
  about 169 kb of linear dsDNA
\item
  \textbf{Early gene expression}: Host RNA polymerase is hijacked;
  phage-encoded anti-sigma factors redirect
  \hyperref[gl:transcription]{\textbf{transcription}}; host DNA is
  degraded by phage nucleases (recycling
  \hyperref[gl:nucleotide]{\textbf{nucleotide}}s)
\item
  \textbf{DNA replication}: Phage DNA replicates using phage-encoded DNA
  polymerase; hydroxymethylcytosine replaces cytosine (protecting phage
  DNA from its own restriction enzymes)
\item
  \textbf{Late gene expression}: Structural proteins (head, tail, tail
  fibers, baseplate) are synthesized
\item
  \textbf{Assembly}: Heads are filled with DNA (headful packaging);
  tails are assembled separately; components join spontaneously
\item
  \textbf{Lysis}: Holin creates pores in the inner membrane; endolysin
  (lysozyme) degrades peptidoglycan; about 100-200 progeny
  phage released per cell; entire cycle takes about 25 minutes
\end{enumerate}

\textbf{Bacteriophage lambda (\texorpdfstring{\ensuremath{\lambda}}{lambda}) lysogeny} represents a molecular
decision switch:

\begin{itemize}
\tightlist
\item
  After infection, \texorpdfstring{\ensuremath{\lambda}}{lambda} integrase catalyzes site-specific recombination
  between phage \emph{attP} and bacterial \emph{attB} sites, inserting
  the phage genome into the host chromosome
\item
  The \textbf{CI repressor} (\texorpdfstring{\ensuremath{\lambda}}{lambda} repressor) binds operator sequences,
  repressing lytic genes and maintaining lysogeny -- the prophage
  replicates passively with the host
\item
  Upon DNA damage (UV exposure, mitomycin C), the SOS response activates
  \textbf{RecA}, which stimulates CI repressor autocleavage,
  derepressing lytic genes and initiating the lytic cycle
\item
  Imprecise excision during induction can produce specialized
  transducing phage carrying \emph{gal} or \emph{bio} genes
\end{itemize}

\subsubsection{Animal Virus Replication}\label{animal-virus-replication}

Animal viruses follow a general replication strategy with virus-specific
variations:

\begin{enumerate}
\def\labelenumi{\arabic{enumi}.}
\tightlist
\item
  \textbf{Attachment}: Viral surface protein binds specific host
  receptor (tropism determinant)

  \begin{itemize}
  \tightlist
  \item
    Influenza: hemagglutinin (HA) binds sialic acid residues
  \item
    HIV: gp120 binds CD4 + CCR5 or CXCR4 coreceptor
  \item
    SARS-CoV-2: spike protein binds ACE2 receptor
  \end{itemize}
\item
  \textbf{Penetration}: Receptor-mediated
  \hyperref[gl:endocytosis]{\textbf{endocytosis}} (most non-enveloped
  viruses) or membrane fusion (enveloped viruses)
\item
  \textbf{Uncoating}: Capsid disassembly releases genome into cytoplasm
  (or nucleus)
\item
  \textbf{Biosynthesis}: Genome replication and protein synthesis
  (location depends on virus type)
\item
  \textbf{Assembly}: New virions assembled in cytoplasm or nucleus
\item
  \textbf{Release}: Budding through host membrane (enveloped viruses --
  acquiring lipid envelope) or cell lysis (non-enveloped viruses)
\end{enumerate}

\subsubsection{The HIV Life Cycle}\label{the-hiv-life-cycle}

HIV (human immunodeficiency virus) is a Class VI retrovirus with a
complex life cycle and multiple drug targets:

\begin{enumerate}
\def\labelenumi{\arabic{enumi}.}
\tightlist
\item
  \textbf{Attachment and fusion}: gp120 binds CD4 on T helper cells,
  then undergoes conformational change to bind coreceptor (CCR5 in early
  infection, CXCR4 in late infection); gp41 mediates membrane fusion
\item
  \textbf{Reverse transcription}: ssRNA genome -\textgreater{} dsDNA via
  reverse transcriptase (RT); RT lacks proofreading (\(\sim 10^{-4}\)
  errors per base per cycle), generating enormous genetic diversity
\item
  \textbf{Nuclear import}: Pre-integration complex enters nucleus
  through nuclear pores
\item
  \textbf{Integration}: Integrase inserts viral dsDNA into host
  chromosome, creating the \textbf{provirus} -- a permanent part of the
  host genome; preferentially integrates into actively transcribed
  regions
\item
  \textbf{Transcription and
  \hyperref[gl:translation]{\textbf{translation}}}: Host RNA Pol II
  transcribes viral mRNA; Tat protein enhances transcription 100-fold;
  Rev protein exports unspliced mRNA from nucleus
\item
  \textbf{Assembly and budding}: Gag and Gag-Pol polyproteins assemble
  at plasma membrane; budding acquires lipid envelope with gp120/gp41
  spikes
\item
  \textbf{Maturation}: Viral protease cleaves Gag-Pol polyprotein into
  functional proteins; immature virion becomes infectious
\end{enumerate}

\textbf{Antiretroviral therapy (ART)} targets multiple steps:

{\def\LTcaptype{none} % do not increment counter
\begin{longtable}[]{@{}
  >{\raggedright\arraybackslash}p{(\linewidth - 4\tabcolsep) * \real{0.3793}}
  >{\raggedright\arraybackslash}p{(\linewidth - 4\tabcolsep) * \real{0.2759}}
  >{\raggedright\arraybackslash}p{(\linewidth - 4\tabcolsep) * \real{0.3448}}@{}}
\toprule\noalign{}
\begin{minipage}[b]{\linewidth}\raggedright
Drug Class
\end{minipage} & \begin{minipage}[b]{\linewidth}\raggedright
Target
\end{minipage} & \begin{minipage}[b]{\linewidth}\raggedright
Examples
\end{minipage} \\
\midrule\noalign{}
\endhead
\bottomrule\noalign{}
\endlastfoot
NRTIs (nucleoside RT inhibitors) & Reverse transcriptase & Tenofovir,
emtricitabine, zidovudine (AZT) \\
NNRTIs (non-nucleoside RT inhibitors) & RT
\hyperref[gl:allosteric]{\textbf{allosteric}} site & Efavirenz,
rilpivirine \\
Protease inhibitors & Viral protease & Ritonavir, darunavir \\
Integrase inhibitors & Integrase & Dolutegravir, raltegravir \\
Entry inhibitors & CCR5 coreceptor & Maraviroc \\
Fusion inhibitors & gp41 & Enfuvirtide \\
\end{longtable}
}

\subsubsection{Prions}\label{prions}

\textbf{Prions} are infectious agents composed entirely of misfolded
protein -- they contain no nucleic acid. The normal cellular prion
protein (PrP\(^C\), predominantly \texorpdfstring{\ensuremath{\alpha}}{alpha}-helical) is converted to the
pathogenic scrapie form (PrP\(^{Sc}\), predominantly \texorpdfstring{\ensuremath{\beta}}{beta}-sheet) through
templated conformational change. PrP\(^{Sc}\) serves as a seed,
converting neighboring PrP\(^C\) molecules in a chain reaction that
produces amyloid fibrils.

Prion diseases (transmissible spongiform encephalopathies) include:

\begin{itemize}
\tightlist
\item
  \textbf{Creutzfeldt-Jakob disease (CJD)}: Sporadic (most common),
  familial, or iatrogenic
\item
  \textbf{Variant CJD (vCJD)}: Transmitted from bovine spongiform
  encephalopathy (BSE/``mad cow disease'') via contaminated beef
\item
  \textbf{Fatal familial insomnia}: Autosomal dominant; progressive
  insomnia leading to death
\item
  \textbf{Kuru}: Transmitted among the Fore people of Papua New Guinea
  through ritual endocannibalism; studied by Gajdusek (Nobel Prize 1976)
\item
  \textbf{Chronic wasting disease (CWD)}: Affects cervids (deer, elk);
  spreading across North America
\end{itemize}

Stanley Prusiner received the Nobel Prize in 1997 for his prion
hypothesis -- initially controversial because it challenged the central
dogma that most infectious agents require nucleic acid for replication.

\begin{quote}
\textbf{Clinical Connection: Bacteriophage Therapy} With the rise of
antibiotic-resistant infections, bacteriophage therapy is experiencing a
renaissance. The UC San Diego Center for Innovative Phage Applications
and Therapeutics (IPATH) has treated patients with life-threatening
multidrug-resistant infections using personalized phage cocktails. In
2016, Tom Patterson was rescued from a pan-resistant \emph{Acinetobacter
baumannii} infection using phages sourced from sewage. Engineered phages
can be designed to target specific bacterial species while sparing
beneficial \hyperref[gl:microbiota]{\textbf{microbiota}} -- a precision
antimicrobial approach unfeasible with conventional antibiotics.
\end{quote}

\subsubsection{Phage Therapy: A Pre-Antibiotic Idea Returns to the
Clinic}\label{phage-therapy-a-pre-antibiotic-idea-returns-to-the-clinic}

Phage therapy was discovered by Félix d'Hérelle in 1917, deployed widely
in the Soviet Union and Eastern Europe through the 20th century, and
largely abandoned in the West after the introduction of penicillin. The
antibiotic-resistance crisis has driven a global resurgence: as of 2024,
more than 30 phage therapy clinical trials are registered worldwide.

\textbf{Why phages are appealing as therapeutics:}

\begin{itemize}
\tightlist
\item
  \textbf{Self-amplification at the site of infection} --- every
  successful infection produces about 50--500 progeny, so dose
  increases where the target bacterium is densest, then declines as the
  infection clears.
\item
  \textbf{Narrow host range} --- most phages target a single species or
  strain, sparing the resident
  \hyperref[gl:microbiota]{\textbf{microbiota}} (a rare property;
  conventional broad-spectrum antibiotics decimate commensals).
\item
  \textbf{Co-evolution with bacteria} --- the same evolutionary force
  that creates antibiotic resistance also creates new phages; phage
  banks can be updated continuously.
\item
  \textbf{Engineerability} --- synthetic biology now enables host-range
  expansion, removal of toxin-encoding genes, addition of CRISPR-Cas
  payloads to specifically target resistance genes, and lytic/lysogenic
  switching.
\end{itemize}

\textbf{Clinical milestones (2016--2026):}

{\def\LTcaptype{none} % do not increment counter
\begin{longtable}[]{@{}
  >{\raggedright\arraybackslash}p{(\linewidth - 6\tabcolsep) * \real{0.1429}}
  >{\raggedright\arraybackslash}p{(\linewidth - 6\tabcolsep) * \real{0.4048}}
  >{\raggedright\arraybackslash}p{(\linewidth - 6\tabcolsep) * \real{0.2381}}
  >{\raggedright\arraybackslash}p{(\linewidth - 6\tabcolsep) * \real{0.2143}}@{}}
\toprule\noalign{}
\begin{minipage}[b]{\linewidth}\raggedright
Year
\end{minipage} & \begin{minipage}[b]{\linewidth}\raggedright
Patient / trial
\end{minipage} & \begin{minipage}[b]{\linewidth}\raggedright
Pathogen
\end{minipage} & \begin{minipage}[b]{\linewidth}\raggedright
Outcome
\end{minipage} \\
\midrule\noalign{}
\endhead
\bottomrule\noalign{}
\endlastfoot
2016 & Tom Patterson (UCSD/IPATH compassionate use) & Pan-resistant
\emph{Acinetobacter baumannii} & Recovery; first US success \\
2019 & Cystic fibrosis lung transplant rescue & \emph{Mycobacterium
abscessus} & Survival; engineered phage cocktail \\
2020--2024 & CYPHY, Phagoburn-2, Tailor-X & \emph{P. aeruginosa} burn /
UTI & Mixed efficacy; safety established \\
2023 & Compassionate-use registry \textgreater{} 200 patients & Multiple
ESKAPE & about 70 \% response in salvage settings \\
2024--26 & Locus Biosciences crPhage (CRISPR-Cas3 phage) & \emph{E.
coli} UTI & Phase II/III readouts \\
\end{longtable}
}

\textbf{Phage pharmacokinetics --- why they break the small-molecule
rules.} Conventional antibiotics follow first-order pharmacokinetics:
\(C(t) = C_0 e^{-k_e t}\), with linear AUC/dose relationships. Phages do
the opposite. The phage population grows where bacterial density is high
(auto-dosing) and crashes when the bacteria are gone (auto-clearance).
The simplest model is a Lotka--Volterra-like coupled ODE:

\[ \frac{dB}{dt} = rB - kBP, \qquad \frac{dP}{dt} = bkBP - mP  \label{eq:unit_VII_bacteria_archaea_viruses_item_8}\]

where \(B\) = bacterial density, \(P\) = phage density, \(r\) =
bacterial growth rate, \(k\) = adsorption rate constant (typically
\(10^{-9}\)--\(10^{-7}\) mL min\texorpdfstring{\textsuperscript{-}}{-}\texorpdfstring{\textsuperscript{1}}{1}), \(b\) = burst size (50--500), \(m\)
= phage clearance rate. The system has a \textbf{proliferation
threshold} at \(B^* = m / (bk)\): above \(B^*\), phage replicate faster
than they are cleared and amplify; below \(B^*\), they are diluted out.
This means low-dose phage \textbf{will not work} against a sub-threshold
infection --- clinicians must dose above the minimum proliferating
concentration. It also means a single phage dose can clear
orders-of-magnitude more bacteria than its initial titer, making
``dose'' a misleading concept.

\textbf{Pharmacokinetic and regulatory challenges that distinguish
phages from small molecules:}

\begin{enumerate}
\def\labelenumi{\arabic{enumi}.}
\tightlist
\item
  \textbf{Bacterial resistance is rapid} --- phage receptors mutate at
  \(10^{-6}\)--\(10^{-8}\) per cell per generation. Cocktails of 3--5
  phages targeting non-overlapping receptors slow this; estimated escape
  probability \(\approx \prod_i 10^{-7} = 10^{-21}\) for a 3-phage
  cocktail with independent receptors.
\item
  \textbf{Immune neutralization} --- IgG against phage capsid emerges
  within about 10 days, lowering blood titres and limiting
  repeated systemic dosing.
\item
  \textbf{Pharmacokinetics are non-linear} --- phage concentration
  depends on bacterial density (replication \textgreater{} clearance
  below threshold; opposite above) --- classical AUC/MIC concepts do not
  apply.
\item
  \textbf{Manufacturing \& regulation} --- purified, endotoxin-free
  phage preparations require GMP processes designed for biologicals; the
  FDA has issued personalized/compassionate-use pathways while broad
  regulatory frameworks are still developing.
\item
  \textbf{Lysogenic phages can transduce virulence/resistance genes}
  (Shiga toxin in STEC came from a lambdoid phage); therapeutic phages
  must be obligately lytic and toxin-free.
\end{enumerate}

The pragmatic clinical role emerging by 2026 is \textbf{adjunctive phage
therapy} --- phages combined with antibiotics for biofilm-associated and
multi-drug-resistant infections (prosthetic joints, CF lungs,
complicated UTIs, infective endocarditis), where phage-driven dispersal
of biofilm matrix re-sensitizes persisters to antibiotics and the
antibiotic suppresses phage-resistant escapees. This combination logic
mirrors HIV ART and TB DOTS --- the lesson that single-agent
antimicrobial therapy invites resistance is now being applied to phages
too.

\begin{quote}
\textbf{Concept Check 4:} HIV reverse transcriptase lacks proofreading
activity, producing approximately one
\hyperref[gl:mutation]{\textbf{mutation}} per genome per replication
cycle. Explain why this high mutation rate is both advantageous for the
virus (immune evasion, drug resistance) and why it also constrains the
maximum genome size of retroviruses compared to DNA viruses like
herpesviruses (about 150 kb genome).
\end{quote}

\begin{quote}
\textbf{Concept Check 4b:} A patient with chronic \emph{Pseudomonas
aeruginosa} lung infection in cystic fibrosis is treated with a phage
cocktail. Within 5 days, the bacterial load drops 100-fold, but by day
14, \emph{P. aeruginosa} has rebounded with mutations in \emph{galU}
(LPS biosynthesis). Explain (a) why galU mutants escape the phage, (b)
why the rebound strain may actually be less fit \emph{in vivo} than the
parent, and (c) why combining phages with tobramycin can suppress this
rebound.
\end{quote}

\begin{center}\rule{0.5\linewidth}{0.5pt}\end{center}

\subsection{Antibiotic Mechanisms and
Resistance}\label{antibiotic-mechanisms-and-resistance}

\subsubsection{Antibiotic Classes by
Mechanism}\label{antibiotic-classes-by-mechanism}

Antibiotics exploit the structural and biochemical differences between
prokaryotic and eukaryotic cells. Each class targets a specific
essential process:

\begin{figure}[htbp]
\centering
\includegraphics[width=0.82\linewidth,height=4.2in,keepaspectratio]{../figures/mermaid_inline/inline_mermaid_0120_6f48da85c781.png}
\caption{Antibiotic targets organised by cellular process (cell-wall synthesis, protein synthesis, nucleic-acid synthesis, membrane disruption) alongside the principal resistance mechanisms.}
\end{figure}

\subsubsection{Selective Toxicity}\label{selective-toxicity}

The foundation of antibiotic therapy is \textbf{selective toxicity} --
targeting structures present in bacteria but absent or sufficiently
different in human cells:

\begin{itemize}
\tightlist
\item
  Peptidoglycan (absent in human cells) -\textgreater{} beta-lactams,
  vancomycin
\item
  70S ribosomes (vs.~human 80S) -\textgreater{} aminoglycosides,
  tetracyclines, macrolides
\item
  Bacterial DNA gyrase (vs.~human topoisomerase II) -\textgreater{}
  fluoroquinolones
\item
  Bacterial RNA polymerase (structurally distinct from human Pol II)
  -\textgreater{} rifampin
\item
  LPS outer membrane (absent in human cells) -\textgreater{} polymyxins
\end{itemize}

\subsubsection{Resistance Mechanisms in
Detail}\label{resistance-mechanisms-in-detail}

\textbf{Beta-lactamases} are the most clinically important resistance
mechanism. These enzymes hydrolyze the beta-lactam ring, inactivating
the antibiotic before it reaches its PBP target. Evolution of
beta-lactamases tracks the history of antibiotic development:

\begin{itemize}
\tightlist
\item
  Penicillinase (TEM-1, SHV-1) -\textgreater{} resistance to penicillins
\item
  Extended-spectrum beta-lactamases (ESBLs: CTX-M, TEM variants)
  -\textgreater{} resistance to 3rd-generation cephalosporins
\item
  Carbapenemases (KPC, NDM-1, OXA-48) -\textgreater{} resistance to
  carbapenems (last-resort beta-lactams)
\end{itemize}

\textbf{MRSA} (\emph{methicillin-resistant Staphylococcus aureus})
carries the \textbf{mecA} or \textbf{mecC} gene encoding PBP2a/PBP2c,
altered penicillin-binding proteins with low affinity for most
traditional beta-lactam antibiotics. The gene resides on SCC\emph{mec}
(staphylococcal cassette chromosome), acquired by horizontal gene
transfer. Treatment depends on syndrome and susceptibility: vancomycin,
daptomycin, linezolid, or ceftaroline/ceftobiprole in settings where
anti-MRSA cephalosporins are appropriate.

\textbf{Efflux pumps} actively transport antibiotics out of the cell.
The AcrAB-TolC system in \emph{E. coli} and MexAB-OprM in
\emph{Pseudomonas aeruginosa} are tripartite pumps spanning the inner
membrane, periplasm, and outer membrane, conferring resistance to
multiple drug classes simultaneously.

\begin{figure}[htbp]
\centering
\includegraphics[width=0.82\linewidth,height=4.2in,keepaspectratio]{../figures/mermaid_inline/inline_mermaid_0121_3ff5e572147d.png}
\caption{AMR horizontal-gene-transfer map. Resistance spreads by both clonal expansion and mobile genetic elements, so stewardship must be paired with infection control, wastewater/environmental surveillance, and organism-resistance-pair reporting \textbackslash{}citep\{who2024bppl,cdc2025antibioticuse\}.}
\end{figure}

\subsubsection{WHO Priority Pathogens
(ESKAPE)}\label{who-priority-pathogens-eskape}

The WHO 2024 Bacterial Priority Pathogens List keeps AMR triage grounded
in public-health burden, resistance trend, transmissibility,
treatability, and pipeline scarcity \citep{who2024bppl}. The
\textbf{ESKAPE} mnemonic remains a useful bedside memory aid, but the
WHO list is broader: it separately prioritises carbapenem-resistant
\emph{Acinetobacter baumannii}, carbapenem- or
third-generation-cephalosporin-resistant Enterobacterales,
drug-resistant \emph{Mycobacterium tuberculosis}, and other
pathogen-resistance pairs.

{\def\LTcaptype{none} % do not increment counter
\begin{longtable}[]{@{}
  >{\raggedright\arraybackslash}p{(\linewidth - 4\tabcolsep) * \real{0.2381}}
  >{\raggedright\arraybackslash}p{(\linewidth - 4\tabcolsep) * \real{0.3571}}
  >{\raggedright\arraybackslash}p{(\linewidth - 4\tabcolsep) * \real{0.4048}}@{}}
\toprule\noalign{}
\begin{minipage}[b]{\linewidth}\raggedright
Pathogen
\end{minipage} & \begin{minipage}[b]{\linewidth}\raggedright
Key Resistance
\end{minipage} & \begin{minipage}[b]{\linewidth}\raggedright
Clinical Setting
\end{minipage} \\
\midrule\noalign{}
\endhead
\bottomrule\noalign{}
\endlastfoot
\emph{Enterococcus faecium} & VRE (vancomycin-resistant) & Bloodstream
infections, UTIs \\
\emph{Staphylococcus aureus} & MRSA (methicillin-resistant) & Skin, soft
tissue, bacteremia \\
\emph{Klebsiella pneumoniae} & ESBL, KPC carbapenemase & Pneumonia,
UTIs, sepsis \\
\emph{Acinetobacter baumannii} & Pan-resistant; OXA carbapenemases &
Ventilator-associated pneumonia \\
\emph{Pseudomonas aeruginosa} & Intrinsic resistance; efflux pumps &
Burn wounds, CF lung infections \\
\emph{Enterobacter} spp. & AmpC beta-lactamase (inducible) & Nosocomial
infections \\
\end{longtable}
}

\begin{quote}
\textbf{Clinical Connection: The Antibiotic Resistance Crisis} The WHO
has declared antimicrobial resistance one of the top 10 global public
health threats. An estimated 1.27 million deaths were directly
attributable to bacterial AMR in 2019, with many more associated deaths
\citep{murray2022amr}. The O'Neill review's 10-million-deaths-per-year
scenario remains a warning about unchecked resistance rather than a
forecast that must occur \citep{oneill2016amr}. Major drivers include
unnecessary human antibiotic prescribing, incomplete access to
diagnostics, antibiotic use in food-animal production, and a thin
discovery pipeline in which truly novel antibacterial classes are rare
rather than absent \citep{cdc2025antibioticuse,who2024bppl}.
\end{quote}

\begin{quote}
\textbf{Concept Check 5:} A hospital isolate of \emph{Klebsiella
pneumoniae} is resistant to most beta-lactams including carbapenems,
aminoglycosides, and fluoroquinolones. Primarily colistin (polymyxin)
remains effective. Explain why colistin resistance, which has recently
emerged via the plasmid-borne mcr-1 gene (encoding a phosphoethanolamine
transferase that modifies lipid A), is particularly alarming from a
public health perspective.
\end{quote}

\begin{center}\rule{0.5\linewidth}{0.5pt}\end{center}

\subsection{CRISPR-Cas: Bacterial Adaptive
Immunity}\label{crispr-cas-bacterial-adaptive-immunity}

\textbf{CRISPR} (Clustered Regularly Interspaced Short Palindromic
Repeats) arrays provide bacteria and archaea with an adaptive, heritable
immune system against phages and plasmids. Discovered in 2007 (Barrangou
\emph{et al.}, \emph{Science} 2007), the system has three functional
stages --- \textbf{adaptation, expression (biogenesis), and
interference} --- that map cleanly onto the three steps of any immune
system: memory formation, surveillance, and effector action.

\begin{figure}[htbp]
\centering
\includegraphics[width=0.82\linewidth,height=4.2in,keepaspectratio]{../figures/mermaid_inline/inline_mermaid_0122_5973a9d90dc7.png}
\caption{The three-stage CRISPR-Cas immune mechanism: adaptation (Cas1-Cas2 captures and integrates new spacers), expression (pre-crRNA processed and loaded onto effector), and interference (target dsDNA cleavage).}
\end{figure}

\subsubsection{Stage 1 --- Spacer Acquisition
(Adaptation)}\label{stage-1-spacer-acquisition-adaptation}

When a bacterium survives phage infection, the \textbf{Cas1-Cas2
integrase complex} captures a short fragment (about 30 bp) of
phage DNA --- called a \textbf{protospacer} --- adjacent to a
\textbf{PAM} (protospacer adjacent motif). Cas1-Cas2 integrates this as
a new \textbf{spacer} at the leader end of the CRISPR array between two
direct repeats. This constitutes an immunological memory: the spacer
sequence is now permanently encoded in the bacterial chromosome and
inherited by most daughter cells.

\begin{equation}
\text{Cas1-Cas2} + \text{protospacer-PAM} \rightarrow \text{pre-spacer} \xrightarrow{\text{integration}} \text{CRISPR array expanded}
\label{eq:bacteria_archaea_viruses_2}
\end{equation}

\textbf{PAM recognition} ensures self vs non-self discrimination: the
PAM sequence is not present in the CRISPR array itself (repeat sequences
flank spacers, not PAMs), so the bacterium does not attack its own
CRISPR locus. The specific PAM depends on the system: NGG (SpCas9, Type
II), 5'-AAT-3' or 5'-CTT-3' (Type I-E), etc.

\subsubsection{Stage 2 --- crRNA
Biogenesis}\label{stage-2-crrna-biogenesis}

The CRISPR array is transcribed as a long pre-CRISPR RNA (pre-crRNA).
Repeat sequences are cleaved by:

\begin{itemize}
\tightlist
\item
  \textbf{Type I and III}: The \emph{Cas6} endoribonuclease (recognizes
  stem-loop structure in repeats)
\item
  \textbf{Type II}: RNase III + \emph{tracrRNA} (trans-activating crRNA)
  hybridises with the repeat; Cas9 is then loaded with a single-guide
  RNA (sgRNA = crRNA + tracrRNA scaffold)
\end{itemize}

\subsubsection{Stage 3 --- Interference
(Targeting)}\label{stage-3-interference-targeting}

The crRNA guides the effector complex to matching sequences in foreign
nucleic acid. The CRISPR-Cas systems are now classified into 6 types
(I--VI) and over 30 subtypes; the three most important types are
summarised below:

{\def\LTcaptype{none} % do not increment counter
\begin{longtable}[]{@{}
  >{\raggedright\arraybackslash}p{(\linewidth - 6\tabcolsep) * \real{0.1750}}
  >{\raggedright\arraybackslash}p{(\linewidth - 6\tabcolsep) * \real{0.2500}}
  >{\raggedright\arraybackslash}p{(\linewidth - 6\tabcolsep) * \real{0.2750}}
  >{\raggedright\arraybackslash}p{(\linewidth - 6\tabcolsep) * \real{0.3000}}@{}}
\toprule\noalign{}
\begin{minipage}[b]{\linewidth}\raggedright
Feature
\end{minipage} & \begin{minipage}[b]{\linewidth}\raggedright
\textbf{Type I}
\end{minipage} & \begin{minipage}[b]{\linewidth}\raggedright
\textbf{Type II}
\end{minipage} & \begin{minipage}[b]{\linewidth}\raggedright
\textbf{Type III}
\end{minipage} \\
\midrule\noalign{}
\endhead
\bottomrule\noalign{}
\endlastfoot
Effector & Multi-subunit \textbf{Cascade} complex (Cas5--Cas8) + Cas3
helicase-nuclease & Single \textbf{Cas9} nuclease &
\textbf{Cas10}-Csm/Cmr complex \\
Target & dsDNA & dsDNA & ssRNA + transcribed DNA \\
Cleavage mechanism & Cas3 unwinds and \textbf{degrades dsDNA
processively} & Cas9 makes blunt-ended \textbf{double-strand break} 3 bp
upstream of PAM & Co-transcriptional ssRNA cleavage by Csm/Cmr;
signalling output \\
PAM requirement & Yes (recognised by Cas8 of Cascade) & Yes (NGG for
SpCas9; varies by ortholog) & None for RNA target;
protospacer-flanking-sequence (PFS) for DNA \\
Guide RNA & crRNA primarily & crRNA + tracrRNA (or fused sgRNA) & crRNA
primarily \\
Signal molecule & --- & --- & \textbf{cyclic oligoadenylate (cA\texorpdfstring{\textsubscript{4}}{4}, cA\texorpdfstring{\textsubscript{6}}{6})}
activates non-specific RNase Csm6 \\
Genome distribution & about  50 \% of CRISPR-bearing genomes
& about  5 \% & about  25 \% \\
Subtypes & I-A through I-G & II-A, II-B, II-C & III-A through III-F \\
Anti-CRISPR susceptibility & Yes (AcrIF, AcrIE families) & Yes (AcrIIA,
AcrIIC families) & Yes (AcrIIIA, anti-cA\texorpdfstring{\textsubscript{4}}{4}) \\
Clinical / biotech use & Cas3 large-deletion editing; phage therapy
design (Locus crPhage) & \textbf{Genome editing (Cas9)}; base editors;
prime editors & Biosensing (cA\texorpdfstring{\textsubscript{4}}{4} readout); SHERLOCK-style RNA
detection \\
\end{longtable}
}

\textbf{Type II (Cas9)} was repurposed for genome editing by Doudna,
Charpentier, and colleagues in 2012 (\emph{Science} 2012), work
recognised with the 2020 Nobel Prize in Chemistry. A synthetic sgRNA
directs Cas9 to any genomic target (specified by a 20-nt spacer sequence
+ NGG PAM), where it makes a site-specific DSB. The cell repairs the
break by NHEJ (creating indels --- gene knockout) or HDR (with a
provided template --- precise edit). Variants now include base editors
(cytosine \texorpdfstring{\ensuremath{\to}}{->} thymine, adenine \texorpdfstring{\ensuremath{\to}}{->} guanine without DSBs), prime editors
(programmable insertions/deletions guided by an extended pegRNA), and
Cas12/Cas13 systems for DNA/RNA targeting respectively. The same
RNA-guided programmability that powers genome editing also underlies the
\textbf{RNA-interference-like} specificity by which bacteria neutralize
phage genomes --- a striking parallel to eukaryotic small-RNA defense
systems first dissected with C. elegans \citep{fire1998}.

\subsubsection{Anti-CRISPR Proteins}\label{anti-crispr-proteins}

Phages have evolved \textbf{anti-CRISPR (Acr)} proteins that counteract
bacterial CRISPR immunity --- an evolutionary arms race. Over 50 Acr
families have been identified (Pawluk \emph{et al.}, \emph{Nature
Microbiology} 2016; Davidson \emph{et al.}, \emph{Science} 2020 cryo-EM
structures):

\begin{itemize}
\tightlist
\item
  \textbf{AcrIF1 (Type I-F inhibitor)}: Binds Cas8f subunit of Cascade;
  blocks crRNA-guided dsDNA binding
\item
  \textbf{AcrIIA2/IIA4 (Type II inhibitors)}: Bind and occlude the
  PAM-recognition domain of Cas9; prevents target engagement
\item
  \textbf{AcrIIIA4}: Degrades the cA\texorpdfstring{\textsubscript{4}}{4} second messenger; blocks Csm6
  activation
\end{itemize}

\begin{quote}
\textbf{Clinical Connection: CRISPR-Based Therapies (2023--2026)} The
first CRISPR-based medicines have now been approved. \textbf{Casgevy}
(exagamglogene autotemcel), developed by Vertex and CRISPR Therapeutics,
received FDA approval in December 2023 for sickle cell disease and
January 2024 approval for transfusion-dependent \texorpdfstring{\ensuremath{\beta}}{beta}-thalassemia; FDA's
2026 product page lists both indications for patients 12 years of age
and older
\citep{fda2023casgevy,fda2024casgevythalassemia,fda2026casgevy}. The ex
vivo approach edits patient HSCs to reactivate fetal haemoglobin (HbF)
by disrupting the BCL11A erythroid enhancer, compensating for defective
adult haemoglobin. Simultaneously, Bluebird Bio's lentiviral gene
therapy Lyfgenia was approved for SCD. By 2026, CRISPR therapies for
transthyretin amyloidosis (intellia NTLA-2001, in vivo hepatic editing)
and \hyperref[gl:heterozygous]{\textbf{heterozygous}} familial
hypercholesterolaemia are in Phase III trials.
\end{quote}

\textbf{Concept Check 6:} \textgreater{} A Type II CRISPR-Cas9 system
fails to cleave a target sequence despite perfect spacer-target
complementarity. List three molecular explanations for this failure,
relating each to a specific step in the CRISPR mechanism (adaptation,
biogenesis, or interference).

\begin{quote}
\textbf{Concept Check 6b:} A bacterial population is challenged with a
phage cocktail. After 5 generations, surviving bacteria are sequenced
--- 70\% have new spacers in their CRISPR array matching the phage; 20\%
have escape mutations in essential phage genes (no new spacer); 10\%
have constitutively expressed Cas9 with broader PAM specificity. Explain
which adaptation strategy is most evolutionary stable, and predict the
population dynamics if a second, immunologically distinct phage is then
introduced.
\end{quote}

\begin{center}\rule{0.5\linewidth}{0.5pt}\end{center}

\subsection{Computational Bridge}\label{computational-bridge-21}

Batch culture kinetics follow the logistic-like curve encoded in
\breaktt{bacterial\_growth\_curve}:

\begin{Shaded}
\begin{Highlighting}[]
\ImportTok{from}\NormalTok{ biology.microbiology }\ImportTok{import}\NormalTok{ bacterial\_growth\_curve}

\NormalTok{curve }\OperatorTok{=}\NormalTok{ bacterial\_growth\_curve(N0}\OperatorTok{=}\FloatTok{1e5}\NormalTok{, doubling\_time\_hr}\OperatorTok{=}\FloatTok{0.5}\NormalTok{, t\_end\_hr}\OperatorTok{=}\FloatTok{6.0}\NormalTok{)}
\BuiltInTok{print}\NormalTok{(curve.populations[}\OperatorTok{{-}}\DecValTok{1}\NormalTok{] }\OperatorTok{\textgreater{}}\NormalTok{ curve.populations[}\DecValTok{0}\NormalTok{])}
\end{Highlighting}
\end{Shaded}

\begin{quote}
\textbf{Clinical / systems note:} Time-kill curves and \(\mathrm{MIC}\)
testing in the clinic are empirical cousins of these growth models;
combination therapy is designed to keep effective populations below
invasion thresholds.
\end{quote}

\begin{center}\rule{0.5\linewidth}{0.5pt}\end{center}

\subsubsection{SARS-CoV-2 Variant Evolution: Real-Time Darwin in a
Global
Population}\label{sars-cov-2-variant-evolution-real-time-darwin-in-a-global-population}

The four-year trajectory of SARS-CoV-2 (late 2019 -- present) is the
best-documented example of \textbf{adaptive viral evolution} in
biological history. Public sequencing through GISAID now archives
\textgreater{} 16 million viral genomes, enabling phylogenetic tracking
at daily resolution. The lineage diversification follows a textbook
evolutionary pattern: \textbf{(1) a narrow founder bottleneck}
(Wuhan-Hu-1 reference, December 2019); \textbf{(2) geographic diffusion}
with local \hyperref[gl:founder-effect]{\textbf{founder effect}}s (D614G
rose to global dominance by June 2020, conferring about 20 \%
higher infectivity via an allosteric receptor-binding-domain (RBD)
``open'' state); \textbf{(3) escape-driven positive selection} as
population immunity accumulated, favouring RBD mutations that evade
neutralising antibodies without compromising ACE2 binding.

Five variants of concern (VOCs) illustrate distinct evolutionary forces.
\textbf{Alpha (B.1.1.7, late 2020)} carried N501Y --- a single RBD
mutation increasing ACE2 affinity about 10\texorpdfstring{\ensuremath{\times}}{x}. \textbf{Beta
(B.1.351)} and \textbf{Gamma (P.1)} added E484K, which partially escapes
antibodies. \textbf{Delta (B.1.617.2, mid-2021)} carried L452R + T478K
and doubled transmissibility. \textbf{Omicron (B.1.1.529, November
2021)} was the punctuated-evolution event of the pandemic --- \textbf{32
spike mutations at once}, not a gradual lineage from Delta but a sudden
saltation, likely from a long-term infection in an immunocompromised
host (a natural ``mutator'' environment). Omicron escaped most
pre-existing neutralising antibodies but attenuated lung tropism via
reduced TMPRSS2-mediated entry, shifting the disease toward upper
respiratory infection. Subsequent sub-lineages (BA.2, BA.5, XBB, JN.1,
KP.3) follow continuous antigenic drift --- the \textbf{``influenza-like
future''} of SARS-CoV-2 is now the expected trajectory.

Quantitatively: \textbf{Ka/Ks ratios} on the RBD exceed 1.5 across the
pandemic (positive selection), versus about 0.1 for
housekeeping ORFs (purifying selection); \textbf{estimated substitution
rate} is about 2 \texorpdfstring{\ensuremath{\times}}{x} 10\texorpdfstring{\textsuperscript{-}}{-}\texorpdfstring{\textsuperscript{3}}{3} per site per year (2\texorpdfstring{\ensuremath{\times}}{x} SARS-CoV-1),
consistent with RNA-virus fidelity limits and increased by
population-scale chronic infections. Public-health surveillance now runs
phylogenetic tools (Nextstrain, Pango) in near-real-time;
\hyperref[gl:vaccine]{\textbf{vaccine}} reformulation (WHO TAG-CO-VAC)
follows the influenza-flu-shot paradigm. The pandemic has transformed
biology teaching: evolution is no longer a fossil-record inference but a
live, observable process with weekly GenBank updates.

\begin{center}\rule{0.5\linewidth}{0.5pt}\end{center}

\subsection{Current Evidence and Frontier
Biology}\label{current-evidence-and-frontier-biology-24}

For \textbf{Bacteria, Archaea, and Viruses}, frontier biology belongs
inside the evidence logic of the chapter. Microbiology and infectious
disease now require One Health reasoning across people, animals,
environments, genomics, and antimicrobial stewardship. The core reading
question is this: microbial claims should identify taxonomy, genome
architecture, metabolism, resistance mechanism, and environment.

\begin{itemize}
\tightlist
\item
  \textbf{What to verify:} identify the observation, model, assay, or
  dataset that would make the claim stronger or weaker.
\item
  \textbf{What to qualify:} state the scale, organism, cell type,
  environmental condition, or population where the claim is expected to
  hold.
\item
  \textbf{What to compare:} test at least one alternative explanation,
  baseline, or null model before treating the pattern as causal.
\item
  \textbf{What to cite:} distinguish primary evidence, review synthesis,
  public dataset, and institutional guidance; for recent or numeric
  claims, prefer the source closest to the measurement and state what
  has changed since it was published.
\end{itemize}

For AMR and pathogen claims, name the organism-resistance pair, the
selection pressure, the transmission route, and the surveillance
evidence that would change triage
\citep{who2024bppl,cdc2025antibioticuse,murray2022amr}.

\textbf{Source practice:} For pathogen, AMR, and intervention claims,
tie statements to organism-resistance pairs, surveillance evidence,
official guidance, and trial/regulatory status
\citep{who2024bppl,who2025tb,who2025malaria,cdc2025lenacapavirprep,cdc2026candidaauris}.

\subsubsection{Current Evidence Map: AMR Movement Across One
Health}\label{current-evidence-map-amr-movement-across-one-health}

\begin{figure}[htbp]
\centering
\includegraphics[width=0.82\linewidth,height=4.2in,keepaspectratio]{../figures/mermaid_inline/inline_mermaid_0123_ddc4b7e43f0b.png}
\caption{WHO priority lists are most useful when students connect the organism-resistance pair to selection pressure, transmission route, and stewardship action \textbackslash{}citep\{who2024bppl,murray2022amr\}.}
\end{figure}

\subsection{Summary}\label{summary-24}

\begin{itemize}
\tightlist
\item
  \textbf{Three domains of life} (Bacteria, Archaea, Eukarya) were
  established by Woese's rRNA phylogeny \citep{woese1977}. Bacteria and
  Archaea share the prokaryotic cell plan but differ fundamentally in
  membrane lipids, cell wall chemistry, and transcriptional machinery.
\item
  \textbf{Bacterial cell walls} contain peptidoglycan, built in an
  11-step pathway (MurA \texorpdfstring{\ensuremath{\to}}{->} MurB \texorpdfstring{\ensuremath{\to}}{->} MurC--F \texorpdfstring{\ensuremath{\to}}{->} MraY \texorpdfstring{\ensuremath{\to}}{->} MurG \texorpdfstring{\ensuremath{\to}}{->} MurJ flipping \texorpdfstring{\ensuremath{\to}}{->}
  transglycosylation \texorpdfstring{\ensuremath{\to}}{->} PBP transpeptidation) that is the target of
  fosfomycin, D-cycloserine, bacitracin, vancomycin, daptomycin, and
  \texorpdfstring{\ensuremath{\beta}}{beta}-lactams. Gram-positive bacteria have thick peptidoglycan with
  teichoic acids; Gram-negative bacteria have thin peptidoglycan plus an
  outer membrane with LPS (endotoxin, activates TLR4, causes septic
  shock).
\item
  \textbf{Bacterial growth equation} \(N(t) = N_0 \cdot 2^{t/t_d}\)
  (\cref{eq:unit_VII_bacterial_growth}) describes exponential phase; the
  natural-log form gives a linear plot whose slope is the specific
  growth rate \(\mu = (\ln 2)/t_d\).
\item
  \textbf{Bacterial metabolic diversity} exceeds most eukaryotes:
  photoautotrophs, chemoautotrophs (lithotrophs), and heterotrophs with
  varied oxygen relationships.
\item
  \textbf{Two-component systems} (sensor histidine kinase + response
  regulator, His\texorpdfstring{\ensuremath{\to}}{->}Asp phosphotransfer) are the dominant prokaryotic
  signal-transduction architecture; \textbf{biofilms} (matrix +
  persisters, c-di-GMP master regulator) and \textbf{quorum sensing}
  (AHL/AIP/AI-2) underpin chronic infections and 65--80 \% of clinical
  infections.
\item
  \textbf{Horizontal gene transfer} (transformation, transduction,
  conjugation) drives bacterial evolution \citep{lawrence1998} and
  antibiotic resistance spread via R-plasmids and transposons.
\item
  \textbf{Archaea} differ from bacteria in ether-linked isoprenoid
  lipids (with G-1-P stereochemistry; tetraether monolayers in
  hyperthermophiles), lack of peptidoglycan, and eukaryote-like
  transcription. \textbf{Asgard archaea} are the closest living
  relatives of eukaryotes and encode actin, ESCRT, and ubiquitin-like
  machinery.
\item
  \textbf{Viruses} (Baltimore classes I-VII) are obligate intracellular
  parasites; bacteriophage lytic/lysogenic cycles and the HIV life cycle
  illustrate replication strategies and drug targets. \textbf{Phage
  therapy} is re-emerging as a clinical adjunct for multi-drug-resistant
  infections, with engineering, cocktail design, and antibiotic
  combination addressing rapid resistance and immune neutralization
  \citep{doub2022phagepseudomonas}. Phage pharmacokinetics are
  non-linear with a proliferation threshold \(B^* = m/(bk)\).
\item
  \textbf{CRISPR-Cas} is a heritable bacterial adaptive immune system
  (adaptation--expression--interference); Types I (Cas3), II (Cas9), and
  III (Cas10) differ in effector architecture, target nucleic acid, and
  PAM requirements. Type II Cas9 has been repurposed as the dominant
  genome-editing technology and now drives FDA-approved therapeutics
  (Casgevy, 2023--2026 FDA records).
\item
  \textbf{Antibiotic resistance} mechanisms (enzymatic inactivation,
  target modification, efflux, reduced permeability) spread via HGT;
  ESKAPE organisms remain a useful mnemonic, but current AMR triage
  should use pathogen-resistance pairs and WHO BPPL priority categories
  rather than a single threat list.
\item
  \textbf{Connections:} See \cref{sec:unit_VI_phylogenetics} for
  phylogeny of mobile elements, \cref{sec:unit_VII_infectious_disease}
  for host immunity, and Unit III for metabolic targets of antibiotics.
\end{itemize}

\begin{center}\rule{0.5\linewidth}{0.5pt}\end{center}

\subsection{Key Terms}\label{key-terms-24}

{\def\LTcaptype{none} % do not increment counter
\begin{longtable}[]{@{}
  >{\raggedright\arraybackslash}p{(\linewidth - 2\tabcolsep) * \real{0.3529}}
  >{\raggedright\arraybackslash}p{(\linewidth - 2\tabcolsep) * \real{0.6471}}@{}}
\toprule\noalign{}
\begin{minipage}[b]{\linewidth}\raggedright
Term
\end{minipage} & \begin{minipage}[b]{\linewidth}\raggedright
Definition
\end{minipage} \\
\midrule\noalign{}
\endhead
\bottomrule\noalign{}
\endlastfoot
\textbf{Peptidoglycan} & Bacterial cell wall polymer of NAG-NAM glycan
strands cross-linked by tetrapeptide chains; target of beta-lactam
antibiotics \\
\textbf{Lipid II} & Lipid-anchored peptidoglycan precursor
(undecaprenyl-PP-MurNAc-pentapeptide-GlcNAc); target of vancomycin,
nisin, daptomycin \\
\textbf{Lipopolysaccharide (LPS)} & Gram-negative outer membrane
component (lipid A + core + O-antigen); endotoxin that activates TLR4
and can cause septic shock \\
\textbf{Gram stain} & Differential staining technique distinguishing
thick-walled (Gram-positive, purple) from thin-walled (Gram-negative,
pink) bacteria \\
\textbf{Two-component system (TCS)} & Sensor histidine kinase + response
regulator pair that converts an environmental stimulus into a
transcriptional output via His\texorpdfstring{\ensuremath{\to}}{->}Asp phosphotransfer \\
\textbf{Biofilm} & Surface-attached, EPS-encased microbial community;
10--1000\texorpdfstring{\ensuremath{\times}}{x} more antibiotic tolerant than planktonic cells \\
\textbf{c-di-GMP} & Cyclic-di-guanosine-monophosphate; master second
messenger that switches between motile (low) and biofilm (high)
lifestyles \\
\textbf{Quorum sensing} & Cell-density-dependent gene regulation via
secreted autoinducer molecules (AHL, AIP, AI-2) \\
\textbf{Persister cells} & Metabolically dormant subpopulations within a
biofilm that survive antibiotic exposure without genetic resistance \\
\textbf{Binary fission} & Asexual bacterial reproduction; chromosome
replication followed by septum formation and cell division \\
\textbf{Specific growth rate (\texorpdfstring{\ensuremath{\mu}}{mu})} & \(\mu = (\ln 2)/t_d\); instantaneous
fractional growth rate in exponential phase, units of
\(\text{time}^{-1}\) \\
\textbf{Horizontal gene transfer} & Transfer of genetic material between
organisms by mechanisms other than vertical (parent-to-offspring)
inheritance \\
\textbf{Transformation} & Uptake of free DNA from the environment by
competent bacteria \\
\textbf{Transduction} & Bacteriophage-mediated transfer of bacterial DNA
between cells \\
\textbf{Conjugation} & Direct cell-to-cell DNA transfer via F pilus;
primary mechanism for R-plasmid spread \\
\textbf{Endospore} & Dormant, highly resistant structure formed by
Bacillus and Clostridium; contains calcium-dipicolinic acid \\
\textbf{Bacteriophage} & Virus that infects bacteria; follows lytic
(cell lysis) or lysogenic (genome integration) cycle \\
\textbf{Lysogeny} & Integration of phage DNA into host chromosome as a
prophage; maintained by CI repressor \\
\textbf{Phage therapy} & Therapeutic use of bacteriophages (often as
cocktails) to treat bacterial infections, especially
multi-drug-resistant ones \\
\textbf{Phage proliferation threshold} & Minimum bacterial density
\(B^* = m/(bk)\) above which a phage population can self-amplify rather
than dilute out \\
\textbf{Baltimore classification} & System classifying viruses into
seven classes based on genome type and replication strategy \\
\textbf{Reverse transcriptase} & RNA-dependent DNA polymerase; used by
retroviruses (HIV) to convert RNA genome to dsDNA \\
\textbf{Prion} & Infectious misfolded protein (PrP\(^{Sc}\)) that
templates conversion of normal PrP\(^C\); no nucleic acid \\
\textbf{MRSA} & Methicillin-resistant \emph{Staphylococcus aureus};
carries mecA gene encoding PBP2a with low beta-lactam affinity \\
\textbf{ESKAPE pathogens} & Six priority antibiotic-resistant organisms:
Enterococcus, S. aureus, Klebsiella, Acinetobacter, Pseudomonas,
Enterobacter \\
\textbf{Asgard archaea} & Deep-branching archaeal lineages (Loki, Thor,
Odin, Heimdall, Hela, Hermod) containing eukaryotic signature proteins;
closest known relatives of eukaryotes \\
\textbf{Eukaryotic signature proteins (ESPs)} & Proteins (actin, ESCRT,
ubiquitin-like, Rab GTPases) encoded by Asgard archaea that were once
thought to be eukaryote-specific \\
\textbf{GDGT lipids} & Glycerol-dialkyl-glycerol-tetraether monolayer
lipids that span the entire archaeal membrane; confer extreme
thermostability \\
\textbf{CRISPR-Cas} & Bacterial/archaeal RNA-guided adaptive immune
system (adaptation--expression--interference); repurposed for genome
editing \\
\textbf{PAM} & Protospacer adjacent motif; short sequence flanking the
target that licenses CRISPR-Cas cleavage and ensures self/non-self
discrimination \\
\end{longtable}
}

\begin{center}\rule{0.5\linewidth}{0.5pt}\end{center}

\subsection{Review Questions}\label{review-questions-24}

\begin{enumerate}
\def\labelenumi{\arabic{enumi}.}
\item
  Compare the cell wall structure of Gram-positive and Gram-negative
  bacteria. Explain why Gram-negative sepsis is often more dangerous
  than Gram-positive sepsis, focusing on the role of LPS and TLR4
  signaling.
\item
  A researcher isolates a bacterium from a hot spring (85 degrees C, pH
  3) that has ether-linked lipids, no peptidoglycan, and an RNA
  polymerase with 12 subunits. To which domain does this organism most
  likely belong? Justify your classification using at least three
  molecular features.
\item
  Explain how a single conjugation event involving an R-plasmid carrying
  genes for beta-lactamase (bla), aminoglycoside acetyltransferase
  (aac), and a tetracycline efflux pump (tet) can convert a susceptible
  bacterium into a multidrug-resistant pathogen in a single step. Why is
  this more concerning than resistance arising through chromosomal
  mutation?
\item
  Compare generalized and specialized transduction. For each, describe
  the molecular error that leads to bacterial DNA being packaged in
  phage particles, and explain why specialized transduction transfers
  primarily specific genes while generalized transduction can transfer
  any gene.
\item
  The HIV provirus can remain latent in resting memory CD4\(^+\) T cells
  for decades, even in patients on suppressive ART. Explain why this
  latent reservoir prevents cure and why the ``shock and kill'' strategy
  (reactivating latent virus with latency-reversing agents while the
  patient is on ART) has been proposed.
\item
  Draw the lytic cycle of bacteriophage T4 and calculate the following:
  if a single T4 phage infects an \emph{E. coli} culture containing
  \(10^6\) cells, the burst size is 150, and the lytic cycle takes 25
  minutes, how many phage particles will be present after 3 complete
  lytic cycles (assuming most released phage successfully infect new
  cells)?
\item
  Explain why \emph{Mycobacterium tuberculosis} is neither truly
  Gram-positive nor Gram-negative despite being classified in the
  Gram-positive phylum Actinobacteria. What cell wall component is
  responsible for this anomaly, and what alternative staining method is
  used?
\item
  The discovery of Asgard archaea has strengthened support for the
  ``two-domain tree'' of life over the traditional ``three-domain
  tree.'' Explain the key evidence supporting this model, including the
  eukaryotic signature proteins found in Asgard genomes and what this
  implies about the origin of eukaryotic cells.
\item
  A hospital laboratory reports that a \emph{Klebsiella pneumoniae}
  isolate is resistant to meropenem (a carbapenem) with an MIC of 32
  \texorpdfstring{\ensuremath{\mu}}{mu}g/mL. The resistance is mediated by a plasmid-borne KPC-2
  carbapenemase. Describe three infection control measures that should
  be implemented immediately, and explain why carbapenem resistance is
  considered a ``sentinel event'' in infectious disease.
\item
  Compare the replication strategies of influenza virus (Class V,
  segmented (-)ssRNA) and SARS-CoV-2 (Class IV, non-segmented (+)ssRNA).
  Explain why influenza can undergo antigenic shift through reassortment
  while coronaviruses cannot, and discuss the implications for pandemic
  preparedness.
\item
  \textbf{Growth-curve calculation.} Apply
  \cref{eq:unit_VII_bacterial_growth}: a culture of \emph{Streptococcus
  pneumoniae} (\(t_d = 30\) min) starts at \(N_0 = 5 \times 10^3\)
  CFU/mL. (a) After how many minutes will the density reach \(10^9\)
  CFU/mL --- the typical density at which sepsis becomes hemodynamically
  significant? (b) If the patient receives a single dose of ceftriaxone
  that reduces the population by 4 log\(_{10}\) instantly, and if the
  residual bacteria resume growth at the same rate 1 hour later, when
  does the density return to \(10^9\)?
\item
  Estimate phage amplification after three rounds if burst size
  \(B=150\) and every virion infects --- relate to the logistic
  assumptions that break down in real cultures.
\item
  Why are archaeal antibiotics rare compared with bacterial antibiotics,
  given clinical need?
\item
  \textbf{Peptidoglycan pathway integration.} Place fosfomycin,
  D-cycloserine, bacitracin, vancomycin, and ampicillin on the
  peptidoglycan biosynthesis pathway and explain why a single resistance
  mutation generally confers resistance to a single of them --- but
  VanA-type vancomycin resistance and PBP2a (mecA) MRSA evade two by
  altering different parts of the same step.
\item
  \textbf{Two-component systems and antibiotic stewardship.} PhoQ-PhoP
  modifies lipid A in response to host antimicrobial peptides. Explain
  how this TCS contributes to colistin resistance and why a PhoQ
  inhibitor would be expected to ``re-sensitize'' \emph{Salmonella} to
  polymyxins.
\item
  \textbf{Biofilm pharmacology.} A patient with a \emph{S. aureus}
  prosthetic-knee infection fails six weeks of vancomycin. Explain four
  distinct biofilm-related reasons for this failure and design a
  combination strategy targeting at least two of them.
\item
  \textbf{Phage therapy resistance.} A clinical phage cocktail uses
  three phages with non-overlapping receptors (LPS O-antigen, OmpC, type
  IV pilus). Estimate the probability that a single bacterium escapes
  most three by independent mutation and explain why cocktails have
  outperformed single-phage therapies in trials.
\item
  \textbf{CRISPR types comparison.} A biotech company wants to develop a
  phage-therapy product that itself targets multi-drug-resistant
  \emph{E. coli} CRISPR-immune to the original phage. Compare Type I
  (Cas3, processive degradation) and Type II (Cas9, blunt DSB)
  approaches for this application, and justify your choice based on (a)
  target size, (b) delivery, and (c) likelihood of bacterial escape.\#\#
  Further Reading and Source Notes
\end{enumerate}

\begin{itemize}
\tightlist
\item
  Woese \& Fox (1977). Phylogenetic structure of the prokaryotic domain:
  The primary kingdoms. \emph{Proceedings of the National Academy of
  Sciences}, 74.
\item
  COVID-19\} (2022). The Lancet Commission on lessons for the future
  from the COVID-19 pandemic. \emph{The Lancet}, 400.
\item
  Sagan (1967). On the origin of mitosing cells. \emph{Journal of
  Theoretical Biology}, 14.
\item
  Madigan et al.~(latest ed.). \emph{Brock Biology of Microorganisms}.
  Pearson.
\item
  Jinek, Chylinski, Fonfara, Hauer, Doudna \& Charpentier (2012). A
  programmable dual-RNA-guided DNA endonuclease in adaptive bacterial
  immunity. \emph{Science}, 337.
\item
  Barrangou et al.~(2007). CRISPR provides acquired resistance against
  viruses in prokaryotes. \emph{Science}, 315.
\item
  Spang et al.~(2015). Complex archaea that bridge the gap between
  prokaryotes and eukaryotes. \emph{Nature}, 521.
\end{itemize}

\begin{center}\rule{0.5\linewidth}{0.5pt}\end{center}

\subsubsection{Companion Source
Module}\label{companion-source-module-24}

\textbf{Bacteria, Archaea, and Viruses} should leave a reproducible
trail from a biological claim to the code, figure, diagram, or
paper-based activity that can test it. Use the surfaces below to inspect
the chapter's assumptions, rerun the relevant model, or compare the
manuscript explanation with companion labs and figures.

{\def\LTcaptype{none} % do not increment counter
\begin{longtable}[]{@{}
  >{\raggedright\arraybackslash}p{(\linewidth - 2\tabcolsep) * \real{0.5000}}
  >{\raggedright\arraybackslash}p{(\linewidth - 2\tabcolsep) * \real{0.5000}}@{}}
\toprule\noalign{}
\begin{minipage}[b]{\linewidth}\raggedright
Surface
\end{minipage} & \begin{minipage}[b]{\linewidth}\raggedright
Use it for
\end{minipage} \\
\midrule\noalign{}
\endhead
\bottomrule\noalign{}
\endlastfoot
\breaktt{src/biology/microbiology/microbiology.py}
(\breaktt{bacterial\_growth\_curve}, \texttt{doubling\_time},
\texttt{ViralReplicationResult}) & Check growth, doubling, and
viral-replication assumptions. \\
\breaktt{src/visualization/plots.py} (\breaktt{plot\_bacterial\_growth}) &
Inspect growth phases and axis scaling. \\
\breaktt{src/mermaid/biology\_diagrams.py}
(\breaktt{viral\_replication\_cycle\_diagram}) & Link genome strategy to
replication cycle. \\
\end{longtable}
}

\textbf{Reproducibility check:} specify taxon, environment, growth
phase, genome type, and measurement method before generalising microbial
claims. \textbf{Cross-reference:} connect with
\cref{sec:unit_VII_microbial_ecology} and
\cref{sec:unit_VII_infectious_disease}.

\newpage

\section{Microbial Ecology and the
Microbiome}\label{microbial-ecology-and-the-microbiome}

\label{sec:unit_VII_microbial_ecology}

\begin{quote}
\textbf{Ch 23} \texorpdfstring{\ensuremath{\cdot}}{.} Level 2/3 \texorpdfstring{\ensuremath{\cdot}}{.} 60 min read \texorpdfstring{\ensuremath{\cdot}}{.} 75 min lecture \texorpdfstring{\ensuremath{\cdot}}{.}
Prerequisites: \cref{sec:unit_VII_bacteria_archaea_viruses}
\end{quote}

\subsection{Learning Objectives}\label{learning-objectives-25}

By the end of this chapter, you should be able to:

\begin{enumerate}
\def\labelenumi{\arabic{enumi}.}
\tightlist
\item
  Explain the ``great plate count anomaly'' and describe
  culture-independent methods for studying microbial communities,
  including 16S rRNA amplicon sequencing and shotgun metagenomics.
\item
  Calculate and interpret alpha and beta diversity metrics for microbial
  communities, including Shannon entropy, Simpson and inverse-Simpson
  indices, and Chao1 richness estimation.
\item
  Distinguish OTU (operational taxonomic unit at 97\% identity) from ASV
  (exact sequence variant) approaches and explain the role of
  rarefaction curves in normalizing diversity comparisons.
\item
  Describe the composition and functional roles of the human
  \hyperref[gl:microbiome]{\textbf{microbiome}} across body sites, with
  emphasis on the gut microbiome and its core/variable components.
\item
  Explain microbiome-host interactions through pattern recognition
  receptors and short-chain fatty acid signalling.
\item
  Explain the relationship between dysbiosis and disease, including
  obesity, inflammatory bowel disease, and colorectal cancer.
\item
  Describe microbial roles in biogeochemical cycles (nitrogen, sulfur,
  carbon), including the gene markers used to track each transformation,
  and explain kill-the-winner dynamics in ocean phage-host systems.
\item
  Describe the roles of microorganisms in
  \hyperref[gl:biofilm]{\textbf{biofilm}} formation, quorum sensing, and
  bioremediation.
\end{enumerate}

\subsubsection{Study Blueprint}\label{study-blueprint-25}

\begin{itemize}
\tightlist
\item
  \textbf{Big idea:} Microbes organize communities, cycles, and host
  physiology through interactions and metabolism.
\item
  \textbf{Core concepts:} microbiomes, symbiosis, diversity,
  biogeochemical cycling.
\item
  \textbf{Framework alignment:} Vision \& Change: Evolution, Systems,
  Structure and function; AP Biology: Evolution, Systems Interactions;
  NGSS-style topics: Structure and Function, Interdependent
  Relationships in Ecosystems.
\item
  \textbf{Model or quantitative lens:} Diversity indices, richness
  estimates, and interaction-network reasoning.
\item
  \textbf{Data skill:} Compute or interpret community metrics from
  abundance data.
\item
  \textbf{Practice cadence:} Questions and Methods, Representing and
  Describing Data, Argumentation.
\item
  \textbf{Common misconception to repair:} A microbiome is not
  automatically beneficial; context determines the effect.
\item
  \textbf{Primary lab:} \cref{sec:lab_unit_VII_microbial_ecology}.
\item
  \textbf{Question bank:} \cref{sec:q_unit_VII_microbial_ecology}.
\item
  \textbf{Transfer task:} Transfer microbial ecology to soils, oceans,
  digestion, disease, and climate feedbacks.
\item
  \textbf{Bridge to computation:}
  \breaktt{biology.ecology.ecology.biodiversity\_indices}.
\end{itemize}

\begin{center}\rule{0.5\linewidth}{0.5pt}\end{center}

\begin{quote}
\textbf{Opening Vignette --- You Are Mostly Not You}

For decades, textbooks cited the figure that the human body contains ten
times more bacterial cells than human cells. In 2016, Ron Milo and Ron
Sender published a careful recalculation in \emph{Cell} --- the ratio is
approximately 1:1, with roughly 38 trillion bacteria and 30 trillion
human cells. But the bacteria contribute only about 0.2 kg of body mass,
compared to 70 kg of self. What changed science wasn't the ratio but the
Human Microbiome Project (2007--2016), which used 16S rRNA sequencing to
catalogue about 10,000 microbial species across 18 body sites
in 300 healthy adults. The project revealed that gut microbiome
composition influences immunity, metabolism, neurotransmitter
production, and drug efficacy --- and differs between individuals as
distinctly as a fingerprint. The microbiome is now recognised as an
additional organ, and its disruption (dysbiosis) is linked to
inflammatory bowel disease, obesity, depression, and
antibiotic-associated colitis. We are ecosystems, not individuals.
\end{quote}

\subsection{Microbial Ecology
Approaches}\label{microbial-ecology-approaches}

\subsubsection{The Great Plate Count
Anomaly}\label{the-great-plate-count-anomaly}

In 1985, Staley and Konopka formalized a long-recognized paradox in
microbiology: when environmental samples are examined by direct
microscopic counts or flow cytometry, cell densities of \(10^6\) per
milliliter are typical in aquatic environments. Yet when the same
samples are plated on standard laboratory media, primarily
\(10^2\)-\(10^3\) colony-forming units (CFU) per milliliter grow -- less
than 1\% of the organisms present. This discrepancy, termed the
\textbf{great plate count anomaly}, revealed that the vast majority of
microbial diversity is invisible to traditional culture-based
microbiology.

The reasons for unculturability are diverse:

\begin{itemize}
\tightlist
\item
  Unknown nutrient requirements or growth factors
\item
  Obligate syntrophic relationships (organisms that depend on
  metabolites from partner species)
\item
  Extremely slow growth rates (generation times of weeks to months)
\item
  Sensitivity to atmospheric oxygen levels (many are strict anaerobes or
  microaerophiles)
\item
  Dormancy or viable-but-not-culturable (VBNC) states
\item
  Requirement for cell-cell signaling or quorum sensing to initiate
  growth
\end{itemize}

\subsubsection{Culture-Independent
Methods}\label{culture-independent-methods}

Modern microbial ecology relies on molecular approaches that bypass
cultivation entirely:

\begin{figure}[htbp]
\centering
\includegraphics[width=0.82\linewidth,height=4.2in,keepaspectratio]{../figures/mermaid_inline/inline_mermaid_0124_4e149edd74cc.png}
\caption{Multi-omic workflow for culture-independent microbial community profiling: 16S amplicon and shotgun metagenomic DNA approaches answer "who is there" and "what could they do," while metatranscriptomics, metaproteomics, and metabolomics answer "what are they actually doing".}
\end{figure}

\subsubsection{The 16S rRNA Amplicon Pipeline: From Sample to ASV
Table}\label{the-16s-rrna-amplicon-pipeline-from-sample-to-asv-table}

\textbf{16S rRNA amplicon sequencing} is the workhorse of microbial
community profiling. Understanding the full pipeline, with its choices
and biases, is essential for interpreting any modern microbiome paper.

\textbf{Why 16S?} The 16S rRNA gene is ideal as a phylogenetic marker
because of three intrinsic properties:

\begin{itemize}
\tightlist
\item
  \textbf{Almost universally present} in most bacteria and
  \hyperref[gl:archaea]{\textbf{archaea}} (about  1500 bp).
\item
  \textbf{Mosaic of conserved and variable regions} --- nine
  hypervariable regions (V1--V9) are sandwiched between highly conserved
  regions that allow universal primer design (e.g., 27F, 515F, 806R,
  1492R).
\item
  \textbf{Slow evolutionary rate} in conserved domains preserves deep
  phylogenetic signal; rapid evolution in V regions enables genus- to
  species-level discrimination.
\end{itemize}

\textbf{Region selection.} Different V regions perform differently:

{\def\LTcaptype{none} % do not increment counter
\begin{longtable}[]{@{}
  >{\raggedright\arraybackslash}p{(\linewidth - 6\tabcolsep) * \real{0.2105}}
  >{\raggedright\arraybackslash}p{(\linewidth - 6\tabcolsep) * \real{0.2105}}
  >{\raggedright\arraybackslash}p{(\linewidth - 6\tabcolsep) * \real{0.2895}}
  >{\raggedright\arraybackslash}p{(\linewidth - 6\tabcolsep) * \real{0.2895}}@{}}
\toprule\noalign{}
\begin{minipage}[b]{\linewidth}\raggedright
Region
\end{minipage} & \begin{minipage}[b]{\linewidth}\raggedright
Length
\end{minipage} & \begin{minipage}[b]{\linewidth}\raggedright
Strengths
\end{minipage} & \begin{minipage}[b]{\linewidth}\raggedright
Weaknesses
\end{minipage} \\
\midrule\noalign{}
\endhead
\bottomrule\noalign{}
\endlastfoot
V1--V3 & about  480 bp & Good genus resolution;
\emph{Staphylococcus}-friendly & Misses some Bifidobacterium \\
\textbf{V3--V4} & about  460 bp & \textbf{De facto standard}
for human microbiome (HMP); good Illumina MiSeq fit & Underestimates
Bifidobacterium and some archaea \\
V4 primarily & about  250 bp & Single-end MiSeq; high read
depth & Lower taxonomic resolution \\
V6--V8 & about  440 bp & Better for environmental archaea &
Less reference-database coverage \\
Full-length (1.5 kb) & 1500 bp & Species- to strain-level resolution;
PacBio HiFi or Nanopore & Higher per-base error rate; lower
throughput \\
\end{longtable}
}

\textbf{Workflow.}

\begin{enumerate}
\def\labelenumi{\arabic{enumi}.}
\tightlist
\item
  \textbf{DNA extraction} --- bead beating (mechanical lysis of
  Gram-positives, fungi, spores) plus a column-based purification (e.g.,
  DNeasy PowerSoil). Extraction kit choice introduces bias of up to \texorpdfstring{\ensuremath{\pm}}{+/-} 30
  \% in relative abundances and is now considered a major
  reproducibility issue.
\item
  \textbf{PCR amplification} --- universal primers (515F/806R for V4,
  341F/805R for V3--V4) plus barcoded adapters allow multiplexing.
  \textbf{PCR cycles should be minimised (\texorpdfstring{\ensuremath{\leq}}{<=} 25)} to reduce chimera
  formation; high-fidelity polymerases (Q5, Phusion) reduce error.
\item
  \textbf{Library preparation and sequencing} --- Illumina MiSeq 2 \texorpdfstring{\ensuremath{\times}}{x} 300
  bp is standard; about  100 000 reads/sample is typical for
  diversity analyses, with about  10 \texorpdfstring{\ensuremath{\times}}{x} that needed to detect
  rare taxa.
\item
  \textbf{Quality filtering} --- trim adapters and primers; discard
  reads with average Q \textless{} 25 or expected errors \textgreater{}
  1.
\item
  \textbf{OTU clustering vs ASV inference} --- see comparison below.
\item
  \textbf{Taxonomic assignment} --- match representative sequences
  against reference databases (SILVA 138, Greengenes2, RDP) using naive
  Bayes (RDP classifier) or k-mer methods (VSEARCH); confidence
  threshold typically 80 \%.
\item
  \textbf{Diversity analyses and statistics} --- alpha (within-sample)
  and beta (between-sample) diversity, with multiple-comparison
  correction.
\end{enumerate}

\subsubsection{OTU vs ASV: A Decade-Long Methodological
Shift}\label{otu-vs-asv-a-decade-long-methodological-shift}

For about  15 years (2007--2017), microbiome studies
clustered reads into \textbf{operational taxonomic units (OTUs)} at a \texorpdfstring{\textgreater{}=}{>=}
97 \% identity threshold (corresponding roughly to genus-level
groupings). Around 2017, a methodological revolution replaced this with
\textbf{amplicon sequence variants (ASVs)} --- exact, error-corrected
sequences resolved at single-nucleotide precision. The contrast:

{\def\LTcaptype{none} % do not increment counter
\begin{longtable}[]{@{}
  >{\raggedright\arraybackslash}p{(\linewidth - 4\tabcolsep) * \real{0.1190}}
  >{\raggedright\arraybackslash}p{(\linewidth - 4\tabcolsep) * \real{0.4524}}
  >{\raggedright\arraybackslash}p{(\linewidth - 4\tabcolsep) * \real{0.4286}}@{}}
\toprule\noalign{}
\begin{minipage}[b]{\linewidth}\raggedright
Property
\end{minipage} & \begin{minipage}[b]{\linewidth}\raggedright
\textbf{OTUs (de novo, closed-reference)}
\end{minipage} & \begin{minipage}[b]{\linewidth}\raggedright
\textbf{ASVs (DADA2, Deblur, UNOISE3)}
\end{minipage} \\
\midrule\noalign{}
\endhead
\bottomrule\noalign{}
\endlastfoot
Clustering & Group reads at 97 \% identity & Statistical denoising; each
unique error-corrected sequence is its own ASV \\
Resolution & Roughly genus-level & Single-nucleotide; potentially
strain-level \\
Reproducibility across studies & \textbf{Poor} (cluster boundaries
depend on dataset) & \textbf{Good} (ASVs are exact sequences) \\
Tolerance to PCR/sequencing errors & High (errors collapsed into
clusters) & Requires explicit error model \\
False-positive rare ASVs & Lower & Higher unless filtered \\
Computational cost & Lower & Higher (DADA2 about  1 hr per
1000 samples) \\
Modern recommendation & Legacy use primarily & \textbf{De facto standard
since about  2018} \\
\end{longtable}
}

The Callahan \emph{et al.} (2017, \emph{ISME Journal}) ``Exact sequence
variants should replace OTUs'' argument has been broadly adopted: ASVs
are reproducible across studies and labs, allow meta-analyses across
datasets, and resolve closely related taxa that 97 \%-clustering would
merge. Closed-reference OTU pipelines (clustering to a fixed reference
database) remain useful for very large meta-analyses but are
increasingly being supplanted.

\textbf{Rarefaction and normalisation.} Different samples are sequenced
to different depths, which artificially inflates richness in deeply
sequenced samples (you see more rare taxa simply because you looked
harder). Two approaches address this:

\begin{itemize}
\tightlist
\item
  \textbf{Rarefaction} --- randomly subsample every sample down to the
  minimum library size before computing diversity. Plotting richness as
  a function of subsample size yields a \textbf{rarefaction curve}: a
  curve that plateaus indicates sufficient sampling depth; one still
  rising indicates undersampling.
\item
  \textbf{Normalisation by total-sum scaling, CSS (cumulative-sum
  scaling), or DESeq2-style geometric means} --- preserves the original
  count data with a statistical correction.
\end{itemize}

Modern best practice: report results from both approaches, confirm
robustness, and rarely compare unrarefied richness across samples of
unequal depth.

\subsubsection{Other Omics Layers}\label{other-omics-layers}

\begin{itemize}
\tightlist
\item
  \textbf{Shotgun metagenomics} sequences most DNA in a sample without
  amplification: no PCR bias; captures most organisms including viruses
  and eukaryotic microbes; enables functional gene profiling; produces
  \textbf{metagenome-assembled genomes (MAGs)} --- reconstructed genomes
  of uncultured organisms from metagenomic reads using coverage depth
  and tetranucleotide frequency binning. The Human Microbiome Project
  (HMP, 2008-2014) characterized the microbiomes of 300 healthy adults
  across 18 body sites, generating \textgreater{} 5.2 Tbp of sequence
  data.
\item
  \textbf{Metatranscriptomics} (total mRNA sequencing) reveals which
  genes are actively expressed -- the functional state of the community
  at the time of sampling. This distinguishes between metabolic
  potential (metagenomics) and metabolic activity.
\item
  \textbf{Metaproteomics} identifies the proteins actually translated,
  capturing post-transcriptional regulation.
\item
  \textbf{Metabolomics} profiles the small molecules present in a sample
  using mass spectrometry (LC-MS/MS) or nuclear magnetic resonance
  (NMR). This captures the chemical outputs and communication molecules
  of the microbial community, including short-chain fatty acids, bile
  acid metabolites, tryptophan derivatives, and quorum sensing signals.
\end{itemize}

\begin{quote}
\textbf{Concept Check 1:} A researcher sequences the 16S rRNA gene from
a soil sample and identifies 500 distinct amplicon sequence variants
(ASVs). She then performs shotgun metagenomics on the same sample and
assembles 200 MAGs. Explain why these numbers differ and what types of
organisms might be captured by one method but not the other.
\end{quote}

\begin{quote}
\textbf{Concept Check 1b:} A reviewer criticises a 2016 microbiome paper
for using 97 \%-identity OTU clustering instead of ASVs. The authors
respond that re-running with DADA2 gives the same overall conclusions.
Explain (a) what the reviewer's underlying concern is, (b) why the
conclusions might still be the same in this case, and (c) one biological
situation where OTU vs ASV could change a paper's conclusions.
\end{quote}

\subsubsection{Diversity Metrics}\label{diversity-metrics}

Quantifying microbial community diversity requires mathematical
frameworks borrowed from classical ecology. Diversity has two components
--- \textbf{richness} (how many distinct taxa?) and \textbf{evenness}
(how uniformly are they distributed?) --- and most metrics are weighted
combinations of these.

\textbf{Alpha diversity} (within-sample diversity):

\begin{itemize}
\tightlist
\item
  \textbf{Observed ASVs/OTUs}: Simple count of distinct taxa; sensitive
  to sequencing depth.
\item
  \textbf{Shannon entropy}: weights both richness and evenness
  logarithmically.
\item
  \textbf{Simpson index} and its \textbf{inverse}: dominance-weighted;
  emphasises common taxa.
\item
  \textbf{Chao1 estimator}: extrapolates from observed
  singletons/doubletons to estimate true richness.
\item
  \textbf{Faith's phylogenetic diversity (PD)}: Sum of branch lengths in
  the phylogenetic tree connecting most observed taxa; incorporates
  evolutionary relationships.
\end{itemize}

\paragraph{Shannon Diversity Index}\label{shannon-diversity-index}

The Shannon entropy is the single most-used alpha-diversity metric. It
is defined as:

\begin{equation}
H' = -\sum_{i=1}^{S} p_i \ln p_i
\label{eq:unit_VII_shannon}
\end{equation}

where \(S\) is the number of species and \(p_i\) is the relative
abundance (fraction) of species \(i\). Higher \(H'\) means greater
diversity. The maximum \(H'\) for \(S\) species is \(\ln S\), achieved
when most species are perfectly even (\(p_i = 1/S\)); the minimum is 0
(single species). A useful intuition: \(e^{H'}\) is the
\textbf{effective number of equally common species} --- a community with
\(H' = 2.30\) has the same Shannon diversity as 10 perfectly even
species.

\paragraph{Worked Example: Computing Shannon Diversity for Five
Species}\label{worked-example-computing-shannon-diversity-for-five-species}

\textbf{Problem:} A gut microbiome sample contains five dominant species
with the following observed counts: \emph{Bacteroides} (500),
\emph{Faecalibacterium} (300), \emph{Roseburia} (100),
\emph{Akkermansia} (60), \emph{Bifidobacterium} (40). Compute the
Shannon diversity index \(H'\) from \cref{eq:unit_VII_shannon}.

\textbf{Solution:}

\begin{enumerate}
\def\labelenumi{\arabic{enumi}.}
\item
  \textbf{Total counts:} \(N = 500 + 300 + 100 + 60 + 40 = 1000\).
\item
  \textbf{Relative abundances:}
  \[ p_1 = 0.500,\ p_2 = 0.300,\ p_3 = 0.100,\ p_4 = 0.060,\ p_5 = 0.040.  \label{eq:unit_VII_microbial_ecology_item_1}\]
\item
  \textbf{Compute \(-p_i \ln p_i\) term-by-term:}

  {\def\LTcaptype{none} % do not increment counter
  \begin{longtable}[]{@{}llll@{}}
  \toprule\noalign{}
  Species & \(p_i\) & \(\ln p_i\) & \(-p_i \ln p_i\) \\
  \midrule\noalign{}
  \endhead
  \bottomrule\noalign{}
  \endlastfoot
  \emph{Bacteroides} & 0.500 & \(-0.693\) & 0.347 \\
  \emph{Faecalibacterium} & 0.300 & \(-1.204\) & 0.361 \\
  \emph{Roseburia} & 0.100 & \(-2.303\) & 0.230 \\
  \emph{Akkermansia} & 0.060 & \(-2.813\) & 0.169 \\
  \emph{Bifidobacterium} & 0.040 & \(-3.219\) & 0.129 \\
  \end{longtable}
  }
\item
  \textbf{Sum:}
  \[ H' = 0.347 + 0.361 + 0.230 + 0.169 + 0.129 \approx 1.236.  \label{eq:unit_VII_microbial_ecology_item_2}\]
\item
  \textbf{Interpretation:} \(e^{H'} = e^{1.236} \approx 3.44\). The
  community has the same Shannon diversity as a perfectly even community
  of about  3.4 species --- capturing the fact that although
  there are 5 species present, the dominance of \emph{Bacteroides} and
  \emph{Faecalibacterium} substantially lowers the effective diversity.
  The maximum possible \(H'\) for 5 species is \(\ln 5 = 1.609\); this
  community has 77 \% of maximum evenness.
\end{enumerate}

\paragraph{Simpson Index and Inverse
Simpson}\label{simpson-index-and-inverse-simpson}

The \textbf{Simpson index} \texorpdfstring{\ensuremath{\lambda}}{lambda} measures dominance --- the probability that
two individuals drawn at random belong to the \emph{same} species:

\[ \lambda = \sum_{i=1}^{S} p_i^2 .  \label{eq:unit_VII_microbial_ecology_item_3}\]

For the example above,
\(\lambda = 0.500^2 + 0.300^2 + 0.100^2 + 0.060^2 + 0.040^2 = 0.250 + 0.090 + 0.010 + 0.0036 + 0.0016 = 0.355\).
Two complementary forms are usually reported:

\begin{itemize}
\tightlist
\item
  \textbf{Gini--Simpson} \(D = 1 - \lambda = 0.645\) --- probability
  that two random individuals are \emph{different} species.
\item
  \textbf{Inverse Simpson} \(1/\lambda = 1/0.355 \approx 2.82\) --- the
  effective number of species under Simpson weighting.
\end{itemize}

Compare with the Shannon effective number (\(e^{H'} = 3.44\)): Simpson
weights are more dominance-sensitive, so its effective number is smaller
because dominant species count more heavily. Reporting both Shannon and
Simpson is now standard because they emphasise different aspects of
community structure: Shannon weights most species roughly
logarithmically; Simpson is dominated by abundant taxa and is less
sensitive to rare taxa or sequencing errors.

\subsubsection{Worked Example: Estimating Species Richness using
Chao1}\label{worked-example-estimating-species-richness-using-chao1}

\textbf{Problem:} A researcher sequences 16S rRNA amplicons from a
deep-sea sediment sample. The initial bioinformatics pipeline identifies
a total of \(S_{obs} = 150\) distinct bacterial species (ASVs).
Examining the abundance tables, the researcher notes that 30 of these
species were observed exactly once (\(f_1 = 30\), singletons), and 10
species were observed exactly twice (\(f_2 = 10\), doubletons). Use the
Chao1 estimator to predict the true species richness of the community.

\textbf{Solution:}

\begin{enumerate}
\def\labelenumi{\arabic{enumi}.}
\item
  \textbf{Identify the variables:}

  \begin{itemize}
  \tightlist
  \item
    Observed species (\(S_{obs}\)) = 150
  \item
    Singletons (\(f_1\)) = 30
  \item
    Doubletons (\(f_2\)) = 10
  \end{itemize}
\item
  \textbf{Apply the Chao1 formula:}
  \[ \hat{S}_{Chao1} = S_{obs} + \frac{f_1^2}{2f_2}  \label{eq:unit_VII_microbial_ecology_item_4}\]

  \[ \hat{S}_{Chao1} = 150 + \frac{30^2}{2(10)}  \label{eq:unit_VII_microbial_ecology_item_5}\]

  \[ \hat{S}_{Chao1} = 150 + \frac{900}{20}  \label{eq:unit_VII_microbial_ecology_item_6}\]

  \[ \hat{S}_{Chao1} = 150 + 45 = 195  \label{eq:unit_VII_microbial_ecology_item_7}\]
\end{enumerate}

The Chao1 estimator predicts that the true species richness is
\textbf{195 species}. This means that despite sequencing deep enough to
find 150 species, approximately 45 rare species remain undetected in the
sample. The high number of singletons relative to doubletons suggests
that the community has a ``long tail'' of rare taxa and requires deeper
sequencing to capture full diversity.

\textbf{Beta diversity} (between-sample diversity):

\begin{itemize}
\tightlist
\item
  \textbf{Bray-Curtis dissimilarity}: Abundance-weighted; ranges from 0
  (identical) to 1 (completely different);
  \(BC_{jk} = 1 - \frac{2C_{jk}}{S_j + S_k}\), where \(C_{jk}\) is the
  sum of lesser abundances shared between samples
\item
  \textbf{UniFrac distance}: Incorporates phylogenetic information;
  unweighted UniFrac considers primarily presence/absence of taxa;
  weighted UniFrac also considers abundance
\item
  \textbf{Jaccard index}: Presence/absence primarily;
  \(J = \frac{|A \cap B|}{|A \cup B|}\)
\end{itemize}

\textbf{Rarefaction curves} plot species richness against sampling
effort (number of sequences). A curve that plateaus indicates sufficient
sampling depth; a curve still rising indicates undersampling.

\begin{quote}
\textbf{Concept Check 2:} Two gut microbiome samples have identical
Shannon entropy values (\(H' = 3.5\)) but very different Bray-Curtis
dissimilarity (0.85). Explain how two samples can have the same
within-sample diversity but be highly dissimilar to each other.
\end{quote}

\begin{quote}
\textbf{Concept Check 2b:} Compute the Shannon and inverse-Simpson
values for these two communities and explain why the rankings differ:

\begin{itemize}
\tightlist
\item
  Community A: 4 species at proportions 0.97, 0.01, 0.01, 0.01
\item
  Community B: 4 species at proportions 0.40, 0.30, 0.20, 0.10
\end{itemize}

Which is more diverse by Shannon? By inverse Simpson? Why?
\end{quote}

\begin{center}\rule{0.5\linewidth}{0.5pt}\end{center}

\subsection{The Human Microbiome}\label{the-human-microbiome}

\subsubsection{Overview}\label{overview}

The human body harbors approximately \(3.8 \times 10^{13}\) bacterial
cells -- close to the number of human somatic cells
(\(3.7 \times 10^{13}\)). The microbial gene catalog is far larger:
approximately 3.3 million unique microbial genes compared to
approximately 20,000 human
\hyperref[gl:protein]{\textbf{protein}}-coding genes, a ratio exceeding
150:1. This ``second genome'' provides metabolic capabilities that human
cells lack, effectively extending our biochemical repertoire.

The total mass of the human microbiome is approximately 200 g,
predominantly residing in the colon. Each body site harbors a distinct
microbial community shaped by local environmental conditions
(\hyperref[gl:ph]{\textbf{pH}}, oxygen, moisture, nutrient availability,
immune surveillance).

\subsubsection{Gut Microbiome}\label{gut-microbiome}

The gastrointestinal tract represents the densest microbial habitat on
the human body, with a gradient from relatively sparse colonization in
the stomach (\(10^1\)-\(10^3\) bacteria/mL, limited by gastric acid) to
extraordinary density in the colon (\(10^{11}\) bacteria/mL --
approaching the theoretical packing limit for bacterial cells).

\textbf{\hyperref[gl:dominant]{Dominant} phyla in healthy adults:}

{\def\LTcaptype{none} % do not increment counter
\begin{longtable}[]{@{}
  >{\raggedright\arraybackslash}p{(\linewidth - 6\tabcolsep) * \real{0.1429}}
  >{\raggedright\arraybackslash}p{(\linewidth - 6\tabcolsep) * \real{0.3393}}
  >{\raggedright\arraybackslash}p{(\linewidth - 6\tabcolsep) * \real{0.2143}}
  >{\raggedright\arraybackslash}p{(\linewidth - 6\tabcolsep) * \real{0.3036}}@{}}
\toprule\noalign{}
\begin{minipage}[b]{\linewidth}\raggedright
Phylum
\end{minipage} & \begin{minipage}[b]{\linewidth}\raggedright
Relative Abundance
\end{minipage} & \begin{minipage}[b]{\linewidth}\raggedright
Key Genera
\end{minipage} & \begin{minipage}[b]{\linewidth}\raggedright
Functional Roles
\end{minipage} \\
\midrule\noalign{}
\endhead
\bottomrule\noalign{}
\endlastfoot
Firmicutes & about 60-65\% & \emph{Faecalibacterium},
\emph{Roseburia}, \emph{Eubacterium}, \emph{Ruminococcus},
\emph{Lactobacillus} & Major butyrate producers; fiber
\hyperref[gl:fermentation]{\textbf{fermentation}} \\
Bacteroidetes & about 20-25\% & \emph{Bacteroides},
\emph{Prevotella} & Complex polysaccharide degradation via PULs \\
Actinobacteria & about 3-5\% & \emph{Bifidobacterium},
\emph{Collinsella} & Dominant in breastfed infants; acetate and lactate
production \\
Proteobacteria & \textless5\% (in health) & \emph{Escherichia},
\emph{Klebsiella} & Expansion is a marker of dysbiosis \\
Verrucomicrobia & about 1-3\% & \emph{Akkermansia muciniphila}
& Mucin degradation; associated with metabolic health \\
Euryarchaeota (Archaea) & about 1\% & \emph{Methanobrevibacter
smithii} & Methanogenesis; removes H\(_2\) to improve fermentation
efficiency \\
\end{longtable}
}

The gut microbiome exhibits considerable inter-individual variation,
influenced by diet (plant-rich diets favor \emph{Prevotella};
animal-rich diets favor \emph{Bacteroides}), geography, age (dramatic
changes in infancy and old age), mode of birth (vaginal delivery
vs.~cesarean section), breastfeeding, and antibiotic exposure.

\subsubsection{Core vs Variable Microbiome and the F/B
Ratio}\label{core-vs-variable-microbiome-and-the-fb-ratio}

Two complementary frameworks describe gut-microbiome composition:

\textbf{Core microbiome} --- taxa present in \texorpdfstring{\textgreater{}=}{>=} 90 \% of healthy
individuals at \texorpdfstring{\textgreater{}=}{>=} 0.1 \% relative abundance. Surprisingly few species
qualify at the species level (5--20 across populations), but at the
\textbf{functional gene} level the core is much larger: \textgreater{}
50 \% of \emph{gene families} are shared across most healthy adults. The
biological inference is that \textbf{functional redundancy is conserved
even when species composition differs} --- the gut requires SCFA
production, bile-acid transformation, and vitamin synthesis, but many
species can supply each function. This is a textbook example of
community-level convergence under functional selection.

\textbf{Variable microbiome} --- the remaining about  80 \%
of taxa, which differ markedly between individuals. Variation is driven
by diet (long-term fiber \texorpdfstring{\ensuremath{\to}}{->} \emph{Prevotella}-dominant; meat-rich Western
\texorpdfstring{\ensuremath{\to}}{->} \emph{Bacteroides}-dominant), geography (the Hadza hunter-gatherers
carry \textgreater{} 50 \% of taxa absent from industrialised
populations), age, and antibiotic history.

\textbf{Firmicutes / Bacteroidetes (F/B) ratio.} The two dominant phyla
in the colon historically defined an enterotype-like axis. The Turnbaugh
\emph{et al.} (2006, \emph{Nature}) finding that obese mice have
elevated F/B prompted thousands of follow-up studies; the relationship
in humans is \textbf{considerably more nuanced} than initially reported.
Key caveats: (1) the ratio is method-sensitive (PCR primer choice
affects Bacteroidetes recovery by \texorpdfstring{\ensuremath{\pm}}{+/-} 30 \%); (2) human F/B in obesity
does not robustly replicate across cohorts (some studies show the
opposite direction); (3) the \emph{functional} signal --- increased
polysaccharide-degradation gene capacity in obese microbiomes --- is
more reproducible than the \emph{taxonomic} F/B ratio. Modern microbiome
research has largely moved away from F/B as a single biomarker and
toward \textbf{functional gene profiles}, \textbf{specific keystone
taxa} (\emph{F. prausnitzii}, \emph{Akkermansia}, \emph{Roseburia}), and
\textbf{microbial metabolites} (SCFAs, secondary bile acids, indole
derivatives).

\textbf{\emph{Akkermansia muciniphila} as a health marker.} \emph{A.
muciniphila} (Verrucomicrobia, 1--4 \% of healthy gut) degrades
intestinal mucin and stimulates host mucin replenishment, paradoxically
\emph{strengthening} the gut barrier. Lower \emph{Akkermansia} abundance
is associated with obesity, type 2 diabetes, IBD, and aging. Pasteurised
\emph{A. muciniphila} (a single non-replicating bacterial supplement,
killed by gentle pasteurisation that preserves the membrane protein
Amuc\_1100 --- which acts on TLR2 to upregulate tight junctions) has
shown improved insulin sensitivity and reduced LDL in early human trials
(Depommier \emph{et al.}, 2019, \emph{Nature Medicine}) --- a striking
case of a defined ``next-generation probiotic'' with a known molecular
mechanism. The European Food Safety Authority approved pasteurised
\emph{A. muciniphila} as a novel food in 2021.

\textbf{Other keystone taxa:}

\begin{itemize}
\tightlist
\item
  \textbf{\emph{Faecalibacterium prausnitzii}} (Firmicutes, 5--15 \% of
  healthy gut) --- major butyrate producer; depleted in IBD; produces a
  15-kDa anti-inflammatory peptide that inhibits NF-\texorpdfstring{\ensuremath{\kappa}}{kappa}B activation in
  colonocytes.
\item
  \textbf{\emph{Bacteroides fragilis}} --- produces polysaccharide A
  (PSA), a TLR2-Treg-inducing immunomodulator.
\item
  \textbf{\emph{Lactobacillus rhamnosus}} GG --- adheres to epithelium,
  secretes p40/p75 proteins that promote tight-junction integrity; the
  most-studied probiotic strain in clinical use.
\item
  \textbf{\emph{Roseburia intestinalis}} --- secondary butyrate producer
  via the butyrate-CoA-transferase pathway; depleted with fiber-poor
  diets.
\end{itemize}

\subsubsection{Oral Microbiome}\label{oral-microbiome}

The oral cavity harbors over 700 species across distinct habitats
(teeth, gingival crevice, tongue dorsum, buccal mucosa, hard palate).
Key species include:

\begin{itemize}
\tightlist
\item
  \emph{Streptococcus mutans}: Primary agent of dental caries;
  metabolizes sucrose to produce lactic acid, lowering local pH below
  the critical threshold (about 5.5) for enamel
  demineralization
\item
  \emph{Porphyromonas gingivalis}: Keystone pathogen in chronic
  periodontitis; produces gingipain proteases that degrade host tissues
  and subvert complement
\item
  \emph{Fusobacterium nucleatum}: Bridge organism in dental plaque
  biofilm formation; recently implicated as enriched in colorectal
  cancer tumors
\end{itemize}

\subsubsection{Skin Microbiome}\label{skin-microbiome}

The skin (about 1.5 m\(^2\) surface area) harbors \(10^4\)-\(10^6\)
bacteria per cm\(^2\), with community composition varying by
microenvironment:

\begin{itemize}
\tightlist
\item
  \textbf{Sebaceous (oily) sites} (forehead, back): Dominated by
  \emph{Cutibacterium acnes} (formerly \emph{Propionibacterium acnes});
  metabolizes sebum triglycerides; implicated in acne vulgaris
\item
  \textbf{Moist sites} (axillae, groin, toe web): \emph{Staphylococcus},
  \emph{Corynebacterium}; body odor results from bacterial metabolism of
  apocrine sweat components
\item
  \textbf{Dry sites} (forearm, leg): Most diverse; \emph{Cutibacterium},
  \emph{Staphylococcus}, \emph{Micrococcus}
\end{itemize}

\subsubsection{Vaginal Microbiome}\label{vaginal-microbiome}

In reproductive-age women, the vaginal microbiome is typically dominated
by \emph{Lactobacillus} species (L. crispatus, L. iners, L. gasseri, L.
jensenii), which maintain a protective low pH (about 3.8-4.5)
through lactic acid production. They also produce hydrogen peroxide and
bacteriocins that inhibit pathogen colonization.

Disruption of \emph{Lactobacillus} dominance -- replacement by anaerobes
such as \emph{Gardnerella vaginalis}, \emph{Atopobium vaginae}, and
\emph{Prevotella} -- characterizes \textbf{bacterial vaginosis (BV)},
associated with preterm birth, increased susceptibility to sexually
transmitted infections including HIV, and pelvic inflammatory disease.

\begin{quote}
\textbf{Clinical Connection: Fecal
\hyperref[gl:microbiota]{\textbf{Microbiota}} Transplant (FMT)}
Recurrent \emph{Clostridioides difficile} infection (rCDI) occurs when
antibiotic-disrupted gut microbiota fails to recover, allowing \emph{C.
difficile} spores to germinate and produce toxins A and B. FMT --
transfer of processed stool from a healthy, screened donor -- restores
colonization resistance and achieves cure rates exceeding 90\%, compared
to approximately 30\% for repeated vancomycin courses. The FDA approved
the first standardized FMT products in 2022-2023 (REBYOTA, a
microbiota-based product, and VOWST, purified \emph{Firmicutes} spores
in oral capsules). The mechanism involves restoration of secondary bile
acid production by commensal \emph{Clostridium scindens}, which inhibits
\emph{C. difficile} spore germination, combined with
\hyperref[gl:competitive-exclusion]{\textbf{competitive exclusion}} and
bacteriocin production by restored commensals.
\end{quote}

\begin{center}\rule{0.5\linewidth}{0.5pt}\end{center}

\subsection{Microbiome Functions and Host
Interactions}\label{microbiome-functions-and-host-interactions}

\subsubsection{Microbiome--Host Molecular
Communication}\label{microbiomehost-molecular-communication}

Host-microbiome interaction is mediated by three classes of molecular
interface: \textbf{(i) pattern-recognition receptors} that surveil
microbe-associated patterns; \textbf{(ii) microbial metabolites} that
act on host receptors; and \textbf{(iii) physical and immune barriers}
that contain the microbiota.

\textbf{Pattern recognition receptors (PRRs).} The intestinal epithelium
and lamina propria express a rich repertoire of PRRs that read the local
microbial milieu in real time:

\begin{itemize}
\tightlist
\item
  \textbf{TLR2} (with TLR1 or TLR6 heterodimers) --- recognises
  Gram-positive lipoteichoic acid, lipopeptides, and \emph{B. fragilis}
  polysaccharide A (PSA). PSA is a remarkable case: it is a microbial
  ligand that activates \emph{anti-inflammatory} Treg responses via
  TLR2, demonstrating that PRR signalling is not uniformly
  pro-inflammatory but tunable by ligand.
\item
  \textbf{TLR4} --- recognises Gram-negative lipid A; in healthy gut
  TLR4 is low-expressed apically and signals are dampened by negative
  regulators (A20, SIGIRR). Loss of this dampening (e.g., increased TLR4
  in IBD) drives chronic inflammation.
\item
  \textbf{TLR5} --- flagellin sensor; intestinal expression is
  \textbf{basolateral primarily}, so commensal flagellated bacteria in
  the lumen do not trigger inflammation, but invading flagellated
  pathogens that cross the epithelium do. A beautiful example of spatial
  PRR localisation as a discrimination strategy.
\item
  \textbf{TLR9} --- unmethylated CpG DNA; engaged by bacterial DNA in
  endosomes.
\item
  \textbf{NOD1, NOD2} --- cytosolic peptidoglycan fragment sensors; NOD2
  mutations are the strongest single genetic risk for Crohn's disease,
  illustrating how loss of microbial sensing predisposes to dysbiosis.
\item
  \textbf{AhR (aryl hydrocarbon receptor)} --- sensor for
  tryptophan-derived microbial metabolites (indoles, IPA,
  indole-3-aldehyde from \emph{Lactobacillus}); upregulates IL-22,
  mucin, and antimicrobial peptide production by intestinal cells.
\end{itemize}

\textbf{Short-chain fatty acid (SCFA) signalling.} SCFAs are not only
fuel but signalling molecules with three identified receptors and one
direct enzymatic mechanism:

{\def\LTcaptype{none} % do not increment counter
\begin{longtable}[]{@{}
  >{\raggedright\arraybackslash}p{(\linewidth - 6\tabcolsep) * \real{0.2500}}
  >{\raggedright\arraybackslash}p{(\linewidth - 6\tabcolsep) * \real{0.2500}}
  >{\raggedright\arraybackslash}p{(\linewidth - 6\tabcolsep) * \real{0.2500}}
  >{\raggedright\arraybackslash}p{(\linewidth - 6\tabcolsep) * \real{0.2500}}@{}}
\toprule\noalign{}
\begin{minipage}[b]{\linewidth}\raggedright
Receptor / target
\end{minipage} & \begin{minipage}[b]{\linewidth}\raggedright
Ligand preference
\end{minipage} & \begin{minipage}[b]{\linewidth}\raggedright
Cell type
\end{minipage} & \begin{minipage}[b]{\linewidth}\raggedright
Effect
\end{minipage} \\
\midrule\noalign{}
\endhead
\bottomrule\noalign{}
\endlastfoot
\textbf{GPR41 (FFAR3)} & Propionate \textgreater{} butyrate \texorpdfstring{\ensuremath{\gg}}{>>} acetate &
Enteroendocrine L-cells; sympathetic neurons & PYY, GLP-1 release;
sympathetic activation \\
\textbf{GPR43 (FFAR2)} & Acetate \texorpdfstring{\ensuremath{\approx}}{~} propionate \textgreater{} butyrate &
L-cells; adipocytes; neutrophils & GLP-1; lipogenesis regulation;
neutrophil chemotaxis \\
\textbf{GPR109A (HCAR2)} & Butyrate (also niacin) & Colonocytes;
macrophages; dendritic cells & Treg induction; anti-inflammatory
cascade \\
\textbf{HDAC inhibition} & Butyrate (mM concentrations in the colon) &
Colonocytes; immune cells & Hyperacetylation of histones; epigenetic
anti-inflammatory programmes \\
\end{longtable}
}

Butyrate has been called a ``molecular bridge'' between microbiome and
host because it can engage four mechanisms simultaneously:
G-protein-coupled signalling (GPR41/43/109A), epigenetic regulation
(HDAC inhibition), mitochondrial fuel (\texorpdfstring{\ensuremath{\beta}}{beta}-oxidation provides
about  70 \% of colonocyte ATP), and Treg differentiation
(via GPR109A and HDAC-dependent FoxP3 upregulation). Butyrate
concentrations in the colonic lumen reach 10--20 mM --- 1000\texorpdfstring{\ensuremath{\times}}{x} higher
than in serum, justifying the local-action interpretation.

\subsubsection{Metabolic Functions}\label{metabolic-functions}

The gut microbiome functions as a metabolic organ, performing
biochemical transformations that human cells cannot:

\begin{figure}[htbp]
\centering
\includegraphics[width=0.82\linewidth,height=4.2in,keepaspectratio]{../figures/mermaid_inline/inline_mermaid_0125_21e6e35bf969.png}
\caption{Microbiome metabolic functions: fermentation of fiber to short-chain fatty acids (butyrate, propionate, acetate) with their specific receptor targets and immune outputs; bile-acid deconjugation; vitamin synthesis; and neuroactive metabolite production.}
\end{figure}

\begin{itemize}
\item
  \textbf{Butyrate} (about 20\% of SCFA): Produced primarily by
  \emph{Faecalibacterium prausnitzii}, \emph{Roseburia intestinalis},
  and \emph{Eubacterium rectale}. Serves as the primary energy source
  for colonocytes (providing about 70\% of their energy needs).
  Acts as a \hyperref[gl:histone]{\textbf{histone}} deacetylase (HDAC)
  inhibitor, promoting anti-inflammatory gene expression. Strengthens
  intestinal barrier function by upregulating tight junction proteins
  (claudins, occludin). Induces colonic regulatory T cell (Treg)
  differentiation via GPR109A, promoting immune tolerance.
\item
  \textbf{Propionate} (about 25\% of SCFA): Transported to the liver via
  portal circulation where it regulates gluconeogenesis. Activates GPR41
  and GPR43 receptors on enteroendocrine L-cells, stimulating release of
  GLP-1 (glucagon-like peptide-1, promotes insulin secretion and
  satiety) and PYY (peptide YY, reduces appetite).
\item
  \textbf{Acetate} (about 55\% of SCFA): The most abundant SCFA; enters
  systemic circulation to serve as fuel for peripheral tissues including
  muscle and brain. Activates GPR43 on adipocytes, regulating lipolysis
  and fat storage.
\end{itemize}

\textbf{Vitamin synthesis}: Gut bacteria produce essential vitamins
including vitamin K2 (menaquinone, synthesized by \emph{Bacteroides} and
other gut anaerobes), vitamin B12 (cobalamin, by
\emph{Propionibacterium}), folate, biotin, and riboflavin. However, most
microbially produced vitamins are absorbed in the colon, where
absorption efficiency is lower than in the small intestine.

\textbf{Bile acid metabolism}: The liver synthesizes primary bile acids
(cholic acid, chenodeoxycholic acid) conjugated with glycine or taurine.
In the colon, gut bacteria deconjugate these bile acids (bile salt
hydrolases) and perform 7\texorpdfstring{\ensuremath{\alpha}}{alpha}-dehydroxylation to produce secondary bile
acids (deoxycholic acid, lithocholic acid). Secondary bile acids are
potent signaling molecules that activate the TGR5 receptor (stimulating
GLP-1 release and energy expenditure) and the farnesoid X receptor (FXR,
regulating lipid and glucose metabolism).

\textbf{Drug metabolism}: The gut microbiome can profoundly affect drug
pharmacokinetics through first-pass microbial biotransformation. The
cardiac glycoside digoxin is inactivated by \emph{Eggerthella lenta};
the anti-cancer prodrug irinotecan is reactivated to its toxic form by
bacterial \texorpdfstring{\ensuremath{\beta}}{beta}-glucuronidase, causing severe diarrhea; the Parkinson's drug
L-DOPA is decarboxylated by gut bacteria, reducing its bioavailability.

\subsubsection{Immune Development}\label{immune-development}

The microbiome is essential for proper immune system development.
Germ-free mice (raised in sterile isolators) exhibit:

\begin{itemize}
\tightlist
\item
  Atrophied Peyer's patches and mesenteric lymph nodes
\item
  Reduced secretory IgA (SIgA) production
\item
  Deficient Th17 cell populations in the lamina propria
\item
  Impaired regulatory T cell function
\item
  Increased susceptibility to infection
\end{itemize}

Colonization of germ-free mice with specific bacteria restores immune
function: segmented filamentous bacteria (SFB) induce Th17 cells in the
small intestine; \emph{Clostridium} clusters IV and XIVa induce colonic
Tregs; \emph{Bacteroides fragilis} polysaccharide A (PSA) activates TLR2
on Tregs, promoting IL-10 production.

In humans, neonatal microbiome composition during a critical
developmental window (first 1-3 years of life) shapes lifelong immune
tone. Disruptions during this window (cesarean delivery, formula
feeding, early antibiotic exposure) are associated with increased risk
of allergic diseases, asthma, and autoimmunity -- consistent with the
``hygiene hypothesis'' and its modern refinement, the ``old friends
hypothesis.''

\subsubsection{Gut-Brain Axis}\label{gut-brain-axis}

The bidirectional communication network between the gut microbiome and
the central nervous system operates through multiple channels:

\begin{itemize}
\tightlist
\item
  \textbf{Vagus nerve}: The tenth cranial nerve provides a direct neural
  pathway from the enteric nervous system (about 100-500
  million \hyperref[gl:neuron]{\textbf{neuron}}s, sometimes called the
  ``second brain'') to the brainstem. Specific bacterial metabolites
  (SCFAs, tryptamine) activate vagal afferents.
\item
  \textbf{Neurotransmitter production}: Approximately 95\% of the body's
  serotonin (5-hydroxytryptamine, 5-HT) is produced by enterochromaffin
  cells in the gut, stimulated by tryptophan metabolites from gut
  bacteria (especially \emph{Clostridiales}). Gut bacteria also produce
  GABA (\emph{Lactobacillus}, \emph{Bacteroides}), dopamine, and
  norepinephrine.
\item
  \textbf{SCFAs}: Cross the blood-brain barrier; modulate microglial
  maturation and function; influence neuroinflammation
\item
  \textbf{Immune mediators}: Microbial-induced cytokines (IL-6, TNF-\texorpdfstring{\ensuremath{\alpha}}{alpha})
  affect brain function via circumventricular organs
\end{itemize}

Experimental evidence includes dramatic behavioral differences in
germ-free mice (increased anxiety-like behavior in some strains,
decreased in others), which can be partially normalized by colonization
with specific probiotic strains (\emph{Lactobacillus rhamnosus} JB-1
reduces anxiety-like behavior via vagal signaling).

\subsubsection{Colonization Resistance}\label{colonization-resistance}

The resident microbiota provides a critical defense against pathogen
colonization through multiple mechanisms:

\begin{itemize}
\tightlist
\item
  \textbf{Competitive exclusion}: Commensals occupy ecological
  \hyperref[gl:niche]{\textbf{niche}}s (nutrient sources, attachment
  sites), preventing pathogen establishment
\item
  \textbf{Bacteriocin production}: Antimicrobial peptides produced by
  commensals directly inhibit pathogen growth
\item
  \textbf{Bile acid-mediated resistance}: Secondary bile acids produced
  by commensals (e.g., \emph{Clostridium scindens} produces deoxycholic
  acid) inhibit \emph{C. difficile} spore germination and vegetative
  growth
\item
  \textbf{Immune priming}: Commensals maintain baseline immune
  surveillance (SIgA, antimicrobial peptide production by Paneth cells)
\end{itemize}

\emph{Salmonella} enterica has evolved strategies to overcome
colonization resistance: it triggers intestinal inflammation (via T3SS
effectors), which disrupts the resident microbiota and generates unique
nutrient sources (tetrathionate, ethanolamine) that \emph{Salmonella}
can metabolize but commensals cannot.

\begin{quote}
\textbf{Concept Check 3:} Explain why broad-spectrum antibiotic
treatment increases the risk of \emph{Clostridioides difficile}
infection, using the concepts of colonization resistance and secondary
bile acid metabolism. Why does FMT work when repeated courses of
vancomycin often fail?
\end{quote}

\begin{quote}
\textbf{Concept Check 3b:} A patient with refractory ulcerative colitis
is treated with butyrate enemas. Symptoms improve over 4 weeks. Explain
the multiple molecular mechanisms by which luminal butyrate can
attenuate colonic inflammation, naming at least three host receptors /
targets engaged. Why might oral butyrate supplements fail where rectal
butyrate succeeds?
\end{quote}

\begin{center}\rule{0.5\linewidth}{0.5pt}\end{center}

\subsection{Dysbiosis and Disease}\label{dysbiosis-and-disease}

\textbf{Dysbiosis} refers to a compositional and functional imbalance in
the microbiome associated with disease. While causation versus
correlation remains a major challenge in microbiome research, several
disease associations have strong mechanistic evidence:

\textbf{Obesity and Metabolic Syndrome}:

The landmark study by Turnbaugh et al.~(2006, \emph{Nature})
demonstrated that the gut microbiome of obese mice (ob/ob) has an
increased Firmicutes/Bacteroidetes (F/B) ratio compared to lean
littermates. When the obese microbiome was transplanted into germ-free
mice, recipients gained 60\% more body fat than recipients of lean
microbiome -- despite consuming identical diets. The mechanism involves
increased energy harvest from dietary polysaccharides (the obese
microbiome has enhanced capacity to ferment complex carbohydrates) and
altered SCFA profiles affecting host energy regulation.

Metabolic endotoxemia provides an additional mechanism: gut barrier
dysfunction (``leaky gut'') allows LPS translocation into the
bloodstream, triggering chronic low-grade inflammation via TLR4
activation on adipose tissue macrophages, contributing to insulin
resistance.

\textbf{Inflammatory Bowel Disease (IBD)}:

Both Crohn's disease and ulcerative colitis are associated with:

\begin{itemize}
\tightlist
\item
  Reduced overall microbial diversity
\item
  Reduced butyrate-producing bacteria (especially \emph{Faecalibacterium
  prausnitzii}, a potent anti-inflammatory species that produces
  butyrate and secretes anti-inflammatory metabolites)
\item
  Increased Proteobacteria (adherent-invasive \emph{E. coli} in ileal
  Crohn's)
\item
  Reduced secondary bile acid production
\item
  Altered tryptophan metabolism favoring the kynurenine pathway
\end{itemize}

\textbf{Colorectal Cancer (CRC)}:

\emph{Fusobacterium nucleatum} is significantly enriched in CRC tumors
compared to adjacent normal tissue. Mechanistic studies demonstrate that
\emph{F. nucleatum}:

\begin{itemize}
\tightlist
\item
  Adheres to cancer cells via its FadA adhesin binding to E-cadherin
\item
  Activates the Wnt/\texorpdfstring{\ensuremath{\beta}}{beta}-catenin signaling pathway, promoting cell
  proliferation
\item
  Recruits tumor-infiltrating myeloid cells that suppress anti-tumor
  immunity
\item
  Correlates with poor prognosis and chemoresistance
\item
  Serves as a potential diagnostic biomarker (detectable in stool
  samples)
\end{itemize}

\textbf{Antibiotic-Induced Dysbiosis}:

Broad-spectrum antibiotics cause profound disruption of the gut
microbiome that can persist for months. Clinical consequences include:

\begin{itemize}
\tightlist
\item
  \emph{Clostridioides difficile} infection (CDI): disrupted
  colonization resistance
\item
  Pediatric antibiotic exposure and later obesity, asthma, and allergies
  (multiple epidemiological studies)
\item
  Selection for antibiotic-resistant organisms within the microbiome
  (resistome expansion)
\item
  Loss of bile acid-metabolizing bacteria, altering host metabolic
  signaling
\end{itemize}

\begin{quote}
\textbf{Clinical Connection: Microbiome-Based Diagnostics and
Therapeutics} The microbiome is emerging as both a diagnostic tool and
therapeutic target. Stool-based tests for \emph{Fusobacterium nucleatum}
DNA are being developed as non-invasive colorectal cancer screening
supplements. Live biotherapeutic products (LBPs) -- defined microbial
consortia -- are in clinical trials for IBD, metabolic syndrome, and
immune checkpoint inhibitor response enhancement. \emph{Akkermansia
muciniphila} supplementation has shown promise in improving metabolic
parameters in overweight individuals (Depommier et al., 2019,
\emph{Nature Medicine}). The field is moving from correlation to
causation through gnotobiotic mouse models, Mendelian randomization
studies, and interventional trials.
\end{quote}

\begin{center}\rule{0.5\linewidth}{0.5pt}\end{center}

\subsection{Environmental
Microbiology}\label{environmental-microbiology}

\subsubsection{Soil Microbiology}\label{soil-microbiology}

Soil is among Earth's most microbially diverse habitats, containing
roughly \(10^8\)-\(10^9\) bacterial cells per gram and thousands to tens
of thousands of amplicon sequence variants per gram depending on soil
type, sequencing depth, and the taxonomic cutoff used. The dominant
phyla often include Proteobacteria, Firmicutes, Actinobacteria,
Bacteroidetes, and Acidobacteria, but relative abundances shift strongly
with pH, moisture, carbon input, and land management.

Soil microorganisms drive global carbon decomposition, nitrogen cycling,
phosphorus solubilization, and organic matter turnover. The
\textbf{rhizosphere} (soil zone immediately surrounding plant roots) is
a hotspot of microbial activity, enriched 10-100 fold compared to bulk
soil, driven by root exudates (sugars, amino acids, organic acids) that
feed microbial communities.

\subsubsection{Marine and Environmental Microbiome
Diversity}\label{marine-and-environmental-microbiome-diversity}

Microbial diversity in environmental systems dwarfs that of
multicellular life. The dominant numerical fact about life on Earth is
that there are about  \(10^{30}\) microbial cells globally
--- \textgreater{} 99 \% of Earth's biological cells, and 70--90 \% of
Earth's biomass nitrogen and phosphorus.

\textbf{Ocean.} The ocean contains approximately \(1.2 \times 10^{29}\)
microbial cells in the photic zone alone. Two organisms dominate the
global ocean microbiome:

\begin{itemize}
\tightlist
\item
  \textbf{SAR11 (\emph{Pelagibacterales})}: The most abundant
  free-living organism on Earth, estimated at
  \(\sim 2.4 \times 10^{28}\) cells globally --- i.e., about 
  25 \% of ocean bacteria. Its genome (about  1.3 Mbp) is
  among the smallest for free-living bacteria, reflecting extreme
  \textbf{genomic streamlining}: the \emph{minimal-cell-cost hypothesis}
  posits that in oligotrophic surface waters where every nitrogen atom
  is precious, selection ruthlessly eliminates non-essential genes.
  SAR11 lacks transcription regulators, sigma factors, sugar-uptake
  systems, and most regulatory pathways --- it lives a high-throughput,
  regulation-free metabolic life.
\item
  \textbf{\emph{Prochlorococcus}} \textbf{and \emph{Synechococcus}}: The
  most abundant photosynthetic organisms on Earth (about 
  \(10^{27}\) cells of \emph{Prochlorococcus} and about 
  \(7 \times 10^{26}\) of \emph{Synechococcus}). \emph{Prochlorococcus}
  was discovered primarily in 1986 (Chisholm), demonstrating how
  recently we have grasped basic ocean ecology. This tiny cyanobacterium
  (\(< 1\) \texorpdfstring{\ensuremath{\mu}}{mu}m diameter) accounts for approximately \textbf{20 \% of
  global ocean primary production} --- it makes about 5 \% of the oxygen
  we breathe. Different ecotypes partition the water column by light
  intensity, with high-light ecotypes near the surface and low-light
  ecotypes extending to 200 m depth.
\end{itemize}

\textbf{Soil.} Containing about  \(10^8\) bacteria per gram
and often about  \(10^4\)--\(10^5\) ASVs per gram in deep
sequencing studies, soil is one of the largest reservoirs of microbial
diversity. Globally, soil holds about  \(2.6 \times 10^{29}\)
microbial cells. The dominant phyla --- Proteobacteria, Acidobacteria,
Actinobacteria, Bacteroidetes, Verrucomicrobia, Firmicutes --- vary
continuously with pH, moisture, and organic matter. Acidobacteria, in
particular, are abundant in molecular surveys but remain under-cultured
relative to their environmental abundance, which is why metagenomics,
single-cell genomics, and enrichment culture are most needed to connect
sequence diversity to physiology.

\textbf{Deep biosphere.} A relatively recent realisation is that the
\textbf{subsurface biosphere} (rocks, sediments, and aquifers below
about  100 m depth) contains an estimated \(10^{29}\) cells
--- comparable to or exceeding most surface marine biomass. These
extremely slow-growing chemolithotrophs (doubling times of years to
centuries) couple H\(_2\) from rock-water reactions to CO\(_2\)
fixation, sustaining ecosystems that have been functionally isolated
from surface life for about  Ma timescales. The deep
biosphere now constitutes one of the largest reservoirs of unstudied
microbial diversity on the planet.

\textbf{Hydrothermal vents.} Chemolithoautotrophic communities at
hydrothermal vents form the base of food webs completely independent of
sunlight. Sulfur-oxidizing bacteria (\emph{Thiomicrospira},
\emph{Beggiatoa}) and hydrogen-oxidizing archaea (\emph{Methanopyrus})
sustain ecosystems including tube worms (\emph{Riftia pachyptila}, which
lack a digestive system and rely entirely on endosymbiotic
sulfur-oxidizing bacteria housed in a specialized organ called the
trophosome).

\subsubsection{Microbial
Biogeochemistry}\label{microbial-biogeochemistry}

Microbes catalyse virtually every redox transformation of nitrogen,
sulfur, and carbon at the global scale. The key enzymatic steps are now
well-characterised, each linked to a marker gene that ecologists use to
track the corresponding flux in environmental DNA.

\textbf{Nitrogen cycle.} The biogeochemical nitrogen cycle is a
quasi-closed loop of microbially-catalysed redox reactions, every step
of which has a diagnostic gene:

{\def\LTcaptype{none} % do not increment counter
\begin{longtable}[]{@{}
  >{\raggedright\arraybackslash}p{(\linewidth - 6\tabcolsep) * \real{0.2500}}
  >{\raggedright\arraybackslash}p{(\linewidth - 6\tabcolsep) * \real{0.2500}}
  >{\raggedright\arraybackslash}p{(\linewidth - 6\tabcolsep) * \real{0.2500}}
  >{\raggedright\arraybackslash}p{(\linewidth - 6\tabcolsep) * \real{0.2500}}@{}}
\toprule\noalign{}
\begin{minipage}[b]{\linewidth}\raggedright
Transformation
\end{minipage} & \begin{minipage}[b]{\linewidth}\raggedright
Reaction
\end{minipage} & \begin{minipage}[b]{\linewidth}\raggedright
Organisms
\end{minipage} & \begin{minipage}[b]{\linewidth}\raggedright
Marker gene
\end{minipage} \\
\midrule\noalign{}
\endhead
\bottomrule\noalign{}
\endlastfoot
\textbf{Nitrogen fixation} & N\(_2\) \texorpdfstring{\ensuremath{\to}}{->} NH\(_3\) & \emph{Rhizobium},
\emph{Azotobacter}, cyanobacteria, \emph{Nostoc} & \emph{nifH} (Fe
protein subunit of nitrogenase) \\
\textbf{Ammonia assimilation} & NH\(_3\) \texorpdfstring{\ensuremath{\to}}{->} glutamate / glutamine & Most
organisms & \emph{gln}A, \emph{gdhA} \\
\textbf{Nitrification step 1 (AOA / AOB)} & NH\(_3\) \texorpdfstring{\ensuremath{\to}}{->} NO\(_2^-\) &
\emph{Nitrosomonas}, \emph{Nitrospira}, ammonia-oxidising archaea (AOA,
\emph{Nitrosopumilus}) & \emph{amoA} (ammonia monooxygenase) \\
\textbf{Nitrification step 2} & NO\(_2^-\) \texorpdfstring{\ensuremath{\to}}{->} NO\(_3^-\) &
\emph{Nitrobacter}, \emph{Nitrospira} & \emph{nxrA} (nitrite
oxidoreductase) \\
\textbf{Comammox} (complete ammonia oxidation, 2015 discovery) &
NH\(_3\) \texorpdfstring{\ensuremath{\to}}{->} NO\(_3^-\) in one organism & \emph{Nitrospira inopinata} &
\emph{amoA} + \emph{nxrA} \\
\textbf{Denitrification} & NO\(_3^-\) \texorpdfstring{\ensuremath{\to}}{->} NO\(_2^-\) \texorpdfstring{\ensuremath{\to}}{->} NO \texorpdfstring{\ensuremath{\to}}{->} N\(_2\)O \texorpdfstring{\ensuremath{\to}}{->}
N\(_2\) & \emph{Pseudomonas}, \emph{Paracoccus}, many facultatives &
\emph{narG, nirK/nirS, norB,} \textbf{\emph{nosZ}} (N\(_2\)O
reductase) \\
\textbf{Anammox} (anaerobic ammonia oxidation, 1995 discovery) &
NH\(_4^+\) + NO\(_2^-\) \texorpdfstring{\ensuremath{\to}}{->} N\(_2\) & \emph{Brocadia}, \emph{Kuenenia}
(\emph{Planctomycetes}) & \emph{hzsA, hzsB} (hydrazine synthase) \\
\textbf{DNRA} (dissimilatory NO\(_3^-\) \texorpdfstring{\ensuremath{\to}}{->} NH\(_4^+\)) & NO\(_3^-\) \texorpdfstring{\ensuremath{\to}}{->}
NH\(_4^+\) & \emph{E. coli}, \emph{Salmonella}, sulfate reducers &
\emph{nrfA} \\
\end{longtable}
}

Two transformations deserve special attention. \textbf{Nitrogen
fixation} by nitrogenase (\(N_2 + 8H^+ + 8e^- + 16\)ATP
\(\rightarrow 2\)NH\$\_3 + \(H\)\_2 + 16\$ADP \(+ 16\)P\(_i\)) is
irreversibly inactivated by oxygen, presenting a paradox for
\hyperref[gl:aerobic]{\textbf{aerobic}} diazotrophs. Solutions include
\emph{Rhizobium}-legume symbiosis (root nodules contain
\textbf{leghemoglobin} delivering O\(_2\) to bacteroid respiration while
keeping free O\(_2\) below the threshold for nitrogenase inactivation,
about  10 nM); cyanobacterial \textbf{heterocysts}
(specialised thick-walled cells lacking photosystem II, expressing
nitrogenase, connected to vegetative cells by microplasmodesmata); and
free-living \emph{Azotobacter vinelandii} using respiratory protection
(extremely high respiration rate consumes O\(_2\) before it reaches
nitrogenase) plus conformational protection (FeSII protein binds
nitrogenase in O\(_2\) presence).

\textbf{Anammox} --- anaerobic ammonia oxidation --- was a stunning 1995
discovery (Mulder \emph{et al.}) showing that \emph{Planctomycetes}
couple NH\(_4^+\) oxidation to NO\(_2^-\) reduction, producing N\(_2\)
via the toxic intermediate hydrazine (N\(_2\)H\(_4\), rocket fuel).
Anammox is responsible for about  30--50 \% of marine N\(_2\)
production globally and is now used industrially in wastewater nitrogen
removal (about  90 \% less aeration energy than conventional
nitrification + denitrification).

\textbf{Sulfur cycle.} Parallel redox cycle:

{\def\LTcaptype{none} % do not increment counter
\begin{longtable}[]{@{}
  >{\raggedright\arraybackslash}p{(\linewidth - 4\tabcolsep) * \real{0.3333}}
  >{\raggedright\arraybackslash}p{(\linewidth - 4\tabcolsep) * \real{0.3333}}
  >{\raggedright\arraybackslash}p{(\linewidth - 4\tabcolsep) * \real{0.3333}}@{}}
\toprule\noalign{}
\begin{minipage}[b]{\linewidth}\raggedright
Transformation
\end{minipage} & \begin{minipage}[b]{\linewidth}\raggedright
Marker gene
\end{minipage} & \begin{minipage}[b]{\linewidth}\raggedright
Organisms
\end{minipage} \\
\midrule\noalign{}
\endhead
\bottomrule\noalign{}
\endlastfoot
\textbf{Sulfate reduction} (SO\(_4^{2-}\) \texorpdfstring{\ensuremath{\to}}{->} H\(_2\)S) & \emph{dsrA}
(dissimilatory sulfite reductase) & \emph{Desulfovibrio},
\emph{Desulfobacter} \\
\textbf{Sulfide oxidation} (H\(_2\)S \texorpdfstring{\ensuremath{\to}}{->} SO\(_4^{2-}\)) & \emph{soxB} &
\emph{Thiobacillus}, \emph{Beggiatoa}, hydrothermal-vent symbionts \\
\textbf{Disproportionation} (S\(^0\) \texorpdfstring{\ensuremath{\to}}{->} SO\(_4^{2-}\) + H\(_2\)S) &
\emph{dsrA} (reverse) & \emph{Desulfocapsa} \\
\end{longtable}
}

\textbf{Carbon cycle.} Microbial carbon transformations include CO\(_2\)
fixation (Calvin cycle: \emph{cbbL/cbbM} RuBisCO; reverse TCA,
\emph{aclB}; Wood--Ljungdahl, \emph{fhs/cooS}), methanogenesis
(\emph{mcrA} gene of methyl-CoM reductase --- diagnostic for most
methanogenic archaea), and anaerobic oxidation of methane (AOM, ANME
archaea + sulfate reducers acting via reverse methanogenesis), which
consumes about  80 \% of methane produced in marine sediments
before it reaches the atmosphere --- a globally important methane
filter.

\subsubsection{Phage--Host Dynamics:
Kill-the-Winner}\label{phagehost-dynamics-kill-the-winner}

Viruses (almost entirely phages) outnumber cellular life in the ocean by
about  10 : 1, with \(\sim 10^{30}\) viral particles
globally. Phages are the dominant agents of bacterial mortality, killing
20--40 \% of marine bacteria per day --- comparable to grazing pressure
--- and the dominant agent of horizontal gene transfer.

The \textbf{kill-the-winner (KTW) hypothesis} (Thingstad, 2000) explains
why marine microbial communities can be enormously diverse despite
competitive exclusion:

\begin{itemize}
\tightlist
\item
  When a bacterial species blooms (becomes the ``winner'' --- high cell
  density), its cognate phage finds it efficiently (encounter rate is
  bilinear in \(B \times P\)).
\item
  Phage proliferation crashes the bloom (the proliferation threshold
  \(B^* = m/(bk)\), derived in the previous chapter, is exceeded once
  \(B\) is high enough).
\item
  Other less-abundant species are not at high enough density to sustain
  their own phages, so they are \textbf{not} killed proportionally ---
  they are released from competitive pressure.
\item
  Diversity is \emph{maintained} because phage predation specifically
  punishes dominance, opening niches for rare taxa.
\end{itemize}

The dynamics are analogous to Lotka--Volterra predator--prey cycles but
with one twist: phage population growth is bursty (each successful
infection produces about  50--500 progeny in
about  30 minutes), giving boom-and-bust cycles superimposed
on the long-term equilibrium. \textbf{Cyanophage-cyanobacteria
oscillations} in the ocean have been directly observed via metagenomics,
with measurable phage-mediated turnover of \emph{Prochlorococcus}
populations every few days. KTW is now considered as fundamental as
competitive exclusion for understanding microbial-community structure:
the predator-mediated maintenance of diversity that ecologists
discovered in macroscopic systems (Paine's keystone-predator
experiments) operates equally --- perhaps more strongly --- in microbial
systems.

\subsubsection{Biofilms}\label{biofilms}

Over 80\% of bacteria in natural environments exist in \textbf{biofilms}
-- structured communities enclosed in a self-produced extracellular
polymeric substance (EPS) matrix of polysaccharides, proteins,
extracellular DNA (eDNA), and lipids.

Biofilm formation proceeds through five stages:

\begin{enumerate}
\def\labelenumi{\arabic{enumi}.}
\tightlist
\item
  \textbf{Reversible attachment}: Planktonic cells contact a surface via
  \hyperref[gl:van-der-waals-forces]{\textbf{van der Waals forces}} and
  hydrophobic interactions
\item
  \textbf{Irreversible attachment}: Specific adhesin-receptor
  interactions; early EPS production; c-di-GMP (cyclic diguanylate
  monophosphate) levels rise, promoting sessile lifestyle
\item
  \textbf{Microcolony formation}: Cell division and EPS accumulation;
  quorum sensing coordinates gene expression
\item
  \textbf{Mature biofilm}: Complex 3D architecture with water channels;
  metabolic gradients (O\(_2\), pH, nutrients) create distinct
  microniches; HGT frequency increases 1,000-fold
\item
  \textbf{Dispersal}: c-di-GMP degradation; EPS-degrading enzymes
  released; motility genes reactivated; planktonic cells seed new
  surfaces
\end{enumerate}

Biofilms are \textbf{10-1,000-fold more antibiotic-resistant} than
planktonic cells due to:

\begin{itemize}
\tightlist
\item
  EPS barrier limiting antibiotic diffusion
\item
  Metabolically dormant persister cells (not susceptible to
  growth-dependent antibiotics)
\item
  Local accumulation of resistance enzymes (\texorpdfstring{\ensuremath{\beta}}{beta}-lactamases in the matrix)
\item
  Microenvironment heterogeneity (anaerobic zones where aminoglycosides
  are ineffective)
\end{itemize}

Clinical impact: about 80\% of chronic bacterial infections
involve biofilms (CDC), including cystic fibrosis lung infections
(\emph{Pseudomonas aeruginosa}), dental plaque, prosthetic joint
infections, catheter-associated UTIs, and endocarditis.

\subsubsection{Quorum Sensing}\label{quorum-sensing}

\textbf{Quorum sensing (QS)} is a cell density-dependent communication
system in which bacteria produce and detect small signaling molecules
(\textbf{autoinducers}) that accumulate as the population grows. When
autoinducer concentration exceeds a threshold, coordinate gene
expression is triggered across the population:

\begin{figure}[htbp]
\centering
\includegraphics[width=0.82\linewidth,height=4.2in,keepaspectratio]{../figures/mermaid_inline/inline_mermaid_0126_10bf714951fb.png}
\caption{Quorum sensing in Vibrio fischeri: at low cell density LuxI-synthesised AHL diffuses away; once a quorum is reached AHL accumulates, binds LuxR, and triggers the luxCDABE operon producing bioluminescence - used for counterillumination camouflage in the bobtail squid.}
\end{figure}

{\def\LTcaptype{none} % do not increment counter
\begin{longtable}[]{@{}
  >{\raggedright\arraybackslash}p{(\linewidth - 6\tabcolsep) * \real{0.2245}}
  >{\raggedright\arraybackslash}p{(\linewidth - 6\tabcolsep) * \real{0.1633}}
  >{\raggedright\arraybackslash}p{(\linewidth - 6\tabcolsep) * \real{0.2245}}
  >{\raggedright\arraybackslash}p{(\linewidth - 6\tabcolsep) * \real{0.3878}}@{}}
\toprule\noalign{}
\begin{minipage}[b]{\linewidth}\raggedright
QS System
\end{minipage} & \begin{minipage}[b]{\linewidth}\raggedright
Signal
\end{minipage} & \begin{minipage}[b]{\linewidth}\raggedright
Organisms
\end{minipage} & \begin{minipage}[b]{\linewidth}\raggedright
Regulated Functions
\end{minipage} \\
\midrule\noalign{}
\endhead
\bottomrule\noalign{}
\endlastfoot
LuxI/LuxR & N-acyl homoserine lactones (AHLs) & Gram-negatives (\emph{V.
fischeri}, \emph{P. aeruginosa}) & Bioluminescence, biofilm,
\hyperref[gl:virulence]{\textbf{virulence}} factors \\
Agr & Autoinducing peptides (AIPs) & Gram-positives (\emph{S. aureus}) &
Toxin production, biofilm dispersal \\
AI-2 (LuxS) & Furanosyl borate diester & Cross-species (most bacteria
with LuxS) & Interspecies communication \\
PQS & 2-heptyl-3-hydroxy-4-quinolone & \emph{P. aeruginosa} & Iron
acquisition, virulence \\
\end{longtable}
}

\subsubsection{Bioremediation}\label{bioremediation}

Microorganisms can be harnessed to degrade environmental pollutants:

\begin{itemize}
\tightlist
\item
  \emph{Pseudomonas putida}: Degrades toluene and xylene via the TOL
  \hyperref[gl:plasmid]{\textbf{plasmid}}-encoded pathways; model
  organism for aromatic hydrocarbon bioremediation
\item
  \emph{Geobacter sulfurreducens}: Reduces Fe\(^{3+}\) and U\(^{6+}\);
  used for uranium immobilization at contaminated sites; generates
  electricity in microbial fuel cells via extracellular electron
  transfer
\item
  \emph{Ideonella sakaiensis}: Discovered in 2016 at a PET bottle
  recycling facility in Japan; produces PETase enzyme that degrades
  polyethylene terephthalate (PET) plastic; engineered variants with
  enhanced activity are being developed for plastic waste bioremediation
\item
  Oil spill bioremediation: \emph{Alcanivorax borkumensis} and
  \emph{Marinobacter} species bloom after oil spills, degrading alkanes
  and polycyclic aromatic hydrocarbons; nutrient addition
  (biostimulation with nitrogen and phosphorus) accelerates degradation
\end{itemize}

\begin{quote}
\textbf{Concept Check 4:} \emph{Pseudomonas aeruginosa} uses quorum
sensing to coordinate virulence factor production. A pharmaceutical
company designs a drug that blocks AHL binding to the LuxR-type receptor
(LasR) without killing the bacteria. Explain the therapeutic rationale
for this ``quorum quenching'' approach compared to traditional
antibiotics, and predict a potential limitation.
\end{quote}

\begin{quote}
\textbf{Concept Check 4b:} Marine \emph{Prochlorococcus} populations
oscillate with their cyanophages on a about  5--7 day cycle,
while their Shannon diversity (about  1.5 nats) is
approximately constant. Use the kill-the-winner framework to (a) explain
why high overall diversity is maintained despite a single species
dominating numerically, (b) predict what happens to phage population
during a \emph{Prochlorococcus} crash, and (c) describe how a
metagenomic time-series experiment could test these predictions.
\end{quote}

\begin{center}\rule{0.5\linewidth}{0.5pt}\end{center}

\subsection{Microbial Biotechnology}\label{microbial-biotechnology}

\subsubsection{Industrial Fermentation}\label{industrial-fermentation}

Microorganisms have been harnessed for food production for millennia,
and modern biotechnology has expanded their applications enormously:

The current frontier is less about replacing fermentation with a new
idea and more about instrumenting it. Genome-scale metabolic models,
automated strain engineering, continuous bioreactors, metagenomic enzyme
discovery, and precision fermentation now let researchers tune yield,
by-product formation, feedstock use, and contamination risk. Claims
about ``sustainable'' biomanufacturing should therefore report substrate
source, energy input, purification burden, waste stream, and life-cycle
context rather than primarily the engineered microbe.

{\def\LTcaptype{none} % do not increment counter
\begin{longtable}[]{@{}
  >{\raggedright\arraybackslash}p{(\linewidth - 6\tabcolsep) * \real{0.2195}}
  >{\raggedright\arraybackslash}p{(\linewidth - 6\tabcolsep) * \real{0.2439}}
  >{\raggedright\arraybackslash}p{(\linewidth - 6\tabcolsep) * \real{0.2195}}
  >{\raggedright\arraybackslash}p{(\linewidth - 6\tabcolsep) * \real{0.3171}}@{}}
\toprule\noalign{}
\begin{minipage}[b]{\linewidth}\raggedright
Product
\end{minipage} & \begin{minipage}[b]{\linewidth}\raggedright
Organism
\end{minipage} & \begin{minipage}[b]{\linewidth}\raggedright
Process
\end{minipage} & \begin{minipage}[b]{\linewidth}\raggedright
Application
\end{minipage} \\
\midrule\noalign{}
\endhead
\bottomrule\noalign{}
\endlastfoot
Ethanol & \emph{Saccharomyces cerevisiae} & Glucose fermentation &
Beverages, biofuel \\
Lactic acid & \emph{Lactobacillus}, \emph{Streptococcus thermophilus} &
Lactose/glucose fermentation & Yogurt, cheese, sauerkraut \\
Acetic acid & \emph{Acetobacter aceti} & Ethanol oxidation & Vinegar \\
Antibiotics & \emph{Streptomyces} spp. & Secondary metabolism &
Streptomycin, tetracycline, erythromycin \\
Amino acids & \emph{Corynebacterium glutamicum} & Engineered
fermentation & L-glutamate (MSG), L-lysine \\
\end{longtable}
}

\subsubsection{Recombinant Protein
Production}\label{recombinant-protein-production}

\emph{E. coli} remains the workhorse for recombinant protein production:

\begin{itemize}
\tightlist
\item
  \textbf{Human insulin (Humulin)}: First recombinant therapeutic
  protein, approved 1982; replaced porcine/bovine insulin
\item
  \textbf{Human growth \hyperref[gl:hormone]{\textbf{hormone}}}:
  Replaced cadaveric pituitary extracts (which transmitted CJD prions)
\item
  \textbf{Erythropoietin (EPO)}: Stimulates red blood cell production;
  produced in CHO (Chinese hamster ovary) cells for proper glycosylation
\item
  \textbf{Hepatitis B \hyperref[gl:vaccine]{\textbf{vaccine}}}:
  Recombinant HBsAg produced in yeast (\emph{Saccharomyces cerevisiae})
  -- the first recombinant vaccine
\end{itemize}

\subsubsection{Environmental
Biotechnology}\label{environmental-biotechnology}

\begin{itemize}
\tightlist
\item
  \textbf{Biogas production}: Anaerobic digestion by methanogenic
  archaea converts organic waste to methane for energy; widely used in
  wastewater treatment and agricultural waste management
\item
  \textbf{Constructed wetlands}: Engineered ecosystems using microbial
  communities for water purification
\item
  \textbf{Synthetic biology}: Engineered microbial biosensors for
  environmental monitoring (arsenic detection, water quality); metabolic
  engineering for production of biofuels, pharmaceuticals, and commodity
  chemicals from renewable feedstocks
\end{itemize}

\begin{quote}
\textbf{Concept Check 5:} The enzyme PETase from \emph{Ideonella
sakaiensis} degrades PET plastic at a rate too slow for industrial
application. Describe two protein engineering approaches (one rational,
one directed evolution) that could be used to improve PETase activity,
and explain why thermostability would be a desirable trait for an
industrial PET-degrading enzyme.
\end{quote}

\begin{quote}
\textbf{Concept Check 5b:} Anammox bacteria couple ammonium oxidation to
nitrite reduction with hydrazine as an intermediate. (a) Why is anammox
more energy-efficient than the conventional nitrification +
denitrification pathway used in sewage treatment plants? (b) Why have
anammox bacteria been so difficult to grow in pure culture (doubling
times of 11+ days)? (c) Predict, in molar terms, the alkalinity change
associated with anammox versus conventional nitrification --- why does
this matter for plant operation?
\end{quote}

\begin{center}\rule{0.5\linewidth}{0.5pt}\end{center}

\subsection{Computational Bridge}\label{computational-bridge-22}

Shannon diversity \(H'\) for tabulated OTU counts matches
\breaktt{biodiversity\_indices}:

\begin{Shaded}
\begin{Highlighting}[]
\ImportTok{from}\NormalTok{ biology.ecology }\ImportTok{import}\NormalTok{ biodiversity\_indices}

\NormalTok{res }\OperatorTok{=}\NormalTok{ biodiversity\_indices([}\DecValTok{120}\NormalTok{, }\DecValTok{80}\NormalTok{, }\DecValTok{40}\NormalTok{, }\DecValTok{10}\NormalTok{])}
\BuiltInTok{print}\NormalTok{(}\BuiltInTok{round}\NormalTok{(res.shannon\_index, }\DecValTok{3}\NormalTok{), res.species\_richness)}
\end{Highlighting}
\end{Shaded}

\begin{quote}
\textbf{Clinical / systems note:} FMT trials stratify donors partly on
diversity and bile-acid-transforming guilds --- operationalisations of
the same richness/evenness concepts.
\end{quote}

\begin{center}\rule{0.5\linewidth}{0.5pt}\end{center}

\subsection{Current Evidence and Frontier
Biology}\label{current-evidence-and-frontier-biology-25}

For \textbf{Microbial Ecology and the Microbiome}, frontier biology
belongs inside the evidence logic of the chapter. Microbiology and
infectious disease now require One Health reasoning across people,
animals, environments, genomics, and antimicrobial stewardship. The core
reading question is this: microbiome claims should distinguish
correlation, mechanism, host context, perturbation, and causality.

\begin{itemize}
\tightlist
\item
  \textbf{What to verify:} identify the observation, model, assay, or
  dataset that would make the claim stronger or weaker.
\item
  \textbf{What to qualify:} state the scale, organism, cell type,
  environmental condition, or population where the claim is expected to
  hold.
\item
  \textbf{What to compare:} test at least one alternative explanation,
  baseline, or null model before treating the pattern as causal.
\item
  \textbf{What to cite:} distinguish primary evidence, review synthesis,
  public dataset, and institutional guidance; for recent or numeric
  claims, prefer the source closest to the measurement and state what
  has changed since it was published.
\end{itemize}

For AMR and pathogen claims, name the organism-resistance pair, the
selection pressure, the transmission route, and the surveillance
evidence that would change triage
\citep{who2024bppl,cdc2025antibioticuse,murray2022amr}.

\textbf{Source practice:} For pathogen, AMR, and intervention claims,
tie statements to organism-resistance pairs, surveillance evidence,
official guidance, and trial/regulatory status
\citep{who2024bppl,who2025tb,who2025malaria,cdc2025lenacapavirprep,cdc2026candidaauris}.

\subsubsection{Current Evidence Map: Microbiome Causality
Ladder}\label{current-evidence-map-microbiome-causality-ladder}

\begin{figure}[htbp]
\centering
\includegraphics[width=0.82\linewidth,height=4.2in,keepaspectratio]{../figures/mermaid_inline/inline_mermaid_0127_31af3b622823.png}
\caption{Microbiome claims become stronger as they move from correlation toward perturbation and rescue evidence; many human associations remain context-dependent.}
\end{figure}

\subsection{Summary}\label{summary-25}

\begin{itemize}
\tightlist
\item
  The \textbf{great plate count anomaly} reveals that \textless1\% of
  environmental microbes are culturable. Culture-independent methods
  (16S amplicon sequencing, shotgun metagenomics, metatranscriptomics,
  metabolomics) have revolutionized microbial ecology, revealing entire
  new phyla and functional capabilities. \textbf{ASVs (exact sequence
  variants)} have replaced 97 \%-clustered OTUs as the de facto standard
  since about  2018.
\item
  \textbf{Diversity metrics}: Alpha diversity quantifies within-sample
  diversity. Shannon entropy \(H' = -\sum p_i \ln p_i\)
  (\cref{eq:unit_VII_shannon}), Simpson index \texorpdfstring{\ensuremath{\lambda}}{lambda} and inverse Simpson
  \(1/\lambda\), Chao1 richness, and Faith's PD weight richness and
  evenness differently. Beta diversity (Bray-Curtis, UniFrac, Jaccard)
  quantifies between-sample differences. Rarefaction curves diagnose
  sampling depth.
\item
  The \textbf{human microbiome} (\(3.8 \times 10^{13}\) cells, 3.3
  million genes) varies by body site. The gut microbiome (Firmicutes,
  Bacteroidetes dominant) has a small \textbf{species-level core} but a
  large \textbf{functional core}; \emph{Akkermansia muciniphila},
  \emph{F. prausnitzii}, \emph{B. fragilis}, and \emph{Roseburia} are
  keystone taxa. Microbial SCFAs (butyrate as colonocyte fuel, HDAC
  inhibitor, GPR109A agonist; propionate via GPR41/43; acetate as
  systemic fuel) signal through three GPCRs and HDAC.
  Pattern-recognition receptors (TLR2/4/5, NOD1/2, AhR) read the
  microbial milieu in real time, with PSA from \emph{B. fragilis} a
  striking example of microbial control of Tregs.
\item
  \textbf{Dysbiosis} is associated with obesity (functional
  polysaccharide-degradation gene shift; F/B ratio is method-sensitive
  and not a reliable single biomarker), IBD (reduced \emph{F.
  prausnitzii}, reduced butyrate), colorectal cancer (\emph{F.
  nucleatum} enrichment), and antibiotic-induced \emph{C. difficile}
  infection (FMT achieves \textgreater90\% cure).
\item
  \textbf{Environmental microbiology}: Soil harbors \(10^8\)-\(10^9\)
  bacteria/g; ocean contains \emph{Prochlorococcus} (about 
  \(10^{27}\) cells, 20 \% of global primary production) and SAR11
  (about  \(10^{28}\) cells); deep biosphere holds another
  \(10^{29}\) cells. \textbf{\hyperref[gl:nitrogen-fixation]{Nitrogen
  fixation}} by nitrogenase requires 16 ATP per N\(_2\) and O\(_2\)
  protection. The microbial nitrogen cycle is tracked by \emph{nifH}
  (fixation), \emph{amoA} (nitrification), \emph{nosZ}
  (denitrification), and \emph{hzsA} (anammox); sulfur by \emph{dsrA}
  and \emph{soxB}; carbon by \emph{mcrA} (methanogenesis) and various
  CO\(_2\)-fixation markers. \textbf{Anammox} accounts for 30--50 \% of
  marine N\(_2\) production.
\item
  \textbf{Phage-host kill-the-winner dynamics} maintain ocean diversity:
  phages preferentially predate dominant bacterial species, releasing
  rare species from competitive exclusion. Phages turn over 20--40 \% of
  marine bacteria daily.
\item
  \textbf{Biofilms} (80\% of natural bacteria) are 10-1,000x more
  antibiotic-resistant; \textbf{quorum sensing} coordinates
  population-level behavior via autoinducers (AHL, AIP, AI-2).
\item
  \textbf{Bioremediation} uses microbial metabolism to degrade
  pollutants including hydrocarbons, heavy metals, and plastics
  (PETase).
\item
  \textbf{Connections:} See \cref{sec:unit_X_ecosystem_ecology} for
  nutrient cycles, \cref{sec:unit_III_bioenergetics_and_respiration} for
  fermentation products, and \cref{sec:unit_VII_infectious_disease} for
  dysbiosis and infection.
\end{itemize}

\begin{center}\rule{0.5\linewidth}{0.5pt}\end{center}

\subsection{Key Terms}\label{key-terms-25}

{\def\LTcaptype{none} % do not increment counter
\begin{longtable}[]{@{}
  >{\raggedright\arraybackslash}p{(\linewidth - 2\tabcolsep) * \real{0.3529}}
  >{\raggedright\arraybackslash}p{(\linewidth - 2\tabcolsep) * \real{0.6471}}@{}}
\toprule\noalign{}
\begin{minipage}[b]{\linewidth}\raggedright
Term
\end{minipage} & \begin{minipage}[b]{\linewidth}\raggedright
Definition
\end{minipage} \\
\midrule\noalign{}
\endhead
\bottomrule\noalign{}
\endlastfoot
\textbf{Great plate count anomaly} & Observation that \textless1\% of
environmental microorganisms grow on standard laboratory culture
media \\
\textbf{16S rRNA sequencing} & Culture-independent method using the
conserved/variable 16S rRNA gene as a phylogenetic marker for bacterial
identification \\
\textbf{V3--V4 region} & The most common 16S amplicon target
(about  460 bp) used in human microbiome studies; provides
genus-level resolution on Illumina MiSeq \\
\textbf{ASV (amplicon sequence variant)} & Error-corrected exact
sequence inferred by DADA2/Deblur/UNOISE3; replaces 97 \%-identity OTUs
as the modern standard \\
\textbf{OTU (operational taxonomic unit)} & Cluster of 16S sequences at
a fixed identity threshold (typically 97 \%); legacy approach,
reproducibility-limited \\
\textbf{Rarefaction curve} & Plot of richness versus sampling depth used
to assess whether sequencing was deep enough \\
\textbf{Metagenomics} & Shotgun sequencing of total environmental DNA;
reveals community composition and functional gene content without
cultivation \\
\textbf{MAG} & Metagenome-assembled genome; computationally
reconstructed genome of an uncultured organism from metagenomic data \\
\textbf{Alpha diversity} & Within-sample diversity measured by richness
(number of taxa) and evenness (relative abundance distribution) \\
\textbf{Beta diversity} & Between-sample diversity; measures
compositional differences between microbial communities \\
\textbf{Shannon entropy} & Alpha diversity metric:
\(H' = -\sum p_i \ln p_i\); accounts for both richness and evenness;
effective species number is \(e^{H'}\) \\
\textbf{Simpson index} & Probability that two random individuals belong
to the same species: \(\lambda = \sum p_i^2\); complement \(1-\lambda\)
is the Gini-Simpson; reciprocal \(1/\lambda\) is the inverse Simpson
effective number \\
\textbf{Chao1} & Estimator of true richness using singletons and
doubletons: \(\hat{S}_{Chao1} = S_{obs} + f_1^2/(2 f_2)\) \\
\textbf{Microbiome} & The collective community of microorganisms (and
their genomes) inhabiting a defined environment \\
\textbf{Core microbiome} & Taxa or genes present in nearly most healthy
individuals; smaller at species level than at functional-gene level \\
\textbf{Akkermansia muciniphila} & Mucin-degrading Verrucomicrobia
species correlated with metabolic health; pasteurised \emph{A.
muciniphila} approved as novel food (EFSA, 2021) \\
\textbf{Dysbiosis} & Compositional and functional imbalance of the
microbiome associated with disease states \\
\textbf{Short-chain fatty acids (SCFAs)} & Bacterial fermentation
products (butyrate, propionate, acetate) with roles in colonocyte
nutrition, immune regulation, and metabolic signaling \\
\textbf{GPR41 / GPR43 / GPR109A} & Three SCFA-sensing GPCRs that
translate microbial fermentation into host hormonal and immune
responses \\
\textbf{Colonization resistance} & Protection against pathogen
establishment provided by resident microbiota through competition,
bacteriocins, and bile acid metabolism \\
\textbf{Fecal microbiota transplant (FMT)} & Transfer of processed donor
stool to restore microbial community in recipients with dysbiosis,
especially recurrent \emph{C. difficile} infection \\
\textbf{Biofilm} & Structured microbial community enclosed in
self-produced EPS matrix; highly antibiotic-resistant \\
\textbf{Quorum sensing} & Cell density-dependent gene regulation via
secreted autoinducer molecules (AHL, AIP, AI-2) \\
\textbf{Nitrogenase} & Enzyme complex (NifH, NifD, NifK) that reduces
N\(_2\) to NH\(_3\); oxygen-sensitive; requires 16 ATP per N\(_2\) \\
\textbf{Anammox} & Anaerobic ammonia oxidation; \emph{Planctomycetes}
couple NH\(_4^+\) + NO\(_2^-\) \texorpdfstring{\ensuremath{\to}}{->} N\(_2\) via hydrazine; key gene
\emph{hzsA} \\
\textbf{Comammox} & Complete ammonia oxidation in a single organism
(\emph{Nitrospira inopinata}, 2015 discovery) carrying both \emph{amoA}
and \emph{nxrA} \\
\textbf{Kill-the-winner} & Phage-driven mechanism that maintains
bacterial diversity by preferentially predating dominant taxa \\
\textbf{Bioremediation} & Use of microbial metabolism to degrade or
immobilize environmental pollutants \\
\textbf{Gut-brain axis} & Bidirectional communication between gut
microbiome and CNS via vagus nerve, neurotransmitters, SCFAs, and immune
mediators \\
\end{longtable}
}

\begin{center}\rule{0.5\linewidth}{0.5pt}\end{center}

\subsection{Review Questions}\label{review-questions-25}

\begin{enumerate}
\def\labelenumi{\arabic{enumi}.}
\item
  Explain the great plate count anomaly and describe three specific
  reasons why environmental microorganisms may fail to grow on standard
  laboratory media. How has metagenomics addressed this limitation?
\item
  A researcher collects gut microbiome samples from 50 healthy
  individuals and 50 patients with Crohn's disease. Describe which alpha
  and beta diversity metrics you would calculate, what statistical test
  you would use to compare communities (e.g., PERMANOVA on Bray-Curtis
  distances), and predict the expected differences.
\item
  Trace the fate of dietary fiber (resistant starch) from ingestion to
  its effects on host metabolism and immunity. Include the specific
  bacteria involved, the SCFAs produced, the host receptors activated,
  and the downstream physiological effects.
\item
  A patient has experienced four recurrences of \emph{Clostridioides
  difficile} infection despite multiple courses of vancomycin. Explain
  the mechanism by which FMT achieves cure rates exceeding 90\%,
  focusing on colonization resistance and secondary bile acid
  metabolism. Why does vancomycin itself perpetuate the cycle?
\item
  Compare the ecological strategies of SAR11 (genome about 1.3
  Mbp) and \emph{Prochlorococcus} (genome about 1.7 Mbp).
  Explain why both organisms have undergone extreme genome streamlining
  despite occupying different ecological niches, and discuss the
  evolutionary trade-offs of minimal genomes.
\item
  Explain why \emph{Pseudomonas aeruginosa} biofilms in the lungs of
  cystic fibrosis patients are resistant to antibiotic therapy. Describe
  at least four distinct resistance mechanisms that operate within the
  biofilm, and suggest a multi-target therapeutic strategy.
\item
  The Firmicutes/Bacteroidetes ratio is often cited as a biomarker for
  obesity. Critically evaluate this claim by discussing: (a) the
  strength of evidence from gnotobiotic mouse experiments, (b)
  confounding factors in human studies, (c) why this ratio alone is an
  oversimplification of the microbiome-obesity relationship, and (d)
  what functional or species-specific biomarkers have replaced F/B in
  modern research.
\item
  Design an experiment using germ-free mice to determine whether the gut
  microbiome influences anxiety-like behavior through the vagus nerve.
  Include appropriate controls, specify which bacterial species you
  would use for colonization, and describe how you would test vagal
  involvement.
\item
  Nitrogenase is irreversibly inactivated by oxygen, yet
  \emph{Azotobacter vinelandii} is an obligate aerobe that fixes
  nitrogen. Explain the molecular mechanisms by which this organism
  protects its nitrogenase from oxygen damage.
\item
  A biofilm-forming \emph{Staphylococcus aureus} strain causes recurrent
  prosthetic joint infection despite appropriate systemic antibiotic
  therapy. Explain why biofilm removal (surgical debridement) is often
  necessary in addition to antibiotic treatment, and describe how
  persister cells within the biofilm contribute to recurrence.
\item
  \textbf{Shannon diversity practice.} Compute \(H'\) from
  \cref{eq:unit_VII_shannon} for the community
  \(\{200, 100, 50, 20, 10, 10, 5, 5\}\). Then compute the inverse
  Simpson index. Interpret which species drive each metric.
\item
  Using the bridge counts, compute Simpson diversity manually from
  \(p_i\) and compare to \texttt{res.simpson\_index}.
\item
  Contrast \textbf{Faith's PD} with Shannon richness for prioritising
  conservation of microbial lineages.
\item
  \textbf{Marker-gene biogeochemistry.} Describe what biogeochemical
  inference you can draw from a soil metagenome that contains: (a)
  abundant \emph{nifH} but very little \emph{amoA}; (b) abundant
  \emph{amoA} and \emph{nxrA}; (c) \emph{dsrA} and \emph{mcrA} in the
  same sample. What environmental setting could each combination
  represent?
\item
  \textbf{Anammox vs denitrification.} Why has anammox replaced
  conventional denitrification in many modern wastewater plants? Compute
  the stoichiometric advantage in O\(_2\) demand per mole of N removed.
\item
  \textbf{Kill-the-winner test.} Design a 30-day metagenomic time-series
  experiment in a marine mesocosm to test whether \emph{Prochlorococcus}
  abundance and cyanophage abundance show predator-prey oscillations.
  Specify sampling frequency, sequencing approach, and the statistical
  analysis (cross-correlation, lagged regression) you would use.\#\#
  Further Reading and Source Notes
\end{enumerate}

\begin{itemize}
\tightlist
\item
  Woese \& Fox (1977). Phylogenetic structure of the prokaryotic domain:
  The primary kingdoms. \emph{Proceedings of the National Academy of
  Sciences}, 74.
\item
  Turnbaugh, Ley, Hamady, Fraser-Liggett, Knight \& Gordon (2007). The
  human microbiome project. \emph{Nature}, 449.
\item
  Falkowski, Fenchel \& Delong (2008). The microbial engines that drive
  Earth's biogeochemical cycles. \emph{Science}, 320.
\item
  Sender, Fuchs \& Milo (2016). Revised estimates for the number of
  human and bacteria cells in the body. \emph{PLOS Biology}, 14.
\item
  Lozupone, Stombaugh, Gordon, Jansson \& Knight (2012). Diversity,
  stability and resilience of the human gut microbiota. \emph{Nature},
  489.
\item
  Madsen (latest ed.). \emph{Environmental Microbiology: From Genomes to
  Biogeochemistry}. Wiley-Blackwell.
\item
  Strous et al.~(1999). Missing lithotroph identified as new
  planctomycete (anammox). \emph{Nature}, 400.
\end{itemize}

\begin{center}\rule{0.5\linewidth}{0.5pt}\end{center}

\subsubsection{Companion Source
Module}\label{companion-source-module-25}

\textbf{Microbial Ecology and the Microbiome} should leave a
reproducible trail from a biological claim to the code, figure, diagram,
or paper-based activity that can test it. Use the surfaces below to
inspect the chapter's assumptions, rerun the relevant model, or compare
the manuscript explanation with companion labs and figures.

{\def\LTcaptype{none} % do not increment counter
\begin{longtable}[]{@{}
  >{\raggedright\arraybackslash}p{(\linewidth - 2\tabcolsep) * \real{0.5000}}
  >{\raggedright\arraybackslash}p{(\linewidth - 2\tabcolsep) * \real{0.5000}}@{}}
\toprule\noalign{}
\begin{minipage}[b]{\linewidth}\raggedright
Surface
\end{minipage} & \begin{minipage}[b]{\linewidth}\raggedright
Use it for
\end{minipage} \\
\midrule\noalign{}
\endhead
\bottomrule\noalign{}
\endlastfoot
\breaktt{src/biology/microbiology/microbiology.py}
(\breaktt{bacterial\_growth\_curve}, \texttt{doubling\_time}) & Quantify
growth constraints for interacting microbial populations. \\
\breaktt{src/biology/ecology/ecology.py} (\texttt{lotka\_volterra},
\texttt{connectance}, \breaktt{biodiversity\_indices}) & Treat
microbiomes as communities with measurable interaction structure. \\
\breaktt{src/mermaid/biology\_diagrams.py} (\texttt{food\_web\_diagram})
& Compare cross-feeding and competition with broader food-web logic. \\
\end{longtable}
}

\textbf{Reproducibility check:} distinguish association, perturbation
response, mechanism, and host/environment context before making a
microbiome-causality claim. \textbf{Cross-reference:} use
\cref{sec:unit_X_community_ecology} and
\cref{sec:unit_VII_infectious_disease}.

\newpage

\section{Infectious Disease and
Immunity}\label{infectious-disease-and-immunity}

\label{sec:unit_VII_infectious_disease}

\begin{quote}
\textbf{Ch 24} \texorpdfstring{\ensuremath{\cdot}}{.} Level 2/3 \texorpdfstring{\ensuremath{\cdot}}{.} 60 min read \texorpdfstring{\ensuremath{\cdot}}{.} 75 min lecture \texorpdfstring{\ensuremath{\cdot}}{.}
Prerequisites: \cref{sec:unit_VII_bacteria_archaea_viruses}
\end{quote}

\subsection{Learning Objectives}\label{learning-objectives-26}

By the end of this chapter, you should be able to:

\begin{enumerate}
\def\labelenumi{\arabic{enumi}.}
\tightlist
\item
  Apply Koch's postulates and molecular Koch's postulates to evaluate
  evidence for microbial causation of disease, and identify their
  limitations.
\item
  Describe \hyperref[gl:virulence]{\textbf{virulence}} factors
  (adhesins, toxins, invasion machinery, immune evasion strategies) used
  by major bacterial, viral, and eukaryotic pathogens.
\item
  Explain innate immune defenses including physical barriers, pattern
  recognition receptors (TLRs, NLRs, RIG-I, cGAS-STING), the three
  complement pathways and their convergence at C3, and cellular
  effectors (neutrophils with NETosis, macrophages, NK cells with
  missing-self recognition).
\item
  Describe adaptive immunity including V(D)J
  \hyperref[gl:recombination]{\textbf{recombination}}, MHC restriction,
  T helper cell subsets, cytotoxic T cell killing mechanisms, B cell
  activation, affinity maturation, and antibody class switching.
\item
  Compare the eight major \hyperref[gl:vaccine]{\textbf{vaccine}}
  platforms (live attenuated, inactivated, subunit, virus-like-particle,
  toxoid, conjugate, mRNA, viral vector) and apply the
  \hyperref[gl:herd-immunity]{\textbf{herd immunity}} equation
  \(p_c = 1 - 1/R_0\) to compute thresholds for measles, polio, and
  COVID-19 variants.
\item
  Describe antigenic variation in influenza (drift vs shift) and HIV
  (reverse-transcriptase quasi-species clouds), and connect this to
  vaccine reformulation and pandemic risk.
\item
  Catalogue the principal antibiotic-resistance mechanisms (\texorpdfstring{\ensuremath{\beta}}{beta}-lactamases
  including ESBLs and carbapenemases, efflux pumps, target modification
  including PBP2a in MRSA and 23S rRNA methylation in MLS-resistant
  streptococci, reduced permeability, bypass pathways) and link each to
  specific drug classes.
\item
  Describe the epidemiology and pathogenesis of major infectious
  diseases including tuberculosis, malaria, HIV/AIDS, and influenza, and
  explain the One Health framework.
\end{enumerate}

\subsubsection{Study Blueprint}\label{study-blueprint-26}

\begin{itemize}
\tightlist
\item
  \textbf{Big idea:} Disease dynamics emerge from host susceptibility,
  pathogen traits, transmission networks, and immunity.
\item
  \textbf{Core concepts:} transmission, R0, immunity, vaccination.
\item
  \textbf{Framework alignment:} Vision \& Change: Evolution, Systems,
  Structure and function; AP Biology: Evolution, Systems Interactions;
  NGSS-style topics: Structure and Function, Interdependent
  Relationships in Ecosystems.
\item
  \textbf{Model or quantitative lens:} R0, herd-immunity threshold, and
  SIR trajectory calculations.
\item
  \textbf{Data skill:} Interpret outbreak curves and intervention
  effects.
\item
  \textbf{Practice cadence:} Questions and Methods, Representing and
  Describing Data, Argumentation.
\item
  \textbf{Common misconception to repair:} R0 is not a fixed property of
  a pathogen alone; it depends on host behaviour and environment.
\item
  \textbf{Primary lab:} \cref{sec:lab_unit_VII_infectious_disease}.
\item
  \textbf{Question bank:} \cref{sec:q_unit_VII_infectious_disease}.
\item
  \textbf{Transfer task:} Transfer disease-dynamic reasoning to
  vaccination, antimicrobial resistance, and public-health policy.
\item
  \textbf{Bridge to computation:}
  \breaktt{biology.microbiology.microbiology.sir\_model}.
\end{itemize}

\begin{center}\rule{0.5\linewidth}{0.5pt}\end{center}

\begin{quote}
\textbf{Opening Vignette --- The Pandemic That Shaped Modern Immunology}

The 1918 influenza pandemic infected an estimated 500 million people ---
one-third of the world's population --- and killed between 50 and 100
million, more than World War I. Most victims were healthy young adults,
a terrifying reversal of the usual mortality pattern, caused by a
cytokine storm in which a vigorous immune response became
catastrophically self-destructive. The pandemic revealed in gruesome
detail the cost of a misdirected immune response and the urgency of
understanding host-pathogen dynamics. From its ashes grew modern
epidemiology, the concept of herd immunity, and the influenza
surveillance networks that today sequence viral
\hyperref[gl:genome]{\textbf{genome}}s in near-real time. When
SARS-CoV-2 emerged in 2019, it was 1918's lessons --- social distancing,
masking, rapid vaccine development --- that shaped the global response.
Infectious disease is not a relic of the past; it is the central
selective pressure on immune system evolution.
\end{quote}

\subsection{Host-Pathogen
Relationships}\label{host-pathogen-relationships}

\subsubsection{Koch's Postulates}\label{kochs-postulates}

In 1884, Robert Koch formalized criteria for establishing that a
specific microorganism causes a specific disease. These four postulates
remain foundational in infectious disease:

\begin{enumerate}
\def\labelenumi{\arabic{enumi}.}
\tightlist
\item
  The microorganism must be found in most cases of the disease but not
  in healthy individuals
\item
  The microorganism must be isolated from the diseased host and grown in
  pure culture
\item
  The cultured microorganism must cause the same disease when inoculated
  into a healthy, susceptible host
\item
  The same microorganism must be re-isolated from the experimentally
  infected host
\end{enumerate}

\textbf{Limitations of Koch's postulates} have become apparent with
advancing knowledge:

\begin{itemize}
\tightlist
\item
  \textbf{Asymptomatic carriers}: \emph{Salmonella typhi} (Typhoid
  Mary), \emph{Neisseria meningitidis} -- healthy carriers exist,
  violating postulate 1
\item
  \textbf{Unculturable organisms}: Viruses (until cell culture),
  \emph{Treponema pallidum}, \emph{Mycobacterium leprae} -- cannot
  fulfill postulate 2
\item
  \textbf{Ethical constraints}: Human experimentation with lethal
  pathogens is prohibited
\item
  \textbf{Polymicrobial disease}: Periodontal disease, bacterial
  vaginosis -- caused by community shifts, not single organisms
\item
  \textbf{Host factors}: \emph{Mycobacterium tuberculosis} causes
  disease in about 5-10\% of infected individuals (the rest maintain
  latent infection)
\end{itemize}

\subsubsection{Molecular Koch's
Postulates}\label{molecular-kochs-postulates}

In 1988, Stanley Falkow proposed a molecular framework for identifying
virulence determinants:

\begin{enumerate}
\def\labelenumi{\arabic{enumi}.}
\tightlist
\item
  The virulence \hyperref[gl:gene]{\textbf{gene}} (or its product)
  should be found in pathogenic strains but not in non-pathogenic
  relatives
\item
  Inactivation of the gene (by \hyperref[gl:mutation]{\textbf{mutation}}
  or deletion) should reduce virulence in an appropriate model
\item
  Complementation (restoration of the gene) should restore virulence
\end{enumerate}

This framework has been essential for identifying virulence factors
through targeted gene knockouts and complementation studies in model
organisms.

\subsubsection{The Infection Continuum}\label{the-infection-continuum}

\begin{figure}[htbp]
\centering
\includegraphics[width=0.82\linewidth,height=4.2in,keepaspectratio]{../figures/mermaid_inline/inline_mermaid_0128_4fb628fe614a.png}
\caption{Decision-tree of pathogen exposure outcomes from initial encounter through clearance, latency, or systemic disease.}
\end{figure}

Not every exposure leads to colonization, not every colonization leads
to infection, and not every infection leads to disease. The outcome
depends on pathogen virulence, inoculum size, route of entry, and host
immune status.

\subsubsection{Virulence Factors}\label{virulence-factors}

Pathogens deploy specific molecular tools to adhere to host tissues,
invade cells, obtain nutrients, and evade immune defenses:

\textbf{Adhesins} mediate initial attachment to host surfaces:

\begin{itemize}
\tightlist
\item
  \textbf{Type IV pili} (\emph{Neisseria gonorrhoeae}, \emph{N.
  meningitidis}): Retractile pili that mediate attachment to epithelial
  surfaces and facilitate twitching motility
\item
  \textbf{FimH} (\emph{E. coli}): Type 1 fimbrial adhesin that binds
  mannose residues on uroepithelial cells, enabling urinary tract
  infection
\item
  \textbf{Fibronectin-binding \hyperref[gl:protein]{\textbf{protein}}s}
  (\emph{Staphylococcus aureus}): Mediate attachment to extracellular
  matrix, enabling wound infections and endocarditis
\end{itemize}

\textbf{Invasion factors}:

\begin{itemize}
\tightlist
\item
  \textbf{Type III secretion system (T3SS)}: \emph{Salmonella enterica}
  SPI-1 (Salmonella Pathogenicity Island 1) encodes a molecular syringe
  that injects effector proteins (SopE, SipA) directly into host
  epithelial cells, triggering \hyperref[gl:actin]{\textbf{actin}}
  \hyperref[gl:cytoskeleton]{\textbf{cytoskeleton}} rearrangement and
  bacterial internalization via membrane ruffling
\item
  \textbf{Internalins} (\emph{Listeria monocytogenes}): InlA binds
  E-cadherin, InlB binds Met receptor; trigger receptor-mediated
  \hyperref[gl:endocytosis]{\textbf{endocytosis}}; once internalized,
  Listeria escapes the phagosome (listeriolysin O) and propels itself
  through the \hyperref[gl:cytoplasm]{\textbf{cytoplasm}} using actin
  polymerization (ActA recruits host Arp2/3 complex)
\end{itemize}

\textbf{Capsules}: Polysaccharide capsules (\emph{Streptococcus
pneumoniae}, \emph{Neisseria meningitidis}, \emph{Haemophilus
influenzae} type b) inhibit phagocytosis by preventing complement C3b
deposition and obscuring surface antigens. The pneumococcal capsule has
\textgreater90 serotypes, forming the basis for conjugate vaccine design
(PCV13, PCV20).

\subsubsection{Toxins}\label{toxins}

\textbf{Exotoxins} are secreted proteins with specific mechanisms of
action. Many have the A-B structure: the B (binding) subunit binds a
host cell receptor, and the A (active) subunit enters the cell to exert
its toxic effect:

{\def\LTcaptype{none} % do not increment counter
\begin{longtable}[]{@{}
  >{\raggedright\arraybackslash}p{(\linewidth - 8\tabcolsep) * \real{0.0986}}
  >{\raggedright\arraybackslash}p{(\linewidth - 8\tabcolsep) * \real{0.1408}}
  >{\raggedright\arraybackslash}p{(\linewidth - 8\tabcolsep) * \real{0.2676}}
  >{\raggedright\arraybackslash}p{(\linewidth - 8\tabcolsep) * \real{0.2676}}
  >{\raggedright\arraybackslash}p{(\linewidth - 8\tabcolsep) * \real{0.2254}}@{}}
\toprule\noalign{}
\begin{minipage}[b]{\linewidth}\raggedright
Toxin
\end{minipage} & \begin{minipage}[b]{\linewidth}\raggedright
Organism
\end{minipage} & \begin{minipage}[b]{\linewidth}\raggedright
B-Subunit Target
\end{minipage} & \begin{minipage}[b]{\linewidth}\raggedright
A-Subunit Activity
\end{minipage} & \begin{minipage}[b]{\linewidth}\raggedright
Clinical Effect
\end{minipage} \\
\midrule\noalign{}
\endhead
\bottomrule\noalign{}
\endlastfoot
Cholera toxin & \emph{V. cholerae} & GM1 ganglioside & ADP-ribosylates
Gs\texorpdfstring{\ensuremath{\alpha}}{alpha} -\textgreater{} constitutive adenylyl cyclase activation
-\textgreater{} cAMP \(\uparrow\) & Cl\(^-\)/H\(_2\)O secretion
-\textgreater{} watery diarrhea (up to 20 L/day) \\
Diphtheria toxin & \emph{C. diphtheriae} & HB-EGF receptor &
ADP-ribosylates EF-2 -\textgreater{} halts protein synthesis &
Pseudomembrane in throat; myocarditis \\
Tetanospasmin & \emph{C. tetani} & Gangliosides (retrograde transport) &
Zinc metalloprotease cleaves VAMP/synaptobrevin & Blocks inhibitory
neurotransmitter release (glycine, GABA) -\textgreater{} spastic
paralysis \\
Botulinum toxin & \emph{C. botulinum} & Gangliosides (peripheral nerve)
& Zinc metalloprotease cleaves SNARE proteins (SNAP-25, syntaxin) &
Blocks ACh release at NMJ -\textgreater{} flaccid paralysis \\
Shiga toxin & \emph{Shigella}, STEC \emph{E. coli} & Gb3
(globotriaosylceramide) & N-glycosidase cleaves 28S rRNA & Halts protein
synthesis; HUS (hemolytic uremic syndrome) \\
\end{longtable}
}

\textbf{Endotoxin (LPS)}: Unlike exotoxins, endotoxin is not secreted
but released upon bacterial lysis. Lipid A activates TLR4 on
macrophages, triggering release of TNF-\texorpdfstring{\ensuremath{\alpha}}{alpha}, IL-1\texorpdfstring{\ensuremath{\beta}}{beta}, and IL-6. At high
concentrations (Gram-negative bacteremia), this causes fever,
hypotension, disseminated intravascular coagulation (DIC), multi-organ
failure, and septic shock.

\textbf{Immune evasion strategies}:

\begin{itemize}
\tightlist
\item
  \textbf{Protein A} (\emph{S. aureus}): Binds the Fc region of IgG in
  the ``wrong'' orientation, preventing opsonization and ADCC
\item
  \textbf{IgA protease} (\emph{N. gonorrhoeae}, \emph{H. influenzae}):
  Cleaves secretory IgA at mucosal surfaces
\item
  \textbf{Antigenic variation}: \emph{Trypanosoma brucei} expresses
  variant surface glycoprotein (VSG) from a library of \textgreater1,000
  \emph{vsg} genes, switching expression before the host mounts an
  effective antibody response
\item
  \textbf{Phagosome arrest}: \emph{M. tuberculosis} inhibits
  phagosome-lysosome fusion via LAM (lipoarabinomannan)
\end{itemize}

\begin{quote}
\textbf{Concept Check 1:} \emph{Vibrio cholerae} produces cholera toxin
that causes massive fluid secretion, yet the bacterium rarely invades
beyond the intestinal epithelium. Explain why this non-invasive strategy
is advantageous for the pathogen's transmission, and predict what would
happen to cholera epidemiology if the toxin were eliminated by gene
deletion.
\end{quote}

\begin{center}\rule{0.5\linewidth}{0.5pt}\end{center}

\subsection{Innate Immunity}\label{innate-immunity}

\subsubsection{Physical and Chemical
Barriers}\label{physical-and-chemical-barriers}

The first line of defense prevents pathogen entry into sterile body
compartments:

A barrier is not just a wall; it is an active ecological and
immunological interface. Skin acidity, mucus flow, antimicrobial
peptides, secretory IgA, iron sequestration, and resident microbiota
most create selection pressures that pathogens must evade. Barrier
failure can therefore come from physical breach, altered chemistry,
disrupted microbial competition, or medical devices that bypass normal
surfaces.

{\def\LTcaptype{none} % do not increment counter
\begin{longtable}[]{@{}
  >{\raggedright\arraybackslash}p{(\linewidth - 4\tabcolsep) * \real{0.2045}}
  >{\raggedright\arraybackslash}p{(\linewidth - 4\tabcolsep) * \real{0.2500}}
  >{\raggedright\arraybackslash}p{(\linewidth - 4\tabcolsep) * \real{0.5455}}@{}}
\toprule\noalign{}
\begin{minipage}[b]{\linewidth}\raggedright
Barrier
\end{minipage} & \begin{minipage}[b]{\linewidth}\raggedright
Mechanism
\end{minipage} & \begin{minipage}[b]{\linewidth}\raggedright
Pathogens That Breach It
\end{minipage} \\
\midrule\noalign{}
\endhead
\bottomrule\noalign{}
\endlastfoot
\textbf{Skin} & Intact keratin layer; low \hyperref[gl:ph]{\textbf{pH}}
(about 5.5); fatty acids; defensins; commensal
\hyperref[gl:microbiota]{\textbf{microbiota}} & Burns, wounds, catheter
insertion bypass skin \\
\textbf{Mucociliary escalator} & Respiratory mucus traps particles;
cilia beat at 12-15 Hz, propelling mucus upward & \emph{P. aeruginosa}
\hyperref[gl:biofilm]{\textbf{biofilm}} in CF; influenza HA cleaves
sialic acid in mucus \\
\textbf{Gastric acid} & pH 1.5-3.5; pepsin activation & \emph{H.
pylori}: urease produces NH\(_3\) to neutralize local pH \\
\textbf{Lysozyme} & Muramidase: cleaves NAG-NAM bonds in peptidoglycan &
Present in tears, saliva, nasal secretions, breast milk \\
\textbf{Lactoferrin} & Iron sequestration (bacteriostatic) & Found at
mucosal surfaces; deprives bacteria of essential iron \\
\textbf{Defensins} & Antimicrobial peptides (AMPs); form pores in
microbial membranes & Produced by epithelial cells and neutrophils;
\texorpdfstring{\ensuremath{\alpha}}{alpha}-defensins (Paneth cells), \texorpdfstring{\ensuremath{\beta}}{beta}-defensins (skin, respiratory) \\
\textbf{Resident microbiota} & Colonization resistance (nutrient
competition, bacteriocins, bile acid metabolism) & \emph{C. difficile}
exploits antibiotic-disrupted microbiota \\
\end{longtable}
}

\subsubsection{Pattern Recognition Receptors
(PRRs)}\label{pattern-recognition-receptors-prrs}

The innate immune system detects conserved microbial structures --
\textbf{pathogen-associated molecular patterns (PAMPs)} -- through
germline-encoded PRRs. This system provides immediate recognition
without prior exposure:

{\def\LTcaptype{none} % do not increment counter
\begin{longtable}[]{@{}
  >{\raggedright\arraybackslash}p{(\linewidth - 8\tabcolsep) * \real{0.0847}}
  >{\raggedright\arraybackslash}p{(\linewidth - 8\tabcolsep) * \real{0.1695}}
  >{\raggedright\arraybackslash}p{(\linewidth - 8\tabcolsep) * \real{0.2712}}
  >{\raggedright\arraybackslash}p{(\linewidth - 8\tabcolsep) * \real{0.3220}}
  >{\raggedright\arraybackslash}p{(\linewidth - 8\tabcolsep) * \real{0.1525}}@{}}
\toprule\noalign{}
\begin{minipage}[b]{\linewidth}\raggedright
PRR
\end{minipage} & \begin{minipage}[b]{\linewidth}\raggedright
Location
\end{minipage} & \begin{minipage}[b]{\linewidth}\raggedright
PAMP Recognized
\end{minipage} & \begin{minipage}[b]{\linewidth}\raggedright
Signaling Pathway
\end{minipage} & \begin{minipage}[b]{\linewidth}\raggedright
Outcome
\end{minipage} \\
\midrule\noalign{}
\endhead
\bottomrule\noalign{}
\endlastfoot
\textbf{TLR4} & Plasma membrane & LPS (lipid A) & MyD88 -\textgreater{}
NF-\texorpdfstring{\ensuremath{\kappa}}{kappa}B; TRIF -\textgreater{} IRF3 & Pro-inflammatory cytokines + type I
IFN \\
\textbf{TLR2} & Plasma membrane & Peptidoglycan, lipoteichoic acid,
zymosan & MyD88 -\textgreater{} NF-\texorpdfstring{\ensuremath{\kappa}}{kappa}B & Pro-inflammatory cytokines \\
\textbf{TLR3} & Endosome & dsRNA & TRIF -\textgreater{} IRF3 & IFN-\texorpdfstring{\ensuremath{\beta}}{beta}
(antiviral state) \\
\textbf{TLR7/8} & Endosome & ssRNA & MyD88 -\textgreater{} NF-\texorpdfstring{\ensuremath{\kappa}}{kappa}B + IRF7
& Type I IFN + cytokines \\
\textbf{TLR9} & Endosome & Unmethylated CpG DNA & MyD88 -\textgreater{}
NF-\texorpdfstring{\ensuremath{\kappa}}{kappa}B & Type I IFN + cytokines \\
\textbf{RIG-I / MDA5} & Cytoplasm & Viral RNA (5'-ppp, long dsRNA) &
MAVS -\textgreater{} IRF3 & IFN-\texorpdfstring{\ensuremath{\beta}}{beta} \\
\textbf{NLRP3} & Cytoplasm & DAMPs (ATP, uric acid, cholesterol
crystals) & ASC -\textgreater{}
\hyperref[gl:caspase]{\textbf{Caspase}}-1 & IL-1\texorpdfstring{\ensuremath{\beta}}{beta}, IL-18, gasdermin D
-\textgreater{} pyroptosis \\
\textbf{cGAS-STING} & Cytoplasm-ER & Cytoplasmic dsDNA & cGAMP
-\textgreater{} STING -\textgreater{} IRF3 & IFN-\texorpdfstring{\ensuremath{\beta}}{beta}; critical for DNA
virus detection \\
\end{longtable}
}

Key signaling outcomes of PRR activation:

\begin{itemize}
\tightlist
\item
  \textbf{NF-\texorpdfstring{\ensuremath{\kappa}}{kappa}B pathway}:
  \hyperref[gl:transcription]{\textbf{Transcription}} of
  pro-inflammatory cytokines (IL-1\texorpdfstring{\ensuremath{\beta}}{beta}, IL-6, TNF-\texorpdfstring{\ensuremath{\alpha}}{alpha}, IL-8/CXCL8) and
  chemokines
\item
  \textbf{IRF3/7 pathway}: Type I interferons (IFN-\texorpdfstring{\ensuremath{\alpha}}{alpha}/\texorpdfstring{\ensuremath{\beta}}{beta}) -\textgreater{}
  JAK-STAT signaling -\textgreater{} interferon-stimulated genes (ISGs:
  OAS/RNase L, Mx proteins, PKR) -\textgreater{} antiviral state in
  neighboring cells
\end{itemize}

\subsubsection{The Complement System: Three Pathways, One
Cascade}\label{the-complement-system-three-pathways-one-cascade}

The complement system is a cascade of about  30 plasma
proteins that opsonize pathogens, recruit inflammation, and lyse
Gram-negative bacteria. Three pathways converge on a common effector
cascade through the central enzyme \textbf{C3 convertase}, which
hydrolyses C3 into C3a (anaphylatoxin) and C3b (opsonin).

\begin{figure}[htbp]
\centering
\includegraphics[width=0.82\linewidth,height=4.2in,keepaspectratio]{../figures/mermaid_inline/inline_mermaid_0129_9049e3a1fb52.png}
\caption{Three pathways of complement activation (classical, lectin, alternative) converge at the C3 convertase, producing C3b (opsonin) and C3a (anaphylatoxin); a second cleavage at C5 generates C5a and the C5b-C9 membrane-attack complex.}
\end{figure}

\begin{itemize}
\tightlist
\item
  \textbf{C3a} --- a 9-kDa anaphylatoxin: binds C3aR on mast cells
  (degranulation, histamine), endothelial cells (vasodilation, vascular
  permeability), and granulocytes (recruitment).
\item
  \textbf{C3b} --- a 175-kDa opsonin: covalently binds pathogen surface;
  recognized by complement receptor 1 (CR1, CD35) on phagocytes,
  dramatically enhancing phagocytosis.
\end{itemize}

A second cleavage step generates:

\begin{itemize}
\tightlist
\item
  \textbf{C5a} --- the strongest neutrophil chemoattractant in plasma
  (effective at picomolar concentrations); binds C5aR1.
\item
  \textbf{C5b} --- initiates the \textbf{membrane attack complex (MAC)}:
  C5b + C6 + C7 + C8 + (10--18 copies of) C9 polymerise into a 10-nm
  transmembrane pore that lyses the target cell osmotically. The MAC is
  most effective against Gram-negative bacteria (which lack the thick
  peptidoglycan barrier of Gram-positives); \emph{Neisseria
  meningitidis} and \emph{N. gonorrhoeae} are particularly
  MAC-susceptible --- terminal-complement deficiencies (C5--C9) cause
  recurrent neisserial infections almost exclusively.
\end{itemize}

\textbf{Complement regulation.} Self cells avoid complement attack via
membrane-bound and soluble regulators:

{\def\LTcaptype{none} % do not increment counter
\begin{longtable}[]{@{}
  >{\raggedright\arraybackslash}p{(\linewidth - 4\tabcolsep) * \real{0.2895}}
  >{\raggedright\arraybackslash}p{(\linewidth - 4\tabcolsep) * \real{0.2105}}
  >{\raggedright\arraybackslash}p{(\linewidth - 4\tabcolsep) * \real{0.5000}}@{}}
\toprule\noalign{}
\begin{minipage}[b]{\linewidth}\raggedright
Regulator
\end{minipage} & \begin{minipage}[b]{\linewidth}\raggedright
Action
\end{minipage} & \begin{minipage}[b]{\linewidth}\raggedright
Defect \texorpdfstring{\ensuremath{\to}}{->} disease
\end{minipage} \\
\midrule\noalign{}
\endhead
\bottomrule\noalign{}
\endlastfoot
\textbf{CD46 (MCP, membrane cofactor protein)} & Cofactor for Factor I
cleavage of C3b/C4b & Atypical haemolytic uraemic syndrome (aHUS) \\
\textbf{CD55 (DAF, decay-accelerating factor)} & Accelerates decay of C3
and C5 convertases & Paroxysmal nocturnal haemoglobinuria (PNH;
CD55+CD59 GPI loss) \\
\textbf{CD59 (protectin)} & Blocks C9 polymerization (MAC) & PNH
(haemolysis) \\
\textbf{Factor H} & Inhibits alternative pathway on self surfaces &
aHUS; age-related macular degeneration (AMD) \\
\textbf{C1 inhibitor} & Blocks C1r/C1s and MASPs & Hereditary angioedema
(HAE) \\
\end{longtable}
}

\emph{Neisseria meningitidis} and \emph{N. gonorrhoeae} have evolved a
remarkable trick --- they bind host Factor H to their surface using
\textbf{factor H--binding protein (fHbp)}, mimicking self and
inactivating alternative-pathway amplification. fHbp is now a target
antigen in two licensed meningococcal-B vaccines (Bexsero, Trumenba).

\subsubsection{Phagocytosis and Cellular
Effectors}\label{phagocytosis-and-cellular-effectors}

\paragraph{Neutrophil Killing: Oxidative Burst and
NETosis}\label{neutrophil-killing-oxidative-burst-and-netosis}

\textbf{Neutrophils} are the most abundant circulating leukocytes
(50--70 \% of blood leukocytes) and the first cells to arrive at sites
of infection (within minutes of chemokine signal). Their lifespan is
extraordinarily short --- about  6--8 hours in circulation,
about  1--4 days at tissue sites --- reflecting a cell
programmed for rapid microbial killing followed by apoptosis to limit
collateral tissue damage.

\textbf{Oxidative burst (respiratory burst).} Upon phagosome formation,
the \textbf{NADPH oxidase complex (NOX2)} assembles at the phagosome
membrane:

\[ 2\text{O}_2 + \text{NADPH} \xrightarrow{\text{NOX2}} 2\text{O}_2^- + \text{NADP}^+ + \text{H}^+  \label{eq:unit_VII_infectious_disease_item_1}\]

Superoxide (\(O_2^-\)) is rapidly converted to hydrogen peroxide
(\(H_2O_2\)) by superoxide dismutase. \textbf{Myeloperoxidase (MPO)} ---
abundant in neutrophil azurophilic granules --- then catalyses the most
potent step of microbial killing:

\[ \text{H}_2\text{O}_2 + \text{Cl}^- \xrightarrow{\text{MPO}} \text{HOCl} + \text{OH}^-  \label{eq:unit_VII_infectious_disease_item_2}\]

\textbf{HOCl (hypochlorous acid, household-bleach chemistry)} is one of
the most potent oxidants in biology --- it chlorinates microbial
proteins, lipids, and DNA, killing many phagocytosed microbes within
minutes. The phagosome is also acidified (pH about  4.5) and
accumulates antimicrobial peptides (defensins, BPI, lactoferrin).
Survival is uncommon for non-adapted microbes but biologically important
pathogens such as \emph{M. tuberculosis}, \emph{Salmonella}, and
\emph{Leishmania} persist by blocking phagosome maturation, resisting
reactive species, or escaping the vacuole.

\textbf{NETosis.} Discovered in 2004 (Brinkmann \emph{et al.},
\emph{Science}), NETosis is a programmed neutrophil death pathway in
which the cell releases its decondensed
\hyperref[gl:chromatin]{\textbf{chromatin}} --- decorated with
citrullinated \hyperref[gl:histone]{\textbf{histone}}s, neutrophil
elastase, MPO, defensins, and granule contents --- to form
\textbf{neutrophil extracellular traps (NETs)} that physically capture
and chemically kill extracellular bacteria, fungi, and large parasites.
The molecular pathway:

\begin{enumerate}
\def\labelenumi{\arabic{enumi}.}
\tightlist
\item
  Activation by PAMPs / immune complexes / cytokines triggers ROS
  production by NOX2.
\item
  \textbf{PAD4 (peptidylarginine deiminase 4)} citrullinates histone H3
  (Arg \texorpdfstring{\ensuremath{\to}}{->} citrulline), neutralizing positive charge and decompacting
  chromatin.
\item
  Nuclear membrane breaks down; chromatin enters cytoplasm; granule
  contents bind to chromatin.
\item
  Plasma membrane lyses, releasing the NET into the extracellular space
  (cell death).
\end{enumerate}

NETs are evolutionarily ancient (zebrafish neutrophils,
\emph{Drosophila} haemocytes, even sea-anemone amoebocytes form similar
traps). Pathological NETosis is now implicated in
\textbf{immunothrombosis} in severe COVID-19 (NETs trap platelets and
trigger coagulation, contributing to microvascular clots), in autoimmune
diseases (citrullinated histones drive ACPA antibodies in rheumatoid
arthritis; anti-NET antibodies in lupus), and in chronic wound failure.

\textbf{Macrophages} are tissue-resident phagocytes with diverse
functions:

\begin{itemize}
\tightlist
\item
  Tissue-specific names: Kupffer cells (liver), alveolar macrophages
  (lung), microglia (brain), osteoclasts (bone), Langerhans cells (skin
  -- actually dendritic cells)
\item
  \textbf{M1 polarization} (IFN-\texorpdfstring{\ensuremath{\gamma}}{gamma} + LPS): Pro-inflammatory; produces
  iNOS -\textgreater{} nitric oxide (NO); secretes TNF-\texorpdfstring{\ensuremath{\alpha}}{alpha}, IL-12;
  activates Th1 responses against intracellular pathogens
\item
  \textbf{M2 polarization} (IL-4, IL-13): Anti-inflammatory; produces
  arginase; secretes IL-10, TGF-\texorpdfstring{\ensuremath{\beta}}{beta}; promotes wound healing, tissue
  repair, and fibrosis
\end{itemize}

\paragraph{NK Cells: Missing-Self and Induced-Self
Recognition}\label{nk-cells-missing-self-and-induced-self-recognition}

\textbf{Natural killer (NK) cells} are innate lymphocytes
(about  5--15 \% of blood lymphocytes) that kill
virus-infected and transformed cells without prior sensitization. They
are the innate immune system's solution to a fundamental problem: how to
detect cells that are \emph{abnormal} from the inside, when no conserved
pathogen pattern is present on the cell surface.

NK-cell killing is governed by an \textbf{integration of inhibitory and
activating receptor signals} --- the ``balance hypothesis.'' A cell is
killed when activating signals exceed inhibitory signals.

\textbf{Missing-self hypothesis (Ljunggren and Karre, 1990).} Healthy
cells display abundant MHC class I on their surface, presenting peptides
for surveillance by CD8\(^+\) T cells. NK cells carry \textbf{inhibitory
KIRs (killer immunoglobulin-like receptors; KIR2DL, KIR3DL)} that
recognise self MHC-I. Engagement delivers an \textbf{inhibitory signal}
through ITIMs (immunoreceptor tyrosine-based inhibitory motifs),
suppressing the NK-cell killing programme. Many viruses (CMV, HIV, KSHV)
and tumours \textbf{downregulate MHC class I} to avoid CD8\(^+\) T cell
detection --- but this loss removes the inhibitory signal to NK cells,
unmasking the cell for NK-mediated killing. The strategy has trade-offs:
CMV has evolved decoy MHC-I-like proteins (UL18) that bind inhibitory NK
receptors to mimic self.

\textbf{Induced-self via NKG2D.} A complementary mechanism detects
cellular stress. \textbf{NKG2D} is an activating receptor on NK cells
(and on subsets of CD8\(^+\) T cells, \texorpdfstring{\ensuremath{\gamma}}{gamma}\texorpdfstring{\ensuremath{\delta}}{delta} T cells, NKT cells) that
recognises \textbf{stress-induced ligands}: \textbf{MICA, MICB} (MHC
class I chain-related), \textbf{ULBP1--6} (UL16-binding proteins). These
ligands are \textbf{absent on healthy cells} but induced by:

\begin{itemize}
\tightlist
\item
  DNA damage (ATR/ATM pathway) --- common in transformed cells.
\item
  Viral infection.
\item
  Heat shock response.
\item
  Oxidative stress.
\end{itemize}

NKG2D ligands are MHC-I-like in fold but lack peptide-binding groove and
\texorpdfstring{\ensuremath{\beta}}{beta}2-microglobulin association. Their induction provides a ``danger
signal'' that bypasses normal MHC-I inhibition. NKG2D engagement signals
through DAP10/DAP12 adaptors to activate cytotoxicity. Many cancers shed
soluble MICA/MICB into circulation as a decoy that downmodulates surface
NKG2D --- an immune escape mechanism now targeted by anti-MICA-shedding
antibodies in cancer immunotherapy.

\textbf{NK-cell killing mechanisms} (shared with CTLs):

\begin{itemize}
\tightlist
\item
  \textbf{Perforin/granzyme} --- perforin polymerizes in target membrane
  (similar to MAC); granzymes (especially granzyme B) enter and cleave
  caspase-3/7 \texorpdfstring{\ensuremath{\to}}{->} \hyperref[gl:apoptosis]{\textbf{apoptosis}}.
\item
  \textbf{Death receptor ligands} --- FasL, TRAIL on NK cell engage Fas,
  DR4/DR5 on target \texorpdfstring{\ensuremath{\to}}{->} DISC \texorpdfstring{\ensuremath{\to}}{->} caspase-8 \texorpdfstring{\ensuremath{\to}}{->} apoptosis.
\item
  \textbf{Antibody-dependent cellular cytotoxicity (ADCC)}: NK cell
  \textbf{CD16 (Fc\texorpdfstring{\ensuremath{\gamma}}{gamma}RIII)} binds IgG bound to target cell surface \texorpdfstring{\ensuremath{\to}}{->}
  directed degranulation. ADCC is the principal mechanism by which
  anti-tumor antibodies (rituximab, trastuzumab) kill cancer cells, and
  a substantial fraction of vaccine-induced antiviral protection.
\end{itemize}

\subsubsection{Inflammation}\label{inflammation}

The cardinal signs of inflammation (rubor, calor, tumor, dolor --
redness, heat, swelling, pain) result from vascular changes triggered by
innate immune activation:

\begin{itemize}
\tightlist
\item
  \textbf{Mast cell degranulation}: Histamine, tryptase, leukotrienes
  -\textgreater{} vasodilation + increased vascular permeability
\item
  \textbf{Cytokine cascade}: IL-1\texorpdfstring{\ensuremath{\beta}}{beta} (fever, endothelial activation);
  TNF-\texorpdfstring{\ensuremath{\alpha}}{alpha} (fever, hypotension, acute-phase response); IL-6 (hepatic
  acute-phase proteins: CRP, fibrinogen, complement, ferritin);
  IL-8/CXCL8 (neutrophil chemotaxis)
\item
  \textbf{Acute-phase response}: Liver produces C-reactive protein
  (CRP), which binds phosphocholine on bacterial surfaces and activates
  complement
\end{itemize}

\begin{quote}
\textbf{Clinical Connection: Sepsis and Cytokine Storm} Sepsis is
defined as life-threatening organ dysfunction caused by a dysregulated
host response to infection (Sepsis-3 criteria, 2016). It affects
approximately 49 million people and causes 11 million deaths annually
worldwide (WHO). In sepsis, the normally protective inflammatory
response becomes pathologically amplified: massive cytokine release
(TNF-\texorpdfstring{\ensuremath{\alpha}}{alpha}, IL-1\texorpdfstring{\ensuremath{\beta}}{beta}, IL-6) causes widespread endothelial dysfunction,
capillary leak, coagulopathy (DIC), and multi-organ failure. The SOFA
(Sequential Organ Failure Assessment) score quantifies organ
dysfunction. Cytokine storm -- a related phenomenon -- was observed in
severe COVID-19 (IL-6 pathway hyperactivation; treated with tocilizumab,
an IL-6 receptor antagonist) and in cytokine release syndrome (CRS)
following CAR-T cell therapy.
\end{quote}

\begin{quote}
\textbf{Concept Check 2:} A child with chronic granulomatous disease
(CGD) has a mutation in the gp91phox subunit of NADPH oxidase. Predict
which types of infections this child will be susceptible to, and explain
why catalase-positive bacteria (\emph{S. aureus}, \emph{Aspergillus})
are particularly dangerous in CGD while catalase-negative bacteria
(\emph{Streptococcus}) are not.
\end{quote}

\begin{quote}
\textbf{Concept Check 2b:} A patient with paroxysmal nocturnal
haemoglobinuria (PNH) has a somatic mutation in \emph{PIGA} causing loss
of GPI-anchored proteins on red blood cells, including CD55 and CD59.
Explain (a) why these red cells are constantly haemolysed by their own
complement system, (b) why the disease is paroxysmal (worse at
night/with sleep), and (c) why eculizumab (anti-C5 monoclonal antibody)
is curative. Which complement step does eculizumab block, and what
infectious-disease vaccination is mandatory before starting it?
\end{quote}

\begin{center}\rule{0.5\linewidth}{0.5pt}\end{center}

\subsection{Adaptive Immunity
Overview}\label{adaptive-immunity-overview}

The adaptive immune system provides antigen-specific responses with
immunological memory. Its defining features include:

\begin{itemize}
\tightlist
\item
  \textbf{Specificity}: Each lymphocyte clone bears a unique antigen
  receptor (TCR or BCR) generated by somatic recombination
\item
  \textbf{Diversity}: \(>10^{12}\) possible TCR specificities; \(>10^8\)
  BCR specificities
\item
  \textbf{Clonal selection}: Antigen selects and expands primarily those
  lymphocyte clones with matching receptors
\item
  \textbf{Memory}: Long-lived memory cells mount faster, stronger
  secondary responses upon re-encounter
\item
  \textbf{Self-tolerance}: Autoreactive clones are eliminated (central
  tolerance) or suppressed (peripheral tolerance)
\end{itemize}

\subsubsection{MHC Molecules}\label{mhc-molecules}

Major histocompatibility complex (MHC) molecules present peptide
fragments to T cells, providing the context for antigen recognition:

\textbf{MHC class I} (HLA-A, HLA-B, HLA-C in humans):

\begin{itemize}
\tightlist
\item
  Expressed on \textbf{most nucleated cells}
\item
  Present \textbf{endogenous peptides} (from cytoplasmic proteins):
  proteasome -\textgreater{} peptide fragments -\textgreater{} TAP
  transporter -\textgreater{} ER -\textgreater{} peptide loading onto
  MHC-I -\textgreater{} cell surface
\item
  Recognized by \textbf{CD8\(^+\) T cells} (cytotoxic T lymphocytes)
\item
  Presents viral antigens and tumor antigens from within the cell
\end{itemize}

\textbf{MHC class II} (HLA-DR, HLA-DP, HLA-DQ):

\begin{itemize}
\tightlist
\item
  Expressed on \textbf{professional antigen-presenting cells (APCs)}:
  dendritic cells, macrophages, B cells
\item
  Present \textbf{exogenous peptides} (from phagocytosed material):
  endosomal/lysosomal degradation -\textgreater{} CLIP removal by HLA-DM
  -\textgreater{} peptide loading onto MHC-II -\textgreater{} cell
  surface
\item
  Recognized by \textbf{CD4\(^+\) T cells} (helper T cells)
\end{itemize}

MHC genes are the most polymorphic in the human genome
(about 20,000 HLA \hyperref[gl:allele]{\textbf{allele}}s in the
population), ensuring that no pathogen peptide can escape presentation
in most individuals -- a population-level defense strategy.

\subsubsection{T Cell Activation and
Differentiation}\label{t-cell-activation-and-differentiation}

\begin{figure}[htbp]
\centering
\includegraphics[width=0.82\linewidth,height=4.2in,keepaspectratio]{../figures/mermaid_inline/inline_mermaid_0130_02191e8ef71e.png}
\caption{Three-signal T cell activation by dendritic cells, with cytokine-context-driven differentiation into Th1/Th2/Th17/Treg/Tfh subsets and downstream activation of cytotoxic CD8\$\textasciicircum{}+\$ T cells and antibody-producing B cells.}
\end{figure}

\begin{itemize}
\tightlist
\item
  \textbf{Signal 1}: TCR recognition of peptide:MHC complex (specificity
  signal)
\item
  \textbf{Signal 2}: Co-stimulatory signal -- CD28 on T cell binds B7
  molecules (CD80/CD86) on APC. Without this signal, the T cell becomes
  \textbf{anergic} (functionally unresponsive) -- a mechanism of
  peripheral tolerance
\item
  \textbf{Signal 3}: Cytokine environment determines effector subset
  differentiation
\end{itemize}

\textbf{T helper cell subsets}:

{\def\LTcaptype{none} % do not increment counter
\begin{longtable}[]{@{}
  >{\raggedright\arraybackslash}p{(\linewidth - 8\tabcolsep) * \real{0.0889}}
  >{\raggedright\arraybackslash}p{(\linewidth - 8\tabcolsep) * \real{0.2111}}
  >{\raggedright\arraybackslash}p{(\linewidth - 8\tabcolsep) * \real{0.3000}}
  >{\raggedright\arraybackslash}p{(\linewidth - 8\tabcolsep) * \real{0.2111}}
  >{\raggedright\arraybackslash}p{(\linewidth - 8\tabcolsep) * \real{0.1889}}@{}}
\toprule\noalign{}
\begin{minipage}[b]{\linewidth}\raggedright
Subset
\end{minipage} & \begin{minipage}[b]{\linewidth}\raggedright
Inducing Cytokines
\end{minipage} & \begin{minipage}[b]{\linewidth}\raggedright
Master Transcription Factor
\end{minipage} & \begin{minipage}[b]{\linewidth}\raggedright
Signature Cytokines
\end{minipage} & \begin{minipage}[b]{\linewidth}\raggedright
Primary Function
\end{minipage} \\
\midrule\noalign{}
\endhead
\bottomrule\noalign{}
\endlastfoot
Th1 & IL-12, IFN-\texorpdfstring{\ensuremath{\gamma}}{gamma} & T-bet & IFN-\texorpdfstring{\ensuremath{\gamma}}{gamma}, TNF-\texorpdfstring{\ensuremath{\alpha}}{alpha} & Macrophage activation;
intracellular pathogens \\
Th2 & IL-4 & GATA3 & IL-4, IL-5, IL-13 & B cell IgE switching;
eosinophils; helminths; allergy \\
Th17 & IL-6 + TGF-\texorpdfstring{\ensuremath{\beta}}{beta} + IL-23 & ROR\texorpdfstring{\ensuremath{\gamma}}{gamma}t & IL-17, IL-22 & Neutrophil
recruitment; mucosal defense; fungi \\
Treg & TGF-\texorpdfstring{\ensuremath{\beta}}{beta} & Foxp3 & IL-10, TGF-\texorpdfstring{\ensuremath{\beta}}{beta} & Suppress excessive inflammation;
maintain tolerance \\
Tfh & IL-6, IL-21 & Bcl6 & IL-21, CXCL13 & B cell help in germinal
centers \\
\end{longtable}
}

\textbf{Cytotoxic T lymphocytes (CTLs, CD8\(^+\))} recognize viral or
tumor antigens on MHC-I and kill target cells through:

\begin{itemize}
\tightlist
\item
  \textbf{Perforin/granzyme pathway}: Perforin polymerizes in the target
  membrane, forming pores; granzyme B enters and activates caspase-3/7
  -\textgreater{} apoptosis
\item
  \textbf{Fas-FasL pathway}: CTL-expressed FasL binds Fas (CD95) on
  target -\textgreater{} DISC formation -\textgreater{} caspase-8
  -\textgreater{} apoptosis
\item
  \textbf{Serial killing}: A single CTL can kill 2-25 target cells
  sequentially via directed degranulation at the immunological synapse
\item
  \textbf{Exhaustion}: In chronic infections (HBV, HCV, HIV), persistent
  antigen exposure leads to upregulation of inhibitory receptors (PD-1,
  TIM-3, LAG-3) and progressive loss of effector function -- the basis
  for PD-1/PD-L1 checkpoint inhibitor therapy in cancer
\end{itemize}

\begin{quote}
\textbf{Concept Check 3:} Explain why a patient with a CD4\(^+\) T cell
count below 200 cells/\texorpdfstring{\ensuremath{\mu}}{mu}L (as in AIDS) is susceptible to opportunistic
infections that healthy individuals easily control. Specifically,
describe which arm of the immune response is most compromised and why
both cellular and humoral immunity are affected despite B cells and
CD8\(^+\) T cells being present.
\end{quote}

\begin{center}\rule{0.5\linewidth}{0.5pt}\end{center}

\subsection{B Cells and Antibodies}\label{b-cells-and-antibodies}

\subsubsection{B Cell Activation}\label{b-cell-activation}

B cell activation pathways differ depending on the nature of the
antigen:

\textbf{T-dependent antigens} (protein antigens): BCR binds and
internalizes antigen -\textgreater{} processes and presents peptides on
MHC-II -\textgreater{} recognized by cognate Tfh cell -\textgreater{}
CD40L-CD40 interaction + IL-21 -\textgreater{} B cell activation,
proliferation, and entry into the germinal center reaction. This pathway
produces high-affinity antibodies, class-switched isotypes, and
long-lived memory.

\textbf{T-independent antigens} (repetitive polysaccharides, LPS):
Cross-link multiple BCRs on the B cell surface -\textgreater{}
activation without T cell help. Produces primarily IgM; limited affinity
maturation; poor memory response. This is why polysaccharide vaccines
(pneumococcal, meningococcal) are poorly immunogenic in children under 2
years (immature T-independent response) and why \textbf{conjugate
vaccines} were developed.

\subsubsection{Germinal Center Reaction}\label{germinal-center-reaction}

The germinal center (GC) in secondary lymphoid organs is the site of
antibody optimization:

\begin{enumerate}
\def\labelenumi{\arabic{enumi}.}
\tightlist
\item
  \textbf{Clonal expansion}: Activated B cells (centroblasts)
  proliferate rapidly in the dark zone
\item
  \textbf{Somatic hypermutation (SHM)}: Activation-induced cytidine
  deaminase (AID) introduces point mutations in the variable (V) regions
  of immunoglobulin genes at a rate of about \(10^{-3}\) per
  base pair per division -- one million-fold higher than the normal
  mutation rate. AID deaminates cytidine to uridine in DNA, which is
  then processed by base excision repair or mismatch repair to generate
  diverse mutations.
\item
  \textbf{Affinity maturation}: Mutant B cells (centrocytes) migrate to
  the light zone, where they compete for antigen displayed on follicular
  dendritic cells (FDCs). Primarily B cells with improved BCR affinity
  receive survival signals; the rest undergo apoptosis. This Darwinian
  selection process progressively increases antibody affinity over
  successive rounds.
\item
  \textbf{Class switch recombination (CSR)}: AID also introduces
  double-strand breaks in switch (S) regions upstream of constant region
  genes, enabling recombination that changes the antibody class from IgM
  to IgG, IgA, or IgE while preserving antigen specificity. Cytokine
  signals determine which class is selected: IL-4 -\textgreater{} IgE;
  IFN-\texorpdfstring{\ensuremath{\gamma}}{gamma} -\textgreater{} IgG1; TGF-\texorpdfstring{\ensuremath{\beta}}{beta} -\textgreater{} IgA.
\end{enumerate}

The GC reaction produces two critical outputs: \textbf{long-lived plasma
cells} (migrate to bone marrow; secrete high-affinity antibody for years
to decades) and \textbf{memory B cells} (reside in secondary lymphoid
organs; rapidly reactivate upon antigen re-encounter).

\subsubsection{Antibody Structure and
Function}\label{antibody-structure-and-function}

An antibody molecule consists of two identical heavy (H) chains and two
identical light (L) chains, each containing variable (V) and constant
(C) domains:

\begin{itemize}
\tightlist
\item
  \textbf{Fab region} (fragment antigen-binding): Contains the variable
  domains (VH + VL) that form the antigen-binding site (paratope). The
  three complementarity-determining regions (CDR1, CDR2, CDR3) within
  each variable domain make direct contact with antigen.
\item
  \textbf{Fc region} (fragment crystallizable): The constant domains
  that determine antibody class and mediate effector functions
  (complement activation, Fc receptor binding on phagocytes and NK
  cells, placental transfer).
\item
  \textbf{Hinge region}: Provides flexibility between Fab and Fc,
  allowing bivalent binding to spatially separated epitopes.
\end{itemize}

\subsubsection{Antibody Classes}\label{antibody-classes}

{\def\LTcaptype{none} % do not increment counter
\begin{longtable}[]{@{}
  >{\raggedright\arraybackslash}p{(\linewidth - 6\tabcolsep) * \real{0.2045}}
  >{\raggedright\arraybackslash}p{(\linewidth - 6\tabcolsep) * \real{0.2500}}
  >{\raggedright\arraybackslash}p{(\linewidth - 6\tabcolsep) * \real{0.2273}}
  >{\raggedright\arraybackslash}p{(\linewidth - 6\tabcolsep) * \real{0.3182}}@{}}
\toprule\noalign{}
\begin{minipage}[b]{\linewidth}\raggedright
Isotype
\end{minipage} & \begin{minipage}[b]{\linewidth}\raggedright
Structure
\end{minipage} & \begin{minipage}[b]{\linewidth}\raggedright
Location
\end{minipage} & \begin{minipage}[b]{\linewidth}\raggedright
Key Functions
\end{minipage} \\
\midrule\noalign{}
\endhead
\bottomrule\noalign{}
\endlastfoot
\textbf{IgM} & Pentamer (10 binding sites) & Blood (primary response) &
First antibody produced; efficient complement activation via classical
pathway; low affinity compensated by high avidity \\
\textbf{IgG} & Monomer (4 subclasses) & Blood, extravascular;
\textbf{crosses placenta} via FcRn & Most abundant serum antibody;
opsonization; ADCC (via CD16 on NK cells); complement; neonatal passive
immunity \\
\textbf{IgA} & Dimer (secretory IgA with secretory component) & Mucosal
surfaces: saliva, tears, breast milk, intestinal lumen & Most produced
antibody class (\textgreater3 g/day); neutralization at mucosal
surfaces; immune exclusion \\
\textbf{IgE} & Monomer & Bound to mast cells and basophils via
high-affinity Fc\texorpdfstring{\ensuremath{\varepsilon}}{epsilon}RI & Very low serum concentration; cross-linking by
allergen -\textgreater{} mast cell degranulation -\textgreater{}
immediate hypersensitivity; anti-helminth defense (eosinophil ADCC) \\
\textbf{IgD} & Monomer & Naive B cell surface (co-expressed with IgM) &
BCR signaling; poorly understood effector function \\
\end{longtable}
}

\textbf{Immunological memory} -- the basis of vaccination -- results
from the persistence of long-lived plasma cells and memory B cells. Upon
re-exposure, the secondary response is:

\begin{itemize}
\tightlist
\item
  \textbf{Faster}: Days instead of 1-2 weeks
\item
  \textbf{Stronger}: 10-100 fold higher antibody titers
\item
  \textbf{Higher affinity}: Memory B cells express somatically
  hypermutated, affinity-matured BCRs
\item
  \textbf{Class-switched}: Predominantly IgG (or IgA at mucosal
  surfaces)
\end{itemize}

\begin{quote}
\textbf{Concept Check 4:} Conjugate vaccines (e.g., PCV13) link
bacterial polysaccharide antigens to protein carriers (e.g., CRM197, a
nontoxic diphtheria toxoid). Explain the immunological rationale: why
does conjugation convert a T-independent antigen into a T-dependent
antigen, and why is this critical for vaccinating infants?
\end{quote}

\begin{center}\rule{0.5\linewidth}{0.5pt}\end{center}

\subsection{Vaccines and Immunological
Memory}\label{vaccines-and-immunological-memory}

\subsubsection{Vaccine Platforms: Eight Strategies for Inducing
Memory}\label{vaccine-platforms-eight-strategies-for-inducing-memory}

Vaccines exploit the same memory mechanisms that protect us after
natural infection, but without the disease. The eight major platform
classes --- each with distinct strengths, limitations, and clinical
exemplars --- represent a 200-year-progressive enrichment of strategies
for safely tricking the immune system into making memory.

\begin{figure}[htbp]
\centering
\includegraphics[width=0.82\linewidth,height=4.2in,keepaspectratio]{../figures/mermaid_inline/inline_mermaid_0131_76eee62eb18a.png}
\caption{Eight vaccine platforms, the immune mechanism each engages, and characteristic clinical responses.}
\end{figure}

{\def\LTcaptype{none} % do not increment counter
\begin{longtable}[]{@{}
  >{\raggedright\arraybackslash}p{(\linewidth - 8\tabcolsep) * \real{0.1724}}
  >{\raggedright\arraybackslash}p{(\linewidth - 8\tabcolsep) * \real{0.1724}}
  >{\raggedright\arraybackslash}p{(\linewidth - 8\tabcolsep) * \real{0.2414}}
  >{\raggedright\arraybackslash}p{(\linewidth - 8\tabcolsep) * \real{0.1897}}
  >{\raggedright\arraybackslash}p{(\linewidth - 8\tabcolsep) * \real{0.2241}}@{}}
\toprule\noalign{}
\begin{minipage}[b]{\linewidth}\raggedright
Platform
\end{minipage} & \begin{minipage}[b]{\linewidth}\raggedright
Examples
\end{minipage} & \begin{minipage}[b]{\linewidth}\raggedright
Antigen form
\end{minipage} & \begin{minipage}[b]{\linewidth}\raggedright
Strengths
\end{minipage} & \begin{minipage}[b]{\linewidth}\raggedright
Limitations
\end{minipage} \\
\midrule\noalign{}
\endhead
\bottomrule\noalign{}
\endlastfoot
\textbf{Live attenuated} & MMR, varicella (Varivax), yellow fever 17D,
oral polio (OPV), BCG, rotavirus & Weakened replicating organism &
Strong T + B memory; often a single dose; mucosal immunity (OPV) &
Reversion to virulence (rare; OPV \texorpdfstring{\ensuremath{\to}}{->} cVDPV); contraindicated in
immunocompromised and pregnancy \\
\textbf{Inactivated / killed} & IPV (Salk), inactivated influenza (IIV),
hepatitis A (Havrix), rabies, whole-cell pertussis (legacy) &
Chemically/heat-killed whole organism & Safe in immunocompromised;
stable & Weaker than live; requires adjuvant + multiple doses \\
\textbf{Subunit / recombinant protein} & Hepatitis B (HBsAg in yeast),
acellular pertussis (aP), zoster (Shingrix gE + AS01) & Purified protein
antigen + adjuvant & Very safe; no infectious risk & Less immunogenic;
adjuvant essential \\
\textbf{Virus-like particles (VLPs)} & HPV (Gardasil, Cervarix), HBV
(also classified as subunit) & Self-assembled capsid proteins (no
genome) & Particulate, highly immunogenic; structurally identical to
virus & Difficult to engineer for some viruses \\
\textbf{Toxoid} & Tetanus (TT), diphtheria (DT, Tdap), botulism toxoid &
Formalin-inactivated toxin & Anti-toxin antibodies neutralise the
\emph{toxin}, not the bacterium & No antibacterial protection; periodic
boosters needed \\
\textbf{Conjugate} & Hib, PCV13, PCV20 (pneumococcal), meningococcal
ACWY (Menactra) & Bacterial polysaccharide covalently linked to a
protein carrier (CRM197, TT) & Converts T-independent \texorpdfstring{\ensuremath{\to}}{->} T-dependent
response; enables infant immunisation; affinity-matured response &
Carrier-specific T cells diverted; serotype replacement after PCV
introduction \\
\textbf{mRNA} & BNT162b2 (Pfizer-BioNTech), mRNA-1273 (Moderna),
mRNA-1345 (RSV, 2024); flu mRNA-1010 (trial) & LNP-encapsulated
nucleoside-modified mRNA encoding the antigen & Rapid design (48 hr from
sequence); strong T + B response; no genome integration & Cold-chain
dependent; rare myocarditis; waning antibody titres \\
\textbf{Viral vector} & rVSV-ZEBOV (Ebola), ChAdOx1 (AstraZeneca COVID),
Ad26.COV2.S (J\&J), gene-therapy AAV vectors & Non-replicating
recombinant virus carrying antigen gene & Strong durable response;
mucosal IgA achievable & Pre-existing vector immunity (Ad5 in adult
populations) reduces efficacy; rare thrombosis (ChAdOx1) \\
\end{longtable}
}

\textbf{mRNA vaccines} are the newest platform and represented a
paradigm shift during the COVID-19 pandemic:

\begin{itemize}
\tightlist
\item
  The spike protein sequence of SARS-CoV-2 was published on January 11,
  2020; Moderna had designed its mRNA-1273 vaccine candidate within 48
  hours.
\item
  The mRNA is modified (N1-methylpseudouridine replaces uridine, Karikó
  and Weissman discovery, Nobel Prize 2023) to reduce TLR7/8
  recognition, block 2'-5'-OAS / RNase L degradation, and increase
  translational efficiency 10--100\texorpdfstring{\ensuremath{\times}}{x}.
\item
  Lipid nanoparticle (LNP) encapsulation protects the mRNA from
  extracellular nucleases and facilitates cellular uptake; the
  four-component LNP (ionisable lipid, phospholipid, cholesterol,
  PEG-lipid) self-assembles into about  80-nm particles.
\item
  Upon injection, host cells (predominantly muscle and dendritic cells
  at the injection site) translate the mRNA into spike protein, which is
  presented on MHC-I (activating CD8\(^+\) CTLs) and
  secreted/surface-displayed (activating B cells and CD4\(^+\) T cells).
\item
  No integration into host DNA occurs --- mRNA is degraded within hours
  to days.
\item
  Spike is locked in the prefusion conformation by two proline
  substitutions (\textbf{2P mutation}) that preserve neutralising
  epitopes.
\end{itemize}

The platform is now extending to \textbf{influenza} (mRNA-1010),
\textbf{RSV} (mRNA-1345, FDA-approved 2024), \textbf{CMV}, and
\textbf{personalised cancer neoantigen vaccines} (Moderna mRNA-4157 in
Phase III for melanoma).

\subsubsection{Antigenic Variation: How Pathogens Outpace Antibody
Responses}\label{antigenic-variation-how-pathogens-outpace-antibody-responses}

The single greatest challenge for vaccine development against many
viruses is that the immune target itself \textbf{evolves under the
selection pressure of the population's own immune response}. Two
paradigmatic examples illustrate the strategies and the resulting
public-health consequences.

\textbf{Influenza A: drift and shift.} Influenza A is a Class V
segmented (-)ssRNA virus with 8 genome segments. Its surface
glycoproteins \textbf{hemagglutinin (HA, 18 subtypes)} and
\textbf{neuraminidase (NA, 11 subtypes)} are the primary targets of
neutralising antibodies. Influenza A undergoes two distinct mechanisms
of antigenic change:

{\def\LTcaptype{none} % do not increment counter
\begin{longtable}[]{@{}
  >{\raggedright\arraybackslash}p{(\linewidth - 4\tabcolsep) * \real{0.3333}}
  >{\raggedright\arraybackslash}p{(\linewidth - 4\tabcolsep) * \real{0.3333}}
  >{\raggedright\arraybackslash}p{(\linewidth - 4\tabcolsep) * \real{0.3333}}@{}}
\toprule\noalign{}
\begin{minipage}[b]{\linewidth}\raggedright
Feature
\end{minipage} & \begin{minipage}[b]{\linewidth}\raggedright
\textbf{Antigenic drift}
\end{minipage} & \begin{minipage}[b]{\linewidth}\raggedright
\textbf{Antigenic shift}
\end{minipage} \\
\midrule\noalign{}
\endhead
\bottomrule\noalign{}
\endlastfoot
Mechanism & Point mutations in HA (and NA) accumulating during
replication & \textbf{Reassortment} of whole genome segments during
co-infection of one cell by two influenza A strains (often in pigs as
the ``mixing vessel'') \\
Rate & Continuous; about  0.5--1 \% per year on HA1 &
Episodic; rare events \\
Consequence & Antigenic mismatch with prior immunity \texorpdfstring{\ensuremath{\to}}{->} \textbf{seasonal
epidemics} & Novel HA subtype unseen by population \texorpdfstring{\ensuremath{\to}}{->} \textbf{pandemic
potential} \\
Vaccine response & Annual reformulation (WHO twice-yearly
recommendation) & Requires entirely new pandemic vaccine \\
Why influenza A primarily? & Influenza B drifts but does not reassort
with animal reservoirs at meaningful frequency; influenza A has a wide
animal host range & Influenza A naturally infects birds, swine, humans,
and other mammals --- sustaining a reservoir of segments for
reassortment \\
\end{longtable}
}

\textbf{Pandemic history:} - \textbf{1918 H1N1} (``Spanish flu'') ---
50--100 million deaths. - \textbf{1957 H2N2} (``Asian flu'') --- 1--2
million deaths; H2 from avian source reassorted with circulating H1N1. -
\textbf{1968 H3N2} (``Hong Kong flu'') --- 1 million deaths; H3 from
avian source. - \textbf{2009 H1N1} (pandemic) --- about 
284,000 deaths; quadruple reassortant (avian + human + two swine
lineages); milder than predicted because of partial cross-immunity from
related H1N1 strains. - \textbf{H5N1 highly pathogenic avian influenza
(HPAI)} --- ongoing outbreaks in dairy cattle and poultry (2024--26);
about  60 \% case-fatality rate in the rare human cases;
intense surveillance for adaptation to mammalian transmission.

\textbf{HIV reverse transcriptase: the quasi-species cloud.} HIV uses an
entirely different mechanism --- extreme intra-host evolution. Reverse
transcriptase has an \textbf{error rate of about  3 \texorpdfstring{\ensuremath{\times}}{x} 10\texorpdfstring{\textsuperscript{-}}{-}\texorpdfstring{\textsuperscript{5}}{5}
per base per replication cycle}, and HIV makes about 
\(10^{10}\) virions per day in an untreated patient. With a 9.7-kb
genome, this produces \textbf{every single point mutation in the viral
genome multiple times every day} and a substantial fraction of possible
double mutations. The resulting population --- a cloud of related
sequences around a consensus --- is termed a \textbf{quasi-species}.

Consequences:

\begin{enumerate}
\def\labelenumi{\arabic{enumi}.}
\tightlist
\item
  \textbf{Drug resistance is essentially preformed}: any single-target
  monotherapy fails within weeks because resistant variants pre-exist in
  the quasi-species cloud. This is why HIV treatment requires
  \textbf{combination ART} (\texorpdfstring{\textgreater{}=}{>=} 3 drugs from \texorpdfstring{\textgreater{}=}{>=} 2 classes); the
  simultaneous probability of a virion carrying mutations conferring
  resistance to most drugs is astronomically low.
\item
  \textbf{Antibody escape is continuous}: the env gene (gp120/gp41)
  evolves rapidly under neutralising-antibody pressure, with the
  variable loops V1/V2/V3 changing every few weeks within an individual.
\item
  \textbf{Vaccine difficulty}: a vaccine that elicits primarily
  narrowly-specific antibodies will fail. Successful HIV vaccine
  candidates aim for \textbf{broadly neutralising antibodies (bNAbs)}
  that target conserved sites (CD4-binding site, fusion peptide) --- but
  these require unusual germline B cells and prolonged affinity
  maturation, which most vaccinees do not develop.
\end{enumerate}

The HIV quasi-species lesson --- that within-host evolution can generate
resistance faster than therapy can clear the virus --- has now been
applied to influenza, hepatitis C, and SARS-CoV-2 (the 32-spike-mutation
Omicron emergence almost certainly arose from chronic infection of an
immunocompromised host, where prolonged replication permitted
accumulation of escape mutations).

\subsubsection{Herd Immunity}\label{herd-immunity}

When a sufficient proportion of a population is immune, the pathogen's
transmission is interrupted even for non-immune individuals. The
\textbf{herd-immunity threshold} is given by:

\begin{equation}
p_c = 1 - \frac{1}{R_0}
\label{eq:unit_VII_herd_immunity}
\end{equation}

where \(p_c\) is the critical fraction of the population that must be
immune and \(R_0\) is the \textbf{basic reproduction number} (the
expected number of secondary cases from one primary case in a fully
susceptible population). The intuition: each case must produce on
average less than one secondary case for the epidemic to die out;
immunity removes a fraction of the contacts an infectious person makes,
so the \textbf{effective reproduction number} is \(R_e = R_0 (1 - p)\).
Solving for \(R_e = 1\) gives \(p = 1 - 1/R_0\).

\subsubsection{Worked Example: Calculating Herd Immunity Threshold
(Three
Pathogens)}\label{worked-example-calculating-herd-immunity-threshold-three-pathogens}

\textbf{Problem:} Compute the herd immunity threshold from
\cref{eq:unit_VII_herd_immunity} for measles (\(R_0 = 15\)), polio
(\(R_0 = 5\)), and original-strain SARS-CoV-2 (\(R_0 = 2.5\)), and
discuss the implications for vaccination policy.

\textbf{Solution:}

{\def\LTcaptype{none} % do not increment counter
\begin{longtable}[]{@{}
  >{\raggedright\arraybackslash}p{(\linewidth - 8\tabcolsep) * \real{0.0962}}
  >{\raggedright\arraybackslash}p{(\linewidth - 8\tabcolsep) * \real{0.0673}}
  >{\raggedright\arraybackslash}p{(\linewidth - 8\tabcolsep) * \real{0.1827}}
  >{\raggedright\arraybackslash}p{(\linewidth - 8\tabcolsep) * \real{0.3077}}
  >{\raggedright\arraybackslash}p{(\linewidth - 8\tabcolsep) * \real{0.3462}}@{}}
\toprule\noalign{}
\begin{minipage}[b]{\linewidth}\raggedright
Pathogen
\end{minipage} & \begin{minipage}[b]{\linewidth}\raggedright
\(R_0\)
\end{minipage} & \begin{minipage}[b]{\linewidth}\raggedright
\(p_c = 1 - 1/R_0\)
\end{minipage} & \begin{minipage}[b]{\linewidth}\raggedright
Realistic vaccine efficacy \(E\)
\end{minipage} & \begin{minipage}[b]{\linewidth}\raggedright
Required coverage \(p_v = p_c / E\)
\end{minipage} \\
\midrule\noalign{}
\endhead
\bottomrule\noalign{}
\endlastfoot
\textbf{Measles} & 15 & \(1 - 1/15 = 0.933\) \texorpdfstring{\ensuremath{\to}}{->} \textbf{93.3 \%} & 0.97
(MMR two-dose) & \(0.933 / 0.97 = 0.962\) \texorpdfstring{\ensuremath{\to}}{->} \textbf{96.2 \%} \\
\textbf{Polio} & 5 & \(1 - 1/5 = 0.800\) \texorpdfstring{\ensuremath{\to}}{->} \textbf{80 \%} & 0.99 (IPV
three-dose) & \(0.800 / 0.99 = 0.808\) \texorpdfstring{\ensuremath{\to}}{->} \textbf{80.8 \%} \\
\textbf{COVID-19 (original)} & 2.5 & \(1 - 1/2.5 = 0.600\) \texorpdfstring{\ensuremath{\to}}{->} \textbf{60
\%} & 0.95 (BNT162b2 against original) & \(0.600 / 0.95 = 0.632\) \texorpdfstring{\ensuremath{\to}}{->}
\textbf{63.2 \%} \\
\textbf{COVID-19 (Omicron)} & 10 & \(1 - 1/10 = 0.900\) \texorpdfstring{\ensuremath{\to}}{->} \textbf{90 \%}
& 0.50 (vs symptomatic infection) & \(0.900 / 0.50 = 1.80\) \texorpdfstring{\ensuremath{\to}}{->}
\textbf{unattainable} \\
\end{longtable}
}

\textbf{Interpretation:}

\begin{enumerate}
\def\labelenumi{\arabic{enumi}.}
\tightlist
\item
  \textbf{Measles is the textbook case for high-coverage vaccination.}
  Even with the excellent MMR vaccine (\(E \approx 97 \%\)),
  \textgreater{} 96 \% of the population must be vaccinated. A drop from
  96 \% to 92 \% (e.g., due to vaccine hesitancy) reduces immunity below
  threshold and allows outbreaks --- exactly what has happened in the
  2018--2019 US/EU measles resurgences.
\item
  \textbf{Polio's lower \(R_0\)} makes it much more tractable: 80 \%
  coverage with IPV is sufficient. This is why polio is on the verge of
  global eradication while measles is not.
\item
  \textbf{COVID-19 Omicron} illustrates the \textbf{vaccine-coverage
  impossibility} result: when efficacy against transmission is moderate
  (about  50 \%) and \(R_0\) is high (about  10
  due to immune escape and inherent transmissibility), the formula
  \(p_v = p_c / E\) exceeds 100 \%, meaning \textbf{herd immunity by
  vaccination alone is unattainable}. Public-health strategy must rely
  on a combination of vaccination (reducing severity),
  non-pharmaceutical interventions (reducing \(R_0\) below the
  achievable coverage), and accepting endemicity.
\end{enumerate}

The herd-immunity equation also explains why \textbf{small reductions in
vaccine coverage produce disproportionate disease resurgences for
high-\(R_0\) pathogens}: at \(R_0 = 15\), a 1-percentage-point drop in
immune coverage raises the effective reproduction number \(R_e\) by
about \(R_0 \times 0.01 = 0.15\). This non-linearity is why measles
outbreaks dominate vaccine-hesitancy clusters.

\subsubsection{\texorpdfstring{R\(_0\) Comparison: Why Pathogens
Differ}{R_0 Comparison: Why Pathogens Differ}}\label{r_0-comparison-why-pathogens-differ}

{\def\LTcaptype{none} % do not increment counter
\begin{longtable}[]{@{}
  >{\raggedright\arraybackslash}p{(\linewidth - 6\tabcolsep) * \real{0.2041}}
  >{\raggedright\arraybackslash}p{(\linewidth - 6\tabcolsep) * \real{0.1429}}
  >{\raggedright\arraybackslash}p{(\linewidth - 6\tabcolsep) * \real{0.5102}}
  >{\raggedright\arraybackslash}p{(\linewidth - 6\tabcolsep) * \real{0.1429}}@{}}
\toprule\noalign{}
\begin{minipage}[b]{\linewidth}\raggedright
Pathogen
\end{minipage} & \begin{minipage}[b]{\linewidth}\raggedright
\(R_0\)
\end{minipage} & \begin{minipage}[b]{\linewidth}\raggedright
Herd-immunity threshold
\end{minipage} & \begin{minipage}[b]{\linewidth}\raggedright
Notes
\end{minipage} \\
\midrule\noalign{}
\endhead
\bottomrule\noalign{}
\endlastfoot
\textbf{Measles} & 12--18 & 92--94 \% & Most contagious common
pathogen \\
\textbf{Mumps} & 4--7 & 75--86 \% & MMR-vaccine-preventable \\
\textbf{Pertussis} & 12--17 & 92--94 \% & Resurgent due to waning
aP-vaccine immunity \\
\textbf{Diphtheria} & 6--7 & 83--86 \% & Toxin-driven;
toxoid-vaccinable \\
\textbf{Polio} & 4--7 & 75--86 \% & Near eradication \\
\textbf{Smallpox} & 5--7 & 80--86 \% & Eradicated 1980 \\
\textbf{Rubella} & 6--7 & 83--86 \% & MMR-vaccine-preventable \\
\textbf{HIV} & 2--5 & 50--80 \% & Sexual / parenteral transmission \\
\textbf{Influenza (seasonal)} & 1.3--2 & 23--50 \% & Annual drift;
vaccine reformulated yearly \\
\textbf{Influenza (1918)} & 2--3 & 50--67 \% & Pandemic \\
\textbf{SARS-CoV-2 (original)} & 2.5--3 & 60--67 \% & Wuhan strain \\
\textbf{SARS-CoV-2 (Delta)} & 5--8 & 80--88 \% & L452R + T478K spike \\
\textbf{SARS-CoV-2 (Omicron, BA.5)} & 8--15 & 88--93 \% & Immune escape
pushes effective threshold higher \\
\textbf{Mpox (clade IIb, 2022)} & 1.0--1.5 & 0--33 \% & Sustained MSM
transmission; vaccine + behaviour change controlled \\
\end{longtable}
}

\subsubsection{Adjuvants}\label{adjuvants}

Adjuvants amplify the innate immune response to vaccine antigens,
improving adaptive immunity:

\begin{itemize}
\tightlist
\item
  \textbf{Aluminum salts (Alum)}: Oldest adjuvant (since 1926);
  activates NLRP3 inflammasome; depot effect (slow antigen release);
  used in most childhood vaccines
\item
  \textbf{AS04} (Alum + monophosphoryl lipid A): TLR4 agonist; used in
  Cervarix (HPV vaccine)
\item
  \textbf{MF59}: Oil-in-water squalene emulsion; enhances antigen uptake
  by APCs; used in some influenza vaccines
\item
  \textbf{CpG oligodeoxynucleotides}: TLR9 agonist; used in Heplisav-B
  (hepatitis B vaccine)
\item
  \textbf{Matrix-M} (saponin-based): Used in Novavax COVID-19 vaccine;
  potent Th1 and CTL response
\item
  \textbf{AS01} (liposomal MPL + QS-21): Used in Shingrix (zoster) and
  RTS,S (malaria); strong Th1 response
\end{itemize}

\begin{quote}
\textbf{Clinical Connection: SARS-CoV-2 and the COVID-19 Pandemic}
SARS-CoV-2 enters cells via its spike protein binding ACE2, with TMPRSS2
protease facilitating membrane fusion. Variants of concern (Alpha,
Delta, Omicron) accumulated mutations in the spike protein's
receptor-binding domain (RBD) that increased transmissibility and/or
enabled immune evasion. Omicron-lineage viruses substantially escaped
neutralising antibodies from both vaccination and prior infection, while
population immunity and viral evolution shifted severe-disease risk
compared with early pandemic waves. mRNA vaccine development
demonstrated the power of platform technology -- Moderna designed its
vaccine sequence within 48 hours of the viral genome publication, though
regulatory approval required 11 months of clinical trials. Long COVID is
now defined operationally as a chronic condition present at least 3
months after SARS-CoV-2 infection; prevalence estimates vary by cohort,
variant era, vaccination status, and case definition. Mechanistic
studies now emphasise persistent immune activation, T-cell exhaustion
signatures, metabolic dysregulation, possible viral-antigen persistence,
autoimmunity, \hyperref[gl:microbiome]{\textbf{microbiome}} disruption,
dysautonomia, and endothelial dysfunction rather than a single cause
\citep{cdc2026longcovid,longcovid2025immune}.
\end{quote}

\begin{figure}[htbp]
\centering
\includegraphics[width=0.82\linewidth,height=4.2in,keepaspectratio]{../figures/mermaid_inline/inline_mermaid_0132_8e0f151220de.png}
\caption{Long COVID mechanism network. Current evidence supports overlapping immune, vascular, neuroautonomic, metabolic, and possible persistence mechanisms; the diagram should be read as a differential-mechanism map, not a single-cause pathway \textbackslash{}citep\{longcovid2026mechanisms\}.}
\end{figure}

\begin{quote}
\textbf{Concept Check 5:} Calculate the effective reproduction number
(\(R_e\)) for measles (\(R_0 = 15\)) in a community where 80\% of the
population is immune. Is an epidemic possible? How many additional
percentage points of immunity are needed to prevent sustained
transmission?
\end{quote}

\begin{quote}
\textbf{Concept Check 5b:} A new COVID-19 vaccine candidate has 60 \%
efficacy against transmission. The Omicron sublineage circulating has
\(R_0 = 10\). Apply \cref{eq:unit_VII_herd_immunity} to determine
whether achievable vaccine coverage can suppress transmission. If not,
what additional public-health measures (and at what magnitude of \(R_0\)
reduction) would make the situation tractable with this vaccine?
\end{quote}

\begin{quote}
\textbf{Concept Check 6:} Influenza A viruses can undergo antigenic
shift through genomic reassortment, but SARS-CoV-2 (a non-segmented RNA
virus) cannot. Despite this, SARS-CoV-2 has generated numerous variants
of concern. Explain the mechanism by which SARS-CoV-2 generates
antigenic diversity and why its approach, while different from
influenza, has still enabled significant immune evasion.
\end{quote}

\begin{center}\rule{0.5\linewidth}{0.5pt}\end{center}

\subsection{Antibiotic Resistance: A Comprehensive Mechanism
Catalogue}\label{antibiotic-resistance-a-comprehensive-mechanism-catalogue}

Antibiotic resistance can be classified into five mechanistic
categories. Most clinically important multi-drug-resistant pathogens
combine multiple mechanisms; an isolate may simultaneously hydrolyse
\texorpdfstring{\ensuremath{\beta}}{beta}-lactams, modify aminoglycosides, methylate ribosomal RNA, and pump out
fluoroquinolones.

\subsubsection{Mechanism 1: Enzymatic
Inactivation}\label{mechanism-1-enzymatic-inactivation}

Direct destruction of the drug by hydrolysis or modification.

\textbf{\texorpdfstring{\ensuremath{\beta}}{beta}-lactamases} are the most clinically important and most diverse
resistance enzymes. They hydrolyse the \texorpdfstring{\ensuremath{\beta}}{beta}-lactam ring (with the
active-site serine --- a tetrahedral covalent intermediate --- or via a
metallo-zinc mechanism) before the drug reaches its PBP target. The
Ambler classification:

{\def\LTcaptype{none} % do not increment counter
\begin{longtable}[]{@{}
  >{\raggedright\arraybackslash}p{(\linewidth - 8\tabcolsep) * \real{0.1129}}
  >{\raggedright\arraybackslash}p{(\linewidth - 8\tabcolsep) * \real{0.1774}}
  >{\raggedright\arraybackslash}p{(\linewidth - 8\tabcolsep) * \real{0.1613}}
  >{\raggedright\arraybackslash}p{(\linewidth - 8\tabcolsep) * \real{0.3871}}
  >{\raggedright\arraybackslash}p{(\linewidth - 8\tabcolsep) * \real{0.1613}}@{}}
\toprule\noalign{}
\begin{minipage}[b]{\linewidth}\raggedright
Class
\end{minipage} & \begin{minipage}[b]{\linewidth}\raggedright
Mechanism
\end{minipage} & \begin{minipage}[b]{\linewidth}\raggedright
Spectrum
\end{minipage} & \begin{minipage}[b]{\linewidth}\raggedright
Inhibitor sensitivity
\end{minipage} & \begin{minipage}[b]{\linewidth}\raggedright
Examples
\end{minipage} \\
\midrule\noalign{}
\endhead
\bottomrule\noalign{}
\endlastfoot
\textbf{Class A} (serine) & Active-site serine (Ser70) & Penicillins;
ESBL extends to 3GC & Inhibited by clavulanate, sulbactam, tazobactam,
avibactam & TEM-1, SHV-1 (penicillinases); CTX-M (the dominant ESBL
globally); KPC (carbapenemase) \\
\textbf{Class B} (metallo-\texorpdfstring{\ensuremath{\beta}}{beta}-lactamase, MBL) & Zn\(^{2+}\) active site;
hydroxide nucleophile & Penicillins, cephalosporins, and carbapenems;
monobactams such as aztreonam are structurally spared by the MBL but can
be hydrolysed by co-produced ESBLs & Inhibited by EDTA in vitro;
\textbf{not inhibited by clinical inhibitors}; aztreonam-avibactam
combinations & \textbf{NDM-1} (New Delhi metallo-\texorpdfstring{\ensuremath{\beta}}{beta}-lactamase, 2008
emergence; global spread); \textbf{VIM}, \textbf{IMP} \\
\textbf{Class C} (AmpC) & Active-site serine & Cephalosporins (including
3GC); poorly inhibited by clavulanate & Cefepime, carbapenems retain
activity & Chromosomal AmpC inducible in \emph{Enterobacter},
\emph{Serratia}, \emph{Citrobacter}, \emph{Pseudomonas}; plasmid-borne
(CMY) \\
\textbf{Class D} (oxacillinase, OXA) & Active-site serine & Variable
(penicillins \texorpdfstring{\ensuremath{\to}}{->} carbapenems for OXA-48) & Variable inhibitor
susceptibility & OXA-48 (Klebsiella, Mediterranean); OXA-23
(Acinetobacter pan-resistance) \\
\end{longtable}
}

\textbf{Other inactivating enzymes:}

\begin{itemize}
\tightlist
\item
  \textbf{Aminoglycoside-modifying enzymes (AMEs)} --- three families:
  acetyltransferases (AAC), phosphotransferases (APH),
  nucleotidyltransferases (ANT). Modify amino or hydroxyl groups on the
  drug, abolishing ribosome binding. Plasmid-encoded; \textgreater{} 50
  enzyme variants known.
\item
  \textbf{Chloramphenicol acetyltransferase (CAT)} --- acetylates the
  drug.
\item
  \textbf{Macrolide esterases / phosphotransferases} (Ere, Mph) --- less
  common than ribosomal methylation.
\item
  \textbf{\texorpdfstring{\ensuremath{\beta}}{beta}-glucuronidases} that activate the cancer drug irinotecan
  into its toxic form (gut bacteria; relevant for drug-drug-microbiome
  interactions).
\end{itemize}

\subsubsection{Mechanism 2: Target
Modification}\label{mechanism-2-target-modification}

Altering the drug target so the antibiotic no longer binds.

{\def\LTcaptype{none} % do not increment counter
\begin{longtable}[]{@{}
  >{\raggedright\arraybackslash}p{(\linewidth - 6\tabcolsep) * \real{0.1739}}
  >{\raggedright\arraybackslash}p{(\linewidth - 6\tabcolsep) * \real{0.2464}}
  >{\raggedright\arraybackslash}p{(\linewidth - 6\tabcolsep) * \real{0.2609}}
  >{\raggedright\arraybackslash}p{(\linewidth - 6\tabcolsep) * \real{0.3188}}@{}}
\toprule\noalign{}
\begin{minipage}[b]{\linewidth}\raggedright
Antibiotic
\end{minipage} & \begin{minipage}[b]{\linewidth}\raggedright
Original target
\end{minipage} & \begin{minipage}[b]{\linewidth}\raggedright
Modified target
\end{minipage} & \begin{minipage}[b]{\linewidth}\raggedright
Resistance mechanism
\end{minipage} \\
\midrule\noalign{}
\endhead
\bottomrule\noalign{}
\endlastfoot
\textbf{\texorpdfstring{\ensuremath{\beta}}{beta}-lactams} & PBP2 (transpeptidase) & \textbf{PBP2a} in MRSA
(encoded by \emph{mecA}) & Low-affinity transpeptidase that performs
cross-linking even when normal PBPs are inhibited; SCC\emph{mec}
cassette mobile element \\
\textbf{Vancomycin} & D-Ala-D-Ala terminus of Lipid II &
\textbf{D-Ala-D-Lac} (VanA, VanB) or \textbf{D-Ala-D-Ser} (VanC) & One
H-bond replaced by an oxygen lone pair; \texorpdfstring{\textgreater{}=}{>=} 1000-fold loss of vancomycin
affinity \\
\textbf{Macrolides, lincosamides, streptogramin B (MLS\(_B\))} & 23S
rRNA peptidyl-transferase center & \textbf{Methylated A2058 of 23S rRNA}
by Erm methylase & Single methyl group blocks most three drug classes
--- cross-resistance \\
\textbf{Aminoglycosides (high-level)} & 16S rRNA A-site & \textbf{16S
rRNA methylation} by ArmA, RmtB, RmtC, RmtD & Plasmid-borne; confers
pan-aminoglycoside resistance \\
\textbf{Linezolid} & 23S rRNA, U2504 & \textbf{G2576T or T2500A 23S rRNA
mutation}; \textbf{cfr methylase} (methylates A2503) & Cfr also confers
cross-resistance to phenicols, lincosamides, streptogramins,
oxazolidinones (PhLOPS\(_A\)) \\
\textbf{Fluoroquinolones} & DNA gyrase (GyrA, GyrB), topoisomerase IV
(ParC, ParE) & \textbf{Mutations in QRDR
(quinolone-resistance-determining region)}: GyrA Ser83, Asp87 & Stepwise
mutations: each step adds resistance; \emph{qnr} genes encode
gyrase-protecting proteins (low-level) \\
\textbf{Rifampin} & RNA polymerase \texorpdfstring{\ensuremath{\beta}}{beta} subunit (RpoB) & \textbf{RpoB
mutations}, especially Ser531Leu, His526Tyr & Single mutations confer
high-level resistance \\
\textbf{Trimethoprim} & DHFR (dihydrofolate reductase) & Acquired
\emph{dfr} genes encoding drug-resistant DHFR; chromosomal mutations &
Bypass through alternative DHFR \\
\textbf{Sulfonamides} & DHPS (dihydropteroate synthase) & Acquired
\emph{sul} genes; chromosomal mutations & Bypass through alternative
DHPS \\
\textbf{Colistin (polymyxin)} & LPS lipid A & \textbf{mcr-1
phosphoethanolamine transferase} (plasmid-borne, 2015 discovery, China);
chromosomal \emph{pmrAB / phoPQ} mutations & Modifies lipid A phosphates
\texorpdfstring{\ensuremath{\to}}{->} reduced affinity for cationic colistin \\
\end{longtable}
}

\textbf{MRSA mechanism in detail.} \emph{S. aureus} resistant to
methicillin and most earlier \texorpdfstring{\ensuremath{\beta}}{beta}-lactams carries the \textbf{mecA gene} on
a mobile genetic element --- \textbf{SCC\emph{mec}} (staphylococcal
cassette chromosome mec). \emph{mecA} encodes \textbf{PBP2a}, an
alternative transpeptidase whose \texorpdfstring{\ensuremath{\beta}}{beta}-lactam-binding pocket has rearranged
to drastically lower affinity. The K\(_i\) for ceftriaxone increases
from about  0.05 \texorpdfstring{\ensuremath{\mu}}{mu}M (PBP2) to about  50 \texorpdfstring{\ensuremath{\mu}}{mu}M
(PBP2a). When normal PBPs are acylated by \texorpdfstring{\ensuremath{\beta}}{beta}-lactam, PBP2a can still
cross-link cell wall, allowing growth. Newer \textbf{5th-generation
cephalosporins} (ceftaroline, ceftobiprole) were specifically engineered
with side chains that bind PBP2a and are now used clinically against
MRSA.

\subsubsection{Mechanism 3: Efflux
Pumps}\label{mechanism-3-efflux-pumps}

Active export of antibiotics out of the cell, often accomplishing
\textbf{multi-drug resistance} through a single pump system.

{\def\LTcaptype{none} % do not increment counter
\begin{longtable}[]{@{}
  >{\raggedright\arraybackslash}p{(\linewidth - 6\tabcolsep) * \real{0.2407}}
  >{\raggedright\arraybackslash}p{(\linewidth - 6\tabcolsep) * \real{0.2778}}
  >{\raggedright\arraybackslash}p{(\linewidth - 6\tabcolsep) * \real{0.2963}}
  >{\raggedright\arraybackslash}p{(\linewidth - 6\tabcolsep) * \real{0.1852}}@{}}
\toprule\noalign{}
\begin{minipage}[b]{\linewidth}\raggedright
Pump family
\end{minipage} & \begin{minipage}[b]{\linewidth}\raggedright
Energy source
\end{minipage} & \begin{minipage}[b]{\linewidth}\raggedright
Drugs effluxed
\end{minipage} & \begin{minipage}[b]{\linewidth}\raggedright
Examples
\end{minipage} \\
\midrule\noalign{}
\endhead
\bottomrule\noalign{}
\endlastfoot
\textbf{RND (Resistance-Nodulation-Division)} & Proton motive force &
\texorpdfstring{\ensuremath{\beta}}{beta}-lactams, fluoroquinolones, tetracyclines, chloramphenicol, macrolides
& \textbf{AcrAB-TolC} (\emph{E. coli}); \textbf{MexAB-OprM},
\textbf{MexCD-OprJ}, \textbf{MexEF-OprN}, \textbf{MexXY-OprM} (\emph{P.
aeruginosa}); \textbf{AdeABC} (\emph{Acinetobacter}) \\
\textbf{MFS (Major Facilitator Superfamily)} & PMF & Tetracyclines
(TetA, TetB), macrolides (Mef), fluoroquinolones & TetA, TetB; MdfA \\
\textbf{MATE} & PMF / Na\(^+\) gradient & Fluoroquinolones,
aminoglycosides & NorM (\emph{V. parahaemolyticus}) \\
\textbf{SMR (Small Multidrug Resistance)} & PMF & QACs, lipophilic
cations & EmrE \\
\textbf{ABC (ATP-Binding Cassette)} & ATP & Macrolides, ketolides &
LmrA \\
\end{longtable}
}

\textbf{RND tripartite pumps} in Gram-negatives are particularly
devastating because they span the entire cell envelope: an
inner-membrane efflux transporter (e.g., AcrB) coupled via a periplasmic
adaptor (AcrA) to an outer-membrane channel (TolC). The pump exports
drugs straight from the cytoplasm or inner membrane into the
extracellular space, bypassing the periplasm. \emph{P. aeruginosa}
MexAB-OprM is the textbook example: constitutive expression confers
intrinsic resistance to multiple drug classes; mutations in regulators
(e.g., \emph{nfxB}, \emph{mexR}) cause overexpression and clinical
failure.

\subsubsection{Mechanism 4: Reduced
Permeability}\label{mechanism-4-reduced-permeability}

Decreasing drug uptake by altering cell-envelope permeability.

\begin{itemize}
\tightlist
\item
  \textbf{Porin loss} in Gram-negatives: deletions or down-regulation of
  OmpF, OmpC, OprD --- particularly important for hydrophilic \texorpdfstring{\ensuremath{\beta}}{beta}-lactams
  and carbapenems. \textbf{OprD loss in \emph{P. aeruginosa}} confers
  imipenem resistance.
\item
  \textbf{LPS modification} in Gram-negatives: reduces hydrophobic-drug
  entry.
\item
  \textbf{Mycobacterial cell wall}: the mycolic-acid layer is
  intrinsically impermeable; \emph{M. tuberculosis} is naturally
  resistant to most antibiotics.
\end{itemize}

\subsubsection{Mechanism 5: Bypass / Target
Overproduction}\label{mechanism-5-bypass-target-overproduction}

Producing more target or alternative target to dilute drug action.

\begin{itemize}
\tightlist
\item
  \textbf{Trimethoprim resistance via} \emph{dfr} genes that encode a
  drug-insensitive DHFR (target replacement).
\item
  \textbf{Sulfonamide resistance via} \emph{sul} genes that encode a
  drug-insensitive DHPS.
\item
  \textbf{Methotrexate resistance} in cancer cells via DHFR
  overexpression --- same principle in eukaryotes.
\end{itemize}

\subsubsection{\texorpdfstring{Resistance in Action: A Pan-Resistant
\emph{Klebsiella
pneumoniae}}{Resistance in Action: A Pan-Resistant Klebsiella pneumoniae}}\label{resistance-in-action-a-pan-resistant-klebsiella-pneumoniae}

A clinical isolate from a 2019 outbreak in a Greek ICU was reported with
the following resistance profile:

{\def\LTcaptype{none} % do not increment counter
\begin{longtable}[]{@{}
  >{\raggedright\arraybackslash}p{(\linewidth - 4\tabcolsep) * \real{0.2075}}
  >{\raggedright\arraybackslash}p{(\linewidth - 4\tabcolsep) * \real{0.3962}}
  >{\raggedright\arraybackslash}p{(\linewidth - 4\tabcolsep) * \real{0.3962}}@{}}
\toprule\noalign{}
\begin{minipage}[b]{\linewidth}\raggedright
Mechanism
\end{minipage} & \begin{minipage}[b]{\linewidth}\raggedright
Genetic determinant
\end{minipage} & \begin{minipage}[b]{\linewidth}\raggedright
Drug class affected
\end{minipage} \\
\midrule\noalign{}
\endhead
\bottomrule\noalign{}
\endlastfoot
\textbf{NDM-1} (Class B MBL) & \emph{bla}\(_\text{NDM-1}\) on plasmid &
Penicillins, cephalosporins, and carbapenems; aztreonam is spared by MBL
chemistry but vulnerable to co-produced ESBLs \\
\textbf{CTX-M-15} ESBL (Class A) & \emph{bla}\(_\text{CTX-M-15}\) &
Cephalosporins (residual) \\
\textbf{OXA-48} (Class D) & \emph{bla}\(_\text{OXA-48}\) & Carbapenems
(residual) \\
\textbf{mcr-1} & \emph{mcr-1} & Colistin \\
\textbf{Aac(6')-Ib-cr} & Plasmid & Aminoglycosides + ciprofloxacin \\
\textbf{rmtB} & Plasmid & Pan-aminoglycoside \\
\textbf{AcrAB-TolC overexpression} & \emph{ramA} mutation & Multiple \\
\textbf{GyrA Ser83Leu} & Chromosome & Fluoroquinolones \\
\textbf{OmpK35 loss} & Chromosome & Reduced \texorpdfstring{\ensuremath{\beta}}{beta}-lactam entry \\
\end{longtable}
}

Such isolates may have no reliable standard single-agent option, so
treatment depends on isolate-specific susceptibility testing, source
control, and expert consultation. Options can include cefiderocol,
ceftazidime-avibactam plus aztreonam for some MBL-producing isolates,
aminoglycoside or fosfomycin combinations when active, and
investigational phage therapy. The case illustrates that resistance is
\textbf{modular and additive}: each mechanism arrives independently,
often on plasmids, and accumulates over time. This is why antibiotic
stewardship, infection control, diagnostics, and new drug development
must operate together.

\begin{quote}
\textbf{Concept Check 6b:} A \emph{K. pneumoniae} isolate is resistant
to the tested \texorpdfstring{\ensuremath{\beta}}{beta}-lactams including carbapenems (KPC-2 carbapenemase),
tigecycline (RamA-dependent AcrAB overexpression), and ciprofloxacin
(GyrA + ParC mutations). The lab reports it as susceptible to colistin
and ceftazidime-avibactam. Explain (a) why ceftazidime-avibactam works
against KPC-2 specifically while ceftazidime alone would not, (b) why
the same combination would fail against an NDM-producer, and (c) how
ESBL-restricted versus carbapenemase-positive resistance is
distinguished by the meropenem-EDTA versus meropenem-boronic-acid
double-disk synergy test.
\end{quote}

\begin{center}\rule{0.5\linewidth}{0.5pt}\end{center}

\subsection{Epidemiology of Infectious
Disease}\label{epidemiology-of-infectious-disease}

\subsubsection{Key Epidemiological
Metrics}\label{key-epidemiological-metrics}

{\def\LTcaptype{none} % do not increment counter
\begin{longtable}[]{@{}
  >{\raggedright\arraybackslash}p{(\linewidth - 4\tabcolsep) * \real{0.2500}}
  >{\raggedright\arraybackslash}p{(\linewidth - 4\tabcolsep) * \real{0.3438}}
  >{\raggedright\arraybackslash}p{(\linewidth - 4\tabcolsep) * \real{0.4062}}@{}}
\toprule\noalign{}
\begin{minipage}[b]{\linewidth}\raggedright
Metric
\end{minipage} & \begin{minipage}[b]{\linewidth}\raggedright
Definition
\end{minipage} & \begin{minipage}[b]{\linewidth}\raggedright
Significance
\end{minipage} \\
\midrule\noalign{}
\endhead
\bottomrule\noalign{}
\endlastfoot
\(R_0\) (basic reproduction number) & Mean secondary cases per primary
case in fully susceptible population & Determines epidemic potential;
\(R_0 > 1\) = epidemic growth \\
\(R_e\) (effective reproduction number) &
\(R_0 \times (1 - \text{fraction immune})\); real-time transmission &
Guides intervention: goal is \(R_e < 1\) \\
CFR (case fatality rate) & Deaths / confirmed cases & Overestimates true
mortality (denominator misses mild cases) \\
IFR (infection fatality rate) & Deaths / most infections (including
asymptomatic) & More accurate; requires seroprevalence data \\
Incubation period & Time from infection to symptom onset & Determines
quarantine duration \\
Serial interval & Time between symptom onset in successive cases &
Determines epidemic speed \\
\end{longtable}
}

\subsubsection{Worked Example: Estimating R0 from an Epidemic
Curve}\label{worked-example-estimating-r0-from-an-epidemic-curve}

\textbf{Problem:} Early in an outbreak, confirmed cases double every
\(T_d = 3.0\) days. The mean serial interval (time between symptom onset
in successive cases) is \(T_s = 5.0\) days. Estimate the basic
reproduction number \(R_0\), then the herd-immunity threshold.

\textbf{Solution:}

\begin{enumerate}
\def\labelenumi{\arabic{enumi}.}
\tightlist
\item
  \textbf{Exponential growth rate \(r\) from the doubling time.} Early
  case counts grow as \(N(t) = N_0\,2^{t/T_d} = N_0\,e^{rt}\), so
  equating exponents gives \(r = \ln 2 / T_d\):
\end{enumerate}

\[ r = \frac{\ln 2}{T_d} = \frac{0.6931}{3.0} = 0.2310 \text{ day}^{-1} \label{eq:unit_VII_infectious_disease_item_3}\]

\begin{enumerate}
\def\labelenumi{\arabic{enumi}.}
\setcounter{enumi}{1}
\tightlist
\item
  \textbf{Convert growth rate to \(R_0\).} For a fixed
  (delta-distributed) serial interval, the renewal equation gives
  \(R_0 = e^{r T_s}\):
\end{enumerate}

\[ R_0 = e^{r T_s} = e^{0.2310 \times 5.0} = e^{1.155} = 3.17 \label{eq:unit_VII_infectious_disease_item_4}\]

(The simpler linear approximation
\(R_0 \approx 1 + r T_s = 1 + 0.2310 \times 5.0 = 2.16\) is a lower
bound that ignores serial-interval dispersion.)

\begin{enumerate}
\def\labelenumi{\arabic{enumi}.}
\setcounter{enumi}{2}
\tightlist
\item
  \textbf{Herd-immunity threshold.} Substituting \(R_0 = 3.17\) into the
  herd-immunity relation \(p_c = 1 - 1/R_0\) from
  \cref{eq:unit_VII_herd_immunity}:
\end{enumerate}

\[ p_c = 1 - \frac{1}{3.17} = 0.685 \;\rightarrow\; 68.5\,\% \label{eq:unit_VII_infectious_disease_item_5}\]

\textbf{Interpretation:} A 3-day doubling time with a 5-day serial
interval implies \(R_0 \approx 3.2\), so roughly 69 \% of the population
must be immune to halt sustained transmission --- a value consistent
with an early pandemic-respiratory pathogen, and one that grows sharply
if the doubling time shortens.

\subsubsection{Major Infectious
Diseases}\label{major-infectious-diseases}

\textbf{Tuberculosis} (\emph{Mycobacterium tuberculosis}):

\begin{itemize}
\tightlist
\item
  Global burden: an estimated 10.7 million people fell ill with TB in
  2024 and 1.23 million died, including 150,000 people with HIV
  \citep{who2025tb}. TB remains the leading cause of death from a single
  infectious agent and a major contributor to AMR-associated mortality.
\item
  Transmission: Airborne droplet nuclei (1-5 \texorpdfstring{\ensuremath{\mu}}{mu}m); can remain suspended
  for hours
\item
  Pathogenesis: Inhaled bacilli are phagocytosed by alveolar macrophages
  but arrest phagosome maturation (LAM inhibits PI3K -\textgreater{}
  blocks Ca\(^{2+}\)/calmodulin -\textgreater{} prevents phagolysosomal
  fusion); survive within macrophages at pH 6.4 instead of lethal pH 4.5
\item
  \textbf{Granuloma}: Hallmark of TB pathology -- organized aggregate of
  infected macrophages, epithelioid cells, multinucleated giant cells,
  and T cells; hypoxic caseous necrotic center; both contains and
  protects the bacilli
\item
  Latent TB: 90-95\% of infected individuals contain the infection as
  LTBI (latent TB infection); 5-10\% lifetime risk of reactivation,
  increased by HIV, immunosuppression, malnutrition
\item
  Treatment: standard drug-susceptible pulmonary TB is treated with a
  multi-drug rifamycin-based regimen, classically 6 months of isoniazid
  + rifampicin + pyrazinamide + ethambutol followed by continuation
  therapy; adherence support matters because metabolically slow bacilli
  and granuloma drug penetration make undertreatment risky.
\item
  \textbf{MDR/RR-TB} (multidrug-resistant or rifampicin-resistant TB),
  \textbf{pre-XDR-TB} (MDR/RR-TB with fluoroquinolone resistance), and
  \textbf{XDR-TB} (MDR/RR-TB plus resistance to a fluoroquinolone and at
  least one additional Group A drug such as bedaquiline or linezolid)
  remain major AMR threats. WHO's consolidated 2025 treatment guidance
  now prioritises shorter most-oral regimens for eligible patients,
  including 6-month BPaLM/BPaL-based options, while reserving longer
  individualized regimens for resistance, toxicity, pregnancy/age
  limits, extrapulmonary disease, or poor early response
  \citep{who2025tb}.
\end{itemize}

\begin{figure}[htbp]
\centering
\includegraphics[width=0.82\linewidth,height=4.2in,keepaspectratio]{../figures/mermaid_inline/inline_mermaid_0133_d23b238a8412.png}
\caption{TB regimen decision schematic. WHO's shorter BPaLM/BPaL options are eligibility-dependent tools for MDR/RR-TB, not a comprehensive replacement for susceptibility testing, toxicity monitoring, or individualized care \textbackslash{}citep\{who2025tb\}.}
\end{figure}

\begin{itemize}
\tightlist
\item
  Global burden: an estimated 282 million cases and 610,000 deaths in
  2024 across 80 countries \citep{who2025malaria}. The WHO African
  Region still accounts for the overwhelming majority of cases and
  deaths, with young children carrying the greatest mortality burden.
\item
  Species: \emph{P. falciparum} (most lethal; 99\% of deaths), \emph{P.
  vivax} (most widespread; hypnozoites cause relapse), \emph{P.
  malariae}, \emph{P. ovale}, \emph{P. knowlesi} (zoonotic, from
  macaques)
\item
  Transmission: Female \emph{Anopheles} mosquito vector; sporozoites
  injected during blood meal
\item
  Life cycle: Sporozoites -\textgreater{} liver (hepatocyte invasion,
  asymptomatic) -\textgreater{} merozoites -\textgreater{} erythrocyte
  invasion (clinical disease: cyclic fever every 48-72 hours
  synchronized with schizogony) -\textgreater{} some become gametocytes
  (taken up by mosquito)
\item
  Immune evasion: \emph{P. falciparum} erythrocyte membrane protein 1
  (PfEMP1) expressed on infected RBC surface mediates cytoadherence
  (sticking to endothelium, avoiding splenic clearance) and rosetting;
  about 60 \emph{var} genes enable antigenic variation
\item
  Vaccine: RTS,S/AS01E (Mosquirix): first approved malaria vaccine
  (2021); targets circumsporozoite protein; protection is partial and
  wanes without booster dosing. R21/Matrix-M is now the second
  WHO-recommended malaria vaccine; both are best understood as additions
  to bed nets, chemoprevention, diagnosis, and vector control rather
  than replacements \citep{who2025malaria}.
\item
  Treatment: Artemisinin-based combination therapy (ACT); artemisinin
  resistance emerging in Southeast Asia (K13 propeller mutations)
\item
  Vector control frontier: WHO issued a 2025 conditional recommendation
  for indoor \textbf{spatial emanators} (also called spatial repellents)
  as an additional tool in areas with ongoing transmission, used
  alongside insecticide-treated nets rather than replacing them. These
  devices release active ingredients such as transfluthrin into indoor
  air to repel, disorient, or kill mosquitoes, with evidence strongest
  for added indoor protection and remaining gaps for standalone use,
  outdoor protection, humanitarian settings, and resistance management
  \citep{who2025spatialemanators}.
\end{itemize}

\begin{figure}[htbp]
\centering
\includegraphics[width=0.82\linewidth,height=4.2in,keepaspectratio]{../figures/mermaid_inline/inline_mermaid_0134_2d5aca90a092.png}
\caption{Malaria vector-control comparison. Spatial emanators add an indoor airspace tool to nets and spraying, but local programmes still need entomological surveillance, insecticide-resistance monitoring, and equity-aware deployment.}
\end{figure}

\textbf{HIV/AIDS}:

\begin{itemize}
\tightlist
\item
  Global burden: an estimated 40.8 million people were living with HIV
  at the end of 2024, with about 1.3 million new infections in 2024 and
  AIDS-related deaths far below their early-2000s peak but still above
  global targets \citep{unaids2025factsheet}.
\item
  Pathogenesis: Progressive depletion of CD4\(^+\) T cells; AIDS defined
  as CD4 count \textless200 cells/\texorpdfstring{\ensuremath{\mu}}{mu}L or AIDS-defining illness
\item
  \textbf{ART} (antiretroviral therapy): Combination of 3+ drugs from
  different classes; suppresses viral load to undetectable (\textless50
  copies/mL); near-normal life expectancy if initiated early; does not
  cure (latent proviral reservoir in resting CD4\(^+\) memory T cells
  persists for decades)
\item
  \textbf{PrEP} (pre-exposure prophylaxis): Oral TDF/FTC or TAF/FTC
  remains highly effective when taken as prescribed; long-acting
  injectable cabotegravir and twice-yearly lenacapavir add
  adherence-sparing options. FDA approved lenacapavir PrEP in 2025, and
  CDC guidance reports trial efficacy of 100\% in a cisgender-female
  trial and 96\% in a primarily male trial over 52 weeks
  \citep{cdc2025lenacapavirprep}.
\item
  \textbf{U=U} (Undetectable = Untransmittable): Individuals with
  sustained undetectable viral load on ART cannot sexually transmit HIV
  (PARTNER study, HPTN 052)
\end{itemize}

\textbf{Influenza}:

\begin{itemize}
\tightlist
\item
  RNA virus (Orthomyxoviridae); 8 segmented (-)ssRNA genome segments
\item
  Surface glycoproteins: hemagglutinin (HA, 18 subtypes) and
  neuraminidase (NA, 11 subtypes); nomenclature: H1N1, H3N2, etc.
\item
  See ``Antigenic Variation'' section above for full discussion of drift
  and shift.
\end{itemize}

\subsubsection{One Health and Emerging
Pathogens}\label{one-health-and-emerging-pathogens}

The \textbf{One Health} framework recognizes the interconnection of
human, animal, and environmental health:

\begin{itemize}
\tightlist
\item
  Over 70\% of emerging infectious diseases are \textbf{zoonotic}
  (originating from animal reservoirs).
\item
  Drivers of emergence: deforestation and habitat destruction
  (increasing human-wildlife contact), wildlife trade (live animal
  markets), climate change (expanding vector ranges), intensive animal
  agriculture (influenza reassortment, AMR selection).
\item
  Examples: SARS-CoV-2 (probable bat origin via intermediate host),
  Ebola (bat reservoir), avian influenza H5N1 (poultry), Nipah
  (bat-to-pig-to-human), Lyme disease (deer-tick-mouse cycle expanding
  with climate change), mpox (rodent reservoirs in West/Central Africa).
\item
  AMR as a One Health issue: antibiotic exposure in people, food
  animals, aquaculture, and shared environments selects for resistance
  genes that can move through water, soil, food chains, plasmids, and
  mobile genetic elements. Stewardship therefore has to cover
  prescribing, veterinary practice, sanitation, surveillance, and access
  to effective treatment \citep{cdc2025antibioticuse,who2024bppl}.
\item
  Fungal disease also belongs in the AMR frame: the WHO fungal priority
  pathogens list highlights \emph{Candida auris}, \emph{Aspergillus
  fumigatus}, \emph{Cryptococcus neoformans}, and other fungi where
  limited diagnostics, few drug classes, immunocompromised hosts, and
  agricultural azole exposure make resistance a One Health problem
  rather than a hospital-confined issue \citep{who2022fungalpriority}.
  CDC guidance treats \emph{Candida auris} as a healthcare-transmissible
  yeast that is often multidrug-resistant, can colonise patients without
  symptoms, persists on surfaces, requires sequencing or mass
  spectrometry for reliable identification, and should be treated
  primarily when causing clinical infection; echinocandins remain
  typical initial adult therapy, but echinocandin-resistant and
  pan-resistant reports are increasing
  \citep{cdc2026candidaauris,cdc2024candidaauristreatment}.
  Mechanistically, the warning is that AMR phenotypes are assembled from
  enzymes, target changes, efflux, permeability shifts, biofilms,
  tolerance states, and mobile genetic elements; the same resistance
  label can hide different treatment constraints.
\end{itemize}

\textbf{Spillover risk factors.} Modern surveillance and modelling have
identified several quantifiable predictors of zoonotic spillover risk:

\begin{enumerate}
\def\labelenumi{\arabic{enumi}.}
\tightlist
\item
  \textbf{Phylogenetic distance} --- viruses from primates spill more
  easily than viruses from rodents than from invertebrates.
\item
  \textbf{Receptor compatibility} --- ACE2 in bats, civets, and humans
  is similar enough to permit SARS-CoV-2 cross-species transmission;
  H5N1 sialic-acid receptor preference shifts (\texorpdfstring{\ensuremath{\alpha}}{alpha}-2,3 \texorpdfstring{\ensuremath{\to}}{->} \texorpdfstring{\ensuremath{\alpha}}{alpha}-2,6) signal
  mammalian adaptation.
\item
  \textbf{Population density at the human-animal interface} --- wet
  markets, intensive agriculture, deforestation edges.
\item
  \textbf{Reservoir host species richness} --- Bat species richness
  correlates with the number of zoonotic viruses; rodents are similarly
  important but underrecognised.
\item
  \textbf{Anthropogenic disruption} --- deforestation, climate change,
  urbanisation increase contact rates.
\end{enumerate}

\textbf{The PREDICT and Global Virome Project} approaches systematically
sample wildlife to catalog viral diversity and identify high-risk
lineages \textbf{before} spillover. By 2026, sequencing \textgreater{}
200,000 viral genomes from \textgreater{} 30,000 animals across
\textgreater{} 30 countries has identified about  10,000
candidate zoonotic viruses, of which about  1700 share key
receptor-binding features with known human pathogens --- a sobering
reservoir.

\begin{quote}
\textbf{Clinical Connection: Antibiotic Stewardship and Real-Time
Resistance Surveillance} The WHO has declared antimicrobial resistance
one of the top 10 global public health threats. An estimated 1.27
million deaths were directly attributable to bacterial AMR in 2019
\citep{murray2022amr}. The O'Neill review's 10-million-deaths-per-year
scenario is a policy warning about uncontrolled resistance, not an
inevitable destiny \citep{oneill2016amr}. Genomic surveillance networks
such as the CDC AR Lab Network and WHO GLASS now track carbapenemases,
\emph{mcr} genes, and emerging plasmids across hospitals and regions.
The lesson from SARS-CoV-2 --- pathogen genomes can be sequenced and
interpreted at population scale --- is being adapted to bacterial AMR,
where slower growth and horizontal gene transfer make the analysis more
complicated.
\end{quote}

\begin{quote}
\textbf{Concept Check 7:} Trace a potential pandemic pathway from
deforestation in Southeast Asia. Identify (a) the ecological factors
that elevate spillover risk, (b) the molecular features (receptor
compatibility, reassortment potential) that determine human
transmissibility, (c) the early-warning indicators that PREDICT-style
surveillance would detect, and (d) the One Health interventions that
could reduce future risk.
\end{quote}

\begin{center}\rule{0.5\linewidth}{0.5pt}\end{center}

\subsection{Computational Bridge}\label{computational-bridge-23}

Early outbreak growth is often sketched with discrete exponential
models:

\begin{Shaded}
\begin{Highlighting}[]
\ImportTok{from}\NormalTok{ biology.ecology }\ImportTok{import}\NormalTok{ exponential\_growth}

\NormalTok{series }\OperatorTok{=}\NormalTok{ exponential\_growth(N0}\OperatorTok{=}\FloatTok{20.0}\NormalTok{, r}\OperatorTok{=}\FloatTok{0.3}\NormalTok{, t\_end}\OperatorTok{=}\FloatTok{12.0}\NormalTok{, steps}\OperatorTok{=}\DecValTok{12}\NormalTok{)}
\BuiltInTok{print}\NormalTok{(}\BuiltInTok{round}\NormalTok{(series.populations[}\OperatorTok{{-}}\DecValTok{1}\NormalTok{], }\DecValTok{2}\NormalTok{))}
\end{Highlighting}
\end{Shaded}

\begin{quote}
\textbf{Clinical / systems note:} \(R_e\) reductions from vaccines
behave like lowering effective \(r\) in these caricatures; heterogeneity
in contact networks breaks the homogeneous assumption.
\end{quote}

\begin{center}\rule{0.5\linewidth}{0.5pt}\end{center}

\subsubsection{mRNA Vaccine Platforms: From Lab Curiosity to Global
Deployment}\label{mrna-vaccine-platforms-from-lab-curiosity-to-global-deployment}

Messenger RNA vaccines --- once a neglected technology because of mRNA's
instability and strong innate-immune activation --- became the
\hyperref[gl:dominant]{\textbf{dominant}} pandemic response platform
between 2020 and 2023. \textbf{BNT162b2 (Pfizer--BioNTech) and mRNA-1273
(Moderna)} demonstrated \texorpdfstring{\textgreater{}=}{>=} 94 \% efficacy against symptomatic COVID-19 in
phase III trials (\textgreater{} 40 000 participants each) within
\textbf{11 months from pathogen identification} --- an unprecedented
pace driven entirely by the modularity of the platform: once a
lipid-nanoparticle delivery system and a nucleoside-modified mRNA
scaffold are validated, swapping the antigen-coding sequence is a matter
of days.

Three molecular innovations make mRNA vaccines work. \textbf{(1)
Nucleoside modification} --- substituting uridines with
\textbf{pseudouridine (\texorpdfstring{\ensuremath{\Psi}}{Psi}) or 1-methyl-pseudouridine (m1\texorpdfstring{\ensuremath{\Psi}}{Psi})} (Karikó \&
Weissman, Nobel Prize 2023) reduces TLR7/8 activation and blocks
2'-5'-oligoadenylate synthetase / RNase L degradation, so
\hyperref[gl:translation]{\textbf{translation}} is 10--100\texorpdfstring{\ensuremath{\times}}{x} more
efficient and innate-immune activation is manageable. \textbf{(2) Lipid
nanoparticle (LNP) delivery} --- a four-component formulation (ionisable
lipid ALC-0315 or SM-102, phospholipid DSPC, cholesterol, PEG-lipid)
self-assembles into about 80 nm particles that endocytose into
dendritic cells; at endosomal pH about 5 the ionisable lipid
becomes cationic, fusing with the endosome and releasing mRNA into the
cytoplasm. \textbf{(3) Antigen design} --- the SARS-CoV-2 spike is
locked in the prefusion conformation by two proline substitutions
(\textbf{2P mutation}) that preserve neutralising epitopes.
Manufacturing uses \emph{in vitro} transcription from linear DNA
templates (T7 polymerase), capping with CleanCap analogue, and
microfluidic LNP assembly --- a single about 20 000 L
bioreactor run can produce \textgreater{} 1 billion doses. Remaining
challenges: cold-chain dependence (though ARCT-154 self-amplifying mRNA
vaccines stable at 2--8 °C are in trials), rare myocarditis risk in
young males (about 1--10 per 100 000 doses, typically mild and
self-limiting), and waning antibody titres (6-month boosting needed).
The platform is now extending to influenza (mRNA-1010), RSV (mRNA-1345,
approved 2024), and personalised cancer neoantigen vaccines (mRNA-4157,
phase III for melanoma).

\subsubsection{cGAS--STING: Innate Cytosolic DNA Sensing and Anti-Tumour
Immunity}\label{cgassting-innate-cytosolic-dna-sensing-and-anti-tumour-immunity}

\textbf{Cyclic GMP--AMP synthase (cGAS)} is a cytosolic DNA sensor
identified in 2013 (Chen lab) that binds double-stranded DNA
non-sequence-specifically and, upon DNA engagement, synthesises the
second messenger \textbf{2'3'-cyclic GMP--AMP (2'3'-cGAMP)} from ATP and
GTP. 2'3'-cGAMP binds the ER-resident adaptor \textbf{STING (Stimulator
of Interferon Genes)}, driving TBK1 phosphorylation of IRF3 and thus
\textbf{Type I interferon (IFN-\texorpdfstring{\ensuremath{\alpha}}{alpha}/\texorpdfstring{\ensuremath{\beta}}{beta})} transcription. Evolutionarily, the
pathway detects pathogen-derived or mis-localised self-DNA --- cytosolic
DNA is otherwise absent in healthy cells.

Three rapid-fire translational applications: \textbf{(1) Cancer
immunotherapy} --- dying tumour cells release DNA that, once dendritic
cells endocytose it, activates cGAS--STING, priming CD8\texorpdfstring{\textsuperscript{+}}{+} T-cell
responses. STING agonists (ADU-S100, MK-1454, diABZI) are in
\textgreater{} 30 oncology clinical trials; \textbf{radiation therapy}
is an accidental STING activator (sub-lethal DNA damage \texorpdfstring{\ensuremath{\to}}{->} micronuclei \texorpdfstring{\ensuremath{\to}}{->}
cytosolic DNA). \textbf{(2) Autoimmunity} --- gain-of-function STING
mutations cause \textbf{SAVI (STING-associated vasculopathy)}, a
paediatric interferonopathy with pulmonary fibrosis and skin
ulcerations; chronic cGAS--STING activation also contributes to
\textbf{age-related inflammation} as mitochondrial and cytosolic DNA
accumulate. \textbf{(3) Anti-viral defence} --- SARS-CoV-2 ORF9b, HSV-1
ICP27, and HIV-1 capsid most suppress cGAS--STING, demonstrating its
evolutionary importance as a front-line sensor. The pathway is now a
canonical example of how fundamental discovery (2013) can translate to
phase III oncology trials in under a decade.

\begin{center}\rule{0.5\linewidth}{0.5pt}\end{center}

\subsection{Current Evidence and Frontier
Biology}\label{current-evidence-and-frontier-biology-26}

For \textbf{Infectious Disease and Immunity}, frontier biology belongs
inside the evidence logic of the chapter. Microbiology and infectious
disease now require One Health reasoning across people, animals,
environments, genomics, and antimicrobial stewardship. The core reading
question is this: infectious-disease reasoning should connect pathogen
biology, transmission, immunity, diagnostics, interventions, and equity.

\begin{itemize}
\tightlist
\item
  \textbf{What to verify:} identify the observation, model, assay, or
  dataset that would make the claim stronger or weaker.
\item
  \textbf{What to qualify:} state the scale, organism, cell type,
  environmental condition, or population where the claim is expected to
  hold.
\item
  \textbf{What to compare:} test at least one alternative explanation,
  baseline, or null model before treating the pattern as causal.
\item
  \textbf{What to cite:} distinguish primary evidence, review synthesis,
  public dataset, and institutional guidance; for recent or numeric
  claims, prefer the source closest to the measurement and state what
  has changed since it was published.
\end{itemize}

For AMR and pathogen claims, name the organism-resistance pair, the
selection pressure, the transmission route, and the surveillance
evidence that would change triage
\citep{who2024bppl,cdc2025antibioticuse,murray2022amr}.

\textbf{Source practice:} For pathogen, AMR, and intervention claims,
tie statements to organism-resistance pairs, surveillance evidence,
official guidance, and trial/regulatory status
\citep{who2024bppl,who2025tb,who2025malaria,cdc2025lenacapavirprep,cdc2026candidaauris}.

\subsubsection{Current Evidence Map: Intervention Choice Across
Pathogens}\label{current-evidence-map-intervention-choice-across-pathogens}

\begin{figure}[htbp]
\centering
\includegraphics[width=0.82\linewidth,height=4.2in,keepaspectratio]{../figures/mermaid_inline/inline_mermaid_0135_1124d195aaca.png}
\caption{TB regimens, malaria spatial emanators, lenacapavir PrEP, Candida auris control, and Long COVID mechanisms require matching intervention to evidence and setting \textbackslash{}citep\{who2025tb,who2025spatialemanators,cdc2025lenacapavirprep,cdc2026candidaauris,longcovid2026mechanisms\}.}
\end{figure}

\subsection{Summary}\label{summary-26}

\begin{itemize}
\tightlist
\item
  \textbf{Koch's postulates} remain foundational for establishing
  disease causation but have important limitations (asymptomatic
  carriers, unculturable organisms, polymicrobial disease). Molecular
  Koch's postulates extend the framework to virulence genes.
\item
  \textbf{Virulence factors} include adhesins (pili, FimH), invasion
  machinery (T3SS), toxins (A-B exotoxins: cholera, diphtheria, tetanus,
  botulinum; endotoxin/LPS), capsules, and immune evasion strategies
  (protein A, antigenic variation, phagosome arrest).
\item
  \textbf{\hyperref[gl:innate-immunity]{Innate immunity}} provides
  immediate, non-specific defense: barriers (skin, mucus, acid,
  lysozyme, defensins), PRRs (TLRs, NLRs, RIG-I, cGAS-STING) \texorpdfstring{\ensuremath{\to}}{->} NF-\texorpdfstring{\ensuremath{\kappa}}{kappa}B +
  IRF3 \texorpdfstring{\ensuremath{\to}}{->} cytokines + type I IFN; \textbf{complement} with three pathways
  (classical via C1q, lectin via MBL/MASPs, alternative via spontaneous
  C3 tickover) converging at C3 convertase, then C5 cleavage and MAC
  formation; cellular effectors (neutrophils with MPO/HOCl oxidative
  burst and NETosis, macrophages M1/M2, NK cells with missing-self via
  inhibitory KIRs and induced-self via NKG2D).
\item
  \textbf{Adaptive immunity}: V(D)J recombination generates diverse TCRs
  and BCRs; MHC-I presents endogenous peptides to CD8\(^+\) CTLs; MHC-II
  presents exogenous peptides to CD4\(^+\) T helpers (Th1, Th2, Th17,
  Treg, Tfh); B cells undergo germinal center reactions (SHM, affinity
  maturation, CSR) producing high-affinity class-switched antibodies and
  long-lived memory.
\item
  \textbf{Vaccines} exploit immunological memory: eight platforms (live
  attenuated, inactivated, subunit, VLP, toxoid, conjugate, mRNA, viral
  vector) span the spectrum from MMR (single-dose lifetime immunity) to
  mRNA (rapid platform with 48-hour design cycle). Herd-immunity
  threshold \(p_c = 1 - 1/R_0\) (\cref{eq:unit_VII_herd_immunity}) ---
  measles requires \textgreater{} 93 \%, polio about  80 \%,
  original COVID-19 about  60 \%, but Omicron-era
  \(R_0 \sim 10\) combined with imperfect vaccine efficacy can exceed
  achievable coverage.
\item
  \textbf{Antigenic variation} in influenza occurs via gradual drift (HA
  point mutations \texorpdfstring{\ensuremath{\to}}{->} seasonal epidemics) and abrupt shift (segment
  reassortment in animal reservoirs \texorpdfstring{\ensuremath{\to}}{->} pandemics: 1918, 1957, 1968,
  2009). HIV reverse transcriptase (3 \texorpdfstring{\ensuremath{\times}}{x} 10\texorpdfstring{\textsuperscript{-}}{-}\texorpdfstring{\textsuperscript{5}}{5} errors/bp/cycle) generates
  a quasi-species cloud that defeats monotherapy and complicates vaccine
  development.
\item
  \textbf{Antibiotic resistance} falls into five categories: enzymatic
  inactivation (\texorpdfstring{\ensuremath{\beta}}{beta}-lactamases --- Class A/B/C/D Ambler; AMEs; CAT; ESBLs
  and carbapenemases KPC/NDM/OXA-48); target modification (PBP2a in
  MRSA; D-Ala-D-Lac in VRE; 23S methylation in MLS; gyrase QRDR
  mutations); efflux (RND tripartite pumps AcrAB-TolC, MexAB-OprM,
  AdeABC); reduced permeability (porin loss); and bypass / target
  overproduction (Dfr, Sul). Pan-resistant ESKAPE isolates combine
  multiple mechanisms simultaneously.
\item
  \textbf{Major infectious diseases}: TB (granuloma, phagosome arrest,
  MDR/XDR crisis), malaria (erythrocyte invasion, PfEMP1 antigenic
  variation, artemisinin resistance), HIV (CD4 depletion, latent
  reservoir, ART, PrEP, U=U), influenza (antigenic drift/shift, pandemic
  potential).
\item
  \textbf{One Health and emerging pathogens}: \textgreater{} 70 \% of
  emerging infections are zoonotic; deforestation, climate change,
  wildlife trade, and agricultural antibiotic use drive emergence and
  AMR. PREDICT-style surveillance now catalogs \textgreater{} 10 000
  candidate zoonotic viruses for early-warning monitoring.
\item
  \textbf{Connections:} See \cref{sec:unit_VII_bacteria_archaea_viruses}
  for pathogen biology, \cref{sec:unit_VII_microbial_ecology} for
  microbiome colonisation resistance, and
  \cref{sec:unit_X_population_ecology} for demographic models.
\end{itemize}

\begin{center}\rule{0.5\linewidth}{0.5pt}\end{center}

\subsection{Key Terms}\label{key-terms-26}

{\def\LTcaptype{none} % do not increment counter
\begin{longtable}[]{@{}
  >{\raggedright\arraybackslash}p{(\linewidth - 2\tabcolsep) * \real{0.3529}}
  >{\raggedright\arraybackslash}p{(\linewidth - 2\tabcolsep) * \real{0.6471}}@{}}
\toprule\noalign{}
\begin{minipage}[b]{\linewidth}\raggedright
Term
\end{minipage} & \begin{minipage}[b]{\linewidth}\raggedright
Definition
\end{minipage} \\
\midrule\noalign{}
\endhead
\bottomrule\noalign{}
\endlastfoot
\textbf{Koch's postulates} & Four criteria for establishing that a
specific microorganism causes a specific disease \\
\textbf{Virulence factor} & Microbial product or strategy that
contributes to pathogenicity (adhesins, toxins, capsules, immune
evasion) \\
\textbf{Exotoxin} & Secreted bacterial protein toxin; often A-B
structure (cholera, diphtheria, tetanus, botulinum) \\
\textbf{Endotoxin (LPS)} & Gram-negative outer membrane component
released upon lysis; activates TLR4; causes fever, septic shock \\
\textbf{PAMP} & Pathogen-Associated Molecular Pattern -- conserved
microbial structure detected by host PRRs \\
\textbf{\hyperref[gl:toll-like-receptor]{Toll-like receptor (TLR)}} &
Membrane-bound PRR that recognizes PAMPs and initiates innate immune
signaling \\
\textbf{Inflammasome} & Intracellular multiprotein complex
(NLRP3-ASC-caspase-1) that processes IL-1\texorpdfstring{\ensuremath{\beta}}{beta} and IL-18 and triggers
pyroptosis \\
\textbf{Complement} & Plasma protein cascade with three pathways
(classical via C1q-antibody, lectin via MBL-carbohydrate, alternative
via spontaneous C3 hydrolysis) converging at the C3 convertase \\
\textbf{C3 convertase} & Central enzyme of complement (C4b2a or C3bBb)
that cleaves C3 into C3a (anaphylatoxin) and C3b (opsonin) \\
\textbf{MAC} & Membrane attack complex (C5b-C9\(_n\)); 10-nm pore that
osmotically lyses Gram-negative bacteria \\
\textbf{Factor H} & Soluble alternative-pathway regulator that binds C3b
on self surfaces; co-opted by \emph{Neisseria} to evade complement \\
\textbf{NETosis} & Programmed neutrophil death pathway in which
decondensed chromatin is extruded as antimicrobial extracellular traps;
PAD4-dependent \\
\textbf{Missing-self hypothesis} & Ljunggren-Karre model: NK cells kill
cells lacking inhibitory MHC-I signals \\
\textbf{NKG2D} & Activating NK-cell receptor recognising stress-induced
ligands (MICA/B, ULBPs); the ``induced-self'' sensor \\
\textbf{MHC class I} & Presents endogenous peptides to CD8\(^+\) CTLs;
expressed on most nucleated cells \\
\textbf{MHC class II} & Presents exogenous peptides to CD4\(^+\) T
helpers; expressed on professional APCs \\
\textbf{Clonal selection} & Antigen selects and expands primarily the
specific lymphocyte clone with matching receptor \\
\textbf{Affinity maturation} & Progressive improvement of antibody
affinity through somatic hypermutation and selection in germinal
centers \\
\textbf{Class switch recombination} & AID-mediated change of antibody
constant region (IgM -\textgreater{} IgG, IgA, or IgE) preserving
antigen specificity \\
\textbf{Herd immunity threshold} & \(p_c = 1 - 1/R_0\); fraction of
population that must be immune to prevent sustained transmission \\
\textbf{\(R_0\)} & Basic reproduction number; mean secondary infections
per primary case in fully susceptible population \\
\textbf{Antigenic drift} & Gradual mutation of surface antigens
(influenza HA/NA) producing seasonal variants \\
\textbf{Antigenic shift} & Reassortment of genome segments between
different influenza strains producing novel pandemic subtypes \\
\textbf{Quasi-species} & The mutant cloud generated by an error-prone
polymerase (HIV reverse transcriptase, RNA-virus RdRp) --- the substrate
of within-host evolution \\
\textbf{VLP (virus-like particle)} & Self-assembled viral capsid lacking
genome; basis of HPV (Gardasil) and HBV vaccines \\
\textbf{mRNA vaccine} & LNP-encapsulated, nucleoside-modified mRNA that
the host translates into antigen; e.g., BNT162b2, mRNA-1273 \\
\textbf{\texorpdfstring{\ensuremath{\beta}}{beta}-lactamase} & Enzyme that hydrolyses \texorpdfstring{\ensuremath{\beta}}{beta}-lactam antibiotics;
Ambler classes A/C/D (serine) and B (metallo); ESBLs, KPC, NDM, OXA \\
\textbf{PBP2a} & Alternative penicillin-binding protein in MRSA, encoded
by \emph{mecA}; low \texorpdfstring{\ensuremath{\beta}}{beta}-lactam affinity \\
\textbf{AcrAB-TolC, MexAB-OprM} & Tripartite RND-family efflux pumps
spanning the entire Gram-negative cell envelope; major multi-drug
resistance mechanism \\
\textbf{mcr-1} & Plasmid-borne phosphoethanolamine transferase; first
transferable colistin resistance gene (China, 2015) \\
\textbf{ART} & Antiretroviral therapy; combination drug regimen that
suppresses HIV viral load to undetectable levels \\
\textbf{One Health} & Framework recognizing the interconnection of
human, animal, and environmental health in infectious disease and AMR \\
\end{longtable}
}

\begin{center}\rule{0.5\linewidth}{0.5pt}\end{center}

\subsection{Review Questions}\label{review-questions-26}

\begin{enumerate}
\def\labelenumi{\arabic{enumi}.}
\item
  Apply Koch's postulates to \emph{Helicobacter pylori} and peptic ulcer
  disease. Barry Marshall famously fulfilled these postulates by
  self-experimentation in 1984. Identify which postulate was most
  difficult to satisfy and explain why the medical establishment was
  initially skeptical.
\item
  Compare the mechanisms of action of cholera toxin and diphtheria
  toxin. Both use A-B structure and ADP-ribosylation, but they target
  different host proteins. Explain how the same enzymatic mechanism
  (ADP-ribosylation) produces completely different clinical outcomes.
\item
  A patient with terminal complement component deficiency (C5-C9)
  presents with recurrent \emph{Neisseria meningitidis} infections.
  Explain why this specific pathogen is problematic in MAC deficiency
  while most other bacterial infections are handled normally. What does
  this tell you about the relative importance of opsonization versus MAC
  lysis for different pathogens?
\item
  Describe the molecular events of T cell activation, including most
  three signals. Explain why Signal 2 (co-stimulation) is critical for
  preventing autoimmunity, and predict what would happen if a
  pharmaceutical agent blocked B7-CD28 interaction globally.
\item
  A mother's IgG antibodies cross the placenta and protect the newborn
  for the first 3-6 months of life. Explain why this passive immunity
  wanes and why active vaccination is necessary starting at 2 months.
  Which antibody class is most important in breast milk, and how does it
  protect the infant's mucosal surfaces?
\item
  \textbf{Herd immunity calculation.} Calculate from
  \cref{eq:unit_VII_herd_immunity} the herd immunity threshold for a
  novel respiratory pathogen with \(R_0 = 8\). If a vaccine with 90\%
  efficacy is available, what percentage of the population must be
  vaccinated to achieve herd immunity? Show your work using the formula
  \(p_v = p_c / E\).
\item
  Explain why tuberculosis treatment requires 6-9 months of multi-drug
  therapy while most bacterial infections are treated with 7-14 days of
  a single antibiotic. Include the concepts of metabolically dormant
  persisters, granuloma pharmacokinetics, and the rationale for
  combination therapy.
\item
  \textbf{Drift vs shift.} Compare antigenic drift and antigenic shift
  in influenza. Explain why antigenic shift can produce pandemics while
  antigenic drift produces seasonal epidemics. Why is influenza A (but
  not influenza B) capable of antigenic shift, and what role do animal
  reservoirs play? How does the HIV quasi-species cloud differ
  mechanistically from influenza, and why does each strategy defeat
  conventional vaccine design?
\item
  The mRNA COVID-19 vaccines (BNT162b2, mRNA-1273) demonstrated
  about 95\% efficacy in clinical trials but this declined over
  time and with new variants. Explain: (a) the immunological mechanism
  of mRNA vaccines, (b) why booster doses are needed (waning immunity
  vs.~antigenic evolution), and (c) how the N1-methylpseudouridine
  modification improves mRNA vaccine performance.
\item
  Using the One Health framework, explain how deforestation in Southeast
  Asia could lead to a novel pandemic. Trace the pathway from habitat
  destruction to zoonotic spillover, identifying the ecological,
  virological, and epidemiological factors at each step.
\item
  Derive the herd immunity threshold \(p_c = 1 - 1/R_0\) from the
  next-generation intuition and relate to Question 6.
\item
  Explain \textbf{antigenic sin} and how it might bias booster responses
  after sequential variant exposure.
\item
  \textbf{Resistance mechanism integration.} A clinical microbiology lab
  reports a \emph{Klebsiella pneumoniae} isolate with MICs: meropenem \texorpdfstring{\textgreater{}=}{>=}
  16 (R), ceftazidime-avibactam 1 (S), cefepime 8 (R), ciprofloxacin \texorpdfstring{\textgreater{}=}{>=} 4
  (R), gentamicin \texorpdfstring{\textgreater{}=}{>=} 16 (R), tigecycline 0.5 (S), colistin 0.5 (S).
  Sequencing reveals: KPC-2 carbapenemase, AAC(6')-Ib-cr, GyrA Ser83Leu,
  ParC Ser80Ile, RamA overexpression. (a) Which mechanism is responsible
  for each resistance phenotype? (b) Why does ceftazidime-avibactam
  still work? (c) Predict whether this strain would respond to
  imipenem-relebactam and meropenem-vaborbactam, and justify.
\item
  \textbf{Vaccine platform comparison.} A new pathogen emerges with
  seasonal transmission and mild illness in adults, severe in infants,
  no available therapy. Compare the trade-offs for a (a)
  live-attenuated, (b) inactivated, (c) subunit-VLP, (d) mRNA, and (e)
  viral-vector vaccine platform decision, given a 9-month timeline to
  first dose.
\item
  \textbf{Innate immunity integration.} A patient with AIDS (CD4
  \textless{} 50) develops disseminated \emph{Mycobacterium avium
  complex} infection. Explain why both adaptive and innate arms of
  immunity have failed despite intact NK cells, neutrophils, complement,
  and antibodies. What is the role of granuloma formation, and why is
  IFN-\texorpdfstring{\ensuremath{\gamma}}{gamma} critical?
\end{enumerate}

\subsection{Further Reading and Source
Notes}\label{further-reading-and-source-notes-21}

\begin{itemize}
\tightlist
\item
  Janeway, Travers, Walport \& Shlomchik (latest ed.). \emph{Janeway's
  Immunobiology}. Garland Science.
\item
  Medzhitov \& Janeway (2000). Innate immunity. \emph{New England
  Journal of Medicine}, 343.
\item
  Plotkin (2010). Correlates of protection induced by vaccination.
  \emph{Clinical and Vaccine Immunology}, 17.
\item
  Anderson \& May (1991). \emph{Infectious Diseases of Humans: Dynamics
  and Control}. Oxford University Press.
\item
  Kermack \& McKendrick (1927). A contribution to the mathematical
  theory of epidemics. \emph{Proceedings of the Royal Society A}, 115.
\item
  Davies, Spagnolo \& Walsh (2010). Origins and evolution of antibiotic
  resistance. \emph{Microbiology and Molecular Biology Reviews}, 74.
\item
  WHO (latest). \emph{Global Antimicrobial Resistance and Use
  Surveillance System (GLASS) Report}. World Health Organization.
\end{itemize}

\begin{center}\rule{0.5\linewidth}{0.5pt}\end{center}

\subsubsection{Companion Source
Module}\label{companion-source-module-26}

\textbf{Infectious Disease and Immunity} should leave a reproducible
trail from a biological claim to the code, figure, diagram, or
paper-based activity that can test it. Use the surfaces below to inspect
the chapter's assumptions, rerun the relevant model, or compare the
manuscript explanation with companion labs and figures.

{\def\LTcaptype{none} % do not increment counter
\begin{longtable}[]{@{}
  >{\raggedright\arraybackslash}p{(\linewidth - 2\tabcolsep) * \real{0.5000}}
  >{\raggedright\arraybackslash}p{(\linewidth - 2\tabcolsep) * \real{0.5000}}@{}}
\toprule\noalign{}
\begin{minipage}[b]{\linewidth}\raggedright
Surface
\end{minipage} & \begin{minipage}[b]{\linewidth}\raggedright
Use it for
\end{minipage} \\
\midrule\noalign{}
\endhead
\bottomrule\noalign{}
\endlastfoot
\breaktt{src/biology/microbiology/microbiology.py}
(\breaktt{basic\_reproduction\_number}, \texttt{sir\_model},
\texttt{mic\_fold\_dilution}) & Reproduce transmission and
antimicrobial-resistance calculations. \\
\breaktt{src/biology/ecology/ecology.py} (\breaktt{exponential\_growth}) &
Compare early outbreak growth with ecological growth models. \\
\breaktt{src/mermaid/biology\_diagrams.py}
(\breaktt{immune\_response\_diagram},
\breaktt{viral\_replication\_cycle\_diagram}) & Connect pathogen life
cycle to host response. \\
\end{longtable}
}

\textbf{Reproducibility check:} identify pathogen, host population,
transmission route, diagnostic window, intervention, and surveillance
source before comparing disease claims. \textbf{Cross-reference:}
connect with \cref{sec:unit_VII_bacteria_archaea_viruses},
\cref{sec:unit_IX_endocrine_and_immune}, and
\cref{sec:unit_X_community_ecology}.

\newpage

\section*{Unit VIII --- Botany --- Plant Biology:
Introduction}\label{unit-viii-botany-plant-biology-introduction}
\addcontentsline{toc}{section}{Unit VIII --- Botany --- Plant Biology:
Introduction}

\subsection*{Why This Unit Matters}\label{why-this-unit-matters-8}
\addcontentsline{toc}{subsection}{Why This Unit Matters}

Plants are the foundation of almost every terrestrial food web, the
source of roughly half the oxygen in the atmosphere, and the origin of
many medicines in human pharmacopoeia --- from aspirin (salicylate,
first identified in willow bark) to taxol (from Pacific yew, \emph{Taxus
brevifolia}), morphine, quinine, and artemisinin. More fundamentally,
plants are solar-energy collectors that translate photons into chemical
bonds, feeding the entire biosphere. Each year, land plants and algae
fix approximately 120 billion tonnes of carbon from atmospheric CO\texorpdfstring{\textsubscript{2}}{2} via
photosynthesis --- a flux so large that it modulates global atmospheric
chemistry.

Plant biology is also a discipline of extremes. The tallest trees (coast
redwoods, \emph{Sequoia sempervirens}, 115 m) move water 115 metres
against gravity using primarily cohesion-tension and solar energy,
achieving flow rates in excess of 1 m/hour. Some desert plants survive
total desiccation and revive on rehydration. Others flower in response
to photoperiod cues measured in the phytochrome system with a resolution
of minutes. The \emph{Arabidopsis thaliana} genome was the first plant
genome sequenced (2000), and this small weed has since become the
"\emph{E. coli} of plant biology" --- a model organism for the molecular
dissection of everything from root development to pathogen resistance.

This unit integrates anatomy (vascular organisation, meristems,
root-shoot polarity), physiology (transpiration, water potential, phloem
loading), reproductive biology (alternation of generations, angiosperm
diversity, seed development), and signalling (phytohormones: auxin,
gibberellin, cytokinin, abscisic acid, ethylene). The quantitative spine
is water potential (\(\Psi = \Psi_s + \Psi_p\)) and transpiration
modelled by the van den Honert equation.

\begin{center}\rule{0.5\linewidth}{0.5pt}\end{center}

\subsection*{Landmark Discoveries}\label{landmark-discoveries-8}
\addcontentsline{toc}{subsection}{Landmark Discoveries}

{\def\LTcaptype{none} % do not increment counter
\begin{longtable}[]{@{}
  >{\raggedright\arraybackslash}p{(\linewidth - 8\tabcolsep) * \real{0.2407}}
  >{\raggedright\arraybackslash}p{(\linewidth - 8\tabcolsep) * \real{0.0741}}
  >{\raggedright\arraybackslash}p{(\linewidth - 8\tabcolsep) * \real{0.2963}}
  >{\raggedright\arraybackslash}p{(\linewidth - 8\tabcolsep) * \real{0.1667}}
  >{\raggedright\arraybackslash}p{(\linewidth - 8\tabcolsep) * \real{0.2222}}@{}}
\toprule\noalign{}
\begin{minipage}[b]{\linewidth}\raggedright
Discoverer(s)
\end{minipage} & \begin{minipage}[b]{\linewidth}\raggedright
Year
\end{minipage} & \begin{minipage}[b]{\linewidth}\raggedright
Journal / Source
\end{minipage} & \begin{minipage}[b]{\linewidth}\raggedright
Discovery
\end{minipage} & \begin{minipage}[b]{\linewidth}\raggedright
Significance
\end{minipage} \\
\midrule\noalign{}
\endhead
\bottomrule\noalign{}
\endlastfoot
Jan Ingenhousz & 1779 & \emph{Experiments upon Vegetables} & Plants
produce oxygen in light, CO\texorpdfstring{\textsubscript{2}}{2} in dark & Established the light-dependency
of photosynthesis \\
Julius von Sachs & 1862 & \emph{Botanische Zeitung} & Starch as first
visible product of photosynthesis & Linked CO\texorpdfstring{\textsubscript{2}}{2} fixation to carbohydrate
storage in chloroplasts \\
Frits Went & 1926 & \emph{Recueil des Travaux Botaniques Néerlandais} &
Isolation of auxin; phototropism mechanism & First plant hormone
identified; explained Darwinian canary grass bending \\
Henry Dixon \& John Joly & 1894 & \emph{Proc. R. Irish Acad.} &
Cohesion-tension theory of water ascent & Explained how trees move water
100+ m with no pump \\
Chailakhyan & 1936 & \emph{C. R. Acad. Sci. URSS} & Florigen hypothesis
(day-length controls flowering) & Proposed a systemic flowering
stimulus; later identified as FT protein \\
\emph{Arabidopsis} Genome Initiative & 2000 & \emph{Nature} & First
plant genome sequence (125 Mb; about 25,500 genes) & Model
organism for molecular plant biology; comparative genomics \\
Gómez-Roldán et al. & 2008 & \emph{Nature} & Strigolactones: new class
of plant hormones & Sixth plant hormone class; regulates branching and
mycorrhizal symbiosis \\
\end{longtable}
}

\begin{center}\rule{0.5\linewidth}{0.5pt}\end{center}

\subsection*{Key Concepts and
Connections}\label{key-concepts-and-connections-8}
\addcontentsline{toc}{subsection}{Key Concepts and Connections}

\begin{figure}[htbp]
\centering
\includegraphics[width=0.82\linewidth,height=4.2in,keepaspectratio]{../figures/mermaid_inline/inline_mermaid_0136_f6243b43cd27.png}
\caption{Flowchart for Key Concepts and Connections: plant Plant Structure (roots, stems, leaves; meristems; vascular: xylem/phloem), water Water Relations (Psi = Psis + Psip; cohesion-tension mechanism), wind Transpiration (stomata; guard cells; transpiration stream), and sunlight Photosynthesis (Unit III to C3, C4, CAM pathways) form the diagram's primary path or branches.}
\end{figure}

\textbf{Unit VIII concept map --- Plant Biology.}

\begin{center}\rule{0.5\linewidth}{0.5pt}\end{center}

\subsection*{Current Evidence Thread}\label{current-evidence-thread-8}
\addcontentsline{toc}{subsection}{Current Evidence Thread}

In this unit, treat plant biology as a body of evidence: each claim
about structure, water transport, environmental response, or
reproduction should rest on physiological measurement, in vivo imaging,
or a field or engineering trial, not on description alone. Plant biology
links molecular regulation to climate stress, water limitation, crop
resilience, phenology, and ecosystem feedbacks. As you move through the
chapters, keep a two-column note: \textbf{claim} on the left,
\textbf{evidence that would change my confidence} on the right. By the
end of the unit, each major idea should be tied to a measurement, model,
citation, or paper-based lab decision.

\subsection*{Chapter Roadmap}\label{chapter-roadmap-9}
\addcontentsline{toc}{subsection}{Chapter Roadmap}

{\def\LTcaptype{none} % do not increment counter
\begin{longtable}[]{@{}
  >{\raggedright\arraybackslash}p{(\linewidth - 6\tabcolsep) * \real{0.1556}}
  >{\raggedright\arraybackslash}p{(\linewidth - 6\tabcolsep) * \real{0.1111}}
  >{\raggedright\arraybackslash}p{(\linewidth - 6\tabcolsep) * \real{0.2889}}
  >{\raggedright\arraybackslash}p{(\linewidth - 6\tabcolsep) * \real{0.4444}}@{}}
\toprule\noalign{}
\begin{minipage}[b]{\linewidth}\raggedright
Chapter
\end{minipage} & \begin{minipage}[b]{\linewidth}\raggedright
Title
\end{minipage} & \begin{minipage}[b]{\linewidth}\raggedright
Core Question
\end{minipage} & \begin{minipage}[b]{\linewidth}\raggedright
Key Equation / Model
\end{minipage} \\
\midrule\noalign{}
\endhead
\bottomrule\noalign{}
\endlastfoot
\textbf{25} & Plant Structure and Water Transport & How is plant body
plan organised, and how does water move from roots to leaves? &
\(\Psi = \Psi_s + \Psi_p\); Poiseuille's law (xylem flow) \\
\textbf{26} & Plant Reproduction & How do plants reproduce sexually and
asexually, and how did flowering plant diversity arise? & Alternation of
generations; polyploidy speciation rates \\
\textbf{27} & Plant Responses to the Environment & How do phytohormones
integrate environmental signals into growth and defence responses? &
Auxin polar transport; ethylene activation kinetics \\
\end{longtable}
}

\begin{center}\rule{0.5\linewidth}{0.5pt}\end{center}

\subsection*{Connections Across the
Textbook}\label{connections-across-the-textbook-8}
\addcontentsline{toc}{subsection}{Connections Across the Textbook}

\begin{itemize}
\tightlist
\item
  \textbf{Photosynthesis} (\cref{sec:unit_III_photosynthesis}) is the
  biochemical foundation of plant carbon metabolism; this unit extends
  it by covering C\texorpdfstring{\textsubscript{4}}{4} and CAM pathways and carbon assimilation in leaves.
\item
  \textbf{Water potential and osmosis}
  (\cref{sec:unit_I_atoms_molecules};
  \cref{sec:unit_II_membrane_transport}) directly underlie the
  transpiration stream model.
\item
  \textbf{Plant hormone receptor signalling} (auxin, ABA) uses the same
  signalling classes covered in Unit II (cell signalling: receptor
  kinases, second messengers).
\item
  \textbf{Plant-pollinator and plant-mycorrhizal interactions} connect
  to Unit X (community ecology, mutualism, co-evolution).
\item
  \textbf{Secondary metabolites as medicines} (alkaloids, terpenoids)
  link to Unit I (functional groups) and Unit VII (antimicrobial
  compounds).
\end{itemize}

\begin{quote}
\textbf{Key vocabulary introduced here:} meristem, apical dominance,
cohesion-tension, water potential (\texorpdfstring{\ensuremath{\Psi}}{Psi}), transpiration, turgor pressure,
phloem, xylem, stomata, guard cell, photoperiodism, phytochrome, auxin,
gibberellin, cytokinin, abscisic acid (ABA), ethylene, alternation of
generations, angiosperm, gymnosperm, double fertilisation.
\end{quote}

\subsection*{Computational Toolbox --- Unit
VIII}\label{computational-toolbox-unit-viii}
\addcontentsline{toc}{subsection}{Computational Toolbox --- Unit VIII}

\begin{Shaded}
\begin{Highlighting}[]
\ImportTok{from}\NormalTok{ biology.botany }\ImportTok{import}\NormalTok{ water\_potential, transpiration\_flux, photosynthesis\_rate}

\CommentTok{\# Water potential: mesophyll cell with 0.30 M solutes and 0.20 MPa turgor}
\NormalTok{psi }\OperatorTok{=}\NormalTok{ water\_potential(solute\_concentration\_M}\OperatorTok{=}\FloatTok{0.30}\NormalTok{, turgor\_pressure\_MPa}\OperatorTok{=}\FloatTok{0.20}\NormalTok{)}
\BuiltInTok{print}\NormalTok{(}\SpecialStringTok{f"Cell water potential Ψ = }\SpecialCharTok{\{}\NormalTok{psi}\SpecialCharTok{.}\NormalTok{water\_potential\_MPa}\SpecialCharTok{:.2f\}}\SpecialStringTok{ MPa"}\NormalTok{)}
\CommentTok{\# Expected: Ψ ≈ {-}0.54 MPa}
\CommentTok{\# Water moves from higher Ψ (soil, ≈ {-}0.03 MPa) to lower Ψ (leaf, ≈ {-}1.5 MPa)}

\CommentTok{\# Stomatal conductance and transpiration}
\CommentTok{\# g = stomatal conductance; concentration gradient is leaf minus air vapour}
\NormalTok{trans }\OperatorTok{=}\NormalTok{ transpiration\_flux(}
\NormalTok{    stomatal\_conductance\_mol\_m2\_s}\OperatorTok{=}\FloatTok{0.20}\NormalTok{,}
\NormalTok{    internal\_vapor\_conc\_mol\_m3}\OperatorTok{=}\FloatTok{1.25}\NormalTok{,}
\NormalTok{    external\_vapor\_conc\_mol\_m3}\OperatorTok{=}\FloatTok{0.85}\NormalTok{,}
\NormalTok{)}
\BuiltInTok{print}\NormalTok{(}\SpecialStringTok{f"Transpiration flux = }\SpecialCharTok{\{}\NormalTok{trans}\SpecialCharTok{.}\NormalTok{flux\_mmol\_m2\_s}\SpecialCharTok{:.1f\}}\SpecialStringTok{ mmol m⁻² s⁻¹"}\NormalTok{)}
\BuiltInTok{print}\NormalTok{(}\SpecialStringTok{f"Net photosynthesis at PAR 1000 = }\SpecialCharTok{\{}\NormalTok{photosynthesis\_rate(}\DecValTok{1000}\NormalTok{)}\SpecialCharTok{:.1f\}}\SpecialStringTok{ µmol CO₂ m⁻² s⁻¹"}\NormalTok{)}
\CommentTok{\# Expected:}
\CommentTok{\# Transpiration flux = 80.0 mmol m⁻² s⁻¹}
\CommentTok{\# Net photosynthesis at PAR 1000 = 9.6 µmol CO₂ m⁻² s⁻¹}
\end{Highlighting}
\end{Shaded}

\begin{quote}
\textbf{Try it yourself:} Compare
\breaktt{stomatal\_conductance\_mol\_m2\_s\ =\ 0.05} (drought stress,
ABA-closed stomata) vs \texttt{0.3} (well-watered). How does the CO\texorpdfstring{\textsubscript{2}}{2}
assimilation rate (linked to photosynthesis) change proportionally?
\end{quote}

\begin{center}\rule{0.5\linewidth}{0.5pt}\end{center}

\emph{Source note: botany helpers support water potential,
transpiration, photosynthesis, and C3/C4/CAM pathway comparisons.}
\emph{Figures: \breaktt{src/visualization/} (water potential diagrams,
transpiration curves); \breaktt{src/mermaid/biology\_diagrams.py} (plant
life cycle, hormone signalling diagrams).}

\newpage

\section{Plant Structure, Water Relations, and
Transport}\label{plant-structure-water-relations-and-transport}

\label{sec:unit_VIII_plant_structure_and_water}

\begin{quote}
\textbf{Ch 25} \texorpdfstring{\ensuremath{\cdot}}{.} Level 2/3 \texorpdfstring{\ensuremath{\cdot}}{.} 55 min read \texorpdfstring{\ensuremath{\cdot}}{.} 75 min lecture \texorpdfstring{\ensuremath{\cdot}}{.}
Prerequisites: \cref{sec:unit_II_membrane_transport}
\end{quote}

\subsection{Learning Objectives}\label{learning-objectives-27}

\begin{enumerate}
\def\labelenumi{\arabic{enumi}.}
\tightlist
\item
  Describe the major organ systems of
  \hyperref[gl:angiosperm]{\textbf{angiosperm}}s: root system (taproot
  vs fibrous), shoot system, and leaf.
\item
  Compare the three plant tissue systems: dermal, ground, and vascular,
  including cell types within each.
\item
  Explain water relations in plants using the four-component water
  potential equation \(\Psi = \Psi_s + \Psi_p + \Psi_m + \Psi_g\).
\item
  Describe \hyperref[gl:transpiration]{\textbf{transpiration}} and the
  cohesion-tension \citep{dixon1894}-adhesion (TACT) mechanism of water
  ascent and critically evaluate its supporting evidence.
\item
  Explain the apoplast vs symplast pathways and the role of the
  Casparian strip and suberin lamella.
\item
  Explain phloem transport and the Münch pressure-flow hypothesis
  \citep{munch1930}, including apoplastic vs symplastic loading.
\item
  Describe stomatal regulation and the role of ABA in drought response
  at the molecular level.
\item
  Explain nutrient uptake strategies including mycorrhizal associations
  and quantitative turgor relations.
\item
  Quantify water-use efficiency at the leaf level and connect it to C3,
  C4, and CAM photosynthetic trade-offs.
\end{enumerate}

\subsubsection{Study Blueprint}\label{study-blueprint-27}

\begin{itemize}
\tightlist
\item
  \textbf{Big idea:} Plant form is an engineering solution to water
  movement, support, gas exchange, and growth.
\item
  \textbf{Core concepts:} xylem, phloem, water potential, transpiration.
\item
  \textbf{Framework alignment:} Vision \& Change: Structure and
  function, Pathways and transformations of energy and matter, Systems;
  AP Biology: Energetics, Systems Interactions; NGSS-style topics:
  Structure and Function, Matter and Energy in Organisms and Ecosystems.
\item
  \textbf{Model or quantitative lens:} Water-potential and
  transpiration-flux calculations.
\item
  \textbf{Data skill:} Interpret plant-water data from pressure, solute,
  and humidity measurements.
\item
  \textbf{Practice cadence:} Visual Representations, Questions and
  Methods, Argumentation.
\item
  \textbf{Common misconception to repair:} Water does not move because
  plants pull with intention; it follows potential gradients and
  cohesion.
\item
  \textbf{Primary lab:}
  \cref{sec:lab_unit_VIII_plant_structure_and_water}.
\item
  \textbf{Question bank:}
  \cref{sec:q_unit_VIII_plant_structure_and_water}.
\item
  \textbf{Transfer task:} Transfer plant-water reasoning to drought,
  irrigation, forest physiology, and crop breeding.
\item
  \textbf{Bridge to computation:}
  \breaktt{biology.botany.botany.water\_potential}.
\end{itemize}

\begin{figure}[htbp]
\centering
\includegraphics[width=0.85\textwidth]{../figures/water_potential_transpiration.png}
\caption{Plant water-potential and transpiration calculations. The left panel shows solute and total water potential as solute concentration changes; the right panel shows Fick-law transpiration flux as the external vapour gradient and stomatal conductance vary.}
\label{fig:unit_VIII_water_potential_transpiration}
\end{figure}

\begin{center}\rule{0.5\linewidth}{0.5pt}\end{center}

\begin{quote}
\textbf{Opening Vignette --- How a Giant Sequoia Pulls Water 100 Metres
Up}

A mature coast redwood (\emph{Sequoia sempervirens}) can stand 115
metres tall. Getting water from soil to the topmost leaves requires
lifting it against a gravitational pull of roughly 1.1 MPa --- while
simultaneously overcoming resistance to flow through millions of
microscopic \hyperref[gl:xylem]{\textbf{xylem}} conduits. How? The
answer, worked out by Henry Dixon and John Joly in 1894
\citep{dixon1894} and validated with pressure-bomb measurements in the
1960s, is cohesion-tension: water molecules cohere to each other and
adhere to xylem walls so strongly that evaporation from leaf
\hyperref[gl:stomata]{\textbf{stomata}} creates a continuous
negative-pressure column (tension) that pulls water up from the roots.
The column can sustain tensions as low as \texorpdfstring{\ensuremath{-}}{-}10 MPa before xylem embolism
shatters it. This entirely passive, thermodynamically driven mechanism
moves trillions of litres of water into Earth's atmosphere annually ---
and is so elegant that engineers studying microfluidics still struggle
to replicate it artificially.
\end{quote}

\subsection{The Plant Body Plan}\label{the-plant-body-plan}

Flowering plants (angiosperms, about 300,000 species) are
organised into two major organ systems that work together to acquire
resources from both soil and atmosphere.

\subsubsection{Root System}\label{root-system}

The root system anchors the plant, absorbs water and minerals, and
stores nutrients. Two major architectural types exist:

\begin{itemize}
\tightlist
\item
  \textbf{Taproot system:} A single
  \hyperref[gl:dominant]{\textbf{dominant}} primary root (derived from
  the embryonic radicle) with smaller lateral roots branching off.
  Common in eudicots (e.g., carrots, dandelions, oaks). Deep taproot
  systems can access water tables several metres below the surface.
\item
  \textbf{Fibrous root system:} A dense mat of similarly sized
  adventitious roots arising from the stem base. The primary root is
  short-lived. Common in monocots (e.g., grasses, wheat, rice).
  Excellent for soil stabilisation and surface nutrient capture.
\end{itemize}

\textbf{Root apical meristem (RAM):} Located near the root tip,
protected by the root cap. The RAM contains slowly dividing cells in the
quiescent centre (QC) surrounded by rapidly dividing initial cells that
produce most root tissues. The root tip is organised into distinct
zones:

\begin{enumerate}
\def\labelenumi{\arabic{enumi}.}
\tightlist
\item
  \textbf{Root cap} -- protects the meristem; secretes mucilage for
  lubrication; columella cells contain starch-filled amyloplasts
  (statoliths) for gravity sensing
\item
  \textbf{Meristematic zone} -- active cell division
\item
  \textbf{Elongation zone} -- cells expand primarily in length (up to
  10-fold); driven by \hyperref[gl:turgor-pressure]{\textbf{turgor
  pressure}} and cell wall loosening
\item
  \textbf{Maturation zone} -- cells differentiate into specialised
  types; root hairs develop here
\end{enumerate}

\textbf{Root hairs} are single elongated epidermal cells that
dramatically amplify surface area (about 600 cm\(^2\) per cm of
root length). Mineral ions are absorbed actively against their
electrochemical gradients by specific transporters (NO\(_3^-\) by
NRT1/NRT2; K\(^+\) by HKT; Fe\(^{2+}\) by IRT1).

\subsubsection{Plant Architecture: Phytomers and
Plastochrons}\label{plant-architecture-phytomers-and-plastochrons}

The shoot is built as a series of repeating modular units called
\textbf{phytomers}. Each phytomer comprises one node, the internode
immediately below it, the leaf attached at the node, and the axillary
bud subtending the leaf. Phytomer iteration generates the entire shoot
--- a tree is, formally, a clonal colony of phytomers expressed by a
single genotype.

The temporal scaffolding is provided by the \textbf{plastochron index
(PI)}, defined as the time interval between successive leaf-primordium
initiations at the SAM (Erickson and Michelini, 1957). At constant
temperature the plastochron is approximately constant for a given
genotype (about 24 h in \emph{Arabidopsis}, about 36 h
in maize), permitting developmental time to be expressed in plastochron
units rather than chronological time. The \textbf{leaf plastochron index
(LPI)} ranks leaves by developmental age --- LPI 0 is the youngest
visible leaf, LPI 1 the next, and so on --- and is the standard
developmental clock for maize and tobacco physiology experiments.

\textbf{Meristematic zones at the SAM} are organised radially and
clonally:

\begin{itemize}
\tightlist
\item
  \textbf{L1 (tunica outer layer):} Strictly anticlinal divisions; gives
  rise to epidermis
\item
  \textbf{L2 (tunica inner layer):} Mostly anticlinal; produces
  mesophyll, gametes
\item
  \textbf{L3 (corpus):} Anticlinal and periclinal divisions; produces
  vascular and ground tissue
\end{itemize}

Genetic chimeras (e.g., variegated geraniums) reveal these clonal layers
as stable, parallel cell lineages whose interactions specify organ
identity.

\subsubsection{Shoot System}\label{shoot-system}

The shoot system comprises the stem and leaves and is responsible for
\hyperref[gl:photosynthesis]{\textbf{photosynthesis}}, reproduction, and
support.

\textbf{Shoot apical meristem (SAM):} Maintains a pool of pluripotent
stem cells at the shoot tip throughout the plant's life. The SAM is
organised into:

\begin{itemize}
\tightlist
\item
  \textbf{Central zone (CZ):} Slowly dividing stem cells maintained by
  the WUS-CLV3 negative feedback loop
\item
  \textbf{Peripheral zone (PZ):} Faster-dividing cells that form lateral
  organ primordia (leaves, flowers)
\item
  \textbf{Rib zone:} Cells below the CZ that contribute to stem
  internodes
\end{itemize}

Leaves are initiated at the PZ where PIN1-mediated auxin maxima form,
producing the characteristic phyllotaxis pattern (often 137.5 degrees
divergence angle).

\subsubsection{Leaf Anatomy}\label{leaf-anatomy}

Leaf cross-section layers (adaxial to abaxial):

\begin{enumerate}
\def\labelenumi{\arabic{enumi}.}
\tightlist
\item
  Upper epidermis + cuticle (waxy layer of cutin; reduces water loss)
\item
  Palisade mesophyll (densely packed columnar cells with abundant
  \hyperref[gl:chloroplast]{\textbf{chloroplast}}s; principal
  photosynthetic layer)
\item
  Spongy mesophyll (loosely packed; CO\(_2\) diffusion through
  intercellular air spaces; about 30-40\% air space)
\item
  Lower epidermis + cuticle
\end{enumerate}

\textbf{Stomata:} Pores bounded by two kidney-shaped
\textbf{\hyperref[gl:guard-cell]{guard cell}s}. Stomatal aperture is
controlled by guard cell turgor:

\begin{itemize}
\tightlist
\item
  Light activates blue-light phototropins and red-light photosynthesis
  in guard cell chloroplasts. K\(^+\) influx (via K\(^+\)
  inward-rectifying channels) lowers osmotic potential. Water enters by
  \hyperref[gl:osmosis]{\textbf{osmosis}}. Guard cells swell and bow
  apart. Stomata open.
\item
  Drought triggers ABA signalling. ABA activates SnRK2 kinases. SLAC1
  anion channels open, releasing Cl\(^-\). Membrane depolarises. K\(^+\)
  leaves via GORK channels. Water follows osmotically. Guard cells
  deflate. Stomata close.
\end{itemize}

Typical stomatal density: about 200 stomata/mm\(^2\) in many
mesophytes. Each stoma can open or close in minutes. Stomata allow
CO\(_2\) in and O\(_2\) plus water vapour out -- the inevitable tradeoff
between photosynthesis and transpiration.

\begin{quote}
\textbf{Clinical Connection:} Understanding stomatal biology is
essential for crop engineering. Overexpression of the SLAC1
\hyperref[gl:gene]{\textbf{gene}} in barley produced plants with more
responsive stomata and 40\% greater water-use efficiency in field
trials, directly relevant to agriculture in drought-prone regions.
\end{quote}

\textbf{Concept Check 1:} Why are most stomata located on the lower
(abaxial) epidermis of leaves rather than the upper surface?

\textbf{Concept Check 2:} A botanist counts an average plastochron of 30
hours in a tomato plant. After 15 days, approximately how many leaf
primordia will have been initiated by the SAM, assuming temperature is
constant?

\begin{center}\rule{0.5\linewidth}{0.5pt}\end{center}

\subsection{Plant Tissue Systems}\label{plant-tissue-systems}

Most plant organs are composed of three tissue systems: dermal, ground,
and vascular.

\begin{figure}[htbp]
\centering
\includegraphics[width=0.82\linewidth,height=4.2in,keepaspectratio]{../figures/mermaid_inline/inline_mermaid_0137_919e1802461e.png}
\caption{Organisation of the three plant tissue systems showing major cell types within each Dermal tissue forms the outer covering, ground tissue fills the interior, and vascular tissue provides transport.}
\end{figure}

\subsubsection{Dermal Tissue System}\label{dermal-tissue-system}

The dermal tissue system is the outer protective covering of the plant:

\begin{itemize}
\tightlist
\item
  \textbf{Epidermis:} A single layer of tightly packed cells covering
  most surfaces. Secretes the cuticle.
\item
  \textbf{Cuticle:} A waxy layer composed of cutin (a polyester of
  hydroxylated fatty acids) and epicuticular waxes. Minimises water loss
  from non-stomatal surfaces. The cuticle can reduce transpiration by
  95\% compared to an unprotected surface.
\item
  \textbf{Trichomes:} Hair-like epidermal appendages. Diverse functions:
  reduce wind speed at leaf surface (boundary layer), reflect excess
  light, secrete defensive compounds (glandular trichomes of tomato
  produce toxic sesquiterpenes), deter
  \hyperref[gl:herbivory]{\textbf{herbivory}}.
\item
  \textbf{Root hairs:} Tubular extensions of single epidermal cells in
  the root maturation zone. Increase absorptive surface area by
  2-10-fold.
\end{itemize}

\subsubsection{Ground Tissue System}\label{ground-tissue-system}

Ground tissue fills the space between dermal and vascular tissues. Three
cell types:

\textbf{Parenchyma cells:} Living at maturity. Thin primary cell walls.
Functions include photosynthesis (mesophyll), storage (starch in potato
tubers, sugar in sugar beet roots), and wound healing (dedifferentiate
to form callus tissue). Parenchyma cells retain the ability to divide,
making them critical for regeneration.

\textbf{Collenchyma cells:} Living at maturity. Unevenly thickened
primary cell walls (thicker at corners). Provide flexible mechanical
support to growing organs. Found in petioles and young stems. Celery
``strings'' are collenchyma strands.

\textbf{Sclerenchyma cells:} Typically dead at maturity. Thick,
lignified secondary cell walls. Provide rigid support. Two types: fibres
(elongated cells in bundles; flax, hemp) and sclereids (short, irregular
cells; the gritty texture in pears; the hard shell of walnuts).

\subsubsection{Vascular Tissue System -- Xylem in
Detail}\label{vascular-tissue-system-xylem-in-detail}

The vascular tissue system is the plant's long-distance transport
network.

\textbf{Xylem} (water and mineral transport, upward) develops in two
temporally distinct phases that yield mechanically distinct conduits:

\textbf{Protoxylem} matures while the surrounding tissue is still
elongating. Its conduits must therefore stretch with the growing organ,
so secondary-wall lignification is restricted to \textbf{annular}
(ring-shaped) or \textbf{helical} (spiral) thickenings. These
thickenings allow elongation but provide limited collapse-resistance
under tension; protoxylem conduits are typically narrow (5--25 \texorpdfstring{\ensuremath{\mu}}{mu}m) and
frequently obliterated as the tissue matures, leaving an air-filled
lacuna.

\textbf{Metaxylem} differentiates after elongation has ceased. With no
constraint on extensibility, metaxylem develops continuous lignified
walls with \textbf{scalariform} (ladder-like), \textbf{reticulate}
(net-like), or \textbf{pitted} thickenings. Metaxylem conduits are wider
(50--500 \texorpdfstring{\ensuremath{\mu}}{mu}m in some lianas), longer-lived, and bear the bulk of mature
transpiration flux.

A simple mnemonic: \textbf{proto = first} (extensible, ephemeral),
\textbf{meta = mature} (durable, mechanically robust). The transition
can be visualised in any longitudinal section of an elongating maize
root: closer to the tip, narrow protoxylem dominates; further back in
the maturation zone, metaxylem vessels of much greater calibre come
on-line.

The two principal cell types of xylem differ markedly in mechanics:

\begin{itemize}
\tightlist
\item
  \textbf{Tracheids:} Long (1-5 mm), tapered cells with lignified
  secondary walls. Dead at maturity. Water moves between tracheids
  through bordered pit pairs (thin regions in the cell wall where
  secondary wall is absent). The pit membrane consists of the unmodified
  primary walls of the two adjacent cells plus the middle lamella; in
  conifers it bears a thickened central \textbf{torus} suspended by a
  porous \textbf{margo}, allowing the torus to seal the pit aperture if
  the adjacent conduit cavitates (an \textbf{aspirated pit}). Tracheids
  are found in most vascular plants and constitute the primary
  conducting element in most gymnosperms.
\item
  \textbf{Vessel elements:} Shorter (0.2-1 mm) but wider (up to 500 \texorpdfstring{\ensuremath{\mu}}{mu}m
  diameter in some lianas) than tracheids. Stacked end-to-end with
  \textbf{perforation plates} (partial or complete dissolution of end
  walls) forming continuous tubes called vessels. Found primarily in
  angiosperms. More efficient but more vulnerable to cavitation than
  tracheids.
\end{itemize}

\textbf{Bordered pit mechanics --- the torus-margo valve:}

The conifer bordered pit is one of the most elegant valve structures in
biology. The pit membrane is differentiated into two regions:

\begin{itemize}
\tightlist
\item
  \textbf{Torus:} A thickened, impermeable central disc
  (about 5 \texorpdfstring{\ensuremath{\mu}}{mu}m diameter) made of cellulose microfibrils embedded
  in a lignified matrix
\item
  \textbf{Margo:} A porous mesh of radial cellulose strands surrounding
  the torus, with pore sizes of 50--200 nm
\end{itemize}

Under normal flow, the torus floats centrally and water moves freely
through the porous margo. When one tracheid embolises, the pressure
differential between the air-filled (atmospheric) and water-filled
(negative pressure) conduits displaces the torus laterally; the torus
seats firmly against the pit aperture, \textbf{aspirating} the pit. The
seal is reinforced by surface tension at the air-water interface inside
the margo pores. Air seeding through this valve requires breaking the
meniscus at the smallest margo pore --- typically requiring tensions of
\(-3\) to \(-10\) MPa, which sets the species' cavitation safety margin.

Angiosperm pit membranes lack a torus; instead they have a homogeneous,
semi-porous structure (typical pore size about 5--20 nm). These
membranes resist air seeding at higher tensions per unit area but lack
the most-or-nothing valving of conifer pits --- once a single membrane
breach opens, the entire vessel cavitates rapidly.

\textbf{Mechanical properties --- the safety-efficiency tradeoff:}

The Hagen-Poiseuille equation predicts that volumetric flow scales as
the \emph{fourth power} of conduit radius:

\begin{equation}
Q = \frac{\pi r^4 \Delta P}{8 \eta L}
\label{eq:unit_VIII_hagen_poiseuille}
\end{equation}

where \(Q\) is volumetric flow (m\texorpdfstring{\textsuperscript{3}}{3}/s), \(r\) is conduit radius (m),
\(\Delta P\) is pressure drop (Pa), \texorpdfstring{\ensuremath{\eta}}{eta} is dynamic viscosity (\texorpdfstring{\ensuremath{\approx}}{~} 10\texorpdfstring{\textsuperscript{-}}{-}\texorpdfstring{\textsuperscript{3}}{3} Pa\texorpdfstring{\ensuremath{\cdot}}{.}s
for water at 20 °C), and \(L\) is conduit length (m). Doubling the
radius therefore increases conductance 16-fold. Wide vessels (oaks,
ring-porous wood) dominate spring transpiration; narrow tracheids
(conifers, drought-adapted woody species) trade efficiency for
resistance to embolism propagation, since narrower conduits maintain
larger surface-tension forces relative to the column volume.

\textbf{Phloem} (sugar and organic compound transport, bidirectional):

\begin{itemize}
\tightlist
\item
  \textbf{Sieve tube elements:} Living at maturity but lack nucleus,
  \hyperref[gl:ribosome]{\textbf{ribosome}}s, and tonoplast at maturity.
  Connected end-to-end at sieve plates (modified cell walls perforated
  by sieve pores, 1-15 \texorpdfstring{\ensuremath{\mu}}{mu}m diameter).
  P-\hyperref[gl:protein]{\textbf{protein}} (phloem protein) plugs sieve
  pores when cells are damaged, preventing phloem sap loss.
\item
  \textbf{Companion cells:} Retain full complement of
  \hyperref[gl:organelle]{\textbf{organelle}}s. Connected to sieve tube
  elements by numerous plasmodesmata. Provide metabolic support
  including ATP for active sugar loading. Transfer cells (specialised
  companion cells with wall ingrowths) increase membrane surface area
  for active transport.
\end{itemize}

\textbf{Concept Check 3:} Why must xylem conducting cells be dead at
maturity, while phloem sieve tube elements must remain alive?

\textbf{Concept Check 4:} Using the Hagen-Poiseuille relation, calculate
the ratio of volumetric flow between a vessel of radius 100 \texorpdfstring{\ensuremath{\mu}}{mu}m and a
tracheid of radius 25 \texorpdfstring{\ensuremath{\mu}}{mu}m, holding pressure gradient and length constant.
Explain why drought-prone species may nevertheless favour the narrower
conduit.

\textbf{Concept Check 5:} A botanist examines a longitudinal section of
an emerging shoot from a deciduous tree. She notes spirally banded
conduits near the apex but reticulate-pitted conduits 5 cm behind.
Identify each conduit type and explain the developmental logic of the
difference.

\begin{center}\rule{0.5\linewidth}{0.5pt}\end{center}

\subsection{Water Relations -- Water
Potential}\label{water-relations-water-potential}

Plants control water movement using \textbf{water potential} (\texorpdfstring{\ensuremath{\Psi}}{Psi}, units:
MPa). The complete four-component formulation captures every
thermodynamically relevant contribution to water's free energy:

\begin{equation}
\Psi = \Psi_s + \Psi_p + \Psi_m + \Psi_g
\label{eq:unit_VIII_water_potential}
\end{equation}

where:

\begin{itemize}
\tightlist
\item
  \textbf{\(\Psi_s\)} = osmotic (solute) potential = \(-iCRT\)
  (typically \(\leq 0\); solutes decrease water potential). \(i\) is the
  van 't Hoff dissociation factor, \(C\) is molar concentration, \(R\)
  is the gas constant, \(T\) is temperature in Kelvin.
\item
  \textbf{\(\Psi_p\)} = pressure (turgor) potential. Positive in turgid
  cells (turgor pressure against cell wall). Negative in xylem under
  tension. Can range from +2 MPa in highly turgid cells to \(-10\) MPa
  or lower in xylem of tall trees.
\item
  \textbf{\(\Psi_m\)} = matric potential. Reflects water binding to
  colloidal surfaces (cell walls, soil clay particles, dry seeds).
  Typically \(\leq 0\). Typically negligible in well-hydrated cells but
  dominates in dry seeds (where \(\Psi_m\) can reach \(-100\) MPa) and
  in dry soils.
\item
  \textbf{\(\Psi_g\)} = gravitational potential = \(\rho g h\).
  Significant primarily in tall trees. At 10 m height, \(\Psi_g = -0.1\)
  MPa. At 100 m (coast redwood), \(\Psi_g = -1.0\) MPa.
\end{itemize}

In well-hydrated, free-flowing systems (most active leaf and root cells)
\(\Psi_m\) is folded into \(\Psi_p\) and the simplified three-component
form \(\Psi = \Psi_s + \Psi_p + \Psi_g\) is used. In dry seeds, lichens,
and very dry soils, \(\Psi_m\) must be retained explicitly. See
\cref{eq:unit_VIII_water_potential}.

\subsubsection{Worked Example: Calculating Water
Potential}\label{worked-example-calculating-water-potential}

\textbf{Problem:} A plant cell has a solute concentration of
\(C = 0.3 \text{ mol L}^{-1}\) sucrose (\(i=1\)) at
\(T = 293 \text{ K}\) (20 °C) and a turgor pressure (\(\Psi_p\)) of
\(0.5 \text{ MPa}\). Assuming normal height where \(\Psi_g = 0\) and the
cell is well hydrated so \(\Psi_m \approx 0\), what is the cell's total
\hyperref[gl:water-potential]{\textbf{water potential (\texorpdfstring{\ensuremath{\Psi}}{Psi})}}? If this
cell is placed in an open beaker of pure water (\(\Psi = 0\)), will
water flow into or out of the cell?

\textbf{Solution:}

\begin{enumerate}
\def\labelenumi{\arabic{enumi}.}
\item
  \textbf{Calculate the osmotic potential (\(\Psi_s\)):} Using the van
  't Hoff equation \(\Psi_s = -iCRT\), where
  \(R = 0.00831 \text{ L MPa mol}^{-1}\text{ K}^{-1}\):
  \[ \Psi_s = -(1) \cdot (0.3 \text{ mol L}^{-1}) \cdot (0.00831 \text{ L MPa mol}^{-1}\text{ K}^{-1}) \cdot (293 \text{ K})  \label{eq:unit_VIII_plant_structure_and_water_item_1}\]

  \[ \Psi_s \approx -0.73 \text{ MPa}  \label{eq:unit_VIII_plant_structure_and_water_item_2}\]
\item
  \textbf{Calculate the total water potential (\texorpdfstring{\ensuremath{\Psi}}{Psi}):}
  \[ \Psi = \Psi_s + \Psi_p + \Psi_m + \Psi_g  \label{eq:unit_VIII_plant_structure_and_water_item_3}\]

  \[ \Psi = -0.73 \text{ MPa} + 0.5 \text{ MPa} + 0 + 0  \label{eq:unit_VIII_plant_structure_and_water_item_4}\]

  \[ \Psi = -0.23 \text{ MPa}  \label{eq:unit_VIII_plant_structure_and_water_item_5}\]
\item
  \textbf{Determine the direction of water flow:} Water typically moves
  from higher water potential to lower water potential. Since the beaker
  of pure water is at \(\Psi = 0\) and the cell is at
  \(\Psi = -0.23 \text{ MPa}\), water will flow \textbf{into the cell}
  until the cell's turgor pressure rises sufficiently to reach
  equilibrium (\(\Psi = 0\)).
\end{enumerate}

Water typically moves from \textbf{higher to lower water potential}
(down the \texorpdfstring{\ensuremath{\Psi}}{Psi} gradient).

\subsubsection{Worked Example: A Leaf Cell at Dawn vs
Midday}\label{worked-example-a-leaf-cell-at-dawn-vs-midday}

\textbf{Problem:} A mesophyll cell at dawn (cool, humid, stomata closed)
has \(\Psi_s = -1.1\) MPa, \(\Psi_p = +0.9\) MPa, \(\Psi_m \approx 0\),
\(\Psi_g \approx 0\). By midday (hot, dry, stomata open and
transpiring), the cell loses water: \(\Psi_s = -1.4\) MPa (solutes
concentrated by about 27\% volume loss), \(\Psi_p = +0.2\) MPa.
Compute \texorpdfstring{\ensuremath{\Psi}}{Psi} in both states and identify which component drives the change.

\textbf{Solution:}

\begin{itemize}
\tightlist
\item
  Dawn: \(\Psi = -1.1 + 0.9 = -0.2\) MPa
\item
  Midday: \(\Psi = -1.4 + 0.2 = -1.2\) MPa
\end{itemize}

Total \texorpdfstring{\ensuremath{\Psi}}{Psi} has fallen by 1.0 MPa. Decomposing the change:

\begin{itemize}
\tightlist
\item
  \(\Delta\Psi_s = -0.3\) MPa (small contribution from solute
  concentration)
\item
  \(\Delta\Psi_p = -0.7\) MPa (dominant contribution)
\end{itemize}

The midday plunge is driven primarily by \textbf{loss of turgor}. The
cell remains alive and below the wilting threshold, but it is
approaching the limit at which growth (which requires \(\Psi_p > Y\) in
the Lockhart equation) ceases. This is the molecular basis of midday
``growth pause'' almost universally observed in field crops.

\begin{figure}[htbp]
\centering
\includegraphics[width=0.82\linewidth,height=4.2in,keepaspectratio]{../figures/mermaid_inline/inline_mermaid_0138_d1e54f8471f6.png}
\caption{Components of water potential and the soil-plant-atmosphere continuum (SPAC) Water moves passively down the water potential gradient from soil to atmosphere. The enormous gradient between leaf and atmosphere drives transpiration.}
\end{figure}

\subsubsection{Quantitative Turgor Relations and
Plasmolysis}\label{quantitative-turgor-relations-and-plasmolysis}

\textbf{Turgor pressure} (\(\Psi_p\)) is generated when a cell with low
internal \(\Psi_s\) takes up water until the cell wall resists further
expansion. The cell-wall yield curve is described by the
\textbf{Lockhart equation} for irreversible expansion:

\begin{equation}
\frac{1}{V}\frac{dV}{dt} = \phi (\Psi_p - Y)
\label{eq:unit_VIII_lockhart}
\end{equation}

where \texorpdfstring{\ensuremath{\varphi}}{phi} is wall extensibility (a measure of how readily the wall
yields), \(\Psi_p\) is current turgor, and \(Y\) is the yield threshold
below which the wall behaves elastically (no growth). Cells grow
primarily when \(\Psi_p > Y\). In the elongation zone, expansin-mediated
wall loosening reduces \(Y\), allowing growth even at modest turgor.

\textbf{Plasmolysis} is the separation of the protoplast from the cell
wall when the cell is placed in a hypertonic solution (external
\(\Psi_s\) more negative than internal). Three phases:

\begin{enumerate}
\def\labelenumi{\arabic{enumi}.}
\tightlist
\item
  \textbf{Incipient plasmolysis:} Turgor reaches zero (\(\Psi_p = 0\));
  at this point \(\Psi = \Psi_s\). The protoplast is just losing contact
  with the wall.
\item
  \textbf{Concave plasmolysis:} Protoplast withdraws inward; plasma
  membrane appears scalloped at the corners.
\item
  \textbf{Convex plasmolysis:} Protoplast becomes a small spherical body
  in the centre of the cell; reversible if rehydrated within minutes,
  but prolonged plasmolysis kills the cell.
\end{enumerate}

\textbf{Worked Example --- Quantitative Plasmolysis:} A cell at
\(\Psi_p = 0.5\) MPa, \(\Psi_s = -0.73\) MPa is placed in a sucrose
solution at \(\Psi_s = -1.0\) MPa. Water leaves the cell. As water
exits, \(\Psi_s\) becomes more negative (solutes concentrate; assume
volume change of 20\% from full turgor to incipient plasmolysis). At
incipient plasmolysis, \(\Psi_p = 0\) and \(\Psi_{cell} = \Psi_s\) must
equal \(\Psi_{external} = -1.0\) MPa. Since \(\Psi_s\) scales inversely
with cell volume (\(\Psi_s \cdot V\) \texorpdfstring{\ensuremath{\approx}}{~} constant for ideal solutes), we
calculate the relative volume at incipient plasmolysis as
\(V_p / V_0 = \Psi_{s,0} / \Psi_{s,p} = -0.73 / -1.0 = 0.73\), i.e., the
cell loses about 27\% of its initial volume before
plasmolysing.

At full turgor: \(\Psi_p = -\Psi_s\), so \(\Psi = 0\) (cell is in
equilibrium with pure water).

\begin{quote}
\textbf{Clinical Connection:} Agricultural irrigation management relies
on measuring soil water potential with tensiometers. When soil \texorpdfstring{\ensuremath{\Psi}}{Psi} falls
below \(-1.5\) MPa (the permanent wilting point), most crop plants
cannot extract water. Precision irrigation triggered at species-specific
\texorpdfstring{\ensuremath{\Psi}}{Psi} thresholds saves 20-40\% of water compared to schedule-based
irrigation.
\end{quote}

\textbf{Concept Check 6:} A seed at storage moisture has
\(\Psi_m \approx -100\) MPa, \(\Psi_s \approx -2\) MPa, and
\(\Psi_p = 0\). When placed on moist filter paper at \(\Psi = -0.05\)
MPa, water is rapidly imbibed. Explain which component dominates the
imbibition gradient and why imbibition is a one-way process.

\begin{center}\rule{0.5\linewidth}{0.5pt}\end{center}

\subsection{Water Transport Pathways}\label{water-transport-pathways}

\subsubsection{Apoplast vs Symplast
Pathways}\label{apoplast-vs-symplast-pathways}

Water and dissolved minerals travel from the root surface to the xylem
via two parallel routes:

\textbf{Apoplast pathway:} Water moves through the continuum of cell
walls and intercellular spaces without crossing any membranes. Fast and
non-selective. Driven by mass flow along pressure gradients.

\textbf{Symplast pathway:} Water moves through the continuum of
\hyperref[gl:cytoplasm]{\textbf{cytoplasm}} connected by plasmodesmata
(channels through cell walls, 40-60 nm diameter, lined by ER). Requires
crossing the plasma membrane at entry. Slower but allows selective
transport.

\textbf{Transmembrane pathway:} Water crosses plasma membranes via
aquaporins (water channel proteins, PIP and TIP families;
about 35 aquaporin genes in \emph{Arabidopsis}). Aquaporin
activity is regulated by phosphorylation, \hyperref[gl:ph]{\textbf{pH}},
and Ca\(^{2+}\).

\subsubsection{Root Anatomy in Depth --- The Endodermis as
Filter}\label{root-anatomy-in-depth-the-endodermis-as-filter}

Radial anatomy of a young dicot root (from outside inward):

\begin{enumerate}
\def\labelenumi{\arabic{enumi}.}
\tightlist
\item
  \textbf{Epidermis} with root hairs --- initial absorption surface
\item
  \textbf{Cortex} --- multiple layers of parenchyma; storage and
  apoplastic transport corridor; rich in plasmodesmata
\item
  \textbf{Endodermis} --- single layer; the critical apoplastic
  checkpoint
\item
  \textbf{Pericycle} --- cell layer giving rise to lateral roots
\item
  \textbf{Vascular cylinder (stele)} --- central xylem and phloem
  strands
\end{enumerate}

\textbf{The Casparian strip --- chemistry and synthesis:}

The Casparian strip is a band of \textbf{lignin and suberin} deposited
in the radial and transverse walls of endodermal cells. The two polymers
play distinct, sequential roles:

\begin{itemize}
\tightlist
\item
  \textbf{Lignin} (a phenylpropanoid polymer of monolignol units
  p-coumaryl, coniferyl, and sinapyl alcohol cross-linked by oxidative
  coupling) is deposited first, forming the \textbf{Casparian strip
  proper} --- a discontinuous belt that fuses cell wall to plasma
  membrane. CASP1--CASP5 transmembrane proteins assemble into a
  hexagonal scaffold that recruits ESB1 (a dirigent protein) and PER64
  (a peroxidase), localising lignification to a precise sub-micron belt.
  \emph{casp1 casp3} double mutants have discontinuous strips and leak
  ions into the stele.
\item
  \textbf{Suberin} (an aliphatic-aromatic polyester of long-chain
  \texorpdfstring{\ensuremath{\omega}}{omega}-hydroxy fatty acids and ferulate esters) is deposited later as the
  \textbf{suberin lamella} on the inner face of the cell wall, after the
  Casparian strip is mature. Suberisation is regulated by SGN3-CIF1/CIF2
  receptor-peptide signalling, which monitors the integrity of the
  Casparian strip via leakage of CIF peptides from stele to cortex.
\end{itemize}

\textbf{Why both polymers?} Lignin is hydrophobic and rigid but porous
to small molecules; suberin is even more hydrophobic and seals against
water and ion flow. The two-stage system gives a fast-acting initial
seal (lignin Casparian strip) followed by a slow-developing
reinforcement (suberin lamella) that further restricts symplastic
transport across the inner endodermal face.

\textbf{Bypass routes --- apoplastic leakage at lateral root junctions:}

Where lateral roots emerge from the pericycle, they must rupture the
endodermis and overlying cortex. During the brief window before the
lateral root re-establishes its own endodermis, an \textbf{apoplastic
bypass} allows ions (and pathogens) to enter the stele without crossing
a plasma membrane. This window can be hours to days long.
Quantitatively, about 70\% of soil-borne Na\(^+\) entering the
stele of rice and wheat does so via these junctional leaks rather than
through the proper symplast. The bypass may explain seasonal salt
accumulation in glycophytes (salt-sensitive species) under saline
irrigation.

\textbf{Symplastic vs transmembrane pathways:}

\begin{itemize}
\tightlist
\item
  \textbf{Pure symplastic transport} uses plasmodesmata to bypass plasma
  membranes entirely. Plasmodesmal aperture is regulated by callose
  deposition at the neck region; under stress, callose deposition closes
  plasmodesmata, isolating cells.
\item
  \textbf{Transmembrane pathway} crosses plasma membranes twice per cell
  (entry + exit) via aquaporins for water and selective transporters for
  ions. Slowest but most selective.
\end{itemize}

In practice, most three pathways operate simultaneously; the relative
contribution shifts with developmental stage, salinity, and mycorrhizal
colonisation.

\textbf{Passage cells:} In older roots, most endodermal cells deposit a
continuous suberin lamella covering the entire inner surface of the cell
wall. \textbf{Passage cells} --- endodermal cells that retain a thin
wall and lack the suberin lamella, located opposite protoxylem poles ---
provide selective transport routes through the otherwise sealed older
endodermis.

In rice, the endodermis is supplemented by an inner \textbf{exodermis}
(a Casparian-strip-bearing layer just inside the epidermis), giving a
double filter. Rice ranks among the most salt-tolerant cereals partly
because the exodermis blocks Na\(^+\) at the root surface itself.

\begin{figure}[htbp]
\centering
\includegraphics[width=0.82\linewidth,height=4.2in,keepaspectratio]{../figures/mermaid_inline/inline_mermaid_0139_f70eb6406088.png}
\caption{Apoplast vs symplast routes through a young root and the Casparian-strip checkpoint at the endodermis Most water and ions reaching the xylem must traverse at least one plasma membrane.}
\end{figure}

\begin{center}\rule{0.5\linewidth}{0.5pt}\end{center}

\subsection{Transpiration and the Cohesion-Tension
Mechanism}\label{transpiration-and-the-cohesion-tension-mechanism}

\subsubsection{\texorpdfstring{Historical Context: The Birth of
Cohesion-Tension Theory
\citep{dixon1894}}{Historical Context: The Birth of Cohesion-Tension Theory}}\label{historical-context-the-birth-of-cohesion-tension-theory}

The question of how water reaches the canopy of a tall tree exercised
plant physiologists throughout the 19th century. Early proposals invoked
\textbf{root pressure} (positive xylem pressure generated by active
mineral secretion overnight; can be measured by exudation from cut
stumps at 0.05--0.5 MPa), but this manifestly cannot lift water more
than about 5 m and is absent during the day in most species.
\textbf{Capillarity} in xylem conduits could lift water about 1 m given
typical conduit radii. \textbf{Atmospheric pressure} can support a
column of about 10.3 m. None of these mechanisms could account for water
rising to 100 m in \emph{Sequoia}.

In 1894, Henry Dixon (Trinity College Dublin) and John Joly proposed a
radical alternative: water in the xylem is under \textbf{tension}
(negative pressure), pulled from above by evaporation rather than pushed
from below \citep{dixon1894}. The idea was counter-intuitive --- bulk
water at \(-1\) MPa is metastable; a single nucleation event should
cause spontaneous boiling --- but the cohesive forces of water (hydrogen
bonds, theoretical tensile strength \(-140\) MPa) and the small radii of
xylem conduits (which suppress nucleation) make the metastable state
biologically sustainable.

\subsubsection{Detailed TACT Mechanism}\label{detailed-tact-mechanism}

Water ascent in tall trees is explained by the
\textbf{cohesion-tension-adhesion (TACT) mechanism}, refined by Renner
(1911), and validated with pressure probes and MRI:

\begin{figure}[htbp]
\centering
\includegraphics[width=0.82\linewidth,height=4.2in,keepaspectratio]{../figures/mermaid_inline/inline_mermaid_0140_36d2b095fc82.png}
\caption{The transpiration-cohesion-tension mechanism of water ascent Evaporation at the leaf surface generates tension (negative pressure) that is transmitted through the continuous water column in the xylem, pulling water upward from the roots. Cohesion between water molecules and adhesion to xylem walls maintain the column.}
\end{figure}

\begin{enumerate}
\def\labelenumi{\arabic{enumi}.}
\item
  \textbf{Transpiration} from leaves drives evaporation at the mesophyll
  cell wall surface. Water evaporates from the thin film coating the
  cell wall into the substomatal air space, then diffuses out through
  open stomata.
\item
  As water evaporates, the air-water interface retreats into the
  nanoscale pores of the cell wall. Surface tension at this curved
  meniscus generates a large negative pressure (tension), following the
  Young-Laplace equation:
\end{enumerate}

\begin{equation}
\Delta P = \frac{2\gamma \cos\theta}{r}
\label{eq:unit_VIII_young_laplace}
\end{equation}

where \texorpdfstring{\ensuremath{\gamma}}{gamma} is surface tension of water (0.0728 N/m), \texorpdfstring{\ensuremath{\theta}}{theta} is contact angle,
and \(r\) is pore radius. For cell wall pores (\(r \approx\) 5-10 nm),
this generates tensions of \(-10\) to \(-30\) MPa.

\begin{enumerate}
\def\labelenumi{\arabic{enumi}.}
\setcounter{enumi}{2}
\item
  This tension propagates through the continuous water columns in xylem
  conduits (tracheids and vessels) because of water's strong
  \textbf{cohesion} (\hyperref[gl:hydrogen-bond]{\textbf{hydrogen
  bond}}s). The theoretical tensile strength of water is approximately
  \(-140\) MPa; in practice, xylem water columns sustain tensions of
  \(-1\) to \(-10\) MPa, with some trees reaching \(-15\) MPa.
\item
  \textbf{Adhesion} of water to the hydrophilic, lignified xylem walls
  prevents the water column from pulling away from the conduit walls.
\item
  The tension at the top of the column draws water from the roots, which
  in turn absorbs water from the soil down the water potential gradient.
\end{enumerate}

\subsubsection{Critical Evaluation of the Cohesion-Tension Theory:
Embolism Debate and Synchrotron
Resolution}\label{critical-evaluation-of-the-cohesion-tension-theory-embolism-debate-and-synchrotron-resolution}

The cohesion-tension model is not without controversy, and its mature
form represents a triumph of integrating ecological, physical, and
microscopic evidence --- together with rebuttals to specific objections.

\textbf{Lines of supporting evidence:}

\begin{enumerate}
\def\labelenumi{\arabic{enumi}.}
\tightlist
\item
  \textbf{Pressure-bomb measurements} (Scholander \emph{et al.}, 1965)
  directly demonstrate that excised stems of transpiring trees release
  water primarily when external pressure is applied, consistent with
  negative xylem pressure prior to excision. Pressure required matches
  predictions of the tension hypothesis (typically \(-1\) to \(-3\) MPa
  for well-watered crops; \(-6\) to \(-10\) MPa under drought).
\item
  \textbf{Magnetic resonance imaging (MRI)} in vivo (Holbrook \emph{et
  al.}, 2001) visualises continuous water columns in mature angiosperm
  trunks, with abrupt transitions to embolised conduits matching
  predictions of tension-induced cavitation.
\item
  \textbf{Pressure probes} (Steudle, 1995) in living root cortex cells
  confirm steep tension gradients between epidermis and stele consistent
  with bulk-flow-driven uptake.
\item
  \textbf{Acoustic emission} of xylem during drought (Tyree and Sperry,
  1989) records sharp clicks corresponding to individual cavitation
  events, the count of which scales with drought severity.
\end{enumerate}

\textbf{The embolism debate:}

Critics (most prominently U. Zimmermann and colleagues, 1990s--2000s)
argued that water columns under tensions of \(-1\) MPa or more should be
metastable and would cavitate continuously, making sustained negative
pressures of \(-10\) MPa physiologically implausible. The Zimmermann
school proposed alternative ``tissue tension'' models in which water
moved by mechanisms other than bulk flow.

This challenge has been progressively resolved by direct visualisation:

\begin{itemize}
\tightlist
\item
  \textbf{Improved pressure probes} (Wei \emph{et al.}, 1999) confirmed
  predicted tensions in living conduits without artefact.
\item
  \textbf{Synthetic-tree experiments} (Wheeler and Stroock,
  \emph{Nature} 2008) constructed transparent microfluidic devices
  replicating the geometry of a leaf's evaporative surface; these
  ``synthetic trees'' sustained water under tension to \(-22\) MPa,
  mirroring biological systems and proving that cell wall pore geometry
  alone is sufficient to support the predicted tensions.
\item
  \textbf{Synchrotron X-ray micro-CT} (Brodersen \emph{et al.},
  \emph{Plant Physiol.} 2010; Choat \emph{et al.}, \emph{Nature} 2012)
  imaged living xylem in real time and showed that under moderate
  drought, embolisms are confined to a small fraction of conduits ---
  settling the long-running argument over how prevalent embolism is in
  actively transpiring plants. Synchrotron imaging also visualised
  \textbf{embolism repair} under daily root pressure: embolised conduits
  refill overnight in \emph{Vitis} and \emph{Laurus}.
\end{itemize}

\textbf{Mechanism of embolism repair:} Two complementary mechanisms are
now well-attested:

\begin{itemize}
\tightlist
\item
  \textbf{Root pressure refilling:} Active mineral secretion into xylem
  at night generates positive xylem pressure (0.05--0.5 MPa) sufficient
  to push water back into embolised conduits. Common in herbaceous
  species and in many trees overnight.
\item
  \textbf{Phloem-driven refilling:} Phloem-derived sugars osmotically
  draw water into embolised conduits via parenchyma cells lining the
  xylem (Salleo \emph{et al.}, 2004; Nardini \emph{et al.}, 2011). This
  mechanism operates against negative bulk xylem pressure and is
  observed during the day.
\end{itemize}

\textbf{Net assessment:} The TACT mechanism is correct in its
essentials. Water in xylem is genuinely under tension; the column is
stabilised by cohesion and adhesion at scales where pore geometry
sustains negative pressures far beyond what bulk water could endure. The
system is metastable but biologically sufficient --- and the cost of
occasional cavitation is mitigated by redundant pathways and active
repair.

\subsubsection{Xylem Transport
Parameters}\label{xylem-transport-parameters}

\textbf{Xylem sap velocity} measured by MRI and heat pulse methods: 3-45
m/h in trees; 15-50 m/h in herbaceous plants. Ring-porous trees (oak,
elm) can reach velocities of 40+ m/h in their large early-wood vessels.

\textbf{Transpiration flux} (Fick's first law for water vapour):

\begin{equation}
E = g_s \cdot \Delta w
\label{eq:unit_VIII_transpiration_flux}
\end{equation}

where \(g_s\) = stomatal conductance (mmol m\(^{-2}\) s\(^{-1}\)) and
\(\Delta w\) = vapour pressure difference between leaf interior and
ambient air (mol mol\(^{-1}\)).

A well-watered broadleaf tree can transpire 200-400 litres of water per
day. A large oak may transpire over 150,000 litres during a growing
season.

\subsubsection{Cavitation and Embolism}\label{cavitation-and-embolism}

When xylem tension exceeds the capacity of the water column, dissolved
gas comes out of solution forming a bubble -- \textbf{cavitation}. The
gas bubble expands to fill the conduit, creating an \textbf{embolism}
(air lock) that blocks water flow.

Plants have evolved multiple strategies to manage cavitation:

\begin{itemize}
\tightlist
\item
  \textbf{Vessel element size:} Smaller conduits are more resistant to
  cavitation (higher surface-tension forces maintain the column) but
  less efficient for transport. This creates a safety-efficiency
  tradeoff.
\item
  \textbf{Pit membranes:} In tracheids, bordered pits have a torus-margo
  structure. When one tracheid cavitates, the torus seals against the
  pit aperture, preventing air from spreading to adjacent conduits.
\item
  \textbf{Embolism repair:} Some plants can refill embolisms overnight
  when transpiration ceases and root pressure (positive xylem pressure
  generated by active mineral secretion into xylem) pushes water back
  into embolised conduits.
\item
  \textbf{Redundant pathways:} Multiple parallel conduits ensure that
  loss of some to embolism does not stop water transport entirely.
\end{itemize}

\begin{quote}
\textbf{Clinical Connection:} The vulnerability of xylem to cavitation
explains why drought-adapted species (e.g., desert shrubs) have narrow
xylem vessels while tropical rainforest trees have wide ones. Climate
change-driven drought is causing widespread forest die-off globally as
tree xylem systems experience unprecedented cavitation stress, termed
``hydraulic failure.'' The 2011 Texas drought killed an estimated 300
million trees.
\end{quote}

\textbf{Concept Check 8:} Modern synchrotron imaging shows that conifer
xylem can withstand tensions of \(-6\) MPa with little embolism, while
many ring-porous angiosperms cavitate at \(-2\) MPa. Explain this
difference in terms of pit membrane structure (torus-margo vs
homogeneous) and conduit dimensions.

\begin{center}\rule{0.5\linewidth}{0.5pt}\end{center}

\subsection{Stomatal Regulation}\label{stomatal-regulation}

Stomata represent the critical control point where plants balance carbon
gain (CO\(_2\) uptake for photosynthesis) against water loss
(transpiration).

\subsubsection{Stomatal Opening --- Guard Cell Ion
Fluxes}\label{stomatal-opening-guard-cell-ion-fluxes}

The mechanism of stomatal opening involves a coordinated cascade of ion
fluxes that generate the osmotic gradient driving water entry. The
molecular logic is one of the best-characterised electrophysiological
systems in plants:

\begin{enumerate}
\def\labelenumi{\arabic{enumi}.}
\tightlist
\item
  \textbf{Blue light (peak 450 nm)} activates phototropin receptors
  (PHOT1/PHOT2) on guard cells via FMN photoreduction
\item
  \textbf{PHOTs activate plasma membrane H\(^+\)-ATPases} (AHA1, AHA2)
  via 14-3-3-mediated phosphorylation of the C-terminal autoinhibitory
  domain
\item
  \textbf{H\(^+\) extrusion hyperpolarises} the membrane (interior
  becomes more negative, \(V_m\) shifts from \(-100\) to \(-180\) mV)
  and acidifies the apoplast (pH 6 \texorpdfstring{\ensuremath{\to}}{->} 5)
\item
  \textbf{K\(^+\) inward-rectifying channels} (KAT1 in
  \emph{Arabidopsis}; voltage-gated, six-transmembrane Shaker family)
  open at hyperpolarised potentials. K\(^+\) flows down its
  electrochemical gradient into the guard cell, accumulating at
  concentrations of 400--800 mM. KAT1 is the rate-limiting K\(^+\) entry
  channel for stomatal opening; \emph{kat1} knockouts have severely
  reduced light-induced opening.
\item
  \textbf{Cl\(^-\) accumulates} as a charge-balancing
  \textbf{counterion} through SLAH3 and CLC channels operating in
  reverse mode under hyperpolarisation
\item
  \textbf{Malate\(^{2-}\)} is synthesised from chloroplast-derived
  starch via PEP carboxylase and converted by NADP-malate dehydrogenase,
  providing additional charge-balancing osmoticum
\item
  \textbf{Sucrose} accumulates in late-afternoon stomata, supplementing
  K\(^+\)/malate as an osmoticum (afternoon stomata may be 50\%
  sucrose-driven)
\item
  \textbf{Decreased osmotic potential} (\(\Psi_s\) more negative, by
  1.5--2 MPa) drives water influx through aquaporins (PIP1, PIP2)
\item
  \textbf{Guard cells swell.} Their differentially thickened walls
  (thicker on the inner pore-facing wall, thinner on the outer wall)
  cause the cells to bow apart asymmetrically
\item
  \textbf{Stomatal aperture} opens by 5--25 \texorpdfstring{\ensuremath{\mu}}{mu}m depending on species
\end{enumerate}

\subsubsection{Stomatal Closure -- ABA Signalling Cascade in
Detail}\label{stomatal-closure-aba-signalling-cascade-in-detail}

Under drought stress, abscisic acid (ABA) triggers rapid stomatal
closure through a molecular relay that has become a textbook example of
plant signal transduction. The pathway can be summarised in five steps:

\begin{quote}
\textbf{ABA \texorpdfstring{\ensuremath{\to}}{->} PYR1 \texorpdfstring{\ensuremath{\to}}{->} PP2C \texorpdfstring{\ensuremath{\to}}{->} SnRK2 \texorpdfstring{\ensuremath{\to}}{->} SLAC1}
\end{quote}

Each step in this cascade has been crystallographically resolved
(PYR1-ABA-PP2C ternary complex, Melcher \emph{et al.}, \emph{Nature}
2009; SnRK2-PP2C complex, Soon \emph{et al.}, \emph{Science} 2012) and
is now textbook signalling.

\begin{enumerate}
\def\labelenumi{\arabic{enumi}.}
\tightlist
\item
  \textbf{ABA synthesis:} Soil drying activates NCED3
  (9-\emph{cis}-epoxycarotenoid dioxygenase) in roots, cleaving
  9'-\emph{cis}-violaxanthin/neoxanthin to xanthoxin (ABA precursor).
  Local ABA biosynthesis also occurs in leaf vasculature and guard cells
  themselves.
\item
  \textbf{ABA distribution:} ABA travels via xylem to guard cells and is
  concentrated by the AIT1/NPF4 family of ABA importers and the ABCG25
  exporter.
\item
  \textbf{Receptor binding:} ABA binds \textbf{PYR/PYL/RCAR} receptors
  (14 paralogs in \emph{Arabidopsis}) at a START-domain pocket. This
  binding ``closes'' the receptor's gate loop over the binding pocket.
\item
  \textbf{PP2C inhibition:} The closed ABA-PYR complex docks onto and
  \textbf{inhibits Type-2C protein phosphatases} (ABI1, ABI2, HAB1).
  PP2Cs normally dephosphorylate (and thereby inactivate) downstream
  kinases.
\item
  \textbf{SnRK2 activation:} Released from PP2C inhibition,
  \textbf{SnRK2 kinases} (especially OST1/SnRK2.6) autophosphorylate and
  become active. OST1 phosphorylates downstream targets including:

  \begin{itemize}
  \tightlist
  \item
    \textbf{SLAC1} anion channel: phosphorylation at Ser120 opens the
    channel; massive Cl\(^-\) and malate\(^{2-}\) efflux
    \textbf{depolarises} the membrane (V\(_m\) shifts to \(-30\) to
    \(-50\) mV)
  \item
    \textbf{GORK} outward-rectifying K\(^+\) channel: opens at
    depolarised potentials; K\(^+\) leaves the cell down its
    electrochemical gradient
  \item
    \textbf{AKT1/KAT1}: inward-rectifying K\(^+\) channels are inhibited
    (preventing K\(^+\) re-entry)
  \item
    \textbf{NADPH oxidase RBOHF}: produces reactive oxygen species (ROS)
    that activate Ca\(^{2+}\)-permeable channels, raising cytosolic
    Ca\(^{2+}\) and amplifying SLAC1/GORK responses
  \end{itemize}
\item
  \textbf{Osmotic deflation:} Loss of K\(^+\), Cl\(^-\), and malate
  raises \(\Psi_s\) by 1.5 MPa; water exits via aquaporins; turgor drops
\item
  \textbf{Aperture closure:} Stomata close within 5-15 minutes
\end{enumerate}

The ABA cascade integrates with a parallel \textbf{CO\(_2\)-sensing
pathway} (HT1 kinase, MPK4/12) and an \textbf{immune pathway} (FLS2
receptor activated by bacterial flagellin) at the level of SLAC1, which
serves as a hub for diverse closure stimuli.

\begin{figure}[htbp]
\centering
\includegraphics[width=0.82\linewidth,height=4.2in,keepaspectratio]{../figures/mermaid_inline/inline_mermaid_0141_4b0dcf4ea02b.png}
\caption{The ABA signalling cascade in guard cells ABA binding releases SnRK2 kinases from PP2C inhibition; SnRK2 then drives the ion-flux changes that osmotically deflate the guard cell.}
\end{figure}

\subsubsection{Other Factors Affecting
Stomata}\label{other-factors-affecting-stomata}

\begin{itemize}
\tightlist
\item
  \textbf{CO\(_2\) concentration:} Elevated internal CO\(_2\) promotes
  closure (reduces need for gas exchange). This is significant under
  rising atmospheric CO\(_2\), where plants partially close stomata,
  reducing transpiration but potentially increasing leaf temperature.
\item
  \textbf{Humidity:} Low humidity (high vapour pressure deficit, VPD)
  promotes closure as a protective response against excessive water loss
\item
  \textbf{Photosynthetically active radiation (PAR):} Drives opening via
  both blue-light specific and photosynthesis-dependent pathways
\item
  \textbf{Circadian clock:} Stomata show anticipatory opening before
  dawn, controlled by the plant circadian clock
\end{itemize}

\subsubsection{Water-Use Efficiency}\label{water-use-efficiency-1}

Stomatal opening trades CO\(_2\) acquisition against water loss. The
intrinsic water-use efficiency captures this trade-off at the leaf
level:

\begin{equation}
iWUE = \frac{A_n}{g_s}
\label{eq:unit_VIII_iwue}
\end{equation}

where \(A_n\) is net CO\(_2\) assimilation rate (\texorpdfstring{\ensuremath{\mu}}{mu}mol CO\(_2\)
m\(^{-2}\) s\(^{-1}\)) and \(g_s\) is stomatal conductance to water
vapour (mol H\(_2\)O m\(^{-2}\) s\(^{-1}\)). Higher iWUE means more
carbon fixed per unit water transpired.

\textbf{C3, C4, and CAM trade-offs:}

\begin{itemize}
\tightlist
\item
  \textbf{C3 plants} (most temperate species; rice, wheat, soybean) have
  iWUE about 50--100 \texorpdfstring{\ensuremath{\mu}}{mu}mol CO\(_2\) mol\(^{-1}\) H\(_2\)O. Their
  primary CO\(_2\)-fixing enzyme RuBisCO operates at near-atmospheric
  {[}CO\(_2\){]} inside the leaf, limiting iWUE.
\item
  \textbf{C4 plants} (maize, sorghum, sugarcane, many tropical grasses)
  concentrate CO\(_2\) in bundle-sheath cells via the C4 pump. Mesophyll
  {[}CO\(_2\){]} is amplified about 10-fold relative to
  ambient, allowing stomata to be partially closed without limiting
  photosynthesis. iWUE about 150--250 \texorpdfstring{\ensuremath{\mu}}{mu}mol mol\(^{-1}\).
\item
  \textbf{CAM plants} (succulents, cacti, agaves, pineapple) open
  stomata primarily at night, when VPD is low. CO\(_2\) is fixed as
  malate and stored in vacuoles; during the day, decarboxylation
  releases CO\(_2\) inside closed leaves for photosynthesis. iWUE can
  exceed 500 \texorpdfstring{\ensuremath{\mu}}{mu}mol mol\(^{-1}\) --- an order of magnitude above C3.
\end{itemize}

The progression C3 \texorpdfstring{\ensuremath{\to}}{->} C4 \texorpdfstring{\ensuremath{\to}}{->} CAM represents progressive specialisation for
water-limited environments, with CAM as the most extreme adaptation. The
iWUE differential explains why C4 maize outperforms C3 wheat in hot,
semi-arid climates and why CAM cacti dominate true deserts.

\textbf{Concept Check 9:} A researcher fumigates a leaf with elevated
CO\(_2\) while simultaneously illuminating it with blue light. Predict
the net effect on stomatal aperture and explain the competing signals at
the level of OST1/SLAC1.

\textbf{Concept Check 10:} An \emph{ost1} loss-of-function
\emph{Arabidopsis} mutant is subjected to drought. Predict its phenotype
relative to wild type. What would happen if the same mutant carried a
constitutively active SLAC1 transgene?

\begin{center}\rule{0.5\linewidth}{0.5pt}\end{center}

\subsection{Phloem Transport -- The Münch Pressure-Flow
Model}\label{phloem-transport-the-muxfcnch-pressure-flow-model}

Unlike xylem (upward, driven by transpiration tension), \textbf{phloem}
transports photosynthate (sucrose at 0.3-0.9 M plus amino acids,
\hyperref[gl:hormone]{\textbf{hormone}}s, mRNAs, and small proteins)
\textbf{bidirectionally} from sources to sinks.

\subsubsection{Phloem Loading --- Apoplastic vs Symplastic
Pathways}\label{phloem-loading-apoplastic-vs-symplastic-pathways}

How sugars enter sieve tubes from photosynthetic mesophyll determines
key physiological properties: loading rate, concentration limits, and
the species-specific tradeoffs between transport efficiency and
protection from herbivores. The choice of loading strategy is one of the
most striking ecophysiological dichotomies in plants.

\textbf{Apoplastic loading} (Type 2; most herbaceous crops, including
most major cereals, \emph{Arabidopsis}, sugar beet, soybean, tomato):

\begin{enumerate}
\def\labelenumi{\arabic{enumi}.}
\tightlist
\item
  Sucrose is exported from mesophyll cells to the apoplast (cell-wall
  space) via \textbf{SWEET11/SWEET12} sucrose efflux carriers
  (uniporters, facilitated diffusion)
\item
  The H\(^+\)-ATPase on the companion-cell plasma membrane creates a
  steep proton gradient (apoplast pH about 5.5; cytoplasm pH
  about 7.5) and a hyperpolarised membrane potential
\item
  \textbf{SUT1} (Sucrose-H\(^+\) symporter; also known as SUC2 in
  \emph{Arabidopsis}) uses the proton gradient to import sucrose against
  its concentration gradient (10--50\texorpdfstring{\ensuremath{\times}}{x} concentration step). The full
  family of \textbf{sucrose transporters (SUTs)} comprises three clades
  with distinct cellular localisations and kinetic properties: SUT1/SUC2
  is the workhorse loader on companion cells (low \(K_m\)
  about 1 mM); SUT4 localises to tonoplasts and mediates
  vacuolar sucrose mobilisation; SUT2 has been proposed as a sucrose
  sensor with high \(K_m\). Knockout of \emph{SUT1} in maize
  (\emph{sut1} mutants) causes sugar accumulation in source leaves and
  starvation in sinks, confirming its role as the rate-limiting loader.
\item
  Imported sucrose moves through plasmodesmata between companion cell
  and sieve element to enter the conducting stream
\item
  Loading rate: 5--15 \texorpdfstring{\ensuremath{\mu}}{mu}mol sucrose per cm\(^2\) of leaf per hour
\end{enumerate}

Apoplastic loading is energetically costly (one ATP per sucrose) but
allows very high sieve-tube concentrations (up to 1 M sucrose) by
``uphill'' transport against the concentration gradient.

\textbf{Symplastic loading} (Type 1; many trees, especially temperate
broadleaves like willow, poplar; also many tropical species):

\begin{enumerate}
\def\labelenumi{\arabic{enumi}.}
\tightlist
\item
  Sucrose moves from mesophyll to companion cells \textbf{entirely
  through plasmodesmata} without crossing any membrane. Plasmodesmal
  density at the mesophyll--companion cell interface is 10--100 \texorpdfstring{\ensuremath{\times}}{x} higher
  in symplastic loaders than in apoplastic loaders. This is the
  diagnostic anatomical feature of symplastic loading --- first
  established by Gamalei (1989) in his comparative survey of
  \textgreater300 species.
\item
  In specialised ``intermediary cells'' (companion-cell variant),
  sucrose is converted to larger oligosaccharides (raffinose, stachyose)
  by \textbf{galactinol synthase + raffinose synthase}. These larger
  oligosaccharides cannot diffuse back through the narrow plasmodesmal
  aperture (about 3 nm size limit at the desmotubule).
\item
  The accumulated raffinose-family oligosaccharides (RFOs) form a
  one-way \textbf{polymer trap} (Turgeon, 1991), lowering \(\Psi_s\) and
  drawing water into the sieve tube.
\end{enumerate}

Symplastic loading is energetically cheap (no membrane transport) but
limited in concentration. Some temperate trees (apple, \emph{Malus}) use
a third hybrid strategy --- passive symplastic loading without polymer
trapping, where the source-sink concentration gradient alone drives
unloading.

\textbf{Plasmodesmatal conductance and regulation:}

Plasmodesmata are dynamic; their aperture is set by the deposition of
callose (\texorpdfstring{\ensuremath{\beta}}{beta}-1,3-glucan) at the neck. Callose synthases (CalS1/3/7/8) and
\texorpdfstring{\ensuremath{\beta}}{beta}-1,3-glucanases (BG\_pap, BG\_ppap) act antagonistically. Under stress
(wounding, pathogen attack), callose deposition closes plasmodesmata
within minutes, isolating affected cells. SUT1-dependent (apoplastic)
loading is unaffected by plasmodesmal closure; symplastic loaders are
critically dependent on open plasmodesmata and are more vulnerable to
wounding-induced flow disruption.

\textbf{Comparison across plant groups:}

{\def\LTcaptype{none} % do not increment counter
\begin{longtable}[]{@{}
  >{\raggedright\arraybackslash}p{(\linewidth - 8\tabcolsep) * \real{0.1188}}
  >{\raggedright\arraybackslash}p{(\linewidth - 8\tabcolsep) * \real{0.1188}}
  >{\raggedright\arraybackslash}p{(\linewidth - 8\tabcolsep) * \real{0.3861}}
  >{\raggedright\arraybackslash}p{(\linewidth - 8\tabcolsep) * \real{0.1782}}
  >{\raggedright\arraybackslash}p{(\linewidth - 8\tabcolsep) * \real{0.1980}}@{}}
\toprule\noalign{}
\begin{minipage}[b]{\linewidth}\raggedright
Loading type
\end{minipage} & \begin{minipage}[b]{\linewidth}\raggedright
Plant groups
\end{minipage} & \begin{minipage}[b]{\linewidth}\raggedright
Plasmodesmata at mesophyll-CC interface
\end{minipage} & \begin{minipage}[b]{\linewidth}\raggedright
Rate-limiting step
\end{minipage} & \begin{minipage}[b]{\linewidth}\raggedright
Phloem sap {[}sucrose{]}
\end{minipage} \\
\midrule\noalign{}
\endhead
\bottomrule\noalign{}
\endlastfoot
Apoplastic (SUT1) & Most herbaceous crops, cereals, \emph{Arabidopsis},
\emph{Solanum} & Few (\textless1 per \texorpdfstring{\ensuremath{\mu}}{mu}m\(^2\)) & SUT1 H\(^+\)-symport &
Up to 1.0 M \\
Symplastic active (polymer trap) & Many temperate broadleaves:
\emph{Salix}, \emph{Populus}, \emph{Cucurbita}, \emph{Coleus} & Many
(\textgreater10 per \texorpdfstring{\ensuremath{\mu}}{mu}m\(^2\)) & Raffinose synthase rate & 0.3--0.8 M \\
Symplastic passive & Some trees (apple, \emph{Quercus}?) & Many & None
(diffusion-limited) & 0.2--0.4 M \\
\end{longtable}
}

\subsubsection{\texorpdfstring{The Münch Pressure-Flow Hypothesis
\citep{munch1930}}{The Münch Pressure-Flow Hypothesis}}\label{the-muxfcnch-pressure-flow-hypothesis}

Ernst Münch's 1930 model remains the consensus explanation for
bidirectional phloem flow, supplemented by about 95 years of
experimental refinement.

\textbf{Mechanism:}

\begin{enumerate}
\def\labelenumi{\arabic{enumi}.}
\tightlist
\item
  \textbf{Sugar loading at source} (mature leaves) raises sucrose
  concentration in sieve tubes to 0.5--1 M, lowering \(\Psi_s\) to
  \(-2\) to \(-3\) MPa
\item
  Water enters by osmosis from adjacent xylem (where \texorpdfstring{\ensuremath{\Psi}}{Psi} is more
  positive). Turgor pressure (\(\Psi_p\)) rises to +1 to +2 MPa at the
  source
\item
  \textbf{Sugar unloading at sink} (roots, developing fruit, growing
  tips, storage organs) removes sucrose. Sucrose may be cleaved by
  invertase (to glucose + fructose) or sucrose synthase. As sucrose is
  removed, \(\Psi_s\) rises, water exits to surrounding tissues, and
  turgor pressure falls to +0.3 to +0.6 MPa
\item
  The \textbf{pressure gradient} between source (high turgor) and sink
  (low turgor) drives bulk flow through the sieve-tube network along the
  path of least resistance
\item
  The unloaded water re-enters the xylem and is recycled to the source
  (closing the loop)
\end{enumerate}

\textbf{Unloading mechanisms} are sink-specific:

\begin{itemize}
\tightlist
\item
  \textbf{Symplastic unloading} (most growing sinks, e.g., root tips,
  developing leaves): sucrose moves through plasmodesmata into
  surrounding cells; intracellular sucrose synthase or invertase cleaves
  it for use.
\item
  \textbf{Apoplastic unloading} (storage sinks isolated by suberin,
  e.g., developing seeds and grains): sucrose is exported by SWEETs into
  the apoplast, then cleaved by cell-wall invertase to glucose +
  fructose, which are imported via hexose transporters. The
  maternal-filial boundary in seeds is typically apoplastic, ensuring
  genetic isolation between mother plant and developing offspring.
\end{itemize}

Flow rates: 0.5-1.5 m/h; rates up to 3 m/h have been measured.

\textbf{Supporting evidence for Münch flow:}

\begin{enumerate}
\def\labelenumi{\arabic{enumi}.}
\tightlist
\item
  \textbf{Aphid-stylet sap collection} (Mittler, 1953): aphids cleanly
  tap individual sieve tubes; collected sap shows sucrose concentrations
  of 10--25\% w/v, consistent with osmotically active loading
\item
  \textbf{Pressure probes} on individual sieve elements directly measure
  positive turgor of +1--2 MPa near the source and lower turgor near
  sinks
\item
  \textbf{Cessation of flow on cooling} the petiole (which inhibits
  H\(^+\)-ATPase but does not block plasmodesmata) confirms active
  loading is required
\item
  \textbf{Source--sink reversal:} Removing the lower leaves of a soybean
  plant reverses flow direction in the stem (downward to upward) within
  hours, consistent with passive bulk flow following the pressure
  gradient
\item
  \textbf{Mathematical modelling} with realistic sieve-tube geometry
  (Thompson and Holbrook, 2003) shows Münch flow can readily account for
  observed transport rates in trees up to about 100 m, provided
  sieve tubes are wide enough
\end{enumerate}

\textbf{Evidence against (and modifications to) Münch flow:}

\begin{enumerate}
\def\labelenumi{\arabic{enumi}.}
\tightlist
\item
  The \textbf{giant-vine objection} (T.W. Mason, 1922): flow rates in
  long vines were measured to exceed Münch's predictions for narrow
  conduits. This was resolved when measurements showed sieve-tube radii
  in such species are actually larger than initially supposed.
\item
  \textbf{Pressure-driven model alone cannot account for observations of
  sieve plate occlusion} (P-protein, callose) in healthy phloem,
  suggesting flow regulation through dynamic plate gating
\item
  \textbf{High-velocity, short-distance flow} in some tropical lianas
  (\textgreater3 m/h) approaches the upper bound of Münch flow and may
  require a supplementary ``relay'' mechanism
\end{enumerate}

The current consensus: Münch pressure-flow is the \textbf{dominant}
mechanism, supplemented by molecular regulators (sieve-plate occlusion,
P-protein gating) that fine-tune flow without overturning the basic
pressure-driven principle.

\begin{quote}
\textbf{Clinical Connection:} Girdling (removal of a ring of bark
including phloem) causes sugar accumulation above the girdle and
starvation below, eventually killing the tree. This principle explains
why bark-stripping by elephants, deer, or porcupines can be lethal.
Conversely, targeted phloem disruption is used as a management tool to
kill invasive tree species.
\end{quote}

\textbf{Concept Check 11:} Predict the consequence for phloem flow if a
chemical inhibitor of the H\(^+\)-ATPase is locally applied to the
source leaf petiole of an apoplastic-loading species. Would the same
treatment have the same effect in a symplastic-loading tree?

\textbf{Concept Check 12:} A maize \emph{sut1} mutant accumulates
sucrose to 200 mM in mesophyll but 10 mM in sieve tubes (wild-type
values: 50 mM and 800 mM, respectively). Trace this phenotype through
the apoplastic loading pathway and predict the appearance of source
leaves and roots.

\begin{center}\rule{0.5\linewidth}{0.5pt}\end{center}

\subsection{Nutrient Uptake and Mineral
Nutrition}\label{nutrient-uptake-and-mineral-nutrition}

\subsubsection{Essential Nutrients}\label{essential-nutrients}

Plants require 17 essential elements, divided into macronutrients and
micronutrients:

\textbf{Macronutrients} (required in large quantities):

{\def\LTcaptype{none} % do not increment counter
\begin{longtable}[]{@{}
  >{\raggedright\arraybackslash}p{(\linewidth - 6\tabcolsep) * \real{0.1429}}
  >{\raggedright\arraybackslash}p{(\linewidth - 6\tabcolsep) * \real{0.1224}}
  >{\raggedright\arraybackslash}p{(\linewidth - 6\tabcolsep) * \real{0.3469}}
  >{\raggedright\arraybackslash}p{(\linewidth - 6\tabcolsep) * \real{0.3878}}@{}}
\toprule\noalign{}
\begin{minipage}[b]{\linewidth}\raggedright
Element
\end{minipage} & \begin{minipage}[b]{\linewidth}\raggedright
Symbol
\end{minipage} & \begin{minipage}[b]{\linewidth}\raggedright
Primary Functions
\end{minipage} & \begin{minipage}[b]{\linewidth}\raggedright
Deficiency Symptoms
\end{minipage} \\
\midrule\noalign{}
\endhead
\bottomrule\noalign{}
\endlastfoot
Nitrogen & N & Amino acids,
\hyperref[gl:nucleotide]{\textbf{nucleotide}}s,
\hyperref[gl:chlorophyll]{\textbf{chlorophyll}} & Chlorosis (yellowing)
of older leaves first \\
Phosphorus & P & ATP, nucleic acids, membranes & Purple/dark green
leaves; stunted growth \\
Potassium & K & Osmotic regulation,
\hyperref[gl:enzyme]{\textbf{enzyme}} activation, stomata & Marginal
leaf necrosis; weak stems \\
Calcium & Ca & Cell wall (middle lamella), signalling & Growing point
death; blossom-end rot \\
Magnesium & Mg & Chlorophyll centre, enzyme cofactor & Interveinal
chlorosis of older leaves \\
Sulphur & S & Cysteine, methionine, coenzyme A & General chlorosis;
stunted growth \\
\end{longtable}
}

\textbf{Micronutrients} (required in trace quantities):

{\def\LTcaptype{none} % do not increment counter
\begin{longtable}[]{@{}lll@{}}
\toprule\noalign{}
Element & Symbol & Primary Functions \\
\midrule\noalign{}
\endhead
\bottomrule\noalign{}
\endlastfoot
Iron & Fe & Electron transport, chlorophyll synthesis \\
Manganese & Mn & Photosystem II (water-splitting complex) \\
Boron & B & Cell wall cross-linking, pollen tube growth \\
Zinc & Zn & Enzyme cofactor (carbonic anhydrase, SOD) \\
Copper & Cu & Plastocyanin, lignin biosynthesis \\
Molybdenum & Mo & Nitrate reductase, nitrogenase cofactor \\
Chlorine & Cl & Photosystem II, osmotic regulation \\
Nickel & Ni & Urease cofactor \\
\end{longtable}
}

\subsubsection{Nitrogen Fixation
Symbiosis}\label{nitrogen-fixation-symbiosis}

Atmospheric N\(_2\) is inaccessible to plants directly. Some plants form
symbiotic associations with nitrogen-fixing bacteria:

\textbf{Rhizobium-legume symbiosis:}

\begin{enumerate}
\def\labelenumi{\arabic{enumi}.}
\tightlist
\item
  Legume roots secrete flavonoids into the rhizosphere
\item
  Flavonoids activate \emph{nod} genes in \emph{Rhizobium}, producing
  Nod factors (lipochitooligosaccharides)
\item
  Nod factors bind NFR1/NFR5 receptors on root hair cells, triggering
  root hair curling and infection thread formation
\item
  Bacteria enter root cortex via infection thread, released into host
  cells, and differentiate into bacteroids within specialised structures
  called nodules
\item
  Bacteroids express \textbf{nitrogenase} (enzyme complex: Fe-protein +
  MoFe-protein) which catalyses:
\end{enumerate}

\begin{equation}
N_2 + 8H^+ + 8e^- + 16\text{ATP} \rightarrow 2NH_3 + H_2 + 16\text{ADP} + 16P_i
\label{eq:unit_VIII_nitrogenase}
\end{equation}

\begin{enumerate}
\def\labelenumi{\arabic{enumi}.}
\setcounter{enumi}{5}
\tightlist
\item
  \textbf{Leghemoglobin} (produced by host, gives nodules pink colour)
  maintains low O\(_2\) concentration (nitrogenase is irreversibly
  inhibited by O\(_2\)) while supplying O\(_2\) for bacteroid
  respiration
\end{enumerate}

The plant provides photosynthate (sucrose, then metabolised to
malate/succinate) to the bacteroid; the bacteroid provides fixed
nitrogen as NH\(_4^+\) or amino acids (alanine, aspartate).

Other nitrogen-fixing associations: \emph{Frankia} (actinobacterium)
with alder, bayberry; \emph{Anabaena} with water fern \emph{Azolla}
(used as green manure in rice paddies for millennia).

\subsubsection{Mycorrhizal Associations}\label{mycorrhizal-associations}

\textbf{Mycorrhizae} (``fungus-root'') are mutualistic associations
between plant roots and soil fungi. Approximately \textbf{80\% of land
plant species} form mycorrhizal associations. Two major types:

\textbf{Arbuscular mycorrhizae (AM, Glomeromycota):} - Obligate plant
symbionts; cannot grow without host - Fungal hyphae penetrate root
cortical cells and form branched structures called \textbf{arbuscules}
(sites of nutrient exchange, turn over every 4-10 days) - Extraradical
hyphae extend several centimetres into soil, accessing a far greater
soil volume than roots alone - Primary benefit: phosphorus delivery
(hyphae secrete phosphatases and organic acids to solubilise bound P) -
Plant provides up to 20-30\% of photosynthate to fungus - Found in
about 80\% of land plant species; ancestral association dating
to \textgreater400 million years ago

\textbf{Ectomycorrhizae (EM, Basidiomycota + Ascomycota):} - Hyphae form
a dense sheath (mantle) around root tips - \textbf{Hartig net:} Hyphae
penetrate between cortical cells but do not enter cells - Provide
nitrogen (via proteolytic enzymes that break down soil organic matter)
and phosphorus - Dominant in temperate and boreal forest trees
(\emph{Pinus}, \emph{Betula}, \emph{Fagus}, \emph{Quercus}) - Many EM
fungi produce mushrooms (e.g., chanterelles, truffles, \emph{Amanita})

\textbf{Common mycorrhizal networks (CMNs):} Mycorrhizal fungi can
connect multiple plants, forming underground networks. Carbon, water,
and nutrient signals have been documented moving between plants via
these networks, though the ecological significance of active plant-plant
signalling through CMNs remains debated.

\textbf{Concept Check 13:} A farmer notices that legume crops planted
alongside non-legume crops improve the growth of both. Explain the
mechanisms by which nitrogen fixed by rhizobia in legume nodules can
become available to neighbouring non-legume plants.

\begin{center}\rule{0.5\linewidth}{0.5pt}\end{center}

\subsection{Adaptations for Water and Nutrient
Acquisition}\label{adaptations-for-water-and-nutrient-acquisition}

\subsubsection{Xerophyte Adaptations}\label{xerophyte-adaptations}

Plants in arid environments have evolved numerous strategies to minimise
water loss and maximise water uptake:

\begin{itemize}
\tightlist
\item
  \textbf{Thick cuticle and epicuticular waxes} (reduce cuticular
  transpiration)
\item
  \textbf{Sunken stomata} in pits or grooves (create humid
  microenvironment; \emph{Nerium oleander})
\item
  \textbf{Reduced leaf surface area} (needles, spines) or deciduousness
  during dry season
\item
  \textbf{CAM photosynthesis} (stomata open at night, close during day;
  CO\(_2\) stored as malate)
\item
  \textbf{Succulence} (water storage in stems or leaves; \emph{Opuntia},
  \emph{Aloe})
\item
  \textbf{Deep or extensive root systems} (mesquite roots can reach 50+
  m depth)
\end{itemize}

\subsubsection{Hydrophyte Adaptations}\label{hydrophyte-adaptations}

Aquatic and wetland plants face the opposite challenge -- too much water
and often \hyperref[gl:anaerobic]{\textbf{anaerobic}} substrates:

\begin{itemize}
\tightlist
\item
  \textbf{Aerenchyma} (large air-filled cortical spaces for gas
  transport to submerged roots)
\item
  \textbf{Thin or absent cuticle} (CO\(_2\) and O\(_2\) diffuse directly
  from water)
\item
  \textbf{Dissected or thin leaves} (maximise surface area for gas
  exchange in water)
\item
  \textbf{Adventitious roots} from stem nodes
\end{itemize}

\subsubsection{Carnivorous Plants}\label{carnivorous-plants}

In nutrient-poor environments (bogs, sandy soils), some plants
supplement mineral nutrition by capturing and digesting insects and
other small animals:

\begin{itemize}
\tightlist
\item
  \textbf{Pitcher plants} (\emph{Nepenthes}, \emph{Sarracenia}):
  modified leaves form pitfall traps with slippery surfaces and
  digestive enzymes
\item
  \textbf{Sundews} (\emph{Drosera}): sticky mucilage on tentacle-like
  trichomes; leaf slowly curls around prey
\item
  \textbf{Venus flytrap} (\emph{Dionaea muscipula}): snap trap triggered
  by touch-sensitive trigger hairs; requires two stimulations within 20
  seconds (counting mechanism involves Ca\(^{2+}\) signalling)
\end{itemize}

Most carnivorous plants derive primarily nitrogen and phosphorus from
prey.

\textbf{Concept Check 14:} Sketch a pressure-flow circuit (source--sink)
for a CAM plant during the night when stomata are open. How does
CAM-driven nighttime malate accumulation interact with phloem loading
dynamics in the same tissue?

\begin{center}\rule{0.5\linewidth}{0.5pt}\end{center}

\subsection{Computational Bridge}\label{computational-bridge-24}

Whole-cell water potential sums osmotic and pressure components in MPa:

\begin{Shaded}
\begin{Highlighting}[]
\ImportTok{from}\NormalTok{ biology.botany }\ImportTok{import}\NormalTok{ water\_potential}

\NormalTok{psi }\OperatorTok{=}\NormalTok{ water\_potential(}\FloatTok{0.3}\NormalTok{, turgor\_pressure\_MPa}\OperatorTok{=}\FloatTok{0.5}\NormalTok{)}
\BuiltInTok{print}\NormalTok{(}\BuiltInTok{round}\NormalTok{(psi.water\_potential\_MPa, }\DecValTok{4}\NormalTok{))}
\end{Highlighting}
\end{Shaded}

\begin{quote}
\textbf{Clinical / systems note:} Crop drought breeding screens often
proxy cellular osmotic adjustment --- the same \(\Psi_s\) shifts you
estimate from solute accumulation.
\end{quote}

\begin{center}\rule{0.5\linewidth}{0.5pt}\end{center}

\subsection{Current Evidence and Frontier
Biology}\label{current-evidence-and-frontier-biology-27}

For \textbf{Plant Structure, Water Relations, and Transport}, frontier
biology belongs inside the evidence logic of the chapter. Plant biology
links molecular regulation to climate stress, water limitation, crop
resilience, phenology, and ecosystem feedbacks. The core reading
question is this: plant-water claims require water potential, hydraulic
pathway, stomatal control, tissue anatomy, and stress context.

\begin{itemize}
\tightlist
\item
  \textbf{What to verify:} identify the observation, model, assay, or
  dataset that would make the claim stronger or weaker.
\item
  \textbf{What to qualify:} state the scale, organism, cell type,
  environmental condition, or population where the claim is expected to
  hold.
\item
  \textbf{What to compare:} test at least one alternative explanation,
  baseline, or null model before treating the pattern as causal.
\item
  \textbf{What to cite:} distinguish primary evidence, review synthesis,
  public dataset, and institutional guidance; for recent or numeric
  claims, prefer the source closest to the measurement and state what
  has changed since it was published.
\end{itemize}

A strong plant explanation names the tissue, signal, environmental
driver, measurable trait, and tradeoff between growth, reproduction,
defence, and water use.

\textbf{Source practice:} For plant-stress and crop claims, name the
tissue, environmental driver, field context, and growth-reproduction
tradeoff; separate laboratory potential from agronomic adoption.

\subsubsection{Current Evidence Map: Hydraulic Safety
Tradeoff}\label{current-evidence-map-hydraulic-safety-tradeoff}

\begin{figure}[htbp]
\centering
\includegraphics[width=0.82\linewidth,height=4.2in,keepaspectratio]{../figures/mermaid_inline/inline_mermaid_0142_2ce273db44e8.png}
\caption{Drought responses should be read as tradeoffs among carbon gain, water loss, hydraulic safety, and growth rather than as simple stress resistance.}
\end{figure}

\subsubsection{Field Translation: Hydraulic Safety
Margins}\label{field-translation-hydraulic-safety-margins}

Water-potential models move from classroom equations to field prediction
when they name the operating range. A hydraulic safety margin compares
the water potential a plant actually reaches during drought with the
water potential that causes damaging loss of xylem conductivity. Choat
and colleagues found that 70\% of 226 sampled forest species from 81
sites operated with margins narrower than 1 MPa against injurious
hydraulic stress \citep{choat2012hydraulicmargins}. That result is a
warning signal, not a one-number forecast: mortality risk also depends
on rooting depth, stored water, non-structural carbohydrate reserves,
stomatal regulation, recovery after embolism, heat load, pests, and
repeated drought history.

When evaluating a drought-tolerant crop or forest-management claim, ask
for the species, tissue, growth stage, vapor-pressure deficit, soil
depth, measured \(\Psi_{\mathrm{min}}\), embolism threshold, and yield
or survival endpoint. A greenhouse water-potential treatment can
demonstrate mechanism, but field adoption requires replicated seasons
and soils because stomatal closure that protects xylem can reduce carbon
gain and reproduction.

\subsection{Summary}\label{summary-27}

\begin{itemize}
\tightlist
\item
  \textbf{Plant body plan:} Root system (taproot or fibrous) + shoot
  system (stem + leaves). Phytomers (node + internode + leaf + axillary
  bud) are the modular building units; plastochron index sets
  developmental time.
\item
  \textbf{Three tissue systems:} Dermal (protection, cuticle, stomata,
  trichomes), ground (parenchyma for photosynthesis/storage, collenchyma
  for flexible support, sclerenchyma for rigid support), vascular
  (xylem: tracheids/vessels for water; phloem: sieve tubes/companion
  cells for sugars).
\item
  \textbf{Xylem development:} Protoxylem (annular/helical, extensible)
  precedes metaxylem (scalariform/pitted, mechanically robust). Vessels
  are wider and more efficient than tracheids but more cavitation-prone
  (Hagen-Poiseuille r\texorpdfstring{\textsuperscript{4}}{4} scaling). Conifer pits with torus-margo valves
  provide most-or-nothing cavitation isolation; angiosperm homogeneous
  pit membranes are more vulnerable.
\item
  \textbf{Water potential:} Full four-component form
  \(\Psi = \Psi_s + \Psi_p + \Psi_m + \Psi_g\)
  (\cref{eq:unit_VIII_water_potential}) governs direction of water flow;
  matric potential dominates in dry seeds and soils. The SPAC has a
  steep gradient driving water from soil (\(-0.01\) to \(-1.5\) MPa) to
  atmosphere (\(-100\) MPa).
\item
  \textbf{Transport pathways:} Apoplast (cell walls, fast,
  non-selective) vs symplast (cytoplasm via plasmodesmata, selective) vs
  transmembrane (aquaporins + transporters, selective). Casparian strip
  (lignin) + suberin lamella in endodermis force water into symplast for
  quality control; passage cells provide selective bypasses; lateral
  root emergence creates transient apoplastic bypass.
\item
  \textbf{TACT mechanism \citep{dixon1894}:} Transpiration creates
  tension in xylem; water's cohesive H-bonds maintain continuous column;
  adhesion to xylem walls supports the column; water rises passively.
  Critical evaluation: pressure probes, MRI, synthetic-tree experiments,
  and synchrotron X-ray micro-CT confirm tensions of \(-1\) to \(-22\)
  MPa; embolism repair via overnight root pressure and phloem-driven
  refilling mitigates cavitation losses.
\item
  \textbf{Stomatal regulation:} Opening via H\(^+\)-ATPase, KAT1, malate
  accumulation (blue light, phototropins); closure via
  ABA-PYR-PP2C-SnRK2-SLAC1/GORK cascade. Balances CO\(_2\) gain vs water
  loss; iWUE = \(A_n/g_s\) (\cref{eq:unit_VIII_iwue}). C4 and CAM plants
  achieve iWUE 2--10\texorpdfstring{\ensuremath{\times}}{x} C3 by concentrating CO\(_2\) behind
  partially-closed stomata.
\item
  \textbf{Phloem:} Münch pressure-flow \citep{munch1930} with apoplastic
  loading (SUT1 cotransporter; herbaceous crops; high sieve-tube
  concentration) or symplastic loading via polymer trap (raffinose; many
  trees; many plasmodesmata; lower concentration) at source.
  Bidirectional flow follows turgor gradient; symplastic vs apoplastic
  unloading varies by sink type.
\item
  \textbf{Nutrient uptake:} 17 essential elements (6 macro, 8 micro,
  plus C/H/O). \hyperref[gl:nitrogen-fixation]{\textbf{Nitrogen
  fixation}} by rhizobia in legume nodules. Mycorrhizae (AM and EM)
  enhance P and N uptake in 80\% of plant species.
\item
  \textbf{Connections:} See \cref{sec:unit_VIII_plant_reproduction} and
  \cref{sec:unit_VIII_plant_responses} for reproduction and signalling,
  \cref{sec:unit_I_water_and_life} for water
  \hyperref[gl:thermodynamics]{\textbf{thermodynamics}}, and
  \cref{sec:unit_X_biomes_and_conservation} for primary production.
\end{itemize}

\begin{center}\rule{0.5\linewidth}{0.5pt}\end{center}

\subsection{Key Terms}\label{key-terms-27}

{\def\LTcaptype{none} % do not increment counter
\begin{longtable}[]{@{}
  >{\raggedright\arraybackslash}p{(\linewidth - 2\tabcolsep) * \real{0.2857}}
  >{\raggedright\arraybackslash}p{(\linewidth - 2\tabcolsep) * \real{0.7143}}@{}}
\toprule\noalign{}
\begin{minipage}[b]{\linewidth}\raggedright
Term
\end{minipage} & \begin{minipage}[b]{\linewidth}\raggedright
Definition
\end{minipage} \\
\midrule\noalign{}
\endhead
\bottomrule\noalign{}
\endlastfoot
\textbf{Water potential (\texorpdfstring{\ensuremath{\Psi}}{Psi})} & Free energy of water per unit volume;
determines direction of water movement; units MPa \\
\textbf{Osmotic potential (\(\Psi_s\))} & Component of water potential
due to solute concentration; typically negative \\
\textbf{Turgor pressure (\(\Psi_p\))} & Positive pressure exerted by
cell contents against the cell wall \\
\textbf{Matric potential (\(\Psi_m\))} & Component due to water binding
on colloidal surfaces; dominant in dry seeds and soils \\
\textbf{Casparian strip} & Lignin (then suberin) band in endodermal cell
walls; forces water through symplast \\
\textbf{Suberin lamella} & Continuous suberin layer in mature endodermal
cells; further restricts apoplast \\
\textbf{Passage cell} & Endodermal cell lacking suberin lamella;
provides selective transport bypass \\
\textbf{Cohesion-tension mechanism} & Transpiration-driven water ascent
relying on hydrogen bond cohesion \\
\textbf{Cavitation} & Formation of gas bubble in xylem under excessive
tension \\
\textbf{Stomata} & Pores bounded by guard cells; regulate gas exchange
and transpiration \\
\textbf{Apoplast} & Continuum of cell walls and intercellular spaces \\
\textbf{Symplast} & Continuum of cytoplasm connected by plasmodesmata \\
\textbf{Tracheids} & Elongated, tapered xylem cells with bordered pits;
found in most vascular plants \\
\textbf{Vessel elements} & Wide xylem cells with perforation plates;
primarily in angiosperms \\
\textbf{Bordered pit (torus-margo)} & Conifer pit with thickened central
torus suspended in porous margo; valves on cavitation \\
\textbf{Protoxylem} & Early xylem with annular/helical thickenings;
extensible during organ growth \\
\textbf{Metaxylem} & Late xylem with scalariform/pitted thickenings;
mechanically robust \\
\textbf{Sieve tube elements} & Enucleate phloem cells connected by sieve
plates; conduct sugars \\
\textbf{Companion cells} & Nucleate cells that metabolically support
sieve tube elements \\
\textbf{SUT1 (SUC2)} & Sucrose-H\(^+\) symporter on companion cells;
rate-limiting loader in apoplastic species \\
\textbf{SWEET11/12} & Sucrose efflux carriers; export sucrose from
mesophyll to apoplast \\
\textbf{Polymer trap} & Symplastic loading mechanism using
raffinose-family oligosaccharides \\
\textbf{Plasmodesmatal conductance} & Bulk transfer rate through
plasmodesmata; regulated by callose deposition \\
\textbf{Münch pressure-flow} & Model of phloem transport driven by
turgor pressure gradient from source to sink \\
\textbf{Phytomer} & Modular shoot unit: node + internode + leaf +
axillary bud \\
\textbf{Plastochron index} & Time interval between successive
leaf-primordium initiations \\
\textbf{Mycorrhizae} & Mutualistic root-fungus associations; enhance
nutrient uptake \\
\textbf{Arbuscular mycorrhizae} & AM fungi that form arbuscules inside
root cortical cells; deliver P \\
\textbf{Ectomycorrhizae} & EM fungi that form external mantle and Hartig
net; deliver N and P \\
\textbf{Nitrogenase} & Enzyme complex in nitrogen-fixing bacteria;
converts N\(_2\) to NH\(_3\) \\
\textbf{Aquaporins} & Water channel proteins (PIP, TIP families) in
plant cell membranes \\
\textbf{OST1/SnRK2.6} & Master kinase of guard-cell ABA response;
phosphorylates SLAC1 \\
\textbf{iWUE (intrinsic WUE)} & \(A_n/g_s\); net CO\(_2\) assimilated
per unit stomatal conductance \\
\end{longtable}
}

\begin{center}\rule{0.5\linewidth}{0.5pt}\end{center}

\subsection{Review Questions}\label{review-questions-27}

\begin{enumerate}
\def\labelenumi{\arabic{enumi}.}
\item
  A plant cell has a solute concentration of 0.4 M, turgor pressure of
  0.3 MPa, and is at ground level. Using \(\Psi_s = -iCRT\) with
  \(i = 1\), \(R = 0.00831\) L MPa mol\(^{-1}\) K\(^{-1}\), and
  \(T = 298\) K, calculate the total water potential. If this cell is
  adjacent to another cell with \(\Psi = -0.5\) MPa, which direction
  will water flow?
\item
  Explain why xylem vessels are more efficient than tracheids for water
  transport but also more vulnerable to cavitation. What structural
  features of tracheids (torus-margo pits) help prevent the spread of
  embolisms? Contrast with angiosperm pit membranes.
\item
  Describe the complete pathway of a water molecule from the soil
  solution to the atmosphere, naming each structure it passes through.
  At which points must the molecule cross a cell membrane?
\item
  Compare apoplastic and symplastic phloem loading. Which is more common
  in herbaceous crop plants? What is the energetic cost of each, and why
  might symplastic loading be advantageous in some environments? Cite
  SUT1 and the polymer-trap mechanism in your answer.
\item
  A farmer observes that phosphorus-deficient plants show improved
  growth when inoculated with arbuscular mycorrhizal fungi but not when
  given additional phosphorus fertiliser alone. Explain this observation
  in terms of the spatial distribution of phosphorus in soil and the
  relative surface area of roots vs mycorrhizal hyphae.
\item
  Explain how the Casparian strip (lignin) and suberin lamella of the
  endodermis act as a quality control checkpoint. What would happen to
  plant ion \hyperref[gl:homeostasis]{\textbf{homeostasis}} if the
  Casparian strip were absent? What is the role of passage cells in
  older roots? Account for apoplastic bypasses at lateral root emergence
  sites.
\item
  During a heat wave, a crop plant simultaneously experiences high
  temperatures (stomatal opening signal from photosynthesis demand) and
  soil drying (ABA-mediated stomatal closure signal). Trace the
  molecular competition between these signals at the level of OST1/SLAC1
  in guard cells and predict which typically dominates.
\item
  Compare and contrast water movement in xylem and phloem in terms of:
  driving force, direction, conducting cell viability, and substances
  transported.
\item
  Describe the nitrogen fixation symbiosis between \emph{Rhizobium} and
  legumes. Why must nitrogenase be protected from oxygen, and how does
  the nodule accomplish this?
\item
  A coast redwood (\emph{Sequoia sempervirens}) is 100 m tall. Calculate
  the gravitational component of water potential at the top of the tree.
  Given that atmospheric \texorpdfstring{\ensuremath{\Psi}}{Psi} can reach \(-100\) MPa at 50\% humidity,
  explain why the tree can still pull water to the canopy.
\item
  Critically evaluate the cohesion-tension theory \citep{dixon1894}.
  Identify two lines of evidence supporting it (including synchrotron
  imaging) and one historical objection that has been resolved. What is
  the air-seeding hypothesis, and how does it apply to torus-margo pits?
\item
  Using \breaktt{transpiration\_flux}, predict whether doubling stomatal
  conductance doubles water loss if the leaf-to-air vapour gradient is
  held fixed. Then compute the iWUE consequence assuming \(A_n\)
  saturates at 80\% of the doubled \(g_s\).
\item
  Contrast \textbf{hydraulic redistribution} via roots with capillary
  rise in soil --- when does each dominate overnight water status?
\item
  Use the Hagen-Poiseuille equation \cref{eq:unit_VIII_hagen_poiseuille}
  to estimate the relative volumetric flow capacity of a 200-\texorpdfstring{\ensuremath{\mu}}{mu}m vessel
  versus a 30-\texorpdfstring{\ensuremath{\mu}}{mu}m tracheid (same length and pressure gradient). How does
  this calculation explain the prevalence of vessels in fast-transpiring
  angiosperms?
\item
  A C3 wheat plant achieves iWUE of 80 \texorpdfstring{\ensuremath{\mu}}{mu}mol CO\(_2\) mol\(^{-1}\)
  H\(_2\)O; an adjacent C4 maize plant achieves iWUE of 200. Both have
  nearly identical \(A_n\) when measured at noon. Calculate the relative
  \(g_s\) of the two species and explain why C4 plants outcompete C3 in
  semi-arid regions.
\end{enumerate}

\begin{center}\rule{0.5\linewidth}{0.5pt}\end{center}

\subsection{Further Reading and Source
Notes}\label{further-reading-and-source-notes-22}

\begin{itemize}
\tightlist
\item
  \citet{dixon1894} --- On the ascent of sap. \emph{Philosophical
  Transactions of the Royal Society B}, 186.
\item
  \citet{munch1930} --- Die Stoffbewegungen in der Pflanze. \emph{Gustav
  Fischer Verlag}.
\end{itemize}

\begin{center}\rule{0.5\linewidth}{0.5pt}\end{center}

\subsubsection{Companion Source
Module}\label{companion-source-module-27}

\textbf{Plant Structure, Water Relations, and Transport} should leave a
reproducible trail from a biological claim to the code, figure, diagram,
or paper-based activity that can test it. Use the surfaces below to
inspect the chapter's assumptions, rerun the relevant model, or compare
the manuscript explanation with companion labs and figures.

{\def\LTcaptype{none} % do not increment counter
\begin{longtable}[]{@{}
  >{\raggedright\arraybackslash}p{(\linewidth - 2\tabcolsep) * \real{0.5000}}
  >{\raggedright\arraybackslash}p{(\linewidth - 2\tabcolsep) * \real{0.5000}}@{}}
\toprule\noalign{}
\begin{minipage}[b]{\linewidth}\raggedright
Surface
\end{minipage} & \begin{minipage}[b]{\linewidth}\raggedright
Use it for
\end{minipage} \\
\midrule\noalign{}
\endhead
\bottomrule\noalign{}
\endlastfoot
\breaktt{src/biology/botany/botany.py} (\texttt{water\_potential},
\breaktt{transpiration\_flux}) & Reproduce plant-water calculations and
hydraulic tradeoffs. \\
\breaktt{src/biology/cell/cell\_biology.py} (\texttt{osmotic\_pressure})
& Connect cellular osmotic pressure to tissue-level water movement. \\
\breaktt{src/visualization/plots.py}
(\breaktt{plot\_light\_response\_curve}) & Practice graph interpretation
for environmental-response curves. \\
\end{longtable}
}

\textbf{Reproducibility check:} list solute potential, pressure
potential, tissue, humidity, temperature, and stomatal state before
predicting water movement. \textbf{Cross-reference:} use
\cref{sec:unit_I_water_and_life} and
\cref{sec:unit_VIII_plant_responses}.

\newpage

\section{Plant Reproduction and
Development}\label{plant-reproduction-and-development}

\label{sec:unit_VIII_plant_reproduction}

\begin{quote}
\textbf{Ch 26} \texorpdfstring{\ensuremath{\cdot}}{.} Level 2/3 \texorpdfstring{\ensuremath{\cdot}}{.} 55 min read \texorpdfstring{\ensuremath{\cdot}}{.} 75 min lecture \texorpdfstring{\ensuremath{\cdot}}{.}
Prerequisites: \cref{sec:unit_VIII_plant_structure_and_water}
\end{quote}

\subsection{Learning Objectives}\label{learning-objectives-28}

By the end of this chapter, you should be able to:

\begin{enumerate}
\def\labelenumi{\arabic{enumi}.}
\tightlist
\item
  Describe the alternation of generations and explain the evolutionary
  trend from gametophyte-\hyperref[gl:dominant]{\textbf{dominant}} to
  sporophyte-dominant life cycles across the plant kingdom.
\item
  Compare reproduction in non-vascular plants, seedless vascular plants,
  gymnosperms, and \hyperref[gl:angiosperm]{\textbf{angiosperm}}s.
\item
  Describe angiosperm flower structure and explain the ABCDE model of
  floral organ identity \citep{coen1991}, including the molecular
  identities of MADS-box transcription factors.
\item
  Describe pollen grain and embryo sac architecture at cellular
  resolution; trace double fertilisation in detail
  \citep{nawaschin1898}, including pollen tube chemotropism, synergid
  degeneration, and sperm cell delivery.
\item
  Explain microsporogenesis, megasporogenesis, and endosperm development
  (free-nuclear, cellular, PEG-pathway) including the 2m:1p genome
  dosage.
\item
  Compare seed dormancy types (physical, physiological, morphological,
  morphophysiological, combinatorial) and the molecular basis of
  germination triggers (stratification, scarification, light via
  phytochrome).
\item
  Describe self-incompatibility (SSI vs GSI), apomixis (sporophytic vs
  gametophytic), polyploidy, and vegetative propagation strategies.
\item
  Explain embryogenesis and meristem organisation (SAM/RAM) including
  the WUS-CLV3 feedback loop, fruit development, and parthenocarpy.
\end{enumerate}

\subsubsection{Study Blueprint}\label{study-blueprint-28}

\begin{itemize}
\tightlist
\item
  \textbf{Big idea:} Plant reproduction integrates life cycles,
  development, dispersal, and environmental timing.
\item
  \textbf{Core concepts:} alternation of generations, flowers, seeds,
  development.
\item
  \textbf{Framework alignment:} Vision \& Change: Structure and
  function, Pathways and transformations of energy and matter, Systems;
  AP Biology: Energetics, Systems Interactions; NGSS-style topics:
  Structure and Function, Matter and Energy in Organisms and Ecosystems.
\item
  \textbf{Model or quantitative lens:} Life-cycle accounting and
  phyllotaxis/growth-pattern calculations.
\item
  \textbf{Data skill:} Track ploidy, tissue origin, and reproductive
  stage from diagrams or observations.
\item
  \textbf{Practice cadence:} Visual Representations, Questions and
  Methods, Argumentation.
\item
  \textbf{Common misconception to repair:} Pollen, spores, seeds, and
  gametes are not interchangeable terms.
\item
  \textbf{Primary lab:} \cref{sec:lab_unit_VIII_plant_reproduction}.
\item
  \textbf{Question bank:} \cref{sec:q_unit_VIII_plant_reproduction}.
\item
  \textbf{Transfer task:} Transfer reproductive reasoning to
  agriculture, pollination ecology, and plant evolution.
\item
  \textbf{Bridge to computation:}
  \breaktt{biology.botany.botany.plant\_biomass\_growth}.
\end{itemize}

\begin{center}\rule{0.5\linewidth}{0.5pt}\end{center}

\begin{quote}
\textbf{Opening Vignette --- The Seeds That Fed a Billion People}

In the 1950s, famines threatened South Asia and Latin America as human
population growth outpaced food production. Norman Borlaug, an American
agronomist working in Mexico with the Rockefeller Foundation, spent a
decade cross-breeding wheat varieties --- shuttling seeds between winter
and summer growing seasons to double selection speed --- until he
produced semi-dwarf, disease-resistant wheat varieties that yielded
three to four times as much grain as traditional tall-stem wheats.
Unlike traditional wheats, the short stems didn't fall over when
supporting heavy grain heads. He introduced these seeds to Pakistan and
India in 1965. By 1970, wheat production had doubled in both countries,
and famine had been averted. Borlaug received the 1970 Nobel Peace
Prize. The same innovations, later extended to rice by the International
Rice Research Institute, constitute the Green Revolution --- a triumph
of applied plant reproductive biology that is estimated to have saved
over a billion lives.
\end{quote}

\subsection{Alternation of
Generations}\label{alternation-of-generations}

Plants have a \textbf{diphasic life cycle} -- alternating between a
diploid \textbf{sporophyte} (spore-producing) generation and a haploid
\textbf{gametophyte} (\hyperref[gl:gamete]{\textbf{gamete}}-producing)
generation. This alternation is unique to plants and some algae; animals
have primarily a brief haploid gametic phase.

\begin{figure}[htbp]
\centering
\includegraphics[width=0.82\linewidth,height=4.2in,keepaspectratio]{../figures/mermaid_inline/inline_mermaid_0143_b2c05fc6accc.png}
\caption{Alternation of generations showing the two phases of the plant life cycle The sporophyte (2n) produces haploid spores by meiosis. Spores develop into the gametophyte (n), which produces gametes by mitosis. Fertilisation restores the diploid sporophyte.}
\end{figure}

{\def\LTcaptype{none} % do not increment counter
\begin{longtable}[]{@{}
  >{\raggedright\arraybackslash}p{(\linewidth - 6\tabcolsep) * \real{0.1410}}
  >{\raggedright\arraybackslash}p{(\linewidth - 6\tabcolsep) * \real{0.2436}}
  >{\raggedright\arraybackslash}p{(\linewidth - 6\tabcolsep) * \real{0.2949}}
  >{\raggedright\arraybackslash}p{(\linewidth - 6\tabcolsep) * \real{0.3205}}@{}}
\toprule\noalign{}
\begin{minipage}[b]{\linewidth}\raggedright
Plant group
\end{minipage} & \begin{minipage}[b]{\linewidth}\raggedright
Dominant generation
\end{minipage} & \begin{minipage}[b]{\linewidth}\raggedright
Gametophyte description
\end{minipage} & \begin{minipage}[b]{\linewidth}\raggedright
Fertilisation requirement
\end{minipage} \\
\midrule\noalign{}
\endhead
\bottomrule\noalign{}
\endlastfoot
Bryophytes (mosses, liverworts) & \textbf{Haploid gametophyte} (the
leafy ``plant'') & Free-living; nutritionally independent & Water
required (flagellated sperm swim) \\
Pteridophytes (ferns, horsetails) & \textbf{Diploid sporophyte} & Small,
free-living prothallus (about 1 cm\(^2\)); needs moisture &
Water required \\
Gymnosperms (\emph{Pinus}, \emph{Ginkgo}, cycads) & Diploid sporophyte &
Reduced; pollen = 4-celled male gametophyte; archegonia in ovule =
female gametophyte & Wind pollination; no water needed \\
Angiosperms (flowering plants) & Diploid sporophyte & Microscopic,
dependent: pollen = 3-celled; embryo sac = 7-celled, 8-nuclei & Wind or
animal pollination \\
\end{longtable}
}

\textbf{Why does dominance shift to the sporophyte?} Diploid cells can
mask deleterious recessive \hyperref[gl:mutation]{\textbf{mutation}}s
(\hyperref[gl:heterozygous]{\textbf{heterozygous}} advantage). A diploid
sporophyte can accumulate more genetic material enabling structural
complexity -- leaves, stems, vascular tissue, seeds -- allowing
colonisation of increasingly dry terrestrial habitats. The seed is the
key innovation: it encases the embryo and gametophyte within maternal
tissue, eliminating dependence on external water for fertilisation.

\subsubsection{Worked Example: Tracking Ploidy Through the Angiosperm
Life
Cycle}\label{worked-example-tracking-ploidy-through-the-angiosperm-life-cycle}

\textbf{Problem:} Maize (\emph{Zea mays}) has a sporophyte chromosome
number of \(2n = 20\). Track the chromosome number at each stage of
alternation of generations, then determine the ploidy of the zygote and
the primary endosperm nucleus produced by double fertilisation.

\textbf{Solution:}

\begin{enumerate}
\def\labelenumi{\arabic{enumi}.}
\item
  \textbf{Establish the haploid number.} The gametophyte (haploid)
  number is half the sporophyte number:
  \[ n = \frac{2n}{2} = \frac{20}{2} = 10 \text{ chromosomes} \label{eq:unit_VIII_plant_reproduction_item_7}\]
\item
  \textbf{Meiosis in the sporophyte.} The diploid megaspore mother cell
  (\(2n = 20\)) undergoes meiosis, producing haploid megaspores with
  \(n = 10\). The surviving megaspore divides \textbf{mitotically} to
  build the embryo sac, so the egg cell and each of the two polar nuclei
  most carry \(n = 10\). The pollen (microgametophyte) likewise delivers
  sperm cells with \(n = 10\).
\item
  \textbf{First fusion --- the zygote.} One sperm fuses with the egg:
  \[ \text{zygote} = n_{\text{egg}} + n_{\text{sperm}} = 10 + 10 = 20 \;\; (2n,\ \text{diploid embryo}) \label{eq:unit_VIII_plant_reproduction_item_8}\]
\item
  \textbf{Second fusion --- the endosperm.} The second sperm fuses with
  the two polar nuclei of the central cell:
  \[ \text{primary endosperm nucleus} = n + n + n = 10 + 10 + 10 = 30 \;\; (3n,\ \text{triploid}) \label{eq:unit_VIII_plant_reproduction_item_9}\]
\item
  \textbf{Genome dosage of the endosperm.} Two of the three genomes are
  maternal (the polar nuclei) and one is paternal (the sperm):
  \[ \frac{\text{maternal}}{\text{paternal}} = \frac{2n}{1n} = \frac{20}{10} = \frac{2}{1} \label{eq:unit_VIII_plant_reproduction_item_10}\]
\end{enumerate}

\textbf{Interpretation:} Meiosis halves \(2n = 20\) to \(n = 10\) for
the gametophyte generation; the two sperm-mediated fusions of double
fertilisation then restore a \(2n = 20\) embryo and create a distinctive
\(3n = 30\) endosperm with a fixed 2:1 maternal:paternal dosage --- the
genetic substrate for the parent-of-origin conflict conserved across
flowering plants.

\begin{center}\rule{0.5\linewidth}{0.5pt}\end{center}

\subsection{Plant Group Reproductive
Strategies}\label{plant-group-reproductive-strategies}

\subsubsection{Non-Vascular Plants
(Bryophytes)}\label{non-vascular-plants-bryophytes}

\textbf{Mosses} (Bryophyta, about 12,000 species),
\textbf{liverworts} (Marchantiophyta), and \textbf{hornworts}
(Anthocerotophyta) represent the ancestral condition:

\begin{itemize}
\tightlist
\item
  The \textbf{gametophyte} is the dominant, photosynthetic plant body
\item
  The sporophyte is small, attached to, and nutritionally dependent on
  the gametophyte
\item
  Flagellated sperm must swim through water films to reach the
  archegonium (female reproductive structure)
\item
  Spores dispersed from capsule (sporangium) at tip of sporophyte seta
\item
  Ecological importance: pioneer colonisers of bare rock, form peat
  (carbon reservoir), retain moisture
\end{itemize}

\textbf{Concept Check 1:} Why are mosses restricted to moist habitats?
What aspect of their reproduction limits their distribution?

\subsubsection{Seedless Vascular Plants (Ferns and
Allies)}\label{seedless-vascular-plants-ferns-and-allies}

\textbf{Ferns} (Polypodiopsida, about 10,500 species) represent
the first major shift to sporophyte dominance:

\begin{itemize}
\tightlist
\item
  The \textbf{sporophyte} is the dominant plant (the familiar fern
  \hyperref[gl:frond]{\textbf{frond}})
\item
  Sporangia are clustered in \textbf{sori} (singular: sorus) on the
  underside of fronds, often protected by an indusium
\item
  Meiosis in sporangia produces haploid spores
\item
  Spores germinate into a small, heart-shaped \textbf{prothallus} (the
  independent gametophyte, about 1 cm)
\item
  The prothallus bears both antheridia (produce flagellated sperm) and
  archegonia (produce eggs)
\item
  Sperm must swim through water to the archegonium -- the limiting step
\item
  After fertilisation, the new sporophyte grows from the prothallus,
  which eventually withers
\end{itemize}

Horsetails (\emph{Equisetum}) and club mosses (\emph{Lycopodium},
\emph{Selaginella}) are also seedless vascular plants.
\emph{Selaginella} shows early heterospory (different-sized spores for
male and female gametophytes), a precursor to the seed habit.

\subsubsection{Gymnosperms}\label{gymnosperms}

Gymnosperms (``naked seed'') include conifers (about 630
species), cycads (about 350), \emph{Ginkgo biloba} (1 species),
and gnetophytes (about 70). Key reproductive features:

\begin{itemize}
\tightlist
\item
  \textbf{Heterospory:} Microsporangia (in pollen cones/microstrobili)
  produce microspores that develop into pollen grains (male
  gametophyte). Megasporangia (in ovulate cones) produce megaspores that
  develop into the female gametophyte (with archegonia).
\item
  \textbf{Pollen eliminates the need for water} for fertilisation.
  Wind-dispersed pollen lands on the ovule (pollination drop mechanism
  in many gymnosperms).
\item
  \textbf{Slow fertilisation:} In pines, 12-18 months elapse between
  pollination and fertilisation. The pollen tube grows slowly through
  the nucellus to reach the archegonium.
\item
  \textbf{Seeds:} The fertilised ovule develops into a seed containing
  the embryo, stored food (female gametophyte tissue), and a protective
  seed coat derived from the integuments. Seeds are ``naked'' (not
  enclosed in a fruit).
\end{itemize}

\subsubsection{Angiosperms (Flowering
Plants)}\label{angiosperms-flowering-plants}

Angiosperms (about 300,000 species) are the most diverse plant
group, with several key innovations:

\begin{itemize}
\tightlist
\item
  \textbf{Flowers:} Specialised reproductive structures that facilitate
  pollination (including animal pollination)
\item
  \textbf{Double fertilisation:} Unique to angiosperms, producing both
  embryo and endosperm
\item
  \textbf{Fruits:} Mature ovary wall enclosing seeds; facilitate
  dispersal
\item
  \textbf{Reduced gametophytes:} Male gametophyte = 3-celled pollen
  grain; female gametophyte = 7-celled, 8-nucleate embryo sac
\end{itemize}

\begin{center}\rule{0.5\linewidth}{0.5pt}\end{center}

\subsection{Angiosperm Flower Structure and the ABCDE
Model}\label{angiosperm-flower-structure-and-the-abcde-model}

\subsubsection{Flower Architecture}\label{flower-architecture}

Angiosperm flowers are \textbf{modified shoots} (shoot lateral organs
with determinate growth). A \textbf{complete flower} contains four
whorls (from outside inward):

{\def\LTcaptype{none} % do not increment counter
\begin{longtable}[]{@{}
  >{\raggedright\arraybackslash}p{(\linewidth - 6\tabcolsep) * \real{0.2273}}
  >{\raggedright\arraybackslash}p{(\linewidth - 6\tabcolsep) * \real{0.1818}}
  >{\raggedright\arraybackslash}p{(\linewidth - 6\tabcolsep) * \real{0.2273}}
  >{\raggedright\arraybackslash}p{(\linewidth - 6\tabcolsep) * \real{0.3636}}@{}}
\toprule\noalign{}
\begin{minipage}[b]{\linewidth}\raggedright
Whorl
\end{minipage} & \begin{minipage}[b]{\linewidth}\raggedright
Name
\end{minipage} & \begin{minipage}[b]{\linewidth}\raggedright
Organ
\end{minipage} & \begin{minipage}[b]{\linewidth}\raggedright
Function
\end{minipage} \\
\midrule\noalign{}
\endhead
\bottomrule\noalign{}
\endlastfoot
1 (outermost) & Calyx & Sepals & Protect flower bud; photosynthetic;
sometimes petaloid \\
2 & Corolla & Petals & Attract pollinators via colour (UV patterns
visible to bees), scent (monoterpenes), oil rewards \\
3 & Androecium & Stamens (filament + anther) & Pollen production (male
gametophyte) \\
4 (innermost) & Gynoecium & Carpels (pistil) = stigma + style + ovary &
Pollen reception; pollen tube guidance; ovule enclosure \\
\end{longtable}
}

\textbf{Floral symmetry:} - \textbf{Actinomorphic (radially symmetric):}
\emph{Rosa}, \emph{Ranunculus}, most basal angiosperms -- generalised
pollinators - \textbf{Zygomorphic (bilaterally symmetric):} Orchidaceae,
Fabaceae, Scrophulariaceae -- specialised pollinators

\textbf{Incomplete flowers} lack one or more whorls. \textbf{Imperfect
flowers} lack either stamens (pistillate/female flowers) or carpels
(staminate/male flowers). Monoecious plants bear both on the same
individual (maize); dioecious plants bear them on separate individuals
(holly, willow).

\subsubsection{\texorpdfstring{The ABCDE Model and MADS-Box Floral
Quartets
\citep{coen1991}}{The ABCDE Model and MADS-Box Floral Quartets}}\label{the-abcde-model-and-mads-box-floral-quartets}

The combinatorial control of floral organ identity by \textbf{MADS-box
transcription factors} is the foundational discovery of plant evo-devo,
and it provides a quantitative framework for understanding floral
diversity. The classical ABC model (Coen and Meyerowitz, 1991) was
extended by the discovery of D and E classes:

{\def\LTcaptype{none} % do not increment counter
\begin{longtable}[]{@{}
  >{\raggedright\arraybackslash}p{(\linewidth - 6\tabcolsep) * \real{0.1111}}
  >{\raggedright\arraybackslash}p{(\linewidth - 6\tabcolsep) * \real{0.1111}}
  >{\raggedright\arraybackslash}p{(\linewidth - 6\tabcolsep) * \real{0.3111}}
  >{\raggedright\arraybackslash}p{(\linewidth - 6\tabcolsep) * \real{0.4667}}@{}}
\toprule\noalign{}
\begin{minipage}[b]{\linewidth}\raggedright
Whorl
\end{minipage} & \begin{minipage}[b]{\linewidth}\raggedright
Organ
\end{minipage} & \begin{minipage}[b]{\linewidth}\raggedright
Classes active
\end{minipage} & \begin{minipage}[b]{\linewidth}\raggedright
\hyperref[gl:gene]{\textbf{Gene}}s (\emph{Arabidopsis})
\end{minipage} \\
\midrule\noalign{}
\endhead
\bottomrule\noalign{}
\endlastfoot
1 & Sepal & A + E & AP1, AP2, SEP1-4 \\
2 & Petal & A + B + E & AP1, AP2, AP3, PI, SEP1-4 \\
3 & Stamen & B + C + E & AP3, PI, AG, SEP1-4 \\
4 & Carpel & C + E & AG, SEP1-4 \\
Inside whorl 4 & Ovule & D + E & SHP1, SHP2, STK, SEP1-4 \\
\end{longtable}
}

\textbf{Molecular identities of the ABC classes:}

\begin{itemize}
\tightlist
\item
  \textbf{A class (AP1, AP2):} AP1 is itself a MADS-box gene; AP2
  belongs to the unrelated AP2/ERF family --- an ``ABC anomaly''
  reflecting the fact that the canonical model was assembled before most
  molecular identities were known. AP1 specifies sepal identity in whorl
  1 and contributes to petal identity in whorl 2.
\item
  \textbf{B class (AP3 and PI):} Both are MADS-box genes; AP3 and PI
  form an obligate heterodimer (AP3-PI). This heterodimer is required
  for petal (with A) and stamen (with C) specification. \emph{ap3} or
  \emph{pi} single mutants lose petals and stamens.
\item
  \textbf{C class (AG = AGAMOUS):} A MADS-box gene that specifies stamen
  identity (with B) and carpel identity (alone, with E). AG also
  enforces floral meristem determinacy: in \emph{ag} mutants, the floral
  meristem proliferates indefinitely, generating extra petals in place
  of stamens and a new flower in place of a carpel --- the genetic basis
  of ``double flowers'' in cultivated roses, carnations, and
  chrysanthemums.
\item
  \textbf{D class (STK, SHP1, SHP2):} MADS-box genes specifying ovule
  identity within whorl 4. Loss of D-class function converts ovules to
  leaf-like structures.
\item
  \textbf{E class (SEP1, SEP2, SEP3, SEP4):} SEPALLATA MADS-box genes;
  required for combinatorial action of A, B, C, D. \emph{sep1 sep2 sep3
  sep4} quadruple mutants produce primarily sepal-like organs in most
  whorls.
\end{itemize}

\textbf{The MADS-box and the floral quartet model:}

The \textbf{MADS-box} is a 56-amino-acid DNA-binding domain (named for
MCM1, AGAMOUS, DEFICIENS, SRF --- the founding members). MADS-box
proteins bind CArG-box DNA motifs (CC{[}A/T{]}\(_6\)GG) and form
\textbf{floral quartets} (Theissen and Saedler, 2001): tetrameric
complexes that combinatorially specify each whorl's identity by binding
two CArG boxes simultaneously, looping the intervening DNA. The
four-protein composition determines target gene specificity:

\begin{itemize}
\tightlist
\item
  AP1-AP1-SEP3-SEP3 \texorpdfstring{\ensuremath{\to}}{->} sepal
\item
  AP3-PI-AP1-SEP3 \texorpdfstring{\ensuremath{\to}}{->} petal
\item
  AP3-PI-AG-SEP3 \texorpdfstring{\ensuremath{\to}}{->} stamen
\item
  AG-AG-SEP3-SEP3 \texorpdfstring{\ensuremath{\to}}{->} carpel
\end{itemize}

\begin{equation}
\text{A} + \text{B} + \text{C} + \text{D} + \text{E} \xrightarrow{\text{combinatorial CArG binding}} \text{floral quartet}
\label{eq:unit_VIII_floral_quartet}
\end{equation}

\textbf{Key predictions and experimental validation:}

\begin{itemize}
\tightlist
\item
  Removing A class (\emph{ap1}/\emph{ap2} mutants) results in carpels in
  whorls 1 and 2 (C class spreads to fill most whorls)
\item
  Removing B class (\emph{ap3}/\emph{pi} mutants): sepals in whorls 1
  and 2; carpels in whorls 3 and 4
\item
  Removing C class (\emph{ag} mutant): petals instead of stamens (whorl
  3) and further petals indefinitely instead of carpel (whorl 4);
  produces double-petalled ``full'' flowers --- the classic phenotype of
  cultivated double roses, carnations, and chrysanthemums
\item
  Removing E class (\emph{sep1 sep2 sep3 sep4} quadruple mutant): most
  whorls produce sepal-like organs, demonstrating SEP genes are required
  for B, C, and D function
\item
  \textbf{Quintuple ABCDE mutant:} Most whorls produce leaf-like organs
  --- confirming the molecular interpretation of Goethe's 1790
  metamorphosis hypothesis that most floral organs are modified leaves
\end{itemize}

\begin{quote}
\textbf{Clinical Connection:} Understanding the ABCDE model has
practical applications in crop breeding. Manipulating floral organ
identity genes can create male-sterile lines (essential for hybrid seed
production in rice and wheat) by silencing the B class to convert
stamens to petals. Conversely, increasing petal number in ornamental
varieties relies on partial loss-of-function \emph{ag} alleles. The
multi-billion dollar cut flower industry relies directly on MADS-box
gene variants for double-flowered roses, carnations, and chrysanthemums.
\end{quote}

See \cref{eq:unit_VIII_floral_quartet} for the combinatorial logic; the
empirical demonstration in \emph{Antirrhinum} and \emph{Arabidopsis} is
the foundational work of \citet{coen1991}.

\begin{center}\rule{0.5\linewidth}{0.5pt}\end{center}

\subsection{Gametophyte Generation -- Cellular
Architecture}\label{gametophyte-generation-cellular-architecture}

\subsubsection{Male Gametophyte: Microsporogenesis and Pollen Grain
Architecture}\label{male-gametophyte-microsporogenesis-and-pollen-grain-architecture}

The mature pollen grain is the most reduced male gametophyte known ---
three cells encased in a desiccation-resistant wall. Its development is
a textbook case of asymmetric cell division.

\begin{figure}[htbp]
\centering
\includegraphics[width=0.82\linewidth,height=4.2in,keepaspectratio]{../figures/mermaid_inline/inline_mermaid_0144_c19e702b611c.png}
\caption{Architecture of an angiosperm pollen grain at maturity The tube cell envelops the generative cell, which divides to yield two sperm cells either before pollen release (tricellular pollen) or during tube growth (bicellular pollen at release).}
\end{figure}

\textbf{Pollen wall layers:}

\begin{itemize}
\tightlist
\item
  \textbf{Intine (inner):} Cellulose + pectin; uniform thickness;
  secreted by the gametophyte itself
\item
  \textbf{Exine (outer):} Sporopollenin (oxidatively polymerised
  phenylpropanoids and fatty acids); deposited by the surrounding
  sporophytic tapetum. Sporopollenin is among the most chemically inert
  biopolymers known --- it is essentially indestructible by acid, base,
  or enzymatic attack, allowing pollen to survive millions of years in
  sediment. Exine sculpturing patterns (colpi, pores, reticulations) are
  diagnostic for plant families and form the basis of
  \textbf{palynology} (fossil pollen analysis).
\item
  \textbf{Tryphine and pollenkitt:} Lipid-rich material in exine
  cavities; contains species-specific recognition factors (S-locus
  proteins, lipidic adhesins) and is essential for hydration on the
  stigma surface
\end{itemize}

\textbf{Pollen cell architecture and microsporogenesis:}

The mature pollen grain is a 3-celled structure: one large
\textbf{vegetative (tube) cell} with a single decondensed nucleus,
enclosing two small \textbf{sperm cells} (the germline). Its development
proceeds through five steps:

\begin{enumerate}
\def\labelenumi{\arabic{enumi}.}
\tightlist
\item
  \textbf{Archesporial cells} within the anther microsporangium
  differentiate into diploid \textbf{microspore mother cells (MMCs)} ---
  also called pollen mother cells
\item
  Each MMC undergoes \textbf{meiosis I + II} to produce \textbf{4
  haploid microspores}, initially held as a tetrad surrounded by a
  \textbf{callose} wall (\texorpdfstring{\ensuremath{\beta}}{beta}-1,3-glucan). The tapetum (a sporophytic
  somatic cell layer surrounding the developing microspores) secretes
  \textbf{callase} (\texorpdfstring{\ensuremath{\beta}}{beta}-1,3-glucanase) at a precisely timed point to
  dissolve callose and release individual microspores.
\item
  \textbf{Microspore mitosis I} is \textbf{asymmetric}: the microspore
  nucleus migrates to one side, and an asymmetric division produces a
  small \textbf{generative cell} and a large \textbf{vegetative cell}.
  The generative cell is engulfed by the vegetative cell cytoplasm --- a
  ``cell within a cell'' architecture unique to plant gametogenesis.
  Asymmetric division depends on the GAMETOPHYTE-DEFECTIVE 1 / DUO1
  transcriptional network.
\item
  \textbf{Microspore mitosis II} divides the generative cell to produce
  \textbf{2 sperm cells}:

  \begin{itemize}
  \tightlist
  \item
    In \textbf{tricellular pollen} (about 30\% of species;
    \emph{Arabidopsis}, \emph{Brassicaceae}, grasses) mitosis II
    completes before pollen is released from the anther
  \item
    In \textbf{bicellular pollen} (about 70\% of species; most
    flowering plants) mitosis II occurs in the growing pollen tube after
    pollination
  \end{itemize}
\item
  The mature pollen grain = 1 vegetative cell + 2 sperm cells = the
  complete \textbf{male gametophyte}
\end{enumerate}

The two sperm cells are not equivalent: one is ``leading'' and tends to
fuse with the egg, the other ``trailing'' with the central cell. The
recognition mechanism uses cell-surface markers --- though preferred
matching is statistical rather than absolute.

\subsubsection{Female Gametophyte: Megasporogenesis and the
Polygonum-Type Embryo
Sac}\label{female-gametophyte-megasporogenesis-and-the-polygonum-type-embryo-sac}

The angiosperm female gametophyte is also drastically reduced --- a
7-celled, 8-nucleate structure embedded in the ovule. The
\textbf{Polygonum-type} embryo sac (about 70\% of angiosperms)
is the canonical form.

\textbf{Cell types and their functions:}

\begin{itemize}
\tightlist
\item
  \textbf{Egg cell} (1, micropylar): Haploid; becomes the zygote upon
  fertilisation. Polarised: nucleus near the chalazal end, large vacuole
  at the micropylar end. Specifies the apical--basal axis of the future
  embryo.
\item
  \textbf{Synergids} (2, micropylar; flank the egg): Haploid. Critical
  functions:

  \begin{itemize}
  \tightlist
  \item
    Secrete \textbf{LURE peptides} (CRP810 family; species-specific)
    through the \textbf{filiform apparatus} (an elaborate cell-wall
    labyrinth at the micropylar pole that increases secretory surface
    area), guiding the pollen tube to the ovule
  \item
    One synergid degenerates (programmed cell death) upon pollen tube
    arrival, providing the entry channel
  \item
    Express the \textbf{FERONIA-LORELEI receptor complex} required for
    pollen tube reception and rupture
  \end{itemize}
\item
  \textbf{Central cell} (1, large, central): Contains \textbf{2 polar
  nuclei} (which fuse before or during fertilisation to form a diploid
  secondary nucleus). Becomes the \textbf{triploid endosperm} upon
  fertilisation with one sperm. The canonical 2:1 maternal:paternal
  genome dosage is widespread across angiosperms and is the substrate
  for genomic imprinting and parental conflict, as discussed in the
  genomic-imprinting section.
\item
  \textbf{Antipodal cells} (3, chalazal): Function in nutritive/transfer
  roles. Highly polyploid (up to 32 n via endoreduplication) in grasses,
  where they are long-lived; in \emph{Arabidopsis} and most dicots,
  antipodals degenerate before fertilisation.
\end{itemize}

\textbf{Megasporogenesis and Embryo Sac Development:}

\begin{enumerate}
\def\labelenumi{\arabic{enumi}.}
\tightlist
\item
  A single \textbf{megaspore mother cell (MMC)} (also called
  megasporocyte) within the ovule's nucellus undergoes meiosis,
  producing 4 haploid megaspores in a linear tetrad along the
  chalazal-micropylar axis
\item
  \textbf{Three megaspores degenerate} via programmed cell death; the
  surviving (typically chalazal) \textbf{functional megaspore} alone
  gives rise to the female gametophyte. The selectivity of which
  megaspore survives is partly genetic (auxin gradients, AGO9-dependent
  siRNA silencing of competing megaspores) and partly stochastic.
\item
  The functional megaspore undergoes \textbf{3 rounds of free nuclear
  mitosis} (without cytokinesis) producing an 8-nucleate coenocyte. The
  8 nuclei migrate to specific positions: 4 nuclei toward each pole.
\item
  \textbf{Cellularisation:} Cell walls form around the 8 nuclei to
  establish the 7 cells: 1 egg + 2 synergids + 1 central cell (with 2
  polar nuclei, one from each pole) + 3 antipodals = the \textbf{mature
  embryo sac}
\end{enumerate}

\textbf{Concept Check 2:} A mutation eliminates the synergids during
embryo sac development (synergid-less \emph{gametophyte mutant1}).
Predict which steps of fertilisation will fail and why double
fertilisation cannot proceed.

\begin{center}\rule{0.5\linewidth}{0.5pt}\end{center}

\subsection{Pollination, Pollen-Pistil Interactions, and Double
Fertilisation}\label{pollination-pollen-pistil-interactions-and-double-fertilisation}

\subsubsection{Pollination Systems}\label{pollination-systems}

\textbf{Transfer of pollen from anther to stigma} occurs by various
vectors, each associated with characteristic floral traits (pollination
syndromes):

{\def\LTcaptype{none} % do not increment counter
\begin{longtable}[]{@{}
  >{\raggedright\arraybackslash}p{(\linewidth - 6\tabcolsep) * \real{0.1333}}
  >{\raggedright\arraybackslash}p{(\linewidth - 6\tabcolsep) * \real{0.1778}}
  >{\raggedright\arraybackslash}p{(\linewidth - 6\tabcolsep) * \real{0.3778}}
  >{\raggedright\arraybackslash}p{(\linewidth - 6\tabcolsep) * \real{0.3111}}@{}}
\toprule\noalign{}
\begin{minipage}[b]{\linewidth}\raggedright
Vector
\end{minipage} & \begin{minipage}[b]{\linewidth}\raggedright
Syndrome
\end{minipage} & \begin{minipage}[b]{\linewidth}\raggedright
Flower morphology
\end{minipage} & \begin{minipage}[b]{\linewidth}\raggedright
Plant examples
\end{minipage} \\
\midrule\noalign{}
\endhead
\bottomrule\noalign{}
\endlastfoot
Bees (\emph{Apis}, \emph{Bombus}) & Melittophily &
Blue/yellow/UV-reflective; landing platform; nectar guide &
\emph{Trifolium}, \emph{Solanum}, \emph{Linaria} \\
Butterflies & Psychophily & Red/pink; narrow tube; nectar platform;
day-opening & \emph{Asclepias}, \emph{Buddleja} \\
Hawkmoths & Sphingophily & White; strong scent dusk-dawn; long narrow
tube; night-opening & \emph{Nicotiana}, \emph{Datura} \\
Bats & Chiropterophily & Large; white/pale; fermented/fruity scent;
nocturnal; robust & \emph{Agave}, \emph{Ceiba}, \emph{Adansonia} \\
Birds (hummingbirds) & Ornithophily & Red/orange (beyond bee vision); no
or little scent; abundant dilute nectar; tubular & \emph{Strelitzia},
\emph{Lobelia}, \emph{Fuchsia} \\
Wind & Anemophily & Small, inconspicuous; no nectar; enormous pollen
quantities; feathery stigma & \emph{Poa}, \emph{Quercus}, grasses,
\emph{Betula} \\
Water & Hydrophily & Variable; pollen at water level &
\emph{Vallisneria}, \emph{Ceratophyllum} \\
\end{longtable}
}

\textbf{80\% of flowering plants} are insect-pollinated (entomophilous).
Global pollination services estimated at about \$577
billion/year (IPBES 2016).

\textbf{Concept Check 3:} A plant species has small, green, scentless
flowers that produce copious dry pollen. What is its likely pollination
vector? What other structural features would you expect?

\subsubsection{Pollen-Pistil Interactions: Self-Incompatibility (SI)
Systems}\label{pollen-pistil-interactions-self-incompatibility-si-systems}

Self-incompatibility (SI) is a widespread genetic mechanism preventing
self-pollination and inbreeding depression. Approximately \textbf{40\%
of angiosperm species} have functional SI. The molecular logic divides
into two paradigms based on whether the pollen's incompatibility
phenotype is determined by its own (haploid) genotype or by the
(diploid) genotype of the pollen-producing parent.

\textbf{Sporophytic SI (SSI)} --- \emph{Brassica} (cabbage, mustard),
\emph{Ipomoea}, Asteraceae

The pollen rejection phenotype is determined by the \textbf{diploid}
sporophytic tissue of the anther (specifically, the tapetum that
deposits proteins onto the pollen wall) --- \emph{not} by the pollen
grain's own haploid genotype.

\begin{itemize}
\tightlist
\item
  \textbf{S-locus genes:}

  \begin{itemize}
  \tightlist
  \item
    \textbf{SRK} (S-receptor kinase): Stigma-expressed; receptor on the
    stigma papilla membrane
  \item
    \textbf{SCR/SP11} (S-locus cysteine-rich protein 11; pollen coat
    protein): Deposited on pollen exine by the diploid tapetum during
    pollen maturation
  \end{itemize}
\item
  \textbf{Mechanism:} When SRK on the stigma encounters its matching SCR
  on the incoming pollen (i.e., self-pollen sharing the same S-haplotype
  as the stigma), SRK autophosphorylates and recruits \textbf{ARC1} (E3
  ubiquitin ligase). ARC1 ubiquitinates \textbf{Exo70A1}, a key
  component of the secretion exocyst required for pollen hydration. The
  pollen fails to hydrate and is rejected at the stigma surface ---
  before pollen tube germination.
\item
  \textbf{Genetic consequence:} Because pollen carries the diploid
  parental S-genotype, most pollen from a heterozygous \emph{S\texorpdfstring{\textsubscript{1}}{1}S\texorpdfstring{\textsubscript{2}}{2}} plant
  is rejected by \emph{S\texorpdfstring{\textsubscript{1}}{1}S\texorpdfstring{\textsubscript{2}}{2}}, \emph{S\texorpdfstring{\textsubscript{1}}{1}S\texorpdfstring{\textsubscript{3}}{3}}, and \emph{S\texorpdfstring{\textsubscript{2}}{2}S\texorpdfstring{\textsubscript{3}}{3}} stigmas.
\end{itemize}

\textbf{Gametophytic SI (GSI)} --- Solanaceae (\emph{Petunia},
\emph{Nicotiana}, tomato), Rosaceae (apple, pear, almond), Papaveraceae

The pollen rejection phenotype is determined by the pollen grain's own
\textbf{haploid} S-genotype.

\begin{itemize}
\tightlist
\item
  \textbf{S-locus genes (S-RNase / SLF system):}

  \begin{itemize}
  \tightlist
  \item
    \textbf{S-RNase}: Pistil-expressed; ribonuclease secreted into the
    style transmitting tract
  \item
    \textbf{SLF (S-locus F-box protein)}: Pollen-expressed; multiple
    SLFs per haplotype, acting collectively
  \end{itemize}
\item
  \textbf{Mechanism (collaborative non-self recognition):} S-RNase
  enters the growing pollen tube non-specifically. Each SLF can ``tag''
  non-self S-RNases (those whose haplotype does not match the pollen
  tube's own) for ubiquitination and proteasomal degradation by the
  SCF\textsuperscript{SLF} complex. \textbf{Self S-RNase escapes
  degradation} because the pollen tube's SLFs cannot recognise the
  matching S-RNase. Surviving S-RNase enters the cytoplasm and degrades
  pollen tube ribosomal RNA, halting tube growth before fertilisation.
\item
  \textbf{Genetic consequence:} A cross between \emph{S\texorpdfstring{\textsubscript{1}}{1}S\texorpdfstring{\textsubscript{2}}{2}} and
  \emph{S\texorpdfstring{\textsubscript{1}}{1}S\texorpdfstring{\textsubscript{3}}{3}} plants: \emph{S\texorpdfstring{\textsubscript{1}}{1}} pollen is rejected by both stigmas;
  \emph{S\texorpdfstring{\textsubscript{2}}{2}} pollen succeeds on the \emph{S\texorpdfstring{\textsubscript{1}}{1}S\texorpdfstring{\textsubscript{3}}{3}} stigma; \emph{S\texorpdfstring{\textsubscript{3}}{3}} pollen
  succeeds on the \emph{S\texorpdfstring{\textsubscript{1}}{1}S\texorpdfstring{\textsubscript{2}}{2}} stigma. Half the pollen is accepted;
  offspring genotypes are \emph{S\texorpdfstring{\textsubscript{1}}{1}S\texorpdfstring{\textsubscript{2}}{2}}, \emph{S\texorpdfstring{\textsubscript{1}}{1}S\texorpdfstring{\textsubscript{3}}{3}}, \emph{S\texorpdfstring{\textsubscript{2}}{2}S\texorpdfstring{\textsubscript{3}}{3}} (no
  homozygotes possible).
\end{itemize}

\textbf{Comparative summary:}

{\def\LTcaptype{none} % do not increment counter
\begin{longtable}[]{@{}
  >{\raggedright\arraybackslash}p{(\linewidth - 4\tabcolsep) * \real{0.1167}}
  >{\raggedright\arraybackslash}p{(\linewidth - 4\tabcolsep) * \real{0.4167}}
  >{\raggedright\arraybackslash}p{(\linewidth - 4\tabcolsep) * \real{0.4667}}@{}}
\toprule\noalign{}
\begin{minipage}[b]{\linewidth}\raggedright
Feature
\end{minipage} & \begin{minipage}[b]{\linewidth}\raggedright
Sporophytic SI (Brassica)
\end{minipage} & \begin{minipage}[b]{\linewidth}\raggedright
Gametophytic SI (Solanaceae)
\end{minipage} \\
\midrule\noalign{}
\endhead
\bottomrule\noalign{}
\endlastfoot
Determinant of pollen phenotype & Diploid parental tissue (tapetum) &
Pollen's own haploid genome \\
Pollen recognition & At stigma surface (before germination) & In style
during tube growth \\
Female factor & SRK (membrane receptor kinase) & S-RNase (secreted
RNase) \\
Male factor & SCR (pollen coat protein) & SLF (F-box) \\
Number of S-haplotypes & Often 50--100 & Often 50--200 \\
\end{longtable}
}

\textbf{Evolutionary significance of SI:} Self-incompatibility maintains
heterozygosity and avoids inbreeding depression in obligate outbreeders.
Its breakdown (typically by loss of S-RNase or SLF function) gives rise
to self-compatible lineages, which often go extinct because of genetic
load --- but occasionally undergo spectacular adaptive radiations (e.g.,
self-compatible \emph{Arabidopsis thaliana} descended from
self-incompatible \emph{A. lyrata}; the loss of SI may have facilitated
\emph{A. thaliana}'s colonisation of new habitats by enabling
single-individual founder events).

\subsubsection{Polyploidy and Speciation in
Plants}\label{polyploidy-and-speciation-in-plants}

Plants exhibit far higher rates of polyploidy than animals;
\textbf{about 70\% of angiosperm species} retain evidence of
one or more whole-genome duplications in their evolutionary history.

\textbf{Modes of polyploid origin:}

\begin{itemize}
\tightlist
\item
  \textbf{Autopolyploidy:} Whole-genome duplication within a single
  species; chromosomes form quadrivalents at meiosis, often causing
  reduced fertility initially
\item
  \textbf{Allopolyploidy:} Hybridisation between two species followed by
  chromosome doubling; chromosomes from each parental genome pair
  separately as bivalents, restoring fertility. Most cultivated
  polyploids are allopolyploids.
\end{itemize}

\textbf{Pathways:}

\begin{enumerate}
\def\labelenumi{\arabic{enumi}.}
\tightlist
\item
  \textbf{Unreduced gamete fusion:} Two 2n gametes fuse to form a 4n
  zygote (one in 1000 gametes is unreduced)
\item
  \textbf{Somatic doubling:} Spontaneous chromosome doubling in a
  meristematic cell during somatic growth
\item
  \textbf{Triploid bridge:} Triploids (3n) produce a wide variety of
  unbalanced gametes; some 4n offspring arise from these
\end{enumerate}

\textbf{Polyploid crops:} Wheat (hexaploid: AABBDD genomes from three
ancestral diploids), cotton (tetraploid), strawberry (octoploid), banana
(often triploid, sterile), oats (hexaploid).

\textbf{Polyploidy promotes speciation} because new polyploids are
reproductively isolated from their parental species (triploid offspring
of a 2n \texorpdfstring{\ensuremath{\times}}{x} 4n cross are sterile), creating an instant reproductive
barrier. Polyploids often exhibit \textbf{heterosis} (hybrid vigour),
enlarged organs (the basis of many crop polyploids), and increased
ecological tolerance.

\subsubsection{\texorpdfstring{Double Fertilisation
\citep{nawaschin1898}}{Double Fertilisation}}\label{double-fertilisation}

Sergei Nawaschin's 1898 discovery of double fertilisation in
\emph{Lilium martagon} and \emph{Fritillaria tenella} --- observing two
simultaneous nuclear fusions in a single ovule --- was transformative
for plant biology and remains the defining synapomorphy of angiosperms.

\begin{figure}[htbp]
\centering
\includegraphics[width=0.82\linewidth,height=4.2in,keepaspectratio]{../figures/mermaid_inline/inline_mermaid_0145_712477dd19ef.png}
\caption{Double fertilisation in angiosperms One sperm fuses with the egg to form the diploid zygote (which develops into the embryo). The second sperm fuses with the central cell's two polar nuclei to form the triploid primary endosperm nucleus (which develops into the nutritive endosperm). This process is unique to angiosperms.}
\end{figure}

\begin{enumerate}
\def\labelenumi{\arabic{enumi}.}
\item
  \textbf{Pollen tube germination on the stigma:} Compatible pollen
  hydrates within minutes on the stigma surface; the \textbf{vegetative
  (tube) cell} extends a callose-walled tube. Tip growth is driven by:

  \begin{itemize}
  \tightlist
  \item
    A \textbf{tip-focused {[}Ca\(^{2+}\){]} gradient}
    (about 1.5--10 \texorpdfstring{\ensuremath{\mu}}{mu}M at the apex; about 150 nM in the
    shank)
  \item
    Vesicle-mediated wall deposition (callose for tube wall, pectin for
    the inner wall)
  \item
    Cytoplasmic streaming carrying the generative/sperm cells to the tip
  \item
    Growth rates: 1--10 mm/h in \emph{Lilium}, 1 cm/h in maize
  \end{itemize}
\item
  \textbf{Spermatogenesis during tube elongation (in bicellular
  species):} The generative cell, carried in the tube cell cytoplasm,
  undergoes mitosis II during pollen tube growth, producing the two
  sperm cells. In tricellular species (\emph{Arabidopsis}, grasses),
  this division has already occurred before pollen release.
\item
  \textbf{Pollen tube guidance --- chemotropism:} The pollen tube
  navigates through the style and ovary toward a single receptive ovule
  via a cascade of guidance cues:

  \begin{itemize}
  \tightlist
  \item
    \textbf{Style cues:} \texorpdfstring{\ensuremath{\gamma}}{gamma}-aminobutyric acid (GABA) gradient in the
    transmitting tract; cysteine-rich peptides
  \item
    \textbf{Funicular guidance:} Short-range cues from ovary tissue
  \item
    \textbf{Micropylar guidance:} \textbf{LURE peptides} (defensin-like,
    cysteine-rich; CRP810 family) released from the synergid filiform
    apparatus; species-specific. \emph{Torenia} LURE1/LURE2 attract
    pollen tubes specifically of the same species --- the molecular
    basis of inter-species pollination barriers. LURE peptides bind PRK6
    receptor kinases on the pollen tube tip, biasing tip growth toward
    the source.
  \end{itemize}
\item
  \textbf{Synergid degeneration:} As the pollen tube approaches the
  ovule, \textbf{one of the two synergids undergoes programmed cell
  death} (within minutes; the receptive synergid). Loss of synergid
  integrity:

  \begin{itemize}
  \tightlist
  \item
    Creates the entry channel into the embryo sac
  \item
    Releases pre-stored signalling molecules
  \item
    Exposes the \textbf{FERONIA (FER)} receptor kinase / \textbf{LORELEI
    (LRE)} GPI-anchored co-receptor complex on the persisting synergid
    surface
  \end{itemize}
\item
  \textbf{Pollen tube reception and rupture:} FER-LRE signalling, via
  reactive oxygen species (ROS) generated by RBOH NADPH oxidases and
  elevated cytosolic {[}Ca\(^{2+}\){]}, triggers explosive rupture of
  the pollen tube tip, releasing the two sperm cells into the embryo
  sac. The signalling depends on RALF peptides binding FER. In
  \emph{fer} mutants, pollen tubes enter the synergid but fail to
  rupture (continued growth and ``supernumerary'' pollen tube delivery
  --- the polyspermy phenotype).
\item
  \textbf{Sperm cell delivery and double fusion:}

  \begin{itemize}
  \tightlist
  \item
    \textbf{Sperm 1 + egg cell} \texorpdfstring{\ensuremath{\to}}{->} karyogamy \texorpdfstring{\ensuremath{\to}}{->} \textbf{2n zygote} \texorpdfstring{\ensuremath{\to}}{->}
    embryo
  \item
    \textbf{Sperm 2 + 2 polar nuclei (central cell)} \texorpdfstring{\ensuremath{\to}}{->} triple fusion \texorpdfstring{\ensuremath{\to}}{->}
    \textbf{3n primary endosperm nucleus} \texorpdfstring{\ensuremath{\to}}{->} endosperm
  \end{itemize}
\end{enumerate}

\begin{equation}
\text{LURE} \xrightarrow{\text{LRE-FER complex}} \text{tube rupture} \rightarrow 2 \text{ sperm released}
\label{eq:unit_VIII_lre_fer}
\end{equation}

\textbf{Significance of double fertilisation:}

\begin{itemize}
\tightlist
\item
  Endosperm provides a ``payment on delivery'' mechanism --- primarily
  fertilised ovules develop endosperm, ensuring optimal parental
  investment
\item
  Endosperm in cereals: starch (70-80\%), storage proteins (gluten in
  wheat, zein in maize), oils, vitamins --- source of
  about 60\% of global human caloric intake
\item
  Triploid endosperm (2 maternal : 1 paternal genome dose) provides a
  unique genetic substrate for \textbf{parental conflict}
  (paternally-imprinted genes promote nutrient transfer;
  maternally-imprinted genes restrict it) --- conserved in ratio across
  most flowering plants
\end{itemize}

\begin{quote}
\textbf{Clinical Connection:} Understanding pollen tube guidance has
implications for crop fertility. In interspecific crosses (e.g., wheat \texorpdfstring{\ensuremath{\times}}{x}
rye to produce triticale), pollen tube guidance often fails because LURE
peptides are species-specific. Genetic engineering of LURE receptors
could enable wider crosses for crop improvement.
\end{quote}

\begin{center}\rule{0.5\linewidth}{0.5pt}\end{center}

\subsection{Endosperm Development and Seed
Biology}\label{endosperm-development-and-seed-biology}

\subsubsection{Endosperm Development: Free-Nuclear, Cellular,
Helobial}\label{endosperm-development-free-nuclear-cellular-helobial}

After triple fusion, the primary endosperm cell (3n; typically 2m:1p
genome ratio) undergoes a stereotyped developmental sequence that
determines the storage architecture of the mature seed. The 2:1
maternal:paternal genome ratio is the canonical angiosperm pattern and
is important for normal seed development, while some lineages show
dosage and developmental variation.

\textbf{Three modes of endosperm development:}

\begin{enumerate}
\def\labelenumi{\arabic{enumi}.}
\tightlist
\item
  \textbf{Nuclear (free-nuclear) endosperm --- most common:} The primary
  endosperm nucleus undergoes repeated mitosis without cytokinesis,
  producing a \textbf{syncytium} of hundreds to thousands of nuclei in a
  common cytoplasm with a large central vacuole. Examples:
  \emph{Arabidopsis}, maize, rice, wheat, coconut (the ``coconut water''
  of an immature coconut is liquid free-nuclear endosperm). After
  about 6--8 nuclear divisions, \textbf{cellularisation}
  initiates from the periphery inward, with cell walls forming radially
  around each nucleus. The mature cellular endosperm then accumulates
  storage products (starch, protein, oil).
\item
  \textbf{Cellular endosperm:} Cytokinesis accompanies every mitosis
  from the start; the endosperm is cellular throughout. Examples:
  \emph{Petunia}, tobacco, magnoliid lineages.
\item
  \textbf{Helobial endosperm:} The first division separates a small
  chalazal cell from a large micropylar cell; subsequent divisions are
  nuclear in the micropylar half and cellular in the chalazal half.
  Examples: many monocots in Alismatales.
\end{enumerate}

\textbf{Seed filling:} During the maturation phase (after
cellularisation), endosperm cells accumulate massive amounts of storage
compounds: - \textbf{Starch:} Synthesised from sucrose imported from
maternal phloem; deposited in plastids (amyloplasts). Cereal grain
starch can reach 60--70\% of dry mass. - \textbf{Storage proteins:}
Family-specific (zeins in maize, glutenins/gliadins in wheat, prolamines
in rice, globulins in legumes). Stored in protein bodies derived from
the ER and vacuole. - \textbf{Oils:} Triacylglycerols stored in oil
bodies (oleosomes); dominant in oilseeds (sunflower, canola, soybean).

\textbf{The PEG/MEG pathway and parental conflict:}

The \textbf{FIS-PRC2} (FERTILIZATION INDEPENDENT SEED -- Polycomb
Repressive Complex 2) complex --- comprising MEDEA (MEA), FIS2, FIE, and
MSI1 --- is the master regulator of endosperm initiation and parental
conflict resolution:

\begin{equation}
\text{Maternal MEA-FIE-FIS2-MSI1} \xrightarrow{\text{H3K27me3}} \text{repression of paternally-imprinted genes (PEGs)}
\label{eq:unit_VIII_peg_pathway}
\end{equation}

\begin{itemize}
\tightlist
\item
  \textbf{PEGs (Paternally Expressed Genes):} Expressed primarily from
  the paternal allele in endosperm (the maternal allele is silenced by
  H3K27me3 deposited by FIS-PRC2). Examples: \emph{PHE1}, \emph{YUC10},
  \emph{AHL10}. PEGs tend to \textbf{promote} endosperm growth and
  nutrient transfer to the embryo.
\item
  \textbf{MEGs (Maternally Expressed Genes):} Expressed primarily from
  the maternal allele. MEGs tend to \textbf{restrict} endosperm growth.
\item
  \textbf{Imbalance in PEG/MEG dosage} (e.g., crosses between species
  with different ploidy) causes seed abortion via the \textbf{endosperm
  balance number (EBN)} mechanism. This is why interploidy crosses fail:
  a 2x \texorpdfstring{\ensuremath{\times}}{x} 4x cross produces triploid embryos with imbalanced PEG/MEG
  expression, triggering seed abortion.
\end{itemize}

The 2m:1p endosperm dosage thus becomes the principal arena for
\textbf{parent-of-origin genetic conflict}, with paternally-imprinted
PEGs evolutionarily favoured to extract more maternal resources for the
offspring, and maternally-imprinted MEGs favoured to ration resources
across multiple offspring.

\subsubsection{Seed Dormancy: Five Classes and
Triggers}\label{seed-dormancy-five-classes-and-triggers}

\textbf{Seed dormancy} delays germination until conditions favour
seedling survival. Baskin and Baskin (2004) classify dormancy into five
types:

{\def\LTcaptype{none} % do not increment counter
\begin{longtable}[]{@{}
  >{\raggedright\arraybackslash}p{(\linewidth - 6\tabcolsep) * \real{0.3095}}
  >{\raggedright\arraybackslash}p{(\linewidth - 6\tabcolsep) * \real{0.1190}}
  >{\raggedright\arraybackslash}p{(\linewidth - 6\tabcolsep) * \real{0.3810}}
  >{\raggedright\arraybackslash}p{(\linewidth - 6\tabcolsep) * \real{0.1905}}@{}}
\toprule\noalign{}
\begin{minipage}[b]{\linewidth}\raggedright
Dormancy type
\end{minipage} & \begin{minipage}[b]{\linewidth}\raggedright
Cause
\end{minipage} & \begin{minipage}[b]{\linewidth}\raggedright
Trigger to break
\end{minipage} & \begin{minipage}[b]{\linewidth}\raggedright
Examples
\end{minipage} \\
\midrule\noalign{}
\endhead
\bottomrule\noalign{}
\endlastfoot
\textbf{Physical (PY)} & Water-impermeable seed coat (dense palisade of
macrosclereids; suberin layers) & Scarification: physical abrasion, fire
heat, gut passage, freeze-thaw & Many legumes (clover, \emph{Acacia});
\emph{Convolvulus}; many Malvaceae \\
\textbf{Physiological (PD)} & Hormonal balance (high ABA / low GA)
blocks germination despite imbibition; embryo viable & Stratification
(cold-moist); after-ripening (dry storage at warm); light (Pfr) &
\emph{Arabidopsis}, lettuce, apple, most temperate trees \\
\textbf{Morphological (MD)} & Embryo immature at dispersal; needs time
to develop & Time + warm-moist conditions & \emph{Ginkgo},
\emph{Magnolia}, parsnips, celery \\
\textbf{Morphophysiological (MPD)} & Combined immature embryo + hormonal
block & Warm followed by cold stratification (or vice versa) & Many
woodland herbaceous spring ephemerals; \emph{Trillium},
\emph{Anemone} \\
\textbf{Combinational (PY+PD)} & Both impermeable coat and hormonal
block & Both scarification and stratification & \emph{Tilia} (linden),
some Rosaceae \\
\end{longtable}
}

\textbf{Germination triggers (sensory inputs):}

\begin{itemize}
\tightlist
\item
  \textbf{Stratification (cold, moist):} Cold treatment (1--5 °C, weeks
  to months) progressively degrades \textbf{ABA} and induces \textbf{GA
  biosynthesis} genes, shifting the ABA/GA balance toward germination.
  Required by many temperate tree seeds. Mechanism involves cold-induced
  demethylation of GA-biosynthesis gene promoters.
\item
  \textbf{Scarification (mechanical or chemical):} Cracks or thins the
  testa, allowing water and oxygen to enter. Can be achieved naturally
  by passage through animal digestive tracts (where stomach acid +
  abrasion abrade the coat), fire (heat scarification of fire-adapted
  species like \emph{Banksia}), or freeze-thaw cycles. Many savanna
  acacia seeds germinate primarily after passage through elephant or
  ungulate guts.
\item
  \textbf{Light (phytochrome-mediated):} Red light (R, 660 nm) converts
  Pr to active Pfr, promoting germination in light-requiring seeds
  (lettuce, \emph{Arabidopsis}, many small-seeded weeds). Far-red light
  (FR, 730 nm) inhibits germination. The R:FR ratio under a leaf canopy
  is low (chlorophyll absorbs R, transmits FR), so seeds under shade
  remain dormant until canopy opening --- an exquisite mechanism for
  detecting gaps in vegetation. The Borthwick-Hendricks experiments
  (1952) on lettuce seeds first demonstrated the R/FR reversibility that
  defined phytochrome.
\item
  \textbf{Temperature fluctuation:} Diurnal temperature swings (e.g., 40
  °C day / 10 °C night, typical of bare soil under sun without
  insulating vegetation) signal a gap in canopy cover. Many desert and
  weed species require such fluctuations.
\item
  \textbf{Smoke and karrikins:} Karrikin compounds (KAR1--KAR6) in
  plant-derived smoke bind the \textbf{KAI2 receptor} (a
  strigolactone-related \texorpdfstring{\ensuremath{\alpha}}{alpha}/\texorpdfstring{\ensuremath{\beta}}{beta} hydrolase). KAR signalling promotes
  germination of fire-adapted species and many weeds (KAR2 promotes
  germination in \emph{Arabidopsis}).
\item
  \textbf{Chemical leaching:} Some desert species require leaching of
  inhibitory compounds (e.g., NaCl in halophytes; phenolic germination
  inhibitors) by sufficient rainfall --- ensuring germination primarily
  after enough rain has fallen to support seedling establishment.
\end{itemize}

\subsubsection{Germination Physiology --- Molecular
Framework}\label{germination-physiology-molecular-framework}

The balance between ABA (dormancy-promoting) and GA
(germination-promoting) signals is the central regulatory axis:

\textbf{ABA pathway (dormancy):}

\textbf{DOG1} (DELAY OF GERMINATION 1): A seed-specific
RNA-binding-domain protein. DOG1 dosage determines dormancy depth by
stabilising ABI3 and ABI5 mRNA --- master transcription factors of the
ABA response. \emph{DOG1} expression peaks at seed maturation,
quantitative trait locus (QTL) responsible for natural variation in
dormancy across \emph{Arabidopsis} accessions (Bentsink \emph{et al.},
\emph{PNAS} 2006).

\begin{equation}
\text{ABA} + \text{PYR/RCAR} \rightleftharpoons \text{ABA-PYR complex} \xrightarrow{} \text{PP2C inhibition (ABI1/ABI2 released)} \rightarrow \text{SnRK2 kinase active}
\label{eq:unit_VIII_aba_germination}
\end{equation}

Active SnRK2 phosphorylates \textbf{ABF/AREB} transcription factors \texorpdfstring{\ensuremath{\to}}{->}
dormancy gene transcription (ABI5 targets: \emph{LEA} proteins, stress
tolerance genes). PP2C inhibition is the key switch: normally PP2C keeps
SnRK2 dephosphorylated (inactive); ABA binding to PYR/RCAR pulls PP2C
off SnRK2.

\textbf{GA pathway (germination):}

\begin{equation}
\text{GA} + \text{GID1} \rightleftharpoons \text{GA-GID1-DELLA complex} \xrightarrow{\text{SCF}^{\text{GID2}}} \text{Ub-DELLA} \xrightarrow{26S} \text{DELLA degradation}
\label{eq:unit_VIII_ga_della}
\end{equation}

DELLA degradation (of RGA, GAI, RGL1-3 in \emph{Arabidopsis}) releases
repression of: - \texorpdfstring{\ensuremath{\alpha}}{alpha}-Amylase \hyperref[gl:promoter]{\textbf{promoter}}s
(mobilise endosperm starch \texorpdfstring{\ensuremath{\to}}{->} maltose \texorpdfstring{\ensuremath{\to}}{->} glucose) - Lipase genes (mobilise
stored triacylglycerols) - Protease genes (mobilise globulins and
albumins)

\textbf{ABA--GA antagonism:} - ABA stabilises DELLA proteins (by
downregulating GID1 and upregulating PP2C) - GA degrades DELLA and
suppresses ABI5 - Environmental cues (light via Pfr, cold via DOG1
degradation, smoke via KAI2) tip the balance toward GA-dominated
germination competence

\textbf{Concept Check 4:} Some desert annual plant seeds require
specific temperature fluctuations (e.g., 40 °C day / 10 °C night) to
break dormancy. Explain why this diurnal fluctuation --- rather than
constant warm temperature --- is required. Connect your answer to the
molecular ABA/GA balance.

\textbf{Concept Check 5:} A \emph{dog1} loss-of-function mutant of
\emph{Arabidopsis} shows severely reduced primary dormancy: seeds
germinate immediately at maturity. Predict the ecological consequence in
(a) a Mediterranean climate with hot dry summers, and (b) a tropical
evergreen forest understory.

\begin{center}\rule{0.5\linewidth}{0.5pt}\end{center}

\subsection{Embryogenesis and Meristem
Organisation}\label{embryogenesis-and-meristem-organisation}

\subsubsection{Early Embryogenesis}\label{early-embryogenesis}

\textbf{Arabidopsis} embryogenesis is the best-studied model:

{\def\LTcaptype{none} % do not increment counter
\begin{longtable}[]{@{}
  >{\raggedright\arraybackslash}p{(\linewidth - 4\tabcolsep) * \real{0.1923}}
  >{\raggedright\arraybackslash}p{(\linewidth - 4\tabcolsep) * \real{0.4231}}
  >{\raggedright\arraybackslash}p{(\linewidth - 4\tabcolsep) * \real{0.3846}}@{}}
\toprule\noalign{}
\begin{minipage}[b]{\linewidth}\raggedright
Stage
\end{minipage} & \begin{minipage}[b]{\linewidth}\raggedright
Description
\end{minipage} & \begin{minipage}[b]{\linewidth}\raggedright
Key events
\end{minipage} \\
\midrule\noalign{}
\endhead
\bottomrule\noalign{}
\endlastfoot
1-cell zygote & Polar asymmetric cell & PIN7 (apical) + PIN1 (basal)
establish auxin gradient \\
2-cell & Asymmetric division: apical cell (embryo) + basal cell
(suspensor) & WOX9 signalling \\
Globular & 8-cell through 32-cell globular stage & WOX2 (apical domain),
WOX9 (basal) pattern embryo \\
Heart & Cotyledon primordia emerge; first visible bilateral symmetry &
YAB/PIN1 separate cotyledons; ARF5/MP establishes vascular axis \\
Torpedo & Elongation; hypocotyl + radicle extend & Protoderm,
procambium, ground meristem differentiate \\
Mature embryo & Desiccation; ABA accumulates; storage proteins deposited
& Two cotyledons + SAM + RAM + hypocotyl + radicle \\
\end{longtable}
}

The \textbf{suspensor} (from the basal cell) anchors the embryo to the
maternal tissue and transfers nutrients. It is programmed for cell death
after the globular stage.

\subsubsection{Shoot Apical Meristem (SAM) -- Stem Cell
Niche}\label{shoot-apical-meristem-sam-stem-cell-niche}

The SAM maintains a pool of pluripotent stem cells throughout the
plant's life:

\begin{itemize}
\tightlist
\item
  \textbf{Central Zone (CZ):} 2-4 slowly dividing stem cells at the
  summit
\item
  \textbf{Peripheral Zone (PZ):} Faster-dividing cells that are
  displaced outward and form lateral organ primordia (leaves, flowers)
\item
  \textbf{Rib Zone:} Beneath the CZ; contributes to stem internodes
\end{itemize}

\textbf{WUS-CLV3 negative feedback loop} (the ``stem cell thermostat''):

\begin{enumerate}
\def\labelenumi{\arabic{enumi}.}
\tightlist
\item
  \textbf{WUSCHEL (WUS)} -- homeodomain transcription factor expressed
  in the organising centre (below CZ). WUS protein moves to CZ via
  plasmodesmata. Activates \textbf{CLAVATA3 (CLV3)} expression.
\item
  \textbf{CLV3} -- a secreted CLE peptide (12 amino acids). CLV3 binds
  the \textbf{CLV1/CLV2/CRN receptor kinase} complex in underlying
  cells, activating a MAPK cascade that represses \textbf{WUS}
  transcription.
\item
  \textbf{Negative feedback:} WUS activates CLV3; CLV3 represses WUS.
  This maintains a constant stem cell pool size.
\end{enumerate}

\textbf{SAM auxin-driven phyllotaxis:} - Each leaf primordium is
initiated where PIN1 (auxin efflux carrier) creates a local auxin
maximum - The established primordium depletes local auxin from
surrounding cells - The next primordium arises at the position of the
next auxin maximum - Mathematical result: primordia arise at
\textbf{137.5 degrees divergence angle} (the ``golden angle'' = 360
degrees \(\times\) (1 \(-\) 1/\(\varphi^2\)) where \(\varphi\) = golden
ratio 1.618) - This produces Fibonacci spirals (1, 1, 2, 3, 5, 8,
13\ldots) visible in pine cones, sunflower heads, and succulent leaf
arrangements

\subsubsection{Worked Example: Calculating the Golden Angle in
Phyllotaxis}\label{worked-example-calculating-the-golden-angle-in-phyllotaxis}

\textbf{Problem:} The arrangement of leaves on a stem (phyllotaxis) is
often determined by the golden angle, which minimises shading of lower
leaves by upper leaves. Calculate the exact value of the golden angle in
degrees using the golden ratio \(\varphi \approx 1.618034\). If a plant
produces a new leaf every 5 days, what will be the total angular
divergence between the first leaf and a fourth leaf (leaf 1 and leaf 4)?

\textbf{Solution:}

\begin{enumerate}
\def\labelenumi{\arabic{enumi}.}
\item
  \textbf{Calculate the golden angle:} Using the formula given: Angle
  \(= 360^\circ \cdot (1 - \frac{1}{\varphi^2})\) First, calculate
  \(\varphi^2\):
  \[ 1.618034^2 \approx 2.618034  \label{eq:unit_VIII_plant_reproduction_item_1}\]

  Calculate the fraction:
  \[ \frac{1}{2.618034} \approx 0.381966  \label{eq:unit_VIII_plant_reproduction_item_2}\]

  Apply to the angle formula:
  \[ \text{Angle} = 360^\circ \cdot (1 - 0.381966) = 360^\circ \cdot 0.618034 \approx 222.49^\circ  \label{eq:unit_VIII_plant_reproduction_item_3}\]

  However, angles are typically measured by the shorter path around the
  circle, so we subtract from \(360^\circ\):
  \[ 360^\circ - 222.49^\circ = 137.51^\circ  \label{eq:unit_VIII_plant_reproduction_item_4}\]

  This is the \textbf{golden angle (about 137.5°)}.
\item
  \textbf{Calculate the total angular divergence:} Between leaf 1 and
  leaf 4, there are 3 developmental intervals (1 to 2, 2 to 3, and 3 to
  4).
  \[ \text{Total Divergence} = 3 \cdot 137.5^\circ = 412.5^\circ  \label{eq:unit_VIII_plant_reproduction_item_5}\]

  To find the apparent angle relative to the first leaf on a single
  \(360^\circ\) circle, take the modulo:
  \[ 412.5^\circ - 360^\circ = 52.5^\circ  \label{eq:unit_VIII_plant_reproduction_item_6}\]

  The fourth leaf will be separated from the first leaf by
  \textbf{\(52.5^\circ\)} on the stem circumference.
\end{enumerate}

\subsubsection{Root Apical Meristem
(RAM)}\label{root-apical-meristem-ram}

RAM structure mirrors SAM with distinct anatomy: - \textbf{Quiescent
Centre (QC):} 4-6 slowly dividing cells (divide about  once per
200 h); WOX5 expression maintains surrounding stem cells -
\textbf{Initial (stem) cells} directly surrounding QC produce clonal
cell files for most root tissues - \textbf{Root cap/columella:} Rapidly
replaced (slough off as root grows); amyloplast-containing statocytes
for gravitropism - \textbf{Casparian strip} in endodermal cells seals
apoplastic pathway for selective mineral transport

\begin{center}\rule{0.5\linewidth}{0.5pt}\end{center}

\subsection{Fruit Development, Vegetative Reproduction, and
Apomixis}\label{fruit-development-vegetative-reproduction-and-apomixis}

\subsubsection{Fruit Development: Hormonal Regulation, Parthenocarpy,
and Climacteric
Ripening}\label{fruit-development-hormonal-regulation-parthenocarpy-and-climacteric-ripening}

Fruit development is initiated by fertilisation and proceeds through
cell division, expansion, and ripening --- each phase under distinct
hormonal control.

\textbf{Phase 1: Fruit set (post-fertilisation initiation):} Auxin from
developing seeds and gibberellins from the maternal pericarp suppress
the abscission programme that would otherwise drop unfertilised flowers.
Without seeds, \textbf{parthenocarpy} can be induced --- fruit
development without fertilisation: - Natural parthenocarpy: cultivated
banana (triploid, sterile), pineapple, navel orange, some \emph{Citrus}
and \emph{Vitis} cultivars - Induced parthenocarpy: exogenous auxin
(NAA, 2,4-D) or GA (GA\(_3\)) sprayed at anthesis. Used commercially in
seedless grapes, watermelons, and tomatoes - Genetic parthenocarpy:
\emph{pin} mutants (pinoid; auxin-overproducing) or \emph{fwf} mutants
(FRUIT WITHOUT FERTILIZATION; tomato) develop fruit without seeds

\textbf{Phase 2: Cell division and expansion:} Cytokinins drive early
cell division in the pericarp; auxin and GA drive subsequent cell
expansion. Cucumbers and watermelons reach mature size by predominantly
cell expansion (each cell can balloon 30-fold).

\textbf{Phase 3: Ripening --- climacteric vs non-climacteric fruits}

The most studied transition is in \textbf{climacteric fruits} (banana,
tomato, apple, avocado, peach, mango), characterised by an
\textbf{autocatalytic burst of ethylene production} at the onset of
ripening:

\begin{equation}
\text{Methionine} \xrightarrow{\text{SAM synthetase}} \text{SAM} \xrightarrow{\text{ACC synthase (ACS)}} \text{ACC} \xrightarrow{\text{ACC oxidase (ACO), O}_2} \text{Ethylene}
\label{eq:unit_VIII_ethylene_pathway}
\end{equation}

\textbf{Two systems of ethylene biosynthesis:}

\begin{itemize}
\tightlist
\item
  \textbf{System 1:} Basal, auto-inhibited ethylene production occurring
  in vegetative tissues and pre-climacteric fruit. Low ethylene levels.
  Maintained at low rates throughout development.
\item
  \textbf{System 2:} Auto-catalytic, ripening-associated ethylene burst.
  Activated by transcription factors (RIN/MADS-RIN, NOR, CNR) that turn
  on a ripening-specific ACS isoform (ACS2) and ACO1. Once initiated,
  the ethylene produced positively feeds back on its own biosynthesis
  (auto-catalysis), generating an exponential rise. Ethylene then
  orchestrates ripening.
\end{itemize}

\textbf{Ripening events under ethylene control:}

\begin{itemize}
\tightlist
\item
  \textbf{Cell wall softening:} Polygalacturonase (PG), pectin
  methylesterase (PME), expansins, \texorpdfstring{\ensuremath{\beta}}{beta}-galactosidase degrade cell wall
  polysaccharides; tissue softens
\item
  \textbf{Sugar accumulation:} Starch \texorpdfstring{\ensuremath{\to}}{->} soluble sugars via amylases;
  sucrose conversion via invertase. Banana goes from 1\% to 18\% sugar
  content during ripening
\item
  \textbf{Pigment changes:} Chlorophyll degradation (Stay-Green protein,
  chlorophyllase) reveals carotenoids (red/orange) and anthocyanins
  (red/purple)
\item
  \textbf{Volatile production:} Esters (banana: isoamyl acetate),
  lactones, terpenes --- aroma compounds that attract dispersers
\item
  \textbf{Acid metabolism:} Malic and citric acids decline, sometimes
  through respiratory consumption
\end{itemize}

\textbf{Non-climacteric fruits} (strawberry, grape, citrus, pineapple,
cherry) do not show the ethylene burst; ABA and auxin appear to
substitute as the master ripening signals. Strawberries in particular
are ABA-responsive: applying ABA accelerates ripening; inhibiting ABA
biosynthesis (NDGA) blocks ripening.

\begin{quote}
\textbf{Clinical Connection:} The commercial fruit industry manipulates
ethylene extensively. Bananas are picked green, shipped under
ethylene-suppressed conditions (using KMnO\(_4\) as an ethylene scrubber
or 1-MCP --- 1-methylcyclopropene --- as an ethylene perception
inhibitor that competitively blocks ETR1), and then ripened on demand by
ethylene gas treatment at distribution centres. The Flavr Savr tomato
(1994) was the first commercial GMO food, with antisense PG suppressing
wall softening; ethylene-resistant tomatoes (silenced ACS or ACO)
preserve fruit quality during long-distance transport.
\end{quote}

\subsubsection{Fruit Types and
Dispersal}\label{fruit-types-and-dispersal}

\begin{figure}[htbp]
\centering
\includegraphics[width=0.82\linewidth,height=4.2in,keepaspectratio]{../figures/mermaid_inline/inline_mermaid_0146_e58d2396aa32.png}
\caption{Fruit development and classification The ovary wall develops into the pericarp (exocarp, mesocarp, endocarp), which differentiates into diverse fruit types adapted for specific dispersal strategies.}
\end{figure}

A \textbf{fruit} is the mature ovary wall (\textbf{pericarp}: exocarp +
mesocarp + endocarp) often incorporating accessory tissues:

{\def\LTcaptype{none} % do not increment counter
\begin{longtable}[]{@{}
  >{\raggedright\arraybackslash}p{(\linewidth - 6\tabcolsep) * \real{0.2174}}
  >{\raggedright\arraybackslash}p{(\linewidth - 6\tabcolsep) * \real{0.2174}}
  >{\raggedright\arraybackslash}p{(\linewidth - 6\tabcolsep) * \real{0.1522}}
  >{\raggedright\arraybackslash}p{(\linewidth - 6\tabcolsep) * \real{0.4130}}@{}}
\toprule\noalign{}
\begin{minipage}[b]{\linewidth}\raggedright
Fruit type
\end{minipage} & \begin{minipage}[b]{\linewidth}\raggedright
Structures
\end{minipage} & \begin{minipage}[b]{\linewidth}\raggedright
Example
\end{minipage} & \begin{minipage}[b]{\linewidth}\raggedright
Dispersal mechanism
\end{minipage} \\
\midrule\noalign{}
\endhead
\bottomrule\noalign{}
\endlastfoot
Drupe & Fleshy mesocarp + hard endocarp (stone) & Cherry, mango, olive,
coconut & Endozoochory (animal ingestion); ocean for coconut \\
Berry (true) & Entire pericarp fleshy & Tomato, grape, blueberry,
capsicum, banana & Endozoochory \\
Achene & Dry, indehiscent; one seed; pericarp free from seed coat &
Sunflower, dandelion (+ pappus) & Wind (dandelion); gravity \\
Samara & Achene-like + wing of pericarp tissue & Maple (\emph{Acer}),
ash (\emph{Fraxinus}) & Wind (autorotation; helicoid flight) \\
Legume & Dehiscent pod; two valves split along two sutures & Pea, bean,
soybean & Explosive hygroscopic dehiscence \\
Capsule & Dry; multiple septa; dehisces by pores or valves & Cotton,
poppy, \emph{Arabidopsis} & Wind-shaken; gravity \\
Bur & Achene with hooks or spines & Burdock (\emph{Arctium}), cocklebur
(\emph{Xanthium}) & Epizoochory (animal fur/feathers) \\
\end{longtable}
}

\subsubsection{Vegetative Reproduction: Mechanisms and Practical
Applications}\label{vegetative-reproduction-mechanisms-and-practical-applications}

Plants have diverse mechanisms for \textbf{vegetative (asexual)
reproduction}, each with distinct anatomical and ecological
characteristics that have been widely exploited in agriculture and
horticulture:

{\def\LTcaptype{none} % do not increment counter
\begin{longtable}[]{@{}
  >{\raggedright\arraybackslash}p{(\linewidth - 6\tabcolsep) * \real{0.1800}}
  >{\raggedright\arraybackslash}p{(\linewidth - 6\tabcolsep) * \real{0.2200}}
  >{\raggedright\arraybackslash}p{(\linewidth - 6\tabcolsep) * \real{0.1600}}
  >{\raggedright\arraybackslash}p{(\linewidth - 6\tabcolsep) * \real{0.4400}}@{}}
\toprule\noalign{}
\begin{minipage}[b]{\linewidth}\raggedright
Mechanism
\end{minipage} & \begin{minipage}[b]{\linewidth}\raggedright
Description
\end{minipage} & \begin{minipage}[b]{\linewidth}\raggedright
Examples
\end{minipage} & \begin{minipage}[b]{\linewidth}\raggedright
Practical applications
\end{minipage} \\
\midrule\noalign{}
\endhead
\bottomrule\noalign{}
\endlastfoot
\textbf{Stolons (runners)} & Horizontal above-ground stems that root and
form daughter plants at nodes & Strawberry (\emph{Fragaria}), spider
plant (\emph{Chlorophytum}) & Commercial strawberry production: each
``mother'' plant produces 5--10 daughter plants per season; entire
fields propagated clonally from elite cultivars \\
\textbf{Rhizomes} & Horizontal underground stems with nodes,
scale-leaves, and adventitious roots & Ginger (\emph{Zingiber}), bamboo,
iris, bracken fern, turmeric & Survives fire, frost, and grazing --- one
rhizome system can persist for centuries (single bracken stand
\textgreater1000 years; \emph{Pando} aspen clone in Utah,
about 80,000 years old). Commercial ginger and turmeric grown
entirely from rhizome divisions \\
\textbf{Bulbs} & Modified shoots with fleshy storage leaves on a short
basal plate; central apical bud & Tulip, garlic, onion, daffodil, lily &
Year-round storage in cold rooms allows seasonal flower forcing (Dutch
tulip industry); garlic propagated entirely by clove (a single bulb
scale) \\
\textbf{Corms} & Solid swollen stem base; superficially bulb-like but no
fleshy storage leaves & Crocus, gladiolus, taro (\emph{Colocasia}) &
Annual replacement: each corm produces ``cormels'' used for the next
planting \\
\textbf{Tubers} & Swollen subterranean stem (stem tubers, e.g., potato)
or root (root tubers, e.g., sweet potato) with axillary buds (eyes) &
Potato (\emph{Solanum tuberosum}), Jerusalem artichoke, sweet potato &
Each ``eye'' can regenerate a complete plant. Global potato cultivation
is essentially clonal --- major cultivars (Russet Burbank, Yukon Gold)
are genetically uniform clones. This uniformity is also their
vulnerability (e.g., the 1845 Irish potato famine, in which a single
\emph{Phytophthora infestans} genotype destroyed a near-monoclonal
crop) \\
\textbf{Adventitious roots} & Roots forming on stem cuttings, leaves, or
other non-root tissue & Willow (\emph{Salix}) cuttings root readily;
many succulents form leaves & Foundation of horticultural cuttings:
roses, grape, citrus, blueberry are commercially propagated by stem
cuttings under mist with rooting hormone (auxin) treatment \\
\textbf{Adventitious plantlets} & Plantlets form spontaneously on leaves
or other organs & \emph{Kalanchoe} (leaf margins), \emph{Bryophyllum} &
Drop from parent and root in soil --- example of natural clonal
propagation \\
\textbf{Fragmentation} & Twig or stem fragments that detach and root &
Willows, many aquatic plants (\emph{Elodea}, \emph{Myriophyllum}) &
Aggressive invasive species often spread by fragmentation (e.g.,
\emph{Elodea canadensis} in European waterways) \\
\textbf{Layering} & Stem touches soil and roots while still attached to
parent & Blackberry, raspberry; horticultural air layering & Common in
commercial production of difficult-to-root species; air layering used
for tropical fruit trees and ornamentals \\
\end{longtable}
}

\textbf{Grafting} (not strictly vegetative reproduction but related):
Joining a scion (desired cultivar) onto a rootstock (provides root
system, often with disease/drought resistance). Used almost universally
in apple, pear, citrus, grape, and stone fruit cultivation. The
19th-century \emph{Phylloxera} aphid pandemic destroyed European
vineyards; the rescue was grafting European \emph{Vitis vinifera} scions
onto American \emph{V. labrusca} rootstocks (which carry natural
resistance) --- every European wine grape today is grown on American
roots.

\subsubsection{Apomixis: Sporophytic and Gametophytic
Pathways}\label{apomixis-sporophytic-and-gametophytic-pathways}

\textbf{Apomixis} is seed formation without fertilisation, producing
genetically identical maternal clones in a seed package. Apomixis is a
powerful evolutionary ``frozen'' genotype dispersal mechanism and the
holy grail of crop breeding.

\textbf{Types of apomixis:}

\begin{enumerate}
\def\labelenumi{\arabic{enumi}.}
\item
  \textbf{Sporophytic (adventitious) apomixis (adventive embryony):}
  Embryos arise directly from sporophytic (somatic) cells of the ovule
  (typically the \textbf{nucellus} or \textbf{integuments}) ---
  bypassing the gametophyte entirely. The embryo is genotypically
  identical to the maternal parent (2n, same as mother). Multiple
  embryos may form per seed (\textbf{polyembryony}). Example: citrus
  (most cultivars produce both zygotic and apomictic embryos in the same
  seed; nucellar embryos eventually outcompete the zygotic embryo,
  preserving the elite genotype). Mango polyembryonic cultivars use the
  same mechanism.
\item
  \textbf{Gametophytic apomixis:} A diploid embryo sac forms without
  meiosis, retaining the maternal genome unreduced. Two sub-types:

  \begin{itemize}
  \tightlist
  \item
    \textbf{Apospory:} Embryo sac arises from a somatic nucellar cell
    that does not undergo meiosis. The MMC degenerates without producing
    megaspores; instead, an adjacent nucellar cell expands and divides
    mitotically to form an unreduced (2n) embryo sac. Examples:
    \emph{Hieracium} (hawkweed), \emph{Pennisetum} (pearl millet
    relatives), \emph{Hypericum}
  \item
    \textbf{Diplospory:} Megaspore mother cell skips or modifies meiosis
    (replacing meiosis with mitosis; ``MiMe''), producing an unreduced
    megaspore that develops into a 2n embryo sac. Examples: dandelion
    (\emph{Taraxacum}; many dandelion species are obligate apomicts),
    \emph{Tripsacum} (relative of maize), some \emph{Boechera}
  \item
    In both, the egg cell is unreduced (2n) and develops
    parthenogenetically into a 2n embryo (matching the mother)
  \item
    \textbf{Endosperm formation in gametophytic apomicts:}

    \begin{itemize}
    \tightlist
    \item
      \textbf{Pseudogamous:} Central cell still requires fertilisation
      by sperm to form endosperm (typical in apomictic grasses like
      \emph{Poa}, \emph{Pennisetum})
    \item
      \textbf{Autonomous:} Central cell develops into endosperm without
      fertilisation (typical of dandelions and \emph{Hieracium})
    \end{itemize}
  \end{itemize}
\end{enumerate}

\textbf{Agricultural significance --- clonal seed and the holy grail of
breeding:}

The dream of apomixis in major crops: capture the \textbf{hybrid vigour
(heterosis)} of an elite F1 hybrid in a self-perpetuating apomictic
seed, eliminating the need to re-create the hybrid each generation.
Currently, F1 hybrid maize, rice, and sorghum require expensive annual
production crosses (typically requiring about 2--3 isolated
hectares of cytoplasmic male sterile mother lines and pollen donor
lines). If apomixis could be engineered into rice or maize, subsistence
farmers could save F1 seed indefinitely and continue to capture hybrid
vigour.

\textbf{Synthetic apomixis breakthrough:} Synthetic apomixis studies in
rice combine apomeiosis with parthenogenesis triggers to produce clonal
seed from hybrid plants. High-frequency systems in hybrid rice reported
many lines with \textgreater80\% clonal seeds and selected lines above
95\%, while subsequent transgenerational work found largely stable
clonal inheritance with rare aneuploid progeny that still require
monitoring
\citep{vernet2022highfrequencyapomixis,liu2023syntheticapomixis}. The
mechanism combines three modifications:

\begin{enumerate}
\def\labelenumi{\arabic{enumi}.}
\tightlist
\item
  \textbf{MiMe} (Mitosis instead of Meiosis): triple knockout of
  \emph{OSD1}, \emph{PAIR1}, \emph{REC8} converts meiosis into mitosis,
  producing unreduced gametes
\item
  \textbf{MTL} (MATRILINEAL): pollen-specific phospholipase knockout
  enabling haploid induction --- sperm cannot complete fusion with the
  egg
\item
  \textbf{BBML} or \textbf{DMP} transgenes inducing parthenogenesis ---
  egg cells initiate embryogenesis without sperm contribution
\end{enumerate}

The combination yields rice plants that produce maternal clones in seed
form across multiple generations --- a foundational step toward
apomictic crops at scale. Field translation remains conditional:
breeders must quantify seed-set frequency, transgene segregation, rare
chromosome loss, fitness under local climates, containment of gene flow
to wild relatives, and whether farmers would gain durable access rather
than new seed-locking arrangements.

\textbf{Risks to wild relatives:} Apomictic hybrids could become
invasive (since each seed is a perfect copy of the elite plant) and
could erode genetic diversity in wild populations through escapes.
Regulatory caution is warranted.

\subsubsection{Plant Biotechnology}\label{plant-biotechnology}

Modern plant biotechnology builds on understanding of plant
reproduction:

\begin{itemize}
\tightlist
\item
  \textbf{Tissue culture and micropropagation:} Totipotent plant cells
  can regenerate whole organisms. The auxin:cytokinin ratio determines
  differentiation: high auxin promotes roots; high cytokinin promotes
  shoots; balanced ratio maintains callus.
\item
  \textbf{Somatic embryogenesis:} Somatic cells (non-reproductive) can
  be induced to form embryo-like structures that develop into complete
  plants. Used commercially for oil palm, conifers, coffee.
\item
  \textbf{Agrobacterium tumefaciens:} Natural genetic engineer. Its Ti
  \hyperref[gl:plasmid]{\textbf{plasmid}} contains T-DNA that integrates
  into the plant \hyperref[gl:genome]{\textbf{genome}}. Modified T-DNA
  (with antibiotic resistance and gene of interest replacing oncogenes)
  is the primary tool for plant transformation.
\item
  \textbf{Applications:} Bt crops (insecticidal crystal protein gene
  from \emph{Bacillus thuringiensis}), golden rice (beta-carotene
  biosynthesis genes), drought-tolerant cultivars (DREB/CBF
  overexpression), herbicide-resistant varieties (modified EPSPS gene).
\end{itemize}

\begin{quote}
\textbf{Clinical Connection:} The Bt toxin proteins (Cry proteins) are
harmless to mammals because they require alkaline gut
\hyperref[gl:ph]{\textbf{pH}} (found in insect midguts but not mammalian
stomachs) to become active. Bt crops have reduced insecticide use by
37\% globally while increasing yields 22\% in developing countries
(meta-analysis, Klumper \& Qaim, 2014).
\end{quote}

\textbf{Concept Check 6:} A breeder wishes to create a triploid (3n)
seedless watermelon for commercial sale. Outline the crossing scheme
starting from diploid (2n) parents. Why are triploid offspring sterile
(no viable seeds), and which cross --- 2n \texorpdfstring{\ensuremath{\times}}{x} 4n or 4n \texorpdfstring{\ensuremath{\times}}{x} 2n --- yields the
best fruit set?

\textbf{Concept Check 7:} \emph{Pyrus communis} (European pear) requires
a pollinator of a different cultivar to set fruit (pollinizer system)
due to gametophytic SI. A grower plants 100\% `Bartlett' (S\texorpdfstring{\textsubscript{1}}{1}S\texorpdfstring{\textsubscript{2}}{2}) pears in
an orchard. Predict the fruit yield. Then describe the genetic logic of
selecting a pollinizer cultivar --- what S-genotype is required?

\textbf{Concept Check 8:} Argue for or against this proposition:
``Apomixis is evolutionarily dead-end because it eliminates genetic
recombination. Hence apomictic species should be short-lived
geologically.'' Cite specific apomictic lineages in your answer (e.g.,
dandelions are apomictic and remain abundant on multiple continents).

\textbf{Concept Check 9:} A tomato cultivar bred for parthenocarpy
carries a constitutively active PIN1 transgene driving auxin
overproduction in the ovary. Predict fruit set in the absence of
pollination, and predict the dormancy and germination behaviour of any
seeds that do form (assuming some pollination occurs by accident).

\begin{center}\rule{0.5\linewidth}{0.5pt}\end{center}

\subsection{Current Evidence and Frontier
Biology}\label{current-evidence-and-frontier-biology-28}

For \textbf{Plant Reproduction and Development}, frontier biology
belongs inside the evidence logic of the chapter. Plant biology links
molecular regulation to climate stress, water limitation, crop
resilience, phenology, and ecosystem feedbacks. The core reading
question is this: plant reproduction links pollination, development,
genetics, phenology, dispersal, and environmental filtering.

\begin{itemize}
\tightlist
\item
  \textbf{What to verify:} identify the observation, model, assay, or
  dataset that would make the claim stronger or weaker.
\item
  \textbf{What to qualify:} state the scale, organism, cell type,
  environmental condition, or population where the claim is expected to
  hold.
\item
  \textbf{What to compare:} test at least one alternative explanation,
  baseline, or null model before treating the pattern as causal.
\item
  \textbf{What to cite:} distinguish primary evidence, review synthesis,
  public dataset, and institutional guidance; for recent or numeric
  claims, prefer the source closest to the measurement and state what
  has changed since it was published.
\end{itemize}

A strong plant explanation names the tissue, signal, environmental
driver, measurable trait, and tradeoff between growth, reproduction,
defence, and water use.

\textbf{Source practice:} For plant-stress and crop claims, name the
tissue, environmental driver, field context, and growth-reproduction
tradeoff; separate laboratory potential from agronomic adoption.

\subsubsection{Field Translation: Fertility, Phenology, and
Containment}\label{field-translation-fertility-phenology-and-containment}

Plant reproductive claims are especially sensitive to the step being
measured. Heat can reduce seed set by damaging tapetum development,
microspores, pollen viability, pollen-tube growth, or ovule/endosperm
development, and these stages differ in timing and temperature threshold
\citep{malik2022heatstressmale}. Phenology claims need the cue model:
spring ephemeral flowering in \emph{Corydalis ambigua} tracked snowmelt
and surface degree-days, while bumblebee emergence tracked soil warming;
seed set declined when flowers opened before bee emergence in a 19-year
field series and snow-removal experiment
\citep{kudo2019phenologicalmismatch}.

The same logic applies to gene drives and engineered reproductive
systems. A plant gene drive or apomictic trait is not just a molecular
construct; it is a population process shaped by dispersal, seed
dormancy, mating system, hybrid fitness, and governance. Responsible
evaluation uses staged containment, ecological risk assessment,
community engagement, and explicit reversibility or mitigation planning
before open release \citep{nasem2016genedrives}.

\subsection{Key Terms}\label{key-terms-28}

{\def\LTcaptype{none} % do not increment counter
\begin{longtable}[]{@{}
  >{\raggedright\arraybackslash}p{(\linewidth - 2\tabcolsep) * \real{0.2857}}
  >{\raggedright\arraybackslash}p{(\linewidth - 2\tabcolsep) * \real{0.7143}}@{}}
\toprule\noalign{}
\begin{minipage}[b]{\linewidth}\raggedright
Term
\end{minipage} & \begin{minipage}[b]{\linewidth}\raggedright
Definition
\end{minipage} \\
\midrule\noalign{}
\endhead
\bottomrule\noalign{}
\endlastfoot
\textbf{Alternation of generations} & Diploid sporophyte
(spore-producing) alternates with haploid gametophyte
(gamete-producing) \\
\textbf{Sporophyte} & \hyperref[gl:diploid]{\textbf{Diploid (2n)}}
generation; produces spores by meiosis; dominant in vascular plants \\
\textbf{Gametophyte} & Haploid (n) generation; produces gametes by
mitosis; dominant in bryophytes \\
\textbf{Sporopollenin} & Highly resistant polymer of pollen exine;
survives millions of years in sediment \\
\textbf{Tube cell} & Vegetative cell of pollen grain; forms the pollen
tube \\
\textbf{Generative cell} & Cell of pollen grain that divides to produce
two sperm cells \\
\textbf{Microsporogenesis} & Meiotic + mitotic divisions producing
pollen grain from microspore mother cell \\
\textbf{Megasporogenesis} & Meiotic divisions producing functional
megaspore from megaspore mother cell \\
\textbf{Synergid} & Two cells flanking the egg in the embryo sac;
secrete LURE peptides; regulate sperm release \\
\textbf{Filiform apparatus} & Cell-wall labyrinth at the synergid
micropylar pole; secretes LURE peptides \\
\textbf{Central cell} & Large embryo sac cell containing two polar
nuclei; gives rise to endosperm (3n; 2m:1p) \\
\textbf{Antipodal cells} & Three cells at the chalazal end of embryo
sac; nutritive function \\
\textbf{Double fertilisation} & Unique to angiosperms: sperm 1 + egg =
embryo (2n); sperm 2 + polar nuclei = endosperm (3n) \\
\textbf{LURE peptides} & Cysteine-rich defensin-like peptides from
synergids; species-specific pollen tube attractants \\
\textbf{FERONIA-LORELEI} & Receptor complex on synergid surface;
required for pollen tube rupture \\
\textbf{ABCDE model} & Model of floral organ identity; combinatorial
MADS-box TF activity specifies each whorl \\
\textbf{MADS-box} & DNA-binding domain in floral organ identity
transcription factors; binds CArG-box motifs \\
\textbf{AP1, AP3, PI, AG, SEP} & MADS-box gene families: A, B, B, C, E
classes respectively \\
\textbf{Floral quartet} & Tetrameric MADS-box complex specifying organ
identity by binding two CArG boxes \\
\textbf{Sporophytic SI} & Self-incompatibility where pollen
\hyperref[gl:phenotype]{\textbf{phenotype}} determined by diploid
tapetum; SRK-SCR system \\
\textbf{Gametophytic SI} & Self-incompatibility where pollen phenotype
determined by its own haploid genotype; S-RNase/SLF \\
\textbf{Polyploidy} & Whole-genome duplication; about 70\% of
angiosperms are paleopolyploids \\
\textbf{Autopolyploid} & Polyploid arising within a single species \\
\textbf{Allopolyploid} & Polyploid arising from interspecific
hybridisation followed by chromosome doubling \\
\textbf{Phyllotaxis} & Arrangement of leaves/organs on stem; golden
angle 137.5 degrees; Fibonacci spirals \\
\textbf{WUS-CLV3 circuit} & Negative feedback loop maintaining stem cell
population in the SAM \\
\textbf{Karrikins} & Smoke-derived butenolides; activate KAI2 receptor;
break seed dormancy after fire \\
\textbf{Apomixis} & Asexual reproduction through seed without
fertilisation; produces maternal clones \\
\textbf{Adventive embryony} & Sporophytic apomixis: embryos from
nucellar/integument cells \\
\textbf{Apospory} & Gametophytic apomixis with embryo sac from somatic
cell \\
\textbf{Diplospory} & Gametophytic apomixis with embryo sac from
modified meiosis \\
\textbf{MiMe} & Mitosis-instead-of-Meiosis triple knockout (OSD1, PAIR1,
REC8); basis of synthetic apomixis \\
\textbf{Endosperm} & Triploid (3n; 2m:1p) nutritive tissue in angiosperm
seeds; stores starch, protein, oils \\
\textbf{Free-nuclear endosperm} & Endosperm form with multiple nuclei in
shared cytoplasm before cellularisation \\
\textbf{PEG (Paternally Expressed Gene)} & Gene expressed primarily from
paternal allele in endosperm \\
\textbf{MEG (Maternally Expressed Gene)} & Gene expressed primarily from
maternal allele in endosperm \\
\textbf{FIS-PRC2} & Polycomb complex (MEA-FIS2-FIE-MSI1) regulating
endosperm imprinting \\
\textbf{Heterospory} & Production of two different spore sizes (micro-
and megaspores); precursor to seed habit \\
\textbf{Pericarp} & Mature ovary wall comprising exocarp, mesocarp, and
endocarp; forms the fruit \\
\textbf{Parthenocarpy} & Fruit development without fertilisation; can be
natural or induced by auxin/GA \\
\textbf{Climacteric fruit} & Fruit type with autocatalytic ethylene
burst at ripening (banana, tomato, apple) \\
\textbf{System 1/System 2 ethylene} & Basal vs auto-catalytic ripening
ethylene biosynthesis \\
\textbf{1-MCP} & Synthetic ethylene perception inhibitor used
commercially to delay fruit ripening \\
\textbf{Stratification} & Cold-moist treatment that breaks seed
dormancy; simulates winter \\
\textbf{Scarification} & Physical or chemical disruption of seed coat
enabling water uptake \\
\textbf{DOG1} & Master regulator of seed dormancy; QTL in
\emph{Arabidopsis} \\
\textbf{\hyperref[gl:totipotency]{Totipotency}} & Ability of a single
plant cell to regenerate an entire organism \\
\end{longtable}
}

\begin{center}\rule{0.5\linewidth}{0.5pt}\end{center}

\subsection{Review Questions}\label{review-questions-28}

\begin{enumerate}
\def\labelenumi{\arabic{enumi}.}
\item
  Compare the male gametophyte of a \emph{Pinus} tree to that of a
  \emph{Solanum lycopersicum} (tomato). How many cells in each? Where is
  mitosis 2 completed? How is sperm delivered to the egg? What is the
  evolutionary significance of the reduction in gametophyte size?
\item
  An \emph{Arabidopsis} triple mutant lacks functional A, B, and C
  domain transcription factors (abc mutant). Predict the identity of
  organs in each floral whorl. Then predict what would happen if, in the
  same mutant, you ectopically expressed B-class genes in whorls 1 and 2
  during flower development. Reference
  \cref{eq:unit_VIII_floral_quartet} in your answer.
\item
  Explain the mechanism of gametophytic self-incompatibility in
  Solanaceae. A researcher crosses two plants with S-genotypes
  \(S_1S_2\) and \(S_1S_3\). What fraction of pollen will be accepted?
  Which S-\hyperref[gl:allele]{\textbf{allele}} combinations appear in
  the viable offspring? Compare with the same scheme under sporophytic
  SI in \emph{Brassica}.
\item
  Compare primary and secondary seed dormancy. For a temperate tree
  species requiring both stratification and light for germination, trace
  the environmental cues from autumn seed fall through spring
  germination, linking each cue to its molecular mechanism (DOG1, ABA/GA
  balance, phytochrome). Identify the dormancy class (PD, PY, MD, MPD,
  combinational).
\item
  The goal of introducing synthetic apomixis into hybrid rice is to
  allow subsistence farmers to save F1 hybrid seed. (a) What genetic
  modifications are needed to create apomictic seed (cite Wang \emph{et
  al.} 2022)? (b) What risks might apomictic crops pose to wild
  relatives via gene flow? (c) Compare the genetic diversity
  implications of apomixis vs vegetative propagation.
\item
  Explain why double fertilisation \citep{nawaschin1898} is considered a
  key evolutionary innovation. What advantage does the triploid
  endosperm provide over the gymnosperm approach where the female
  gametophyte tissue serves as the nutritive tissue? Discuss the
  parental conflict implications via the PEG/MEG mechanism. Why is the
  canonical 2m:1p genome dosage common, and what kinds of exceptions
  would challenge an over-simple rule?
\item
  A fruit biologist examines a mystery fruit and finds: fleshy mesocarp,
  hard endocarp containing a single seed, and thin exocarp. Classify
  this fruit type. Name three plants that produce this type of fruit and
  describe the most likely dispersal mechanism. Is this likely a
  climacteric or non-climacteric fruit?
\item
  Describe the WUS-CLV3 feedback loop. Predict the phenotype of: (a) a
  loss-of-function \emph{clv3} mutant, (b) a gain-of-function \emph{WUS}
  overexpression line, and (c) a double mutant lacking both WUS and
  CLV3.
\item
  Run \breaktt{plant\_biomass\_growth} from 1 g to 40 g capacity over 60
  days. How does the curve differ from unrestricted exponential growth?
\item
  Why might \textbf{endosperm ploidy} (3n) stabilise parent-offspring
  conflict compared with purely maternal provisioning? Use the FIS-PRC2
  mechanism in your answer.
\item
  A breeder wants to create double-flowered carnations (extra petals).
  Which class of MADS-box gene should be partially inactivated and why?
  Cite the relevant Coen and Meyerowitz framework \citep{coen1991} and
  the molecular identity of the gene (AG).
\item
  Explain the commercial pre-harvest treatment with 1-MCP for apples in
  long-distance export. Which step of the ethylene pathway shown in
  \cref{eq:unit_VIII_ethylene_pathway} does 1-MCP block, and how does
  this preserve fruit firmness during shipping? Distinguish System 1 and
  System 2 ethylene biosynthesis.
\item
  A plant exhibits adventive embryony (sporophytic apomixis) in some
  seeds and zygotic embryos in others. From a single cross, predict the
  genotypic distribution of offspring and the agricultural value of this
  dual-mode reproduction.
\end{enumerate}

\begin{center}\rule{0.5\linewidth}{0.5pt}\end{center}

\subsection{Further Reading and Source
Notes}\label{further-reading-and-source-notes-23}

\begin{itemize}
\tightlist
\item
  \citet{coen1991} --- The war of the whorls: Genetic interactions
  controlling flower development. \emph{Nature}, 353.
\item
  \citet{nawaschin1898} --- Resultate einer Revision der
  Befruchtungsvorgänge bei \emph{Lilium martagon} und \emph{Fritillaria
  tenella}. \emph{Bulletin de l'Académie Impériale des Sciences de
  Saint-Pétersbourg}, 9.
\end{itemize}

\begin{center}\rule{0.5\linewidth}{0.5pt}\end{center}

\subsection{Computational Bridge}\label{computational-bridge-25}

Sporophyte growth phases are approximated logistically in
\breaktt{plant\_biomass\_growth}:

\begin{Shaded}
\begin{Highlighting}[]
\ImportTok{from}\NormalTok{ biology.botany }\ImportTok{import}\NormalTok{ plant\_biomass\_growth}

\NormalTok{out }\OperatorTok{=}\NormalTok{ plant\_biomass\_growth(}\FloatTok{1.0}\NormalTok{, }\FloatTok{0.15}\NormalTok{, }\FloatTok{80.0}\NormalTok{, }\FloatTok{40.0}\NormalTok{, steps}\OperatorTok{=}\DecValTok{40}\NormalTok{)}
\BuiltInTok{print}\NormalTok{(}\BuiltInTok{round}\NormalTok{(out.biomass\_g[}\OperatorTok{{-}}\DecValTok{1}\NormalTok{], }\DecValTok{2}\NormalTok{))}
\end{Highlighting}
\end{Shaded}

\begin{quote}
\textbf{Clinical / systems note:} Apomixis and clonal crops change
epidemiological risk profiles (uniform genetics) much like monocultures
in human agriculture policy.
\end{quote}

\begin{center}\rule{0.5\linewidth}{0.5pt}\end{center}

\subsection{Summary}\label{summary-28}

\begin{itemize}
\tightlist
\item
  \textbf{Alternation of generations:} Sporophyte (2n) produces spores
  by meiosis; gametophyte (n) produces gametes by mitosis. Evolutionary
  trend: sporophyte dominant, gametophyte reduced.
\item
  \textbf{Plant group diversity:} Bryophytes (gametophyte dominant,
  water-dependent fertilisation) through gymnosperms (pollen, seeds,
  wind pollination) to angiosperms (flowers, double fertilisation,
  fruits, animal pollination).
\item
  \textbf{Floral ABCDE model \citep{coen1991}:} A+E = sepal; A+B+E =
  petal; B+C+E = stamen; C+E = carpel; D+E = ovule. MADS-box TFs (AP1,
  AP3-PI heterodimer, AG, SEP1-4) form floral quartets binding two CArG
  boxes simultaneously.
\item
  \textbf{Gametophytes:} Male = 3-cell pollen grain (tube cell + 2
  sperm) inside sporopollenin exine; female = 7-cell, 8-nucleate
  Polygonum-type embryo sac (egg + 2 synergids + central cell with 2
  polar nuclei + 3 antipodals). Microsporogenesis: MMC \texorpdfstring{\ensuremath{\to}}{->} meiosis \texorpdfstring{\ensuremath{\to}}{->} 4
  microspores \texorpdfstring{\ensuremath{\to}}{->} asymmetric mitosis I \texorpdfstring{\ensuremath{\to}}{->} 2-cell pollen; mitosis II yields
  2 sperm. Megasporogenesis: MMC \texorpdfstring{\ensuremath{\to}}{->} meiosis \texorpdfstring{\ensuremath{\to}}{->} 1 functional + 3
  degenerating megaspores \texorpdfstring{\ensuremath{\to}}{->} 3 free nuclear divisions \texorpdfstring{\ensuremath{\to}}{->} 8 nuclei \texorpdfstring{\ensuremath{\to}}{->}
  cellularisation \texorpdfstring{\ensuremath{\to}}{->} 7-cell embryo sac.
\item
  \textbf{Double fertilisation \citep{nawaschin1898}:} LURE-LRE-FER
  signalling triggers pollen tube rupture in synergid; sperm 1 + egg =
  2n zygote \texorpdfstring{\ensuremath{\to}}{->} embryo; sperm 2 + 2 polar nuclei = 3n primary endosperm.
  Pollen tube guidance: tip-focused Ca\(^{2+}\) gradient + LURE
  chemotropism; spermatogenesis during tube elongation in bicellular
  pollen species.
\item
  \textbf{Endosperm (typically 3n; 2m:1p):} Free-nuclear \texorpdfstring{\ensuremath{\to}}{->} cellular
  development; PEG/MEG imprinting via FIS-PRC2 drives parent-of-origin
  expression and parental-conflict regulation. Dosage is canonically
  2m:1p but varies in some lineages and developmental contexts.
\item
  \textbf{Self-incompatibility:} Sporophytic SI (SRK-SCR-ARC1;
  \emph{Brassica}) acts at stigma surface; gametophytic SI (S-RNase/SLF;
  Solanaceae, Rosaceae) acts in style by collaborative non-self
  recognition.
\item
  \textbf{Polyploidy:} Auto- vs allopolyploids; 70\% of angiosperms have
  paleopolyploid history; major crop genomes (wheat, cotton, strawberry)
  are polyploid.
\item
  \textbf{Fruit types:} Diverse adaptations for dispersal (endozoochory,
  anemochory, explosive dehiscence, epizoochory). Climacteric fruits
  (banana, tomato, apple) ripen via System 2 autocatalytic ethylene
  burst --- controlled commercially with 1-MCP. Parthenocarpy can be
  induced by auxin/GA (seedless grapes, watermelons).
\item
  \textbf{Seed dormancy:} Five classes (PD, PY, MD, MPD, combinational);
  broken by stratification, scarification, light (Pfr), smoke-karrikins
  (KAI2), or temperature fluctuation. Germination: GA degrades DELLA,
  releasing \texorpdfstring{\ensuremath{\alpha}}{alpha}-amylase transcription; ABA-DOG1-PP2C-SnRK2 axis maintains
  dormancy.
\item
  \textbf{Vegetative reproduction:} Stolons, rhizomes, bulbs, corms,
  tubers, adventitious roots/plantlets, fragmentation, layering,
  grafting. Foundation of clonal crop production (banana, potato,
  garlic, grape).
\item
  \textbf{Apomixis:} Sporophytic (adventive embryony, e.g., citrus
  nucellar) vs gametophytic (apospory in \emph{Hieracium}, diplospory in
  \emph{Taraxacum}; pseudogamous vs autonomous endosperm). Wang \emph{et
  al.} 2022 synthetic apomixis in rice via MiMe + MTL/BBML.
\item
  \textbf{Meristems:} SAM (WUS-CLV3 feedback, PIN1-auxin phyllotaxis at
  golden angle); RAM (QC + initials). Both maintain lifelong growth.
\item
  \textbf{Biotechnology:} Tissue culture exploits totipotency;
  \emph{Agrobacterium} Ti plasmid enables genetic transformation; Bt
  crops, golden rice.
\item
  \textbf{Connections:} See Unit IV for meiosis and life cycle,
  \cref{sec:unit_VIII_plant_structure_and_water} for water and growth,
  and \cref{sec:unit_III_photosynthesis} for photosynthate partitioning.
\end{itemize}

\begin{center}\rule{0.5\linewidth}{0.5pt}\end{center}

\begin{center}\rule{0.5\linewidth}{0.5pt}\end{center}

\subsubsection{Companion Source
Module}\label{companion-source-module-28}

\textbf{Plant Reproduction and Development} should leave a reproducible
trail from a biological claim to the code, figure, diagram, or
paper-based activity that can test it. Use the surfaces below to inspect
the chapter's assumptions, rerun the relevant model, or compare the
manuscript explanation with companion labs and figures.

{\def\LTcaptype{none} % do not increment counter
\begin{longtable}[]{@{}
  >{\raggedright\arraybackslash}p{(\linewidth - 2\tabcolsep) * \real{0.5000}}
  >{\raggedright\arraybackslash}p{(\linewidth - 2\tabcolsep) * \real{0.5000}}@{}}
\toprule\noalign{}
\begin{minipage}[b]{\linewidth}\raggedright
Surface
\end{minipage} & \begin{minipage}[b]{\linewidth}\raggedright
Use it for
\end{minipage} \\
\midrule\noalign{}
\endhead
\bottomrule\noalign{}
\endlastfoot
\breaktt{src/biology/botany/botany.py} (\breaktt{plant\_biomass\_growth})
& Explore growth allocation and reproductive tradeoffs. \\
\breaktt{src/biology/genetics/genetics.py} (\texttt{punnett\_square},
\texttt{chi\_squared\_test}) & Connect inheritance evidence to breeding
and reproductive outcomes. \\
\breaktt{src/mermaid/biology\_diagrams.py}
(\breaktt{hormone\_signaling\_diagram}) & Link developmental timing to
hormone signalling. \\
\end{longtable}
}

\textbf{Reproducibility check:} state pollination mechanism,
developmental stage, genetic model, and environmental filter before
interpreting reproductive success. \textbf{Cross-reference:} connect
with \cref{sec:unit_V_mendelian_genetics} and
\cref{sec:unit_VIII_plant_responses}.

\newpage

\section{Plant Responses to the
Environment}\label{plant-responses-to-the-environment}

\label{sec:unit_VIII_plant_responses}

\begin{quote}
\textbf{Ch 27} \texorpdfstring{\ensuremath{\cdot}}{.} Level 2/3 \texorpdfstring{\ensuremath{\cdot}}{.} 55 min read \texorpdfstring{\ensuremath{\cdot}}{.} 75 min lecture \texorpdfstring{\ensuremath{\cdot}}{.}
Prerequisites: \cref{sec:unit_VIII_plant_reproduction}
\end{quote}

\subsection{Learning Objectives}\label{learning-objectives-29}

By the end of this chapter, you should be able to:

\begin{enumerate}
\def\labelenumi{\arabic{enumi}.}
\tightlist
\item
  Describe tropisms and how plants sense and respond to directional
  stimuli, including the molecular pathway for phototropism
  \citep{went1926}.
\item
  Explain the major plant \hyperref[gl:hormone]{\textbf{hormone}}s
  (auxin, gibberellin, cytokinin, ABA, ethylene, brassinosteroids,
  jasmonates, salicylate, strigolactones) and their signalling pathways.
\item
  Describe photomorphogenesis: phytochrome biochemistry (Pr/Pfr
  photoconversion), cryptochrome FAD signalling, and shade avoidance.
\item
  Explain gravitropism in roots and shoots, including amyloplast
  settling and PIN \hyperref[gl:protein]{\textbf{protein}}
  relocalisation.
\item
  Describe \hyperref[gl:photoperiodism]{\textbf{photoperiodism}}, the
  role of \hyperref[gl:phytochrome]{\textbf{phytochrome}}, florigen
  (FT), and the molecular basis of flowering time control via
  vernalisation (FLC silencing).
\item
  Explain the plant circadian clock molecular architecture (morning loop
  CCA1/LHY/TOC1; evening loop PRR5/7/9) and how it gates growth,
  flowering, and \hyperref[gl:photosynthesis]{\textbf{photosynthesis}}.
\item
  Describe drought (ABA pathway), cold (CBF-COR pathway), heat (HSP
  response), salt, flooding, and
  \hyperref[gl:herbivory]{\textbf{herbivory}} (jasmonate) stress
  responses.
\item
  Describe plant immunity (PTI, ETI), the hypersensitive response,
  systemic acquired resistance (SAR), and JA/SA antagonism.
\item
  Evaluate agricultural applications of plant hormone biology: ethylene
  inhibitors, gibberellin dwarfing genes, and auxin herbicides.
\end{enumerate}

\subsubsection{Study Blueprint}\label{study-blueprint-29}

\begin{itemize}
\tightlist
\item
  \textbf{Big idea:} Plants sense environmental signals and respond
  through growth, hormones, and physiological regulation.
\item
  \textbf{Core concepts:} tropisms, hormones, photoperiodism, stress
  responses.
\item
  \textbf{Framework alignment:} Vision \& Change: Structure and
  function, Pathways and transformations of energy and matter, Systems;
  AP Biology: Energetics, Systems Interactions; NGSS-style topics:
  Structure and Function, Matter and Energy in Organisms and Ecosystems.
\item
  \textbf{Model or quantitative lens:} Dose-response,
  water-use-efficiency, and hormone-interaction reasoning.
\item
  \textbf{Data skill:} Interpret plant response data across light,
  gravity, water, and hormone treatments.
\item
  \textbf{Practice cadence:} Visual Representations, Questions and
  Methods, Argumentation.
\item
  \textbf{Common misconception to repair:} A plant response is not
  passive; plants actively regulate development and physiology without
  neurons.
\item
  \textbf{Primary lab:} \cref{sec:lab_unit_VIII_plant_responses}.
\item
  \textbf{Question bank:} \cref{sec:q_unit_VIII_plant_responses}.
\item
  \textbf{Transfer task:} Transfer response logic to shade avoidance,
  drought, flowering, and crop management.
\item
  \textbf{Bridge to computation:}
  \breaktt{biology.botany.botany.photosynthesis\_rate}.
\end{itemize}

\begin{center}\rule{0.5\linewidth}{0.5pt}\end{center}

\begin{quote}
\textbf{Opening Vignette --- Darwin's Love Affair with Carnivorous
Plants}

``I care more about Drosera than the origin of species in the world,''
Darwin wrote to his friend Asa Gray in 1860. For fifteen years he
systematically experimented on sundews and Venus flytraps, weighing
morsels of meat to 1/78,000 of a grain, timing trap closure, and testing
with hundreds of chemical solutions --- establishing that carnivorous
plants actively respond to nitrogenous stimuli, not merely mechanical
touch. His 1875 book \emph{Insectivorous Plants} described hormone-like
chemical signals in plants decades before the word ``hormone'' existed.
In those experiments lurked the seed of plant hormone biology: auxin,
gibberellins, abscisic acid, cytokinin, ethylene --- the five classical
plant hormones whose discovery unfolded over the next century. Darwin
rarely found them by name, but his insight that plants sense and respond
to their environment through diffusible chemical signals laid much of
the groundwork that followed.
\end{quote}

\subsection{Plant Hormones Overview}\label{plant-hormones-overview}

Plant hormones (phytohormones) are small organic molecules produced in
low concentrations that regulate virtually every aspect of plant growth,
development, and stress responses. Unlike animal hormones, plant
hormones are often produced at many sites throughout the plant and can
act locally.

\begin{figure}[htbp]
\centering
\includegraphics[width=0.82\linewidth,height=4.2in,keepaspectratio]{../figures/mermaid_inline/inline_mermaid_0147_94efd40a8d97.png}
\caption{Overview of major plant hormone signalling pathways from biosynthesis through perception to cellular response Each pathway follows a general scheme: stimulus triggers hormone biosynthesis, the hormone is perceived by a specific receptor, and downstream signalling leads to gene expression changes.}
\end{figure}

\subsubsection{Auxin (Indole-3-Acetic Acid,
IAA)}\label{auxin-indole-3-acetic-acid-iaa}

\textbf{Biosynthesis:} Primarily from tryptophan via the TAA1/YUCCA
pathway. Produced in young leaves, shoot apical meristem, and developing
embryos.

\textbf{Transport:} Auxin moves by \textbf{polar auxin transport} (PAT),
unique among plant hormones: - \textbf{PIN efflux carriers} (PIN1-8):
Localised to specific faces of the plasma membrane; their polar
distribution determines the direction of auxin flow. PIN proteins cycle
between plasma membrane and endosomal compartments; their localisation
is regulated by phosphorylation (PID kinase) and vesicle trafficking. -
\textbf{AUX1/LAX influx carriers:} Import auxin into cells; less polar
than PIN proteins - Polar transport rate: about 1 cm/h in stems
(basipetal in shoots, acropetal in roots)

\textbf{Signalling mechanism (TIR1 pathway):} 1. Auxin binds TIR1/AFB
F-box proteins (nuclear receptors) 2. Auxin acts as ``molecular glue''
between TIR1 and Aux/IAA transcriptional repressors 3. Aux/IAA proteins
are ubiquitinated by SCF\(^{TIR1}\) complex and degraded by 26S
proteasome 4. \textbf{ARF (Auxin Response Factor)}
\hyperref[gl:transcription]{\textbf{transcription}} factors are released
from Aux/IAA repression 5. ARF activates auxin-responsive genes

\textbf{Key functions:} - \textbf{Apical dominance:} High auxin from the
apex suppresses lateral bud outgrowth; removal of apex (decapitation)
releases lateral buds - \textbf{Phototropism:} Asymmetric auxin
distribution bends shoots toward light. - \textbf{Gravitropism:}
Asymmetric auxin distribution in response to gravity - \textbf{Root
initiation:} High auxin promotes adventitious and lateral root formation
- \textbf{Fruit development:} Auxin from developing seeds promotes fruit
growth

\subsubsection{Gibberellins (GAs)}\label{gibberellins-gas}

\textbf{Biosynthesis:} From the MEP pathway in plastids; active forms
include GA\(_1\), GA\(_3\), GA\(_4\). Produced in young leaves, root
tips, and developing seeds.

\textbf{Signalling mechanism:} 1. GA binds the soluble receptor
\textbf{GID1} 2. GA-GID1 complex interacts with \textbf{DELLA proteins}
(transcriptional repressors: GAI, RGA, RGL1-3 in \emph{Arabidopsis}) 3.
DELLA proteins are ubiquitinated by SCF\(^{SLY1}\) E3 ligase and
degraded 4. Release of DELLA repression allows growth-promoting
transcription factors to act

\textbf{Key functions:} - \textbf{Stem elongation:} Promotes internodal
elongation; GA-deficient mutants are dwarf plants (basis of ``Green
Revolution'' semi-dwarf wheat: \emph{Rht} gene encodes GA-insensitive
DELLA) - \textbf{Seed germination:} GA induces \texorpdfstring{\ensuremath{\alpha}}{alpha}-amylase synthesis in
cereal aleurone layer, mobilising starch reserves - \textbf{Fruit
development:} Can induce parthenocarpy (seedless fruit production
without fertilisation; seedless grapes) - \textbf{Flowering:} In some
long-day plants, GA can substitute for the long-day photoperiod
requirement

\subsubsection{Cytokinins}\label{cytokinins}

\textbf{Biosynthesis:} From adenine derivatives via IPT
(isopentenyltransferase). Primarily synthesised in \textbf{root tips}
and transported to shoots via \hyperref[gl:xylem]{\textbf{xylem}}.

\textbf{Key functions:} - \textbf{Cell division:} Essential for
cytokinesis; cytokinin + auxin together drive cell proliferation -
\textbf{Shoot promotion:} Antagonise auxin's apical dominance; promote
lateral bud outgrowth - \textbf{Delay of senescence:} Maintain
\hyperref[gl:chlorophyll]{\textbf{chlorophyll}} and protein levels in
leaves; SAUR (senescence-associated ubiquitin-related) pathway -
\textbf{Tissue culture:} Auxin:cytokinin ratio determines organogenesis
-- high auxin = roots; high cytokinin = shoots; equal = callus
(undifferentiated cell mass)

\subsubsection{Abscisic Acid (ABA)}\label{abscisic-acid-aba}

\textbf{Biosynthesis:} From carotenoid cleavage (NCED3
\hyperref[gl:enzyme]{\textbf{enzyme}}, 9-\emph{cis}-epoxycarotenoid
dioxygenase). Produced in roots under soil drying and in leaves under
dehydration stress.

\textbf{Key functions:} - \textbf{Stomatal closure:} The primary drought
response hormone, detailed in the abscisic-acid discussion. -
\textbf{Seed dormancy:} Maintains dormancy until conditions favour
germination; antagonised by GA - \textbf{Stress gene expression:}
Activates LEA (Late Embryogenesis Abundant) protein genes; dehydrins;
osmolyte synthesis

\textbf{Signalling:} ABA binds PYR/PYL/RCAR receptors, which inhibit
PP2C phosphatases, activating SnRK2 kinases. SnRK2 phosphorylates SLAC1
(stomatal closure), AREB/ABF transcription factors (stress genes), and
ion channels.

\subsubsection{Ethylene}\label{ethylene}

\textbf{Biosynthesis:} From methionine via SAM (S-adenosylmethionine) to
ACC (1-aminocyclopropane-1-carboxylic acid) by ACC synthase, then to
ethylene by ACC oxidase. Ethylene is a \textbf{gas} -- it diffuses
freely between cells and even between plants.

\textbf{Key functions:} - \textbf{Fruit ripening:} Climacteric fruits
(banana, tomato, apple, avocado) show a burst of ethylene production
triggering ripening. Non-climacteric fruits (strawberry, grape, citrus)
do not show this burst. - \textbf{Senescence:} Promotes leaf yellowing
and abscission (leaf drop); petal senescence in cut flowers -
\textbf{Triple response:} In dark-grown seedlings, ethylene causes: (1)
shortened, thickened hypocotyl, (2) exaggerated apical hook, (3)
horizontal growth. This protects the seedling apex as it pushes through
soil. - \textbf{Flooding response:} Submerged roots accumulate ACC
(cannot oxidise to ethylene without O\(_2\)); when ACC reaches shoots,
aerenchyma formation is triggered (programmed cell death creates air
channels) - \textbf{Abscission:} Ethylene activates cellulase and
pectinase in the abscission zone at the leaf petiole base

\begin{quote}
\textbf{Clinical Connection:} The commercial fruit industry manipulates
ethylene extensively. Bananas are picked green, shipped under
ethylene-suppressed conditions (using KMnO\(_4\) as an ethylene scrubber
or 1-MCP as an ethylene perception inhibitor), and then ripened on
demand by ethylene gas treatment at distribution centres. This enables
global tropical fruit supply chains.
\end{quote}

\subsubsection{Jasmonates (JA)}\label{jasmonates-ja}

\textbf{Biosynthesis:} From linolenic acid (membrane lipid) via the LOX
(lipoxygenase) pathway. The active form is \textbf{JA-Ile}
(jasmonyl-isoleucine conjugate, produced by JAR1).

\textbf{Signalling:} JA-Ile binds \textbf{COI1} (F-box protein of the
SCF\textsuperscript{COI1} E3 ubiquitin ligase), which recruits
\textbf{JAZ} (Jasmonate ZIM-domain) repressor proteins into the SCF
complex:

\begin{equation}
\text{JA-Ile} + \text{COI1-JAZ} \xrightarrow{\text{SCF}^{\text{COI1}}} \text{Ub-JAZ} \xrightarrow{\text{26S proteasome}} \text{MYC2/3/4 activation}
\label{eq:plant_responses_1}
\end{equation}

Pre-wounding, JAZ proteins recruit \textbf{TOPLESS} corepressor via the
\textbf{NINJA} adaptor, forming a trimeric JAZ-NINJA-TOPLESS complex
that represses MYC2 through \hyperref[gl:histone]{\textbf{histone}}
deacetylase recruitment. Post-wounding JA-Ile surge displaces this
complex by directing SCF\textsuperscript{COI1} to ubiquitinate JAZ via
the \textbf{SLY degron} motif (Jas domain). 2024 cryo-EM structures
(\hyperref[gl:allosteric]{\textbf{allosteric}} lid closure model) show
COI1 undergoes a conformational switch upon JA-Ile binding, with
ASK1-CUL1-RBX1 forming the catalytic core (Sheard \emph{et al.},
\emph{Nature} 2010; Xu \emph{et al.}, \emph{bioRxiv} 2024). COI1
paralogs (COI1a/b) fine-tune JAZ subtype specificity across
defence/development contexts.

\textbf{Key functions:} - \textbf{Wound/herbivore defence:} Protease
inhibitors, alkaloids (nicotine in tobacco), glucosinolates, terpene
volatiles - \textbf{UV damage response:} Accumulation of UV-absorbing
flavonoids - \textbf{Mycorrhizal signalling:} JA promotes mycorrhizal
colonisation - \textbf{Pollen development:} Required for stamen
maturation and pollen release

\subsubsection{Salicylic Acid (SA)}\label{salicylic-acid-sa}

\textbf{Biosynthesis:} From phenylalanine (PAL pathway) or chorismate
(ICS pathway; \hyperref[gl:dominant]{\textbf{dominant}} in
\emph{Arabidopsis}).

\textbf{Key functions:} - \textbf{Systemic Acquired Resistance (SAR):}
After local pathogen infection, SA signals spread systemically, priming
distant tissues for enhanced defence in the plant-immunity response. -
\textbf{PR protein induction:} Pathogenesis-related proteins (PR-1,
\texorpdfstring{\ensuremath{\beta}}{beta}-1,3-glucanase, thaumatin-like proteins)

\textbf{Concept Check:} Aspirin (acetylsalicylic acid) is derived from
salicylic acid, originally extracted from willow bark (\emph{Salix}).
Why might plants produce SA in response to pathogens, and what does this
tell us about the evolutionary age of this defence pathway?

\subsubsection{Strigolactones}\label{strigolactones}

\textbf{Biosynthesis:} From carotenoid precursors via D27 isomerase and
CCD7/CCD8 cleavage. Produced primarily in roots.

\textbf{Key functions:} - \textbf{Branch suppression:} Inhibit lateral
bud outgrowth; work with auxin in apical dominance - \textbf{Mycorrhizal
recruitment:} Secreted into rhizosphere to attract AM fungi -
\textbf{Root architecture:} Promote lateral root elongation -
\textbf{Parasitic plant germination:} \emph{Striga} (witchweed) seeds
germinate in response to strigolactones from host roots, a devastating
crop \hyperref[gl:parasite]{\textbf{parasite}} in Africa

\begin{center}\rule{0.5\linewidth}{0.5pt}\end{center}

\subsection{Phototropism}\label{phototropism}

The molecular pathway for phototropism in \emph{Arabidopsis thaliana} is
one of the best-characterised plant signalling cascades:

\begin{enumerate}
\def\labelenumi{\arabic{enumi}.}
\tightlist
\item
  \textbf{Blue light} (390-500 nm) illuminates the shoot asymmetrically
\item
  \textbf{PHOT1} or \textbf{PHOT2} (LOV1/LOV2-domain photoreceptors; LOV
  = Light Oxygen Voltage) absorbs blue light. Photochemical covalent
  FMN-cysteine adduct forms, causing conformational change and
  autophosphorylation (Ser849 of PHOT1, kinase domain released).
\item
  PHOT1/2 phosphorylate \textbf{NPH3} (Non-Phototropic Hypocotyl 3).
  NPH3 interacts with ARF7/ARF19 and reorganises.
\item
  \textbf{Auxin lateral redistribution} (the Cholodny-Went hypothesis,
  quantified by \citet{went1926} \emph{Avena} coleoptile bioassay):

  \begin{itemize}
  \tightlist
  \item
    PHOT1/2 triggers \textbf{PIN3} (auxin efflux carrier) relocalisation
    from the bottom to the lateral face of shoot endodermal cells
  \item
    More auxin (IAA) is transported to the shaded side; less to the
    illuminated side
  \end{itemize}
\item
  \textbf{Acid growth hypothesis} \citep{hager1971}:

  \begin{itemize}
  \tightlist
  \item
    High {[}IAA{]} on shaded side activates plasma membrane
    \textbf{H\(^+\)-ATPase} (AHA1/2). Proton extrusion acidifies the
    cell wall (\hyperref[gl:ph]{\textbf{pH}} 5.5 to 4.5).
  \item
    Acidic pH activates \textbf{expansins} (proteins that disrupt
    H-bonding between cellulose microfibrils and hemicellulose) causing
    wall loosening. Turgor-driven elongation follows.
  \end{itemize}
\item
  Differential elongation (shaded side elongates more than lit side)
  causes the shoot to bend toward light. Phototropic curvature is
  completed within 30-60 minutes.
\end{enumerate}

\begin{quote}
Historical note: Darwin and Darwin (\emph{Power of Movement in Plants},
1880) first demonstrated that the phototropic signal originates in the
coleoptile tip. \citet{went1926} isolated auxin using agar blocks. The
molecular receptors (PHOTs) were identified in 1997 by Briggs and
colleagues.
\end{quote}

\subsubsection{Worked Example: Phototropic Curvature from a Lateral
Auxin
Gradient}\label{worked-example-phototropic-curvature-from-a-lateral-auxin-gradient}

\textbf{Problem:} A coleoptile segment is illuminated from one side.
Following Cholodny-Went redistribution, the shaded flank receives more
auxin and, via acid growth, elongates 30\% while the illuminated flank
elongates 10\%. Each flank starts at length \(L_0 = 10\) mm and the
coleoptile diameter (distance between the two flanks) is \(w = 1.5\) mm.
Calculate the phototropic bending angle.

\textbf{Solution:}

\begin{enumerate}
\def\labelenumi{\arabic{enumi}.}
\item
  \textbf{Elongation of each flank.} Apply the relative growth fractions
  to the initial length:
  \[ \Delta s_{\text{shaded}} = L_0 \times 0.30 = 10 \times 0.30 = 3.0 \text{ mm}, \qquad \Delta s_{\text{lit}} = L_0 \times 0.10 = 10 \times 0.10 = 1.0 \text{ mm} \label{eq:unit_VIII_plant_responses_item_5}\]
\item
  \textbf{Differential elongation.} The bend is driven by the difference
  in arc length between the two flanks:
  \[ \Delta s_{\text{shaded}} - \Delta s_{\text{lit}} = 3.0 - 1.0 = 2.0 \text{ mm} \label{eq:unit_VIII_plant_responses_item_6}\]
\item
  \textbf{Bending angle.} When an organ of width \(w\) bends through
  angle \texorpdfstring{\ensuremath{\theta}}{theta} (radians), the arc-length difference between its two surfaces
  equals \(w\theta\), so
  \(\theta = (\Delta s_{\text{shaded}} - \Delta s_{\text{lit}})/w\):
  \[ \theta = \frac{2.0 \text{ mm}}{1.5 \text{ mm}} = 1.33 \text{ rad} = 1.33 \times \frac{180^\circ}{\pi} \approx 76.4^\circ \label{eq:unit_VIII_plant_responses_item_7}\]
\end{enumerate}

\textbf{Interpretation:} A 3:1 ratio of differential flank elongation
over a 1.5 mm-wide coleoptile produces about a 76 degree bend toward the
light --- quantitatively illustrating how a modest lateral auxin
gradient, amplified by acid-growth wall loosening, generates the large
directional curvature characteristic of phototropism.

\begin{center}\rule{0.5\linewidth}{0.5pt}\end{center}

\subsection{Gravitropism}\label{gravitropism}

\begin{figure}[htbp]
\centering
\includegraphics[width=0.82\linewidth,height=4.2in,keepaspectratio]{../figures/mermaid_inline/inline_mermaid_0148_234141e89d08.png}
\caption{Gravitropism in roots and shoots In both organs, amyloplast settling triggers asymmetric auxin redistribution via PIN protein relocalisation. The critical difference: root cells are INHIBITED by high auxin concentrations (causing the lower side to grow less, bending the root downward), while shoot cells are PROMOTED by high auxin (causing the lower side to grow more, bending the shoot upward).}
\end{figure}

\textbf{Key molecular details --- statoliths and PIN asymmetry:}

\begin{itemize}
\tightlist
\item
  \textbf{Statocytes:} Specialised gravity-sensing cells. In roots,
  these are the \textbf{columella cells} of the root cap (a few central
  cells just behind the meristem). In shoots, they are the
  \textbf{endodermal cells} (starch sheath surrounding the vascular
  cylinder).
\item
  \textbf{Amyloplasts (statoliths):} Starch-filled plastids (1--10 \texorpdfstring{\ensuremath{\mu}}{mu}m
  diameter) that sediment under gravity within seconds of reorientation.
  Their settling pulls on the actin cytoskeleton and ER membranes,
  generating mechanical signals that are transduced via \textbf{TOC
  (translocon of outer chloroplast) complex} to the LAZY1/IGT proteins
  on the new lower face of the cell.
\item
  \textbf{PIN auxin transporter asymmetry:}

  \begin{itemize}
  \tightlist
  \item
    In root columella cells, \textbf{PIN3} is initially symmetrically
    localised at the plasma membrane. Within minutes of gravity
    stimulation, PIN3 relocalises to the \textbf{new lower (now
    downward-facing) face} of the cell.
  \item
    In the elongation zone, \textbf{PIN2} is permanently localised on
    the upper face of cortical and epidermal cells; PIN2 is critical for
    transporting auxin from the root tip back along the upper side of
    the elongation zone toward the shoot. Relocalisation of PIN3 in
    columella cells biases auxin export toward the lower side, where
    PIN2 then carries it acropetally, accumulating auxin on the lower
    side of the elongation zone.
  \end{itemize}
\item
  \textbf{Lateral auxin redistribution:} In \emph{Arabidopsis} roots,
  the asymmetry develops within about 5 min of gravity
  stimulation; bending is visible by 20 min and complete within 1--2
  hours.
\item
  \textbf{LAZY1 protein:} Required for normal gravitropism. \emph{lazy1}
  mutants show reversed gravitropic responses (shoots grow down,
  branches grow at unusual angles). LAZY1 relocalises to the lower side
  of the cell after gravity stimulation and facilitates PIN
  relocalisation.
\item
  \textbf{Root vs shoot dose-response difference:} The same auxin
  redistribution (more auxin on the lower side) has opposite effects
  because root cells have a much \textbf{lower optimal auxin
  concentration} (about \(10^{-9}\) M for root elongation)
  than shoot cells (about \(10^{-5}\) M for stem elongation).
  Above the root optimum, auxin \textbf{inhibits} elongation via the
  stabilisation of Aux/IAA repressors and ethylene production; thus the
  lower side of the gravity-stimulated root grows less than the upper,
  and the root bends downward.
\end{itemize}

\begin{center}\rule{0.5\linewidth}{0.5pt}\end{center}

\subsection{Photomorphogenesis: Phytochromes and
Cryptochromes}\label{photomorphogenesis-phytochromes-and-cryptochromes}

Plants integrate diverse light signals --- wavelength, intensity,
direction, duration, timing --- through a constellation of
photoreceptors. The two most studied families are phytochromes
(red/far-red sensors) and cryptochromes (blue/UV-A sensors).

\subsubsection{Phytochrome Biochemistry}\label{phytochrome-biochemistry}

Phytochromes are dimeric about 125 kDa chromoproteins. Each
monomer carries a bilin chromophore (phytochromobilin) covalently
attached to a conserved cysteine in the photosensory module. The
chromophore exists in two photointerconvertible forms:

\begin{itemize}
\tightlist
\item
  \textbf{Pr (red-absorbing form, \(\lambda_{max} = 660\) nm):} Inactive
  form; synthesised in the dark; primarily cytoplasmic
\item
  \textbf{Pfr (far-red-absorbing form, \(\lambda_{max} = 730\) nm):}
  Active form; translocates to the nucleus where it interacts with PIFs
  (Phytochrome Interacting Factors); triggers physiological responses
\end{itemize}

\textbf{Photoconversion:}

\begin{equation}
\text{Pr} \xrightarrow{\text{red 660 nm}} \text{Pfr}
\qquad
\text{Pfr} \xrightarrow{\text{far-red 730 nm or dark reversion}} \text{Pr}
\label{eq:unit_VIII_phytochrome_photoconversion}
\end{equation}

Red light (660 nm) converts Pr to Pfr; far-red light (730 nm) converts
Pfr back to Pr. In darkness, Pfr slowly reverts to Pr (thermal
reversion, half-life about  hours) and is also degraded by the
COP1 ubiquitin-proteasome pathway. The R/FR reversibility of
physiological responses (Borthwick and Hendricks, 1952) is the
diagnostic signature of phytochrome involvement.

\textbf{Phytochrome diversity:}

\emph{Arabidopsis} has five phytochromes (PhyA--PhyE), each with
distinct light-stability and physiological roles:

\begin{itemize}
\tightlist
\item
  \textbf{PhyA:} Light-labile (degraded rapidly in light). Mediates
  very-low-fluence and far-red-high-irradiance responses (FR-HIR). PhyA
  is the principal sensor for \textbf{shade-avoidance under deep shade}
  (where FR predominates) and for \textbf{de-etiolation in dim light}.
\item
  \textbf{PhyB:} Light-stable. Mediates the classic R/FR reversible
  (low-fluence) responses. PhyB is the \textbf{principal sensor of
  canopy shade} via the \textbf{R:FR ratio} (high R:FR = open sun; low
  R:FR = under leaf canopy because chlorophyll absorbs red and transmits
  far-red). Loss-of-function \emph{phyB} mutants are constitutively
  shade-avoiding (elongated, pale, early-flowering).
\item
  \textbf{PhyC, PhyD, PhyE:} Specialised roles in flowering time, leaf
  angle, and modulating the responses of PhyA/B.
\end{itemize}

\textbf{Shade avoidance via R:FR ratio:} Sun-grown plants typically
experience R:FR ratios of about 1.2; under a leaf canopy, the
ratio drops to 0.1--0.4 because chlorophyll preferentially absorbs red.
PhyB Pfr is destabilised under low R:FR (returning to Pr); this releases
PIF4/5/7 from PhyB-mediated degradation, activating YUCCA-driven auxin
synthesis and triggering hypocotyl elongation, petiole elongation, and
accelerated flowering --- the \textbf{shade avoidance syndrome}.

\subsubsection{Cryptochrome Signalling --- Blue Light and
FAD}\label{cryptochrome-signalling-blue-light-and-fad}

\textbf{Cryptochromes (CRY1/CRY2)} are blue/UV-A photoreceptors with two
cofactors: a \textbf{FAD} (flavin adenine dinucleotide) chromophore for
blue light absorption (about 450 nm) and an MTHF antenna
pigment that broadens absorption.

\textbf{Photoactivation mechanism:} 1. Blue light reduces FAD from
oxidised (FAD\(_{ox}\)) to neutral semiquinone (FADH•) 2. Reduction
triggers conformational change in the C-terminal extension of CRY 3.
Activated CRY interacts with multiple downstream targets: - \textbf{COP1
(CONSTITUTIVE PHOTOMORPHOGENIC 1):} an E3 ubiquitin ligase; CRY binding
inhibits COP1, thereby stabilising COP1 substrates (HY5, HFR1, CO, LAF1)
--- bZIP and bHLH transcription factors that drive photomorphogenic gene
expression - \textbf{SPA proteins:} COP1 partners; CRY binding disrupts
COP1-SPA interaction - \textbf{PIFs:} CRY can also directly modulate PIF
activity

\textbf{Key blue-light responses:} - \textbf{Hypocotyl growth
inhibition:} Light-grown seedlings have short hypocotyls; \emph{cry1
cry2} double mutants show etiolated (long-hypocotyl) phenotype in blue
light despite normal red light response - \textbf{Anthocyanin
biosynthesis:} Blue light induces phenylpropanoid pathway genes (PAL,
CHS, CHI) producing UV-protective and pigment compounds -
\textbf{Photoperiodic flowering:} CRY2 stabilises CONSTANS protein in
long days - \textbf{Stomatal opening:} Cooperatively with phototropins
(covered in \cref{sec:unit_VIII_plant_structure_and_water})

\textbf{CRY-COP1-HY5 module:} The COP1/SPA E3 ubiquitin ligase complex
is the master ``darkness'' enforcer. In darkness, COP1-SPA ubiquitinates
HY5 and other photomorphogenic TFs, targeting them for degradation.
Light activates CRYs and PhyB, which physically interact with COP1 and
disrupt the COP1-SPA complex; HY5 accumulates and drives
photomorphogenic gene expression (chlorophyll biosynthesis, anthocyanin
synthesis, hypocotyl growth arrest). Loss-of-function \emph{cop1}
mutants are constitutively photomorphogenic in darkness --- confirming
the central role of this regulator.

\begin{center}\rule{0.5\linewidth}{0.5pt}\end{center}

\subsection{The Plant Circadian Clock --- Molecular
Architecture}\label{the-plant-circadian-clock-molecular-architecture}

Plants maintain an endogenous \textbf{about 24-hour oscillator}
(circadian clock) that anticipates dawn/dusk and gates downstream
processes including photosynthesis, growth, flowering, and stress
responses. The plant clock comprises interlocked transcriptional
feedback loops with morning, evening, and intermediate components.

\subsubsection{Morning Loop: CCA1/LHY Repressors and TOC1
Activator}\label{morning-loop-cca1lhy-repressors-and-toc1-activator}

\begin{itemize}
\tightlist
\item
  \textbf{CCA1 (Circadian Clock Associated 1) and LHY (Late Elongated
  Hypocotyl):} Single-MYB transcription factors expressed at dawn (peak
  about  ZT0, where ZT0 is dawn). They function as a
  heterodimer and act as \textbf{repressors} of evening-expressed genes
  (TOC1, ELF3, ELF4, LUX, GI).
\item
  \textbf{TOC1 (Timing of CAB expression 1; pseudo-response regulator
  PRR1):} Expressed at dusk (peak about  ZT12). TOC1 binds and
  represses the \emph{CCA1} and \emph{LHY} promoters, closing the
  negative feedback loop. TOC1 also activates morning genes after dusk
  through PIF interactions.
\end{itemize}

The morning loop is the ``central oscillator'': CCA1/LHY repress TOC1 in
the morning; as CCA1/LHY decay through the day, TOC1 rises; TOC1 then
represses CCA1/LHY through the night until they rise again at dawn.

\subsubsection{Evening Loop: PRR9 \texorpdfstring{\ensuremath{\to}}{->} PRR7 \texorpdfstring{\ensuremath{\to}}{->} PRR5 \texorpdfstring{\ensuremath{\to}}{->}
TOC1}\label{evening-loop-prr9-prr7-prr5-toc1}

A cascade of \textbf{pseudo-response regulators (PRRs)} sequentially
repress CCA1/LHY through the day:

\begin{itemize}
\tightlist
\item
  \textbf{PRR9:} Peaks about 4 h after dawn (ZT4); repressed by
  CCA1/LHY; represses CCA1/LHY transcription
\item
  \textbf{PRR7:} Peaks about 8 h after dawn (ZT8); repressed by
  PRR9 partially; represses CCA1/LHY
\item
  \textbf{PRR5:} Peaks at midday (ZT10--12); activates downstream
  targets including CO
\item
  \textbf{TOC1/PRR1:} Peaks at dusk (ZT12); closes morning loop
\end{itemize}

This PRR cascade creates a wave of repression sweeping across the day,
with each successive PRR repressing the previous one's targets. The
molecular result is a robust about 24 h oscillation despite
individual gene half-lives of primarily hours.

\subsubsection{Evening Complex:
ELF3-ELF4-LUX}\label{evening-complex-elf3-elf4-lux}

The \textbf{Evening Complex (EC)} comprises ELF3, ELF4, and LUX (LUX
ARRHYTHMO; a Myb-domain TF). ELF3 is a scaffold protein; LUX provides
DNA binding. The EC peaks at ZT12--16 and represses morning genes (PRR9,
PIF4, GI) --- gating early-night growth. \emph{elf3} mutants are
arrhythmic.

\subsubsection{Temperature Compensation}\label{temperature-compensation}

A defining property of the circadian clock is \textbf{temperature
compensation}: the period remains about 24 h across a wide
temperature range (10--30 °C) despite the underlying biochemical
reactions accelerating with temperature (Q\(_{10}\) about 
2). Mechanisms include: - \textbf{Alternative splicing of CCA1:} at
lower temperature, an alternative splice form (CCA1\texorpdfstring{\ensuremath{\beta}}{beta}) acts as a
repressor of CCA1\texorpdfstring{\ensuremath{\alpha}}{alpha}, slowing the loop - \textbf{Temperature-dependent
protein stability of PRR proteins:} PRR9/7 are more rapidly degraded at
higher temperatures - \textbf{JMJ30 demethylase:} counters temperature
effects on chromatin marks at clock loci

\subsubsection{Clock Outputs}\label{clock-outputs}

The clock gates numerous physiological processes by regulating
downstream gene expression: - \textbf{Stomatal opening:} Anticipatory
opening before dawn; controlled by clock-gated H\(^+\)-ATPase activity -
\textbf{Photosynthesis genes:} CAB (chlorophyll a/b binding protein),
Rubisco activase peak at dawn - \textbf{Flowering:} GI-CO-FT pathway is
clock-gated through photoperiodic hormone signalling. - \textbf{Growth:}
Hypocotyl elongation peaks in the late night/early morning (gated by
PIF4-driven YUCCA expression) - \textbf{Defence:} Many pathogen
resistance genes (e.g., \emph{PR1}) show circadian regulation,
anticipating bacterial infections that peak at midday

\begin{center}\rule{0.5\linewidth}{0.5pt}\end{center}

\subsection{Photoperiodism, Flowering Time, and
Vernalisation}\label{photoperiodism-flowering-time-and-vernalisation}

\subsubsection{Photoperiodic Categories}\label{photoperiodic-categories}

Plants are classified by their flowering response to day/night length:

\begin{itemize}
\tightlist
\item
  \textbf{Short-day plants} (SDPs): Flower when nights are
  \textbf{longer} than a critical length. Examples: chrysanthemum, rice,
  \emph{Pharbitis} (Japanese morning glory), tobacco (\emph{Nicotiana
  tabacum} Maryland Mammoth), soybean. Despite the name, it is
  \textbf{night length} that matters -- a brief flash of red light
  during the long night breaks the inductive dark period (the
  \textbf{night-break experiment}).
\item
  \textbf{Long-day plants} (LDPs): Flower when nights are
  \textbf{shorter} than a critical length. Examples: spinach,
  \emph{Arabidopsis}, barley, wheat, radish.
\item
  \textbf{Day-neutral plants:} Flower after sufficient vegetative growth
  regardless of photoperiod. Examples: tomato, cucumber, maize.
\end{itemize}

\subsubsection{Night-Length Measurement and Circadian
Gating}\label{night-length-measurement-and-circadian-gating}

\textbf{Phytochrome} is the primary photoreceptor for photoperiodism.
During a long night, Pfr levels decline below a threshold via thermal
reversion and proteolysis. In short-day plants, this decline is required
for flowering. A brief red-light pulse during the night converts Pr back
to Pfr, resetting the clock and preventing flowering. This \textbf{night
break} effect can be reversed by an immediately subsequent far-red light
pulse --- classic evidence for phytochrome involvement.

Photoperiod is measured not by simply timing dark intervals but by
\textbf{circadian gating} --- the clock provides a window during which
light is permissive for flowering. The molecular mechanism (the
``external coincidence model'') is: CO mRNA is rhythmically expressed by
the clock (peaks in late afternoon), but CO protein is stable primarily
in light. Coincidence of light + CO protein \texorpdfstring{\ensuremath{\to}}{->} FT activation \texorpdfstring{\ensuremath{\to}}{->} flowering.

\subsubsection{The CONSTANS-FT (Florigen)
Pathway}\label{the-constans-ft-florigen-pathway}

\begin{figure}[htbp]
\centering
\includegraphics[width=0.82\linewidth,height=4.2in,keepaspectratio]{../figures/mermaid_inline/inline_mermaid_0149_87b101c8c23e.png}
\caption{Molecular pathway of photoperiodic flowering in long-day plants The circadian clock gates CONSTANS (CO) protein stability so that in long days, CO is present in the light (when it is stabilised), activating FLOWERING LOCUS T (FT). FT protein (florigen) moves from leaves to the SAM, triggering flower development.}
\end{figure}

\textbf{Florigen --- FT as the mobile flowering signal:}

\begin{enumerate}
\def\labelenumi{\arabic{enumi}.}
\tightlist
\item
  In long-day plants, the circadian clock ensures that \textbf{CONSTANS
  (CO)} mRNA peaks in the late afternoon
\item
  CO protein is stabilised by light (blue light via CRY2; far-red light
  via PhyA) but degraded in darkness via COP1 ubiquitin ligase
\item
  In long days: CO mRNA peaks coincide with daylight, CO protein is
  stable, CO binds the \textbf{FT (FLOWERING LOCUS T)}
  \hyperref[gl:promoter]{\textbf{promoter}}, activating transcription
\item
  \textbf{FT protein} --- the long-sought ``florigen'' first
  hypothesised by Chailakhyan in 1937 --- is produced in
  \textbf{companion cells of leaf phloem} and moves systemically via
  sieve tubes to the SAM. FT is a small (about 20 kDa) globular
  protein structurally related to mammalian RAF kinase inhibitors; it is
  loaded into phloem via plasmodesmata and unloaded at the SAM.
\item
  At the SAM, FT interacts with \textbf{FD} (bZIP transcription factor)
  to activate \textbf{AP1} (floral meristem identity) and \textbf{SOC1}
  (flowering pathway integrator). The FT-FD-14-3-3 floral activation
  complex (FAC) was crystallographically resolved in 2011 (Taoka
  \emph{et al.}, \emph{Nature} 2011).
\item
  The vegetative meristem transitions irreversibly to a floral meristem
  --- the moment of ``evocation''
\end{enumerate}

\textbf{Diversity of FT homologues:} Most plants have multiple FT
paralogues with distinct functions. Some act as floral activators (FT in
\emph{Arabidopsis}, Hd3a in rice --- the rice florigen); some as floral
repressors (TFL1, BFT). The balance of activator and repressor FT-family
proteins fine-tunes flowering time.

\subsubsection{Vernalisation: Epigenetic Silencing of
FLC}\label{vernalisation-epigenetic-silencing-of-flc}

Prolonged cold exposure (weeks at 4 °C) promotes flowering in winter
annuals (\emph{Arabidopsis} ecotype Ler, accessions from northern
latitudes) and biennials (carrots, foxglove, sugar beet). The molecular
mechanism is a paradigmatic example of \textbf{epigenetic memory}:

\textbf{The FLC repressor (MADS-box):}

\textbf{FLC (FLOWERING LOCUS C)} is a MADS-box transcription factor that
represses \emph{FT} and \emph{SOC1}, blocking the floral transition. In
winter annuals before vernalisation, FLC is highly expressed, preventing
flowering. Vernalisation silences FLC, releasing FT/SOC1 and permitting
flowering.

\textbf{Mechanism of cold-induced FLC silencing --- the VRN2-PRC2
H3K27me3 mark:}

\begin{enumerate}
\def\labelenumi{\arabic{enumi}.}
\tightlist
\item
  Extended cold (\textgreater4 weeks at 4 °C) induces \textbf{VIN3}
  (VERNALISATION INSENSITIVE 3), a PHD-domain protein. VIN3 expression
  depends on a cold-memory element (CME) upstream of the \emph{VIN3}
  promoter; VIN3 transcript accumulates primarily after sustained cold
  (not pulses).
\item
  VIN3 joins \textbf{PRC2} (Polycomb Repressive Complex 2: catalytic
  subunit CLF/SWN, MSI1, FIE, VRN5) at the \textbf{FLC}
  \hyperref[gl:chromatin]{\textbf{chromatin}} via the \textbf{VRN2}
  subunit (a vernalisation-specific PRC2 partner).
\item
  PRC2 catalyses the trimethylation of histone H3 at lysine 27:
\end{enumerate}

\begin{equation}
\text{H3} \xrightarrow{\text{PRC2 (CLF/SWN, VRN2, FIE, MSI1, VIN3)}} \text{H3K27me3 at FLC}
\label{eq:unit_VIII_vernalisation}
\end{equation}

\begin{enumerate}
\def\labelenumi{\arabic{enumi}.}
\setcounter{enumi}{3}
\tightlist
\item
  H3K27me3 marks initially appear at a discrete nucleation region within
  the \emph{FLC} locus (a 1.5 kb cold-memory region spanning exon 1 and
  intron 1) during cold exposure; mark spreading across the entire
  \emph{FLC} locus occurs after return to warm.
\item
  \textbf{After return to warm:} VIN3 expression fades, but H3K27me3 is
  \textbf{maintained} by VRN2-PRC2 (which lacks VIN3 in the warm but
  retains the catalytic core), creating a \textbf{mitotically stable
  memory} of cold exposure (Finnegan \& Dennis, \emph{Plant J.} 2007;
  Sheldon \emph{et al.}, 2009).
\item
  \textbf{Mitotic inheritance:} As cells divide during spring growth,
  H3K27me3 is faithfully copied to daughter chromatin, propagating FLC
  silencing throughout the vernalised plant. The memory is gradually
  erased by demethylases in the offspring generation, resetting
  \emph{FLC} for the next winter.
\item
  \textbf{Result:} FLC (MADS-box repressor) can no longer repress
  \textbf{FT} and \textbf{SOC1} \texorpdfstring{\ensuremath{\to}}{->} flowering proceeds in spring
\end{enumerate}

\begin{equation}
\text{Cold} \xrightarrow{\text{VIN3}} \text{PRC2-VRN2} \xrightarrow{\text{H3K27me3 at FLC}} \downarrow \text{FLC} \rightarrow \uparrow \text{FT} \rightarrow \text{Flowering}
\label{eq:unit_VIII_vernalisation_summary}
\end{equation}

\textbf{Concept Check:} Why does vernalisation require weeks rather than
days of cold? What happens if prolonged cold is given in pulses of a few
days each, separated by warm intervals?

\begin{center}\rule{0.5\linewidth}{0.5pt}\end{center}

\subsection{Gravitropism, Touch Responses, and Other
Tropisms}\label{gravitropism-touch-responses-and-other-tropisms}

\subsubsection{Gravitropism --- Statolith and PIN Asymmetry (deep
dive)}\label{gravitropism-statolith-and-pin-asymmetry-deep-dive}

The detailed molecular mechanism of root gravitropism integrates
statolith sedimentation, PIN protein relocalisation, and the
dose-response asymmetry between root and shoot:

\begin{enumerate}
\def\labelenumi{\arabic{enumi}.}
\tightlist
\item
  \textbf{Gravity perception} by amyloplast sedimentation in columella
  cells (about 30 s after reorientation)
\item
  \textbf{Mechanical signal transduction} via actin cytoskeleton and ER
  membrane interactions; LAZY1 family proteins relocalise to the new
  lower face
\item
  \textbf{PIN3 asymmetric relocalisation} to the new lower face of
  columella cells (about 3--5 min)
\item
  \textbf{Lateral auxin transport} out the new lower face into the
  cortex below
\item
  \textbf{PIN2 acropetal transport} of auxin in the elongation zone,
  biased toward the lower side
\item
  \textbf{Differential elongation:} auxin inhibits elongation on the
  lower side (root); the upper side grows faster, bending the root
  downward
\item
  \textbf{Bending complete} in 1--2 hours; PIN proteins reset to
  symmetric distribution
\end{enumerate}

\subsubsection{Touch Responses
(Thigmomorphogenesis)}\label{touch-responses-thigmomorphogenesis}

Plants respond to mechanical stimuli --- wind, rain, animal contact ---
through \textbf{thigmomorphogenesis}, often resulting in shorter,
sturdier growth. The molecular pathway:

\begin{enumerate}
\def\labelenumi{\arabic{enumi}.}
\tightlist
\item
  \textbf{TCH (TOUCH) genes} --- first identified in \emph{Arabidopsis}
  by Janet Braam (1990). The four classical TCH genes are:

  \begin{itemize}
  \tightlist
  \item
    \textbf{TCH1 (CAM2):} calmodulin
  \item
    \textbf{TCH2 (CML24):} calmodulin-like
  \item
    \textbf{TCH3 (CML12):} calmodulin-like
  \item
    \textbf{TCH4 (XTH22):} xyloglucan endotransglucosylase/hydrolase
    (cell-wall modifier)
  \end{itemize}
\item
  \textbf{Mechanism:} Mechanical perturbation triggers Ca\(^{2+}\)
  influx via mechanosensitive channels (MCA1, MSL2, OSCA1; piezo-like).
  Cytosolic {[}Ca\(^{2+}\){]} rises within seconds.
\item
  \textbf{Calmodulin (TCH1) and CMLs (TCH2/3)} bind Ca\(^{2+}\) and
  translate the calcium signal to downstream targets --- kinases,
  transcription factors, cell-wall enzymes.
\item
  \textbf{Jasmonate production:} Mechanical wounding activates JA
  biosynthesis via the LOX pathway, mobilising defence responses in the
  herbivory and plant-immunity network.
\item
  \textbf{Phenotypic consequences:} Shorter internodes, thicker stems,
  increased lignin, smaller leaves. Brushed \emph{Arabidopsis} plants
  are 30--40\% shorter than unbrushed controls; the response is adaptive
  in windy environments.
\end{enumerate}

\textbf{Tendril coiling:} Climbing plants like peas, grapes, and
cucumbers produce \textbf{tendrils} that exhibit \textbf{thigmonastic}
coiling on contact with a support. The mechanism involves: -
Differential elongation: cells on the contact side stop elongating;
cells on the opposite side continue - Jasmonate signalling: JA mediates
the coiling response; \emph{Bryonia} tendrils require JA for coiling -
Within minutes of contact, tendril coils a complete 360° around the
support

\subsubsection{Other Tropisms}\label{other-tropisms}

\begin{itemize}
\tightlist
\item
  \textbf{Hydrotropism:} Roots grow toward moisture; mediated by ABA and
  MIZ1 (mizu-kussei 1) protein in root cap. \emph{miz1} mutants lack
  hydrotropism but retain gravitropism.
\item
  \textbf{Thigmotropism in roots:} Roots growing along a substrate
  exhibit obstacle avoidance, mediated by mechanosensitive channels.
\item
  \textbf{Heliotropism (solar tracking):} Sunflowers track the sun
  across the sky (East to West during the day; reverse to East at
  night). Driven by differential growth rates on east vs west sides of
  the stem, gated by the circadian clock.
\end{itemize}

\begin{center}\rule{0.5\linewidth}{0.5pt}\end{center}

\subsection{Abiotic Stress Responses}\label{abiotic-stress-responses}

\subsubsection{Drought -- ABA Signalling
Cascade}\label{drought-aba-signalling-cascade}

ABA is the primary drought stress hormone:

\begin{enumerate}
\def\labelenumi{\arabic{enumi}.}
\tightlist
\item
  Soil water deficit detected. Root ABA synthesis increases (NCED3
  upregulated). ABA moves via xylem to leaves.
\item
  \hyperref[gl:guard-cell]{\textbf{Guard cell}}: ABA binds
  \textbf{PYR/PYL/RCAR receptors} (14 paralogs in \emph{Arabidopsis})
\item
  ABA-PYR complex inhibits \textbf{PP2C phosphatases} (ABI1, ABI2, HAB1)
\item
  \textbf{SnRK2 kinases} (especially OST1/SnRK2.6) are freed from PP2C
  inhibition and phosphorylate:

  \begin{itemize}
  \tightlist
  \item
    \textbf{SLAC1} anion channel: Cl\(^-\) efflux depolarises guard cell
    membrane
  \item
    \textbf{GORK} K\(^+\) outward-rectifying channel opens: K\(^+\)
    efflux
  \item
    Water follows osmotically: guard cells shrink,
    \hyperref[gl:stomata]{\textbf{stomata}} close
  \item
    \textbf{AREB/ABF} transcription factors: activate drought-responsive
    genes (dehydrins, LEA proteins, osmolyte synthesis)
  \end{itemize}
\end{enumerate}

\textbf{Water use efficiency (WUE):}

\begin{equation}
WUE = \frac{A_{net}}{g_s} = \frac{\text{Net CO}_2 \text{ assimilation}}{\text{stomatal conductance}}
\label{eq:plant_responses_4}
\end{equation}

Plants under water stress maximise WUE by partially closing stomata.

\subsubsection{Worked Example: Calculating Water Use Efficiency
(WUE)}\label{worked-example-calculating-water-use-efficiency-wue}

\textbf{Problem:} An agronomist is evaluating a new drought-tolerant
wheat cultivar. Using a portable photosynthesis system, the researcher
measures a net CO\(_2\) assimilation rate (\(A_{net}\)) of
\(15 \mu\text{mol CO}_2 \text{ m}^{-2} \text{ s}^{-1}\). At the same
time, the stomatal conductance to water vapour (\(g_s\)) is measured at
\(0.15 \text{ mol H}_2\text{O m}^{-2} \text{ s}^{-1}\). 1. Calculate the
intrinsic Water Use Efficiency (WUE). 2. If ABA treatment causes the
stomata to partially close, decreasing \(g_s\) by 50\% but primarily
decreasing \(A_{net}\) by 20\%, what is the new WUE?

\textbf{Solution:}

\begin{enumerate}
\def\labelenumi{\arabic{enumi}.}
\item
  \textbf{Calculate the initial WUE:} It is standard to express \(g_s\)
  in \(\mu\text{mol}\) or \(A_{net}\) in \(\mu\text{mol}\) so units
  match. First, convert \(g_s\) from mol to \(\mu\text{mol}\):
  \[ g_s = 0.15 \text{ mol} = 150,000 \mu\text{mol H}_2\text{O m}^{-2} \text{ s}^{-1}  \label{eq:unit_VIII_plant_responses_item_1}\]

  Now, apply the WUE formula:
  \[ WUE_{initial} = \frac{A_{net}}{g_s} = \frac{15}{150,000} = 0.0001 \frac{\mu\text{mol CO}_2}{\mu\text{mol H}_2\text{O}}  \label{eq:unit_VIII_plant_responses_item_2}\]

  Alternatively, keeping \(g_s\) in mol:
  \[ WUE = \frac{15 \mu\text{mol CO}_2 \text{ m}^{-2} \text{ s}^{-1}}{0.15 \text{ mol H}_2\text{O m}^{-2} \text{ s}^{-1}} = 100 \frac{\mu\text{mol CO}_2}{\text{mol H}_2\text{O}}  \label{eq:unit_VIII_plant_responses_item_3}\]
\item
  \textbf{Calculate the new rates after ABA treatment:} New \(g_s\):
  \(0.15 \cdot (1 - 0.50) = 0.075 \text{ mol H}_2\text{O m}^{-2} \text{ s}^{-1}\)
  New \(A_{net}\):
  \(15 \cdot (1 - 0.20) = 12 \mu\text{mol CO}_2 \text{ m}^{-2} \text{ s}^{-1}\)
\item
  \textbf{Calculate the new WUE:}
  \[ WUE_{new} = \frac{12}{0.075} = 160 \frac{\mu\text{mol CO}_2}{\text{mol H}_2\text{O}}  \label{eq:unit_VIII_plant_responses_item_4}\]

  The partial stomatal closure significantly improved the water use
  efficiency (from 100 to 160), demonstrating how plants optimize carbon
  gain against water loss during drought conditions.
\end{enumerate}

\textbf{Root architecture changes under drought:} ABA and auxin promote
deeper root growth. Lateral root emergence is suppressed while primary
root elongation continues, shifting the root system to access deeper
soil water.

\subsubsection{Cold Stress -- CBF-COR
Pathway}\label{cold-stress-cbf-cor-pathway}

\textbf{Cold acclimation} (hardening at 4 degrees C for 1-3 weeks)
increases freezing tolerance:

\begin{enumerate}
\def\labelenumi{\arabic{enumi}.}
\tightlist
\item
  Temperature drop triggers membrane rigidification (gel-state
  transition) and Ca\(^{2+}\) influx (MCA1/MCA2 mechanosensitive
  channels)
\item
  \textbf{ICE1} (Inducer of CBF Expression 1; bHLH transcription factor)
  activates \textbf{CBF1/CBF2/CBF3} (C-repeat Binding Factors; AP2/ERF
  TFs) within minutes
\item
  CBFs bind CRT/DRE elements in promoters, activating \textbf{COR
  (Cold-Regulated) genes:}

  \begin{itemize}
  \tightlist
  \item
    \textbf{COR15:} \hyperref[gl:chloroplast]{\textbf{Chloroplast}}
    outer envelope protein; prevents membrane vesiculation
  \item
    \textbf{KIN1/KIN2:} Dehydrins; cryoprotectants that prevent protein
    aggregation
  \item
    \textbf{FAD2/FAD3:} Fatty acid desaturases that increase membrane
    unsaturated fatty acid content, maintaining membrane fluidity
  \end{itemize}
\item
  \textbf{Extracellular ice management:} Ice nucleates in the apoplast
  (lower solute concentration than
  \hyperref[gl:cytoplasm]{\textbf{cytoplasm}}). Water is drawn
  osmotically from cells, producing controlled dehydration rather than
  lethal intracellular ice formation.
\end{enumerate}

\textbf{Antifreeze proteins:} Some cold-tolerant plants (rye, carrots)
express proteins that bind ice crystal surfaces and inhibit crystal
growth (thermal hysteresis), protecting membranes.

\subsubsection{Heat Stress -- HSP
Response}\label{heat-stress-hsp-response}

Above 35-42 degrees C, the heat shock response is activated: -
\textbf{HSFA1} transcription factors activate \textbf{HSP70, HSP90,
HSP101} chaperones that prevent protein aggregation and refold misfolded
proteins - \textbf{Small HSPs} (sHSPs, 15-40 kDa): ATP-independent
holdases that accumulate in cytoplasmic aggregates -
\textbf{Thermomemory:} H3K4me3 epigenetic marks at HSP loci persist for
several days, allowing faster response to subsequent heat episodes

\subsubsection{Salt Stress}\label{salt-stress}

High soil salinity (NaCl) causes both osmotic stress and ion toxicity: -
\textbf{SOS pathway} (Salt Overly Sensitive): SOS3 (Ca\(^{2+}\) sensor)
activates SOS2 (kinase), which activates \textbf{SOS1} (Na\(^+\)/H\(^+\)
antiporter) to expel Na\(^+\) from roots - \textbf{Vacuolar
sequestration:} NHX1 antiporter pumps Na\(^+\) into vacuole -
\textbf{Compatible osmolytes:} Proline, glycine betaine, mannitol
accumulate to lower \(\psi_s\) without enzyme inhibition

\subsubsection{Flooding Response}\label{flooding-response}

Submergence deprives roots of O\(_2\): - \textbf{ERFVII transcription
factors} (RAP2.2, RAP2.12): Stabilised under hypoxia (normally degraded
by the N-end rule pathway when O\(_2\) is present -- direct O\(_2\)
sensing) - ERFVIIs activate \textbf{\hyperref[gl:anaerobic]{anaerobic}
genes:} ADH (alcohol dehydrogenase) for
\hyperref[gl:fermentation]{\textbf{fermentation}}, PDC (pyruvate
decarboxylase) - \textbf{Aerenchyma formation:} Ethylene (accumulates
because ACC oxidase requires O\(_2\)) triggers programmed cell death in
cortical cells, creating gas-filled channels for O\(_2\) transport from
aerial parts to submerged roots - \textbf{Escape strategy:} Deepwater
rice varieties elongate internodes rapidly (up to 25 cm/day) via GA and
ethylene signalling to keep leaves above water

\begin{center}\rule{0.5\linewidth}{0.5pt}\end{center}

\subsection{Plant Immunity --- A Two-Layer Defence
System}\label{plant-immunity-a-two-layer-defence-system}

Plants lack circulating immune cells but possess a sophisticated
cell-autonomous immune system organised into two layers. The
Jones--Dangl ``zigzag model'' (2006) is the canonical framework.

\subsubsection{Layer 1: PAMP-Triggered Immunity
(PTI)}\label{layer-1-pamp-triggered-immunity-pti}

Plants recognise broadly conserved molecular signatures of pathogens ---
\textbf{PAMPs} (Pathogen-Associated Molecular Patterns) or
\textbf{MAMPs} (Microbe-Associated Molecular Patterns) --- via
cell-surface \textbf{pattern recognition receptors (PRRs)}:

\begin{itemize}
\tightlist
\item
  \textbf{FLS2} (Flagellin Sensitive 2): A leucine-rich repeat receptor
  kinase (LRR-RK) that binds \textbf{flg22} (a 22 amino acid epitope of
  bacterial flagellin). FLS2-flg22 binding recruits \textbf{BAK1}
  co-receptor, activating intracellular kinase signalling.
\item
  \textbf{EFR}: Recognises bacterial elongation factor Tu (epitope
  \textbf{elf18})
\item
  \textbf{CERK1} and \textbf{LYK} receptors: Recognise fungal chitin
  oligomers (mostly tetra- and pentamers of N-acetylglucosamine)
\item
  \textbf{PEPR1/2}: Recognise plant-derived danger signals (PEPs)
\end{itemize}

\textbf{PTI signalling cascade:} 1. PRR ligand binding \texorpdfstring{\ensuremath{\to}}{->} BAK1
co-receptor recruitment 2. \textbf{BIK1} receptor-like cytoplasmic
kinase activated 3. \textbf{RBOHD} NADPH oxidase produces apoplastic ROS
burst (reactive oxygen species; H\(_2\)O\(_2\) within minutes) 4.
\textbf{MPK3/MPK6} MAP kinase cascade activated; transcription of
defence genes (PR1, WRKY TFs) 5. \textbf{Stomatal closure} via SLAC1
(preventing bacterial entry through stomata) 6. \textbf{Callose
deposition} at cell wall (papillae) limiting pathogen ingress

PTI provides broad-spectrum, modest resistance to most non-specialist
pathogens.

\subsubsection{Layer 2: Effector-Triggered Immunity
(ETI)}\label{layer-2-effector-triggered-immunity-eti}

Successful pathogens deliver \textbf{effector proteins} into host cells
(via type III secretion systems in bacteria, haustoria in
fungi/oomycetes) to suppress PTI. Plants in turn evolved intracellular
\textbf{NBS-LRR proteins} (also called R proteins or NLRs) that
recognise these effectors:

\begin{itemize}
\tightlist
\item
  \textbf{NBS-LRR} (Nucleotide-Binding Site, Leucine-Rich Repeat):
  Intracellular immune receptors. \emph{Arabidopsis} has
  about 150 NBS-LRR genes. Two structural classes: TIR-NLRs
  (with N-terminal Toll/Interleukin-1 Receptor domain) and CC-NLRs
  (coiled-coil).
\item
  \textbf{Recognition modes:}

  \begin{itemize}
  \tightlist
  \item
    \textbf{Direct binding:} R protein binds effector directly (e.g.,
    RPP1 binds \emph{Hyaloperonospora} ATR1)
  \item
    \textbf{Indirect ``guard'' model:} R protein guards a host target of
    effector; effector modification of the target triggers R activation
    (e.g., RPS2 guards RIN4; \emph{Pseudomonas} AvrRpt2 cleaves RIN4,
    activating RPS2)
  \item
    \textbf{Decoy model:} Plant evolves a non-functional mimic of an
    effector target as a ``decoy'' to bait R protein activation
  \end{itemize}
\end{itemize}

\textbf{ETI signalling and the hypersensitive response (HR):}

\begin{enumerate}
\def\labelenumi{\arabic{enumi}.}
\tightlist
\item
  NLR activation \texorpdfstring{\ensuremath{\to}}{->} oligomerisation (resistosome formation; ZAR1 forms a
  pentameric ``resistosome'' channel; Wang \emph{et al.}, \emph{Science}
  2019)
\item
  Resistosome inserts into plasma membrane and conducts Ca\(^{2+}\) \texorpdfstring{\ensuremath{\to}}{->}
  cytosolic {[}Ca\(^{2+}\){]} surge
\item
  \textbf{EDS1/PAD4 hub} activates downstream signalling
\item
  \textbf{Hypersensitive response (HR):} Programmed cell death at the
  infection site; sacrificial necrotic lesion that prevents pathogen
  spread. HR is morphologically similar to animal apoptosis but
  molecularly distinct (involves vacuolar processing enzymes, autophagy
  components, and ROS).
\item
  \textbf{Salicylic acid (SA) accumulation} locally and systemically
\end{enumerate}

ETI is faster and stronger than PTI --- typically resulting in the
visible \textbf{hypersensitive lesion} within 24--48 hours of infection.

\subsubsection{Systemic Acquired Resistance
(SAR)}\label{systemic-acquired-resistance-sar}

Local infection triggers a long-lasting, systemic defence priming
throughout the plant --- \textbf{systemic acquired resistance (SAR)}.
The mobile signal includes: - \textbf{Methyl salicylate (MeSA)}:
volatile, can travel through phloem and air - \textbf{Pipecolic acid
(Pip)} and N-hydroxy-pipecolic acid (NHP): essential systemic mobile
signals (Bernsdorff \emph{et al.}, \emph{PNAS} 2016) -
\textbf{Glycerol-3-phosphate (G3P)} and azelaic acid

In distal tissues, SA accumulates and activates \textbf{NPR1}
(Non-expressor of PR genes 1):

\begin{enumerate}
\def\labelenumi{\arabic{enumi}.}
\tightlist
\item
  NPR1 in the cytoplasm exists as oligomers (inactive) due to disulfide
  bonds
\item
  SA accumulation alters cellular redox; NPR1 is reduced to monomers
\item
  Monomeric NPR1 enters nucleus and binds \textbf{TGA} transcription
  factors
\item
  NPR1-TGA complex activates \textbf{PR1 (Pathogenesis-Related 1)} and
  other PR genes --- antimicrobial proteins, \texorpdfstring{\ensuremath{\beta}}{beta}-1,3-glucanases,
  chitinases
\item
  SAR provides 1--2 weeks of broad-spectrum resistance to subsequent
  infections
\end{enumerate}

\subsubsection{JA/SA Antagonism}\label{jasa-antagonism}

The two major immune hormones --- \textbf{salicylic acid (SA)} for
biotrophic pathogens and \textbf{jasmonic acid (JA)} for chewing
herbivores and necrotrophic pathogens --- exhibit pervasive
\textbf{antagonism}:

\begin{itemize}
\tightlist
\item
  \textbf{SA suppresses JA:} SA-induced NPR1 represses \emph{MYC2}
  (master JA TF) and JA-responsive genes; high SA \texorpdfstring{\ensuremath{\to}}{->} low JA defences
\item
  \textbf{JA suppresses SA:} JA-induced \emph{WRKY70} and \emph{MYC2}
  repress SA biosynthesis (ICS1) and PR1; high JA \texorpdfstring{\ensuremath{\to}}{->} low SA defences
\end{itemize}

\textbf{Why antagonism?} Different pathogens require different defence
chemistries: biotrophs (powdery mildew, rust fungi) feed on living
tissue and are best controlled by SA-induced HR; necrotrophs
(\emph{Botrytis}, \emph{Pythium}) and herbivores benefit from cell death
and are best controlled by JA-induced toxins/proteases. Allocating to
both is metabolically expensive; antagonism enables prioritisation.

\textbf{Pathogen exploitation of antagonism:} \emph{Pseudomonas
syringae} pv. \emph{tomato} produces \textbf{coronatine} (COR), a
structural mimic of JA-Ile. COR activates the JA pathway, suppressing
the SA pathway and stomatal-closure-mediated immunity. The bacterium
uses this to enter through normally-closed stomata. \emph{Arabidopsis
coi1} mutants (insensitive to JA) are resistant to COR-mediated stomatal
reopening.

\begin{center}\rule{0.5\linewidth}{0.5pt}\end{center}

\subsection{Herbivory and Wound Defence -- Jasmonate
Signalling}\label{herbivory-and-wound-defence-jasmonate-signalling}

Physical insect damage triggers the JA signalling cascade:

\begin{enumerate}
\def\labelenumi{\arabic{enumi}.}
\tightlist
\item
  Membrane damage releases linolenic acid
\item
  \textbf{13-LOX} (lipoxygenase) pathway in plastids produces 12-OPDA
\item
  Peroxisomal \texorpdfstring{\ensuremath{\beta}}{beta}-oxidation produces JA
\item
  \textbf{JAR1} conjugates JA with isoleucine to form \textbf{JA-Ile}
\item
  JA-Ile binds \textbf{COI1} (F-box protein), recruiting JAZ repressors
  for ubiquitin-mediated degradation
\item
  \textbf{MYC2/MYC3/MYC4} transcription factors are released,
  activating:

  \begin{itemize}
  \tightlist
  \item
    \textbf{Protease inhibitors:} Reduce digestibility of plant proteins
    in insect gut
  \item
    \textbf{Alkaloids:} Nicotine in tobacco, glucosinolates in
    Brassicaceae
  \item
    \textbf{Volatile terpenes:} (E)-\texorpdfstring{\ensuremath{\beta}}{beta}-ocimene, linalool attract
    predatory insects and parasitic wasps (indirect defence)
  \end{itemize}
\end{enumerate}

\begin{quote}
\textbf{Clinical Connection:} The wound response in plants is remarkably
fast. Within 30 seconds of \emph{Manduca sexta} (tobacco hornworm)
feeding on \emph{Nicotiana attenuata}, JA-Ile accumulates. Within 1
hour, systemic protease inhibitor expression is detected. Within 24
hours, volatile terpenoids attract predatory \emph{Geocoris} bugs. This
multi-layered defence system has inspired ``push-pull'' agricultural
strategies where companion plants emit volatiles to repel pests and
attract their natural enemies.
\end{quote}

\begin{center}\rule{0.5\linewidth}{0.5pt}\end{center}

\subsection{Shade Avoidance Syndrome}\label{shade-avoidance-syndrome}

In canopy shade, the R:FR (red:far-red) ratio decreases because
chlorophyll absorbs red light but transmits far-red:

\begin{itemize}
\tightlist
\item
  \textbf{PhyB} converts from active Pfr to inactive Pr form
\item
  \textbf{PIF4, PIF5, PIF7} (Phytochrome Interacting Factors) freed from
  PhyB repression
\item
  PIF4/5 activate \textbf{YUCCA} genes, increasing auxin synthesis in
  hypocotyls
\item
  Rapid hypocotyl and petiole elongation (shade avoidance)
\item
  Accelerated flowering (via PIF4 activation of FT promoter)
\item
  Reduced branching, leaf area, and root investment
\end{itemize}

This shade avoidance syndrome represents a competitive strategy: grow
taller to overtop neighbours and reproduce before being shaded out.
However, it comes at a cost of reduced root investment and mechanical
stability.

\begin{center}\rule{0.5\linewidth}{0.5pt}\end{center}

\subsection{Epigenetic Stress Memory in
Plants}\label{epigenetic-stress-memory-in-plants}

Plants lack a germline-soma separation and can transmit epigenetic marks
trans-generationally:

\begin{itemize}
\tightlist
\item
  \textbf{Somatic priming:} SA treatment deposits H3K4me3 marks at WRKY
  promoters, maintained for weeks. Subsequent pathogen attack triggers
  faster and stronger PR gene induction.
\item
  \textbf{Trans-generational immunity:} Pathogen-stressed parents
  produce offspring with enhanced resistance via DNA methylation changes
  at transposable element loci (RdDM pathway -- RNA-directed DNA
  methylation).
\item
  \textbf{Heat stress memory:} H3K4me3 at HSP genes persists for several
  days after heat episode.
\end{itemize}

\textbf{Concept Check:} Why might trans-generational epigenetic stress
memory be particularly advantageous in plants compared to animals?

\begin{center}\rule{0.5\linewidth}{0.5pt}\end{center}

\subsection{Agricultural Applications of Plant Hormone
Biology}\label{agricultural-applications-of-plant-hormone-biology}

Decades of basic research on plant hormones have generated
transformative agricultural technologies:

The strongest applications now combine hormone biology with genotype,
microbiome, water status, and field context. Drought tolerance, heat
tolerance, salinity tolerance, and pathogen resistance often trade off
against growth or yield, so breeders increasingly stack traits and test
them across environments rather than searching for a single
comprehensive ``stress gene.'' Hormone interventions should be evaluated
by phenotype, resource cost, reproductive outcome, and performance under
realistic combined stresses.

\subsubsection{Ethylene Inhibitors --- Postharvest
Storage}\label{ethylene-inhibitors-postharvest-storage}

\textbf{1-MCP (1-methylcyclopropene):} A small gaseous molecule that
competitively binds the ethylene receptor ETR1, blocking ethylene
perception. Commercial product \emph{SmartFresh} applied to apples,
pears, kiwifruit, and avocados: - Apples treated with 1-MCP can be
stored 6--9 months (vs 2--3 months untreated) with retained firmness,
sugar content, and flavour - The North American apple industry is
essentially built on 1-MCP storage - Other ethylene management:
\textbf{AVG} (aminoethoxyvinylglycine; ACS inhibitor) used pre-harvest
in stone fruits to delay ripening; \textbf{silver thiosulfate (STS)} in
cut flowers to preserve vase life

\subsubsection{Gibberellin Biology --- The Green
Revolution}\label{gibberellin-biology-the-green-revolution}

The semi-dwarf wheat varieties that catalysed the Green Revolution carry
mutations in \textbf{gibberellin signalling}: - \textbf{Rht-B1b and
Rht-D1b} in wheat: encode gain-of-function alleles of \emph{DELLA} genes
that produce stable, GA-insensitive DELLA proteins. Plants do not
respond to GA \texorpdfstring{\ensuremath{\to}}{->} reduced internodal elongation \texorpdfstring{\ensuremath{\to}}{->} semi-dwarf stature \texorpdfstring{\ensuremath{\to}}{->} no
lodging under heavy grain heads. - \textbf{sd1} in rice (the IR8
``miracle rice'' of 1966): loss-of-function mutation in \emph{GA20ox-2}
(a GA biosynthesis enzyme) \texorpdfstring{\ensuremath{\to}}{->} reduced GA levels \texorpdfstring{\ensuremath{\to}}{->} semi-dwarf - Both
mutations enable plants to support large grain heads, harvest more
nitrogen-fertiliser into seed, and resist storm damage. Together they
account for \textgreater50\% of global wheat and rice acreage.

\textbf{GA application:} Exogenous GA\(_3\) is sprayed on: - Seedless
table grapes (induces parthenocarpy and elongates clusters) - Sugarcane
(increases stem length and sugar yield) - Hybrid rice seed production
(activates female lines)

\subsubsection{Auxin Herbicides}\label{auxin-herbicides}

\textbf{2,4-D (2,4-dichlorophenoxyacetic acid)} and related compounds
(dicamba, MCPA) are synthetic auxin analogues used as selective
broadleaf herbicides: - Mechanism: synthetic auxins activate the same
TIR1-Aux/IAA-ARF pathway as natural IAA, but at much higher and
sustained concentrations - Sustained activation triggers ethylene
biosynthesis, ABA accumulation, and ultimately growth abnormalities and
death in dicots - Monocots (grasses, including cereal crops) tolerate
2,4-D because they have lower TIR1 binding affinity, faster auxin
metabolism, and morphology that limits leaf retention of spray - 2,4-D
was developed in WWII as part of biological warfare research;
declassified post-war and adopted for agricultural use in 1946. It
became the most-used herbicide globally for decades. -
Glyphosate-resistance now dominates, but 2,4-D and dicamba have
re-emerged in stacked-trait crops (Enlist E3 soybean, Xtend cotton)

\subsubsection{Other Hormone-Based
Technologies}\label{other-hormone-based-technologies}

\begin{itemize}
\tightlist
\item
  \textbf{Cytokinin sprays:} maintain green leaves in cut flowers; delay
  senescence
\item
  \textbf{ABA analogues:} in development for drought-priming of crop
  seedlings
\item
  \textbf{Strigolactone analogues (e.g., GR24):} induce suicidal
  germination of \emph{Striga} parasitic weeds in absence of host,
  reducing field infestation
\item
  \textbf{Plant growth retardants:} uniconazole, paclobutrazol --- GA
  biosynthesis inhibitors used to keep ornamentals compact
\end{itemize}

\begin{quote}
\textbf{Clinical Connection:} Plant hormone biology supplies an
extraordinary fraction of agricultural innovation. The Green Revolution
(semi-dwarf GA-insensitive wheat and rice), modern fruit storage
(1-MCP), and the largest herbicide markets (2,4-D, dicamba --- synthetic
auxins) most derive directly from understanding hormone signalling
pathways characterised in \emph{Arabidopsis} and other model species.
\end{quote}

\begin{center}\rule{0.5\linewidth}{0.5pt}\end{center}

\subsection{Current Evidence and Frontier
Biology}\label{current-evidence-and-frontier-biology-29}

For \textbf{Plant Responses to the Environment}, frontier biology
belongs inside the evidence logic of the chapter. Plant biology links
molecular regulation to climate stress, water limitation, crop
resilience, phenology, and ecosystem feedbacks. The core reading
question is this: plant response claims should connect signal
perception, hormone network, gene expression, phenotype, and tradeoff.

\begin{itemize}
\tightlist
\item
  \textbf{What to verify:} identify the observation, model, assay, or
  dataset that would make the claim stronger or weaker.
\item
  \textbf{What to qualify:} state the scale, organism, cell type,
  environmental condition, or population where the claim is expected to
  hold.
\item
  \textbf{What to compare:} test at least one alternative explanation,
  baseline, or null model before treating the pattern as causal.
\item
  \textbf{What to cite:} distinguish primary evidence, review synthesis,
  public dataset, and institutional guidance; for recent or numeric
  claims, prefer the source closest to the measurement and state what
  has changed since it was published.
\end{itemize}

A strong plant explanation names the tissue, signal, environmental
driver, measurable trait, and tradeoff between growth, reproduction,
defence, and water use.

\textbf{Source practice:} For plant-stress and crop claims, name the
tissue, environmental driver, field context, and growth-reproduction
tradeoff; separate laboratory potential from agronomic adoption.

\subsubsection{Current Evidence Map: Plant Stress Response
Decision}\label{current-evidence-map-plant-stress-response-decision}

\begin{figure}[htbp]
\centering
\includegraphics[width=0.82\linewidth,height=4.2in,keepaspectratio]{../figures/mermaid_inline/inline_mermaid_0150_9c1b03a0028a.png}
\caption{Stress tolerance is not free: the same response that improves survival can reduce growth or reproduction depending on timing and environment.}
\end{figure}

\subsubsection{Field Translation: Stress Circuits Are Not Field
Proof}\label{field-translation-stress-circuits-are-not-field-proof}

Synthetic response circuits and hormone interventions should be read as
hypotheses about a controllable decision point, not as proof of
agronomic benefit. A moisture-responsive stomatal circuit must show that
closure occurs before damaging xylem tension but does not depress carbon
gain enough to reduce yield. An inducible defence promoter must show low
leakiness, fast activation under natural herbivory, and minimal cost to
growth or beneficial insects. A flowering-time circuit must be tested
against local chilling, photoperiod, heat, and pollinator windows. The
field claim becomes credible when mechanism, expression stability,
off-target phenotypes, gene flow, management context, and multi-season
yield or fitness data point in the same direction.

\subsection{Key Terms}\label{key-terms-29}

{\def\LTcaptype{none} % do not increment counter
\begin{longtable}[]{@{}
  >{\raggedright\arraybackslash}p{(\linewidth - 2\tabcolsep) * \real{0.2857}}
  >{\raggedright\arraybackslash}p{(\linewidth - 2\tabcolsep) * \real{0.7143}}@{}}
\toprule\noalign{}
\begin{minipage}[b]{\linewidth}\raggedright
Term
\end{minipage} & \begin{minipage}[b]{\linewidth}\raggedright
Definition
\end{minipage} \\
\midrule\noalign{}
\endhead
\bottomrule\noalign{}
\endlastfoot
\textbf{Tropism} & Directional growth response to an environmental
stimulus \\
\textbf{Phototropin (PHOT1/2)} & Blue-light photoreceptor kinases; LOV
domains; autophosphorylate upon blue light absorption \\
\textbf{Auxin (IAA)} & Indole-3-acetic acid; major growth hormone; polar
transport via PIN efflux carriers \\
\textbf{PIN proteins} & Auxin efflux carriers; polar localisation
determines transport direction; PIN2 for root gravitropism \\
\textbf{Acid growth hypothesis} & Auxin activates H\(^+\)-ATPase,
acidifying cell wall, activating expansins, enabling elongation \\
\textbf{DELLA proteins} & Transcriptional repressors degraded by GA-GID1
signalling; growth suppressors \\
\textbf{Phytochrome} & Dual R/FR photoreceptors (PhyA-E); Pr (inactive)
and Pfr (active) conformers \\
\textbf{Cryptochrome} & Blue/UV-A photoreceptor; FAD chromophore;
CRY1/CRY2 in \emph{Arabidopsis} \\
\textbf{COP1} & E3 ubiquitin ligase; ``darkness enforcer''; targets HY5,
CO, HFR1 for degradation \\
\textbf{PIFs} & Phytochrome Interacting Factors; bHLH TFs degraded by
Pfr; mediate shade avoidance \\
\textbf{Photoperiodism} & Regulation of developmental events (especially
flowering) by day/night length \\
\textbf{CONSTANS (CO)} & bZIP transcription factor; activator of FT in
long-day flowering \\
\textbf{Florigen (FT)} & Mobile protein produced in leaf companion
cells; travels to SAM via phloem to trigger flowering \\
\textbf{FLC} & MADS-box repressor of flowering; silenced by
vernalisation (H3K27me3 by VRN2-PRC2) \\
\textbf{Vernalisation} & Cold-mediated epigenetic silencing of FLC
enabling flowering \\
\textbf{Circadian clock} & Endogenous about 24 h oscillator;
CCA1/LHY-TOC1-PRR interlocked loops \\
\textbf{CCA1/LHY} & Morning-expressed MYB TFs; repress evening genes \\
\textbf{TOC1/PRR1} & Evening-expressed pseudo-response regulator;
represses CCA1/LHY \\
\textbf{PRR9/7/5} & Sequential daytime PRRs in evening loop \\
\textbf{Evening Complex (ELF3-ELF4-LUX)} & Repressor complex peaking at
ZT12-16; gates early-night growth \\
\textbf{Statolith} & Starch-filled amyloplast in columella cells;
sediments under gravity \\
\textbf{LAZY1} & Protein required for normal gravitropism; relocalises
to lower face after gravity stimulation \\
\textbf{TCH genes} & Touch-induced genes (CAM2, CML24, CML12, XTH22);
calmodulin-related \\
\textbf{Thigmomorphogenesis} & Mechanical-stimulus-induced morphology
change (shorter, sturdier plants) \\
\textbf{ABA} & Abscisic acid; stress hormone; PYR/PYL receptors activate
SnRK2 kinases \\
\textbf{CBF/DREB} & C-repeat Binding Factors; AP2/ERF TFs; cold-induced,
activate COR genes \\
\textbf{JA-Ile} & Jasmonyl-isoleucine conjugate; active jasmonate;
ligand for COI1 receptor \\
\textbf{COI1} & Coronatine Insensitive 1; F-box protein; JA-Ile
receptor; degrades JAZ repressors \\
\textbf{PTI} & PAMP-Triggered Immunity; first-layer defence via
cell-surface PRRs \\
\textbf{ETI} & Effector-Triggered Immunity; second-layer defence via
intracellular NBS-LRR proteins \\
\textbf{NBS-LRR (NLR)} & Nucleotide-Binding Site, Leucine-Rich Repeat;
intracellular immune receptors \\
\textbf{HR (Hypersensitive Response)} & Programmed cell death at
infection site; ETI hallmark \\
\textbf{SAR} & Systemic Acquired Resistance; SA-mediated whole-plant
defence priming \\
\textbf{NPR1} & SA-induced master regulator of PR gene expression \\
\textbf{Strigolactone} & Carotenoid-derived hormone; suppresses
branching; recruits mycorrhizae \\
\textbf{1-MCP} & 1-methylcyclopropene; ethylene receptor antagonist;
postharvest storage \\
\textbf{2,4-D} & Synthetic auxin herbicide; selective for dicots \\
\end{longtable}
}

\begin{center}\rule{0.5\linewidth}{0.5pt}\end{center}

\subsection{Review Questions}\label{review-questions-29}

\begin{enumerate}
\def\labelenumi{\arabic{enumi}.}
\item
  An \emph{Arabidopsis phot1 phot2} double mutant is grown under
  asymmetric blue light. Predict the
  \hyperref[gl:phenotype]{\textbf{phenotype}}. Would this mutant show
  any phototropic response under high-intensity white light? (Hint:
  consider cryptochromes.)
\item
  Explain why the same auxin redistribution (accumulation on the lower
  side) causes opposite growth responses in roots vs shoots during
  gravitropism. Include the concepts of optimal auxin concentration,
  dose-response curves, and PIN2 vs PIN3 localisation.
\item
  A farmer notices that long-day wheat varieties flower prematurely when
  grown near street lights that extend the day length. Explain the
  molecular mechanism, starting from phytochrome and proceeding through
  the CO-FT pathway. Include the role of CRY2 in stabilising CO protein.
\item
  Compare the ABA-mediated drought response with the CBF-COR cold
  acclimation pathway. What do they share in common? How do the final
  cellular protective mechanisms differ?
\item
  \emph{Pseudomonas syringae} produces coronatine (COR), a JA-Ile mimic.
  Explain how COR suppresses stomatal immunity (PTI) and promotes
  bacterial entry. Design a crop plant with COR-insensitive COI1 that
  resists this strategy without compromising wound response.
\item
  A researcher discovers that a particular \emph{Arabidopsis} accession
  shows enhanced cold tolerance without prior cold acclimation. The line
  has constitutively high expression of CBF3. What are the likely growth
  penalties of constitutive CBF expression, and why?
\item
  Explain the molecular basis of the ``night break'' experiment in
  short-day plants. Why does a brief pulse of red light in the middle of
  a long night prevent flowering, and why does an immediately subsequent
  far-red pulse restore the flowering response?
\item
  Compare ethylene's role in fruit ripening vs flooding response. How is
  the same hormone adapted for such different functions?
\item
  Sample three points from \breaktt{light\_response\_curve()} and
  identify roughly where the curve saturates for default C3-like
  parameters.
\item
  Explain how \textbf{phyB} thermosensing can decouple flowering time
  from true photoperiod in warming climates.
\item
  Describe the molecular memory of vernalisation in \emph{Arabidopsis}.
  How is the VRN2-PRC2-mediated H3K27me3 mark at \emph{FLC} maintained
  through cell division and erased between generations?
\item
  Distinguish PTI and ETI in the Jones--Dangl zigzag model. What is the
  role of resistosome formation? How does the SA-NPR1 pathway translate
  local infection into systemic acquired resistance?
\item
  Explain why 2,4-D selectively kills broadleaf weeds in cereal crops.
  At the molecular level, how does sustained synthetic auxin signalling
  cause death in dicots but not monocots?
\item
  Trace the \emph{Arabidopsis} circadian clock through one full day
  starting at dawn (ZT0). Identify which genes are expressed at ZT0,
  ZT4, ZT8, ZT12, and ZT16, and explain the negative-feedback
  architecture.
\item
  Tendril coiling in cucumber requires touch and JA. A \emph{coi1}
  loss-of-function mutant cucumber is challenged with a thin string.
  Predict the response, and explain why JA signalling is downstream of
  mechanical perception.
\end{enumerate}

\begin{center}\rule{0.5\linewidth}{0.5pt}\end{center}

\subsection{Further Reading and Source
Notes}\label{further-reading-and-source-notes-24}

\begin{itemize}
\tightlist
\item
  Hager (1971). Acid growth theory of auxin action. \emph{Planta}, 100.
\item
  \citet{went1926} --- On diffusion and the stimulation of growth by
  auxins. \emph{Recueil des Travaux Botaniques Néerlandais}, 33.
\end{itemize}

\begin{center}\rule{0.5\linewidth}{0.5pt}\end{center}

\subsection{Computational Bridge}\label{computational-bridge-26}

Photosynthetic light response is tabulated for any PAR grid:

\begin{Shaded}
\begin{Highlighting}[]
\ImportTok{from}\NormalTok{ biology.botany }\ImportTok{import}\NormalTok{ light\_response\_curve}

\NormalTok{curve }\OperatorTok{=}\NormalTok{ light\_response\_curve(n\_points}\OperatorTok{=}\DecValTok{8}\NormalTok{, max\_par}\OperatorTok{=}\FloatTok{1600.0}\NormalTok{)}
\BuiltInTok{print}\NormalTok{(}\BuiltInTok{round}\NormalTok{(curve[}\OperatorTok{{-}}\DecValTok{1}\NormalTok{][}\DecValTok{1}\NormalTok{], }\DecValTok{2}\NormalTok{))}
\end{Highlighting}
\end{Shaded}

\begin{quote}
\textbf{Clinical / systems note:} Shade-intolerant crops under
hospital-city light pollution mirror the long-day wheat problem ---
photoperiod pathways wired to yield stability.
\end{quote}

\begin{center}\rule{0.5\linewidth}{0.5pt}\end{center}

\subsection{Summary}\label{summary-29}

\begin{itemize}
\tightlist
\item
  \textbf{Plant hormones:} Nine major classes (auxin, GA, cytokinin,
  ABA, ethylene, JA, SA, brassinosteroids, strigolactones) regulate most
  aspects of growth, development, and stress responses through
  receptor-mediated signalling pathways.
\item
  \textbf{Phototropism \citep{went1926}:} PHOT1/2 (LOV-domain kinases)
  trigger asymmetric PIN3 relocalisation, creating a lateral auxin
  gradient. Acid growth (H\(^+\)-ATPase to expansins) drives
  differential elongation \citep{hager1971}.
\item
  \textbf{Photomorphogenesis:} Phytochromes (Pr/Pfr photoconversion at
  660/730 nm) sense R:FR ratio for shade avoidance; cryptochromes
  (CRY1/CRY2 with FAD) sense blue light for photomorphogenesis. PhyA
  mediates FR-HIR; PhyB is the canopy-shade sensor. CRY-COP1-HY5 module
  enforces light/dark choice.
\item
  \textbf{Circadian clock molecular architecture:} Morning loop
  (CCA1/LHY repressors of evening genes; TOC1 activator); evening loop
  (PRR9/7/5/TOC1 cascade repressing CCA1/LHY); evening complex
  (ELF3-ELF4-LUX). Temperature-compensated.
\item
  \textbf{Photoperiodism:} Night length is the true signal. Phytochrome
  (Pr/Pfr) measures dark duration. The CO-FT-AP1 pathway integrates
  circadian clock and light signals. FT protein moves from leaf
  companion cells to SAM as the long-sought florigen.
\item
  \textbf{Vernalisation:} Cold-induced VIN3 joins VRN2-PRC2; H3K27me3
  deposition at \emph{FLC} silences this MADS-box repressor of
  flowering. Memory is mitotically stable but reset between generations.
\item
  \textbf{Gravitropism:} Amyloplast settling in statocytes triggers PIN3
  relocalisation. Same auxin redistribution causes opposite responses in
  roots (PIN2 carries auxin acropetally; high auxin inhibits root cells)
  vs shoots (high auxin promotes shoot cells).
\item
  \textbf{Touch responses:} TCH genes (calmodulins, XTH22) mediate
  Ca\(^{2+}\)-driven thigmomorphogenesis; jasmonate links wound
  perception to defence and tendril coiling.
\item
  \textbf{Drought:} ABA activates PYR-PP2C-SnRK2-SLAC1 cascade for
  stomatal closure plus AREB/ABF for stress gene expression.
\item
  \textbf{Cold:} CBF1/2/3 activated by ICE1 within minutes of cold
  exposure; COR genes produce dehydrins, desaturases, and cryoprotective
  proteins.
\item
  \textbf{Flooding:} ERFVII transcription factors sense O\(_2\) via the
  N-end rule; ethylene promotes aerenchyma formation.
\item
  \textbf{Plant immunity (Jones--Dangl zigzag):} PTI via cell-surface
  PRRs (FLS2/BAK1, CERK1) recognising PAMPs. ETI via intracellular
  NBS-LRR (NLR) proteins recognising effectors; resistosome formation
  triggers HR. NPR1-mediated SAR provides systemic broad-spectrum
  priming via SA. JA-SA antagonism partitions defences;
  \emph{Pseudomonas} coronatine exploits this.
\item
  \textbf{Herbivory:} JA-Ile triggers COI1-JAZ degradation, releasing
  MYC2 to activate protease inhibitors, alkaloids, and volatile-mediated
  indirect defence.
\item
  \textbf{Shade avoidance:} Low R:FR frees PIF4/5 from PhyB repression,
  increasing auxin and accelerating flowering.
\item
  \textbf{Agricultural applications:} Ethylene inhibitors (1-MCP for
  postharvest); GA-insensitive DELLA mutations (\emph{Rht}, \emph{sd1})
  underlie Green Revolution semi-dwarf wheat and rice; synthetic auxin
  herbicides (2,4-D, dicamba) selectively kill broadleaf weeds.
\item
  \textbf{Connections:} See \cref{sec:unit_III_photosynthesis} for light
  reactions, \cref{sec:unit_II_cell_signaling} for receptor signalling
  parallels, and Unit X for phenology under climate change.
\end{itemize}

\begin{center}\rule{0.5\linewidth}{0.5pt}\end{center}

\subsubsection{Companion Source
Module}\label{companion-source-module-29}

\textbf{Plant Responses to the Environment} should leave a reproducible
trail from a biological claim to the code, figure, diagram, or
paper-based activity that can test it. Use the surfaces below to inspect
the chapter's assumptions, rerun the relevant model, or compare the
manuscript explanation with companion labs and figures.

{\def\LTcaptype{none} % do not increment counter
\begin{longtable}[]{@{}
  >{\raggedright\arraybackslash}p{(\linewidth - 2\tabcolsep) * \real{0.5000}}
  >{\raggedright\arraybackslash}p{(\linewidth - 2\tabcolsep) * \real{0.5000}}@{}}
\toprule\noalign{}
\begin{minipage}[b]{\linewidth}\raggedright
Surface
\end{minipage} & \begin{minipage}[b]{\linewidth}\raggedright
Use it for
\end{minipage} \\
\midrule\noalign{}
\endhead
\bottomrule\noalign{}
\endlastfoot
\breaktt{src/biology/botany/botany.py} (\breaktt{photosynthesis\_rate},
\breaktt{light\_response\_curve}, \breaktt{transpiration\_flux}) &
Quantify how light, CO2, water, and temperature shape response
curves. \\
\breaktt{src/visualization/plots.py}
(\breaktt{plot\_light\_response\_curve}) & Check saturation and stress
interpretation visually. \\
\breaktt{src/mermaid/biology\_diagrams.py}
(\breaktt{photosynthesis\_light\_dark\_diagram},
\breaktt{hormone\_signaling\_diagram}) & Link environmental sensing to
pathway response. \\
\end{longtable}
}

\textbf{Reproducibility check:} name the signal, receptor/tissue,
hormone network, phenotype, and tradeoff before claiming adaptive
response. \textbf{Cross-reference:} use
\cref{sec:unit_VIII_plant_structure_and_water},
\cref{sec:unit_VIII_plant_reproduction}, and
\cref{sec:unit_III_photosynthesis}.

\newpage

\section*{Unit IX --- Zoology and Systems Physiology:
Introduction}\label{unit-ix-zoology-and-systems-physiology-introduction}
\addcontentsline{toc}{section}{Unit IX --- Zoology and Systems
Physiology: Introduction}

\subsection*{Why This Unit Matters}\label{why-this-unit-matters-9}
\addcontentsline{toc}{subsection}{Why This Unit Matters}

In 1628, William Harvey published \emph{Exercitatio Anatomica de Motu
Cordis et Sanguinis in Animalibus} (On the Motion of the Heart and Blood
in Animals) and overturned nearly 1,500 years of Galenic medicine.
Harvey demonstrated, through quantitative reasoning alone, that the
heart must be a pump recirculating the same blood --- not a furnace
generating new blood from food. He calculated that the heart pushes
approximately 1.8 litres of blood per minute; if blood flowed outward
and were not returned, the entire blood volume (about 5 litres) would be
exhausted in under 3 minutes. Therefore, it must circulate. This was
physiology as computational biology --- numbers proving mechanism before
any understanding of cellular biology existed.

Modern systems physiology is even more mathematical. The Hodgkin-Huxley
model (1952) describes the action potential as a set of four coupled
nonlinear ordinary differential equations and was solved numerically
before computers capable of doing so efficiently were widely available.
Today, the same mathematical framework underlies cardiac
electrophysiology, pharmacological ion channel modelling, and the design
of neural prosthetics. Homeostasis --- the maintenance of stable
internal conditions --- is formalised as a negative feedback control
system: a detector, comparator, and effector maintaining a set-point,
mathematically identical to engineering control systems.

This unit covers four major physiological systems --- circulatory,
nervous, endocrine, and immune --- with quantitative models throughout.
You will apply the Hodgkin-Huxley equations to the action potential,
calculate Starling forces for capillary filtration, model hormonal
negative feedback with Hill equations, and contrast innate and adaptive
immune responses as computational pattern-recognition systems. Clinical
connections appear in every section: from cardiac arrhythmias to
diabetes mellitus, from anaphylaxis to Parkinson's disease.

\begin{center}\rule{0.5\linewidth}{0.5pt}\end{center}

\subsection*{Landmark Discoveries}\label{landmark-discoveries-9}
\addcontentsline{toc}{subsection}{Landmark Discoveries}

{\def\LTcaptype{none} % do not increment counter
\begin{longtable}[]{@{}
  >{\raggedright\arraybackslash}p{(\linewidth - 8\tabcolsep) * \real{0.2407}}
  >{\raggedright\arraybackslash}p{(\linewidth - 8\tabcolsep) * \real{0.0741}}
  >{\raggedright\arraybackslash}p{(\linewidth - 8\tabcolsep) * \real{0.2963}}
  >{\raggedright\arraybackslash}p{(\linewidth - 8\tabcolsep) * \real{0.1667}}
  >{\raggedright\arraybackslash}p{(\linewidth - 8\tabcolsep) * \real{0.2222}}@{}}
\toprule\noalign{}
\begin{minipage}[b]{\linewidth}\raggedright
Discoverer(s)
\end{minipage} & \begin{minipage}[b]{\linewidth}\raggedright
Year
\end{minipage} & \begin{minipage}[b]{\linewidth}\raggedright
Journal / Source
\end{minipage} & \begin{minipage}[b]{\linewidth}\raggedright
Discovery
\end{minipage} & \begin{minipage}[b]{\linewidth}\raggedright
Significance
\end{minipage} \\
\midrule\noalign{}
\endhead
\bottomrule\noalign{}
\endlastfoot
William Harvey & 1628 & \emph{Exercitatio Anatomica de Motu Cordis} &
Systemic circulation: heart as pump & Quantitative physiology; disproved
Galenic humoral theory \\
Claude Bernard & 1854--1879 & \emph{Leçons de physiologie expérimentale}
& Internal milieu and homeostasis & Defined the concept of regulated
internal environment \\
Ernest Starling & 1918 & \emph{J. Physiol.} & Frank-Starling law of the
heart & Cardiac output proportional to venous return; basis of heart
failure physiology \\
Alan Hodgkin \& Andrew Huxley & 1952 & \emph{J. Physiol.} & Mathematical
model of nerve action potential & Four ODEs describing Na\texorpdfstring{\textsuperscript{+}}{+}/K\texorpdfstring{\textsuperscript{+}}{+}
conductance; Nobel Prize 1963 \\
Barry Marshall \& Robin Warren & 1984 & \emph{The Lancet} &
\emph{Helicobacter pylori} causes peptic ulcers & Disproved "acid
stress" dogma; Nobel Prize 2005; key discovery in gastrointestinal
physiology \\
Rosalyn Yalow \& Solomon Berson & 1959 & \emph{J. Clin. Invest.} &
Radioimmunoassay (RIA) for insulin & Enabled measurement of hormone
concentrations at picomolar levels; Nobel Prize 1977 \\
César Milstein \& Georges Köhler & 1975 & \emph{Nature} & Monoclonal
antibody production (hybridoma) & Foundation of modern immunotherapy,
diagnostics, and targeted cancer therapy; Nobel Prize 1984 \\
\end{longtable}
}

\begin{center}\rule{0.5\linewidth}{0.5pt}\end{center}

\subsection*{Key Concepts and
Connections}\label{key-concepts-and-connections-9}
\addcontentsline{toc}{subsection}{Key Concepts and Connections}

\begin{figure}[htbp]
\centering
\includegraphics[width=0.82\linewidth,height=4.2in,keepaspectratio]{../figures/mermaid_inline/inline_mermaid_0151_39f9cb216cd7.png}
\caption{Flowchart for Key Concepts and Connections: balance Homeostasis (negative feedback; set-point; effector), Circulatory System (cardiac output; BP; Starling forces), Respiratory System (gas exchange; ventilation-perfusion), and Endocrine System (hormones; anterior pituitary; feedback loops) form the diagram's primary path or branches.}
\end{figure}

\textbf{Unit IX concept map --- Zoology and Systems Physiology.}

\begin{center}\rule{0.5\linewidth}{0.5pt}\end{center}

\subsection*{Current Evidence Thread}\label{current-evidence-thread-9}
\addcontentsline{toc}{subsection}{Current Evidence Thread}

Read this unit as physiology that is known because it has been measured:
every claim about action potentials, circulation, hormones, immunity,
and neural circuits rests on a recording, an image, a perturbation, or a
clinical measurement, and the strength of the claim is no better than
the method behind it. Physiology now blends mechanism with allostasis,
immune-endocrine-neural coupling, wearable data, and individualized risk
without reducing bodies to simple machines. As you move through the
chapters, keep a two-column note: \textbf{claim} on the left,
\textbf{evidence that would change my confidence} on the right. By the
end of the unit, each major idea should be tied to a measurement, model,
citation, or paper-based lab decision.

\subsection*{Chapter Roadmap}\label{chapter-roadmap-10}
\addcontentsline{toc}{subsection}{Chapter Roadmap}

{\def\LTcaptype{none} % do not increment counter
\begin{longtable}[]{@{}
  >{\raggedright\arraybackslash}p{(\linewidth - 6\tabcolsep) * \real{0.1556}}
  >{\raggedright\arraybackslash}p{(\linewidth - 6\tabcolsep) * \real{0.1111}}
  >{\raggedright\arraybackslash}p{(\linewidth - 6\tabcolsep) * \real{0.2889}}
  >{\raggedright\arraybackslash}p{(\linewidth - 6\tabcolsep) * \real{0.4444}}@{}}
\toprule\noalign{}
\begin{minipage}[b]{\linewidth}\raggedright
Chapter
\end{minipage} & \begin{minipage}[b]{\linewidth}\raggedright
Title
\end{minipage} & \begin{minipage}[b]{\linewidth}\raggedright
Core Question
\end{minipage} & \begin{minipage}[b]{\linewidth}\raggedright
Key Equation / Model
\end{minipage} \\
\midrule\noalign{}
\endhead
\bottomrule\noalign{}
\endlastfoot
\textbf{28} & Circulation, Respiration, and Homeostasis & How do the
cardiovascular and respiratory systems maintain homeostasis? & Fick's
principle; Starling equation; \(P_{O_2}\)/Hb saturation (Hill eq.) \\
\textbf{29} & Nervous System Organisation & How is the nervous system
organised from molecular to systems level? & Cable equation; receptor
field models \\
\textbf{30} & Action Potentials and Synapses & How do neurons generate
and transmit electrical signals? & Hodgkin-Huxley:
\(I = C_m dV/dt + I_{Na} + I_K + I_L\) \\
\textbf{31} & Endocrine and Immune Systems & How do hormones and immune
cells maintain internal defence and stability? & Negative feedback
oscillator; MHC diversity; clonal selection \\
\end{longtable}
}

\begin{center}\rule{0.5\linewidth}{0.5pt}\end{center}

\subsection*{Connections Across the
Textbook}\label{connections-across-the-textbook-9}
\addcontentsline{toc}{subsection}{Connections Across the Textbook}

\begin{itemize}
\tightlist
\item
  \textbf{Ion channels and membrane potential} (this unit) build
  directly on membrane transport (\cref{sec:unit_II_membrane_transport})
  and electrochemistry (Unit I).
\item
  \textbf{Haemoglobin cooperativity} (Hill equation, \(n \approx 2.8\))
  connects to enzyme allostery (\cref{sec:unit_I_enzymes_and_kinetics})
  and protein structure (\cref{sec:unit_I_macromolecules}).
\item
  \textbf{Immune system} connects to Unit VII (innate response to
  bacteria/viruses) and Unit IV (recombination of V(D)J gene segments in
  antibody diversity).
\item
  \textbf{Endocrine regulation} (steroid hormones acting on nuclear
  receptors) links to Unit IV (gene regulation) and Unit III (metabolic
  regulation by insulin/glucagon vs.~AMPK).
\item
  \textbf{Negative feedback homeostasis} is mathematically equivalent to
  the control theory underlying Unit X (population regulation) and Unit
  III (metabolic feedback).
\end{itemize}

\begin{quote}
\textbf{Key vocabulary introduced here:} homeostasis, action potential,
resting membrane potential, ion channel, neurotransmitter, synapse,
GPCR, second messenger, cardiac output, stroke volume, preload,
afterload, Frank-Starling law, haemoglobin cooperativity, Bohr effect,
hormone, receptor tyrosine kinase, antigen, MHC, clonal selection,
innate immunity, adaptive immunity.
\end{quote}

\subsection*{Computational Toolbox --- Unit
IX}\label{computational-toolbox-unit-ix}
\addcontentsline{toc}{subsection}{Computational Toolbox --- Unit IX}

\begin{Shaded}
\begin{Highlighting}[]
\ImportTok{from}\NormalTok{ biology.physiology }\ImportTok{import}\NormalTok{ oxygen\_saturation}
\ImportTok{from}\NormalTok{ biology.neuroscience }\ImportTok{import}\NormalTok{ action\_potential\_hh}

\CommentTok{\# Haemoglobin oxygen saturation (Hill equation): arterial vs venous}
\CommentTok{\# Hill coefficient n=2.7; P50=26 mmHg for adult HbA at 37°C, pH 7.4}
\NormalTok{arterial }\OperatorTok{=}\NormalTok{ oxygen\_saturation(pO2\_mmHg}\OperatorTok{=}\DecValTok{100}\NormalTok{, p50\_mmHg}\OperatorTok{=}\DecValTok{26}\NormalTok{, hill\_coefficient}\OperatorTok{=}\FloatTok{2.7}\NormalTok{)}
\NormalTok{venous }\OperatorTok{=}\NormalTok{ oxygen\_saturation(pO2\_mmHg}\OperatorTok{=}\DecValTok{40}\NormalTok{, p50\_mmHg}\OperatorTok{=}\DecValTok{26}\NormalTok{, hill\_coefficient}\OperatorTok{=}\FloatTok{2.7}\NormalTok{)}
\BuiltInTok{print}\NormalTok{(}\SpecialStringTok{f"Arterial SO2  (PO2=100 mmHg): }\SpecialCharTok{\{}\NormalTok{arterial}\SpecialCharTok{.}\NormalTok{saturation}\SpecialCharTok{:.1\%\}}\SpecialStringTok{"}\NormalTok{)}
\BuiltInTok{print}\NormalTok{(}\SpecialStringTok{f"Venous   SO2  (PO2=40  mmHg): }\SpecialCharTok{\{}\NormalTok{venous}\SpecialCharTok{.}\NormalTok{saturation}\SpecialCharTok{:.1\%\}}\SpecialStringTok{"}\NormalTok{)}
\BuiltInTok{print}\NormalTok{(}\SpecialStringTok{f"O2 extraction ratio: }\SpecialCharTok{\{}\NormalTok{(arterial.saturation }\OperatorTok{{-}}\NormalTok{ venous.saturation) }\OperatorTok{/}\NormalTok{ arterial}\SpecialCharTok{.}\NormalTok{saturation}\SpecialCharTok{:.1\%\}}\SpecialStringTok{"}\NormalTok{)}
\CommentTok{\# Expected:}
\CommentTok{\# Arterial SO2:  97.0\%}
\CommentTok{\# Venous   SO2:  76.2\%}
\CommentTok{\# O2 extraction: 21.7\%  (about 5 mL O2 per 100 mL blood at rest)}

\CommentTok{\# Hodgkin{-}Huxley action potential simulation (10 µA/cm² stimulus, 50 ms)}
\NormalTok{hh }\OperatorTok{=}\NormalTok{ action\_potential\_hh(stimulus\_current\_µA}\OperatorTok{=}\FloatTok{10.0}\NormalTok{, t\_end\_ms}\OperatorTok{=}\FloatTok{50.0}\NormalTok{)}
\NormalTok{peak\_V }\OperatorTok{=} \BuiltInTok{max}\NormalTok{(hh.voltage\_mV)}
\BuiltInTok{print}\NormalTok{(}\SpecialStringTok{f"Peak membrane voltage: }\SpecialCharTok{\{}\NormalTok{peak\_V}\SpecialCharTok{:.1f\}}\SpecialStringTok{ mV"}\NormalTok{)}
\BuiltInTok{print}\NormalTok{(}\SpecialStringTok{f"Firing threshold crossed: }\SpecialCharTok{\{}\StringTok{\textquotesingle{}yes\textquotesingle{}} \ControlFlowTok{if}\NormalTok{ peak\_V }\OperatorTok{\textgreater{}} \DecValTok{0} \ControlFlowTok{else} \StringTok{\textquotesingle{}no\textquotesingle{}}\SpecialCharTok{\}}\SpecialStringTok{"}\NormalTok{)}
\CommentTok{\# Expected:}
\CommentTok{\# Peak membrane voltage: about +48.0 mV (overshoot above 0 mV)}
\CommentTok{\# Firing threshold crossed: yes}
\end{Highlighting}
\end{Shaded}

\begin{quote}
\textbf{Try it yourself:} Reduce \texttt{stimulus\_current\_\texorpdfstring{\ensuremath{\mu}}{mu}A} to 5
\texorpdfstring{\ensuremath{\mu}}{mu}A/cm\texorpdfstring{\textsuperscript{2}}{2} --- the membrane depolarises but may not reach threshold. This
illustrates the \textbf{most-or-none principle}.
\end{quote}

\begin{center}\rule{0.5\linewidth}{0.5pt}\end{center}

\emph{Source modules: \breaktt{src/biology/physiology/} and
\breaktt{src/biology/neuroscience/} --- \breaktt{oxygen\_saturation()},
\texttt{poiseuille\_flow()}, \breaktt{homeostasis\_response()},
\breaktt{action\_potential\_hh()}.} \emph{Figures:
\breaktt{src/visualization/} (action potential traces, haemoglobin
saturation curves, hormone feedback oscillators);
\breaktt{src/mermaid/biology\_diagrams.py} (nervous system diagrams,
immune response cascades).}

\newpage

\section{Circulation, Respiration, and
Homeostasis}\label{circulation-respiration-and-homeostasis}

\label{sec:unit_IX_circulation_respiration_homeostasis}

\begin{quote}
\textbf{Ch 28} \texorpdfstring{\ensuremath{\cdot}}{.} Level 3/3 \texorpdfstring{\ensuremath{\cdot}}{.} 60 min read \texorpdfstring{\ensuremath{\cdot}}{.} 100 min lecture \texorpdfstring{\ensuremath{\cdot}}{.}
Prerequisites: \cref{sec:unit_II_membrane_transport},
\cref{sec:unit_III_bioenergetics_and_respiration}
\end{quote}

\subsection{Learning Objectives}\label{learning-objectives-30}

\begin{enumerate}
\def\labelenumi{\arabic{enumi}.}
\tightlist
\item
  Describe the evolutionary progression of cardiovascular systems from
  open to closed, and from 2-chamber to 4-chamber hearts.
\item
  Explain the cardiac cycle including electrical conduction, mechanical
  events, and the ECG.
\item
  Apply Poiseuille's law to blood flow and explain arteriolar regulation
  of vascular resistance.
\item
  Describe oxygen transport by haemoglobin using the Hill equation,
  including the Bohr and Haldane effects.
\item
  Explain CO\(_2\) transport mechanisms and the chloride shift.
\item
  Describe pulmonary ventilation mechanics, lung volumes, and
  spirometry.
\item
  Apply Fick's law to alveolar gas exchange.
\item
  Explain respiratory control by medullary centres and chemoreceptors.
\item
  Describe homeostatic control systems,
  \hyperref[gl:thermoregulation]{\textbf{thermoregulation}}, and fluid
  balance.
\end{enumerate}

\subsubsection{Study Blueprint}\label{study-blueprint-30}

\begin{itemize}
\tightlist
\item
  \textbf{Big idea:} Animal transport systems maintain gradients that
  let cells exchange gases, nutrients, heat, and wastes.
\item
  \textbf{Core concepts:} cardiac output, gas exchange, homeostasis,
  feedback.
\item
  \textbf{Framework alignment:} Vision \& Change: Structure and
  function, Systems; AP Biology: Systems Interactions, Energetics;
  NGSS-style topics: Structure and Function.
\item
  \textbf{Model or quantitative lens:} Cardiac-output, diffusion,
  oxygen-saturation, and feedback calculations.
\item
  \textbf{Data skill:} Interpret physiological data from pressure, flow,
  saturation, or set-point changes.
\item
  \textbf{Practice cadence:} Visual Representations, Statistical Tests
  and Data Analysis, Argumentation.
\item
  \textbf{Common misconception to repair:} Homeostasis is dynamic
  regulation, not an unchanging internal state.
\item
  \textbf{Primary lab:}
  \cref{sec:lab_unit_IX_circulation_respiration_homeostasis}.
\item
  \textbf{Question bank:}
  \cref{sec:q_unit_IX_circulation_respiration_homeostasis}.
\item
  \textbf{Transfer task:} Transfer homeostatic reasoning to exercise,
  altitude, hemorrhage, fever, and shock.
\item
  \textbf{Bridge to computation:}
  \breaktt{biology.physiology.physiology.oxygen\_saturation}.
\end{itemize}

\begin{center}\rule{0.5\linewidth}{0.5pt}\end{center}

\begin{quote}
\textbf{Opening Vignette --- The Doctor Who Proved the Heart Was a Pump}

In 1628, William Harvey published \emph{Exercitatio Anatomica de Motu
Cordis et Sanguinis} --- a slim 72-page monograph that demolished 1,400
years of medical dogma. The \hyperref[gl:dominant]{\textbf{dominant}}
Galenic view held that blood was continuously manufactured from food in
the liver and consumed by tissues. Harvey's insight was quantitative: he
calculated the volume expelled by the heart with each beat (about 60 mL)
and multiplied by heart rate. The heart expels roughly 4.3 kg of blood
per minute --- far more than the body could possibly generate from food.
The primarily explanation was that blood circulates, continuously
recirculated by the heart. Harvey's argument was revolutionary in method
as much as finding: he used mathematical reasoning, vivisection, and
comparative anatomy across many species rather than received authority.
He was ridiculed widely and lost patients; he died in 1657. Modern
cardiology, cardiac surgery, and pharmacology of the cardiovascular
system most stand on the foundation he built alone.
\end{quote}

\subsection{Evolution of Cardiovascular
Systems}\label{evolution-of-cardiovascular-systems}

The cardiovascular system evolved to overcome the limitations of
diffusion for O\(_2\) and nutrient delivery in large organisms:

\textbf{Open circulatory systems} (arthropods, most molluscs): Heart
pumps haemolymph into open sinuses (haemocoel). Haemolymph bathes
tissues directly. Low pressure, slow flow. Adequate for small ectotherms
with low metabolic rates.

\textbf{Closed circulatory systems} (annelids, cephalopods, most
vertebrates): Blood confined within vessels. Higher pressure enables
faster, more directed flow. Allows separate regulation of blood flow to
different organs.

\textbf{Vertebrate heart evolution:}

{\def\LTcaptype{none} % do not increment counter
\begin{longtable}[]{@{}
  >{\raggedright\arraybackslash}p{(\linewidth - 6\tabcolsep) * \real{0.2564}}
  >{\raggedright\arraybackslash}p{(\linewidth - 6\tabcolsep) * \real{0.2051}}
  >{\raggedright\arraybackslash}p{(\linewidth - 6\tabcolsep) * \real{0.2308}}
  >{\raggedright\arraybackslash}p{(\linewidth - 6\tabcolsep) * \real{0.3077}}@{}}
\toprule\noalign{}
\begin{minipage}[b]{\linewidth}\raggedright
Heart Type
\end{minipage} & \begin{minipage}[b]{\linewidth}\raggedright
Chambers
\end{minipage} & \begin{minipage}[b]{\linewidth}\raggedright
Organisms
\end{minipage} & \begin{minipage}[b]{\linewidth}\raggedright
Key Features
\end{minipage} \\
\midrule\noalign{}
\endhead
\bottomrule\noalign{}
\endlastfoot
2-chamber & 1 atrium + 1 ventricle & Fish & Single circuit: heart to
gills to body to heart \\
3-chamber & 2 atria + 1 ventricle & Amphibians, most reptiles & Double
circuit with some mixing in ventricle \\
3.5-chamber & Partial septum & Crocodilians & Nearly complete
ventricular separation \\
4-chamber & 2 atria + 2 ventricles & Birds, mammals & Complete
separation of pulmonary and systemic circuits \\
\end{longtable}
}

The 4-chamber heart enables complete separation of oxygenated and
deoxygenated blood, supporting the high metabolic rates required for
endothermy.

\begin{center}\rule{0.5\linewidth}{0.5pt}\end{center}

\subsection{Heart Anatomy and the Cardiac
Cycle}\label{heart-anatomy-and-the-cardiac-cycle}

\subsubsection{Cardiac Conduction
System}\label{cardiac-conduction-system}

The heart is a self-exciting organ. Its electrical conduction system
ensures coordinated contraction:

\begin{figure}[htbp]
\centering
\includegraphics[width=0.82\linewidth,height=4.2in,keepaspectratio]{../figures/mermaid_inline/inline_mermaid_0152_78bc1cfd0164.png}
\caption{Cardiac conduction system and corresponding ECG waves. The SA node initiates depolarisation, which spreads through the atria (P wave), is delayed at the AV node, then rapidly conducted through the bundle of His and Purkinje fibres to the ventricles (QRS complex). Ventricular repolarisation produces the T wave.}
\end{figure}

\textbf{ECG waves and their significance:} - \textbf{P wave:} Atrial
depolarisation and contraction - \textbf{PR interval:} Time from atrial
depolarisation to ventricular depolarisation (0.12-0.20 s). Prolonged PR
= AV block. - \textbf{QRS complex:} Ventricular depolarisation
(0.06-0.10 s). Wide QRS = conduction defect or ventricular origin. -
\textbf{T wave:} Ventricular repolarisation - \textbf{QT interval:}
Total ventricular electrical activity. Prolonged QT = risk of torsade de
pointes arrhythmia.

\subsubsection{The Cardiac Cycle}\label{the-cardiac-cycle}

At a heart rate of 70 bpm, one cardiac cycle lasts about 0.8
seconds:

\begin{itemize}
\tightlist
\item
  \textbf{Systole} (contraction, about 0.3 s):

  \begin{itemize}
  \tightlist
  \item
    Isovolumetric contraction: Most valves closed, ventricular pressure
    rises rapidly
  \item
    Ejection phase: Ventricular pressure exceeds aortic pressure
    (about 80 mmHg), aortic valve opens, blood ejected
  \item
    Peak systolic pressure: about 120 mmHg in left ventricle
  \end{itemize}
\item
  \textbf{Diastole} (relaxation, about 0.5 s):

  \begin{itemize}
  \tightlist
  \item
    Isovolumetric relaxation: Most valves closed, ventricular pressure
    falls
  \item
    Rapid filling: Ventricular pressure falls below atrial pressure, AV
    valve opens, blood flows passively (accounts for about 80\%
    of filling)
  \item
    Atrial systole (``atrial kick''): Final about 20\% of
    ventricular filling
  \end{itemize}
\end{itemize}

\textbf{Starling's Law of the Heart (Frank-Starling \citep{starling1914}
mechanism):} Ventricular stroke volume increases with end-diastolic
volume (preload). Increased venous return stretches ventricular
cardiomyocytes, optimising \hyperref[gl:actin]{\textbf{actin}}-myosin
overlap and increasing Ca\(^{2+}\) sensitivity of troponin C. This
intrinsic mechanism ensures that left and right cardiac outputs remain
matched.

\begin{equation}
SV \propto EDV \quad \text{(within physiological range)}
\label{eq:circulation_respiration_homeostasis_1}
\end{equation}

\begin{quote}
\textbf{Clinical Connection:} Heart failure occurs when the
Frank-Starling mechanism can no longer compensate for reduced
contractility. In systolic heart failure (HFrEF, ejection fraction
\textless40\%), the weakened ventricle operates on a depressed Starling
curve. Treatments include ACE inhibitors (reduce afterload),
beta-blockers (reduce myocardial demand), and SGLT2 inhibitors (reduce
preload). The 2024 ESC guidelines emphasise early combination therapy.
\end{quote}

\begin{center}\rule{0.5\linewidth}{0.5pt}\end{center}

\subsection{Cardiac Output and Blood
Pressure}\label{cardiac-output-and-blood-pressure}

\subsubsection{Cardiac Output}\label{cardiac-output}

\begin{equation}
CO = HR \times SV
\label{eq:unit_IX_cardiac_output}
\end{equation}

\begin{itemize}
\tightlist
\item
  At rest:
  \(CO = 70\;\text{bpm} \times 70\;\text{mL} \approx 5\;\text{L/min}\)
\item
  Maximum exercise:
  \(CO = 200\;\text{bpm} \times 130\;\text{mL} \approx 26\;\text{L/min}\)
  (elite athletes --- up to 35 L/min documented)
\end{itemize}

\textbf{The Frank-Starling mechanism} \citep{starling1914} is the
heart's intrinsic capacity to match its output to venous return on a
beat-to-beat basis. Within physiological limits, increased end-diastolic
volume (preload) stretches ventricular cardiomyocytes, optimising
actin-myosin overlap (length-tension relationship) and increasing the
Ca\(^{2+}\)-sensitivity of troponin C. The result is an increased stroke
volume \emph{without} any extrinsic neural or hormonal command. Because
the right and left hearts are in series, this intrinsic matching is
essential --- any sustained mismatch would cause pulmonary or systemic
congestion. Acutely transfusing fluid or rising from supine to legs-up
posture both expand venous return, and the Frank-Starling mechanism
translates that into matched increases in both ventricles.

\textbf{Regulation of cardiac output:}

\begin{itemize}
\tightlist
\item
  \textbf{Sympathetic stimulation} (\(\beta_1\) adrenergic receptors):
  Increases both HR (positive chronotropy via SA node) and SV (positive
  inotropy via increased Ca\(^{2+}\) entry through L-type Ca\(^{2+}\)
  channels and SR Ca\(^{2+}\) release)
\item
  \textbf{Parasympathetic stimulation} (M\(_2\) muscarinic receptors via
  vagus nerve): Decreases HR primarily (negative chronotropy). Vagal
  tone dominates at rest, which is why resting HR is about 70
  bpm rather than the SA node's intrinsic 100 bpm.
\item
  \textbf{Frank-Starling mechanism:} Intrinsic matching of SV to venous
  return (see above)
\end{itemize}

\subsubsection{Blood Pressure}\label{blood-pressure}

\begin{equation}
MAP = DP + \frac{1}{3}PP = DP + \frac{1}{3}(SP - DP)
\label{eq:circulation_respiration_homeostasis_3}
\end{equation}

where MAP = mean arterial pressure, DP = diastolic pressure, SP =
systolic pressure, PP = pulse pressure.

Normal: SP/DP = 120/80 mmHg, MAP about 93 mmHg.

\begin{equation}
MAP = CO \times TPR
\label{eq:circulation_respiration_homeostasis_4}
\end{equation}

where TPR = total peripheral resistance. Blood pressure is regulated by:
- Cardiac output (HR, SV, blood volume) - Total peripheral resistance
(arteriolar diameter -- the dominant factor)

\subsubsection{Blood Flow -- Poiseuille's
Law}\label{blood-flow-poiseuilles-law}

\begin{equation}
Q = \frac{\pi r^4 \Delta P}{8 \eta L}
\label{eq:circulation_respiration_homeostasis_5}
\end{equation}

The \textbf{r\(^4\) dependence} is the most physiologically critical
feature: halving vessel radius reduces flow by 16-fold. Arterioles
(radius about 50-100 um) serve as the primary resistance
vessels.

\textbf{Regulation of arteriolar tone:} - \textbf{Local autoregulation:}
- Myogenic response: Smooth muscle contracts in response to stretch
(maintains constant flow despite pressure changes) - Metabolic:
CO\(_2\), H\(^+\), K\(^+\), adenosine from active tissues cause
vasodilation - Endothelial: Shear stress stimulates NO (nitric oxide)
production via eNOS. NO diffuses to smooth muscle, activates guanylyl
cyclase to cGMP to PKG to smooth muscle relaxation. - \textbf{Neural
(sympathetic):} Noradrenaline acts on \(\alpha_1\) receptors causing
vasoconstriction in most vascular beds. Exceptions: skeletal muscle
during exercise (cholinergic vasodilation and metabolic override). -
\textbf{Hormonal:} Angiotensin II (potent vasoconstrictor), ANP
(vasodilator), ADH/vasopressin (vasoconstrictor at high doses).

\textbf{Concept Check:} During exercise, cardiac output increases from 5
to 25 L/min, but blood pressure increases primarily modestly (e.g.,
120/80 to 160/80). How is this possible? (Consider what happens to TPR.)

\subsubsection{Worked Example: Cardiac
Reserve}\label{worked-example-cardiac-reserve}

A healthy 40-year-old has resting cardiac output of 5 L/min (HR 70 bpm \texorpdfstring{\ensuremath{\times}}{x}
SV 70 mL). The \textbf{age-predicted maximum heart rate} is given by:

\begin{equation}
HR_{\max} \approx 220 - \text{age (years)}
\label{eq:unit_IX_hrmax}
\end{equation}

(The Tanaka 2001 formula \(HR_{\max} = 208 - 0.7 \cdot \text{age}\) is
more accurate for older adults and is now preferred in cardiology.)

For our 40-year-old, \(HR_{\max} \approx 180\) bpm. With
exercise-induced sympathetic activation increasing stroke volume to
about 120 mL via the \textbf{Frank-Starling mechanism}
\citep{starling1914} and \texorpdfstring{\ensuremath{\uparrow}}{up}contractility (\texorpdfstring{\ensuremath{\beta}}{beta}\(_1\)):

\[CO_{\max} = 180 \times 120 = 21{,}600\;\text{mL/min} \approx 21.6\;\text{L/min} \label{eq:unit_IX_circulation_respiration_homeostasis_item_1}\]

This represents a \textbf{about 4.3-fold increase} over rest
--- typical for an untrained adult. Endurance training increases the SV
reserve more than the HR reserve, allowing elite athletes to reach
30--35 L/min cardiac output despite \emph{lower} maximum heart rates
than untrained subjects.

\textbf{Cardiac reserve
\(= CO_{\max} - CO_{\text{rest}} \approx 16.6\;\text{L/min}\)} in this
example. Cardiac reserve is the master variable that limits exercise
tolerance: heart-failure patients with reduced ejection fraction may
have a reserve of 3--5 L/min and become symptomatic with mild activity.

\subsubsection{Baroreceptor Reflex --- Beat-to-Beat Blood Pressure
Control}\label{baroreceptor-reflex-beat-to-beat-blood-pressure-control}

The arterial \textbf{baroreceptor reflex} is the body's primary
short-term blood pressure stabiliser, operating on a beat-to-beat
timescale (seconds). Stretch-sensitive mechanoreceptors in the
\textbf{carotid sinus} (innervated by the glossopharyngeal nerve, CN IX)
and \textbf{aortic arch} (innervated by the vagus, CN X) increase their
firing rate when arterial wall stretch increases.

\begin{figure}[htbp]
\centering
\includegraphics[width=0.82\linewidth,height=4.2in,keepaspectratio]{../figures/mermaid_inline/inline_mermaid_0153_b805ffa795c5.png}
\caption{Baroreceptor reflex Pressure rise stretches arterial baroreceptors, increasing firing to the NTS in the medulla. The NTS inhibits the rostral ventrolateral medulla (RVLM, sympathetic premotor neurons) and excites the nucleus ambiguus (parasympathetic). Reduced sympathetic and increased parasympathetic outflow lower HR, contractility, and TPR - restoring normal MAP within seconds.}
\end{figure}

The reflex is \textbf{bidirectional}: low MAP unloads baroreceptors \texorpdfstring{\ensuremath{\to}}{->}
reduced firing \texorpdfstring{\ensuremath{\to}}{->} disinhibition of sympathetic outflow \texorpdfstring{\ensuremath{\to}}{->} tachycardia,
vasoconstriction, increased contractility. \textbf{Orthostatic
hypotension} results from impaired baroreflex (autonomic neuropathy in
diabetes, ageing, or pure autonomic failure).

\textbf{Baroreceptor adaptation:} With sustained hypertension over hours
to days, baroreceptors \textbf{reset} to the elevated pressure --- they
fire less for a given pressure than in a normotensive person. This is
why chronic hypertension is not corrected by the reflex: the system has
adapted to defend the new (elevated) set point. Long-term blood pressure
control requires the renin-angiotensin-aldosterone system (RAAS) and
renal pressure-natriuresis, not the baroreflex.

\subsubsection{RAAS Cascade --- Long-term Blood Pressure and Volume
Control}\label{raas-cascade-long-term-blood-pressure-and-volume-control}

The \textbf{renin-angiotensin-aldosterone system} is the dominant
long-term BP regulator, operating over hours to days through changes in
vascular tone, sodium retention, and blood volume.

\begin{equation}
\text{Angiotensinogen} \xrightarrow{\text{renin}} \text{Angiotensin I} \xrightarrow{\text{ACE}} \text{Angiotensin II}
\label{eq:circulation_respiration_homeostasis_raas}
\end{equation}

\textbf{Trigger signals for renin release} from juxtaglomerular cells in
afferent arterioles: 1. \textbf{Reduced renal perfusion pressure}
(sensed by JG cells as stretch \texorpdfstring{\ensuremath{\downarrow}}{down}) 2. \textbf{Reduced NaCl delivery} to
the macula densa (sensed in distal tubule via NKCC2) 3.
\textbf{Sympathetic activation} (\texorpdfstring{\ensuremath{\beta}}{beta}\(_1\) adrenergic stimulation of JG
cells)

\textbf{Angiotensin II effects (multi-organ):} -
\textbf{Vasoconstriction} at AT\(_1\) receptors (potent direct effect on
arterioles \texorpdfstring{\ensuremath{\to}}{->} \texorpdfstring{\ensuremath{\uparrow}}{up}TPR \texorpdfstring{\ensuremath{\to}}{->} \texorpdfstring{\ensuremath{\uparrow}}{up}MAP) - \textbf{Aldosterone release} from adrenal
zona glomerulosa \texorpdfstring{\ensuremath{\to}}{->} renal Na\(^+\) reabsorption (ENaC) and K\(^+\)
secretion \texorpdfstring{\ensuremath{\to}}{->} \texorpdfstring{\ensuremath{\uparrow}}{up}blood volume - \textbf{ADH release} from posterior
pituitary \texorpdfstring{\ensuremath{\to}}{->} renal water reabsorption (aquaporin-2) \texorpdfstring{\ensuremath{\to}}{->} \texorpdfstring{\ensuremath{\uparrow}}{up}blood volume -
\textbf{Sympathetic facilitation} (central and peripheral) -
\textbf{Thirst} (subfornical organ) - \textbf{Direct renal effects:}
preferential constriction of efferent arteriole (preserves GFR during
hypoperfusion); proximal tubule Na\(^+\) reabsorption

\begin{quote}
\textbf{Clinical Connection --- RAAS pharmacology:} ACE inhibitors
(lisinopril, enalapril) block conversion of Ang I to Ang II. ARBs
(losartan, valsartan) block AT\(_1\) receptors. Direct renin inhibitors
(aliskiren) block at the top. Aldosterone antagonists (spironolactone,
eplerenone) block the mineralocorticoid receptor. These four classes are
foundational for hypertension, heart failure (HFrEF), diabetic
nephropathy, and post-MI ventricular remodelling. The 2024 PARADIGM-HF
trials established \textbf{ARNI (angiotensin receptor-neprilysin
inhibitor; sacubitril/valsartan)} as superior to ACE inhibition alone in
HFrEF --- neprilysin inhibition prolongs natriuretic peptides (ANP/BNP)
while ARB blocks AT\(_1\).
\end{quote}

\subsection{Digestive, Nutritional, Renal, and Excretory
Integration}\label{digestive-nutritional-renal-and-excretory-integration}

Animal homeostasis also depends on how the gut, liver, kidney, and
excretory surfaces couple intake to internal composition. Digestion is
not just ``food breakdown''; it is staged chemical and mechanical
processing: stomach acid denatures proteins and limits microbes,
pancreatic enzymes and bile complete macromolecule hydrolysis in the
small intestine, and villi plus microvilli expand absorptive surface for
monosaccharides, amino acids, lipids, vitamins, and ions
\citep{niddk2024digestivesystem}. Nutritional status then becomes a
systems variable. A high-protein meal changes hepatic urea production
and renal nitrogen excretion; a low-salt or dehydrating environment
recruits RAAS and ADH; and malnutrition weakens immune barriers, wound
repair, growth, and reproductive function \citep{fao2025sofi}.

The vertebrate kidney stabilises plasma volume, osmolality, pH,
electrolytes, and nitrogen balance by filtering plasma, reclaiming
useful solutes, secreting selected wastes, and concentrating urine
\citep{niddk2024kidneys}. Filtration at the glomerulus is
pressure-driven; proximal tubules reclaim most filtered Na\(^+\),
glucose, amino acids, and bicarbonate; the loop of Henle builds the
corticomedullary osmotic gradient; distal nephron segments fine-tune
Na\(^+\), K\(^+\), acid-base, and water balance under aldosterone and
ADH. Across animals, excretory designs solve the same problem with
different constraints: aquatic fishes can excrete ammonia directly,
mammals convert nitrogen to urea, and birds/reptiles conserve water by
excreting uric acid. The unifying claim is comparative and mechanistic:
excretion trades ATP, water, and toxicity risk against the environment
an organism inhabits.

\begin{center}\rule{0.5\linewidth}{0.5pt}\end{center}

\subsection{Capillary Fluid Exchange and the Lymphatic
System}\label{capillary-fluid-exchange-and-the-lymphatic-system}

\subsubsection{Starling Forces}\label{starling-forces}

At the capillary level, O\(_2\) and nutrients are delivered to tissues
while CO\(_2\) and waste products are removed. Fluid exchange across the
capillary wall is governed by the \textbf{Starling forces} (Ernest
Starling, 1896):

\begin{equation}
J_v = L_p \left[(P_c - P_{if}) - \sigma(\pi_c - \pi_{if})\right]
\label{eq:circulation_respiration_homeostasis_6}
\end{equation}

where: - \(J_v\) = net fluid flux (positive = filtration out of
capillary) - \(L_p\) = hydraulic conductivity of the capillary wall -
\(P_c\) = capillary hydrostatic pressure - \(P_{if}\) = interstitial
fluid hydrostatic pressure - \texorpdfstring{\ensuremath{\sigma}}{sigma} = reflection coefficient (0 = fully
permeable to \hyperref[gl:protein]{\textbf{protein}}; 1 = impermeable) -
\(\pi_c\) = capillary oncotic pressure (colloid osmotic pressure, from
plasma proteins, mainly albumin) - \(\pi_{if}\) = interstitial oncotic
pressure

\textbf{Typical values at the arteriolar end of a skeletal muscle
capillary:}

{\def\LTcaptype{none} % do not increment counter
\begin{longtable}[]{@{}
  >{\raggedright\arraybackslash}p{(\linewidth - 4\tabcolsep) * \real{0.1852}}
  >{\raggedright\arraybackslash}p{(\linewidth - 4\tabcolsep) * \real{0.4815}}
  >{\raggedright\arraybackslash}p{(\linewidth - 4\tabcolsep) * \real{0.3333}}@{}}
\toprule\noalign{}
\begin{minipage}[b]{\linewidth}\raggedright
Force
\end{minipage} & \begin{minipage}[b]{\linewidth}\raggedright
Value (mmHg)
\end{minipage} & \begin{minipage}[b]{\linewidth}\raggedright
Direction
\end{minipage} \\
\midrule\noalign{}
\endhead
\bottomrule\noalign{}
\endlastfoot
Capillary hydrostatic pressure (\(P_c\)) & 35 & Out (filtration) \\
Interstitial hydrostatic pressure (\(P_{if}\)) & \(-3\) & Out
(filtration) \\
Capillary oncotic pressure (\(\pi_c\)) & 26 & In (reabsorption) \\
Interstitial oncotic pressure (\(\pi_{if}\)) & 1 & Out (filtration) \\
\end{longtable}
}

\textbf{Net filtration pressure (arteriolar end):}

\begin{equation}
\text{NFP} = (P_c + \pi_{if}) - (\pi_c + P_{if}) = (35 + 1) - (26 + (-3)) = 36 - 23 = +13\;\text{mmHg}
\label{eq:circulation_respiration_homeostasis_7}
\end{equation}

Positive NFP: fluid filters \textbf{out} of capillary at the arteriolar
end.

\textbf{At the venular end} (\(P_c\) drops to about 15 mmHg due
to resistance loss):

\begin{equation}
\text{NFP} = (15 + 1) - (26 + (-3)) = 16 - 23 = -7\;\text{mmHg}
\label{eq:circulation_respiration_homeostasis_8}
\end{equation}

Negative NFP: fluid is \textbf{reabsorbed} into capillary at the venular
end.

\textbf{Net result:} about 90\% of filtered fluid is reabsorbed
at the venular end. The remaining about 10\% (about 2--3 L/day)
must be returned via the \textbf{lymphatic system}.

\subsubsection{The Lymphatic System}\label{the-lymphatic-system}

The lymphatic system drains excess interstitial fluid and returned
plasma proteins that leak out of capillaries:

\begin{itemize}
\tightlist
\item
  \textbf{Blind-ended lymphatic capillaries} in tissues are highly
  permeable; interstitial fluid + proteins enter by bulk flow
\item
  Lymph flows through lymph vessels (propelled by skeletal muscle
  compression and one-way valves) into lymph nodes
\item
  Lymph nodes filter lymph, removing pathogens and debris, and add
  lymphocytes
\item
  Right lymphatic duct (drains right chest, arm, head) and thoracic duct
  (drains everything else) return lymph to venous circulation
  (subclavian veins)
\item
  Flow: about 2--3 L/day; lymph protein concentration
  about 2 g/dL (compared with plasma about 7 g/dL)
\end{itemize}

\begin{quote}
\textbf{Clinical Connection:} \textbf{Lymphoedema} occurs when lymphatic
return is impaired, causing protein-rich fluid to accumulate in the
interstitium. Causes: filariasis (\emph{Wuchereria bancrofti}, the
leading worldwide cause), surgical lymph node dissection (common after
breast cancer surgery), or radiation damage. Unlike simple oedema
(protein-poor), lymphoedema is high-protein, making it prone to skin
fibrosis, infection (cellulitis), and functional impairment.
\textbf{Pitting oedema} (non-lymphoedema) results from elevated \(P_c\)
(heart failure, portal hypertension) or reduced \(\pi_c\)
(hypoalbuminaemia from cirrhosis or nephrotic syndrome).
\end{quote}

\subsection{Worked Example: Hypoalbuminaemia and Net Filtration
Pressure}\label{worked-example-hypoalbuminaemia-and-net-filtration-pressure}

A patient has serum albumin of 2 g/dL (normal 4 g/dL), reducing
\(\pi_c\) from 26 mmHg to 13 mmHg (oncotic pressure \(\propto\) protein
concentration). At the arteriolar end (\(P_c\) = 35 mmHg):

\begin{equation}
\text{NFP} = (35 + 1) - (13 + (-3)) = 36 - 10 = +26\;\text{mmHg}
\label{eq:circulation_respiration_homeostasis_9}
\end{equation}

And at the venular end (\(P_c\) = 15 mmHg):

\begin{equation}
\text{NFP} = (15 + 1) - (13 + (-3)) = 16 - 10 = +6\;\text{mmHg}
\label{eq:circulation_respiration_homeostasis_10}
\end{equation}

Both ends now show net filtration --- fluid cannot be reabsorbed,
leading to progressive oedema. The lymphatics are overwhelmed. This
explains the \textbf{anasarca} (generalised oedema) seen in severe
hypoalbuminaemia from nephrotic syndrome or liver failure.

\textbf{Concept Check:} Using the Starling equation with arteriolar
\(P_c = 35\) mmHg, \(P_{if} = -3\) mmHg, \(\pi_{if} = 1\) mmHg, predict
the sign of the net filtration pressure when plasma albumin halves
(\(\pi_c\): 26 -> 13 mmHg), and explain why nephrotic-syndrome oedema
arises even though capillary hydrostatic pressure is unchanged. (Hint:
compute NFP at both the arteriolar and venular ends before and after the
\(\pi_c\) fall.)

\begin{center}\rule{0.5\linewidth}{0.5pt}\end{center}

\subsection{Oxygen Transport}\label{oxygen-transport}

\begin{figure}[htbp]
\centering
\includegraphics[width=0.85\textwidth]{../figures/oxygen_dissociation_curve.png}
\caption{Oxygen--haemoglobin dissociation curves showing percent saturation versus $p\mathrm{O}_2$ for normal adult haemoglobin, a right-shifted fever/exercise condition, and a left-shifted fetal-haemoglobin condition.}
\label{fig:unit_IX_oxygen_dissociation}
\end{figure}

\subsubsection{Haemoglobin Structure}\label{haemoglobin-structure}

Haemoglobin (Hb) is an \(\alpha_2\beta_2\) tetramer. Each subunit
contains a \textbf{haeme} group (iron-porphyrin; Fe\(^{2+}\)) that binds
one O\(_2\). Oxygen-carrying capacity: 1.34 mL O\(_2\) per gram Hb;
typical Hb concentration: 150 g/L blood. Total: about 201 mL
O\(_2\)/L blood bound to Hb. Plasma dissolved O\(_2\): 0.003 \(\times\)
PO\(_2\) (mL/L/mmHg), giving about 0.3 mL/L at PO\(_2\) = 100 mmHg. Hb
carries about 670\(\times\) more O\(_2\) than plasma alone.

\subsubsection{\texorpdfstring{O\(_2\)-Hb Dissociation Curve and the
Hill
Equation}{O_2-Hb Dissociation Curve and the Hill Equation}}\label{o_2-hb-dissociation-curve-and-the-hill-equation}

\begin{figure}[htbp]
\centering
\includegraphics[width=0.82\linewidth,height=4.2in,keepaspectratio]{../figures/mermaid_inline/inline_mermaid_0154_aa7954aaccab.png}
\caption{Oxygen delivery from lung to tissue and the Bohr effect. In the lungs (high PO\$\_2\$), haemoglobin loads O\$\_2\$ to 97\% saturation. In tissues (lower PO\$\_2\$), O\$\_2\$ is released. The Bohr effect rightward shift of the dissociation curve in metabolically active tissues (low pH, high CO\$\_2\$, high temperature) enhances O\$\_2\$ delivery where it is most needed.}
\end{figure}

Hb-O\(_2\) binding is \textbf{cooperative} (T-state to R-state
conformational switch). The resulting sigmoidal binding curve, together
with physiological affinity shifts
(\cref{fig:unit_IX_oxygen_dissociation}), follows the Hill equation:

\begin{equation}
SO_2 = \frac{(PO_2 / P_{50})^n}{1 + (PO_2 / P_{50})^n}
\label{eq:circulation_respiration_homeostasis_11}
\end{equation}

where \textbf{P\(_{50}\)} = PO\(_2\) at 50\% saturation (about 26 mmHg
for human HbA at 37 degrees C, pH 7.4) and \textbf{n} is about 2.7 (Hill
coefficient; n = 1 would be non-cooperative; n = 4 would be perfectly
cooperative for a tetramer).

\textbf{Bohr effect:} The curve shifts rightward (P\(_{50}\) increases,
affinity decreases) with: - Increased {[}H\(^+\){]} (lower pH): Protons
bind to histidine residues on Hb, stabilising T-state - Increased
CO\(_2\): Forms carbaminohaemoglobin and generates H\(^+\) - Increased
temperature: Weakens Hb-O\(_2\) bond - Increased 2,3-DPG
(2,3-diphosphoglycerate): Binds between \texorpdfstring{\ensuremath{\beta}}{beta} chains, stabilising T-state

\textbf{Physiological significance:} In metabolically active tissues
(low pH, high CO\(_2\), high temperature), Hb releases more O\(_2\)
exactly where it is needed.

\subsection{Worked Example: 28.Z --- Bohr Effect and 2,3-BPG in
Exercise}\label{worked-example-28.z-bohr-effect-and-23-bpg-in-exercise}

A trained cyclist's vastus lateralis at maximal exercise has tissue
\(P\mathrm{O}_2 = 18\) mmHg, pH = 7.20 (lactic acidosis), and
temperature = 39 °C. At rest, the same muscle has \(P\mathrm{O}_2 = 40\)
mmHg, pH = 7.40, and 37 °C. Estimate the resulting \textbf{shift in
P\(_{50}\)} and the resulting \textbf{fractional change in O\(_2\)
unloading per cycle of blood}.

\textbf{Step 1.} Quantify the rightward shift in P\(_{50}\) from
exercise conditions. Empirical relations (Severinghaus correction):

\[\Delta \log P_{50} = -0.48(\Delta \text{pH}) + 0.024(\Delta T) + 0.061 \log\left(\frac{[\text{2,3-BPG}]}{4.6\;\text{mM}}\right) \label{eq:unit_IX_circulation_respiration_homeostasis_item_2}\]

For pH change 7.40 \texorpdfstring{\ensuremath{\to}}{->} 7.20 (\texorpdfstring{\ensuremath{\Delta}}{Delta}pH = \texorpdfstring{\ensuremath{-}}{-}0.20) and \(T\) change 37 \texorpdfstring{\ensuremath{\to}}{->} 39 °C
(\(\Delta T\) = +2):

\[\Delta \log P_{50} = -0.48(-0.20) + 0.024(2) = 0.096 + 0.048 = 0.144 \label{eq:unit_IX_circulation_respiration_homeostasis_item_3}\]

\[P_{50}^{\text{exercise}} = P_{50}^{\text{rest}} \times 10^{0.144} = 26 \times 1.39 \approx 36\;\text{mmHg} \label{eq:unit_IX_circulation_respiration_homeostasis_item_4}\]

\textbf{Step 2.} Compute saturation at tissue O\(_2\) tension under each
condition using Hill (\(n = 2.7\)):

At rest (P\(_{50}\) = 26 mmHg, \(P\mathrm{O}_2 = 40\)):
\[SO_2 = \frac{(40/26)^{2.7}}{1 + (40/26)^{2.7}} = \frac{3.32}{4.32} \approx 0.77 \label{eq:unit_IX_circulation_respiration_homeostasis_item_5}\]

At exercise (P\(_{50}\) = 36 mmHg, \(P\mathrm{O}_2 = 18\)):
\[SO_2 = \frac{(18/36)^{2.7}}{1 + (18/36)^{2.7}} = \frac{0.154}{1.154} \approx 0.13 \label{eq:unit_IX_circulation_respiration_homeostasis_item_6}\]

\textbf{Step 3.} Compare arterial-to-venous O\(_2\) unloading. Arterial
saturation in both cases \texorpdfstring{\ensuremath{\approx}}{~} 0.97 (lung Bohr/Haldane effect already
accounted for there). Then:

\begin{itemize}
\tightlist
\item
  Rest extraction: \(\Delta SO_2 = 0.97 - 0.77 = 0.20\) (20\% of bound
  O\(_2\) released)
\item
  Exercise extraction: \(\Delta SO_2 = 0.97 - 0.13 = 0.84\) (84\% of
  bound O\(_2\) released)
\end{itemize}

\textbf{Result.} Exercise increases per-pass O\(_2\) extraction more
than 4-fold. Combined with about 5-fold increase in cardiac
output, total O\(_2\) delivery to working muscle rises
about 20-fold --- fully accounting for the increased VO\(_2\)
during maximal aerobic exercise.

\textbf{2,3-BPG adaptation.} Chronic hypoxia (high altitude, anaemia,
chronic lung disease) elevates RBC 2,3-BPG within hours, further
right-shifting the curve to enhance peripheral O\(_2\) delivery. Stored
bank blood (\textgreater2 weeks) loses 2,3-BPG \texorpdfstring{\ensuremath{\to}}{->} left-shifted curve \texorpdfstring{\ensuremath{\to}}{->}
poor O\(_2\) delivery despite normal Hb concentration.

\textbf{Haldane effect:} Deoxygenated Hb binds CO\(_2\) and H\(^+\) more
readily than oxygenated Hb. In tissues, O\(_2\) release promotes
CO\(_2\) and H\(^+\) binding. In lungs, O\(_2\) loading promotes
CO\(_2\) release. This is quantitatively as important as the Bohr effect
for gas exchange.

\textbf{Special haemoglobins:} - \textbf{Foetal Hb (HbF,
\(\alpha_2\gamma_2\)):} P\(_{50}\) about 19 mmHg; higher O\(_2\)
affinity than adult HbA because \texorpdfstring{\ensuremath{\gamma}}{gamma} chains bind 2,3-DPG poorly. Foetus
extracts O\(_2\) from maternal blood across placenta. -
\textbf{Myoglobin:} Monomeric (no cooperativity); P\(_{50}\) about 2
mmHg; O\(_2\) storage in muscle; delivers O\(_2\) during contraction
when blood flow is compressed. - \textbf{HbS (sickle cell):} Glu6Val
\hyperref[gl:mutation]{\textbf{mutation}} in \texorpdfstring{\ensuremath{\beta}}{beta}-globin. HbS polymerises
when deoxygenated, distorting RBCs into sickle shapes.
\hyperref[gl:heterozygous]{\textbf{Heterozygous}} advantage: resistance
to \emph{Plasmodium falciparum} malaria.

\subsubsection{\texorpdfstring{CO\(_2\)
Transport}{CO_2 Transport}}\label{co_2-transport}

CO\(_2\) is transported from tissues to lungs by three mechanisms: -
\textbf{Dissolved CO\(_2\):} 8\% (physically dissolved in plasma;
proportional to PCO\(_2\)) - \textbf{Carbaminohaemoglobin:} 27\%
(CO\(_2\) binds to terminal amino groups of Hb) - \textbf{Bicarbonate
(HCO\(_3^-\)):} 65\% (the dominant mechanism)

\begin{equation}
\mathrm{CO}_2 + \mathrm{H}_2\mathrm{O} \rightleftharpoons
\mathrm{H}_2\mathrm{CO}_3 \rightarrow \mathrm{H}^+ + \mathrm{HCO}_3^-
\label{eq:circulation_respiration_homeostasis_12}
\end{equation}

Carbonic anhydrase catalyses the rapid reversible hydration step.
\textbf{Chloride shift:} As HCO\(_3^-\) is produced inside RBCs, it is
exported via the Band 3 (AE1) Cl\(^-\)/HCO\(_3^-\) exchanger. Cl\(^-\)
enters the RBC in exchange, maintaining electroneutrality. The process
reverses in the lungs.

\begin{quote}
\textbf{Clinical Connection:} Carbon monoxide (CO) poisoning occurs
because CO binds haemoglobin with 200-250x higher affinity than O\(_2\),
forming carboxyhaemoglobin (COHb). Even low CO levels (e.g., 10\% COHb)
shift the O\(_2\)-Hb curve leftward (remaining Hb holds O\(_2\) more
tightly), reducing O\(_2\) delivery to tissues. Treatment: 100\% O\(_2\)
or hyperbaric O\(_2\) to competitively displace CO.
\end{quote}

\textbf{Concept Check:} Exercising muscle is hotter, more acidic, and
higher in CO\(_2\) and 2,3-BPG than resting muscle. State whether each
change shifts the O\(_2\)-Hb dissociation curve left or right, and
explain why a curve that shifts in the \emph{opposite} direction (e.g.,
in stored bank blood depleted of 2,3-BPG) impairs tissue O\(_2\)
delivery even when haemoglobin concentration and arterial saturation are
normal. (Hint: track what each shift does to \(P_{50}\) and to the
O\(_2\) released between arterial and tissue \(P\mathrm{O}_2\).)

\begin{center}\rule{0.5\linewidth}{0.5pt}\end{center}

\subsection{Respiratory System and
Ventilation}\label{respiratory-system-and-ventilation}

\subsubsection{Anatomy}\label{anatomy}

\textbf{Conducting zone} (dead space, about 150 mL): Nasal
passages (warm, humidify, filter) to pharynx to larynx to trachea
(C-shaped cartilage rings) to bronchi (right and left mainstem) to
bronchioles. Bronchioles have smooth muscle (regulated by autonomic NS:
sympathetic \(\beta_2\) = bronchodilation; parasympathetic M\(_3\) =
bronchoconstriction).

\textbf{Respiratory zone:} Respiratory bronchioles to alveolar ducts to
alveolar sacs. about 300 million alveoli in adult human lungs.
Total alveolar surface area: about 70 m\(^2\) (half a tennis
court). Alveolar wall thickness: about 0.5 um. Type I
pneumocytes (gas exchange, 95\% of surface area). Type II pneumocytes
(produce \hyperref[gl:surfactant]{\textbf{surfactant}}, 5\% of surface
area; also stem cells that regenerate Type I cells).

\subsubsection{Lung Mechanics}\label{lung-mechanics}

\textbf{Inspiration (active at rest):} Diaphragm + external intercostals
contract. Thorax volume increases. Intrapleural pressure decreases (from
\(-5\) to \(-8\) cmH\(_2\)O). Alveolar pressure falls about 1
cmH\(_2\)O below atmospheric. Air flows in.

\textbf{Expiration (passive at rest):} Elastic recoil of lung tissue and
chest wall. Alveolar pressure rises about 1 cmH\(_2\)O above
atmospheric. Air flows out.

\subsubsection{Pulmonary Surfactant --- Laplace's Law and Alveolar
Stability}\label{pulmonary-surfactant-laplaces-law-and-alveolar-stability}

For a spherical air-liquid interface, Laplace's law gives the collapse
pressure due to surface tension:

\begin{equation}
P = \frac{2\gamma}{r}
\label{eq:unit_IX_laplace}
\end{equation}

where \texorpdfstring{\ensuremath{\gamma}}{gamma} is the surface tension (N/m) and \(r\) is the alveolar radius.
Using \(T\) as the symbol for surface tension is also common in
textbooks; both conventions appear below.

where \(T\) = surface tension (N/m) and \(r\) = radius. With the surface
tension of water (about 70 mN/m), a 50 \texorpdfstring{\ensuremath{\mu}}{mu}m alveolus would
collapse at \(P \approx 2(0.07)/(2.5 \times 10^{-5}) = 5{,}600\) Pa = 56
cmH\(_2\)O --- far exceeding inspiratory pressure. Without surfactant,
small alveoli (r\(_{\min}\) during expiration) would empty into larger
ones, producing massive atelectasis.

\textbf{Surfactant} is a complex of phospholipids (about 80\%,
dominated by \textbf{DPPC, dipalmitoylphosphatidylcholine}) and four
surfactant-specific proteins, secreted by \textbf{Type II pneumocytes}
as lamellar bodies that unwind into a monolayer at the air-liquid
interface.

{\def\LTcaptype{none} % do not increment counter
\begin{longtable}[]{@{}
  >{\raggedright\arraybackslash}p{(\linewidth - 2\tabcolsep) * \real{0.4667}}
  >{\raggedright\arraybackslash}p{(\linewidth - 2\tabcolsep) * \real{0.5333}}@{}}
\toprule\noalign{}
\begin{minipage}[b]{\linewidth}\raggedright
Protein
\end{minipage} & \begin{minipage}[b]{\linewidth}\raggedright
Function
\end{minipage} \\
\midrule\noalign{}
\endhead
\bottomrule\noalign{}
\endlastfoot
\textbf{SP-A} & Hydrophilic; innate immunity (opsonin for pathogens);
regulates surfactant turnover; mutations rare \\
\textbf{SP-B} & Hydrophobic; essential for lamellar body biogenesis and
surface film formation. \textbf{SP-B deficiency is uniformly lethal in
newborns} (severe RDS unresponsive to therapy) \\
\textbf{SP-C} & Hydrophobic; stabilises surface film during
compression-expansion. Mutations cause familial interstitial lung
disease \\
\textbf{SP-D} & Hydrophilic; innate immunity (similar role to SP-A);
collectin family \\
\end{longtable}
}

\textbf{Surfactant action.} DPPC molecules align at the air-liquid
interface, displacing water molecules and reducing surface tension from
about 70 mN/m to about 5--10 mN/m at full lung volume
--- and to \textbf{near zero} during expiration when alveoli are
smallest and the surfactant film is most compressed. This
\textbf{non-linear} behaviour (lower tension at smaller radius) is
critical: it inverts the destabilising effect of Laplace's law, making
small alveoli \textbf{stable} rather than collapse-prone.

\textbf{Surfactant turnover} is rapid (about 10 h half-life).
Type II cells synthesise, secrete, internalise, and recycle surfactant
continuously. Deep breaths (sighs) extend and refresh the surface film;
absence of sighs (e.g., during anaesthesia, prolonged shallow breathing)
leads to atelectasis.

\begin{quote}
\textbf{Clinical Connection:} Premature infants (\textless34 weeks
gestation) lack mature Type II pneumocytes and adequate surfactant,
leading to \textbf{neonatal respiratory distress syndrome (NRDS /
hyaline membrane disease)} --- the leading cause of death in premature
infants before surfactant therapy. Treatment: intratracheal
administration of exogenous surfactant (beractant from bovine lung;
poractant alfa from porcine lung). \textbf{Antenatal corticosteroids}
(betamethasone) given to mothers in preterm labour 24--48 h before
delivery accelerate fetal Type II cell maturation and reduce NRDS
incidence by about 50\% --- one of the most cost-effective
interventions in modern medicine. \textbf{Adult ARDS} has surfactant
inactivation by inflammatory exudates as a contributing mechanism;
replacement therapy in adults has been less successful than in neonates.
\end{quote}

\subsubsection{Lung Volumes and
Spirometry}\label{lung-volumes-and-spirometry}

{\def\LTcaptype{none} % do not increment counter
\begin{longtable}[]{@{}
  >{\raggedright\arraybackslash}p{(\linewidth - 4\tabcolsep) * \real{0.3947}}
  >{\raggedright\arraybackslash}p{(\linewidth - 4\tabcolsep) * \real{0.2632}}
  >{\raggedright\arraybackslash}p{(\linewidth - 4\tabcolsep) * \real{0.3421}}@{}}
\toprule\noalign{}
\begin{minipage}[b]{\linewidth}\raggedright
Volume/Capacity
\end{minipage} & \begin{minipage}[b]{\linewidth}\raggedright
Definition
\end{minipage} & \begin{minipage}[b]{\linewidth}\raggedright
Typical Value
\end{minipage} \\
\midrule\noalign{}
\endhead
\bottomrule\noalign{}
\endlastfoot
Tidal volume (TV) & Normal breath & 500 mL \\
Inspiratory reserve (IRV) & Maximum additional inspiration & 3,000 mL \\
Expiratory reserve (ERV) & Maximum additional expiration & 1,200 mL \\
Residual volume (RV) & Air remaining after maximal expiration & 1,200
mL \\
Vital capacity (VC) & TV + IRV + ERV & 4,700 mL \\
Total lung capacity (TLC) & VC + RV & 5,900 mL \\
Functional residual capacity (FRC) & ERV + RV & 2,400 mL \\
\end{longtable}
}

\textbf{Spirometry and pulmonary function} \citep{graham2019spirometry}:
- \textbf{FEV\(_1\):} Volume expired in first second of forced
expiration - \textbf{FVC:} Total volume during forced expiration -
\textbf{FEV\(_1\)/FVC ratio:} - Normal: above the age-dependent lower
limit of normal; a fixed at least 70\% threshold is a useful screening
simplification but can over-call obstruction in older adults and
under-call it in younger adults. - \textbf{Obstructive disease} (COPD,
asthma): reduced FEV\(_1\)/FVC (air trapping, increased resistance) -
\textbf{Restrictive disease} (pulmonary fibrosis): preserved or high
FEV\(_1\)/FVC with reduced FVC/TLC (stiff lungs, reduced compliance)

\begin{center}\rule{0.5\linewidth}{0.5pt}\end{center}

\subsection{Gas Exchange}\label{gas-exchange}

\begin{figure}[htbp]
\centering
\includegraphics[width=0.82\linewidth,height=4.2in,keepaspectratio]{../figures/mermaid_inline/inline_mermaid_0155_8d32c9d6a8bf.png}
\caption{Respiratory control and gas exchange. The medullary respiratory centre generates the breathing rhythm. Chemoreceptors (central and peripheral) provide feedback to adjust ventilation rate. Gas exchange at the alveolar-capillary interface follows Fick's law, with equilibration occurring within 0.3 seconds.}
\end{figure}

\textbf{Fick's law of diffusion for gas exchange:}

\begin{equation}
\dot{V}_{gas} = \frac{D \cdot A \cdot \Delta P}{T}
\label{eq:circulation_respiration_homeostasis_13}
\end{equation}

where D = diffusion coefficient (solubility/molecular weight), A =
alveolar surface area (about 70 m\(^2\)), \(\Delta P\) =
partial pressure difference, T = barrier thickness (about 0.5
um).

CO\(_2\) diffuses 20x faster than O\(_2\) (despite similar molecular
weight) because CO\(_2\) is much more soluble in aqueous tissue.

\subsection{Worked Example: Alveolar Oxygen
Diffusion}\label{worked-example-alveolar-oxygen-diffusion}

Calculate the rate of O\(_2\) diffusion across the alveolar-capillary
barrier using Fick's law with these values: - \(D_{O_2}\) (in tissue at
37°C) \(\approx 1.5 \times 10^{-5}\) cm\(^2\)/s - \(A\) = 70 m\(^2\) =
\(70 \times 10^4\) cm\(^2\) - \(\Delta P_{O_2}\) = alveolar \(PO_2\) 100
mmHg \(-\) venous \(PO_2\) 40 mmHg = 60 mmHg \(= 60 \times 133.3\) Pa
\(\approx 8{,}000\) Pa; for gas diffusion use \(\Delta P\) in mmHg
directly with a tissue-specific coefficient - \(T\) = 0.5 \texorpdfstring{\ensuremath{\mu}}{mu}m
\(= 0.5 \times 10^{-4}\) cm

Using Krogh's diffusion constant \(K_{O_2} \approx 1.73 \times 10^{-9}\)
mL O\(_2\)/(cm\(\cdot\)s\(\cdot\)mmHg):

\begin{equation}
\dot{V}_{O_2} = K_{O_2} \cdot A \cdot \frac{\Delta P}{T} = 1.73 \times 10^{-9} \times 70 \times 10^4 \times \frac{60}{0.5 \times 10^{-4}}
\label{eq:circulation_respiration_homeostasis_14}
\end{equation}

\begin{equation}
= 1.73 \times 10^{-9} \times 7 \times 10^5 \times 1.2 \times 10^6 \approx 1.45\;\text{mL O}_2/\text{s} = 87\;\text{mL O}_2/\text{min}
\label{eq:circulation_respiration_homeostasis_15}
\end{equation}

At rest, whole-body O\(_2\) consumption is about 250 mL/min ---
so the diffusing capacity, even under resting conditions, is ample.
During maximal exercise, increased cardiac output, capillary
recruitment, and a wider \(\Delta P\) (tissues extract more O\(_2\),
reducing venous \(PO_2\) to about 20 mmHg) push diffusion to
the limit in elite athletes.

\begin{quote}
\textbf{Clinical Connection:} In \textbf{diffusion-limited} lung disease
(pulmonary fibrosis, interstitial lung disease), the alveolar-capillary
membrane is thickened (increased \(T\)) and the surface area is reduced.
Fick's law predicts reduced O\(_2\) transfer: patients become hypoxaemic
especially during exercise when transit time through capillaries is
shortened and the blood cannot equilibrate with alveolar O\(_2\).
\end{quote}

\subsubsection{Respiratory Control}\label{respiratory-control}

\textbf{Central pattern generator:} Medullary respiratory centres
generate the basic breathing rhythm: - \textbf{Dorsal respiratory group
(DRG):} Primarily inspiratory \hyperref[gl:neuron]{\textbf{neuron}}s;
drives diaphragm via phrenic nerve - \textbf{Ventral respiratory group
(VRG):} Active during forced breathing; contains both inspiratory and
expiratory neurons

\textbf{Pontine centres:} - \textbf{Pneumotaxic centre:} Limits
inspiration duration; increases respiratory rate - \textbf{Apneustic
centre:} Promotes prolonged inspiration (normally inhibited by
pneumotaxic centre)

\textbf{Chemoreceptors:} - \textbf{Central chemoreceptors} (ventral
medullary surface): Respond to H\(^+\) in CSF (which reflects arterial
PCO\(_2\) because CO\(_2\) freely crosses the blood-brain barrier and is
converted to H\(^+\) by carbonic anhydrase). This is the \textbf{primary
driver} of normal ventilation. - \textbf{Peripheral chemoreceptors}
(carotid bodies at carotid bifurcation; aortic bodies on aortic arch):
Respond to decreased PO\(_2\) (\textless60 mmHg primarily), increased
PCO\(_2\), and decreased pH. The carotid bodies are the primary sensors
of arterial PO\(_2\).

\textbf{Concept Check:} In chronic COPD, patients may retain CO\(_2\)
chronically (hypercapnia). Over time, central chemoreceptors reset to
the elevated CO\(_2\) level. What becomes the primary stimulus for
breathing in these patients, and why is high-flow oxygen potentially
dangerous?

\begin{center}\rule{0.5\linewidth}{0.5pt}\end{center}

\subsection{Exercise Physiology --- Integrated Cardiovascular and
Respiratory
Adjustments}\label{exercise-physiology-integrated-cardiovascular-and-respiratory-adjustments}

Exercise is the most demanding physiological challenge for the
cardiovascular and respiratory systems. The match between O\(_2\)
delivery and demand is precise and rapid. The \textbf{Fick principle}
quantifies O\(_2\) consumption:

\begin{equation}
\dot{V}\mathrm{O}_2 = CO \times (Ca\mathrm{O}_2 - Cv\mathrm{O}_2)
\label{eq:circulation_respiration_homeostasis_fick}
\end{equation}

where \(Ca\mathrm{O}_2 - Cv\mathrm{O}_2\) is the arteriovenous O\(_2\)
content difference. At rest, an average adult has
\(\dot{V}\mathrm{O}_2 \approx 250\) mL/min = 5 L/min \texorpdfstring{\ensuremath{\times}}{x} (200 \texorpdfstring{\ensuremath{-}}{-} 150) mL/L
= 5 L/min \texorpdfstring{\ensuremath{\times}}{x} 50 mL/L. Maximal O\(_2\) uptake (V̇O\(_2\)max) in elite
endurance athletes can exceed 80 mL/kg/min (vs about 40
mL/kg/min in untrained adults).

\subsubsection{Cardiovascular
adjustments}\label{cardiovascular-adjustments}

{\def\LTcaptype{none} % do not increment counter
\begin{longtable}[]{@{}
  >{\raggedright\arraybackslash}p{(\linewidth - 6\tabcolsep) * \real{0.2162}}
  >{\raggedright\arraybackslash}p{(\linewidth - 6\tabcolsep) * \real{0.1081}}
  >{\raggedright\arraybackslash}p{(\linewidth - 6\tabcolsep) * \real{0.4324}}
  >{\raggedright\arraybackslash}p{(\linewidth - 6\tabcolsep) * \real{0.2432}}@{}}
\toprule\noalign{}
\begin{minipage}[b]{\linewidth}\raggedright
Variable
\end{minipage} & \begin{minipage}[b]{\linewidth}\raggedright
Rest
\end{minipage} & \begin{minipage}[b]{\linewidth}\raggedright
Maximal exercise
\end{minipage} & \begin{minipage}[b]{\linewidth}\raggedright
Mechanism
\end{minipage} \\
\midrule\noalign{}
\endhead
\bottomrule\noalign{}
\endlastfoot
\textbf{Heart rate} & 70 bpm & 200 bpm (220\texorpdfstring{\ensuremath{-}}{-}age rule) & Sympathetic \texorpdfstring{\ensuremath{\uparrow}}{up} +
vagal withdrawal \\
\textbf{Stroke volume} & 70 mL & 130--160 mL (untrained);
\textgreater200 mL (elite) & Frank-Starling (\texorpdfstring{\ensuremath{\uparrow}}{up}venous return) +
\texorpdfstring{\ensuremath{\uparrow}}{up}contractility (\texorpdfstring{\ensuremath{\beta}}{beta}\(_1\)) \\
\textbf{Cardiac output} & 5 L/min & 26--35 L/min (5--7\texorpdfstring{\ensuremath{\times}}{x}) & HR \texorpdfstring{\ensuremath{\times}}{x} SV \\
\textbf{Systolic BP} & 120 mmHg & 200 mmHg & \texorpdfstring{\ensuremath{\uparrow}}{up}CO \texorpdfstring{\ensuremath{\uparrow}}{up} TPR (running muscle
excepted) \\
\textbf{Diastolic BP} & 80 mmHg & 80 mmHg (unchanged) & Active muscle
vasodilation offsets sympathetic vasoconstriction \\
\textbf{MAP} & 93 mmHg & 120 mmHg & Pulse pressure widens \\
\textbf{Skeletal muscle blood flow} & 1 L/min (20\% of CO) & 22 L/min
(84\% of CO) & Local metabolic vasodilation (adenosine, K\(^+\),
lactate, H\(^+\), CO\(_2\)) overrides sympathetic vasoconstriction \\
\textbf{Splanchnic blood flow} & 1.4 L/min & 0.3 L/min & Sympathetic
vasoconstriction redistributes flow \\
\textbf{Skin blood flow} & 0.4 L/min & 1.5 L/min (heat dissipation) &
Hypothalamic thermoregulation \\
\end{longtable}
}

The \textbf{a-vO\(_2\) difference} widens from 50 mL/L at rest to
\textgreater150 mL/L at maximal exercise --- combined with
about 5-fold \texorpdfstring{\ensuremath{\uparrow}}{up}CO, produces about 20-fold \texorpdfstring{\ensuremath{\uparrow}}{up}V̇O\(_2\).

\subsubsection{Respiratory adjustments}\label{respiratory-adjustments}

\begin{itemize}
\tightlist
\item
  \textbf{Minute ventilation} (\(\dot{V}_E\)): rises from 6 L/min at
  rest to 100--200 L/min at maximal exercise (about 20--30\texorpdfstring{\ensuremath{\times}}{x}
  increase) --- driven primarily by \texorpdfstring{\ensuremath{\uparrow}}{up}tidal volume initially (up to
  about 50\% of vital capacity) and then \texorpdfstring{\ensuremath{\uparrow}}{up}frequency (up to
  about 50/min).
\item
  \textbf{Anaerobic threshold:} about 50--80\% of V̇O\(_2\)max
  in untrained vs trained subjects; lactate production exceeds
  clearance, blood lactate rises, ventilation rises disproportionately
  to drive off CO\(_2\) and buffer H\(^+\) (respiratory compensation for
  metabolic acidosis).
\item
  \textbf{Diffusion capacity:} Pulmonary capillary recruitment and
  dilation increase the area for gas exchange about 3-fold,
  maintaining alveolar-end-capillary equilibration even with shortened
  transit time (0.75 s at rest \texorpdfstring{\ensuremath{\to}}{->} 0.25 s at exercise).
\end{itemize}

\textbf{Training adaptations:} Endurance training increases stroke
volume (cardiac hypertrophy, eccentric), capillary density,
mitochondrial volume, and oxidative enzyme content (citrate synthase,
\texorpdfstring{\ensuremath{\beta}}{beta}-HAD). V̇O\(_2\)max can increase 15--25\% with training. Cardiac output
increases at any work rate; resting heart rate falls (about 50
bpm or lower in trained athletes --- eccentric hypertrophy and \texorpdfstring{\ensuremath{\uparrow}}{up}vagal
tone).

\textbf{Concept Check:} An athlete has \(\dot{V}\mathrm{O}_2 = 4{,}000\)
mL/min, arterial O\(_2\) content \(Ca\mathrm{O}_2 = 200\) mL/L, and
mixed-venous content \(Cv\mathrm{O}_2 = 40\) mL/L at maximal exercise.
Use the Fick principle to solve for cardiac output, and explain why a
trained endurance athlete reaches a higher \(\dot{V}\mathrm{O}_2\)max
than an untrained person of identical maximal cardiac output. (Hint:
rearrange
\(\dot{V}\mathrm{O}_2 = CO \times (Ca\mathrm{O}_2 - Cv\mathrm{O}_2)\)
and consider the a-vO\(_2\) difference.)

\begin{center}\rule{0.5\linewidth}{0.5pt}\end{center}

\subsection{Homeostasis}\label{homeostasis}

\subsubsection{Principles}\label{principles}

\textbf{\hyperref[gl:homeostasis]{Homeostasis}} \citep{cannon1932}:
Maintenance of a relatively stable internal environment (milieu
interieur, Claude Bernard 1865) despite external fluctuations.

\textbf{Negative feedback:} The dominant control mechanism. A change in
a regulated variable is detected by a sensor, compared to a set point by
a control centre, and corrected by an effector that opposes the change.

\textbf{Positive feedback:} Amplifies a change rather than opposing it.
Rare in physiology; examples include: - Oxytocin and uterine
contractions during labour - Platelet activation in haemostasis -
\hyperref[gl:action-potential]{\textbf{Action potential}} upstroke
(Na\(^+\) channel positive feedback)

\subsubsection{Temperature Regulation}\label{temperature-regulation}

\textbf{Ectotherms} (fish, amphibians, reptiles, invertebrates): Body
temperature tracks environmental temperature. Behavioural
thermoregulation (basking, seeking shade).

\textbf{Endotherms} (mammals, birds): Maintain body temperature near a
set point (about 37 degrees C in humans) via metabolic heat
production and physiological regulation.

The \textbf{hypothalamus} serves as the thermostat:

\textbf{Heat dissipation (when too hot):} - Cutaneous vasodilation
(increased blood flow to skin surface) - Sweating (evaporative cooling;
about 580 kcal/L evaporated) - Behavioural responses (seeking
shade, reducing activity)

\textbf{Heat conservation/generation (when too cold):} - Cutaneous
vasoconstriction (reduced heat loss from skin) - Piloerection (limited
effectiveness in humans) - Shivering thermogenesis (skeletal muscle
contraction without useful work; increases metabolic rate 5-fold) -
\textbf{Non-shivering thermogenesis:} Brown adipose tissue (BAT)
expresses UCP1 (uncoupling protein 1), which dissipates the
mitochondrial proton gradient as heat rather than ATP. Significant in
neonates and hibernating mammals; recently confirmed active in adult
humans by PET-CT imaging.

\subsubsection{Countercurrent Heat
Exchange}\label{countercurrent-heat-exchange}

\textbf{Countercurrent heat exchange} is a passive but elegant mechanism
that conserves core heat in cold-stressed appendages. Arteries supplying
the limb run alongside paired veins (the \textbf{vena comitans}
arrangement). Warm arterial blood transfers heat to the cooler venous
blood returning from the cold extremity. By the time arterial blood
reaches the distal limb, it has cooled substantially; conversely,
returning venous blood is rewarmed before reaching the core.

\begin{figure}[htbp]
\centering
\includegraphics[width=0.82\linewidth,height=4.2in,keepaspectratio]{../figures/mermaid_inline/inline_mermaid_0156_6603d0e995af.png}
\caption{Countercurrent heat exchange in a cold limb Heat flows continuously from warmer arterial to cooler venous blood along the entire vessel length. Warm arterial blood is "pre-cooled" before reaching the distal extremity; cold venous blood is "rewarmed" before re-entering the core. The result is dramatic conservation of core body temperature despite cold exposure of the appendage.}
\end{figure}

The same principle operates in \textbf{fish gills} (countercurrent
O\(_2\) extraction --- water flow opposite to blood flow extracts up to
80--90\% of dissolved O\(_2\)), the \textbf{renal medulla}
(countercurrent multiplication for urine concentration), and the
\textbf{placenta} (maternal-fetal exchange).

\subsubsection{Hibernation, Torpor, and
Heterothermy}\label{hibernation-torpor-and-heterothermy}

Some endotherms can dramatically lower their body temperature and
metabolic rate, entering states of \textbf{torpor} (short-term, hours to
days) or \textbf{hibernation} (seasonal, weeks to months):

\begin{itemize}
\tightlist
\item
  \textbf{Daily torpor} in small mammals (hummingbirds, mice) and many
  bats: nightly drops in T\(_b\) to about 10--20 °C, saving
  about 50\% of energy expenditure.
\item
  \textbf{Multi-day torpor} in some marsupials (sugar gliders) and
  primates (lemurs).
\item
  \textbf{Hibernation} in ground squirrels, marmots, bears: T\(_b\)
  falls to near 0 °C in some species; metabolic rate drops to
  \textless2\% of basal; heart rate falls from about 300 to
  \textless10 bpm. Periodic arousals (every 1--3 weeks) raise T\(_b\)
  briefly to enable function then return to torpor.
\end{itemize}

\textbf{Mechanisms:} - \textbf{Hypothalamic set point reduction:} During
torpor entry, the thermal set point is actively lowered (not a failure
of regulation). - \textbf{Pre-fattening:} Animals deposit large lipid
stores (white and brown) in autumn. - \textbf{Metabolic suppression:}
Mitochondrial respiration and ATPase activity are suppressed beyond what
temperature alone predicts (Q\(_{10}\) exceeds typical values). -
\textbf{Adaptive cold tolerance:} Membrane composition shifts to
maintain fluidity; SERCA pumps are modified to maintain Ca\(^{2+}\)
handling. - \textbf{Arousal via BAT:} BAT thermogenesis (UCP1,
sympathetic input) plus shivering raises T\(_b\) to normothermic levels
for arousal.

\begin{quote}
\textbf{Clinical Connection:} Therapeutic hypothermia (32--36 °C for
24--72 h) is now standard care after cardiac arrest with return of
spontaneous circulation, and after neonatal hypoxic-ischaemic
encephalopathy. The mechanism --- slowed metabolic demand, reduced
excitotoxicity, suppressed apoptotic pathways --- parallels natural
hibernation. Pharmacological induction of torpor-like states
(``synthetic torpor'') is an active area of research for trauma, stroke,
and even spaceflight.
\end{quote}

\begin{quote}
\textbf{Clinical Connection:} Fever is not a failure of thermoregulation
but a regulated elevation of the hypothalamic set point. Pyrogens (IL-1,
IL-6, TNF-alpha from macrophages) stimulate hypothalamic COX-2 to
produce PGE\(_2\), which raises the set point. The body then uses normal
heat-generating mechanisms (vasoconstriction, shivering) to reach the
new, higher set point. NSAIDs (aspirin, ibuprofen) reduce fever by
inhibiting COX-2.
\end{quote}

\subsubsection{Fluid and Electrolyte
Homeostasis}\label{fluid-and-electrolyte-homeostasis}

\textbf{Renal regulation:} The kidney filters about 180 L of
plasma per day but excretes about 1.5 L urine.

The nephron performs: 1. \textbf{Glomerular filtration:}
about 125 mL/min GFR. Pressure-driven ultrafiltration. Proteins
and blood cells excluded. 2. \textbf{Tubular reabsorption:} PCT
reabsorbs about 65\% of filtrate; Loop of Henle creates
corticomedullary concentration gradient by countercurrent
multiplication; DCT + collecting duct adjust water (ADH) and Na\(^+\)
(aldosterone). 3. \textbf{Tubular secretion:} Organic acids, drugs,
K\(^+\), H\(^+\) secreted.

\textbf{Key \hyperref[gl:hormone]{\textbf{hormone}}s in fluid balance:}
- \textbf{ADH (vasopressin):} Released from posterior pituitary when
plasma osmolality rises. Inserts aquaporin-2 channels into collecting
duct, increasing water reabsorption. Diabetes insipidus: deficiency of
ADH (central) or renal insensitivity (nephrogenic) causes massive urine
output (up to 20 L/day). - \textbf{Aldosterone:} Released from adrenal
cortex zona glomerulosa (stimulated by angiotensin II and high K\(^+\)).
Increases ENaC Na\(^+\) channels in collecting duct, increasing Na\(^+\)
and water reabsorption. - \textbf{ANP (atrial natriuretic peptide):}
Released from atrial cardiomyocytes when stretched. Promotes Na\(^+\)
and water excretion. Antagonises RAAS.

\textbf{RAAS (Renin-Angiotensin-Aldosterone System):} 1. Low renal
perfusion pressure detected by juxtaglomerular cells 2. Renin secreted,
cleaves angiotensinogen (liver) to angiotensin I 3. ACE (lung
endothelium) converts angiotensin I to angiotensin II 4. Angiotensin II:
vasoconstriction + aldosterone release + ADH release + thirst
stimulation 5. Net effect: blood pressure and volume restoration

\begin{center}\rule{0.5\linewidth}{0.5pt}\end{center}

\subsection{Worked Example}\label{worked-example-1}

\textbf{Problem:} A patient has a resting heart rate (\(HR\)) of
\(75\text{ beats/min}\) and a stroke volume (\(SV\)) of
\(70\text{ mL/beat}\). During exercise, the heart rate increases to
\(150\text{ beats/min}\) and stroke volume increases to
\(110\text{ mL/beat}\). Calculate the resting cardiac output
(\(CO_{rest}\)) and the exercise cardiac output (\(CO_{exercise}\)) in
L/min. What is the fold-increase in cardiac output?

\textbf{Solution:}

\textbf{Step 1. Calculate resting cardiac output.}
\[CO = HR \times SV \label{eq:unit_IX_circulation_respiration_homeostasis_item_7}\]

\[CO_{rest} = 75\text{ beats/min} \times 70\text{ mL/beat} = 5,250\text{ mL/min} = 5.25\text{ L/min} \label{eq:unit_IX_circulation_respiration_homeostasis_item_8}\]

\textbf{Step 2. Calculate exercise cardiac output.}
\[CO_{exercise} = 150\text{ beats/min} \times 110\text{ mL/beat} = 16,500\text{ mL/min} = 16.5\text{ L/min} \label{eq:unit_IX_circulation_respiration_homeostasis_item_9}\]

\textbf{Step 3. Calculate the fold-increase.}
\[\text{Fold-increase} = \frac{CO_{exercise}}{CO_{rest}} \label{eq:unit_IX_circulation_respiration_homeostasis_item_10}\]

\[\text{Fold-increase} = \frac{16.5\text{ L/min}}{5.25\text{ L/min}} \approx 3.14 \label{eq:unit_IX_circulation_respiration_homeostasis_item_11}\]

\textbf{Answer:} The patient's cardiac output increases from
\textbf{\(5.25\text{ L/min}\)} at rest to \textbf{\(16.5\text{ L/min}\)}
during exercise, representing a \textbf{3.14-fold} increase.

\begin{center}\rule{0.5\linewidth}{0.5pt}\end{center}

\subsection{Computational Bridge}\label{computational-bridge-27}

The Hb--O\(_2\) curve used in problems is materialised as discrete
samples:

\begin{Shaded}
\begin{Highlighting}[]
\ImportTok{from}\NormalTok{ biology.physiology }\ImportTok{import}\NormalTok{ oxygen\_dissociation\_curve}

\NormalTok{curve }\OperatorTok{=}\NormalTok{ oxygen\_dissociation\_curve(p50\_mmHg}\OperatorTok{=}\FloatTok{26.0}\NormalTok{, n\_points}\OperatorTok{=}\DecValTok{8}\NormalTok{)}
\BuiltInTok{print}\NormalTok{(}\BuiltInTok{round}\NormalTok{(curve[}\DecValTok{4}\NormalTok{].saturation, }\DecValTok{3}\NormalTok{))}
\end{Highlighting}
\end{Shaded}

\begin{quote}
\textbf{Clinical / systems note:} Pulse oximeters estimate SpO\(_2\)
along the same saturation curve; right shifts from fever or acidosis
explain why patients can be ``happy hypoxaemic'' until they abruptly
decompensate.
\end{quote}

\begin{center}\rule{0.5\linewidth}{0.5pt}\end{center}

\subsection{Current Evidence and Frontier
Biology}\label{current-evidence-and-frontier-biology-30}

For \textbf{Circulation, Respiration, and Homeostasis}, frontier biology
belongs inside the evidence logic of the chapter. Physiology now blends
mechanism with allostasis, immune-endocrine-neural coupling, wearable
data, and individualized risk without reducing bodies to simple
machines. The core reading question is this: homeostasis claims should
connect flow, diffusion, control loops, reserve capacity, and
measurement limits.

\begin{itemize}
\tightlist
\item
  \textbf{What to verify:} identify the observation, model, assay, or
  dataset that would make the claim stronger or weaker.
\item
  \textbf{What to qualify:} state the scale, organism, cell type,
  environmental condition, or population where the claim is expected to
  hold.
\item
  \textbf{What to compare:} test at least one alternative explanation,
  baseline, or null model before treating the pattern as causal.
\item
  \textbf{What to cite:} distinguish primary evidence, review synthesis,
  public dataset, and institutional guidance; for recent or numeric
  claims, prefer the source closest to the measurement and state what
  has changed since it was published.
\end{itemize}

Interpret physiological data by separating baseline variation,
perturbation response, compensation, and the threshold where
compensation becomes pathology.

\textbf{Source practice:} For physiology claims, cite the measurement
context and distinguish baseline variation, compensation,
pathophysiology, and treatment evidence.

\subsection{Summary}\label{summary-30}

\begin{itemize}
\tightlist
\item
  \textbf{Cardiovascular evolution:} Open to closed; 2-chamber (fish) to
  4-chamber (mammals/birds). Complete separation enables high-pressure
  systemic and low-pressure pulmonary circuits.
\item
  \textbf{Cardiac cycle:} SA node (pacemaker) to AV node (delay) to
  bundle of His to Purkinje fibres. ECG: P (atrial), QRS (ventricular
  depolarisation), T (repolarisation). Starling's law matches SV to
  preload.
\item
  \textbf{Cardiac output:} CO = HR \(\times\) SV; about 5 L/min
  at rest. Regulated by sympathetic (\(\beta_1\)), parasympathetic
  (M\(_2\)), and Starling mechanism.
\item
  \textbf{Blood flow:} Poiseuille's law (\(Q \propto r^4\)); arteriolar
  tone is the key regulator. Local autoregulation (myogenic, metabolic,
  endothelial NO) and neural/hormonal control.
\item
  \textbf{O\(_2\) transport:} Hb cooperatively binds O\(_2\) (Hill
  equation; n about 2.7; P\(_{50}\) about 26 mmHg). Bohr effect unloads
  O\(_2\) in metabolically active tissues. Haldane effect facilitates
  CO\(_2\) transport.
\item
  \textbf{CO\(_2\) transport:} 65\% as bicarbonate (carbonic anhydrase),
  27\% carbaminohaemoglobin, 8\% dissolved.
\item
  \textbf{Ventilation:} Diaphragm contraction creates negative
  intrapulmonary pressure. Surfactant prevents alveolar collapse.
  FEV\(_1\)/FVC ratio distinguishes obstructive from restrictive
  disease.
\item
  \textbf{Gas exchange:} Fick's law governs diffusion across the
  alveolar-capillary membrane (70 m\(^2\) surface area, 0.5 um
  thickness).
\item
  \textbf{Respiratory control:} Medullary DRG/VRG generate rhythm.
  Central chemoreceptors (H\(^+\)/CO\(_2\) in CSF) are the primary
  driver; peripheral chemoreceptors (carotid/aortic bodies) detect
  PO\(_2\).
\item
  \textbf{Homeostasis:} Negative feedback dominates. Thermoregulation
  via hypothalamus (vasomotor, sweating, shivering, BAT). Fluid balance
  via RAAS, ADH, aldosterone, ANP.
\item
  \textbf{Connections:} See \cref{sec:unit_IX_nervous_system} and
  \cref{sec:unit_IX_action_potential_synapses} for neural control of
  breathing and autonomic output,
  \cref{sec:unit_III_metabolic_integration} for metabolic heat load, and
  \cref{sec:unit_I_water_and_life} for blood osmolarity.
\end{itemize}

\begin{center}\rule{0.5\linewidth}{0.5pt}\end{center}

\subsection{Key Terms}\label{key-terms-30}

{\def\LTcaptype{none} % do not increment counter
\begin{longtable}[]{@{}
  >{\raggedright\arraybackslash}p{(\linewidth - 2\tabcolsep) * \real{0.2857}}
  >{\raggedright\arraybackslash}p{(\linewidth - 2\tabcolsep) * \real{0.7143}}@{}}
\toprule\noalign{}
\begin{minipage}[b]{\linewidth}\raggedright
Term
\end{minipage} & \begin{minipage}[b]{\linewidth}\raggedright
Definition
\end{minipage} \\
\midrule\noalign{}
\endhead
\bottomrule\noalign{}
\endlastfoot
\textbf{\hyperref[gl:cardiac-output]{Cardiac output (CO)}} & Volume of
blood pumped per minute; CO = HR \(\times\) SV \\
\textbf{Stroke volume (SV)} & Volume ejected per heartbeat;
about 70 mL at rest \\
\textbf{\hyperref[gl:frank-starling-law]{Frank-Starling law}
(Equationabout \eqref{eq:circulation_respiration_homeostasis_1})}
& Increased preload (EDV) causes increased SV \\
\textbf{Poiseuille's law} & \(Q = \pi r^4 \Delta P / (8\eta L)\); flow
depends on radius to the 4th power \\
\textbf{Mean arterial pressure} & MAP = DP + 1/3(SP - DP);
about 93 mmHg normally \\
\textbf{Haemoglobin} & \(\alpha_2\beta_2\) tetramer; cooperative O\(_2\)
binding; 4 haeme groups \\
\textbf{Bohr effect} & Rightward shift of O\(_2\)-Hb curve with
decreased pH / increased CO\(_2\) \\
\textbf{Haldane effect} & Deoxygenated Hb binds CO\(_2\) and H\(^+\)
more readily than oxygenated Hb \\
\textbf{P\(_{50}\)} & PO\(_2\) at 50\% Hb saturation; about 26
mmHg for adult HbA \\
\textbf{Surfactant} & DPPC from Type II pneumocytes; reduces alveolar
surface tension \\
\textbf{FEV\(_1\)/FVC} & Spirometric ratio interpreted against the
age-dependent lower limit of normal; obstruction lowers the ratio, while
restriction usually preserves or raises it with reduced absolute
volumes \\
\textbf{Chemoreceptors} & Central (medullary, respond to
H\(^+\)/CO\(_2\)) and peripheral (carotid/aortic, respond to
PO\(_2\)) \\
\textbf{Homeostasis} & Maintenance of stable internal environment via
negative feedback \\
\textbf{RAAS} & Renin-Angiotensin-Aldosterone System; restores blood
pressure and volume \\
\textbf{UCP1} & Uncoupling protein 1 in brown adipose tissue; generates
heat \\
\textbf{ADH (vasopressin)} & Posterior pituitary hormone; inserts
aquaporin-2 in collecting duct \\
\end{longtable}
}

\begin{center}\rule{0.5\linewidth}{0.5pt}\end{center}

\subsection{Review Questions}\label{review-questions-30}

\begin{enumerate}
\def\labelenumi{\arabic{enumi}.}
\item
  A patient's ECG shows a prolonged PR interval (0.28 s) but normal QRS
  complexes. Name the type of heart block and explain which structure in
  the conduction system is likely affected.
\item
  Using the Hill equation with P\(_{50}\) = 26 mmHg and n = 2.7,
  calculate SO\(_2\) at PO\(_2\) = 100 mmHg (arterial) and PO\(_2\) = 40
  mmHg (venous). What is the O\(_2\) extraction ratio?
\item
  Explain why CO\(_2\) diffuses 20 times faster than O\(_2\) across the
  alveolar-capillary membrane despite having a similar molecular weight.
  What property accounts for this difference?
\item
  An arteriole supplying skeletal muscle constricts from radius 50 um to
  40 um. Using Poiseuille's law, calculate the fold-change in blood flow
  (assume most other variables constant). Why does this make arteriolar
  tone the dominant regulator of organ blood flow?
\item
  A mountain climber at 5,500 m altitude breathes air with PO\(_2\) = 75
  mmHg. Describe the acute and chronic compensatory responses involving
  the respiratory system, cardiovascular system, and blood.
\item
  Compare central and peripheral chemoreceptors in terms of location,
  stimulus, and response latency. Why is the central chemoreceptor
  considered the primary driver of normal ventilation?
\item
  Explain why administering high-flow oxygen to a patient with chronic
  COPD and CO\(_2\) retention could suppress their ventilatory drive.
  What is the safer approach?
\item
  A patient with diabetes insipidus produces 15 L of dilute urine per
  day. Explain the molecular defect (central vs nephrogenic), and
  describe how ADH normally regulates water reabsorption at the
  molecular level.
\item
  During intense exercise, skeletal muscle produces large amounts of
  CO\(_2\), H\(^+\), and heat. Explain how each of these factors
  enhances O\(_2\) delivery to the exercising muscle via the Bohr
  effect.
\item
  Compare shivering thermogenesis and non-shivering thermogenesis
  (BAT/UCP1) in terms of mechanism, efficiency, tissue involved, and
  developmental significance.
\item
  A 45-year-old patient with nephrotic syndrome (severe proteinuria,
  serum albumin 1.8 g/dL) presents with bilateral pitting oedema to the
  knees. Using the Starling equation, explain: (a) why plasma oncotic
  pressure is reduced; (b) how this shifts the net filtration pressure
  at both arteriolar and venular capillary ends; (c) why the oedema is
  pitting rather than non-pitting. How does treatment with albumin
  infusion immediately reduce oedema?
\item
  Using Fick's law, predict what happens to alveolar O\(_2\) diffusion
  in: (a) a patient with idiopathic pulmonary fibrosis where membrane
  thickness doubles and surface area halves; (b) a healthy person
  exercising at altitude (4,000 m, alveolar \(PO_2\) = 60 mmHg, tissue
  \(PO_2\) = 30 mmHg). Which is more likely to cause frank hypoxaemia at
  rest vs exercise?
\end{enumerate}

\begin{center}\rule{0.5\linewidth}{0.5pt}\end{center}

\subsection{Further Reading and Source
Notes}\label{further-reading-and-source-notes-25}

\begin{itemize}
\tightlist
\item
  Starling (1918). The Linacre Lecture on the Law of the Heart.
  \emph{Longmans, Green and Co.}.
\item
  Cannon (1932). \emph{The Wisdom of the Body}. W. W. Norton.
\end{itemize}

\begin{center}\rule{0.5\linewidth}{0.5pt}\end{center}

\subsubsection{Companion Source
Module}\label{companion-source-module-30}

\textbf{Circulation, Respiration, and Homeostasis} should leave a
reproducible trail from a biological claim to the code, figure, diagram,
or paper-based activity that can test it. Use the surfaces below to
inspect the chapter's assumptions, rerun the relevant model, or compare
the manuscript explanation with companion labs and figures.

{\def\LTcaptype{none} % do not increment counter
\begin{longtable}[]{@{}
  >{\raggedright\arraybackslash}p{(\linewidth - 2\tabcolsep) * \real{0.5000}}
  >{\raggedright\arraybackslash}p{(\linewidth - 2\tabcolsep) * \real{0.5000}}@{}}
\toprule\noalign{}
\begin{minipage}[b]{\linewidth}\raggedright
Surface
\end{minipage} & \begin{minipage}[b]{\linewidth}\raggedright
Use it for
\end{minipage} \\
\midrule\noalign{}
\endhead
\bottomrule\noalign{}
\endlastfoot
\breaktt{src/biology/physiology/physiology.py}
(\texttt{poiseuille\_flow}, \texttt{oxygen\_saturation},
\breaktt{oxygen\_dissociation\_curve}, \breaktt{homeostasis\_response}) &
Reproduce flow, gas transport, and regulatory response claims. \\
\breaktt{src/visualization/plots.py}
(\breaktt{plot\_oxygen\_dissociation}) & Inspect shifts in oxygen loading
and unloading. \\
\end{longtable}
}

\textbf{Reproducibility check:} state vessel radius, pressure gradient,
haemoglobin state, tissue demand, and feedback variable before
interpreting homeostasis. \textbf{Cross-reference:} connect with
\cref{sec:unit_IX_endocrine_and_immune} and
\cref{sec:unit_III_bioenergetics_and_respiration}.

\newpage

\section{Nervous System and Neural
Signalling}\label{nervous-system-and-neural-signalling}

\label{sec:unit_IX_nervous_system}

\begin{quote}
\textbf{Ch 29} \texorpdfstring{\ensuremath{\cdot}}{.} Level 3/3 \texorpdfstring{\ensuremath{\cdot}}{.} 55 min read \texorpdfstring{\ensuremath{\cdot}}{.} 75 min lecture \texorpdfstring{\ensuremath{\cdot}}{.}
Prerequisites: \cref{sec:unit_IX_circulation_respiration_homeostasis}
\end{quote}

\subsection{Learning Objectives}\label{learning-objectives-31}

\begin{enumerate}
\def\labelenumi{\arabic{enumi}.}
\tightlist
\item
  Describe the organisation of the vertebrate nervous system (CNS, PNS)
  and the roles of glial cells.
\item
  Explain \hyperref[gl:neuron]{\textbf{neuron}} structure and classify
  neuron types by function and morphology.
\item
  \textbf{Analyse} how the \hyperref[gl:axon-hillock]{\textbf{axon
  hillock}} integrates competing excitatory and inhibitory synaptic
  inputs to determine whether the neuron fires.
\item
  Describe graded potentials, synaptic summation, and how the
  \hyperref[gl:axon-hillock]{\textbf{axon hillock}} integrates inputs.
\item
  Compare the sympathetic and parasympathetic divisions of the autonomic
  nervous system.
\item
  Describe the major brain regions and their functions.
\item
  Describe sensory systems and neural plasticity.
\item
  \textbf{Calculate} the resting membrane potential from ion
  concentrations and relative permeabilities using the
  Goldman-Hodgkin-Katz equation.
\item
  \textbf{Predict} the direction and approximate magnitude of the
  \(V_m\) shift produced by a stated change in K\(^+\) permeability.
\end{enumerate}

\subsubsection{Study Blueprint}\label{study-blueprint-31}

\begin{itemize}
\tightlist
\item
  \textbf{Big idea:} Nervous systems compute with cells whose structure
  shapes information flow and behaviour.
\item
  \textbf{Core concepts:} neurons, glia, circuits, sensory processing.
\item
  \textbf{Framework alignment:} Vision \& Change: Structure and
  function, Systems; AP Biology: Systems Interactions, Energetics;
  NGSS-style topics: Structure and Function.
\item
  \textbf{Model or quantitative lens:} Cable-length, conduction, and
  simple circuit calculations.
\item
  \textbf{Data skill:} Interpret neural data from anatomy, timing, or
  lesion evidence.
\item
  \textbf{Practice cadence:} Visual Representations, Statistical Tests
  and Data Analysis, Argumentation.
\item
  \textbf{Common misconception to repair:} The brain is not only
  neurons; glia and circuit context are essential to function.
\item
  \textbf{Primary lab:} \cref{sec:lab_unit_IX_nervous_system}.
\item
  \textbf{Question bank:} \cref{sec:q_unit_IX_nervous_system}.
\item
  \textbf{Transfer task:} Transfer circuit reasoning to reflexes,
  sensory systems, learning, and disease.
\item
  \textbf{Bridge to computation:}
  \breaktt{biology.neuroscience.neuroscience.cable\_voltage\_attenuation}.
\end{itemize}

\begin{center}\rule{0.5\linewidth}{0.5pt}\end{center}

\begin{quote}
\textbf{Opening Vignette --- The Reflex That Revealed the Synapse}

In 1897, Charles Scott Sherrington coined the word ``synapse'' --- from
the Greek for ``clasp'' --- to describe the morphological gap between
neurons. Over four decades, working with decerebrate cats and dogs whose
spinal cord reflexes could be studied in isolation, Sherrington showed
that reflexes were not simple mechanical responses but involved
integration: competing signals converged on motor neurons, inhibiting
some pathways while exciting others, with the outcome depending on the
algebraic summation of hundreds of inputs. His 1906 masterwork,
\emph{The Integrative Action \citep{sherrington1906} of the Nervous
System}, established the principles of convergence, divergence,
summation, inhibition, and the motor neuron as the ``final common
pathway.'' The electrical basis of these phenomena ---
\hyperref[gl:action-potential]{\textbf{action potential}}s, EPSPs, IPSPs
--- was not known until Hodgkin, Huxley, and Eccles worked it out in the
1950s. Sherrington and Adrian shared the 1932 Nobel Prize; Eccles,
Hodgkin, and Huxley shared the 1963 Nobel Prize. The synapse Sherrington
named launched two Nobel Prizes' worth of discoveries.
\end{quote}

\subsection{Nervous System
Organisation}\label{nervous-system-organisation}

The human nervous system comprises approximately \textbf{86 billion
neurons} and about 85 billion non-neuronal cells (Azevedo et
al., 2009). It is organised hierarchically into two major divisions.

\begin{figure}[htbp]
\centering
\includegraphics[width=0.82\linewidth,height=4.2in,keepaspectratio]{../figures/mermaid_inline/inline_mermaid_0157_4539c84bcc27.png}
\caption{Organisation of the vertebrate nervous system. The CNS (brain and spinal cord) integrates information. The PNS (somatic and autonomic divisions) connects the CNS to the body. The autonomic NS has sympathetic, parasympathetic, and enteric subdivisions.}
\end{figure}

\begin{itemize}
\tightlist
\item
  Blood-brain barrier (BBB): Tight junctions between endothelial cells +
  astrocyte endfeet
\item
  Three meningeal layers: dura mater (tough outer), arachnoid mater
  (middle; CSF-filled subarachnoid space), pia mater (thin; adheres to
  brain surface)
\item
  Bone: Cranium and vertebral column
\end{itemize}

\textbf{Peripheral nervous system (PNS):} Most nervous tissue outside
the CNS. Includes 12 cranial nerves and 31 spinal nerve pairs.

\begin{center}\rule{0.5\linewidth}{0.5pt}\end{center}

\subsection{Neuron Types and
Structure}\label{neuron-types-and-structure}

\subsubsection{Neuron Classification}\label{neuron-classification}

\textbf{By function:} - \textbf{Sensory (afferent) neurons:} Carry
information from sensory receptors toward CNS. Most are pseudounipolar
(cell body off to one side of a single process that bifurcates). -
\textbf{Motor (efferent) neurons:} Carry commands from CNS to effectors
(muscles, glands). Multipolar with long axons. - \textbf{Interneurons:}
Integration and processing within the CNS. Most numerous
(about 99\% of neurons). Diverse morphologies.

\textbf{By morphology:} - \textbf{Multipolar:} Multiple dendrites + one
axon (most CNS neurons, most motor neurons) - \textbf{Bipolar:} One
dendrite + one axon (retinal bipolar cells, olfactory receptor neurons)
- \textbf{Pseudounipolar:} Single process that splits into peripheral
and central branches (most sensory neurons in DRG) - \textbf{Anaxonic:}
No distinguishable axon (amacrine cells of retina; some interneurons)

\subsubsection{Neuron Structure}\label{neuron-structure}

A typical neuron comprises:

\begin{itemize}
\tightlist
\item
  \textbf{Soma (cell body):} Contains nucleus, rough ER (Nissl
  substance), Golgi apparatus. Site of most
  \hyperref[gl:protein]{\textbf{protein}} synthesis and metabolic
  activity. Diameter: 5-100 um.
\item
  \textbf{Dendrites:} Receive synaptic inputs. Dendritic spines (0.5-2
  um protrusions) increase surface area and compartmentalise Ca\(^{2+}\)
  signals. A single cortical pyramidal neuron may have 10,000-30,000
  dendritic spines.
\item
  \textbf{Axon hillock:} Junction between soma and axon. Highest density
  of voltage-gated Na\(^+\) channels (about 800-1200/um\(^2\)).
  Lowest threshold for action potential initiation. The
  ``decision-making'' node.
\item
  \textbf{Axon:} Signal conduction pathway. Myelinated or unmyelinated.
  Can extend \textgreater1 m (motor neurons to foot muscles). Axoplasmic
  transport: fast anterograde (kinesin, 200-400 mm/day for vesicles),
  slow anterograde (0.5-5 mm/day for cytoskeletal proteins), fast
  retrograde (dynein, 200-300 mm/day for recycled materials and trophic
  signals).
\item
  \textbf{Axon terminal (bouton):} Contains synaptic vesicles (40-50 nm
  diameter, each containing about 5,000-10,000 neurotransmitter
  molecules). Active zone: specialised region of presynaptic membrane
  where vesicles dock and fuse.
\end{itemize}

\begin{center}\rule{0.5\linewidth}{0.5pt}\end{center}

\subsection{Glial Cells}\label{glial-cells}

{\def\LTcaptype{none} % do not increment counter
\begin{longtable}[]{@{}
  >{\raggedright\arraybackslash}p{(\linewidth - 4\tabcolsep) * \real{0.3333}}
  >{\raggedright\arraybackslash}p{(\linewidth - 4\tabcolsep) * \real{0.2667}}
  >{\raggedright\arraybackslash}p{(\linewidth - 4\tabcolsep) * \real{0.4000}}@{}}
\toprule\noalign{}
\begin{minipage}[b]{\linewidth}\raggedright
Glial type
\end{minipage} & \begin{minipage}[b]{\linewidth}\raggedright
Location
\end{minipage} & \begin{minipage}[b]{\linewidth}\raggedright
Key functions
\end{minipage} \\
\midrule\noalign{}
\endhead
\bottomrule\noalign{}
\endlastfoot
\textbf{Astrocytes} & CNS & Metabolic support (lactate shuttle to
neurons); BBB construction; K\(^+\) spatial buffering;
glutamate-glutamine recycling; tripartite synapse (modulate synaptic
transmission); scar formation after injury \\
\textbf{Oligodendrocytes} & CNS & Myelination (one cell can myelinate
segments on up to 50 axons simultaneously) \\
\textbf{Schwann cells} & PNS & Myelination (one cell per axon segment,
about 1 mm internode); nerve repair via Bands of Bungner \\
\textbf{Microglia} & CNS & Resident macrophages; synaptic pruning
(complement C1q/C3-tagged synapses phagocytosed); neuroinflammation;
activated in Alzheimer's, TBI, MS \\
\textbf{Ependymal cells} & CNS & Line ventricles; ciliated surface moves
CSF; produce CSF (together with choroid plexus; 450-500 mL/day; total
CSF volume about 150 mL, replaced about 3
times/day) \\
\textbf{Radial glia} & Developing CNS & Neural stem cells; scaffold for
neuronal migration (inside-out cortical layering); give rise to both
neurons and glia \\
\textbf{Satellite cells} & PNS & Surround cell bodies in ganglia;
analogous to astrocytes \\
\end{longtable}
}

\begin{quote}
\textbf{Clinical Connection:} Microglia-mediated synaptic pruning is
essential during development (eliminating about 50\% of
synapses during adolescence). Dysregulated pruning has been linked to
schizophrenia: the complement component C4A
\hyperref[gl:gene]{\textbf{gene}} is the strongest genetic risk factor
for schizophrenia (Sekar et al., 2016, \emph{Nature}). Individuals with
high C4A expression show excessive synaptic pruning in prefrontal cortex
during adolescence, correlating with symptom onset.
\end{quote}

\subsubsection{Astrocytes --- More Than
Glue}\label{astrocytes-more-than-glue}

Astrocytes were once considered passive support cells but are now
recognised as active partners in synaptic transmission, brain
metabolism, and homeostasis. A single cortical astrocyte contacts
about 100,000--2,000,000 synapses through its fine processes
(``perisynaptic astrocyte processes'' or PAPs) --- forming the third
element of the \textbf{tripartite synapse} (presynaptic terminal +
postsynaptic spine + astrocyte process).

\textbf{Key astrocyte functions:}

\begin{itemize}
\tightlist
\item
  \textbf{K\(^+\) spatial buffering:} Following synaptic activity,
  extracellular K\(^+\) rises locally. Astrocytes express high densities
  of inwardly-rectifying Kir4.1 channels and aquaporin-4 (AQP4) at
  endfeet on capillaries. K\(^+\) enters astrocytes locally and exits at
  sites of low K\(^+\) via the panastrocytic syncytium (gap junctions
  formed by connexin 43/30) --- preventing extracellular K\(^+\)
  accumulation that would depolarise neurons and disrupt firing.
\item
  \textbf{Glutamate clearance:} EAAT1 and EAAT2 transporters (the
  ``glutamate-aspartate transporters'') on astrocytes clear
  about 90\% of synaptically released glutamate. EAAT2/GLT-1
  dysfunction is implicated in ALS (mutations) and excitotoxicity in
  stroke.
\item
  \textbf{Glutamate-glutamine cycle:} Astrocytes convert glutamate to
  glutamine via glutamine synthetase, then export glutamine to neurons,
  where neurons reconvert it back to glutamate via glutaminase. This
  metabolically traps glutamate within neurons (glutamate is itself a
  metabolic intermediate that cannot be allowed to accumulate in
  extracellular space).
\item
  \textbf{Astrocyte-neuron lactate shuttle:} Astrocytes preferentially
  perform glycolysis (express PFK1 and LDH-A) and export lactate. Active
  neurons take up lactate via MCT2 and use it as a major energy
  substrate during sustained activity. This ``ANLS hypothesis''
  (Pellerin and Magistretti, 1994) explains the glucose-lactate
  metabolic coupling underlying fMRI BOLD signals.
\item
  \textbf{Gliotransmission:} Activated astrocytes release glutamate,
  ATP, D-serine (an NMDA receptor co-agonist), and other transmitters
  via Ca\(^{2+}\)-dependent vesicle exocytosis or channel-mediated
  release --- modulating synaptic transmission and plasticity.
\item
  \textbf{Reactive astrogliosis:} After CNS injury, astrocytes
  hypertrophy, upregulate GFAP, and form glial scars. This both contains
  damage and inhibits axonal regeneration (chondroitin sulphate
  proteoglycans). Modulating reactive astrogliosis is a target for
  spinal cord injury therapeutics.
\end{itemize}

\subsubsection{Oligodendrocytes and
Myelin}\label{oligodendrocytes-and-myelin}

Oligodendrocytes wrap layers of plasma membrane around axons in the CNS.
A single oligodendrocyte myelinates segments on up to \textbf{50
different axons} simultaneously (contrast with Schwann cells in PNS: one
cell, one axon). The myelin sheath is up to \textbf{100 layers} thick,
formed by repeated wrapping with extrusion of cytoplasm. The major
proteins are: - \textbf{Myelin basic protein (MBP):} Compaction of
cytoplasmic faces. Major autoantigen in multiple sclerosis. -
\textbf{Proteolipid protein (PLP):} Compaction of extracellular faces.
Mutations cause Pelizaeus-Merzbacher disease. - \textbf{Myelin
oligodendrocyte glycoprotein (MOG):} Outer surface; minor protein but
major autoantigen in some demyelinating diseases (e.g., MOG antibody
disease).

Myelin development continues into the third decade of life, particularly
in prefrontal cortex --- coinciding with maturation of executive
function. Adult oligodendrocyte progenitor cells (OPCs, NG2+ cells)
maintain a regenerative capacity, but it is incomplete after
demyelinating injury.

\subsubsection{Microglia --- Resident Macrophages of the
CNS}\label{microglia-resident-macrophages-of-the-cns}

Microglia derive from yolk-sac primitive macrophages that colonised the
developing CNS before formation of the blood-brain barrier. They are
\textbf{self-renewing} (independent of blood monocytes) and constitute
about 10\% of CNS cells.

\textbf{States and functions:} - \textbf{Resting (surveillance):} Highly
ramified processes constantly sample the parenchyma; turnover of most
brain volume every few hours. - \textbf{Activated (M1-like,
pro-inflammatory):} Upon detection of pathogens (TLR ligands), tissue
damage (DAMPs: ATP, HMGB1), or aggregated proteins (A\texorpdfstring{\ensuremath{\beta}}{beta}, \texorpdfstring{\ensuremath{\alpha}}{alpha}-synuclein) ---
retract processes, become amoeboid, secrete TNF-\texorpdfstring{\ensuremath{\alpha}}{alpha}, IL-1\texorpdfstring{\ensuremath{\beta}}{beta}, IL-6, NO;
phagocytose debris. - \textbf{Alternatively activated (M2-like,
repair):} Secrete IL-10, TGF-\texorpdfstring{\ensuremath{\beta}}{beta}; promote tissue repair. -
\textbf{Synaptic pruning:} Microglia tag underused synapses with
complement (C1q \texorpdfstring{\ensuremath{\to}}{->} C3) for phagocytic engulfment. Critical during
development; pathological in Alzheimer's (excessive pruning in
hippocampus) and schizophrenia (excessive pruning in prefrontal cortex
during adolescence).

\subsubsection{Blood-Brain Barrier --- Molecular
Architecture}\label{blood-brain-barrier-molecular-architecture}

The \textbf{blood-brain barrier (BBB)} is a selective barrier between
blood and brain parenchyma that excludes most polar molecules, ions,
large molecules, and pathogens, while allowing essential nutrients to
enter. It is formed by:

\begin{enumerate}
\def\labelenumi{\arabic{enumi}.}
\tightlist
\item
  \textbf{Tight junctions} between brain capillary endothelial cells,
  sealed by claudin-5, occludin, JAM-1, and zonula occludens (ZO-1,
  ZO-2). Transendothelial electrical resistance is about 1500
  \texorpdfstring{\ensuremath{\Omega}}{Omega}\texorpdfstring{\ensuremath{\cdot}}{.}cm\(^2\) --- orders of magnitude higher than peripheral capillaries
  (about 10 \texorpdfstring{\ensuremath{\Omega}}{Omega}\texorpdfstring{\ensuremath{\cdot}}{.}cm\(^2\)).
\item
  \textbf{Pericytes} ensheathing the abluminal endothelial surface ---
  regulate capillary diameter, BBB integrity, and angiogenesis.
\item
  \textbf{Astrocyte endfeet} with AQP4 channels --- form a
  near-continuous covering of capillaries and signal to maintain
  endothelial tight junction integrity.
\item
  \textbf{Basement membrane} between endothelial cells and astrocyte
  endfeet.
\end{enumerate}

\textbf{Transport mechanisms across the BBB:} - \textbf{Diffusion:}
Small lipophilic molecules (O\(_2\), CO\(_2\), ethanol, anaesthetics)
cross freely. - \textbf{Carrier-mediated transport:} GLUT1 (glucose);
LAT1 (large neutral amino acids --- competition explains why
high-protein meals can affect L-DOPA delivery); MCT1 (lactate, ketone
bodies). - \textbf{Receptor-mediated transcytosis:} Insulin, transferrin
(iron), leptin --- receptors mediate endocytosis on the luminal side and
exocytosis on the abluminal side. This is exploited to deliver
therapeutics: a ``Trojan horse'' strategy fuses drugs to
anti-transferrin-receptor antibodies. - \textbf{Active efflux pumps:}
\textbf{P-glycoprotein (P-gp/ABCB1)} and \textbf{BCRP/ABCG2} on the
luminal surface actively pump substrates back into blood. Many
lipophilic drugs that should enter the brain by passive diffusion are
pumped out --- explaining the limited brain penetration of many
chemotherapies and the role of P-gp polymorphisms in inter-individual
drug response.

\textbf{BBB-bypassing structures (circumventricular organs):}
Subfornical organ, area postrema (chemotrigger zone), median eminence,
neurohypophysis, OVLT, pineal gland --- lack tight junctions and have
fenestrated capillaries to allow hormones and circulating signals to
reach specialised neurons (e.g., AT\(_1\) on subfornical neurons \texorpdfstring{\ensuremath{\to}}{->}
thirst).

\begin{quote}
\textbf{Clinical Connection:} BBB breakdown contributes to many CNS
diseases. In multiple sclerosis, autoreactive T cells must first cross
the BBB (via VLA-4/VCAM-1 interaction); the monoclonal antibody
\textbf{natalizumab} blocks VLA-4 and dramatically reduces relapses. In
Alzheimer's, BBB pericyte dysfunction precedes overt neurodegeneration.
In ischaemic stroke, BBB breakdown produces vasogenic oedema and
haemorrhagic transformation. Pharmacologically, the BBB is the major
obstacle to CNS drug delivery --- about 2\% of small molecules and 0\%
of biologics readily cross.
\end{quote}

\begin{center}\rule{0.5\linewidth}{0.5pt}\end{center}

\subsection{Resting Membrane
Potential}\label{resting-membrane-potential}

\begin{figure}[htbp]
\centering
\includegraphics[width=0.85\textwidth]{../figures/action_potential.png}
\caption{Hodgkin--Huxley action-potential simulation showing membrane voltage over time, resting potential, and threshold.}
\label{fig:unit_IX_action_potential}
\end{figure}

\subsubsection{Ion Distribution and the Nernst
Equation}\label{ion-distribution-and-the-nernst-equation}

The resting membrane potential (about \(-70\) mV) arises from
unequal ion distributions maintained by the Na\(^+\)/K\(^+\)-ATPase and
selective membrane permeability.

\textbf{Typical neuronal ion concentrations:}

{\def\LTcaptype{none} % do not increment counter
\begin{longtable}[]{@{}
  >{\raggedright\arraybackslash}p{(\linewidth - 6\tabcolsep) * \real{0.0492}}
  >{\raggedright\arraybackslash}p{(\linewidth - 6\tabcolsep) * \real{0.3115}}
  >{\raggedright\arraybackslash}p{(\linewidth - 6\tabcolsep) * \real{0.3115}}
  >{\raggedright\arraybackslash}p{(\linewidth - 6\tabcolsep) * \real{0.3279}}@{}}
\toprule\noalign{}
\begin{minipage}[b]{\linewidth}\raggedright
Ion
\end{minipage} & \begin{minipage}[b]{\linewidth}\raggedright
Intracellular (mM)
\end{minipage} & \begin{minipage}[b]{\linewidth}\raggedright
Extracellular (mM)
\end{minipage} & \begin{minipage}[b]{\linewidth}\raggedright
Equilibrium Potential
\end{minipage} \\
\midrule\noalign{}
\endhead
\bottomrule\noalign{}
\endlastfoot
K\(^+\) & 140 & 5 & \(-89\) mV \\
Na\(^+\) & 12 & 145 & \(+67\) mV \\
Cl\(^-\) & 7 & 110 & \(-74\) mV \\
Ca\(^{2+}\) & 0.0001 & 2 & \(+132\) mV \\
\end{longtable}
}

\textbf{Nernst equation} for a single ion:

\begin{equation}
E_X = \frac{RT}{zF} \ln\frac{[X]_o}{[X]_i} = \frac{61.5 \text{ mV}}{z} \log_{10}\frac{[X]_o}{[X]_i} \quad \text{(at 37°C)}
\label{eq:nervous_system_1}
\end{equation}

\subsubsection{Goldman-Hodgkin-Katz (GHK)
Equation}\label{goldman-hodgkin-katz-ghk-equation}

The \hyperref[gl:resting-potential]{\textbf{resting potential}} is a
weighted average of equilibrium potentials, determined by relative
membrane permeability to each ion:

\begin{equation}
V_m = \frac{RT}{F} \ln\frac{P_K[K^+]_o + P_{Na}[Na^+]_o + P_{Cl}[Cl^-]_i}{P_K[K^+]_i + P_{Na}[Na^+]_i + P_{Cl}[Cl^-]_o}
\label{eq:nervous_system_2}
\end{equation}

At rest, \(P_K : P_{Na} : P_{Cl} \approx 1 : 0.04 : 0.45\), giving
\(V_m \approx -68\) mV.

The resting potential is dominated by \textbf{K\(^+\)} (highest resting
conductance via KCNK leak channels). The Na\(^+\)/K\(^+\)-ATPase (3
Na\(^+\) out, 2 K\(^+\) in per ATP) is electrogenic, contributing
about \(-3\) mV directly, but its main role is maintaining
the concentration gradients.

\textbf{Concept Check:} If you suddenly doubled the extracellular
K\(^+\) concentration from 5 mM to 10 mM, what would happen to E\(_K\)
and the resting membrane potential? Calculate using the Nernst equation.

\begin{center}\rule{0.5\linewidth}{0.5pt}\end{center}

\subsection{Cable Properties and Passive Signal
Spread}\label{cable-properties-and-passive-signal-spread}

Before a graded potential can trigger an action potential at the axon
hillock, it must travel electrotonically from the synapse. This passive
spread is governed by the \textbf{cable properties} of the dendrite,
modelled as a leaky electrical cable.

\subsubsection{Length Constant (\texorpdfstring{\ensuremath{\lambda}}{lambda})}\label{length-constant-ux3bb}

The \textbf{space constant} (or length constant) \texorpdfstring{\ensuremath{\lambda}}{lambda} determines how far a
steady-state voltage change spreads along a dendrite or axon:

\begin{equation}
\lambda = \sqrt{\frac{r_m}{r_i}}
\label{eq:nervous_system_3}
\end{equation}

where \(r_m\) = membrane resistance per unit length (\texorpdfstring{\ensuremath{\Omega}}{Omega}\texorpdfstring{\ensuremath{\cdot}}{.}cm) and \(r_i\) =
axial (cytoplasmic) resistance per unit length (\texorpdfstring{\ensuremath{\Omega}}{Omega}/cm).

\begin{equation}
V(x) = V_0 \, e^{-x/\lambda}
\label{eq:nervous_system_4}
\end{equation}

At distance \(x = \lambda\), the voltage has decayed to
\(V_0/e \approx 37\%\) of its peak value.

\textbf{Typical values:}

{\def\LTcaptype{none} % do not increment counter
\begin{longtable}[]{@{}
  >{\raggedright\arraybackslash}p{(\linewidth - 6\tabcolsep) * \real{0.1915}}
  >{\raggedright\arraybackslash}p{(\linewidth - 6\tabcolsep) * \real{0.2766}}
  >{\raggedright\arraybackslash}p{(\linewidth - 6\tabcolsep) * \real{0.2979}}
  >{\raggedright\arraybackslash}p{(\linewidth - 6\tabcolsep) * \real{0.2340}}@{}}
\toprule\noalign{}
\begin{minipage}[b]{\linewidth}\raggedright
Structure
\end{minipage} & \begin{minipage}[b]{\linewidth}\raggedright
Diameter (\texorpdfstring{\ensuremath{\mu}}{mu}m)
\end{minipage} & \begin{minipage}[b]{\linewidth}\raggedright
\texorpdfstring{\ensuremath{\lambda}}{lambda} (mm)
\end{minipage} & \begin{minipage}[b]{\linewidth}\raggedright
Implication
\end{minipage} \\
\midrule\noalign{}
\endhead
\bottomrule\noalign{}
\endlastfoot
Thin dendrite & 0.5 & about 0.1--0.2 & Distal synapses
attenuate strongly before reaching soma \\
Large dendrite & 5 & about 0.5--1.0 & Better passive
integration \\
Squid giant axon & 500 & about 5--7 & Enables long-distance
spread in unmyelinated axon \\
Myelinated axon & 10 & about 2--5 & Myelin increases \(r_m\),
dramatically increasing \texorpdfstring{\ensuremath{\lambda}}{lambda} \\
\end{longtable}
}

\textbf{Key determinants of \texorpdfstring{\ensuremath{\lambda}}{lambda}:} - \textbf{Myelin} increases \(r_m\)
about 100-fold (reduces current leak) \texorpdfstring{\ensuremath{\to}}{->} \texorpdfstring{\ensuremath{\lambda}}{lambda} increases
about 10-fold - \textbf{Larger diameter} reduces \(r_i\) (more
conducting cross-section) \texorpdfstring{\ensuremath{\to}}{->} \texorpdfstring{\ensuremath{\lambda}}{lambda} increases - \textbf{GABA\(_A\) receptor
activation} opens Cl\(^-\) channels, reducing \(r_m\) \texorpdfstring{\ensuremath{\to}}{->} \texorpdfstring{\ensuremath{\lambda}}{lambda} decreases ---
a quantitative mechanism for \textbf{shunting inhibition} (nearby
inhibitory synapses drastically reduce the effective length constant,
terminating EPSP integration before reaching the hillock)

\subsubsection{Time Constant (\texorpdfstring{\ensuremath{\tau}}{tau}\_m)}\label{time-constant-ux3c4_m}

The \textbf{membrane time constant} determines how fast the membrane
voltage changes in response to a current:

\begin{equation}
\tau_m = r_m \cdot C_m
\label{eq:nervous_system_5}
\end{equation}

where \(C_m \approx 1\) \texorpdfstring{\ensuremath{\mu}}{mu}F/cm\texorpdfstring{\textsuperscript{2}}{2} (specific membrane capacitance --- nearly
constant across cell types). Typical \(\tau_m \approx 10\)--\(20\) ms
for neurons.

\begin{equation}
V(t) = V_\infty \left(1 - e^{-t/\tau_m}\right)\quad \text{(voltage rising in response to a step current)}
\label{eq:nervous_system_6}
\end{equation}

\textbf{Significance:} \(\tau_m\) governs the time window for
\textbf{temporal summation} in neuronal membrane integration. If two
EPSPs arrive within \(\tau_m\) of each other at the same synapse, the
second EPSP begins before the membrane has fully repolarized from the
first, and the \hyperref[gl:depolarisation]{\textbf{depolarisation}}s
summate.

\subsection{Worked Example: Dendritic Length
Constant}\label{worked-example-dendritic-length-constant}

A cortical pyramidal neuron has \(r_m = 40{,}000\;\Omega\text{\texorpdfstring{\ensuremath{\cdot}}{.}cm}\) and
\(r_i = 200\;\Omega/\text{cm}\) for a typical basal dendrite (diameter
about 1 \texorpdfstring{\ensuremath{\mu}}{mu}m). What is the length constant?

\begin{equation}
\lambda = \sqrt{\frac{40{,}000}{200}} = \sqrt{200} \approx 14.1\;\text{cm?}
\label{eq:nervous_system_7}
\end{equation}

Wait --- the formula uses \emph{specific} membrane resistance per unit
\textbf{length} for a cylinder. For a dendrite of radius \(a\):

\begin{equation}
r_m(\text{per unit length}) = \frac{R_m}{2\pi a}, \quad r_i(\text{per unit length}) = \frac{R_i}{\pi a^2}
\label{eq:nervous_system_8}
\end{equation}

where \(R_m\) = specific membrane resistance (about 20\{,\}000
\texorpdfstring{\ensuremath{\Omega}}{Omega}\texorpdfstring{\ensuremath{\cdot}}{.}cm\texorpdfstring{\textsuperscript{2}}{2}) and \(R_i\) = cytoplasmic resistivity (about 100 \texorpdfstring{\ensuremath{\Omega}}{Omega}\texorpdfstring{\ensuremath{\cdot}}{.}cm).
For \(a = 0.5\;\mu\)m \(= 0.5 \times 10^{-4}\) cm:

\begin{equation}
r_m = \frac{20{,}000}{2\pi \times 5 \times 10^{-5}} \approx 6.37 \times 10^7\;\Omega/\text{cm}
\label{eq:nervous_system_9}
\end{equation}

\begin{equation}
r_i = \frac{100}{\pi \times (5 \times 10^{-5})^2} \approx 1.27 \times 10^{10}\;\Omega/\text{cm}
\label{eq:nervous_system_10}
\end{equation}

\begin{equation}
\lambda = \sqrt{\frac{r_m}{r_i}} = \sqrt{\frac{6.37 \times 10^7}{1.27 \times 10^{10}}} = \sqrt{0.00501} \approx 0.071\;\text{cm} = 0.71\;\text{mm}
\label{eq:nervous_system_11}
\end{equation}

This means a synapse on a thin distal dendrite 1.4 mm from the soma
delivers primarily \(e^{-2} \approx 13.5\%\) of its peak voltage to the
hillock --- illustrating how \textbf{dendritic location profoundly
affects synaptic efficacy}, and why proximal synapses (closer to the
hillock) have disproportionate influence.

\begin{quote}
\textbf{Clinical Connection:} Peripheral neuropathy (e.g., from
diabetes, chemotherapy, or Guillain-Barré syndrome) reduces \(r_m\) in
peripheral axons by disrupting myelin or causing axonal damage. The
resulting decrease in \texorpdfstring{\ensuremath{\lambda}}{lambda} causes voltage signals to decay faster along the
axon, slowing conduction and reducing EPSP efficacy at nerve-muscle
junctions --- producing the characteristic distal-to-proximal weakness
and sensory loss of a length-dependent neuropathy.
\end{quote}

\begin{center}\rule{0.5\linewidth}{0.5pt}\end{center}

\subsection{Graded Potentials and Synaptic
Integration}\label{graded-potentials-and-synaptic-integration}

\subsubsection{Graded Potentials}\label{graded-potentials}

\textbf{Excitatory postsynaptic potentials (EPSPs):} Depolarising events
(typically caused by Na\(^+\) or mixed cation influx through ionotropic
receptors like AMPA). Amplitude proportional to stimulus strength. Decay
with distance (electrotonic spread governed by the length constant \texorpdfstring{\ensuremath{\lambda}}{lambda}).

\textbf{Inhibitory postsynaptic potentials (IPSPs):} Hyperpolarising
events (Cl\(^-\) influx through GABA\(_A\) receptors or K\(^+\) efflux
through GABA\(_B\)-activated K\(^+\) channels). Shunting inhibition:
even if IPSP does not change \(V_m\) much, opening Cl\(^-\) channels
increases membrane conductance, reducing the effectiveness of nearby
EPSPs.

\subsubsection{Summation at the Axon
Hillock}\label{summation-at-the-axon-hillock}

\begin{figure}[htbp]
\centering
\includegraphics[width=0.82\linewidth,height=4.2in,keepaspectratio]{../figures/mermaid_inline/inline_mermaid_0158_e6c1f83cf706.png}
\caption{Neural circuit types and synaptic integration. Convergence allows integration of multiple inputs; divergence allows signal distribution. At the axon hillock, temporal and spatial summation of EPSPs and IPSPs determine whether threshold is reached.}
\end{figure}

If the net depolarisation at the axon hillock exceeds threshold
(about \(-55\) mV), an most-or-none action potential is
generated, with the stereotyped Na\(^+\)- and K\(^+\)-driven waveform of
rise, peak, and hyperpolarisation shown in
\cref{fig:unit_IX_action_potential}. The axon hillock has the highest
density of Nav1.2 and Nav1.6 channels, making it the lowest-threshold
site.

\textbf{Strategic placement of inhibition:} Inhibitory synapses
(GABA\(_A\)) are preferentially located on the soma and proximal
dendrites, positioned to maximally dampen axon hillock depolarisation.

\begin{center}\rule{0.5\linewidth}{0.5pt}\end{center}

\subsection{Autonomic Nervous System}\label{autonomic-nervous-system}

The \hyperref[gl:autonomic-nervous-system]{\textbf{autonomic nervous
system (ANS)}} regulates involuntary functions through two antagonistic
divisions:

\subsubsection{Sympathetic Division (``Fight or
Flight'')}\label{sympathetic-division-fight-or-flight}

\begin{itemize}
\tightlist
\item
  \textbf{Preganglionic neurons:} T1-L2 spinal cord (thoracolumbar
  outflow). Short preganglionic fibres to paravertebral (sympathetic
  chain) or prevertebral ganglia.
\item
  \textbf{Preganglionic neurotransmitter:} ACh at nicotinic receptors
\item
  \textbf{Postganglionic neurotransmitter:} Norepinephrine (NE) at \texorpdfstring{\ensuremath{\alpha}}{alpha} and
  \texorpdfstring{\ensuremath{\beta}}{beta} adrenergic receptors on target organs
\item
  \textbf{Exception:} Adrenal medulla -- preganglionic sympathetic
  fibres synapse directly on chromaffin cells (modified postganglionic
  neurons) that release epinephrine (80\%) and NE (20\%) into blood as
  \hyperref[gl:hormone]{\textbf{hormone}}s
\end{itemize}

\textbf{Sympathetic effects:} Increased HR (\(\beta_1\)),
bronchodilation (\(\beta_2\)), pupil dilation (mydriasis), inhibited GI
motility, glycogenolysis in liver (\(\beta_2\)), vasoconstriction
(\(\alpha_1\)) in skin/viscera, vasodilation (\(\beta_2\)) in skeletal
muscle.

\subsubsection{Parasympathetic Division (``Rest and
Digest'')}\label{parasympathetic-division-rest-and-digest}

\begin{itemize}
\tightlist
\item
  \textbf{Preganglionic neurons:} Brainstem (CN III, VII, IX, X) and
  sacral spinal cord (S2-S4; craniosacral outflow). Long preganglionic
  fibres to terminal ganglia near or within target organs.
\item
  \textbf{Both neurotransmitters:} ACh. Preganglionic: nicotinic
  receptors. Postganglionic: muscarinic receptors (M1-M5) on target
  organs.
\item
  \textbf{Vagus nerve (CN X):} Carries about 75\% of
  parasympathetic fibres; innervates heart, lungs, GI tract to splenic
  flexure.
\end{itemize}

\textbf{Parasympathetic effects:} Decreased HR (M\(_2\)),
bronchoconstriction (M\(_3\)), pupil constriction (miosis), increased GI
motility and secretion, bladder contraction.

\subsubsection{Enteric Nervous System}\label{enteric-nervous-system}

\begin{itemize}
\tightlist
\item
  about 500 million neurons in the gut wall (more than in the
  spinal cord)
\item
  \textbf{Myenteric (Auerbach's) plexus:} Between longitudinal and
  circular muscle layers; controls motility
\item
  \textbf{Submucosal (Meissner's) plexus:} Controls secretion and blood
  flow
\item
  Uses many neurotransmitters: ACh, NO, serotonin, substance P, VIP
\item
  Can function autonomously but is modulated by sympathetic (generally
  inhibitory) and parasympathetic (generally excitatory) input
\end{itemize}

\begin{quote}
\textbf{Clinical Connection:} Understanding ANS pharmacology is
fundamental to medicine. Beta-blockers (\(\beta_1\) antagonists:
metoprolol, atenolol) are first-line treatments for hypertension, heart
failure, and arrhythmias. Atropine (muscarinic antagonist) treats
bradycardia. Prazosin (\(\alpha_1\) antagonist) treats hypertension and
PTSD nightmares. Pilocarpine (muscarinic agonist) treats glaucoma by
promoting aqueous humour drainage.
\end{quote}

\begin{center}\rule{0.5\linewidth}{0.5pt}\end{center}

\subsection{Brain Anatomy and
Function}\label{brain-anatomy-and-function}

\begin{figure}[htbp]
\centering
\includegraphics[width=0.82\linewidth,height=4.2in,keepaspectratio]{../figures/mermaid_inline/inline_mermaid_0159_7b9f1e7b1bee.png}
\caption{Major brain regions and their functions. The cerebral cortex is divided into four lobes with specialised functions. Subcortical structures handle motor control (basal ganglia), sensory relay (thalamus), homeostasis (hypothalamus), memory (hippocampus), and emotion (amygdala). The cerebellum coordinates movement, and the brainstem controls vital functions.}
\end{figure}

\subsubsection{Cerebral Cortex}\label{cerebral-cortex}

The cortex is about 3 mm thick, contains about 20
billion neurons, and has a surface area of about 2,500 cm\(^2\)
(increased by folding into gyri and sulci).

\begin{itemize}
\tightlist
\item
  \textbf{Motor homunculus:} Topographic representation of body parts on
  the primary motor cortex (precentral gyrus). Hands and face have
  disproportionately large representations (fine motor control).
\item
  \textbf{Sensory homunculus:} Topographic representation on the primary
  somatosensory cortex (postcentral gyrus). Hands, lips, and tongue are
  overrepresented (high receptor density).
\item
  \textbf{Broca's area} (left inferior frontal gyrus): Speech
  production. Damage causes non-fluent (expressive) aphasia --
  comprehension intact but speech production impaired.
\item
  \textbf{Wernicke's area} (left posterior superior temporal gyrus):
  Language comprehension. Damage causes fluent (receptive) aphasia --
  speech flows but is nonsensical.
\end{itemize}

\subsubsection{Basal Ganglia}\label{basal-ganglia}

The basal ganglia (caudate, putamen, globus pallidus, subthalamic
nucleus, substantia nigra) modulate cortical motor output:

\begin{itemize}
\tightlist
\item
  \textbf{Direct pathway (Go):} Cortex excites striatum; striatum
  inhibits GPi/SNr (GABAergic); GPi/SNr normally inhibit thalamus. Net
  effect: disinhibition of thalamus, promoting movement.
\item
  \textbf{Indirect pathway (NoGo):} Cortex excites striatum; striatum
  inhibits GPe; GPe normally inhibits STN; STN excites GPi/SNr. Net
  effect: increased GPi/SNr inhibition of thalamus, suppressing
  movement.
\item
  \textbf{Dopamine:} Substantia nigra pars compacta (SNc) neurons
  release dopamine onto striatal neurons. D1 receptors excite direct
  pathway neurons; D2 receptors inhibit indirect pathway neurons. Net:
  dopamine facilitates movement.
\end{itemize}

\textbf{Parkinson's disease:} Degeneration of SNc dopaminergic neurons
(\textgreater60\% loss before symptoms appear). Loss of dopamine removes
facilitation of movement, causing: bradykinesia (slow movement),
rigidity, resting tremor, postural instability. Treatment: L-DOPA
(dopamine precursor crosses BBB); deep brain stimulation of STN.

\subsubsection{Other Key Structures}\label{other-key-structures}

\textbf{Cerebellum:} Contains more neurons than the rest of the brain
combined (about 70 billion granule cells). Functions: motor
coordination, error correction (compares intended vs actual movement),
timing, motor learning, balance. Damage causes ataxia (uncoordinated
movement) but not paralysis.

\textbf{Hippocampus:} Essential for consolidating declarative (episodic
and semantic) memories from short-term to long-term storage. Contains
place cells (fire at specific locations; O'Keefe, Nobel 2014) and grid
cells (fire in a hexagonal spatial pattern; Moser and Moser, Nobel
2014). One of few brain regions with confirmed adult neurogenesis
(dentate gyrus).

\textbf{Hypothalamus:} Homeostatic control centre with 13+ nuclei
controlling: temperature, hunger/satiety, thirst, circadian rhythms
(suprachiasmatic nucleus, SCN), autonomic output, and endocrine function
(connects to pituitary via the hypothalamo-hypophyseal portal system).

\textbf{Spinal cord:} Grey matter (butterfly-shaped, contains neuronal
cell bodies): dorsal horn (sensory processing), ventral horn (motor
neuron cell bodies), lateral horn (T1-L2: sympathetic preganglionic
neurons). White matter (surrounding, myelinated tracts): ascending
(sensory, e.g., spinothalamic tract for pain/temperature) and descending
(motor, e.g., corticospinal tract for voluntary movement).

\begin{center}\rule{0.5\linewidth}{0.5pt}\end{center}

\subsection{Sensory Systems}\label{sensory-systems}

\subsubsection{Somatosensation}\label{somatosensation}

{\def\LTcaptype{none} % do not increment counter
\begin{longtable}[]{@{}
  >{\raggedright\arraybackslash}p{(\linewidth - 6\tabcolsep) * \real{0.3333}}
  >{\raggedright\arraybackslash}p{(\linewidth - 6\tabcolsep) * \real{0.2051}}
  >{\raggedright\arraybackslash}p{(\linewidth - 6\tabcolsep) * \real{0.2564}}
  >{\raggedright\arraybackslash}p{(\linewidth - 6\tabcolsep) * \real{0.2051}}@{}}
\toprule\noalign{}
\begin{minipage}[b]{\linewidth}\raggedright
Receptor Type
\end{minipage} & \begin{minipage}[b]{\linewidth}\raggedright
Stimulus
\end{minipage} & \begin{minipage}[b]{\linewidth}\raggedright
Adaptation
\end{minipage} & \begin{minipage}[b]{\linewidth}\raggedright
Modality
\end{minipage} \\
\midrule\noalign{}
\endhead
\bottomrule\noalign{}
\endlastfoot
Meissner's corpuscle & Light touch, texture & Rapidly adapting & Fine
touch discrimination \\
Merkel's disc & Sustained pressure & Slowly adapting & Shape, edge
detection \\
Pacinian corpuscle & Vibration, deep pressure & Rapidly adapting &
Vibration detection \\
Ruffini ending & Skin stretch & Slowly adapting & Joint position,
stretch \\
Free nerve endings & Pain, temperature & Variable & Nociception,
thermoreception \\
\end{longtable}
}

\textbf{Nociception:} Two fibre types carry pain signals: -
\textbf{A-delta fibres:} Thinly myelinated (5-30 m/s). Sharp,
well-localised ``first pain'' - \textbf{C fibres:} Unmyelinated (0.5-2
m/s). Dull, diffuse ``second pain''

\subsubsection{Visual Pathway}\label{visual-pathway}

Retina (photoreceptors: rods for dim light/peripheral vision; cones for
colour/acuity) to optic nerve to optic chiasm (nasal fibres cross) to
\textbf{lateral geniculate nucleus (LGN)} of thalamus to primary visual
cortex (V1, striate cortex) in occipital lobe.

Beyond V1, visual processing splits into: - \textbf{Dorsal stream}
(``where/how''): V1 to posterior parietal cortex. Motion, spatial
relationships, visually guided action. - \textbf{Ventral stream}
(``what''): V1 to inferotemporal cortex. Object recognition, face
recognition, colour.

\subsubsection{Proprioception}\label{proprioception}

Muscle spindles (detect muscle length and stretch velocity), Golgi
tendon organs (detect muscle tension), and joint receptors provide
unconscious awareness of body position. Information travels via dorsal
column-medial lemniscal pathway (proprioception, fine touch) or
spinocerebellar tracts (to cerebellum for motor coordination).

\subsubsection{Musculoskeletal Control and
Behaviour}\label{musculoskeletal-control-and-behaviour}

Skeletal movement is a loop, not a one-way command. Alpha motor neurons
release acetylcholine at the neuromuscular junction; muscle fibres
depolarise, release Ca\(^{2+}\) from the sarcoplasmic reticulum, and
contract by the sliding-filament mechanism in which myosin heads
cyclically bind actin, pull, detach, and reset
\citep{huxley1954sliding}. Sensory feedback closes the loop: muscle
spindles report length, Golgi tendon organs report tension, cutaneous
mechanoreceptors report contact, and vestibular inputs report head
acceleration. Reflexes are therefore local control policies embedded in
a larger behavioural system, not primitive leftovers.

Behavioural biology adds four levels of explanation that should not be
collapsed: mechanism (the neural and hormonal circuit), development (how
the behaviour changes across life), function (what fitness problem it
solves), and evolutionary history (how related species differ)
\citep{tinbergen1963aims}. A startle reflex, bird song, courtship
display, or human reaching movement can be read through those four
levels. The organismal habit is to ask which sensory cue, motor
effector, motivational state, developmental window, and ecological
payoff are actually supported by the evidence.

\begin{center}\rule{0.5\linewidth}{0.5pt}\end{center}

\subsection{Brain Imaging --- Reading Activity Through Indirect
Signals}\label{brain-imaging-reading-activity-through-indirect-signals}

Modern neuroscience and clinical neurology rely on a portfolio of
imaging modalities, each measuring a different physical correlate of
neural activity.

\subsubsection{fMRI and the BOLD signal}\label{fmri-and-the-bold-signal}

\textbf{Functional MRI (fMRI)} measures the
\textbf{blood-oxygen-level-dependent (BOLD)} signal, which exploits a
magnetic-resonance peculiarity of haemoglobin: oxyhaemoglobin is
\textbf{diamagnetic}, while deoxyhaemoglobin is \textbf{paramagnetic}
and distorts the local T2\emph{-weighted MR signal. Active brain regions
paradoxically }increase* local oxyhaemoglobin (and decrease
deoxyhaemoglobin) within seconds, because the \textbf{neurovascular
coupling} response over-supplies blood relative to the metabolic demand.
The mismatch between blood-flow increase (about 50\%) and
O\(_2\) extraction increase (about 5--20\%) produces the
\textbf{BOLD signal rise}, peaking about 5 s after activity
onset.

The neurovascular coupling that drives BOLD is mediated largely by
astrocytes (recall the astrocyte tripartite synapse): synaptic glutamate
elevates astrocytic Ca\(^{2+}\), which releases vasoactive arachidonic
acid metabolites (PGE\(_2\), EETs) onto local arterioles. The
astrocyte-neuron lactate shuttle (ANLS) is the metabolic coupling
underlying fMRI BOLD signals.

BOLD does \emph{not} measure spike rates directly; it measures the
haemodynamic response to recent local-field-potential activity. fMRI
excels at spatial localisation (about 1--3 mm) but trades
temporal resolution (seconds) --- the inverse trade-off compared with
EEG.

\subsubsection{EEG and oscillation
bands}\label{eeg-and-oscillation-bands}

\textbf{Electroencephalography (EEG)} measures the summed extracellular
electrical potentials produced by synchronised cortical pyramidal-neuron
currents. EEG has millisecond temporal resolution but limited spatial
localisation (about 1--2 cm). Brain rhythms occupy distinct
frequency bands tied to behavioural states:

{\def\LTcaptype{none} % do not increment counter
\begin{longtable}[]{@{}
  >{\raggedright\arraybackslash}p{(\linewidth - 6\tabcolsep) * \real{0.0769}}
  >{\raggedright\arraybackslash}p{(\linewidth - 6\tabcolsep) * \real{0.2692}}
  >{\raggedright\arraybackslash}p{(\linewidth - 6\tabcolsep) * \real{0.2692}}
  >{\raggedright\arraybackslash}p{(\linewidth - 6\tabcolsep) * \real{0.3846}}@{}}
\toprule\noalign{}
\begin{minipage}[b]{\linewidth}\raggedright
Band
\end{minipage} & \begin{minipage}[b]{\linewidth}\raggedright
Frequency (Hz)
\end{minipage} & \begin{minipage}[b]{\linewidth}\raggedright
Dominant state
\end{minipage} & \begin{minipage}[b]{\linewidth}\raggedright
Generator / function
\end{minipage} \\
\midrule\noalign{}
\endhead
\bottomrule\noalign{}
\endlastfoot
\textbf{Delta (\texorpdfstring{\ensuremath{\delta}}{delta})} & 0.5--4 & Deep NREM sleep (N3); coma; infancy &
Thalamocortical bursts; slow-wave sleep; memory consolidation \\
\textbf{Theta (\texorpdfstring{\ensuremath{\theta}}{theta})} & 4--8 & Drowsiness; REM; hippocampal navigation &
Hippocampal ``theta rhythm''; encoding of new memories \\
\textbf{Alpha (\texorpdfstring{\ensuremath{\alpha}}{alpha})} & 8--13 & Awake, eyes-closed, relaxed (occipital) &
Posterior cortical idling; suppressed by visual attention \\
\textbf{Beta (\texorpdfstring{\ensuremath{\beta}}{beta})} & 13--30 & Awake, focused; anxious; motor tasks &
Cortical activation; motor planning (sensorimotor cortex) \\
\textbf{Gamma (\texorpdfstring{\ensuremath{\gamma}}{gamma})} & 30--80+ & Sensory binding; conscious perception &
Local cortical inhibitory networks (PV interneurons) \\
\end{longtable}
}

Clinical EEG is the bedside tool for diagnosing \textbf{epilepsy}
(interictal spikes; rhythmic ictal patterns), \textbf{encephalopathy}
(generalised slowing), and \textbf{brain death} (electrocerebral
silence). MEG (magnetoencephalography) measures the same neuronal
currents via tiny magnetic fields with better localisation than EEG.

\subsubsection{Connectivity, MEG, NIRS,
PET}\label{connectivity-meg-nirs-pet}

\begin{itemize}
\tightlist
\item
  \textbf{Diffusion tensor imaging (DTI):} Maps white-matter tracts by
  measuring water diffusion anisotropy along axons. Used in stroke,
  traumatic axonal injury, neurosurgical planning.
\item
  \textbf{Resting-state fMRI:} Identifies networks of co-fluctuating
  brain regions (default mode network, salience network, frontoparietal
  control). Disturbances correlate with depression, schizophrenia,
  Alzheimer's.
\item
  \textbf{PET (positron emission tomography):} Tracer-based metabolic or
  receptor-binding imaging; FDG-PET shows glucose uptake (epilepsy
  localisation, dementia patterns); amyloid-PET (florbetapir) and
  tau-PET enable Alzheimer's disease characterisation.
\item
  \textbf{NIRS (near-infrared spectroscopy):} Optical analogue of fMRI;
  portable; used in neonatal monitoring and bedside cerebral oximetry.
\end{itemize}

\subsubsection{Neural Prosthetics and Brain-Computer
Interfaces}\label{neural-prosthetics-and-brain-computer-interfaces}

Brain-computer interfaces (BCIs) and neural prosthetics convert neural
activity into an external action or stimulation pattern. The strongest
current speech demonstrations are still experimental implants, not
consumer mind-reading devices: 2023 intracortical speech neuroprostheses
decoded attempted speech into text or avatar control at high performance
in single-participant research settings
\citep{willett2023speechneuroprosthesis,metzger2023avatarneuroprosthesis}.
The scientific claim is therefore bounded: motor and speech cortices
contain decodable intention signals, but clinical translation requires
stable implants, calibration, infection and tissue-response management,
privacy safeguards, and evidence that performance generalises beyond
highly supported trials.

\begin{center}\rule{0.5\linewidth}{0.5pt}\end{center}

\subsection{Pain Pathways}\label{pain-pathways}

Pain is the body's most clinically important sensory modality and serves
both protective and pathological roles.

\subsubsection{Nociceptors and
transduction}\label{nociceptors-and-transduction}

Free nerve endings of A-delta and C fibres express \textbf{molecular
sensors} that transduce noxious stimuli into membrane depolarisation:

\begin{itemize}
\tightlist
\item
  \textbf{TRPV1:} Capsaicin- and heat-activated (\textgreater43 °C);
  also activated by acidic pH and inflammatory lipids. Resiniferatoxin
  and lidocaine-soaked TRPV1 agonist patches deplete nociceptors of
  substance P.
\item
  \textbf{TRPM8:} Cold- and menthol-activated (15--25 °C).
\item
  \textbf{TRPA1:} Mustard oil, cinnamaldehyde, environmental irritants;
  cold-activated in some species.
\item
  \textbf{ASICs (acid-sensing ion channels):} Activated by tissue
  acidosis (ischaemia, inflammation).
\item
  \textbf{Nav1.7 (SCN9A):} Voltage-gated Na\(^+\) channel essential for
  nociceptor excitability. \emph{Loss-of-function} mutations cause
  \textbf{congenital insensitivity to pain (CIP)} --- patients fail to
  detect injuries and self-mutilate. \emph{Gain-of-function} mutations
  cause \textbf{inherited erythromelalgia} (severe burning pain).
  Nav1.7-selective inhibitors are in clinical development as non-opioid
  analgesics.
\end{itemize}

\subsubsection{Ascending pathway --- the spinothalamic
tract}\label{ascending-pathway-the-spinothalamic-tract}

A-delta and C fibres release glutamate \textbf{and substance P} into the
dorsal horn of the spinal cord (lamina I, II, V). Second-order neurons
cross the midline (anterior white commissure), ascend in the
contralateral \textbf{spinothalamic tract}, and synapse in the ventral
posterolateral (VPL) nucleus of the thalamus. Third-order neurons
project to primary somatosensory cortex (S1), insula (interoceptive),
and anterior cingulate (affective dimension of pain). Substance P binds
NK1 receptors on second-order neurons; the combined glutamate-Substance
P transmission distinguishes pain from innocuous somatosensation.

\subsubsection{Descending modulation}\label{descending-modulation}

The brain actively \textbf{suppresses} ongoing pain signalling through
descending circuits:

\begin{itemize}
\tightlist
\item
  \textbf{Periaqueductal grey (PAG)} in the midbrain receives input from
  cortex, hypothalamus, and amygdala
\item
  PAG activates the \textbf{rostral ventromedial medulla (RVM)}
\item
  RVM serotonergic neurons project to the \textbf{dorsal horn (DH)},
  where they release serotonin, noradrenaline, and endogenous opioids
  onto pain-transmission neurons
\end{itemize}

This \textbf{PAG \texorpdfstring{\ensuremath{\to}}{->} RVM \texorpdfstring{\ensuremath{\to}}{->} DH} axis is the substrate for
\textbf{stress-induced analgesia} (hand-on-flame reflex without pain in
life-or-death situations) and for the placebo effect. Endogenous opioids
(\texorpdfstring{\ensuremath{\beta}}{beta}-endorphin, enkephalin, dynorphin) acting on \texorpdfstring{\ensuremath{\mu}}{mu}-, \texorpdfstring{\ensuremath{\delta}}{delta}-, and \texorpdfstring{\ensuremath{\kappa}}{kappa}-opioid
receptors are released along this axis. \textbf{Morphine and other
opioid analgesics} exploit the same descending pathway.

\textbf{Gate control theory} (Melzack \& Wall, 1965): A-\texorpdfstring{\ensuremath{\beta}}{beta} fibres
(mechanoreceptors, ``innocuous touch'') activate inhibitory interneurons
in the dorsal horn that ``gate'' C-fibre transmission --- the basis for
\textbf{rubbing a stubbed toe} to reduce pain, and the rationale for
\textbf{TENS (transcutaneous electrical nerve stimulation)} therapy.

\begin{quote}
\textbf{Clinical Connection:} Chronic neuropathic pain (post-herpetic
neuralgia, diabetic neuropathy) reflects maladaptive central
sensitisation --- peripheral nerve injury upregulates Nav1.7/1.8 in
nociceptors, expands receptive fields, and reduces descending
inhibition. Treatments include gabapentinoids (target \texorpdfstring{\ensuremath{\alpha}}{alpha}\(_2\delta\)
subunits of voltage-gated Ca\(^{2+}\) channels), SNRIs (boost descending
serotonin/noradrenaline), and topical lidocaine (Nav blocker).
\end{quote}

\begin{center}\rule{0.5\linewidth}{0.5pt}\end{center}

\subsection{Sleep --- Two-Process Model}\label{sleep-two-process-model}

Sleep is not the absence of brain activity but a structured succession
of states with distinct functions.

\subsubsection{Architecture}\label{architecture}

Adult sleep cycles between \textbf{NREM} (further divided into N1, N2,
N3) and \textbf{REM} stages, each cycle about 90 minutes:

\begin{itemize}
\tightlist
\item
  \textbf{N1 (light):} Theta activity; transition from waking
\item
  \textbf{N2:} Sleep spindles (12--14 Hz thalamic bursts) and
  K-complexes; about 50\% of total sleep
\item
  \textbf{N3 (slow-wave / deep):} Delta-dominant; ``homeostatic'' sleep
  restoring tissue and consolidating declarative memory;
  \textbf{glymphatic clearance} of brain metabolites including amyloid-\texorpdfstring{\ensuremath{\beta}}{beta}
\item
  \textbf{REM (paradoxical):} Low-amplitude desynchronised EEG
  resembling waking; rapid eye movements; \textbf{muscle atonia}
  (mediated by glycinergic inhibition of motor neurons by sublateral
  dorsal nucleus); vivid dreams; consolidation of procedural and
  emotional memory
\end{itemize}

REM proportion is maximal in infancy (about 50\% of sleep) and
declines with age. Sleep deprivation produces a ``REM rebound'' --- REM
disproportionately recovers after deprivation, suggesting it serves a
critical function.

\subsubsection{The two-process model (Borbély,
1982)}\label{the-two-process-model-borbuxe9ly-1982}

Sleep timing is regulated by two interacting processes:

\begin{itemize}
\tightlist
\item
  \textbf{Process C (Circadian):} Sinusoidal drive from the
  suprachiasmatic nucleus (SCN) of the hypothalamus, entrained by light
  via the retinohypothalamic tract. Promotes wakefulness during the day
  and sleep at night; mediated by hypothalamic circuits including the
  ventrolateral preoptic nucleus (VLPO, sleep-promoting) and
  orexin/hypocretin neurons of the lateral hypothalamus
  (wake-promoting).
\item
  \textbf{Process S (Sleep homeostasis):} Increases monotonically during
  waking and dissipates exponentially during sleep. Its molecular
  substrate is largely \textbf{adenosine}, which accumulates in basal
  forebrain during sustained activity (from local ATP catabolism).
  Adenosine acting on A\(_1\) receptors inhibits wake-promoting neurons;
  acting on A\(_{2A}\) receptors in the ventral striatum, it promotes
  sleep. \textbf{Caffeine} (an adenosine A\(_1\)/A\(_{2A}\) antagonist)
  blocks Process S --- this is its mechanism for reducing sleep
  pressure.
\end{itemize}

Sleep occurs when Process C (low) coincides with Process S (high).
Disruption of either process produces sleep disorders: SCN damage (or
shift work) disrupts Process C, producing irregular sleep timing;
insomnia often reflects elevated arousal that overrides Process S.

\subsubsection{Functions of sleep}\label{functions-of-sleep}

\begin{itemize}
\tightlist
\item
  \textbf{Memory consolidation:} Slow-wave sleep replays hippocampal
  place-cell sequences; REM consolidates emotional and procedural
  memory.
\item
  \textbf{Glymphatic clearance:} During NREM, the brain's glymphatic
  (glial-lymphatic) system expands extracellular space
  about 60\% and accelerates CSF--ISF exchange, clearing
  amyloid-\texorpdfstring{\ensuremath{\beta}}{beta}, tau, and other metabolites. Chronic sleep deprivation is a
  recognised risk factor for Alzheimer's disease.
\item
  \textbf{Synaptic homeostasis (Tononi-Cirelli hypothesis):} Wake-time
  learning produces net synaptic potentiation; sleep allows broad
  synaptic downscaling, restoring information capacity.
\item
  \textbf{Metabolic and immune restoration:} GH and prolactin peak
  during deep sleep; cytokine cycling supports immune surveillance.
\end{itemize}

\begin{quote}
\textbf{Clinical Connection:} Narcolepsy type 1 results from loss of
orexin/hypocretin neurons (autoimmune destruction in genetically
susceptible individuals). Without orexin, Process C cannot stably
maintain wakefulness, producing daytime sleepiness, REM intrusion into
wake (cataplexy, sleep paralysis, hallucinations). Treatment includes
modafinil (wake-promoting), sodium oxybate (consolidates sleep), and
emerging orexin receptor agonists.
\end{quote}

\begin{center}\rule{0.5\linewidth}{0.5pt}\end{center}

\subsection{Neural Plasticity}\label{neural-plasticity}

\textbf{Hebbian learning:} ``Neurons that fire together wire together''
\citep{hebb1949}. Correlated pre- and postsynaptic activity strengthens
synaptic connections.

\textbf{Long-term potentiation (LTP):} Activity-dependent, long-lasting
increase in synaptic efficacy. At Schaffer collateral to CA1 synapses in
hippocampus: 1. High-frequency stimulation depolarises postsynaptic
membrane enough to relieve Mg\(^{2+}\) block of NMDA receptors 2.
Ca\(^{2+}\) influx through NMDA receptors activates CaMKII 3. CaMKII
phosphorylates AMPA receptors (increased conductance) and promotes AMPA
receptor insertion into synapse 4. Structural plasticity: Dendritic
spine enlargement, new spine formation

\textbf{Long-term depression (LTD):} Opposite of LTP. Low-frequency
stimulation activates protein phosphatases (PP1, calcineurin) that
dephosphorylate AMPA receptors and promote their
\hyperref[gl:endocytosis]{\textbf{endocytosis}}. Weakens synaptic
connections.

\textbf{Homeostatic plasticity:} Synaptic scaling adjusts the overall
strength of most synapses on a neuron to maintain a stable firing rate.
If a neuron's activity is chronically reduced, it scales up most
synaptic strengths; if chronically increased, it scales them down. This
prevents runaway excitation or silencing.

\textbf{Critical periods \citep{beggs2003}:} Windows during development
when experience has maximal impact on neural circuit formation.
Monocular deprivation during the visual critical period
(about  birth to age 5 in humans) permanently reduces cortical
responsiveness to the deprived eye (amblyopia). Molecular brakes on
plasticity include perineuronal nets (PNNs), myelin-associated
inhibitors, and GABA maturation.

\textbf{Concept Check:} Explain why the NMDA receptor is called a
``coincidence detector'' and why this property makes it ideal for
Hebbian learning.

\begin{center}\rule{0.5\linewidth}{0.5pt}\end{center}

\subsection{Worked Example}\label{worked-example-2}

\textbf{Problem:} Calculate the equilibrium potential for potassium
(\(E_K\)) in a mammalian neuron at \(37^\circ\text{C}\), given an
intracellular potassium concentration \([K^+]_i\) of \(140\text{ mM}\)
and an extracellular concentration \([K^+]_o\) of \(5\text{ mM}\).

The Nernst equation for a monovalent cation at \(37^\circ\text{C}\) is:
\[E_K = 61.5\text{ mV} \times \log_{10}\left(\frac{[K^+]_o}{[K^+]_i}\right) \label{eq:unit_IX_nervous_system_item_1}\]

\textbf{Solution:}

\textbf{Step 1. Identify the given variables.} -
\([K^+]_o = 5\text{ mM}\) - \([K^+]_i = 140\text{ mM}\)

\textbf{Step 2. Substitute the values into the Nernst equation.}
\[E_K = 61.5\text{ mV} \times \log_{10}\left(\frac{5}{140}\right) \label{eq:unit_IX_nervous_system_item_2}\]

\textbf{Step 3. Calculate the logarithm.}
\[\frac{5}{140} \approx 0.0357 \label{eq:unit_IX_nervous_system_item_3}\]

\[\log_{10}(0.0357) \approx -1.447 \label{eq:unit_IX_nervous_system_item_4}\]

\textbf{Step 4. Determine the final potential.}
\[E_K = 61.5\text{ mV} \times (-1.447) \approx -89\text{ mV} \label{eq:unit_IX_nervous_system_item_5}\]

\textbf{Answer:} The equilibrium potential for potassium is
approximately \textbf{\(-89\text{ mV}\)}.

\begin{center}\rule{0.5\linewidth}{0.5pt}\end{center}

\subsection{Worked Example: Resting Potential from the
Goldman-Hodgkin-Katz
Equation}\label{worked-example-resting-potential-from-the-goldman-hodgkin-katz-equation}

\textbf{Problem:} Calculate the resting membrane potential of a
mammalian neuron at \(37^\circ\text{C}\) using the Goldman-Hodgkin-Katz
(GHK) equation, given the ion concentrations \([K^+]_o = 5\text{ mM}\),
\([K^+]_i = 140\text{ mM}\), \([Na^+]_o = 145\text{ mM}\),
\([Na^+]_i = 12\text{ mM}\), \([Cl^-]_o = 110\text{ mM}\),
\([Cl^-]_i = 7\text{ mM}\), and the resting permeability ratio
\(P_K : P_{Na} : P_{Cl} = 1 : 0.04 : 0.45\).

The GHK voltage equation, written in base-10 form for a
\(37^\circ\text{C}\) neuron, is:
\[V_m = 61.5\text{ mV} \times \log_{10}\left(\frac{P_K[K^+]_o + P_{Na}[Na^+]_o + P_{Cl}[Cl^-]_i}{P_K[K^+]_i + P_{Na}[Na^+]_i + P_{Cl}[Cl^-]_o}\right)\]

Note that the Cl\(^-\) terms are inverted relative to the cations (a
consequence of the \(z = -1\) charge): the outside-flux numerator
carries \([Cl^-]_i\) and the inside-flux denominator carries
\([Cl^-]_o\).

\textbf{Solution:}

\textbf{Step 1. Form the numerator (outward-driving terms), in mM.}
\[P_K[K^+]_o + P_{Na}[Na^+]_o + P_{Cl}[Cl^-]_i = (1)(5) + (0.04)(145) + (0.45)(7) = 5 + 5.8 + 3.15 = 13.95\]

\textbf{Step 2. Form the denominator (inward-driving terms), in mM.}
\[P_K[K^+]_i + P_{Na}[Na^+]_i + P_{Cl}[Cl^-]_o = (1)(140) + (0.04)(12) + (0.45)(110) = 140 + 0.48 + 49.5 = 189.98\]

\textbf{Step 3. Take the ratio and its base-10 logarithm.}
\[\frac{13.95}{189.98} \approx 0.0734, \qquad \log_{10}(0.0734) \approx -1.134\]

\textbf{Step 4. Scale by the \(37^\circ\text{C}\) Nernst slope.}
\[V_m = 61.5\text{ mV} \times (-1.134) \approx -69.7\text{ mV}\]

\textbf{Answer:} The resting membrane potential is approximately
\textbf{\(-70\text{ mV}\)} --- close to \(E_K\) (\(-89\) mV) but
depolarised relative to it because the finite Na\(^+\) and Cl\(^-\)
permeabilities pull \(V_m\) toward their (more positive) equilibrium
potentials. The K\(^+\) term dominates because \(P_K\) is the largest
permeability.

\textbf{Concept Check:} If \(P_{Na}\) rose tenfold (from 0.04 to 0.40)
while \(P_K\) and \(P_{Cl}\) were unchanged, predict the new sign of the
numerator-vs-denominator imbalance and whether \(V_m\) moves toward
\(E_{Na}\) or \(E_K\).

\textbf{Concept Check:} At rest, which single term in the GHK numerator
and denominator contributes the most to \(V_m\), and predict how \(V_m\)
changes if \(P_K\) is doubled while most concentrations are held fixed.

\begin{center}\rule{0.5\linewidth}{0.5pt}\end{center}

\subsection{Current Evidence and Frontier
Biology}\label{current-evidence-and-frontier-biology-31}

For \textbf{Nervous System and Neural Signalling}, frontier biology
belongs inside the evidence logic of the chapter. Physiology now blends
mechanism with allostasis, immune-endocrine-neural coupling, wearable
data, and individualized risk without reducing bodies to simple
machines. The core reading question is this: neural explanations should
separate circuit architecture, glial support, plasticity, behaviour, and
evidence scale.

\begin{itemize}
\tightlist
\item
  \textbf{What to verify:} identify the observation, model, assay, or
  dataset that would make the claim stronger or weaker.
\item
  \textbf{What to qualify:} state the scale, organism, cell type,
  environmental condition, or population where the claim is expected to
  hold.
\item
  \textbf{What to compare:} test at least one alternative explanation,
  baseline, or null model before treating the pattern as causal.
\item
  \textbf{What to cite:} distinguish primary evidence, review synthesis,
  public dataset, and institutional guidance; for recent or numeric
  claims, prefer the source closest to the measurement and state what
  has changed since it was published.
\end{itemize}

Interpret physiological data by separating baseline variation,
perturbation response, compensation, and the threshold where
compensation becomes pathology.

\textbf{Source practice:} For physiology claims, cite the measurement
context and distinguish baseline variation, compensation,
pathophysiology, and treatment evidence.

\subsection{Key Terms}\label{key-terms-31}

{\def\LTcaptype{none} % do not increment counter
\begin{longtable}[]{@{}
  >{\raggedright\arraybackslash}p{(\linewidth - 2\tabcolsep) * \real{0.2857}}
  >{\raggedright\arraybackslash}p{(\linewidth - 2\tabcolsep) * \real{0.7143}}@{}}
\toprule\noalign{}
\begin{minipage}[b]{\linewidth}\raggedright
Term
\end{minipage} & \begin{minipage}[b]{\linewidth}\raggedright
Definition
\end{minipage} \\
\midrule\noalign{}
\endhead
\bottomrule\noalign{}
\endlastfoot
\textbf{Resting membrane potential} & Electrochemical potential across
neuronal membrane at rest; about \(-70\) mV \\
\textbf{Nernst potential} & Voltage at which net flux of a single ion
species = 0 \\
\textbf{GHK equation} & Multi-ion extension of Nernst; weights each ion
by membrane permeability \\
\textbf{EPSP} & Excitatory postsynaptic potential: depolarisation toward
threshold \\
\textbf{IPSP} & Inhibitory postsynaptic potential: hyperpolarisation or
shunting inhibition \\
\textbf{Axon hillock} & Lowest threshold region; highest Nav channel
density; AP initiation site \\
\textbf{Tripartite synapse} & Presynaptic + postsynaptic + astrocyte
glial process \\
\textbf{Sympathetic NS} & Thoracolumbar outflow; fight or flight; NE on
adrenergic receptors \\
\textbf{Parasympathetic NS} & Craniosacral outflow; rest and digest; ACh
on muscarinic receptors \\
\textbf{Basal ganglia} & Subcortical nuclei modulating motor output;
direct (Go) and indirect (NoGo) pathways \\
\textbf{LTP} & Long-term potentiation; Hebbian strengthening of synaptic
connections \\
\textbf{LTD} & Long-term depression; weakening of synaptic
connections \\
\textbf{Broca's area} & Left inferior frontal gyrus; speech
production \\
\textbf{Wernicke's area} & Left superior temporal gyrus; language
comprehension \\
\textbf{Place cells} & Hippocampal neurons firing at specific spatial
locations \\
\textbf{Microglia} & CNS resident macrophages; synaptic pruning;
neuroinflammation \\
\textbf{Myelin} & Lipid-rich insulation produced by oligodendrocytes
(CNS) or Schwann cells (PNS) \\
\textbf{Na\(^+\)/K\(^+\)-ATPase} & Pump maintaining ion gradients; 3
Na\(^+\) out, 2 K\(^+\) in per ATP \\
\end{longtable}
}

\begin{center}\rule{0.5\linewidth}{0.5pt}\end{center}

\subsection{Review Questions}\label{review-questions-31}

\begin{enumerate}
\def\labelenumi{\arabic{enumi}.}
\item
  Calculate E\(_K\) using the Nernst equation for a neuron with
  {[}K\(^+\){]}\(_{in}\) = 150 mM and {[}K\(^+\){]}\(_{out}\) = 4 mM at
  37 degrees C. How does this compare to the resting potential of
  \(-70\) mV, and what accounts for the difference?
\item
  Compare the sympathetic and parasympathetic divisions in terms of: (a)
  spinal cord origin, (b) preganglionic fibre length, (c)
  neurotransmitters at ganglia and target organs, (d) receptor types,
  and (e) effects on heart rate.
\item
  A patient presents with non-fluent speech but intact comprehension.
  Which brain region is likely damaged? What would you expect on MRI?
  How does this differ from damage to Wernicke's area?
\item
  Explain the direct and indirect pathways of the basal ganglia. How
  does loss of dopaminergic input from SNc (as in Parkinson's disease)
  shift the balance between these pathways, and why does this cause
  bradykinesia?
\item
  Explain synaptic pruning and its relationship to microglia and
  complement proteins. What evidence links synaptic pruning
  dysregulation to schizophrenia risk? (Sekar et al.~2016,
  \emph{Nature}: C4A).
\item
  A researcher applies a blocker of KCNK leak potassium channels to a
  neuron. Predict the effect on (a) resting membrane potential, (b) the
  relative contribution of Na\(^+\) to the resting potential, and (c)
  neuronal excitability.
\item
  Compare ionotropic and metabotropic receptors in terms of structure,
  speed, ionic selectivity, and role in modulating synaptic gain. Give
  one clinical example for each class.
\item
  Explain why inhibitory synapses are strategically placed on the soma
  and proximal dendrites rather than on distal dendrites. How does the
  concept of shunting inhibition differ from hyperpolarising inhibition?
\item
  A dendrite has specific membrane resistance
  \(R_m = 25{,}000\;\Omega\text{\texorpdfstring{\ensuremath{\cdot}}{.}cm}^2\), cytoplasmic resistivity
  \(R_i = 80\;\Omega\text{\texorpdfstring{\ensuremath{\cdot}}{.}cm}\), and radius \(a = 1\;\mu\)m. Calculate
  the length constant \texorpdfstring{\ensuremath{\lambda}}{lambda}. If a synapse is located 2 mm from the soma,
  what fraction of the original EPSP amplitude reaches the soma? How
  would myelination of the dendrite change your answer qualitatively?
\item
  A patient with chronic inflammatory demyelinating
  polyradiculoneuropathy (CIDP) has reduced conduction velocity in
  peripheral sensory nerves and complains of ``glove-and-stocking''
  numbness. Using the cable properties framework, explain: (a) why
  demyelination slows conduction velocity, (b) why sensory deficits are
  worse distally (length-dependent), and (c) why IVIG treatment
  (suppressing antibody-mediated myelin attack) can restore sensation.
\item
  Estimate \texorpdfstring{\ensuremath{\lambda}}{lambda} from \breaktt{cable\_voltage\_attenuation} defaults and
  compare to the hand-calculated \texorpdfstring{\ensuremath{\lambda}}{lambda} in Question 9.
\item
  Why does \textbf{ephaptic coupling} complicate the independence
  assumption of parallel dendritic inputs?
\end{enumerate}

\begin{center}\rule{0.5\linewidth}{0.5pt}\end{center}

\subsection{Further Reading and Source
Notes}\label{further-reading-and-source-notes-26}

\begin{itemize}
\tightlist
\item
  Sherrington (1906). \emph{The Integrative Action of the Nervous
  System}. Yale University Press.
\item
  Hebb (1949). \emph{The Organization of Behavior: A Neuropsychological
  Theory}. Wiley.
\item
  Beggs \& Plenz (2003). Neuronal avalanches in neocortical circuits.
  \emph{Journal of Neuroscience}, 23.
\end{itemize}

\begin{center}\rule{0.5\linewidth}{0.5pt}\end{center}

\subsection{Computational Bridge}\label{computational-bridge-28}

Passive cable length constant is returned with the attenuation profile:

\begin{Shaded}
\begin{Highlighting}[]
\ImportTok{from}\NormalTok{ biology.neuroscience }\ImportTok{import}\NormalTok{ cable\_voltage\_attenuation}

\NormalTok{c }\OperatorTok{=}\NormalTok{ cable\_voltage\_attenuation(}\FloatTok{2.0}\NormalTok{, max\_distance\_µm}\OperatorTok{=}\FloatTok{800.0}\NormalTok{, n\_points}\OperatorTok{=}\DecValTok{20}\NormalTok{)}
\BuiltInTok{print}\NormalTok{(}\BuiltInTok{round}\NormalTok{(c.lambda\_µm, }\DecValTok{2}\NormalTok{))}
\end{Highlighting}
\end{Shaded}

\begin{quote}
\textbf{Clinical / systems note:} Local anaesthetics shorten the
depolarisation segment that electrotonic spread must bridge ---
functionally shrinking \texorpdfstring{\ensuremath{\lambda}}{lambda} until synaptic input fails to reach threshold.
\end{quote}

\begin{center}\rule{0.5\linewidth}{0.5pt}\end{center}

\subsubsection{Optogenetics: Light-Gated Ion Channels as Tools for
Mapping Neural
Circuits}\label{optogenetics-light-gated-ion-channels-as-tools-for-mapping-neural-circuits}

Until 2005, neuroscience had two experimental tools for probing neural
circuits: \textbf{electrical stimulation} (fast, but indiscriminate ---
excites every cell and axon within a about 100 \texorpdfstring{\ensuremath{\mu}}{mu}m radius) and
\textbf{pharmacological inactivation} (specific, but slow and diffuse).
\textbf{Optogenetics} --- introduced by Boyden, Deisseroth, and
colleagues (\emph{Nat. Neurosci.} 2005) --- merged genetic specificity
with millisecond temporal precision by transplanting light-gated ion
channels from microbes into mammalian neurons.

The core tools: \textbf{Channelrhodopsin-2 (ChR2)}, a blue-light-gated
non-selective cation channel from the green alga \emph{Chlamydomonas
reinhardtii}, opens in about 1 ms on 470 nm illumination and
drives neurons to fire action potentials with up to 40 Hz fidelity.
\textbf{Halorhodopsin (NpHR)} from \emph{Natronomonas pharaonis} is a
yellow-light-driven inward Cl\texorpdfstring{\textsuperscript{-}}{-} pump that hyperpolarises neurons and
silences firing. \textbf{Archaerhodopsin (Arch)} is a green-light-driven
outward H\texorpdfstring{\textsuperscript{+}}{+} pump with similar inhibitory effect. Transgene expression via
AAV with cell-type-specific \hyperref[gl:promoter]{\textbf{promoter}}s
(Thy1 for cortical pyramidal neurons, PV-Cre \texorpdfstring{\ensuremath{\times}}{x} ChR2 for parvalbumin
interneurons, DAT-Cre for dopamine neurons) delivers the channel
primarily to genetically defined populations. Combined with implanted
optical fibres (for deep structures) or surface-mounted LEDs (for
cortex), the experimenter can turn specific cell types on or off in
milliseconds.

Landmark findings enabled by optogenetics: \textbf{(1) Memory engrams}
--- Tonegawa and colleagues (\emph{Nature} 2012, \emph{Science} 2014)
showed that optogenetic reactivation of the specific hippocampal neurons
active during fear conditioning is sufficient to recall the fear memory
in a neutral context, proving engram cells are \emph{the} physical
substrate of memory. \textbf{(2) Parkinson's disease circuit dissection}
--- Kravitz et al.~(\emph{Nature} 2010) showed that direct-pathway (D1)
striatal neuron optogenetic activation rescues motor deficits in
parkinsonian mice, while indirect-pathway (D2) activation mimics them,
confirming the 1990s pharmacological model. \textbf{(3) Basis of deep
brain stimulation} --- optogenetic DBS mimics in rodents has clarified
which STN cell types produce the therapeutic benefit in Parkinson's
patients. \textbf{Clinical
\hyperref[gl:translation]{\textbf{translation}}}: RST-001 (jSAM for
retinitis pigmentosa, GenSight Biologics) uses ChrimsonR delivered by
AAV to retinal ganglion cells and specialised goggles to convert images
to patterned light. A 2021 \emph{Nature Medicine} report documented
partial restoration of visual function in a blind patient ---
optogenetics' first human therapeutic success.

\begin{center}\rule{0.5\linewidth}{0.5pt}\end{center}

\subsection{Summary}\label{summary-31}

\begin{itemize}
\tightlist
\item
  \textbf{NS organisation:} CNS (brain + spinal cord) + PNS (somatic +
  autonomic + enteric). 86 billion neurons; glia provide myelin, BBB,
  metabolic support, immune surveillance.
\item
  \textbf{Neuron types:} Sensory (afferent, pseudounipolar), motor
  (efferent, multipolar), interneurons (integration). Axon hillock has
  highest Nav channel density.
\item
  \textbf{Resting potential:} E\(_K\) = \(-89\) mV dominates;
  Na\(^+\)/K\(^+\)-ATPase maintains gradients; GHK equation weights by
  permeabilities; \(V_m \approx -70\) mV.
\item
  \textbf{Synaptic integration:} EPSPs (depolarising) and IPSPs
  (hyperpolarising/shunting) summate temporally and spatially at the
  axon hillock. Threshold at about \(-55\) mV triggers
  most-or-none AP.
\item
  \textbf{ANS:} Sympathetic (T1-L2, NE, fight-or-flight) vs
  parasympathetic (craniosacral, ACh, rest-and-digest). Enteric NS:
  semi-autonomous gut control.
\item
  \textbf{Brain:} Cortex (4 lobes, motor/sensory homunculi, language
  areas); basal ganglia (direct/indirect pathways, dopamine); thalamus
  (sensory relay); hippocampus (memory, place cells); hypothalamus
  (homeostasis); cerebellum (coordination); brainstem (vital functions).
\item
  \textbf{Sensory systems:} Mechanoreceptors (Meissner, Pacinian,
  Ruffini, Merkel), nociceptors (A-delta, C fibres), proprioceptors;
  visual pathway (retina to LGN to V1 to dorsal/ventral streams).
\item
  \textbf{Neural plasticity:} LTP (NMDA-Ca\(^{2+}\)-CaMKII-AMPA), LTD,
  homeostatic scaling, critical periods, synaptic pruning.
\item
  \textbf{Connections:} See \cref{sec:unit_IX_action_potential_synapses}
  for active propagation, \cref{sec:unit_II_cell_signaling} for second
  messengers, and \cref{sec:unit_0_active_inference} for active
  inference themes.
\end{itemize}

\begin{center}\rule{0.5\linewidth}{0.5pt}\end{center}

\subsubsection{Companion Source
Module}\label{companion-source-module-31}

\textbf{Nervous System and Neural Signalling} should leave a
reproducible trail from a biological claim to the code, figure, diagram,
or paper-based activity that can test it. Use the surfaces below to
inspect the chapter's assumptions, rerun the relevant model, or compare
the manuscript explanation with companion labs and figures.

{\def\LTcaptype{none} % do not increment counter
\begin{longtable}[]{@{}
  >{\raggedright\arraybackslash}p{(\linewidth - 2\tabcolsep) * \real{0.5000}}
  >{\raggedright\arraybackslash}p{(\linewidth - 2\tabcolsep) * \real{0.5000}}@{}}
\toprule\noalign{}
\begin{minipage}[b]{\linewidth}\raggedright
Surface
\end{minipage} & \begin{minipage}[b]{\linewidth}\raggedright
Use it for
\end{minipage} \\
\midrule\noalign{}
\endhead
\bottomrule\noalign{}
\endlastfoot
\breaktt{src/biology/neuroscience/neuroscience.py}
(\breaktt{cable\_voltage\_attenuation}, \breaktt{hebbian\_weight\_update},
\breaktt{action\_potential\_hh}) & Connect circuit architecture, passive
spread, spiking, and plasticity. \\
\breaktt{src/visualization/plots.py} (\breaktt{plot\_action\_potential}) &
Check timing and amplitude of neural signals. \\
\breaktt{src/mermaid/biology\_diagrams.py}
(\breaktt{nervous\_system\_reflex\_diagram}) & Keep stimulus,
integration, motor output, and feedback distinct. \\
\end{longtable}
}

\textbf{Reproducibility check:} specify cell type, circuit level,
recording method, and behavioural readout before linking neural
mechanism to outcome. \textbf{Cross-reference:} use
\cref{sec:unit_IX_action_potential_synapses} and
\cref{sec:unit_0_active_inference}.

\newpage

\section{Action Potentials and Synaptic
Transmission}\label{action-potentials-and-synaptic-transmission}

\label{sec:unit_IX_action_potential_synapses}

\begin{quote}
\textbf{Ch 30} \texorpdfstring{\ensuremath{\cdot}}{.} Level 3/3 \texorpdfstring{\ensuremath{\cdot}}{.} 55 min read \texorpdfstring{\ensuremath{\cdot}}{.} 100 min lecture \texorpdfstring{\ensuremath{\cdot}}{.}
Prerequisites: \cref{sec:unit_IX_nervous_system}
\end{quote}

\subsection{Learning Objectives}\label{learning-objectives-32}

\begin{enumerate}
\def\labelenumi{\arabic{enumi}.}
\tightlist
\item
  Write and explain the Hodgkin-Huxley equations for the
  \hyperref[gl:action-potential]{\textbf{action potential}}.
\item
  Explain the kinetics of Na\(^+\) activation (m), inactivation (h), and
  K\(^+\) activation (n) gating variables.
\item
  Describe the absolute and relative refractory periods in terms of
  channel states.
\item
  Explain myelination,
  \hyperref[gl:saltatory-conduction]{\textbf{saltatory conduction}}, and
  how axon diameter affects velocity.
\item
  Describe the molecular machinery of chemical synaptic transmission
  including SNARE \hyperref[gl:protein]{\textbf{protein}}s and
  Ca\(^{2+}\) triggering.
\item
  Compare major neurotransmitter systems and their receptor
  pharmacology.
\item
  Describe long-term potentiation \citep{frey1997} (LTP) and long-term
  depression (LTD) and their roles in learning and memory.
\item
  Explain the mechanisms of action of major neuroactive drugs.
\end{enumerate}

\subsubsection{Study Blueprint}\label{study-blueprint-32}

\begin{itemize}
\tightlist
\item
  \textbf{Big idea:} Electrical excitability and synaptic transmission
  convert ion gradients into rapid communication.
\item
  \textbf{Core concepts:} action potentials, ion channels, synapses,
  plasticity.
\item
  \textbf{Framework alignment:} Vision \& Change: Structure and
  function, Systems; AP Biology: Systems Interactions, Energetics;
  NGSS-style topics: Structure and Function.
\item
  \textbf{Model or quantitative lens:} Nernst/Goldman, Hodgkin-Huxley,
  and synaptic-current reasoning.
\item
  \textbf{Data skill:} Interpret voltage traces, conductance changes,
  and synaptic perturbations.
\item
  \textbf{Practice cadence:} Visual Representations, Statistical Tests
  and Data Analysis, Argumentation.
\item
  \textbf{Common misconception to repair:} An action potential is not
  electricity flowing like a wire; it is regenerated ion-channel
  dynamics.
\item
  \textbf{Primary lab:}
  \cref{sec:lab_unit_IX_action_potential_synapses}.
\item
  \textbf{Question bank:}
  \cref{sec:q_unit_IX_action_potential_synapses}.
\item
  \textbf{Transfer task:} Transfer excitability reasoning to anesthesia,
  toxins, epilepsy, and neuromuscular disease.
\item
  \textbf{Bridge to computation:}
  \breaktt{biology.neuroscience.neuroscience.action\_potential\_hh}.
\end{itemize}

\begin{center}\rule{0.5\linewidth}{0.5pt}\end{center}

\begin{quote}
\textbf{Opening Vignette --- The Squid That Explained the Brain}

The squid \emph{Loligo} possesses a giant axon up to 1 mm in diameter
--- large enough for researchers to insert a glass micropipette
electrode into the interior. In 1939, Alan Hodgkin and Andrew Huxley
exploited this remarkable anatomy to make the first intracellular
recordings of an action potential. By 1952, using voltage-clamp
apparatus of their own construction, they had measured the time courses
of sodium and potassium conductances across the membrane with stunning
precision and fitted them to a set of differential equations now known
as the Hodgkin-Huxley model. The equations predicted the shape,
amplitude, threshold, refractory period, and propagation velocity of
action potentials --- most from first principles of ionic movement. The
1963 Nobel Prize in Physiology or Medicine rewarded what remains one of
the most quantitatively rigorous pieces of biological science ever
accomplished. Every antiepileptic, local anaesthetic, cardiac
antiarrhythmic, and membrane biophysics textbook stems ultimately from
that one peculiarly large squid axon.
\end{quote}

\subsection{The Hodgkin-Huxley Model}\label{the-hodgkin-huxley-model}

\begin{figure}[htbp]
\centering
\includegraphics[width=0.85\textwidth]{../figures/nernst_potentials.png}
\caption{Nernst equilibrium potentials for major physiological ions (Na$^+$, K$^+$, Cl$^-$, Ca$^{2+}$) calculated from their inside/outside concentration gradients.}
\label{fig:unit_IX_nernst_potentials}
\end{figure}

Alan Hodgkin and Andrew Huxley recorded action potentials from the
\textbf{squid giant axon} (diameter about 0.5 mm) using voltage
clamp (Hodgkin \& Huxley, 1952, \emph{J Physiol} -- Nobel Prize 1963).
Each ionic current is driven by the difference between \(V\) and that
ion's Nernst equilibrium potential, the gradient-dependent values of
which are plotted in \cref{fig:unit_IX_nernst_potentials}. They
described membrane current as:

\begin{equation}
C_m \frac{dV}{dt} = -I_{Na} - I_K - I_L + I_{ext}
\label{eq:action_potential_synapses_1}
\end{equation}

where: - \(C_m\) = membrane capacitance about 1 uF/cm\(^2\) -
\(I_{Na} = \bar{g}_{Na} m^3 h (V - E_{Na})\); \(\bar{g}_{Na}\) = 120
mS/cm\(^2\); \(E_{Na}\) = +50 mV - \(I_K = \bar{g}_K n^4 (V - E_K)\);
\(\bar{g}_K\) = 36 mS/cm\(^2\); \(E_K\) = \(-77\) mV -
\(I_L = g_L (V - E_L)\); \(g_L\) = 0.3 mS/cm\(^2\); \(E_L\) = \(-54.4\)
mV

\subsubsection{Gating Variables}\label{gating-variables}

Each gating variable (\(m\), \(h\), \(n\)) evolves according to
first-order Markov kinetics:

\begin{equation}
\frac{dx}{dt} = \alpha_x(V)(1-x) - \beta_x(V)x
\label{eq:action_potential_synapses_2}
\end{equation}

At steady state: \(x_\infty(V) = \alpha_x / (\alpha_x + \beta_x)\); time
constant: \(\tau_x(V) = 1/(\alpha_x + \beta_x)\)

\begin{figure}[htbp]
\centering
\includegraphics[width=0.82\linewidth,height=4.2in,keepaspectratio]{../figures/mermaid_inline/inline_mermaid_0160_6aa448c0dd72.png}
\caption{Voltage-gated Na\$\textasciicircum{}+\$ channel states. The channel transitions through three states: Closed (resting, ready to open), Open (conducting Na\$\textasciicircum{}+\$, brief about 0.5 ms), and Inactivated (blocked by the inactivation gate, cannot reopen until repolarisation restores the resting state). The inactivated state underlies the absolute refractory period.}
\end{figure}

{\def\LTcaptype{none} % do not increment counter
\begin{longtable}[]{@{}
  >{\raggedright\arraybackslash}p{(\linewidth - 4\tabcolsep) * \real{0.4167}}
  >{\raggedright\arraybackslash}p{(\linewidth - 4\tabcolsep) * \real{0.3056}}
  >{\raggedright\arraybackslash}p{(\linewidth - 4\tabcolsep) * \real{0.2778}}@{}}
\toprule\noalign{}
\begin{minipage}[b]{\linewidth}\raggedright
Gating variable
\end{minipage} & \begin{minipage}[b]{\linewidth}\raggedright
\(\alpha(V)\)
\end{minipage} & \begin{minipage}[b]{\linewidth}\raggedright
\(\beta(V)\)
\end{minipage} \\
\midrule\noalign{}
\endhead
\bottomrule\noalign{}
\endlastfoot
m (Na\(^+\) activation) & \(0.1(V+40) / [1 - \exp(-(V+40)/10)]\) &
\(4 \exp(-(V+65)/18)\) \\
h (Na\(^+\) inactivation) & \(0.07 \exp(-(V+65)/20)\) &
\(1 / [1 + \exp(-(V+35)/10)]\) \\
n (K\(^+\) activation) & \(0.01(V+55) / [1 - \exp(-(V+55)/10)]\) &
\(0.125 \exp(-(V+65)/80)\) \\
\end{longtable}
}

\subsubsection{Biophysical Interpretation of the Action
Potential}\label{biophysical-interpretation-of-the-action-potential}

\begin{itemize}
\tightlist
\item
  \textbf{Threshold} (about \(-55\) mV): Point where inward
  Na\(^+\) current exceeds outward K\(^+\) leak. Positive feedback
  begins (more \hyperref[gl:depolarisation]{\textbf{depolarisation}}
  opens more Na\(^+\) channels).
\item
  \textbf{Upstroke (depolarisation):} Rapid m-gate opening
  (\(\tau_m \approx 0.5\) ms). Na\(^+\) rushes in. Membrane potential
  approaches E\(_{Na}\) (+50 mV) but typically peaks at +30 to +40 mV.
\item
  \textbf{Repolarisation:} h-gate closes (Na\(^+\) inactivation,
  \(\tau_h \approx 5\) ms) AND n-gate opens (K\(^+\) delayed rectifier,
  \(\tau_n \approx 5\) ms). Na\(^+\) influx stops; K\(^+\) efflux drives
  membrane back toward E\(_K\).
\item
  \textbf{Afterhyperpolarisation (undershoot):} n-gate slow to close.
  K\(^+\) efflux continues briefly, hyperpolarising membrane below
  \hyperref[gl:resting-potential]{\textbf{resting potential}}
  (about \(-80\) mV). This is the relative refractory period.
\end{itemize}

\subsubsection{Channel Pharmacology}\label{channel-pharmacology}

\textbf{Tetrodotoxin (TTX):} Puffer fish toxin. Blocks Nav channels from
the extracellular side by occluding the selectivity filter pore. Blocks
most AP generation. LD\(_{50}\) in mice: about 10 ug/kg. Used
extensively as a research tool to study Na\(^+\)-dependent processes.

\textbf{4-Aminopyridine (4-AP):} Blocks Kv (voltage-gated K\(^+\))
channels. Prolongs AP duration, enhances neurotransmitter release. Used
clinically as dalfampridine (Ampyra) for multiple sclerosis to improve
conduction in demyelinated axons.

\textbf{Nav channel subtypes:} Nav1.1-Nav1.9 encoded by SCN1A-SCN11A
\hyperref[gl:gene]{\textbf{gene}}s. Different subtypes have distinct
tissue distributions, voltage sensitivities, and pharmacological
profiles: - Nav1.7 (SCN9A): Pain perception. Loss-of-function:
congenital insensitivity to pain. Gain-of-function: erythromelalgia
(burning pain). - Nav1.5 (SCN5A): Cardiac muscle.
\hyperref[gl:mutation]{\textbf{Mutation}}s cause Long QT syndrome,
Brugada syndrome.

\textbf{Concept Check:} TTX blocks Na\(^+\) channels and abolishes
action potentials. Local anaesthetics (lidocaine) also block Na\(^+\)
channels. Why does lidocaine preferentially block pain fibres rather
than motor fibres? (Hint: consider use-dependence and fibre diameter.)

\begin{center}\rule{0.5\linewidth}{0.5pt}\end{center}

\subsection{Refractory Periods}\label{refractory-periods}

\textbf{Absolute refractory period} (about 1-2 ms after spike
peak): Na\(^+\) channels are in the inactivated state (h about 0). They
cannot be re-opened regardless of stimulus strength. This supports: -
Unidirectional AP propagation (the region just behind the AP cannot be
re-excited) - Maximum firing rate (about 500-1000 Hz)

\textbf{Relative refractory period} (about 5-10 ms following
absolute): Na\(^+\) channels are recovering from inactivation (h
recovering), but K\(^+\) channels remain open. The membrane is more
negative than resting potential (afterhyperpolarisation). A
stronger-than-normal stimulus is needed to trigger a second spike. APs
generated during this period have reduced amplitude.

\begin{quote}
\textbf{Clinical Connection:} The cardiac refractory period is much
longer than in \hyperref[gl:neuron]{\textbf{neuron}}s
(about 200-300 ms vs about 2 ms) because cardiac AP
duration is prolonged by L-type Ca\(^{2+}\) channels (plateau phase).
This prevents tetanic contraction of the heart (which would be fatal --
the heart must relax to fill). The ``vulnerable period'' near the end of
the T wave (relative refractory period of ventricles) is when a
premature stimulus can trigger ventricular fibrillation.
\end{quote}

\textbf{Concept Check:} During the falling phase of the action potential
the m, h, and n gates are most changing at once. Which single
Hodgkin-Huxley gating variable is the \emph{direct} cause of
repolarisation, and why does Na\(^+\) inactivation (h) alone fail to
repolarise the membrane without it? (Hint: compare what each conductance
does to the driving force toward \(E_K\) versus \(E_{Na}\).)

\begin{center}\rule{0.5\linewidth}{0.5pt}\end{center}

\subsection{Myelination and Conduction
Velocity}\label{myelination-and-conduction-velocity}

\subsubsection{Unmyelinated Conduction}\label{unmyelinated-conduction}

In continuous conduction, the AP generates local circuits along the axon
membrane:

\begin{equation}
\theta \propto \sqrt{a/r_i}
\label{eq:action_potential_synapses_3}
\end{equation}

where a = axon radius and \(r_i\) = axial resistance per unit length.
Increasing diameter reduces axial resistance, increasing velocity. Squid
giant axon (500 um diameter) achieves about 25 m/s by
brute-force gigantism.

\subsubsection{Myelinated Axons -- Saltatory
Conduction}\label{myelinated-axons-saltatory-conduction}

Myelin dramatically increases the \textbf{space constant} (\texorpdfstring{\ensuremath{\lambda}}{lambda}):

\begin{equation}
\lambda = \sqrt{r_m / r_i}
\label{eq:action_potential_synapses_4}
\end{equation}

Myelin increases \(r_m\) (membrane resistance, preventing current leak)
and reduces \(C_m\) (membrane capacitance, reducing charge needed to
change voltage). The result: depolarisation spreads much further without
decay.

APs can primarily arise at \textbf{nodes of Ranvier} (1-2 um gaps
between myelin segments, every about 1 mm). Nodes have
concentrated voltage-gated Na\(^+\) channels
(about 1,000/um\(^2\) at nodes vs about 25/um\(^2\)
under myelin).

\textbf{Saltatory conduction:} The AP appears to ``jump'' between nodes.
In myelinated axons, velocity scales linearly with diameter
(\(\theta \propto d\)), rather than as the square root:

{\def\LTcaptype{none} % do not increment counter
\begin{longtable}[]{@{}lllll@{}}
\toprule\noalign{}
Fibre type & Diameter (um) & Velocity (m/s) & Myelinated? & Function \\
\midrule\noalign{}
\endhead
\bottomrule\noalign{}
\endlastfoot
A-alpha (Ia) & 12-20 & 70-120 & Yes & Muscle spindle afferents \\
A-beta & 6-12 & 30-70 & Yes & Touch, pressure \\
A-delta & 1-5 & 5-30 & Thinly & Sharp pain, temperature \\
B & 1-3 & 3-15 & Yes & Preganglionic autonomic \\
C & 0.2-1.5 & 0.5-2 & No & Slow pain, itch, postganglionic autonomic \\
\end{longtable}
}

\subsubsection{Worked Example: Myelination Gain --- Saltatory vs
Continuous
Conduction}\label{worked-example-myelination-gain-saltatory-vs-continuous-conduction}

Suppose two unmyelinated axons of equal small diameter (\(d = 1\;\mu\)m)
carry the same action potential. Continuous conduction velocity scales
as \(\theta_\text{unmyel} \propto \sqrt{d}\). From empirical data, an
unmyelinated 1 \texorpdfstring{\ensuremath{\mu}}{mu}m axon conducts at about 0.5 m/s.

Now myelinate one axon with internodes about 1 mm apart and 100
wraps of myelin. Myelin increases membrane resistance (\(r_m\))
about 100-fold and reduces capacitance (\(C_m\))
about 50-fold:

\begin{itemize}
\tightlist
\item
  Length constant \(\lambda = \sqrt{r_m / r_i}\) increases
  about 10-fold
\item
  Time constant at the node \(\tau_m = r_m C_m\) --- although both
  increase, the \emph{time to reach threshold} at the next node depends
  primarily on the rapid charging of node membrane (which has low
  capacitance) by the spreading depolarisation
\item
  Effective conduction velocity in myelinated fibre:
  \(\theta_\text{myel} \propto d\) (linear, not square root)
\end{itemize}

For the same 1 \texorpdfstring{\ensuremath{\mu}}{mu}m fibre with myelination, conduction approaches
about 6 m/s (\textgreater10\texorpdfstring{\ensuremath{\times}}{x} faster). For larger myelinated
fibres (10--20 \texorpdfstring{\ensuremath{\mu}}{mu}m A-\texorpdfstring{\ensuremath{\alpha}}{alpha}), the gain is even more striking:

{\def\LTcaptype{none} % do not increment counter
\begin{longtable}[]{@{}lllll@{}}
\toprule\noalign{}
Fibre & Diameter & Conduction velocity & Myelinated? & Velocity gain \\
\midrule\noalign{}
\endhead
\bottomrule\noalign{}
\endlastfoot
C fibre & 1 \texorpdfstring{\ensuremath{\mu}}{mu}m & 0.5 m/s & No & (baseline) \\
A-\texorpdfstring{\ensuremath{\delta}}{delta} & 1--5 \texorpdfstring{\ensuremath{\mu}}{mu}m & 5--30 m/s & Thinly & about 10--60\texorpdfstring{\ensuremath{\times}}{x} \\
A-\texorpdfstring{\ensuremath{\alpha}}{alpha} (Ia) & 12--20 \texorpdfstring{\ensuremath{\mu}}{mu}m & 70--120 m/s & Yes (heavily) &
about 150--250\texorpdfstring{\ensuremath{\times}}{x} \\
\end{longtable}
}

\textbf{The fundamental insight:} Saltatory conduction provides the
\textbf{same speed at about 50\texorpdfstring{\ensuremath{\times}}{x} lower diameter} compared to
brute-force gigantism (squid giant axon, 500 \texorpdfstring{\ensuremath{\mu}}{mu}m). The vertebrate
solution to fast conduction --- myelinate small axons rather than grow
giant ones --- is enormously more space-efficient and energetically
cheaper. The about \(10^{12}\) axons in the human white
matter would not fit in the cranium without myelin.

The \textbf{energetic cost} of action potentials is also dramatically
reduced: in unmyelinated axons, Na\(^+\)/K\(^+\)-ATPase must
continuously pump out Na\(^+\) along the entire length. In myelinated
axons, AP-associated ion fluxes occur primarily at nodes
(about 0.1\% of axon surface), reducing the metabolic cost of
spike propagation by a similar factor.

\subsubsection{Multiple Sclerosis}\label{multiple-sclerosis}

\textbf{Multiple sclerosis (MS):} Autoimmune demyelination of CNS white
matter. T cells recognise myelin basic protein (MBP) and proteolipid
protein (PLP) as autoantigens. Inflammatory attack destroys
oligodendrocyte myelin sheaths.

Consequences: Slowed or blocked conduction in demyelinated segments.
Symptoms depend on lesion location: optic neuritis (visual loss),
weakness, sensory loss, fatigue, cognitive impairment.
Relapsing-remitting course in about 85\% of patients.

Treatment: Disease-modifying therapies (interferons, natalizumab,
ocrelizumab) reduce relapse frequency. 4-aminopyridine (blocks exposed
K\(^+\) channels in demyelinated segments) improves conduction.

\textbf{Concept Check:} A 1 \texorpdfstring{\ensuremath{\mu}}{mu}m unmyelinated axon conducts at
about 0.5 m/s, yet a 1 \texorpdfstring{\ensuremath{\mu}}{mu}m myelinated axon conducts at
about 6 m/s while a 20 \texorpdfstring{\ensuremath{\mu}}{mu}m unmyelinated axon would need a far
larger diameter to match it. Using the velocity scalings
\(\theta \propto \sqrt{d}\) (unmyelinated) versus \(\theta \propto d\)
(myelinated), explain why myelination --- not gigantism --- is the
space- and energy-efficient solution, and predict what happens to
conduction when a node of Ranvier is demyelinated. (Hint: think about
what myelin does to \(r_m\), \(C_m\), and the length constant \texorpdfstring{\ensuremath{\lambda}}{lambda}.)

\begin{center}\rule{0.5\linewidth}{0.5pt}\end{center}

\subsection{Chemical Synaptic
Transmission}\label{chemical-synaptic-transmission}

\subsubsection{Synapse Types}\label{synapse-types}

\textbf{Electrical synapses} (gap junctions): Connexin/innexin
hemichannels directly connect
\hyperref[gl:cytoplasm]{\textbf{cytoplasm}}s. Bidirectional, nearly
instantaneous (\textless0.1 ms). Found in: cardiac pacemaker cells
(connexin 43), retinal processing (rod-cone coupling), hippocampal
interneuron synchronisation. Connexin 26/30 mutations cause hereditary
deafness.

\textbf{Chemical synapses:} Unidirectional. Modifiable (plasticity).
Synaptic delay: 0.5-2 ms. The \hyperref[gl:dominant]{\textbf{dominant}}
synapse type in the CNS.

\subsubsection{Presynaptic Machinery}\label{presynaptic-machinery}

\begin{figure}[htbp]
\centering
\includegraphics[width=0.82\linewidth,height=4.2in,keepaspectratio]{../figures/mermaid_inline/inline_mermaid_0161_670c2dedea6d.png}
\caption{Molecular machinery of synaptic vesicle fusion. Depolarisation opens Ca\$\textasciicircum{}\{2+\}\$ channels. Ca\$\textasciicircum{}\{2+\}\$ binds synaptotagmin-1, which displaces the complexin fusion clamp and triggers full SNARE complex zippering, forcing the vesicle membrane to fuse with the plasma membrane and release neurotransmitter into the cleft.}
\end{figure}

\textbf{Key proteins:} - \textbf{SNARE complex:} The minimal fusion
machinery. \textbf{Synaptobrevin} (v-SNARE, on vesicle membrane) +
\textbf{syntaxin} + \textbf{SNAP-25} (both t-SNAREs, on plasma
membrane). The four \texorpdfstring{\ensuremath{\alpha}}{alpha}-helices zipper together N-to-C terminally, pulling
membranes into contact. - \textbf{Synaptotagmin-1:} Ca\(^{2+}\) sensor.
Its C2A and C2B domains bind 4-5 Ca\(^{2+}\) ions, triggering membrane
insertion and SNARE complex activation. Cooperativity of Ca\(^{2+}\)
binding gives the steep Ca\(^{2+}\)-release relationship
(about 4th power). - \textbf{Munc18 / Munc13:} Priming factors
that prepare SNARE complexes for fast fusion - \textbf{RIM:} Active zone
scaffold protein. Positions vesicles within nanometres of Ca\(^{2+}\)
channels. - \textbf{Complexin:} ``Fusion clamp.'' Binds partially
zippered SNARE complex and prevents spontaneous fusion.
Ca\(^{2+}\)/synaptotagmin displaces complexin to trigger release. -
\textbf{NSF + alpha-SNAP:} After fusion, these AAA\(^+\) ATPases
disassemble SNARE complexes for recycling.

\textbf{Quantal release (Katz, Nobel 1970):} Neurotransmitter is
released in discrete packets (quanta), each corresponding to the
contents of one vesicle. The end-plate potential (EPP) at the NMJ = n
\(\times\) p \(\times\) q, where n = number of release-ready vesicles, p
= release probability per vesicle, q = quantal size (postsynaptic
response to one vesicle).

Miniature EPPs (mEPPs) represent spontaneous release of single vesicles
and have a constant amplitude (about 0.5 mV at the NMJ). The
evoked EPP is an integer multiple of the mEPP amplitude.

\subsubsection{Synaptic Vesicle Pools}\label{synaptic-vesicle-pools}

A presynaptic terminal does \textbf{not} treat most its vesicles
equivalently. Decades of imaging, electrophysiology, and FM dye
experiments establish three functionally distinct pools that differ in
their release-readiness, location, and mobilisation kinetics:

{\def\LTcaptype{none} % do not increment counter
\begin{longtable}[]{@{}
  >{\raggedright\arraybackslash}p{(\linewidth - 8\tabcolsep) * \real{0.0548}}
  >{\raggedright\arraybackslash}p{(\linewidth - 8\tabcolsep) * \real{0.4384}}
  >{\raggedright\arraybackslash}p{(\linewidth - 8\tabcolsep) * \real{0.1096}}
  >{\raggedright\arraybackslash}p{(\linewidth - 8\tabcolsep) * \real{0.2192}}
  >{\raggedright\arraybackslash}p{(\linewidth - 8\tabcolsep) * \real{0.1781}}@{}}
\toprule\noalign{}
\begin{minipage}[b]{\linewidth}\raggedright
Pool
\end{minipage} & \begin{minipage}[b]{\linewidth}\raggedright
Approximate size per active zone
\end{minipage} & \begin{minipage}[b]{\linewidth}\raggedright
Location
\end{minipage} & \begin{minipage}[b]{\linewidth}\raggedright
Release kinetics
\end{minipage} & \begin{minipage}[b]{\linewidth}\raggedright
Replenishment
\end{minipage} \\
\midrule\noalign{}
\endhead
\bottomrule\noalign{}
\endlastfoot
\textbf{Readily releasable pool (RRP)} & about 10--20 vesicles
(CNS); about 50--100 (NMJ) & Docked at active zone; primed &
Released within about 1 ms of Ca\(^{2+}\) entry; depleted
within first few APs of a high-frequency train & Refilled from recycling
pool over 1--10 s \\
\textbf{Recycling pool} & about 100--200 vesicles & Near active
zone; fused with PM during sustained activity then reformed locally &
Sustains release during moderate-frequency activity & Cycle time
about 30--60 s \\
\textbf{Reserve pool} & about 300--500+ vesicles & Tethered to
actin/synapsin in the cytoplasm & Mobilised primarily during intense or
prolonged activity (typically when synapsin is phosphorylated by
CaMKII/PKA) & Slow refilling (minutes) from synthesis and endocytosis \\
\end{longtable}
}

The three-pool model elegantly explains the kinetics of synaptic
transmission during different activity patterns. \textbf{Brief,
low-frequency stimulation} uses primarily the RRP and the system rapidly
returns to baseline. \textbf{Sustained moderate activity} engages the
recycling pool. \textbf{Tetanic stimulation} mobilises the reserve pool
by phosphorylation of synapsin (which normally clamps reserve vesicles
to the actin cytoskeleton), releasing them to refill the active zone.

\subsubsection{Short-term plasticity --- facilitation and
depression}\label{short-term-plasticity-facilitation-and-depression}

Synaptic strength is dynamic across milliseconds-to-seconds, even before
any LTP/LTD is induced. Two opposing short-term phenomena dominate:

\textbf{Synaptic facilitation} (paired-pulse facilitation; tens to
hundreds of ms): If two APs arrive in close succession, the second
EPSC/IPSC is \emph{larger} than the first. Mechanism: residual
presynaptic Ca\(^{2+}\) from the first AP has not yet been pumped out or
buffered; when the second AP triggers Ca\(^{2+}\) entry, the residual
Ca\(^{2+}\) adds to the new flux, super-linearly increasing release
probability \(p\). Facilitation is most prominent at synapses with low
initial \(p\) (e.g., parallel-fibre to Purkinje cell synapses;
about \(p \approx 0.05\)); high-\(p\) synapses have less room
for further increase.

\textbf{Synaptic depression} (paired-pulse depression; tens of ms to
seconds): If two APs arrive in close succession at a high-\(p\) synapse,
the second EPSC is \emph{smaller} than the first. Mechanism: the RRP is
partially depleted by the first AP and has not had time to refill.
Depression is most prominent at synapses with high initial \(p\) (e.g.,
calyx of Held auditory synapses; \(p \approx 0.4\)).

Whether a given synapse facilitates or depresses depends on its starting
release probability, vesicle pool dynamics, and Ca\(^{2+}\) handling.
Some synapses transition: facilitate at low frequency, then depress as
frequency rises and pools deplete.

\subsubsection{Tsodyks-Markram model}\label{tsodyks-markram-model}

The classical phenomenological model of short-term plasticity (Tsodyks
\& Markram, 1997) treats each synapse as having two state variables:

\begin{itemize}
\tightlist
\item
  \textbf{\(x\)}: the fraction of available (un-depleted) resources
  (vesicles in the RRP), with \(0 \le x \le 1\)
\item
  \textbf{\(u\)}: the utilisation parameter (effective release
  probability), influenced by residual Ca\(^{2+}\)
\end{itemize}

After each spike, \(x \to x \cdot (1 - u)\) (resources depleted) and
\(u \to u + U \cdot (1 - u)\) (release probability incremented by
Ca\(^{2+}\) accumulation). Between spikes, both variables relax back to
baseline with time constants \(\tau_d\) (recovery, about 500
ms--2 s) and \(\tau_f\) (facilitation, about 100--500 ms):

\[\frac{dx}{dt} = \frac{1 - x}{\tau_d}; \qquad \frac{du}{dt} = -\frac{u - U}{\tau_f} \label{eq:unit_IX_action_potential_synapses_item_1}\]

The post-spike EPSC amplitude is proportional to \(u \cdot x\). This
single model captures the full repertoire of facilitating, depressing,
and mixed synapses --- by varying the parameters \(U\), \(\tau_d\),
\(\tau_f\) --- and is the standard building block for biologically
realistic spiking neural network simulations.

\begin{quote}
\textbf{Clinical Connection:} Lambert-Eaton myasthenic syndrome (LEMS)
is an autoimmune disorder targeting presynaptic P/Q-type Ca\(^{2+}\)
channels at the NMJ, dramatically reducing initial release probability.
The hallmark on bedside testing is \textbf{post-tetanic facilitation}
--- strength briefly improves after sustained voluntary contraction,
reflecting accumulated Ca\(^{2+}\) partially compensating for reduced
channel number. Treatment with \textbf{3,4-diaminopyridine} (a K\(^+\)
channel blocker that prolongs presynaptic AP duration \texorpdfstring{\ensuremath{\to}}{->} more Ca\(^{2+}\)
entry per spike) increases initial \(p\) and restores neuromuscular
transmission.
\end{quote}

\subsubsection{Postsynaptic Receptors}\label{postsynaptic-receptors}

{\def\LTcaptype{none} % do not increment counter
\begin{longtable}[]{@{}
  >{\raggedright\arraybackslash}p{(\linewidth - 8\tabcolsep) * \real{0.2800}}
  >{\raggedright\arraybackslash}p{(\linewidth - 8\tabcolsep) * \real{0.1400}}
  >{\raggedright\arraybackslash}p{(\linewidth - 8\tabcolsep) * \real{0.3000}}
  >{\raggedright\arraybackslash}p{(\linewidth - 8\tabcolsep) * \real{0.1000}}
  >{\raggedright\arraybackslash}p{(\linewidth - 8\tabcolsep) * \real{0.1800}}@{}}
\toprule\noalign{}
\begin{minipage}[b]{\linewidth}\raggedright
Receptor class
\end{minipage} & \begin{minipage}[b]{\linewidth}\raggedright
Example
\end{minipage} & \begin{minipage}[b]{\linewidth}\raggedright
Ion selectivity
\end{minipage} & \begin{minipage}[b]{\linewidth}\raggedright
Speed
\end{minipage} & \begin{minipage}[b]{\linewidth}\raggedright
Mechanism
\end{minipage} \\
\midrule\noalign{}
\endhead
\bottomrule\noalign{}
\endlastfoot
\textbf{Ionotropic} & AMPA (GluA1-4) & Na\(^+\)/K\(^+\) (some
Ca\(^{2+}\)) & Fast (ms) & Ligand-gated ion channel \\
Ionotropic & NMDA (GluN1/2A-D) & Na\(^+\)/K\(^+\)/Ca\(^{2+}\) & Slower
(10s of ms) & Voltage + ligand gated; Mg\(^{2+}\) block \\
Ionotropic & GABA\(_A\) & Cl\(^-\) & Fast (ms) & Ligand-gated anion
channel \\
Ionotropic & nAChR & Na\(^+\)/K\(^+\) & Fast (ms) & Pentameric cation
channel \\
\textbf{Metabotropic} & mGluR1-8 & N/A & Slow (100s ms-min) & GPCR to
G\(_q\)/G\(_i\) to IP\(_3\)/cAMP \\
Metabotropic & GABA\(_B\) & K\(^+\) (indirect) & Slow & G\(_i\) to
K\(^+\) channel activation \\
Metabotropic & D1/D2 & N/A & Slow & G\(_s\)/G\(_i\) to cAMP \\
Metabotropic & 5-HT\(_{1A-7}\) & N/A & Slow & Various G protein
coupling \\
\end{longtable}
}

\textbf{NMDA receptor -- ``coincidence detector'':} The NMDA receptor
requires BOTH glutamate binding AND postsynaptic depolarisation (to
relieve Mg\(^{2+}\) block) to conduct. This makes it a detector of
coincident pre- and postsynaptic activity -- exactly the condition Hebb
proposed should strengthen synapses. Ca\(^{2+}\) entry through NMDA
receptors triggers LTP.

\textbf{Concept Check:} At a synapse with \(n = 10\) release-ready
vesicles, release probability \(p = 0.2\), and quantal size \(q = 0.5\)
mV, a single presynaptic spike must drive the postsynaptic membrane from
\(-70\) mV to a \(-55\) mV threshold. Using the \(n \cdot p \cdot q\)
model, decide whether one spike reaches threshold and explain why a drug
that raises \(p\) (e.g., an aminopyridine prolonging the presynaptic AP)
can convert a sub-threshold synapse to a supra-threshold one without
adding any new vesicles. (Hint: compute the expected EPSP amplitude and
compare it to the 15 mV gap.)

\begin{center}\rule{0.5\linewidth}{0.5pt}\end{center}

\subsection{Neurotransmitter Systems}\label{neurotransmitter-systems}

\subsubsection{Glutamate (Excitatory)}\label{glutamate-excitatory}

The major excitatory neurotransmitter in the CNS (about 90\% of
excitatory synapses).

\textbf{Receptors:} - \textbf{AMPA} (GluA1-4): Fast Na\(^+\)/K\(^+\)
current. Ca\(^{2+}\)-permeable if lacking GluA2 (GluA2 RNA editing of
Q/R site prevents Ca\(^{2+}\) permeability in most adult neurons). -
\textbf{NMDA} (GluN1/GluN2A-D): Coincidence detector. Requires glutamate
+ glycine/D-serine co-agonist + depolarisation to relieve Mg\(^{2+}\)
block. High Ca\(^{2+}\) permeability. Critical for LTP. -
\textbf{Kainate} (GluK1-5): Similar to AMPA; modulatory roles at some
synapses. - \textbf{mGluR} (mGluR1-8): Metabotropic. Group I (mGluR1/5):
G\(_q\)-coupled, excitatory. Group II/III: G\(_i\)-coupled, inhibitory
(presynaptic autoreceptors).

\textbf{\hyperref[gl:excitotoxicity]{Excitotoxicity}:} Excessive
glutamate release (e.g., during stroke/ischaemia) overactivates NMDA
receptors. Massive Ca\(^{2+}\) influx activates calpains, endonucleases,
and nitric oxide synthase, leading to neuronal death. This is a major
mechanism of ischaemic brain damage.

\subsubsection{GABA (Inhibitory)}\label{gaba-inhibitory}

The major inhibitory neurotransmitter (about 30\% of CNS
synapses). Synthesised from glutamate by glutamic acid decarboxylase
(GAD65/GAD67).

\textbf{Receptors:} - \textbf{GABA\(_A\)} (ionotropic): Cl\(^-\)
channel. Pentameric (typically \(\alpha_1\beta_2\gamma_2\)). Binding
sites for: GABA (agonist), benzodiazepines (positive
\hyperref[gl:allosteric]{\textbf{allosteric}} modulators at \texorpdfstring{\ensuremath{\alpha}}{alpha}/\texorpdfstring{\ensuremath{\gamma}}{gamma}
interface), barbiturates, neurosteroids, ethanol, anaesthetics
(propofol, isoflurane). - \textbf{GABA\(_B\)} (metabotropic):
G\(_i\)-coupled GPCR. Activates K\(^+\) channels (postsynaptic
hyperpolarisation) and inhibits Ca\(^{2+}\) channels (presynaptic,
reduces NT release). Baclofen is a GABA\(_B\) agonist (used for
spasticity).

\subsubsection{Acetylcholine (ACh)}\label{acetylcholine-ach}

\textbf{Receptors:} - \textbf{Nicotinic (nAChR):} Ionotropic pentamer.
Muscle type (\(\alpha_1^2\beta_1\delta\epsilon\)) at NMJ. Neuronal types
(\(\alpha_4\beta_2\), \(\alpha_7\)) in CNS. Na\(^+\)/K\(^+\) channels. -
\textbf{Muscarinic (mAChR, M1-M5):} GPCRs. M1/M3/M5: G\(_q\)
(excitatory). M2/M4: G\(_i\) (inhibitory, e.g., vagal slowing of heart
rate).

\textbf{CNS cholinergic system:} Basal forebrain (nucleus basalis of
Meynert) projects to cortex. Critical for attention and memory.
Degeneration in Alzheimer's disease.

\subsubsection{Dopamine}\label{dopamine}

\textbf{Receptors:} D1-D5 (most GPCRs). D1-like (D1, D5): G\(_s\),
increase cAMP. D2-like (D2, D3, D4): G\(_i\), decrease cAMP.

\textbf{Major pathways:} - \textbf{Nigrostriatal} (SNc to striatum):
Motor control. Loss causes Parkinson's disease. - \textbf{Mesolimbic}
(VTA to nucleus accumbens): Reward prediction error, motivation.
Hyperactivity linked to positive symptoms of schizophrenia. -
\textbf{Mesocortical} (VTA to prefrontal cortex): Working memory,
executive function. Hypoactivity linked to negative symptoms of
schizophrenia and ADHD. - \textbf{Tuberoinfundibular} (hypothalamus to
pituitary): Inhibits prolactin release.

\subsubsection{Serotonin (5-HT)}\label{serotonin-5-ht}

Synthesised from tryptophan by tryptophan hydroxylase (TPH2 in CNS).
Raphe nuclei in brainstem project widely. 14 receptor subtypes
(5-HT\(_1\)A through 5-HT\(_7\)).

\textbf{Functions:} Mood regulation, sleep/wake cycle, appetite, pain
modulation. about 90\% of body's serotonin is in
enterochromaffin cells of the GI tract.

\subsubsection{Norepinephrine (NE)}\label{norepinephrine-ne}

Synthesised from dopamine by dopamine \texorpdfstring{\ensuremath{\beta}}{beta}-hydroxylase (DBH). Locus
coeruleus (brainstem) projects to entire cortex. Functions: arousal,
attention, vigilance, stress response.

\textbf{Receptors:} \(\alpha_1\) (G\(_q\)), \(\alpha_2\) (G\(_i\),
presynaptic autoreceptor), \(\beta_1\) (G\(_s\), heart), \(\beta_2\)
(G\(_s\), bronchial smooth muscle, blood vessels).

\begin{center}\rule{0.5\linewidth}{0.5pt}\end{center}

\subsection{AI and Computational Neuroscience: Modeling the Brain with
Data and
Algorithms}\label{ai-and-computational-neuroscience-modeling-the-brain-with-data-and-algorithms}

\begin{figure}[htbp]
\centering
\includegraphics[width=0.82\linewidth,height=4.2in,keepaspectratio]{../figures/mermaid_inline/inline_mermaid_0162_36db1762ef16.png}
\caption{The convergence of computational neuroscience and AI. Early neural models (Hodgkin-Huxley, integrate-and-fire) inspired artificial neurons, while modern deep learning architectures are now used to model brain function and analyze neural data.}
\end{figure}

\subsubsection{Computational Neuroscience
Models}\label{computational-neuroscience-models}

\textbf{Biophysical models}:

\begin{itemize}
\item
  \textbf{Hodgkin-Huxley model} (1952): A set of nonlinear differential
  equations that describe how action potentials in neurons are initiated
  and propagated. It remains one of the most important mathematical
  models in neuroscience.
\item
  \textbf{FitzHugh-Nagumo model} (1961): A simplified two-variable model
  that captures the essential dynamics of the Hodgkin-Huxley model while
  being more mathematically tractable.
\item
  \textbf{Izhikevich model} (2003): A hybrid model that combines the
  biological plausibility of Hodgkin-Huxley with the computational
  efficiency of integrate-and-fire models.
\end{itemize}

\textbf{Large-scale brain models}:

\begin{itemize}
\tightlist
\item
  \textbf{Blue Brain Project} (EPFL): Aims to create a digital
  reconstruction of the rodent brain and ultimately the human brain
  using detailed biophysical models.
\item
  \textbf{Human Brain Project} (EU): A €1 billion initiative to simulate
  the entire human brain using supercomputers.
\item
  \textbf{Spaun} (University of Waterloo): A large-scale functional
  brain model that can perform multiple cognitive tasks.
\end{itemize}

\subsubsection{Machine Learning for Neural Data
Analysis}\label{machine-learning-for-neural-data-analysis}

Modern neuroscience generates massive datasets (e.g., calcium imaging,
EEG, fMRI, intracranial recordings). Machine learning is essential for
extracting meaningful information:

\begin{itemize}
\item
  \textbf{Spike sorting}: Unsupervised learning algorithms (e.g., PCA +
  clustering, deep learning) to separate the activity of individual
  neurons from extracellular recordings.
\item
  \textbf{Calcium imaging analysis}: Deconvolution algorithms (e.g.,
  OASIS, CaImAn) to extract spike trains from fluorescence signals.
\item
  \textbf{Neural decoding}: Supervised learning to predict behavior or
  perception from neural activity (e.g., BMI control, speech decoding).
\item
  \textbf{fMRI pattern analysis}: Multivariate pattern analysis (MVPA)
  and deep learning to identify mental states from brain activity
  patterns.
\item
  \textbf{Connectomics}: Computer vision and graph neural networks to
  map neural connectivity from electron microscopy data.
\end{itemize}

\begin{Shaded}
\begin{Highlighting}[]
\ImportTok{import}\NormalTok{ numpy }\ImportTok{as}\NormalTok{ np}
\ImportTok{import}\NormalTok{ torch}
\ImportTok{import}\NormalTok{ torch.nn }\ImportTok{as}\NormalTok{ nn}
\ImportTok{from}\NormalTok{ scipy.signal }\ImportTok{import}\NormalTok{ butter, filtfilt}

\CommentTok{\# Example: Simple neural decoder for BMI control}
\KeywordTok{class}\NormalTok{ NeuralDecoder(nn.Module):}
    \CommentTok{"""A simple LSTM{-}based decoder for predicting movement from neural spikes."""}
    \KeywordTok{def} \FunctionTok{\_\_init\_\_}\NormalTok{(}\VariableTok{self}\NormalTok{, input\_dim, hidden\_dim, output\_dim, num\_layers}\OperatorTok{=}\DecValTok{2}\NormalTok{):}
        \BuiltInTok{super}\NormalTok{().}\FunctionTok{\_\_init\_\_}\NormalTok{()}
        \VariableTok{self}\NormalTok{.lstm }\OperatorTok{=}\NormalTok{ nn.LSTM(input\_dim, hidden\_dim, num\_layers, batch\_first}\OperatorTok{=}\VariableTok{True}\NormalTok{)}
        \VariableTok{self}\NormalTok{.fc }\OperatorTok{=}\NormalTok{ nn.Linear(hidden\_dim, output\_dim)}

    \KeywordTok{def}\NormalTok{ forward(}\VariableTok{self}\NormalTok{, x):}
        \CommentTok{\# x shape: (batch, seq\_len, input\_dim)}
\NormalTok{        lstm\_out, \_ }\OperatorTok{=} \VariableTok{self}\NormalTok{.lstm(x)}
        \CommentTok{\# Take the last time step output}
\NormalTok{        out }\OperatorTok{=} \VariableTok{self}\NormalTok{.fc(lstm\_out[:, }\OperatorTok{{-}}\DecValTok{1}\NormalTok{, :])}
        \ControlFlowTok{return}\NormalTok{ out}

\CommentTok{\# Example: Spike sorting with PCA + clustering}
\KeywordTok{def}\NormalTok{ spike\_sorting(waveforms):}
    \CommentTok{"""Simple spike sorting using PCA and k{-}means."""}
    \CommentTok{\# Perform PCA}
    \ImportTok{from}\NormalTok{ sklearn.decomposition }\ImportTok{import}\NormalTok{ PCA}
\NormalTok{    pca }\OperatorTok{=}\NormalTok{ PCA(n\_components}\OperatorTok{=}\DecValTok{3}\NormalTok{)}
\NormalTok{    pca\_result }\OperatorTok{=}\NormalTok{ pca.fit\_transform(waveforms)}
    \CommentTok{\# Cluster}
    \ImportTok{from}\NormalTok{ sklearn.cluster }\ImportTok{import}\NormalTok{ KMeans}
\NormalTok{    kmeans }\OperatorTok{=}\NormalTok{ KMeans(n\_clusters}\OperatorTok{=}\DecValTok{4}\NormalTok{, random\_state}\OperatorTok{=}\DecValTok{0}\NormalTok{).fit(pca\_result)}
    \ControlFlowTok{return}\NormalTok{ kmeans.labels\_}

\CommentTok{\# Example: Calcium imaging deconvolution}
\KeywordTok{def}\NormalTok{ oasis\_deconvolution(fluor\_signal, dt}\OperatorTok{=}\FloatTok{0.1}\NormalTok{, lambda\_}\OperatorTok{=}\DecValTok{10}\NormalTok{):}
    \CommentTok{"""OASIS algorithm for spike inference from calcium imaging."""}
    \ImportTok{from}\NormalTok{ oasis }\ImportTok{import}\NormalTok{ oasisAR1N}
    \CommentTok{\# Convert to AR1N model}
\NormalTok{    result }\OperatorTok{=}\NormalTok{ oasisAR1N(fluor\_signal, dt}\OperatorTok{=}\NormalTok{dt, lambda\_}\OperatorTok{=}\NormalTok{lambda\_, g}\OperatorTok{=}\FloatTok{0.95}\NormalTok{)}
    \ControlFlowTok{return}\NormalTok{ result[}\StringTok{"spikes"}\NormalTok{]  }\CommentTok{\# Estimated spike times}
\end{Highlighting}
\end{Shaded}

\subsubsection{AI for Drug Discovery in
Neurology}\label{ai-for-drug-discovery-in-neurology}

Traditional drug discovery takes 10-15 years and billions of dollars. AI
is accelerating the process:

\begin{itemize}
\item
  \textbf{Target identification}: Machine learning models analyze
  genomic, transcriptomic, and proteomic data to identify promising drug
  targets.
\item
  \textbf{Compound screening}: Deep learning models predict the activity
  of millions of compounds against specific targets, enabling virtual
  screening.
\item
  \textbf{De novo drug design}: Generative models (GANs, VAEs,
  transformers) can design novel molecules with desired properties.
\item
  \textbf{Clinical trial optimization}: AI helps identify suitable
  patient populations, predict outcomes, and optimize trial design.
\end{itemize}

\textbf{Case study}: Insilico Medicine used AI to identify a novel
target for idiopathic pulmonary fibrosis, design a novel molecule, and
complete preclinical experiments in under 18 months (vs.~typical 4-6
years).

\subsubsection{Brain-Computer Interfaces
(BCIs)}\label{brain-computer-interfaces-bcis}

BCIs translate neural activity into commands for external devices,
offering hope for paralysis, locked-in syndrome, and other neurological
conditions:

\begin{itemize}
\item
  \textbf{Signal acquisition}: Microelectrode arrays (Utah array,
  Blackrock), ECoG, EEG, or optical imaging.
\item
  \textbf{Signal processing}: Filtering, spike detection, feature
  extraction.
\item
  \textbf{Decoding algorithms}: Linear models (Kalman filter, Wiener
  filter), neural networks, and reinforcement learning to map neural
  signals to actions.
\item
  \textbf{Applications}:

  \begin{itemize}
  \tightlist
  \item
    \textbf{Communication}: Typing at 20+ words per minute using
    imagined handwriting.
  \item
    \textbf{Mobility}: Controlling robotic limbs or exoskeletons.
  \item
    \textbf{Sensory restoration}: Cortical visual prostheses for the
    blind.
  \end{itemize}
\end{itemize}

\textbf{Neuralink} and other companies are developing high-bandwidth,
minimally invasive BCIs with thousands of channels, aiming to treat
paralysis and eventually enhance human cognition.

\subsubsection{Ethical Considerations}\label{ethical-considerations-1}

As AI becomes more integrated with neuroscience, important ethical
questions arise:

\begin{itemize}
\tightlist
\item
  \textbf{Privacy}: Neural data is the ultimate private information. How
  do we protect brain data from misuse?
\item
  \textbf{Agency and responsibility}: If a BCI-controlled prosthetic
  acts unintentionally, who is responsible?
\item
  \textbf{Cognitive enhancement}: Should AI-assisted cognitive
  enhancement be allowed? What are the societal implications?
\item
  \textbf{Dual-use}: The same technology that helps paralyzed patients
  could be used for military applications or surveillance.
\item
  \textbf{Bias in AI}: If training data lacks diversity, AI systems may
  perform poorly for underrepresented groups, exacerbating health
  disparities.
\end{itemize}

\begin{Shaded}
\begin{Highlighting}[]
\CommentTok{\# Example: Simple neural encoding model}
\KeywordTok{class}\NormalTok{ NeuralEncoder(nn.Module):}
    \CommentTok{"""A simple model that predicts neural response to visual stimuli."""}
    \KeywordTok{def} \FunctionTok{\_\_init\_\_}\NormalTok{(}\VariableTok{self}\NormalTok{, stimulus\_dim, hidden\_dim, neural\_dim):}
        \BuiltInTok{super}\NormalTok{().}\FunctionTok{\_\_init\_\_}\NormalTok{()}
        \VariableTok{self}\NormalTok{.conv }\OperatorTok{=}\NormalTok{ nn.Conv2d(stimulus\_dim, }\DecValTok{64}\NormalTok{, kernel\_size}\OperatorTok{=}\DecValTok{7}\NormalTok{, stride}\OperatorTok{=}\DecValTok{2}\NormalTok{)}
        \VariableTok{self}\NormalTok{.fc }\OperatorTok{=}\NormalTok{ nn.Linear(}\DecValTok{64} \OperatorTok{*} \DecValTok{7} \OperatorTok{*} \DecValTok{7}\NormalTok{, hidden\_dim)}
        \VariableTok{self}\NormalTok{.output }\OperatorTok{=}\NormalTok{ nn.Linear(hidden\_dim, neural\_dim)}

    \KeywordTok{def}\NormalTok{ forward(}\VariableTok{self}\NormalTok{, stimulus):}
\NormalTok{        x }\OperatorTok{=}\NormalTok{ torch.relu(}\VariableTok{self}\NormalTok{.conv(stimulus))}
\NormalTok{        x }\OperatorTok{=}\NormalTok{ x.view(x.size(}\DecValTok{0}\NormalTok{), }\OperatorTok{{-}}\DecValTok{1}\NormalTok{)}
\NormalTok{        x }\OperatorTok{=}\NormalTok{ torch.relu(}\VariableTok{self}\NormalTok{.fc(x))}
        \ControlFlowTok{return} \VariableTok{self}\NormalTok{.output(x)}

\CommentTok{\# Example: Reinforcement learning for adaptive BCIs}
\ImportTok{import}\NormalTok{ gym}
\ImportTok{from}\NormalTok{ stable\_baselines3 }\ImportTok{import}\NormalTok{ PPO}

\CommentTok{\# Create a custom environment for BCI control}
\KeywordTok{class}\NormalTok{ BCIEnv(gym.Env):}
    \KeywordTok{def} \FunctionTok{\_\_init\_\_}\NormalTok{(}\VariableTok{self}\NormalTok{):}
        \BuiltInTok{super}\NormalTok{().}\FunctionTok{\_\_init\_\_}\NormalTok{()}
        \VariableTok{self}\NormalTok{.observation\_space }\OperatorTok{=}\NormalTok{ gym.spaces.Box(low}\OperatorTok{={-}}\DecValTok{1}\NormalTok{, high}\OperatorTok{=}\DecValTok{1}\NormalTok{, shape}\OperatorTok{=}\NormalTok{(}\DecValTok{100}\NormalTok{,))}
        \VariableTok{self}\NormalTok{.action\_space }\OperatorTok{=}\NormalTok{ gym.spaces.Box(low}\OperatorTok{={-}}\DecValTok{1}\NormalTok{, high}\OperatorTok{=}\DecValTok{1}\NormalTok{, shape}\OperatorTok{=}\NormalTok{(}\DecValTok{2}\NormalTok{,))}

    \KeywordTok{def}\NormalTok{ step(}\VariableTok{self}\NormalTok{, action):}
        \CommentTok{\# Simplified: reward based on how close action is to target}
\NormalTok{        reward }\OperatorTok{=} \OperatorTok{{-}}\NormalTok{np.}\BuiltInTok{sum}\NormalTok{((action }\OperatorTok{{-}} \VariableTok{self}\NormalTok{.target) }\OperatorTok{**} \DecValTok{2}\NormalTok{)}
        \VariableTok{self}\NormalTok{.steps }\OperatorTok{+=} \DecValTok{1}
\NormalTok{        done }\OperatorTok{=} \VariableTok{self}\NormalTok{.steps }\OperatorTok{\textgreater{}=} \DecValTok{200} \KeywordTok{or}\NormalTok{ reward }\OperatorTok{\textless{}} \OperatorTok{{-}}\DecValTok{10}
        \ControlFlowTok{return} \VariableTok{self}\NormalTok{.\_get\_observation(), reward, done, \{\}}

    \KeywordTok{def}\NormalTok{ reset(}\VariableTok{self}\NormalTok{):}
        \VariableTok{self}\NormalTok{.steps }\OperatorTok{=} \DecValTok{0}
        \VariableTok{self}\NormalTok{.target }\OperatorTok{=}\NormalTok{ np.random.randn(}\DecValTok{2}\NormalTok{)}
        \ControlFlowTok{return} \VariableTok{self}\NormalTok{.\_get\_observation()}

    \KeywordTok{def}\NormalTok{ \_get\_observation(}\VariableTok{self}\NormalTok{):}
        \CommentTok{\# Simulated neural data}
        \ControlFlowTok{return}\NormalTok{ np.random.randn(}\DecValTok{100}\NormalTok{)}

\CommentTok{\# Train a policy}
\NormalTok{env }\OperatorTok{=}\NormalTok{ BCIEnv()}
\NormalTok{model }\OperatorTok{=}\NormalTok{ PPO(}\StringTok{"MlpPolicy"}\NormalTok{, env).learn(total\_timesteps}\OperatorTok{=}\DecValTok{10000}\NormalTok{)}
\end{Highlighting}
\end{Shaded}

\textbf{Concept Check 6.1}

\begin{quote}
\begin{enumerate}
\def\labelenumi{\arabic{enumi}.}
\tightlist
\item
  What is the difference between Hodgkin-Huxley and Izhikevich models?
  When would you use each?
\item
  How can machine learning be used to analyze calcium imaging data?
\item
  What are some ethical concerns with brain-computer interfaces?
\item
  How might AI accelerate drug discovery for neurological disorders?
\end{enumerate}
\end{quote}

\begin{center}\rule{0.5\linewidth}{0.5pt}\end{center}

\subsubsection{Neuropeptides}\label{neuropeptides}

\textbf{Characteristics:} Larger molecules (3-40 amino acids);
synthesised in soma (not terminals); stored in large dense-core
vesicles; released by higher-frequency stimulation; diffuse farther
(volume transmission); slower, longer-lasting effects.

\textbf{Examples:} - \textbf{Substance P:} Pain transmission (dorsal
horn); co-released with glutamate from nociceptors - \textbf{Enkephalins
/ endorphins:} Endogenous opioids; bind mu (\texorpdfstring{\ensuremath{\mu}}{mu}), delta (\texorpdfstring{\ensuremath{\delta}}{delta}), kappa (\texorpdfstring{\ensuremath{\kappa}}{kappa})
opioid receptors; inhibit pain pathways - \textbf{NPY:} Appetite
stimulation; anxiolysis; one of the most abundant peptides in the brain
- \textbf{Orexin/hypocretin:} Wakefulness promotion; loss causes
narcolepsy type 1

\begin{center}\rule{0.5\linewidth}{0.5pt}\end{center}

\subsection{Synaptic Plasticity and
Memory}\label{synaptic-plasticity-and-memory}

\subsubsection{Long-Term Potentiation
(LTP)}\label{long-term-potentiation-ltp}

\begin{figure}[htbp]
\centering
\includegraphics[width=0.82\linewidth,height=4.2in,keepaspectratio]{../figures/mermaid_inline/inline_mermaid_0163_fc7d28b2a9c2.png}
\caption{Molecular mechanisms of LTP and neurotransmitter recycling. LTP induction requires NMDA receptor opening (coincidence detection) and Ca\$\textasciicircum{}\{2+\}\$ influx. Early LTP involves CaMKII-mediated AMPA receptor modification. Late LTP requires gene expression (CREB) and structural synaptic changes. Neurotransmitters are cleared by reuptake, enzymatic degradation, or diffusion.}
\end{figure}

\textbf{LTP at Schaffer collateral to CA1 synapse (hippocampus):}

\begin{enumerate}
\def\labelenumi{\arabic{enumi}.}
\tightlist
\item
  High-frequency stimulation (100 Hz, 1 s) produces large EPSPs via AMPA
  receptors
\item
  Sufficient depolarisation removes Mg\(^{2+}\) block from NMDA
  receptors
\item
  Ca\(^{2+}\) floods into the postsynaptic spine via NMDA receptors
\item
  Ca\(^{2+}\) activates \textbf{CaMKII} (calcium/calmodulin-dependent
  protein kinase II):

  \begin{itemize}
  \tightlist
  \item
    Autophosphorylation at T286 makes CaMKII constitutively active
    (Ca\(^{2+}\)-independent)
  \item
    Phosphorylates AMPA receptors at Ser831, increasing single-channel
    conductance
  \end{itemize}
\item
  AMPA receptor trafficking: New AMPA receptors inserted into the
  synapse from recycling endosomes
\item
  \textbf{Silent synapses:} Some synapses have NMDA receptors but no
  AMPA receptors. LTP ``unmasks'' these by inserting AMPA receptors,
  converting them to functional synapses.
\item
  \textbf{Late LTP} (\textgreater3 h): Requires new gene expression.
  Ca\(^{2+}\) activates PKA, which phosphorylates CREB. CREB-mediated
  \hyperref[gl:transcription]{\textbf{transcription}} produces Arc,
  BDNF, c-Fos. Structural changes include dendritic spine enlargement
  and formation of new spines.
\end{enumerate}

\subsubsection{Long-Term Depression
(LTD)}\label{long-term-depression-ltd}

Low-frequency stimulation (1 Hz, 15 min) produces modest Ca\(^{2+}\)
entry via NMDA receptors: - Lower Ca\(^{2+}\) levels preferentially
activate protein phosphatases (PP1, calcineurin/PP2B) rather than CaMKII
- Phosphatases dephosphorylate AMPA receptors - AMPA receptors are
endocytosed (clathrin-mediated) from the synapse - Synaptic strength
decreases

\textbf{Bidirectional plasticity:} The same synapse can undergo LTP or
LTD depending on the pattern of Ca\(^{2+}\) entry. This is captured by
the \textbf{BCM (Bienenstock-Cooper-Munro) rule:} there exists a
modification threshold (\(\theta_M\)) -- activity above \(\theta_M\)
causes LTP, below causes LTD. The threshold itself slides based on
recent postsynaptic activity (metaplasticity).

\begin{quote}
\textbf{Concept Check 2:} A neuroscientist finds that pharmacologically
blocking the NMDA receptor abolishes LTP at hippocampal
Schaffer-collateral \texorpdfstring{\ensuremath{\to}}{->} CA1 synapses but does not abolish LTP at
mossy-fibre \texorpdfstring{\ensuremath{\to}}{->} CA3 synapses. Propose a mechanism that explains why these
two forms of LTP differ. Given that the CA3 form is presynaptic
(cAMP-PKA-dependent), which pharmacological manipulation would
selectively block it without affecting NMDA-receptor-dependent LTP?
\end{quote}

\begin{center}\rule{0.5\linewidth}{0.5pt}\end{center}

\subsection{Drug Mechanisms}\label{drug-mechanisms}

{\def\LTcaptype{none} % do not increment counter
\begin{longtable}[]{@{}
  >{\raggedright\arraybackslash}p{(\linewidth - 6\tabcolsep) * \real{0.0976}}
  >{\raggedright\arraybackslash}p{(\linewidth - 6\tabcolsep) * \real{0.3902}}
  >{\raggedright\arraybackslash}p{(\linewidth - 6\tabcolsep) * \real{0.2195}}
  >{\raggedright\arraybackslash}p{(\linewidth - 6\tabcolsep) * \real{0.2927}}@{}}
\toprule\noalign{}
\begin{minipage}[b]{\linewidth}\raggedright
Drug
\end{minipage} & \begin{minipage}[b]{\linewidth}\raggedright
Molecular Target
\end{minipage} & \begin{minipage}[b]{\linewidth}\raggedright
Mechanism
\end{minipage} & \begin{minipage}[b]{\linewidth}\raggedright
Clinical Use
\end{minipage} \\
\midrule\noalign{}
\endhead
\bottomrule\noalign{}
\endlastfoot
\textbf{Cocaine} & DAT (dopamine transporter) & Blocks DA reuptake; DA
accumulates in synapse & Abuse potential; local anaesthetic (Na\(^+\)
channel block) \\
\textbf{Amphetamine} & DAT, VMAT & Reverses DAT (DA efflux); releases
vesicular DA & ADHD; narcolepsy; abuse \\
\textbf{SSRIs} (fluoxetine, sertraline) & SERT (serotonin transporter) &
Block 5-HT reuptake & Depression; anxiety; OCD \\
\textbf{SNRIs} (venlafaxine, duloxetine) & SERT + NET & Block both 5-HT
and NE reuptake & Depression; neuropathic pain; generalised anxiety
disorder \\
\textbf{Morphine/heroin} & \texorpdfstring{\ensuremath{\mu}}{mu}-opioid receptor & Agonist; inhibits GABA
interneurons (disinhibition of DA) in VTA & Pain; abuse \\
\textbf{Naloxone} & \texorpdfstring{\ensuremath{\mu}}{mu}-opioid receptor & Antagonist; reverses opioid
effects & Opioid overdose reversal \\
\textbf{Buprenorphine} & \texorpdfstring{\ensuremath{\mu}}{mu}-opioid receptor & Partial agonist; ceiling
effect; long half-life & Opioid use disorder (maintenance treatment) \\
\textbf{Benzodiazepines} (diazepam) & GABA\(_A\) (\texorpdfstring{\ensuremath{\alpha}}{alpha}/\texorpdfstring{\ensuremath{\gamma}}{gamma} interface) &
Positive allosteric modulator; increases Cl\(^-\) channel open
\textbf{frequency} & Anxiety; seizures; insomnia \\
\textbf{Barbiturates} (phenobarbital) & GABA\(_A\) (transmembrane
domain) & Positive allosteric modulator; increases Cl\(^-\) channel open
\textbf{duration}; direct activation at high doses & Seizures;
anaesthesia; narrow therapeutic index \\
\textbf{Ketamine} & NMDA receptor & Non-competitive antagonist
(open-channel block) & Anaesthesia; rapid-acting antidepressant
(sub-anaesthetic dose) \\
\textbf{Ethanol} & GABA\(_A\) + NMDA & Potentiates GABA\(_A\) (increased
Cl\(^-\) conductance); inhibits NMDA & Recreational; anxiolytic \\
\textbf{Caffeine} & Adenosine A\(_1\)/A\(_{2A}\) receptors & Competitive
antagonist; blocks adenosine's sleep-promoting effect & Wakefulness;
headache treatment \\
\textbf{L-DOPA} & AADC (aromatic amino acid decarboxylase) & Dopamine
precursor; crosses BBB; converted to DA in brain & Parkinson's
disease \\
\textbf{Levodopa + carbidopa} & AADC (peripheral inhibition by
carbidopa) & Prevents peripheral conversion; maximises CNS L-DOPA
delivery & Parkinson's disease \\
\textbf{Haloperidol} & D2 receptor & Competitive D2 antagonist &
First-generation antipsychotic (schizophrenia); risk of tardive
dyskinesia \\
\textbf{Clozapine} & D4, 5-HT\(_{2A}\), H1, M1, \(\alpha_1\) &
``Atypical'' antipsychotic --- low D2 affinity, high 5-HT\(_{2A}\) block
& Treatment-resistant schizophrenia; risk of agranulocytosis requires
WBC monitoring \\
\textbf{Curare (d-tubocurarine)} & Muscle nAChR & Competitive
antagonist; blocks NMJ transmission & Muscle relaxant (anaesthesia) \\
\textbf{Botulinum toxin} & SNARE (SNAP-25) & Cleaves SNAP-25; blocks
vesicle fusion & Focal dystonia; cosmetic (Botox); therapeutic
spasticity \\
\textbf{Gabapentin; pregabalin} & Ca-channel \(\alpha_2\delta\) subunit
& Reduces dorsal-horn neurotransmitter release & Neuropathic pain;
epilepsy; fibromyalgia; anxiety \\
\textbf{Levetiracetam} & SV2A (synaptic vesicle glycoprotein) & Binds
SV2A; reduces Ca\(^{2+}\)-evoked vesicle release & Broad-spectrum
antiepileptic; minimal drug interactions \\
\textbf{Lithium} & IMPase, GSK-3\texorpdfstring{\ensuremath{\beta}}{beta}, and others & Inhibits inositol
monophosphatase (depletes PI signalling substrate); inhibits GSK-3\texorpdfstring{\ensuremath{\beta}}{beta}
(neuroprotection) & Bipolar disorder (mood stabiliser); narrow
therapeutic index (Li\(^+\) 0.6--1.2 mEq/L) \\
\end{longtable}
}

\begin{quote}
\textbf{Clinical Connection:} The opioid epidemic has killed over
500,000 Americans since 1999. Understanding the molecular pharmacology
is essential: opioids activate mu receptors on GABAergic interneurons in
the VTA, removing tonic inhibition of dopamine neurons. The resulting
dopamine surge in the nucleus accumbens produces euphoria and
reinforcement. Tolerance develops as mu receptors are desensitised and
downregulated. Withdrawal occurs because the compensatory upregulation
of cAMP signalling is unmasked when the drug is removed. Naloxone
(Narcan) competitively displaces opioids from mu receptors and can
reverse respiratory depression within minutes.
\end{quote}

\begin{center}\rule{0.5\linewidth}{0.5pt}\end{center}

\subsection{Worked Example}\label{worked-example-3}

\textbf{Problem:} You voltage-clamp a frog neuromuscular junction and
record spontaneous miniature end-plate currents (mEPCs). From 500
events, the mean mEPC amplitude is \(q = 1.2\) nA (quantal size). You
then stimulate the motor nerve with single pulses and measure evoked
EPCs.

Over 200 evoked trials: - 8 trials produce no response (failures) - 72
trials produce a single quantal event (about 1.2 nA) - 68
trials produce two quantal events (about 2.4 nA) - 36 trials
produce three quantal events (about 3.6 nA) - 12 trials produce
four quantal events (about 4.8 nA) - 4 trials produce five
quantal events (about 6.0 nA)

\textbf{(a)} Calculate the mean quantal content \(\bar{m}\) = mean
number of quanta released per stimulus.

\begin{equation}
\bar{m} = \frac{0(8) + 1(72) + 2(68) + 3(36) + 4(12) + 5(4)}{200} = \frac{0 + 72 + 136 + 108 + 48 + 20}{200} = \frac{384}{200} = 1.92
\label{eq:action_potential_synapses_5}
\end{equation}

\textbf{(b)} Use the Poisson approximation to estimate the number of
failures expected if \(\bar{m}\) = 1.92:

For a Poisson process, the fraction of failures
\(P(0) = e^{-\bar{m}} = e^{-1.92} = 0.146\)

Expected failures = \(0.146 \times 200 = 29.3\). Observed failures = 8.
\textbf{The discrepancy} suggests the release is \textbf{not} strictly
Poisson --- there may be heterogeneity in vesicle release probabilities
across active zones, or the preparation has an unusually low failure
rate due to high quantal content.

\textbf{(c)} Use the method of failures to give an alternative estimate
of \(\bar{m}\):

\begin{equation}
\bar{m} = \ln\left(\frac{N}{n_0}\right) = \ln\left(\frac{200}{8}\right) = \ln(25) = 3.22
\label{eq:action_potential_synapses_6}
\end{equation}

The discrepancy between the two estimates (1.92 by direct count vs 3.22
by method of failures) indicates that the \textbf{binomial model} (with
heterogeneous release probabilities) fits better than the simple
Poisson. This is a classic result: real synapses have multiple active
zones with heterogeneous \(p\) values, violating the simple Poisson
assumption (Zucker 1973; Bhumbra \& Bhatt 2022).

\textbf{(d)} Given \(\bar{m} = 1.92\) and the number of release-ready
vesicles \(n = 10\) (estimated by high-osmotic sucrose stimulation),
estimate the mean release probability \(\bar{p}\):

\begin{equation}
\bar{p} = \frac{\bar{m}}{n} = \frac{1.92}{10} = 0.192
\label{eq:action_potential_synapses_7}
\end{equation}

This is a typical resting release probability for a vertebrate
neuromuscular junction (0.1--0.4 at physiological Ca\(^{2+}\)). The
\(n \cdot p \cdot q\) synaptic model predicts that drugs which increase
\(p\) (e.g., by prolonging the AP and increasing Ca\(^{2+}\) influx,
such as aminopyridines) will linearly increase \(\bar{m}\) and thus EPP
amplitude --- a pharmacological basis for treatments in Lambert-Eaton
myasthenic syndrome.

\begin{center}\rule{0.5\linewidth}{0.5pt}\end{center}

\subsection{Current Evidence and Frontier
Biology}\label{current-evidence-and-frontier-biology-32}

For \textbf{Action Potentials and Synaptic Transmission}, frontier
biology belongs inside the evidence logic of the chapter. Physiology now
blends mechanism with allostasis, immune-endocrine-neural coupling,
wearable data, and individualized risk without reducing bodies to simple
machines. The core reading question is this: synaptic claims require
ion-channel timing, driving force, transmitter release, receptor
dynamics, and plasticity.

\begin{itemize}
\tightlist
\item
  \textbf{What to verify:} identify the observation, model, assay, or
  dataset that would make the claim stronger or weaker.
\item
  \textbf{What to qualify:} state the scale, organism, cell type,
  environmental condition, or population where the claim is expected to
  hold.
\item
  \textbf{What to compare:} test at least one alternative explanation,
  baseline, or null model before treating the pattern as causal.
\item
  \textbf{What to cite:} distinguish primary evidence, review synthesis,
  public dataset, and institutional guidance; for recent or numeric
  claims, prefer the source closest to the measurement and state what
  has changed since it was published.
\end{itemize}

Interpret physiological data by separating baseline variation,
perturbation response, compensation, and the threshold where
compensation becomes pathology.

\textbf{Source practice:} For physiology claims, cite the measurement
context and distinguish baseline variation, compensation,
pathophysiology, and treatment evidence.

\subsection{Key Terms}\label{key-terms-32}

{\def\LTcaptype{none} % do not increment counter
\begin{longtable}[]{@{}
  >{\raggedright\arraybackslash}p{(\linewidth - 2\tabcolsep) * \real{0.2857}}
  >{\raggedright\arraybackslash}p{(\linewidth - 2\tabcolsep) * \real{0.7143}}@{}}
\toprule\noalign{}
\begin{minipage}[b]{\linewidth}\raggedright
Term
\end{minipage} & \begin{minipage}[b]{\linewidth}\raggedright
Definition
\end{minipage} \\
\midrule\noalign{}
\endhead
\bottomrule\noalign{}
\endlastfoot
\textbf{Hodgkin-Huxley model} & Mathematical description of AP using
gating variables m, h, n for Na\(^+\) and K\(^+\) conductances \\
\textbf{Threshold} & Membrane potential (about \(-55\) mV) at
which inward Na\(^+\) current exceeds outward K\(^+\) current \\
\textbf{Absolute refractory period} & about 1-2 ms when
Na\(^+\) channels are inactivated; no AP possible regardless of
stimulus \\
\textbf{Relative refractory period} & about 5-10 ms when
Na\(^+\) channels recovering; stronger stimulus needed for AP \\
\textbf{Saltatory conduction} & AP jumps between nodes of Ranvier in
myelinated axons \\
\textbf{SNARE complex} & Synaptobrevin + syntaxin + SNAP-25: core
vesicle fusion machinery \\
\textbf{Synaptotagmin} & Ca\(^{2+}\) sensor that triggers vesicle
fusion; C2 domains bind 4-5 Ca\(^{2+}\) ions \\
\textbf{Quantal release} & Neurotransmitter released in discrete
vesicle-sized packets (Katz, 1950s) \\
\textbf{NMDA receptor} & Coincidence detector: requires glutamate +
depolarisation; Ca\(^{2+}\) permeable \\
\textbf{CaMKII} & Calcium/calmodulin kinase II; autophosphorylation at
T286; key LTP effector \\
\textbf{LTP} & Long-term potentiation: persistent increase in synaptic
strength \\
\textbf{LTD} & Long-term depression: persistent decrease in synaptic
strength \\
\textbf{Excitotoxicity} & Neuronal death from excessive
glutamate/Ca\(^{2+}\) (stroke, epilepsy) \\
\textbf{Tetrodotoxin (TTX)} & Puffer fish toxin blocking Nav channels \\
\textbf{BCM rule} & Sliding modification threshold determining LTP vs
LTD \\
\textbf{Silent synapse} & Synapse with NMDA but no AMPA receptors;
``unmasked'' during LTP \\
\end{longtable}
}

\begin{center}\rule{0.5\linewidth}{0.5pt}\end{center}

\subsection{Review Questions}\label{review-questions-32}

\begin{enumerate}
\def\labelenumi{\arabic{enumi}.}
\item
  Using the Hodgkin-Huxley equations, explain why the peak of the action
  potential does not reach E\(_{Na}\) (+50 mV) but instead peaks at
  approximately +30 to +40 mV. Which gating variable is responsible?
\item
  Explain the molecular basis of the absolute refractory period. Why is
  this essential for unidirectional AP propagation?
\item
  Compare saltatory conduction in myelinated axons with continuous
  conduction in unmyelinated axons. Why does velocity scale as
  \(\sqrt{d}\) in unmyelinated axons but as \(d\) in myelinated axons?
\item
  Describe the roles of the five key presynaptic proteins:
  synaptobrevin, syntaxin, SNAP-25, synaptotagmin-1, and complexin. What
  happens if each is eliminated?
\item
  The NMDA receptor is often called the ``Hebbian synapse detector.''
  Explain this concept, linking the receptor's biophysical properties to
  Hebb's postulate and to LTP induction.
\item
  A patient presents with progressive muscle weakness. Antibodies
  against the nicotinic ACh receptor are detected. Name the disease,
  explain the pathophysiology, and describe two treatment strategies
  based on cholinergic pharmacology.
\item
  Compare the mechanisms of cocaine and amphetamine at dopaminergic
  synapses. Both increase synaptic dopamine, but by different
  mechanisms. Why might their clinical effects and addiction profiles
  differ?
\item
  Explain why benzodiazepines are safer in overdose than barbiturates,
  despite both acting on GABA\(_A\) receptors. (Hint: consider the
  difference between increasing channel open \textbf{frequency} vs open
  \textbf{duration}, and the concept of ceiling effect.)
\item
  Design a hypothetical drug that could enhance memory formation by
  targeting the LTP pathway. Identify the molecular target, predict the
  therapeutic effect, and discuss potential side effects (excitotoxicity
  risk).
\item
  Why does multiple sclerosis cause such diverse neurological symptoms?
  Explain the relationship between lesion location in white matter
  tracts and the specific symptoms produced.
\item
  At a synapse, you observe mEPP amplitude of 0.5 mV and mean EPP
  amplitude of 4.5 mV (low-Ca\(^{2+}\) conditions to allow quantal
  resolution). (a) Estimate the mean quantal content \(\bar{m}\). (b)
  Use the Poisson failure method to calculate the expected probability
  of a failure. (c) If you increase extracellular Ca\(^{2+}\) from 1 mM
  to 2 mM, predict whether \(\bar{m}\) will increase linearly,
  quadratically, or via a 4th-power relationship (based on synaptotagmin
  cooperativity), and explain the molecular basis.
\item
  Clozapine is described as an ``atypical'' antipsychotic with lower
  risk of tardive dyskinesia compared to haloperidol, despite both
  treating positive symptoms of schizophrenia. Using the drug table,
  explain: (a) their differing receptor binding profiles; (b) why high
  D2 occupancy causes tardive dyskinesia; (c) why clozapine's unique
  receptor profile (including 5-HT\(_{2A}\) blockade) may explain its
  superior efficacy in treatment-resistant schizophrenia.
\end{enumerate}

\begin{center}\rule{0.5\linewidth}{0.5pt}\end{center}

\subsection{Further Reading and Source
Notes}\label{further-reading-and-source-notes-27}

\begin{itemize}
\tightlist
\item
  Frey \& Morris (1997). Synaptic tagging and long-term potentiation.
  \emph{Nature}, 385.
\item
  Hodgkin \& Huxley (1952). A quantitative description of membrane
  current and its application to conduction and excitation in nerve.
  \emph{Journal of Physiology}, 117.
\item
  Katz (1969). \emph{The Release of Neural Transmitter Substances}.
  Liverpool University Press.
\item
  Bliss \& Lømo (1973). Long-lasting potentiation of synaptic
  transmission in the dentate area of the anaesthetized rabbit.
  \emph{Journal of Physiology}, 232.
\item
  Kandel et al.~(latest ed.). \emph{Principles of Neural Science}
  (action potential and synaptic transmission chapters). McGraw-Hill.
\item
  Neher \& Sakmann (1976). Single-channel currents recorded from
  membrane of denervated frog muscle fibres. \emph{Nature}, 260.
\end{itemize}

\begin{center}\rule{0.5\linewidth}{0.5pt}\end{center}

\subsection{Computational Bridge}\label{computational-bridge-29}

Driving-force synaptic current in pA follows
\(I = g\,(V_m - E_\mathrm{rev})\):

\begin{Shaded}
\begin{Highlighting}[]
\ImportTok{from}\NormalTok{ biology.neuroscience }\ImportTok{import}\NormalTok{ synaptic\_current}

\NormalTok{exc }\OperatorTok{=}\NormalTok{ synaptic\_current(}\FloatTok{0.0}\NormalTok{, }\OperatorTok{{-}}\FloatTok{65.0}\NormalTok{, }\FloatTok{12.0}\NormalTok{)}
\NormalTok{inh }\OperatorTok{=}\NormalTok{ synaptic\_current(}\OperatorTok{{-}}\FloatTok{70.0}\NormalTok{, }\OperatorTok{{-}}\FloatTok{65.0}\NormalTok{, }\FloatTok{20.0}\NormalTok{)}
\BuiltInTok{print}\NormalTok{(}\BuiltInTok{round}\NormalTok{(exc.peak\_current\_pA, }\DecValTok{2}\NormalTok{), }\BuiltInTok{round}\NormalTok{(inh.peak\_current\_pA, }\DecValTok{2}\NormalTok{))}
\end{Highlighting}
\end{Shaded}

\begin{quote}
\textbf{Clinical / systems note:} Benzodiazepines increase GABA\(_A\)
\textbf{frequency} without changing \(E_\mathrm{rev}\); barbiturates
prolong open time --- both shift inhibitory current but with different
overdose ceilings.
\end{quote}

\begin{center}\rule{0.5\linewidth}{0.5pt}\end{center}

\subsection{Summary}\label{summary-32}

\begin{itemize}
\tightlist
\item
  \textbf{Hodgkin-Huxley:} \(I_{Na} = \bar{g}_{Na} m^3 h (V - E_{Na})\),
  \(I_K = \bar{g}_K n^4 (V - E_K)\). Gating variables obey first-order
  kinetics. m (fast activation), h (slow inactivation), n (slow
  activation).
\item
  \textbf{Refractory periods:} Absolute (Na\(^+\) inactivated, h about
  0, 1-2 ms): cannot fire. Relative (K\(^+\) channels open, membrane
  hyperpolarised, 5-10 ms): stronger stimulus needed.
\item
  \textbf{Na\(^+\) channel states:} Closed (ready) to Open (conducting,
  about 0.5 ms) to Inactivated (blocked from inside, cannot
  reopen). TTX and local anaesthetics block the pore.
\item
  \textbf{Myelination:} Saltatory conduction between nodes of Ranvier.
  Velocity proportional to diameter (linear). A-alpha fibres: 70-120
  m/s; C fibres: 0.5-2 m/s. MS: autoimmune demyelination.
\item
  \textbf{Chemical synapse:} AP opens Cav2.1/2.2 channels. Ca\(^{2+}\)
  binds synaptotagmin-1. Complexin displaced. SNARE complex zippers.
  Vesicle fuses. NT released in \textless1 ms. Quantal release (Katz).
\item
  \textbf{Neurotransmitters:} Glutamate (AMPA/NMDA/mGluR, excitatory);
  GABA (GABA\(_A\)/GABA\(_B\), inhibitory); ACh (nicotinic/muscarinic);
  DA (D1-D5, reward/motor); 5-HT (mood/sleep); NE (arousal);
  neuropeptides (volume transmission, pain, reward).
\item
  \textbf{LTP:} NMDA-Ca\(^{2+}\)-CaMKII-AMPA pathway. Early:
  phosphorylation and insertion of AMPA receptors. Late: CREB-dependent
  gene expression, structural spine changes. LTD: opposite; phosphatases
  remove AMPA receptors.
\item
  \textbf{Drug mechanisms:} Cocaine (DAT block), morphine (mu agonist),
  SSRIs (SERT block), benzodiazepines (GABA\(_A\) PAM), ketamine (NMDA
  block), botulinum toxin (SNAP-25 cleavage).
\item
  \textbf{Connections:} See \cref{sec:unit_IX_nervous_system} for
  resting potential and integration,
  \cref{sec:unit_VII_infectious_disease} for synaptic pathogens and
  toxins, and \cref{sec:unit_II_membrane_transport} for ion channels.
\end{itemize}

\begin{center}\rule{0.5\linewidth}{0.5pt}\end{center}

\subsubsection{Companion Source
Module}\label{companion-source-module-32}

\textbf{Action Potentials and Synaptic Transmission} should leave a
reproducible trail from a biological claim to the code, figure, diagram,
or paper-based activity that can test it. Use the surfaces below to
inspect the chapter's assumptions, rerun the relevant model, or compare
the manuscript explanation with companion labs and figures.

{\def\LTcaptype{none} % do not increment counter
\begin{longtable}[]{@{}
  >{\raggedright\arraybackslash}p{(\linewidth - 2\tabcolsep) * \real{0.5000}}
  >{\raggedright\arraybackslash}p{(\linewidth - 2\tabcolsep) * \real{0.5000}}@{}}
\toprule\noalign{}
\begin{minipage}[b]{\linewidth}\raggedright
Surface
\end{minipage} & \begin{minipage}[b]{\linewidth}\raggedright
Use it for
\end{minipage} \\
\midrule\noalign{}
\endhead
\bottomrule\noalign{}
\endlastfoot
\breaktt{src/biology/neuroscience/neuroscience.py}
(\breaktt{action\_potential\_hh}, \texttt{synaptic\_current},
\breaktt{cable\_voltage\_attenuation}) & Reproduce spike timing,
postsynaptic currents, and passive spread. \\
\breaktt{src/biology/cell/cell\_biology.py} (\texttt{nernst\_potential},
\texttt{goldman\_equation}) & Check ion gradients and membrane-voltage
assumptions. \\
\breaktt{src/visualization/plots.py} (\breaktt{plot\_action\_potential},
\breaktt{plot\_nernst\_potentials}) & Compare calculated voltages with
plotted signals. \\
\end{longtable}
}

\textbf{Reproducibility check:} list ion concentrations, conductances,
reversal potentials, synaptic delay, and receptor type before
interpreting excitability. \textbf{Cross-reference:} connect with
\cref{sec:unit_II_membrane_transport} and
\cref{sec:unit_IX_nervous_system}.

\newpage

\section{Endocrine and Immune
Systems}\label{endocrine-and-immune-systems}

\label{sec:unit_IX_endocrine_and_immune}

\begin{quote}
\textbf{Ch 31} \texorpdfstring{\ensuremath{\cdot}}{.} Level 2/3 \texorpdfstring{\ensuremath{\cdot}}{.} 55 min read \texorpdfstring{\ensuremath{\cdot}}{.} 75 min lecture \texorpdfstring{\ensuremath{\cdot}}{.}
Prerequisites: \cref{sec:unit_IX_circulation_respiration_homeostasis}
\end{quote}

\subsection{Learning Objectives}\label{learning-objectives-33}

\begin{enumerate}
\def\labelenumi{\arabic{enumi}.}
\tightlist
\item
  Compare endocrine and nervous system signalling in terms of speed,
  duration, and specificity.
\item
  Classify \hyperref[gl:hormone]{\textbf{hormone}}s by chemical class
  (peptide, steroid, amine, eicosanoid) and describe their synthesis,
  transport, receptor location, and signalling duration.
\item
  Explain the HPA, HPT, and HPG axes with feedback regulation, including
  detailed steroidogenesis and the circadian profile of cortisol.
\item
  Trace insulin and glucagon signalling in glucose
  \hyperref[gl:homeostasis]{\textbf{homeostasis}}, including
  IR/IRS/PI3K/Akt/GLUT4, crosstalk with leptin and GLP-1, and the
  pathophysiology of diabetes.
\item
  Describe the adrenal gland structure and function (cortex and
  medulla), including cortisol synthesis from cholesterol.
\item
  Describe thyroid hormone synthesis, T4\texorpdfstring{\ensuremath{\to}}{->}T3 conversion, the nuclear
  receptor mechanism, and the Wolff-Chaikoff effect.
\item
  Describe the growth hormone axis and IGF-1.
\item
  Explain prostaglandin and eicosanoid synthesis from arachidonic acid;
  differentiate non-selective NSAIDs from COX-2 selective inhibitors.
\item
  Describe endocrine disruption by xenoestrogens, BPA, phthalates, and
  PFAS.
\item
  Distinguish innate and adaptive immunity, including PRRs (TLRs, NLRs,
  RLRs) and their downstream signalling pathways (MyD88/NF-\texorpdfstring{\ensuremath{\kappa}}{kappa}B vs
  TRIF/IRF3 vs cGAS-STING).
\item
  Explain the complement system (classical, lectin, alternative
  pathways), the C3/C5 convertases, the MAC, and amplification dynamics.
\item
  Describe T cell development through the DN1--DN4 stages, positive and
  negative thymic selection, and TCR diversity generation.
\item
  Describe B cell activation, T-dependent vs T-independent responses,
  germinal centre reactions, somatic hypermutation, and class switching.
\item
  Tabulate the cytokine network and key cytokines (IL-1, IL-6, TNF,
  IFN-\texorpdfstring{\ensuremath{\gamma}}{gamma}, IL-10, IL-4, IL-17, etc.).
\item
  Describe immunological memory formation and explain why memory cells
  respond faster.
\item
  Explain mechanisms of central and peripheral tolerance and how their
  failure causes autoimmunity (molecular mimicry, bystander activation,
  epitope spreading).
\item
  Describe the Type I--IV classification of hypersensitivity reactions
  and their treatments.
\item
  Explain immunotherapy: checkpoint inhibitors (PD-1/PD-L1, CTLA-4) and
  CAR-T cell therapy.
\end{enumerate}

\subsubsection{Study Blueprint}\label{study-blueprint-33}

\begin{itemize}
\tightlist
\item
  \textbf{Big idea:} Long-range signaling and immune recognition
  coordinate body-wide adaptation and defense.
\item
  \textbf{Core concepts:} hormones, receptors, innate immunity, adaptive
  immunity.
\item
  \textbf{Framework alignment:} Vision \& Change: Structure and
  function, Systems; AP Biology: Systems Interactions, Energetics;
  NGSS-style topics: Structure and Function.
\item
  \textbf{Model or quantitative lens:} Hormone feedback, dose-response,
  and immune-memory reasoning.
\item
  \textbf{Data skill:} Interpret endocrine or immune data from time
  courses, titers, or perturbations.
\item
  \textbf{Practice cadence:} Visual Representations, Statistical Tests
  and Data Analysis, Argumentation.
\item
  \textbf{Common misconception to repair:} Immunity is not just attack;
  recognition, tolerance, memory, and regulation are equally central.
\item
  \textbf{Primary lab:} \cref{sec:lab_unit_IX_endocrine_and_immune}.
\item
  \textbf{Question bank:} \cref{sec:q_unit_IX_endocrine_and_immune}.
\item
  \textbf{Transfer task:} Transfer signaling and immunity reasoning to
  vaccination, autoimmunity, stress, and metabolism.
\item
  \textbf{Bridge to computation:}
  \breaktt{biology.physiology.physiology.homeostasis\_response}.
\end{itemize}

\begin{center}\rule{0.5\linewidth}{0.5pt}\end{center}

\begin{quote}
\textbf{Opening Vignette --- The Hormone That Changed Medicine Forever}

Before 1921, Type 1 diabetes was a death sentence. Children diagnosed
with it were placed on starvation diets --- sometimes eating fewer than
500 calories per day --- which extended their lives by months while
slowly wasting them. Then Frederick Banting, a young Canadian surgeon,
persuaded the University of Toronto to give him laboratory space and a
few dogs. Working with student Charles Best and biochemist J.B. Collip
to purify the extract, they isolated the pancreatic secretion that
controlled blood glucose --- insulin. The first human injection was
given to 14-year-old Leonard Thompson on January 11, 1922. He had been
near death; within days his blood glucose normalised and he survived.
Banting and John Macleod received the Nobel Prize in 1923. Insulin was
the first hormone to be purified, the first to be sequenced (by
Frederick Sanger, 1951), and the first to be produced by recombinant DNA
technology (1982). No single molecule has had a more direct impact on
human survival.
\end{quote}

\subsection{Endocrine System Overview}\label{endocrine-system-overview}

\subsubsection{Endocrine vs Nervous
System}\label{endocrine-vs-nervous-system}

{\def\LTcaptype{none} % do not increment counter
\begin{longtable}[]{@{}
  >{\raggedright\arraybackslash}p{(\linewidth - 4\tabcolsep) * \real{0.1892}}
  >{\raggedright\arraybackslash}p{(\linewidth - 4\tabcolsep) * \real{0.3784}}
  >{\raggedright\arraybackslash}p{(\linewidth - 4\tabcolsep) * \real{0.4324}}@{}}
\toprule\noalign{}
\begin{minipage}[b]{\linewidth}\raggedright
Feature
\end{minipage} & \begin{minipage}[b]{\linewidth}\raggedright
Nervous System
\end{minipage} & \begin{minipage}[b]{\linewidth}\raggedright
Endocrine System
\end{minipage} \\
\midrule\noalign{}
\endhead
\bottomrule\noalign{}
\endlastfoot
Signal type & Electrical + chemical (neurotransmitter) & Chemical
(hormone) \\
Speed & Fast (ms) & Slow (seconds to hours) \\
Duration & Brief (ms) & Prolonged (minutes to days) \\
Target & Specific (synapse) & Widespread (most cells with receptor) \\
Distance & Short (synaptic cleft, 20 nm) & Long (via blood
circulation) \\
\end{longtable}
}

\textbf{Mixed signalling:} Neuroendocrine cells (e.g., hypothalamic
\hyperref[gl:neuron]{\textbf{neuron}}s secreting releasing hormones,
adrenal medulla chromaffin cells) bridge both systems. Sterling and
Eyer's concept of \hyperref[gl:allostasis]{\textbf{allostasis}}
\citep{sterling1988,sterling2015} extends classical homeostasis: the
brain anticipates physiological needs and adjusts setpoints
predictively, integrating endocrine and autonomic outputs.

\subsubsection{Hormone Classes --- Synthesis, Transport, and
Mechanism}\label{hormone-classes-synthesis-transport-and-mechanism}

{\def\LTcaptype{none} % do not increment counter
\begin{longtable}[]{@{}
  >{\raggedright\arraybackslash}p{(\linewidth - 14\tabcolsep) * \real{0.0549}}
  >{\raggedright\arraybackslash}p{(\linewidth - 14\tabcolsep) * \real{0.1099}}
  >{\raggedright\arraybackslash}p{(\linewidth - 14\tabcolsep) * \real{0.0989}}
  >{\raggedright\arraybackslash}p{(\linewidth - 14\tabcolsep) * \real{0.1978}}
  >{\raggedright\arraybackslash}p{(\linewidth - 14\tabcolsep) * \real{0.1868}}
  >{\raggedright\arraybackslash}p{(\linewidth - 14\tabcolsep) * \real{0.1758}}
  >{\raggedright\arraybackslash}p{(\linewidth - 14\tabcolsep) * \real{0.0879}}
  >{\raggedright\arraybackslash}p{(\linewidth - 14\tabcolsep) * \real{0.0879}}@{}}
\toprule\noalign{}
\begin{minipage}[b]{\linewidth}\raggedright
Class
\end{minipage} & \begin{minipage}[b]{\linewidth}\raggedright
Solubility
\end{minipage} & \begin{minipage}[b]{\linewidth}\raggedright
Synthesis
\end{minipage} & \begin{minipage}[b]{\linewidth}\raggedright
Transport in Blood
\end{minipage} & \begin{minipage}[b]{\linewidth}\raggedright
Receptor Location
\end{minipage} & \begin{minipage}[b]{\linewidth}\raggedright
Signalling Speed
\end{minipage} & \begin{minipage}[b]{\linewidth}\raggedright
Duration
\end{minipage} & \begin{minipage}[b]{\linewidth}\raggedright
Examples
\end{minipage} \\
\midrule\noalign{}
\endhead
\bottomrule\noalign{}
\endlastfoot
\textbf{Peptide/\hyperref[gl:protein]{\textbf{protein}}} & Water-soluble
& Ribosomal synthesis as preprohormones; cleaved in ER/Golgi; stored in
secretory granules & Free in plasma & Plasma membrane (RTKs, GPCRs) &
Seconds to minutes & Minutes to hours & Insulin, glucagon, GH, ACTH,
ADH, PTH, prolactin \\
\textbf{Steroid} & Lipid-soluble & Synthesised on demand from
cholesterol (no storage); enzymatic cascades in mitochondria/SER & Bound
to carrier proteins (CBG, SHBG, albumin) --- about 5\% free and
bioactive & Nuclear receptors (intracellular) & Hours (transcription
required) & Hours to days & \hyperref[gl:cortisol]{\textbf{Cortisol}},
aldosterone, oestrogen, testosterone, vitamin D \\
\textbf{Amine --- Catecholamines} & Water-soluble & Tyrosine \texorpdfstring{\ensuremath{\to}}{->} DOPA \texorpdfstring{\ensuremath{\to}}{->}
dopamine \texorpdfstring{\ensuremath{\to}}{->} norepinephrine \texorpdfstring{\ensuremath{\to}}{->} epinephrine; stored in chromaffin granules &
Free in plasma; very short half-life & Plasma membrane (\texorpdfstring{\ensuremath{\alpha}}{alpha}/\texorpdfstring{\ensuremath{\beta}}{beta} adrenergic
GPCRs) & Seconds & Seconds to minutes & Epinephrine, norepinephrine,
dopamine \\
\textbf{Amine --- Thyroid hormones} & Lipid-soluble & Synthesised on
iodinated thyroglobulin scaffold in colloid; T4 prohormone converted to
T3 peripherally & 99.97\% bound to TBG, transthyretin, albumin & Nuclear
receptors (TR\texorpdfstring{\ensuremath{\alpha}}{alpha}, TR\texorpdfstring{\ensuremath{\beta}}{beta}) & Hours (transcription required) & Days (T4 t½
about 7 d) & T3, T4 \\
\textbf{Eicosanoids} & Lipid-soluble & Generated locally on demand from
arachidonic acid; not stored & Diffuse to neighbouring cells (paracrine)
& GPCRs and nuclear receptors (PPARs) & Seconds to minutes & Seconds to
minutes (rapid degradation) & PGE\(_2\), PGI\(_2\), TXA\(_2\),
leukotrienes \\
\end{longtable}
}

\paragraph{Comprehensive hormone table --- synthesis location,
transport, receptor,
half-life}\label{comprehensive-hormone-table-synthesis-location-transport-receptor-half-life}

{\def\LTcaptype{none} % do not increment counter
\begin{longtable}[]{@{}
  >{\raggedright\arraybackslash}p{(\linewidth - 12\tabcolsep) * \real{0.0897}}
  >{\raggedright\arraybackslash}p{(\linewidth - 12\tabcolsep) * \real{0.0641}}
  >{\raggedright\arraybackslash}p{(\linewidth - 12\tabcolsep) * \real{0.2308}}
  >{\raggedright\arraybackslash}p{(\linewidth - 12\tabcolsep) * \real{0.2051}}
  >{\raggedright\arraybackslash}p{(\linewidth - 12\tabcolsep) * \real{0.1026}}
  >{\raggedright\arraybackslash}p{(\linewidth - 12\tabcolsep) * \real{0.1923}}
  >{\raggedright\arraybackslash}p{(\linewidth - 12\tabcolsep) * \real{0.1154}}@{}}
\toprule\noalign{}
\begin{minipage}[b]{\linewidth}\raggedright
Hormone
\end{minipage} & \begin{minipage}[b]{\linewidth}\raggedright
Class
\end{minipage} & \begin{minipage}[b]{\linewidth}\raggedright
Synthesis location
\end{minipage} & \begin{minipage}[b]{\linewidth}\raggedright
Plasma transport
\end{minipage} & \begin{minipage}[b]{\linewidth}\raggedright
Receptor
\end{minipage} & \begin{minipage}[b]{\linewidth}\raggedright
Signal duration
\end{minipage} & \begin{minipage}[b]{\linewidth}\raggedright
Half-life
\end{minipage} \\
\midrule\noalign{}
\endhead
\bottomrule\noalign{}
\endlastfoot
\textbf{Insulin} & Peptide (51 aa) & \texorpdfstring{\ensuremath{\beta}}{beta}-cells, islets of Langerhans &
Free & IR (RTK, plasma membrane) & Minutes & about 5 min \\
\textbf{Glucagon} & Peptide (29 aa) & \texorpdfstring{\ensuremath{\alpha}}{alpha}-cells, islets of Langerhans &
Free & GcgR (G\(_s\) GPCR, hepatocyte) & Minutes & about 5
min \\
\textbf{Growth hormone (GH)} & Peptide (191 aa) & Anterior pituitary
somatotrophs & Free; some bound to GHBP & GHR (JAK2/STAT5 cytokine
receptor, liver) & Hours (via IGF-1) & about 20 min \\
\textbf{PTH} & Peptide (84 aa) & Parathyroid chief cells & Free & PTH1R
(G\(_s\)/G\(_q\) GPCR; bone, kidney) & Minutes & about 4 min \\
\textbf{ADH (vasopressin)} & Peptide (9 aa) & Hypothalamic SON/PVN;
stored in posterior pituitary & Free & V\(_2\)R (G\(_s\), kidney);
V\(_1\)R (G\(_q\), vessels) & Minutes & about 15 min \\
\textbf{Cortisol} & Steroid (C\(_{21}\)) & Adrenal cortex (zona
fasciculata) & 90\% CBG; 5\% albumin; 5\% free & GR (NR3C1,
cytoplasmic\texorpdfstring{\ensuremath{\to}}{->}nuclear) & Hours to days & about 60--90 min \\
\textbf{Aldosterone} & Steroid (C\(_{21}\)) & Adrenal cortex (zona
glomerulosa) & about 50\% bound to CBG/albumin & MR (NR3C2,
cytoplasmic\texorpdfstring{\ensuremath{\to}}{->}nuclear) & Hours & about 20 min \\
\textbf{Testosterone} & Steroid (C\(_{19}\)) & Leydig cells (testis);
zona reticularis (adrenal) & 60\% SHBG; 38\% albumin; 2\% free & AR
(NR3C4, cytoplasmic\texorpdfstring{\ensuremath{\to}}{->}nuclear) & Hours to days & about 70 min \\
\textbf{Oestradiol} & Steroid (C\(_{18}\)) & Ovarian granulosa cells
(aromatase from androgens) & SHBG, albumin; 1--2\% free & ER\texorpdfstring{\ensuremath{\alpha}}{alpha}/ER\texorpdfstring{\ensuremath{\beta}}{beta}
(nuclear); GPER (membrane) & Hours to days & about 13 h \\
\textbf{Progesterone} & Steroid (C\(_{21}\)) & Corpus luteum, placenta &
CBG, albumin & PR (NR3C3, nuclear) & Hours & about 5 min \\
\textbf{Vitamin D (1,25-(OH)\(_2\)D\(_3\))} & Steroid (secosteroid) &
Kidney (CYP27B1 hydroxylation of 25-OH-D) & DBP (vitamin D binding
protein) & VDR (nuclear) & Days & about 15 h \\
\textbf{Epinephrine} & Catecholamine & Adrenal medulla chromaffin cells
& Free & \texorpdfstring{\ensuremath{\alpha}}{alpha}/\texorpdfstring{\ensuremath{\beta}}{beta} adrenergic (GPCRs) & Seconds--minutes & about 2
min \\
\textbf{Norepinephrine} & Catecholamine & Adrenal medulla; sympathetic
neurons & Free & \texorpdfstring{\ensuremath{\alpha}}{alpha}/\texorpdfstring{\ensuremath{\beta}}{beta} adrenergic (GPCRs) & Seconds--minutes &
about 2 min \\
\textbf{T4 (thyroxine)} & Iodinated tyrosine & Thyroid follicular cells
& 99.97\% TBG/TTR/albumin; 0.03\% free & TR\texorpdfstring{\ensuremath{\alpha}}{alpha}/\texorpdfstring{\ensuremath{\beta}}{beta} (nuclear) after D2
conversion to T3 & Days--weeks & about 6--7 d \\
\textbf{T3 (triiodothyronine)} & Iodinated tyrosine & Thyroid
(about 20\%); peripheral D1/D2 (about 80\%) & 99.7\%
bound; 0.3\% free & TR\texorpdfstring{\ensuremath{\alpha}}{alpha}/\texorpdfstring{\ensuremath{\beta}}{beta} (nuclear) & Days & about 1 d \\
\textbf{Melatonin} & Indolamine & Pineal gland (from tryptophan via
serotonin) & Albumin & MT\(_1\), MT\(_2\) (G\(_i\) GPCRs) &
Minutes--hours & about 30--50 min \\
\textbf{PGE\(_2\)} & Eicosanoid & Most cells (COX-2 inducible) & Local
(paracrine) & EP\(_1\)--EP\(_4\) (GPCRs) & Seconds & about 30
s \\
\textbf{TXA\(_2\)} & Eicosanoid & Platelets (TXAS) & Local & TP receptor
(G\(_q\)) & Seconds & about 30 s \\
\textbf{Leptin} & Peptide (167 aa) & Adipocytes (proportional to fat
mass) & Free; some bound & LepR (JAK2/STAT3, hypothalamus) & Hours &
about 25 min \\
\textbf{Adiponectin} & Peptide (244 aa) & Adipocytes & Free;
trimers/hexamers/HMW & AdipoR1, AdipoR2 (AMPK pathway) & Hours &
about 14 h \\
\end{longtable}
}

\textbf{Key contrasts:}

\begin{itemize}
\tightlist
\item
  \textbf{Speed vs duration trade-off.} Membrane receptor signalling
  (peptides, catecholamines) is fast but transient: cAMP, IP\(_3\), and
  Ca\(^{2+}\) signals decay in seconds when ligand washes away. Nuclear
  receptor signalling (steroids, T3) is slow but persistent: transcribed
  mRNAs and translated proteins last hours to days.
\item
  \textbf{Storage vs synthesis.} Peptide hormones are pre-synthesised
  and stored in granules ready for rapid release --- insulin can be
  released within seconds of glucose elevation. Steroids cannot be
  stored (lipid-soluble, would diffuse away); they are synthesised on
  demand, which limits acute response speed.
\item
  \textbf{Carrier proteins as a reservoir.} CBG (cortisol binding
  globulin), SHBG (sex hormone binding globulin), and TBG (thyroxine
  binding globulin) bind 95--99\% of their target hormones. The bound
  fraction is biologically inactive but provides a circulating
  reservoir, buffering plasma levels and extending half-life. Free
  hormone is in equilibrium with bound; pregnancy elevates oestrogen
  which raises CBG, increasing total cortisol while free cortisol
  remains normal.
\end{itemize}

\subsubsection{Quantitative
Endocrinology}\label{quantitative-endocrinology}

\textbf{Hormone--receptor binding.} The dissociation constant \(K_d\)
characterises receptor affinity (lower \(K_d\) = higher affinity):

\begin{equation}
K_d = \frac{[\text{H}][\text{R}]}{[\text{HR}]}; \qquad \text{fractional occupancy} = \frac{[\text{H}]}{[\text{H}] + K_d}
\label{eq:unit_IX_kd}
\end{equation}

Typical \(K_d\) values: \hyperref[gl:insulin-receptor]{\textbf{insulin
receptor}} about 0.1 nM; glucocorticoid receptor
about 5 nM; epinephrine at \(\beta_2\) receptor
about 1 \texorpdfstring{\ensuremath{\mu}}{mu}M. The enormous affinity difference explains why
insulin circulates at picomolar concentrations while catecholamines
require nanomolar--micromolar levels for effect.

\textbf{Hormone half-life and clearance.} Plasma hormone concentration
decays exponentially after secretion ceases:

\begin{equation}
C(t) = C_0 \, e^{-kt}, \qquad k = \frac{\ln 2}{t_{1/2}}
\label{eq:unit_IX_halflife}
\end{equation}

Half-lives vary enormously: catecholamines about 1--2 min
(rapid enzymatic degradation by MAO/COMT); cortisol
about 60--90 min (hepatic metabolism); thyroid hormone (T4)
about 6--7 days (protein binding extends clearance).
Protein-bound hormones (e.g., TBG-bound T4, CBG-bound cortisol) are
protected from degradation and renal filtration.

\textbf{Dose--response relationship.} The sigmoidal (Hill)
dose--response curve relates hormone concentration to biological effect:

\begin{equation}
E = E_{\max} \cdot \frac{[\text{H}]^{n}}{EC_{50}^{n} + [\text{H}]^{n}}
\label{eq:unit_IX_doseresponse}
\end{equation}

where \(EC_{50}\) is the concentration producing 50\% of maximal effect
and \(n\) is the Hill coefficient (cooperativity). Spare receptors shift
the \(EC_{50}\) leftward from \(K_d\): insulin requires occupation of
about 5\% of receptors for maximal glucose uptake.

\textbf{\hyperref[gl:herd-immunity]{Herd immunity} threshold.} The
fraction of the population that must be immune to prevent epidemic
spread is:

\begin{equation}
p_c = 1 - \frac{1}{R_0}
\label{eq:unit_IX_herd}
\end{equation}

where \(R_0\) is the basic reproductive number. For measles
(\(R_0 \approx 12\)--18), \(p_c \approx 92\)--95\%, explaining why even
small drops in vaccination coverage trigger outbreaks.

\begin{center}\rule{0.5\linewidth}{0.5pt}\end{center}

\subsection{Hypothalamic-Pituitary
Axes}\label{hypothalamic-pituitary-axes}

The hypothalamus serves as the master integrator, linking the nervous
and endocrine systems. It produces releasing and inhibiting hormones
that control the anterior pituitary, which in turn controls target
glands throughout the body.

\textbf{Anterior pituitary hormones:} ACTH (corticotrophs), TSH
(thyrotrophs), GH (somatotrophs), LH and FSH (gonadotrophs), prolactin
(lactotrophs).

\textbf{Posterior pituitary:} Stores and releases ADH (vasopressin) and
oxytocin, which are synthesised in hypothalamic supraoptic and
paraventricular nuclei and transported down axons to the posterior
pituitary for release.

\begin{figure}[htbp]
\centering
\includegraphics[width=0.82\linewidth,height=4.2in,keepaspectratio]{../figures/mermaid_inline/inline_mermaid_0164_3472280e63b0.png}
\caption{The three major hypothalamic-pituitary axes Each follows the same hierarchical pattern: hypothalamic releasing hormone stimulates anterior pituitary tropic hormone, which stimulates target gland hormone production. Negative feedback at multiple levels prevents overstimulation. The HPG axis uniquely features positive feedback during the mid-cycle LH surge.}
\end{figure}

\subsubsection{HPA Axis (Hypothalamic-Pituitary-Adrenal) --- Detailed
Cortisol
Physiology}\label{hpa-axis-hypothalamic-pituitary-adrenal-detailed-cortisol-physiology}

\paragraph{Cortisol synthesis from cholesterol
(steroidogenesis)}\label{cortisol-synthesis-from-cholesterol-steroidogenesis}

Most adrenal steroids derive from cholesterol via a series of cytochrome
P450-mediated hydroxylations and oxidations. The rate-limiting step is
the transport of cholesterol from the outer to the inner mitochondrial
membrane, mediated by \textbf{StAR (steroidogenic acute regulatory
protein)} and stimulated by ACTH via cAMP/PKA.

\begin{figure}[htbp]
\centering
\includegraphics[width=0.82\linewidth,height=4.2in,keepaspectratio]{../figures/mermaid_inline/inline_mermaid_0165_9ce0bbb6ae5d.png}
\caption{Adrenal steroidogenesis The pathway diverges from common precursor pregnenolone into three parallel routes producing aldosterone (zona glomerulosa), cortisol (zona fasciculata), and androgens (zona reticularis). Each zone expresses zone-specific enzymes; for example, primarily the zona glomerulosa expresses aldosterone synthase (CYP11B2), and primarily the zona fasciculata expresses 11beta-hydroxylase (CYP11B1).}
\end{figure}

\paragraph{Glucocorticoid receptor (GR)
mechanism}\label{glucocorticoid-receptor-gr-mechanism}

Cortisol diffuses freely across the plasma membrane (lipid-soluble) and
binds the cytoplasmic GR (NR3C1). In the unliganded state, GR is held in
an inactive complex with HSP90, HSP70, and immunophilins. Cortisol
binding induces a conformational change, releasing the chaperones and
exposing nuclear localisation signals. The activated GR translocates to
the nucleus, dimerises, and binds \textbf{glucocorticoid response
elements (GREs)} in target gene promoters with the consensus sequence
5'-AGAACAnnnTGTTCT-3'.

GR has \textbf{two transcriptional modes:}

\begin{itemize}
\tightlist
\item
  \textbf{Transactivation} at positive GREs: directly upregulates
  gluconeogenic enzymes (PEPCK, G6Pase) and anti-inflammatory proteins
  (annexin A1, MKP-1, I\texorpdfstring{\ensuremath{\kappa}}{kappa}B\texorpdfstring{\ensuremath{\alpha}}{alpha}, IL-10).
\item
  \textbf{Transrepression} by tethering to NF-\texorpdfstring{\ensuremath{\kappa}}{kappa}B or AP-1: GR binds these
  pro-inflammatory transcription factors and prevents their activation
  of cytokine genes (IL-2, IL-6, TNF-\texorpdfstring{\ensuremath{\alpha}}{alpha}, COX-2). Transrepression mediates
  much of the anti-inflammatory action of glucocorticoid drugs.
\end{itemize}

The dual mechanism is therapeutically central. Most classical
glucocorticoids (prednisone, dexamethasone) drive both modes --- the
unwanted metabolic effects (hyperglycaemia, osteoporosis, muscle
wasting) come predominantly from transactivation, while the
anti-inflammatory benefit comes mostly from transrepression.
Pharmaceutical efforts to design ``selective GR agonists'' (SEGRAs) that
drive transrepression preferentially have been an enduring (but largely
unrealised) goal.

\paragraph{Cortisol effects (anti-inflammatory gene
programme)}\label{cortisol-effects-anti-inflammatory-gene-programme}

\begin{itemize}
\tightlist
\item
  \textbf{Gluconeogenesis} (transactivation of PEPCK, G6Pase via GR;
  mobilises amino acids from muscle protein catabolism)
\item
  \textbf{Immunosuppression} at multiple levels:

  \begin{itemize}
  \tightlist
  \item
    Direct: NF-\texorpdfstring{\ensuremath{\kappa}}{kappa}B inhibition; lymphocyte
    \hyperref[gl:apoptosis]{\textbf{apoptosis}}; suppression of IL-2
    transcription; mast-cell stabilisation
  \item
    Indirect: induction of IL-10, TGF-\texorpdfstring{\ensuremath{\beta}}{beta}, annexin A1; suppression of
    leukocyte trafficking by reducing E-selectin and ICAM-1
  \item
    Net cytokine effect: \textbf{down} TNF-\texorpdfstring{\ensuremath{\alpha}}{alpha}, IL-1\texorpdfstring{\ensuremath{\beta}}{beta}, IL-2, IL-6, IFN-\texorpdfstring{\ensuremath{\gamma}}{gamma};
    \textbf{up} IL-10
  \end{itemize}
\item
  \textbf{Lipolysis} in peripheral adipose tissue (with central fat
  redistribution from chronic excess)
\item
  \textbf{Anti-inflammatory} (inhibits COX-2, phospholipase A\(_2\) via
  annexin A1 induction; reduces histamine, leukotriene production)
\item
  \textbf{Permissive effects} on catecholamine action (\texorpdfstring{\ensuremath{\beta}}{beta}-adrenergic
  receptor expression maintenance)
\item
  \textbf{Bone resorption} (chronic excess: osteoporosis via direct
  osteoblast suppression and indirect calcium handling)
\item
  \textbf{CNS effects} (mood --- chronic excess: depression, psychosis,
  cognitive impairment via hippocampal atrophy)
\end{itemize}

\paragraph{Circadian rhythm and
feedback}\label{circadian-rhythm-and-feedback}

Cortisol secretion follows a robust \textbf{circadian rhythm} driven by
the suprachiasmatic nucleus (SCN): peak at approximately 08:00 (cortisol
awakening response, CAR), nadir at midnight. Superimposed are
\textbf{ultradian pulses} every 60--90 minutes. The pulsatile pattern is
essential for proper GR function --- continuous (non-pulsatile) cortisol
exposure desensitises target tissues.

Schematically, plasma cortisol \(\approx 18\,\mu\text{g/dL}\) at 08:00,
drops through the day to \(\approx 8\,\mu\text{g/dL}\) at 16:00, and is
\(<5\,\mu\text{g/dL}\) at midnight. Disrupted circadian rhythm (shift
work, depression, Cushing's) loses the diurnal variation;
\textbf{midnight salivary cortisol} is one of the most sensitive
screening tests for endogenous Cushing's syndrome.

\textbf{Negative feedback} operates at two levels:

\begin{itemize}
\tightlist
\item
  \textbf{Fast feedback} (minutes): non-genomic; cortisol acts on
  hypothalamic CRH neurons via membrane GRs and presynaptic
  endocannabinoid release to suppress CRH secretion within minutes.
\item
  \textbf{Slow feedback} (hours): genomic; cortisol downregulates POMC
  gene transcription (the ACTH precursor) in pituitary corticotrophs and
  CRH transcription in hypothalamus.
\end{itemize}

\paragraph{Stress vs basal regulation}\label{stress-vs-basal-regulation}

Under \textbf{basal (non-stress)} conditions, the HPA axis runs on its
circadian/ultradian rhythm with tight negative feedback. The day's
cortisol output is roughly 10--20 mg/24 h.

Under \textbf{acute stress} (trauma, hypoglycaemia, infection, surgery),
hypothalamic PVN neurons receive convergent inputs (limbic, brainstem,
circulating cytokines, baroreceptors) and override basal restraint:

\begin{itemize}
\tightlist
\item
  CRH secretion rises 5--10 fold
\item
  ACTH plasma concentration rises within 5--15 min
\item
  Cortisol peaks about 30 min after stress onset
\item
  Cortisol output can increase 3--10 fold over basal (50--100 mg/24 h in
  major surgery)
\item
  Negative feedback is partially suppressed, allowing sustained
  elevation
\end{itemize}

Chronic psychosocial stress engages the same axis but with
\textbf{glucocorticoid receptor downregulation} and \textbf{flattened
diurnal rhythm} --- features now linked to metabolic syndrome,
depression, hippocampal atrophy, and accelerated cognitive decline.

\textbf{Pathology:}

\begin{itemize}
\tightlist
\item
  \textbf{Cushing's syndrome:} Cortisol excess. Causes: pituitary
  adenoma (Cushing's disease, 70\%), ectopic ACTH (lung small cell
  carcinoma), adrenal adenoma, iatrogenic (chronic glucocorticoid
  therapy). Features: central obesity, moon face, purple striae,
  hyperglycaemia, osteoporosis, immunosuppression, hypertension
  (cortisol cross-reacts with mineralocorticoid receptor).
\item
  \textbf{Addison's disease:} Primary adrenal insufficiency (cortisol
  deficiency). Most common cause: autoimmune adrenalitis. Features:
  hypotension, hyperpigmentation (elevated ACTH drives
  melanocyte-stimulating hormone activity, since both are cleaved from
  the same precursor POMC), fatigue, salt craving. Can cause
  life-threatening adrenal crisis.
\end{itemize}

\begin{quote}
\textbf{Clinical Connection:} The dexamethasone suppression test
exploits negative feedback to diagnose Cushing's syndrome. Dexamethasone
is a synthetic glucocorticoid that normally suppresses CRH and ACTH
secretion. In Cushing's disease (pituitary adenoma), low-dose
dexamethasone fails to suppress cortisol, but high-dose dexamethasone
does. In ectopic ACTH production, neither dose suppresses cortisol. This
distinguishes the source of excess cortisol.
\end{quote}

\subsubsection{HPT Axis (Hypothalamic-Pituitary-Thyroid) --- Detailed
Mechanism}\label{hpt-axis-hypothalamic-pituitary-thyroid-detailed-mechanism}

\paragraph{Thyroid hormone synthesis}\label{thyroid-hormone-synthesis}

Follicular cells of the thyroid gland surround a colloid-filled follicle
and execute a multi-step iodination process:

\begin{enumerate}
\def\labelenumi{\arabic{enumi}.}
\tightlist
\item
  \textbf{Iodide uptake.} The basolateral Na\(^+\)/I\(^-\) symporter
  (NIS) actively concentrates iodide about 25--100-fold above
  plasma. Pertechnetate (\(^{99m}\)TcO\(_4^-\)) is also transported by
  NIS and is used in thyroid scintigraphy.
\item
  \textbf{Thyroglobulin synthesis.} Follicular cells produce
  thyroglobulin (Tg, about 660 kDa glycoprotein with
  about 120 tyrosine residues), packaged into vesicles and
  exocytosed into the colloid lumen.
\item
  \textbf{Iodide oxidation and tyrosine iodination.} At the apical
  membrane, \textbf{thyroid peroxidase (TPO)} oxidises I\(^-\) to a
  reactive iodinating species (I\(_2\) or I\(^+\)) using H\(_2\)O\(_2\)
  as cofactor. Tyrosines on Tg are iodinated to form monoiodotyrosine
  (MIT) and diiodotyrosine (DIT).
\item
  \textbf{Coupling reaction.} Still on the Tg scaffold, TPO catalyses
  oxidative coupling: DIT + DIT \texorpdfstring{\ensuremath{\to}}{->} T4 (thyroxine, 4 iodines); MIT + DIT \texorpdfstring{\ensuremath{\to}}{->}
  T3 (3 iodines).
\item
  \textbf{Endocytosis and proteolysis.} TSH stimulation triggers
  endocytosis of Tg into the follicular cell, where lysosomal proteases
  cleave the iodothyronines from Tg. T4 and T3 are released into the
  bloodstream.
\item
  \textbf{Recycling.} MIT and DIT not coupled into T3/T4 are deiodinated
  by intracellular dehalogenase, recycling iodide.
\end{enumerate}

\paragraph{Wolff-Chaikoff effect}\label{wolff-chaikoff-effect}

Acutely high iodide (e.g., from radiocontrast, amiodarone, or seaweed
binge) paradoxically \textbf{inhibits} thyroid hormone synthesis: excess
iodide downregulates NIS expression and impairs TPO-mediated
organification. Most healthy individuals ``escape'' within 7--10 days as
NIS resets. In patients with autoimmune thyroid disease, escape may fail
--- producing \textbf{iodine-induced hypothyroidism}. Conversely, in
nodular goitres with autonomous follicles, iodine load can drive
\textbf{iodine-induced hyperthyroidism (Jod-Basedow)}. The
Wolff-Chaikoff effect is exploited therapeutically: \textbf{potassium
iodide} (Lugol's solution, SSKI) is given pre-operatively in Graves'
disease to reduce thyroid vascularity and slow hormone release before
thyroidectomy.

\paragraph{T4 \texorpdfstring{\ensuremath{\to}}{->} T3 conversion
(deiodinases)}\label{t4-t3-conversion-deiodinases}

T4 is the predominant secreted form (about 80\%) but is largely
a \textbf{prohormone}. Three iodothyronine deiodinases activate or
inactivate T4 in target tissues:

{\def\LTcaptype{none} % do not increment counter
\begin{longtable}[]{@{}
  >{\raggedright\arraybackslash}p{(\linewidth - 6\tabcolsep) * \real{0.3125}}
  >{\raggedright\arraybackslash}p{(\linewidth - 6\tabcolsep) * \real{0.2500}}
  >{\raggedright\arraybackslash}p{(\linewidth - 6\tabcolsep) * \real{0.1875}}
  >{\raggedright\arraybackslash}p{(\linewidth - 6\tabcolsep) * \real{0.2500}}@{}}
\toprule\noalign{}
\begin{minipage}[b]{\linewidth}\raggedright
Deiodinase
\end{minipage} & \begin{minipage}[b]{\linewidth}\raggedright
Reaction
\end{minipage} & \begin{minipage}[b]{\linewidth}\raggedright
Tissue
\end{minipage} & \begin{minipage}[b]{\linewidth}\raggedright
Function
\end{minipage} \\
\midrule\noalign{}
\endhead
\bottomrule\noalign{}
\endlastfoot
\textbf{D1} & T4 \texorpdfstring{\ensuremath{\to}}{->} T3 (5'-deiodination) & Liver, kidney, thyroid & Bulk
plasma T3 production \\
\textbf{D2} & T4 \texorpdfstring{\ensuremath{\to}}{->} T3 (5'-deiodination) & Brain, pituitary, brown
adipose & Local T3 generation; pituitary feedback set-point \\
\textbf{D3} & T4 \texorpdfstring{\ensuremath{\to}}{->} reverse T3 (rT3); T3 \texorpdfstring{\ensuremath{\to}}{->} T2 & Placenta, fetal tissues &
Inactivation; protects fetus from maternal T3 \\
\end{longtable}
}

T3 is \textbf{about 4\texorpdfstring{\ensuremath{\times}}{x} more potent} than T4 at the nuclear
receptor. During fasting, illness, or stress (non-thyroidal illness), D1
activity drops and D3 activity rises, shifting balance toward inactive
rT3 --- adaptive metabolic slowing.

\paragraph{Nuclear thyroid receptor
mechanism}\label{nuclear-thyroid-receptor-mechanism}

T3 enters cells via MCT8/MCT10 transporters and binds nuclear
\textbf{thyroid hormone receptors TR\texorpdfstring{\ensuremath{\alpha}}{alpha} (cardiac, neural)} and \textbf{TR\texorpdfstring{\ensuremath{\beta}}{beta}
(liver, pituitary)}. TRs heterodimerise with \textbf{RXR (retinoid X
receptor)} and bind \textbf{thyroid response elements (TREs)} in target
gene promoters --- even in the absence of T3.

\begin{itemize}
\tightlist
\item
  \textbf{Without T3:} Unliganded TR/RXR recruits corepressors (NCoR,
  SMRT, HDACs) and \textbf{represses} target gene transcription.
\item
  \textbf{With T3:} Conformational change releases corepressors and
  recruits coactivators (SRC-1, p300/CBP, histone acetyltransferases) \texorpdfstring{\ensuremath{\to}}{->}
  \textbf{active transcription}.
\end{itemize}

This dual mechanism explains why hypothyroidism causes such pronounced
symptoms: target genes are not merely ``not turned on'' but are actively
repressed below baseline.

\paragraph{Effects on metabolism}\label{effects-on-metabolism}

\begin{itemize}
\tightlist
\item
  \textbf{Increased basal metabolic rate (BMR)} through upregulation of
  Na\(^+\)/K\(^+\)-ATPase (each cell uses about 25\% of ATP for
  this pump; T3 increases pump density).
\item
  \textbf{Mitochondrial uncoupling} via UCP-mediated proton leak (heat
  production, calorigenic effect).
\item
  \textbf{Increased cardiac output} (\(\beta_1\)-adrenergic receptor
  upregulation; positive chronotropic and inotropic effects).
\item
  \textbf{Lipolysis and cholesterol degradation} (hyperthyroidism lowers
  LDL; hypothyroidism causes hypercholesterolaemia).
\item
  \textbf{Neurodevelopment:} essential for myelination, synaptogenesis,
  neuronal migration in fetus and infant. Severe untreated congenital
  hypothyroidism produces cretinism (intellectual disability, growth
  retardation) --- preventable by neonatal screening.
\end{itemize}

\textbf{Pathology:}

\begin{itemize}
\tightlist
\item
  \textbf{Hypothyroidism:} Low T3/T4, elevated TSH. Hashimoto's
  thyroiditis (autoimmune; anti-TPO and anti-thyroglobulin antibodies).
  Features: fatigue, weight gain, cold intolerance, bradycardia,
  constipation, myxoedema.
\item
  \textbf{Hyperthyroidism:} High T3/T4, suppressed TSH. Graves' disease
  (autoimmune; TSH receptor-stimulating antibodies ---
  thyroid-stimulating immunoglobulins, TSI). Features: weight loss, heat
  intolerance, tachycardia/atrial fibrillation, tremor, exophthalmos
  (due to retro-orbital inflammation), diffuse goitre.
\item
  \textbf{Iodine deficiency:} Most common preventable cause of
  intellectual disability worldwide; affects approximately 1.9 billion
  people. Addressed by salt iodisation programmes.
\end{itemize}

\subsubsection{HPG Axis and Reproductive
Hormones}\label{hpg-axis-and-reproductive-hormones}

\textbf{GnRH pulsatility is critical:} High-frequency pulses (every
60--90 min) favour LH secretion; low-frequency pulses (every 2--4 h)
favour FSH. Continuous GnRH paradoxically suppresses the axis
(downregulates GnRH receptors), which is the basis of GnRH agonist
therapy for prostate cancer, endometriosis, and precocious puberty.

\textbf{Male reproductive endocrinology:} LH stimulates Leydig cells
(testosterone production). FSH + testosterone stimulate Sertoli cells
(support spermatogenesis; produce inhibin B, which provides negative
feedback on FSH). Testosterone effects: muscle mass, bone density,
secondary sexual characteristics, spermatogenesis, libido.

\textbf{Female menstrual cycle (28-day average):}

\begin{itemize}
\tightlist
\item
  \textbf{Follicular phase} (days 1--14): FSH promotes follicular
  growth. \hyperref[gl:dominant]{\textbf{Dominant}} follicle produces
  rising oestradiol. Oestradiol initially provides negative feedback.
\item
  \textbf{Ovulation} (day approximately 14): \textbf{Positive feedback}
  --- sustained high oestradiol above a threshold for \textgreater36 h
  switches from negative to positive feedback, triggering an LH surge.
  The LH surge induces ovulation (follicular rupture, oocyte release).
\item
  \textbf{Luteal phase} (days 14--28): Corpus luteum (remnant of
  ovulated follicle) produces progesterone + oestradiol. Progesterone
  prepares the endometrium for implantation (secretory phase). If no
  implantation occurs, corpus luteum degenerates (luteolysis),
  progesterone and oestradiol fall, endometrium sheds (menstruation).
\end{itemize}

\textbf{Development and reproduction as endocrine timing problems:}
Fertilisation, implantation, placentation, fetal growth, birth,
lactation, puberty, and reproductive senescence depend on timed
endocrine signals interacting with local tissue cues. Early embryos are
initially regulated by maternal transcripts and local morphogen
gradients; later development adds fetal-placental endocrine exchange,
thyroid-hormone-dependent neurodevelopment, glucocorticoid-driven lung
maturation, and sex-steroid-dependent reproductive tract
differentiation. Clinically, infertility is diagnosed after 12 months of
regular unprotected intercourse without pregnancy (or after 6 months
when the female partner is 35 years or older), because reproductive
physiology depends on ovulation, sperm production, tubal transport,
implantation, uterine receptivity, endocrine timing, and age-dependent
gamete quality \citep{cdc2024reproductivehealth}. The organismal point
is that reproduction is not one organ system; it is a coordinated
life-history transition linking gonads, hypothalamus, pituitary,
placenta, metabolism, immune tolerance, and behaviour.

\begin{quote}
\textbf{Concept Check 1:} Why does continuous GnRH administration
paradoxically suppress the HPG axis rather than stimulate it? How is
this exploited clinically in prostate cancer treatment, where the goal
is to lower testosterone?
\end{quote}

\begin{center}\rule{0.5\linewidth}{0.5pt}\end{center}

\subsection{Pancreatic Hormones, Glucose Homeostasis, and Energy
Balance}\label{pancreatic-hormones-glucose-homeostasis-and-energy-balance}

\textbf{Normal fasting blood glucose:} 4.0--5.5 mmol/L (72--99 mg/dL).
Post-prandial peak: \textless7.8 mmol/L (\textless140 mg/dL).

\subsubsection{Insulin signalling --- molecular
detail}\label{insulin-signalling-molecular-detail}

\textbf{After a meal (high glucose):}

\begin{enumerate}
\def\labelenumi{\arabic{enumi}.}
\tightlist
\item
  Glucose enters \texorpdfstring{\ensuremath{\beta}}{beta}-cells via GLUT2 transporter (high-K\(_m\) ``glucose
  sensor'')
\item
  Glucose metabolism increases the ATP/ADP ratio
\item
  ATP-sensitive K\(^+\) channels (K\(_{ATP}\), SUR1/Kir6.2 subunits)
  close \texorpdfstring{\ensuremath{\to}}{->} membrane depolarisation
\item
  Voltage-gated L-type Ca\(^{2+}\) channels open \texorpdfstring{\ensuremath{\to}}{->} Ca\(^{2+}\) influx
\item
  Ca\(^{2+}\) triggers \textbf{insulin exocytosis} from dense-core
  granules (mature insulin = A and B chains held by disulphide bonds;
  C-peptide co-released as marker)
\item
  \textbf{Insulin signalling in target cells:}

  \begin{itemize}
  \tightlist
  \item
    Insulin binds \textbf{insulin receptor (IR)} --- an
    \(\alpha_2\beta_2\) tetrameric receptor tyrosine kinase. Each
    \texorpdfstring{\ensuremath{\alpha}}{alpha}-subunit is extracellular and binds insulin; \texorpdfstring{\ensuremath{\beta}}{beta}-subunits span the
    membrane and contain intracellular tyrosine kinase domains.
  \item
    Binding causes conformational change \texorpdfstring{\ensuremath{\to}}{->}
    \textbf{trans-autophosphorylation} of \texorpdfstring{\ensuremath{\beta}}{beta}-subunit tyrosine residues
    (Y1158, Y1162, Y1163 in the activation loop)
  \item
    Phosphotyrosines recruit \textbf{IRS1/IRS2} scaffold proteins \texorpdfstring{\ensuremath{\to}}{->} IRS
    phosphorylation
  \item
    IRS phosphotyrosines recruit \textbf{PI3K} (p85 regulatory + p110
    catalytic) via SH2 domains
  \item
    PI3K converts PIP\(_2\) \texorpdfstring{\ensuremath{\to}}{->} \textbf{PIP\(_3\)}; PDK1 and mTORC2 dock
    on PIP\(_3\) and activate \textbf{Akt (PKB)} by phosphorylation at
    T308 and S473
  \item
    \textbf{Akt phosphorylates AS160 (TBC1D4)}, releasing Rab-GAP
    inhibition of GLUT4-vesicle exocytosis \texorpdfstring{\ensuremath{\to}}{->}
    \textbf{\hyperref[gl:glut4]{GLUT4} translocation to plasma membrane}
    in muscle and adipose (the dominant mechanism for postprandial
    glucose disposal)
  \item
    \textbf{Akt inhibits GSK3\texorpdfstring{\ensuremath{\beta}}{beta}} \texorpdfstring{\ensuremath{\to}}{->} glycogen synthase activation \texorpdfstring{\ensuremath{\to}}{->}
    \textbf{glycogen synthesis}
  \item
    \textbf{Akt activates mTORC1} \texorpdfstring{\ensuremath{\to}}{->} ribosomal S6K1 and 4E-BP1
    phosphorylation \texorpdfstring{\ensuremath{\to}}{->} \textbf{protein synthesis}
  \item
    \textbf{Akt phosphorylates FoxO1} \texorpdfstring{\ensuremath{\to}}{->} nuclear export \texorpdfstring{\ensuremath{\to}}{->} represses
    gluconeogenic gene transcription (PEPCK, G6Pase)
  \end{itemize}
\end{enumerate}

\begin{figure}[htbp]
\centering
\includegraphics[width=0.82\linewidth,height=4.2in,keepaspectratio]{../figures/mermaid_inline/inline_mermaid_0166_a076687f9cc5.png}
\caption{Insulin signalling cascade Insulin binding to IR triggers trans-autophosphorylation, recruitment of IRS, activation of PI3K, generation of PIP\$\_3\$, and activation of Akt. Akt then drives four parallel branches: GLUT4 translocation, glycogen synthesis, protein synthesis, and suppression of gluconeogenic gene transcription.}
\end{figure}

\paragraph{Glucagon signalling}\label{glucagon-signalling}

\textbf{During fasting (low glucose):}

\begin{enumerate}
\def\labelenumi{\arabic{enumi}.}
\tightlist
\item
  \texorpdfstring{\ensuremath{\alpha}}{alpha}-cells release \textbf{glucagon} (triggered by low glucose,
  sympathetic activation, amino acids)
\item
  Glucagon binds G\(_s\)-coupled receptor on hepatocytes
\item
  cAMP-PKA pathway: phosphorylase kinase activates \textbf{glycogen
  phosphorylase} for \textbf{glycogenolysis}
\item
  cAMP also activates \textbf{CREB} \texorpdfstring{\ensuremath{\to}}{->} PEPCK and G6Pase gene
  transcription \texorpdfstring{\ensuremath{\to}}{->} \textbf{gluconeogenesis}
\item
  PKA phosphorylates and inhibits PFK-2/FBPase-2 (PFKFB1), lowering
  F2,6BP \texorpdfstring{\ensuremath{\to}}{->} favours fructose-1,6-bisphosphatase over PFK-1 \texorpdfstring{\ensuremath{\to}}{->}
  gluconeogenesis dominates
\end{enumerate}

\paragraph{Crosstalk with leptin, adiponectin, and
GLP-1}\label{crosstalk-with-leptin-adiponectin-and-glp-1}

\textbf{Leptin} is secreted by adipocytes in proportion to fat mass and
acts on hypothalamic neurons (arcuate nucleus POMC and AgRP/NPY neurons)
via the leptin receptor (LepR, JAK2/STAT3 signalling) to suppress
appetite and increase energy expenditure. Leptin enhances central
insulin sensitivity and provides a long-term signal of energy stores.
Most obese individuals have high leptin but show \textbf{leptin
resistance} --- impaired hypothalamic LepR signalling and JAK2/STAT3
attenuation, partly via SOCS3 induction.

\textbf{Adiponectin} is also secreted by adipocytes --- but
\textbf{inversely} with fat mass (higher in lean states). It acts via
AdipoR1 and AdipoR2, activating \textbf{AMPK} (AMP-activated protein
kinase, the cellular ``energy sensor''). AMPK activation:

\begin{itemize}
\tightlist
\item
  Increases fatty acid oxidation (phosphorylates ACC, decreasing
  malonyl-CoA, releasing CPT-1 from inhibition)
\item
  Decreases hepatic gluconeogenesis (decreases CREB-regulated
  transcription)
\item
  Increases skeletal muscle glucose uptake (independent of insulin)
\item
  Decreases mTORC1 and protein synthesis (reciprocal to insulin/Akt)
\end{itemize}

Adiponectin therefore acts as an \textbf{insulin sensitiser}. Plasma
adiponectin is reduced in T2DM, obesity, and metabolic syndrome;
pioglitazone (a thiazolidinedione, PPAR\texorpdfstring{\ensuremath{\gamma}}{gamma} agonist) raises adiponectin and
improves insulin sensitivity.

\textbf{GLP-1 (glucagon-like peptide 1)} is an \textbf{incretin}
released from intestinal L-cells in response to nutrient ingestion. It
acts via GLP-1R (G\(_s\)-coupled GPCR) to:

\begin{itemize}
\tightlist
\item
  Enhance glucose-dependent insulin secretion (primarily when glucose
  elevated \texorpdfstring{\ensuremath{\to}}{->} low hypoglycaemia risk)
\item
  Suppress glucagon release from \texorpdfstring{\ensuremath{\alpha}}{alpha}-cells
\item
  Slow gastric emptying
\item
  Activate hypothalamic anorexigenic circuits \texorpdfstring{\ensuremath{\to}}{->} reduce food intake
\item
  Promote \texorpdfstring{\ensuremath{\beta}}{beta}-cell survival and proliferation (animal studies)
\end{itemize}

GLP-1 has a very short half-life (about 2 min) due to
degradation by \textbf{DPP-4} (dipeptidyl peptidase-4). Drug strategies:
GLP-1 analogues with DPP-4--resistant modifications (semaglutide,
liraglutide), or DPP-4 inhibitors (sitagliptin) that prolong endogenous
GLP-1.

\subsubsection{Diabetes Mellitus}\label{diabetes-mellitus}

\textbf{Type 1 diabetes mellitus (T1DM):}

\begin{itemize}
\tightlist
\item
  Autoimmune destruction of \texorpdfstring{\ensuremath{\beta}}{beta}-cells (anti-GAD65, anti-IA-2,
  anti-insulin, anti-ZnT8 antibodies)
\item
  Complete insulin deficiency
\item
  Onset usually in childhood/adolescence (but can occur at any age ---
  ``LADA'' in adults)
\item
  Without insulin: hyperglycaemia, ketoacidosis (lipolysis produces FFA,
  hepatic \texorpdfstring{\ensuremath{\beta}}{beta}-oxidation produces ketone bodies, metabolic acidosis)
\item
  Treatment: exogenous insulin (basal-bolus regimen), increasingly
  closed-loop insulin pump systems
\end{itemize}

\textbf{Type 2 diabetes mellitus (T2DM):}

\begin{itemize}
\tightlist
\item
  Peripheral insulin resistance: ceramide and pro-inflammatory cytokines
  (TNF-\texorpdfstring{\ensuremath{\alpha}}{alpha}, IL-6 from visceral adipose tissue) cause serine
  phosphorylation of IRS1, blocking normal tyrosine phosphorylation. ER
  stress and mitochondrial dysfunction also contribute.
\item
  Compensatory \texorpdfstring{\ensuremath{\beta}}{beta}-cell hyperinsulinaemia initially maintains
  normoglycaemia
\item
  Progressive \texorpdfstring{\ensuremath{\beta}}{beta}-cell failure (glucotoxicity, lipotoxicity, amyloid
  deposition) leads to overt hyperglycaemia
\item
  Treatment: lifestyle modification + metformin (AMPK activation,
  reduced hepatic glucose production) + GLP-1 receptor agonists
  (semaglutide) + SGLT2 inhibitors (empagliflozin: blocks renal glucose
  reabsorption)
\end{itemize}

\begin{quote}
\textbf{Clinical Connection:} Semaglutide (Ozempic/Wegovy) has
transformed T2DM and obesity management. Clinical trials show 15--20\%
body weight reduction with semaglutide, plus cardiovascular and renal
benefits, by mimicking natural GLP-1 signalling at hypothalamic appetite
centres and pancreatic \texorpdfstring{\ensuremath{\beta}}{beta}-cells.
\end{quote}

\begin{quote}
\textbf{Concept Check 2:} A type-1 diabetic receives a long-acting
insulin analogue (glargine) once daily. Blood glucose is well controlled
during the day but the patient develops reactive hyperglycaemia every
morning (``dawn phenomenon''). Given that cortisol, GH, and glucagon
most rise in the pre-waking hours, explain \emph{which}
counter-regulatory hormones are responsible for the morning rise and
\emph{why} glargine's 24-hour profile is insufficient to cover it. What
feature of a more modern analogue (e.g.~degludec, with
about 42-h half-life) addresses this?
\end{quote}

\subsubsection{Adrenal Gland --- Zonal
Architecture}\label{adrenal-gland-zonal-architecture}

\textbf{Adrenal cortex} (mesodermal origin, steroid hormones):

\begin{itemize}
\tightlist
\item
  \textbf{Zona glomerulosa:} Aldosterone (mineralocorticoid). Regulated
  by RAAS and K\(^+\). Promotes Na\(^+\) reabsorption, K\(^+\) secretion
  in collecting duct.
\item
  \textbf{Zona fasciculata:} Cortisol (glucocorticoid). Regulated by
  ACTH.
\item
  \textbf{Zona reticularis:} DHEA, androstenedione (adrenal androgens).
  Puberty (adrenarche).
\end{itemize}

\textbf{Adrenal medulla} (neural crest origin, modified sympathetic
ganglion):

\begin{itemize}
\tightlist
\item
  Chromaffin cells release epinephrine (80\%) and norepinephrine (20\%)
  directly into blood
\item
  Epinephrine effects: heart rate increase (\(\beta_1\)),
  bronchodilation (\(\beta_2\)), glycogenolysis (\(\beta_2\) in liver),
  lipolysis (\(\beta_3\) in adipose), vasoconstriction (\(\alpha_1\)) in
  skin/viscera, vasodilation (\(\beta_2\)) in skeletal muscle
\item
  Duration: seconds to minutes (rapid metabolic clearance by MAO and
  COMT)
\end{itemize}

\subsubsection{Growth Hormone Axis and
IGF-1}\label{growth-hormone-axis-and-igf-1}

\begin{figure}[htbp]
\centering
\includegraphics[width=0.82\linewidth,height=4.2in,keepaspectratio]{../figures/mermaid_inline/inline_mermaid_0167_3c9c6ef56171.png}
\caption{Growth hormone axis GHRH stimulates and somatostatin inhibits pituitary GH release. GH acts directly on tissues (lipolysis, insulin antagonism - "diabetogenic") and indirectly via hepatic IGF-1, which mediates linear growth and protein synthesis. IGF-1 provides negative feedback at hypothalamus (somatostatin) and pituitary (GH transcription).}
\end{figure}

GH from anterior pituitary somatotrophs: pulsatile release (peaks during
deep sleep and exercise). GH binds the GH receptor (a JAK2-coupled
cytokine receptor) on hepatocytes, activating STAT5, which drives
\textbf{IGF-1 (insulin-like growth factor 1)} transcription and
secretion. Circulating IGF-1 is bound to IGFBP-3 (extends half-life from
minutes to hours).

\textbf{GH/IGF-1 effects:}

\begin{itemize}
\tightlist
\item
  Linear bone growth (epiphyseal plate chondrogenesis driven by IGF-1)
\item
  Protein synthesis (positive nitrogen balance)
\item
  Lipolysis (direct GH effect; ``diabetogenic'')
\item
  Reduced peripheral glucose uptake (GH antagonises insulin)
\end{itemize}

\textbf{Pathology:}

\begin{itemize}
\tightlist
\item
  \textbf{Acromegaly:} Excess GH in adults (pituitary adenoma).
  Characteristic features: enlarged hands, feet, jaw; organomegaly;
  impaired glucose tolerance. Diagnosed by failure to suppress GH during
  oral glucose tolerance test.
\item
  \textbf{Gigantism:} Excess GH before epiphyseal plate closure
  (childhood). Tall stature.
\item
  \textbf{GH deficiency in children:} Short stature; treated with
  recombinant human GH (somatropin).
\item
  \textbf{Laron syndrome:} GH receptor mutation; high GH but no IGF-1
  response. Short stature; remarkably, low IGF-1 confers resistance to
  cancer and diabetes.
\end{itemize}

\begin{center}\rule{0.5\linewidth}{0.5pt}\end{center}

\subsection{Prostaglandins and
Eicosanoids}\label{prostaglandins-and-eicosanoids}

\textbf{Eicosanoids} are 20-carbon paracrine signalling lipids derived
from membrane phospholipids. Unlike conventional hormones, they are not
stored, are synthesised on demand, act locally (paracrine/autocrine),
and are rapidly inactivated.

\subsubsection{Synthesis from arachidonic
acid}\label{synthesis-from-arachidonic-acid}

\begin{figure}[htbp]
\centering
\includegraphics[width=0.82\linewidth,height=4.2in,keepaspectratio]{../figures/mermaid_inline/inline_mermaid_0168_2ed715e1da8a.png}
\caption{Eicosanoid biosynthesis pathways Phospholipase A\$\_2\$ liberates arachidonic acid from membrane phospholipids. Three branches generate distinct families: cyclooxygenase (COX) to prostaglandins and thromboxanes; 5-lipoxygenase to leukotrienes; cytochrome P450 to epoxyeicosatrienoic acids. Glucocorticoids inhibit at the PLA\$\_2\$ step; NSAIDs inhibit at COX.}
\end{figure}

\subsubsection{COX-1 vs COX-2 --- distinct
physiology}\label{cox-1-vs-cox-2-distinct-physiology}

{\def\LTcaptype{none} % do not increment counter
\begin{longtable}[]{@{}
  >{\raggedright\arraybackslash}p{(\linewidth - 4\tabcolsep) * \real{0.2800}}
  >{\raggedright\arraybackslash}p{(\linewidth - 4\tabcolsep) * \real{0.3600}}
  >{\raggedright\arraybackslash}p{(\linewidth - 4\tabcolsep) * \real{0.3600}}@{}}
\toprule\noalign{}
\begin{minipage}[b]{\linewidth}\raggedright
Feature
\end{minipage} & \begin{minipage}[b]{\linewidth}\raggedright
\textbf{COX-1}
\end{minipage} & \begin{minipage}[b]{\linewidth}\raggedright
\textbf{COX-2}
\end{minipage} \\
\midrule\noalign{}
\endhead
\bottomrule\noalign{}
\endlastfoot
Expression & Constitutive (most tissues) & Inducible (inflammation,
growth factors); constitutive in kidney, brain, vascular endothelium \\
Function & Gastric mucosa protection (PGE\(_2\), PGI\(_2\)); platelet
TXA\(_2\); renal autoregulation & Inflammatory PGE\(_2\)/PGI\(_2\);
pain, fever; renal salt/water handling \\
Knockout phenotype & Gastric ulcers; reduced platelet aggregation &
Reduced inflammation; renal abnormalities; fertility defects \\
Selective inhibitor & (none in clinical use) & Celecoxib, etoricoxib,
parecoxib \\
\end{longtable}
}

\textbf{Selective COX-2 inhibitors (coxibs)} were developed to spare
COX-1 (preserving gastric prostaglandins and reducing GI bleeding).
Initial successes (Vioxx/rofecoxib, Bextra/valdecoxib) were tempered by
cardiovascular concerns: selective COX-2 inhibition reduces endothelial
PGI\(_2\) (antiplatelet, vasodilator) without reducing platelet
TXA\(_2\) (made by COX-1) --- shifting the haemostatic balance toward
thrombosis. Rofecoxib was withdrawn in 2004; celecoxib remains in use
with cardiovascular labelling.

\subsubsection{Pharmacological targets}\label{pharmacological-targets}

{\def\LTcaptype{none} % do not increment counter
\begin{longtable}[]{@{}
  >{\raggedright\arraybackslash}p{(\linewidth - 6\tabcolsep) * \real{0.1818}}
  >{\raggedright\arraybackslash}p{(\linewidth - 6\tabcolsep) * \real{0.2727}}
  >{\raggedright\arraybackslash}p{(\linewidth - 6\tabcolsep) * \real{0.4091}}
  >{\raggedright\arraybackslash}p{(\linewidth - 6\tabcolsep) * \real{0.1364}}@{}}
\toprule\noalign{}
\begin{minipage}[b]{\linewidth}\raggedright
Drug
\end{minipage} & \begin{minipage}[b]{\linewidth}\raggedright
Target
\end{minipage} & \begin{minipage}[b]{\linewidth}\raggedright
Mechanism
\end{minipage} & \begin{minipage}[b]{\linewidth}\raggedright
Use
\end{minipage} \\
\midrule\noalign{}
\endhead
\bottomrule\noalign{}
\endlastfoot
\textbf{Glucocorticoids} & Phospholipase A\(_2\) (indirectly via annexin
A1) & Block most eicosanoid synthesis at the source & Inflammation
(broad effect) \\
\textbf{Aspirin} & COX-1, COX-2 & \textbf{Irreversible} acetylation of
Ser529; permanently inactivates platelet COX-1 (no nucleus \texorpdfstring{\ensuremath{\to}}{->} cannot
resynthesise) & Antiplatelet, anti-inflammatory, analgesic,
antipyretic \\
\textbf{Ibuprofen, naproxen} & COX-1, COX-2 & Reversible competitive
inhibition & Anti-inflammatory, analgesic \\
\textbf{Celecoxib} & COX-2 selective & Reversible & Reduced GI toxicity
(COX-1 spared in gastric mucosa); slight \texorpdfstring{\ensuremath{\uparrow}}{up} thrombotic risk (PGI\(_2\) \texorpdfstring{\ensuremath{\downarrow}}{down}
without TXA\(_2\) \texorpdfstring{\ensuremath{\downarrow}}{down}) \\
\textbf{Montelukast} & CysLT\(_1\) receptor & LT receptor antagonist &
Asthma, allergic rhinitis \\
\textbf{Zileuton} & 5-lipoxygenase & Direct enzyme inhibition &
Asthma \\
\textbf{Misoprostol} & PGE\(_1\) analogue & Synthetic prostaglandin &
Gastric protection, induction of labour \\
\textbf{Latanoprost} & PGF\(_{2\alpha}\) analogue & Increases
uveoscleral outflow & Glaucoma (lowers IOP) \\
\end{longtable}
}

\begin{quote}
\textbf{Clinical Connection --- The aspirin paradox:} Low-dose aspirin
(75--100 mg/day) selectively inhibits platelet COX-1 (anucleate
platelets cannot regenerate the enzyme; 7-day duration) but primarily
transiently inhibits endothelial COX-2 (nucleated cells continuously
resynthesise). Result: net antiplatelet effect with preserved
endothelial PGI\(_2\) \texorpdfstring{\ensuremath{\to}}{->} reduced thrombotic risk. High-dose aspirin loses
this selectivity.
\end{quote}

\begin{center}\rule{0.5\linewidth}{0.5pt}\end{center}

\subsection{Endocrine Disruption}\label{endocrine-disruption}

\textbf{Endocrine-disrupting chemicals (EDCs)} are exogenous substances
that interfere with hormone synthesis, secretion, transport, binding, or
elimination. They are particularly concerning for fetal and neonatal
development, when small hormone perturbations can have lifelong
consequences.

\subsubsection{Mechanisms of disruption}\label{mechanisms-of-disruption}

\begin{enumerate}
\def\labelenumi{\arabic{enumi}.}
\tightlist
\item
  \textbf{Receptor agonism (mimicry):} Exogenous compound binds and
  activates a hormone receptor, mimicking the endogenous ligand (e.g.,
  xenoestrogens activate ER).
\item
  \textbf{Receptor antagonism:} Block native hormone binding (e.g., DDE
  blocks AR; flutamide-like effect).
\item
  \textbf{Altered hormone synthesis:} Inhibit steroidogenic enzymes
  (e.g., conazole fungicides inhibit aromatase; perchlorate blocks NIS
  at the thyroid).
\item
  \textbf{Altered transport/clearance:} Disrupt carrier protein binding
  or hepatic metabolism (e.g., PCBs displace T4 from transthyretin).
\item
  \textbf{Altered receptor expression:} Epigenetic modification of
  hormone receptor genes.
\end{enumerate}

\subsubsection{Bisphenol A (BPA) and other
xenoestrogens}\label{bisphenol-a-bpa-and-other-xenoestrogens}

\textbf{Bisphenol A} is a high-volume industrial monomer used in
polycarbonate plastics, epoxy resins, and thermal receipts. Estimated
production exceeds 6 million tonnes/year. BPA leaches from food
containers, especially when heated or contacting acidic foods.
Detectable BPA is found in \textgreater90\% of US adults' urine.

\textbf{Mechanism:}

\begin{itemize}
\tightlist
\item
  Weak agonist at classical nuclear oestrogen receptors (ER\texorpdfstring{\ensuremath{\alpha}}{alpha}, ER\texorpdfstring{\ensuremath{\beta}}{beta};
  affinity about 10,000\texorpdfstring{\ensuremath{\times}}{x} lower than oestradiol)
\item
  High-affinity agonist at the membrane-bound G-protein--coupled
  oestrogen receptor (GPER, formerly GPR30)
\item
  Binds androgen receptor as antagonist
\item
  Binds thyroid hormone receptor as antagonist
\item
  Activates pregnane X receptor (xenobiotic metabolism)
\end{itemize}

The combination of multiple low-affinity but high-prevalence
interactions makes BPA a ``low-dose'' disruptor --- its non-monotonic
dose--response curve (effects at very low doses absent at higher doses)
violates the classical toxicological assumption that ``the dose makes
the poison.''

\textbf{Effects} (animal and observational human studies): altered
pubertal timing, reduced sperm count and quality, increased risk of
breast and prostate cancers, metabolic dysfunction (obesity, diabetes),
neurodevelopmental and behavioural effects in children. Regulatory
responses have lowered BPA exposure limits and led to bans in baby
bottles in the EU, Canada, and US --- though substitutes (BPS, BPF)
appear to share similar disrupting profiles (``regrettable
substitution'').

\subsubsection{Phthalates, PFAS, and other major
EDCs}\label{phthalates-pfas-and-other-major-edcs}

\textbf{Phthalates} (DEHP, DBP, BBzP) are plasticisers added to PVC to
confer flexibility; also used in personal-care products as
solvents/fragrance carriers. Dietary intake is the main route.
Mechanism: act as \textbf{anti-androgens} via reduced testosterone
synthesis (suppression of StAR and CYP17A1) and AR antagonism in the
developing male reproductive tract. Animal exposures during the male
sex-differentiation window produce the \textbf{``phthalate syndrome''}:
cryptorchidism, hypospadias, reduced anogenital distance, decreased
sperm count. Human cohort studies link prenatal phthalate exposure to
similar genitourinary endpoints.

\textbf{PFAS (per- and polyfluoroalkyl substances)} --- the so-called
``\textbf{forever chemicals}'' (e.g., PFOA, PFOS, GenX) --- are
characterised by C--F bonds that resist environmental and biological
breakdown. Half-lives in humans range from years to decades. Sources:
non-stick cookware (Teflon), water-repellent textiles, firefighting
foams (AFFF), food packaging. Mechanisms include PPAR\texorpdfstring{\ensuremath{\alpha}}{alpha} activation
(fatty-acid metabolism disruption), thyroid hormone displacement from
carrier proteins, and dose-dependent immunosuppression (reduced antibody
response to childhood vaccines documented in Faroe Islands cohort
studies). The C8 Health Project (West Virginia) linked high-dose
occupational PFOA exposure to elevated risks of testicular cancer,
kidney cancer, ulcerative colitis, thyroid disease,
hypercholesterolaemia, and pregnancy-induced hypertension.

{\def\LTcaptype{none} % do not increment counter
\begin{longtable}[]{@{}
  >{\raggedright\arraybackslash}p{(\linewidth - 6\tabcolsep) * \real{0.1304}}
  >{\raggedright\arraybackslash}p{(\linewidth - 6\tabcolsep) * \real{0.2609}}
  >{\raggedright\arraybackslash}p{(\linewidth - 6\tabcolsep) * \real{0.3043}}
  >{\raggedright\arraybackslash}p{(\linewidth - 6\tabcolsep) * \real{0.3043}}@{}}
\toprule\noalign{}
\begin{minipage}[b]{\linewidth}\raggedright
EDC
\end{minipage} & \begin{minipage}[b]{\linewidth}\raggedright
Source
\end{minipage} & \begin{minipage}[b]{\linewidth}\raggedright
Targets
\end{minipage} & \begin{minipage}[b]{\linewidth}\raggedright
Effects
\end{minipage} \\
\midrule\noalign{}
\endhead
\bottomrule\noalign{}
\endlastfoot
\textbf{Phthalates} (DEHP, DBP) & Plasticisers in PVC, cosmetics & AR
antagonism; PPAR\texorpdfstring{\ensuremath{\alpha}}{alpha}, PPAR\texorpdfstring{\ensuremath{\gamma}}{gamma} activation; \texorpdfstring{\ensuremath{\downarrow}}{down} testosterone synthesis &
Anti-androgenic effects on male reproductive development \\
\textbf{DDT/DDE} & Banned pesticide; persistent & ER agonism; AR
antagonism & Eggshell thinning (raptors); breast cancer association \\
\textbf{PCBs} & Banned but persistent industrial chemicals & Bind
transthyretin; AhR agonism & Thyroid disruption; neurodevelopmental
harm \\
\textbf{Atrazine} & Herbicide & Aromatase induction & Feminisation of
male amphibians \\
\textbf{PFAS} (PFOA, PFOS) & Non-stick coatings, firefighting foam,
water & PPAR\texorpdfstring{\ensuremath{\alpha}}{alpha}; thyroid; immune & ``Forever chemicals''; thyroid
disruption, immunosuppression, cancer \\
\textbf{Phytoestrogens} (genistein, daidzein) & Soy, legumes & Weak ER
agonist (preferentially ER\texorpdfstring{\ensuremath{\beta}}{beta}) & Modest health effects; debate about
infant soy formula \\
\textbf{Tributyltin} & Marine antifouling paint & RXR/PPAR\texorpdfstring{\ensuremath{\gamma}}{gamma} agonism
(obesogen) & Imposex in molluscs; potential adipogenesis \\
\textbf{Pesticides (vinclozolin, linuron)} & Agriculture & AR antagonism
& Anti-androgenic, transgenerational epigenetic effects \\
\end{longtable}
}

\subsubsection{Vulnerability of
development}\label{vulnerability-of-development}

Fetal and early-life development is uniquely vulnerable to EDCs:

\begin{itemize}
\tightlist
\item
  \textbf{Hormone-dependent organogenesis:} Sex differentiation
  (Wolffian/Müllerian fates), brain dimorphism, and thyroid-driven CNS
  myelination depend on tightly timed hormone windows. A small
  perturbation during these windows can permanently rewire the
  architecture in ways that would not occur in adults.
\item
  \textbf{Limited metabolic clearance:} Fetal liver enzymes (CYPs, UGTs)
  are immature; placenta does not always exclude lipophilic toxicants.
\item
  \textbf{Exponential cell division:} Programming events, including DNA
  methylation and histone modifications, are particularly susceptible to
  epigenetic disruption during rapid cell division.
\item
  \textbf{Critical periods are non-recoverable:} Once a developmental
  window closes, the missing signal cannot be replaced by later
  supplementation.
\end{itemize}

\begin{quote}
\textbf{Concept Check 3:} Why might the developmental period (fetal life
through puberty) be uniquely vulnerable to endocrine disruptors compared
to adult exposures? Consider hormone-dependent organ development and the
absence of redundant pathways during organogenesis.
\end{quote}

\begin{center}\rule{0.5\linewidth}{0.5pt}\end{center}

\newpage

\subsection{Immune System}\label{immune-system}

The immune system protects against pathogens and tumour cells while
preserving tolerance to self. It comprises two integrated arms:
\textbf{innate immunity} (rapid, non-specific, germline-encoded) and
\textbf{adaptive immunity} (slow, antigen-specific, somatically
generated).

\subsubsection{Innate Immunity}\label{innate-immunity-1}

\hyperref[gl:innate-immunity]{\textbf{Innate immunity}} provides
immediate (seconds to hours), non-specific protection:

\begin{figure}[htbp]
\centering
\includegraphics[width=0.82\linewidth,height=4.2in,keepaspectratio]{../figures/mermaid_inline/inline_mermaid_0169_631db71c65c6.png}
\caption{Components of innate immunity Physical barriers form the first line of defence. When breached, inflammation recruits phagocytes, complement activates, NK cells kill infected cells, and interferons establish an antiviral state. Antigen-presenting cells bridge innate to adaptive immunity.}
\end{figure}

\textbf{Key innate immune cells:}

{\def\LTcaptype{none} % do not increment counter
\begin{longtable}[]{@{}
  >{\raggedright\arraybackslash}p{(\linewidth - 4\tabcolsep) * \real{0.3103}}
  >{\raggedright\arraybackslash}p{(\linewidth - 4\tabcolsep) * \real{0.2759}}
  >{\raggedright\arraybackslash}p{(\linewidth - 4\tabcolsep) * \real{0.4138}}@{}}
\toprule\noalign{}
\begin{minipage}[b]{\linewidth}\raggedright
Cell Type
\end{minipage} & \begin{minipage}[b]{\linewidth}\raggedright
Function
\end{minipage} & \begin{minipage}[b]{\linewidth}\raggedright
Key Features
\end{minipage} \\
\midrule\noalign{}
\endhead
\bottomrule\noalign{}
\endlastfoot
\textbf{Neutrophils} & First responders; phagocytosis; NETs; oxidative
burst & Most abundant WBC (60--70\%); short-lived (hours) \\
\textbf{Macrophages} & Phagocytosis; antigen presentation; cytokine
production & Tissue-resident (Kupffer cells in liver, microglia in
brain, alveolar macrophages in lung) \\
\textbf{Dendritic cells} & Professional APCs; bridge innate and adaptive
& Most potent antigen presenters \\
\textbf{NK cells} & Kill virus-infected and tumour cells & ``Missing
self'' detection via KIR receptors \\
\textbf{Mast cells} & Histamine release; IgE-mediated degranulation &
Allergy; \hyperref[gl:parasite]{\textbf{parasite}} defence \\
\textbf{Eosinophils} & Parasite defence; allergic inflammation & Major
basic protein toxic to helminths \\
\textbf{Basophils} & Histamine; IL-4 production & Rarest WBC
(\textless1\%) \\
\end{longtable}
}

\subsubsection{Pattern Recognition Receptors
(PRRs)}\label{pattern-recognition-receptors-prrs-1}

Innate immune cells detect pathogens through germline-encoded
\textbf{PRRs} that recognise conserved molecular signatures unique to
pathogens --- \textbf{pathogen-associated molecular patterns (PAMPs)}
--- or signals of cellular damage --- \textbf{damage-associated
molecular patterns (DAMPs)}. PRRs fall into four major families based on
cellular location and ligand class.

\paragraph{Toll-like receptors (TLRs)}\label{toll-like-receptors-tlrs}

Membrane-bound (plasma membrane or endosomal). Humans express 10 TLRs.

Species specificity matters here. Human TLR1--TLR10 are not a
comprehensive mammalian template: mice lack a direct functional
equivalent of human TLR10 but retain TLR11--TLR13, which detect
microbial ligands such as profilin-like proteins and bacterial RNA.
TLR10 itself remains less mechanistically settled than TLR4, TLR7/8, or
TLR9. When comparing innate-immunity experiments across humans, mice,
and cell lines, students should ask whether the receptor repertoire and
ligand preparation actually match the claimed pathogen-sensing pathway.

{\def\LTcaptype{none} % do not increment counter
\begin{longtable}[]{@{}
  >{\raggedright\arraybackslash}p{(\linewidth - 6\tabcolsep) * \real{0.0968}}
  >{\raggedright\arraybackslash}p{(\linewidth - 6\tabcolsep) * \real{0.2581}}
  >{\raggedright\arraybackslash}p{(\linewidth - 6\tabcolsep) * \real{0.1935}}
  >{\raggedright\arraybackslash}p{(\linewidth - 6\tabcolsep) * \real{0.4516}}@{}}
\toprule\noalign{}
\begin{minipage}[b]{\linewidth}\raggedright
TLR
\end{minipage} & \begin{minipage}[b]{\linewidth}\raggedright
Location
\end{minipage} & \begin{minipage}[b]{\linewidth}\raggedright
Ligand
\end{minipage} & \begin{minipage}[b]{\linewidth}\raggedright
Pathogen class
\end{minipage} \\
\midrule\noalign{}
\endhead
\bottomrule\noalign{}
\endlastfoot
TLR1/2 & Plasma membrane & Triacyl lipopeptides & Bacteria
(mycobacteria) \\
TLR2/6 & Plasma membrane & Diacyl lipopeptides, peptidoglycan & Gram+
bacteria, fungi \\
TLR3 & Endosomal & dsRNA & Viruses \\
\textbf{TLR4} & Plasma membrane & \textbf{LPS (lipopolysaccharide)} &
Gram\texorpdfstring{\ensuremath{-}}{-} bacteria \\
TLR5 & Plasma membrane & Flagellin & Motile bacteria \\
TLR7/8 & Endosomal & ssRNA & RNA viruses \\
TLR9 & Endosomal & Unmethylated CpG DNA & Bacteria, DNA viruses \\
\end{longtable}
}

\paragraph{TLR4 \texorpdfstring{\ensuremath{\to}}{->} MyD88 \texorpdfstring{\ensuremath{\to}}{->} NF-\texorpdfstring{\ensuremath{\kappa}}{kappa}B pathway (bacterial LPS
response)}\label{tlr4-myd88-nf-ux3bab-pathway-bacterial-lps-response}

\begin{figure}[htbp]
\centering
\includegraphics[width=0.82\linewidth,height=4.2in,keepaspectratio]{../figures/mermaid_inline/inline_mermaid_0170_f94978e2386b.png}
\caption{TLR4/MyD88/NF-kappaB pathway. LBP and CD14 deliver bacterial LPS to TLR4/MD-2, TLR4 dimerization recruits TIRAP and MyD88, IRAK kinases and TRAF6 activate TAK1 and IKK, and IkappaBalpha degradation releases NF-kappaB to induce inflammatory genes.}
\end{figure}

\paragraph{TLR3/TRIF \texorpdfstring{\ensuremath{\to}}{->} IRF3 \texorpdfstring{\ensuremath{\to}}{->} IFN-\texorpdfstring{\ensuremath{\beta}}{beta} pathway (antiviral
response)}\label{tlr3trif-irf3-ifn-ux3b2-pathway-antiviral-response}

A parallel branch is engaged by TLR3 (endosomal dsRNA) and the
late-endosome pool of TLR4. The adaptor \textbf{TRIF} recruits TBK1,
which phosphorylates \textbf{IRF3}. Phospho-IRF3 dimerises, enters the
nucleus, and drives transcription of \textbf{type I interferons
(IFN-\texorpdfstring{\ensuremath{\alpha}}{alpha}/\texorpdfstring{\ensuremath{\beta}}{beta})}. IFN-\texorpdfstring{\ensuremath{\beta}}{beta} released into the extracellular space binds IFNAR on
neighbouring cells, activating JAK1/TYK2 \texorpdfstring{\ensuremath{\to}}{->} STAT1/STAT2 \texorpdfstring{\ensuremath{\to}}{->} ISGF3 \texorpdfstring{\ensuremath{\to}}{->}
induction of hundreds of \textbf{interferon-stimulated genes (ISGs)}
that establish an antiviral state. The MyD88/NF-\texorpdfstring{\ensuremath{\kappa}}{kappa}B vs TRIF/IRF3
dichotomy explains why bacterial LPS produces fever and inflammation
while viral dsRNA produces an interferon-driven antiviral state.

\paragraph{NOD-like receptors (NLRs) and the NLRP3
inflammasome}\label{nod-like-receptors-nlrs-and-the-nlrp3-inflammasome}

Cytosolic. Detect intracellular bacterial components (peptidoglycan
derivatives) and danger signals.

\begin{itemize}
\tightlist
\item
  \textbf{NOD1} detects iE-DAP (Gram\texorpdfstring{\ensuremath{-}}{-}); \textbf{NOD2} detects MDP
  (comprehensive). Both activate NF-\texorpdfstring{\ensuremath{\kappa}}{kappa}B. NOD2 mutations cause Crohn's
  disease (impaired mucosal immunity \texorpdfstring{\ensuremath{\to}}{->} dysbiosis \texorpdfstring{\ensuremath{\to}}{->} inflammation).
\end{itemize}

The \textbf{NLRP3 inflammasome} illustrates the \textbf{two-signal}
model:

\begin{itemize}
\tightlist
\item
  \textbf{Signal 1 (priming):} TLR or cytokine engagement \texorpdfstring{\ensuremath{\to}}{->} NF-\texorpdfstring{\ensuremath{\kappa}}{kappa}B \texorpdfstring{\ensuremath{\to}}{->}
  upregulates NLRP3 and pro-IL-1\texorpdfstring{\ensuremath{\beta}}{beta} transcription. Without this, no
  inflammasome assembly.
\item
  \textbf{Signal 2 (activation):} Diverse triggers --- extracellular ATP
  (P2X7), K\(^+\) efflux, lysosomal rupture (urate crystals, cholesterol
  crystals, silica, alum), mitochondrial ROS, mitochondrial DNA in
  cytosol --- activate NLRP3.
\item
  \textbf{Assembly:} NLRP3 oligomerises via its NACHT domain; recruits
  ASC adaptor via PYD--PYD interactions; ASC nucleates pro-caspase-1 via
  CARD--CARD; pro-caspase-1 self-cleaves to active caspase-1.
\item
  \textbf{Output:} Caspase-1 cleaves pro-IL-1\texorpdfstring{\ensuremath{\beta}}{beta} \texorpdfstring{\ensuremath{\to}}{->} IL-1\texorpdfstring{\ensuremath{\beta}}{beta} (released via
  gasdermin D pores) and pro-IL-18 \texorpdfstring{\ensuremath{\to}}{->} IL-18; cleaves gasdermin D, whose
  N-terminal fragment forms 10--20 nm pores in the plasma membrane
  causing \textbf{pyroptosis} (lytic cell death with massive cytokine
  release).
\end{itemize}

NLRP3 mutations cause cryopyrin-associated periodic syndromes (CAPS);
chronic NLRP3 activity drives gout (urate crystals), atherosclerosis
(cholesterol crystals), Alzheimer's-related neuroinflammation, and type
2 diabetes. \textbf{Anakinra} (recombinant IL-1Ra) and
\textbf{canakinumab} (anti-IL-1\texorpdfstring{\ensuremath{\beta}}{beta}) target the inflammasome output.

\paragraph{RIG-I-like receptors (RLRs)}\label{rig-i-like-receptors-rlrs}

Cytosolic RNA sensors detecting viral replication.

\begin{itemize}
\tightlist
\item
  \textbf{RIG-I:} detects 5'-triphosphate RNA (host RNA is capped; viral
  RNA is not).
\item
  \textbf{MDA5:} detects long dsRNA.
\item
  Signal via \textbf{MAVS} (mitochondrial antiviral signalling protein)
  \texorpdfstring{\ensuremath{\to}}{->} IRF3/IRF7 \texorpdfstring{\ensuremath{\to}}{->} type I interferons.
\end{itemize}

\paragraph{cGAS--STING pathway (cytosolic DNA
sensing)}\label{cgassting-pathway-cytosolic-dna-sensing}

The \textbf{cGAS--STING} axis is the principal sensor for cytosolic DNA
--- a hallmark of intracellular bacterial or viral infection (and,
problematically, mislocalised mitochondrial or self DNA).

\begin{itemize}
\tightlist
\item
  \textbf{cGAS (cyclic GMP--AMP synthase)} binds dsDNA
  non-sequence-specifically through a phase-separation-like
  condensation. Activated cGAS catalyses synthesis of the cyclic
  dinucleotide \textbf{2'3'-cGAMP} from ATP and GTP.
\item
  \textbf{STING (stimulator of interferon genes)}, an ER-resident
  transmembrane protein, binds 2'3'-cGAMP, undergoes a major
  conformational change, traffics from ER to ERGIC/Golgi, and recruits
  \textbf{TBK1}, which phosphorylates \textbf{IRF3} \texorpdfstring{\ensuremath{\to}}{->} type I interferon
  transcription. STING also activates NF-\texorpdfstring{\ensuremath{\kappa}}{kappa}B via a parallel branch.
\end{itemize}

The cGAS--STING pathway is essential for control of HSV-1, vaccinia, and
many cytosolic bacteria. \textbf{Dysregulation drives autoimmunity:}
mutations causing constitutive STING activation produce \textbf{SAVI
(STING-associated vasculopathy with onset in infancy)}, an
interferonopathy. Aicardi-Goutières syndrome arises when defective
DNases (TREX1, RNASEH2) cannot clear cytoplasmic nucleic acids,
chronically engaging cGAS-STING. Pharmacologically, STING agonists
(ADU-S100) are being trialled as cancer adjuvants because tumour-induced
type I IFN can boost antitumour immunity.

\begin{figure}[htbp]
\centering
\includegraphics[width=0.82\linewidth,height=4.2in,keepaspectratio]{../figures/mermaid_inline/inline_mermaid_0171_1e90bd6d08e2.png}
\caption{PRR signalling pathways TLRs, NLRs, RLRs, and cGAS converge on transcription factors NF-kappaB (inflammation), IRF3/7 (interferons), and the inflammasome (IL-1beta, pyroptosis). Different pathogen classes preferentially engage different sensors.}
\end{figure}

\subsubsection{Complement System --- Comprehensive
Overview}\label{complement-system-comprehensive-overview}

The complement system comprises about 30 plasma proteins that
amplify innate responses through enzymatic cascades. There are three
pathways of activation, most converging on a common terminal pathway.

\paragraph{Three activation pathways}\label{three-activation-pathways}

{\def\LTcaptype{none} % do not increment counter
\begin{longtable}[]{@{}
  >{\raggedright\arraybackslash}p{(\linewidth - 6\tabcolsep) * \real{0.1750}}
  >{\raggedright\arraybackslash}p{(\linewidth - 6\tabcolsep) * \real{0.1750}}
  >{\raggedright\arraybackslash}p{(\linewidth - 6\tabcolsep) * \real{0.3750}}
  >{\raggedright\arraybackslash}p{(\linewidth - 6\tabcolsep) * \real{0.2750}}@{}}
\toprule\noalign{}
\begin{minipage}[b]{\linewidth}\raggedright
Pathway
\end{minipage} & \begin{minipage}[b]{\linewidth}\raggedright
Trigger
\end{minipage} & \begin{minipage}[b]{\linewidth}\raggedright
Initiation step
\end{minipage} & \begin{minipage}[b]{\linewidth}\raggedright
Convergence
\end{minipage} \\
\midrule\noalign{}
\endhead
\bottomrule\noalign{}
\endlastfoot
\textbf{Classical} & Antibody (IgM, IgG) bound to antigen on pathogen &
C1q binds Fc \texorpdfstring{\ensuremath{\to}}{->} activates C1r \texorpdfstring{\ensuremath{\to}}{->} C1s \texorpdfstring{\ensuremath{\to}}{->} cleaves C4 + C2 & C3 convertase =
C4b2a \\
\textbf{Lectin} & Pathogen surface mannose / GlcNAc & MBL or ficolins
bind sugars \texorpdfstring{\ensuremath{\to}}{->} MASP1/MASP2 (analogous to C1r/C1s) \texorpdfstring{\ensuremath{\to}}{->} cleave C4 + C2 & C3
convertase = C4b2a \\
\textbf{Alternative} & Spontaneous ``tick-over'' hydrolysis of C3;
amplified on pathogen surfaces lacking complement regulators &
C3(H\(_2\)O) + factor B + factor D \texorpdfstring{\ensuremath{\to}}{->} C3(H\(_2\)O)Bb (initial fluid-phase
convertase) \texorpdfstring{\ensuremath{\to}}{->} deposits C3b on surface \texorpdfstring{\ensuremath{\to}}{->} C3bBb (surface convertase,
stabilised by properdin) & C3 convertase = C3bBb \\
\end{longtable}
}

\paragraph{C3 convertase, C5 convertase,
MAC}\label{c3-convertase-c5-convertase-mac}

Most three pathways generate a \textbf{C3 convertase} (C4b2a or C3bBb)
that cleaves \textbf{C3 \texorpdfstring{\ensuremath{\to}}{->} C3a + C3b}. C3b is deposited on the pathogen
surface; binding of an additional C3b to the existing C3 convertase
yields the \textbf{C5 convertase} (C4b2aC3b or C3bBbC3b) that cleaves
\textbf{C5 \texorpdfstring{\ensuremath{\to}}{->} C5a + C5b}.

C5b initiates the \textbf{terminal pathway}: C5b \texorpdfstring{\ensuremath{\to}}{->} C5b-C6 \texorpdfstring{\ensuremath{\to}}{->} C5b-C6-C7
(membrane-inserting) \texorpdfstring{\ensuremath{\to}}{->} C5b-C6-C7-C8 \texorpdfstring{\ensuremath{\to}}{->} addition of multiple C9 monomers
polymerising into the \textbf{membrane attack complex (MAC, C5b-9)}. The
MAC forms a 10 nm transmembrane pore that lyses the target cell.
Gram-negative bacteria are particularly vulnerable; encapsulated
bacteria (\emph{Neisseria}) require complement for clearance, which is
why C5--C9 deficiencies present with recurrent meningococcal infection.

\paragraph{Effector functions}\label{effector-functions}

\begin{itemize}
\tightlist
\item
  \textbf{Opsonisation:} C3b coats pathogen \texorpdfstring{\ensuremath{\to}}{->} recognised by phagocyte
  receptors \textbf{CR1} (C3b/C4b), \textbf{CR3} (iC3b), \textbf{CR4}.
  Opsonised particles are 1000-fold more efficiently phagocytosed.
\item
  \textbf{Membrane attack complex (MAC):} C5b-9 polymerises in target
  membrane, forming a 10 nm pore.
\item
  \textbf{Anaphylatoxins (chemotaxis and inflammation):} C3a and
  \textbf{C5a} (the most potent) recruit neutrophils, activate mast
  cells, increase vascular permeability, and amplify local inflammation.
\item
  \textbf{Immune complex clearance:} CR1 on erythrocytes binds
  C3b-coated immune complexes and ferries them to liver/spleen for
  disposal.
\item
  \textbf{B cell co-stimulation:} C3d coupled to antigen lowers the BCR
  signalling threshold about 10,000 fold via CR2 (CD21).
\end{itemize}

\paragraph{Amplification dynamics and
regulation}\label{amplification-dynamics-and-regulation}

The cascade is intrinsically amplifying because each enzyme cleaves many
substrates. If one C3 convertase cleaves N copies of C3 per second, with
decay rate \(k_d\), the steady-state C3b concentration scales as

\begin{equation}
[\text{C3b}] = \frac{[\text{C3conv}] \cdot N}{k_d}
\label{eq:unit_IX_amplification}
\end{equation}

The \textbf{alternative pathway amplification loop} is positive: each
new C3b binds factor B \texorpdfstring{\ensuremath{\to}}{->} C3 convertase \texorpdfstring{\ensuremath{\to}}{->} cleaves more C3 \texorpdfstring{\ensuremath{\to}}{->} more C3b.
Without regulators, this loop would consume most plasma C3 within
minutes.

Regulators confine the cascade to pathogen surfaces:

{\def\LTcaptype{none} % do not increment counter
\begin{longtable}[]{@{}
  >{\raggedright\arraybackslash}p{(\linewidth - 4\tabcolsep) * \real{0.3600}}
  >{\raggedright\arraybackslash}p{(\linewidth - 4\tabcolsep) * \real{0.3200}}
  >{\raggedright\arraybackslash}p{(\linewidth - 4\tabcolsep) * \real{0.3200}}@{}}
\toprule\noalign{}
\begin{minipage}[b]{\linewidth}\raggedright
Regulator
\end{minipage} & \begin{minipage}[b]{\linewidth}\raggedright
Location
\end{minipage} & \begin{minipage}[b]{\linewidth}\raggedright
Function
\end{minipage} \\
\midrule\noalign{}
\endhead
\bottomrule\noalign{}
\endlastfoot
\textbf{DAF (CD55)} & GPI-anchored on host cells & Accelerates C3/C5
convertase decay \\
\textbf{CD59} & GPI-anchored on host cells & Blocks MAC assembly (C9
incorporation) \\
\textbf{Factor H} & Soluble plasma protein & Binds host-specific sialic
acid; cofactor for factor I cleavage of C3b \\
\textbf{C4BP (C4b binding protein)} & Soluble & Inactivates C4b \\
\textbf{C1-INH (C1 esterase inhibitor)} & Soluble & Inhibits C1r/C1s and
MASPs \\
\textbf{CR1 (CD35)} & Erythrocytes, lymphocytes & Cofactor for factor I;
immune complex clearance \\
\end{longtable}
}

Patients with \textbf{paroxysmal nocturnal hemoglobinuria (PNH)} lack
the GPI anchor that tethers DAF and CD59 to RBCs \texorpdfstring{\ensuremath{\to}}{->} uncontrolled
complement activation \texorpdfstring{\ensuremath{\to}}{->} haemolysis. Treated with \textbf{eculizumab}
(anti-C5 monoclonal antibody; blocks C5 cleavage and MAC formation).
\textbf{C1-INH deficiency} causes \textbf{hereditary angioedema}
(uncontrolled C1 \texorpdfstring{\ensuremath{\to}}{->} bradykinin generation via the kinin--kallikrein
system \texorpdfstring{\ensuremath{\to}}{->} sudden tissue swelling).

\subsubsection{Adaptive Immunity
Overview}\label{adaptive-immunity-overview-1}

\hyperref[gl:adaptive-immunity]{\textbf{Adaptive immunity}} is slower
(7--10 days for primary response) but provides \textbf{specificity}
(each lymphocyte recognises a unique antigen) and \textbf{memory}
(faster, stronger response on re-exposure --- basis of vaccination).

\textbf{Two arms:}

\begin{itemize}
\tightlist
\item
  \textbf{Cell-mediated:} T cells (CD8+ cytotoxic kill infected cells;
  CD4+ helper coordinate response)
\item
  \textbf{Humoral:} B cells produce antibodies that neutralise
  pathogens, opsonise, activate complement
\end{itemize}

\textbf{Antigen presentation via MHC:}

\begin{itemize}
\tightlist
\item
  \textbf{MHC class I} (on most nucleated cells): presents endogenous
  peptides (8--10 aa) from cytosolic proteins. Pathway: cytosolic
  protein \texorpdfstring{\ensuremath{\to}}{->} proteasome \texorpdfstring{\ensuremath{\to}}{->} TAP transporter \texorpdfstring{\ensuremath{\to}}{->} ER \texorpdfstring{\ensuremath{\to}}{->} loaded onto MHC I \texorpdfstring{\ensuremath{\to}}{->}
  surface. Presented to \textbf{CD8+ T cells}.
\item
  \textbf{MHC class II} (on professional APCs: dendritic cells,
  macrophages, B cells): presents exogenous peptides (12--25 aa) from
  internalised pathogens. Pathway: phagocytosis \texorpdfstring{\ensuremath{\to}}{->} endolysosome \texorpdfstring{\ensuremath{\to}}{->}
  cathepsin cleavage \texorpdfstring{\ensuremath{\to}}{->} loaded onto MHC II (after CLIP removal by HLA-DM)
  \texorpdfstring{\ensuremath{\to}}{->} surface. Presented to \textbf{CD4+ T cells}.
\item
  \textbf{Cross-presentation:} Dendritic cells can also load exogenous
  antigens onto MHC I, important for activating CD8+ responses against
  viruses that don't directly infect APCs.
\end{itemize}

\subsubsection{T Cell Development and
Selection}\label{t-cell-development-and-selection}

T cell precursors leave the bone marrow as immature CD4\(^-\)CD8\(^-\)
``double-negative'' thymocytes and migrate to the thymus, where they
undergo somatic recombination and selection.

\paragraph{Double-negative (DN1--DN4)
staging}\label{double-negative-dn1dn4-staging}

Within the thymic cortex, double-negative thymocytes pass through four
sequential stages defined by surface CD44 and CD25 expression:

{\def\LTcaptype{none} % do not increment counter
\begin{longtable}[]{@{}
  >{\raggedright\arraybackslash}p{(\linewidth - 6\tabcolsep) * \real{0.2083}}
  >{\raggedright\arraybackslash}p{(\linewidth - 6\tabcolsep) * \real{0.1667}}
  >{\raggedright\arraybackslash}p{(\linewidth - 6\tabcolsep) * \real{0.1667}}
  >{\raggedright\arraybackslash}p{(\linewidth - 6\tabcolsep) * \real{0.4583}}@{}}
\toprule\noalign{}
\begin{minipage}[b]{\linewidth}\raggedright
Stage
\end{minipage} & \begin{minipage}[b]{\linewidth}\raggedright
CD44
\end{minipage} & \begin{minipage}[b]{\linewidth}\raggedright
CD25
\end{minipage} & \begin{minipage}[b]{\linewidth}\raggedright
Major event
\end{minipage} \\
\midrule\noalign{}
\endhead
\bottomrule\noalign{}
\endlastfoot
\textbf{DN1} & + & \texorpdfstring{\ensuremath{-}}{-} & Early thymic progenitor; multipotent
(T/NK/myeloid) \\
\textbf{DN2} & + & + & Lineage commitment; TCR \texorpdfstring{\ensuremath{\beta}}{beta}/\texorpdfstring{\ensuremath{\gamma}}{gamma}/\texorpdfstring{\ensuremath{\delta}}{delta} rearrangement
begins \\
\textbf{DN3} & \texorpdfstring{\ensuremath{-}}{-} & + & TCR \texorpdfstring{\ensuremath{\beta}}{beta} rearrangement complete; \textbf{\texorpdfstring{\ensuremath{\beta}}{beta}-selection
checkpoint} (primarily cells with productive TCR \texorpdfstring{\ensuremath{\beta}}{beta} survive, via
signalling from pre-TCR with surrogate \texorpdfstring{\ensuremath{\alpha}}{alpha} chain pT\texorpdfstring{\ensuremath{\alpha}}{alpha}) \\
\textbf{DN4} & \texorpdfstring{\ensuremath{-}}{-} & \texorpdfstring{\ensuremath{-}}{-} & Proliferative burst; transition to
double-positive (CD4+CD8+) \\
\end{longtable}
}

Cells then become \textbf{double-positive (CD4+CD8+)} and rearrange TCR
\texorpdfstring{\ensuremath{\alpha}}{alpha}. With both TCR chains expressed, they undergo \textbf{positive} then
\textbf{negative} selection.

\paragraph{TCR diversity generation}\label{tcr-diversity-generation}

Like immunoglobulin loci, the T cell receptor (TCR) loci undergo
\textbf{V(D)J recombination} mediated by RAG1/RAG2 endonucleases:

\begin{itemize}
\tightlist
\item
  TCR \texorpdfstring{\ensuremath{\beta}}{beta} chain: V--D--J recombination (about 52 V \texorpdfstring{\ensuremath{\times}}{x} 2 D \texorpdfstring{\ensuremath{\times}}{x} 13 J
  segments \texorpdfstring{\ensuremath{\to}}{->} about 1,400 combinations)
\item
  TCR \texorpdfstring{\ensuremath{\alpha}}{alpha} chain: V--J recombination (about 70 V \texorpdfstring{\ensuremath{\times}}{x}
  about 61 J segments \texorpdfstring{\ensuremath{\to}}{->} about 4,300 combinations)
\item
  \textbf{Junctional diversity:} TdT (terminal deoxynucleotidyl
  transferase) adds non-templated N-nucleotides at junctions
\item
  \textbf{Combinatorial diversity:} \texorpdfstring{\ensuremath{\alpha}}{alpha}\texorpdfstring{\ensuremath{\beta}}{beta} pairing creates about 6
  \texorpdfstring{\ensuremath{\times}}{x} 10\(^6\) unique receptors before junctional diversity
\item
  \textbf{With junctional diversity:} the theoretical TCR repertoire
  exceeds 10\(^{18}\), far larger than the about 10\(^{11}\) T
  cells in the human body --- most TCRs are rarely realised.
\end{itemize}

Rough calculation: \(1{,}400 \times 4{,}300 \approx 6 \times 10^6\)
V(D)J combinations \texorpdfstring{\ensuremath{\times}}{x} about 10\(^{12}\) junctional possibilities
\texorpdfstring{\ensuremath{\approx}}{~} 10\(^{18}\) theoretical receptors.

Defects in this recombination machinery cause severe combined
immunodeficiencies (RAG1/RAG2 SCID --- ``bubble boy'' disease);
ataxia-telangiectasia (ATM mutation) causes radiosensitivity and
lymphoid malignancy.

\paragraph{Positive selection (cortex)}\label{positive-selection-cortex}

In the thymic cortex, double-positive (CD4+CD8+) thymocytes encounter
\textbf{cortical thymic epithelial cells (cTECs)} displaying
self-peptide--MHC complexes. T cells whose TCR engages MHC with
sufficient (but not excessive) affinity receive a survival signal. T
cells with no MHC affinity die by neglect (about 95\% of
thymocytes).

\begin{itemize}
\tightlist
\item
  TCR engagement of MHC I \texorpdfstring{\ensuremath{\to}}{->} CD8+ single-positive T cell
\item
  TCR engagement of MHC II \texorpdfstring{\ensuremath{\to}}{->} CD4+ single-positive T cell
\end{itemize}

Positive selection ensures the surviving repertoire is
\textbf{MHC-restricted} --- primarily recognises antigen presented in
the context of self MHC.

\paragraph{Negative selection (medulla) --- AIRE and Treg
generation}\label{negative-selection-medulla-aire-and-treg-generation}

Surviving thymocytes migrate to the medulla and encounter
\textbf{medullary thymic epithelial cells (mTECs)} and dendritic cells.
mTECs express the autoimmune regulator \textbf{AIRE} transcription
factor, which drives ectopic expression of thousands of tissue-specific
antigens (insulin, thyroglobulin, myelin proteins) normally restricted
to peripheral tissues. T cells whose TCR binds self-peptide--MHC
complexes with \textbf{high affinity} undergo apoptosis (clonal
deletion). A small fraction with intermediate self-reactivity become
\textbf{regulatory T cells (Tregs)} --- CD4+CD25+FoxP3+ cells that
police self-tolerance in the periphery.

\textbf{AIRE mutations} cause autoimmune polyendocrinopathy syndrome
type 1 (APS-1, APECED) --- patients fail to delete autoreactive T cells
against multiple endocrine organs, developing chronic mucocutaneous
candidiasis, hypoparathyroidism, and adrenal insufficiency.
\textbf{FoxP3 mutations} cause IPEX syndrome --- fatal multi-organ
autoimmunity in infancy from absent Tregs.

\subsubsection{B Cell Activation and Antibody
Diversification}\label{b-cell-activation-and-antibody-diversification}

\paragraph{Stepwise B cell activation}\label{stepwise-b-cell-activation}

\begin{enumerate}
\def\labelenumi{\arabic{enumi}.}
\tightlist
\item
  \textbf{Antigen encounter.} Naïve B cell in lymphoid follicle
  encounters cognate antigen (in soluble form, or displayed on
  subcapsular sinus macrophages and follicular dendritic cells).
\item
  \textbf{BCR cross-linking.} Multivalent antigen cross-links several
  BCRs \texorpdfstring{\ensuremath{\to}}{->} tyrosine phosphorylation of Ig\texorpdfstring{\ensuremath{\alpha}}{alpha}/Ig\texorpdfstring{\ensuremath{\beta}}{beta} ITAMs by Lyn/Fyn \texorpdfstring{\ensuremath{\to}}{->}
  recruitment of Syk \texorpdfstring{\ensuremath{\to}}{->} activation of PI3K, PLC\texorpdfstring{\ensuremath{\gamma}}{gamma}2, Ras/MAPK cascades.
\item
  \textbf{Antigen internalisation.} B cell internalises antigen via BCR,
  processes peptides in MHC class II compartment.
\item
  \textbf{Migration to T-B border.} Activated B cell upregulates CCR7,
  migrates to T-cell zone of lymph node.
\item
  \textbf{T cell help.} Cognate Tfh (follicular helper T cell)
  recognises B cell-presented peptide on MHC II. Engagement of
  \textbf{CD40L (Tfh) -- CD40 (B cell)} plus cytokines (IL-4, IL-21)
  provides ``second signal.''
\item
  \textbf{Outcome:} B cells either differentiate into
  \textbf{short-lived extrafollicular plasmablasts} (rapid IgM, low
  affinity) or enter the germinal centre.
\end{enumerate}

\paragraph{B cell activation modes}\label{b-cell-activation-modes}

\begin{itemize}
\tightlist
\item
  \textbf{T-independent (TI) responses:} Polysaccharide and repetitive
  antigens cross-link many BCRs simultaneously. Produces predominantly
  IgM, no germinal centre, no affinity maturation, weak memory.
  Important for encapsulated bacteria (\emph{Streptococcus pneumoniae},
  \emph{Haemophilus influenzae}); explains why polysaccharide vaccines
  work poorly in children \textless2 years.
\item
  \textbf{T-dependent (TD) responses:} Protein antigens. B cell
  internalises antigen via BCR, processes it, presents peptide on MHC
  II, is recognised by cognate CD4+ T helper cell (specifically Tfh ---
  follicular helper T cell). T:B interaction at the T-B border activates
  the B cell to enter the \textbf{germinal centre reaction}. Produces
  high-affinity, class-switched antibodies and long-lived memory.
\end{itemize}

\paragraph{Germinal centre reaction --- somatic hypermutation and
affinity
maturation}\label{germinal-centre-reaction-somatic-hypermutation-and-affinity-maturation}

In secondary lymphoid organs (lymph nodes, spleen, Peyer's patches),
activated B cells form \textbf{germinal centres} with two zones:

\begin{itemize}
\tightlist
\item
  \textbf{Dark zone (centroblasts):} Rapid proliferation.
  \textbf{Activation-induced cytidine deaminase (AID)} introduces
  somatic point mutations into the variable regions of immunoglobulin
  genes --- \textbf{somatic hypermutation (SHM)},
  about 10\(^{-3}\) mutations per base per generation
  (about 10\(^6\) \texorpdfstring{\ensuremath{\times}}{x} normal mutation rate).
\item
  \textbf{Light zone (centrocytes):} B cells re-encounter antigen
  displayed on follicular dendritic cells (FDCs). Cells whose mutated
  BCR has improved antigen affinity capture more antigen, internalise
  it, present more peptide on MHC II, and receive stronger Tfh help \texorpdfstring{\ensuremath{\to}}{->}
  survival and re-cycling. Cells with reduced affinity die by apoptosis.
  This is \textbf{affinity maturation} --- Darwinian selection at the
  cellular level, driving 1000-fold increases in antibody affinity over
  weeks.
\end{itemize}

Light-zone cells differentiate into:

\begin{itemize}
\tightlist
\item
  \textbf{Plasma cells} (long-lived in bone marrow, secrete antibodies
  for years/decades)
\item
  \textbf{Memory B cells} (rapidly mobilised on antigen re-encounter)
\end{itemize}

\paragraph{Class switch recombination
(CSR)}\label{class-switch-recombination-csr}

Initially B cells produce IgM (default isotype). Cytokines from Tfh
cells direct \textbf{class switching} to IgG, IgA, or IgE by recombining
the heavy chain constant region (the variable region --- and thus
antigen specificity --- is preserved). At the molecular level,
\textbf{AID} deaminates cytidines in switch (S) regions upstream of each
constant-region gene; subsequent base-excision repair generates
double-strand breaks that are joined by NHEJ to produce switch
recombination, deleting intervening DNA.

{\def\LTcaptype{none} % do not increment counter
\begin{longtable}[]{@{}
  >{\raggedright\arraybackslash}p{(\linewidth - 4\tabcolsep) * \real{0.3200}}
  >{\raggedright\arraybackslash}p{(\linewidth - 4\tabcolsep) * \real{0.3600}}
  >{\raggedright\arraybackslash}p{(\linewidth - 4\tabcolsep) * \real{0.3200}}@{}}
\toprule\noalign{}
\begin{minipage}[b]{\linewidth}\raggedright
Cytokine
\end{minipage} & \begin{minipage}[b]{\linewidth}\raggedright
Switch to
\end{minipage} & \begin{minipage}[b]{\linewidth}\raggedright
Function
\end{minipage} \\
\midrule\noalign{}
\endhead
\bottomrule\noalign{}
\endlastfoot
IFN-\texorpdfstring{\ensuremath{\gamma}}{gamma} & IgG1, IgG3 & Opsonisation, complement, intracellular
pathogens \\
IL-4 & IgG4, IgE & Allergy, helminth defence \\
TGF-\texorpdfstring{\ensuremath{\beta}}{beta} & IgA & Mucosal immunity \\
IL-21 & IgG1, IgG3 & Synergises with other switches \\
\end{longtable}
}

Defects in AID cause \textbf{hyper-IgM syndrome} (HIGM2 --- failure of
class switching and somatic hypermutation; primarily IgM is produced).
CD40L mutations cause an X-linked form (HIGM1) --- failure of cognate
T-cell help.

\paragraph{Antibody isotypes}\label{antibody-isotypes}

{\def\LTcaptype{none} % do not increment counter
\begin{longtable}[]{@{}
  >{\raggedright\arraybackslash}p{(\linewidth - 6\tabcolsep) * \real{0.1923}}
  >{\raggedright\arraybackslash}p{(\linewidth - 6\tabcolsep) * \real{0.1538}}
  >{\raggedright\arraybackslash}p{(\linewidth - 6\tabcolsep) * \real{0.3462}}
  >{\raggedright\arraybackslash}p{(\linewidth - 6\tabcolsep) * \real{0.3077}}@{}}
\toprule\noalign{}
\begin{minipage}[b]{\linewidth}\raggedright
Class
\end{minipage} & \begin{minipage}[b]{\linewidth}\raggedright
Form
\end{minipage} & \begin{minipage}[b]{\linewidth}\raggedright
Half-life
\end{minipage} & \begin{minipage}[b]{\linewidth}\raggedright
Function
\end{minipage} \\
\midrule\noalign{}
\endhead
\bottomrule\noalign{}
\endlastfoot
\textbf{IgM} & Pentamer & about 5 d & Primary response;
complement activation \\
\textbf{IgG} & Monomer & about 21 d & Secondary response;
opsonisation; placental transfer; complement \\
\textbf{IgA} & Dimer (mucosal) & about 6 d & Mucosal immunity
(gut, respiratory, breast milk) \\
\textbf{IgE} & Monomer & about 2 d (3 weeks bound to Fc\texorpdfstring{\ensuremath{\varepsilon}}{epsilon}R on
mast cells) & Allergy; helminth defence \\
\textbf{IgD} & Monomer & about 3 d & B cell receptor (function
unclear) \\
\end{longtable}
}

\subsubsection{Cytokine Network --- Comprehensive
Reference}\label{cytokine-network-comprehensive-reference}

Cytokines are small (about 15--25 kDa) signalling proteins that
coordinate immune cell function. They act locally (paracrine/autocrine)
at very low concentrations (pM--nM) via JAK/STAT or other receptor
families.

{\def\LTcaptype{none} % do not increment counter
\begin{longtable}[]{@{}
  >{\raggedright\arraybackslash}p{(\linewidth - 8\tabcolsep) * \real{0.1231}}
  >{\raggedright\arraybackslash}p{(\linewidth - 8\tabcolsep) * \real{0.1846}}
  >{\raggedright\arraybackslash}p{(\linewidth - 8\tabcolsep) * \real{0.2000}}
  >{\raggedright\arraybackslash}p{(\linewidth - 8\tabcolsep) * \real{0.3231}}
  >{\raggedright\arraybackslash}p{(\linewidth - 8\tabcolsep) * \real{0.1692}}@{}}
\toprule\noalign{}
\begin{minipage}[b]{\linewidth}\raggedright
Cytokine
\end{minipage} & \begin{minipage}[b]{\linewidth}\raggedright
Major Source
\end{minipage} & \begin{minipage}[b]{\linewidth}\raggedright
Major Targets
\end{minipage} & \begin{minipage}[b]{\linewidth}\raggedright
Receptor / Signalling
\end{minipage} & \begin{minipage}[b]{\linewidth}\raggedright
Key effects
\end{minipage} \\
\midrule\noalign{}
\endhead
\bottomrule\noalign{}
\endlastfoot
\textbf{IL-1 (\texorpdfstring{\ensuremath{\alpha}}{alpha}/\texorpdfstring{\ensuremath{\beta}}{beta})} & Macrophages, DCs (NLRP3 inflammasome for IL-1\texorpdfstring{\ensuremath{\beta}}{beta}) &
Endothelium, hypothalamus, T cells & IL-1R / MyD88 \texorpdfstring{\ensuremath{\to}}{->} NF-\texorpdfstring{\ensuremath{\kappa}}{kappa}B & Fever
(PGE\(_2\)), endothelial activation, T cell co-stimulation \\
\textbf{IL-2} & Activated CD4+ T cells & T cells, NK cells, Tregs &
IL-2R (\texorpdfstring{\ensuremath{\alpha}}{alpha}\texorpdfstring{\ensuremath{\beta}}{beta}\texorpdfstring{\ensuremath{\gamma}}{gamma}) / JAK1/3 \texorpdfstring{\ensuremath{\to}}{->} STAT5 & T cell proliferation; Treg survival (Treg
uses IL-2 as `sink') \\
\textbf{IL-4} & Th2 cells, mast cells, basophils & B cells, Th2
polarisation & IL-4R / JAK1/3 \texorpdfstring{\ensuremath{\to}}{->} STAT6 & IgE class switch; Th2
differentiation; allergy \\
\textbf{IL-5} & Th2 cells, ILC2 & Eosinophils & IL-5R / JAK2 \texorpdfstring{\ensuremath{\to}}{->} STAT5 &
Eosinophil growth/activation (helminths, allergy) \\
\textbf{IL-6} & Macrophages, T cells, hepatocytes & Liver, B cells, T
cells & IL-6R/gp130 / JAK1/2 \texorpdfstring{\ensuremath{\to}}{->} STAT3 & Acute phase response (CRP,
fibrinogen); Th17 differentiation; B cell maturation \\
\textbf{IL-7} & Stromal cells (BM, thymus) & Naïve and memory T cells &
IL-7R / JAK1/3 \texorpdfstring{\ensuremath{\to}}{->} STAT5 & T cell development and homeostatic survival \\
\textbf{IL-8 (CXCL8)} & Macrophages, endothelium & Neutrophils &
CXCR1/CXCR2 (G\(_i\) GPCRs) & Neutrophil chemotaxis \\
\textbf{IL-10} & Tregs, macrophages, B cells & Macrophages, T cells &
IL-10R / JAK1/TYK2 \texorpdfstring{\ensuremath{\to}}{->} STAT3 & \textbf{Anti-inflammatory} --- suppresses
Th1 and macrophage activation \\
\textbf{IL-12} & Macrophages, DCs & NK cells, T cells & IL-12R /
TYK2/JAK2 \texorpdfstring{\ensuremath{\to}}{->} STAT4 & Th1 differentiation; IFN-\texorpdfstring{\ensuremath{\gamma}}{gamma} induction \\
\textbf{IL-13} & Th2, ILC2, mast cells & B cells, smooth muscle,
epithelium & IL-13R\texorpdfstring{\ensuremath{\alpha}}{alpha}1/IL-4R\texorpdfstring{\ensuremath{\alpha}}{alpha} / STAT6 & Allergy/asthma; goblet cell
mucus; tissue remodelling; target of dupilumab \\
\textbf{IL-15} & DCs, monocytes & NK, memory CD8+ T cells & IL-15R /
JAK1/3 \texorpdfstring{\ensuremath{\to}}{->} STAT5 & NK and memory T cell maintenance \\
\textbf{IL-17} & Th17 cells, \texorpdfstring{\ensuremath{\gamma}}{gamma}\texorpdfstring{\ensuremath{\delta}}{delta} T cells & Epithelium, neutrophils &
IL-17R / Act1 \texorpdfstring{\ensuremath{\to}}{->} NF-\texorpdfstring{\ensuremath{\kappa}}{kappa}B & Mucocutaneous defence (fungi, extracellular
bacteria); psoriasis, IBD when dysregulated \\
\textbf{IL-21} & Tfh cells & B cells, T cells & IL-21R / JAK1/3 \texorpdfstring{\ensuremath{\to}}{->} STAT3
& Germinal centre reactions, class switching \\
\textbf{IL-22} & Th17, Th22, ILC3 & Epithelium & IL-22R / JAK1/2 \texorpdfstring{\ensuremath{\to}}{->} STAT3
& Mucosal barrier defence; antimicrobial peptide production \\
\textbf{IL-23} & Macrophages, DCs & Th17 cells & IL-23R / JAK2/TYK2 \texorpdfstring{\ensuremath{\to}}{->}
STAT3 & Th17 stabilisation; target of ustekinumab/risankizumab in
psoriasis/IBD \\
\textbf{TNF-\texorpdfstring{\ensuremath{\alpha}}{alpha}} & Macrophages, T cells & Endothelium, tumour cells &
TNFR1/2 \texorpdfstring{\ensuremath{\to}}{->} NF-\texorpdfstring{\ensuremath{\kappa}}{kappa}B / apoptosis & Inflammation, septic shock, apoptosis;
target of anti-TNF biologics \\
\textbf{IFN-\texorpdfstring{\ensuremath{\alpha}}{alpha}/\texorpdfstring{\ensuremath{\beta}}{beta} (Type I)} & Most cells (esp.~plasmacytoid DCs) & Most
nucleated cells & IFNAR1/2 / JAK1/TYK2 \texorpdfstring{\ensuremath{\to}}{->} STAT1/2 \texorpdfstring{\ensuremath{\to}}{->} ISGF3 & Antiviral
state (induces hundreds of ISGs); MHC I upregulation \\
\textbf{IFN-\texorpdfstring{\ensuremath{\gamma}}{gamma} (Type II)} & Th1 cells, NK cells, CD8+ & Macrophages, most
cells & IFNGR1/2 / JAK1/2 \texorpdfstring{\ensuremath{\to}}{->} STAT1 & Macrophage activation; MHC
upregulation; intracellular pathogen defence \\
\textbf{TGF-\texorpdfstring{\ensuremath{\beta}}{beta}} & Tregs, many cells & Many & T\texorpdfstring{\ensuremath{\beta}}{beta}R-I/II / SMAD2/3 & Treg
induction; tissue repair; wound healing; oncogenic in late cancer \\
\textbf{GM-CSF} & T cells, macrophages, endothelium & Myeloid
progenitors & CSF2R / JAK2 \texorpdfstring{\ensuremath{\to}}{->} STAT5 & Granulocyte/macrophage growth; DC
differentiation; therapeutic in neutropenia \\
\end{longtable}
}

\textbf{Th polarisation:}

\begin{itemize}
\tightlist
\item
  \textbf{Th1} (driven by IL-12; produces IFN-\texorpdfstring{\ensuremath{\gamma}}{gamma}): intracellular
  pathogens, viruses, autoimmunity (MS, T1DM)
\item
  \textbf{Th2} (driven by IL-4; produces IL-4/5/13): helminths, allergy,
  asthma
\item
  \textbf{Th17} (driven by IL-6 + TGF-\texorpdfstring{\ensuremath{\beta}}{beta}; produces IL-17): extracellular
  bacteria, fungi; psoriasis, IBD
\item
  \textbf{Treg} (driven by TGF-\texorpdfstring{\ensuremath{\beta}}{beta}; FoxP3+): peripheral tolerance,
  suppression
\item
  \textbf{Tfh} (in germinal centres; produces IL-21): B cell help
\end{itemize}

\subsubsection{Immunological Memory --- Why Memory Cells Respond
Faster}\label{immunological-memory-why-memory-cells-respond-faster}

After antigen clearance, about 95\% of effector T and B cells
die by apoptosis (contraction phase). The remaining 5\% become
long-lived memory cells. Memory cells respond to re-exposure 100--1000\texorpdfstring{\ensuremath{\times}}{x}
faster and stronger than naïve cells:

{\def\LTcaptype{none} % do not increment counter
\begin{longtable}[]{@{}
  >{\raggedright\arraybackslash}p{(\linewidth - 4\tabcolsep) * \real{0.2759}}
  >{\raggedright\arraybackslash}p{(\linewidth - 4\tabcolsep) * \real{0.3448}}
  >{\raggedright\arraybackslash}p{(\linewidth - 4\tabcolsep) * \real{0.3793}}@{}}
\toprule\noalign{}
\begin{minipage}[b]{\linewidth}\raggedright
Property
\end{minipage} & \begin{minipage}[b]{\linewidth}\raggedright
Naïve cell
\end{minipage} & \begin{minipage}[b]{\linewidth}\raggedright
Memory cell
\end{minipage} \\
\midrule\noalign{}
\endhead
\bottomrule\noalign{}
\endlastfoot
\textbf{Frequency for given antigen} & about 1 in 10\(^6\) &
100--1000-fold higher (clonal expansion preserved) \\
\textbf{Activation threshold} & High (needs 24--48 h priming, full
co-stimulation) & Low (rapid proliferation on re-encounter) \\
\textbf{Effector repertoire} & Naïve, must differentiate & Pre-armed;
rapid cytokine production / class-switched antibody \\
\textbf{Tissue distribution} & Recirculate through 2° lymphoid organs &
T\(_{CM}\) recirculate; T\(_{EM}\)/T\(_{RM}\) patrol peripheral and
barrier tissues \\
\textbf{Homeostatic maintenance} & None & Tonic IL-7 + IL-15; long
lifespan (decades for some clones) \\
\end{longtable}
}

The signalling determinants of effector vs memory differentiation
include:

\begin{itemize}
\tightlist
\item
  \textbf{Strong, prolonged TCR signalling + IL-2} \texorpdfstring{\ensuremath{\to}}{->} effector
  (terminally differentiated, short-lived, KLRG1\(^{hi}\))
\item
  \textbf{Brief TCR signalling + IL-7, IL-15} \texorpdfstring{\ensuremath{\to}}{->} memory precursors
  (CD127\(^{hi}\), CD62L\(^+\))
\item
  \textbf{Memory subsets:}

  \begin{itemize}
  \tightlist
  \item
    \textbf{Central memory (T\(_{CM}\)):} CCR7\(^+\), CD62L\(^+\);
    recirculate through lymph nodes; high proliferative potential
  \item
    \textbf{Effector memory (T\(_{EM}\)):} CCR7\(^-\); patrol peripheral
    tissues; immediate effector function
  \item
    \textbf{Tissue-resident memory (T\(_{RM}\)):} CD69\(^+\),
    CD103\(^+\); reside in barrier tissues (skin, gut, lung); first to
    detect re-infection
  \end{itemize}
\end{itemize}

Memory persistence depends on tonic IL-7 and IL-15 cytokine signalling
and homeostatic proliferation, not antigen re-stimulation. This is why
vaccine immunity can last decades despite no reinfection --- and why
immunosuppressive therapies that cripple homeostatic cytokines (e.g.,
IL-7 axis blockade) reduce memory.

\subsubsection{Tolerance and
Autoimmunity}\label{tolerance-and-autoimmunity}

Self-tolerance is maintained at two levels:

Autoimmunity is rarely explained by a single broken checkpoint.
Susceptibility alleles, sex-biased immune regulation, molecular mimicry,
epitope spreading, tissue damage, microbiome state, and timing of
infection can most shift the threshold between protective immunity and
self-damage. A precise answer should therefore distinguish the
initiating trigger, the self-antigen or tissue compartment, the effector
mechanism, and the evidence that tolerance failed rather than merely
inflammation increased.

\paragraph{Central tolerance (thymus and bone
marrow)}\label{central-tolerance-thymus-and-bone-marrow}

\begin{itemize}
\tightlist
\item
  T cells: Negative selection in thymic medulla (AIRE-driven
  self-antigen expression)
\item
  B cells: Receptor editing and clonal deletion in bone marrow upon BCR
  engagement of self-antigen
\end{itemize}

Failure modes: AIRE mutations (APS-1); central tolerance is incomplete
because not most self-antigens can be expressed in primary lymphoid
organs.

\paragraph{Peripheral tolerance}\label{peripheral-tolerance}

\begin{itemize}
\tightlist
\item
  \textbf{Anergy:} TCR engagement without co-stimulation (CD28 by B7)
  inactivates rather than activates the T cell.
\item
  \textbf{Treg suppression:} FoxP3+ regulatory T cells suppress
  conventional T cell activation by IL-2 consumption, IL-10/TGF-\texorpdfstring{\ensuremath{\beta}}{beta}
  secretion, CTLA-4 outcompeting CD28, and granzyme killing.
\item
  \textbf{Inhibitory checkpoints:} \textbf{CTLA-4} and \textbf{PD-1}
  restrain T cell responses (see immunotherapy below).
\item
  \textbf{Activation-induced cell death (AICD):} Repeatedly activated T
  cells upregulate Fas \texorpdfstring{\ensuremath{\to}}{->} undergo apoptosis on Fas ligation.
\end{itemize}

\paragraph{Mechanisms of autoimmunity}\label{mechanisms-of-autoimmunity}

Failures of tolerance produce autoimmunity through several archetypal
mechanisms:

\begin{itemize}
\tightlist
\item
  \textbf{Molecular mimicry:} A pathogen antigen shares structural
  similarity with self-antigen. Memory cells generated to fight the
  pathogen cross-react with self.

  \begin{itemize}
  \tightlist
  \item
    \emph{Example:} Group A \emph{Streptococcus} M protein mimics
    cardiac myosin \texorpdfstring{\ensuremath{\to}}{->} \textbf{rheumatic heart disease} after
    streptococcal pharyngitis.
  \item
    \emph{Example:} \emph{Campylobacter jejuni} lipo-oligosaccharide
    mimics peripheral nerve gangliosides \texorpdfstring{\ensuremath{\to}}{->} \textbf{Guillain-Barré
    syndrome}.
  \end{itemize}
\item
  \textbf{Bystander activation:} Local infection/inflammation drives
  massive cytokine release, MHC upregulation, and APC maturation.
  Autoreactive T cells that escaped thymic deletion are activated by
  their cognate self-antigen now presented in a strongly stimulatory
  context.
\item
  \textbf{Epitope spreading:} Initial response to one self-antigen
  broadens over time as tissue damage releases additional self-antigens
  that are taken up by APCs and presented to new autoreactive T cell
  clones.

  \begin{itemize}
  \tightlist
  \item
    \emph{Example:} In \textbf{Type 1 diabetes}, autoantibodies often
    appear in stereotyped sequence (insulin \texorpdfstring{\ensuremath{\to}}{->} GAD65 \texorpdfstring{\ensuremath{\to}}{->} IA-2 \texorpdfstring{\ensuremath{\to}}{->} ZnT8) over
    months to years before clinical onset, reflecting progressive
    epitope spreading within the islet.
  \item
    \emph{Example:} In \textbf{multiple sclerosis}, initial myelin basic
    protein (MBP) response can spread to PLP, MOG, and other myelin
    antigens.
  \end{itemize}
\item
  \textbf{Loss of regulatory T cells:} FoxP3 mutations cause IPEX
  syndrome (Immunodysregulation Polyendocrinopathy Enteropathy X-linked)
  --- fatal multi-organ autoimmunity in infancy.
\item
  \textbf{Defective apoptosis:} Fas / FasL mutations (autoimmune
  lymphoproliferative syndrome, ALPS) prevent AICD \texorpdfstring{\ensuremath{\to}}{->} accumulation of
  self-reactive lymphocytes.
\end{itemize}

{\def\LTcaptype{none} % do not increment counter
\begin{longtable}[]{@{}
  >{\raggedright\arraybackslash}p{(\linewidth - 6\tabcolsep) * \real{0.3953}}
  >{\raggedright\arraybackslash}p{(\linewidth - 6\tabcolsep) * \real{0.3023}}
  >{\raggedright\arraybackslash}p{(\linewidth - 6\tabcolsep) * \real{0.0930}}
  >{\raggedright\arraybackslash}p{(\linewidth - 6\tabcolsep) * \real{0.2093}}@{}}
\toprule\noalign{}
\begin{minipage}[b]{\linewidth}\raggedright
Autoimmune disease
\end{minipage} & \begin{minipage}[b]{\linewidth}\raggedright
Target tissue
\end{minipage} & \begin{minipage}[b]{\linewidth}\raggedright
Type
\end{minipage} & \begin{minipage}[b]{\linewidth}\raggedright
Mechanism
\end{minipage} \\
\midrule\noalign{}
\endhead
\bottomrule\noalign{}
\endlastfoot
Type 1 diabetes & \texorpdfstring{\ensuremath{\beta}}{beta}-cells & Cell-mediated (CD8+) & Anti-GAD65, anti-IA-2
antibodies; T cell destruction; epitope spreading \\
Multiple sclerosis & CNS myelin & Cell-mediated (Th1/Th17) & Anti-MBP T
cells; molecular mimicry to EBV nuclear antigen 1 implicated \\
Rheumatoid arthritis & Synovial joints & Mixed & Anti-CCP, RF
antibodies; TNF-driven \\
SLE & Multi-organ & Antibody/IC & Anti-dsDNA, anti-Sm; immune complex
deposition \\
Hashimoto's & Thyroid & Antibody & Anti-TPO, anti-Tg \\
Graves' & Thyroid & Antibody (stimulating) & Anti-TSH-R (agonist) \\
Myasthenia gravis & NMJ & Antibody & Anti-AChR antibodies block
receptors \\
Goodpasture's & Kidney, lung & Antibody (Type II) & Anti-\texorpdfstring{\ensuremath{\alpha}}{alpha}3(IV)
collagen \\
\end{longtable}
}

\subsubsection{Hypersensitivity Reactions (Gell-Coombs
Classification)}\label{hypersensitivity-reactions-gell-coombs-classification}

The four-tier Gell-Coombs classification organises immune-mediated
hypersensitivity by mechanism, time course, and treatment.

{\def\LTcaptype{none} % do not increment counter
\begin{longtable}[]{@{}
  >{\raggedright\arraybackslash}p{(\linewidth - 8\tabcolsep) * \real{0.1143}}
  >{\raggedright\arraybackslash}p{(\linewidth - 8\tabcolsep) * \real{0.2571}}
  >{\raggedright\arraybackslash}p{(\linewidth - 8\tabcolsep) * \real{0.1429}}
  >{\raggedright\arraybackslash}p{(\linewidth - 8\tabcolsep) * \real{0.2286}}
  >{\raggedright\arraybackslash}p{(\linewidth - 8\tabcolsep) * \real{0.2571}}@{}}
\toprule\noalign{}
\begin{minipage}[b]{\linewidth}\raggedright
Type
\end{minipage} & \begin{minipage}[b]{\linewidth}\raggedright
Mechanism
\end{minipage} & \begin{minipage}[b]{\linewidth}\raggedright
Onset
\end{minipage} & \begin{minipage}[b]{\linewidth}\raggedright
Examples
\end{minipage} & \begin{minipage}[b]{\linewidth}\raggedright
Treatment
\end{minipage} \\
\midrule\noalign{}
\endhead
\bottomrule\noalign{}
\endlastfoot
\textbf{Type I (immediate)} & IgE-mediated; mast cell degranulation;
histamine, leukotrienes, tryptase & Minutes & Anaphylaxis, allergic
asthma, food allergy, hay fever & Epinephrine (anaphylaxis),
antihistamines (H1 blockers --- diphenhydramine, cetirizine), inhaled
corticosteroids, cromolyn (mast cell stabiliser), allergen immunotherapy
(gradual desensitisation), anti-IgE (omalizumab) \\
\textbf{Type II (cytotoxic)} & IgG/IgM antibody against cell-surface
antigen \texorpdfstring{\ensuremath{\to}}{->} complement, ADCC & Hours & Autoimmune haemolytic anaemia,
transfusion reaction, Goodpasture's, Graves' (stimulating) &
Plasmapheresis, IVIG, immunosuppression (rituximab anti-CD20 for
ITP/AIHA) \\
\textbf{Type III (immune complex)} & Soluble antigen-antibody complexes
deposit in tissues \texorpdfstring{\ensuremath{\to}}{->} complement, neutrophils & Hours to days & SLE,
serum sickness, post-streptococcal glomerulonephritis, hypersensitivity
pneumonitis & Corticosteroids, cyclophosphamide, removal of antigen
source, plasmapheresis \\
\textbf{Type IV (delayed/cell-mediated)} & T cell-mediated; no antibody
& 24--72 h & Contact dermatitis (poison ivy, nickel), tuberculin/PPD
test, Type 1 diabetes onset, MS, contact allergy to nickel &
Topical/systemic corticosteroids, calcineurin inhibitors (tacrolimus),
allergen avoidance \\
\end{longtable}
}

\paragraph{Type I --- molecular detail}\label{type-i-molecular-detail}

Sensitisation: First encounter with allergen \texorpdfstring{\ensuremath{\to}}{->} APCs prime Th2 \texorpdfstring{\ensuremath{\to}}{->} IL-4
drives B cells to class-switch to \textbf{IgE}. IgE binds \textbf{Fc\texorpdfstring{\ensuremath{\varepsilon}}{epsilon}RI}
on mast cells and basophils with extraordinarily high affinity (K\(_d\)
about 10\(^{-10}\) M) and remains bound for weeks.

Effector: Re-encounter with allergen cross-links mast cell-bound IgE \texorpdfstring{\ensuremath{\to}}{->}
degranulation within seconds: \textbf{histamine} (vasodilation, vascular
leak, smooth muscle contraction, pruritus), \textbf{tryptase} (tissue
damage), \textbf{leukotrienes (C4, D4, E4)} (delayed
bronchoconstriction), \textbf{prostaglandins}, \textbf{TNF},
\textbf{IL-4/IL-5/IL-13} (late phase, 4--8 h: eosinophil and Th2
recruitment).

Clinical spectrum: localised (allergic rhinitis, urticaria) \texorpdfstring{\ensuremath{\to}}{->} systemic
(anaphylaxis with hypotension, bronchospasm, laryngeal oedema).
Treatment of anaphylaxis: \textbf{intramuscular epinephrine} (\texorpdfstring{\ensuremath{\alpha}}{alpha}\(_1\)
vasoconstriction reverses hypotension; \texorpdfstring{\ensuremath{\beta}}{beta}\(_2\) bronchodilates; \texorpdfstring{\ensuremath{\beta}}{beta}\(_2\)
stabilises mast cells), oxygen, fluids, antihistamines, corticosteroids,
possible airway management. \textbf{Allergen immunotherapy (AIT)}
gradually shifts the response from Th2/IgE toward Th1/IgG4 (``blocking
antibodies'') and induces Tregs --- the primary disease-modifying
therapy for allergy.

\subsubsection{Immunotherapy}\label{immunotherapy}

Modern oncology has been transformed by therapies that release brakes on
antitumour immunity (\textbf{checkpoint inhibitors}) or redirect T cells
(\textbf{CAR-T}). Because the product roster and safety language change,
treatment tables below should be read with current FDA product pages and
labels, not as a permanent formulary
\citep{fda2026cellulargeneproducts,fda2024cartwarning}.

\paragraph{Immune checkpoint inhibitors --- PD-1/PD-L1 and CTLA-4
biology}\label{immune-checkpoint-inhibitors-pd-1pd-l1-and-ctla-4-biology}

\textbf{PD-1 (programmed death 1)} is an inhibitory receptor expressed
on activated T cells (and B cells, NK cells, Tregs). Its ligand
\textbf{PD-L1} is widely expressed under inflammatory conditions and
constitutively on many tumours; \textbf{PD-L2} is more restricted to
APCs. PD-1 engagement recruits the SHP-2 phosphatase to the immune
synapse, dephosphorylating proximal TCR-signalling components (CD3\texorpdfstring{\ensuremath{\zeta}}{zeta},
ZAP70) --- this functionally ``exhausts'' chronically activated T cells
and limits collateral damage. Tumours hijack this brake by upregulating
PD-L1 (often in response to local IFN-\texorpdfstring{\ensuremath{\gamma}}{gamma} from infiltrating T cells --- a
paradoxical ``adaptive resistance''), tonically inhibiting
tumour-infiltrating T cells. Current drug examples and approved-use
language should be checked against NCI or regulatory pages before
publication because checkpoint indications change by tumour type,
biomarker, and line of therapy \citep{nci2026checkpointinhibitors}.

\textbf{CTLA-4 (cytotoxic T-lymphocyte-associated 4, CD152)} is also
inhibitory but acts at the priming stage in lymph nodes. Resting T cells
receive co-stimulation when CD28 binds B7-1/B7-2 (CD80/CD86) on APCs.
Upon activation, CTLA-4 is upregulated and translocates to the immune
synapse, where it binds B7 with much higher affinity than CD28 and
trans-endocytoses B7 from the APC --- depriving CD28 of its co-stimulus.
CTLA-4 is also constitutively expressed on Tregs, contributing to their
suppressive function. CTLA-4 blockade therefore acts earlier and more
broadly than PD-1 blockade, with stronger autoimmune toxicity.

\begin{figure}[htbp]
\centering
\includegraphics[width=0.82\linewidth,height=4.2in,keepaspectratio]{../figures/mermaid_inline/inline_mermaid_0172_2aed8e7e5564.png}
\caption{Immune checkpoint blockade Normal T cells use CTLA-4 and PD-1 as brakes to prevent over-activation. Tumours exploit PD-L1 to suppress local T cells. Anti-PD-1 antibodies (pembrolizumab) and anti-CTLA-4 antibodies (ipilimumab) block these interactions, releasing the immune brake on antitumour T cells.}
\end{figure}

{\def\LTcaptype{none} % do not increment counter
\begin{longtable}[]{@{}
  >{\raggedright\arraybackslash}p{(\linewidth - 6\tabcolsep) * \real{0.0784}}
  >{\raggedright\arraybackslash}p{(\linewidth - 6\tabcolsep) * \real{0.1176}}
  >{\raggedright\arraybackslash}p{(\linewidth - 6\tabcolsep) * \real{0.3922}}
  >{\raggedright\arraybackslash}p{(\linewidth - 6\tabcolsep) * \real{0.4118}}@{}}
\toprule\noalign{}
\begin{minipage}[b]{\linewidth}\raggedright
Drug
\end{minipage} & \begin{minipage}[b]{\linewidth}\raggedright
Target
\end{minipage} & \begin{minipage}[b]{\linewidth}\raggedright
Approved indications
\end{minipage} & \begin{minipage}[b]{\linewidth}\raggedright
Approx. response rate
\end{minipage} \\
\midrule\noalign{}
\endhead
\bottomrule\noalign{}
\endlastfoot
\textbf{Ipilimumab} & CTLA-4 & Melanoma, RCC (combined with nivolumab) &
about 10--15\% monotherapy, about 50\% with
anti-PD-1 \\
\textbf{Nivolumab, pembrolizumab} & PD-1 & Melanoma, NSCLC, RCC,
head/neck, Hodgkin, MMR-deficient cancers & about 20--40\%
across indications \\
\textbf{Atezolizumab, durvalumab} & PD-L1 & NSCLC, urothelial, TNBC &
about 15--30\% \\
\textbf{Relatlimab} & LAG-3 & Melanoma (combined with nivolumab) &
Combination effect; LAG-3 alone modest \\
\textbf{Tislelizumab, cemiplimab} & PD-1 & NSCLC, cutaneous SCC &
Variable \\
\end{longtable}
}

The 2018 Nobel Prize in Physiology or Medicine was awarded to James
Allison (CTLA-4) and Tasuku Honjo (PD-1) for the discoveries underlying
this field. Response rates vary widely by tumour type, biomarker status,
line of therapy, and combination regimen; the table's approximate ranges
are orientation values, not patient-level predictions. A key qualitative
pattern remains \textbf{durable remissions} in some responders (the
``tail of the curve''). Side effects: \textbf{immune-related adverse
events (irAEs)} --- colitis, pneumonitis, hypophysitis, thyroiditis,
hepatitis, dermatitis --- reflecting the breaking of self-tolerance.
Management: corticosteroids; if severe, infliximab (anti-TNF) for
colitis.

\paragraph{CAR-T cell therapy --- design and clinical
outcomes}\label{car-t-cell-therapy-design-and-clinical-outcomes}

\textbf{Chimeric antigen receptor T cells:} Patient T cells are removed
by leukapheresis, genetically engineered ex vivo (lentiviral or
retroviral vector) to express a synthetic receptor combining:

\begin{itemize}
\tightlist
\item
  \textbf{Extracellular antigen-binding domain:} Single-chain variable
  fragment (scFv) derived from a tumour-specific monoclonal antibody
  (e.g., anti-CD19 scFv from FMC63 antibody).
\item
  \textbf{Hinge and transmembrane:} From CD8\texorpdfstring{\ensuremath{\alpha}}{alpha} or IgG.
\item
  \textbf{Co-stimulatory domain:} CD28 (rapid effector function, shorter
  persistence) or 4-1BB/CD137 (slower expansion, longer persistence).
\item
  \textbf{Activation domain:} CD3\texorpdfstring{\ensuremath{\zeta}}{zeta} (TCR-derived ITAMs).
\end{itemize}

Modified T cells are expanded in vitro and infused back into the
patient. The CAR-T cells recognise the surface antigen (e.g., CD19 on
B-cell malignancies) \textbf{without MHC restriction}, which enables
recognition of MHC-low tumours that escape conventional T cells.

{\def\LTcaptype{none} % do not increment counter
\begin{longtable}[]{@{}
  >{\raggedright\arraybackslash}p{(\linewidth - 8\tabcolsep) * \real{0.1296}}
  >{\raggedright\arraybackslash}p{(\linewidth - 8\tabcolsep) * \real{0.1111}}
  >{\raggedright\arraybackslash}p{(\linewidth - 8\tabcolsep) * \real{0.3519}}
  >{\raggedright\arraybackslash}p{(\linewidth - 8\tabcolsep) * \real{0.1296}}
  >{\raggedright\arraybackslash}p{(\linewidth - 8\tabcolsep) * \real{0.2778}}@{}}
\toprule\noalign{}
\begin{minipage}[b]{\linewidth}\raggedright
Product
\end{minipage} & \begin{minipage}[b]{\linewidth}\raggedright
Target
\end{minipage} & \begin{minipage}[b]{\linewidth}\raggedright
Approved indication
\end{minipage} & \begin{minipage}[b]{\linewidth}\raggedright
Co-stim
\end{minipage} & \begin{minipage}[b]{\linewidth}\raggedright
Remission rate
\end{minipage} \\
\midrule\noalign{}
\endhead
\bottomrule\noalign{}
\endlastfoot
\textbf{Tisagenlecleucel (Kymriah)} & CD19 & Pediatric/young adult
B-cell acute lymphoblastic leukaemia & 4-1BB & about 80\% at 3
months \\
\textbf{Axicabtagene ciloleucel (Yescarta)} & CD19 & DLBCL (3rd line+,
then earlier) & CD28 & about 50--60\% complete response \\
\textbf{Brexucabtagene autoleucel (Tecartus)} & CD19 & Mantle cell
lymphoma; B-cell acute lymphoblastic leukaemia & CD28 &
about 67\% in MCL \\
\textbf{Idecabtagene vicleucel (Abecma)} & BCMA & Multiple myeloma &
4-1BB & about 30\% complete response \\
\textbf{Ciltacabtagene autoleucel (Carvykti)} & BCMA & Multiple myeloma
& 4-1BB & about 80\% overall response \\
\end{longtable}
}

\textbf{Toxicity:}

\begin{itemize}
\tightlist
\item
  \textbf{Cytokine release syndrome (CRS):} fever, hypotension, hypoxia
  within hours-to-days from massive cytokine release (esp.~IL-6, IFN-\texorpdfstring{\ensuremath{\gamma}}{gamma},
  TNF) --- graded 1--4. Treatment: \textbf{tocilizumab} (anti-IL-6R),
  corticosteroids if severe.
\item
  \textbf{Immune effector cell-associated neurotoxicity syndrome
  (ICANS):} confusion, aphasia, seizures, cerebral oedema; treated with
  corticosteroids.
\item
  \textbf{B cell aplasia:} Anti-CD19 CAR-T kills normal B cells too;
  long-term IVIG replacement may be needed.
\item
  \textbf{Second-malignancy surveillance:} FDA required boxed-warning
  language for T-cell malignancies after BCMA-directed or CD19-directed
  autologous CAR-T products in 2024, making long-term follow-up part of
  the source-governance burden for any current CAR-T table
  \citep{fda2024cartwarning}.
\end{itemize}

\begin{quote}
\textbf{Concept Check 4:} A patient on pembrolizumab (anti-PD-1) for
metastatic melanoma develops new-onset hyperthyroidism with elevated
free T4 and undetectable TSH. Explain mechanistically why checkpoint
blockade can trigger autoimmune endocrinopathy. What does this teach us
about the normal role of PD-1 in maintaining peripheral tolerance to
thyroid antigens?
\end{quote}

\begin{quote}
\textbf{Concept Check 5:} Why do polysaccharide vaccines (e.g., the
original 23-valent pneumococcal vaccine) fail in children younger than 2
years, while conjugate vaccines (PCV13 --- same polysaccharides
covalently linked to a carrier protein) work well? Connect your answer
to T-dependent vs T-independent B cell activation and germinal centre
formation.
\end{quote}

\begin{quote}
\textbf{Concept Check 6:} A patient with Type 1 diabetes is found to
have autoantibodies against multiple islet antigens (GAD65, IA-2, ZnT8)
in addition to the originally targeted insulin. Explain this observation
in terms of \textbf{epitope spreading} and discuss why early
immunomodulatory therapy (before extensive \texorpdfstring{\ensuremath{\beta}}{beta}-cell destruction) might be
more effective than later intervention.
\end{quote}

\begin{quote}
\textbf{Concept Check 7:} During a CAR-T infusion for B-cell acute
lymphoblastic leukaemia, the patient develops high fever, hypotension,
and elevated IL-6. Why does CAR-T therapy frequently cause cytokine
release syndrome, and why is anti-IL-6R (tocilizumab) effective without
abrogating the antitumour response?
\end{quote}

\begin{quote}
\textbf{Concept Check 8:} Compare central and peripheral tolerance. A
patient with FoxP3 mutation (IPEX syndrome) develops autoimmunity in
infancy despite intact thymic negative selection. What does this teach
us about the redundancy and complementarity of these two tolerance
mechanisms?
\end{quote}

\begin{quote}
\textbf{Concept Check 9:} A patient with hereditary angioedema has
frequent episodes of facial/laryngeal swelling. Genetic testing reveals
a C1-INH (C1-inhibitor) deficiency. Why does deficiency of this
complement regulator cause swelling --- and why is the swelling
primarily mediated by \textbf{bradykinin} rather than complement-derived
anaphylatoxins?
\end{quote}

\begin{quote}
\textbf{Concept Check 10:} Why do anti-IL-17 (secukinumab) and
anti-IL-23 (risankizumab) biologics treat psoriasis but exacerbate
Crohn's disease in some patients? Connect your answer to the dual role
of Th17 in barrier defence and pathological inflammation.
\end{quote}

\begin{center}\rule{0.5\linewidth}{0.5pt}\end{center}

\subsection{Worked Example}\label{worked-example-4}

\textbf{Problem:} A 35-year-old woman is found to have an elevated total
cortisol level of 28 \texorpdfstring{\ensuremath{\mu}}{mu}g/dL on a morning blood draw. CBG
(cortisol-binding globulin) is also elevated at 60 \texorpdfstring{\ensuremath{\mu}}{mu}g/mL (normal:
about 30 \texorpdfstring{\ensuremath{\mu}}{mu}g/mL). She is six months pregnant. Estimate her
\textbf{free (bioactive) cortisol} assuming free cortisol normally
accounts for about 5\% of total at baseline CBG. Should she be
diagnosed with Cushing's syndrome?

\textbf{Solution:}

\textbf{Step 1.} Recognise that pregnancy elevates oestrogen, which
increases hepatic CBG synthesis about 2-fold. Total cortisol
rises in parallel with CBG, but \textbf{free cortisol} --- the
biologically active fraction --- remains in the normal range.

\textbf{Step 2.} Estimate baseline free cortisol fraction. With normal
CBG, 5\% of total cortisol is free. With doubled CBG, the free fraction
drops because more cortisol is sequestered:

\[\text{free fraction} \propto \frac{1}{1 + [\text{CBG}]/K_d^{\text{CBG}}} \label{eq:unit_IX_endocrine_and_immune_item_1}\]

\[\text{free fraction (pregnancy)} \approx \frac{5\%}{2} = 2.5\% \label{eq:unit_IX_endocrine_and_immune_item_2}\]

\textbf{Step 3.} Calculate free cortisol:

\[[\text{cortisol}]_{\text{free}} = 0.025 \times 28\;\mu\text{g/dL} = 0.7\;\mu\text{g/dL} \label{eq:unit_IX_endocrine_and_immune_item_3}\]

This is within the normal range (free cortisol about 0.5--1.0
\texorpdfstring{\ensuremath{\mu}}{mu}g/dL).

\textbf{Answer:} The elevated total cortisol reflects increased CBG, not
pathological cortisol excess. \textbf{Free cortisol} is normal. The
patient should not be diagnosed with Cushing's syndrome based on total
cortisol alone. The appropriate test in pregnancy is \textbf{24-h
urinary free cortisol} (which measures the unbound fraction excreted by
the kidney) or salivary cortisol (free cortisol equilibrates with
saliva).

This illustrates a fundamental principle: \textbf{for protein-bound
hormones, total plasma concentration is misleading; primarily the free
fraction is biologically active.} The same applies to thyroid hormones
(TBG changes during pregnancy), testosterone (SHBG), and many others.

\begin{center}\rule{0.5\linewidth}{0.5pt}\end{center}

\subsection{Worked Example: Clonal
Expansion}\label{worked-example-clonal-expansion}

\textbf{Problem:} During a T-dependent germinal-centre response, an
antigen-specific B-cell clone proliferates in the dark zone with a
doubling time of \(t_d = 6\) h. Starting from \(N_0 = 100\)
antigen-specific cells, calculate the fold expansion and the absolute
clone size after \(t = 72\) h (3 days). Compare the result with the
100--1000-fold higher precursor frequency that defines an established
memory pool.

\textbf{Solution:}

\textbf{Step 1.} State the exponential-growth relation for a clone
dividing with a fixed doubling time. Each doubling multiplies cell
number by 2, and the number of doublings in time \(t\) is \(t/t_d\):

\[N = N_0 \cdot 2^{\,t/t_d} \label{eq:unit_IX_endocrine_and_immune_clonal_1}\]

\textbf{Step 2.} Substitute the measured values. With \(t = 72\) h and
\(t_d = 6\) h, the number of doublings is \(t/t_d = 72/6 = 12\). The
fold expansion is therefore \(2^{12}\):

\[\frac{N}{N_0} = 2^{72/6} = 2^{12} = 4096 \label{eq:unit_IX_endocrine_and_immune_clonal_2}\]

\textbf{Step 3.} Compute the absolute clone size:
\(N = 100 \times 4096 = 4.10 \times 10^{5}\) cells --- a roughly
4000-fold expansion in three days.

\textbf{Answer:} A 6-hour doubling time yields 12 doublings in 72 h,
expanding the clone about 4096-fold (from 100 to
about 4.1 \texorpdfstring{\ensuremath{\times}}{x} 10\(^5\) cells). This magnitude matches the
chapter's account of dark-zone proliferation and explains how a naive
precursor frequency of about 1 in 10\(^6\) can be driven up to
the 100--1000-fold higher frequency that characterises an established
memory pool. The same \(N = N_0 \cdot 2^{\,t/t_d}\) relation predicts
that even a one-hour increase in doubling time markedly slows the
response --- which is why memory cells, with their lower activation
threshold and shorter effective doubling time, clear a re-challenge far
faster than naive cells.

\begin{center}\rule{0.5\linewidth}{0.5pt}\end{center}

\subsection{Computational Bridge}\label{computational-bridge-30}

Generic negative-feedback controllers map errors to corrective outputs:

\begin{Shaded}
\begin{Highlighting}[]
\ImportTok{from}\NormalTok{ biology.physiology }\ImportTok{import}\NormalTok{ homeostasis\_response}

\NormalTok{h }\OperatorTok{=}\NormalTok{ homeostasis\_response(}\FloatTok{37.0}\NormalTok{, }\FloatTok{38.5}\NormalTok{, gain}\OperatorTok{=}\FloatTok{0.8}\NormalTok{)}
\BuiltInTok{print}\NormalTok{(}\BuiltInTok{round}\NormalTok{(h.corrective\_response, }\DecValTok{3}\NormalTok{), h.is\_within\_tolerance)}
\end{Highlighting}
\end{Shaded}

\begin{quote}
\textbf{Clinical / systems note:} Glucose clamp studies in humans
literally estimate controller gain for insulin secretion --- the same
proportional picture in a different domain.
\end{quote}

\begin{center}\rule{0.5\linewidth}{0.5pt}\end{center}

\subsection{Current Evidence and Frontier
Biology}\label{current-evidence-and-frontier-biology-33}

For \textbf{Endocrine and Immune Systems}, frontier biology belongs
inside the evidence logic of the chapter. Physiology now blends
mechanism with allostasis, immune-endocrine-neural coupling, wearable
data, and individualized risk without reducing bodies to simple
machines. The core reading question is this: endocrine-immune claims
should include feedback, timing, receptor sensitivity, inflammation, and
allostatic load.

\begin{itemize}
\tightlist
\item
  \textbf{What to verify:} identify the observation, model, assay, or
  dataset that would make the claim stronger or weaker.
\item
  \textbf{What to qualify:} state the scale, organism, cell type,
  environmental condition, or population where the claim is expected to
  hold.
\item
  \textbf{What to compare:} test at least one alternative explanation,
  baseline, or null model before treating the pattern as causal.
\item
  \textbf{What to cite:} distinguish primary evidence, review synthesis,
  public dataset, and institutional guidance; for recent or numeric
  claims, prefer the source closest to the measurement and state what
  has changed since it was published.
\end{itemize}

Interpret physiological data by separating baseline variation,
perturbation response, compensation, and the threshold where
compensation becomes pathology.

\textbf{Source practice:} For physiology claims, cite the measurement
context and distinguish baseline variation, compensation,
pathophysiology, and treatment evidence.

\subsubsection{Current Evidence Map: Allostasis and Immune-Endocrine
Coupling}\label{current-evidence-map-allostasis-and-immune-endocrine-coupling}

\begin{figure}[htbp]
\centering
\includegraphics[width=0.82\linewidth,height=4.2in,keepaspectratio]{../figures/mermaid_inline/inline_mermaid_0173_32b78aac9da5.png}
\caption{Physiology is often adaptive over short time scales and costly over long time scales, so baseline, perturbation, compensation, and pathology must be distinguished.}
\end{figure}

\subsection{Summary}\label{summary-33}

\begin{itemize}
\tightlist
\item
  \textbf{Endocrine system:} Hierarchical hypothalamic-pituitary-target
  gland axes with negative feedback. Three hormone classes: peptide
  (surface receptors, second messengers), steroid (nuclear receptors,
  transcription), amino acid derivatives (variable). Carrier proteins
  (CBG, SHBG, TBG) buffer plasma levels of lipid-soluble hormones.
\item
  \textbf{HPA axis:} Stress \texorpdfstring{\ensuremath{\to}}{->} CRH \texorpdfstring{\ensuremath{\to}}{->} ACTH \texorpdfstring{\ensuremath{\to}}{->} cortisol. Cortisol
  synthesised from cholesterol via StAR-mediated transport and CYP
  enzymes. GR acts via transactivation (gluconeogenic genes) and
  transrepression (NF-\texorpdfstring{\ensuremath{\kappa}}{kappa}B). Circadian (peak 08:00) and pulsatile.
  Cushing's (excess) vs Addison's (deficiency).
\item
  \textbf{HPT axis:} TRH \texorpdfstring{\ensuremath{\to}}{->} TSH \texorpdfstring{\ensuremath{\to}}{->} T4/T3. T4 synthesised on iodinated
  thyroglobulin scaffold by TPO. T4 is prohormone; D1/D2 deiodinases
  convert to active T3. T3 binds nuclear TR/RXR heterodimers at TREs
  (active repression in absence of T3). Increases BMR via
  Na\(^+\)/K\(^+\)-ATPase, mitochondrial uncoupling. Hashimoto's (hypo)
  vs Graves' (hyper). Wolff-Chaikoff effect on iodine load.
\item
  \textbf{HPG axis:} Pulsatile GnRH \texorpdfstring{\ensuremath{\to}}{->} LH/FSH \texorpdfstring{\ensuremath{\to}}{->} sex steroids. Continuous
  GnRH suppresses (basis of agonist therapy). Menstrual cycle:
  follicular (oestrogen), ovulation (positive feedback LH surge), luteal
  (progesterone).
\item
  \textbf{Glucose homeostasis:} Insulin (\texorpdfstring{\ensuremath{\beta}}{beta}-cells) \texorpdfstring{\ensuremath{\to}}{->} IR (RTK) \texorpdfstring{\ensuremath{\to}}{->} IRS \texorpdfstring{\ensuremath{\to}}{->}
  PI3K \texorpdfstring{\ensuremath{\to}}{->} Akt \texorpdfstring{\ensuremath{\to}}{->} GLUT4 translocation, glycogen synthesis, FoxO
  inactivation. Glucagon (\texorpdfstring{\ensuremath{\alpha}}{alpha}-cells) \texorpdfstring{\ensuremath{\to}}{->} cAMP/PKA \texorpdfstring{\ensuremath{\to}}{->}
  glycogenolysis/gluconeogenesis. Crosstalk with leptin (long-term
  energy stores), adiponectin (AMPK insulin sensitiser), and GLP-1
  (postprandial incretin). T1DM: autoimmune. T2DM: insulin resistance
  then \texorpdfstring{\ensuremath{\beta}}{beta}-cell failure. GLP-1 agonists (semaglutide) and SGLT2 inhibitors
  transforming treatment.
\item
  \textbf{Adrenal medulla:} Chromaffin cells release epinephrine/NE on
  sympathetic stimulation; brief, fight-or-flight effects.
\item
  \textbf{GH axis:} Pulsatile GH \texorpdfstring{\ensuremath{\to}}{->} JAK2/STAT5 \texorpdfstring{\ensuremath{\to}}{->} hepatic IGF-1. Linear
  bone growth, protein synthesis, lipolysis. Acromegaly (excess),
  gigantism, GH deficiency.
\item
  \textbf{Eicosanoids:} PLA\(_2\) liberates arachidonic acid from
  membrane phospholipids. COX \texorpdfstring{\ensuremath{\to}}{->} prostaglandins/thromboxanes. 5-LOX \texorpdfstring{\ensuremath{\to}}{->}
  leukotrienes. Glucocorticoids inhibit at PLA\(_2\); aspirin
  irreversibly acetylates COX-1. COX-2 selective inhibitors (celecoxib)
  reduce GI toxicity but raise thrombotic risk.
\item
  \textbf{Endocrine disruption:} BPA, phthalates, PFAS interact with
  hormone receptors and enzymes. Most concerning during fetal/neonatal
  development. Non-monotonic dose-responses violate classical toxicology
  assumptions.
\item
  \textbf{Innate immunity:} Immediate, non-specific. Physical barriers,
  complement (opsonisation, MAC, chemotaxis with positive amplification
  feedback), phagocytes (neutrophils, macrophages, DCs), NK cells
  (missing self), inflammation, interferons, fever. PRRs: TLRs
  (membrane; TLR4\texorpdfstring{\ensuremath{\to}}{->}MyD88\texorpdfstring{\ensuremath{\to}}{->}NF-\texorpdfstring{\ensuremath{\kappa}}{kappa}B or TLR3\texorpdfstring{\ensuremath{\to}}{->}TRIF\texorpdfstring{\ensuremath{\to}}{->}IRF3\texorpdfstring{\ensuremath{\to}}{->}IFN-\texorpdfstring{\ensuremath{\beta}}{beta}); NLRs (cytosolic;
  NLRP3 inflammasome \texorpdfstring{\ensuremath{\to}}{->} caspase-1 \texorpdfstring{\ensuremath{\to}}{->} IL-1\texorpdfstring{\ensuremath{\beta}}{beta} + pyroptosis); RLRs (viral
  RNA); cGAS-STING (cytosolic DNA \texorpdfstring{\ensuremath{\to}}{->} IFN-\texorpdfstring{\ensuremath{\beta}}{beta}).
\item
  \textbf{Complement:} Three pathways converge on C3 convertase \texorpdfstring{\ensuremath{\to}}{->} C5
  convertase \texorpdfstring{\ensuremath{\to}}{->} MAC. Effector functions: opsonisation (C3b), MAC (C5b-9),
  anaphylatoxins (C3a, C5a). Regulators (DAF, CD59, factor H, C1-INH)
  confine to pathogen surfaces. PNH, hereditary angioedema, eculizumab.
\item
  \textbf{Adaptive immunity:} Specific, slow (7--10 days), memory. MHC I
  (endogenous) \texorpdfstring{\ensuremath{\to}}{->} CD8+ CTL. MHC II (exogenous) \texorpdfstring{\ensuremath{\to}}{->} CD4+ Th. T cell
  development DN1\texorpdfstring{\ensuremath{\to}}{->}DN4\texorpdfstring{\ensuremath{\to}}{->}DP\texorpdfstring{\ensuremath{\to}}{->}positive selection (cTECs)\texorpdfstring{\ensuremath{\to}}{->}negative selection
  (mTECs/AIRE/Tregs). Th polarisation: Th1, Th2, Th17, Treg, Tfh.
\item
  \textbf{B cells:} T-dependent vs T-independent activation. Germinal
  centre reaction: SHM (AID) \texorpdfstring{\ensuremath{\to}}{->} affinity maturation; CSR \texorpdfstring{\ensuremath{\to}}{->} IgG/IgA/IgE.
  Antibody isotypes: IgM (primary), IgG (secondary, opsonisation,
  placenta), IgA (mucosal), IgE (allergy).
\item
  \textbf{Cytokine network:} IL-1/IL-6/TNF (inflammation), IL-2 (T cell
  expansion), IL-4 (Th2/IgE), IL-17 (Th17/barrier), IFN-\texorpdfstring{\ensuremath{\gamma}}{gamma}
  (Th1/macrophage activation), IL-10/TGF-\texorpdfstring{\ensuremath{\beta}}{beta} (anti-inflammatory).
\item
  \textbf{Memory:} T\(_{CM}\) (recirculating), T\(_{EM}\) (peripheral),
  T\(_{RM}\) (tissue-resident); maintained by IL-7/IL-15. Pre-existing
  antigen-specific cells, lower activation threshold, pre-armed effector
  machinery explain the rapid recall response.
\item
  \textbf{Tolerance and autoimmunity:} Central (thymic/BM deletion,
  AIRE) and peripheral (Treg, anergy, PD-1, CTLA-4). Mechanisms of
  failure: molecular mimicry, bystander activation, epitope spreading,
  defective Tregs/apoptosis. Diseases: SLE, RA, T1DM, MS.
  Hypersensitivity I--IV.
\item
  \textbf{Immunotherapy:} Checkpoint inhibitors (anti-PD-1, anti-CTLA-4,
  anti-LAG-3) --- 2018 Nobel Prize; response rates vary by tumour,
  biomarker, line of therapy, and combination strategy, with durable
  ``tail of the curve'' remissions in some responders. CAR-T for
  haematological malignancies (anti-CD19 for B-cell acute lymphoblastic
  leukaemia and DLBCL; anti-BCMA for myeloma); CRS managed by
  tocilizumab.
\item
  \textbf{Connections:} See \cref{sec:unit_III_metabolic_integration}
  for insulin/glucagon integration,
  \cref{sec:unit_VII_infectious_disease} for PAMP recognition, and
  \cref{sec:unit_IX_circulation_respiration_homeostasis} for stress
  axes.
\end{itemize}

\begin{center}\rule{0.5\linewidth}{0.5pt}\end{center}

\subsection{Key Terms}\label{key-terms-33}

{\def\LTcaptype{none} % do not increment counter
\begin{longtable}[]{@{}
  >{\raggedright\arraybackslash}p{(\linewidth - 2\tabcolsep) * \real{0.2857}}
  >{\raggedright\arraybackslash}p{(\linewidth - 2\tabcolsep) * \real{0.7143}}@{}}
\toprule\noalign{}
\begin{minipage}[b]{\linewidth}\raggedright
Term
\end{minipage} & \begin{minipage}[b]{\linewidth}\raggedright
Definition
\end{minipage} \\
\midrule\noalign{}
\endhead
\bottomrule\noalign{}
\endlastfoot
\textbf{Hormone} & Chemical signal secreted into blood; acts at distant
target via specific receptor \\
\textbf{HPA axis} & Hypothalamus--pituitary--adrenal; stress response;
cortisol \\
\textbf{HPT axis} & Hypothalamus--pituitary--thyroid; metabolism;
T3/T4 \\
\textbf{HPG axis} & Hypothalamus--pituitary--gonad; reproduction; sex
steroids \\
\textbf{StAR} & Steroidogenic acute regulatory protein; rate-limiting
cholesterol transport \\
\textbf{GR} & Glucocorticoid receptor; nuclear receptor with
transactivation and transrepression modes \\
\textbf{Deiodinase} & T4\texorpdfstring{\ensuremath{\to}}{->}T3 conversion enzymes (D1, D2 activate; D3
inactivates) \\
\textbf{Wolff-Chaikoff effect} & Acute iodide excess transiently
inhibits thyroid hormone synthesis \\
\textbf{Insulin receptor} & \(\alpha_2\beta_2\) RTK; phosphorylates IRS
\texorpdfstring{\ensuremath{\to}}{->} PI3K \texorpdfstring{\ensuremath{\to}}{->} Akt \texorpdfstring{\ensuremath{\to}}{->} GLUT4 \\
\textbf{Leptin} & Adipocyte hormone; satiety; long-term energy signal;
LepR/JAK2/STAT3 \\
\textbf{Adiponectin} & Adipocyte hormone; insulin sensitiser via AMPK \\
\textbf{GLP-1} & Incretin from intestinal L-cells; potentiates
glucose-stimulated insulin \\
\textbf{PRR} & Pattern recognition receptor (TLR, NLR, RLR, cGAS) \\
\textbf{TLR4} & Recognises LPS; signals via MyD88 \texorpdfstring{\ensuremath{\to}}{->} NF-\texorpdfstring{\ensuremath{\kappa}}{kappa}B \\
\textbf{NLRP3} & Cytosolic inflammasome; activates caspase-1 \texorpdfstring{\ensuremath{\to}}{->} IL-1\texorpdfstring{\ensuremath{\beta}}{beta} +
pyroptosis \\
\textbf{cGAS-STING} & Cytosolic DNA sensor \texorpdfstring{\ensuremath{\to}}{->} type I interferons \\
\textbf{Complement} & Cascade producing opsonisation (C3b), MAC (C5b-9),
anaphylatoxins (C3a, C5a) \\
\textbf{C3 convertase / C5 convertase} & Central enzyme amplifying
complement (C4b2a or C3bBb / C4b2aC3b) \\
\textbf{MAC} & Membrane attack complex (C5b-9); 10 nm pore lyses target
cell \\
\textbf{MHC I/II} & Antigen-presenting molecules to CD8+/CD4+ T cells \\
\textbf{AIRE} & Thymic transcription factor enabling expression of
tissue-specific antigens for negative selection \\
\textbf{AID} & Activation-induced cytidine deaminase; somatic
hypermutation and class switching \\
\textbf{Germinal centre} & Lymphoid microenvironment for affinity
maturation and class switching \\
\textbf{Treg} & FoxP3+ regulatory T cell; peripheral tolerance \\
\textbf{PD-1 / PD-L1} & Inhibitory checkpoint exploited by tumours;
target of pembrolizumab \\
\textbf{CTLA-4} & Inhibitory checkpoint on activated T cells; target of
ipilimumab \\
\textbf{CAR-T cell} & Engineered T cell with synthetic antigen receptor
for tumour killing \\
\textbf{Hypersensitivity I--IV} & Allergy/anaphylaxis (IgE), cytotoxic,
immune complex, delayed cell-mediated \\
\textbf{Molecular mimicry} & Pathogen antigen resembles self; triggers
cross-reactive autoimmunity \\
\textbf{Epitope spreading} & Autoimmune response broadens from one
self-antigen to many over time \\
\end{longtable}
}

\begin{center}\rule{0.5\linewidth}{0.5pt}\end{center}

\subsection{Review Questions}\label{review-questions-33}

\begin{enumerate}
\def\labelenumi{\arabic{enumi}.}
\item
  Compare peptide and steroid hormone signalling in terms of synthesis,
  transport, receptor location, signalling speed, and duration. Why do
  steroid hormones have longer-lasting effects despite similar
  half-lives?
\item
  Explain cortisol synthesis from cholesterol, including the role of
  StAR. Which enzyme deficiency causes the most common form of
  congenital adrenal hyperplasia, and why does it produce both
  glucocorticoid deficiency and androgen excess?
\item
  Compare the two transcriptional modes of the glucocorticoid receptor
  (transactivation vs transrepression). Which mediates the
  anti-inflammatory effect of glucocorticoid drugs, and why does this
  dual mechanism complicate the development of ``dissociated''
  glucocorticoids with anti-inflammatory action but no metabolic side
  effects?
\item
  Explain why T4 is called a prohormone. Describe the role of D1, D2,
  and D3 deiodinases. Why is rT3 elevated during severe illness (``low
  T3 syndrome'')?
\item
  A patient has elevated TSH but normal free T4 (``subclinical
  hypothyroidism''). Explain the mechanism (think about the slope of the
  TSH--T4 feedback curve) and the clinical decision about treatment.
\item
  Trace the molecular pathway from insulin binding to GLUT4
  translocation in muscle. Identify three nodes where Type 2
  diabetes-associated defects occur.
\item
  Explain how leptin signals body energy stores to the brain and why
  most obese individuals have leptin resistance. How does adiponectin
  act as a complementary signal? Why doesn't exogenous leptin work for
  obesity?
\item
  Describe TLR4 signalling in response to bacterial LPS. Compare with
  the TRIF-IRF3 branch and the cGAS-STING pathway. Why does septic shock
  result from massive cytokine release, and what is the mechanism of
  action of ``anti-TNF'' therapeutics like infliximab?
\item
  Explain how the alternative complement pathway amplifies its own
  activation. Describe the formation of C3 convertase, C5 convertase,
  and the MAC. Why is complement activity restricted to pathogen
  surfaces and not host cells? What goes wrong in PNH, and how does
  eculizumab work?
\item
  Describe the DN1--DN4 stages of T cell development. What is
  \texorpdfstring{\ensuremath{\beta}}{beta}-selection? Then describe positive and negative T cell selection in
  the thymus. What is the role of AIRE, and what disease results from
  its mutation?
\item
  Compare T-dependent and T-independent B cell responses. Why are
  conjugate vaccines (PCV13, Hib) more effective in young children than
  pure polysaccharide vaccines? Trace the steps from naïve B cell
  antigen encounter to long-lived plasma cell.
\item
  Explain the molecular basis of affinity maturation in the germinal
  centre. How does AID-mediated somatic hypermutation combined with
  antigen-driven selection produce 1000-fold increases in antibody
  affinity?
\item
  Compare the four types of hypersensitivity reactions. Identify a
  clinical example and treatment for each. Explain the difference in
  time course between Type I and Type IV.
\item
  Explain the difference between molecular mimicry, bystander
  activation, and epitope spreading as mechanisms of autoimmunity. Give
  a specific clinical example for each.
\item
  A patient on pembrolizumab develops new-onset hypothyroidism,
  vitiligo, and colitis. Explain mechanistically why anti-PD-1 therapy
  causes these immune-related adverse events. What does this reveal
  about the normal role of PD-1?
\item
  Describe the design of an anti-CD19 CAR-T cell --- antigen recognition
  domain, hinge, transmembrane, co-stimulatory, and activation domains.
  Why does CAR-T therapy frequently cause cytokine release syndrome, and
  how does anti-IL-6R tocilizumab treat CRS without compromising tumour
  killing?
\item
  Bisphenol A binds oestrogen receptors with much lower affinity than
  oestradiol. Explain the concept of non-monotonic dose-response curves
  and why low-dose endocrine disruption can have greater developmental
  impact than high-dose adult exposure.
\item
  Aspirin and ibuprofen both inhibit cyclooxygenase, but primarily
  aspirin permanently inhibits platelet COX-1. Explain the molecular and
  cellular basis of this difference and its therapeutic significance.
  Why did selective COX-2 inhibitors like rofecoxib increase
  cardiovascular risk?
\item
  Explain why memory T and B cells respond more rapidly and strongly to
  re-exposure than naïve cells. Identify at least four mechanistic
  differences.
\item
  A patient develops anaphylaxis after a wasp sting. Trace the cellular
  and molecular events from prior sensitisation through the acute
  reaction. Explain why intramuscular epinephrine is the first-line
  treatment and how each of its receptor effects (\texorpdfstring{\ensuremath{\alpha}}{alpha}\(_1\), \texorpdfstring{\ensuremath{\beta}}{beta}\(_1\),
  \texorpdfstring{\ensuremath{\beta}}{beta}\(_2\)) addresses different components of anaphylaxis.
\end{enumerate}

\begin{center}\rule{0.5\linewidth}{0.5pt}\end{center}

\subsection{Further Reading and Source
Notes}\label{further-reading-and-source-notes-28}

\begin{itemize}
\tightlist
\item
  Sterling \& Eyer (1988). Allostasis: a new paradigm to explain arousal
  pathology. In \emph{Handbook of Life Stress, Cognition and Health}
  \citep{sterling1988}.
\item
  Sterling (2015). Principles of allostasis: optimal design, predictive
  regulation, pathophysiology, and rational therapeutics
  \citep{sterling2015}.
\item
  Cannon (1932). \emph{The Wisdom of the Body}. W. W. Norton.
\end{itemize}

\begin{center}\rule{0.5\linewidth}{0.5pt}\end{center}

\subsubsection{Companion Source
Module}\label{companion-source-module-33}

\textbf{Endocrine and Immune Systems} should leave a reproducible trail
from a biological claim to the code, figure, diagram, or paper-based
activity that can test it. Use the surfaces below to inspect the
chapter's assumptions, rerun the relevant model, or compare the
manuscript explanation with companion labs and figures.

{\def\LTcaptype{none} % do not increment counter
\begin{longtable}[]{@{}
  >{\raggedright\arraybackslash}p{(\linewidth - 2\tabcolsep) * \real{0.5000}}
  >{\raggedright\arraybackslash}p{(\linewidth - 2\tabcolsep) * \real{0.5000}}@{}}
\toprule\noalign{}
\begin{minipage}[b]{\linewidth}\raggedright
Surface
\end{minipage} & \begin{minipage}[b]{\linewidth}\raggedright
Use it for
\end{minipage} \\
\midrule\noalign{}
\endhead
\bottomrule\noalign{}
\endlastfoot
\breaktt{src/biology/physiology/physiology.py}
(\breaktt{homeostasis\_response}) & Compare hormone feedback and
inflammatory regulation as control problems. \\
\breaktt{src/biology/cell/cell\_biology.py}
(\breaktt{receptor\_occupancy}, \breaktt{signal\_amplification}) &
Quantify receptor sensitivity and cascade gain. \\
\breaktt{src/mermaid/biology\_diagrams.py}
(\breaktt{immune\_response\_diagram},
\breaktt{hormone\_signaling\_diagram}) & Connect endocrine and immune
sequence logic. \\
\end{longtable}
}

\textbf{Reproducibility check:} specify ligand/cytokine, receptor,
timing, tissue, feedback loop, and readout before calling a response
adaptive or pathological. \textbf{Cross-reference:} use
\cref{sec:unit_II_cell_signaling},
\cref{sec:unit_IX_circulation_respiration_homeostasis}, and
\cref{sec:unit_VII_infectious_disease}.

\newpage

\section*{Unit X --- Ecology:
Introduction}\label{unit-x-ecology-introduction}
\addcontentsline{toc}{section}{Unit X --- Ecology: Introduction}

\subsection*{Why This Unit Matters}\label{why-this-unit-matters-10}
\addcontentsline{toc}{subsection}{Why This Unit Matters}

In 1966, ecologist Robert Paine removed most sea stars (\emph{Pisaster
ochraceus}) from a stretch of rocky intertidal shore in Washington State
and waited. Within 18 months, the mussel population exploded, crowding
out barnacles, limpets, chitons, and anemones. What had been a
15-species community collapsed to a near-monoculture of \emph{Mytilus
californianus}. The sea star was a \textbf{keystone species}: a species
whose removal from the community caused a change in biodiversity
disproportionate to its own abundance. Paine's experiment created the
concept and the measurement --- the ratio of effect to biomass --- and
it changed how biologists think about the architecture of ecosystems
(Paine, 1966, \emph{American Naturalist}).

Ecology is the study of how organisms interact with each other and with
their physical environment. It operates across scales spanning six
orders of magnitude: from the foraging decision of a single aphid on a
single leaf, to the population dynamics of wolves in Yellowstone, to the
cycling of carbon through the entire Earth system. Its models --- the
logistic growth equation, the Lotka-Volterra predator-prey system, the
species-area relationship --- are among the most mathematically
tractable in biology, yet they govern phenomena as urgent as climate
change, mass extinction, and pandemic emergence.

The 21st century has made ecology a discipline under crisis and urgent
relevance. The current mass extinction event --- the sixth in Earth's
history --- is driven by habitat loss, overexploitation, invasive
species, pollution, and climate change. Biodiversity is losing the
equivalent of tens of thousands of species per year. Understanding
population dynamics, community structure, and ecosystem function is no
longer merely academic: it is an existential scientific priority for the
conservation of life on Earth. This unit will give you the quantitative
tools to understand why, and to participate in the solutions.

\begin{center}\rule{0.5\linewidth}{0.5pt}\end{center}

\subsection*{Landmark Discoveries}\label{landmark-discoveries-10}
\addcontentsline{toc}{subsection}{Landmark Discoveries}

{\def\LTcaptype{none} % do not increment counter
\begin{longtable}[]{@{}
  >{\raggedright\arraybackslash}p{(\linewidth - 8\tabcolsep) * \real{0.2407}}
  >{\raggedright\arraybackslash}p{(\linewidth - 8\tabcolsep) * \real{0.0741}}
  >{\raggedright\arraybackslash}p{(\linewidth - 8\tabcolsep) * \real{0.2963}}
  >{\raggedright\arraybackslash}p{(\linewidth - 8\tabcolsep) * \real{0.1667}}
  >{\raggedright\arraybackslash}p{(\linewidth - 8\tabcolsep) * \real{0.2222}}@{}}
\toprule\noalign{}
\begin{minipage}[b]{\linewidth}\raggedright
Discoverer(s)
\end{minipage} & \begin{minipage}[b]{\linewidth}\raggedright
Year
\end{minipage} & \begin{minipage}[b]{\linewidth}\raggedright
Journal / Source
\end{minipage} & \begin{minipage}[b]{\linewidth}\raggedright
Discovery
\end{minipage} & \begin{minipage}[b]{\linewidth}\raggedright
Significance
\end{minipage} \\
\midrule\noalign{}
\endhead
\bottomrule\noalign{}
\endlastfoot
Alfred Lotka \& Vito Volterra & 1925--1926 & \emph{Elements of Physical
Biology}; \emph{Mem. Accad. Lincei} & Lotka-Volterra predator-prey ODEs
& First quantitative model of species interactions; linked ecology to
mathematics \\
Eugene Odum & 1953 & \emph{Fundamentals of Ecology} & Ecosystem as unit
of ecology with energy flow & Introduced energy-flow and trophic-level
thinking into mainstream ecology \\
Robert Paine & 1966 & \emph{American Naturalist} & Keystone species
concept (sea star experiment) & Showed top-down regulation; single
species can structure entire communities \\
Robert MacArthur \& E.O. Wilson & 1967 & \emph{The Theory of Island
Biogeography} & Island biogeography: \(S = cA^z\) & Connected species
richness to area and isolation; foundational for conservation \\
Paul Ehrlich \& Peter Raven & 1964 & \emph{Evolution} & Coevolution
between plants and butterflies & Launched chemical ecology and
co-evolutionary theory \\
Intergovernmental Science-Policy Platform on Biodiversity (IPBES) & 2019
& IPBES Global Assessment & 1 million species threatened with extinction
& Quantified the biodiversity crisis; policy-science interface \\
Ripple et al. & 2014 & \emph{Science} & Trophic cascades from large
predators & Wolves reintroduced to Yellowstone reshaped rivers via
trophic cascade \\
\end{longtable}
}

\begin{center}\rule{0.5\linewidth}{0.5pt}\end{center}

\subsection*{Key Concepts and
Connections}\label{key-concepts-and-connections-10}
\addcontentsline{toc}{subsection}{Key Concepts and Connections}

\begin{figure}[htbp]
\centering
\includegraphics[width=0.82\linewidth,height=4.2in,keepaspectratio]{../figures/mermaid_inline/inline_mermaid_0174_bf7655c0a300.png}
\caption{Unit X concept map - Ecology. Green = population; blue = community; orange = ecosystem; red = conservation.}
\end{figure}

\begin{center}\rule{0.5\linewidth}{0.5pt}\end{center}

\subsection*{Current Evidence Thread}\label{current-evidence-thread-10}
\addcontentsline{toc}{subsection}{Current Evidence Thread}

Read this unit as a chain of ecological evidence rather than a list of
topics: population, community, ecosystem, and conservation claims are
each backed by a characteristic kind of measurement --- field sampling
and mark-recapture, manipulative and observational experiments, remote
sensing and flux networks, dynamical and statistical models, and
long-term monitoring programs. Ecology and conservation decisions
increasingly combine field data, remote sensing, community knowledge,
model uncertainty, and explicit values. As you move through the
chapters, keep a two-column note: \textbf{claim} on the left,
\textbf{evidence that would change my confidence} on the right. By the
end of the unit, each major idea should be tied to a measurement, model,
citation, or paper-based lab decision.

\subsection*{Chapter Roadmap}\label{chapter-roadmap-11}
\addcontentsline{toc}{subsection}{Chapter Roadmap}

{\def\LTcaptype{none} % do not increment counter
\begin{longtable}[]{@{}
  >{\raggedright\arraybackslash}p{(\linewidth - 6\tabcolsep) * \real{0.1556}}
  >{\raggedright\arraybackslash}p{(\linewidth - 6\tabcolsep) * \real{0.1111}}
  >{\raggedright\arraybackslash}p{(\linewidth - 6\tabcolsep) * \real{0.2889}}
  >{\raggedright\arraybackslash}p{(\linewidth - 6\tabcolsep) * \real{0.4444}}@{}}
\toprule\noalign{}
\begin{minipage}[b]{\linewidth}\raggedright
Chapter
\end{minipage} & \begin{minipage}[b]{\linewidth}\raggedright
Title
\end{minipage} & \begin{minipage}[b]{\linewidth}\raggedright
Core Question
\end{minipage} & \begin{minipage}[b]{\linewidth}\raggedright
Key Equation / Model
\end{minipage} \\
\midrule\noalign{}
\endhead
\bottomrule\noalign{}
\endlastfoot
\textbf{32} & Population Ecology & How do populations grow, and what
factors regulate population size? & Logistic: \(dN/dt = rN(1-N/K)\);
life tables; \(\lambda = e^r\) \\
\textbf{33} & Community Ecology & How do species interactions structure
biological communities? & Lotka-Volterra:
\(dN/dt = \alpha N - \beta NP\); Shannon \(H'\) \\
\textbf{34} & Ecosystem Ecology & How does energy flow and how do
nutrients cycle through ecosystems? & Trophic efficiency
about 10\%; \(NEP = GPP - R_{total}\); biogeochemical cycle
models \\
\textbf{35} & Biomes and Conservation Biology & What determines global
patterns of biodiversity, and how do we protect them? & Species-area:
\(S = cA^z\); minimum viable population (MVP); extinction rate models \\
\end{longtable}
}

\begin{center}\rule{0.5\linewidth}{0.5pt}\end{center}

\subsection*{Connections Across the
Textbook}\label{connections-across-the-textbook-10}
\addcontentsline{toc}{subsection}{Connections Across the Textbook}

\begin{itemize}
\tightlist
\item
  \textbf{Population models} (logistic growth, age structure) use the
  same differential equation framework as Unit VII (SIR epidemiological
  models) --- ecology and epidemiology are mathematically unified.
\item
  \textbf{Species diversity indices} (Shannon \(H'\), Simpson \(D\)) are
  computed using tools from Unit VII (microbial communities) and help
  evaluate ecosystem health.
\item
  \textbf{Nutrient cycling} (nitrogen, carbon, phosphorus) directly ties
  to Unit III (carbon fixation by photosynthesis) and Unit VII
  (microbial decomposers, nitrogen fixation by bacteria).
\item
  \textbf{Climate change impacts} on biomes connect to Unit VIII (plant
  responses to drought, CO\texorpdfstring{\textsubscript{2}}{2} enrichment) and Unit VI (evolution of
  species under selection by climate).
\item
  \textbf{Conservation biology} uses evolutionary genetics (Unit VI:
  effective population size, genetic rescue) and macroecology (island
  biogeography, keystone species).
\end{itemize}

\begin{quote}
\textbf{Key vocabulary introduced here:} intrinsic rate of increase
(\(r\)), carrying capacity (\(K\)), logistic growth, density dependence,
keystone species, trophic cascade, competitive exclusion, fundamental
vs.~realised niche, coevolution, mutualism, parasitism, primary
productivity (GPP, NPP), trophic efficiency, biogeochemical cycle,
biome, species-area relationship, island biogeography, minimum viable
population (MVP), extinction debt.
\end{quote}

\subsection*{Computational Toolbox --- Unit
X}\label{computational-toolbox-unit-x}
\addcontentsline{toc}{subsection}{Computational Toolbox --- Unit X}

\begin{Shaded}
\begin{Highlighting}[]
\ImportTok{from}\NormalTok{ biology.ecology }\ImportTok{import}\NormalTok{ logistic\_growth, lotka\_volterra}

\CommentTok{\# Logistic growth: white{-}tailed deer reintroduction (r = 0.35/yr, K = 500, N0 = 10)}
\NormalTok{lg }\OperatorTok{=}\NormalTok{ logistic\_growth(N0}\OperatorTok{=}\FloatTok{10.0}\NormalTok{, r}\OperatorTok{=}\FloatTok{0.35}\NormalTok{, K}\OperatorTok{=}\FloatTok{500.0}\NormalTok{, t\_end}\OperatorTok{=}\FloatTok{30.0}\NormalTok{, steps}\OperatorTok{=}\DecValTok{300}\NormalTok{)}
\BuiltInTok{print}\NormalTok{(}\SpecialStringTok{f"N(0) ≈ }\SpecialCharTok{\{}\NormalTok{lg}\SpecialCharTok{.}\NormalTok{populations[}\DecValTok{0}\NormalTok{]}\SpecialCharTok{:.0f\}}\SpecialStringTok{; N(t\_end) ≈ }\SpecialCharTok{\{}\NormalTok{lg}\SpecialCharTok{.}\NormalTok{populations[}\OperatorTok{{-}}\DecValTok{1}\NormalTok{]}\SpecialCharTok{:.0f\}}\SpecialStringTok{"}\NormalTok{)}

\CommentTok{\# Lotka{-}Volterra predator{-}prey (same ODEs used in community ecology)}
\NormalTok{lv }\OperatorTok{=}\NormalTok{ lotka\_volterra(}
    \FloatTok{100.0}\NormalTok{, }\FloatTok{10.0}\NormalTok{, }\FloatTok{0.5}\NormalTok{, }\FloatTok{0.02}\NormalTok{, }\FloatTok{0.01}\NormalTok{, }\FloatTok{0.2}\NormalTok{, t\_end}\OperatorTok{=}\FloatTok{80.0}\NormalTok{, steps}\OperatorTok{=}\DecValTok{800}\NormalTok{,}
\NormalTok{)}
\BuiltInTok{print}\NormalTok{(}\SpecialStringTok{f"Prey range: }\SpecialCharTok{\{}\BuiltInTok{min}\NormalTok{(lv.prey)}\SpecialCharTok{:.1f\}}\SpecialStringTok{–}\SpecialCharTok{\{}\BuiltInTok{max}\NormalTok{(lv.prey)}\SpecialCharTok{:.1f\}}\SpecialStringTok{"}\NormalTok{)}
\BuiltInTok{print}\NormalTok{(}\SpecialStringTok{f"Predator range: }\SpecialCharTok{\{}\BuiltInTok{min}\NormalTok{(lv.predator)}\SpecialCharTok{:.1f\}}\SpecialStringTok{–}\SpecialCharTok{\{}\BuiltInTok{max}\NormalTok{(lv.predator)}\SpecialCharTok{:.1f\}}\SpecialStringTok{"}\NormalTok{)}
\end{Highlighting}
\end{Shaded}

\begin{quote}
\textbf{Try it yourself:} Increase \texttt{beta} (predation) and watch
oscillation amplitude widen in \breaktt{plot\_lotka\_volterra} output
from \breaktt{scripts/generate\_figures.py}.
\end{quote}

\begin{center}\rule{0.5\linewidth}{0.5pt}\end{center}

\emph{Source note: ecology helpers support logistic and Allee growth,
predator-prey cycles, species-area curves, biodiversity indices, and
biome data.} \emph{Figures: \breaktt{src/visualization/} (Lotka-Volterra
cycles, logistic growth curves);
\breaktt{src/mermaid/biology\_diagrams.py} (food web diagrams, nutrient
cycle diagrams).}

\newpage

\section{Population Ecology and Growth
Models}\label{population-ecology-and-growth-models}

\label{sec:unit_X_population_ecology}

\begin{quote}
\textbf{Ch 32} \texorpdfstring{\ensuremath{\cdot}}{.} Level 3/3 \texorpdfstring{\ensuremath{\cdot}}{.} 75 min read \texorpdfstring{\ensuremath{\cdot}}{.} 100 min lecture \texorpdfstring{\ensuremath{\cdot}}{.}
Prerequisites: \cref{sec:unit_V_population_genetics}
\end{quote}

\subsection{Learning Objectives}\label{learning-objectives-34}

\begin{enumerate}
\def\labelenumi{\arabic{enumi}.}
\tightlist
\item
  Describe the key demographic parameters (birth rate, death rate,
  survivorship, net reproductive rate) of a population.
\item
  Derive and interpret the exponential and
  \hyperref[gl:logistic-growth]{\textbf{logistic growth}} equations and
  apply them to real populations.
\item
  Construct and interpret \hyperref[gl:life-table]{\textbf{life table}}s
  and \hyperref[gl:survivorship-curve]{\textbf{survivorship curve}}s
  (Types I, II, III).
\item
  Explain the Allee effect and describe how it alters population
  dynamics in small populations.
\item
  Apply the Lotka \citep{lotka1925}-Volterra \citep{volterra1926}
  competition and predator-prey equations and interpret phase-plane
  outcomes.
\item
  Describe r vs.~K selection strategies and their adaptive ecological
  contexts.
\item
  Explain metapopulation dynamics and source-sink models for fragmented
  habitats.
\item
  Apply mark-recapture methods and distance sampling to estimate
  population size.
\item
  Construct a Leslie matrix from a life table, compute \(\lambda_1\) and
  the stable age distribution, and use sensitivity/elasticity analysis
  to identify management-critical demographic rates.
\item
  Distinguish individual-based models from deterministic models, and
  outline a population viability analysis (PVA) workflow for a small
  population.
\end{enumerate}

\begin{figure}[htbp]
\centering
\includegraphics[width=0.85\textwidth]{../figures/logistic_growth.png}
\caption{Logistic growth of a population with carrying capacity K. Population size follows dN/dt = rN(1 - N/K); the S-shaped curve asymptotes at K after an inflection at N = K/2.}
\label{fig:unit_X_logistic_growth}
\end{figure}

\subsubsection{Study Blueprint}\label{study-blueprint-34}

\begin{itemize}
\tightlist
\item
  \textbf{Big idea:} Population change reflects births, deaths,
  movement, age structure, and density dependence.
\item
  \textbf{Core concepts:} exponential growth, logistic growth, life
  tables, population viability.
\item
  \textbf{Framework alignment:} Vision \& Change: Systems, Evolution,
  Pathways and transformations of energy and matter; AP Biology: Systems
  Interactions, Evolution, Energetics; NGSS-style topics: Interdependent
  Relationships in Ecosystems, Matter and Energy in Organisms and
  Ecosystems, Natural Selection and Evolution.
\item
  \textbf{Model or quantitative lens:} Exponential/logistic growth,
  mark-recapture, and matrix projection.
\item
  \textbf{Data skill:} Use abundance or age-structure data to estimate
  growth and risk.
\item
  \textbf{Practice cadence:} Representing and Describing Data,
  Statistical Tests and Data Analysis, Argumentation.
\item
  \textbf{Common misconception to repair:} Carrying capacity is not a
  fixed magic number; it changes with resources, interactions, and
  disturbance.
\item
  \textbf{Primary lab:} \cref{sec:lab_unit_X_population_ecology}.
\item
  \textbf{Question bank:} \cref{sec:q_unit_X_population_ecology}.
\item
  \textbf{Transfer task:} Transfer population models to fisheries,
  invasive species, epidemiology, and endangered species.
\item
  \textbf{Bridge to computation:}
  \breaktt{biology.ecology.ecology.logistic\_growth}.
\end{itemize}

\begin{center}\rule{0.5\linewidth}{0.5pt}\end{center}

\begin{quote}
\textbf{Opening Vignette --- The Equation That Explains Why Populations
Don't Grow Forever}

In 1838, Belgian mathematician Pierre François Verhulst looked at census
data and asked: why don't populations simply grow exponentially forever?
His answer was the logistic equation --- dN/dt = rN(1 \texorpdfstring{\ensuremath{-}}{-} N/K) --- where
K, the carrying capacity, captures the idea that resources limit growth.
For nearly a century, the logistic model was theoretical. Then in 1934,
Russian ecologist Georgy Gause grew two \emph{Paramecium} species in
test tubes and observed the sigmoid growth curve Verhulst had predicted:
slow growth, rapid growth, leveling off at carrying capacity. His
experiments also showed
\hyperref[gl:competitive-exclusion]{\textbf{competitive exclusion}} ---
two species competing for the same \hyperref[gl:niche]{\textbf{niche}}
cannot coexist indefinitely. Gause published this in \emph{The Struggle
for Existence} at age 23. The logistic model now underlies fisheries
management, epidemiology (the SIR model is a direct descendant),
conservation minimum viable population estimates, and pandemic
projections. Verhulst's equation, written in a 4-page paper with no
data, may have saved more lives than any equation since Newton's second
law.
\end{quote}

\subsubsection{Chapter Roadmap}\label{chapter-roadmap-12}

The chapter spans the full pipeline from individuals to metapopulations.
Read it as three nested scales:

\begin{itemize}
\tightlist
\item
  \textbf{The individual scale.} Population attributes, life tables,
  survivorship. What you need before any growth equation makes sense.
\item
  \textbf{The population scale.} Exponential growth, logistic growth,
  Allee effects, Lotka-Volterra interactions --- the canonical equations
  and their named modifications.
\item
  \textbf{The spatial, structured, and human scale.} r/K life histories,
  age-structured matrix models, individual-based simulation, population
  viability analysis, metapopulation dynamics, field-estimation
  techniques, and human demographic transitions.
\end{itemize}

For a brief treatment, focus on the core growth mathematics and the
applied field-estimation and demographic-transition material. For a full
quantitative course, work through most sections including the
matrix-model and PVA toolkit, which underlie modern conservation
practice.

\subsection{What Is a Population?}\label{what-is-a-population}

A \textbf{population} consists of individuals of a single species
(conspecifics) inhabiting a defined area at a given time. Population
ecology investigates the factors that regulate population size, density,
distribution, and growth over time. Understanding population dynamics is
fundamental to wildlife management, conservation biology, epidemiology,
and sustainable resource harvest.

\subsubsection{Population Attributes}\label{population-attributes}

Populations are characterised by several measurable attributes:

{\def\LTcaptype{none} % do not increment counter
\begin{longtable}[]{@{}
  >{\raggedright\arraybackslash}p{(\linewidth - 6\tabcolsep) * \real{0.3000}}
  >{\raggedright\arraybackslash}p{(\linewidth - 6\tabcolsep) * \real{0.2000}}
  >{\raggedright\arraybackslash}p{(\linewidth - 6\tabcolsep) * \real{0.3333}}
  >{\raggedright\arraybackslash}p{(\linewidth - 6\tabcolsep) * \real{0.1667}}@{}}
\toprule\noalign{}
\begin{minipage}[b]{\linewidth}\raggedright
Attribute
\end{minipage} & \begin{minipage}[b]{\linewidth}\raggedright
Symbol
\end{minipage} & \begin{minipage}[b]{\linewidth}\raggedright
Definition
\end{minipage} & \begin{minipage}[b]{\linewidth}\raggedright
Units
\end{minipage} \\
\midrule\noalign{}
\endhead
\bottomrule\noalign{}
\endlastfoot
Census size & \(N\) & Total number of individuals & Individuals \\
Population density & \(D\) & Individuals per unit area or volume &
ind/km\(^2\) or ind/L \\
Per capita birth rate & \(b\) & Births per individual per time unit &
births/ind/time \\
Per capita death rate & \(d\) & Deaths per individual per time unit &
deaths/ind/time \\
Intrinsic rate of increase & \(r\) & \(r = b - d\) (can be negative) &
time\(^{-1}\) \\
Finite rate of increase & \texorpdfstring{\ensuremath{\lambda}}{lambda} & \(\lambda = e^r\); \(\lambda > 1\)
growing; \(\lambda < 1\) declining & dimensionless \\
Net reproductive rate & \(R_0\) & \(R_0 = \sum l_x m_x\); female
offspring per female per lifetime & offspring/female \\
Generation time & \(T\) & \(T = \sum x \cdot l_x \cdot m_x / R_0\) &
time units \\
\end{longtable}
}

The fundamental relationship between \(R_0\), \(T\), and \(r\) is the
\textbf{Euler-Lotka equation}:

\begin{equation}
\sum_{x=0}^{\omega} e^{-rx} l_x m_x = 1
\label{eq:population_ecology_1}
\end{equation}

For approximate computation: \(r \approx \ln(R_0) / T\)

\subsubsection{Dispersion Patterns}\label{dispersion-patterns}

Individuals within a population are distributed in one of three spatial
patterns:

{\def\LTcaptype{none} % do not increment counter
\begin{longtable}[]{@{}
  >{\raggedright\arraybackslash}p{(\linewidth - 6\tabcolsep) * \real{0.2059}}
  >{\raggedright\arraybackslash}p{(\linewidth - 6\tabcolsep) * \real{0.3235}}
  >{\raggedright\arraybackslash}p{(\linewidth - 6\tabcolsep) * \real{0.2647}}
  >{\raggedright\arraybackslash}p{(\linewidth - 6\tabcolsep) * \real{0.2059}}@{}}
\toprule\noalign{}
\begin{minipage}[b]{\linewidth}\raggedright
Pattern
\end{minipage} & \begin{minipage}[b]{\linewidth}\raggedright
Description
\end{minipage} & \begin{minipage}[b]{\linewidth}\raggedright
Mechanism
\end{minipage} & \begin{minipage}[b]{\linewidth}\raggedright
Example
\end{minipage} \\
\midrule\noalign{}
\endhead
\bottomrule\noalign{}
\endlastfoot
\textbf{Clumped} & Aggregated in patches & Resource heterogeneity,
social behaviour, limited dispersal & Schooling fish, herding ungulates,
fungi near rotting logs \\
\textbf{Uniform} & Evenly spaced & Territoriality, allelopathy,
intraspecific competition & Nesting penguins, creosote bush in
deserts \\
\textbf{Random} & No predictable pattern & Absence of strong attraction
or repulsion & Wind-dispersed seeds in homogeneous habitat \\
\end{longtable}
}

The \textbf{variance-to-mean ratio} (\(\sigma^2 / \mu\)) of quadrat
counts distinguishes these patterns: ratio = 1 (random, Poisson),
\textgreater1 (clumped, negative binomial), \textless1 (uniform).

\begin{quote}
\textbf{Concept Check:} A researcher counts organisms in 20 quadrats and
calculates a variance-to-mean ratio of 3.7. What dispersion pattern does
this suggest, and what ecological mechanisms might produce it?
\end{quote}

\begin{center}\rule{0.5\linewidth}{0.5pt}\end{center}

\subsection{Life Tables and Survivorship
Curves}\label{life-tables-and-survivorship-curves}

A \textbf{cohort life table} (also called a horizontal or dynamic life
table) tracks a birth cohort from birth until most members die,
recording survival and fecundity at each age:

{\def\LTcaptype{none} % do not increment counter
\begin{longtable}[]{@{}
  >{\raggedright\arraybackslash}p{(\linewidth - 14\tabcolsep) * \real{0.0703}}
  >{\raggedright\arraybackslash}p{(\linewidth - 14\tabcolsep) * \real{0.1094}}
  >{\raggedright\arraybackslash}p{(\linewidth - 14\tabcolsep) * \real{0.1641}}
  >{\raggedright\arraybackslash}p{(\linewidth - 14\tabcolsep) * \real{0.1172}}
  >{\raggedright\arraybackslash}p{(\linewidth - 14\tabcolsep) * \real{0.1719}}
  >{\raggedright\arraybackslash}p{(\linewidth - 14\tabcolsep) * \real{0.1406}}
  >{\raggedright\arraybackslash}p{(\linewidth - 14\tabcolsep) * \real{0.0781}}
  >{\raggedright\arraybackslash}p{(\linewidth - 14\tabcolsep) * \real{0.1484}}@{}}
\toprule\noalign{}
\begin{minipage}[b]{\linewidth}\raggedright
Age (\(x\))
\end{minipage} & \begin{minipage}[b]{\linewidth}\raggedright
\(n_x\) (alive)
\end{minipage} & \begin{minipage}[b]{\linewidth}\raggedright
\(l_x\) (survivorship)
\end{minipage} & \begin{minipage}[b]{\linewidth}\raggedright
\(d_x\) (deaths)
\end{minipage} & \begin{minipage}[b]{\linewidth}\raggedright
\(q_x\) (mortality rate)
\end{minipage} & \begin{minipage}[b]{\linewidth}\raggedright
\(m_x\) (fecundity)
\end{minipage} & \begin{minipage}[b]{\linewidth}\raggedright
\(l_x m_x\)
\end{minipage} & \begin{minipage}[b]{\linewidth}\raggedright
\(x \cdot l_x m_x\)
\end{minipage} \\
\midrule\noalign{}
\endhead
\bottomrule\noalign{}
\endlastfoot
0 & 1000 & 1.000 & 200 & 0.200 & 0 & 0 & 0 \\
1 & 800 & 0.800 & 200 & 0.250 & 0.5 & 0.400 & 0.400 \\
2 & 600 & 0.600 & 400 & 0.667 & 1.2 & 0.720 & 1.440 \\
3 & 200 & 0.200 & 180 & 0.900 & 1.0 & 0.200 & 0.600 \\
4 & 20 & 0.020 & 20 & 1.000 & 0 & 0 & 0 \\
\end{longtable}
}

From this table:

\begin{equation}
R_0 = \sum l_x m_x = 0 + 0.400 + 0.720 + 0.200 + 0 = 1.320
\label{eq:population_ecology_2}
\end{equation}

\begin{equation}
T = \frac{\sum x \cdot l_x m_x}{R_0} = \frac{2.440}{1.320} = 1.848 \text{ generations}
\label{eq:population_ecology_3}
\end{equation}

\begin{equation}
r \approx \frac{\ln(R_0)}{T} = \frac{\ln(1.320)}{1.848} = \frac{0.278}{1.848} = 0.150 \text{ per time unit}
\label{eq:population_ecology_4}
\end{equation}

Since \(R_0 > 1\), this population is \textbf{growing}.

A \textbf{static life table} (vertical life table) uses age structure
data from a single time point --- useful when tracking cohorts is
impractical (e.g., long-lived species like elephants or trees). It
assumes stable age distribution.

\subsubsection{Survivorship Curves}\label{survivorship-curves}

Three canonical survivorship curve types (Pearl 1928; Deevey 1947):

\begin{figure}[htbp]
\centering
\includegraphics[width=0.82\linewidth,height=4.2in,keepaspectratio]{../figures/mermaid_inline/inline_mermaid_0175_d4c12273ba7b.png}
\caption{Flowchart for Survivorship Curves: Convex curve on log plot, Straight line on log plot, Concave curve on log plot, and High parental investment, few offspring, K-selected traits"] form the diagram's primary path or branches.}
\end{figure}

\textbf{Mathematical representation of survivorship:}

For Type I: \(l_x \approx e^{-bx^n}\) where \(n > 1\) (Gompertz-Makeham
law for human mortality)

For Type II: \(l_x = e^{-\mu x}\) (constant hazard rate \texorpdfstring{\ensuremath{\mu}}{mu})

For Type III: \(l_x \approx e^{-bx^n}\) where \(n < 1\)

\begin{quote}
\texorpdfstring{}{} \textbf{Clinical Connection --- Actuarial Science and Human
Survivorship:} Human survivorship curves have shifted dramatically over
the past 200 years. In 1800, human survivorship resembled a Type II
curve due to high infant and childhood mortality from infectious
disease. Modern medicine, sanitation, and nutrition have transformed the
human curve to an extreme Type I, with about 99\% survival to
age 50 in high-income countries. The \textbf{Gompertz law of mortality}
states that human death rate doubles approximately every 8 years after
age 30: \(\mu(x) = \alpha e^{\beta x}\). This principle underlies life
insurance actuarial tables and pension fund projections.
\end{quote}

\begin{quote}
\textbf{Concept Check:} Sea turtles lay 50-200 eggs per nesting event
but about 1 in 1,000 hatchlings survives to reproductive age. Which
survivorship curve type does this represent, and what would a
conservation program need to achieve (in terms of stage-specific
survival rates) to increase population growth?
\end{quote}

\subsubsection{Reproductive Value}\label{reproductive-value}

Fisher's \textbf{reproductive value} (\(v_x\)) quantifies the expected
future contribution of an individual of age \(x\) to population growth:

\begin{equation}
v_x = \frac{e^{rx}}{l_x} \sum_{y=x}^{\omega} e^{-ry} l_y m_y
\label{eq:population_ecology_5}
\end{equation}

Reproductive value peaks at the age of first reproduction in growing
populations and declines thereafter. This concept is critical for
conservation: protecting age classes with highest \(v_x\) yields the
greatest impact on population recovery. For sea turtles, subadult and
adult females have the highest reproductive value --- hence
\textbf{Turtle Excluder Devices (TEDs)} in fishing nets target this life
stage.

\begin{center}\rule{0.5\linewidth}{0.5pt}\end{center}

\subsection{Exponential Growth}\label{exponential-growth}

When resources are unlimited and the environment exerts no
\hyperref[gl:density-dependent-regulation]{\textbf{density-dependent
regulation}}, population growth is \textbf{exponential} (or geometric in
discrete-time models):

\textbf{Continuous-time model:}

\begin{equation}
\frac{dN}{dt} = rN \quad \Rightarrow \quad N(t) = N_0 e^{rt}
\label{eq:population_ecology_6}
\end{equation}

\textbf{Discrete-time model (non-overlapping generations):}

\begin{equation}
N_{t+1} = \lambda N_t \quad \Rightarrow \quad N_t = N_0 \lambda^t
\label{eq:population_ecology_7}
\end{equation}

where \(\lambda = e^r\) is the \textbf{finite rate of increase}.

\textbf{Doubling time:}
\(t_2 = \frac{\ln 2}{r} \approx \frac{0.693}{r}\)

\subsubsection{Real-World Examples of Near-Exponential
Growth}\label{real-world-examples-of-near-exponential-growth}

{\def\LTcaptype{none} % do not increment counter
\begin{longtable}[]{@{}
  >{\raggedright\arraybackslash}p{(\linewidth - 6\tabcolsep) * \real{0.3030}}
  >{\raggedright\arraybackslash}p{(\linewidth - 6\tabcolsep) * \real{0.0909}}
  >{\raggedright\arraybackslash}p{(\linewidth - 6\tabcolsep) * \real{0.3939}}
  >{\raggedright\arraybackslash}p{(\linewidth - 6\tabcolsep) * \real{0.2121}}@{}}
\toprule\noalign{}
\begin{minipage}[b]{\linewidth}\raggedright
Population
\end{minipage} & \begin{minipage}[b]{\linewidth}\raggedright
\(r\)
\end{minipage} & \begin{minipage}[b]{\linewidth}\raggedright
Doubling time
\end{minipage} & \begin{minipage}[b]{\linewidth}\raggedright
Context
\end{minipage} \\
\midrule\noalign{}
\endhead
\bottomrule\noalign{}
\endlastfoot
\emph{E. coli} (optimal) & 1.7/hr & 24.5 min & Binary fission; unlimited
glucose \\
World human population (1650--1800) & 0.005/yr & 139 yr &
Pre-industrial, low density \\
COVID-19 early spread (March 2020) & 0.23/day & 3 days & Unmitigated, no
immunity \\
Ring-necked pheasant, Protection Island (1937--42) & 1.02/yr & 0.68 yr &
8 birds introduced; no predators \\
Reindeer on St.~Matthew Island (1944--63) & 0.34/yr & 2.0 yr & 29
animals introduced; no wolves \\
\end{longtable}
}

The St.~Matthew Island reindeer population illustrates the
\textbf{catastrophic crash} that follows unsustainable exponential
growth: 29 reindeer introduced in 1944 grew to 6,000 by 1963, then
crashed to 42 by 1966 due to overgrazing of lichen habitat. This case
study is a classic illustration of \textbf{ecological overshoot}.

\begin{Shaded}
\begin{Highlighting}[]
\ImportTok{from}\NormalTok{ biology.ecology }\ImportTok{import}\NormalTok{ logistic\_growth}

\CommentTok{\# Exponential phase (K → ∞ approximation)}
\NormalTok{result }\OperatorTok{=}\NormalTok{ logistic\_growth(}
\NormalTok{    initial\_population}\OperatorTok{=}\DecValTok{100}\NormalTok{,}
\NormalTok{    growth\_rate}\OperatorTok{=}\FloatTok{0.5}\NormalTok{,}
\NormalTok{    carrying\_capacity}\OperatorTok{=}\FloatTok{1e12}\NormalTok{,   }\CommentTok{\# effectively unlimited}
\NormalTok{    time\_steps}\OperatorTok{=}\DecValTok{20}
\NormalTok{)}
\BuiltInTok{print}\NormalTok{(}\SpecialStringTok{f"After 20 time steps: N = }\SpecialCharTok{\{}\NormalTok{result}\SpecialCharTok{.}\NormalTok{trajectory[}\OperatorTok{{-}}\DecValTok{1}\NormalTok{]}\SpecialCharTok{:.0f\}}\SpecialStringTok{"}\NormalTok{)}
\BuiltInTok{print}\NormalTok{(}\SpecialStringTok{f"Predicted (analytic): N = }\SpecialCharTok{\{}\DecValTok{100} \OperatorTok{*} \FloatTok{2.718}\OperatorTok{**}\NormalTok{( }\FloatTok{0.5}\OperatorTok{*}\DecValTok{20}\NormalTok{ )}\SpecialCharTok{:.0f\}}\SpecialStringTok{"}\NormalTok{)}
\end{Highlighting}
\end{Shaded}

\begin{quote}
\textbf{Concept Check:} If a bacterial culture starts with 500 cells and
has \(r = 1.4\)/hr, how many cells will be present after 6 hours? What
assumptions does this calculation require?
\end{quote}

\begin{center}\rule{0.5\linewidth}{0.5pt}\end{center}

\subsection{Logistic Growth and Density
Dependence}\label{logistic-growth-and-density-dependence}

As \(N\) approaches the \textbf{carrying capacity} (\(K\)),
intraspecific competition for resources reduces the per capita growth
rate. The \textbf{logistic growth equation} \citep{verhulst1838}
incorporates this density dependence:

\begin{equation}
\frac{dN}{dt} = rN\left(1 - \frac{N}{K}\right)
\label{eq:population_ecology_8}
\end{equation}

The term \((1 - N/K)\) is the \textbf{density-dependence factor} (or
unused portion of carrying capacity), and it bends exponential growth
into the S-shaped trajectory of \cref{fig:unit_X_logistic_growth} with
the following key properties:

\begin{itemize}
\tightlist
\item
  When \(N \ll K\): growth is approximately exponential
  (\(dN/dt \approx rN\))
\item
  When \(N = K/2\): \(dN/dt\) is \textbf{maximised} (inflection point of
  the sigmoid curve)
\item
  When \(N = K\): \(dN/dt = 0\) (population is at equilibrium)
\item
  When \(N > K\): \(dN/dt < 0\) (population declines toward \(K\))
\end{itemize}

\textbf{Analytical solution:}

\begin{equation}
N(t) = \frac{K}{1 + \left(\frac{K - N_0}{N_0}\right)e^{-rt}}
\label{eq:population_ecology_9}
\end{equation}

\begin{figure}[htbp]
\centering
\includegraphics[width=0.82\linewidth,height=4.2in,keepaspectratio]{../figures/mermaid_inline/inline_mermaid_0176_f42f98ef3168.png}
\caption{State diagram for Logistic Growth and Density Dependence showing transitions among N0 << K, Growth accelerates, N approaches K/2, and N = K/2, max dN/dt.}
\end{figure}

\subsubsection{Density-Dependent
Factors}\label{density-dependent-factors}

{\def\LTcaptype{none} % do not increment counter
\begin{longtable}[]{@{}
  >{\raggedright\arraybackslash}p{(\linewidth - 4\tabcolsep) * \real{0.2727}}
  >{\raggedright\arraybackslash}p{(\linewidth - 4\tabcolsep) * \real{0.4091}}
  >{\raggedright\arraybackslash}p{(\linewidth - 4\tabcolsep) * \real{0.3182}}@{}}
\toprule\noalign{}
\begin{minipage}[b]{\linewidth}\raggedright
Factor
\end{minipage} & \begin{minipage}[b]{\linewidth}\raggedright
Mechanism
\end{minipage} & \begin{minipage}[b]{\linewidth}\raggedright
Example
\end{minipage} \\
\midrule\noalign{}
\endhead
\bottomrule\noalign{}
\endlastfoot
\textbf{Intraspecific competition} & Scramble or contest competition for
food, space, mates & Flour beetles (\emph{Tribolium}): cannibalism
increases at high density \\
\textbf{Disease transmission} & Contact rate increases with density &
Distemper in Serengeti lions; COVID-19 in mink farms \\
\textbf{Predation} & Predators aggregate at high-density prey patches &
Functional response Type III (sigmoidal) \\
\textbf{Territoriality} & \hyperref[gl:dominant]{\textbf{Dominant}}
individuals exclude subordinates & Red grouse: territorial males hold
heather patches \\
\textbf{Toxic waste accumulation} & Metabolic byproducts inhibit growth
& Yeast: ethanol accumulation limits
\hyperref[gl:fermentation]{\textbf{fermentation}} \\
\end{longtable}
}

\subsubsection{Density-Independent
Factors}\label{density-independent-factors}

Some factors affect populations regardless of density:

\begin{itemize}
\tightlist
\item
  \textbf{Weather extremes} (frost, drought, hurricanes)
\item
  \textbf{Volcanic eruptions}, wildfires, floods
\item
  \textbf{Pesticide application}
\end{itemize}

In practice, most populations are regulated by a combination of
density-dependent and density-independent factors.

\subsubsection{Maximum Sustainable Yield
(MSY)}\label{maximum-sustainable-yield-msy}

From the logistic equation, the maximum absolute growth rate occurs at
\(N = K/2\):

\begin{equation}
\text{MSY} = \frac{rK}{4}
\label{eq:population_ecology_10}
\end{equation}

Fisheries and wildlife management use MSY to set harvest quotas.
However, MSY has significant limitations:

\begin{enumerate}
\def\labelenumi{\arabic{enumi}.}
\tightlist
\item
  \textbf{Assumes logistic growth} --- real populations often show
  non-logistic dynamics
\item
  \textbf{K and r are difficult to estimate} --- environmental
  stochasticity changes both
\item
  \textbf{Ignores age/size structure} --- harvesting large,
  reproductively valuable individuals disproportionately impacts growth
\item
  \textbf{No safety margin} --- harvesting exactly at MSY leaves no
  \hyperref[gl:buffer]{\textbf{buffer}} for environmental variation
\end{enumerate}

\begin{quote}
\texorpdfstring{}{} \textbf{Clinical Connection --- Pacific Sardine Collapse:} The
Pacific sardine (\emph{Sardinops sagax}) fishery collapsed in the
1940s-1950s, with the population crashing from 3 million tonnes to
near-zero. The collapse resulted from harvesting above MSY during a
period when ocean conditions (Pacific Decadal Oscillation) shifted
unfavorably, reducing \(K\). The combined effect of over-exploitation
and environmental change was catastrophic. This case led to the
development of \textbf{precautionary reference points} in modern
fisheries management, where harvest targets are set below MSY (typically
at 0.5-0.8 MSY) to provide a buffer.
\end{quote}

\begin{Shaded}
\begin{Highlighting}[]
\ImportTok{from}\NormalTok{ biology.ecology }\ImportTok{import}\NormalTok{ logistic\_growth}

\NormalTok{result }\OperatorTok{=}\NormalTok{ logistic\_growth(}
\NormalTok{    initial\_population}\OperatorTok{=}\DecValTok{10}\NormalTok{,}
\NormalTok{    growth\_rate}\OperatorTok{=}\FloatTok{0.8}\NormalTok{,}
\NormalTok{    carrying\_capacity}\OperatorTok{=}\DecValTok{1000}\NormalTok{,}
\NormalTok{    time\_steps}\OperatorTok{=}\DecValTok{30}
\NormalTok{)}
\CommentTok{\# Find maximum per{-}capita growth (should be near K/2 = 500)}
\NormalTok{growth\_rates }\OperatorTok{=}\NormalTok{ [result.trajectory[i}\OperatorTok{+}\DecValTok{1}\NormalTok{] }\OperatorTok{{-}}\NormalTok{ result.trajectory[i]}
                \ControlFlowTok{for}\NormalTok{ i }\KeywordTok{in} \BuiltInTok{range}\NormalTok{(}\BuiltInTok{len}\NormalTok{(result.trajectory)}\OperatorTok{{-}}\DecValTok{1}\NormalTok{)]}
\NormalTok{peak\_N }\OperatorTok{=}\NormalTok{ result.trajectory[growth\_rates.index(}\BuiltInTok{max}\NormalTok{(growth\_rates))]}
\BuiltInTok{print}\NormalTok{(}\SpecialStringTok{f"Peak growth at N ≈ }\SpecialCharTok{\{}\NormalTok{peak\_N}\SpecialCharTok{:.0f\}}\SpecialStringTok{ (expect K/2 = 500)"}\NormalTok{)}
\end{Highlighting}
\end{Shaded}

\subsubsection{Extensions of the Logistic
Model}\label{extensions-of-the-logistic-model}

\textbf{Theta-logistic model} --- allows flexible density dependence:

\begin{equation}
\frac{dN}{dt} = rN\left[1 - \left(\frac{N}{K}\right)^\theta\right]
\label{eq:population_ecology_11}
\end{equation}

When \(\theta = 1\): standard logistic. When \(\theta > 1\): density
effects are weak until \(N\) is close to \(K\) (concave per-capita
growth). When \(\theta < 1\): density effects are strong even at low
\(N\) (convex). Empirical estimates for large mammals typically give
\(\theta \approx 2-7\) (Sibly et al.~2005, \emph{Science}), meaning
density dependence is weaker than the standard logistic assumes.

\textbf{Time-lagged logistic model:}

\begin{equation}
\frac{dN}{dt} = rN\left(1 - \frac{N(t-\tau)}{K}\right)
\label{eq:population_ecology_12}
\end{equation}

When the time lag \texorpdfstring{\ensuremath{\tau}}{tau} is large relative to \(1/r\), the population can
\textbf{overshoot} \(K\) and exhibit damped oscillations, limit cycles,
or even chaos (May 1976, \emph{Nature}).

\begin{quote}
\textbf{Concept Check:} A population of deer has \(r = 0.3\)/yr and
\(K = 5000\). Calculate MSY and the population size at which it occurs.
If the current population is 4000, how many deer can be sustainably
harvested per year?
\end{quote}

\begin{center}\rule{0.5\linewidth}{0.5pt}\end{center}

\subsection{The Allee Effect}\label{the-allee-effect}

The \textbf{Allee effect} \citep{allee1931} occurs when per capita
fitness \emph{decreases} at low population density --- the opposite of
standard logistic density-dependence. This creates a positive feedback
loop where small populations become increasingly vulnerable.

\subsubsection{Strong vs.~Weak Allee
Effects}\label{strong-vs.-weak-allee-effects}

\textbf{Strong Allee effect} --- per capita growth rate becomes negative
below a threshold:

\begin{equation}
\frac{dN}{dt} = rN\left(\frac{N}{A} - 1\right)\left(1 - \frac{N}{K}\right)
\label{eq:population_ecology_13}
\end{equation}

where \(A\) = \textbf{Allee threshold}. Below \(A\), the population
declines deterministically to extinction. There are three equilibria:
\(N = 0\) (stable), \(N = A\) (unstable), and \(N = K\) (stable).

\begin{figure}[htbp]
\centering
\includegraphics[width=0.85\textwidth]{../figures/allee_threshold_dynamics.png}
\caption{Strong Allee-effect threshold dynamics. Initial population sizes below, at, and above the Allee threshold show extinction, unstable threshold persistence, and recovery toward carrying capacity.}
\label{fig:unit_X_allee_threshold_dynamics}
\end{figure}

\textbf{Weak Allee effect} --- per capita growth rate decreases at low
density but remains positive. No threshold exists; the population can
recover from any positive size, but growth is slower than expected at
low \(N\).

\subsubsection{Mechanisms of Allee
Effects}\label{mechanisms-of-allee-effects}

{\def\LTcaptype{none} % do not increment counter
\begin{longtable}[]{@{}
  >{\raggedright\arraybackslash}p{(\linewidth - 4\tabcolsep) * \real{0.2727}}
  >{\raggedright\arraybackslash}p{(\linewidth - 4\tabcolsep) * \real{0.2121}}
  >{\raggedright\arraybackslash}p{(\linewidth - 4\tabcolsep) * \real{0.5152}}@{}}
\toprule\noalign{}
\begin{minipage}[b]{\linewidth}\raggedright
Mechanism
\end{minipage} & \begin{minipage}[b]{\linewidth}\raggedright
Example
\end{minipage} & \begin{minipage}[b]{\linewidth}\raggedright
Effect on fitness
\end{minipage} \\
\midrule\noalign{}
\endhead
\bottomrule\noalign{}
\endlastfoot
\textbf{Mate finding} & Ivory-billed woodpecker; many marine
invertebrates & Low density \texorpdfstring{\ensuremath{\to}}{->} failed mate encounter \texorpdfstring{\ensuremath{\to}}{->} reduced
reproduction \\
\textbf{Predator satiation} & Mast-seeding trees (oaks, beeches) & Below
threshold seed output \texorpdfstring{\ensuremath{\to}}{->} disproportionate predation loss \\
\textbf{Cooperative breeding/hunting} & African wild dogs (pack size \texorpdfstring{\textgreater{}=}{>=} 5
for efficient hunting) & Small packs \texorpdfstring{\ensuremath{\to}}{->} insufficient prey capture \texorpdfstring{\ensuremath{\to}}{->}
starvation \\
\textbf{Cooperative defence} & Muskoxen forming defensive circles & Too
few adults \texorpdfstring{\ensuremath{\to}}{->} predators penetrate defence \\
\textbf{Genetic diversity} & Florida panther (pre-1995,
\(N_e \approx 25\)) & Small \(N\) \texorpdfstring{\ensuremath{\to}}{->} inbreeding depression \texorpdfstring{\ensuremath{\to}}{->} reduced
fitness \\
\textbf{Pollination failure} & Rare plants in fragmented meadows & Low
density \texorpdfstring{\ensuremath{\to}}{->} pollinators don't visit \texorpdfstring{\ensuremath{\to}}{->} seed set fails \\
\textbf{Environmental conditioning} & Soil
\hyperref[gl:microbiome]{\textbf{microbiome}} enrichment by plant roots
& Few plants \texorpdfstring{\ensuremath{\to}}{->} soil biota depauperate \texorpdfstring{\ensuremath{\to}}{->} poor seedling establishment \\
\end{longtable}
}

Stephens et al.~(1999, \emph{Trends Ecol. Evol.}) distinguished the
\textbf{component Allee effect} (reduction in any fitness component at
low density) from the \textbf{demographic Allee effect} (reduction in
per capita population growth rate). A species may experience component
Allee effects in reproduction without a demographic Allee effect if
compensating survival increases at low density offset the reproductive
reduction.

\begin{Shaded}
\begin{Highlighting}[]
\ImportTok{from}\NormalTok{ biology.ecology }\ImportTok{import}\NormalTok{ allee\_strong\_growth}

\CommentTok{\# Strong Allee: dN/dt = r N (N/A {-} 1)(1 {-} N/K) — same form as Eq. above}
\NormalTok{below }\OperatorTok{=}\NormalTok{ allee\_strong\_growth(N0}\OperatorTok{=}\FloatTok{45.0}\NormalTok{, r}\OperatorTok{=}\FloatTok{0.5}\NormalTok{, A}\OperatorTok{=}\FloatTok{50.0}\NormalTok{, K}\OperatorTok{=}\FloatTok{1000.0}\NormalTok{, t\_end}\OperatorTok{=}\FloatTok{25.0}\NormalTok{, steps}\OperatorTok{=}\DecValTok{2500}\NormalTok{)}
\NormalTok{above }\OperatorTok{=}\NormalTok{ allee\_strong\_growth(N0}\OperatorTok{=}\FloatTok{55.0}\NormalTok{, r}\OperatorTok{=}\FloatTok{0.5}\NormalTok{, A}\OperatorTok{=}\FloatTok{50.0}\NormalTok{, K}\OperatorTok{=}\FloatTok{1000.0}\NormalTok{, t\_end}\OperatorTok{=}\FloatTok{40.0}\NormalTok{, steps}\OperatorTok{=}\DecValTok{4000}\NormalTok{)}
\BuiltInTok{print}\NormalTok{(}\SpecialStringTok{f"Below A: N\_final ≈ }\SpecialCharTok{\{}\NormalTok{below}\SpecialCharTok{.}\NormalTok{populations[}\OperatorTok{{-}}\DecValTok{1}\NormalTok{]}\SpecialCharTok{:.1f\}}\SpecialStringTok{"}\NormalTok{)}
\BuiltInTok{print}\NormalTok{(}\SpecialStringTok{f"Above A: N\_final ≈ }\SpecialCharTok{\{}\NormalTok{above}\SpecialCharTok{.}\NormalTok{populations[}\OperatorTok{{-}}\DecValTok{1}\NormalTok{]}\SpecialCharTok{:.1f\}}\SpecialStringTok{"}\NormalTok{)}
\end{Highlighting}
\end{Shaded}

\begin{quote}
\texorpdfstring{}{} \textbf{Clinical Connection --- Northern White Rhinoceros:} The
northern white rhinoceros (\emph{Ceratotherium simum cottoni})
represents a tragic Allee effect in action. As of 2024, two females
remain (Najin and Fatu at Ol Pejeta Conservancy, Kenya). Even with
preserved sperm and IVF technology, the population has fallen below
plausible demographic and genetic recovery thresholds for conventional
breeding. Stem cell-derived \hyperref[gl:gamete]{\textbf{gamete}}s from
banked fibroblasts represent a highly experimental rescue pathway rather
than a routine conservation tool. This case illustrates that for species
with strong Allee effects, conservation intervention must occur while
demographic options and genetic variation are still large enough to
matter.
\end{quote}

\begin{quote}
\textbf{Concept Check:} A whale population has an Allee threshold of
\(A = 200\) individuals. Current population is \(N = 180\). The
intrinsic growth rate (in the absence of Allee effects) is
\(r = 0.04\)/yr. What will happen to this population without
intervention? What minimum number of individuals would need to be
translocated to push the population above the threshold?
\end{quote}

\begin{center}\rule{0.5\linewidth}{0.5pt}\end{center}

\subsection{Lotka-Volterra Competition and
Predator-Prey}\label{lotka-volterra-competition-and-predator-prey}

\begin{figure}[htbp]
\centering
\includegraphics[width=0.85\textwidth]{../figures/lotka_volterra.png}
\caption{Lotka--Volterra predator--prey dynamics: oscillating population cycles of prey and predator plotted over time and as a phase-plane portrait.}
\label{fig:unit_X_lotka_volterra}
\end{figure}

\subsubsection{Interspecific
Competition}\label{interspecific-competition}

Two species sharing a limiting resource compete. The Lotka-Volterra
competition model:

\begin{equation}
\frac{dN_1}{dt} = r_1 N_1 \left(1 - \frac{N_1 + \alpha_{12} N_2}{K_1}\right)
\label{eq:population_ecology_14}
\end{equation}

\begin{equation}
\frac{dN_2}{dt} = r_2 N_2 \left(1 - \frac{N_2 + \alpha_{21} N_1}{K_2}\right)
\label{eq:population_ecology_15}
\end{equation}

The \textbf{competition coefficient} \(\alpha_{12}\) represents the
per-capita effect of species 2 on species 1, expressed in units of
species-1 equivalents. If \(\alpha_{12} = 0.5\), each individual of
species 2 has half the competitive impact of one individual of species
1.

\textbf{Zero-growth isoclines} (setting \(dN/dt = 0\)): - Species 1
isocline: \(N_1 = K_1 - \alpha_{12}N_2\) (line from \(K_1\) on
\(N_1\)-axis to \(K_1/\alpha_{12}\) on \(N_2\)-axis) - Species 2
isocline: \(N_2 = K_2 - \alpha_{21}N_1\) (line from \(K_2\) on
\(N_2\)-axis to \(K_2/\alpha_{21}\) on \(N_1\)-axis)

\textbf{Phase-plane outcomes:}

{\def\LTcaptype{none} % do not increment counter
\begin{longtable}[]{@{}
  >{\raggedright\arraybackslash}p{(\linewidth - 4\tabcolsep) * \real{0.2195}}
  >{\raggedright\arraybackslash}p{(\linewidth - 4\tabcolsep) * \real{0.1707}}
  >{\raggedright\arraybackslash}p{(\linewidth - 4\tabcolsep) * \real{0.6098}}@{}}
\toprule\noalign{}
\begin{minipage}[b]{\linewidth}\raggedright
Condition
\end{minipage} & \begin{minipage}[b]{\linewidth}\raggedright
Outcome
\end{minipage} & \begin{minipage}[b]{\linewidth}\raggedright
Biological interpretation
\end{minipage} \\
\midrule\noalign{}
\endhead
\bottomrule\noalign{}
\endlastfoot
\(K_1 > K_2/\alpha_{12}\) and \(K_2 < K_1/\alpha_{21}\) & Species 1
wins; species 2 excluded & Species 1 tolerates competition better \\
\(K_2 > K_1/\alpha_{21}\) and \(K_1 < K_2/\alpha_{12}\) & Species 2
wins; species 1 excluded & Species 2 tolerates competition better \\
\(K_1 < K_2/\alpha_{12}\) and \(K_2 < K_1/\alpha_{21}\) & \textbf{Stable
coexistence} & Intraspecific \textgreater{} interspecific competition
(niche differentiation) \\
\(K_1 > K_2/\alpha_{12}\) and \(K_2 > K_1/\alpha_{21}\) &
\textbf{Unstable equilibrium} (priority effect) & Interspecific
\textgreater{} intraspecific; winner depends on initial \(N\) \\
\end{longtable}
}

\textbf{Chesson's modern coexistence theory} (Chesson 2000, \emph{Annu.
Rev.~Ecol. Syst.}): Stable coexistence requires that intraspecific
competition exceed interspecific competition:
\(\alpha_{11}\alpha_{22} > \alpha_{12}\alpha_{21}\), or equivalently,
niche overlap \(< 1\). Two mechanisms enable this:

\begin{enumerate}
\def\labelenumi{\arabic{enumi}.}
\tightlist
\item
  \textbf{Stabilising mechanisms} --- increase niche differentiation
  (negative frequency dependence; rare species advantage)
\item
  \textbf{Equalising mechanisms} --- reduce fitness differences between
  species (similar competitive abilities)
\end{enumerate}

\subsubsection{Lotka-Volterra
Predator-Prey}\label{lotka-volterra-predator-prey}

\begin{equation}
\frac{dV}{dt} = \alpha V - \beta VP \quad \text{(prey: grow in absence of predator; decline with encounters)}
\label{eq:population_ecology_16}
\end{equation}

\begin{equation}
\frac{dP}{dt} = \delta VP - \gamma P \quad \text{(predator: grow via prey; decline without prey)}
\label{eq:population_ecology_17}
\end{equation}

This system produces \textbf{neutrally stable oscillations} with period
\(\approx 2\pi/\sqrt{\alpha\gamma}\) and amplitude that depends on
initial conditions, tracing the offset prey--predator cycles and closed
phase-plane orbit shown in \cref{fig:unit_X_lotka_volterra}.

\textbf{Functional responses} \citep{holling1959} describe the per
capita predation rate as a function of prey density:

{\def\LTcaptype{none} % do not increment counter
\begin{longtable}[]{@{}
  >{\raggedright\arraybackslash}p{(\linewidth - 6\tabcolsep) * \real{0.1667}}
  >{\raggedright\arraybackslash}p{(\linewidth - 6\tabcolsep) * \real{0.3333}}
  >{\raggedright\arraybackslash}p{(\linewidth - 6\tabcolsep) * \real{0.2083}}
  >{\raggedright\arraybackslash}p{(\linewidth - 6\tabcolsep) * \real{0.2917}}@{}}
\toprule\noalign{}
\begin{minipage}[b]{\linewidth}\raggedright
Type
\end{minipage} & \begin{minipage}[b]{\linewidth}\raggedright
Equation
\end{minipage} & \begin{minipage}[b]{\linewidth}\raggedright
Shape
\end{minipage} & \begin{minipage}[b]{\linewidth}\raggedright
Example
\end{minipage} \\
\midrule\noalign{}
\endhead
\bottomrule\noalign{}
\endlastfoot
\textbf{Type I} & \(f(V) = aV\) & Linear & Filter feeders (whale sharks,
baleen whales) \\
\textbf{Type II} & \(f(V) = \frac{aV}{1 + ahV}\) & Decelerating
(hyperbolic) & Most predators; handling time \(h\) limits intake \\
\textbf{Type III} & \(f(V) = \frac{aV^2}{1 + ahV^2}\) & Sigmoidal &
Predators that switch prey; search image formation \\
\end{longtable}
}

The Type II functional response is also known as the
\textbf{Michaelis-Menten} or \textbf{disc equation}. The parameter \(a\)
= attack rate and \(h\) = handling time per prey item. Maximum
consumption rate = \(1/h\).

Type III responses are stabilising because at low prey density,
per-capita predation risk is low (predator switching to alternative
prey), creating a \textbf{low-density refuge}.

\textbf{Numerical response:} predator population growth rate as a
function of prey density --- driven by the functional response but with
a time lag (gestation, maturation).

\begin{Shaded}
\begin{Highlighting}[]
\ImportTok{from}\NormalTok{ biology.ecology }\ImportTok{import}\NormalTok{ lotka\_volterra}

\NormalTok{result }\OperatorTok{=}\NormalTok{ lotka\_volterra(}
    \FloatTok{100.0}\NormalTok{, }\FloatTok{20.0}\NormalTok{, }\FloatTok{1.0}\NormalTok{, }\FloatTok{0.1}\NormalTok{, }\FloatTok{0.075}\NormalTok{, }\FloatTok{1.5}\NormalTok{, t\_end}\OperatorTok{=}\FloatTok{30.0}\NormalTok{, steps}\OperatorTok{=}\DecValTok{300}\NormalTok{,}
\NormalTok{)}
\BuiltInTok{print}\NormalTok{(}\SpecialStringTok{f"Prey range: }\SpecialCharTok{\{}\BuiltInTok{min}\NormalTok{(result.prey)}\SpecialCharTok{:.0f\}}\SpecialStringTok{–}\SpecialCharTok{\{}\BuiltInTok{max}\NormalTok{(result.prey)}\SpecialCharTok{:.0f\}}\SpecialStringTok{"}\NormalTok{)}
\BuiltInTok{print}\NormalTok{(}\SpecialStringTok{f"Predator range: }\SpecialCharTok{\{}\BuiltInTok{min}\NormalTok{(result.predator)}\SpecialCharTok{:.0f\}}\SpecialStringTok{–}\SpecialCharTok{\{}\BuiltInTok{max}\NormalTok{(result.predator)}\SpecialCharTok{:.0f\}}\SpecialStringTok{"}\NormalTok{)}
\end{Highlighting}
\end{Shaded}

\subsubsection{Lynx-Hare Cycle: A Case
Study}\label{lynx-hare-cycle-a-case-study}

The snowshoe hare (\emph{Lepus americanus}) and Canada lynx (\emph{Lynx
canadensis}) display coupled about 10-year population cycles
across boreal Canada, documented from Hudson's Bay Company fur trading
records (1845-1935; Keith 1963). Modern analysis reveals the cycle is
more complex than simple Lotka-Volterra:

\begin{enumerate}
\def\labelenumi{\arabic{enumi}.}
\tightlist
\item
  \textbf{Predation} accounts for about 60\% of the cycle
  driver (Krebs et al.~2001, \emph{Science})
\item
  \textbf{Food quality} --- at high hare density, overgrazing induces
  plant chemical defences (phenolic compounds in willow/birch browse),
  reducing hare food quality (about 20\% of cycle)
\item
  \textbf{Stress physiology} --- high density and predation risk elevate
  \hyperref[gl:cortisol]{\textbf{cortisol}}, suppressing reproduction
  (about 20\% of cycle)
\item
  The \textbf{tri-trophic interaction} (vegetation-hare-lynx) is
  necessary to explain the observed amplitude and period
\end{enumerate}

\begin{quote}
\textbf{Concept Check:} In the basic Lotka-Volterra predator-prey model,
the oscillations are neutrally stable. What modification to the model
would produce damped oscillations converging to a stable equilibrium?
(Hint: consider what happens when you add density dependence to the prey
equation.)
\end{quote}

\begin{center}\rule{0.5\linewidth}{0.5pt}\end{center}

\subsection{r vs.~K Selection and Life History
Theory}\label{r-vs.-k-selection-and-life-history-theory}

\citet{macarthur1967} proposed that
\hyperref[gl:natural-selection]{\textbf{natural selection}} shapes life
histories along an r-K continuum:

{\def\LTcaptype{none} % do not increment counter
\begin{longtable}[]{@{}
  >{\raggedright\arraybackslash}p{(\linewidth - 4\tabcolsep) * \real{0.1613}}
  >{\raggedright\arraybackslash}p{(\linewidth - 4\tabcolsep) * \real{0.4194}}
  >{\raggedright\arraybackslash}p{(\linewidth - 4\tabcolsep) * \real{0.4194}}@{}}
\toprule\noalign{}
\begin{minipage}[b]{\linewidth}\raggedright
Trait
\end{minipage} & \begin{minipage}[b]{\linewidth}\raggedright
\hyperref[gl:r-strategist]{\textbf{r-strategist}}s
\end{minipage} & \begin{minipage}[b]{\linewidth}\raggedright
\hyperref[gl:k-strategist]{\textbf{K-strategist}}s
\end{minipage} \\
\midrule\noalign{}
\endhead
\bottomrule\noalign{}
\endlastfoot
Body size & Small & Large \\
Life span & Short & Long \\
Maturation time & Fast & Slow \\
Offspring number & Many; small & Few; large \\
Parental investment & Minimal & Extensive \\
Population dynamics & Exponential; boom/bust; opportunistic & Near K;
stable; competitive \\
Environmental tolerance & High variability; unpredictable & Low
variability; stable niches \\
Survivorship curve & Type III & Type I \\
Examples & Dandelions, flies, bacteria, aphids & Elephants, whales,
great apes, condors \\
\end{longtable}
}

\begin{figure}[htbp]
\centering
\includegraphics[width=0.82\linewidth,height=4.2in,keepaspectratio]{../figures/mermaid_inline/inline_mermaid_0177_c0b67a78890a.png}
\caption{Flowchart for r vs. K Selection and Life History Theory: r-selected High fecundity Small body Short lifespan Type III survivorship Colonisers, K-selected Low fecundity Large body Long lifespan Type I survivorship Competitors, Bacteria r about  40/day, and Insects r about  10/yr form the diagram's primary path or branches.}
\end{figure}

\subsubsection{Beyond r-K: Modern Life History
Theory}\label{beyond-r-k-modern-life-history-theory}

The r-K framework has been largely superseded by more nuanced models:

\textbf{Stearns' (1992) demographic classification} uses age-specific
mortality patterns rather than a single r-K axis, predicting that: -
High juvenile mortality \texorpdfstring{\ensuremath{\to}}{->} early maturation, high fecundity - High adult
mortality \texorpdfstring{\ensuremath{\to}}{->} delayed maturation, high reproductive effort per event -
Variable mortality \texorpdfstring{\ensuremath{\to}}{->} bet-hedging strategies

\textbf{Bet-hedging} (Philippi \& Seger 1989, \emph{Trends Ecol.
Evol.}): In variable environments, intermediate-risk reproductive
strategies outperform extreme approaches by reducing \textbf{variance}
in fitness rather than maximising mean fitness: - \textbf{Conservative
bet-hedging:} Low-variance strategy (e.g., typically produce moderate
number of offspring) - \textbf{Diversified bet-hedging:} Produce
variable \hyperref[gl:phenotype]{\textbf{phenotype}}s (e.g., desert
annual plants with variable seed dormancy --- don't germinate most seeds
in a single good year)

\textbf{Grime's CSR triangle} (plants): Competitors (large, slow),
Stress-tolerators (tough, slow), Ruderals (fast, weedy) --- a three-way
classification rather than binary r-K.

\begin{quote}
\texorpdfstring{}{} \textbf{Clinical Connection --- Human Life History and Demographic
Transition:} Human populations have undergone a dramatic
\textbf{demographic transition} as societies industrialise: from high
birth rates and high death rates (Stage 1) through declining death rates
(Stage 2, rapid population growth), declining birth rates (Stage 3), to
low birth/death rates (Stage 4). Some developed nations are now in Stage
5 with birth rates below replacement (\(R_0 < 1\)). Japan's population
declined by about 800,000 in 2023 alone
(\(r \approx -0.006\)/yr). The demographic transition illustrates a
shift from r-selected to K-selected reproductive strategy within a
single species as environmental conditions change.
\end{quote}

\begin{center}\rule{0.5\linewidth}{0.5pt}\end{center}

\subsection{7B Age-Structured Models and the Leslie
Matrix}\label{b-age-structured-models-and-the-leslie-matrix}

The scalar growth models treat most individuals as demographically
identical. Real populations are structured by age (or stage), and
demographic rates are highly age-specific: a one-year-old salmon and a
four-year-old salmon contribute very differently to population growth.
\textbf{Matrix population models} \citep{lotka1925} make this structure
explicit and underlie virtually every modern population viability
analysis, \hyperref[gl:fisheries]{\textbf{fisheries}} stock assessment,
and conservation triage exercise.

\subsubsection{Leslie Matrix
Construction}\label{leslie-matrix-construction}

A \textbf{Leslie matrix} \(\mathbf{L}\) is a square non-negative matrix
that projects an age-classified population vector
\(\mathbf{n}(t) = (n_0, n_1, \ldots, n_\omega)^T\) from time \(t\) to
\(t+1\):

\begin{equation}
\mathbf{n}(t+1) = \mathbf{L}\,\mathbf{n}(t)
\label{eq:unit_X_leslie_projection}
\end{equation}

For a population with maximum age \texorpdfstring{\ensuremath{\omega}}{omega}, the Leslie matrix has age-specific
fecundities \(F_x = l_x m_x / l_{x-1}\) on the top row and survival
probabilities \(P_x = l_{x+1}/l_x\) on the sub-diagonal:

\begin{equation}
\mathbf{L} = \begin{pmatrix}
F_0 & F_1 & F_2 & \cdots & F_{\omega-1} & F_\omega \\
P_0 & 0   & 0   & \cdots & 0           & 0 \\
0   & P_1 & 0   & \cdots & 0           & 0 \\
0   & 0   & P_2 & \cdots & 0           & 0 \\
\vdots & \vdots & \vdots & \ddots & \vdots & \vdots \\
0   & 0   & 0   & \cdots & P_{\omega-1} & 0
\end{pmatrix}
\label{eq:unit_X_leslie_matrix}
\end{equation}

By the \textbf{Perron--Frobenius theorem}, \(\mathbf{L}\) has a unique
positive dominant eigenvalue \(\lambda_1\) --- this is the
\textbf{asymptotic finite rate of increase}. The associated right
eigenvector \(\mathbf{w}\) gives the \textbf{stable age distribution}
(the proportions \(n_x/N\) to which the population converges) and the
left eigenvector \(\mathbf{v}\) gives the \textbf{age-specific
reproductive value} \(v_x\) (\cref{eq:population_ecology_5}, but now
read off the matrix directly).

\begin{equation}
\mathbf{L}\mathbf{w} = \lambda_1 \mathbf{w}, \qquad \mathbf{v}^T \mathbf{L} = \lambda_1 \mathbf{v}^T
\label{eq:unit_X_leslie_eigen}
\end{equation}

The Euler-Lotka equation (\cref{eq:population_ecology_1}) is the
characteristic polynomial of \(\mathbf{L}\) --- every demographic
statement made in the life-table and survivorship discussion is
recoverable from the matrix.

\subsubsection{Sensitivity and Elasticity
Analysis}\label{sensitivity-and-elasticity-analysis}

Conservation managers rarely have the luxury of changing every
demographic rate. Which rate, if perturbed by 1\%, would change
\(\lambda_1\) the most? \textbf{Sensitivity} \(s_{ij}\) and
\textbf{elasticity} \(e_{ij}\) answer this:

\begin{equation}
s_{ij} = \frac{\partial \lambda_1}{\partial a_{ij}} = \frac{v_i w_j}{\langle \mathbf{v}, \mathbf{w} \rangle}, \qquad e_{ij} = \frac{a_{ij}}{\lambda_1}\,s_{ij}
\label{eq:unit_X_sensitivity_elasticity}
\end{equation}

where \(a_{ij}\) is the matrix entry,
\(\langle \mathbf{v}, \mathbf{w}\rangle = \sum_x v_x w_x\), and
elasticities sum to one (\(\sum_{ij} e_{ij} = 1\)), making them directly
interpretable as proportional contributions.

{\def\LTcaptype{none} % do not increment counter
\begin{longtable}[]{@{}
  >{\raggedright\arraybackslash}p{(\linewidth - 4\tabcolsep) * \real{0.3200}}
  >{\raggedright\arraybackslash}p{(\linewidth - 4\tabcolsep) * \real{0.5600}}
  >{\raggedright\arraybackslash}p{(\linewidth - 4\tabcolsep) * \real{0.1200}}@{}}
\toprule\noalign{}
\begin{minipage}[b]{\linewidth}\raggedright
Quantity
\end{minipage} & \begin{minipage}[b]{\linewidth}\raggedright
Interpretation
\end{minipage} & \begin{minipage}[b]{\linewidth}\raggedright
Use
\end{minipage} \\
\midrule\noalign{}
\endhead
\bottomrule\noalign{}
\endlastfoot
Sensitivity \(s_{ij}\) & Absolute change in \(\lambda_1\) per unit
change in \(a_{ij}\) & Rates measured on different scales (survival
0--1, fecundity 0--\(\infty\)) --- direct comparison misleading \\
Elasticity \(e_{ij}\) & Proportional change in \(\lambda_1\) per
proportional change in \(a_{ij}\) & Comparing across rates; identifying
highest-leverage management target \\
\end{longtable}
}

\textbf{Worked Example --- Loggerhead sea turtle.} Crouse, Crowder \&
Caswell (1987, \emph{Ecology}) built a 7-stage Leslie-style matrix for
\emph{Caretta caretta} in the southeastern USA. Hatchling protection
(the dominant 1980s strategy) had elasticity \(\approx 0.04\), whereas
adult survival had elasticity \(\approx 0.20-0.30\). A 1 \% increase in
adult survival raised \(\lambda_1\) five to seven times more than a 1 \%
increase in egg-stage survival. This elasticity result drove the
regulatory adoption of \textbf{Turtle Excluder Devices} in shrimp trawls
--- a direct policy outcome of matrix sensitivity analysis. The same
logic explains why protecting reproductive-aged females beats protecting
eggs, echoing the reproductive-value argument developed earlier.

\begin{quote}
\textbf{Concept Check:} A 3-stage matrix has elasticities
\(e_{\text{juvenile survival}} = 0.05\),
\(e_{\text{subadult survival}} = 0.35\),
\(e_{\text{adult survival}} = 0.50\),
\(e_{\text{adult fecundity}} = 0.10\). Which life-stage transition
should a recovery plan target first, and why is ``save the babies''
intuition misleading here?
\end{quote}

\subsubsection{Stage-Structured
Generalisations}\label{stage-structured-generalisations}

Many organisms are better classified by \textbf{stage} (size, instar,
reproductive status) than by age --- trees, corals, and
indeterminate-growth fish are obvious examples. The \textbf{Lefkovitch
matrix} generalises Leslie by allowing positive entries on the diagonal
(probability of remaining in the same stage) and on multiple
sub-diagonals (skipping stages, regression). The eigen-analysis
machinery is shared with Leslie models, while the biological
interpretation of each entry changes.

\begin{Shaded}
\begin{Highlighting}[]
\ImportTok{import}\NormalTok{ numpy }\ImportTok{as}\NormalTok{ np}

\CommentTok{\# Loggerhead{-}style 5{-}stage matrix (illustrative)}
\NormalTok{L }\OperatorTok{=}\NormalTok{ np.array([}
\NormalTok{    [}\FloatTok{0.00}\NormalTok{, }\FloatTok{0.00}\NormalTok{, }\FloatTok{0.00}\NormalTok{, }\FloatTok{4.66}\NormalTok{, }\FloatTok{61.9}\NormalTok{],   }\CommentTok{\# fecundities (eggs, hatchlings)}
\NormalTok{    [}\FloatTok{0.68}\NormalTok{, }\FloatTok{0.74}\NormalTok{, }\FloatTok{0.00}\NormalTok{, }\FloatTok{0.00}\NormalTok{, }\FloatTok{0.00}\NormalTok{],   }\CommentTok{\# juvenile survival}
\NormalTok{    [}\FloatTok{0.00}\NormalTok{, }\FloatTok{0.05}\NormalTok{, }\FloatTok{0.66}\NormalTok{, }\FloatTok{0.00}\NormalTok{, }\FloatTok{0.00}\NormalTok{],   }\CommentTok{\# subadult transition}
\NormalTok{    [}\FloatTok{0.00}\NormalTok{, }\FloatTok{0.00}\NormalTok{, }\FloatTok{0.02}\NormalTok{, }\FloatTok{0.68}\NormalTok{, }\FloatTok{0.00}\NormalTok{],   }\CommentTok{\# novice breeder}
\NormalTok{    [}\FloatTok{0.00}\NormalTok{, }\FloatTok{0.00}\NormalTok{, }\FloatTok{0.00}\NormalTok{, }\FloatTok{0.06}\NormalTok{, }\FloatTok{0.81}\NormalTok{],   }\CommentTok{\# mature breeder}
\NormalTok{])}
\NormalTok{eigvals, eigvecs }\OperatorTok{=}\NormalTok{ np.linalg.eig(L)}
\NormalTok{lam1 }\OperatorTok{=} \BuiltInTok{max}\NormalTok{(eigvals.real)}
\NormalTok{w }\OperatorTok{=}\NormalTok{ eigvecs[:, np.argmax(eigvals.real)].real}
\NormalTok{w }\OperatorTok{=}\NormalTok{ w }\OperatorTok{/}\NormalTok{ w.}\BuiltInTok{sum}\NormalTok{()}
\BuiltInTok{print}\NormalTok{(}\SpecialStringTok{f"lambda\_1 = }\SpecialCharTok{\{}\NormalTok{lam1}\SpecialCharTok{:.3f\}}\SpecialStringTok{"}\NormalTok{)}
\BuiltInTok{print}\NormalTok{(}\SpecialStringTok{f"stable stage distribution = }\SpecialCharTok{\{}\NormalTok{np}\SpecialCharTok{.}\BuiltInTok{round}\NormalTok{(w, }\DecValTok{3}\NormalTok{)}\SpecialCharTok{\}}\SpecialStringTok{"}\NormalTok{)}
\CommentTok{\# elasticity of adult survival — the loggerhead conservation lever}
\end{Highlighting}
\end{Shaded}

\begin{center}\rule{0.5\linewidth}{0.5pt}\end{center}

\subsection{7C Individual-Based Models and Stochastic
Simulation}\label{c-individual-based-models-and-stochastic-simulation}

Matrix models are deterministic and assume infinite, well-mixed
populations. Two regimes break those assumptions and demand a different
tool:

\begin{enumerate}
\def\labelenumi{\arabic{enumi}.}
\tightlist
\item
  \textbf{Small populations} (conservation contexts) where demographic
  stochasticity dominates --- a population of 30 condors with mean
  \(\lambda_1 = 1.05\) can still go extinct from an unlucky run of
  male-heavy chicks.
\item
  \textbf{Heterogeneous individuals} --- variation in body size,
  behaviour, territory quality, or local conditions that cannot be
  averaged into a mean rate without losing the dynamics.
\end{enumerate}

\textbf{Individual-based models (IBMs)}, also called agent-based models,
simulate each organism explicitly. Births, deaths, dispersal, and
interactions are drawn as random events from probability distributions
parameterised by individual attributes. Across many replicate runs, IBMs
produce \emph{distributions} of outcomes (extinction probability, time
to recovery) rather than single trajectories.

{\def\LTcaptype{none} % do not increment counter
\begin{longtable}[]{@{}
  >{\raggedright\arraybackslash}p{(\linewidth - 2\tabcolsep) * \real{0.3214}}
  >{\raggedright\arraybackslash}p{(\linewidth - 2\tabcolsep) * \real{0.6786}}@{}}
\toprule\noalign{}
\begin{minipage}[b]{\linewidth}\raggedright
Advantage
\end{minipage} & \begin{minipage}[b]{\linewidth}\raggedright
Why IBMs deliver it
\end{minipage} \\
\midrule\noalign{}
\endhead
\bottomrule\noalign{}
\endlastfoot
Demographic stochasticity & Discrete birth/death events --- no
fractional individuals \\
Individual heterogeneity & Each agent carries its own age, size,
genotype, location \\
Spatial structure & Landscape grids, dispersal kernels, local
interactions \\
Emergent dynamics & Allee effects,
\hyperref[gl:self-organization]{\textbf{self-organization}}, threshold
behaviours arise naturally \\
Complex behaviours & Mate choice, learning, territoriality without
analytical closed form \\
\end{longtable}
}

\textbf{Limitations:} IBMs trade analytical clarity for realism. They
are computationally expensive (millions of replicates needed for
tail-probability estimates), parameter-hungry, and validation is harder
than for closed-form models. Best practice: build the simplest model
that captures the focal mechanism, then add complexity primarily when
justified by data (the \textbf{Pattern-Oriented Modelling} philosophy of
Grimm \& Railsback 2005).

\begin{quote}
\textbf{Concept Check:} Two black-footed ferret populations both have
mean \(\lambda_1 = 1.02\). Population A has \(N = 500\); population B
has \(N = 25\). Why does the IBM extinction probability differ vastly
between them despite identical mean growth rates? What would a
deterministic model predict?
\end{quote}

\begin{center}\rule{0.5\linewidth}{0.5pt}\end{center}

\subsection{7D Population Viability Analysis
(PVA)}\label{d-population-viability-analysis-pva}

\textbf{Population viability analysis} integrates matrix models and
stochastic simulation into a quantitative extinction-risk forecast. PVA
is the formal answer to ``How long does this species have, and what
intervention buys the most time?'' --- required by the U.S. Endangered
Species Act, the IUCN Red List Criterion E, and most national recovery
plans.

\subsubsection{Inputs to a PVA}\label{inputs-to-a-pva}

{\def\LTcaptype{none} % do not increment counter
\begin{longtable}[]{@{}
  >{\raggedright\arraybackslash}p{(\linewidth - 4\tabcolsep) * \real{0.2000}}
  >{\raggedright\arraybackslash}p{(\linewidth - 4\tabcolsep) * \real{0.2400}}
  >{\raggedright\arraybackslash}p{(\linewidth - 4\tabcolsep) * \real{0.5600}}@{}}
\toprule\noalign{}
\begin{minipage}[b]{\linewidth}\raggedright
Input
\end{minipage} & \begin{minipage}[b]{\linewidth}\raggedright
Source
\end{minipage} & \begin{minipage}[b]{\linewidth}\raggedright
Why it matters
\end{minipage} \\
\midrule\noalign{}
\endhead
\bottomrule\noalign{}
\endlastfoot
Stage- or age-specific vital rates & Mark-recapture, life table & Mean
dynamics \\
Variance in vital rates & Long-term monitoring & Environmental
stochasticity \\
Catastrophe frequency and severity & Disturbance history & Tail risk
(fire, disease, storm) \\
Density dependence (often theta-logistic,
\cref{eq:population_ecology_11}) & Time-series fits & Carrying-capacity
feedback \\
Initial age/stage distribution & Field census & Starting condition \\
Habitat capacity \(K\) & Resource estimates & Ceiling on population \\
Inbreeding cost on vital rates & Pedigree or ROH & Genetic stochasticity
(link to \cref{sec:unit_X_biomes_and_conservation}) \\
\end{longtable}
}

\subsubsection{Outputs and Decision
Rules}\label{outputs-and-decision-rules}

The two headline outputs are \textbf{extinction probability}
\(P_{ext}(T)\) over a horizon \(T\) (commonly 100 years or 100
generations) and \textbf{mean time to extinction} \(\bar{T}_{ext}\).

\begin{equation}
P_{ext}(T) = \Pr\bigl(\,N(t) < N_{\text{crit}}\ \text{for some } t \le T\bigr)
\label{eq:unit_X_extinction_probability}
\end{equation}

The threshold \(N_{\text{crit}}\) (often called the
\textbf{quasi-extinction threshold}) is set above zero --- typically the
Allee threshold \(A\) (\cref{eq:population_ecology_13}) or a number
below which the population is ``functionally extinct.'' Hitting
\(N_{\text{crit}}\) counts as failure even if the simulation has not
bottomed out.

The \textbf{minimum viable population (MVP)} in standard usage is the
smallest \(N_0\) such that \(P_{ext}(100\,\text{yr}) \le 0.05\). Reed et
al.~(2003, \emph{Biol. Conserv.}) reviewed PVAs across 102 vertebrates
and found MVPs clustering at 5,000--7,000 adults --- far above the
classical \(N_e = 500\) rule of thumb (which ignores demographic and
environmental stochasticity) and a key empirical result that pushed
conservation toward the \textbf{100/1000 rule} discussed in
\cref{sec:unit_X_biomes_and_conservation}.

\subsubsection{Worked Example --- California Condor
PVA}\label{worked-example-california-condor-pva}

In 1982 the known California condor (\emph{Gymnogyps californianus})
population had fallen to 22 birds. A PVA built around captive-breeding
fecundity, juvenile mortality from lead poisoning, and stochastic
catastrophe (wildfire) gave \(P_{ext}(50\,\text{yr}) > 0.99\) for the
no-intervention scenario. The same model showed
\(P_{ext}(50\,\text{yr}) < 0.10\) if (i) the remaining birds were
captured for breeding, and (ii) released birds were protected from lead
ammunition. The captive-breeding decision, controversial at the time,
was justified directly by PVA elasticities --- adult survival and
juvenile lead mortality had the largest leverage on \(\lambda_1\). The
wild population is now \textgreater{} 250 (see
\cref{sec:unit_X_biomes_and_conservation}).

\subsubsection{Limitations and Honest
Use}\label{limitations-and-honest-use}

PVA outputs depend on the model structure, input data, and uncertainty
assumptions; small-population data are often sparse and biased toward
observed survivors. Best practice (Beissinger \& Westphal 1998, \emph{J.
Wildl. Manage.}):

\begin{itemize}
\tightlist
\item
  Treat absolute \(P_{ext}\) values with skepticism; use \emph{relative}
  comparisons across management scenarios
\item
  Report sensitivity to parameter uncertainty (Monte Carlo over input
  distributions)
\item
  Update PVAs as monitoring data accumulate --- they are living tools,
  not one-off reports
\item
  Avoid the temptation to over-parameterise; structural uncertainty
  often exceeds parameter uncertainty
\end{itemize}

\begin{quote}
\textbf{Concept Check:} A PVA gives \(P_{ext}(100\,\text{yr}) = 0.30\)
under current management and 0.05 under a proposed habitat-corridor
intervention. Should you trust the absolute numbers, the difference, or
both? What three sources of uncertainty would you investigate before
recommending the corridor?
\end{quote}

\begin{center}\rule{0.5\linewidth}{0.5pt}\end{center}

\subsection{Metapopulation Dynamics}\label{metapopulation-dynamics}

A \textbf{metapopulation} (Levins 1969; Hanski 1994) is a ``population
of populations'' --- a set of spatially separated subpopulations
(patches) connected by dispersal. Local extinctions are balanced by
recolonisation from other patches.

In conservation practice, the important question is usually whether
dispersal is functional, not whether habitat patches are merely drawn on
a map. Roads, dams, fences, heat waves, disease, and behavioural
avoidance can lower effective colonisation even when patches appear
close. Corridor design therefore needs demographic monitoring, genetic
connectivity, mortality-risk mapping, and sensitivity analysis for \(c\)
and \(e\), because a corridor that increases movement can still fail if
it also increases mortality or spreads disease.

\subsubsection{The Levins Model}\label{the-levins-model}

\begin{equation}
\frac{dp}{dt} = cp(1-p) - ep
\label{eq:population_ecology_18}
\end{equation}

where: - \(p\) = proportion of occupied patches - \(c\) = colonisation
rate - \(e\) = extinction rate per patch

At equilibrium: \(\hat{p} = 1 - e/c\)

The metapopulation persists if \(e/c < 1\) (colonisation rate exceeds
extinction rate). If habitat patches are destroyed (reducing total
patches available), the effective \(c\) decreases, potentially pushing
\(e/c > 1\) and causing \textbf{metapopulation collapse}.

\subsubsection{\texorpdfstring{Source-Sink Dynamics
\citep{pulliam1988}}{Source-Sink Dynamics}}\label{source-sink-dynamics}

Patches differ in leverage:

{\def\LTcaptype{none} % do not increment counter
\begin{longtable}[]{@{}
  >{\raggedright\arraybackslash}p{(\linewidth - 4\tabcolsep) * \real{0.4000}}
  >{\raggedright\arraybackslash}p{(\linewidth - 4\tabcolsep) * \real{0.4400}}
  >{\raggedright\arraybackslash}p{(\linewidth - 4\tabcolsep) * \real{0.1600}}@{}}
\toprule\noalign{}
\begin{minipage}[b]{\linewidth}\raggedright
Patch type
\end{minipage} & \begin{minipage}[b]{\linewidth}\raggedright
Growth rate
\end{minipage} & \begin{minipage}[b]{\linewidth}\raggedright
Role
\end{minipage} \\
\midrule\noalign{}
\endhead
\bottomrule\noalign{}
\endlastfoot
\textbf{Source} & \(\lambda > 1\) & Produces surplus individuals; net
emigration \\
\textbf{Sink} & \(\lambda < 1\) & Cannot sustain itself; relies on
immigration from sources \\
\end{longtable}
}

A population can persist in sink habitats indefinitely if immigration
from sources is sufficient. This has profound conservation implications:
protecting source habitats is critical, while sink habitats may appear
to support populations but are ``ecological traps'' without the source.

\subsubsection{Rescue Effect and
Connectivity}\label{rescue-effect-and-connectivity}

The \textbf{rescue effect} (Brown \& Kodric-Brown 1977): immigration
from nearby patches prevents or delays local extinction. This means: -
Closely spaced patches have lower extinction rates - Habitat corridors
increase connectivity and reduce local extinction - Isolation increases
both extinction probability and time to recolonisation

\begin{quote}
\textbf{Concept Check:} A metapopulation has 100 habitat patches with
\(c = 0.3\) and \(e = 0.1\). Calculate the equilibrium proportion of
occupied patches. If 40\% of patches are destroyed (reducing total to 60
patches), what happens to the metapopulation? (Consider how reduced
patch number affects colonisation rate.)
\end{quote}

\begin{center}\rule{0.5\linewidth}{0.5pt}\end{center}

\subsubsection{Visualizing Metapopulation and Source-Sink
Dynamics}\label{visualizing-metapopulation-and-source-sink-dynamics}

\begin{figure}[htbp]
\centering
\includegraphics[width=0.82\linewidth,height=4.2in,keepaspectratio]{../figures/mermaid_inline/inline_mermaid_0178_477737b69ed1.png}
\caption{Flowchart for Visualizing Metapopulation and Source-Sink Dynamics: Source Patch A lambda > 1 High quality habitat Net Emigration, Source Patch B lambda > 1 High quality habitat Net Emigration, Sink Patch C lambda Poor quality habitat Net Immigration, and Sink Patch D lambda Ecological Trap? Net Immigration form the diagram's primary path or branches.}
\end{figure}

\begin{center}\rule{0.5\linewidth}{0.5pt}\end{center}

\subsection{Estimating Population
Size}\label{estimating-population-size}

\subsubsection{Mark-Recapture Methods}\label{mark-recapture-methods}

The \textbf{Lincoln-Petersen method} (simplest mark-recapture) estimates
population size from two sampling events:

\begin{equation}
\hat{N} = \frac{M \cdot C}{R}
\label{eq:population_ecology_19}
\end{equation}

where: - \(M\) = number marked and released in first sample - \(C\) =
total captured in second sample - \(R\) = number of marked individuals
recaptured

\textbf{Assumptions:} (1) Closed population (no births, deaths,
immigration, emigration between samples), (2) marks are not lost, (3)
marked and unmarked individuals have equal capture probability, (4)
marks do not affect survival.

\textbf{Chapman's correction} for small samples:

\begin{equation}
\hat{N} = \frac{(M+1)(C+1)}{R+1} - 1
\label{eq:population_ecology_20}
\end{equation}

\textbf{Jolly-Seber method} --- for open populations with multiple
sampling occasions. Estimates \(N\), survival rate (\texorpdfstring{\ensuremath{\varphi}}{phi}), and recruitment
(\(B\)) at each time point. Computationally intensive; requires at least
3 capture occasions.

\subsubsection{Distance Sampling}\label{distance-sampling}

\textbf{Line transect distance sampling} estimates density from the
perpendicular distances of detected objects from the transect line:

\begin{equation}
\hat{D} = \frac{n \cdot \hat{f}(0)}{2L}
\label{eq:population_ecology_21}
\end{equation}

where \(n\) = number detected, \(L\) = total transect length,
\(\hat{f}(0)\) = estimated probability density function of detection
distances evaluated at zero distance. The key assumption is that most
objects on the line are detected (\(g(0) = 1\)).

\begin{quote}
\texorpdfstring{}{} \textbf{Clinical Connection --- Epidemiological Surveillance:}
Mark-recapture methodology has been adapted for human epidemiology as
\textbf{capture-recapture estimation} to estimate the completeness of
disease registries. For example, comparing hospital records with death
certificates as two independent ``captures'' allows estimation of the
true number of opioid overdose deaths, accounting for under-reporting.
The WHO uses this approach for HIV prevalence estimation in populations
where direct surveys are impractical.
\end{quote}

\begin{center}\rule{0.5\linewidth}{0.5pt}\end{center}

\subsection{Human Population
Demographics}\label{human-population-demographics}

The global human population reached 8 billion in November 2022. The UN
World Population Prospects 2024 projects continued growth to about 10.3
billion in the mid-2080s, followed by a slight decline toward about 10.2
billion by 2100 \citep{un2024population}. Treat this as a scenario-based
projection rather than a fixed destiny: fertility, mortality, migration,
policy, and environmental stress can shift regional trajectories.

\subsubsection{Age Structure and Population
Pyramids}\label{age-structure-and-population-pyramids}

The \textbf{age structure} of a population determines its future growth
trajectory:

\begin{itemize}
\tightlist
\item
  \textbf{Expansive pyramid} (broad base): high proportion of
  pre-reproductive individuals; rapid growth (e.g., Nigeria,
  \(r \approx 0.025\)/yr)
\item
  \textbf{Stationary pyramid} (column shape): roughly equal proportions
  across age classes; near-zero growth (e.g., USA,
  \(r \approx 0.002\)/yr)
\item
  \textbf{Constrictive pyramid} (narrow base): higher proportion of
  older individuals; declining population (e.g., Japan,
  \(r \approx -0.006\)/yr)
\end{itemize}

\textbf{Population momentum:} Even after fertility drops to replacement
level (\(R_0 = 1\)), a population with an expansive age structure
continues growing for decades because the large cohort of young people
has not yet reproduced. This explains why the world population will
continue growing despite declining global fertility rates (from 5.3 in
1963 to 2.3 in 2023).

\subsubsection{Total Fertility Rate (TFR) and Replacement
Level}\label{total-fertility-rate-tfr-and-replacement-level}

The \textbf{total fertility rate} is the average number of children born
to a woman over her lifetime. \textbf{Replacement-level fertility} is
TFR \(\approx 2.1\) in developed nations (slightly above 2 to account
for child mortality). In 2023, global TFR was 2.3 --- barely above
replacement.

\textbf{Demographic dividend:} Countries in mid-transition (declining
fertility, large working-age cohort, low dependency ratio) experience
accelerated economic growth. East Asia (1965-1990) and sub-Saharan
Africa (projected 2030-2060) demonstrate this phenomenon.

\begin{center}\rule{0.5\linewidth}{0.5pt}\end{center}

\subsection{Worked Examples}\label{worked-examples-3}

\subsubsection{Worked Example 1: Leslie Matrix
Projection}\label{worked-example-1-leslie-matrix-projection}

\textbf{Problem:} A 3-age-class population (juvenile, subadult, adult)
has Leslie matrix
\[\mathbf{L} = \begin{pmatrix} 0 & 1.2 & 3.0 \\ 0.4 & 0 & 0 \\ 0 & 0.6 & 0.7 \end{pmatrix} \label{eq:unit_X_population_ecology_item_1}\]

with initial vector \(\mathbf{n}(0) = (40, 20, 10)^T\). (a) Project the
population one time step. (b) Compute \(\lambda_1\) and decide whether
the population is growing, stable, or declining.

\textbf{Solution:}

\textbf{Step 1. One-step projection.}
\(\mathbf{n}(1) = \mathbf{L}\mathbf{n}(0)\) -
\(n_0(1) = 0 \cdot 40 + 1.2 \cdot 20 + 3.0 \cdot 10 = 54\) -
\(n_1(1) = 0.4 \cdot 40 + 0 + 0 = 16\) -
\(n_2(1) = 0 + 0.6 \cdot 20 + 0.7 \cdot 10 = 19\) - Total: \(N(1) = 89\)
vs.~\(N(0) = 70\) --- apparent growth, but transient.

\textbf{Step 2. Characteristic polynomial.} The asymptotic rate
\(\lambda_1\) solves
\(\det(\mathbf{L} - \lambda \mathbf{I}) = 0 \;\Rightarrow\; -\lambda^3 + 0.7\lambda^2 + 0.48\lambda + 0.72 = 0\)

Numerical root-finding gives \(\lambda_1 \approx 1.094\).

\textbf{Step 3. Interpretation.} Because \(\lambda_1 > 1\), the
population grows asymptotically at \(\sim 9.4\%\) per time step once it
converges to its stable age distribution. The transient ratio
\(89/70 = 1.27\) overstates long-run growth because the initial vector
is juvenile-heavy.

\textbf{Answer:} \(\lambda_1 \approx 1.09\); the population is
\textbf{growing} in the long run.

\begin{center}\rule{0.5\linewidth}{0.5pt}\end{center}

\subsubsection{Worked Example 2 --- Mark-Recapture
Estimation}\label{worked-example-2-mark-recapture-estimation}

\textbf{Problem:} A population ecologist uses the mark-recapture method
to estimate the size of a grasshopper population in a meadow. On the
first day, she captures, marks, and releases 120 grasshoppers (\(M\)).
On the second day, she captures a total of 150 grasshoppers (\(C\)), of
which 30 are marked (\(R\)). Calculate the estimated population size
(\(N\)).

\textbf{Solution:}

\textbf{Step 1. Identify the given variables.} - Number marked in first
sample (\(M\)) = 120 - Total number in second sample (\(C\)) = 150 -
Number of marked recaptures (\(R\)) = 30

\textbf{Step 2. Use the Lincoln-Petersen estimator.} The formula for the
total population size (\(N\)) is:
\[N = \frac{M \times C}{R} \label{eq:unit_X_population_ecology_item_2}\]

\textbf{Step 3. Calculate \(N\).}
\[N = \frac{120 \times 150}{30} \label{eq:unit_X_population_ecology_item_3}\]

\[N = \frac{18,000}{30} \label{eq:unit_X_population_ecology_item_4}\]

\[N = 600 \label{eq:unit_X_population_ecology_item_5}\]

\textbf{Answer:} The estimated size of the grasshopper population is
\textbf{600 individuals}.

\begin{center}\rule{0.5\linewidth}{0.5pt}\end{center}

\subsection{Current Evidence and Frontier
Biology}\label{current-evidence-and-frontier-biology-34}

For \textbf{Population Ecology and Growth Models}, frontier biology
belongs inside the evidence logic of the chapter. Ecology and
conservation decisions increasingly combine field data, remote sensing,
community knowledge, model uncertainty, and explicit values. The core
reading question is this: population claims require density dependence,
demographic stochasticity, dispersal, age structure, and management
objective.

\begin{itemize}
\tightlist
\item
  \textbf{What to verify:} identify the observation, model, assay, or
  dataset that would make the claim stronger or weaker.
\item
  \textbf{What to qualify:} state the scale, organism, cell type,
  environmental condition, or population where the claim is expected to
  hold.
\item
  \textbf{What to compare:} test at least one alternative explanation,
  baseline, or null model before treating the pattern as causal.
\item
  \textbf{What to cite:} distinguish primary evidence, review synthesis,
  public dataset, and institutional guidance; for recent or numeric
  claims, prefer the source closest to the measurement and state what
  has changed since it was published.
\end{itemize}

Use biodiversity metrics carefully: population indices, extinction risk
categories, ecosystem services, and management targets answer different
questions
\citep{ipbes2019global,ipbes2024transformative,wwf2024livingplanet,iucn2025redlist,fao2024sofia}.

\textbf{Source practice:} For conservation claims, cite assessment
sources and state whether the evidence is a population index,
extinction-risk assessment, ecosystem-service valuation, satellite
product, or policy synthesis
\citep{ipbes2024transformative,noaa2025coralbleaching,fao2025sofi}.

\subsection{Key Terms}\label{key-terms-34}

{\def\LTcaptype{none} % do not increment counter
\begin{longtable}[]{@{}
  >{\raggedright\arraybackslash}p{(\linewidth - 2\tabcolsep) * \real{0.2857}}
  >{\raggedright\arraybackslash}p{(\linewidth - 2\tabcolsep) * \real{0.7143}}@{}}
\toprule\noalign{}
\begin{minipage}[b]{\linewidth}\raggedright
Term
\end{minipage} & \begin{minipage}[b]{\linewidth}\raggedright
Definition
\end{minipage} \\
\midrule\noalign{}
\endhead
\bottomrule\noalign{}
\endlastfoot
\textbf{Carrying capacity (\(K\))} & Maximum population size supported
by a given environment \\
\textbf{Intrinsic rate of increase (\(r\))} & Per capita growth rate in
ideal conditions (\(= b - d\)) \\
\textbf{Net reproductive rate (\(R_0\))} & Mean lifetime offspring per
individual; \(R_0 > 1\) means growing \\
\textbf{Allee effect} & Positive density-dependence; fitness decreases
at low \(N\) \\
\textbf{Allee threshold (\(A\))} & \(N\) below which population declines
deterministically (strong Allee effect) \\
\textbf{Competition coefficient (\texorpdfstring{\ensuremath{\alpha}}{alpha})} & Per-capita interspecific
competition effect \\
\textbf{Competitive exclusion} & One species eliminates another when
niches overlap completely \\
\textbf{Maximum sustainable yield (MSY)} & \(rK/4\); optimal harvest at
\(N = K/2\) \\
\textbf{\hyperref[gl:lotka-volterra-equations]{Lotka-Volterra
equations}} & ODEs describing competition and predator-prey dynamics \\
\textbf{Bet-hedging} & Variable reproductive strategy reducing variance
in fitness \\
\textbf{Functional response} & Per-capita predation rate as a function
of prey density (Types I, II, III) \\
\textbf{Metapopulation} & Network of subpopulations connected by
dispersal; local extinction balanced by recolonisation \\
\textbf{Source-sink dynamics} & Source patches (\(\lambda > 1\)) sustain
sink patches (\(\lambda < 1\)) via emigration \\
\textbf{Mark-recapture} & Method for estimating population size using
marked individuals \\
\textbf{Reproductive value (\(v_x\))} & Expected future contribution of
age-\(x\) individual to population growth \\
\textbf{Demographic transition} & Shift from high birth/death rates to
low birth/death rates during development \\
\textbf{Theta-logistic model} & Extension of logistic growth with
flexible density-dependence exponent \texorpdfstring{\ensuremath{\theta}}{theta} \\
\textbf{Survivorship curve} & Plot of proportion surviving (\(l_x\))
vs.~age; Types I (late loss), II (constant), III (early loss) \\
\textbf{Leslie matrix} & Age-structured projection matrix; dominant
eigenvalue \(= \lambda_1\) \\
\textbf{Sensitivity / Elasticity} &
\(\partial \lambda_1/\partial a_{ij}\) (absolute) and proportional
analogue; identify highest-leverage rate \\
\textbf{Stable age distribution} & Right eigenvector of Leslie matrix;
long-run age proportions \\
\textbf{Individual-based model (IBM)} & Stochastic simulation tracking
each organism; captures demographic stochasticity and heterogeneity \\
\textbf{Population viability analysis (PVA)} & Stochastic projection of
extinction probability over a fixed horizon \\
\textbf{Quasi-extinction threshold} & \(N_{\text{crit}}\) below which a
population is treated as functionally extinct in PVA \\
\end{longtable}
}

\begin{center}\rule{0.5\linewidth}{0.5pt}\end{center}

\subsection{Review Questions}\label{review-questions-34}

\begin{enumerate}
\def\labelenumi{\arabic{enumi}.}
\item
  A whale population has \(R_0 = 0.95\) and \(T = 20\) years. Estimate
  \(r\). Is this population growing or declining? Calculate how many
  years until the population halves. What conservation interventions
  could change the trajectory?
\item
  Explain why maximum sustainable yield \(\text{MSY} = rK/4\). A
  fisheries manager observes that the current catch equals MSY but the
  population continues to decline. Identify at least three reasons why
  this might occur despite apparently sustainable harvest levels.
\item
  A large predator population is reintroduced into a prey-depleted
  ecosystem with 40 prey individuals remaining. The predator's Allee
  threshold is \(A = 100\). Quantitatively predict whether the predator
  population will recover. What two conditions must simultaneously be
  met for recovery?
\item
  Two grass species compete: \(\alpha_{12} = 0.8\),
  \(\alpha_{21} = 0.7\), \(K_1 = 1000\), \(K_2 = 900\). Use phase-plane
  analysis to determine whether stable coexistence, competitive
  exclusion, or priority effect is predicted. Show your work by
  comparing isocline intercepts.
\item
  Compare the lynx-hare cycle with a simple Lotka-Volterra prediction.
  What additional factors (beyond \texorpdfstring{\ensuremath{\alpha}}{alpha}, \texorpdfstring{\ensuremath{\beta}}{beta}, \texorpdfstring{\ensuremath{\delta}}{delta}, \texorpdfstring{\ensuremath{\gamma}}{gamma} parameters) are needed to
  explain the observed 10-year periodicity? Cite at least one primary
  study.
\item
  A metapopulation of a rare butterfly occupies 25 of 100 available
  habitat patches. The colonisation rate is \(c = 0.15\) and extinction
  rate is \(e = 0.05\). (a) Calculate the expected equilibrium
  occupancy. (b) A highway project will destroy 30 patches. Will the
  metapopulation persist? (c) What colonisation rate would be needed to
  maintain the same equilibrium occupancy after the loss, with 70
  patches remaining?
\item
  You capture 200 salamanders from a pond, mark them, and release them.
  Two weeks later, you capture 250 salamanders and find that 40 are
  marked. (a) Estimate the population size using the Lincoln-Petersen
  method. (b) List three assumptions of this method and describe how
  violating each would bias your estimate.
\item
  Using the concept of reproductive value, explain why protecting adult
  female sea turtles (through Turtle Excluder Devices in fishing nets)
  is more effective for population recovery than protecting eggs on
  nesting beaches.
\item
  Japan has a current population of 125 million and
  \(r \approx -0.006\)/yr. Assuming constant \(r\), project the
  population in 2050, 2075, and 2100. What demographic, economic, and
  social challenges does population decline create?
\item
  The St.~Matthew Island reindeer population grew from 29 to 6,000 in 19
  years, then crashed to 42. (a) Calculate \(r\) during the growth
  phase. (b) What caused the crash in ecological terms? (c) How does
  this case study illustrate the limitations of exponential growth
  models?
\item
  From \texttt{logistic\_growth}, read off \(N(t)\) near \(K/2\) and
  relate to MSY management targets.
\item
  When does the Levins metapopulation model \(\hat p = 1 - e/c\) predict
  \textbf{rescue} vs.~extinction?
\item
  A 4-stage Leslie matrix yields elasticities \(e_{F_3} = 0.08\) (adult
  fecundity), \(e_{P_1} = 0.32\) (juvenile\texorpdfstring{\ensuremath{\to}}{->}subadult survival),
  \(e_{P_2} = 0.34\) (subadult\texorpdfstring{\ensuremath{\to}}{->}adult survival), \(e_{P_3} = 0.26\)
  (adult survival). A wildlife agency proposes (i) doubling clutch size
  via supplemental feeding or (ii) a 5\% absolute increase in subadult
  survival. Which delivers more growth, and what does the elasticity
  comparison imply about the limits of headcount-centered conservation?
\item
  Explain why an individual-based model and a deterministic Leslie
  matrix model can give the same \(\lambda_1\) yet completely different
  extinction predictions for a small population. Under what conditions
  do the two approaches converge?
\item
  Sketch the workflow of a PVA for a critically endangered amphibian:
  (a) what data must be collected, (b) how is the quasi-extinction
  threshold \(N_{\text{crit}}\) chosen, (c) why are \emph{relative}
  comparisons across management scenarios more defensible than absolute
  extinction probabilities?
\end{enumerate}

\begin{center}\rule{0.5\linewidth}{0.5pt}\end{center}

\subsection{Further Reading and Source
Notes}\label{further-reading-and-source-notes-29}

\begin{itemize}
\tightlist
\item
  Lotka (1925). \emph{Elements of Physical Biology}. Williams \&
  Wilkins.
\item
  Volterra (1926). Variazioni e fluttuazioni del numero d'individui in
  specie animali conviventi. \emph{Memorie dell'Accademia dei Lincei},
  2.
\item
  Verhulst (1838). Notice sur la loi que la population poursuit dans son
  accroissement. \emph{Correspondance Math\{\textquotesingle e\}matique
  et Physique}, 10.
\item
  Allee (1931). \emph{Animal Aggregations: A Study in General
  Sociology}. University of Chicago Press.
\item
  Holling (1959). Some characteristics of simple types of predation and
  parasitism. \emph{Canadian Entomologist}, 91.
\item
  MacArthur \& Wilson (1967). \emph{The Theory of
  \hyperref[gl:island-biogeography]{\textbf{Island biogeography}}}.
  Princeton University Press.
\end{itemize}

\begin{center}\rule{0.5\linewidth}{0.5pt}\end{center}

\subsection{Computational Bridge}\label{computational-bridge-31}

Logistic time series are integrated deterministically:

\begin{Shaded}
\begin{Highlighting}[]
\ImportTok{from}\NormalTok{ biology.ecology }\ImportTok{import}\NormalTok{ logistic\_growth}

\NormalTok{ser }\OperatorTok{=}\NormalTok{ logistic\_growth(N0}\OperatorTok{=}\FloatTok{40.0}\NormalTok{, r}\OperatorTok{=}\FloatTok{0.2}\NormalTok{, K}\OperatorTok{=}\FloatTok{500.0}\NormalTok{, t\_end}\OperatorTok{=}\FloatTok{60.0}\NormalTok{, steps}\OperatorTok{=}\DecValTok{240}\NormalTok{)}
\BuiltInTok{print}\NormalTok{(}\BuiltInTok{round}\NormalTok{(ser.populations[}\OperatorTok{{-}}\DecValTok{1}\NormalTok{], }\DecValTok{2}\NormalTok{))}
\end{Highlighting}
\end{Shaded}

\begin{quote}
\textbf{Clinical / systems note:} Harvest models for fisheries and
forest stands reuse the same \(r\), \(K\), and MSY logic; collapse
usually means \(F > F_\mathrm{MSY}\) for too long.
\end{quote}

\begin{center}\rule{0.5\linewidth}{0.5pt}\end{center}

\subsection{Summary}\label{summary-34}

\begin{itemize}
\tightlist
\item
  \textbf{Demographics:} \(r = b - d\); \(R_0 = \sum l_x m_x\); doubling
  time \(= \ln 2/r\); generation time
  \(T = \sum x \cdot l_x m_x / R_0\). Euler-Lotka equation provides
  exact relationship.
\item
  \textbf{Dispersion patterns:} clumped (most common), uniform, random
  --- detected by variance-to-mean ratio.
\item
  \textbf{Survivorship curves:} Type I (late loss, K-selected), Type II
  (constant), Type III (early loss, r-selected). Gompertz law for human
  mortality.
\item
  \textbf{Logistic growth:} \(dN/dt = rN(1 - N/K)\); inflection at
  \(K/2\); MSY \(= rK/4\). Extensions: theta-logistic, time-lagged
  logistic.
\item
  \textbf{Allee effect:} \(dN/dt < 0\) below threshold \(A\); threatens
  small surviving populations. Component vs.~demographic Allee effects.
\item
  \textbf{LV competition:} stable coexistence requires niche
  differentiation (\(\alpha_{12}\alpha_{21} < 1\)); otherwise exclusion
  or priority effects. Chesson's theory: stabilising + equalising
  mechanisms.
\item
  \textbf{LV predator-prey:} neutral oscillations; period
  \(\approx 2\pi/\sqrt{\alpha\gamma}\); lynx-hare 10-yr cycle requires
  tri-trophic explanation.
\item
  \textbf{Functional responses:} Type I (linear), Type II (decelerating,
  Holling disc equation), Type III (sigmoidal, stabilising).
\item
  \textbf{Age-structured models:} Leslie matrix \(\mathbf{L}\) with
  fecundities on top row and survivorships on sub-diagonal; dominant
  eigenvalue \(\lambda_1\) = asymptotic growth rate; right and left
  eigenvectors give stable age distribution and reproductive value.
  \textbf{Sensitivity} \(s_{ij} = \partial \lambda_1/\partial a_{ij}\)
  and \textbf{elasticity} \(e_{ij}\) identify highest-leverage
  demographic rates (loggerhead turtle adult survival \texorpdfstring{\ensuremath{\to}}{->} TEDs).
\item
  \textbf{Individual-based models (IBMs):} stochastic agent-based
  simulations capture demographic stochasticity, individual
  heterogeneity, and emergent dynamics that mean-field models miss.
\item
  \textbf{Population viability analysis (PVA):} integrates matrix and
  stochastic machinery into extinction-probability forecasts;
  quasi-extinction threshold \(N_{\text{crit}}\); empirical MVPs cluster
  at 5,000--7,000 adults (Reed et al.~2003), well above the classical
  \(N_e = 500\) rule.
\item
  \textbf{r vs.~K continuum:} life-history trade-offs shaped by habitat
  stability; superseded by Stearns' and Grime's frameworks; bet-hedging
  in variable environments.
\item
  \textbf{Metapopulations:} Levins model (\(\hat{p} = 1 - e/c\));
  source-sink dynamics; rescue effect. Habitat corridors increase
  connectivity.
\item
  \textbf{Population estimation:} mark-recapture (Lincoln-Petersen,
  Jolly-Seber), distance sampling.
\item
  \textbf{Human demographics:} 8 billion (2022); demographic transition;
  population momentum; TFR declining globally.
\item
  \textbf{Connections:} See \cref{sec:unit_X_community_ecology} for
  species interactions, \cref{sec:unit_V_population_genetics} for
  \hyperref[gl:allele]{\textbf{allele}} frequencies in subdivided
  populations, and Unit VI for life-history evolution.
\end{itemize}

\begin{center}\rule{0.5\linewidth}{0.5pt}\end{center}

\subsubsection{Companion Source
Module}\label{companion-source-module-34}

\textbf{Population Ecology and Growth Models} should leave a
reproducible trail from a biological claim to the code, figure, diagram,
or paper-based activity that can test it. Use the surfaces below to
inspect the chapter's assumptions, rerun the relevant model, or compare
the manuscript explanation with companion labs and figures.

{\def\LTcaptype{none} % do not increment counter
\begin{longtable}[]{@{}
  >{\raggedright\arraybackslash}p{(\linewidth - 2\tabcolsep) * \real{0.5000}}
  >{\raggedright\arraybackslash}p{(\linewidth - 2\tabcolsep) * \real{0.5000}}@{}}
\toprule\noalign{}
\begin{minipage}[b]{\linewidth}\raggedright
Surface
\end{minipage} & \begin{minipage}[b]{\linewidth}\raggedright
Use it for
\end{minipage} \\
\midrule\noalign{}
\endhead
\bottomrule\noalign{}
\endlastfoot
\breaktt{src/biology/ecology/ecology.py} (\breaktt{exponential\_growth},
\texttt{logistic\_growth}, \breaktt{allee\_strong\_growth}) & Reproduce
density-independent, density-dependent, and Allee-effect scenarios. \\
\breaktt{src/visualization/plots.py} (\breaktt{plot\_logistic\_growth}) &
Inspect carrying capacity and growth-rate assumptions. \\
\breaktt{src/mermaid/biology\_diagrams.py}
(\breaktt{population\_growth\_stages\_diagram}) & Link model phases to
visual summaries. \\
\end{longtable}
}

\textbf{Reproducibility check:} state time step, units, density
dependence, stochasticity, and management objective before forecasting a
population. \textbf{Cross-reference:} compare with
\cref{sec:unit_V_population_genetics} and
\cref{sec:unit_X_community_ecology}.

\newpage

\section{Community Ecology and Species
Interactions}\label{community-ecology-and-species-interactions}

\label{sec:unit_X_community_ecology}

\begin{quote}
\textbf{Ch 33} \texorpdfstring{\ensuremath{\cdot}}{.} Level 2/3 \texorpdfstring{\ensuremath{\cdot}}{.} 80 min read \texorpdfstring{\ensuremath{\cdot}}{.} 100 min lecture \texorpdfstring{\ensuremath{\cdot}}{.}
Prerequisites: \cref{sec:unit_X_population_ecology}
\end{quote}

\subsection{Learning Objectives}\label{learning-objectives-35}

By the end of this chapter, you should be able to:

\begin{enumerate}
\def\labelenumi{\arabic{enumi}.}
\tightlist
\item
  Define a community and categorise the six types of biotic interactions
  with examples.
\item
  Apply Lotka-Volterra competition equations to predict competitive
  outcomes and explain the
  \hyperref[gl:competitive-exclusion]{\textbf{competitive exclusion}}
  principle and \hyperref[gl:niche]{\textbf{niche}} theory.
\item
  Explain \hyperref[gl:trophic-cascade]{\textbf{trophic cascade}}s and
  keystone \citep{paine1966} species with quantitative examples.
\item
  Compare primary and
  \hyperref[gl:secondary-succession]{\textbf{secondary succession}} and
  explain the intermediate disturbance \citep{connell1978} hypothesis.
\item
  Calculate Shannon diversity (H'), Simpson index, and species evenness
  and explain what each measures.
\item
  Apply \hyperref[gl:island-biogeography]{\textbf{island biogeography}}
  theory, the
  \hyperref[gl:species-area-relationship]{\textbf{species-area
  relationship}}, and SLOSS debate to conservation design.
\item
  Describe the role of disturbance, facilitation, and alternative stable
  states in shaping communities.
\item
  Explain food web topology, connectance, network robustness, and the
  relationship between complexity and stability.
\item
  Compare niche-based and neutral theories of biodiversity, and explain
  when each null model is appropriate.
\item
  Use functional traits (CSR strategies, leaf economics spectrum) to
  predict community assembly and ecosystem function beyond species
  lists.
\item
  Distinguish classical, augmentative, and conservation biological
  control, and evaluate their risk profiles using community ecology
  principles.
\end{enumerate}

\subsubsection{Study Blueprint}\label{study-blueprint-35}

\begin{itemize}
\tightlist
\item
  \textbf{Big idea:} Communities are structured by interactions among
  species and by the context of those interactions.
\item
  \textbf{Core concepts:} competition, predation, mutualism, succession.
\item
  \textbf{Framework alignment:} Vision \& Change: Systems, Evolution,
  Pathways and transformations of energy and matter; AP Biology: Systems
  Interactions, Evolution, Energetics; NGSS-style topics: Interdependent
  Relationships in Ecosystems, Matter and Energy in Organisms and
  Ecosystems, Natural Selection and Evolution.
\item
  \textbf{Model or quantitative lens:} Lotka-Volterra-style interaction
  and diversity-pattern reasoning.
\item
  \textbf{Data skill:} Interpret abundance, interaction, or disturbance
  data from communities.
\item
  \textbf{Practice cadence:} Representing and Describing Data,
  Statistical Tests and Data Analysis, Argumentation.
\item
  \textbf{Common misconception to repair:} A species interaction is not
  permanently good or bad; the sign can change with context.
\item
  \textbf{Primary lab:} \cref{sec:lab_unit_X_community_ecology}.
\item
  \textbf{Question bank:} \cref{sec:q_unit_X_community_ecology}.
\item
  \textbf{Transfer task:} Transfer interaction reasoning to restoration,
  disease ecology, agriculture, and invasion biology.
\item
  \textbf{Bridge to computation:}
  \breaktt{biology.ecology.ecology.biodiversity\_indices}.
\end{itemize}

\begin{center}\rule{0.5\linewidth}{0.5pt}\end{center}

\begin{quote}
\textbf{Opening Vignette --- How Wolves Changed Rivers}

In 1995, 14 gray wolves were reintroduced to Yellowstone National Park
after a 70-year absence. What happened next has become one of ecology's
most vivid demonstrations of a trophic cascade. With predators back, elk
avoided grazing in valleys and riverside areas where they were
vulnerable. Vegetation in those areas --- willows, aspens, cottonwoods
--- rebounded within years. With trees stabilising river banks, erosion
slowed. Beaver colonies, dependent on willows, increased sixfold. Beaver
dams created wetlands that supported fish, otters, ducks, and
amphibians. River channels narrowed and meandered, becoming more
complex. The wolves, through fear alone --- the ``landscape of fear''
effect documented by William Ripple --- had changed the physical
geography of the park. The Yellowstone study has been cited thousands of
times, popularised by a viral YouTube video with 40 million views, and
debated (some effects took decades to show). But it remains the
canonical example that removing or restoring apex predators cascades
through every \hyperref[gl:trophic-level]{\textbf{trophic level}} of a
community.
\end{quote}

\subsubsection{Chapter Roadmap}\label{chapter-roadmap-13}

This is a long chapter that covers eight closely-related but
distinguishable topics. Read it as \emph{two halves}:

\begin{itemize}
\tightlist
\item
  \textbf{Part A --- Local interactions between species:} what happens
  when two or a few species meet. The chapter begins with community
  definitions, then develops competition, predation and trophic
  cascades, \hyperref[gl:mutualism]{\textbf{mutualism}}, and parasitism.
\item
  \textbf{Part B --- Community-scale patterns and assembly:} what
  emerges when you scale up. Succession, diversity measurement, island
  biogeography, and food-web network structure become the organising
  themes.
\item
  \textbf{Part C --- Process-level theory and applications:} neutral
  theory supplies a null model for biodiversity; trait-based ecology
  (CSR, leaf economics) predicts ecosystem function; biological control
  applies the entire chapter to pests and disease vectors.
\end{itemize}

If you are reading for a one-semester survey, Part A supplies the
mechanistic vocabulary and Part B supplies the integrative patterns.
Part C provides the modern frameworks and applied translations.
Instructors wanting to split the chapter over two lectures can use the
Part-A/Part-B boundary; a third lecture can cover Part C.

\subsection{Community Ecology
Fundamentals}\label{community-ecology-fundamentals}

An \textbf{ecological community} is an assemblage of populations of
different species occupying the same region and time, interacting with
each other and their \hyperref[gl:abiotic]{\textbf{abiotic}}
environment. Community ecology analyses the \textbf{biotic interactions}
among species and their effects on community structure --- species
composition and relative abundance --- as well as the processes driving
community assembly and succession.

\subsubsection{Emergent Properties of
Communities}\label{emergent-properties-of-communities}

Communities possess properties that cannot be predicted from studying
individual species in isolation:

{\def\LTcaptype{none} % do not increment counter
\begin{longtable}[]{@{}
  >{\raggedright\arraybackslash}p{(\linewidth - 4\tabcolsep) * \real{0.2759}}
  >{\raggedright\arraybackslash}p{(\linewidth - 4\tabcolsep) * \real{0.3448}}
  >{\raggedright\arraybackslash}p{(\linewidth - 4\tabcolsep) * \real{0.3793}}@{}}
\toprule\noalign{}
\begin{minipage}[b]{\linewidth}\raggedright
Property
\end{minipage} & \begin{minipage}[b]{\linewidth}\raggedright
Definition
\end{minipage} & \begin{minipage}[b]{\linewidth}\raggedright
Measurement
\end{minipage} \\
\midrule\noalign{}
\endhead
\bottomrule\noalign{}
\endlastfoot
\textbf{Species richness (\(S\))} & Number of species present & Count of
unique species \\
\textbf{Species evenness} & How equal are species abundances? & Pielou's
\(J'\) \\
\textbf{Diversity} & Combined richness and evenness & Shannon \(H'\),
Simpson \(1-D\) \\
\textbf{Community structure} & Food web topology, trophic levels,
functional groups & Network analysis, connectance \\
\textbf{Resistance} & Ability to withstand perturbation without change &
Deviation from baseline after disturbance \\
\textbf{Resilience} & Speed of return to baseline after perturbation &
Recovery time \\
\textbf{Rank-abundance distribution} & Pattern of relative abundance
across species & Log-normal, geometric series, broken stick \\
\end{longtable}
}

\subsubsection{Rank-Abundance Models}\label{rank-abundance-models}

The distribution of individuals among species in a community follows
predictable patterns:

{\def\LTcaptype{none} % do not increment counter
\begin{longtable}[]{@{}
  >{\raggedright\arraybackslash}p{(\linewidth - 6\tabcolsep) * \real{0.1163}}
  >{\raggedright\arraybackslash}p{(\linewidth - 6\tabcolsep) * \real{0.1628}}
  >{\raggedright\arraybackslash}p{(\linewidth - 6\tabcolsep) * \real{0.3256}}
  >{\raggedright\arraybackslash}p{(\linewidth - 6\tabcolsep) * \real{0.3953}}@{}}
\toprule\noalign{}
\begin{minipage}[b]{\linewidth}\raggedright
Model
\end{minipage} & \begin{minipage}[b]{\linewidth}\raggedright
Pattern
\end{minipage} & \begin{minipage}[b]{\linewidth}\raggedright
Interpretation
\end{minipage} & \begin{minipage}[b]{\linewidth}\raggedright
Typical community
\end{minipage} \\
\midrule\noalign{}
\endhead
\bottomrule\noalign{}
\endlastfoot
\textbf{Geometric series} & Steep, convex & Strong dominance; niche
pre-emption & Species-poor, harsh environments \\
\textbf{Log-series} & Moderate slope & Many rare species, few common &
Island faunas, successional communities \\
\textbf{Log-normal} & Moderate, symmetric bell curve on log scale & Most
natural communities & Large, undisturbed communities \\
\textbf{Broken stick} & Shallow, even & Resources divided equally &
Species-poor, saturated communities \\
\end{longtable}
}

\textbf{Preston's canonical log-normal} (1962): In most large
communities, when species are binned by abundance in octaves
(doublings), the distribution is approximately log-normal. This has deep
connections to the species-area relationship.

\subsubsection{Types of Biotic
Interactions}\label{types-of-biotic-interactions}

\begin{figure}[htbp]
\centering
\includegraphics[width=0.82\linewidth,height=4.2in,keepaspectratio]{../figures/mermaid_inline/inline_mermaid_0179_75b8c647214a.png}
\caption{Flowchart for Types of Biotic Interactions: Mutualism (+/+) Both benefit Mycorrhizae, fig-wasp, cleaner fish, coral-zooxanthellae, Parasitism (+/) Parasite benefits, host harmed Tapeworm, mistletoe, Toxoplasma, malaria, Predation (+/) Predator benefits, prey killed Wolf-moose, owl-mouse, spider-fly, and Competition (/) Both harmed Warblers in same tree, plants competing for light form the diagram's primary path or branches.}
\end{figure}

{\def\LTcaptype{none} % do not increment counter
\begin{longtable}[]{@{}
  >{\raggedright\arraybackslash}p{(\linewidth - 8\tabcolsep) * \real{0.1930}}
  >{\raggedright\arraybackslash}p{(\linewidth - 8\tabcolsep) * \real{0.2632}}
  >{\raggedright\arraybackslash}p{(\linewidth - 8\tabcolsep) * \real{0.2632}}
  >{\raggedright\arraybackslash}p{(\linewidth - 8\tabcolsep) * \real{0.1579}}
  >{\raggedright\arraybackslash}p{(\linewidth - 8\tabcolsep) * \real{0.1228}}@{}}
\toprule\noalign{}
\begin{minipage}[b]{\linewidth}\raggedright
Interaction
\end{minipage} & \begin{minipage}[b]{\linewidth}\raggedright
Species A effect
\end{minipage} & \begin{minipage}[b]{\linewidth}\raggedright
Species B effect
\end{minipage} & \begin{minipage}[b]{\linewidth}\raggedright
Mechanism
\end{minipage} & \begin{minipage}[b]{\linewidth}\raggedright
Example
\end{minipage} \\
\midrule\noalign{}
\endhead
\bottomrule\noalign{}
\endlastfoot
Mutualism (+/+) & Benefits & Benefits & Direct reciprocal benefit &
Mycorrhizae (+plant, +fungus); fig-wasp pollination; cleaner fish \\
Commensalism (+/0) & Benefits & Neutral & One benefits, other unaffected
& Cattle egret + cattle (cattle disturb insects); barnacles on whale \\
Parasitism (+/-) & Benefits (\hyperref[gl:parasite]{\textbf{parasite}})
& Harmed (host) & Partial exploitation; usually not lethal & Tapeworm +
human; mistletoe + tree; \emph{Toxoplasma} + rodent \\
Predation (+/-) & Benefits (predator) & Harmed (prey) & Prey consumed;
drives prey adaptation & Wolf + moose; \emph{Daphnia} +
\hyperref[gl:phytoplankton]{\textbf{phytoplankton}} \\
Competition (-/-) & Harmed & Harmed & Shared resource demand & Two
warbler species in same niche; plants competing for light \\
Amensalism (0/-) & Neutral & Harmed & Inhibitory compounds; physical
suppression & Juglone from black walnut;
\hyperref[gl:biofilm]{\textbf{biofilm}} quorum quenching \\
\end{longtable}
}

\begin{quote}
\textbf{Concept Check:} A remora fish attaches to a shark, feeding on
scraps from the shark's meals. Is this mutualism, commensalism, or
parasitism? What additional information would you need to determine the
exact interaction type?
\end{quote}

\begin{center}\rule{0.5\linewidth}{0.5pt}\end{center}

\subsection{Competition Theory}\label{competition-theory}

\subsubsection{Lotka-Volterra Interspecific
Competition}\label{lotka-volterra-interspecific-competition}

Two competing species \(N_1\) and \(N_2\) with shared resources are
modelled by:

\begin{equation}
\frac{dN_1}{dt} = r_1 N_1 \left(1 - \frac{N_1 + \alpha_{12} N_2}{K_1}\right)
\label{eq:community_ecology_1}
\end{equation}

\begin{equation}
\frac{dN_2}{dt} = r_2 N_2 \left(1 - \frac{N_2 + \alpha_{21} N_1}{K_2}\right)
\label{eq:community_ecology_2}
\end{equation}

where: - \(\alpha_{12}\) = competitive effect of species 2 on species 1
(per individual) - \(\alpha_{21}\) = competitive effect of species 1 on
species 2 - \(K_1\), \(K_2\) = carrying capacities of each species

\textbf{Isocline analysis:} Four outcomes depending on isocline
intersections: 1. \textbf{Species 1 wins:} \(K_1 > K_2/\alpha_{12}\) AND
\(K_1/\alpha_{21} > K_2\) 2. \textbf{Species 2 wins:}
\(K_2 > K_1/\alpha_{21}\) AND \(K_2/\alpha_{12} > K_1\) (symmetric case)
3. \textbf{Unstable equilibrium (priority effect):}
\(K_1 > K_2/\alpha_{12}\) AND \(K_2 > K_1/\alpha_{21}\) --- whoever
starts higher wins 4. \textbf{Stable coexistence:}
\(K_1 < K_2/\alpha_{12}\) AND \(K_2 < K_1/\alpha_{21}\) ---
interspecific competition weaker than intraspecific

\subsubsection{Competitive Exclusion Principle (Gause's
Law)}\label{competitive-exclusion-principle-gauses-law}

\textbf{\citet{gause1934}:} Two species competing for the same limiting
resource cannot stably coexist; the superior competitor excludes the
other. Demonstrated with \emph{Paramecium aurelia} vs.~\emph{P.
caudatum} grown together on a single bacterial food source --- \emph{P.
aurelia} consistently drove \emph{P. caudatum} to extinction within 20
days.

\textbf{Important qualifications:} Competitive exclusion requires (1)
complete niche overlap, (2) constant environment, (3) no spatial
refugia, and (4) sufficient time. In nature, these conditions are rarely
fully met.

\subsubsection{Hutchinson's Niche
Concept}\label{hutchinsons-niche-concept}

\textbf{Fundamental niche} \citep{hutchinson1957}: the n-dimensional
hypervolume of environmental conditions and resources permitting a
species to maintain \(r \geq 0\).

\textbf{Realised niche:} the subset of the fundamental niche actually
occupied after accounting for interspecific competition, predation, and
other biotic interactions. Typically smaller than or equal to the
fundamental niche.

\begin{equation}
\text{Realised niche} = \text{Fundamental niche} - \text{Competitive exclusion zone}
\label{eq:community_ecology_3}
\end{equation}

\textbf{Hutchinson's paradox of the plankton} (1961): Why do hundreds of
phytoplankton species coexist in the apparently homogeneous water
column, most competing for the same light and nutrients? Violates
competitive exclusion. Solutions:

\begin{enumerate}
\def\labelenumi{\arabic{enumi}.}
\tightlist
\item
  \textbf{Temporal niche partitioning} --- seasonal turnover,
  disturbance prevents equilibrium
\item
  \textbf{Spatial heterogeneity} --- microscale gradients in light,
  nutrients, turbulence
\item
  \textbf{Predation (selective grazing)} --- zooplankton preferentially
  graze \hyperref[gl:dominant]{\textbf{dominant}} species
\item
  \textbf{Allelopathy} --- chemical warfare among phytoplankton
\item
  \textbf{Non-equilibrium dynamics} --- the system rarely reaches
  competitive exclusion
\end{enumerate}

\subsubsection{\texorpdfstring{Modern Coexistence Theory
\citep{chesson2000}}{Modern Coexistence Theory}}\label{modern-coexistence-theory}

Coexistence requires that \textbf{intraspecific competition
\textgreater{} interspecific competition} for both species. This is
enabled by two classes of mechanisms:

\textbf{Stabilising mechanisms} (niche differences): - \textbf{Resource
partitioning:} different foods, microhabitats, activity times -
\textbf{Janzen-Connell effect:} species-specific enemies concentrate
near conspecifics, giving heterospecifics an advantage - \textbf{Storage
effect:} temporal environmental variation favours different species at
different times; long-lived adults ``store'' good years -
\textbf{Relative nonlinearity:} species respond differently to resource
fluctuations

\textbf{Equalising mechanisms} (fitness similarity): - Trade-offs
between competitive ability and other traits (dispersal, stress
tolerance) - Similar per-capita growth rates at low density

The mathematical condition for coexistence is:

\begin{equation}
\rho < \frac{\kappa_j}{\kappa_i} < \frac{1}{\rho}
\label{eq:community_ecology_4}
\end{equation}

where \texorpdfstring{\ensuremath{\rho}}{rho} = niche overlap (0 to 1) and \(\kappa_i/\kappa_j\) = fitness
ratio. Lower \texorpdfstring{\ensuremath{\rho}}{rho} (more niche differentiation) allows greater fitness
inequality.

\begin{Shaded}
\begin{Highlighting}[]
\ImportTok{from}\NormalTok{ biology.ecology.ecology }\ImportTok{import}\NormalTok{ lotka\_volterra}

\NormalTok{result }\OperatorTok{=}\NormalTok{ lotka\_volterra(}
\NormalTok{    N1\_0}\OperatorTok{=}\DecValTok{100}\NormalTok{, N2\_0}\OperatorTok{=}\DecValTok{100}\NormalTok{,}
\NormalTok{    r1}\OperatorTok{=}\FloatTok{0.5}\NormalTok{, r2}\OperatorTok{=}\FloatTok{0.5}\NormalTok{,}
\NormalTok{    K1}\OperatorTok{=}\DecValTok{500}\NormalTok{, K2}\OperatorTok{=}\DecValTok{500}\NormalTok{,}
\NormalTok{    alpha12}\OperatorTok{=}\FloatTok{0.6}\NormalTok{,   }\CommentTok{\# species 2 slightly suppresses species 1}
\NormalTok{    alpha21}\OperatorTok{=}\FloatTok{0.8}\NormalTok{,   }\CommentTok{\# species 1 strongly suppresses species 2}
\NormalTok{    t\_end}\OperatorTok{=}\DecValTok{100}
\NormalTok{)}
\BuiltInTok{print}\NormalTok{(}\SpecialStringTok{f"Final N1: }\SpecialCharTok{\{}\NormalTok{result}\SpecialCharTok{.}\NormalTok{N1[}\OperatorTok{{-}}\DecValTok{1}\NormalTok{]}\SpecialCharTok{:.0f\}}\SpecialStringTok{, N2: }\SpecialCharTok{\{}\NormalTok{result}\SpecialCharTok{.}\NormalTok{N2[}\OperatorTok{{-}}\DecValTok{1}\NormalTok{]}\SpecialCharTok{:.0f\}}\SpecialStringTok{"}\NormalTok{)}
\CommentTok{\# Expected: species 1 wins (lower alpha12 means less impact from competitor)}
\end{Highlighting}
\end{Shaded}

\begin{quote}
\texorpdfstring{}{} \textbf{Clinical Connection --- Competitive Exclusion in the
\hyperref[gl:microbiome]{\textbf{Microbiome}}:} The competitive
exclusion principle operates within the human gut microbiome.
\textbf{Clostridium difficile} infection (CDI) typically occurs after
antibiotic treatment eliminates normal gut flora, removing competitive
exclusion and allowing \emph{C. difficile} to proliferate unopposed.
\textbf{Fecal \hyperref[gl:microbiota]{\textbf{microbiota}}
transplantation (FMT)} restores competitive exclusion by reintroducing a
diverse microbial community, achieving about 90\% cure rates
for recurrent CDI. This is Gause's principle applied to clinical
medicine.
\end{quote}

\begin{quote}
\textbf{Concept Check:} Two species of Paramecium (\emph{P. aurelia} and
\emph{P. bursaria}) coexist in the same pond. \emph{P. aurelia} feeds on
bacteria in open water; \emph{P. bursaria} harbours symbiotic algae and
feeds near the bottom. Explain their coexistence using Chesson's
framework: what is the stabilising mechanism?
\end{quote}

\begin{center}\rule{0.5\linewidth}{0.5pt}\end{center}

\subsection{Predation, Keystone Species, and Trophic
Cascades}\label{predation-keystone-species-and-trophic-cascades}

\subsubsection{Predator-Prey Arms Races}\label{predator-prey-arms-races}

\hyperref[gl:coevolution]{\textbf{Coevolution}} \citep{ehrlich1964}
between predators and prey drives a Red Queen dynamic of escalating
adaptations:

\textbf{Anti-predator defences:}

{\def\LTcaptype{none} % do not increment counter
\begin{longtable}[]{@{}
  >{\raggedright\arraybackslash}p{(\linewidth - 4\tabcolsep) * \real{0.3333}}
  >{\raggedright\arraybackslash}p{(\linewidth - 4\tabcolsep) * \real{0.3750}}
  >{\raggedright\arraybackslash}p{(\linewidth - 4\tabcolsep) * \real{0.2917}}@{}}
\toprule\noalign{}
\begin{minipage}[b]{\linewidth}\raggedright
Strategy
\end{minipage} & \begin{minipage}[b]{\linewidth}\raggedright
Mechanism
\end{minipage} & \begin{minipage}[b]{\linewidth}\raggedright
Example
\end{minipage} \\
\midrule\noalign{}
\endhead
\bottomrule\noalign{}
\endlastfoot
\textbf{Crypsis} & Match background appearance & \emph{Biston betularia}
(peppered moth) --- industrial melanism \\
\textbf{Aposematism} & Bright warning colours advertise toxicity &
Poison dart frogs (\emph{Dendrobates}) --- alkaloid warning \\
\textbf{Mullerian mimicry} & Two toxic species converge on same warning
signal & \emph{Heliconius} butterflies sharing wing patterns \\
\textbf{Batesian mimicry} & Palatable species mimics toxic model &
Viceroy (\emph{Limenitis archippus}) mimics monarch \\
\textbf{Chemical defence} & Toxic compounds deter predators & Monarch
butterflies sequester cardenolide glycosides \\
\textbf{Startle display} & Sudden reveal of eyespots or bright colours &
Io moth (\emph{Automeris io}) eyespot flash \\
\textbf{Behavioural} & Alarm calls, mobbing, confusion effect & Starling
murmurations confuse raptors \\
\textbf{Morphological} & Spines, shells, armour & Porcupine quills;
turtle shells; hedgehog spines \\
\end{longtable}
}

\textbf{Counter-adaptations in predators:} - \textbf{Kingsnakes}
(\emph{Lampropeltis}) evolved resistance to rattlesnake venom -
\textbf{Rough-skinned newt} (\emph{Taricha granulosa})
vs.~\textbf{garter snake} (\emph{Thamnophis sirtalis}): escalating toxin
(tetrodotoxin) and resistance --- one of the best-documented
coevolutionary arms races \citep{brodie1999} - \textbf{Cuckoo}
(\emph{Cuculus canorus}) egg mimicry vs.~host egg discrimination ---
classic parasitic arms race

\subsubsection{Hare-Lynx Cycle}\label{hare-lynx-cycle}

\emph{Lepus americanus} (snowshoe hare) and \emph{Lynx canadensis} show
coupled about 10-year population cycles in boreal Canada.
Hudson's Bay Company fur trading records from 1736-1940 reveal the
cycle. Modern analysis:

\begin{equation}
\frac{dH}{dt} = r_{max}H - aPH \qquad \frac{dP}{dt} = b \cdot a \cdot P \cdot H - dP
\label{eq:community_ecology_5}
\end{equation}

The cycle is driven by multiple feedback loops: 1. \textbf{Predation}
(about 60\% of cycle) --- lynx consumption drives hare decline
2. \textbf{Food quality} (about 20\%) --- hare overgrazing
induces phenolic toughening of willow/birch browse 3. \textbf{Stress
physiology} (about 20\%) --- predation risk elevates
\hyperref[gl:cortisol]{\textbf{cortisol}}, suppressing reproduction 4.
True cycles require \textbf{tri-trophic} feedbacks; purely
Lotka-Volterra predation alone is insufficient (Krebs et al.~2001,
\emph{Science})

\subsubsection{Keystone Species}\label{keystone-species}

A \textbf{\hyperref[gl:keystone-species]{keystone species}} (Paine 1966,
1969) has disproportionately large effects on community structure
relative to its abundance. Removal of the keystone causes dramatic
community reorganisation.

\textbf{Types of keystone effects:}

{\def\LTcaptype{none} % do not increment counter
\begin{longtable}[]{@{}
  >{\raggedright\arraybackslash}p{(\linewidth - 4\tabcolsep) * \real{0.2000}}
  >{\raggedright\arraybackslash}p{(\linewidth - 4\tabcolsep) * \real{0.4500}}
  >{\raggedright\arraybackslash}p{(\linewidth - 4\tabcolsep) * \real{0.3500}}@{}}
\toprule\noalign{}
\begin{minipage}[b]{\linewidth}\raggedright
Type
\end{minipage} & \begin{minipage}[b]{\linewidth}\raggedright
Mechanism
\end{minipage} & \begin{minipage}[b]{\linewidth}\raggedright
Example
\end{minipage} \\
\midrule\noalign{}
\endhead
\bottomrule\noalign{}
\endlastfoot
\textbf{Keystone predator} & Prevents competitive exclusion by
suppressing dominant prey & \emph{Pisaster ochraceus} (sea star) \texorpdfstring{\ensuremath{\to}}{->}
prevents mussel monopoly \\
\textbf{Keystone herbivore} & Controls dominant plant, maintaining
diversity & African elephant \texorpdfstring{\ensuremath{\to}}{->} prevents woodland encroachment on
savanna \\
\textbf{Keystone mutualist} & Supports many other species through
interaction network & Fig trees \texorpdfstring{\ensuremath{\to}}{->} fruit for 1,200+ vertebrate species in
tropical forests \\
\textbf{Ecosystem engineer} & Physically modifies habitat & Beaver dams
create wetland habitat for diverse communities \\
\end{longtable}
}

\subsubsection{Trophic Cascades}\label{trophic-cascades}

A \textbf{trophic cascade} occurs when changes at one trophic level
ripple through the food web to affect non-adjacent levels:

\textbf{Top-down cascade} (predator-driven): - Sea otter \texorpdfstring{\ensuremath{\to}}{->} sea urchin \texorpdfstring{\ensuremath{\to}}{->}
kelp cascade \citep{estes1974} - Wolf \texorpdfstring{\ensuremath{\to}}{->} elk \texorpdfstring{\ensuremath{\to}}{->} willow cascade in
Yellowstone \citep{ripple2012}

\textbf{Bottom-up cascade} (resource-driven): - Nutrient enrichment \texorpdfstring{\ensuremath{\to}}{->}
phytoplankton \texorpdfstring{\ensuremath{\to}}{->} zooplankton \texorpdfstring{\ensuremath{\to}}{->} fish

\begin{figure}[htbp]
\centering
\includegraphics[width=0.82\linewidth,height=4.2in,keepaspectratio]{../figures/mermaid_inline/inline_mermaid_0180_18b43723222a.png}
\caption{Sequence diagram for Trophic Cascades showing ordered interaction among Wolves (reintroduced 1995), Elk, Willows/Cottonwoods, and Beavers.}
\end{figure}

\textbf{Quantitative trophic cascade --- Sea Otter example:} - 1 sea
otter consumes about 10 kg sea urchin/day - Each urchin removed
protects about 1 m\(^2\) of kelp - Kelp productivity:
about 300 g C/m\(^2\)/yr - Therefore: 1 otter \texorpdfstring{\ensuremath{\to}}{->}
about 3,650 m\(^2\) kelp protected \texorpdfstring{\ensuremath{\to}}{->} about 1,095 kg
C/yr sequestered - Kelp forests also attenuate wave energy by 60-70\%,
providing coastal protection valued at about \$10,000/km/yr

\textbf{Mesopredator release:} removal of apex predators \texorpdfstring{\ensuremath{\to}}{->} explosion of
medium-sized predators \texorpdfstring{\ensuremath{\to}}{->} disproportionate prey decline. Example: coyote
increase after wolf removal \texorpdfstring{\ensuremath{\to}}{->} ground-nesting bird decline in North
America.

\begin{quote}
\texorpdfstring{}{} \textbf{Clinical Connection --- Trophic Cascades and Lyme Disease:}
The decline of apex predators in eastern North America has contributed
to a trophic cascade affecting human health. Wolf and cougar removal \texorpdfstring{\ensuremath{\to}}{->}
deer population explosion \texorpdfstring{\ensuremath{\to}}{->} increased tick-deer encounters \texorpdfstring{\ensuremath{\to}}{->} increased
\emph{Borrelia burgdorferi} transmission \texorpdfstring{\ensuremath{\to}}{->} Lyme disease incidence
increased 25-fold from 1990 to 2020 in the northeastern USA. Deer also
browse forest understory, reducing small mammal habitat diversity, which
concentrates ticks on the most competent reservoir hosts (white-footed
mice), further amplifying transmission. Predator restoration could
reduce disease burden --- a health-relevant trophic cascade.
\end{quote}

\begin{quote}
\textbf{Concept Check:} In Yellowstone, wolves primarily cause a
``behavioural trophic cascade'' rather than a purely numerical one. What
is the difference? How does the ``landscape of fear'' concept explain
why elk behaviour change (avoiding riverbanks) may be more important
than elk population reduction?
\end{quote}

\begin{center}\rule{0.5\linewidth}{0.5pt}\end{center}

\subsection{Mutualism, Parasitism, and
Facilitation}\label{mutualism-parasitism-and-facilitation}

\subsubsection{Mutualism}\label{mutualism}

\textbf{Obligate mutualism:} neither partner can survive without the
other (e.g., fig-wasp pollination; lichens = fungus +
alga/cyanobacterium).

\textbf{Facultative mutualism:} both benefit but can survive
independently (e.g., seed-dispersing birds and fruiting trees).

\textbf{Mycorrhizal networks:} - \textbf{Arbuscular mycorrhizae (AM,
Glomeromycota):} obligate symbionts; hyphae form arbuscules inside root
cortex cells; deliver P and micronutrients; receive up to 30\% of plant
photosynthate; present in about 80\% of land plant species -
\textbf{Ectomycorrhizae (EM, Basidiomycota + Ascomycota):} hyphae form
Hartig net around root cells; dominant in temperate/boreal forests
(\emph{Pinus}, \emph{Betula}, \emph{Fagus}); provide N via proteolytic
\hyperref[gl:enzyme]{\textbf{enzyme}}s - \textbf{Common mycorrhizal
networks (CMNs) --- ``Wood Wide Web'':} CMNs transfer carbon, water, and
nutrient signals between plants of same and different species; large
``mother trees'' supply carbon to understory seedlings (Simard et
al.~1997, \emph{Nature}). Contested: whether transfer is truly adaptive
signalling or passive diffusion remains debated (Karst et al.~2023,
\emph{Nature Ecology \& Evolution})

\subsubsection{Parasitism and Disease
Ecology}\label{parasitism-and-disease-ecology}

Parasites regulate host populations and can function as keystone
species:

{\def\LTcaptype{none} % do not increment counter
\begin{longtable}[]{@{}
  >{\raggedright\arraybackslash}p{(\linewidth - 4\tabcolsep) * \real{0.3514}}
  >{\raggedright\arraybackslash}p{(\linewidth - 4\tabcolsep) * \real{0.1892}}
  >{\raggedright\arraybackslash}p{(\linewidth - 4\tabcolsep) * \real{0.4595}}@{}}
\toprule\noalign{}
\begin{minipage}[b]{\linewidth}\raggedright
Parasite type
\end{minipage} & \begin{minipage}[b]{\linewidth}\raggedright
Example
\end{minipage} & \begin{minipage}[b]{\linewidth}\raggedright
Ecological effect
\end{minipage} \\
\midrule\noalign{}
\endhead
\bottomrule\noalign{}
\endlastfoot
\textbf{Macroparasite} & Intestinal helminths, ticks, lice & Reduce host
fitness; can regulate population size \\
\textbf{Microparasite} & \emph{Batrachochytrium dendrobatidis} (Bd) &
Chytrid fungus: \textgreater90 amphibian species extinctions since
1970 \\
\textbf{Parasitoid} & \emph{Cotesia glomerata} (braconid wasp) & Lays
eggs inside caterpillars; larvae consume host \\
\textbf{Social parasite} & Cuckoo (\emph{Cuculus canorus}) & Brood
parasitism; reduces host reproductive output \\
\textbf{Manipulative parasite} & \emph{Toxoplasma gondii} & Alters
rodent behaviour (decreased fear of cats) to facilitate transmission \\
\end{longtable}
}

\textbf{The parasite-mediated competition hypothesis:} parasites can
determine competitive outcomes between host species, effectively
functioning as ``hidden'' keystone species.

\subsubsection{Facilitation}\label{facilitation}

\textbf{Facilitation} occurs when one species improves the survival or
reproduction of another. Unlike mutualism, facilitation can be
unidirectional:

\begin{itemize}
\tightlist
\item
  \textbf{Nurse plants} in deserts: cacti establish under the shade of
  nurse shrubs (\emph{Larrea}, \emph{Ambrosia})
\item
  \textbf{Foundation species:} create habitat structure (e.g., corals,
  kelp, \emph{Spartina} grass in salt marshes)
\item
  \textbf{Nitrogen fixers:} \emph{Lupinus} colonises volcanic
  substrates, enriching soil N for later successional species (Mount
  St.~Helens \hyperref[gl:primary-succession]{\textbf{primary
  succession}})
\end{itemize}

\begin{quote}
\textbf{Concept Check:} The relationship between clownfish and sea
anemones is often described as mutualism. The clownfish gains protection
from predators; the anemone may benefit from clownfish defending against
butterfly fish that eat anemone tentacles. How would you design an
experiment to determine whether this interaction is truly mutualism
(+/+) or commensalism (+/0)?
\end{quote}

\begin{center}\rule{0.5\linewidth}{0.5pt}\end{center}

\subsection{Ecological Succession}\label{ecological-succession}

\textbf{Succession} = directional, generally predictable change in
community composition over time after disturbance or on newly available
substrate.

Contemporary succession ecology is less deterministic than the classical
``march to climax'' story. Recovery trajectories depend on surviving
legacies, seed banks, dispersal corridors, disturbance severity,
herbivory, invasive species, soil microbes, and climate conditions
during recovery. The practical question is not only which stage comes
next, but which intervention would change the trajectory: protecting
refuges, adding propagules, removing barriers, or accepting a novel
stable state.

{\def\LTcaptype{none} % do not increment counter
\begin{longtable}[]{@{}
  >{\raggedright\arraybackslash}p{(\linewidth - 4\tabcolsep) * \real{0.1591}}
  >{\raggedright\arraybackslash}p{(\linewidth - 4\tabcolsep) * \real{0.3864}}
  >{\raggedright\arraybackslash}p{(\linewidth - 4\tabcolsep) * \real{0.4545}}@{}}
\toprule\noalign{}
\begin{minipage}[b]{\linewidth}\raggedright
Feature
\end{minipage} & \begin{minipage}[b]{\linewidth}\raggedright
Primary succession
\end{minipage} & \begin{minipage}[b]{\linewidth}\raggedright
Secondary succession
\end{minipage} \\
\midrule\noalign{}
\endhead
\bottomrule\noalign{}
\endlastfoot
Starting substrate & Bare rock/sterile substrate (no soil) & Disturbed
community with soil intact \\
Typical rate & Centuries to millennia & Decades to centuries \\
Pioneer species & Cyanobacteria, lichens, mosses & Annual weeds,
grasses \\
Soil development & Must form from scratch (weathering + organic
accumulation) & Already present; seed bank may survive \\
Example & Krakatoa (1883); Surtsey, Iceland (1963); Mount St.~Helens
(1980) & Old-field succession (abandoned farmland); post-fire forest
regrowth \\
\end{longtable}
}

\subsubsection{Mechanisms of Succession}\label{mechanisms-of-succession}

\textbf{\citet{connell1977}} proposed three mechanisms:

{\def\LTcaptype{none} % do not increment counter
\begin{longtable}[]{@{}
  >{\raggedright\arraybackslash}p{(\linewidth - 4\tabcolsep) * \real{0.2381}}
  >{\raggedright\arraybackslash}p{(\linewidth - 4\tabcolsep) * \real{0.4286}}
  >{\raggedright\arraybackslash}p{(\linewidth - 4\tabcolsep) * \real{0.3333}}@{}}
\toprule\noalign{}
\begin{minipage}[b]{\linewidth}\raggedright
Model
\end{minipage} & \begin{minipage}[b]{\linewidth}\raggedright
Mechanism
\end{minipage} & \begin{minipage}[b]{\linewidth}\raggedright
Example
\end{minipage} \\
\midrule\noalign{}
\endhead
\bottomrule\noalign{}
\endlastfoot
\textbf{Facilitation} & Early species modify environment to favour later
species & Nitrogen-fixing \emph{Alnus} (alder) enriches soil, enabling
spruce colonisation (Glacier Bay, Alaska) \\
\textbf{Tolerance} & Later species can establish regardless of early
species but grow more slowly & Shade-tolerant species slowly replace
shade-intolerant pioneers \\
\textbf{Inhibition} & Early colonists resist replacement; succession
proceeds primarily when pioneers die & \emph{Cladonia} lichen crusts
inhibit vascular plant establishment on sand dunes \\
\end{longtable}
}

\begin{figure}[htbp]
\centering
\includegraphics[width=0.82\linewidth,height=4.2in,keepaspectratio]{../figures/mermaid_inline/inline_mermaid_0181_5b60733bc96e.png}
\caption{Ecological succession pathways. Primary succession begins on bare substrate and takes centuries; secondary succession begins with intact soil and proceeds in decades. Three mechanisms (facilitation, tolerance, inhibition) operate at transition points. Modern ecology views the endpoint as a dynamic equilibrium rather than a fixed climax state.}
\end{figure}

In reality, most successional sequences involve most three mechanisms
operating simultaneously at different stages and spatial scales.

\subsubsection{Climax Community Concept}\label{climax-community-concept}

\textbf{Classical view \citep{clements1916}:} Succession proceeds toward
a single, deterministic climax community determined by regional climate
(the ``organismal'' model --- community as superorganism).

\textbf{Modern view (Gleason 1926; Whittaker 1953):} Communities are
assemblages of individually distributed species (the ``individualistic''
model). Multiple stable endpoints possible; disturbance history,
stochastic colonisation, and priority effects most influence trajectory.
The concept of a single climax is largely abandoned.

\subsubsection{Intermediate Disturbance Hypothesis
(IDH)}\label{intermediate-disturbance-hypothesis-idh}

\textbf{Connell (1978):} Communities at intermediate levels of
disturbance frequency and intensity show maximum species diversity:

\begin{itemize}
\tightlist
\item
  \textbf{Low disturbance} \texorpdfstring{\ensuremath{\to}}{->} competitive exclusion \texorpdfstring{\ensuremath{\to}}{->} low diversity
  (dominant competitor wins)
\item
  \textbf{High disturbance} \texorpdfstring{\ensuremath{\to}}{->} primarily r-selected pioneers survive \texorpdfstring{\ensuremath{\to}}{->}
  low diversity
\item
  \textbf{Intermediate} \texorpdfstring{\ensuremath{\to}}{->} prevents competitive dominance while enabling
  diverse colonisation
\end{itemize}

\begin{equation}
H' = f(\text{disturbance frequency, intensity})
\label{eq:community_ecology_6}
\end{equation}

\textbf{Evidence:} Coral reefs, tropical forests, stream invertebrate
communities --- moderate disturbance (hurricanes, floods, fires)
increases diversity.

\textbf{Criticisms:} The IDH has been challenged as overly simplistic
(Fox 2013, \emph{Ecology}). Some communities show monotonic
diversity-disturbance relationships. The hypothesis also assumes a
competition-colonisation trade-off that is not comprehensive.

\subsubsection{Alternative Stable States and Regime
Shifts}\label{alternative-stable-states-and-regime-shifts}

\textbf{Alternative stable states} (Lewontin 1969; Scheffer et
al.~2001): Some ecosystems can exist in multiple stable configurations
under the same environmental conditions. Transitions between states
(\textbf{regime shifts}) can be triggered by small perturbations near
\textbf{tipping points}:

{\def\LTcaptype{none} % do not increment counter
\begin{longtable}[]{@{}
  >{\raggedright\arraybackslash}p{(\linewidth - 6\tabcolsep) * \real{0.1622}}
  >{\raggedright\arraybackslash}p{(\linewidth - 6\tabcolsep) * \real{0.1892}}
  >{\raggedright\arraybackslash}p{(\linewidth - 6\tabcolsep) * \real{0.1892}}
  >{\raggedright\arraybackslash}p{(\linewidth - 6\tabcolsep) * \real{0.4595}}@{}}
\toprule\noalign{}
\begin{minipage}[b]{\linewidth}\raggedright
System
\end{minipage} & \begin{minipage}[b]{\linewidth}\raggedright
State 1
\end{minipage} & \begin{minipage}[b]{\linewidth}\raggedright
State 2
\end{minipage} & \begin{minipage}[b]{\linewidth}\raggedright
Tipping mechanism
\end{minipage} \\
\midrule\noalign{}
\endhead
\bottomrule\noalign{}
\endlastfoot
Shallow lake & Clear water (macrophytes dominant) & Turbid (algal bloom)
& Nutrient loading exceeds P threshold \\
Coral reef & Coral dominated & Macroalgae dominated & Overfishing of
herbivorous fish \\
Savanna & Grassland with scattered trees & Closed-canopy forest & Fire
suppression \\
Arctic & Perennial sea ice & Seasonal/ice-free & Ocean warming
threshold \\
\end{longtable}
}

\textbf{Early warning signals} of approaching tipping points: -
Increased temporal variance (flickering) - \textbf{Critical slowing
down} --- recovery from perturbation becomes slower - Increased spatial
correlation - Skewed distribution of system state

These early warning indicators are being applied to monitor reef, lake,
and climate system stability.

\begin{quote}
\texorpdfstring{}{} \textbf{Clinical Connection --- Microbiome Regime Shifts:} The human
gut microbiome exhibits alternative stable states analogous to
ecological regime shifts. A healthy, diverse microbiome represents one
stable state. Antibiotic treatment or \emph{C. difficile} infection can
trigger a regime shift to an impoverished, pathogenic state. Once
established, the unhealthy state resists return to the diverse state
(hysteresis) --- explaining why probiotics alone are often insufficient
and why fecal microbiota transplantation (which provides a complete
microbial community) is more effective.
\end{quote}

\begin{quote}
\textbf{Concept Check:} After the 1980 eruption of Mount St.~Helens,
ecologists observed primary succession proceeding much faster than
predicted. Surviving pocket gophers, ants, and lupine plants in patches
of surviving soil created nuclei of rapid recovery. Which successional
mechanism(s) does this illustrate, and why does it challenge Clements'
classical model?
\end{quote}

\begin{center}\rule{0.5\linewidth}{0.5pt}\end{center}

\subsection{Measuring Biodiversity}\label{measuring-biodiversity}

\subsubsection{Alpha, Beta, and Gamma
Diversity}\label{alpha-beta-and-gamma-diversity}

\textbf{\citet{whittaker1960}} distinguished three scales of diversity:

These scales also clarify what modern biodiversity tools can and cannot
show. Environmental DNA, acoustic monitoring, camera traps, remote
sensing, and citizen-science records can reveal turnover across space
faster than classical plots alone, but each method has detection bias,
taxonomic gaps, and scale limits. A credible diversity comparison states
the sampling unit, detection method, taxonomic resolution, and whether
the result is richness, evenness, composition, or functional change.

{\def\LTcaptype{none} % do not increment counter
\begin{longtable}[]{@{}
  >{\raggedright\arraybackslash}p{(\linewidth - 4\tabcolsep) * \real{0.2381}}
  >{\raggedright\arraybackslash}p{(\linewidth - 4\tabcolsep) * \real{0.4762}}
  >{\raggedright\arraybackslash}p{(\linewidth - 4\tabcolsep) * \real{0.2857}}@{}}
\toprule\noalign{}
\begin{minipage}[b]{\linewidth}\raggedright
Scale
\end{minipage} & \begin{minipage}[b]{\linewidth}\raggedright
Definition
\end{minipage} & \begin{minipage}[b]{\linewidth}\raggedright
Metric
\end{minipage} \\
\midrule\noalign{}
\endhead
\bottomrule\noalign{}
\endlastfoot
\textbf{Alpha (\texorpdfstring{\ensuremath{\alpha}}{alpha}) diversity} & Species richness within a single
community/habitat & \(H'\), Simpson's, species count \\
\textbf{Beta (\texorpdfstring{\ensuremath{\beta}}{beta}) diversity} & Turnover in species composition between
communities & Jaccard index, Bray-Curtis dissimilarity \\
\textbf{Gamma (\texorpdfstring{\ensuremath{\gamma}}{gamma}) diversity} & Total diversity across most communities
in a landscape & \(\gamma = \bar{\alpha} \times \beta\)
(multiplicative) \\
\end{longtable}
}

\subsubsection{\texorpdfstring{Shannon-Wiener Diversity Index
(\(H'\))}{Shannon-Wiener Diversity Index (Htextquotesingle)}}\label{shannon-wiener-diversity-index-h}

\begin{equation}
H' = -\sum_{i=1}^{S} p_i \ln p_i
\label{eq:community_ecology_7}
\end{equation}

where \(p_i\) = proportion of individuals belonging to species \(i\).

\begin{itemize}
\tightlist
\item
  Range: \(H' = 0\) (monoculture) to \(H' = \ln(S)\) (perfectly even)
\item
  Sensitive to rare species
\item
  Typical values: 1.5-3.5 for most communities; \textgreater4.0 for
  extremely diverse tropical communities
\end{itemize}

\subsubsection{\texorpdfstring{Simpson Diversity Index
(\(1-D\))}{Simpson Diversity Index (1-D)}}\label{simpson-diversity-index-1-d}

\begin{equation}
D = \sum_{i=1}^{S} p_i^2 \qquad \text{Simpson's diversity} = 1 - D
\label{eq:community_ecology_8}
\end{equation}

\begin{itemize}
\tightlist
\item
  Range: 0-1; probability that two randomly selected individuals are
  from \textbf{different} species
\item
  Weighted toward dominant species (less affected by rare species than
  \(H'\))
\item
  Also expressed as reciprocal: Simpson's reciprocal \(= 1/D\)
  (effective number of species)
\end{itemize}

\subsubsection{\texorpdfstring{Pielou's Evenness
(\(J'\))}{Pielou's Evenness (Jtextquotesingle)}}\label{pielous-evenness-j}

\begin{equation}
J' = \frac{H'}{\ln S} \in [0, 1]
\label{eq:community_ecology_9}
\end{equation}

\(J' = 1\) means most species equally abundant; \(J' \to 0\) means
dominated by one species.

\subsubsection{Worked Example}\label{worked-example-5}

{\def\LTcaptype{none} % do not increment counter
\begin{longtable}[]{@{}lllll@{}}
\toprule\noalign{}
Species & Abundance & \(p_i\) & \(p_i \ln p_i\) & \(p_i^2\) \\
\midrule\noalign{}
\endhead
\bottomrule\noalign{}
\endlastfoot
A & 50 & 0.500 & -0.347 & 0.250 \\
B & 30 & 0.300 & -0.361 & 0.090 \\
C & 15 & 0.150 & -0.285 & 0.023 \\
D & 4 & 0.040 & -0.129 & 0.002 \\
E & 1 & 0.010 & -0.046 & 0.000 \\
\textbf{Total} & \textbf{100} & \textbf{1.000} & \textbf{-1.167} &
\textbf{0.364} \\
\end{longtable}
}

\(H' = 1.167\); \(H'_{max} = \ln 5 = 1.609\);
\(J' = 1.167/1.609 = 0.725\)

Simpson's \(D = 0.364\); Simpson's diversity \(= 1 - 0.364 = 0.636\)

\subsubsection{Beta Diversity: Measuring
Turnover}\label{beta-diversity-measuring-turnover}

\textbf{Jaccard similarity index:}

\begin{equation}
J = \frac{|A \cap B|}{|A \cup B|}
\label{eq:community_ecology_10}
\end{equation}

where \(A\) and \(B\) are species sets from two communities. \(J = 1\)
means identical composition; \(J = 0\) means no shared species.

\textbf{Bray-Curtis dissimilarity} (incorporates abundance):

\begin{equation}
BC = 1 - \frac{2 \sum_i \min(a_i, b_i)}{\sum_i (a_i + b_i)}
\label{eq:community_ecology_11}
\end{equation}

\begin{Shaded}
\begin{Highlighting}[]
\ImportTok{from}\NormalTok{ biology.ecology.ecology }\ImportTok{import}\NormalTok{ biodiversity\_indices}

\CommentTok{\# Community counts (array of species abundances)}
\NormalTok{counts }\OperatorTok{=}\NormalTok{ [}\DecValTok{50}\NormalTok{, }\DecValTok{30}\NormalTok{, }\DecValTok{15}\NormalTok{, }\DecValTok{4}\NormalTok{, }\DecValTok{1}\NormalTok{]    }\CommentTok{\# 5 species; very uneven}
\NormalTok{result }\OperatorTok{=}\NormalTok{ biodiversity\_indices(counts)}
\BuiltInTok{print}\NormalTok{(}\SpecialStringTok{f"Species richness: }\SpecialCharTok{\{}\NormalTok{result}\SpecialCharTok{.}\NormalTok{species\_richness}\SpecialCharTok{\}}\SpecialStringTok{"}\NormalTok{)}
\BuiltInTok{print}\NormalTok{(}\SpecialStringTok{f"Shannon H\textquotesingle{}: }\SpecialCharTok{\{}\NormalTok{result}\SpecialCharTok{.}\NormalTok{shannon\_index}\SpecialCharTok{:.3f\}}\SpecialStringTok{"}\NormalTok{)}
\BuiltInTok{print}\NormalTok{(}\SpecialStringTok{f"Simpson 1{-}D: }\SpecialCharTok{\{}\NormalTok{result}\SpecialCharTok{.}\NormalTok{simpson\_index}\SpecialCharTok{:.3f\}}\SpecialStringTok{"}\NormalTok{)}
\BuiltInTok{print}\NormalTok{(}\SpecialStringTok{f"Evenness J\textquotesingle{}: }\SpecialCharTok{\{}\NormalTok{result}\SpecialCharTok{.}\NormalTok{evenness}\SpecialCharTok{:.3f\}}\SpecialStringTok{"}\NormalTok{)}
\end{Highlighting}
\end{Shaded}

\begin{quote}
\textbf{Concept Check:} Community X has 4 species with abundances
{[}100, 100, 100, 100{]}. Community Y has 4 species with abundances
{[}394, 2, 2, 2{]}. Both have \(S = 4\). Calculate \(H'\) and \(J'\) for
each. Which community is more ``diverse'' and why does evenness matter
as much as richness?
\end{quote}

\begin{center}\rule{0.5\linewidth}{0.5pt}\end{center}

\subsection{Island Biogeography
Theory}\label{island-biogeography-theory}

\subsubsection{MacArthur-Wilson Equilibrium
Model}\label{macarthur-wilson-equilibrium-model}

\textbf{\citet{macarthur1967}:} \emph{The Theory of Island Biogeography}
--- equilibrium species richness on islands balances immigration
(colonisation) rate and extinction rate:

\begin{itemize}
\tightlist
\item
  \textbf{Immigration rate} decreases as \(S\) increases (fewer
  uncolonised species remain in the mainland pool)
\item
  \textbf{Extinction rate} increases as \(S\) increases (more species
  present means more potential extinctions)
\item
  At \textbf{equilibrium} (\(\hat{S}\)): immigration rate = extinction
  rate
\end{itemize}

\textbf{Effects of island characteristics:}

{\def\LTcaptype{none} % do not increment counter
\begin{longtable}[]{@{}
  >{\raggedright\arraybackslash}p{(\linewidth - 4\tabcolsep) * \real{0.1765}}
  >{\raggedright\arraybackslash}p{(\linewidth - 4\tabcolsep) * \real{0.5588}}
  >{\raggedright\arraybackslash}p{(\linewidth - 4\tabcolsep) * \real{0.2647}}@{}}
\toprule\noalign{}
\begin{minipage}[b]{\linewidth}\raggedright
Factor
\end{minipage} & \begin{minipage}[b]{\linewidth}\raggedright
Effect on \(\hat{S}\)
\end{minipage} & \begin{minipage}[b]{\linewidth}\raggedright
Mechanism
\end{minipage} \\
\midrule\noalign{}
\endhead
\bottomrule\noalign{}
\endlastfoot
\textbf{Larger area} & Higher \(\hat{S}\) & Lower extinction rate
(larger populations, more habitats) \\
\textbf{Closer to mainland} & Higher \(\hat{S}\) & Higher immigration
rate (easier colonisation) \\
\textbf{Small + far} & Lowest \(\hat{S}\) & High extinction, low
immigration \\
\textbf{Large + near} & Highest \(\hat{S}\) & Low extinction, high
immigration \\
\end{longtable}
}

\begin{center}\rule{0.5\linewidth}{0.5pt}\end{center}

\subsubsection{Visualizing Food Web Connectance and
Modularity}\label{visualizing-food-web-connectance-and-modularity}

\begin{figure}[htbp]
\centering
\includegraphics[width=0.82\linewidth,height=4.2in,keepaspectratio]{../figures/mermaid_inline/inline_mermaid_0182_8a2e02816705.png}
\caption{Flowchart for Visualizing Food Web Connectance and Modularity: High C = L/S2 Robust to single extinctions Rapid perturbation spread and High Modularity Compartmentalised links Contains perturbations Typical of large ecosystems form the diagram's primary path or branches.}
\end{figure}

\begin{center}\rule{0.5\linewidth}{0.5pt}\end{center}

\subsubsection{Species-Area
Relationship}\label{species-area-relationship}

\begin{equation}
S = cA^z
\label{eq:community_ecology_12}
\end{equation}

Log-transformed: \(\log S = \log c + z \log A\)

{\def\LTcaptype{none} % do not increment counter
\begin{longtable}[]{@{}
  >{\raggedright\arraybackslash}p{(\linewidth - 4\tabcolsep) * \real{0.2188}}
  >{\raggedright\arraybackslash}p{(\linewidth - 4\tabcolsep) * \real{0.3438}}
  >{\raggedright\arraybackslash}p{(\linewidth - 4\tabcolsep) * \real{0.4375}}@{}}
\toprule\noalign{}
\begin{minipage}[b]{\linewidth}\raggedright
Context
\end{minipage} & \begin{minipage}[b]{\linewidth}\raggedright
Typical \(z\)
\end{minipage} & \begin{minipage}[b]{\linewidth}\raggedright
Interpretation
\end{minipage} \\
\midrule\noalign{}
\endhead
\bottomrule\noalign{}
\endlastfoot
True islands & 0.20-0.35 & Higher \(z\) because islands are more
isolated \\
Mainland habitat fragments & 0.12-0.17 & Lower \(z\) due to rescue
effect from surrounding matrix \\
Archipelagoes & 0.25-0.33 & Standard island prediction \\
\end{longtable}
}

\textbf{Practical calculation:} Doubling area increases \(S\) by
\(\approx 19\%\) (when \(z = 0.25\), since \(2^{0.25} - 1 = 0.189\)).

\begin{Shaded}
\begin{Highlighting}[]
\ImportTok{import}\NormalTok{ math}

\KeywordTok{def}\NormalTok{ species\_area(c, A, z):}
    \ControlFlowTok{return}\NormalTok{ c }\OperatorTok{*}\NormalTok{ (A }\OperatorTok{**}\NormalTok{ z)}

\CommentTok{\# Old{-}growth forest fragment richness predictions}
\ControlFlowTok{for}\NormalTok{ area\_ha }\KeywordTok{in}\NormalTok{ [}\DecValTok{1}\NormalTok{, }\DecValTok{10}\NormalTok{, }\DecValTok{100}\NormalTok{, }\DecValTok{1000}\NormalTok{, }\DecValTok{10000}\NormalTok{]:}
\NormalTok{    S }\OperatorTok{=}\NormalTok{ species\_area(c}\OperatorTok{=}\DecValTok{5}\NormalTok{, A}\OperatorTok{=}\NormalTok{area\_ha, z}\OperatorTok{=}\FloatTok{0.25}\NormalTok{)}
    \BuiltInTok{print}\NormalTok{(}\SpecialStringTok{f"Area = }\SpecialCharTok{\{}\NormalTok{area\_ha}\SpecialCharTok{:\textgreater{}6\}}\SpecialStringTok{ ha → S = }\SpecialCharTok{\{}\NormalTok{S}\SpecialCharTok{:.1f\}}\SpecialStringTok{ bird species"}\NormalTok{)}
\end{Highlighting}
\end{Shaded}

\subsubsection{Empirical Tests}\label{empirical-tests}

\begin{itemize}
\tightlist
\item
  \textbf{Post-Krakatoa recolonisation} (1883-1983): After complete
  sterilisation by eruption, species accumulated toward equilibrium;
  initial overshoot then relaxation \citep{whittaker1975}
\item
  \textbf{Florida Keys experiment} \citep{simberloff1969}: Fumigated
  small mangrove islands \texorpdfstring{\ensuremath{\to}}{->} recolonisation to predicted equilibrium
  within 2 years; confirmed immigration-extinction dynamics
\item
  \textbf{Habitat fragments as islands:} Biological Dynamics of Forest
  Fragments Project (BDFFP, Amazonia; Laurance et al.~2011) --- 40 years
  of data on isolated 1-100 ha fragments showing predictable species
  loss following species-area relationship
\end{itemize}

\subsubsection{SLOSS Debate}\label{sloss-debate}

\textbf{Single Large Or Several Small} --- for conservation, which
reserve design maximises biodiversity?

{\def\LTcaptype{none} % do not increment counter
\begin{longtable}[]{@{}
  >{\raggedright\arraybackslash}p{(\linewidth - 4\tabcolsep) * \real{0.2069}}
  >{\raggedright\arraybackslash}p{(\linewidth - 4\tabcolsep) * \real{0.3448}}
  >{\raggedright\arraybackslash}p{(\linewidth - 4\tabcolsep) * \real{0.4483}}@{}}
\toprule\noalign{}
\begin{minipage}[b]{\linewidth}\raggedright
Design
\end{minipage} & \begin{minipage}[b]{\linewidth}\raggedright
Advantages
\end{minipage} & \begin{minipage}[b]{\linewidth}\raggedright
Disadvantages
\end{minipage} \\
\midrule\noalign{}
\endhead
\bottomrule\noalign{}
\endlastfoot
\textbf{Single large} & Lower edge-to-interior ratio; larger MVPs;
protects area-demanding species & Vulnerability to single catastrophe;
may miss regional habitat diversity \\
\textbf{Several small} & Redundancy against disaster; covers more
habitat types; protects local endemics & Edge effects dominant; lower
\(S\) per patch; connectivity problems \\
\end{longtable}
}

\textbf{Modern consensus: Core-corridor-matrix design:} - Large core
reserves (the ``single large'' benefit) - Habitat corridors or stepping
stones connecting reserves - Matrix land uses compatible with wildlife
movement (agroforestry, wildlife-friendly farming) - Example
initiatives: Yellowstone-to-Yukon (Y2Y, 3,200 km); Florida Wildlife
Corridor (1,000+ miles)

\subsubsection{Extinction Debt}\label{extinction-debt}

\textbf{Extinction debt} (Tilman et al.~1994, \emph{Nature}): Species
currently present in degraded/fragmented habitats may be committed to
future extinction --- but the extinction is \textbf{delayed} by decades
to centuries as existing individuals slowly die without replacement.

\begin{equation}
\text{Species committed to extinction} = S_{\text{pre-fragmentation}} - c \cdot A_{\text{remaining}}^z
\label{eq:community_ecology_13}
\end{equation}

\textbf{Implications:} - Present-day species richness
\textbf{overestimates} long-term viability - Conservation assessments
based on current presence may be dangerously optimistic - European
calcareous grasslands: 20-50\% of plant species face committed
extinction due to historical fragmentation \citep{lindborg2004} - Amazon
deforestation: 25-year extinction debt means current biodiversity
surveys underestimate eventual losses

\subsubsection{Metapopulation Dynamics (Levins, 1969; Hanski,
1994)}\label{metapopulation-dynamics-levins-1969-hanski-1994}

A patchwork of habitat fragments connected by dispersal --- local
extinctions offset by recolonisation. Metapopulation persistence
requires fraction of occupied patches to exceed extinction threshold:

\begin{equation}
\hat{p} = 1 - \frac{e}{c}
\label{eq:community_ecology_14}
\end{equation}

where \(e\) = extinction rate per patch, \(c\) = colonisation rate;
metapopulation persists if \(e/c < 1\).

\textbf{Hanski's incidence function model} extends this by incorporating
patch area (affects extinction rate) and isolation (affects colonisation
rate):

\begin{equation}
J_i = \frac{1}{1 + (e_i/c_i)^2}
\label{eq:community_ecology_15}
\end{equation}

where \(J_i\) = probability that patch \(i\) is occupied,
\(e_i \propto A_i^{-x}\), and \(c_i\) depends on distance to occupied
patches.

\begin{quote}
\texorpdfstring{}{} \textbf{Clinical Connection --- Habitat Fragmentation and Zoonotic
Disease:} As forests are fragmented into small patches, edge effects
increase human-wildlife contact, facilitating zoonotic pathogen
spillover. The emergence of Nipah virus in Malaysia (1998-99) was linked
to bat habitat loss: displaced fruit bats (\emph{Pteropus} spp.) moved
to pig farms near forest fragments, facilitating transmission to
livestock and then humans. Similarly, Ebola outbreaks correlate
spatially with deforestation frontiers. Island biogeography theory
predicts that small fragments lose apex predators first (mesopredator
release), increasing populations of rodents and bats that serve as
pathogen reservoirs.
\end{quote}

\begin{quote}
\textbf{Concept Check:} A national park of 10,000 km\(^2\) is split by a
highway into two fragments of 6,000 km\(^2\) and 4,000 km\(^2\). Using
\(S = 50A^{0.25}\), calculate species richness for: (a) the intact park,
(b) each fragment separately, (c) both fragments combined assuming no
species overlap. How does a wildlife corridor change your prediction?
\end{quote}

\begin{center}\rule{0.5\linewidth}{0.5pt}\end{center}

\subsection{Food Web Structure and Network
Ecology}\label{food-web-structure-and-network-ecology}

\subsubsection{Food Web Topology}\label{food-web-topology}

Food webs are network representations of feeding relationships in
communities --- a graph where nodes are species (or trophic units) and
directed edges encode ``\(i\) eats \(j\).'' Treating the community as a
graph unlocks the entire toolbox of network science (degree
distributions, clustering, modularity, motif analysis), which has
reshaped community ecology over the past two decades. Key metrics:

{\def\LTcaptype{none} % do not increment counter
\begin{longtable}[]{@{}
  >{\raggedright\arraybackslash}p{(\linewidth - 4\tabcolsep) * \real{0.2069}}
  >{\raggedright\arraybackslash}p{(\linewidth - 4\tabcolsep) * \real{0.3448}}
  >{\raggedright\arraybackslash}p{(\linewidth - 4\tabcolsep) * \real{0.4483}}@{}}
\toprule\noalign{}
\begin{minipage}[b]{\linewidth}\raggedright
Metric
\end{minipage} & \begin{minipage}[b]{\linewidth}\raggedright
Definition
\end{minipage} & \begin{minipage}[b]{\linewidth}\raggedright
Typical value
\end{minipage} \\
\midrule\noalign{}
\endhead
\bottomrule\noalign{}
\endlastfoot
\textbf{Connectance (\(C\))} & \(C = L/S^2\) (or \(L/[S(S-1)/2]\));
fraction of possible links realised & 0.05-0.30 \\
\textbf{Links per species (\(L/S\))} & Average number of trophic links
per species & about 2 for most food webs \\
\textbf{Chain length} & Number of links from base to top & 3-5
typically \\
\textbf{Omnivory} & Fraction of species feeding at \textgreater1 trophic
level & about 50\% in many webs \\
\textbf{Mean trophic level} & Average path length from primary producers
& 2.5--3.5 typical \\
\textbf{Generality} & Mean number of resources per consumer & Skewed:
most consumers specialised \\
\textbf{Vulnerability} & Mean number of consumers per resource & Skewed:
most prey have few predators \\
\textbf{Nestedness} & Specialists' diets are subsets of generalists'
diets & High in mutualistic networks \\
\end{longtable}
}

\begin{Shaded}
\begin{Highlighting}[]
\CommentTok{\# Minimal food{-}web statistics from an adjacency dictionary}
\KeywordTok{def}\NormalTok{ web\_stats(web):}
\NormalTok{    species }\OperatorTok{=} \BuiltInTok{set}\NormalTok{(web).union(\{p }\ControlFlowTok{for}\NormalTok{ prey }\KeywordTok{in}\NormalTok{ web.values() }\ControlFlowTok{for}\NormalTok{ p }\KeywordTok{in}\NormalTok{ prey\})}
\NormalTok{    S }\OperatorTok{=} \BuiltInTok{len}\NormalTok{(species)}
\NormalTok{    L }\OperatorTok{=} \BuiltInTok{sum}\NormalTok{(}\BuiltInTok{len}\NormalTok{(v) }\ControlFlowTok{for}\NormalTok{ v }\KeywordTok{in}\NormalTok{ web.values())}
\NormalTok{    C }\OperatorTok{=}\NormalTok{ L }\OperatorTok{/}\NormalTok{ (S }\OperatorTok{*}\NormalTok{ S)}
    \ControlFlowTok{return}\NormalTok{ \{}\StringTok{"S"}\NormalTok{: S, }\StringTok{"L"}\NormalTok{: L, }\StringTok{"L/S"}\NormalTok{: L }\OperatorTok{/}\NormalTok{ S, }\StringTok{"connectance"}\NormalTok{: C\}}

\NormalTok{example }\OperatorTok{=}\NormalTok{ \{}
    \StringTok{"phyto"}\NormalTok{: [], }\StringTok{"zoo"}\NormalTok{: [}\StringTok{"phyto"}\NormalTok{], }\StringTok{"shrimp"}\NormalTok{: [}\StringTok{"phyto"}\NormalTok{],}
    \StringTok{"minnow"}\NormalTok{: [}\StringTok{"zoo"}\NormalTok{, }\StringTok{"shrimp"}\NormalTok{], }\StringTok{"bass"}\NormalTok{: [}\StringTok{"minnow"}\NormalTok{, }\StringTok{"shrimp"}\NormalTok{],}
\NormalTok{\}}
\BuiltInTok{print}\NormalTok{(web\_stats(example))}
\end{Highlighting}
\end{Shaded}

\subsubsection{Cascade Effects and
Robustness}\label{cascade-effects-and-robustness}

Food webs are not just descriptive: their topology predicts how
perturbations propagate. Removal experiments (real or simulated) reveal
three classes of network response:

{\def\LTcaptype{none} % do not increment counter
\begin{longtable}[]{@{}
  >{\raggedright\arraybackslash}p{(\linewidth - 4\tabcolsep) * \real{0.2105}}
  >{\raggedright\arraybackslash}p{(\linewidth - 4\tabcolsep) * \real{0.2368}}
  >{\raggedright\arraybackslash}p{(\linewidth - 4\tabcolsep) * \real{0.5526}}@{}}
\toprule\noalign{}
\begin{minipage}[b]{\linewidth}\raggedright
Response
\end{minipage} & \begin{minipage}[b]{\linewidth}\raggedright
Mechanism
\end{minipage} & \begin{minipage}[b]{\linewidth}\raggedright
Topological signature
\end{minipage} \\
\midrule\noalign{}
\endhead
\bottomrule\noalign{}
\endlastfoot
\textbf{Bottom-up cascade} & Loss of a producer ripples upward through
consumers & Few generalist consumers; long chains \\
\textbf{Top-down cascade} (\cref{eq:community_ecology_5} system) & Loss
of a top predator releases mesopredators or herbivores & Strong vertical
interaction strengths \\
\textbf{Horizontal cascade} & Loss of one prey shifts predator to
alternative prey, harming it indirectly & Apparent competition; shared
predators \\
\end{longtable}
}

\textbf{Robustness analysis} (Dunne et al.~2002, \emph{Ecol. Lett.})
sequentially removes species (random order, most-connected first, or
rarest first) and tracks the fraction of secondary extinctions. Results
across 16 well-resolved food webs:

\begin{itemize}
\tightlist
\item
  Random removal: about 50\% of species can disappear before
  web collapse
\item
  \textbf{Most-connected-first removal: collapse after about 20\% of
  species removed} --- hubs (highly connected generalists) are
  disproportionately important
\item
  Rarest-first: minimal cascade --- rare species are weakly embedded
\end{itemize}

This is a striking parallel to the \textbf{scale-free network}
robustness literature (Barabási \& Albert 1999): tolerance to random
failure but vulnerability to targeted attack. For conservation, hub
species are ecological-network analogues of keystone species,
identifiable from topology alone.

\subsubsection{Complexity-Stability
Debate}\label{complexity-stability-debate}

\textbf{May (1972, \emph{Nature}):} Mathematical analysis of random
community matrices showed that complexity (high \(S\), high \(C\),
strong interaction strengths) \textbf{destabilises} communities ---
contradicting the intuitive ``diversity begets stability'' hypothesis.
The May criterion: a random Jacobian is locally stable primarily if
\(\sigma\sqrt{SC} < 1\) (with \texorpdfstring{\ensuremath{\sigma}}{sigma} = standard deviation of interaction
strengths).

\textbf{Resolution (McCann 2000, \emph{Nature}; Allesina \& Tang 2012,
\emph{Nature}):} - Weak interactions stabilise food webs (dampen
oscillations) - Non-random interaction structure (modularity,
nestedness) promotes stability - Real food webs are not random --- they
have specific architectures that enhance stability -
\textbf{Weak-interaction dominance:} most interactions are weak; few are
strong. The many weak links act as stabilisers.

\textbf{Modularity:} food webs are organised into modules (compartments)
with strong within-module interactions and weak between-module
interactions. This structure limits the spread of perturbations --- a
forest pathogen outbreak in one module rarely cascades into the aquatic
module of the same landscape food web.

\begin{center}\rule{0.5\linewidth}{0.5pt}\end{center}

\subsection{Neutral Theory of
Biodiversity}\label{neutral-theory-of-biodiversity}

The competition-coexistence theory developed earlier explains community
structure through \textbf{niches} --- species coexist because they
differ. \textbf{Hubbell's unified neutral theory of biodiversity and
biogeography (UNTB; Hubbell 2001)} proposes the radical alternative: at
the scale of trophically similar species (e.g., canopy trees in a
tropical forest), \textbf{most individuals are demographically
equivalent} regardless of species, and observed community patterns
emerge from random birth--death--dispersal--speciation alone.

\subsubsection{The Neutral Assumption}\label{the-neutral-assumption}

Each individual, regardless of species, has the same per-capita
probabilities of: - Birth (replacement of a vacancy by an offspring) -
Death (creation of a vacancy) - Immigration \(m\) from a regional
\textbf{metacommunity} (replaces local community via dispersal) -
Speciation (rate \texorpdfstring{\ensuremath{\nu}}{nu} in the metacommunity)

The local community has fixed size \(J\) (zero-sum constraint --- every
death is replaced); dynamics are pure ecological drift, mathematically
isomorphic to neutral \hyperref[gl:genetic-drift]{\textbf{genetic
drift}} (\cref{sec:unit_V_population_genetics}). Two parameters do most
of the work:

\begin{equation}
\theta = 2 J_M \nu \qquad \text{(fundamental biodiversity number)}
\label{eq:unit_X_neutral_theta}
\end{equation}

\begin{equation}
m = \text{immigration probability per local death}
\label{eq:unit_X_neutral_m}
\end{equation}

The metacommunity size is \(J_M\). Together, \texorpdfstring{\ensuremath{\theta}}{theta} and \(m\) predict species
abundance distributions, species--area curves, and beta diversity ---
without invoking niche differences at most.

\subsubsection{Predictions and Empirical
Successes}\label{predictions-and-empirical-successes}

UNTB predicts: - \textbf{Species abundance distribution} follows a
zero-sum multinomial (very close to the rank-abundance pattern
introduced earlier, but with a heavier tail of rare species) -
\textbf{Species--area relationships} with \(z \approx 0.20\)--0.30
(matching empirical island values, \cref{eq:community_ecology_12}) -
\textbf{Beta diversity} declines with geographic distance even in
homogeneous habitat (purely from limited dispersal)

Hubbell's analysis of the Barro Colorado Island 50-ha forest plot
(Panama, \textgreater{} 200 tree species) showed UNTB fit the species
abundance distribution as well as niche-based models. Volkov et
al.~(2003, \emph{Nature}) confirmed neutral fits across forests in
Panama, Ecuador, India, and Malaysia.

\subsubsection{Limits and Synthesis}\label{limits-and-synthesis}

UNTB has been heavily criticised: real species are \emph{not} equivalent
(massive trait variation; clear niche differences in measured
demographic rates), and dynamic predictions (e.g., extinction times)
often fail. The modern view, following Adler, HilleRisLambers \& Levine
(2007, \emph{Ecol. Lett.}), is that \textbf{niche and neutral processes
co-occur}: Chesson-style niche differences stabilise coexistence, while
neutral drift adds stochastic variation in relative abundance. UNTB is
the null model --- if your data look neutral, you have not yet found the
ecology that distinguishes species.

\begin{quote}
\textbf{Concept Check:} Two tropical forest plots have nearly identical
species-abundance distributions. The first sits in a homogeneous lowland
habitat; the second spans a steep elevation gradient. Why is the
\emph{neutral} explanation of the species-abundance pattern more
plausible for the first plot than the second, and what additional data
would discriminate the two hypotheses?
\end{quote}

\begin{center}\rule{0.5\linewidth}{0.5pt}\end{center}

\subsection{Trait-Based Community
Ecology}\label{trait-based-community-ecology}

Counting species ignores that a 1 mm aphid and a 10 m oak both add
``+1'' to richness. \textbf{Trait-based ecology} replaces (or
complements) species lists with measured \textbf{functional traits} ---
morphological, physiological, or phenological attributes that influence
performance. This shifts community ecology from a taxonomic enterprise
to a quantitative, predictive one.

\subsubsection{Grime's CSR Triangle}\label{grimes-csr-triangle}

Grime's CSR scheme (1977, 2001) classifies plant strategies along two
stress axes --- disturbance and stress (resource limitation) ---
yielding three primary strategies and the gradients between them:

{\def\LTcaptype{none} % do not increment counter
\begin{longtable}[]{@{}
  >{\raggedright\arraybackslash}p{(\linewidth - 6\tabcolsep) * \real{0.2000}}
  >{\raggedright\arraybackslash}p{(\linewidth - 6\tabcolsep) * \real{0.2500}}
  >{\raggedright\arraybackslash}p{(\linewidth - 6\tabcolsep) * \real{0.3500}}
  >{\raggedright\arraybackslash}p{(\linewidth - 6\tabcolsep) * \real{0.2000}}@{}}
\toprule\noalign{}
\begin{minipage}[b]{\linewidth}\raggedright
Strategy
\end{minipage} & \begin{minipage}[b]{\linewidth}\raggedright
Conditions
\end{minipage} & \begin{minipage}[b]{\linewidth}\raggedright
Trait syndrome
\end{minipage} & \begin{minipage}[b]{\linewidth}\raggedright
Examples
\end{minipage} \\
\midrule\noalign{}
\endhead
\bottomrule\noalign{}
\endlastfoot
\textbf{C --- Competitor} & Low stress, low disturbance & Tall,
fast-growing, high biomass, deep roots; leaves long-lived but expensive
& Mature forest dominants (oak, beech, \emph{Eucalyptus}) \\
\textbf{S --- Stress-tolerator} & High stress, low disturbance & Slow
growth, evergreen, sclerophyllous leaves, conservative resource use &
Desert succulents, alpine cushion plants, lichens, conifers in
nutrient-poor soils \\
\textbf{R --- Ruderal} & Low stress, high disturbance & Short-lived,
high seed output, rapid maturation, weak competitors & Annual weeds
(\emph{Senecio}, \emph{Capsella}); pioneer trees \\
\end{longtable}
}

Most species fall along intermediates (CR, CS, SR, CSR). Crucially, the
CSR axes capture independent gradients to the
\hyperref[gl:r-strategist]{\textbf{r-strategist}}/\hyperref[gl:k-strategist]{\textbf{K-strategist}}
axis of population ecology --- a stress-tolerator is K-like in longevity
but slow-growing, not high-fecundity.

\subsubsection{The Leaf Economics
Spectrum}\label{the-leaf-economics-spectrum}

Wright et al.~(2004, \emph{Nature}) analysed \textgreater{} 2,500 plant
species globally and uncovered a stunningly tight axis along which leaf
traits co-vary --- the \textbf{leaf economics spectrum (LES)}:

{\def\LTcaptype{none} % do not increment counter
\begin{longtable}[]{@{}
  >{\raggedright\arraybackslash}p{(\linewidth - 2\tabcolsep) * \real{0.4889}}
  >{\raggedright\arraybackslash}p{(\linewidth - 2\tabcolsep) * \real{0.5111}}@{}}
\toprule\noalign{}
\begin{minipage}[b]{\linewidth}\raggedright
Acquisitive (fast) end
\end{minipage} & \begin{minipage}[b]{\linewidth}\raggedright
Conservative (slow) end
\end{minipage} \\
\midrule\noalign{}
\endhead
\bottomrule\noalign{}
\endlastfoot
High specific leaf area (SLA, m\(^2\) kg\(^{-1}\)) & Low SLA (thick,
dense leaves) \\
High mass-based photosynthetic rate (\(A_{\text{mass}}\)) & Low
\(A_{\text{mass}}\) \\
High leaf N concentration (Rubisco-rich) & Low leaf N \\
Short leaf lifespan (months) & Long leaf lifespan (years; evergreen) \\
Low construction cost per unit area & High construction cost;
defended \\
\end{longtable}
}

The LES is widely observed across biomes --- tropical pioneer trees and
desert annuals can occupy the same acquisitive end of the axis despite
very different evolutionary histories. The axis captures a major
trade-off between rapid resource capture and resource conservation, but
local exceptions occur because water stress, herbivory, nutrient
limitation, and phylogeny can bend the relationship. Analogous spectra
exist for roots (Bergmann et al.~2020, \emph{Sci. Adv.}) and for
whole-plant strategy axes (Diaz et al.~2016, \emph{Nature}).

\subsubsection{Functional Diversity and Community
Assembly}\label{functional-diversity-and-community-assembly}

Trait data enable \textbf{functional diversity} indices (functional
richness, evenness, divergence; Villéger, Mason \& Mouillot 2008,
\emph{Ecology}) that detect community assembly mechanisms invisible to
species lists:

\begin{itemize}
\tightlist
\item
  \textbf{Functional clustering} (traits more similar than expected by
  chance) \texorpdfstring{\ensuremath{\to}}{->} \emph{environmental filtering} (primarily some strategies
  tolerate the local abiotic conditions)
\item
  \textbf{Functional overdispersion} (traits more different than
  expected) \texorpdfstring{\ensuremath{\to}}{->} \emph{competitive exclusion} via limiting similarity and
  Chesson-style stabilising mechanisms.
\end{itemize}

A community of 30 alpine cushion plants and a community of 30 tropical
canopy trees both have \(S = 30\), but their trait spreads on the LES
are wildly different --- and that difference predicts ecosystem function
(productivity, decomposition, drought response) far better than richness
alone.

\begin{quote}
\textbf{Concept Check:} Two grassland plots have identical Shannon
diversity \(H' = 2.5\). Plot A's plants cluster tightly on the
conservative end of the LES; plot B's plants span the entire LES. Which
plot is more productive on average, which is more drought-resilient, and
which would you predict has stronger competitive interactions?
\end{quote}

\begin{center}\rule{0.5\linewidth}{0.5pt}\end{center}

\subsection{Biological Control}\label{biological-control}

The competition, predation, and parasitism theory developed above has
direct translational application: \textbf{biological control} uses
natural enemies to suppress pest populations, replacing or reducing
chemical pesticides. It is community ecology deployed for agriculture,
public health, and invasive-species management.

\subsubsection{Three Classical
Strategies}\label{three-classical-strategies}

{\def\LTcaptype{none} % do not increment counter
\begin{longtable}[]{@{}
  >{\raggedright\arraybackslash}p{(\linewidth - 6\tabcolsep) * \real{0.2105}}
  >{\raggedright\arraybackslash}p{(\linewidth - 6\tabcolsep) * \real{0.2105}}
  >{\raggedright\arraybackslash}p{(\linewidth - 6\tabcolsep) * \real{0.2632}}
  >{\raggedright\arraybackslash}p{(\linewidth - 6\tabcolsep) * \real{0.3158}}@{}}
\toprule\noalign{}
\begin{minipage}[b]{\linewidth}\raggedright
Strategy
\end{minipage} & \begin{minipage}[b]{\linewidth}\raggedright
Approach
\end{minipage} & \begin{minipage}[b]{\linewidth}\raggedright
Time scale
\end{minipage} & \begin{minipage}[b]{\linewidth}\raggedright
Risk profile
\end{minipage} \\
\midrule\noalign{}
\endhead
\bottomrule\noalign{}
\endlastfoot
\textbf{Classical (importation)} & Import a specialist natural enemy
from the pest's native range; expect it to establish and self-perpetuate
& Years to decades; one-shot release & High --- non-target effects if
enemy is not specialist enough \\
\textbf{Augmentative} & Mass-rear and release natural enemies repeatedly
to suppress current outbreak & Weeks to months; greenhouse, glasshouse &
Low --- released organisms typically fail to overwinter \\
\textbf{Conservation} & Modify habitat to favour resident natural
enemies (hedgerows, beetle banks, insectary plants, reduced pesticide
use) & Continuous & Lowest --- uses native species \\
\end{longtable}
}

\subsubsection{Classical Successes and
Failures}\label{classical-successes-and-failures}

\textbf{Cottony cushion scale (\emph{Icerya purchasi}) and the vedalia
beetle (\emph{Rodolia cardinalis}).} Imported in 1888 from Australia to
California citrus groves --- within two years the citrus industry was
saved. This is the textbook success and the prototype for most
subsequent classical biocontrol programs.

\textbf{Prickly pear (\emph{Opuntia stricta}) and \emph{Cactoblastis
cactorum}.} Released in 1925 in Australia, the moth larvae cleared 25
million ha of invasive cactus by 1940 --- one of the largest-scale
ecological interventions in history.

\textbf{Cane toad (\emph{Rhinella marina}).} Released in 1935 in
Queensland to control sugarcane beetles; the toads ignored the beetles,
ate everything else, and are now an invasive plague across northern
Australia. Classical biocontrol's most-cited cautionary tale: a
generalist natural enemy will eat non-target species, and the lesson
cost ecologists decades of credibility.

\subsubsection{Risk Assessment and Modern
Practice}\label{risk-assessment-and-modern-practice}

After the cane-toad era, regulatory frameworks (e.g., the FAO
\textbf{Code of Conduct for the Import and Release of Exotic Biological
Control Agents}, 1996) require:

\begin{enumerate}
\def\labelenumi{\arabic{enumi}.}
\tightlist
\item
  \textbf{Host-specificity testing} --- quarantine trials against native
  non-target species
\item
  \textbf{Climate matching} --- does the agent's native climate predict
  its spread in the target range?
\item
  \textbf{Population modelling} --- Lotka--Volterra or matrix projection
  of agent and target dynamics; predicted suppression vs.~escape
\item
  \textbf{Reversibility analysis} --- can the agent be eradicated if it
  causes damage? (Almost typically: no.)
\end{enumerate}

Modern programs increasingly favour \textbf{conservation biological
control}, which avoids introducing non-native species entirely. Field
margins of insectary plants (e.g., \emph{Phacelia}, \emph{Fagopyrum})
increase syrphid and parasitoid wasp populations that suppress aphids in
adjacent crops by 30--50 \% (Tscharntke et al.~2007, \emph{Biol.
Control}). This connects directly to ecosystem-services valuation
(\cref{sec:unit_X_ecosystem_ecology}).

\begin{quote}
\texorpdfstring{}{} \textbf{Clinical Connection --- \emph{Wolbachia} and Mosquito-Borne
Disease.} Conservation biological control extends to public health.
\emph{Wolbachia} is an intracellular bacterium that, when introduced
into \emph{Aedes aegypti} mosquitoes, blocks dengue, Zika, and
chikungunya virus transmission. Field trials in Yogyakarta (Indonesia)
and Niterói (Brazil) released \emph{Wolbachia}-carrying mosquitoes that
spread the bacterium through the wild population via cytoplasmic
incompatibility. Within 27 months of release in Yogyakarta, dengue
incidence dropped 77\% in treated neighbourhoods (Utarini et al.~2021,
\emph{N. Engl. J. Med.}). This is biological control of disease vectors
via a manipulative endosymbiont --- community ecology in service of
pandemic prevention.
\end{quote}

\begin{quote}
\textbf{Concept Check:} A regulator must approve or reject a proposed
classical biocontrol release of a parasitoid wasp against an invasive
moth pest. (a) What three host-specificity tests would you require? (b)
Why is ``the agent attacks the target moth in the lab'' insufficient
evidence for safety? (c) How does a sound elasticity analysis
(\cref{eq:unit_X_sensitivity_elasticity}) of the \emph{target's} matrix
model help you predict whether the release will succeed?
\end{quote}

\begin{center}\rule{0.5\linewidth}{0.5pt}\end{center}

\subsection{Worked Example}\label{worked-example-6}

\textbf{Problem:} An ecologist surveys a small forest patch and
identifies three species of trees: - Species A: 50 individuals - Species
B: 30 individuals - Species C: 20 individuals

Calculate the Shannon Diversity Index (\(H\)) for this tree community.
The formula is:
\[H = -\sum (p_i \ln p_i) \label{eq:unit_X_community_ecology_item_1}\]

where \(p_i\) is the proportion of total individuals belonging to the
\(i\)-th species.

\textbf{Solution:}

\textbf{Step 1. Calculate the total number of individuals (\(N\)).}
\[N = 50 + 30 + 20 = 100 \label{eq:unit_X_community_ecology_item_2}\]

\textbf{Step 2. Calculate the proportion (\(p_i\)) for each species.} -
\(p_A = 50 / 100 = 0.50\) - \(p_B = 30 / 100 = 0.30\) -
\(p_C = 20 / 100 = 0.20\)

\textbf{Step 3. Calculate \(p_i \ln p_i\) for each species.} - Species
A: \(0.50 \times \ln(0.50) \approx 0.50 \times (-0.693) \approx -0.347\)
- Species B:
\(0.30 \times \ln(0.30) \approx 0.30 \times (-1.204) \approx -0.361\) -
Species C:
\(0.20 \times \ln(0.20) \approx 0.20 \times (-1.609) \approx -0.322\)

\textbf{Step 4. Sum the values and multiply by \(-1\).}
\[H = - (-0.347 + -0.361 + -0.322) = - (-1.030) = 1.030 \label{eq:unit_X_community_ecology_item_3}\]

\textbf{Answer:} The Shannon Diversity Index for this tree community is
approximately \textbf{1.03}.

\begin{center}\rule{0.5\linewidth}{0.5pt}\end{center}

\subsection{Current Evidence and Frontier
Biology}\label{current-evidence-and-frontier-biology-35}

For \textbf{Community Ecology and Species Interactions}, frontier
biology belongs inside the evidence logic of the chapter. Ecology and
conservation decisions increasingly combine field data, remote sensing,
community knowledge, model uncertainty, and explicit values. The core
reading question is this: community claims should identify interaction
type, network position, disturbance regime, and observational limits.

\begin{itemize}
\tightlist
\item
  \textbf{What to verify:} identify the observation, model, assay, or
  dataset that would make the claim stronger or weaker.
\item
  \textbf{What to qualify:} state the scale, organism, cell type,
  environmental condition, or population where the claim is expected to
  hold.
\item
  \textbf{What to compare:} test at least one alternative explanation,
  baseline, or null model before treating the pattern as causal.
\item
  \textbf{What to cite:} distinguish primary evidence, review synthesis,
  public dataset, and institutional guidance; for recent or numeric
  claims, prefer the source closest to the measurement and state what
  has changed since it was published.
\end{itemize}

Use biodiversity metrics carefully: population indices, extinction risk
categories, ecosystem services, and management targets answer different
questions
\citep{ipbes2019global,ipbes2024transformative,wwf2024livingplanet,iucn2025redlist,fao2024sofia}.

\textbf{Source practice:} For conservation claims, cite assessment
sources and state whether the evidence is a population index,
extinction-risk assessment, ecosystem-service valuation, satellite
product, or policy synthesis
\citep{ipbes2024transformative,noaa2025coralbleaching,fao2025sofi}.

\subsection{Key Terms}\label{key-terms-35}

{\def\LTcaptype{none} % do not increment counter
\begin{longtable}[]{@{}
  >{\raggedright\arraybackslash}p{(\linewidth - 2\tabcolsep) * \real{0.2857}}
  >{\raggedright\arraybackslash}p{(\linewidth - 2\tabcolsep) * \real{0.7143}}@{}}
\toprule\noalign{}
\begin{minipage}[b]{\linewidth}\raggedright
Term
\end{minipage} & \begin{minipage}[b]{\linewidth}\raggedright
Definition
\end{minipage} \\
\midrule\noalign{}
\endhead
\bottomrule\noalign{}
\endlastfoot
\textbf{Competitive exclusion principle} & Two species cannot stably
coexist if they occupy the exact same niche; the superior competitor
excludes the other \\
\textbf{Keystone species} & Species with disproportionately large
effects on community structure relative to its abundance \\
\textbf{Trophic cascade} & Indirect effect of apex predator on plant
biomass via suppressing herbivores; top-down regulation \\
\textbf{Succession} & Directional change in community composition over
time; primary (bare substrate) or secondary (after disturbance) \\
\textbf{Intermediate Disturbance Hypothesis} & Moderate disturbance
frequency/intensity maximises species diversity by preventing
competitive exclusion \\
\textbf{Shannon diversity (\(H'\))} & \(H' = -\sum p_i \ln p_i\);
quantifies species richness and evenness simultaneously \\
\textbf{Simpson index (\(1-D\))} & Probability that two random
individuals are from different species; weighted toward dominant
species \\
\textbf{Species-area relationship} & \(S = cA^z\); species richness
increases with area; \(z \approx 0.25\) for habitat islands \\
\textbf{SLOSS} & Single Large Or Several Small reserves debate in
conservation biology \\
\textbf{Extinction debt} & Species present today but predicted to go
extinct based on current habitat loss; time-delayed response \\
\textbf{Metapopulation} & Network of semi-isolated subpopulations
connected by dispersal; local extinctions offset by recolonisation \\
\textbf{Alternative stable states} & Multiple stable community
configurations under the same environmental conditions; regime shifts
between states \\
\textbf{Niche differentiation} & Process by which competing species
evolve to use different resources, reducing niche overlap \\
\textbf{Fundamental niche} & Full range of environmental conditions
where a species can maintain \(r \geq 0\) (Hutchinson) \\
\textbf{Realised niche} & Subset of fundamental niche actually occupied
after biotic interactions \\
\textbf{Connectance} & Proportion of possible trophic links that are
realised in a food web \\
\textbf{Facilitation} & One species improves survival or reproduction of
another; mechanism of succession \\
\textbf{Beta diversity} & Species turnover between communities; measured
by Jaccard or Bray-Curtis indices \\
\textbf{Neutral theory (UNTB)} & Hubbell's framework: per-capita
demographic equivalence; biodiversity arises from drift, dispersal, and
speciation \\
\textbf{Fundamental biodiversity number (\texorpdfstring{\ensuremath{\theta}}{theta})} & \(\theta = 2J_M\nu\);
controls species richness in the metacommunity under UNTB \\
\textbf{CSR strategies} & Grime's competitor / stress-tolerator /
ruderal classification of plant strategies \\
\textbf{Leaf economics spectrum (LES)} & Comprehensive axis from
acquisitive (high SLA, short-lived leaves) to conservative (thick,
long-lived) leaves \\
\textbf{Functional diversity} & Trait-based diversity metric; clustering
vs.~overdispersion diagnoses assembly mechanism \\
\textbf{Classical biological control} & Importation of a specialist
natural enemy from a pest's native range \\
\textbf{Conservation biological control} & Habitat modification to
favour resident natural enemies; lowest-risk strategy \\
\textbf{Network robustness} & Resilience of a food web to species
removal; collapses fastest under most-connected-first deletion \\
\end{longtable}
}

\begin{center}\rule{0.5\linewidth}{0.5pt}\end{center}

\subsection{Review Questions}\label{review-questions-35}

\begin{enumerate}
\def\labelenumi{\arabic{enumi}.}
\item
  Two warbler species (\emph{Dendroica castanea} and \emph{D. fusca})
  share a boreal spruce tree habitat. MacArthur (1958) showed they
  coexist by partitioning the tree into feeding zones. Using Chesson's
  modern coexistence theory: (a) Which of Chesson's two mechanisms
  (equalising vs.~stabilising) enables their coexistence? (b) Write
  \hyperref[gl:lotka-volterra-equations]{\textbf{Lotka-Volterra
  equations}} and specify the condition for stable coexistence in terms
  of \(\alpha_{12}\) and \(\alpha_{21}\).
\item
  A trophic cascade operates in Yellowstone: wolves \texorpdfstring{\ensuremath{\to}}{->} elk \texorpdfstring{\ensuremath{\to}}{->} willows \texorpdfstring{\ensuremath{\to}}{->}
  beavers \texorpdfstring{\ensuremath{\to}}{->} altered stream morphology. (a) Is this a bottom-up or
  top-down cascade? (b) Predict what happens to each level if chronic
  wasting disease eliminates 90\% of the elk population. (c) How does
  the ``landscape of fear'' concept modify purely density-mediated
  cascade predictions?
\item
  A tropical rainforest island (area = 500 km\(^2\)) currently has
  \(S = 200\) bird species. Deforestation reduces area to 125 km\(^2\)
  (25\% of original). (a) Using the species-area relationship with
  \(z = 0.25\) and appropriate \(c\), predict the final equilibrium
  species richness. (b) If the extinction debt takes 50 years to fully
  manifest, what is the expected species number at year 10 assuming
  linear debt decay?
\item
  Compare the Shannon diversity index (\(H'\)) for three communities:

  \begin{itemize}
  \tightlist
  \item
    Community A: 4 species, abundances = {[}25, 25, 25, 25{]}
  \item
    Community B: 4 species, abundances = {[}97, 1, 1, 1{]}
  \item
    Community C: 8 species, abundances = {[}50, 20, 10, 8, 5, 4, 2, 1{]}
    Calculate \(H'\) and \(J'\) for each. Which community shows: highest
    richness? highest evenness? highest \(H'\)?
  \end{itemize}
\item
  Explain Hutchinson's ``paradox of the plankton.'' Why does the
  coexistence of hundreds of phytoplankton species apparently violate
  competitive exclusion, and what are three mechanisms that resolve the
  paradox?
\item
  The Biological Dynamics of Forest Fragments Project (BDFFP) in
  Amazonia found that 1-ha fragments lost 50\% of their bird species
  within 15 years, while 100-ha fragments lost 10\%. (a) Calculate \(z\)
  from these two data points. (b) Predict the species richness of a
  10-ha fragment relative to a 100-ha fragment. (c) What role do edge
  effects play beyond the simple species-area prediction?
\item
  A shallow lake in the Netherlands exists in a clear-water state with
  abundant macrophytes. Nutrient loading from agriculture gradually
  increases. Describe the regime shift to the turbid, algal-dominated
  state using the framework of alternative stable states. Why is
  reducing nutrients back to the original level insufficient to restore
  the clear-water state (hysteresis)?
\item
  Design a reserve system for a large mammal species (home range = 100
  km\(^2\), \(N_e\) requirement = 500) in a landscape of forest
  fragments. Use island biogeography principles to specify: (a) minimum
  total reserve area, (b) number and size of core reserves, (c) corridor
  design, (d) matrix management. Justify each design element.
\item
  The rough-skinned newt and common garter snake represent one of the
  best-documented coevolutionary arms races. Describe the escalation:
  what toxin does the newt produce, what resistance mechanism has the
  snake evolved, and what geographic mosaic pattern is observed? How
  does this relate to Red Queen dynamics?
\item
  Explain how the competitive exclusion principle applies to
  \emph{Clostridium difficile} infection and its treatment by fecal
  microbiota transplantation. What ecological principles make FMT more
  effective than antibiotic treatment for recurrent CDI?
\item
  Using \texttt{lotka\_volterra}, estimate whether predator peak
  \textbf{lags} prey peak for default parameters (inspect time series
  qualitatively).
\item
  How does \textbf{connectance} \(C=L/S^2\) relate to stability
  arguments in diverse food webs?
\item
  Hubbell's neutral theory predicts a species-abundance distribution
  from drift and dispersal alone. (a) Why is UNTB best regarded as a
  \emph{null model} rather than a competitor to niche theory? (b) If a
  community fits UNTB perfectly, what does that \emph{not} prove about
  the underlying ecology? Cite the Adler, HilleRisLambers \& Levine
  (2007) synthesis.
\item
  Two oak forests share identical species lists (\(S = 30\),
  \(H' = 2.8\)) but differ markedly in leaf economics: forest A's
  species cluster on the conservative end of the LES; forest B's species
  span the entire spectrum. (a) Which forest is more functionally
  diverse? (b) Predict relative productivity, decomposition rate, and
  drought-resilience for each. (c) Why do two communities with identical
  Shannon diversity behave differently?
\item
  A food-web robustness analysis simulates extinctions in random
  vs.~most-connected-first orders. The web collapses after 50\% removal
  under random deletion but after 18\% removal under
  most-connected-first deletion. (a) Explain the asymmetry using network
  topology. (b) Translate the result into a conservation prioritisation
  rule for hub species (analogous to keystone identification but
  topological).
\item
  A government agency proposes releasing a parasitoid wasp from East
  Asia to control an invasive aphid species in California. (a) Outline
  the host-specificity testing required before approval. (b) Why is the
  cane toad disaster relevant precedent? (c) Sketch how the agent and
  aphid populations would interact in a coupled Lotka-Volterra model ---
  and what condition supports long-term aphid suppression rather than
  coexistence?
\end{enumerate}

\begin{center}\rule{0.5\linewidth}{0.5pt}\end{center}

\subsection{Further Reading and Source
Notes}\label{further-reading-and-source-notes-30}

\begin{itemize}
\tightlist
\item
  Paine (1966). Food Web Complexity and Species Diversity. \emph{The
  American Naturalist}, 100.
\item
  Connell (1978). Diversity in tropical rain forests and coral reefs.
  \emph{Science}, 199.
\item
  Gause (1934). \emph{The Struggle for Existence}. Williams \& Wilkins.
\item
  Hutchinson (1957). Concluding remarks. \emph{Cold Spring Harbor
  Symposia on Quantitative Biology}, 22.
\item
  Chesson (2000). Mechanisms of maintenance of species diversity.
  \emph{Annual Review of Ecology and Systematics}, 31.
\item
  Ehrlich \& Raven (1964). Butterflies and plants: A study in
  coevolution. \emph{Evolution}, 18.
\end{itemize}

\begin{center}\rule{0.5\linewidth}{0.5pt}\end{center}

\subsection{Computational Bridge}\label{computational-bridge-32}

Predator--prey cycles from the chapter are reproduced by numerical
integration:

\begin{Shaded}
\begin{Highlighting}[]
\ImportTok{from}\NormalTok{ biology.ecology }\ImportTok{import}\NormalTok{ lotka\_volterra}

\NormalTok{lv }\OperatorTok{=}\NormalTok{ lotka\_volterra(}\FloatTok{40.0}\NormalTok{, }\FloatTok{9.0}\NormalTok{, }\FloatTok{0.5}\NormalTok{, }\FloatTok{0.02}\NormalTok{, }\FloatTok{0.01}\NormalTok{, }\FloatTok{0.2}\NormalTok{, t\_end}\OperatorTok{=}\FloatTok{80.0}\NormalTok{)}
\BuiltInTok{print}\NormalTok{(}\BuiltInTok{len}\NormalTok{(lv.times), }\BuiltInTok{round}\NormalTok{(lv.prey[}\OperatorTok{{-}}\DecValTok{1}\NormalTok{], }\DecValTok{2}\NormalTok{))}
\end{Highlighting}
\end{Shaded}

\begin{quote}
\textbf{Clinical / systems note:} FMT restores diversity and
colonisation resistance --- an ecological intervention for a microbiome
community treated as a competitive network.
\end{quote}

\begin{center}\rule{0.5\linewidth}{0.5pt}\end{center}

\subsection{Summary}\label{summary-35}

\begin{itemize}
\tightlist
\item
  \textbf{Biotic interactions:} mutualism (+/+), commensalism (+/0),
  parasitism (+/-), predation (+/-), competition (-/-), amensalism (0/-)
  --- shape community composition and evolution.
\item
  \textbf{Competition:} Lotka-Volterra equations; competitive exclusion
  if niches overlap completely; stable coexistence if intraspecific
  \textgreater{} interspecific effects. Modern coexistence theory
  \citep{chesson2000}: stabilising + equalising mechanisms.
\item
  \textbf{Niche theory:} Hutchinson's n-dimensional hypervolume;
  fundamental vs.~realised niche; paradox of the plankton resolved by
  non-equilibrium dynamics and niche partitioning.
\item
  \textbf{Trophic cascades:} keystone predators suppress herbivores \texorpdfstring{\ensuremath{\to}}{->}
  plant biomass increases (wolves\texorpdfstring{\ensuremath{\to}}{->}elk\texorpdfstring{\ensuremath{\to}}{->}willows; sea otters\texorpdfstring{\ensuremath{\to}}{->}urchins\texorpdfstring{\ensuremath{\to}}{->}kelp).
  Behavioural cascades via ``landscape of fear.''
\item
  \textbf{Succession:} primary (bare substrate, slow) vs.~secondary
  (disturbed community, faster); facilitation/tolerance/inhibition
  mechanisms; IDH: intermediate disturbance maximises diversity;
  alternative stable states and regime shifts with hysteresis.
\item
  \textbf{Diversity indices:} \(H'\) (Shannon) combines richness +
  evenness; Simpson \(1-D\) = probability of inter-specific encounter;
  \(J'\) (evenness) = \(H'/\ln S\). Alpha, beta, gamma diversity at
  different spatial scales.
\item
  \textbf{Island biogeography:} \(S = cA^z\); equilibrium balances
  colonisation and extinction; SLOSS debate resolved toward
  core-corridor-matrix design; extinction debt warns of delayed species
  loss.
\item
  \textbf{Food web structure:} connectance, modularity, weak-interaction
  dominance stabilise complex communities. \textbf{Network robustness}
  under targeted (most-connected-first) deletion collapses faster than
  under random deletion --- hub species are topological keystones.
\item
  \textbf{Neutral theory (UNTB):} Hubbell's per-capita-equivalence
  framework; two parameters \texorpdfstring{\ensuremath{\theta}}{theta} and \(m\) predict species-abundance
  distributions and species--area relationships from drift, dispersal,
  and speciation alone. Best treated as a null model that niche-based
  theory must beat.
\item
  \textbf{Trait-based ecology:} Grime's CSR triangle; the leaf economics
  spectrum (LES) is a widely observed acquisitive\texorpdfstring{\ensuremath{\leftrightarrow}}{<->}conservative axis with
  local exceptions. Functional diversity diagnoses assembly mechanisms
  (clustering = environmental filtering; overdispersion = limiting
  similarity) that species lists miss.
\item
  \textbf{Biological control:} classical importation (vedalia beetle
  success; cane toad disaster), augmentative releases, and conservation
  biological control. \emph{Wolbachia}-loaded \emph{Aedes aegypti}
  extends the framework to vector-borne disease.
\item
  \textbf{Connections:} See \cref{sec:unit_X_population_ecology} for
  consumer-resource oscillations, \cref{sec:unit_X_ecosystem_ecology}
  for energy flux, and \cref{sec:unit_VII_microbial_ecology} for
  microbial communities.
\end{itemize}

\begin{center}\rule{0.5\linewidth}{0.5pt}\end{center}

\subsubsection{Companion Source
Module}\label{companion-source-module-35}

\textbf{Community Ecology and Species Interactions} should leave a
reproducible trail from a biological claim to the code, figure, diagram,
or paper-based activity that can test it. Use the surfaces below to
inspect the chapter's assumptions, rerun the relevant model, or compare
the manuscript explanation with companion labs and figures.

{\def\LTcaptype{none} % do not increment counter
\begin{longtable}[]{@{}
  >{\raggedright\arraybackslash}p{(\linewidth - 2\tabcolsep) * \real{0.5000}}
  >{\raggedright\arraybackslash}p{(\linewidth - 2\tabcolsep) * \real{0.5000}}@{}}
\toprule\noalign{}
\begin{minipage}[b]{\linewidth}\raggedright
Surface
\end{minipage} & \begin{minipage}[b]{\linewidth}\raggedright
Use it for
\end{minipage} \\
\midrule\noalign{}
\endhead
\bottomrule\noalign{}
\endlastfoot
\breaktt{src/biology/ecology/ecology.py} (\texttt{lotka\_volterra},
\texttt{connectance}, \breaktt{biodiversity\_indices}) & Quantify
interactions, network structure, and community diversity. \\
\breaktt{src/visualization/plots.py} (\breaktt{plot\_lotka\_volterra},
\breaktt{plot\_species\_area\_relationship}) & Inspect dynamics and
richness-area patterns. \\
\breaktt{src/mermaid/biology\_diagrams.py} (\texttt{food\_web\_diagram})
& Keep trophic links and interaction signs explicit. \\
\end{longtable}
}

\textbf{Reproducibility check:} define interaction sign, spatial scale,
sampling effort, disturbance history, and network boundary before
interpreting community patterns. \textbf{Cross-reference:} use
\cref{sec:unit_X_population_ecology},
\cref{sec:unit_X_ecosystem_ecology}, and
\cref{sec:unit_VII_microbial_ecology}.

\newpage

\section{Ecosystem Ecology and Biogeochemical
Cycles}\label{ecosystem-ecology-and-biogeochemical-cycles}

\label{sec:unit_X_ecosystem_ecology}

\begin{quote}
\textbf{Ch 34} \texorpdfstring{\ensuremath{\cdot}}{.} Level 2/3 \texorpdfstring{\ensuremath{\cdot}}{.} 65 min read \texorpdfstring{\ensuremath{\cdot}}{.} 75 min lecture \texorpdfstring{\ensuremath{\cdot}}{.}
Prerequisites: \cref{sec:unit_X_community_ecology},
\cref{sec:unit_III_photosynthesis}
\end{quote}

\subsection{Learning Objectives}\label{learning-objectives-36}

By the end of this chapter, you should be able to:

\begin{enumerate}
\def\labelenumi{\arabic{enumi}.}
\tightlist
\item
  Define an ecosystem \citep{levin1998} and distinguish
  \hyperref[gl:abiotic]{\textbf{abiotic}} from biotic components;
  categorise consumers by trophic role.
\item
  Explain energy flow through
  \hyperref[gl:trophic-level]{\textbf{trophic level}}s, calculate
  ecological efficiency, and explain why food chains are short.
\item
  Describe the carbon, nitrogen, phosphorus, and sulphur cycles and
  their key anthropogenic disruptions.
\item
  Calculate gross primary production (GPP), net primary production
  (NPP), and net ecosystem production (NEP).
\item
  Describe \hyperref[gl:eutrophication]{\textbf{eutrophication}} and
  hypoxia as consequences of nutrient cycle disruption.
\item
  Explain ecosystem services and their economic valuation.
\item
  Describe climate change feedbacks involving biogeochemical cycles
  (permafrost carbon, ocean acidification).
\item
  Explain the biological pump and its role in ocean carbon
  sequestration.
\item
  Compare methods for measuring NPP (harvest, eddy covariance,
  \(^{14}\)C, remote sensing) and explain why ANPP/BNPP partitioning
  matters for global carbon accounting.
\item
  Use the Hubbard Brook watershed experiment to explain how live
  vegetation regulates nutrient retention, and contrast open vs.~closed
  nutrient cycling regimes.
\item
  Describe ocean productivity zones (upwelling, gyres, polar) and the
  Martin curve, and connect soil pedogenesis (CLORPT, horizons) to
  long-term carbon storage.
\item
  Use the Redfield ratio to diagnose nutrient limitation, and evaluate
  the planetary boundaries and Gaia hypothesis as frameworks for Earth
  system thinking.
\end{enumerate}

\subsubsection{Study Blueprint}\label{study-blueprint-36}

\begin{itemize}
\tightlist
\item
  \textbf{Big idea:} Ecosystems couple energy flow and matter cycling
  across organisms, environments, and time.
\item
  \textbf{Core concepts:} primary productivity, trophic efficiency,
  nutrient cycles, decomposition.
\item
  \textbf{Framework alignment:} Vision \& Change: Systems, Evolution,
  Pathways and transformations of energy and matter; AP Biology: Systems
  Interactions, Evolution, Energetics; NGSS-style topics: Interdependent
  Relationships in Ecosystems, Matter and Energy in Organisms and
  Ecosystems, Natural Selection and Evolution.
\item
  \textbf{Model or quantitative lens:} Energy-flow, productivity, and
  nutrient-budget calculations.
\item
  \textbf{Data skill:} Trace matter and energy through food webs and
  biogeochemical cycles.
\item
  \textbf{Practice cadence:} Representing and Describing Data,
  Statistical Tests and Data Analysis, Argumentation.
\item
  \textbf{Common misconception to repair:} Energy flows through
  ecosystems, but matter cycles; confusing the two breaks many
  explanations.
\item
  \textbf{Primary lab:} \cref{sec:lab_unit_X_ecosystem_ecology}.
\item
  \textbf{Question bank:} \cref{sec:q_unit_X_ecosystem_ecology}.
\item
  \textbf{Transfer task:} Transfer ecosystem reasoning to
  eutrophication, carbon budgets, agriculture, and climate feedback.
\item
  \textbf{Bridge to computation:}
  \breaktt{biology.ecology.ecology.food\_web\_trophic\_levels}.
\end{itemize}

\begin{center}\rule{0.5\linewidth}{0.5pt}\end{center}

\begin{quote}
\textbf{Opening Vignette --- Measuring What an Ecosystem Actually Runs
On}

In 1953, Howard T. Odum waded into Silver Springs, Florida --- a
crystal-clear, constant-temperature spring --- and spent months
measuring oxygen concentrations at stations upstream and downstream to
calculate the primary productivity of an entire ecosystem. His method,
upstream-downstream dissolved oxygen change, was brilliantly simple: if
an ecosystem produces oxygen during the day and consumes it at night,
measuring both rates gives net and gross production. From his data, Odum
constructed the first complete energy flow diagram of a natural
ecosystem, quantifying how much energy entered via
\hyperref[gl:photosynthesis]{\textbf{photosynthesis}}, how much was lost
at each trophic level, and how much left the system as heat. The 10\%
rule --- roughly 10\% of energy transfers between trophic levels ---
came directly from Odum's data. The concept that ecosystems are
energy-processing machines, quantifiable and modelable, born in Silver
Springs, became foundational to biogeochemistry, conservation biology,
fisheries management, and climate modelling. Silver Springs is still one
of the most-studied aquatic ecosystems on Earth.
\end{quote}

\subsection{Ecosystem Concepts}\label{ecosystem-concepts}

An \textbf{ecosystem} comprises most organisms (biotic component) and
the physical/chemical environment (abiotic component) in a defined area,
interacting through energy flow and nutrient cycling
\citep{bormann1967}. Ecosystems vary from a few litres (rock pool) to
thousands of km\(^2\) (boreal forest). Their defining features are:

\begin{enumerate}
\def\labelenumi{\arabic{enumi}.}
\tightlist
\item
  \textbf{Energy input} (solar or chemical) drives biological activity
\item
  \textbf{Nutrient cycling} --- atoms cycle repeatedly between biotic
  and abiotic phases (unlike energy, which flows linearly through and
  out of the system)
\item
  \textbf{Emergent properties} --- productivity, stability, resistance
  --- cannot be predicted from individual species properties alone
\end{enumerate}

\subsubsection{The Two Fundamental
Processes}\label{the-two-fundamental-processes}

The distinction between energy \textbf{flow} and nutrient
\textbf{cycling} is fundamental:

{\def\LTcaptype{none} % do not increment counter
\begin{longtable}[]{@{}
  >{\raggedright\arraybackslash}p{(\linewidth - 4\tabcolsep) * \real{0.3182}}
  >{\raggedright\arraybackslash}p{(\linewidth - 4\tabcolsep) * \real{0.2727}}
  >{\raggedright\arraybackslash}p{(\linewidth - 4\tabcolsep) * \real{0.4091}}@{}}
\toprule\noalign{}
\begin{minipage}[b]{\linewidth}\raggedright
Process
\end{minipage} & \begin{minipage}[b]{\linewidth}\raggedright
Energy
\end{minipage} & \begin{minipage}[b]{\linewidth}\raggedright
Nutrients
\end{minipage} \\
\midrule\noalign{}
\endhead
\bottomrule\noalign{}
\endlastfoot
Direction & One-way (flows through system) & Cyclic (atoms recycled
indefinitely) \\
Source & Solar radiation (99.9\%) or chemosynthesis & Geological and
atmospheric reservoirs \\
Efficiency & about 1-2\% of solar energy captured by
photosynthesis & about 100\% recycled over geological time \\
Limiting? & Rarely limiting (except deep sea, caves) & Often limiting
(N, P, Fe, water) \\
Conservation law & Energy conserved but degraded (2nd law) & Mass
conserved (atoms neither created nor destroyed) \\
\end{longtable}
}

\subsubsection{Ecosystem Services}\label{ecosystem-services}

\textbf{Costanza et al.~(1997, \emph{Nature})} estimated global
ecosystem services at \textgreater\$33 trillion/year (updated to
about \$125 trillion/year by Costanza et al.~2014). The four
categories:

{\def\LTcaptype{none} % do not increment counter
\begin{longtable}[]{@{}
  >{\raggedright\arraybackslash}p{(\linewidth - 6\tabcolsep) * \real{0.1905}}
  >{\raggedright\arraybackslash}p{(\linewidth - 6\tabcolsep) * \real{0.1905}}
  >{\raggedright\arraybackslash}p{(\linewidth - 6\tabcolsep) * \real{0.1905}}
  >{\raggedright\arraybackslash}p{(\linewidth - 6\tabcolsep) * \real{0.4286}}@{}}
\toprule\noalign{}
\begin{minipage}[b]{\linewidth}\raggedright
Category
\end{minipage} & \begin{minipage}[b]{\linewidth}\raggedright
Services
\end{minipage} & \begin{minipage}[b]{\linewidth}\raggedright
Examples
\end{minipage} & \begin{minipage}[b]{\linewidth}\raggedright
Valuation approach
\end{minipage} \\
\midrule\noalign{}
\endhead
\bottomrule\noalign{}
\endlastfoot
\textbf{Provisioning} & Materials and energy & Food, freshwater, timber,
fuel, medicine, genetic resources & Market prices \\
\textbf{Regulating} & Environmental processes & Climate regulation (C
sequestration), flood control, pollination, disease regulation, air
purification, water purification & Avoided cost; replacement cost \\
\textbf{Cultural} & Non-material benefits & Recreation, ecotourism,
spiritual value, aesthetic beauty, education & Willingness-to-pay
surveys \\
\textbf{Supporting} & Foundation for most others & Soil formation,
nutrient cycling, primary productivity, water cycling & Cost of
replication \\
\end{longtable}
}

The \textbf{Millennium Ecosystem Assessment} (2005; 1,300 scientists
from 95 countries) concluded that 60\% of Earth's ecosystem services are
being degraded or used unsustainably. The \textbf{IPBES Global
Assessment} (2019) expanded this to 75\% of terrestrial and 66\% of
marine ecosystems significantly altered.

\begin{quote}
\textbf{Concept Check:} A wetland provides flood attenuation, water
purification, carbon sequestration, recreation, and biodiversity
habitat. If this wetland is drained for agriculture, which ecosystem
services are lost? How would you calculate the economic cost of losing
flood attenuation alone?
\end{quote}

\begin{center}\rule{0.5\linewidth}{0.5pt}\end{center}

\subsection{Trophic Structure and Energy
Flow}\label{trophic-structure-and-energy-flow}

\textbf{Trophic levels} (from Greek \emph{trophe} = nourishment)
describe feeding position in the food web:

{\def\LTcaptype{none} % do not increment counter
\begin{longtable}[]{@{}
  >{\raggedright\arraybackslash}p{(\linewidth - 4\tabcolsep) * \real{0.3714}}
  >{\raggedright\arraybackslash}p{(\linewidth - 4\tabcolsep) * \real{0.2571}}
  >{\raggedright\arraybackslash}p{(\linewidth - 4\tabcolsep) * \real{0.3714}}@{}}
\toprule\noalign{}
\begin{minipage}[b]{\linewidth}\raggedright
Trophic level
\end{minipage} & \begin{minipage}[b]{\linewidth}\raggedright
Organisms
\end{minipage} & \begin{minipage}[b]{\linewidth}\raggedright
Energy source
\end{minipage} \\
\midrule\noalign{}
\endhead
\bottomrule\noalign{}
\endlastfoot
1 --- Producers (\hyperref[gl:autotroph]{\textbf{autotroph}}s) & Plants,
algae, cyanobacteria, chemolithotrophs & Solar energy (GPP) or chemical
energy \\
2 --- Primary consumers (herbivores) & Insects, zooplankton, ungulates,
small mammals & NPP consumed \\
3 --- Secondary consumers (carnivores) & Spiders, small fish, small
birds & Primary consumer biomass \\
4 --- Tertiary consumers & Large fish, raptors & Secondary consumer
biomass \\
Apex --- Top predators & Orca, wolf, lion, large sharks & Top-down
regulation \\
\hyperref[gl:decomposer]{\textbf{Decomposer}}/detritivore (parallel) &
Bacteria, fungi, earthworms, millipedes & Dead organic matter (detrital
pathway) \\
\end{longtable}
}

\subsubsection{Ecological Efficiency}\label{ecological-efficiency}

\begin{figure}[htbp]
\centering
\includegraphics[width=0.82\linewidth,height=4.2in,keepaspectratio]{../figures/mermaid_inline/inline_mermaid_0183_9db01333935b.png}
\caption{Flowchart for Ecological Efficiency: Solar Energy 1,000,000 kJ, Producers 10,000 kJ NPP, Herbivores 1,000 kJ, and Secondary Consumers 100 kJ form the diagram's primary path or branches.}
\end{figure}

The \textbf{trophic efficiency} (10\% rule, Lindeman 1942,
\emph{Ecology}) describes the fraction of energy at trophic level \(n\)
transferred to level \(n+1\):

\begin{equation}
\text{Trophic Efficiency} = \frac{\text{Production}_{n+1}}{\text{Production}_n} \times 100\% \approx 5\text{-}20\%
\label{eq:ecosystem_ecology_1}
\end{equation}

\textbf{Decomposition of efficiency:}

\begin{equation}
\text{Trophic efficiency} = \text{Consumption efficiency} \times \text{Assimilation efficiency} \times \text{Production efficiency}
\label{eq:ecosystem_ecology_2}
\end{equation}

{\def\LTcaptype{none} % do not increment counter
\begin{longtable}[]{@{}
  >{\raggedright\arraybackslash}p{(\linewidth - 4\tabcolsep) * \real{0.2812}}
  >{\raggedright\arraybackslash}p{(\linewidth - 4\tabcolsep) * \real{0.3125}}
  >{\raggedright\arraybackslash}p{(\linewidth - 4\tabcolsep) * \real{0.4062}}@{}}
\toprule\noalign{}
\begin{minipage}[b]{\linewidth}\raggedright
Component
\end{minipage} & \begin{minipage}[b]{\linewidth}\raggedright
Definition
\end{minipage} & \begin{minipage}[b]{\linewidth}\raggedright
Typical range
\end{minipage} \\
\midrule\noalign{}
\endhead
\bottomrule\noalign{}
\endlastfoot
\textbf{Consumption efficiency} & Fraction of available production
ingested & 20-50\% (herbivores); 50-100\% (carnivores) \\
\textbf{Assimilation efficiency} & Fraction of ingested energy absorbed
across gut & 20-60\% (herbivores on plants); 60-90\% (carnivores on
animal tissue) \\
\textbf{Production efficiency} & Fraction of assimilated energy
converted to new biomass & 1-3\% (endotherms); 20-60\% (ectotherms) \\
\end{longtable}
}

\subsubsection{Production Efficiency by Organism
Group}\label{production-efficiency-by-organism-group}

{\def\LTcaptype{none} % do not increment counter
\begin{longtable}[]{@{}
  >{\raggedright\arraybackslash}p{(\linewidth - 4\tabcolsep) * \real{0.1613}}
  >{\raggedright\arraybackslash}p{(\linewidth - 4\tabcolsep) * \real{0.6452}}
  >{\raggedright\arraybackslash}p{(\linewidth - 4\tabcolsep) * \real{0.1935}}@{}}
\toprule\noalign{}
\begin{minipage}[b]{\linewidth}\raggedright
Group
\end{minipage} & \begin{minipage}[b]{\linewidth}\raggedright
Production efficiency
\end{minipage} & \begin{minipage}[b]{\linewidth}\raggedright
Reason
\end{minipage} \\
\midrule\noalign{}
\endhead
\bottomrule\noalign{}
\endlastfoot
Insects (ectotherms) & 40-60\% & No
\hyperref[gl:thermoregulation]{\textbf{thermoregulation}} cost \\
Fish (ectotherms) & 20-30\% & Moderate metabolic costs \\
Amphibians (ectotherms) & 30-40\% & Low metabolic rate \\
Birds (endotherms) & 1-3\% & about 97\% of assimilated energy
on thermoregulation \\
Mammals (endotherms) & 1-3\% & High metabolic overhead \\
\end{longtable}
}

\textbf{Consequence for food chain length:} 10\% efficiency means energy
decreases by an order of magnitude per trophic level. Chains longer than
4-5 levels are energetically unsustainable in most ecosystems.

\textbf{Human dietary implications:} Producing 1 kg beef requires
about 7-8 kg grain and about 15,000 L water. Shifting
from trophic level 3 (carnivore) to level 2 (herbivore) reduces land and
water footprint by 5-10x. Insect \hyperref[gl:protein]{\textbf{protein}}
(40-60\% production efficiency) is about 20x more efficient
than beef protein (1-3\%).

\begin{quote}
\texorpdfstring{}{} \textbf{Clinical Connection --- Biomagnification and Human Health:}
The 10\% energy transfer rule has a dangerous corollary: persistent
pollutants that are not metabolised \textbf{biomagnify} through the food
chain. DDT (dichlorodiphenyltrichloroethane) concentrates
about 10x per trophic level. In Lake Michigan in the 1960s:
water = 0.014 ppm DDT \texorpdfstring{\ensuremath{\to}}{->}
\hyperref[gl:phytoplankton]{\textbf{phytoplankton}} = 5 ppm \texorpdfstring{\ensuremath{\to}}{->}
zooplankton = 10 ppm \texorpdfstring{\ensuremath{\to}}{->} small fish = 50 ppm \texorpdfstring{\ensuremath{\to}}{->} large fish = 200 ppm \texorpdfstring{\ensuremath{\to}}{->} bald
eagle eggs = 2,000 ppm. At 2,000 ppm, DDT metabolite DDE inhibits
carbonic anhydrase in the shell gland, producing thin eggshells \texorpdfstring{\ensuremath{\to}}{->}
reproductive failure \texorpdfstring{\ensuremath{\to}}{->} near-extinction. Mercury (methylmercury) follows
the same pattern: tuna and swordfish accumulate mercury to levels that
prompt FDA advisories for pregnant women (\textgreater0.5 ppm \texorpdfstring{\ensuremath{\to}}{->}
neurological risk to developing fetus).
\end{quote}

\begin{quote}
\textbf{Concept Check:} A savanna ecosystem has NPP = 500 g
C/m\(^2\)/yr. If herbivores consume 20\% of NPP with 10\% trophic
efficiency, and a lion (tertiary consumer) eats herbivores with 10\%
efficiency, how much energy (g C/m\(^2\)/yr) is available to sustain the
lion population?
\end{quote}

\begin{center}\rule{0.5\linewidth}{0.5pt}\end{center}

\subsection{Primary Productivity}\label{primary-productivity}

\textbf{Gross Primary Production (GPP):} Total rate of CO\(_2\) fixation
by most autotrophs (\texorpdfstring{\ensuremath{\mu}}{mu}mol CO\(_2\)/m\(^2\)/s or g C/m\(^2\)/year).

\textbf{Net Primary Production (NPP):} GPP minus autotrophic respiration
(\(R_a\)):

\begin{equation}
NPP = GPP - R_a
\label{eq:ecosystem_ecology_3}
\end{equation}

\textbf{Net Ecosystem Production (NEP):} NPP minus heterotrophic
respiration (\(R_h\), decomposers):

\begin{equation}
NEP = NPP - R_h
\label{eq:ecosystem_ecology_4}
\end{equation}

\textbf{Net \hyperref[gl:biome]{\textbf{Biome}} Production (NBP):} NEP
minus disturbance losses (fire, harvest, land-use change):

\begin{equation}
NBP = NEP - D
\label{eq:ecosystem_ecology_5}
\end{equation}

\begin{quote}
NBP determines whether an ecosystem is a net carbon \textbf{sink}
(\(NBP > 0\)) or \textbf{source} (\(NBP < 0\)).
\end{quote}

\subsubsection{The Productivity
Hierarchy}\label{the-productivity-hierarchy}

\begin{equation}
\text{Solar radiation} \xrightarrow{\sim 1\text{-}2\%} GPP \xrightarrow{-R_a (\sim 50\%)} NPP \xrightarrow{-R_h} NEP \xrightarrow{-D} NBP
\label{eq:ecosystem_ecology_6}
\end{equation}

Only about 1--2\% of incident solar energy is captured
photosynthetically as \textbf{gross primary production} (GPP).
Autotrophic respiration (\(R_a\)) consumes roughly half of GPP, leaving
\textbf{net primary production} (NPP) --- the biomass actually available
to consumers. Subtracting heterotrophic respiration (\(R_h\), by
animals, fungi, and microbes) gives \textbf{net ecosystem production}
(NEP); subtracting further non-respiratory losses such as fire, harvest,
and erosion (\(D\)) gives \textbf{net biome production} (NBP), the
quantity that determines whether a region is a long-term carbon sink or
source. Each arrow is a quantitatively large loss, which is why food
chains are short and why NBP is small relative to GPP.

\subsubsection{Global Primary
Production}\label{global-primary-production}

{\def\LTcaptype{none} % do not increment counter
\begin{longtable}[]{@{}
  >{\raggedright\arraybackslash}p{(\linewidth - 6\tabcolsep) * \real{0.1343}}
  >{\raggedright\arraybackslash}p{(\linewidth - 6\tabcolsep) * \real{0.2836}}
  >{\raggedright\arraybackslash}p{(\linewidth - 6\tabcolsep) * \real{0.2985}}
  >{\raggedright\arraybackslash}p{(\linewidth - 6\tabcolsep) * \real{0.2836}}@{}}
\toprule\noalign{}
\begin{minipage}[b]{\linewidth}\raggedright
Ecosystem
\end{minipage} & \begin{minipage}[b]{\linewidth}\raggedright
NPP (g C/m\(^2\)/yr)
\end{minipage} & \begin{minipage}[b]{\linewidth}\raggedright
Area (10\(^6\) km\(^2\))
\end{minipage} & \begin{minipage}[b]{\linewidth}\raggedright
Global contribution
\end{minipage} \\
\midrule\noalign{}
\endhead
\bottomrule\noalign{}
\endlastfoot
Tropical rainforests & 900-1,700 & 17 & about 30\% of
terrestrial NPP \\
Temperate forests & 400-900 & 12 & \\
Boreal forests & 100-400 & 12 & \\
Tropical grasslands/savanna & 200-600 & 22 & \\
Desert & 10-50 & 42 & Minimal \\
Tundra & 50-150 & 9 & \\
Open ocean & 50-150 & 332 & about 45\% of global NPP \\
Coastal/upwelling & 200-600 & 27 & \\
Coral reefs & 500-2,000 & 0.6 & High per area, small total \\
\textbf{Global total} & --- & --- & \textbf{about 120 Pg C/yr (land) +
90 Pg C/yr (ocean)} \\
\end{longtable}
}

\subsubsection{Factors Limiting GPP}\label{factors-limiting-gpp}

{\def\LTcaptype{none} % do not increment counter
\begin{longtable}[]{@{}
  >{\raggedright\arraybackslash}p{(\linewidth - 4\tabcolsep) * \real{0.2609}}
  >{\raggedright\arraybackslash}p{(\linewidth - 4\tabcolsep) * \real{0.4783}}
  >{\raggedright\arraybackslash}p{(\linewidth - 4\tabcolsep) * \real{0.2609}}@{}}
\toprule\noalign{}
\begin{minipage}[b]{\linewidth}\raggedright
Factor
\end{minipage} & \begin{minipage}[b]{\linewidth}\raggedright
Terrestrial
\end{minipage} & \begin{minipage}[b]{\linewidth}\raggedright
Marine
\end{minipage} \\
\midrule\noalign{}
\endhead
\bottomrule\noalign{}
\endlastfoot
Light & Yes (canopy shading, LAI) & Yes (depth; turbidity; latitude) \\
Temperature & Yes (\hyperref[gl:enzyme]{\textbf{enzyme}} rates; freeze)
& Yes (Arctic/Antarctic) \\
Water & Primary limitation in many biomes & N/A \\
Nitrogen & Limiting in most ecosystems & Primary limiter in most open
ocean \\
Phosphorus & Co-limiting, especially in old weathered soils &
Co-limiting \\
Iron & Not typically & Yes; 30\% of ocean is iron-limited (HNLC
regions) \\
CO\(_2\) & Potentially (CO\(_2\) fertilisation effect) & Dissolved
CO\(_2\) rarely limiting \\
\end{longtable}
}

\subsubsection{Measuring NPP}\label{measuring-npp}

Different methods access different components of NPP and apply across
different scales. No single method captures everything; the modern
practice combines them.

{\def\LTcaptype{none} % do not increment counter
\begin{longtable}[]{@{}
  >{\raggedright\arraybackslash}p{(\linewidth - 8\tabcolsep) * \real{0.1818}}
  >{\raggedright\arraybackslash}p{(\linewidth - 8\tabcolsep) * \real{0.1515}}
  >{\raggedright\arraybackslash}p{(\linewidth - 8\tabcolsep) * \real{0.2424}}
  >{\raggedright\arraybackslash}p{(\linewidth - 8\tabcolsep) * \real{0.2424}}
  >{\raggedright\arraybackslash}p{(\linewidth - 8\tabcolsep) * \real{0.1818}}@{}}
\toprule\noalign{}
\begin{minipage}[b]{\linewidth}\raggedright
Method
\end{minipage} & \begin{minipage}[b]{\linewidth}\raggedright
Scale
\end{minipage} & \begin{minipage}[b]{\linewidth}\raggedright
Approach
\end{minipage} & \begin{minipage}[b]{\linewidth}\raggedright
Captures
\end{minipage} & \begin{minipage}[b]{\linewidth}\raggedright
Misses
\end{minipage} \\
\midrule\noalign{}
\endhead
\bottomrule\noalign{}
\endlastfoot
\textbf{Harvest method} & Plot (m\(^2\)--ha) & Clip aboveground biomass
at peak; correct for litterfall, herbivory, mortality & ANPP directly &
BNPP; rapid turnover; non-peak production \\
\textbf{Eddy covariance} & Hectares (canopy footprint) & Tower-mounted
sonic anemometer + IRGA measure CO\(_2\) flux 10--20 Hz; partition NEP \texorpdfstring{\ensuremath{\to}}{->}
GPP, \(R_e\) & NEP at half-hourly resolution; 24/7/365 & Primarily the
tower footprint; tall tower needed for forests \\
\textbf{Remote sensing} & Regional / global & MODIS NDVI \texorpdfstring{\ensuremath{\to}}{->} fAPAR
(fraction absorbed PAR) \texorpdfstring{\ensuremath{\to}}{->} light-use-efficiency NPP models (e.g., MOD17)
& Continental and global scales & Hidden BNPP; saturation in dense
canopies; cloud occlusion \\
\textbf{\(^{14}\)C uptake} & Aquatic, bottle & Add NaH\(^{14}\)CO\(_3\);
incubate; filter and count \(^{14}\)C in cells & Phytoplankton primary
production & Mostly net (some respired \(^{14}\)C lost); bottle
effects \\
\textbf{O\(_2\) light/dark bottles} & Aquatic & Incubate paired bottles;
light = NPP, dark = respiration; sum = GPP & GPP and NPP simultaneously
& Sensitive to incubation duration \\
\textbf{Litterfall traps + stem increment} & Forest plot & Annual leaf
and twig fall + dendrometer band; sum = ANPP & Long-term forest dynamics
& Belowground; understory \\
\textbf{Minirhizotrons / ingrowth cores} & Plot, root-specific & Image
roots through transparent tubes or harvest from buried mesh cores & BNPP
estimates & Disturbance artefacts; coarse roots underrepresented \\
\end{longtable}
}

\textbf{Eddy covariance} is the current gold standard for
ecosystem-level carbon flux measurement. A global network of
about 800 flux towers (FLUXNET) provides continuous GPP,
\(R_e\), and NEP measurements across biomes; the data drive every modern
terrestrial carbon model and are a primary input to the IPCC reports.

\textbf{Cross-validation matters.} Eddy covariance NEP at the Harvard
Forest tower agreed with biometric (harvest + litterfall + dendrometer)
NEP within 15\% over 12 years (Barford et al.~2001, \emph{Science}) ---
independent confirmation that the flux community's numbers are real, not
artefacts of filtering or gap-filling.

\subsubsection{ANPP, BNPP, and Their
Ratio}\label{anpp-bnpp-and-their-ratio}

Total NPP partitions into \textbf{aboveground} (ANPP --- leaves, stems,
reproductive organs) and \textbf{belowground} (BNPP --- roots, root
exudates, mycorrhizal carbon transfer). Belowground production is
notoriously hard to measure but is often the larger fraction:

{\def\LTcaptype{none} % do not increment counter
\begin{longtable}[]{@{}
  >{\raggedright\arraybackslash}p{(\linewidth - 8\tabcolsep) * \real{0.1316}}
  >{\raggedright\arraybackslash}p{(\linewidth - 8\tabcolsep) * \real{0.2105}}
  >{\raggedright\arraybackslash}p{(\linewidth - 8\tabcolsep) * \real{0.2105}}
  >{\raggedright\arraybackslash}p{(\linewidth - 8\tabcolsep) * \real{0.2895}}
  >{\raggedright\arraybackslash}p{(\linewidth - 8\tabcolsep) * \real{0.1579}}@{}}
\toprule\noalign{}
\begin{minipage}[b]{\linewidth}\raggedright
Biome
\end{minipage} & \begin{minipage}[b]{\linewidth}\raggedright
ANPP/NPP
\end{minipage} & \begin{minipage}[b]{\linewidth}\raggedright
BNPP/NPP
\end{minipage} & \begin{minipage}[b]{\linewidth}\raggedright
BNPP / ANPP
\end{minipage} & \begin{minipage}[b]{\linewidth}\raggedright
Driver
\end{minipage} \\
\midrule\noalign{}
\endhead
\bottomrule\noalign{}
\endlastfoot
Tropical rainforest & 0.65--0.75 & 0.25--0.35 & 0.4 & Resource-rich;
allocation favours canopy competition \\
Temperate deciduous forest & 0.55--0.65 & 0.35--0.45 & 0.7 & Seasonal,
balanced \\
Temperate grassland & 0.30--0.40 & 0.60--0.70 & 1.5--2.5 & Drought,
fire, grazing favour root storage \\
Tundra & 0.20--0.30 & 0.70--0.80 & 3--4 & Cold soils; nutrient
acquisition expensive; root protection from freezing \\
Boreal forest & 0.40--0.55 & 0.45--0.60 & 1.0--1.3 & Slow N
mineralisation; mycorrhizal demand high \\
\end{longtable}
}

Globally, BNPP is roughly half of total NPP --- meaning \emph{half} of
the planet's photosynthate is invested below the soil surface, where
carbon turnover times can stretch into millennia. Ignoring BNPP in
carbon-stock estimates underestimates terrestrial sequestration by
30--50 \% for grasslands and tundra.

\begin{quote}
\textbf{Concept Check:} A grassland and a forest both report ANPP = 400
g C/m\(^2\)/yr. Which ecosystem has higher \emph{total} NPP, and why
does this matter for policy (e.g., carbon-credit accounting under
REDD+)?
\end{quote}

\begin{quote}
\textbf{Concept Check:} A tropical rainforest has GPP = 3,000 g
C/m\(^2\)/yr. If autotrophic respiration consumes 50\% and heterotrophic
respiration consumes 80\% of NPP, calculate NPP and NEP. Is this
ecosystem a net carbon sink?
\end{quote}

\begin{center}\rule{0.5\linewidth}{0.5pt}\end{center}

\subsection{The Carbon Cycle and Climate
Change}\label{the-carbon-cycle-and-climate-change}

\subsubsection{Carbon Reservoirs}\label{carbon-reservoirs}

{\def\LTcaptype{none} % do not increment counter
\begin{longtable}[]{@{}
  >{\raggedright\arraybackslash}p{(\linewidth - 6\tabcolsep) * \real{0.2368}}
  >{\raggedright\arraybackslash}p{(\linewidth - 6\tabcolsep) * \real{0.2895}}
  >{\raggedright\arraybackslash}p{(\linewidth - 6\tabcolsep) * \real{0.3421}}
  >{\raggedright\arraybackslash}p{(\linewidth - 6\tabcolsep) * \real{0.1316}}@{}}
\toprule\noalign{}
\begin{minipage}[b]{\linewidth}\raggedright
Reservoir
\end{minipage} & \begin{minipage}[b]{\linewidth}\raggedright
Size (Pg C)
\end{minipage} & \begin{minipage}[b]{\linewidth}\raggedright
Turnover time
\end{minipage} & \begin{minipage}[b]{\linewidth}\raggedright
Notes
\end{minipage} \\
\midrule\noalign{}
\endhead
\bottomrule\noalign{}
\endlastfoot
Atmosphere & about 860 (2024) & Years & 421 ppm CO\(_2\)
(Keeling Curve) \\
Vegetation & about 560 & Decades & Mainly tropical forests \\
Soil organic carbon & 1,500-2,400 & Decades-millennia & Largest
terrestrial reservoir \\
Permafrost & about 1,700 & Millennia & Vulnerable to thaw \\
Ocean surface & about 900 & Months & Rapid exchange with
atmosphere \\
Ocean deep & about 37,000 & Centuries-millennia & Largest
active reservoir \\
Lithosphere (fossil fuels) & about 3,700 (recoverable) &
Geological timescales & Being transferred to atmosphere \\
Marine sediments (carbonate) & \textgreater60,000,000 & Millions of
years & Long-term geological storage \\
\end{longtable}
}

\subsubsection{Annual Carbon Fluxes
(2020s)}\label{annual-carbon-fluxes-2020s}

\begin{figure}[htbp]
\centering
\includegraphics[width=0.82\linewidth,height=4.2in,keepaspectratio]{../figures/mermaid_inline/inline_mermaid_0184_2c3c7914855c.png}
\caption{Net atmospheric accumulation:. about 4.7 Pg C/yr to CO\$\_2\$ rising at about 2.4 ppm/yr.}
\end{figure}

\subsubsection{The Biological Pump}\label{the-biological-pump}

The \textbf{biological pump} transfers carbon from ocean surface to deep
water via biological processes:

\begin{enumerate}
\def\labelenumi{\arabic{enumi}.}
\tightlist
\item
  \textbf{Photosynthesis} by phytoplankton in euphotic zone fixes
  CO\(_2\) into organic carbon
\item
  \textbf{Sinking particles} --- dead phytoplankton, fecal pellets,
  aggregates (marine snow) sink to depth
\item
  \textbf{Active transport} --- zooplankton diel vertical migration
  moves carbon downward (feeding at surface at night, metabolising at
  depth during day)
\item
  \textbf{Dissolution} at depth --- organic carbon remineralised by
  bacteria, releasing CO\(_2\) into deep water
\item
  \textbf{Carbonate pump} --- CaCO\(_3\) shells (foraminifera,
  coccolithophores) sink and dissolve below the carbonate compensation
  depth (about 4,000 m)
\end{enumerate}

The biological pump exports about 10-15 Pg C/yr to deep water,
sequestering it for centuries to millennia. Without the pump,
atmospheric CO\(_2\) would be about 200 ppm higher.

\subsubsection{Key Climate Feedbacks}\label{key-climate-feedbacks}

\textbf{1. Permafrost carbon feedback (positive):} Arctic/subarctic
permafrost contains about 1,700 Pg C (twice the atmospheric
reservoir). Warming \texorpdfstring{\ensuremath{\to}}{->} thaw \texorpdfstring{\ensuremath{\to}}{->} CO\(_2\) + CH\(_4\) release \texorpdfstring{\ensuremath{\to}}{->} more warming.

\begin{itemize}
\tightlist
\item
  CH\(_4\) from \hyperref[gl:anaerobic]{\textbf{anaerobic}}
  decomposition in thermokarst lakes is 80x more potent than CO\(_2\)
  over 20 years
\item
  Estimated +1.5-2.5\(^\circ\)C additional warming if permafrost fully
  thaws (IPCC AR6 2021)
\item
  \textbf{Abrupt thaw} (collapse of ice-rich permafrost) may be more
  important than gradual thaw; largely omitted from current climate
  models (Turetsky et al.~2020, \emph{Nature Geoscience})
\end{itemize}

\textbf{2. Ocean acidification:}

\begin{equation}
CO_2 + H_2O \rightarrow H_2CO_3 \rightarrow HCO_3^- + H^+
\label{eq:ecosystem_ecology_7}
\end{equation}

Ocean \hyperref[gl:ph]{\textbf{pH}} has decreased from
about 8.2 to about 8.05 since industrialisation (30\%
increase in {[}H\(^+\){]}). Consequences:

\begin{itemize}
\tightlist
\item
  Decreased {[}CO\(_3^{2-}\){]} \texorpdfstring{\ensuremath{\to}}{->} carbonate undersaturation \texorpdfstring{\ensuremath{\to}}{->}
  dissolution of calcareous shells
\item
  \textbf{Aragonite saturation state:}
  \(\Omega_{aragonite} = \frac{[Ca^{2+}][CO_3^{2-}]}{K_{sp,aragonite}}\)
\item
  When \(\Omega < 1\): aragonite dissolves
\item
  Tropical coral reefs projected to be in net erosion by 2050 at current
  emission trajectories (Hoegh-Guldberg et al.~2007, \emph{Science})
\item
  Pteropods (sea butterflies) --- thin aragonite shells already
  dissolving in Southern Ocean
\item
  Ocean acidification also impairs fish olfaction (affecting predator
  avoidance behaviour)
\end{itemize}

\textbf{3. Methane hydrates (clathrate gun hypothesis):}
about 500-2,500 Pg C locked in methane clathrates (ice-like
structures) in ocean sediments and permafrost. Destabilisation during
rapid warming \texorpdfstring{\ensuremath{\to}}{->} mass marine CH\(_4\) release \texorpdfstring{\ensuremath{\to}}{->} catastrophic warming.
Possible role in the Paleocene-Eocene Thermal Maximum (PETM,
about 56 Ma): 5-8\(^\circ\)C warming over about 10,000
years, mass marine extinction.

\textbf{4. CO\(_2\) fertilisation effect (negative feedback,
partially):} Higher atmospheric CO\(_2\) increases photosynthetic rate
(especially in C3 plants), partially offsetting emissions. However, this
effect is limited by nitrogen and phosphorus availability, and saturates
at high CO\(_2\) levels. FACE (Free-Air CO\(_2\) Enrichment) experiments
show about 15-25\% NPP increase at 550 ppm, but diminishing
returns beyond this.

\begin{quote}
\texorpdfstring{}{} \textbf{Clinical Connection --- Ocean Acidification and Food
Security:} Ocean acidification threatens the \$100 billion/year global
shellfish industry. Oyster larvae in Pacific Northwest hatcheries
experienced mass mortality beginning about 2005, traced to
corrosive upwelled water with low \(\Omega_{aragonite}\). Hatcheries now
monitor real-time carbonate chemistry and
\hyperref[gl:buffer]{\textbf{buffer}} intake water with sodium
carbonate. Wild shellfish populations cannot be similarly managed. By
2100, oyster calcification rates are projected to decline 25-40\% under
RCP 8.5, with cascading effects on coastal economies and nutrition ---
particularly in developing nations dependent on shellfish protein.
\end{quote}

\begin{quote}
\textbf{Concept Check:} The ocean currently absorbs about 2.9
Pg C/yr from the atmosphere. If ocean warming reduces the solubility
pump by 20\%, how much additional CO\(_2\) would remain in the
atmosphere annually? What would be the approximate additional warming
contribution over 50 years?
\end{quote}

\begin{center}\rule{0.5\linewidth}{0.5pt}\end{center}

\subsection{The Nitrogen Cycle}\label{the-nitrogen-cycle}

Nitrogen (N) is often the most limiting macronutrient in terrestrial and
freshwater ecosystems. Despite comprising 78\% of the atmosphere,
\(N_2\) is biologically inert --- its triple bond (\(\equiv\)) requires
enormous energy to break.

\subsubsection{Key Nitrogen
Transformations}\label{key-nitrogen-transformations}

\begin{equation}
\text{N}_2 \xrightarrow{\text{nitrogenase}} \text{NH}_3 \xrightarrow{\text{GS/GOGAT}} \text{Organic N} \xrightarrow{\text{ammonification}} \text{NH}_4^+ \xrightarrow{\text{Nitrosomonas}} \text{NO}_2^- \xrightarrow{\text{Nitrobacter}} \text{NO}_3^- \xrightarrow{\text{denitrification}} \text{N}_2\text{O} \to \text{N}_2
\label{eq:ecosystem_ecology_8}
\end{equation}

{\def\LTcaptype{none} % do not increment counter
\begin{longtable}[]{@{}
  >{\raggedright\arraybackslash}p{(\linewidth - 8\tabcolsep) * \real{0.1842}}
  >{\raggedright\arraybackslash}p{(\linewidth - 8\tabcolsep) * \real{0.2368}}
  >{\raggedright\arraybackslash}p{(\linewidth - 8\tabcolsep) * \real{0.2105}}
  >{\raggedright\arraybackslash}p{(\linewidth - 8\tabcolsep) * \real{0.2105}}
  >{\raggedright\arraybackslash}p{(\linewidth - 8\tabcolsep) * \real{0.1579}}@{}}
\toprule\noalign{}
\begin{minipage}[b]{\linewidth}\raggedright
Process
\end{minipage} & \begin{minipage}[b]{\linewidth}\raggedright
Organisms
\end{minipage} & \begin{minipage}[b]{\linewidth}\raggedright
Reaction
\end{minipage} & \begin{minipage}[b]{\linewidth}\raggedright
Location
\end{minipage} & \begin{minipage}[b]{\linewidth}\raggedright
Oxygen
\end{minipage} \\
\midrule\noalign{}
\endhead
\bottomrule\noalign{}
\endlastfoot
\textbf{N\(_2\) fixation} & \emph{Rhizobium}, \emph{Azotobacter},
\emph{Anabaena}, \emph{Frankia} & N\(_2\) + 8H\(^+\) + 8e\(^-\) + 16ATP
\texorpdfstring{\ensuremath{\to}}{->} 2NH\(_3\) + H\(_2\) & Root nodules; soil; ocean & O\(_2\)-sensitive
(nitrogenase) \\
\textbf{Ammonification} & Most decomposers & Organic N \texorpdfstring{\ensuremath{\to}}{->} NH\(_4^+\) &
Ubiquitous & \hyperref[gl:aerobic]{\textbf{Aerobic}} or anaerobic \\
\textbf{\hyperref[gl:nitrification]{Nitrification} (step 1)} &
\emph{Nitrosomonas}, \emph{Nitrososphaera} (AOA) & NH\(_3\) \texorpdfstring{\ensuremath{\to}}{->} NO\(_2^-\)
& Aerobic soil/water & Required \\
\textbf{Nitrification (step 2)} & \emph{Nitrobacter}, \emph{Nitrospira}
& NO\(_2^-\) \texorpdfstring{\ensuremath{\to}}{->} NO\(_3^-\) & Aerobic soil/water & Required \\
\textbf{Comammox} & \emph{Nitrospira inopinata} & NH\(_3\) \texorpdfstring{\ensuremath{\to}}{->} NO\(_3^-\)
(complete) & Soil, engineered systems & Required \\
\textbf{Denitrification} & \emph{Pseudomonas}, \emph{Paracoccus} &
NO\(_3^-\) \texorpdfstring{\ensuremath{\to}}{->} N\(_2\)O \texorpdfstring{\ensuremath{\to}}{->} N\(_2\) & Waterlogged soil; sediments &
Anaerobic \\
\textbf{Anammox} & \emph{Candidatus Kuenenia} & NH\(_4^+\) + NO\(_2^-\)
\texorpdfstring{\ensuremath{\to}}{->} N\(_2\) + 2H\(_2\)O & Marine sediments & Anaerobic \\
\textbf{DNRA} & Various bacteria & NO\(_3^-\) \texorpdfstring{\ensuremath{\to}}{->} NH\(_4^+\) & Low C:N
sediments & Anaerobic \\
\end{longtable}
}

\begin{figure}[htbp]
\centering
\includegraphics[width=0.82\linewidth,height=4.2in,keepaspectratio]{../figures/mermaid_inline/inline_mermaid_0185_4fcc4cb37f87.png}
\caption{The nitrogen cycle showing biological transformations (fixation, ammonification, nitrification, denitrification) and anthropogenic disruption via the Haber-Bosch process. Excess reactive nitrogen cascades through ecosystems, causing eutrophication, hypoxia, and N2O-mediated climate warming and ozone depletion.}
\end{figure}

\subsubsection{Biological Nitrogen
Fixation}\label{biological-nitrogen-fixation}

The \textbf{nitrogenase enzyme complex} (Fe-Mo cofactor) catalyses the
most energetically expensive biological reaction:

\begin{equation}
N_2 + 8H^+ + 8e^- + 16ATP \rightarrow 2NH_3 + H_2 + 16ADP + 16P_i
\label{eq:ecosystem_ecology_9}
\end{equation}

Nitrogenase is \textbf{irreversibly inactivated by O\(_2\)}. Strategies
for O\(_2\) protection: - \textbf{Rhizobium} in legume root nodules:
leghemoglobin (O\(_2\)-scavenging pink protein) maintains low O\(_2\) -
\textbf{Cyanobacteria} (\emph{Anabaena}): specialised heterocyst cells
lack photosystem II (no O\(_2\) production) - \textbf{Azotobacter}: very
high respiration rate consumes O\(_2\) before it reaches nitrogenase

Global biological N fixation: about 130 Tg N/yr (natural) +
about 130 Tg N/yr (Haber-Bosch) = about 260 Tg N/yr
total.

\subsubsection{Anthropogenic Nitrogen
Disruption}\label{anthropogenic-nitrogen-disruption}

\textbf{Haber-Bosch process} (1909; Fritz Haber, Nobel Prize 1918):

\begin{equation}
N_2 + 3H_2 \xrightarrow{400\text{-}500°C, 150\text{-}300\text{ atm, Fe catalyst}} 2NH_3
\label{eq:ecosystem_ecology_10}
\end{equation}

This single industrial process has \textbf{doubled global reactive
nitrogen} and enabled the feeding of about 4 billion additional
people. It is arguably the most important chemical invention of the 20th
century --- and one of the most environmentally damaging.

\textbf{The reactive nitrogen cascade:} Anthropogenic N travels through
multiple compartments, causing damage at each step:

\begin{enumerate}
\def\labelenumi{\arabic{enumi}.}
\tightlist
\item
  \textbf{Fertiliser} \texorpdfstring{\ensuremath{\to}}{->} crop N (typically 30-50\% taken up by plants)
\item
  \textbf{Runoff} \texorpdfstring{\ensuremath{\to}}{->} freshwater eutrophication \texorpdfstring{\ensuremath{\to}}{->} algal blooms
\item
  \textbf{Coastal zones} \texorpdfstring{\ensuremath{\to}}{->} marine eutrophication \texorpdfstring{\ensuremath{\to}}{->} hypoxia (``dead
  zones'')
\item
  \textbf{Denitrification} \texorpdfstring{\ensuremath{\to}}{->} N\(_2\)O (potent greenhouse gas, 298x
  CO\(_2\) warming potential over 100 years)
\item
  \textbf{N\(_2\)O} \texorpdfstring{\ensuremath{\to}}{->} stratospheric ozone depletion (now the primary
  ozone-depleting substance; Ravishankara et al.~2009, \emph{Science})
\end{enumerate}

\textbf{Major hypoxic ``dead zones'':} |{} Location |{}
Area (km\(^2\)) |{} Primary nutrient source |{} |{}
-------- |{} ------------- |{} ----------------------
|{} |{} Gulf of Mexico |{} \textgreater15,000
|{} Mississippi River agricultural runoff |{} |{}
Baltic Sea |{} about 60,000 |{} European
agricultural + industrial |{} |{} Chesapeake Bay
|{} about 7,000 |{} Agricultural + urban
|{} |{} East China Sea |{} about 20,000
|{} Yangtze River |{}

\begin{quote}
\texorpdfstring{}{} \textbf{Clinical Connection --- Nitrate Contamination of Drinking
Water:} Agricultural nitrogen runoff contaminates groundwater with
nitrate (NO\(_3^-\)). In infants, gut bacteria reduce NO\(_3^-\) to
NO\(_2^-\), which oxidises hemoglobin to methemoglobin (Fe\(^{3+}\)),
incapable of carrying O\(_2\). This causes \textbf{methemoglobinemia}
(``blue baby syndrome'') --- cyanosis and potentially fatal hypoxia. The
EPA maximum contaminant level for nitrate is 10 mg/L as N. In the US
Corn Belt, \textgreater20\% of private wells exceed this limit. Chronic
low-level nitrate exposure is also epidemiologically linked to
colorectal cancer (Ward et al.~2018, \emph{Int. J. Cancer}).
\end{quote}

\begin{quote}
\textbf{Concept Check:} The Haber-Bosch process has doubled global
reactive nitrogen. Trace the fate of nitrogen applied as ammonium
fertiliser to a cornfield: what percentage is taken up by the crop, what
happens to the rest, and what environmental consequences arise at each
step?
\end{quote}

\begin{center}\rule{0.5\linewidth}{0.5pt}\end{center}

\subsection{The Phosphorus Cycle}\label{the-phosphorus-cycle}

Phosphorus has \textbf{no significant gaseous phase} --- its cycle is
sedimentary (unlike C and N):

\subsubsection{Phosphorus Reservoirs and
Fluxes}\label{phosphorus-reservoirs-and-fluxes}

{\def\LTcaptype{none} % do not increment counter
\begin{longtable}[]{@{}
  >{\raggedright\arraybackslash}p{(\linewidth - 4\tabcolsep) * \real{0.3600}}
  >{\raggedright\arraybackslash}p{(\linewidth - 4\tabcolsep) * \real{0.1600}}
  >{\raggedright\arraybackslash}p{(\linewidth - 4\tabcolsep) * \real{0.4800}}@{}}
\toprule\noalign{}
\begin{minipage}[b]{\linewidth}\raggedright
Reservoir
\end{minipage} & \begin{minipage}[b]{\linewidth}\raggedright
Size
\end{minipage} & \begin{minipage}[b]{\linewidth}\raggedright
Key features
\end{minipage} \\
\midrule\noalign{}
\endhead
\bottomrule\noalign{}
\endlastfoot
\textbf{Continental crust} & Primary source & P in phosphate rock
(apatite: Ca\(_5\)(PO\(_4\))\(_3\)(OH,F,Cl)); mined for fertiliser \\
\textbf{Soil} & Variable & P sorbed to Fe\(^{3+}\) and Al\(^{3+}\)
oxides (makes P unavailable at pH \textless{} 5 or \textgreater{} 8) \\
\textbf{Freshwater} & Very low & Dissolved inorganic P (DIP) often
\textless{} 10 \texorpdfstring{\ensuremath{\mu}}{mu}g/L \\
\textbf{Ocean} & Low & DIP at nanomolar concentrations in surface;
higher at depth (remineralisation) \\
\textbf{Marine sediments} & Very large & Long-term geological sink \\
\end{longtable}
}

\subsubsection{P Limitation in
Ecosystems}\label{p-limitation-in-ecosystems}

\begin{itemize}
\tightlist
\item
  \textbf{Freshwater lakes:} primarily P-limited (Schindler 1977 ---
  whole-lake experiment in Ontario demonstrated P controls algal growth;
  this finding led to P-detergent bans)
\item
  \textbf{Open ocean:} co-limitation by N and P; N in short-term, P in
  long-term (Redfield ratio: 106C:16N:1P in marine phytoplankton)
\item
  \textbf{Tropical soils:} severely P-limited (old weathered laterite
  soils; P leached over millions of years; mycorrhizal associations
  critical for P acquisition)
\end{itemize}

\subsubsection{The Redfield Ratio}\label{the-redfield-ratio}

\textbf{\citet{redfield1958}:} Marine phytoplankton consistently
maintain an elemental ratio of:

\begin{equation}
106C : 16N : 1P
\label{eq:ecosystem_ecology_11}
\end{equation}

This stoichiometric constraint means: - N:P ratio \textless{} 16 \texorpdfstring{\ensuremath{\to}}{->}
N-limiting - N:P ratio \textgreater{} 16 \texorpdfstring{\ensuremath{\to}}{->} P-limiting - Deviation from
Redfield ratio indicates which nutrient limits growth

\subsubsection{Phosphorus Crisis}\label{phosphorus-crisis}

Global phosphate rock reserves are concentrated in about 6
countries (Morocco controls \textgreater70\% of reserves). Depletion
projections vary widely (50-400 years at current extraction rates).
\textbf{Phosphorus recovery} from wastewater (struvite precipitation:
MgNH\(_4\)PO\(_4 \cdot 6H_2O\)) is critical for long-term food security.
Unlike nitrogen, phosphorus cannot be synthesised --- it must be mined
or recycled.

\begin{quote}
\textbf{Concept Check:} Schindler's whole-lake experiment added P to one
half and N+C to the other half. Primarily the P-enriched half developed
algal blooms. Why is P the primary limiting nutrient in freshwater but N
is often more limiting in the ocean?
\end{quote}

\begin{center}\rule{0.5\linewidth}{0.5pt}\end{center}

\subsection{The Sulphur Cycle}\label{the-sulphur-cycle}

Sulphur (S) cycling connects biological, atmospheric, and geological
processes:

Sulphur is a useful corrective to overly simple nutrient-cycle diagrams
because it links metabolism, redox gradients, aerosols, odours, mining,
acid rain, and ocean-atmosphere exchange. In sediments, sulphate
reduction can dominate anaerobic respiration after oxygen and nitrate
are depleted; at oxic-anoxic interfaces, sulphide oxidisers recycle
reduced sulphur back toward sulphate. The biological interpretation
therefore depends on redox state, electron donors, pH, and whether the
system is microbial mat, wetland, hydrothermal vent, soil, or ocean
surface.

\subsubsection{Key Sulphur
Transformations}\label{key-sulphur-transformations}

{\def\LTcaptype{none} % do not increment counter
\begin{longtable}[]{@{}
  >{\raggedright\arraybackslash}p{(\linewidth - 6\tabcolsep) * \real{0.2000}}
  >{\raggedright\arraybackslash}p{(\linewidth - 6\tabcolsep) * \real{0.2571}}
  >{\raggedright\arraybackslash}p{(\linewidth - 6\tabcolsep) * \real{0.2286}}
  >{\raggedright\arraybackslash}p{(\linewidth - 6\tabcolsep) * \real{0.3143}}@{}}
\toprule\noalign{}
\begin{minipage}[b]{\linewidth}\raggedright
Process
\end{minipage} & \begin{minipage}[b]{\linewidth}\raggedright
Organisms
\end{minipage} & \begin{minipage}[b]{\linewidth}\raggedright
Reaction
\end{minipage} & \begin{minipage}[b]{\linewidth}\raggedright
Environment
\end{minipage} \\
\midrule\noalign{}
\endhead
\bottomrule\noalign{}
\endlastfoot
\textbf{Assimilatory S reduction} & Plants, bacteria & SO\(_4^{2-}\) \texorpdfstring{\ensuremath{\to}}{->}
organic S (cysteine, methionine) & Ubiquitous \\
\textbf{Decomposition} & Decomposers & Organic S \texorpdfstring{\ensuremath{\to}}{->} H\(_2\)S & Anaerobic
sediments \\
\textbf{Dissimilatory S reduction} & \emph{Desulfovibrio},
\emph{Desulfobacter} & SO\(_4^{2-}\) + H\(_2\) \texorpdfstring{\ensuremath{\to}}{->} H\(_2\)S & Anaerobic
(below N-reduction zone) \\
\textbf{S oxidation} & \emph{Thiobacillus}, \emph{Beggiatoa} & H\(_2\)S
\texorpdfstring{\ensuremath{\to}}{->} S\(^0\) \texorpdfstring{\ensuremath{\to}}{->} SO\(_4^{2-}\) & Chemolithotrophic; oxic-anoxic interface \\
\textbf{DMS production} & Marine phytoplankton & DMSP \texorpdfstring{\ensuremath{\to}}{->} DMS
(dimethylsulphide) & Ocean surface \\
\end{longtable}
}

\subsubsection{The CLAW Hypothesis}\label{the-claw-hypothesis}

\textbf{Charlson-Lovelock-Andreae-Warren (1987):} Marine phytoplankton
emit DMS \texorpdfstring{\ensuremath{\to}}{->} DMS oxidises to SO\(_4^{2-}\) aerosols in the atmosphere \texorpdfstring{\ensuremath{\to}}{->}
these aerosols serve as \textbf{cloud condensation nuclei (CCN)} \texorpdfstring{\ensuremath{\to}}{->} more
clouds \texorpdfstring{\ensuremath{\to}}{->} increased albedo \texorpdfstring{\ensuremath{\to}}{->} cooling. This represents a potential
\textbf{planetary thermostat} (negative feedback).

However, the strength of this feedback remains debated. Satellite
studies show correlations between phytoplankton blooms and cloud
properties, but the magnitude of the DMS-climate link is uncertain.

\subsubsection{Acid Rain}\label{acid-rain}

Fossil fuel combustion releases SO\(_2\) and NO\(_x\):

\begin{equation}
SO_2 + H_2O + \frac{1}{2}O_2 \rightarrow H_2SO_4
\label{eq:ecosystem_ecology_12}
\end{equation}

\begin{equation}
2NO_2 + H_2O \rightarrow HNO_3 + HNO_2
\label{eq:ecosystem_ecology_13}
\end{equation}

\textbf{Environmental effects:} pH \textless{} 4.2 in sensitive lakes \texorpdfstring{\ensuremath{\to}}{->}
fish kills, forest dieback (1970s-1990s). The success of sulphur
emission regulations (US Clean Air Act 1990) in reversing acid rain
damage demonstrates that environmental recovery is possible with
political action --- a major environmental success story.

\begin{center}\rule{0.5\linewidth}{0.5pt}\end{center}

\subsection{7B Nutrient Cycling Models: Open vs.~Closed
Systems}\label{b-nutrient-cycling-models-open-vs.-closed-systems}

Whether an ecosystem is a net \textbf{sink} or \textbf{source} for a
nutrient depends on the relative rates of input, internal recycling, and
output (leakage). Two idealised limits are useful:

{\def\LTcaptype{none} % do not increment counter
\begin{longtable}[]{@{}
  >{\raggedright\arraybackslash}p{(\linewidth - 8\tabcolsep) * \real{0.1136}}
  >{\raggedright\arraybackslash}p{(\linewidth - 8\tabcolsep) * \real{0.1364}}
  >{\raggedright\arraybackslash}p{(\linewidth - 8\tabcolsep) * \real{0.1591}}
  >{\raggedright\arraybackslash}p{(\linewidth - 8\tabcolsep) * \real{0.4091}}
  >{\raggedright\arraybackslash}p{(\linewidth - 8\tabcolsep) * \real{0.1818}}@{}}
\toprule\noalign{}
\begin{minipage}[b]{\linewidth}\raggedright
Model
\end{minipage} & \begin{minipage}[b]{\linewidth}\raggedright
Inputs
\end{minipage} & \begin{minipage}[b]{\linewidth}\raggedright
Outputs
\end{minipage} & \begin{minipage}[b]{\linewidth}\raggedright
Internal recycling
\end{minipage} & \begin{minipage}[b]{\linewidth}\raggedright
Examples
\end{minipage} \\
\midrule\noalign{}
\endhead
\bottomrule\noalign{}
\endlastfoot
\textbf{Closed} & \(\to 0\) & \(\to 0\) & Dominates & Mature tropical
rainforest on ancient soil; coral reef interior \\
\textbf{Open} & Large & Large & Small relative to fluxes & Floodplain
wetland; nutrient-loaded estuary; agricultural field \\
\end{longtable}
}

For a single nutrient pool \(X\) (e.g., soil-bound P, kg/ha) with input
rate \(I\) (e.g., atmospheric deposition + weathering) and output rate
\(O\) (leaching + erosion + harvest):

\begin{equation}
\frac{dX}{dt} = I - O = I - k X
\label{eq:ecosystem_ecology_nutrient_pool}
\end{equation}

where \(k\) is the leakage rate (yr\(^{-1}\)). At steady state
(\(dX/dt = 0\)): \(\hat X = I/k\), and the \textbf{mean residence time}
is \(\tau = 1/k = X/O\). A closed system has \(\tau \gg 1\) (nutrients
cycle internally many times before leaving); an open system has
\(\tau \le 1\) (nutrients flow through quickly).

\subsubsection{Hubbard Brook: The Watershed That Defined the
Field}\label{hubbard-brook-the-watershed-that-defined-the-field}

The \textbf{Hubbard Brook Experimental Forest} (New Hampshire, USA)
hosts the longest-running ecosystem-scale watershed experiment in the
world \citep{bormann1967}. The design is brilliant: a small forested
watershed has a single stream outlet, so weighing rainfall inputs
against streamwater outputs gives the entire ecosystem nutrient budget
--- no plot extrapolation needed.

In 1965--1966, Bormann, Likens, and colleagues clear-cut Watershed 2 and
suppressed regrowth with herbicides for three years, then compared
inputs and outputs against an undisturbed reference watershed. The
results overturned the dogma that mature forests ``leak'' nutrients in
proportion to inputs:

{\def\LTcaptype{none} % do not increment counter
\begin{longtable}[]{@{}
  >{\raggedright\arraybackslash}p{(\linewidth - 6\tabcolsep) * \real{0.3692}}
  >{\raggedright\arraybackslash}p{(\linewidth - 6\tabcolsep) * \real{0.2769}}
  >{\raggedright\arraybackslash}p{(\linewidth - 6\tabcolsep) * \real{0.2769}}
  >{\raggedright\arraybackslash}p{(\linewidth - 6\tabcolsep) * \real{0.0769}}@{}}
\toprule\noalign{}
\begin{minipage}[b]{\linewidth}\raggedright
Stream output (kg/ha/yr)
\end{minipage} & \begin{minipage}[b]{\linewidth}\raggedright
Reference (intact)
\end{minipage} & \begin{minipage}[b]{\linewidth}\raggedright
Clear-cut (year 2)
\end{minipage} & \begin{minipage}[b]{\linewidth}\raggedright
Ratio
\end{minipage} \\
\midrule\noalign{}
\endhead
\bottomrule\noalign{}
\endlastfoot
\textbf{NO\(_3^-\)--N} & 2.0 & 53 & 26\texorpdfstring{\ensuremath{\times}}{x} \\
\textbf{Ca\(^{2+}\)} & 14 & 78 & 5.6\texorpdfstring{\ensuremath{\times}}{x} \\
\textbf{K\(^+\)} & 1.7 & 36 & 21\texorpdfstring{\ensuremath{\times}}{x} \\
\textbf{Streamflow (cm/yr)} & 73 & 100 & 1.4\texorpdfstring{\ensuremath{\times}}{x} (no transpiration) \\
\end{longtable}
}

The clear-cut watershed lost more N in two years than it had accumulated
over decades. The lessons:

\begin{enumerate}
\def\labelenumi{\arabic{enumi}.}
\tightlist
\item
  \textbf{Live vegetation is the single largest control on nutrient
  retention.} Roots take up dissolved N-\textgreater{} build biomass
  -\textgreater{} N stays in the system.
\item
  \textbf{Transpiration moves water out of soils} and prevents
  nitrate-leaching events. Removing it raises water tables and flushes
  nutrients to streams.
\item
  \textbf{Mature forests are quasi-closed for most nutrients}, retaining
  80--95 \% of inputs through tight biological cycling.
\item
  The same logic explains why agricultural watersheds (no perennial
  deep-rooted vegetation) leak nitrate and drive hypoxic dead zones in
  downstream receiving waters.
\end{enumerate}

Hubbard Brook also detected \textbf{acid rain} in the 1970s by tracking
declining streamwater Ca\(^{2+}\) concentration over decades; that
signal contributed to the science underlying the U.S. Clean Air Act
amendments of 1990. Few ecosystem experiments have had comparable
scientific or policy impact.

\begin{quote}
\textbf{Concept Check:} Why does a clear-cut forest lose 20\texorpdfstring{\ensuremath{\times}}{x} more
nitrate in streamflow than an intact forest, even though the standing
soil-N pool is the same? Trace the mechanism through (a) plant uptake,
(b) decomposition rate, (c) hydrology.
\end{quote}

\begin{center}\rule{0.5\linewidth}{0.5pt}\end{center}

\subsection{7C Ocean Ecosystem Ecology}\label{c-ocean-ecosystem-ecology}

Marine ecosystems contribute roughly half of global NPP
(about 50 Pg C/yr) on just two-thirds of the planet's surface
--- but the productivity is enormously heterogeneous, structured by
light, nutrients, and physical mixing.

\subsubsection{Productivity Zones of the
Ocean}\label{productivity-zones-of-the-ocean}

{\def\LTcaptype{none} % do not increment counter
\begin{longtable}[]{@{}
  >{\raggedright\arraybackslash}p{(\linewidth - 8\tabcolsep) * \real{0.1000}}
  >{\raggedright\arraybackslash}p{(\linewidth - 8\tabcolsep) * \real{0.1250}}
  >{\raggedright\arraybackslash}p{(\linewidth - 8\tabcolsep) * \real{0.1250}}
  >{\raggedright\arraybackslash}p{(\linewidth - 8\tabcolsep) * \real{0.4750}}
  >{\raggedright\arraybackslash}p{(\linewidth - 8\tabcolsep) * \real{0.1750}}@{}}
\toprule\noalign{}
\begin{minipage}[b]{\linewidth}\raggedright
Zone
\end{minipage} & \begin{minipage}[b]{\linewidth}\raggedright
Depth
\end{minipage} & \begin{minipage}[b]{\linewidth}\raggedright
Light
\end{minipage} & \begin{minipage}[b]{\linewidth}\raggedright
NPP (g C/m\(^2\)/yr)
\end{minipage} & \begin{minipage}[b]{\linewidth}\raggedright
Drivers
\end{minipage} \\
\midrule\noalign{}
\endhead
\bottomrule\noalign{}
\endlastfoot
\textbf{Euphotic} & 0--100 m & \textgreater{} 1 \% surface PAR & 50--600
& Phytoplankton growth; temperature- and nutrient-limited \\
\textbf{Mesopelagic (twilight)} & 100--1,000 m & \textless{} 1 \% & 0
(heterotrophic) & Diel vertical migrators; remineralisation \\
\textbf{Bathypelagic / abyssopelagic} & 1,000--6,000 m & None & 0
(heterotrophic + chemo) & Marine snow rains down; deep-sea food webs \\
\textbf{Coastal upwelling} & 0--200 m & High & 500--2,000 & Wind-driven
upwelling brings cold, nutrient-rich water (Peruvian, Benguela,
Californian, Canary currents) \\
\textbf{Subtropical gyres (``ocean deserts'')} & 0--200 m & High &
25--100 & Permanent thermal stratification; nutrients trapped at
depth \\
\textbf{Polar oceans} & 0--200 m (seasonal) & Strong seasonality &
50--300 (intense in summer) & Diatom blooms after ice retreat \\
\textbf{Continental shelves} & 0--200 m & High & 100--600 & River
nutrient inputs; tidal mixing \\
\end{longtable}
}

The four major \textbf{coastal upwelling systems} (eastern boundaries of
ocean basins) cover \textless{} 1 \% of ocean area but produce
about 20 \% of global fisheries catch --- a dramatic
illustration of how local mixing physics can dwarf area-based
productivity expectations.

\subsubsection{Biological Pump and Its
Efficiency}\label{biological-pump-and-its-efficiency}

The biological pump operates at three scales:

\begin{equation}
\text{Export production} = \text{NPP} \times \text{e-ratio}
\label{eq:ecosystem_ecology_export_ratio}
\end{equation}

where the \textbf{e-ratio} is the fraction of NPP that escapes the
surface layer as sinking particles. Typical values: 0.1 in oligotrophic
gyres (most carbon is recycled in surface), 0.3--0.5 in upwelling zones
(large diatoms sink fast), and 0.5+ during diatom blooms in polar
oceans.

\begin{equation}
\text{Sequestration efficiency} = \exp(-z/z_{1/2})
\label{eq:ecosystem_ecology_martin_curve}
\end{equation}

The \textbf{Martin curve} (Martin et al.~1987, \emph{Deep-Sea Research})
describes how export flux attenuates with depth as bacteria remineralise
sinking organic matter. Carbon that escapes below about 1,000 m
is sequestered for centuries to millennia; carbon remineralised in the
upper ocean returns to the atmosphere within years.

Climate change is reorganising both pumps. Warming \textbf{stratifies
the surface ocean}, reducing nutrient supply from depth and lowering
NPP; acidification reduces calcareous-shell production, weakening the
carbonate pump; and oxygen minimum zones are expanding, shifting
remineralisation pathways. The net effect on the future ocean carbon
sink is one of the largest open questions in Earth-system science.

\begin{quote}
\texorpdfstring{}{} \textbf{Clinical Connection --- Harmful Algal Blooms.}
Nutrient-loaded coastal waters increasingly experience \textbf{harmful
algal blooms} (HABs) by \emph{Karenia brevis} (red tide, Florida ---
produces brevetoxins \texorpdfstring{\ensuremath{\to}}{->} respiratory distress, fish kills),
\emph{Pseudo-nitzschia} (Pacific coast --- produces domoic acid \texorpdfstring{\ensuremath{\to}}{->}
amnesic shellfish poisoning, lethal in birds and humans), and
cyanobacterial blooms in Lake Erie (microcystin \texorpdfstring{\ensuremath{\to}}{->} liver toxicity; the
2014 Toledo, Ohio drinking-water shutdown affected 500,000 people). HABs
are open-system biogeochemistry meeting public health: agricultural
runoff loads N and P to coastal waters, warm summers stratify and warm
the surface, and the resulting blooms produce neurotoxins that move up
the food chain via shellfish and finfish to humans. Each of the carbon,
nitrogen, and phosphorus cycle disruptions in this chapter contributes.
\end{quote}

\begin{center}\rule{0.5\linewidth}{0.5pt}\end{center}

\subsection{7D Soil Formation and
Pedogenesis}\label{d-soil-formation-and-pedogenesis}

Soils are the central terrestrial reservoir for organic carbon
(1,500--2,400 Pg C, more than vegetation and atmosphere combined) and
the medium through which nutrient cycles are mediated.
\textbf{Pedogenesis} (soil formation) is itself an ecosystem-scale
process operating over centuries to millennia.

\subsubsection{Five State Factors (Jenny
1941)}\label{five-state-factors-jenny-1941}

\begin{equation}
\text{Soil} = f\,(\text{climate, organisms, relief, parent material, time})
\label{eq:ecosystem_ecology_clorpt}
\end{equation}

This \textbf{CLORPT} framework remains the foundational organising
principle of pedology. Climate sets the rates of weathering and
decomposition; organisms (especially plants and microbes) drive organic
matter accumulation; relief (topography) routes water and erosion;
parent material (bedrock) supplies the mineralogy; and time integrates
everything.

\subsubsection{Soil Horizons}\label{soil-horizons}

A mature soil typically displays vertical layering reflecting
depth-dependent processes:

{\def\LTcaptype{none} % do not increment counter
\begin{longtable}[]{@{}
  >{\raggedright\arraybackslash}p{(\linewidth - 6\tabcolsep) * \real{0.2188}}
  >{\raggedright\arraybackslash}p{(\linewidth - 6\tabcolsep) * \real{0.1562}}
  >{\raggedright\arraybackslash}p{(\linewidth - 6\tabcolsep) * \real{0.2812}}
  >{\raggedright\arraybackslash}p{(\linewidth - 6\tabcolsep) * \real{0.3438}}@{}}
\toprule\noalign{}
\begin{minipage}[b]{\linewidth}\raggedright
Horizon
\end{minipage} & \begin{minipage}[b]{\linewidth}\raggedright
Depth
\end{minipage} & \begin{minipage}[b]{\linewidth}\raggedright
Processes
\end{minipage} & \begin{minipage}[b]{\linewidth}\raggedright
Composition
\end{minipage} \\
\midrule\noalign{}
\endhead
\bottomrule\noalign{}
\endlastfoot
\textbf{O (organic)} & Top 0--10 cm & Litter accumulation; decomposition
& Leaf litter; partially decomposed organic matter \\
\textbf{A (topsoil)} & 5--30 cm & Bioturbation; humification; root
activity & Dark, organic-matter-rich; high biological activity \\
\textbf{E (eluvial)} & 0--50 cm (boreal/podzol) & Leaching of clays, Fe,
Al downward & Pale, depleted; sandy texture \\
\textbf{B (subsoil)} & 30--150 cm & Illuviation: accumulation of leached
materials & Clay, Fe, Al oxides; often reddish or yellowish \\
\textbf{C (parent material)} & Below B & Slow weathering of bedrock &
Rock fragments transitioning to soil \\
\textbf{R (bedrock)} & Bottom & None & Unweathered rock \\
\end{longtable}
}

\subsubsection{Three Stages of
Pedogenesis}\label{three-stages-of-pedogenesis}

\begin{enumerate}
\def\labelenumi{\arabic{enumi}.}
\tightlist
\item
  \textbf{Weathering and primary mineral breakdown.} Physical
  fragmentation and chemical hydrolysis of bedrock release mineral
  nutrients (Ca, Mg, K, P, Fe). Cyanobacteria and lichens accelerate
  weathering through carbonic and oxalic acid secretion. Net rate:
  0.01--0.1 mm/yr (slow!).
\item
  \textbf{Organic matter accumulation.} Pioneer plants colonise; their
  litter fuels decomposers; humus forms (recalcitrant polymeric organic
  matter). Net soil organic carbon (SOC) increases at 5--50 g
  C/m\(^2\)/yr in temperate ecosystems for centuries.
\item
  \textbf{Horizon differentiation and steady-state.} Vertical
  translocation of clays, Fe/Al, and dissolved organic matter creates
  the A--E--B horizon sequence. Eventually inputs balance losses
  (leaching + decomposition + erosion) and SOC reaches a
  climate-determined ceiling.
\end{enumerate}

Tropical soils on ancient Gondwanan parent material (Brazilian shield,
Australian outback) are at stage 3 \emph{and ancient} --- most P has
leached over millions of years, leaving extremely low-fertility
laterites; this explains why tropical rainforest productivity is high
(rapid recycling) but cleared tropical soils are infertile (no
recyclable pool left). Glacier-Bay primary succession
(\cref{sec:unit_X_community_ecology}) is the same process running at the
centennial timescale.

\subsubsection{Soil Carbon and Climate
Feedback}\label{soil-carbon-and-climate-feedback}

Warming accelerates microbial decomposition more than it accelerates
plant productivity, so a warmer world is expected to release soil carbon
to the atmosphere --- a positive feedback estimated at +30 to +200 Pg C
over the 21st century (Crowther et al.~2016, \emph{Nature}). Permafrost
soils contain about 1,700 Pg C, twice the atmospheric pool, and
are one of the largest uncertainties in 21st-century climate
trajectories.

\begin{center}\rule{0.5\linewidth}{0.5pt}\end{center}

\subsection{7E Coupled Biogeochemical Cycles and the Redfield
Stoichiometry}\label{e-coupled-biogeochemical-cycles-and-the-redfield-stoichiometry}

The C, N, and P cycles do not operate independently --- biological
demand stitches them into a coupled system whose stoichiometry is one of
the most powerful diagnostic tools in ecosystem science.

\subsubsection{The Redfield Ratio
Revisited}\label{the-redfield-ratio-revisited}

Marine phytoplankton consistently maintain an elemental composition of
\(106\,\text{C} : 16\,\text{N} : 1\,\text{P}\) \citep{redfield1958}.
This ratio also describes the average composition of dissolved nutrients
in deep ocean water --- strong evidence that \textbf{biology actively
controls ocean chemistry} (the Redfield insight). Deviations from
Redfield diagnose limitation:

{\def\LTcaptype{none} % do not increment counter
\begin{longtable}[]{@{}ll@{}}
\toprule\noalign{}
Ambient N : P & Implication \\
\midrule\noalign{}
\endhead
\bottomrule\noalign{}
\endlastfoot
\textless{} 16 (e.g., 8) & N-limiting; nitrogen-fixers favoured \\
\texorpdfstring{\ensuremath{\approx}}{~} 16 & Balanced; both potentially limiting \\
\textgreater{} 16 (e.g., 30) & P-limiting; high-P-affinity species
favoured \\
\end{longtable}
}

Terrestrial systems show similar but more variable ratios reflecting
differences between woody (high C : N : P) and herbaceous tissues.
Forest canopy: about \(1200 : 28 : 1\). Soil microbial
biomass: about \(60 : 7 : 1\) (N- and P-rich relative to
plants --- explaining why microbes immobilise mineral nutrients).

\subsubsection{Stoichiometric Coupling}\label{stoichiometric-coupling}

Ecosystem-scale stoichiometry generates predictable feedbacks:

\begin{itemize}
\tightlist
\item
  \textbf{CO\(_2\) fertilisation has limits.} Higher atmospheric
  CO\(_2\) raises plant C : N ratios; without proportional N input,
  growth saturates (the \textbf{progressive nitrogen limitation}
  hypothesis confirmed by long-term FACE experiments).
\item
  \textbf{N deposition has P costs.} Anthropogenic N enrichment shifts
  ecosystems from N- to P-limited; soil and freshwater P demand rises
  sharply.
\item
  \textbf{Anoxia shifts cycles.} Hypoxic dead zones suppress
  nitrate-dependent respiration pathways but enable iron and sulphate
  reduction; the N : P : Fe : S coupling reorganises.
\end{itemize}

\subsubsection{Planetary Boundaries}\label{planetary-boundaries}

\citet{richardson2023earth} and the Stockholm Resilience Centre frame
anthropogenic biogeochemical disruption as transgression of
\textbf{planetary boundaries}. In the current assessment, six of nine
boundaries are crossed: climate change, biosphere integrity,
biogeochemical flows (N and P), land-system change, freshwater change,
and novel entities (chemical pollution). Stratospheric ozone remains a
regulatory success story, while ocean acidification is close to its
boundary. Ecosystem ecology provides the quantitative basis for each of
these limits --- and for designing interventions to return inside them.

\begin{center}\rule{0.5\linewidth}{0.5pt}\end{center}

\subsection{7F Ecosystem Services Quantification and Earth System
Science}\label{f-ecosystem-services-quantification-and-earth-system-science}

\subsubsection{Quantifying and Valuing
Services}\label{quantifying-and-valuing-services}

Beyond the opening qualitative categories, ecosystem services are
increasingly quantified using three valuation approaches:

{\def\LTcaptype{none} % do not increment counter
\begin{longtable}[]{@{}
  >{\raggedright\arraybackslash}p{(\linewidth - 6\tabcolsep) * \real{0.1765}}
  >{\raggedright\arraybackslash}p{(\linewidth - 6\tabcolsep) * \real{0.2353}}
  >{\raggedright\arraybackslash}p{(\linewidth - 6\tabcolsep) * \real{0.2647}}
  >{\raggedright\arraybackslash}p{(\linewidth - 6\tabcolsep) * \real{0.3235}}@{}}
\toprule\noalign{}
\begin{minipage}[b]{\linewidth}\raggedright
Method
\end{minipage} & \begin{minipage}[b]{\linewidth}\raggedright
Approach
\end{minipage} & \begin{minipage}[b]{\linewidth}\raggedright
Strengths
\end{minipage} & \begin{minipage}[b]{\linewidth}\raggedright
Limitations
\end{minipage} \\
\midrule\noalign{}
\endhead
\bottomrule\noalign{}
\endlastfoot
\textbf{Market price} & Use observed prices for traded ecosystem outputs
(timber, fish) & Direct; defensible & Primarily captures provisioning
services \\
\textbf{Replacement cost} & Estimate cost of engineering substitutes
(water treatment plants for wetland filtration) & Tangible & Assumes
substitutes exist; ignores non-substitutables \\
\textbf{Willingness-to-pay} & Survey residents on what they would pay
for clean air, biodiversity, etc. & Captures cultural and existence
value & Hypothetical bias; income-dependent \\
\end{longtable}
}

Costanza et al.~(2014, \emph{Ecosyst. Serv.}) updated their landmark
1997 estimate: global ecosystem services are worth about \$125
trillion/yr (vs.~global GDP of about \$100 trillion/yr in
2024). The TEEB initiative (The Economics of Ecosystems and
Biodiversity, 2010) extended these methods to corporate and government
decision-making. Ecosystem-service trade-offs (e.g., cropland conversion
gains provisioning food but loses regulating climate and water
purification) are now standard in environmental impact assessments.

\subsubsection{The Gaia Hypothesis: Critical
Evaluation}\label{the-gaia-hypothesis-critical-evaluation}

\citet{lovelock1974atmospheric} proposed the \textbf{Gaia hypothesis}:
life and the abiotic environment form a coupled, self-regulating system.
Modern Earth System Science keeps the coupling but rejects teleology;
the key question is which feedbacks stabilise the Earth system and which
amplify perturbations. Core claims and current standing:

{\def\LTcaptype{none} % do not increment counter
\begin{longtable}[]{@{}
  >{\raggedright\arraybackslash}p{(\linewidth - 2\tabcolsep) * \real{0.6250}}
  >{\raggedright\arraybackslash}p{(\linewidth - 2\tabcolsep) * \real{0.3750}}@{}}
\toprule\noalign{}
\begin{minipage}[b]{\linewidth}\raggedright
Gaia claim
\end{minipage} & \begin{minipage}[b]{\linewidth}\raggedright
Status
\end{minipage} \\
\midrule\noalign{}
\endhead
\bottomrule\noalign{}
\endlastfoot
Atmospheric composition (O\(_2\), CH\(_4\), N\(_2\)O) is biologically
driven & \textbf{Strongly supported} --- biology is essential \\
Earth's climate is biologically buffered (CLAW, biological pump,
vegetation albedo) & \textbf{Partially supported} --- feedbacks exist
but are not always stabilising \\
Biota act \emph{teleologically} to maintain habitability &
\textbf{Rejected} --- no mechanism for group/planetary selection
(Doolittle 1981); selection acts on individuals \\
The biosphere is a \textbf{superorganism} & \textbf{Metaphorical} ---
useful for systems thinking but not a literal claim \\
\end{longtable}
}

The modern synthesis treats Gaia as \textbf{Earth System Science}:
biology, geochemistry, hydrology, and atmospheric science form a single
coupled system, with feedbacks that can be stabilising (negative) or
destabilising (positive) depending on the perturbation. The IPCC
reports, the Anthropocene Working Group, and the Planetary Boundaries
framework are the descendants of this line of thinking.

Per \citet{dobzhansky1973}'s dictum that ``nothing in biology makes
sense except in the light of evolution,'' nothing in 21st-century
ecology makes sense except in the light of human-driven Earth system
change. The ecosystem ecology of this chapter --- energy flow,
biogeochemical cycles, ecosystem services --- is now the operating
manual for managing the planet itself.

\begin{quote}
\textbf{Concept Check:} The Gaia hypothesis predicts that the biosphere
stabilises Earth's habitability. List two examples where biological
feedbacks stabilise the climate system and two where they destabilise
it. What does this tell you about the difference between ``Gaia as
superorganism'' and ``Earth as a coupled biogeochemical system''?
\end{quote}

\begin{center}\rule{0.5\linewidth}{0.5pt}\end{center}

\subsection{The Water Cycle (Hydrological
Cycle)}\label{the-water-cycle-hydrological-cycle}

The water cycle connects most biogeochemical cycles and drives nutrient
transport:

For biology, the water cycle is also a vegetation and land-use cycle.
Roots, stomata, leaf area, soil organic matter, and microbial crusts
influence infiltration, transpiration, runoff, and groundwater recharge;
deforestation, drainage, compaction, and irrigation alter those same
fluxes. Climate-change impacts on drought, flood risk, and food security
therefore cannot be read from precipitation totals alone; residence
time, timing, storage, and plant access to water matter just as much.

\subsubsection{Key Fluxes}\label{key-fluxes}

{\def\LTcaptype{none} % do not increment counter
\begin{longtable}[]{@{}
  >{\raggedright\arraybackslash}p{(\linewidth - 4\tabcolsep) * \real{0.1707}}
  >{\raggedright\arraybackslash}p{(\linewidth - 4\tabcolsep) * \real{0.5610}}
  >{\raggedright\arraybackslash}p{(\linewidth - 4\tabcolsep) * \real{0.2683}}@{}}
\toprule\noalign{}
\begin{minipage}[b]{\linewidth}\raggedright
Process
\end{minipage} & \begin{minipage}[b]{\linewidth}\raggedright
Rate (10\(^3\) km\(^3\)/yr)
\end{minipage} & \begin{minipage}[b]{\linewidth}\raggedright
Description
\end{minipage} \\
\midrule\noalign{}
\endhead
\bottomrule\noalign{}
\endlastfoot
\textbf{Evaporation} (ocean) & 434 & Largest single flux \\
\textbf{Precipitation} (ocean) & 398 & \\
\textbf{Evapotranspiration} (land) & 71 & Plant
\hyperref[gl:transpiration]{\textbf{transpiration}} = 60-90\% of
terrestrial evaporation \\
\textbf{Precipitation} (land) & 107 & \\
\textbf{Runoff} (rivers to ocean) & 36 & Balances net ocean\texorpdfstring{\ensuremath{\to}}{->}land
atmospheric transport \\
\end{longtable}
}

\textbf{Plant transpiration} is a massively underappreciated flux: a
single large oak tree transpires about 400 L/day. Tropical
forests generate about 50\% of their own rainfall through
transpiration recycling (the ``flying rivers'' of the Amazon).
Deforestation disrupts this cycle, potentially triggering tipping points
toward savannification.

\begin{center}\rule{0.5\linewidth}{0.5pt}\end{center}

\subsection{Ecosystem Services and Conservation
Economics}\label{ecosystem-services-and-conservation-economics}

\subsubsection{Biodiversity-Ecosystem Functioning
(BEF)}\label{biodiversity-ecosystem-functioning-bef}

\textbf{Tilman et al.~(1996, 2001):} Species richness generally
increases ecosystem productivity, stability, and nutrient retention.

Mechanisms: - \textbf{Complementarity:} different species use different
resource pools/times; \hyperref[gl:niche]{\textbf{niche}} partitioning
increases total resource capture - \textbf{Sampling effect:} more
diverse communities more likely to contain highly productive species -
\textbf{Facilitation:} some species improve conditions for others
(nitrogen fixers enrich soil for grasses)

\textbf{The insurance hypothesis} \citep{yachi1999}: Biodiversity
provides insurance against environmental fluctuations. In diverse
communities, when one species declines due to unfavourable conditions,
others compensate. This stabilises ecosystem function over time.

\subsubsection{Payment for Ecosystem Services
(PES)}\label{payment-for-ecosystem-services-pes}

{\def\LTcaptype{none} % do not increment counter
\begin{longtable}[]{@{}
  >{\raggedright\arraybackslash}p{(\linewidth - 4\tabcolsep) * \real{0.3333}}
  >{\raggedright\arraybackslash}p{(\linewidth - 4\tabcolsep) * \real{0.4286}}
  >{\raggedright\arraybackslash}p{(\linewidth - 4\tabcolsep) * \real{0.2381}}@{}}
\toprule\noalign{}
\begin{minipage}[b]{\linewidth}\raggedright
Program
\end{minipage} & \begin{minipage}[b]{\linewidth}\raggedright
Mechanism
\end{minipage} & \begin{minipage}[b]{\linewidth}\raggedright
Scale
\end{minipage} \\
\midrule\noalign{}
\endhead
\bottomrule\noalign{}
\endlastfoot
\textbf{REDD+} & Carbon market payments to prevent deforestation &
International \\
\textbf{Costa Rica PES} & Government payments to landowners for forest
conservation & National \\
\textbf{Pollination services} & Beekeepers paid by orchardists &
Local \\
\textbf{Wetland banking} & Developers purchase credits for wetland
destruction & National (US) \\
\end{longtable}
}

\textbf{Challenges:} Additionality (would conservation have happened
anyway?), leakage (deforestation displaced elsewhere), permanence
(forest may be cut later), measurement (how much carbon is actually
stored?).

\begin{quote}
\textbf{Concept Check:} The Millennium Ecosystem Assessment found that
60\% of ecosystem services are degraded. Give three specific examples of
degraded ecosystem services and explain how the degradation of each
affects human wellbeing.
\end{quote}

\begin{center}\rule{0.5\linewidth}{0.5pt}\end{center}

\subsection{Worked Example}\label{worked-example-7}

\textbf{Problem:} In a lake ecosystem, the gross primary productivity
(GPP) of phytoplankton is \(10,000\text{ J/m}^2\text{/yr}\). The
phytoplankton expend \(4,000\text{ J/m}^2\text{/yr}\) via cellular
respiration (\(R\)). Primary consumers (zooplankton) subsequently ingest
the phytoplankton with a trophic efficiency of 10\%. Calculate the net
primary productivity (NPP) of the phytoplankton and the energy available
to the secondary consumers.

\textbf{Solution:}

\textbf{Step 1. Calculate Net Primary Productivity (NPP).} NPP is the
energy that remains in the primary producers after respiration.
\[NPP = GPP - R \label{eq:unit_X_ecosystem_ecology_item_1}\]

\[NPP = 10,000\text{ J} - 4,000\text{ J} = 6,000\text{ J/m}^2\text{/yr} \label{eq:unit_X_ecosystem_ecology_item_2}\]

\textbf{Step 2. Calculate the energy transferred to primary consumers.}
The zooplankton consume with a 10\% trophic efficiency.
\[\text{Zooplankton Production} = NPP \times 0.10 \label{eq:unit_X_ecosystem_ecology_item_3}\]

\[\text{Zooplankton Production} = 6,000\text{ J} \times 0.10 = 600\text{ J/m}^2\text{/yr} \label{eq:unit_X_ecosystem_ecology_item_4}\]

\textbf{Step 3. Calculate the energy available to secondary consumers.}
If secondary consumers (e.g., small fish) also have a 10\% trophic
efficiency:
\[\text{Fish Production} = \text{Zooplankton Production} \times 0.10 \label{eq:unit_X_ecosystem_ecology_item_5}\]

\[\text{Fish Production} = 600\text{ J} \times 0.10 = 60\text{ J/m}^2\text{/yr} \label{eq:unit_X_ecosystem_ecology_item_6}\]

\textbf{Answer:} The NPP of the phytoplankton is
\textbf{\(6,000\text{ J/m}^2\text{/yr}\)}, and the energy available to
the secondary consumers (produced by primary consumers) is
\textbf{\(60\text{ J/m}^2\text{/yr}\)}.

\begin{center}\rule{0.5\linewidth}{0.5pt}\end{center}

\subsection{Worked Example: Nutrient Pool Mass Balance and Residence
Time}\label{worked-example-nutrient-pool-mass-balance-and-residence-time}

\textbf{Problem:} A forested watershed holds a soil-available nitrogen
pool of \(X = 800\text{ kg N/ha}\). Atmospheric deposition plus
weathering supply an input of \(I = 8\text{ kg N/ha/yr}\), and the pool
is at steady state. Calculate the stream-plus-erosion output flux \(O\),
the mean residence time \texorpdfstring{\ensuremath{\tau}}{tau}, and the leakage rate \(k\), then verify the
steady-state pool size and interpret whether the system is open or
closed.

\textbf{Solution:}

\textbf{Step 1. Identify the given variables.} - Nitrogen pool (\(X\)) =
\(800\text{ kg N/ha}\) - Input flux (\(I\)) = \(8\text{ kg N/ha/yr}\) -
Steady state: \(dX/dt = 0\), so output equals input.

\textbf{Step 2. Determine the output flux from the steady-state mass
balance.}
\[\frac{dX}{dt} = I - O = 0 \;\Rightarrow\; O = I = 8\text{ kg N/ha/yr} \label{eq:unit_X_ecosystem_ecology_item_7}\]

\textbf{Step 3. Compute the mean residence time.}
\[\tau = \frac{X}{O} = \frac{800\text{ kg N/ha}}{8\text{ kg N/ha/yr}} = 100\text{ yr} \label{eq:unit_X_ecosystem_ecology_item_8}\]

\textbf{Step 4. Compute the leakage rate and check the steady-state
pool.}
\[k = \frac{1}{\tau} = \frac{1}{100\text{ yr}} = 0.01\text{ yr}^{-1} \label{eq:unit_X_ecosystem_ecology_item_9}\]

\[\hat{X} = \frac{I}{k} = \frac{8\text{ kg N/ha/yr}}{0.01\text{ yr}^{-1}} = 800\text{ kg N/ha} \label{eq:unit_X_ecosystem_ecology_item_10}\]

The recovered \(\hat{X}\) equals the stated pool, confirming the system
is genuinely at steady state.

\textbf{Answer:} \(O = 8\text{ kg N/ha/yr}\), \(\tau = 100\text{ yr}\),
and \(k = 0.01\text{ yr}^{-1}\). Because \(\tau = 100\text{ yr} \gg 1\),
an average nitrogen atom is recycled internally about 100 times before
it leaves the watershed, so this is a tightly retaining, quasi-closed
system; a disturbance that strips deep-rooted vegetation would raise
\(O\), collapse \texorpdfstring{\ensuremath{\tau}}{tau}, and rapidly flush the nitrogen pool to streamwater.

\begin{center}\rule{0.5\linewidth}{0.5pt}\end{center}

\subsection{Current Evidence and Frontier
Biology}\label{current-evidence-and-frontier-biology-36}

For \textbf{Ecosystem Ecology and Biogeochemical Cycles}, frontier
biology belongs inside the evidence logic of the chapter. Ecology and
conservation decisions increasingly combine field data, remote sensing,
community knowledge, model uncertainty, and explicit values. The core
reading question is this: ecosystem claims should track stocks, fluxes,
residence times, boundaries, and coupled cycles.

\begin{itemize}
\tightlist
\item
  \textbf{What to verify:} identify the observation, model, assay, or
  dataset that would make the claim stronger or weaker.
\item
  \textbf{What to qualify:} state the scale, organism, cell type,
  environmental condition, or population where the claim is expected to
  hold.
\item
  \textbf{What to compare:} test at least one alternative explanation,
  baseline, or null model before treating the pattern as causal.
\item
  \textbf{What to cite:} distinguish primary evidence, review synthesis,
  public dataset, and institutional guidance; for recent or numeric
  claims, prefer the source closest to the measurement and state what
  has changed since it was published.
\end{itemize}

Use biodiversity metrics carefully: population indices, extinction risk
categories, ecosystem services, and management targets answer different
questions
\citep{ipbes2019global,ipbes2024transformative,wwf2024livingplanet,iucn2025redlist,fao2024sofia}.

\textbf{Source practice:} For conservation claims, cite assessment
sources and state whether the evidence is a population index,
extinction-risk assessment, ecosystem-service valuation, satellite
product, or policy synthesis
\citep{ipbes2024transformative,noaa2025coralbleaching,fao2025sofi}.

\subsubsection{Current Evidence Map: Agroecology as Coupled
Fluxes}\label{current-evidence-map-agroecology-as-coupled-fluxes}

\begin{figure}[htbp]
\centering
\includegraphics[width=0.82\linewidth,height=4.2in,keepaspectratio]{../figures/mermaid_inline/inline_mermaid_0186_30023f5c24bf.png}
\caption{Food-security claims should connect ecological mechanisms to access, resilience, livelihoods, and tradeoffs rather than equating yield alone with nutrition \textbackslash{}citep\{fao2025sofi\}.}
\end{figure}

\subsection{Key Terms}\label{key-terms-36}

{\def\LTcaptype{none} % do not increment counter
\begin{longtable}[]{@{}
  >{\raggedright\arraybackslash}p{(\linewidth - 2\tabcolsep) * \real{0.2857}}
  >{\raggedright\arraybackslash}p{(\linewidth - 2\tabcolsep) * \real{0.7143}}@{}}
\toprule\noalign{}
\begin{minipage}[b]{\linewidth}\raggedright
Term
\end{minipage} & \begin{minipage}[b]{\linewidth}\raggedright
Definition
\end{minipage} \\
\midrule\noalign{}
\endhead
\bottomrule\noalign{}
\endlastfoot
\textbf{GPP} & Gross Primary Production; total CO\(_2\) fixation rate in
an ecosystem \\
\textbf{NPP} & Net Primary Production; GPP - autotrophic respiration;
biomass available to consumers \\
\textbf{NEP} & Net Ecosystem Production; NPP - heterotrophic
respiration; ecosystem carbon balance \\
\textbf{NBP} & Net Biome Production; NEP - disturbance losses; regional
carbon balance \\
\textbf{Trophic efficiency} & Fraction of energy at trophic level \(n\)
transferred to level \(n+1\); typically about 10\% \\
\textbf{Biomagnification} & Increasing concentration of persistent
pollutants at higher trophic levels \\
\textbf{Eutrophication} & Nutrient enrichment of water body \texorpdfstring{\ensuremath{\to}}{->} algal
bloom \texorpdfstring{\ensuremath{\to}}{->} low O\(_2\) \texorpdfstring{\ensuremath{\to}}{->} hypoxia \\
\textbf{Biogeochemical cycle} & Movement of an element through biotic
and abiotic compartments \\
\textbf{Haber-Bosch process} & Industrial NH\(_3\) synthesis; doubles
global reactive N; drives eutrophication \\
\textbf{Biological pump} & Transfer of carbon from ocean surface to deep
water via biological processes \\
\textbf{Permafrost feedback} & Warming \texorpdfstring{\ensuremath{\to}}{->} permafrost thaw \texorpdfstring{\ensuremath{\to}}{->} CO\(_2\) +
CH\(_4\) release \texorpdfstring{\ensuremath{\to}}{->} further warming \\
\textbf{Ocean acidification} & CO\(_2\) absorption \texorpdfstring{\ensuremath{\to}}{->} carbonic acid \texorpdfstring{\ensuremath{\to}}{->} pH
decrease \texorpdfstring{\ensuremath{\to}}{->} carbonate undersaturation \\
\textbf{DMS (dimethylsulphide)} & Volatile S compound from
phytoplankton; \texorpdfstring{\ensuremath{\to}}{->} cloud condensation nuclei \\
\textbf{Redfield ratio} & 106C:16N:1P stoichiometry of marine
phytoplankton \\
\textbf{Reactive nitrogen cascade} & Sequential environmental damage as
anthropogenic N moves through compartments \\
\textbf{REDD+} & International mechanism for carbon payments to prevent
deforestation \\
\textbf{Eddy covariance} & Micrometeorological method for measuring
ecosystem carbon flux \\
\textbf{CLAW hypothesis} & DMS-cloud-albedo feedback as planetary
thermostat \\
\textbf{ANPP / BNPP} & Aboveground / belowground net primary production;
ratio varies from 0.4 (rainforest) to 3+ (tundra) \\
\textbf{Eddy covariance} & Continuous tower-based measurement of
ecosystem CO\(_2\) flux; gold standard for NEP \\
\textbf{Martin curve} & Exponential attenuation of organic-matter flux
with depth in the ocean \\
\textbf{Hubbard Brook} & Long-running watershed experiment that
demonstrated the dominant role of vegetation in nutrient retention \\
\textbf{CLORPT} & Jenny's five soil-forming state factors: climate,
organisms, relief, parent material, time \\
\textbf{Soil horizons} & Vertical layers (O, A, E, B, C, R) reflecting
depth-dependent pedogenic processes \\
\textbf{Planetary boundaries} & Quantitative limits on anthropogenic
perturbation across nine Earth-system processes \\
\textbf{Earth System Science} & Coupled study of biology, geochemistry,
hydrology, and atmospheric science as a single planetary system \\
\textbf{Export production / e-ratio} & Fraction of NPP that escapes the
surface ocean as sinking particles \\
\textbf{Open vs.~closed cycle} & Whether a nutrient cycle is dominated
by external fluxes (open) or internal recycling (closed) \\
\textbf{Mean residence time (\texorpdfstring{\ensuremath{\tau}}{tau})} & \(X/O\); average time a nutrient atom
spends in a pool before leaving \\
\end{longtable}
}

\begin{center}\rule{0.5\linewidth}{0.5pt}\end{center}

\subsection{Review Questions}\label{review-questions-36}

\begin{enumerate}
\def\labelenumi{\arabic{enumi}.}
\item
  An estuary receives large inputs of fertiliser-derived N and P from
  agricultural runoff each spring. Trace the \textbf{cascade of events}
  from nutrient loading to hypoxia, naming each ecological and microbial
  process involved. Which specific microbial process creates the anoxic
  condition?
\item
  A tropical rainforest has GPP = 3,000 g C/m\(^2\)/yr and autotrophic
  respiration = 1,500 g C/m\(^2\)/yr. Heterotrophic respiration = 1,200
  g C/m\(^2\)/yr. Calculate: (a) NPP, (b) NEP, (c) Is this ecosystem a
  net carbon sink or source? (d) If deforestation converts this to
  farmland with NEP = -200 g C/m\(^2\)/yr, what is the additional carbon
  released per hectare per year?
\item
  Iron fertilisation experiments (Martin 1990: ``Give me half a tanker
  of iron and I will give you an ice age'') showed that adding Fe to
  HNLC ocean regions triggers phytoplankton blooms. (a) Explain why
  these regions are high in N and P but low in phytoplankton. (b) Why
  doesn't iron fertilisation cause long-term net CO\(_2\) drawdown? (c)
  What are the ecological risks of large-scale iron fertilisation?
\item
  The permafrost carbon feedback is estimated to release 1-1.5 Pg
  CO\(_2\)-equivalent per year for each degree of additional warming.
  Current warming projections suggest 2-4\(^\circ\)C above
  pre-industrial by 2100. Calculate the range of additional
  CO\(_2\)-equivalent that permafrost might contribute over 80 years,
  assuming linear temperature rise. Compare with annual fossil fuel
  emissions of about 10 Pg C/yr.
\item
  Compare the nitrogen and phosphorus cycles in terms of: (a) whether
  they have a significant atmospheric phase, (b) the primary
  anthropogenic disruption, (c) the main environmental consequence, and
  (d) whether the element can be industrially synthesised from
  atmospheric or geological sources.
\item
  Using the concept of trophic efficiency, calculate how many kilograms
  of phytoplankton are needed to produce 1 kg of tuna (assume tuna feeds
  at trophic level 4 with 10\% efficiency at each step). How does this
  relate to the environmental footprint of eating high-trophic-level
  seafood?
\item
  Explain the biological pump in detail. Why is it important for climate
  regulation? If ocean warming stratifies the surface layer more
  strongly (reducing nutrient upwelling), what happens to the efficiency
  of the biological pump and atmospheric CO\(_2\)?
\item
  The Keeling Curve shows atmospheric CO\(_2\) rising from 315 ppm in
  1958 to 421 ppm in 2024. Given that about 55\% of
  anthropogenic CO\(_2\) emissions are absorbed by land and ocean sinks:
  (a) What would atmospheric CO\(_2\) be without these sinks? (b) If
  ocean acidification reduces the ocean's capacity to absorb CO\(_2\) by
  15\%, how much additional CO\(_2\) would remain in the atmosphere per
  year?
\item
  A farmer applies 200 kg N/ha as ammonium nitrate fertiliser. Primarily
  40\% is taken up by the corn crop. Trace the fate of the remaining 120
  kg N/ha through the nitrogen cycle: what fraction is likely
  denitrified, leached as nitrate, volatilised as NH\(_3\), and
  converted to N\(_2\)O? What is the total environmental cost of this
  ``excess'' nitrogen?
\item
  Compare the CLAW hypothesis (DMS-cloud feedback) with the permafrost
  carbon feedback. One is a negative feedback and one is positive.
  Explain which is which, their relative magnitudes, and their
  implications for Earth's climate trajectory.
\item
  Assign trophic levels with \breaktt{food\_web\_trophic\_levels} for a
  three-species chain; verify tuna-style pyramids require multiplying
  efficiencies.
\item
  Why is \textbf{Redfield} stoichiometry a diagnostic for nutrient
  limitation rather than a comprehensive constant?
\item
  The Hubbard Brook clear-cut experiment increased streamwater
  NO\(_3^-\) losses 26-fold. Trace the mechanism through (a) cessation
  of plant uptake, (b) accelerated decomposition without canopy cooling,
  (c) elevated water tables and reduced transpiration. Why is this
  evidence that mature forests function as quasi-closed nutrient
  systems?
\item
  Compare four methods for measuring NPP (harvest, eddy covariance,
  \(^{14}\)C, MODIS remote sensing). For each, give one ecosystem where
  it is the \emph{most} defensible method and one where it would
  systematically underestimate true NPP.
\item
  A coastal upwelling zone has surface NPP = 1,500 g C/m\(^2\)/yr and an
  e-ratio of 0.4; a subtropical gyre has NPP = 50 g C/m\(^2\)/yr and an
  e-ratio of 0.1. (a) Calculate export production for each. (b) If both
  export equal \emph{fractions} of carbon below 1,000 m via the Martin
  curve (\(z_{1/2} = 200\) m), which contributes more to long-term
  carbon sequestration per unit area? (c) How does climate-driven
  stratification of the surface ocean change both numbers?
\item
  Explain the ANPP/BNPP partition for a temperate grassland (BNPP/ANPP \texorpdfstring{\ensuremath{\approx}}{~}
  2) vs.~a tropical rainforest (BNPP/ANPP \texorpdfstring{\ensuremath{\approx}}{~} 0.4). Why does this
  partition affect (a) carbon-credit accounting under REDD+, (b)
  ecosystem response to drought, (c) long-term soil organic carbon
  accumulation?
\item
  Apply the Redfield ratio to diagnose three scenarios: (a) ambient N :
  P = 6, (b) N : P = 16, (c) N : P = 35. For each, predict which
  nutrient is limiting and which functional groups (e.g., diatoms,
  cyanobacteria, dinoflagellates) you expect to dominate.
\item
  The Gaia hypothesis claims Earth's biosphere actively maintains
  habitability. Evaluate this claim by listing two stabilising
  biological feedbacks (negative) and two destabilising ones (positive).
  What is the difference between ``Gaia as a metaphor for Earth System
  Science'' and ``Gaia as a literal superorganism''? Why was Doolittle
  (1981)'s critique decisive?
\end{enumerate}

\begin{center}\rule{0.5\linewidth}{0.5pt}\end{center}

\subsection{Further Reading and Source
Notes}\label{further-reading-and-source-notes-31}

\begin{itemize}
\tightlist
\item
  Levin (1998). Ecosystems and the Biosphere as Complex Adaptive
  Systems. \emph{Ecosystems}, 1.
\item
  Bormann \& Likens (1967). Nutrient cycling. \emph{Science}, 155.
\item
  Redfield (1958). The biological control of chemical factors in the
  environment. \emph{American Scientist}, 46.
\item
  Yachi \& Loreau (1999). Biodiversity and ecosystem productivity in a
  fluctuating environment: The insurance hypothesis. \emph{Proceedings
  of the National Academy of Sciences}, 96.
\end{itemize}

\begin{center}\rule{0.5\linewidth}{0.5pt}\end{center}

\subsection{Computational Bridge}\label{computational-bridge-33}

Trophic position can be read off an adjacency list of ``eats''
relations:

\begin{Shaded}
\begin{Highlighting}[]
\ImportTok{from}\NormalTok{ biology.ecology }\ImportTok{import}\NormalTok{ food\_web\_trophic\_levels}

\NormalTok{web }\OperatorTok{=}\NormalTok{ \{}\StringTok{"phyto"}\NormalTok{: [], }\StringTok{"zooplankton"}\NormalTok{: [}\StringTok{"phyto"}\NormalTok{], }\StringTok{"fish"}\NormalTok{: [}\StringTok{"zooplankton"}\NormalTok{]\}}
\BuiltInTok{print}\NormalTok{(food\_web\_trophic\_levels(web))}
\end{Highlighting}
\end{Shaded}

\begin{quote}
\textbf{Clinical / systems note:} Biomagnification of methylmercury
follows the same stepped energy flow you quantify with trophic level
assignments.
\end{quote}

\begin{center}\rule{0.5\linewidth}{0.5pt}\end{center}

\subsection{Summary}\label{summary-36}

\begin{itemize}
\tightlist
\item
  \textbf{Ecosystems} = biotic + abiotic components; energy flows
  through trophic levels (about 10\% efficiency; Lindeman
  1942); nutrient atoms cycle.
\item
  \textbf{GPP \texorpdfstring{\ensuremath{\to}}{->} NPP \texorpdfstring{\ensuremath{\to}}{->} NEP \texorpdfstring{\ensuremath{\to}}{->} NBP} hierarchy; global terrestrial NPP
  about 120 Pg C/yr; tropical rainforests dominate.
\item
  \textbf{Energy flow:} one-way through system with about 10\%
  trophic efficiency; biomagnification of persistent pollutants
  concentrates toxins at higher levels.
\item
  \textbf{Carbon cycle:} 10 Pg C/yr fossil fuel input; atmospheric
  CO\(_2\) at 421 ppm (2024); biological pump exports 10-15 Pg C/yr to
  deep ocean; permafrost (1,700 Pg) vulnerable; ocean acidification
  (\(\Omega < 1\) \texorpdfstring{\ensuremath{\to}}{->} coral dissolution).
\item
  \textbf{Nitrogen cycle:} N\(_2\) fixation (nitrogenase; Haber-Bosch
  doubles global reactive N) \texorpdfstring{\ensuremath{\to}}{->} nitrification \texorpdfstring{\ensuremath{\to}}{->} denitrification /
  anammox. Reactive N cascade causes eutrophication, hypoxia, N\(_2\)O
  (298x CO\(_2\) warming potential).
\item
  \textbf{Phosphorus cycle:} sedimentary (no gas phase); P-limiting in
  freshwater; Redfield ratio 106C:16N:1P; globally finite reserves.
\item
  \textbf{Sulphur cycle:} DMS \texorpdfstring{\ensuremath{\to}}{->} cloud condensation nuclei \texorpdfstring{\ensuremath{\to}}{->} planetary
  thermostat (CLAW hypothesis); acid rain from SO\(_2\) (reversed by
  regulation).
\item
  \textbf{Measuring NPP:} harvest, eddy covariance (FLUXNET), \(^{14}\)C
  uptake, MODIS remote sensing --- each captures different components;
  ANPP/BNPP partition varies from 0.4 (rainforest) to \textgreater3
  (tundra) and shifts conclusions about global carbon storage.
\item
  \textbf{Hubbard Brook \citep{bormann1967}:} the clear-cut experiment
  showed live vegetation retains 80--95\% of nutrient inputs; mature
  forests are quasi-closed for most nutrients. Open vs.~closed nutrient
  regimes governed by mean residence time \(\tau = X/O\).
\item
  \textbf{Ocean ecosystem ecology:} half of global NPP; productivity
  zones (euphotic, mesopelagic, abyssopelagic, upwelling, gyres, polar);
  biological pump exports about 10--15 Pg C/yr; Martin curve
  attenuates flux with depth.
\item
  \textbf{Soil formation:} Jenny's CLORPT factors; horizon
  differentiation (O--A--E--B--C--R); soils store 1,500--2,400 Pg C;
  warming-driven SOC release is a major positive climate feedback.
\item
  \textbf{Coupled cycles:} Redfield
  \(106\,\text{C}:16\,\text{N}:1\,\text{P}\) diagnoses nutrient
  limitation; the planetary boundaries framework places six of nine
  Earth-system limits already crossed \citep{richardson2023earth}.
\item
  \textbf{Water cycle:} plant transpiration drives 60-90\% of
  terrestrial evaporation; tropical forests generate their own rainfall.
\item
  \textbf{Ecosystem services} (about \$125T/yr; Costanza 2014);
  60\% degraded (MEA 2005); biodiversity increases ecosystem function
  via complementarity and insurance.
\item
  \textbf{Earth System Science / Gaia:} modern descendant of the
  Lovelock--Margulis hypothesis; biology and abiotic systems are
  coupled, but the ``superorganism'' claim is rejected --- feedbacks can
  stabilise \emph{or} destabilise.
\item
  \textbf{Connections:} See \cref{sec:unit_III_photosynthesis} for
  photosynthetic GPP/NPP concepts, Unit VIII for transpiration in the
  water cycle, and \cref{sec:unit_X_biomes_and_conservation} for biome
  NPP benchmarks.
\end{itemize}

\begin{center}\rule{0.5\linewidth}{0.5pt}\end{center}

\subsubsection{Companion Source
Module}\label{companion-source-module-36}

\textbf{Ecosystem Ecology and Biogeochemical Cycles} should leave a
reproducible trail from a biological claim to the code, figure, diagram,
or paper-based activity that can test it. Use the surfaces below to
inspect the chapter's assumptions, rerun the relevant model, or compare
the manuscript explanation with companion labs and figures.

{\def\LTcaptype{none} % do not increment counter
\begin{longtable}[]{@{}
  >{\raggedright\arraybackslash}p{(\linewidth - 2\tabcolsep) * \real{0.5000}}
  >{\raggedright\arraybackslash}p{(\linewidth - 2\tabcolsep) * \real{0.5000}}@{}}
\toprule\noalign{}
\begin{minipage}[b]{\linewidth}\raggedright
Surface
\end{minipage} & \begin{minipage}[b]{\linewidth}\raggedright
Use it for
\end{minipage} \\
\midrule\noalign{}
\endhead
\bottomrule\noalign{}
\endlastfoot
\breaktt{src/biology/ecology/ecology.py}
(\breaktt{food\_web\_trophic\_levels}, \texttt{connectance},
\breaktt{biodiversity\_indices}) & Translate ecosystem structure into
trophic, network, and diversity calculations. \\
\breaktt{src/biology/botany/botany.py} (\breaktt{photosynthesis\_rate},
\texttt{water\_potential}) & Link primary production and plant-water
constraints to ecosystem fluxes. \\
\breaktt{src/mermaid/biology\_diagrams.py}
(\breaktt{nutrient\_cycle\_diagram}, \texttt{food\_web\_diagram}) & Keep
stocks, fluxes, and boundaries visually explicit. \\
\end{longtable}
}

\textbf{Reproducibility check:} state system boundary, stock, flux,
residence time, unit, and time window before comparing ecosystem
budgets. \textbf{Cross-reference:} connect with
\cref{sec:unit_III_photosynthesis}, \cref{sec:unit_X_community_ecology},
and \cref{sec:unit_X_biomes_and_conservation}.

\newpage

\section{Biomes and Conservation
Biology}\label{biomes-and-conservation-biology}

\label{sec:unit_X_biomes_and_conservation}

\begin{quote}
\textbf{Ch 35} \texorpdfstring{\ensuremath{\cdot}}{.} Level 2/3 \texorpdfstring{\ensuremath{\cdot}}{.} 70 min read \texorpdfstring{\ensuremath{\cdot}}{.} 75 min lecture \texorpdfstring{\ensuremath{\cdot}}{.}
Prerequisites: \cref{sec:unit_X_ecosystem_ecology}
\end{quote}

\subsection{Learning Objectives}\label{learning-objectives-37}

\begin{enumerate}
\def\labelenumi{\arabic{enumi}.}
\tightlist
\item
  Describe the 9 major terrestrial \hyperref[gl:biome]{\textbf{biome}}s
  and 5 aquatic biome types with climate, NPP, and indicator species.
\item
  Explain climate-driven biome boundary shifts and tipping points from
  2023-2025 scholarship.
\item
  Define the IUCN Red List categories,
  \hyperref[gl:minimum-viable-population]{\textbf{minimum viable
  population (MVP)}}, and the 50/500 rule.
\item
  Explain extinction debt and its implications for conservation
  timescales.
\item
  Design nature reserves using
  \hyperref[gl:island-biogeography]{\textbf{island biogeography}} and
  SLOSS principles with quantitative examples.
\item
  Describe rewilding and climate-adaptive management as emerging
  conservation strategies.
\item
  Explain conservation genomics and assisted
  \hyperref[gl:gene]{\textbf{gene}} flow as tools for managing
  threatened populations.
\item
  Evaluate the effectiveness of international conservation frameworks
  (CBD, 30x30, CITES).
\item
  Quantify IPBES findings (extinction rates, decline of vertebrate
  populations, drivers of biodiversity loss) and use them to evaluate
  conservation priorities.
\item
  Compute climate velocity for different landscape gradients and
  identify topographic, hydrological, microclimatic, and elevational
  refugia.
\item
  Compare ecosystem-based adaptation strategies (mangroves, urban
  greening, wetlands, coral reefs) with hard infrastructure on
  cost-benefit and co-benefit dimensions.
\end{enumerate}

\begin{figure}[htbp]
\centering
\includegraphics[width=0.85\textwidth]{../figures/biome_distribution.png}
\caption{Representative biome positions in temperature-precipitation space. Point area scales with net primary productivity, making tropical rainforest, temperate forest, savanna/grassland, desert, tundra, and marine reference conditions visually comparable.}
\label{fig:unit_X_biome_distribution}
\end{figure}

\subsubsection{Study Blueprint}\label{study-blueprint-37}

\begin{itemize}
\tightlist
\item
  \textbf{Big idea:} Biodiversity patterns and conservation decisions
  emerge from climate, history, disturbance, and human choice.
\item
  \textbf{Core concepts:} biomes, biodiversity, extinction risk,
  conservation planning.
\item
  \textbf{Framework alignment:} Vision \& Change: Systems, Evolution,
  Pathways and transformations of energy and matter; AP Biology: Systems
  Interactions, Evolution, Energetics; NGSS-style topics: Interdependent
  Relationships in Ecosystems, Matter and Energy in Organisms and
  Ecosystems, Natural Selection and Evolution.
\item
  \textbf{Model or quantitative lens:} Species-area, risk, and
  prioritization calculations.
\item
  \textbf{Data skill:} Use maps, trend data, and threat categories to
  justify conservation priorities.
\item
  \textbf{Practice cadence:} Representing and Describing Data,
  Statistical Tests and Data Analysis, Argumentation.
\item
  \textbf{Common misconception to repair:} Conservation is
  decision-making under constraints and uncertainty, not only preserving
  untouched nature.
\item
  \textbf{Primary lab:} \cref{sec:lab_unit_X_biomes_and_conservation}.
\item
  \textbf{Question bank:} \cref{sec:q_unit_X_biomes_and_conservation}.
\item
  \textbf{Transfer task:} Transfer conservation reasoning to land use,
  climate corridors, restoration, and environmental justice.
\item
  \textbf{Bridge to computation:}
  \breaktt{biology.ecology.ecology.species\_area\_relationship}.
\end{itemize}

\begin{center}\rule{0.5\linewidth}{0.5pt}\end{center}

\begin{quote}
\textbf{Opening Vignette --- The Experiment That Launched Conservation
Biology}

In 1966, ecologist E.O. Wilson and his student Daniel Simberloff tested
a bold prediction of island biogeography theory --- that small islands
far from the mainland support fewer species, and that a disturbed island
will bounce back to a predictable species count. Simberloff hired
professional exterminators to fumigate small mangrove islands off the
Florida Keys with methyl bromide, killing most arthropods. He then
monitored recolonisation. Within a year, species counts had returned to
near the pre-defaunation levels predicted by island area and distance
equations Wilson had developed with Robert MacArthur in 1967. The
MacArthur-Wilson Theory of Island Biogeography --- predicting species
number from area \(S = cA^z\) --- became the mathematical foundation for
reserve design in conservation biology. The SLOSS debate (Single Large
Or Several Small reserves), minimum viable populations, habitat
fragmentation analysis, and biodiversity hotspot identification most
derive from Wilson and MacArthur's equations. And it most started with
methyl bromide on tiny Keys mangroves.
\end{quote}

\subsection{Biome Concepts and
Classification}\label{biome-concepts-and-classification}

A \textbf{biome} is a large-scale ecological zone characterised by
climate (principally temperature and precipitation) and the type of
vegetation and animal communities adapted to those conditions. Biome
distribution follows Whittaker's (1975) climate-vegetation gradient map,
subsequently confirmed and refined by MODIS satellite NPP data (Running
et al.~2004, \emph{Science}) and updated with machine-learning-enhanced
Koppen-Geiger classifications (Monzon-Alvarado et al.~2023, \emph{Glob.
Change Biol.}).

\subsubsection{Climate as the Master
Variable}\label{climate-as-the-master-variable}

The distribution of biomes is determined primarily by two climate
variables:

\begin{enumerate}
\def\labelenumi{\arabic{enumi}.}
\tightlist
\item
  \textbf{Mean annual temperature (MAT)} --- determines the length of
  the growing season and the type of vegetation (tropical vs.~temperate
  vs.~boreal)
\item
  \textbf{Mean annual precipitation (MAP)} --- determines vegetation
  structure (forest vs.~grassland vs.~desert)
\end{enumerate}

\textbf{Whittaker's biome diagram}
(\cref{fig:unit_X_biome_distribution}) plots MAT vs.~MAP and shows that
biome boundaries are remarkably predictable from these two variables
alone. However, other factors also matter: seasonality, soil type, fire
frequency, and evolutionary history.

\subsubsection{The Nine Major Terrestrial
Biomes}\label{the-nine-major-terrestrial-biomes}

{\def\LTcaptype{none} % do not increment counter
\begin{longtable}[]{@{}
  >{\raggedright\arraybackslash}p{(\linewidth - 8\tabcolsep) * \real{0.0676}}
  >{\raggedright\arraybackslash}p{(\linewidth - 8\tabcolsep) * \real{0.1486}}
  >{\raggedright\arraybackslash}p{(\linewidth - 8\tabcolsep) * \real{0.1757}}
  >{\raggedright\arraybackslash}p{(\linewidth - 8\tabcolsep) * \real{0.3514}}
  >{\raggedright\arraybackslash}p{(\linewidth - 8\tabcolsep) * \real{0.2568}}@{}}
\toprule\noalign{}
\begin{minipage}[b]{\linewidth}\raggedright
Biome
\end{minipage} & \begin{minipage}[b]{\linewidth}\raggedright
Temperature
\end{minipage} & \begin{minipage}[b]{\linewidth}\raggedright
Precipitation
\end{minipage} & \begin{minipage}[b]{\linewidth}\raggedright
Characteristic adaptations
\end{minipage} & \begin{minipage}[b]{\linewidth}\raggedright
NPP (g C/m\(^2\)/yr)
\end{minipage} \\
\midrule\noalign{}
\endhead
\bottomrule\noalign{}
\endlastfoot
\textbf{Tropical rainforest} & 25-30\(^\circ\)C year-round &
\textgreater2,000 mm/yr & Buttress roots,
\hyperref[gl:epiphyte]{\textbf{epiphyte}}s, drip tips, fig-wasp
\hyperref[gl:mutualism]{\textbf{mutualism}}, stratified canopy
(emergent/canopy/understory/forest floor) & 1,000-1,750 \\
\textbf{Tropical seasonal forest / savanna} & 25\(^\circ\)C; distinct
dry season & 1,000-2,000 mm/yr & Deciduousness (dry season);
fire-adapted grasses (resprouting); ungulate herding; C4 grasses &
450-1,000 \\
\textbf{Desert} & Extreme diurnal range; hot or cold & \textless250
mm/yr & CAM \hyperref[gl:photosynthesis]{\textbf{photosynthesis}};
nocturnal behaviour; waxy cuticles; deep taproots; water storage &
45-125 \\
\textbf{Temperate grassland / prairie} & 0-25\(^\circ\)C; large annual
range & 300-1,500 mm/yr & Deep root systems (\textgreater2 m); C4
grasses in warm areas; burrowing animals; fire-maintained & 250-450 \\
\textbf{Temperate deciduous forest} & 0-20\(^\circ\)C; 4 distinct
seasons & 750-1,500 mm/yr & Autumn leaf senescence (ABA-mediated);
arbuscular mycorrhizae; spring ephemerals & 500-1,250 \\
\textbf{Temperate rainforest} & 5-15\(^\circ\)C; ocean-moderated &
\textgreater1,500 mm/yr & Tallest trees on Earth (coastal redwood
\emph{Sequoia sempervirens}, 115 m); heavy epiphytic moss cover &
750-1,500 \\
\textbf{Boreal forest (taiga)} & -10 to 10\(^\circ\)C & 300-850 mm/yr
(snow) & Conical crown shape (snow shedding); antifreeze
\hyperref[gl:protein]{\textbf{protein}}s; permafrost soils; low
decomposition & 300-600 \\
\textbf{Tundra (arctic)} & -30 to 10\(^\circ\)C; short growing season &
\textless250 mm/yr & Cushion plants; freeze-tolerance (cryoprotectants);
permafrost (active layer \textless1 m); lichens & 50-200 \\
\textbf{Mediterranean shrubland (chaparral)} & Mild wet winters; hot dry
summers & 300-800 mm/yr & Sclerophyllous leaves (thick cuticle, high
fibre); fire-serotinous cones; resprouting lignotubers & 225-650 \\
\end{longtable}
}

\subsubsection{Five Major Aquatic Biome
Types}\label{five-major-aquatic-biome-types}

{\def\LTcaptype{none} % do not increment counter
\begin{longtable}[]{@{}
  >{\raggedright\arraybackslash}p{(\linewidth - 6\tabcolsep) * \real{0.0943}}
  >{\raggedright\arraybackslash}p{(\linewidth - 6\tabcolsep) * \real{0.2264}}
  >{\raggedright\arraybackslash}p{(\linewidth - 6\tabcolsep) * \real{0.3208}}
  >{\raggedright\arraybackslash}p{(\linewidth - 6\tabcolsep) * \real{0.3585}}@{}}
\toprule\noalign{}
\begin{minipage}[b]{\linewidth}\raggedright
Biome
\end{minipage} & \begin{minipage}[b]{\linewidth}\raggedright
Key features
\end{minipage} & \begin{minipage}[b]{\linewidth}\raggedright
Indicator species
\end{minipage} & \begin{minipage}[b]{\linewidth}\raggedright
NPP (g C/m\(^2\)/yr)
\end{minipage} \\
\midrule\noalign{}
\endhead
\bottomrule\noalign{}
\endlastfoot
\textbf{Freshwater lakes \& rivers} & Salinity \textless0.5 ppt; thermal
stratification (epilimnion/hypolimnion); lotic vs.~lentic & Trout,
cattails, caddisfly larvae & 50-750 \\
\textbf{Wetlands} & High organic load;
\hyperref[gl:anaerobic]{\textbf{anaerobic}} sediment; CH\(_4\)
production; water table at or near surface & Herons, \emph{Sphagnum}
moss, alligators & 400-1,750 \\
\textbf{Coastal/neritic marine} & Shallow (\textless200 m); upwelling
nutrients; high turbidity; continental shelf & Kelp, cod, sea otters,
oysters & 100-1,250 \\
\textbf{Open ocean (pelagic)} & Nutrient-poor subtropical gyres;
stratified; ``blue desert'' &
\hyperref[gl:phytoplankton]{\textbf{Phytoplankton}}, bluefin tuna, sperm
whale & 25-200 \\
\textbf{Coral reefs} & Tropical; 20-30\(^\circ\)C; high light;
mutualistic \emph{Symbiodiniaceae} (zooxanthellae) & Brain coral,
parrotfish, reef sharks & 500-2,000 \\
\end{longtable}
}

\subsubsection{Lake Stratification and
Turnover}\label{lake-stratification-and-turnover}

Temperate lakes exhibit seasonal stratification:

{\def\LTcaptype{none} % do not increment counter
\begin{longtable}[]{@{}
  >{\raggedright\arraybackslash}p{(\linewidth - 6\tabcolsep) * \real{0.1316}}
  >{\raggedright\arraybackslash}p{(\linewidth - 6\tabcolsep) * \real{0.2895}}
  >{\raggedright\arraybackslash}p{(\linewidth - 6\tabcolsep) * \real{0.1842}}
  >{\raggedright\arraybackslash}p{(\linewidth - 6\tabcolsep) * \real{0.3947}}@{}}
\toprule\noalign{}
\begin{minipage}[b]{\linewidth}\raggedright
Layer
\end{minipage} & \begin{minipage}[b]{\linewidth}\raggedright
Temperature
\end{minipage} & \begin{minipage}[b]{\linewidth}\raggedright
Density
\end{minipage} & \begin{minipage}[b]{\linewidth}\raggedright
Characteristics
\end{minipage} \\
\midrule\noalign{}
\endhead
\bottomrule\noalign{}
\endlastfoot
\textbf{Epilimnion} & Warm (20-25\(^\circ\)C) & Low & Surface; mixed by
wind; O\(_2\)-rich, nutrient-depleted \\
\textbf{Thermocline} (metalimnion) & Rapid change & Transition & Density
barrier prevents mixing \\
\textbf{Hypolimnion} & Cold (4\(^\circ\)C) & High & Deep;
O\(_2\)-depleted; nutrient-rich from decomposition \\
\end{longtable}
}

\textbf{Turnover} occurs in spring and autumn when surface temperature
passes through 4\(^\circ\)C (maximum density of water), eliminating the
thermocline and allowing complete mixing. This redistributes O\(_2\) and
nutrients throughout the water column.

\subsubsection{Coral Reef Bleaching}\label{coral-reef-bleaching}

\textbf{Mechanism:} Thermal stress (anomaly \textgreater1\(^\circ\)C
sustained \textgreater4 weeks, measured as degree heating weeks or DHW)
triggers expulsion of symbiotic zooxanthellae (\emph{Symbiodiniaceae})
from coral tissue. Without their photosynthetic partners, corals lose
90\% of energy supply and appear white (``bleached''). If thermal stress
persists \textgreater8 weeks, coral mortality follows.

\textbf{Recent mass bleaching events:} The Great Barrier Reef
experienced mass bleaching in 2016, 2017, 2020, 2022, 2024, and
\textbf{2025} (six events in ten years). NOAA and the International
Coral Reef Initiative confirmed the fourth global bleaching event
beginning in 2023; by late 2025, bleaching-level heat stress had
affected about 84\% of the world's coral reef area
\citep{noaa2025coralbleaching}. At \textgreater2\(^\circ\)C global
warming, \textgreater99\% of coral reefs are projected to experience
annual severe bleaching \citep{ipcc2021ar6wg1}.

\textbf{Assisted evolution caveat:} Conservation tools now include
selective breeding, assisted gene flow, symbiont manipulation, and
microbiome conditioning for heat tolerance. These interventions can
protect specific reefs or buy time, but they do not remove the thermal
driver; without emissions reduction and local water-quality management,
heat-tolerant genotypes still face repeated bleaching, acidification,
disease, and storm damage.

\begin{figure}[htbp]
\centering
\includegraphics[width=0.82\linewidth,height=4.2in,keepaspectratio]{../figures/mermaid_inline/inline_mermaid_0187_fb7fd0b64f1f.png}
\caption{Coral bleaching and assisted-evolution response. Heat-tolerant hosts, symbionts, or microbiomes can improve local resilience in some settings, but they are complements to climate mitigation and reef-water-quality management, not substitutes \textbackslash{}citep\{noaa2025coralbleaching,strader2022coralheat\}.}
\end{figure}

\begin{quote}
\texorpdfstring{}{} \textbf{Clinical Connection --- Reef-Derived Pharmaceuticals:} Coral
reef ecosystems are a source of bioactive compounds for drug
development. \textbf{Pseudopterosin} (from the Caribbean sea whip
\emph{Pseudopterogorgia elisabethae}) has potent anti-inflammatory and
analgesic properties and is used in commercial wound-healing
formulations. \textbf{Discodermolide} (from the deep-sea sponge
\emph{Discodermia dissoluta}) shows anti-cancer properties by
stabilising microtubules, similar to taxol. The cone snail toxin
\textbf{ziconotide} (Prialt) is an FDA-approved non-opioid analgesic for
severe chronic pain. Reef destruction through bleaching and
acidification threatens the loss of undiscovered pharmaceutical
compounds before they can be identified --- a form of pharmacological
extinction debt.
\end{quote}

\begin{quote}
\textbf{Concept Check:} What climatic variables determine whether a
region supports tropical rainforest vs.~tropical savanna? How does fire
interact with precipitation to maintain the savanna biome?
\end{quote}

\begin{center}\rule{0.5\linewidth}{0.5pt}\end{center}

\subsection{Climate Change and Biome Boundary
Shifts}\label{climate-change-and-biome-boundary-shifts}

\begin{figure}[htbp]
\centering
\includegraphics[width=0.82\linewidth,height=4.2in,keepaspectratio]{../figures/mermaid_inline/inline_mermaid_0188_654bba2ed7a3.png}
\caption{Flowchart for Climate Change and Biome Boundary Shifts: Global Warming +1.1 degrees C since pre-industrial, Poleward Shift 5-10 km/decade (Chen et al. 2023), Upslope Shift 6 m/decade elevation (mountain ecosystems), and Tipping Points (non-linear transitions) form the diagram's primary path or branches.}
\end{figure}

\subsubsection{Key Findings (2023-2025)}\label{key-findings-2023-2025}

{\def\LTcaptype{none} % do not increment counter
\begin{longtable}[]{@{}
  >{\raggedright\arraybackslash}p{(\linewidth - 4\tabcolsep) * \real{0.3182}}
  >{\raggedright\arraybackslash}p{(\linewidth - 4\tabcolsep) * \real{0.2273}}
  >{\raggedright\arraybackslash}p{(\linewidth - 4\tabcolsep) * \real{0.4545}}@{}}
\toprule\noalign{}
\begin{minipage}[b]{\linewidth}\raggedright
Finding
\end{minipage} & \begin{minipage}[b]{\linewidth}\raggedright
Study
\end{minipage} & \begin{minipage}[b]{\linewidth}\raggedright
Key result
\end{minipage} \\
\midrule\noalign{}
\endhead
\bottomrule\noalign{}
\endlastfoot
Biome boundaries shifting poleward & Chen et al.~(2023, \emph{Nat. Clim.
Change}) & 5-10 km/decade poleward; 6 m/decade upslope \\
Amazon dieback tipping point & Lovejoy \& Nobre (2024, \emph{Sci. Adv.})
& 20-25\% deforestation \texorpdfstring{\ensuremath{\to}}{->} runaway savannification; currently at
about 17\% \\
Boreal \texorpdfstring{\ensuremath{\to}}{->} tundra greening & Elmendorf et al.~(2023, \emph{Nat. Ecol.
Evol.}) & Shrub cover +50\% in 20 years; permafrost destabilisation \\
Novel biomes (no-analog climates) & Williams et al.~(2023, \emph{Proc.
R. Soc. B}) & about 20\% of Earth's surface may develop novel
biome types by 2100 \\
Permafrost thaw methane pulse & Turetsky et al.~(2024, \emph{Nat.
Geosci.}) & Abrupt thaw adds 2-5 Gt C/yr by 2050; omitted from most
models \\
Tree mortality acceleration & McDowell et al.~(2024, \emph{Science}) &
Global tree mortality rates doubled since 1970s; drought/heat/insects \\
\end{longtable}
}

\subsubsection{Biome Suitability
Modelling}\label{biome-suitability-modelling}

\begin{equation}
S = \sum_i w_i \cdot f_i(T, P)
\label{eq:biomes_and_conservation_1}
\end{equation}

where \(w_i\) = variable weights (temperature, precipitation,
seasonality index), \(f_i\) = Gaussian response function centred on
optimal conditions.

\textbf{Species distribution models (SDMs)} project future range shifts
by correlating current distributions with climate variables and
projecting onto future climate scenarios:

{\def\LTcaptype{none} % do not increment counter
\begin{longtable}[]{@{}
  >{\raggedright\arraybackslash}p{(\linewidth - 6\tabcolsep) * \real{0.2632}}
  >{\raggedright\arraybackslash}p{(\linewidth - 6\tabcolsep) * \real{0.2105}}
  >{\raggedright\arraybackslash}p{(\linewidth - 6\tabcolsep) * \real{0.2368}}
  >{\raggedright\arraybackslash}p{(\linewidth - 6\tabcolsep) * \real{0.2895}}@{}}
\toprule\noalign{}
\begin{minipage}[b]{\linewidth}\raggedright
Model type
\end{minipage} & \begin{minipage}[b]{\linewidth}\raggedright
Approach
\end{minipage} & \begin{minipage}[b]{\linewidth}\raggedright
Strengths
\end{minipage} & \begin{minipage}[b]{\linewidth}\raggedright
Limitations
\end{minipage} \\
\midrule\noalign{}
\endhead
\bottomrule\noalign{}
\endlastfoot
\textbf{Bioclimatic envelope} & Correlative (climate \texorpdfstring{\ensuremath{\to}}{->} presence) &
Simple; good for many species & Ignores dispersal, biotic
interactions \\
\textbf{MaxEnt} & Machine learning; maximum entropy & High predictive
power; handles sparse data & Overfitting risk; correlative primarily \\
\textbf{Process-based (e.g., LPJ-GUESS)} & Mechanistic plant physiology
& Captures mechanisms; projects novel climates & Data-intensive;
computationally expensive \\
\end{longtable}
}

\subsubsection{Climate Velocity}\label{climate-velocity}

\textbf{Climate velocity} = the speed at which climate conditions shift
geographically:

\begin{equation}
v_{climate} = \frac{\text{rate of temperature change (°C/yr)}}{\text{spatial temperature gradient (°C/km)}}
\label{eq:biomes_and_conservation_2}
\end{equation}

In flat landscapes (e.g., Amazon basin), spatial gradients are shallow \texorpdfstring{\ensuremath{\to}}{->}
climate velocity is fast (about 10 km/yr). In mountains, steep
gradients \texorpdfstring{\ensuremath{\to}}{->} slow climate velocity (about 1 km/yr). Species must
disperse at least as fast as climate velocity to track suitable
conditions.

\textbf{Migration lag:} Many species cannot track climate velocity: -
Trees: typical dispersal rate = 0.1-0.5 km/yr via seed dispersal -
Climate velocity in temperate lowlands: about 4-10 km/yr - Gap:
10-100x faster movement needed than most trees can achieve

\begin{quote}
\textbf{Concept Check:} The Amazon rainforest is at about 17\%
deforestation, with a tipping point estimated at 20-25\%. Why is the
tipping point non-linear (i.e., why doesn't the forest simply lose area
proportionally to deforestation)? What role does
\hyperref[gl:transpiration]{\textbf{transpiration}} recycling play in
the tipping mechanism?
\end{quote}

\begin{center}\rule{0.5\linewidth}{0.5pt}\end{center}

\subsection{Conservation Biology
Fundamentals}\label{conservation-biology-fundamentals}

\subsubsection{IPBES Global Assessment: The Quantified
Crisis}\label{ipbes-global-assessment-the-quantified-crisis}

The \textbf{Intergovernmental Science-Policy Platform on Biodiversity
and Ecosystem Services (IPBES)} --- analogous to the IPCC for climate
--- released its first Global Assessment in 2019 and its Transformative
Change Assessment in 2024
\citep{ipbes2019global,ipbes2024transformative}. Its findings frame
contemporary conservation:

{\def\LTcaptype{none} % do not increment counter
\begin{longtable}[]{@{}
  >{\raggedright\arraybackslash}p{(\linewidth - 2\tabcolsep) * \real{0.4062}}
  >{\raggedright\arraybackslash}p{(\linewidth - 2\tabcolsep) * \real{0.5938}}@{}}
\toprule\noalign{}
\begin{minipage}[b]{\linewidth}\raggedright
IPBES finding
\end{minipage} & \begin{minipage}[b]{\linewidth}\raggedright
Quantitative result
\end{minipage} \\
\midrule\noalign{}
\endhead
\bottomrule\noalign{}
\endlastfoot
Species threatened with extinction & about 1,000,000 in the
IPBES synthesis; the IUCN Red List provides the tracked assessed-species
subset \citep{iucn2025redlist} \\
Average abundance of native species in major terrestrial biomes &
Declined \texorpdfstring{\textgreater{}=}{>=} 20\% (mostly since 1900) \\
Monitored vertebrate populations & 73\% average decline since 1970 in
the Living Planet Index; this is a population-index trend, not a count
of individual animals \citep{wwf2024livingplanet} \\
Wetlands & 85\% lost since 1700 \\
Forest cover & 32\% lost (vs.~pre-industrial) \\
Coral reef cover & 50\% lost since 1870 \\
Pollinator populations & 40\% of insect pollinators threatened \\
Local livestock breeds & 9\% extinct; another 1,000 threatened (genetic
erosion of food security) \\
\end{longtable}
}

\subsubsection{Quantified Drivers (IPBES
2019)}\label{quantified-drivers-ipbes-2019}

IPBES ranked the relative importance of the five direct drivers of
biodiversity loss in \textbf{terrestrial and freshwater systems}:

{\def\LTcaptype{none} % do not increment counter
\begin{longtable}[]{@{}
  >{\raggedright\arraybackslash}p{(\linewidth - 4\tabcolsep) * \real{0.1176}}
  >{\raggedright\arraybackslash}p{(\linewidth - 4\tabcolsep) * \real{0.1765}}
  >{\raggedright\arraybackslash}p{(\linewidth - 4\tabcolsep) * \real{0.7059}}@{}}
\toprule\noalign{}
\begin{minipage}[b]{\linewidth}\raggedright
Rank
\end{minipage} & \begin{minipage}[b]{\linewidth}\raggedright
Driver
\end{minipage} & \begin{minipage}[b]{\linewidth}\raggedright
Approximate contribution
\end{minipage} \\
\midrule\noalign{}
\endhead
\bottomrule\noalign{}
\endlastfoot
1 & Land-use change (agriculture, urbanisation, infrastructure) &
about 30\% \\
2 & Direct exploitation (hunting, fishing, logging) &
about 23\% \\
3 & Climate change & about 14\% (rising fast, projected to
dominate by 2050) \\
4 & Pollution (nutrients, plastics, persistent organics) &
about 14\% \\
5 & Invasive alien species & about 11\% \\
\end{longtable}
}

For \textbf{marine systems}, the ranking changes: direct exploitation
\textgreater{} climate change \textgreater{} pollution \textgreater{}
land-/sea-use change \textgreater{} invasive species. Climate change is
projected to overtake direct exploitation as the dominant marine driver
by 2050.

The Transformative Change Assessment (IPBES 2024) goes further:
\textbf{incremental conservation alone cannot meet 30x30 or biodiversity
targets}. Systemic changes to economic, financial, and governance
systems are required --- including reforming environmentally harmful
subsidies, shifting incentives toward ecosystem stewardship, and
integrating biodiversity into most policy domains (the
``whole-of-government'' approach) \citep{ipbes2024transformative}.

Food systems are a useful place to make ``transformative change''
concrete. FAO, IFAD, UNICEF, WFP, and WHO report that high food-price
inflation continues to undermine access to healthy diets, especially for
low-income populations \citep{fao2025sofi}. Conservation biology
therefore cannot treat agricultural landscapes as outside nature:
agroecology, diversified crop rotations, soil-carbon restoration,
pollinator habitat, reduced food waste, sustainable fisheries, and
equitable market access are biodiversity strategies and food-security
strategies at the same time. The claim is not that agroecology alone
solves hunger; it is that biodiversity, nutrition, farm livelihoods, and
climate resilience are coupled enough that single-objective policies
often fail.

\subsubsection{The Sixth Mass
Extinction}\label{the-sixth-mass-extinction}

Earth currently experiences the \textbf{Sixth Mass Extinction} ---
driven primarily by human activities:

{\def\LTcaptype{none} % do not increment counter
\begin{longtable}[]{@{}
  >{\raggedright\arraybackslash}p{(\linewidth - 4\tabcolsep) * \real{0.1714}}
  >{\raggedright\arraybackslash}p{(\linewidth - 4\tabcolsep) * \real{0.2571}}
  >{\raggedright\arraybackslash}p{(\linewidth - 4\tabcolsep) * \real{0.5714}}@{}}
\toprule\noalign{}
\begin{minipage}[b]{\linewidth}\raggedright
Driver
\end{minipage} & \begin{minipage}[b]{\linewidth}\raggedright
Mechanism
\end{minipage} & \begin{minipage}[b]{\linewidth}\raggedright
Relative contribution
\end{minipage} \\
\midrule\noalign{}
\endhead
\bottomrule\noalign{}
\endlastfoot
\textbf{Habitat destruction} & Land-use change, deforestation,
urbanisation & about 30\% of threatened species \\
\textbf{Overexploitation} & Hunting, fishing, wildlife trade &
about 20\% \\
\textbf{Invasive species} & Competition, predation, disease introduction
& about 15\% \\
\textbf{Pollution} & Pesticides, plastics, nutrient loading &
about 10\% \\
\textbf{Climate change} & Range shifts, phenological mismatch, extreme
events & about 8\% (increasing rapidly) \\
\textbf{Disease} & Chytrid fungus (Bd), white-nose syndrome &
about 5\% \\
\end{longtable}
}

\textbf{Key statistics:} - Current extinction rate: 100-1,000x
background rate (Ceballos et al.~2017, \emph{PNAS}) - IPBES 2019 Global
Assessment: about 1 million species threatened with extinction
- \textbf{Living Planet Index} (WWF 2024): 73\% average decline in
monitored vertebrate population sizes since 1970; this measures a
weighted set of monitored populations, not every vertebrate individual
\citep{wwf2024livingplanet} - 75\% of terrestrial ecosystems
significantly altered by humans - 50\% of global wetland area lost since
1900 - 33\% of marine fish stocks overfished in FAO's recent global
assessments \citep{fao2024sofia}

\subsubsection{IUCN Red List Categories (2025
Update)}\label{iucn-red-list-categories-2025-update}

{\def\LTcaptype{none} % do not increment counter
\begin{longtable}[]{@{}
  >{\raggedright\arraybackslash}p{(\linewidth - 4\tabcolsep) * \real{0.2353}}
  >{\raggedright\arraybackslash}p{(\linewidth - 4\tabcolsep) * \real{0.1176}}
  >{\raggedright\arraybackslash}p{(\linewidth - 4\tabcolsep) * \real{0.6471}}@{}}
\toprule\noalign{}
\begin{minipage}[b]{\linewidth}\raggedright
Category
\end{minipage} & \begin{minipage}[b]{\linewidth}\raggedright
Code
\end{minipage} & \begin{minipage}[b]{\linewidth}\raggedright
Quantitative threshold
\end{minipage} \\
\midrule\noalign{}
\endhead
\bottomrule\noalign{}
\endlastfoot
\textbf{Extinct} & EX & No reasonable doubt that last individual has
died \\
\textbf{Extinct in the Wild} & EW & Primarily in cultivation or
captivity \\
\textbf{Critically Endangered} & CR & at least 80\% reduction over 10 yr
or 3 generations \\
\textbf{Endangered} & EN & at least 50\% reduction over 10 yr or 3
generations \\
\textbf{Vulnerable} & VU & at least 30\% reduction over 10 yr or 3
generations \\
\textbf{Near Threatened} & NT & Close to qualifying for VU; being
monitored \\
\textbf{Least Concern} & LC & Widespread and abundant \\
\textbf{Data Deficient} & DD & Insufficient data to assess \\
\end{longtable}
}

The IUCN Red List is a tracked subset of assessed species, not a census
of global biodiversity. Version 2025-1 listed 169,420 assessed species,
with 47,187 classified as threatened (Vulnerable, Endangered, or
Critically Endangered), about 28\% of assessed species
\citep{iucn2025redlist}. Counts change with new assessments and
reassessments, so the category logic is more durable than any single
annual total.

\textbf{Criteria for listing} (IUCN uses 5 criteria, any one
sufficient): - \textbf{A:} Population reduction (quantitative decline
thresholds) - \textbf{B:} Geographic range (extent of occurrence or area
of occupancy) with fragmentation/decline - \textbf{C:} Small population
size (\(N < 10,000\) for VU; \(N < 2,500\) for EN; \(N < 250\) for CR)
with decline - \textbf{D:} Very small population (\(N < 1,000\) for VU;
\(N < 250\) for EN; \(N < 50\) for CR) - \textbf{E:} Quantitative
analysis (PVA) showing extinction probability \textgreater10\% within
100 years (VU)

\begin{quote}
\textbf{Concept Check:} A bird species has a global population of 8,000
mature individuals with a documented 35\% decline over the past 15 years
(approximately 3 generations). Under which IUCN criteria and category
would this species be listed?
\end{quote}

\begin{figure}[htbp]
\centering
\includegraphics[width=0.82\linewidth,height=4.2in,keepaspectratio]{../figures/mermaid_inline/inline_mermaid_0189_700b94aae4af.png}
\caption{IUCN Red List assessment framework and conservation triage decision tree. Species are assessed against five criteria (population decline, geographic range, small population, very small population, quantitative analysis). Threatened species receive conservation interventions based on urgency and feasibility.}
\end{figure}

\begin{center}\rule{0.5\linewidth}{0.5pt}\end{center}

\subsection{Population Viability and the 50/500
Rule}\label{population-viability-and-the-50500-rule}

\subsubsection{Minimum Viable Population
(MVP)}\label{minimum-viable-population-mvp}

The \textbf{minimum viable population} is the smallest population with
at least 95\% probability of persistence for at least 100 generations,
accounting for four types of stochasticity:

{\def\LTcaptype{none} % do not increment counter
\begin{longtable}[]{@{}
  >{\raggedright\arraybackslash}p{(\linewidth - 4\tabcolsep) * \real{0.3396}}
  >{\raggedright\arraybackslash}p{(\linewidth - 4\tabcolsep) * \real{0.1698}}
  >{\raggedright\arraybackslash}p{(\linewidth - 4\tabcolsep) * \real{0.4906}}@{}}
\toprule\noalign{}
\begin{minipage}[b]{\linewidth}\raggedright
Stochasticity type
\end{minipage} & \begin{minipage}[b]{\linewidth}\raggedright
Mechanism
\end{minipage} & \begin{minipage}[b]{\linewidth}\raggedright
Effect on small populations
\end{minipage} \\
\midrule\noalign{}
\endhead
\bottomrule\noalign{}
\endlastfoot
\textbf{Demographic} & Random variation in births/deaths & Even when
\(r > 0\), small \(N\) \texorpdfstring{\ensuremath{\to}}{->} high extinction probability \\
\textbf{Environmental} & Weather, food fluctuations & Entire population
experiences same bad conditions \\
\textbf{Genetic} & Inbreeding depression, drift & \(F\) per generation
\(= 1/(2N_e)\); \(F > 0.10\) \texorpdfstring{\ensuremath{\to}}{->} fitness costs \\
\textbf{Catastrophes} & Fire, disease epidemic, storm & Sudden
\textgreater50\% mortality events; frequency matters \\
\end{longtable}
}

\subsubsection{Franklin's 50/500 Rule}\label{franklins-50500-rule}

\textbf{Franklin (1980) and \citet{franklinsoule1980}:}

\begin{itemize}
\tightlist
\item
  \textbf{\(N_e \geq 50\):} Short-term viability --- prevents rapid
  inbreeding depression

  \begin{itemize}
  \tightlist
  \item
    Inbreeding rate: \(\Delta F = \frac{1}{2N_e} \leq 1\%\) per
    generation
  \item
    At \(F > 0.10\): measurable inbreeding depression (reduced litter
    size, increased juvenile mortality)
  \end{itemize}
\item
  \textbf{\(N_e \geq 500\):} Long-term viability --- maintains additive
  genetic variance for adaptive evolution

  \begin{itemize}
  \tightlist
  \item
    Genetic variance lost per generation: \(\frac{1}{2N_e} = 0.1\%\)
  \item
    Balanced by new \hyperref[gl:mutation]{\textbf{mutation}} if
    \(N_e > 500\) \citep{frankham1995}
  \end{itemize}
\end{itemize}

\textbf{Revised 100/1000 rule} (Frankham et al.~2014, \emph{Biol.
Conserv.}): More conservative; accounts for reduced purging
effectiveness and updated genomic data.

\textbf{\(N_e\) vs.~census \(N\):} Effective population size is
typically much smaller than census size:

\begin{equation}
N_e = \frac{4N_m N_f}{N_m + N_f}
\label{eq:biomes_and_conservation_3}
\end{equation}

Typical \(N_e/N\) ratio: 0.10-0.25 for vertebrates. Therefore, to
achieve \(N_e = 500\), census population needs to be
\(N = 2,000-5,000\).

\subsubsection{Population Viability Analysis
(PVA)}\label{population-viability-analysis-pva}

\textbf{PVA} uses species-specific demographic data and stochastic
simulation models (VORTEX, RAMAS) to project extinction probability:

\textbf{Inputs:} age/sex-specific survival, fecundity, variance in vital
rates, carrying capacity, catastrophe frequency, inbreeding
coefficients, habitat area.

\textbf{Outputs:} extinction probability over specified time horizon,
mean time to extinction, quasi-extinction threshold probability.

\textbf{Case study --- Florida panther} (\emph{Puma concolor coryi}): -
1994 PVA \texorpdfstring{\ensuremath{\to}}{->} predicted extinction within 25 years without intervention -
\(N \approx 30\), \(N_e \approx 10\) \texorpdfstring{\ensuremath{\to}}{->} severe inbreeding (kinked tails,
cryptorchidism, heart defects) - 1995: 8 female Texas cougars introduced
\texorpdfstring{\ensuremath{\to}}{->} \textbf{genetic rescue} - 2023: population \textgreater200
individuals; heterozygosity increased, inbreeding depression reversed -
One of conservation biology's greatest success stories

\begin{quote}
\texorpdfstring{}{} \textbf{Clinical Connection --- Conservation Genetics and Inbreeding
Depression in Human Populations:} The principles of MVP and inbreeding
depression apply to human populations. Small, isolated populations
(e.g., the Amish community of Lancaster County, Pennsylvania) show
elevated frequencies of rare autosomal recessive disorders including
Ellis-van Creveld syndrome (6-fingered dwarfism; about 5\%
carrier frequency vs.~about 0.001\% in general population),
maple syrup urine disease, and glutaric aciduria. The
\textbf{\hyperref[gl:founder-effect]{founder effect}} --- establishment
of a population by a small number of individuals carrying a
non-representative sample of the gene pool --- mirrors the
\hyperref[gl:genetic-drift]{\textbf{genetic drift}} effects that reduce
\(N_e\) in endangered wildlife. Genomic screening programs in such
communities follow the same logic as conservation genomics: identifying
deleterious \hyperref[gl:allele]{\textbf{allele}}s and managing their
frequency.
\end{quote}

\begin{center}\rule{0.5\linewidth}{0.5pt}\end{center}

\subsection{Extinction Debt and Its Conservation
Implications}\label{extinction-debt-and-its-conservation-implications}

\textbf{Extinction debt} (Tilman et al.~1994, \emph{Nature}): Species
present in degraded/fragmented habitats are committed to future
extinction because habitat can no longer support viable populations ---
but extinction is \textbf{delayed} (50-500 years) by existing long-lived
individuals.

\subsubsection{Mathematical Framework}\label{mathematical-framework}

\begin{equation}
\text{Species committed to extinction} = S_{\text{pre-fragmentation}} - c \cdot A_{\text{remaining}}^z
\label{eq:biomes_and_conservation_4}
\end{equation}

where \(S_{\text{pre-fragmentation}}\) is species richness before
habitat loss, \(A_{\text{remaining}}\) is remaining habitat area, and
\(c\) and \(z\) are the species-area parameters whose log--log slope
steepens from contiguous continental biomes toward isolated oceanic
islands (\cref{fig:unit_X_species_area}).

\begin{figure}[htbp]
\centering
\includegraphics[width=0.85\textwidth]{../figures/species_area_relationship.png}
\caption{Species--area relationship on log--log axes. Mainland-fragment and oceanic-island curves show contrasting slopes ($z about 0.15$ versus $z about 0.35$), with a dotted reference curve at $z = 0.25$.}
\label{fig:unit_X_species_area}
\end{figure}

\subsubsection{Relaxation Time}\label{relaxation-time}

\textbf{Relaxation time} = the period between habitat loss and final
extinction of committed species. Determined by: - Longevity of
individuals (trees: centuries; insects: months) - Generation time
(longer generation = longer debt) - Habitat quality of remaining
fragments - Connectivity between patches (rescue effect)

\subsubsection{Case Studies}\label{case-studies}

{\def\LTcaptype{none} % do not increment counter
\begin{longtable}[]{@{}
  >{\raggedright\arraybackslash}p{(\linewidth - 6\tabcolsep) * \real{0.1395}}
  >{\raggedright\arraybackslash}p{(\linewidth - 6\tabcolsep) * \real{0.3023}}
  >{\raggedright\arraybackslash}p{(\linewidth - 6\tabcolsep) * \real{0.3488}}
  >{\raggedright\arraybackslash}p{(\linewidth - 6\tabcolsep) * \real{0.2093}}@{}}
\toprule\noalign{}
\begin{minipage}[b]{\linewidth}\raggedright
System
\end{minipage} & \begin{minipage}[b]{\linewidth}\raggedright
Debt estimate
\end{minipage} & \begin{minipage}[b]{\linewidth}\raggedright
Relaxation time
\end{minipage} & \begin{minipage}[b]{\linewidth}\raggedright
Reference
\end{minipage} \\
\midrule\noalign{}
\endhead
\bottomrule\noalign{}
\endlastfoot
European calcareous grasslands & 20-50\% of plant species & 50-200 years
& Lindborg \& Eriksson 2004 \\
Amazon forest fragments & 30-50\% of canopy trees & 100-500 years &
Laurance et al.~2011 \\
African tropical forests & 10-30\% of endemic species & Ongoing & Brooks
et al.~1999 \\
Belgian forest herbs & 30-40\% in fragments \textless1 ha & 100+ years &
Vellend et al.~2006 \\
\end{longtable}
}

\textbf{Conservation implication:} Present-day species richness
\textbf{overestimates} long-term viability. Conservation assessments
based on current presence may be dangerously optimistic. Extinction debt
means that the full consequences of today's deforestation will not be
realised for decades to centuries.

\begin{quote}
\textbf{Concept Check:} A 10,000-ha forest is fragmented into ten 100-ha
patches. Using \(S = 20A^{0.25}\), calculate species richness for the
intact forest and for each fragment. How many species are ``committed''
to extinction via extinction debt? How would wildlife corridors between
patches modify this prediction?
\end{quote}

\begin{center}\rule{0.5\linewidth}{0.5pt}\end{center}

\subsection{Reserve Design}\label{reserve-design}

\subsubsection{Island Biogeography Applied to
Conservation}\label{island-biogeography-applied-to-conservation}

MacArthur-Wilson island biogeography theory provides the theoretical
foundation for reserve design:

\textbf{\hyperref[gl:species-area-relationship]{Species-area
relationship}:} \(S = cA^z\) (see Community Ecology chapter for
derivation)

\begin{Shaded}
\begin{Highlighting}[]
\ImportTok{from}\NormalTok{ biology.ecology }\ImportTok{import}\NormalTok{ species\_area\_relationship}

\CommentTok{\# Compare large vs. small reserves}
\NormalTok{S\_large }\OperatorTok{=}\NormalTok{ species\_area\_relationship(area\_ha}\OperatorTok{=}\DecValTok{10\_000}\NormalTok{, c}\OperatorTok{=}\DecValTok{7}\NormalTok{, z}\OperatorTok{=}\FloatTok{0.26}\NormalTok{)}
\NormalTok{S\_small }\OperatorTok{=}\NormalTok{ species\_area\_relationship(area\_ha}\OperatorTok{=}\DecValTok{100}\NormalTok{, c}\OperatorTok{=}\DecValTok{7}\NormalTok{, z}\OperatorTok{=}\FloatTok{0.26}\NormalTok{)}
\BuiltInTok{print}\NormalTok{(}\SpecialStringTok{f"Large reserve (10,000 ha): }\SpecialCharTok{\{}\NormalTok{S\_large}\SpecialCharTok{:.0f\}}\SpecialStringTok{ species"}\NormalTok{)}
\BuiltInTok{print}\NormalTok{(}\SpecialStringTok{f"Small reserve (100 ha):    }\SpecialCharTok{\{}\NormalTok{S\_small}\SpecialCharTok{:.0f\}}\SpecialStringTok{ species"}\NormalTok{)}
\BuiltInTok{print}\NormalTok{(}\SpecialStringTok{f"Species lost from fragmentation: }\SpecialCharTok{\{}\NormalTok{((S\_large}\OperatorTok{{-}}\NormalTok{S\_small)}\OperatorTok{/}\NormalTok{S\_large)}\OperatorTok{*}\DecValTok{100}\SpecialCharTok{:.0f\}}\SpecialStringTok{\%"}\NormalTok{)}
\end{Highlighting}
\end{Shaded}

\subsubsection{SLOSS Debate Resolution}\label{sloss-debate-resolution}

{\def\LTcaptype{none} % do not increment counter
\begin{longtable}[]{@{}
  >{\raggedright\arraybackslash}p{(\linewidth - 4\tabcolsep) * \real{0.4167}}
  >{\raggedright\arraybackslash}p{(\linewidth - 4\tabcolsep) * \real{0.3889}}
  >{\raggedright\arraybackslash}p{(\linewidth - 4\tabcolsep) * \real{0.1944}}@{}}
\toprule\noalign{}
\begin{minipage}[b]{\linewidth}\raggedright
Design principle
\end{minipage} & \begin{minipage}[b]{\linewidth}\raggedright
When it applies
\end{minipage} & \begin{minipage}[b]{\linewidth}\raggedright
Example
\end{minipage} \\
\midrule\noalign{}
\endhead
\bottomrule\noalign{}
\endlastfoot
\textbf{Single large} preferred & Area-demanding species; interior
habitat specialists; low landscape connectivity & Amazon: large reserves
protect jaguar, harpy eagle \\
\textbf{Several small} preferred & Different habitat types in different
patches; high beta diversity; patches protect local endemics & Caribbean
islands: each island has unique endemic reptiles \\
\textbf{Core-corridor-matrix} (modern synthesis) & Most real-world
situations & Y2Y (Yellowstone to Yukon, 3,200 km) \\
\end{longtable}
}

\subsubsection{Diamond's Design Principles (1975,
updated)}\label{diamonds-design-principles-1975-updated}

{\def\LTcaptype{none} % do not increment counter
\begin{longtable}[]{@{}
  >{\raggedright\arraybackslash}p{(\linewidth - 4\tabcolsep) * \real{0.2903}}
  >{\raggedright\arraybackslash}p{(\linewidth - 4\tabcolsep) * \real{0.2903}}
  >{\raggedright\arraybackslash}p{(\linewidth - 4\tabcolsep) * \real{0.4194}}@{}}
\toprule\noalign{}
\begin{minipage}[b]{\linewidth}\raggedright
Principle
\end{minipage} & \begin{minipage}[b]{\linewidth}\raggedright
Reasoning
\end{minipage} & \begin{minipage}[b]{\linewidth}\raggedright
Modern status
\end{minipage} \\
\midrule\noalign{}
\endhead
\bottomrule\noalign{}
\endlastfoot
Larger is better than smaller & Species-area relationship & Supported \\
One large \textgreater{} several small of equal total area & More
interior habitat; lower extinction & Context-dependent \\
Close together better than far apart & Higher colonisation rate &
Supported \\
Clustered better than linear arrangement & More connectivity &
Supported \\
Circular better than elongated & Less edge effect & Supported \\
Connected by corridors better than isolated & Rescue effect; gene flow &
Strongly supported \\
\end{longtable}
}

\subsubsection{Edge Effects}\label{edge-effects}

Habitat edges create microclimatic gradients that penetrate 100-300 m
into forest fragments:

{\def\LTcaptype{none} % do not increment counter
\begin{longtable}[]{@{}
  >{\raggedright\arraybackslash}p{(\linewidth - 4\tabcolsep) * \real{0.3056}}
  >{\raggedright\arraybackslash}p{(\linewidth - 4\tabcolsep) * \real{0.5278}}
  >{\raggedright\arraybackslash}p{(\linewidth - 4\tabcolsep) * \real{0.1667}}@{}}
\toprule\noalign{}
\begin{minipage}[b]{\linewidth}\raggedright
Edge effect
\end{minipage} & \begin{minipage}[b]{\linewidth}\raggedright
Penetration distance
\end{minipage} & \begin{minipage}[b]{\linewidth}\raggedright
Impact
\end{minipage} \\
\midrule\noalign{}
\endhead
\bottomrule\noalign{}
\endlastfoot
Increased light & 50-200 m & Favours pioneer species; dries forest
interior \\
Wind damage & 100-300 m & Increased tree mortality \\
Elevated temperature & 50-100 m & Shifts microclimate; stresses shade
species \\
Nest predation/parasitism & 200-600 m & Cowbird parasitism and corvid
predation increase \\
Invasive species & 100-500 m & Exotic plants establish in disturbed
edges \\
\end{longtable}
}

For a circular reserve:

\begin{equation}
\text{Interior area} = \pi(r - d)^2
\label{eq:biomes_and_conservation_5}
\end{equation}

where \(r\) = reserve radius and \(d\) = edge penetration distance. For
a 10-ha circular reserve (\(r = 178\) m) with 200-m edge effects:
\textbf{zero interior habitat.} This illustrates why small fragments are
ecological traps for interior-dependent species.

\subsubsection{Habitat Corridors}\label{habitat-corridors}

\textbf{Evidence for corridor effectiveness:} - Tewksbury et al.~(2002,
\emph{PNAS}): Connected patches had 25-100\% more pollinator movement
than unconnected - Gilbert-Norton et al.~(2010, meta-analysis): Movement
between patches 50\% higher with corridors - \textbf{Florida Wildlife
Corridor} (legislated 2021): 1,000+ mile protected corridor from
Everglades to Okefenokee Swamp

\textbf{Stepping stones:} For species that can cross moderate distances
(birds, bats, strong-flying insects), stepping stone patches may be as
effective as continuous corridors at lower cost.

\begin{quote}
\texorpdfstring{}{} \textbf{Clinical Connection --- Urban Green Corridors and Human
Health:} The principles of habitat corridors extend to urban planning
and human health. Studies in Philadelphia, Chicago, and Singapore
demonstrate that connected urban greenways (tree-lined corridors
connecting parks) reduce urban heat island effects by 2-5\(^\circ\)C,
decrease particulate matter (PM\(_{2.5}\)) by 15-25\%, and increase
physical activity by 30-40\% in nearby residents. Green corridors also
reduce stress biomarkers (salivary
\hyperref[gl:cortisol]{\textbf{cortisol}} decreased 20\% after 30-minute
walks in urban green corridors vs.~streets; Roe et al.~2013).
Conservation biology's core-corridor-matrix design principles thus have
direct applications to One Health approaches integrating biodiversity
conservation, climate adaptation, and human wellbeing.
\end{quote}

\begin{quote}
\textbf{Concept Check:} A conservation planner has a budget to protect
1,000 ha of tropical forest. Option A: one 1,000-ha reserve. Option B:
ten 100-ha reserves spread across different habitat types. Option C:
five 100-ha reserves connected by 100-m-wide corridors (total 500 ha
reserves + 500 ha corridors). Which design best protects: (a) a
large-bodied apex predator? (b) maximum regional species richness? (c)
genetic connectivity? Justify each answer using island biogeography
theory.
\end{quote}

\begin{center}\rule{0.5\linewidth}{0.5pt}\end{center}

\subsection{Emerging Conservation
Strategies}\label{emerging-conservation-strategies}

\subsubsection{Rewilding}\label{rewilding}

\textbf{Rewilding} (Soule \& Noss 1998; updated Jorgensen 2024,
\emph{Nat. Ecol. Evol.}): Restore ecological processes --- especially
trophic interactions --- by reintroducing lost species, particularly
apex predators and megaherbivores.

\textbf{Three types of rewilding:}

{\def\LTcaptype{none} % do not increment counter
\begin{longtable}[]{@{}
  >{\raggedright\arraybackslash}p{(\linewidth - 4\tabcolsep) * \real{0.2105}}
  >{\raggedright\arraybackslash}p{(\linewidth - 4\tabcolsep) * \real{0.4211}}
  >{\raggedright\arraybackslash}p{(\linewidth - 4\tabcolsep) * \real{0.3684}}@{}}
\toprule\noalign{}
\begin{minipage}[b]{\linewidth}\raggedright
Type
\end{minipage} & \begin{minipage}[b]{\linewidth}\raggedright
Approach
\end{minipage} & \begin{minipage}[b]{\linewidth}\raggedright
Example
\end{minipage} \\
\midrule\noalign{}
\endhead
\bottomrule\noalign{}
\endlastfoot
\textbf{Trophic rewilding} & Reintroduce apex predators/herbivores to
restore cascades & Wolves in Yellowstone (1995); lynx in UK
(proposed) \\
\textbf{Pleistocene rewilding} & Use ecological proxies for extinct
megafauna & Pleistocene Park, Siberia (bison, horses to create mammoth
steppe) \\
\textbf{Passive rewilding} & Remove human disturbance; allow natural
regeneration & European farmland abandonment \texorpdfstring{\ensuremath{\to}}{->} forest recovery \\
\end{longtable}
}

\textbf{Quantitative outcomes of wolf reintroduction in Yellowstone:} -
Elk: population reduced about 50\% (12,000 \texorpdfstring{\ensuremath{\to}}{->} 6,000); behaviour
change (avoid riverbanks) - Willows: height increased 5-fold in riparian
areas within 15 years - Beavers: 0 colonies in 1995 \texorpdfstring{\ensuremath{\to}}{->} 12 colonies by
2012 - Songbirds: 5 species returned to recovering riparian habitat -
Stream morphology: channels narrowed and deepened as bank vegetation
stabilised soil - Berry production: increased \texorpdfstring{\ensuremath{\to}}{->} grizzly bear diet
shifted (consuming 10x more berries)

\textbf{Pleistocene Park hypothesis:} Zimov et al.~(2012,
\emph{Science}): Reintroducing large herbivores to Arctic tundra
compacts snow in winter, allowing cold air to freeze soil deeper,
\textbf{slowing permafrost thaw} by about 80\% compared to
snow-insulated control plots.

\subsubsection{Conservation Genomics}\label{conservation-genomics}

\textbf{Whole-\hyperref[gl:genome]{\textbf{genome}} sequencing} of
threatened populations enables:

Conservation genomics is strongest when sequencing is tied to a
management decision rather than treated as an expensive description.
Runs of homozygosity estimate recent inbreeding, allele-frequency
outliers can nominate local adaptation, pangenome or long-read
approaches can reveal structural variation, and pedigree-aware sampling
can separate ancestry from recent bottlenecks. The ethical layer is
inseparable: sampling permits, Indigenous data sovereignty,
benefit-sharing, and the risk of over-prioritising ``genetically
interesting'' populations most affect whether the genomic evidence
should guide action.

{\def\LTcaptype{none} % do not increment counter
\begin{longtable}[]{@{}
  >{\raggedright\arraybackslash}p{(\linewidth - 4\tabcolsep) * \real{0.4583}}
  >{\raggedright\arraybackslash}p{(\linewidth - 4\tabcolsep) * \real{0.2500}}
  >{\raggedright\arraybackslash}p{(\linewidth - 4\tabcolsep) * \real{0.2917}}@{}}
\toprule\noalign{}
\begin{minipage}[b]{\linewidth}\raggedright
Application
\end{minipage} & \begin{minipage}[b]{\linewidth}\raggedright
Method
\end{minipage} & \begin{minipage}[b]{\linewidth}\raggedright
Example
\end{minipage} \\
\midrule\noalign{}
\endhead
\bottomrule\noalign{}
\endlastfoot
\textbf{Inbreeding assessment} & Runs of homozygosity (ROH) analysis &
Cheetah: about 10 Mb average ROH; severe bottleneck \\
\textbf{Adaptive variation} & Genome-wide association (GWAS) for
climate-relevant loci & Coral heat tolerance alleles in
\emph{Acropora} \\
\textbf{Management units} & Population structure (FST, admixture
analysis) & Tiger: 6 subspecies \texorpdfstring{\ensuremath{\to}}{->} conservation priorities \\
\textbf{Genetic rescue planning} & Identify source populations
maximising heterozygosity & Florida panther \texorpdfstring{\ensuremath{\times}}{x} Texas cougar \\
\textbf{Ancient DNA} & Reconstruct historical diversity; identify lost
alleles & Woolly mammoth \texorpdfstring{\ensuremath{\to}}{->} mammoth de-extinction project \\
\end{longtable}
}

\begin{figure}[htbp]
\centering
\includegraphics[width=0.82\linewidth,height=4.2in,keepaspectratio]{../figures/mermaid_inline/inline_mermaid_0190_3a6a9842f9c3.png}
\caption{Sequence diagram for Conservation Genomics showing ordered interaction among Endangered Population, Whole-Genome Sequencing, Genomic Analysis, and Management Decision.}
\end{figure}

\subsubsection{Assisted Gene Flow and Assisted
Migration}\label{assisted-gene-flow-and-assisted-migration}

\textbf{Assisted gene flow:} Intentionally moving alleles from adapted
to non-adapted populations of the same species: - Coral: transplanting
heat-tolerant \emph{Acropora} genotypes to vulnerable reefs (Palumbi et
al.~2014, \emph{Science}) - Trees: moving southern seed sources to
northern planting sites to pre-adapt to warming (assisted migration of
\emph{Pinus contorta} in British Columbia)

\textbf{Assisted migration (managed relocation):} Moving entire species
to new locations outside their historical range where future climate is
suitable. Controversial due to invasion risk.

\textbf{De-extinction:} Using genomic engineering to resurrect
functional analogues of extinct species: - \textbf{Woolly mammoth
project} (Colossal Biosciences): Editing Asian elephant genome with
mammoth cold-adaptation genes (TRPV3, haemoglobin variants, fat
deposition); goal: produce cold-adapted elephants for Arctic rewilding -
\textbf{Thylacine project} (Colossal): Reconstructing thylacine genome
from museum specimens - \textbf{Ethical debates:} Should resources be
spent on de-extinction rather than preventing current extinctions?

\subsubsection{Climate-Adaptive
Management}\label{climate-adaptive-management}

Traditional conservation aims to preserve historical baselines.
\textbf{Climate-adaptive management} accepts that ecosystems are
changing and manages for future conditions:

The hard part is deciding when to resist change, when to help a system
transition, and when to redirect effort toward ecosystem function rather
than historical composition. Coral reefs illustrate the tradeoff:
assisted gene flow, symbiont management, shading, and local
water-quality improvement can raise local resilience, but they do not
remove the external heat and acidification drivers. Climate-adaptive
management is therefore a portfolio problem, pairing local stress
reduction with monitoring, emissions context, social feasibility, and
explicit thresholds for changing strategy
\citep{noaa2025coralbleaching,ipbes2024transformative}.

{\def\LTcaptype{none} % do not increment counter
\begin{longtable}[]{@{}
  >{\raggedright\arraybackslash}p{(\linewidth - 4\tabcolsep) * \real{0.3478}}
  >{\raggedright\arraybackslash}p{(\linewidth - 4\tabcolsep) * \real{0.3478}}
  >{\raggedright\arraybackslash}p{(\linewidth - 4\tabcolsep) * \real{0.3043}}@{}}
\toprule\noalign{}
\begin{minipage}[b]{\linewidth}\raggedright
Approach
\end{minipage} & \begin{minipage}[b]{\linewidth}\raggedright
Strategy
\end{minipage} & \begin{minipage}[b]{\linewidth}\raggedright
Example
\end{minipage} \\
\midrule\noalign{}
\endhead
\bottomrule\noalign{}
\endlastfoot
\textbf{Protect climate refugia} & Identify and protect areas where
climate change impacts are minimised & Deep canyons, north-facing
slopes, groundwater-fed wetlands \\
\textbf{Increase connectivity} & Corridors for range shifts &
Continental-scale corridors (Y2Y, Florida Wildlife Corridor) \\
\textbf{Managed retreat} & Accept loss of some habitats; invest in new
ones & Managed realignment of coastal defences; creating new
marshland \\
\textbf{Reduce non-climate stressors} & Remove pollution, invasive
species to increase resilience & Reducing sediment runoff on coral
reefs \\
\end{longtable}
}

\begin{center}\rule{0.5\linewidth}{0.5pt}\end{center}

\subsubsection{Visualizing Reserve Design
Principles}\label{visualizing-reserve-design-principles}

\begin{figure}[htbp]
\centering
\includegraphics[width=0.82\linewidth,height=4.2in,keepaspectratio]{../figures/mermaid_inline/inline_mermaid_0191_b8f5c21f48b4.png}
\caption{Flowchart for Visualizing Reserve Design Principles: Connected \& Large High gene flow Lower extinction risk Resilient to shifts, Isolated \& Small Edge effects dominate Inbreeding risk Vulnerable to stochasticity, Single Large Protects wide-ranging apex predators, and Several Small Protects rare plants in unique micro-habitats form the diagram's primary path or branches.}
\end{figure}

\begin{center}\rule{0.5\linewidth}{0.5pt}\end{center}

\subsection{7B Climate Velocity, Refugia, and Range
Shifts}\label{b-climate-velocity-refugia-and-range-shifts}

\citet{ripple2012} and Loarie et al.~(2009, \emph{Nature}) established
\textbf{climate velocity} as the central organising metric for
climate-adaptive conservation. The velocity is the speed at which a
given thermal envelope moves across the landscape, and species must
disperse at least this fast to track their climatic niche.

\begin{equation}
v_{\text{climate}} = \frac{\partial T/\partial t}{|\nabla T|}
\label{eq:unit_X_climate_velocity}
\end{equation}

The numerator is the rate of climate change at a point (°C/yr); the
denominator is the local spatial temperature gradient (°C/km).
Rearranging gives velocity in km/yr.

{\def\LTcaptype{none} % do not increment counter
\begin{longtable}[]{@{}
  >{\raggedright\arraybackslash}p{(\linewidth - 4\tabcolsep) * \real{0.1875}}
  >{\raggedright\arraybackslash}p{(\linewidth - 4\tabcolsep) * \real{0.3333}}
  >{\raggedright\arraybackslash}p{(\linewidth - 4\tabcolsep) * \real{0.4792}}@{}}
\toprule\noalign{}
\begin{minipage}[b]{\linewidth}\raggedright
Landscape
\end{minipage} & \begin{minipage}[b]{\linewidth}\raggedright
Spatial gradient
\end{minipage} & \begin{minipage}[b]{\linewidth}\raggedright
Velocity at +0.03 °C/yr
\end{minipage} \\
\midrule\noalign{}
\endhead
\bottomrule\noalign{}
\endlastfoot
Mountain (e.g., Andes) & 5 °C/km vertically & 0.006 km/yr (6 m/yr ---
slow) \\
Boreal foothills & 0.5 °C/km & 0.06 km/yr \\
Temperate plains & 0.05 °C/km & 0.6 km/yr \\
Amazon basin (flat) & 0.005 °C/km & 6 km/yr \\
Subtropical plains & 0.003 °C/km & 10 km/yr \\
\end{longtable}
}

The flattest landscapes generate the fastest velocities --- exactly
where species cannot escape upslope. \textbf{Migration lags} therefore
cluster in lowland tropical and temperate plains: trees often disperse
at roughly 0.1--0.5 km/yr, which can be 10--100\texorpdfstring{\ensuremath{\times}}{x} too slow for tropical
lowlands but adequate for some mountain settings.

\subsubsection{Refugia and Their
Identification}\label{refugia-and-their-identification}

\textbf{Climate refugia} are locations where the local climate changes
more slowly than the regional average, providing temporary or permanent
havens for species unable to track climate velocity. Recognised refugia
types:

{\def\LTcaptype{none} % do not increment counter
\begin{longtable}[]{@{}
  >{\raggedright\arraybackslash}p{(\linewidth - 4\tabcolsep) * \real{0.1905}}
  >{\raggedright\arraybackslash}p{(\linewidth - 4\tabcolsep) * \real{0.4286}}
  >{\raggedright\arraybackslash}p{(\linewidth - 4\tabcolsep) * \real{0.3810}}@{}}
\toprule\noalign{}
\begin{minipage}[b]{\linewidth}\raggedright
Type
\end{minipage} & \begin{minipage}[b]{\linewidth}\raggedright
Mechanism
\end{minipage} & \begin{minipage}[b]{\linewidth}\raggedright
Examples
\end{minipage} \\
\midrule\noalign{}
\endhead
\bottomrule\noalign{}
\endlastfoot
\textbf{Topographic} & Cold-air pooling, north-facing slopes, deep
canyons buffer temperature & Klamath-Siskiyou (NW USA); Mediterranean
canyons \\
\textbf{Hydrological} & Groundwater-fed streams, springs, lakes maintain
cool temperatures & Spring-fed cold-water salmonid habitat \\
\textbf{Microclimatic} & Forest interior buffering (--2 to --5 °C
vs.~open canopy) & Cathedral-grove old growth; tropical understory \\
\textbf{Elevational} & Mountains hold residual cool habitat above
warming lowlands & Alpine ecosystems globally \\
\textbf{Coastal upwelling} & Cold-water upwelling buffers warming
surface oceans & California Current; Benguela; Peru \\
\textbf{Cryptic / paleo} & Sites that remained inhabitable through
Pleistocene glaciations & Many high-biodiversity hotspots map onto past
refugia \\
\end{longtable}
}

Refugia identification is a major effort of climate-adaptive
conservation planning: high-resolution climate downscaling (\texorpdfstring{\ensuremath{\leq}}{<=} 1 km
grids) plus species distribution modelling reveals locations where
target species can persist longest, prioritising them for protection.
The \textbf{Adapt West} initiative (Conservation Biology Institute) maps
climate refugia for North American species; analogous tools exist in
Australia, Europe, and parts of South America.

\begin{quote}
\textbf{Concept Check:} Two endangered birds occupy similar habitat:
species A in the high Andes (5 °C/km gradient) and species B on the
Brazilian cerrado (flat). With identical thermal requirements and
warming rates, which species faces the greater extinction risk from
climate change, and what reserve-design implications follow?
\end{quote}

\begin{center}\rule{0.5\linewidth}{0.5pt}\end{center}

\subsection{7C Ecosystem-Based Adaptation: Using Nature for Climate
Resilience}\label{c-ecosystem-based-adaptation-using-nature-for-climate-resilience}

\textbf{Ecosystem-Based Adaptation (EbA)} uses biodiversity and
ecosystem services as part of an overall adaptation strategy to help
people adapt to climate change. The framework --- formalised by the CBD
and IUCN in 2009 --- sits at the intersection of conservation and human
resilience: protecting and restoring ecosystems to buffer communities
against floods, storms, droughts, heat waves, and rising seas, while
simultaneously sequestering carbon and conserving biodiversity.

\subsubsection{Major EbA Strategies}\label{major-eba-strategies}

{\def\LTcaptype{none} % do not increment counter
\begin{longtable}[]{@{}
  >{\raggedright\arraybackslash}p{(\linewidth - 4\tabcolsep) * \real{0.2963}}
  >{\raggedright\arraybackslash}p{(\linewidth - 4\tabcolsep) * \real{0.3333}}
  >{\raggedright\arraybackslash}p{(\linewidth - 4\tabcolsep) * \real{0.3704}}@{}}
\toprule\noalign{}
\begin{minipage}[b]{\linewidth}\raggedright
Strategy
\end{minipage} & \begin{minipage}[b]{\linewidth}\raggedright
Mechanism
\end{minipage} & \begin{minipage}[b]{\linewidth}\raggedright
Co-benefit
\end{minipage} \\
\midrule\noalign{}
\endhead
\bottomrule\noalign{}
\endlastfoot
\textbf{Mangrove restoration} & Wave attenuation; coastal sediment
stabilisation; storm surge buffering (each 100 m of mangrove reduces
wave height about 20 \%) & Carbon sequestration (3--5 \texorpdfstring{\ensuremath{\times}}{x}
tropical forest density); fisheries nursery \\
\textbf{Coral reef restoration} & Wave dissipation (reefs reduce wave
energy 70--95 \%); coastal protection valued at
about \$10,000/km/yr & Tourism; fisheries; biodiversity \\
\textbf{Urban green infrastructure} & Tree canopy reduces urban heat
island by 2--5 °C; permeable surfaces reduce flooding & Air quality;
mental health; physical activity \\
\textbf{Wetland conservation} & Floodplain water storage; nutrient
filtration & Wildlife habitat; recreation \\
\textbf{Reforestation / agroforestry} & Soil retention; rainfall
recycling; microclimate stabilisation & Carbon storage; smallholder
livelihoods \\
\textbf{Drought-tolerant agroecosystems} & Diverse cropping, perennial
roots, soil organic matter & Food security; reduced N\(_2\)O \\
\textbf{Living shorelines} & Marsh and oyster reef construction
replacing seawalls & Habitat; better long-term storm performance than
hard infrastructure \\
\end{longtable}
}

\subsubsection{Quantified Performance}\label{quantified-performance}

EbA projects often outperform engineered alternatives on cost and
resilience:

\begin{itemize}
\tightlist
\item
  \textbf{Mangroves vs.~seawalls.} Coastal protection from mangroves in
  the Philippines averages \$94/m of coastline; equivalent seawall
  construction averages \$2,000/m, with no co-benefits and high
  maintenance (Beck et al.~2018, \emph{Nat. Commun.}).
\item
  \textbf{Hurricane Katrina (2005).} Wetland loss in the Mississippi
  Delta increased storm surge by about 25\%; restoring 100
  km\(^2\) of marsh would reduce surge damage in vulnerable parishes by
  about 10--15 \% (Day et al.~2007, \emph{Science}).
\item
  \textbf{2003 European heat wave.} Cities with dense urban green
  infrastructure (e.g., Vienna) experienced about 30 \% fewer
  heat-related deaths than comparable cities with less greening
  (Smargiassi et al.~2009).
\end{itemize}

\subsubsection{Limits of EbA}\label{limits-of-eba}

EbA is not a panacea. Limitations:

\begin{itemize}
\tightlist
\item
  \textbf{Cannot substitute for emissions reduction.} No amount of
  coastal restoration prevents the catastrophic outcomes of 4 °C
  warming.
\item
  \textbf{Climate change itself is undermining the adapters.} Coral
  bleaching kills the reefs we want to use for coastal protection;
  mangroves face sea-level rise faster than they can build sediment.
\item
  \textbf{Equity and tenure issues.} EbA projects fail when they ignore
  Indigenous land rights or displace communities.
\end{itemize}

EbA is best understood as one component of a multi-layered adaptation
portfolio that also includes hard infrastructure, social systems, and
aggressive emissions cuts.

\begin{quote}
\texorpdfstring{}{} \textbf{Clinical Connection --- Heat Action Plans and Public Health.}
Following the 2003 European heat wave (about 70,000 excess
deaths) and the 2010 Russian heat wave (about 55,000 excess
deaths), cities increasingly couple ecological strategies (urban tree
planting, cool roofs, restored wetlands and parks) with public health
interventions (heat warning systems, cooling shelters,
vulnerable-population outreach). Ahmedabad, India implemented South
Asia's first Heat Action Plan in 2013 after a 2010 heat wave killed
\textgreater{} 1,300 residents; subsequent years saw about 30\%
reductions in heat-related mortality despite continued warming. Mexico
City, Phoenix, and Chicago have similar plans. These One Health
interventions integrate biodiversity conservation, urban ecology, and
human health --- operationalising the EbA framework.
\end{quote}

\begin{quote}
\textbf{Concept Check:} A coastal municipality must spend its limited
adaptation budget on protecting 20 km of shoreline from increasing storm
surge. (a) Compare the cost-benefit profile of seawall construction
vs.~mangrove restoration. (b) Why is EbA not always the best choice ---
what conditions favour hard infrastructure? (c) How does coral bleaching
complicate the math?
\end{quote}

\begin{center}\rule{0.5\linewidth}{0.5pt}\end{center}

\subsection{International Conservation
Frameworks}\label{international-conservation-frameworks}

\subsubsection{Key Agreements}\label{key-agreements}

{\def\LTcaptype{none} % do not increment counter
\begin{longtable}[]{@{}
  >{\raggedright\arraybackslash}p{(\linewidth - 6\tabcolsep) * \real{0.2812}}
  >{\raggedright\arraybackslash}p{(\linewidth - 6\tabcolsep) * \real{0.1250}}
  >{\raggedright\arraybackslash}p{(\linewidth - 6\tabcolsep) * \real{0.1875}}
  >{\raggedright\arraybackslash}p{(\linewidth - 6\tabcolsep) * \real{0.4062}}@{}}
\toprule\noalign{}
\begin{minipage}[b]{\linewidth}\raggedright
Framework
\end{minipage} & \begin{minipage}[b]{\linewidth}\raggedright
Year
\end{minipage} & \begin{minipage}[b]{\linewidth}\raggedright
Target
\end{minipage} & \begin{minipage}[b]{\linewidth}\raggedright
Status (2024)
\end{minipage} \\
\midrule\noalign{}
\endhead
\bottomrule\noalign{}
\endlastfoot
\textbf{CITES} & 1975 & Regulate international wildlife trade & Parties
and listed taxa change as appendices are amended; cite the live
appendices for current counts \citep{cites2026appendices} \\
\textbf{CBD} (Convention on Biological Diversity) & 1992 & Conservation,
sustainable use, equitable sharing & 196 parties (US has not
ratified) \\
\textbf{Aichi Targets} & 2010 & 20 targets for 2020 & Most targets
missed \\
\textbf{Kunming-Montreal Framework} (30x30) & 2022 & Protect 30\% of
land and ocean by 2030 & Currently about 17\% land,
about 8\% ocean protected \\
\textbf{Paris Agreement} & 2015 & Limit warming to 1.5-2\(^\circ\)C &
Tracking toward 2.5-3\(^\circ\)C \\
\end{longtable}
}

\subsubsection{The 30x30 Target}\label{the-30x30-target}

The Kunming-Montreal Global Biodiversity Framework (GBF, 2022) set a
target to protect 30\% of land and 30\% of ocean by 2030
\citep{cbd2022kunmingmontreal}. Current status:

\begin{equation}
\text{Terrestrial: } 16.7\% \xrightarrow{\text{need}} 30\% \text{ (add 20 million km}^2\text{)}
\label{eq:biomes_and_conservation_6}
\end{equation}

\begin{equation}
\text{Marine: } 8.3\% \xrightarrow{\text{need}} 30\% \text{ (add 78 million km}^2\text{)}
\label{eq:biomes_and_conservation_7}
\end{equation}

\textbf{Challenges:} - Quality vs.~quantity (many protected areas are
``paper parks'' with no enforcement) - Indigenous and local community
rights (30x30 must not displace indigenous peoples) - Representativeness
(current protected areas biased toward ``rock and ice'' ---
high-altitude, low-productivity areas) - Financial gap: estimated \$700
billion/yr needed; current spending about \$100 billion/yr

\subsubsection{Payments for Ecosystem
Services}\label{payments-for-ecosystem-services}

\textbf{REDD+} (Reducing Emissions from Deforestation and Forest
Degradation): International carbon market payments incentivise
forest-rich developing nations to avoid deforestation. \$5 billion in
pledges as of 2024, but implementation faces challenges of
additionality, leakage, permanence, and measurement.

\begin{quote}
\textbf{Concept Check:} The Aichi Targets (2010-2020) were largely
unmet. What factors contributed to their failure, and how does the
Kunming-Montreal Framework attempt to address these shortcomings?
\end{quote}

\begin{center}\rule{0.5\linewidth}{0.5pt}\end{center}

\subsection{Worked Example}\label{worked-example-8}

\textbf{Problem:} An ecologist studies bird species on an archipelago.
The species-area relationship for this region follows the equation
\(S = cA^z\), where \(c = 10\) and \(z = 0.3\). Island X has an area
(\(A\)) of \(1,000\text{ km}^2\). Calculate the expected number of bird
species (\(S\)) on Island X.

\textbf{Solution:}

\textbf{Step 1. Identify the given variables.} - Area (\(A\)) =
\(1,000\text{ km}^2\) - Constant (\(c\)) = 10 - Exponent (\(z\)) = 0.3

\textbf{Step 2. Substitute the values into the species-area equation.}
\[S = cA^z \label{eq:unit_X_biomes_and_conservation_item_1}\]

\[S = 10 \times (1000)^{0.3} \label{eq:unit_X_biomes_and_conservation_item_2}\]

\textbf{Step 3. Calculate the exponent.} Notice that \(1000 = 10^3\).
Therefore:
\[(10^3)^{0.3} = 10^{(3 \times 0.3)} = 10^{0.9} \label{eq:unit_X_biomes_and_conservation_item_3}\]

Let's approximate \(10^{0.9}\). Since \(10^1 = 10\) and
\(10^{0.5} \approx 3.16\), \(10^{0.9} \approx 7.94\). (Exact value:
\(10^{0.9} \approx 7.943\))

\textbf{Step 4. Solve for \(S\).}
\[S = 10 \times 7.943 \label{eq:unit_X_biomes_and_conservation_item_4}\]

\[S \approx 79.43 \label{eq:unit_X_biomes_and_conservation_item_5}\]

\textbf{Answer:} Island X is expected to support approximately
\textbf{79 bird species}.

\begin{center}\rule{0.5\linewidth}{0.5pt}\end{center}

\subsection{Worked Example: Effective Population Size and the 50/500
Rule}\label{worked-example-effective-population-size-and-the-50500-rule}

\textbf{Problem:} A captive recovery herd of an endangered ungulate has
a census size of \(N = 120\) individuals, but primarily \(N_m = 20\)
males and \(N_f = 100\) females actually breed. Calculate the effective
population size \(N_e\), the per-generation inbreeding rate
\(\Delta F\), and evaluate the herd against the 50/500 rule for
short-term and long-term viability.

\textbf{Solution:}

\textbf{Step 1. Identify the given variables.} - Breeding males
(\(N_m\)) = 20 - Breeding females (\(N_f\)) = 100 - Census size (\(N\))
= 120

\textbf{Step 2. Substitute into the unequal-sex-ratio effective-size
equation.}
\[N_e = \frac{4 N_m N_f}{N_m + N_f} \label{eq:unit_X_biomes_and_conservation_item_6}\]

\[N_e = \frac{4 \times 20 \times 100}{20 + 100} = \frac{8{,}000}{120} \approx 66.7 \label{eq:unit_X_biomes_and_conservation_item_7}\]

\textbf{Step 3. Compute the per-generation inbreeding rate.}
\[\Delta F = \frac{1}{2 N_e} = \frac{1}{2 \times 66.7} \approx 0.0075 = 0.75\%\text{ per generation} \label{eq:unit_X_biomes_and_conservation_item_8}\]

\textbf{Step 4. Evaluate against the 50/500 rule.} With
\(N_e \approx 67\), the herd clears the short-term threshold
(\(N_e \geq 50\)) that guards against rapid inbreeding depression, but
it falls far short of the long-term threshold (\(N_e \geq 500\)) needed
to retain additive genetic variance for future adaptation. The skewed
sex ratio is the cause: the effective size is barely over half the
census size, even though 120 animals are present.

\textbf{Answer:} \(N_e \approx 67\) and \(\Delta F \approx 0.75\%\) per
generation. The herd is demographically defensible in the short term but
is genetically committed to losing roughly 0.75\% of its heterozygosity
each generation; without managed gene flow or population growth it
remains at elevated long-term extinction risk because it cannot meet the
\(N_e \geq 500\) adaptive-potential standard.

\begin{center}\rule{0.5\linewidth}{0.5pt}\end{center}

\subsubsection{Minimum Viable Population and Genetic Rescue: Case
Studies from Conservation
Biology}\label{minimum-viable-population-and-genetic-rescue-case-studies-from-conservation-biology}

\textbf{Minimum viable population (MVP)} is defined as the smallest
population size with \(\geq 95\ \%\) probability of persistence over 100
generations, incorporating demographic stochasticity, environmental
variance, and genetic deterioration from inbreeding depression and loss
of adaptive variation. The classical rule-of-thumb is the \textbf{``50 /
500 rule''} (Franklin 1980): \(N_e \geq 50\) to avoid short-term
inbreeding depression, \(N_e \geq 500\) to retain long-term adaptive
potential against mutation-drift balance. Modern refinements (Frankham,
Brook \& Bradshaw 2014) revise these upward --- \(N_e \geq 100\) and
\(N_e \geq 1000\) --- as the more defensible thresholds given empirical
estimates of mutational-variance and inbreeding-load effects.

\textbf{Case 1 --- The Arabian Oryx (\emph{Oryx leucoryx})}. Extinct in
the wild by 1972 after hunting reduced the Arabian population to zero,
the species survived primarily in a 9-founder ``World Herd'' assembled
at Phoenix Zoo. Careful pedigree management and reintroduction programs
(Oman 1982, Saudi Arabia 1989, Jordan 1997) have grown the wild
population to about 1250 (IUCN 2017, Vulnerable). The genetic
cost: current population heterozygosity is about 70 \% of the
pre-bottleneck estimate, and most 6 microsatellite loci have at most 2
alleles --- a cautionary example that demographic recovery does not
erase the genetic signature of a bottleneck.

\textbf{Case 2 --- The California Condor (\emph{Gymnogyps
californianus})}. From a 1982 low of 22 individuals (the entire species
in captivity by 1987), a captive-breeding / release program has grown
the population to \textgreater{} 500 (half wild, half captive). Genomic
sequencing (Ralls et al.~2020) detected \textbf{23 known deleterious
alleles} in the founders, two fixed during the bottleneck. The ongoing
challenge is not census size but \textbf{lead toxicosis} from ingestion
of lead ammunition in carcasses --- a clear example of \emph{why}
habitat and regulatory intervention remain necessary even when
demographic targets are met.

\textbf{Genetic rescue} is the deliberate introduction of unrelated
individuals into a small, inbred population to restore heterozygosity.
The textbook case is the \textbf{Florida panther (\emph{Puma concolor
coryi})}: by the 1990s, about 30 individuals remained, with striking
inbreeding-depression \hyperref[gl:phenotype]{\textbf{phenotype}}s
(cryptorchidism, cardiac defects). Introducing 8 female Texas cougars in
1995 doubled heterozygosity within one generation; population is now
about 200 with reduced deformities and increased reproductive
success (Johnson et al., \emph{Science} 2010). The Isle Royale wolves
(introduction of 19 wolves 2018--2021) represent the latest live
experiment. Genetic rescue is \emph{not} a comprehensive solution --- it
can disrupt locally adapted gene complexes, introduce maladapted
alleles, and dilute evolutionary uniqueness (the Florida panther now
bears about 5 \% Texas-cougar ancestry). The decision is
typically a balance of short-term demographic benefit against long-term
evolutionary coherence.

\begin{center}\rule{0.5\linewidth}{0.5pt}\end{center}

\subsection{Current Evidence and Frontier
Biology}\label{current-evidence-and-frontier-biology-37}

For \textbf{Biomes and Conservation Biology}, frontier biology belongs
inside the evidence logic of the chapter. Ecology and conservation
decisions increasingly combine field data, remote sensing, community
knowledge, model uncertainty, and explicit values. The core reading
question is this: conservation claims must separate ecological evidence,
social values, feasibility, and uncertainty in tradeoffs.

\begin{itemize}
\tightlist
\item
  \textbf{What to verify:} identify the observation, model, assay, or
  dataset that would make the claim stronger or weaker.
\item
  \textbf{What to qualify:} state the scale, organism, cell type,
  environmental condition, or population where the claim is expected to
  hold.
\item
  \textbf{What to compare:} test at least one alternative explanation,
  baseline, or null model before treating the pattern as causal.
\item
  \textbf{What to cite:} distinguish primary evidence, review synthesis,
  public dataset, and institutional guidance; for recent or numeric
  claims, prefer the source closest to the measurement and state what
  has changed since it was published.
\end{itemize}

Use biodiversity metrics carefully: population indices, extinction risk
categories, ecosystem services, and management targets answer different
questions
\citep{ipbes2019global,ipbes2024transformative,wwf2024livingplanet,iucn2025redlist,fao2024sofia}.

\textbf{Source practice:} For conservation claims, cite assessment
sources and state whether the evidence is a population index,
extinction-risk assessment, ecosystem-service valuation, satellite
product, or policy synthesis
\citep{ipbes2024transformative,noaa2025coralbleaching,fao2025sofi}.

\subsubsection{Current Evidence Map: Conservation Decision Evidence
Chain}\label{current-evidence-map-conservation-decision-evidence-chain}

\begin{figure}[htbp]
\centering
\includegraphics[width=0.82\linewidth,height=4.2in,keepaspectratio]{../figures/mermaid_inline/inline_mermaid_0192_45e01ed0a495.png}
\caption{Coral assisted evolution and IPBES transformative-change examples show why conservation decisions need evidence, values, feasibility, and monitoring in the same frame \textbackslash{}citep\{ipbes2024transformative,noaa2025coralbleaching,strader2022coralheat\}.}
\end{figure}

\subsection{Key Terms}\label{key-terms-37}

{\def\LTcaptype{none} % do not increment counter
\begin{longtable}[]{@{}
  >{\raggedright\arraybackslash}p{(\linewidth - 2\tabcolsep) * \real{0.2857}}
  >{\raggedright\arraybackslash}p{(\linewidth - 2\tabcolsep) * \real{0.7143}}@{}}
\toprule\noalign{}
\begin{minipage}[b]{\linewidth}\raggedright
Term
\end{minipage} & \begin{minipage}[b]{\linewidth}\raggedright
Definition
\end{minipage} \\
\midrule\noalign{}
\endhead
\bottomrule\noalign{}
\endlastfoot
\textbf{Biome} & Large-scale ecological zone defined by climate and
\hyperref[gl:dominant]{\textbf{dominant}} vegetation type \\
\textbf{Net primary productivity (NPP)} & Carbon fixed by plants minus
autotrophic respiration; g C/m\(^2\)/yr \\
\textbf{IUCN Red List} & Global system for evaluating extinction risk
using quantitative criteria \\
\textbf{Minimum viable population (MVP)} & Smallest population with at
least 95\% persistence probability for at least 100 generations \\
\textbf{50/500 rule} & \(N_e \geq 50\) short-term; \(N_e \geq 500\)
long-term viability (Franklin/Soule) \\
\textbf{Effective population size (\(N_e\))} & Genetically effective
number of breeding individuals; typically \(0.1-0.25 \times N\) \\
\textbf{Extinction debt} & Delayed extinction of species committed to
disappear due to past habitat loss \\
\textbf{SLOSS debate} & Trade-off between single large vs.~several small
reserves \\
\textbf{Rewilding} & Restoration of ecological processes via
reintroduction of lost species \\
\textbf{Population Viability Analysis (PVA)} & Stochastic simulation of
extinction probability \\
\textbf{Conservation genomics} & Application of whole-genome analysis to
conservation management \\
\textbf{Assisted gene flow} & Intentional translocation of alleles from
adapted to non-adapted populations \\
\textbf{Genetic rescue} & Introduction of unrelated individuals to
reverse inbreeding depression \\
\textbf{Climate velocity} & Speed at which climate conditions shift
geographically (km/yr) \\
\textbf{Edge effect} & Microclimatic and biotic changes at habitat
fragment boundaries \\
\textbf{Coral bleaching} & Thermal stress-induced expulsion of
zooxanthellae; loss of coral energy supply \\
\textbf{30x30} & Target to protect 30\% of land and ocean by 2030
(Kunming-Montreal GBF) \\
\textbf{De-extinction} & Using genomic engineering to resurrect
functional analogues of extinct species \\
\textbf{Novel biome} & Ecosystem with no historical analog, arising from
climate change \\
\textbf{IPBES} & Intergovernmental Science-Policy Platform on
Biodiversity and Ecosystem Services; ``IPCC for biodiversity'' \\
\textbf{Climate refugia} & Locations where local climate changes more
slowly than the regional average; havens for slow-dispersing species \\
\textbf{Ecosystem-based adaptation (EbA)} & Use of biodiversity and
ecosystem services to help people adapt to climate change \\
\textbf{Living shorelines} & Marsh, oyster reef, and mangrove
infrastructure replacing seawalls for coastal protection \\
\textbf{Living Planet Index} & WWF metric of mean change in monitored
vertebrate populations since 1970 \\
\textbf{Heat Action Plan} & Coupled public-health and ecological
response to extreme heat events \\
\textbf{Genetic load} & Accumulated deleterious alleles; heightened in
bottlenecked populations; affects translocation strategies \\
\textbf{Runs of homozygosity (ROH)} & Long stretches of homozygous DNA
indicating recent inbreeding; key conservation-genomic metric \\
\end{longtable}
}

\begin{center}\rule{0.5\linewidth}{0.5pt}\end{center}

\subsection{Review Questions}\label{review-questions-37}

\begin{enumerate}
\def\labelenumi{\arabic{enumi}.}
\item
  Using the species-area relationship \(S = cA^z\) with \(c = 6\) and
  \(z = 0.28\), calculate species richness for a 50,000 ha intact forest
  and after fragmentation into 5 patches of 1,000 ha each. What fraction
  of species are committed to extinction via the fragmentation-driven
  extinction debt? How long might this debt take to be ``paid'' for
  long-lived tree species?
\item
  A population of sea turtles has \(N_e = 35\) and current
  heterozygosity \(H = 0.60\). (a) Apply the 50/500 rule to evaluate its
  short-term and long-term viability. (b) Calculate how many generations
  until \(H < 0.40\) using \(H_t = H_0(1 - 1/2N_e)^t\). (c) What genetic
  rescue strategy could improve the prognosis?
\item
  The Amazon has currently lost about 17\% of its original
  forest cover. Lovejoy and Nobre's tipping point threshold is
  about 20-25\%. (a) Explain the tipping mechanism (why is it
  non-linear?). (b) What conservation interventions are needed in terms
  of scale, urgency, and geographic targeting? (c) How does the
  transpiration recycling mechanism create a positive feedback loop
  during deforestation?
\item
  Explain the difference between demographic, environmental, genetic,
  and catastrophic stochasticity in an MVP context. Which was most
  important in the Florida panther case, and how was it mitigated?
  Calculate the inbreeding rate (\(\Delta F\)) for the panther
  population at \(N_e = 10\) vs.~after genetic rescue at \(N_e = 50\).
\item
  Compare rewilding and assisted gene flow as conservation strategies
  for climate change adaptation. For each, give one specific
  quantitative example (species, location, outcome metric, timescale).
\item
  A conservation biologist must design a reserve system for a large
  carnivore with home range = 200 km\(^2\) and an MVP of \(N_e = 500\)
  (\(N/N_e\) ratio = 4). Calculate: (a) minimum total census population
  needed, (b) minimum reserve area assuming 50\% habitat overlap between
  individuals, (c) whether a single reserve or corridor-connected system
  is more feasible. Design the reserve system.
\item
  The Great Barrier Reef has experienced 6 mass bleaching events in 10
  years. Using the degree heating weeks (DHW) metric, explain the
  mechanism of bleaching. At what level of global warming does coral
  reef survival become highly improbable? What is the economic value at
  stake (cite Deloitte or similar valuation studies)?
\item
  Evaluate the 30x30 target (Kunming-Montreal Framework). Calculate the
  additional area of land and ocean that must be protected to reach 30\%
  from current levels. What are the three most significant
  implementation challenges, and how might they be addressed?
\item
  The woolly mammoth de-extinction project aims to create cold-adapted
  elephants for Arctic tundra. (a) What ecological function would they
  serve? (b) How would they slow permafrost thaw? (c) What are the
  strongest ecological and ethical arguments for and against this
  approach?
\item
  Using the concept of climate velocity, explain why montane species are
  at lower risk of range contraction than lowland species. Calculate
  climate velocity for a flat landscape with temperature change =
  0.03\(^\circ\)C/yr and spatial gradient = 0.005\(^\circ\)C/km, vs.~a
  mountain with spatial gradient = 5\(^\circ\)C/km. What are the
  implications for conservation prioritisation?
\item
  Compute \breaktt{species\_area\_relationship} for
  \(A \in \{1, 10, 100\}\) km\(^2\) with \(c=2\), \(z=0.25\); discuss
  reserve sizing.
\item
  Contrast \textbf{inbreeding depression} vs.~\textbf{outbreeding
  depression} in translocation programs.
\item
  The IPBES 2019 Global Assessment ranked land-use change as the
  dominant terrestrial driver of biodiversity loss (about 30
  \%), with climate change at about 14 \% but rising. (a) Why
  is this \emph{current} ranking expected to invert by 2050? (b) How
  should conservation priority-setting respond to this projected shift?
  (c) The Transformative Change Assessment (IPBES 2024) argues that
  incremental conservation cannot meet biodiversity targets. What three
  systemic reforms does it recommend, and why are they outside
  traditional conservation biology?
\item
  Compare two endangered birds: species A in the high Andes (5 °C/km
  gradient) and species B on the Brazilian cerrado (0.005 °C/km
  gradient), with both facing 0.03 °C/yr warming. (a) Calculate climate
  velocity for each. (b) Predict their relative extinction risks. (c)
  How does this calculation re-prioritise conservation funding between
  the two regions?
\item
  A coastal municipality is choosing between (a) a 20-km concrete
  seawall (\$2,000/m, 50-yr lifetime, no co-benefits, requires periodic
  rebuilding) and (b) a 20-km mangrove restoration project (\$94/m
  initial, low maintenance, multiple co-benefits, but vulnerable to
  sea-level rise faster than mangrove sediment accretion). Calculate
  total costs over 50 years for each, and identify under what
  climate-change scenario each option is preferred.
\item
  Explain how the 50/500 rule, conservation genomics (runs of
  homozygosity, genetic load), and population viability analysis
  converge to inform genetic rescue decisions for the Florida panther.
  Why was the 1995 Texas-cougar introduction successful, and what risks
  (e.g., outbreeding depression) had to be balanced?
\item
  Define a climate refugium and list four types (topographic,
  hydrological, microclimatic, elevational). Why is identifying refugia
  critical for 21st-century reserve design? What datasets (climate
  models, species distribution models, paleoecology) does refugium
  identification require?
\item
  Critically evaluate the 30x30 target. (a) Calculate the additional
  area of land and ocean needed to reach 30 \% from current levels. (b)
  Why is \emph{quality} of protection (effective management, no ``paper
  parks'') at least as important as quantity? (c) How can 30x30 be
  implemented without violating Indigenous land rights, given that
  about 80 \% of remaining biodiversity is on
  Indigenous-managed lands?
\end{enumerate}

\begin{center}\rule{0.5\linewidth}{0.5pt}\end{center}

\subsection{Further Reading and Source
Notes}\label{further-reading-and-source-notes-32}

\begin{itemize}
\tightlist
\item
  Soule (1980). Thresholds for Survival: Maintaining Fitness and
  Evolutionary Potential. \emph{Conservation Biology: An
  Evolutionary-Ecological Perspective}.
\item
  Frankham (1995). Effective population size/adult population size
  ratios in wildlife: A review. \emph{Genetical Research}, 66.
\item
  Frankham (2014). Conservation genetics: Setting the scene.
  \emph{Philosophical Transactions of the Royal Society B}, 369.
\end{itemize}

\begin{center}\rule{0.5\linewidth}{0.5pt}\end{center}

\subsection{Computational Bridge}\label{computational-bridge-34}

The power-law S--A curve is a one-liner:

\begin{Shaded}
\begin{Highlighting}[]
\ImportTok{from}\NormalTok{ biology.ecology }\ImportTok{import}\NormalTok{ species\_area\_relationship}

\ControlFlowTok{for}\NormalTok{ a }\KeywordTok{in}\NormalTok{ (}\FloatTok{1.0}\NormalTok{, }\FloatTok{10.0}\NormalTok{, }\FloatTok{100.0}\NormalTok{):}
    \BuiltInTok{print}\NormalTok{(a, }\BuiltInTok{round}\NormalTok{(species\_area\_relationship(a, c}\OperatorTok{=}\FloatTok{2.0}\NormalTok{, z}\OperatorTok{=}\FloatTok{0.25}\NormalTok{), }\DecValTok{2}\NormalTok{))}
\end{Highlighting}
\end{Shaded}

\begin{quote}
\textbf{Clinical / systems note:} Habitat fragmentation in zoonotic
emergence maps onto the same area-dependent extinction risk curves used
in terrestrial conservation.
\end{quote}

\begin{center}\rule{0.5\linewidth}{0.5pt}\end{center}

\subsection{Summary}\label{summary-37}

\begin{itemize}
\tightlist
\item
  \textbf{Nine terrestrial biomes:} tropical rainforest highest NPP
  (about 1,750 g C/m\(^2\)/yr); desert lowest
  (about 45-125); boreal forest largest land biome by area.
  Five aquatic biome types including coral reefs (most biodiverse marine
  ecosystem, threatened by bleaching).
\item
  \textbf{Climate change:} biome boundaries shifting 5-10 km/decade
  poleward; Amazon at about 17\% deforestation approaching
  20-25\% tipping point; permafrost thaw adds 2-5 Gt C/yr; novel biomes
  emerging on about 20\% of Earth's surface by 2100.
\item
  \textbf{Sixth Mass Extinction:} 100-1,000x background extinction rate;
  73\% vertebrate population decline since 1970; 1 million species
  threatened.
\item
  \textbf{IPBES findings:} about 1 million species threatened
  (12.5\%); 85\% of wetlands lost since 1700; 50\% coral reef loss since
  1870. Drivers ranked: land-use change \textgreater{} exploitation
  \textgreater{} climate change (rising fast) \textgreater{} pollution
  \textgreater{} invasives. Transformative Change Assessment (2024)
  argues incremental conservation insufficient.
\item
  \textbf{Climate velocity} \(v = (\partial T/\partial t)/|\nabla T|\):
  fast in flat tropical and temperate plains (6--10 km/yr), slow in
  mountains (6 m/yr); refugia (topographic, hydrological, microclimatic,
  elevational) are the key targets for climate-adaptive reserve design.
\item
  \textbf{Ecosystem-based adaptation (EbA):} mangroves, coral reefs,
  urban green infrastructure, wetlands, agroforestry. Often outperforms
  hard infrastructure on cost (\$94/m vs.~\$2,000/m for coastal
  protection) and provides co-benefits, but cannot substitute for
  emissions reduction.
\item
  \textbf{IUCN Red List:} CR = at least 80\% population decline under
  Criterion A; version 2025-1 listed 47,187 threatened species among
  169,420 assessed, and totals change as assessments expand
  \citep{iucn2025redlist}.
\item
  \textbf{MVP:} 50/500 rule (Franklin); revised 100/1000
  \citep{frankham2014}; PVA accounts for most stochasticity types.
  \(N_e/N\) ratio typically 0.10-0.25.
\item
  \textbf{Extinction debt:} fragmented species committed to extinction
  with 50-500 yr lags; current surveys overestimate long-term viability.
\item
  \textbf{Reserve design:} core-corridor-matrix; edge effects penetrate
  100-300 m; corridors increase movement 50\%; 30x30 target.
\item
  \textbf{Emerging strategies:} rewilding
  (\hyperref[gl:trophic-cascade]{\textbf{trophic cascade}}s),
  conservation genomics (ROH, adaptive variation), assisted gene flow
  (coral heat tolerance), de-extinction (mammoth project),
  climate-adaptive management.
\item
  \textbf{International frameworks:} CBD, CITES, 30x30 target; \$700
  billion/yr funding gap; quality and representativeness challenges.
\item
  \textbf{Connections:} See \cref{sec:unit_VI_phylogenetics} for
  phylogenetic diversity, \cref{sec:unit_X_community_ecology} for
  trophic rewilding, and \cref{sec:unit_V_population_genetics} for
  small-\hyperref[gl:population-genetics]{\textbf{population genetics}}.
\end{itemize}

\begin{center}\rule{0.5\linewidth}{0.5pt}\end{center}

\begin{center}\rule{0.5\linewidth}{0.5pt}\end{center}

\subsubsection{Companion Source
Module}\label{companion-source-module-37}

\textbf{Biomes and Conservation Biology} should leave a reproducible
trail from a biological claim to the code, figure, diagram, or
paper-based activity that can test it. Use the surfaces below to inspect
the chapter's assumptions, rerun the relevant model, or compare the
manuscript explanation with companion labs and figures.

{\def\LTcaptype{none} % do not increment counter
\begin{longtable}[]{@{}
  >{\raggedright\arraybackslash}p{(\linewidth - 2\tabcolsep) * \real{0.5000}}
  >{\raggedright\arraybackslash}p{(\linewidth - 2\tabcolsep) * \real{0.5000}}@{}}
\toprule\noalign{}
\begin{minipage}[b]{\linewidth}\raggedright
Surface
\end{minipage} & \begin{minipage}[b]{\linewidth}\raggedright
Use it for
\end{minipage} \\
\midrule\noalign{}
\endhead
\bottomrule\noalign{}
\endlastfoot
\breaktt{src/biology/ecology/ecology.py}
(\breaktt{species\_area\_relationship}, \breaktt{biodiversity\_indices},
\texttt{connectance}) & Reproduce conservation metrics and
tradeoff-sensitive summaries. \\
\breaktt{src/visualization/plots.py}
(\breaktt{plot\_species\_area\_relationship},
\breaktt{plot\_biome\_distribution}) & Inspect species-area assumptions
and biome comparisons. \\
\breaktt{src/mermaid/biology\_diagrams.py} (\texttt{food\_web\_diagram},
\breaktt{nutrient\_cycle\_diagram}) & Connect conservation action to
ecological pathways. \\
\end{longtable}
}

\textbf{Reproducibility check:} separate ecological evidence, social
objective, feasibility, uncertainty, and monitoring indicator before
choosing a conservation action. \textbf{Cross-reference:} use
\cref{sec:unit_X_population_ecology},
\cref{sec:unit_X_community_ecology}, and
\cref{sec:unit_X_ecosystem_ecology}.

\newpage

\section*{Lab A --- Systems Science and the Logic of
Emergence}\label{lab-a-systems-science-and-the-logic-of-emergence}
\addcontentsline{toc}{section}{Lab A --- Systems Science and the Logic
of Emergence}

\label{sec:lab_unit_0_systems_science}

\emph{This activity accompanies \cref{sec:unit_0_systems_science} of the
textbook --- review that chapter before attempting the exercises below.}

\textbf{Unit 0 \texorpdfstring{\ensuremath{\cdot}}{.} Lab A}

\begin{center}\rule{0.5\linewidth}{0.5pt}\end{center}

\subsection*{Learning Objectives}\label{learning-objectives-38}
\addcontentsline{toc}{subsection}{Learning Objectives}

\begin{itemize}
\tightlist
\item
  Identify the three canonical elements of a system (components,
  interactions, boundary) in concrete biological examples.
\item
  Diagram negative- and positive-feedback loops for real physiological
  systems and predict their stabilising or amplifying behaviour.
\item
  Compute response curves from the Hill equation and relate the Hill
  coefficient \(n\) to cooperativity and switch-like behaviour.
\item
  Reason about the difference between a complicated machine and a
  complex system using a worked example.
\end{itemize}

\begin{center}\rule{0.5\linewidth}{0.5pt}\end{center}

\subsubsection*{Alignment and Rubric
Map}\label{alignment-and-rubric-map}
\addcontentsline{toc}{subsubsection}{Alignment and Rubric Map}

\begin{itemize}
\tightlist
\item
  \textbf{Outcome 1 (LO1):} Interpret the supplied evidence or model
  output for Systems Science and the Logic of Emergence.
\item
  \textbf{Outcome 2 (LO2):} Identify controls and comparison groups that
  make the claim testable.
\item
  \textbf{Outcome 3 (LO3):} Quantify uncertainty, boundary conditions,
  or alternative explanations before concluding.
\item
  \textbf{Outcome 4 (LO4):} Transfer the mechanism to a new biological
  case or public-facing decision.
\item
  \textbf{Chapter LO coverage:} LO1, LO2, LO3, LO4
\item
  \textbf{Rubric dimensions:} evidence; controls; uncertainty;
  mechanism; transfer. \#\# Background \{.unnumbered\}
\end{itemize}

A \textbf{system} is a set of interacting components, bounded against
its surroundings, whose behaviour depends on the interactions as much as
on the parts. The central mathematical tool for stability is the
\textbf{Hill equation}:

\begin{equation}
\theta \;=\; \frac{[L]^n}{K_d^n + [L]^n}
\label{eq:lab_a_hill}
\end{equation}

where \texorpdfstring{\ensuremath{\theta}}{theta} is the fraction of receptors bound, \([L]\) the ligand
concentration, \(K_d\) the dissociation constant, and \(n\) the Hill
coefficient. For \(n = 1\) the curve is hyperbolic (graded response);
for \(n = 4\) it is steep and switch-like (as in haemoglobin's O\(_2\)
binding).

Today's lab trains three intuitions: what counts as ``the system'', how
feedback loops behave, and how cooperativity converts graded chemistry
into binary biology.

\begin{center}\rule{0.5\linewidth}{0.5pt}\end{center}

\subsection*{Pre-Lab (to complete before
class)}\label{pre-lab-to-complete-before-class}
\addcontentsline{toc}{subsection}{Pre-Lab (to complete before class)}

\begin{enumerate}
\def\labelenumi{\arabic{enumi}.}
\tightlist
\item
  Copy \cref{eq:lab_a_hill} into your notebook. Compute \texorpdfstring{\ensuremath{\theta}}{theta} for
  \(n = 1, 2, 4\) at \([L] = K_d\), \(2K_d\), \(5K_d\). Tabulate.
\item
  Define in your own words: \textbf{system}, \textbf{emergence},
  \textbf{feedback}, \textbf{homeostasis}, \textbf{modularity}.
\item
  Bring one example of a biological positive feedback loop and one
  example of a biological negative feedback loop (not from the
  textbook).
\end{enumerate}

\begin{center}\rule{0.5\linewidth}{0.5pt}\end{center}

\subsection*{Paper-Based Materials}\label{paper-based-materials}
\addcontentsline{toc}{subsection}{Paper-Based Materials}

\begin{itemize}
\tightlist
\item
  Graph paper (2 sheets per student).
\item
  Calculator.
\item
  Coloured pens.
\item
  Worksheet with blank feedback diagrams and Hill-curve axes.
\end{itemize}

\begin{center}\rule{0.5\linewidth}{0.5pt}\end{center}

\subsection*{Paper-Based Investigation}\label{paper-based-investigation}
\addcontentsline{toc}{subsection}{Paper-Based Investigation}

\subsubsection*{Part 1 --- Identifying
Systems}\label{part-1-identifying-systems}
\addcontentsline{toc}{subsubsection}{Part 1 --- Identifying Systems}

For each of the following, identify: (a) the components, (b) the key
interactions, (c) the boundary separating system from environment, and
(d) one input and one output crossing the boundary.

\begin{enumerate}
\def\labelenumi{\arabic{enumi}.}
\tightlist
\item
  A single mitochondrion.
\item
  The glucose-insulin-glucagon regulatory axis.
\item
  A coral reef.
\item
  A bacterial biofilm on a medical catheter.
\item
  A neuron's axon hillock.
\end{enumerate}

\subsubsection*{Part 2 --- Feedback Loop
Diagrams}\label{part-2-feedback-loop-diagrams}
\addcontentsline{toc}{subsubsection}{Part 2 --- Feedback Loop Diagrams}

For each system below, diagram the feedback loop and classify it as
positive or negative. Predict the steady-state behaviour (stable fixed
point, oscillation, or runaway).

\begin{enumerate}
\def\labelenumi{\arabic{enumi}.}
\tightlist
\item
  Blood pressure \texorpdfstring{\ensuremath{\to}}{->} baroreceptor firing \texorpdfstring{\ensuremath{\to}}{->} vagal tone \texorpdfstring{\ensuremath{\to}}{->} heart rate \texorpdfstring{\ensuremath{\to}}{->} blood
  pressure.
\item
  Oxytocin \texorpdfstring{\ensuremath{\to}}{->} uterine contraction \texorpdfstring{\ensuremath{\to}}{->} stretch receptor \texorpdfstring{\ensuremath{\to}}{->} oxytocin (during
  labour).
\item
  Luteinising hormone \texorpdfstring{\ensuremath{\to}}{->} ovarian oestradiol \texorpdfstring{\ensuremath{\to}}{->} hypothalamic GnRH
  (follicular phase).
\item
  Luteinising hormone \texorpdfstring{\ensuremath{\to}}{->} ovarian oestradiol \texorpdfstring{\ensuremath{\to}}{->} hypothalamic GnRH
  (ovulatory phase).
\item
  Depolarising Na\(^+\) current \texorpdfstring{\ensuremath{\to}}{->} more open Na\(^+\) channels \texorpdfstring{\ensuremath{\to}}{->}
  depolarisation (action potential upstroke).
\end{enumerate}

Why do loops 3 and 4 have different signs, and what does the switch
imply for the dynamics of the menstrual cycle?

\subsubsection*{Part 3 --- Hill Curve by
Hand}\label{part-3-hill-curve-by-hand}
\addcontentsline{toc}{subsubsection}{Part 3 --- Hill Curve by Hand}

\begin{enumerate}
\def\labelenumi{\arabic{enumi}.}
\tightlist
\item
  Using \(K_d = 10\,\mu\mathrm{M}\), compute \texorpdfstring{\ensuremath{\theta}}{theta} at
  \([L] = 1, 2, 5, 10, 20, 50, 100\,\mu\mathrm{M}\) for three values of
  \(n\): 1, 2.8 (physiological Hb), and 4.
\item
  Plot most three curves on the same axes. Label the ``switch zone''
  where small changes in \([L]\) produce large changes in \texorpdfstring{\ensuremath{\theta}}{theta}.
\item
  For which \(n\) is the transition sharpest? How does this
  cooperativity benefit an oxygen-carrying pigment?
\end{enumerate}

\subsubsection*{Part 4 --- Complicated
vs.~Complex}\label{part-4-complicated-vs.-complex}
\addcontentsline{toc}{subsubsection}{Part 4 --- Complicated vs.~Complex}

A wristwatch and a live cell both contain many interacting components.
Yet one is \textbf{complicated} and the other \textbf{complex}. In a
table, contrast:

{\def\LTcaptype{none} % do not increment counter
\begin{longtable}[]{@{}
  >{\raggedright\arraybackslash}p{(\linewidth - 4\tabcolsep) * \real{0.2692}}
  >{\raggedright\arraybackslash}p{(\linewidth - 4\tabcolsep) * \real{0.3846}}
  >{\raggedright\arraybackslash}p{(\linewidth - 4\tabcolsep) * \real{0.3462}}@{}}
\toprule\noalign{}
\begin{minipage}[b]{\linewidth}\raggedright
Feature
\end{minipage} & \begin{minipage}[b]{\linewidth}\raggedright
Wristwatch
\end{minipage} & \begin{minipage}[b]{\linewidth}\raggedright
Live cell
\end{minipage} \\
\midrule\noalign{}
\endhead
\bottomrule\noalign{}
\endlastfoot
Component count & about 100 & about 10\texorpdfstring{\textsuperscript{9}}{9} molecules \\
Predictability & High (mechanistic) & Statistical \\
Repairability & External & Self-repair \\
Adaptation & None & Evolutionary \\
Failure mode & Component breakage & Network dysregulation \\
Emergent properties & None (function designed in) & Many (metabolism,
identity, division) \\
\end{longtable}
}

Discuss in your group one additional feature that distinguishes complex
from complicated, and one case where the distinction is blurry (e.g., a
modern autopilot with adaptive learning).

\begin{center}\rule{0.5\linewidth}{0.5pt}\end{center}

\subsection*{Data Tables}\label{data-tables}
\addcontentsline{toc}{subsection}{Data Tables}

{\def\LTcaptype{none} % do not increment counter
\begin{longtable}[]{@{}llllll@{}}
\toprule\noalign{}
System & Components & Interactions & Boundary & Input & Output \\
\midrule\noalign{}
\endhead
\bottomrule\noalign{}
\endlastfoot
Mitochondrion & & & & & \\
Glucose axis & & & & & \\
Coral reef & & & & & \\
Biofilm & & & & & \\
Axon hillock & & & & & \\
\end{longtable}
}

{\def\LTcaptype{none} % do not increment counter
\begin{longtable}[]{@{}lll@{}}
\toprule\noalign{}
Feedback loop & +/\texorpdfstring{\ensuremath{-}}{-} & Predicted behaviour \\
\midrule\noalign{}
\endhead
\bottomrule\noalign{}
\endlastfoot
Baroreceptor reflex & & \\
Oxytocin-uterus & & \\
GnRH follicular & & \\
GnRH ovulatory & & \\
Na\(^+\) action potential & & \\
\end{longtable}
}

{\def\LTcaptype{none} % do not increment counter
\begin{longtable}[]{@{}
  >{\raggedright\arraybackslash}p{(\linewidth - 6\tabcolsep) * \real{0.2933}}
  >{\raggedright\arraybackslash}p{(\linewidth - 6\tabcolsep) * \real{0.2267}}
  >{\raggedright\arraybackslash}p{(\linewidth - 6\tabcolsep) * \real{0.2533}}
  >{\raggedright\arraybackslash}p{(\linewidth - 6\tabcolsep) * \real{0.2267}}@{}}
\toprule\noalign{}
\begin{minipage}[b]{\linewidth}\raggedright
\([L]\,(\mu\mathrm{M})\)
\end{minipage} & \begin{minipage}[b]{\linewidth}\raggedright
\texorpdfstring{\ensuremath{\theta}}{theta} at \(n=1\)
\end{minipage} & \begin{minipage}[b]{\linewidth}\raggedright
\texorpdfstring{\ensuremath{\theta}}{theta} at \(n=2.8\)
\end{minipage} & \begin{minipage}[b]{\linewidth}\raggedright
\texorpdfstring{\ensuremath{\theta}}{theta} at \(n=4\)
\end{minipage} \\
\midrule\noalign{}
\endhead
\bottomrule\noalign{}
\endlastfoot
1 & & & \\
2 & & & \\
5 & & & \\
10 & & & \\
20 & & & \\
50 & & & \\
100 & & & \\
\end{longtable}
}

\begin{center}\rule{0.5\linewidth}{0.5pt}\end{center}

\subsection*{Evidence and Reproducibility
Checklist}\label{evidence-and-reproducibility-checklist}
\addcontentsline{toc}{subsection}{Evidence and Reproducibility
Checklist}

\begin{itemize}
\tightlist
\item
  \textbf{Primary evidence goal:} Build and critique a feedback model
  for a familiar living system.
\item
  \textbf{Data skill to practice:} Translate a verbal biological system
  into variables, links, and testable predictions.
\item
  \textbf{BioSkills emphasis:} Modeling and simulation, Process of
  science, Quantitative reasoning.
\item
  \textbf{Control logic:} identify at least one positive control, one
  negative control, or one baseline comparison before interpreting
  results.
\item
  \textbf{Measurement discipline:} record units, uncertainty, sample
  size, and any discarded observation with a reason.
\item
  \textbf{Mechanistic link:} connect one result directly to the parent
  chapter's big idea before writing the conclusion.
\item
  \textbf{Reproducibility check:} state one procedural detail that
  another group would need in order to reproduce the result.
\end{itemize}

\subsection*{Paper-Based Evidence
Upgrade}\label{paper-based-evidence-upgrade}
\addcontentsline{toc}{subsection}{Paper-Based Evidence Upgrade}

Before answering the analysis questions, annotate the paper dataset for
\textbf{Systems Science and the Logic of Emergence} with a
reproducibility pass:

{\def\LTcaptype{none} % do not increment counter
\begin{longtable}[]{@{}
  >{\raggedright\arraybackslash}p{(\linewidth - 2\tabcolsep) * \real{0.5000}}
  >{\raggedright\arraybackslash}p{(\linewidth - 2\tabcolsep) * \real{0.5000}}@{}}
\toprule\noalign{}
\begin{minipage}[b]{\linewidth}\raggedright
Evidence check
\end{minipage} & \begin{minipage}[b]{\linewidth}\raggedright
Student action
\end{minipage} \\
\midrule\noalign{}
\endhead
\bottomrule\noalign{}
\endlastfoot
Control logic & Mark the comparison that functions as the baseline,
negative control, or reference case. \\
Uncertainty & Circle the row, card, diagram feature, or model assumption
most likely to change the conclusion. \\
Model comparison & State whether a simpler rule, null model, or
alternative mechanism could explain the same pattern. \\
Decision threshold & Write the minimum evidence that would make you
revise the interpretation. \\
Reproducibility & Record the exact scoring rule another group would need
to reproduce your classification. \\
\end{longtable}
}

Focus note: For the feedback-diagram and Hill-curve datasets, the
reproducible quantities are the sign assigned to each loop edge and the
Hill coefficient \(n\) used for the saturation calculation: state the
system boundary you drew, justify every \(+\) or \(-\) sign from the
dataset, and report the O\(_2\)-saturation numbers another group must be
able to regenerate from your stated \(n\). Keep required work
paper-based; any material-handling or equipment version belongs only in
an optional extension.

\subsection*{Analysis Questions}\label{analysis-questions}
\addcontentsline{toc}{subsection}{Analysis Questions}

\begin{enumerate}
\def\labelenumi{\arabic{enumi}.}
\tightlist
\item
  For haemoglobin (\(n \approx 2.8\)), compute the difference in O\(_2\)
  saturation between lung capillaries
  (\([\mathrm{O}_2] = 100\,\mathrm{mmHg}\)) and active muscle
  capillaries (\([\mathrm{O}_2] = 20\,\mathrm{mmHg}\)). Why is the
  steepness of the Hill curve a life-preserving property?
\item
  In Part 2, the ovarian-pituitary loop reverses sign across the
  menstrual cycle. What does this imply for the stability of the overall
  system, and what biological mechanism produces the sign reversal?
\item
  Modular design makes biological systems easier to evolve (new modules
  can be added without breaking old ones). Give one counter-example ---
  a domain where modularity \emph{constrains} evolution rather than
  freeing it.
\item
  The textbook notes that biological delays destabilise feedback. A
  surgical patient on a mechanical ventilator has a CO\(_2\)-sensing
  feedback loop with a built-in 10-second measurement lag. Predict what
  happens as you increase the loop gain; suggest one design fix.
\item
  Design a paper-based experiment to test whether a bacterial-population
  model exhibits emergent collective behaviour distinct from
  isolated-cell behaviour. What is the null hypothesis and what
  observation card would refute it?
\end{enumerate}

\begin{center}\rule{0.5\linewidth}{0.5pt}\end{center}

\subsection*{Discussion Questions}\label{discussion-questions-5}
\addcontentsline{toc}{subsection}{Discussion Questions}

\begin{enumerate}
\def\labelenumi{\arabic{enumi}.}
\tightlist
\item
  A reductionist insists that knowing everything about every molecule in
  a cell is sufficient to understand cellular behaviour. Where does
  systems science push back?
\item
  Are there systems in biology that are \emph{not} adaptive? Propose a
  candidate and defend your choice.
\item
  The Covid-19 pandemic exposed how hospital systems behave under
  extreme input. Which systems-science principles explain why ICU
  capacity failed in some cities and not others?
\end{enumerate}

\begin{center}\rule{0.5\linewidth}{0.5pt}\end{center}

\subsection*{Safety and Ethics Notes}\label{safety-and-ethics-notes}
\addcontentsline{toc}{subsection}{Safety and Ethics Notes}

Paper-based lab --- no reagents or instruments, no physical risk.
Discussion of the menstrual cycle, labour, and ventilator management
should be inclusive and accurate; seek clarification from the instructor
if terminology feels unfamiliar. The Covid-19 example should be grounded
in public-health data rather than anecdote.

\emph{Module: \breaktt{src/biology/cell/cell\_biology.py}
(\texttt{hill\_equation}, \breaktt{receptor\_occupancy});
\breaktt{src/biology/ecology/ecology.py} (\texttt{logistic\_growth}).}

\newpage

\section*{Lab B --- Complex Adaptive
Systems}\label{lab-b-complex-adaptive-systems}
\addcontentsline{toc}{section}{Lab B --- Complex Adaptive Systems}

\label{sec:lab_unit_0_complex_adaptive_systems}

\emph{This activity accompanies
\cref{sec:unit_0_complex_adaptive_systems} of the textbook --- review
that chapter before attempting the exercises below.}

\textbf{Unit 0 \texorpdfstring{\ensuremath{\cdot}}{.} Lab B}

\begin{center}\rule{0.5\linewidth}{0.5pt}\end{center}

\subsection*{Learning Objectives}\label{learning-objectives-39}
\addcontentsline{toc}{subsection}{Learning Objectives}

\begin{itemize}
\tightlist
\item
  Simulate a simple agent-based system on paper and observe emergent
  global behaviour.
\item
  Compute the period of a logistic map for selected reproductive rates
  and identify the onset of chaos.
\item
  Sketch a phase diagram for a two-basin attractor landscape and reason
  about tipping points.
\item
  Distinguish \emph{robustness} from \emph{redundancy} using a worked
  biological example.
\end{itemize}

\begin{center}\rule{0.5\linewidth}{0.5pt}\end{center}

\subsubsection*{Alignment and Rubric
Map}\label{alignment-and-rubric-map-1}
\addcontentsline{toc}{subsubsection}{Alignment and Rubric Map}

\begin{itemize}
\tightlist
\item
  \textbf{Outcome 1 (LO1):} Interpret the supplied evidence or model
  output for Complex Adaptive Systems.
\item
  \textbf{Outcome 2 (LO2):} Identify controls and comparison groups that
  make the claim testable.
\item
  \textbf{Outcome 3 (LO3):} Quantify uncertainty, boundary conditions,
  or alternative explanations before concluding.
\item
  \textbf{Outcome 4 (LO4):} Transfer the mechanism to a new biological
  case or public-facing decision.
\item
  \textbf{Chapter LO coverage:} LO1, LO2, LO3, LO4
\item
  \textbf{Rubric dimensions:} evidence; controls; uncertainty;
  mechanism; transfer. \#\# Background \{.unnumbered\}
\end{itemize}

A \textbf{complex adaptive system (CAS)} is a population of agents whose
collective behaviour emerges from local rules rather than central
control. Examples run from neurons in a cortical circuit to ants in a
colony to species in an ecosystem. Two mathematical tools recur across
scales:

\begin{enumerate}
\def\labelenumi{\arabic{enumi}.}
\tightlist
\item
  \textbf{The logistic map} --- a discrete-time model
  \(x_{n+1} = r x_n (1 - x_n)\) that passes from stable fixed point \texorpdfstring{\ensuremath{\to}}{->}
  period-2 \texorpdfstring{\ensuremath{\to}}{->} period-4 \texorpdfstring{\ensuremath{\to}}{->} chaos as the parameter \(r\) increases. It
  captures the essence of bifurcation theory in a single line of
  algebra.
\item
  \textbf{Phase-space attractor landscapes} --- a potential-well picture
  that shows how systems can occupy alternative stable states with a
  ``tipping point'' between them, and how hysteresis delays recovery.
\end{enumerate}

This lab is paper-based; a calculator and graph paper are the only tools
needed.

\begin{center}\rule{0.5\linewidth}{0.5pt}\end{center}

\subsection*{Pre-Lab (to complete before
class)}\label{pre-lab-to-complete-before-class-1}
\addcontentsline{toc}{subsection}{Pre-Lab (to complete before class)}

\begin{enumerate}
\def\labelenumi{\arabic{enumi}.}
\tightlist
\item
  Write down the logistic map and draw its form for \(0 \le x \le 1\),
  \(r = 2.5\). Mark the intersection with the identity line
  \(x_{n+1} = x_n\) --- the fixed point.
\item
  Briefly define: \textbf{attractor}, \textbf{bifurcation},
  \textbf{hysteresis}, \textbf{self-organised criticality}.
\item
  Bring three examples of bistable biological systems to class (for Part
  3 discussion).
\end{enumerate}

\begin{center}\rule{0.5\linewidth}{0.5pt}\end{center}

\subsection*{Paper-Based Materials}\label{paper-based-materials-1}
\addcontentsline{toc}{subsection}{Paper-Based Materials}

\begin{itemize}
\tightlist
\item
  Graph paper (2 sheets per student).
\item
  Calculator.
\item
  Coloured pens.
\item
  Worksheet with a blank logistic-map table and a blank phase landscape.
\end{itemize}

\begin{center}\rule{0.5\linewidth}{0.5pt}\end{center}

\subsection*{Paper-Based
Investigation}\label{paper-based-investigation-1}
\addcontentsline{toc}{subsection}{Paper-Based Investigation}

\subsubsection*{Part 1 --- The Logistic Map by
Hand}\label{part-1-the-logistic-map-by-hand}
\addcontentsline{toc}{subsubsection}{Part 1 --- The Logistic Map by
Hand}

\begin{enumerate}
\def\labelenumi{\arabic{enumi}.}
\tightlist
\item
  Starting from \(x_0 = 0.2\), compute 20 iterations of the logistic map
  for each of these \(r\) values: \(r = 2.5,\; 3.2,\; 3.5,\; 3.8\).
\item
  For each \(r\), plot \(x_n\) against iteration number \(n\) on the
  same axes (use a different colour for each \(r\)).
\item
  Identify the long-term behaviour: fixed point, period-2, period-4, or
  chaotic.
\item
  Bonus: compute \(x_{100}\) for \(r = 3.8\) with \(x_0 = 0.2\) and
  \(x_0 = 0.20001\). How different are they? Relate to ``sensitive
  dependence on initial conditions.''
\end{enumerate}

\subsubsection*{Part 2 --- Attractor Landscape and Tipping
Points}\label{part-2-attractor-landscape-and-tipping-points}
\addcontentsline{toc}{subsubsection}{Part 2 --- Attractor Landscape and
Tipping Points}

A shallow lake has two alternative stable states: \textbf{clear} (low
algae, clear water) and \textbf{turbid} (high algae, murky). Phosphorus
loading \(P\) shifts which state is stable.

\begin{enumerate}
\def\labelenumi{\arabic{enumi}.}
\tightlist
\item
  Sketch the potential landscape \(U(A)\) (where \(A\) = algal biomass)
  for three values of \(P\): low (primarily clear is stable),
  intermediate (both states stable), and high (primarily turbid is
  stable).
\item
  Mark the \textbf{tipping point} (\(P_1\)) where the clear state loses
  stability as \(P\) increases.
\item
  Mark the \textbf{restoration point} (\(P_2\)) where the turbid state
  loses stability as \(P\) decreases.
\item
  Compare \(P_1\) and \(P_2\) --- which is smaller? This gap is the
  hysteresis loop and explains why eutrophication is hard to reverse.
\end{enumerate}

\subsubsection*{Part 3 --- Robustness Versus
Redundancy}\label{part-3-robustness-versus-redundancy}
\addcontentsline{toc}{subsubsection}{Part 3 --- Robustness Versus
Redundancy}

Consider the mammalian immune system as a CAS. It has five categories of
robustness:

\begin{itemize}
\tightlist
\item
  \textbf{Redundancy} --- multiple B-cell clones specific to the same
  antigen.
\item
  \textbf{Degeneracy} --- T- and B-cells with overlapping but distinct
  function.
\item
  \textbf{Modularity} --- innate vs.~adaptive compartments.
\item
  \textbf{Negative feedback} --- regulatory T cells suppress overactive
  responders.
\item
  \textbf{Diversity} --- V(D)J recombination generates
  about \(10^{11}\) receptor sequences.
\end{itemize}

\begin{enumerate}
\def\labelenumi{\arabic{enumi}.}
\tightlist
\item
  For each category, give one concrete example from the immune system.
\item
  Then list one \emph{fragility} that the same architectural feature
  introduces. (E.g., immune redundancy also means multiple places
  autoimmunity can arise.)
\item
  Discuss in your group: would a maximally robust system also be
  maximally evolvable?
\end{enumerate}

\begin{center}\rule{0.5\linewidth}{0.5pt}\end{center}

\subsection*{Data Tables}\label{data-tables-1}
\addcontentsline{toc}{subsection}{Data Tables}

{\def\LTcaptype{none} % do not increment counter
\begin{longtable}[]{@{}
  >{\raggedright\arraybackslash}p{(\linewidth - 8\tabcolsep) * \real{0.1688}}
  >{\raggedright\arraybackslash}p{(\linewidth - 8\tabcolsep) * \real{0.2078}}
  >{\raggedright\arraybackslash}p{(\linewidth - 8\tabcolsep) * \real{0.2078}}
  >{\raggedright\arraybackslash}p{(\linewidth - 8\tabcolsep) * \real{0.2078}}
  >{\raggedright\arraybackslash}p{(\linewidth - 8\tabcolsep) * \real{0.2078}}@{}}
\toprule\noalign{}
\begin{minipage}[b]{\linewidth}\raggedright
Iteration \(n\)
\end{minipage} & \begin{minipage}[b]{\linewidth}\raggedright
\(x_n\) at \(r=2.5\)
\end{minipage} & \begin{minipage}[b]{\linewidth}\raggedright
\(x_n\) at \(r=3.2\)
\end{minipage} & \begin{minipage}[b]{\linewidth}\raggedright
\(x_n\) at \(r=3.5\)
\end{minipage} & \begin{minipage}[b]{\linewidth}\raggedright
\(x_n\) at \(r=3.8\)
\end{minipage} \\
\midrule\noalign{}
\endhead
\bottomrule\noalign{}
\endlastfoot
0 & 0.2 & 0.2 & 0.2 & 0.2 \\
1 & & & & \\
\ldots{} & & & & \\
20 & & & & \\
\end{longtable}
}

Long-term regime at each \(r\): \(r=2.5\): \textbf{\emph{, \(r=3.2\):
}}, \(r=3.5\): \textbf{\emph{, \(r=3.8\): }}.

{\def\LTcaptype{none} % do not increment counter
\begin{longtable}[]{@{}lll@{}}
\toprule\noalign{}
CAS property & Immune-system example & Associated fragility \\
\midrule\noalign{}
\endhead
\bottomrule\noalign{}
\endlastfoot
Redundancy & & \\
Degeneracy & & \\
Modularity & & \\
Negative feedback & & \\
Diversity & & \\
\end{longtable}
}

\begin{center}\rule{0.5\linewidth}{0.5pt}\end{center}

\subsection*{Evidence and Reproducibility
Checklist}\label{evidence-and-reproducibility-checklist-1}
\addcontentsline{toc}{subsection}{Evidence and Reproducibility
Checklist}

\begin{itemize}
\tightlist
\item
  \textbf{Primary evidence goal:} Explore how small rule changes alter
  collective behaviour.
\item
  \textbf{Data skill to practice:} Distinguish deterministic trends from
  stochastic variation in repeated simulations.
\item
  \textbf{BioSkills emphasis:} Modeling and simulation, Process of
  science, Quantitative reasoning.
\item
  \textbf{Control logic:} identify at least one positive control, one
  negative control, or one baseline comparison before interpreting
  results.
\item
  \textbf{Measurement discipline:} record units, uncertainty, sample
  size, and any discarded observation with a reason.
\item
  \textbf{Mechanistic link:} connect one result directly to the parent
  chapter's big idea before writing the conclusion.
\item
  \textbf{Reproducibility check:} state one procedural detail that
  another group would need in order to reproduce the result.
\end{itemize}

\subsection*{Paper-Based Evidence
Upgrade}\label{paper-based-evidence-upgrade-1}
\addcontentsline{toc}{subsection}{Paper-Based Evidence Upgrade}

Before answering the analysis questions, annotate the paper dataset for
\textbf{Complex Adaptive Systems} with a reproducibility pass:

{\def\LTcaptype{none} % do not increment counter
\begin{longtable}[]{@{}
  >{\raggedright\arraybackslash}p{(\linewidth - 2\tabcolsep) * \real{0.5000}}
  >{\raggedright\arraybackslash}p{(\linewidth - 2\tabcolsep) * \real{0.5000}}@{}}
\toprule\noalign{}
\begin{minipage}[b]{\linewidth}\raggedright
Evidence check
\end{minipage} & \begin{minipage}[b]{\linewidth}\raggedright
Student action
\end{minipage} \\
\midrule\noalign{}
\endhead
\bottomrule\noalign{}
\endlastfoot
Control logic & Mark the comparison that functions as the baseline,
negative control, or reference case. \\
Uncertainty & Circle the row, card, diagram feature, or model assumption
most likely to change the conclusion. \\
Model comparison & State whether a simpler rule, null model, or
alternative mechanism could explain the same pattern. \\
Decision threshold & Write the minimum evidence that would make you
revise the interpretation. \\
Reproducibility & Record the exact scoring rule another group would need
to reproduce your classification. \\
\end{longtable}
}

Focus note: Here the reproducibility pass turns on the logistic-map
iteration and the hysteresis-loop dataset: record the exact value of
\(r\) and the seed/initial condition for every trajectory, and for the
tipping-point loop note which branch (clear-water versus turbid) you
read and the phosphorus level at which the jump occurred, since a
different starting point silently changes the conclusion. Keep required
work paper-based; any material-handling or equipment version belongs
only in an optional extension.

\subsection*{Analysis Questions}\label{analysis-questions-1}
\addcontentsline{toc}{subsection}{Analysis Questions}

\begin{enumerate}
\def\labelenumi{\arabic{enumi}.}
\tightlist
\item
  The logistic map transitions from periodic to chaotic behaviour as
  \(r\) passes through the Feigenbaum point (about 3.57). Why
  does this transition occur even though the underlying equation is
  completely deterministic?
\item
  In Part 2, you drew a hysteresis loop. Why can reducing \(P\) to the
  level \emph{at which} eutrophication occurred fail to restore the
  clear-water state? What mechanism sustains the turbid attractor?
\item
  Power-law distributions appear in earthquake magnitudes, neural
  avalanches, and extinction events. What common CAS mechanism could
  explain most three?
\item
  Cancer has been described as a ``failed CAS'' where normal tissue
  homeostasis breaks down. Identify two CAS features that go wrong in
  malignancy.
\item
  Propose a paper simulation or card-based experiment using a
  bacterial-biofilm model to demonstrate one CAS principle (e.g.,
  stigmergic self-organisation, threshold transitions, or robustness to
  perturbation). Describe the set-up, observable, and expected result.
\end{enumerate}

\begin{center}\rule{0.5\linewidth}{0.5pt}\end{center}

\subsection*{Discussion Questions}\label{discussion-questions-6}
\addcontentsline{toc}{subsection}{Discussion Questions}

\begin{enumerate}
\def\labelenumi{\arabic{enumi}.}
\tightlist
\item
  The SARS-CoV-2 pandemic generated striking geographic and temporal
  heterogeneity. Which features of a CAS (noise, nonlinearity, network
  topology) explain the emergence of variants?
\item
  Is a single cell a CAS? Defend or refute using at least three of the
  Unit 0.2 criteria.
\item
  What additional scientific ingredients would be needed to ``build life
  from scratch'' in a synthetic CAS? What does the field of synthetic
  biology currently get right, and what do its failures reveal about the
  difficulty of engineering life?
\end{enumerate}

\begin{center}\rule{0.5\linewidth}{0.5pt}\end{center}

\subsection*{Safety and Ethics Notes}\label{safety-and-ethics-notes-1}
\addcontentsline{toc}{subsection}{Safety and Ethics Notes}

Paper-based lab --- no reagents, no risk. Discussion of SARS-CoV-2
should be current-affairs aware; respect divergent lived experiences of
the pandemic. Discussion of cancer as a ``failed CAS'' should be framed
carefully if anyone in the class has personal experience --- use the
system-level language, not the blame-oriented ``cancer cells are
selfish'' framing.

\emph{Module: \breaktt{src/biology/ecology/ecology.py}
(\texttt{logistic\_growth}, \texttt{lotka\_volterra});
\breaktt{src/biology/evolution/evolution.py} (\texttt{simulate\_drift}).}

\newpage

\section*{Lab C --- Active Inference and the Free Energy
Principle}\label{lab-c-active-inference-and-the-free-energy-principle}
\addcontentsline{toc}{section}{Lab C --- Active Inference and the Free
Energy Principle}

\label{sec:lab_unit_0_active_inference}

\emph{This activity accompanies \cref{sec:unit_0_active_inference} of
the textbook --- review that chapter before attempting the exercises
below.}

\textbf{Unit 0 \texorpdfstring{\ensuremath{\cdot}}{.} Lab C}

\begin{center}\rule{0.5\linewidth}{0.5pt}\end{center}

\subsection*{Learning Objectives}\label{learning-objectives-40}
\addcontentsline{toc}{subsection}{Learning Objectives}

\begin{itemize}
\tightlist
\item
  Apply core principles of active inference to concrete biological
  observations.
\item
  Compute a Bayesian posterior by hand for a one-dimensional inference
  problem and interpret the result in terms of prior, likelihood, and
  evidence.
\item
  Distinguish \emph{perceptual} inference (updating beliefs) from
  \emph{active} inference (acting to satisfy beliefs) using everyday
  examples.
\item
  Quantify how sensor precision changes the trade-off between updating a
  belief and acting to make a belief true.
\end{itemize}

\begin{center}\rule{0.5\linewidth}{0.5pt}\end{center}

\subsubsection*{Alignment and Rubric
Map}\label{alignment-and-rubric-map-2}
\addcontentsline{toc}{subsubsection}{Alignment and Rubric Map}

\begin{itemize}
\tightlist
\item
  \textbf{Outcome 1 (LO1):} Interpret the supplied evidence or model
  output for Active Inference and the Free Energy Principle.
\item
  \textbf{Outcome 2 (LO2):} Identify controls and comparison groups that
  make the claim testable.
\item
  \textbf{Outcome 3 (LO3):} Quantify uncertainty, boundary conditions,
  or alternative explanations before concluding.
\item
  \textbf{Outcome 4 (LO4):} Transfer the mechanism to a new biological
  case or public-facing decision.
\item
  \textbf{Chapter LO coverage:} LO1, LO2, LO3, LO4
\item
  \textbf{Rubric dimensions:} evidence; controls; uncertainty;
  mechanism; transfer. \#\# Background \{.unnumbered\}
\end{itemize}

Active inference proposes that living agents minimise a quantity called
\textbf{variational free energy} --- an upper bound on the ``surprise''
of observed sensory data given the agent's internal generative model.
The core equation (simplified):

\begin{equation}
F \;=\; \mathbb{E}_{Q}[\,\log Q(\mathbf{z}) - \log P(\mathbf{o}, \mathbf{z})\,]
\label{eq:lab_ai_fe}
\end{equation}

where \(Q(\mathbf{z})\) is the agent's current belief about hidden
causes \(\mathbf{z}\), and \(P(\mathbf{o}, \mathbf{z})\) is the agent's
model of how those causes generate observations \(\mathbf{o}\).
Minimising \(F\) is achieved either by \textbf{updating beliefs}
(perceptual inference) or \textbf{changing the world through action}
(active inference).

Today's lab is paper-based: you will work through two hand-computed
examples and one behaviour-card classification exercise. No computer is
required, though a calculator (or phone app) is useful for the Gaussian
arithmetic.

\begin{center}\rule{0.5\linewidth}{0.5pt}\end{center}

\subsection*{Pre-Lab (to complete before
class)}\label{pre-lab-to-complete-before-class-2}
\addcontentsline{toc}{subsection}{Pre-Lab (to complete before class)}

\begin{enumerate}
\def\labelenumi{\arabic{enumi}.}
\tightlist
\item
  Copy \cref{eq:lab_ai_fe} into your notebook and annotate each symbol.
\item
  For a Gaussian prior \(\mathcal{N}(\mu_0, \sigma_0^2)\) and Gaussian
  likelihood \(\mathcal{N}(o, \sigma_o^2)\), write the closed-form
  posterior mean and variance.
\item
  Choose one everyday example of \emph{perceptual inference} and one of
  \emph{active inference} that you could describe in two sentences each.
  You will compare them with the scenario cards during Part 3.
\end{enumerate}

\begin{center}\rule{0.5\linewidth}{0.5pt}\end{center}

\subsection*{Paper-Based Materials}\label{paper-based-materials-2}
\addcontentsline{toc}{subsection}{Paper-Based Materials}

\begin{itemize}
\tightlist
\item
  Printed worksheet with two Bayesian inference problems (provided).
\item
  Behaviour scenario-card set with ten animal and human examples.
\item
  Calculator or phone.
\item
  Coloured pens.
\end{itemize}

\begin{center}\rule{0.5\linewidth}{0.5pt}\end{center}

\subsection*{Paper-Based
Investigation}\label{paper-based-investigation-2}
\addcontentsline{toc}{subsection}{Paper-Based Investigation}

\subsubsection*{Part 1 --- Hand-computed Bayesian
Update}\label{part-1-hand-computed-bayesian-update}
\addcontentsline{toc}{subsubsection}{Part 1 --- Hand-computed Bayesian
Update}

A bird forages for seeds hidden under two snow patches. Its prior belief
about the location of food is a Gaussian centred at position
\(\mu_0 = 10\) m (where food was yesterday) with variance
\(\sigma_0^2 = 4\,\mathrm{m}^2\). A partial snow melt gives a noisy
visual cue that food is at position \(o = 16\) m with likelihood
variance \(\sigma_o^2 = 12\,\mathrm{m}^2\).

\begin{enumerate}
\def\labelenumi{\arabic{enumi}.}
\tightlist
\item
  Compute the \textbf{prediction error}: \(\varepsilon = o - \mu_0\).
\item
  Compute the \textbf{Kalman gain}:
  \(K = \sigma_0^2 / (\sigma_0^2 + \sigma_o^2)\).
\item
  Compute the \textbf{posterior mean}:
  \(\mu_1 = \mu_0 + K \, \varepsilon\).
\item
  Compute the \textbf{posterior variance}:
  \(\sigma_1^2 = (1 - K) \, \sigma_0^2\).
\item
  Re-compute \(\mu_1\) under the assumption that the cue is \emph{very}
  noisy (\(\sigma_o^2 = 100\)) and \emph{very} precise
  (\(\sigma_o^2 = 0.5\)). Explain in words how the bird's final belief
  depends on sensor precision.
\end{enumerate}

\subsubsection*{Part 2 --- Perceptual vs.~Active
Inference}\label{part-2-perceptual-vs.-active-inference}
\addcontentsline{toc}{subsubsection}{Part 2 --- Perceptual vs.~Active
Inference}

A person wakes up cold. Their generative model has a tight prior on body
temperature (\(\mu = 37\) °C, \(\sigma = 0.5\)). Their interoceptive
sensor reports 34.5 °C with \(\sigma = 1.0\).

\begin{enumerate}
\def\labelenumi{\arabic{enumi}.}
\tightlist
\item
  Compute the posterior belief about their body temperature (same
  formulae as Part 1).
\item
  The posterior free energy is reduced but still positive. List
  \textbf{three actions} the person could take that would drive free
  energy further toward zero by \emph{changing the observation} rather
  than the belief.
\item
  Argue in one paragraph why active inference is more effective than
  perceptual inference when the prior precision is tight (small
  \(\sigma_0\)) --- i.e., when the organism is ``committed'' to a
  particular internal state.
\end{enumerate}

\subsubsection*{Part 3 --- Behaviour Scenario
Classification}\label{part-3-behaviour-scenario-classification}
\addcontentsline{toc}{subsubsection}{Part 3 --- Behaviour Scenario
Classification}

Use the ten printed behaviour scenario cards. For each behaviour,
annotate whether it is perceptual (sampling information) or active
(imposing prediction on the environment). Record the evidence phrase
that made your classification reproducible. As an optional extension
outside class, you may compare the cards with a respectful observation
of an animal or consenting human, but the scenario cards are the
required dataset.

\begin{center}\rule{0.5\linewidth}{0.5pt}\end{center}

\subsection*{Data Tables}\label{data-tables-2}
\addcontentsline{toc}{subsection}{Data Tables}

{\def\LTcaptype{none} % do not increment counter
\begin{longtable}[]{@{}lllll@{}}
\toprule\noalign{}
Step & Quantity & Symbol & Computed value & Units \\
\midrule\noalign{}
\endhead
\bottomrule\noalign{}
\endlastfoot
1 & Prediction error & \(\varepsilon\) & & m \\
2 & Kalman gain & \(K\) & & --- \\
3 & Posterior mean & \(\mu_1\) & & m \\
4 & Posterior variance & \(\sigma_1^2\) & & m\texorpdfstring{\textsuperscript{2}}{2} \\
5a & \(\mu_1\) with very noisy cue & & & m \\
5b & \(\mu_1\) with very precise cue & & & m \\
\end{longtable}
}

{\def\LTcaptype{none} % do not increment counter
\begin{longtable}[]{@{}
  >{\raggedright\arraybackslash}p{(\linewidth - 6\tabcolsep) * \real{0.1912}}
  >{\raggedright\arraybackslash}p{(\linewidth - 6\tabcolsep) * \real{0.2794}}
  >{\raggedright\arraybackslash}p{(\linewidth - 6\tabcolsep) * \real{0.3088}}
  >{\raggedright\arraybackslash}p{(\linewidth - 6\tabcolsep) * \real{0.2206}}@{}}
\toprule\noalign{}
\begin{minipage}[b]{\linewidth}\raggedright
Observation \#
\end{minipage} & \begin{minipage}[b]{\linewidth}\raggedright
Behaviour described
\end{minipage} & \begin{minipage}[b]{\linewidth}\raggedright
Perceptual or active?
\end{minipage} & \begin{minipage}[b]{\linewidth}\raggedright
Brief reasoning
\end{minipage} \\
\midrule\noalign{}
\endhead
\bottomrule\noalign{}
\endlastfoot
1 & & & \\
2 & & & \\
3 & & & \\
4 & & & \\
5 & & & \\
\end{longtable}
}

\begin{center}\rule{0.5\linewidth}{0.5pt}\end{center}

\subsection*{Evidence and Reproducibility
Checklist}\label{evidence-and-reproducibility-checklist-2}
\addcontentsline{toc}{subsection}{Evidence and Reproducibility
Checklist}

\begin{itemize}
\tightlist
\item
  \textbf{Primary evidence goal:} Model sensing and action as coupled
  inference in a simple organism.
\item
  \textbf{Data skill to practice:} Read a small probability table and
  update a prediction after new evidence.
\item
  \textbf{BioSkills emphasis:} Modeling and simulation, Process of
  science, Quantitative reasoning.
\item
  \textbf{Control logic:} identify at least one positive control, one
  negative control, or one baseline comparison before interpreting
  results.
\item
  \textbf{Measurement discipline:} record units, uncertainty, sample
  size, and any discarded observation with a reason.
\item
  \textbf{Mechanistic link:} connect one result directly to the parent
  chapter's big idea before writing the conclusion.
\item
  \textbf{Reproducibility check:} state one procedural detail that
  another group would need in order to reproduce the result.
\end{itemize}

\subsection*{Paper-Based Evidence
Upgrade}\label{paper-based-evidence-upgrade-2}
\addcontentsline{toc}{subsection}{Paper-Based Evidence Upgrade}

Before answering the analysis questions, annotate the paper dataset for
\textbf{Active Inference and the Free Energy Principle} with a
reproducibility pass:

{\def\LTcaptype{none} % do not increment counter
\begin{longtable}[]{@{}
  >{\raggedright\arraybackslash}p{(\linewidth - 2\tabcolsep) * \real{0.5000}}
  >{\raggedright\arraybackslash}p{(\linewidth - 2\tabcolsep) * \real{0.5000}}@{}}
\toprule\noalign{}
\begin{minipage}[b]{\linewidth}\raggedright
Evidence check
\end{minipage} & \begin{minipage}[b]{\linewidth}\raggedright
Student action
\end{minipage} \\
\midrule\noalign{}
\endhead
\bottomrule\noalign{}
\endlastfoot
Control logic & Mark the comparison that functions as the baseline,
negative control, or reference case. \\
Uncertainty & Circle the row, card, diagram feature, or model assumption
most likely to change the conclusion. \\
Model comparison & State whether a simpler rule, null model, or
alternative mechanism could explain the same pattern. \\
Decision threshold & Write the minimum evidence that would make you
revise the interpretation. \\
Reproducibility & Record the exact scoring rule another group would need
to reproduce your classification. \\
\end{longtable}
}

Focus note: When you annotate the Bayesian-update and
perceptual-versus-active-inference datasets, the load-bearing numbers
are the prior variance, the sensory variance, and the resulting Kalman
gain: check that the posterior shift is reproducible from those three
values alone, and that each behaviour-classification verdict cites the
specific prediction-error a sceptic could recompute. Keep required work
paper-based; any material-handling or equipment version belongs only in
an optional extension.

\subsection*{Analysis Questions}\label{analysis-questions-2}
\addcontentsline{toc}{subsection}{Analysis Questions}

\begin{enumerate}
\def\labelenumi{\arabic{enumi}.}
\tightlist
\item
  Explain in one paragraph how the Kalman gain can be read as
  ``precision-weighted uncertainty reduction'' --- why \(K\) depends on
  the ratio of prior variance to total variance.
\item
  The generative model in Part 2 uses a \emph{tight} prior (low
  variance) on body temperature. How would the calculation change if the
  prior were loose (\(\sigma = 5\) °C)? Which regime corresponds to a
  pre-homeostatic organism; which to a strictly regulated endotherm?
\item
  In humans, chronic anxiety is sometimes described as a ``tight prior
  on safety.'' Using your computation from Part 2, explain why normal
  day-to-day sensory noise becomes persistently surprising to an anxious
  agent and what ``widening the prior'' (via exposure therapy or
  pharmacology) is doing mathematically.
\item
  Scientists have attempted to build artificial life --- robots that
  minimise free energy with respect to a battery-level sensor. Describe
  what the agent's \textbf{prior} on battery state and \textbf{action
  repertoire} would need to be so that the robot ``feels hungry'' at low
  charge and seeks a charging pad.
\item
  The textbook chapter notes that evolution itself can be cast as
  free-energy minimisation at the population scale: organisms whose
  generative models fail are eliminated. Argue for or against this
  framing by contrasting it with the Modern Synthesis's emphasis on
  fitness differentials.
\end{enumerate}

\begin{center}\rule{0.5\linewidth}{0.5pt}\end{center}

\subsection*{Discussion Questions}\label{discussion-questions-7}
\addcontentsline{toc}{subsection}{Discussion Questions}

\begin{enumerate}
\def\labelenumi{\arabic{enumi}.}
\tightlist
\item
  Is ``surprise'' a purely mathematical quantity, or does it have
  phenomenal correlates? Consider how a newborn infant's surprise
  differs from an adult's when encountering the same stimulus.
\item
  Could two agents with identical sensory input hold different beliefs
  and both be rational under active inference? If yes, what determines
  each?
\item
  Does the free-energy principle predict that living things should seek
  complexity or simplicity in their environments? Justify using the
  trade-off between epistemic and instrumental value.
\end{enumerate}

\begin{center}\rule{0.5\linewidth}{0.5pt}\end{center}

\subsection*{Safety and Ethics Notes}\label{safety-and-ethics-notes-2}
\addcontentsline{toc}{subsection}{Safety and Ethics Notes}

This lab is entirely paper-based and computational --- no reagents, no
instruments, no risk. Optional observations outside class must not
disturb animals. If recording human behaviour, obtain consent and
preserve anonymity.

\emph{Module: \breaktt{src/biology/neuroscience/neuroscience.py}
(Bayesian update illustrations).}

\newpage

\section*{Lab 1 --- Atoms, Molecules, and Chemical
Bonds}\label{lab-1-atoms-molecules-and-chemical-bonds}
\addcontentsline{toc}{section}{Lab 1 --- Atoms, Molecules, and Chemical
Bonds}

\label{sec:lab_unit_I_atoms_molecules}

\subsection*{Learning Objectives}\label{learning-objectives-41}
\addcontentsline{toc}{subsection}{Learning Objectives}

\emph{This activity accompanies \cref{sec:unit_I_atoms_molecules} of the
textbook --- review that chapter before attempting the exercises below.}

\begin{itemize}
\tightlist
\item
  Identify polar and non-polar molecular regions using electronegativity
  reasoning
\item
  Interpret hydrogen bonding properties using printed observations and
  molecular cards
\item
  Distinguish ionic, covalent, and hydrogen bonds from structural and
  data evidence
\item
  Apply molecular polarity concepts to explain biological solubility
  patterns
\end{itemize}

\subsubsection*{Alignment and Rubric
Map}\label{alignment-and-rubric-map-3}
\addcontentsline{toc}{subsubsection}{Alignment and Rubric Map}

\begin{itemize}
\tightlist
\item
  \textbf{Outcome 1 (LO1):} Interpret the supplied evidence or model
  output for Atoms, Molecules, and Chemical Bonds.
\item
  \textbf{Outcome 2 (LO2):} Identify controls and comparison groups that
  make the claim testable.
\item
  \textbf{Outcome 3 (LO3):} Quantify uncertainty, boundary conditions,
  or alternative explanations before concluding.
\item
  \textbf{Outcome 4 (LO4):} Transfer the mechanism to a new biological
  case or public-facing decision.
\item
  \textbf{Chapter LO coverage:} LO1, LO2, LO3, LO4
\item
  \textbf{Rubric dimensions:} evidence; controls; uncertainty;
  mechanism; transfer. \#\# Background \{.unnumbered\}
\end{itemize}

Chemistry underpins every biological process. Atoms bond through
electronegativity differences: when atoms share electrons unequally,
they form polar covalent bonds (as in O--H), creating partial charges
(\texorpdfstring{\ensuremath{\delta}}{delta}+ and \texorpdfstring{\ensuremath{\delta}}{delta}\texorpdfstring{\ensuremath{-}}{-}). These partial charges enable hydrogen bonding --- the weak
but pervasive force that gives water its remarkable properties and
stabilises the three-dimensional shapes of proteins and DNA.

In this lab, you will use a paper dataset of solubility and
surface-tension observations as proxies for molecular polarity and
hydrogen bonding capacity: ``like dissolves like'' --- polar solutes
dissolve in polar solvents (water), while non-polar solutes dissolve
primarily in non-polar solvents (oil). You will also visualise electron
sharing differences using electronegativity tables and molecule cards.

\subsection*{Paper-Based Materials}\label{paper-based-materials-3}
\addcontentsline{toc}{subsection}{Paper-Based Materials}

{\def\LTcaptype{none} % do not increment counter
\begin{longtable}[]{@{}
  >{\raggedright\arraybackslash}p{(\linewidth - 2\tabcolsep) * \real{0.3333}}
  >{\raggedright\arraybackslash}p{(\linewidth - 2\tabcolsep) * \real{0.6667}}@{}}
\toprule\noalign{}
\begin{minipage}[b]{\linewidth}\raggedright
Item
\end{minipage} & \begin{minipage}[b]{\linewidth}\raggedright
Quantity
\end{minipage} \\
\midrule\noalign{}
\endhead
\bottomrule\noalign{}
\endlastfoot
Molecule cards: water, oil-like hydrocarbon, NaCl, sucrose, dye A, dye
B, soap & 1 set \\
Printed solubility observation dataset & 1 \\
Printed surface-tension observation card & 1 \\
Electronegativity table (provided) & 1 \\
Graph paper or notebook paper & 2 sheets \\
Calculator & 1 \\
\end{longtable}
}

\subsection*{Paper-Based
Investigation}\label{paper-based-investigation-3}
\addcontentsline{toc}{subsection}{Paper-Based Investigation}

\begin{enumerate}
\def\labelenumi{\arabic{enumi}.}
\tightlist
\item
  Sort the molecule cards into polar, non-polar, ionic, and amphipathic
  categories. Record one structural feature that justifies each
  placement.
\item
  Use the printed observation dataset to classify each solute/solvent
  pair as soluble, partially soluble, or insoluble.
\item
  Identify one positive control for polar solubility, one negative
  control, and one ambiguous case that would need more evidence.
\item
  Using the electronegativity table, calculate the electronegativity
  difference (\texorpdfstring{\ensuremath{\Delta}}{Delta}EN) for N--H, O--H, C--H, and C--C bonds. Record which
  are polar and which are non-polar.
\item
  Read the surface-tension card showing a paperclip floating before soap
  and sinking after soap. Explain the before/after pattern using
  hydrogen bonding and amphipathic disruption.
\item
  Re-score one card after a peer group challenges your classification.
  Note whether the change improved consistency with the evidence.
\end{enumerate}

\subsection*{Data Recording}\label{data-recording}
\addcontentsline{toc}{subsection}{Data Recording}

{\def\LTcaptype{none} % do not increment counter
\begin{longtable}[]{@{}lllll@{}}
\toprule\noalign{}
Cup & Solvent & Solute/Dye & Dissolves? (Y/N) & Observations \\
\midrule\noalign{}
\endhead
\bottomrule\noalign{}
\endlastfoot
A & Water & NaCl & & \\
B & Water & Sucrose & & \\
C & Oil & NaCl & & \\
D & Oil & Sucrose & & \\
E & Water & Water dye & & \\
F & Oil & Turmeric & & \\
\end{longtable}
}

{\def\LTcaptype{none} % do not increment counter
\begin{longtable}[]{@{}lll@{}}
\toprule\noalign{}
Bond & \texorpdfstring{\ensuremath{\Delta}}{Delta}EN & Polar? \\
\midrule\noalign{}
\endhead
\bottomrule\noalign{}
\endlastfoot
O--H & & \\
N--H & & \\
C--H & & \\
C--C & & \\
\end{longtable}
}

Paperclip float: \_\_\_\_\_\_\_\_ After soap: \_\_\_\_\_\_\_\_

\subsection*{Evidence and Reproducibility
Checklist}\label{evidence-and-reproducibility-checklist-3}
\addcontentsline{toc}{subsection}{Evidence and Reproducibility
Checklist}

\begin{itemize}
\tightlist
\item
  \textbf{Primary evidence goal:} Classify bonds and solvent behaviour
  using simple observations.
\item
  \textbf{Data skill to practice:} Use tabular atomic data to predict
  polarity, solubility, and biological reactivity.
\item
  \textbf{BioSkills emphasis:} Quantitative reasoning, Interdisciplinary
  nature of science.
\item
  \textbf{Control logic:} identify at least one positive control, one
  negative control, or one baseline comparison before interpreting
  results.
\item
  \textbf{Measurement discipline:} record units, uncertainty, sample
  size, and any discarded observation with a reason.
\item
  \textbf{Mechanistic link:} connect one result directly to the parent
  chapter's big idea before writing the conclusion.
\item
  \textbf{Reproducibility check:} state one procedural detail that
  another group would need in order to reproduce the result.
\end{itemize}

\subsection*{Paper-Based Evidence
Upgrade}\label{paper-based-evidence-upgrade-3}
\addcontentsline{toc}{subsection}{Paper-Based Evidence Upgrade}

Before answering the analysis questions, annotate the paper dataset for
\textbf{Atoms, Molecules, and Chemical Bonds} with a reproducibility
pass:

{\def\LTcaptype{none} % do not increment counter
\begin{longtable}[]{@{}
  >{\raggedright\arraybackslash}p{(\linewidth - 2\tabcolsep) * \real{0.5000}}
  >{\raggedright\arraybackslash}p{(\linewidth - 2\tabcolsep) * \real{0.5000}}@{}}
\toprule\noalign{}
\begin{minipage}[b]{\linewidth}\raggedright
Evidence check
\end{minipage} & \begin{minipage}[b]{\linewidth}\raggedright
Student action
\end{minipage} \\
\midrule\noalign{}
\endhead
\bottomrule\noalign{}
\endlastfoot
Control logic & Mark the comparison that functions as the baseline,
negative control, or reference case. \\
Uncertainty & Circle the row, card, diagram feature, or model assumption
most likely to change the conclusion. \\
Model comparison & State whether a simpler rule, null model, or
alternative mechanism could explain the same pattern. \\
Decision threshold & Write the minimum evidence that would make you
revise the interpretation. \\
Reproducibility & Record the exact scoring rule another group would need
to reproduce your classification. \\
\end{longtable}
}

Focus note: read each solubility result as a test of bond polarity --
treat the water-versus-oil pairing as the control contrast, and decide
in advance how strong a colour or layering change counts as
``dissolved'' for an ionic versus a non-polar solute. Keep required work
paper-based; any material-handling or equipment version belongs only in
an optional extension.

\subsection*{Analysis Questions}\label{analysis-questions-3}
\addcontentsline{toc}{subsection}{Analysis Questions}

\begin{enumerate}
\def\labelenumi{\arabic{enumi}.}
\tightlist
\item
  Which solutes dissolved in water but not oil? Which dissolved in oil
  but not water? Explain in terms of bond polarity.
\item
  Why did NaCl dissolve in water but not in vegetable oil? What type of
  bond holds NaCl together, and how does water break it apart?
\item
  Based on \texorpdfstring{\ensuremath{\Delta}}{Delta}EN values, rank the four bonds from most to least polar. How
  does bond polarity influence molecular behaviour in biological
  systems?
\item
  The paperclip floated due to surface tension. What property of water
  creates surface tension, and why did soap destroy it? What is the
  biological significance of surfactants?
\item
  Predict whether cholesterol (a steroid with a polar hydroxyl group and
  a large non-polar body) would dissolve in blood plasma (aqueous) or in
  a membrane (lipid). Justify your answer using polarity principles.
\end{enumerate}

\subsection*{Optional Hands-On
Extension}\label{optional-hands-on-extension}
\addcontentsline{toc}{subsection}{Optional Hands-On Extension}

If the instructor approves a demonstration, the printed solubility and
surface-tension cases may be reproduced with food-grade materials. The
paper dataset remains the required evidence source; any demonstration
should be treated as a check on the model, not as a prerequisite for
completing the lab.

\subsection*{Safety and Ethics Notes}\label{safety-and-ethics-notes-3}
\addcontentsline{toc}{subsection}{Safety and Ethics Notes}

Paper-based lab --- no reagents, no required wet materials, and no
physical risk. For optional demonstrations, avoid ingesting samples and
clean any spills immediately.

\subsection*{Debrief and Reflection}\label{debrief-and-reflection}
\addcontentsline{toc}{subsection}{Debrief and Reflection}

After you finish the practical work, spend 5--10 minutes in your small
group comparing results and discussing the following prompts. Each
member should contribute at least one observation before moving to the
next prompt:

\begin{enumerate}
\def\labelenumi{\arabic{enumi}.}
\tightlist
\item
  \textbf{What did your measurements show} --- compare the group's
  results to the textbook's predictions. Where they diverge, suggest at
  least one mechanistic explanation before concluding ``experimental
  error.''
\item
  \textbf{What would change the outcome} --- propose one modification to
  the procedure that would sharpen the measurement or extend the result
  to a new biological context, and predict what you would observe.
\item
  \textbf{One-sentence headline} --- each student composes a single
  sentence summarising the lab's take-home message, suitable for a
  tweet. Compare sentences across groups; good headlines are short,
  quantitative, and mechanistic.
\item
  \textbf{Connection back to the textbook} --- identify one section of
  \cref{sec:unit_I_atoms_molecules} that your data either confirmed or
  complicated. Cite the specific passage.
\end{enumerate}

\subsection*{Further Reading (Lab)}\label{further-reading-lab}
\addcontentsline{toc}{subsection}{Further Reading (Lab)}

\begin{itemize}
\tightlist
\item
  Revisit the parent chapter \cref{sec:unit_I_atoms_molecules} for the
  theoretical foundations on which this lab is built.
\item
  Look up any \textbf{bolded glossary term} introduced in the textbook
  chapter (each has a \texttt{\#gl:term-slug} link in the text) --- its
  master definition is in \breaktt{manuscript/glossary.md}.
\item
  Explore the appended \texttt{src/} module that implements the
  corresponding quantitative model (when applicable) --- referenced in
  the parent chapter's ``Bridge to Computation'' subsection.
\end{itemize}

\emph{Module footer: parent chapter
\texttt{\textbackslash{}cref\{sec:unit\_I\_atoms\_molecules\}}; all
numerical quantities in this lab use SI units --- see Appendix D of the
textbook for unit conversions and biological-scale reference values.}

\newpage

\section*{Lab 2 --- Water --- The Molecule of
Life}\label{lab-2-water-the-molecule-of-life}
\addcontentsline{toc}{section}{Lab 2 --- Water --- The Molecule of Life}

\label{sec:lab_unit_I_water_and_life}

\subsection*{Learning Objectives}\label{learning-objectives-42}
\addcontentsline{toc}{subsection}{Learning Objectives}

\emph{This activity accompanies \cref{sec:unit_I_water_and_life} of the
textbook --- review that chapter before attempting the exercises below.}

\begin{itemize}
\tightlist
\item
  Compare specific heat capacity patterns for water and oil using a
  printed dataset
\item
  Explain cohesion, adhesion, and capillarity from diagrams and model
  observations
\item
  Use solute concentration data to reason about water potential and
  osmosis
\item
  Relate water's physical properties to biological functions
\end{itemize}

\subsubsection*{Alignment and Rubric
Map}\label{alignment-and-rubric-map-4}
\addcontentsline{toc}{subsubsection}{Alignment and Rubric Map}

\begin{itemize}
\tightlist
\item
  \textbf{Outcome 1 (LO1):} Interpret the supplied evidence or model
  output for Water --- The Molecule of Life.
\item
  \textbf{Outcome 2 (LO2):} Identify controls and comparison groups that
  make the claim testable.
\item
  \textbf{Outcome 3 (LO3):} Quantify uncertainty, boundary conditions,
  or alternative explanations before concluding.
\item
  \textbf{Outcome 4 (LO4):} Transfer the mechanism to a new biological
  case or public-facing decision.
\item
  \textbf{Chapter LO coverage:} LO1, LO2, LO3, LO4
\item
  \textbf{Rubric dimensions:} evidence; controls; uncertainty;
  mechanism; transfer. \#\# Background \{.unnumbered\}
\end{itemize}

Water's extraordinary properties arise from its bent molecular geometry
(104.5° H--O--H bond angle) and the resulting network of hydrogen bonds.
Each water molecule can form up to four hydrogen bonds --- two through
lone pairs (O as H-bond acceptor) and two through O--H groups (H-bond
donor). This network gives water its high specific heat (4.18 J/g\texorpdfstring{\ensuremath{\cdot}}{.}°C),
high surface tension (72 mN/m), strong cohesion/adhesion, and the
ability to dissolve ionic and polar solutes. Most of these properties
are exploited by living organisms.

In this lab you will analyse heat-absorption data for water and oil,
interpret capillary-rise diagrams, and track osmotic movement in a model
membrane dataset. The required work is paper-based so the scientific
reasoning is centred on controls, units, gradients, and reproducible
interpretation.

\subsection*{Paper-Based Materials}\label{paper-based-materials-4}
\addcontentsline{toc}{subsection}{Paper-Based Materials}

{\def\LTcaptype{none} % do not increment counter
\begin{longtable}[]{@{}
  >{\raggedright\arraybackslash}p{(\linewidth - 2\tabcolsep) * \real{0.3333}}
  >{\raggedright\arraybackslash}p{(\linewidth - 2\tabcolsep) * \real{0.6667}}@{}}
\toprule\noalign{}
\begin{minipage}[b]{\linewidth}\raggedright
Item
\end{minipage} & \begin{minipage}[b]{\linewidth}\raggedright
Quantity
\end{minipage} \\
\midrule\noalign{}
\endhead
\bottomrule\noalign{}
\endlastfoot
Heat-capacity dataset: water and oil temperature over time & 1 \\
Capillary-rise diagram packet for water, oil, and xylem & 1 \\
Osmosis model dataset for membrane bags in different solute conditions &
1 \\
Control-identification cards & 1 set \\
Ruler (mm) & 1 \\
Calculator and graph paper & 1 each \\
\end{longtable}
}

\subsection*{Paper-Based
Investigation}\label{paper-based-investigation-4}
\addcontentsline{toc}{subsection}{Paper-Based Investigation}

\textbf{Part A --- Specific Heat}

\begin{enumerate}
\def\labelenumi{\arabic{enumi}.}
\tightlist
\item
  Graph the provided temperature data for equal masses of water and oil
  under the same heat input.
\item
  Calculate \texorpdfstring{\ensuremath{\Delta}}{Delta}T per 2-minute interval and identify which substance
  buffered temperature change more strongly.
\item
  Mark the baseline, the controlled variables, and one possible
  measurement uncertainty on the graph.
\end{enumerate}

\textbf{Part B --- Capillarity}

\begin{enumerate}
\def\labelenumi{\arabic{enumi}.}
\setcounter{enumi}{3}
\tightlist
\item
  Use the capillary-rise diagrams to measure the height of model water
  and oil columns in millimetres.
\item
  Compare the xylem diagram with the capillary tube model. Identify
  where adhesion, cohesion, and tube diameter enter the explanation.
\item
  Write one prediction for how capillary height would change if tube
  radius doubled, then check it against the model card.
\end{enumerate}

\textbf{Part C --- Osmosis}

\begin{enumerate}
\def\labelenumi{\arabic{enumi}.}
\setcounter{enumi}{6}
\tightlist
\item
  Use the membrane-bag dataset to calculate \texorpdfstring{\ensuremath{\Delta}}{Delta}mass over time for
  hypotonic, isotonic, and hypertonic surroundings.
\item
  Draw the water-potential gradient for each condition and identify
  which case functions as the best negative control for net water
  movement.
\end{enumerate}

\subsection*{Data Recording}\label{data-recording-1}
\addcontentsline{toc}{subsection}{Data Recording}

{\def\LTcaptype{none} % do not increment counter
\begin{longtable}[]{@{}lll@{}}
\toprule\noalign{}
Time (min) & Water Temp (°C) & Oil Temp (°C) \\
\midrule\noalign{}
\endhead
\bottomrule\noalign{}
\endlastfoot
0 & & \\
2 & & \\
4 & & \\
6 & & \\
8 & & \\
10 & & \\
\end{longtable}
}

Capillary rise --- water: \_\_\_\_\_\_ mm; oil: \_\_\_\_\_\_ mm

{\def\LTcaptype{none} % do not increment counter
\begin{longtable}[]{@{}lll@{}}
\toprule\noalign{}
Time (min) & Dialysis tubing mass (g) & \texorpdfstring{\ensuremath{\Delta}}{Delta}Mass (g) \\
\midrule\noalign{}
\endhead
\bottomrule\noalign{}
\endlastfoot
0 & & --- \\
5 & & \\
10 & & \\
20 & & \\
\end{longtable}
}

\subsection*{Evidence and Reproducibility
Checklist}\label{evidence-and-reproducibility-checklist-4}
\addcontentsline{toc}{subsection}{Evidence and Reproducibility
Checklist}

\begin{itemize}
\tightlist
\item
  \textbf{Primary evidence goal:} Measure surface tension, polarity, and
  buffering behaviour.
\item
  \textbf{Data skill to practice:} Convert between pH, hydrogen ion
  concentration, and buffer ratios.
\item
  \textbf{BioSkills emphasis:} Quantitative reasoning, Interdisciplinary
  nature of science.
\item
  \textbf{Control logic:} identify at least one positive control, one
  negative control, or one baseline comparison before interpreting
  results.
\item
  \textbf{Measurement discipline:} record units, uncertainty, sample
  size, and any discarded observation with a reason.
\item
  \textbf{Mechanistic link:} connect one result directly to the parent
  chapter's big idea before writing the conclusion.
\item
  \textbf{Reproducibility check:} state one procedural detail that
  another group would need in order to reproduce the result.
\end{itemize}

\subsection*{Paper-Based Evidence
Upgrade}\label{paper-based-evidence-upgrade-4}
\addcontentsline{toc}{subsection}{Paper-Based Evidence Upgrade}

Before answering the analysis questions, annotate the paper dataset for
\textbf{Water --- The Molecule of Life} with a reproducibility pass:

{\def\LTcaptype{none} % do not increment counter
\begin{longtable}[]{@{}
  >{\raggedright\arraybackslash}p{(\linewidth - 2\tabcolsep) * \real{0.5000}}
  >{\raggedright\arraybackslash}p{(\linewidth - 2\tabcolsep) * \real{0.5000}}@{}}
\toprule\noalign{}
\begin{minipage}[b]{\linewidth}\raggedright
Evidence check
\end{minipage} & \begin{minipage}[b]{\linewidth}\raggedright
Student action
\end{minipage} \\
\midrule\noalign{}
\endhead
\bottomrule\noalign{}
\endlastfoot
Control logic & Mark the comparison that functions as the baseline,
negative control, or reference case. \\
Uncertainty & Circle the row, card, diagram feature, or model assumption
most likely to change the conclusion. \\
Model comparison & State whether a simpler rule, null model, or
alternative mechanism could explain the same pattern. \\
Decision threshold & Write the minimum evidence that would make you
revise the interpretation. \\
Reproducibility & Record the exact scoring rule another group would need
to reproduce your classification. \\
\end{longtable}
}

Focus note: this lab compares a thermal-property measurement (water
versus oil heating rate, reporting specific heat) against a separate
interfacial measurement (capillary rise from adhesion and cohesion);
keep the two evidence chains distinct and state the temperature or
height resolution at which a difference becomes convincing. Keep
required work paper-based; any material-handling or equipment version
belongs only in an optional extension.

\subsection*{Analysis Questions}\label{analysis-questions-4}
\addcontentsline{toc}{subsection}{Analysis Questions}

\begin{enumerate}
\def\labelenumi{\arabic{enumi}.}
\tightlist
\item
  Which substance heated more quickly --- water or oil? How does this
  reflect differences in specific heat, and what is the ecological
  significance for oceans moderating coastal climates?
\item
  Why did water rise higher than oil in the capillary tube? Name the two
  adhesive forces responsible and identify the molecules involved.
\item
  Explain what you observed in the celery xylem. Which water property
  --- cohesion, adhesion, or both --- most directly drives this upward
  movement?
\item
  In the osmosis experiment, did the dialysis tubing gain or lose mass?
  Draw and explain the water potential gradient that caused this
  direction of movement.
\item
  Patients receiving intravenous fluids receive 0.9\% NaCl (normal
  saline) rather than pure water. Predict what would happen if red blood
  cells were placed in pure water vs 3\% NaCl solution. Explain using
  water potential (\(\Psi = \Psi_s + \Psi_p\)).
\end{enumerate}

\subsection*{Optional Hands-On
Extension}\label{optional-hands-on-extension-1}
\addcontentsline{toc}{subsection}{Optional Hands-On Extension}

An instructor may demonstrate capillary rise or osmosis with household
materials, but the printed dataset is the required source for analysis.
Treat any demonstration as an illustration that can be compared with the
model, not as required evidence.

\subsection*{Safety and Ethics Notes}\label{safety-and-ethics-notes-4}
\addcontentsline{toc}{subsection}{Safety and Ethics Notes}

Paper-based lab --- no hot water baths, glassware, membranes, or wet
materials are required. Optional demonstrations should use low-risk
materials and ordinary spill precautions.

\subsection*{Debrief and Reflection}\label{debrief-and-reflection-1}
\addcontentsline{toc}{subsection}{Debrief and Reflection}

After you finish the practical work, spend 5--10 minutes in your small
group comparing results and discussing the following prompts. Each
member should contribute at least one observation before moving to the
next prompt:

\begin{enumerate}
\def\labelenumi{\arabic{enumi}.}
\tightlist
\item
  \textbf{What did your measurements show} --- compare the group's
  results to the textbook's predictions. Where they diverge, suggest at
  least one mechanistic explanation before concluding ``experimental
  error.''
\item
  \textbf{What would change the outcome} --- propose one modification to
  the procedure that would sharpen the measurement or extend the result
  to a new biological context, and predict what you would observe.
\item
  \textbf{One-sentence headline} --- each student composes a single
  sentence summarising the lab's take-home message, suitable for a
  tweet. Compare sentences across groups; good headlines are short,
  quantitative, and mechanistic.
\item
  \textbf{Connection back to the textbook} --- identify one section of
  \cref{sec:unit_I_water_and_life} that your data either confirmed or
  complicated. Cite the specific passage.
\end{enumerate}

\subsection*{Further Reading (Lab)}\label{further-reading-lab-1}
\addcontentsline{toc}{subsection}{Further Reading (Lab)}

\begin{itemize}
\tightlist
\item
  Revisit the parent chapter \cref{sec:unit_I_water_and_life} for the
  theoretical foundations on which this lab is built.
\item
  Look up any \textbf{bolded glossary term} introduced in the textbook
  chapter (each has a \texttt{\#gl:term-slug} link in the text) --- its
  master definition is in \breaktt{manuscript/glossary.md}.
\item
  Explore the appended \texttt{src/} module that implements the
  corresponding quantitative model (when applicable) --- referenced in
  the parent chapter's ``Bridge to Computation'' subsection.
\end{itemize}

\emph{Module footer: parent chapter
\texttt{\textbackslash{}cref\{sec:unit\_I\_water\_and\_life\}}; all
numerical quantities in this lab use SI units --- see Appendix D of the
textbook for unit conversions and biological-scale reference values.}

\newpage

\section*{Lab 3 --- Biological
Macromolecules}\label{lab-3-biological-macromolecules}
\addcontentsline{toc}{section}{Lab 3 --- Biological Macromolecules}

\label{sec:lab_unit_I_macromolecules}

\subsection*{Learning Objectives}\label{learning-objectives-43}
\addcontentsline{toc}{subsection}{Learning Objectives}

\emph{This activity accompanies \cref{sec:unit_I_macromolecules} of the
textbook --- review that chapter before attempting the exercises below.}

\begin{itemize}
\tightlist
\item
  Apply simulated colorimetric-test evidence to identify carbohydrates,
  proteins, and lipids
\item
  Distinguish reducing sugars from non-reducing sugars using structural
  logic
\item
  Identify proteins and lipids from reaction cards, labels, and model
  results
\item
  Interpret qualitative colour evidence with controls, uncertainty, and
  reproducibility checks
\end{itemize}

\subsubsection*{Alignment and Rubric
Map}\label{alignment-and-rubric-map-5}
\addcontentsline{toc}{subsubsection}{Alignment and Rubric Map}

\begin{itemize}
\tightlist
\item
  \textbf{Outcome 1 (LO1):} Interpret the supplied evidence or model
  output for Biological Macromolecules.
\item
  \textbf{Outcome 2 (LO2):} Identify controls and comparison groups that
  make the claim testable.
\item
  \textbf{Outcome 3 (LO3):} Quantify uncertainty, boundary conditions,
  or alternative explanations before concluding.
\item
  \textbf{Outcome 4 (LO4):} Transfer the mechanism to a new biological
  case or public-facing decision.
\item
  \textbf{Chapter LO coverage:} LO1, LO2, LO3, LO4
\item
  \textbf{Rubric dimensions:} evidence; controls; uncertainty;
  mechanism; transfer. \#\# Pre-Lab (to complete before class)
  \{.unnumbered\}
\end{itemize}

\begin{enumerate}
\def\labelenumi{\arabic{enumi}.}
\tightlist
\item
  For each of the four macromolecule classes, write one sentence that
  names its monomer, its primary covalent bond, and one biological
  function.
\item
  classic reducing-sugar assay detects \textbf{aldehyde} groups on
  reducing sugars. Which of glucose, fructose, sucrose, and starch have
  a free aldehyde? For those without, explain why (look up hemiketal and
  glycosidic bonds).
\item
  Compute the energy density of pure triglyceride (37 kJ g\texorpdfstring{\textsuperscript{-}}{-}\texorpdfstring{\textsuperscript{1}}{1}) vs pure
  starch (17 kJ g\texorpdfstring{\textsuperscript{-}}{-}\texorpdfstring{\textsuperscript{1}}{1}) vs pure protein (17 kJ g\texorpdfstring{\textsuperscript{-}}{-}\texorpdfstring{\textsuperscript{1}}{1}). Why do organisms store
  long-term energy as fat rather than glycogen?
\item
  Predict which of the following will yield a positive classic
  peptide-bond assay test at room temperature: a single amino acid
  (glycine), a dipeptide, albumin. Justify in one sentence (classic
  peptide-bond assay requires \texorpdfstring{\textgreater{}=}{>=}2 peptide bonds \texorpdfstring{\ensuremath{\to}}{->} at least tripeptide).
\end{enumerate}

\subsection*{Background}\label{background}
\addcontentsline{toc}{subsection}{Background}

Macromolecules --- carbohydrates, lipids, proteins, and nucleic acids
--- form the molecular architecture of cells. Each class has distinctive
chemical groups that produce diagnostic patterns in classic colorimetric
tests. Reducing sugars (glucose, fructose, maltose) donate electrons to
Cu\texorpdfstring{\textsuperscript{2}}{2}\texorpdfstring{\textsuperscript{+}}{+} in classic reducing-sugar assay, producing a colour change from
blue to yellow/orange/red; sucrose (non-reducing) does not react unless
hydrolysed. The classic peptide-bond assay reaction turns purple in the
presence of peptide bonds (\texorpdfstring{\textgreater{}=}{>=}2). classic lipid-partition assay, an
oil-soluble dye, partitions into lipids and stains them red-orange.

In this lab you will interpret a blinded paper dataset for five unknown
samples (A--E), compare it with positive and negative controls, and
identify the macromolecules present in each. The chemistry is treated as
evidence to reason from, not as a required wet procedure.

\subsection*{Paper-Based Materials}\label{paper-based-materials-5}
\addcontentsline{toc}{subsection}{Paper-Based Materials}

{\def\LTcaptype{none} % do not increment counter
\begin{longtable}[]{@{}
  >{\raggedright\arraybackslash}p{(\linewidth - 2\tabcolsep) * \real{0.3333}}
  >{\raggedright\arraybackslash}p{(\linewidth - 2\tabcolsep) * \real{0.6667}}@{}}
\toprule\noalign{}
\begin{minipage}[b]{\linewidth}\raggedright
Item
\end{minipage} & \begin{minipage}[b]{\linewidth}\raggedright
Quantity
\end{minipage} \\
\midrule\noalign{}
\endhead
\bottomrule\noalign{}
\endlastfoot
Blinded macromolecule reaction dataset for samples A--E & 1 \\
Molecular-structure cards: glucose, fructose, sucrose, starch, albumin,
triglyceride & 1 set \\
Positive-control and negative-control cards & 1 set \\
Hydrolysis time-course graph packet & 1 \\
AI protein-model evidence packet: confidence map, predicted interface,
ligand-contact table, and validation notes & 1 \\
Decision-matrix worksheet & 1 \\
Calculator and graph paper & 1 each \\
\end{longtable}
}

\subsection*{Paper-Based
Investigation}\label{paper-based-investigation-5}
\addcontentsline{toc}{subsection}{Paper-Based Investigation}

\begin{enumerate}
\def\labelenumi{\arabic{enumi}.}
\tightlist
\item
  Read the blinded dataset for samples A--E. For each simulated test,
  mark whether the result is positive, negative, weak positive, or
  ambiguous.
\item
  Use the control cards to verify that each test behaved as expected
  before interpreting unknowns.
\item
  Build a decision matrix: reducing sugar evidence, peptide-bond
  evidence, lipid-partition evidence, and structural clue evidence.
\item
  Assign a macromolecule profile to each unknown and rate your
  confidence as high, medium, or low.
\item
  Identify one result that would need replication or an additional test
  before a confident claim.
\end{enumerate}

\subsubsection*{Part D --- Enzyme Hydrolysis
Data}\label{part-d-enzyme-hydrolysis-data}
\addcontentsline{toc}{subsubsection}{Part D --- Enzyme Hydrolysis Data}

\begin{enumerate}
\def\labelenumi{\arabic{enumi}.}
\setcounter{enumi}{5}
\tightlist
\item
  Interpret the provided starch-hydrolysis time course for active
  amylase, denatured enzyme, and no-enzyme controls.
\item
  Estimate a crude hydrolysis rate (mg starch per minute) from the graph
  and explain why the denatured-enzyme line functions as a negative
  control.
\item
  Decide whether the evidence supports enzyme specificity, temperature
  sensitivity, or both.
\end{enumerate}

\subsubsection*{Part E --- Sucrose Hydrolysis
Model}\label{part-e-sucrose-hydrolysis-model}
\addcontentsline{toc}{subsubsection}{Part E --- Sucrose Hydrolysis
Model}

\begin{enumerate}
\def\labelenumi{\arabic{enumi}.}
\setcounter{enumi}{8}
\tightlist
\item
  Compare the structural cards for untreated sucrose and hydrolysed
  sucrose.
\item
  Predict the simulated Benedict pattern before and after hydrolysis,
  then check it against the provided model result. Explain why cleavage
  into glucose and fructose changes the reducing-sugar evidence.
\end{enumerate}

\subsubsection*{Part F --- AI Structure Evidence
Check}\label{part-f-ai-structure-evidence-check}
\addcontentsline{toc}{subsubsection}{Part F --- AI Structure Evidence
Check}

\begin{enumerate}
\def\labelenumi{\arabic{enumi}.}
\setcounter{enumi}{10}
\tightlist
\item
  Read the AI protein-model packet for a hypothetical enzyme-substrate
  interaction. Mark each predicted contact as high, medium, or low
  confidence using the confidence map and interface notes.
\item
  Decide which claim is justified: ``this residue is worth mutating,''
  ``this ligand probably binds,'' or ``this interaction is
  established.'' Explain why the last claim requires independent assay
  evidence.
\item
  Propose a paper-based validation plan: one conservation comparison,
  one predicted-disruptive mutation, one binding or kinetic readout, and
  one result that would make you reject the predicted interface.
\end{enumerate}

\subsection*{Data Recording}\label{data-recording-2}
\addcontentsline{toc}{subsection}{Data Recording}

{\def\LTcaptype{none} % do not increment counter
\begin{longtable}[]{@{}
  >{\raggedright\arraybackslash}p{(\linewidth - 8\tabcolsep) * \real{0.0750}}
  >{\raggedright\arraybackslash}p{(\linewidth - 8\tabcolsep) * \real{0.2375}}
  >{\raggedright\arraybackslash}p{(\linewidth - 8\tabcolsep) * \real{0.1875}}
  >{\raggedright\arraybackslash}p{(\linewidth - 8\tabcolsep) * \real{0.2000}}
  >{\raggedright\arraybackslash}p{(\linewidth - 8\tabcolsep) * \real{0.3000}}@{}}
\toprule\noalign{}
\begin{minipage}[b]{\linewidth}\raggedright
Sample
\end{minipage} & \begin{minipage}[b]{\linewidth}\raggedright
Reducing-sugar evidence
\end{minipage} & \begin{minipage}[b]{\linewidth}\raggedright
Peptide-bond evidence
\end{minipage} & \begin{minipage}[b]{\linewidth}\raggedright
Lipid-partition evidence
\end{minipage} & \begin{minipage}[b]{\linewidth}\raggedright
Macromolecule(s) present
\end{minipage} \\
\midrule\noalign{}
\endhead
\bottomrule\noalign{}
\endlastfoot
A & & & & \\
B & & & & \\
C & & & & \\
D & & & & \\
E & & & & \\
(+) Glucose & & --- & --- & Reducing sugar \\
(+) Albumin & --- & & --- & Protein \\
(\texorpdfstring{\ensuremath{-}}{-}) Water & & & & None \\
\end{longtable}
}

AI model evidence: highest-confidence contact = \_\_\_\_\_\_\_;
lowest-confidence region = \_\_\_\_\_\_\_; validation assay proposed =
\_\_\_\_\_\_\_; claim justified now = \_\_\_\_\_\_\_

\subsection*{Evidence and Reproducibility
Checklist}\label{evidence-and-reproducibility-checklist-5}
\addcontentsline{toc}{subsection}{Evidence and Reproducibility
Checklist}

\begin{itemize}
\tightlist
\item
  \textbf{Primary evidence goal:} Use chemical tests and structural
  clues to identify biological macromolecules.
\item
  \textbf{Data skill to practice:} Classify molecules from structural
  evidence rather than names alone.
\item
  \textbf{BioSkills emphasis:} Quantitative reasoning, Interdisciplinary
  nature of science.
\item
  \textbf{Control logic:} identify at least one positive control, one
  negative control, or one baseline comparison before interpreting
  results.
\item
  \textbf{Measurement discipline:} record units, uncertainty, sample
  size, and any discarded observation with a reason.
\item
  \textbf{Mechanistic link:} connect one result directly to the parent
  chapter's big idea before writing the conclusion.
\item
  \textbf{Reproducibility check:} state one procedural detail that
  another group would need in order to reproduce the result.
\end{itemize}

\subsection*{Paper-Based Evidence
Upgrade}\label{paper-based-evidence-upgrade-5}
\addcontentsline{toc}{subsection}{Paper-Based Evidence Upgrade}

Before answering the analysis questions, annotate the paper dataset for
\textbf{Biological Macromolecules} with a reproducibility pass:

{\def\LTcaptype{none} % do not increment counter
\begin{longtable}[]{@{}
  >{\raggedright\arraybackslash}p{(\linewidth - 2\tabcolsep) * \real{0.5000}}
  >{\raggedright\arraybackslash}p{(\linewidth - 2\tabcolsep) * \real{0.5000}}@{}}
\toprule\noalign{}
\begin{minipage}[b]{\linewidth}\raggedright
Evidence check
\end{minipage} & \begin{minipage}[b]{\linewidth}\raggedright
Student action
\end{minipage} \\
\midrule\noalign{}
\endhead
\bottomrule\noalign{}
\endlastfoot
Control logic & Mark the comparison that functions as the baseline,
negative control, or reference case. \\
Uncertainty & Circle the row, card, diagram feature, or model assumption
most likely to change the conclusion. \\
Model comparison & State whether a simpler rule, null model, or
alternative mechanism could explain the same pattern. \\
Decision threshold & Write the minimum evidence that would make you
revise the interpretation. \\
Reproducibility & Record the exact scoring rule another group would need
to reproduce your classification. \\
\end{longtable}
}

Focus note: each colour-change assay (reducing-sugar, peptide-bond, and
the lipid grease-spot test) is a separate qualitative detector with its
own false positives -- record the positive and negative reference
samples and the minimum colour shift you will accept before scoring a
macromolecule class as present. Keep required work paper-based; any
material-handling or equipment version belongs only in an optional
extension.

\subsection*{Analysis Questions}\label{analysis-questions-5}
\addcontentsline{toc}{subsection}{Analysis Questions}

\begin{enumerate}
\def\labelenumi{\arabic{enumi}.}
\tightlist
\item
  Which samples tested positive for reducing sugars? Were any positive
  for sucrose (hint: sucrose does NOT turn classic reducing-sugar assay
  red without hydrolysis)? How would you distinguish glucose from
  sucrose using primarily classic reducing-sugar assay test?
\item
  A sample of milk tested positive for both classic reducing-sugar assay
  (weakly) and classic peptide-bond assay (strongly) and showed a grease
  spot. Name the three macromolecule classes present in milk.
\item
  Why does classic reducing-sugar assay require heating to detect
  reducing sugars, while classic peptide-bond assay works at room
  temperature? What does this reveal about the reaction mechanisms?
\item
  Dietary fibre (cellulose) is a polysaccharide. Predict whether it
  would give a positive classic reducing-sugar assay test. Explain in
  terms of glycosidic bond type (\texorpdfstring{\ensuremath{\alpha}}{alpha} vs \texorpdfstring{\ensuremath{\beta}}{beta} linkage) and digestibility.
\item
  A patient has phenylketonuria (PKU) --- inability to metabolise
  phenylalanine. Which macromolecule class and which specific monomer is
  relevant? How does dietary management (low-phenylalanine diet) relate
  to the biochemistry of this lab?
\item
  An AlphaFold-derived packet predicts a ligand contact in a
  low-confidence loop. What additional evidence would move the claim
  from ``hypothesis'' to ``supported mechanism''? Name at least one
  sequence, structural, and biochemical check.
\end{enumerate}

\subsection*{Extension Analysis
Questions}\label{extension-analysis-questions}
\addcontentsline{toc}{subsection}{Extension Analysis Questions}

\begin{enumerate}
\def\labelenumi{\arabic{enumi}.}
\setcounter{enumi}{6}
\tightlist
\item
  Galactosaemia patients cannot metabolise galactose (a component of
  lactose). Using your understanding of glycosidic bonds, explain why
  such a patient must avoid milk but can safely consume glucose or
  fructose sweetened foods.
\item
  The Keto diet shifts metabolism to fat oxidation. Predict which
  macromolecule category would dominate your test results on a piece of
  bacon vs a slice of bread. Discuss why each is energy-dense but in
  biochemically different ways (immediate vs sustained release).
\item
  Protein-losing enteropathy causes albumin leakage into the gut. Given
  classic peptide-bond assay's sensitivity (about 0.1 mg/mL),
  would you expect a positive classic peptide-bond assay reaction on a
  faecal sample from such a patient? Why might faecal protein
  quantification via classic peptide-bond assay be unreliable (hint:
  digestive proteolysis)?
\end{enumerate}

\subsection*{Optional Wet-Lab
Extension}\label{optional-wet-lab-extension}
\addcontentsline{toc}{subsection}{Optional Wet-Lab Extension}

An instructor may run classic Benedict, classic peptide-bond assay,
Sudan, iodine, or hydrolysis demonstrations in a properly equipped
laboratory. These are optional demonstrations primarily; the required
activity is the paper dataset and decision matrix.

\subsection*{Safety and Ethics Notes}\label{safety-and-ethics-notes-5}
\addcontentsline{toc}{subsection}{Safety and Ethics Notes}

Paper-based lab --- no reagents, heating, saliva sampling, or unknown
foods are required. Optional wet demonstrations require instructor
supervision, approved chemical handling, and appropriate waste disposal.

\subsection*{Debrief and Reflection}\label{debrief-and-reflection-2}
\addcontentsline{toc}{subsection}{Debrief and Reflection}

After you finish the practical work, spend 5--10 minutes in your small
group comparing results and discussing the following prompts. Each
member should contribute at least one observation before moving to the
next prompt:

\begin{enumerate}
\def\labelenumi{\arabic{enumi}.}
\tightlist
\item
  \textbf{What did your measurements show} --- compare the group's
  results to the textbook's predictions. Where they diverge, suggest at
  least one mechanistic explanation before concluding ``experimental
  error.''
\item
  \textbf{What would change the outcome} --- propose one modification to
  the procedure that would sharpen the measurement or extend the result
  to a new biological context, and predict what you would observe.
\item
  \textbf{One-sentence headline} --- each student composes a single
  sentence summarising the lab's take-home message, suitable for a
  tweet. Compare sentences across groups; good headlines are short,
  quantitative, and mechanistic.
\item
  \textbf{Connection back to the textbook} --- identify one section of
  \cref{sec:unit_I_macromolecules} that your data either confirmed or
  complicated. Cite the specific passage.
\end{enumerate}

\subsection*{Further Reading (Lab)}\label{further-reading-lab-2}
\addcontentsline{toc}{subsection}{Further Reading (Lab)}

\begin{itemize}
\tightlist
\item
  Revisit the parent chapter \cref{sec:unit_I_macromolecules} for the
  theoretical foundations on which this lab is built.
\item
  Look up any \textbf{bolded glossary term} introduced in the textbook
  chapter (each has a \texttt{\#gl:term-slug} link in the text) --- its
  master definition is in \breaktt{manuscript/glossary.md}.
\item
  Explore the appended \texttt{src/} module that implements the
  corresponding quantitative model (when applicable) --- referenced in
  the parent chapter's ``Bridge to Computation'' subsection.
\end{itemize}

\emph{Module footer: parent chapter
\texttt{\textbackslash{}cref\{sec:unit\_I\_macromolecules\}}; all
numerical quantities in this lab use SI units --- see Appendix D of the
textbook for unit conversions and biological-scale reference values.}

\newpage

\section*{Lab 4 --- Enzymes and the Kinetics of
Catalysis}\label{lab-4-enzymes-and-the-kinetics-of-catalysis}
\addcontentsline{toc}{section}{Lab 4 --- Enzymes and the Kinetics of
Catalysis}

\label{sec:lab_unit_I_enzymes_and_kinetics}

\subsection*{Learning Objectives}\label{learning-objectives-44}
\addcontentsline{toc}{subsection}{Learning Objectives}

\emph{This activity accompanies \cref{sec:unit_I_enzymes_and_kinetics}
of the textbook --- review that chapter before attempting the exercises
below.}

\begin{itemize}
\tightlist
\item
  Interpret enzyme activity data by tracking O\texorpdfstring{\textsubscript{2}}{2} production as a proxy
  for reaction rate
\item
  Determine how temperature, pH, and substrate concentration affect
  enzyme rate (Vmax approximation)
\item
  Construct a rate vs \hyperref[s]{S} curve and identify qualitative Km
\item
  Interpret enzyme inhibition from a competitive inhibitor treatment
\end{itemize}

\subsubsection*{Alignment and Rubric
Map}\label{alignment-and-rubric-map-6}
\addcontentsline{toc}{subsubsection}{Alignment and Rubric Map}

\begin{itemize}
\tightlist
\item
  \textbf{Outcome 1 (LO1):} Interpret the supplied evidence or model
  output for Enzymes and the Kinetics of Catalysis.
\item
  \textbf{Outcome 2 (LO2):} Identify controls and comparison groups that
  make the claim testable.
\item
  \textbf{Outcome 3 (LO3):} Quantify uncertainty, boundary conditions,
  or alternative explanations before concluding.
\item
  \textbf{Outcome 4 (LO4):} Transfer the mechanism to a new biological
  case or public-facing decision.
\item
  \textbf{Chapter LO coverage:} LO1, LO2, LO3, LO4
\item
  \textbf{Rubric dimensions:} evidence; controls; uncertainty;
  mechanism; transfer. \#\# Background \{.unnumbered\}
\end{itemize}

Enzymes are biological catalysts that lower activation energy. The
enzyme catalase converts hydrogen peroxide (H\texorpdfstring{\textsubscript{2}}{2}O\texorpdfstring{\textsubscript{2}}{2}) to water and oxygen:
\textbf{2 H\texorpdfstring{\textsubscript{2}}{2}O\texorpdfstring{\textsubscript{2}}{2} \texorpdfstring{\ensuremath{\to}}{->} 2 H\texorpdfstring{\textsubscript{2}}{2}O + O\texorpdfstring{\textsubscript{2}}{2}}. This reaction is often monitored by
counting O\texorpdfstring{\textsubscript{2}}{2} bubbles per minute or by measuring the height of foam
produced. Catalase is one of the fastest known enzymes (k\_cat
about 40,000,000 s\texorpdfstring{\textsuperscript{-}}{-}\texorpdfstring{\textsuperscript{1}}{1}) and is found in nearly most aerobic
organisms.

In this lab you will analyse a paper dataset showing how temperature,
pH, substrate concentration, and inhibition affect catalase activity.
The goal is to infer mechanism from controlled comparisons rather than
to run a wet reaction.

\subsection*{Paper-Based Materials}\label{paper-based-materials-6}
\addcontentsline{toc}{subsection}{Paper-Based Materials}

{\def\LTcaptype{none} % do not increment counter
\begin{longtable}[]{@{}ll@{}}
\toprule\noalign{}
Item & Quantity \\
\midrule\noalign{}
\endhead
\bottomrule\noalign{}
\endlastfoot
Catalase activity dataset with three replicates per condition & 1 \\
Enzyme-structure and active-site cards & 1 set \\
Michaelis-Menten curve worksheet & 1 \\
Inhibitor scenario cards & 1 set \\
Ruler (mm) & 1 \\
Calculator and graph paper & 1 each \\
\end{longtable}
}

\subsection*{Paper-Based
Investigation}\label{paper-based-investigation-6}
\addcontentsline{toc}{subsection}{Paper-Based Investigation}

\textbf{Part A --- Temperature Effect}

\begin{enumerate}
\def\labelenumi{\arabic{enumi}.}
\tightlist
\item
  Plot the provided catalase-rate data for 0°C, 22°C, 37°C, 50°C, and
  denatured enzyme conditions.
\item
  Calculate the mean and range for each condition. Identify the
  denatured enzyme as a negative control and the 37°C condition as a
  reference comparison.
\item
  Use the active-site cards to explain why low temperature and
  denaturation reduce activity for different reasons.
\end{enumerate}

\textbf{Part B --- pH Effect}

\begin{enumerate}
\def\labelenumi{\arabic{enumi}.}
\setcounter{enumi}{3}
\tightlist
\item
  Graph catalase rate across pH 2.5, 5, 7, 8.5, and 11. Mark the pH
  optimum and one condition where enzyme structure is likely disrupted.
\item
  Decide whether pH changes primarily affect enzyme shape, substrate
  charge, or both. Support the claim with one pattern in the dataset.
\end{enumerate}

\textbf{Part C --- Substrate Concentration Effect}

\begin{enumerate}
\def\labelenumi{\arabic{enumi}.}
\setcounter{enumi}{5}
\tightlist
\item
  Plot rate versus {[}H\texorpdfstring{\textsubscript{2}}{2}O\texorpdfstring{\textsubscript{2}}{2}{]} for 0\%, 1\%, 3\%, and 6\% substrate.
\item
  Estimate the concentration where the curve begins to plateau and label
  it as a qualitative Vmax region. Use the inhibitor cards to predict
  how a competitive inhibitor would change apparent Km and Vmax.
\end{enumerate}

\subsection*{Data Recording}\label{data-recording-3}
\addcontentsline{toc}{subsection}{Data Recording}

{\def\LTcaptype{none} % do not increment counter
\begin{longtable}[]{@{}
  >{\raggedright\arraybackslash}p{(\linewidth - 8\tabcolsep) * \real{0.1406}}
  >{\raggedright\arraybackslash}p{(\linewidth - 8\tabcolsep) * \real{0.2656}}
  >{\raggedright\arraybackslash}p{(\linewidth - 8\tabcolsep) * \real{0.0312}}
  >{\raggedright\arraybackslash}p{(\linewidth - 8\tabcolsep) * \real{0.1719}}
  >{\raggedright\arraybackslash}p{(\linewidth - 8\tabcolsep) * \real{0.3906}}@{}}
\toprule\noalign{}
\begin{minipage}[b]{\linewidth}\raggedright
Condition
\end{minipage} & \begin{minipage}[b]{\linewidth}\raggedright
Temperature (°C)
\end{minipage} & \begin{minipage}[b]{\linewidth}\raggedright
pH
\end{minipage} & \begin{minipage}[b]{\linewidth}\raggedright
{[}H\texorpdfstring{\textsubscript{2}}{2}O\texorpdfstring{\textsubscript{2}}{2}{]} (\%)
\end{minipage} & \begin{minipage}[b]{\linewidth}\raggedright
Simulated O\texorpdfstring{\textsubscript{2}}{2} proxy at 60 s
\end{minipage} \\
\midrule\noalign{}
\endhead
\bottomrule\noalign{}
\endlastfoot
Ice bath & 0 & 7 & 3 & \\
Room temp & 22 & 7 & 3 & \\
Warm bath & 50 & 7 & 3 & \\
Denatured enzyme & 22 & 7 & 3 & \\
Acid (pH 2.5) & 22 & 2.5 & 3 & \\
Neutral (pH 7) & 22 & 7 & 3 & \\
Alkaline (pH 8.5) & 22 & 8.5 & 3 & \\
\hyperref[s]{S} 0\% & 22 & 7 & 0 & \\
\hyperref[s]{S} 1\% & 22 & 7 & 1 & \\
\hyperref[s]{S} 3\% & 22 & 7 & 3 & \\
\hyperref[s]{S} 6\% & 22 & 7 & 6 & \\
\end{longtable}
}

\subsection*{Evidence and Reproducibility
Checklist}\label{evidence-and-reproducibility-checklist-6}
\addcontentsline{toc}{subsection}{Evidence and Reproducibility
Checklist}

\begin{itemize}
\tightlist
\item
  \textbf{Primary evidence goal:} Measure reaction-rate proxies and
  compare inhibition scenarios.
\item
  \textbf{Data skill to practice:} Fit or interpret enzyme-rate data and
  identify which parameter changed.
\item
  \textbf{BioSkills emphasis:} Quantitative reasoning, Interdisciplinary
  nature of science.
\item
  \textbf{Control logic:} identify at least one positive control, one
  negative control, or one baseline comparison before interpreting
  results.
\item
  \textbf{Measurement discipline:} record units, uncertainty, sample
  size, and any discarded observation with a reason.
\item
  \textbf{Mechanistic link:} connect one result directly to the parent
  chapter's big idea before writing the conclusion.
\item
  \textbf{Reproducibility check:} state one procedural detail that
  another group would need in order to reproduce the result.
\end{itemize}

\subsection*{Paper-Based Evidence
Upgrade}\label{paper-based-evidence-upgrade-6}
\addcontentsline{toc}{subsection}{Paper-Based Evidence Upgrade}

Before answering the analysis questions, annotate the paper dataset for
\textbf{Enzymes and the Kinetics of Catalysis} with a reproducibility
pass:

{\def\LTcaptype{none} % do not increment counter
\begin{longtable}[]{@{}
  >{\raggedright\arraybackslash}p{(\linewidth - 2\tabcolsep) * \real{0.5000}}
  >{\raggedright\arraybackslash}p{(\linewidth - 2\tabcolsep) * \real{0.5000}}@{}}
\toprule\noalign{}
\begin{minipage}[b]{\linewidth}\raggedright
Evidence check
\end{minipage} & \begin{minipage}[b]{\linewidth}\raggedright
Student action
\end{minipage} \\
\midrule\noalign{}
\endhead
\bottomrule\noalign{}
\endlastfoot
Control logic & Mark the comparison that functions as the baseline,
negative control, or reference case. \\
Uncertainty & Circle the row, card, diagram feature, or model assumption
most likely to change the conclusion. \\
Model comparison & State whether a simpler rule, null model, or
alternative mechanism could explain the same pattern. \\
Decision threshold & Write the minimum evidence that would make you
revise the interpretation. \\
Reproducibility & Record the exact scoring rule another group would need
to reproduce your classification. \\
\end{longtable}
}

Focus note: treat the activity-versus-temperature series as a
dose-response curve whose optimum and denaturation endpoints are the
claims; the denatured-enzyme and no-enzyme conditions are the controls
that decide whether a low signal means lost catalysis or a failed assay.
Keep required work paper-based; any material-handling or equipment
version belongs only in an optional extension.

\subsection*{Analysis Questions}\label{analysis-questions-6}
\addcontentsline{toc}{subsection}{Analysis Questions}

\begin{enumerate}
\def\labelenumi{\arabic{enumi}.}
\tightlist
\item
  Plot the simulated O\texorpdfstring{\textsubscript{2}}{2} proxy vs temperature. Identify the approximate
  optimal temperature. Explain the molecular basis for activity loss at
  0°C and above 50°C.
\item
  Explain why the denatured-enzyme condition showed no activity using
  the concept of protein denaturation and active site geometry.
\item
  Plot the simulated O\texorpdfstring{\textsubscript{2}}{2} proxy vs {[}H\texorpdfstring{\textsubscript{2}}{2}O\texorpdfstring{\textsubscript{2}}{2}{]}. Describe the shape of the
  curve. At what concentration does activity appear to plateau? How does
  this relate to the concept of Vmax and enzyme saturation?
\item
  If you added a competitive inhibitor (e.g., a molecule structurally
  similar to H\texorpdfstring{\textsubscript{2}}{2}O\texorpdfstring{\textsubscript{2}}{2}), predict how the rate vs \hyperref[s]{S} curve would
  change in apparent Km and Vmax relative to uninhibited enzyme.
\item
  Catalase deficiency is a human genetic condition (acatalasia)
  associated with recurrent oral infections. Using your experimental
  results, explain why catalase is biologically important in protecting
  cells from H\texorpdfstring{\textsubscript{2}}{2}O\texorpdfstring{\textsubscript{2}}{2} produced during oxidative metabolism.
\end{enumerate}

\subsection*{Optional Wet-Lab
Extension}\label{optional-wet-lab-extension-1}
\addcontentsline{toc}{subsection}{Optional Wet-Lab Extension}

An instructor may demonstrate catalase activity with approved materials
after a safety briefing. The required lab remains the paper dataset and
the mechanistic interpretation of controls, rate curves, and inhibition.

\subsection*{Safety and Ethics Notes}\label{safety-and-ethics-notes-6}
\addcontentsline{toc}{subsection}{Safety and Ethics Notes}

Paper-based lab --- no hydrogen peroxide, heating, or biological samples
are required. Optional wet demonstrations require instructor supervision
and appropriate eye and skin protection.

\subsection*{Debrief and Reflection}\label{debrief-and-reflection-3}
\addcontentsline{toc}{subsection}{Debrief and Reflection}

After you finish the practical work, spend 5--10 minutes in your small
group comparing results and discussing the following prompts. Each
member should contribute at least one observation before moving to the
next prompt:

\begin{enumerate}
\def\labelenumi{\arabic{enumi}.}
\tightlist
\item
  \textbf{What did your measurements show} --- compare the group's
  results to the textbook's predictions. Where they diverge, suggest at
  least one mechanistic explanation before concluding ``experimental
  error.''
\item
  \textbf{What would change the outcome} --- propose one modification to
  the procedure that would sharpen the measurement or extend the result
  to a new biological context, and predict what you would observe.
\item
  \textbf{One-sentence headline} --- each student composes a single
  sentence summarising the lab's take-home message, suitable for a
  tweet. Compare sentences across groups; good headlines are short,
  quantitative, and mechanistic.
\item
  \textbf{Connection back to the textbook} --- identify one section of
  \cref{sec:unit_I_enzymes_and_kinetics} that your data either confirmed
  or complicated. Cite the specific passage.
\end{enumerate}

\subsection*{Further Reading (Lab)}\label{further-reading-lab-3}
\addcontentsline{toc}{subsection}{Further Reading (Lab)}

\begin{itemize}
\tightlist
\item
  Revisit the parent chapter \cref{sec:unit_I_enzymes_and_kinetics} for
  the theoretical foundations on which this lab is built.
\item
  Look up any \textbf{bolded glossary term} introduced in the textbook
  chapter (each has a \texttt{\#gl:term-slug} link in the text) --- its
  master definition is in \breaktt{manuscript/glossary.md}.
\item
  Explore the appended \texttt{src/} module that implements the
  corresponding quantitative model (when applicable) --- referenced in
  the parent chapter's ``Bridge to Computation'' subsection.
\end{itemize}

\emph{Module footer: parent chapter
\texttt{\textbackslash{}cref\{sec:unit\_I\_enzymes\_and\_kinetics\}};
all numerical quantities in this lab use SI units --- see Appendix D of
the textbook for unit conversions and biological-scale reference
values.}

\newpage

\section*{Lab 5 --- Cell Theory and Cell
Types}\label{lab-5-cell-theory-and-cell-types}
\addcontentsline{toc}{section}{Lab 5 --- Cell Theory and Cell Types}

\label{sec:lab_unit_II_cell_theory}

\subsection*{Learning Objectives}\label{learning-objectives-45}
\addcontentsline{toc}{subsection}{Learning Objectives}

\emph{This activity accompanies \cref{sec:unit_II_cell_theory} of the
textbook --- review that chapter before attempting the exercises below.}

\begin{itemize}
\tightlist
\item
  Interpret micrograph packets of prokaryotic and eukaryotic cells
\item
  Apply cell theory to evaluate cases from printed microscopy evidence
\item
  Measure cell dimensions using scale bars and reference data
\item
  Compare structural features of prokaryotic and plant/animal eukaryotic
  cells
\end{itemize}

\subsubsection*{Alignment and Rubric
Map}\label{alignment-and-rubric-map-7}
\addcontentsline{toc}{subsubsection}{Alignment and Rubric Map}

\begin{itemize}
\tightlist
\item
  \textbf{Outcome 1 (LO1):} Interpret the supplied evidence or model
  output for Cell Theory and Cell Types.
\item
  \textbf{Outcome 2 (LO2):} Identify controls and comparison groups that
  make the claim testable.
\item
  \textbf{Outcome 3 (LO3):} Quantify uncertainty, boundary conditions,
  or alternative explanations before concluding.
\item
  \textbf{Outcome 4 (LO4):} Transfer the mechanism to a new biological
  case or public-facing decision.
\item
  \textbf{Chapter LO coverage:} LO1, LO2, LO3, LO4
\item
  \textbf{Rubric dimensions:} evidence; controls; uncertainty;
  mechanism; transfer. \#\# Pre-Lab (to complete before class)
  \{.unnumbered\}
\end{itemize}

\begin{enumerate}
\def\labelenumi{\arabic{enumi}.}
\tightlist
\item
  Define in one sentence each: \textbf{resolution},
  \textbf{magnification}, \textbf{field of view (FOV)}. Note which
  depends on optics and which on apertures.
\item
  Compute the FOV diameter at 400\texorpdfstring{\ensuremath{\times}}{x} if the FOV at 40\texorpdfstring{\ensuremath{\times}}{x} is 4.5 mm. (Answer:
  4.5 mm \texorpdfstring{\ensuremath{\times}}{x} 40/400 = 0.45 mm = 450 \texorpdfstring{\ensuremath{\mu}}{mu}m.)
\item
  Look up Robert Hooke (1665) and Antoni van Leeuwenhoek (1673). In 2--3
  sentences, explain who saw what first and why both contributed to cell
  theory.
\item
  Predict: at 400\texorpdfstring{\ensuremath{\times}}{x} total magnification, how many \emph{Bacillus
  subtilis} cells (2 \texorpdfstring{\ensuremath{\mu}}{mu}m long) would fit end-to-end across the diameter
  of the FOV? (Answer: 450 \texorpdfstring{\ensuremath{\mu}}{mu}m / 2 \texorpdfstring{\ensuremath{\mu}}{mu}m \texorpdfstring{\ensuremath{\approx}}{~} 225 cells.)
\end{enumerate}

\subsection*{Background}\label{background-1}
\addcontentsline{toc}{subsection}{Background}

Cell theory, formalised by Schleiden, Schwann, and Virchow (1838--1855),
states that: (1) most living things are composed of cells; (2) the cell
is the basic unit of life; and (3) most cells arise from pre-existing
cells. Two fundamental cell types exist: \textbf{prokaryotes} (Bacteria,
Archaea --- no nucleus, typically 0.5--5 \texorpdfstring{\ensuremath{\mu}}{mu}m) and \textbf{eukaryotes}
(Eukarya --- membrane-bounded nucleus, 10--100 \texorpdfstring{\ensuremath{\mu}}{mu}m). Light microscopy can
resolve features above about 200 nm, allowing clear
visualisation of cell walls, chloroplasts, and the nucleus, though the
electron microscope is needed for ribosomes and membranes.

\subsection*{Paper-Based Materials}\label{paper-based-materials-7}
\addcontentsline{toc}{subsection}{Paper-Based Materials}

{\def\LTcaptype{none} % do not increment counter
\begin{longtable}[]{@{}
  >{\raggedright\arraybackslash}p{(\linewidth - 2\tabcolsep) * \real{0.3333}}
  >{\raggedright\arraybackslash}p{(\linewidth - 2\tabcolsep) * \real{0.6667}}@{}}
\toprule\noalign{}
\begin{minipage}[b]{\linewidth}\raggedright
Item
\end{minipage} & \begin{minipage}[b]{\linewidth}\raggedright
Quantity
\end{minipage} \\
\midrule\noalign{}
\endhead
\bottomrule\noalign{}
\endlastfoot
Micrograph packet: \emph{Bacillus}, onion epidermis, leaf mesophyll,
epithelial cell & 1 \\
Scale-bar worksheet & 1 \\
Cell-feature evidence cards & 1 set \\
Historical cell-theory source cards & 1 set \\
Ruler (mm) & 1 \\
Calculator & 1 \\
\end{longtable}
}

\subsection*{Paper-Based
Investigation}\label{paper-based-investigation-7}
\addcontentsline{toc}{subsection}{Paper-Based Investigation}

\begin{enumerate}
\def\labelenumi{\arabic{enumi}.}
\tightlist
\item
  \textbf{Prokaryote micrograph:} Inspect the \emph{Bacillus} image.
  Sketch cells; note shape, arrangement, and size relative to the scale
  bar.
\item
  \textbf{Plant epidermis micrograph:} Inspect the onion image. Identify
  cell wall, nucleus where visible, and vacuole evidence.
\item
  \textbf{Leaf mesophyll micrograph:} Inspect the chloroplast-rich
  plant-cell image. Identify chloroplasts and explain why this cell type
  differs from onion epidermis.
\item
  \textbf{Animal epithelial micrograph:} Inspect the epithelial-cell
  image. Identify nucleus, plasma membrane boundary, and lack of cell
  wall.
\item
  For each case, record estimated cell diameter, presence/absence of
  nucleus, cell wall, chloroplasts, and large vacuole. Flag any feature
  that cannot be resolved from the image.
\end{enumerate}

\subsubsection*{Part D --- Microbial Diversity Image
Packet}\label{part-d-microbial-diversity-image-packet}
\addcontentsline{toc}{subsubsection}{Part D --- Microbial Diversity
Image Packet}

\begin{enumerate}
\def\labelenumi{\arabic{enumi}.}
\setcounter{enumi}{5}
\tightlist
\item
  Use the printed microbial-diversity packet to compare chain-forming
  cocci and rod-shaped cells. Sketch both and identify the evidence for
  prokaryotic organisation.
\item
  Estimate cells per field of view from the annotated image grid, then
  calculate how the estimate would change under a 1:10 dilution. State
  one counting assumption that affects reproducibility.
\end{enumerate}

\subsubsection*{Part E --- Scale Bar
Arithmetic}\label{part-e-scale-bar-arithmetic}
\addcontentsline{toc}{subsubsection}{Part E --- Scale Bar Arithmetic}

\begin{enumerate}
\def\labelenumi{\arabic{enumi}.}
\setcounter{enumi}{7}
\item
  The instructor will distribute a printed micrograph of a eukaryotic
  cell with a labelled scale bar of 10 \texorpdfstring{\ensuremath{\mu}}{mu}m. Measure the longest cell
  diameter with a ruler (mm); convert to true \texorpdfstring{\ensuremath{\mu}}{mu}m using the scale-bar
  ratio. This is the core arithmetic for every published micrograph:
  \emph{true size = (measured distance \texorpdfstring{\ensuremath{\div}}{/} scale bar length) \texorpdfstring{\ensuremath{\times}}{x} scale bar
  value}.
\item
  Using the same formula, estimate the total surface area of a spherical
  epithelial-cell image (diameter \texorpdfstring{\ensuremath{\approx}}{~} 60 \texorpdfstring{\ensuremath{\mu}}{mu}m):
  \(A = 4\pi r^2 = 4\pi(30)^2 \approx 1.13 \times 10^4\) \texorpdfstring{\ensuremath{\mu}}{mu}m\texorpdfstring{\textsuperscript{2}}{2}. Compare
  with a spherical \emph{Bacillus} (diameter 2 \texorpdfstring{\ensuremath{\mu}}{mu}m, \(A \approx 12.6\)
  \texorpdfstring{\ensuremath{\mu}}{mu}m\texorpdfstring{\textsuperscript{2}}{2}). Note the about 1000\texorpdfstring{\ensuremath{\times}}{x} difference --- which cell has a
  larger surface-to-volume ratio, and why is that biologically
  important?
\end{enumerate}

\subsection*{Data Recording}\label{data-recording-4}
\addcontentsline{toc}{subsection}{Data Recording}

{\def\LTcaptype{none} % do not increment counter
\begin{longtable}[]{@{}
  >{\raggedright\arraybackslash}p{(\linewidth - 12\tabcolsep) * \real{0.0964}}
  >{\raggedright\arraybackslash}p{(\linewidth - 12\tabcolsep) * \real{0.2289}}
  >{\raggedright\arraybackslash}p{(\linewidth - 12\tabcolsep) * \real{0.2530}}
  >{\raggedright\arraybackslash}p{(\linewidth - 12\tabcolsep) * \real{0.0843}}
  >{\raggedright\arraybackslash}p{(\linewidth - 12\tabcolsep) * \real{0.1084}}
  >{\raggedright\arraybackslash}p{(\linewidth - 12\tabcolsep) * \real{0.1446}}
  >{\raggedright\arraybackslash}p{(\linewidth - 12\tabcolsep) * \real{0.0843}}@{}}
\toprule\noalign{}
\begin{minipage}[b]{\linewidth}\raggedright
Case
\end{minipage} & \begin{minipage}[b]{\linewidth}\raggedright
Cell type (Pro/Euk)
\end{minipage} & \begin{minipage}[b]{\linewidth}\raggedright
Approx. diameter (\texorpdfstring{\ensuremath{\mu}}{mu}m)
\end{minipage} & \begin{minipage}[b]{\linewidth}\raggedright
Nucleus
\end{minipage} & \begin{minipage}[b]{\linewidth}\raggedright
Cell wall
\end{minipage} & \begin{minipage}[b]{\linewidth}\raggedright
Chloroplasts
\end{minipage} & \begin{minipage}[b]{\linewidth}\raggedright
Vacuole
\end{minipage} \\
\midrule\noalign{}
\endhead
\bottomrule\noalign{}
\endlastfoot
Bacillus & & & & & & \\
Onion & & & & & & \\
Leaf mesophyll & & & & & & \\
Epithelial image & & & & & & \\
\end{longtable}
}

Sketch each cell type in the boxes below:

\texttt{{[}Sketch\ boxes\ —\ label\ most\ visible\ structures{]}}

\subsection*{Evidence and Reproducibility
Checklist}\label{evidence-and-reproducibility-checklist-7}
\addcontentsline{toc}{subsection}{Evidence and Reproducibility
Checklist}

\begin{itemize}
\tightlist
\item
  \textbf{Primary evidence goal:} Compare cell types and quantify
  surface-area limits.
\item
  \textbf{Data skill to practice:} Infer cellular constraints from
  measurements, micrographs, and scale bars.
\item
  \textbf{BioSkills emphasis:} Modeling and simulation, Process of
  science.
\item
  \textbf{Control logic:} identify at least one positive control, one
  negative control, or one baseline comparison before interpreting
  results.
\item
  \textbf{Measurement discipline:} record units, uncertainty, sample
  size, and any discarded observation with a reason.
\item
  \textbf{Mechanistic link:} connect one result directly to the parent
  chapter's big idea before writing the conclusion.
\item
  \textbf{Reproducibility check:} state one procedural detail that
  another group would need in order to reproduce the result.
\end{itemize}

\subsection*{Paper-Based Evidence
Upgrade}\label{paper-based-evidence-upgrade-7}
\addcontentsline{toc}{subsection}{Paper-Based Evidence Upgrade}

Before answering the analysis questions, annotate the paper dataset for
\textbf{Cell Theory and Cell Types} with a reproducibility pass:

{\def\LTcaptype{none} % do not increment counter
\begin{longtable}[]{@{}
  >{\raggedright\arraybackslash}p{(\linewidth - 2\tabcolsep) * \real{0.5000}}
  >{\raggedright\arraybackslash}p{(\linewidth - 2\tabcolsep) * \real{0.5000}}@{}}
\toprule\noalign{}
\begin{minipage}[b]{\linewidth}\raggedright
Evidence check
\end{minipage} & \begin{minipage}[b]{\linewidth}\raggedright
Student action
\end{minipage} \\
\midrule\noalign{}
\endhead
\bottomrule\noalign{}
\endlastfoot
Control logic & Mark the comparison that functions as the baseline,
negative control, or reference case. \\
Uncertainty & Circle the row, card, diagram feature, or model assumption
most likely to change the conclusion. \\
Model comparison & State whether a simpler rule, null model, or
alternative mechanism could explain the same pattern. \\
Decision threshold & Write the minimum evidence that would make you
revise the interpretation. \\
Reproducibility & Record the exact scoring rule another group would need
to reproduce your classification. \\
\end{longtable}
}

Focus note: when you assign a cell to a type, write down the operational
definition you used and whether it is a discrete category or a position
on a continuum; treat a single-cell or spatial atlas, a lineage trace,
or a minimal-cell construct as the evidence that would change your call.
Keep required work paper-based; any material-handling or equipment
version belongs only in an optional extension.

\subsection*{Analysis Questions}\label{analysis-questions-7}
\addcontentsline{toc}{subsection}{Analysis Questions}

\begin{enumerate}
\def\labelenumi{\arabic{enumi}.}
\tightlist
\item
  How does the size of \emph{Bacillus} compare to the onion epidermal
  cell? Express as a ratio. What structural feature allows eukaryotic
  cells to be so much larger (hint: internal compartmentalisation)?
\item
  The onion image packet includes a starch-localisation annotation. What
  molecule was detected? Which organelle does this implicate, and why is
  this unexpected in onion cells that are not green?
\item
  Both onion and leaf mesophyll cells are plant cells, but mesophyll
  cells have abundant chloroplasts and onion epidermal cells typically
  do not. Explain using cell theory and differential gene expression.
\item
  A printed pond-water image shows cells with a nucleus but without a
  cell wall. Is this organism a plant, animal, or fungus? What
  additional image feature or molecular marker would narrow the
  identification?
\item
  Devise a paper-based study using time-series image cards to test
  whether temperature affects the rate of cytoplasmic streaming.
  Identify your independent, dependent, and controlled variables.
\end{enumerate}

\subsection*{Extension Analysis
Questions}\label{extension-analysis-questions-1}
\addcontentsline{toc}{subsection}{Extension Analysis Questions}

\begin{enumerate}
\def\labelenumi{\arabic{enumi}.}
\setcounter{enumi}{5}
\tightlist
\item
  Hooke's 1665 \emph{Micrographia} described ``cells'' in cork --- but
  what he saw were primarily the \textbf{walls} of dead plant cells, not
  living matter. Why does this still count as the founding observation
  of cell theory? How does it illustrate the difference between
  \textbf{observing} a cell and \textbf{understanding} it?
\item
  Mitochondria and chloroplasts have their own DNA, 70S ribosomes (like
  bacteria, not the 80S eukaryotic ribosomes in cytoplasm), and divide
  independently. Connect these facts to the \textbf{endosymbiotic
  theory} (Margulis, 1967) --- why is maternal inheritance of
  mitochondrial DNA in humans a direct prediction of this theory?
\item
  A single mammalian cell contains about 10\texorpdfstring{\textsuperscript{7}}{7} proteins. Given a cell
  volume of 1000 \texorpdfstring{\ensuremath{\mu}}{mu}m\texorpdfstring{\textsuperscript{3}}{3}, compute the approximate bulk protein concentration
  (mass/volume). Discuss why ``macromolecular crowding''
  (about 30\% by volume) changes effective reaction kinetics
  relative to dilute test-tube conditions.
\end{enumerate}

\subsection*{Optional Microscopy
Extension}\label{optional-microscopy-extension}
\addcontentsline{toc}{subsection}{Optional Microscopy Extension}

Prepared slides or classroom microscopy may be used as an optional
extension. Do not collect human cells or prepare wet mounts for the
required activity; the micrograph packet supplies most evidence needed
for the lab.

\subsection*{Safety and Ethics Notes}\label{safety-and-ethics-notes-7}
\addcontentsline{toc}{subsection}{Safety and Ethics Notes}

Paper-based lab --- no swabs, stains, wet mounts, cultures, glass
slides, or microscopes are required. Optional microscopy work requires
instructor supervision and standard glassware precautions.

\subsection*{Debrief and Reflection}\label{debrief-and-reflection-4}
\addcontentsline{toc}{subsection}{Debrief and Reflection}

After you finish the practical work, spend 5--10 minutes in your small
group comparing results and discussing the following prompts. Each
member should contribute at least one observation before moving to the
next prompt:

\begin{enumerate}
\def\labelenumi{\arabic{enumi}.}
\tightlist
\item
  \textbf{What did your measurements show} --- compare the group's
  results to the textbook's predictions. Where they diverge, suggest at
  least one mechanistic explanation before concluding ``experimental
  error.''
\item
  \textbf{What would change the outcome} --- propose one modification to
  the procedure that would sharpen the measurement or extend the result
  to a new biological context, and predict what you would observe.
\item
  \textbf{One-sentence headline} --- each student composes a single
  sentence summarising the lab's take-home message, suitable for a
  tweet. Compare sentences across groups; good headlines are short,
  quantitative, and mechanistic.
\item
  \textbf{Connection back to the textbook} --- identify one section of
  \cref{sec:unit_II_cell_theory} that your data either confirmed or
  complicated. Cite the specific passage.
\end{enumerate}

\subsection*{Further Reading (Lab)}\label{further-reading-lab-4}
\addcontentsline{toc}{subsection}{Further Reading (Lab)}

\begin{itemize}
\tightlist
\item
  Revisit the parent chapter \cref{sec:unit_II_cell_theory} for the
  theoretical foundations on which this lab is built.
\item
  Look up any \textbf{bolded glossary term} introduced in the textbook
  chapter (each has a \texttt{\#gl:term-slug} link in the text) --- its
  master definition is in \breaktt{manuscript/glossary.md}.
\item
  Explore the appended \texttt{src/} module that implements the
  corresponding quantitative model (when applicable) --- referenced in
  the parent chapter's ``Bridge to Computation'' subsection.
\end{itemize}

\emph{Module footer: parent chapter
\texttt{\textbackslash{}cref\{sec:unit\_II\_cell\_theory\}}; all
numerical quantities in this lab use SI units --- see Appendix D of the
textbook for unit conversions and biological-scale reference values.}

\newpage

\section*{Lab 6 --- Cell Structure and
Organelles}\label{lab-6-cell-structure-and-organelles}
\addcontentsline{toc}{section}{Lab 6 --- Cell Structure and Organelles}

\label{sec:lab_unit_II_cell_structure}

\subsection*{Learning Objectives}\label{learning-objectives-46}
\addcontentsline{toc}{subsection}{Learning Objectives}

\emph{This activity accompanies \cref{sec:unit_II_cell_structure} of the
textbook --- review that chapter before attempting the exercises below.}

\begin{itemize}
\tightlist
\item
  Identify major organelles in electron micrograph and model-image
  packets
\item
  Relate organelle structure to function using the endomembrane system
\item
  Estimate the magnification and real size of organelles from scale bars
\item
  Compare organelle complements across cell types (liver, muscle, leaf
  mesophyll)
\end{itemize}

\subsubsection*{Alignment and Rubric
Map}\label{alignment-and-rubric-map-8}
\addcontentsline{toc}{subsubsection}{Alignment and Rubric Map}

\begin{itemize}
\tightlist
\item
  \textbf{Outcome 1 (LO1):} Interpret the supplied evidence or model
  output for Cell Structure and Organelles.
\item
  \textbf{Outcome 2 (LO2):} Identify controls and comparison groups that
  make the claim testable.
\item
  \textbf{Outcome 3 (LO3):} Quantify uncertainty, boundary conditions,
  or alternative explanations before concluding.
\item
  \textbf{Outcome 4 (LO4):} Transfer the mechanism to a new biological
  case or public-facing decision.
\item
  \textbf{Chapter LO coverage:} LO1, LO2, LO3, LO4
\item
  \textbf{Rubric dimensions:} evidence; controls; uncertainty;
  mechanism; transfer. \#\# Pre-Lab (to complete before class)
  \{.unnumbered\}
\end{itemize}

\begin{enumerate}
\def\labelenumi{\arabic{enumi}.}
\tightlist
\item
  For each organelle, write one sentence matching it to its dominant
  function: nucleus, ER (rough and smooth), Golgi, mitochondrion,
  chloroplast, lysosome, peroxisome.
\item
  Define in one line: \textbf{scale bar}, \textbf{magnification},
  \textbf{endomembrane system}. Which three compartments make up the
  endomembrane system?
\item
  Compute total mitochondrial volume in a hepatocyte if each
  mitochondrion is a 1 \texorpdfstring{\ensuremath{\mu}}{mu}m \texorpdfstring{\ensuremath{\times}}{x} 3 \texorpdfstring{\ensuremath{\mu}}{mu}m cylinder and the cell contains 2000 of
  them. Express as a fraction of a cell with total volume 5000 \texorpdfstring{\ensuremath{\mu}}{mu}m\texorpdfstring{\textsuperscript{3}}{3}.
\item
  Predict three cell types where you would expect mitochondrial content
  to approach 40\% of cell volume (answer: cardiac muscle, BAT,
  hepatocytes).
\end{enumerate}

\subsection*{Background}\label{background-2}
\addcontentsline{toc}{subsection}{Background}

Eukaryotic cells are organised into membrane-bounded compartments
(organelles) that divide labour: the nucleus stores and expresses DNA;
the rough ER synthesises membrane and secretory proteins; the Golgi
receives, modifies, and sorts vesicle cargo; mitochondria generate ATP
via oxidative phosphorylation. The type and abundance of organelles
reflects a cell's function --- liver hepatocytes are rich in rough ER
(secretory proteins) and smooth ER (lipid metabolism); muscle cells are
packed with mitochondria; leaf mesophyll cells contain hundreds of
chloroplasts.

\subsection*{Paper-Based Materials}\label{paper-based-materials-8}
\addcontentsline{toc}{subsection}{Paper-Based Materials}

{\def\LTcaptype{none} % do not increment counter
\begin{longtable}[]{@{}
  >{\raggedright\arraybackslash}p{(\linewidth - 2\tabcolsep) * \real{0.3333}}
  >{\raggedright\arraybackslash}p{(\linewidth - 2\tabcolsep) * \real{0.6667}}@{}}
\toprule\noalign{}
\begin{minipage}[b]{\linewidth}\raggedright
Item
\end{minipage} & \begin{minipage}[b]{\linewidth}\raggedright
Quantity
\end{minipage} \\
\midrule\noalign{}
\endhead
\bottomrule\noalign{}
\endlastfoot
Transmission electron micrograph (TEM) print packet: plant cell, animal
cell, secretory cell, mitochondrion close-up, chloroplast close-up & 1
set \\
Ruler (mm) & 1 \\
Scale bar calculator worksheet & 1 \\
Coloured pencils (10 colours) & 1 set \\
Printed light-microscopy image cards: onion root tip and pancreas & 1
set \\
Cell function card deck (instructor-supplied: 20 cards) & 1 per group \\
\end{longtable}
}

\subsection*{Paper-Based
Investigation}\label{paper-based-investigation-8}
\addcontentsline{toc}{subsection}{Paper-Based Investigation}

\begin{enumerate}
\def\labelenumi{\arabic{enumi}.}
\tightlist
\item
  \textbf{TEM analysis:} For each of the 5 TEM prints, identify and
  label (with arrows and a key): nucleus, nuclear envelope, nucleolus,
  rough ER, smooth ER, Golgi apparatus, mitochondria, chloroplast (where
  present), plasma membrane, cell wall (where present), vacuole (where
  present), ribosomes, lysosomes.
\item
  \textbf{Size estimation:} Using the scale bar provided on each TEM,
  measure (mm) and calculate the real size of: (a) the nucleus diameter
  in the animal cell image; (b) one mitochondrion; (c) one chloroplast.
\item
  \textbf{Light microscopy image cards:} Use the onion root tip and
  pancreas images to identify cells in interphase versus mitosis and
  secretory acinar cells with abundant granular cytoplasm.
\item
  \textbf{Card sort:} Deal the function cards (e.g., ``exports
  proteins'', ``detoxifies drugs'', ``contains DNA'') and assign each to
  the correct organelle. Record your assignments.
\end{enumerate}

\subsubsection*{Part D --- Cell-Type Organelle
Profile}\label{part-d-cell-type-organelle-profile}
\addcontentsline{toc}{subsubsection}{Part D --- Cell-Type Organelle
Profile}

\begin{enumerate}
\def\labelenumi{\arabic{enumi}.}
\setcounter{enumi}{4}
\tightlist
\item
  Using the provided TEM images, count or estimate the abundance of each
  major organelle (rough ER, smooth ER, mitochondria, Golgi, vacuole,
  chloroplast, peroxisome) for each of: hepatocyte, cardiac myocyte,
  pancreatic acinar cell, palisade mesophyll cell, skin fibroblast. Fill
  in the table below (use 0/+/++/+++ rather than precise counts).
\end{enumerate}

\subsubsection*{Part E --- Scale-Bar Arithmetic in
Action}\label{part-e-scale-bar-arithmetic-in-action}
\addcontentsline{toc}{subsubsection}{Part E --- Scale-Bar Arithmetic in
Action}

\begin{enumerate}
\def\labelenumi{\arabic{enumi}.}
\setcounter{enumi}{5}
\tightlist
\item
  For TEM image \#3 (chloroplast), the scale bar reads 1 \texorpdfstring{\ensuremath{\mu}}{mu}m and measures
  8 mm on the printout. Measure the thylakoid stack width (disk
  diameter) in mm, then compute its real size. A plausible answer is
  400--500 nm per grana disk.
\end{enumerate}

\subsection*{Data Recording}\label{data-recording-5}
\addcontentsline{toc}{subsection}{Data Recording}

{\def\LTcaptype{none} % do not increment counter
\begin{longtable}[]{@{}
  >{\raggedright\arraybackslash}p{(\linewidth - 6\tabcolsep) * \real{0.1154}}
  >{\raggedright\arraybackslash}p{(\linewidth - 6\tabcolsep) * \real{0.2051}}
  >{\raggedright\arraybackslash}p{(\linewidth - 6\tabcolsep) * \real{0.2949}}
  >{\raggedright\arraybackslash}p{(\linewidth - 6\tabcolsep) * \real{0.3846}}@{}}
\toprule\noalign{}
\begin{minipage}[b]{\linewidth}\raggedright
Organelle
\end{minipage} & \begin{minipage}[b]{\linewidth}\raggedright
TEM image (tick)
\end{minipage} & \begin{minipage}[b]{\linewidth}\raggedright
Function (in own words)
\end{minipage} & \begin{minipage}[b]{\linewidth}\raggedright
Cell types where most abundant
\end{minipage} \\
\midrule\noalign{}
\endhead
\bottomrule\noalign{}
\endlastfoot
Nucleus & & & \\
Rough ER & & & \\
Golgi & & & \\
Mitochondrion & & & \\
Chloroplast & & & \\
Lysosome & & & \\
Vacuole & & & \\
Ribosome & & & \\
\end{longtable}
}

Size calculations:

{\def\LTcaptype{none} % do not increment counter
\begin{longtable}[]{@{}
  >{\raggedright\arraybackslash}p{(\linewidth - 6\tabcolsep) * \real{0.1429}}
  >{\raggedright\arraybackslash}p{(\linewidth - 6\tabcolsep) * \real{0.3333}}
  >{\raggedright\arraybackslash}p{(\linewidth - 6\tabcolsep) * \real{0.3016}}
  >{\raggedright\arraybackslash}p{(\linewidth - 6\tabcolsep) * \real{0.2222}}@{}}
\toprule\noalign{}
\begin{minipage}[b]{\linewidth}\raggedright
Structure
\end{minipage} & \begin{minipage}[b]{\linewidth}\raggedright
Scale bar (mm) = (\texorpdfstring{\ensuremath{\mu}}{mu}m)
\end{minipage} & \begin{minipage}[b]{\linewidth}\raggedright
Measured size (mm)
\end{minipage} & \begin{minipage}[b]{\linewidth}\raggedright
Real size (\texorpdfstring{\ensuremath{\mu}}{mu}m)
\end{minipage} \\
\midrule\noalign{}
\endhead
\bottomrule\noalign{}
\endlastfoot
Nucleus & & & \\
Mitochondrion & & & \\
Chloroplast & & & \\
\end{longtable}
}

\subsection*{Evidence and Reproducibility
Checklist}\label{evidence-and-reproducibility-checklist-8}
\addcontentsline{toc}{subsection}{Evidence and Reproducibility
Checklist}

\begin{itemize}
\tightlist
\item
  \textbf{Primary evidence goal:} Map organelle function with
  microscopy-style observations.
\item
  \textbf{Data skill to practice:} Trace a molecule through compartments
  using evidence from labels or perturbations.
\item
  \textbf{BioSkills emphasis:} Modeling and simulation, Process of
  science.
\item
  \textbf{Control logic:} identify at least one positive control, one
  negative control, or one baseline comparison before interpreting
  results.
\item
  \textbf{Measurement discipline:} record units, uncertainty, sample
  size, and any discarded observation with a reason.
\item
  \textbf{Mechanistic link:} connect one result directly to the parent
  chapter's big idea before writing the conclusion.
\item
  \textbf{Reproducibility check:} state one procedural detail that
  another group would need in order to reproduce the result.
\end{itemize}

\subsection*{Paper-Based Evidence
Upgrade}\label{paper-based-evidence-upgrade-8}
\addcontentsline{toc}{subsection}{Paper-Based Evidence Upgrade}

Before answering the analysis questions, annotate the paper dataset for
\textbf{Cell Structure and Organelles} with a reproducibility pass:

{\def\LTcaptype{none} % do not increment counter
\begin{longtable}[]{@{}
  >{\raggedright\arraybackslash}p{(\linewidth - 2\tabcolsep) * \real{0.5000}}
  >{\raggedright\arraybackslash}p{(\linewidth - 2\tabcolsep) * \real{0.5000}}@{}}
\toprule\noalign{}
\begin{minipage}[b]{\linewidth}\raggedright
Evidence check
\end{minipage} & \begin{minipage}[b]{\linewidth}\raggedright
Student action
\end{minipage} \\
\midrule\noalign{}
\endhead
\bottomrule\noalign{}
\endlastfoot
Control logic & Mark the comparison that functions as the baseline,
negative control, or reference case. \\
Uncertainty & Circle the row, card, diagram feature, or model assumption
most likely to change the conclusion. \\
Model comparison & State whether a simpler rule, null model, or
alternative mechanism could explain the same pattern. \\
Decision threshold & Write the minimum evidence that would make you
revise the interpretation. \\
Reproducibility & Record the exact scoring rule another group would need
to reproduce your classification. \\
\end{longtable}
}

Focus note: when you classify an organelle observation, state the
imaging method and resolution, whether a membrane contact site or
condensate is involved, and the cellular context; treat a cryo-electron
tomogram of the native structure as the evidence that would change your
call. Keep required work paper-based; any material-handling or equipment
version belongs only in an optional extension.

\subsection*{Analysis Questions}\label{analysis-questions-8}
\addcontentsline{toc}{subsection}{Analysis Questions}

\begin{enumerate}
\def\labelenumi{\arabic{enumi}.}
\tightlist
\item
  Describe the path taken by a newly synthesised digestive enzyme (e.g.,
  lipase) from its gene in the nucleus to its secretion from a
  pancreatic acinar cell. Name every organelle involved in sequence.
\item
  A cell type has an unusually large number of mitochondria and
  virtually no rough ER. Propose two cell types that fit this
  description and explain the functional logic.
\item
  Using your size calculations, how many mitochondria could fit
  side-by-side across the diameter of the nucleus in your animal cell
  TEM? Show your arithmetic.
\item
  Chloroplasts and mitochondria are both double-membrane organelles with
  their own DNA. What does this suggest about their evolutionary origin?
  Name and explain the theory.
\item
  A student examines a mystery cell and finds: large central vacuole,
  cell wall, chloroplasts, but no centrioles. Is this a plant, animal,
  or fungal cell? Justify each structural clue.
\end{enumerate}

\subsection*{Extension Analysis
Questions}\label{extension-analysis-questions-2}
\addcontentsline{toc}{subsection}{Extension Analysis Questions}

\begin{enumerate}
\def\labelenumi{\arabic{enumi}.}
\setcounter{enumi}{5}
\tightlist
\item
  Cells in brown adipose tissue (BAT) are dedicated to
  \textbf{non-shivering thermogenesis} --- uncoupling electron transport
  from ATP synthesis via UCP1. Predict which organelle complement would
  be dramatically amplified relative to white adipose tissue. Why?
\item
  Peroxisomes carry out very-long-chain fatty-acid \texorpdfstring{\ensuremath{\beta}}{beta}-oxidation and
  hydrogen-peroxide metabolism. Zellweger syndrome is a lethal
  peroxisomal biogenesis disorder. Given your understanding of the
  endomembrane system, why can't the ER or mitochondrion simply take
  over peroxisomal function?
\item
  A TEM-resolved nuclear envelope shows pores about 120 nm
  wide, with a central transporter. Calculate (approximately) the mass
  of mRNA (about 500 kDa for an average mRNP) that can transit
  through one pore per second, given that active transport is limited to
  about 1000 events/pore/s. Compare to the about 2000
  pores per mammalian nucleus.
\end{enumerate}

\subsection*{Optional Microscopy
Extension}\label{optional-microscopy-extension-1}
\addcontentsline{toc}{subsection}{Optional Microscopy Extension}

Prepared slides may be observed if microscopes are available, but the
required analysis uses printed image packets and scale bars. Optional
work should not require live cultures or animal tissues.

\subsection*{Safety and Ethics Notes}\label{safety-and-ethics-notes-8}
\addcontentsline{toc}{subsection}{Safety and Ethics Notes}

Paper-based lab --- no live cultures, animal tissues, glass slides, or
microscopes are required. Electron micrograph prints are shared
resources; do not write on them without protective paper underneath.

\subsection*{Debrief and Reflection}\label{debrief-and-reflection-5}
\addcontentsline{toc}{subsection}{Debrief and Reflection}

After you finish the practical work, spend 5--10 minutes in your small
group comparing results and discussing the following prompts. Each
member should contribute at least one observation before moving to the
next prompt:

\begin{enumerate}
\def\labelenumi{\arabic{enumi}.}
\tightlist
\item
  \textbf{What did your measurements show} --- compare the group's
  results to the textbook's predictions. Where they diverge, suggest at
  least one mechanistic explanation before concluding ``experimental
  error.''
\item
  \textbf{What would change the outcome} --- propose one modification to
  the procedure that would sharpen the measurement or extend the result
  to a new biological context, and predict what you would observe.
\item
  \textbf{One-sentence headline} --- each student composes a single
  sentence summarising the lab's take-home message, suitable for a
  tweet. Compare sentences across groups; good headlines are short,
  quantitative, and mechanistic.
\item
  \textbf{Connection back to the textbook} --- identify one section of
  \cref{sec:unit_II_cell_structure} that your data either confirmed or
  complicated. Cite the specific passage.
\end{enumerate}

\subsection*{Further Reading (Lab)}\label{further-reading-lab-5}
\addcontentsline{toc}{subsection}{Further Reading (Lab)}

\begin{itemize}
\tightlist
\item
  Revisit the parent chapter \cref{sec:unit_II_cell_structure} for the
  theoretical foundations on which this lab is built.
\item
  Look up any \textbf{bolded glossary term} introduced in the textbook
  chapter (each has a \texttt{\#gl:term-slug} link in the text) --- its
  master definition is in \breaktt{manuscript/glossary.md}.
\item
  Explore the appended \texttt{src/} module that implements the
  corresponding quantitative model (when applicable) --- referenced in
  the parent chapter's ``Bridge to Computation'' subsection.
\end{itemize}

\emph{Module footer: parent chapter
\texttt{\textbackslash{}cref\{sec:unit\_II\_cell\_structure\}}; all
numerical quantities in this lab use SI units --- see Appendix D of the
textbook for unit conversions and biological-scale reference values.}

\newpage

\section*{Lab 7 --- Membrane Structure and
Transport}\label{lab-7-membrane-structure-and-transport}
\addcontentsline{toc}{section}{Lab 7 --- Membrane Structure and
Transport}

\label{sec:lab_unit_II_membrane_transport}

\subsection*{Learning Objectives}\label{learning-objectives-47}
\addcontentsline{toc}{subsection}{Learning Objectives}

\emph{This activity accompanies \cref{sec:unit_II_membrane_transport} of
the textbook --- review that chapter before attempting the exercises
below.}

\begin{itemize}
\tightlist
\item
  Interpret osmosis data and estimate osmotic potential from model
  tissues
\item
  Quantify plasmolysis using annotated image packets across NaCl
  concentrations
\item
  Relate membrane lipid bilayer fluidity to temperature through
  simulated colour-leakage data
\item
  Apply the fluid-mosaic model to explain observed transport phenomena
\end{itemize}

\subsubsection*{Alignment and Rubric
Map}\label{alignment-and-rubric-map-9}
\addcontentsline{toc}{subsubsection}{Alignment and Rubric Map}

\begin{itemize}
\tightlist
\item
  \textbf{Outcome 1 (LO1):} Interpret the supplied evidence or model
  output for Membrane Structure and Transport.
\item
  \textbf{Outcome 2 (LO2):} Identify controls and comparison groups that
  make the claim testable.
\item
  \textbf{Outcome 3 (LO3):} Quantify uncertainty, boundary conditions,
  or alternative explanations before concluding.
\item
  \textbf{Outcome 4 (LO4):} Transfer the mechanism to a new biological
  case or public-facing decision.
\item
  \textbf{Chapter LO coverage:} LO1, LO2, LO3, LO4
\item
  \textbf{Rubric dimensions:} evidence; controls; uncertainty;
  mechanism; transfer. \#\# Background \{.unnumbered\}
\end{itemize}

The plasma membrane is a selectively permeable phospholipid bilayer with
embedded proteins. Water moves by osmosis across membranes from regions
of high water potential (\texorpdfstring{\ensuremath{\Psi}}{Psi}) to low \texorpdfstring{\ensuremath{\Psi}}{Psi} (high solute concentration). In
plant cells, plasmolysis occurs when the cell is placed in hypertonic
solution --- the cytoplasm shrinks away from the rigid cell wall. In
animal cells, a hypertonic solution causes crenation (cell shrinking).
Membrane fluidity depends on phospholipid composition (unsaturated fatty
acids \texorpdfstring{\ensuremath{\to}}{->} more fluid) and temperature (higher temperature \texorpdfstring{\ensuremath{\to}}{->} more fluid).

In this lab you will use printed datasets to investigate how temperature
damages membranes (simulated betacyanin leakage) and how plasmolysis
changes across a NaCl concentration gradient. The required work is
paper-based and emphasises controls, graphing, and model interpretation.

\subsection*{Paper-Based Materials}\label{paper-based-materials-9}
\addcontentsline{toc}{subsection}{Paper-Based Materials}

{\def\LTcaptype{none} % do not increment counter
\begin{longtable}[]{@{}ll@{}}
\toprule\noalign{}
Item & Quantity \\
\midrule\noalign{}
\endhead
\bottomrule\noalign{}
\endlastfoot
Membrane-leakage dataset at 20°C, 40°C, 60°C, 80°C, and 100°C & 1 \\
Annotated onion plasmolysis image packet & 1 \\
NaCl gradient and water-potential cards & 1 set \\
Control and confound worksheet & 1 \\
Calculator and graph paper & 1 each \\
\end{longtable}
}

\subsection*{Paper-Based
Investigation}\label{paper-based-investigation-9}
\addcontentsline{toc}{subsection}{Paper-Based Investigation}

\textbf{Part A --- Temperature and Membrane Integrity}

\begin{enumerate}
\def\labelenumi{\arabic{enumi}.}
\tightlist
\item
  Plot the provided absorbance-at-530-nm dataset for model membrane
  leakage across temperatures.
\item
  Identify the baseline condition, the sharpest change in leakage, and
  one confound that a real wet assay would need to control.
\item
  Use the membrane-composition cards to predict how cholesterol or
  unsaturated fatty acids would shift the curve.
\end{enumerate}

\textbf{Part B --- Osmosis and Plasmolysis in Onion}

\begin{enumerate}
\def\labelenumi{\arabic{enumi}.}
\setcounter{enumi}{3}
\tightlist
\item
  Use the annotated image packet for 0\%, 0.5\%, 1\%, 1.5\%, and 2\%
  NaCl.
\item
  Count the marked cells in each image and classify each as plasmolysed
  or not plasmolysed using the rubric.
\item
  Calculate percent plasmolysis and estimate the concentration where
  50\% plasmolysis occurs.
\end{enumerate}

\subsection*{Data Recording}\label{data-recording-6}
\addcontentsline{toc}{subsection}{Data Recording}

{\def\LTcaptype{none} % do not increment counter
\begin{longtable}[]{@{}
  >{\raggedright\arraybackslash}p{(\linewidth - 4\tabcolsep) * \real{0.2576}}
  >{\raggedright\arraybackslash}p{(\linewidth - 4\tabcolsep) * \real{0.3182}}
  >{\raggedright\arraybackslash}p{(\linewidth - 4\tabcolsep) * \real{0.4242}}@{}}
\toprule\noalign{}
\begin{minipage}[b]{\linewidth}\raggedright
Temperature (°C)
\end{minipage} & \begin{minipage}[b]{\linewidth}\raggedright
Absorbance at 530 nm
\end{minipage} & \begin{minipage}[b]{\linewidth}\raggedright
Qualitative colour of water
\end{minipage} \\
\midrule\noalign{}
\endhead
\bottomrule\noalign{}
\endlastfoot
20 & & \\
40 & & \\
60 & & \\
80 & & \\
100 & & \\
\end{longtable}
}

{\def\LTcaptype{none} % do not increment counter
\begin{longtable}[]{@{}lll@{}}
\toprule\noalign{}
{[}NaCl{]} (\%) & Cells plasmolysed (of 20) & \% Plasmolysis \\
\midrule\noalign{}
\endhead
\bottomrule\noalign{}
\endlastfoot
0 & & \\
0.5 & & \\
1.0 & & \\
1.5 & & \\
2.0 & & \\
\end{longtable}
}

\subsection*{Evidence and Reproducibility
Checklist}\label{evidence-and-reproducibility-checklist-9}
\addcontentsline{toc}{subsection}{Evidence and Reproducibility
Checklist}

\begin{itemize}
\tightlist
\item
  \textbf{Primary evidence goal:} Test osmosis and transport mechanisms
  with simple cell or model-membrane observations.
\item
  \textbf{Data skill to practice:} Interpret transport data from
  gradients, rates, and membrane potentials.
\item
  \textbf{BioSkills emphasis:} Modeling and simulation, Process of
  science.
\item
  \textbf{Control logic:} identify at least one positive control, one
  negative control, or one baseline comparison before interpreting
  results.
\item
  \textbf{Measurement discipline:} record units, uncertainty, sample
  size, and any discarded observation with a reason.
\item
  \textbf{Mechanistic link:} connect one result directly to the parent
  chapter's big idea before writing the conclusion.
\item
  \textbf{Reproducibility check:} state one procedural detail that
  another group would need in order to reproduce the result.
\end{itemize}

\subsection*{Paper-Based Evidence
Upgrade}\label{paper-based-evidence-upgrade-9}
\addcontentsline{toc}{subsection}{Paper-Based Evidence Upgrade}

Before answering the analysis questions, annotate the paper dataset for
\textbf{Membrane Structure and Transport} with a reproducibility pass:

{\def\LTcaptype{none} % do not increment counter
\begin{longtable}[]{@{}
  >{\raggedright\arraybackslash}p{(\linewidth - 2\tabcolsep) * \real{0.5000}}
  >{\raggedright\arraybackslash}p{(\linewidth - 2\tabcolsep) * \real{0.5000}}@{}}
\toprule\noalign{}
\begin{minipage}[b]{\linewidth}\raggedright
Evidence check
\end{minipage} & \begin{minipage}[b]{\linewidth}\raggedright
Student action
\end{minipage} \\
\midrule\noalign{}
\endhead
\bottomrule\noalign{}
\endlastfoot
Control logic & Mark the comparison that functions as the baseline,
negative control, or reference case. \\
Uncertainty & Circle the row, card, diagram feature, or model assumption
most likely to change the conclusion. \\
Model comparison & State whether a simpler rule, null model, or
alternative mechanism could explain the same pattern. \\
Decision threshold & Write the minimum evidence that would make you
revise the interpretation. \\
Reproducibility & Record the exact scoring rule another group would need
to reproduce your classification. \\
\end{longtable}
}

Focus note: when you judge a transport claim, specify the concentration
gradient, the electrochemical driving force, the gating state, and the
energy coupling assumed; treat a cryo-EM conformational state or a
single-channel flux recording as the evidence that would change your
call. Keep required work paper-based; any material-handling or equipment
version belongs only in an optional extension.

\subsection*{Analysis Questions}\label{analysis-questions-9}
\addcontentsline{toc}{subsection}{Analysis Questions}

\begin{enumerate}
\def\labelenumi{\arabic{enumi}.}
\tightlist
\item
  Plot absorbance vs temperature. At approximately what temperature did
  membrane integrity fail dramatically (sharp increase in absorbance)?
  What happened to the phospholipid bilayer at this temperature?
\item
  Plot \% plasmolysed cells vs {[}NaCl{]}. At what concentration does
  50\% plasmolysis occur (incipient plasmolysis)? This approximates the
  osmotic potential of the onion cell cytoplasm --- what is your
  estimate in osmol/L?
\item
  In a real beetroot assay, why would rinsing cut discs matter? What
  confound would exist if the dataset mixed surface-cut leakage with
  temperature-induced leakage?
\item
  Predict how results would differ if you used a saturated fat membrane
  (e.g., modelled by adding cholesterol to a liposome) vs a
  polyunsaturated fat membrane. Which would show damage at a lower
  temperature?
\item
  Red blood cells placed in 2\% NaCl would crenate. Compare this to the
  plant cell response in 2\% NaCl. Why do plant cells plasmolysis
  (shrink away from the cell wall) instead of simply shrinking like
  animal cells?
\end{enumerate}

\subsection*{Optional Wet-Lab
Extension}\label{optional-wet-lab-extension-2}
\addcontentsline{toc}{subsection}{Optional Wet-Lab Extension}

An instructor may demonstrate membrane leakage or plasmolysis with
approved materials, but the paper dataset is the required evidence
source. Optional work should be framed as a validation exercise for the
model.

\subsection*{Safety and Ethics Notes}\label{safety-and-ethics-notes-9}
\addcontentsline{toc}{subsection}{Safety and Ethics Notes}

Paper-based lab --- no beetroot, hot water baths, slides, microscopes,
or solutions are required. Optional wet demonstrations require
instructor supervision, stain precautions, and heat-safety procedures.

\subsection*{Debrief and Reflection}\label{debrief-and-reflection-6}
\addcontentsline{toc}{subsection}{Debrief and Reflection}

After you finish the practical work, spend 5--10 minutes in your small
group comparing results and discussing the following prompts. Each
member should contribute at least one observation before moving to the
next prompt:

\begin{enumerate}
\def\labelenumi{\arabic{enumi}.}
\tightlist
\item
  \textbf{What did your measurements show} --- compare the group's
  results to the textbook's predictions. Where they diverge, suggest at
  least one mechanistic explanation before concluding ``experimental
  error.''
\item
  \textbf{What would change the outcome} --- propose one modification to
  the procedure that would sharpen the measurement or extend the result
  to a new biological context, and predict what you would observe.
\item
  \textbf{One-sentence headline} --- each student composes a single
  sentence summarising the lab's take-home message, suitable for a
  tweet. Compare sentences across groups; good headlines are short,
  quantitative, and mechanistic.
\item
  \textbf{Connection back to the textbook} --- identify one section of
  \cref{sec:unit_II_membrane_transport} that your data either confirmed
  or complicated. Cite the specific passage.
\end{enumerate}

\subsection*{Further Reading (Lab)}\label{further-reading-lab-6}
\addcontentsline{toc}{subsection}{Further Reading (Lab)}

\begin{itemize}
\tightlist
\item
  Revisit the parent chapter \cref{sec:unit_II_membrane_transport} for
  the theoretical foundations on which this lab is built.
\item
  Look up any \textbf{bolded glossary term} introduced in the textbook
  chapter (each has a \texttt{\#gl:term-slug} link in the text) --- its
  master definition is in \breaktt{manuscript/glossary.md}.
\item
  Explore the appended \texttt{src/} module that implements the
  corresponding quantitative model (when applicable) --- referenced in
  the parent chapter's ``Bridge to Computation'' subsection.
\end{itemize}

\emph{Module footer: parent chapter
\texttt{\textbackslash{}cref\{sec:unit\_II\_membrane\_transport\}}; all
numerical quantities in this lab use SI units --- see Appendix D of the
textbook for unit conversions and biological-scale reference values.}

\newpage

\section*{Lab 8 --- Cell Signalling and
Communication}\label{lab-8-cell-signalling-and-communication}
\addcontentsline{toc}{section}{Lab 8 --- Cell Signalling and
Communication}

\label{sec:lab_unit_II_cell_signaling}

\subsection*{Learning Objectives}\label{learning-objectives-48}
\addcontentsline{toc}{subsection}{Learning Objectives}

\emph{This activity accompanies \cref{sec:unit_II_cell_signaling} of the
textbook --- review that chapter before attempting the exercises below.}

\begin{itemize}
\tightlist
\item
  Model the three-stage cell signalling cascade (reception \texorpdfstring{\ensuremath{\to}}{->}
  transduction \texorpdfstring{\ensuremath{\to}}{->} response) through a physical simulation
\item
  Analyse how signal amplification occurs through enzymatic cascades
\item
  Evaluate the role of second messengers using a classroom analogy and
  data interpretation
\item
  Design an experiment to test how a signalling inhibitor disrupts a
  cellular response
\end{itemize}

\subsubsection*{Alignment and Rubric
Map}\label{alignment-and-rubric-map-10}
\addcontentsline{toc}{subsubsection}{Alignment and Rubric Map}

\begin{itemize}
\tightlist
\item
  \textbf{Outcome 1 (LO1):} Interpret the supplied evidence or model
  output for Cell Signalling and Communication.
\item
  \textbf{Outcome 2 (LO2):} Identify controls and comparison groups that
  make the claim testable.
\item
  \textbf{Outcome 3 (LO3):} Quantify uncertainty, boundary conditions,
  or alternative explanations before concluding.
\item
  \textbf{Outcome 4 (LO4):} Transfer the mechanism to a new biological
  case or public-facing decision.
\item
  \textbf{Chapter LO coverage:} LO1, LO2, LO3, LO4
\item
  \textbf{Rubric dimensions:} evidence; controls; uncertainty;
  mechanism; transfer. \#\# Background \{.unnumbered\}
\end{itemize}

Cell signalling enables cells to detect and respond to extracellular
information. The three stages are: (1) \textbf{Reception} --- a ligand
binds its receptor (GPCR, RTK, nuclear receptor); (2)
\textbf{Transduction} --- signal is converted and amplified through a
cascade of molecular switches; (3) \textbf{Response} --- altered gene
expression, enzyme activity, or cytoskeletal change. Amplification is a
key feature: one activated receptor can activate hundreds of G proteins,
each activating dozens of adenylate cyclase molecules, each producing
thousands of cAMP molecules --- giving a cascade amplification of
\textgreater10\texorpdfstring{\textsuperscript{6}}{6}.

In this lab you will simulate and calculate signal amplification through
a cascade model, analyse a published cell signalling dose-response
dataset, and design an inhibitor experiment.

\subsection*{Paper-Based Materials}\label{paper-based-materials-10}
\addcontentsline{toc}{subsection}{Paper-Based Materials}

{\def\LTcaptype{none} % do not increment counter
\begin{longtable}[]{@{}
  >{\raggedright\arraybackslash}p{(\linewidth - 2\tabcolsep) * \real{0.3333}}
  >{\raggedright\arraybackslash}p{(\linewidth - 2\tabcolsep) * \real{0.6667}}@{}}
\toprule\noalign{}
\begin{minipage}[b]{\linewidth}\raggedright
Item
\end{minipage} & \begin{minipage}[b]{\linewidth}\raggedright
Quantity
\end{minipage} \\
\midrule\noalign{}
\endhead
\bottomrule\noalign{}
\endlastfoot
Signal cascade worksheet (instructor-provided: table with step values) &
1 \\
Calculator & 1 \\
Graph paper or laptop (for dose-response curve) & 1 \\
Coloured card tokens (representing signalling molecules): 1 ``ligand''
card, 10 ``G protein'' cards, 50 ``cAMP'' cards, 500 ``PKA'' cards & 1
set \\
12-well plate or paper diagram of receptor \texorpdfstring{\ensuremath{\to}}{->} cascade \texorpdfstring{\ensuremath{\to}}{->} response & 1 \\
Sample dataset: \texorpdfstring{\ensuremath{\beta}}{beta}-adrenergic response (heart rate vs {[}adrenaline{]})
(instructor-provided) & 1 \\
\end{longtable}
}

\subsection*{Paper-Based
Investigation}\label{paper-based-investigation-10}
\addcontentsline{toc}{subsection}{Paper-Based Investigation}

\textbf{Part A --- Signal Amplification Simulation}

\begin{enumerate}
\def\labelenumi{\arabic{enumi}.}
\tightlist
\item
  One student is the ``ligand'': activate one ``receptor'' card. Each
  activated receptor activates 10 G protein cards. Each G protein
  activates 10 ``adenylate cyclase'' steps (represented by producing 10
  cAMP tokens). Each cAMP activates 5 PKA tokens.
\item
  Count total active tokens at each level. Record in the amplification
  table.
\item
  Calculate total fold amplification from ligand to PKA response.
\end{enumerate}

\textbf{Part B --- Dose-Response Curve Analysis}

\begin{enumerate}
\def\labelenumi{\arabic{enumi}.}
\setcounter{enumi}{3}
\tightlist
\item
  Using the provided dataset of heart rate (bpm) vs {[}adrenaline{]}
  (nM), plot a dose-response curve on semilog paper (log
  {[}adrenaline{]} on x-axis, heart rate on y).
\item
  Identify the EC\texorpdfstring{\textsubscript{5}}{5}\texorpdfstring{\textsubscript{0}}{0} (concentration producing 50\% of maximal response)
  from your curve.
\item
  A second curve shows the same response in the presence of propranolol
  (a \texorpdfstring{\ensuremath{\beta}}{beta}-blocker). Plot both curves; compare EC\texorpdfstring{\textsubscript{5}}{5}\texorpdfstring{\textsubscript{0}}{0} values and maximum
  responses.
\end{enumerate}

\textbf{Part C --- Inhibitor Experiment Design}

\begin{enumerate}
\def\labelenumi{\arabic{enumi}.}
\setcounter{enumi}{6}
\tightlist
\item
  Design a controlled paper experiment to test whether blocking
  adenylate cyclase with SQ 22536 would prevent an adrenaline-mediated
  response. Use provided pathway cards and a simulated response dataset
  rather than cells or tissue.
\end{enumerate}

\subsection*{Data Recording}\label{data-recording-7}
\addcontentsline{toc}{subsection}{Data Recording}

{\def\LTcaptype{none} % do not increment counter
\begin{longtable}[]{@{}
  >{\raggedright\arraybackslash}p{(\linewidth - 4\tabcolsep) * \real{0.2254}}
  >{\raggedright\arraybackslash}p{(\linewidth - 4\tabcolsep) * \real{0.2676}}
  >{\raggedright\arraybackslash}p{(\linewidth - 4\tabcolsep) * \real{0.5070}}@{}}
\toprule\noalign{}
\begin{minipage}[b]{\linewidth}\raggedright
Signalling level
\end{minipage} & \begin{minipage}[b]{\linewidth}\raggedright
Molecules activated
\end{minipage} & \begin{minipage}[b]{\linewidth}\raggedright
Fold amplification vs previous step
\end{minipage} \\
\midrule\noalign{}
\endhead
\bottomrule\noalign{}
\endlastfoot
Ligand & 1 & --- \\
Receptors activated & & \\
G proteins activated & & \\
cAMP produced & & \\
PKA catalytic subunits activated & & \\
\textbf{Total amplification (ligand \texorpdfstring{\ensuremath{\to}}{->} PKA)} & & \\
\end{longtable}
}

EC\texorpdfstring{\textsubscript{5}}{5}\texorpdfstring{\textsubscript{0}}{0} without propranolol: \_\_\_\_\_\_\_ nM EC\texorpdfstring{\textsubscript{5}}{5}\texorpdfstring{\textsubscript{0}}{0} with propranolol:
\_\_\_\_\_\_\_ nM Shift in EC\texorpdfstring{\textsubscript{5}}{5}\texorpdfstring{\textsubscript{0}}{0}: \_\_\_\_\_\_\_ fold

Your inhibitor experiment design (brief):

\subsection*{Evidence and Reproducibility
Checklist}\label{evidence-and-reproducibility-checklist-10}
\addcontentsline{toc}{subsection}{Evidence and Reproducibility
Checklist}

\begin{itemize}
\tightlist
\item
  \textbf{Primary evidence goal:} Model signaling response from receptor
  activation to cellular output.
\item
  \textbf{Data skill to practice:} Read pathway diagrams and infer the
  effect of agonists, antagonists, or mutations.
\item
  \textbf{BioSkills emphasis:} Modeling and simulation, Process of
  science.
\item
  \textbf{Control logic:} identify at least one positive control, one
  negative control, or one baseline comparison before interpreting
  results.
\item
  \textbf{Measurement discipline:} record units, uncertainty, sample
  size, and any discarded observation with a reason.
\item
  \textbf{Mechanistic link:} connect one result directly to the parent
  chapter's big idea before writing the conclusion.
\item
  \textbf{Reproducibility check:} state one procedural detail that
  another group would need in order to reproduce the result.
\end{itemize}

\subsection*{Paper-Based Evidence
Upgrade}\label{paper-based-evidence-upgrade-10}
\addcontentsline{toc}{subsection}{Paper-Based Evidence Upgrade}

Before answering the analysis questions, annotate the paper dataset for
\textbf{Cell Signalling and Communication} with a reproducibility pass:

{\def\LTcaptype{none} % do not increment counter
\begin{longtable}[]{@{}
  >{\raggedright\arraybackslash}p{(\linewidth - 2\tabcolsep) * \real{0.5000}}
  >{\raggedright\arraybackslash}p{(\linewidth - 2\tabcolsep) * \real{0.5000}}@{}}
\toprule\noalign{}
\begin{minipage}[b]{\linewidth}\raggedright
Evidence check
\end{minipage} & \begin{minipage}[b]{\linewidth}\raggedright
Student action
\end{minipage} \\
\midrule\noalign{}
\endhead
\bottomrule\noalign{}
\endlastfoot
Control logic & Mark the comparison that functions as the baseline,
negative control, or reference case. \\
Uncertainty & Circle the row, card, diagram feature, or model assumption
most likely to change the conclusion. \\
Model comparison & State whether a simpler rule, null model, or
alternative mechanism could explain the same pattern. \\
Decision threshold & Write the minimum evidence that would make you
revise the interpretation. \\
Reproducibility & Record the exact scoring rule another group would need
to reproduce your classification. \\
\end{longtable}
}

Focus note: when you score a signalling case, name the receptor, the
dose or stimulus level, the timing of the readout, and any feedback or
crosstalk you assumed; treat a live-cell biosensor or optogenetic
perturbation trace as the evidence that would change your call. Keep
required work paper-based; any material-handling or equipment version
belongs only in an optional extension.

\subsection*{Analysis Questions}\label{analysis-questions-10}
\addcontentsline{toc}{subsection}{Analysis Questions}

\begin{enumerate}
\def\labelenumi{\arabic{enumi}.}
\tightlist
\item
  Calculate the fold amplification your cascade produced. How does this
  compare to the theoretical amplification cited in the chapter
  (about 10\texorpdfstring{\textsuperscript{6}}{6})? What real-world factors account for the
  difference?
\item
  Propranolol is a competitive antagonist of the \texorpdfstring{\ensuremath{\beta}}{beta}-adrenergic receptor.
  Based on your dose-response curves, does it shift the EC\texorpdfstring{\textsubscript{5}}{5}\texorpdfstring{\textsubscript{0}}{0}, reduce
  Emax, or both? What does this pattern indicate about mechanism?
\item
  Steroid hormones such as oestrogen can enter cells and bind nuclear
  receptors directly. Why is this signalling pathway typically slower
  (minutes to hours) than GPCR cascades (seconds)? What is the cellular
  destination of the hormone-receptor complex?
\item
  Tamoxifen is a selective oestrogen receptor modulator (SERM) that acts
  as an antagonist in breast tissue but a partial agonist in bone.
  Explain how the same molecule can produce opposite effects in
  different tissues.
\item
  A mutation makes the Gs protein GTPase-inactive (cannot hydrolyse GTP
  to GDP). Predict the effect on adenylate cyclase activity, cAMP
  levels, and PKA activity in the affected cell. What disease would this
  phenotype resemble?
\end{enumerate}

\subsection*{Safety and Ethics Notes}\label{safety-and-ethics-notes-10}
\addcontentsline{toc}{subsection}{Safety and Ethics Notes}

No hazardous chemicals in this lab. Card and paper components primarily.
If a computer simulation module is used, follow data handling and
privacy guidelines for any shared datasets.

\subsection*{Debrief and Reflection}\label{debrief-and-reflection-7}
\addcontentsline{toc}{subsection}{Debrief and Reflection}

After you finish the practical work, spend 5--10 minutes in your small
group comparing results and discussing the following prompts. Each
member should contribute at least one observation before moving to the
next prompt:

\begin{enumerate}
\def\labelenumi{\arabic{enumi}.}
\tightlist
\item
  \textbf{What did your measurements show} --- compare the group's
  results to the textbook's predictions. Where they diverge, suggest at
  least one mechanistic explanation before concluding ``experimental
  error.''
\item
  \textbf{What would change the outcome} --- propose one modification to
  the procedure that would sharpen the measurement or extend the result
  to a new biological context, and predict what you would observe.
\item
  \textbf{One-sentence headline} --- each student composes a single
  sentence summarising the lab's take-home message, suitable for a
  tweet. Compare sentences across groups; good headlines are short,
  quantitative, and mechanistic.
\item
  \textbf{Connection back to the textbook} --- identify one section of
  \cref{sec:unit_II_cell_signaling} that your data either confirmed or
  complicated. Cite the specific passage.
\end{enumerate}

\subsection*{Further Reading (Lab)}\label{further-reading-lab-7}
\addcontentsline{toc}{subsection}{Further Reading (Lab)}

\begin{itemize}
\tightlist
\item
  Revisit the parent chapter \cref{sec:unit_II_cell_signaling} for the
  theoretical foundations on which this lab is built.
\item
  Look up any \textbf{bolded glossary term} introduced in the textbook
  chapter (each has a \texttt{\#gl:term-slug} link in the text) --- its
  master definition is in \breaktt{manuscript/glossary.md}.
\item
  Explore the appended \texttt{src/} module that implements the
  corresponding quantitative model (when applicable) --- referenced in
  the parent chapter's ``Bridge to Computation'' subsection.
\end{itemize}

\emph{Module footer: parent chapter
\texttt{\textbackslash{}cref\{sec:unit\_II\_cell\_signaling\}}; all
numerical quantities in this lab use SI units --- see Appendix D of the
textbook for unit conversions and biological-scale reference values.}

\newpage

\section*{Lab 9 --- Bioenergetics and Cellular
Respiration}\label{lab-9-bioenergetics-and-cellular-respiration}
\addcontentsline{toc}{section}{Lab 9 --- Bioenergetics and Cellular
Respiration}

\label{sec:lab_unit_III_bioenergetics_and_respiration}

\subsection*{Learning Objectives}\label{learning-objectives-49}
\addcontentsline{toc}{subsection}{Learning Objectives}

\emph{This activity accompanies
\cref{sec:unit_III_bioenergetics_and_respiration} of the textbook ---
review that chapter before attempting the exercises below.}

\begin{itemize}
\tightlist
\item
  Interpret CO\texorpdfstring{\textsubscript{2}}{2} and O\texorpdfstring{\textsubscript{2}}{2} datasets as indicators of cellular respiration
  rate
\item
  Compare aerobic respiration, fermentation, and the effect of substrate
  type on respiration rate
\item
  Calculate respiration rate (mL CO\texorpdfstring{\textsubscript{2}}{2} per minute per gram of tissue)
\item
  Distinguish the contributions of glycolysis, pyruvate oxidation, and
  the Krebs cycle to overall ATP yield
\end{itemize}

\subsubsection*{Alignment and Rubric
Map}\label{alignment-and-rubric-map-11}
\addcontentsline{toc}{subsubsection}{Alignment and Rubric Map}

\begin{itemize}
\tightlist
\item
  \textbf{Outcome 1 (LO1):} Interpret the supplied evidence or model
  output for Bioenergetics and Cellular Respiration.
\item
  \textbf{Outcome 2 (LO2):} Identify controls and comparison groups that
  make the claim testable.
\item
  \textbf{Outcome 3 (LO3):} Quantify uncertainty, boundary conditions,
  or alternative explanations before concluding.
\item
  \textbf{Outcome 4 (LO4):} Transfer the mechanism to a new biological
  case or public-facing decision.
\item
  \textbf{Chapter LO coverage:} LO1, LO2, LO3, LO4
\item
  \textbf{Rubric dimensions:} evidence; controls; uncertainty;
  mechanism; transfer. \#\# Background \{.unnumbered\}
\end{itemize}

Cellular respiration converts the chemical energy in organic molecules
to ATP. The complete aerobic pathway --- glycolysis \texorpdfstring{\ensuremath{\to}}{->} pyruvate oxidation
\texorpdfstring{\ensuremath{\to}}{->} Krebs cycle \texorpdfstring{\ensuremath{\to}}{->} oxidative phosphorylation --- yields a theoretical
maximum of about 30 ATP per glucose. The overall equation is:
\textbf{C\texorpdfstring{\textsubscript{6}}{6}H\texorpdfstring{\textsubscript{1}}{1}\texorpdfstring{\textsubscript{2}}{2}O\texorpdfstring{\textsubscript{6}}{6} + 6 O\texorpdfstring{\textsubscript{2}}{2} \texorpdfstring{\ensuremath{\to}}{->} 6 CO\texorpdfstring{\textsubscript{2}}{2} + 6 H\texorpdfstring{\textsubscript{2}}{2}O + about 30 ATP}. CO\texorpdfstring{\textsubscript{2}}{2}
production rate and O\texorpdfstring{\textsubscript{2}}{2} consumption rate can serve as proxies for
metabolic rate. In this lab, printed datasets provide model evidence for
comparing aerobic respiration, fermentation, and substrate effects.

\subsection*{Paper-Based Materials}\label{paper-based-materials-11}
\addcontentsline{toc}{subsection}{Paper-Based Materials}

{\def\LTcaptype{none} % do not increment counter
\begin{longtable}[]{@{}
  >{\raggedright\arraybackslash}p{(\linewidth - 2\tabcolsep) * \real{0.3333}}
  >{\raggedright\arraybackslash}p{(\linewidth - 2\tabcolsep) * \real{0.6667}}@{}}
\toprule\noalign{}
\begin{minipage}[b]{\linewidth}\raggedright
Item
\end{minipage} & \begin{minipage}[b]{\linewidth}\raggedright
Quantity
\end{minipage} \\
\midrule\noalign{}
\endhead
\bottomrule\noalign{}
\endlastfoot
Fermentation CO\texorpdfstring{\textsubscript{2}}{2} dataset by substrate type & 1 \\
Respirometry O\texorpdfstring{\textsubscript{2}}{2}-consumption dataset for active, dormant, and
volume-control cases & 1 \\
Pathway carbon-and-electron tracing cards & 1 set \\
Calibration and control worksheet & 1 \\
Calculator and graph paper & 1 each \\
\end{longtable}
}

\subsection*{Paper-Based
Investigation}\label{paper-based-investigation-11}
\addcontentsline{toc}{subsection}{Paper-Based Investigation}

\textbf{Part A --- Yeast Fermentation by Substrate Type}

\begin{enumerate}
\def\labelenumi{\arabic{enumi}.}
\tightlist
\item
  Plot the provided CO\texorpdfstring{\textsubscript{2}}{2} proxy data for four model conditions: sucrose,
  glucose, starch, and no-enzyme negative control.
\item
  Calculate the rate of CO\texorpdfstring{\textsubscript{2}}{2} increase for each condition between 0 and 40
  minutes.
\item
  Identify the negative control, the fastest substrate condition, and
  one alternative explanation that the design rules out.
\end{enumerate}

\textbf{Part B --- Active vs Dormant Tissue Model}

\begin{enumerate}
\def\labelenumi{\arabic{enumi}.}
\setcounter{enumi}{3}
\tightlist
\item
  Use the printed O\texorpdfstring{\textsubscript{2}}{2}-consumption dataset for active seed tissue, dormant
  seed tissue, and an inert volume control.
\item
  Correct the active and dormant readings by subtracting the
  volume-control displacement.
\item
  Convert the corrected values to mL O\texorpdfstring{\textsubscript{2}}{2} per minute per gram and connect
  the result to oxidative phosphorylation.
\end{enumerate}

\subsection*{Data Recording}\label{data-recording-8}
\addcontentsline{toc}{subsection}{Data Recording}

{\def\LTcaptype{none} % do not increment counter
\begin{longtable}[]{@{}lllll@{}}
\toprule\noalign{}
Time (min) & Bottle A & Bottle B & Bottle C & Bottle D \\
\midrule\noalign{}
\endhead
\bottomrule\noalign{}
\endlastfoot
0 & & & & \\
10 & & & & \\
20 & & & & \\
30 & & & & \\
40 & & & & \\
\end{longtable}
}

O\texorpdfstring{\textsubscript{2}}{2} consumed (germinating peas): \_\_\_\_\_\_ mL / 10 min O\texorpdfstring{\textsubscript{2}}{2} consumed
(dry peas): \_\_\_\_\_\_ mL / 10 min O\texorpdfstring{\textsubscript{2}}{2} consumed (glass beads):
\_\_\_\_\_\_ mL / 10 min

\subsection*{Evidence and Reproducibility
Checklist}\label{evidence-and-reproducibility-checklist-11}
\addcontentsline{toc}{subsection}{Evidence and Reproducibility
Checklist}

\begin{itemize}
\tightlist
\item
  \textbf{Primary evidence goal:} Estimate energy yield under aerobic
  and anaerobic conditions.
\item
  \textbf{Data skill to practice:} Track carbon, electrons, and ATP
  across a pathway.
\item
  \textbf{BioSkills emphasis:} Quantitative reasoning, Modeling and
  simulation.
\item
  \textbf{Control logic:} identify at least one positive control, one
  negative control, or one baseline comparison before interpreting
  results.
\item
  \textbf{Measurement discipline:} record units, uncertainty, sample
  size, and any discarded observation with a reason.
\item
  \textbf{Mechanistic link:} connect one result directly to the parent
  chapter's big idea before writing the conclusion.
\item
  \textbf{Reproducibility check:} state one procedural detail that
  another group would need in order to reproduce the result.
\end{itemize}

\subsection*{Paper-Based Evidence
Upgrade}\label{paper-based-evidence-upgrade-11}
\addcontentsline{toc}{subsection}{Paper-Based Evidence Upgrade}

Before answering the analysis questions, annotate the paper dataset for
\textbf{Bioenergetics and Cellular Respiration} with a reproducibility
pass:

{\def\LTcaptype{none} % do not increment counter
\begin{longtable}[]{@{}
  >{\raggedright\arraybackslash}p{(\linewidth - 2\tabcolsep) * \real{0.5000}}
  >{\raggedright\arraybackslash}p{(\linewidth - 2\tabcolsep) * \real{0.5000}}@{}}
\toprule\noalign{}
\begin{minipage}[b]{\linewidth}\raggedright
Evidence check
\end{minipage} & \begin{minipage}[b]{\linewidth}\raggedright
Student action
\end{minipage} \\
\midrule\noalign{}
\endhead
\bottomrule\noalign{}
\endlastfoot
Control logic & Mark the comparison that functions as the baseline,
negative control, or reference case. \\
Uncertainty & Circle the row, card, diagram feature, or model assumption
most likely to change the conclusion. \\
Model comparison & State whether a simpler rule, null model, or
alternative mechanism could explain the same pattern. \\
Decision threshold & Write the minimum evidence that would make you
revise the interpretation. \\
Reproducibility & Record the exact scoring rule another group would need
to reproduce your classification. \\
\end{longtable}
}

Focus note: in a respiration dataset, follow the electrons and protons
explicitly --- ask which step the measured rate (O\texorpdfstring{\textsubscript{2}}{2} uptake, ATP, \texorpdfstring{\ensuremath{\Delta}}{Delta}p)
actually constrains, and whether the reported P/O or ATP yield holds
once shuttle use and proton leak are stated. Keep required work
paper-based; any material-handling or equipment version belongs only in
an optional extension.

\subsection*{Analysis Questions}\label{analysis-questions-11}
\addcontentsline{toc}{subsection}{Analysis Questions}

\begin{enumerate}
\def\labelenumi{\arabic{enumi}.}
\tightlist
\item
  Which substrate model produced the fastest CO\texorpdfstring{\textsubscript{2}}{2} increase --- sucrose or
  glucose? Explain using the concept of hydrolysis and substrate entry
  into glycolysis.
\item
  Why did the starch treatment model produce less CO\texorpdfstring{\textsubscript{2}}{2} than sucrose? How
  does this relate to enzyme availability for substrate breakdown?
\item
  The active-tissue model showed higher O\texorpdfstring{\textsubscript{2}}{2} consumption than the
  dormant-tissue model. Explain at which stage of respiration this O\texorpdfstring{\textsubscript{2}}{2} is
  consumed and write the half-reaction involved.
\item
  Facultative anaerobe models can continue producing CO\texorpdfstring{\textsubscript{2}}{2} with little O\texorpdfstring{\textsubscript{2}}{2}
  available. Predict whether CO\texorpdfstring{\textsubscript{2}}{2} production would stop, slow, or
  continue in a no-O\texorpdfstring{\textsubscript{2}}{2} scenario. Explain using fermentation pathways and
  the role of NAD\texorpdfstring{\textsuperscript{+}}{+} regeneration.
\item
  A marathon runner transitions from aerobic to anaerobic metabolism
  past about 80\% VO\texorpdfstring{\textsubscript{2}}{2}max. Explain what metabolic shift occurs
  in muscle cells, what molecule accumulates (not lactic acid --- what
  is the actual molecule?), and why fatigue results.
\end{enumerate}

\subsection*{Optional Demonstration
Extension}\label{optional-demonstration-extension}
\addcontentsline{toc}{subsection}{Optional Demonstration Extension}

An instructor may demonstrate fermentation or respirometry with approved
materials, but the required lab uses printed datasets. Treat any
demonstration as a calibration or validation exercise rather than the
source of required evidence.

\subsection*{Safety and Ethics Notes}\label{safety-and-ethics-notes-11}
\addcontentsline{toc}{subsection}{Safety and Ethics Notes}

Paper-based lab --- no live cultures, balloons, indicators, seeds, or
sealed respirometers are required. Optional demonstrations require
instructor supervision and standard classroom safety procedures.

\subsection*{Debrief and Reflection}\label{debrief-and-reflection-8}
\addcontentsline{toc}{subsection}{Debrief and Reflection}

After you finish the practical work, spend 5--10 minutes in your small
group comparing results and discussing the following prompts. Each
member should contribute at least one observation before moving to the
next prompt:

\begin{enumerate}
\def\labelenumi{\arabic{enumi}.}
\tightlist
\item
  \textbf{What did your measurements show} --- compare the group's
  results to the textbook's predictions. Where they diverge, suggest at
  least one mechanistic explanation before concluding ``experimental
  error.''
\item
  \textbf{What would change the outcome} --- propose one modification to
  the procedure that would sharpen the measurement or extend the result
  to a new biological context, and predict what you would observe.
\item
  \textbf{One-sentence headline} --- each student composes a single
  sentence summarising the lab's take-home message, suitable for a
  tweet. Compare sentences across groups; good headlines are short,
  quantitative, and mechanistic.
\item
  \textbf{Connection back to the textbook} --- identify one section of
  \cref{sec:unit_III_bioenergetics_and_respiration} that your data
  either confirmed or complicated. Cite the specific passage.
\end{enumerate}

\subsection*{Further Reading (Lab)}\label{further-reading-lab-8}
\addcontentsline{toc}{subsection}{Further Reading (Lab)}

\begin{itemize}
\tightlist
\item
  Revisit the parent chapter
  \cref{sec:unit_III_bioenergetics_and_respiration} for the theoretical
  foundations on which this lab is built.
\item
  Look up any \textbf{bolded glossary term} introduced in the textbook
  chapter (each has a \texttt{\#gl:term-slug} link in the text) --- its
  master definition is in \breaktt{manuscript/glossary.md}.
\item
  Explore the appended \texttt{src/} module that implements the
  corresponding quantitative model (when applicable) --- referenced in
  the parent chapter's ``Bridge to Computation'' subsection.
\end{itemize}

\emph{Module footer: parent chapter
\texttt{\textbackslash{}cref\{sec:unit\_III\_bioenergetics\_and\_respiration\}};
all numerical quantities in this lab use SI units --- see Appendix D of
the textbook for unit conversions and biological-scale reference
values.}

\newpage

\section*{Lab 10 --- Photosynthesis}\label{lab-10-photosynthesis}
\addcontentsline{toc}{section}{Lab 10 --- Photosynthesis}

\label{sec:lab_unit_III_photosynthesis}

\emph{This activity accompanies \cref{sec:unit_III_photosynthesis} of
the textbook --- review that chapter before attempting the exercises
below.}

\subsection*{Learning Objectives}\label{learning-objectives-50}
\addcontentsline{toc}{subsection}{Learning Objectives}

\begin{itemize}
\tightlist
\item
  Design and analyse an experiment to quantify photosynthesis rates
  using a simulated leaf-disc dataset
\item
  Formulate testable hypotheses about how light wavelength and intensity
  affect photosynthesis
\item
  Identify independent, dependent, and controlled variables in
  experimental designs
\item
  Analyze experimental data using basic statistical methods and data
  visualization
\item
  Write a brief scientific report with hypothesis, methods, results, and
  conclusion
\end{itemize}

\subsubsection*{Alignment and Rubric
Map}\label{alignment-and-rubric-map-12}
\addcontentsline{toc}{subsubsection}{Alignment and Rubric Map}

\begin{itemize}
\tightlist
\item
  \textbf{Outcome 1 (LO1):} Interpret the supplied evidence or model
  output for Photosynthesis.
\item
  \textbf{Outcome 2 (LO2):} Identify controls and comparison groups that
  make the claim testable.
\item
  \textbf{Outcome 3 (LO3):} Quantify uncertainty, boundary conditions,
  or alternative explanations before concluding.
\item
  \textbf{Outcome 4 (LO4):} Transfer the mechanism to a new biological
  case or public-facing decision.
\item
  \textbf{Chapter LO coverage:} LO1, LO2, LO3, LO4
\item
  \textbf{Rubric dimensions:} evidence; controls; uncertainty;
  mechanism; transfer. \#\# Pre-Lab Inquiry Questions \{.unnumbered\}
\end{itemize}

\emph{Complete these questions before coming to lab. Use your textbook
and additional research to inform your hypotheses.}

\begin{enumerate}
\def\labelenumi{\arabic{enumi}.}
\item
  \textbf{Hypothesis Formation}: Based on the absorption spectra of
  chlorophyll a and b, predict which color of light (red, blue, green,
  or clear) will result in the fastest photosynthesis rate. Write a
  clear, testable hypothesis statement.
\item
  \textbf{Experimental Design}: What is the independent variable in this
  experiment? What is the dependent variable? List at least 5 variables
  that should be controlled to ensure a fair test.
\item
  \textbf{Scientific Context}: The Ruben \& Kamen (1941) experiment used
  \(^{18}\)O-labeled water to prove that oxygen released during
  photosynthesis comes from water, not carbon dioxide. Why was this
  experiment groundbreaking? How did it change our understanding of
  photosynthesis?
\item
  \textbf{Quantitative Reasoning}: Calculate the energy per photon for
  red light (\(\lambda = 700\) nm) and blue light (\(\lambda = 450\)
  nm). Which has more energy? How might this affect photosynthesis
  efficiency?
\item
  \textbf{Real-World Application}: Commercial greenhouses often use LED
  grow lights. Based on your knowledge of photosynthesis, what light
  spectrum would you recommend for maximizing plant growth? Consider
  both energy efficiency and photosynthetic efficiency.
\end{enumerate}

\subsection*{Background}\label{background-3}
\addcontentsline{toc}{subsection}{Background}

Photosynthesis converts light energy to chemical energy: \textbf{6 CO\texorpdfstring{\textsubscript{2}}{2} +
6 H\texorpdfstring{\textsubscript{2}}{2}O + light \texorpdfstring{\ensuremath{\to}}{->} C\texorpdfstring{\textsubscript{6}}{6}H\texorpdfstring{\textsubscript{1}}{1}\texorpdfstring{\textsubscript{2}}{2}O\texorpdfstring{\textsubscript{6}}{6} + 6 O\texorpdfstring{\textsubscript{2}}{2}}. The light-dependent reactions in the
thylakoid membrane split water (photolysis) and generate NADPH and ATP;
the Calvin cycle in the stroma uses NADPH and ATP to fix CO\texorpdfstring{\textsubscript{2}}{2} into G3P
and ultimately glucose. A floating leaf-disc assay can estimate
photosynthesis because O\texorpdfstring{\textsubscript{2}}{2} accumulation increases buoyancy; here, an
image-and-data packet provides the same ET\texorpdfstring{\textsubscript{5}}{5}\texorpdfstring{\textsubscript{0}}{0} evidence without requiring
plant material, bicarbonate solution, lamps, or cutting tools.

\subsection*{Paper-Based Materials}\label{paper-based-materials-12}
\addcontentsline{toc}{subsection}{Paper-Based Materials}

{\def\LTcaptype{none} % do not increment counter
\begin{longtable}[]{@{}ll@{}}
\toprule\noalign{}
Item & Quantity \\
\midrule\noalign{}
\endhead
\bottomrule\noalign{}
\endlastfoot
Simulated leaf-disc time-course dataset by light colour & 1 \\
Light-intensity and inverse-square worksheet & 1 \\
Chlorophyll absorption spectrum packet & 1 \\
Experimental-design card set: controls, variables, confounds & 1 \\
Calculator or optional Python REPL with this project installed & 1 \\
Spreadsheet software (Excel/Google Sheets) & 1 \\
\end{longtable}
}

\subsection*{Paper-Based
Investigation}\label{paper-based-investigation-12}
\addcontentsline{toc}{subsection}{Paper-Based Investigation}

\subsubsection*{Part 1: Experimental Design and Hypothesis
Testing}\label{part-1-experimental-design-and-hypothesis-testing}
\addcontentsline{toc}{subsubsection}{Part 1: Experimental Design and
Hypothesis Testing}

\begin{enumerate}
\def\labelenumi{\arabic{enumi}.}
\item
  \textbf{Formulate Your Hypothesis}: Based on your pre-lab research,
  write a clear hypothesis about how light color will affect
  photosynthesis rate. Example: ``If photosynthesis is driven by
  chlorophyll absorption, then red light will result in a faster
  photosynthesis rate (lower ET\texorpdfstring{\textsubscript{5}}{5}\texorpdfstring{\textsubscript{0}}{0}) than green light because chlorophyll
  absorbs red wavelengths more efficiently.''
\item
  \textbf{Identify Variables}:

  \begin{itemize}
  \tightlist
  \item
    Independent variable: Light color (red, blue, green, clear)
  \item
    Dependent variable: ET\texorpdfstring{\textsubscript{5}}{5}\texorpdfstring{\textsubscript{0}}{0} (time for 50\% of discs to float)
  \item
    Controlled variables: Temperature, CO\texorpdfstring{\textsubscript{2}}{2} proxy concentration, disc
    size in the model, initial disc status, light-source distance, etc.
  \end{itemize}
\item
  \textbf{Set Up Controls}: Why is it important to have a clear filter
  (white light) control? What other controls could you include?
\item
  \textbf{Audit the Simulated Assay}: Use the method card to check
  whether disc size, starting status, CO\texorpdfstring{\textsubscript{2}}{2} availability, temperature, and
  light distance are controlled across most conditions.
\item
  \textbf{Data Collection}:

  \begin{itemize}
  \tightlist
  \item
    Read the number of floating discs at 1-minute intervals for 20
    minutes from the dataset
  \item
    Use the three provided replicate runs for each condition
  \item
    Calculate ET\texorpdfstring{\textsubscript{5}}{5}\texorpdfstring{\textsubscript{0}}{0} for each replicate using linear interpolation or the
    formula: ET\texorpdfstring{\textsubscript{5}}{5}\texorpdfstring{\textsubscript{0}}{0} = (t\texorpdfstring{\textsubscript{2}}{2} - t\texorpdfstring{\textsubscript{1}}{1}) \texorpdfstring{\ensuremath{\times}}{x} (5 - n\texorpdfstring{\textsubscript{1}}{1}) / (n\texorpdfstring{\textsubscript{2}}{2} - n\texorpdfstring{\textsubscript{1}}{1}) + t\texorpdfstring{\textsubscript{1}}{1}
  \end{itemize}
\end{enumerate}

\subsubsection*{Part 2: Computational Biology Exercise - Data Analysis
with
Python}\label{part-2-computational-biology-exercise---data-analysis-with-python}
\addcontentsline{toc}{subsubsection}{Part 2: Computational Biology
Exercise - Data Analysis with Python}

\emph{Optional computational check: run this self-contained Python
snippet from the project root. It uses tested \texttt{src/biology}
modules and requires no external notebook or CSV file.}

Compare your light-response table with the project photosynthesis model.

\begin{Shaded}
\begin{Highlighting}[]
\ImportTok{from}\NormalTok{ biology.botany }\ImportTok{import}\NormalTok{ photosynthesis\_rate, light\_response\_curve}

\ControlFlowTok{for}\NormalTok{ light }\KeywordTok{in}\NormalTok{ (}\DecValTok{0}\NormalTok{, }\DecValTok{100}\NormalTok{, }\DecValTok{500}\NormalTok{, }\DecValTok{1000}\NormalTok{):}
\NormalTok{    rate }\OperatorTok{=}\NormalTok{ photosynthesis\_rate(light, max\_rate\_µmol\_CO2\_m2\_s}\OperatorTok{=}\FloatTok{22.0}\NormalTok{)}
    \BuiltInTok{print}\NormalTok{(light, }\BuiltInTok{round}\NormalTok{(rate, }\DecValTok{2}\NormalTok{))}

\NormalTok{curve }\OperatorTok{=}\NormalTok{ light\_response\_curve(n\_points}\OperatorTok{=}\DecValTok{5}\NormalTok{)}
\BuiltInTok{print}\NormalTok{(}\StringTok{"first/last model points:"}\NormalTok{, curve[}\DecValTok{0}\NormalTok{], curve[}\OperatorTok{{-}}\DecValTok{1}\NormalTok{])}
\end{Highlighting}
\end{Shaded}

\subsubsection*{Part 3: Light Intensity and Inverse-Square
Law}\label{part-3-light-intensity-and-inverse-square-law}
\addcontentsline{toc}{subsubsection}{Part 3: Light Intensity and
Inverse-Square Law}

\begin{enumerate}
\def\labelenumi{\arabic{enumi}.}
\setcounter{enumi}{5}
\item
  \textbf{Intensity Experiment}: Use the provided white-light dataset at
  three distances: 10 cm, 20 cm, and 40 cm. Record ET\texorpdfstring{\textsubscript{5}}{5}\texorpdfstring{\textsubscript{0}}{0} at each
  distance.
\item
  \textbf{Data Analysis}: Plot ET\texorpdfstring{\textsubscript{5}}{5}\texorpdfstring{\textsubscript{0}}{0} against 1/d\texorpdfstring{\textsuperscript{2}}{2} (relative intensity).
  Is the relationship linear? The linearity should hold primarily up to
  light saturation; beyond that, rate is limited by Calvin-cycle
  throughput, not light.
\item
  \textbf{Calculate Photon Flux}: If a point light source emits 1500
  \texorpdfstring{\ensuremath{\mu}}{mu}mol photons/m\texorpdfstring{\textsuperscript{2}}{2}/s at 20 cm, what is the photon flux at 10 cm and 40
  cm? (Assume inverse-square law applies to point sources.)
\end{enumerate}

\subsubsection*{Part 4: CAM vs C3 Comparison (Extended
Investigation)}\label{part-4-cam-vs-c3-comparison-extended-investigation}
\addcontentsline{toc}{subsubsection}{Part 4: CAM vs C3 Comparison
(Extended Investigation)}

\begin{enumerate}
\def\labelenumi{\arabic{enumi}.}
\setcounter{enumi}{7}
\item
  \textbf{Design an Alternative Experiment}: Design a paper protocol to
  compare photosynthesis rates between a C3 plant model and a CAM plant
  model. Consider how CAM plants fix CO\texorpdfstring{\textsubscript{2}}{2} at night and release it during
  the day. What simulated measurements would you request? How would you
  control variables?
\item
  \textbf{Hypothesis}: Predict what you would observe if you measured O\texorpdfstring{\textsubscript{2}}{2}
  production or CO\texorpdfstring{\textsubscript{2}}{2} uptake in light vs dark for both plant types.
\end{enumerate}

\subsection*{Data Recording}\label{data-recording-9}
\addcontentsline{toc}{subsection}{Data Recording}

{\def\LTcaptype{none} % do not increment counter
\begin{longtable}[]{@{}lllll@{}}
\toprule\noalign{}
Time (min) & Clear (white) & Red filter & Blue filter & Green filter \\
\midrule\noalign{}
\endhead
\bottomrule\noalign{}
\endlastfoot
0 & 0/10 & 0/10 & 0/10 & 0/10 \\
2 & & & & \\
4 & & & & \\
6 & & & & \\
8 & & & & \\
10 & & & & \\
15 & & & & \\
20 & & & & \\
\end{longtable}
}

ET\texorpdfstring{\textsubscript{5}}{5}\texorpdfstring{\textsubscript{0}}{0} (time for 5 discs to float): Clear \_\_ Red \_\_ Blue \_\_ Green
\_\_

Light intensity experiment: |{} Distance (cm) |{} Relative
Intensity (1/d\texorpdfstring{\textsuperscript{2}}{2}) |{} ET\texorpdfstring{\textsubscript{5}}{5}\texorpdfstring{\textsubscript{0}}{0} |{} |{} -------------
|{} ------------------------- |{} ---- |{}
|{} 10 |{} 1.00 |{} |{} |{} 20
|{} 0.25 |{} |{} |{} 40 |{} 0.0625
|{} |{}

\subsection*{Evidence and Reproducibility
Checklist}\label{evidence-and-reproducibility-checklist-12}
\addcontentsline{toc}{subsection}{Evidence and Reproducibility
Checklist}

\begin{itemize}
\tightlist
\item
  \textbf{Primary evidence goal:} Measure or model photosynthetic rate
  under changing light or carbon conditions.
\item
  \textbf{Data skill to practice:} Interpret light-response and
  carbon-fixation data.
\item
  \textbf{BioSkills emphasis:} Quantitative reasoning, Modeling and
  simulation.
\item
  \textbf{Control logic:} identify at least one positive control, one
  negative control, or one baseline comparison before interpreting
  results.
\item
  \textbf{Measurement discipline:} record units, uncertainty, sample
  size, and any discarded observation with a reason.
\item
  \textbf{Mechanistic link:} connect one result directly to the parent
  chapter's big idea before writing the conclusion.
\item
  \textbf{Reproducibility check:} state one procedural detail that
  another group would need in order to reproduce the result.
\end{itemize}

\subsection*{Paper-Based Evidence
Upgrade}\label{paper-based-evidence-upgrade-12}
\addcontentsline{toc}{subsection}{Paper-Based Evidence Upgrade}

Before answering the analysis questions, annotate the paper dataset for
\textbf{Photosynthesis} with a reproducibility pass:

{\def\LTcaptype{none} % do not increment counter
\begin{longtable}[]{@{}
  >{\raggedright\arraybackslash}p{(\linewidth - 2\tabcolsep) * \real{0.5000}}
  >{\raggedright\arraybackslash}p{(\linewidth - 2\tabcolsep) * \real{0.5000}}@{}}
\toprule\noalign{}
\begin{minipage}[b]{\linewidth}\raggedright
Evidence check
\end{minipage} & \begin{minipage}[b]{\linewidth}\raggedright
Student action
\end{minipage} \\
\midrule\noalign{}
\endhead
\bottomrule\noalign{}
\endlastfoot
Control logic & Mark the comparison that functions as the baseline,
negative control, or reference case. \\
Uncertainty & Circle the row, card, diagram feature, or model assumption
most likely to change the conclusion. \\
Model comparison & State whether a simpler rule, null model, or
alternative mechanism could explain the same pattern. \\
Decision threshold & Write the minimum evidence that would make you
revise the interpretation. \\
Reproducibility & Record the exact scoring rule another group would need
to reproduce your classification. \\
\end{longtable}
}

Focus note: in a photosynthesis dataset, name the limiting factor ---
ask whether a rate change reflects light, CO\texorpdfstring{\textsubscript{2}}{2} supply/stomata, or
RuBisCO/photorespiration, and test whether a leaf- or chloroplast-scale
gain would survive whole-plant water and sink limits. Keep required work
paper-based; any material-handling or equipment version belongs only in
an optional extension.

\subsection*{Analysis Questions}\label{analysis-questions-12}
\addcontentsline{toc}{subsection}{Analysis Questions}

\begin{enumerate}
\def\labelenumi{\arabic{enumi}.}
\item
  Which wavelength produced the fastest disc floating (lowest ET\texorpdfstring{\textsubscript{5}}{5}\texorpdfstring{\textsubscript{0}}{0})?
  Which produced the slowest? How does this correlate with the
  absorption spectrum of chlorophyll a and b? Was your hypothesis
  supported?
\item
  Why were discs first made to sink? What gas was removed during vacuum
  infiltration, and why does replacing it with CO\texorpdfstring{\textsubscript{2}}{2}-solution cause
  sinking?
\item
  As light intensity increases, photosynthesis rate increases to a
  plateau (light saturation point). Identify the limiting factor at each
  side of this curve and explain at the biochemical level.
\item
  The Calvin cycle does not directly require light. Why would removing
  light quickly halt the Calvin cycle reactions? (Hint: consider the
  fate of NADPH and ATP between the two stages.)
\item
  A crop scientist wants to maximise yield using artificial grow lights.
  Using your results, recommend the optimal light wavelength(s) and
  explain. What other environmental factors (CO\texorpdfstring{\textsubscript{2}}{2}, temperature, water)
  should she simultaneously optimise, and what law of limiting factors
  applies?
\end{enumerate}

\subsection*{Extension Analysis
Questions}\label{extension-analysis-questions-3}
\addcontentsline{toc}{subsection}{Extension Analysis Questions}

\begin{enumerate}
\def\labelenumi{\arabic{enumi}.}
\setcounter{enumi}{5}
\item
  Rising atmospheric CO\texorpdfstring{\textsubscript{2}}{2} (from about 280 ppm pre-industrial to
  \textgreater420 ppm today) is predicted to increase C3 photosynthesis
  more than C4. Explain why, using your understanding of RuBisCO's
  oxygenation reaction (photorespiration) and the CO\texorpdfstring{\textsubscript{2}}{2}-concentrating
  mechanism in C4 plants.
\item
  Plants grown under monochromatic green light develop abnormally
  (etiolated). Given that green light is weakly absorbed by chlorophyll,
  explain why green leaves appear green and why gardeners with LED grow
  lights generally select red + blue mixtures rather than
  ``full-spectrum'' (which wastes energy in the green).
\item
  The \textbf{Hill reaction} (Robin Hill, 1939) demonstrated that
  isolated chloroplasts can reduce artificial electron acceptors (e.g.,
  DCPIP) in light, without CO\texorpdfstring{\textsubscript{2}}{2} fixation. Describe what this result
  proved about the separability of the light and dark reactions --- and
  why it was historically transformative for photosynthesis research.
\item
  \textbf{Statistical Analysis}: Calculate the standard error of the
  mean (SEM) for each light condition. Construct 95\% confidence
  intervals. Are the differences between light conditions statistically
  significant based on your ANOVA?
\item
  \textbf{Experimental Error Analysis}: What are three potential sources
  of error in this experiment? How could you modify the procedure to
  reduce each source of error?
\end{enumerate}

\subsection*{Group Project Extension
(Multi-Session)}\label{group-project-extension-multi-session}
\addcontentsline{toc}{subsection}{Group Project Extension
(Multi-Session)}

\textbf{Design a Data Study to Test the Effect of Temperature on
Photosynthesis}

This project will span 2-3 lab sessions:

\textbf{Session 1}: Research Question and Hypothesis - How does
temperature affect the rate of photosynthesis? - What enzymes are
involved? What is the optimal temperature range for these enzymes? -
Write a detailed paper protocol to test temperature effects using
provided ET\texorpdfstring{\textsubscript{5}}{5}\texorpdfstring{\textsubscript{0}}{0} datasets or simulation cards

\textbf{Session 2}: Data Collection - Analyse the provided ET\texorpdfstring{\textsubscript{5}}{5}\texorpdfstring{\textsubscript{0}}{0} data at
different temperatures - Include proper controls and replicates

\textbf{Session 3}: Data Analysis and Presentation - Analyze your data
using statistical methods - Create a scientific poster or presentation -
Compare your results to the Q\texorpdfstring{\textsubscript{1}}{1}\texorpdfstring{\textsubscript{0}}{0} rule for temperature effects on
biochemical reactions

\subsection*{Real-World Problem Solving: Climate Change and
Photosynthesis}\label{real-world-problem-solving-climate-change-and-photosynthesis}
\addcontentsline{toc}{subsection}{Real-World Problem Solving: Climate
Change and Photosynthesis}

\textbf{Case Study: Rising CO\texorpdfstring{\textsubscript{2}}{2} and Plant Productivity}

Very high-emissions scenarios can push atmospheric CO\texorpdfstring{\textsubscript{2}}{2} toward roughly
800-1000 ppm by 2100, while lower-emissions pathways remain far below
that range \citep{ipcc2021ar6wg1}. While C3 plants may benefit from CO\texorpdfstring{\textsubscript{2}}{2}
fertilization under controlled conditions, temperature, water
availability, nutrient limitation, ozone exposure, pests, and sink
demand often constrain the effect in real ecosystems and farms.

\begin{enumerate}
\def\labelenumi{\arabic{enumi}.}
\item
  \textbf{Research Task}: Investigate the concept of CO\texorpdfstring{\textsubscript{2}}{2} fertilization.
  Which plant types (C3 vs C4) benefit more? What are the limitations of
  this effect in natural ecosystems?
\item
  \textbf{Ethical Consideration}: Should we rely on enhanced plant
  growth to offset carbon emissions? What are the potential ecological
  consequences of artificially increasing CO\texorpdfstring{\textsubscript{2}}{2} levels in forests or
  agricultural systems?
\item
  \textbf{Policy Proposal}: Design a policy that balances the need for
  carbon sequestration through plant growth with the preservation of
  biodiversity and ecosystem function.
\end{enumerate}

\subsection*{Optional Wet-Lab
Extension}\label{optional-wet-lab-extension-3}
\addcontentsline{toc}{subsection}{Optional Wet-Lab Extension}

An instructor may run a floating leaf-disc demonstration with approved
materials, but the required lab uses the simulated dataset and design
audit. Optional work should be framed as validating the paper model.

\subsection*{Safety and Ethics Notes}\label{safety-and-ethics-notes-12}
\addcontentsline{toc}{subsection}{Safety and Ethics Notes}

Paper-based lab --- no plant material, cutting tools, bicarbonate
solution, cups, hot lamps, or wet materials are required. Optional
demonstrations require instructor supervision, eye-comfort precautions
around bright lights, and ordinary tool safety.

\subsection*{Debrief and Reflection}\label{debrief-and-reflection-9}
\addcontentsline{toc}{subsection}{Debrief and Reflection}

After you finish the practical work, spend 5--10 minutes in your small
group comparing results and discussing the following prompts. Each
member should contribute at least one observation before moving to the
next prompt:

\begin{enumerate}
\def\labelenumi{\arabic{enumi}.}
\item
  \textbf{What did your measurements show} --- compare the group's
  results to the textbook's predictions. Where they diverge, suggest at
  least one mechanistic explanation before concluding ``experimental
  error.''
\item
  \textbf{What would change the outcome} --- propose one modification to
  the procedure that would sharpen the measurement or extend the result
  to a new biological context, and predict what you would observe.
\item
  \textbf{One-sentence headline} --- each student composes a single
  sentence summarising the lab's take-home message, suitable for a
  tweet. Compare sentences across groups; good headlines are short,
  quantitative, and mechanistic.
\item
  \textbf{Connection back to the textbook} --- identify one section of
  \cref{sec:unit_III_photosynthesis} that your data either confirmed or
  complicated. Cite the specific passage.
\item
  \textbf{Experimental Design Reflection}: What was the strongest aspect
  of your experimental design? What would you change if you could repeat
  the experiment?
\end{enumerate}

\subsection*{Further Reading (Lab)}\label{further-reading-lab-9}
\addcontentsline{toc}{subsection}{Further Reading (Lab)}

\begin{itemize}
\tightlist
\item
  Revisit the parent chapter \cref{sec:unit_III_photosynthesis} for the
  theoretical foundations on which this lab is built.
\item
  Look up any \textbf{bolded glossary term} introduced in the textbook
  chapter (each has a \texttt{\#gl:term-slug} link in the text) --- its
  master definition is in \breaktt{manuscript/glossary.md}.
\item
  Explore the Python code in
  \breaktt{src/biology/biochemistry/biochemistry.py} for
  photosynthesis-related calculations.
\item
  Use the self-contained Part 2 snippet as the computational template;
  it runs against tested project modules without external notebooks or
  CSV files.
\end{itemize}

\newpage

\section*{Lab 11 --- Metabolic Integration and
Regulation}\label{lab-11-metabolic-integration-and-regulation}
\addcontentsline{toc}{section}{Lab 11 --- Metabolic Integration and
Regulation}

\label{sec:lab_unit_III_metabolic_integration}

\emph{This activity accompanies
\cref{sec:unit_III_metabolic_integration} of the textbook --- review
that chapter before attempting the exercises below.}

\subsection*{Learning Objectives}\label{learning-objectives-51}
\addcontentsline{toc}{subsection}{Learning Objectives}

\begin{itemize}
\tightlist
\item
  Design and execute a glucose tolerance test (GTT) simulation to
  analyze metabolic regulation
\item
  Formulate testable hypotheses about how insulin and glucagon regulate
  blood glucose
\item
  Identify independent, dependent, and controlled variables in metabolic
  experiments
\item
  Analyze experimental data using area under the curve (AUC) and
  statistical methods
\item
  Write a brief scientific report with hypothesis, methods, results, and
  conclusion
\end{itemize}

\subsubsection*{Alignment and Rubric
Map}\label{alignment-and-rubric-map-13}
\addcontentsline{toc}{subsubsection}{Alignment and Rubric Map}

\begin{itemize}
\tightlist
\item
  \textbf{Outcome 1 (LO1):} Interpret the supplied evidence or model
  output for Metabolic Integration and Regulation.
\item
  \textbf{Outcome 2 (LO2):} Identify controls and comparison groups that
  make the claim testable.
\item
  \textbf{Outcome 3 (LO3):} Quantify uncertainty, boundary conditions,
  or alternative explanations before concluding.
\item
  \textbf{Outcome 4 (LO4):} Transfer the mechanism to a new biological
  case or public-facing decision.
\item
  \textbf{Chapter LO coverage:} LO1, LO2, LO3, LO4
\item
  \textbf{Rubric dimensions:} evidence; controls; uncertainty;
  mechanism; transfer. \#\# Pre-Lab Inquiry Questions \{.unnumbered\}
\end{itemize}

\emph{Complete these questions before coming to lab. Use your textbook
and additional research to inform your hypotheses.}

\begin{enumerate}
\def\labelenumi{\arabic{enumi}.}
\item
  \textbf{Hypothesis Formation}: Based on your understanding of insulin
  and glucagon, predict how a healthy individual's blood glucose would
  respond to a glucose challenge compared to a type 2 diabetic. Write a
  clear, testable hypothesis.
\item
  \textbf{Experimental Design}: What is the independent variable in a
  glucose tolerance test? What is the dependent variable? List at least
  5 variables that should be controlled to ensure a fair test.
\item
  \textbf{Scientific Context}: The glucose tolerance test is a clinical
  diagnostic tool. Why is the area under the curve (AUC) an important
  metric? What does a larger AUC indicate about metabolic health?
\item
  \textbf{Quantitative Reasoning}: If a normal person's blood glucose
  peaks at 30 minutes and returns to baseline by 120 minutes, while a
  diabetic's peaks at 60 minutes and remains elevated at 120 minutes,
  what does this tell you about insulin secretion and sensitivity?
\item
  \textbf{Real-World Application}: Metformin is a first-line drug for
  type 2 diabetes. Based on your knowledge of metabolic regulation, how
  might metformin work? Consider its effects on AMP:ATP ratio and
  downstream targets.
\end{enumerate}

\subsection*{Background}\label{background-4}
\addcontentsline{toc}{subsection}{Background}

Metabolic integration means that pathways (glycolysis, gluconeogenesis,
\texorpdfstring{\ensuremath{\beta}}{beta}-oxidation, glycogen synthesis) are coordinated by hormonal signals and
allosteric regulators, not running independently. Insulin (released
postprandially) activates GLUT4 translocation, glycogen synthase, and
fatty acid synthesis; it inhibits gluconeogenesis and lipolysis.
Glucagon (released during fasting) does the reverse. PFK-1 is the key
allosteric valve of glycolysis: activated by AMP, ADP, and
fructose-2,6-bisphosphate; inhibited by ATP and citrate.

In this data-analysis lab you will interpret simulated glucose tolerance
test data from three subjects (normal, pre-diabetic, type 2 diabetic)
and model the allosteric regulation of PFK-1.

\subsection*{Paper-Based Materials}\label{paper-based-materials-13}
\addcontentsline{toc}{subsection}{Paper-Based Materials}

{\def\LTcaptype{none} % do not increment counter
\begin{longtable}[]{@{}
  >{\raggedright\arraybackslash}p{(\linewidth - 2\tabcolsep) * \real{0.3333}}
  >{\raggedright\arraybackslash}p{(\linewidth - 2\tabcolsep) * \real{0.6667}}@{}}
\toprule\noalign{}
\begin{minipage}[b]{\linewidth}\raggedright
Item
\end{minipage} & \begin{minipage}[b]{\linewidth}\raggedright
Quantity
\end{minipage} \\
\midrule\noalign{}
\endhead
\bottomrule\noalign{}
\endlastfoot
GTT dataset (instructor-provided, printed or digital) & 1 per student \\
Graph paper or laptop with spreadsheet software & 1 \\
Calculator & 1 \\
PFK-1 activity vs {[}fructose-2,6-BP{]} graph (provided) & 1 \\
Coloured pencils (3 colours) & 1 set \\
Calculator or optional Python REPL with this project installed & 1 \\
Ruler & 1 \\
\end{longtable}
}

\subsection*{Paper-Based
Investigation}\label{paper-based-investigation-13}
\addcontentsline{toc}{subsection}{Paper-Based Investigation}

\subsubsection*{Part 1: Glucose Tolerance Test
Analysis}\label{part-1-glucose-tolerance-test-analysis}
\addcontentsline{toc}{subsubsection}{Part 1: Glucose Tolerance Test
Analysis}

\begin{enumerate}
\def\labelenumi{\arabic{enumi}.}
\item
  \textbf{Formulate Your Hypothesis}: Based on your pre-lab research,
  write a hypothesis about how the GTT curves will differ between
  normal, pre-diabetic, and diabetic subjects.
\item
  \textbf{Identify Variables}:

  \begin{itemize}
  \tightlist
  \item
    Independent variable: Subject group (normal, pre-diabetic, type 2
    diabetic)
  \item
    Dependent variable: Blood glucose concentration (mM) over time
  \item
    Controlled variables: Glucose dose, timing of measurements, subject
    age/fasting status, etc.
  \end{itemize}
\item
  \textbf{Set Up Controls}: Why is it important to have a normal control
  group? What other controls could you include in a clinical GTT study?
\item
  \textbf{Data Analysis}:

  \begin{itemize}
  \tightlist
  \item
    Plot blood glucose (mM) vs time (min, 0--120 min) for most three
    subjects on the same graph
  \item
    Annotate key events: glucose ingestion (t = 0), peak blood glucose,
    return to baseline
  \item
    Calculate AUC for each subject using the trapezoid rule at 30-minute
    intervals
  \end{itemize}
\end{enumerate}

\subsubsection*{Part 2: Computational Biology Exercise - Statistical
Analysis with
Python}\label{part-2-computational-biology-exercise---statistical-analysis-with-python}
\addcontentsline{toc}{subsubsection}{Part 2: Computational Biology
Exercise - Statistical Analysis with Python}

\emph{Optional computational check: run this self-contained Python
snippet from the project root. It uses tested \texttt{src/biology}
modules and requires no external notebook or CSV file.}

Use the biochemical models to anchor the class discussion in measured
quantities.

\begin{Shaded}
\begin{Highlighting}[]
\ImportTok{from}\NormalTok{ biology.biochemistry }\ImportTok{import}\NormalTok{ atp\_free\_energy, glycolysis\_summary, michaelis\_menten}

\NormalTok{glycolysis }\OperatorTok{=}\NormalTok{ glycolysis\_summary()}
\NormalTok{rate }\OperatorTok{=}\NormalTok{ michaelis\_menten(substrate\_conc}\OperatorTok{=}\FloatTok{2.0}\NormalTok{, Vmax}\OperatorTok{=}\FloatTok{10.0}\NormalTok{, Km}\OperatorTok{=}\FloatTok{2.0}\NormalTok{)}
\NormalTok{atp }\OperatorTok{=}\NormalTok{ atp\_free\_energy()}

\BuiltInTok{print}\NormalTok{(}\StringTok{"net ATP:"}\NormalTok{, glycolysis.net\_atp)}
\BuiltInTok{print}\NormalTok{(}\StringTok{"half{-}saturation rate:"}\NormalTok{, }\BuiltInTok{round}\NormalTok{(rate.reaction\_rate, }\DecValTok{2}\NormalTok{))}
\BuiltInTok{print}\NormalTok{(}\StringTok{"ATP hydrolysis ΔG:"}\NormalTok{, }\BuiltInTok{round}\NormalTok{(atp, }\DecValTok{2}\NormalTok{))}
\end{Highlighting}
\end{Shaded}

\subsubsection*{Part 3: Allosteric Regulation
Modeling}\label{part-3-allosteric-regulation-modeling}
\addcontentsline{toc}{subsubsection}{Part 3: Allosteric Regulation
Modeling}

\begin{enumerate}
\def\labelenumi{\arabic{enumi}.}
\setcounter{enumi}{4}
\item
  \textbf{Design an Alternative Data Study}: Instead of just analysing
  provided data, design a paper experiment to test how different
  concentrations of fructose-2,6-bisphosphate affect PFK-1 activity.
  What hypothesis would you test? What concentrations would you request
  in the dataset?
\item
  \textbf{Hypothesis}: Predict the shape of the dose-response curve.
  Would it be linear, sigmoidal, or something else? Why?
\end{enumerate}

\subsection*{Data Recording}\label{data-recording-10}
\addcontentsline{toc}{subsection}{Data Recording}

{\def\LTcaptype{none} % do not increment counter
\begin{longtable}[]{@{}
  >{\raggedright\arraybackslash}p{(\linewidth - 6\tabcolsep) * \real{0.1639}}
  >{\raggedright\arraybackslash}p{(\linewidth - 6\tabcolsep) * \real{0.2787}}
  >{\raggedright\arraybackslash}p{(\linewidth - 6\tabcolsep) * \real{0.2787}}
  >{\raggedright\arraybackslash}p{(\linewidth - 6\tabcolsep) * \real{0.2787}}@{}}
\toprule\noalign{}
\begin{minipage}[b]{\linewidth}\raggedright
Time (min)
\end{minipage} & \begin{minipage}[b]{\linewidth}\raggedright
Subject A BG (mM)
\end{minipage} & \begin{minipage}[b]{\linewidth}\raggedright
Subject B BG (mM)
\end{minipage} & \begin{minipage}[b]{\linewidth}\raggedright
Subject C BG (mM)
\end{minipage} \\
\midrule\noalign{}
\endhead
\bottomrule\noalign{}
\endlastfoot
0 & 4.5 & 5.7 & 7.0 \\
30 & 6.8 & 9.2 & 12.4 \\
60 & 5.9 & 8.5 & 13.1 \\
90 & 5.1 & 7.3 & 11.8 \\
120 & 4.6 & 6.4 & 10.5 \\
\end{longtable}
}

AUC Subject A: \_\_\_\_\_ mM\texorpdfstring{\ensuremath{\cdot}}{.}min\\
AUC Subject B: \_\_\_\_\_ mM\texorpdfstring{\ensuremath{\cdot}}{.}min\\
AUC Subject C: \_\_\_\_\_ mM\texorpdfstring{\ensuremath{\cdot}}{.}min

PFK-1 at 0 \texorpdfstring{\ensuremath{\mu}}{mu}M fructose-2,6-BP: \_\_\_\_\_; K\texorpdfstring{\textsubscript{0}}{0}.\texorpdfstring{\textsubscript{5}}{5}: \_\_\_\_\_; at 10 \texorpdfstring{\ensuremath{\mu}}{mu}M:
\_\_\_\_\_

\subsection*{Evidence and Reproducibility
Checklist}\label{evidence-and-reproducibility-checklist-13}
\addcontentsline{toc}{subsection}{Evidence and Reproducibility
Checklist}

\begin{itemize}
\tightlist
\item
  \textbf{Primary evidence goal:} Compare metabolic responses across
  nutrient and hormone conditions.
\item
  \textbf{Data skill to practice:} Use pathway evidence to infer which
  metabolic state or tissue is active.
\item
  \textbf{BioSkills emphasis:} Quantitative reasoning, Modeling and
  simulation.
\item
  \textbf{Control logic:} identify at least one positive control, one
  negative control, or one baseline comparison before interpreting
  results.
\item
  \textbf{Measurement discipline:} record units, uncertainty, sample
  size, and any discarded observation with a reason.
\item
  \textbf{Mechanistic link:} connect one result directly to the parent
  chapter's big idea before writing the conclusion.
\item
  \textbf{Reproducibility check:} state one procedural detail that
  another group would need in order to reproduce the result.
\end{itemize}

\subsection*{Paper-Based Evidence
Upgrade}\label{paper-based-evidence-upgrade-13}
\addcontentsline{toc}{subsection}{Paper-Based Evidence Upgrade}

Before answering the analysis questions, annotate the paper dataset for
\textbf{Metabolic Integration and Regulation} with a reproducibility
pass:

{\def\LTcaptype{none} % do not increment counter
\begin{longtable}[]{@{}
  >{\raggedright\arraybackslash}p{(\linewidth - 2\tabcolsep) * \real{0.5000}}
  >{\raggedright\arraybackslash}p{(\linewidth - 2\tabcolsep) * \real{0.5000}}@{}}
\toprule\noalign{}
\begin{minipage}[b]{\linewidth}\raggedright
Evidence check
\end{minipage} & \begin{minipage}[b]{\linewidth}\raggedright
Student action
\end{minipage} \\
\midrule\noalign{}
\endhead
\bottomrule\noalign{}
\endlastfoot
Control logic & Mark the comparison that functions as the baseline,
negative control, or reference case. \\
Uncertainty & Circle the row, card, diagram feature, or model assumption
most likely to change the conclusion. \\
Model comparison & State whether a simpler rule, null model, or
alternative mechanism could explain the same pattern. \\
Decision threshold & Write the minimum evidence that would make you
revise the interpretation. \\
Reproducibility & Record the exact scoring rule another group would need
to reproduce your classification. \\
\end{longtable}
}

Focus note: in an integrated-metabolism dataset, separate flux from pool
size --- ask whether isotope-tracer or rate evidence supports each
directional claim, and pin every result to a fed/fasted state, tissue,
and time window before generalizing. Keep required work paper-based; any
material-handling or equipment version belongs only in an optional
extension.

\subsection*{Analysis Questions}\label{analysis-questions-13}
\addcontentsline{toc}{subsection}{Analysis Questions}

\begin{enumerate}
\def\labelenumi{\arabic{enumi}.}
\item
  Based on the GTT data, which AUC threshold might a clinician use to
  diagnose type 2 diabetes? How does the shape of Subject C's curve
  differ from Subject A's, and what does this indicate about \texorpdfstring{\ensuremath{\beta}}{beta}-cell
  insulin secretion capacity?
\item
  Subject B is pre-diabetic. Suggest two lifestyle interventions that
  would shift Subject B's GTT curve toward Subject A's profile. For
  each, identify the specific metabolic pathway or hormone axis
  restored.
\item
  In the post-meal state (high insulin), what happens to hepatic
  fructose-2,6-bisphosphate levels? Trace: insulin \texorpdfstring{\ensuremath{\to}}{->} PFK-2
  phosphorylation state \texorpdfstring{\ensuremath{\to}}{->} {[}fructose-2,6-BP{]} \texorpdfstring{\ensuremath{\to}}{->} PFK-1 activity \texorpdfstring{\ensuremath{\to}}{->}
  glycolysis rate.
\item
  During intense exercise, AMP levels spike. Using the allosteric model,
  predict and explain the combined effect of elevated AMP and elevated
  fructose-2,6-BP on glycolytic flux in exercising muscle.
\item
  Metformin, the first-line drug for type 2 diabetes, inhibits Complex I
  of the mitochondrial electron transport chain, increasing the AMP:ATP
  ratio. Using your metabolic knowledge, predict at least three
  downstream effects of elevated AMP: (a) on PFK-1; (b) on glycogen
  synthase; (c) on gluconeogenesis. Does this make metabolic sense as a
  diabetes treatment?
\end{enumerate}

\subsection*{Extension Analysis
Questions}\label{extension-analysis-questions-4}
\addcontentsline{toc}{subsection}{Extension Analysis Questions}

\begin{enumerate}
\def\labelenumi{\arabic{enumi}.}
\setcounter{enumi}{5}
\item
  \textbf{Statistical Analysis}: Calculate the standard error of the
  mean (SEM) for each subject group. Construct 95\% confidence
  intervals. Are the differences between groups statistically
  significant based on your ANOVA?
\item
  \textbf{Experimental Error Analysis}: What are three potential sources
  of error in this data analysis? How could you modify the procedure to
  reduce each source of error?
\item
  \textbf{Experimental Design}: If you were to design a follow-up
  experiment to distinguish between insulin resistance and \texorpdfstring{\ensuremath{\beta}}{beta}-cell
  dysfunction in the pre-diabetic subject, what would you measure? How
  would you interpret the results?
\end{enumerate}

\subsection*{Group Project Extension
(Multi-Session)}\label{group-project-extension-multi-session-1}
\addcontentsline{toc}{subsection}{Group Project Extension
(Multi-Session)}

\textbf{Design a Metabolic Syndrome Risk Assessment Tool}

This project will span 2-3 lab sessions:

\textbf{Session 1}: Research Question and Hypothesis - What factors
contribute to metabolic syndrome (obesity, hypertension, high blood
glucose, dyslipidemia)? - How can we use simple measurements to assess
risk? - Write a detailed protocol for assessing metabolic health in a
human population

\textbf{Session 2}: Data Collection and Analysis - Use de-identified
public datasets or fully simulated classroom data - Calculate risk
scores based on waist circumference, blood pressure, fasting glucose,
triglycerides, HDL - Analyze correlations between different risk factors

\textbf{Session 3}: Tool Development and Presentation - Develop a simple
risk assessment tool (questionnaire or calculator) - Validate it against
known metabolic syndrome criteria - Create a scientific poster or
presentation

\subsection*{Real-World Problem Solving: Metabolic
Diseases}\label{real-world-problem-solving-metabolic-diseases}
\addcontentsline{toc}{subsection}{Real-World Problem Solving: Metabolic
Diseases}

\textbf{Case Study: Type 2 Diabetes Prevention}

Type 2 diabetes is largely preventable through lifestyle modifications,
yet incidence continues to rise globally.

\begin{enumerate}
\def\labelenumi{\arabic{enumi}.}
\item
  \textbf{Research Task}: Investigate the effectiveness of different
  interventions (diet, exercise, medication) in preventing or delaying
  type 2 diabetes in high-risk individuals.
\item
  \textbf{Ethical Consideration}: Should employers or insurance
  companies offer incentives for healthy behaviors? What are the
  potential benefits and risks of such programs?
\item
  \textbf{Policy Proposal}: Design a community-based program to reduce
  the incidence of type 2 diabetes. Consider factors like access to
  healthy food, safe spaces for exercise, and education.
\end{enumerate}

\subsection*{Safety and Ethics Notes}\label{safety-and-ethics-notes-13}
\addcontentsline{toc}{subsection}{Safety and Ethics Notes}

Paper-based data-analysis lab --- no chemical hazards, biological
samples, or volunteer data collection are required. If using real public
patient data, most data must be de-identified and handled with
privacy-preserving practices.

\subsection*{Debrief and Reflection}\label{debrief-and-reflection-10}
\addcontentsline{toc}{subsection}{Debrief and Reflection}

After you finish the practical work, spend 5--10 minutes in your small
group comparing results and discussing the following prompts. Each
member should contribute at least one observation before moving to the
next prompt:

\begin{enumerate}
\def\labelenumi{\arabic{enumi}.}
\item
  \textbf{What did your measurements show} --- compare the group's
  results to the textbook's predictions. Where they diverge, suggest at
  least one mechanistic explanation before concluding ``experimental
  error.''
\item
  \textbf{What would change the outcome} --- propose one modification to
  the procedure that would sharpen the measurement or extend the result
  to a new biological context, and predict what you would observe.
\item
  \textbf{One-sentence headline} --- each student composes a single
  sentence summarising the lab's take-home message, suitable for a
  tweet. Compare sentences across groups; good headlines are short,
  quantitative, and mechanistic.
\item
  \textbf{Connection back to the textbook} --- identify one section of
  \cref{sec:unit_III_metabolic_integration} that your data either
  confirmed or complicated. Cite the specific passage.
\item
  \textbf{Experimental Design Reflection}: What was the strongest aspect
  of your experimental design? What would you change if you could repeat
  the experiment?
\end{enumerate}

\subsection*{Further Reading (Lab)}\label{further-reading-lab-10}
\addcontentsline{toc}{subsection}{Further Reading (Lab)}

\begin{itemize}
\tightlist
\item
  Revisit the parent chapter \cref{sec:unit_III_metabolic_integration}
  for the theoretical foundations on which this lab is built.
\item
  Look up any \textbf{bolded glossary term} introduced in the textbook
  chapter (each has a \texttt{\#gl:term-slug} link in the text) --- its
  master definition is in \breaktt{manuscript/glossary.md}.
\item
  Explore the Python code in
  \breaktt{src/biology/biochemistry/biochemistry.py} for metabolic
  calculations.
\item
  Use the self-contained Part 2 snippet as the computational template;
  it runs against tested project modules without external notebooks or
  CSV files.
\end{itemize}

\newpage

\section*{Lab 12 --- DNA Replication and the Cell
Cycle}\label{lab-12-dna-replication-and-the-cell-cycle}
\addcontentsline{toc}{section}{Lab 12 --- DNA Replication and the Cell
Cycle}

\label{sec:lab_unit_IV_dna_replication_and_cell_cycle}

\emph{This activity accompanies
\cref{sec:unit_IV_dna_replication_and_cell_cycle} of the textbook ---
review that chapter before attempting the exercises below.}

\subsection*{Learning Objectives}\label{learning-objectives-52}
\addcontentsline{toc}{subsection}{Learning Objectives}

\begin{itemize}
\tightlist
\item
  Design and evaluate a mitotic staging dataset to calculate mitotic
  index
\item
  Formulate testable hypotheses about how cell cycle progression varies
  across tissue types
\item
  Identify independent, dependent, and controlled variables in cell
  biology experiments
\item
  Analyze experimental data using statistical methods and cell biology
  concepts
\item
  Write a brief scientific report with hypothesis, methods, results, and
  conclusion
\end{itemize}

\subsubsection*{Alignment and Rubric
Map}\label{alignment-and-rubric-map-14}
\addcontentsline{toc}{subsubsection}{Alignment and Rubric Map}

\begin{itemize}
\tightlist
\item
  \textbf{Outcome 1 (LO1):} Interpret the supplied evidence or model
  output for DNA Replication and the Cell Cycle.
\item
  \textbf{Outcome 2 (LO2):} Identify controls and comparison groups that
  make the claim testable.
\item
  \textbf{Outcome 3 (LO3):} Quantify uncertainty, boundary conditions,
  or alternative explanations before concluding.
\item
  \textbf{Outcome 4 (LO4):} Transfer the mechanism to a new biological
  case or public-facing decision.
\item
  \textbf{Chapter LO coverage:} LO1, LO2, LO3, LO4
\item
  \textbf{Rubric dimensions:} evidence; controls; uncertainty;
  mechanism; transfer. \#\# Pre-Lab Inquiry Questions \{.unnumbered\}
\end{itemize}

\emph{Complete these questions before coming to lab. Use your textbook
and additional research to inform your hypotheses.}

\begin{enumerate}
\def\labelenumi{\arabic{enumi}.}
\item
  \textbf{Hypothesis Formation}: Based on your understanding of the cell
  cycle, predict which stage of mitosis will take the longest time. What
  molecular events occur during this stage that might explain its
  duration?
\item
  \textbf{Experimental Design}: What is the independent variable in a
  mitotic index experiment? What is the dependent variable? List at
  least 5 variables that should be controlled to ensure accurate
  staging.
\item
  \textbf{Scientific Context}: The mitotic index is used in cancer
  diagnosis. Why is a high mitotic index indicative of malignant tumors?
  What are the limitations of this measure?
\item
  \textbf{Quantitative Reasoning}: If a tissue has a mitotic index of
  5\%, approximately what percentage of cells are in S phase? (Assume S
  phase duration is about 6-8 hours in a 24-hour cycle.)
\item
  \textbf{Real-World Application}: Chemotherapy drugs often target
  rapidly dividing cells. How does the cell cycle explain why these
  drugs affect cancer cells more than normal cells? What are the side
  effects related to this mechanism?
\end{enumerate}

\subsection*{Background}\label{background-5}
\addcontentsline{toc}{subsection}{Background}

DNA replication is semiconservative (Meselson--Stahl, 1958): each
daughter double helix retains one original strand and one newly
synthesised strand. Replication proceeds bidirectionally from origins of
replication; DNA polymerase extends DNA in the 5'\texorpdfstring{\ensuremath{\to}}{->}3' direction,
requiring the lagging strand to be synthesised discontinuously as
Okazaki fragments. The cell cycle (G\texorpdfstring{\textsubscript{1}}{1} \texorpdfstring{\ensuremath{\to}}{->} S \texorpdfstring{\ensuremath{\to}}{->} G\texorpdfstring{\textsubscript{2}}{2} \texorpdfstring{\ensuremath{\to}}{->} M) is driven by
cyclin-CDK complexes and monitored by checkpoints (G\texorpdfstring{\textsubscript{1}}{1}/S, G\texorpdfstring{\textsubscript{2}}{2}/M, spindle
assembly checkpoint). The mitotic index (fraction of cells in mitosis)
reflects proliferative activity and is elevated in tumours.

\subsection*{Paper-Based Materials}\label{paper-based-materials-14}
\addcontentsline{toc}{subsection}{Paper-Based Materials}

{\def\LTcaptype{none} % do not increment counter
\begin{longtable}[]{@{}
  >{\raggedright\arraybackslash}p{(\linewidth - 2\tabcolsep) * \real{0.3333}}
  >{\raggedright\arraybackslash}p{(\linewidth - 2\tabcolsep) * \real{0.6667}}@{}}
\toprule\noalign{}
\begin{minipage}[b]{\linewidth}\raggedright
Item
\end{minipage} & \begin{minipage}[b]{\linewidth}\raggedright
Quantity
\end{minipage} \\
\midrule\noalign{}
\endhead
\bottomrule\noalign{}
\endlastfoot
Printed cell-cycle image cards from root tip, embryo, and tumour tissue
datasets & 1 set \\
Mitotic stage decision key with example features & 1 \\
Meselson-Stahl simulation dataset (printed) & 1 \\
Cell cycle timeline diagram (printed) & 1 per student \\
Checkpoint perturbation cards (untreated, DNA polymerase inhibitor,
spindle poison) & 1 set \\
Transparent grid overlay or printed counting frame & 1 \\
Calculator & 1 \\
Calculator or optional Python REPL with this project installed & 1 \\
\end{longtable}
}

\subsection*{Paper-Based
Investigation}\label{paper-based-investigation-14}
\addcontentsline{toc}{subsection}{Paper-Based Investigation}

\subsubsection*{Part 1: Mitotic Index Determination from Image
Cards}\label{part-1-mitotic-index-determination-from-image-cards}
\addcontentsline{toc}{subsubsection}{Part 1: Mitotic Index Determination
from Image Cards}

\begin{enumerate}
\def\labelenumi{\arabic{enumi}.}
\item
  \textbf{Formulate Your Hypothesis}: Based on your pre-lab research,
  write a hypothesis about how the mitotic index might differ among root
  tip, embryo, and tumour image datasets.
\item
  \textbf{Identify Variables}:

  \begin{itemize}
  \tightlist
  \item
    Independent variable: Tissue dataset type
  \item
    Dependent variable: Mitotic index (\% cells in mitosis)
  \item
    Controlled variables: image source, magnification label, counting
    frame size, staging key, scorer training, and exclusion rules.
  \end{itemize}
\item
  \textbf{Set Up Controls}: Use the instructor-provided answer-key
  subset as a positive control for each mitotic stage and a non-dividing
  tissue card as a negative control. Before scoring the unknown cards,
  each group must reach at least 80\% agreement on the control cards.
\item
  \textbf{Apply the Staging Key}:

  \begin{itemize}
  \tightlist
  \item
    Place the transparent grid over each printed image card.
  \item
    Classify exactly 100 consecutive visible cells per card as
    interphase, prophase, metaphase, anaphase, telophase, or
    ``uncertain.''
  \item
    Record uncertain cells separately; do not force a stage call unless
    the chromosome features match the decision key.
  \end{itemize}
\item
  \textbf{Data Collection}:

  \begin{itemize}
  \tightlist
  \item
    Score three cards per tissue dataset, with two students
    independently scoring one shared card to estimate inter-rater
    reliability.
  \item
    Calculate mitotic index = (cells in mitosis \texorpdfstring{\ensuremath{\div}}{/} total cells counted) \texorpdfstring{\ensuremath{\times}}{x}
    100\%.
  \item
    Report the mean, range, and standard error across replicate cards;
    flag any card whose uncertain-cell count exceeds 10\%.
  \end{itemize}
\end{enumerate}

\subsubsection*{Part 2: Computational Biology Exercise - Cell Cycle
Modeling with
Python}\label{part-2-computational-biology-exercise---cell-cycle-modeling-with-python}
\addcontentsline{toc}{subsubsection}{Part 2: Computational Biology
Exercise - Cell Cycle Modeling with Python}

\emph{Optional computational check: run this self-contained Python
snippet from the project root. It uses tested \texttt{src/biology}
modules and requires no external notebook or CSV file.}

Check template-strand logic before interpreting replication and
cell-cycle data.

\begin{Shaded}
\begin{Highlighting}[]
\ImportTok{from}\NormalTok{ biology.genetics }\ImportTok{import}\NormalTok{ cpg\_methylation\_remaining, dna\_complement, transcribe\_dna\_to\_mrna}

\NormalTok{template }\OperatorTok{=} \StringTok{"TACGGA"}
\BuiltInTok{print}\NormalTok{(}\StringTok{"complement:"}\NormalTok{, dna\_complement(template))}
\BuiltInTok{print}\NormalTok{(}\StringTok{"mRNA:"}\NormalTok{, transcribe\_dna\_to\_mrna(template))}
\BuiltInTok{print}\NormalTok{(}\StringTok{"methylation after 3 divisions:"}\NormalTok{, }\BuiltInTok{round}\NormalTok{(cpg\_methylation\_remaining(}\FloatTok{0.9}\NormalTok{, }\DecValTok{3}\NormalTok{, }\FloatTok{0.85}\NormalTok{), }\DecValTok{3}\NormalTok{))}
\end{Highlighting}
\end{Shaded}

\subsubsection*{Part 3: Semiconservative Replication
Analysis}\label{part-3-semiconservative-replication-analysis}
\addcontentsline{toc}{subsubsection}{Part 3: Semiconservative
Replication Analysis}

\begin{enumerate}
\def\labelenumi{\arabic{enumi}.}
\setcounter{enumi}{5}
\item
  \textbf{Design an Alternative Investigation}: Using the checkpoint
  perturbation cards, design a paper simulation to test the effect of a
  drug that inhibits DNA polymerase on cell cycle progression. What
  hypothesis would you test? Which stage proportions would be your
  evidence?
\item
  \textbf{Hypothesis}: Predict how inhibiting DNA polymerase would
  affect the duration of S phase and the overall cell cycle time.
\end{enumerate}

\subsection*{Data Recording}\label{data-recording-11}
\addcontentsline{toc}{subsection}{Data Recording}

{\def\LTcaptype{none} % do not increment counter
\begin{longtable}[]{@{}
  >{\raggedright\arraybackslash}p{(\linewidth - 6\tabcolsep) * \real{0.0781}}
  >{\raggedright\arraybackslash}p{(\linewidth - 6\tabcolsep) * \real{0.2344}}
  >{\raggedright\arraybackslash}p{(\linewidth - 6\tabcolsep) * \real{0.2188}}
  >{\raggedright\arraybackslash}p{(\linewidth - 6\tabcolsep) * \real{0.4688}}@{}}
\toprule\noalign{}
\begin{minipage}[b]{\linewidth}\raggedright
Stage
\end{minipage} & \begin{minipage}[b]{\linewidth}\raggedright
Number of cells
\end{minipage} & \begin{minipage}[b]{\linewidth}\raggedright
Proportion (\%)
\end{minipage} & \begin{minipage}[b]{\linewidth}\raggedright
Estimated time in stage (min)*
\end{minipage} \\
\midrule\noalign{}
\endhead
\bottomrule\noalign{}
\endlastfoot
Interphase & & & \\
Prophase & & & \\
Metaphase & & & \\
Anaphase & & & \\
Telophase & & & \\
\textbf{Total} & 100 & 100\% & 24 h total \\
\end{longtable}
}

Mitotic index: \_\_\_\_\_\_\_ \%\\
Inter-rater agreement on shared card: \_\_\_\_\_\_\_ \%\\
Cards excluded or flagged, with reason:
\_\_\_\_\_\_\_\_\_\_\_\_\_\_\_\_\_\_\_\_

Generation 2 Meselson-Stahl prediction (sketch):

\subsection*{Evidence and Reproducibility
Checklist}\label{evidence-and-reproducibility-checklist-14}
\addcontentsline{toc}{subsection}{Evidence and Reproducibility
Checklist}

\begin{itemize}
\tightlist
\item
  \textbf{Primary evidence goal:} Model replication dynamics and
  checkpoint outcomes.
\item
  \textbf{Data skill to practice:} Interpret replication or cell-cycle
  data from timing, labeling, or checkpoint perturbations.
\item
  \textbf{BioSkills emphasis:} Process of science, Science and society,
  Quantitative reasoning.
\item
  \textbf{Control logic:} identify at least one positive control, one
  negative control, or one baseline comparison before interpreting
  results.
\item
  \textbf{Measurement discipline:} record units, uncertainty, sample
  size, and any discarded observation with a reason.
\item
  \textbf{Mechanistic link:} connect one result directly to the parent
  chapter's big idea before writing the conclusion.
\item
  \textbf{Reproducibility check:} state one procedural detail that
  another group would need in order to reproduce the result.
\end{itemize}

\subsection*{Paper-Based Evidence
Upgrade}\label{paper-based-evidence-upgrade-14}
\addcontentsline{toc}{subsection}{Paper-Based Evidence Upgrade}

Before answering the analysis questions, annotate the paper dataset for
\textbf{DNA Replication and the Cell Cycle} with a reproducibility pass:

{\def\LTcaptype{none} % do not increment counter
\begin{longtable}[]{@{}
  >{\raggedright\arraybackslash}p{(\linewidth - 2\tabcolsep) * \real{0.5000}}
  >{\raggedright\arraybackslash}p{(\linewidth - 2\tabcolsep) * \real{0.5000}}@{}}
\toprule\noalign{}
\begin{minipage}[b]{\linewidth}\raggedright
Evidence check
\end{minipage} & \begin{minipage}[b]{\linewidth}\raggedright
Student action
\end{minipage} \\
\midrule\noalign{}
\endhead
\bottomrule\noalign{}
\endlastfoot
Control logic & Mark the comparison that functions as the baseline,
negative control, or reference case. \\
Uncertainty & Circle the row, card, diagram feature, or model assumption
most likely to change the conclusion. \\
Model comparison & State whether a simpler rule, null model, or
alternative mechanism could explain the same pattern. \\
Decision threshold & Write the minimum evidence that would make you
revise the interpretation. \\
Reproducibility & Record the exact scoring rule another group would need
to reproduce your classification. \\
\end{longtable}
}

Focus note: When scoring mitotic stages, Meselson-Stahl bands, or
replication-fidelity cards, tie every classification to the molecular
event it claims to show and to a stated counting or banding rule another
group could reproduce. Keep required work paper-based; any
material-handling or equipment version belongs only in an optional
extension.

\subsection*{Analysis Questions}\label{analysis-questions-14}
\addcontentsline{toc}{subsection}{Analysis Questions}

\begin{enumerate}
\def\labelenumi{\arabic{enumi}.}
\item
  Which stage of mitosis took the longest (highest proportion)? Which
  was shortest? Explain in terms of the molecular events occurring in
  each stage.
\item
  Your mitotic index suggests what fraction of cells are actively
  dividing. Would you expect a higher or lower mitotic index in: (a) a
  tumour; (b) a mature neuron; (c) a healing wound? Explain for each.
\item
  In the Meselson-Stahl simulation, after two rounds of replication in
  light medium: (a) what fraction of DNA molecules are hybrid (\texorpdfstring{\textsuperscript{1}}{1}\texorpdfstring{\textsuperscript{5}}{5}N/\texorpdfstring{\textsuperscript{1}}{1}\texorpdfstring{\textsuperscript{4}}{4}N)?
  (b) What fraction are light-light? How do these proportions
  distinguish semiconservative from conservative replication?
\item
  DNA polymerase requires a primer and extends DNA in the 5'\texorpdfstring{\ensuremath{\to}}{->}3'
  direction. Explain why the lagging strand is initiated multiple times
  as Okazaki fragments, and name the enzyme that seals these fragments.
\item
  A drug (e.g., taxol) stabilises microtubules and prevents spindle
  disassembly. At which checkpoint would cells arrest, and what specific
  checkpoint protein detects unattached kinetochores? Why is this drug
  effective against rapidly dividing cancer cells?
\end{enumerate}

\subsection*{Extension Analysis
Questions}\label{extension-analysis-questions-5}
\addcontentsline{toc}{subsection}{Extension Analysis Questions}

\begin{enumerate}
\def\labelenumi{\arabic{enumi}.}
\setcounter{enumi}{5}
\item
  \textbf{Statistical Analysis}: Calculate the standard error of the
  mean (SEM) for your mitotic index across replicate cards. Construct a
  95\% confidence interval. How precise is your estimate?
\item
  \textbf{Experimental Error Analysis}: What are three potential sources
  of error in image-based mitotic staging? How could you modify the
  scoring rules or dataset design to reduce each source of error?
\item
  \textbf{Experimental Design}: If you wanted to test the effect of a
  new chemotherapy drug on cell cycle progression, what would you
  measure? How would you distinguish between drugs that affect DNA
  synthesis vs.~those that affect mitosis?
\end{enumerate}

\subsection*{Group Project Extension
(Multi-Session)}\label{group-project-extension-multi-session-2}
\addcontentsline{toc}{subsection}{Group Project Extension
(Multi-Session)}

\textbf{Design a Cancer Cell Detection Tool Using Machine Learning}

This project will span 2-3 lab sessions:

\textbf{Session 1}: Research Question and Hypothesis - Can we
distinguish cancer cells from normal cells based on cell cycle
abnormalities? - What features (size, shape, staining intensity, mitotic
index) are most predictive? - Write a reproducible sampling plan for
choosing public image datasets and metadata fields

\textbf{Session 2}: Data Collection and Analysis - Use public image
datasets of normal and cancer cells - Extract features using image
analysis software - Train a simple machine learning classifier (e.g.,
logistic regression)

\textbf{Session 3}: Tool Development and Testing - Test your classifier
on unseen data - Calculate sensitivity, specificity, and accuracy -
Create a scientific poster or presentation

\subsection*{Real-World Problem Solving: Cancer
Diagnostics}\label{real-world-problem-solving-cancer-diagnostics}
\addcontentsline{toc}{subsection}{Real-World Problem Solving: Cancer
Diagnostics}

\textbf{Case Study: Liquid Biopsy for Cancer Detection}

Liquid biopsies detect circulating tumor DNA (ctDNA) or circulating
tumor cells (CTCs) in blood samples.

\begin{enumerate}
\def\labelenumi{\arabic{enumi}.}
\item
  \textbf{Research Task}: Investigate the advantages and limitations of
  liquid biopsies compared to traditional tissue biopsies.
\item
  \textbf{Ethical Consideration}: Should liquid biopsy screening be
  offered to asymptomatic individuals? What are the potential benefits
  and risks of early detection?
\item
  \textbf{Policy Proposal}: Design a screening program for a specific
  cancer type using liquid biopsy technology. Consider
  cost-effectiveness, accessibility, and follow-up procedures.
\end{enumerate}

\subsection*{Safety and Ethics Notes}\label{safety-and-ethics-notes-14}
\addcontentsline{toc}{subsection}{Safety and Ethics Notes}

Default lab work uses printed datasets primarily, with no biological or
chemical hazards. When using machine learning on human data, ensure
privacy, use de-identified public datasets, and obtain proper consent
for any non-public images. Any wet-lab root-tip squash should be treated
as an optional instructor-supervised extension with local safety
approval, not as a required activity.

\subsection*{Debrief and Reflection}\label{debrief-and-reflection-11}
\addcontentsline{toc}{subsection}{Debrief and Reflection}

After you finish the practical work, spend 5--10 minutes in your small
group comparing results and discussing the following prompts. Each
member should contribute at least one observation before moving to the
next prompt:

\begin{enumerate}
\def\labelenumi{\arabic{enumi}.}
\item
  \textbf{What did your measurements show} --- compare the group's
  results to the textbook's predictions. Where they diverge, suggest at
  least one mechanistic explanation before concluding ``experimental
  error.''
\item
  \textbf{What would change the outcome} --- propose one modification to
  the procedure that would sharpen the measurement or extend the result
  to a new biological context, and predict what you would observe.
\item
  \textbf{One-sentence headline} --- each student composes a single
  sentence summarising the lab's take-home message, suitable for a
  tweet. Compare sentences across groups; good headlines are short,
  quantitative, and mechanistic.
\item
  \textbf{Connection back to the textbook} --- identify one section of
  \cref{sec:unit_IV_dna_replication_and_cell_cycle} that your data
  either confirmed or complicated. Cite the specific passage.
\item
  \textbf{Experimental Design Reflection}: What was the strongest aspect
  of your experimental design? What would you change if you could repeat
  the experiment?
\end{enumerate}

\subsection*{Further Reading (Lab)}\label{further-reading-lab-11}
\addcontentsline{toc}{subsection}{Further Reading (Lab)}

\begin{itemize}
\tightlist
\item
  Revisit the parent chapter
  \cref{sec:unit_IV_dna_replication_and_cell_cycle} for the theoretical
  foundations on which this lab is built.
\item
  Look up any \textbf{bolded glossary term} introduced in the textbook
  chapter (each has a \texttt{\#gl:term-slug} link in the text) --- its
  master definition is in \breaktt{manuscript/glossary.md}.
\item
  Explore the Python code in \breaktt{src/biology/genetics/genetics.py}
  for DNA replication calculations.
\item
  Use the self-contained Part 2 snippet as the computational template;
  it runs against tested project modules without external notebooks or
  CSV files.
\end{itemize}

\newpage

\section*{Lab 13 --- Gene Expression --- Transcription and
Translation}\label{lab-13-gene-expression-transcription-and-translation}
\addcontentsline{toc}{section}{Lab 13 --- Gene Expression ---
Transcription and Translation}

\label{sec:lab_unit_IV_gene_expression}

\subsection*{Learning Objectives}\label{learning-objectives-53}
\addcontentsline{toc}{subsection}{Learning Objectives}

\emph{This activity accompanies \cref{sec:unit_IV_gene_expression} of
the textbook --- review that chapter before attempting the exercises
below.}

\begin{itemize}
\tightlist
\item
  Transcribe and translate a provided DNA template sequence manually to
  produce an amino acid sequence
\item
  Identify the effect of point mutations (silent, missense, nonsense) on
  the final protein
\item
  Analyse a polysome electron micrograph to determine relative
  transcription and translation rates
\item
  Interpret a Northern blot and a Western blot to determine gene
  expression changes
\end{itemize}

\subsubsection*{Alignment and Rubric
Map}\label{alignment-and-rubric-map-15}
\addcontentsline{toc}{subsubsection}{Alignment and Rubric Map}

\begin{itemize}
\tightlist
\item
  \textbf{Outcome 1 (LO1):} Interpret the supplied evidence or model
  output for Gene Expression --- Transcription and Translation.
\item
  \textbf{Outcome 2 (LO2):} Identify controls and comparison groups that
  make the claim testable.
\item
  \textbf{Outcome 3 (LO3):} Quantify uncertainty, boundary conditions,
  or alternative explanations before concluding.
\item
  \textbf{Outcome 4 (LO4):} Transfer the mechanism to a new biological
  case or public-facing decision.
\item
  \textbf{Chapter LO coverage:} LO1, LO2, LO3, LO4
\item
  \textbf{Rubric dimensions:} evidence; controls; uncertainty;
  mechanism; transfer. \#\# Background \{.unnumbered\}
\end{itemize}

Gene expression flows from DNA \texorpdfstring{\ensuremath{\to}}{->} mRNA (transcription, in the nucleus) \texorpdfstring{\ensuremath{\to}}{->}
protein (translation, by ribosomes in the cytoplasm). The genetic code
is triplet (3 nucleotides = 1 codon), non-overlapping, comprehensive,
and degenerate (multiple codons = one amino acid). A \textbf{polysome}
(polyribosome) consists of multiple ribosomes simultaneously translating
the same mRNA; the number of ribosomes per mRNA is proportional to
protein production demand and mRNA length.

In this lab you will perform manual transcription/translation using
codon tables, predict the effect of mutations, and interpret blot data
to determine whether gene expression changes are occurring at the
transcriptional or translational level.

\subsection*{Paper-Based Materials}\label{paper-based-materials-15}
\addcontentsline{toc}{subsection}{Paper-Based Materials}

{\def\LTcaptype{none} % do not increment counter
\begin{longtable}[]{@{}
  >{\raggedright\arraybackslash}p{(\linewidth - 2\tabcolsep) * \real{0.3333}}
  >{\raggedright\arraybackslash}p{(\linewidth - 2\tabcolsep) * \real{0.6667}}@{}}
\toprule\noalign{}
\begin{minipage}[b]{\linewidth}\raggedright
Item
\end{minipage} & \begin{minipage}[b]{\linewidth}\raggedright
Quantity
\end{minipage} \\
\midrule\noalign{}
\endhead
\bottomrule\noalign{}
\endlastfoot
DNA template sequence cards (instructor-provided: 3 sequences ---
original, missense mutant, nonsense mutant) & 1 set \\
mRNA codon table (standard genetic code) & 1 per student \\
Polysome TEM image (printed) & 1 \\
Northern blot image (printed, 2 lanes: untreated vs treated sample) &
1 \\
Western blot image (printed, 2 lanes: untreated vs treated sample) &
1 \\
Ruler & 1 \\
Coloured pens & 3 \\
\end{longtable}
}

\subsection*{Paper-Based
Investigation}\label{paper-based-investigation-15}
\addcontentsline{toc}{subsection}{Paper-Based Investigation}

\textbf{Part A --- Transcription and Translation Simulation}

\begin{enumerate}
\def\labelenumi{\arabic{enumi}.}
\tightlist
\item
  Using Cards A (original), B (missense), C (nonsense): Transcribe each
  DNA template to mRNA (3'\texorpdfstring{\ensuremath{\to}}{->}5' template \texorpdfstring{\ensuremath{\to}}{->} 5'\texorpdfstring{\ensuremath{\to}}{->}3' mRNA), then translate
  using the codon table to produce an amino acid sequence.
\item
  For Cards B and C, identify: which codon was changed, what type of
  mutation (silent/missense/nonsense/frameshift), and what effect on the
  final protein.
\end{enumerate}

\textbf{Part B --- Polysome Analysis}

\begin{enumerate}
\def\labelenumi{\arabic{enumi}.}
\setcounter{enumi}{2}
\tightlist
\item
  Examine the polysome TEM image. Count the number of ribosomes on the
  longest visible mRNA strand. Measure the approximate length of the
  mRNA strand (mm). Calculate ribosome density (ribosomes per mm of
  mRNA).
\end{enumerate}

\textbf{Part C --- Blot Interpretation}

\begin{enumerate}
\def\labelenumi{\arabic{enumi}.}
\setcounter{enumi}{3}
\tightlist
\item
  The Northern blot shows mRNA levels (bands). The Western blot shows
  protein levels (bands). Compare band intensities between untreated
  (lane 1) and treated (lane 2) samples. Complete the interpretation
  table.
\end{enumerate}

\subsection*{Data Recording}\label{data-recording-12}
\addcontentsline{toc}{subsection}{Data Recording}

{\def\LTcaptype{none} % do not increment counter
\begin{longtable}[]{@{}
  >{\raggedright\arraybackslash}p{(\linewidth - 8\tabcolsep) * \real{0.0625}}
  >{\raggedright\arraybackslash}p{(\linewidth - 8\tabcolsep) * \real{0.2031}}
  >{\raggedright\arraybackslash}p{(\linewidth - 8\tabcolsep) * \real{0.2031}}
  >{\raggedright\arraybackslash}p{(\linewidth - 8\tabcolsep) * \real{0.2656}}
  >{\raggedright\arraybackslash}p{(\linewidth - 8\tabcolsep) * \real{0.2656}}@{}}
\toprule\noalign{}
\begin{minipage}[b]{\linewidth}\raggedright
Card
\end{minipage} & \begin{minipage}[b]{\linewidth}\raggedright
Mutation type
\end{minipage} & \begin{minipage}[b]{\linewidth}\raggedright
Mutated codon
\end{minipage} & \begin{minipage}[b]{\linewidth}\raggedright
Amino acid change
\end{minipage} & \begin{minipage}[b]{\linewidth}\raggedright
Effect on protein
\end{minipage} \\
\midrule\noalign{}
\endhead
\bottomrule\noalign{}
\endlastfoot
A & --- & --- & --- & Normal \\
B & & & & \\
C & & & & \\
\end{longtable}
}

mRNA length: \_\_\_\_\_\_\_\_ mm; Ribosomes counted: \_\_\_\_\_\_\_\_;
Ribosome density: \_\_\_\_\_\_\_\_ per mm

{\def\LTcaptype{none} % do not increment counter
\begin{longtable}[]{@{}
  >{\raggedright\arraybackslash}p{(\linewidth - 6\tabcolsep) * \real{0.0896}}
  >{\raggedright\arraybackslash}p{(\linewidth - 6\tabcolsep) * \real{0.5224}}
  >{\raggedright\arraybackslash}p{(\linewidth - 6\tabcolsep) * \real{0.1791}}
  >{\raggedright\arraybackslash}p{(\linewidth - 6\tabcolsep) * \real{0.2090}}@{}}
\toprule\noalign{}
\begin{minipage}[b]{\linewidth}\raggedright
Sample
\end{minipage} & \begin{minipage}[b]{\linewidth}\raggedright
Northern band (strong/weak/absent)
\end{minipage} & \begin{minipage}[b]{\linewidth}\raggedright
Western band
\end{minipage} & \begin{minipage}[b]{\linewidth}\raggedright
Interpretation
\end{minipage} \\
\midrule\noalign{}
\endhead
\bottomrule\noalign{}
\endlastfoot
Untreated & & & \\
Treated & & & \\
\end{longtable}
}

\subsection*{Evidence and Reproducibility
Checklist}\label{evidence-and-reproducibility-checklist-15}
\addcontentsline{toc}{subsection}{Evidence and Reproducibility
Checklist}

\begin{itemize}
\tightlist
\item
  \textbf{Primary evidence goal:} Trace expression from DNA template to
  protein or regulatory RNA.
\item
  \textbf{Data skill to practice:} Convert DNA/RNA sequence data into
  predicted molecular products.
\item
  \textbf{BioSkills emphasis:} Process of science, Science and society,
  Quantitative reasoning.
\item
  \textbf{Control logic:} identify at least one positive control, one
  negative control, or one baseline comparison before interpreting
  results.
\item
  \textbf{Measurement discipline:} record units, uncertainty, sample
  size, and any discarded observation with a reason.
\item
  \textbf{Mechanistic link:} connect one result directly to the parent
  chapter's big idea before writing the conclusion.
\item
  \textbf{Reproducibility check:} state one procedural detail that
  another group would need in order to reproduce the result.
\end{itemize}

\subsection*{Paper-Based Evidence
Upgrade}\label{paper-based-evidence-upgrade-15}
\addcontentsline{toc}{subsection}{Paper-Based Evidence Upgrade}

Before answering the analysis questions, annotate the paper dataset for
\textbf{Gene Expression --- Transcription and Translation} with a
reproducibility pass:

{\def\LTcaptype{none} % do not increment counter
\begin{longtable}[]{@{}
  >{\raggedright\arraybackslash}p{(\linewidth - 2\tabcolsep) * \real{0.5000}}
  >{\raggedright\arraybackslash}p{(\linewidth - 2\tabcolsep) * \real{0.5000}}@{}}
\toprule\noalign{}
\begin{minipage}[b]{\linewidth}\raggedright
Evidence check
\end{minipage} & \begin{minipage}[b]{\linewidth}\raggedright
Student action
\end{minipage} \\
\midrule\noalign{}
\endhead
\bottomrule\noalign{}
\endlastfoot
Control logic & Mark the comparison that functions as the baseline,
negative control, or reference case. \\
Uncertainty & Circle the row, card, diagram feature, or model assumption
most likely to change the conclusion. \\
Model comparison & State whether a simpler rule, null model, or
alternative mechanism could explain the same pattern. \\
Decision threshold & Write the minimum evidence that would make you
revise the interpretation. \\
Reproducibility & Record the exact scoring rule another group would need
to reproduce your classification. \\
\end{longtable}
}

Focus note: When reading transcription, splicing, or translation
datasets, keep the regulatory steps separate and ask which step the
evidence actually measures, since mRNA level, isoform choice, and
protein output can move independently. Keep required work paper-based;
any material-handling or equipment version belongs only in an optional
extension.

\subsection*{Analysis Questions}\label{analysis-questions-15}
\addcontentsline{toc}{subsection}{Analysis Questions}

\begin{enumerate}
\def\labelenumi{\arabic{enumi}.}
\tightlist
\item
  Card B (missense mutation): the amino acid sequence changed but the
  protein was still functional. Explain how this is possible using the
  concept of functionally conservative amino acid substitutions (e.g.,
  Leu \texorpdfstring{\ensuremath{\to}}{->} Ile).
\item
  Card C introduced a premature stop codon. How does the cell detect and
  degrade mRNAs with premature stop codons? Name this quality-control
  pathway.
\item
  If Northern blot shows no change in mRNA level, but Western blot shows
  dramatically reduced protein, at what level is gene expression being
  regulated? Name two possible regulatory mechanisms.
\item
  A gene has 8 ribosomes per mRNA; a second gene's mRNA has 40 ribosomes
  present. Which gene is being translated more actively? What are two
  reasons a cell might load one mRNA with more ribosomes than another?
\item
  The antibiotic streptomycin causes misreading of the A-site codon by
  destabilising the decoding centre of the 30S ribosomal subunit.
  Predict the effect on translation fidelity and which kingdom of
  organisms would be most affected (hint: consider that streptomycin
  targets 30S vs 60S subunits).
\end{enumerate}

\subsection*{Safety and Ethics Notes}\label{safety-and-ethics-notes-15}
\addcontentsline{toc}{subsection}{Safety and Ethics Notes}

No chemical hazards in this simulation/analysis lab. Handle TEM prints
with care. If using online databases for real sequence data, comply with
institutional access agreements.

\subsection*{Debrief and Reflection}\label{debrief-and-reflection-12}
\addcontentsline{toc}{subsection}{Debrief and Reflection}

After you finish the practical work, spend 5--10 minutes in your small
group comparing results and discussing the following prompts. Each
member should contribute at least one observation before moving to the
next prompt:

\begin{enumerate}
\def\labelenumi{\arabic{enumi}.}
\tightlist
\item
  \textbf{What did your measurements show} --- compare the group's
  results to the textbook's predictions. Where they diverge, suggest at
  least one mechanistic explanation before concluding ``experimental
  error.''
\item
  \textbf{What would change the outcome} --- propose one modification to
  the procedure that would sharpen the measurement or extend the result
  to a new biological context, and predict what you would observe.
\item
  \textbf{One-sentence headline} --- each student composes a single
  sentence summarising the lab's take-home message, suitable for a
  tweet. Compare sentences across groups; good headlines are short,
  quantitative, and mechanistic.
\item
  \textbf{Connection back to the textbook} --- identify one section of
  \cref{sec:unit_IV_gene_expression} that your data either confirmed or
  complicated. Cite the specific passage.
\end{enumerate}

\subsection*{Further Reading (Lab)}\label{further-reading-lab-12}
\addcontentsline{toc}{subsection}{Further Reading (Lab)}

\begin{itemize}
\tightlist
\item
  Revisit the parent chapter \cref{sec:unit_IV_gene_expression} for the
  theoretical foundations on which this lab is built.
\item
  Look up any \textbf{bolded glossary term} introduced in the textbook
  chapter (each has a \texttt{\#gl:term-slug} link in the text) --- its
  master definition is in \breaktt{manuscript/glossary.md}.
\item
  Explore the appended \texttt{src/} module that implements the
  corresponding quantitative model (when applicable) --- referenced in
  the parent chapter's ``Bridge to Computation'' subsection.
\end{itemize}

\emph{Module footer: parent chapter
\texttt{\textbackslash{}cref\{sec:unit\_IV\_gene\_expression\}}; all
numerical quantities in this lab use SI units --- see Appendix D of the
textbook for unit conversions and biological-scale reference values.}

\newpage

\section*{Lab 14 --- Mutations, CRISPR, and
Genomics}\label{lab-14-mutations-crispr-and-genomics}
\addcontentsline{toc}{section}{Lab 14 --- Mutations, CRISPR, and
Genomics}

\label{sec:lab_unit_IV_mutations_and_genomics}

\emph{This activity accompanies
\cref{sec:unit_IV_mutations_and_genomics} of the textbook --- review
that chapter before attempting the exercises below.}

\subsection*{Learning Objectives}\label{learning-objectives-54}
\addcontentsline{toc}{subsection}{Learning Objectives}

\begin{itemize}
\tightlist
\item
  Design and evaluate a CRISPR-Cas9 editing scenario using sequence
  cards and gel images
\item
  Formulate testable hypotheses about mutation effects on protein
  function
\item
  Identify independent, dependent, and controlled variables in genome
  editing experiments
\item
  Analyze experimental data using gel electrophoresis and statistical
  methods
\item
  Write a brief scientific report with hypothesis, methods, results, and
  conclusion
\end{itemize}

\subsubsection*{Alignment and Rubric
Map}\label{alignment-and-rubric-map-16}
\addcontentsline{toc}{subsubsection}{Alignment and Rubric Map}

\begin{itemize}
\tightlist
\item
  \textbf{Outcome 1 (LO1):} Interpret the supplied evidence or model
  output for Mutations, CRISPR, and Genomics.
\item
  \textbf{Outcome 2 (LO2):} Identify controls and comparison groups that
  make the claim testable.
\item
  \textbf{Outcome 3 (LO3):} Quantify uncertainty, boundary conditions,
  or alternative explanations before concluding.
\item
  \textbf{Outcome 4 (LO4):} Transfer the mechanism to a new biological
  case or public-facing decision.
\item
  \textbf{Chapter LO coverage:} LO1, LO2, LO3, LO4
\item
  \textbf{Rubric dimensions:} evidence; controls; uncertainty;
  mechanism; transfer. \#\# Pre-Lab Inquiry Questions \{.unnumbered\}
\end{itemize}

\emph{Complete these questions before coming to lab. Use your textbook
and additional research to inform your hypotheses.}

\begin{enumerate}
\def\labelenumi{\arabic{enumi}.}
\item
  \textbf{Hypothesis Formation}: Based on your understanding of mutation
  types, predict which class of mutation (synonymous, missense,
  nonsense, frameshift) would be most likely to cause a loss of protein
  function. Explain your reasoning.
\item
  \textbf{Experimental Design}: What is the independent variable in a
  CRISPR genome-editing data analysis? What is the dependent variable?
  List at least 5 variables that should be controlled to ensure specific
  interpretation.
\item
  \textbf{Scientific Context}: The T7E1 assay is commonly used to detect
  CRISPR-induced mutations. How does this assay work? What are its
  limitations?
\item
  \textbf{Quantitative Reasoning}: If a CRISPR editing dataset reports
  an editing efficiency of 30\%, how many cells would need to be
  represented to infer 10 successfully edited cells? What factors might
  affect this estimate?
\item
  \textbf{Real-World Application}: Germline genome editing raises
  significant ethical concerns. What are the main arguments for and
  against editing human embryos to eliminate genetic diseases? Where do
  you stand on this issue?
\end{enumerate}

\subsection*{Background}\label{background-6}
\addcontentsline{toc}{subsection}{Background}

Mutations --- heritable changes in DNA sequence --- can be neutral
(synonymous), harmful (loss-of-function missense or nonsense), or
beneficial (gain-of-function). CRISPR-Cas9 (Clustered Regularly
Interspaced Short Palindromic Repeats) is a programmable nuclease: the
guide RNA (gRNA; about 20 nt complementary to the target)
directs Cas9 to create a blunt-ended double-strand break 3 bp upstream
of the PAM sequence (5'-NGG-3'). Repair by NHEJ creates indel mutations
(disruptive); HDR with a template creates precise edits.

\subsection*{Paper-Based Materials}\label{paper-based-materials-16}
\addcontentsline{toc}{subsection}{Paper-Based Materials}

{\def\LTcaptype{none} % do not increment counter
\begin{longtable}[]{@{}
  >{\raggedright\arraybackslash}p{(\linewidth - 2\tabcolsep) * \real{0.3333}}
  >{\raggedright\arraybackslash}p{(\linewidth - 2\tabcolsep) * \real{0.6667}}@{}}
\toprule\noalign{}
\begin{minipage}[b]{\linewidth}\raggedright
Item
\end{minipage} & \begin{minipage}[b]{\linewidth}\raggedright
Quantity
\end{minipage} \\
\midrule\noalign{}
\endhead
\bottomrule\noalign{}
\endlastfoot
Wild-type and mutant DNA sequences (printed cards: WT, SNP1, SNP2,
frameshift) & 1 set \\
Codon table & 1 \\
CRISPR guide RNA design worksheet (target region provided) & 1 \\
Gel electrophoresis image (T7E1 assay: uncut control, CRISPR-treated)
(printed) & 1 \\
Off-target candidate sequence cards with mismatch positions marked & 1
set \\
Editing-efficiency replicate dataset & 1 \\
Pangenome path cards: linear reference, alternate insertion path,
inversion path, and read-alignment summaries & 1 set \\
Prime-editing gamma-globin promoter case card with pegRNA, intended
edits, donor variability, and HbF readout & 1 \\
Ruler (mm) & 1 \\
Ethical debate scenario cards (4 scenarios) & 1 set \\
Calculator or optional Python REPL with this project installed & 1 \\
DNA ladder image (printed) & 1 \\
\end{longtable}
}

\subsection*{Paper-Based
Investigation}\label{paper-based-investigation-16}
\addcontentsline{toc}{subsection}{Paper-Based Investigation}

\subsubsection*{Part 1: Mutation Classification and
Analysis}\label{part-1-mutation-classification-and-analysis}
\addcontentsline{toc}{subsubsection}{Part 1: Mutation Classification and
Analysis}

\begin{enumerate}
\def\labelenumi{\arabic{enumi}.}
\item
  \textbf{Formulate Your Hypothesis}: Based on your pre-lab research,
  write a hypothesis about how different mutation types (missense vs
  nonsense vs frameshift) affect protein function in a specific disease
  context (e.g., cystic fibrosis).
\item
  \textbf{Identify Variables}:

  \begin{itemize}
  \tightlist
  \item
    Independent variable: Mutation type (synonymous, missense, nonsense,
    frameshift)
  \item
    Dependent variable: Predicted effect on protein function (neutral,
    mild, severe, lethal)
  \item
    Controlled variables: Gene context, protein domain, evolutionary
    conservation, etc.
  \end{itemize}
\item
  \textbf{Set Up Controls}: Why is it important to have a wild-type (WT)
  control sequence? What other controls could you include in a mutation
  classification experiment?
\item
  \textbf{Data Collection}:

  \begin{itemize}
  \tightlist
  \item
    Compare SNP1, SNP2, and frameshift sequences with the WT. For each:
    identify the nucleotide change, classify the mutation type,
    determine the amino acid change, and assess likely functional
    impact.
  \item
    Use the provided codon table and protein domain information to
    inform your assessment.
  \end{itemize}
\end{enumerate}

\subsubsection*{Part 2: Computational Biology Exercise - CRISPR gRNA
Design and Efficiency
Analysis}\label{part-2-computational-biology-exercise---crispr-grna-design-and-efficiency-analysis}
\addcontentsline{toc}{subsubsection}{Part 2: Computational Biology
Exercise - CRISPR gRNA Design and Efficiency Analysis}

\emph{Optional computational check: run this self-contained Python
snippet from the project root. It uses tested \texttt{src/biology}
modules and requires no external notebook or CSV file.}

Quantify sequence change with the same distance functions used in the
chapter.

\begin{Shaded}
\begin{Highlighting}[]
\ImportTok{from}\NormalTok{ biology.genetics }\ImportTok{import}\NormalTok{ gc\_content, hamming\_distance, jukes\_cantor\_distance}

\NormalTok{reference }\OperatorTok{=} \StringTok{"ATGCGTAC"}
\NormalTok{variant }\OperatorTok{=} \StringTok{"ATGAGTTC"}
\NormalTok{p\_distance }\OperatorTok{=}\NormalTok{ hamming\_distance(reference, variant) }\OperatorTok{/} \BuiltInTok{len}\NormalTok{(reference)}

\BuiltInTok{print}\NormalTok{(}\StringTok{"GC reference:"}\NormalTok{, }\BuiltInTok{round}\NormalTok{(gc\_content(reference), }\DecValTok{2}\NormalTok{))}
\BuiltInTok{print}\NormalTok{(}\StringTok{"Hamming distance:"}\NormalTok{, hamming\_distance(reference, variant))}
\BuiltInTok{print}\NormalTok{(}\StringTok{"Jukes{-}Cantor distance:"}\NormalTok{, }\BuiltInTok{round}\NormalTok{(jukes\_cantor\_distance(p\_distance), }\DecValTok{3}\NormalTok{))}
\end{Highlighting}
\end{Shaded}

\subsubsection*{Part 3: Gel Electrophoresis
Interpretation}\label{part-3-gel-electrophoresis-interpretation}
\addcontentsline{toc}{subsubsection}{Part 3: Gel Electrophoresis
Interpretation}

\begin{enumerate}
\def\labelenumi{\arabic{enumi}.}
\setcounter{enumi}{4}
\item
  \textbf{Design an Alternative Investigation}: Instead of proposing a
  live genome-editing experiment, design a blinded analysis of printed
  gRNA-design scenarios. What hypothesis would you test about
  seed-region mismatches, GC content, and editing efficiency? Which
  sequence-card controls would you include?
\item
  \textbf{Hypothesis}: Predict the relationship between gRNA binding
  affinity and editing efficiency in the dataset. Would it be linear,
  saturating, or something else? Why?
\end{enumerate}

\subsubsection*{Part 4: Pangenome and Prime-Editing Decision
Cards}\label{part-4-pangenome-and-prime-editing-decision-cards}
\addcontentsline{toc}{subsubsection}{Part 4: Pangenome and Prime-Editing
Decision Cards}

\begin{enumerate}
\def\labelenumi{\arabic{enumi}.}
\setcounter{enumi}{6}
\tightlist
\item
  Sort the pangenome path cards into three evidence classes: confidently
  represented by a single linear reference, better represented by an
  alternate graph path, and unresolved without long-read or phased
  evidence.
\item
  For the gamma-globin promoter prime-editing card, identify the
  intended regulatory effect, the difference between editing a coding
  sequence and editing an enhancer/promoter motif, and the
  donor-to-donor variability that would limit a comprehensive claim.
\item
  Build a decision matrix comparing Casgevy-style BCL11A enhancer
  disruption, base editing of an \emph{HBB} point mutation, and prime
  editing of \emph{HBG1/HBG2} promoter motifs. Score each for target
  logic, double-strand-break risk, validation burden, and current
  clinical maturity.
\end{enumerate}

\subsection*{Data Recording}\label{data-recording-13}
\addcontentsline{toc}{subsection}{Data Recording}

{\def\LTcaptype{none} % do not increment counter
\begin{longtable}[]{@{}
  >{\raggedright\arraybackslash}p{(\linewidth - 8\tabcolsep) * \real{0.1067}}
  >{\raggedright\arraybackslash}p{(\linewidth - 8\tabcolsep) * \real{0.2267}}
  >{\raggedright\arraybackslash}p{(\linewidth - 8\tabcolsep) * \real{0.1733}}
  >{\raggedright\arraybackslash}p{(\linewidth - 8\tabcolsep) * \real{0.1200}}
  >{\raggedright\arraybackslash}p{(\linewidth - 8\tabcolsep) * \real{0.3733}}@{}}
\toprule\noalign{}
\begin{minipage}[b]{\linewidth}\raggedright
Sequence
\end{minipage} & \begin{minipage}[b]{\linewidth}\raggedright
Nucleotide change
\end{minipage} & \begin{minipage}[b]{\linewidth}\raggedright
Mutation type
\end{minipage} & \begin{minipage}[b]{\linewidth}\raggedright
AA change
\end{minipage} & \begin{minipage}[b]{\linewidth}\raggedright
Functional impact prediction
\end{minipage} \\
\midrule\noalign{}
\endhead
\bottomrule\noalign{}
\endlastfoot
SNP1 & & & & \\
SNP2 & & & & \\
Frameshift & & & & \\
\end{longtable}
}

gRNA spacer sequences identified: (list most with PAM)

Gel band sizes: Control: \_\_\_\_\_\_\_ bp; Lane 2 bands: \_\_\_\_\_\_\_
bp, \_\_\_\_\_\_\_ bp\\
Indel efficiency estimate = (cleavage bands / total DNA) \texorpdfstring{\ensuremath{\times}}{x} 100\%:
\_\_\_\_\_\_\_\% Off-target risk score for selected gRNA:
\_\_\_\_\_\_\_; highest-risk mismatch pattern: \_\_\_\_\_\_\_ Replicate
editing efficiencies: \textbf{\emph{\%, }}\%, \textbf{\emph{\%; mean \texorpdfstring{\ensuremath{\pm}}{+/-}
SEM: }}\_\textbf{\emph{ Pangenome path selected: }}\_\textbf{\emph{;
graph evidence needed: }}\_\textbf{\emph{; prime-editing validation
readout: }}\_\_\_\_

Ethics summary (2--3 sentences per group position):

\subsection*{Evidence and Reproducibility
Checklist}\label{evidence-and-reproducibility-checklist-16}
\addcontentsline{toc}{subsection}{Evidence and Reproducibility
Checklist}

\begin{itemize}
\tightlist
\item
  \textbf{Primary evidence goal:} Evaluate genome-editing or
  variant-calling scenarios.
\item
  \textbf{Data skill to practice:} Classify variants and predict likely
  molecular effect from sequence evidence.
\item
  \textbf{BioSkills emphasis:} Process of science, Science and society,
  Quantitative reasoning.
\item
  \textbf{Control logic:} identify at least one positive control, one
  negative control, or one baseline comparison before interpreting
  results.
\item
  \textbf{Measurement discipline:} record units, uncertainty, sample
  size, and any discarded observation with a reason.
\item
  \textbf{Mechanistic link:} connect one result directly to the parent
  chapter's big idea before writing the conclusion.
\item
  \textbf{Reproducibility check:} state one procedural detail that
  another group would need in order to reproduce the result.
\end{itemize}

\subsection*{Paper-Based Evidence
Upgrade}\label{paper-based-evidence-upgrade-16}
\addcontentsline{toc}{subsection}{Paper-Based Evidence Upgrade}

Before answering the analysis questions, annotate the paper dataset for
\textbf{Mutations, CRISPR, and Genomics} with a reproducibility pass:

{\def\LTcaptype{none} % do not increment counter
\begin{longtable}[]{@{}
  >{\raggedright\arraybackslash}p{(\linewidth - 2\tabcolsep) * \real{0.5000}}
  >{\raggedright\arraybackslash}p{(\linewidth - 2\tabcolsep) * \real{0.5000}}@{}}
\toprule\noalign{}
\begin{minipage}[b]{\linewidth}\raggedright
Evidence check
\end{minipage} & \begin{minipage}[b]{\linewidth}\raggedright
Student action
\end{minipage} \\
\midrule\noalign{}
\endhead
\bottomrule\noalign{}
\endlastfoot
Control logic & Mark the comparison that functions as the baseline,
negative control, or reference case. \\
Uncertainty & Circle the row, card, diagram feature, or model assumption
most likely to change the conclusion. \\
Model comparison & State whether a simpler rule, null model, or
alternative mechanism could explain the same pattern. \\
Decision threshold & Write the minimum evidence that would make you
revise the interpretation. \\
Reproducibility & Record the exact scoring rule another group would need
to reproduce your classification. \\
\end{longtable}
}

Focus note: When classifying variants from the paper dataset, separate
the act of calling a variant from interpreting its pathogenicity, and
record the ancestry coverage and evidence threshold behind any clinical
claim. Keep required work paper-based; any material-handling or
equipment version belongs only in an optional extension.

\subsection*{Analysis Questions}\label{analysis-questions-16}
\addcontentsline{toc}{subsection}{Analysis Questions}

\begin{enumerate}
\def\labelenumi{\arabic{enumi}.}
\item
  Why does a frameshift mutation typically have a more severe effect on
  protein function than a missense mutation? Under what circumstance
  could a frameshift mutation be silent (no effect on protein)?
\item
  A gRNA with sequence mismatches at positions 1--4 from the PAM distal
  end is less likely to cause cutting than mismatches at the seed region
  (positions 1--12 from PAM-proximal end). Why? What does this imply
  about gRNA design strategy to minimise off-target effects?
\item
  The T7E1 enzyme cleaves primarily at mismatched positions in
  heteroduplex DNA. Explain why NHEJ-induced indels would create
  mismatches and thus be detectable by T7E1, whereas HDR edits with
  perfect template would not be.
\item
  In 2018, He Jiankui used CRISPR to edit the CCR5 gene in human
  embryos, which were then implanted and brought to term. The stated
  goal was to confer HIV resistance. Using your ethical framework from
  Part D: name two distinct ethical objections to this experiment beyond
  the scientific accuracy concerns.
\item
  A patient with Duchenne muscular dystrophy (DMD) has a frameshift in
  exon 51 causing a premature stop codon. Propose a CRISPR-based
  strategy (exon skipping via splice site disruption) to restore the
  reading frame. What are the limitations of this approach?
\item
  A short-read sample maps poorly to a disease locus because one
  haplotype contains a 3-kb insertion absent from the linear reference.
  Explain how a pangenome graph could represent the evidence more
  fairly, and name one validation step before using the call clinically.
\end{enumerate}

\subsection*{Extension Analysis
Questions}\label{extension-analysis-questions-6}
\addcontentsline{toc}{subsection}{Extension Analysis Questions}

\begin{enumerate}
\def\labelenumi{\arabic{enumi}.}
\setcounter{enumi}{6}
\item
  \textbf{Statistical Analysis}: If the printed dataset includes three
  replicate editing-efficiency estimates for each gRNA, how would you
  analyze the variation? What statistical test would you use to
  determine if differences between gRNAs are significant?
\item
  \textbf{Experimental Error Analysis}: What are three potential sources
  of error in T7E1 gel analysis? How could you modify the procedure to
  reduce each source of error?
\item
  \textbf{Experimental Design}: If you wanted to test the effect of a
  specific mutation on protein function, how would you design a
  paper-based evidence plan before any wet work? What sequence controls,
  phenotype readouts, and off-target checks would be required?
\end{enumerate}

\subsection*{Group Project Extension
(Multi-Session)}\label{group-project-extension-multi-session-3}
\addcontentsline{toc}{subsection}{Group Project Extension
(Multi-Session)}

\textbf{Design a CRISPR-based Diagnostic Tool for Infectious Disease}

This project will span 2-3 lab sessions:

\textbf{Session 1}: Research Question and Hypothesis - Can CRISPR-Cas12a
(Cpf1) be used to detect viral DNA (e.g., SARS-CoV-2) with high
specificity? - What are the advantages of CRISPR diagnostics over PCR? -
Write a detailed decision workflow for a CRISPR-based diagnostic assay
using provided sequences and fluorescence traces

\textbf{Session 2}: Data Collection and Analysis - Design gRNAs specific
for a viral target from printed sequence alignments - Simulate detection
by classifying provided fluorescence-output traces - Analyze detection
thresholds, false positives, and false negatives over time

\textbf{Session 3}: Tool Development and Testing - Optimize assay
conditions in a decision matrix (gRNA concentration, reaction time,
detection threshold) - Test specificity against related viral strains -
Create a scientific poster or presentation

\subsection*{Real-World Problem Solving: Genome Editing
Ethics}\label{real-world-problem-solving-genome-editing-ethics}
\addcontentsline{toc}{subsection}{Real-World Problem Solving: Genome
Editing Ethics}

\textbf{Case Study: Somatic vs Germline Editing}

\begin{enumerate}
\def\labelenumi{\arabic{enumi}.}
\item
  \textbf{Research Task}: Investigate the current regulatory landscape
  for somatic vs germline genome editing in two different countries
  (e.g., USA vs China).
\item
  \textbf{Ethical Consideration}: Should there be a global moratorium on
  germline editing? What are the arguments for and against international
  regulation?
\item
  \textbf{Policy Proposal}: Design a policy framework for the
  responsible use of germline genome editing. Consider: which diseases
  would qualify, who would have access, how would long-term effects be
  monitored?
\end{enumerate}

\subsection*{Safety and Ethics Notes}\label{safety-and-ethics-notes-16}
\addcontentsline{toc}{subsection}{Safety and Ethics Notes}

No biological or chemical hazards in this lab (paper/card based). Most
CRISPR discussion is theoretical. When discussing patient/human
scenarios, maintain respectful, evidence-based discourse. Follow
institutional guidelines for ethical discussions.

\subsection*{Debrief and Reflection}\label{debrief-and-reflection-13}
\addcontentsline{toc}{subsection}{Debrief and Reflection}

After you finish the practical work, spend 5--10 minutes in your small
group comparing results and discussing the following prompts. Each
member should contribute at least one observation before moving to the
next prompt:

\begin{enumerate}
\def\labelenumi{\arabic{enumi}.}
\item
  \textbf{What did your measurements show} --- compare the group's
  results to the textbook's predictions. Where they diverge, suggest at
  least one mechanistic explanation before concluding ``experimental
  error.''
\item
  \textbf{What would change the outcome} --- propose one modification to
  the procedure that would sharpen the measurement or extend the result
  to a new biological context, and predict what you would observe.
\item
  \textbf{One-sentence headline} --- each student composes a single
  sentence summarising the lab's take-home message, suitable for a
  tweet. Compare sentences across groups; good headlines are short,
  quantitative, and mechanistic.
\item
  \textbf{Connection back to the textbook} --- identify one section of
  \cref{sec:unit_IV_mutations_and_genomics} that your data either
  confirmed or complicated. Cite the specific passage.
\item
  \textbf{Experimental Design Reflection}: What was the strongest aspect
  of your experimental design? What would you change if you could repeat
  the experiment?
\end{enumerate}

\subsection*{Further Reading (Lab)}\label{further-reading-lab-13}
\addcontentsline{toc}{subsection}{Further Reading (Lab)}

\begin{itemize}
\tightlist
\item
  Revisit the parent chapter \cref{sec:unit_IV_mutations_and_genomics}
  for the theoretical foundations on which this lab is built.
\item
  Look up any \textbf{bolded glossary term} introduced in the textbook
  chapter (each has a \texttt{\#gl:term-slug} link in the text) --- its
  master definition is in \breaktt{manuscript/glossary.md}.
\item
  Explore the Python code in \breaktt{src/biology/genetics/genetics.py}
  for mutation analysis calculations.
\item
  Use the self-contained Part 2 snippet as the computational template;
  it runs against tested project modules without external notebooks or
  CSV files.
\end{itemize}

\newpage

\section*{Lab 15 --- Epigenetics and Gene
Regulation}\label{lab-15-epigenetics-and-gene-regulation}
\addcontentsline{toc}{section}{Lab 15 --- Epigenetics and Gene
Regulation}

\label{sec:lab_unit_IV_epigenetics_and_gene_regulation}

\emph{This activity accompanies
\cref{sec:unit_IV_epigenetics_and_gene_regulation} of the textbook ---
review that chapter before attempting the exercises below.}

\subsection*{Learning Objectives}\label{learning-objectives-55}
\addcontentsline{toc}{subsection}{Learning Objectives}

\begin{itemize}
\tightlist
\item
  Design and evaluate an epigenetic data-analysis scenario to determine
  gene regulation
\item
  Formulate testable hypotheses about how environmental factors affect
  epigenetic marks
\item
  Identify independent, dependent, and controlled variables in
  epigenetic studies
\item
  Analyze experimental data using methylation analysis and ChIP-seq
  interpretation
\item
  Write a brief scientific report with hypothesis, methods, results, and
  conclusion
\end{itemize}

\subsubsection*{Alignment and Rubric
Map}\label{alignment-and-rubric-map-17}
\addcontentsline{toc}{subsubsection}{Alignment and Rubric Map}

\begin{itemize}
\tightlist
\item
  \textbf{Outcome 1 (LO1):} Interpret the supplied evidence or model
  output for Epigenetics and Gene Regulation.
\item
  \textbf{Outcome 2 (LO2):} Identify controls and comparison groups that
  make the claim testable.
\item
  \textbf{Outcome 3 (LO3):} Quantify uncertainty, boundary conditions,
  or alternative explanations before concluding.
\item
  \textbf{Outcome 4 (LO4):} Transfer the mechanism to a new biological
  case or public-facing decision.
\item
  \textbf{Chapter LO coverage:} LO1, LO2, LO3, LO4
\item
  \textbf{Rubric dimensions:} evidence; controls; uncertainty;
  mechanism; transfer. \#\# Pre-Lab Inquiry Questions \{.unnumbered\}
\end{itemize}

\emph{Complete these questions before coming to lab. Use your textbook
and additional research to inform your hypotheses.}

\begin{enumerate}
\def\labelenumi{\arabic{enumi}.}
\item
  \textbf{Hypothesis Formation}: Based on your understanding of DNA
  methylation, predict how a gene promoter's methylation status might
  change in response to a high-sugar diet. Would you expect
  hypermethylation or hypomethylation? Why?
\item
  \textbf{Experimental Design}: What is the independent variable in an
  epigenetic diet study? What is the dependent variable? List at least 5
  variables that should be controlled to ensure valid results.
\item
  \textbf{Scientific Context}: The Horvath epigenetic clock uses DNA
  methylation patterns to predict biological age. Why might biological
  age differ from chronological age? What factors can accelerate
  epigenetic ageing?
\item
  \textbf{Quantitative Reasoning}: If a methylation-sensitive PCR shows
  a 200 bp band in the untreated sample and a 400 bp band in the
  bisulphite-treated sample, what does this tell you about the
  methylation status of the gene?
\item
  \textbf{Real-World Application}: Epigenetic drugs like azacitidine are
  used to treat certain cancers. How do these drugs work? What are the
  potential side effects of globally altering epigenetic marks?
\end{enumerate}

\subsection*{Background}\label{background-7}
\addcontentsline{toc}{subsection}{Background}

Epigenetic regulation controls gene expression without altering DNA
sequence: DNA methylation (CpG islands; methylation = gene silencing,
generally), histone modification (H3K4me3 = activation; H3K27me3 =
silencing; H3K9ac = activation), and non-coding RNA (miRNA, siRNA). The
lac operon is the classical model for prokaryotic transcription
regulation: the lac repressor blocks transcription in the absence of
lactose; allolactose (a lactose isomer) binds the repressor and releases
it from the operator.

\subsection*{Paper-Based Materials}\label{paper-based-materials-17}
\addcontentsline{toc}{subsection}{Paper-Based Materials}

{\def\LTcaptype{none} % do not increment counter
\begin{longtable}[]{@{}
  >{\raggedright\arraybackslash}p{(\linewidth - 2\tabcolsep) * \real{0.3333}}
  >{\raggedright\arraybackslash}p{(\linewidth - 2\tabcolsep) * \real{0.6667}}@{}}
\toprule\noalign{}
\begin{minipage}[b]{\linewidth}\raggedright
Item
\end{minipage} & \begin{minipage}[b]{\linewidth}\raggedright
Quantity
\end{minipage} \\
\midrule\noalign{}
\endhead
\bottomrule\noalign{}
\endlastfoot
MS-PCR gel image (printed: two lanes per sample,
bisulphite-treated/untreated) & 1 \\
ChIP-seq track printout (H3K4me3 signal at BRCA1 promoter: tumour vs
normal) & 1 \\
Lac operon gene-expression logic table (worksheet) & 1 per student \\
``Agouti mouse'' case study reading (1 page) & 1 per student \\
Coloured pens & 3 \\
Calculator or optional Python REPL with this project installed & 1 \\
Epigenetic age calculator worksheet & 1 \\
\end{longtable}
}

\subsection*{Paper-Based
Investigation}\label{paper-based-investigation-17}
\addcontentsline{toc}{subsection}{Paper-Based Investigation}

\subsubsection*{Part 1: Epigenetic Data
Analysis}\label{part-1-epigenetic-data-analysis}
\addcontentsline{toc}{subsubsection}{Part 1: Epigenetic Data Analysis}

\begin{enumerate}
\def\labelenumi{\arabic{enumi}.}
\item
  \textbf{Formulate Your Hypothesis}: Based on your pre-lab research,
  write a hypothesis about how a specific environmental factor (diet,
  stress, toxins) might affect DNA methylation patterns at a particular
  gene.
\item
  \textbf{Identify Variables}:

  \begin{itemize}
  \tightlist
  \item
    Independent variable: Environmental exposure (e.g., high-fat diet,
    pollutant)
  \item
    Dependent variable: DNA methylation level at specific CpG sites
  \item
    Controlled variables: Age, sex, genetic background, sample
    processing, etc.
  \end{itemize}
\item
  \textbf{Set Up Controls}: Why is it important to include both treated
  and untreated control samples in epigenetic studies? What other
  controls could you include to validate your results?
\item
  \textbf{Data Collection}:

  \begin{itemize}
  \tightlist
  \item
    Interpret the MS-PCR gel: determine which gene is methylated in
    cancer vs normal cells.
  \item
    Analyze the ChIP-seq track: compare H3K4me3 peak height at the BRCA1
    promoter between tumour and normal. Does BRCA1 appear silenced in
    the tumour?
  \item
    Complete the lac operon logic table for four conditions.
  \item
    Read and analyze the agouti mouse case study.
  \end{itemize}
\end{enumerate}

\subsubsection*{Part 2: Computational Biology Exercise - Epigenetic Age
Prediction with
Python}\label{part-2-computational-biology-exercise---epigenetic-age-prediction-with-python}
\addcontentsline{toc}{subsubsection}{Part 2: Computational Biology
Exercise - Epigenetic Age Prediction with Python}

\emph{Optional computational check: run this self-contained Python
snippet from the project root. It uses tested \texttt{src/biology}
modules and requires no external notebook or CSV file.}

Model maintenance methylation and classify histone marks without
external files.

\begin{Shaded}
\begin{Highlighting}[]
\ImportTok{from}\NormalTok{ biology.genetics }\ImportTok{import}\NormalTok{ cpg\_methylation\_remaining, histone\_modification\_state}

\ControlFlowTok{for}\NormalTok{ efficiency }\KeywordTok{in}\NormalTok{ (}\FloatTok{0.95}\NormalTok{, }\FloatTok{0.85}\NormalTok{, }\FloatTok{0.60}\NormalTok{):}
\NormalTok{    remaining }\OperatorTok{=}\NormalTok{ cpg\_methylation\_remaining(}\FloatTok{0.8}\NormalTok{, divisions}\OperatorTok{=}\DecValTok{4}\NormalTok{, maintenance\_efficiency}\OperatorTok{=}\NormalTok{efficiency)}
    \BuiltInTok{print}\NormalTok{(efficiency, }\BuiltInTok{round}\NormalTok{(remaining, }\DecValTok{3}\NormalTok{))}

\ControlFlowTok{for}\NormalTok{ mark }\KeywordTok{in}\NormalTok{ (}\StringTok{"H3K27me3"}\NormalTok{, }\StringTok{"H3K27ac"}\NormalTok{, }\StringTok{"H3K4me2"}\NormalTok{):}
    \BuiltInTok{print}\NormalTok{(mark, histone\_modification\_state(mark))}
\end{Highlighting}
\end{Shaded}

\subsubsection*{Part 3: Advanced Epigenetic
Analysis}\label{part-3-advanced-epigenetic-analysis}
\addcontentsline{toc}{subsubsection}{Part 3: Advanced Epigenetic
Analysis}

\begin{enumerate}
\def\labelenumi{\arabic{enumi}.}
\setcounter{enumi}{4}
\item
  \textbf{Design an Alternative Investigation}: Instead of just
  analyzing provided data, design a paper-based study plan to test how a
  specific nutrient (e.g., folate, B12) affects DNA methylation patterns
  in a model organism. What hypothesis would you test? Which tissue
  datasets, controls, and metadata fields would you require?
\item
  \textbf{Hypothesis}: Predict the direction and magnitude of
  methylation changes you might observe with nutrient supplementation.
\end{enumerate}

\subsection*{Data Recording}\label{data-recording-14}
\addcontentsline{toc}{subsection}{Data Recording}

{\def\LTcaptype{none} % do not increment counter
\begin{longtable}[]{@{}
  >{\raggedright\arraybackslash}p{(\linewidth - 6\tabcolsep) * \real{0.1154}}
  >{\raggedright\arraybackslash}p{(\linewidth - 6\tabcolsep) * \real{0.3590}}
  >{\raggedright\arraybackslash}p{(\linewidth - 6\tabcolsep) * \real{0.2692}}
  >{\raggedright\arraybackslash}p{(\linewidth - 6\tabcolsep) * \real{0.2564}}@{}}
\toprule\noalign{}
\begin{minipage}[b]{\linewidth}\raggedright
Condition
\end{minipage} & \begin{minipage}[b]{\linewidth}\raggedright
Repressor bound to operator?
\end{minipage} & \begin{minipage}[b]{\linewidth}\raggedright
CAP-cAMP at promoter?
\end{minipage} & \begin{minipage}[b]{\linewidth}\raggedright
lac genes expressed?
\end{minipage} \\
\midrule\noalign{}
\endhead
\bottomrule\noalign{}
\endlastfoot
No lactose, no glucose & Yes & No & No \\
Lactose present, no glucose & & & \\
No lactose, glucose present & & & \\
Lactose + glucose & & & \\
\end{longtable}
}

MS-PCR result --- BRCA1: Methylated? (Y/N) Normal: \_\_\_ Cancer:
\_\_\_\\
ChIP-seq --- BRCA1 H3K4me3 peak height: Normal: \_\_\_ Cancer: \_\_\_
Interpretation: \_\_\_

Agouti mouse case study: Methyl donor involved: \textbf{\emph{;
Epigenetic enzyme: }}; Phenotypic consequence: \_\_\_

\subsection*{Evidence and Reproducibility
Checklist}\label{evidence-and-reproducibility-checklist-17}
\addcontentsline{toc}{subsection}{Evidence and Reproducibility
Checklist}

\begin{itemize}
\tightlist
\item
  \textbf{Primary evidence goal:} Compare regulatory perturbations and
  predicted expression outputs.
\item
  \textbf{Data skill to practice:} Interpret chromatin or expression
  evidence from simple regulatory datasets.
\item
  \textbf{BioSkills emphasis:} Process of science, Science and society,
  Quantitative reasoning.
\item
  \textbf{Control logic:} identify at least one positive control, one
  negative control, or one baseline comparison before interpreting
  results.
\item
  \textbf{Measurement discipline:} record units, uncertainty, sample
  size, and any discarded observation with a reason.
\item
  \textbf{Mechanistic link:} connect one result directly to the parent
  chapter's big idea before writing the conclusion.
\item
  \textbf{Reproducibility check:} state one procedural detail that
  another group would need in order to reproduce the result.
\end{itemize}

\subsection*{Paper-Based Evidence
Upgrade}\label{paper-based-evidence-upgrade-17}
\addcontentsline{toc}{subsection}{Paper-Based Evidence Upgrade}

Before answering the analysis questions, annotate the paper dataset for
\textbf{Epigenetics and Gene Regulation} with a reproducibility pass:

{\def\LTcaptype{none} % do not increment counter
\begin{longtable}[]{@{}
  >{\raggedright\arraybackslash}p{(\linewidth - 2\tabcolsep) * \real{0.5000}}
  >{\raggedright\arraybackslash}p{(\linewidth - 2\tabcolsep) * \real{0.5000}}@{}}
\toprule\noalign{}
\begin{minipage}[b]{\linewidth}\raggedright
Evidence check
\end{minipage} & \begin{minipage}[b]{\linewidth}\raggedright
Student action
\end{minipage} \\
\midrule\noalign{}
\endhead
\bottomrule\noalign{}
\endlastfoot
Control logic & Mark the comparison that functions as the baseline,
negative control, or reference case. \\
Uncertainty & Circle the row, card, diagram feature, or model assumption
most likely to change the conclusion. \\
Model comparison & State whether a simpler rule, null model, or
alternative mechanism could explain the same pattern. \\
Decision threshold & Write the minimum evidence that would make you
revise the interpretation. \\
Reproducibility & Record the exact scoring rule another group would need
to reproduce your classification. \\
\end{longtable}
}

Focus note: When interpreting methylation, chromatin, or gene-silencing
datasets, demand a causal perturbation, a matched cell type, and a
control for timing or inheritance before reading a pattern as a
regulatory mechanism. Keep required work paper-based; any
material-handling or equipment version belongs only in an optional
extension.

\subsection*{Analysis Questions}\label{analysis-questions-17}
\addcontentsline{toc}{subsection}{Analysis Questions}

\begin{enumerate}
\def\labelenumi{\arabic{enumi}.}
\item
  The lac repressor is an allosteric protein. How does allolactose
  binding change the shape of the repressor so that it can no longer
  bind the operator? Relate this to the induced fit model.
\item
  In catabolite repression, high glucose causes low cAMP because glucose
  inhibits adenylate cyclase. Trace: high glucose \texorpdfstring{\ensuremath{\to}}{->} {[}cAMP{]} \texorpdfstring{\ensuremath{\to}}{->}
  CAP-cAMP binding \texorpdfstring{\ensuremath{\to}}{->} lac promoter activity. Why does the cell
  preferentially use glucose over lactose?
\item
  A cancer cell has methylation of a tumour suppressor gene promoter but
  no mutation in the coding sequence. Is this a genetic or epigenetic
  mutation? Why is this clinically important for therapy selection
  (hint: methyltransferase inhibitors vs traditional chemotherapy)?
\item
  The agouti mouse model showed that maternal diet (methyl-group-rich
  foods) changes offspring coat colour and obesity risk through
  epigenetic marks. Name the methyl donors involved (from diet) and the
  epigenetic writing enzyme that places the mark on CpG islands.
\item
  Design an experiment using induced pluripotent stem cell (iPSC)
  technology to test whether erasing epigenetic marks (global
  demethylation) is sufficient to reprogram a differentiated skin
  fibroblast. What control groups would you include?
\end{enumerate}

\subsection*{Extension Analysis
Questions}\label{extension-analysis-questions-7}
\addcontentsline{toc}{subsection}{Extension Analysis Questions}

\begin{enumerate}
\def\labelenumi{\arabic{enumi}.}
\setcounter{enumi}{5}
\item
  Methylation-sensitive restriction enzymes (e.g., HpaII) and
  methyl-insensitive isoschizomers (e.g., MspI) cut the same sequence
  but primarily MspI cuts methylated DNA. Describe one experimental use
  of this enzyme pair for mapping genome-wide methylation without
  bisulphite conversion.
\item
  Transgenerational epigenetic inheritance in humans remains
  controversial. What \emph{specific} molecular evidence would convince
  you it is real? (Address confounding: shared environment, germline
  reprogramming.)
\item
  The textbook introduces CRISPR-dCas9 epigenome editing. Design a
  dCas9-based experiment to re-activate the silenced tumour-suppressor
  gene you identified in the ChIP-seq analysis, and predict one
  off-target risk.
\item
  \textbf{Statistical Analysis}: If you performed three replicates of an
  MS-PCR experiment, how would you analyze the variation in methylation
  status? What statistical test would you use to determine if
  differences between groups are significant?
\item
  \textbf{Experimental Error Analysis}: What are three potential sources
  of error in epigenetic data analysis? How could you modify the
  procedure to reduce each source of error?
\end{enumerate}

\subsection*{Group Project Extension
(Multi-Session)}\label{group-project-extension-multi-session-4}
\addcontentsline{toc}{subsection}{Group Project Extension
(Multi-Session)}

\textbf{Design an Epigenetic Diet Intervention Study}

This project will span 2-3 lab sessions:

\textbf{Session 1}: Research Question and Hypothesis - Can dietary
methyl donors (folate, B12, choline) alter DNA methylation patterns in
humans? - What genes might be most responsive to such changes? - Write a
detailed protocol for a dietary intervention study

\textbf{Session 2}: Data Collection and Analysis - Collect dietary
intake data and DNA samples (or use public datasets) - Analyze
methylation at candidate genes using bisulphite sequencing - Compare
methylation changes between intervention and control groups

\textbf{Session 3}: Data Interpretation and Presentation - Correlate
methylation changes with health outcomes - Create a scientific poster or
presentation - Discuss limitations and future directions

\subsection*{Real-World Problem Solving: Epigenetic
Ethics}\label{real-world-problem-solving-epigenetic-ethics}
\addcontentsline{toc}{subsection}{Real-World Problem Solving: Epigenetic
Ethics}

\textbf{Case Study: Epigenetic Discrimination}

\begin{enumerate}
\def\labelenumi{\arabic{enumi}.}
\item
  \textbf{Research Task}: Investigate the current legal protections
  against genetic discrimination (e.g., GINA in the US). Do these laws
  cover epigenetic information? Why or why not?
\item
  \textbf{Ethical Consideration}: Should employers or insurance
  companies be allowed to use epigenetic age estimates to make
  decisions? What are the potential benefits and risks?
\item
  \textbf{Policy Proposal}: Design a policy framework for the ethical
  use of epigenetic information. Consider: consent, privacy, access, and
  potential for discrimination.
\end{enumerate}

\subsection*{Safety and Ethics Notes}\label{safety-and-ethics-notes-17}
\addcontentsline{toc}{subsection}{Safety and Ethics Notes}

Data-analysis primarily lab --- no reagents, no risk. Discussions of
patient cancer data should respect privacy principles
(de-identification; aggregate reporting). Agouti mouse studies involve
animal experimentation --- students should briefly discuss the 3Rs
framework (Replacement, Reduction, Refinement) and when mouse models are
(and are not) scientifically justified. Epigenetic age clocks have
attracted commercial wellness-industry interest; discuss the limits of
current evidence before drawing lifestyle conclusions from a single
patient's clock.

\subsection*{Debrief and Reflection}\label{debrief-and-reflection-14}
\addcontentsline{toc}{subsection}{Debrief and Reflection}

After you finish the practical work, spend 5--10 minutes in your small
group comparing results and discussing the following prompts. Each
member should contribute at least one observation before moving to the
next prompt:

\begin{enumerate}
\def\labelenumi{\arabic{enumi}.}
\item
  \textbf{What did your measurements show} --- compare the group's
  results to the textbook's predictions. Where they diverge, suggest at
  least one mechanistic explanation before concluding ``experimental
  error.''
\item
  \textbf{What would change the outcome} --- propose one modification to
  the procedure that would sharpen the measurement or extend the result
  to a new biological context, and predict what you would observe.
\item
  \textbf{One-sentence headline} --- each student composes a single
  sentence summarising the lab's take-home message, suitable for a
  tweet. Compare sentences across groups; good headlines are short,
  quantitative, and mechanistic.
\item
  \textbf{Connection back to the textbook} --- identify one section of
  \cref{sec:unit_IV_epigenetics_and_gene_regulation} that your data
  either confirmed or complicated. Cite the specific passage.
\item
  \textbf{Experimental Design Reflection}: What was the strongest aspect
  of your experimental design? What would you change if you could repeat
  the experiment?
\end{enumerate}

\subsection*{Further Reading (Lab)}\label{further-reading-lab-14}
\addcontentsline{toc}{subsection}{Further Reading (Lab)}

\begin{itemize}
\tightlist
\item
  Revisit the parent chapter
  \cref{sec:unit_IV_epigenetics_and_gene_regulation} for the theoretical
  foundations on which this lab is built.
\item
  Look up any \textbf{bolded glossary term} introduced in the textbook
  chapter (each has a \texttt{\#gl:term-slug} link in the text) --- its
  master definition is in \breaktt{manuscript/glossary.md}.
\item
  Explore the Python code in \breaktt{src/biology/genetics/genetics.py}
  for epigenetic calculations.
\item
  Use the self-contained Part 2 snippet as the computational template;
  it runs against tested project modules without external notebooks or
  CSV files.
\end{itemize}

\newpage

\section*{Lab 16 --- Mendelian Genetics and
Heredity}\label{lab-16-mendelian-genetics-and-heredity}
\addcontentsline{toc}{section}{Lab 16 --- Mendelian Genetics and
Heredity}

\label{sec:lab_unit_V_mendelian_genetics}

\emph{This activity accompanies \cref{sec:unit_V_mendelian_genetics} of
the textbook --- review that chapter before attempting the exercises
below.}

\subsection*{Learning Objectives}\label{learning-objectives-56}
\addcontentsline{toc}{subsection}{Learning Objectives}

\begin{itemize}
\tightlist
\item
  Design and evaluate genetic cross simulations to determine inheritance
  patterns
\item
  Formulate testable hypotheses about dominance, recessiveness, and
  linkage
\item
  Identify independent, dependent, and controlled variables in genetic
  experiments
\item
  Analyze experimental data using chi-square tests and statistical
  methods
\item
  Write a brief scientific report with hypothesis, methods, results, and
  conclusion
\end{itemize}

\subsubsection*{Alignment and Rubric
Map}\label{alignment-and-rubric-map-18}
\addcontentsline{toc}{subsubsection}{Alignment and Rubric Map}

\begin{itemize}
\tightlist
\item
  \textbf{Outcome 1 (LO1):} Interpret the supplied evidence or model
  output for Mendelian Genetics and Heredity.
\item
  \textbf{Outcome 2 (LO2):} Identify controls and comparison groups that
  make the claim testable.
\item
  \textbf{Outcome 3 (LO3):} Quantify uncertainty, boundary conditions,
  or alternative explanations before concluding.
\item
  \textbf{Outcome 4 (LO4):} Transfer the mechanism to a new biological
  case or public-facing decision.
\item
  \textbf{Chapter LO coverage:} LO1, LO2, LO3, LO4
\item
  \textbf{Rubric dimensions:} evidence; controls; uncertainty;
  mechanism; transfer. \#\# Pre-Lab Inquiry Questions \{.unnumbered\}
\end{itemize}

\emph{Complete these questions before coming to lab. Use your textbook
and additional research to inform your hypotheses.}

\begin{enumerate}
\def\labelenumi{\arabic{enumi}.}
\item
  \textbf{Hypothesis Formation}: Based on your understanding of
  Mendelian inheritance, predict the outcome of a test cross between a
  dominant phenotype individual (unknown genotype) and a homozygous
  recessive individual. What phenotypic ratio would you expect if the
  dominant individual is homozygous vs heterozygous?
\item
  \textbf{Experimental Design}: What is the independent variable in a
  Mendelian inheritance experiment? What is the dependent variable? List
  at least 5 variables that should be controlled to ensure valid
  results.
\item
  \textbf{Scientific Context}: The chi-square test is used to determine
  if observed ratios fit expected Mendelian ratios. Why is this test
  important in genetics? What does a significant deviation from expected
  ratios tell you?
\item
  \textbf{Quantitative Reasoning}: If you perform a monohybrid cross (Aa
  \texorpdfstring{\ensuremath{\times}}{x} Aa) and observe 450 round seeds and 130 wrinkled seeds (total 580),
  what is the \texorpdfstring{\ensuremath{\chi}}{chi}\texorpdfstring{\textsuperscript{2}}{2} value? With 1 degree of freedom, is this significantly
  different from the expected 3:1 ratio at p = 0.05?
\item
  \textbf{Real-World Application}: Many human genetic disorders are
  recessive. Why are recessive disorders more common than dominant
  disorders? What are the implications for genetic counseling?
\end{enumerate}

\subsection*{Background}\label{background-8}
\addcontentsline{toc}{subsection}{Background}

Gregor Mendel's experiments with pea plants (1856--1863) established the
laws of segregation (allele pairs separate during gamete formation) and
independent assortment (non-linked genes assort independently).
Monohybrid crosses between heterozygotes (Aa \texorpdfstring{\ensuremath{\times}}{x} Aa) produce 3:1 phenotype
ratios; dihybrid crosses produce 9:3:3:1 ratios. Deviations from these
ratios suggest linkage, incomplete dominance, codominance, or epistasis.

\subsubsection*{Example Pedigree --- Autosomal Recessive Trait (Cystic
Fibrosis)}\label{example-pedigree-autosomal-recessive-trait-cystic-fibrosis}
\addcontentsline{toc}{subsubsection}{Example Pedigree --- Autosomal
Recessive Trait (Cystic Fibrosis)}

Below is the three-generation pedigree you will analyse in Part B.
Filled squares/circles are affected individuals; half-filled are
obligate carriers; empty are unaffected. Use this as a reference for how
pedigree symbols encode genotypes.

\subsection*{Paper-Based Materials}\label{paper-based-materials-18}
\addcontentsline{toc}{subsection}{Paper-Based Materials}

{\def\LTcaptype{none} % do not increment counter
\begin{longtable}[]{@{}
  >{\raggedright\arraybackslash}p{(\linewidth - 2\tabcolsep) * \real{0.3333}}
  >{\raggedright\arraybackslash}p{(\linewidth - 2\tabcolsep) * \real{0.6667}}@{}}
\toprule\noalign{}
\begin{minipage}[b]{\linewidth}\raggedright
Item
\end{minipage} & \begin{minipage}[b]{\linewidth}\raggedright
Quantity
\end{minipage} \\
\midrule\noalign{}
\endhead
\bottomrule\noalign{}
\endlastfoot
Punnett square worksheets (3: monohybrid, dihybrid, incomplete
dominance) & 1 set \\
Pedigree chart (printed: 3-generation autosomal recessive trait) & 1 \\
Coin or random number table (to simulate gamete formation) & 2 \\
Chi-square table (critical values) & 1 \\
Calculator & 1 \\
Coloured counters (2 colours \texorpdfstring{\ensuremath{\times}}{x} 20) for simulating gametes & 1 set \\
Calculator or optional Python REPL with this project installed & 1 \\
Statistical analysis worksheet & 1 \\
\end{longtable}
}

\subsection*{Paper-Based
Investigation}\label{paper-based-investigation-18}
\addcontentsline{toc}{subsection}{Paper-Based Investigation}

\subsubsection*{Part 1: Punnett Square
Analysis}\label{part-1-punnett-square-analysis}
\addcontentsline{toc}{subsubsection}{Part 1: Punnett Square Analysis}

\begin{enumerate}
\def\labelenumi{\arabic{enumi}.}
\item
  \textbf{Formulate Your Hypothesis}: Based on your pre-lab research,
  write a hypothesis about how the inheritance pattern of a specific
  trait (e.g., pea seed shape) follows Mendelian ratios.
\item
  \textbf{Identify Variables}:

  \begin{itemize}
  \tightlist
  \item
    Independent variable: Parental genotypes
  \item
    Dependent variable: Offspring genotype/phenotype ratios
  \item
    Controlled variables: Cross design, counting method, environmental
    conditions, etc.
  \end{itemize}
\item
  \textbf{Set Up Controls}: Why is it important to use true-breeding
  parental lines? What other controls could you include in a Mendelian
  inheritance experiment?
\item
  \textbf{Data Collection}:

  \begin{itemize}
  \tightlist
  \item
    Complete monohybrid crosses: Aa \texorpdfstring{\ensuremath{\times}}{x} Aa; AA \texorpdfstring{\ensuremath{\times}}{x} Aa; Aa \texorpdfstring{\ensuremath{\times}}{x} aa. Record
    genotype and phenotype ratios.
  \item
    Complete a dihybrid cross: AaBb \texorpdfstring{\ensuremath{\times}}{x} AaBb. Record the 9:3:3:1 phenotype
    ratio.
  \item
    Complete an incomplete dominance cross: R\texorpdfstring{\textsubscript{1}}{1}R\texorpdfstring{\textsubscript{1}}{1} (red) \texorpdfstring{\ensuremath{\times}}{x} R\texorpdfstring{\textsubscript{2}}{2}R\texorpdfstring{\textsubscript{2}}{2} (white).
    What is the F\texorpdfstring{\textsubscript{1}}{1} phenotype? F\texorpdfstring{\textsubscript{2}}{2} ratios?
  \end{itemize}
\end{enumerate}

\subsubsection*{Part 2: Computational Biology Exercise - Statistical
Analysis with
Python}\label{part-2-computational-biology-exercise---statistical-analysis-with-python-1}
\addcontentsline{toc}{subsubsection}{Part 2: Computational Biology
Exercise - Statistical Analysis with Python}

\emph{Optional computational check: run this self-contained Python
snippet from the project root. It uses tested \texttt{src/biology}
modules and requires no external notebook or CSV file.}

Validate the paper Punnett-square and chi-square calculations.

\begin{Shaded}
\begin{Highlighting}[]
\ImportTok{from}\NormalTok{ biology.genetics }\ImportTok{import}\NormalTok{ chi\_squared\_test, punnett\_square}

\NormalTok{cross }\OperatorTok{=}\NormalTok{ punnett\_square(}\StringTok{"Aa"}\NormalTok{, }\StringTok{"Aa"}\NormalTok{)}
\NormalTok{observed }\OperatorTok{=}\NormalTok{ [}\FloatTok{450.0}\NormalTok{, }\FloatTok{130.0}\NormalTok{]}
\NormalTok{expected }\OperatorTok{=}\NormalTok{ [}\FloatTok{0.75} \OperatorTok{*} \BuiltInTok{sum}\NormalTok{(observed), }\FloatTok{0.25} \OperatorTok{*} \BuiltInTok{sum}\NormalTok{(observed)]}
\NormalTok{chi }\OperatorTok{=}\NormalTok{ chi\_squared\_test(observed, expected)}

\BuiltInTok{print}\NormalTok{(}\StringTok{"genotype ratios:"}\NormalTok{, cross.genotype\_ratios)}
\BuiltInTok{print}\NormalTok{(}\StringTok{"phenotype ratios:"}\NormalTok{, cross.phenotype\_ratios)}
\BuiltInTok{print}\NormalTok{(}\StringTok{"χ²:"}\NormalTok{, }\BuiltInTok{round}\NormalTok{(chi.chi\_squared, }\DecValTok{2}\NormalTok{), }\StringTok{"reject?"}\NormalTok{, chi.reject\_null)}
\end{Highlighting}
\end{Shaded}

\subsubsection*{Part 3: Pedigree
Analysis}\label{part-3-pedigree-analysis}
\addcontentsline{toc}{subsubsection}{Part 3: Pedigree Analysis}

\begin{enumerate}
\def\labelenumi{\arabic{enumi}.}
\setcounter{enumi}{4}
\item
  \textbf{Design an Alternative Investigation}: Instead of just
  analyzing the provided pedigree, design a paper-based study to
  determine the inheritance pattern of a human genetic trait in a small
  population. What hypothesis would you test? What consent, privacy, and
  pedigree metadata would be required?
\item
  \textbf{Hypothesis}: Predict the inheritance pattern (autosomal
  dominant, autosomal recessive, X-linked) for a trait that appears in
  every generation, affects both males and females equally, and shows no
  father-to-son transmission.
\end{enumerate}

\subsection*{Data Recording}\label{data-recording-15}
\addcontentsline{toc}{subsection}{Data Recording}

{\def\LTcaptype{none} % do not increment counter
\begin{longtable}[]{@{}lll@{}}
\toprule\noalign{}
Cross & Offspring genotype ratio & Offspring phenotype ratio \\
\midrule\noalign{}
\endhead
\bottomrule\noalign{}
\endlastfoot
Aa \texorpdfstring{\ensuremath{\times}}{x} Aa & & \\
AaBb \texorpdfstring{\ensuremath{\times}}{x} AaBb & & \\
R\texorpdfstring{\textsubscript{1}}{1}R\texorpdfstring{\textsubscript{2}}{2} \texorpdfstring{\ensuremath{\times}}{x} R\texorpdfstring{\textsubscript{1}}{1}R\texorpdfstring{\textsubscript{2}}{2} (incomplete dom.) & & \\
\end{longtable}
}

Pedigree genotypes:

{\def\LTcaptype{none} % do not increment counter
\begin{longtable}[]{@{}llll@{}}
\toprule\noalign{}
& Gen I & Gen II & Gen III \\
\midrule\noalign{}
\endhead
\bottomrule\noalign{}
\endlastfoot
Individual & & & \\
\end{longtable}
}

Chi-square result: \texorpdfstring{\ensuremath{\chi}}{chi}\texorpdfstring{\textsuperscript{2}}{2} = \textbf{\emph{; df = }}; p \textless{} 0.05?
(Y/N); Conclusion:

\subsection*{Evidence and Reproducibility
Checklist}\label{evidence-and-reproducibility-checklist-18}
\addcontentsline{toc}{subsection}{Evidence and Reproducibility
Checklist}

\begin{itemize}
\tightlist
\item
  \textbf{Primary evidence goal:} Solve inheritance patterns and test
  expected ratios.
\item
  \textbf{Data skill to practice:} Use family or cross data to infer
  genotype probabilities.
\item
  \textbf{BioSkills emphasis:} Quantitative reasoning, Process of
  science.
\item
  \textbf{Control logic:} identify at least one positive control, one
  negative control, or one baseline comparison before interpreting
  results.
\item
  \textbf{Measurement discipline:} record units, uncertainty, sample
  size, and any discarded observation with a reason.
\item
  \textbf{Mechanistic link:} connect one result directly to the parent
  chapter's big idea before writing the conclusion.
\item
  \textbf{Reproducibility check:} state one procedural detail that
  another group would need in order to reproduce the result.
\end{itemize}

\subsection*{Paper-Based Evidence
Upgrade}\label{paper-based-evidence-upgrade-18}
\addcontentsline{toc}{subsection}{Paper-Based Evidence Upgrade}

Before answering the analysis questions, annotate the paper dataset for
\textbf{Mendelian Genetics and Heredity} with a reproducibility pass:

{\def\LTcaptype{none} % do not increment counter
\begin{longtable}[]{@{}
  >{\raggedright\arraybackslash}p{(\linewidth - 2\tabcolsep) * \real{0.5000}}
  >{\raggedright\arraybackslash}p{(\linewidth - 2\tabcolsep) * \real{0.5000}}@{}}
\toprule\noalign{}
\begin{minipage}[b]{\linewidth}\raggedright
Evidence check
\end{minipage} & \begin{minipage}[b]{\linewidth}\raggedright
Student action
\end{minipage} \\
\midrule\noalign{}
\endhead
\bottomrule\noalign{}
\endlastfoot
Control logic & Mark the comparison that functions as the baseline,
negative control, or reference case. \\
Uncertainty & Circle the row, card, diagram feature, or model assumption
most likely to change the conclusion. \\
Equity and consent & Mark any genetic information that could reveal risk
to relatives, change insurance or employment concerns, or require
counseling before disclosure. \\
Model comparison & State whether a simpler rule, null model, or
alternative mechanism could explain the same pattern. \\
Decision threshold & Write the minimum evidence that would make you
revise the interpretation. \\
Reproducibility & Record the exact scoring rule another group would need
to reproduce your classification. \\
\end{longtable}
}

Focus note: in this lab the cross results and offspring counts are the
evidence --- use a chi-square goodness-of-fit test against the expected
ratio before deciding whether a deviation signals linkage, penetrance,
or sampling noise. Keep required work paper-based; any material-handling
or equipment version belongs only in an optional extension.

\subsection*{Analysis Questions}\label{analysis-questions-18}
\addcontentsline{toc}{subsection}{Analysis Questions}

\begin{enumerate}
\def\labelenumi{\arabic{enumi}.}
\item
  In the AaBb \texorpdfstring{\ensuremath{\times}}{x} AaBb dihybrid cross, what fraction of offspring are aabb
  (homozygous recessive for both traits)? Show your Punnett square
  logic.
\item
  The pedigree shows two unaffected parents who produce an affected
  child. What does this tell you about dominance? Can you determine
  whether the trait is autosomal or X-linked from this generation alone?
  What additional information would resolve this?
\item
  Your chi-square test: did the observed ratio significantly deviate
  from 3:1? What biological reasons could cause deviation from the
  expected ratio (give at least two sources of biological deviation
  other than statistical sampling error)?
\item
  Blood types in humans follow codominance (A and B antigens) and simple
  recessive inheritance (O). A type AB mother and type O father have two
  children --- one type A, one type B. What genotype does each family
  member have? Is this consistent with Mendelian inheritance?
\item
  A gene controls tail length in cats: T (tailless, Manx) is dominant
  but homozygous TT is lethal. Set up the cross Tt \texorpdfstring{\ensuremath{\times}}{x} Tt. What are the
  observed genotype and phenotype ratios among the LIVING offspring? How
  does this illustrate the concept of a lethal allele modifying a
  Mendelian ratio?
\end{enumerate}

\subsection*{Extension Analysis
Questions}\label{extension-analysis-questions-8}
\addcontentsline{toc}{subsection}{Extension Analysis Questions}

\begin{enumerate}
\def\labelenumi{\arabic{enumi}.}
\setcounter{enumi}{5}
\item
  \textbf{Statistical Analysis}: If you performed three replicates of a
  genetic cross, how would you analyze the variation in offspring
  ratios? What statistical test would you use to determine if
  differences between crosses are significant?
\item
  \textbf{Experimental Error Analysis}: What are three potential sources
  of error in Mendelian inheritance experiments? How could you modify
  the procedure to reduce each source of error?
\item
  \textbf{Experimental Design}: If you wanted to test whether two genes
  are linked or assort independently, how would you design the
  experiment? What data would you collect? How would you analyze it?
\end{enumerate}

\subsection*{Group Project Extension
(Multi-Session)}\label{group-project-extension-multi-session-5}
\addcontentsline{toc}{subsection}{Group Project Extension
(Multi-Session)}

\textbf{Design a Population Genetics Study of a Human Trait}

This project will span 2-3 lab sessions:

\textbf{Session 1}: Research Question and Hypothesis - Choose a human
Mendelian trait (e.g., attached earlobes, tongue rolling, hitchhiker's
thumb) - Formulate a hypothesis about its inheritance pattern in your
population - Design a survey to collect pedigree data from families

\textbf{Session 2}: Data Collection and Analysis - Collect pedigree data
from volunteers (with consent) - Analyze inheritance patterns using
Punnett squares and chi-square tests - Compare your results to expected
Mendelian ratios

\textbf{Session 3}: Data Interpretation and Presentation - Create a
report with your findings - Discuss any deviations from expected ratios
and possible explanations - Create a scientific poster or presentation

\subsection*{Real-World Problem Solving: Genetic
Counseling}\label{real-world-problem-solving-genetic-counseling}
\addcontentsline{toc}{subsection}{Real-World Problem Solving: Genetic
Counseling}

\textbf{Case Study: Carrier Screening for Cystic Fibrosis}

\begin{enumerate}
\def\labelenumi{\arabic{enumi}.}
\item
  \textbf{Research Task}: Investigate the prevalence of cystic fibrosis
  carrier status in different populations. Why is the carrier rate
  higher in people of Northern European descent?
\item
  \textbf{Ethical Consideration}: Should carrier screening for recessive
  disorders be routinely offered to prospective parents in
  high-prevalence or high-risk contexts? What are the potential
  benefits, risks, consent requirements, and equity concerns?
\item
  \textbf{Policy Proposal}: Design a policy for responsible use of
  genetic screening in reproductive healthcare. Consider: access,
  counseling, privacy, data retention, secondary use, family
  implications, and follow-up options.
\end{enumerate}

\subsection*{Safety and Ethics Notes}\label{safety-and-ethics-notes-18}
\addcontentsline{toc}{subsection}{Safety and Ethics Notes}

This lab involves primarily paper-based activities and discussions. When
discussing patient genetic data, maintain confidentiality and respect
for individual privacy. Treat pedigrees and PRS examples as potentially
identifiable family information: do not infer ancestry, parentage, or
disease status beyond the evidence given, and separate biological risk
estimates from policy decisions about disclosure, insurance, employment,
or family notification.

\subsection*{Debrief and Reflection}\label{debrief-and-reflection-15}
\addcontentsline{toc}{subsection}{Debrief and Reflection}

After you finish the practical work, spend 5--10 minutes in your small
group comparing results and discussing the following prompts. Each
member should contribute at least one observation before moving to the
next prompt:

\begin{enumerate}
\def\labelenumi{\arabic{enumi}.}
\item
  \textbf{What did your measurements show} --- compare the group's
  results to the textbook's predictions. Where they diverge, suggest at
  least one mechanistic explanation before concluding ``experimental
  error.''
\item
  \textbf{What would change the outcome} --- propose one modification to
  the procedure that would sharpen the measurement or extend the result
  to a new biological context, and predict what you would observe.
\item
  \textbf{One-sentence headline} --- each student composes a single
  sentence summarising the lab's take-home message, suitable for a
  tweet. Compare sentences across groups; good headlines are short,
  quantitative, and mechanistic.
\item
  \textbf{Connection back to the textbook} --- identify one section of
  \cref{sec:unit_V_mendelian_genetics} that your data either confirmed
  or complicated. Cite the specific passage.
\item
  \textbf{Experimental Design Reflection}: What was the strongest aspect
  of your experimental design? What would you change if you could repeat
  the experiment?
\end{enumerate}

\subsection*{Further Reading (Lab)}\label{further-reading-lab-15}
\addcontentsline{toc}{subsection}{Further Reading (Lab)}

\begin{itemize}
\tightlist
\item
  Revisit the parent chapter \cref{sec:unit_V_mendelian_genetics} for
  the theoretical foundations on which this lab is built.
\item
  Look up any \textbf{bolded glossary term} introduced in the textbook
  chapter (each has a \texttt{\#gl:term-slug} link in the text) --- its
  master definition is in \breaktt{manuscript/glossary.md}.
\item
  Explore the Python code in \breaktt{src/biology/genetics/genetics.py}
  for Mendelian genetics calculations.
\item
  Use the self-contained Part 2 snippet as the computational template;
  it runs against tested project modules without external notebooks or
  CSV files.
\end{itemize}

\newpage

\section*{Lab 17 --- Chromosomal Inheritance and
Linkage}\label{lab-17-chromosomal-inheritance-and-linkage}
\addcontentsline{toc}{section}{Lab 17 --- Chromosomal Inheritance and
Linkage}

\label{sec:lab_unit_V_chromosomal_inheritance}

\emph{This activity accompanies
\cref{sec:unit_V_chromosomal_inheritance} of the textbook --- review
that chapter before attempting the exercises below.}

\subsection*{Learning Objectives}\label{learning-objectives-57}
\addcontentsline{toc}{subsection}{Learning Objectives}

\begin{itemize}
\tightlist
\item
  Design and evaluate genetic linkage mapping datasets to determine gene
  order
\item
  Formulate testable hypotheses about chromosomal abnormalities from
  karyotype data
\item
  Identify independent, dependent, and controlled variables in
  cytogenetic studies
\item
  Analyze experimental data using recombination frequencies and
  statistical methods
\item
  Write a brief scientific report with hypothesis, methods, results, and
  conclusion
\end{itemize}

\subsubsection*{Alignment and Rubric
Map}\label{alignment-and-rubric-map-19}
\addcontentsline{toc}{subsubsection}{Alignment and Rubric Map}

\begin{itemize}
\tightlist
\item
  \textbf{Outcome 1 (LO1):} Interpret the supplied evidence or model
  output for Chromosomal Inheritance and Linkage.
\item
  \textbf{Outcome 2 (LO2):} Identify controls and comparison groups that
  make the claim testable.
\item
  \textbf{Outcome 3 (LO3):} Quantify uncertainty, boundary conditions,
  or alternative explanations before concluding.
\item
  \textbf{Outcome 4 (LO4):} Transfer the mechanism to a new biological
  case or public-facing decision.
\item
  \textbf{Chapter LO coverage:} LO1, LO2, LO3, LO4
\item
  \textbf{Rubric dimensions:} evidence; controls; uncertainty;
  mechanism; transfer. \#\# Pre-Lab Inquiry Questions \{.unnumbered\}
\end{itemize}

\emph{Complete these questions before coming to lab. Use your textbook
and additional research to inform your hypotheses.}

\begin{enumerate}
\def\labelenumi{\arabic{enumi}.}
\item
  \textbf{Hypothesis Formation}: Based on your understanding of genetic
  linkage, predict how the recombination frequency between two genes
  would change if they were located very close together on the same
  chromosome versus far apart.
\item
  \textbf{Experimental Design}: What is the independent variable in a
  linkage mapping experiment? What is the dependent variable? List at
  least 5 variables that should be controlled to ensure accurate
  mapping.
\item
  \textbf{Scientific Context}: The chi-square test can be used to
  determine if observed recombination frequencies differ significantly
  from expected. Why is this test important in linkage analysis? What
  does a significant deviation tell you?
\item
  \textbf{Quantitative Reasoning}: If genes A and B have a recombination
  frequency of 12 cM, and genes B and C have 8 cM, what is the expected
  recombination frequency between A and C if there is no interference?
  How does interference affect this?
\item
  \textbf{Real-World Application}: Chromosomal abnormalities like Down
  syndrome are often caused by nondisjunction. Why does maternal age
  increase the risk of nondisjunction? What cellular mechanisms normally
  prevent this?
\end{enumerate}

\subsection*{Background}\label{background-9}
\addcontentsline{toc}{subsection}{Background}

Genes on the same chromosome are linked and tend to be inherited
together (violating independent assortment). The \textbf{recombination
frequency} between linked genes --- measured as the percentage of
recombinant offspring among total --- is proportional to map distance
(1\% recombination = 1 centimorgan, cM). X-linked traits are carried on
the X chromosome; since males are hemizygous (XY), a single recessive
allele is expressed. Nondisjunction during meiosis (failure of
chromosomes to separate) produces aneuploid gametes, leading to
conditions such as trisomy 21 (Down syndrome).

\subsection*{Paper-Based Materials}\label{paper-based-materials-19}
\addcontentsline{toc}{subsection}{Paper-Based Materials}

{\def\LTcaptype{none} % do not increment counter
\begin{longtable}[]{@{}
  >{\raggedright\arraybackslash}p{(\linewidth - 2\tabcolsep) * \real{0.3333}}
  >{\raggedright\arraybackslash}p{(\linewidth - 2\tabcolsep) * \real{0.6667}}@{}}
\toprule\noalign{}
\begin{minipage}[b]{\linewidth}\raggedright
Item
\end{minipage} & \begin{minipage}[b]{\linewidth}\raggedright
Quantity
\end{minipage} \\
\midrule\noalign{}
\endhead
\bottomrule\noalign{}
\endlastfoot
Linkage mapping dataset (3 genes, pairwise recombination frequencies)
(printed) & 1 \\
Karyotype images (4 cases: normal, trisomy 21, Turner syndrome 45,X,
translocation) (printed) & 1 set \\
Scissors and paste (for karyotype sorting exercise) & 1 set \\
X-linked pedigrees (2 cases: colour blindness, haemophilia) (printed) &
1 \\
Ruler & 1 \\
Chromosome pair cards (22 homologous pairs, numbered) & 1 set \\
Calculator or optional Python REPL with this project installed & 1 \\
Statistical analysis worksheet & 1 \\
\end{longtable}
}

\subsection*{Paper-Based
Investigation}\label{paper-based-investigation-19}
\addcontentsline{toc}{subsection}{Paper-Based Investigation}

\subsubsection*{Part 1: Genetic Linkage
Mapping}\label{part-1-genetic-linkage-mapping}
\addcontentsline{toc}{subsubsection}{Part 1: Genetic Linkage Mapping}

\begin{enumerate}
\def\labelenumi{\arabic{enumi}.}
\item
  \textbf{Formulate Your Hypothesis}: Based on your pre-lab research,
  write a hypothesis about the order of three genes given their pairwise
  recombination frequencies.
\item
  \textbf{Identify Variables}:

  \begin{itemize}
  \tightlist
  \item
    Independent variable: Gene order (which gene is in the middle)
  \item
    Dependent variable: Recombination frequencies between gene pairs
  \item
    Controlled variables: Population size, crossover detection method,
    etc.
  \end{itemize}
\item
  \textbf{Set Up Controls}: Why is it important to have a large sample
  size in linkage mapping? What other controls could you include to
  ensure accurate recombination frequency measurements?
\item
  \textbf{Data Collection}:

  \begin{itemize}
  \tightlist
  \item
    Given pairwise recombination frequencies: genes A--B = 12 cM, B--C =
    8 cM, A--C = 20 cM.
  \item
    Draw the linear linkage map. Which gene is in the middle? Indicate
    the distances between each adjacent pair.
  \item
    Calculate the expected coefficient of coincidence (observed double
    crossovers / expected double crossovers) and interference (1 -
    coefficient of coincidence).
  \end{itemize}
\end{enumerate}

\subsubsection*{Part 2: Computational Biology Exercise - Linkage
Analysis with
Python}\label{part-2-computational-biology-exercise---linkage-analysis-with-python}
\addcontentsline{toc}{subsubsection}{Part 2: Computational Biology
Exercise - Linkage Analysis with Python}

\emph{Optional computational check: run this self-contained Python
snippet from the project root. It uses tested \texttt{src/biology}
modules and requires no external notebook or CSV file.}

Infer a simple three-point map from pairwise recombination distances.

\begin{Shaded}
\begin{Highlighting}[]
\ImportTok{from}\NormalTok{ biology.genetics }\ImportTok{import}\NormalTok{ genetic\_distance, infer\_three\_point\_order}

\NormalTok{distances }\OperatorTok{=}\NormalTok{ \{(}\StringTok{"A"}\NormalTok{, }\StringTok{"B"}\NormalTok{): }\FloatTok{12.0}\NormalTok{, (}\StringTok{"B"}\NormalTok{, }\StringTok{"C"}\NormalTok{): }\FloatTok{8.0}\NormalTok{, (}\StringTok{"A"}\NormalTok{, }\StringTok{"C"}\NormalTok{): }\FloatTok{20.0}\NormalTok{\}}
\NormalTok{order }\OperatorTok{=}\NormalTok{ infer\_three\_point\_order(distances)}

\BuiltInTok{print}\NormalTok{(}\StringTok{"gene order:"}\NormalTok{, }\StringTok{" {-} "}\NormalTok{.join(order.order))}
\BuiltInTok{print}\NormalTok{(}\StringTok{"adjacent distances:"}\NormalTok{, order.adjacent\_distances\_cM)}
\BuiltInTok{print}\NormalTok{(}\StringTok{"24 recombinants among 200 progeny:"}\NormalTok{, genetic\_distance(}\DecValTok{24}\NormalTok{, }\DecValTok{200}\NormalTok{), }\StringTok{"cM"}\NormalTok{)}
\end{Highlighting}
\end{Shaded}

\subsubsection*{Part 3: Karyotype
Analysis}\label{part-3-karyotype-analysis}
\addcontentsline{toc}{subsubsection}{Part 3: Karyotype Analysis}

\begin{enumerate}
\def\labelenumi{\arabic{enumi}.}
\setcounter{enumi}{4}
\item
  \textbf{Design an Alternative Investigation}: Instead of just
  analyzing provided karyotype images, design a paper-based case-control
  study using archived karyotype datasets to investigate the effect of a
  spindle-disrupting condition on chromosome segregation. What
  hypothesis would you test? How would you analyze the resulting
  karyotype calls?
\item
  \textbf{Hypothesis}: Predict how exposure to a microtubule-disrupting
  agent (like colchicine) would affect chromosome number and structure
  in dividing cells.
\end{enumerate}

\subsection*{Data Recording}\label{data-recording-16}
\addcontentsline{toc}{subsection}{Data Recording}

Linkage map sketch:

\texttt{A\ —{[}\_\_\_cM{]}—\ B\ —{[}\_\_\_cM{]}—\ C}

{\def\LTcaptype{none} % do not increment counter
\begin{longtable}[]{@{}llll@{}}
\toprule\noalign{}
Karyotype Case & \# chromosomes & Abnormality & Syndrome (if any) \\
\midrule\noalign{}
\endhead
\bottomrule\noalign{}
\endlastfoot
1 & & & \\
2 & & & \\
3 & & & \\
4 & & & \\
\end{longtable}
}

X-linked pedigree genotypes:

{\def\LTcaptype{none} % do not increment counter
\begin{longtable}[]{@{}lll@{}}
\toprule\noalign{}
Individual & Phenotype & Genotype \\
\midrule\noalign{}
\endhead
\bottomrule\noalign{}
\endlastfoot
Gen I male & Unaffected & X\^{}B Y \\
Gen I female & & \\
Gen II affected male & & \\
Gen III carrier female & & \\
\end{longtable}
}

\subsection*{Evidence and Reproducibility
Checklist}\label{evidence-and-reproducibility-checklist-19}
\addcontentsline{toc}{subsection}{Evidence and Reproducibility
Checklist}

\begin{itemize}
\tightlist
\item
  \textbf{Primary evidence goal:} Map linked genes and reason from
  chromosomal anomalies.
\item
  \textbf{Data skill to practice:} Infer gene order or chromosomal
  mechanism from offspring counts.
\item
  \textbf{BioSkills emphasis:} Quantitative reasoning, Process of
  science.
\item
  \textbf{Control logic:} identify at least one positive control, one
  negative control, or one baseline comparison before interpreting
  results.
\item
  \textbf{Measurement discipline:} record units, uncertainty, sample
  size, and any discarded observation with a reason.
\item
  \textbf{Mechanistic link:} connect one result directly to the parent
  chapter's big idea before writing the conclusion.
\item
  \textbf{Reproducibility check:} state one procedural detail that
  another group would need in order to reproduce the result.
\end{itemize}

\subsection*{Paper-Based Evidence
Upgrade}\label{paper-based-evidence-upgrade-19}
\addcontentsline{toc}{subsection}{Paper-Based Evidence Upgrade}

Before answering the analysis questions, annotate the paper dataset for
\textbf{Chromosomal Inheritance and Linkage} with a reproducibility
pass:

{\def\LTcaptype{none} % do not increment counter
\begin{longtable}[]{@{}
  >{\raggedright\arraybackslash}p{(\linewidth - 2\tabcolsep) * \real{0.5000}}
  >{\raggedright\arraybackslash}p{(\linewidth - 2\tabcolsep) * \real{0.5000}}@{}}
\toprule\noalign{}
\begin{minipage}[b]{\linewidth}\raggedright
Evidence check
\end{minipage} & \begin{minipage}[b]{\linewidth}\raggedright
Student action
\end{minipage} \\
\midrule\noalign{}
\endhead
\bottomrule\noalign{}
\endlastfoot
Control logic & Mark the comparison that functions as the baseline,
negative control, or reference case. \\
Uncertainty & Circle the row, card, diagram feature, or model assumption
most likely to change the conclusion. \\
Model comparison & State whether a simpler rule, null model, or
alternative mechanism could explain the same pattern. \\
Decision threshold & Write the minimum evidence that would make you
revise the interpretation. \\
Reproducibility & Record the exact scoring rule another group would need
to reproduce your classification. \\
\end{longtable}
}

Focus note: in this lab the paper karyotypes, three-point cross data,
and pedigrees are the evidence --- score recombination frequencies, gene
order, and nondisjunction division from the dataset itself rather than
from the expected ratio. Keep required work paper-based; any
material-handling or equipment version belongs only in an optional
extension.

\subsection*{Analysis Questions}\label{analysis-questions-19}
\addcontentsline{toc}{subsection}{Analysis Questions}

\begin{enumerate}
\def\labelenumi{\arabic{enumi}.}
\item
  In linkage mapping, why does the A--C recombination frequency (20 cM)
  not simply equal A--B + B--C (12 + 8 = 20 cM) in real experiments?
  Explain the concept of double crossovers and interference.
\item
  A woman is a carrier of haemophilia A (X\^{}H X\^{}h) and is also
  heterozygous for an autosomal trait (Aa). She marries a normal man
  (X\^{}H Y; AA). Calculate the probability their first-born son is
  affected with haemophilia. Show most steps using the product rule.
\item
  Karyotype Case 3 showed trisomy 21. During which meiotic division is
  nondisjunction most commonly occurring in cases of Down syndrome, and
  what is the evidence supporting this conclusion (hint: consider
  maternal age effects and the stage vulnerable to cohesin loss)?
\item
  An individual with Turner syndrome (45,X) is phenotypically female.
  Explain why the presence of a single X chromosome leads to infertility
  and why males (46,XY) with a single X chromosome are not similarly
  affected.
\item
  Philadelphia chromosome (a translocation between chromosomes 9 and 22,
  creating BCR-ABL fusion oncogene) causes chronic myeloid leukaemia
  (CML). Explain: (a) why a translocation creates a fusion gene; (b)
  what the BCR-ABL kinase does biochemically; (c) why imatinib (a
  tyrosine kinase inhibitor) is effective.
\end{enumerate}

\subsection*{Extension Analysis
Questions}\label{extension-analysis-questions-9}
\addcontentsline{toc}{subsection}{Extension Analysis Questions}

\begin{enumerate}
\def\labelenumi{\arabic{enumi}.}
\setcounter{enumi}{5}
\item
  \textbf{Statistical Analysis}: If you performed three replicates of a
  linkage mapping experiment, how would you analyze the variation in
  recombination frequencies? What statistical test would you use to
  determine if differences between gene pairs are significant?
\item
  \textbf{Experimental Error Analysis}: What are three potential sources
  of error in karyotype analysis? How could you modify the procedure to
  reduce each source of error?
\item
  \textbf{Experimental Design}: If you wanted to test whether two genes
  are linked or assort independently, how would you design the
  experiment? What data would you collect? How would you analyze it?
\end{enumerate}

\subsection*{Group Project Extension
(Multi-Session)}\label{group-project-extension-multi-session-6}
\addcontentsline{toc}{subsection}{Group Project Extension
(Multi-Session)}

\textbf{Design a Cytogenetic Study of Environmental Mutagens}

This project will span 2-3 lab sessions:

\textbf{Session 1}: Research Question and Hypothesis - Choose an
environmental mutagen (e.g., tobacco smoke, pesticides, radiation) -
Formulate a hypothesis about its effect on chromosome structure or
number - Design an experiment using model organisms (e.g., onion root
tip, fruit fly)

\textbf{Session 2}: Data Collection and Analysis - Expose organisms to
different concentrations of mutagen - Prepare and analyze chromosome
spreads - Quantify chromosome abnormalities (breaks, fragments,
aneuploidy)

\textbf{Session 3}: Data Interpretation and Presentation - Compare
abnormality rates between control and treated groups - Calculate
statistical significance - Create a scientific poster or presentation

\subsection*{Real-World Problem Solving: Genetic
Counseling}\label{real-world-problem-solving-genetic-counseling-1}
\addcontentsline{toc}{subsection}{Real-World Problem Solving: Genetic
Counseling}

\textbf{Case Study: Prenatal Screening for Chromosomal Abnormalities}

\begin{enumerate}
\def\labelenumi{\arabic{enumi}.}
\item
  \textbf{Research Task}: Investigate the different methods for prenatal
  screening and diagnosis (ultrasound, serum markers, NIPT,
  amniocentesis). What are the advantages and limitations of each?
\item
  \textbf{Ethical Consideration}: Should most pregnant women be offered
  prenatal screening for chromosomal abnormalities? What are the
  potential psychological and social implications of such testing?
\item
  \textbf{Policy Proposal}: Design a policy for responsible use of
  prenatal genetic testing. Consider: informed consent, access,
  counseling, and follow-up options.
\end{enumerate}

\subsection*{Safety and Ethics Notes}\label{safety-and-ethics-notes-19}
\addcontentsline{toc}{subsection}{Safety and Ethics Notes}

No chemical hazards. Karyotype analysis using real patient images: most
should be de-identified. Discussion of chromosomal conditions should use
person-first language and be conducted respectfully. Consider the
ethical implications of genetic testing and counseling.

\subsection*{Debrief and Reflection}\label{debrief-and-reflection-16}
\addcontentsline{toc}{subsection}{Debrief and Reflection}

After you finish the practical work, spend 5--10 minutes in your small
group comparing results and discussing the following prompts. Each
member should contribute at least one observation before moving to the
next prompt:

\begin{enumerate}
\def\labelenumi{\arabic{enumi}.}
\item
  \textbf{What did your measurements show} --- compare the group's
  results to the textbook's predictions. Where they diverge, suggest at
  least one mechanistic explanation before concluding ``experimental
  error.''
\item
  \textbf{What would change the outcome} --- propose one modification to
  the procedure that would sharpen the measurement or extend the result
  to a new biological context, and predict what you would observe.
\item
  \textbf{One-sentence headline} --- each student composes a single
  sentence summarising the lab's take-home message, suitable for a
  tweet. Compare sentences across groups; good headlines are short,
  quantitative, and mechanistic.
\item
  \textbf{Connection back to the textbook} --- identify one section of
  \cref{sec:unit_V_chromosomal_inheritance} that your data either
  confirmed or complicated. Cite the specific passage.
\item
  \textbf{Experimental Design Reflection}: What was the strongest aspect
  of your experimental design? What would you change if you could repeat
  the experiment?
\end{enumerate}

\subsection*{Further Reading (Lab)}\label{further-reading-lab-16}
\addcontentsline{toc}{subsection}{Further Reading (Lab)}

\begin{itemize}
\tightlist
\item
  Revisit the parent chapter \cref{sec:unit_V_chromosomal_inheritance}
  for the theoretical foundations on which this lab is built.
\item
  Look up any \textbf{bolded glossary term} introduced in the textbook
  chapter (each has a \texttt{\#gl:term-slug} link in the text) --- its
  master definition is in \breaktt{manuscript/glossary.md}.
\item
  Explore the Python code in \breaktt{src/biology/genetics/genetics.py}
  for chromosomal inheritance calculations.
\item
  Use the self-contained Part 2 snippet as the computational template;
  it runs against tested project modules without external notebooks or
  CSV files.
\end{itemize}

\newpage

\section*{Lab 18 --- Population Genetics and Hardy-Weinberg
Equilibrium}\label{lab-18-population-genetics-and-hardy-weinberg-equilibrium}
\addcontentsline{toc}{section}{Lab 18 --- Population Genetics and
Hardy-Weinberg Equilibrium}

\label{sec:lab_unit_V_population_genetics}

\subsection*{Learning Objectives}\label{learning-objectives-58}
\addcontentsline{toc}{subsection}{Learning Objectives}

\emph{This activity accompanies \cref{sec:unit_V_population_genetics} of
the textbook --- review that chapter before attempting the exercises
below.}

\begin{itemize}
\tightlist
\item
  Apply Hardy-Weinberg equations to calculate allele and genotype
  frequencies from population data
\item
  Test for deviations from HWE using chi-square analysis
\item
  Simulate the effect of genetic drift on small vs large populations
  using a bead model
\item
  Evaluate the effect of natural selection on allele frequencies across
  generations
\end{itemize}

\subsubsection*{Alignment and Rubric
Map}\label{alignment-and-rubric-map-20}
\addcontentsline{toc}{subsubsection}{Alignment and Rubric Map}

\begin{itemize}
\tightlist
\item
  \textbf{Outcome 1 (LO1):} Interpret the supplied evidence or model
  output for Population Genetics and Hardy-Weinberg Equilibrium.
\item
  \textbf{Outcome 2 (LO2):} Identify controls and comparison groups that
  make the claim testable.
\item
  \textbf{Outcome 3 (LO3):} Quantify uncertainty, boundary conditions,
  or alternative explanations before concluding.
\item
  \textbf{Outcome 4 (LO4):} Transfer the mechanism to a new biological
  case or public-facing decision.
\item
  \textbf{Chapter LO coverage:} LO1, LO2, LO3, LO4
\item
  \textbf{Rubric dimensions:} evidence; controls; uncertainty;
  mechanism; transfer. \#\# Background \{.unnumbered\}
\end{itemize}

The Hardy-Weinberg principle states that in a large, randomly mating
population with no selection, mutation, migration, or drift, allele
frequencies (p, q) and genotype frequencies (p\texorpdfstring{\textsuperscript{2}}{2}, 2pq, q\texorpdfstring{\textsuperscript{2}}{2}) remain
constant across generations. Deviations from HWE indicate one or more of
these forces is acting. Genetic drift --- random changes in allele
frequency --- is most powerful in small populations, causing loss of
variation and potentially fixing or eliminating alleles by chance alone.

\subsection*{Paper-Based Materials}\label{paper-based-materials-20}
\addcontentsline{toc}{subsection}{Paper-Based Materials}

{\def\LTcaptype{none} % do not increment counter
\begin{longtable}[]{@{}
  >{\raggedright\arraybackslash}p{(\linewidth - 2\tabcolsep) * \real{0.3333}}
  >{\raggedright\arraybackslash}p{(\linewidth - 2\tabcolsep) * \real{0.6667}}@{}}
\toprule\noalign{}
\begin{minipage}[b]{\linewidth}\raggedright
Item
\end{minipage} & \begin{minipage}[b]{\linewidth}\raggedright
Quantity
\end{minipage} \\
\midrule\noalign{}
\endhead
\bottomrule\noalign{}
\endlastfoot
Two colours of beads (representing A and a alleles): 50 of each colour &
100 total \\
Opaque bag or cup (for random sampling) & 2 \\
Population genetics data table (printed: real MN blood group data) &
1 \\
Calculator & 1 \\
Chi-square critical value table & 1 \\
Graph paper & 1 \\
\end{longtable}
}

\subsection*{Paper-Based
Investigation}\label{paper-based-investigation-20}
\addcontentsline{toc}{subsection}{Paper-Based Investigation}

\textbf{Part A --- Hardy-Weinberg Calculation}

\begin{enumerate}
\def\labelenumi{\arabic{enumi}.}
\tightlist
\item
  Using the provided MN blood group data (genotype counts \textbf{MM =
  298, MN = 489, NN = 213}, from a sample of \emph{n} = 1000; classic
  combined US-Caucasian MN blood-group survey, as reproduced in standard
  population-genetics texts): calculate allele frequencies \emph{p} (M)
  and \emph{q} (N). \emph{(Check: p(M) \texorpdfstring{\ensuremath{\approx}}{~} 0.5425, q(N) \texorpdfstring{\ensuremath{\approx}}{~} 0.4575.)}
\item
  Calculate expected genotype frequencies under HWE (p\texorpdfstring{\textsuperscript{2}}{2}, 2pq, q\texorpdfstring{\textsuperscript{2}}{2}) and
  expected counts.
\item
  Perform chi-square test to determine whether the population is in HWE.
\end{enumerate}

\textbf{Part B --- Genetic Drift Simulation}

\begin{enumerate}
\def\labelenumi{\arabic{enumi}.}
\setcounter{enumi}{3}
\tightlist
\item
  Place 50 red (A) and 50 blue (a) beads in the bag: this represents a
  population of 100 with p = q = 0.5.
\item
  Without looking, draw 10 beads (simulating a bottleneck to N=10).
  Record the frequency of red (A) beads drawn. Then \textbf{discard the
  old bag and rebuild a fresh 100-bead pool at the drawn allele ratio}:
  multiply the observed red fraction by 100, round to the nearest whole
  bead, and use that many red beads with the remainder blue (e.g.~6/10
  red \texorpdfstring{\ensuremath{\to}}{->} 60 red + 40 blue). Repeat the draw--rebuild cycle for 5
  generations, recording \emph{p} (red fraction) each generation.
\item
  A second group runs the same simulation with N=50 (draw 50 beads, then
  rebuild a 100-bead pool at the drawn allele ratio by the same rounding
  rule) for 5 generations.
\end{enumerate}

\textbf{Part C --- Selection Simulation (Optional)}

\begin{enumerate}
\def\labelenumi{\arabic{enumi}.}
\setcounter{enumi}{6}
\tightlist
\item
  In generation 1 of the drift simulation, designate aa (2 blue drawn)
  as lethal --- remove any drawn aa pairs and resample. Track p over 3
  generations.
\end{enumerate}

\subsection*{Data Recording}\label{data-recording-17}
\addcontentsline{toc}{subsection}{Data Recording}

MN blood group HWE test (observed counts from the \emph{n} = 1000
sample: MM = 298, MN = 489, NN = 213; source: classic combined
US-Caucasian MN blood-group survey reproduced in standard
population-genetics texts):

{\def\LTcaptype{none} % do not increment counter
\begin{longtable}[]{@{}lllll@{}}
\toprule\noalign{}
Genotype & Observed & p or q & Expected & (O\texorpdfstring{\ensuremath{-}}{-}E)\texorpdfstring{\textsuperscript{2}}{2}/E \\
\midrule\noalign{}
\endhead
\bottomrule\noalign{}
\endlastfoot
MM & 298 & p\texorpdfstring{\textsuperscript{2}}{2} = & & \\
NN & 213 & q\texorpdfstring{\textsuperscript{2}}{2} = & & \\
MN & 489 & 2pq = & & \\
\textbf{Total} & \textbf{1000} & --- & \textbf{1000} & \texorpdfstring{\ensuremath{\chi}}{chi}\texorpdfstring{\textsuperscript{2}}{2} = \_\_\_ \\
\end{longtable}
}

Fill the \emph{Expected} column using your calculated \emph{p} and
\emph{q} (Part A), then \texorpdfstring{\ensuremath{\chi}}{chi}\texorpdfstring{\textsuperscript{2}}{2} = \texorpdfstring{\ensuremath{\Sigma}}{Sigma} (O\texorpdfstring{\ensuremath{-}}{-}E)\texorpdfstring{\textsuperscript{2}}{2}/E.

\texorpdfstring{\ensuremath{\chi}}{chi}\texorpdfstring{\textsuperscript{2}}{2} = \_\_\_; df = 1; HWE at \texorpdfstring{\ensuremath{\alpha}}{alpha} = 0.05? (Y/N) --- compare \texorpdfstring{\ensuremath{\chi}}{chi}\texorpdfstring{\textsuperscript{2}}{2} to the
critical value 3.841.

Drift simulation:

{\def\LTcaptype{none} % do not increment counter
\begin{longtable}[]{@{}lll@{}}
\toprule\noalign{}
Generation & p (N=10 group) & p (N=50 group) \\
\midrule\noalign{}
\endhead
\bottomrule\noalign{}
\endlastfoot
0 & 0.5 & 0.5 \\
1 & & \\
2 & & \\
3 & & \\
4 & & \\
5 & & \\
\end{longtable}
}

\subsection*{Evidence and Reproducibility
Checklist}\label{evidence-and-reproducibility-checklist-20}
\addcontentsline{toc}{subsection}{Evidence and Reproducibility
Checklist}

\begin{itemize}
\tightlist
\item
  \textbf{Primary evidence goal:} Test equilibrium and model
  allele-frequency change.
\item
  \textbf{Data skill to practice:} Estimate genotype or allele
  frequencies from population data.
\item
  \textbf{BioSkills emphasis:} Quantitative reasoning, Process of
  science.
\item
  \textbf{Control logic:} identify at least one positive control, one
  negative control, or one baseline comparison before interpreting
  results.
\item
  \textbf{Measurement discipline:} record units, uncertainty, sample
  size, and any discarded observation with a reason.
\item
  \textbf{Mechanistic link:} connect one result directly to the parent
  chapter's big idea before writing the conclusion.
\item
  \textbf{Reproducibility check:} state one procedural detail that
  another group would need in order to reproduce the result.
\end{itemize}

\subsection*{Paper-Based Evidence
Upgrade}\label{paper-based-evidence-upgrade-20}
\addcontentsline{toc}{subsection}{Paper-Based Evidence Upgrade}

Before answering the analysis questions, annotate the paper dataset for
\textbf{Population Genetics and Hardy-Weinberg Equilibrium} with a
reproducibility pass:

{\def\LTcaptype{none} % do not increment counter
\begin{longtable}[]{@{}
  >{\raggedright\arraybackslash}p{(\linewidth - 2\tabcolsep) * \real{0.5000}}
  >{\raggedright\arraybackslash}p{(\linewidth - 2\tabcolsep) * \real{0.5000}}@{}}
\toprule\noalign{}
\begin{minipage}[b]{\linewidth}\raggedright
Evidence check
\end{minipage} & \begin{minipage}[b]{\linewidth}\raggedright
Student action
\end{minipage} \\
\midrule\noalign{}
\endhead
\bottomrule\noalign{}
\endlastfoot
Control logic & Mark the comparison that functions as the baseline,
negative control, or reference case. \\
Uncertainty & Circle the row, card, diagram feature, or model assumption
most likely to change the conclusion. \\
Model comparison & State whether a simpler rule, null model, or
alternative mechanism could explain the same pattern. \\
Decision threshold & Write the minimum evidence that would make you
revise the interpretation. \\
Reproducibility & Record the exact scoring rule another group would need
to reproduce your classification. \\
\end{longtable}
}

Focus note: in this lab the simulated allele-frequency series across
generations are the evidence --- compare observed genotype counts
against the Hardy-Weinberg expectation and contrast small versus large N
to attribute change to drift, selection, or sampling. Keep required work
paper-based; any material-handling or equipment version belongs only in
an optional extension.

\subsection*{Analysis Questions}\label{analysis-questions-20}
\addcontentsline{toc}{subsection}{Analysis Questions}

\begin{enumerate}
\def\labelenumi{\arabic{enumi}.}
\tightlist
\item
  Was the MN blood group population in HWE? Which of the HWE assumptions
  might be violated in a real human population, and how would each
  violation shift allele frequencies over generations?
\item
  Was there more variation in allele frequency across generations in the
  N=10 or N=50 group? Define genetic drift in terms of sampling error
  and explain why population size affects drift magnitude.
\item
  Founder effect: a group of 20 individuals colonises a new island. By
  chance, the founding group contains no aa individuals. Even without
  selection against aa, will this condition persist? Explain using HWE
  and the concept that mutation-selection equilibrium is absent if
  frequency starts at zero.
\item
  Heterozygote advantage: sickle-cell anaemia (aa) is lethal, but
  carriers (Aa) are more resistant to malaria than AA homozygotes. With
  the heterozygote Aa as the fittest genotype, let \emph{s} be the
  selection coefficient against AA and \emph{t} the selection
  coefficient against aa (relative fitnesses 1\texorpdfstring{\ensuremath{-}}{-}\emph{s} : 1 :
  1\texorpdfstring{\ensuremath{-}}{-}\emph{t}). The stable equilibrium allele frequency of \emph{a} is
  \textbf{q̂ = s / (s + t)}. For a malarious region with \emph{s} = 0.2
  (AA disadvantage) and \emph{t} = 1 (aa lethal), compute q̂ and
  interpret why the \emph{a} allele persists in the population despite
  being lethal in homozygotes.
\item
  A conservation biologist studying a population of 25 cheetahs finds
  they are almost completely genetically homogeneous (very low
  heterozygosity). Explain two consequences of this for: (a) short-term
  disease resistance; (b) long-term adaptive capacity. What management
  intervention would restore genetic variation?
\end{enumerate}

\subsection*{Safety and Ethics Notes}\label{safety-and-ethics-notes-20}
\addcontentsline{toc}{subsection}{Safety and Ethics Notes}

No chemical hazards. Beads present a swallowing hazard --- do not use
with young children. When discussing human population genetics data
(e.g., blood groups by ethnicity), maintain scientific objectivity and
avoid discriminatory interpretation.

\subsection*{Debrief and Reflection}\label{debrief-and-reflection-17}
\addcontentsline{toc}{subsection}{Debrief and Reflection}

After you finish the practical work, spend 5--10 minutes in your small
group comparing results and discussing the following prompts. Each
member should contribute at least one observation before moving to the
next prompt:

\begin{enumerate}
\def\labelenumi{\arabic{enumi}.}
\tightlist
\item
  \textbf{What did your measurements show} --- compare the group's
  results to the textbook's predictions. Where they diverge, suggest at
  least one mechanistic explanation before concluding ``experimental
  error.''
\item
  \textbf{What would change the outcome} --- propose one modification to
  the procedure that would sharpen the measurement or extend the result
  to a new biological context, and predict what you would observe.
\item
  \textbf{One-sentence headline} --- each student composes a single
  sentence summarising the lab's take-home message, suitable for a
  tweet. Compare sentences across groups; good headlines are short,
  quantitative, and mechanistic.
\item
  \textbf{Connection back to the textbook} --- identify one section of
  \cref{sec:unit_V_population_genetics} that your data either confirmed
  or complicated. Cite the specific passage.
\end{enumerate}

\subsection*{Further Reading (Lab)}\label{further-reading-lab-17}
\addcontentsline{toc}{subsection}{Further Reading (Lab)}

\begin{itemize}
\tightlist
\item
  Revisit the parent chapter \cref{sec:unit_V_population_genetics} for
  the theoretical foundations on which this lab is built.
\item
  Look up any \textbf{bolded glossary term} introduced in the textbook
  chapter (each has a \texttt{\#gl:term-slug} link in the text) --- its
  master definition is in \breaktt{manuscript/glossary.md}.
\item
  Explore the appended \texttt{src/} module that implements the
  corresponding quantitative model (when applicable) --- referenced in
  the parent chapter's ``Bridge to Computation'' subsection.
\end{itemize}

\emph{Module footer: parent chapter
\texttt{\textbackslash{}cref\{sec:unit\_V\_population\_genetics\}}; all
numerical quantities in this lab use SI units --- see Appendix D of the
textbook for unit conversions and biological-scale reference values.}

\newpage

\section*{Lab 19 --- Evolution --- Theory, Natural Selection, and
Adaptation}\label{lab-19-evolution-theory-natural-selection-and-adaptation}
\addcontentsline{toc}{section}{Lab 19 --- Evolution --- Theory, Natural
Selection, and Adaptation}

\label{sec:lab_unit_VI_evolution_and_selection}

\emph{This activity accompanies
\cref{sec:unit_VI_evolution_and_selection} of the textbook --- review
that chapter before attempting the exercises below.}

\subsection*{Learning Objectives}\label{learning-objectives-59}
\addcontentsline{toc}{subsection}{Learning Objectives}

\begin{itemize}
\tightlist
\item
  Design and evaluate a natural selection simulation to measure
  selection coefficients
\item
  Formulate testable hypotheses about how environmental change affects
  allele frequencies
\item
  Identify independent, dependent, and controlled variables in
  evolutionary experiments
\item
  Analyze experimental data using fitness calculations and statistical
  methods
\item
  Write a brief scientific report with hypothesis, methods, results, and
  conclusion
\end{itemize}

\subsubsection*{Alignment and Rubric
Map}\label{alignment-and-rubric-map-21}
\addcontentsline{toc}{subsubsection}{Alignment and Rubric Map}

\begin{itemize}
\tightlist
\item
  \textbf{Outcome 1 (LO1):} Interpret the supplied evidence or model
  output for Evolution --- Theory, Natural Selection, and Adaptation.
\item
  \textbf{Outcome 2 (LO2):} Identify controls and comparison groups that
  make the claim testable.
\item
  \textbf{Outcome 3 (LO3):} Quantify uncertainty, boundary conditions,
  or alternative explanations before concluding.
\item
  \textbf{Outcome 4 (LO4):} Transfer the mechanism to a new biological
  case or public-facing decision.
\item
  \textbf{Chapter LO coverage:} LO1, LO2, LO3, LO4
\item
  \textbf{Rubric dimensions:} evidence; controls; uncertainty;
  mechanism; transfer. \#\# Pre-Lab Inquiry Questions \{.unnumbered\}
\end{itemize}

\emph{Complete these questions before coming to lab. Use your textbook
and additional research to inform your hypotheses.}

\begin{enumerate}
\def\labelenumi{\arabic{enumi}.}
\item
  \textbf{Hypothesis Formation}: Based on your understanding of natural
  selection, predict how a change in environment (e.g., from green to
  brown background) would affect the survival of different colored prey.
  Write a clear, testable hypothesis.
\item
  \textbf{Experimental Design}: What is the independent variable in a
  natural selection simulation? What is the dependent variable? List at
  least 5 variables that should be controlled to ensure valid results.
\item
  \textbf{Scientific Context}: The selection coefficient (s) quantifies
  the fitness cost of a genotype. Why is this measure important in
  evolutionary biology? How does it relate to the rate of evolutionary
  change?
\item
  \textbf{Quantitative Reasoning}: If a genotype has a relative fitness
  of 0.8, what is its selection coefficient? If the initial frequency of
  this genotype is 0.3, what would you expect its frequency to be after
  one generation of selection?
\item
  \textbf{Real-World Application}: Antibiotic resistance is a classic
  example of natural selection. How does the overuse of antibiotics
  create selection pressure for resistant bacteria? What can be done to
  slow this process?
\end{enumerate}

\subsection*{Background}\label{background-10}
\addcontentsline{toc}{subsection}{Background}

Natural selection operates when three conditions are met: (1) heritable
variation exists; (2) variation affects survival/reproduction (fitness);
(3) traits are inherited. Over generations, favourable alleles increase
in frequency (directional selection), multiple phenotypes are maintained
(balancing/disruptive selection), or variation is reduced around an
optimum (stabilising selection). The selection coefficient \emph{s}
quantifies fitness cost: a lethal allele has s = 1.0; a neutral allele
has s = 0.

\subsection*{Paper-Based Materials}\label{paper-based-materials-21}
\addcontentsline{toc}{subsection}{Paper-Based Materials}

{\def\LTcaptype{none} % do not increment counter
\begin{longtable}[]{@{}
  >{\raggedright\arraybackslash}p{(\linewidth - 2\tabcolsep) * \real{0.3333}}
  >{\raggedright\arraybackslash}p{(\linewidth - 2\tabcolsep) * \real{0.6667}}@{}}
\toprule\noalign{}
\begin{minipage}[b]{\linewidth}\raggedright
Item
\end{minipage} & \begin{minipage}[b]{\linewidth}\raggedright
Quantity
\end{minipage} \\
\midrule\noalign{}
\endhead
\bottomrule\noalign{}
\endlastfoot
Coloured paper squares (30 each of: green, red, brown, on green and
brown backgrounds) & 1 set \\
Forceps (one per student, simulating different beak types) & 4 \\
60-second timer & 1 \\
Ruler & 1 \\
Darwin's finch beak data table (printed: beak depth vs seed hardness
across 14 species) & 1 \\
Selection data table (survival rates by genotype in two environments) &
1 \\
AMR allele-frequency dataset across simulated antibiotic gradients &
1 \\
Resistance policy decision matrix & 1 \\
Calculator & 1 \\
Calculator or optional Python REPL with this project installed & 1 \\
Graph paper & 1 \\
\end{longtable}
}

\subsection*{Paper-Based
Investigation}\label{paper-based-investigation-21}
\addcontentsline{toc}{subsection}{Paper-Based Investigation}

\subsubsection*{Part 1: Predator-Prey Selection
Simulation}\label{part-1-predator-prey-selection-simulation}
\addcontentsline{toc}{subsubsection}{Part 1: Predator-Prey Selection
Simulation}

\begin{enumerate}
\def\labelenumi{\arabic{enumi}.}
\item
  \textbf{Formulate Your Hypothesis}: Based on your pre-lab research,
  write a hypothesis about how different colored prey will survive on
  different colored backgrounds.
\item
  \textbf{Identify Variables}:

  \begin{itemize}
  \tightlist
  \item
    Independent variable: Prey color and background color
  \item
    Dependent variable: Survival rate (\% of prey consumed or surviving)
  \item
    Controlled variables: Foraging time, number of prey, predator type,
    etc.
  \end{itemize}
\item
  \textbf{Set Up Controls}: Why is it important to have a control
  background (e.g., green prey on green background)? What other controls
  could you include in a predation simulation?
\item
  \textbf{Data Collection}:

  \begin{itemize}
  \tightlist
  \item
    Scatter 90 coloured squares randomly across a green (grassy)
    background table.
  \item
    Students ``forage'' for 30 seconds using forceps, prioritising
    whatever they can pick up fastest (competitive foraging).
  \item
    Count survivors in each colour category. Calculate \% survival per
    colour on the green background.
  \item
    Repeat on a brown paper background.
  \item
    Perform at least 3 replicates for each background type.
  \end{itemize}
\end{enumerate}

\subsubsection*{Part 2: Computational Biology Exercise - Selection
Coefficient Analysis with
Python}\label{part-2-computational-biology-exercise---selection-coefficient-analysis-with-python}
\addcontentsline{toc}{subsubsection}{Part 2: Computational Biology
Exercise - Selection Coefficient Analysis with Python}

\emph{Optional computational check: run this self-contained Python
snippet from the project root. It uses tested \texttt{src/biology}
modules and requires no external notebook or CSV file.}

Compare a hand-calculated allele-frequency change with the selection
simulator.

\begin{Shaded}
\begin{Highlighting}[]
\ImportTok{from}\NormalTok{ biology.evolution }\ImportTok{import}\NormalTok{ Population, simulate\_selection}

\NormalTok{start }\OperatorTok{=}\NormalTok{ Population(name}\OperatorTok{=}\StringTok{"demo"}\NormalTok{, p}\OperatorTok{=}\FloatTok{0.3}\NormalTok{, q}\OperatorTok{=}\FloatTok{0.7}\NormalTok{, fitness\_AA}\OperatorTok{=}\FloatTok{1.0}\NormalTok{, fitness\_Aa}\OperatorTok{=}\FloatTok{0.9}\NormalTok{, fitness\_aa}\OperatorTok{=}\FloatTok{0.6}\NormalTok{)}
\NormalTok{history }\OperatorTok{=}\NormalTok{ simulate\_selection(start, generations}\OperatorTok{=}\DecValTok{10}\NormalTok{)}

\BuiltInTok{print}\NormalTok{(}\StringTok{"starting p:"}\NormalTok{, start.p)}
\BuiltInTok{print}\NormalTok{(}\StringTok{"final p:"}\NormalTok{, }\BuiltInTok{round}\NormalTok{(history[}\OperatorTok{{-}}\DecValTok{1}\NormalTok{].p, }\DecValTok{3}\NormalTok{))}
\BuiltInTok{print}\NormalTok{(}\StringTok{"generations recorded:"}\NormalTok{, }\BuiltInTok{len}\NormalTok{(history))}
\end{Highlighting}
\end{Shaded}

\subsubsection*{Part 3: Finch Beak
Analysis}\label{part-3-finch-beak-analysis}
\addcontentsline{toc}{subsubsection}{Part 3: Finch Beak Analysis}

\begin{enumerate}
\def\labelenumi{\arabic{enumi}.}
\setcounter{enumi}{4}
\item
  \textbf{Design an Alternative Experiment}: Instead of just analyzing
  provided data, design an experiment to test how seed availability
  affects beak morphology in a bird population over multiple
  generations. What hypothesis would you test? How would you measure
  beak depth?
\item
  \textbf{Hypothesis}: Predict the relationship between seed hardness
  and average beak depth. Would you expect a linear relationship, or
  something else? Why?
\end{enumerate}

\subsection*{Data Recording}\label{data-recording-18}
\addcontentsline{toc}{subsection}{Data Recording}

{\def\LTcaptype{none} % do not increment counter
\begin{longtable}[]{@{}
  >{\raggedright\arraybackslash}p{(\linewidth - 10\tabcolsep) * \real{0.0652}}
  >{\raggedright\arraybackslash}p{(\linewidth - 10\tabcolsep) * \real{0.1087}}
  >{\raggedright\arraybackslash}p{(\linewidth - 10\tabcolsep) * \real{0.2065}}
  >{\raggedright\arraybackslash}p{(\linewidth - 10\tabcolsep) * \real{0.2065}}
  >{\raggedright\arraybackslash}p{(\linewidth - 10\tabcolsep) * \real{0.2065}}
  >{\raggedright\arraybackslash}p{(\linewidth - 10\tabcolsep) * \real{0.2065}}@{}}
\toprule\noalign{}
\begin{minipage}[b]{\linewidth}\raggedright
Colour
\end{minipage} & \begin{minipage}[b]{\linewidth}\raggedright
Starting \#
\end{minipage} & \begin{minipage}[b]{\linewidth}\raggedright
\# Surviving (green)
\end{minipage} & \begin{minipage}[b]{\linewidth}\raggedright
\% Survival (green)
\end{minipage} & \begin{minipage}[b]{\linewidth}\raggedright
\# Surviving (brown)
\end{minipage} & \begin{minipage}[b]{\linewidth}\raggedright
\% Survival (brown)
\end{minipage} \\
\midrule\noalign{}
\endhead
\bottomrule\noalign{}
\endlastfoot
Green & 30 & & & & \\
Red & 30 & & & & \\
Brown & 30 & & & & \\
\end{longtable}
}

Selection coefficients:

{\def\LTcaptype{none} % do not increment counter
\begin{longtable}[]{@{}
  >{\raggedright\arraybackslash}p{(\linewidth - 8\tabcolsep) * \real{0.1333}}
  >{\raggedright\arraybackslash}p{(\linewidth - 8\tabcolsep) * \real{0.3500}}
  >{\raggedright\arraybackslash}p{(\linewidth - 8\tabcolsep) * \real{0.0833}}
  >{\raggedright\arraybackslash}p{(\linewidth - 8\tabcolsep) * \real{0.3500}}
  >{\raggedright\arraybackslash}p{(\linewidth - 8\tabcolsep) * \real{0.0833}}@{}}
\toprule\noalign{}
\begin{minipage}[b]{\linewidth}\raggedright
Genotype
\end{minipage} & \begin{minipage}[b]{\linewidth}\raggedright
Habitat A fitness (w)
\end{minipage} & \begin{minipage}[b]{\linewidth}\raggedright
s (A)
\end{minipage} & \begin{minipage}[b]{\linewidth}\raggedright
Habitat B fitness (w)
\end{minipage} & \begin{minipage}[b]{\linewidth}\raggedright
s (B)
\end{minipage} \\
\midrule\noalign{}
\endhead
\bottomrule\noalign{}
\endlastfoot
AA & & & & \\
Aa & & & & \\
aa & & & & \\
\end{longtable}
}

Finch graph trend description:

\subsection*{Evidence and Reproducibility
Checklist}\label{evidence-and-reproducibility-checklist-21}
\addcontentsline{toc}{subsection}{Evidence and Reproducibility
Checklist}

\begin{itemize}
\tightlist
\item
  \textbf{Primary evidence goal:} Simulate selection across generations
  under alternative fitness assumptions.
\item
  \textbf{Data skill to practice:} Interpret fitness data and
  distinguish selection from other forces.
\item
  \textbf{BioSkills emphasis:} Modeling and simulation, Quantitative
  reasoning, Communication and collaboration.
\item
  \textbf{Control logic:} identify at least one positive control, one
  negative control, or one baseline comparison before interpreting
  results.
\item
  \textbf{Measurement discipline:} record units, uncertainty, sample
  size, and any discarded observation with a reason.
\item
  \textbf{Mechanistic link:} connect one result directly to the parent
  chapter's big idea before writing the conclusion.
\item
  \textbf{Reproducibility check:} state one procedural detail that
  another group would need in order to reproduce the result.
\end{itemize}

\subsection*{Paper-Based Evidence
Upgrade}\label{paper-based-evidence-upgrade-21}
\addcontentsline{toc}{subsection}{Paper-Based Evidence Upgrade}

Before answering the analysis questions, annotate the paper dataset for
\textbf{Evolution --- Theory, Natural Selection, and Adaptation} with a
reproducibility pass:

{\def\LTcaptype{none} % do not increment counter
\begin{longtable}[]{@{}
  >{\raggedright\arraybackslash}p{(\linewidth - 2\tabcolsep) * \real{0.5000}}
  >{\raggedright\arraybackslash}p{(\linewidth - 2\tabcolsep) * \real{0.5000}}@{}}
\toprule\noalign{}
\begin{minipage}[b]{\linewidth}\raggedright
Evidence check
\end{minipage} & \begin{minipage}[b]{\linewidth}\raggedright
Student action
\end{minipage} \\
\midrule\noalign{}
\endhead
\bottomrule\noalign{}
\endlastfoot
Control logic & Mark the comparison that functions as the baseline,
negative control, or reference case. \\
Uncertainty & Circle the row, card, diagram feature, or model assumption
most likely to change the conclusion. \\
Model comparison & State whether a simpler rule, null model, or
alternative mechanism could explain the same pattern. \\
Decision threshold & Write the minimum evidence that would make you
revise the interpretation. \\
Reproducibility & Record the exact scoring rule another group would need
to reproduce your classification. \\
\end{longtable}
}

Focus note: in this background-survival and finch-beak lab, treat each
generation's surviving colour or beak counts as a measured fitness
differential, and separate selection acting on heritable variation from
within-generation plasticity before claiming adaptation. Keep required
work paper-based; any material-handling or equipment version belongs
only in an optional extension.

\subsection*{Analysis Questions}\label{analysis-questions-21}
\addcontentsline{toc}{subsection}{Analysis Questions}

\begin{enumerate}
\def\labelenumi{\arabic{enumi}.}
\item
  On the green background, which colour survived best? Explain in terms
  of camouflage and predator visual system. How does this simulate
  directional selection in a real habitat?
\item
  If the green background was then ``burned'' and replaced with brown
  substrate (environmental change), predict how allele frequencies for
  ``green colour'' would change over generations. Which type of
  selection would now operate?
\item
  The finch beak data showed that species with deep beaks eat harder
  seeds. Explain how this pattern could arise through: (a) individual
  phenotypic plasticity within one generation; (b) natural selection
  across many generations. How would you distinguish these mechanisms
  with an experimental test?
\item
  Antibiotics are not ``mutagens'' --- they do not create resistant
  mutants. They primarily \emph{select for} pre-existing resistant
  mutants. Explain why this distinction matters for how we should use
  antibiotics clinically, and how the selection coefficient for
  resistance changes with antibiotic presence vs absence.
\item
  Compare the evidence supporting the gradual-change model of evolution
  (phyletic gradualism) vs the punctuated equilibrium model. What does
  the fossil record actually show, and what is the current consensus
  about the pace of evolutionary change?
\end{enumerate}

\subsection*{Extension Analysis
Questions}\label{extension-analysis-questions-10}
\addcontentsline{toc}{subsection}{Extension Analysis Questions}

\begin{enumerate}
\def\labelenumi{\arabic{enumi}.}
\setcounter{enumi}{5}
\item
  \textbf{Statistical Analysis}: If you performed three replicates of
  the selection simulation, how would you analyze the variation in
  survival rates? What statistical test would you use to determine if
  differences between prey colors are significant?
\item
  \textbf{Experimental Error Analysis}: What are three potential sources
  of error in the predator-prey simulation? How could you modify the
  procedure to reduce each source of error?
\item
  \textbf{Experimental Design}: If you wanted to test whether a
  particular trait is under directional, stabilizing, or disruptive
  selection, how would you design the experiment? What data would you
  collect? How would you analyze it?
\end{enumerate}

\subsection*{Group Project Extension
(Multi-Session)}\label{group-project-extension-multi-session-7}
\addcontentsline{toc}{subsection}{Group Project Extension
(Multi-Session)}

\textbf{Design an Evolutionary Study of Antibiotic Resistance}

This project will span 2-3 lab sessions:

\textbf{Session 1}: Research Question and Hypothesis - How does
antibiotic concentration affect the rate of resistance evolution in
bacteria? - Formulate a hypothesis about the relationship between
selection pressure and resistance development - Design a simulation
using provided allele-frequency and MIC datasets across different
antibiotic concentrations

\textbf{Session 2}: Data Collection and Analysis - Analyze printed
time-series data showing resistant and sensitive genotypes over
generations - Estimate resistance levels from provided minimum
inhibitory concentration tables - Compare allele-frequency changes
against a no-antibiotic baseline and a high-dose scenario

\textbf{Session 3}: Data Interpretation and Presentation - Compare
resistance evolution across different antibiotic concentrations -
Calculate selection coefficients for resistant vs sensitive bacteria -
Evaluate whether each scenario is reproducible by checking sample size,
replicate count, and starting allele frequency - Create a scientific
poster or presentation

\subsection*{Real-World Problem Solving: Evolution and
Medicine}\label{real-world-problem-solving-evolution-and-medicine}
\addcontentsline{toc}{subsection}{Real-World Problem Solving: Evolution
and Medicine}

\textbf{Case Study: Evolution of Virulence}

\begin{enumerate}
\def\labelenumi{\arabic{enumi}.}
\item
  \textbf{Research Task}: Investigate the trade-off hypothesis for the
  evolution of virulence. Under what conditions would a pathogen evolve
  to be more or less virulent?
\item
  \textbf{Ethical Consideration}: Should we use anti-evolution drugs
  (e.g., ``anti-evolution'' compounds that reduce mutation rates or
  horizontal gene transfer) to slow the evolution of resistance? What
  are the potential benefits and risks?
\item
  \textbf{Policy Proposal}: Design a policy for the responsible use of
  antibiotics in agriculture. Consider: growth promotion vs disease
  prevention, alternatives to antibiotics, and monitoring of resistance.
\end{enumerate}

\subsection*{Safety and Ethics Notes}\label{safety-and-ethics-notes-21}
\addcontentsline{toc}{subsection}{Safety and Ethics Notes}

Forceps (plastic or blunt-tip) are low risk. Ensure fairness in
simulation --- most students forage simultaneously. Small coloured
squares: keep away from young children (swallowing hazard).
Antibiotic-resistance work in this lab is paper-based primarily; do not
culture bacteria or handle antibiotics. When discussing antibiotic
resistance, emphasize following current prescriber instructions,
avoiding unnecessary antibiotic starts, using diagnostics when
available, and protecting access to effective treatment.

\subsection*{Debrief and Reflection}\label{debrief-and-reflection-18}
\addcontentsline{toc}{subsection}{Debrief and Reflection}

After you finish the practical work, spend 5--10 minutes in your small
group comparing results and discussing the following prompts. Each
member should contribute at least one observation before moving to the
next prompt:

\begin{enumerate}
\def\labelenumi{\arabic{enumi}.}
\item
  \textbf{What did your measurements show} --- compare the group's
  results to the textbook's predictions. Where they diverge, suggest at
  least one mechanistic explanation before concluding ``experimental
  error.''
\item
  \textbf{What would change the outcome} --- propose one modification to
  the procedure that would sharpen the measurement or extend the result
  to a new biological context, and predict what you would observe.
\item
  \textbf{One-sentence headline} --- each student composes a single
  sentence summarising the lab's take-home message, suitable for a
  tweet. Compare sentences across groups; good headlines are short,
  quantitative, and mechanistic.
\item
  \textbf{Connection back to the textbook} --- identify one section of
  \cref{sec:unit_VI_evolution_and_selection} that your data either
  confirmed or complicated. Cite the specific passage.
\item
  \textbf{Experimental Design Reflection}: What was the strongest aspect
  of your experimental design? What would you change if you could repeat
  the experiment?
\end{enumerate}

\subsection*{Further Reading (Lab)}\label{further-reading-lab-18}
\addcontentsline{toc}{subsection}{Further Reading (Lab)}

\begin{itemize}
\tightlist
\item
  Revisit the parent chapter \cref{sec:unit_VI_evolution_and_selection}
  for the theoretical foundations on which this lab is built.
\item
  Look up any \textbf{bolded glossary term} introduced in the textbook
  chapter (each has a \texttt{\#gl:term-slug} link in the text) --- its
  master definition is in \breaktt{manuscript/glossary.md}.
\item
  Explore the Python code in \breaktt{src/biology/evolution/evolution.py}
  for evolutionary calculations.
\item
  Use the self-contained Part 2 snippet as the computational template;
  it runs against tested project modules without external notebooks or
  CSV files.
\end{itemize}

\newpage

\section*{Lab 20 --- Genetic Drift, Gene Flow, and
Speciation}\label{lab-20-genetic-drift-gene-flow-and-speciation}
\addcontentsline{toc}{section}{Lab 20 --- Genetic Drift, Gene Flow, and
Speciation}

\label{sec:lab_unit_VI_genetic_drift_and_speciation}

\emph{This activity accompanies
\cref{sec:unit_VI_genetic_drift_and_speciation} of the textbook ---
review that chapter before attempting the exercises below.}

\subsection*{Learning Objectives}\label{learning-objectives-60}
\addcontentsline{toc}{subsection}{Learning Objectives}

\begin{itemize}
\tightlist
\item
  Design and evaluate a founder effect simulation to measure genetic
  drift
\item
  Formulate testable hypotheses about how population size affects
  genetic diversity
\item
  Identify independent, dependent, and controlled variables in
  population genetics experiments
\item
  Analyze experimental data using allele frequency calculations and Fst
  statistics
\item
  Write a brief scientific report with hypothesis, methods, results, and
  conclusion
\end{itemize}

\subsubsection*{Alignment and Rubric
Map}\label{alignment-and-rubric-map-22}
\addcontentsline{toc}{subsubsection}{Alignment and Rubric Map}

\begin{itemize}
\tightlist
\item
  \textbf{Outcome 1 (LO1):} Interpret the supplied evidence or model
  output for Genetic Drift, Gene Flow, and Speciation.
\item
  \textbf{Outcome 2 (LO2):} Identify controls and comparison groups that
  make the claim testable.
\item
  \textbf{Outcome 3 (LO3):} Quantify uncertainty, boundary conditions,
  or alternative explanations before concluding.
\item
  \textbf{Outcome 4 (LO4):} Transfer the mechanism to a new biological
  case or public-facing decision.
\item
  \textbf{Chapter LO coverage:} LO1, LO2, LO3, LO4
\item
  \textbf{Rubric dimensions:} evidence; controls; uncertainty;
  mechanism; transfer. \#\# Pre-Lab Inquiry Questions \{.unnumbered\}
\end{itemize}

\emph{Complete these questions before coming to lab. Use your textbook
and additional research to inform your hypotheses.}

\begin{enumerate}
\def\labelenumi{\arabic{enumi}.}
\item
  \textbf{Hypothesis Formation}: Based on your understanding of genetic
  drift, predict how the allele frequency in a small founding population
  (e.g., 6 individuals) would compare to the source population (100
  individuals). Write a clear, testable hypothesis.
\item
  \textbf{Experimental Design}: What is the independent variable in a
  genetic drift simulation? What is the dependent variable? List at
  least 5 variables that should be controlled to ensure valid results.
\item
  \textbf{Scientific Context}: The Fixation Index (Fst) quantifies
  genetic differentiation between populations. Why is this measure
  important in evolutionary biology? How does it relate to gene flow and
  speciation?
\item
  \textbf{Quantitative Reasoning}: If two populations have allele
  frequencies of 0.8 and 0.2 for the same gene, what would their Fst be?
  Show your calculation.
\item
  \textbf{Real-World Application}: The founder effect is thought to have
  played a role in the high prevalence of certain genetic disorders in
  isolated populations (e.g., Tay-Sachs in Ashkenazi Jews). How does the
  founder effect explain this pattern? What are the implications for
  genetic counseling?
\end{enumerate}

\subsection*{Background}\label{background-11}
\addcontentsline{toc}{subsection}{Background}

Even without natural selection, populations evolve through random
genetic drift, gene flow, and mutation. The \textbf{founder effect} (a
special case of bottleneck) occurs when a small number of individuals
colonise a new area, carrying primarily a subset of the original gene
pool's variation. \textbf{Speciation} --- the formation of new species
--- occurs when reproductive isolation evolves between populations,
typically after geographic (allopatric) or ecological (sympatric)
separation. The Fixation Index \textbf{Fst = (H\_T \texorpdfstring{\ensuremath{-}}{-} H\_S) / H\_T}
quantifies genetic differentiation between populations (0 = no
differentiation; 1 = complete differentiation).

\subsection*{Paper-Based Materials}\label{paper-based-materials-22}
\addcontentsline{toc}{subsection}{Paper-Based Materials}

{\def\LTcaptype{none} % do not increment counter
\begin{longtable}[]{@{}
  >{\raggedright\arraybackslash}p{(\linewidth - 2\tabcolsep) * \real{0.3333}}
  >{\raggedright\arraybackslash}p{(\linewidth - 2\tabcolsep) * \real{0.6667}}@{}}
\toprule\noalign{}
\begin{minipage}[b]{\linewidth}\raggedright
Item
\end{minipage} & \begin{minipage}[b]{\linewidth}\raggedright
Quantity
\end{minipage} \\
\midrule\noalign{}
\endhead
\bottomrule\noalign{}
\endlastfoot
100 beads (50 red, 50 blue --- representing two alleles) in a bag & 1 \\
Second bag for ``new island'' & 1 \\
Speciation case study cards (allopatric: Hawaiian honeycreepers;
sympatric: cichlid fish in Lake Victoria) & 1 set \\
Galapagos mockingbird morphological data (printed: beak length, body
mass per island) & 1 \\
Calculator & 1 \\
Graph paper & 1 \\
Calculator or optional Python REPL with this project installed & 1 \\
Population genetics worksheet & 1 \\
\end{longtable}
}

\subsection*{Paper-Based
Investigation}\label{paper-based-investigation-22}
\addcontentsline{toc}{subsection}{Paper-Based Investigation}

\subsubsection*{Part 1: Founder Effect
Simulation}\label{part-1-founder-effect-simulation}
\addcontentsline{toc}{subsubsection}{Part 1: Founder Effect Simulation}

\begin{enumerate}
\def\labelenumi{\arabic{enumi}.}
\item
  \textbf{Formulate Your Hypothesis}: Based on your pre-lab research,
  write a hypothesis about how the size of a founding population affects
  the rate of genetic drift.
\item
  \textbf{Identify Variables}:

  \begin{itemize}
  \tightlist
  \item
    Independent variable: Founder population size (6 vs 20 individuals)
  \item
    Dependent variable: Change in allele frequency from source
    population
  \item
    Controlled variables: Source population allele frequency, number of
    generations, etc.
  \end{itemize}
\item
  \textbf{Set Up Controls}: Why is it important to start with a known
  source population allele frequency? What other controls could you
  include in a genetic drift simulation?
\item
  \textbf{Data Collection}:

  \begin{itemize}
  \tightlist
  \item
    Start: 100 beads at p = 0.5. Draw 6 beads randomly (founder
    population) --- do not replace.
  \item
    Record allele frequencies in the founder group; use these to
    repopulate to 100 (multiply founder frequencies \texorpdfstring{\ensuremath{\times}}{x} 100).
  \item
    Repeat sampling: draw 20 from the new population to simulate a
    moderate founding event.
  \item
    Compare p values: original, 6-founder, 20-founder.
  \item
    Perform at least 3 replicates for each founder size.
  \end{itemize}
\end{enumerate}

\subsubsection*{Part 2: Computational Biology Exercise - Population
Genetics Analysis with
Python}\label{part-2-computational-biology-exercise---population-genetics-analysis-with-python}
\addcontentsline{toc}{subsubsection}{Part 2: Computational Biology
Exercise - Population Genetics Analysis with Python}

\emph{Optional computational check: run this self-contained Python
snippet from the project root. It uses tested \texttt{src/biology}
modules and requires no external notebook or CSV file.}

Run deterministic drift simulations by fixing the random seed.

\begin{Shaded}
\begin{Highlighting}[]
\ImportTok{from}\NormalTok{ biology.evolution }\ImportTok{import}\NormalTok{ isolation\_index, simulate\_drift}

\NormalTok{small }\OperatorTok{=}\NormalTok{ simulate\_drift(p}\OperatorTok{=}\FloatTok{0.5}\NormalTok{, N}\OperatorTok{=}\DecValTok{20}\NormalTok{, generations}\OperatorTok{=}\DecValTok{10}\NormalTok{, rng\_seed}\OperatorTok{=}\DecValTok{7}\NormalTok{)}
\NormalTok{large }\OperatorTok{=}\NormalTok{ simulate\_drift(p}\OperatorTok{=}\FloatTok{0.5}\NormalTok{, N}\OperatorTok{=}\DecValTok{500}\NormalTok{, generations}\OperatorTok{=}\DecValTok{10}\NormalTok{, rng\_seed}\OperatorTok{=}\DecValTok{7}\NormalTok{)}

\BuiltInTok{print}\NormalTok{(}\StringTok{"small{-}population final p:"}\NormalTok{, }\BuiltInTok{round}\NormalTok{(small[}\OperatorTok{{-}}\DecValTok{1}\NormalTok{], }\DecValTok{3}\NormalTok{))}
\BuiltInTok{print}\NormalTok{(}\StringTok{"large{-}population final p:"}\NormalTok{, }\BuiltInTok{round}\NormalTok{(large[}\OperatorTok{{-}}\DecValTok{1}\NormalTok{], }\DecValTok{3}\NormalTok{))}
\BuiltInTok{print}\NormalTok{(}\StringTok{"isolation example:"}\NormalTok{, }\BuiltInTok{round}\NormalTok{(isolation\_index(gene\_flow\_rate}\OperatorTok{=}\FloatTok{0.01}\NormalTok{, mutation\_rate}\OperatorTok{=}\FloatTok{0.001}\NormalTok{), }\DecValTok{3}\NormalTok{))}
\end{Highlighting}
\end{Shaded}

\subsubsection*{Part 3: Morphological Divergence
Analysis}\label{part-3-morphological-divergence-analysis}
\addcontentsline{toc}{subsubsection}{Part 3: Morphological Divergence
Analysis}

\begin{enumerate}
\def\labelenumi{\arabic{enumi}.}
\setcounter{enumi}{4}
\item
  \textbf{Design an Alternative Investigation}: Instead of just
  analyzing provided morphological data, design a paper-based comparison
  to test how geographic isolation affects trait divergence in a natural
  population. What hypothesis would you test? Which archived
  measurements, sample sizes, and map distances would you need?
\item
  \textbf{Hypothesis}: Predict the relationship between geographic
  distance and morphological divergence. Would you expect a linear
  relationship, or something else? Why?
\end{enumerate}

\subsection*{Data Recording}\label{data-recording-19}
\addcontentsline{toc}{subsection}{Data Recording}

{\def\LTcaptype{none} % do not increment counter
\begin{longtable}[]{@{}llll@{}}
\toprule\noalign{}
Founding event & Sample size & p (red allele) & q (blue allele) \\
\midrule\noalign{}
\endhead
\bottomrule\noalign{}
\endlastfoot
Original & 100 & 0.5 & 0.5 \\
6-founder & 6 & & \\
20-founder & 20 & & \\
\end{longtable}
}

Mockingbird data summary (fill from provided table):

{\def\LTcaptype{none} % do not increment counter
\begin{longtable}[]{@{}
  >{\raggedright\arraybackslash}p{(\linewidth - 6\tabcolsep) * \real{0.0882}}
  >{\raggedright\arraybackslash}p{(\linewidth - 6\tabcolsep) * \real{0.3088}}
  >{\raggedright\arraybackslash}p{(\linewidth - 6\tabcolsep) * \real{0.2794}}
  >{\raggedright\arraybackslash}p{(\linewidth - 6\tabcolsep) * \real{0.3235}}@{}}
\toprule\noalign{}
\begin{minipage}[b]{\linewidth}\raggedright
Island
\end{minipage} & \begin{minipage}[b]{\linewidth}\raggedright
Mean beak length (mm)
\end{minipage} & \begin{minipage}[b]{\linewidth}\raggedright
Mean body mass (g)
\end{minipage} & \begin{minipage}[b]{\linewidth}\raggedright
Distinct from others?
\end{minipage} \\
\midrule\noalign{}
\endhead
\bottomrule\noalign{}
\endlastfoot
Santa Cruz & & & \\
Genovesa & & & \\
Española & & & \\
Fernandina & & & \\
\end{longtable}
}

Fst calculation: H\_T = \textbf{\emph{; H\_S = }}; Fst = \_\_\_;
Interpretation:

\subsection*{Evidence and Reproducibility
Checklist}\label{evidence-and-reproducibility-checklist-22}
\addcontentsline{toc}{subsection}{Evidence and Reproducibility
Checklist}

\begin{itemize}
\tightlist
\item
  \textbf{Primary evidence goal:} Model drift and speciation scenarios
  with repeated trials.
\item
  \textbf{Data skill to practice:} Distinguish stochastic from
  directional change in allele-frequency data.
\item
  \textbf{BioSkills emphasis:} Modeling and simulation, Quantitative
  reasoning, Communication and collaboration.
\item
  \textbf{Control logic:} identify at least one positive control, one
  negative control, or one baseline comparison before interpreting
  results.
\item
  \textbf{Measurement discipline:} record units, uncertainty, sample
  size, and any discarded observation with a reason.
\item
  \textbf{Mechanistic link:} connect one result directly to the parent
  chapter's big idea before writing the conclusion.
\item
  \textbf{Reproducibility check:} state one procedural detail that
  another group would need in order to reproduce the result.
\end{itemize}

\subsection*{Paper-Based Evidence
Upgrade}\label{paper-based-evidence-upgrade-22}
\addcontentsline{toc}{subsection}{Paper-Based Evidence Upgrade}

Before answering the analysis questions, annotate the paper dataset for
\textbf{Genetic Drift, Gene Flow, and Speciation} with a reproducibility
pass:

{\def\LTcaptype{none} % do not increment counter
\begin{longtable}[]{@{}
  >{\raggedright\arraybackslash}p{(\linewidth - 2\tabcolsep) * \real{0.5000}}
  >{\raggedright\arraybackslash}p{(\linewidth - 2\tabcolsep) * \real{0.5000}}@{}}
\toprule\noalign{}
\begin{minipage}[b]{\linewidth}\raggedright
Evidence check
\end{minipage} & \begin{minipage}[b]{\linewidth}\raggedright
Student action
\end{minipage} \\
\midrule\noalign{}
\endhead
\bottomrule\noalign{}
\endlastfoot
Control logic & Mark the comparison that functions as the baseline,
negative control, or reference case. \\
Uncertainty & Circle the row, card, diagram feature, or model assumption
most likely to change the conclusion. \\
Model comparison & State whether a simpler rule, null model, or
alternative mechanism could explain the same pattern. \\
Decision threshold & Write the minimum evidence that would make you
revise the interpretation. \\
Reproducibility & Record the exact scoring rule another group would need
to reproduce your classification. \\
\end{longtable}
}

Focus note: in this founder-event and \(F_{ST}\) lab, distinguish
allele-frequency change caused by sampling in small colonising groups
(drift) from differentiation maintained by restricted gene flow, and
state which reproductive-isolation barrier each cichlid or mockingbird
comparison would require before calling the populations distinct
species. Keep required work paper-based; any material-handling or
equipment version belongs only in an optional extension.

\subsection*{Analysis Questions}\label{analysis-questions-22}
\addcontentsline{toc}{subsection}{Analysis Questions}

\begin{enumerate}
\def\labelenumi{\arabic{enumi}.}
\item
  How did the allele frequencies in the 6-founder event compare to the
  original population? What does this demonstrate about why small
  colonising populations often show low genetic diversity compared to
  mainland populations?
\item
  The mockingbird populations on different Galapagos islands show
  morphological divergence. What mechanism prevented gene flow between
  islands, and what type of speciation does this represent? What
  reproductive isolation mechanisms would need to evolve for them to
  become separate species?
\item
  Fst = 0 indicates populations are identical (freely exchanging genes).
  Fst = 1 indicates complete isolation. How would you interpret Fst =
  0.15 --- is there gene flow? Use your calculated Fst to describe the
  degree of isolation between cichlid populations A, B, and C.
\item
  Lake Victoria cichlids underwent adaptive radiation in
  about 15,000 years, producing \textgreater500 species. What
  does this rapid speciation suggest about the speed at which
  reproductive isolation can evolve? What kind of speciation
  (allopatric, parapatric, or sympatric) is most consistent with this
  pattern?
\item
  Ring species (like Larus gulls circling the Arctic) demonstrate
  speciation in action: adjacent populations can interbreed, but the two
  ends of the ring cannot. Explain what this demonstrates about the
  continuity of the speciation process and why ring species challenge a
  strict definition of ``species.''
\end{enumerate}

\subsection*{Extension Analysis
Questions}\label{extension-analysis-questions-11}
\addcontentsline{toc}{subsection}{Extension Analysis Questions}

\begin{enumerate}
\def\labelenumi{\arabic{enumi}.}
\setcounter{enumi}{5}
\item
  \textbf{Statistical Analysis}: If you performed three replicates of
  the founder effect simulation, how would you analyze the variation in
  allele frequency change? What statistical test would you use to
  determine if differences between founder sizes are significant?
\item
  \textbf{Experimental Error Analysis}: What are three potential sources
  of error in the bead simulation? How could you modify the procedure to
  reduce each source of error?
\item
  \textbf{Experimental Design}: If you wanted to test whether two
  populations are diverging due to genetic drift or natural selection,
  how would you design the experiment? What data would you collect? How
  would you analyze it?
\end{enumerate}

\subsection*{Group Project Extension
(Multi-Session)}\label{group-project-extension-multi-session-8}
\addcontentsline{toc}{subsection}{Group Project Extension
(Multi-Session)}

\textbf{Design a Study of Human Population Differentiation}

This project will span 2-3 lab sessions:

\textbf{Session 1}: Research Question and Hypothesis - How does
geographic distance affect genetic differentiation in human populations?
- Formulate a hypothesis about the relationship between distance and Fst
- Design a study using publicly available genetic data (e.g., 1000
Genomes Project)

\textbf{Session 2}: Data Collection and Analysis - Download genetic data
for multiple populations - Calculate Fst between population pairs - Plot
Fst against geographic distance

\textbf{Session 3}: Data Interpretation and Presentation - Analyze the
correlation between distance and genetic differentiation - Discuss the
implications for understanding human migration history - Create a
scientific poster or presentation

\subsection*{Real-World Problem Solving: Conservation
Genetics}\label{real-world-problem-solving-conservation-genetics}
\addcontentsline{toc}{subsection}{Real-World Problem Solving:
Conservation Genetics}

\textbf{Case Study: Genetic Rescue of Endangered Populations}

\begin{enumerate}
\def\labelenumi{\arabic{enumi}.}
\item
  \textbf{Research Task}: Investigate the Florida panther bottleneck.
  How did genetic drift and inbreeding depression affect the population?
  What was the outcome of the genetic rescue program?
\item
  \textbf{Ethical Consideration}: Should we introduce individuals from
  other populations to increase genetic diversity in endangered species?
  What are the potential benefits and risks (e.g., outbreeding
  depression)?
\item
  \textbf{Policy Proposal}: Design a policy for managing genetic
  diversity in small, isolated populations. Consider: when to intervene,
  how to select individuals for translocation, and how to monitor
  success.
\end{enumerate}

\subsection*{Safety and Ethics Notes}\label{safety-and-ethics-notes-22}
\addcontentsline{toc}{subsection}{Safety and Ethics Notes}

Bead simulation: ensure beads are not ingested. Discussions of
speciation based on genetic data should be careful not to conflate
genetic population differences with outdated racial taxonomic thinking.
Consider the ethical implications of conservation genetics.

\subsection*{Debrief and Reflection}\label{debrief-and-reflection-19}
\addcontentsline{toc}{subsection}{Debrief and Reflection}

After you finish the practical work, spend 5--10 minutes in your small
group comparing results and discussing the following prompts. Each
member should contribute at least one observation before moving to the
next prompt:

\begin{enumerate}
\def\labelenumi{\arabic{enumi}.}
\item
  \textbf{What did your measurements show} --- compare the group's
  results to the textbook's predictions. Where they diverge, suggest at
  least one mechanistic explanation before concluding ``experimental
  error.''
\item
  \textbf{What would change the outcome} --- propose one modification to
  the procedure that would sharpen the measurement or extend the result
  to a new biological context, and predict what you would observe.
\item
  \textbf{One-sentence headline} --- each student composes a single
  sentence summarising the lab's take-home message, suitable for a
  tweet. Compare sentences across groups; good headlines are short,
  quantitative, and mechanistic.
\item
  \textbf{Connection back to the textbook} --- identify one section of
  \cref{sec:unit_VI_genetic_drift_and_speciation} that your data either
  confirmed or complicated. Cite the specific passage.
\item
  \textbf{Experimental Design Reflection}: What was the strongest aspect
  of your experimental design? What would you change if you could repeat
  the experiment?
\end{enumerate}

\subsection*{Further Reading (Lab)}\label{further-reading-lab-19}
\addcontentsline{toc}{subsection}{Further Reading (Lab)}

\begin{itemize}
\tightlist
\item
  Revisit the parent chapter
  \cref{sec:unit_VI_genetic_drift_and_speciation} for the theoretical
  foundations on which this lab is built.
\item
  Look up any \textbf{bolded glossary term} introduced in the textbook
  chapter (each has a \texttt{\#gl:term-slug} link in the text) --- its
  master definition is in \breaktt{manuscript/glossary.md}.
\item
  Explore the Python code in \breaktt{src/biology/evolution/evolution.py}
  for population genetics calculations.
\item
  Use the self-contained Part 2 snippet as the computational template;
  it runs against tested project modules without external notebooks or
  CSV files.
\end{itemize}

\newpage

\section*{Lab 21 --- Phylogenetics and the Tree of
Life}\label{lab-21-phylogenetics-and-the-tree-of-life}
\addcontentsline{toc}{section}{Lab 21 --- Phylogenetics and the Tree of
Life}

\label{sec:lab_unit_VI_phylogenetics}

\subsection*{Learning Objectives}\label{learning-objectives-61}
\addcontentsline{toc}{subsection}{Learning Objectives}

\emph{This activity accompanies \cref{sec:unit_VI_phylogenetics} of the
textbook --- review that chapter before attempting the exercises below.}

\begin{itemize}
\tightlist
\item
  Construct a cladogram from a morphological character matrix using
  parsimony
\item
  Interpret a molecular phylogeny to determine divergence times and
  evolutionary relationships
\item
  Identify synapomorphies and plesiomorphies from a character matrix
\item
  Evaluate conflicting phylogenetic hypotheses using bootstrap support
  values
\end{itemize}

\subsubsection*{Alignment and Rubric
Map}\label{alignment-and-rubric-map-23}
\addcontentsline{toc}{subsubsection}{Alignment and Rubric Map}

\begin{itemize}
\tightlist
\item
  \textbf{Outcome 1 (LO1):} Interpret the supplied evidence or model
  output for Phylogenetics and the Tree of Life.
\item
  \textbf{Outcome 2 (LO2):} Identify controls and comparison groups that
  make the claim testable.
\item
  \textbf{Outcome 3 (LO3):} Quantify uncertainty, boundary conditions,
  or alternative explanations before concluding.
\item
  \textbf{Outcome 4 (LO4):} Transfer the mechanism to a new biological
  case or public-facing decision.
\item
  \textbf{Chapter LO coverage:} LO1, LO2, LO3, LO4
\item
  \textbf{Rubric dimensions:} evidence; controls; uncertainty;
  mechanism; transfer. \#\# Background \{.unnumbered\}
\end{itemize}

Phylogenetics reconstructs evolutionary history by identifying shared
derived characters (\textbf{synapomorphies}) that define clades. Maximum
parsimony selects the tree requiring the fewest evolutionary changes.
Molecular phylogenies use DNA/protein sequence differences, calibrated
with fossil dates, to construct ultrametric trees with branch lengths
proportional to substitution rate or time. Bootstrap values (percentage
of resampled datasets recovering the same clade) indicate support for
individual nodes.

\subsection*{Paper-Based Materials}\label{paper-based-materials-23}
\addcontentsline{toc}{subsection}{Paper-Based Materials}

{\def\LTcaptype{none} % do not increment counter
\begin{longtable}[]{@{}
  >{\raggedright\arraybackslash}p{(\linewidth - 2\tabcolsep) * \real{0.3333}}
  >{\raggedright\arraybackslash}p{(\linewidth - 2\tabcolsep) * \real{0.6667}}@{}}
\toprule\noalign{}
\begin{minipage}[b]{\linewidth}\raggedright
Item
\end{minipage} & \begin{minipage}[b]{\linewidth}\raggedright
Quantity
\end{minipage} \\
\midrule\noalign{}
\endhead
\bottomrule\noalign{}
\endlastfoot
Character matrix (12 taxa \texorpdfstring{\ensuremath{\times}}{x} 10 morphological characters; 0=absent,
1=present) & 1 printed \\
Molecular phylogeny printout (5 taxa; with branch lengths and bootstrap
values) & 1 \\
Cladogram construction guide & 1 \\
Scissors; paste; paper & 1 set \\
Calculator & 1 \\
\end{longtable}
}

\subsection*{Paper-Based
Investigation}\label{paper-based-investigation-23}
\addcontentsline{toc}{subsection}{Paper-Based Investigation}

\textbf{Part A --- Manual Cladogram Construction}

\begin{enumerate}
\def\labelenumi{\arabic{enumi}.}
\tightlist
\item
  Using the character matrix, identify which characters are shared among
  subsets of taxa (synapomorphies) vs ancestral (plesiomorphies shared
  by outgroup).
\item
  Use the principle of parsimony to construct the most parsimonious
  cladogram for 6 selected taxa (the other 6 are ``outgroups'' --- use
  to root the tree).
\item
  Count the total number of character-state changes on your tree vs an
  alternative arrangement. Identify which tree requires fewer changes
  (most parsimonious).
\end{enumerate}

\textbf{Part B --- Molecular Phylogeny Interpretation}

\begin{enumerate}
\def\labelenumi{\arabic{enumi}.}
\setcounter{enumi}{3}
\tightlist
\item
  On the provided molecular phylogeny: identify (a) the most recent
  common ancestor of taxa A and B; (b) which taxon diverged first from
  the others; (c) any polytomies (unresolved nodes); (d) the clade with
  the highest bootstrap support.
\end{enumerate}

\subsection*{Data Recording}\label{data-recording-20}
\addcontentsline{toc}{subsection}{Data Recording}

Character matrix (excerpt --- 5 taxa):

{\def\LTcaptype{none} % do not increment counter
\begin{longtable}[]{@{}
  >{\raggedright\arraybackslash}p{(\linewidth - 10\tabcolsep) * \real{0.0926}}
  >{\raggedright\arraybackslash}p{(\linewidth - 10\tabcolsep) * \real{0.1667}}
  >{\raggedright\arraybackslash}p{(\linewidth - 10\tabcolsep) * \real{0.1852}}
  >{\raggedright\arraybackslash}p{(\linewidth - 10\tabcolsep) * \real{0.1481}}
  >{\raggedright\arraybackslash}p{(\linewidth - 10\tabcolsep) * \real{0.2593}}
  >{\raggedright\arraybackslash}p{(\linewidth - 10\tabcolsep) * \real{0.1481}}@{}}
\toprule\noalign{}
\begin{minipage}[b]{\linewidth}\raggedright
Taxon
\end{minipage} & \begin{minipage}[b]{\linewidth}\raggedright
Vertebrae
\end{minipage} & \begin{minipage}[b]{\linewidth}\raggedright
Four limbs
\end{minipage} & \begin{minipage}[b]{\linewidth}\raggedright
Feathers
\end{minipage} & \begin{minipage}[b]{\linewidth}\raggedright
Mammary glands
\end{minipage} & \begin{minipage}[b]{\linewidth}\raggedright
Placenta
\end{minipage} \\
\midrule\noalign{}
\endhead
\bottomrule\noalign{}
\endlastfoot
Lamprey & 1 & 0 & 0 & 0 & 0 \\
Salamander & 1 & 1 & 0 & 0 & 0 \\
Crocodile & 1 & 1 & 0 & 0 & 0 \\
Robin & 1 & 1 & 1 & 0 & 0 \\
Human & 1 & 1 & 0 & 1 & 1 \\
\end{longtable}
}

Your cladogram sketch:

Most parsimonious tree --- total changes: \textbf{\emph{ Alternative
tree --- total changes: }}

Molecular phylogeny notes: - Taxon diverging first: \textbf{\emph{ -
Highest bootstrap node: }} (\textbf{\emph{\% support) - Any polytomies?
}}

\subsection*{Evidence and Reproducibility
Checklist}\label{evidence-and-reproducibility-checklist-23}
\addcontentsline{toc}{subsection}{Evidence and Reproducibility
Checklist}

\begin{itemize}
\tightlist
\item
  \textbf{Primary evidence goal:} Infer relationships from character or
  sequence evidence.
\item
  \textbf{Data skill to practice:} Read trees correctly and map traits
  or sequences onto branches.
\item
  \textbf{BioSkills emphasis:} Modeling and simulation, Quantitative
  reasoning, Communication and collaboration.
\item
  \textbf{Control logic:} identify at least one positive control, one
  negative control, or one baseline comparison before interpreting
  results.
\item
  \textbf{Measurement discipline:} record units, uncertainty, sample
  size, and any discarded observation with a reason.
\item
  \textbf{Mechanistic link:} connect one result directly to the parent
  chapter's big idea before writing the conclusion.
\item
  \textbf{Reproducibility check:} state one procedural detail that
  another group would need in order to reproduce the result.
\end{itemize}

\subsection*{Paper-Based Evidence
Upgrade}\label{paper-based-evidence-upgrade-23}
\addcontentsline{toc}{subsection}{Paper-Based Evidence Upgrade}

Before answering the analysis questions, annotate the paper dataset for
\textbf{Phylogenetics and the Tree of Life} with a reproducibility pass:

{\def\LTcaptype{none} % do not increment counter
\begin{longtable}[]{@{}
  >{\raggedright\arraybackslash}p{(\linewidth - 2\tabcolsep) * \real{0.5000}}
  >{\raggedright\arraybackslash}p{(\linewidth - 2\tabcolsep) * \real{0.5000}}@{}}
\toprule\noalign{}
\begin{minipage}[b]{\linewidth}\raggedright
Evidence check
\end{minipage} & \begin{minipage}[b]{\linewidth}\raggedright
Student action
\end{minipage} \\
\midrule\noalign{}
\endhead
\bottomrule\noalign{}
\endlastfoot
Control logic & Mark the comparison that functions as the baseline,
negative control, or reference case. \\
Uncertainty & Circle the row, card, diagram feature, or model assumption
most likely to change the conclusion. \\
Model comparison & State whether a simpler rule, null model, or
alternative mechanism could explain the same pattern. \\
Decision threshold & Write the minimum evidence that would make you
revise the interpretation. \\
Reproducibility & Record the exact scoring rule another group would need
to reproduce your classification. \\
\end{longtable}
}

Focus note: in this character-matrix and tree-building lab, justify each
clade with a shared derived character (synapomorphy) rather than overall
similarity, treat bootstrap percentages as resampling support rather
than proof, and flag where horizontal gene transfer or hybridisation
would make a strictly bifurcating tree the wrong model. Keep required
work paper-based; any material-handling or equipment version belongs
only in an optional extension.

\subsection*{Analysis Questions}\label{analysis-questions-23}
\addcontentsline{toc}{subsection}{Analysis Questions}

\begin{enumerate}
\def\labelenumi{\arabic{enumi}.}
\tightlist
\item
  In your character matrix, which character (vertebrae, four limbs,
  feathers, etc.) defines the broadest clade (most inclusive group)?
  Which is the most restricted synapomorphy (smallest clade)? Explain
  the concept of nested hierarchy.
\item
  Crocodiles share a more recent common ancestor with birds than with
  lizards, despite looking more like lizards. What does this illustrate
  about the limits of morphological similarity as an indicator of
  evolutionary relationship?
\item
  A bootstrap value of 95\% means that 95\% of resampled datasets
  recover that node. How would you interpret a node with 42\% bootstrap
  support? Should it be included in the published tree?
\item
  Horizontal gene transfer (HGT) in bacteria allows genes to jump
  between unrelated lineages --- violating the vertical
  (parent-to-offspring) assumption of phylogenetics. Explain how HGT
  would appear as a misleading signal in a molecular phylogeny, and
  describe one method used to detect it.
\item
  The ``Tree of Life'' analogy implies a bifurcating, hierarchical
  structure. However, hybridisation (interspecies mating, e.g., in
  plants) produces reticulate evolution --- a network rather than a
  tree. Name two plant genera where reticulate evolution occurs and
  describe how allopolyploidy contributes to speciation.
\end{enumerate}

\subsection*{Safety and Ethics Notes}\label{safety-and-ethics-notes-23}
\addcontentsline{toc}{subsection}{Safety and Ethics Notes}

No chemical hazards. When discussing evolutionary relationships
involving humans and other primates, maintain scientific respect for
most organisms and evolutionary histories.

\subsection*{Debrief and Reflection}\label{debrief-and-reflection-20}
\addcontentsline{toc}{subsection}{Debrief and Reflection}

After you finish the practical work, spend 5--10 minutes in your small
group comparing results and discussing the following prompts. Each
member should contribute at least one observation before moving to the
next prompt:

\begin{enumerate}
\def\labelenumi{\arabic{enumi}.}
\tightlist
\item
  \textbf{What did your measurements show} --- compare the group's
  results to the textbook's predictions. Where they diverge, suggest at
  least one mechanistic explanation before concluding ``experimental
  error.''
\item
  \textbf{What would change the outcome} --- propose one modification to
  the procedure that would sharpen the measurement or extend the result
  to a new biological context, and predict what you would observe.
\item
  \textbf{One-sentence headline} --- each student composes a single
  sentence summarising the lab's take-home message, suitable for a
  tweet. Compare sentences across groups; good headlines are short,
  quantitative, and mechanistic.
\item
  \textbf{Connection back to the textbook} --- identify one section of
  \cref{sec:unit_VI_phylogenetics} that your data either confirmed or
  complicated. Cite the specific passage.
\end{enumerate}

\subsection*{Further Reading (Lab)}\label{further-reading-lab-20}
\addcontentsline{toc}{subsection}{Further Reading (Lab)}

\begin{itemize}
\tightlist
\item
  Revisit the parent chapter \cref{sec:unit_VI_phylogenetics} for the
  theoretical foundations on which this lab is built.
\item
  Look up any \textbf{bolded glossary term} introduced in the textbook
  chapter (each has a \texttt{\#gl:term-slug} link in the text) --- its
  master definition is in \breaktt{manuscript/glossary.md}.
\item
  Explore the appended \texttt{src/} module that implements the
  corresponding quantitative model (when applicable) --- referenced in
  the parent chapter's ``Bridge to Computation'' subsection.
\end{itemize}

\emph{Module footer: parent chapter
\texttt{\textbackslash{}cref\{sec:unit\_VI\_phylogenetics\}}; all
numerical quantities in this lab use SI units --- see Appendix D of the
textbook for unit conversions and biological-scale reference values.}

\newpage

\section*{Lab 22 --- Bacteria, Archaea, and
Viruses}\label{lab-22-bacteria-archaea-and-viruses}
\addcontentsline{toc}{section}{Lab 22 --- Bacteria, Archaea, and
Viruses}

\label{sec:lab_unit_VII_bacteria_archaea_viruses}

\emph{This activity accompanies
\cref{sec:unit_VII_bacteria_archaea_viruses} of the textbook --- review
that chapter before attempting the exercises below.}

\subsection*{Learning Objectives}\label{learning-objectives-62}
\addcontentsline{toc}{subsection}{Learning Objectives}

\begin{itemize}
\tightlist
\item
  Design and evaluate a bacterial identification case packet using
  Gram-stain images
\item
  Formulate testable hypotheses about how environmental factors affect
  simulated bacterial growth
\item
  Identify independent, dependent, and controlled variables in
  microbiology experiments
\item
  Analyze experimental data using colony counting and statistical
  methods
\item
  Write a brief scientific report with hypothesis, methods, results, and
  conclusion
\end{itemize}

\subsubsection*{Alignment and Rubric
Map}\label{alignment-and-rubric-map-24}
\addcontentsline{toc}{subsubsection}{Alignment and Rubric Map}

\begin{itemize}
\tightlist
\item
  \textbf{Outcome 1 (LO1):} Interpret the supplied evidence or model
  output for Bacteria, Archaea, and Viruses.
\item
  \textbf{Outcome 2 (LO2):} Identify controls and comparison groups that
  make the claim testable.
\item
  \textbf{Outcome 3 (LO3):} Quantify uncertainty, boundary conditions,
  or alternative explanations before concluding.
\item
  \textbf{Outcome 4 (LO4):} Transfer the mechanism to a new biological
  case or public-facing decision.
\item
  \textbf{Chapter LO coverage:} LO1, LO2, LO3, LO4
\item
  \textbf{Rubric dimensions:} evidence; controls; uncertainty;
  mechanism; transfer. \#\# Pre-Lab Inquiry Questions \{.unnumbered\}
\end{itemize}

\emph{Complete these questions before coming to lab. Use your textbook
and additional research to inform your hypotheses.}

\begin{enumerate}
\def\labelenumi{\arabic{enumi}.}
\item
  \textbf{Hypothesis Formation}: Based on your understanding of Gram
  staining, predict which type of bacteria (Gram-positive or
  Gram-negative) would be more resistant to antibiotics that target cell
  wall synthesis. Explain your reasoning.
\item
  \textbf{Experimental Design}: What is the independent variable in a
  bacterial growth dataset? What is the dependent variable? List at
  least 5 variables that should be controlled to ensure valid
  conclusions.
\item
  \textbf{Scientific Context}: The Gram stain is a fundamental tool in
  microbiology. Why is it important to classify bacteria as
  Gram-positive or Gram-negative? How does this classification inform
  treatment decisions?
\item
  \textbf{Quantitative Reasoning}: A printed dilution dataset shows 150
  colonies on a 10\texorpdfstring{\textsuperscript{-}}{-}\texorpdfstring{\textsuperscript{5}}{5} dilution plate, with 0.1 mL plated. What is the
  concentration of bacteria in the original culture? Show your
  calculation.
\item
  \textbf{Real-World Application}: Antibiotic resistance is a growing
  problem. How does the misuse of antibiotics in agriculture contribute
  to the evolution of resistant bacteria? What policies could help
  address this issue?
\end{enumerate}

\subsection*{Background}\label{background-12}
\addcontentsline{toc}{subsection}{Background}

Prokaryotes --- Bacteria and Archaea --- are the most abundant life
forms on Earth (about 10\texorpdfstring{\textsuperscript{3}}{3}\texorpdfstring{\textsuperscript{1}}{1} cells). Gram staining differentiates
bacteria based on cell wall composition: Gram-positive (thick
peptidoglycan layer retains crystal violet --- purple) vs Gram-negative
(thin peptidoglycan + outer membrane --- loses crystal violet, appears
pink after safranin counterstain). Viruses are acellular entities that
replicate primarily inside host cells; bacteriophages infect bacteria
and follow either lytic (immediate replication and lysis) or lysogenic
(DNA integration, dormancy) cycles.

\subsection*{Paper-Based Materials}\label{paper-based-materials-24}
\addcontentsline{toc}{subsection}{Paper-Based Materials}

{\def\LTcaptype{none} % do not increment counter
\begin{longtable}[]{@{}
  >{\raggedright\arraybackslash}p{(\linewidth - 2\tabcolsep) * \real{0.3333}}
  >{\raggedright\arraybackslash}p{(\linewidth - 2\tabcolsep) * \real{0.6667}}@{}}
\toprule\noalign{}
\begin{minipage}[b]{\linewidth}\raggedright
Item
\end{minipage} & \begin{minipage}[b]{\linewidth}\raggedright
Quantity
\end{minipage} \\
\midrule\noalign{}
\endhead
\bottomrule\noalign{}
\endlastfoot
Gram-stain image cards with metadata removed & 1 set \\
Known Gram-positive and Gram-negative control image cards & 1 set \\
Printed serial-dilution plate images (10\texorpdfstring{\textsuperscript{-}}{-}\texorpdfstring{\textsuperscript{4}}{4} to 10\texorpdfstring{\textsuperscript{-}}{-}\texorpdfstring{\textsuperscript{6}}{6}, three replicates) &
1 set \\
Bacterial growth curve data under temperature, pH, and salt scenarios &
1 \\
Phage life-cycle decision cards & 1 set \\
AMR treatment decision matrix & 1 \\
Horizontal-gene-transfer map cards: plasmid, transposon, integron,
phage, transformation, wastewater, and hospital ward & 1 set \\
\emph{Candida auris} clinical triage card: colonization vs infection,
identification method, surface persistence, susceptibility result & 1 \\
Calculator or optional Python REPL with this project installed & 1 \\
Ruler or transparent counting grid & 1 \\
\end{longtable}
}

\subsection*{Paper-Based
Investigation}\label{paper-based-investigation-24}
\addcontentsline{toc}{subsection}{Paper-Based Investigation}

\subsubsection*{Part 1: Gram-Stain Image
Classification}\label{part-1-gram-stain-image-classification}
\addcontentsline{toc}{subsubsection}{Part 1: Gram-Stain Image
Classification}

\begin{enumerate}
\def\labelenumi{\arabic{enumi}.}
\item
  \textbf{Formulate Your Hypothesis}: Based on your pre-lab research,
  write a hypothesis about how the Gram stain results will differ
  between \emph{E. coli} and \emph{Bacillus} species.
\item
  \textbf{Identify Variables}:

  \begin{itemize}
  \tightlist
  \item
    Independent variable: Bacterial species
  \item
    Dependent variable: Gram stain result (purple vs pink) and cell
    morphology
  \item
    Controlled variables: image source, colour-calibration note, scale
    label, scoring key, scorer training, and blind-card order.
  \end{itemize}
\item
  \textbf{Set Up Controls}: Score the known Gram-positive and
  Gram-negative control cards first. What other controls could you
  include to ensure accurate identification, such as a mixed-species
  image, a low-quality image, or a blinded duplicate?
\item
  \textbf{Data Collection}:

  \begin{itemize}
  \tightlist
  \item
    Classify each blinded image card as Gram-positive, Gram-negative,
    mixed, or uninterpretable.
  \item
    Record: colour (purple = Gram+, pink = Gram-), cell shape (coccus,
    bacillus, spirillum), arrangement (single, pairs, chains, clusters),
    and confidence score (1-5).
  \item
    Compare independent scorers. Resolve disagreements primarily after
    recording the original calls.
  \end{itemize}
\end{enumerate}

\subsubsection*{Part 2: Computational Biology Exercise - Bacterial
Growth Analysis with
Python}\label{part-2-computational-biology-exercise---bacterial-growth-analysis-with-python}
\addcontentsline{toc}{subsubsection}{Part 2: Computational Biology
Exercise - Bacterial Growth Analysis with Python}

\emph{Optional computational check: run this self-contained Python
snippet from the project root. It uses tested \texttt{src/biology}
modules and requires no external notebook or CSV file.}

Use the microbiology module to check growth and MIC calculations.

\begin{Shaded}
\begin{Highlighting}[]
\ImportTok{from}\NormalTok{ biology.microbiology }\ImportTok{import}\NormalTok{ bacterial\_growth\_curve, doubling\_time, mic\_fold\_dilution}

\NormalTok{growth }\OperatorTok{=}\NormalTok{ bacterial\_growth\_curve(N0}\OperatorTok{=}\DecValTok{1\_000}\NormalTok{, doubling\_time\_hr}\OperatorTok{=}\FloatTok{0.5}\NormalTok{, t\_end\_hr}\OperatorTok{=}\FloatTok{4.0}\NormalTok{)}
\BuiltInTok{print}\NormalTok{(}\StringTok{"final population:"}\NormalTok{, }\BuiltInTok{int}\NormalTok{(growth.populations[}\OperatorTok{{-}}\DecValTok{1}\NormalTok{]))}
\BuiltInTok{print}\NormalTok{(}\StringTok{"doubling time from counts:"}\NormalTok{, }\BuiltInTok{round}\NormalTok{(doubling\_time(}\DecValTok{1\_000}\NormalTok{, }\DecValTok{16\_000}\NormalTok{, }\FloatTok{2.0}\NormalTok{), }\DecValTok{2}\NormalTok{))}
\BuiltInTok{print}\NormalTok{(}\StringTok{"MIC dilution series:"}\NormalTok{, mic\_fold\_dilution(}\FloatTok{128.0}\NormalTok{, dilution\_factor}\OperatorTok{=}\DecValTok{2}\NormalTok{, n\_tubes}\OperatorTok{=}\DecValTok{5}\NormalTok{))}
\end{Highlighting}
\end{Shaded}

\subsubsection*{Part 3: Serial Dilution and Colony Counting from Plate
Images}\label{part-3-serial-dilution-and-colony-counting-from-plate-images}
\addcontentsline{toc}{subsubsection}{Part 3: Serial Dilution and Colony
Counting from Plate Images}

\begin{enumerate}
\def\labelenumi{\arabic{enumi}.}
\setcounter{enumi}{4}
\item
  \textbf{Design an Alternative Investigation}: Instead of performing a
  standard dilution series, design a data-card investigation to test how
  different environmental conditions (temperature, pH, salt
  concentration) affect bacterial growth rates. What hypothesis would
  you test? Which printed curve or plate-image evidence would measure
  growth?
\item
  \textbf{Hypothesis}: Predict how increasing salt concentration would
  affect the growth of \emph{E. coli}. Would you expect a linear
  decrease in growth rate, or a threshold effect? Why?
\end{enumerate}

\subsubsection*{Part 4: AMR Network and Fungal Triage
Cards}\label{part-4-amr-network-and-fungal-triage-cards}
\addcontentsline{toc}{subsubsection}{Part 4: AMR Network and Fungal
Triage Cards}

\begin{enumerate}
\def\labelenumi{\arabic{enumi}.}
\setcounter{enumi}{6}
\tightlist
\item
  Build an AMR transfer map from the provided cards. For each arrow,
  label whether it represents selection, horizontal gene transfer,
  environmental persistence, or surveillance.
\item
  Use the \emph{Candida auris} card to decide whether the scenario calls
  for treatment, infection-control action, both, or neither. Justify
  your answer using colonization status, clinical infection evidence,
  identification method, and susceptibility data.
\item
  Compare the bacterial AMR map with the fungal triage card. Identify
  one shared One Health principle and one difference between bacterial
  resistance-gene transfer and healthcare-associated fungal persistence.
\end{enumerate}

\subsection*{Data Recording}\label{data-recording-21}
\addcontentsline{toc}{subsection}{Data Recording}

{\def\LTcaptype{none} % do not increment counter
\begin{longtable}[]{@{}llll@{}}
\toprule\noalign{}
Species & Gram stain result & Cell shape & Cell arrangement \\
\midrule\noalign{}
\endhead
\bottomrule\noalign{}
\endlastfoot
\emph{E. coli} & & & \\
\emph{Bacillus sp.} & & & \\
\end{longtable}
}

{\def\LTcaptype{none} % do not increment counter
\begin{longtable}[]{@{}lll@{}}
\toprule\noalign{}
Dilution & Colonies counted & CFU/mL \\
\midrule\noalign{}
\endhead
\bottomrule\noalign{}
\endlastfoot
10\texorpdfstring{\textsuperscript{-}}{-}\texorpdfstring{\textsuperscript{4}}{4} & & \\
10\texorpdfstring{\textsuperscript{-}}{-}\texorpdfstring{\textsuperscript{5}}{5} & & \\
10\texorpdfstring{\textsuperscript{-}}{-}\texorpdfstring{\textsuperscript{6}}{6} & & \\
\end{longtable}
}

Best estimate of original concentration: \_\_\_\_\_\_\_ CFU/mL Blinded
duplicate agreement: \_\_\_\_\_\_\_\%; plate images excluded, with
reason: \_\_\_\_\_\_\_ AMR transfer pathway most supported:
\_\_\_\_\_\_\_; \emph{Candida auris} action selected: \_\_\_\_\_\_\_;
surveillance evidence needed: \_\_\_\_\_\_\_

\subsection*{Evidence and Reproducibility
Checklist}\label{evidence-and-reproducibility-checklist-24}
\addcontentsline{toc}{subsection}{Evidence and Reproducibility
Checklist}

\begin{itemize}
\tightlist
\item
  \textbf{Primary evidence goal:} Measure growth and compare microbial
  life strategies.
\item
  \textbf{Data skill to practice:} Interpret microbial observations from
  growth, sequence, or structural evidence.
\item
  \textbf{BioSkills emphasis:} Science and society, Process of science,
  Modeling and simulation.
\item
  \textbf{Control logic:} identify at least one positive control, one
  negative control, or one baseline comparison before interpreting
  results.
\item
  \textbf{Measurement discipline:} record units, uncertainty, sample
  size, and any discarded observation with a reason.
\item
  \textbf{Mechanistic link:} connect one result directly to the parent
  chapter's big idea before writing the conclusion.
\item
  \textbf{Reproducibility check:} state one procedural detail that
  another group would need in order to reproduce the result.
\end{itemize}

\subsection*{Paper-Based Evidence
Upgrade}\label{paper-based-evidence-upgrade-24}
\addcontentsline{toc}{subsection}{Paper-Based Evidence Upgrade}

Before answering the analysis questions, annotate the paper dataset for
\textbf{Bacteria, Archaea, and Viruses} with a reproducibility pass:

{\def\LTcaptype{none} % do not increment counter
\begin{longtable}[]{@{}
  >{\raggedright\arraybackslash}p{(\linewidth - 2\tabcolsep) * \real{0.5000}}
  >{\raggedright\arraybackslash}p{(\linewidth - 2\tabcolsep) * \real{0.5000}}@{}}
\toprule\noalign{}
\begin{minipage}[b]{\linewidth}\raggedright
Evidence check
\end{minipage} & \begin{minipage}[b]{\linewidth}\raggedright
Student action
\end{minipage} \\
\midrule\noalign{}
\endhead
\bottomrule\noalign{}
\endlastfoot
Control logic & Mark the comparison that functions as the baseline,
negative control, or reference case. \\
Uncertainty & Circle the row, card, diagram feature, or model assumption
most likely to change the conclusion. \\
Model comparison & State whether a simpler rule, null model, or
alternative mechanism could explain the same pattern. \\
Decision threshold & Write the minimum evidence that would make you
revise the interpretation. \\
Reproducibility & Record the exact scoring rule another group would need
to reproduce your classification. \\
\end{longtable}
}

Focus note: in this lab, treat each microbe as an evidence claim built
from a Gram-stain phenotype, a colony or plaque count, and a
wall-structure or phage-genetics inference, and state for every
conclusion whether it rests on a cultured isolate or on a
sequence/diagram on paper. Keep required work paper-based; any
material-handling or equipment version belongs only in an optional
extension.

\subsection*{Analysis Questions}\label{analysis-questions-24}
\addcontentsline{toc}{subsection}{Analysis Questions}

\begin{enumerate}
\def\labelenumi{\arabic{enumi}.}
\item
  Why does \emph{E. coli} stain Gram-negative? Describe the structural
  difference between the cell walls of Gram-positive and Gram-negative
  bacteria, and explain why ethanol decolourises Gram-negative but not
  Gram-positive cells.
\item
  Calculate the CFU/mL from your best printed plate count. Assuming
  exponential growth at a doubling time of 20 minutes, estimate how many
  cells would be present after 4 hours if starting from 100 CFU/mL.
\item
  Penicillin inhibits transpeptidase (cross-links peptidoglycan). Based
  on cell wall structure, predict whether penicillin is more effective
  against \emph{E. coli} or \emph{Bacillus subtilis}, and why
  Gram-negative bacteria often require higher doses.
\item
  A bacteriophage follows a lysogenic cycle for 200 generations, then
  enters the lytic cycle. What molecular signal might trigger the switch
  from lysogeny to lysis? Name the specific protease involved and the CI
  repressor it degrades.
\item
  SARS-CoV-2 is a single-stranded, positive-sense RNA virus with an
  envelope. Compare its replication cycle to a lytic bacteriophage:
  identify two key differences in where replication occurs, the host
  enzyme used, and the method of exiting the host cell.
\item
  A carbapenemase gene appears in related \emph{Klebsiella} isolates on
  two hospital wards and in a wastewater sample. Use plasmid, integron,
  and clonal-spread evidence to decide what additional data would
  distinguish horizontal transfer from one expanding clone.
\end{enumerate}

\subsection*{Extension Analysis
Questions}\label{extension-analysis-questions-12}
\addcontentsline{toc}{subsection}{Extension Analysis Questions}

\begin{enumerate}
\def\labelenumi{\arabic{enumi}.}
\setcounter{enumi}{6}
\item
  \textbf{Statistical Analysis}: If the dataset includes three replicate
  plate images, how would you analyze the variation in CFU/mL estimates?
  What statistical test would you use to determine if differences
  between species are significant?
\item
  \textbf{Experimental Error Analysis}: What are three potential sources
  of error in bacterial colony counting from images? How could you
  modify the scoring rules or image set to reduce each source of error?
\item
  \textbf{Experimental Design}: If you wanted to test the effectiveness
  of different antibiotics on bacterial growth without culturing, how
  would you design the dataset? What controls would you include? How
  would you measure the effect from growth curves, MIC tables, or
  inhibition-zone images?
\end{enumerate}

\subsection*{Group Project Extension
(Multi-Session)}\label{group-project-extension-multi-session-9}
\addcontentsline{toc}{subsection}{Group Project Extension
(Multi-Session)}

\textbf{Design a Study of Antibiotic Resistance Evolution}

This project will span 2-3 lab sessions:

\textbf{Session 1}: Research Question and Hypothesis - How does
sub-lethal antibiotic concentration affect the rate of resistance
evolution in bacteria? - Formulate a hypothesis about the relationship
between antibiotic concentration and resistance development - Design a
simulation using provided \emph{E. coli} growth curves and different
antibiotic concentrations

\textbf{Session 2}: Data Collection and Analysis - Analyze printed
growth curves collected under antibiotic gradients - Measure resistance
levels from minimum inhibitory concentration tables - Analyze
allele-frequency changes using provided PCR or sequencing summaries

\textbf{Session 3}: Data Interpretation and Presentation - Compare
resistance evolution across different antibiotic concentrations -
Calculate selection coefficients for resistant vs sensitive bacteria -
Create a scientific poster or presentation

\subsection*{Real-World Problem Solving: Microbial
Ethics}\label{real-world-problem-solving-microbial-ethics}
\addcontentsline{toc}{subsection}{Real-World Problem Solving: Microbial
Ethics}

\textbf{Case Study: Antibiotic Use in Agriculture}

\begin{enumerate}
\def\labelenumi{\arabic{enumi}.}
\item
  \textbf{Research Task}: Investigate the use of antibiotics as growth
  promoters in livestock. What percentage of antibiotics sold in the US
  are used for this purpose? What are the most common antibiotics used?
\item
  \textbf{Ethical Consideration}: Should the routine use of antibiotics
  for growth promotion in agriculture be banned? What are the potential
  economic impacts on farmers vs the public health benefits?
\item
  \textbf{Policy Proposal}: Design a policy to reduce antibiotic use in
  agriculture while ensuring animal health and farmer livelihoods.
  Consider: alternatives to antibiotics, monitoring systems, and
  transition assistance for farmers.
\end{enumerate}

\subsection*{Safety and Ethics Notes}\label{safety-and-ethics-notes-24}
\addcontentsline{toc}{subsection}{Safety and Ethics Notes}

Default lab work uses printed image cards, growth curves, and decision
matrices primarily. No cultures, plates, antibiotic discs, or
environmental samples are required. Optional wet microbiology
demonstrations must use approved BSL-1 organisms under instructor
supervision and local disposal rules. When discussing antibiotic
resistance, emphasize following current prescriber instructions,
avoiding unnecessary antibiotic starts, using diagnostics when
available, and protecting access to effective treatment.

\subsection*{Debrief and Reflection}\label{debrief-and-reflection-21}
\addcontentsline{toc}{subsection}{Debrief and Reflection}

After you finish the practical work, spend 5--10 minutes in your small
group comparing results and discussing the following prompts. Each
member should contribute at least one observation before moving to the
next prompt:

\begin{enumerate}
\def\labelenumi{\arabic{enumi}.}
\item
  \textbf{What did your measurements show} --- compare the group's
  results to the textbook's predictions. Where they diverge, suggest at
  least one mechanistic explanation before concluding ``experimental
  error.''
\item
  \textbf{What would change the outcome} --- propose one modification to
  the procedure that would sharpen the measurement or extend the result
  to a new biological context, and predict what you would observe.
\item
  \textbf{One-sentence headline} --- each student composes a single
  sentence summarising the lab's take-home message, suitable for a
  tweet. Compare sentences across groups; good headlines are short,
  quantitative, and mechanistic.
\item
  \textbf{Connection back to the textbook} --- identify one section of
  \cref{sec:unit_VII_bacteria_archaea_viruses} that your data either
  confirmed or complicated. Cite the specific passage.
\item
  \textbf{Experimental Design Reflection}: What was the strongest aspect
  of your experimental design? What would you change if you could repeat
  the experiment?
\end{enumerate}

\subsection*{Further Reading (Lab)}\label{further-reading-lab-21}
\addcontentsline{toc}{subsection}{Further Reading (Lab)}

\begin{itemize}
\tightlist
\item
  Revisit the parent chapter
  \cref{sec:unit_VII_bacteria_archaea_viruses} for the theoretical
  foundations on which this lab is built.
\item
  Look up any \textbf{bolded glossary term} introduced in the textbook
  chapter (each has a \texttt{\#gl:term-slug} link in the text) --- its
  master definition is in \breaktt{manuscript/glossary.md}.
\item
  Explore the Python code in
  \breaktt{src/biology/microbiology/microbiology.py} for bacterial growth
  calculations.
\item
  Use the self-contained Part 2 snippet as the computational template;
  it runs against tested project modules without external notebooks or
  CSV files.
\end{itemize}

\newpage

\section*{Lab 23 --- Microbial Ecology and the
Microbiome}\label{lab-23-microbial-ecology-and-the-microbiome}
\addcontentsline{toc}{section}{Lab 23 --- Microbial Ecology and the
Microbiome}

\label{sec:lab_unit_VII_microbial_ecology}

\emph{This activity accompanies \cref{sec:unit_VII_microbial_ecology} of
the textbook --- review that chapter before attempting the exercises
below.}

\subsection*{Learning Objectives}\label{learning-objectives-63}
\addcontentsline{toc}{subsection}{Learning Objectives}

\begin{itemize}
\tightlist
\item
  Design and evaluate a microbial community analysis using printed
  culture-independent datasets
\item
  Formulate testable hypotheses about how environmental factors affect
  microbial diversity
\item
  Identify independent, dependent, and controlled variables in microbial
  ecology experiments
\item
  Analyze experimental data using diversity indices and network analysis
\item
  Write a brief scientific report with hypothesis, methods, results, and
  conclusion
\end{itemize}

\subsubsection*{Alignment and Rubric
Map}\label{alignment-and-rubric-map-25}
\addcontentsline{toc}{subsubsection}{Alignment and Rubric Map}

\begin{itemize}
\tightlist
\item
  \textbf{Outcome 1 (LO1):} Interpret the supplied evidence or model
  output for Microbial Ecology and the Microbiome.
\item
  \textbf{Outcome 2 (LO2):} Identify controls and comparison groups that
  make the claim testable.
\item
  \textbf{Outcome 3 (LO3):} Quantify uncertainty, boundary conditions,
  or alternative explanations before concluding.
\item
  \textbf{Outcome 4 (LO4):} Transfer the mechanism to a new biological
  case or public-facing decision.
\item
  \textbf{Chapter LO coverage:} LO1, LO2, LO3, LO4
\item
  \textbf{Rubric dimensions:} evidence; controls; uncertainty;
  mechanism; transfer. \#\# Pre-Lab Inquiry Questions \{.unnumbered\}
\end{itemize}

\emph{Complete these questions before coming to lab. Use your textbook
and additional research to inform your hypotheses.}

\begin{enumerate}
\def\labelenumi{\arabic{enumi}.}
\item
  \textbf{Hypothesis Formation}: Based on your understanding of
  microbial ecology, predict which provided community dataset (soil,
  skin, tap water, or gut) will have the highest microbial diversity.
  Write a clear, testable hypothesis.
\item
  \textbf{Experimental Design}: What is the independent variable in a
  microbial diversity study? What is the dependent variable? List at
  least 5 variables that should be controlled to ensure valid results.
\item
  \textbf{Scientific Context}: The Shannon diversity index (H')
  incorporates both richness and evenness. Why is this measure important
  in microbial ecology? How does it differ from simply counting species?
\item
  \textbf{Quantitative Reasoning}: If a community has three species with
  relative abundances of 0.5, 0.3, and 0.2, what is its Shannon
  diversity index? Show your calculation.
\item
  \textbf{Real-World Application}: The human gut microbiome is linked to
  health and disease. How might antibiotic treatment reduce microbial
  diversity, and what are the potential health consequences of this
  reduction?
\end{enumerate}

\subsection*{Background}\label{background-13}
\addcontentsline{toc}{subsection}{Background}

The human gut microbiome contains about 3.8 \texorpdfstring{\ensuremath{\times}}{x} 10\texorpdfstring{\textsuperscript{1}}{1}\texorpdfstring{\textsuperscript{3}}{3} bacteria
from hundreds of species; the ratio of microbial to human cells is
approximately 1:1 (revised from earlier overestimates). Culture-based
methods recover about 1\% of environmental microbes (the ``great plate
count anomaly''). Metagenomics and 16S rRNA gene sequencing (amplifying
the hypervariable V3--V4 region) now allow culture-independent
assessment of community composition. The \textbf{Shannon diversity index
H' = \texorpdfstring{\ensuremath{-}}{-}\texorpdfstring{\ensuremath{\Sigma}}{Sigma}(p\_i \texorpdfstring{\ensuremath{\times}}{x} ln p\_i)} integrates species richness and evenness.

\subsection*{Paper-Based Materials}\label{paper-based-materials-25}
\addcontentsline{toc}{subsection}{Paper-Based Materials}

{\def\LTcaptype{none} % do not increment counter
\begin{longtable}[]{@{}
  >{\raggedright\arraybackslash}p{(\linewidth - 2\tabcolsep) * \real{0.3333}}
  >{\raggedright\arraybackslash}p{(\linewidth - 2\tabcolsep) * \real{0.6667}}@{}}
\toprule\noalign{}
\begin{minipage}[b]{\linewidth}\raggedright
Item
\end{minipage} & \begin{minipage}[b]{\linewidth}\raggedright
Quantity
\end{minipage} \\
\midrule\noalign{}
\endhead
\bottomrule\noalign{}
\endlastfoot
Printed community case packet (soil, skin, tap water, gut) with metadata
cards & 1 \\
Printed colony-image panel from archived BSL-1 teaching plates & 1 \\
Printed 16S rRNA OTU table (two samples: healthy gut vs
antibiotic-treated gut --- 10 taxa, abundance data) & 1 \\
Calculator & 1 \\
Coloured pencils & 3 \\
Interaction network diagram template (printed) & 1 \\
Calculator or optional Python REPL with this project installed & 1 \\
Replicate read-count table and rarefaction checklist & 1 \\
\end{longtable}
}

\subsection*{Paper-Based
Investigation}\label{paper-based-investigation-25}
\addcontentsline{toc}{subsection}{Paper-Based Investigation}

\subsubsection*{Part 1: Community Diversity from Case
Packets}\label{part-1-community-diversity-from-case-packets}
\addcontentsline{toc}{subsubsection}{Part 1: Community Diversity from
Case Packets}

\begin{enumerate}
\def\labelenumi{\arabic{enumi}.}
\item
  \textbf{Formulate Your Hypothesis}: Based on your pre-lab research,
  write a hypothesis about how the provided community datasets will
  differ in microbial diversity.
\item
  \textbf{Identify Variables}:

  \begin{itemize}
  \tightlist
  \item
    Independent variable: Community source represented in the case
    packet
  \item
    Dependent variable: Taxon richness, Shannon diversity index, and
    network connectance
  \item
    Controlled variables: sequencing depth, primer region, read-quality
    threshold, rarefaction rule, metadata completeness, and replicate
    count.
  \end{itemize}
\item
  \textbf{Set Up Controls}: Why is it important to use the same
  sequencing-depth rule for most samples? What other controls could you
  include, such as a mock community, a blank extraction control, or a
  duplicate sample?
\item
  \textbf{Data Collection}:

  \begin{itemize}
  \tightlist
  \item
    Use the printed OTU table to calculate total reads, relative
    abundance, richness, Shannon diversity, and evenness for each
    sample.
  \item
    Use the archived colony-image panel primarily as a comparison for
    why morphology underestimates molecular diversity.
  \item
    Check whether each sample has enough reads for fair comparison; if
    not, apply the rarefaction rule provided in the case packet.
  \item
    Compare at least three replicate rows per sample type before drawing
    conclusions.
  \end{itemize}
\end{enumerate}

\subsubsection*{Part 2: Computational Biology Exercise - Diversity
Analysis with
Python}\label{part-2-computational-biology-exercise---diversity-analysis-with-python}
\addcontentsline{toc}{subsubsection}{Part 2: Computational Biology
Exercise - Diversity Analysis with Python}

\emph{Optional computational check: run this self-contained Python
snippet from the project root. It uses tested \texttt{src/biology}
modules and requires no external notebook or CSV file.}

Calculate diversity from table counts and compare named reference
organisms.

\begin{Shaded}
\begin{Highlighting}[]
\ImportTok{from}\NormalTok{ biology.ecology }\ImportTok{import}\NormalTok{ biodiversity\_indices}
\ImportTok{from}\NormalTok{ biology.microbiology }\ImportTok{import}\NormalTok{ REFERENCE\_ORGANISMS}

\NormalTok{counts }\OperatorTok{=}\NormalTok{ [}\DecValTok{42}\NormalTok{, }\DecValTok{18}\NormalTok{, }\DecValTok{15}\NormalTok{, }\DecValTok{9}\NormalTok{, }\DecValTok{6}\NormalTok{]}
\NormalTok{diversity }\OperatorTok{=}\NormalTok{ biodiversity\_indices(counts)}
\BuiltInTok{print}\NormalTok{(}\StringTok{"Shannon:"}\NormalTok{, }\BuiltInTok{round}\NormalTok{(diversity.shannon\_index, }\DecValTok{3}\NormalTok{))}
\BuiltInTok{print}\NormalTok{(}\StringTok{"evenness:"}\NormalTok{, }\BuiltInTok{round}\NormalTok{(diversity.evenness, }\DecValTok{3}\NormalTok{))}
\BuiltInTok{print}\NormalTok{(}\StringTok{"reference domains:"}\NormalTok{, }\BuiltInTok{sorted}\NormalTok{(\{organism.domain }\ControlFlowTok{for}\NormalTok{ organism }\KeywordTok{in}\NormalTok{ REFERENCE\_ORGANISMS\}))}
\end{Highlighting}
\end{Shaded}

\subsubsection*{Part 3: Microbial Interaction Network
Analysis}\label{part-3-microbial-interaction-network-analysis}
\addcontentsline{toc}{subsubsection}{Part 3: Microbial Interaction
Network Analysis}

\begin{enumerate}
\def\labelenumi{\arabic{enumi}.}
\setcounter{enumi}{4}
\item
  \textbf{Design an Alternative Investigation}: Instead of collecting
  microbiome samples, design a paper-based analysis to test how
  antibiotic treatment affects microbial interaction networks. What
  hypothesis would you test? How would you construct the network from
  co-occurrence or time-series data?
\item
  \textbf{Hypothesis}: Predict how the connectance (number of
  interactions) in a microbial network would change with reduced
  diversity. Would you expect more or fewer interactions? Why?
\end{enumerate}

\subsection*{Data Recording}\label{data-recording-22}
\addcontentsline{toc}{subsection}{Data Recording}

{\def\LTcaptype{none} % do not increment counter
\begin{longtable}[]{@{}lllll@{}}
\toprule\noalign{}
Sample source & Taxon richness & Shannon H' & Evenness & Read depth \\
\midrule\noalign{}
\endhead
\bottomrule\noalign{}
\endlastfoot
Soil & & & & \\
Skin & & & & \\
Tap water & & & & \\
Gut & & & & \\
\end{longtable}
}

{\def\LTcaptype{none} % do not increment counter
\begin{longtable}[]{@{}
  >{\raggedright\arraybackslash}p{(\linewidth - 10\tabcolsep) * \real{0.0549}}
  >{\raggedright\arraybackslash}p{(\linewidth - 10\tabcolsep) * \real{0.2198}}
  >{\raggedright\arraybackslash}p{(\linewidth - 10\tabcolsep) * \real{0.1429}}
  >{\raggedright\arraybackslash}p{(\linewidth - 10\tabcolsep) * \real{0.1429}}
  >{\raggedright\arraybackslash}p{(\linewidth - 10\tabcolsep) * \real{0.2527}}
  >{\raggedright\arraybackslash}p{(\linewidth - 10\tabcolsep) * \real{0.1868}}@{}}
\toprule\noalign{}
\begin{minipage}[b]{\linewidth}\raggedright
Taxon
\end{minipage} & \begin{minipage}[b]{\linewidth}\raggedright
Healthy sample reads
\end{minipage} & \begin{minipage}[b]{\linewidth}\raggedright
p\_i (healthy)
\end{minipage} & \begin{minipage}[b]{\linewidth}\raggedright
p\_i \texorpdfstring{\ensuremath{\times}}{x} ln(p\_i)
\end{minipage} & \begin{minipage}[b]{\linewidth}\raggedright
Antibiotic sample reads
\end{minipage} & \begin{minipage}[b]{\linewidth}\raggedright
p\_i (antibiotic)
\end{minipage} \\
\midrule\noalign{}
\endhead
\bottomrule\noalign{}
\endlastfoot
& & & & & \\
& & & & & \\
& & & & & \\
& & & & & \\
& & & & & \\
& & & & & \\
& & & & & \\
& & & & & \\
& & & & & \\
& & & & & \\
\end{longtable}
}

H' (healthy): \_\_\_\_; H' (antibiotic): \_\_\_\_

Keystone taxon in network: \textbf{\emph{; Rationale: }} Rarefaction or
exclusion decision: \textbf{\emph{; Reproducibility concern: }}

\subsection*{Evidence and Reproducibility
Checklist}\label{evidence-and-reproducibility-checklist-25}
\addcontentsline{toc}{subsection}{Evidence and Reproducibility
Checklist}

\begin{itemize}
\tightlist
\item
  \textbf{Primary evidence goal:} Compare microbial communities and
  infer interaction hypotheses.
\item
  \textbf{Data skill to practice:} Compute or interpret community
  metrics from abundance data.
\item
  \textbf{BioSkills emphasis:} Science and society, Process of science,
  Modeling and simulation.
\item
  \textbf{Control logic:} identify at least one positive control, one
  negative control, or one baseline comparison before interpreting
  results.
\item
  \textbf{Measurement discipline:} record units, uncertainty, sample
  size, and any discarded observation with a reason.
\item
  \textbf{Mechanistic link:} connect one result directly to the parent
  chapter's big idea before writing the conclusion.
\item
  \textbf{Reproducibility check:} state one procedural detail that
  another group would need in order to reproduce the result.
\end{itemize}

\subsection*{Paper-Based Evidence
Upgrade}\label{paper-based-evidence-upgrade-25}
\addcontentsline{toc}{subsection}{Paper-Based Evidence Upgrade}

Before answering the analysis questions, annotate the paper dataset for
\textbf{Microbial Ecology and the Microbiome} with a reproducibility
pass:

{\def\LTcaptype{none} % do not increment counter
\begin{longtable}[]{@{}
  >{\raggedright\arraybackslash}p{(\linewidth - 2\tabcolsep) * \real{0.5000}}
  >{\raggedright\arraybackslash}p{(\linewidth - 2\tabcolsep) * \real{0.5000}}@{}}
\toprule\noalign{}
\begin{minipage}[b]{\linewidth}\raggedright
Evidence check
\end{minipage} & \begin{minipage}[b]{\linewidth}\raggedright
Student action
\end{minipage} \\
\midrule\noalign{}
\endhead
\bottomrule\noalign{}
\endlastfoot
Control logic & Mark the comparison that functions as the baseline,
negative control, or reference case. \\
Uncertainty & Circle the row, card, diagram feature, or model assumption
most likely to change the conclusion. \\
Model comparison & State whether a simpler rule, null model, or
alternative mechanism could explain the same pattern. \\
Decision threshold & Write the minimum evidence that would make you
revise the interpretation. \\
Reproducibility & Record the exact scoring rule another group would need
to reproduce your classification. \\
\end{longtable}
}

Focus note: in this lab, separate community-composition correlations
from mechanistic claims by recording the diversity index or count, the
host or environmental context, and the perturbation, and state the
experiment that would be needed to move a correlation toward causation.
Keep required work paper-based; any material-handling or equipment
version belongs only in an optional extension.

\subsection*{Analysis Questions}\label{analysis-questions-25}
\addcontentsline{toc}{subsection}{Analysis Questions}

\begin{enumerate}
\def\labelenumi{\arabic{enumi}.}
\item
  The soil dataset likely shows higher richness than the skin or
  tap-water dataset. Why are soils among the most microbially diverse
  environments on Earth? What physical and chemical habitat
  heterogeneity drives this?
\item
  Antibiotic treatment usually reduces H'. What does low H' indicate
  about community stability? Name one specific clinical consequence of
  reduced microbiome diversity following broad-spectrum antibiotic use.
\item
  The ``great plate count anomaly'' refers to the fact that about 1\% of
  environmental microbes form colonies on standard media. Propose three
  reasons why the remaining 99\% are ``unculturable'' and describe one
  molecular technique that can characterise them without culturing.
\item
  \emph{Akkermansia muciniphila} (a gut commensal) degrades mucin and
  produces acetate and propionate consumed by other bacteria. If
  \emph{Akkermansia} is lost, predict the cascade effects on adjacent
  taxa in your interaction network.
\item
  A patient receives a faecal microbiota transplant (FMT) to treat
  recurrent \emph{C. difficile} infection. Using the concept of
  competitive exclusion and microbiome diversity, explain the mechanism
  by which FMT suppresses \emph{C. difficile} and describe what the
  ideal donor microbiome looks like.
\end{enumerate}

\subsection*{Extension Analysis
Questions}\label{extension-analysis-questions-13}
\addcontentsline{toc}{subsection}{Extension Analysis Questions}

\begin{enumerate}
\def\labelenumi{\arabic{enumi}.}
\setcounter{enumi}{5}
\item
  \textbf{Statistical Analysis}: If the dataset includes three replicate
  OTU tables per source, how would you analyze the variation in Shannon
  diversity? What statistical test would you use to determine if
  differences between sample types are significant?
\item
  \textbf{Experimental Error Analysis}: What are three potential sources
  of error in microbial diversity measurements? How could you modify the
  procedure to reduce each source of error?
\item
  \textbf{Experimental Design}: If you wanted to test whether a specific
  environmental factor (e.g., pH, moisture) affects microbial diversity,
  how would you design the experiment? What data would you collect? How
  would you analyze it?
\end{enumerate}

\subsection*{Group Project Extension
(Multi-Session)}\label{group-project-extension-multi-session-10}
\addcontentsline{toc}{subsection}{Group Project Extension
(Multi-Session)}

\textbf{Design a Study of Microbiome and Health}

This project will span 2-3 lab sessions:

\textbf{Session 1}: Research Question and Hypothesis - How does diet
affect gut microbiome diversity? - Formulate a hypothesis about the
relationship between fiber intake and microbial diversity - Design a
study plan using dietary data and public microbiome datasets

\textbf{Session 2}: Data Collection and Analysis - Use dietary surveys
and public microbiome datasets - Analyze 16S rRNA sequencing data to
calculate diversity indices - Compare diversity across different diet
groups

\textbf{Session 3}: Data Interpretation and Presentation - Correlate
diversity with specific dietary components - Discuss the implications
for human health - Create a scientific poster or presentation

\subsection*{Real-World Problem Solving: Microbial
Ethics}\label{real-world-problem-solving-microbial-ethics-1}
\addcontentsline{toc}{subsection}{Real-World Problem Solving: Microbial
Ethics}

\textbf{Case Study: Microbiome Biobanks}

\begin{enumerate}
\def\labelenumi{\arabic{enumi}.}
\item
  \textbf{Research Task}: Investigate the ethical, legal, and social
  implications (ELSI) of collecting and storing human microbiome
  samples. What privacy concerns exist? How is microbiome data different
  from genetic data?
\item
  \textbf{Ethical Consideration}: Should individuals have property
  rights over their microbiome? What are the potential benefits and
  risks of commercialising microbiome-based therapies?
\item
  \textbf{Policy Proposal}: Design a policy for the ethical collection
  and use of human microbiome samples. Consider: informed consent, data
  privacy, benefit-sharing, and commercialisation.
\end{enumerate}

\subsection*{Safety and Ethics Notes}\label{safety-and-ethics-notes-25}
\addcontentsline{toc}{subsection}{Safety and Ethics Notes}

Default lab work uses printed OTU tables, archived colony images,
metadata cards, and network diagrams primarily. No environmental
sampling, swabbing, agar plates, incubation, or culturing is required.
Optional sampling demonstrations must be instructor-supervised and
follow local biosafety rules. When discussing microbiome data, respect
participant privacy and follow data protection guidelines.

\subsection*{Debrief and Reflection}\label{debrief-and-reflection-22}
\addcontentsline{toc}{subsection}{Debrief and Reflection}

After you finish the practical work, spend 5--10 minutes in your small
group comparing results and discussing the following prompts. Each
member should contribute at least one observation before moving to the
next prompt:

\begin{enumerate}
\def\labelenumi{\arabic{enumi}.}
\item
  \textbf{What did your measurements show} --- compare the group's
  results to the textbook's predictions. Where they diverge, suggest at
  least one mechanistic explanation before concluding ``experimental
  error.''
\item
  \textbf{What would change the outcome} --- propose one modification to
  the procedure that would sharpen the measurement or extend the result
  to a new biological context, and predict what you would observe.
\item
  \textbf{One-sentence headline} --- each student composes a single
  sentence summarising the lab's take-home message, suitable for a
  tweet. Compare sentences across groups; good headlines are short,
  quantitative, and mechanistic.
\item
  \textbf{Connection back to the textbook} --- identify one section of
  \cref{sec:unit_VII_microbial_ecology} that your data either confirmed
  or complicated. Cite the specific passage.
\item
  \textbf{Experimental Design Reflection}: What was the strongest aspect
  of your experimental design? What would you change if you could repeat
  the experiment?
\end{enumerate}

\subsection*{Further Reading (Lab)}\label{further-reading-lab-22}
\addcontentsline{toc}{subsection}{Further Reading (Lab)}

\begin{itemize}
\tightlist
\item
  Revisit the parent chapter \cref{sec:unit_VII_microbial_ecology} for
  the theoretical foundations on which this lab is built.
\item
  Look up any \textbf{bolded glossary term} introduced in the textbook
  chapter (each has a \texttt{\#gl:term-slug} link in the text) --- its
  master definition is in \breaktt{manuscript/glossary.md}.
\item
  Explore the Python code in
  \breaktt{src/biology/microbiology/microbiology.py} for microbial
  ecology calculations.
\item
  Use the self-contained Part 2 snippet as the computational template;
  it runs against tested project modules without external notebooks or
  CSV files.
\end{itemize}

\newpage

\section*{Lab 24 --- Infectious Disease and
Immunity}\label{lab-24-infectious-disease-and-immunity}
\addcontentsline{toc}{section}{Lab 24 --- Infectious Disease and
Immunity}

\label{sec:lab_unit_VII_infectious_disease}

\emph{This activity accompanies \cref{sec:unit_VII_infectious_disease}
of the textbook --- review that chapter before attempting the exercises
below.}

\subsection*{Learning Objectives}\label{learning-objectives-64}
\addcontentsline{toc}{subsection}{Learning Objectives}

\begin{itemize}
\tightlist
\item
  Design and evaluate an epidemic investigation dataset to determine
  transmission dynamics
\item
  Formulate testable hypotheses about how immunity affects disease
  spread
\item
  Identify independent, dependent, and controlled variables in
  epidemiological studies
\item
  Analyze experimental data using epidemic curves and herd immunity
  calculations
\item
  Write a brief scientific report with hypothesis, methods, results, and
  conclusion
\end{itemize}

\subsubsection*{Alignment and Rubric
Map}\label{alignment-and-rubric-map-26}
\addcontentsline{toc}{subsubsection}{Alignment and Rubric Map}

\begin{itemize}
\tightlist
\item
  \textbf{Outcome 1 (LO1):} Interpret the supplied evidence or model
  output for Infectious Disease and Immunity.
\item
  \textbf{Outcome 2 (LO2):} Identify controls and comparison groups that
  make the claim testable.
\item
  \textbf{Outcome 3 (LO3):} Quantify uncertainty, boundary conditions,
  or alternative explanations before concluding.
\item
  \textbf{Outcome 4 (LO4):} Transfer the mechanism to a new biological
  case or public-facing decision.
\item
  \textbf{Chapter LO coverage:} LO1, LO2, LO3, LO4
\item
  \textbf{Rubric dimensions:} evidence; controls; uncertainty;
  mechanism; transfer. \#\# Pre-Lab Inquiry Questions \{.unnumbered\}
\end{itemize}

\emph{Complete these questions before coming to lab. Use your textbook
and additional research to inform your hypotheses.}

\begin{enumerate}
\def\labelenumi{\arabic{enumi}.}
\item
  \textbf{Hypothesis Formation}: Based on your understanding of epidemic
  curves, predict what shape the curve would take for a point-source
  exposure (e.g., contaminated food) versus person-to-person
  transmission (e.g., influenza). Write a clear, testable hypothesis.
\item
  \textbf{Experimental Design}: What is the independent variable in an
  epidemic investigation? What is the dependent variable? List at least
  5 variables that should be controlled to ensure valid results.
\item
  \textbf{Scientific Context}: The basic reproduction number (R\texorpdfstring{\textsubscript{0}}{0}) is a
  key metric in epidemiology. Why is R\texorpdfstring{\textsubscript{0}}{0} important for understanding
  disease spread? How does it relate to the concept of herd immunity?
\item
  \textbf{Quantitative Reasoning}: If a disease has an R\texorpdfstring{\textsubscript{0}}{0} of 5, what
  fraction of the population needs to be immune to achieve herd
  immunity? Show your calculation using the formula p\_c = 1 - (1/R\texorpdfstring{\textsubscript{0}}{0}).
\item
  \textbf{Real-World Application}: Antibiotic resistance is a major
  public health concern. How does the overuse of antibiotics in medicine
  and agriculture contribute to the evolution of resistant bacteria?
  What policies could help address this issue?
\end{enumerate}

\subsection*{Background}\label{background-14}
\addcontentsline{toc}{subsection}{Background}

Epidemiology uses epidemic curves (cases vs time), attack rates,
secondary attack rates, and reproductive numbers (R\texorpdfstring{\textsubscript{0}}{0}) to characterise
outbreaks. R\texorpdfstring{\textsubscript{0}}{0} is the average number of secondary cases generated by one
infected individual in a fully susceptible population; herd immunity is
achieved when the fraction vaccinated exceeds the threshold \textbf{p\_c
= 1 \texorpdfstring{\ensuremath{-}}{-} (1/R\texorpdfstring{\textsubscript{0}}{0})}. Antimicrobial susceptibility datasets often report
disc-diffusion zones of inhibition from Kirby-Bauer assays --- larger
zone = more susceptible pathogen --- but this lab interprets printed
data primarily.

\subsection*{Paper-Based Materials}\label{paper-based-materials-26}
\addcontentsline{toc}{subsection}{Paper-Based Materials}

{\def\LTcaptype{none} % do not increment counter
\begin{longtable}[]{@{}
  >{\raggedright\arraybackslash}p{(\linewidth - 2\tabcolsep) * \real{0.3333}}
  >{\raggedright\arraybackslash}p{(\linewidth - 2\tabcolsep) * \real{0.6667}}@{}}
\toprule\noalign{}
\begin{minipage}[b]{\linewidth}\raggedright
Item
\end{minipage} & \begin{minipage}[b]{\linewidth}\raggedright
Quantity
\end{minipage} \\
\midrule\noalign{}
\endhead
\bottomrule\noalign{}
\endlastfoot
Epidemic curve dataset (printed: hypothetical norovirus outbreak,
72-hour window, 150 cases) & 1 \\
Disc diffusion data (printed: zones of inhibition for 5 antibiotics
against 2 bacterial strains) & 1 \\
Printed disc-diffusion plate images with blinded replicate labels & 1
set \\
AMR mechanism cards (enzyme, target modification, efflux, permeability)
& 1 set \\
TB regimen decision cards: drug-susceptible, MDR/RR-TB,
fluoroquinolone-resistant, toxicity-risk, and adherence-support
scenarios & 1 set \\
HIV PrEP comparison card: oral PrEP, cabotegravir, and twice-yearly
lenacapavir attributes & 1 \\
Malaria vector-control comparison card: bed nets, indoor residual
spraying, spatial emanators, vaccine, and chemoprevention & 1 \\
Calculator & 1 \\
Ruler (mm) & 1 \\
Herd immunity R\texorpdfstring{\textsubscript{0}}{0} worksheet & 1 \\
Calculator or optional Python REPL with this project installed & 1 \\
Epidemic simulation worksheet & 1 \\
\end{longtable}
}

\subsection*{Paper-Based
Investigation}\label{paper-based-investigation-26}
\addcontentsline{toc}{subsection}{Paper-Based Investigation}

\subsubsection*{Part 1: Epidemic Curve
Analysis}\label{part-1-epidemic-curve-analysis}
\addcontentsline{toc}{subsubsection}{Part 1: Epidemic Curve Analysis}

\begin{enumerate}
\def\labelenumi{\arabic{enumi}.}
\item
  \textbf{Formulate Your Hypothesis}: Based on your pre-lab research,
  write a hypothesis about the transmission dynamics of the norovirus
  outbreak.
\item
  \textbf{Identify Variables}:

  \begin{itemize}
  \tightlist
  \item
    Independent variable: Time (hours since exposure)
  \item
    Dependent variable: Number of new cases
  \item
    Controlled variables: Population size, exposure source, reporting
    methods, etc.
  \end{itemize}
\item
  \textbf{Set Up Controls}: Why is it important to have accurate case
  reporting? What other controls could you include in an epidemic
  investigation to ensure valid conclusions?
\item
  \textbf{Data Collection}:

  \begin{itemize}
  \tightlist
  \item
    Plot the epidemic curve (cases per 2-hour interval).
  \item
    Identify the single peak (suggesting point-source exposure) or
    multiple waves (suggesting person-to-person spread).
  \item
    Calculate: (a) attack rate = (cases / total exposed) \texorpdfstring{\ensuremath{\times}}{x} 100\%; (b)
    approximate incubation period from the peak; (c) R\texorpdfstring{\textsubscript{0}}{0} from the growth
    rate during exponential phase using the provided formula.
  \item
    Compare two blinded disc-diffusion image replicates, measure zone
    diameters from the printed scale, and classify each antibiotic as
    susceptible, intermediate, or resistant using the provided
    breakpoint table.
  \end{itemize}
\end{enumerate}

\subsubsection*{Part 2: Computational Biology Exercise - Epidemic
Modeling with
Python}\label{part-2-computational-biology-exercise---epidemic-modeling-with-python}
\addcontentsline{toc}{subsubsection}{Part 2: Computational Biology
Exercise - Epidemic Modeling with Python}

\emph{Optional computational check: run this self-contained Python
snippet from the project root. It uses tested \texttt{src/biology}
modules and requires no external notebook or CSV file.}

Model outbreak curves with the tested SIR implementation.

\begin{Shaded}
\begin{Highlighting}[]
\ImportTok{from}\NormalTok{ biology.microbiology }\ImportTok{import}\NormalTok{ sir\_model}

\NormalTok{baseline }\OperatorTok{=}\NormalTok{ sir\_model(population}\OperatorTok{=}\DecValTok{10\_000}\NormalTok{, initial\_infected}\OperatorTok{=}\DecValTok{10}\NormalTok{, beta\_per\_day}\OperatorTok{=}\FloatTok{0.35}\NormalTok{, gamma\_per\_day}\OperatorTok{=}\FloatTok{0.1}\NormalTok{, days}\OperatorTok{=}\DecValTok{30}\NormalTok{)}
\NormalTok{distancing }\OperatorTok{=}\NormalTok{ sir\_model(population}\OperatorTok{=}\DecValTok{10\_000}\NormalTok{, initial\_infected}\OperatorTok{=}\DecValTok{10}\NormalTok{, beta\_per\_day}\OperatorTok{=}\FloatTok{0.18}\NormalTok{, gamma\_per\_day}\OperatorTok{=}\FloatTok{0.1}\NormalTok{, days}\OperatorTok{=}\DecValTok{30}\NormalTok{)}

\BuiltInTok{print}\NormalTok{(}\StringTok{"baseline R0:"}\NormalTok{, baseline.r0, }\StringTok{"peak infected:"}\NormalTok{, }\BuiltInTok{round}\NormalTok{(}\BuiltInTok{max}\NormalTok{(baseline.infected)))}
\BuiltInTok{print}\NormalTok{(}\StringTok{"distancing R0:"}\NormalTok{, distancing.r0, }\StringTok{"peak infected:"}\NormalTok{, }\BuiltInTok{round}\NormalTok{(}\BuiltInTok{max}\NormalTok{(distancing.infected)))}
\end{Highlighting}
\end{Shaded}

\subsubsection*{Part 3: Herd Immunity
Calculation}\label{part-3-herd-immunity-calculation}
\addcontentsline{toc}{subsubsection}{Part 3: Herd Immunity Calculation}

\begin{enumerate}
\def\labelenumi{\arabic{enumi}.}
\setcounter{enumi}{4}
\item
  \textbf{Design an Alternative Experiment}: Instead of just calculating
  herd immunity thresholds, design an experiment to test how different
  vaccination rates affect disease spread in a simulated population.
  What hypothesis would you test? How would you model transmission?
\item
  \textbf{Hypothesis}: Predict the relationship between vaccination
  coverage and peak epidemic size. Would you expect a linear
  relationship, or something else? Why?
\end{enumerate}

\subsubsection*{Part 4: Current Intervention Decision
Matrices}\label{part-4-current-intervention-decision-matrices}
\addcontentsline{toc}{subsubsection}{Part 4: Current Intervention
Decision Matrices}

\begin{enumerate}
\def\labelenumi{\arabic{enumi}.}
\setcounter{enumi}{6}
\tightlist
\item
  Use the TB decision cards to classify each case as standard
  drug-susceptible therapy, BPaLM/BPaL eligibility review, or
  individualized specialist regimen. Record the evidence that makes you
  confident or uncertain.
\item
  Use the HIV PrEP comparison card to match patient scenarios to a
  plausible PrEP option. Separate biological efficacy, adherence
  feasibility, contraindications, testing needs, and access constraints.
\item
  Use the malaria vector-control card to design a layered prevention
  package for a household with night-time exposure, daytime indoor
  exposure, and insecticide-resistance concerns. Explain why no single
  tool replaces surveillance.
\end{enumerate}

\subsection*{Data Recording}\label{data-recording-23}
\addcontentsline{toc}{subsection}{Data Recording}

Epidemic curve: peak at hour \textbf{\emph{; attack rate: }}\%;
estimated incubation: \_\_\_ h; R\texorpdfstring{\textsubscript{0}}{0} \texorpdfstring{\ensuremath{\approx}}{~} \_\_\_

{\def\LTcaptype{none} % do not increment counter
\begin{longtable}[]{@{}lll@{}}
\toprule\noalign{}
Antibiotic & Zone (mm) & Interpretation (S/I/R) \\
\midrule\noalign{}
\endhead
\bottomrule\noalign{}
\endlastfoot
Ampicillin & & \\
Kanamycin & & \\
Tetracycline & & \\
Chloramphenicol & & \\
\end{longtable}
}

Replicate zone agreement within \texorpdfstring{\ensuremath{\pm}}{+/-}2 mm? \_\_\_\_\_\_\_; AMR mechanism
card selected: \_\_\_\_\_\_\_ TB regimen path selected: \_\_\_\_\_\_\_;
PrEP option justified: \_\_\_\_\_\_\_; malaria control package:
\_\_\_\_\_\_\_

{\def\LTcaptype{none} % do not increment counter
\begin{longtable}[]{@{}lll@{}}
\toprule\noalign{}
Disease & R\texorpdfstring{\textsubscript{0}}{0} & Herd immunity threshold (p\_c) \\
\midrule\noalign{}
\endhead
\bottomrule\noalign{}
\endlastfoot
Measles & 15 & \\
Influenza & 2.5 & \\
COVID-19 & 2.9 & \\
Mumps & 5 & \\
Polio & 6 & \\
\end{longtable}
}

\subsection*{Evidence and Reproducibility
Checklist}\label{evidence-and-reproducibility-checklist-26}
\addcontentsline{toc}{subsection}{Evidence and Reproducibility
Checklist}

\begin{itemize}
\tightlist
\item
  \textbf{Primary evidence goal:} Model transmission and compare
  prevention strategies.
\item
  \textbf{Data skill to practice:} Interpret outbreak curves and
  intervention effects.
\item
  \textbf{BioSkills emphasis:} Science and society, Process of science,
  Modeling and simulation.
\item
  \textbf{Control logic:} identify at least one positive control, one
  negative control, or one baseline comparison before interpreting
  results.
\item
  \textbf{Measurement discipline:} record units, uncertainty, sample
  size, and any discarded observation with a reason.
\item
  \textbf{Mechanistic link:} connect one result directly to the parent
  chapter's big idea before writing the conclusion.
\item
  \textbf{Reproducibility check:} state one procedural detail that
  another group would need in order to reproduce the result.
\end{itemize}

\subsection*{Paper-Based Evidence
Upgrade}\label{paper-based-evidence-upgrade-26}
\addcontentsline{toc}{subsection}{Paper-Based Evidence Upgrade}

Before answering the analysis questions, annotate the paper dataset for
\textbf{Infectious Disease and Immunity} with a reproducibility pass:

{\def\LTcaptype{none} % do not increment counter
\begin{longtable}[]{@{}
  >{\raggedright\arraybackslash}p{(\linewidth - 2\tabcolsep) * \real{0.5000}}
  >{\raggedright\arraybackslash}p{(\linewidth - 2\tabcolsep) * \real{0.5000}}@{}}
\toprule\noalign{}
\begin{minipage}[b]{\linewidth}\raggedright
Evidence check
\end{minipage} & \begin{minipage}[b]{\linewidth}\raggedright
Student action
\end{minipage} \\
\midrule\noalign{}
\endhead
\bottomrule\noalign{}
\endlastfoot
Control logic & Mark the comparison that functions as the baseline,
negative control, or reference case. \\
Uncertainty & Circle the row, card, diagram feature, or model assumption
most likely to change the conclusion. \\
Model comparison & State whether a simpler rule, null model, or
alternative mechanism could explain the same pattern. \\
Decision threshold & Write the minimum evidence that would make you
revise the interpretation. \\
Reproducibility & Record the exact scoring rule another group would need
to reproduce your classification. \\
\end{longtable}
}

Focus note: in this lab, trace every outbreak conclusion back to its
evidence chain (case definition, transmission parameter, and the
diagnostic or surveillance signal it rests on), and name the alternative
epidemic explanation you ruled out before reporting an \(R_0\) or attack
rate. Keep required work paper-based; any material-handling or equipment
version belongs only in an optional extension.

\subsection*{Analysis Questions}\label{analysis-questions-26}
\addcontentsline{toc}{subsection}{Analysis Questions}

\begin{enumerate}
\def\labelenumi{\arabic{enumi}.}
\item
  If the epidemic curve shows a single sharp peak within a window
  shorter than the incubation period, what type of source does this
  suggest? How would a propagated (person-to-person) epidemic curve look
  different?
\item
  \emph{E. coli} showed resistance to ampicillin but susceptibility to
  kanamycin. Explain the likely molecular mechanism of ampicillin
  resistance (naming the specific enzyme) and why one \texorpdfstring{\ensuremath{\beta}}{beta}-lactamase can
  create cross-resistance to several related \texorpdfstring{\ensuremath{\beta}}{beta}-lactams while still
  leaving some \texorpdfstring{\ensuremath{\beta}}{beta}-lactam/\texorpdfstring{\ensuremath{\beta}}{beta}-lactamase-inhibitor combinations, carbapenems,
  or non-\texorpdfstring{\ensuremath{\beta}}{beta}-lactam drugs active depending on the enzyme and permeability
  context.
\item
  Measles has an R\texorpdfstring{\textsubscript{0}}{0} of about 15, requiring about 93\%
  vaccination for herd immunity. Your calculation shows that even high
  measles vaccination coverage is needed. What happens to herd immunity
  when 5\% of a vaccinated population declines vaccination? Relate this
  to observed measles outbreaks in Europe (2019).
\item
  An immunocompromised patient (CD4 T cell count \textless{} 200/\texorpdfstring{\ensuremath{\mu}}{mu}L,
  HIV+ AIDS) contracts measles despite prior vaccination. Explain: (a)
  why the vaccine may be insufficient; (b) which branch of adaptive
  immunity (humoral or cellular) is most deficient; (c) what specific
  viral clearance mechanism is impaired.
\item
  The O'Neill Review's 10-million-deaths-per-year by 2050 AMR estimate
  is a historical warning scenario, not a current burden count
  \citep{oneill2016amr}. Describe three specific mechanisms by which
  bacteria acquire resistance genes, and evaluate two distinct policy
  interventions that would slow the spread of AMR at a population level.
\item
  A malaria programme already distributes insecticide-treated nets, but
  household surveys show substantial indoor activity before bedtime. Use
  the vector-control matrix to decide whether spatial emanators could be
  an added tool, and state two evidence gaps that would still require
  local monitoring.
\end{enumerate}

\subsection*{Extension Analysis
Questions}\label{extension-analysis-questions-14}
\addcontentsline{toc}{subsection}{Extension Analysis Questions}

\begin{enumerate}
\def\labelenumi{\arabic{enumi}.}
\setcounter{enumi}{6}
\item
  \textbf{Statistical Analysis}: If the printed dataset includes three
  replicate disc-diffusion images, how would you analyze the variation
  in zone diameters? What statistical test would you use to determine if
  differences between antibiotics are significant?
\item
  \textbf{Experimental Error Analysis}: What are three potential sources
  of error in epidemic curve analysis? How could you modify the
  procedure to reduce each source of error?
\item
  \textbf{Experimental Design}: If you wanted to test the effectiveness
  of different public health interventions (e.g., vaccination, social
  distancing) in controlling an outbreak, how would you design the
  experiment? What data would you collect? How would you analyze it?
\end{enumerate}

\subsection*{Group Project Extension
(Multi-Session)}\label{group-project-extension-multi-session-11}
\addcontentsline{toc}{subsection}{Group Project Extension
(Multi-Session)}

\textbf{Design an Epidemiological Study of Vaccine Hesitancy}

This project will span 2-3 lab sessions:

\textbf{Session 1}: Research Question and Hypothesis - How does vaccine
hesitancy affect disease outbreak size? - Formulate a hypothesis about
the relationship between vaccine coverage and outbreak probability -
Design a study using simulation or real outbreak data

\textbf{Session 2}: Data Collection and Analysis - Collect data on
vaccine coverage and outbreak sizes from public health records - Analyze
the correlation between coverage and outbreak size - Model the impact of
increasing vaccine hesitancy on outbreak risk

\textbf{Session 3}: Data Interpretation and Presentation - Calculate the
critical vaccination threshold for different diseases - Discuss the
public health implications of declining vaccine coverage - Create a
scientific poster or presentation

\subsection*{Real-World Problem Solving: Epidemic
Ethics}\label{real-world-problem-solving-epidemic-ethics}
\addcontentsline{toc}{subsection}{Real-World Problem Solving: Epidemic
Ethics}

\textbf{Case Study: Quarantine and Civil Liberties}

\begin{enumerate}
\def\labelenumi{\arabic{enumi}.}
\item
  \textbf{Research Task}: Investigate the legal and ethical basis for
  quarantine during infectious disease outbreaks. What are the key
  principles that balance public health needs with individual rights?
\item
  \textbf{Ethical Consideration}: Should governments have the power to
  enforce quarantine on individuals who refuse vaccination? What are the
  potential benefits and risks of such policies?
\item
  \textbf{Policy Proposal}: Design a policy for implementing quarantine
  measures during an outbreak that respects civil liberties while
  protecting public health. Consider: due process, compensation for lost
  income, and transparent decision-making.
\end{enumerate}

\subsection*{Safety and Ethics Notes}\label{safety-and-ethics-notes-26}
\addcontentsline{toc}{subsection}{Safety and Ethics Notes}

Default lab work uses printed outbreak datasets, susceptibility images,
and AMR mechanism cards primarily. No bacterial cultures, agar plates,
or antibiotic discs are required. Optional wet susceptibility
demonstrations must use approved BSL-1 organisms under instructor
supervision and local disposal rules. Discuss epidemic data without
stigmatising affected groups. When discussing antibiotics, emphasize
following current prescriber instructions, avoiding unnecessary
antibiotic starts, using diagnostics when available, and protecting
access to effective treatment.

\subsection*{Debrief and Reflection}\label{debrief-and-reflection-23}
\addcontentsline{toc}{subsection}{Debrief and Reflection}

After you finish the practical work, spend 5--10 minutes in your small
group comparing results and discussing the following prompts. Each
member should contribute at least one observation before moving to the
next prompt:

\begin{enumerate}
\def\labelenumi{\arabic{enumi}.}
\item
  \textbf{What did your measurements show} --- compare the group's
  results to the textbook's predictions. Where they diverge, suggest at
  least one mechanistic explanation before concluding ``experimental
  error.''
\item
  \textbf{What would change the outcome} --- propose one modification to
  the procedure that would sharpen the measurement or extend the result
  to a new biological context, and predict what you would observe.
\item
  \textbf{One-sentence headline} --- each student composes a single
  sentence summarising the lab's take-home message, suitable for a
  tweet. Compare sentences across groups; good headlines are short,
  quantitative, and mechanistic.
\item
  \textbf{Connection back to the textbook} --- identify one section of
  \cref{sec:unit_VII_infectious_disease} that your data either confirmed
  or complicated. Cite the specific passage.
\item
  \textbf{Experimental Design Reflection}: What was the strongest aspect
  of your experimental design? What would you change if you could repeat
  the experiment?
\end{enumerate}

\subsection*{Further Reading (Lab)}\label{further-reading-lab-23}
\addcontentsline{toc}{subsection}{Further Reading (Lab)}

\begin{itemize}
\tightlist
\item
  Revisit the parent chapter \cref{sec:unit_VII_infectious_disease} for
  the theoretical foundations on which this lab is built.
\item
  Look up any \textbf{bolded glossary term} introduced in the textbook
  chapter (each has a \texttt{\#gl:term-slug} link in the text) --- its
  master definition is in \breaktt{manuscript/glossary.md}.
\item
  Explore the Python code in
  \breaktt{src/biology/microbiology/microbiology.py} for epidemiological
  calculations.
\item
  Use the self-contained Part 2 snippet as the computational template;
  it runs against tested project modules without external notebooks or
  CSV files.
\end{itemize}

\newpage

\section*{Lab 25 --- Plant Structure, Water Relations, and
Transport}\label{lab-25-plant-structure-water-relations-and-transport}
\addcontentsline{toc}{section}{Lab 25 --- Plant Structure, Water
Relations, and Transport}

\label{sec:lab_unit_VIII_plant_structure_and_water}

\emph{This activity accompanies
\cref{sec:unit_VIII_plant_structure_and_water} of the textbook ---
review that chapter before attempting the exercises below.}

\subsection*{Learning Objectives}\label{learning-objectives-65}
\addcontentsline{toc}{subsection}{Learning Objectives}

\begin{itemize}
\tightlist
\item
  Design a paper-based transpiration investigation using printed
  plant-water datasets
\item
  Formulate testable hypotheses about how environmental factors affect
  transpiration rate
\item
  Identify independent, dependent, and controlled variables in plant
  physiology experiments
\item
  Analyze experimental data using water potential calculations and
  statistical methods
\item
  Write a brief scientific report with hypothesis, methods, results, and
  conclusion
\end{itemize}

\subsubsection*{Alignment and Rubric
Map}\label{alignment-and-rubric-map-27}
\addcontentsline{toc}{subsubsection}{Alignment and Rubric Map}

\begin{itemize}
\tightlist
\item
  \textbf{Outcome 1 (LO1):} Interpret the supplied evidence or model
  output for Plant Structure, Water Relations, and Transport.
\item
  \textbf{Outcome 2 (LO2):} Identify controls and comparison groups that
  make the claim testable.
\item
  \textbf{Outcome 3 (LO3):} Quantify uncertainty, boundary conditions,
  or alternative explanations before concluding.
\item
  \textbf{Outcome 4 (LO4):} Transfer the mechanism to a new biological
  case or public-facing decision.
\item
  \textbf{Chapter LO coverage:} LO1, LO2, LO3, LO4
\item
  \textbf{Rubric dimensions:} evidence; controls; uncertainty;
  mechanism; transfer. \#\# Pre-Lab Inquiry Questions \{.unnumbered\}
\end{itemize}

\emph{Complete these questions before coming to lab. Use your textbook
and additional research to inform your hypotheses.}

\begin{enumerate}
\def\labelenumi{\arabic{enumi}.}
\item
  \textbf{Hypothesis Formation}: Based on your understanding of
  transpiration, predict which environmental condition (bright light,
  wind, high humidity, dim light) will result in the highest
  transpiration rate. Write a clear, testable hypothesis.
\item
  \textbf{Experimental Design}: What is the independent variable in a
  transpiration experiment? What is the dependent variable? List at
  least 5 variables that should be controlled to ensure valid results.
\item
  \textbf{Scientific Context}: The cohesion-tension theory explains how
  water moves upward in plants against gravity. Why is this mechanism
  dependent on transpiration? What would happen to water transport if
  transpiration were completely stopped?
\item
  \textbf{Quantitative Reasoning}: If a plant loses 5 grams of water in
  30 minutes, what is its transpiration rate in g/hr? If the leaf
  surface area is 100 cm\texorpdfstring{\textsuperscript{2}}{2}, what is the transpiration rate per unit area
  (g/m\texorpdfstring{\textsuperscript{2}}{2}/hr)?
\item
  \textbf{Real-World Application}: Drought stress affects crop yields.
  How does understanding transpiration help in developing
  drought-resistant crops? What traits would you select for?
\end{enumerate}

\subsection*{Background}\label{background-15}
\addcontentsline{toc}{subsection}{Background}

Vascular plants have three tissue systems: dermal (epidermis, cuticle,
guard cells), ground (parenchyma, collenchyma, sclerenchyma), and
vascular (xylem for water/mineral transport upward; phloem for
sugar/assimilate transport, bidirectional). Transpiration ---
evaporation of water from stomata --- drives the cohesion-tension
mechanism that pulls water from root to leaf. Environmental factors
(light, humidity, temperature, wind) modulate transpiration rate by
affecting stomatal aperture.

\subsection*{Paper-Based Materials}\label{paper-based-materials-27}
\addcontentsline{toc}{subsection}{Paper-Based Materials}

{\def\LTcaptype{none} % do not increment counter
\begin{longtable}[]{@{}
  >{\raggedright\arraybackslash}p{(\linewidth - 2\tabcolsep) * \real{0.3333}}
  >{\raggedright\arraybackslash}p{(\linewidth - 2\tabcolsep) * \real{0.6667}}@{}}
\toprule\noalign{}
\begin{minipage}[b]{\linewidth}\raggedright
Item
\end{minipage} & \begin{minipage}[b]{\linewidth}\raggedright
Quantity
\end{minipage} \\
\midrule\noalign{}
\endhead
\bottomrule\noalign{}
\endlastfoot
Printed dicot stem and root cross-section diagrams with scale bars & 1
packet \\
Transpiration dataset for four environmental treatments, with three
replicates each & 1 \\
Leaf-area and stomatal-density data cards for broadleaf, needleleaf, and
succulent plants & 1 set \\
Water-potential reference sheet and worked example & 1 \\
Crop water-use case packet with drought, yield, and irrigation data &
1 \\
Graph paper or spreadsheet template & 1 \\
Calculator or optional Python REPL with this project installed & 1 \\
Optional extension card describing a potometer or celery-dye
demonstration & 1 \\
\end{longtable}
}

\subsection*{Paper-Based
Investigation}\label{paper-based-investigation-27}
\addcontentsline{toc}{subsection}{Paper-Based Investigation}

\subsubsection*{Part 1: Anatomical Diagram and Tissue
Identification}\label{part-1-anatomical-diagram-and-tissue-identification}
\addcontentsline{toc}{subsubsection}{Part 1: Anatomical Diagram and
Tissue Identification}

\begin{enumerate}
\def\labelenumi{\arabic{enumi}.}
\item
  \textbf{Formulate Your Hypothesis}: Based on your pre-lab research,
  write a hypothesis about how stem and root anatomy differ in their
  vascular tissue arrangement.
\item
  \textbf{Identify Variables}:

  \begin{itemize}
  \tightlist
  \item
    Independent variable: Plant organ (stem vs root)
  \item
    Dependent variable: Tissue arrangement and identification
  \item
    Controlled variables: image source, scale-bar calibration, diagram
    resolution, tissue orientation, scoring key, and blind-card order.
  \end{itemize}
\item
  \textbf{Set Up Controls}: Why is it important for the diagram packet
  to include a labelled reference image and an unlabelled image from the
  same species? What other controls could you include to ensure accurate
  tissue identification?
\item
  \textbf{Data Collection}:

  \begin{itemize}
  \tightlist
  \item
    Use the printed stem cross-section diagram to identify and label:
    epidermis, cortex (parenchyma), vascular bundles (xylem = larger,
    thick-walled; phloem = smaller, thin-walled), pith.
  \item
    Use the printed root cross-section diagram to identify: epidermis
    (root hairs), cortex, endodermis (Casparian strip), pericycle,
    vascular cylinder (xylem star pattern; phloem between xylem arms).
  \item
    Record at least three independent tissue-identification decisions
    per organ, citing the visual feature used for each decision.
  \end{itemize}
\end{enumerate}

\subsubsection*{Part 2: Computational Biology Exercise - Transpiration
Analysis with
Python}\label{part-2-computational-biology-exercise---transpiration-analysis-with-python}
\addcontentsline{toc}{subsubsection}{Part 2: Computational Biology
Exercise - Transpiration Analysis with Python}

\emph{Optional computational check: run this self-contained Python
snippet from the project root. It uses tested \texttt{src/biology}
modules and requires no external notebook or CSV file.}

Check water-potential and transpiration calculations from your data
table.

\begin{Shaded}
\begin{Highlighting}[]
\ImportTok{from}\NormalTok{ biology.botany }\ImportTok{import}\NormalTok{ transpiration\_flux, water\_potential}

\NormalTok{leaf }\OperatorTok{=}\NormalTok{ water\_potential(solute\_concentration\_M}\OperatorTok{=}\FloatTok{0.3}\NormalTok{, turgor\_pressure\_MPa}\OperatorTok{=}\FloatTok{0.4}\NormalTok{)}
\NormalTok{flux }\OperatorTok{=}\NormalTok{ transpiration\_flux(}
\NormalTok{    stomatal\_conductance\_mol\_m2\_s}\OperatorTok{=}\FloatTok{0.2}\NormalTok{,}
\NormalTok{    internal\_vapor\_conc\_mol\_m3}\OperatorTok{=}\FloatTok{0.5}\NormalTok{,}
\NormalTok{    external\_vapor\_conc\_mol\_m3}\OperatorTok{=}\FloatTok{0.3}\NormalTok{,}
\NormalTok{)}

\BuiltInTok{print}\NormalTok{(}\StringTok{"water potential:"}\NormalTok{, }\BuiltInTok{round}\NormalTok{(leaf.water\_potential\_MPa, }\DecValTok{3}\NormalTok{), }\StringTok{"MPa"}\NormalTok{)}
\BuiltInTok{print}\NormalTok{(}\StringTok{"transpiration flux:"}\NormalTok{, }\BuiltInTok{round}\NormalTok{(flux.flux\_mmol\_m2\_s, }\DecValTok{2}\NormalTok{), }\StringTok{"mmol m\^{}{-}2 s\^{}{-}1"}\NormalTok{)}
\end{Highlighting}
\end{Shaded}

\subsubsection*{Part 3: Transpiration Dataset and Model
Check}\label{part-3-transpiration-dataset-and-model-check}
\addcontentsline{toc}{subsubsection}{Part 3: Transpiration Dataset and
Model Check}

\begin{enumerate}
\def\labelenumi{\arabic{enumi}.}
\setcounter{enumi}{4}
\item
  \textbf{Design an Alternative Dataset}: Instead of just comparing
  environmental treatments, design a dataset to test how plant species
  with different leaf types (e.g., broadleaf vs needleleaf) differ in
  their transpiration rates. What hypothesis would you test? Which
  variables would you hold constant?
\item
  \textbf{Hypothesis}: Predict the relationship between leaf surface
  area and transpiration rate. Would you expect a linear relationship,
  or would stomatal density also play a role? Why?
\item
  \textbf{Reproducibility Check}: For each treatment, calculate the
  mean, range, and coefficient of variation across the three printed
  replicates. Flag any replicate that would change the treatment ranking
  if removed, and decide whether it should be retained.
\end{enumerate}

\subsection*{Data Recording}\label{data-recording-24}
\addcontentsline{toc}{subsection}{Data Recording}

{\def\LTcaptype{none} % do not increment counter
\begin{longtable}[]{@{}llll@{}}
\toprule\noalign{}
Structure & Observed in stem? & Observed in root? & Function \\
\midrule\noalign{}
\endhead
\bottomrule\noalign{}
\endlastfoot
Epidermis & & & \\
Xylem & & & \\
Phloem & & & \\
Endodermis & & & \\
Pith & & & \\
\end{longtable}
}

{\def\LTcaptype{none} % do not increment counter
\begin{longtable}[]{@{}
  >{\raggedright\arraybackslash}p{(\linewidth - 8\tabcolsep) * \real{0.1098}}
  >{\raggedright\arraybackslash}p{(\linewidth - 8\tabcolsep) * \real{0.1707}}
  >{\raggedright\arraybackslash}p{(\linewidth - 8\tabcolsep) * \real{0.1220}}
  >{\raggedright\arraybackslash}p{(\linewidth - 8\tabcolsep) * \real{0.3049}}
  >{\raggedright\arraybackslash}p{(\linewidth - 8\tabcolsep) * \real{0.2927}}@{}}
\toprule\noalign{}
\begin{minipage}[b]{\linewidth}\raggedright
Condition
\end{minipage} & \begin{minipage}[b]{\linewidth}\raggedright
Mean water loss (g)
\end{minipage} & \begin{minipage}[b]{\linewidth}\raggedright
Time (min)
\end{minipage} & \begin{minipage}[b]{\linewidth}\raggedright
Transpiration rate (g/min)
\end{minipage} & \begin{minipage}[b]{\linewidth}\raggedright
Transpiration (g/m\texorpdfstring{\textsuperscript{2}}{2}/min)
\end{minipage} \\
\midrule\noalign{}
\endhead
\bottomrule\noalign{}
\endlastfoot
Dim light, still & & & & \\
Bright light, still & & & & \\
Dim light, fan & & & & \\
Bag (humid, still) & & & & \\
\end{longtable}
}

Diagram evidence for xylem pathway: feature 1: \textbf{\emph{; feature
2: }}; feature 3: \_\_\_

\subsection*{Evidence and Reproducibility
Checklist}\label{evidence-and-reproducibility-checklist-27}
\addcontentsline{toc}{subsection}{Evidence and Reproducibility
Checklist}

\begin{itemize}
\tightlist
\item
  \textbf{Primary evidence goal:} Measure or model water movement
  through plant tissues.
\item
  \textbf{Data skill to practice:} Interpret plant-water data from
  pressure, solute, and humidity measurements.
\item
  \textbf{BioSkills emphasis:} Interdisciplinary nature of science,
  Process of science.
\item
  \textbf{Control logic:} identify at least one positive control, one
  negative control, or one baseline comparison before interpreting
  results.
\item
  \textbf{Measurement discipline:} record units, uncertainty, sample
  size, and any discarded observation with a reason.
\item
  \textbf{Mechanistic link:} connect one result directly to the parent
  chapter's big idea before writing the conclusion.
\item
  \textbf{Reproducibility check:} state one procedural detail that
  another group would need in order to reproduce the result.
\end{itemize}

\subsection*{Paper-Based Evidence
Upgrade}\label{paper-based-evidence-upgrade-27}
\addcontentsline{toc}{subsection}{Paper-Based Evidence Upgrade}

Before answering the analysis questions, annotate the paper dataset for
\textbf{Plant Structure, Water Relations, and Transport} with a
reproducibility pass:

{\def\LTcaptype{none} % do not increment counter
\begin{longtable}[]{@{}
  >{\raggedright\arraybackslash}p{(\linewidth - 2\tabcolsep) * \real{0.5000}}
  >{\raggedright\arraybackslash}p{(\linewidth - 2\tabcolsep) * \real{0.5000}}@{}}
\toprule\noalign{}
\begin{minipage}[b]{\linewidth}\raggedright
Evidence check
\end{minipage} & \begin{minipage}[b]{\linewidth}\raggedright
Student action
\end{minipage} \\
\midrule\noalign{}
\endhead
\bottomrule\noalign{}
\endlastfoot
Control logic & Mark the comparison that functions as the baseline,
negative control, or reference case. \\
Uncertainty & Circle the row, card, diagram feature, or model assumption
most likely to change the conclusion. \\
Model comparison & State whether a simpler rule, null model, or
alternative mechanism could explain the same pattern. \\
Decision threshold & Write the minimum evidence that would make you
revise the interpretation. \\
Reproducibility & Record the exact scoring rule another group would need
to reproduce your classification. \\
\end{longtable}
}

Focus note: in this water-relations lab, anchor every claim to a
water-potential value or gradient and the hydraulic pathway it acts on
(soil-root-xylem-leaf-air), and distinguish stomatal regulation, tissue
anatomy, and stress context as separate, testable controls on transport.
Keep required work paper-based; any material-handling or equipment
version belongs only in an optional extension.

\subsection*{Analysis Questions}\label{analysis-questions-27}
\addcontentsline{toc}{subsection}{Analysis Questions}

\begin{enumerate}
\def\labelenumi{\arabic{enumi}.}
\item
  The diagram packet traces water movement through xylem rather than
  phloem. Why does bulk water transport follow xylem pathways, and what
  would you predict for sugar movement in phloem under the pressure-flow
  hypothesis?
\item
  Which condition produced the highest transpiration rate --- bright
  light or fan? Explain why wind increases transpiration using the
  concept of the diffusion boundary layer.
\item
  Stomata open when guard cells swell (K\texorpdfstring{\textsuperscript{+}}{+} influx followed by water via
  osmosis). Explain how light triggers stomatal opening via the proton
  pump (H\texorpdfstring{\textsuperscript{+}}{+}-ATPase) and K\texorpdfstring{\textsuperscript{+}}{+} channels, and why abscisic acid (ABA) causes
  stomatal closure during drought.
\item
  The Casparian strip forces water from the cortex to enter the
  endodermis symplastically (through cells), not apoplastically (through
  cell walls). Why is this selectivity important for plant mineral
  uptake?
\item
  A plant biologist clips the tip of a phloem sieve tube using a laser
  scalpel. Predict what happens to sugar transport above and below the
  cut within 30 seconds, using your knowledge of the pressure-flow
  hypothesis for phloem loading and unloading.
\end{enumerate}

\subsection*{Extension Analysis
Questions}\label{extension-analysis-questions-15}
\addcontentsline{toc}{subsection}{Extension Analysis Questions}

\begin{enumerate}
\def\labelenumi{\arabic{enumi}.}
\setcounter{enumi}{5}
\item
  \textbf{Statistical Analysis}: Using the three printed replicates per
  treatment, how would you analyze variation in transpiration rates?
  What statistical test would you use to determine if differences
  between conditions are significant?
\item
  \textbf{Experimental Error Analysis}: What are three potential sources
  of error in transpiration measurements? How could you modify the
  procedure to reduce each source of error?
\item
  \textbf{Experimental Design}: If you wanted to test the effect of a
  specific hormone (e.g., ABA) on stomatal closure using a paper
  dataset, what concentrations would you include, what response variable
  would you measure, and what baseline comparison would make the result
  interpretable?
\end{enumerate}

\subsection*{Group Project Extension
(Multi-Session)}\label{group-project-extension-multi-session-12}
\addcontentsline{toc}{subsection}{Group Project Extension
(Multi-Session)}

\textbf{Design a Study of Plant Water Use Efficiency}

This project will span 2-3 lab sessions:

\textbf{Session 1}: Research Question and Hypothesis - How does water
use efficiency (WUE) vary among plant species from different
environments? - Formulate a hypothesis about the relationship between
WUE and drought tolerance - Design an experiment measuring transpiration
and photosynthesis rates

\textbf{Session 2}: Data Extraction and Analysis - Extract transpiration
rates from printed potometer or gravimetric datasets - Compare
photosynthesis rates from CO\texorpdfstring{\textsubscript{2}}{2} exchange or O\texorpdfstring{\textsubscript{2}}{2} evolution data tables -
Calculate intrinsic water use efficiency (photosynthesis/transpiration)

\textbf{Session 3}: Data Interpretation and Presentation - Compare WUE
across species - Correlate WUE with environmental adaptations (e.g.,
leaf thickness, stomatal density) - Create a scientific poster or
presentation

\subsection*{Real-World Problem Solving: Plant Water Relations and
Agriculture}\label{real-world-problem-solving-plant-water-relations-and-agriculture}
\addcontentsline{toc}{subsection}{Real-World Problem Solving: Plant
Water Relations and Agriculture}

\textbf{Case Study: Irrigation Management}

\begin{enumerate}
\def\labelenumi{\arabic{enumi}.}
\item
  \textbf{Research Task}: Investigate the concept of crop water use
  efficiency (WUE). What is the typical range for major crops like
  wheat, rice, and corn? How does WUE vary with environmental
  conditions?
\item
  \textbf{Ethical Consideration}: Should farmers be encouraged or
  required to use water-saving irrigation techniques (e.g., drip
  irrigation) in water-scarce regions? What are the economic and social
  implications of such policies?
\item
  \textbf{Policy Proposal}: Design a policy to promote water-efficient
  agriculture in drought-prone areas. Consider: incentives for adopting
  efficient irrigation, water pricing, and monitoring systems.
\end{enumerate}

\subsection*{Optional Wet-Lab
Demonstration}\label{optional-wet-lab-demonstration}
\addcontentsline{toc}{subsection}{Optional Wet-Lab Demonstration}

An instructor may demonstrate celery dye movement or potometer data
collection if local safety rules and materials allow. This is not
required for the lab; the default investigation uses printed diagrams
and datasets.

\subsection*{Safety and Ethics Notes}\label{safety-and-ethics-notes-27}
\addcontentsline{toc}{subsection}{Safety and Ethics Notes}

No wet materials are required. If an optional demonstration is used,
handle blades, glassware, and food dye primarily under instructor
supervision. When discussing agricultural water use, consider the needs
of different stakeholders.

\subsection*{Debrief and Reflection}\label{debrief-and-reflection-24}
\addcontentsline{toc}{subsection}{Debrief and Reflection}

After you finish the practical work, spend 5--10 minutes in your small
group comparing results and discussing the following prompts. Each
member should contribute at least one observation before moving to the
next prompt:

\begin{enumerate}
\def\labelenumi{\arabic{enumi}.}
\item
  \textbf{What did your measurements show} --- compare the group's
  results to the textbook's predictions. Where they diverge, suggest at
  least one mechanistic explanation before concluding ``experimental
  error.''
\item
  \textbf{What would change the outcome} --- propose one modification to
  the procedure that would sharpen the measurement or extend the result
  to a new biological context, and predict what you would observe.
\item
  \textbf{One-sentence headline} --- each student composes a single
  sentence summarising the lab's take-home message, suitable for a
  tweet. Compare sentences across groups; good headlines are short,
  quantitative, and mechanistic.
\item
  \textbf{Connection back to the textbook} --- identify one section of
  \cref{sec:unit_VIII_plant_structure_and_water} that your data either
  confirmed or complicated. Cite the specific passage.
\item
  \textbf{Experimental Design Reflection}: What was the strongest aspect
  of your experimental design? What would you change if you could repeat
  the experiment?
\end{enumerate}

\subsection*{Further Reading (Lab)}\label{further-reading-lab-24}
\addcontentsline{toc}{subsection}{Further Reading (Lab)}

\begin{itemize}
\tightlist
\item
  Revisit the parent chapter
  \cref{sec:unit_VIII_plant_structure_and_water} for the theoretical
  foundations on which this lab is built.
\item
  Look up any \textbf{bolded glossary term} introduced in the textbook
  chapter (each has a \texttt{\#gl:term-slug} link in the text) --- its
  master definition is in \breaktt{manuscript/glossary.md}.
\item
  Explore the Python code in \breaktt{src/biology/botany/botany.py} for
  plant water relations calculations.
\item
  Use the self-contained Part 2 snippet as the computational template;
  it runs against tested project modules without external notebooks or
  CSV files.
\end{itemize}

\newpage

\section*{Lab 26 --- Plant Reproduction and
Development}\label{lab-26-plant-reproduction-and-development}
\addcontentsline{toc}{section}{Lab 26 --- Plant Reproduction and
Development}

\label{sec:lab_unit_VIII_plant_reproduction}

\emph{This activity accompanies \cref{sec:unit_VIII_plant_reproduction}
of the textbook --- review that chapter before attempting the exercises
below.}

\subsection*{Learning Objectives}\label{learning-objectives-66}
\addcontentsline{toc}{subsection}{Learning Objectives}

\begin{itemize}
\tightlist
\item
  Design a paper-based pollination investigation using anatomical
  diagrams and pollen-tube datasets
\item
  Formulate testable hypotheses about how environmental factors affect
  pollen germination data
\item
  Identify independent, dependent, and controlled variables in plant
  reproduction experiments
\item
  Analyze experimental data using growth rate calculations and
  statistical methods
\item
  Write a brief scientific report with hypothesis, methods, results, and
  conclusion
\end{itemize}

\subsubsection*{Alignment and Rubric
Map}\label{alignment-and-rubric-map-28}
\addcontentsline{toc}{subsubsection}{Alignment and Rubric Map}

\begin{itemize}
\tightlist
\item
  \textbf{Outcome 1 (LO1):} Interpret the supplied evidence or model
  output for Plant Reproduction and Development.
\item
  \textbf{Outcome 2 (LO2):} Identify controls and comparison groups that
  make the claim testable.
\item
  \textbf{Outcome 3 (LO3):} Quantify uncertainty, boundary conditions,
  or alternative explanations before concluding.
\item
  \textbf{Outcome 4 (LO4):} Transfer the mechanism to a new biological
  case or public-facing decision.
\item
  \textbf{Chapter LO coverage:} LO1, LO2, LO3, LO4
\item
  \textbf{Rubric dimensions:} evidence; controls; uncertainty;
  mechanism; transfer. \#\# Pre-Lab Inquiry Questions \{.unnumbered\}
\end{itemize}

\emph{Complete these questions before coming to lab. Use your textbook
and additional research to inform your hypotheses.}

\begin{enumerate}
\def\labelenumi{\arabic{enumi}.}
\item
  \textbf{Hypothesis Formation}: Based on your understanding of pollen
  germination, predict how temperature might affect pollen tube growth
  rate. Write a clear, testable hypothesis.
\item
  \textbf{Experimental Design}: What is the independent variable in a
  pollen germination experiment? What is the dependent variable? List at
  least 5 variables that should be controlled to ensure valid results.
\item
  \textbf{Scientific Context}: Double fertilization is a key innovation
  in angiosperms. Why is double fertilization advantageous compared to
  single fertilization in gymnosperms? What evolutionary pressures might
  have selected for this trait?
\item
  \textbf{Quantitative Reasoning}: If pollen tubes grow at an average
  rate of 1 mm/hour, how long would it take for a pollen tube to reach
  an ovule 5 cm away from the stigma? Show your calculation.
\item
  \textbf{Real-World Application}: Climate change is affecting plant
  reproduction. How might rising temperatures impact pollen viability
  and fertilization success? What are the implications for crop
  production?
\end{enumerate}

\subsection*{Background}\label{background-16}
\addcontentsline{toc}{subsection}{Background}

Angiosperms (flowering plants) undergo \textbf{double fertilisation}:
one sperm nucleus (n) fuses with the egg (n) to form the zygote (2n); a
second sperm fuses with the two polar nuclei (2\texorpdfstring{\ensuremath{\times}}{x}n = 2n) to form the
primary endosperm nucleus (3n), which nourishes the developing embryo.
Pollen grains (male gametophytes, n) germinate on the stigma and grow a
pollen tube through the style to deliver sperm to the embryo sac (female
gametophyte) in the ovule.

\subsection*{Paper-Based Materials}\label{paper-based-materials-28}
\addcontentsline{toc}{subsection}{Paper-Based Materials}

{\def\LTcaptype{none} % do not increment counter
\begin{longtable}[]{@{}
  >{\raggedright\arraybackslash}p{(\linewidth - 2\tabcolsep) * \real{0.3333}}
  >{\raggedright\arraybackslash}p{(\linewidth - 2\tabcolsep) * \real{0.6667}}@{}}
\toprule\noalign{}
\begin{minipage}[b]{\linewidth}\raggedright
Item
\end{minipage} & \begin{minipage}[b]{\linewidth}\raggedright
Quantity
\end{minipage} \\
\midrule\noalign{}
\endhead
\bottomrule\noalign{}
\endlastfoot
Flower anatomy diagram packet with labelled and unlabelled angiosperm
structures & 1 \\
Pollination syndrome card set (wind, bee, moth, bird, and generalized
flowers) & 1 set \\
Pollen-tube growth dataset across temperature and sucrose treatments &
1 \\
Printed micrograph panels or schematic pollen-tube images with scale
bars & 1 set \\
Double-fertilisation ploidy worksheet & 1 \\
Climate-risk case packet for crop flowering, heat stress, and pollinator
service data & 1 \\
Graph paper or spreadsheet template & 1 \\
Calculator or optional Python REPL with this project installed & 1 \\
\end{longtable}
}

\subsection*{Paper-Based
Investigation}\label{paper-based-investigation-28}
\addcontentsline{toc}{subsection}{Paper-Based Investigation}

\subsubsection*{Part 1: Flower Anatomy and Pollination Strategy
Cards}\label{part-1-flower-anatomy-and-pollination-strategy-cards}
\addcontentsline{toc}{subsubsection}{Part 1: Flower Anatomy and
Pollination Strategy Cards}

\begin{enumerate}
\def\labelenumi{\arabic{enumi}.}
\item
  \textbf{Formulate Your Hypothesis}: Based on your pre-lab research,
  write a hypothesis about how flower morphology relates to pollination
  strategy.
\item
  \textbf{Identify Variables}:

  \begin{itemize}
  \tightlist
  \item
    Independent variable: Pollination syndrome (wind vs insect)
  \item
    Dependent variable: Flower morphological traits (petal colour,
    fragrance, pollen size, stigma surface)
  \item
    Controlled variables: Plant species, growing conditions, measurement
    methods, etc.
  \end{itemize}
\item
  \textbf{Set Up Controls}: Why is it important for each
  pollination-syndrome card to include both a labelled reference image
  and a blinded comparison image? What other controls could you include
  to ensure accurate comparisons?
\item
  \textbf{Data Collection}:

  \begin{itemize}
  \tightlist
  \item
    Identify external parts on the printed diagram: sepals, petals,
    stamens (anther + filament), pistil (stigma + style + ovary).
  \item
    Use the ovary cross-section diagram to locate ovules and trace the
    path a pollen tube would follow.
  \item
    Use micrograph panels or schematic cards to record pollen
    characteristics such as size, surface texture, and tube length.
  \item
    Compare wind-pollinated vs insect-pollinated flower cards. Fill in
    the comparison table and cite the trait evidence behind each
    classification.
  \end{itemize}
\end{enumerate}

\subsubsection*{Part 2: Computational Biology Exercise - Pollen Tube
Growth Analysis with
Python}\label{part-2-computational-biology-exercise---pollen-tube-growth-analysis-with-python}
\addcontentsline{toc}{subsubsection}{Part 2: Computational Biology
Exercise - Pollen Tube Growth Analysis with Python}

\emph{Optional computational check: run this self-contained Python
snippet from the project root. It uses tested \texttt{src/biology}
modules and requires no external notebook or CSV file.}

Use a growth curve as a quantitative analogue for pollen-tube extension.

\begin{Shaded}
\begin{Highlighting}[]
\ImportTok{from}\NormalTok{ biology.botany }\ImportTok{import}\NormalTok{ plant\_biomass\_growth}

\NormalTok{growth }\OperatorTok{=}\NormalTok{ plant\_biomass\_growth(}
\NormalTok{    initial\_biomass\_g}\OperatorTok{=}\FloatTok{0.2}\NormalTok{,}
\NormalTok{    relative\_growth\_rate}\OperatorTok{=}\FloatTok{0.35}\NormalTok{,}
\NormalTok{    carrying\_capacity\_g}\OperatorTok{=}\FloatTok{10.0}\NormalTok{,}
\NormalTok{    duration\_days}\OperatorTok{=}\DecValTok{10}\NormalTok{,}
\NormalTok{)}
\BuiltInTok{print}\NormalTok{(}\StringTok{"initial biomass:"}\NormalTok{, growth.biomass\_g[}\DecValTok{0}\NormalTok{])}
\BuiltInTok{print}\NormalTok{(}\StringTok{"final biomass:"}\NormalTok{, }\BuiltInTok{round}\NormalTok{(growth.biomass\_g[}\OperatorTok{{-}}\DecValTok{1}\NormalTok{], }\DecValTok{2}\NormalTok{))}
\BuiltInTok{print}\NormalTok{(}\StringTok{"days recorded:"}\NormalTok{, }\BuiltInTok{len}\NormalTok{(growth.times\_days))}
\end{Highlighting}
\end{Shaded}

\subsubsection*{Part 3: Pollen Germination
Dataset}\label{part-3-pollen-germination-dataset}
\addcontentsline{toc}{subsubsection}{Part 3: Pollen Germination Dataset}

\begin{enumerate}
\def\labelenumi{\arabic{enumi}.}
\setcounter{enumi}{4}
\item
  \textbf{Design an Alternative Dataset}: Instead of just measuring
  pollen tube growth at different temperatures, design a dataset to test
  how sugar concentration affects pollen tube growth. What hypothesis
  would you test? What sugar concentrations would you include, and what
  treatment would serve as the baseline?
\item
  \textbf{Hypothesis}: Predict the relationship between sucrose
  concentration and pollen tube growth rate. Would you expect a linear
  increase, a saturating curve, or something else? Why?
\item
  \textbf{Reproducibility Check}: For the printed pollen-tube dataset,
  calculate mean growth rate and standard deviation for each treatment.
  Decide whether the temperature effect is larger than within-treatment
  variation.
\end{enumerate}

\subsection*{Data Recording}\label{data-recording-25}
\addcontentsline{toc}{subsection}{Data Recording}

Flower anatomy diagram (label most structures):

\texttt{{[}Annotate\ the\ printed\ flower\ diagram\ here{]}}

{\def\LTcaptype{none} % do not increment counter
\begin{longtable}[]{@{}lll@{}}
\toprule\noalign{}
Structure & Present? & Description (colour, size, texture) \\
\midrule\noalign{}
\endhead
\bottomrule\noalign{}
\endlastfoot
Sepals & & \\
Petals & & \\
Anthers & & \\
Stigma & & \\
Style & & \\
Ovary & & \\
Ovules & & \\
\end{longtable}
}

{\def\LTcaptype{none} % do not increment counter
\begin{longtable}[]{@{}
  >{\raggedright\arraybackslash}p{(\linewidth - 4\tabcolsep) * \real{0.1515}}
  >{\raggedright\arraybackslash}p{(\linewidth - 4\tabcolsep) * \real{0.3485}}
  >{\raggedright\arraybackslash}p{(\linewidth - 4\tabcolsep) * \real{0.5000}}@{}}
\toprule\noalign{}
\begin{minipage}[b]{\linewidth}\raggedright
Time (min)
\end{minipage} & \begin{minipage}[b]{\linewidth}\raggedright
Pollen tube length (\texorpdfstring{\ensuremath{\mu}}{mu}m)
\end{minipage} & \begin{minipage}[b]{\linewidth}\raggedright
Pollen tube growth rate (\texorpdfstring{\ensuremath{\mu}}{mu}m/min)
\end{minipage} \\
\midrule\noalign{}
\endhead
\bottomrule\noalign{}
\endlastfoot
0 & 0 & --- \\
30 & & \\
60 & & \\
90 & & \\
120 & & \\
\end{longtable}
}

{\def\LTcaptype{none} % do not increment counter
\begin{longtable}[]{@{}lll@{}}
\toprule\noalign{}
Character & Wind-pollinated (grass) & Insect-pollinated (lily) \\
\midrule\noalign{}
\endhead
\bottomrule\noalign{}
\endlastfoot
Petal colour & & \\
Fragrance & & \\
Pollen size & & \\
Pollen surface & & \\
Stigma surface & & \\
\end{longtable}
}

\subsection*{Evidence and Reproducibility
Checklist}\label{evidence-and-reproducibility-checklist-28}
\addcontentsline{toc}{subsection}{Evidence and Reproducibility
Checklist}

\begin{itemize}
\tightlist
\item
  \textbf{Primary evidence goal:} Compare reproductive structures and
  developmental outcomes.
\item
  \textbf{Data skill to practice:} Track ploidy, tissue origin, and
  reproductive stage from diagrams or observations.
\item
  \textbf{BioSkills emphasis:} Interdisciplinary nature of science,
  Process of science.
\item
  \textbf{Control logic:} identify at least one positive control, one
  negative control, or one baseline comparison before interpreting
  results.
\item
  \textbf{Measurement discipline:} record units, uncertainty, sample
  size, and any discarded observation with a reason.
\item
  \textbf{Mechanistic link:} connect one result directly to the parent
  chapter's big idea before writing the conclusion.
\item
  \textbf{Reproducibility check:} state one procedural detail that
  another group would need in order to reproduce the result.
\end{itemize}

\subsection*{Paper-Based Evidence
Upgrade}\label{paper-based-evidence-upgrade-28}
\addcontentsline{toc}{subsection}{Paper-Based Evidence Upgrade}

Before answering the analysis questions, annotate the paper dataset for
\textbf{Plant Reproduction and Development} with a reproducibility pass:

{\def\LTcaptype{none} % do not increment counter
\begin{longtable}[]{@{}
  >{\raggedright\arraybackslash}p{(\linewidth - 2\tabcolsep) * \real{0.5000}}
  >{\raggedright\arraybackslash}p{(\linewidth - 2\tabcolsep) * \real{0.5000}}@{}}
\toprule\noalign{}
\begin{minipage}[b]{\linewidth}\raggedright
Evidence check
\end{minipage} & \begin{minipage}[b]{\linewidth}\raggedright
Student action
\end{minipage} \\
\midrule\noalign{}
\endhead
\bottomrule\noalign{}
\endlastfoot
Control logic & Mark the comparison that functions as the baseline,
negative control, or reference case. \\
Uncertainty & Circle the row, card, diagram feature, or model assumption
most likely to change the conclusion. \\
Field translation & Note whether the evidence comes from a diagram,
growth chamber, greenhouse, field plot, or long-term observation, and
state what changes when the scale changes. \\
Model comparison & State whether a simpler rule, null model, or
alternative mechanism could explain the same pattern. \\
Decision threshold & Write the minimum evidence that would make you
revise the interpretation. \\
Reproducibility & Record the exact scoring rule another group would need
to reproduce your classification. \\
\end{longtable}
}

Focus note: in this reproduction lab, tie every classification to the
reproductive structure or stage it depends on (sporophyte versus
gametophyte, pollen versus ovule, self- versus cross-pollination) and
state the pollination, development, or dispersal mechanism that links
genotype to the observed trait. Keep required work paper-based; any
material-handling or equipment version belongs only in an optional
extension.

\subsection*{Analysis Questions}\label{analysis-questions-28}
\addcontentsline{toc}{subsection}{Analysis Questions}

\begin{enumerate}
\def\labelenumi{\arabic{enumi}.}
\item
  Pollen tube growth is driven by cytoplasmic streaming and tip-directed
  vesicle fusion. What is the role of Ca\texorpdfstring{\textsuperscript{2}}{2}\texorpdfstring{\textsuperscript{+}}{+} ions (specifically the
  tip-focused Ca\texorpdfstring{\textsuperscript{2}}{2}\texorpdfstring{\textsuperscript{+}}{+} gradient) in regulating pollen tube polarity and
  growth direction?
\item
  Describe double fertilisation: where does each sperm go, what does
  each produce, and what are the ploidy levels of the zygote and
  endosperm?
\item
  A bee visits a fragrant, brightly coloured flower with sticky pollen.
  A grass flower has no petals, no fragrance, and produces vast
  quantities of dry pollen. Explain the evolutionary logic behind each
  strategy using cost-benefit trade-offs.
\item
  Self-incompatibility (SI) systems prevent self-fertilisation in many
  plant species. The S-RNase system in \emph{Petunia} involves the
  pistil secreting an RNase that degrades pollen tube RNA if the pollen
  S-allele matches the pistil S-allele. How does this mechanism
  selectively destroy ``self'' pollen without harming ``foreign''
  pollen?
\item
  Seedless watermelons are triploid (3n). Explain the steps used to
  produce them (starting from a colchicine treatment of diploid cells),
  and why triploid plants are sterile --- using the concept of
  homologous chromosome pairing during meiosis.
\end{enumerate}

\subsection*{Extension Analysis
Questions}\label{extension-analysis-questions-16}
\addcontentsline{toc}{subsection}{Extension Analysis Questions}

\begin{enumerate}
\def\labelenumi{\arabic{enumi}.}
\setcounter{enumi}{5}
\item
  \textbf{Statistical Analysis}: If the pollen germination dataset
  includes three replicates per treatment, how would you analyze
  variation in growth rates? What statistical test would you use to
  determine if differences between temperatures are significant?
\item
  \textbf{Experimental Error Analysis}: What are three potential sources
  of error in pollen tube growth measurements? How could you modify the
  procedure to reduce each source of error?
\item
  \textbf{Experimental Design}: If you wanted to test the effect of a
  specific hormone (e.g., auxin) on pollen tube growth using archived
  image measurements, what concentrations would you test, how would the
  images be blinded, and how would you measure the effect?
\end{enumerate}

\subsection*{Group Project Extension
(Multi-Session)}\label{group-project-extension-multi-session-13}
\addcontentsline{toc}{subsection}{Group Project Extension
(Multi-Session)}

\textbf{Design a Study of Plant Reproductive Strategies}

This project will span 2-3 lab sessions:

\textbf{Session 1}: Research Question and Hypothesis - How do flower
traits correlate with pollination syndrome? - Formulate a hypothesis
about the relationship between specific traits (e.g., petal colour,
nectar production) and pollinator type - Design a study comparing
multiple plant species

\textbf{Session 2}: Data Extraction and Analysis - Extract flower
morphology data from species cards, image databases, or published
datasets - Score pollinator visits from provided video stills or
summarized observation tables - Analyze correlations between traits and
pollinator type

\textbf{Session 3}: Data Interpretation and Presentation - Perform
statistical analysis (e.g., principal component analysis) - Discuss the
evolutionary implications of your findings - Create a scientific poster
or presentation

\subsection*{Real-World Problem Solving: Plant Reproduction and
Agriculture}\label{real-world-problem-solving-plant-reproduction-and-agriculture}
\addcontentsline{toc}{subsection}{Real-World Problem Solving: Plant
Reproduction and Agriculture}

\textbf{Case Study: Pollinator Decline}

\begin{enumerate}
\def\labelenumi{\arabic{enumi}.}
\item
  \textbf{Research Task}: Investigate the causes and consequences of
  pollinator decline (bees, butterflies, etc.). What crop groups depend
  on animal pollination, and how do dependence estimates differ between
  staple calories, fruits, nuts, and specialty crops?
\item
  \textbf{Field-Translation Consideration}: A warming treatment advances
  flowering in a chamber experiment, but a field dataset shows variable
  pollinator emergence. What evidence would let you decide whether the
  crop faces pollen heat failure, pollinator mismatch, or both?
\item
  \textbf{Policy Proposal}: Design a policy to protect pollinator
  populations while supporting agricultural productivity. Consider:
  incentives for habitat restoration, pesticide timing, heat-wave
  monitoring during flowering, and seed-set surveillance.
\end{enumerate}

\subsection*{Optional Wet-Lab
Demonstration}\label{optional-wet-lab-demonstration-1}
\addcontentsline{toc}{subsection}{Optional Wet-Lab Demonstration}

An instructor may demonstrate flower dissection or pollen germination
with appropriate allergy precautions and approved materials. This is not
required; the default lab uses diagrams, cards, and archived pollen-tube
measurements.

\subsection*{Safety and Ethics Notes}\label{safety-and-ethics-notes-28}
\addcontentsline{toc}{subsection}{Safety and Ethics Notes}

No wet materials are required. Pollen can be an allergen in optional
demonstrations, and colchicine is discussed conceptually primarily. When
discussing pollinator decline, consider the impacts on both ecosystems
and human food security.

\subsection*{Debrief and Reflection}\label{debrief-and-reflection-25}
\addcontentsline{toc}{subsection}{Debrief and Reflection}

After you finish the practical work, spend 5--10 minutes in your small
group comparing results and discussing the following prompts. Each
member should contribute at least one observation before moving to the
next prompt:

\begin{enumerate}
\def\labelenumi{\arabic{enumi}.}
\item
  \textbf{What did your measurements show} --- compare the group's
  results to the textbook's predictions. Where they diverge, suggest at
  least one mechanistic explanation before concluding ``experimental
  error.''
\item
  \textbf{What would change the outcome} --- propose one modification to
  the procedure that would sharpen the measurement or extend the result
  to a new biological context, and predict what you would observe.
\item
  \textbf{One-sentence headline} --- each student composes a single
  sentence summarising the lab's take-home message, suitable for a
  tweet. Compare sentences across groups; good headlines are short,
  quantitative, and mechanistic.
\item
  \textbf{Connection back to the textbook} --- identify one section of
  \cref{sec:unit_VIII_plant_reproduction} that your data either
  confirmed or complicated. Cite the specific passage.
\item
  \textbf{Experimental Design Reflection}: What was the strongest aspect
  of your experimental design? What would you change if you could repeat
  the experiment?
\end{enumerate}

\subsection*{Further Reading (Lab)}\label{further-reading-lab-25}
\addcontentsline{toc}{subsection}{Further Reading (Lab)}

\begin{itemize}
\tightlist
\item
  Revisit the parent chapter \cref{sec:unit_VIII_plant_reproduction} for
  the theoretical foundations on which this lab is built.
\item
  Look up any \textbf{bolded glossary term} introduced in the textbook
  chapter (each has a \texttt{\#gl:term-slug} link in the text) --- its
  master definition is in \breaktt{manuscript/glossary.md}.
\item
  Explore the Python code in \breaktt{src/biology/botany/botany.py} for
  plant reproduction calculations.
\item
  Use the self-contained Part 2 snippet as the computational template;
  it runs against tested project modules without external notebooks or
  CSV files.
\end{itemize}

\newpage

\section*{Lab 27 --- Plant Responses to the
Environment}\label{lab-27-plant-responses-to-the-environment}
\addcontentsline{toc}{section}{Lab 27 --- Plant Responses to the
Environment}

\label{sec:lab_unit_VIII_plant_responses}

\subsection*{Learning Objectives}\label{learning-objectives-67}
\addcontentsline{toc}{subsection}{Learning Objectives}

\emph{This activity accompanies \cref{sec:unit_VIII_plant_responses} of
the textbook --- review that chapter before attempting the exercises
below.}

\begin{itemize}
\tightlist
\item
  Interpret phototropism and gravitropism from time-series images and
  angle data
\item
  Measure the effect of exogenous auxin concentration from a printed
  dose-response dataset
\item
  Quantify the effect of gibberellin on dwarf pea elongation from
  replicated growth tables
\item
  Evaluate the role of phytochrome in seed germination through
  red/far-red treatment data
\end{itemize}

\subsubsection*{Alignment and Rubric
Map}\label{alignment-and-rubric-map-29}
\addcontentsline{toc}{subsubsection}{Alignment and Rubric Map}

\begin{itemize}
\tightlist
\item
  \textbf{Outcome 1 (LO1):} Interpret the supplied evidence or model
  output for Plant Responses to the Environment.
\item
  \textbf{Outcome 2 (LO2):} Identify controls and comparison groups that
  make the claim testable.
\item
  \textbf{Outcome 3 (LO3):} Quantify uncertainty, boundary conditions,
  or alternative explanations before concluding.
\item
  \textbf{Outcome 4 (LO4):} Transfer the mechanism to a new biological
  case or public-facing decision.
\item
  \textbf{Chapter LO coverage:} LO1, LO2, LO3, LO4
\item
  \textbf{Rubric dimensions:} evidence; controls; uncertainty;
  mechanism; transfer. \#\# Background \{.unnumbered\}
\end{itemize}

Plants integrate environmental signals and respond through
photoreceptors and hormone signalling. Auxin (IAA) is produced in the
shoot apex and promotes elongation at low concentrations but inhibits at
high concentrations (biphasic response). Gibberellins (GAs) promote stem
elongation and seed germination by targeting DELLA protein degradation.
Phototropism (bending toward light) results from asymmetric auxin
distribution driven by the PIN auxin efflux carriers. Phytochrome Pr
(inactive, absorbs red 660 nm) \texorpdfstring{\ensuremath{\leftrightarrow}}{<->} Pfr (active, absorbs far-red 730 nm)
provides a light quality sensor --- Pfr promotes seed germination and
flowering.

\subsection*{Paper-Based Materials}\label{paper-based-materials-29}
\addcontentsline{toc}{subsection}{Paper-Based Materials}

{\def\LTcaptype{none} % do not increment counter
\begin{longtable}[]{@{}
  >{\raggedright\arraybackslash}p{(\linewidth - 2\tabcolsep) * \real{0.3333}}
  >{\raggedright\arraybackslash}p{(\linewidth - 2\tabcolsep) * \real{0.6667}}@{}}
\toprule\noalign{}
\begin{minipage}[b]{\linewidth}\raggedright
Item
\end{minipage} & \begin{minipage}[b]{\linewidth}\raggedright
Quantity
\end{minipage} \\
\midrule\noalign{}
\endhead
\bottomrule\noalign{}
\endlastfoot
Phototropism and gravitropism image sequence with angle scale & 1
packet \\
Auxin dose-response dataset with replicated curvature measurements &
1 \\
Gibberellin dwarf-pea stem-length dataset with control and GA-treated
groups & 1 \\
Phytochrome red/far-red germination treatment table & 1 \\
Environmental cue cards for canopy shade, drought, salinity, and warming
& 1 set \\
Graph paper or spreadsheet template & 1 \\
Calculator & 1 \\
\end{longtable}
}

\subsection*{Paper-Based
Investigation}\label{paper-based-investigation-29}
\addcontentsline{toc}{subsection}{Paper-Based Investigation}

\textbf{Part A --- Phototropism Demonstration}

\begin{enumerate}
\def\labelenumi{\arabic{enumi}.}
\tightlist
\item
  Use the image sequence to compare unilateral light, darkness, and
  rotated-orientation treatments. Measure the angle of curvature from
  the printed scale and identify the baseline control.
\end{enumerate}

\textbf{Part B --- Auxin Dose-Response}

\begin{enumerate}
\def\labelenumi{\arabic{enumi}.}
\setcounter{enumi}{1}
\tightlist
\item
  Plot the replicated IAA curvature dataset on a log concentration axis.
  Mark the concentration with maximum response and calculate the
  coefficient of variation for each treatment.
\end{enumerate}

\textbf{Part C --- Gibberellin Effect on Dwarfism}

\begin{enumerate}
\def\labelenumi{\arabic{enumi}.}
\setcounter{enumi}{2}
\tightlist
\item
  Analyze dwarf pea stem-length data in two conditions: (a) control; (b)
  GA\texorpdfstring{\textsubscript{3}}{3}-treated. Calculate daily growth rate and the percent change
  relative to control.
\end{enumerate}

\textbf{Part D --- Phytochrome and Seed Germination}

\begin{enumerate}
\def\labelenumi{\arabic{enumi}.}
\setcounter{enumi}{3}
\tightlist
\item
  Use the phytochrome table for five treatment groups: (a) dark
  primarily; (b) red light pulse then dark; (c) far-red then dark; (d)
  red -\textgreater{} far-red -\textgreater{} dark; (e) far-red
  -\textgreater{} red -\textgreater{} dark.
\item
  Calculate percent germination and identify which final light pulse
  best predicts the response.
\end{enumerate}

\subsection*{Data Recording}\label{data-recording-26}
\addcontentsline{toc}{subsection}{Data Recording}

Phototropism: angle of curvature toward light = \textbf{\emph{°; dark
control = }}°

{\def\LTcaptype{none} % do not increment counter
\begin{longtable}[]{@{}ll@{}}
\toprule\noalign{}
IAA concentration (\texorpdfstring{\ensuremath{\mu}}{mu}M) & Angle of curvature (°) \\
\midrule\noalign{}
\endhead
\bottomrule\noalign{}
\endlastfoot
0 (lanolin primarily) & \\
0.1 & \\
1 & \\
10 & \\
100 & \\
\end{longtable}
}

Gibberellin --- dwarf pea stem length (mm):

{\def\LTcaptype{none} % do not increment counter
\begin{longtable}[]{@{}lll@{}}
\toprule\noalign{}
Day & Water control & +GA\texorpdfstring{\textsubscript{3}}{3} \\
\midrule\noalign{}
\endhead
\bottomrule\noalign{}
\endlastfoot
0 & & \\
2 & & \\
5 & & \\
\end{longtable}
}

Phytochrome germination results:

{\def\LTcaptype{none} % do not increment counter
\begin{longtable}[]{@{}lll@{}}
\toprule\noalign{}
Treatment & Germinated (of 20) & \% germination \\
\midrule\noalign{}
\endhead
\bottomrule\noalign{}
\endlastfoot
Dark primarily & & \\
Red \texorpdfstring{\ensuremath{\to}}{->} dark & & \\
Far-red \texorpdfstring{\ensuremath{\to}}{->} dark & & \\
Red \texorpdfstring{\ensuremath{\to}}{->} far-red \texorpdfstring{\ensuremath{\to}}{->} dark & & \\
Far-red \texorpdfstring{\ensuremath{\to}}{->} red \texorpdfstring{\ensuremath{\to}}{->} dark & & \\
\end{longtable}
}

Reproducibility check: treatment with largest mean response =
\textbf{\emph{; treatment with largest variation = }}; conclusion most
sensitive to one replicate = \_\_\_

\subsection*{Evidence and Reproducibility
Checklist}\label{evidence-and-reproducibility-checklist-29}
\addcontentsline{toc}{subsection}{Evidence and Reproducibility
Checklist}

\begin{itemize}
\tightlist
\item
  \textbf{Primary evidence goal:} Test plant response predictions under
  controlled environmental cues.
\item
  \textbf{Data skill to practice:} Interpret plant response data across
  light, gravity, water, and hormone treatments.
\item
  \textbf{BioSkills emphasis:} Interdisciplinary nature of science,
  Process of science.
\item
  \textbf{Control logic:} identify at least one positive control, one
  negative control, or one baseline comparison before interpreting
  results.
\item
  \textbf{Measurement discipline:} record units, uncertainty, sample
  size, and any discarded observation with a reason.
\item
  \textbf{Mechanistic link:} connect one result directly to the parent
  chapter's big idea before writing the conclusion.
\item
  \textbf{Reproducibility check:} state one procedural detail that
  another group would need in order to reproduce the result.
\end{itemize}

\subsection*{Paper-Based Evidence
Upgrade}\label{paper-based-evidence-upgrade-29}
\addcontentsline{toc}{subsection}{Paper-Based Evidence Upgrade}

Before answering the analysis questions, annotate the paper dataset for
\textbf{Plant Responses to the Environment} with a reproducibility pass:

{\def\LTcaptype{none} % do not increment counter
\begin{longtable}[]{@{}
  >{\raggedright\arraybackslash}p{(\linewidth - 2\tabcolsep) * \real{0.5000}}
  >{\raggedright\arraybackslash}p{(\linewidth - 2\tabcolsep) * \real{0.5000}}@{}}
\toprule\noalign{}
\begin{minipage}[b]{\linewidth}\raggedright
Evidence check
\end{minipage} & \begin{minipage}[b]{\linewidth}\raggedright
Student action
\end{minipage} \\
\midrule\noalign{}
\endhead
\bottomrule\noalign{}
\endlastfoot
Control logic & Mark the comparison that functions as the baseline,
negative control, or reference case. \\
Uncertainty & Circle the row, card, diagram feature, or model assumption
most likely to change the conclusion. \\
Field translation & State whether the response was measured in a
chamber, greenhouse, or field setting, and what environmental variable
is missing from that setting. \\
Model comparison & State whether a simpler rule, null model, or
alternative mechanism could explain the same pattern. \\
Decision threshold & Write the minimum evidence that would make you
revise the interpretation. \\
Reproducibility & Record the exact scoring rule another group would need
to reproduce your classification. \\
\end{longtable}
}

Focus note: in this environmental-response lab, trace each conclusion
along the perception-to-phenotype path --- name the stimulus, the
photoreceptor or hormone, the downstream gene-expression or signalling
step, and the growth-versus-defence tradeoff it implies --- before
treating any response as causal. Keep required work paper-based; any
material-handling or equipment version belongs only in an optional
extension.

\subsection*{Analysis Questions}\label{analysis-questions-29}
\addcontentsline{toc}{subsection}{Analysis Questions}

\begin{enumerate}
\def\labelenumi{\arabic{enumi}.}
\tightlist
\item
  In the phototropism experiment, which side of the shoot had longer
  cells? Explain how asymmetric auxin distribution (mediated by PIN
  proteins) creates differential elongation, and name the specific plant
  hormone involved.
\item
  Your auxin dose-response curve should show a bell shape. At what
  concentration did you see maximum curvature? Explain why very high
  auxin concentrations actually inhibit shoot elongation (relate to the
  acidic growth hypothesis and receptor downregulation).
\item
  GA-treated dwarf peas grew taller than controls. Explain the molecular
  mechanism: how does GA binding its receptor (GID1) lead to degradation
  of DELLA proteins and activate cell elongation genes?
\item
  Interpret your phytochrome germination results: which treatment(s)
  promoted germination? Which inhibited it? Explain why the red/far-red
  ratio in natural shade (under a canopy) would suppress germination of
  shade-intolerant species.
\item
  Long-day plants flower when nights are shorter than a critical length
  (they actually measure night length via phytochrome). Explain why a
  brief red light pulse in the middle of an otherwise long night would
  prevent flowering in a short-day plant, but a subsequent far-red pulse
  would restore it --- and what this tells us about the form of
  phytochrome that promotes flowering inhibition.
\item
  Choose one environmental cue card (drought, canopy shade, salinity, or
  warming). What field measurement would you need before turning the
  hormone-response result into a crop-management recommendation?
\end{enumerate}

\subsection*{Safety and Ethics Notes}\label{safety-and-ethics-notes-29}
\addcontentsline{toc}{subsection}{Safety and Ethics Notes}

No wet materials, live seedlings, hormones, or lamps are required.
Optional instructor demonstrations with seedlings, IAA, GA\texorpdfstring{\textsubscript{3}}{3}, or
red/far-red lights should follow local chemical and heat-safety rules
and remain separate from the required paper-based investigation.

\subsection*{Debrief and Reflection}\label{debrief-and-reflection-26}
\addcontentsline{toc}{subsection}{Debrief and Reflection}

After you finish the practical work, spend 5--10 minutes in your small
group comparing results and discussing the following prompts. Each
member should contribute at least one observation before moving to the
next prompt:

\begin{enumerate}
\def\labelenumi{\arabic{enumi}.}
\tightlist
\item
  \textbf{What did your measurements show} --- compare the group's
  results to the textbook's predictions. Where they diverge, suggest at
  least one mechanistic explanation before concluding ``experimental
  error.''
\item
  \textbf{What would change the outcome} --- propose one modification to
  the procedure that would sharpen the measurement or extend the result
  to a new biological context, and predict what you would observe.
\item
  \textbf{One-sentence headline} --- each student composes a single
  sentence summarising the lab's take-home message, suitable for a
  tweet. Compare sentences across groups; good headlines are short,
  quantitative, and mechanistic.
\item
  \textbf{Connection back to the textbook} --- identify one section of
  \cref{sec:unit_VIII_plant_responses} that your data either confirmed
  or complicated. Cite the specific passage.
\end{enumerate}

\subsection*{Further Reading (Lab)}\label{further-reading-lab-26}
\addcontentsline{toc}{subsection}{Further Reading (Lab)}

\begin{itemize}
\tightlist
\item
  Revisit the parent chapter \cref{sec:unit_VIII_plant_responses} for
  the theoretical foundations on which this lab is built.
\item
  Look up any \textbf{bolded glossary term} introduced in the textbook
  chapter (each has a \texttt{\#gl:term-slug} link in the text) --- its
  master definition is in \breaktt{manuscript/glossary.md}.
\item
  Explore the appended \texttt{src/} module that implements the
  corresponding quantitative model (when applicable) --- referenced in
  the parent chapter's ``Bridge to Computation'' subsection.
\end{itemize}

\emph{Module footer: parent chapter
\texttt{\textbackslash{}cref\{sec:unit\_VIII\_plant\_responses\}}; all
numerical quantities in this lab use SI units --- see Appendix D of the
textbook for unit conversions and biological-scale reference values.}

\newpage

\section*{Lab 28 --- Circulation, Respiration, and
Homeostasis}\label{lab-28-circulation-respiration-and-homeostasis}
\addcontentsline{toc}{section}{Lab 28 --- Circulation, Respiration, and
Homeostasis}

\label{sec:lab_unit_IX_circulation_respiration_homeostasis}

\subsection*{Learning Objectives}\label{learning-objectives-68}
\addcontentsline{toc}{subsection}{Learning Objectives}

\emph{This activity accompanies
\cref{sec:unit_IX_circulation_respiration_homeostasis} of the textbook
--- review that chapter before attempting the exercises below.}

\begin{itemize}
\tightlist
\item
  Analyze resting and post-exercise heart rate and blood pressure
  datasets and apply the Frank-Starling law
\item
  Calculate cardiac output and stroke volume from anonymized case data
  using CO = HR \texorpdfstring{\ensuremath{\times}}{x} SV
\item
  Apply Fick's principle to estimate oxygen consumption from heart rate
  data
\item
  Analyse spirometry data to determine tidal volume, IRV, ERV, and FVC
\end{itemize}

\subsubsection*{Alignment and Rubric
Map}\label{alignment-and-rubric-map-30}
\addcontentsline{toc}{subsubsection}{Alignment and Rubric Map}

\begin{itemize}
\tightlist
\item
  \textbf{Outcome 1 (LO1):} Interpret the supplied evidence or model
  output for Circulation, Respiration, and Homeostasis.
\item
  \textbf{Outcome 2 (LO2):} Identify controls and comparison groups that
  make the claim testable.
\item
  \textbf{Outcome 3 (LO3):} Quantify uncertainty, boundary conditions,
  or alternative explanations before concluding.
\item
  \textbf{Outcome 4 (LO4):} Transfer the mechanism to a new biological
  case or public-facing decision.
\item
  \textbf{Chapter LO coverage:} LO1, LO2, LO3, LO4
\item
  \textbf{Rubric dimensions:} evidence; controls; uncertainty;
  mechanism; transfer. \#\# Background \{.unnumbered\}
\end{itemize}

The cardiovascular system delivers oxygen and removes CO\texorpdfstring{\textsubscript{2}}{2} at a rate
matched to tissue demand. Cardiac output (CO) = stroke volume (SV) \texorpdfstring{\ensuremath{\times}}{x}
heart rate (HR). The Frank-Starling law states that greater ventricular
filling (preload) stretches sarcomeres to a more optimal myosin--actin
overlap, increasing stroke volume. Spirometry measures lung volumes:
tidal volume (TV about 500 mL), inspiratory reserve (IRV
about 3,000 mL), expiratory reserve (ERV about 1,200
mL), and FVC = (TV + IRV + ERV) for healthy adults.

\subsection*{Paper-Based Materials}\label{paper-based-materials-30}
\addcontentsline{toc}{subsection}{Paper-Based Materials}

{\def\LTcaptype{none} % do not increment counter
\begin{longtable}[]{@{}
  >{\raggedright\arraybackslash}p{(\linewidth - 2\tabcolsep) * \real{0.3333}}
  >{\raggedright\arraybackslash}p{(\linewidth - 2\tabcolsep) * \real{0.6667}}@{}}
\toprule\noalign{}
\begin{minipage}[b]{\linewidth}\raggedright
Item
\end{minipage} & \begin{minipage}[b]{\linewidth}\raggedright
Quantity
\end{minipage} \\
\midrule\noalign{}
\endhead
\bottomrule\noalign{}
\endlastfoot
Calculator & 1 \\
Graph paper or spreadsheet template & 1 \\
Cardiovascular response dataset with rest, exercise, recovery, and
heart-failure cases & 1 \\
Spirometry dataset printout (3 subjects: normal, asthmatic, COPD) & 1 \\
Altitude acclimatization case card with PaO2, ventilation, hematocrit,
and symptom data & 1 \\
Digestive-renal balance cards: meal composition, water intake, plasma
osmolality, urea, and urine concentration & 1 set \\
\end{longtable}
}

\subsection*{Paper-Based
Investigation}\label{paper-based-investigation-30}
\addcontentsline{toc}{subsection}{Paper-Based Investigation}

\textbf{Part A --- Cardiovascular Measurements}

\begin{enumerate}
\def\labelenumi{\arabic{enumi}.}
\tightlist
\item
  Use the cardiovascular response dataset to compare resting HR,
  estimated stroke volume, and recovery time across cases.
\item
  Plot HR at rest, 0, 1, 2, 5, and 10 minutes post-exercise for the
  printed subjects. Identify the baseline control and the slowest
  recovery profile.
\item
  Calculate CO = HR \texorpdfstring{\ensuremath{\times}}{x} SV at rest and immediately post-exercise. Check
  whether the data are consistent with the Frank-Starling mechanism and
  autonomic regulation.
\end{enumerate}

\textbf{Part B --- Spirometry Analysis}

\begin{enumerate}
\def\labelenumi{\arabic{enumi}.}
\setcounter{enumi}{3}
\tightlist
\item
  Use the printed spirometry dataset to calculate TV, FVC, FEV1, and
  FEV1/FVC where available.
\item
  Compare FVC and FEV1/FVC ratios for normal, asthmatic, and COPD
  subjects.
\item
  Use the altitude case card to connect PaO2 changes to ventilation,
  erythropoietin, hematocrit, and acclimatization.
\end{enumerate}

\textbf{Part C --- Digestive-Renal Homeostasis}

\begin{enumerate}
\def\labelenumi{\arabic{enumi}.}
\setcounter{enumi}{6}
\tightlist
\item
  Use the digestive-renal cards to trace how a high-protein salty meal
  changes gut absorption, hepatic urea production, plasma osmolality,
  ADH release, and urine concentration.
\item
  Compare a freshwater fish, a desert mammal, and a bird/reptile
  excretory strategy. Identify whether ammonia, urea, or uric acid is
  the dominant nitrogen-waste solution and what water trade-off that
  solution implies.
\end{enumerate}

\subsection*{Data Recording}\label{data-recording-27}
\addcontentsline{toc}{subsection}{Data Recording}

{\def\LTcaptype{none} % do not increment counter
\begin{longtable}[]{@{}llll@{}}
\toprule\noalign{}
Time (min) & HR (bpm) & Estimated SV (mL) & CO (L/min) \\
\midrule\noalign{}
\endhead
\bottomrule\noalign{}
\endlastfoot
Rest & & 70 & \\
0 (immediately post-exercise) & & 90 & \\
1 & & & \\
2 & & & \\
5 & & & \\
10 & & & \\
\end{longtable}
}

Spirometry (from dataset):

{\def\LTcaptype{none} % do not increment counter
\begin{longtable}[]{@{}lllll@{}}
\toprule\noalign{}
Subject & FVC (L) & FEV\texorpdfstring{\textsubscript{1}}{1} (L) & FEV\texorpdfstring{\textsubscript{1}}{1}/FVC (\%) & Pattern \\
\midrule\noalign{}
\endhead
\bottomrule\noalign{}
\endlastfoot
Normal & 4.8 & 4.3 & & \\
Asthmatic & 4.6 & 2.8 & & \\
COPD & 3.1 & 1.5 & & \\
\end{longtable}
}

Reproducibility check: sampling interval used for recovery graph =
\textbf{\emph{; rule for classifying obstructive pattern = }}

\subsection*{Evidence and Reproducibility
Checklist}\label{evidence-and-reproducibility-checklist-30}
\addcontentsline{toc}{subsection}{Evidence and Reproducibility
Checklist}

\begin{itemize}
\tightlist
\item
  \textbf{Primary evidence goal:} Model circulation, respiration, and
  feedback compensation.
\item
  \textbf{Data skill to practice:} Interpret physiological data from
  pressure, flow, saturation, or set-point changes.
\item
  \textbf{BioSkills emphasis:} Modeling and simulation, Quantitative
  reasoning, Science and society.
\item
  \textbf{Control logic:} identify at least one positive control, one
  negative control, or one baseline comparison before interpreting
  results.
\item
  \textbf{Measurement discipline:} record units, uncertainty, sample
  size, and any discarded observation with a reason.
\item
  \textbf{Mechanistic link:} connect one result directly to the parent
  chapter's big idea before writing the conclusion.
\item
  \textbf{Reproducibility check:} state one procedural detail that
  another group would need in order to reproduce the result.
\end{itemize}

\subsection*{Paper-Based Evidence
Upgrade}\label{paper-based-evidence-upgrade-30}
\addcontentsline{toc}{subsection}{Paper-Based Evidence Upgrade}

Before answering the analysis questions, annotate the paper dataset for
\textbf{Circulation, Respiration, and Homeostasis} with a
reproducibility pass:

{\def\LTcaptype{none} % do not increment counter
\begin{longtable}[]{@{}
  >{\raggedright\arraybackslash}p{(\linewidth - 2\tabcolsep) * \real{0.5000}}
  >{\raggedright\arraybackslash}p{(\linewidth - 2\tabcolsep) * \real{0.5000}}@{}}
\toprule\noalign{}
\begin{minipage}[b]{\linewidth}\raggedright
Evidence check
\end{minipage} & \begin{minipage}[b]{\linewidth}\raggedright
Student action
\end{minipage} \\
\midrule\noalign{}
\endhead
\bottomrule\noalign{}
\endlastfoot
Control logic & Mark the comparison that functions as the baseline,
negative control, or reference case. \\
Uncertainty & Circle the row, card, diagram feature, or model assumption
most likely to change the conclusion. \\
Model comparison & State whether a simpler rule, null model, or
alternative mechanism could explain the same pattern. \\
Decision threshold & Write the minimum evidence that would make you
revise the interpretation. \\
Reproducibility & Record the exact scoring rule another group would need
to reproduce your classification. \\
\end{longtable}
}

Focus note: in this lab, tie every cardiorespiratory claim to a flow or
gas-exchange measurement and the control loop it implies: state which
baseline, perturbation, and recovery values you used, and whether the
apparent regulation could instead reflect reserve capacity or the
resolution limit of the recording. Keep required work paper-based; any
material-handling or equipment version belongs only in an optional
extension.

\subsection*{Analysis Questions}\label{analysis-questions-30}
\addcontentsline{toc}{subsection}{Analysis Questions}

\begin{enumerate}
\def\labelenumi{\arabic{enumi}.}
\tightlist
\item
  Why did HR increase during exercise, and why did it remain elevated
  for several minutes post-exercise? Name the autonomic branch primarily
  responsible for the exercise-induced HR increase.
\item
  A reduced FEV\texorpdfstring{\textsubscript{1}}{1}/FVC ratio below the age-dependent lower limit of normal
  supports an obstructive pattern; a fixed 70\% cutoff is a screening
  shortcut, not a comprehensive diagnostic boundary. The asthmatic shows
  reduced FEV\texorpdfstring{\textsubscript{1}}{1} but relatively preserved FVC. Explain the structural
  difference in the airways that causes obstructive vs restrictive
  spirometry patterns.
\item
  Calculate the difference in cardiac output between rest and peak
  exercise in your data. Using Fick's principle (VO\texorpdfstring{\textsubscript{2}}{2} = CO \texorpdfstring{\ensuremath{\times}}{x} (CaO\texorpdfstring{\textsubscript{2}}{2} \texorpdfstring{\ensuremath{-}}{-}
  CvO\texorpdfstring{\textsubscript{2}}{2})), estimate oxygen consumption if the arterio-venous O\texorpdfstring{\textsubscript{2}}{2}
  difference increases from 50 mL/L at rest to 150 mL/L during maximal
  exercise.
\item
  A patient with congestive heart failure has chronically elevated
  ventricular end-diastolic pressure. Explain how this shifts the
  Frank-Starling curve and contributes to pulmonary oedema (connect to
  Starling capillary forces).
\item
  At high altitude (3,500 m), PaO\texorpdfstring{\textsubscript{2}}{2} drops from 100 to about 60
  mmHg. Identify three physiological acclimatisation responses (acute
  and chronic) and for each, specify the sensing mechanism and the
  effector system involved.
\item
  A dehydrated student eats a salty, high-protein meal. Predict the
  direction of change for plasma osmolality, ADH, urea production, and
  urine concentration. Which measurements would distinguish ``not enough
  water intake'' from impaired kidney concentrating ability?
\end{enumerate}

\subsection*{Safety and Ethics Notes}\label{safety-and-ethics-notes-30}
\addcontentsline{toc}{subsection}{Safety and Ethics Notes}

No exercise testing, blood pressure measurement, or spirometry is
required. Optional demonstrations involving human participants require
consent, screening, single-use mouthpiece filters where relevant, and
immediate stopping for dizziness, chest pain, or discomfort.

\subsection*{Debrief and Reflection}\label{debrief-and-reflection-27}
\addcontentsline{toc}{subsection}{Debrief and Reflection}

After you finish the practical work, spend 5--10 minutes in your small
group comparing results and discussing the following prompts. Each
member should contribute at least one observation before moving to the
next prompt:

\begin{enumerate}
\def\labelenumi{\arabic{enumi}.}
\tightlist
\item
  \textbf{What did your measurements show} --- compare the group's
  results to the textbook's predictions. Where they diverge, suggest at
  least one mechanistic explanation before concluding ``experimental
  error.''
\item
  \textbf{What would change the outcome} --- propose one modification to
  the procedure that would sharpen the measurement or extend the result
  to a new biological context, and predict what you would observe.
\item
  \textbf{One-sentence headline} --- each student composes a single
  sentence summarising the lab's take-home message, suitable for a
  tweet. Compare sentences across groups; good headlines are short,
  quantitative, and mechanistic.
\item
  \textbf{Connection back to the textbook} --- identify one section of
  \cref{sec:unit_IX_circulation_respiration_homeostasis} that your data
  either confirmed or complicated. Cite the specific passage.
\end{enumerate}

\subsection*{Further Reading (Lab)}\label{further-reading-lab-27}
\addcontentsline{toc}{subsection}{Further Reading (Lab)}

\begin{itemize}
\tightlist
\item
  Revisit the parent chapter
  \cref{sec:unit_IX_circulation_respiration_homeostasis} for the
  theoretical foundations on which this lab is built.
\item
  Look up any \textbf{bolded glossary term} introduced in the textbook
  chapter (each has a \texttt{\#gl:term-slug} link in the text) --- its
  master definition is in \breaktt{manuscript/glossary.md}.
\item
  Explore the appended \texttt{src/} module that implements the
  corresponding quantitative model (when applicable) --- referenced in
  the parent chapter's ``Bridge to Computation'' subsection.
\end{itemize}

\emph{Module footer: parent chapter
\texttt{\textbackslash{}cref\{sec:unit\_IX\_circulation\_respiration\_homeostasis\}};
all numerical quantities in this lab use SI units --- see Appendix D of
the textbook for unit conversions and biological-scale reference
values.}

\newpage

\section*{Lab 29 --- Nervous System and Neural
Signalling}\label{lab-29-nervous-system-and-neural-signalling}
\addcontentsline{toc}{section}{Lab 29 --- Nervous System and Neural
Signalling}

\label{sec:lab_unit_IX_nervous_system}

\subsection*{Learning Objectives}\label{learning-objectives-69}
\addcontentsline{toc}{subsection}{Learning Objectives}

\emph{This activity accompanies \cref{sec:unit_IX_nervous_system} of the
textbook --- review that chapter before attempting the exercises below.}

\begin{itemize}
\tightlist
\item
  Interpret reflex arc function using a clinical timing dataset and map
  receptive fields from paper data
\item
  Analyze two-point discrimination thresholds from anonymized class
  datasets
\item
  Analyse an action potential recording and identify threshold, peak,
  undershoot, and refractory periods
\item
  Apply the length constant equation to predict how signal decay differs
  in myelinated vs unmyelinated axons
\end{itemize}

\subsubsection*{Alignment and Rubric
Map}\label{alignment-and-rubric-map-31}
\addcontentsline{toc}{subsubsection}{Alignment and Rubric Map}

\begin{itemize}
\tightlist
\item
  \textbf{Outcome 1 (LO1):} Interpret the supplied evidence or model
  output for Nervous System and Neural Signalling.
\item
  \textbf{Outcome 2 (LO2):} Identify controls and comparison groups that
  make the claim testable.
\item
  \textbf{Outcome 3 (LO3):} Quantify uncertainty, boundary conditions,
  or alternative explanations before concluding.
\item
  \textbf{Outcome 4 (LO4):} Transfer the mechanism to a new biological
  case or public-facing decision.
\item
  \textbf{Chapter LO coverage:} LO1, LO2, LO3, LO4
\item
  \textbf{Rubric dimensions:} evidence; controls; uncertainty;
  mechanism; transfer. \#\# Background \{.unnumbered\}
\end{itemize}

Neural signalling depends on the passive spread of graded potentials
and, for long-distance signalling, the most-or-nothing action potential
(AP). The \textbf{two-point discrimination threshold} (Weber's two-point
test) reveals the density of somatosensory receptors in the skin:
fingertips (2--3 mm threshold) vs upper back (40--70 mm). The patellar
reflex is monosynaptic: stretching the patellar tendon activates Ia
afferents in the quadriceps \texorpdfstring{\ensuremath{\to}}{->} synapse on \texorpdfstring{\ensuremath{\alpha}}{alpha} motor neurons in the spinal
cord \texorpdfstring{\ensuremath{\to}}{->} quadriceps contraction. The entire arc takes \textless50 ms.

\subsection*{Paper-Based Materials}\label{paper-based-materials-31}
\addcontentsline{toc}{subsection}{Paper-Based Materials}

{\def\LTcaptype{none} % do not increment counter
\begin{longtable}[]{@{}
  >{\raggedright\arraybackslash}p{(\linewidth - 2\tabcolsep) * \real{0.3333}}
  >{\raggedright\arraybackslash}p{(\linewidth - 2\tabcolsep) * \real{0.6667}}@{}}
\toprule\noalign{}
\begin{minipage}[b]{\linewidth}\raggedright
Item
\end{minipage} & \begin{minipage}[b]{\linewidth}\raggedright
Quantity
\end{minipage} \\
\midrule\noalign{}
\endhead
\bottomrule\noalign{}
\endlastfoot
Patellar reflex case dataset with baseline, Jendrassik, and neuropathy
scenarios & 1 \\
Two-point discrimination dataset for fingertip, palm, forearm, and upper
back & 1 \\
Action potential recording trace (printed: one AP at scale with time
axis in ms, voltage axis in mV) & 1 \\
Cable properties worksheet (with \texorpdfstring{\ensuremath{\lambda}}{lambda} formula: \texorpdfstring{\ensuremath{\lambda}}{lambda} = \texorpdfstring{\(\sqrt{\;}\)}{sqrt}(r\_m/r\_i)) & 1 \\
Calculator & 1 \\
Body surface map (printed outline of human figure) & 1 \\
Behavioural ethogram and neural-prosthesis evidence cards & 1 set \\
\end{longtable}
}

\subsection*{Paper-Based
Investigation}\label{paper-based-investigation-31}
\addcontentsline{toc}{subsection}{Paper-Based Investigation}

\textbf{Part A --- Patellar Reflex}

\begin{enumerate}
\def\labelenumi{\arabic{enumi}.}
\tightlist
\item
  Use the reflex dataset to compare baseline, Jendrassik manoeuvre,
  peripheral neuropathy, and upper motor neuron scenarios. Record reflex
  magnitude and latency.
\item
  Identify the baseline control and the positive comparison that
  confirms the arc can increase with descending facilitation.
\item
  Map the likely lesion location for each abnormal case using the reflex
  arc diagram.
\end{enumerate}

\textbf{Part B --- Two-Point Discrimination}

\begin{enumerate}
\def\labelenumi{\arabic{enumi}.}
\setcounter{enumi}{3}
\tightlist
\item
  Analyze the anonymized two-point dataset at various spacings (1, 2, 5,
  10, 20, 40 mm) for the fingertip, palm, forearm, and upper back.
  Record the minimum spacing perceived as two points on each body area.
\end{enumerate}

\textbf{Part C --- Action Potential Trace Analysis}

\begin{enumerate}
\def\labelenumi{\arabic{enumi}.}
\setcounter{enumi}{4}
\tightlist
\item
  From the printed AP trace: measure (a) resting membrane potential
  (mV); (b) threshold voltage; (c) amplitude (peak mV); (d) duration of
  absolute refractory period (ms); (e) relative refractory period (ms);
  (f) undershoot (afterhyperpolarisation) level (mV).
\end{enumerate}

\textbf{Part D --- Cable Properties Calculation}

\begin{enumerate}
\def\labelenumi{\arabic{enumi}.}
\setcounter{enumi}{5}
\tightlist
\item
  Calculate the length constant \texorpdfstring{\ensuremath{\lambda}}{lambda} for two axons: (a) unmyelinated 0.5 \texorpdfstring{\ensuremath{\mu}}{mu}m
  diameter axon (r\_m = 10,000 \texorpdfstring{\ensuremath{\Omega}}{Omega}\texorpdfstring{\ensuremath{\cdot}}{.}cm, r\_i = 80 M\texorpdfstring{\ensuremath{\Omega}}{Omega}/cm); (b) myelinated 10
  \texorpdfstring{\ensuremath{\mu}}{mu}m diameter axon (r\_m = 800,000 \texorpdfstring{\ensuremath{\Omega}}{Omega}\texorpdfstring{\ensuremath{\cdot}}{.}cm, r\_i = 4 M\texorpdfstring{\ensuremath{\Omega}}{Omega}/cm). Determine
  V(x)/V\texorpdfstring{\textsubscript{0}}{0} at a distance of 1 mm.
\end{enumerate}

\textbf{Part E --- Behaviour and Neural Prosthesis Evidence}

\begin{enumerate}
\def\labelenumi{\arabic{enumi}.}
\setcounter{enumi}{6}
\tightlist
\item
  Sort each behaviour card into Tinbergen's four questions: mechanism,
  development, function, and evolutionary history.
\item
  Use the neural-prosthesis cards to separate what implanted recordings
  directly show (decodable neural activity) from the clinical claims
  that require longer-term evidence (implant stability, calibration,
  safety, and generalisation).
\end{enumerate}

\subsection*{Data Recording}\label{data-recording-28}
\addcontentsline{toc}{subsection}{Data Recording}

Patellar reflex: extension observed? (Y/N) \textbf{\emph{; Jendrassik
effect: }}

Two-point discrimination thresholds:

{\def\LTcaptype{none} % do not increment counter
\begin{longtable}[]{@{}ll@{}}
\toprule\noalign{}
Body area & Minimum two-point distance (mm) \\
\midrule\noalign{}
\endhead
\bottomrule\noalign{}
\endlastfoot
Fingertip & \\
Palm & \\
Forearm & \\
Upper back & \\
\end{longtable}
}

Action potential measurements from trace:

{\def\LTcaptype{none} % do not increment counter
\begin{longtable}[]{@{}ll@{}}
\toprule\noalign{}
Parameter & Value \\
\midrule\noalign{}
\endhead
\bottomrule\noalign{}
\endlastfoot
Resting membrane potential & mV \\
Threshold & mV \\
Peak & mV \\
AP amplitude (threshold\texorpdfstring{\ensuremath{\to}}{->}peak) & mV \\
Absolute refractory period & ms \\
Relative refractory period & ms \\
Undershoot & mV \\
\end{longtable}
}

Cable properties --- length constant \texorpdfstring{\ensuremath{\lambda}}{lambda}: Unmyelinated: \texorpdfstring{\ensuremath{\lambda}}{lambda} = \_\_\_ mm;
V(1mm)/V\texorpdfstring{\textsubscript{0}}{0} = \textbf{\emph{ Myelinated: \texorpdfstring{\ensuremath{\lambda}}{lambda} = }} mm; V(1mm)/V\texorpdfstring{\textsubscript{0}}{0} = \_\_\_

Reproducibility check: number of trials per body area = \textbf{\emph{;
threshold rule used (for example, two correct responses in a row) = }}

\subsection*{Evidence and Reproducibility
Checklist}\label{evidence-and-reproducibility-checklist-31}
\addcontentsline{toc}{subsection}{Evidence and Reproducibility
Checklist}

\begin{itemize}
\tightlist
\item
  \textbf{Primary evidence goal:} Map circuit structure onto predicted
  information flow.
\item
  \textbf{Data skill to practice:} Interpret neural data from anatomy,
  timing, or lesion evidence.
\item
  \textbf{BioSkills emphasis:} Modeling and simulation, Quantitative
  reasoning, Science and society.
\item
  \textbf{Control logic:} identify at least one positive control, one
  negative control, or one baseline comparison before interpreting
  results.
\item
  \textbf{Measurement discipline:} record units, uncertainty, sample
  size, and any discarded observation with a reason.
\item
  \textbf{Mechanistic link:} connect one result directly to the parent
  chapter's big idea before writing the conclusion.
\item
  \textbf{Reproducibility check:} state one procedural detail that
  another group would need in order to reproduce the result.
\end{itemize}

\subsection*{Paper-Based Evidence
Upgrade}\label{paper-based-evidence-upgrade-31}
\addcontentsline{toc}{subsection}{Paper-Based Evidence Upgrade}

Before answering the analysis questions, annotate the paper dataset for
\textbf{Nervous System and Neural Signalling} with a reproducibility
pass:

{\def\LTcaptype{none} % do not increment counter
\begin{longtable}[]{@{}
  >{\raggedright\arraybackslash}p{(\linewidth - 2\tabcolsep) * \real{0.5000}}
  >{\raggedright\arraybackslash}p{(\linewidth - 2\tabcolsep) * \real{0.5000}}@{}}
\toprule\noalign{}
\begin{minipage}[b]{\linewidth}\raggedright
Evidence check
\end{minipage} & \begin{minipage}[b]{\linewidth}\raggedright
Student action
\end{minipage} \\
\midrule\noalign{}
\endhead
\bottomrule\noalign{}
\endlastfoot
Control logic & Mark the comparison that functions as the baseline,
negative control, or reference case. \\
Uncertainty & Circle the row, card, diagram feature, or model assumption
most likely to change the conclusion. \\
Model comparison & State whether a simpler rule, null model, or
alternative mechanism could explain the same pattern. \\
Decision threshold & Write the minimum evidence that would make you
revise the interpretation. \\
Reproducibility & Record the exact scoring rule another group would need
to reproduce your classification. \\
\end{longtable}
}

Focus note: in this lab, keep the level of evidence explicit for every
neural claim: separate what the structural or anatomical data show from
what the recorded activity or behaviour shows, and state the scale
(cell, circuit, system) at which the inference is licensed before
generalising. Keep required work paper-based; any material-handling or
equipment version belongs only in an optional extension.

\subsection*{Analysis Questions}\label{analysis-questions-31}
\addcontentsline{toc}{subsection}{Analysis Questions}

\begin{enumerate}
\def\labelenumi{\arabic{enumi}.}
\tightlist
\item
  The fingertip had a much shorter two-point discrimination threshold
  than the upper back. Relate this to somatotopic mapping in the
  somatosensory cortex --- which area of the sensory homunculus
  represents fingers vs trunk, and what is the neural basis for this
  high-density representation?
\item
  The Jendrassik manoeuvre increased the patellar reflex. Explain this
  using the concept of motor neuron excitability and descending
  facilitation from the cortex/brainstem.
\item
  The action potential's absolute refractory period prevents reversed
  signal propagation (backward re-excitation). Explain how sodium
  channel inactivation (h-gate closed) creates this period and why this
  is essential for unidirectional propagation.
\item
  Your cable property calculation showed much higher \texorpdfstring{\ensuremath{\lambda}}{lambda} in the myelinated
  axon. Relate this to saltatory conduction: why does an action
  potential jump between nodes of Ranvier rather than regenerate
  continuously along the axon?
\item
  In multiple sclerosis, the myelin sheath is destroyed. Using your \texorpdfstring{\ensuremath{\lambda}}{lambda}
  calculations, predict what happens to action potential propagation in
  demyelinated axons, and explain why patients might experience
  increased sensitivity to heat (Uhthoff's phenomenon).
\item
  A speech BCI card reports high decoding accuracy from one participant
  during attempted speech. Which conclusion is directly supported, and
  which conclusions would require additional controls or longitudinal
  data?
\end{enumerate}

\subsection*{Safety and Ethics Notes}\label{safety-and-ethics-notes-31}
\addcontentsline{toc}{subsection}{Safety and Ethics Notes}

No touch-based neurological testing is required. If an instructor offers
an optional demonstration, student consent is required, taps must be
gentle, discriminator tips must be blunted, and students may opt out
without explanation.

\subsection*{Debrief and Reflection}\label{debrief-and-reflection-28}
\addcontentsline{toc}{subsection}{Debrief and Reflection}

After you finish the practical work, spend 5--10 minutes in your small
group comparing results and discussing the following prompts. Each
member should contribute at least one observation before moving to the
next prompt:

\begin{enumerate}
\def\labelenumi{\arabic{enumi}.}
\tightlist
\item
  \textbf{What did your measurements show} --- compare the group's
  results to the textbook's predictions. Where they diverge, suggest at
  least one mechanistic explanation before concluding ``experimental
  error.''
\item
  \textbf{What would change the outcome} --- propose one modification to
  the procedure that would sharpen the measurement or extend the result
  to a new biological context, and predict what you would observe.
\item
  \textbf{One-sentence headline} --- each student composes a single
  sentence summarising the lab's take-home message, suitable for a
  tweet. Compare sentences across groups; good headlines are short,
  quantitative, and mechanistic.
\item
  \textbf{Connection back to the textbook} --- identify one section of
  \cref{sec:unit_IX_nervous_system} that your data either confirmed or
  complicated. Cite the specific passage.
\end{enumerate}

\subsection*{Further Reading (Lab)}\label{further-reading-lab-28}
\addcontentsline{toc}{subsection}{Further Reading (Lab)}

\begin{itemize}
\tightlist
\item
  Revisit the parent chapter \cref{sec:unit_IX_nervous_system} for the
  theoretical foundations on which this lab is built.
\item
  Look up any \textbf{bolded glossary term} introduced in the textbook
  chapter (each has a \texttt{\#gl:term-slug} link in the text) --- its
  master definition is in \breaktt{manuscript/glossary.md}.
\item
  Explore the appended \texttt{src/} module that implements the
  corresponding quantitative model (when applicable) --- referenced in
  the parent chapter's ``Bridge to Computation'' subsection.
\end{itemize}

\emph{Module footer: parent chapter
\texttt{\textbackslash{}cref\{sec:unit\_IX\_nervous\_system\}}; all
numerical quantities in this lab use SI units --- see Appendix D of the
textbook for unit conversions and biological-scale reference values.}

\newpage

\section*{Lab 30 --- Action Potentials and Synaptic
Transmission}\label{lab-30-action-potentials-and-synaptic-transmission}
\addcontentsline{toc}{section}{Lab 30 --- Action Potentials and Synaptic
Transmission}

\label{sec:lab_unit_IX_action_potential_synapses}

\subsection*{Learning Objectives}\label{learning-objectives-70}
\addcontentsline{toc}{subsection}{Learning Objectives}

\emph{This activity accompanies
\cref{sec:unit_IX_action_potential_synapses} of the textbook --- review
that chapter before attempting the exercises below.}

\begin{itemize}
\tightlist
\item
  Analyse a compound action potential (CAP) recording from a nerve
  preparation to distinguish A, B, and C fibre types
\item
  Simulate quantal neurotransmitter release and calculate mean quantal
  content
\item
  Model the effect of GABA-A receptor modulation by benzodiazepines vs
  barbiturates
\item
  Evaluate synaptic plasticity mechanisms using an LTP induction
  protocol dataset
\end{itemize}

\subsubsection*{Alignment and Rubric
Map}\label{alignment-and-rubric-map-32}
\addcontentsline{toc}{subsubsection}{Alignment and Rubric Map}

\begin{itemize}
\tightlist
\item
  \textbf{Outcome 1 (LO1):} Interpret the supplied evidence or model
  output for Action Potentials and Synaptic Transmission.
\item
  \textbf{Outcome 2 (LO2):} Identify controls and comparison groups that
  make the claim testable.
\item
  \textbf{Outcome 3 (LO3):} Quantify uncertainty, boundary conditions,
  or alternative explanations before concluding.
\item
  \textbf{Outcome 4 (LO4):} Transfer the mechanism to a new biological
  case or public-facing decision.
\item
  \textbf{Chapter LO coverage:} LO1, LO2, LO3, LO4
\item
  \textbf{Rubric dimensions:} evidence; controls; uncertainty;
  mechanism; transfer. \#\# Background \{.unnumbered\}
\end{itemize}

Synaptic transmission involves vesicle fusion (Ca\texorpdfstring{\textsuperscript{2}}{2}\texorpdfstring{\textsuperscript{+}}{+}-triggered SNARE
complex assembly), neurotransmitter diffusion across the synaptic cleft,
and receptor binding on the postsynaptic membrane. Release is quantal:
the minimum unit is the vesicle (one quantum, about 5,000
molecules). The mean quantal content (m̄) can be estimated by the Poisson
failure method: m̄ = ln(N/F\texorpdfstring{\textsubscript{0}}{0}) where N = total trials and F\texorpdfstring{\textsubscript{0}}{0} = failures.
Long-term potentiation (LTP) at Schaffer collateral \texorpdfstring{\ensuremath{\to}}{->} CA1 synapses
requires NMDA receptor-mediated Ca\texorpdfstring{\textsuperscript{2}}{2}\texorpdfstring{\textsuperscript{+}}{+} influx and CaMKII
autophosphorylation.

\subsection*{Paper-Based Materials}\label{paper-based-materials-32}
\addcontentsline{toc}{subsection}{Paper-Based Materials}

{\def\LTcaptype{none} % do not increment counter
\begin{longtable}[]{@{}
  >{\raggedright\arraybackslash}p{(\linewidth - 2\tabcolsep) * \real{0.3333}}
  >{\raggedright\arraybackslash}p{(\linewidth - 2\tabcolsep) * \real{0.6667}}@{}}
\toprule\noalign{}
\begin{minipage}[b]{\linewidth}\raggedright
Item
\end{minipage} & \begin{minipage}[b]{\linewidth}\raggedright
Quantity
\end{minipage} \\
\midrule\noalign{}
\endhead
\bottomrule\noalign{}
\endlastfoot
Compound action potential trace (printed: 3 peaks corresponding to A\texorpdfstring{\ensuremath{\alpha}}{alpha},
A\texorpdfstring{\ensuremath{\delta}}{delta}, C fibres with time and conduction velocity axes) & 1 \\
Quantal release dataset (printed: 100 stimulus trials, number of quanta
released each trial, count of failures) & 1 \\
GABA-A receptor pharmacology worksheet & 1 \\
LTP induction dataset (EPSP slope before/after high-frequency
stimulation; two conditions: control and AP5-treated) & 1 \\
Calculator & 1 \\
Graph paper or spreadsheet template & 1 \\
\end{longtable}
}

\subsection*{Paper-Based
Investigation}\label{paper-based-investigation-32}
\addcontentsline{toc}{subsection}{Paper-Based Investigation}

\textbf{Part A --- Compound Action Potential Analysis}

\begin{enumerate}
\def\labelenumi{\arabic{enumi}.}
\tightlist
\item
  From the CAP trace, identify the three peaks (A\texorpdfstring{\ensuremath{\alpha}}{alpha}, A\texorpdfstring{\ensuremath{\delta}}{delta}, C fibre).
  Measure the latency of each peak from the stimulus artefact.
\item
  Given the distance from electrode to nerve (30 mm), calculate
  conduction velocity for each peak (velocity = distance/time).
\item
  Compare your calculated values to known conduction velocities (A\texorpdfstring{\ensuremath{\alpha}}{alpha}:
  70--120 m/s; A\texorpdfstring{\ensuremath{\delta}}{delta}: 5--30 m/s; C: 0.5--2 m/s).
\end{enumerate}

\textbf{Part B --- Quantal Analysis Simulation}

\begin{enumerate}
\def\labelenumi{\arabic{enumi}.}
\setcounter{enumi}{3}
\tightlist
\item
  From the dataset (100 trials, F\texorpdfstring{\textsubscript{0}}{0} = 8 failures): calculate m̄ by Poisson
  failure method (m̄ = ln(100/8) = ln(12.5)).
\item
  Also calculate m̄ by direct count method (m̄ = total quanta released /
  100 trials).
\item
  Estimate release probability: p̄ = m̄/n (where n = total available
  release sites, given as 15).
\end{enumerate}

\textbf{Part C --- GABA-A Pharmacology}

\begin{enumerate}
\def\labelenumi{\arabic{enumi}.}
\setcounter{enumi}{6}
\tightlist
\item
  From the worksheet, benzodiazepines increase the \textbf{frequency} of
  Cl\texorpdfstring{\textsuperscript{-}}{-} channel opening; barbiturates increase \textbf{duration} of Cl\texorpdfstring{\textsuperscript{-}}{-}
  channel opening (at low concentration) and open channels directly (at
  high concentration). On the I-V plot provided, annotate the effect of
  each drug class on Cl\texorpdfstring{\textsuperscript{-}}{-} current.
\end{enumerate}

\textbf{Part D --- LTP Dataset}

\begin{enumerate}
\def\labelenumi{\arabic{enumi}.}
\setcounter{enumi}{7}
\tightlist
\item
  Plot EPSP slope (\% baseline) vs time (minutes) for control and
  AP5-treated groups before and after high-frequency stimulation (HFS at
  t=0).
\item
  Calculate the LTP magnitude: (mean EPSP slope 30--60 min post-HFS /
  mean EPSP slope 10 min pre-HFS) \texorpdfstring{\ensuremath{\times}}{x} 100\%.
\item
  Reproducibility check: record the baseline window, post-HFS window,
  and any outlier rule before calculating LTP magnitude.
\end{enumerate}

\subsection*{Data Recording}\label{data-recording-29}
\addcontentsline{toc}{subsection}{Data Recording}

{\def\LTcaptype{none} % do not increment counter
\begin{longtable}[]{@{}
  >{\raggedright\arraybackslash}p{(\linewidth - 6\tabcolsep) * \real{0.1613}}
  >{\raggedright\arraybackslash}p{(\linewidth - 6\tabcolsep) * \real{0.1935}}
  >{\raggedright\arraybackslash}p{(\linewidth - 6\tabcolsep) * \real{0.4194}}
  >{\raggedright\arraybackslash}p{(\linewidth - 6\tabcolsep) * \real{0.2258}}@{}}
\toprule\noalign{}
\begin{minipage}[b]{\linewidth}\raggedright
Fibre type
\end{minipage} & \begin{minipage}[b]{\linewidth}\raggedright
Latency (ms)
\end{minipage} & \begin{minipage}[b]{\linewidth}\raggedright
Conduction velocity (m/s)
\end{minipage} & \begin{minipage}[b]{\linewidth}\raggedright
Expected range
\end{minipage} \\
\midrule\noalign{}
\endhead
\bottomrule\noalign{}
\endlastfoot
Peak 1 (A\texorpdfstring{\ensuremath{\alpha}}{alpha}) & & & 70--120 m/s \\
Peak 2 (A\texorpdfstring{\ensuremath{\delta}}{delta}) & & & 5--30 m/s \\
Peak 3 (C) & & & 0.5--2 m/s \\
\end{longtable}
}

Quantal analysis: F\texorpdfstring{\textsubscript{0}}{0} = 8; N = 100 m̄ (Poisson) = ln(100/8) =
\textbf{\emph{ m̄ (direct count) from data = }} p̄ = m̄/n (n=15) = \_\_\_

LTP magnitude: Control = \textbf{\emph{\%; AP5-treated = }}\%

Reproducibility check: baseline window = \textbf{\emph{; post-HFS window
= }}; outlier rule = \_\_\_

\subsection*{Evidence and Reproducibility
Checklist}\label{evidence-and-reproducibility-checklist-32}
\addcontentsline{toc}{subsection}{Evidence and Reproducibility
Checklist}

\begin{itemize}
\tightlist
\item
  \textbf{Primary evidence goal:} Model excitability and synaptic
  response under changed ion or channel conditions.
\item
  \textbf{Data skill to practice:} Interpret voltage traces, conductance
  changes, and synaptic perturbations.
\item
  \textbf{BioSkills emphasis:} Modeling and simulation, Quantitative
  reasoning, Science and society.
\item
  \textbf{Control logic:} identify at least one positive control, one
  negative control, or one baseline comparison before interpreting
  results.
\item
  \textbf{Measurement discipline:} record units, uncertainty, sample
  size, and any discarded observation with a reason.
\item
  \textbf{Mechanistic link:} connect one result directly to the parent
  chapter's big idea before writing the conclusion.
\item
  \textbf{Reproducibility check:} state one procedural detail that
  another group would need in order to reproduce the result.
\end{itemize}

\subsection*{Paper-Based Evidence
Upgrade}\label{paper-based-evidence-upgrade-32}
\addcontentsline{toc}{subsection}{Paper-Based Evidence Upgrade}

Before answering the analysis questions, annotate the paper dataset for
\textbf{Action Potentials and Synaptic Transmission} with a
reproducibility pass:

{\def\LTcaptype{none} % do not increment counter
\begin{longtable}[]{@{}
  >{\raggedright\arraybackslash}p{(\linewidth - 2\tabcolsep) * \real{0.5000}}
  >{\raggedright\arraybackslash}p{(\linewidth - 2\tabcolsep) * \real{0.5000}}@{}}
\toprule\noalign{}
\begin{minipage}[b]{\linewidth}\raggedright
Evidence check
\end{minipage} & \begin{minipage}[b]{\linewidth}\raggedright
Student action
\end{minipage} \\
\midrule\noalign{}
\endhead
\bottomrule\noalign{}
\endlastfoot
Control logic & Mark the comparison that functions as the baseline,
negative control, or reference case. \\
Uncertainty & Circle the row, card, diagram feature, or model assumption
most likely to change the conclusion. \\
Model comparison & State whether a simpler rule, null model, or
alternative mechanism could explain the same pattern. \\
Decision threshold & Write the minimum evidence that would make you
revise the interpretation. \\
Reproducibility & Record the exact scoring rule another group would need
to reproduce your classification. \\
\end{longtable}
}

Focus note: in this lab, treat each excitability or transmission claim
as a timing argument: ask which trace, latency, or quantal statistic
supports it, and whether the conduction-velocity or release estimate
would survive a change in fibre type, temperature, or
release-probability heterogeneity. Keep required work paper-based; any
material-handling or equipment version belongs only in an optional
extension.

\subsection*{Analysis Questions}\label{analysis-questions-32}
\addcontentsline{toc}{subsection}{Analysis Questions}

\begin{enumerate}
\def\labelenumi{\arabic{enumi}.}
\tightlist
\item
  C fibres (unmyelinated) have much slower conduction velocities than A\texorpdfstring{\ensuremath{\alpha}}{alpha}
  fibres (myelinated). Using your cable property knowledge from Lab 29,
  explain the molecular basis for this difference.
\item
  The two methods for calculating m̄ gave similar but not identical
  answers. Which method is more reliable when there is heterogeneity in
  release probabilities across synapses? Explain the statistical
  assumption violated by each method.
\item
  A patient given a benzodiazepine (e.g., diazepam) overdose can be
  treated with flumazenil (a competitive benzodiazepine antagonist).
  Explain why flumazenil doesn't cause seizures on its own even though
  it fully blocks GABA-A modulation, while complete GABA-A inhibition by
  bicuculline does cause seizures.
\item
  AP5 (an NMDA receptor antagonist) blocked LTP induction. Explain the
  molecular events from NMDA receptor opening to CaMKII
  autophosphorylation to AMPA receptor insertion (at least 5 steps).
\item
  Ketamine is an NMDA receptor open-channel blocker used as an
  anaesthetic and, at sub-anaesthetic doses, as a rapid antidepressant.
  Propose a mechanism by which blocking NMDA receptors could relieve
  depression --- consider what happens to synaptic strength in GABAergic
  interneurons when their NMDA receptors are blocked.
\end{enumerate}

\subsection*{Safety and Ethics Notes}\label{safety-and-ethics-notes-32}
\addcontentsline{toc}{subsection}{Safety and Ethics Notes}

Data-analysis and simulation primarily --- no chemical or biological
hazards. Discussions of neuropharmacology (drug mechanisms) should be
factual and clinical in tone.

\subsection*{Debrief and Reflection}\label{debrief-and-reflection-29}
\addcontentsline{toc}{subsection}{Debrief and Reflection}

After you finish the practical work, spend 5--10 minutes in your small
group comparing results and discussing the following prompts. Each
member should contribute at least one observation before moving to the
next prompt:

\begin{enumerate}
\def\labelenumi{\arabic{enumi}.}
\tightlist
\item
  \textbf{What did your measurements show} --- compare the group's
  results to the textbook's predictions. Where they diverge, suggest at
  least one mechanistic explanation before concluding ``experimental
  error.''
\item
  \textbf{What would change the outcome} --- propose one modification to
  the procedure that would sharpen the measurement or extend the result
  to a new biological context, and predict what you would observe.
\item
  \textbf{One-sentence headline} --- each student composes a single
  sentence summarising the lab's take-home message, suitable for a
  tweet. Compare sentences across groups; good headlines are short,
  quantitative, and mechanistic.
\item
  \textbf{Connection back to the textbook} --- identify one section of
  \cref{sec:unit_IX_action_potential_synapses} that your data either
  confirmed or complicated. Cite the specific passage.
\end{enumerate}

\subsection*{Further Reading (Lab)}\label{further-reading-lab-29}
\addcontentsline{toc}{subsection}{Further Reading (Lab)}

\begin{itemize}
\tightlist
\item
  Revisit the parent chapter
  \cref{sec:unit_IX_action_potential_synapses} for the theoretical
  foundations on which this lab is built.
\item
  Look up any \textbf{bolded glossary term} introduced in the textbook
  chapter (each has a \texttt{\#gl:term-slug} link in the text) --- its
  master definition is in \breaktt{manuscript/glossary.md}.
\item
  Explore the appended \texttt{src/} module that implements the
  corresponding quantitative model (when applicable) --- referenced in
  the parent chapter's ``Bridge to Computation'' subsection.
\end{itemize}

\emph{Module footer: parent chapter
\texttt{\textbackslash{}cref\{sec:unit\_IX\_action\_potential\_synapses\}};
all numerical quantities in this lab use SI units --- see Appendix D of
the textbook for unit conversions and biological-scale reference
values.}

\newpage

\section*{Lab 31 --- Endocrine and Immune
Systems}\label{lab-31-endocrine-and-immune-systems}
\addcontentsline{toc}{section}{Lab 31 --- Endocrine and Immune Systems}

\label{sec:lab_unit_IX_endocrine_and_immune}

\subsection*{Learning Objectives}\label{learning-objectives-71}
\addcontentsline{toc}{subsection}{Learning Objectives}

\emph{This activity accompanies \cref{sec:unit_IX_endocrine_and_immune}
of the textbook --- review that chapter before attempting the exercises
below.}

\begin{itemize}
\tightlist
\item
  Interpret blood hormone panels to diagnose endocrine disorders (hypo-
  vs hyperthyroidism, diabetes, Cushing's)
\item
  Analyse a cytokine profile to determine the type and severity of an
  inflammatory response
\item
  Model the HPT axis feedback loop using a paper-based regulatory
  diagram
\item
  Evaluate innate immune pattern recognition using TLR ligand-response
  data
\end{itemize}

\subsubsection*{Alignment and Rubric
Map}\label{alignment-and-rubric-map-33}
\addcontentsline{toc}{subsubsection}{Alignment and Rubric Map}

\begin{itemize}
\tightlist
\item
  \textbf{Outcome 1 (LO1):} Interpret the supplied evidence or model
  output for Endocrine and Immune Systems.
\item
  \textbf{Outcome 2 (LO2):} Identify controls and comparison groups that
  make the claim testable.
\item
  \textbf{Outcome 3 (LO3):} Quantify uncertainty, boundary conditions,
  or alternative explanations before concluding.
\item
  \textbf{Outcome 4 (LO4):} Transfer the mechanism to a new biological
  case or public-facing decision.
\item
  \textbf{Chapter LO coverage:} LO1, LO2, LO3, LO4
\item
  \textbf{Rubric dimensions:} evidence; controls; uncertainty;
  mechanism; transfer. \#\# Background \{.unnumbered\}
\end{itemize}

The endocrine system uses hormones (chemical messengers released into
the bloodstream) to coordinate physiology across long timescales.
Negative feedback loops maintain homeostasis: in the HPT axis, TRH
(hypothalamus) \texorpdfstring{\ensuremath{\to}}{->} TSH (pituitary) \texorpdfstring{\ensuremath{\to}}{->} T3/T4 (thyroid) \texorpdfstring{\ensuremath{\to}}{->} inhibit TRH and
TSH. In the innate immune system, Toll-like receptors (TLRs) detect
conserved pathogen-associated molecular patterns (PAMPs; e.g., LPS,
flagellin) and activate NF-\texorpdfstring{\ensuremath{\kappa}}{kappa}B, triggering pro-inflammatory cytokine
production (IL-6, TNF-\texorpdfstring{\ensuremath{\alpha}}{alpha}, IL-1\texorpdfstring{\ensuremath{\beta}}{beta}).

\subsection*{Paper-Based Materials}\label{paper-based-materials-33}
\addcontentsline{toc}{subsection}{Paper-Based Materials}

{\def\LTcaptype{none} % do not increment counter
\begin{longtable}[]{@{}
  >{\raggedright\arraybackslash}p{(\linewidth - 2\tabcolsep) * \real{0.3333}}
  >{\raggedright\arraybackslash}p{(\linewidth - 2\tabcolsep) * \real{0.6667}}@{}}
\toprule\noalign{}
\begin{minipage}[b]{\linewidth}\raggedright
Item
\end{minipage} & \begin{minipage}[b]{\linewidth}\raggedright
Quantity
\end{minipage} \\
\midrule\noalign{}
\endhead
\bottomrule\noalign{}
\endlastfoot
Endocrine case study data cards (4 cases: TSH, T4, cortisol, fasting
glucose, HbA1c values) & 1 set \\
HPT axis feedback diagram template (blank: boxes for hypothalamus,
pituitary, thyroid) & 1 \\
Cytokine data table (printed: IL-6, TNF-\texorpdfstring{\ensuremath{\alpha}}{alpha}, CRP, PCT by condition:
sepsis, mild infection, healthy) & 1 \\
TLR ligand-response worksheet (TLR4-LPS \texorpdfstring{\ensuremath{\to}}{->} NF-\texorpdfstring{\ensuremath{\kappa}}{kappa}B \texorpdfstring{\ensuremath{\to}}{->} cytokines; effects of
TAK-242 inhibitor) & 1 \\
Calculator & 1 \\
Reference ranges table & 1 \\
Reproductive-development timing cards: puberty, ovulation, implantation,
pregnancy, lactation, and infertility triage & 1 set \\
\end{longtable}
}

\subsection*{Paper-Based
Investigation}\label{paper-based-investigation-33}
\addcontentsline{toc}{subsection}{Paper-Based Investigation}

\textbf{Part A --- Endocrine Diagnosis}

\begin{enumerate}
\def\labelenumi{\arabic{enumi}.}
\tightlist
\item
  For each of 4 case study patients, compare provided hormone values to
  reference ranges. Diagnose: (a) primary hypothyroidism; (b)
  hyperthyroidism (Graves' disease); (c) type 2 diabetes; (d) Cushing's
  syndrome (hypercortisolism).
\item
  For each case, state which hormone is abnormal, whether it is elevated
  or suppressed, and the feedback loop explanation.
\end{enumerate}

\textbf{Part B --- HPT Axis Feedback Model}

\begin{enumerate}
\def\labelenumi{\arabic{enumi}.}
\setcounter{enumi}{2}
\tightlist
\item
  Complete the blank HPT axis diagram: draw arrows showing stimulatory
  (\texorpdfstring{\ensuremath{\to}}{->}) and inhibitory (\texorpdfstring{\ensuremath{\dashv}}{-|}) connections. Annotate what happens to each
  hormone level in primary hypothyroidism (thyroid fails) vs secondary
  hypothyroidism (pituitary fails).
\end{enumerate}

\textbf{Part C --- Cytokine Analysis}

\begin{enumerate}
\def\labelenumi{\arabic{enumi}.}
\setcounter{enumi}{3}
\tightlist
\item
  From the cytokine data table, calculate the ratio of each cytokine
  (sepsis) / (healthy). Rank the three markers by fold-change to
  determine which is most useful as a biomarker for severe sepsis.
\item
  Mark the clinical action threshold for PCT (procalcitonin
  \textgreater{} 2 ng/mL \texorpdfstring{\ensuremath{\to}}{->} antibiotic treatment recommended) on a bar
  chart.
\end{enumerate}

\textbf{Part D --- TLR4/NF-\texorpdfstring{\ensuremath{\kappa}}{kappa}B Pathway}

\begin{enumerate}
\def\labelenumi{\arabic{enumi}.}
\setcounter{enumi}{5}
\tightlist
\item
  Using the TLR ligand-response worksheet: trace LPS \texorpdfstring{\ensuremath{\to}}{->} TLR4 \texorpdfstring{\ensuremath{\to}}{->} MyD88 \texorpdfstring{\ensuremath{\to}}{->}
  IRAK \texorpdfstring{\ensuremath{\to}}{->} TRAF6 \texorpdfstring{\ensuremath{\to}}{->} IKK \texorpdfstring{\ensuremath{\to}}{->} I\texorpdfstring{\ensuremath{\kappa}}{kappa}B phosphorylation \texorpdfstring{\ensuremath{\to}}{->} NF-\texorpdfstring{\ensuremath{\kappa}}{kappa}B nuclear translocation
  \texorpdfstring{\ensuremath{\to}}{->} IL-6, TNF-\texorpdfstring{\ensuremath{\alpha}}{alpha} expression. Indicate where TAK-242 blocks the pathway.
\item
  Reproducibility check: state the reference range source, diagnostic
  threshold, and fold-change rule used for every case classification.
\end{enumerate}

\textbf{Part E --- Reproductive and Developmental Endocrinology}

\begin{enumerate}
\def\labelenumi{\arabic{enumi}.}
\setcounter{enumi}{7}
\tightlist
\item
  Sort the reproductive-development timing cards into upstream signal,
  target tissue, critical window, and measurable evidence.
\item
  For an infertility triage card, identify whether the first evidence
  gap is ovulation timing, sperm production, tubal transport, uterine
  receptivity, endocrine feedback, or age-dependent gamete quality.
\end{enumerate}

\subsection*{Data Recording}\label{data-recording-30}
\addcontentsline{toc}{subsection}{Data Recording}

{\def\LTcaptype{none} % do not increment counter
\begin{longtable}[]{@{}llllll@{}}
\toprule\noalign{}
Case & TSH & T4 & Cortisol & Glucose/HbA1c & Diagnosis \\
\midrule\noalign{}
\endhead
\bottomrule\noalign{}
\endlastfoot
A & 12 mIU/L \texorpdfstring{\ensuremath{\uparrow}}{up} & 7 pmol/L \texorpdfstring{\ensuremath{\downarrow}}{down} & normal & normal & \\
B & 0.02 mIU/L \texorpdfstring{\ensuremath{\downarrow}}{down} & 42 pmol/L \texorpdfstring{\ensuremath{\uparrow}}{up} & normal & normal & \\
C & normal & normal & 890 nmol/L \texorpdfstring{\ensuremath{\uparrow}}{up} & 8.1 mmol/L \texorpdfstring{\ensuremath{\uparrow}}{up} & \\
D & normal & normal & normal & 11 mmol/L \texorpdfstring{\ensuremath{\uparrow}}{up}; HbA1c 9.2\% \texorpdfstring{\ensuremath{\uparrow}}{up} & \\
\end{longtable}
}

Cytokine fold-changes:

{\def\LTcaptype{none} % do not increment counter
\begin{longtable}[]{@{}lllll@{}}
\toprule\noalign{}
Marker & Healthy & Sepsis & Fold-change & Rank \\
\midrule\noalign{}
\endhead
\bottomrule\noalign{}
\endlastfoot
IL-6 (pg/mL) & 5 & 1,200 & & \\
TNF-\texorpdfstring{\ensuremath{\alpha}}{alpha} (pg/mL) & 10 & 380 & & \\
PCT (ng/mL) & 0.05 & 18 & & \\
\end{longtable}
}

HPT axis diagram annotations (describe what happens in primary
hypothyroidism):

Reproducibility check: reference range source = \textbf{\emph{; cytokine
fold-change rule = }}; uncertain case needing more data = \_\_\_

\subsection*{Evidence and Reproducibility
Checklist}\label{evidence-and-reproducibility-checklist-33}
\addcontentsline{toc}{subsection}{Evidence and Reproducibility
Checklist}

\begin{itemize}
\tightlist
\item
  \textbf{Primary evidence goal:} Compare feedback loops and immune
  activation scenarios.
\item
  \textbf{Data skill to practice:} Interpret endocrine or immune data
  from time courses, titers, or perturbations.
\item
  \textbf{BioSkills emphasis:} Modeling and simulation, Quantitative
  reasoning, Science and society.
\item
  \textbf{Control logic:} identify at least one positive control, one
  negative control, or one baseline comparison before interpreting
  results.
\item
  \textbf{Measurement discipline:} record units, uncertainty, sample
  size, and any discarded observation with a reason.
\item
  \textbf{Mechanistic link:} connect one result directly to the parent
  chapter's big idea before writing the conclusion.
\item
  \textbf{Reproducibility check:} state one procedural detail that
  another group would need in order to reproduce the result.
\end{itemize}

\subsection*{Paper-Based Evidence
Upgrade}\label{paper-based-evidence-upgrade-33}
\addcontentsline{toc}{subsection}{Paper-Based Evidence Upgrade}

Before answering the analysis questions, annotate the paper dataset for
\textbf{Endocrine and Immune Systems} with a reproducibility pass:

{\def\LTcaptype{none} % do not increment counter
\begin{longtable}[]{@{}
  >{\raggedright\arraybackslash}p{(\linewidth - 2\tabcolsep) * \real{0.5000}}
  >{\raggedright\arraybackslash}p{(\linewidth - 2\tabcolsep) * \real{0.5000}}@{}}
\toprule\noalign{}
\begin{minipage}[b]{\linewidth}\raggedright
Evidence check
\end{minipage} & \begin{minipage}[b]{\linewidth}\raggedright
Student action
\end{minipage} \\
\midrule\noalign{}
\endhead
\bottomrule\noalign{}
\endlastfoot
Control logic & Mark the comparison that functions as the baseline,
negative control, or reference case. \\
Uncertainty & Circle the row, card, diagram feature, or model assumption
most likely to change the conclusion. \\
Model comparison & State whether a simpler rule, null model, or
alternative mechanism could explain the same pattern. \\
Decision threshold & Write the minimum evidence that would make you
revise the interpretation. \\
Reproducibility & Record the exact scoring rule another group would need
to reproduce your classification. \\
\end{longtable}
}

Focus note: in this lab, read each endocrine or immune claim as a
feedback-and-timing argument: identify the dose, latency, and
receptor-sensitivity assumptions behind it, and check whether an
inflammatory confound or shifted set-point would change the conclusion
before treating the response as specific. Keep required work
paper-based; any material-handling or equipment version belongs only in
an optional extension.

\subsection*{Analysis Questions}\label{analysis-questions-33}
\addcontentsline{toc}{subsection}{Analysis Questions}

\begin{enumerate}
\def\labelenumi{\arabic{enumi}.}
\tightlist
\item
  In Case A (primary hypothyroidism), TSH is elevated but T4 is low.
  Trace the negative feedback loop to explain why TSH is high despite
  low thyroid function. What drug would you prescribe, and which axis
  level does it target?
\item
  Graves' disease (Case B) involves thyroid-stimulating immunoglobulins
  (TSI) that mimic TSH. Why does the pituitary produce less TSH even
  though T4 is high? Can negative feedback suppress TSI production?
\item
  The cytokine data showed PCT was elevated 360-fold in sepsis while
  TNF-\texorpdfstring{\ensuremath{\alpha}}{alpha} was elevated 38-fold. Why is PCT particularly specific for
  bacterial infection over viral infection, and why does this make it a
  better guide for antibiotic prescribing?
\item
  NF-\texorpdfstring{\ensuremath{\kappa}}{kappa}B is activated during sepsis and produces the cytokine storm that
  contributes to multi-organ failure. A drug (TAK-242) blocks TLR4
  signalling. Predict the effect on: (a) cytokine levels; (b) pathogen
  clearance. What is the risk of fully blocking TLR4 signalling?
\item
  An athlete uses anabolic steroids (supraphysiological testosterone
  dose). Predict the effect on LH and FSH levels via the HPG axis
  feedback. After stopping steroids, how long does hormonal recovery
  take, and which cells in the testes are most affected during prolonged
  suppression?
\item
  A couple has not conceived after 12 months of regular unprotected
  intercourse. Build a first-pass evidence map that separates ovulation,
  sperm production, tubal transport, implantation, uterine receptivity,
  endocrine feedback, and age-dependent gamete quality.
\end{enumerate}

\subsection*{Safety and Ethics Notes}\label{safety-and-ethics-notes-33}
\addcontentsline{toc}{subsection}{Safety and Ethics Notes}

Data-analysis primarily. Most patient case data should be
anonymised/fictional. Discuss real clinical conditions (Cushing's,
diabetes) with sensitivity to students who may have personal connections
to these diagnoses.

\subsection*{Debrief and Reflection}\label{debrief-and-reflection-30}
\addcontentsline{toc}{subsection}{Debrief and Reflection}

After you finish the practical work, spend 5--10 minutes in your small
group comparing results and discussing the following prompts. Each
member should contribute at least one observation before moving to the
next prompt:

\begin{enumerate}
\def\labelenumi{\arabic{enumi}.}
\tightlist
\item
  \textbf{What did your measurements show} --- compare the group's
  results to the textbook's predictions. Where they diverge, suggest at
  least one mechanistic explanation before concluding ``experimental
  error.''
\item
  \textbf{What would change the outcome} --- propose one modification to
  the procedure that would sharpen the measurement or extend the result
  to a new biological context, and predict what you would observe.
\item
  \textbf{One-sentence headline} --- each student composes a single
  sentence summarising the lab's take-home message, suitable for a
  tweet. Compare sentences across groups; good headlines are short,
  quantitative, and mechanistic.
\item
  \textbf{Connection back to the textbook} --- identify one section of
  \cref{sec:unit_IX_endocrine_and_immune} that your data either
  confirmed or complicated. Cite the specific passage.
\end{enumerate}

\subsection*{Further Reading (Lab)}\label{further-reading-lab-30}
\addcontentsline{toc}{subsection}{Further Reading (Lab)}

\begin{itemize}
\tightlist
\item
  Revisit the parent chapter \cref{sec:unit_IX_endocrine_and_immune} for
  the theoretical foundations on which this lab is built.
\item
  Look up any \textbf{bolded glossary term} introduced in the textbook
  chapter (each has a \texttt{\#gl:term-slug} link in the text) --- its
  master definition is in \breaktt{manuscript/glossary.md}.
\item
  Explore the appended \texttt{src/} module that implements the
  corresponding quantitative model (when applicable) --- referenced in
  the parent chapter's ``Bridge to Computation'' subsection.
\end{itemize}

\emph{Module footer: parent chapter
\texttt{\textbackslash{}cref\{sec:unit\_IX\_endocrine\_and\_immune\}};
all numerical quantities in this lab use SI units --- see Appendix D of
the textbook for unit conversions and biological-scale reference
values.}

\newpage

\section*{Lab 32 --- Population Ecology and Growth
Models}\label{lab-32-population-ecology-and-growth-models}
\addcontentsline{toc}{section}{Lab 32 --- Population Ecology and Growth
Models}

\label{sec:lab_unit_X_population_ecology}

\subsection*{Learning Objectives}\label{learning-objectives-72}
\addcontentsline{toc}{subsection}{Learning Objectives}

\emph{This activity accompanies \cref{sec:unit_X_population_ecology} of
the textbook --- review that chapter before attempting the exercises
below.}

\begin{itemize}
\tightlist
\item
  Fit logistic and exponential growth equations to real or simulated
  population data
\item
  Calculate intrinsic rate of increase (r) and carrying capacity (K)
  from population time series
\item
  Apply the mark-recapture method (Lincoln-Petersen) to estimate
  population size
\item
  Evaluate density-dependent vs density-independent limiting factors
  from case study data
\end{itemize}

\subsubsection*{Alignment and Rubric
Map}\label{alignment-and-rubric-map-34}
\addcontentsline{toc}{subsubsection}{Alignment and Rubric Map}

\begin{itemize}
\tightlist
\item
  \textbf{Outcome 1 (LO1):} Interpret the supplied evidence or model
  output for Population Ecology and Growth Models.
\item
  \textbf{Outcome 2 (LO2):} Identify controls and comparison groups that
  make the claim testable.
\item
  \textbf{Outcome 3 (LO3):} Quantify uncertainty, boundary conditions,
  or alternative explanations before concluding.
\item
  \textbf{Outcome 4 (LO4):} Transfer the mechanism to a new biological
  case or public-facing decision.
\item
  \textbf{Chapter LO coverage:} LO1, LO2, LO3, LO4
\item
  \textbf{Rubric dimensions:} evidence; controls; uncertainty;
  mechanism; transfer. \#\# Background \{.unnumbered\}
\end{itemize}

Population growth follows exponential growth (dN/dt = rN) when resources
are unlimited, and logistic growth (dN/dt = rN(1 \texorpdfstring{\ensuremath{-}}{-} N/K)) as populations
approach the carrying capacity K. Real populations fluctuate around K
due to density-dependent factors (competition, predation, disease ---
intensify as N increases) or may crash due to density-independent
factors (storms, droughts --- unrelated to N). The Lincoln-Petersen
mark-recapture estimate: N̂ = (M \texorpdfstring{\ensuremath{\times}}{x} C) / R, where M = initial marked
individuals, C = recapture sample size, R = recaptured marked
individuals.

\subsection*{Paper-Based Materials}\label{paper-based-materials-34}
\addcontentsline{toc}{subsection}{Paper-Based Materials}

{\def\LTcaptype{none} % do not increment counter
\begin{longtable}[]{@{}
  >{\raggedright\arraybackslash}p{(\linewidth - 2\tabcolsep) * \real{0.3333}}
  >{\raggedright\arraybackslash}p{(\linewidth - 2\tabcolsep) * \real{0.6667}}@{}}
\toprule\noalign{}
\begin{minipage}[b]{\linewidth}\raggedright
Item
\end{minipage} & \begin{minipage}[b]{\linewidth}\raggedright
Quantity
\end{minipage} \\
\midrule\noalign{}
\endhead
\bottomrule\noalign{}
\endlastfoot
Population time series data (printed: 3 species --- moth (exponential),
snowshoe hare (oscillating), logistic yeast) & 1 \\
Graph paper or spreadsheet template & 1 \\
Calculator & 1 \\
Mark-recapture card deck or printed capture-history table & 1 \\
Management scenario cards for invasive species, harvested fisheries, and
climate-driven range shifts & 1 set \\
\end{longtable}
}

\subsection*{Paper-Based
Investigation}\label{paper-based-investigation-34}
\addcontentsline{toc}{subsection}{Paper-Based Investigation}

\textbf{Part A --- Growth Curve Fitting}

\begin{enumerate}
\def\labelenumi{\arabic{enumi}.}
\tightlist
\item
  Plot most three population time series on graph paper. For the yeast
  data, identify: lag phase, exponential phase, plateau (K).
\item
  Estimate K from the plateau of the logistic curve.
\item
  Estimate r from the exponential phase: r \texorpdfstring{\ensuremath{\approx}}{~} (ln N\texorpdfstring{\textsubscript{2}}{2} \texorpdfstring{\ensuremath{-}}{-} ln N\texorpdfstring{\textsubscript{1}}{1}) / (t\texorpdfstring{\textsubscript{2}}{2} \texorpdfstring{\ensuremath{-}}{-} t\texorpdfstring{\textsubscript{1}}{1})
  during the period of fastest growth.
\end{enumerate}

\textbf{Part B --- Mark-Recapture}

\begin{enumerate}
\def\labelenumi{\arabic{enumi}.}
\setcounter{enumi}{3}
\tightlist
\item
  Use the printed capture-history table: M = 50 initially marked
  individuals and C = 30 individuals in the recapture sample. Count the
  marked recaptures (R) in the table and calculate Nhat = (50 x 30) / R.
\item
  Repeat the calculation for the second printed sampling event. Compare
  both estimates with the known simulated population size and calculate
  percent error.
\end{enumerate}

\textbf{Part C --- Density Dependence Analysis}

\begin{enumerate}
\def\labelenumi{\arabic{enumi}.}
\setcounter{enumi}{5}
\tightlist
\item
  Given data on white-tailed deer population growth rate (r) vs density
  in 4 regions: plot r vs N/K (relative density). Identify whether r
  decreases with increasing density (density-dependent regulation) ---
  this confirms a logistic pattern.
\end{enumerate}

\subsection*{Data Recording}\label{data-recording-31}
\addcontentsline{toc}{subsection}{Data Recording}

Logistic yeast data:

{\def\LTcaptype{none} % do not increment counter
\begin{longtable}[]{@{}lll@{}}
\toprule\noalign{}
Time (h) & Population (cells/mL) & ln N \\
\midrule\noalign{}
\endhead
\bottomrule\noalign{}
\endlastfoot
0 & 10 & \\
6 & 80 & \\
12 & 410 & \\
18 & 1,200 & \\
24 & 2,800 & \\
36 & 3,900 & \\
48 & 4,100 & \\
\end{longtable}
}

Estimated K = \textbf{\emph{; Estimated r = }}

Mark-recapture: M = 50; C = 30; R = \textbf{\emph{; N̂ = }} Second
sampling estimate: R = \textbf{\emph{; N̂ = }}; percent error from known
simulated N = \_\_\_\%

\subsection*{Evidence and Reproducibility
Checklist}\label{evidence-and-reproducibility-checklist-34}
\addcontentsline{toc}{subsection}{Evidence and Reproducibility
Checklist}

\begin{itemize}
\tightlist
\item
  \textbf{Primary evidence goal:} Model population trajectories and
  compare sampling methods.
\item
  \textbf{Data skill to practice:} Use abundance or age-structure data
  to estimate growth and risk.
\item
  \textbf{BioSkills emphasis:} Modeling and simulation, Science and
  society, Communication and collaboration.
\item
  \textbf{Control logic:} identify at least one positive control, one
  negative control, or one baseline comparison before interpreting
  results.
\item
  \textbf{Measurement discipline:} record units, uncertainty, sample
  size, and any discarded observation with a reason.
\item
  \textbf{Mechanistic link:} connect one result directly to the parent
  chapter's big idea before writing the conclusion.
\item
  \textbf{Reproducibility check:} state one procedural detail that
  another group would need in order to reproduce the result.
\end{itemize}

\subsection*{Paper-Based Evidence
Upgrade}\label{paper-based-evidence-upgrade-34}
\addcontentsline{toc}{subsection}{Paper-Based Evidence Upgrade}

Before answering the analysis questions, annotate the paper dataset for
\textbf{Population Ecology and Growth Models} with a reproducibility
pass:

{\def\LTcaptype{none} % do not increment counter
\begin{longtable}[]{@{}
  >{\raggedright\arraybackslash}p{(\linewidth - 2\tabcolsep) * \real{0.5000}}
  >{\raggedright\arraybackslash}p{(\linewidth - 2\tabcolsep) * \real{0.5000}}@{}}
\toprule\noalign{}
\begin{minipage}[b]{\linewidth}\raggedright
Evidence check
\end{minipage} & \begin{minipage}[b]{\linewidth}\raggedright
Student action
\end{minipage} \\
\midrule\noalign{}
\endhead
\bottomrule\noalign{}
\endlastfoot
Control logic & Mark the comparison that functions as the baseline,
negative control, or reference case. \\
Uncertainty & Circle the row, card, diagram feature, or model assumption
most likely to change the conclusion. \\
Model comparison & State whether a simpler rule, null model, or
alternative mechanism could explain the same pattern. \\
Decision threshold & Write the minimum evidence that would make you
revise the interpretation. \\
Reproducibility & Record the exact scoring rule another group would need
to reproduce your classification. \\
\end{longtable}
}

Focus note: for this population lab, state the demographic model behind
each printed estimate (exponential, logistic, or mark-recapture), check
whether its assumptions about density dependence, closure, and detection
hold for the simulated capture histories, and tie any harvest or
recovery conclusion to an explicit management objective. Keep required
work paper-based; any material-handling or equipment version belongs
only in an optional extension.

\subsection*{Analysis Questions}\label{analysis-questions-34}
\addcontentsline{toc}{subsection}{Analysis Questions}

\begin{enumerate}
\def\labelenumi{\arabic{enumi}.}
\tightlist
\item
  Why does the yeast population growth slow as it approaches K? Use the
  term ``intraspecific competition'' and identify two specific
  density-dependent resources that become limiting.
\item
  Snowshoe hare populations cycle with about 10-year
  periodicity (tied to lynx cycles). Is this best explained as: (a)
  density-dependent predation; (b) density-independent weather events;
  or (c) a two-species coupled oscillator? Describe the evidence.
\item
  Your two mark-recapture estimates probably differed slightly. List
  three assumptions of the Lincoln-Petersen method and explain which
  assumption is most likely violated in the printed capture-history
  simulation.
\item
  Invasive species often exhibit exponential growth in new habitats.
  Name a specific invasive species, identify what density-dependent
  factor is absent in the new habitat, and predict when/how logistic
  behaviour will eventually emerge.
\item
  Apply a fisheries management scenario: a cod population has K =
  500,000 tonnes and r = 0.2/year. Using the maximum sustainable yield
  formula (MSY = rK/4), calculate the MSY and the population size at MSY
  (N = K/2). If fishing reduces the population below K/4, explain why
  the population will continue to decline even if fishing stops.
\end{enumerate}

\subsection*{Safety and Ethics Notes}\label{safety-and-ethics-notes-34}
\addcontentsline{toc}{subsection}{Safety and Ethics Notes}

No hazardous materials or living specimens are required. When discussing
harvested or invasive species, maintain ecological objectivity and
identify uncertainty in population estimates.

\subsection*{Debrief and Reflection}\label{debrief-and-reflection-31}
\addcontentsline{toc}{subsection}{Debrief and Reflection}

After you finish the practical work, spend 5--10 minutes in your small
group comparing results and discussing the following prompts. Each
member should contribute at least one observation before moving to the
next prompt:

\begin{enumerate}
\def\labelenumi{\arabic{enumi}.}
\tightlist
\item
  \textbf{What did your measurements show} --- compare the group's
  results to the textbook's predictions. Where they diverge, suggest at
  least one mechanistic explanation before concluding ``experimental
  error.''
\item
  \textbf{What would change the outcome} --- propose one modification to
  the procedure that would sharpen the measurement or extend the result
  to a new biological context, and predict what you would observe.
\item
  \textbf{One-sentence headline} --- each student composes a single
  sentence summarising the lab's take-home message, suitable for a
  tweet. Compare sentences across groups; good headlines are short,
  quantitative, and mechanistic.
\item
  \textbf{Connection back to the textbook} --- identify one section of
  \cref{sec:unit_X_population_ecology} that your data either confirmed
  or complicated. Cite the specific passage.
\end{enumerate}

\subsection*{Further Reading (Lab)}\label{further-reading-lab-31}
\addcontentsline{toc}{subsection}{Further Reading (Lab)}

\begin{itemize}
\tightlist
\item
  Revisit the parent chapter \cref{sec:unit_X_population_ecology} for
  the theoretical foundations on which this lab is built.
\item
  Look up any \textbf{bolded glossary term} introduced in the textbook
  chapter (each has a \texttt{\#gl:term-slug} link in the text) --- its
  master definition is in \breaktt{manuscript/glossary.md}.
\item
  Explore the appended \texttt{src/} module that implements the
  corresponding quantitative model (when applicable) --- referenced in
  the parent chapter's ``Bridge to Computation'' subsection.
\end{itemize}

\emph{Module footer: parent chapter
\texttt{\textbackslash{}cref\{sec:unit\_X\_population\_ecology\}}; all
numerical quantities in this lab use SI units --- see Appendix D of the
textbook for unit conversions and biological-scale reference values.}

\newpage

\section*{Lab 33 --- Community Ecology and Species
Interactions}\label{lab-33-community-ecology-and-species-interactions}
\addcontentsline{toc}{section}{Lab 33 --- Community Ecology and Species
Interactions}

\label{sec:lab_unit_X_community_ecology}

\subsection*{Learning Objectives}\label{learning-objectives-73}
\addcontentsline{toc}{subsection}{Learning Objectives}

\emph{This activity accompanies \cref{sec:unit_X_community_ecology} of
the textbook --- review that chapter before attempting the exercises
below.}

\begin{itemize}
\tightlist
\item
  Interpret a paper-based competition dataset between two species using
  Lotka-Volterra predictions
\item
  Analyse food web data to determine trophic levels and calculate energy
  transfer efficiency
\item
  Evaluate keystone predator effects using a removal/addition case study
\item
  Apply successional theory to interpret vegetation data from different
  chronosequences
\end{itemize}

\subsubsection*{Alignment and Rubric
Map}\label{alignment-and-rubric-map-35}
\addcontentsline{toc}{subsubsection}{Alignment and Rubric Map}

\begin{itemize}
\tightlist
\item
  \textbf{Outcome 1 (LO1):} Interpret the supplied evidence or model
  output for Community Ecology and Species Interactions.
\item
  \textbf{Outcome 2 (LO2):} Identify controls and comparison groups that
  make the claim testable.
\item
  \textbf{Outcome 3 (LO3):} Quantify uncertainty, boundary conditions,
  or alternative explanations before concluding.
\item
  \textbf{Outcome 4 (LO4):} Transfer the mechanism to a new biological
  case or public-facing decision.
\item
  \textbf{Chapter LO coverage:} LO1, LO2, LO3, LO4
\item
  \textbf{Rubric dimensions:} evidence; controls; uncertainty;
  mechanism; transfer. \#\# Background \{.unnumbered\}
\end{itemize}

Community ecology studies how species interact (competition, predation,
mutualism, commensalism, parasitism) and how those interactions shape
community composition. The Lotka-Volterra competition model predicts
stable coexistence primarily when intraspecific competition
\textgreater{} interspecific competition for both species (Gause's
principle: competitive exclusion eliminates one species if niches
overlap completely). Energy flows through trophic levels with
about 10\% efficiency: about 10\% of energy at level N reaches
level N+1 (ecological pyramid).

\subsection*{Paper-Based Materials}\label{paper-based-materials-35}
\addcontentsline{toc}{subsection}{Paper-Based Materials}

{\def\LTcaptype{none} % do not increment counter
\begin{longtable}[]{@{}
  >{\raggedright\arraybackslash}p{(\linewidth - 2\tabcolsep) * \real{0.3333}}
  >{\raggedright\arraybackslash}p{(\linewidth - 2\tabcolsep) * \real{0.6667}}@{}}
\toprule\noalign{}
\begin{minipage}[b]{\linewidth}\raggedright
Item
\end{minipage} & \begin{minipage}[b]{\linewidth}\raggedright
Quantity
\end{minipage} \\
\midrule\noalign{}
\endhead
\bottomrule\noalign{}
\endlastfoot
Lotka-Volterra outcome table (4 cases: stable coexistence, species 1
wins, species 2 wins, unstable equilibrium) & 1 \\
Food web diagram: temperate forest ecosystem (15 species, 3 trophic
levels) & 1 \\
Species competition dataset (N1 and N2 over 14 days in mono- and mixed
treatments) & 1 \\
Successional chronosequence and remote-sensing cards (bare substrate
-\textgreater{} pioneer -\textgreater{} scrub -\textgreater{} forest) &
1 set \\
Ecosystem-service decision matrix for pollination, carbon storage, flood
buffering, and recreation & 1 \\
Calculator & 1 \\
\end{longtable}
}

\subsection*{Paper-Based
Investigation}\label{paper-based-investigation-35}
\addcontentsline{toc}{subsection}{Paper-Based Investigation}

\textbf{Part A --- Competition Simulation/Analysis}

\begin{enumerate}
\def\labelenumi{\arabic{enumi}.}
\tightlist
\item
  Using the competition dataset, plot N1 and N2 over time in mixed
  treatment on one graph.
\item
  Predict the competitive outcome from the Lotka-Volterra equilibrium
  conditions (\texorpdfstring{\ensuremath{\alpha}}{alpha}\texorpdfstring{\textsubscript{1}}{1}\texorpdfstring{\textsubscript{2}}{2} and \texorpdfstring{\ensuremath{\alpha}}{alpha}\texorpdfstring{\textsubscript{2}}{2}\texorpdfstring{\textsubscript{1}}{1} values provided).
\item
  Compare the mixed-treatment trajectory with both monoculture
  baselines. Identify whether interspecific competition is stronger,
  weaker, or similar to intraspecific competition for each species.
\end{enumerate}

\textbf{Part B --- Food Web Analysis}

\begin{enumerate}
\def\labelenumi{\arabic{enumi}.}
\setcounter{enumi}{3}
\tightlist
\item
  From the food web diagram: identify most producers, primary consumers,
  secondary consumers, and apex predators. List most possible food
  chains containing 4 links.
\item
  If a producer level has 10,000 kJ available, calculate the energy at
  each trophic level assuming 10\% efficiency. Calculate how many humans
  (apex predators) can be supported vs if they ate at the primary
  consumer level.
\end{enumerate}

\textbf{Part C --- Keystone Predator Analysis}

\begin{enumerate}
\def\labelenumi{\arabic{enumi}.}
\setcounter{enumi}{5}
\tightlist
\item
  Read the starfish removal experiment data (Paine, 1966, summarised):
  count species richness before and after sea star (\emph{Pisaster})
  removal from an intertidal zone. Explain how removing one species
  reduced diversity.
\end{enumerate}

\textbf{Part D --- Succession Mapping}

\begin{enumerate}
\def\labelenumi{\arabic{enumi}.}
\setcounter{enumi}{6}
\tightlist
\item
  For each chronosequence or remote-sensing card: record dominant
  vegetation type, estimated time since disturbance, key abiotic
  factors, and one ecosystem service expected to change during
  succession.
\end{enumerate}

\subsection*{Data Recording}\label{data-recording-32}
\addcontentsline{toc}{subsection}{Data Recording}

Species competition (from dataset):

{\def\LTcaptype{none} % do not increment counter
\begin{longtable}[]{@{}lllll@{}}
\toprule\noalign{}
Day & N\texorpdfstring{\textsubscript{1}}{1} (monoculture) & N\texorpdfstring{\textsubscript{1}}{1} (mixed) & N\texorpdfstring{\textsubscript{2}}{2} (monoculture) & N\texorpdfstring{\textsubscript{2}}{2} (mixed) \\
\midrule\noalign{}
\endhead
\bottomrule\noalign{}
\endlastfoot
0 & & & & \\
4 & & & & \\
8 & & & & \\
14 & & & & \\
\end{longtable}
}

Food web energy pyramid:

{\def\LTcaptype{none} % do not increment counter
\begin{longtable}[]{@{}lll@{}}
\toprule\noalign{}
Trophic level & Energy (kJ) & \# Humans supported \\
\midrule\noalign{}
\endhead
\bottomrule\noalign{}
\endlastfoot
Producers & 10,000 & --- \\
Primary consumers & & \\
Secondary consumers & & \\
Apex predators & & \\
\end{longtable}
}

Species richness before sea star removal: \textbf{\emph{; After: }}

Ecosystem-service tradeoff most improved by succession: \textbf{\emph{;
tradeoff most reduced by disturbance: }}

\subsection*{Evidence and Reproducibility
Checklist}\label{evidence-and-reproducibility-checklist-35}
\addcontentsline{toc}{subsection}{Evidence and Reproducibility
Checklist}

\begin{itemize}
\tightlist
\item
  \textbf{Primary evidence goal:} Compare species-interaction outcomes
  under changing conditions.
\item
  \textbf{Data skill to practice:} Interpret abundance, interaction, or
  disturbance data from communities.
\item
  \textbf{BioSkills emphasis:} Modeling and simulation, Science and
  society, Communication and collaboration.
\item
  \textbf{Control logic:} identify at least one positive control, one
  negative control, or one baseline comparison before interpreting
  results.
\item
  \textbf{Measurement discipline:} record units, uncertainty, sample
  size, and any discarded observation with a reason.
\item
  \textbf{Mechanistic link:} connect one result directly to the parent
  chapter's big idea before writing the conclusion.
\item
  \textbf{Reproducibility check:} state one procedural detail that
  another group would need in order to reproduce the result.
\end{itemize}

\subsection*{Paper-Based Evidence
Upgrade}\label{paper-based-evidence-upgrade-35}
\addcontentsline{toc}{subsection}{Paper-Based Evidence Upgrade}

Before answering the analysis questions, annotate the paper dataset for
\textbf{Community Ecology and Species Interactions} with a
reproducibility pass:

{\def\LTcaptype{none} % do not increment counter
\begin{longtable}[]{@{}
  >{\raggedright\arraybackslash}p{(\linewidth - 2\tabcolsep) * \real{0.5000}}
  >{\raggedright\arraybackslash}p{(\linewidth - 2\tabcolsep) * \real{0.5000}}@{}}
\toprule\noalign{}
\begin{minipage}[b]{\linewidth}\raggedright
Evidence check
\end{minipage} & \begin{minipage}[b]{\linewidth}\raggedright
Student action
\end{minipage} \\
\midrule\noalign{}
\endhead
\bottomrule\noalign{}
\endlastfoot
Control logic & Mark the comparison that functions as the baseline,
negative control, or reference case. \\
Uncertainty & Circle the row, card, diagram feature, or model assumption
most likely to change the conclusion. \\
Model comparison & State whether a simpler rule, null model, or
alternative mechanism could explain the same pattern. \\
Decision threshold & Write the minimum evidence that would make you
revise the interpretation. \\
Reproducibility & Record the exact scoring rule another group would need
to reproduce your classification. \\
\end{longtable}
}

Focus note: for this community-interaction lab, label every link in the
printed network by interaction type (competition, predation, mutualism)
and by whether it was observed directly or inferred from co-occurrence,
and state how the sampling effort and disturbance history could bias
which interactions appear. Keep required work paper-based; any
material-handling or equipment version belongs only in an optional
extension.

\subsection*{Analysis Questions}\label{analysis-questions-35}
\addcontentsline{toc}{subsection}{Analysis Questions}

\begin{enumerate}
\def\labelenumi{\arabic{enumi}.}
\tightlist
\item
  In the Lotka-Volterra framework, stable coexistence requires \texorpdfstring{\ensuremath{\alpha}}{alpha}\texorpdfstring{\textsubscript{1}}{1}\texorpdfstring{\textsubscript{2}}{2}
  \textless{} K\texorpdfstring{\textsubscript{1}}{1}/K\texorpdfstring{\textsubscript{2}}{2} and \texorpdfstring{\ensuremath{\alpha}}{alpha}\texorpdfstring{\textsubscript{2}}{2}\texorpdfstring{\textsubscript{1}}{1} \textless{} K\texorpdfstring{\textsubscript{2}}{2}/K\texorpdfstring{\textsubscript{1}}{1}. What do these
  inequalities mean biologically in terms of intra- vs interspecific
  competition strength?
\item
  The pyramid of energy shows that far more humans can be supported by
  eating at lower trophic levels. Calculate the maximum number of humans
  at each trophic level from your energy data. What are three reasons
  why human food systems often involve higher trophic level consumption?
\item
  When \emph{Pisaster} (sea star) was removed, mussel populations
  exploded and crowded out other species, reducing diversity. Define
  keystone species and explain why removing an apex predator can reduce
  community diversity (trophic cascade).
\item
  Primary succession begins on bare rock (no soil) while secondary
  succession follows disturbance on existing soil. Explain how pioneer
  species (lichens, mosses) facilitate succession on bare rock --- which
  physical and chemical properties of soil do they create? Name Connell
  and Slatyer's facilitation model.
\item
  A student claims ``climax forests are the most stable and should
  generally be our conservation goal.'' Critique this statement using
  the concept of intermediate disturbance hypothesis (Connell 1978) and
  give an example of an ecosystem that requires periodic disturbance to
  maintain high diversity.
\end{enumerate}

\subsection*{Safety and Ethics Notes}\label{safety-and-ethics-notes-35}
\addcontentsline{toc}{subsection}{Safety and Ethics Notes}

No live organisms or ecological specimens are required. Optional
living-culture demonstrations should be instructor managed, disposed of
without release to waterways, and treated as enrichment rather than
required evidence.

\subsection*{Debrief and Reflection}\label{debrief-and-reflection-32}
\addcontentsline{toc}{subsection}{Debrief and Reflection}

After you finish the practical work, spend 5--10 minutes in your small
group comparing results and discussing the following prompts. Each
member should contribute at least one observation before moving to the
next prompt:

\begin{enumerate}
\def\labelenumi{\arabic{enumi}.}
\tightlist
\item
  \textbf{What did your measurements show} --- compare the group's
  results to the textbook's predictions. Where they diverge, suggest at
  least one mechanistic explanation before concluding ``experimental
  error.''
\item
  \textbf{What would change the outcome} --- propose one modification to
  the procedure that would sharpen the measurement or extend the result
  to a new biological context, and predict what you would observe.
\item
  \textbf{One-sentence headline} --- each student composes a single
  sentence summarising the lab's take-home message, suitable for a
  tweet. Compare sentences across groups; good headlines are short,
  quantitative, and mechanistic.
\item
  \textbf{Connection back to the textbook} --- identify one section of
  \cref{sec:unit_X_community_ecology} that your data either confirmed or
  complicated. Cite the specific passage.
\end{enumerate}

\subsection*{Further Reading (Lab)}\label{further-reading-lab-32}
\addcontentsline{toc}{subsection}{Further Reading (Lab)}

\begin{itemize}
\tightlist
\item
  Revisit the parent chapter \cref{sec:unit_X_community_ecology} for the
  theoretical foundations on which this lab is built.
\item
  Look up any \textbf{bolded glossary term} introduced in the textbook
  chapter (each has a \texttt{\#gl:term-slug} link in the text) --- its
  master definition is in \breaktt{manuscript/glossary.md}.
\item
  Explore the appended \texttt{src/} module that implements the
  corresponding quantitative model (when applicable) --- referenced in
  the parent chapter's ``Bridge to Computation'' subsection.
\end{itemize}

\emph{Module footer: parent chapter
\texttt{\textbackslash{}cref\{sec:unit\_X\_community\_ecology\}}; all
numerical quantities in this lab use SI units --- see Appendix D of the
textbook for unit conversions and biological-scale reference values.}

\newpage

\section*{Lab 34 --- Ecosystem Ecology and Biogeochemical
Cycles}\label{lab-34-ecosystem-ecology-and-biogeochemical-cycles}
\addcontentsline{toc}{section}{Lab 34 --- Ecosystem Ecology and
Biogeochemical Cycles}

\label{sec:lab_unit_X_ecosystem_ecology}

\subsection*{Learning Objectives}\label{learning-objectives-74}
\addcontentsline{toc}{subsection}{Learning Objectives}

\emph{This activity accompanies \cref{sec:unit_X_ecosystem_ecology} of
the textbook --- review that chapter before attempting the exercises
below.}

\begin{itemize}
\tightlist
\item
  Calculate primary productivity of an aquatic ecosystem from light-dark
  bottle datasets
\item
  Model nitrogen cycle transformations using an annotated diagram
\item
  Calculate carbon flux between reservoirs using published global carbon
  budget data
\item
  Evaluate the impact of nutrient loading on eutrophication using
  scenario data
\end{itemize}

\subsubsection*{Alignment and Rubric
Map}\label{alignment-and-rubric-map-36}
\addcontentsline{toc}{subsubsection}{Alignment and Rubric Map}

\begin{itemize}
\tightlist
\item
  \textbf{Outcome 1 (LO1):} Interpret the supplied evidence or model
  output for Ecosystem Ecology and Biogeochemical Cycles.
\item
  \textbf{Outcome 2 (LO2):} Identify controls and comparison groups that
  make the claim testable.
\item
  \textbf{Outcome 3 (LO3):} Quantify uncertainty, boundary conditions,
  or alternative explanations before concluding.
\item
  \textbf{Outcome 4 (LO4):} Transfer the mechanism to a new biological
  case or public-facing decision.
\item
  \textbf{Chapter LO coverage:} LO1, LO2, LO3, LO4
\item
  \textbf{Rubric dimensions:} evidence; controls; uncertainty;
  mechanism; transfer. \#\# Background \{.unnumbered\}
\end{itemize}

Ecosystems are characterised by flows of energy (one-way, lost as heat
at each step) and cycles of matter (recycled through biogeochemical
pathways). \textbf{Gross primary productivity (GPP)} is total
photosynthesis; \textbf{net primary productivity (NPP) = GPP \texorpdfstring{\ensuremath{-}}{-}
respiration}. The nitrogen cycle includes fixation (N\texorpdfstring{\textsubscript{2}}{2} \texorpdfstring{\ensuremath{\to}}{->} NH\texorpdfstring{\textsubscript{3}}{3} by
nitrogenase), nitrification (NH\texorpdfstring{\textsubscript{3}}{3} \texorpdfstring{\ensuremath{\to}}{->} NO\texorpdfstring{\textsubscript{2}}{2}\texorpdfstring{\textsuperscript{-}}{-} \texorpdfstring{\ensuremath{\to}}{->} NO\texorpdfstring{\textsubscript{3}}{3}\texorpdfstring{\textsuperscript{-}}{-} by \emph{Nitrosomonas},
\emph{Nitrobacter}), denitrification (NO\texorpdfstring{\textsubscript{3}}{3}\texorpdfstring{\textsuperscript{-}}{-} \texorpdfstring{\ensuremath{\to}}{->} N\texorpdfstring{\textsubscript{2}}{2} by anaerobic bacteria),
and assimilation (NH\texorpdfstring{\textsubscript{4}}{4}\texorpdfstring{\textsuperscript{+}}{+}/NO\texorpdfstring{\textsubscript{3}}{3}\texorpdfstring{\textsuperscript{-}}{-} \texorpdfstring{\ensuremath{\to}}{->} amino acids by plants).

\subsection*{Paper-Based Materials}\label{paper-based-materials-36}
\addcontentsline{toc}{subsection}{Paper-Based Materials}

{\def\LTcaptype{none} % do not increment counter
\begin{longtable}[]{@{}
  >{\raggedright\arraybackslash}p{(\linewidth - 2\tabcolsep) * \real{0.3333}}
  >{\raggedright\arraybackslash}p{(\linewidth - 2\tabcolsep) * \real{0.6667}}@{}}
\toprule\noalign{}
\begin{minipage}[b]{\linewidth}\raggedright
Item
\end{minipage} & \begin{minipage}[b]{\linewidth}\raggedright
Quantity
\end{minipage} \\
\midrule\noalign{}
\endhead
\bottomrule\noalign{}
\endlastfoot
Light-dark bottle dissolved oxygen dataset with three replicates per
treatment & 1 \\
Nitrogen cycle blank diagram (for annotation) & 1 \\
Global carbon budget data (IPCC AR6 summary values, printed) & 1 \\
Eutrophication case packet with nutrient loading, turbidity,
chlorophyll-a, and fish-kill observations & 1 \\
Ecosystem-service card set for wetlands, forests, croplands, and urban
streams & 1 \\
Graph paper or spreadsheet template & 1 \\
Calculator & 1 \\
\end{longtable}
}

\subsection*{Paper-Based
Investigation}\label{paper-based-investigation-36}
\addcontentsline{toc}{subsection}{Paper-Based Investigation}

\textbf{Part A --- Primary Productivity (Light-Dark Bottle)}

\begin{enumerate}
\def\labelenumi{\arabic{enumi}.}
\tightlist
\item
  Use the provided initial, light-bottle, and dark-bottle dissolved
  oxygen data. Confirm that each treatment has three replicates and the
  same incubation time.
\item
  Calculate mean DO and range for each bottle type.
\item
  Compare the light and dark treatments to the initial baseline, then
  flag any replicate that strongly changes the mean.
\item
  Calculate: NPP = DO\_light \texorpdfstring{\ensuremath{-}}{-} DO\_initial; Respiration = DO\_initial \texorpdfstring{\ensuremath{-}}{-}
  DO\_dark; GPP = NPP + Respiration.
\end{enumerate}

\textbf{Part B --- Nitrogen Cycle Annotation}

\begin{enumerate}
\def\labelenumi{\arabic{enumi}.}
\setcounter{enumi}{4}
\tightlist
\item
  Complete the blank nitrogen cycle diagram: label most 5 processes
  (fixation, nitrification, denitrification, assimilation,
  decomposition/ammonification). Name one key organism for each process.
\end{enumerate}

\textbf{Part C --- Carbon Budget Calculation}

\begin{enumerate}
\def\labelenumi{\arabic{enumi}.}
\setcounter{enumi}{5}
\tightlist
\item
  Using IPCC data (approximate 2022 values): GPP (terrestrial) = \texorpdfstring{\ensuremath{-}}{-}120 Gt
  C/yr (uptake); plant respiration = +60 Gt C/yr; FF emissions = +10 Gt
  C/yr; ocean uptake = \texorpdfstring{\ensuremath{-}}{-}3 Gt C/yr; land use change = +2 Gt C/yr.
  Calculate net atmospheric accumulation of CO\texorpdfstring{\textsubscript{2}}{2} per year. Compare to
  measured atmospheric growth rate (about 5 Gt C/yr).
\end{enumerate}

\textbf{Part D --- Eutrophication Experiment}

\begin{enumerate}
\def\labelenumi{\arabic{enumi}.}
\setcounter{enumi}{6}
\tightlist
\item
  Use the eutrophication case packet to compare control and
  nutrient-addition scenarios over 5 days. Interpret turbidity,
  chlorophyll-a, dissolved oxygen, and fish-kill observations as
  indicators of algal bloom risk.
\end{enumerate}

\subsection*{Data Recording}\label{data-recording-33}
\addcontentsline{toc}{subsection}{Data Recording}

{\def\LTcaptype{none} % do not increment counter
\begin{longtable}[]{@{}llll@{}}
\toprule\noalign{}
Bottle type & Initial DO (mg/L) & Final DO (mg/L) & Change \\
\midrule\noalign{}
\endhead
\bottomrule\noalign{}
\endlastfoot
Dark & & & \\
Light & & & \\
Initial (t=0) & & --- & --- \\
\end{longtable}
}

NPP = \textbf{\emph{; Respiration = }}; GPP = \_\_\_

Carbon budget: Total sources (FF + land use) = \textbf{\emph{; Total
sinks (terrestrial GPP - respiration + ocean) = }}; Net accumulation =
\_\_\_ Gt C/yr

Eutrophication indicators (day 3-5): chlorophyll-a change:
\textbf{\emph{; DO change: }}; turbidity change: \textbf{\emph{;
fish-kill risk: }}

\subsection*{Evidence and Reproducibility
Checklist}\label{evidence-and-reproducibility-checklist-36}
\addcontentsline{toc}{subsection}{Evidence and Reproducibility
Checklist}

\begin{itemize}
\tightlist
\item
  \textbf{Primary evidence goal:} Quantify ecosystem fluxes and identify
  limiting pools.
\item
  \textbf{Data skill to practice:} Trace matter and energy through food
  webs and biogeochemical cycles.
\item
  \textbf{BioSkills emphasis:} Modeling and simulation, Science and
  society, Communication and collaboration.
\item
  \textbf{Control logic:} identify at least one positive control, one
  negative control, or one baseline comparison before interpreting
  results.
\item
  \textbf{Measurement discipline:} record units, uncertainty, sample
  size, and any discarded observation with a reason.
\item
  \textbf{Mechanistic link:} connect one result directly to the parent
  chapter's big idea before writing the conclusion.
\item
  \textbf{Reproducibility check:} state one procedural detail that
  another group would need in order to reproduce the result.
\end{itemize}

\subsection*{Paper-Based Evidence
Upgrade}\label{paper-based-evidence-upgrade-36}
\addcontentsline{toc}{subsection}{Paper-Based Evidence Upgrade}

Before answering the analysis questions, annotate the paper dataset for
\textbf{Ecosystem Ecology and Biogeochemical Cycles} with a
reproducibility pass:

{\def\LTcaptype{none} % do not increment counter
\begin{longtable}[]{@{}
  >{\raggedright\arraybackslash}p{(\linewidth - 2\tabcolsep) * \real{0.5000}}
  >{\raggedright\arraybackslash}p{(\linewidth - 2\tabcolsep) * \real{0.5000}}@{}}
\toprule\noalign{}
\begin{minipage}[b]{\linewidth}\raggedright
Evidence check
\end{minipage} & \begin{minipage}[b]{\linewidth}\raggedright
Student action
\end{minipage} \\
\midrule\noalign{}
\endhead
\bottomrule\noalign{}
\endlastfoot
Control logic & Mark the comparison that functions as the baseline,
negative control, or reference case. \\
Uncertainty & Circle the row, card, diagram feature, or model assumption
most likely to change the conclusion. \\
Model comparison & State whether a simpler rule, null model, or
alternative mechanism could explain the same pattern. \\
Decision threshold & Write the minimum evidence that would make you
revise the interpretation. \\
Reproducibility & Record the exact scoring rule another group would need
to reproduce your classification. \\
\end{longtable}
}

Focus note: for this biogeochemical lab, keep stocks and fluxes in
distinct columns, check that the printed budget balances within its
stated uncertainty, and identify the system boundary and residence time
implied before comparing one pool or pathway against another. Keep
required work paper-based; any material-handling or equipment version
belongs only in an optional extension.

\subsection*{Analysis Questions}\label{analysis-questions-36}
\addcontentsline{toc}{subsection}{Analysis Questions}

\begin{enumerate}
\def\labelenumi{\arabic{enumi}.}
\tightlist
\item
  The dark bottle showed decreased DO while the light bottle showed
  increased DO. Explain why: which process dominates in the light
  bottle, and which process in the dark bottle? Why is O\texorpdfstring{\textsubscript{2}}{2} a useful proxy
  for carbon fixation here?
\item
  Your carbon budget calculation: does the simple sum of sources and
  sinks balance with the measured atmospheric accumulation
  (about 5 Gt C/yr)? What are the largest sources of
  uncertainty in the global carbon budget?
\item
  Nitrogen fixation by \emph{Rhizobium} in legume root nodules requires
  enormous amounts of ATP and anaerobic conditions (nitrogenase is
  inactivated by O\texorpdfstring{\textsubscript{2}}{2}). Explain why this process is energetically costly
  and how the leghaemoglobin protein maintains low O\texorpdfstring{\textsubscript{2}}{2} in the nodule
  while supporting aerobic bacterial respiration.
\item
  In the eutrophication case packet, the nutrient-addition scenario
  shows increased chlorophyll-a followed by bottom-water hypoxia. Trace
  the ecological mechanism from nutrient addition -\textgreater{} algal
  bloom -\textgreater{} bacterial decomposition -\textgreater{} hypoxia
  -\textgreater{} fish kill.
\item
  Wetlands are sometimes called ``kidneys of the landscape.'' Explain
  their role in nitrogen cycling (using specific N-cycle steps),
  phosphorus immobilisation, and carbon sequestration. Why is wetland
  drainage for agriculture a double ecological concern?
\end{enumerate}

\subsection*{Safety and Ethics Notes}\label{safety-and-ethics-notes-36}
\addcontentsline{toc}{subsection}{Safety and Ethics Notes}

No pond water, fertiliser, bottles, probes, or biological specimens are
required. Optional demonstrations using water samples or nutrients
should follow local disposal rules and should rarely release enriched
water into drains or outdoor waterways.

\subsection*{Debrief and Reflection}\label{debrief-and-reflection-33}
\addcontentsline{toc}{subsection}{Debrief and Reflection}

After you finish the practical work, spend 5--10 minutes in your small
group comparing results and discussing the following prompts. Each
member should contribute at least one observation before moving to the
next prompt:

\begin{enumerate}
\def\labelenumi{\arabic{enumi}.}
\tightlist
\item
  \textbf{What did your measurements show} --- compare the group's
  results to the textbook's predictions. Where they diverge, suggest at
  least one mechanistic explanation before concluding ``experimental
  error.''
\item
  \textbf{What would change the outcome} --- propose one modification to
  the procedure that would sharpen the measurement or extend the result
  to a new biological context, and predict what you would observe.
\item
  \textbf{One-sentence headline} --- each student composes a single
  sentence summarising the lab's take-home message, suitable for a
  tweet. Compare sentences across groups; good headlines are short,
  quantitative, and mechanistic.
\item
  \textbf{Connection back to the textbook} --- identify one section of
  \cref{sec:unit_X_ecosystem_ecology} that your data either confirmed or
  complicated. Cite the specific passage.
\end{enumerate}

\subsection*{Further Reading (Lab)}\label{further-reading-lab-33}
\addcontentsline{toc}{subsection}{Further Reading (Lab)}

\begin{itemize}
\tightlist
\item
  Revisit the parent chapter \cref{sec:unit_X_ecosystem_ecology} for the
  theoretical foundations on which this lab is built.
\item
  Look up any \textbf{bolded glossary term} introduced in the textbook
  chapter (each has a \texttt{\#gl:term-slug} link in the text) --- its
  master definition is in \breaktt{manuscript/glossary.md}.
\item
  Explore the appended \texttt{src/} module that implements the
  corresponding quantitative model (when applicable) --- referenced in
  the parent chapter's ``Bridge to Computation'' subsection.
\end{itemize}

\emph{Module footer: parent chapter
\texttt{\textbackslash{}cref\{sec:unit\_X\_ecosystem\_ecology\}}; all
numerical quantities in this lab use SI units --- see Appendix D of the
textbook for unit conversions and biological-scale reference values.}

\newpage

\section*{Lab 35 --- Biomes and Conservation
Biology}\label{lab-35-biomes-and-conservation-biology}
\addcontentsline{toc}{section}{Lab 35 --- Biomes and Conservation
Biology}

\label{sec:lab_unit_X_biomes_and_conservation}

\subsection*{Learning Objectives}\label{learning-objectives-75}
\addcontentsline{toc}{subsection}{Learning Objectives}

\emph{This activity accompanies
\cref{sec:unit_X_biomes_and_conservation} of the textbook --- review
that chapter before attempting the exercises below.}

\begin{itemize}
\tightlist
\item
  Classify biomes from climate data (mean annual temperature,
  precipitation, Whittaker diagram)
\item
  Calculate species-area relationship parameters and predict extinction
  risk from habitat loss
\item
  Evaluate minimum viable population (MVP) and reserve design principles
  for a case species
\item
  Apply IUCN Red List criteria to classify conservation status of a
  provided species profile
\end{itemize}

\subsubsection*{Alignment and Rubric
Map}\label{alignment-and-rubric-map-37}
\addcontentsline{toc}{subsubsection}{Alignment and Rubric Map}

\begin{itemize}
\tightlist
\item
  \textbf{Outcome 1 (LO1):} Interpret the supplied evidence or model
  output for Biomes and Conservation Biology.
\item
  \textbf{Outcome 2 (LO2):} Identify controls and comparison groups that
  make the claim testable.
\item
  \textbf{Outcome 3 (LO3):} Quantify uncertainty, boundary conditions,
  or alternative explanations before concluding.
\item
  \textbf{Outcome 4 (LO4):} Transfer the mechanism to a new biological
  case or public-facing decision.
\item
  \textbf{Chapter LO coverage:} LO1, LO2, LO3, LO4
\item
  \textbf{Rubric dimensions:} evidence; controls; uncertainty;
  mechanism; transfer. \#\# Background \{.unnumbered\}
\end{itemize}

Biomes are large-scale communities determined primarily by climate
(temperature and precipitation). The species-area relationship
\textbf{\(S = cA^z\)} (log-linear: log S = log c + z\texorpdfstring{\ensuremath{\cdot}}{.}log A) predicts
that decreasing habitat area reduces species richness; the exponent z
(typically 0.2--0.35 for oceanic islands) can be used to predict
extinction rates from deforestation. Minimum viable population (MVP)
analyses determine the population size needed for a defined probability
of persistence over a defined time frame. IUCN Red List criteria assign
species to categories (LC, NT, VU, EN, CR, EW, EX) based on population
decline rates, geographic range, and quantitative extinction risk.

\subsection*{Paper-Based Materials}\label{paper-based-materials-37}
\addcontentsline{toc}{subsection}{Paper-Based Materials}

{\def\LTcaptype{none} % do not increment counter
\begin{longtable}[]{@{}
  >{\raggedright\arraybackslash}p{(\linewidth - 2\tabcolsep) * \real{0.3333}}
  >{\raggedright\arraybackslash}p{(\linewidth - 2\tabcolsep) * \real{0.6667}}@{}}
\toprule\noalign{}
\begin{minipage}[b]{\linewidth}\raggedright
Item
\end{minipage} & \begin{minipage}[b]{\linewidth}\raggedright
Quantity
\end{minipage} \\
\midrule\noalign{}
\endhead
\bottomrule\noalign{}
\endlastfoot
Whittaker biome diagram (printed with temperature and precipitation
axes) & 1 \\
Climate data cards (10 locations: MAT and MAP values) & 1 set \\
Species-area relationship dataset (6 habitat fragments: area and species
count) & 1 \\
Log-log graph paper & 1 \\
IUCN criteria worksheet (table of criteria A--E) & 1 \\
Species profile card: fictional vaquita-like cetacean (population data,
trend, range) & 1 \\
Reserve-design decision matrix with habitat area, corridor cost, carbon
storage, and community co-benefits & 1 \\
Climate-shift scenario cards for biome boundaries and conservation
priorities & 1 set \\
Coral bleaching degree-heating-week cards with recovery, mortality, and
symbiont-shuffling scenarios & 1 set \\
IPBES/food-security tradeoff card: agroecology, yield, pollinator
habitat, waste reduction, and equity indicators & 1 \\
Calculator & 1 \\
\end{longtable}
}

\subsection*{Paper-Based
Investigation}\label{paper-based-investigation-37}
\addcontentsline{toc}{subsection}{Paper-Based Investigation}

\textbf{Part A --- Biome Classification}

\begin{enumerate}
\def\labelenumi{\arabic{enumi}.}
\tightlist
\item
  For each of 10 climate data cards (MAT in °C, MAP in mm), plot the
  location on the Whittaker diagram. Identify the biome for each
  location.
\end{enumerate}

\textbf{Part B --- Species-Area Relationship}

\begin{enumerate}
\def\labelenumi{\arabic{enumi}.}
\setcounter{enumi}{1}
\tightlist
\item
  Plot log(S) vs log(A) for the 6 habitat fragments on log-log paper.
  Fit a line; calculate the slope (z exponent) and the y-intercept (log
  c).
\item
  Using \(S = cA^z\), predict species richness in a fragment one-tenth
  the area of the largest.
\item
  If a deforestation event reduces a forest from 10,000 ha to 100 ha,
  predict the percentage of species lost using the same equation.
\end{enumerate}

\textbf{Part C --- MVP Analysis}

\begin{enumerate}
\def\labelenumi{\arabic{enumi}.}
\setcounter{enumi}{4}
\tightlist
\item
  Given MVP data for the provided species: 50/500 rule (N\_e
  \textgreater{} 50 for short-term genetic viability; N\_e
  \textgreater{} 500 for long-term). Current population: 35 individuals.
  Is the population below the short-term MVP? Design a 10-year recovery
  plan with specific targets.
\end{enumerate}

\textbf{Part D --- IUCN Classification}

\begin{enumerate}
\def\labelenumi{\arabic{enumi}.}
\setcounter{enumi}{5}
\tightlist
\item
  Apply IUCN criteria to the species profile (population decline 80\% in
  10 years; \textless{} 250 mature individuals; geographic range
  \textless{} 100 km\texorpdfstring{\textsuperscript{2}}{2}). Determine IUCN status.
\item
  Use the reserve-design matrix to compare a single large reserve,
  several small reserves, and a connected-corridor plan. Score each
  option for extinction risk, climate refugia, ecosystem services, and
  governance tradeoffs.
\end{enumerate}

\textbf{Part E --- Coral and Food-System Decision Cards}

\begin{enumerate}
\def\labelenumi{\arabic{enumi}.}
\setcounter{enumi}{7}
\tightlist
\item
  Classify coral scenarios by degree-heating-week exposure, recovery
  interval, and local stressors. Decide whether the evidence supports
  likely recovery, repeated bleaching risk, or high mortality risk.
\item
  Use the symbiont-response cards to distinguish symbiont shuffling,
  assisted gene flow, selective breeding, and microbiome conditioning.
  State one benefit, one tradeoff, and one missing validation step for
  each.
\item
  Use the IPBES/food-security card to compare two conservation plans:
  one maximizing protected area alone and one combining habitat
  protection with agroecological transition, pollinator corridors, and
  food-waste reduction. Score biodiversity, nutrition, equity, and
  feasibility separately.
\end{enumerate}

\subsection*{Data Recording}\label{data-recording-34}
\addcontentsline{toc}{subsection}{Data Recording}

Biome classifications:

{\def\LTcaptype{none} % do not increment counter
\begin{longtable}[]{@{}llll@{}}
\toprule\noalign{}
Location & MAT (°C) & MAP (mm) & Biome \\
\midrule\noalign{}
\endhead
\bottomrule\noalign{}
\endlastfoot
1 & 28 & 2,400 & \\
2 & \texorpdfstring{\ensuremath{-}}{-}5 & 300 & \\
3 & 15 & 600 & \\
4 & 22 & 100 & \\
5 & 0 & 50 & \\
\end{longtable}
}

(continue for 10 locations)

Species-area relationship: z = \textbf{\emph{; c = }}; predicted S at
1/10 area = \textbf{\emph{; \% species lost = }}\%

MVP assessment: current N = 35; MVP\_short = 50; MVP\_long = 500 Status:
BELOW/ABOVE short-term MVP; IUCN status: \_\_\_

Reserve-design score: extinction risk = \textbf{\emph{; climate refugia
= }}; ecosystem services = \textbf{\emph{; governance risk = }} Coral
risk class: \_\_\_\_\_\_\_; assisted-evolution option: \_\_\_\_\_\_\_;
food-system tradeoff score: biodiversity \_\_\_ / nutrition \_\_\_ /
equity \_\_\_ / feasibility \_\_\_

\subsection*{Evidence and Reproducibility
Checklist}\label{evidence-and-reproducibility-checklist-37}
\addcontentsline{toc}{subsection}{Evidence and Reproducibility
Checklist}

\begin{itemize}
\tightlist
\item
  \textbf{Primary evidence goal:} Compare biome patterns and
  conservation tradeoffs.
\item
  \textbf{Data skill to practice:} Use maps, trend data, and threat
  categories to justify conservation priorities.
\item
  \textbf{BioSkills emphasis:} Modeling and simulation, Science and
  society, Communication and collaboration.
\item
  \textbf{Control logic:} identify at least one positive control, one
  negative control, or one baseline comparison before interpreting
  results.
\item
  \textbf{Measurement discipline:} record units, uncertainty, sample
  size, and any discarded observation with a reason.
\item
  \textbf{Mechanistic link:} connect one result directly to the parent
  chapter's big idea before writing the conclusion.
\item
  \textbf{Reproducibility check:} state one procedural detail that
  another group would need in order to reproduce the result.
\end{itemize}

\subsection*{Paper-Based Evidence
Upgrade}\label{paper-based-evidence-upgrade-37}
\addcontentsline{toc}{subsection}{Paper-Based Evidence Upgrade}

Before answering the analysis questions, annotate the paper dataset for
\textbf{Biomes and Conservation Biology} with a reproducibility pass:

{\def\LTcaptype{none} % do not increment counter
\begin{longtable}[]{@{}
  >{\raggedright\arraybackslash}p{(\linewidth - 2\tabcolsep) * \real{0.5000}}
  >{\raggedright\arraybackslash}p{(\linewidth - 2\tabcolsep) * \real{0.5000}}@{}}
\toprule\noalign{}
\begin{minipage}[b]{\linewidth}\raggedright
Evidence check
\end{minipage} & \begin{minipage}[b]{\linewidth}\raggedright
Student action
\end{minipage} \\
\midrule\noalign{}
\endhead
\bottomrule\noalign{}
\endlastfoot
Control logic & Mark the comparison that functions as the baseline,
negative control, or reference case. \\
Uncertainty & Circle the row, card, diagram feature, or model assumption
most likely to change the conclusion. \\
Model comparison & State whether a simpler rule, null model, or
alternative mechanism could explain the same pattern. \\
Decision threshold & Write the minimum evidence that would make you
revise the interpretation. \\
Reproducibility & Record the exact scoring rule another group would need
to reproduce your classification. \\
\end{longtable}
}

Focus note: for this conservation lab, separate the ecological signal in
the printed dataset (a biodiversity index, range map, or protected-area
comparison) from the value judgement and the feasibility constraint, and
name explicitly which counterfactual or unprotected reference the
decision is being measured against. Keep required work paper-based; any
material-handling or equipment version belongs only in an optional
extension.

\subsection*{Analysis Questions}\label{analysis-questions-37}
\addcontentsline{toc}{subsection}{Analysis Questions}

\begin{enumerate}
\def\labelenumi{\arabic{enumi}.}
\tightlist
\item
  Two locations had the same MAT but very different biomes because of
  different MAP values. What does this demonstrate about the relative
  importance of temperature vs precipitation in biome determination?
  Name an example pair from the Whittaker diagram where this contrast is
  striking.
\item
  Your log-log plot of the species-area relationship should be linear.
  What does the slope (z) represent ecologically? Why do oceanic islands
  have higher z values (about 0.35) than mainland habitat
  patches (about 0.12)?
\item
  Predicted species loss from deforestation assumes equilibrium. Explain
  the concept of ``extinction debt'' --- why the full extinction toll of
  habitat reduction is primarily realised decades after the habitat
  loss, and what this implies for urgency in conservation action.
\item
  The 50/500 rule gives threshold population sizes for genetic
  viability. Explain: (a) why populations below N\_e = 50 suffer
  inbreeding depression; (b) why populations below N\_e = 500 lose the
  variation needed to adapt to future environmental change. What is the
  difference between census population size (N\_c) and effective
  population size (N\_e)?
\item
  Design a conservation reserve for a large mammal (e.g., tiger, Amur
  leopard) with a home range of 200 km\texorpdfstring{\textsuperscript{2}}{2} that requires a minimum viable
  population of 100 individuals. Calculate the minimum reserve area,
  explain whether a single large reserve (SLOSS debate: Single Large or
  Several Small) is preferable in this case, and identify two human
  activities outside the reserve boundary that would most threaten
  population persistence.
\item
  A reef has high DHW exposure but a recovery window long enough for
  partial symbiont reshuffling, while a nearby watershed continues to
  deliver nutrient pollution. Which intervention should be treated as
  climate adaptation, which as local-stressor reduction, and why would
  assisted evolution alone be insufficient?
\item
  A food-security policy increases short-term yield by removing
  hedgerows and reducing crop diversity. Use the IPBES/food-system
  tradeoff card to evaluate why yield, pollination, soil stability,
  nutrition, and equity must be scored separately rather than collapsed
  into one ``more food'' metric.
\end{enumerate}

\subsection*{Safety and Ethics Notes}\label{safety-and-ethics-notes-37}
\addcontentsline{toc}{subsection}{Safety and Ethics Notes}

No hazardous materials. Conservation discussions involving indigenous
land rights and protected area design should be approached with cultural
sensitivity and recognition of land sovereignty principles. IUCN
assessments require reliable species data --- discuss uncertainty and
the precautionary principle.

\subsection*{Debrief and Reflection}\label{debrief-and-reflection-34}
\addcontentsline{toc}{subsection}{Debrief and Reflection}

After you finish the practical work, spend 5--10 minutes in your small
group comparing results and discussing the following prompts. Each
member should contribute at least one observation before moving to the
next prompt:

\begin{enumerate}
\def\labelenumi{\arabic{enumi}.}
\tightlist
\item
  \textbf{What did your measurements show} --- compare the group's
  results to the textbook's predictions. Where they diverge, suggest at
  least one mechanistic explanation before concluding ``experimental
  error.''
\item
  \textbf{What would change the outcome} --- propose one modification to
  the procedure that would sharpen the measurement or extend the result
  to a new biological context, and predict what you would observe.
\item
  \textbf{One-sentence headline} --- each student composes a single
  sentence summarising the lab's take-home message, suitable for a
  tweet. Compare sentences across groups; good headlines are short,
  quantitative, and mechanistic.
\item
  \textbf{Connection back to the textbook} --- identify one section of
  \cref{sec:unit_X_biomes_and_conservation} that your data either
  confirmed or complicated. Cite the specific passage.
\end{enumerate}

\subsection*{Further Reading (Lab)}\label{further-reading-lab-34}
\addcontentsline{toc}{subsection}{Further Reading (Lab)}

\begin{itemize}
\tightlist
\item
  Revisit the parent chapter \cref{sec:unit_X_biomes_and_conservation}
  for the theoretical foundations on which this lab is built.
\item
  Look up any \textbf{bolded glossary term} introduced in the textbook
  chapter (each has a \texttt{\#gl:term-slug} link in the text) --- its
  master definition is in \breaktt{manuscript/glossary.md}.
\item
  Explore the appended \texttt{src/} module that implements the
  corresponding quantitative model (when applicable) --- referenced in
  the parent chapter's ``Bridge to Computation'' subsection.
\end{itemize}

\emph{Module footer: parent chapter
\texttt{\textbackslash{}cref\{sec:unit\_X\_biomes\_and\_conservation\}};
all numerical quantities in this lab use SI units --- see Appendix D of
the textbook for unit conversions and biological-scale reference
values.}

\newpage

\section*{Questions --- Ch 0.1: Systems Science and the Logic of
Emergence}\label{questions-ch-0.1-systems-science-and-the-logic-of-emergence}
\addcontentsline{toc}{section}{Questions --- Ch 0.1: Systems Science and
the Logic of Emergence}

\label{sec:q_unit_0_systems_science}

\emph{This activity accompanies \cref{sec:unit_0_systems_science} of the
textbook --- review that chapter before attempting the exercises below.}

\begin{center}\rule{0.5\linewidth}{0.5pt}\end{center}

\subsection*{Recall Questions (1 mark
each)}\label{recall-questions-1-mark-each}
\addcontentsline{toc}{subsection}{Recall Questions (1 mark each)}

\subsection*{Instructor Use and Coverage
Notes}\label{instructor-use-and-coverage-notes}
\addcontentsline{toc}{subsection}{Instructor Use and Coverage Notes}

\begin{itemize}
\tightlist
\item
  \textbf{Coverage target:} Given a biological scenario, identify the
  system boundary, feedback loop, and missing measurement.
\item
  \textbf{Model/data emphasis:} Box-and-arrow causal models with
  explicit inputs, outputs, and feedback signs.
\item
  \textbf{Assessment alignment:} Questions and Methods, Representing and
  Describing Data, Argumentation.
\item
  \textbf{Misconception probe:} A system is not just a list of parts;
  the interactions are part of the explanation.
\item
  \textbf{Transfer product:} Apply the same feedback map to a cell,
  organism, and ecosystem, then name what changes at each scale.
\item
  \textbf{Grading focus:} award full credit for mechanism, evidence,
  boundary conditions, and units when a calculation is required.
\item
  \textbf{Suggested use:} draw one recall item, one application item,
  and one synthesis item when building a short quiz from this bank.
\end{itemize}

\begin{enumerate}
\def\labelenumi{\arabic{enumi}.}
\tightlist
\item
  Define `systems science' in your own words and provide a biological
  example.
\end{enumerate}

\begin{enumerate}
\def\labelenumi{\arabic{enumi}.}
\setcounter{enumi}{1}
\tightlist
\item
  What distinguishes a complex adaptive system from a simple system?
\end{enumerate}

\begin{enumerate}
\def\labelenumi{\arabic{enumi}.}
\setcounter{enumi}{2}
\tightlist
\item
  What is the free energy principle? What does free energy measure?
\end{enumerate}

\begin{enumerate}
\def\labelenumi{\arabic{enumi}.}
\setcounter{enumi}{3}
\tightlist
\item
  Describe one biological example of negative feedback and one of
  positive feedback.
\end{enumerate}

\begin{enumerate}
\def\labelenumi{\arabic{enumi}.}
\setcounter{enumi}{4}
\tightlist
\item
  What is the difference between homeostasis and allostasis?
\end{enumerate}

\begin{center}\rule{0.5\linewidth}{0.5pt}\end{center}

\subsection*{Application Questions (2 marks
each)}\label{application-questions-2-marks-each}
\addcontentsline{toc}{subsection}{Application Questions (2 marks each)}

\begin{enumerate}
\def\labelenumi{\arabic{enumi}.}
\setcounter{enumi}{5}
\tightlist
\item
  A drug blocks a key kinase in a signalling cascade, but the tumour
  eventually regrows. Explain this outcome using the concept of feedback
  compensation.
\end{enumerate}

\begin{enumerate}
\def\labelenumi{\arabic{enumi}.}
\setcounter{enumi}{6}
\tightlist
\item
  Explain why the lac operon switch is described as bistable. What
  advantage does bistability provide over a graded (linear) response?
\end{enumerate}

\begin{enumerate}
\def\labelenumi{\arabic{enumi}.}
\setcounter{enumi}{7}
\tightlist
\item
  How does the concept of attractors in phase space relate to cell fate?
  Give an example.
\end{enumerate}

\begin{enumerate}
\def\labelenumi{\arabic{enumi}.}
\setcounter{enumi}{8}
\tightlist
\item
  Apply the active inference framework to explain how a fever might be
  adaptive rather than merely a pathological response.
\end{enumerate}

\begin{enumerate}
\def\labelenumi{\arabic{enumi}.}
\setcounter{enumi}{9}
\tightlist
\item
  In evolutionary terms, what is niche construction, and how does it
  relate to the free energy principle?
\end{enumerate}

\begin{center}\rule{0.5\linewidth}{0.5pt}\end{center}

\subsection*{Synthesis Questions (4 marks
each)}\label{synthesis-questions-4-marks-each}
\addcontentsline{toc}{subsection}{Synthesis Questions (4 marks each)}

\begin{enumerate}
\def\labelenumi{\arabic{enumi}.}
\setcounter{enumi}{10}
\tightlist
\item
  Compare and contrast the immune system and an evolving tumour as
  complex adaptive systems. In each case, identify (a) the agents, (b)
  the selection pressures, (c) an emergent property, and (d) a critical
  vulnerability.
\end{enumerate}

\begin{enumerate}
\def\labelenumi{\arabic{enumi}.}
\setcounter{enumi}{11}
\tightlist
\item
  A patient with chronic anxiety has persistent overactivation of the
  sympathetic nervous system even in the absence of objective threat.
  Using predictive coding and active inference, construct a mechanistic
  account of how prior beliefs about threat might drive this pattern.
  How might therapeutic interventions (CBT, interoceptive exposure)
  alter the generative model?
\end{enumerate}

\begin{enumerate}
\def\labelenumi{\arabic{enumi}.}
\setcounter{enumi}{12}
\tightlist
\item
  Describe how ecosystem collapse (e.g., coral bleaching, fishery
  collapse) can be framed as a phase transition in a complex adaptive
  system. What warning signals might precede the transition? Why is
  recovery often much harder than the original decline?
\end{enumerate}

\subsection*{Data and Model Interpretation Questions (3 marks
each)}\label{data-and-model-interpretation-questions-3-marks-each}
\addcontentsline{toc}{subsection}{Data and Model Interpretation
Questions (3 marks each)}

\begin{enumerate}
\def\labelenumi{\arabic{enumi}.}
\setcounter{enumi}{13}
\tightlist
\item
  A homeostatic variable is measured at 98, 101, 105, 112, and 130 units
  after a perturbation. Identify whether the response appears damped,
  amplified, or unstable, and name one missing measurement needed before
  drawing a systems-level conclusion.
\end{enumerate}

\begin{enumerate}
\def\labelenumi{\arabic{enumi}.}
\setcounter{enumi}{14}
\tightlist
\item
  Draw a signed feedback diagram for blood glucose regulation that
  includes glucose, insulin, glucagon, liver glycogen breakdown, and
  muscle glucose uptake. Which links are negative feedback links?
\end{enumerate}

\begin{enumerate}
\def\labelenumi{\arabic{enumi}.}
\setcounter{enumi}{15}
\tightlist
\item
  In a predator-prey system, prey abundance rises first, predator
  abundance rises after a delay, and prey then falls. Explain why the
  delay matters for oscillation rather than immediate stabilization.
\end{enumerate}

\begin{enumerate}
\def\labelenumi{\arabic{enumi}.}
\setcounter{enumi}{16}
\tightlist
\item
  A pathway has two parallel routes from receptor activation to
  transcription. Blocking one route reduces output by 20\%. What systems
  property explains this robustness, and what experiment would test it?
\end{enumerate}

\begin{enumerate}
\def\labelenumi{\arabic{enumi}.}
\setcounter{enumi}{17}
\tightlist
\item
  A network diagram has many nodes with two links and three nodes with
  hundreds of links. What kind of network does this suggest, and why is
  it robust to random node loss but vulnerable to targeted hub loss?
\end{enumerate}

\begin{enumerate}
\def\labelenumi{\arabic{enumi}.}
\setcounter{enumi}{18}
\tightlist
\item
  A cell-fate model has two stable attractors, A and B. A weak transient
  signal moves the cell partway toward B, but it returns to A; a
  stronger signal switches it to B permanently. Explain this in terms of
  basins of attraction.
\end{enumerate}

\begin{enumerate}
\def\labelenumi{\arabic{enumi}.}
\setcounter{enumi}{19}
\tightlist
\item
  A researcher reports that deleting one enzyme in a pathway has no
  phenotype, so the enzyme is unimportant. Critique this conclusion
  using compensation and context dependence.
\end{enumerate}

\begin{enumerate}
\def\labelenumi{\arabic{enumi}.}
\setcounter{enumi}{20}
\tightlist
\item
  A negative feedback loop has a sensor that saturates at high input
  values. Predict what happens when the input exceeds the sensor's
  dynamic range.
\end{enumerate}

\begin{enumerate}
\def\labelenumi{\arabic{enumi}.}
\setcounter{enumi}{21}
\tightlist
\item
  Explain how modularity can make a system both more evolvable and more
  fragile.
\end{enumerate}

\begin{enumerate}
\def\labelenumi{\arabic{enumi}.}
\setcounter{enumi}{22}
\tightlist
\item
  Compare a thermostat, blood calcium regulation, and quorum sensing.
  Which parts of the analogy are useful, and where does the analogy
  break down?
\end{enumerate}

\subsection*{Evaluation and Design Questions (4 marks
each)}\label{evaluation-and-design-questions-4-marks-each}
\addcontentsline{toc}{subsection}{Evaluation and Design Questions (4
marks each)}

\begin{enumerate}
\def\labelenumi{\arabic{enumi}.}
\setcounter{enumi}{23}
\tightlist
\item
  Design a minimal experiment to distinguish direct causation from
  correlation in a proposed feedback loop between a hormone and a
  behaviour.
\end{enumerate}

\begin{enumerate}
\def\labelenumi{\arabic{enumi}.}
\setcounter{enumi}{24}
\tightlist
\item
  A model of fever treats body temperature as a fixed set point.
  Evaluate why this model fails during infection and propose a better
  systems description.
\end{enumerate}

\begin{enumerate}
\def\labelenumi{\arabic{enumi}.}
\setcounter{enumi}{25}
\tightlist
\item
  Propose two early-warning indicators that a lake ecosystem is
  approaching a tipping point. Explain why each indicator is
  mechanistically plausible.
\end{enumerate}

\begin{enumerate}
\def\labelenumi{\arabic{enumi}.}
\setcounter{enumi}{26}
\tightlist
\item
  A clinician wants a single biomarker for sepsis severity. Use systems
  thinking to argue for or against relying on one marker.
\end{enumerate}

\begin{enumerate}
\def\labelenumi{\arabic{enumi}.}
\setcounter{enumi}{27}
\tightlist
\item
  Build a modular explanation of wound healing with at least three
  modules and one feedback link between modules.
\end{enumerate}

\begin{enumerate}
\def\labelenumi{\arabic{enumi}.}
\setcounter{enumi}{28}
\tightlist
\item
  A student says, ``If we know every molecule in a cell, we know the
  cell.'' Respond using emergence, scale, and model choice.
\end{enumerate}

\begin{enumerate}
\def\labelenumi{\arabic{enumi}.}
\setcounter{enumi}{29}
\tightlist
\item
  Choose one biological system from another unit of the textbook and
  outline how you would model it as a system with state variables,
  inputs, outputs, and feedback.
\end{enumerate}

\begin{center}\rule{0.5\linewidth}{0.5pt}\end{center}

\emph{Module: \texttt{src/biology/} (general).}

\newpage

\section*{Questions --- Ch 0.2: Complex Adaptive
Systems}\label{questions-ch-0.2-complex-adaptive-systems}
\addcontentsline{toc}{section}{Questions --- Ch 0.2: Complex Adaptive
Systems}

\label{sec:q_unit_0_complex_adaptive_systems}

\emph{This activity accompanies
\cref{sec:unit_0_complex_adaptive_systems} of the textbook --- review
that chapter before attempting the exercises below.}

\begin{center}\rule{0.5\linewidth}{0.5pt}\end{center}

\subsection*{Recall Questions (1 mark
each)}\label{recall-questions-1-mark-each-1}
\addcontentsline{toc}{subsection}{Recall Questions (1 mark each)}

\subsection*{Instructor Use and Coverage
Notes}\label{instructor-use-and-coverage-notes-1}
\addcontentsline{toc}{subsection}{Instructor Use and Coverage Notes}

\begin{itemize}
\tightlist
\item
  \textbf{Coverage target:} Explain an observed population pattern from
  individual-level rules and identify one falsifiable prediction.
\item
  \textbf{Model/data emphasis:} Agent-rule and scaling arguments for
  emergent biological patterns.
\item
  \textbf{Assessment alignment:} Questions and Methods, Representing and
  Describing Data, Argumentation.
\item
  \textbf{Misconception probe:} Emergence is not mysterious; it is a
  reproducible consequence of interactions plus constraints.
\item
  \textbf{Transfer product:} Compare flocking, immune activation, and
  microbial biofilms as adaptive systems.
\item
  \textbf{Grading focus:} award full credit for mechanism, evidence,
  boundary conditions, and units when a calculation is required.
\item
  \textbf{Suggested use:} draw one recall item, one application item,
  and one synthesis item when building a short quiz from this bank.
\end{itemize}

\begin{enumerate}
\def\labelenumi{\arabic{enumi}.}
\tightlist
\item
  Define `complex adaptive systems' in your own words and provide a
  biological example.
\end{enumerate}

\begin{enumerate}
\def\labelenumi{\arabic{enumi}.}
\setcounter{enumi}{1}
\tightlist
\item
  What distinguishes a complex adaptive system from a simple system?
\end{enumerate}

\begin{enumerate}
\def\labelenumi{\arabic{enumi}.}
\setcounter{enumi}{2}
\tightlist
\item
  What is the free energy principle? What does free energy measure?
\end{enumerate}

\begin{enumerate}
\def\labelenumi{\arabic{enumi}.}
\setcounter{enumi}{3}
\tightlist
\item
  Describe one biological example of negative feedback and one of
  positive feedback.
\end{enumerate}

\begin{enumerate}
\def\labelenumi{\arabic{enumi}.}
\setcounter{enumi}{4}
\tightlist
\item
  What is the difference between homeostasis and allostasis?
\end{enumerate}

\begin{center}\rule{0.5\linewidth}{0.5pt}\end{center}

\subsection*{Application Questions (2 marks
each)}\label{application-questions-2-marks-each-1}
\addcontentsline{toc}{subsection}{Application Questions (2 marks each)}

\begin{enumerate}
\def\labelenumi{\arabic{enumi}.}
\setcounter{enumi}{5}
\tightlist
\item
  A drug blocks a key kinase in a signalling cascade, but the tumour
  eventually regrows. Explain this outcome using the concept of feedback
  compensation.
\end{enumerate}

\begin{enumerate}
\def\labelenumi{\arabic{enumi}.}
\setcounter{enumi}{6}
\tightlist
\item
  Explain why the lac operon switch is described as bistable. What
  advantage does bistability provide over a graded (linear) response?
\end{enumerate}

\begin{enumerate}
\def\labelenumi{\arabic{enumi}.}
\setcounter{enumi}{7}
\tightlist
\item
  How does the concept of attractors in phase space relate to cell fate?
  Give an example.
\end{enumerate}

\begin{enumerate}
\def\labelenumi{\arabic{enumi}.}
\setcounter{enumi}{8}
\tightlist
\item
  Apply the active inference framework to explain how a fever might be
  adaptive rather than merely a pathological response.
\end{enumerate}

\begin{enumerate}
\def\labelenumi{\arabic{enumi}.}
\setcounter{enumi}{9}
\tightlist
\item
  In evolutionary terms, what is niche construction, and how does it
  relate to the free energy principle?
\end{enumerate}

\begin{center}\rule{0.5\linewidth}{0.5pt}\end{center}

\subsection*{Synthesis Questions (4 marks
each)}\label{synthesis-questions-4-marks-each-1}
\addcontentsline{toc}{subsection}{Synthesis Questions (4 marks each)}

\begin{enumerate}
\def\labelenumi{\arabic{enumi}.}
\setcounter{enumi}{10}
\tightlist
\item
  Compare and contrast the immune system and an evolving tumour as
  complex adaptive systems. In each case, identify (a) the agents, (b)
  the selection pressures, (c) an emergent property, and (d) a critical
  vulnerability.
\end{enumerate}

\begin{enumerate}
\def\labelenumi{\arabic{enumi}.}
\setcounter{enumi}{11}
\tightlist
\item
  A patient with chronic anxiety has persistent overactivation of the
  sympathetic nervous system even in the absence of objective threat.
  Using predictive coding and active inference, construct a mechanistic
  account of how prior beliefs about threat might drive this pattern.
  How might therapeutic interventions (CBT, interoceptive exposure)
  alter the generative model?
\end{enumerate}

\begin{enumerate}
\def\labelenumi{\arabic{enumi}.}
\setcounter{enumi}{12}
\tightlist
\item
  Describe how ecosystem collapse (e.g., coral bleaching, fishery
  collapse) can be framed as a phase transition in a complex adaptive
  system. What warning signals might precede the transition? Why is
  recovery often much harder than the original decline?
\end{enumerate}

\subsection*{Data and Model Interpretation Questions (3 marks
each)}\label{data-and-model-interpretation-questions-3-marks-each-1}
\addcontentsline{toc}{subsection}{Data and Model Interpretation
Questions (3 marks each)}

\begin{enumerate}
\def\labelenumi{\arabic{enumi}.}
\setcounter{enumi}{13}
\tightlist
\item
  In an agent-based simulation, changing one local rule from ``align
  with three neighbours'' to ``align with six neighbours'' causes global
  order to appear faster. Explain why this is an emergent result rather
  than a command from a central controller.
\end{enumerate}

\begin{enumerate}
\def\labelenumi{\arabic{enumi}.}
\setcounter{enumi}{14}
\tightlist
\item
  A microbial biofilm becomes resistant to an antibiotic even though
  many isolated cells from the same biofilm remain sensitive. Give two
  complex-adaptive explanations.
\end{enumerate}

\begin{enumerate}
\def\labelenumi{\arabic{enumi}.}
\setcounter{enumi}{15}
\tightlist
\item
  A fitness landscape has one tall narrow peak and one lower broad peak.
  Which peak is easier for a population to maintain under mutation and
  environmental noise, and why?
\end{enumerate}

\begin{enumerate}
\def\labelenumi{\arabic{enumi}.}
\setcounter{enumi}{16}
\tightlist
\item
  Explain how a phase transition differs from a gradual linear change.
  Use one biological example.
\end{enumerate}

\begin{enumerate}
\def\labelenumi{\arabic{enumi}.}
\setcounter{enumi}{17}
\tightlist
\item
  A population repeatedly evolves the same resistance mutation in
  separate cultures. Does this prove evolution is deterministic? Give a
  careful answer.
\end{enumerate}

\begin{enumerate}
\def\labelenumi{\arabic{enumi}.}
\setcounter{enumi}{18}
\tightlist
\item
  Compare self-organization in embryonic patterning and in ant foraging.
  What is shared, and what differs?
\end{enumerate}

\begin{enumerate}
\def\labelenumi{\arabic{enumi}.}
\setcounter{enumi}{19}
\tightlist
\item
  A complex system model reproduces one observed pattern but fails under
  a new perturbation. What does this tell you about model validation?
\end{enumerate}

\begin{enumerate}
\def\labelenumi{\arabic{enumi}.}
\setcounter{enumi}{20}
\tightlist
\item
  Explain why heterogeneity among agents can make a population more
  resilient.
\end{enumerate}

\begin{enumerate}
\def\labelenumi{\arabic{enumi}.}
\setcounter{enumi}{21}
\tightlist
\item
  A cancer therapy kills 99\% of tumour cells but leaves a small
  resistant subpopulation. Explain why this can select for relapse.
\end{enumerate}

\begin{enumerate}
\def\labelenumi{\arabic{enumi}.}
\setcounter{enumi}{22}
\tightlist
\item
  Distinguish adaptation by natural selection from adaptation by
  learning within an organism. Why can both be described as adaptive
  processes?
\end{enumerate}

\subsection*{Evaluation and Design Questions (4 marks
each)}\label{evaluation-and-design-questions-4-marks-each-1}
\addcontentsline{toc}{subsection}{Evaluation and Design Questions (4
marks each)}

\begin{enumerate}
\def\labelenumi{\arabic{enumi}.}
\setcounter{enumi}{23}
\tightlist
\item
  Design a simple agent-based model for immune-cell recruitment to an
  infection site. Specify agents, local rules, and one emergent output.
\end{enumerate}

\begin{enumerate}
\def\labelenumi{\arabic{enumi}.}
\setcounter{enumi}{24}
\tightlist
\item
  Evaluate the claim: ``Because complex systems are unpredictable,
  modeling them is useless.''
\end{enumerate}

\begin{enumerate}
\def\labelenumi{\arabic{enumi}.}
\setcounter{enumi}{25}
\tightlist
\item
  Propose a way to detect whether a community pattern is produced by
  local interactions or by a shared external gradient.
\end{enumerate}

\begin{enumerate}
\def\labelenumi{\arabic{enumi}.}
\setcounter{enumi}{26}
\tightlist
\item
  A student says emergence means ``anything can happen.'' Correct this
  using constraints and repeated patterns.
\end{enumerate}

\begin{enumerate}
\def\labelenumi{\arabic{enumi}.}
\setcounter{enumi}{27}
\tightlist
\item
  Build a complex-adaptive explanation of antibiotic resistance that
  includes mutation, selection, horizontal gene transfer, and clinical
  practice.
\end{enumerate}

\begin{enumerate}
\def\labelenumi{\arabic{enumi}.}
\setcounter{enumi}{28}
\tightlist
\item
  Choose one chapter later in the textbook and explain how its topic
  could be reframed as a complex adaptive system.
\end{enumerate}

\begin{enumerate}
\def\labelenumi{\arabic{enumi}.}
\setcounter{enumi}{29}
\tightlist
\item
  Design a falsifiable prediction that would distinguish two competing
  rule sets in an agent-based model of cell migration.
\end{enumerate}

\begin{center}\rule{0.5\linewidth}{0.5pt}\end{center}

\emph{Module: \texttt{src/biology/} (general).}

\newpage

\section*{Questions --- Ch 0.3: Active Inference and the Free Energy
Principle}\label{questions-ch-0.3-active-inference-and-the-free-energy-principle}
\addcontentsline{toc}{section}{Questions --- Ch 0.3: Active Inference
and the Free Energy Principle}

\label{sec:q_unit_0_active_inference}

\emph{This activity accompanies \cref{sec:unit_0_active_inference} of
the textbook --- review that chapter before attempting the exercises
below.}

\begin{center}\rule{0.5\linewidth}{0.5pt}\end{center}

\subsection*{Recall Questions (1 mark
each)}\label{recall-questions-1-mark-each-2}
\addcontentsline{toc}{subsection}{Recall Questions (1 mark each)}

\subsection*{Instructor Use and Coverage
Notes}\label{instructor-use-and-coverage-notes-2}
\addcontentsline{toc}{subsection}{Instructor Use and Coverage Notes}

\begin{itemize}
\tightlist
\item
  \textbf{Coverage target:} Calculate a posterior belief and explain
  what action would reduce uncertainty or restore a set point.
\item
  \textbf{Model/data emphasis:} Bayesian belief updating and
  expected-free-energy-style policy comparison.
\item
  \textbf{Assessment alignment:} Questions and Methods, Representing and
  Describing Data, Argumentation.
\item
  \textbf{Misconception probe:} Active inference is not passive
  prediction; action changes the sensory data that arrive next.
\item
  \textbf{Transfer product:} Map prediction-error reasoning onto
  chemotaxis, thermoregulation, and attention.
\item
  \textbf{Grading focus:} award full credit for mechanism, evidence,
  boundary conditions, and units when a calculation is required.
\item
  \textbf{Suggested use:} draw one recall item, one application item,
  and one synthesis item when building a short quiz from this bank.
\end{itemize}

\begin{enumerate}
\def\labelenumi{\arabic{enumi}.}
\tightlist
\item
  Define `active inference' in your own words and provide a biological
  example.
\end{enumerate}

\begin{enumerate}
\def\labelenumi{\arabic{enumi}.}
\setcounter{enumi}{1}
\tightlist
\item
  What distinguishes a complex adaptive system from a simple system?
\end{enumerate}

\begin{enumerate}
\def\labelenumi{\arabic{enumi}.}
\setcounter{enumi}{2}
\tightlist
\item
  What is the free energy principle? What does free energy measure?
\end{enumerate}

\begin{enumerate}
\def\labelenumi{\arabic{enumi}.}
\setcounter{enumi}{3}
\tightlist
\item
  Describe one biological example of negative feedback and one of
  positive feedback.
\end{enumerate}

\begin{enumerate}
\def\labelenumi{\arabic{enumi}.}
\setcounter{enumi}{4}
\tightlist
\item
  What is the difference between homeostasis and allostasis?
\end{enumerate}

\begin{center}\rule{0.5\linewidth}{0.5pt}\end{center}

\subsection*{Application Questions (2 marks
each)}\label{application-questions-2-marks-each-2}
\addcontentsline{toc}{subsection}{Application Questions (2 marks each)}

\begin{enumerate}
\def\labelenumi{\arabic{enumi}.}
\setcounter{enumi}{5}
\tightlist
\item
  A drug blocks a key kinase in a signalling cascade, but the tumour
  eventually regrows. Explain this outcome using the concept of feedback
  compensation.
\end{enumerate}

\begin{enumerate}
\def\labelenumi{\arabic{enumi}.}
\setcounter{enumi}{6}
\tightlist
\item
  Explain why the lac operon switch is described as bistable. What
  advantage does bistability provide over a graded (linear) response?
\end{enumerate}

\begin{enumerate}
\def\labelenumi{\arabic{enumi}.}
\setcounter{enumi}{7}
\tightlist
\item
  How does the concept of attractors in phase space relate to cell fate?
  Give an example.
\end{enumerate}

\begin{enumerate}
\def\labelenumi{\arabic{enumi}.}
\setcounter{enumi}{8}
\tightlist
\item
  Apply the active inference framework to explain how a fever might be
  adaptive rather than merely a pathological response.
\end{enumerate}

\begin{enumerate}
\def\labelenumi{\arabic{enumi}.}
\setcounter{enumi}{9}
\tightlist
\item
  In evolutionary terms, what is niche construction, and how does it
  relate to the free energy principle?
\end{enumerate}

\begin{center}\rule{0.5\linewidth}{0.5pt}\end{center}

\subsection*{Synthesis Questions (4 marks
each)}\label{synthesis-questions-4-marks-each-2}
\addcontentsline{toc}{subsection}{Synthesis Questions (4 marks each)}

\begin{enumerate}
\def\labelenumi{\arabic{enumi}.}
\setcounter{enumi}{10}
\tightlist
\item
  Compare and contrast the immune system and an evolving tumour as
  complex adaptive systems. In each case, identify (a) the agents, (b)
  the selection pressures, (c) an emergent property, and (d) a critical
  vulnerability.
\end{enumerate}

\begin{enumerate}
\def\labelenumi{\arabic{enumi}.}
\setcounter{enumi}{11}
\tightlist
\item
  A patient with chronic anxiety has persistent overactivation of the
  sympathetic nervous system even in the absence of objective threat.
  Using predictive coding and active inference, construct a mechanistic
  account of how prior beliefs about threat might drive this pattern.
  How might therapeutic interventions (CBT, interoceptive exposure)
  alter the generative model?
\end{enumerate}

\begin{enumerate}
\def\labelenumi{\arabic{enumi}.}
\setcounter{enumi}{12}
\tightlist
\item
  Describe how ecosystem collapse (e.g., coral bleaching, fishery
  collapse) can be framed as a phase transition in a complex adaptive
  system. What warning signals might precede the transition? Why is
  recovery often much harder than the original decline?
\end{enumerate}

\subsection*{Data and Model Interpretation Questions (3 marks
each)}\label{data-and-model-interpretation-questions-3-marks-each-2}
\addcontentsline{toc}{subsection}{Data and Model Interpretation
Questions (3 marks each)}

\begin{enumerate}
\def\labelenumi{\arabic{enumi}.}
\setcounter{enumi}{13}
\tightlist
\item
  A prior predicts body temperature should be 37.0 degrees C, but
  sensory data report 36.2 degrees C with low uncertainty. Predict
  whether belief updating or action is more likely to reduce prediction
  error, and justify your answer.
\end{enumerate}

\begin{enumerate}
\def\labelenumi{\arabic{enumi}.}
\setcounter{enumi}{14}
\tightlist
\item
  Explain the difference between prediction error and surprise in an
  active-inference account.
\end{enumerate}

\begin{enumerate}
\def\labelenumi{\arabic{enumi}.}
\setcounter{enumi}{15}
\tightlist
\item
  A patient with panic disorder interprets benign heart-rate increases
  as evidence of danger. Identify the prior, the sensory evidence, and
  the prediction-error loop.
\end{enumerate}

\begin{enumerate}
\def\labelenumi{\arabic{enumi}.}
\setcounter{enumi}{16}
\tightlist
\item
  Compare homeostasis and allostasis from an active-inference
  perspective.
\end{enumerate}

\begin{enumerate}
\def\labelenumi{\arabic{enumi}.}
\setcounter{enumi}{17}
\tightlist
\item
  A chemotactic bacterium tumbles more often when attractant
  concentration falls. Describe this as action changing future sensory
  input.
\end{enumerate}

\begin{enumerate}
\def\labelenumi{\arabic{enumi}.}
\setcounter{enumi}{18}
\tightlist
\item
  What does it mean to assign high precision to a sensory signal? Give
  one adaptive and one maladaptive example.
\end{enumerate}

\begin{enumerate}
\def\labelenumi{\arabic{enumi}.}
\setcounter{enumi}{19}
\tightlist
\item
  A model predicts that reducing sensory precision should reduce
  self-generated tickle sensations. Explain why.
\end{enumerate}

\begin{enumerate}
\def\labelenumi{\arabic{enumi}.}
\setcounter{enumi}{20}
\tightlist
\item
  Explain why active inference is not the same as reinforcement
  learning, even though both can describe adaptive behaviour.
\end{enumerate}

\begin{enumerate}
\def\labelenumi{\arabic{enumi}.}
\setcounter{enumi}{21}
\tightlist
\item
  Use active inference to explain why inflammation can be both
  protective and harmful.
\end{enumerate}

\begin{enumerate}
\def\labelenumi{\arabic{enumi}.}
\setcounter{enumi}{22}
\tightlist
\item
  A generative model predicts predator risk from smell and sound. If
  smell is blocked, what should happen to reliance on sound, and what
  experiment would test this?
\end{enumerate}

\subsection*{Evaluation and Design Questions (4 marks
each)}\label{evaluation-and-design-questions-4-marks-each-2}
\addcontentsline{toc}{subsection}{Evaluation and Design Questions (4
marks each)}

\begin{enumerate}
\def\labelenumi{\arabic{enumi}.}
\setcounter{enumi}{23}
\tightlist
\item
  Design a behavioural experiment to test whether interoceptive
  precision differs between anxious and non-anxious participants.
\end{enumerate}

\begin{enumerate}
\def\labelenumi{\arabic{enumi}.}
\setcounter{enumi}{24}
\tightlist
\item
  Evaluate the claim: ``Active inference explains everything, so it
  explains nothing.''
\end{enumerate}

\begin{enumerate}
\def\labelenumi{\arabic{enumi}.}
\setcounter{enumi}{25}
\tightlist
\item
  Propose an active-inference account of placebo analgesia.
\end{enumerate}

\begin{enumerate}
\def\labelenumi{\arabic{enumi}.}
\setcounter{enumi}{26}
\tightlist
\item
  A clinician wants to reduce chronic dyspnea by changing primarily lung
  mechanics. Use active inference to explain why cognitive or
  interoceptive therapy might also help.
\end{enumerate}

\begin{enumerate}
\def\labelenumi{\arabic{enumi}.}
\setcounter{enumi}{27}
\tightlist
\item
  Build an active-inference explanation of immune tolerance.
\end{enumerate}

\begin{enumerate}
\def\labelenumi{\arabic{enumi}.}
\setcounter{enumi}{28}
\tightlist
\item
  Compare an organism seeking food and a scientist designing an
  experiment as epistemic agents.
\end{enumerate}

\begin{enumerate}
\def\labelenumi{\arabic{enumi}.}
\setcounter{enumi}{29}
\tightlist
\item
  Choose one later textbook topic and formulate an active-inference
  version of its central mechanism.
\end{enumerate}

\begin{center}\rule{0.5\linewidth}{0.5pt}\end{center}

\emph{Module: \texttt{src/biology/} (general).}

\newpage

\section*{Questions --- Ch 1: Atoms, Molecules, and Chemical
Bonds}\label{questions-ch-1-atoms-molecules-and-chemical-bonds}
\addcontentsline{toc}{section}{Questions --- Ch 1: Atoms, Molecules, and
Chemical Bonds}

\label{sec:q_unit_I_atoms_molecules}

\subsection*{Instructor Use and Coverage
Notes}\label{instructor-use-and-coverage-notes-3}
\addcontentsline{toc}{subsection}{Instructor Use and Coverage Notes}

\begin{itemize}
\tightlist
\item
  \textbf{Coverage target:} Predict bond polarity, formal charge, or
  isotope behaviour and connect it to a biological consequence.
\item
  \textbf{Model/data emphasis:} Formal charge and
  electronegativity-difference reasoning.
\item
  \textbf{Assessment alignment:} Concept Explanation, Statistical Tests
  and Data Analysis, Argumentation.
\item
  \textbf{Misconception probe:} Weak interactions are not unimportant;
  many weak interactions together dominate structure and specificity.
\item
  \textbf{Transfer product:} Use atomic reasoning to explain a medical
  tracer, enzyme cofactor, or membrane-solubility problem.
\item
  \textbf{Grading focus:} award full credit for mechanism, evidence,
  boundary conditions, and units when a calculation is required.
\item
  \textbf{Suggested use:} draw one recall item, one application item,
  and one synthesis item when building a short quiz from this bank.
\end{itemize}

\subsection*{Questions 1--10: Recall and
Comprehension}\label{questions-110-recall-and-comprehension}
\addcontentsline{toc}{subsection}{Questions 1--10: Recall and
Comprehension}

\emph{This activity accompanies \cref{sec:unit_I_atoms_molecules} of the
textbook --- review that chapter before attempting the exercises below.}

\begin{enumerate}
\def\labelenumi{\arabic{enumi}.}
\tightlist
\item
  Define an atom and identify its three subatomic particles. Which
  determines the element's identity?
\end{enumerate}

\begin{enumerate}
\def\labelenumi{\arabic{enumi}.}
\setcounter{enumi}{1}
\tightlist
\item
  What is the atomic number of carbon? How many valence electrons does
  carbon have?
\end{enumerate}

\begin{enumerate}
\def\labelenumi{\arabic{enumi}.}
\setcounter{enumi}{2}
\tightlist
\item
  Distinguish between covalent and ionic bonds. Give one biological
  example of each.
\end{enumerate}

\begin{enumerate}
\def\labelenumi{\arabic{enumi}.}
\setcounter{enumi}{3}
\tightlist
\item
  Define electronegativity. Which element has the highest
  electronegativity?
\end{enumerate}

\begin{enumerate}
\def\labelenumi{\arabic{enumi}.}
\setcounter{enumi}{4}
\tightlist
\item
  What makes a covalent bond polar? What partial charges (\texorpdfstring{\ensuremath{\delta}}{delta}+ and \texorpdfstring{\ensuremath{\delta}}{delta}\texorpdfstring{\ensuremath{-}}{-})
  arise in an O--H bond?
\end{enumerate}

\begin{enumerate}
\def\labelenumi{\arabic{enumi}.}
\setcounter{enumi}{5}
\tightlist
\item
  Define a hydrogen bond. Which atoms typically act as hydrogen bond
  donors and acceptors?
\end{enumerate}

\begin{enumerate}
\def\labelenumi{\arabic{enumi}.}
\setcounter{enumi}{6}
\tightlist
\item
  What is a van der Waals interaction? In which large molecules are van
  der Waals forces collectively significant?
\end{enumerate}

\begin{enumerate}
\def\labelenumi{\arabic{enumi}.}
\setcounter{enumi}{7}
\tightlist
\item
  List four properties of water that arise directly from hydrogen
  bonding.
\end{enumerate}

\begin{enumerate}
\def\labelenumi{\arabic{enumi}.}
\setcounter{enumi}{8}
\tightlist
\item
  What is pH? What hydrogen ion concentration corresponds to pH 7 and pH
  4?
\end{enumerate}

\begin{enumerate}
\def\labelenumi{\arabic{enumi}.}
\setcounter{enumi}{9}
\tightlist
\item
  Define an acid and a base using the Brønsted-Lowry definition. Give
  one biological example of each.
\end{enumerate}

\subsection*{Questions 11--20: Application and
Analysis}\label{questions-1120-application-and-analysis}
\addcontentsline{toc}{subsection}{Questions 11--20: Application and
Analysis}

\begin{enumerate}
\def\labelenumi{\arabic{enumi}.}
\setcounter{enumi}{10}
\tightlist
\item
  Carbon forms four covalent bonds. Explain how this tetravalency allows
  carbon to form diverse structures including chains, branches, and
  rings. Why is carbon the basis of biological molecules rather than
  silicon?
\end{enumerate}

\begin{enumerate}
\def\labelenumi{\arabic{enumi}.}
\setcounter{enumi}{11}
\tightlist
\item
  The electronegativity difference for O--H is 1.4, for N--H is 0.9, and
  for C--H is 0.4. Rank these bonds from most to least polar and predict
  which can form hydrogen bonds.
\end{enumerate}

\begin{enumerate}
\def\labelenumi{\arabic{enumi}.}
\setcounter{enumi}{12}
\tightlist
\item
  A buffer solution contains equal concentrations of carbonic acid
  (H\texorpdfstring{\textsubscript{2}}{2}CO\texorpdfstring{\textsubscript{3}}{3}) and bicarbonate (HCO\texorpdfstring{\textsubscript{3}}{3}\texorpdfstring{\textsuperscript{-}}{-}). At what pH does the buffer have
  maximum capacity? Express this using the Henderson-Hasselbalch
  equation (pKa of carbonic acid = 6.1).
\end{enumerate}

\begin{enumerate}
\def\labelenumi{\arabic{enumi}.}
\setcounter{enumi}{13}
\tightlist
\item
  Explain why non-polar molecules such as cholesterol are insoluble in
  water but soluble in lipids. What thermodynamic principle governs this
  behaviour (hint: entropy of hydrophobic interactions)?
\end{enumerate}

\begin{enumerate}
\def\labelenumi{\arabic{enumi}.}
\setcounter{enumi}{14}
\tightlist
\item
  A molecule has the formula C\texorpdfstring{\textsubscript{1}}{1}\texorpdfstring{\textsubscript{8}}{8}H\texorpdfstring{\textsubscript{3}}{3}\texorpdfstring{\textsubscript{6}}{6}O\texorpdfstring{\textsubscript{2}}{2}. Predict whether it is saturated
  or unsaturated based on the degree-of-unsaturation (DoU) formula.
  Identify it as a likely carbohydrate, lipid, or protein monomer.
\end{enumerate}

\begin{enumerate}
\def\labelenumi{\arabic{enumi}.}
\setcounter{enumi}{15}
\tightlist
\item
  Compare the properties of strong acids (e.g., HCl) vs weak acids
  (e.g., acetic acid, pKa = 4.76) at the same initial concentration. At
  pH 5, what fraction of acetic acid molecules are dissociated? (Use
  Henderson-Hasselbalch.)
\end{enumerate}

\begin{enumerate}
\def\labelenumi{\arabic{enumi}.}
\setcounter{enumi}{16}
\tightlist
\item
  A cell's cytoplasmic pH drops from 7.4 to 6.9 during intense exercise.
  By what factor does {[}H\texorpdfstring{\textsuperscript{+}}{+}{]} increase? What buffers in the cytoplasm
  resist this pH change?
\end{enumerate}

\begin{enumerate}
\def\labelenumi{\arabic{enumi}.}
\setcounter{enumi}{17}
\tightlist
\item
  Explain why CO\texorpdfstring{\textsubscript{2}}{2} dissolved in blood lowers pH. Write the relevant
  chemical equations, naming each species formed.
\end{enumerate}

\begin{enumerate}
\def\labelenumi{\arabic{enumi}.}
\setcounter{enumi}{18}
\tightlist
\item
  The peptide bond in proteins involves C--N with partial double-bond
  character due to resonance. Explain how this resonance prevents
  rotation around the peptide bond and what structural constraint this
  places on protein backbone geometry.
\end{enumerate}

\begin{enumerate}
\def\labelenumi{\arabic{enumi}.}
\setcounter{enumi}{19}
\tightlist
\item
  Magnesium (Mg\texorpdfstring{\textsuperscript{2}}{2}\texorpdfstring{\textsuperscript{+}}{+}) is a common enzyme cofactor. Using your understanding
  of atomic structure, explain why Mg\texorpdfstring{\textsuperscript{2}}{2}\texorpdfstring{\textsuperscript{+}}{+} can stabilise negatively charged
  phosphate groups in ATP. What orbital property of Mg\texorpdfstring{\textsuperscript{2}}{2}\texorpdfstring{\textsuperscript{+}}{+} makes it a
  Lewis acid?
\end{enumerate}

\subsection*{Questions 21--30: Synthesis and
Evaluation}\label{questions-2130-synthesis-and-evaluation}
\addcontentsline{toc}{subsection}{Questions 21--30: Synthesis and
Evaluation}

\begin{enumerate}
\def\labelenumi{\arabic{enumi}.}
\setcounter{enumi}{20}
\tightlist
\item
  Evaluate the statement: ``Life could have evolved using silicon
  instead of carbon as its molecular backbone.'' Using comparative
  chemistry, identify at least two key chemical differences between the
  bonding properties of C and Si that make carbon uniquely suited for
  biological macromolecules.
\end{enumerate}

\begin{enumerate}
\def\labelenumi{\arabic{enumi}.}
\setcounter{enumi}{21}
\tightlist
\item
  Design an experiment to test whether the hydrophobic effect is
  entropy-driven. Describe what measurements you would take, what
  temperature you would vary, and what thermodynamic parameters you
  would calculate. Predict the outcome.
\end{enumerate}

\begin{enumerate}
\def\labelenumi{\arabic{enumi}.}
\setcounter{enumi}{22}
\tightlist
\item
  A drug molecule must cross the blood-brain barrier (which is composed
  of lipid bilayer membranes). Using the structure-activity relationship
  principle, explain what molecular properties (MW, logP, H-bond
  donors/acceptors) favour CNS penetration. How would you modify a
  highly polar drug to improve CNS penetration?
\end{enumerate}

\begin{enumerate}
\def\labelenumi{\arabic{enumi}.}
\setcounter{enumi}{23}
\tightlist
\item
  Carbonic anhydrase catalyses CO\texorpdfstring{\textsubscript{2}}{2} + H\texorpdfstring{\textsubscript{2}}{2}O \texorpdfstring{\ensuremath{\rightleftharpoons}}{<=>} H\texorpdfstring{\textsubscript{2}}{2}CO\texorpdfstring{\textsubscript{3}}{3} with a rate constant
  about 10\texorpdfstring{\textsuperscript{6}}{6}\texorpdfstring{\ensuremath{\times}}{x} faster than the uncatalysed reaction. The active
  site uses a Zn\texorpdfstring{\textsuperscript{2}}{2}\texorpdfstring{\textsuperscript{+}}{+} ion to polarise the O--H bond of a water molecule.
  Explain the mechanistic logic: why does metal coordination of water
  make it a better nucleophile in this context?
\end{enumerate}

\begin{enumerate}
\def\labelenumi{\arabic{enumi}.}
\setcounter{enumi}{24}
\tightlist
\item
  Compare the evolutionary origins of the genetic code (DNA/RNA using N-
  and O-containing bases) with a hypothetical alternative code using
  entirely non-polar building blocks. What biochemical properties would
  be lost or gained? Evaluate why polarity and hydrogen bonding in
  nucleobases are essential for genetic fidelity.
\end{enumerate}

\begin{enumerate}
\def\labelenumi{\arabic{enumi}.}
\setcounter{enumi}{25}
\tightlist
\item
  A newly discovered antibiotic (X) functions by chelating Mg\texorpdfstring{\textsuperscript{2}}{2}\texorpdfstring{\textsuperscript{+}}{+} ions in
  bacterial ribosomes. Predict: (a) which step of protein synthesis
  would be disrupted; (b) whether X would be active vs gram-positive or
  gram-negative bacteria; (c) possible mechanisms of resistance the
  bacterium might evolve.
\end{enumerate}

\begin{enumerate}
\def\labelenumi{\arabic{enumi}.}
\setcounter{enumi}{26}
\tightlist
\item
  The Henderson-Hasselbalch equation assumes that pKa is constant.
  However, the pKa of His residues in proteins is strongly
  context-dependent (ranging from 4.0 to 8.0). Explain how protein
  microenvironment (electrostatics, hydrophobicity, hydrogen bonding)
  shifts the pKa of an ionisable side chain, and why this is
  functionally important for enzyme catalysis.
\end{enumerate}

\begin{enumerate}
\def\labelenumi{\arabic{enumi}.}
\setcounter{enumi}{27}
\tightlist
\item
  Bohr effect: CO\texorpdfstring{\textsubscript{2}}{2} and H\texorpdfstring{\textsuperscript{+}}{+} reduce haemoglobin's O\texorpdfstring{\textsubscript{2}}{2} affinity by
  stabilising the T-state (deoxy form). Using atomic-level reasoning:
  (a) which specific residues of haemoglobin are protonated at low pH;
  (b) how does protonation alter quaternary structure; (c) why is this
  physiologically beneficial in exercising tissues?
\end{enumerate}

\begin{enumerate}
\def\labelenumi{\arabic{enumi}.}
\setcounter{enumi}{28}
\tightlist
\item
  A patient presents with severe metabolic acidosis (blood pH 7.1, HCO\texorpdfstring{\textsubscript{3}}{3}\texorpdfstring{\textsuperscript{-}}{-}
  = 12 mEq/L, pCO\texorpdfstring{\textsubscript{2}}{2} = 30 mmHg). Using Henderson-Hasselbalch (pH = 6.1 +
  log{[}HCO\texorpdfstring{\textsubscript{3}}{3}\texorpdfstring{\textsuperscript{-}}{-}{]}/{[}0.03 \texorpdfstring{\ensuremath{\times}}{x} pCO\texorpdfstring{\textsubscript{2}}{2}{]}), verify the expected pH vs measured.
  Is there respiratory compensation? What primary disorder is this
  (metabolic or respiratory) and how would you distinguish a primary
  from a compensatory change?
\end{enumerate}

\begin{enumerate}
\def\labelenumi{\arabic{enumi}.}
\setcounter{enumi}{29}
\tightlist
\item
  Propose a molecular mechanism by which a histone deacetylase (HDAC)
  inhibitor (e.g., vorinostat) could be used to treat cancer. Trace from
  HDAC inhibition \texorpdfstring{\ensuremath{\to}}{->} altered histone acetylation state \texorpdfstring{\ensuremath{\to}}{->} chromatin
  opening \texorpdfstring{\ensuremath{\to}}{->} transcription of tumour suppressor genes \texorpdfstring{\ensuremath{\to}}{->} reduced tumour
  cell survival. Identify the specific bond type removed/restored at
  each step.
\end{enumerate}

\newpage

\section*{Questions --- Ch 2: Water --- The Molecule of
Life}\label{questions-ch-2-water-the-molecule-of-life}
\addcontentsline{toc}{section}{Questions --- Ch 2: Water --- The
Molecule of Life}

\label{sec:q_unit_I_water_and_life}

\subsection*{Instructor Use and Coverage
Notes}\label{instructor-use-and-coverage-notes-4}
\addcontentsline{toc}{subsection}{Instructor Use and Coverage Notes}

\begin{itemize}
\tightlist
\item
  \textbf{Coverage target:} Calculate a pH or buffer ratio and interpret
  what it means for a biological system.
\item
  \textbf{Model/data emphasis:} Henderson-Hasselbalch and pH-scale
  calculations.
\item
  \textbf{Assessment alignment:} Concept Explanation, Statistical Tests
  and Data Analysis, Argumentation.
\item
  \textbf{Misconception probe:} Water is not an inert background; it is
  an active participant in structure and reaction chemistry.
\item
  \textbf{Transfer product:} Transfer water-property reasoning to blood
  buffering, plant transport, or protein folding.
\item
  \textbf{Grading focus:} award full credit for mechanism, evidence,
  boundary conditions, and units when a calculation is required.
\item
  \textbf{Suggested use:} draw one recall item, one application item,
  and one synthesis item when building a short quiz from this bank.
\end{itemize}

\subsection*{Questions 1--10: Recall and
Comprehension}\label{questions-110-recall-and-comprehension-1}
\addcontentsline{toc}{subsection}{Questions 1--10: Recall and
Comprehension}

\emph{This activity accompanies \cref{sec:unit_I_water_and_life} of the
textbook --- review that chapter before attempting the exercises below.}

\begin{enumerate}
\def\labelenumi{\arabic{enumi}.}
\tightlist
\item
  Draw the water molecule; label the bond angle (104.5°), partial
  charges, and hydrogen bonds with adjacent molecules.
\end{enumerate}

\begin{enumerate}
\def\labelenumi{\arabic{enumi}.}
\setcounter{enumi}{1}
\tightlist
\item
  Define specific heat capacity. Why does water have a high specific
  heat compared to most liquids?
\end{enumerate}

\begin{enumerate}
\def\labelenumi{\arabic{enumi}.}
\setcounter{enumi}{2}
\tightlist
\item
  What is the heat of vaporisation of water (cal/g) and what biological
  process exploits this?
\end{enumerate}

\begin{enumerate}
\def\labelenumi{\arabic{enumi}.}
\setcounter{enumi}{3}
\tightlist
\item
  Define cohesion and adhesion. Which property allows water to form a
  meniscus in a glass tube?
\end{enumerate}

\begin{enumerate}
\def\labelenumi{\arabic{enumi}.}
\setcounter{enumi}{4}
\tightlist
\item
  What is osmosis? In which direction does water move --- toward higher
  or lower solute concentration?
\end{enumerate}

\begin{enumerate}
\def\labelenumi{\arabic{enumi}.}
\setcounter{enumi}{5}
\tightlist
\item
  Define water potential (\texorpdfstring{\ensuremath{\Psi}}{Psi}). Write the equation \texorpdfstring{\ensuremath{\Psi}}{Psi} = \texorpdfstring{\ensuremath{\Psi}}{Psi}\_s + \texorpdfstring{\ensuremath{\Psi}}{Psi}\_p and
  define each term.
\end{enumerate}

\begin{enumerate}
\def\labelenumi{\arabic{enumi}.}
\setcounter{enumi}{6}
\tightlist
\item
  What is tonicity? Distinguish hypertonic, hypotonic, and isotonic
  solutions relative to a red blood cell.
\end{enumerate}

\begin{enumerate}
\def\labelenumi{\arabic{enumi}.}
\setcounter{enumi}{7}
\tightlist
\item
  Define capillarity. Name two forces responsible for capillary rise in
  a tube.
\end{enumerate}

\begin{enumerate}
\def\labelenumi{\arabic{enumi}.}
\setcounter{enumi}{8}
\tightlist
\item
  What is the pH of pure water at 25°C? What is the Kw of water at this
  temperature?
\end{enumerate}

\begin{enumerate}
\def\labelenumi{\arabic{enumi}.}
\setcounter{enumi}{9}
\tightlist
\item
  Describe what happens to a freshwater fish placed in seawater
  (hypertonic solution) in terms of osmosis.
\end{enumerate}

\subsection*{Questions 11--20: Application and
Analysis}\label{questions-1120-application-and-analysis-1}
\addcontentsline{toc}{subsection}{Questions 11--20: Application and
Analysis}

\begin{enumerate}
\def\labelenumi{\arabic{enumi}.}
\setcounter{enumi}{10}
\tightlist
\item
  Calculate: if the osmotic potential of a plant cell is \texorpdfstring{\ensuremath{-}}{-}0.8 MPa and
  pressure potential is +0.3 MPa, what is the water potential? Will
  water enter or leave the cell if it is placed in a solution with \texorpdfstring{\ensuremath{\Psi}}{Psi} =
  \texorpdfstring{\ensuremath{-}}{-}0.4 MPa?
\end{enumerate}

\begin{enumerate}
\def\labelenumi{\arabic{enumi}.}
\setcounter{enumi}{11}
\tightlist
\item
  During sweating, 1 litre of sweat evaporated removes approximately 580
  kcal of heat from the body. Calculate: if a marathon runner produces 2
  L/h of sweat, how many watts of heat are removed? (1 watt = 1 J/s)
\end{enumerate}

\begin{enumerate}
\def\labelenumi{\arabic{enumi}.}
\setcounter{enumi}{12}
\tightlist
\item
  Sea ice insulates frigid arctic waters and protects aquatic organisms.
  Explain, at the molecular level, why ice is less dense than liquid
  water and why this unusual property has profound ecological
  consequences.
\end{enumerate}

\begin{enumerate}
\def\labelenumi{\arabic{enumi}.}
\setcounter{enumi}{13}
\tightlist
\item
  An onion cell with osmotic potential \texorpdfstring{\ensuremath{-}}{-}0.6 MPa is placed in a 0.3 M
  NaCl solution. The NaCl solution has \texorpdfstring{\ensuremath{\Psi}}{Psi}\_s = \texorpdfstring{\ensuremath{-}}{-}0.75 MPa. Predict the
  direction of water movement and whether the cell will plasmolyse or
  swell.
\end{enumerate}

\begin{enumerate}
\def\labelenumi{\arabic{enumi}.}
\setcounter{enumi}{14}
\tightlist
\item
  Compare the water requirements of a plant using C\texorpdfstring{\textsubscript{4}}{4} photosynthesis
  (e.g., maize) vs C\texorpdfstring{\textsubscript{3}}{3} photosynthesis (e.g., wheat) in a hot, dry
  environment. Explain in terms of stomatal conductance and water-use
  efficiency.
\end{enumerate}

\begin{enumerate}
\def\labelenumi{\arabic{enumi}.}
\setcounter{enumi}{15}
\tightlist
\item
  Describe three specific properties of water that make it essential for
  the cohesion-tension mechanism of xylem transport. What evidence
  refutes the alternative ``root pressure'' hypothesis as the primary
  force driving water to the tops of tall trees?
\end{enumerate}

\begin{enumerate}
\def\labelenumi{\arabic{enumi}.}
\setcounter{enumi}{16}
\tightlist
\item
  A student proposes using ethanol instead of water as a biological
  solvent. Identify three properties of water (not of ethanol) that
  would be lost, and predict the consequences for enzyme function, ion
  transport, and pH buffering.
\end{enumerate}

\begin{enumerate}
\def\labelenumi{\arabic{enumi}.}
\setcounter{enumi}{17}
\tightlist
\item
  During hypothermia, core body temperature drops to 32°C. Using
  specific heat capacity principles, calculate the heat (kJ) that must
  be extracted from a 70 kg person (composed of about 60\%
  water) to lower their temperature from 37°C to 32°C (specific heat of
  the body \texorpdfstring{\ensuremath{\approx}}{~} 3.47 kJ/kg\texorpdfstring{\ensuremath{\cdot}}{.}°C).
\end{enumerate}

\begin{enumerate}
\def\labelenumi{\arabic{enumi}.}
\setcounter{enumi}{18}
\tightlist
\item
  Explain the role of aquaporins in transmembrane water movement. How
  does the selectivity filter exclude protons (H\texorpdfstring{\textsubscript{3}}{3}O\texorpdfstring{\textsuperscript{+}}{+}) despite water
  molecules sharing similar size? What disease results from aquaporin-2
  deficiency?
\end{enumerate}

\begin{enumerate}
\def\labelenumi{\arabic{enumi}.}
\setcounter{enumi}{19}
\tightlist
\item
  The mangrove tree lives in salty seawater. Its roots must extract
  fresh water against an extremely negative osmotic potential. Explain
  what root structure adaptations mangroves have developed, and
  calculate the minimum root pressure required if external seawater has
  \texorpdfstring{\ensuremath{\Psi}}{Psi}\_s = \texorpdfstring{\ensuremath{-}}{-}2.5 MPa and leaf water loss creates \texorpdfstring{\ensuremath{\Psi}}{Psi}\_p = \texorpdfstring{\ensuremath{-}}{-}1.0 MPa.
\end{enumerate}

\subsection*{Questions 21--30: Synthesis and
Evaluation}\label{questions-2130-synthesis-and-evaluation-1}
\addcontentsline{toc}{subsection}{Questions 21--30: Synthesis and
Evaluation}

\begin{enumerate}
\def\labelenumi{\arabic{enumi}.}
\setcounter{enumi}{20}
\tightlist
\item
  Water is simultaneously the product and reactant of countless
  biochemical reactions (hydrolysis, condensation, and hydration).
  Evaluate whether increasing intracellular water concentration would
  shift equilibrium toward condensation reactions (protein synthesis) or
  hydrolysis reactions. What does Le Chatelier's principle predict?
\end{enumerate}

\begin{enumerate}
\def\labelenumi{\arabic{enumi}.}
\setcounter{enumi}{21}
\tightlist
\item
  Propose an experiment to test whether the anomalous expansion of water
  upon freezing is due to hydrogen bond structure. What control
  conditions and measurements would you use? What would falsify the
  hydrogen bond hypothesis?
\end{enumerate}

\begin{enumerate}
\def\labelenumi{\arabic{enumi}.}
\setcounter{enumi}{22}
\tightlist
\item
  The blood-brain barrier restricts water flux into the CNS. Cerebral
  oedema (excess water accumulation in brain tissue) can be fatal. Using
  water potential principles, explain: (a) why hyponatraemia (low blood
  {[}Na\texorpdfstring{\textsuperscript{+}}{+}{]}) causes cerebral oedema; (b) why mannitol infusion reduces
  oedema; (c) how AQP4 water channels contribute to normal fluid
  homeostasis.
\end{enumerate}

\begin{enumerate}
\def\labelenumi{\arabic{enumi}.}
\setcounter{enumi}{23}
\tightlist
\item
  Compare the water potential of: (a) pure water; (b) 0.1 M sucrose
  solution; (c) wilted plant cell (\texorpdfstring{\ensuremath{-}}{-}1.2 MPa osmotic, +0.1 MPa pressure);
  (d) turgid plant cell (\texorpdfstring{\ensuremath{-}}{-}1.2 MPa osmotic, +0.9 MPa pressure). Rank them
  from highest to lowest \texorpdfstring{\ensuremath{\Psi}}{Psi} and predict the direction of water flux
  between each pair.
\end{enumerate}

\begin{enumerate}
\def\labelenumi{\arabic{enumi}.}
\setcounter{enumi}{24}
\tightlist
\item
  Desiccation-tolerant (resurrection) plants can survive to \textless{}
  5\% water content and revive when rehydrated. Propose two cellular
  mechanisms that allow membranes and proteins to remain functional in
  the nearly anhydrous state. Consider the role of trehalose as a
  vitrifying agent.
\end{enumerate}

\begin{enumerate}
\def\labelenumi{\arabic{enumi}.}
\setcounter{enumi}{25}
\tightlist
\item
  Reverse osmosis (RO) desalination uses high pressure to push seawater
  through semi-permeable membranes. Calculate: if seawater has an
  osmotic pressure of approximately 27 atm (2.7 MPa), what minimum
  external pressure must the RO system exert to produce fresh water?
  What energy (kJ) is required per litre, given 1 atm \texorpdfstring{\ensuremath{\approx}}{~} 101.3 kPa and
  volume = 1 L?
\end{enumerate}

\begin{enumerate}
\def\labelenumi{\arabic{enumi}.}
\setcounter{enumi}{26}
\tightlist
\item
  During ischemia (reduced blood flow), intracellular Na\texorpdfstring{\textsuperscript{+}}{+} rises due to
  Na/K-ATPase failure. Predict the effect of rising intracellular Na\texorpdfstring{\textsuperscript{+}}{+}
  on: (a) intracellular osmotic potential; (b) water flux across the
  plasma membrane; (c) cell volume. How does this contribute to
  ischaemic cell death?
\end{enumerate}

\begin{enumerate}
\def\labelenumi{\arabic{enumi}.}
\setcounter{enumi}{27}
\tightlist
\item
  Antarctic ice fish have no haemoglobin and survive at \texorpdfstring{\ensuremath{-}}{-}1.9°C seawater
  by relying on the unusual fluidity of water at low temperatures.
  Explain: (a) how antifreeze proteins (AFPs) prevent ice crystal
  formation at subzero temperatures (using the concepts of ice crystal
  growth inhibition); (b) how AFP ice binding differs from osmotic
  freezing protection (as in glycerol).
\end{enumerate}

\begin{enumerate}
\def\labelenumi{\arabic{enumi}.}
\setcounter{enumi}{28}
\tightlist
\item
  Water's high surface tension (72 mN/m at 20°C) is biologically
  critical. Describe two specific biological examples of surface tension
  exploitation (one arthropod, one mammalian respiratory) and explain
  the molecular basis of surface tension in each, and what would happen
  if surface tension were reduced by half in each case.
\end{enumerate}

\begin{enumerate}
\def\labelenumi{\arabic{enumi}.}
\setcounter{enumi}{29}
\tightlist
\item
  Critique the role of water activity (a\_w) vs water potential in
  determining microbial growth. Explain why halophilic bacteria can grow
  at a\_w = 0.75 (equivalent to about 5 M NaCl), while most
  bacteria cannot survive below a\_w = 0.95. How does reduced water
  activity inhibit cellular processes even when cells are fully
  hydrated? What field of food science exploits this principle and how?
\end{enumerate}

\newpage

\section*{Questions --- Ch 3: Biological
Macromolecules}\label{questions-ch-3-biological-macromolecules}
\addcontentsline{toc}{section}{Questions --- Ch 3: Biological
Macromolecules}

\label{sec:q_unit_I_macromolecules}

\subsection*{Instructor Use and Coverage
Notes}\label{instructor-use-and-coverage-notes-5}
\addcontentsline{toc}{subsection}{Instructor Use and Coverage Notes}

\begin{itemize}
\tightlist
\item
  \textbf{Coverage target:} Given a molecular feature, infer likely
  solubility, interaction partners, or cellular role.
\item
  \textbf{Model/data emphasis:} Polymerization energetics and
  sequence-to-structure reasoning.
\item
  \textbf{Assessment alignment:} Concept Explanation, Statistical Tests
  and Data Analysis, Argumentation.
\item
  \textbf{Misconception probe:} Structure is not decoration; small
  chemical changes can redirect function and recognition.
\item
  \textbf{Transfer product:} Explain how a mutation, lipid substitution,
  or glycosylation change propagates to phenotype.
\item
  \textbf{Grading focus:} award full credit for mechanism, evidence,
  boundary conditions, and units when a calculation is required.
\item
  \textbf{Suggested use:} draw one recall item, one application item,
  and one synthesis item when building a short quiz from this bank.
\end{itemize}

\subsection*{Questions 1--10: Recall and
Comprehension}\label{questions-110-recall-and-comprehension-2}
\addcontentsline{toc}{subsection}{Questions 1--10: Recall and
Comprehension}

\emph{This activity accompanies \cref{sec:unit_I_macromolecules} of the
textbook --- review that chapter before attempting the exercises below.}

\begin{enumerate}
\def\labelenumi{\arabic{enumi}.}
\tightlist
\item
  Name the four classes of biological macromolecules and list the
  monomers of each.
\end{enumerate}

\begin{enumerate}
\def\labelenumi{\arabic{enumi}.}
\setcounter{enumi}{1}
\tightlist
\item
  What is a condensation (dehydration) reaction? What bond is formed
  when two amino acids condense?
\end{enumerate}

\begin{enumerate}
\def\labelenumi{\arabic{enumi}.}
\setcounter{enumi}{2}
\tightlist
\item
  Describe the four levels of protein structure. What type of bond holds
  the primary structure together?
\end{enumerate}

\begin{enumerate}
\def\labelenumi{\arabic{enumi}.}
\setcounter{enumi}{3}
\tightlist
\item
  What is the difference between \texorpdfstring{\ensuremath{\alpha}}{alpha}-helix and \texorpdfstring{\ensuremath{\beta}}{beta}-sheet secondary
  structures? What forces stabilise each?
\end{enumerate}

\begin{enumerate}
\def\labelenumi{\arabic{enumi}.}
\setcounter{enumi}{4}
\tightlist
\item
  Draw a phospholipid and label the hydrophilic head and hydrophobic
  tail. What does ``amphipathic'' mean?
\end{enumerate}

\begin{enumerate}
\def\labelenumi{\arabic{enumi}.}
\setcounter{enumi}{5}
\tightlist
\item
  Distinguish saturated from unsaturated fatty acids. How does the
  degree of saturation affect membrane fluidity?
\end{enumerate}

\begin{enumerate}
\def\labelenumi{\arabic{enumi}.}
\setcounter{enumi}{6}
\tightlist
\item
  Write the general formula for a monosaccharide. Name three
  monosaccharides important in biology.
\end{enumerate}

\begin{enumerate}
\def\labelenumi{\arabic{enumi}.}
\setcounter{enumi}{7}
\tightlist
\item
  What is the difference between \texorpdfstring{\ensuremath{\alpha}}{alpha}- and \texorpdfstring{\ensuremath{\beta}}{beta}-glycosidic bonds? Give one
  example of each.
\end{enumerate}

\begin{enumerate}
\def\labelenumi{\arabic{enumi}.}
\setcounter{enumi}{8}
\tightlist
\item
  Define a nucleotide and name its three components. Which nucleotides
  carry energy?
\end{enumerate}

\begin{enumerate}
\def\labelenumi{\arabic{enumi}.}
\setcounter{enumi}{9}
\tightlist
\item
  What reagents are used to test for: (a) reducing sugars; (b) starch;
  (c) protein; (d) lipids?
\end{enumerate}

\subsection*{Questions 11--20: Application and
Analysis}\label{questions-1120-application-and-analysis-2}
\addcontentsline{toc}{subsection}{Questions 11--20: Application and
Analysis}

\begin{enumerate}
\def\labelenumi{\arabic{enumi}.}
\setcounter{enumi}{10}
\tightlist
\item
  Cellulose and starch are both glucose polymers, yet starch is
  digestible by humans and cellulose is not. Explain this difference in
  terms of glycosidic bond type (\texorpdfstring{\ensuremath{\alpha}}{alpha}-1,4 vs \texorpdfstring{\ensuremath{\beta}}{beta}-1,4), three-dimensional
  structure, and the specific enzyme (cellulase) that humans lack.
\end{enumerate}

\begin{enumerate}
\def\labelenumi{\arabic{enumi}.}
\setcounter{enumi}{11}
\tightlist
\item
  A protein has Tm (melting temperature) of 55°C in the presence of 8 M
  urea and 70°C in the absence of denaturant. Identify which types of
  bonds are disrupted by urea (that are absent from primary structure),
  and predict whether a disulphide-bonded protein would be denatured by
  urea alone.
\end{enumerate}

\begin{enumerate}
\def\labelenumi{\arabic{enumi}.}
\setcounter{enumi}{12}
\tightlist
\item
  Collagen is a fibrous protein with a repeating Gly-X-Y motif (X =
  often Pro; Y = often Hyp). Explain: (a) why Gly must occur every third
  residue; (b) what modification produces hydroxyproline (Hyp) and which
  vitamin is required; (c) what disease results from Hyp deficiency.
\end{enumerate}

\begin{enumerate}
\def\labelenumi{\arabic{enumi}.}
\setcounter{enumi}{13}
\tightlist
\item
  Phospholipids form bilayers in water. Explain: (a) the thermodynamic
  basis for bilayer formation (hydrophobic effect); (b) why cholesterol
  increases membrane rigidity at high temperatures but increases
  fluidity at low temperatures; (c) how membrane asymmetry (different
  inner vs outer leaflet composition) is maintained.
\end{enumerate}

\begin{enumerate}
\def\labelenumi{\arabic{enumi}.}
\setcounter{enumi}{14}
\tightlist
\item
  Compare the structures of DNA and RNA --- list at least four
  differences (sugar, bases, strand number, stability). Explain why RNA
  is less stable than DNA chemically (relate to the 2'-hydroxyl group).
\end{enumerate}

\begin{enumerate}
\def\labelenumi{\arabic{enumi}.}
\setcounter{enumi}{15}
\tightlist
\item
  Calculate the number of possible distinct dipeptides from 20 standard
  amino acids. Then calculate the number of different 100-residue
  proteins possible. What does this astronomical number imply for
  protein sequence space and directed evolution?
\end{enumerate}

\begin{enumerate}
\def\labelenumi{\arabic{enumi}.}
\setcounter{enumi}{16}
\tightlist
\item
  Haemoglobin (quaternary structure, four subunits) loses cooperative O\texorpdfstring{\textsubscript{2}}{2}
  binding when disaggregated into monomers. Explain how quaternary
  structure enables cooperativity (Hill coefficient n \textgreater{} 1),
  and why this is physiologically advantageous.
\end{enumerate}

\begin{enumerate}
\def\labelenumi{\arabic{enumi}.}
\setcounter{enumi}{17}
\tightlist
\item
  A glycoprotein has oligosaccharide chains attached to Asn residues
  (N-linked glycosylation). Describe: (a) where in the cell
  N-glycosylation occurs; (b) what enzyme adds the glycan; (c) three
  functions of glycosylation on cell-surface proteins.
\end{enumerate}

\begin{enumerate}
\def\labelenumi{\arabic{enumi}.}
\setcounter{enumi}{18}
\tightlist
\item
  Lipoproteins (LDL, HDL) are macromolecular complexes of lipid and
  protein. Explain how their amphipathic structure allows transport of
  hydrophobic lipids in the aqueous bloodstream. Why does elevated LDL
  cholesterol increase atherosclerosis risk while elevated HDL is
  protective?
\end{enumerate}

\begin{enumerate}
\def\labelenumi{\arabic{enumi}.}
\setcounter{enumi}{19}
\tightlist
\item
  A 10-residue peptide contains one disulphide bond between Cys
  residues. How many possible disulphide bond arrangements are there?
  Under what cellular conditions do disulphide bonds form (redox
  potential), and why are they rare in intracellular proteins?
\end{enumerate}

\subsection*{Questions 21--30: Synthesis and
Evaluation}\label{questions-2130-synthesis-and-evaluation-2}
\addcontentsline{toc}{subsection}{Questions 21--30: Synthesis and
Evaluation}

\begin{enumerate}
\def\labelenumi{\arabic{enumi}.}
\setcounter{enumi}{20}
\tightlist
\item
  Prions are infectious proteins (misfolded PrP\^{}Sc) that propagate by
  templating the misfolding of normal PrP\^{}C. Evaluate the hypothesis
  that protein conformation can carry infectivity without nucleic acid:
  what evidence supports it, and what challenges remain for explaining
  how a protein can act as an infectious agent?
\end{enumerate}

\begin{enumerate}
\def\labelenumi{\arabic{enumi}.}
\setcounter{enumi}{21}
\tightlist
\item
  Intrinsically disordered proteins (IDPs) lack stable tertiary
  structure under physiological conditions but are biologically
  functional. Evaluate the traditional sequence-structure-function
  paradigm in light of IDPs. What advantages might disorder confer for
  signalling and regulatory proteins?
\end{enumerate}

\begin{enumerate}
\def\labelenumi{\arabic{enumi}.}
\setcounter{enumi}{22}
\tightlist
\item
  Silk fibroin is a \texorpdfstring{\ensuremath{\beta}}{beta}-sheet-rich protein with extraordinary material
  strength (tensile strength 0.5 GPa). Explain, at the inter-molecular
  level, why \texorpdfstring{\ensuremath{\beta}}{beta}-sheet stacking produces this strength. Compare to
  collagen's triple helix and synthetic Kevlar (also a \texorpdfstring{\ensuremath{\beta}}{beta}-sheet polymer).
  Design a biotechnology experiment to engineer silk with improved
  elasticity by incorporating glycine-rich loops.
\end{enumerate}

\begin{enumerate}
\def\labelenumi{\arabic{enumi}.}
\setcounter{enumi}{23}
\tightlist
\item
  Deep-sea bacteria living at 110°C require heat-stable enzymes
  (thermozymes). Propose three molecular-level strategies these
  organisms use to maintain enzyme structure at extreme temperatures
  (consider ion pairs, proline residues, hydrophobic packing, and
  disulphide bonds). How might you engineer a mesophilic enzyme for
  thermostability using directed evolution?
\end{enumerate}

\begin{enumerate}
\def\labelenumi{\arabic{enumi}.}
\setcounter{enumi}{24}
\tightlist
\item
  Glycolipids in the outer leaflet of the plasma membrane act as
  cell-surface recognition molecules (blood type antigens, pathogen
  receptors). Evaluate the claim that carbohydrate diversity makes
  carbohydrates better suited than proteins for cell-surface diversity.
  Consider information content per unit mass (number of possible
  trisaccharides vs tripeptides from their respective monomer sets).
\end{enumerate}

\begin{enumerate}
\def\labelenumi{\arabic{enumi}.}
\setcounter{enumi}{25}
\tightlist
\item
  Amyloid plaques (\texorpdfstring{\ensuremath{\beta}}{beta}-sheet-rich aggregates of A\texorpdfstring{\ensuremath{\beta}}{beta} peptide) are associated
  with Alzheimer's disease. Evaluate the ``amyloid cascade hypothesis''
  citing: (a) the evidence for A\texorpdfstring{\ensuremath{\beta}}{beta} oligomers (not fibrils) as the toxic
  species; (b) why therapeutic targeting of A\texorpdfstring{\ensuremath{\beta}}{beta} aggregation has largely
  failed clinically; (c) an alternative hypothesis involving tau
  protein.
\end{enumerate}

\begin{enumerate}
\def\labelenumi{\arabic{enumi}.}
\setcounter{enumi}{26}
\tightlist
\item
  Design a protein engineering strategy to convert an enzyme that uses
  NADH as cofactor to one that uses NADPH. Identify the specific
  residues in the coenzyme-binding pocket that differ between
  NADH-specific and NADPH-specific enzymes, and describe what mutations
  you would introduce. What selection pressure would you use in directed
  evolution?
\end{enumerate}

\begin{enumerate}
\def\labelenumi{\arabic{enumi}.}
\setcounter{enumi}{27}
\tightlist
\item
  RNA aptamers can fold into three-dimensional structures that
  specifically bind proteins (like antibodies but made of RNA). Compare
  RNA and protein as molecular recognition scaffolds: evaluate the
  structural repertoire (functional groups available, backbone rigidity,
  available secondary structures) of each, and explain why protein
  antibodies have been preferred over RNA aptamers clinically, and what
  advantages aptamers offer in terms of synthesis and immunogenicity.
\end{enumerate}

\begin{enumerate}
\def\labelenumi{\arabic{enumi}.}
\setcounter{enumi}{28}
\tightlist
\item
  Polyhydroxyalkanoates (PHAs) are bacterial polyesters (biological
  plastics). Compare their biosynthesis pathway and polymer chemistry to
  triacylglycerides (TAGs) --- which also function as carbon storage
  molecules in animals. Evaluate whether PHAs could replace
  petroleum-derived plastics industrially, considering degradation
  rates, production cost, and material properties.
\end{enumerate}

\begin{enumerate}
\def\labelenumi{\arabic{enumi}.}
\setcounter{enumi}{29}
\tightlist
\item
  The hydrophobic core of most globular proteins contains residues with
  about 15--30\% exposed surface area. Propose and critically evaluate
  three evolutionary forces that have driven the formation of the
  hydrophobic core: (a) thermodynamic stability; (b) protection of
  reactive centres; (c) evolutionary constraint on folding kinetics. Use
  a specific example (e.g., cytochrome c, lysozyme) to ground your
  argument.
\end{enumerate}

\newpage

\section*{Questions --- Ch 4: Enzymes and the Kinetics of
Catalysis}\label{questions-ch-4-enzymes-and-the-kinetics-of-catalysis}
\addcontentsline{toc}{section}{Questions --- Ch 4: Enzymes and the
Kinetics of Catalysis}

\label{sec:q_unit_I_enzymes_and_kinetics}

\subsection*{Instructor Use and Coverage
Notes}\label{instructor-use-and-coverage-notes-6}
\addcontentsline{toc}{subsection}{Instructor Use and Coverage Notes}

\begin{itemize}
\tightlist
\item
  \textbf{Coverage target:} Infer Vmax, Km, or inhibition type from data
  and justify the mechanistic interpretation.
\item
  \textbf{Model/data emphasis:} Michaelis-Menten, Lineweaver-Burk, and
  inhibition-pattern calculations.
\item
  \textbf{Assessment alignment:} Concept Explanation, Statistical Tests
  and Data Analysis, Argumentation.
\item
  \textbf{Misconception probe:} A catalyst changes rate, not the
  equilibrium constant or the sign of delta G.
\item
  \textbf{Transfer product:} Apply enzyme-kinetic reasoning to drug
  dosing, metabolic control, or diagnostic assays.
\item
  \textbf{Grading focus:} award full credit for mechanism, evidence,
  boundary conditions, and units when a calculation is required.
\item
  \textbf{Suggested use:} draw one recall item, one application item,
  and one synthesis item when building a short quiz from this bank.
\end{itemize}

\subsection*{Questions 1--10: Recall and
Comprehension}\label{questions-110-recall-and-comprehension-3}
\addcontentsline{toc}{subsection}{Questions 1--10: Recall and
Comprehension}

\emph{This activity accompanies \cref{sec:unit_I_enzymes_and_kinetics}
of the textbook --- review that chapter before attempting the exercises
below.}

\begin{enumerate}
\def\labelenumi{\arabic{enumi}.}
\tightlist
\item
  Define an enzyme. What type of macromolecule are most enzymes? Name
  one exception.
\end{enumerate}

\begin{enumerate}
\def\labelenumi{\arabic{enumi}.}
\setcounter{enumi}{1}
\tightlist
\item
  What is the active site? Explain the induced fit model.
\end{enumerate}

\begin{enumerate}
\def\labelenumi{\arabic{enumi}.}
\setcounter{enumi}{2}
\tightlist
\item
  Define activation energy (Ea). How does an enzyme reduce Ea without
  altering \texorpdfstring{\ensuremath{\Delta}}{Delta}G of the reaction?
\end{enumerate}

\begin{enumerate}
\def\labelenumi{\arabic{enumi}.}
\setcounter{enumi}{3}
\tightlist
\item
  What are Michaelis-Menten kinetics? Define Km and Vmax.
\end{enumerate}

\begin{enumerate}
\def\labelenumi{\arabic{enumi}.}
\setcounter{enumi}{4}
\tightlist
\item
  Write the Michaelis-Menten equation. At \hyperref[s]{S} = Km, what
  fraction of Vmax is achieved?
\end{enumerate}

\begin{enumerate}
\def\labelenumi{\arabic{enumi}.}
\setcounter{enumi}{5}
\tightlist
\item
  Distinguish competitive, non-competitive, and uncompetitive
  inhibition. Which type changes Km? Which changes Vmax?
\end{enumerate}

\begin{enumerate}
\def\labelenumi{\arabic{enumi}.}
\setcounter{enumi}{6}
\tightlist
\item
  What is a coenzyme? Name two vitamins that serve as coenzymes.
\end{enumerate}

\begin{enumerate}
\def\labelenumi{\arabic{enumi}.}
\setcounter{enumi}{7}
\tightlist
\item
  What is an allosteric enzyme? How does cooperativity affect the v vs
  \hyperref[s]{S} curve shape?
\end{enumerate}

\begin{enumerate}
\def\labelenumi{\arabic{enumi}.}
\setcounter{enumi}{8}
\tightlist
\item
  Define feedback inhibition (product inhibition) and give one example
  from a metabolic pathway.
\end{enumerate}

\begin{enumerate}
\def\labelenumi{\arabic{enumi}.}
\setcounter{enumi}{9}
\tightlist
\item
  How does temperature affect enzyme activity? Name two effects of
  extreme temperatures on enzyme function.
\end{enumerate}

\subsection*{Questions 11--20: Application and
Analysis}\label{questions-1120-application-and-analysis-3}
\addcontentsline{toc}{subsection}{Questions 11--20: Application and
Analysis}

\begin{enumerate}
\def\labelenumi{\arabic{enumi}.}
\setcounter{enumi}{10}
\tightlist
\item
  An enzyme has Km = 2 mM and Vmax = 10 \texorpdfstring{\ensuremath{\mu}}{mu}mol/min. Calculate the reaction
  rate at \hyperref[s]{S} = 1 mM, 2 mM, and 10 mM. What does the
  saturation behaviour tell you about enzyme kinetics?
\end{enumerate}

\begin{enumerate}
\def\labelenumi{\arabic{enumi}.}
\setcounter{enumi}{11}
\tightlist
\item
  In a competitive inhibition experiment, Km increases from 2 mM to 8
  mM, but Vmax is unchanged (10 \texorpdfstring{\ensuremath{\mu}}{mu}mol/min). Calculate the inhibitor
  constant Ki if \hyperref[i]{I} = 50 \texorpdfstring{\ensuremath{\mu}}{mu}M (use the relationship
  K\_m\^{}app = Km(1 + \hyperref[i]{I}/Ki)).
\end{enumerate}

\begin{enumerate}
\def\labelenumi{\arabic{enumi}.}
\setcounter{enumi}{12}
\tightlist
\item
  ATCase (aspartate transcarbamoylase) shows sigmoidal kinetics (Hill
  coefficient n = 2) rather than hyperbolic Michaelis-Menten kinetics.
  Explain what the Hill coefficient indicates about subunit
  cooperativity and how CTP (the end product) acts as a heterotropic
  allosteric inhibitor.
\end{enumerate}

\begin{enumerate}
\def\labelenumi{\arabic{enumi}.}
\setcounter{enumi}{13}
\tightlist
\item
  Temperature doubles reaction rates approximately every 10°C (Q\texorpdfstring{\textsubscript{1}}{1}\texorpdfstring{\textsubscript{0}}{0} \texorpdfstring{\ensuremath{\approx}}{~} 2).
  However, above an optimal temperature, enzyme rate drops sharply.
  Explain: (a) why the Arrhenius relationship predicts increasing rate
  with temperature; (b) why denaturation overrides this above Topt; (c)
  how thermophilic enzymes (Tm \textgreater{} 80°C) avoid denaturation.
\end{enumerate}

\begin{enumerate}
\def\labelenumi{\arabic{enumi}.}
\setcounter{enumi}{14}
\tightlist
\item
  A Lineweaver-Burk (double-reciprocal) plot for enzyme X gives a
  y-intercept of 0.1 \texorpdfstring{\ensuremath{\mu}}{mu}mol\texorpdfstring{\textsuperscript{-}}{-}\texorpdfstring{\textsuperscript{1}}{1}\texorpdfstring{\ensuremath{\cdot}}{.}min and x-intercept of \texorpdfstring{\ensuremath{-}}{-}0.5 mM\texorpdfstring{\textsuperscript{-}}{-}\texorpdfstring{\textsuperscript{1}}{1}. Calculate
  Vmax and Km. If an inhibitor shifts the x-intercept to \texorpdfstring{\ensuremath{-}}{-}0.25 mM\texorpdfstring{\textsuperscript{-}}{-}\texorpdfstring{\textsuperscript{1}}{1}
  without changing the y-intercept, what type of inhibition is this?
\end{enumerate}

\begin{enumerate}
\def\labelenumi{\arabic{enumi}.}
\setcounter{enumi}{15}
\tightlist
\item
  Aconitase, an enzyme in the TCA cycle, uses Fe\texorpdfstring{\textsuperscript{2}}{2}\texorpdfstring{\textsuperscript{+}}{+} as a cofactor. When
  \emph{E. coli} is iron-starved, aconitase activity is abolished.
  Explain: (a) the likely role of Fe\texorpdfstring{\textsuperscript{2}}{2}\texorpdfstring{\textsuperscript{+}}{+} in the active site mechanism; (b)
  why iron starvation globally affects energy metabolism; (c) how some
  bacteria produce iron-chelating siderophores to restore enzyme
  function.
\end{enumerate}

\begin{enumerate}
\def\labelenumi{\arabic{enumi}.}
\setcounter{enumi}{16}
\tightlist
\item
  Isoenzymes (isozymes) are distinct enzymes that catalyse the same
  reaction but have different kinetic properties. Hexokinase
  (Km\_glucose = 0.1 mM) and glucokinase/hexokinase IV (Km\_glucose = 10
  mM) both catalyse glucose phosphorylation. Explain why the liver
  expresses high-Km glucokinase rather than low-Km hexokinase and how
  this enables hepatic glucose buffering post-meal.
\end{enumerate}

\begin{enumerate}
\def\labelenumi{\arabic{enumi}.}
\setcounter{enumi}{17}
\tightlist
\item
  Serine proteases (trypsin, chymotrypsin, elastase) use the catalytic
  triad Ser-His-Asp. Explain the mechanism: how does the His residue act
  as a proton acceptor/donor, how does Ser become a nucleophile, and how
  does the oxyanion hole stabilise the tetrahedral intermediate?
\end{enumerate}

\begin{enumerate}
\def\labelenumi{\arabic{enumi}.}
\setcounter{enumi}{18}
\tightlist
\item
  Ribozymes are RNA enzymes that catalyse phosphodiester bond cleavage.
  Compare ribozyme catalytic mechanisms to protein enzyme mechanisms:
  what functional group strategies can RNA use that proteins cannot (and
  vice versa)? What does the existence of ribozymes tell us about the
  origin of life?
\end{enumerate}

\begin{enumerate}
\def\labelenumi{\arabic{enumi}.}
\setcounter{enumi}{19}
\tightlist
\item
  Suicide (mechanism-based) inhibitors permanently inactivate enzymes by
  forming covalent bonds after being processed by the enzyme's active
  site. Give a specific clinical example of a suicide inhibitor and
  explain, step-by-step, how it inactivates the target enzyme. Why are
  suicide inhibitors more selective than reversible inhibitors?
\end{enumerate}

\subsection*{Questions 21--30: Synthesis and
Evaluation}\label{questions-2130-synthesis-and-evaluation-3}
\addcontentsline{toc}{subsection}{Questions 21--30: Synthesis and
Evaluation}

\begin{enumerate}
\def\labelenumi{\arabic{enumi}.}
\setcounter{enumi}{20}
\tightlist
\item
  Aspirin (acetylsalicylate) irreversibly inhibits COX-1 and COX-2 by
  acetylating Ser530 in the active site. Evaluate: (a) why aspirin's
  anti-platelet effect lasts 7--10 days (the lifetime of a platelet);
  (b) how selective COX-2 inhibitors (celecoxib) were designed to reduce
  GI side effects; (c) the unexpected cardiovascular risk of COX-2
  inhibitors detected primarily after clinical deployment.
\end{enumerate}

\begin{enumerate}
\def\labelenumi{\arabic{enumi}.}
\setcounter{enumi}{21}
\tightlist
\item
  Design an enzyme engineering experiment to improve the Km of an
  industrial enzyme for a non-natural substrate. Describe: the rational
  design approach using structural data; the directed evolution
  approach; how you would screen large variant libraries; and how you
  would verify that the improved Km reflects genuine active-site
  affinity changes vs changes in protein stability.
\end{enumerate}

\begin{enumerate}
\def\labelenumi{\arabic{enumi}.}
\setcounter{enumi}{22}
\tightlist
\item
  Enzyme kinetics assumes a well-mixed, dilute environment. However,
  intracellular enzymes operate in a crowded, compartmentalised
  environment with local \hyperref[s]{S} quite different from bulk
  cytoplasm. Evaluate three ways that molecular crowding alters
  effective Km and Vmax in vivo vs in vitro measurements, and propose an
  experimental strategy to measure kinetics in a crowded environment.
\end{enumerate}

\begin{enumerate}
\def\labelenumi{\arabic{enumi}.}
\setcounter{enumi}{23}
\tightlist
\item
  Zymogen activation (e.g., pepsinogen \texorpdfstring{\ensuremath{\to}}{->} pepsin) is a mechanism for
  controlling enzyme activity. Evaluate the physiological logic of
  synthesising inactive proenzymes, using three specific examples from
  digestion and blood coagulation. What happens when zymogen activation
  goes wrong (e.g., pancreatitis)?
\end{enumerate}

\begin{enumerate}
\def\labelenumi{\arabic{enumi}.}
\setcounter{enumi}{24}
\tightlist
\item
  Methotrexate inhibits dihydrofolate reductase (DHFR) by competitive
  inhibition. In cancer chemotherapy, it exploits rapidly dividing
  cells' high demand for folate-derived nucleotides. Evaluate: (a) why
  cancer cells cannot simply overcome competitive inhibition by
  upregulating substrate concentration; (b) why resistance to
  methotrexate often involves DHFR gene amplification; (c) the trade-off
  between therapeutic index and resistance evolution.
\end{enumerate}

\begin{enumerate}
\def\labelenumi{\arabic{enumi}.}
\setcounter{enumi}{25}
\tightlist
\item
  PCSK9 (proprotein convertase subtilisin/kexin type 9) is a serine
  protease that degrades LDL receptors, raising LDL cholesterol.
  Monoclonal antibody inhibitors of PCSK9 dramatically lower LDL.
  Evaluate why an enzyme inhibitor (rather than a simple receptor
  agonist) is necessary here, and explain how the
  enzyme-substrate-inhibitor kinetics differ from conventional
  small-molecule competitive inhibitors targeting PCSK9's active site.
\end{enumerate}

\begin{enumerate}
\def\labelenumi{\arabic{enumi}.}
\setcounter{enumi}{26}
\tightlist
\item
  Enzymes that catalyse the same reaction in different organisms often
  have different optimal pH values (e.g., pepsin pH 2 vs trypsin pH 8).
  Propose three molecular reasons why the optimal pH of an enzyme is
  tuned to its cellular compartment, and describe how you would
  experimentally shift the pH optimum of pepsin toward neutral pH by
  rational mutation.
\end{enumerate}

\begin{enumerate}
\def\labelenumi{\arabic{enumi}.}
\setcounter{enumi}{27}
\tightlist
\item
  Single-molecule enzyme kinetics (SMEK) experiments reveal that
  individual enzyme molecules show fluctuating activity --- some are
  momentarily fast, some slow --- a phenomenon called ``dynamic
  disorder.'' Evaluate what SMEK data reveals about the assumptions
  underlying the traditional Michaelis-Menten model, and propose what
  these fluctuations imply for drug design targeting enzymes in disease.
\end{enumerate}

\begin{enumerate}
\def\labelenumi{\arabic{enumi}.}
\setcounter{enumi}{28}
\tightlist
\item
  Metabolic channelling occurs when the product of enzyme A is directly
  transferred to enzyme B without diffusing into the bulk solution
  (e.g., tryptophan synthase, CAD complex, purinosome). Evaluate the
  advantages of metabolic channelling over free diffusion for: (a)
  reaction efficiency; (b) protection of unstable intermediates; (c)
  prevention of inhibitory product accumulation. Describe one
  experimental test that distinguishes channelling from free diffusion.
\end{enumerate}

\begin{enumerate}
\def\labelenumi{\arabic{enumi}.}
\setcounter{enumi}{29}
\tightlist
\item
  Criticise the following hypothesis: ``The optimal enzyme evolves to
  have a Km equal to the in vivo substrate concentration.'' Using
  evolutionary arguments, real enzyme kinetic data, and the concept of
  ``catalytic perfection'' (kcat/Km approaching the diffusion limit
  about 10\texorpdfstring{\textsuperscript{8}}{8}--10\texorpdfstring{\textsuperscript{9}}{9} M\texorpdfstring{\textsuperscript{-}}{-}\texorpdfstring{\textsuperscript{1}}{1}s\texorpdfstring{\textsuperscript{-}}{-}\texorpdfstring{\textsuperscript{1}}{1}), evaluate whether enzymes are actually
  at evolutionary optimum, or whether constraints (protein stability,
  pleiotropic effects) prevent optimisation. Cite specific examples.
\end{enumerate}

\newpage

\section*{Questions --- Ch 5: Cell Theory and Cell
Types}\label{questions-ch-5-cell-theory-and-cell-types}
\addcontentsline{toc}{section}{Questions --- Ch 5: Cell Theory and Cell
Types}

\label{sec:q_unit_II_cell_theory}

\subsection*{Instructor Use and Coverage
Notes}\label{instructor-use-and-coverage-notes-7}
\addcontentsline{toc}{subsection}{Instructor Use and Coverage Notes}

\begin{itemize}
\tightlist
\item
  \textbf{Coverage target:} Use size or microscopy evidence to classify
  a cell and predict a constraint on exchange.
\item
  \textbf{Model/data emphasis:} Surface-area-to-volume scaling.
\item
  \textbf{Assessment alignment:} Visual Representations, Questions and
  Methods, Argumentation.
\item
  \textbf{Misconception probe:} Cells are not just small bags of fluid;
  boundaries and internal organization create function.
\item
  \textbf{Transfer product:} Transfer scale reasoning to eggs, neurons,
  root hairs, and microbial colonies.
\item
  \textbf{Grading focus:} award full credit for mechanism, evidence,
  boundary conditions, and units when a calculation is required.
\item
  \textbf{Suggested use:} draw one recall item, one application item,
  and one synthesis item when building a short quiz from this bank.
\end{itemize}

\subsection*{Questions 1--10: Recall and
Comprehension}\label{questions-110-recall-and-comprehension-4}
\addcontentsline{toc}{subsection}{Questions 1--10: Recall and
Comprehension}

\emph{This activity accompanies \cref{sec:unit_II_cell_theory} of the
textbook --- review that chapter before attempting the exercises below.}

\begin{enumerate}
\def\labelenumi{\arabic{enumi}.}
\tightlist
\item
  State the three principles of cell theory and name the three
  scientists who established it.
\end{enumerate}

\begin{enumerate}
\def\labelenumi{\arabic{enumi}.}
\setcounter{enumi}{1}
\tightlist
\item
  List five structural differences between a prokaryote and a eukaryote.
\end{enumerate}

\begin{enumerate}
\def\labelenumi{\arabic{enumi}.}
\setcounter{enumi}{2}
\tightlist
\item
  What is the smallest known living cell? What is the largest known
  single cell?
\end{enumerate}

\begin{enumerate}
\def\labelenumi{\arabic{enumi}.}
\setcounter{enumi}{3}
\tightlist
\item
  What is the endomembrane system? List most organelles that belong to
  it.
\end{enumerate}

\begin{enumerate}
\def\labelenumi{\arabic{enumi}.}
\setcounter{enumi}{4}
\tightlist
\item
  Define resolution in microscopy. What is the resolution limit of a
  light microscope?
\end{enumerate}

\begin{enumerate}
\def\labelenumi{\arabic{enumi}.}
\setcounter{enumi}{5}
\tightlist
\item
  What is the cell theory exception posed by viruses? Are viruses
  ``alive''?
\end{enumerate}

\begin{enumerate}
\def\labelenumi{\arabic{enumi}.}
\setcounter{enumi}{6}
\tightlist
\item
  Describe the inner membrane of a mitochondrion. What are cristae and
  why are they important?
\end{enumerate}

\begin{enumerate}
\def\labelenumi{\arabic{enumi}.}
\setcounter{enumi}{7}
\tightlist
\item
  What is the nucleolus and what is assembled there?
\end{enumerate}

\begin{enumerate}
\def\labelenumi{\arabic{enumi}.}
\setcounter{enumi}{8}
\tightlist
\item
  Distinguish cell walls of bacteria, plants, and fungi (composition).
\end{enumerate}

\begin{enumerate}
\def\labelenumi{\arabic{enumi}.}
\setcounter{enumi}{9}
\tightlist
\item
  What is meant by ``compartmentalisation'' and why is it advantageous
  for eukaryotes?
\end{enumerate}

\subsection*{Questions 11--20: Application and
Analysis}\label{questions-1120-application-and-analysis-4}
\addcontentsline{toc}{subsection}{Questions 11--20: Application and
Analysis}

\begin{enumerate}
\def\labelenumi{\arabic{enumi}.}
\setcounter{enumi}{10}
\tightlist
\item
  Calculate the surface area-to-volume ratio (SA:V) for spherical cells
  with radii of 1 \texorpdfstring{\ensuremath{\mu}}{mu}m, 5 \texorpdfstring{\ensuremath{\mu}}{mu}m, and 10 \texorpdfstring{\ensuremath{\mu}}{mu}m. Explain why SA:V limits the
  maximum size of cells.
\end{enumerate}

\begin{enumerate}
\def\labelenumi{\arabic{enumi}.}
\setcounter{enumi}{11}
\tightlist
\item
  Colchicine destroys microtubules. Predict its effects on: (a) mitosis;
  (b) vesicle trafficking; (c) cilia; (d) cell shape.
\end{enumerate}

\begin{enumerate}
\def\labelenumi{\arabic{enumi}.}
\setcounter{enumi}{12}
\tightlist
\item
  \emph{Thiomargarita namibiensis} is a bacterium \textgreater750 \texorpdfstring{\ensuremath{\mu}}{mu}m in
  diameter --- larger than many eukaryotes. How does it solve the SA:V
  problem?
\end{enumerate}

\begin{enumerate}
\def\labelenumi{\arabic{enumi}.}
\setcounter{enumi}{13}
\tightlist
\item
  Describe the experimental basis for the discovery that cells arise
  from pre-existing cells (Pasteur's swan-neck flask experiment).
\end{enumerate}

\begin{enumerate}
\def\labelenumi{\arabic{enumi}.}
\setcounter{enumi}{14}
\tightlist
\item
  Explain five lines of molecular evidence supporting the endosymbiotic
  origin of mitochondria.
\end{enumerate}

\begin{enumerate}
\def\labelenumi{\arabic{enumi}.}
\setcounter{enumi}{15}
\tightlist
\item
  A secretory cell is treated with brefeldin A (disrupts the Golgi).
  Predict the consequences for enzyme secretion.
\end{enumerate}

\begin{enumerate}
\def\labelenumi{\arabic{enumi}.}
\setcounter{enumi}{16}
\tightlist
\item
  Explain why red blood cells lack a nucleus and mitochondria. What
  metabolic consequences does this create?
\end{enumerate}

\begin{enumerate}
\def\labelenumi{\arabic{enumi}.}
\setcounter{enumi}{17}
\tightlist
\item
  A mutation eliminates the peroxisomal targeting signal (PTS1) from
  catalase. Where will catalase end up? What biochemical consequence
  results?
\end{enumerate}

\begin{enumerate}
\def\labelenumi{\arabic{enumi}.}
\setcounter{enumi}{18}
\tightlist
\item
  Ribosome composition differs between prokaryotes (70S = 50S + 30S) and
  eukaryotes (80S = 60S + 40S). How are antibiotics exploited to target
  primarily bacterial ribosomes?
\end{enumerate}

\begin{enumerate}
\def\labelenumi{\arabic{enumi}.}
\setcounter{enumi}{19}
\tightlist
\item
  A student uses STED super-resolution microscopy and finds that nuclear
  pore complexes are distributed non-randomly on the nuclear envelope,
  clustering near transcriptionally active chromatin. Interpret this
  finding in terms of mRNA export and nuclear architecture.
\end{enumerate}

\subsection*{Questions 21--30: Synthesis and
Evaluation}\label{questions-2130-synthesis-and-evaluation-4}
\addcontentsline{toc}{subsection}{Questions 21--30: Synthesis and
Evaluation}

\begin{enumerate}
\def\labelenumi{\arabic{enumi}.}
\setcounter{enumi}{20}
\tightlist
\item
  Evaluate the claim: ``Most life evolved from a single comprehensive
  common ancestor (LUCA).'' What molecular evidence supports LUCA, and
  what evidence (e.g., horizontal gene transfer) complicates the simple
  tree-of-life narrative?
\end{enumerate}

\begin{enumerate}
\def\labelenumi{\arabic{enumi}.}
\setcounter{enumi}{21}
\tightlist
\item
  Synthetic biology aims to build minimal artificial cells (protocells).
  Evaluate three fundamental hurdles: (a) stable lipid encapsulation;
  (b) genetic self-replication; (c) energy supply. At what stage are
  current (2024--2025) efforts?
\end{enumerate}

\begin{enumerate}
\def\labelenumi{\arabic{enumi}.}
\setcounter{enumi}{22}
\tightlist
\item
  Organelles proliferate or contract depending on cellular demand.
  Design an experiment to determine whether smooth ER proliferation in
  response to lipid overload is controlled by transcriptional
  upregulation of ER resident proteins or by a different mechanism
  (e.g., membrane lipid incorporation itself driving ER expansion).
\end{enumerate}

\begin{enumerate}
\def\labelenumi{\arabic{enumi}.}
\setcounter{enumi}{23}
\tightlist
\item
  Ageing cells accumulate dysfunctional mitochondria. Evaluate the
  ``mitochondrial theory of ageing'': does mitochondrial dysfunction
  cause ageing, or result from it? Cite evidence from caloric
  restriction and NAD+ supplementation studies.
\end{enumerate}

\begin{enumerate}
\def\labelenumi{\arabic{enumi}.}
\setcounter{enumi}{24}
\tightlist
\item
  Some unicellular organisms (e.g., \emph{Euglena}) contain both
  mitochondria and chloroplasts. Evaluate how a single cell coordinates
  two endosymbiotic organelles with conflicting demands on the cell's
  redox state.
\end{enumerate}

\begin{enumerate}
\def\labelenumi{\arabic{enumi}.}
\setcounter{enumi}{25}
\tightlist
\item
  Evaluate the role of the centrosome as a signalling hub beyond its
  microtubule-organising function. What kinases localise to the
  centrosome, and how does centrosome duplication failure contribute to
  cancer?
\end{enumerate}

\begin{enumerate}
\def\labelenumi{\arabic{enumi}.}
\setcounter{enumi}{26}
\tightlist
\item
  Apicoplasts in malaria parasites (\emph{Plasmodium}) are
  non-photosynthetic plastids derived from a secondary endosymbiosis.
  Evaluate them as drug targets --- what essential biosynthetic pathways
  occur exclusively in the apicoplast, and why does targeting apicoplast
  translation (e.g., with doxycycline) kill the parasite but spare human
  cells?
\end{enumerate}

\begin{enumerate}
\def\labelenumi{\arabic{enumi}.}
\setcounter{enumi}{27}
\tightlist
\item
  Liquid-liquid phase separation (LLPS) creates membraneless organelles
  (stress granules, P-bodies). Evaluate the evidence that these
  condensates are functionally equivalent to membrane-bounded organelles
  in concentrating reaction components.
\end{enumerate}

\begin{enumerate}
\def\labelenumi{\arabic{enumi}.}
\setcounter{enumi}{28}
\tightlist
\item
  Compare the evolutionary trajectories of a symbiotic bacterium
  becoming a mitochondrion vs becoming a lysosome. What structural and
  genomic changes would you expect in each pathway?
\end{enumerate}

\begin{enumerate}
\def\labelenumi{\arabic{enumi}.}
\setcounter{enumi}{29}
\tightlist
\item
  Critically assess the hypothesis that cancer is a disease of cell
  identity --- a reversion to a primordial unicellular state where
  individual cell fitness (unlimited proliferation, resistance to death)
  overrides multicellular cooperation. What evidence from cancer
  evolutionary biology supports or contradicts this view?
\end{enumerate}

\newpage

\section*{Questions --- Ch 6: Cell Structure and
Organelles}\label{questions-ch-6-cell-structure-and-organelles}
\addcontentsline{toc}{section}{Questions --- Ch 6: Cell Structure and
Organelles}

\label{sec:q_unit_II_cell_structure}

\subsection*{Instructor Use and Coverage
Notes}\label{instructor-use-and-coverage-notes-8}
\addcontentsline{toc}{subsection}{Instructor Use and Coverage Notes}

\begin{itemize}
\tightlist
\item
  \textbf{Coverage target:} Predict which organelle or cytoskeletal
  process is disrupted from a phenotype.
\item
  \textbf{Model/data emphasis:} Compartment-flow and motor-transport
  calculations.
\item
  \textbf{Assessment alignment:} Visual Representations, Questions and
  Methods, Argumentation.
\item
  \textbf{Misconception probe:} Organelles are dynamic process hubs, not
  static textbook icons.
\item
  \textbf{Transfer product:} Apply compartment logic to secretion,
  apoptosis, infection, or cell division.
\item
  \textbf{Grading focus:} award full credit for mechanism, evidence,
  boundary conditions, and units when a calculation is required.
\item
  \textbf{Suggested use:} draw one recall item, one application item,
  and one synthesis item when building a short quiz from this bank.
\end{itemize}

\subsection*{Questions 1--10: Recall and
Comprehension}\label{questions-110-recall-and-comprehension-5}
\addcontentsline{toc}{subsection}{Questions 1--10: Recall and
Comprehension}

\emph{This activity accompanies \cref{sec:unit_II_cell_structure} of the
textbook --- review that chapter before attempting the exercises below.}

\begin{enumerate}
\def\labelenumi{\arabic{enumi}.}
\tightlist
\item
  Which organelle is responsible for: (a) protein synthesis; (b) lipid
  synthesis; (c) protein glycosylation; (d) intracellular digestion; (e)
  ATP synthesis?
\end{enumerate}

\begin{enumerate}
\def\labelenumi{\arabic{enumi}.}
\setcounter{enumi}{1}
\tightlist
\item
  What is the difference between rough ER and smooth ER in structure and
  function?
\end{enumerate}

\begin{enumerate}
\def\labelenumi{\arabic{enumi}.}
\setcounter{enumi}{2}
\tightlist
\item
  Describe the cis-to-trans polarity of the Golgi apparatus and the
  direction of cargo flow.
\end{enumerate}

\begin{enumerate}
\def\labelenumi{\arabic{enumi}.}
\setcounter{enumi}{3}
\tightlist
\item
  What is the optimal pH inside a lysosome and why?
\end{enumerate}

\begin{enumerate}
\def\labelenumi{\arabic{enumi}.}
\setcounter{enumi}{4}
\tightlist
\item
  Name the three components of the cytoskeleton and one function of
  each.
\end{enumerate}

\begin{enumerate}
\def\labelenumi{\arabic{enumi}.}
\setcounter{enumi}{5}
\tightlist
\item
  What are cilia and flagella made of (protein and structure)? What
  drives their beating?
\end{enumerate}

\begin{enumerate}
\def\labelenumi{\arabic{enumi}.}
\setcounter{enumi}{6}
\tightlist
\item
  Describe the nuclear envelope. How many membranes does it have?
\end{enumerate}

\begin{enumerate}
\def\labelenumi{\arabic{enumi}.}
\setcounter{enumi}{7}
\tightlist
\item
  What two reactions does the peroxisome catalyse involving hydrogen
  peroxide?
\end{enumerate}

\begin{enumerate}
\def\labelenumi{\arabic{enumi}.}
\setcounter{enumi}{8}
\tightlist
\item
  Distinguish the large central vacuole of plants from contractile
  vacuoles of protists.
\end{enumerate}

\begin{enumerate}
\def\labelenumi{\arabic{enumi}.}
\setcounter{enumi}{9}
\tightlist
\item
  What is the centrosome? What happens to centrosomes during mitosis?
\end{enumerate}

\subsection*{Questions 11--20: Application and
Analysis}\label{questions-1120-application-and-analysis-5}
\addcontentsline{toc}{subsection}{Questions 11--20: Application and
Analysis}

\begin{enumerate}
\def\labelenumi{\arabic{enumi}.}
\setcounter{enumi}{10}
\tightlist
\item
  Trace the complete path of a secretory protein from its gene to the
  extracellular space, naming every organelle and vesicle type it passes
  through.
\end{enumerate}

\begin{enumerate}
\def\labelenumi{\arabic{enumi}.}
\setcounter{enumi}{11}
\tightlist
\item
  A mutation disables the signal recognition particle (SRP). Predict the
  fate of a lysosomal enzyme without functional SRP.
\end{enumerate}

\begin{enumerate}
\def\labelenumi{\arabic{enumi}.}
\setcounter{enumi}{12}
\tightlist
\item
  Colchicine prevents microtubule polymerisation. Would it disrupt: (a)
  chromosome segregation; (b) ER-Golgi trafficking; (c) cell migration;
  (d) axonal transport? Justify each.
\end{enumerate}

\begin{enumerate}
\def\labelenumi{\arabic{enumi}.}
\setcounter{enumi}{13}
\tightlist
\item
  Gaucher disease results from a deficiency of the lysosomal enzyme
  glucocerebrosidase. What accumulates, in which cells, and why does
  this cause splenomegaly and bone marrow infiltration?
\end{enumerate}

\begin{enumerate}
\def\labelenumi{\arabic{enumi}.}
\setcounter{enumi}{14}
\tightlist
\item
  If a cell's mitochondria were removed, which metabolic pathways would
  be abolished and which would remain functional? Calculate the maximum
  ATP yield from 1 glucose in a cell without mitochondria.
\end{enumerate}

\begin{enumerate}
\def\labelenumi{\arabic{enumi}.}
\setcounter{enumi}{15}
\tightlist
\item
  Clathrin-coated vesicles mediate receptor-mediated endocytosis.
  Describe the sequence from receptor-ligand binding to delivery of
  cargo to endosomes. Name two other coat protein systems and their
  functions.
\end{enumerate}

\begin{enumerate}
\def\labelenumi{\arabic{enumi}.}
\setcounter{enumi}{16}
\tightlist
\item
  Motor proteins kinesin (moves to + end) and dynein (moves to \texorpdfstring{\ensuremath{-}}{-} end)
  both use microtubule tracks. How does a cell ensure organelles are
  delivered to the correct destination --- axon terminal vs cell body?
\end{enumerate}

\begin{enumerate}
\def\labelenumi{\arabic{enumi}.}
\setcounter{enumi}{17}
\tightlist
\item
  The rough ER membrane is studded with ribosomes, while the smooth ER
  is ribosome-free. A mutation eliminates the ribosome-binding site on
  the ER translocon (Sec61). Predict the consequence for: (a) membrane
  protein insertion; (b) secretory protein signal peptide cleavage.
\end{enumerate}

\begin{enumerate}
\def\labelenumi{\arabic{enumi}.}
\setcounter{enumi}{18}
\tightlist
\item
  Zellweger syndrome (peroxisome biogenesis disorder) causes progressive
  neurological deterioration. Identify two specific metabolic functions
  of peroxisomes in neurons that would be disrupted, and explain the
  mechanism of neural damage.
\end{enumerate}

\begin{enumerate}
\def\labelenumi{\arabic{enumi}.}
\setcounter{enumi}{19}
\tightlist
\item
  A researcher fuses a nuclear localisation signal (NLS) to a
  cytoplasmic protein (GFP). Describe the molecular events leading to
  the protein's nuclear import, including the roles of importin-\texorpdfstring{\ensuremath{\alpha}}{alpha},
  importin-\texorpdfstring{\ensuremath{\beta}}{beta}, RanGTP, and the nuclear pore complex.
\end{enumerate}

\subsection*{Questions 21--30: Synthesis and
Evaluation}\label{questions-2130-synthesis-and-evaluation-5}
\addcontentsline{toc}{subsection}{Questions 21--30: Synthesis and
Evaluation}

\begin{enumerate}
\def\labelenumi{\arabic{enumi}.}
\setcounter{enumi}{20}
\tightlist
\item
  Primary cilia are now recognised as signalling organelles for
  Hedgehog, Wnt, and PDGFR pathways. Evaluate why concentrating
  receptors in a small ciliary compartment (0.1\% of total cell surface)
  enhances signalling sensitivity compared to distributing receptors
  over the entire plasma membrane.
\end{enumerate}

\begin{enumerate}
\def\labelenumi{\arabic{enumi}.}
\setcounter{enumi}{21}
\tightlist
\item
  Design an experiment to determine whether a newly discovered organelle
  in a marine protist is related to the peroxisome or the lysosome,
  using: (a) proteomic analysis; (b) membrane lipid analysis; (c)
  fluorescent marker trafficking; (d) inhibitor sensitivity.
\end{enumerate}

\begin{enumerate}
\def\labelenumi{\arabic{enumi}.}
\setcounter{enumi}{22}
\tightlist
\item
  Autophagy (macro, micro, chaperone-mediated) degrades cytosolic
  material. Evaluate the selective vs non-selective aspects of each
  type, and propose why cells might preferentially use CMA
  (chaperone-mediated autophagy) rather than macroautophagy for
  long-lived soluble proteins.
\end{enumerate}

\begin{enumerate}
\def\labelenumi{\arabic{enumi}.}
\setcounter{enumi}{23}
\tightlist
\item
  Reticulophagy (selective ER autophagy) is triggered by ER stress.
  Trace the molecular pathway from misfolded protein accumulation \texorpdfstring{\ensuremath{\to}}{->}
  IRE1\texorpdfstring{\ensuremath{\alpha}}{alpha}/PERK/ATF6 activation \texorpdfstring{\ensuremath{\to}}{->} reticulophagy receptor induction \texorpdfstring{\ensuremath{\to}}{->}
  ER-phagy. What would happen if reticulophagy were completely blocked?
\end{enumerate}

\begin{enumerate}
\def\labelenumi{\arabic{enumi}.}
\setcounter{enumi}{24}
\tightlist
\item
  The Golgi apparatus disperses into ministacks during mitosis and
  reassembles in each daughter cell after division. Evaluate the
  mechanism of Golgi fragmentation (role of mitotic kinases) and
  reassembly (role of p97 ATPase), and explain why Golgi inheritance
  must be coordinated with chromosome segregation.
\end{enumerate}

\begin{enumerate}
\def\labelenumi{\arabic{enumi}.}
\setcounter{enumi}{25}
\tightlist
\item
  Compare endocytosis mechanisms: clathrin-mediated endocytosis,
  caveolae-mediated endocytosis, macropinocytosis, and phagocytosis. For
  each: what cargo is internalised, what signalling events initiate it,
  and how is cargo processed after internalisation?
\end{enumerate}

\begin{enumerate}
\def\labelenumi{\arabic{enumi}.}
\setcounter{enumi}{26}
\tightlist
\item
  COPI-coated vesicles mediate retrograde transport from Golgi back to
  ER (retrieval of escaped ER-resident proteins). How do KDEL-bearing
  ER-resident proteins get retrieved? What is the KDEL receptor, and
  what happens to its affinity for KDEL at Golgi pH vs ER pH?
\end{enumerate}

\begin{enumerate}
\def\labelenumi{\arabic{enumi}.}
\setcounter{enumi}{27}
\tightlist
\item
  The nuclear lamina (lamin intermediate filaments) underlies the inner
  nuclear membrane. Mutations in lamin A cause laminopathies including
  Hutchinson-Gilford progeria syndrome (HGPS) --- a premature ageing
  disease. Evaluate how nuclear structural failure leads to an
  accelerated-ageing phenotype at the cellular level.
\end{enumerate}

\begin{enumerate}
\def\labelenumi{\arabic{enumi}.}
\setcounter{enumi}{28}
\tightlist
\item
  Cells in developing lungs must convert from secreting surfactant
  (lipid-rich, packaged in lamellar bodies --- lysosome-related
  organelles) to absorbing alveolar fluid at birth. Evaluate the
  signalling mechanisms controlling this switch and what organelle
  remodelling is required.
\end{enumerate}

\begin{enumerate}
\def\labelenumi{\arabic{enumi}.}
\setcounter{enumi}{29}
\tightlist
\item
  Evaluate the concept of ``organelle crosstalk'': mitochondria-ER
  contact sites (MAM --- mitochondria-associated membranes) facilitate
  Ca\texorpdfstring{\textsuperscript{2}}{2}\texorpdfstring{\textsuperscript{+}}{+} and lipid exchange. How does disruption of MAM contact alter: (a)
  mitochondrial Ca\texorpdfstring{\textsuperscript{2}}{2}\texorpdfstring{\textsuperscript{+}}{+} buffering; (b) phosphatidylserine transfer; (c)
  mitochondrial morphology? Cite evidence from neurodegenerative disease
  models.
\end{enumerate}

\newpage

\section*{Questions --- Ch 7: Membrane Structure and
Transport}\label{questions-ch-7-membrane-structure-and-transport}
\addcontentsline{toc}{section}{Questions --- Ch 7: Membrane Structure
and Transport}

\label{sec:q_unit_II_membrane_transport}

\subsection*{Instructor Use and Coverage
Notes}\label{instructor-use-and-coverage-notes-9}
\addcontentsline{toc}{subsection}{Instructor Use and Coverage Notes}

\begin{itemize}
\tightlist
\item
  \textbf{Coverage target:} Calculate or predict ion movement and
  explain which force dominates.
\item
  \textbf{Model/data emphasis:} Nernst, Goldman, osmotic, and
  facilitated-transport calculations.
\item
  \textbf{Assessment alignment:} Visual Representations, Questions and
  Methods, Argumentation.
\item
  \textbf{Misconception probe:} Equilibrium does not mean equal
  concentration; charge and permeability matter.
\item
  \textbf{Transfer product:} Transfer gradient logic to neurons,
  kidneys, roots, and mitochondrial membranes.
\item
  \textbf{Grading focus:} award full credit for mechanism, evidence,
  boundary conditions, and units when a calculation is required.
\item
  \textbf{Suggested use:} draw one recall item, one application item,
  and one synthesis item when building a short quiz from this bank.
\end{itemize}

\subsection*{Questions 1--10: Recall and
Comprehension}\label{questions-110-recall-and-comprehension-6}
\addcontentsline{toc}{subsection}{Questions 1--10: Recall and
Comprehension}

\emph{This activity accompanies \cref{sec:unit_II_membrane_transport} of
the textbook --- review that chapter before attempting the exercises
below.}

\begin{enumerate}
\def\labelenumi{\arabic{enumi}.}
\tightlist
\item
  Describe the fluid-mosaic model of the plasma membrane
  (Singer-Nicolson, 1972). What are the two major components?
\end{enumerate}

\begin{enumerate}
\def\labelenumi{\arabic{enumi}.}
\setcounter{enumi}{1}
\tightlist
\item
  What is a phospholipid bilayer? Explain why it forms spontaneously in
  an aqueous environment.
\end{enumerate}

\begin{enumerate}
\def\labelenumi{\arabic{enumi}.}
\setcounter{enumi}{2}
\tightlist
\item
  Define selective permeability. Which types of molecules cross
  membranes freely and which require assistance?
\end{enumerate}

\begin{enumerate}
\def\labelenumi{\arabic{enumi}.}
\setcounter{enumi}{3}
\tightlist
\item
  Distinguish simple diffusion, facilitated diffusion, and active
  transport.
\end{enumerate}

\begin{enumerate}
\def\labelenumi{\arabic{enumi}.}
\setcounter{enumi}{4}
\tightlist
\item
  What is the sodium-potassium ATPase (Na\texorpdfstring{\textsuperscript{+}}{+}/K\texorpdfstring{\textsuperscript{+}}{+}-ATPase)? How many ions
  does it move per ATP?
\end{enumerate}

\begin{enumerate}
\def\labelenumi{\arabic{enumi}.}
\setcounter{enumi}{5}
\tightlist
\item
  Define osmosis. In which direction does water move --- toward high or
  low solute concentration?
\end{enumerate}

\begin{enumerate}
\def\labelenumi{\arabic{enumi}.}
\setcounter{enumi}{6}
\tightlist
\item
  What is a membrane potential? What is the typical resting membrane
  potential of a mammalian cell?
\end{enumerate}

\begin{enumerate}
\def\labelenumi{\arabic{enumi}.}
\setcounter{enumi}{7}
\tightlist
\item
  Distinguish primary from secondary active transport. Give one example
  of each.
\end{enumerate}

\begin{enumerate}
\def\labelenumi{\arabic{enumi}.}
\setcounter{enumi}{8}
\tightlist
\item
  What is endocytosis? Distinguish phagocytosis, pinocytosis, and
  receptor-mediated endocytosis.
\end{enumerate}

\begin{enumerate}
\def\labelenumi{\arabic{enumi}.}
\setcounter{enumi}{9}
\tightlist
\item
  What is the role of cholesterol in membrane structure? How does it
  affect fluidity at high and low temperatures?
\end{enumerate}

\subsection*{Questions 11--20: Application and
Analysis}\label{questions-1120-application-and-analysis-6}
\addcontentsline{toc}{subsection}{Questions 11--20: Application and
Analysis}

\begin{enumerate}
\def\labelenumi{\arabic{enumi}.}
\setcounter{enumi}{10}
\tightlist
\item
  A cell is placed in a hypertonic solution. Describe, at the molecular
  level, what happens to: (a) the direction of water movement; (b) cell
  volume; (c) the cytoskeleton if shrinkage is severe; (d) membrane
  potential due to increased ion concentration.
\end{enumerate}

\begin{enumerate}
\def\labelenumi{\arabic{enumi}.}
\setcounter{enumi}{11}
\tightlist
\item
  The Na\texorpdfstring{\textsuperscript{+}}{+}/K\texorpdfstring{\textsuperscript{+}}{+}-ATPase pumps 3 Na\texorpdfstring{\textsuperscript{+}}{+} out and 2 K\texorpdfstring{\textsuperscript{+}}{+} in per ATP hydrolysed.
  Explain why this creates an electrogenic current. Calculate the charge
  moved per pump cycle and predict how ouabain (a pump inhibitor) would
  affect membrane potential.
\end{enumerate}

\begin{enumerate}
\def\labelenumi{\arabic{enumi}.}
\setcounter{enumi}{12}
\tightlist
\item
  Glucose enters intestinal epithelial cells via SGLT1 (Na\texorpdfstring{\textsuperscript{+}}{+}-glucose
  co-transporter) on the apical membrane and exits via GLUT2 (passive
  facilitated diffusion) on the basolateral membrane. Explain: (a) which
  transport is primary active and which is secondary active; (b) why
  this arrangement allows glucose to be absorbed against its
  concentration gradient.
\end{enumerate}

\begin{enumerate}
\def\labelenumi{\arabic{enumi}.}
\setcounter{enumi}{13}
\tightlist
\item
  CFTR is a Cl\texorpdfstring{\textsuperscript{-}}{-} channel mutated in cystic fibrosis. The \texorpdfstring{\ensuremath{\Delta}}{Delta}F508 deletion
  causes CFTR misfolding and ER retention. Explain how this affects: (a)
  epithelial Cl\texorpdfstring{\textsuperscript{-}}{-} secretion; (b) mucus viscosity; (c) bacterial
  colonisation of airways.
\end{enumerate}

\begin{enumerate}
\def\labelenumi{\arabic{enumi}.}
\setcounter{enumi}{14}
\tightlist
\item
  A cell biologist measures the fluorescence recovery after
  photobleaching (FRAP) of a fluorescently labelled membrane protein.
  Recovery takes 30 seconds in normal cells but 5 minutes in cells with
  disrupted actin cortex. Interpret this result using the fluid-mosaic
  model and the concept of membrane protein diffusion constraints.
\end{enumerate}

\begin{enumerate}
\def\labelenumi{\arabic{enumi}.}
\setcounter{enumi}{15}
\tightlist
\item
  Explain how the membrane potential arises from differences in ion
  permeability (Goldman equation). If a cell becomes 10\texorpdfstring{\ensuremath{\times}}{x} more permeable
  to Na\texorpdfstring{\textsuperscript{+}}{+} (e.g., during an action potential), predict the direction of
  change in membrane potential and apply the Nernst equation to estimate
  the equilibrium potential for Na\texorpdfstring{\textsuperscript{+}}{+}.
\end{enumerate}

\begin{enumerate}
\def\labelenumi{\arabic{enumi}.}
\setcounter{enumi}{16}
\tightlist
\item
  Vesicle fusion requires SNARE proteins. Explain the zipper model of
  SNARE complex assembly and how this drives membrane fusion. Why does
  botulinum toxin (which cleaves SNARE proteins) block neurotransmitter
  release?
\end{enumerate}

\begin{enumerate}
\def\labelenumi{\arabic{enumi}.}
\setcounter{enumi}{17}
\tightlist
\item
  Aquaporin-1 (AQP1) is absent in red blood cells from certain
  individuals, yet these individuals are healthy. AQP1 knockout mice
  have impaired urinary concentrating ability. Explain this discrepancy
  --- what does it reveal about the relative importance of aquaporin vs
  lipid bilayer water permeability in different tissues?
\end{enumerate}

\begin{enumerate}
\def\labelenumi{\arabic{enumi}.}
\setcounter{enumi}{18}
\tightlist
\item
  LDL receptor-mediated endocytosis concentrates LDL receptors in
  clathrin-coated pits. FH (familial hypercholesterolaemia) patients
  lack functional LDL receptors. Explain: (a) why plasma LDL rises
  dramatically; (b) why heterozygous FH is less severe than homozygous;
  (c) why statins partially compensate for LDL receptor deficiency.
\end{enumerate}

\begin{enumerate}
\def\labelenumi{\arabic{enumi}.}
\setcounter{enumi}{19}
\tightlist
\item
  Membrane asymmetry: phosphatidylserine (PS) is normally on the inner
  leaflet but is exposed on the outer leaflet of apoptotic cells.
  Explain: (a) the enzyme (flippase) that normally maintains PS on the
  inner leaflet; (b) the enzyme (scramblase) that destroys asymmetry
  during apoptosis; (c) the physiological consequence of outer-leaflet
  PS exposure (phagocyte recognition).
\end{enumerate}

\subsection*{Questions 21--30: Synthesis and
Evaluation}\label{questions-2130-synthesis-and-evaluation-6}
\addcontentsline{toc}{subsection}{Questions 21--30: Synthesis and
Evaluation}

\begin{enumerate}
\def\labelenumi{\arabic{enumi}.}
\setcounter{enumi}{20}
\tightlist
\item
  Evaluate the thermodynamic basis of lipid raft formation. What causes
  certain lipids (sphingomyelin, cholesterol, saturated phospholipids)
  to phase-separate into liquid-ordered domains? Evaluate the evidence
  that lipid rafts are real functional domains vs artefacts of cold
  detergent extraction.
\end{enumerate}

\begin{enumerate}
\def\labelenumi{\arabic{enumi}.}
\setcounter{enumi}{21}
\tightlist
\item
  The membrane potential (typically \texorpdfstring{\ensuremath{-}}{-}70 mV, inside negative) represents
  stored electrochemical energy. Calculate the energy available per mole
  of Na\texorpdfstring{\textsuperscript{+}}{+} flowing inward across a neuron membrane using \texorpdfstring{\ensuremath{\Delta}}{Delta}G = zF\texorpdfstring{\ensuremath{\Delta}}{Delta}\texorpdfstring{\ensuremath{\Psi}}{Psi} + RT
  ln({[}Na\texorpdfstring{\textsuperscript{+}}{+}{]}\_in/{[}Na\texorpdfstring{\textsuperscript{+}}{+}{]}\_out). How does this compare to the energy
  of ATP hydrolysis (about \texorpdfstring{\ensuremath{-}}{-}57 kJ/mol at physiological
  concentrations)?
\end{enumerate}

\begin{enumerate}
\def\labelenumi{\arabic{enumi}.}
\setcounter{enumi}{22}
\tightlist
\item
  Ionophores are small molecules that carry specific ions across lipid
  bilayers. Valinomycin is a K\texorpdfstring{\textsuperscript{+}}{+}-selective ionophore; gramicidin is a
  non-selective cation channel. Compare their effects on: (a) membrane
  potential; (b) mitochondrial function; (c) cell viability at
  therapeutic vs toxic concentrations. How are they used experimentally
  as tools?
\end{enumerate}

\begin{enumerate}
\def\labelenumi{\arabic{enumi}.}
\setcounter{enumi}{23}
\tightlist
\item
  Membrane curvature is actively generated during endocytosis, vesicle
  budding, and tubule formation. Evaluate the mechanisms: (a) BAR domain
  protein scaffolding; (b) asymmetric lipid composition; (c) protein
  insertion into the cytoplasmic leaflet. Which mechanism is most
  critical for clathrin-mediated endocytosis at the cell's flat plasma
  membrane?
\end{enumerate}

\begin{enumerate}
\def\labelenumi{\arabic{enumi}.}
\setcounter{enumi}{24}
\tightlist
\item
  Osmotic stress responses in bacteria involve membrane sensors (e.g.,
  EnvZ osmosensor) that activate two-component signalling systems.
  Compare bacterial osmosensing to mammalian cell osmosensing via
  volume-sensitive ion channels (VSOR). Which system responds faster and
  why?
\end{enumerate}

\begin{enumerate}
\def\labelenumi{\arabic{enumi}.}
\setcounter{enumi}{25}
\tightlist
\item
  Fatty acid composition of membrane phospholipids changes with
  temperature (homeoviscous adaptation). Design an experiment to
  determine whether \emph{E. coli} adjusts membrane fluidity in response
  to temperature via: (a) transcriptional regulation of fatty acid
  desaturases; (b) selective incorporation of exogenous fatty acids; or
  (c) lipid flip-flop remodelling. What lipid analysis technique would
  you use?
\end{enumerate}

\begin{enumerate}
\def\labelenumi{\arabic{enumi}.}
\setcounter{enumi}{26}
\tightlist
\item
  P-glycoprotein (MDR1) is an ABC transporter that pumps hydrophobic
  drugs out of cancer cells, causing multidrug resistance. Evaluate: (a)
  why hydrophobic drugs are its substrates; (b) why inhibiting MDR1 has
  largely failed clinically; (c) an alternative strategy (e.g., prodrug
  activation inside the cell) that bypasses MDR1.
\end{enumerate}

\begin{enumerate}
\def\labelenumi{\arabic{enumi}.}
\setcounter{enumi}{27}
\tightlist
\item
  The mitochondrial inner membrane maintains a proton gradient (\texorpdfstring{\ensuremath{\Delta}}{Delta}\texorpdfstring{\ensuremath{\mu}}{mu}H\texorpdfstring{\textsuperscript{+}}{+} =
  \texorpdfstring{\ensuremath{-}}{-}200 mV, \texorpdfstring{\ensuremath{\Delta}}{Delta}pH = 1 unit = \texorpdfstring{\ensuremath{-}}{-}59 mV, electrical = \texorpdfstring{\ensuremath{-}}{-}141 mV). Using the
  chemiosmotic equation \texorpdfstring{\ensuremath{\Delta}}{Delta}G = nF\texorpdfstring{\ensuremath{\Delta}}{Delta}\texorpdfstring{\ensuremath{\Psi}}{Psi} + 2.303RT\texorpdfstring{\ensuremath{\Delta}}{Delta}pH, calculate the total
  proton-motive force. How does this compare to the energy needed to
  synthesise ATP from ADP + Pi (about +37 kJ/mol at
  physiological {[}ATP{]}/{[}ADP{]})? Assume 3 H\texorpdfstring{\textsuperscript{+}}{+} per ATP (c-ring H\texorpdfstring{\textsuperscript{+}}{+}/ATP
  coupling ratio).
\end{enumerate}

\begin{enumerate}
\def\labelenumi{\arabic{enumi}.}
\setcounter{enumi}{28}
\tightlist
\item
  Nanoparticle drug delivery systems exploit endocytosis for
  intracellular drug delivery. Evaluate the barriers nanoparticles face:
  (a) being taken up via endocytosis rather than direct membrane fusion;
  (b) endosomal escape before lysosomal degradation; (c) targeting
  specificity via surface ligands. What design features address each
  barrier?
\end{enumerate}

\begin{enumerate}
\def\labelenumi{\arabic{enumi}.}
\setcounter{enumi}{29}
\tightlist
\item
  Stretch-activated ion channels (SACs) open when the plasma membrane is
  mechanically deformed. Design a complete research programme to: (a)
  identify the molecular identity of a SAC in cardiac myocytes; (b)
  determine its role in cardiac arrhythmia during myocardial stretch;
  (c) develop a selective SAC inhibitor and evaluate its anti-arrhythmic
  potential in a relevant animal model.
\end{enumerate}

\newpage

\section*{Questions --- Ch 8: Cell Signalling and
Communication}\label{questions-ch-8-cell-signalling-and-communication}
\addcontentsline{toc}{section}{Questions --- Ch 8: Cell Signalling and
Communication}

\label{sec:q_unit_II_cell_signaling}

\subsection*{Instructor Use and Coverage
Notes}\label{instructor-use-and-coverage-notes-10}
\addcontentsline{toc}{subsection}{Instructor Use and Coverage Notes}

\begin{itemize}
\tightlist
\item
  \textbf{Coverage target:} Predict the outcome of perturbing one
  pathway component and justify the causal chain.
\item
  \textbf{Model/data emphasis:} Ligand-receptor occupancy and
  dose-response reasoning.
\item
  \textbf{Assessment alignment:} Visual Representations, Questions and
  Methods, Argumentation.
\item
  \textbf{Misconception probe:} A pathway diagram is a causal model, not
  a memorization chart.
\item
  \textbf{Transfer product:} Apply signaling logic to hormones,
  neurotransmitters, immune receptors, or cancer mutations.
\item
  \textbf{Grading focus:} award full credit for mechanism, evidence,
  boundary conditions, and units when a calculation is required.
\item
  \textbf{Suggested use:} draw one recall item, one application item,
  and one synthesis item when building a short quiz from this bank.
\end{itemize}

\subsection*{Questions 1--10: Recall and
Comprehension}\label{questions-110-recall-and-comprehension-7}
\addcontentsline{toc}{subsection}{Questions 1--10: Recall and
Comprehension}

\emph{This activity accompanies \cref{sec:unit_II_cell_signaling} of the
textbook --- review that chapter before attempting the exercises below.}

\begin{enumerate}
\def\labelenumi{\arabic{enumi}.}
\tightlist
\item
  List the three stages of cell signalling. Give a brief description of
  each stage.
\end{enumerate}

\begin{enumerate}
\def\labelenumi{\arabic{enumi}.}
\setcounter{enumi}{1}
\tightlist
\item
  Name four chemical classes of signalling molecules. Give one
  biological example of each.
\end{enumerate}

\begin{enumerate}
\def\labelenumi{\arabic{enumi}.}
\setcounter{enumi}{2}
\tightlist
\item
  Distinguish autocrine, paracrine, endocrine, and contact-dependent
  signalling.
\end{enumerate}

\begin{enumerate}
\def\labelenumi{\arabic{enumi}.}
\setcounter{enumi}{3}
\tightlist
\item
  What is a G protein-coupled receptor (GPCR)? Name the three subunits
  of the heterotrimeric G protein.
\end{enumerate}

\begin{enumerate}
\def\labelenumi{\arabic{enumi}.}
\setcounter{enumi}{4}
\tightlist
\item
  Describe the second messengers cAMP and IP\texorpdfstring{\textsubscript{3}}{3}/DAG and name the enzymes
  that produce each.
\end{enumerate}

\begin{enumerate}
\def\labelenumi{\arabic{enumi}.}
\setcounter{enumi}{5}
\tightlist
\item
  What is a receptor tyrosine kinase (RTK)? Name two growth factor
  receptors of this type.
\end{enumerate}

\begin{enumerate}
\def\labelenumi{\arabic{enumi}.}
\setcounter{enumi}{6}
\tightlist
\item
  Define signal amplification. Give a numerical example from a GPCR
  cascade.
\end{enumerate}

\begin{enumerate}
\def\labelenumi{\arabic{enumi}.}
\setcounter{enumi}{7}
\tightlist
\item
  What is a MAP kinase cascade? List the three kinases in the Ras-MAPK
  pathway.
\end{enumerate}

\begin{enumerate}
\def\labelenumi{\arabic{enumi}.}
\setcounter{enumi}{8}
\tightlist
\item
  Define receptor downregulation. What cellular mechanism drives it?
\end{enumerate}

\begin{enumerate}
\def\labelenumi{\arabic{enumi}.}
\setcounter{enumi}{9}
\tightlist
\item
  What is the role of protein phosphatases in signal termination?
\end{enumerate}

\subsection*{Questions 11--20: Application and
Analysis}\label{questions-1120-application-and-analysis-7}
\addcontentsline{toc}{subsection}{Questions 11--20: Application and
Analysis}

\begin{enumerate}
\def\labelenumi{\arabic{enumi}.}
\setcounter{enumi}{10}
\tightlist
\item
  Adrenaline binds \texorpdfstring{\ensuremath{\beta}}{beta}-adrenergic receptors (Gs-coupled). Trace the
  complete signalling cascade from ligand binding to glycogen
  phosphorylase activation, naming every molecule. How many molecules of
  glycogen phosphorylase are activated per initial adrenaline molecule?
\end{enumerate}

\begin{enumerate}
\def\labelenumi{\arabic{enumi}.}
\setcounter{enumi}{11}
\tightlist
\item
  Steroid hormones (e.g., cortisol) exert their effects via
  intracellular receptors that act as transcription factors. Compare the
  timescale, reversibility, and specificity of steroid signalling vs
  adrenaline signalling via GPCR. Which is more suitable for rapid
  physiological responses, and why?
\end{enumerate}

\begin{enumerate}
\def\labelenumi{\arabic{enumi}.}
\setcounter{enumi}{12}
\tightlist
\item
  Explain how Ca\texorpdfstring{\textsuperscript{2}}{2}\texorpdfstring{\textsuperscript{+}}{+} acts as a second messenger. What keeps resting
  cytoplasmic {[}Ca\texorpdfstring{\textsuperscript{2}}{2}\texorpdfstring{\textsuperscript{+}}{+}{]} so low (about 100 nM)? What is the
  role of calmodulin as a Ca\texorpdfstring{\textsuperscript{2}}{2}\texorpdfstring{\textsuperscript{+}}{+} sensor, and which kinase does
  Ca\texorpdfstring{\textsuperscript{2}}{2}\texorpdfstring{\textsuperscript{+}}{+}-calmodulin activate?
\end{enumerate}

\begin{enumerate}
\def\labelenumi{\arabic{enumi}.}
\setcounter{enumi}{13}
\tightlist
\item
  Trastuzumab (Herceptin) is a monoclonal antibody that targets HER2 (an
  RTK overexpressed in about 25\% of breast cancers). Explain
  its mechanism of action: how does receptor blockade inhibit downstream
  MAPK and PI3K-Akt signalling, and why does HER2 dimerisation (with
  HER3) make it particularly potent?
\end{enumerate}

\begin{enumerate}
\def\labelenumi{\arabic{enumi}.}
\setcounter{enumi}{14}
\tightlist
\item
  An inhibitor blocks PLC-\texorpdfstring{\ensuremath{\beta}}{beta} (phospholipase C-\texorpdfstring{\ensuremath{\beta}}{beta}). Predict the effects on:
  (a) IP\texorpdfstring{\textsubscript{3}}{3} production; (b) DAG levels; (c) PKC activation; (d)
  intracellular Ca\texorpdfstring{\textsuperscript{2}}{2}\texorpdfstring{\textsuperscript{+}}{+} release from the ER; (e) smooth muscle contraction
  downstream of M3 muscarinic receptor activation.
\end{enumerate}

\begin{enumerate}
\def\labelenumi{\arabic{enumi}.}
\setcounter{enumi}{15}
\tightlist
\item
  Cross-talk between signal pathways allows integration of multiple
  inputs. Describe two specific examples where activation of one
  signalling pathway inhibits or potentiates the response of another
  pathway. How might this cross-talk produce distinct cellular outcomes
  depending on context?
\end{enumerate}

\begin{enumerate}
\def\labelenumi{\arabic{enumi}.}
\setcounter{enumi}{16}
\tightlist
\item
  Wnt signalling inhibits the destruction complex (APC/Axin/CK1/GSK3\texorpdfstring{\ensuremath{\beta}}{beta}),
  allowing \texorpdfstring{\ensuremath{\beta}}{beta}-catenin to accumulate and enter the nucleus. Explain how
  mutations in APC (found in \textgreater80\% of colorectal cancers)
  constitutively activate Wnt signalling and promote cell proliferation.
\end{enumerate}

\begin{enumerate}
\def\labelenumi{\arabic{enumi}.}
\setcounter{enumi}{17}
\tightlist
\item
  Nitric oxide (NO) is produced by nitric oxide synthase (NOS) and
  diffuses freely across membranes to activate soluble guanylate cyclase
  (sGC) in smooth muscle cells, producing cGMP. Explain: (a) why NO is
  unusual as a paracrine signal; (b) how sildenafil (Viagra) potentiates
  the NO-cGMP pathway; (c) what happens to cGMP without PDE5 inhibition.
\end{enumerate}

\begin{enumerate}
\def\labelenumi{\arabic{enumi}.}
\setcounter{enumi}{18}
\tightlist
\item
  Notch signalling requires direct cell-cell contact (juxtacrine). When
  Notch receptor (on receiving cell) binds Delta ligand (on sending
  cell), \texorpdfstring{\ensuremath{\gamma}}{gamma}-secretase cleaves the Notch intracellular domain (NICD),
  which enters the nucleus. Explain: (a) why this signalling cannot be
  mediated by a diffusible messenger; (b) the role of Notch in lateral
  inhibition during development.
\end{enumerate}

\begin{enumerate}
\def\labelenumi{\arabic{enumi}.}
\setcounter{enumi}{19}
\tightlist
\item
  TGF-\texorpdfstring{\ensuremath{\beta}}{beta} receptors (type I and type II serine/threonine kinase receptors)
  signal via SMAD proteins. Explain: (a) how receptor dimerisation
  activates SMAD2/3 phosphorylation; (b) how SMAD4 forms a complex with
  phospho-SMAD2/3; (c) how this complex regulates target gene
  transcription in the context of tumour suppression vs EMT
  (epithelial-mesenchymal transition).
\end{enumerate}

\subsection*{Questions 21--30: Synthesis and
Evaluation}\label{questions-2130-synthesis-and-evaluation-7}
\addcontentsline{toc}{subsection}{Questions 21--30: Synthesis and
Evaluation}

\begin{enumerate}
\def\labelenumi{\arabic{enumi}.}
\setcounter{enumi}{20}
\tightlist
\item
  Evaluate the concept of ``signal integration'': a cell simultaneously
  receives growth factor (MAPK activation), anchorage signals
  (integrin-FAK signalling), and survival signals (insulin \texorpdfstring{\ensuremath{\to}}{->} PI3K-Akt).
  How does the cell integrate these three pathways to make a
  proliferation vs death decision? Name at least one molecular node
  where most three signals converge.
\end{enumerate}

\begin{enumerate}
\def\labelenumi{\arabic{enumi}.}
\setcounter{enumi}{21}
\tightlist
\item
  Design a CRISPR screen to identify novel negative regulators of the
  EGFR-MAPK pathway in non-small-cell lung cancer cells. Describe the
  library design, the cellular readout, and how you would validate the
  top hits as therapeutically relevant targets.
\end{enumerate}

\begin{enumerate}
\def\labelenumi{\arabic{enumi}.}
\setcounter{enumi}{22}
\tightlist
\item
  Receptor endocytosis and downregulation reduce cell sensitivity to a
  signal. However, some endosomal signalling (e.g., from
  endosome-localised EGFR) is qualitatively different from plasma
  membrane signalling. Evaluate the concept of ``spatial encoding'' of
  signal quality by receptor location. What specific downstream
  effectors are activated preferentially from endosomal vs surface
  receptors?
\end{enumerate}

\begin{enumerate}
\def\labelenumi{\arabic{enumi}.}
\setcounter{enumi}{23}
\tightlist
\item
  Intracellular Ca\texorpdfstring{\textsuperscript{2}}{2}\texorpdfstring{\textsuperscript{+}}{+} oscillations are widely used for signal encoding.
  Evaluate how the frequency and amplitude of Ca\texorpdfstring{\textsuperscript{2}}{2}\texorpdfstring{\textsuperscript{+}}{+} oscillations can
  encode different cellular responses (NF-\texorpdfstring{\ensuremath{\kappa}}{kappa}B vs NFAT activation). What
  molecular mechanism (IP\texorpdfstring{\textsubscript{3}}{3}R sensitivity to Ca\texorpdfstring{\textsuperscript{2}}{2}\texorpdfstring{\textsuperscript{+}}{+} and inhibition by Ca\texorpdfstring{\textsuperscript{2}}{2}\texorpdfstring{\textsuperscript{+}}{+})
  generates oscillations?
\end{enumerate}

\begin{enumerate}
\def\labelenumi{\arabic{enumi}.}
\setcounter{enumi}{24}
\tightlist
\item
  Synaptic signalling relies on glutamate receptor subtypes (AMPA, NMDA,
  mGluR) with vastly different ion conductances, Ca\texorpdfstring{\textsuperscript{2}}{2}\texorpdfstring{\textsuperscript{+}}{+} permeabilities,
  and kinetics. Evaluate how the combination of fast ionotropic (AMPA)
  and slow modulatory (NMDA, mGluR) signalling enables both rapid
  synaptic transmission and long-term plasticity at the same synapse.
\end{enumerate}

\begin{enumerate}
\def\labelenumi{\arabic{enumi}.}
\setcounter{enumi}{25}
\tightlist
\item
  Drug resistance in targeted therapy often results from mutations that
  reactivate downstream signalling. In BRAF\^{}V600E melanoma treated
  with vemurafenib (BRAF inhibitor), resistance arises via: (a) RAS
  mutation bypassing BRAF; (b) MEK1 mutation making it insensitive; (c)
  RAF1 dimerisation activating MEK. Evaluate why combination therapy
  (BRAF + MEK inhibition) reduces resistance emergence.
\end{enumerate}

\begin{enumerate}
\def\labelenumi{\arabic{enumi}.}
\setcounter{enumi}{26}
\tightlist
\item
  The PI3K-Akt-mTOR pathway is a master regulator of cell growth and
  metabolism. Evaluate the consequences of complete mTOR inhibition in
  cancer vs healthy cells: why do rapalogs (mTORC1 inhibitors) produce
  cytostatic rather than cytotoxic effects in most cancers, and what
  feedback mechanisms restore Akt activity?
\end{enumerate}

\begin{enumerate}
\def\labelenumi{\arabic{enumi}.}
\setcounter{enumi}{27}
\tightlist
\item
  Feedback loops in signalling networks can generate bistability
  (switch-like most-or-none responses) or oscillations (e.g., NF-\texorpdfstring{\ensuremath{\kappa}}{kappa}B
  oscillations). Evaluate the mathematical conditions (using a two-state
  model) under which positive feedback creates bistability vs
  oscillations. Give one biological example of each.
\end{enumerate}

\begin{enumerate}
\def\labelenumi{\arabic{enumi}.}
\setcounter{enumi}{28}
\tightlist
\item
  Exosomes (30--150 nm extracellular vesicles) carry signalling
  molecules (proteins, miRNAs, lipids) between cells over long
  distances. Evaluate the evidence that exosomes constitute a bona fide
  signalling mechanism distinct from classical paracrine signalling, and
  propose experimental criteria to distinguish ``exosomal'' from
  conventional secreted-protein-mediated communication.
\end{enumerate}

\begin{enumerate}
\def\labelenumi{\arabic{enumi}.}
\setcounter{enumi}{29}
\tightlist
\item
  Critically evaluate the ``oncogene addiction'' hypothesis: cancer
  cells become ``addicted'' to a single dominant oncogene signal (e.g.,
  BCR-ABL in CML), making targeted inhibition highly effective. Under
  what conditions does addiction arise, and why do multiple oncogenic
  mutations (e.g., in pancreatic cancer) prevent effective single-agent
  targeted therapy?
\end{enumerate}

\newpage

\section*{Questions --- Ch 9: Bioenergetics and Cellular
Respiration}\label{questions-ch-9-bioenergetics-and-cellular-respiration}
\addcontentsline{toc}{section}{Questions --- Ch 9: Bioenergetics and
Cellular Respiration}

\label{sec:q_unit_III_bioenergetics_and_respiration}

\subsection*{Instructor Use and Coverage
Notes}\label{instructor-use-and-coverage-notes-11}
\addcontentsline{toc}{subsection}{Instructor Use and Coverage Notes}

\begin{itemize}
\tightlist
\item
  \textbf{Coverage target:} Calculate an energy yield or redox
  consequence and connect it to cellular work.
\item
  \textbf{Model/data emphasis:} Delta G, ATP yield, and electron-carrier
  accounting.
\item
  \textbf{Assessment alignment:} Representing and Describing Data,
  Statistical Tests and Data Analysis.
\item
  \textbf{Misconception probe:} ATP is not stored energy in a vague
  sense; it is a coupling currency with defined reaction chemistry.
\item
  \textbf{Transfer product:} Transfer energy accounting to exercise,
  fermentation, hypoxia, and mitochondrial disease.
\item
  \textbf{Grading focus:} award full credit for mechanism, evidence,
  boundary conditions, and units when a calculation is required.
\item
  \textbf{Suggested use:} draw one recall item, one application item,
  and one synthesis item when building a short quiz from this bank.
\end{itemize}

\subsection*{Questions 1--10: Recall and
Comprehension}\label{questions-110-recall-and-comprehension-8}
\addcontentsline{toc}{subsection}{Questions 1--10: Recall and
Comprehension}

\emph{This activity accompanies
\cref{sec:unit_III_bioenergetics_and_respiration} of the textbook ---
review that chapter before attempting the exercises below.}

\begin{enumerate}
\def\labelenumi{\arabic{enumi}.}
\tightlist
\item
  Define free energy (G). What does a negative \texorpdfstring{\ensuremath{\Delta}}{Delta}G indicate about a
  reaction?
\end{enumerate}

\begin{enumerate}
\def\labelenumi{\arabic{enumi}.}
\setcounter{enumi}{1}
\tightlist
\item
  Write the overall equation for aerobic cellular respiration.
\end{enumerate}

\begin{enumerate}
\def\labelenumi{\arabic{enumi}.}
\setcounter{enumi}{2}
\tightlist
\item
  What is ATP? What bonds are broken during ATP hydrolysis, and why is
  this energy release large?
\end{enumerate}

\begin{enumerate}
\def\labelenumi{\arabic{enumi}.}
\setcounter{enumi}{3}
\tightlist
\item
  Name the four stages of aerobic respiration and state where each
  occurs in the eukaryotic cell.
\end{enumerate}

\begin{enumerate}
\def\labelenumi{\arabic{enumi}.}
\setcounter{enumi}{4}
\tightlist
\item
  How many ATP are produced by substrate-level phosphorylation in
  glycolysis?
\end{enumerate}

\begin{enumerate}
\def\labelenumi{\arabic{enumi}.}
\setcounter{enumi}{5}
\tightlist
\item
  What is the purpose of the preparatory steps (pyruvate
  decarboxylation) before the Krebs cycle?
\end{enumerate}

\begin{enumerate}
\def\labelenumi{\arabic{enumi}.}
\setcounter{enumi}{6}
\tightlist
\item
  What is the electron transport chain? Name the four major protein
  complexes.
\end{enumerate}

\begin{enumerate}
\def\labelenumi{\arabic{enumi}.}
\setcounter{enumi}{7}
\tightlist
\item
  What is the chemiosmotic theory? Who proposed it and when?
\end{enumerate}

\begin{enumerate}
\def\labelenumi{\arabic{enumi}.}
\setcounter{enumi}{8}
\tightlist
\item
  What is fermentation? Name two types and the products of each.
\end{enumerate}

\begin{enumerate}
\def\labelenumi{\arabic{enumi}.}
\setcounter{enumi}{9}
\tightlist
\item
  What is the difference between substrate-level phosphorylation and
  oxidative phosphorylation?
\end{enumerate}

\subsection*{Questions 11--20: Application and
Analysis}\label{questions-1120-application-and-analysis-8}
\addcontentsline{toc}{subsection}{Questions 11--20: Application and
Analysis}

\begin{enumerate}
\def\labelenumi{\arabic{enumi}.}
\setcounter{enumi}{10}
\tightlist
\item
  Calculate the theoretical maximum ATP yield from one glucose molecule
  in aerobic respiration. How does the actual yield (about 30
  ATP) differ from the textbook maximum (about 38), and why?
\end{enumerate}

\begin{enumerate}
\def\labelenumi{\arabic{enumi}.}
\setcounter{enumi}{11}
\tightlist
\item
  In anaerobic fermentation, yeast converts pyruvate to ethanol + CO\texorpdfstring{\textsubscript{2}}{2}.
  Explain why this regeneration of NAD\texorpdfstring{\textsuperscript{+}}{+} is essential for continued
  glycolysis, even though no ATP is produced in the process.
\end{enumerate}

\begin{enumerate}
\def\labelenumi{\arabic{enumi}.}
\setcounter{enumi}{12}
\tightlist
\item
  Cyanide (CN\texorpdfstring{\textsuperscript{-}}{-}) blocks Complex IV (cytochrome c oxidase). Predict the
  consequences for: (a) proton gradient; (b) NADH/NAD\texorpdfstring{\textsuperscript{+}}{+} ratio; (c) ATP
  production; (d) O\texorpdfstring{\textsubscript{2}}{2} consumption; (e) lactic acid production.
\end{enumerate}

\begin{enumerate}
\def\labelenumi{\arabic{enumi}.}
\setcounter{enumi}{13}
\tightlist
\item
  Uncoupling proteins (UCPs) such as thermogenin in brown adipose tissue
  dissipate the proton gradient as heat rather than driving ATP
  synthesis. Explain the physiological role of this process and
  calculate: if 50\% of the proton gradient is uncoupled, what happens
  to O\texorpdfstring{\textsubscript{2}}{2} consumption vs ATP output?
\end{enumerate}

\begin{enumerate}
\def\labelenumi{\arabic{enumi}.}
\setcounter{enumi}{14}
\tightlist
\item
  Glycolysis includes an irreversible phosphofructokinase-1 (PFK-1)
  step. PFK-1 is allosterically activated by AMP. ADP, and F-2,6-BP, and
  inhibited by ATP and citrate. Explain how this regulation matches
  glycolytic rate to cellular energy status.
\end{enumerate}

\begin{enumerate}
\def\labelenumi{\arabic{enumi}.}
\setcounter{enumi}{15}
\tightlist
\item
  Brown adipose tissue (BAT) is thermogenic. Compare the metabolic
  activity of BAT vs white adipose tissue (WAT): what substrate do they
  preferentially oxidise, how do they regulate UCP1 expression
  (adrenergic \texorpdfstring{\ensuremath{\to}}{->} cAMP \texorpdfstring{\ensuremath{\to}}{->} PKA \texorpdfstring{\ensuremath{\to}}{->} lipolysis \texorpdfstring{\ensuremath{\to}}{->} FA activation of UCP1), and
  what clinical interest does BAT activation generate?
\end{enumerate}

\begin{enumerate}
\def\labelenumi{\arabic{enumi}.}
\setcounter{enumi}{16}
\tightlist
\item
  The proton-motive force (pmf) has two components: electrical (\texorpdfstring{\ensuremath{\Delta}}{Delta}\texorpdfstring{\ensuremath{\Psi}}{Psi} \texorpdfstring{\ensuremath{\approx}}{~}
  \texorpdfstring{\ensuremath{-}}{-}140 mV) and chemical (\texorpdfstring{\ensuremath{\Delta}}{Delta}pH \texorpdfstring{\ensuremath{\approx}}{~} 1 unit = \texorpdfstring{\ensuremath{-}}{-}59 mV at 37°C; note sign
  convention for force driving H\texorpdfstring{\textsuperscript{+}}{+} into the matrix). Calculate the total
  pmf in mV and the Gibbs energy per mole of protons entering the matrix
  (\texorpdfstring{\ensuremath{\Delta}}{Delta}G = F\texorpdfstring{\ensuremath{\Delta}}{Delta}\texorpdfstring{\ensuremath{\Psi}}{Psi} + 2.303RT\texorpdfstring{\ensuremath{\cdot}}{.}\texorpdfstring{\ensuremath{\Delta}}{Delta}pH).
\end{enumerate}

\begin{enumerate}
\def\labelenumi{\arabic{enumi}.}
\setcounter{enumi}{17}
\tightlist
\item
  The malate-aspartate shuttle transfers NADH electrons from the
  cytoplasm into the mitochondrial matrix. Why is this shuttle needed,
  and why does a cell using the glycerol-3-phosphate shuttle instead
  produce slightly less ATP? (Hint: which NADH-equivalent is delivered
  to which complex?)
\end{enumerate}

\begin{enumerate}
\def\labelenumi{\arabic{enumi}.}
\setcounter{enumi}{18}
\tightlist
\item
  Cancer cells preferentially use aerobic glycolysis (Warburg effect)
  even when oxygen is available. Propose three mechanistic explanations
  for why rapidly proliferating cells would benefit from this apparently
  inefficient strategy.
\end{enumerate}

\begin{enumerate}
\def\labelenumi{\arabic{enumi}.}
\setcounter{enumi}{19}
\tightlist
\item
  Chloramphenicol inhibits bacterial 50S ribosomes and mitochondrial
  ribosomes. Explain why chloramphenicol can cause myelosuppression
  (bone marrow failure) after prolonged use, and why it is still used
  clinically when other antibiotics fail.
\end{enumerate}

\subsection*{Questions 21--30: Synthesis and
Evaluation}\label{questions-2130-synthesis-and-evaluation-8}
\addcontentsline{toc}{subsection}{Questions 21--30: Synthesis and
Evaluation}

\begin{enumerate}
\def\labelenumi{\arabic{enumi}.}
\setcounter{enumi}{20}
\tightlist
\item
  Evaluate the ATP synthase as a molecular motor. Describe the rotary
  mechanism (F\texorpdfstring{\textsubscript{0}}{0} rotation driven by H\texorpdfstring{\textsuperscript{+}}{+} flux \texorpdfstring{\ensuremath{\to}}{->} F\texorpdfstring{\textsubscript{1}}{1} conformational change \texorpdfstring{\ensuremath{\to}}{->}
  alternating binding-change mechanism). Calculate the torque generated
  if F\texorpdfstring{\textsubscript{0}}{0} rotates at 100 rpm with a step size of 120° per ATP.
\end{enumerate}

\begin{enumerate}
\def\labelenumi{\arabic{enumi}.}
\setcounter{enumi}{21}
\tightlist
\item
  Evaluate the efficiency of oxidative phosphorylation: how close is the
  actual ATP/O ratio to the theoretical maximum, and what factors cause
  inefficiency (proton leak, incomplete coupling)? What impact does
  membrane composition (cardiolipin content) have on coupling
  efficiency?
\end{enumerate}

\begin{enumerate}
\def\labelenumi{\arabic{enumi}.}
\setcounter{enumi}{22}
\tightlist
\item
  Design an experiment to test whether a newly discovered
  mitochondrially targeted compound acts as an uncoupler, a chain
  blocker, or an ATP synthase inhibitor. What measurements (O\texorpdfstring{\textsubscript{2}}{2}
  consumption rate, membrane potential, ATP levels, pH) would you take,
  and in what order?
\end{enumerate}

\begin{enumerate}
\def\labelenumi{\arabic{enumi}.}
\setcounter{enumi}{23}
\tightlist
\item
  Cells can switch between aerobic respiration, fermentation, and fatty
  acid oxidation. Design a metabolic flux experiment using \texorpdfstring{\textsuperscript{1}}{1}\texorpdfstring{\textsuperscript{3}}{3}C-labelled
  glucose and isotope mass spectrometry (metabolomics) to map which
  pathways are active under normoxia, hypoxia, and caloric restriction.
  What isotopomer patterns would distinguish glycolytic from TCA-derived
  carbon?
\end{enumerate}

\begin{enumerate}
\def\labelenumi{\arabic{enumi}.}
\setcounter{enumi}{24}
\tightlist
\item
  The pentose phosphate pathway (PPP) branches from glycolysis at G6P
  and produces NADPH and ribose-5-phosphate. Evaluate the relative
  importance of PPP in: (a) red blood cells (where NADPH protects
  against oxidative damage via glutathione); (b) rapidly proliferating
  cancer cells (where ribose is needed for nucleotide synthesis). Design
  a therapeutic strategy targeting the PPP specifically in cancer cells.
\end{enumerate}

\begin{enumerate}
\def\labelenumi{\arabic{enumi}.}
\setcounter{enumi}{25}
\tightlist
\item
  Mitochondria-derived reactive oxygen species (ROS) cause oxidative
  stress. Evaluate the paradox that low ROS function in redox signalling
  (activating NF-\texorpdfstring{\ensuremath{\kappa}}{kappa}B, Nrf2, and HIF-1\texorpdfstring{\ensuremath{\alpha}}{alpha}) while high ROS cause
  macromolecular damage. What antioxidant systems exist in mitochondria
  (SOD2, thioredoxin, peroxiredoxin), and how might targeted
  mitochondrial antioxidants (MitoQ, SS-31) prevent disease?
\end{enumerate}

\begin{enumerate}
\def\labelenumi{\arabic{enumi}.}
\setcounter{enumi}{26}
\tightlist
\item
  Metabolic reprogramming in hypoxia (HIF-1\texorpdfstring{\ensuremath{\alpha}}{alpha} pathway) switches cells
  from oxidative phosphorylation to anaerobic glycolysis. Evaluate the
  molecular logic: why would a cell transcriptionally upregulate
  glycolytic enzymes AND simultaneously downregulate Complex IV (via
  COX4-2 subunit switch) AND activate BNIP3-mediated mitophagy
  (selective mitochondrial degradation) in response to O\texorpdfstring{\textsubscript{2}}{2} lack?
\end{enumerate}

\begin{enumerate}
\def\labelenumi{\arabic{enumi}.}
\setcounter{enumi}{27}
\tightlist
\item
  The ketone bodies (\texorpdfstring{\ensuremath{\beta}}{beta}-hydroxybutyrate, acetoacetate) produced during
  fasting and ketogenic diets are increasingly recognised as signalling
  molecules (not just fuel). Evaluate the evidence that
  \texorpdfstring{\ensuremath{\beta}}{beta}-hydroxybutyrate inhibits histone deacetylases (HDACs), thereby
  acting as an epigenetic regulator. What genes are induced, and how
  might this explain the neuroprotective and anti-inflammatory
  properties attributed to ketogenic diets?
\end{enumerate}

\begin{enumerate}
\def\labelenumi{\arabic{enumi}.}
\setcounter{enumi}{28}
\tightlist
\item
  Thermodynamic coupling in oxidative phosphorylation can be described
  by non-equilibrium thermodynamics. Evaluate: when the ATP/ADP ratio is
  very high, does flow of H\texorpdfstring{\textsuperscript{+}}{+} through ATP synthase reverse (ATP \texorpdfstring{\ensuremath{\to}}{->} ADP +
  Pi, using pmf to rotate in reverse)? Under what physiological
  conditions does this reverse mode actually occur, and what is its role
  in maintaining mitochondrial membrane potential when Complex I/III
  activity is inhibited?
\end{enumerate}

\begin{enumerate}
\def\labelenumi{\arabic{enumi}.}
\setcounter{enumi}{29}
\tightlist
\item
  Critically evaluate the ``two-hit'' model of mitochondrial disease:
  patients with heteroplasmic mtDNA mutations are often asymptomatic
  until a critical threshold (often about 80\% mutated mtDNA)
  is exceeded. What cellular mechanisms (quality control,
  complementation in fused networks) maintain function below the
  threshold? At what threshold does bioenergetic failure occur, and how
  does tissue-specific metabolism (brain vs muscle vs liver) determine
  which tissue is most sensitive?
\end{enumerate}

\newpage

\section*{Questions --- Ch 10:
Photosynthesis}\label{questions-ch-10-photosynthesis}
\addcontentsline{toc}{section}{Questions --- Ch 10: Photosynthesis}

\label{sec:q_unit_III_photosynthesis}

\subsection*{Instructor Use and Coverage
Notes}\label{instructor-use-and-coverage-notes-12}
\addcontentsline{toc}{subsection}{Instructor Use and Coverage Notes}

\begin{itemize}
\tightlist
\item
  \textbf{Coverage target:} Compute carbon-fixation requirements and
  explain a limiting factor.
\item
  \textbf{Model/data emphasis:} Photon, ATP/NADPH, and Calvin-cycle
  stoichiometry.
\item
  \textbf{Assessment alignment:} Representing and Describing Data,
  Statistical Tests and Data Analysis.
\item
  \textbf{Misconception probe:} Plants do not eat sunlight; they use
  light energy to reduce carbon using electrons and enzymes.
\item
  \textbf{Transfer product:} Transfer photosynthetic constraints to
  crops, algae, climate, or ecosystem productivity.
\item
  \textbf{Grading focus:} award full credit for mechanism, evidence,
  boundary conditions, and units when a calculation is required.
\item
  \textbf{Suggested use:} draw one recall item, one application item,
  and one synthesis item when building a short quiz from this bank.
\end{itemize}

\subsection*{Questions 1--10: Recall and
Comprehension}\label{questions-110-recall-and-comprehension-9}
\addcontentsline{toc}{subsection}{Questions 1--10: Recall and
Comprehension}

\emph{This activity accompanies \cref{sec:unit_III_photosynthesis} of
the textbook --- review that chapter before attempting the exercises
below.}

\begin{enumerate}
\def\labelenumi{\arabic{enumi}.}
\tightlist
\item
  Write the overall equation for photosynthesis. Which reactant provides
  O\texorpdfstring{\textsubscript{2}}{2}?
\end{enumerate}

\begin{enumerate}
\def\labelenumi{\arabic{enumi}.}
\setcounter{enumi}{1}
\tightlist
\item
  Locate photosynthesis: in which organelle and in which specific
  membrane compartment do the light reactions occur?
\end{enumerate}

\begin{enumerate}
\def\labelenumi{\arabic{enumi}.}
\setcounter{enumi}{2}
\tightlist
\item
  What are photosystems I and II? What is the primary electron donor in
  each?
\end{enumerate}

\begin{enumerate}
\def\labelenumi{\arabic{enumi}.}
\setcounter{enumi}{3}
\tightlist
\item
  Define the Z-scheme. Trace the path of electrons from water to NADP\texorpdfstring{\textsuperscript{+}}{+}.
\end{enumerate}

\begin{enumerate}
\def\labelenumi{\arabic{enumi}.}
\setcounter{enumi}{4}
\tightlist
\item
  What are the products of the light reactions passed to the Calvin
  cycle?
\end{enumerate}

\begin{enumerate}
\def\labelenumi{\arabic{enumi}.}
\setcounter{enumi}{5}
\tightlist
\item
  Name the three stages of the Calvin cycle (also called the C3 cycle or
  light-independent reactions).
\end{enumerate}

\begin{enumerate}
\def\labelenumi{\arabic{enumi}.}
\setcounter{enumi}{6}
\tightlist
\item
  What is RuBisCO? What reaction does it catalyse and what is its
  substrate?
\end{enumerate}

\begin{enumerate}
\def\labelenumi{\arabic{enumi}.}
\setcounter{enumi}{7}
\tightlist
\item
  Distinguish C\texorpdfstring{\textsubscript{3}}{3}, C\texorpdfstring{\textsubscript{4}}{4}, and CAM plants. Give one example species of each.
\end{enumerate}

\begin{enumerate}
\def\labelenumi{\arabic{enumi}.}
\setcounter{enumi}{8}
\tightlist
\item
  What is photorespiration? Under what conditions is it maximised?
\end{enumerate}

\begin{enumerate}
\def\labelenumi{\arabic{enumi}.}
\setcounter{enumi}{9}
\tightlist
\item
  What is the wavelength absorbed by chlorophyll a and chlorophyll b?
\end{enumerate}

\subsection*{Questions 11--20: Application and
Analysis}\label{questions-1120-application-and-analysis-9}
\addcontentsline{toc}{subsection}{Questions 11--20: Application and
Analysis}

\begin{enumerate}
\def\labelenumi{\arabic{enumi}.}
\setcounter{enumi}{10}
\tightlist
\item
  DCMU (herbicide, kills plants) blocks the electron acceptor site of
  Photosystem II. Predict its effects on: (a) O\texorpdfstring{\textsubscript{2}}{2} evolution; (b) NADPH
  production; (c) ATP (photophosphorylation); (d) the Calvin cycle.
\end{enumerate}

\begin{enumerate}
\def\labelenumi{\arabic{enumi}.}
\setcounter{enumi}{11}
\tightlist
\item
  Non-cyclic photophosphorylation produces both ATP and NADPH. Cyclic
  photophosphorylation produces primarily ATP. Explain the mechanism of
  cyclic flow (electron recirculates through Fd \texorpdfstring{\ensuremath{\to}}{->} PQ \texorpdfstring{\ensuremath{\to}}{->} Cyt b\texorpdfstring{\textsubscript{6}}{6}f \texorpdfstring{\ensuremath{\to}}{->} PSI)
  and when a cell uses cyclic vs non-cyclic flow.
\end{enumerate}

\begin{enumerate}
\def\labelenumi{\arabic{enumi}.}
\setcounter{enumi}{12}
\tightlist
\item
  A student growing \emph{Arabidopsis thaliana} increases light
  intensity 10-fold. Initially photosynthetic rate rises, then plateaus.
  Explain the two-phase response: which reaction limits at low light
  (light reactions) and which limits at high light (Calvin cycle enzyme
  capacity)?
\end{enumerate}

\begin{enumerate}
\def\labelenumi{\arabic{enumi}.}
\setcounter{enumi}{13}
\tightlist
\item
  C\texorpdfstring{\textsubscript{4}}{4} plants (e.g., maize) separate initial CO\texorpdfstring{\textsubscript{2}}{2} fixation (HCO\texorpdfstring{\textsubscript{3}}{3}\texorpdfstring{\textsuperscript{-}}{-} \texorpdfstring{\ensuremath{\to}}{->}
  oxaloacetate, in mesophyll) from RuBisCO (in bundle sheath). Explain:
  (a) how this concentrates CO\texorpdfstring{\textsubscript{2}}{2} around RuBisCO; (b) why it reduces
  photorespiration; (c) why C\texorpdfstring{\textsubscript{4}}{4} plants outcompete C\texorpdfstring{\textsubscript{3}}{3} plants in hot, sunny
  conditions.
\end{enumerate}

\begin{enumerate}
\def\labelenumi{\arabic{enumi}.}
\setcounter{enumi}{14}
\tightlist
\item
  The quantum yield of O\texorpdfstring{\textsubscript{2}}{2} production is about 1 O\texorpdfstring{\textsubscript{2}}{2} per 8
  photons absorbed (theoretical: 1 O\texorpdfstring{\textsubscript{2}}{2} per 8 photons absorbed). Why is
  this the minimum, and how does the Z-scheme energetically justify why
  4 photons per PSII and 4 per PSI are needed to produce 1 O\texorpdfstring{\textsubscript{2}}{2} and 2
  NADPH?
\end{enumerate}

\begin{enumerate}
\def\labelenumi{\arabic{enumi}.}
\setcounter{enumi}{15}
\tightlist
\item
  Photorespiration occurs when RuBisCO oxygenates RuBP (instead of
  carboxylating it), producing phosphoglycolate, which must be recycled
  through the glycolate pathway. Calculate: if 25\% of RuBisCO reactions
  are oxygenation rather than carboxylation, what fraction of carbon
  fixed is immediately lost? How does this affect net carbon gain?
\end{enumerate}

\begin{enumerate}
\def\labelenumi{\arabic{enumi}.}
\setcounter{enumi}{16}
\tightlist
\item
  Antenna complexes (LHCI, LHCII) harvest photons and funnel excitation
  energy to reaction centres by resonance energy transfer (Förster
  transfer). Explain: (a) why antenna complexes increase efficiency; (b)
  why there is a risk of photodamage (singlet O\texorpdfstring{\textsubscript{2}}{2} production) at high
  light; (c) what non-photochemical quenching (NPQ) does to protect the
  reaction centres.
\end{enumerate}

\begin{enumerate}
\def\labelenumi{\arabic{enumi}.}
\setcounter{enumi}{17}
\tightlist
\item
  Rubisco activase unfolds the tight inhibitory conformation of RuBisCO
  when CO\texorpdfstring{\textsubscript{2}}{2} is absent. Explain: (a) why RuBisCO requires activation
  before catalysis; (b) how activase uses ATP hydrolysis to open the
  active site; (c) why rubisco activase is itself heat-sensitive and
  limits photosynthesis at high temperatures.
\end{enumerate}

\begin{enumerate}
\def\labelenumi{\arabic{enumi}.}
\setcounter{enumi}{18}
\tightlist
\item
  CAM (Crassulacean acid metabolism) plants open stomata primarily at
  night. Trace: (a) nocturnal CO\texorpdfstring{\textsubscript{2}}{2} fixation into malate (by PEPC); (b)
  daytime malate decarboxylation to CO\texorpdfstring{\textsubscript{2}}{2}; (c) CO\texorpdfstring{\textsubscript{2}}{2} fixation by RuBisCO;
  (d) the advantage in arid environments. Compare water-use efficiency
  (WUE) of CAM vs C\texorpdfstring{\textsubscript{3}}{3} plants.
\end{enumerate}

\begin{enumerate}
\def\labelenumi{\arabic{enumi}.}
\setcounter{enumi}{19}
\tightlist
\item
  The chloroplast ATP synthase (CF\texorpdfstring{\textsubscript{1}}{1}F\texorpdfstring{\textsubscript{0}}{0}) is structurally similar to the
  mitochondrial ATP synthase. Compare: the direction of proton gradient
  (lumen to stroma vs IMS to matrix); whether CF\texorpdfstring{\textsubscript{1}}{1}F\texorpdfstring{\textsubscript{0}}{0} can run in reverse
  (hydrolyse ATP to pump H\texorpdfstring{\textsuperscript{+}}{+}); and the regulatory role of the
  thiol-regulated \texorpdfstring{\ensuremath{\gamma}}{gamma}-subunit in the dark vs light.
\end{enumerate}

\subsection*{Questions 21--30: Synthesis and
Evaluation}\label{questions-2130-synthesis-and-evaluation-9}
\addcontentsline{toc}{subsection}{Questions 21--30: Synthesis and
Evaluation}

\begin{enumerate}
\def\labelenumi{\arabic{enumi}.}
\setcounter{enumi}{20}
\tightlist
\item
  Evaluate the artificial leaf as a technology for solar fuel
  production: splitting water to produce H\texorpdfstring{\textsubscript{2}}{2} photochemically. What
  catalysts mimic the Mn cluster of the oxygen-evolving complex (OEC)?
  What are the current efficiency limits, and how do they compare to
  photosynthetic efficiency (about 3\% vs photovoltaic
  about 25\%)?
\end{enumerate}

\begin{enumerate}
\def\labelenumi{\arabic{enumi}.}
\setcounter{enumi}{21}
\tightlist
\item
  Rising CO\texorpdfstring{\textsubscript{2}}{2} atmospheric concentrations (currently about 420
  ppm) should theoretically increase C\texorpdfstring{\textsubscript{3}}{3} photosynthesis (CO\texorpdfstring{\textsubscript{2}}{2}
  fertilisation effect) but FACE (Free Air CO\texorpdfstring{\textsubscript{2}}{2} Enrichment) experiments
  show attenuated responses. Evaluate three mechanisms of
  ``photosynthetic downregulation'' under elevated CO\texorpdfstring{\textsubscript{2}}{2} and assess the
  net impact on crop yields.
\end{enumerate}

\begin{enumerate}
\def\labelenumi{\arabic{enumi}.}
\setcounter{enumi}{22}
\tightlist
\item
  RuBisCO is notoriously inefficient (kcat about 3 s\texorpdfstring{\textsuperscript{-}}{-}\texorpdfstring{\textsuperscript{1}}{1},
  compared to \textgreater1,000 s\texorpdfstring{\textsuperscript{-}}{-}\texorpdfstring{\textsuperscript{1}}{1} for most enzymes). Evaluate why
  evolution has not improved RuBisCO, and assess recent synthetic
  biology approaches to engineer an improved RuBisCO (e.g., Form II
  RuBisCO from bacteria, directed evolution) for crop improvement.
\end{enumerate}

\begin{enumerate}
\def\labelenumi{\arabic{enumi}.}
\setcounter{enumi}{23}
\tightlist
\item
  Shade-adapted plants have higher LHCII content and lower PSII/PSI
  ratios than sun plants. Evaluate the molecular mechanisms of this
  acclimation: how does excess excitation pressure in low light signal
  LHCII assembly vs LHCII disassembly in high light? What role does STN7
  kinase (LHCII phosphorylation) play in state transitions?
\end{enumerate}

\begin{enumerate}
\def\labelenumi{\arabic{enumi}.}
\setcounter{enumi}{24}
\tightlist
\item
  Oxygenic photosynthesis evolved about 2.7 billion years ago
  and fundamentally transformed Earth's atmosphere. Evaluate the
  geochemical evidence (banded iron formations, the Great Oxidation
  Event) for the timing of oxygenic photosynthesis. What was the
  ecological impact on anaerobic life, and how did some lineages evolve
  oxygen tolerance?
\end{enumerate}

\begin{enumerate}
\def\labelenumi{\arabic{enumi}.}
\setcounter{enumi}{25}
\tightlist
\item
  Synthetic biology has produced a ``photosynthetic organism'' where an
  entire C\texorpdfstring{\textsubscript{4}}{4} photosynthetic apparatus is introduced into a C\texorpdfstring{\textsubscript{3}}{3} plant (C\texorpdfstring{\textsubscript{4}}{4}
  rice project). Evaluate what genetic modifications are required, which
  modifications have been achieved, and what the bottlenecks are for
  realising this in a crop plant.
\end{enumerate}

\begin{enumerate}
\def\labelenumi{\arabic{enumi}.}
\setcounter{enumi}{26}
\tightlist
\item
  Photoinhibition occurs when PSII is irreversibly damaged by excess
  light. The D1 protein of PSII is a site of photodamage and must be
  replaced every about 30 minutes under high light. Evaluate
  the molecular mechanism of D1 turnover (degradation by FtsH protease,
  ribosomal de novo synthesis, PSII reassembly in thylakoid membrane)
  and explain why this turnover rate limits maximal photosynthetic
  productivity.
\end{enumerate}

\begin{enumerate}
\def\labelenumi{\arabic{enumi}.}
\setcounter{enumi}{27}
\tightlist
\item
  Horizontal gene transfer of photosynthesis genes between cyanobacteria
  and early eukaryotes (primary endosymbiosis) versus between
  non-photosynthetic organisms and photosynthetic ones (secondary
  endosymbiosis) produced the diverse eukaryotic algal groups. Evaluate
  the complexity of gene transfer that was required to set up a
  functional chloroplast, focusing on the challenge of targeting
  nuclear-encoded proteins to the chloroplast (TIC/TOC translocon
  import).
\end{enumerate}

\begin{enumerate}
\def\labelenumi{\arabic{enumi}.}
\setcounter{enumi}{28}
\tightlist
\item
  Phytoplankton (mainly diatoms, dinoflagellates, cyanobacteria) account
  for about 50\% of global CO\texorpdfstring{\textsubscript{2}}{2} fixation despite having about
  1\% of terrestrial plant biomass. Evaluate what properties of aquatic
  photosynthesis explain this high per-mass productivity, and assess the
  feasibility of enhancing ocean photosynthesis (e.g., iron
  fertilisation) as a climate change mitigation strategy.
\end{enumerate}

\begin{enumerate}
\def\labelenumi{\arabic{enumi}.}
\setcounter{enumi}{29}
\tightlist
\item
  Critically assess whether an organism could be engineered to
  simultaneously perform photosynthesis AND oxidative phosphorylation in
  the same compartment. What major thermodynamic and biochemical
  obstacles would need to be overcome? (Consider: O\texorpdfstring{\textsubscript{2}}{2} produced by
  photosynthesis would support oxidative phosphorylation, but the
  competing demands for NAD(P)H, the different evolutionary origins of
  the two systems, and the compartmentation challenges in a single
  organelle context.)
\end{enumerate}

\newpage

\section*{Questions --- Ch 11: Metabolic Integration and
Regulation}\label{questions-ch-11-metabolic-integration-and-regulation}
\addcontentsline{toc}{section}{Questions --- Ch 11: Metabolic
Integration and Regulation}

\label{sec:q_unit_III_metabolic_integration}

\subsection*{Instructor Use and Coverage
Notes}\label{instructor-use-and-coverage-notes-13}
\addcontentsline{toc}{subsection}{Instructor Use and Coverage Notes}

\begin{itemize}
\tightlist
\item
  \textbf{Coverage target:} Predict how a change in substrate, hormone,
  or enzyme activity redirects flux.
\item
  \textbf{Model/data emphasis:} Energy charge, control points, and
  pathway-flux comparisons.
\item
  \textbf{Assessment alignment:} Representing and Describing Data,
  Statistical Tests and Data Analysis.
\item
  \textbf{Misconception probe:} A pathway is not a one-way assembly
  line; reversibility and regulation define the real route.
\item
  \textbf{Transfer product:} Apply metabolic network reasoning to
  diabetes, fasting, exercise, or cancer metabolism.
\item
  \textbf{Grading focus:} award full credit for mechanism, evidence,
  boundary conditions, and units when a calculation is required.
\item
  \textbf{Suggested use:} draw one recall item, one application item,
  and one synthesis item when building a short quiz from this bank.
\end{itemize}

\subsection*{Questions 1--10: Recall and
Comprehension}\label{questions-110-recall-and-comprehension-10}
\addcontentsline{toc}{subsection}{Questions 1--10: Recall and
Comprehension}

\emph{This activity accompanies
\cref{sec:unit_III_metabolic_integration} of the textbook --- review
that chapter before attempting the exercises below.}

\begin{enumerate}
\def\labelenumi{\arabic{enumi}.}
\tightlist
\item
  What does ``metabolic integration'' mean? Give one example of how
  carbohydrate and lipid metabolism are connected.
\end{enumerate}

\begin{enumerate}
\def\labelenumi{\arabic{enumi}.}
\setcounter{enumi}{1}
\tightlist
\item
  Define anabolism and catabolism. Give one example of each.
\end{enumerate}

\begin{enumerate}
\def\labelenumi{\arabic{enumi}.}
\setcounter{enumi}{2}
\tightlist
\item
  What is the role of AMP:ATP ratio as a metabolic signal? Which enzyme
  is the key AMP sensor?
\end{enumerate}

\begin{enumerate}
\def\labelenumi{\arabic{enumi}.}
\setcounter{enumi}{3}
\tightlist
\item
  What is AMPK? Describe two metabolic processes it activates and two it
  inhibits.
\end{enumerate}

\begin{enumerate}
\def\labelenumi{\arabic{enumi}.}
\setcounter{enumi}{4}
\tightlist
\item
  What is mTOR? What signals activate it, and what anabolic processes
  does it promote?
\end{enumerate}

\begin{enumerate}
\def\labelenumi{\arabic{enumi}.}
\setcounter{enumi}{5}
\tightlist
\item
  Name the organs that play a primary role in glucose homeostasis after
  a meal.
\end{enumerate}

\begin{enumerate}
\def\labelenumi{\arabic{enumi}.}
\setcounter{enumi}{6}
\tightlist
\item
  What is the metabolic role of the liver in the fed vs fasted state?
\end{enumerate}

\begin{enumerate}
\def\labelenumi{\arabic{enumi}.}
\setcounter{enumi}{7}
\tightlist
\item
  Define the Cori cycle. Which two organs are involved, and which
  molecule cycles between them?
\end{enumerate}

\begin{enumerate}
\def\labelenumi{\arabic{enumi}.}
\setcounter{enumi}{8}
\tightlist
\item
  What is the glucose-fatty acid cycle (Randle cycle)? What is the
  primary substrate competition it describes?
\end{enumerate}

\begin{enumerate}
\def\labelenumi{\arabic{enumi}.}
\setcounter{enumi}{9}
\tightlist
\item
  Define metabolic syndrome. List four diagnostic criteria.
\end{enumerate}

\subsection*{Questions 11--20: Application and
Analysis}\label{questions-1120-application-and-analysis-10}
\addcontentsline{toc}{subsection}{Questions 11--20: Application and
Analysis}

\begin{enumerate}
\def\labelenumi{\arabic{enumi}.}
\setcounter{enumi}{10}
\tightlist
\item
  After a high-carbohydrate meal, blood glucose rises. Trace the
  hormonal and metabolic responses in: (a) the pancreas; (b) the liver;
  (c) skeletal muscle; (d) adipose tissue --- over the 2 hours
  post-meal.
\end{enumerate}

\begin{enumerate}
\def\labelenumi{\arabic{enumi}.}
\setcounter{enumi}{11}
\tightlist
\item
  During prolonged fasting (48 hours), the body must maintain blood
  glucose for the brain. List in order the fuels used: (a) which depot
  is depleted first; (b) when and where does gluconeogenesis begin; (c)
  when do ketone bodies become the primary brain fuel?
\end{enumerate}

\begin{enumerate}
\def\labelenumi{\arabic{enumi}.}
\setcounter{enumi}{12}
\tightlist
\item
  Insulin stimulates GLUT4 translocation to the plasma membrane of
  muscle and adipose cells via PI3K-Akt-AS160 pathway. Explain the
  molecular cascade from insulin receptor activation to GLUT4 vesicle
  exocytosis. What happens to GLUT4 when insulin is removed?
\end{enumerate}

\begin{enumerate}
\def\labelenumi{\arabic{enumi}.}
\setcounter{enumi}{13}
\tightlist
\item
  PFK-1 (phosphofructokinase-1) is a key regulatory enzyme of
  glycolysis. It is activated by AMP, ADP, and fructose-2,6-bisphosphate
  (F-2,6-BP) and inhibited by ATP, citrate, and H\texorpdfstring{\textsuperscript{+}}{+}. Explain how each
  regulator makes physiological sense in terms of energy status
  signalling.
\end{enumerate}

\begin{enumerate}
\def\labelenumi{\arabic{enumi}.}
\setcounter{enumi}{14}
\tightlist
\item
  A patient with Type 2 Diabetes has high fasting blood glucose and high
  insulin (insulin resistance). Explain: (a) why cells are resistant to
  insulin; (b) why the liver continues gluconeogenesis despite high
  insulin; (c) what happens to HbA1c in chronic hyperglycaemia; (d) why
  metformin (activates AMPK) is effective.
\end{enumerate}

\begin{enumerate}
\def\labelenumi{\arabic{enumi}.}
\setcounter{enumi}{15}
\tightlist
\item
  The malate-aspartate shuttle moves NADH equivalents from the cytoplasm
  to the mitochondrial matrix. Explain: (a) why direct NADH cannot cross
  the inner mitochondrial membrane; (b) the carrier proteins involved;
  (c) why cardiac muscle prefers this shuttle over the
  glycerol-3-phosphate shuttle.
\end{enumerate}

\begin{enumerate}
\def\labelenumi{\arabic{enumi}.}
\setcounter{enumi}{16}
\tightlist
\item
  Brown adipose tissue (BAT) thermogenesis is mediated by UCP1. Explain
  the adrenergic activation cascade: cold \texorpdfstring{\ensuremath{\to}}{->} sympathetic nervous system \texorpdfstring{\ensuremath{\to}}{->}
  \texorpdfstring{\ensuremath{\beta}}{beta}\texorpdfstring{\textsubscript{3}}{3}-adrenergic receptor \texorpdfstring{\ensuremath{\to}}{->} cAMP \texorpdfstring{\ensuremath{\to}}{->} PKA \texorpdfstring{\ensuremath{\to}}{->} hormone-sensitive lipase \texorpdfstring{\ensuremath{\to}}{->} free
  fatty acid (FFA) release \texorpdfstring{\ensuremath{\to}}{->} FFA activate/open UCP1 \texorpdfstring{\ensuremath{\to}}{->} H\texorpdfstring{\textsuperscript{+}}{+} leak \texorpdfstring{\ensuremath{\to}}{->} heat. At
  which step is UCP1 expression also increased transcriptionally?
\end{enumerate}

\begin{enumerate}
\def\labelenumi{\arabic{enumi}.}
\setcounter{enumi}{17}
\tightlist
\item
  Glucagon activates gluconeogenesis in liver cells. Trace from glucagon
  binding \texorpdfstring{\ensuremath{\to}}{->} Gs \texorpdfstring{\ensuremath{\to}}{->} cAMP \texorpdfstring{\ensuremath{\to}}{->} PKA \texorpdfstring{\ensuremath{\to}}{->} phosphorylation of CREB \texorpdfstring{\ensuremath{\to}}{->} activation of
  PEPCK and G6Pase gene transcription. How does insulin oppose this
  pathway via Akt (inactivating FOXO1 transcription factor)?
\end{enumerate}

\begin{enumerate}
\def\labelenumi{\arabic{enumi}.}
\setcounter{enumi}{18}
\tightlist
\item
  The hypothalamus integrates hunger and satiety signals from peripheral
  hormones (leptin, ghrelin, PYY, insulin). Explain: (a) which neurons
  (NPY/AgRP vs POMC/CART) are activated by fasting vs feeding; (b) how
  melanocortin-4 receptor (MC4R) controls food intake; (c) why leptin
  resistance in obesity prevents normalisation of energy balance.
\end{enumerate}

\begin{enumerate}
\def\labelenumi{\arabic{enumi}.}
\setcounter{enumi}{19}
\tightlist
\item
  Fatty acid synthesis and \texorpdfstring{\ensuremath{\beta}}{beta}-oxidation cannot occur simultaneously ---
  they would create a futile cycle. Explain how the cell prevents this:
  (a) what molecule inhibits carnitine palmitoyltransferase I (CPTI)
  during fatty acid synthesis; (b) in which compartment is malonyl-CoA
  produced; (c) why this system specifically blocks mitochondrial FA
  import while allowing FA synthesis in the cytoplasm.
\end{enumerate}

\subsection*{Questions 21--30: Synthesis and
Evaluation}\label{questions-2130-synthesis-and-evaluation-10}
\addcontentsline{toc}{subsection}{Questions 21--30: Synthesis and
Evaluation}

\begin{enumerate}
\def\labelenumi{\arabic{enumi}.}
\setcounter{enumi}{20}
\tightlist
\item
  Evaluate the ``fuel partitioning'' hypothesis in obesity: does obesity
  arise primarily from excess carbohydrate or excess fat intake, and
  does macronutrient composition matter for energy balance beyond
  caloric excess? Cite the results of isocaloric controlled diet trials.
\end{enumerate}

\begin{enumerate}
\def\labelenumi{\arabic{enumi}.}
\setcounter{enumi}{21}
\tightlist
\item
  Evaluate the ketogenic diet as a treatment for drug-resistant
  epilepsy. Propose at least three molecular mechanisms by which
  elevated ketone bodies (\texorpdfstring{\ensuremath{\beta}}{beta}-hydroxybutyrate primarily as an HDAC
  inhibitor, adenosine release, or KATP channel opening) could reduce
  neuronal excitability.
\end{enumerate}

\begin{enumerate}
\def\labelenumi{\arabic{enumi}.}
\setcounter{enumi}{22}
\tightlist
\item
  Design an experiment to determine whether AMPK activation in liver
  cells reduces hepatic glucose output primarily by: (a) decreasing
  gluconeogenesis gene expression; (b) acutely phosphorylating and
  inactivating TORC2 (CREB co-activator); (c) decreasing hepatic ATP to
  slow enzymatic gluconeogenesis. Distinguish these mechanisms using
  selective inhibitors and reporter assays.
\end{enumerate}

\begin{enumerate}
\def\labelenumi{\arabic{enumi}.}
\setcounter{enumi}{23}
\tightlist
\item
  Evaluate the evidence for ``metabolic flexibility'' --- the ability of
  healthy cells to switch between glucose and fat oxidation based on
  substrate availability. How does this differ in Type 2 Diabetes,
  insulin resistance, and cancer cells? What biochemical assay (Seahorse
  XF analyser) measures metabolic flexibility?
\end{enumerate}

\begin{enumerate}
\def\labelenumi{\arabic{enumi}.}
\setcounter{enumi}{24}
\tightlist
\item
  The metabolic role of the gut microbiome is increasingly recognised:
  SCFAs (acetate, propionate, butyrate) produced by bacterial
  fermentation of dietary fibre signal to host cells. Evaluate how
  propionate signals to the liver (inhibiting gluconeogenesis) and
  butyrate signals to colonocytes (primary fuel) and immune cells (HDAC
  inhibition, regulatory T cell induction).
\end{enumerate}

\begin{enumerate}
\def\labelenumi{\arabic{enumi}.}
\setcounter{enumi}{25}
\tightlist
\item
  Evaluate the metabolic rationale for intermittent fasting (IF)
  regimens (16:8 time-restricted eating, 5:2 alternate-day fasting).
  Identify three metabolic benefits of IF beyond simple caloric
  restriction, and propose the molecular mechanisms driving each
  benefit.
\end{enumerate}

\begin{enumerate}
\def\labelenumi{\arabic{enumi}.}
\setcounter{enumi}{26}
\tightlist
\item
  Cancer cells depend on glutamine (not glucose) for anaplerosis ---
  replenishing TCA cycle intermediates for biosynthesis. Evaluate: (a)
  which TCA intermediate is most depleted in rapidly proliferating
  cells; (b) how glutamine enters the TCA cycle (via glutaminase \texorpdfstring{\ensuremath{\to}}{->}
  glutamate \texorpdfstring{\ensuremath{\to}}{->} \texorpdfstring{\ensuremath{\alpha}}{alpha}-ketoglutarate); (c) why targeting glutamine metabolism
  (CB-839, a glutaminase inhibitor) might be selectively toxic to cancer
  cells.
\end{enumerate}

\begin{enumerate}
\def\labelenumi{\arabic{enumi}.}
\setcounter{enumi}{27}
\tightlist
\item
  Evaluate the concept of ``metaflammation'' --- the chronic low-grade
  inflammation seen in obesity. Propose a mechanistic link from
  adipocyte lipid overload \texorpdfstring{\ensuremath{\to}}{->} ER stress \texorpdfstring{\ensuremath{\to}}{->} NF-\texorpdfstring{\ensuremath{\kappa}}{kappa}B activation \texorpdfstring{\ensuremath{\to}}{->}
  pro-inflammatory cytokine (TNF-\texorpdfstring{\ensuremath{\alpha}}{alpha}, IL-6) secretion \texorpdfstring{\ensuremath{\to}}{->} insulin resistance
  in muscle and liver. At which step would you intervene
  pharmacologically?
\end{enumerate}

\begin{enumerate}
\def\labelenumi{\arabic{enumi}.}
\setcounter{enumi}{28}
\tightlist
\item
  NAFLD (non-alcoholic fatty liver disease) progresses from steatosis \texorpdfstring{\ensuremath{\to}}{->}
  NASH (steatohepatitis) \texorpdfstring{\ensuremath{\to}}{->} cirrhosis. Evaluate the ``two-hit
  hypothesis'': what is the first hit (lipid accumulation), what is the
  second hit (oxidative stress, inflammatory cytokines from adipose
  tissue or gut microbiome), and why targeting a single hit is
  insufficient for treatment?
\end{enumerate}

\begin{enumerate}
\def\labelenumi{\arabic{enumi}.}
\setcounter{enumi}{29}
\tightlist
\item
  Critically evaluate the use of GLP-1 receptor agonists (semaglutide,
  tirzepatide) for obesity treatment. What are the distinct mechanisms
  by which GLP-1 agonists reduce appetite (hypothalamic signalling),
  slow gastric emptying, and promote pancreatic \texorpdfstring{\ensuremath{\beta}}{beta}-cell survival? What
  are the metabolic consequences of rapid weight loss induced by these
  drugs (muscle loss, bone density) and how should therapy be managed?
\end{enumerate}

\newpage

\section*{Questions --- Ch 12: DNA Replication and the Cell
Cycle}\label{questions-ch-12-dna-replication-and-the-cell-cycle}
\addcontentsline{toc}{section}{Questions --- Ch 12: DNA Replication and
the Cell Cycle}

\label{sec:q_unit_IV_dna_replication_and_cell_cycle}

\subsection*{Instructor Use and Coverage
Notes}\label{instructor-use-and-coverage-notes-14}
\addcontentsline{toc}{subsection}{Instructor Use and Coverage Notes}

\begin{itemize}
\tightlist
\item
  \textbf{Coverage target:} Calculate replication timing or error burden
  and infer checkpoint consequences.
\item
  \textbf{Model/data emphasis:} Replication-fork speed, error-rate, and
  cell-cycle timing calculations.
\item
  \textbf{Assessment alignment:} Concept Explanation, Questions and
  Methods, Argumentation.
\item
  \textbf{Misconception probe:} High fidelity is not automatic; it is
  built from multiple partially redundant safeguards.
\item
  \textbf{Transfer product:} Transfer replication logic to cancer,
  aging, viral replication, or antibiotic targets.
\item
  \textbf{Grading focus:} award full credit for mechanism, evidence,
  boundary conditions, and units when a calculation is required.
\item
  \textbf{Suggested use:} draw one recall item, one application item,
  and one synthesis item when building a short quiz from this bank.
\end{itemize}

\subsection*{Questions 1--10: Recall and
Comprehension}\label{questions-110-recall-and-comprehension-11}
\addcontentsline{toc}{subsection}{Questions 1--10: Recall and
Comprehension}

\emph{This activity accompanies
\cref{sec:unit_IV_dna_replication_and_cell_cycle} of the textbook ---
review that chapter before attempting the exercises below.}

\begin{enumerate}
\def\labelenumi{\arabic{enumi}.}
\tightlist
\item
  Describe the semiconservative model of DNA replication. What
  experiment confirmed it (Meselson-Stahl, 1958)?
\end{enumerate}

\begin{enumerate}
\def\labelenumi{\arabic{enumi}.}
\setcounter{enumi}{1}
\tightlist
\item
  List the enzymes required for DNA replication and state one function
  of each.
\end{enumerate}

\begin{enumerate}
\def\labelenumi{\arabic{enumi}.}
\setcounter{enumi}{2}
\tightlist
\item
  Why does DNA polymerase require a primer? Who synthesises the primer?
\end{enumerate}

\begin{enumerate}
\def\labelenumi{\arabic{enumi}.}
\setcounter{enumi}{3}
\tightlist
\item
  Distinguish the leading strand from the lagging strand during
  replication. Why does the lagging strand form Okazaki fragments?
\end{enumerate}

\begin{enumerate}
\def\labelenumi{\arabic{enumi}.}
\setcounter{enumi}{4}
\tightlist
\item
  What is the replication fork? Name the enzyme that separates the
  double helix at the fork.
\end{enumerate}

\begin{enumerate}
\def\labelenumi{\arabic{enumi}.}
\setcounter{enumi}{5}
\tightlist
\item
  Name the four phases of the eukaryotic cell cycle. What occurs in
  each?
\end{enumerate}

\begin{enumerate}
\def\labelenumi{\arabic{enumi}.}
\setcounter{enumi}{6}
\tightlist
\item
  What are cyclins and cyclin-dependent kinases (CDKs)? Give one example
  of a cyclin-CDK pair and its function.
\end{enumerate}

\begin{enumerate}
\def\labelenumi{\arabic{enumi}.}
\setcounter{enumi}{7}
\tightlist
\item
  What is a checkpoint? Name three checkpoints in the cell cycle and
  their functions.
\end{enumerate}

\begin{enumerate}
\def\labelenumi{\arabic{enumi}.}
\setcounter{enumi}{8}
\tightlist
\item
  What does the tumour suppressor protein p53 do when DNA damage is
  detected?
\end{enumerate}

\begin{enumerate}
\def\labelenumi{\arabic{enumi}.}
\setcounter{enumi}{9}
\tightlist
\item
  What is telomerase? Why do most somatic cells lack telomerase
  activity?
\end{enumerate}

\subsection*{Questions 11--20: Application and
Analysis}\label{questions-1120-application-and-analysis-11}
\addcontentsline{toc}{subsection}{Questions 11--20: Application and
Analysis}

\begin{enumerate}
\def\labelenumi{\arabic{enumi}.}
\setcounter{enumi}{10}
\tightlist
\item
  Using the Meselson-Stahl data: if \texorpdfstring{\textsuperscript{1}}{1}\texorpdfstring{\textsuperscript{5}}{5}N/\texorpdfstring{\textsuperscript{1}}{1}\texorpdfstring{\textsuperscript{5}}{5}N-labelled DNA is allowed to
  replicate for two rounds in \texorpdfstring{\textsuperscript{1}}{1}\texorpdfstring{\textsuperscript{4}}{4}N medium, what densities of DNA are
  expected and in what proportions? Show your reasoning with a diagram.
\end{enumerate}

\begin{enumerate}
\def\labelenumi{\arabic{enumi}.}
\setcounter{enumi}{11}
\tightlist
\item
  DNA polymerase III (prokaryotic) and DNA polymerase \texorpdfstring{\ensuremath{\delta}}{delta}/\texorpdfstring{\ensuremath{\varepsilon}}{epsilon} (eukaryotic)
  require a processivity clamp (\texorpdfstring{\ensuremath{\beta}}{beta}-clamp in prokaryotes; PCNA in
  eukaryotes). Explain how the sliding clamp increases processivity, and
  why this matters for replication speed of the 3 \texorpdfstring{\ensuremath{\times}}{x} 10\texorpdfstring{\textsuperscript{9}}{9} bp human genome.
\end{enumerate}

\begin{enumerate}
\def\labelenumi{\arabic{enumi}.}
\setcounter{enumi}{12}
\tightlist
\item
  Okazaki fragments in eukaryotes are about 200 nucleotides
  (roughly one nucleosome length). After primase removes the primer, DNA
  pol \texorpdfstring{\ensuremath{\delta}}{delta} fills the gap using the downstream Okazaki fragment as a flap.
  Explain the role of FEN1 (flap endonuclease) and DNA ligase in
  completing lagging strand synthesis.
\end{enumerate}

\begin{enumerate}
\def\labelenumi{\arabic{enumi}.}
\setcounter{enumi}{13}
\tightlist
\item
  The G\texorpdfstring{\textsubscript{1}}{1}/S checkpoint (restriction point) ensures cells don't enter S
  phase with insufficient resources or DNA damage. Explain the molecular
  mechanism: how does mitogen (growth factor) signalling lead to Rb
  phosphorylation \texorpdfstring{\ensuremath{\to}}{->} release of E2F transcription factor \texorpdfstring{\ensuremath{\to}}{->} S-phase gene
  expression? At what step does p16 (CDKN2A tumour suppressor) block
  this pathway?
\end{enumerate}

\begin{enumerate}
\def\labelenumi{\arabic{enumi}.}
\setcounter{enumi}{14}
\tightlist
\item
  Aphidicolin blocks DNA polymerase \texorpdfstring{\ensuremath{\alpha}}{alpha} and \texorpdfstring{\ensuremath{\delta}}{delta}, arresting cells at the G\texorpdfstring{\textsubscript{1}}{1}/S
  boundary. Hydroxyurea (HU) depletes dNTP pools. Compare the effects of
  each drug on: (a) replication fork speed; (b) fork stalling vs
  collapse; (c) activation of ATR checkpoint kinase; (d) cell cycle
  arrest phase.
\end{enumerate}

\begin{enumerate}
\def\labelenumi{\arabic{enumi}.}
\setcounter{enumi}{15}
\tightlist
\item
  Telomere shortening occurs at a rate of about 100 bp per cell
  division in most somatic cells. Calculate how many divisions a cell
  with initial telomere length of 15 kb could undergo before reaching
  the critical length of 6 kb. What happens to cells that continue
  dividing past this limit (crisis)?
\end{enumerate}

\begin{enumerate}
\def\labelenumi{\arabic{enumi}.}
\setcounter{enumi}{16}
\tightlist
\item
  The spindle assembly checkpoint (SAC) detects kinetochores not under
  tension (unattached to spindle microtubules). Explain the molecular
  mechanism: how does an unattached kinetochore produce the wait
  anaphase signal (MCC complex), how does this inhibit APC/C
  (anaphase-promoting complex), and what happens to separase and cohesin
  when SAC is satisfied?
\end{enumerate}

\begin{enumerate}
\def\labelenumi{\arabic{enumi}.}
\setcounter{enumi}{17}
\tightlist
\item
  CDK1-Cyclin B drives entry into mitosis. CDK1-Cyclin B phosphorylates
  nuclear lamins (dissolving the nuclear envelope), activates condensin
  (chromosome condensation), and activates APC/C (through intermediate
  steps). Explain why Cyclin B must be abruptly degraded at the
  metaphase-anaphase transition.
\end{enumerate}

\begin{enumerate}
\def\labelenumi{\arabic{enumi}.}
\setcounter{enumi}{18}
\tightlist
\item
  Oncogenes often bypass cell cycle checkpoints. Compare how Ras
  (constitutively active mutant) vs Myc (overexpressed) each drive cell
  cycle progression: through which cyclin-CDK pair(s) do they act, and
  why do cells with oncogenic mutations often also need to lose p53 or
  Rb to become fully cancerous?
\end{enumerate}

\begin{enumerate}
\def\labelenumi{\arabic{enumi}.}
\setcounter{enumi}{19}
\tightlist
\item
  Yeast two-hybrid screens identified that geminin (overexpressed in
  many cancers) inhibits Cdt1, preventing re-licensing of origins of
  replication. Explain: (a) what origin licensing means; (b) why
  re-replication would be catastrophic; (c) how overexpression of
  geminin PREVENTS re-replication by preventing MCM loading even after S
  phase begins.
\end{enumerate}

\subsection*{Questions 21--30: Synthesis and
Evaluation}\label{questions-2130-synthesis-and-evaluation-11}
\addcontentsline{toc}{subsection}{Questions 21--30: Synthesis and
Evaluation}

\begin{enumerate}
\def\labelenumi{\arabic{enumi}.}
\setcounter{enumi}{20}
\tightlist
\item
  Evaluate the precision of the S-phase checkpoint: ATR kinase is
  activated by RPA-coated single-stranded DNA (from stalled forks). How
  does this trigger CHK1 activation \texorpdfstring{\ensuremath{\to}}{->} Cdc25A phosphorylation \texorpdfstring{\ensuremath{\to}}{->}
  degradation \texorpdfstring{\ensuremath{\to}}{->} CDK2 inactivation \texorpdfstring{\ensuremath{\to}}{->} fork stabilisation? What happens if
  ATR is constitutively inactive (as in a mutation)?
\end{enumerate}

\begin{enumerate}
\def\labelenumi{\arabic{enumi}.}
\setcounter{enumi}{21}
\tightlist
\item
  Design an experiment to test whether telomerase reactivation in a
  telomere-shortened cancer cell line would allow unlimited
  proliferation or would require additional changes. Include: how you
  would deliver hTERT (the catalytic subunit), how you would measure
  telomere length over generations, and what growth advantage you would
  expect and under what assay conditions.
\end{enumerate}

\begin{enumerate}
\def\labelenumi{\arabic{enumi}.}
\setcounter{enumi}{22}
\tightlist
\item
  DNA replication is not 100\% accurate --- polymerase error rate is
  about 1 in 10\texorpdfstring{\textsuperscript{7}}{7}. After mismatch repair (MMR), the rate drops
  to about 1 in 10\texorpdfstring{\textsuperscript{9}}{9}. Evaluate: (a) the molecular mechanism of
  MMR (MutS, MutL, excision, resynthesis); (b) why MMR deficiency causes
  microsatellite instability (MSI); (c) why MSI cancers respond better
  to immune checkpoint inhibitors (PD-1 blockade).
\end{enumerate}

\begin{enumerate}
\def\labelenumi{\arabic{enumi}.}
\setcounter{enumi}{23}
\tightlist
\item
  Evaluate origin-of-replication usage in cancer: cancer cells often
  show replication stress (RFS --- replication fork stalling) due to
  premature S-phase entry driven by oncogenes. Propose how elevated RFS
  leads to: (a) DSB (double-strand break) formation at stalled forks;
  (b) inappropriate NHEJ repair causing chromosomal translocations; (c)
  selection pressure for p53 loss (cells that could normally
  checkpoint-arrest now instead die, so p53 loss allows survival).
\end{enumerate}

\begin{enumerate}
\def\labelenumi{\arabic{enumi}.}
\setcounter{enumi}{24}
\tightlist
\item
  The pre-replication complex (pre-RC) assembly (ORC \texorpdfstring{\ensuremath{\to}}{->} Cdc6 \texorpdfstring{\ensuremath{\to}}{->} Cdt1 \texorpdfstring{\ensuremath{\to}}{->}
  MCM2-7) licenses origins from late mitosis to G\texorpdfstring{\textsubscript{1}}{1}. Evaluate two
  licensing control systems: (a) Cdt1 degradation (Cullin4-Cdt2 E3
  ligase triggered by PCNA on DNA); (b) Geminin inhibition of Cdt1.
  Propose why two redundant mechanisms are needed, and what would happen
  if both were inactivated simultaneously.
\end{enumerate}

\begin{enumerate}
\def\labelenumi{\arabic{enumi}.}
\setcounter{enumi}{25}
\tightlist
\item
  Replication stress is thought to be the earliest event in
  pre-cancerous lesions (before DNA mutation). Evaluate the evidence
  from oncogene-induced senescence (OIS): how does RAS\^{}V12
  overexpression initially stall replication forks and activate the DNA
  damage response (p53-p21 axis, producing cellular senescence) before
  genomic instability occurs?
\end{enumerate}

\begin{enumerate}
\def\labelenumi{\arabic{enumi}.}
\setcounter{enumi}{26}
\tightlist
\item
  Centromere identity in eukaryotes is determined not by DNA sequence
  but by the presence of CENP-A (a histone H3 variant). Evaluate the
  epigenetic model of centromere inheritance: how is CENP-A loaded (by
  HJURP chaperone) at each S phase, and how can a ``centromere'' be
  moved to a new chromosomal location (neocentromere formation)?
\end{enumerate}

\begin{enumerate}
\def\labelenumi{\arabic{enumi}.}
\setcounter{enumi}{27}
\tightlist
\item
  Multi-fork replication in prokaryotes: \emph{E. coli} can initiate
  replication at the same origin before the previous round is complete
  (multi-fork replication), allowing doubling times faster than the
  S-phase duration. Evaluate whether eukaryotes could employ multi-fork
  replication and what structural constraints (nuclear envelope,
  nucleosomes, cohesin) would prevent this.
\end{enumerate}

\begin{enumerate}
\def\labelenumi{\arabic{enumi}.}
\setcounter{enumi}{28}
\tightlist
\item
  Evaluate the hypothesis that cancer is fundamentally a disease of the
  cell cycle: most hallmarks of cancer (unlimited proliferation, evasion
  of apoptosis, invasion) result from failure of cell cycle checkpoints
  rather than from any particular oncogenic mutation. Construct the best
  argument for and against this view using specific molecular examples.
\end{enumerate}

\begin{enumerate}
\def\labelenumi{\arabic{enumi}.}
\setcounter{enumi}{29}
\tightlist
\item
  Critically evaluate the kinetics of DNA repair: NHEJ (non-homologous
  end joining, fast, about 30 min) and HR (homologous
  recombination, slow, requires sister chromatid, restricted to S/G2).
  Predict the consequences for genomic integrity if NHEJ were the
  primary repair pathway available in a cell. Why is NHEJ
  ``error-prone'' at the molecular level, and what types of mutations
  does it introduce?
\end{enumerate}

\newpage

\section*{Questions --- Ch 13: Gene Expression --- Transcription and
Translation}\label{questions-ch-13-gene-expression-transcription-and-translation}
\addcontentsline{toc}{section}{Questions --- Ch 13: Gene Expression ---
Transcription and Translation}

\label{sec:q_unit_IV_gene_expression}

\subsection*{Instructor Use and Coverage
Notes}\label{instructor-use-and-coverage-notes-15}
\addcontentsline{toc}{subsection}{Instructor Use and Coverage Notes}

\begin{itemize}
\tightlist
\item
  \textbf{Coverage target:} Translate or annotate a sequence and explain
  how a regulatory change alters output.
\item
  \textbf{Model/data emphasis:} Reading-frame, codon, and
  expression-output calculations.
\item
  \textbf{Assessment alignment:} Concept Explanation, Questions and
  Methods, Argumentation.
\item
  \textbf{Misconception probe:} A gene is not simply a protein recipe;
  context controls when, where, and how much product appears.
\item
  \textbf{Transfer product:} Apply expression logic to mutations,
  biotechnology, development, and disease diagnostics.
\item
  \textbf{Grading focus:} award full credit for mechanism, evidence,
  boundary conditions, and units when a calculation is required.
\item
  \textbf{Suggested use:} draw one recall item, one application item,
  and one synthesis item when building a short quiz from this bank.
\end{itemize}

\subsection*{Questions 1--10: Recall and
Comprehension}\label{questions-110-recall-and-comprehension-12}
\addcontentsline{toc}{subsection}{Questions 1--10: Recall and
Comprehension}

\emph{This activity accompanies \cref{sec:unit_IV_gene_expression} of
the textbook --- review that chapter before attempting the exercises
below.}

\begin{enumerate}
\def\labelenumi{\arabic{enumi}.}
\tightlist
\item
  Define gene expression. What are the two major steps?
\end{enumerate}

\begin{enumerate}
\def\labelenumi{\arabic{enumi}.}
\setcounter{enumi}{1}
\tightlist
\item
  Describe the structure of a eukaryotic mRNA from 5' to 3', including
  most regulatory elements.
\end{enumerate}

\begin{enumerate}
\def\labelenumi{\arabic{enumi}.}
\setcounter{enumi}{2}
\tightlist
\item
  What is the role of RNA polymerase II in transcription? Where does it
  bind?
\end{enumerate}

\begin{enumerate}
\def\labelenumi{\arabic{enumi}.}
\setcounter{enumi}{3}
\tightlist
\item
  Define a transcription factor. Distinguish general vs specific
  transcription factors.
\end{enumerate}

\begin{enumerate}
\def\labelenumi{\arabic{enumi}.}
\setcounter{enumi}{4}
\tightlist
\item
  What is RNA splicing? What molecular machinery performs it?
\end{enumerate}

\begin{enumerate}
\def\labelenumi{\arabic{enumi}.}
\setcounter{enumi}{5}
\tightlist
\item
  Define the genetic code. What is a codon? How many codons encode amino
  acids?
\end{enumerate}

\begin{enumerate}
\def\labelenumi{\arabic{enumi}.}
\setcounter{enumi}{6}
\tightlist
\item
  What is the ribosome composed of? Name the three ribosomal sites (A,
  P, E sites).
\end{enumerate}

\begin{enumerate}
\def\labelenumi{\arabic{enumi}.}
\setcounter{enumi}{7}
\tightlist
\item
  Describe initiation, elongation, and termination of translation in
  eukaryotes.
\end{enumerate}

\begin{enumerate}
\def\labelenumi{\arabic{enumi}.}
\setcounter{enumi}{8}
\tightlist
\item
  What is tRNA? What is the anticodon, and what is aminoacyl-tRNA
  synthetase?
\end{enumerate}

\begin{enumerate}
\def\labelenumi{\arabic{enumi}.}
\setcounter{enumi}{9}
\tightlist
\item
  What is the wobble hypothesis? Which position of the codon allows
  degeneracy?
\end{enumerate}

\subsection*{Questions 11--20: Application and
Analysis}\label{questions-1120-application-and-analysis-12}
\addcontentsline{toc}{subsection}{Questions 11--20: Application and
Analysis}

\begin{enumerate}
\def\labelenumi{\arabic{enumi}.}
\setcounter{enumi}{10}
\tightlist
\item
  Pre-mRNA splicing must be precise to within one nucleotide. Explain
  the two-step transesterification mechanism of splicing and the role of
  branch point A, 5' splice site, and 3' splice site in spliceosome
  assembly.
\end{enumerate}

\begin{enumerate}
\def\labelenumi{\arabic{enumi}.}
\setcounter{enumi}{11}
\tightlist
\item
  Alternative splicing generates multiple protein isoforms from one
  gene. \emph{Drosophila Dscam} gene can produce \textgreater38,000
  isoforms through mutually exclusive exon splicing. Explain: (a) the
  general mechanism of alternative splicing; (b) why \emph{Dscam}
  diversity is important for axon guidance in neural wiring; (c) one
  human disease caused by aberrant alternative splicing.
\end{enumerate}

\begin{enumerate}
\def\labelenumi{\arabic{enumi}.}
\setcounter{enumi}{12}
\tightlist
\item
  The 5' cap (m\texorpdfstring{\textsuperscript{7}}{7}G) and poly-A tail of eukaryotic mRNA serve similar
  functions --- protecting mRNA from exonuclease attack and promoting
  translation initiation. Explain: (a) how the 5' cap is added; (b) how
  eIF4E binding to the cap recruits the 43S pre-initiation complex; (c)
  why some viruses (e.g., poliovirus) cleave eIF4G to direct ribosomes
  to IRES sequences instead.
\end{enumerate}

\begin{enumerate}
\def\labelenumi{\arabic{enumi}.}
\setcounter{enumi}{13}
\tightlist
\item
  The ribosome proof-reads aminoacyl-tRNA selection at the A site.
  Explain: (a) initial selection based on codon-anticodon base pairing
  (induced fit); (b) proofreading step (kinetokinetics --- GTPase
  activation primarily if correct aa-tRNA is bound); (c) why errors
  occur at about 1 in 10,000 vs DNA polymerase errors of 1 in
  10\texorpdfstring{\textsuperscript{9}}{9} (why is translation less accurate?).
\end{enumerate}

\begin{enumerate}
\def\labelenumi{\arabic{enumi}.}
\setcounter{enumi}{14}
\tightlist
\item
  Cycloheximide blocks the elongation step of eukaryotic translation
  (not prokaryotic), because it specifically inhibits the 60S subunit
  translocation step. Explain: (a) what translocation means
  mechanistically; (b) why prokaryotic ribosomes (70S) are not blocked;
  (c) how this selectivity is exploited in laboratory cell biology
  experiments.
\end{enumerate}

\begin{enumerate}
\def\labelenumi{\arabic{enumi}.}
\setcounter{enumi}{15}
\tightlist
\item
  Nonsense-mediated mRNA decay (NMD) degrades mRNAs with premature stop
  codons. Explain: (a) how the cell recognises a premature vs normal
  stop codon (role of exon-junction complexes); (b) the UPF1/SMG
  cascade; (c) why NMD is a quality control mechanism that also
  regulates normal gene expression for about 10\% of genes.
\end{enumerate}

\begin{enumerate}
\def\labelenumi{\arabic{enumi}.}
\setcounter{enumi}{16}
\tightlist
\item
  Polyribosomes (polysomes) consist of multiple ribosomes simultaneously
  translating the same mRNA. If a 1,500-nt coding sequence is occupied
  by ribosomes spaced on average every 80 nt, how many ribosomes are
  simultaneously translating this mRNA? Calculate the rate of protein
  synthesis per mRNA if each ribosome adds 2 amino acids/sec.
\end{enumerate}

\begin{enumerate}
\def\labelenumi{\arabic{enumi}.}
\setcounter{enumi}{17}
\tightlist
\item
  Internal ribosome entry sites (IRES) allow cap-independent translation
  initiation, used by viruses (Hepatitis C, poliovirus) and in cellular
  stress responses, when eIF4E activity is reduced. Compare
  IRES-mediated vs cap-dependent initiation: which cellular proteins are
  replaced by viral ITAF proteins, and under what stress conditions does
  cellular IRES-mediated translation occur?
\end{enumerate}

\begin{enumerate}
\def\labelenumi{\arabic{enumi}.}
\setcounter{enumi}{18}
\tightlist
\item
  The signal recognition particle (SRP) recognises hydrophobic signal
  peptides emerging from ribosomes translating secretory proteins.
  Explain: (a) how SRP stalls translation while targeting the ribosome
  to the rough ER; (b) how signal peptidase cleaves the signal peptide
  in the ER lumen; (c) what happens to cytosolic misfolded proteins
  lacking a signal peptide vs ER-targeted ones.
\end{enumerate}

\begin{enumerate}
\def\labelenumi{\arabic{enumi}.}
\setcounter{enumi}{19}
\tightlist
\item
  Transcription factors bind specific DNA sequences (motifs). A
  transcriptional activator contains a DNA-binding domain (DBD) and an
  activation domain (AD). Explain: (a) what structural motifs are common
  in DBDs (zinc finger, helix-turn-helix, leucine zipper); (b) how the
  activation domain recruits the Mediator complex; (c) how enhancers
  (located kb away from promoters) loop to contact the preinitiation
  complex through cohesin-mediated chromosome looping.
\end{enumerate}

\subsection*{Questions 21--30: Synthesis and
Evaluation}\label{questions-2130-synthesis-and-evaluation-12}
\addcontentsline{toc}{subsection}{Questions 21--30: Synthesis and
Evaluation}

\begin{enumerate}
\def\labelenumi{\arabic{enumi}.}
\setcounter{enumi}{20}
\tightlist
\item
  Evaluate the dual role of SR proteins (serine-arginine-rich splicing
  factors) in both constitutive and alternative splicing and mRNA
  export. What does the dual role of SR proteins reveal about the
  evolutionary coupling between nuclear RNA processing and cytoplasmic
  translation?
\end{enumerate}

\begin{enumerate}
\def\labelenumi{\arabic{enumi}.}
\setcounter{enumi}{21}
\tightlist
\item
  Design a comprehensive experiment to map most enhancers controlling a
  specific gene in a given cell type using: (a) H3K27ac ChIP-seq (active
  enhancer mark); (b) ATAC-seq (open chromatin); (c) HiC chromatin
  conformation capture; (d) CRISPR deletion of candidate enhancers and
  transcriptional output measurement. How would you prioritise which
  enhancers are most important?
\end{enumerate}

\begin{enumerate}
\def\labelenumi{\arabic{enumi}.}
\setcounter{enumi}{22}
\tightlist
\item
  The ``RNA world'' hypothesis proposes that early life was based on RNA
  molecules that could both carry genetic information and catalyse
  reactions. Evaluate: (a) the evidence supporting RNA world (ribozymes,
  ribosome as a ribozyme, universality of the genetic code); (b) the
  transition from RNA world to the DNA-protein world; (c) what modern
  ribozymes tell us about the chemistry of early catalysis.
\end{enumerate}

\begin{enumerate}
\def\labelenumi{\arabic{enumi}.}
\setcounter{enumi}{23}
\tightlist
\item
  Transcription and translation are physically coupled in prokaryotes
  (translation begins before transcription is complete). Evaluate the
  consequences of this coupling: how does it allow attenuation
  (regulation of transcription termination by ribosome speed), and why
  would coupling like this have evolved in prokaryotes but not
  eukaryotes (where transcription and translation are spatially
  separated)?
\end{enumerate}

\begin{enumerate}
\def\labelenumi{\arabic{enumi}.}
\setcounter{enumi}{24}
\tightlist
\item
  Ribosome biogenesis occurs in the nucleolus and is rate-limiting for
  cell growth. Evaluate the regulatory relationship between mTORC1, Pol
  I transcription of rDNA, and ribosome biogenesis. Why does mTOR
  inhibition suppress cell growth even in cells with unlimited amino
  acid supply?
\end{enumerate}

\begin{enumerate}
\def\labelenumi{\arabic{enumi}.}
\setcounter{enumi}{25}
\tightlist
\item
  Evaluate the molecular mechanism of translation by the ribosome in the
  context of the ``peptidyl transferase'' ribosome activity:
  specifically, is the catalysis chemical (RNA catalysis) or entropic
  (positioning effect)? Cite the key experiment (Nobel Prize 2009) that
  settled this debate.
\end{enumerate}

\begin{enumerate}
\def\labelenumi{\arabic{enumi}.}
\setcounter{enumi}{26}
\tightlist
\item
  Single-molecule translation experiments using fluorescently labelled
  ribosomes and mRNAs have revealed ``pausing'' during elongation
  opposite rare codons, and ``synchronised'' waves of ribosome density
  on mRNA. Evaluate what these observations reveal about codon usage
  optimisation and how synonymous mutations (same amino acid, different
  codon) can dramatically affect protein folding by altering translation
  speed.
\end{enumerate}

\begin{enumerate}
\def\labelenumi{\arabic{enumi}.}
\setcounter{enumi}{27}
\tightlist
\item
  Design a gene therapy strategy for a patient with haemophilia A
  (lacking factor VIII). Compare: mRNA therapy (synthetic mRNA with
  pseudo-uridine modification for stability); AAV vector gene therapy
  (liver-tropism, long-term expression); and base editing (correct the
  point mutation in situ). Evaluate the trade-offs in immunogenicity,
  duration of effect, and delivery efficiency.
\end{enumerate}

\begin{enumerate}
\def\labelenumi{\arabic{enumi}.}
\setcounter{enumi}{28}
\tightlist
\item
  RNA interference (RNAi) was discovered when dsRNA triggered gene
  silencing in \emph{C. elegans} (Fire and Mello, 1998). Evaluate the
  mechanism (Dicer, RISC, Ago2-mediated cleavage), the evolutionary
  origin of RNAi (antiviral defence), and the clinical development of
  siRNA therapeutics (inclisiran for PCSK9 silencing). What advantages
  does siRNA have over antisense oligonucleotides (ASOs) for gene
  silencing?
\end{enumerate}

\begin{enumerate}
\def\labelenumi{\arabic{enumi}.}
\setcounter{enumi}{29}
\tightlist
\item
  Critically evaluate the concept of the ``central dogma'' (DNA \texorpdfstring{\ensuremath{\to}}{->} RNA \texorpdfstring{\ensuremath{\to}}{->}
  protein). Identify three biological exceptions that violate a strict
  reading of the central dogma (reverse transcription, RNA editing,
  prion replication), explain their mechanisms, and assess whether these
  exceptions require a fundamental revision of the central dogma or
  merely represent elaborations within the same framework.
\end{enumerate}

\newpage

\section*{Questions --- Ch 14: Mutations, CRISPR, and
Genomics}\label{questions-ch-14-mutations-crispr-and-genomics}
\addcontentsline{toc}{section}{Questions --- Ch 14: Mutations, CRISPR,
and Genomics}

\label{sec:q_unit_IV_mutations_and_genomics}

\subsection*{Instructor Use and Coverage
Notes}\label{instructor-use-and-coverage-notes-16}
\addcontentsline{toc}{subsection}{Instructor Use and Coverage Notes}

\begin{itemize}
\tightlist
\item
  \textbf{Coverage target:} Predict the impact of a variant and justify
  the evidence needed to validate it.
\item
  \textbf{Model/data emphasis:} Mutation-rate, edit-efficiency, and
  sequence-comparison calculations.
\item
  \textbf{Assessment alignment:} Concept Explanation, Questions and
  Methods, Argumentation.
\item
  \textbf{Misconception probe:} Not every mutation is harmful, and not
  every harmful mutation changes a protein sequence.
\item
  \textbf{Transfer product:} Transfer variant reasoning to cancer
  genomics, ancestry, gene therapy, or microbial evolution.
\item
  \textbf{Grading focus:} award full credit for mechanism, evidence,
  boundary conditions, and units when a calculation is required.
\item
  \textbf{Suggested use:} draw one recall item, one application item,
  and one synthesis item when building a short quiz from this bank.
\end{itemize}

\subsection*{Questions 1--10: Recall and
Comprehension}\label{questions-110-recall-and-comprehension-13}
\addcontentsline{toc}{subsection}{Questions 1--10: Recall and
Comprehension}

\emph{This activity accompanies
\cref{sec:unit_IV_mutations_and_genomics} of the textbook --- review
that chapter before attempting the exercises below.}

\begin{enumerate}
\def\labelenumi{\arabic{enumi}.}
\tightlist
\item
  Define a mutation. Classify mutations by: (a) type of base change; (b)
  effect on the protein.
\end{enumerate}

\begin{enumerate}
\def\labelenumi{\arabic{enumi}.}
\setcounter{enumi}{1}
\tightlist
\item
  What is a frameshift mutation? How does an insertion vs deletion cause
  this?
\end{enumerate}

\begin{enumerate}
\def\labelenumi{\arabic{enumi}.}
\setcounter{enumi}{2}
\tightlist
\item
  What are transposons? Distinguish replicative (copy-and-paste) from
  non-replicative (cut-and-paste) transposons.
\end{enumerate}

\begin{enumerate}
\def\labelenumi{\arabic{enumi}.}
\setcounter{enumi}{3}
\tightlist
\item
  What is the Ames test? What does a positive result indicate?
\end{enumerate}

\begin{enumerate}
\def\labelenumi{\arabic{enumi}.}
\setcounter{enumi}{4}
\tightlist
\item
  Describe the three types of DNA damage repair: (a) nucleotide excision
  repair (NER); (b) base excision repair (BER); (c) mismatch repair
  (MMR).
\end{enumerate}

\begin{enumerate}
\def\labelenumi{\arabic{enumi}.}
\setcounter{enumi}{5}
\tightlist
\item
  What is CRISPR-Cas9? Where does CRISPR originate in bacteria?
\end{enumerate}

\begin{enumerate}
\def\labelenumi{\arabic{enumi}.}
\setcounter{enumi}{6}
\tightlist
\item
  What is gRNA (guide RNA) and what does it do in the Cas9 complex?
\end{enumerate}

\begin{enumerate}
\def\labelenumi{\arabic{enumi}.}
\setcounter{enumi}{7}
\tightlist
\item
  Define the PAM sequence. Why is it required for CRISPR cutting?
\end{enumerate}

\begin{enumerate}
\def\labelenumi{\arabic{enumi}.}
\setcounter{enumi}{8}
\tightlist
\item
  What is whole-genome sequencing (WGS)? What sequencing technology is
  most commonly used?
\end{enumerate}

\begin{enumerate}
\def\labelenumi{\arabic{enumi}.}
\setcounter{enumi}{9}
\tightlist
\item
  What is a single nucleotide polymorphism (SNP)? How many SNPs does the
  average human genome contain?
\end{enumerate}

\subsection*{Questions 11--20: Application and
Analysis}\label{questions-1120-application-and-analysis-13}
\addcontentsline{toc}{subsection}{Questions 11--20: Application and
Analysis}

\begin{enumerate}
\def\labelenumi{\arabic{enumi}.}
\setcounter{enumi}{10}
\tightlist
\item
  A chemotherapy drug causes G:C \texorpdfstring{\ensuremath{\to}}{->} A:T transitions. What type of DNA
  damage mechanism does this suggest (alkylation, intercalation,
  cross-linking, or depurination)? What repair pathway would normally
  correct this damage?
\end{enumerate}

\begin{enumerate}
\def\labelenumi{\arabic{enumi}.}
\setcounter{enumi}{11}
\tightlist
\item
  Xeroderma pigmentosum (XP) patients have mutations in NER pathway
  genes (XPA-XPG). Explain: (a) what lesion NER repairs (CPDs from
  UV-B); (b) why XP patients have dramatically increased skin cancer
  risk; (c) why the nervous system is also affected in some XP subtypes.
\end{enumerate}

\begin{enumerate}
\def\labelenumi{\arabic{enumi}.}
\setcounter{enumi}{12}
\tightlist
\item
  A researcher designs a CRISPR experiment to correct a point mutation
  (G\texorpdfstring{\ensuremath{\to}}{->}A) in the \emph{HBB} gene causing sickle-cell disease. List: (a)
  the guide RNA design rules; (b) the delivery method for cells (RNP
  electroporation vs viral vector); (c) why homology-directed repair
  (HDR, requiring a donor template) is needed rather than NHEJ; (d) how
  to increase HDR efficiency.
\end{enumerate}

\begin{enumerate}
\def\labelenumi{\arabic{enumi}.}
\setcounter{enumi}{13}
\tightlist
\item
  A patient's tumour has microsatellite instability (MSI-high) due to
  MMR deficiency. Explain: (a) what kinds of mutations accumulate in
  MSI-high tumours; (b) why MSI tumours produce many neoantigens; (c)
  why pembrolizumab (anti-PD-1) is effective in MSI-high tumours
  regardless of cancer type (tissue-agnostic treatment).
\end{enumerate}

\begin{enumerate}
\def\labelenumi{\arabic{enumi}.}
\setcounter{enumi}{14}
\tightlist
\item
  Transposaon Alu (SINE family) comprise about 11\% of the
  human genome. Explain: (a) how Alu elements replicate via RNA
  intermediates (SINE retrotransposition); (b) how an Alu insertion in
  exon 16 of BRCA2 could cause breast cancer; (c) how the host genome
  suppresses transposon activity in the germline (piRNA pathway).
\end{enumerate}

\begin{enumerate}
\def\labelenumi{\arabic{enumi}.}
\setcounter{enumi}{15}
\tightlist
\item
  GWAS (genome-wide association studies) have identified thousands of
  SNPs associated with disease risk, but most are in non-coding regions.
  Explain: (a) why non-coding SNPs are biologically meaningful (enhancer
  activity, splicing regulation); (b) how eQTL analysis links SNPs to
  gene expression levels; (c) why a SNP with odds ratio 1.2 for Type 2
  Diabetes is clinically uninformative for individual risk prediction.
\end{enumerate}

\begin{enumerate}
\def\labelenumi{\arabic{enumi}.}
\setcounter{enumi}{16}
\tightlist
\item
  The P53 gene is mutated in \textgreater50\% of all human cancers.
  Explain: (a) the ``gain-of-function'' mutations in TP53 (R248W, R175H)
  that not only eliminate tumour suppression but actively promote
  invasion; (b) how wild-type p53 regulates CDK inhibitor p21 and
  pro-apoptotic PUMA; (c) what therapeutic strategy (restoring wild-type
  TP53 conformation using small molecules --- APR-246) is being tested.
\end{enumerate}

\begin{enumerate}
\def\labelenumi{\arabic{enumi}.}
\setcounter{enumi}{17}
\tightlist
\item
  Gene therapy for Duchenne muscular dystrophy (DMD, caused by large
  deletions in \emph{Dystrophin}) uses exon skipping with antisense
  oligonucleotides (AONs). Explain: (a) why some deletions cause severe
  DMD while frameshifted deletions are less severe (in-frame vs
  out-of-frame); (b) how AON targeting exon 51 restores the reading
  frame; (c) why this produces ``Becker-like'' dystrophin rather than
  fully corrected dystrophin.
\end{enumerate}

\begin{enumerate}
\def\labelenumi{\arabic{enumi}.}
\setcounter{enumi}{18}
\tightlist
\item
  Epigenome-wide association studies (EWAS) examine DNA methylation
  differences between cases and controls. Explain: (a) why methylation
  differences (epimutations) are harder to interpret than genetic
  mutations; (b) why cell-type composition confounds EWAS analysis; (c)
  how MR (Mendelian randomisation) can be used to determine whether
  methylation changes are causal or merely correlative with disease.
\end{enumerate}

\begin{enumerate}
\def\labelenumi{\arabic{enumi}.}
\setcounter{enumi}{19}
\tightlist
\item
  Oxford Nanopore sequencing reads individual DNA/RNA molecules by
  measuring ionic current changes as they pass through a protein pore.
  Compare this technology to Illumina short-read sequencing: advantages
  (read length, direct RNA, base modifications, real-time),
  disadvantages (raw error rate about 5\% vs \textless0.1\%),
  and the application where each is most suitable.
\end{enumerate}

\subsection*{Questions 21--30: Synthesis and
Evaluation}\label{questions-2130-synthesis-and-evaluation-13}
\addcontentsline{toc}{subsection}{Questions 21--30: Synthesis and
Evaluation}

\begin{enumerate}
\def\labelenumi{\arabic{enumi}.}
\setcounter{enumi}{20}
\tightlist
\item
  Evaluate the safety of germline CRISPR editing (modifying embryo DNA,
  producing heritable changes) as attempted by He Jiankui in 2018.
  Identify three specific scientific risks (off-target edits, mosaicism,
  unintended immune consequences) and three ethical concerns (informed
  consent, enhancement vs therapy, equity of access). What international
  governance frameworks should govern germline editing?
\end{enumerate}

\begin{enumerate}
\def\labelenumi{\arabic{enumi}.}
\setcounter{enumi}{21}
\tightlist
\item
  A patient carries a disease-associated haplotype in which short reads
  map poorly because a 3-kb insertion and nearby promoter variant are
  absent from the linear reference. Design a graph-aware interpretation
  workflow: compare a single-reference call, a pangenome graph path, and
  long-read validation; then explain when base editing, prime editing,
  or Casgevy-style enhancer disruption would be the most defensible
  therapeutic logic.
\end{enumerate}

\begin{enumerate}
\def\labelenumi{\arabic{enumi}.}
\setcounter{enumi}{22}
\tightlist
\item
  Evaluate the ``cancer genome atlas'' concept: sequencing thousands of
  tumour genomes has generated a mutation landscape for most cancer
  types, identifying ``drivers'' vs ``passengers.'' What statistical
  criteria distinguish drivers from passengers, and why have most cancer
  driver mutations proven difficult to target therapeutically?
\end{enumerate}

\begin{enumerate}
\def\labelenumi{\arabic{enumi}.}
\setcounter{enumi}{23}
\tightlist
\item
  Synthetic chromosomes have been designed and built in yeast (Sc2.0,
  synthetic chromosome project). Evaluate the design principles (removal
  of transposons, recombination tags added, inducible evolution mode via
  SCRaMbLE), and discuss whether a fully synthetic human chromosome
  could be built and what it would require.
\end{enumerate}

\begin{enumerate}
\def\labelenumi{\arabic{enumi}.}
\setcounter{enumi}{24}
\tightlist
\item
  Transposons in the germline are silenced by piRNAs (PIWI-interacting
  RNAs). The piRNA pathway (PIWI proteins, ping-pong amplification)
  produces a feed-forward silencing loop. Evaluate: (a) how piRNAs
  recognise transposons without classical Watson-Crick base pairing; (b)
  why depletion of PIWI proteins causes germline transposon mobilisation
  and infertility in mice; (c) whether piRNAs can also target
  non-transposon mRNAs.
\end{enumerate}

\begin{enumerate}
\def\labelenumi{\arabic{enumi}.}
\setcounter{enumi}{25}
\tightlist
\item
  Precision oncology based on tumour sequencing identifies actionable
  mutations. However, intratumoral heterogeneity (ITH) means that some
  cells within the tumour may lack the targeted mutation. Evaluate: (a)
  how clonal vs subclonal mutations are distinguished by variant allele
  frequency; (b) why targeting a clonal mutation (present in most cells)
  is more likely to be effective; (c) how liquid biopsy (circulating
  tumour DNA, ctDNA) can detect resistance mechanisms before clinical
  progression.
\end{enumerate}

\begin{enumerate}
\def\labelenumi{\arabic{enumi}.}
\setcounter{enumi}{26}
\tightlist
\item
  Evaluate whether ``de-extinction'' (reviving extinct species via
  CRISPR genome editing of a close relative) is scientifically feasible
  for: (a) the woolly mammoth (editing \emph{Loxodonta africana} genome
  about 50 key cold-adaptation genes); (b) the passenger pigeon
  (\emph{Ectopistes migratorius}, using band-tailed pigeon as template).
  What are the ecological risks of reintroducing de-extinct animals?
\end{enumerate}

\begin{enumerate}
\def\labelenumi{\arabic{enumi}.}
\setcounter{enumi}{27}
\tightlist
\item
  Retroviral integration (HIV integrates into host genome) is both a
  challenge for cure research and a tool for gene therapy vectors.
  Evaluate the ``semi-random'' integration pattern of HIV (preference
  for active transcription units) and how this differs from AAV
  (predominantly episomal in non-dividing cells, rare integration at
  AAVS1 safe harbour). Why does insertional mutagenesis risk differ
  between retroviral and AAV gene therapies?
\end{enumerate}

\begin{enumerate}
\def\labelenumi{\arabic{enumi}.}
\setcounter{enumi}{28}
\tightlist
\item
  tRNA suppressors can read through stop codons, restoring protein
  production from genes with premature stop mutations. Evaluate the
  therapeutic use of aminoglycoside antibiotics (gentamicin) as
  readthrough agents for nonsense mutations (e.g., in \emph{Dystrophin})
  --- what are the molecular mechanisms, the challenges of specificity
  (also read through normal stop codons), and the clinical results to
  date?
\end{enumerate}

\begin{enumerate}
\def\labelenumi{\arabic{enumi}.}
\setcounter{enumi}{29}
\tightlist
\item
  Critically evaluate the technical and ethical implications of building
  a ``minimal synthetic cell'' with a fully defined genome (as in the
  Venter Institute's JCVI-syn3.0, about 473 genes). What does
  the set of essential genes reveal about the minimal requirements for
  cellular life? If this minimal genome could be commercially
  manufactured and sold, what biosafety and biosecurity regulations
  should apply?
\end{enumerate}

\newpage

\section*{Questions --- Ch 15: Epigenetics and Gene
Regulation}\label{questions-ch-15-epigenetics-and-gene-regulation}
\addcontentsline{toc}{section}{Questions --- Ch 15: Epigenetics and Gene
Regulation}

\label{sec:q_unit_IV_epigenetics_and_gene_regulation}

\subsection*{Instructor Use and Coverage
Notes}\label{instructor-use-and-coverage-notes-17}
\addcontentsline{toc}{subsection}{Instructor Use and Coverage Notes}

\begin{itemize}
\tightlist
\item
  \textbf{Coverage target:} Explain how a regulatory mark changes
  accessibility and predict downstream expression.
\item
  \textbf{Model/data emphasis:} Regulatory-state and expression-ratio
  reasoning.
\item
  \textbf{Assessment alignment:} Concept Explanation, Questions and
  Methods, Argumentation.
\item
  \textbf{Misconception probe:} Epigenetic does not mean independent of
  DNA sequence or permanently inherited.
\item
  \textbf{Transfer product:} Apply regulation logic to differentiation,
  imprinting, cancer, or environmental responses.
\item
  \textbf{Grading focus:} award full credit for mechanism, evidence,
  boundary conditions, and units when a calculation is required.
\item
  \textbf{Suggested use:} draw one recall item, one application item,
  and one synthesis item when building a short quiz from this bank.
\end{itemize}

\subsection*{Questions 1--10: Recall and
Comprehension}\label{questions-110-recall-and-comprehension-14}
\addcontentsline{toc}{subsection}{Questions 1--10: Recall and
Comprehension}

\emph{This activity accompanies
\cref{sec:unit_IV_epigenetics_and_gene_regulation} of the textbook ---
review that chapter before attempting the exercises below.}

\begin{enumerate}
\def\labelenumi{\arabic{enumi}.}
\tightlist
\item
  Define epigenetics. Give two examples of epigenetic modifications.
\end{enumerate}

\begin{enumerate}
\def\labelenumi{\arabic{enumi}.}
\setcounter{enumi}{1}
\tightlist
\item
  What is DNA methylation? Which nucleotide is typically methylated in
  mammals, and what effect does methylation at gene promoters generally
  have on transcription?
\end{enumerate}

\begin{enumerate}
\def\labelenumi{\arabic{enumi}.}
\setcounter{enumi}{2}
\tightlist
\item
  What is a histone? Name four types of histone modification and one
  enzyme family for each.
\end{enumerate}

\begin{enumerate}
\def\labelenumi{\arabic{enumi}.}
\setcounter{enumi}{3}
\tightlist
\item
  What does H3K4me3 mark indicate? What about H3K27me3?
\end{enumerate}

\begin{enumerate}
\def\labelenumi{\arabic{enumi}.}
\setcounter{enumi}{4}
\tightlist
\item
  What is chromatin remodelling? Name one chromatin remodelling complex.
\end{enumerate}

\begin{enumerate}
\def\labelenumi{\arabic{enumi}.}
\setcounter{enumi}{5}
\tightlist
\item
  What is the lac operon? Which molecules control it?
\end{enumerate}

\begin{enumerate}
\def\labelenumi{\arabic{enumi}.}
\setcounter{enumi}{6}
\tightlist
\item
  What is a repressor protein? What is an activator protein?
\end{enumerate}

\begin{enumerate}
\def\labelenumi{\arabic{enumi}.}
\setcounter{enumi}{7}
\tightlist
\item
  Define allosteric regulation in the context of transcription. Give an
  example.
\end{enumerate}

\begin{enumerate}
\def\labelenumi{\arabic{enumi}.}
\setcounter{enumi}{8}
\tightlist
\item
  What are non-coding RNAs? Name three biologically important types.
\end{enumerate}

\begin{enumerate}
\def\labelenumi{\arabic{enumi}.}
\setcounter{enumi}{9}
\tightlist
\item
  What is genomic imprinting? Give one example of an imprinted gene.
\end{enumerate}

\subsection*{Questions 11--20: Application and
Analysis}\label{questions-1120-application-and-analysis-14}
\addcontentsline{toc}{subsection}{Questions 11--20: Application and
Analysis}

\begin{enumerate}
\def\labelenumi{\arabic{enumi}.}
\setcounter{enumi}{10}
\tightlist
\item
  A CpG island in the CDKN2A (p16) promoter is hypermethylated in a lung
  cancer cell. Trace the molecular consequences: (a) which protein is
  silenced; (b) what happens to CDK4/6 activity; (c) what happens to
  cell cycle progression; (d) why this is an alternative mechanism to
  p16 deletion for bypassing the G\texorpdfstring{\textsubscript{1}}{1} checkpoint.
\end{enumerate}

\begin{enumerate}
\def\labelenumi{\arabic{enumi}.}
\setcounter{enumi}{11}
\tightlist
\item
  Polycomb repressive complex 2 (PRC2) trimethylates H3K27 and is
  overexpressed in many cancers. Explain: (a) the molecular mechanism of
  H3K27me3-mediated gene silencing; (b) how PRC2 spreads repression
  across large chromosomal domains; (c) why EZH2 (the catalytic subunit
  of PRC2) is a therapeutic target (tazemetostat).
\end{enumerate}

\begin{enumerate}
\def\labelenumi{\arabic{enumi}.}
\setcounter{enumi}{12}
\tightlist
\item
  In the lac operon, explain the logic of dual control by the lac
  repressor AND catabolite activator protein (CAP). Under which
  combination of conditions is the operon most highly transcribed? Under
  which is it completely silent?
\end{enumerate}

\begin{enumerate}
\def\labelenumi{\arabic{enumi}.}
\setcounter{enumi}{13}
\tightlist
\item
  RNA interference (RNAi) via siRNA silences gene expression
  post-transcriptionally. Compare siRNA (perfectly complementary \texorpdfstring{\ensuremath{\to}}{->} mRNA
  cleavage by Ago2) vs miRNA (partially complementary \texorpdfstring{\ensuremath{\to}}{->} translational
  repression and mRNA destabilisation). How does each affect the
  steady-state abundance of target mRNA?
\end{enumerate}

\begin{enumerate}
\def\labelenumi{\arabic{enumi}.}
\setcounter{enumi}{14}
\tightlist
\item
  Enhancers can activate promoters over distances of \textgreater1 Mb in
  the genome by chromatin looping mediated by cohesin and CTCF. Explain:
  (a) the role of CTCF as a chromatin barrier/insulator; (b) how loss of
  a CTCF binding site at the IGF2/H19 imprinting control region (ICR)
  causes Beckwith-Wiedemann syndrome; (c) how 4C and HiC experiments
  confirm enhancer-promoter looping.
\end{enumerate}

\begin{enumerate}
\def\labelenumi{\arabic{enumi}.}
\setcounter{enumi}{15}
\tightlist
\item
  The X-inactivation centre (Xic) on the inactive X chromosome produces
  Xist lncRNA, which coats the inactive X in cis. Explain: (a) how Xist
  RNA spreads along the chromosome; (b) how it recruits PRC2 (H3K27me3)
  and PRC1 (H2AK119ub1) to compact the Xi; (c) why a small region
  (pseudo-autosomal region, PAR) escapes X-inactivation.
\end{enumerate}

\begin{enumerate}
\def\labelenumi{\arabic{enumi}.}
\setcounter{enumi}{16}
\tightlist
\item
  A researcher uses Cas9 fused to a DNMT3A (DNA methyltransferase)
  domain (dCas9-DNMT3A) to specifically methylate a target promoter.
  Predict: (a) the epigenetic effect at the target gene; (b) whether the
  silencing will be maintained through DNA replication (yes, DNMT1
  maintenance methyltransferase copies the pattern); (c) how you would
  verify that methylation, not dCas9 binding alone, is responsible for
  silencing.
\end{enumerate}

\begin{enumerate}
\def\labelenumi{\arabic{enumi}.}
\setcounter{enumi}{17}
\tightlist
\item
  Transgenerational epigenetic inheritance refers to the transmission of
  epigenetic marks across generations without DNA sequence change. The
  agouti mouse model showed that maternal methyl-donor diet changes
  offspring coat colour and obesity risk via CpG methylation of the Avy
  allele. What does this imply about the erasure of epigenetic marks
  between generations (which normally occurs in primordial germ cells)?
\end{enumerate}

\begin{enumerate}
\def\labelenumi{\arabic{enumi}.}
\setcounter{enumi}{18}
\tightlist
\item
  Long non-coding RNAs (lncRNAs) number \textgreater100,000 in the human
  transcriptome but most have unknown functions. Evaluate the challenges
  of studying lncRNA function: (a) conservation is low across species
  (cannot use model organisms easily); (b) knockdown may affect the
  genomic region itself (enhancer function of the locus); (c) lncRNAs
  have complex secondary structure but no ORFs. Propose a systematic
  approach to determine function.
\end{enumerate}

\begin{enumerate}
\def\labelenumi{\arabic{enumi}.}
\setcounter{enumi}{19}
\tightlist
\item
  Chromatin accessibility (measured by ATAC-seq) is strongly correlated
  with gene expression. Explain: (a) how nucleosome positioning
  regulates transcription factor accessibility; (b) the role of pioneer
  transcription factors (e.g., FOXA1) in opening closed chromatin; (c)
  how ATAC-seq works mechanistically (Tn5 transposase cuts accessible
  DNA and adds sequencing adapters).
\end{enumerate}

\subsection*{Questions 21--30: Synthesis and
Evaluation}\label{questions-2130-synthesis-and-evaluation-14}
\addcontentsline{toc}{subsection}{Questions 21--30: Synthesis and
Evaluation}

\begin{enumerate}
\def\labelenumi{\arabic{enumi}.}
\setcounter{enumi}{20}
\tightlist
\item
  Evaluate the claim that cancer is largely an epigenetic disease. Cite
  evidence from: (a) global hypomethylation and CpG island
  hypermethylation in cancer; (b) frequent mutations in chromatin
  remodelling enzymes (ARID1A, SMARCA4, KMT2D) in cancer; (c) the
  success of HDAC inhibitors and DNMT inhibitors clinically. Under what
  conditions are purely epigenetic therapies (without targeted mutation
  correction) appropriate?
\end{enumerate}

\begin{enumerate}
\def\labelenumi{\arabic{enumi}.}
\setcounter{enumi}{21}
\tightlist
\item
  Design a pharmaceutical strategy to reactivate a silenced tumour
  suppressor gene in a cancer with promoter hypermethylation. Compare
  DNMT inhibitors (5-azacitidine, broad demethylation), targeted
  dCas9-TET1 (site-specific demethylation), and HDAC inhibitors. What
  are the risks of broad epigenome reprogramming vs targeted approaches?
\end{enumerate}

\begin{enumerate}
\def\labelenumi{\arabic{enumi}.}
\setcounter{enumi}{22}
\tightlist
\item
  Super-enhancers are clusters of enhancers (about 25--50 kb)
  bound by an unusually high density of transcription factors (BRD4,
  Mediator), associated with strong transcription of oncogenes (MYC,
  BCL2) in cancer. Evaluate the evidence for super-enhancer addiction in
  cancer cells, and explain why BET bromodomain inhibitors (JQ1)
  preferentially disengage from super-enhancers over typical enhancers.
\end{enumerate}

\begin{enumerate}
\def\labelenumi{\arabic{enumi}.}
\setcounter{enumi}{23}
\tightlist
\item
  Evaluate the concept of ``epigenetic clock'' (Horvath clock):
  methylation of specific CpG sites in blood DNA correlates with
  biological age more accurately than chronological age. What do
  deviations from the epigenetic clock predict? Can lifestyle
  interventions (caloric restriction, exercise) reduce epigenetic age,
  and what experimental evidence supports this?
\end{enumerate}

\begin{enumerate}
\def\labelenumi{\arabic{enumi}.}
\setcounter{enumi}{24}
\tightlist
\item
  3D chromatin organisation (TADs --- topologically associating domains)
  is largely maintained by CTCF and cohesin. Evaluate what would happen
  to gene regulation if cohesin loading were abolished globally: (a)
  which enhancer-promoter pairs would become aberrantly activated; (b)
  what new gene expression states would emerge; (c) is this consistent
  with the observation that cohesin mutations cause Cornelia de Lange
  syndrome?
\end{enumerate}

\begin{enumerate}
\def\labelenumi{\arabic{enumi}.}
\setcounter{enumi}{25}
\tightlist
\item
  CRISPR interference (CRISPRi) uses dCas9-KRAB (a transcriptional
  repressor) to silences target genes without cutting DNA. Compare
  CRISPRi with ASO (antisense oligonucleotide) knockdown at the mRNA
  level and classical promoter deletion by CRISPR. What are the
  therapeutic advantages of CRISPRi (temporal control, reversibility, no
  off-target DSBs)?
\end{enumerate}

\begin{enumerate}
\def\labelenumi{\arabic{enumi}.}
\setcounter{enumi}{26}
\tightlist
\item
  Maternal nutrition during peri-conception affects offspring health
  across the life course (Barker/DOHaD hypothesis). Using three specific
  nutrients (folate, methionine, choline) and their roles in one-carbon
  metabolism \texorpdfstring{\ensuremath{\to}}{->} SAM production \texorpdfstring{\ensuremath{\to}}{->} methyltransferase substrate supply,
  evaluate the molecular mechanism by which poor maternal nutrition
  could cause persistent epigenetic changes in offspring that increase
  adult diabetes or cardiovascular disease risk.
\end{enumerate}

\begin{enumerate}
\def\labelenumi{\arabic{enumi}.}
\setcounter{enumi}{27}
\tightlist
\item
  Evaluate the concept of ``bivalent chromatin domains'' in embryonic
  stem cells: developmental gene promoters carry both active (H3K4me3)
  and repressive (H3K27me3) marks simultaneously, maintaining them in a
  ``poised'' state. How does this enable rapid differentiation in either
  direction, and what happens to bivalent domains as cells commit to a
  specific lineage?
\end{enumerate}

\begin{enumerate}
\def\labelenumi{\arabic{enumi}.}
\setcounter{enumi}{28}
\tightlist
\item
  Fragile X syndrome is caused by CGG repeat expansion (\textgreater200
  repeats) in the \emph{FMR1} gene's 5' UTR, leading to hypermethylation
  and gene silencing. Evaluate the molecular relationship between repeat
  expansion and hypermethylation: which comes first, what protein
  (MBNL1/DNMT3A) reads the repeat expansion, and why does the
  ``premutation'' (55--200 repeats) have a partially different phenotype
  (FXTAS) compared to full mutation?
\end{enumerate}

\begin{enumerate}
\def\labelenumi{\arabic{enumi}.}
\setcounter{enumi}{29}
\tightlist
\item
  Critically evaluate the potential of epigenetic editing as a permanent
  therapeutic intervention. Specifically: if dCas9-DNMT3A is used to
  methylate and silence an oncogene, will the silencing persist through
  cell division (via DNMT1 maintenance methylation), and what are the
  risks of epigenetic drift or passenger methylation at off-target
  sites? Propose experimental criteria to determine whether a
  therapeutic epigenetic edit is ``safe and durable.''
\end{enumerate}

\newpage

\section*{Questions --- Ch 16: Mendelian Genetics and
Heredity}\label{questions-ch-16-mendelian-genetics-and-heredity}
\addcontentsline{toc}{section}{Questions --- Ch 16: Mendelian Genetics
and Heredity}

\label{sec:q_unit_V_mendelian_genetics}

\subsection*{Instructor Use and Coverage
Notes}\label{instructor-use-and-coverage-notes-18}
\addcontentsline{toc}{subsection}{Instructor Use and Coverage Notes}

\begin{itemize}
\tightlist
\item
  \textbf{Coverage target:} Compute a genetic probability or chi-square
  result and explain the biological assumption.
\item
  \textbf{Model/data emphasis:} Punnett, binomial, pedigree, and
  chi-square calculations.
\item
  \textbf{Assessment alignment:} Statistical Tests and Data Analysis,
  Representing and Describing Data.
\item
  \textbf{Misconception probe:} Dominant does not mean common, stronger,
  or better.
\item
  \textbf{Transfer product:} Transfer inheritance reasoning to
  counseling, breeding, model organisms, and risk prediction.
\item
  \textbf{Grading focus:} award full credit for mechanism, evidence,
  boundary conditions, and units when a calculation is required.
\item
  \textbf{Suggested use:} draw one recall item, one application item,
  and one synthesis item when building a short quiz from this bank.
\end{itemize}

\subsection*{Questions 1--10: Recall and
Comprehension}\label{questions-110-recall-and-comprehension-15}
\addcontentsline{toc}{subsection}{Questions 1--10: Recall and
Comprehension}

\emph{This activity accompanies \cref{sec:unit_V_mendelian_genetics} of
the textbook --- review that chapter before attempting the exercises
below.}

\begin{enumerate}
\def\labelenumi{\arabic{enumi}.}
\tightlist
\item
  State Mendel's two laws. Give the biological basis for each in terms
  of meiosis.
\end{enumerate}

\begin{enumerate}
\def\labelenumi{\arabic{enumi}.}
\setcounter{enumi}{1}
\tightlist
\item
  Define allele, locus, genotype, and phenotype.
\end{enumerate}

\begin{enumerate}
\def\labelenumi{\arabic{enumi}.}
\setcounter{enumi}{2}
\tightlist
\item
  What is a Punnett square? Draw and solve: Aa \texorpdfstring{\ensuremath{\times}}{x} Aa.
\end{enumerate}

\begin{enumerate}
\def\labelenumi{\arabic{enumi}.}
\setcounter{enumi}{3}
\tightlist
\item
  What are the expected phenotype ratios for a monohybrid cross between
  two heterozygotes?
\end{enumerate}

\begin{enumerate}
\def\labelenumi{\arabic{enumi}.}
\setcounter{enumi}{4}
\tightlist
\item
  What are the expected phenotype ratios for a dihybrid cross between
  two double heterozygotes?
\end{enumerate}

\begin{enumerate}
\def\labelenumi{\arabic{enumi}.}
\setcounter{enumi}{5}
\tightlist
\item
  Define codominance. Give the human blood type system as an example.
\end{enumerate}

\begin{enumerate}
\def\labelenumi{\arabic{enumi}.}
\setcounter{enumi}{6}
\tightlist
\item
  What is incomplete dominance? Give one example and show the F\texorpdfstring{\textsubscript{1}}{1} and F\texorpdfstring{\textsubscript{2}}{2}
  ratios.
\end{enumerate}

\begin{enumerate}
\def\labelenumi{\arabic{enumi}.}
\setcounter{enumi}{7}
\tightlist
\item
  What is a testcross? For what is it used?
\end{enumerate}

\begin{enumerate}
\def\labelenumi{\arabic{enumi}.}
\setcounter{enumi}{8}
\tightlist
\item
  What is epistasis? Give one example where two loci interact to modify
  phenotype.
\end{enumerate}

\begin{enumerate}
\def\labelenumi{\arabic{enumi}.}
\setcounter{enumi}{9}
\tightlist
\item
  What is pleiotropy? Give one example of a single gene that affects
  multiple traits.
\end{enumerate}

\subsection*{Questions 11--20: Application and
Analysis}\label{questions-1120-application-and-analysis-15}
\addcontentsline{toc}{subsection}{Questions 11--20: Application and
Analysis}

\begin{enumerate}
\def\labelenumi{\arabic{enumi}.}
\setcounter{enumi}{10}
\tightlist
\item
  A black cat (BB) is crossed with a white cat (bb). What is the F\texorpdfstring{\textsubscript{1}}{1}
  phenotype? What is the F\texorpdfstring{\textsubscript{2}}{2} ratio if F\texorpdfstring{\textsubscript{1}}{1} cats are interbred? If black is
  incompletely dominant, what is the F\texorpdfstring{\textsubscript{1}}{1} phenotype?
\end{enumerate}

\begin{enumerate}
\def\labelenumi{\arabic{enumi}.}
\setcounter{enumi}{11}
\tightlist
\item
  A chi-square analysis of 580 pea seeds (450 round:130 wrinkled) gives
  \texorpdfstring{\ensuremath{\chi}}{chi}\texorpdfstring{\textsuperscript{2}}{2} = 2.6 (critical value 3.84, p = 0.05). Should the null hypothesis
  of 3:1 ratio be rejected? What does this conclusion mean biologically?
\end{enumerate}

\begin{enumerate}
\def\labelenumi{\arabic{enumi}.}
\setcounter{enumi}{12}
\tightlist
\item
  Two true-breeding strains of wheat produce red and white kernels. The
  cross produces medium-red F\texorpdfstring{\textsubscript{1}}{1}, and F\texorpdfstring{\textsubscript{2}}{2} shows a continuous range of
  colours with frequencies approximately 1:4:6:4:1. This is consistent
  with two genes each contributing an additive unit of colour. What mode
  of inheritance does this illustrate?
\end{enumerate}

\begin{enumerate}
\def\labelenumi{\arabic{enumi}.}
\setcounter{enumi}{13}
\tightlist
\item
  A woman is phenotypically normal but her father had phenylketonuria
  (PKU, autosomal recessive). What is the probability that she is a
  carrier? If she marries a known carrier, what is the probability their
  first child will have PKU?
\end{enumerate}

\begin{enumerate}
\def\labelenumi{\arabic{enumi}.}
\setcounter{enumi}{14}
\tightlist
\item
  In coat colour genetics of Labrador retrievers: B is dominant for
  black pigment; b is recessive (brown/chocolate). E is dominant for
  coat expression; ee gives yellow regardless of B/b. Set up the
  dihybrid cross BbEe \texorpdfstring{\ensuremath{\times}}{x} BbEe and calculate the expected ratio of
  black:chocolate:yellow.
\end{enumerate}

\begin{enumerate}
\def\labelenumi{\arabic{enumi}.}
\setcounter{enumi}{15}
\tightlist
\item
  ABO blood type shows codominance (A and B antigens) and simple
  recessive (O, no antigen). A man of unknown blood type has a type O
  child with a type A woman (confirmed genotype IA\^{}i). What is the
  man's blood type? What are most possible blood types of this union's
  children?
\end{enumerate}

\begin{enumerate}
\def\labelenumi{\arabic{enumi}.}
\setcounter{enumi}{16}
\tightlist
\item
  Maternal effect genes in \emph{Drosophila} (e.g., \emph{bicoid})
  determine the embryo's polarity based on the maternal genotype, not
  the embryo's genotype. Explain how this violates the simple Mendelian
  expectation that offspring phenotype reflects offspring genotype. What
  does this tell us about developmental genetics?
\end{enumerate}

\begin{enumerate}
\def\labelenumi{\arabic{enumi}.}
\setcounter{enumi}{17}
\tightlist
\item
  Haemophilia A is X-linked recessive. A carrier woman (X\^{}H X\^{}h)
  has children with a normal man (X\^{}H Y). What are the expected
  phenotype ratios for daughters? For sons? What is the probability that
  the first four children include at least two affected sons (use
  binomial distribution)?
\end{enumerate}

\begin{enumerate}
\def\labelenumi{\arabic{enumi}.}
\setcounter{enumi}{18}
\tightlist
\item
  Set up and interpret the F\texorpdfstring{\textsubscript{2}}{2} results of a dihybrid cross (AaBb \texorpdfstring{\ensuremath{\times}}{x} AaBb)
  if the two genes show complementary epistasis (A-B- = one phenotype;
  A-bb, aaB-, aabb = albino). What is the expected ratio? What does this
  imply about the biochemical pathway?
\end{enumerate}

\begin{enumerate}
\def\labelenumi{\arabic{enumi}.}
\setcounter{enumi}{19}
\tightlist
\item
  Three pea plants of unknown genotype most produce round seeds. Two
  plants crossed to a wrinkled plant (rr) produce: plant A = most round;
  plant B = 50\% round:50\% wrinkled; plant C = 75\% round:25\%
  wrinkled. Determine the genotype of each plant and explain the
  crosses.
\end{enumerate}

\subsection*{Questions 21--30: Synthesis and
Evaluation}\label{questions-2130-synthesis-and-evaluation-15}
\addcontentsline{toc}{subsection}{Questions 21--30: Synthesis and
Evaluation}

\begin{enumerate}
\def\labelenumi{\arabic{enumi}.}
\setcounter{enumi}{20}
\tightlist
\item
  Evaluate the concept of ``genetic heterogeneity'': the same phenotype
  (e.g., congenital deafness) can result from mutations in any of
  \textgreater100 different genes, each segregating in different
  families. How does genetic heterogeneity complicate: (a) pedigree
  analysis; (b) population frequency estimates; (c) the development of a
  comprehensive gene therapy?
\end{enumerate}

\begin{enumerate}
\def\labelenumi{\arabic{enumi}.}
\setcounter{enumi}{21}
\tightlist
\item
  Design a genetic screen in \emph{C. elegans} to identify most genes
  involved in a specific pathway (e.g., vulva development). Describe the
  EMS (ethyl methanesulfonate) mutagenesis strategy, the phenotypic
  screen, and how you would determine whether two mutations affect the
  same gene (complementation test) or the same pathway (epistasis test).
\end{enumerate}

\begin{enumerate}
\def\labelenumi{\arabic{enumi}.}
\setcounter{enumi}{22}
\tightlist
\item
  Quantitative traits (height, intelligence, blood pressure) are
  determined by many genes plus environment. Evaluate the methodology of
  heritability estimation from twin studies: what does h\texorpdfstring{\textsuperscript{2}}{2} = 0.8 actually
  mean, and why does high heritability NOT mean that environmental
  interventions are ineffective? Use PKU as a concrete counter-example.
\end{enumerate}

\begin{enumerate}
\def\labelenumi{\arabic{enumi}.}
\setcounter{enumi}{23}
\tightlist
\item
  Evaluate Mendel's choice of pea traits: most seven traits happened to
  be either not linked or on different chromosomes. If two of Mendel's
  traits had been tightly linked (say, flower colour and pod colour on
  the same chromosome), how would the F\texorpdfstring{\textsubscript{2}}{2} data have deviated from the
  9:3:3:1 ratio, and would Mendel still have discovered independent
  assortment?
\end{enumerate}

\begin{enumerate}
\def\labelenumi{\arabic{enumi}.}
\setcounter{enumi}{24}
\tightlist
\item
  Inbreeding increases homozygosity and can reveal recessive deleterious
  alleles. Calculate the inbreeding coefficient (F) for offspring of a
  first-cousin mating, and predict the increase in: (a) homozygosity for
  a rare allele (q = 0.01); (b) probability of autosomal recessive
  disease compared to random mating. Why do many societies have legal or
  cultural prohibitions against close-relative marriages?
\end{enumerate}

\begin{enumerate}
\def\labelenumi{\arabic{enumi}.}
\setcounter{enumi}{25}
\tightlist
\item
  Anticipation (increasing severity across generations) occurs in
  trinucleotide repeat disorders (Huntington's disease, myotonic
  dystrophy). Explain the molecular mechanism by which repeats expand
  between generations (slippage during DNA replication) and why larger
  repeats cause earlier onset and greater severity. Why is repeat
  expansion more common through the paternal germline in HD but the
  maternal germline in DM1?
\end{enumerate}

\begin{enumerate}
\def\labelenumi{\arabic{enumi}.}
\setcounter{enumi}{26}
\tightlist
\item
  Evaluate the concept of ``genomic imprinting conflict'' (kinship
  theory): paternally imprinted genes (expressed primarily from paternal
  allele) are generally growth-promoting (\emph{IGF2}), while maternally
  imprinted genes (expressed primarily from maternal allele) are
  generally growth-restraining (\emph{H19}). Evaluate the evolutionary
  logic that predicts this based on parent-offspring conflict theory.
  What happens when this balance is disrupted (Silver-Russell syndrome
  vs Beckwith-Wiedemann syndrome)?
\end{enumerate}

\begin{enumerate}
\def\labelenumi{\arabic{enumi}.}
\setcounter{enumi}{27}
\tightlist
\item
  Multifactorial diseases (Type 1 Diabetes, schizophrenia) have both
  genetic and environmental contributions. Evaluate the strengths and
  limitations of GWAS and polygenic risk scores: (a) minimum detectable
  effect size at genome-wide significance (p \textless{} 5\texorpdfstring{\ensuremath{\times}}{x}10\texorpdfstring{\textsuperscript{-}}{-}\texorpdfstring{\textsuperscript{8}}{8}); (b)
  the ``missing heritability'' problem; (c) why PRS transferability
  across ancestries is a scientific and equity concern; (d) what consent
  and privacy information belongs in a clinical PRS report.
\end{enumerate}

\begin{enumerate}
\def\labelenumi{\arabic{enumi}.}
\setcounter{enumi}{28}
\tightlist
\item
  Phenocopies are environmentally caused conditions that resemble
  genetic conditions. Thalidomide-induced limb malformations in
  1950s--60s phenocopied the autosomal recessive condition phocomelia.
  Evaluate: (a) how phenocopies complicate pedigree analysis; (b) how
  you would distinguish a phenocopy from a genetic condition in a
  clinical or research context; (c) the ethical implications of
  confusing genetic and environmental causation for public health
  policy.
\end{enumerate}

\begin{enumerate}
\def\labelenumi{\arabic{enumi}.}
\setcounter{enumi}{29}
\tightlist
\item
  Critically evaluate Mendel's pea experiments, considering: (a) sample
  sizes (he used thousands of plants, giving high statistical power);
  (b) evidence that Mendel may have ``improved'' his data (Weldon's 1902
  re-analysis and Ronald Fisher's 1936 statistical critique suggesting
  perfect ratios are too good to be true); (c) whether Mendel's results
  are reproducible, and what this debate reveals about the role of
  confirmation bias in science.
\end{enumerate}

\newpage

\section*{Questions --- Ch 17: Chromosomal Inheritance and
Linkage}\label{questions-ch-17-chromosomal-inheritance-and-linkage}
\addcontentsline{toc}{section}{Questions --- Ch 17: Chromosomal
Inheritance and Linkage}

\label{sec:q_unit_V_chromosomal_inheritance}

\subsection*{Instructor Use and Coverage
Notes}\label{instructor-use-and-coverage-notes-19}
\addcontentsline{toc}{subsection}{Instructor Use and Coverage Notes}

\begin{itemize}
\tightlist
\item
  \textbf{Coverage target:} Use progeny classes to calculate map
  distances and infer the most likely gene order.
\item
  \textbf{Model/data emphasis:} Recombination frequency and three-point
  mapping.
\item
  \textbf{Assessment alignment:} Statistical Tests and Data Analysis,
  Representing and Describing Data.
\item
  \textbf{Misconception probe:} Independent assortment applies to
  unlinked loci, not to every pair of genes.
\item
  \textbf{Transfer product:} Transfer linkage reasoning to disease
  mapping, breeding, and genome assemblies.
\item
  \textbf{Grading focus:} award full credit for mechanism, evidence,
  boundary conditions, and units when a calculation is required.
\item
  \textbf{Suggested use:} draw one recall item, one application item,
  and one synthesis item when building a short quiz from this bank.
\end{itemize}

\subsection*{Questions 1--10: Recall and
Comprehension}\label{questions-110-recall-and-comprehension-16}
\addcontentsline{toc}{subsection}{Questions 1--10: Recall and
Comprehension}

\emph{This activity accompanies
\cref{sec:unit_V_chromosomal_inheritance} of the textbook --- review
that chapter before attempting the exercises below.}

\begin{enumerate}
\def\labelenumi{\arabic{enumi}.}
\tightlist
\item
  What is a chromosome? How many chromosomes does a normal human somatic
  cell contain?
\end{enumerate}

\begin{enumerate}
\def\labelenumi{\arabic{enumi}.}
\setcounter{enumi}{1}
\tightlist
\item
  Define linkage. Why do linked genes violate Mendel's law of
  independent assortment?
\end{enumerate}

\begin{enumerate}
\def\labelenumi{\arabic{enumi}.}
\setcounter{enumi}{2}
\tightlist
\item
  What is a centimorgan (cM)? What recombination frequency corresponds
  to 1 cM?
\end{enumerate}

\begin{enumerate}
\def\labelenumi{\arabic{enumi}.}
\setcounter{enumi}{3}
\tightlist
\item
  Describe the process of crossing over. At what stage of meiosis does
  it occur?
\end{enumerate}

\begin{enumerate}
\def\labelenumi{\arabic{enumi}.}
\setcounter{enumi}{4}
\tightlist
\item
  What is a genetic map? How is it constructed from recombination
  frequency data?
\end{enumerate}

\begin{enumerate}
\def\labelenumi{\arabic{enumi}.}
\setcounter{enumi}{5}
\tightlist
\item
  What is X-linked inheritance? Give two examples of X-linked recessive
  disorders.
\end{enumerate}

\begin{enumerate}
\def\labelenumi{\arabic{enumi}.}
\setcounter{enumi}{6}
\tightlist
\item
  Define nondisjunction. What aneuploid conditions result from
  nondisjunction at meiosis I vs meiosis II?
\end{enumerate}

\begin{enumerate}
\def\labelenumi{\arabic{enumi}.}
\setcounter{enumi}{7}
\tightlist
\item
  What is trisomy 21? What maternal-age-related risk factor is
  associated with it?
\end{enumerate}

\begin{enumerate}
\def\labelenumi{\arabic{enumi}.}
\setcounter{enumi}{8}
\tightlist
\item
  What is a chromosomal translocation? How does the Philadelphia
  chromosome arise?
\end{enumerate}

\begin{enumerate}
\def\labelenumi{\arabic{enumi}.}
\setcounter{enumi}{9}
\tightlist
\item
  What is a karyotype? How is one prepared?
\end{enumerate}

\subsection*{Questions 11--20: Application and
Analysis}\label{questions-1120-application-and-analysis-16}
\addcontentsline{toc}{subsection}{Questions 11--20: Application and
Analysis}

\begin{enumerate}
\def\labelenumi{\arabic{enumi}.}
\setcounter{enumi}{10}
\tightlist
\item
  Three linked genes A, B, C have pairwise recombination frequencies:
  A--B = 12\%, B--C = 7\%, A--C = 19\%. Does map distance A--C equal
  A--B + B--C? If not, why, and what genetic phenomenon accounts for the
  discrepancy?
\end{enumerate}

\begin{enumerate}
\def\labelenumi{\arabic{enumi}.}
\setcounter{enumi}{11}
\tightlist
\item
  Morgan's experiments showed that white eye and miniature wing in
  \emph{Drosophila} are X-linked and linked to each other (recombination
  frequency about 36\%). Explain how these results violate both
  Mendel's laws simultaneously, and how recombination still produces
  some parental and recombinant gamete classes.
\end{enumerate}

\begin{enumerate}
\def\labelenumi{\arabic{enumi}.}
\setcounter{enumi}{12}
\tightlist
\item
  A man with colour blindness (X-linked recessive) who has normal vision
  has a sister who is a confirmed carrier. What is the probability that
  his daughter will: (a) be a carrier; (b) be colour blind; (c) have
  normal vision? What is the probability any son of his will be colour
  blind?
\end{enumerate}

\begin{enumerate}
\def\labelenumi{\arabic{enumi}.}
\setcounter{enumi}{13}
\tightlist
\item
  Barr bodies (condensed inactive X chromosomes) are present in cells
  with \textgreater1 X chromosome: 1 Barr body per n+1 X chromosomes.
  How many Barr bodies would be present in cells from: (a) 46,XX female;
  (b) 47,XXX female (triple X); (c) 47,XXY male (Klinefelter); (d) 45,X
  female (Turner)?
\end{enumerate}

\begin{enumerate}
\def\labelenumi{\arabic{enumi}.}
\setcounter{enumi}{14}
\tightlist
\item
  Down syndrome (trisomy 21) risk rises with maternal age. Propose a
  mechanistic explanation: which proteins hold meiosis I arrested in the
  human oocyte, how does cohesin decline with age, and how does
  premature chromatid separation \texorpdfstring{\ensuremath{\to}}{->} nondisjunction \texorpdfstring{\ensuremath{\to}}{->} trisomy? Name the
  specific cohesin subunit implicated.
\end{enumerate}

\begin{enumerate}
\def\labelenumi{\arabic{enumi}.}
\setcounter{enumi}{15}
\tightlist
\item
  Robertsonian translocation Down syndrome: a chromosome 21 fuses to
  chromosome 14. A carrier parent (45 chromosomes, with 14;21
  translocation) produces offspring via meiosis. Using a diagram, show
  the three viable gamete types and predict the proportion of offspring
  who have translocation Down syndrome, normal karyotype, and carrier
  status.
\end{enumerate}

\begin{enumerate}
\def\labelenumi{\arabic{enumi}.}
\setcounter{enumi}{16}
\tightlist
\item
  The Fst (fixation index) of a particular autosomal gene between two
  human populations is 0.08. Interpret this value: how much of total
  human genetic variation lies between populations vs within
  populations? Compare this to the claim that ``race explains most human
  genetic variation.''
\end{enumerate}

\begin{enumerate}
\def\labelenumi{\arabic{enumi}.}
\setcounter{enumi}{17}
\tightlist
\item
  Huntington's disease is autosomal dominant. The causative allele (HTT
  with \textgreater36 CAG repeats) shows near-100\% penetrance. If an
  affected parent (heterozygous, Hh) has five children, what fraction
  are expected to be affected? Why does penetrance differ from simple
  Mendelian dominant inheritance in this context?
\end{enumerate}

\begin{enumerate}
\def\labelenumi{\arabic{enumi}.}
\setcounter{enumi}{18}
\tightlist
\item
  Genomic imprinting: if the paternal allele of IGF2 is expressed and
  the maternal allele is silenced, predict the phenotype of: (a) a child
  who inherits a deletion on the paternal chromosome 11
  (Beckwith-Wiedemann: overgrowth); (b) a child inheriting a deletion on
  the maternal chromosome 11 (Silver-Russell: growth restriction). Use
  the concept of parent-of-origin-specific expression to explain both.
\end{enumerate}

\begin{enumerate}
\def\labelenumi{\arabic{enumi}.}
\setcounter{enumi}{19}
\tightlist
\item
  Inversion heterozygosity in \emph{Drosophila} suppresses recombination
  in the inverted region because crossing over within an inversion
  produces unbalanced chromosomes (duplications and deletions), which
  are lethal. Explain the mechanism and why evolution often preserves
  inversions if they capture coadapted gene complexes.
\end{enumerate}

\subsection*{Questions 21--30: Synthesis and
Evaluation}\label{questions-2130-synthesis-and-evaluation-16}
\addcontentsline{toc}{subsection}{Questions 21--30: Synthesis and
Evaluation}

\begin{enumerate}
\def\labelenumi{\arabic{enumi}.}
\setcounter{enumi}{20}
\tightlist
\item
  Evaluate the role of the synaptonemal complex (SC) in meiotic
  recombination: how does the SC align homologous chromosomes, what
  components does it consist of, and how does SC disassembly correlate
  with chiasma formation? What happens when SC formation fails?
\end{enumerate}

\begin{enumerate}
\def\labelenumi{\arabic{enumi}.}
\setcounter{enumi}{21}
\tightlist
\item
  Comparative genomic hybridisation (CGH) and SNP arrays detect copy
  number variation (CNVs) across the genome. Evaluate the clinical
  utility of chromosomal microarray (CMA) analysis in: (a) prenatal
  diagnosis (replacing karyotyping?); (b) diagnosing children with
  unexplained intellectual disability; (c) identifying germline vs
  somatic CNVs in cancer. What is the main limitation of CMA vs
  long-read sequencing for CNV detection?
\end{enumerate}

\begin{enumerate}
\def\labelenumi{\arabic{enumi}.}
\setcounter{enumi}{22}
\tightlist
\item
  Design an experiment to map the XIC (X-inactivation centre) on the
  human X chromosome using: (a) X-autosome translocation cell lines; (b)
  CRISPR deletion of candidate regions; (c) Xist RNA FISH in human iPSC
  differentiation. What molecular result would confirm the XIC has been
  correctly identified?
\end{enumerate}

\begin{enumerate}
\def\labelenumi{\arabic{enumi}.}
\setcounter{enumi}{23}
\tightlist
\item
  Evaluate the genomic consequences of chromosomal rearrangements that
  bring proto-oncogenes near active enhancers (``enhancer hijacking'').
  Give two specific examples (BCL2 and IGH in follicular lymphoma; MYC
  and IGH in Burkitt lymphoma) and explain why the same translocation in
  different individuals produces the same cancer phenotype.
\end{enumerate}

\begin{enumerate}
\def\labelenumi{\arabic{enumi}.}
\setcounter{enumi}{24}
\tightlist
\item
  The genetic and epigenetic mechanisms of centromere identity are
  distinct from those of euchromatin. Evaluate: (a) why centromeres
  require CENP-A (not canonical H3) for proper kinetochore assembly; (b)
  why centromere sequences are highly repetitive (\texorpdfstring{\ensuremath{\alpha}}{alpha}-satellite) despite
  being functionally defined by CENP-A deposition; (c) what would happen
  to chromosome segregation if CENP-A deposition were globally
  abolished.
\end{enumerate}

\begin{enumerate}
\def\labelenumi{\arabic{enumi}.}
\setcounter{enumi}{25}
\tightlist
\item
  Position effects in Drosophila (position effect variegation, PEV)
  occur when a gene is relocated near heterochromatin via chromosomal
  inversion, causing mosaic silencing. Evaluate the molecular basis of
  heterochromatic silencing spreading (HP1-H3K9me3 loop), what
  suppresses PEV (loss of heterochromatin factors), and how this relates
  to genes at the chromosome 16p11.2 deletion locus in autism.
\end{enumerate}

\begin{enumerate}
\def\labelenumi{\arabic{enumi}.}
\setcounter{enumi}{26}
\tightlist
\item
  Evaluate the selective sweep as an evolutionary force producing
  chromosomal linkage disequilibrium (LD). When a beneficial mutation
  sweeps to fixation, flanking variants are dragged along (``genetic
  hitchhiking''), producing extended LD. How is extended LD detected
  (iHS statistic), and what gene regions in the human genome show
  evidence of recent selective sweeps?
\end{enumerate}

\begin{enumerate}
\def\labelenumi{\arabic{enumi}.}
\setcounter{enumi}{27}
\tightlist
\item
  Chromosome conformation capture (3C/4C/HiC) reveals that chromosomes
  occupy specific territories in the nucleus and that A (active) and B
  (inactive) compartments are spatially separated. Evaluate how
  compartment identity is maintained through cell division (given that
  nuclear organisation must be re-established after mitosis), and what
  proteins mark the A vs B compartment boundaries.
\end{enumerate}

\begin{enumerate}
\def\labelenumi{\arabic{enumi}.}
\setcounter{enumi}{28}
\tightlist
\item
  Meiotic drive (non-Mendelian segregation) is a selfish genetic element
  that biases its own transmission into gametes at \textgreater50\%
  frequency. Evaluate the molecular mechanisms of drive (t-allele in
  mice, SD system in \emph{Drosophila}) and the population genetic
  consequences: under what conditions does a driving allele reach
  fixation vs form a stable polymorphism? What restricts drive from
  eliminating the drive-suppressor population?
\end{enumerate}

\begin{enumerate}
\def\labelenumi{\arabic{enumi}.}
\setcounter{enumi}{29}
\tightlist
\item
  Critically evaluate the concept of ``chromosomal instability'' (CIN)
  in cancer, where cells gain and lose whole chromosomes or large
  chromosome arms with each division. Is CIN a cause or consequence of
  cancer? Evaluate the evidence that CIN generates the aneuploidy needed
  for tumour evolution vs the evidence that CIN is incompatible with
  cell viability and is therefore tolerated primarily within a narrow
  window of ``just right'' instability.
\end{enumerate}

\newpage

\section*{Questions --- Ch 18: Population Genetics and Hardy-Weinberg
Equilibrium}\label{questions-ch-18-population-genetics-and-hardy-weinberg-equilibrium}
\addcontentsline{toc}{section}{Questions --- Ch 18: Population Genetics
and Hardy-Weinberg Equilibrium}

\label{sec:q_unit_V_population_genetics}

\subsection*{Instructor Use and Coverage
Notes}\label{instructor-use-and-coverage-notes-20}
\addcontentsline{toc}{subsection}{Instructor Use and Coverage Notes}

\begin{itemize}
\tightlist
\item
  \textbf{Coverage target:} Calculate expected frequencies and identify
  which evolutionary force could explain deviations.
\item
  \textbf{Model/data emphasis:} Hardy-Weinberg and allele-frequency
  recurrence calculations.
\item
  \textbf{Assessment alignment:} Statistical Tests and Data Analysis,
  Representing and Describing Data.
\item
  \textbf{Misconception probe:} Hardy-Weinberg is a null model, not a
  claim that populations do not evolve.
\item
  \textbf{Transfer product:} Apply population-genetic reasoning to
  screening, conservation, and pathogen evolution.
\item
  \textbf{Grading focus:} award full credit for mechanism, evidence,
  boundary conditions, and units when a calculation is required.
\item
  \textbf{Suggested use:} draw one recall item, one application item,
  and one synthesis item when building a short quiz from this bank.
\end{itemize}

\subsection*{Questions 1--10: Recall and
Comprehension}\label{questions-110-recall-and-comprehension-17}
\addcontentsline{toc}{subsection}{Questions 1--10: Recall and
Comprehension}

\emph{This activity accompanies \cref{sec:unit_V_population_genetics} of
the textbook --- review that chapter before attempting the exercises
below.}

\begin{enumerate}
\def\labelenumi{\arabic{enumi}.}
\tightlist
\item
  State the Hardy-Weinberg principle and list the five conditions
  required for equilibrium.
\end{enumerate}

\begin{enumerate}
\def\labelenumi{\arabic{enumi}.}
\setcounter{enumi}{1}
\tightlist
\item
  Write the Hardy-Weinberg equations for a two-allele system. What do p,
  q, p\texorpdfstring{\textsuperscript{2}}{2}, 2pq, and q\texorpdfstring{\textsuperscript{2}}{2} represent?
\end{enumerate}

\begin{enumerate}
\def\labelenumi{\arabic{enumi}.}
\setcounter{enumi}{2}
\tightlist
\item
  What is allele frequency? What is genotype frequency? How do they
  relate?
\end{enumerate}

\begin{enumerate}
\def\labelenumi{\arabic{enumi}.}
\setcounter{enumi}{3}
\tightlist
\item
  Define genetic drift. Why is it more pronounced in small populations?
\end{enumerate}

\begin{enumerate}
\def\labelenumi{\arabic{enumi}.}
\setcounter{enumi}{4}
\tightlist
\item
  What is the founder effect? Give one human example.
\end{enumerate}

\begin{enumerate}
\def\labelenumi{\arabic{enumi}.}
\setcounter{enumi}{5}
\tightlist
\item
  Define gene flow (migration). How does it affect allele frequencies in
  two populations?
\end{enumerate}

\begin{enumerate}
\def\labelenumi{\arabic{enumi}.}
\setcounter{enumi}{6}
\tightlist
\item
  What is mutation rate? What is the typical per-base-pair mutation rate
  in humans per generation?
\end{enumerate}

\begin{enumerate}
\def\labelenumi{\arabic{enumi}.}
\setcounter{enumi}{7}
\tightlist
\item
  Define natural selection. What is the selection coefficient (s)?
\end{enumerate}

\begin{enumerate}
\def\labelenumi{\arabic{enumi}.}
\setcounter{enumi}{8}
\tightlist
\item
  What is heterozygote advantage (overdominance)? Give the sickle-cell
  anaemia example.
\end{enumerate}

\begin{enumerate}
\def\labelenumi{\arabic{enumi}.}
\setcounter{enumi}{9}
\tightlist
\item
  What is the neutral theory of molecular evolution? Who proposed it?
\end{enumerate}

\subsection*{Questions 11--20: Application and
Analysis}\label{questions-1120-application-and-analysis-17}
\addcontentsline{toc}{subsection}{Questions 11--20: Application and
Analysis}

\begin{enumerate}
\def\labelenumi{\arabic{enumi}.}
\setcounter{enumi}{10}
\tightlist
\item
  A population has 500 individuals: 200 AA, 240 Aa, 60 aa. Calculate
  allele frequencies p and q. Calculate expected HWE genotype
  frequencies. Use chi-square to test if the population is in HWE.
\end{enumerate}

\begin{enumerate}
\def\labelenumi{\arabic{enumi}.}
\setcounter{enumi}{11}
\tightlist
\item
  After a natural disaster reduces a population from 10,000 to 100
  individuals, what are the genetic consequences? Calculate how much
  genetic variation is expected to be lost (as measured by
  heterozygosity) in one generation: H\_t = H\texorpdfstring{\textsubscript{0}}{0}(1 \texorpdfstring{\ensuremath{-}}{-} 1/(2N\_e)).
\end{enumerate}

\begin{enumerate}
\def\labelenumi{\arabic{enumi}.}
\setcounter{enumi}{12}
\tightlist
\item
  In a population of N\_e = 50 individuals, an allele at frequency q =
  0.2 is selectively neutral. Calculate the probability it is lost by
  drift in the first generation (approximate using 1 \texorpdfstring{\ensuremath{-}}{-} q = 0.8). Over
  many generations, what is the probability it eventually reaches
  fixation (= q\texorpdfstring{\textsubscript{0}}{0})?
\end{enumerate}

\begin{enumerate}
\def\labelenumi{\arabic{enumi}.}
\setcounter{enumi}{13}
\tightlist
\item
  Selection against a recessive allele (s = 1, lethal; aa dies before
  reproduction): starting from q\texorpdfstring{\textsubscript{0}}{0} = 0.5, calculate the allele frequency
  after 1, 5, and 10 generations using q\_t = q\texorpdfstring{\textsubscript{0}}{0}/(1 + t\texorpdfstring{\ensuremath{\cdot}}{.}s\texorpdfstring{\ensuremath{\cdot}}{.}q\texorpdfstring{\textsubscript{0}}{0}). How many
  generations to reduce q from 0.5 to 0.01?
\end{enumerate}

\begin{enumerate}
\def\labelenumi{\arabic{enumi}.}
\setcounter{enumi}{14}
\tightlist
\item
  Cystic fibrosis has a frequency of 1/2,500 in European populations.
  Calculate: (a) q (frequency of CF allele); (b) the carrier frequency
  (2pq); (c) what Hardy-Weinberg assumption must hold for this
  calculation to be valid; (d) why the CF allele is maintained at
  relatively high frequency (propose two hypotheses).
\end{enumerate}

\begin{enumerate}
\def\labelenumi{\arabic{enumi}.}
\setcounter{enumi}{15}
\tightlist
\item
  Fst between two populations is 0.15 for an autosomal locus. Interpret
  this: what fraction of total variation is between vs within
  populations? Calculate the estimated gene flow (Nm) using the
  approximation Fst \texorpdfstring{\ensuremath{\approx}}{~} 1/(1 + 4Nm).
\end{enumerate}

\begin{enumerate}
\def\labelenumi{\arabic{enumi}.}
\setcounter{enumi}{16}
\tightlist
\item
  A bottleneck reduces a population to 10 individuals for one
  generation, then expands. If initial heterozygosity was 0.6, calculate
  expected heterozygosity after the bottleneck (H = H\texorpdfstring{\textsubscript{0}}{0}(1 \texorpdfstring{\ensuremath{-}}{-}
  1/(2N\_bottleneck))). After 10 generations of recovery with large N,
  how does heterozygosity change?
\end{enumerate}

\begin{enumerate}
\def\labelenumi{\arabic{enumi}.}
\setcounter{enumi}{17}
\tightlist
\item
  The neutral theory predicts that the rate of neutral molecular
  evolution equals the mutation rate (\texorpdfstring{\ensuremath{\mu}}{mu}), independent of population
  size. Explain the logic: if advantageous mutations are rare and
  deleterious mutations are removed by selection, then substitutions in
  the fossil record should primarily reflect neutral changes. What types
  of sequences evolve fastest (synonymous sites) and slowest (amino
  acid-changing in essential proteins)?
\end{enumerate}

\begin{enumerate}
\def\labelenumi{\arabic{enumi}.}
\setcounter{enumi}{18}
\tightlist
\item
  Balancing selection maintains two or more alleles at a locus above
  neutral expectations. Name three types of balancing selection and one
  example of each: (a) heterozygote advantage; (b) frequency-dependent
  selection; (c) fluctuating selection (spatially or temporally
  varying). Which type is operative at MHC loci?
\end{enumerate}

\begin{enumerate}
\def\labelenumi{\arabic{enumi}.}
\setcounter{enumi}{19}
\tightlist
\item
  Sexual selection is a form of natural selection acting on mating
  success rather than survival. Distinguish intrasexual selection
  (competition between same-sex individuals) from intersexual selection
  (mate choice). Give one example of each in animals. How might sexual
  selection drive the rapid evolution of secondary sexual
  characteristics even if they reduce survival (Zahavian handicap
  model)?
\end{enumerate}

\subsection*{Questions 21--30: Synthesis and
Evaluation}\label{questions-2130-synthesis-and-evaluation-17}
\addcontentsline{toc}{subsection}{Questions 21--30: Synthesis and
Evaluation}

\begin{enumerate}
\def\labelenumi{\arabic{enumi}.}
\setcounter{enumi}{20}
\tightlist
\item
  Evaluate the molecular clock hypothesis: if the neutral substitution
  rate equals \texorpdfstring{\ensuremath{\mu}}{mu} (the mutation rate), then the number of differences
  between sequences accumulates linearly with time. What calibrates the
  molecular clock (fossil dates), how has the clock been used to date
  the human-chimpanzee split, and what are the sources of rate variation
  (generation time, effective population size, repair efficiency) that
  make the clock ``sloppy''?
\end{enumerate}

\begin{enumerate}
\def\labelenumi{\arabic{enumi}.}
\setcounter{enumi}{21}
\tightlist
\item
  The 1000 Genomes Project revealed that most human genetic variants are
  rare (minor allele frequency \textless{} 1\%). Evaluate what the site
  frequency spectrum (SFS) of human variants tells us about recent
  population expansion (bottleneck \texorpdfstring{\ensuremath{\to}}{->} colonisation of new territories \texorpdfstring{\ensuremath{\to}}{->}
  rapid expansion), and how deviations from neutral SFS predictions
  (Tajima's D) detect selection.
\end{enumerate}

\begin{enumerate}
\def\labelenumi{\arabic{enumi}.}
\setcounter{enumi}{22}
\tightlist
\item
  Evaluate the evidence that human populations show significantly
  different allele frequencies at loci controlling pigmentation (MC1R,
  SLC24A5, HERC2/OCA2) and that these differences reflect recent local
  adaptation (UV radiation, vitamin D synthesis) rather than random
  drift. What evolutionary mechanism drove rapid adaptation at these
  loci in the past 10,000--20,000 years?
\end{enumerate}

\begin{enumerate}
\def\labelenumi{\arabic{enumi}.}
\setcounter{enumi}{23}
\tightlist
\item
  Design an analysis to identify signatures of selection in a
  genome-wide SNP dataset from 1,000 individuals in two populations.
  Describe the following statistical tests you would use: Fst outlier
  analysis, XP-EHH (cross-population extended haplotype homozygosity),
  and Tajima's D, stating what each detects. What results would indicate
  recent selective sweeps, and how would you validate a candidate gene
  using functional data?
\end{enumerate}

\begin{enumerate}
\def\labelenumi{\arabic{enumi}.}
\setcounter{enumi}{24}
\tightlist
\item
  Evaluate ``negative frequency-dependent selection'' at the ABO blood
  group locus. Propose why O (producing no antigen) remains at
  \textgreater60\% frequency in many human populations despite the
  disadvantage of lacking most antigens (O individuals cannot receive A,
  B, or AB blood). What pathogens preferentially target specific blood
  group antigens, and how does this drive balancing selection?
\end{enumerate}

\begin{enumerate}
\def\labelenumi{\arabic{enumi}.}
\setcounter{enumi}{25}
\tightlist
\item
  Effective population size (N\_e) in humans is estimated at
  about 10,000--20,000 (from genetic diversity), far smaller
  than the current census population (about 8 billion).
  Evaluate three reasons why N\_e \textless\textless{} N\_c: (a) highly
  unequal reproductive success; (b) past bottlenecks; (c) sex ratio
  inequality. How does a small historical N\_e affect our genome (e.g.,
  low heterozygosity, linkage disequilibrium)?
\end{enumerate}

\begin{enumerate}
\def\labelenumi{\arabic{enumi}.}
\setcounter{enumi}{26}
\tightlist
\item
  Genetic load is the reduction in population fitness due to deleterious
  alleles. Evaluate the ``mutation load paradox'' in humans: each human
  carries about 70--100 de novo mutations per generation, yet
  the population continues to grow. How is deleterious load purged by
  selection (purifying selection), and what analytical framework
  (fitness decline per mutation, calculated from GWAS) resolves the
  paradox?
\end{enumerate}

\begin{enumerate}
\def\labelenumi{\arabic{enumi}.}
\setcounter{enumi}{27}
\tightlist
\item
  Evaluate the role of sexual reproduction vs asexual reproduction in
  maintaining genetic variation under selection. The Red Queen
  hypothesis proposes that sexual reproduction is advantageous because
  it generates novel combinations to keep pace with rapidly evolving
  parasites. Evaluate: (a) the ecological evidence (host-parasite
  coevolution in snails, \emph{Potamopyrgus antipodarum}); (b) the
  genetic evidence (rate of Muller's ratchet in asexual vs sexual
  populations); (c) the cost of sex (only about 50\% of offspring bear
  offspring themselves).
\end{enumerate}

\begin{enumerate}
\def\labelenumi{\arabic{enumi}.}
\setcounter{enumi}{28}
\tightlist
\item
  Population stratification in GWAS occurs when ancestry-related allele
  frequency differences confound disease-association signals. Evaluate
  how principal component analysis (PCA) corrects for stratification in
  GWAS, why failing to correct can generate spurious associations, and
  why the same caution applies to forensic or outbreak
  population-inference claims.
\end{enumerate}

\begin{enumerate}
\def\labelenumi{\arabic{enumi}.}
\setcounter{enumi}{29}
\tightlist
\item
  Critically evaluate the ``mutation-selection-drift'' balance framework
  as a complete explanation for the maintenance of genetic variation in
  natural populations. Where does this framework work well (e.g., for
  clearly deleterious alleles), and where does it break down (complex
  traits with many loci of small effect, phenotypic plasticity,
  developmental canalization, environmental gene \texorpdfstring{\ensuremath{\times}}{x} environment
  interactions)? Propose how the modern extended synthesis incorporates
  concepts beyond classical population genetics.
\end{enumerate}

\newpage

\section*{Questions --- Ch 19: Evolution --- Theory, Natural Selection,
and
Adaptation}\label{questions-ch-19-evolution-theory-natural-selection-and-adaptation}
\addcontentsline{toc}{section}{Questions --- Ch 19: Evolution ---
Theory, Natural Selection, and Adaptation}

\label{sec:q_unit_VI_evolution_and_selection}

\subsection*{Instructor Use and Coverage
Notes}\label{instructor-use-and-coverage-notes-21}
\addcontentsline{toc}{subsection}{Instructor Use and Coverage Notes}

\begin{itemize}
\tightlist
\item
  \textbf{Coverage target:} Estimate selection strength from data and
  explain why the result is context dependent.
\item
  \textbf{Model/data emphasis:} Selection coefficient and
  allele-frequency trajectory calculations.
\item
  \textbf{Assessment alignment:} Visual Representations, Statistical
  Tests and Data Analysis, Argumentation.
\item
  \textbf{Misconception probe:} Evolution is not goal-directed progress;
  it is local change in populations under constraints.
\item
  \textbf{Transfer product:} Transfer selection reasoning to
  antibiotics, pesticide resistance, cancer, or climate adaptation.
\item
  \textbf{Grading focus:} award full credit for mechanism, evidence,
  boundary conditions, and units when a calculation is required.
\item
  \textbf{Suggested use:} draw one recall item, one application item,
  and one synthesis item when building a short quiz from this bank.
\end{itemize}

\subsection*{Questions 1--10: Recall and
Comprehension}\label{questions-110-recall-and-comprehension-18}
\addcontentsline{toc}{subsection}{Questions 1--10: Recall and
Comprehension}

\emph{This activity accompanies
\cref{sec:unit_VI_evolution_and_selection} of the textbook --- review
that chapter before attempting the exercises below.}

\begin{enumerate}
\def\labelenumi{\arabic{enumi}.}
\tightlist
\item
  State Darwin's mechanism of natural selection in four steps.
\end{enumerate}

\begin{enumerate}
\def\labelenumi{\arabic{enumi}.}
\setcounter{enumi}{1}
\tightlist
\item
  What is fitness in evolutionary biology? Does it relate to physical
  strength?
\end{enumerate}

\begin{enumerate}
\def\labelenumi{\arabic{enumi}.}
\setcounter{enumi}{2}
\tightlist
\item
  Define adaptation. Give one structural, one physiological, and one
  behavioural example.
\end{enumerate}

\begin{enumerate}
\def\labelenumi{\arabic{enumi}.}
\setcounter{enumi}{3}
\tightlist
\item
  What is directional selection? What is stabilising selection? Give one
  example of each.
\end{enumerate}

\begin{enumerate}
\def\labelenumi{\arabic{enumi}.}
\setcounter{enumi}{4}
\tightlist
\item
  What is the difference between microevolution and macroevolution?
\end{enumerate}

\begin{enumerate}
\def\labelenumi{\arabic{enumi}.}
\setcounter{enumi}{5}
\tightlist
\item
  Name four lines of evidence for evolution by natural selection.
\end{enumerate}

\begin{enumerate}
\def\labelenumi{\arabic{enumi}.}
\setcounter{enumi}{6}
\tightlist
\item
  What is the fossil record? What is a transitional fossil? Give one
  example.
\end{enumerate}

\begin{enumerate}
\def\labelenumi{\arabic{enumi}.}
\setcounter{enumi}{7}
\tightlist
\item
  What is homology? Distinguish structural homology from molecular
  homology.
\end{enumerate}

\begin{enumerate}
\def\labelenumi{\arabic{enumi}.}
\setcounter{enumi}{8}
\tightlist
\item
  What are vestigial structures? Give two examples in humans.
\end{enumerate}

\begin{enumerate}
\def\labelenumi{\arabic{enumi}.}
\setcounter{enumi}{9}
\tightlist
\item
  What is convergent evolution? Give one example.
\end{enumerate}

\subsection*{Questions 11--20: Application and
Analysis}\label{questions-1120-application-and-analysis-18}
\addcontentsline{toc}{subsection}{Questions 11--20: Application and
Analysis}

\begin{enumerate}
\def\labelenumi{\arabic{enumi}.}
\setcounter{enumi}{10}
\tightlist
\item
  Rock pocket mice in New Mexico exist on both dark lava and light soil
  habitats. Mice with dark coats survive better on lava; light mice
  survive better on soil. Explain: (a) the type of selection
  (disruptive); (b) the molecular basis of dark coat colour (MC1R
  mutations); (c) what the independent evolution of dark coats in
  different geographic populations tells us about convergent evolution.
\end{enumerate}

\begin{enumerate}
\def\labelenumi{\arabic{enumi}.}
\setcounter{enumi}{11}
\tightlist
\item
  The industrial melanism of the peppered moth (\emph{Biston betularia})
  is a textbook example of directional selection. Calculate: if initial
  light morph frequency = 0.99 and dark morph = 0.01, and selection
  coefficient against light morph = 0.33 (on industrial trees), after 50
  years (50 generations), predict the approximate frequency of the dark
  morph using the directional selection equation.
\end{enumerate}

\begin{enumerate}
\def\labelenumi{\arabic{enumi}.}
\setcounter{enumi}{12}
\tightlist
\item
  Sexual selection can drive traits to fixation even when they reduce
  survival (peacock tail). Use the Zahavian handicap principle to
  explain: (a) why primarily males with good ``condition'' can afford to
  develop large ornaments; (b) how female preference for ornaments
  evolves (Fisherian runaway vs good genes models); (c) evidence
  distinguishing the two models.
\end{enumerate}

\begin{enumerate}
\def\labelenumi{\arabic{enumi}.}
\setcounter{enumi}{13}
\tightlist
\item
  Kin selection allows altruistic behaviour to evolve when altruism
  benefits genetic relatives. Apply Hamilton's rule: altruism evolves
  when rb \textgreater{} c (r = relatedness, b = benefit, c = cost). A
  worker honey bee (r = 0.75 to sisters) helps raise 10 sisters (b = 10)
  at the cost of 0 own offspring (due to sterility). Is altruism
  favoured? What would rB need to be to just exceed c = 3.0?
\end{enumerate}

\begin{enumerate}
\def\labelenumi{\arabic{enumi}.}
\setcounter{enumi}{14}
\tightlist
\item
  Selection on quantitative traits can be modelled by the breeder's
  equation: R = h\texorpdfstring{\textsuperscript{2}}{2}S (R = response to selection, h\texorpdfstring{\textsuperscript{2}}{2} = heritability, S =
  selection differential). If height heritability is 0.80 and we select
  parents who are 5 cm taller than the population mean, predict R. How
  many generations of selection would be needed to shift mean height by
  10 cm?
\end{enumerate}

\begin{enumerate}
\def\labelenumi{\arabic{enumi}.}
\setcounter{enumi}{15}
\tightlist
\item
  Gould and Lewontin's ``spandrel'' argument suggests that not most
  traits are adaptations --- some are architectural by-products of other
  adaptations (spandrels in the context of evolutionary biology). Give
  two specific biological examples of putative spandrels and evaluate
  the methodological criteria needed to distinguish a spandrel from an
  adaptation.
\end{enumerate}

\begin{enumerate}
\def\labelenumi{\arabic{enumi}.}
\setcounter{enumi}{16}
\tightlist
\item
  Antibiotic resistance evolution offers real-time evidence of natural
  selection. Describe the Lenski long-term evolution experiment with
  \emph{E. coli} (\textgreater75,000 generations). What phenotypic
  innovations have evolved, and how does the experiment distinguish
  ``chance'' events from adaptive improvement?
\end{enumerate}

\begin{enumerate}
\def\labelenumi{\arabic{enumi}.}
\setcounter{enumi}{17}
\tightlist
\item
  Coevolution occurs when two species reciprocally affect each other's
  evolution. Describe the Red Queen model of host-parasite coevolution:
  (a) why does the host need to evolve continuously; (b) what does the
  snail (\emph{Potamopyrgus antipodarum}) -- trematode parasite system
  teach us; (c) what arms race dynamics emerge between plant alkaloids
  and insect detoxification enzymes?
\end{enumerate}

\begin{enumerate}
\def\labelenumi{\arabic{enumi}.}
\setcounter{enumi}{18}
\tightlist
\item
  Punctuated equilibrium (Gould and Eldredge, 1972) argues that most
  evolutionary change is concentrated in rapid bursts associated with
  speciation, not in gradual anagenesis. Evaluate the fossil evidence
  for and against punctuated equilibrium using the Cerion land snail
  data from the Bahamas and the evolution of horse lineages.
\end{enumerate}

\begin{enumerate}
\def\labelenumi{\arabic{enumi}.}
\setcounter{enumi}{19}
\tightlist
\item
  Biogeography provides evidence for evolution: organisms on islands
  often resemble those on the nearest mainland but have evolved unique
  adaptations. Apply the concept of adaptive radiation to Hawaiian
  honeycreepers (about 50 species from one ancestral finch-like
  coloniser): which resources were exploited by different lineages, and
  what ecological opportunity allowed such rapid diversification?
\end{enumerate}

\subsection*{Questions 21--30: Synthesis and
Evaluation}\label{questions-2130-synthesis-and-evaluation-18}
\addcontentsline{toc}{subsection}{Questions 21--30: Synthesis and
Evaluation}

\begin{enumerate}
\def\labelenumi{\arabic{enumi}.}
\setcounter{enumi}{20}
\tightlist
\item
  Evaluate the extended evolutionary synthesis (EES) relative to the
  modern synthesis. Which additional mechanisms does EES incorporate
  (developmental plasticity, epigenetic inheritance, niche construction,
  evolvability)? Choose one mechanism and evaluate the evidence that it
  plays a role comparable in importance to natural selection in at least
  one biological system.
\end{enumerate}

\begin{enumerate}
\def\labelenumi{\arabic{enumi}.}
\setcounter{enumi}{21}
\tightlist
\item
  Design an experiment to test the ``good genes'' hypothesis for female
  preference in peacocks (\emph{Pavo cristatus}). Describe: the relevant
  measurements (eyespot count, parasite load, extra-pair offspring
  survival), statistical analyses needed, and a falsification criterion.
  What result would confirm vs reject the hypothesis?
\end{enumerate}

\begin{enumerate}
\def\labelenumi{\arabic{enumi}.}
\setcounter{enumi}{22}
\tightlist
\item
  Evaluate the evolutionary arms race between cuckoos and their hosts
  (\emph{Cuculus canorus} vs reed warblers). Describe: (a) forgery (egg
  mimicry); (b) host counter-adaptations (pattern recognition, egg
  ejection); (c) what would have to happen for the host to ``win'' and
  cuckoos to go extinct. What does the coevolutionary balance tell us
  about the role of Red Queen dynamics in parasite-host coevolution?
\end{enumerate}

\begin{enumerate}
\def\labelenumi{\arabic{enumi}.}
\setcounter{enumi}{23}
\tightlist
\item
  Molecular convergence (independent evolution of the same molecular
  change in different lineages) provides the strongest evidence for
  adaptation at the sequence level. Evaluate three examples: (a) same
  amino acid substitution in haemoglobin for high-altitude adaptation in
  birds; (b) convergent microbial drug resistance mutations; (c)
  convergent VLDLR mutations for flight in birds. What does molecular
  convergence tell us about the size of genotypic variation space
  accessible by natural selection?
\end{enumerate}

\begin{enumerate}
\def\labelenumi{\arabic{enumi}.}
\setcounter{enumi}{24}
\tightlist
\item
  The ``fitness landscape'' model represents most possible genotypes as
  a surface where height corresponds to fitness. Evaluate the
  limitations of this model for understanding: (a) epistasis
  (non-additive interactions between loci); (b) adaptive valleys (can
  evolution cross fitness valleys?); (c) frequency-dependent fitness
  (landscape changes with allele frequency). Propose an alternative
  model that addresses these limitations.
\end{enumerate}

\begin{enumerate}
\def\labelenumi{\arabic{enumi}.}
\setcounter{enumi}{25}
\tightlist
\item
  Group selection (selection acting on groups rather than individuals)
  was largely dismissed after Williams (1966), but has been revived as
  multilevel selection (MLS) theory. Evaluate: (a) the conditions under
  which group selection overrides individual selection; (b) the evidence
  from slime mould (\emph{Dictyostelium}) cooperation; (c) whether human
  cooperation can be explained by MLS or requires cultural group
  selection.
\end{enumerate}

\begin{enumerate}
\def\labelenumi{\arabic{enumi}.}
\setcounter{enumi}{26}
\tightlist
\item
  Evaluate the claim that evolutionary theory is a scientific fact while
  specific evolutionary narratives (the origin of feathers, the
  evolution of eukaryotes) are provisional hypotheses. What does
  ``theory'' mean in science compared to colloquial usage? How does the
  misunderstanding of ``theory'' fuel creationist objections, and what
  is the appropriate scientific response?
\end{enumerate}

\begin{enumerate}
\def\labelenumi{\arabic{enumi}.}
\setcounter{enumi}{27}
\tightlist
\item
  Evolvability --- the capacity to evolve --- may itself be a product of
  natural selection. Evaluate: (a) the evidence that the modular
  architecture of proteins (protein domains) evolved to enhance
  evolvability; (b) whether mutation rate itself can be subject to
  natural selection; (c) the concept of robustness (canalization) and
  its relationship to evolvability.
\end{enumerate}

\begin{enumerate}
\def\labelenumi{\arabic{enumi}.}
\setcounter{enumi}{28}
\tightlist
\item
  The ``major evolutionary transitions'' (Szathmáry and Maynard Smith,
  1995) describe eight key events in the history of life (e.g., origin
  of the eukaryotic cell, origin of multicellularity). Evaluate what
  distinguishes each transition from ordinary speciation, and propose
  what the next major evolutionary transition might be (e.g., a digital
  information-based evolving system?).
\end{enumerate}

\begin{enumerate}
\def\labelenumi{\arabic{enumi}.}
\setcounter{enumi}{29}
\tightlist
\item
  Critically evaluate whether evolutionary theory can be directly
  applied to artificial systems (machine learning algorithms, memes,
  cultural evolution). What structural features must any system have to
  undergo Darwinian evolution (heritable variation, differential
  reproduction, competition)? Does cultural evolution satisfy these
  criteria, and how does it differ from biological evolution in terms of
  inheritance mechanisms and selection units?
\end{enumerate}

\newpage

\section*{Questions --- Ch 20: Genetic Drift, Gene Flow, and
Speciation}\label{questions-ch-20-genetic-drift-gene-flow-and-speciation}
\addcontentsline{toc}{section}{Questions --- Ch 20: Genetic Drift, Gene
Flow, and Speciation}

\label{sec:q_unit_VI_genetic_drift_and_speciation}

\subsection*{Instructor Use and Coverage
Notes}\label{instructor-use-and-coverage-notes-22}
\addcontentsline{toc}{subsection}{Instructor Use and Coverage Notes}

\begin{itemize}
\tightlist
\item
  \textbf{Coverage target:} Predict how population size or isolation
  changes divergence and fixation risk.
\item
  \textbf{Model/data emphasis:} Fixation probability, effective
  population size, and migration-selection balance.
\item
  \textbf{Assessment alignment:} Visual Representations, Statistical
  Tests and Data Analysis, Argumentation.
\item
  \textbf{Misconception probe:} Random does not mean patternless;
  stochastic processes have predictable distributions.
\item
  \textbf{Transfer product:} Apply drift reasoning to endangered
  populations, founder effects, and island radiations.
\item
  \textbf{Grading focus:} award full credit for mechanism, evidence,
  boundary conditions, and units when a calculation is required.
\item
  \textbf{Suggested use:} draw one recall item, one application item,
  and one synthesis item when building a short quiz from this bank.
\end{itemize}

\subsection*{Questions 1--10: Recall and
Comprehension}\label{questions-110-recall-and-comprehension-19}
\addcontentsline{toc}{subsection}{Questions 1--10: Recall and
Comprehension}

\emph{This activity accompanies
\cref{sec:unit_VI_genetic_drift_and_speciation} of the textbook ---
review that chapter before attempting the exercises below.}

\begin{enumerate}
\def\labelenumi{\arabic{enumi}.}
\tightlist
\item
  Define genetic drift. How does it differ from natural selection?
\end{enumerate}

\begin{enumerate}
\def\labelenumi{\arabic{enumi}.}
\setcounter{enumi}{1}
\tightlist
\item
  What is the founder effect? What is the bottleneck effect? How do they
  differ?
\end{enumerate}

\begin{enumerate}
\def\labelenumi{\arabic{enumi}.}
\setcounter{enumi}{2}
\tightlist
\item
  Define gene flow. What effect does gene flow have on genetic
  differentiation between populations?
\end{enumerate}

\begin{enumerate}
\def\labelenumi{\arabic{enumi}.}
\setcounter{enumi}{3}
\tightlist
\item
  Define biological species concept (BSC). What is a reproductive
  isolation barrier?
\end{enumerate}

\begin{enumerate}
\def\labelenumi{\arabic{enumi}.}
\setcounter{enumi}{4}
\tightlist
\item
  Distinguish pre-zygotic from post-zygotic reproductive isolation
  barriers. Give two examples of each.
\end{enumerate}

\begin{enumerate}
\def\labelenumi{\arabic{enumi}.}
\setcounter{enumi}{5}
\tightlist
\item
  What is allopatric speciation? What is sympatric speciation? Give one
  natural example of each.
\end{enumerate}

\begin{enumerate}
\def\labelenumi{\arabic{enumi}.}
\setcounter{enumi}{6}
\tightlist
\item
  What is parapatric speciation? Give one example from plant
  populations.
\end{enumerate}

\begin{enumerate}
\def\labelenumi{\arabic{enumi}.}
\setcounter{enumi}{7}
\tightlist
\item
  Define peripatric speciation. How does it relate to the founder
  effect?
\end{enumerate}

\begin{enumerate}
\def\labelenumi{\arabic{enumi}.}
\setcounter{enumi}{8}
\tightlist
\item
  Define Fst (fixation index). What does Fst = 0 mean? What does Fst = 1
  mean?
\end{enumerate}

\begin{enumerate}
\def\labelenumi{\arabic{enumi}.}
\setcounter{enumi}{9}
\tightlist
\item
  What is adaptive radiation? Give one famous example.
\end{enumerate}

\subsection*{Questions 11--20: Application and
Analysis}\label{questions-1120-application-and-analysis-19}
\addcontentsline{toc}{subsection}{Questions 11--20: Application and
Analysis}

\begin{enumerate}
\def\labelenumi{\arabic{enumi}.}
\setcounter{enumi}{10}
\tightlist
\item
  A population of 1,000 individuals is reduced to 5 for three
  generations, then recovers to 1,000. Calculate the harmonic mean of
  the effective population size over the three-generation bottleneck,
  and use H\_t = H\texorpdfstring{\textsubscript{0}}{0}(1 \texorpdfstring{\ensuremath{-}}{-} 1/(2N\_e))\^{}t to estimate the heterozygosity
  remaining after recovery.
\end{enumerate}

\begin{enumerate}
\def\labelenumi{\arabic{enumi}.}
\setcounter{enumi}{11}
\tightlist
\item
  The Amish population (descended from 200 founders in the 1700s) has a
  1,000-fold higher frequency of Ellis-van Creveld syndrome (short-limb
  dwarfism) than the general population. Explain: (a) which mechanism is
  responsible; (b) whether natural selection for or against the
  condition is required to explain its frequency; (c) what this implies
  for genetic counselling in this community.
\end{enumerate}

\begin{enumerate}
\def\labelenumi{\arabic{enumi}.}
\setcounter{enumi}{12}
\tightlist
\item
  Two sister populations of fish separated 10,000 years ago are now
  parapatrically distributed along a river system. They differ in colour
  pattern and mating call. When brought together in the lab, they
  hybridise successfully but hybrids show reduced fertility. Classify:
  (a) what type of speciation occurred; (b) what type of reproductive
  isolation is demonstrated; (c) whether they are currently separate
  species under the BSC.
\end{enumerate}

\begin{enumerate}
\def\labelenumi{\arabic{enumi}.}
\setcounter{enumi}{13}
\tightlist
\item
  Calculate the estimated gene flow (Nm = number of migrants per
  generation) between two populations with Fst = 0.08, using Fst \texorpdfstring{\ensuremath{\approx}}{~} 1/(1
  + 4Nm). Is this a high or low level of gene flow? What migration rate
  (Nm \textgreater{} 1) is needed to prevent significant divergence?
\end{enumerate}

\begin{enumerate}
\def\labelenumi{\arabic{enumi}.}
\setcounter{enumi}{14}
\tightlist
\item
  In cichlid fish adaptive radiation in Lake Victoria, \textgreater{}
  500 species evolved in \textless{} 15,000 years --- an extremely rapid
  speciation event. Propose: (a) what ecological opportunity promoted
  rapid radiation; (b) what role mate choice based on colour pattern
  played in sympatric speciation; (c) how light transmission in
  different lake depths may have driven divergence.
\end{enumerate}

\begin{enumerate}
\def\labelenumi{\arabic{enumi}.}
\setcounter{enumi}{15}
\tightlist
\item
  Ring species provide a spatial illustration of the speciation
  continuum. The Larus gull chain circles the Arctic: populations A
  interbreed with B, B with C, C with D, D with E, but A and E do not
  interbreed. Explain: (a) why A and E are or are not separate species
  under BSC; (b) what this shows about speciation as a gradual,
  continuous process; (c) how climate change and habitat connectivity
  changes could affect this ring.
\end{enumerate}

\begin{enumerate}
\def\labelenumi{\arabic{enumi}.}
\setcounter{enumi}{16}
\tightlist
\item
  Hybridisation between species can produce evolutionary novelty.
  Allopolyploidy (hybrid + chromosome doubling) has produced many crop
  species: bread wheat (AABBDD, from three diploid ancestors). Explain:
  (a) how polyploidy immediately produces reproductive isolation from
  both ancestors; (b) why polyploid plants are often phenotypically
  superior (hybrid vigour/heterosis); (c) what the estimated
  contribution of allopolyploidy to flowering plant diversity is.
\end{enumerate}

\begin{enumerate}
\def\labelenumi{\arabic{enumi}.}
\setcounter{enumi}{17}
\tightlist
\item
  ``Magic traits'' are traits that simultaneously cause both ecological
  divergence (adaptation to different niches) and assortative mating
  (non-random mating based on the same trait). Give one example of a
  magic trait and explain why this simplifies the theory of sympatric
  speciation.
\end{enumerate}

\begin{enumerate}
\def\labelenumi{\arabic{enumi}.}
\setcounter{enumi}{18}
\tightlist
\item
  The biological species concept has been criticised as inapplicable to:
  (a) asexual organisms; (b) freely hybridising plant species. Evaluate
  two alternative species concepts (phylogenetic species concept,
  ecological species concept) and explain which is most useful for
  biodiversity conservation.
\end{enumerate}

\begin{enumerate}
\def\labelenumi{\arabic{enumi}.}
\setcounter{enumi}{19}
\tightlist
\item
  Hybrid zones are regions where two distinct populations interbreed.
  Using the tension zone model: (a) explain why a hybrid zone maintains
  a stable width; (b) what happens to hybrid fitness at the centre of
  the zone; (c) how hybrid zone data can estimate the selection against
  hybrids and the rate of gene flow.
\end{enumerate}

\subsection*{Questions 21--30: Synthesis and
Evaluation}\label{questions-2130-synthesis-and-evaluation-19}
\addcontentsline{toc}{subsection}{Questions 21--30: Synthesis and
Evaluation}

\begin{enumerate}
\def\labelenumi{\arabic{enumi}.}
\setcounter{enumi}{20}
\tightlist
\item
  Evaluate the evidence that speciation can occur without geographic
  separation (sympatric speciation) using the apple maggot fly
  (\emph{Rhagoletis pomonella}) host-race formation as a case study.
  What conditions (assortative mating by habitat, additive gene effects,
  limited gene flow) are required, and what genetic evidence suggests
  incipient sympatric speciation is occurring?
\end{enumerate}

\begin{enumerate}
\def\labelenumi{\arabic{enumi}.}
\setcounter{enumi}{21}
\tightlist
\item
  Design a population genomics study to determine whether two
  morphologically distinct bird populations on isolated islands should
  be classified as separate species. Describe: the genomic data you
  would collect (RAD-seq or whole-genome resequencing), the analyses you
  would perform (Fst, phylogeny, ABBA-BABA test for gene flow), and the
  threshold values you would use to make the classification decision.
\end{enumerate}

\begin{enumerate}
\def\labelenumi{\arabic{enumi}.}
\setcounter{enumi}{22}
\tightlist
\item
  Evaluate the ``glacial refugia'' hypothesis for European biodiversity:
  during Pleistocene glaciations, populations were isolated in southern
  refugia (Iberia, Italy, Balkans), underwent allopatric evolution, then
  expanded northward after glaciers retreated. What phylogeographic data
  would test this hypothesis, and what consequences does secondary
  contact of previously isolated populations have for hybrid zone
  formation?
\end{enumerate}

\begin{enumerate}
\def\labelenumi{\arabic{enumi}.}
\setcounter{enumi}{23}
\tightlist
\item
  Gene drive technology (e.g., CRISPR-based homing drives) can spread a
  gene (e.g., infertility) through a wild population at rates faster
  than Mendelian inheritance. Evaluate the ecological risks of releasing
  a gene drive into a wild mosquito population to suppress malaria
  transmission: (a) off-target spread to non-target species; (b)
  irreversibility; (c) international governance challenges (the drive
  does not respect national borders).
\end{enumerate}

\begin{enumerate}
\def\labelenumi{\arabic{enumi}.}
\setcounter{enumi}{24}
\tightlist
\item
  Inbreeding depression (reduced fitness in highly homozygous
  individuals) limits the viability of small, isolated populations.
  Evaluate: (a) the mechanisms of inbreeding depression (overdominance
  vs partial dominance of deleterious recessives); (b) the threshold
  F\_cr (critical inbreeding coefficient) below which genetic rescue is
  needed; (c) policy implications for reintroduction of individuals from
  genetically distinct populations.
\end{enumerate}

\begin{enumerate}
\def\labelenumi{\arabic{enumi}.}
\setcounter{enumi}{25}
\tightlist
\item
  The concept of ``reproductive character displacement'' predicts that
  the signal used for mate recognition must diverge faster between
  sympatric sister species than between allopatric ones, to prevent
  costly hybridisation. Evaluate the evidence for reproductive character
  displacement in: (a) \emph{Ficedula} flycatchers in Scandinavia; (b)
  cricket frog calling frequencies; (c) \emph{Phlox drummondii} flower
  colour contrast.
\end{enumerate}

\begin{enumerate}
\def\labelenumi{\arabic{enumi}.}
\setcounter{enumi}{26}
\tightlist
\item
  Evaluate the genomics of speciation: ``genomic islands of divergence''
  are regions of elevated Fst interspersed with regions of low Fst in
  partially reproductively isolated populations. Interpret these
  patterns: do islands represent targets of selection, or artefacts of
  reduced recombination around centromeres? What sequencing/analysis
  strategy would distinguish selective sweeps from reduced recombination
  as drivers of Fst peaks?
\end{enumerate}

\begin{enumerate}
\def\labelenumi{\arabic{enumi}.}
\setcounter{enumi}{27}
\tightlist
\item
  Evaluate the concept of the ``speciation continuum'': from partially
  isolated populations (e.g., Fst = 0.1) through incipient species
  (partial reproductive isolation, fertile hybrids) to full species
  (complete reproductive isolation). At what point on this continuum
  should conservation efforts classify distinct evolutionary significant
  units (ESUs) for legal protection?
\end{enumerate}

\begin{enumerate}
\def\labelenumi{\arabic{enumi}.}
\setcounter{enumi}{28}
\tightlist
\item
  Mass extinction events (e.g., end-Cretaceous, 66 Mya) dramatically
  restructured biodiversity. Evaluate: (a) what evolutionary patterns
  characterise the burst of diversification after mass extinction
  (adaptive radiation); (b) whether the sixth mass extinction (current
  human-driven biodiversity crisis) is comparable in rate to previous
  ones; (c) whether survivors of previous extinctions provide models for
  predicting which present-day species will persist.
\end{enumerate}

\begin{enumerate}
\def\labelenumi{\arabic{enumi}.}
\setcounter{enumi}{29}
\tightlist
\item
  Critically evaluate the ``two-speed genome'' concept in fungi
  (\emph{Zymoseptoria tritici}, wheat pathogen): the genome is divided
  into ``core'' (gene-rich, slowly evolving) and ``accessory''
  (transposon-rich, rapidly evolving) chromosomes. How does this
  architecture accelerate host-pathogen coevolution, and how might it
  apply to our understanding of evolutionary rate variation in other
  organisms, including humans (conserved regions vs rapidly evolving
  regions)?
\end{enumerate}

\newpage

\section*{Questions --- Ch 21: Phylogenetics and the Tree of
Life}\label{questions-ch-21-phylogenetics-and-the-tree-of-life}
\addcontentsline{toc}{section}{Questions --- Ch 21: Phylogenetics and
the Tree of Life}

\label{sec:q_unit_VI_phylogenetics}

\subsection*{Instructor Use and Coverage
Notes}\label{instructor-use-and-coverage-notes-23}
\addcontentsline{toc}{subsection}{Instructor Use and Coverage Notes}

\begin{itemize}
\tightlist
\item
  \textbf{Coverage target:} Use a tree to infer shared ancestry,
  character transitions, or divergence time.
\item
  \textbf{Model/data emphasis:} Tree-distance, parsimony, and
  molecular-clock calculations.
\item
  \textbf{Assessment alignment:} Visual Representations, Statistical
  Tests and Data Analysis, Argumentation.
\item
  \textbf{Misconception probe:} Living species are cousins, not
  ancestors of one another.
\item
  \textbf{Transfer product:} Transfer tree thinking to pathogens,
  conservation units, gene families, and development.
\item
  \textbf{Grading focus:} award full credit for mechanism, evidence,
  boundary conditions, and units when a calculation is required.
\item
  \textbf{Suggested use:} draw one recall item, one application item,
  and one synthesis item when building a short quiz from this bank.
\end{itemize}

\subsection*{Questions 1--10: Recall and
Comprehension}\label{questions-110-recall-and-comprehension-20}
\addcontentsline{toc}{subsection}{Questions 1--10: Recall and
Comprehension}

\emph{This activity accompanies \cref{sec:unit_VI_phylogenetics} of the
textbook --- review that chapter before attempting the exercises below.}

\begin{enumerate}
\def\labelenumi{\arabic{enumi}.}
\tightlist
\item
  What is a phylogenetic tree? What does each node, branch, and tip
  represent?
\end{enumerate}

\begin{enumerate}
\def\labelenumi{\arabic{enumi}.}
\setcounter{enumi}{1}
\tightlist
\item
  What is a clade (monophyletic group)? Distinguish monophyletic,
  paraphyletic, and polyphyletic groups.
\end{enumerate}

\begin{enumerate}
\def\labelenumi{\arabic{enumi}.}
\setcounter{enumi}{2}
\tightlist
\item
  Define synapomorphy. How does it differ from a plesiomorphy?
\end{enumerate}

\begin{enumerate}
\def\labelenumi{\arabic{enumi}.}
\setcounter{enumi}{3}
\tightlist
\item
  What is parsimony in phylogenetics? What does the most parsimonious
  tree minimise?
\end{enumerate}

\begin{enumerate}
\def\labelenumi{\arabic{enumi}.}
\setcounter{enumi}{4}
\tightlist
\item
  What is a molecular phylogeny? What type of data is used to construct
  one?
\end{enumerate}

\begin{enumerate}
\def\labelenumi{\arabic{enumi}.}
\setcounter{enumi}{5}
\tightlist
\item
  What is a bootstrap value? What does a bootstrap of 90\% indicate?
\end{enumerate}

\begin{enumerate}
\def\labelenumi{\arabic{enumi}.}
\setcounter{enumi}{6}
\tightlist
\item
  What is an outgroup? How is it used to root a phylogenetic tree?
\end{enumerate}

\begin{enumerate}
\def\labelenumi{\arabic{enumi}.}
\setcounter{enumi}{7}
\tightlist
\item
  What is horizontal gene transfer (HGT)? In which domains of life is it
  most common?
\end{enumerate}

\begin{enumerate}
\def\labelenumi{\arabic{enumi}.}
\setcounter{enumi}{8}
\tightlist
\item
  Define convergent evolution. How can it mislead phylogenetic analysis?
\end{enumerate}

\begin{enumerate}
\def\labelenumi{\arabic{enumi}.}
\setcounter{enumi}{9}
\tightlist
\item
  What are molecular clocks? How are they calibrated?
\end{enumerate}

\subsection*{Questions 11--20: Application and
Analysis}\label{questions-1120-application-and-analysis-20}
\addcontentsline{toc}{subsection}{Questions 11--20: Application and
Analysis}

\begin{enumerate}
\def\labelenumi{\arabic{enumi}.}
\setcounter{enumi}{10}
\tightlist
\item
  Given the character matrix below (1=present, 0=absent): determine the
  synapomorphies that define the clade (birds + crocodiles) relative to
  lizards and snakes, using the parsimony principle.
\end{enumerate}

{\def\LTcaptype{none} % do not increment counter
\begin{longtable}[]{@{}lllll@{}}
\toprule\noalign{}
Taxon & Archosaur skull & Feathers & 4-chambered heart & Endothermy \\
\midrule\noalign{}
\endhead
\bottomrule\noalign{}
\endlastfoot
Lizard & 0 & 0 & 0 & 0 \\
Crocodile & 1 & 0 & 1 & 0 \\
Bird & 1 & 1 & 1 & 1 \\
\end{longtable}
}

\begin{enumerate}
\def\labelenumi{\arabic{enumi}.}
\setcounter{enumi}{11}
\tightlist
\item
  Two competing phylogenies are proposed for the Paleocene mammals: (A)
  requires 12 evolutionary changes; (B) requires 15 changes. By
  parsimony, which tree is preferred? Explain what assumption is made
  when choosing the more parsimonious tree.
\end{enumerate}

\begin{enumerate}
\def\labelenumi{\arabic{enumi}.}
\setcounter{enumi}{12}
\tightlist
\item
  Dolphins and ichthyosaurs (extinct reptiles) both have streamlined
  bodies and fin-like forelimbs. Explain: (a) why a morphological
  phylogeny based solely on body shape would incorrectly group them; (b)
  what more reliable character would avoid this error; (c) what type of
  evolution produced the similarity.
\end{enumerate}

\begin{enumerate}
\def\labelenumi{\arabic{enumi}.}
\setcounter{enumi}{13}
\tightlist
\item
  A molecular phylogeny based on 16S rRNA sequences places Archaea as a
  sister clade to Eukaryotes, while Bacteria branch more basely. Explain
  what this means for the ``three-domain of life'' vs ``two-domain of
  life'' (Eocyte hypothesis) debate.
\end{enumerate}

\begin{enumerate}
\def\labelenumi{\arabic{enumi}.}
\setcounter{enumi}{14}
\tightlist
\item
  The molecular clock for mitochondrial cytochrome b sequences in
  mammals ticks at about 2\% divergence per million years
  (Mya). Two sister species differ by 8\% in their cytochrome b
  sequences. Estimate when they diverged, and list two sources of
  uncertainty in this estimate.
\end{enumerate}

\begin{enumerate}
\def\labelenumi{\arabic{enumi}.}
\setcounter{enumi}{15}
\tightlist
\item
  ABBA-BABA (D-statistics) tests detect gene flow between populations.
  If D is significantly positive, what does it imply about the
  relationships between the four taxa (P1, P2, P3, and outgroup)? How
  did this test reveal Neanderthal--modern human gene flow?
\end{enumerate}

\begin{enumerate}
\def\labelenumi{\arabic{enumi}.}
\setcounter{enumi}{16}
\tightlist
\item
  Biogeographic analysis shows that the freshwater fish genus
  \emph{Nannoperca} is distributed across southwestern Australia
  (isolated from other parts by arid zones). How would you distinguish
  between: (a) vicariance (ancient populations separated by geological
  changes); (b) long-distance dispersal (a single colonisation event
  across an unsuitable barrier)?
\end{enumerate}

\begin{enumerate}
\def\labelenumi{\arabic{enumi}.}
\setcounter{enumi}{17}
\tightlist
\item
  Maximum likelihood (ML) phylogenetics uses an explicit model of
  sequence evolution (e.g., GTR+\texorpdfstring{\ensuremath{\Gamma}}{Gamma}) to evaluate the likelihood of
  observing the data given a tree. Why is ML generally preferred over
  parsimony for molecular data, especially for distantly related taxa?
\end{enumerate}

\begin{enumerate}
\def\labelenumi{\arabic{enumi}.}
\setcounter{enumi}{18}
\tightlist
\item
  Gene trees vs species trees: individual gene trees can differ from the
  species tree due to incomplete lineage sorting (ILS) or horizontal
  gene transfer. What is ILS, and why is it more likely for genes with
  short internodal branches (rapid speciation)?
\end{enumerate}

\begin{enumerate}
\def\labelenumi{\arabic{enumi}.}
\setcounter{enumi}{19}
\tightlist
\item
  The ``long branch attraction'' (LBA) artefact occurs when rapidly
  evolving lineages with many changes falsely cluster together in
  parsimony or distance analysis. Explain the mechanism of LBA and
  propose two strategies to avoid it (taxon sampling, model-based
  methods).
\end{enumerate}

\subsection*{Questions 21--30: Synthesis and
Evaluation}\label{questions-2130-synthesis-and-evaluation-20}
\addcontentsline{toc}{subsection}{Questions 21--30: Synthesis and
Evaluation}

\begin{enumerate}
\def\labelenumi{\arabic{enumi}.}
\setcounter{enumi}{20}
\tightlist
\item
  Evaluate the evidence that Archaean clade Asgard is the closest
  prokaryotic relative of eukaryotes. What specific eukaryotic features
  (endomembrane, cytoskeleton homologues, ubiquitin) were identified in
  their genomes? What does this suggest about the origin of eukaryotic
  complexity?
\end{enumerate}

\begin{enumerate}
\def\labelenumi{\arabic{enumi}.}
\setcounter{enumi}{21}
\tightlist
\item
  Design a phylogenomic study (using hundreds of genes rather than one
  marker) to resolve a ``soft polytomy'' (unresolved node) in the avian
  tree of life. Describe: coalescent methods (ASTRAL) vs concatenation
  (ML on combined alignment), what sources of conflict between gene
  trees need to be accounted for, and how you would determine whether
  the polytomy reflects rapid radiation rather than incomplete data.
\end{enumerate}

\begin{enumerate}
\def\labelenumi{\arabic{enumi}.}
\setcounter{enumi}{22}
\tightlist
\item
  The ribosomal RNA (rRNA) phylogeny (Woese, 1977) unified most life
  into three domains. Evaluate why rRNA was chosen as a comprehensive
  phylogenetic marker: (a) ubiquitous, almost universally homologous;
  (b) highly conserved but containing variable regions; (c) unlikely to
  undergo HGT. Are there limitations to rRNA as a ``true'' species tree
  marker?
\end{enumerate}

\begin{enumerate}
\def\labelenumi{\arabic{enumi}.}
\setcounter{enumi}{23}
\tightlist
\item
  Bayesian phylogenetics (MrBayes, BEAST) uses prior probabilities and
  MCMC sampling to estimate the posterior distribution of trees. Compare
  Bayesian vs maximum likelihood approaches: (a) what is the difference
  in interpretation of the ``support'' values (posterior probabilities
  vs bootstrap); (b) when does Bayesian inference provide overconfident
  node support; (c) how is divergence time estimation performed in
  BEAST?
\end{enumerate}

\begin{enumerate}
\def\labelenumi{\arabic{enumi}.}
\setcounter{enumi}{24}
\tightlist
\item
  Ancient DNA (aDNA) sequencing has transformed our understanding of
  human history by adding extinct hominin and historical populations to
  phylogenies. Evaluate the technical challenges of aDNA analysis: (a)
  DNA degradation (short fragments, C\texorpdfstring{\ensuremath{\to}}{->}U deamination); (b) contamination
  from environmental microbes and modern human DNA; (c) how mapDamage
  software and UDG treatment are used to address these.
\end{enumerate}

\begin{enumerate}
\def\labelenumi{\arabic{enumi}.}
\setcounter{enumi}{25}
\tightlist
\item
  Evaluate the ``bush of life'' metaphor vs the traditional ``tree of
  life'' for describing microbial evolution. How prevalent is HGT among
  bacteria and archaea (what fraction of the typical prokaryote genome
  has been acquired by HGT), and at what level (genus, family, phylum)
  does HGT obscure true vertical inheritance?
\end{enumerate}

\begin{enumerate}
\def\labelenumi{\arabic{enumi}.}
\setcounter{enumi}{26}
\tightlist
\item
  TransposableElements have been used as phylogenetic markers
  (SINEs/LINEs as synapomorphies): a SINE insertion shared by multiple
  taxa must have occurred once in a common ancestor. Evaluate the
  advantages of transposon-based markers over sequence-based
  phylogenetics: (a) irreversibility of insertion; (b) independence of
  molecular evolution model assumptions; (c) limitations (element
  excision, parallel insertions in hotspot regions).
\end{enumerate}

\begin{enumerate}
\def\labelenumi{\arabic{enumi}.}
\setcounter{enumi}{27}
\tightlist
\item
  The ``total evidence dating'' approach (tip calibration) uses
  morphological characters from fossils as direct terminal taxa (not
  just calibration points) in Bayesian time-tree analysis. Evaluate the
  advantages of total evidence dating over node calibration, and
  describe what happens to morphological branch length estimates when
  the model of trait evolution (e.g., Mk model) is misspecified.
\end{enumerate}

\begin{enumerate}
\def\labelenumi{\arabic{enumi}.}
\setcounter{enumi}{28}
\tightlist
\item
  Horizontal gene transfer in cancer (somatic HGT from tumour
  microenvironment bacteria into human cells) has been reported but
  remains controversial. Evaluate the evidence: (a) what sequencing
  artefacts can mimic bacterial HGT; (b) what controls are needed to
  rule out contamination; (c) if real, what functional consequences
  might cancer-incorporated bacterial genes have?
\end{enumerate}

\begin{enumerate}
\def\labelenumi{\arabic{enumi}.}
\setcounter{enumi}{29}
\tightlist
\item
  Critically assess whether a ``tree of life'' can be reconstructed from
  extant biodiversity alone, without fossil calibration. What
  information is intrinsically lost when primarily the tips (extant
  species) and not the internal nodes (extinct ancestors) are known? Use
  the concept of incomplete ancestry and the ``big-bang'' topology
  problem of rapid early diversification as a specific example.
\end{enumerate}

\newpage

\section*{Questions --- Ch 22: Bacteria, Archaea, and
Viruses}\label{questions-ch-22-bacteria-archaea-and-viruses}
\addcontentsline{toc}{section}{Questions --- Ch 22: Bacteria, Archaea,
and Viruses}

\label{sec:q_unit_VII_bacteria_archaea_viruses}

\subsection*{Instructor Use and Coverage
Notes}\label{instructor-use-and-coverage-notes-24}
\addcontentsline{toc}{subsection}{Instructor Use and Coverage Notes}

\begin{itemize}
\tightlist
\item
  \textbf{Coverage target:} Classify a microbe or virus from evidence
  and predict a likely vulnerability.
\item
  \textbf{Model/data emphasis:} Growth-rate, genome-size, and dilution
  calculations.
\item
  \textbf{Assessment alignment:} Questions and Methods, Representing and
  Describing Data, Argumentation.
\item
  \textbf{Misconception probe:} Viruses are not simply tiny bacteria;
  they use fundamentally different replication logic.
\item
  \textbf{Transfer product:} Apply microbial diversity reasoning to
  antibiotics, biotechnology, ecology, and outbreaks.
\item
  \textbf{Grading focus:} award full credit for mechanism, evidence,
  boundary conditions, and units when a calculation is required.
\item
  \textbf{Suggested use:} draw one recall item, one application item,
  and one synthesis item when building a short quiz from this bank.
\end{itemize}

\subsection*{Questions 1--10: Recall and
Comprehension}\label{questions-110-recall-and-comprehension-21}
\addcontentsline{toc}{subsection}{Questions 1--10: Recall and
Comprehension}

\emph{This activity accompanies
\cref{sec:unit_VII_bacteria_archaea_viruses} of the textbook --- review
that chapter before attempting the exercises below.}

\begin{enumerate}
\def\labelenumi{\arabic{enumi}.}
\tightlist
\item
  List three differences between bacteria and archaea in cell wall
  chemistry and membrane lipids.
\end{enumerate}

\begin{enumerate}
\def\labelenumi{\arabic{enumi}.}
\setcounter{enumi}{1}
\tightlist
\item
  What is the Gram stain? What does Gram-positive vs Gram-negative
  indicate about cell wall structure?
\end{enumerate}

\begin{enumerate}
\def\labelenumi{\arabic{enumi}.}
\setcounter{enumi}{2}
\tightlist
\item
  Describe the replication cycle of a bacteriophage. Distinguish lytic
  from lysogenic cycles.
\end{enumerate}

\begin{enumerate}
\def\labelenumi{\arabic{enumi}.}
\setcounter{enumi}{3}
\tightlist
\item
  What is a prophage? What signal triggers induction of the lytic cycle
  from lysogeny?
\end{enumerate}

\begin{enumerate}
\def\labelenumi{\arabic{enumi}.}
\setcounter{enumi}{4}
\tightlist
\item
  Name three types of bacterial reproduction. Which is the fastest?
\end{enumerate}

\begin{enumerate}
\def\labelenumi{\arabic{enumi}.}
\setcounter{enumi}{5}
\tightlist
\item
  What is binary fission? How does the doubling time relate to bacterial
  growth rate?
\end{enumerate}

\begin{enumerate}
\def\labelenumi{\arabic{enumi}.}
\setcounter{enumi}{6}
\tightlist
\item
  Describe the structure of a typical bacterium: nucleoid, ribosomes,
  pili, flagella, capsule.
\end{enumerate}

\begin{enumerate}
\def\labelenumi{\arabic{enumi}.}
\setcounter{enumi}{7}
\tightlist
\item
  What is conjugation? What is an F plasmid?
\end{enumerate}

\begin{enumerate}
\def\labelenumi{\arabic{enumi}.}
\setcounter{enumi}{8}
\tightlist
\item
  Distinguish DNA viruses from RNA viruses. Give one example of each.
\end{enumerate}

\begin{enumerate}
\def\labelenumi{\arabic{enumi}.}
\setcounter{enumi}{9}
\tightlist
\item
  What is a retrovirus? Describe its replication cycle.
\end{enumerate}

\subsection*{Questions 11--20: Application and
Analysis}\label{questions-1120-application-and-analysis-21}
\addcontentsline{toc}{subsection}{Questions 11--20: Application and
Analysis}

\begin{enumerate}
\def\labelenumi{\arabic{enumi}.}
\setcounter{enumi}{10}
\tightlist
\item
  A bacterial culture starts with 100 cells and has a doubling time of
  20 minutes. Calculate: (a) the number of cells after 2 hours; (b) the
  approximate mass of bacteria (assuming 10\texorpdfstring{\textsuperscript{-}}{-}\texorpdfstring{\textsuperscript{1}}{1}\texorpdfstring{\textsuperscript{5}}{5} g/cell) after 24 hours.
  Why does exponential growth not continue indefinitely?
\end{enumerate}

\begin{enumerate}
\def\labelenumi{\arabic{enumi}.}
\setcounter{enumi}{11}
\tightlist
\item
  Design a serial dilution experiment to count colonies from a bacterial
  sample expected to have about 10\texorpdfstring{\textsuperscript{6}}{6} CFU/mL. Which dilutions
  would give countable plates (30--300 colonies)? Calculate CFU/mL from
  a plate showing 175 colonies at 10\texorpdfstring{\textsuperscript{-}}{-}\texorpdfstring{\textsuperscript{5}}{5} dilution.
\end{enumerate}

\begin{enumerate}
\def\labelenumi{\arabic{enumi}.}
\setcounter{enumi}{12}
\tightlist
\item
  Compare the lytic and lysogenic cycles: what are the advantages of
  each for the phage? Under what environmental/physiological conditions
  does the lambda phage switch from lysogeny to the lytic cycle? (Name
  the specific protease.)
\end{enumerate}

\begin{enumerate}
\def\labelenumi{\arabic{enumi}.}
\setcounter{enumi}{13}
\tightlist
\item
  MRSA (methicillin-resistant \emph{Staphylococcus aureus}) carries
  \emph{mecA} or \emph{mecC}, encoding a modified penicillin-binding
  protein (PBP2a/PBP2c) with low affinity for most traditional \texorpdfstring{\ensuremath{\beta}}{beta}-lactam
  antibiotics. Explain: (a) how \emph{mec} genes are acquired by
  horizontal gene transfer on SCC\emph{mec} elements; (b) why many
  \texorpdfstring{\ensuremath{\beta}}{beta}-lactams fail but anti-MRSA cephalosporins can be exceptions in
  susceptible isolates; (c) what alternative antibiotics are used and
  how their mechanisms bypass PBP2a.
\end{enumerate}

\begin{enumerate}
\def\labelenumi{\arabic{enumi}.}
\setcounter{enumi}{14}
\tightlist
\item
  Archaea are found in extreme environments (hyperthermophiles,
  halophiles, methanogens). Explain how: (a) archaeal lipid membranes
  (isoprene ether-linked, rather than ester-linked) withstand high
  temperatures; (b) halophilic archaea stabilise proteins at high salt
  by having an excess of surface-facing acidic residues; (c)
  methanogenic archaea use H\texorpdfstring{\textsubscript{2}}{2} and CO\texorpdfstring{\textsubscript{2}}{2} as energy sources via a unique
  archaeal pathway.
\end{enumerate}

\begin{enumerate}
\def\labelenumi{\arabic{enumi}.}
\setcounter{enumi}{15}
\tightlist
\item
  The CRISPR-Cas system in bacteria provides adaptive immunity against
  phage. Explain: (a) spacer acquisition from phage DNA; (b)
  crRNA-guided Cas9 surveillance; (c) cleavage of matching invader DNA;
  (d) why the phage can evolve to escape CRISPR immunity (PAM sequence
  mutation). How did bacterial CRISPR inspire the molecular biology
  tool?
\end{enumerate}

\begin{enumerate}
\def\labelenumi{\arabic{enumi}.}
\setcounter{enumi}{16}
\tightlist
\item
  Viral tropism (which cells a virus can infect) depends on
  receptor-ligand binding. HIV infects cells expressing both CD4 and
  CCR5/CXCR4 (co-receptors). Individuals homozygous for a 32-bp deletion
  in \emph{CCR5} (\texorpdfstring{\ensuremath{\Delta}}{Delta}32 allele) are largely resistant to HIV infection.
  Explain why this is and why CCR5 antagonists (maraviroc) alone do not
  fully cure HIV.
\end{enumerate}

\begin{enumerate}
\def\labelenumi{\arabic{enumi}.}
\setcounter{enumi}{17}
\tightlist
\item
  Prions are misfolded proteins (PrP\^{}Sc) that convert normally folded
  proteins (PrP\^{}C) to the pathogenic conformation. Explain: (a) why
  prion diseases (CJD, BSE) are transmissible without nucleic acid; (b)
  the proposed templating mechanism; (c) why prions are nearly extremely
  difficult to decontaminate by standard autoclaving.
\end{enumerate}

\begin{enumerate}
\def\labelenumi{\arabic{enumi}.}
\setcounter{enumi}{18}
\tightlist
\item
  Bacteriophage therapy (phage therapy) is being reconsidered as an
  alternative to antibiotics for drug-resistant infections. Compare
  phage therapy vs antibiotics: (a) host range (broad-spectrum
  antibiotics vs phage-host specificity); (b) resistance evolution rate;
  (c) phage ability to co-evolve with bacterial host. Give one clinical
  case report of successful phage therapy.
\end{enumerate}

\begin{enumerate}
\def\labelenumi{\arabic{enumi}.}
\setcounter{enumi}{19}
\tightlist
\item
  SARS-CoV-2 is a +ssRNA coronavirus. Explain: (a) how it uses its spike
  protein to bind ACE2 and enter cells; (b) why mRNA vaccines encoding
  spike protein induce protective immunity; (c) why variants (Omicron)
  with many spike mutations can partially evade vaccine-induced
  antibodies.
\end{enumerate}

\subsection*{Questions 21--30: Synthesis and
Evaluation}\label{questions-2130-synthesis-and-evaluation-21}
\addcontentsline{toc}{subsection}{Questions 21--30: Synthesis and
Evaluation}

\begin{enumerate}
\def\labelenumi{\arabic{enumi}.}
\setcounter{enumi}{20}
\tightlist
\item
  Evaluate the hypotheses for the origin of viruses: (a) regression
  hypothesis (viruses are degenerate cells); (b) escape hypothesis
  (mobile genetic elements that gained capsids); (c) primordial virion
  hypothesis (viruses predate cells). Which hypothesis is most
  consistent with current phylogenomic data (nucleocytoplasmic large DNA
  viruses and their relatives)?
\end{enumerate}

\begin{enumerate}
\def\labelenumi{\arabic{enumi}.}
\setcounter{enumi}{21}
\tightlist
\item
  Design a strategy to engineer a phage to specifically kill a
  Gram-negative ESKAPE pathogen (\emph{Klebsiella pneumoniae}) that is
  resistant to multiple first-line antibiotics and has few standard
  options. Describe: (a) the host-range engineering approach (tail fibre
  protein swapping); (b) the anti-biofilm component you would add; (c)
  the regulatory and safety evidence required before use in humans.
\end{enumerate}

\begin{enumerate}
\def\labelenumi{\arabic{enumi}.}
\setcounter{enumi}{22}
\tightlist
\item
  Quorum sensing in \emph{Pseudomonas aeruginosa} controls virulence
  gene expression via N-acyl-homoserine lactone (AHL) signalling (las
  and rhl systems). Evaluate: (a) how quorum sensing allows coordinated
  virulence (biofilm formation, exotoxin production) primarily at high
  cell density; (b) why quorum sensing inhibitors (QSIs) are attractive
  anti-virulence targets (less resistance selection pressure); (c) why
  QSI clinical development has stalled.
\end{enumerate}

\begin{enumerate}
\def\labelenumi{\arabic{enumi}.}
\setcounter{enumi}{23}
\tightlist
\item
  Evaluate the evolutionary biology of rapid viral evolution (HIV,
  influenza). HIV's high mutation rate (error-prone RT,
  about 3\texorpdfstring{\ensuremath{\times}}{x}10\texorpdfstring{\textsuperscript{-}}{-}\texorpdfstring{\textsuperscript{5}}{5}/base/replication) generates enormous diversity.
  How does this diversity allow immune escape and drug resistance?
  Evaluate why combination antiretroviral therapy (ART) --- targeting
  multiple steps simultaneously --- dramatically suppresses resistance
  evolution and how ``cure'' strategies (shock and kill) aim to
  eliminate latent reservoirs.
\end{enumerate}

\begin{enumerate}
\def\labelenumi{\arabic{enumi}.}
\setcounter{enumi}{24}
\tightlist
\item
  Giant viruses (Mimivirus, Pandoravirus) have genomes \textgreater1 Mb
  --- larger than some bacteria --- and contain genes for
  translation-related functions unprecedented in viruses. Evaluate
  whether giant viruses blur the boundary between viruses and cells, and
  whether they support the regression hypothesis for viral origin.
\end{enumerate}

\begin{enumerate}
\def\labelenumi{\arabic{enumi}.}
\setcounter{enumi}{25}
\tightlist
\item
  Horizontal gene transfer (HGT) in bacteria can be transduction,
  transformation, or conjugation. Use a One Health AMR map to compare
  how an ESBL or carbapenemase gene could move among hospital patients,
  wastewater, livestock, soil microbes, and wildlife. Which transfer
  mechanism is most plausible in each setting, and what surveillance
  evidence would distinguish clonal spread from mobile-element spread?
\end{enumerate}

\begin{enumerate}
\def\labelenumi{\arabic{enumi}.}
\setcounter{enumi}{26}
\tightlist
\item
  The microbiome of the human gut contains \textgreater1,000 bacterial
  species that collectively perform essential metabolic functions (SCFA
  production, vitamin synthesis, immune training). Evaluate whether the
  gut microbiome should be considered an ``accessory genome'' of the
  human host, and what criteria would need to be met for this ``extended
  genome'' concept to be scientifically defensible.
\end{enumerate}

\begin{enumerate}
\def\labelenumi{\arabic{enumi}.}
\setcounter{enumi}{27}
\tightlist
\item
  Bacteriophage encode auxiliary metabolic genes (AMGs) --- metabolic
  genes that boost host metabolism during infection (e.g., phage-encoded
  photosynthesis genes in cyanophages). Evaluate the role of phage AMGs
  in global carbon and nitrogen cycling, and propose how phage-encoded
  metabolic capacity shapes marine primary productivity.
\end{enumerate}

\begin{enumerate}
\def\labelenumi{\arabic{enumi}.}
\setcounter{enumi}{28}
\tightlist
\item
  The virome --- the complete collection of viruses in an ecosystem ---
  is far larger than the known virome catalogued to date. Evaluate
  metagenomics approaches (shotgun metagenomics, viromics with CRISPR
  spacer-based enrichment) for discovering novel viruses, and assess
  what fraction of global viral diversity may remain undiscovered.
\end{enumerate}

\begin{enumerate}
\def\labelenumi{\arabic{enumi}.}
\setcounter{enumi}{29}
\tightlist
\item
  Critically evaluate the claim that viruses are ``the dark matter of
  the biosphere'' --- underappreciated forces shaping ecosystems, gene
  flows, and evolution. What evidence from marine phage ecology, the
  pan-genome of bacteria, and viral integration into eukaryotic genomes
  supports this claim? How does the expanding virology of the
  2010s--2020s change our understanding of life's diversity?
\end{enumerate}

\newpage

\section*{Questions --- Ch 23: Microbial Ecology and the
Microbiome}\label{questions-ch-23-microbial-ecology-and-the-microbiome}
\addcontentsline{toc}{section}{Questions --- Ch 23: Microbial Ecology
and the Microbiome}

\label{sec:q_unit_VII_microbial_ecology}

\subsection*{Instructor Use and Coverage
Notes}\label{instructor-use-and-coverage-notes-25}
\addcontentsline{toc}{subsection}{Instructor Use and Coverage Notes}

\begin{itemize}
\tightlist
\item
  \textbf{Coverage target:} Use community data to infer diversity,
  dominance, and a plausible ecological mechanism.
\item
  \textbf{Model/data emphasis:} Diversity indices, richness estimates,
  and interaction-network reasoning.
\item
  \textbf{Assessment alignment:} Questions and Methods, Representing and
  Describing Data, Argumentation.
\item
  \textbf{Misconception probe:} A microbiome is not automatically
  beneficial; context determines the effect.
\item
  \textbf{Transfer product:} Transfer microbial ecology to soils,
  oceans, digestion, disease, and climate feedbacks.
\item
  \textbf{Grading focus:} award full credit for mechanism, evidence,
  boundary conditions, and units when a calculation is required.
\item
  \textbf{Suggested use:} draw one recall item, one application item,
  and one synthesis item when building a short quiz from this bank.
\end{itemize}

\subsection*{Questions 1--10: Recall and
Comprehension}\label{questions-110-recall-and-comprehension-22}
\addcontentsline{toc}{subsection}{Questions 1--10: Recall and
Comprehension}

\emph{This activity accompanies \cref{sec:unit_VII_microbial_ecology} of
the textbook --- review that chapter before attempting the exercises
below.}

\begin{enumerate}
\def\labelenumi{\arabic{enumi}.}
\tightlist
\item
  Define the microbiome. What is the difference between microbiome and
  microbiota?
\end{enumerate}

\begin{enumerate}
\def\labelenumi{\arabic{enumi}.}
\setcounter{enumi}{1}
\tightlist
\item
  What is 16S rRNA sequencing? Why is the 16S gene used to classify
  bacteria?
\end{enumerate}

\begin{enumerate}
\def\labelenumi{\arabic{enumi}.}
\setcounter{enumi}{2}
\tightlist
\item
  Define the Shannon diversity index. What does a high value indicate?
\end{enumerate}

\begin{enumerate}
\def\labelenumi{\arabic{enumi}.}
\setcounter{enumi}{3}
\tightlist
\item
  What are the four major phyla of the human gut microbiome?
\end{enumerate}

\begin{enumerate}
\def\labelenumi{\arabic{enumi}.}
\setcounter{enumi}{4}
\tightlist
\item
  What is a biofilm? Name two biological advantages biofilm growth
  confers on bacteria.
\end{enumerate}

\begin{enumerate}
\def\labelenumi{\arabic{enumi}.}
\setcounter{enumi}{5}
\tightlist
\item
  Define mutualism, commensalism, and parasitism. Give one example of
  each in microbial ecology.
\end{enumerate}

\begin{enumerate}
\def\labelenumi{\arabic{enumi}.}
\setcounter{enumi}{6}
\tightlist
\item
  What is the rhizosphere? Name one bacterium that benefits plant growth
  there.
\end{enumerate}

\begin{enumerate}
\def\labelenumi{\arabic{enumi}.}
\setcounter{enumi}{7}
\tightlist
\item
  What are short-chain fatty acids (SCFAs)? Name three and state which
  gut bacteria produce them.
\end{enumerate}

\begin{enumerate}
\def\labelenumi{\arabic{enumi}.}
\setcounter{enumi}{8}
\tightlist
\item
  Define dysbiosis. Associate dysbiosis with two human diseases.
\end{enumerate}

\begin{enumerate}
\def\labelenumi{\arabic{enumi}.}
\setcounter{enumi}{9}
\tightlist
\item
  What is the global nitrogen cycle? Name two types of bacteria that fix
  atmospheric nitrogen.
\end{enumerate}

\subsection*{Questions 11--20: Application and
Analysis}\label{questions-1120-application-and-analysis-22}
\addcontentsline{toc}{subsection}{Questions 11--20: Application and
Analysis}

\begin{enumerate}
\def\labelenumi{\arabic{enumi}.}
\setcounter{enumi}{10}
\tightlist
\item
  Fecal microbiome transplant (FMT) resolves recurrent
  \emph{Clostridioides difficile} infection in about 90\% of
  cases. Explain: (a) why antibiotic use predisposes patients to C.
  diff; (b) how FMT restores colonisation resistance; (c) why FMT is
  less likely to be effective in IBD (Crohn's, ulcerative colitis)?
\end{enumerate}

\begin{enumerate}
\def\labelenumi{\arabic{enumi}.}
\setcounter{enumi}{11}
\tightlist
\item
  A researcher performs 16S metagenomics on soil samples from fertile
  farmland, degraded land, and old-growth forest. Predict: which
  environment would show highest Shannon diversity; whether any taxa
  would serve as indicators of soil health; and what functional genes
  you would look for in shotgun metagenomics data.
\end{enumerate}

\begin{enumerate}
\def\labelenumi{\arabic{enumi}.}
\setcounter{enumi}{12}
\tightlist
\item
  The ``hygiene hypothesis'' (now refined as the ``old friends''
  hypothesis) proposes that reduced microbial exposure in early life
  increases allergic and autoimmune disease risk. Evaluate the evidence:
  (a) epidemiological (higher allergy rates in developed countries); (b)
  experimental (germ-free mice have hyperactive immune systems); (c)
  specific microbial species (Lactobacillus, Faecalibacterium
  prausnitzii) that prevent allergy.
\end{enumerate}

\begin{enumerate}
\def\labelenumi{\arabic{enumi}.}
\setcounter{enumi}{13}
\tightlist
\item
  Nitrogen cycle: N-fixing bacteria (\emph{Rhizobium} in legume nodules)
  convert N\texorpdfstring{\textsubscript{2}}{2} \texorpdfstring{\ensuremath{\to}}{->} NH\texorpdfstring{\textsubscript{3}}{3}. Nitrifying bacteria convert NH\texorpdfstring{\textsubscript{3}}{3} \texorpdfstring{\ensuremath{\to}}{->} NO\texorpdfstring{\textsubscript{2}}{2}\texorpdfstring{\textsuperscript{-}}{-} \texorpdfstring{\ensuremath{\to}}{->} NO\texorpdfstring{\textsubscript{3}}{3}\texorpdfstring{\textsuperscript{-}}{-}.
  Denitrifying bacteria convert NO\texorpdfstring{\textsubscript{3}}{3}\texorpdfstring{\textsuperscript{-}}{-} \texorpdfstring{\ensuremath{\to}}{->} N\texorpdfstring{\textsubscript{2}}{2}. Trace a nitrogen atom from
  atmospheric N\texorpdfstring{\textsubscript{2}}{2} through these steps into a plant protein, and then back
  to N\texorpdfstring{\textsubscript{2}}{2} in soil. Why is this cycle vulnerable to disruption by synthetic
  fertilisers?
\end{enumerate}

\begin{enumerate}
\def\labelenumi{\arabic{enumi}.}
\setcounter{enumi}{14}
\tightlist
\item
  Marine microbial ecology: \emph{Prochlorococcus} is the most abundant
  photosynthetic organism on Earth (about 10\texorpdfstring{\textsuperscript{2}}{2}\texorpdfstring{\textsuperscript{7}}{7} cells). Explain:
  (a) its adaptations for low-nutrient, high-light (and low-light) ocean
  zones; (b) how phage predation of Prochlorococcus releases dissolved
  organic carbon (DOC) fuelling heterotrophic bacteria; (c) what would
  happen to ocean productivity if Prochlorococcus were eliminated.
\end{enumerate}

\begin{enumerate}
\def\labelenumi{\arabic{enumi}.}
\setcounter{enumi}{15}
\tightlist
\item
  The human oral microbiome (about 700 species) produces dental
  caries via \emph{Streptococcus mutans} acidifying enamel via lactic
  acid fermentation. Explain: (a) how sucrose is a uniquely cariogenic
  sugar (forms glucan adhesion matrix via glucosyltransferase); (b) what
  pH triggers enamel demineralisation (\textless5.5); (c) how fluoride
  in drinking water chemically protects enamel.
\end{enumerate}

\begin{enumerate}
\def\labelenumi{\arabic{enumi}.}
\setcounter{enumi}{16}
\tightlist
\item
  Quorum sensing coordinates population-level behaviour in bacteria
  (biofilm formation, bioluminescence, sporulation). Explain in
  \emph{Vibrio harveyi}: (a) how two quorum sensing circuits (LuxN and
  LuxPQ) integrate autoinducer signals at different cell densities; (b)
  why \emph{Vibrio fischeri} in the squid light organ shows
  population-density-dependent luminescence; (c) how quorum quenching
  enzymes (AHL lactonase) disrupt this signalling.
\end{enumerate}

\begin{enumerate}
\def\labelenumi{\arabic{enumi}.}
\setcounter{enumi}{17}
\tightlist
\item
  Antibiotic resistance spreads horizontally via mobile genetic elements
  (plasmids, integrons, transposons). Trace how resistance to
  trimethoprim and ampicillin on a co-integrated resistance plasmid
  would affect \emph{E. coli} evolution in a hospital drain environment:
  (a) does selection for one resistance indirectly maintain the other?
  (b) how does this create a co-selection problem when antibiotics are
  used?
\end{enumerate}

\begin{enumerate}
\def\labelenumi{\arabic{enumi}.}
\setcounter{enumi}{18}
\tightlist
\item
  Extremophile archaea in deep-sea hydrothermal vents oxidise H\texorpdfstring{\textsubscript{2}}{2}S as
  electron donor, fixing CO\texorpdfstring{\textsubscript{2}}{2} via the reverse TCA cycle (without
  photosynthesis). Compare this chemolithoautotrophic primary production
  to photosynthesis: (a) what is the electron donor; (b) what energy
  source; (c) how much biomass is produced per unit energy; (d) why this
  is important for understanding the origin of life.
\end{enumerate}

\begin{enumerate}
\def\labelenumi{\arabic{enumi}.}
\setcounter{enumi}{19}
\tightlist
\item
  Meta-analysis of 10 gut microbiome studies in Type 2 Diabetes (T2D)
  patients vs healthy controls identified depletion of \emph{Akkermansia
  muciniphila} and \emph{Faecalibacterium prausnitzii} in T2D. Design a
  causal inference test: (a) what correlation measure would you use; (b)
  what germ-free mouse experiment would establish causality; (c) what
  confounders (diet, medication) need controlling?
\end{enumerate}

\subsection*{Questions 21--30: Synthesis and
Evaluation}\label{questions-2130-synthesis-and-evaluation-22}
\addcontentsline{toc}{subsection}{Questions 21--30: Synthesis and
Evaluation}

\begin{enumerate}
\def\labelenumi{\arabic{enumi}.}
\setcounter{enumi}{20}
\tightlist
\item
  Evaluate the ``microbiome-brain axis'' (gut-brain axis): propose the
  mechanistic routes (vagus nerve, SCFAs, immune mediators, enteric
  serotonin) by which gut microbial composition can influence mood and
  cognition. What evidence from: (a) germ-free mice (increased anxiety,
  HPA axis changes); (b) faecal transfer from depressed patients to
  rats; (c) probiotic clinical trials in humans --- supports a causal
  role?
\end{enumerate}

\begin{enumerate}
\def\labelenumi{\arabic{enumi}.}
\setcounter{enumi}{21}
\tightlist
\item
  Design a metagenomics study to characterise the resistome (complete
  antibiotic resistance gene complement) of wastewater treatment plant
  effluent. Describe: (a) DNA extraction strategy; (b) how you would
  annotate resistance genes; (c) how you would determine whether
  identified genes are in pathogens vs harmless commensals; (d) how you
  would track the resistome over a full seasonal cycle.
\end{enumerate}

\begin{enumerate}
\def\labelenumi{\arabic{enumi}.}
\setcounter{enumi}{22}
\tightlist
\item
  Evaluate the concept of the ``core microbiome'' --- a set of microbial
  taxa present in most healthy individuals of a species. What criteria
  would define a ``core'' species, how would you distinguish core from
  accessory microbiome species, and why the concept is complicated by
  inter-individual variation (about 20\% of gut species may be
  ``comprehensive core'').
\end{enumerate}

\begin{enumerate}
\def\labelenumi{\arabic{enumi}.}
\setcounter{enumi}{23}
\tightlist
\item
  Microbial succession in natural environments follows predictable
  patterns (e.g., pioneer \texorpdfstring{\ensuremath{\to}}{->} climax community in forest soils;
  successional microbiome during decomposition of whale falls). Evaluate
  the mechanisms (facilitation, inhibition, tolerance) controlling
  succession, and propose an experiment to test whether succession is
  driven by deterministic (niche-based) or stochastic processes.
\end{enumerate}

\begin{enumerate}
\def\labelenumi{\arabic{enumi}.}
\setcounter{enumi}{24}
\tightlist
\item
  Evaluate the concept of ``ecosystem services'' provided by soil
  microbiomes: nitrogen cycling, carbon sequestration, plant pathogen
  suppression. How does intensive agriculture (tillage, pesticide,
  monoculture) reduce soil microbial diversity and function? Propose a
  restorative soil management strategy that would re-establish lost
  ecosystem services.
\end{enumerate}

\begin{enumerate}
\def\labelenumi{\arabic{enumi}.}
\setcounter{enumi}{25}
\tightlist
\item
  The human virome includes about 6 \texorpdfstring{\ensuremath{\times}}{x} 10\texorpdfstring{\textsuperscript{1}}{1}\texorpdfstring{\textsuperscript{2}}{2} phage particles,
  largely matching the bacteriome. Evaluate: (a) how the virome
  influences bacteriome composition (phage predation driving community
  dynamics); (b) the evidence that an abnormal virome (e.g., expanded
  Caudovirales in Crohn's disease) can drive inflammation independently
  of the bacteriome; (c) whether the virome should be considered in the
  design of microbiome therapeutics.
\end{enumerate}

\begin{enumerate}
\def\labelenumi{\arabic{enumi}.}
\setcounter{enumi}{26}
\tightlist
\item
  Syntrophic communities (where no single organism can grow on a
  substrate alone --- requires metabolic cross-feeding with a partner)
  are central to anaerobic digestion of organic matter. Using the
  sulfate-reducing and methanogenic guild competition as an example,
  evaluate how thermodynamic constraints (Gibbs energy for each
  reaction) determine community composition, and why hydrogen partial
  pressure is the key variable.
\end{enumerate}

\begin{enumerate}
\def\labelenumi{\arabic{enumi}.}
\setcounter{enumi}{27}
\tightlist
\item
  Evaluate the potential of engineered synthetic microbial consortia for
  biotechnology (e.g., consortium for complete cellulose-to-ethanol
  conversion, or coordinated tumour-targeting drug delivery). What
  design principles govern consortium stability, and how can synthetic
  quorum sensing circuits be used to coordinate the activities of
  different strains?
\end{enumerate}

\begin{enumerate}
\def\labelenumi{\arabic{enumi}.}
\setcounter{enumi}{28}
\tightlist
\item
  The ocean microbiome is structured by ``niche partitioning'' on
  individual dissolved organic carbon (DOC) molecules: the ``competitive
  exclusion/keystone species'' model predicts that competitive exclusion
  dominates in low-diversity environments, while microbes coexist in
  complex high-DOC environments by partitioning onto different
  compounds. Evaluate how ``genome streamlining'' (loss of non-essential
  genes) in marine bacteria (e.g., SAR11) reflects adaptation to this
  competitive, nutrient-limited environment.
\end{enumerate}

\begin{enumerate}
\def\labelenumi{\arabic{enumi}.}
\setcounter{enumi}{29}
\tightlist
\item
  Critically evaluate the therapeutic potential of ``designer
  probiotics'' --- genetically engineered bacteria with specific
  metabolic or signalling functions. Using the example of a genetically
  engineered \emph{E. coli} Nissle expressing a GLP-1 precursor for
  treating obesity (clinical trial, 2025), evaluate: the safety concerns
  (containment of GMO bacteria, gene transfer to resident microbiome),
  the evidence of efficacy, and the regulatory framework needed.
\end{enumerate}

\newpage

\section*{Questions --- Ch 24: Infectious Disease and
Immunity}\label{questions-ch-24-infectious-disease-and-immunity}
\addcontentsline{toc}{section}{Questions --- Ch 24: Infectious Disease
and Immunity}

\label{sec:q_unit_VII_infectious_disease}

\subsection*{Instructor Use and Coverage
Notes}\label{instructor-use-and-coverage-notes-26}
\addcontentsline{toc}{subsection}{Instructor Use and Coverage Notes}

\begin{itemize}
\tightlist
\item
  \textbf{Coverage target:} Calculate a threshold or outbreak trajectory
  and identify the intervention lever.
\item
  \textbf{Model/data emphasis:} R0, herd-immunity threshold, and SIR
  trajectory calculations.
\item
  \textbf{Assessment alignment:} Questions and Methods, Representing and
  Describing Data, Argumentation.
\item
  \textbf{Misconception probe:} R0 is not a fixed property of a pathogen
  alone; it depends on host behaviour and environment.
\item
  \textbf{Transfer product:} Transfer disease-dynamic reasoning to
  vaccination, antimicrobial resistance, and public-health policy.
\item
  \textbf{Grading focus:} award full credit for mechanism, evidence,
  boundary conditions, and units when a calculation is required.
\item
  \textbf{Suggested use:} draw one recall item, one application item,
  and one synthesis item when building a short quiz from this bank.
\end{itemize}

\subsection*{Questions 1--10: Recall and
Comprehension}\label{questions-110-recall-and-comprehension-23}
\addcontentsline{toc}{subsection}{Questions 1--10: Recall and
Comprehension}

\emph{This activity accompanies \cref{sec:unit_VII_infectious_disease}
of the textbook --- review that chapter before attempting the exercises
below.}

\begin{enumerate}
\def\labelenumi{\arabic{enumi}.}
\tightlist
\item
  Distinguish innate from adaptive immunity. Name two cell types
  involved in each.
\end{enumerate}

\begin{enumerate}
\def\labelenumi{\arabic{enumi}.}
\setcounter{enumi}{1}
\tightlist
\item
  What is a pathogen-associated molecular pattern (PAMP)? Name one
  example and its receptor.
\end{enumerate}

\begin{enumerate}
\def\labelenumi{\arabic{enumi}.}
\setcounter{enumi}{2}
\tightlist
\item
  What is an antibody? Name its four structural components.
\end{enumerate}

\begin{enumerate}
\def\labelenumi{\arabic{enumi}.}
\setcounter{enumi}{3}
\tightlist
\item
  What are the five classes of immunoglobulin? State one function of IgG
  and one of IgA.
\end{enumerate}

\begin{enumerate}
\def\labelenumi{\arabic{enumi}.}
\setcounter{enumi}{4}
\tightlist
\item
  Describe the clonal selection theory of lymphocyte activation.
\end{enumerate}

\begin{enumerate}
\def\labelenumi{\arabic{enumi}.}
\setcounter{enumi}{5}
\tightlist
\item
  What is the MHC (major histocompatibility complex)? Distinguish MHC
  class I from class II.
\end{enumerate}

\begin{enumerate}
\def\labelenumi{\arabic{enumi}.}
\setcounter{enumi}{6}
\tightlist
\item
  What are cytotoxic T lymphocytes (CTLs)? What do they recognise and
  how do they kill target cells?
\end{enumerate}

\begin{enumerate}
\def\labelenumi{\arabic{enumi}.}
\setcounter{enumi}{7}
\tightlist
\item
  What is an epidemic? What is a pandemic? Define basic reproduction
  number R\texorpdfstring{\textsubscript{0}}{0}.
\end{enumerate}

\begin{enumerate}
\def\labelenumi{\arabic{enumi}.}
\setcounter{enumi}{8}
\tightlist
\item
  What is herd immunity? How does vaccination achieve herd immunity?
\end{enumerate}

\begin{enumerate}
\def\labelenumi{\arabic{enumi}.}
\setcounter{enumi}{9}
\tightlist
\item
  Name three mechanisms of antibiotic resistance evolution.
\end{enumerate}

\subsection*{Questions 11--20: Application and
Analysis}\label{questions-1120-application-and-analysis-23}
\addcontentsline{toc}{subsection}{Questions 11--20: Application and
Analysis}

\begin{enumerate}
\def\labelenumi{\arabic{enumi}.}
\setcounter{enumi}{10}
\tightlist
\item
  SARS-CoV-2 infected about 700 million people globally.
  Estimate the herd immunity threshold (H = 1 \texorpdfstring{\ensuremath{-}}{-} 1/R\texorpdfstring{\textsubscript{0}}{0}) if R\texorpdfstring{\textsubscript{0}}{0} for the
  original strain was 2.5. For the Omicron variant (R\texorpdfstring{\textsubscript{0}}{0} \texorpdfstring{\ensuremath{\approx}}{~} 10),
  recalculate the herd immunity threshold and explain the implications
  for vaccination coverage needed.
\end{enumerate}

\begin{enumerate}
\def\labelenumi{\arabic{enumi}.}
\setcounter{enumi}{11}
\tightlist
\item
  The inflammasome (NLRP3) is activated by uric acid crystals (gout),
  cholesterol crystals (atherosclerosis), and pathogens. Explain: (a)
  how NLRP3 assembles and activates caspase-1; (b) how caspase-1 cleaves
  pro-IL-1\texorpdfstring{\ensuremath{\beta}}{beta} and pro-IL-18; (c) why NLRP3 inhibitors (MCC950) are under
  development for inflammatory diseases.
\end{enumerate}

\begin{enumerate}
\def\labelenumi{\arabic{enumi}.}
\setcounter{enumi}{12}
\tightlist
\item
  Influenza can shift from H1N1 to H3N2 by antigenic shift (segment
  reassortment) or antigenic drift (point mutations in haemagglutinin).
  Explain: (a) why antigenic drift requires annual vaccine
  reformulation; (b) why antigenic shift can cause pandemics; (c) what
  ``original antigenic sin'' means for immune memory.
\end{enumerate}

\begin{enumerate}
\def\labelenumi{\arabic{enumi}.}
\setcounter{enumi}{13}
\tightlist
\item
  The complement system (C1-C9) is a cascade of proteolytic activation
  that opsonises pathogens and lyses bacteria directly (MAC --- membrane
  attack complex). Describe the three complement activation pathways
  (classical, lectin, alternative), naming the central converging point
  (C3 convertase), and explain how bacteria evade complement (factor H
  mimicry, capsule).
\end{enumerate}

\begin{enumerate}
\def\labelenumi{\arabic{enumi}.}
\setcounter{enumi}{14}
\tightlist
\item
  Natural killer (NK) cells kill virus-infected cells that downregulate
  MHC class I (a strategy used by many viruses to evade CTLs). Explain
  the ``missing self'' hypothesis: how do inhibitory KIR receptors
  normally inhibit NK killing, and how does MHC class I downregulation
  remove inhibition and trigger NK-mediated lysis?
\end{enumerate}

\begin{enumerate}
\def\labelenumi{\arabic{enumi}.}
\setcounter{enumi}{15}
\tightlist
\item
  mRNA vaccines (Pfizer-BioNTech, Moderna) for SARS-CoV-2 required
  solving two key problems: (a) mRNA instability in vivo (solution:
  pseudouridine substitution reduces TLR7/8 sensing and increases
  stability); (b) delivery to antigen-presenting cells in lymph nodes
  (solution: LNP --- lipid nanoparticle encapsulation). Trace the
  cascade from IM injection \texorpdfstring{\ensuremath{\to}}{->} LNP uptake by dendritic cells \texorpdfstring{\ensuremath{\to}}{->} spike
  protein expression \texorpdfstring{\ensuremath{\to}}{->} naive B and T cell activation \texorpdfstring{\ensuremath{\to}}{->} class switch
  recombination \texorpdfstring{\ensuremath{\to}}{->} IgG antibody production.
\end{enumerate}

\begin{enumerate}
\def\labelenumi{\arabic{enumi}.}
\setcounter{enumi}{16}
\tightlist
\item
  Regulatory T cells (Tregs, FOXP3\texorpdfstring{\textsuperscript{+}}{+}) suppress immune responses and
  prevent autoimmunity. Explain: (a) how Tregs develop (thymic selection
  vs peripheral induction); (b) their mechanisms of suppression (IL-10,
  TGF-\texorpdfstring{\ensuremath{\beta}}{beta}, CTLA-4-mediated antigen-presenting cell suppression); (c) why
  depleting Tregs in cancer (anti-CTLA-4, ipilimumab) can restore
  anti-tumour immunity at the cost of autoimmunity.
\end{enumerate}

\begin{enumerate}
\def\labelenumi{\arabic{enumi}.}
\setcounter{enumi}{17}
\tightlist
\item
  Malaria (\emph{Plasmodium falciparum}) has complex immune evasion
  through antigenic variation of PfEMP1 (erythrocyte membrane protein 1)
  encoded by about 60 \emph{var} genes, a single active at a
  time. Explain: (a) why the immune system struggles to clear malaria;
  (b) how epigenetic switching between var genes works; (c) what
  approach (conserved antigen targets) vaccine developers use to
  circumvent this evasion.
\end{enumerate}

\begin{enumerate}
\def\labelenumi{\arabic{enumi}.}
\setcounter{enumi}{18}
\tightlist
\item
  Calculate basic epidemiological parameters: if an influenza outbreak
  shows 15 secondary cases per 10 infected individuals over a mean
  infectious period of 3 days (generation time 3 days), estimate R\texorpdfstring{\textsubscript{0}}{0} and
  the doubling time of the epidemic. What control measure would need to
  be implemented to bring R to \textless{} 1?
\end{enumerate}

\begin{enumerate}
\def\labelenumi{\arabic{enumi}.}
\setcounter{enumi}{19}
\tightlist
\item
  Intestinal IgA is produced in the lamina propria by plasma cells and
  specifically transported into the gut lumen (via polymeric
  immunoglobulin receptor, pIgR). Explain: (a) how secretory IgA (sIgA)
  differs from serum IgA structurally; (b) the role of the secretory
  component; (c) how sIgA provides ``immune exclusion'' of pathogens
  without triggering inflammation.
\end{enumerate}

\subsection*{Questions 21--30: Synthesis and
Evaluation}\label{questions-2130-synthesis-and-evaluation-23}
\addcontentsline{toc}{subsection}{Questions 21--30: Synthesis and
Evaluation}

\begin{enumerate}
\def\labelenumi{\arabic{enumi}.}
\setcounter{enumi}{20}
\tightlist
\item
  Evaluate the ``inverted-U'' relationship between immune activation and
  pathogen clearance, using sepsis as a model: (a) why severe sepsis
  kills by cytokine storm (uncontrolled innate immune activation) rather
  than the pathogen itself; (b) the failure of anti-TNF-\texorpdfstring{\ensuremath{\alpha}}{alpha} therapy in
  sepsis trials despite success in rheumatoid arthritis; (c) how the
  ``immunological hysteresis'' concept explains why the same cytokine
  blockade is beneficial in autoimmunity but harmful in sepsis.
\end{enumerate}

\begin{enumerate}
\def\labelenumi{\arabic{enumi}.}
\setcounter{enumi}{21}
\tightlist
\item
  Build a paper-based infectious-disease decision matrix for three
  current interventions: (a) BPaL/BPaLM-style most-oral treatment for
  rifampicin-resistant tuberculosis; (b) twice-yearly lenacapavir PrEP
  for people with substantial HIV exposure risk; (c) indoor spatial
  emanators added to insecticide-treated nets for malaria control. For
  each, state the biological target, the surveillance or eligibility
  evidence required, one equity constraint, and one reason the
  intervention should not be treated as a standalone solution.
\end{enumerate}

\begin{enumerate}
\def\labelenumi{\arabic{enumi}.}
\setcounter{enumi}{22}
\tightlist
\item
  Evaluate the evidence for ``trained innate immunity'': after exposure
  to \texorpdfstring{\ensuremath{\beta}}{beta}-glucan (from \emph{Candida}) or BCG vaccination, innate immune
  cells (macrophages, NK cells) show epigenetically reprogrammed,
  enhanced responses 3--6 months later. What epigenetic changes (H3K4me1
  at enhancers, mTOR-dependent metabolic rewiring) mediate this
  ``immunological memory'' in non-adaptive immune cells, and what are
  the implications for innate immune-based vaccines (BCG against
  COVID-19)?
\end{enumerate}

\begin{enumerate}
\def\labelenumi{\arabic{enumi}.}
\setcounter{enumi}{23}
\tightlist
\item
  Autoimmune diseases arise when self-tolerance breaks down. Evaluate
  three mechanisms of self-tolerance failure in Type 1 Diabetes (T1D):
  (a) central tolerance failure (inadequate thymic deletion of
  autoreactive T cells directed at insulin); (b) peripheral tolerance
  failure (Treg dysfunction); (c) molecular mimicry (GAD65-Coxsackie B
  virus homology). What evidence supports each mechanism in human T1D?
\end{enumerate}

\begin{enumerate}
\def\labelenumi{\arabic{enumi}.}
\setcounter{enumi}{24}
\tightlist
\item
  During COVID-19 pandemic, several SARS-CoV-2 variants of concern
  (Alpha, Delta, Omicron) arose sequentially. Evaluate the evolutionary
  pressures (immune evasion vs transmissibility) that drove variant
  emergence, and why Omicron showed more ACE2-binding mutations
  (increased transmissibility) but more spike mutations outside the RBD
  (immune evasion from vaccine-induced antibodies).
\end{enumerate}

\begin{enumerate}
\def\labelenumi{\arabic{enumi}.}
\setcounter{enumi}{25}
\tightlist
\item
  Evaluate checkpoint immunotherapy for cancer: PD-1/PD-L1 blockade,
  CTLA-4 blockade, LAG-3 blockade. For what biomarker(s) (TMB, PD-L1
  expression, MSI-H) are these drugs most effective, and why do about
  20--30\% of eligible cancer patients benefit? Propose a combination
  therapy strategy to increase response rates while managing autoimmune
  toxicity (immune-related adverse events, irAEs).
\end{enumerate}

\begin{enumerate}
\def\labelenumi{\arabic{enumi}.}
\setcounter{enumi}{26}
\tightlist
\item
  Evaluate the concept of the ``immunological homunculus'' --- that the
  normal immune system retains low-frequency self-reactive clones that
  perform surveillance and immunoregulation (not only pathology).
  Critically assess the evidence that low-affinity self-reactive T and B
  cells are essential for shaping the immune repertoire, and what
  happens when these regulatory self-reactive cells are eliminated by
  lymphodepletion.
\end{enumerate}

\begin{enumerate}
\def\labelenumi{\arabic{enumi}.}
\setcounter{enumi}{27}
\tightlist
\item
  The ``two-signal model'' of T cell activation requires: (1) TCR
  engagement with MHC-peptide and (2) a co-stimulatory signal (CD28-B7).
  Without signal 2, T cells become anergic (unresponsive). Evaluate how
  tumours exploit this model to suppress anti-tumour T cells by: (a)
  expressing PD-L1 (sending an anergy-inducing signal); (b) lacking
  co-stimulatory molecules; (c) expressing IDO (depleting tryptophan
  needed for T cell activation). How does CAR-T cell therapy bypass this
  requirement for co-stimulation?
\end{enumerate}

\begin{enumerate}
\def\labelenumi{\arabic{enumi}.}
\setcounter{enumi}{28}
\tightlist
\item
  Evaluate the emerging field of ``onco-immunology'' --- the mechanisms
  by which tumours edit the immune landscape (immunoediting: elimination
  \texorpdfstring{\ensuremath{\to}}{->} equilibrium \texorpdfstring{\ensuremath{\to}}{->} escape). Give specific molecular examples of immune
  editing at each phase, and assess whether the immune equilibrium phase
  can be maintained indefinitely in any cancer type, or whether escape
  is inevitable.
\end{enumerate}

\begin{enumerate}
\def\labelenumi{\arabic{enumi}.}
\setcounter{enumi}{29}
\tightlist
\item
  Critically evaluate the public health ethics of mandatory vaccination
  programmes. Using the framework of herd immunity thresholds,
  individual liberty, social contract theory, and harm principle (Mill),
  construct a utilitarian argument for and a libertarian argument
  against mandatory vaccination. What vaccination coverage is minimally
  required for herd immunity against measles (R\texorpdfstring{\textsubscript{0}}{0} \texorpdfstring{\ensuremath{\approx}}{~} 18), and how do
  vaccine hesitancy and anti-vaccine movements affect public health?
\end{enumerate}

\newpage

\section*{Questions --- Ch 25: Plant Structure, Water Relations, and
Transport}\label{questions-ch-25-plant-structure-water-relations-and-transport}
\addcontentsline{toc}{section}{Questions --- Ch 25: Plant Structure,
Water Relations, and Transport}

\label{sec:q_unit_VIII_plant_structure_and_water}

\subsection*{Instructor Use and Coverage
Notes}\label{instructor-use-and-coverage-notes-27}
\addcontentsline{toc}{subsection}{Instructor Use and Coverage Notes}

\begin{itemize}
\tightlist
\item
  \textbf{Coverage target:} Calculate water potential and predict the
  direction of water movement.
\item
  \textbf{Model/data emphasis:} Water-potential and transpiration-flux
  calculations.
\item
  \textbf{Assessment alignment:} Visual Representations, Questions and
  Methods, Argumentation.
\item
  \textbf{Misconception probe:} Water does not move because plants pull
  with intention; it follows potential gradients and cohesion.
\item
  \textbf{Transfer product:} Transfer plant-water reasoning to drought,
  irrigation, forest physiology, and crop breeding.
\item
  \textbf{Grading focus:} award full credit for mechanism, evidence,
  boundary conditions, and units when a calculation is required.
\item
  \textbf{Suggested use:} draw one recall item, one application item,
  and one synthesis item when building a short quiz from this bank.
\end{itemize}

\subsection*{Questions 1--10: Recall and
Comprehension}\label{questions-110-recall-and-comprehension-24}
\addcontentsline{toc}{subsection}{Questions 1--10: Recall and
Comprehension}

\emph{This activity accompanies
\cref{sec:unit_VIII_plant_structure_and_water} of the textbook ---
review that chapter before attempting the exercises below.}

\begin{enumerate}
\def\labelenumi{\arabic{enumi}.}
\tightlist
\item
  Name the three major tissue systems in plants. Give one cell type
  belonging to each.
\end{enumerate}

\begin{enumerate}
\def\labelenumi{\arabic{enumi}.}
\setcounter{enumi}{1}
\tightlist
\item
  What is the difference between dermal, ground, and vascular tissue?
\end{enumerate}

\begin{enumerate}
\def\labelenumi{\arabic{enumi}.}
\setcounter{enumi}{2}
\tightlist
\item
  Define xylem and phloem. Which transports water and which transports
  sugars?
\end{enumerate}

\begin{enumerate}
\def\labelenumi{\arabic{enumi}.}
\setcounter{enumi}{3}
\tightlist
\item
  What are tracheids and vessel elements? How do they differ
  structurally?
\end{enumerate}

\begin{enumerate}
\def\labelenumi{\arabic{enumi}.}
\setcounter{enumi}{4}
\tightlist
\item
  Define water potential (\texorpdfstring{\ensuremath{\Psi}}{Psi}). Write the equation \texorpdfstring{\ensuremath{\Psi}}{Psi} = \texorpdfstring{\ensuremath{\Psi}}{Psi}\_s + \texorpdfstring{\ensuremath{\Psi}}{Psi}\_p.
\end{enumerate}

\begin{enumerate}
\def\labelenumi{\arabic{enumi}.}
\setcounter{enumi}{5}
\tightlist
\item
  What is the Casparian strip? In which cell layer is it found and what
  does it do?
\end{enumerate}

\begin{enumerate}
\def\labelenumi{\arabic{enumi}.}
\setcounter{enumi}{6}
\tightlist
\item
  Define transpiration. What drives transpiration in leaves?
\end{enumerate}

\begin{enumerate}
\def\labelenumi{\arabic{enumi}.}
\setcounter{enumi}{7}
\tightlist
\item
  What are stomata? What cells regulate their opening and closing?
\end{enumerate}

\begin{enumerate}
\def\labelenumi{\arabic{enumi}.}
\setcounter{enumi}{8}
\tightlist
\item
  What is root pressure? Under what conditions does it occur?
\end{enumerate}

\begin{enumerate}
\def\labelenumi{\arabic{enumi}.}
\setcounter{enumi}{9}
\tightlist
\item
  Describe the pressure-flow (Münch) hypothesis for phloem transport.
\end{enumerate}

\subsection*{Questions 11--20: Application and
Analysis}\label{questions-1120-application-and-analysis-24}
\addcontentsline{toc}{subsection}{Questions 11--20: Application and
Analysis}

\begin{enumerate}
\def\labelenumi{\arabic{enumi}.}
\setcounter{enumi}{10}
\tightlist
\item
  A plant cell has \texorpdfstring{\ensuremath{\Psi}}{Psi}\_s = \texorpdfstring{\ensuremath{-}}{-}0.9 MPa and \texorpdfstring{\ensuremath{\Psi}}{Psi}\_p = +0.4 MPa. Calculate its \texorpdfstring{\ensuremath{\Psi}}{Psi}.
  If placed in a solution with \texorpdfstring{\ensuremath{\Psi}}{Psi} = \texorpdfstring{\ensuremath{-}}{-}0.3 MPa, will water enter or leave
  the cell? What will change as the cell reaches equilibrium?
\end{enumerate}

\begin{enumerate}
\def\labelenumi{\arabic{enumi}.}
\setcounter{enumi}{11}
\tightlist
\item
  The cohesion-tension theory explains xylem water transport. Explain
  each component: (a) cohesion (H-bonds between water molecules); (b)
  adhesion (attraction to xylem walls); (c) tension (negative pressure
  from leaf transpiration). How does this mechanism pull water to 100 m
  in tall trees?
\end{enumerate}

\begin{enumerate}
\def\labelenumi{\arabic{enumi}.}
\setcounter{enumi}{12}
\tightlist
\item
  A researcher applies ABA (abscisic acid) to stomatal guard cells.
  Predict: (a) what happens to K\texorpdfstring{\textsuperscript{+}}{+} channels (efflux activated); (b) what
  happens to guard cell turgor; (c) what happens to stomatal aperture;
  (d) how CO\texorpdfstring{\textsubscript{2}}{2} supply to mesophyll cells changes.
\end{enumerate}

\begin{enumerate}
\def\labelenumi{\arabic{enumi}.}
\setcounter{enumi}{13}
\tightlist
\item
  Compare apoplastic vs symplastic water transport from root epidermis
  to the xylem. What role does the Casparian strip play in forcing water
  through the symplastic pathway? Why is this gatekeeping important for
  ion selection?
\end{enumerate}

\begin{enumerate}
\def\labelenumi{\arabic{enumi}.}
\setcounter{enumi}{14}
\tightlist
\item
  Phloem loading concentrates sucrose in sieve tubes at the source.
  Explain: (a) the apoplastic loading pathway (H\texorpdfstring{\textsuperscript{+}}{+}-sucrose
  co-transporter); (b) how increased solute concentration in sieve tubes
  lowers \texorpdfstring{\ensuremath{\Psi}}{Psi}\_s; (c) how water follows by osmosis, generating turgor
  pressure that drives bulk flow toward the sink.
\end{enumerate}

\begin{enumerate}
\def\labelenumi{\arabic{enumi}.}
\setcounter{enumi}{15}
\tightlist
\item
  Cavitation (air bubble formation) in xylem vessels can block water
  transport. Explain: (a) how tension increases cavitation risk during
  drought; (b) how pit membranes between vessel elements limit
  cavitation spread; (c) how some trees (ring-porous oaks) repair
  cavitated vessels at night via root pressure.
\end{enumerate}

\begin{enumerate}
\def\labelenumi{\arabic{enumi}.}
\setcounter{enumi}{16}
\tightlist
\item
  A potometer measures transpiration rate by tracking water uptake. If a
  leafy shoot uptakes 4.5 mL/h at 25°C and 2.0 mL/h at 15°C, calculate:
  (a) the Q\texorpdfstring{\textsubscript{1}}{1}\texorpdfstring{\textsubscript{0}}{0} for transpiration; (b) what factor associated with
  temperature change drives this difference; (c) why humidity must also
  be controlled.
\end{enumerate}

\begin{enumerate}
\def\labelenumi{\arabic{enumi}.}
\setcounter{enumi}{17}
\tightlist
\item
  Mycorrhizal fungi dramatically increase root surface area and provide
  phosphorus to plants in exchange for photosynthetic carbon. Compare
  arbuscular mycorrhizae (AM, intracellular fungal hyphae) vs
  ectomycorrhizal (ECM, extracellular mantle). Estimate the increase in
  effective root absorptive area (10\texorpdfstring{\ensuremath{\times}}{x} to 1,000\texorpdfstring{\ensuremath{\times}}{x}) and explain why nearly
  90\% of plant species have mycorrhizal associations.
\end{enumerate}

\begin{enumerate}
\def\labelenumi{\arabic{enumi}.}
\setcounter{enumi}{18}
\tightlist
\item
  Halophytes (salt-tolerant plants) grow in saline soils (\texorpdfstring{\ensuremath{\Psi}}{Psi}\_s
  \textless{} \texorpdfstring{\ensuremath{-}}{-}1.0 MPa). Explain three strategies: (a) accumulation of
  compatible solutes (proline) to lower \texorpdfstring{\ensuremath{\Psi}}{Psi}\_s; (b) ion
  compartmentalisation into vacuoles (NHX1 Na\texorpdfstring{\textsuperscript{+}}{+}/H\texorpdfstring{\textsuperscript{+}}{+} antiporter); (c) salt
  secretion via specialised salt glands.
\end{enumerate}

\begin{enumerate}
\def\labelenumi{\arabic{enumi}.}
\setcounter{enumi}{19}
\tightlist
\item
  A plant anatomist sections a dicot stem and a monocot stem. List four
  differences in vascular bundle arrangement, cambium presence,
  cortex/pith proportions, and secondary growth capability.
\end{enumerate}

\subsection*{Questions 21--30: Synthesis and
Evaluation}\label{questions-2130-synthesis-and-evaluation-24}
\addcontentsline{toc}{subsection}{Questions 21--30: Synthesis and
Evaluation}

\begin{enumerate}
\def\labelenumi{\arabic{enumi}.}
\setcounter{enumi}{20}
\tightlist
\item
  Evaluate the limits of the cohesion-tension theory: the maximum
  theoretical xylem tension is approximately \texorpdfstring{\ensuremath{-}}{-}20 MPa, but measured
  tensions rarely exceed \texorpdfstring{\ensuremath{-}}{-}3 MPa. Does this discrepancy refute the
  theory? Discuss the ``vulnerability curve'' (loss of hydraulic
  conductivity vs tension) and how different species show different
  cavitation thresholds. What evolutionary strategies balance xylem
  safety vs efficiency?
\end{enumerate}

\begin{enumerate}
\def\labelenumi{\arabic{enumi}.}
\setcounter{enumi}{21}
\tightlist
\item
  Design an experiment to measure the hydraulic conductivity of a stem
  segment. Describe the apparatus (pressure head, flow rate
  measurement), the calculation (K\_h = flow rate \texorpdfstring{\ensuremath{\times}}{x} segment length /
  pressure difference), and how you would normalise for sapwood area.
  What would you expect for a conifer (tracheids primarily) vs a
  diffuse-porous angiosperm (vessel elements)?
\end{enumerate}

\begin{enumerate}
\def\labelenumi{\arabic{enumi}.}
\setcounter{enumi}{22}
\tightlist
\item
  Evaluate the evidence that aquaporins regulate short-term water
  transport in roots. How does the abundance and phosphorylation state
  of root aquaporins (PIP1, PIP2 subfamilies) change in response to: (a)
  drought; (b) flooding; (c) diurnal light cycles? What does this imply
  about the relative importance of the apoplastic vs symplastic pathway
  under different environmental conditions?
\end{enumerate}

\begin{enumerate}
\def\labelenumi{\arabic{enumi}.}
\setcounter{enumi}{23}
\tightlist
\item
  Propose a bioengineering strategy to improve drought tolerance in crop
  plants by manipulating water transport. Compare three approaches: (a)
  engineering deeper root systems (overexpressing DRO1); (b) increasing
  xylem vessel diameter for greater hydraulic conductivity; (c)
  constitutively closing stomata (overexpressing OST1 kinase). For each,
  evaluate the trade-off between water conservation and carbon
  assimilation. Which strategy is most promising for wheat in arid
  environments?
\end{enumerate}

\begin{enumerate}
\def\labelenumi{\arabic{enumi}.}
\setcounter{enumi}{24}
\tightlist
\item
  Phloem transport speed is about 1 m/h. Calculate: if a tree
  is 30 m tall and assimilate loading occurs in the canopy, how long
  does it take for sucrose to reach root tips? Compare this to xylem
  transport speed (about 10--40 m/h). What limits phloem
  transport speed (viscosity of concentrated sucrose solution, sieve
  plate resistance)?
\end{enumerate}

\begin{enumerate}
\def\labelenumi{\arabic{enumi}.}
\setcounter{enumi}{25}
\tightlist
\item
  Evaluate the evolutionary transition from tracheids (present in most
  vascular plants) to vessel elements (present in angiosperms and some
  gymnosperms). What anatomical changes (perforation plates, wider
  diameter) conferred increased hydraulic efficiency, and what
  trade-offs (increased vulnerability to cavitation) accompanied this
  transition? How does this relate to the ecological dominance of
  angiosperms?
\end{enumerate}

\begin{enumerate}
\def\labelenumi{\arabic{enumi}.}
\setcounter{enumi}{26}
\tightlist
\item
  The pressure-flow model (Münch) for phloem transport makes predictions
  about: (a) turgor in source vs sink phloem; (b) the osmotic potential
  gradient between source and sink; (c) the effect of girdling (removing
  phloem) on sugar accumulation above the girdle. Evaluate each
  prediction using experimental evidence, and assess the main criticism
  of the model (can passive flow alone explain transport over
  \textgreater50 m?).
\end{enumerate}

\begin{enumerate}
\def\labelenumi{\arabic{enumi}.}
\setcounter{enumi}{27}
\tightlist
\item
  Evaluate the role of plasmodesmata in cell-to-cell communication in
  plants. How do plasmodesmata allow selective symplastic transport
  (size exclusion limit about 1 kDa), what signals move through
  them (transcription factors like KNOTTED1, siRNAs), and how do viruses
  exploit plasmodesmata for cell-to-cell spread (movement proteins)?
\end{enumerate}

\begin{enumerate}
\def\labelenumi{\arabic{enumi}.}
\setcounter{enumi}{28}
\tightlist
\item
  Climate change increases atmospheric CO\texorpdfstring{\textsubscript{2}}{2} and temperature, both of
  which affect transpiration, stomatal conductance, and water-use
  efficiency. Evaluate the predicted effects on forest water balance:
  (a) increased CO\texorpdfstring{\textsubscript{2}}{2} \texorpdfstring{\ensuremath{\to}}{->} partial stomatal closure \texorpdfstring{\ensuremath{\to}}{->} reduced transpiration \texorpdfstring{\ensuremath{\to}}{->}
  increased runoff; (b) increased temperature \texorpdfstring{\ensuremath{\to}}{->} increased VPD \texorpdfstring{\ensuremath{\to}}{->}
  increased transpiration \texorpdfstring{\ensuremath{\to}}{->} drought stress. Which effect dominates in
  tropical forests vs boreal forests?
\end{enumerate}

\begin{enumerate}
\def\labelenumi{\arabic{enumi}.}
\setcounter{enumi}{29}
\tightlist
\item
  Critically evaluate whether the ``hydraulic architecture'' of a tree
  (the specific pattern of vessel diameters, vessel density, and branch
  tapering from base to tip) is optimised by natural selection for
  efficient water delivery. Apply Murray's law (branching geometry) and
  the ``safety vs efficiency'' trade-off framework. What does damage
  from extreme drought events reveal about the limits of optimisation?
\end{enumerate}

\newpage

\section*{Questions --- Ch 26: Plant Reproduction and
Development}\label{questions-ch-26-plant-reproduction-and-development}
\addcontentsline{toc}{section}{Questions --- Ch 26: Plant Reproduction
and Development}

\label{sec:q_unit_VIII_plant_reproduction}

\subsection*{Instructor Use and Coverage
Notes}\label{instructor-use-and-coverage-notes-28}
\addcontentsline{toc}{subsection}{Instructor Use and Coverage Notes}

\begin{itemize}
\tightlist
\item
  \textbf{Coverage target:} Trace a plant life cycle and predict the
  consequence of disrupting one stage.
\item
  \textbf{Model/data emphasis:} Life-cycle accounting and
  phyllotaxis/growth-pattern calculations.
\item
  \textbf{Assessment alignment:} Visual Representations, Questions and
  Methods, Argumentation.
\item
  \textbf{Misconception probe:} Pollen, spores, seeds, and gametes are
  not interchangeable terms.
\item
  \textbf{Transfer product:} Transfer reproductive reasoning to
  agriculture, pollination ecology, and plant evolution.
\item
  \textbf{Grading focus:} award full credit for mechanism, evidence,
  boundary conditions, and units when a calculation is required.
\item
  \textbf{Suggested use:} draw one recall item, one application item,
  and one synthesis item when building a short quiz from this bank.
\end{itemize}

\subsection*{Questions 1--10: Recall and
Comprehension}\label{questions-110-recall-and-comprehension-25}
\addcontentsline{toc}{subsection}{Questions 1--10: Recall and
Comprehension}

\emph{This activity accompanies \cref{sec:unit_VIII_plant_reproduction}
of the textbook --- review that chapter before attempting the exercises
below.}

\begin{enumerate}
\def\labelenumi{\arabic{enumi}.}
\tightlist
\item
  Describe the alternation of generations in plant life cycles. What is
  the gametophyte? What is the sporophyte?
\end{enumerate}

\begin{enumerate}
\def\labelenumi{\arabic{enumi}.}
\setcounter{enumi}{1}
\tightlist
\item
  Name the male and female parts of a flower and state their functions.
\end{enumerate}

\begin{enumerate}
\def\labelenumi{\arabic{enumi}.}
\setcounter{enumi}{2}
\tightlist
\item
  What is pollination? Distinguish self-pollination from
  cross-pollination.
\end{enumerate}

\begin{enumerate}
\def\labelenumi{\arabic{enumi}.}
\setcounter{enumi}{3}
\tightlist
\item
  What is double fertilisation in angiosperms? What two structures
  result?
\end{enumerate}

\begin{enumerate}
\def\labelenumi{\arabic{enumi}.}
\setcounter{enumi}{4}
\tightlist
\item
  What is a seed? Name its three components.
\end{enumerate}

\begin{enumerate}
\def\labelenumi{\arabic{enumi}.}
\setcounter{enumi}{5}
\tightlist
\item
  What is a fruit? How does it develop after fertilisation?
\end{enumerate}

\begin{enumerate}
\def\labelenumi{\arabic{enumi}.}
\setcounter{enumi}{6}
\tightlist
\item
  Define germination. Name three environmental factors required.
\end{enumerate}

\begin{enumerate}
\def\labelenumi{\arabic{enumi}.}
\setcounter{enumi}{7}
\tightlist
\item
  What is vegetative reproduction? Give two examples.
\end{enumerate}

\begin{enumerate}
\def\labelenumi{\arabic{enumi}.}
\setcounter{enumi}{8}
\tightlist
\item
  What is meristem tissue? Distinguish apical meristems from lateral
  meristems.
\end{enumerate}

\begin{enumerate}
\def\labelenumi{\arabic{enumi}.}
\setcounter{enumi}{9}
\tightlist
\item
  What is photoperiodism? Distinguish long-day, short-day, and
  day-neutral plants.
\end{enumerate}

\subsection*{Questions 11--20: Application and
Analysis}\label{questions-1120-application-and-analysis-25}
\addcontentsline{toc}{subsection}{Questions 11--20: Application and
Analysis}

\begin{enumerate}
\def\labelenumi{\arabic{enumi}.}
\setcounter{enumi}{10}
\tightlist
\item
  Double fertilisation: one sperm fuses with the egg (2n zygote),
  another fuses with the central cell (3n endosperm). Explain the
  evolutionary advantage of triploid endosperm vs a haploid nutritive
  tissue. Why is endosperm essential for seed development in cereals?
\end{enumerate}

\begin{enumerate}
\def\labelenumi{\arabic{enumi}.}
\setcounter{enumi}{11}
\tightlist
\item
  SI (self-incompatibility) mechanisms prevent self-pollination. In
  gametophytic SI (e.g., \emph{Nicotiana}), pollen S-RNases degrade RNA
  in incompatible (self) pollen tubes. Explain: (a) the molecular logic
  --- why would a plant destroy its own pollen; (b) the advantage for
  genetic diversity; (c) how SI is overcome in commercial apple orchards
  (polliniser trees).
\end{enumerate}

\begin{enumerate}
\def\labelenumi{\arabic{enumi}.}
\setcounter{enumi}{12}
\tightlist
\item
  Seed dormancy allows seeds to delay germination until conditions are
  favourable. Explain: (a) the role of ABA (prevents germination) vs
  gibberellin (promotes germination); (b) why stratification (cold
  treatment) breaks dormancy in temperate-region seeds; (c) the
  ecological advantage of dormancy for fire-adapted species
  (heat-triggered germination).
\end{enumerate}

\begin{enumerate}
\def\labelenumi{\arabic{enumi}.}
\setcounter{enumi}{13}
\tightlist
\item
  Fruit types: distinguish fleshy fruits (berry, drupe, pome) from dry
  fruits (achene, capsule, legume). For each type, name an example
  species and explain how the fruit structure promotes seed dispersal
  (animal, wind, water, or ballistic).
\end{enumerate}

\begin{enumerate}
\def\labelenumi{\arabic{enumi}.}
\setcounter{enumi}{14}
\tightlist
\item
  Arabidopsis ABCE model of floral organ identity: A genes \texorpdfstring{\ensuremath{\to}}{->} sepals; A+B
  \texorpdfstring{\ensuremath{\to}}{->} petals; B+C \texorpdfstring{\ensuremath{\to}}{->} stamens; C \texorpdfstring{\ensuremath{\to}}{->} carpels; E is required in most whorls.
  Predict the floral phenotype when: (a) the B gene is knocked out; (b)
  the A gene is knocked out; (c) both A and C are knocked out. What
  would each mutant flower look like?
\end{enumerate}

\begin{enumerate}
\def\labelenumi{\arabic{enumi}.}
\setcounter{enumi}{15}
\tightlist
\item
  Pollen tube navigation: pollen tubes grow through the style following
  chemical gradients (e.g., LURE peptides from the synergid cells
  flanking the egg cell). Explain: (a) how Ca\texorpdfstring{\textsuperscript{2}}{2}\texorpdfstring{\textsuperscript{+}}{+} oscillations drive
  directional tip growth; (b) the role of the filiform apparatus in
  synergid cells; (c) what happens if a pollen tube arrives at an
  already-fertilised ovule (polytubey block).
\end{enumerate}

\begin{enumerate}
\def\labelenumi{\arabic{enumi}.}
\setcounter{enumi}{16}
\tightlist
\item
  Apomixis (asexual seed production without fertilisation) occurs
  naturally in some grasses and dandelions. Explain: (a) the mechanism
  (unreduced embryo sac, parthenogenesis); (b) why agriculture would
  benefit from engineered apomixis in hybrid crops (fixing heterosis
  permanently); (c) the rice synthetic-apomixis approach using MiMe plus
  a parthenogenesis trigger; (d) what field-translation risks remain.
\end{enumerate}

\begin{enumerate}
\def\labelenumi{\arabic{enumi}.}
\setcounter{enumi}{17}
\tightlist
\item
  Compare seed dispersal strategies: epizoochory (burs attaching to
  animal fur), endozoochory (fleshy fruits eaten by birds), anemochory
  (wind, e.g., dandelion pappus), hydrochory (water, coconut). For each,
  predict which seed/fruit trait is maximised (adhesion strength,
  palatability, surface area, buoyancy).
\end{enumerate}

\begin{enumerate}
\def\labelenumi{\arabic{enumi}.}
\setcounter{enumi}{18}
\tightlist
\item
  The transition from vegetative to reproductive growth (floral
  transition) is controlled by the florigen protein FT (FLOWERING LOCUS
  T) produced in leaves and transported via phloem to the shoot apical
  meristem. Explain: (a) how photoperiod (sensed by CONSTANS protein in
  leaves) triggers FT expression; (b) how FT interacts with FD
  transcription factor at the meristem to activate floral identity genes
  (AP1); (c) why vernalisation (prolonged cold) is required for winter
  annuals.
\end{enumerate}

\begin{enumerate}
\def\labelenumi{\arabic{enumi}.}
\setcounter{enumi}{19}
\tightlist
\item
  Embryogenesis in plants establishes the apical-basal axis (shoot-root)
  and radial axis (epidermis-vascular tissue-ground tissue). Using
  Arabidopsis as a model, describe: (a) the role of the asymmetric first
  division; (b) the signalling between the suspensor and the embryo
  proper; (c) the auxin gradient that specifies the root pole
  (PIN-mediated polar transport).
\end{enumerate}

\subsection*{Questions 21--30: Synthesis and
Evaluation}\label{questions-2130-synthesis-and-evaluation-25}
\addcontentsline{toc}{subsection}{Questions 21--30: Synthesis and
Evaluation}

\begin{enumerate}
\def\labelenumi{\arabic{enumi}.}
\setcounter{enumi}{20}
\tightlist
\item
  Coevolution between flowering plants and their pollinators has shaped
  morphological diversity. Evaluate the evidence for pollination
  syndromes (e.g., red tubular flowers \texorpdfstring{\ensuremath{\to}}{->} hummingbird; white fragrant
  nocturnal \texorpdfstring{\ensuremath{\to}}{->} moth; UV-patterned \texorpdfstring{\ensuremath{\to}}{->} bee). Are syndromes perfect
  predictors? What does the ``most effective pollinator'' principle say
  about syndrome fidelity vs pollinator generalisation?
\end{enumerate}

\begin{enumerate}
\def\labelenumi{\arabic{enumi}.}
\setcounter{enumi}{21}
\tightlist
\item
  Design a breeding programme to introduce self-incompatibility from a
  wild relative into a self-compatible crop species (e.g., tomato).
  Describe the genetic crosses required, how you would screen for SI
  expression, and the trade-off between enforced outcrossing (greater
  heterosis) and practical challenges (pollinator dependence).
\end{enumerate}

\begin{enumerate}
\def\labelenumi{\arabic{enumi}.}
\setcounter{enumi}{22}
\tightlist
\item
  Evaluate the hypothesis that the origin of the flower was a single
  evolutionary event (monophyletic origin of angiosperms). Cite evidence
  from: (a) molecular phylogenies (most angiosperms share a common
  ancestor); (b) fossil data (earliest undisputed flower fossils
  about 132 Mya); (c) the ``evo-devo'' evidence that the
  MADS-box gene family was duplicated and diversified at the base of
  angiosperms. What is ``Darwin's abominable mystery''?
\end{enumerate}

\begin{enumerate}
\def\labelenumi{\arabic{enumi}.}
\setcounter{enumi}{23}
\tightlist
\item
  Parthenocarpy (seedless fruit production) can be induced by GA\texorpdfstring{\textsubscript{3}}{3}
  application (seedless grapes) or by gene engineering (auxin- or
  GA-overproduction in ovary). Evaluate: (a) why seedless fruits are
  commercially valuable; (b) what transcriptional changes GA\texorpdfstring{\textsubscript{3}}{3} induces in
  the ovary wall (cell expansion genes); (c) whether GM parthenocarpy
  would require approval as a GMO in the EU.
\end{enumerate}

\begin{enumerate}
\def\labelenumi{\arabic{enumi}.}
\setcounter{enumi}{24}
\tightlist
\item
  Evaluate the evolutionary transition from water-dependent reproduction
  (bryophytes --- motile sperm requiring a water film) to
  water-independent reproduction (angiosperms --- pollen grain, pollen
  tube, sealed ovule). What were the key innovations at each transition
  (vascular plants \texorpdfstring{\ensuremath{\to}}{->} seeds \texorpdfstring{\ensuremath{\to}}{->} flowers \texorpdfstring{\ensuremath{\to}}{->} enclosed ovules), and how did
  each innovation expand environmental range?
\end{enumerate}

\begin{enumerate}
\def\labelenumi{\arabic{enumi}.}
\setcounter{enumi}{25}
\tightlist
\item
  The seed coat (testa) provides physical and chemical protection
  against pathogens, UV, and desiccation. Evaluate the molecular basis
  of seed coat impermeability (suberin and lignin deposits) and how seed
  coat-imposed dormancy is released by scarification (physical damage)
  or microbial degradation. What trade-off exists between strong
  dormancy (survival) and rapid germination (competitive advantage)?
\end{enumerate}

\begin{enumerate}
\def\labelenumi{\arabic{enumi}.}
\setcounter{enumi}{26}
\tightlist
\item
  Epigenetic regulation of seed development: the MEDEA/FIE Polycomb
  complex prevents autonomous endosperm development before
  fertilisation. After fertilisation, paternal genome expression
  overcomes MEDEA repression (genomic imprinting). Evaluate: (a) why the
  paternal genome specifically is needed to initiate endosperm; (b) the
  parent-of-origin conflict hypothesis (Haig theory) applied to seed
  provisioning; (c) what happens in interploidy crosses (2x \texorpdfstring{\ensuremath{\times}}{x} 4x) where
  parental dosage is disrupted.
\end{enumerate}

\begin{enumerate}
\def\labelenumi{\arabic{enumi}.}
\setcounter{enumi}{27}
\tightlist
\item
  Evaluate the potential of synthetic seeds (somatic embryos
  encapsulated in alginate beads) for clonal propagation of elite crop
  genotypes. What are the technical challenges (desiccation tolerance of
  somatic embryos, germination rate, production at scale), and how do
  synthetic seeds compare to traditional seed in terms of cost, genetic
  uniformity, and epigenetic stability?
\end{enumerate}

\begin{enumerate}
\def\labelenumi{\arabic{enumi}.}
\setcounter{enumi}{28}
\tightlist
\item
  Climate change alters flowering phenology (timing of first bloom).
  Evaluate: (a) how warming temperatures can alter vernalisation,
  florigen release, and pollen fertility; (b) what happens when plant
  flowering phenology shifts but pollinator phenology does not; (c) what
  field evidence would distinguish pollen heat failure from
  plant-pollinator phenological mismatch.
\end{enumerate}

\begin{enumerate}
\def\labelenumi{\arabic{enumi}.}
\setcounter{enumi}{29}
\tightlist
\item
  Critically evaluate CRISPR-based gene drives in wild plant
  populations: if a drive were engineered to spread herbicide resistance
  from a crop to its wild relatives, evaluate the ecological risks
  (invasive hybrid vigour, competitive dominance), the containment
  strategies (temporal/spatial/genetic containment), and the regulatory
  framework needed. Is there ever a scenario where a gene drive in
  plants would be ethically justified?
\end{enumerate}

\newpage

\section*{Questions --- Ch 27: Plant Responses to the
Environment}\label{questions-ch-27-plant-responses-to-the-environment}
\addcontentsline{toc}{section}{Questions --- Ch 27: Plant Responses to
the Environment}

\label{sec:q_unit_VIII_plant_responses}

\subsection*{Instructor Use and Coverage
Notes}\label{instructor-use-and-coverage-notes-29}
\addcontentsline{toc}{subsection}{Instructor Use and Coverage Notes}

\begin{itemize}
\tightlist
\item
  \textbf{Coverage target:} Predict a plant response from a signal and
  identify the likely hormone or pathway.
\item
  \textbf{Model/data emphasis:} Dose-response, water-use-efficiency, and
  hormone-interaction reasoning.
\item
  \textbf{Assessment alignment:} Visual Representations, Questions and
  Methods, Argumentation.
\item
  \textbf{Misconception probe:} A plant response is not passive; plants
  actively regulate development and physiology without neurons.
\item
  \textbf{Transfer product:} Transfer response logic to shade avoidance,
  drought, flowering, and crop management.
\item
  \textbf{Grading focus:} award full credit for mechanism, evidence,
  boundary conditions, and units when a calculation is required.
\item
  \textbf{Suggested use:} draw one recall item, one application item,
  and one synthesis item when building a short quiz from this bank.
\end{itemize}

\subsection*{Questions 1--10: Recall and
Comprehension}\label{questions-110-recall-and-comprehension-26}
\addcontentsline{toc}{subsection}{Questions 1--10: Recall and
Comprehension}

\emph{This activity accompanies \cref{sec:unit_VIII_plant_responses} of
the textbook --- review that chapter before attempting the exercises
below.}

\begin{enumerate}
\def\labelenumi{\arabic{enumi}.}
\tightlist
\item
  What is a tropism? Name three types and state the stimulus for each.
\end{enumerate}

\begin{enumerate}
\def\labelenumi{\arabic{enumi}.}
\setcounter{enumi}{1}
\tightlist
\item
  What is auxin? Where is it produced and how does it move through the
  plant?
\end{enumerate}

\begin{enumerate}
\def\labelenumi{\arabic{enumi}.}
\setcounter{enumi}{2}
\tightlist
\item
  What is phototropism? What photoreceptor detects blue light
  (phototropin)?
\end{enumerate}

\begin{enumerate}
\def\labelenumi{\arabic{enumi}.}
\setcounter{enumi}{3}
\tightlist
\item
  What is gravitropism? How do amyloplasts (statoliths) detect gravity?
\end{enumerate}

\begin{enumerate}
\def\labelenumi{\arabic{enumi}.}
\setcounter{enumi}{4}
\tightlist
\item
  Name four plant hormones other than auxin and state one major function
  of each.
\end{enumerate}

\begin{enumerate}
\def\labelenumi{\arabic{enumi}.}
\setcounter{enumi}{5}
\tightlist
\item
  What is ethylene? What role does it play in fruit ripening?
\end{enumerate}

\begin{enumerate}
\def\labelenumi{\arabic{enumi}.}
\setcounter{enumi}{6}
\tightlist
\item
  What is the phytochrome system? Distinguish Pr and Pfr forms.
\end{enumerate}

\begin{enumerate}
\def\labelenumi{\arabic{enumi}.}
\setcounter{enumi}{7}
\tightlist
\item
  What are plant defence responses to herbivory? Name two chemical
  defences.
\end{enumerate}

\begin{enumerate}
\def\labelenumi{\arabic{enumi}.}
\setcounter{enumi}{8}
\tightlist
\item
  Define thigmotropism. Give one example.
\end{enumerate}

\begin{enumerate}
\def\labelenumi{\arabic{enumi}.}
\setcounter{enumi}{9}
\tightlist
\item
  What is vernalisation? Which gene (FLC) is epigenetically silenced by
  prolonged cold?
\end{enumerate}

\subsection*{Questions 11--20: Application and
Analysis}\label{questions-1120-application-and-analysis-26}
\addcontentsline{toc}{subsection}{Questions 11--20: Application and
Analysis}

\begin{enumerate}
\def\labelenumi{\arabic{enumi}.}
\setcounter{enumi}{10}
\tightlist
\item
  The Cholodny-Went model explains phototropism: auxin redistributes
  laterally to the shaded side, causing differential elongation.
  Describe the evidence: (a) agar block experiments; (b) PIN3 auxin
  efflux carrier redistribution; (c) auxin-responsive reporter
  (DR5::GFP). What alternative explanations (differential auxin
  sensitivity, flavonoid phototropins) have been proposed?
\end{enumerate}

\begin{enumerate}
\def\labelenumi{\arabic{enumi}.}
\setcounter{enumi}{11}
\tightlist
\item
  Gibberellins promote stem elongation by degrading DELLA growth
  repressors. Explain: (a) the molecular pathway (GA \texorpdfstring{\ensuremath{\to}}{->} GID1 receptor \texorpdfstring{\ensuremath{\to}}{->}
  DELLA ubiquitination \texorpdfstring{\ensuremath{\to}}{->} proteasome degradation); (b) how dwarf mutants
  (GA-deficient or DELLA gain-of-function) fail to elongate; (c) why the
  Green Revolution wheat varieties (Rht-B1b, Rht-D1b) carry
  semi-dominant DELLA mutations that reduce height and lodging.
\end{enumerate}

\begin{enumerate}
\def\labelenumi{\arabic{enumi}.}
\setcounter{enumi}{12}
\tightlist
\item
  Ethylene triggers the triple response in dark-grown seedlings
  (shortened stem, thickened stem, horizontal growth). Design an
  experiment using: (a) the ethylene biosynthesis inhibitor AVG; (b) the
  ethylene-insensitive mutant \emph{ein2}; (c) an ethylene-overproducing
  mutant (\emph{eto1}) to confirm that each component of the triple
  response is ethylene-dependent.
\end{enumerate}

\begin{enumerate}
\def\labelenumi{\arabic{enumi}.}
\setcounter{enumi}{13}
\tightlist
\item
  Jasmonic acid (JA) mediates herbivore-induced defences. When a
  caterpillar chews a leaf, systemin (a peptide signal) activates JA
  biosynthesis, which induces protease inhibitors in distant unwounded
  leaves. Trace the entire signalling cascade: wounding \texorpdfstring{\ensuremath{\to}}{->} systemin
  release \texorpdfstring{\ensuremath{\to}}{->} systemin receptor \texorpdfstring{\ensuremath{\to}}{->} JA biosynthesis (from linolenic acid in
  chloroplast) \texorpdfstring{\ensuremath{\to}}{->} COI1-JAZ interaction \texorpdfstring{\ensuremath{\to}}{->} MYC2 transcription factor
  activation \texorpdfstring{\ensuremath{\to}}{->} PI gene expression.
\end{enumerate}

\begin{enumerate}
\def\labelenumi{\arabic{enumi}.}
\setcounter{enumi}{14}
\tightlist
\item
  Phytochrome interconversion: Pr (absorbs red 660 nm, inactive) \texorpdfstring{\ensuremath{\rightleftharpoons}}{<=>} Pfr
  (absorbs far-red 730 nm, active). In a long-day plant, flowering is
  promoted by long photoperiods (short nights). Explain: (a) how the
  phytochrome system measures night length (Pfr \texorpdfstring{\ensuremath{\to}}{->} Pr reversion rate in
  the dark); (b) how a brief red flash in the middle of a long night
  promotes flowering; (c) how a subsequent far-red flash reverses this.
\end{enumerate}

\begin{enumerate}
\def\labelenumi{\arabic{enumi}.}
\setcounter{enumi}{15}
\tightlist
\item
  Abscisic acid (ABA) promotes stomatal closure under drought stress.
  Explain the guard cell signalling cascade: ABA \texorpdfstring{\ensuremath{\to}}{->} RCAR/PYL receptor \texorpdfstring{\ensuremath{\to}}{->}
  PP2C inactivation \texorpdfstring{\ensuremath{\to}}{->} SnRK2 activation \texorpdfstring{\ensuremath{\to}}{->} SLAC1 anion channel opening \texorpdfstring{\ensuremath{\to}}{->}
  Cl\texorpdfstring{\textsuperscript{-}}{-} and malate\texorpdfstring{\textsuperscript{2}}{2}\texorpdfstring{\textsuperscript{-}}{-} efflux \texorpdfstring{\ensuremath{\to}}{->} K\texorpdfstring{\textsuperscript{+}}{+} efflux \texorpdfstring{\ensuremath{\to}}{->} turgor loss \texorpdfstring{\ensuremath{\to}}{->} stomata close.
\end{enumerate}

\begin{enumerate}
\def\labelenumi{\arabic{enumi}.}
\setcounter{enumi}{16}
\tightlist
\item
  Auxin dose-response curves differ between stems and roots: stems show
  maximum elongation at about 10\texorpdfstring{\textsuperscript{-}}{-}\texorpdfstring{\textsuperscript{5}}{5} M auxin, while roots show
  maximum elongation at about 10\texorpdfstring{\textsuperscript{-}}{-}\texorpdfstring{\textsuperscript{1}}{1}\texorpdfstring{\textsuperscript{0}}{0} M auxin and are inhibited
  at concentrations that stimulate stems. Explain: (a) how this
  difference underlies gravitropism (auxin accumulates on the lower
  side, inhibiting root elongation while stimulating shoot elongation);
  (b) why 2,4-D (a synthetic auxin) is an effective herbicide for
  broadleaf plants.
\end{enumerate}

\begin{enumerate}
\def\labelenumi{\arabic{enumi}.}
\setcounter{enumi}{17}
\tightlist
\item
  Plants produce volatile organic compounds (VOCs) in response to
  herbivory that attract predators of the herbivore (``crying for
  help''). The lima bean-spider mite-predatory mite tritrophic system is
  a classic example. Explain: (a) which VOCs (terpenoids, green leaf
  volatiles) are emitted; (b) how neighbouring undamaged plants
  ``eavesdrop'' via VOC-induced defensive gene expression; (c) whether
  this constitutes plant communication or merely an information-rich
  cue.
\end{enumerate}

\begin{enumerate}
\def\labelenumi{\arabic{enumi}.}
\setcounter{enumi}{18}
\tightlist
\item
  Circadian clock in plants: the core oscillator (CCA1/LHY \texorpdfstring{\ensuremath{\to}}{->} TOC1 \texorpdfstring{\ensuremath{\to}}{->}
  CCA1/LHY feedback loop) regulates about 30\% of plant gene
  expression. Explain: (a) how the clock anticipates dawn (pre-dawn
  expression of photosynthesis genes); (b) how jet-lag experiments
  (sudden photoperiod shift) reveal resynchronisation kinetics; (c) why
  clock mutants (TOC1-ox) have reduced fitness under natural light-dark
  cycles.
\end{enumerate}

\begin{enumerate}
\def\labelenumi{\arabic{enumi}.}
\setcounter{enumi}{19}
\tightlist
\item
  SAR (systemic acquired resistance) provides long-lasting,
  broad-spectrum pathogen protection. Following a localised pathogen
  infection, the signal molecule salicylic acid (SA) activates the SA
  pathway \texorpdfstring{\ensuremath{\to}}{->} NPR1 monomerisation \texorpdfstring{\ensuremath{\to}}{->} nuclear import \texorpdfstring{\ensuremath{\to}}{->} TGA transcription
  factor activation \texorpdfstring{\ensuremath{\to}}{->} PR (pathogenesis-related) gene expression. How
  does this confer resistance to secondary infections in distal tissues?
\end{enumerate}

\subsection*{Questions 21--30: Synthesis and
Evaluation}\label{questions-2130-synthesis-and-evaluation-26}
\addcontentsline{toc}{subsection}{Questions 21--30: Synthesis and
Evaluation}

\begin{enumerate}
\def\labelenumi{\arabic{enumi}.}
\setcounter{enumi}{20}
\tightlist
\item
  Evaluate the evidence for ``plant intelligence'' or ``plant
  neurobiology'' --- the concept that plant signalling networks
  (electrical action potentials, calcium waves, hormone gradients)
  constitute a form of information processing analogous to animal
  brains. Critically assess: (a) what criteria define intelligence; (b)
  whether plants satisfy any of these; (c) the boundaries of legitimate
  analogy vs anthropomorphism.
\end{enumerate}

\begin{enumerate}
\def\labelenumi{\arabic{enumi}.}
\setcounter{enumi}{21}
\tightlist
\item
  Design a CRISPR knockout experiment in Arabidopsis to determine
  whether the gravitropic response requires the starch-statolith
  mechanism (amyloplast sedimentation) or an alternative membrane-based
  mechanosensing pathway. What gene(s) would you knock out (PGM1 for
  starchless mutants), what phenotype would you measure, and what result
  would distinguish the two hypotheses?
\end{enumerate}

\begin{enumerate}
\def\labelenumi{\arabic{enumi}.}
\setcounter{enumi}{22}
\tightlist
\item
  Evaluate the molecular mechanisms of plant adaptation to flooding
  stress (waterlogging and submergence). Compare: (a) the SUB1A strategy
  in rice (quiescence --- ethylene \texorpdfstring{\ensuremath{\to}}{->} SUB1A \texorpdfstring{\ensuremath{\to}}{->} DELLA stabilisation \texorpdfstring{\ensuremath{\to}}{->}
  growth arrest); (b) the SNORKEL strategy (escape --- ethylene \texorpdfstring{\ensuremath{\to}}{->}
  SK1/SK2 \texorpdfstring{\ensuremath{\to}}{->} GA signalling \texorpdfstring{\ensuremath{\to}}{->} rapid internode elongation to escape
  floodwater). Under what flooding regimes is each strategy
  advantageous?
\end{enumerate}

\begin{enumerate}
\def\labelenumi{\arabic{enumi}.}
\setcounter{enumi}{23}
\tightlist
\item
  Climate change increases UV-B radiation at Earth's surface due to
  ozone depletion. Evaluate plant UV-B responses: (a) UVR8
  photoreceptor-mediated signalling \texorpdfstring{\ensuremath{\to}}{->} COP1-SPA \texorpdfstring{\ensuremath{\to}}{->} HY5 transcription
  factor \texorpdfstring{\ensuremath{\to}}{->} flavonoid/anthocyanin biosynthesis (sunscreen); (b) UV-B
  effects on DNA (CPD lesion formation) and repair (photolyase); (c)
  predicted impacts on crop productivity.
\end{enumerate}

\begin{enumerate}
\def\labelenumi{\arabic{enumi}.}
\setcounter{enumi}{24}
\tightlist
\item
  Evaluate the evolutionary arms race between plants and herbivores.
  Using the example of glucosinolate-myrosinase system in Brassicaceae
  (the ``mustard oil bomb''): (a) how the system works (glucosinolates
  are stored separately from myrosinase; upon tissue damage, they mix to
  produce toxic isothiocyanates); (b) how specialised herbivores
  (diamondback moth) detoxify isothiocyanates; (c) what countermeasure
  plant lineages have evolved (novel glucosinolates).
\end{enumerate}

\begin{enumerate}
\def\labelenumi{\arabic{enumi}.}
\setcounter{enumi}{25}
\tightlist
\item
  Symbiotic nitrogen fixation (Rhizobium-legume): evaluate the molecular
  cross-talk leading to nodule formation: (a) plant-derived flavonoids \texorpdfstring{\ensuremath{\to}}{->}
  bacterial NodD \texorpdfstring{\ensuremath{\to}}{->} Nod factor production \texorpdfstring{\ensuremath{\to}}{->} root hair curling \texorpdfstring{\ensuremath{\to}}{->} infection
  thread \texorpdfstring{\ensuremath{\to}}{->} nodule primordia; (b) the oxygen paradox (nitrogenase is
  O\texorpdfstring{\textsubscript{2}}{2}-sensitive but Rhizobium needs aerobic respiration \texorpdfstring{\ensuremath{\to}}{->} leghaemoglobin
  O\texorpdfstring{\textsubscript{2}}{2} buffering); (c) the autoregulation of nodulation (too many nodules
  are costly \texorpdfstring{\ensuremath{\to}}{->} CLE peptide negative feedback).
\end{enumerate}

\begin{enumerate}
\def\labelenumi{\arabic{enumi}.}
\setcounter{enumi}{26}
\tightlist
\item
  Evaluate the role of epigenetics in plant stress memory. After a heat
  stress episode, some plants show enhanced heat tolerance for weeks
  (priming). What molecular marks (H3K4me3 at HSP loci, small RNA
  populations, chromatin remodelling) mediate stress memory, and does
  this memory have transgenerational stability?
\end{enumerate}

\begin{enumerate}
\def\labelenumi{\arabic{enumi}.}
\setcounter{enumi}{27}
\tightlist
\item
  Parasitic plants (Striga, Orobanche, dodder) exploit host plant
  signals to locate and attach to hosts. Evaluate: (a) how Striga seeds
  detect host-derived strigolactones as germination cues; (b) how the
  haustorium forms and penetrates host xylem/phloem; (c) why Striga is
  the most economically devastating weed in sub-Saharan Africa; (d) what
  control strategies (resistant crop varieties, suicidal germination
  stimulants) are being deployed.
\end{enumerate}

\begin{enumerate}
\def\labelenumi{\arabic{enumi}.}
\setcounter{enumi}{28}
\tightlist
\item
  Evaluate the potential of ``smart crops'' engineered with synthetic
  hormone signalling circuits: e.g., a crop that closes stomata in
  response to soil moisture sensors (synthetic ABA-independent
  signalling), activates JA defences primarily when herbivores are
  detected (synthetic promoter responsive to caterpillar oral
  secretions), and accelerates flowering under climate-variable
  conditions (synthetic florigen circuit). What are the engineering
  feasibility, field-validation, regulatory, and ecological risks?
\end{enumerate}

\begin{enumerate}
\def\labelenumi{\arabic{enumi}.}
\setcounter{enumi}{29}
\tightlist
\item
  Critically assess the claim that plant volatile communication
  represents a form of ``language'': (a) define what communication
  requires (sender, signal, receiver, response, evolved function); (b)
  evaluate whether plant VOC emission meets these criteria; (c) compare
  plant volatile signalling to animal pheromone communication; (d)
  assess the fitness benefit to the emitter (is it a true signal or an
  incidental cue exploited by neighbours?).
\end{enumerate}

\newpage

\section*{Questions --- Ch 28: Circulation, Respiration, and
Homeostasis}\label{questions-ch-28-circulation-respiration-and-homeostasis}
\addcontentsline{toc}{section}{Questions --- Ch 28: Circulation,
Respiration, and Homeostasis}

\label{sec:q_unit_IX_circulation_respiration_homeostasis}

\subsection*{Instructor Use and Coverage
Notes}\label{instructor-use-and-coverage-notes-30}
\addcontentsline{toc}{subsection}{Instructor Use and Coverage Notes}

\begin{itemize}
\tightlist
\item
  \textbf{Coverage target:} Calculate a transport or feedback variable
  and explain the compensatory response.
\item
  \textbf{Model/data emphasis:} Cardiac-output, diffusion,
  oxygen-saturation, and feedback calculations.
\item
  \textbf{Assessment alignment:} Visual Representations, Statistical
  Tests and Data Analysis, Argumentation.
\item
  \textbf{Misconception probe:} Homeostasis is dynamic regulation, not
  an unchanging internal state.
\item
  \textbf{Transfer product:} Transfer homeostatic reasoning to exercise,
  altitude, hemorrhage, fever, and shock.
\item
  \textbf{Grading focus:} award full credit for mechanism, evidence,
  boundary conditions, and units when a calculation is required.
\item
  \textbf{Suggested use:} draw one recall item, one application item,
  and one synthesis item when building a short quiz from this bank.
\end{itemize}

\subsection*{Questions 1--10: Recall and
Comprehension}\label{questions-110-recall-and-comprehension-27}
\addcontentsline{toc}{subsection}{Questions 1--10: Recall and
Comprehension}

\emph{This activity accompanies
\cref{sec:unit_IX_circulation_respiration_homeostasis} of the textbook
--- review that chapter before attempting the exercises below.}

\begin{enumerate}
\def\labelenumi{\arabic{enumi}.}
\tightlist
\item
  What is homeostasis? Give two examples of physiological variables
  maintained by homeostasis.
\end{enumerate}

\begin{enumerate}
\def\labelenumi{\arabic{enumi}.}
\setcounter{enumi}{1}
\tightlist
\item
  Describe the structure of the human heart. Name the four chambers and
  four valves.
\end{enumerate}

\begin{enumerate}
\def\labelenumi{\arabic{enumi}.}
\setcounter{enumi}{2}
\tightlist
\item
  What is the cardiac cycle? Define systole and diastole.
\end{enumerate}

\begin{enumerate}
\def\labelenumi{\arabic{enumi}.}
\setcounter{enumi}{3}
\tightlist
\item
  What is cardiac output (CO)? Write the formula CO = HR \texorpdfstring{\ensuremath{\times}}{x} SV.
\end{enumerate}

\begin{enumerate}
\def\labelenumi{\arabic{enumi}.}
\setcounter{enumi}{4}
\tightlist
\item
  Describe the structure and function of arteries, veins, and
  capillaries.
\end{enumerate}

\begin{enumerate}
\def\labelenumi{\arabic{enumi}.}
\setcounter{enumi}{5}
\tightlist
\item
  What is blood pressure? What are typical systolic and diastolic
  values?
\end{enumerate}

\begin{enumerate}
\def\labelenumi{\arabic{enumi}.}
\setcounter{enumi}{6}
\tightlist
\item
  What is tidal volume? What is vital capacity?
\end{enumerate}

\begin{enumerate}
\def\labelenumi{\arabic{enumi}.}
\setcounter{enumi}{7}
\tightlist
\item
  Describe the mechanics of ventilation (breathing): inhalation and
  exhalation.
\end{enumerate}

\begin{enumerate}
\def\labelenumi{\arabic{enumi}.}
\setcounter{enumi}{8}
\tightlist
\item
  What is the oxygen-haemoglobin dissociation curve? What does its
  sigmoidal shape indicate?
\end{enumerate}

\begin{enumerate}
\def\labelenumi{\arabic{enumi}.}
\setcounter{enumi}{9}
\tightlist
\item
  Define thermoregulation. Distinguish endotherms from ectotherms.
\end{enumerate}

\subsection*{Questions 11--20: Application and
Analysis}\label{questions-1120-application-and-analysis-27}
\addcontentsline{toc}{subsection}{Questions 11--20: Application and
Analysis}

\begin{enumerate}
\def\labelenumi{\arabic{enumi}.}
\setcounter{enumi}{10}
\tightlist
\item
  During exercise, HR increases from 70 to 160 bpm and SV from 70 mL to
  120 mL. Calculate the CO at rest and during exercise. What is the
  fold-increase? What mechanisms increase HR (sympathetic \texorpdfstring{\ensuremath{\to}}{->} SA node \texorpdfstring{\ensuremath{\to}}{->}
  faster depolarisation) and SV (Frank-Starling mechanism)?
\end{enumerate}

\begin{enumerate}
\def\labelenumi{\arabic{enumi}.}
\setcounter{enumi}{11}
\tightlist
\item
  The Bohr effect: increased CO\texorpdfstring{\textsubscript{2}}{2} and H\texorpdfstring{\textsuperscript{+}}{+} in exercising muscle shift the
  oxygen-haemoglobin dissociation curve rightward, decreasing Hb's O\texorpdfstring{\textsubscript{2}}{2}
  affinity. Explain: (a) the molecular basis (protonation of His146
  stabilising the T-state); (b) the physiological benefit (enhanced O\texorpdfstring{\textsubscript{2}}{2}
  unloading in active tissues); (c) what happens at the lungs (low CO\texorpdfstring{\textsubscript{2}}{2},
  high pH \texorpdfstring{\ensuremath{\to}}{->} leftward shift \texorpdfstring{\ensuremath{\to}}{->} enhanced O\texorpdfstring{\textsubscript{2}}{2} loading).
\end{enumerate}

\begin{enumerate}
\def\labelenumi{\arabic{enumi}.}
\setcounter{enumi}{12}
\tightlist
\item
  Ventilation-perfusion (V/Q) matching ensures that regions of the lung
  receiving airflow also receive blood flow. Explain: (a) what happens
  when V/Q is high (dead space) vs low (shunt); (b) how local hypoxic
  pulmonary vasoconstriction (HPV) diverts blood from poorly ventilated
  alveoli; (c) why chronic lung disease (COPD) causes V/Q mismatch and
  hypoxaemia.
\end{enumerate}

\begin{enumerate}
\def\labelenumi{\arabic{enumi}.}
\setcounter{enumi}{13}
\tightlist
\item
  Blood pressure regulation: the baroreceptor reflex senses carotid
  sinus stretch. When BP drops: (a) baroreceptor firing decreases; (b)
  medullary cardiovascular centre increases sympathetic output; (c) HR,
  contractility, and peripheral resistance increase; (d) BP is restored.
  What happens if the baroreceptor reflex fails (autonomic neuropathy)?
\end{enumerate}

\begin{enumerate}
\def\labelenumi{\arabic{enumi}.}
\setcounter{enumi}{14}
\tightlist
\item
  Surfactant (produced by type II alveolar cells) reduces alveolar
  surface tension. Without surfactant: (a) what happens to alveolar
  compliance (Laplace's law: P = 2T/r \texorpdfstring{\ensuremath{\to}}{->} smaller alveoli collapse into
  larger ones); (b) what condition results in premature infants (NRDS
  --- neonatal respiratory distress syndrome); (c) how is antenatal
  corticosteroid treatment used to mature fetal lungs?
\end{enumerate}

\begin{enumerate}
\def\labelenumi{\arabic{enumi}.}
\setcounter{enumi}{15}
\tightlist
\item
  The kidney maintains blood osmolality homeostasis. ADH (vasopressin)
  is released from the posterior pituitary in response to high blood
  osmolality. Explain: (a) how ADH regulates AQP2 insertion into the
  collecting duct apical membrane; (b) why ADH deficiency produces
  dilute, voluminous urine (diabetes insipidus); (c) how alcohol
  inhibits ADH secretion and causes diuresis.
\end{enumerate}

\begin{enumerate}
\def\labelenumi{\arabic{enumi}.}
\setcounter{enumi}{16}
\tightlist
\item
  Counter-current exchange in the vasa recta (renal medullary blood
  vessels) maintains the osmotic gradient established by the loop of
  Henle. Explain: (a) why the vasa recta loop parallel to the loop of
  Henle; (b) how ascending and descending limbs exchange solutes; (c)
  what would happen to medullary osmolality if the vasa recta were
  replaced by a single straight vessel.
\end{enumerate}

\begin{enumerate}
\def\labelenumi{\arabic{enumi}.}
\setcounter{enumi}{17}
\tightlist
\item
  During cold exposure, thermoregulation maintains core temperature at
  37°C. Trace the response: peripheral thermoreceptors \texorpdfstring{\ensuremath{\to}}{->} hypothalamus \texorpdfstring{\ensuremath{\to}}{->}
  sympathetic nervous system \texorpdfstring{\ensuremath{\to}}{->} cutaneous vasoconstriction \texorpdfstring{\ensuremath{\to}}{->} reduced heat
  loss + shivering \texorpdfstring{\ensuremath{\to}}{->} increased heat production. What is the set-point
  concept, and how does fever alter the set-point?
\end{enumerate}

\begin{enumerate}
\def\labelenumi{\arabic{enumi}.}
\setcounter{enumi}{18}
\tightlist
\item
  The oxygen cascade describes the progressive decrease in PO\texorpdfstring{\textsubscript{2}}{2} from
  atmosphere (160 mmHg) to mitochondria (about 1--3 mmHg).
  Identify each step: inspired air \texorpdfstring{\ensuremath{\to}}{->} alveolar air (104 mmHg, diluted by
  CO\texorpdfstring{\textsubscript{2}}{2} and water vapour) \texorpdfstring{\ensuremath{\to}}{->} arterial blood (100 mmHg, some V/Q mismatch) \texorpdfstring{\ensuremath{\to}}{->}
  capillary blood \texorpdfstring{\ensuremath{\to}}{->} tissue \texorpdfstring{\ensuremath{\to}}{->} mitochondria. At which step is the largest
  PO\texorpdfstring{\textsubscript{2}}{2} drop, and what determines its magnitude?
\end{enumerate}

\begin{enumerate}
\def\labelenumi{\arabic{enumi}.}
\setcounter{enumi}{19}
\tightlist
\item
  Heart failure with reduced ejection fraction (HFrEF) results in an EF
  \textless{} 40\% (normal 55--70\%). Explain: (a) why reduced
  contractility leads to fluid retention (RAAS activation); (b) why
  diuretics provide symptomatic relief but do not improve survival; (c)
  how ACE inhibitors + beta-blockers + ARNI (sacubitril/valsartan) have
  been shown to reduce mortality.
\end{enumerate}

\subsection*{Questions 21--30: Synthesis and
Evaluation}\label{questions-2130-synthesis-and-evaluation-27}
\addcontentsline{toc}{subsection}{Questions 21--30: Synthesis and
Evaluation}

\begin{enumerate}
\def\labelenumi{\arabic{enumi}.}
\setcounter{enumi}{20}
\tightlist
\item
  Evaluate the Frank-Starling mechanism at the molecular level:
  increased preload \texorpdfstring{\ensuremath{\to}}{->} increased sarcomere stretch \texorpdfstring{\ensuremath{\to}}{->} increased Ca\texorpdfstring{\textsuperscript{2}}{2}\texorpdfstring{\textsuperscript{+}}{+}
  sensitivity of troponin C \texorpdfstring{\ensuremath{\to}}{->} increased force production. What molecular
  mechanism converts sarcomere length into increased Ca\texorpdfstring{\textsuperscript{2}}{2}\texorpdfstring{\textsuperscript{+}}{+} sensitivity,
  and why does this relationship break down in dilated cardiomyopathy
  (DCM)?
\end{enumerate}

\begin{enumerate}
\def\labelenumi{\arabic{enumi}.}
\setcounter{enumi}{21}
\tightlist
\item
  Design an experiment to measure the maximal aerobic capacity (VO\texorpdfstring{\textsubscript{2}}{2} max)
  of a human subject. Describe: the incremental exercise protocol
  (treadmill or cycle), the gas analysis measurements (O\texorpdfstring{\textsubscript{2}}{2} consumption,
  CO\texorpdfstring{\textsubscript{2}}{2} production, RER), and the criteria for reaching true VO\texorpdfstring{\textsubscript{2}}{2} max. How
  does VO\texorpdfstring{\textsubscript{2}}{2} max differ between trained athletes and sedentary
  individuals, and what factors (cardiac output, O\texorpdfstring{\textsubscript{2}}{2} extraction, muscle
  mitochondrial density) limit performance?
\end{enumerate}

\begin{enumerate}
\def\labelenumi{\arabic{enumi}.}
\setcounter{enumi}{22}
\tightlist
\item
  Evaluate the roles of the sympathetic and parasympathetic nervous
  systems in cardiovascular regulation via their respective
  neurotransmitters (norepinephrine \texorpdfstring{\ensuremath{\to}}{->} \texorpdfstring{\ensuremath{\beta}}{beta}\texorpdfstring{\textsubscript{1}}{1} receptors \texorpdfstring{\ensuremath{\to}}{->} increased HR vs
  acetylcholine \texorpdfstring{\ensuremath{\to}}{->} M\texorpdfstring{\textsubscript{2}}{2} receptors \texorpdfstring{\ensuremath{\to}}{->} decreased HR). What is ``vagal tone''
  at rest, and why does denervation of the heart (cardiac transplant
  recipients) result in a resting HR of about 100 bpm?
\end{enumerate}

\begin{enumerate}
\def\labelenumi{\arabic{enumi}.}
\setcounter{enumi}{23}
\tightlist
\item
  Evaluate the physiological trade-off between capillary density and
  distance-limited O\texorpdfstring{\textsubscript{2}}{2} diffusion: (a) use Fick's law of diffusion (J = D
  \texorpdfstring{\ensuremath{\times}}{x} A \texorpdfstring{\ensuremath{\times}}{x} \texorpdfstring{\ensuremath{\Delta}}{Delta}P/d) to explain why capillary density must increase in
  metabolically active tissues; (b) calculate the maximum O\texorpdfstring{\textsubscript{2}}{2} diffusion
  distance (about 100 \texorpdfstring{\ensuremath{\mu}}{mu}m); (c) explain why tumours develop
  necrotic cores when they grow beyond 200 \texorpdfstring{\ensuremath{\mu}}{mu}m without angiogenesis.
\end{enumerate}

\begin{enumerate}
\def\labelenumi{\arabic{enumi}.}
\setcounter{enumi}{24}
\tightlist
\item
  High altitude acclimatisation involves acute and chronic responses:
  acute (hyperventilation, increased HR); chronic (EPO \texorpdfstring{\ensuremath{\to}}{->} erythrocytosis,
  2,3-BPG increase, vascular remodelling). Evaluate: (a) how the carotid
  body senses hypoxia and triggers hyperventilation; (b) why respiratory
  alkalosis results; (c) how renal bicarbonate excretion compensates;
  (d) why chronic mountain sickness (Monge's disease) can result from
  excessive erythropoiesis.
\end{enumerate}

\begin{enumerate}
\def\labelenumi{\arabic{enumi}.}
\setcounter{enumi}{25}
\tightlist
\item
  Evaluate the ``oxygen paradox'' in ischaemia-reperfusion injury:
  restored oxygen supply to ischaemic tissue generates ROS (via xanthine
  oxidase and Complex I) that cause greater damage than the ischaemia
  itself. How do ROS cause injury (membrane lipid peroxidation,
  mitochondrial permeability transition pore opening), and what
  therapeutic strategies (hypothermia, ischaemic pre-conditioning)
  mitigate reperfusion injury?
\end{enumerate}

\begin{enumerate}
\def\labelenumi{\arabic{enumi}.}
\setcounter{enumi}{26}
\tightlist
\item
  Evaluate the hypothesis that endothermy evolved from aerobic capacity
  (the ``aerobic capacity model''): selection for sustained locomotion \texorpdfstring{\ensuremath{\to}}{->}
  increased mitochondrial density \texorpdfstring{\ensuremath{\to}}{->} increased resting metabolic rate \texorpdfstring{\ensuremath{\to}}{->}
  internal heat production \texorpdfstring{\ensuremath{\to}}{->} endothermy. Compare this to the ``parental
  care model'' (selection for incubation of eggs \texorpdfstring{\ensuremath{\to}}{->} endothermy). Which
  model is better supported by paleontological and comparative
  physiological evidence?
\end{enumerate}

\begin{enumerate}
\def\labelenumi{\arabic{enumi}.}
\setcounter{enumi}{27}
\tightlist
\item
  The lymphatic system returns interstitial fluid to the bloodstream.
  Evaluate: (a) how lymph is formed (Starling's forces \texorpdfstring{\ensuremath{\to}}{->} net filtration
  at the arterial end of capillaries); (b) what happens when lymphatic
  drainage fails (lymphoedema --- elephantiasis caused by parasitic
  nematode \emph{Wuchereria bancrofti}); (c) the dual role of lymphatics
  in immune surveillance (dendritic cell migration to lymph nodes).
\end{enumerate}

\begin{enumerate}
\def\labelenumi{\arabic{enumi}.}
\setcounter{enumi}{28}
\tightlist
\item
  Diving mammals (Weddell seal, sperm whale) can hold their breath for
  \textgreater60 minutes while diving to \textgreater1,000 m. Evaluate
  four physiological adaptations: (a) bradycardia and peripheral
  vasoconstriction (dive reflex); (b) elevated myoglobin content (O\texorpdfstring{\textsubscript{2}}{2}
  storage); (c) splenic contraction (autotransfusion of O\texorpdfstring{\textsubscript{2}}{2}-carrying
  RBCs); (d) tolerance of lactate (anaerobic metabolism during deep
  dives). How do these adaptations differ between pinnipeds and
  cetaceans?
\end{enumerate}

\begin{enumerate}
\def\labelenumi{\arabic{enumi}.}
\setcounter{enumi}{29}
\tightlist
\item
  Critically evaluate the concept of ``physiological reserve'' --- the
  excess capacity of organ systems beyond what is needed at rest (e.g.,
  cardiac output can increase 5\texorpdfstring{\ensuremath{\times}}{x} during exercise, kidneys can maintain
  GFR with a single kidney). How does ageing erode physiological
  reserve, and at what threshold does age-related reserve depletion
  produce clinical disease? Apply the concept to explain why the elderly
  are disproportionately vulnerable to COVID-19, heat waves, and
  surgical stress. Then transfer the same homeostatic-reserve logic to a
  salty high-protein meal with low water intake: predict plasma
  osmolality, ADH/aldosterone direction, hepatic urea production, urine
  concentration, and why desert mammals and freshwater fish handle
  nitrogen differently.
\end{enumerate}

\newpage

\section*{Questions --- Ch 29: Nervous System and Neural
Signalling}\label{questions-ch-29-nervous-system-and-neural-signalling}
\addcontentsline{toc}{section}{Questions --- Ch 29: Nervous System and
Neural Signalling}

\label{sec:q_unit_IX_nervous_system}

\subsection*{Instructor Use and Coverage
Notes}\label{instructor-use-and-coverage-notes-31}
\addcontentsline{toc}{subsection}{Instructor Use and Coverage Notes}

\begin{itemize}
\tightlist
\item
  \textbf{Coverage target:} Predict behavioural or signal consequences
  from a circuit or structural perturbation.
\item
  \textbf{Model/data emphasis:} Cable-length, conduction, and simple
  circuit calculations.
\item
  \textbf{Assessment alignment:} Visual Representations, Statistical
  Tests and Data Analysis, Argumentation.
\item
  \textbf{Misconception probe:} The brain is not only neurons; glia and
  circuit context are essential to function.
\item
  \textbf{Transfer product:} Transfer circuit reasoning to reflexes,
  sensory systems, learning, and disease.
\item
  \textbf{Grading focus:} award full credit for mechanism, evidence,
  boundary conditions, and units when a calculation is required.
\item
  \textbf{Suggested use:} draw one recall item, one application item,
  and one synthesis item when building a short quiz from this bank.
\end{itemize}

\subsection*{Questions 1--10: Recall and
Comprehension}\label{questions-110-recall-and-comprehension-28}
\addcontentsline{toc}{subsection}{Questions 1--10: Recall and
Comprehension}

\emph{This activity accompanies \cref{sec:unit_IX_nervous_system} of the
textbook --- review that chapter before attempting the exercises below.}

\begin{enumerate}
\def\labelenumi{\arabic{enumi}.}
\tightlist
\item
  Name the two divisions of the nervous system. What does each contain?
\end{enumerate}

\begin{enumerate}
\def\labelenumi{\arabic{enumi}.}
\setcounter{enumi}{1}
\tightlist
\item
  What is a neuron? Draw and label the soma, dendrites, axon, and
  synaptic terminal.
\end{enumerate}

\begin{enumerate}
\def\labelenumi{\arabic{enumi}.}
\setcounter{enumi}{2}
\tightlist
\item
  What is the resting membrane potential of a typical neuron? Which ion
  is most permeable at rest?
\end{enumerate}

\begin{enumerate}
\def\labelenumi{\arabic{enumi}.}
\setcounter{enumi}{3}
\tightlist
\item
  What is an action potential? Describe the four phases (resting,
  depolarisation, repolarisation, hyperpolarisation).
\end{enumerate}

\begin{enumerate}
\def\labelenumi{\arabic{enumi}.}
\setcounter{enumi}{4}
\tightlist
\item
  What is the myelin sheath? Name the cells that produce myelin in the
  CNS and PNS.
\end{enumerate}

\begin{enumerate}
\def\labelenumi{\arabic{enumi}.}
\setcounter{enumi}{5}
\tightlist
\item
  What is saltatory conduction? How does it increase conduction
  velocity?
\end{enumerate}

\begin{enumerate}
\def\labelenumi{\arabic{enumi}.}
\setcounter{enumi}{6}
\tightlist
\item
  Define a synapse. Distinguish chemical from electrical synapses.
\end{enumerate}

\begin{enumerate}
\def\labelenumi{\arabic{enumi}.}
\setcounter{enumi}{7}
\tightlist
\item
  Name three neurotransmitters and state one function of each.
\end{enumerate}

\begin{enumerate}
\def\labelenumi{\arabic{enumi}.}
\setcounter{enumi}{8}
\tightlist
\item
  What is the reflex arc? Name its five components.
\end{enumerate}

\begin{enumerate}
\def\labelenumi{\arabic{enumi}.}
\setcounter{enumi}{9}
\tightlist
\item
  Distinguish the somatic from the autonomic nervous system.
\end{enumerate}

\subsection*{Questions 11--20: Application and
Analysis}\label{questions-1120-application-and-analysis-28}
\addcontentsline{toc}{subsection}{Questions 11--20: Application and
Analysis}

\begin{enumerate}
\def\labelenumi{\arabic{enumi}.}
\setcounter{enumi}{10}
\tightlist
\item
  The Goldman equation predicts resting membrane potential from the
  relative permeabilities of Na\texorpdfstring{\textsuperscript{+}}{+}, K\texorpdfstring{\textsuperscript{+}}{+}, and Cl\texorpdfstring{\textsuperscript{-}}{-}. If a mutation doubles Na\texorpdfstring{\textsuperscript{+}}{+}
  permeability at rest (P\_Na increases 2\texorpdfstring{\ensuremath{\times}}{x}), predict the new resting
  potential direction (depolarisation toward E\_Na = +60 mV). Calculate
  using simplified Goldman equation.
\end{enumerate}

\begin{enumerate}
\def\labelenumi{\arabic{enumi}.}
\setcounter{enumi}{11}
\tightlist
\item
  Multiple sclerosis (MS) destroys myelin in the CNS. Predict: (a) the
  effect on conduction velocity; (b) which neuronal functions are
  impaired first; (c) why remyelination therapies are being
  investigated; (d) why the PNS is relatively spared (different
  myelinating cells).
\end{enumerate}

\begin{enumerate}
\def\labelenumi{\arabic{enumi}.}
\setcounter{enumi}{12}
\tightlist
\item
  Local anaesthetics (lidocaine) block voltage-gated Na\texorpdfstring{\textsuperscript{+}}{+} channels.
  Explain: (a) why this prevents action potential generation; (b) why
  pain fibres (small, unmyelinated C fibres) are blocked before motor
  fibres (large, myelinated A\texorpdfstring{\ensuremath{\alpha}}{alpha} fibres); (c) is this use-dependent block
  (more block with higher firing frequency)?
\end{enumerate}

\begin{enumerate}
\def\labelenumi{\arabic{enumi}.}
\setcounter{enumi}{13}
\tightlist
\item
  The patellar (knee-jerk) reflex involves a monosynaptic stretch reflex
  arc. Trace the pathway: muscle spindle \texorpdfstring{\ensuremath{\to}}{->} Ia afferent \texorpdfstring{\ensuremath{\to}}{->} spinal cord
  (monosynaptic excitation of \texorpdfstring{\ensuremath{\alpha}}{alpha}-motor neuron) \texorpdfstring{\ensuremath{\to}}{->} quadriceps contraction.
  Explain: (a) why the reflex is monosynaptic (no interneuron); (b) the
  role of reciprocal inhibition of the hamstring (via an inhibitory
  interneuron).
\end{enumerate}

\begin{enumerate}
\def\labelenumi{\arabic{enumi}.}
\setcounter{enumi}{14}
\tightlist
\item
  Sensory adaptation: olfactory receptors rapidly adapt to constant
  odour exposure (within minutes). Explain: (a) the molecular mechanism
  (cAMP \texorpdfstring{\ensuremath{\to}}{->} Ca\texorpdfstring{\textsuperscript{2}}{2}\texorpdfstring{\textsuperscript{+}}{+} \texorpdfstring{\ensuremath{\to}}{->} calmodulin \texorpdfstring{\ensuremath{\to}}{->} reduced CNG channel sensitivity); (b) the
  difference between phasic (rapidly adapting) and tonic (slowly
  adapting) receptors; (c) one example of a tonic receptor.
\end{enumerate}

\begin{enumerate}
\def\labelenumi{\arabic{enumi}.}
\setcounter{enumi}{15}
\tightlist
\item
  The blood-brain barrier (BBB) restricts drug entry into the CNS.
  Explain: (a) the cellular basis (tight junctions between endothelial
  cells, pericytes, astrocyte end-feet); (b) which molecules can cross
  (small, lipophilic, MW \textless{} 400); (c) how L-DOPA (prodrug for
  dopamine) exploits amino acid transporters to bypass the BBB.
\end{enumerate}

\begin{enumerate}
\def\labelenumi{\arabic{enumi}.}
\setcounter{enumi}{16}
\tightlist
\item
  Neuroplasticity: after a stroke destroys the motor cortex controlling
  the right hand, the left hand often improves in dexterity over weeks.
  Explain: (a) cortical remapping (adjacent cortical areas take over);
  (b) the role of long-term potentiation (LTP) in this process; (c) why
  intensive rehabilitation (constraint-induced movement therapy)
  accelerates recovery.
\end{enumerate}

\begin{enumerate}
\def\labelenumi{\arabic{enumi}.}
\setcounter{enumi}{17}
\tightlist
\item
  The autonomic nervous system has sympathetic (fight-or-flight) and
  parasympathetic (rest-and-digest) divisions. Compare: (a)
  preganglionic neurotransmitter (both ACh); (b) postganglionic
  neurotransmitter (sympathetic = NE, parasympathetic = ACh); (c)
  effects on heart rate, bronchioles, and GI motility.
\end{enumerate}

\begin{enumerate}
\def\labelenumi{\arabic{enumi}.}
\setcounter{enumi}{18}
\tightlist
\item
  Pain processing: nociceptors detect tissue damage and transmit signals
  via A\texorpdfstring{\ensuremath{\delta}}{delta} (sharp, fast) and C (dull, slow) fibres. Explain the gate
  control theory: how does activation of large-diameter A\texorpdfstring{\ensuremath{\beta}}{beta} touch fibres
  ``close the gate'' in the spinal cord dorsal horn, reducing pain
  transmission to the brain?
\end{enumerate}

\begin{enumerate}
\def\labelenumi{\arabic{enumi}.}
\setcounter{enumi}{19}
\tightlist
\item
  Glial cells (astrocytes, oligodendrocytes, microglia) outnumber
  neurons. Describe two functional roles of astrocytes beyond structural
  support: (a) glutamate uptake at synapses (preventing excitotoxicity);
  (b) regulation of extracellular K\texorpdfstring{\textsuperscript{+}}{+}; (c) the ``tripartite synapse''
  concept (astrocyte modulates synaptic strength via gliotransmitters).
\end{enumerate}

\subsection*{Questions 21--30: Synthesis and
Evaluation}\label{questions-2130-synthesis-and-evaluation-28}
\addcontentsline{toc}{subsection}{Questions 21--30: Synthesis and
Evaluation}

\begin{enumerate}
\def\labelenumi{\arabic{enumi}.}
\setcounter{enumi}{20}
\tightlist
\item
  Evaluate the ``integrate-and-fire'' model of neuronal computation: a
  neuron summates thousands of excitatory and inhibitory postsynaptic
  potentials (EPSPs and IPSPs) and fires when threshold is reached. What
  are the limitations of this model (dendritic computation, non-linear
  integration, backpropagating APs), and how do modern computational
  neuroscience models address them?
\end{enumerate}

\begin{enumerate}
\def\labelenumi{\arabic{enumi}.}
\setcounter{enumi}{21}
\tightlist
\item
  Design an experiment to determine whether a newly discovered brain
  region is involved in spatial memory rather than procedural memory.
  Describe: (a) lesion studies (selective ablation); (b)
  electrophysiology (place cells, grid cells); (c) optogenetic
  activation/inhibition during a behavioural task; (d) what behavioural
  paradigms (Morris water maze, T-maze) you would use.
\end{enumerate}

\begin{enumerate}
\def\labelenumi{\arabic{enumi}.}
\setcounter{enumi}{22}
\tightlist
\item
  Evaluate the molecular basis of Alzheimer's disease: (a) A\texorpdfstring{\ensuremath{\beta}}{beta} oligomer
  toxicity hypothesis; (b) tau hyperphosphorylation and neurofibrillary
  tangles; (c) the failure of anti-A\texorpdfstring{\ensuremath{\beta}}{beta} antibody therapies (lecanemab
  showed modest benefit, 27\% slowing of decline). What does the modest
  efficacy of anti-A\texorpdfstring{\ensuremath{\beta}}{beta} antibodies tell us about the ``amyloid cascade''
  as a therapeutic target?
\end{enumerate}

\begin{enumerate}
\def\labelenumi{\arabic{enumi}.}
\setcounter{enumi}{23}
\tightlist
\item
  Evaluate the role of adult neurogenesis in the hippocampus: (a)
  evidence that new neurons are continuously generated in the
  subgranular zone of the dentate gyrus; (b) the functional role of
  newborn neurons in pattern separation (distinguishing similar
  memories); (c) the controversial evidence for adult neurogenesis in
  humans (Spalding et al., \texorpdfstring{\textsuperscript{1}}{1}\texorpdfstring{\textsuperscript{4}}{4}C birth-dating vs Sorrells et al.,
  immunohistochemistry showing no new neurons after age 13).
\end{enumerate}

\begin{enumerate}
\def\labelenumi{\arabic{enumi}.}
\setcounter{enumi}{24}
\tightlist
\item
  Evaluate the mechanism of general anaesthesia: despite \textgreater150
  years of use, the molecular targets of volatile anaesthetics
  (sevoflurane, isoflurane) are incompletely understood. Current
  evidence implicates: (a) GABA\_A receptor potentiation; (b) NMDA
  receptor inhibition; (c) two-pore domain K\texorpdfstring{\textsuperscript{+}}{+} channel (TREK/TASK)
  activation. How does the combination of these effects produce
  unconsciousness, and why is the ``unitary'' theory of anaesthesia (one
  target) now rejected?
\end{enumerate}

\begin{enumerate}
\def\labelenumi{\arabic{enumi}.}
\setcounter{enumi}{25}
\tightlist
\item
  Evaluate the concept of ``neural oscillations'' --- rhythmic brain
  activity at specific frequency bands (delta 1--4 Hz, theta 4--8 Hz,
  alpha 8--13 Hz, beta 13--30 Hz, gamma 30--100 Hz). What cognitive
  processes are associated with each band, and what evidence from
  EEG/MEG studies supports the idea that frequency-band-specific
  synchronisation enables ``binding'' of distributed cortical
  representations into a unified percept?
\end{enumerate}

\begin{enumerate}
\def\labelenumi{\arabic{enumi}.}
\setcounter{enumi}{26}
\tightlist
\item
  Chronic pain involves maladaptive plasticity (central sensitisation)
  in the spinal cord dorsal horn. Evaluate: (a) the molecular mechanism
  (NMDA receptor-dependent wind-up, glial activation, BDNF release); (b)
  why opioids provide temporary relief but worsen chronic pain through
  opioid-induced hyperalgesia (OIH); (c) what non-opioid alternatives
  (gabapentinoids, SNRI antidepressants, spinal cord stimulation) target
  central sensitisation.
\end{enumerate}

\begin{enumerate}
\def\labelenumi{\arabic{enumi}.}
\setcounter{enumi}{27}
\tightlist
\item
  The enteric nervous system (ENS, ``second brain'') contains
  about 500 million neurons controlling GI motility, secretion,
  and blood flow independently of the CNS. Evaluate: (a) why the ENS is
  sometimes called autonomous (it functions after vagotomy); (b) the
  evidence linking ENS dysfunction to Parkinson's disease (\texorpdfstring{\ensuremath{\alpha}}{alpha}-synuclein
  aggregation begins in the ENS years before brain symptoms); (c) how
  the gut-brain axis connects the ENS to cognition and mood.
\end{enumerate}

\begin{enumerate}
\def\labelenumi{\arabic{enumi}.}
\setcounter{enumi}{28}
\tightlist
\item
  Optogenetics uses light-activated ion channels (channelrhodopsin-2,
  halorhodopsin) to activate or inhibit specific neurons with
  millisecond precision. Evaluate: (a) the technical requirements (viral
  vector delivery of opsin gene, implanted fibre optic); (b) an
  experiment that optogenetics uniquely enables (e.g., activating and
  silencing the same neurons to confirm necessity and sufficiency for a
  behaviour); (c) the path to clinical translation (optogenetic gene
  therapy for retinitis pigmentosa).
\end{enumerate}

\begin{enumerate}
\def\labelenumi{\arabic{enumi}.}
\setcounter{enumi}{29}
\tightlist
\item
  Critically evaluate the ``connectome'' project --- the effort to map
  every synaptic connection in the human brain (about 86
  billion neurons, about 150 trillion synapses). Is the
  connectome sufficient to understand brain function, or is dynamic
  activity (electrophysiology, neuromodulation) equally important?
  Compare the completed \emph{C. elegans} connectome (302 neurons, fully
  mapped) to our understanding of worm behaviour --- does knowing the
  wiring diagram explain the animal's behaviour? Then apply the same
  evidence-boundary reasoning to a high-performance speech
  neuroprosthesis reported in one participant with paralysis.
\end{enumerate}

\newpage

\section*{Questions --- Ch 30: Action Potentials and Synaptic
Transmission}\label{questions-ch-30-action-potentials-and-synaptic-transmission}
\addcontentsline{toc}{section}{Questions --- Ch 30: Action Potentials
and Synaptic Transmission}

\label{sec:q_unit_IX_action_potential_synapses}

\subsection*{Instructor Use and Coverage
Notes}\label{instructor-use-and-coverage-notes-32}
\addcontentsline{toc}{subsection}{Instructor Use and Coverage Notes}

\begin{itemize}
\tightlist
\item
  \textbf{Coverage target:} Predict how changing an ion gradient or
  channel conductance alters excitability.
\item
  \textbf{Model/data emphasis:} Nernst/Goldman, Hodgkin-Huxley, and
  synaptic-current reasoning.
\item
  \textbf{Assessment alignment:} Visual Representations, Statistical
  Tests and Data Analysis, Argumentation.
\item
  \textbf{Misconception probe:} An action potential is not electricity
  flowing like a wire; it is regenerated ion-channel dynamics.
\item
  \textbf{Transfer product:} Transfer excitability reasoning to
  anesthesia, toxins, epilepsy, and neuromuscular disease.
\item
  \textbf{Grading focus:} award full credit for mechanism, evidence,
  boundary conditions, and units when a calculation is required.
\item
  \textbf{Suggested use:} draw one recall item, one application item,
  and one synthesis item when building a short quiz from this bank.
\end{itemize}

\subsection*{Questions 1--10: Recall and
Comprehension}\label{questions-110-recall-and-comprehension-29}
\addcontentsline{toc}{subsection}{Questions 1--10: Recall and
Comprehension}

\emph{This activity accompanies
\cref{sec:unit_IX_action_potential_synapses} of the textbook --- review
that chapter before attempting the exercises below.}

\begin{enumerate}
\def\labelenumi{\arabic{enumi}.}
\tightlist
\item
  Define an action potential. What are the threshold, peak, and
  undershoot voltages?
\end{enumerate}

\begin{enumerate}
\def\labelenumi{\arabic{enumi}.}
\setcounter{enumi}{1}
\tightlist
\item
  What are voltage-gated Na\texorpdfstring{\textsuperscript{+}}{+} channels? Describe their three states
  (closed, open, inactivated).
\end{enumerate}

\begin{enumerate}
\def\labelenumi{\arabic{enumi}.}
\setcounter{enumi}{2}
\tightlist
\item
  What is the absolute refractory period? What is the relative
  refractory period?
\end{enumerate}

\begin{enumerate}
\def\labelenumi{\arabic{enumi}.}
\setcounter{enumi}{3}
\tightlist
\item
  What is a compound action potential (CAP)? How does it differ from a
  single-neuron AP?
\end{enumerate}

\begin{enumerate}
\def\labelenumi{\arabic{enumi}.}
\setcounter{enumi}{4}
\tightlist
\item
  Define a synapse. Name the pre- and postsynaptic components.
\end{enumerate}

\begin{enumerate}
\def\labelenumi{\arabic{enumi}.}
\setcounter{enumi}{5}
\tightlist
\item
  What is a neurotransmitter vesicle? How is it released (Ca\texorpdfstring{\textsuperscript{2}}{2}\texorpdfstring{\textsuperscript{+}}{+}-dependent
  exocytosis)?
\end{enumerate}

\begin{enumerate}
\def\labelenumi{\arabic{enumi}.}
\setcounter{enumi}{6}
\tightlist
\item
  What is an EPSP? What is an IPSP? How do they differ in ionic
  mechanism?
\end{enumerate}

\begin{enumerate}
\def\labelenumi{\arabic{enumi}.}
\setcounter{enumi}{7}
\tightlist
\item
  Define temporal and spatial summation of postsynaptic potentials.
\end{enumerate}

\begin{enumerate}
\def\labelenumi{\arabic{enumi}.}
\setcounter{enumi}{8}
\tightlist
\item
  What is a neuromuscular junction? What neurotransmitter is released
  there?
\end{enumerate}

\begin{enumerate}
\def\labelenumi{\arabic{enumi}.}
\setcounter{enumi}{9}
\tightlist
\item
  What is long-term potentiation (LTP)? In which brain region was it
  first described?
\end{enumerate}

\subsection*{Questions 11--20: Application and
Analysis}\label{questions-1120-application-and-analysis-29}
\addcontentsline{toc}{subsection}{Questions 11--20: Application and
Analysis}

\begin{enumerate}
\def\labelenumi{\arabic{enumi}.}
\setcounter{enumi}{10}
\tightlist
\item
  Using the Nernst equation (E = (RT/zF)
  ln({[}ion{]}\_out/{[}ion{]}\_in)), calculate the equilibrium potential
  for K\texorpdfstring{\textsuperscript{+}}{+} given {[}K\texorpdfstring{\textsuperscript{+}}{+}{]}\_out = 5 mM and {[}K\texorpdfstring{\textsuperscript{+}}{+}{]}\_in = 150 mM at 37°C.
  Then calculate E\_Na with {[}Na\texorpdfstring{\textsuperscript{+}}{+}{]}\_out = 145 mM and {[}Na\texorpdfstring{\textsuperscript{+}}{+}{]}\_in =
  15 mM. Which equilibrium potential is closest to the resting membrane
  potential?
\end{enumerate}

\begin{enumerate}
\def\labelenumi{\arabic{enumi}.}
\setcounter{enumi}{11}
\tightlist
\item
  Tetrodotoxin (TTX, from pufferfish) blocks voltage-gated Na\texorpdfstring{\textsuperscript{+}}{+} channels.
  Tetraethylammonium (TEA) blocks voltage-gated K\texorpdfstring{\textsuperscript{+}}{+} channels. Predict the
  effects on the action potential waveform when: (a) TTX alone is
  applied; (b) TEA alone; (c) both together.
\end{enumerate}

\begin{enumerate}
\def\labelenumi{\arabic{enumi}.}
\setcounter{enumi}{12}
\tightlist
\item
  Quantal analysis of neurotransmitter release: at the neuromuscular
  junction, each synaptic vesicle contains about 10,000 ACh
  molecules (one ``quantum''). The miniature end-plate potential (mEPP)
  = 0.4 mV. A normal EPSP = 40 mV. Estimate the quantal content (m =
  EPSP/mEPP). What parameter does the quantal content depend on (n =
  number of release sites, p = release probability)?
\end{enumerate}

\begin{enumerate}
\def\labelenumi{\arabic{enumi}.}
\setcounter{enumi}{13}
\tightlist
\item
  NMDA receptors are both ligand-gated (glutamate) AND voltage-gated
  (Mg\texorpdfstring{\textsuperscript{2}}{2}\texorpdfstring{\textsuperscript{+}}{+} block removed at depolarised potentials). Explain: (a) why NMDA
  receptors function as ``coincidence detectors'' (both presynaptic
  glutamate and postsynaptic depolarisation are required); (b) why this
  property is considered the molecular basis for Hebbian learning
  (``neurons that fire together, wire together''); (c) the Ca\texorpdfstring{\textsuperscript{2}}{2}\texorpdfstring{\textsuperscript{+}}{+} signal
  that enters via NMDA receptors and triggers LTP.
\end{enumerate}

\begin{enumerate}
\def\labelenumi{\arabic{enumi}.}
\setcounter{enumi}{14}
\tightlist
\item
  Myasthenia gravis (MG) is an autoimmune disease where antibodies
  target nicotinic ACh receptors at the NMJ. Predict: (a) whether the
  EPSP amplitude at the NMJ is increased or decreased; (b) the clinical
  manifestation (fatigable muscle weakness); (c) why
  acetylcholinesterase inhibitors (pyridostigmine) temporarily improve
  symptoms.
\end{enumerate}

\begin{enumerate}
\def\labelenumi{\arabic{enumi}.}
\setcounter{enumi}{15}
\tightlist
\item
  GABAergic inhibition: GABA\_A receptors are Cl\texorpdfstring{\textsuperscript{-}}{-} channels. In adult
  neurons E\_Cl \texorpdfstring{\ensuremath{\approx}}{~} \texorpdfstring{\ensuremath{-}}{-}70 mV (inhibitory). In neonatal neurons E\_Cl \texorpdfstring{\ensuremath{\approx}}{~} \texorpdfstring{\ensuremath{-}}{-}40
  mV (excitatory because resting potential is about \texorpdfstring{\ensuremath{-}}{-}55 mV).
  Explain: (a) why GABA is excitatory in neonatal neurons; (b) the
  developmental switch (upregulation of KCC2 Cl\texorpdfstring{\textsuperscript{-}}{-} transporter); (c) why
  this has implications for neonatal seizure treatment.
\end{enumerate}

\begin{enumerate}
\def\labelenumi{\arabic{enumi}.}
\setcounter{enumi}{16}
\tightlist
\item
  Short-term plasticity: paired-pulse facilitation (PPF) occurs when the
  second of two closely-spaced stimuli evokes a larger EPSP. Explain the
  residual calcium hypothesis: Ca\texorpdfstring{\textsuperscript{2}}{2}\texorpdfstring{\textsuperscript{+}}{+} from the first stimulus remains and
  adds to Ca\texorpdfstring{\textsuperscript{2}}{2}\texorpdfstring{\textsuperscript{+}}{+} from the second stimulus, producing greater vesicle
  release. Predict: would PPF be increased or decreased by a Ca\texorpdfstring{\textsuperscript{2}}{2}\texorpdfstring{\textsuperscript{+}}{+}
  chelator (BAPTA)?
\end{enumerate}

\begin{enumerate}
\def\labelenumi{\arabic{enumi}.}
\setcounter{enumi}{17}
\tightlist
\item
  Excitotoxicity: excessive glutamate release (e.g., during stroke)
  causes neuronal death via NMDA receptor overactivation \texorpdfstring{\ensuremath{\to}}{->} excessive
  Ca\texorpdfstring{\textsuperscript{2}}{2}\texorpdfstring{\textsuperscript{+}}{+} influx \texorpdfstring{\ensuremath{\to}}{->} mitochondrial dysfunction \texorpdfstring{\ensuremath{\to}}{->} ROS production \texorpdfstring{\ensuremath{\to}}{->} cell death.
  Explain: (a) why NMDA-blocking drugs (memantine, at low doses) are
  neuroprotective in clinical trials; (b) why high-dose NMDA block
  causes psychotic symptoms (dissociative anaesthetics --- ketamine);
  (c) what the narrow therapeutic window implies for stroke
  intervention.
\end{enumerate}

\begin{enumerate}
\def\labelenumi{\arabic{enumi}.}
\setcounter{enumi}{18}
\tightlist
\item
  Synaptic vesicle recycling occurs via three modes: kiss-and-run
  (transient fusion pore, rapid), full collapse/clathrin-mediated
  endocytosis (5--20 s), and bulk endocytosis (during intense
  stimulation). Explain: (a) how each mode is suited to different firing
  rates; (b) what happens when recycling capacity is exceeded (synaptic
  depression); (c) the role of dynamin GTPase in clathrin-mediated
  endocytosis.
\end{enumerate}

\begin{enumerate}
\def\labelenumi{\arabic{enumi}.}
\setcounter{enumi}{19}
\tightlist
\item
  Long-term depression (LTD) at cerebellar parallel fibre-Purkinje cell
  synapses is triggered by climbing fibre co-activation \texorpdfstring{\ensuremath{\to}}{->} Ca\texorpdfstring{\textsuperscript{2}}{2}\texorpdfstring{\textsuperscript{+}}{+} + PKC \texorpdfstring{\ensuremath{\to}}{->}
  AMPAR internalisation. Explain: (a) how LTD differs from LTP
  mechanistically (removal vs insertion of receptors); (b) the role of
  LTD in motor learning (error correction signals from climbing fibres);
  (c) how cerebellar LTD produces adaptation in the vestibulo-ocular
  reflex.
\end{enumerate}

\subsection*{Questions 21--30: Synthesis and
Evaluation}\label{questions-2130-synthesis-and-evaluation-29}
\addcontentsline{toc}{subsection}{Questions 21--30: Synthesis and
Evaluation}

\begin{enumerate}
\def\labelenumi{\arabic{enumi}.}
\setcounter{enumi}{20}
\tightlist
\item
  Evaluate the Hodgkin-Huxley (HH) model of the action potential (1952,
  Nobel Prize): using voltage-clamp data, they described Na\texorpdfstring{\textsuperscript{+}}{+} and K\texorpdfstring{\textsuperscript{+}}{+}
  conductances as m\texorpdfstring{\textsuperscript{3}}{3}h and n\texorpdfstring{\textsuperscript{4}}{4} gating variables. Assess: (a) the key
  insight (ion-specific conductances with distinct kinetics); (b) the
  predictive power of the model (propagation velocity, frequency
  coding); (c) the limitations (no molecular detail, does not account
  for ion channel diversity, subthreshold oscillations).
\end{enumerate}

\begin{enumerate}
\def\labelenumi{\arabic{enumi}.}
\setcounter{enumi}{21}
\tightlist
\item
  Design an experiment to determine whether a specific synapse exhibits
  Hebbian LTP, anti-Hebbian LTP, or homeostatic plasticity. Describe:
  the stimulation protocols (high-frequency \texorpdfstring{\ensuremath{\to}}{->} Hebbian, low-frequency \texorpdfstring{\ensuremath{\to}}{->}
  homeostatic), the electrophysiological measurements (field EPSP
  slope), and the pharmacological tools (NMDA antagonist, AMPAR
  antagonist) needed to distinguish these forms.
\end{enumerate}

\begin{enumerate}
\def\labelenumi{\arabic{enumi}.}
\setcounter{enumi}{22}
\tightlist
\item
  Evaluate the ``synaptic tagging and capture'' hypothesis (Frey and
  Morris, 1997): early LTP is protein synthesis-independent (lasts
  about 1--2 h), while late LTP requires new protein synthesis.
  A weak stimulus sets a ``tag'' at activated synapses; the tag captures
  plasticity-related proteins (PRPs) synthesised in response to a strong
  stimulus at another input. Evaluate: (a) the evidence for this model;
  (b) how it explains input-specificity of LTP; (c) implications for
  memory consolidation.
\end{enumerate}

\begin{enumerate}
\def\labelenumi{\arabic{enumi}.}
\setcounter{enumi}{23}
\tightlist
\item
  Evaluate the role of astrocytic gliotransmission (release of
  glutamate, D-serine, ATP from astrocytes) in modulating synaptic
  plasticity. What evidence supports the ``tripartite synapse'' model,
  and what criticisms have been raised about the specificity of
  astrocytic Ca\texorpdfstring{\textsuperscript{2}}{2}\texorpdfstring{\textsuperscript{+}}{+} signals and the difficulty of distinguishing neuronal
  from glial glutamate release?
\end{enumerate}

\begin{enumerate}
\def\labelenumi{\arabic{enumi}.}
\setcounter{enumi}{24}
\tightlist
\item
  Channelopathies: mutations in ion channels cause neurological
  diseases. Evaluate three examples: (a) SCN1A (Nav1.1) mutations \texorpdfstring{\ensuremath{\to}}{->}
  Dravet syndrome (severe childhood epilepsy); (b) KCNQ2/3 mutations \texorpdfstring{\ensuremath{\to}}{->}
  benign familial neonatal seizures; (c) CACNA1A mutations \texorpdfstring{\ensuremath{\to}}{->} familial
  hemiplegic migraine. For each, describe the channel, the mutation's
  effect on neuronal excitability, and the therapeutic approach.
\end{enumerate}

\begin{enumerate}
\def\labelenumi{\arabic{enumi}.}
\setcounter{enumi}{25}
\tightlist
\item
  Evaluate the dopamine hypothesis of schizophrenia: (a) traditional
  version (excess D\texorpdfstring{\textsubscript{2}}{2} receptor signalling in mesolimbic pathway); (b)
  revised version (hypo-dopaminergia in prefrontal cortex,
  hyper-dopaminergia in striatum); (c) the glutamate hypothesis (NMDA
  receptor hypofunction). How do these hypotheses explain the positive
  symptoms (hallucinations, delusions) and negative symptoms (flat
  affect, cognitive deficits) of schizophrenia?
\end{enumerate}

\begin{enumerate}
\def\labelenumi{\arabic{enumi}.}
\setcounter{enumi}{26}
\tightlist
\item
  Spike-timing-dependent plasticity (STDP): if the presynaptic neuron
  fires before the postsynaptic neuron (within about 20 ms),
  the synapse is strengthened (LTP). If after, the synapse is weakened
  (LTD). Evaluate: (a) the molecular mechanism (STDP timing depends on
  NMDA receptor/Ca\texorpdfstring{\textsuperscript{2}}{2}\texorpdfstring{\textsuperscript{+}}{+} dynamics); (b) the computational significance (STDP
  allows precise temporal learning of sequences); (c) how STDP differs
  in inhibitory vs excitatory synapses.
\end{enumerate}

\begin{enumerate}
\def\labelenumi{\arabic{enumi}.}
\setcounter{enumi}{27}
\tightlist
\item
  Evaluate the molecular mechanism of presynaptic neurotransmitter
  release: Ca\texorpdfstring{\textsuperscript{2}}{2}\texorpdfstring{\textsuperscript{+}}{+} entry \texorpdfstring{\ensuremath{\to}}{->} synaptotagmin-1 (Ca\texorpdfstring{\textsuperscript{2}}{2}\texorpdfstring{\textsuperscript{+}}{+} sensor) \texorpdfstring{\ensuremath{\to}}{->} conformational
  change \texorpdfstring{\ensuremath{\to}}{->} interaction with SNARE complex \texorpdfstring{\ensuremath{\to}}{->} fusion pore opening \texorpdfstring{\ensuremath{\to}}{->}
  transmitter release. Specifically: why is synaptotagmin-1 the timing
  sensor (low Ca\texorpdfstring{\textsuperscript{2}}{2}\texorpdfstring{\textsuperscript{+}}{+} affinity, fast kinetics, about 200 \texorpdfstring{\ensuremath{\mu}}{mu}s from
  Ca\texorpdfstring{\textsuperscript{2}}{2}\texorpdfstring{\textsuperscript{+}}{+} binding to fusion), and what does the about 4th power
  Ca\texorpdfstring{\textsuperscript{2}}{2}\texorpdfstring{\textsuperscript{+}}{+} dependence of release tell us about the number of Ca\texorpdfstring{\textsuperscript{2}}{2}\texorpdfstring{\textsuperscript{+}}{+} ions
  needed to trigger a single vesicle fusion event?
\end{enumerate}

\begin{enumerate}
\def\labelenumi{\arabic{enumi}.}
\setcounter{enumi}{28}
\tightlist
\item
  Evaluate the clinical utility of transcranial magnetic stimulation
  (TMS) for modulating neural circuits non-invasively. How does TMS
  induce action potentials in cortical neurons, what stimulation
  patterns produce LTP-like (theta-burst) vs LTD-like (continuous
  low-frequency) effects, and what is the evidence for therapeutic
  efficacy in treatment-resistant depression?
\end{enumerate}

\begin{enumerate}
\def\labelenumi{\arabic{enumi}.}
\setcounter{enumi}{29}
\tightlist
\item
  Critically evaluate the concept of the ``synaptome'' --- the catalogue
  of synapse types and their molecular composition across the brain.
  Using proteomics of isolated synaptosomes, researchers have identified
  \textgreater1,000 proteins at excitatory synapses. Evaluate: (a) what
  the molecular diversity of synapses implies for circuit function; (b)
  whether the synaptome concept is more informative than the connectome
  for understanding brain computation; (c) how synaptome-level mutations
  (SHANK3, PSD-95) cause neurodevelopmental disorders.
\end{enumerate}

\newpage

\section*{Questions --- Ch 31: Endocrine and Immune
Systems}\label{questions-ch-31-endocrine-and-immune-systems}
\addcontentsline{toc}{section}{Questions --- Ch 31: Endocrine and Immune
Systems}

\label{sec:q_unit_IX_endocrine_and_immune}

\subsection*{Instructor Use and Coverage
Notes}\label{instructor-use-and-coverage-notes-33}
\addcontentsline{toc}{subsection}{Instructor Use and Coverage Notes}

\begin{itemize}
\tightlist
\item
  \textbf{Coverage target:} Predict the consequence of changing a
  hormone, receptor, antigen, or immune-cell population.
\item
  \textbf{Model/data emphasis:} Hormone feedback, dose-response, and
  immune-memory reasoning.
\item
  \textbf{Assessment alignment:} Visual Representations, Statistical
  Tests and Data Analysis, Argumentation.
\item
  \textbf{Misconception probe:} Immunity is not just attack;
  recognition, tolerance, memory, and regulation are equally central.
\item
  \textbf{Transfer product:} Transfer signaling and immunity reasoning
  to vaccination, autoimmunity, stress, and metabolism.
\item
  \textbf{Grading focus:} award full credit for mechanism, evidence,
  boundary conditions, and units when a calculation is required.
\item
  \textbf{Suggested use:} draw one recall item, one application item,
  and one synthesis item when building a short quiz from this bank.
\end{itemize}

\subsection*{Questions 1--10: Recall and
Comprehension}\label{questions-110-recall-and-comprehension-30}
\addcontentsline{toc}{subsection}{Questions 1--10: Recall and
Comprehension}

\emph{This activity accompanies \cref{sec:unit_IX_endocrine_and_immune}
of the textbook --- review that chapter before attempting the exercises
below.}

\begin{enumerate}
\def\labelenumi{\arabic{enumi}.}
\tightlist
\item
  What is the endocrine system? How does it differ from the nervous
  system in signalling speed and duration?
\end{enumerate}

\begin{enumerate}
\def\labelenumi{\arabic{enumi}.}
\setcounter{enumi}{1}
\tightlist
\item
  Name the major endocrine glands and one hormone secreted by each.
\end{enumerate}

\begin{enumerate}
\def\labelenumi{\arabic{enumi}.}
\setcounter{enumi}{2}
\tightlist
\item
  Distinguish peptide hormones from steroid hormones in terms of
  synthesis, transport, and receptor location.
\end{enumerate}

\begin{enumerate}
\def\labelenumi{\arabic{enumi}.}
\setcounter{enumi}{3}
\tightlist
\item
  What is the hypothalamic-pituitary axis? Name one releasing hormone
  and one tropic hormone.
\end{enumerate}

\begin{enumerate}
\def\labelenumi{\arabic{enumi}.}
\setcounter{enumi}{4}
\tightlist
\item
  What is negative feedback? Give an example involving the thyroid
  gland.
\end{enumerate}

\begin{enumerate}
\def\labelenumi{\arabic{enumi}.}
\setcounter{enumi}{5}
\tightlist
\item
  What is cortisol? Name its target tissues and three physiological
  effects.
\end{enumerate}

\begin{enumerate}
\def\labelenumi{\arabic{enumi}.}
\setcounter{enumi}{6}
\tightlist
\item
  What are the components of innate immunity? Name four.
\end{enumerate}

\begin{enumerate}
\def\labelenumi{\arabic{enumi}.}
\setcounter{enumi}{7}
\tightlist
\item
  What is a Toll-like receptor (TLR)? Name a ligand for TLR4.
\end{enumerate}

\begin{enumerate}
\def\labelenumi{\arabic{enumi}.}
\setcounter{enumi}{8}
\tightlist
\item
  What is the complement system? What are its three outcomes
  (opsonisation, MAC, inflammation)?
\end{enumerate}

\begin{enumerate}
\def\labelenumi{\arabic{enumi}.}
\setcounter{enumi}{9}
\tightlist
\item
  Define phagocytosis. Name two professional phagocytes.
\end{enumerate}

\subsection*{Questions 11--20: Application and
Analysis}\label{questions-1120-application-and-analysis-30}
\addcontentsline{toc}{subsection}{Questions 11--20: Application and
Analysis}

\begin{enumerate}
\def\labelenumi{\arabic{enumi}.}
\setcounter{enumi}{10}
\tightlist
\item
  The HPA axis: CRH (hypothalamus) \texorpdfstring{\ensuremath{\to}}{->} ACTH (anterior pituitary) \texorpdfstring{\ensuremath{\to}}{->}
  cortisol (adrenal cortex). Cortisol feeds back negatively on CRH and
  ACTH secretion. If a patient has an adrenal tumour autonomously
  producing cortisol, predict: (a) ACTH levels; (b) CRH levels; (c)
  contralateral adrenal gland size; (d) response to the dexamethasone
  suppression test.
\end{enumerate}

\begin{enumerate}
\def\labelenumi{\arabic{enumi}.}
\setcounter{enumi}{11}
\tightlist
\item
  Insulin (from pancreatic \texorpdfstring{\ensuremath{\beta}}{beta}-cells) and glucagon (from \texorpdfstring{\ensuremath{\alpha}}{alpha}-cells) have
  opposing metabolic effects. Compare: (a) the stimulus for secretion of
  each; (b) the target tissue responses (liver, muscle, adipose); (c)
  the blood glucose level at which each dominates. What happens in Type
  1 Diabetes (complete \texorpdfstring{\ensuremath{\beta}}{beta}-cell destruction)?
\end{enumerate}

\begin{enumerate}
\def\labelenumi{\arabic{enumi}.}
\setcounter{enumi}{12}
\tightlist
\item
  Thyroid hormones (T\texorpdfstring{\textsubscript{3}}{3}/T\texorpdfstring{\textsubscript{4}}{4}) are unusual peptide-derived hormones that act
  like steroids (entering cells, binding nuclear receptors). Explain:
  (a) why T\texorpdfstring{\textsubscript{3}}{3}/T\texorpdfstring{\textsubscript{4}}{4} can cross membranes (lipophilic due to iodination); (b)
  how T\texorpdfstring{\textsubscript{3}}{3} binds thyroid hormone receptors (TR\texorpdfstring{\ensuremath{\alpha}}{alpha}/TR\texorpdfstring{\ensuremath{\beta}}{beta}) to regulate gene
  transcription; (c) why hypothyroidism causes bradycardia, fatigue, and
  weight gain.
\end{enumerate}

\begin{enumerate}
\def\labelenumi{\arabic{enumi}.}
\setcounter{enumi}{13}
\tightlist
\item
  The inflammasome (NLRP3) is activated by diverse danger signals (K\texorpdfstring{\textsuperscript{+}}{+}
  efflux, ROS, crystal particles). Trace: (a) NLRP3 oligomerisation; (b)
  ASC adaptor recruitment; (c) caspase-1 activation; (d) IL-1\texorpdfstring{\ensuremath{\beta}}{beta} and IL-18
  maturation and release; (e) pyroptosis (inflammatory cell death). Name
  two diseases where NLRP3 is inappropriately activated.
\end{enumerate}

\begin{enumerate}
\def\labelenumi{\arabic{enumi}.}
\setcounter{enumi}{14}
\tightlist
\item
  Neutrophils are the first responders to bacterial infection. Describe
  their killing mechanisms: (a) phagocytosis and phagolysosome fusion;
  (b) respiratory burst (NADPH oxidase \texorpdfstring{\ensuremath{\to}}{->} superoxide \texorpdfstring{\ensuremath{\to}}{->} HOCl via
  myeloperoxidase); (c) NETosis (neutrophil extracellular trap
  formation). What disease results from NADPH oxidase deficiency (CGD)?
\end{enumerate}

\begin{enumerate}
\def\labelenumi{\arabic{enumi}.}
\setcounter{enumi}{15}
\tightlist
\item
  Adrenal medulla catecholamines (epinephrine, norepinephrine) are
  released in response to sympathetic stimulation. Compare: (a) the
  receptor subtypes activated by each (\texorpdfstring{\ensuremath{\alpha}}{alpha}\texorpdfstring{\textsubscript{1}}{1}, \texorpdfstring{\ensuremath{\alpha}}{alpha}\texorpdfstring{\textsubscript{2}}{2}, \texorpdfstring{\ensuremath{\beta}}{beta}\texorpdfstring{\textsubscript{1}}{1}, \texorpdfstring{\ensuremath{\beta}}{beta}\texorpdfstring{\textsubscript{2}}{2}); (b) the effects
  on heart rate, blood vessels, and bronchioles; (c) why
  phaeochromocytoma (adrenal medullary tumour) causes paroxysmal
  hypertension.
\end{enumerate}

\begin{enumerate}
\def\labelenumi{\arabic{enumi}.}
\setcounter{enumi}{16}
\tightlist
\item
  Growth hormone (GH) from the anterior pituitary has both direct and
  indirect (via IGF-1 from the liver) effects. Explain: (a) the GH pulse
  pattern (nocturnal secretion regulated by GHRH and somatostatin); (b)
  the metabolic effects of GH (lipolysis, anti-insulin, protein
  synthesis); (c) why acromegaly (GH excess in adults) causes the
  symptoms it does.
\end{enumerate}

\begin{enumerate}
\def\labelenumi{\arabic{enumi}.}
\setcounter{enumi}{17}
\tightlist
\item
  Dendritic cells (DCs) bridge innate and adaptive immunity. Explain:
  (a) how immature DCs in tissues sample antigens via macropinocytosis
  and TLR activation; (b) how DC maturation involves upregulation of MHC
  class II and co-stimulatory molecules (CD80/86); (c) how mature DCs
  migrate to lymph nodes to activate naïve T cells (antigen
  presentation). Why are DCs the most potent professional
  antigen-presenting cells?
\end{enumerate}

\begin{enumerate}
\def\labelenumi{\arabic{enumi}.}
\setcounter{enumi}{18}
\tightlist
\item
  Parathyroid hormone (PTH) and calcitonin regulate blood calcium. PTH
  increases Ca\texorpdfstring{\textsuperscript{2}}{2}\texorpdfstring{\textsuperscript{+}}{+} by: (a) stimulating osteoclasts (bone resorption); (b)
  increasing renal Ca\texorpdfstring{\textsuperscript{2}}{2}\texorpdfstring{\textsuperscript{+}}{+} reabsorption; (c) activating 1\texorpdfstring{\ensuremath{\alpha}}{alpha}-hydroxylase
  (vitamin D activation \texorpdfstring{\ensuremath{\to}}{->} intestinal Ca\texorpdfstring{\textsuperscript{2}}{2}\texorpdfstring{\textsuperscript{+}}{+} absorption). Calcitonin
  decreases Ca\texorpdfstring{\textsuperscript{2}}{2}\texorpdfstring{\textsuperscript{+}}{+} by inhibiting osteoclasts. What disease results from
  PTH excess (hyperparathyroidism)?
\end{enumerate}

\begin{enumerate}
\def\labelenumi{\arabic{enumi}.}
\setcounter{enumi}{19}
\tightlist
\item
  Natural killer (NK) cells use the ``missing self'' mechanism:
  inhibitory KIR receptors bind MHC class I on normal cells and suppress
  killing. Virus-infected cells and tumour cells often downregulate MHC
  I, releasing the NK cell from inhibition. Explain: (a) the activating
  receptors (NKG2D, recognising stress ligands MICA/MICB); (b) how the
  balance of activating vs inhibiting signals determines the NK cell
  response; (c) how cancer immunotherapy exploits NK cells (bispecific
  engagers).
\end{enumerate}

\subsection*{Questions 21--30: Synthesis and
Evaluation}\label{questions-2130-synthesis-and-evaluation-30}
\addcontentsline{toc}{subsection}{Questions 21--30: Synthesis and
Evaluation}

\begin{enumerate}
\def\labelenumi{\arabic{enumi}.}
\setcounter{enumi}{20}
\tightlist
\item
  Evaluate the integration of endocrine and immune systems: cortisol (a
  glucocorticoid) suppresses inflammation by: (a) inhibiting NF-\texorpdfstring{\ensuremath{\kappa}}{kappa}B
  transcription factor; (b) reducing prostaglandin synthesis (inhibiting
  PLA\texorpdfstring{\textsubscript{2}}{2} and COX-2); (c) inducing lymphocyte apoptosis. Why do critically
  ill ICU patients sometimes develop adrenal insufficiency, and why does
  physiological stress-dose hydrocortisone improve sepsis survival?
\end{enumerate}

\begin{enumerate}
\def\labelenumi{\arabic{enumi}.}
\setcounter{enumi}{21}
\tightlist
\item
  Evaluate the molecular mechanisms of endocrine disruptors (BPA,
  phthalates, pesticides): how do these chemicals interfere with hormone
  signalling (estrogenic, anti-androgenic, thyroid-disrupting), at what
  environmental concentrations are effects seen (non-monotonic
  dose-response), and what evidence from wildlife populations
  (feminisation of male fish, reproductive abnormalities in alligators)
  supports endocrine disruption in the environment?
\end{enumerate}

\begin{enumerate}
\def\labelenumi{\arabic{enumi}.}
\setcounter{enumi}{22}
\tightlist
\item
  Design an experiment to determine whether a newly discovered peptide
  is a hormone: (a) demonstrate that it is secreted into the bloodstream
  in response to a physiological stimulus; (b) show that it has a target
  tissue with specific receptors; (c) prove that removing the source
  gland eliminates the response and that hormone replacement restores
  it; (d) show that the response is dose-dependent.
\end{enumerate}

\begin{enumerate}
\def\labelenumi{\arabic{enumi}.}
\setcounter{enumi}{23}
\tightlist
\item
  Evaluate the concept of ``trained immunity'' --- epigenetic
  reprogramming of innate immune cells (monocytes/macrophages) by prior
  pathogen exposure, resulting in enhanced responses to subsequent
  unrelated infections. How does BCG vaccination induce trained immunity
  (H3K4me3 at inflammatory gene promoters, mTOR-dependent metabolic
  shift from OXPHOS to glycolysis), and what are the implications for
  broad-spectrum vaccine development?
\end{enumerate}

\begin{enumerate}
\def\labelenumi{\arabic{enumi}.}
\setcounter{enumi}{24}
\tightlist
\item
  Evaluate the evolutionary conservation and divergence of the HPA axis:
  fish have an HPI axis (hypothalamic-pituitary-interrenal), amphibians
  have an HPA axis with cortisol, and birds have corticosterone. What
  does this conservation tell us about the fundamental biological
  requirement for glucocorticoid stress response across vertebrates?
\end{enumerate}

\begin{enumerate}
\def\labelenumi{\arabic{enumi}.}
\setcounter{enumi}{25}
\tightlist
\item
  Evaluate the therapeutic potential of checkpoint immunotherapy
  combined with endocrine modulation: tumour-associated macrophages
  (TAMs) are often immunosuppressive (M2-polarised). Propose a strategy
  combining: (a) CSF1R inhibitor to repolarise TAMs; (b) anti-PD-1 to
  release T cell checkpoint brakes; (c) low-dose dexamethasone to reduce
  treatment-related inflammation without suppressing anti-tumour
  immunity. What preclinical evidence supports this combination?
\end{enumerate}

\begin{enumerate}
\def\labelenumi{\arabic{enumi}.}
\setcounter{enumi}{26}
\tightlist
\item
  The pineal gland produces melatonin in response to darkness (via the
  suprachiasmatic nucleus \texorpdfstring{\ensuremath{\to}}{->} sympathetic \texorpdfstring{\ensuremath{\to}}{->} pinealocyte). Evaluate: (a)
  the circadian rhythm of melatonin secretion; (b) the evidence for
  melatonin's role in sleep-wake regulation vs a ``darkness signal'';
  (c) the efficacy of exogenous melatonin for jet lag and shift work
  disorder; (d) the immunomodulatory effects of melatonin
  (anti-inflammatory, antioxidant, NK cell activation).
\end{enumerate}

\begin{enumerate}
\def\labelenumi{\arabic{enumi}.}
\setcounter{enumi}{27}
\tightlist
\item
  Evaluate autoinflammatory diseases (distinct from autoimmune):
  conditions caused by inappropriate innate immune activation without
  autoantibodies. Using familial Mediterranean fever (FMF, pyrin
  mutations) as a case study: (a) how does mutant pyrin activate the
  inflammasome; (b) why does colchicine (inhibiting neutrophil
  migration) work as treatment; (c) how does IL-1 blockade (anakinra,
  canakinumab) provide definitive treatment?
\end{enumerate}

\begin{enumerate}
\def\labelenumi{\arabic{enumi}.}
\setcounter{enumi}{28}
\tightlist
\item
  The gut-brain-immune axis: enteroendocrine cells in the gut produce
  \textgreater30 hormones (GLP-1, PYY, CCK, serotonin) that signal to
  the brain and immune system. Evaluate: (a) how gut serotonin
  (\textgreater90\% of body's serotonin) modulates gut motility and
  immune cell function; (b) the evidence linking gut endocrine
  dysfunction to IBS (irritable bowel syndrome); (c) how GLP-1 receptor
  agonists (semaglutide) affect not just glucose metabolism but also
  neuroinflammation.
\end{enumerate}

\begin{enumerate}
\def\labelenumi{\arabic{enumi}.}
\setcounter{enumi}{29}
\tightlist
\item
  Critically evaluate the concept of ``immunometabolism'' --- the idea
  that metabolic reprogramming of immune cells determines their
  functional polarisation. Using the M1 (pro-inflammatory, glycolytic)
  vs M2 (anti-inflammatory, OXPHOS) macrophage paradigm: (a) how does
  TCA cycle disruption in M1 macrophages lead to itaconate and succinate
  accumulation; (b) how does succinate stabilise HIF-1\texorpdfstring{\ensuremath{\alpha}}{alpha} and promote
  IL-1\texorpdfstring{\ensuremath{\beta}}{beta} production; (c) why targeting immunometabolism (e.g., dimethyl
  fumarate activating Nrf2) is a therapeutic frontier? Then apply the
  chapter's evidence-map habit to reproductive physiology: what evidence
  separates ovulation timing, sperm production, tubal transport, uterine
  receptivity, endocrine feedback, implantation, and age-dependent
  gamete quality after 12 months without conception?
\end{enumerate}

\newpage

\section*{Questions --- Ch 32: Population Ecology and Growth
Models}\label{questions-ch-32-population-ecology-and-growth-models}
\addcontentsline{toc}{section}{Questions --- Ch 32: Population Ecology
and Growth Models}

\label{sec:q_unit_X_population_ecology}

\subsection*{Instructor Use and Coverage
Notes}\label{instructor-use-and-coverage-notes-34}
\addcontentsline{toc}{subsection}{Instructor Use and Coverage Notes}

\begin{itemize}
\tightlist
\item
  \textbf{Coverage target:} Calculate a growth or sampling parameter and
  interpret its conservation meaning.
\item
  \textbf{Model/data emphasis:} Exponential/logistic growth,
  mark-recapture, and matrix projection.
\item
  \textbf{Assessment alignment:} Representing and Describing Data,
  Statistical Tests and Data Analysis, Argumentation.
\item
  \textbf{Misconception probe:} Carrying capacity is not a fixed magic
  number; it changes with resources, interactions, and disturbance.
\item
  \textbf{Transfer product:} Transfer population models to fisheries,
  invasive species, epidemiology, and endangered species.
\item
  \textbf{Grading focus:} award full credit for mechanism, evidence,
  boundary conditions, and units when a calculation is required.
\item
  \textbf{Suggested use:} draw one recall item, one application item,
  and one synthesis item when building a short quiz from this bank.
\end{itemize}

\subsection*{Questions 1--10: Recall and
Comprehension}\label{questions-110-recall-and-comprehension-31}
\addcontentsline{toc}{subsection}{Questions 1--10: Recall and
Comprehension}

\emph{This activity accompanies \cref{sec:unit_X_population_ecology} of
the textbook --- review that chapter before attempting the exercises
below.}

\begin{enumerate}
\def\labelenumi{\arabic{enumi}.}
\tightlist
\item
  Define population ecology. What are the four demographic parameters
  that influence population size?
\end{enumerate}

\begin{enumerate}
\def\labelenumi{\arabic{enumi}.}
\setcounter{enumi}{1}
\tightlist
\item
  What is the exponential growth model? Write its equation.
\end{enumerate}

\begin{enumerate}
\def\labelenumi{\arabic{enumi}.}
\setcounter{enumi}{2}
\tightlist
\item
  What is the logistic growth model? Write its equation and define K
  (carrying capacity).
\end{enumerate}

\begin{enumerate}
\def\labelenumi{\arabic{enumi}.}
\setcounter{enumi}{3}
\tightlist
\item
  What are r-selected and K-selected species? Give one example of each.
\end{enumerate}

\begin{enumerate}
\def\labelenumi{\arabic{enumi}.}
\setcounter{enumi}{4}
\tightlist
\item
  Define survivorship curve. Describe Types I, II, and III.
\end{enumerate}

\begin{enumerate}
\def\labelenumi{\arabic{enumi}.}
\setcounter{enumi}{5}
\tightlist
\item
  What is the mark-recapture method? Write the Lincoln-Petersen
  equation.
\end{enumerate}

\begin{enumerate}
\def\labelenumi{\arabic{enumi}.}
\setcounter{enumi}{6}
\tightlist
\item
  What is population density? Distinguish crude density from ecological
  density.
\end{enumerate}

\begin{enumerate}
\def\labelenumi{\arabic{enumi}.}
\setcounter{enumi}{7}
\tightlist
\item
  Define natality, mortality, immigration, and emigration.
\end{enumerate}

\begin{enumerate}
\def\labelenumi{\arabic{enumi}.}
\setcounter{enumi}{8}
\tightlist
\item
  What is a life table? What parameters does it include?
\end{enumerate}

\begin{enumerate}
\def\labelenumi{\arabic{enumi}.}
\setcounter{enumi}{9}
\tightlist
\item
  What is the difference between density-dependent and
  density-independent limiting factors?
\end{enumerate}

\subsection*{Questions 11--20: Application and
Analysis}\label{questions-1120-application-and-analysis-31}
\addcontentsline{toc}{subsection}{Questions 11--20: Application and
Analysis}

\begin{enumerate}
\def\labelenumi{\arabic{enumi}.}
\setcounter{enumi}{10}
\tightlist
\item
  An E. coli population starts with N\texorpdfstring{\textsubscript{0}}{0} = 100 cells. The intrinsic rate
  of increase r = 0.693 h\texorpdfstring{\textsuperscript{-}}{-}\texorpdfstring{\textsuperscript{1}}{1} (doubling time = 1 hour). Calculate N\_t at
  t = 5 hours using N\_t = N\texorpdfstring{\textsubscript{0}}{0}e\^{}(rt). At what time would the
  population reach 10\texorpdfstring{\textsuperscript{6}}{6} cells?
\end{enumerate}

\begin{enumerate}
\def\labelenumi{\arabic{enumi}.}
\setcounter{enumi}{11}
\tightlist
\item
  A deer population (K = 1,000) starts at N = 50 and grows logistically
  with r = 0.3/year. Calculate dN/dt at N = 50, N = 500, and N = 900. At
  which population size is growth rate highest? What happens to growth
  rate as N approaches K?
\end{enumerate}

\begin{enumerate}
\def\labelenumi{\arabic{enumi}.}
\setcounter{enumi}{12}
\tightlist
\item
  In a mark-recapture study: 200 butterflies are captured, marked, and
  released. A week later, 250 butterflies are captured, of which 40 are
  marked. Estimate population size using the Lincoln-Petersen formula.
  What assumptions must hold for this estimate to be valid?
\end{enumerate}

\begin{enumerate}
\def\labelenumi{\arabic{enumi}.}
\setcounter{enumi}{13}
\tightlist
\item
  A population crash occurs when a population overshoots K (exceeds
  carrying capacity) and resources are depleted, causing K to decline.
  Using the logistic model, explain: (a) why the population exceeds K
  (time lag in density-dependent response); (b) the population
  trajectory after overshooting (damped oscillations vs extinction); (c)
  a real-world example (reindeer on St.~Matthew Island).
\end{enumerate}

\begin{enumerate}
\def\labelenumi{\arabic{enumi}.}
\setcounter{enumi}{14}
\tightlist
\item
  Compare r-selected (small, many offspring, high mortality, short
  lifespan) and K-selected (large, few offspring, low mortality, long
  lifespan) strategies. Place the following organisms on the r-K
  continuum: elephant, dandelion, oak tree, mosquito, whale, annual
  grass. Explain why the r/K model is now considered an
  oversimplification.
\end{enumerate}

\begin{enumerate}
\def\labelenumi{\arabic{enumi}.}
\setcounter{enumi}{15}
\tightlist
\item
  A Type I survivorship curve (e.g., humans) shows high survival until
  old age. A Type III curve (e.g., sea turtles) shows high juvenile
  mortality with high adult survival. Calculate: if a sea turtle lays
  100 eggs per year and 2\% survive to adulthood, what is the expected
  number of adult recruits per year? How many years of reproduction are
  needed to replace a dying adult (assume one-for-one replacement)?
\end{enumerate}

\begin{enumerate}
\def\labelenumi{\arabic{enumi}.}
\setcounter{enumi}{16}
\tightlist
\item
  The Allee effect describes a positive relationship between population
  density and per capita growth rate at low densities (below a
  threshold, the population cannot recover). Explain: (a) three
  mechanisms (mate finding difficulty, loss of cooperative defense,
  genetic problems from inbreeding); (b) how the Allee effect creates an
  extinction vortex for small populations; (c) what conservation
  strategies address the Allee effect.
\end{enumerate}

\begin{enumerate}
\def\labelenumi{\arabic{enumi}.}
\setcounter{enumi}{17}
\tightlist
\item
  Age-structured population models (Leslie matrix): a population with
  two age classes (juvenile and adult) has the following transition
  rates: juvenile survival = 0.5, adult survival = 0.8, adult fecundity
  = 3 juveniles/adult. Write the Leslie matrix and project the
  population for two time steps from an initial vector of {[}100
  juveniles, 50 adults{]}.
\end{enumerate}

\begin{enumerate}
\def\labelenumi{\arabic{enumi}.}
\setcounter{enumi}{18}
\tightlist
\item
  Source-sink metapopulation dynamics: a ``source'' habitat (birth rate
  \textgreater{} death rate) produces emigrants that maintain a ``sink''
  population (death rate \textgreater{} birth rate). Using a specific
  example (e.g., wetland birds using good-quality vs poor-quality
  wetlands), explain: (a) why the sink population would go extinct
  without immigration from the source; (b) what happens if the source
  habitat is destroyed; (c) implications for habitat conservation
  priority.
\end{enumerate}

\begin{enumerate}
\def\labelenumi{\arabic{enumi}.}
\setcounter{enumi}{19}
\tightlist
\item
  Climate change is shifting species' ranges poleward and upslope.
  Evaluate: (a) the evidence from range shift data (butterflies, birds,
  trees); (b) why some species cannot keep pace (dispersal limitation,
  habitat fragmentation); (c) how population viability analysis (PVA) is
  used to predict extinction risk under different warming scenarios.
\end{enumerate}

\subsection*{Questions 21--30: Synthesis and
Evaluation}\label{questions-2130-synthesis-and-evaluation-31}
\addcontentsline{toc}{subsection}{Questions 21--30: Synthesis and
Evaluation}

\begin{enumerate}
\def\labelenumi{\arabic{enumi}.}
\setcounter{enumi}{20}
\tightlist
\item
  Evaluate the applicability of the logistic growth model to real
  populations: where does it succeed (general S-shaped growth curves in
  laboratory conditions) and where does it fail (time lags,
  environmental stochasticity, age structure, Allee effects)? Propose a
  more realistic model that incorporates at least two of these omitted
  features.
\end{enumerate}

\begin{enumerate}
\def\labelenumi{\arabic{enumi}.}
\setcounter{enumi}{21}
\tightlist
\item
  Design a long-term population monitoring programme for a threatened
  species (e.g., African wild dog, \emph{Lycaon pictus}). Describe: (a)
  the sampling methodology (camera traps, aerialsurveys, genetic
  mark-recapture using faecal DNA); (b) the demographic parameters to
  estimate (\texorpdfstring{\ensuremath{\lambda}}{lambda} = finite rate of increase, survival by age class,
  reproductive output); (c) how you would build a PVA model and define
  minimum viable population (MVP) size.
\end{enumerate}

\begin{enumerate}
\def\labelenumi{\arabic{enumi}.}
\setcounter{enumi}{22}
\tightlist
\item
  Evaluate invasive species population dynamics: the brown tree snake
  (\emph{Boiga irregularis}) invaded Guam after WWII and caused the
  extinction of most native forest birds. Using population ecology
  theory: (a) what growth model best describes the initial invasion
  (exponential growth in the absence of natural enemies); (b) what new
  carrying capacity was reached; (c) what control strategies (traps,
  acetaminophen-laced dead mice) have been attempted, and why are they
  insufficient?
\end{enumerate}

\begin{enumerate}
\def\labelenumi{\arabic{enumi}.}
\setcounter{enumi}{23}
\tightlist
\item
  Evaluate the concept of ``ecological traps'' --- habitats that appear
  attractive to organisms but result in reduced fitness (e.g., sea
  turtles attracted to coastal lights, birds nesting in mowed
  agricultural fields). How do ecological traps arise (evolutionary
  mismatch between environmental cues and habitat quality), and how
  might they contribute to population decline in rapidly changing
  environments?
\end{enumerate}

\begin{enumerate}
\def\labelenumi{\arabic{enumi}.}
\setcounter{enumi}{24}
\tightlist
\item
  The maximum sustainable yield (MSY) of a fishery occurs at N = K/2 in
  the logistic model. Evaluate: (a) why harvesting at MSY maximises
  catch; (b) why MSY is inherently unstable (small perturbations can
  push the population toward collapse); (c) why modern fisheries
  management uses precautionary approaches (reference points below MSY,
  no-take zones) rather than attempting to maintain stocks at MSY.
\end{enumerate}

\begin{enumerate}
\def\labelenumi{\arabic{enumi}.}
\setcounter{enumi}{25}
\tightlist
\item
  Evaluate the role of demographic stochasticity vs environmental
  stochasticity in small population extinction risk. Using Monte Carlo
  simulations: (a) at what population size does demographic
  stochasticity dominate (N \textless{} 50); (b) how does environmental
  stochasticity affect populations of any size (droughts, disease
  outbreaks); (c) what does the combination of both stochastic forces
  imply for the ``50/500 rule'' in conservation genetics (Ne
  \textgreater{} 50 to avoid inbreeding, \textgreater{} 500 to maintain
  evolutionary potential)?
\end{enumerate}

\begin{enumerate}
\def\labelenumi{\arabic{enumi}.}
\setcounter{enumi}{26}
\tightlist
\item
  Human population grew exponentially from about 1 billion
  (1800) to about 8 billion (2024). Evaluate: (a) which
  technological advances removed carrying capacity limits (agriculture,
  medicine, energy); (b) the demographic transition model (stages I--IV)
  and where different countries are; (c) whether global population will
  stabilise (UN WPP 2024 median projection: peak near 10.3 billion in
  the mid-2080s, then about 10.2 billion by 2100) or decline,
  and the assumptions underlying different projections.
\end{enumerate}

\begin{enumerate}
\def\labelenumi{\arabic{enumi}.}
\setcounter{enumi}{27}
\tightlist
\item
  Evaluate the ``landscape of fear'' concept: predator presence alters
  prey behaviour (reduced foraging time, habitat avoidance), which
  reduces prey population growth rate even without direct predation.
  Using the wolf-elk-aspen trophic cascade in Yellowstone as a case
  study, evaluate the evidence that wolf reintroduction changed elk
  behaviour enough to allow aspen regeneration, and whether this
  qualifies as a ``landscape of fear'' effect or simply reduced elk
  density.
\end{enumerate}

\begin{enumerate}
\def\labelenumi{\arabic{enumi}.}
\setcounter{enumi}{28}
\tightlist
\item
  Evaluate the mathematical basis for chaos in ecological populations:
  discrete-time logistic models (N\_\{t+1\} = rN\_t(1 \texorpdfstring{\ensuremath{-}}{-} N\_t/K)) with
  high r (r \textgreater{} 2.57) produce chaotic dynamics where
  population size is deterministic but unpredictable. What does chaos
  imply for long-term population prediction, and how can evidence of
  chaos be detected in real population time series (e.g., flour beetles,
  \emph{Tribolium})?
\end{enumerate}

\begin{enumerate}
\def\labelenumi{\arabic{enumi}.}
\setcounter{enumi}{29}
\tightlist
\item
  Critically evaluate the concept of ``planetary boundaries'' (Stockholm
  Resilience Centre, 2009) in relation to human population ecology. Of
  the nine identified boundaries (climate change, biodiversity loss,
  nitrogen cycling, etc.), which have been transgressed, and how do
  population density, consumption patterns, and technology interact to
  determine whether human civilisation is operating within sustainable
  limits?
\end{enumerate}

\newpage

\section*{Questions --- Ch 33: Community Ecology and Species
Interactions}\label{questions-ch-33-community-ecology-and-species-interactions}
\addcontentsline{toc}{section}{Questions --- Ch 33: Community Ecology
and Species Interactions}

\label{sec:q_unit_X_community_ecology}

\subsection*{Instructor Use and Coverage
Notes}\label{instructor-use-and-coverage-notes-35}
\addcontentsline{toc}{subsection}{Instructor Use and Coverage Notes}

\begin{itemize}
\tightlist
\item
  \textbf{Coverage target:} Predict how changing one species or resource
  alters community structure.
\item
  \textbf{Model/data emphasis:} Lotka-Volterra-style interaction and
  diversity-pattern reasoning.
\item
  \textbf{Assessment alignment:} Representing and Describing Data,
  Statistical Tests and Data Analysis, Argumentation.
\item
  \textbf{Misconception probe:} A species interaction is not permanently
  good or bad; the sign can change with context.
\item
  \textbf{Transfer product:} Transfer interaction reasoning to
  restoration, disease ecology, agriculture, and invasion biology.
\item
  \textbf{Grading focus:} award full credit for mechanism, evidence,
  boundary conditions, and units when a calculation is required.
\item
  \textbf{Suggested use:} draw one recall item, one application item,
  and one synthesis item when building a short quiz from this bank.
\end{itemize}

\subsection*{Questions 1--10: Recall and
Comprehension}\label{questions-110-recall-and-comprehension-32}
\addcontentsline{toc}{subsection}{Questions 1--10: Recall and
Comprehension}

\emph{This activity accompanies \cref{sec:unit_X_community_ecology} of
the textbook --- review that chapter before attempting the exercises
below.}

\begin{enumerate}
\def\labelenumi{\arabic{enumi}.}
\tightlist
\item
  Define a biological community. How does it differ from a population?
\end{enumerate}

\begin{enumerate}
\def\labelenumi{\arabic{enumi}.}
\setcounter{enumi}{1}
\tightlist
\item
  Name five types of species interactions. Give one example of each.
\end{enumerate}

\begin{enumerate}
\def\labelenumi{\arabic{enumi}.}
\setcounter{enumi}{2}
\tightlist
\item
  What is interspecific competition? State the competitive exclusion
  principle (Gause's law).
\end{enumerate}

\begin{enumerate}
\def\labelenumi{\arabic{enumi}.}
\setcounter{enumi}{3}
\tightlist
\item
  What is a niche? Distinguish fundamental niche from realised niche.
\end{enumerate}

\begin{enumerate}
\def\labelenumi{\arabic{enumi}.}
\setcounter{enumi}{4}
\tightlist
\item
  What is predation? Name two anti-predator adaptations.
\end{enumerate}

\begin{enumerate}
\def\labelenumi{\arabic{enumi}.}
\setcounter{enumi}{5}
\tightlist
\item
  Define mutualism. Give one obligate and one facultative example.
\end{enumerate}

\begin{enumerate}
\def\labelenumi{\arabic{enumi}.}
\setcounter{enumi}{6}
\tightlist
\item
  What is a food web? How does it differ from a food chain?
\end{enumerate}

\begin{enumerate}
\def\labelenumi{\arabic{enumi}.}
\setcounter{enumi}{7}
\tightlist
\item
  What is a keystone species? Give one classic example.
\end{enumerate}

\begin{enumerate}
\def\labelenumi{\arabic{enumi}.}
\setcounter{enumi}{8}
\tightlist
\item
  Define ecological succession. Distinguish primary from secondary
  succession.
\end{enumerate}

\begin{enumerate}
\def\labelenumi{\arabic{enumi}.}
\setcounter{enumi}{9}
\tightlist
\item
  What is species richness? How does it differ from species diversity
  (Shannon index)?
\end{enumerate}

\subsection*{Questions 11--20: Application and
Analysis}\label{questions-1120-application-and-analysis-32}
\addcontentsline{toc}{subsection}{Questions 11--20: Application and
Analysis}

\begin{enumerate}
\def\labelenumi{\arabic{enumi}.}
\setcounter{enumi}{10}
\tightlist
\item
  The Lotka-Volterra competition model uses two equations: dN\texorpdfstring{\textsubscript{1}}{1}/dt =
  r\texorpdfstring{\textsubscript{1}}{1}N\texorpdfstring{\textsubscript{1}}{1}(K\texorpdfstring{\textsubscript{1}}{1} \texorpdfstring{\ensuremath{-}}{-} N\texorpdfstring{\textsubscript{1}}{1} \texorpdfstring{\ensuremath{-}}{-} \texorpdfstring{\ensuremath{\alpha}}{alpha}\texorpdfstring{\textsubscript{1}}{1}\texorpdfstring{\textsubscript{2}}{2}N\texorpdfstring{\textsubscript{2}}{2})/K\texorpdfstring{\textsubscript{1}}{1} and dN\texorpdfstring{\textsubscript{2}}{2}/dt = r\texorpdfstring{\textsubscript{2}}{2}N\texorpdfstring{\textsubscript{2}}{2}(K\texorpdfstring{\textsubscript{2}}{2} \texorpdfstring{\ensuremath{-}}{-} N\texorpdfstring{\textsubscript{2}}{2} \texorpdfstring{\ensuremath{-}}{-} \texorpdfstring{\ensuremath{\alpha}}{alpha}\texorpdfstring{\textsubscript{2}}{2}\texorpdfstring{\textsubscript{1}}{1}N\texorpdfstring{\textsubscript{1}}{1})/K\texorpdfstring{\textsubscript{2}}{2}. If K\texorpdfstring{\textsubscript{1}}{1}
  = 1,000, K\texorpdfstring{\textsubscript{2}}{2} = 800, \texorpdfstring{\ensuremath{\alpha}}{alpha}\texorpdfstring{\textsubscript{1}}{1}\texorpdfstring{\textsubscript{2}}{2} = 0.5, \texorpdfstring{\ensuremath{\alpha}}{alpha}\texorpdfstring{\textsubscript{2}}{2}\texorpdfstring{\textsubscript{1}}{1} = 1.2, predict the outcome of
  competition. Under what conditions do two species coexist (K\texorpdfstring{\textsubscript{1}}{1}/\texorpdfstring{\ensuremath{\alpha}}{alpha}\texorpdfstring{\textsubscript{1}}{1}\texorpdfstring{\textsubscript{2}}{2}
  \textgreater{} K\texorpdfstring{\textsubscript{2}}{2} AND K\texorpdfstring{\textsubscript{2}}{2}/\texorpdfstring{\ensuremath{\alpha}}{alpha}\texorpdfstring{\textsubscript{2}}{2}\texorpdfstring{\textsubscript{1}}{1} \textgreater{} K\texorpdfstring{\textsubscript{1}}{1})?
\end{enumerate}

\begin{enumerate}
\def\labelenumi{\arabic{enumi}.}
\setcounter{enumi}{11}
\tightlist
\item
  Paine's classic experiment (1966) removed the sea star \emph{Pisaster
  ochraceus} from intertidal rocks. The result was dominance by the
  mussel \emph{Mytilus californianus}, which outcompeted most other
  sessile species. Explain: (a) why \emph{Pisaster} is a keystone
  predator; (b) how its removal reduced species diversity; (c) the
  concept of ``keystone species'' and why biomass alone does not predict
  ecological importance.
\end{enumerate}

\begin{enumerate}
\def\labelenumi{\arabic{enumi}.}
\setcounter{enumi}{12}
\tightlist
\item
  Character displacement: Darwin's finches on islands where two species
  coexist show greater beak size divergence than on islands where they
  occur alone. Explain: (a) the mechanism (sympatric competition selects
  for niche divergence); (b) the evidence from \emph{Geospiza fortis}
  and \emph{G. fuliginosa} beak depth distributions; (c) how this
  relates to niche partitioning.
\end{enumerate}

\begin{enumerate}
\def\labelenumi{\arabic{enumi}.}
\setcounter{enumi}{13}
\tightlist
\item
  Trophic cascades: when wolves were reintroduced to Yellowstone, elk
  herbivory decreased, allowing aspen and willow regeneration, which
  stabilised stream banks and increased beaver populations. Trace this
  top-down trophic cascade: wolves \texorpdfstring{\ensuremath{\to}}{->} elk behaviour change \texorpdfstring{\ensuremath{\to}}{->} vegetation
  recovery \texorpdfstring{\ensuremath{\to}}{->} stream geomorphology change \texorpdfstring{\ensuremath{\to}}{->} beaver habitat expansion. Is
  this primarily a density-mediated or trait-mediated cascade?
\end{enumerate}

\begin{enumerate}
\def\labelenumi{\arabic{enumi}.}
\setcounter{enumi}{14}
\tightlist
\item
  Mycorrhizal networks (``wood-wide web'') connect trees underground,
  potentially allowing nutrient and carbon transfer between individuals.
  Evaluate the evidence: (a) \texorpdfstring{\textsuperscript{1}}{1}\texorpdfstring{\textsuperscript{3}}{3}C/\texorpdfstring{\textsuperscript{1}}{1}\texorpdfstring{\textsuperscript{5}}{5}N labelling experiments showing
  carbon transfer between Douglas fir and paper birch; (b) whether
  transfer is mutualistic (mother trees provisioning seedlings) or
  incidental (fungal self-interest); (c) the evidence from Simard's
  studies and the controversy around the ``mother tree'' narrative.
\end{enumerate}

\begin{enumerate}
\def\labelenumi{\arabic{enumi}.}
\setcounter{enumi}{15}
\tightlist
\item
  Island biogeography theory (MacArthur and Wilson, 1967): species
  richness on an island results from a balance between immigration
  (decreasing with distance from mainland) and extinction (increasing
  with island size decrease). Apply this to habitat islands: predict
  species richness on forest fragments of 1 ha, 10 ha, and 1,000 ha that
  are 5 km from a large continuous forest. What does the ``species-area
  relationship'' (\(S = cA^z\)) predict?
\end{enumerate}

\begin{enumerate}
\def\labelenumi{\arabic{enumi}.}
\setcounter{enumi}{16}
\tightlist
\item
  Herbivore-plant coevolution has produced elaborate chemical defence
  systems. The milkweed (\emph{Asclepias}) -- monarch butterfly
  (\emph{Danaus plexippus}) system illustrates escalation: (a) milkweed
  produces cardenolides (toxic Na\texorpdfstring{\textsuperscript{+}}{+}/K\texorpdfstring{\textsuperscript{+}}{+}-ATPase inhibitors); (b) monarch
  larvae sequester cardenolides for their own defence; (c) monarch
  Na\texorpdfstring{\textsuperscript{+}}{+}/K\texorpdfstring{\textsuperscript{+}}{+}-ATPase has evolved resistance mutations. Trace this
  coevolutionary arms race.
\end{enumerate}

\begin{enumerate}
\def\labelenumi{\arabic{enumi}.}
\setcounter{enumi}{17}
\tightlist
\item
  Primary succession on volcanic islands (e.g., Surtsey, Iceland, formed
  1963): pioneer species (lichens, mosses) \texorpdfstring{\ensuremath{\to}}{->} grasses \texorpdfstring{\ensuremath{\to}}{->} shrubs \texorpdfstring{\ensuremath{\to}}{->}
  woodland. Explain the facilitation model: how do early colonisers
  modify the environment (soil formation, nitrogen fixation) to enable
  later successional stages?
\end{enumerate}

\begin{enumerate}
\def\labelenumi{\arabic{enumi}.}
\setcounter{enumi}{18}
\tightlist
\item
  Calculate the Shannon diversity index (H' = \texorpdfstring{\ensuremath{-}}{-}\texorpdfstring{\ensuremath{\Sigma}}{Sigma}p\texorpdfstring{\textsubscript{i}}{i} ln p\texorpdfstring{\textsubscript{i}}{i}) for a
  community with 4 species at proportions: 0.50, 0.25, 0.15, 0.10. Then
  calculate for a community with 4 species at equal proportions (0.25
  each). Which community is more diverse? What does evenness contribute
  to diversity?
\end{enumerate}

\begin{enumerate}
\def\labelenumi{\arabic{enumi}.}
\setcounter{enumi}{19}
\tightlist
\item
  Parasitoids (e.g., parasitoid wasps) lay eggs inside or on host
  insects; the larva consumes the host from inside. Compare
  parasitoid-host dynamics to classical predator-prey dynamics: (a) why
  parasitoid search efficiency matters (Nicholson-Bailey model); (b) why
  parasitoid-host systems tend toward unstable oscillations in simple
  models; (c) what stabilising mechanisms exist in nature (spatial
  refugia, host immunity).
\end{enumerate}

\subsection*{Questions 21--30: Synthesis and
Evaluation}\label{questions-2130-synthesis-and-evaluation-32}
\addcontentsline{toc}{subsection}{Questions 21--30: Synthesis and
Evaluation}

\begin{enumerate}
\def\labelenumi{\arabic{enumi}.}
\setcounter{enumi}{20}
\tightlist
\item
  Evaluate the intermediate disturbance hypothesis (IDH, Connell 1978):
  species diversity is maximised at intermediate disturbance
  frequency/intensity. Evaluate: (a) the mechanism (low disturbance \texorpdfstring{\ensuremath{\to}}{->}
  competitive exclusion; high disturbance \texorpdfstring{\ensuremath{\to}}{->} primarily ruderals survive);
  (b) empirical support from coral reef and tropical forest studies; (c)
  criticisms (meta-analyses showing weak or inconsistent support for
  IDH).
\end{enumerate}

\begin{enumerate}
\def\labelenumi{\arabic{enumi}.}
\setcounter{enumi}{21}
\tightlist
\item
  Design an experiment to test whether a suspected keystone species is
  truly a keystone (disproportionate effect on community structure
  relative to its biomass). Describe: (a) the removal experiment design;
  (b) what community metrics (species richness, evenness, biomass
  distribution) you would measure before and after removal; (c) how you
  would distinguish keystone effects from dominant species effects.
\end{enumerate}

\begin{enumerate}
\def\labelenumi{\arabic{enumi}.}
\setcounter{enumi}{22}
\tightlist
\item
  Evaluate the ``biodiversity-ecosystem function'' (BEF) debate: does
  higher species diversity lead to greater ecosystem productivity? Cite:
  (a) the Cedar Creek grassland experiments (Tilman); (b) the ``sampling
  effect'' (more species \texorpdfstring{\ensuremath{\to}}{->} more likely to include a highly productive
  species) vs ``complementarity effect'' (niche partitioning \texorpdfstring{\ensuremath{\to}}{->} more
  complete resource use); (c) implications for conservation (``every
  species matters'' vs ``some species are redundant'').
\end{enumerate}

\begin{enumerate}
\def\labelenumi{\arabic{enumi}.}
\setcounter{enumi}{23}
\tightlist
\item
  Evaluate the role of facilitation (positive interactions) vs
  competition in structuring communities under environmental stress. The
  stress gradient hypothesis (SGH) predicts facilitation dominates in
  harsh environments and competition in benign ones. Evaluate the
  evidence from: (a) alpine cushion plants facilitating other species;
  (b) coral reef structure facilitating fish diversity; (c) exceptions
  where competition dominates even under stress.
\end{enumerate}

\begin{enumerate}
\def\labelenumi{\arabic{enumi}.}
\setcounter{enumi}{24}
\tightlist
\item
  Biological invasions: invasive species can restructure entire
  communities. Evaluate three hypotheses for why invasive species
  succeed: (a) enemy release hypothesis (no predators/parasites in new
  range); (b) novel weapons hypothesis (allelochemicals that native
  species cannot tolerate); (c) empty niche hypothesis (exploiting
  unused resources). For each, cite one case study.
\end{enumerate}

\begin{enumerate}
\def\labelenumi{\arabic{enumi}.}
\setcounter{enumi}{25}
\tightlist
\item
  Evaluate the concept of ``ecological networks'' beyond food webs:
  mutualistic networks (plant-pollinator, plant-seed disperser) show a
  nested structure where specialist species interact with a subset of
  the partners of generalist species. What does nestedness imply for:
  (a) network stability (resilience to species loss); (b) the order of
  species extinction under habitat loss; (c) the importance of highly
  connected ``hub'' species for network integrity?
\end{enumerate}

\begin{enumerate}
\def\labelenumi{\arabic{enumi}.}
\setcounter{enumi}{26}
\tightlist
\item
  Evaluate Hubbell's Neutral Theory of Biodiversity (2001): species are
  ecologically equivalent, and community composition is determined by
  stochastic birth, death, and migration rather than niche differences.
  Compare neutral theory predictions with niche-based models: (a) what
  patterns can neutral theory explain (species-abundance distributions,
  species-area relationships); (b) what it cannot (functional trait
  differences, competitive exclusion); (c) is neutral theory ``wrong but
  useful''?
\end{enumerate}

\begin{enumerate}
\def\labelenumi{\arabic{enumi}.}
\setcounter{enumi}{27}
\tightlist
\item
  Climate change is causing phenological mismatches between interacting
  species (e.g., caterpillar peak occurring before oak leaf-out, but
  after great tit breeding initiation). Evaluate: (a) the evidence from
  long-term phenological datasets; (b) whether species can adapt quickly
  enough (evolutionary rescue vs extinction); (c) what community-level
  consequences emerge from cascading mismatches across multiple trophic
  levels.
\end{enumerate}

\begin{enumerate}
\def\labelenumi{\arabic{enumi}.}
\setcounter{enumi}{28}
\tightlist
\item
  Evaluate the concept of ``alternative stable states'' in ecological
  communities: an ecosystem can exist in multiple stable configurations
  (e.g., clear-water lake vs turbid-water lake), and transitions between
  states (regime shifts) are driven by perturbations that push the
  system past a tipping point. Using the shallow lake eutrophication
  example: (a) what drives the regime shift; (b) why is the shift often
  irreversible (hysteresis); (c) what early warning signals predict an
  approaching tipping point?
\end{enumerate}

\begin{enumerate}
\def\labelenumi{\arabic{enumi}.}
\setcounter{enumi}{29}
\tightlist
\item
  Critically evaluate the use of ecological network analysis for
  predicting the cascading effects of species extinction. Using
  simulated extinction sequences (removing species from most to least
  connected): (a) how do co-extinction cascades propagate through food
  webs; (b) why do mutualistic networks (plant-pollinator) collapse more
  rapidly than antagonistic networks (predator-prey) under targeted
  extinction; (c) what does this imply for setting conservation
  priorities in species-rich ecosystems?
\end{enumerate}

\newpage

\section*{Questions --- Ch 34: Ecosystem Ecology and Biogeochemical
Cycles}\label{questions-ch-34-ecosystem-ecology-and-biogeochemical-cycles}
\addcontentsline{toc}{section}{Questions --- Ch 34: Ecosystem Ecology
and Biogeochemical Cycles}

\label{sec:q_unit_X_ecosystem_ecology}

\subsection*{Instructor Use and Coverage
Notes}\label{instructor-use-and-coverage-notes-36}
\addcontentsline{toc}{subsection}{Instructor Use and Coverage Notes}

\begin{itemize}
\tightlist
\item
  \textbf{Coverage target:} Calculate a flux or efficiency and explain
  which pool or process limits the system.
\item
  \textbf{Model/data emphasis:} Energy-flow, productivity, and
  nutrient-budget calculations.
\item
  \textbf{Assessment alignment:} Representing and Describing Data,
  Statistical Tests and Data Analysis, Argumentation.
\item
  \textbf{Misconception probe:} Energy flows through ecosystems, but
  matter cycles; confusing the two breaks many explanations.
\item
  \textbf{Transfer product:} Transfer ecosystem reasoning to
  eutrophication, carbon budgets, agriculture, and climate feedback.
\item
  \textbf{Grading focus:} award full credit for mechanism, evidence,
  boundary conditions, and units when a calculation is required.
\item
  \textbf{Suggested use:} draw one recall item, one application item,
  and one synthesis item when building a short quiz from this bank.
\end{itemize}

\subsection*{Questions 1--10: Recall and
Comprehension}\label{questions-110-recall-and-comprehension-33}
\addcontentsline{toc}{subsection}{Questions 1--10: Recall and
Comprehension}

\emph{This activity accompanies \cref{sec:unit_X_ecosystem_ecology} of
the textbook --- review that chapter before attempting the exercises
below.}

\begin{enumerate}
\def\labelenumi{\arabic{enumi}.}
\tightlist
\item
  Define an ecosystem. What are its biotic and abiotic components?
\end{enumerate}

\begin{enumerate}
\def\labelenumi{\arabic{enumi}.}
\setcounter{enumi}{1}
\tightlist
\item
  What is primary productivity? Distinguish GPP, NPP, and NEP.
\end{enumerate}

\begin{enumerate}
\def\labelenumi{\arabic{enumi}.}
\setcounter{enumi}{2}
\tightlist
\item
  What is a trophic level? Name four trophic levels in a food chain.
\end{enumerate}

\begin{enumerate}
\def\labelenumi{\arabic{enumi}.}
\setcounter{enumi}{3}
\tightlist
\item
  What is ecological efficiency? What is the typical transfer efficiency
  between trophic levels?
\end{enumerate}

\begin{enumerate}
\def\labelenumi{\arabic{enumi}.}
\setcounter{enumi}{4}
\tightlist
\item
  Define the carbon cycle. Name the major carbon reservoirs on Earth.
\end{enumerate}

\begin{enumerate}
\def\labelenumi{\arabic{enumi}.}
\setcounter{enumi}{5}
\tightlist
\item
  What is the nitrogen cycle? Name the four major transformations.
\end{enumerate}

\begin{enumerate}
\def\labelenumi{\arabic{enumi}.}
\setcounter{enumi}{6}
\tightlist
\item
  What is the phosphorus cycle? How does it differ from the nitrogen and
  carbon cycles?
\end{enumerate}

\begin{enumerate}
\def\labelenumi{\arabic{enumi}.}
\setcounter{enumi}{7}
\tightlist
\item
  What is the water cycle (hydrological cycle)? Name the major fluxes.
\end{enumerate}

\begin{enumerate}
\def\labelenumi{\arabic{enumi}.}
\setcounter{enumi}{8}
\tightlist
\item
  What is eutrophication? What nutrients cause it and why?
\end{enumerate}

\begin{enumerate}
\def\labelenumi{\arabic{enumi}.}
\setcounter{enumi}{9}
\tightlist
\item
  Define decomposition. What organisms are primary decomposers?
\end{enumerate}

\subsection*{Questions 11--20: Application and
Analysis}\label{questions-1120-application-and-analysis-33}
\addcontentsline{toc}{subsection}{Questions 11--20: Application and
Analysis}

\begin{enumerate}
\def\labelenumi{\arabic{enumi}.}
\setcounter{enumi}{10}
\tightlist
\item
  In a grassland ecosystem, GPP = 5,000 kcal/m\texorpdfstring{\textsuperscript{2}}{2}/yr and plant respiration
  = 3,000 kcal/m\texorpdfstring{\textsuperscript{2}}{2}/yr. Calculate NPP. If herbivores consume 20\% of NPP
  and their assimilation efficiency is 40\%, calculate secondary
  production at the herbivore level. How much energy reaches the top
  predator (3rd consumer) assuming 10\% transfer efficiency?
\end{enumerate}

\begin{enumerate}
\def\labelenumi{\arabic{enumi}.}
\setcounter{enumi}{11}
\tightlist
\item
  The inverted biomass pyramid of the open ocean: phytoplankton biomass
  is much lower than zooplankton biomass at any instant, yet
  phytoplankton NPP exceeds zooplankton production. Explain this paradox
  using turnover rate (phytoplankton divide daily; zooplankton live
  weeks to months).
\end{enumerate}

\begin{enumerate}
\def\labelenumi{\arabic{enumi}.}
\setcounter{enumi}{12}
\tightlist
\item
  Nitrogen fixation by \emph{Rhizobium} in legume root nodules adds
  about 200 kg N/ha/yr to agricultural soils. Haber-Bosch
  industrial fixation produces about 150 Tg N/yr globally.
  Compare: (a) the energy cost of each; (b) the environmental
  consequences of excess reactive nitrogen (N\texorpdfstring{\textsubscript{2}}{2}O greenhouse gas, aquatic
  eutrophication, acid rain); (c) why reducing fertiliser overuse is
  critical for planetary health.
\end{enumerate}

\begin{enumerate}
\def\labelenumi{\arabic{enumi}.}
\setcounter{enumi}{13}
\tightlist
\item
  The biological pump in the ocean: surface phytoplankton fix CO\texorpdfstring{\textsubscript{2}}{2} \texorpdfstring{\ensuremath{\to}}{->}
  organic particles sink \texorpdfstring{\ensuremath{\to}}{->} deep ocean decomposition releases CO\texorpdfstring{\textsubscript{2}}{2} \texorpdfstring{\ensuremath{\to}}{->} CO\texorpdfstring{\textsubscript{2}}{2}
  is sequestered for centuries. Calculate: if the biological pump
  sequesters about 11 Gt C/yr and anthropogenic emissions are
  about 10 Gt C/yr, what fraction of emissions is the
  biological pump offsetting? What would happen to atmospheric CO\texorpdfstring{\textsubscript{2}}{2} if
  the biological pump weakened by 20\%?
\end{enumerate}

\begin{enumerate}
\def\labelenumi{\arabic{enumi}.}
\setcounter{enumi}{14}
\tightlist
\item
  A mesocosm experiment measures primary productivity using the
  light-dark bottle method. The light bottle shows dissolved O\texorpdfstring{\textsubscript{2}}{2} increase
  of 3 mg/L over 24 hours. The dark bottle shows O\texorpdfstring{\textsubscript{2}}{2} decrease of 1.5
  mg/L. Calculate: (a) NPP (net change in light bottle); (b) respiration
  (O\texorpdfstring{\textsubscript{2}}{2} loss in dark bottle); (c) GPP (NPP + R). Convert to carbon units
  assuming 1 mg O\texorpdfstring{\textsubscript{2}}{2} = 0.375 mg C.
\end{enumerate}

\begin{enumerate}
\def\labelenumi{\arabic{enumi}.}
\setcounter{enumi}{15}
\tightlist
\item
  Phosphorus is the limiting nutrient in most freshwater ecosystems
  (unlike marine systems where nitrogen often limits). Explain: (a) why
  P availability limits productivity (no atmospheric reservoir; P cycles
  primarily through rock weathering); (b) how agricultural runoff causes
  lake eutrophication (Schindler's whole-lake experiments); (c) why
  banning phosphate detergents was effective in lake recovery.
\end{enumerate}

\begin{enumerate}
\def\labelenumi{\arabic{enumi}.}
\setcounter{enumi}{16}
\tightlist
\item
  Decomposition rates vary with temperature and moisture. Using the ``k
  constant'' for leaf litter decomposition (mass remaining =
  M\texorpdfstring{\textsubscript{0}}{0}e\^{}(\texorpdfstring{\ensuremath{-}}{-}kt)): if k = 0.5/yr in a temperate forest and k = 2.0/yr in a
  tropical forest, calculate the time for 90\% decomposition in each
  ecosystem. What does this imply for nutrient cycling rates?
\end{enumerate}

\begin{enumerate}
\def\labelenumi{\arabic{enumi}.}
\setcounter{enumi}{17}
\tightlist
\item
  The sulfur cycle includes atmospheric SO\texorpdfstring{\textsubscript{2}}{2} (from volcanic eruptions and
  fossil fuel combustion), which forms sulfuric acid (H\texorpdfstring{\textsubscript{2}}{2}SO\texorpdfstring{\textsubscript{4}}{4}) in clouds.
  Trace: SO\texorpdfstring{\textsubscript{2}}{2} \texorpdfstring{\ensuremath{\to}}{->} sulfate aerosol \texorpdfstring{\ensuremath{\to}}{->} acid rain \texorpdfstring{\ensuremath{\to}}{->} soil and lake acidification
  \texorpdfstring{\ensuremath{\to}}{->} biological damage. Name the specific impacts on: (a) fish
  populations; (b) forest soil chemistry (Al\texorpdfstring{\textsuperscript{3}}{3}\texorpdfstring{\textsuperscript{+}}{+} toxicity); (c) building
  materials (limestone dissolution).
\end{enumerate}

\begin{enumerate}
\def\labelenumi{\arabic{enumi}.}
\setcounter{enumi}{18}
\tightlist
\item
  Carbon budget: the global carbon budget (2024) estimates: fossil fuel
  emissions = 9.9 Gt C/yr; land-use change emissions = 1.0 Gt C/yr;
  ocean sink = 2.8 Gt C/yr; land sink = 3.3 Gt C/yr; atmospheric growth
  = 4.8 Gt C/yr. Verify: do sources = sinks + atmospheric growth? What
  is the ``residual'' and what does it reveal about our understanding of
  the carbon cycle?
\end{enumerate}

\begin{enumerate}
\def\labelenumi{\arabic{enumi}.}
\setcounter{enumi}{19}
\tightlist
\item
  Wetlands serve as both carbon sinks and methane (CH\texorpdfstring{\textsubscript{4}}{4}) sources.
  Evaluate: (a) how waterlogged, anoxic conditions promote
  methanogenesis (aceticlastic vs hydrogenotrophic pathways); (b) the
  100-year global warming potential of CH\texorpdfstring{\textsubscript{4}}{4} (about 27--30\texorpdfstring{\ensuremath{\times}}{x} CO\texorpdfstring{\textsubscript{2}}{2} depending
  on fossil vs biogenic accounting in IPCC AR6); (c) the net climate
  effect of wetlands --- do they cool or warm the planet considering
  both C sequestration and CH\texorpdfstring{\textsubscript{4}}{4} emission?
\end{enumerate}

\subsection*{Questions 21--30: Synthesis and
Evaluation}\label{questions-2130-synthesis-and-evaluation-33}
\addcontentsline{toc}{subsection}{Questions 21--30: Synthesis and
Evaluation}

\begin{enumerate}
\def\labelenumi{\arabic{enumi}.}
\setcounter{enumi}{20}
\tightlist
\item
  Evaluate the concept of ``planetary boundaries'' for biogeochemical
  cycles (Rockström et al., 2009). Which cycles have transgressed safe
  operating spaces (nitrogen, phosphorus, climate change), and what does
  this mean for Earth system stability? Propose specific policy
  interventions for the nitrogen cycle that would bring humanity back
  within the safe boundary.
\end{enumerate}

\begin{enumerate}
\def\labelenumi{\arabic{enumi}.}
\setcounter{enumi}{21}
\tightlist
\item
  Design an ecosystem-scale experiment (like Biosphere 2) to test
  whether a simplified ecosystem (5 trophic levels, 50 species) can
  maintain biogeochemical cycling in a sealed environment. What are the
  critical species you would include, what cycling processes must be
  maintained (C, N, P, H\texorpdfstring{\textsubscript{2}}{2}O), and what failure modes would you predict?
\end{enumerate}

\begin{enumerate}
\def\labelenumi{\arabic{enumi}.}
\setcounter{enumi}{22}
\tightlist
\item
  Evaluate the ``metabolic theory of ecology'' (MTE, Brown et al.,
  2004): body mass and temperature determine metabolic rate, which
  scales ecosystem processes (productivity, decomposition, nutrient
  flux). What does MTE predict about the response of tropical vs boreal
  ecosystem productivity to climate warming, and where does MTE break
  down (nutrient limitation, biodiversity effects)?
\end{enumerate}

\begin{enumerate}
\def\labelenumi{\arabic{enumi}.}
\setcounter{enumi}{23}
\tightlist
\item
  Evaluate the role of fire as an ecosystem process: (a) how fire
  releases nutrients (P, K) from biomass; (b) how fire-adapted
  ecosystems (savanna, chaparral, boreal forest) depend on periodic fire
  for regeneration; (c) how fire suppression in fire-dependent
  ecosystems leads to fuel accumulation and catastrophic megafires; (d)
  how indigenous fire management (``cool burns'') can be integrated into
  modern land management.
\end{enumerate}

\begin{enumerate}
\def\labelenumi{\arabic{enumi}.}
\setcounter{enumi}{24}
\tightlist
\item
  Evaluate the evidence that ocean acidification (decreasing pH from CO\texorpdfstring{\textsubscript{2}}{2}
  absorption, currently pH 8.1 \texorpdfstring{\ensuremath{\to}}{->} projected 7.8 by 2100) threatens marine
  ecosystems: (a) what organisms are most affected (calcifiers ---
  corals, pteropods, coccolithophores); (b) how does reduced aragonite
  saturation state (\texorpdfstring{\ensuremath{\Omega}}{Omega}\_ar \textless{} 1) dissolve CaCO\texorpdfstring{\textsubscript{3}}{3} shells; (c) what
  cascading ecosystem effects are predicted; (d) is there evidence of
  evolutionary adaptation to lower pH in any marine organisms?
\end{enumerate}

\begin{enumerate}
\def\labelenumi{\arabic{enumi}.}
\setcounter{enumi}{25}
\tightlist
\item
  Evaluate the concept of ``natural capital'' --- valuing ecosystem
  services (pollination, water purification, carbon storage) in monetary
  terms. Using Costanza et al.~(1997, updated 2014) estimates of global
  ecosystem services at \$125 trillion/yr: (a) is monetary valuation of
  nature scientifically valid; (b) does it effectively change policy;
  (c) what ethical concerns arise from commodifying nature?
\end{enumerate}

\begin{enumerate}
\def\labelenumi{\arabic{enumi}.}
\setcounter{enumi}{26}
\tightlist
\item
  Evaluate the role of permafrost thaw in the global carbon cycle:
  about 1,500 Gt C is stored in permafrost (twice atmospheric
  carbon). As Arctic warming thaws permafrost: (a) what forms of carbon
  are released (CO\texorpdfstring{\textsubscript{2}}{2} from aerobic decomposition, CH\texorpdfstring{\textsubscript{4}}{4} from anaerobic
  decomposition); (b) what is the ``permafrost carbon feedback'' and how
  much warming does it add to existing projections; (c) is permafrost
  thaw a tipping point (irreversible once initiated)?
\end{enumerate}

\begin{enumerate}
\def\labelenumi{\arabic{enumi}.}
\setcounter{enumi}{27}
\tightlist
\item
  Evaluate the concept of ``novel ecosystems'' --- ecosystems with no
  historical analogue, arising from human activity (urban ecosystems,
  post-mining landscapes, invasive-dominated habitats). Should
  conservation efforts attempt to restore historical ecosystems or
  manage novel ecosystems for desired functions? Use a specific example
  to argue both sides.
\end{enumerate}

\begin{enumerate}
\def\labelenumi{\arabic{enumi}.}
\setcounter{enumi}{28}
\tightlist
\item
  Evaluate geoengineering approaches targeting biogeochemical cycles:
  (a) ocean iron fertilisation (stimulating phytoplankton growth to
  increase biological pump CO\texorpdfstring{\textsubscript{2}}{2} sequestration); (b) enhanced weathering
  (grinding silicate rocks to accelerate CO\texorpdfstring{\textsubscript{2}}{2} consumption); (c) bioenergy
  with carbon capture and storage (BECCS). For each, evaluate the scale
  of CO\texorpdfstring{\textsubscript{2}}{2} removal, the ecological risks, and the technological readiness.
\end{enumerate}

\begin{enumerate}
\def\labelenumi{\arabic{enumi}.}
\setcounter{enumi}{29}
\tightlist
\item
  Critically evaluate whether Earth's biogeochemical cycles constitute a
  self-regulating system (Gaia hypothesis, Lovelock 1979). Specifically:
  (a) what evidence supports biological regulation of atmospheric O\texorpdfstring{\textsubscript{2}}{2} at
  about 21\%; (b) what evidence supports biological regulation
  of ocean salinity; (c) what is the difference between ``strong Gaia''
  (Earth as a superorganism) and ``weak Gaia'' (life influences but does
  not control the environment); (d) which version, if either, is
  scientifically defensible?
\end{enumerate}

\newpage

\section*{Questions --- Ch 35: Biomes and Conservation
Biology}\label{questions-ch-35-biomes-and-conservation-biology}
\addcontentsline{toc}{section}{Questions --- Ch 35: Biomes and
Conservation Biology}

\label{sec:q_unit_X_biomes_and_conservation}

\subsection*{Instructor Use and Coverage
Notes}\label{instructor-use-and-coverage-notes-37}
\addcontentsline{toc}{subsection}{Instructor Use and Coverage Notes}

\begin{itemize}
\tightlist
\item
  \textbf{Coverage target:} Evaluate a conservation scenario using
  evidence, tradeoffs, and uncertainty.
\item
  \textbf{Model/data emphasis:} Species-area, risk, and prioritization
  calculations.
\item
  \textbf{Assessment alignment:} Representing and Describing Data,
  Statistical Tests and Data Analysis, Argumentation.
\item
  \textbf{Misconception probe:} Conservation is decision-making under
  constraints and uncertainty, not only preserving untouched nature.
\item
  \textbf{Transfer product:} Transfer conservation reasoning to land
  use, climate corridors, restoration, and environmental justice.
\item
  \textbf{Grading focus:} award full credit for mechanism, evidence,
  boundary conditions, and units when a calculation is required.
\item
  \textbf{Suggested use:} draw one recall item, one application item,
  and one synthesis item when building a short quiz from this bank.
\end{itemize}

\subsection*{Questions 1--10: Recall and
Comprehension}\label{questions-110-recall-and-comprehension-34}
\addcontentsline{toc}{subsection}{Questions 1--10: Recall and
Comprehension}

\emph{This activity accompanies
\cref{sec:unit_X_biomes_and_conservation} of the textbook --- review
that chapter before attempting the exercises below.}

\begin{enumerate}
\def\labelenumi{\arabic{enumi}.}
\tightlist
\item
  Define a biome. What two abiotic factors primarily determine biome
  distribution?
\end{enumerate}

\begin{enumerate}
\def\labelenumi{\arabic{enumi}.}
\setcounter{enumi}{1}
\tightlist
\item
  Name six major terrestrial biomes and list their characteristic
  temperature and precipitation ranges.
\end{enumerate}

\begin{enumerate}
\def\labelenumi{\arabic{enumi}.}
\setcounter{enumi}{2}
\tightlist
\item
  What is the difference between weather and climate?
\end{enumerate}

\begin{enumerate}
\def\labelenumi{\arabic{enumi}.}
\setcounter{enumi}{3}
\tightlist
\item
  What is biodiversity? Distinguish genetic, species, and ecosystem
  diversity.
\end{enumerate}

\begin{enumerate}
\def\labelenumi{\arabic{enumi}.}
\setcounter{enumi}{4}
\tightlist
\item
  Define an endangered species. What organisation maintains the Red
  List?
\end{enumerate}

\begin{enumerate}
\def\labelenumi{\arabic{enumi}.}
\setcounter{enumi}{5}
\tightlist
\item
  What is habitat fragmentation? Why does it threaten biodiversity?
\end{enumerate}

\begin{enumerate}
\def\labelenumi{\arabic{enumi}.}
\setcounter{enumi}{6}
\tightlist
\item
  Name three causes of the current biodiversity crisis.
\end{enumerate}

\begin{enumerate}
\def\labelenumi{\arabic{enumi}.}
\setcounter{enumi}{7}
\tightlist
\item
  What is a biodiversity hotspot? Name two examples.
\end{enumerate}

\begin{enumerate}
\def\labelenumi{\arabic{enumi}.}
\setcounter{enumi}{8}
\tightlist
\item
  What is an in situ conservation strategy? What is an ex situ strategy?
\end{enumerate}

\begin{enumerate}
\def\labelenumi{\arabic{enumi}.}
\setcounter{enumi}{9}
\tightlist
\item
  Define sustainability. What are the three pillars of sustainability?
\end{enumerate}

\subsection*{Questions 11--20: Application and
Analysis}\label{questions-1120-application-and-analysis-34}
\addcontentsline{toc}{subsection}{Questions 11--20: Application and
Analysis}

\begin{enumerate}
\def\labelenumi{\arabic{enumi}.}
\setcounter{enumi}{10}
\tightlist
\item
  The species-area relationship (\(S = cA^z\), where z \texorpdfstring{\ensuremath{\approx}}{~} 0.25 for
  oceanic islands) predicts that reducing habitat area by 90\% reduces
  species number by about 50\%. Calculate: if a 10,000 km\texorpdfstring{\textsuperscript{2}}{2}
  forest is reduced to 1,000 km\texorpdfstring{\textsuperscript{2}}{2}, and the original forest had 200 bird
  species, how many bird species are expected in the remaining fragment?
  What assumptions does this model make?
\end{enumerate}

\begin{enumerate}
\def\labelenumi{\arabic{enumi}.}
\setcounter{enumi}{11}
\tightlist
\item
  EDGE scores (Evolutionarily Distinct and Globally Endangered) are used
  to prioritise species for conservation investment. The aye-aye
  (\emph{Daubentonia madagascariensis}) --- the sole living member of
  its family --- has a very high EDGE score. Explain: (a) what
  evolutionary distinctiveness means; (b) why losing an ``edge species''
  results in greater phylogenetic diversity loss than losing a species
  in a species-rich clade; (c) how EDGE scoring changes conservation
  priorities compared to a focus on charisma.
\end{enumerate}

\begin{enumerate}
\def\labelenumi{\arabic{enumi}.}
\setcounter{enumi}{12}
\tightlist
\item
  Minimum viable population (MVP) analysis: a PVA model estimates that a
  population of 500 individuals has a 95\% probability of persisting for
  100 years, but a population of 50 has a 20\% chance. Identify three
  factors incorporated in the model (demographic stochasticity,
  environmental stochasticity, inbreeding depression) and explain how
  each reduces persistence probability.
\end{enumerate}

\begin{enumerate}
\def\labelenumi{\arabic{enumi}.}
\setcounter{enumi}{13}
\tightlist
\item
  Climate change is shifting biome boundaries: the boreal-temperate
  boundary is moving northward, and the Sahel-Sahara boundary is
  fluctuating. Predict the consequences for: (a) species that cannot
  migrate fast enough (boreal specialists trapped in shrinking habitat);
  (b) agricultural zones (expansion of growing season in formerly cold
  regions); (c) novel community assemblages (no-analogue communities
  formed by differential range shifts).
\end{enumerate}

\begin{enumerate}
\def\labelenumi{\arabic{enumi}.}
\setcounter{enumi}{14}
\tightlist
\item
  Coral reef bleaching occurs when accumulated heat stress disrupts the
  coral-Symbiodiniaceae partnership. Explain: (a) how
  degree-heating-week evidence links heat exposure to symbiont
  photosystem stress, reactive oxygen species, and symbiont loss; (b)
  why recovery depends on repeated heat exposure, water quality,
  recruitment, and symbiont community composition rather than a fixed
  clock; (c) how assisted evolution or symbiont turnover could improve
  local resilience while still failing as a substitute for climate
  mitigation.
\end{enumerate}

\begin{enumerate}
\def\labelenumi{\arabic{enumi}.}
\setcounter{enumi}{15}
\tightlist
\item
  Protected areas currently cover about 17\% of land and
  about 8\% of ocean. The Kunming-Montreal Global Biodiversity
  Framework (2022) targets 30\% by 2030 (``30 by 30''). Evaluate: (a)
  whether area alone is sufficient (effective management, connectivity,
  representativeness); (b) the concept of ``half-Earth'' (Wilson, 2016
  --- protecting 50\% of Earth's surface); (c) the social justice
  concerns (indigenous displacement, land rights).
\end{enumerate}

\begin{enumerate}
\def\labelenumi{\arabic{enumi}.}
\setcounter{enumi}{16}
\tightlist
\item
  Invasive species economic damage: the global estimated cost of
  invasive species is \textgreater\$423 billion/yr (2019). Using the
  example of the invasive cane toad (\emph{Rhinella marina}) in
  Australia: (a) what was the original purpose of introduction; (b) what
  ecological impacts have occurred (native predator mortality from toad
  toxin); (c) what control strategies have been attempted (biological
  control, trapping, gene drive research).
\end{enumerate}

\begin{enumerate}
\def\labelenumi{\arabic{enumi}.}
\setcounter{enumi}{17}
\tightlist
\item
  Payment for ecosystem services (PES): Costa Rica's PES programme pays
  landowners to conserve forest (carbon storage, water regulation).
  Evaluate: (a) how PES monetises ecosystem services; (b) evidence of
  effectiveness (Costa Rica reversed deforestation); (c) criticisms
  (additionality --- would forests have been conserved anyway? Equity
  --- do poor landowners benefit or are payments captured by wealthy
  elites?).
\end{enumerate}

\begin{enumerate}
\def\labelenumi{\arabic{enumi}.}
\setcounter{enumi}{18}
\tightlist
\item
  The IPBES Global Assessment estimates that about one million species
  are threatened, while the IPBES Transformative Change Assessment
  argues that biodiversity loss is driven by economic, governance, and
  social systems as well as local habitat conversion. Evaluate: (a) how
  threatened-species estimates combine assessed taxa, extrapolation, and
  habitat trends; (b) why food systems, healthy diets, agroecology,
  fisheries, and commodity supply chains belong inside conservation
  biology; (c) how to distinguish evidence-based transformative change
  from vague sustainability branding.
\end{enumerate}

\begin{enumerate}
\def\labelenumi{\arabic{enumi}.}
\setcounter{enumi}{19}
\tightlist
\item
  Rewilding involves reintroducing apex predators and megafauna to
  restore ecosystem function. Evaluate the Eurasian rewilding movement
  (European bison in Romania, wolves in Germany, lynx in UK): (a) what
  ecological functions are restored; (b) what human-wildlife conflict
  arises; (c) the evidence from existing rewilding projects
  (Oostvaardersplassen, Netherlands) --- successes and failures.
\end{enumerate}

\subsection*{Questions 21--30: Synthesis and
Evaluation}\label{questions-2130-synthesis-and-evaluation-34}
\addcontentsline{toc}{subsection}{Questions 21--30: Synthesis and
Evaluation}

\begin{enumerate}
\def\labelenumi{\arabic{enumi}.}
\setcounter{enumi}{20}
\tightlist
\item
  Evaluate the ``extinction debt'' concept: habitat destruction in the
  past has committed species to future extinction even if no further
  habitat is lost, because small fragments cannot sustain viable
  populations over decades. Using the example of tropical forest
  fragmentation (Ferraz et al., 2003): (a) how is extinction debt
  estimated; (b) how long is the time lag (about 20--100
  years); (c) what does this imply for current biodiversity assessments
  (we are underestimating the true impact of past deforestation)?
\end{enumerate}

\begin{enumerate}
\def\labelenumi{\arabic{enumi}.}
\setcounter{enumi}{21}
\tightlist
\item
  Design a conservation genetics programme for a critically endangered
  species (e.g., California condor, about 500 individuals):
  describe (a) the genetic assessment (genomic sequencing of sampled
  individuals, measurement of heterozygosity and inbreeding
  coefficient); (b) the breeding management strategy (maximising N\_e,
  using mean kinship to minimise inbreeding); (c) the genetic rescue
  strategy (introducing individuals from a genetically distinct
  population to restore lost alleles). What are the risks of genetic
  rescue (outbreeding depression)?
\end{enumerate}

\begin{enumerate}
\def\labelenumi{\arabic{enumi}.}
\setcounter{enumi}{22}
\tightlist
\item
  Evaluate climate-smart conservation: traditional conservation aims to
  protect current biodiversity in current locations. Climate change
  requires protecting future habitats and enabling species movement.
  Evaluate three climate-smart strategies: (a) connectivity corridors
  between protected areas; (b) assisted migration (moving species to
  suitable future habitat); (c) protecting climate refugia (areas
  predicted to remain climatically stable). What are the risks of each?
\end{enumerate}

\begin{enumerate}
\def\labelenumi{\arabic{enumi}.}
\setcounter{enumi}{23}
\tightlist
\item
  Evaluate the role of indigenous and traditional ecological knowledge
  (ITEK) in conservation: (a) the evidence that indigenous-managed lands
  have higher biodiversity than conventional protected areas; (b) the UN
  IPBES recognition of ITEK as an evidence base equal to scientific
  knowledge; (c) the challenges of integrating ITEK with Western science
  (different epistemologies, power dynamics, intellectual property).
\end{enumerate}

\begin{enumerate}
\def\labelenumi{\arabic{enumi}.}
\setcounter{enumi}{24}
\tightlist
\item
  Evaluate the concept of ``de-extinction'' in the context of
  conservation biology: if woolly mammoths could be de-extincted (via
  CRISPR editing of Asian elephant genome), would this be a conservation
  priority? Discuss: (a) the opportunity cost (diverting funds from
  existing species); (b) the ecological benefit (restoring mammoth
  steppe ecosystem, permafrost preservation); (c) the ethical questions
  (welfare of first-generation animals, ecological risk).
\end{enumerate}

\begin{enumerate}
\def\labelenumi{\arabic{enumi}.}
\setcounter{enumi}{25}
\tightlist
\item
  Evaluate the impact of plastic pollution on marine biomes: (a)
  microplastic ingestion by marine organisms across trophic levels; (b)
  bioaccumulation and biomagnification of associated chemical pollutants
  (phthalates, BPA); (c) the ``plastisphere'' --- microbial communities
  colonising marine plastic that can include pathogens; (d) policy
  interventions (single-use plastic bans, Extended Producer
  Responsibility).
\end{enumerate}

\begin{enumerate}
\def\labelenumi{\arabic{enumi}.}
\setcounter{enumi}{26}
\tightlist
\item
  Evaluate Nature-based Solutions (NbS) to climate change: (a)
  reforestation and afforestation (potential about 10 Gt
  CO\texorpdfstring{\textsubscript{2}}{2}/yr); (b) peatland and mangrove restoration; (c) regenerative
  agriculture (soil carbon sequestration). For each, evaluate the
  permanence of carbon storage, the biodiversity co-benefits, and the
  risk of ``greenwashing'' (using NbS to justify continued fossil fuel
  use).
\end{enumerate}

\begin{enumerate}
\def\labelenumi{\arabic{enumi}.}
\setcounter{enumi}{27}
\tightlist
\item
  Evaluate the role of zoos and botanical gardens (ex situ conservation)
  in preventing extinction: (a) the success stories (Arabian oryx,
  black-footed ferret, California condor --- saved from extinction by
  captive breeding); (b) the limitations (genetic bottlenecks,
  behavioural domestication, space constraints); (c) the shift from
  ``Noah's Ark'' (saving species) to ``conservation hubs'' (research,
  education, advocacy).
\end{enumerate}

\begin{enumerate}
\def\labelenumi{\arabic{enumi}.}
\setcounter{enumi}{28}
\tightlist
\item
  Evaluate the ``triage'' approach to conservation: given limited
  resources, should we focus on: (a) species most likely to be saved
  (pragmatic triage); (b) species with highest evolutionary
  distinctiveness (phylogenetic triage); (c) species with greatest
  functional importance (ecosystem triage)? Compare these approaches
  using a specific threatened community (e.g., coral reef, tropical
  forest), and discuss the ethical implications of explicitly choosing
  to let some species go extinct.
\end{enumerate}

\begin{enumerate}
\def\labelenumi{\arabic{enumi}.}
\setcounter{enumi}{29}
\tightlist
\item
  Critically evaluate whether the concept of ``ecological footprint''
  (measuring human demand on Earth's biocapacity in global hectares) is
  a scientifically rigorous and policy-useful metric. Compare ecological
  footprint to alternative sustainability metrics (planetary boundaries,
  natural capital accounting, genuine progress indicator). What does
  each capture that others miss, and which is most actionable for policy
  change?
\end{enumerate}

\newpage

\section*{Appendix A --- Curriculum
Map}\label{appendix-a-curriculum-map}
\addcontentsline{toc}{section}{Appendix A --- Curriculum Map}

\label{sec:appendix_curriculum_map}

This appendix is generated from \breaktt{src/biology/curriculum.py} and
\breaktt{src/biology/alignment.py}. It links each chapter to its
companion lab, question bank, model/data skill, misconception probe,
transfer task, and standards/skills alignment so the textbook can be
composed into different course formats without losing pedagogical
coherence.

\subsection*{Unit 0}\label{unit-0}
\addcontentsline{toc}{subsection}{Unit 0}

\subsubsection*{Systems Science and the Logic of
Emergence}\label{systems-science-and-the-logic-of-emergence-1}
\addcontentsline{toc}{subsubsection}{Systems Science and the Logic of
Emergence}

\begin{itemize}
\tightlist
\item
  \textbf{Chapter:} \cref{sec:unit_0_systems_science}.
\item
  \textbf{Lab:} \cref{sec:lab_unit_0_systems_science}.
\item
  \textbf{Question bank:} \cref{sec:q_unit_0_systems_science}.
\item
  \textbf{Big idea:} Biological explanation improves when parts,
  interactions, feedback, and scale are kept in view together.
\item
  \textbf{Core concepts:} systems boundaries, feedback, emergence, state
  variables.
\item
  \textbf{Framework alignment:} Vision \& Change: Systems, Structure and
  function; AP Biology: Systems Interactions; NGSS-style topics:
  Structure and Function, Interdependent Relationships in Ecosystems.
\item
  \textbf{Practice cadence:} Questions and Methods, Representing and
  Describing Data, Argumentation.
\item
  \textbf{Model/data skill:} Box-and-arrow causal models with explicit
  inputs, outputs, and feedback signs. Translate a verbal biological
  system into variables, links, and testable predictions.
\item
  \textbf{Misconception probe:} A system is not just a list of parts;
  the interactions are part of the explanation.
\item
  \textbf{Transfer task:} Apply the same feedback map to a cell,
  organism, and ecosystem, then name what changes at each scale.
\end{itemize}

\subsubsection*{Complex Adaptive
Systems}\label{complex-adaptive-systems-1}
\addcontentsline{toc}{subsubsection}{Complex Adaptive Systems}

\begin{itemize}
\tightlist
\item
  \textbf{Chapter:} \cref{sec:unit_0_complex_adaptive_systems}.
\item
  \textbf{Lab:} \cref{sec:lab_unit_0_complex_adaptive_systems}.
\item
  \textbf{Question bank:} \cref{sec:q_unit_0_complex_adaptive_systems}.
\item
  \textbf{Big idea:} Local rules can generate population-level patterns
  without a central controller.
\item
  \textbf{Core concepts:} nonlinearity, self-organization, attractors,
  power laws.
\item
  \textbf{Framework alignment:} Vision \& Change: Systems, Structure and
  function; AP Biology: Systems Interactions; NGSS-style topics:
  Structure and Function, Interdependent Relationships in Ecosystems.
\item
  \textbf{Practice cadence:} Questions and Methods, Representing and
  Describing Data, Argumentation.
\item
  \textbf{Model/data skill:} Agent-rule and scaling arguments for
  emergent biological patterns. Distinguish deterministic trends from
  stochastic variation in repeated simulations.
\item
  \textbf{Misconception probe:} Emergence is not mysterious; it is a
  reproducible consequence of interactions plus constraints.
\item
  \textbf{Transfer task:} Compare flocking, immune activation, and
  microbial biofilms as adaptive systems.
\end{itemize}

\subsubsection*{Active Inference and the Free Energy
Principle}\label{active-inference-and-the-free-energy-principle-1}
\addcontentsline{toc}{subsubsection}{Active Inference and the Free
Energy Principle}

\begin{itemize}
\tightlist
\item
  \textbf{Chapter:} \cref{sec:unit_0_active_inference}.
\item
  \textbf{Lab:} \cref{sec:lab_unit_0_active_inference}.
\item
  \textbf{Question bank:} \cref{sec:q_unit_0_active_inference}.
\item
  \textbf{Big idea:} Living systems reduce uncertainty by acting on the
  world as well as by updating internal models.
\item
  \textbf{Core concepts:} prediction error, Bayesian updating, policy
  selection, homeostasis.
\item
  \textbf{Framework alignment:} Vision \& Change: Systems, Structure and
  function; AP Biology: Systems Interactions; NGSS-style topics:
  Structure and Function, Interdependent Relationships in Ecosystems.
\item
  \textbf{Practice cadence:} Questions and Methods, Representing and
  Describing Data, Argumentation.
\item
  \textbf{Model/data skill:} Bayesian belief updating and
  expected-free-energy-style policy comparison. Read a small probability
  table and update a prediction after new evidence.
\item
  \textbf{Misconception probe:} Active inference is not passive
  prediction; action changes the sensory data that arrive next.
\item
  \textbf{Transfer task:} Map prediction-error reasoning onto
  chemotaxis, thermoregulation, and attention.
\end{itemize}

\subsection*{Unit I}\label{unit-i}
\addcontentsline{toc}{subsection}{Unit I}

\subsubsection*{Atoms, Molecules, and Chemical
Bonds}\label{atoms-molecules-and-chemical-bonds-1}
\addcontentsline{toc}{subsubsection}{Atoms, Molecules, and Chemical
Bonds}

\begin{itemize}
\tightlist
\item
  \textbf{Chapter:} \cref{sec:unit_I_atoms_molecules}.
\item
  \textbf{Lab:} \cref{sec:lab_unit_I_atoms_molecules}.
\item
  \textbf{Question bank:} \cref{sec:q_unit_I_atoms_molecules}.
\item
  \textbf{Big idea:} Atomic structure constrains the bonds, charges, and
  geometries that make biological molecules possible.
\item
  \textbf{Core concepts:} valence, electronegativity, formal charge,
  isotopes.
\item
  \textbf{Framework alignment:} Vision \& Change: Structure and
  function, Pathways and transformations of energy and matter; AP
  Biology: Energetics, Systems Interactions; NGSS-style topics: Matter
  and Energy in Organisms and Ecosystems, Structure and Function.
\item
  \textbf{Practice cadence:} Concept Explanation, Statistical Tests and
  Data Analysis, Argumentation.
\item
  \textbf{Model/data skill:} Formal charge and
  electronegativity-difference reasoning. Use tabular atomic data to
  predict polarity, solubility, and biological reactivity.
\item
  \textbf{Misconception probe:} Weak interactions are not unimportant;
  many weak interactions together dominate structure and specificity.
\item
  \textbf{Transfer task:} Use atomic reasoning to explain a medical
  tracer, enzyme cofactor, or membrane-solubility problem.
\end{itemize}

\subsubsection*{Water --- The Molecule of
Life}\label{water-the-molecule-of-life-1}
\addcontentsline{toc}{subsubsection}{Water --- The Molecule of Life}

\begin{itemize}
\tightlist
\item
  \textbf{Chapter:} \cref{sec:unit_I_water_and_life}.
\item
  \textbf{Lab:} \cref{sec:lab_unit_I_water_and_life}.
\item
  \textbf{Question bank:} \cref{sec:q_unit_I_water_and_life}.
\item
  \textbf{Big idea:} Water's polarity, hydrogen bonding, and ionization
  make cells physically possible.
\item
  \textbf{Core concepts:} hydrogen bonding, cohesion, pH, buffers.
\item
  \textbf{Framework alignment:} Vision \& Change: Structure and
  function, Pathways and transformations of energy and matter; AP
  Biology: Energetics, Systems Interactions; NGSS-style topics: Matter
  and Energy in Organisms and Ecosystems, Structure and Function.
\item
  \textbf{Practice cadence:} Concept Explanation, Statistical Tests and
  Data Analysis, Argumentation.
\item
  \textbf{Model/data skill:} Henderson-Hasselbalch and pH-scale
  calculations. Convert between pH, hydrogen ion concentration, and
  buffer ratios.
\item
  \textbf{Misconception probe:} Water is not an inert background; it is
  an active participant in structure and reaction chemistry.
\item
  \textbf{Transfer task:} Transfer water-property reasoning to blood
  buffering, plant transport, or protein folding.
\end{itemize}

\subsubsection*{Biological
Macromolecules}\label{biological-macromolecules-1}
\addcontentsline{toc}{subsubsection}{Biological Macromolecules}

\begin{itemize}
\tightlist
\item
  \textbf{Chapter:} \cref{sec:unit_I_macromolecules}.
\item
  \textbf{Lab:} \cref{sec:lab_unit_I_macromolecules}.
\item
  \textbf{Question bank:} \cref{sec:q_unit_I_macromolecules}.
\item
  \textbf{Big idea:} Macromolecule structure links monomer chemistry to
  biological function and information storage.
\item
  \textbf{Core concepts:} polymers, dehydration synthesis, protein
  structure, nucleic acids.
\item
  \textbf{Framework alignment:} Vision \& Change: Structure and
  function, Pathways and transformations of energy and matter; AP
  Biology: Energetics, Systems Interactions; NGSS-style topics: Matter
  and Energy in Organisms and Ecosystems, Structure and Function.
\item
  \textbf{Practice cadence:} Concept Explanation, Statistical Tests and
  Data Analysis, Argumentation.
\item
  \textbf{Model/data skill:} Polymerization energetics and
  sequence-to-structure reasoning. Classify molecules from structural
  evidence rather than names alone.
\item
  \textbf{Misconception probe:} Structure is not decoration; small
  chemical changes can redirect function and recognition.
\item
  \textbf{Transfer task:} Explain how a mutation, lipid substitution, or
  glycosylation change propagates to phenotype.
\end{itemize}

\subsubsection*{Enzymes and the Kinetics of
Catalysis}\label{enzymes-and-the-kinetics-of-catalysis-1}
\addcontentsline{toc}{subsubsection}{Enzymes and the Kinetics of
Catalysis}

\begin{itemize}
\tightlist
\item
  \textbf{Chapter:} \cref{sec:unit_I_enzymes_and_kinetics}.
\item
  \textbf{Lab:} \cref{sec:lab_unit_I_enzymes_and_kinetics}.
\item
  \textbf{Question bank:} \cref{sec:q_unit_I_enzymes_and_kinetics}.
\item
  \textbf{Big idea:} Enzymes accelerate reactions by stabilizing
  transition states without changing thermodynamic endpoints.
\item
  \textbf{Core concepts:} activation energy, Michaelis-Menten kinetics,
  inhibition, allostery.
\item
  \textbf{Framework alignment:} Vision \& Change: Structure and
  function, Pathways and transformations of energy and matter; AP
  Biology: Energetics, Systems Interactions; NGSS-style topics: Matter
  and Energy in Organisms and Ecosystems, Structure and Function.
\item
  \textbf{Practice cadence:} Concept Explanation, Statistical Tests and
  Data Analysis, Argumentation.
\item
  \textbf{Model/data skill:} Michaelis-Menten, Lineweaver-Burk, and
  inhibition-pattern calculations. Fit or interpret enzyme-rate data and
  identify which parameter changed.
\item
  \textbf{Misconception probe:} A catalyst changes rate, not the
  equilibrium constant or the sign of delta G.
\item
  \textbf{Transfer task:} Apply enzyme-kinetic reasoning to drug dosing,
  metabolic control, or diagnostic assays.
\end{itemize}

\subsection*{Unit II}\label{unit-ii}
\addcontentsline{toc}{subsection}{Unit II}

\subsubsection*{Cell Theory and Cell
Types}\label{cell-theory-and-cell-types-1}
\addcontentsline{toc}{subsubsection}{Cell Theory and Cell Types}

\begin{itemize}
\tightlist
\item
  \textbf{Chapter:} \cref{sec:unit_II_cell_theory}.
\item
  \textbf{Lab:} \cref{sec:lab_unit_II_cell_theory}.
\item
  \textbf{Question bank:} \cref{sec:q_unit_II_cell_theory}.
\item
  \textbf{Big idea:} Cells are bounded, evolving systems whose size and
  organization reflect physical constraints.
\item
  \textbf{Core concepts:} cell theory, surface area, microscopy,
  prokaryote/eukaryote comparison.
\item
  \textbf{Framework alignment:} Vision \& Change: Structure and
  function, Systems, Information flow, exchange, and storage; AP
  Biology: Systems Interactions, Information Storage and Transmission;
  NGSS-style topics: Structure and Function.
\item
  \textbf{Practice cadence:} Visual Representations, Questions and
  Methods, Argumentation.
\item
  \textbf{Model/data skill:} Surface-area-to-volume scaling. Infer
  cellular constraints from measurements, micrographs, and scale bars.
\item
  \textbf{Misconception probe:} Cells are not just small bags of fluid;
  boundaries and internal organization create function.
\item
  \textbf{Transfer task:} Transfer scale reasoning to eggs, neurons,
  root hairs, and microbial colonies.
\end{itemize}

\subsubsection*{Cell Structure and
Organelles}\label{cell-structure-and-organelles-1}
\addcontentsline{toc}{subsubsection}{Cell Structure and Organelles}

\begin{itemize}
\tightlist
\item
  \textbf{Chapter:} \cref{sec:unit_II_cell_structure}.
\item
  \textbf{Lab:} \cref{sec:lab_unit_II_cell_structure}.
\item
  \textbf{Question bank:} \cref{sec:q_unit_II_cell_structure}.
\item
  \textbf{Big idea:} Organelle structure partitions work, traffic,
  information, and energy inside eukaryotic cells.
\item
  \textbf{Core concepts:} organelles, endomembrane system, cytoskeleton,
  motor proteins.
\item
  \textbf{Framework alignment:} Vision \& Change: Structure and
  function, Systems, Information flow, exchange, and storage; AP
  Biology: Systems Interactions, Information Storage and Transmission;
  NGSS-style topics: Structure and Function.
\item
  \textbf{Practice cadence:} Visual Representations, Questions and
  Methods, Argumentation.
\item
  \textbf{Model/data skill:} Compartment-flow and motor-transport
  calculations. Trace a molecule through compartments using evidence
  from labels or perturbations.
\item
  \textbf{Misconception probe:} Organelles are dynamic process hubs, not
  static textbook icons.
\item
  \textbf{Transfer task:} Apply compartment logic to secretion,
  apoptosis, infection, or cell division.
\end{itemize}

\subsubsection*{Membrane Structure and
Transport}\label{membrane-structure-and-transport-1}
\addcontentsline{toc}{subsubsection}{Membrane Structure and Transport}

\begin{itemize}
\tightlist
\item
  \textbf{Chapter:} \cref{sec:unit_II_membrane_transport}.
\item
  \textbf{Lab:} \cref{sec:lab_unit_II_membrane_transport}.
\item
  \textbf{Question bank:} \cref{sec:q_unit_II_membrane_transport}.
\item
  \textbf{Big idea:} Membranes convert gradients, permeability, and
  selective transport into cellular physiology.
\item
  \textbf{Core concepts:} bilayers, diffusion, osmosis, electrochemical
  gradients.
\item
  \textbf{Framework alignment:} Vision \& Change: Structure and
  function, Systems, Information flow, exchange, and storage; AP
  Biology: Systems Interactions, Information Storage and Transmission;
  NGSS-style topics: Structure and Function.
\item
  \textbf{Practice cadence:} Visual Representations, Questions and
  Methods, Argumentation.
\item
  \textbf{Model/data skill:} Nernst, Goldman, osmotic, and
  facilitated-transport calculations. Interpret transport data from
  gradients, rates, and membrane potentials.
\item
  \textbf{Misconception probe:} Equilibrium does not mean equal
  concentration; charge and permeability matter.
\item
  \textbf{Transfer task:} Transfer gradient logic to neurons, kidneys,
  roots, and mitochondrial membranes.
\end{itemize}

\subsubsection*{Cell Signalling and
Communication}\label{cell-signalling-and-communication-1}
\addcontentsline{toc}{subsubsection}{Cell Signalling and Communication}

\begin{itemize}
\tightlist
\item
  \textbf{Chapter:} \cref{sec:unit_II_cell_signaling}.
\item
  \textbf{Lab:} \cref{sec:lab_unit_II_cell_signaling}.
\item
  \textbf{Question bank:} \cref{sec:q_unit_II_cell_signaling}.
\item
  \textbf{Big idea:} Cells communicate by converting external signals
  into regulated intracellular decisions.
\item
  \textbf{Core concepts:} receptors, second messengers, signal
  amplification, feedback.
\item
  \textbf{Framework alignment:} Vision \& Change: Structure and
  function, Systems, Information flow, exchange, and storage; AP
  Biology: Systems Interactions, Information Storage and Transmission;
  NGSS-style topics: Structure and Function.
\item
  \textbf{Practice cadence:} Visual Representations, Questions and
  Methods, Argumentation.
\item
  \textbf{Model/data skill:} Ligand-receptor occupancy and dose-response
  reasoning. Read pathway diagrams and infer the effect of agonists,
  antagonists, or mutations.
\item
  \textbf{Misconception probe:} A pathway diagram is a causal model, not
  a memorization chart.
\item
  \textbf{Transfer task:} Apply signaling logic to hormones,
  neurotransmitters, immune receptors, or cancer mutations.
\end{itemize}

\subsection*{Unit III}\label{unit-iii}
\addcontentsline{toc}{subsection}{Unit III}

\subsubsection*{Bioenergetics and Cellular
Respiration}\label{bioenergetics-and-cellular-respiration-1}
\addcontentsline{toc}{subsubsection}{Bioenergetics and Cellular
Respiration}

\begin{itemize}
\tightlist
\item
  \textbf{Chapter:} \cref{sec:unit_III_bioenergetics_and_respiration}.
\item
  \textbf{Lab:} \cref{sec:lab_unit_III_bioenergetics_and_respiration}.
\item
  \textbf{Question bank:}
  \cref{sec:q_unit_III_bioenergetics_and_respiration}.
\item
  \textbf{Big idea:} Cells harvest free energy by coupling redox
  chemistry to phosphoryl transfer and ion gradients.
\item
  \textbf{Core concepts:} free energy, redox, glycolysis, oxidative
  phosphorylation.
\item
  \textbf{Framework alignment:} Vision \& Change: Pathways and
  transformations of energy and matter, Systems; AP Biology: Energetics,
  Systems Interactions; NGSS-style topics: Matter and Energy in
  Organisms and Ecosystems.
\item
  \textbf{Practice cadence:} Representing and Describing Data,
  Statistical Tests and Data Analysis.
\item
  \textbf{Model/data skill:} Delta G, ATP yield, and electron-carrier
  accounting. Track carbon, electrons, and ATP across a pathway.
\item
  \textbf{Misconception probe:} ATP is not stored energy in a vague
  sense; it is a coupling currency with defined reaction chemistry.
\item
  \textbf{Transfer task:} Transfer energy accounting to exercise,
  fermentation, hypoxia, and mitochondrial disease.
\end{itemize}

\subsubsection*{Photosynthesis}\label{photosynthesis-1}
\addcontentsline{toc}{subsubsection}{Photosynthesis}

\begin{itemize}
\tightlist
\item
  \textbf{Chapter:} \cref{sec:unit_III_photosynthesis}.
\item
  \textbf{Lab:} \cref{sec:lab_unit_III_photosynthesis}.
\item
  \textbf{Question bank:} \cref{sec:q_unit_III_photosynthesis}.
\item
  \textbf{Big idea:} Photosynthesis couples light-driven electron flow
  to carbon fixation and planetary productivity.
\item
  \textbf{Core concepts:} photosystems, electron transport, Calvin
  cycle, photorespiration.
\item
  \textbf{Framework alignment:} Vision \& Change: Pathways and
  transformations of energy and matter, Systems; AP Biology: Energetics,
  Systems Interactions; NGSS-style topics: Matter and Energy in
  Organisms and Ecosystems.
\item
  \textbf{Practice cadence:} Representing and Describing Data,
  Statistical Tests and Data Analysis.
\item
  \textbf{Model/data skill:} Photon, ATP/NADPH, and Calvin-cycle
  stoichiometry. Interpret light-response and carbon-fixation data.
\item
  \textbf{Misconception probe:} Plants do not eat sunlight; they use
  light energy to reduce carbon using electrons and enzymes.
\item
  \textbf{Transfer task:} Transfer photosynthetic constraints to crops,
  algae, climate, or ecosystem productivity.
\end{itemize}

\subsubsection*{Metabolic Integration and
Regulation}\label{metabolic-integration-and-regulation-1}
\addcontentsline{toc}{subsubsection}{Metabolic Integration and
Regulation}

\begin{itemize}
\tightlist
\item
  \textbf{Chapter:} \cref{sec:unit_III_metabolic_integration}.
\item
  \textbf{Lab:} \cref{sec:lab_unit_III_metabolic_integration}.
\item
  \textbf{Question bank:} \cref{sec:q_unit_III_metabolic_integration}.
\item
  \textbf{Big idea:} Metabolism is a regulated network that reallocates
  flux across tissues, time, and nutrient states.
\item
  \textbf{Core concepts:} flux, energy charge, hormonal control,
  fed/fasted states.
\item
  \textbf{Framework alignment:} Vision \& Change: Pathways and
  transformations of energy and matter, Systems; AP Biology: Energetics,
  Systems Interactions; NGSS-style topics: Matter and Energy in
  Organisms and Ecosystems.
\item
  \textbf{Practice cadence:} Representing and Describing Data,
  Statistical Tests and Data Analysis.
\item
  \textbf{Model/data skill:} Energy charge, control points, and
  pathway-flux comparisons. Use pathway evidence to infer which
  metabolic state or tissue is active.
\item
  \textbf{Misconception probe:} A pathway is not a one-way assembly
  line; reversibility and regulation define the real route.
\item
  \textbf{Transfer task:} Apply metabolic network reasoning to diabetes,
  fasting, exercise, or cancer metabolism.
\end{itemize}

\subsection*{Unit IV}\label{unit-iv}
\addcontentsline{toc}{subsection}{Unit IV}

\subsubsection*{DNA Replication and the Cell
Cycle}\label{dna-replication-and-the-cell-cycle-1}
\addcontentsline{toc}{subsubsection}{DNA Replication and the Cell Cycle}

\begin{itemize}
\tightlist
\item
  \textbf{Chapter:} \cref{sec:unit_IV_dna_replication_and_cell_cycle}.
\item
  \textbf{Lab:} \cref{sec:lab_unit_IV_dna_replication_and_cell_cycle}.
\item
  \textbf{Question bank:}
  \cref{sec:q_unit_IV_dna_replication_and_cell_cycle}.
\item
  \textbf{Big idea:} Genome copying is accurate because chemistry,
  enzyme proofreading, and checkpoints cooperate.
\item
  \textbf{Core concepts:} semiconservative replication, polymerase,
  proofreading, checkpoints.
\item
  \textbf{Framework alignment:} Vision \& Change: Information flow,
  exchange, and storage, Structure and function; AP Biology: Information
  Storage and Transmission, Systems Interactions; NGSS-style topics:
  Inheritance and Variation of Traits, Structure and Function.
\item
  \textbf{Practice cadence:} Concept Explanation, Questions and Methods,
  Argumentation.
\item
  \textbf{Model/data skill:} Replication-fork speed, error-rate, and
  cell-cycle timing calculations. Interpret replication or cell-cycle
  data from timing, labeling, or checkpoint perturbations.
\item
  \textbf{Misconception probe:} High fidelity is not automatic; it is
  built from multiple partially redundant safeguards.
\item
  \textbf{Transfer task:} Transfer replication logic to cancer, aging,
  viral replication, or antibiotic targets.
\end{itemize}

\subsubsection*{Gene Expression --- Transcription and
Translation}\label{gene-expression-transcription-and-translation-1}
\addcontentsline{toc}{subsubsection}{Gene Expression --- Transcription
and Translation}

\begin{itemize}
\tightlist
\item
  \textbf{Chapter:} \cref{sec:unit_IV_gene_expression}.
\item
  \textbf{Lab:} \cref{sec:lab_unit_IV_gene_expression}.
\item
  \textbf{Question bank:} \cref{sec:q_unit_IV_gene_expression}.
\item
  \textbf{Big idea:} Gene expression translates sequence information
  into regulated RNA and protein products.
\item
  \textbf{Core concepts:} transcription, translation, codons, RNA
  processing.
\item
  \textbf{Framework alignment:} Vision \& Change: Information flow,
  exchange, and storage, Structure and function; AP Biology: Information
  Storage and Transmission, Systems Interactions; NGSS-style topics:
  Inheritance and Variation of Traits, Structure and Function.
\item
  \textbf{Practice cadence:} Concept Explanation, Questions and Methods,
  Argumentation.
\item
  \textbf{Model/data skill:} Reading-frame, codon, and expression-output
  calculations. Convert DNA/RNA sequence data into predicted molecular
  products.
\item
  \textbf{Misconception probe:} A gene is not simply a protein recipe;
  context controls when, where, and how much product appears.
\item
  \textbf{Transfer task:} Apply expression logic to mutations,
  biotechnology, development, and disease diagnostics.
\end{itemize}

\subsubsection*{Mutations, CRISPR, and
Genomics}\label{mutations-crispr-and-genomics-1}
\addcontentsline{toc}{subsubsection}{Mutations, CRISPR, and Genomics}

\begin{itemize}
\tightlist
\item
  \textbf{Chapter:} \cref{sec:unit_IV_mutations_and_genomics}.
\item
  \textbf{Lab:} \cref{sec:lab_unit_IV_mutations_and_genomics}.
\item
  \textbf{Question bank:} \cref{sec:q_unit_IV_mutations_and_genomics}.
\item
  \textbf{Big idea:} Genomic variation becomes biological consequence
  through sequence context, repair, and selection.
\item
  \textbf{Core concepts:} mutation classes, DNA repair, CRISPR, genome
  analysis.
\item
  \textbf{Framework alignment:} Vision \& Change: Information flow,
  exchange, and storage, Structure and function; AP Biology: Information
  Storage and Transmission, Systems Interactions; NGSS-style topics:
  Inheritance and Variation of Traits, Structure and Function.
\item
  \textbf{Practice cadence:} Concept Explanation, Questions and Methods,
  Argumentation.
\item
  \textbf{Model/data skill:} Mutation-rate, edit-efficiency, and
  sequence-comparison calculations. Classify variants and predict likely
  molecular effect from sequence evidence.
\item
  \textbf{Misconception probe:} Not every mutation is harmful, and not
  every harmful mutation changes a protein sequence.
\item
  \textbf{Transfer task:} Transfer variant reasoning to cancer genomics,
  ancestry, gene therapy, or microbial evolution.
\end{itemize}

\subsubsection*{Epigenetics and Gene
Regulation}\label{epigenetics-and-gene-regulation-1}
\addcontentsline{toc}{subsubsection}{Epigenetics and Gene Regulation}

\begin{itemize}
\tightlist
\item
  \textbf{Chapter:} \cref{sec:unit_IV_epigenetics_and_gene_regulation}.
\item
  \textbf{Lab:} \cref{sec:lab_unit_IV_epigenetics_and_gene_regulation}.
\item
  \textbf{Question bank:}
  \cref{sec:q_unit_IV_epigenetics_and_gene_regulation}.
\item
  \textbf{Big idea:} Cells create stable yet reversible expression
  states through chromatin, DNA marks, and regulatory circuits.
\item
  \textbf{Core concepts:} chromatin, methylation, histone modification,
  enhancers.
\item
  \textbf{Framework alignment:} Vision \& Change: Information flow,
  exchange, and storage, Structure and function; AP Biology: Information
  Storage and Transmission, Systems Interactions; NGSS-style topics:
  Inheritance and Variation of Traits, Structure and Function.
\item
  \textbf{Practice cadence:} Concept Explanation, Questions and Methods,
  Argumentation.
\item
  \textbf{Model/data skill:} Regulatory-state and expression-ratio
  reasoning. Interpret chromatin or expression evidence from simple
  regulatory datasets.
\item
  \textbf{Misconception probe:} Epigenetic does not mean independent of
  DNA sequence or permanently inherited.
\item
  \textbf{Transfer task:} Apply regulation logic to differentiation,
  imprinting, cancer, or environmental responses.
\end{itemize}

\subsection*{Unit V}\label{unit-v}
\addcontentsline{toc}{subsection}{Unit V}

\subsubsection*{Mendelian Genetics and
Heredity}\label{mendelian-genetics-and-heredity-1}
\addcontentsline{toc}{subsubsection}{Mendelian Genetics and Heredity}

\begin{itemize}
\tightlist
\item
  \textbf{Chapter:} \cref{sec:unit_V_mendelian_genetics}.
\item
  \textbf{Lab:} \cref{sec:lab_unit_V_mendelian_genetics}.
\item
  \textbf{Question bank:} \cref{sec:q_unit_V_mendelian_genetics}.
\item
  \textbf{Big idea:} Inheritance patterns arise from chromosome
  behaviour, probability, and genotype-to-phenotype mapping.
\item
  \textbf{Core concepts:} segregation, independent assortment,
  pedigrees, epistasis.
\item
  \textbf{Framework alignment:} Vision \& Change: Information flow,
  exchange, and storage, Evolution; AP Biology: Information Storage and
  Transmission, Evolution; NGSS-style topics: Inheritance and Variation
  of Traits, Natural Selection and Evolution.
\item
  \textbf{Practice cadence:} Statistical Tests and Data Analysis,
  Representing and Describing Data.
\item
  \textbf{Model/data skill:} Punnett, binomial, pedigree, and chi-square
  calculations. Use family or cross data to infer genotype
  probabilities.
\item
  \textbf{Misconception probe:} Dominant does not mean common, stronger,
  or better.
\item
  \textbf{Transfer task:} Transfer inheritance reasoning to counseling,
  breeding, model organisms, and risk prediction.
\end{itemize}

\subsubsection*{Chromosomal Inheritance and
Linkage}\label{chromosomal-inheritance-and-linkage-1}
\addcontentsline{toc}{subsubsection}{Chromosomal Inheritance and
Linkage}

\begin{itemize}
\tightlist
\item
  \textbf{Chapter:} \cref{sec:unit_V_chromosomal_inheritance}.
\item
  \textbf{Lab:} \cref{sec:lab_unit_V_chromosomal_inheritance}.
\item
  \textbf{Question bank:} \cref{sec:q_unit_V_chromosomal_inheritance}.
\item
  \textbf{Big idea:} Genes travel on chromosomes, so linkage,
  recombination, and chromosome structure shape inheritance.
\item
  \textbf{Core concepts:} linkage, recombination, sex linkage,
  chromosomal rearrangements.
\item
  \textbf{Framework alignment:} Vision \& Change: Information flow,
  exchange, and storage, Evolution; AP Biology: Information Storage and
  Transmission, Evolution; NGSS-style topics: Inheritance and Variation
  of Traits, Natural Selection and Evolution.
\item
  \textbf{Practice cadence:} Statistical Tests and Data Analysis,
  Representing and Describing Data.
\item
  \textbf{Model/data skill:} Recombination frequency and three-point
  mapping. Infer gene order or chromosomal mechanism from offspring
  counts.
\item
  \textbf{Misconception probe:} Independent assortment applies to
  unlinked loci, not to every pair of genes.
\item
  \textbf{Transfer task:} Transfer linkage reasoning to disease mapping,
  breeding, and genome assemblies.
\end{itemize}

\subsubsection*{Population Genetics and Hardy-Weinberg
Equilibrium}\label{population-genetics-and-hardy-weinberg-equilibrium-1}
\addcontentsline{toc}{subsubsection}{Population Genetics and
Hardy-Weinberg Equilibrium}

\begin{itemize}
\tightlist
\item
  \textbf{Chapter:} \cref{sec:unit_V_population_genetics}.
\item
  \textbf{Lab:} \cref{sec:lab_unit_V_population_genetics}.
\item
  \textbf{Question bank:} \cref{sec:q_unit_V_population_genetics}.
\item
  \textbf{Big idea:} Allele frequencies change when assumptions about
  random mating, population size, and fitness are violated.
\item
  \textbf{Core concepts:} Hardy-Weinberg, selection, mutation, genetic
  drift.
\item
  \textbf{Framework alignment:} Vision \& Change: Information flow,
  exchange, and storage, Evolution; AP Biology: Information Storage and
  Transmission, Evolution; NGSS-style topics: Inheritance and Variation
  of Traits, Natural Selection and Evolution.
\item
  \textbf{Practice cadence:} Statistical Tests and Data Analysis,
  Representing and Describing Data.
\item
  \textbf{Model/data skill:} Hardy-Weinberg and allele-frequency
  recurrence calculations. Estimate genotype or allele frequencies from
  population data.
\item
  \textbf{Misconception probe:} Hardy-Weinberg is a null model, not a
  claim that populations do not evolve.
\item
  \textbf{Transfer task:} Apply population-genetic reasoning to
  screening, conservation, and pathogen evolution.
\end{itemize}

\subsection*{Unit VI}\label{unit-vi}
\addcontentsline{toc}{subsection}{Unit VI}

\subsubsection*{Evolution --- Theory, Natural Selection, and
Adaptation}\label{evolution-theory-natural-selection-and-adaptation-1}
\addcontentsline{toc}{subsubsection}{Evolution --- Theory, Natural
Selection, and Adaptation}

\begin{itemize}
\tightlist
\item
  \textbf{Chapter:} \cref{sec:unit_VI_evolution_and_selection}.
\item
  \textbf{Lab:} \cref{sec:lab_unit_VI_evolution_and_selection}.
\item
  \textbf{Question bank:} \cref{sec:q_unit_VI_evolution_and_selection}.
\item
  \textbf{Big idea:} Natural selection is differential reproductive
  success acting on heritable variation in context.
\item
  \textbf{Core concepts:} variation, fitness, adaptation, selection
  coefficient.
\item
  \textbf{Framework alignment:} Vision \& Change: Evolution, Systems; AP
  Biology: Evolution, Systems Interactions; NGSS-style topics: Natural
  Selection and Evolution, Interdependent Relationships in Ecosystems.
\item
  \textbf{Practice cadence:} Visual Representations, Statistical Tests
  and Data Analysis, Argumentation.
\item
  \textbf{Model/data skill:} Selection coefficient and allele-frequency
  trajectory calculations. Interpret fitness data and distinguish
  selection from other forces.
\item
  \textbf{Misconception probe:} Evolution is not goal-directed progress;
  it is local change in populations under constraints.
\item
  \textbf{Transfer task:} Transfer selection reasoning to antibiotics,
  pesticide resistance, cancer, or climate adaptation.
\end{itemize}

\subsubsection*{Genetic Drift, Gene Flow, and
Speciation}\label{genetic-drift-gene-flow-and-speciation-1}
\addcontentsline{toc}{subsubsection}{Genetic Drift, Gene Flow, and
Speciation}

\begin{itemize}
\tightlist
\item
  \textbf{Chapter:} \cref{sec:unit_VI_genetic_drift_and_speciation}.
\item
  \textbf{Lab:} \cref{sec:lab_unit_VI_genetic_drift_and_speciation}.
\item
  \textbf{Question bank:}
  \cref{sec:q_unit_VI_genetic_drift_and_speciation}.
\item
  \textbf{Big idea:} Chance, population structure, and barriers to gene
  flow can generate divergence even without adaptive change.
\item
  \textbf{Core concepts:} drift, effective population size, gene flow,
  speciation.
\item
  \textbf{Framework alignment:} Vision \& Change: Evolution, Systems; AP
  Biology: Evolution, Systems Interactions; NGSS-style topics: Natural
  Selection and Evolution, Interdependent Relationships in Ecosystems.
\item
  \textbf{Practice cadence:} Visual Representations, Statistical Tests
  and Data Analysis, Argumentation.
\item
  \textbf{Model/data skill:} Fixation probability, effective population
  size, and migration-selection balance. Distinguish stochastic from
  directional change in allele-frequency data.
\item
  \textbf{Misconception probe:} Random does not mean patternless;
  stochastic processes have predictable distributions.
\item
  \textbf{Transfer task:} Apply drift reasoning to endangered
  populations, founder effects, and island radiations.
\end{itemize}

\subsubsection*{Phylogenetics and the Tree of
Life}\label{phylogenetics-and-the-tree-of-life-1}
\addcontentsline{toc}{subsubsection}{Phylogenetics and the Tree of Life}

\begin{itemize}
\tightlist
\item
  \textbf{Chapter:} \cref{sec:unit_VI_phylogenetics}.
\item
  \textbf{Lab:} \cref{sec:lab_unit_VI_phylogenetics}.
\item
  \textbf{Question bank:} \cref{sec:q_unit_VI_phylogenetics}.
\item
  \textbf{Big idea:} Phylogenies are evidence-based hypotheses about
  ancestry, not ladders of progress.
\item
  \textbf{Core concepts:} homology, tree topology, parsimony, molecular
  clocks.
\item
  \textbf{Framework alignment:} Vision \& Change: Evolution, Systems; AP
  Biology: Evolution, Systems Interactions; NGSS-style topics: Natural
  Selection and Evolution, Interdependent Relationships in Ecosystems.
\item
  \textbf{Practice cadence:} Visual Representations, Statistical Tests
  and Data Analysis, Argumentation.
\item
  \textbf{Model/data skill:} Tree-distance, parsimony, and
  molecular-clock calculations. Read trees correctly and map traits or
  sequences onto branches.
\item
  \textbf{Misconception probe:} Living species are cousins, not
  ancestors of one another.
\item
  \textbf{Transfer task:} Transfer tree thinking to pathogens,
  conservation units, gene families, and development.
\end{itemize}

\subsection*{Unit VII}\label{unit-vii}
\addcontentsline{toc}{subsection}{Unit VII}

\subsubsection*{Bacteria, Archaea, and
Viruses}\label{bacteria-archaea-and-viruses-1}
\addcontentsline{toc}{subsubsection}{Bacteria, Archaea, and Viruses}

\begin{itemize}
\tightlist
\item
  \textbf{Chapter:} \cref{sec:unit_VII_bacteria_archaea_viruses}.
\item
  \textbf{Lab:} \cref{sec:lab_unit_VII_bacteria_archaea_viruses}.
\item
  \textbf{Question bank:}
  \cref{sec:q_unit_VII_bacteria_archaea_viruses}.
\item
  \textbf{Big idea:} Microbial diversity reflects different cell
  architectures, genomes, metabolisms, and evolutionary histories.
\item
  \textbf{Core concepts:} prokaryotic structure, archaea, viruses,
  horizontal gene transfer.
\item
  \textbf{Framework alignment:} Vision \& Change: Evolution, Systems,
  Structure and function; AP Biology: Evolution, Systems Interactions;
  NGSS-style topics: Structure and Function, Interdependent
  Relationships in Ecosystems.
\item
  \textbf{Practice cadence:} Questions and Methods, Representing and
  Describing Data, Argumentation.
\item
  \textbf{Model/data skill:} Growth-rate, genome-size, and dilution
  calculations. Interpret microbial observations from growth, sequence,
  or structural evidence.
\item
  \textbf{Misconception probe:} Viruses are not simply tiny bacteria;
  they use fundamentally different replication logic.
\item
  \textbf{Transfer task:} Apply microbial diversity reasoning to
  antibiotics, biotechnology, ecology, and outbreaks.
\end{itemize}

\subsubsection*{Microbial Ecology and the
Microbiome}\label{microbial-ecology-and-the-microbiome-1}
\addcontentsline{toc}{subsubsection}{Microbial Ecology and the
Microbiome}

\begin{itemize}
\tightlist
\item
  \textbf{Chapter:} \cref{sec:unit_VII_microbial_ecology}.
\item
  \textbf{Lab:} \cref{sec:lab_unit_VII_microbial_ecology}.
\item
  \textbf{Question bank:} \cref{sec:q_unit_VII_microbial_ecology}.
\item
  \textbf{Big idea:} Microbes organize communities, cycles, and host
  physiology through interactions and metabolism.
\item
  \textbf{Core concepts:} microbiomes, symbiosis, diversity,
  biogeochemical cycling.
\item
  \textbf{Framework alignment:} Vision \& Change: Evolution, Systems,
  Structure and function; AP Biology: Evolution, Systems Interactions;
  NGSS-style topics: Structure and Function, Interdependent
  Relationships in Ecosystems.
\item
  \textbf{Practice cadence:} Questions and Methods, Representing and
  Describing Data, Argumentation.
\item
  \textbf{Model/data skill:} Diversity indices, richness estimates, and
  interaction-network reasoning. Compute or interpret community metrics
  from abundance data.
\item
  \textbf{Misconception probe:} A microbiome is not automatically
  beneficial; context determines the effect.
\item
  \textbf{Transfer task:} Transfer microbial ecology to soils, oceans,
  digestion, disease, and climate feedbacks.
\end{itemize}

\subsubsection*{Infectious Disease and
Immunity}\label{infectious-disease-and-immunity-1}
\addcontentsline{toc}{subsubsection}{Infectious Disease and Immunity}

\begin{itemize}
\tightlist
\item
  \textbf{Chapter:} \cref{sec:unit_VII_infectious_disease}.
\item
  \textbf{Lab:} \cref{sec:lab_unit_VII_infectious_disease}.
\item
  \textbf{Question bank:} \cref{sec:q_unit_VII_infectious_disease}.
\item
  \textbf{Big idea:} Disease dynamics emerge from host susceptibility,
  pathogen traits, transmission networks, and immunity.
\item
  \textbf{Core concepts:} transmission, R0, immunity, vaccination.
\item
  \textbf{Framework alignment:} Vision \& Change: Evolution, Systems,
  Structure and function; AP Biology: Evolution, Systems Interactions;
  NGSS-style topics: Structure and Function, Interdependent
  Relationships in Ecosystems.
\item
  \textbf{Practice cadence:} Questions and Methods, Representing and
  Describing Data, Argumentation.
\item
  \textbf{Model/data skill:} R0, herd-immunity threshold, and SIR
  trajectory calculations. Interpret outbreak curves and intervention
  effects.
\item
  \textbf{Misconception probe:} R0 is not a fixed property of a pathogen
  alone; it depends on host behaviour and environment.
\item
  \textbf{Transfer task:} Transfer disease-dynamic reasoning to
  vaccination, antimicrobial resistance, and public-health policy.
\end{itemize}

\subsection*{Unit VIII}\label{unit-viii}
\addcontentsline{toc}{subsection}{Unit VIII}

\subsubsection*{Plant Structure, Water Relations, and
Transport}\label{plant-structure-water-relations-and-transport-1}
\addcontentsline{toc}{subsubsection}{Plant Structure, Water Relations,
and Transport}

\begin{itemize}
\tightlist
\item
  \textbf{Chapter:} \cref{sec:unit_VIII_plant_structure_and_water}.
\item
  \textbf{Lab:} \cref{sec:lab_unit_VIII_plant_structure_and_water}.
\item
  \textbf{Question bank:}
  \cref{sec:q_unit_VIII_plant_structure_and_water}.
\item
  \textbf{Big idea:} Plant form is an engineering solution to water
  movement, support, gas exchange, and growth.
\item
  \textbf{Core concepts:} xylem, phloem, water potential, transpiration.
\item
  \textbf{Framework alignment:} Vision \& Change: Structure and
  function, Pathways and transformations of energy and matter, Systems;
  AP Biology: Energetics, Systems Interactions; NGSS-style topics:
  Structure and Function, Matter and Energy in Organisms and Ecosystems.
\item
  \textbf{Practice cadence:} Visual Representations, Questions and
  Methods, Argumentation.
\item
  \textbf{Model/data skill:} Water-potential and transpiration-flux
  calculations. Interpret plant-water data from pressure, solute, and
  humidity measurements.
\item
  \textbf{Misconception probe:} Water does not move because plants pull
  with intention; it follows potential gradients and cohesion.
\item
  \textbf{Transfer task:} Transfer plant-water reasoning to drought,
  irrigation, forest physiology, and crop breeding.
\end{itemize}

\subsubsection*{Plant Reproduction and
Development}\label{plant-reproduction-and-development-1}
\addcontentsline{toc}{subsubsection}{Plant Reproduction and Development}

\begin{itemize}
\tightlist
\item
  \textbf{Chapter:} \cref{sec:unit_VIII_plant_reproduction}.
\item
  \textbf{Lab:} \cref{sec:lab_unit_VIII_plant_reproduction}.
\item
  \textbf{Question bank:} \cref{sec:q_unit_VIII_plant_reproduction}.
\item
  \textbf{Big idea:} Plant reproduction integrates life cycles,
  development, dispersal, and environmental timing.
\item
  \textbf{Core concepts:} alternation of generations, flowers, seeds,
  development.
\item
  \textbf{Framework alignment:} Vision \& Change: Structure and
  function, Pathways and transformations of energy and matter, Systems;
  AP Biology: Energetics, Systems Interactions; NGSS-style topics:
  Structure and Function, Matter and Energy in Organisms and Ecosystems.
\item
  \textbf{Practice cadence:} Visual Representations, Questions and
  Methods, Argumentation.
\item
  \textbf{Model/data skill:} Life-cycle accounting and
  phyllotaxis/growth-pattern calculations. Track ploidy, tissue origin,
  and reproductive stage from diagrams or observations.
\item
  \textbf{Misconception probe:} Pollen, spores, seeds, and gametes are
  not interchangeable terms.
\item
  \textbf{Transfer task:} Transfer reproductive reasoning to
  agriculture, pollination ecology, and plant evolution.
\end{itemize}

\subsubsection*{Plant Responses to the
Environment}\label{plant-responses-to-the-environment-1}
\addcontentsline{toc}{subsubsection}{Plant Responses to the Environment}

\begin{itemize}
\tightlist
\item
  \textbf{Chapter:} \cref{sec:unit_VIII_plant_responses}.
\item
  \textbf{Lab:} \cref{sec:lab_unit_VIII_plant_responses}.
\item
  \textbf{Question bank:} \cref{sec:q_unit_VIII_plant_responses}.
\item
  \textbf{Big idea:} Plants sense environmental signals and respond
  through growth, hormones, and physiological regulation.
\item
  \textbf{Core concepts:} tropisms, hormones, photoperiodism, stress
  responses.
\item
  \textbf{Framework alignment:} Vision \& Change: Structure and
  function, Pathways and transformations of energy and matter, Systems;
  AP Biology: Energetics, Systems Interactions; NGSS-style topics:
  Structure and Function, Matter and Energy in Organisms and Ecosystems.
\item
  \textbf{Practice cadence:} Visual Representations, Questions and
  Methods, Argumentation.
\item
  \textbf{Model/data skill:} Dose-response, water-use-efficiency, and
  hormone-interaction reasoning. Interpret plant response data across
  light, gravity, water, and hormone treatments.
\item
  \textbf{Misconception probe:} A plant response is not passive; plants
  actively regulate development and physiology without neurons.
\item
  \textbf{Transfer task:} Transfer response logic to shade avoidance,
  drought, flowering, and crop management.
\end{itemize}

\subsection*{Unit IX}\label{unit-ix}
\addcontentsline{toc}{subsection}{Unit IX}

\subsubsection*{Circulation, Respiration, and
Homeostasis}\label{circulation-respiration-and-homeostasis-1}
\addcontentsline{toc}{subsubsection}{Circulation, Respiration, and
Homeostasis}

\begin{itemize}
\tightlist
\item
  \textbf{Chapter:}
  \cref{sec:unit_IX_circulation_respiration_homeostasis}.
\item
  \textbf{Lab:}
  \cref{sec:lab_unit_IX_circulation_respiration_homeostasis}.
\item
  \textbf{Question bank:}
  \cref{sec:q_unit_IX_circulation_respiration_homeostasis}.
\item
  \textbf{Big idea:} Animal transport systems maintain gradients that
  let cells exchange gases, nutrients, heat, and wastes.
\item
  \textbf{Core concepts:} cardiac output, gas exchange, homeostasis,
  feedback.
\item
  \textbf{Framework alignment:} Vision \& Change: Structure and
  function, Systems; AP Biology: Systems Interactions, Energetics;
  NGSS-style topics: Structure and Function.
\item
  \textbf{Practice cadence:} Visual Representations, Statistical Tests
  and Data Analysis, Argumentation.
\item
  \textbf{Model/data skill:} Cardiac-output, diffusion,
  oxygen-saturation, and feedback calculations. Interpret physiological
  data from pressure, flow, saturation, or set-point changes.
\item
  \textbf{Misconception probe:} Homeostasis is dynamic regulation, not
  an unchanging internal state.
\item
  \textbf{Transfer task:} Transfer homeostatic reasoning to exercise,
  altitude, hemorrhage, fever, and shock.
\end{itemize}

\subsubsection*{Nervous System and Neural
Signalling}\label{nervous-system-and-neural-signalling-1}
\addcontentsline{toc}{subsubsection}{Nervous System and Neural
Signalling}

\begin{itemize}
\tightlist
\item
  \textbf{Chapter:} \cref{sec:unit_IX_nervous_system}.
\item
  \textbf{Lab:} \cref{sec:lab_unit_IX_nervous_system}.
\item
  \textbf{Question bank:} \cref{sec:q_unit_IX_nervous_system}.
\item
  \textbf{Big idea:} Nervous systems compute with cells whose structure
  shapes information flow and behaviour.
\item
  \textbf{Core concepts:} neurons, glia, circuits, sensory processing.
\item
  \textbf{Framework alignment:} Vision \& Change: Structure and
  function, Systems; AP Biology: Systems Interactions, Energetics;
  NGSS-style topics: Structure and Function.
\item
  \textbf{Practice cadence:} Visual Representations, Statistical Tests
  and Data Analysis, Argumentation.
\item
  \textbf{Model/data skill:} Cable-length, conduction, and simple
  circuit calculations. Interpret neural data from anatomy, timing, or
  lesion evidence.
\item
  \textbf{Misconception probe:} The brain is not only neurons; glia and
  circuit context are essential to function.
\item
  \textbf{Transfer task:} Transfer circuit reasoning to reflexes,
  sensory systems, learning, and disease.
\end{itemize}

\subsubsection*{Action Potentials and Synaptic
Transmission}\label{action-potentials-and-synaptic-transmission-1}
\addcontentsline{toc}{subsubsection}{Action Potentials and Synaptic
Transmission}

\begin{itemize}
\tightlist
\item
  \textbf{Chapter:} \cref{sec:unit_IX_action_potential_synapses}.
\item
  \textbf{Lab:} \cref{sec:lab_unit_IX_action_potential_synapses}.
\item
  \textbf{Question bank:}
  \cref{sec:q_unit_IX_action_potential_synapses}.
\item
  \textbf{Big idea:} Electrical excitability and synaptic transmission
  convert ion gradients into rapid communication.
\item
  \textbf{Core concepts:} action potentials, ion channels, synapses,
  plasticity.
\item
  \textbf{Framework alignment:} Vision \& Change: Structure and
  function, Systems; AP Biology: Systems Interactions, Energetics;
  NGSS-style topics: Structure and Function.
\item
  \textbf{Practice cadence:} Visual Representations, Statistical Tests
  and Data Analysis, Argumentation.
\item
  \textbf{Model/data skill:} Nernst/Goldman, Hodgkin-Huxley, and
  synaptic-current reasoning. Interpret voltage traces, conductance
  changes, and synaptic perturbations.
\item
  \textbf{Misconception probe:} An action potential is not electricity
  flowing like a wire; it is regenerated ion-channel dynamics.
\item
  \textbf{Transfer task:} Transfer excitability reasoning to anesthesia,
  toxins, epilepsy, and neuromuscular disease.
\end{itemize}

\subsubsection*{Endocrine and Immune
Systems}\label{endocrine-and-immune-systems-1}
\addcontentsline{toc}{subsubsection}{Endocrine and Immune Systems}

\begin{itemize}
\tightlist
\item
  \textbf{Chapter:} \cref{sec:unit_IX_endocrine_and_immune}.
\item
  \textbf{Lab:} \cref{sec:lab_unit_IX_endocrine_and_immune}.
\item
  \textbf{Question bank:} \cref{sec:q_unit_IX_endocrine_and_immune}.
\item
  \textbf{Big idea:} Long-range signaling and immune recognition
  coordinate body-wide adaptation and defense.
\item
  \textbf{Core concepts:} hormones, receptors, innate immunity, adaptive
  immunity.
\item
  \textbf{Framework alignment:} Vision \& Change: Structure and
  function, Systems; AP Biology: Systems Interactions, Energetics;
  NGSS-style topics: Structure and Function.
\item
  \textbf{Practice cadence:} Visual Representations, Statistical Tests
  and Data Analysis, Argumentation.
\item
  \textbf{Model/data skill:} Hormone feedback, dose-response, and
  immune-memory reasoning. Interpret endocrine or immune data from time
  courses, titers, or perturbations.
\item
  \textbf{Misconception probe:} Immunity is not just attack;
  recognition, tolerance, memory, and regulation are equally central.
\item
  \textbf{Transfer task:} Transfer signaling and immunity reasoning to
  vaccination, autoimmunity, stress, and metabolism.
\end{itemize}

\subsection*{Unit X}\label{unit-x}
\addcontentsline{toc}{subsection}{Unit X}

\subsubsection*{Population Ecology and Growth
Models}\label{population-ecology-and-growth-models-1}
\addcontentsline{toc}{subsubsection}{Population Ecology and Growth
Models}

\begin{itemize}
\tightlist
\item
  \textbf{Chapter:} \cref{sec:unit_X_population_ecology}.
\item
  \textbf{Lab:} \cref{sec:lab_unit_X_population_ecology}.
\item
  \textbf{Question bank:} \cref{sec:q_unit_X_population_ecology}.
\item
  \textbf{Big idea:} Population change reflects births, deaths,
  movement, age structure, and density dependence.
\item
  \textbf{Core concepts:} exponential growth, logistic growth, life
  tables, population viability.
\item
  \textbf{Framework alignment:} Vision \& Change: Systems, Evolution,
  Pathways and transformations of energy and matter; AP Biology: Systems
  Interactions, Evolution, Energetics; NGSS-style topics: Interdependent
  Relationships in Ecosystems, Matter and Energy in Organisms and
  Ecosystems, Natural Selection and Evolution.
\item
  \textbf{Practice cadence:} Representing and Describing Data,
  Statistical Tests and Data Analysis, Argumentation.
\item
  \textbf{Model/data skill:} Exponential/logistic growth,
  mark-recapture, and matrix projection. Use abundance or age-structure
  data to estimate growth and risk.
\item
  \textbf{Misconception probe:} Carrying capacity is not a fixed magic
  number; it changes with resources, interactions, and disturbance.
\item
  \textbf{Transfer task:} Transfer population models to fisheries,
  invasive species, epidemiology, and endangered species.
\end{itemize}

\subsubsection*{Community Ecology and Species
Interactions}\label{community-ecology-and-species-interactions-1}
\addcontentsline{toc}{subsubsection}{Community Ecology and Species
Interactions}

\begin{itemize}
\tightlist
\item
  \textbf{Chapter:} \cref{sec:unit_X_community_ecology}.
\item
  \textbf{Lab:} \cref{sec:lab_unit_X_community_ecology}.
\item
  \textbf{Question bank:} \cref{sec:q_unit_X_community_ecology}.
\item
  \textbf{Big idea:} Communities are structured by interactions among
  species and by the context of those interactions.
\item
  \textbf{Core concepts:} competition, predation, mutualism, succession.
\item
  \textbf{Framework alignment:} Vision \& Change: Systems, Evolution,
  Pathways and transformations of energy and matter; AP Biology: Systems
  Interactions, Evolution, Energetics; NGSS-style topics: Interdependent
  Relationships in Ecosystems, Matter and Energy in Organisms and
  Ecosystems, Natural Selection and Evolution.
\item
  \textbf{Practice cadence:} Representing and Describing Data,
  Statistical Tests and Data Analysis, Argumentation.
\item
  \textbf{Model/data skill:} Lotka-Volterra-style interaction and
  diversity-pattern reasoning. Interpret abundance, interaction, or
  disturbance data from communities.
\item
  \textbf{Misconception probe:} A species interaction is not permanently
  good or bad; the sign can change with context.
\item
  \textbf{Transfer task:} Transfer interaction reasoning to restoration,
  disease ecology, agriculture, and invasion biology.
\end{itemize}

\subsubsection*{Ecosystem Ecology and Biogeochemical
Cycles}\label{ecosystem-ecology-and-biogeochemical-cycles-1}
\addcontentsline{toc}{subsubsection}{Ecosystem Ecology and
Biogeochemical Cycles}

\begin{itemize}
\tightlist
\item
  \textbf{Chapter:} \cref{sec:unit_X_ecosystem_ecology}.
\item
  \textbf{Lab:} \cref{sec:lab_unit_X_ecosystem_ecology}.
\item
  \textbf{Question bank:} \cref{sec:q_unit_X_ecosystem_ecology}.
\item
  \textbf{Big idea:} Ecosystems couple energy flow and matter cycling
  across organisms, environments, and time.
\item
  \textbf{Core concepts:} primary productivity, trophic efficiency,
  nutrient cycles, decomposition.
\item
  \textbf{Framework alignment:} Vision \& Change: Systems, Evolution,
  Pathways and transformations of energy and matter; AP Biology: Systems
  Interactions, Evolution, Energetics; NGSS-style topics: Interdependent
  Relationships in Ecosystems, Matter and Energy in Organisms and
  Ecosystems, Natural Selection and Evolution.
\item
  \textbf{Practice cadence:} Representing and Describing Data,
  Statistical Tests and Data Analysis, Argumentation.
\item
  \textbf{Model/data skill:} Energy-flow, productivity, and
  nutrient-budget calculations. Trace matter and energy through food
  webs and biogeochemical cycles.
\item
  \textbf{Misconception probe:} Energy flows through ecosystems, but
  matter cycles; confusing the two breaks many explanations.
\item
  \textbf{Transfer task:} Transfer ecosystem reasoning to
  eutrophication, carbon budgets, agriculture, and climate feedback.
\end{itemize}

\subsubsection*{Biomes and Conservation
Biology}\label{biomes-and-conservation-biology-1}
\addcontentsline{toc}{subsubsection}{Biomes and Conservation Biology}

\begin{itemize}
\tightlist
\item
  \textbf{Chapter:} \cref{sec:unit_X_biomes_and_conservation}.
\item
  \textbf{Lab:} \cref{sec:lab_unit_X_biomes_and_conservation}.
\item
  \textbf{Question bank:} \cref{sec:q_unit_X_biomes_and_conservation}.
\item
  \textbf{Big idea:} Biodiversity patterns and conservation decisions
  emerge from climate, history, disturbance, and human choice.
\item
  \textbf{Core concepts:} biomes, biodiversity, extinction risk,
  conservation planning.
\item
  \textbf{Framework alignment:} Vision \& Change: Systems, Evolution,
  Pathways and transformations of energy and matter; AP Biology: Systems
  Interactions, Evolution, Energetics; NGSS-style topics: Interdependent
  Relationships in Ecosystems, Matter and Energy in Organisms and
  Ecosystems, Natural Selection and Evolution.
\item
  \textbf{Practice cadence:} Representing and Describing Data,
  Statistical Tests and Data Analysis, Argumentation.
\item
  \textbf{Model/data skill:} Species-area, risk, and prioritization
  calculations. Use maps, trend data, and threat categories to justify
  conservation priorities.
\item
  \textbf{Misconception probe:} Conservation is decision-making under
  constraints and uncertainty, not only preserving untouched nature.
\item
  \textbf{Transfer task:} Transfer conservation reasoning to land use,
  climate corridors, restoration, and environmental justice.
\end{itemize}

\newpage

\section*{Appendix B --- Instructor Orchestration
Guide}\label{appendix-b-instructor-orchestration-guide}
\addcontentsline{toc}{section}{Appendix B --- Instructor Orchestration
Guide}

\label{sec:appendix_instructor_orchestration}

This appendix is generated from the same curriculum and alignment
records that populate chapter Study Blueprints. It gives instructors a
compact sequence for turning the textbook into lectures, labs, quizzes,
and transfer tasks without hand-rebuilding a course map.

\subsection*{Reusable Teaching Loop}\label{reusable-teaching-loop}
\addcontentsline{toc}{subsection}{Reusable Teaching Loop}

\begin{enumerate}
\def\labelenumi{\arabic{enumi}.}
\tightlist
\item
  \textbf{Launch from a phenomenon.} Begin with a case, dataset, model
  failure, or observed pattern that makes the chapter question
  necessary.
\item
  \textbf{Model the mechanism.} Ask students to draw, compute, simulate,
  or annotate the causal structure before reading the full explanation.
\item
  \textbf{Gather evidence.} Use the companion lab as the measurement and
  reproducibility surface for the chapter.
\item
  \textbf{Check transfer.} Use the question bank to sample recall,
  application, data interpretation, and argumentation.
\item
  \textbf{Close the loop.} Return to the transfer task and require a
  claim, evidence, reasoning, and uncertainty statement.
\end{enumerate}

\subsection*{Chapter Orchestration
Matrix}\label{chapter-orchestration-matrix}
\addcontentsline{toc}{subsection}{Chapter Orchestration Matrix}

\subsubsection*{Unit 0}\label{unit-0-1}
\addcontentsline{toc}{subsubsection}{Unit 0}

\paragraph*{Systems Science and the Logic of
Emergence}\label{systems-science-and-the-logic-of-emergence-2}
\addcontentsline{toc}{paragraph}{Systems Science and the Logic of
Emergence}

\begin{itemize}
\tightlist
\item
  \textbf{Core thread:} Use systems diagrams, uncertainty, and feedback
  as recurring sense-making tools.
\item
  \textbf{Instructor move:} Launch from a phenomenon, have students
  model it with Box-and-arrow causal models with explicit inputs,
  outputs, and feedback signs, then test the model through data.
\item
  \textbf{Formative check:} Ask students to translate a verbal
  biological system into variables, links, and testable predictions and
  then answer: given a biological scenario, identify the system
  boundary, feedback loop, and missing measurement.
\item
  \textbf{Summative product:} Apply the same feedback map to a cell,
  organism, and ecosystem, then name what changes at each scale.
\item
  \textbf{Lab/question pair:} \cref{sec:lab_unit_0_systems_science};
  \cref{sec:q_unit_0_systems_science}.
\item
  \textbf{External alignment:} Vision \& Change: Systems, Structure and
  function; AP Biology: Systems Interactions; NGSS-style topics:
  Structure and Function, Interdependent Relationships in Ecosystems.
\end{itemize}

\paragraph*{Complex Adaptive Systems}\label{complex-adaptive-systems-2}
\addcontentsline{toc}{paragraph}{Complex Adaptive Systems}

\begin{itemize}
\tightlist
\item
  \textbf{Core thread:} Use systems diagrams, uncertainty, and feedback
  as recurring sense-making tools.
\item
  \textbf{Instructor move:} Launch from a phenomenon, have students
  model it with Agent-rule and scaling arguments for emergent biological
  patterns, then test the model through data.
\item
  \textbf{Formative check:} Ask students to distinguish deterministic
  trends from stochastic variation in repeated simulations and then
  answer: explain an observed population pattern from individual-level
  rules and identify one falsifiable prediction.
\item
  \textbf{Summative product:} Compare flocking, immune activation, and
  microbial biofilms as adaptive systems.
\item
  \textbf{Lab/question pair:}
  \cref{sec:lab_unit_0_complex_adaptive_systems};
  \cref{sec:q_unit_0_complex_adaptive_systems}.
\item
  \textbf{External alignment:} Vision \& Change: Systems, Structure and
  function; AP Biology: Systems Interactions; NGSS-style topics:
  Structure and Function, Interdependent Relationships in Ecosystems.
\end{itemize}

\paragraph*{Active Inference and the Free Energy
Principle}\label{active-inference-and-the-free-energy-principle-2}
\addcontentsline{toc}{paragraph}{Active Inference and the Free Energy
Principle}

\begin{itemize}
\tightlist
\item
  \textbf{Core thread:} Use systems diagrams, uncertainty, and feedback
  as recurring sense-making tools.
\item
  \textbf{Instructor move:} Launch from a phenomenon, have students
  model it with Bayesian belief updating and expected-free-energy-style
  policy comparison, then test the model through data.
\item
  \textbf{Formative check:} Ask students to read a small probability
  table and update a prediction after new evidence and then answer:
  calculate a posterior belief and explain what action would reduce
  uncertainty or restore a set point.
\item
  \textbf{Summative product:} Map prediction-error reasoning onto
  chemotaxis, thermoregulation, and attention.
\item
  \textbf{Lab/question pair:} \cref{sec:lab_unit_0_active_inference};
  \cref{sec:q_unit_0_active_inference}.
\item
  \textbf{External alignment:} Vision \& Change: Systems, Structure and
  function; AP Biology: Systems Interactions; NGSS-style topics:
  Structure and Function, Interdependent Relationships in Ecosystems.
\end{itemize}

\subsubsection*{Unit I}\label{unit-i-1}
\addcontentsline{toc}{subsubsection}{Unit I}

\paragraph*{Atoms, Molecules, and Chemical
Bonds}\label{atoms-molecules-and-chemical-bonds-2}
\addcontentsline{toc}{paragraph}{Atoms, Molecules, and Chemical Bonds}

\begin{itemize}
\tightlist
\item
  \textbf{Core thread:} Connect chemical structure, water, polymers, and
  catalysis to biological mechanism.
\item
  \textbf{Instructor move:} Launch from a phenomenon, have students
  model it with Formal charge and electronegativity-difference
  reasoning, then test the model through data.
\item
  \textbf{Formative check:} Ask students to use tabular atomic data to
  predict polarity, solubility, and biological reactivity and then
  answer: predict bond polarity, formal charge, or isotope behaviour and
  connect it to a biological consequence.
\item
  \textbf{Summative product:} Use atomic reasoning to explain a medical
  tracer, enzyme cofactor, or membrane-solubility problem.
\item
  \textbf{Lab/question pair:} \cref{sec:lab_unit_I_atoms_molecules};
  \cref{sec:q_unit_I_atoms_molecules}.
\item
  \textbf{External alignment:} Vision \& Change: Structure and function,
  Pathways and transformations of energy and matter; AP Biology:
  Energetics, Systems Interactions; NGSS-style topics: Matter and Energy
  in Organisms and Ecosystems, Structure and Function.
\end{itemize}

\paragraph*{Water --- The Molecule of
Life}\label{water-the-molecule-of-life-2}
\addcontentsline{toc}{paragraph}{Water --- The Molecule of Life}

\begin{itemize}
\tightlist
\item
  \textbf{Core thread:} Connect chemical structure, water, polymers, and
  catalysis to biological mechanism.
\item
  \textbf{Instructor move:} Launch from a phenomenon, have students
  model it with Henderson-Hasselbalch and pH-scale calculations, then
  test the model through data.
\item
  \textbf{Formative check:} Ask students to convert between pH, hydrogen
  ion concentration, and buffer ratios and then answer: calculate a pH
  or buffer ratio and interpret what it means for a biological system.
\item
  \textbf{Summative product:} Transfer water-property reasoning to blood
  buffering, plant transport, or protein folding.
\item
  \textbf{Lab/question pair:} \cref{sec:lab_unit_I_water_and_life};
  \cref{sec:q_unit_I_water_and_life}.
\item
  \textbf{External alignment:} Vision \& Change: Structure and function,
  Pathways and transformations of energy and matter; AP Biology:
  Energetics, Systems Interactions; NGSS-style topics: Matter and Energy
  in Organisms and Ecosystems, Structure and Function.
\end{itemize}

\paragraph*{Biological
Macromolecules}\label{biological-macromolecules-2}
\addcontentsline{toc}{paragraph}{Biological Macromolecules}

\begin{itemize}
\tightlist
\item
  \textbf{Core thread:} Connect chemical structure, water, polymers, and
  catalysis to biological mechanism.
\item
  \textbf{Instructor move:} Launch from a phenomenon, have students
  model it with Polymerization energetics and sequence-to-structure
  reasoning, then test the model through data.
\item
  \textbf{Formative check:} Ask students to classify molecules from
  structural evidence rather than names alone and then answer: given a
  molecular feature, infer likely solubility, interaction partners, or
  cellular role.
\item
  \textbf{Summative product:} Explain how a mutation, lipid
  substitution, or glycosylation change propagates to phenotype.
\item
  \textbf{Lab/question pair:} \cref{sec:lab_unit_I_macromolecules};
  \cref{sec:q_unit_I_macromolecules}.
\item
  \textbf{External alignment:} Vision \& Change: Structure and function,
  Pathways and transformations of energy and matter; AP Biology:
  Energetics, Systems Interactions; NGSS-style topics: Matter and Energy
  in Organisms and Ecosystems, Structure and Function.
\end{itemize}

\paragraph*{Enzymes and the Kinetics of
Catalysis}\label{enzymes-and-the-kinetics-of-catalysis-2}
\addcontentsline{toc}{paragraph}{Enzymes and the Kinetics of Catalysis}

\begin{itemize}
\tightlist
\item
  \textbf{Core thread:} Connect chemical structure, water, polymers, and
  catalysis to biological mechanism.
\item
  \textbf{Instructor move:} Launch from a phenomenon, have students
  model it with Michaelis-Menten, Lineweaver-Burk, and
  inhibition-pattern calculations, then test the model through data.
\item
  \textbf{Formative check:} Ask students to fit or interpret enzyme-rate
  data and identify which parameter changed and then answer: infer Vmax,
  Km, or inhibition type from data and justify the mechanistic
  interpretation.
\item
  \textbf{Summative product:} Apply enzyme-kinetic reasoning to drug
  dosing, metabolic control, or diagnostic assays.
\item
  \textbf{Lab/question pair:}
  \cref{sec:lab_unit_I_enzymes_and_kinetics};
  \cref{sec:q_unit_I_enzymes_and_kinetics}.
\item
  \textbf{External alignment:} Vision \& Change: Structure and function,
  Pathways and transformations of energy and matter; AP Biology:
  Energetics, Systems Interactions; NGSS-style topics: Matter and Energy
  in Organisms and Ecosystems, Structure and Function.
\end{itemize}

\subsubsection*{Unit II}\label{unit-ii-1}
\addcontentsline{toc}{subsubsection}{Unit II}

\paragraph*{Cell Theory and Cell
Types}\label{cell-theory-and-cell-types-2}
\addcontentsline{toc}{paragraph}{Cell Theory and Cell Types}

\begin{itemize}
\tightlist
\item
  \textbf{Core thread:} Revisit boundaries, compartments, and signals as
  causal models of cellular work.
\item
  \textbf{Instructor move:} Launch from a phenomenon, have students
  model it with Surface-area-to-volume scaling, then test the model
  through data.
\item
  \textbf{Formative check:} Ask students to infer cellular constraints
  from measurements, micrographs, and scale bars and then answer: use
  size or microscopy evidence to classify a cell and predict a
  constraint on exchange.
\item
  \textbf{Summative product:} Transfer scale reasoning to eggs, neurons,
  root hairs, and microbial colonies.
\item
  \textbf{Lab/question pair:} \cref{sec:lab_unit_II_cell_theory};
  \cref{sec:q_unit_II_cell_theory}.
\item
  \textbf{External alignment:} Vision \& Change: Structure and function,
  Systems, Information flow, exchange, and storage; AP Biology: Systems
  Interactions, Information Storage and Transmission; NGSS-style topics:
  Structure and Function.
\end{itemize}

\paragraph*{Cell Structure and
Organelles}\label{cell-structure-and-organelles-2}
\addcontentsline{toc}{paragraph}{Cell Structure and Organelles}

\begin{itemize}
\tightlist
\item
  \textbf{Core thread:} Revisit boundaries, compartments, and signals as
  causal models of cellular work.
\item
  \textbf{Instructor move:} Launch from a phenomenon, have students
  model it with Compartment-flow and motor-transport calculations, then
  test the model through data.
\item
  \textbf{Formative check:} Ask students to trace a molecule through
  compartments using evidence from labels or perturbations and then
  answer: predict which organelle or cytoskeletal process is disrupted
  from a phenotype.
\item
  \textbf{Summative product:} Apply compartment logic to secretion,
  apoptosis, infection, or cell division.
\item
  \textbf{Lab/question pair:} \cref{sec:lab_unit_II_cell_structure};
  \cref{sec:q_unit_II_cell_structure}.
\item
  \textbf{External alignment:} Vision \& Change: Structure and function,
  Systems, Information flow, exchange, and storage; AP Biology: Systems
  Interactions, Information Storage and Transmission; NGSS-style topics:
  Structure and Function.
\end{itemize}

\paragraph*{Membrane Structure and
Transport}\label{membrane-structure-and-transport-2}
\addcontentsline{toc}{paragraph}{Membrane Structure and Transport}

\begin{itemize}
\tightlist
\item
  \textbf{Core thread:} Revisit boundaries, compartments, and signals as
  causal models of cellular work.
\item
  \textbf{Instructor move:} Launch from a phenomenon, have students
  model it with Nernst, Goldman, osmotic, and facilitated-transport
  calculations, then test the model through data.
\item
  \textbf{Formative check:} Ask students to interpret transport data
  from gradients, rates, and membrane potentials and then answer:
  calculate or predict ion movement and explain which force dominates.
\item
  \textbf{Summative product:} Transfer gradient logic to neurons,
  kidneys, roots, and mitochondrial membranes.
\item
  \textbf{Lab/question pair:} \cref{sec:lab_unit_II_membrane_transport};
  \cref{sec:q_unit_II_membrane_transport}.
\item
  \textbf{External alignment:} Vision \& Change: Structure and function,
  Systems, Information flow, exchange, and storage; AP Biology: Systems
  Interactions, Information Storage and Transmission; NGSS-style topics:
  Structure and Function.
\end{itemize}

\paragraph*{Cell Signalling and
Communication}\label{cell-signalling-and-communication-2}
\addcontentsline{toc}{paragraph}{Cell Signalling and Communication}

\begin{itemize}
\tightlist
\item
  \textbf{Core thread:} Revisit boundaries, compartments, and signals as
  causal models of cellular work.
\item
  \textbf{Instructor move:} Launch from a phenomenon, have students
  model it with Ligand-receptor occupancy and dose-response reasoning,
  then test the model through data.
\item
  \textbf{Formative check:} Ask students to read pathway diagrams and
  infer the effect of agonists, antagonists, or mutations and then
  answer: predict the outcome of perturbing one pathway component and
  justify the causal chain.
\item
  \textbf{Summative product:} Apply signaling logic to hormones,
  neurotransmitters, immune receptors, or cancer mutations.
\item
  \textbf{Lab/question pair:} \cref{sec:lab_unit_II_cell_signaling};
  \cref{sec:q_unit_II_cell_signaling}.
\item
  \textbf{External alignment:} Vision \& Change: Structure and function,
  Systems, Information flow, exchange, and storage; AP Biology: Systems
  Interactions, Information Storage and Transmission; NGSS-style topics:
  Structure and Function.
\end{itemize}

\subsubsection*{Unit III}\label{unit-iii-1}
\addcontentsline{toc}{subsubsection}{Unit III}

\paragraph*{Bioenergetics and Cellular
Respiration}\label{bioenergetics-and-cellular-respiration-2}
\addcontentsline{toc}{paragraph}{Bioenergetics and Cellular Respiration}

\begin{itemize}
\tightlist
\item
  \textbf{Core thread:} Track matter, electrons, and free energy through
  pathways and ecosystems.
\item
  \textbf{Instructor move:} Launch from a phenomenon, have students
  model it with Delta G, ATP yield, and electron-carrier accounting,
  then test the model through data.
\item
  \textbf{Formative check:} Ask students to track carbon, electrons, and
  ATP across a pathway and then answer: calculate an energy yield or
  redox consequence and connect it to cellular work.
\item
  \textbf{Summative product:} Transfer energy accounting to exercise,
  fermentation, hypoxia, and mitochondrial disease.
\item
  \textbf{Lab/question pair:}
  \cref{sec:lab_unit_III_bioenergetics_and_respiration};
  \cref{sec:q_unit_III_bioenergetics_and_respiration}.
\item
  \textbf{External alignment:} Vision \& Change: Pathways and
  transformations of energy and matter, Systems; AP Biology: Energetics,
  Systems Interactions; NGSS-style topics: Matter and Energy in
  Organisms and Ecosystems.
\end{itemize}

\paragraph*{Photosynthesis}\label{photosynthesis-2}
\addcontentsline{toc}{paragraph}{Photosynthesis}

\begin{itemize}
\tightlist
\item
  \textbf{Core thread:} Track matter, electrons, and free energy through
  pathways and ecosystems.
\item
  \textbf{Instructor move:} Launch from a phenomenon, have students
  model it with Photon, ATP/NADPH, and Calvin-cycle stoichiometry, then
  test the model through data.
\item
  \textbf{Formative check:} Ask students to interpret light-response and
  carbon-fixation data and then answer: compute carbon-fixation
  requirements and explain a limiting factor.
\item
  \textbf{Summative product:} Transfer photosynthetic constraints to
  crops, algae, climate, or ecosystem productivity.
\item
  \textbf{Lab/question pair:} \cref{sec:lab_unit_III_photosynthesis};
  \cref{sec:q_unit_III_photosynthesis}.
\item
  \textbf{External alignment:} Vision \& Change: Pathways and
  transformations of energy and matter, Systems; AP Biology: Energetics,
  Systems Interactions; NGSS-style topics: Matter and Energy in
  Organisms and Ecosystems.
\end{itemize}

\paragraph*{Metabolic Integration and
Regulation}\label{metabolic-integration-and-regulation-2}
\addcontentsline{toc}{paragraph}{Metabolic Integration and Regulation}

\begin{itemize}
\tightlist
\item
  \textbf{Core thread:} Track matter, electrons, and free energy through
  pathways and ecosystems.
\item
  \textbf{Instructor move:} Launch from a phenomenon, have students
  model it with Energy charge, control points, and pathway-flux
  comparisons, then test the model through data.
\item
  \textbf{Formative check:} Ask students to use pathway evidence to
  infer which metabolic state or tissue is active and then answer:
  predict how a change in substrate, hormone, or enzyme activity
  redirects flux.
\item
  \textbf{Summative product:} Apply metabolic network reasoning to
  diabetes, fasting, exercise, or cancer metabolism.
\item
  \textbf{Lab/question pair:}
  \cref{sec:lab_unit_III_metabolic_integration};
  \cref{sec:q_unit_III_metabolic_integration}.
\item
  \textbf{External alignment:} Vision \& Change: Pathways and
  transformations of energy and matter, Systems; AP Biology: Energetics,
  Systems Interactions; NGSS-style topics: Matter and Energy in
  Organisms and Ecosystems.
\end{itemize}

\subsubsection*{Unit IV}\label{unit-iv-1}
\addcontentsline{toc}{subsubsection}{Unit IV}

\paragraph*{DNA Replication and the Cell
Cycle}\label{dna-replication-and-the-cell-cycle-2}
\addcontentsline{toc}{paragraph}{DNA Replication and the Cell Cycle}

\begin{itemize}
\tightlist
\item
  \textbf{Core thread:} Connect sequence, regulation, genome integrity,
  and phenotype through evidence.
\item
  \textbf{Instructor move:} Launch from a phenomenon, have students
  model it with Replication-fork speed, error-rate, and cell-cycle
  timing calculations, then test the model through data.
\item
  \textbf{Formative check:} Ask students to interpret replication or
  cell-cycle data from timing, labeling, or checkpoint perturbations and
  then answer: calculate replication timing or error burden and infer
  checkpoint consequences.
\item
  \textbf{Summative product:} Transfer replication logic to cancer,
  aging, viral replication, or antibiotic targets.
\item
  \textbf{Lab/question pair:}
  \cref{sec:lab_unit_IV_dna_replication_and_cell_cycle};
  \cref{sec:q_unit_IV_dna_replication_and_cell_cycle}.
\item
  \textbf{External alignment:} Vision \& Change: Information flow,
  exchange, and storage, Structure and function; AP Biology: Information
  Storage and Transmission, Systems Interactions; NGSS-style topics:
  Inheritance and Variation of Traits, Structure and Function.
\end{itemize}

\paragraph*{Gene Expression --- Transcription and
Translation}\label{gene-expression-transcription-and-translation-2}
\addcontentsline{toc}{paragraph}{Gene Expression --- Transcription and
Translation}

\begin{itemize}
\tightlist
\item
  \textbf{Core thread:} Connect sequence, regulation, genome integrity,
  and phenotype through evidence.
\item
  \textbf{Instructor move:} Launch from a phenomenon, have students
  model it with Reading-frame, codon, and expression-output
  calculations, then test the model through data.
\item
  \textbf{Formative check:} Ask students to convert DNA/RNA sequence
  data into predicted molecular products and then answer: translate or
  annotate a sequence and explain how a regulatory change alters output.
\item
  \textbf{Summative product:} Apply expression logic to mutations,
  biotechnology, development, and disease diagnostics.
\item
  \textbf{Lab/question pair:} \cref{sec:lab_unit_IV_gene_expression};
  \cref{sec:q_unit_IV_gene_expression}.
\item
  \textbf{External alignment:} Vision \& Change: Information flow,
  exchange, and storage, Structure and function; AP Biology: Information
  Storage and Transmission, Systems Interactions; NGSS-style topics:
  Inheritance and Variation of Traits, Structure and Function.
\end{itemize}

\paragraph*{Mutations, CRISPR, and
Genomics}\label{mutations-crispr-and-genomics-2}
\addcontentsline{toc}{paragraph}{Mutations, CRISPR, and Genomics}

\begin{itemize}
\tightlist
\item
  \textbf{Core thread:} Connect sequence, regulation, genome integrity,
  and phenotype through evidence.
\item
  \textbf{Instructor move:} Launch from a phenomenon, have students
  model it with Mutation-rate, edit-efficiency, and sequence-comparison
  calculations, then test the model through data.
\item
  \textbf{Formative check:} Ask students to classify variants and
  predict likely molecular effect from sequence evidence and then
  answer: predict the impact of a variant and justify the evidence
  needed to validate it.
\item
  \textbf{Summative product:} Transfer variant reasoning to cancer
  genomics, ancestry, gene therapy, or microbial evolution.
\item
  \textbf{Lab/question pair:}
  \cref{sec:lab_unit_IV_mutations_and_genomics};
  \cref{sec:q_unit_IV_mutations_and_genomics}.
\item
  \textbf{External alignment:} Vision \& Change: Information flow,
  exchange, and storage, Structure and function; AP Biology: Information
  Storage and Transmission, Systems Interactions; NGSS-style topics:
  Inheritance and Variation of Traits, Structure and Function.
\end{itemize}

\paragraph*{Epigenetics and Gene
Regulation}\label{epigenetics-and-gene-regulation-2}
\addcontentsline{toc}{paragraph}{Epigenetics and Gene Regulation}

\begin{itemize}
\tightlist
\item
  \textbf{Core thread:} Connect sequence, regulation, genome integrity,
  and phenotype through evidence.
\item
  \textbf{Instructor move:} Launch from a phenomenon, have students
  model it with Regulatory-state and expression-ratio reasoning, then
  test the model through data.
\item
  \textbf{Formative check:} Ask students to interpret chromatin or
  expression evidence from simple regulatory datasets and then answer:
  explain how a regulatory mark changes accessibility and predict
  downstream expression.
\item
  \textbf{Summative product:} Apply regulation logic to differentiation,
  imprinting, cancer, or environmental responses.
\item
  \textbf{Lab/question pair:}
  \cref{sec:lab_unit_IV_epigenetics_and_gene_regulation};
  \cref{sec:q_unit_IV_epigenetics_and_gene_regulation}.
\item
  \textbf{External alignment:} Vision \& Change: Information flow,
  exchange, and storage, Structure and function; AP Biology: Information
  Storage and Transmission, Systems Interactions; NGSS-style topics:
  Inheritance and Variation of Traits, Structure and Function.
\end{itemize}

\subsubsection*{Unit V}\label{unit-v-1}
\addcontentsline{toc}{subsubsection}{Unit V}

\paragraph*{Mendelian Genetics and
Heredity}\label{mendelian-genetics-and-heredity-2}
\addcontentsline{toc}{paragraph}{Mendelian Genetics and Heredity}

\begin{itemize}
\tightlist
\item
  \textbf{Core thread:} Use probability and population reasoning to
  connect inheritance to evolution.
\item
  \textbf{Instructor move:} Launch from a phenomenon, have students
  model it with Punnett, binomial, pedigree, and chi-square
  calculations, then test the model through data.
\item
  \textbf{Formative check:} Ask students to use family or cross data to
  infer genotype probabilities and then answer: compute a genetic
  probability or chi-square result and explain the biological
  assumption.
\item
  \textbf{Summative product:} Transfer inheritance reasoning to
  counseling, breeding, model organisms, and risk prediction.
\item
  \textbf{Lab/question pair:} \cref{sec:lab_unit_V_mendelian_genetics};
  \cref{sec:q_unit_V_mendelian_genetics}.
\item
  \textbf{External alignment:} Vision \& Change: Information flow,
  exchange, and storage, Evolution; AP Biology: Information Storage and
  Transmission, Evolution; NGSS-style topics: Inheritance and Variation
  of Traits, Natural Selection and Evolution.
\end{itemize}

\paragraph*{Chromosomal Inheritance and
Linkage}\label{chromosomal-inheritance-and-linkage-2}
\addcontentsline{toc}{paragraph}{Chromosomal Inheritance and Linkage}

\begin{itemize}
\tightlist
\item
  \textbf{Core thread:} Use probability and population reasoning to
  connect inheritance to evolution.
\item
  \textbf{Instructor move:} Launch from a phenomenon, have students
  model it with Recombination frequency and three-point mapping, then
  test the model through data.
\item
  \textbf{Formative check:} Ask students to infer gene order or
  chromosomal mechanism from offspring counts and then answer: use
  progeny classes to calculate map distances and infer the most likely
  gene order.
\item
  \textbf{Summative product:} Transfer linkage reasoning to disease
  mapping, breeding, and genome assemblies.
\item
  \textbf{Lab/question pair:}
  \cref{sec:lab_unit_V_chromosomal_inheritance};
  \cref{sec:q_unit_V_chromosomal_inheritance}.
\item
  \textbf{External alignment:} Vision \& Change: Information flow,
  exchange, and storage, Evolution; AP Biology: Information Storage and
  Transmission, Evolution; NGSS-style topics: Inheritance and Variation
  of Traits, Natural Selection and Evolution.
\end{itemize}

\paragraph*{Population Genetics and Hardy-Weinberg
Equilibrium}\label{population-genetics-and-hardy-weinberg-equilibrium-2}
\addcontentsline{toc}{paragraph}{Population Genetics and Hardy-Weinberg
Equilibrium}

\begin{itemize}
\tightlist
\item
  \textbf{Core thread:} Use probability and population reasoning to
  connect inheritance to evolution.
\item
  \textbf{Instructor move:} Launch from a phenomenon, have students
  model it with Hardy-Weinberg and allele-frequency recurrence
  calculations, then test the model through data.
\item
  \textbf{Formative check:} Ask students to estimate genotype or allele
  frequencies from population data and then answer: calculate expected
  frequencies and identify which evolutionary force could explain
  deviations.
\item
  \textbf{Summative product:} Apply population-genetic reasoning to
  screening, conservation, and pathogen evolution.
\item
  \textbf{Lab/question pair:} \cref{sec:lab_unit_V_population_genetics};
  \cref{sec:q_unit_V_population_genetics}.
\item
  \textbf{External alignment:} Vision \& Change: Information flow,
  exchange, and storage, Evolution; AP Biology: Information Storage and
  Transmission, Evolution; NGSS-style topics: Inheritance and Variation
  of Traits, Natural Selection and Evolution.
\end{itemize}

\subsubsection*{Unit VI}\label{unit-vi-1}
\addcontentsline{toc}{subsubsection}{Unit VI}

\paragraph*{Evolution --- Theory, Natural Selection, and
Adaptation}\label{evolution-theory-natural-selection-and-adaptation-2}
\addcontentsline{toc}{paragraph}{Evolution --- Theory, Natural
Selection, and Adaptation}

\begin{itemize}
\tightlist
\item
  \textbf{Core thread:} Return to variation, evidence, and historical
  inference across evolutionary scales.
\item
  \textbf{Instructor move:} Launch from a phenomenon, have students
  model it with Selection coefficient and allele-frequency trajectory
  calculations, then test the model through data.
\item
  \textbf{Formative check:} Ask students to interpret fitness data and
  distinguish selection from other forces and then answer: estimate
  selection strength from data and explain why the result is context
  dependent.
\item
  \textbf{Summative product:} Transfer selection reasoning to
  antibiotics, pesticide resistance, cancer, or climate adaptation.
\item
  \textbf{Lab/question pair:}
  \cref{sec:lab_unit_VI_evolution_and_selection};
  \cref{sec:q_unit_VI_evolution_and_selection}.
\item
  \textbf{External alignment:} Vision \& Change: Evolution, Systems; AP
  Biology: Evolution, Systems Interactions; NGSS-style topics: Natural
  Selection and Evolution, Interdependent Relationships in Ecosystems.
\end{itemize}

\paragraph*{Genetic Drift, Gene Flow, and
Speciation}\label{genetic-drift-gene-flow-and-speciation-2}
\addcontentsline{toc}{paragraph}{Genetic Drift, Gene Flow, and
Speciation}

\begin{itemize}
\tightlist
\item
  \textbf{Core thread:} Return to variation, evidence, and historical
  inference across evolutionary scales.
\item
  \textbf{Instructor move:} Launch from a phenomenon, have students
  model it with Fixation probability, effective population size, and
  migration-selection balance, then test the model through data.
\item
  \textbf{Formative check:} Ask students to distinguish stochastic from
  directional change in allele-frequency data and then answer: predict
  how population size or isolation changes divergence and fixation risk.
\item
  \textbf{Summative product:} Apply drift reasoning to endangered
  populations, founder effects, and island radiations.
\item
  \textbf{Lab/question pair:}
  \cref{sec:lab_unit_VI_genetic_drift_and_speciation};
  \cref{sec:q_unit_VI_genetic_drift_and_speciation}.
\item
  \textbf{External alignment:} Vision \& Change: Evolution, Systems; AP
  Biology: Evolution, Systems Interactions; NGSS-style topics: Natural
  Selection and Evolution, Interdependent Relationships in Ecosystems.
\end{itemize}

\paragraph*{Phylogenetics and the Tree of
Life}\label{phylogenetics-and-the-tree-of-life-2}
\addcontentsline{toc}{paragraph}{Phylogenetics and the Tree of Life}

\begin{itemize}
\tightlist
\item
  \textbf{Core thread:} Return to variation, evidence, and historical
  inference across evolutionary scales.
\item
  \textbf{Instructor move:} Launch from a phenomenon, have students
  model it with Tree-distance, parsimony, and molecular-clock
  calculations, then test the model through data.
\item
  \textbf{Formative check:} Ask students to read trees correctly and map
  traits or sequences onto branches and then answer: use a tree to infer
  shared ancestry, character transitions, or divergence time.
\item
  \textbf{Summative product:} Transfer tree thinking to pathogens,
  conservation units, gene families, and development.
\item
  \textbf{Lab/question pair:} \cref{sec:lab_unit_VI_phylogenetics};
  \cref{sec:q_unit_VI_phylogenetics}.
\item
  \textbf{External alignment:} Vision \& Change: Evolution, Systems; AP
  Biology: Evolution, Systems Interactions; NGSS-style topics: Natural
  Selection and Evolution, Interdependent Relationships in Ecosystems.
\end{itemize}

\subsubsection*{Unit VII}\label{unit-vii-1}
\addcontentsline{toc}{subsubsection}{Unit VII}

\paragraph*{Bacteria, Archaea, and
Viruses}\label{bacteria-archaea-and-viruses-2}
\addcontentsline{toc}{paragraph}{Bacteria, Archaea, and Viruses}

\begin{itemize}
\tightlist
\item
  \textbf{Core thread:} Use microbes to integrate evolution, ecology,
  host response, and public decisions.
\item
  \textbf{Instructor move:} Launch from a phenomenon, have students
  model it with Growth-rate, genome-size, and dilution calculations,
  then test the model through data.
\item
  \textbf{Formative check:} Ask students to interpret microbial
  observations from growth, sequence, or structural evidence and then
  answer: classify a microbe or virus from evidence and predict a likely
  vulnerability.
\item
  \textbf{Summative product:} Apply microbial diversity reasoning to
  antibiotics, biotechnology, ecology, and outbreaks.
\item
  \textbf{Lab/question pair:}
  \cref{sec:lab_unit_VII_bacteria_archaea_viruses};
  \cref{sec:q_unit_VII_bacteria_archaea_viruses}.
\item
  \textbf{External alignment:} Vision \& Change: Evolution, Systems,
  Structure and function; AP Biology: Evolution, Systems Interactions;
  NGSS-style topics: Structure and Function, Interdependent
  Relationships in Ecosystems.
\end{itemize}

\paragraph*{Microbial Ecology and the
Microbiome}\label{microbial-ecology-and-the-microbiome-2}
\addcontentsline{toc}{paragraph}{Microbial Ecology and the Microbiome}

\begin{itemize}
\tightlist
\item
  \textbf{Core thread:} Use microbes to integrate evolution, ecology,
  host response, and public decisions.
\item
  \textbf{Instructor move:} Launch from a phenomenon, have students
  model it with Diversity indices, richness estimates, and
  interaction-network reasoning, then test the model through data.
\item
  \textbf{Formative check:} Ask students to compute or interpret
  community metrics from abundance data and then answer: use community
  data to infer diversity, dominance, and a plausible ecological
  mechanism.
\item
  \textbf{Summative product:} Transfer microbial ecology to soils,
  oceans, digestion, disease, and climate feedbacks.
\item
  \textbf{Lab/question pair:} \cref{sec:lab_unit_VII_microbial_ecology};
  \cref{sec:q_unit_VII_microbial_ecology}.
\item
  \textbf{External alignment:} Vision \& Change: Evolution, Systems,
  Structure and function; AP Biology: Evolution, Systems Interactions;
  NGSS-style topics: Structure and Function, Interdependent
  Relationships in Ecosystems.
\end{itemize}

\paragraph*{Infectious Disease and
Immunity}\label{infectious-disease-and-immunity-2}
\addcontentsline{toc}{paragraph}{Infectious Disease and Immunity}

\begin{itemize}
\tightlist
\item
  \textbf{Core thread:} Use microbes to integrate evolution, ecology,
  host response, and public decisions.
\item
  \textbf{Instructor move:} Launch from a phenomenon, have students
  model it with R0, herd-immunity threshold, and SIR trajectory
  calculations, then test the model through data.
\item
  \textbf{Formative check:} Ask students to interpret outbreak curves
  and intervention effects and then answer: calculate a threshold or
  outbreak trajectory and identify the intervention lever.
\item
  \textbf{Summative product:} Transfer disease-dynamic reasoning to
  vaccination, antimicrobial resistance, and public-health policy.
\item
  \textbf{Lab/question pair:}
  \cref{sec:lab_unit_VII_infectious_disease};
  \cref{sec:q_unit_VII_infectious_disease}.
\item
  \textbf{External alignment:} Vision \& Change: Evolution, Systems,
  Structure and function; AP Biology: Evolution, Systems Interactions;
  NGSS-style topics: Structure and Function, Interdependent
  Relationships in Ecosystems.
\end{itemize}

\subsubsection*{Unit VIII}\label{unit-viii-1}
\addcontentsline{toc}{subsubsection}{Unit VIII}

\paragraph*{Plant Structure, Water Relations, and
Transport}\label{plant-structure-water-relations-and-transport-2}
\addcontentsline{toc}{paragraph}{Plant Structure, Water Relations, and
Transport}

\begin{itemize}
\tightlist
\item
  \textbf{Core thread:} Read plant form as an engineered compromise
  among transport, energy, and reproduction.
\item
  \textbf{Instructor move:} Launch from a phenomenon, have students
  model it with Water-potential and transpiration-flux calculations,
  then test the model through data.
\item
  \textbf{Formative check:} Ask students to interpret plant-water data
  from pressure, solute, and humidity measurements and then answer:
  calculate water potential and predict the direction of water movement.
\item
  \textbf{Summative product:} Transfer plant-water reasoning to drought,
  irrigation, forest physiology, and crop breeding.
\item
  \textbf{Lab/question pair:}
  \cref{sec:lab_unit_VIII_plant_structure_and_water};
  \cref{sec:q_unit_VIII_plant_structure_and_water}.
\item
  \textbf{External alignment:} Vision \& Change: Structure and function,
  Pathways and transformations of energy and matter, Systems; AP
  Biology: Energetics, Systems Interactions; NGSS-style topics:
  Structure and Function, Matter and Energy in Organisms and Ecosystems.
\end{itemize}

\paragraph*{Plant Reproduction and
Development}\label{plant-reproduction-and-development-2}
\addcontentsline{toc}{paragraph}{Plant Reproduction and Development}

\begin{itemize}
\tightlist
\item
  \textbf{Core thread:} Read plant form as an engineered compromise
  among transport, energy, and reproduction.
\item
  \textbf{Instructor move:} Launch from a phenomenon, have students
  model it with Life-cycle accounting and phyllotaxis/growth-pattern
  calculations, then test the model through data.
\item
  \textbf{Formative check:} Ask students to track ploidy, tissue origin,
  and reproductive stage from diagrams or observations and then answer:
  trace a plant life cycle and predict the consequence of disrupting one
  stage.
\item
  \textbf{Summative product:} Transfer reproductive reasoning to
  agriculture, pollination ecology, and plant evolution.
\item
  \textbf{Lab/question pair:}
  \cref{sec:lab_unit_VIII_plant_reproduction};
  \cref{sec:q_unit_VIII_plant_reproduction}.
\item
  \textbf{External alignment:} Vision \& Change: Structure and function,
  Pathways and transformations of energy and matter, Systems; AP
  Biology: Energetics, Systems Interactions; NGSS-style topics:
  Structure and Function, Matter and Energy in Organisms and Ecosystems.
\end{itemize}

\paragraph*{Plant Responses to the
Environment}\label{plant-responses-to-the-environment-2}
\addcontentsline{toc}{paragraph}{Plant Responses to the Environment}

\begin{itemize}
\tightlist
\item
  \textbf{Core thread:} Read plant form as an engineered compromise
  among transport, energy, and reproduction.
\item
  \textbf{Instructor move:} Launch from a phenomenon, have students
  model it with Dose-response, water-use-efficiency, and
  hormone-interaction reasoning, then test the model through data.
\item
  \textbf{Formative check:} Ask students to interpret plant response
  data across light, gravity, water, and hormone treatments and then
  answer: predict a plant response from a signal and identify the likely
  hormone or pathway.
\item
  \textbf{Summative product:} Transfer response logic to shade
  avoidance, drought, flowering, and crop management.
\item
  \textbf{Lab/question pair:} \cref{sec:lab_unit_VIII_plant_responses};
  \cref{sec:q_unit_VIII_plant_responses}.
\item
  \textbf{External alignment:} Vision \& Change: Structure and function,
  Pathways and transformations of energy and matter, Systems; AP
  Biology: Energetics, Systems Interactions; NGSS-style topics:
  Structure and Function, Matter and Energy in Organisms and Ecosystems.
\end{itemize}

\subsubsection*{Unit IX}\label{unit-ix-1}
\addcontentsline{toc}{subsubsection}{Unit IX}

\paragraph*{Circulation, Respiration, and
Homeostasis}\label{circulation-respiration-and-homeostasis-2}
\addcontentsline{toc}{paragraph}{Circulation, Respiration, and
Homeostasis}

\begin{itemize}
\tightlist
\item
  \textbf{Core thread:} Spiral homeostasis from organ systems to neural,
  endocrine, and immune decisions.
\item
  \textbf{Instructor move:} Launch from a phenomenon, have students
  model it with Cardiac-output, diffusion, oxygen-saturation, and
  feedback calculations, then test the model through data.
\item
  \textbf{Formative check:} Ask students to interpret physiological data
  from pressure, flow, saturation, or set-point changes and then answer:
  calculate a transport or feedback variable and explain the
  compensatory response.
\item
  \textbf{Summative product:} Transfer homeostatic reasoning to
  exercise, altitude, hemorrhage, fever, and shock.
\item
  \textbf{Lab/question pair:}
  \cref{sec:lab_unit_IX_circulation_respiration_homeostasis};
  \cref{sec:q_unit_IX_circulation_respiration_homeostasis}.
\item
  \textbf{External alignment:} Vision \& Change: Structure and function,
  Systems; AP Biology: Systems Interactions, Energetics; NGSS-style
  topics: Structure and Function.
\end{itemize}

\paragraph*{Nervous System and Neural
Signalling}\label{nervous-system-and-neural-signalling-2}
\addcontentsline{toc}{paragraph}{Nervous System and Neural Signalling}

\begin{itemize}
\tightlist
\item
  \textbf{Core thread:} Spiral homeostasis from organ systems to neural,
  endocrine, and immune decisions.
\item
  \textbf{Instructor move:} Launch from a phenomenon, have students
  model it with Cable-length, conduction, and simple circuit
  calculations, then test the model through data.
\item
  \textbf{Formative check:} Ask students to interpret neural data from
  anatomy, timing, or lesion evidence and then answer: predict
  behavioural or signal consequences from a circuit or structural
  perturbation.
\item
  \textbf{Summative product:} Transfer circuit reasoning to reflexes,
  sensory systems, learning, and disease.
\item
  \textbf{Lab/question pair:} \cref{sec:lab_unit_IX_nervous_system};
  \cref{sec:q_unit_IX_nervous_system}.
\item
  \textbf{External alignment:} Vision \& Change: Structure and function,
  Systems; AP Biology: Systems Interactions, Energetics; NGSS-style
  topics: Structure and Function.
\end{itemize}

\paragraph*{Action Potentials and Synaptic
Transmission}\label{action-potentials-and-synaptic-transmission-2}
\addcontentsline{toc}{paragraph}{Action Potentials and Synaptic
Transmission}

\begin{itemize}
\tightlist
\item
  \textbf{Core thread:} Spiral homeostasis from organ systems to neural,
  endocrine, and immune decisions.
\item
  \textbf{Instructor move:} Launch from a phenomenon, have students
  model it with Nernst/Goldman, Hodgkin-Huxley, and synaptic-current
  reasoning, then test the model through data.
\item
  \textbf{Formative check:} Ask students to interpret voltage traces,
  conductance changes, and synaptic perturbations and then answer:
  predict how changing an ion gradient or channel conductance alters
  excitability.
\item
  \textbf{Summative product:} Transfer excitability reasoning to
  anesthesia, toxins, epilepsy, and neuromuscular disease.
\item
  \textbf{Lab/question pair:}
  \cref{sec:lab_unit_IX_action_potential_synapses};
  \cref{sec:q_unit_IX_action_potential_synapses}.
\item
  \textbf{External alignment:} Vision \& Change: Structure and function,
  Systems; AP Biology: Systems Interactions, Energetics; NGSS-style
  topics: Structure and Function.
\end{itemize}

\paragraph*{Endocrine and Immune
Systems}\label{endocrine-and-immune-systems-2}
\addcontentsline{toc}{paragraph}{Endocrine and Immune Systems}

\begin{itemize}
\tightlist
\item
  \textbf{Core thread:} Spiral homeostasis from organ systems to neural,
  endocrine, and immune decisions.
\item
  \textbf{Instructor move:} Launch from a phenomenon, have students
  model it with Hormone feedback, dose-response, and immune-memory
  reasoning, then test the model through data.
\item
  \textbf{Formative check:} Ask students to interpret endocrine or
  immune data from time courses, titers, or perturbations and then
  answer: predict the consequence of changing a hormone, receptor,
  antigen, or immune-cell population.
\item
  \textbf{Summative product:} Transfer signaling and immunity reasoning
  to vaccination, autoimmunity, stress, and metabolism.
\item
  \textbf{Lab/question pair:}
  \cref{sec:lab_unit_IX_endocrine_and_immune};
  \cref{sec:q_unit_IX_endocrine_and_immune}.
\item
  \textbf{External alignment:} Vision \& Change: Structure and function,
  Systems; AP Biology: Systems Interactions, Energetics; NGSS-style
  topics: Structure and Function.
\end{itemize}

\subsubsection*{Unit X}\label{unit-x-1}
\addcontentsline{toc}{subsubsection}{Unit X}

\paragraph*{Population Ecology and Growth
Models}\label{population-ecology-and-growth-models-2}
\addcontentsline{toc}{paragraph}{Population Ecology and Growth Models}

\begin{itemize}
\tightlist
\item
  \textbf{Core thread:} Use ecological models to connect populations,
  interactions, matter, climate, and decisions.
\item
  \textbf{Instructor move:} Launch from a phenomenon, have students
  model it with Exponential/logistic growth, mark-recapture, and matrix
  projection, then test the model through data.
\item
  \textbf{Formative check:} Ask students to use abundance or
  age-structure data to estimate growth and risk and then answer:
  calculate a growth or sampling parameter and interpret its
  conservation meaning.
\item
  \textbf{Summative product:} Transfer population models to fisheries,
  invasive species, epidemiology, and endangered species.
\item
  \textbf{Lab/question pair:} \cref{sec:lab_unit_X_population_ecology};
  \cref{sec:q_unit_X_population_ecology}.
\item
  \textbf{External alignment:} Vision \& Change: Systems, Evolution,
  Pathways and transformations of energy and matter; AP Biology: Systems
  Interactions, Evolution, Energetics; NGSS-style topics: Interdependent
  Relationships in Ecosystems, Matter and Energy in Organisms and
  Ecosystems, Natural Selection and Evolution.
\end{itemize}

\paragraph*{Community Ecology and Species
Interactions}\label{community-ecology-and-species-interactions-2}
\addcontentsline{toc}{paragraph}{Community Ecology and Species
Interactions}

\begin{itemize}
\tightlist
\item
  \textbf{Core thread:} Use ecological models to connect populations,
  interactions, matter, climate, and decisions.
\item
  \textbf{Instructor move:} Launch from a phenomenon, have students
  model it with Lotka-Volterra-style interaction and diversity-pattern
  reasoning, then test the model through data.
\item
  \textbf{Formative check:} Ask students to interpret abundance,
  interaction, or disturbance data from communities and then answer:
  predict how changing one species or resource alters community
  structure.
\item
  \textbf{Summative product:} Transfer interaction reasoning to
  restoration, disease ecology, agriculture, and invasion biology.
\item
  \textbf{Lab/question pair:} \cref{sec:lab_unit_X_community_ecology};
  \cref{sec:q_unit_X_community_ecology}.
\item
  \textbf{External alignment:} Vision \& Change: Systems, Evolution,
  Pathways and transformations of energy and matter; AP Biology: Systems
  Interactions, Evolution, Energetics; NGSS-style topics: Interdependent
  Relationships in Ecosystems, Matter and Energy in Organisms and
  Ecosystems, Natural Selection and Evolution.
\end{itemize}

\paragraph*{Ecosystem Ecology and Biogeochemical
Cycles}\label{ecosystem-ecology-and-biogeochemical-cycles-2}
\addcontentsline{toc}{paragraph}{Ecosystem Ecology and Biogeochemical
Cycles}

\begin{itemize}
\tightlist
\item
  \textbf{Core thread:} Use ecological models to connect populations,
  interactions, matter, climate, and decisions.
\item
  \textbf{Instructor move:} Launch from a phenomenon, have students
  model it with Energy-flow, productivity, and nutrient-budget
  calculations, then test the model through data.
\item
  \textbf{Formative check:} Ask students to trace matter and energy
  through food webs and biogeochemical cycles and then answer: calculate
  a flux or efficiency and explain which pool or process limits the
  system.
\item
  \textbf{Summative product:} Transfer ecosystem reasoning to
  eutrophication, carbon budgets, agriculture, and climate feedback.
\item
  \textbf{Lab/question pair:} \cref{sec:lab_unit_X_ecosystem_ecology};
  \cref{sec:q_unit_X_ecosystem_ecology}.
\item
  \textbf{External alignment:} Vision \& Change: Systems, Evolution,
  Pathways and transformations of energy and matter; AP Biology: Systems
  Interactions, Evolution, Energetics; NGSS-style topics: Interdependent
  Relationships in Ecosystems, Matter and Energy in Organisms and
  Ecosystems, Natural Selection and Evolution.
\end{itemize}

\paragraph*{Biomes and Conservation
Biology}\label{biomes-and-conservation-biology-2}
\addcontentsline{toc}{paragraph}{Biomes and Conservation Biology}

\begin{itemize}
\tightlist
\item
  \textbf{Core thread:} Use ecological models to connect populations,
  interactions, matter, climate, and decisions.
\item
  \textbf{Instructor move:} Launch from a phenomenon, have students
  model it with Species-area, risk, and prioritization calculations,
  then test the model through data.
\item
  \textbf{Formative check:} Ask students to use maps, trend data, and
  threat categories to justify conservation priorities and then answer:
  evaluate a conservation scenario using evidence, tradeoffs, and
  uncertainty.
\item
  \textbf{Summative product:} Transfer conservation reasoning to land
  use, climate corridors, restoration, and environmental justice.
\item
  \textbf{Lab/question pair:}
  \cref{sec:lab_unit_X_biomes_and_conservation};
  \cref{sec:q_unit_X_biomes_and_conservation}.
\item
  \textbf{External alignment:} Vision \& Change: Systems, Evolution,
  Pathways and transformations of energy and matter; AP Biology: Systems
  Interactions, Evolution, Energetics; NGSS-style topics: Interdependent
  Relationships in Ecosystems, Matter and Energy in Organisms and
  Ecosystems, Natural Selection and Evolution.
\end{itemize}

\newpage

\section*{Appendix C --- Mathematical Review for
Biology}\label{appendix-c-mathematical-review-for-biology}
\addcontentsline{toc}{section}{Appendix C --- Mathematical Review for
Biology}

\label{sec:appendix_math_review}

\begin{quote}
\textbf{Appendix C} \texorpdfstring{\ensuremath{\cdot}}{.} Level 1/3 \texorpdfstring{\ensuremath{\cdot}}{.} 45 min read \texorpdfstring{\ensuremath{\cdot}}{.} Prerequisites: none \texorpdfstring{\ensuremath{\cdot}}{.}
Use as reference
\end{quote}

This appendix collects the mathematics a reader needs to work through
the quantitative chapters of this textbook. It is not a course in
mathematics --- it is a compact reference for the specific tools used in
the chapters. Each section is self-contained; dip in as needed.

\begin{center}\rule{0.5\linewidth}{0.5pt}\end{center}

\subsection*{C.1 Logarithms and the Natural
Log}\label{c.1-logarithms-and-the-natural-log}
\addcontentsline{toc}{subsection}{C.1 Logarithms and the Natural Log}

\subsubsection*{Definition}\label{definition-1}
\addcontentsline{toc}{subsubsection}{Definition}

\[ \log_b(x) = y \iff b^y = x  \label{eq:appendices_appendix_math_review_item_1}\]

Biology uses two bases heavily:

\begin{itemize}
\tightlist
\item
  \textbf{Base 10} (\(\log_{10}\)) --- pH, sound intensity, earthquake
  magnitude, powers-of-ten units (nM, \texorpdfstring{\ensuremath{\mu}}{mu}M, mM).
\item
  \textbf{Base \(e\)} (\(\ln\), natural log) --- most first-order
  kinetics (exponential growth, decay, enzyme kinetics), where
  \(e \approx 2.71828\).
\end{itemize}

Conversion: \(\log_{10}(x) = \ln(x) / \ln(10) \approx \ln(x) / 2.303\).

\subsubsection*{Key identities}\label{key-identities}
\addcontentsline{toc}{subsubsection}{Key identities}

{\def\LTcaptype{none} % do not increment counter
\begin{longtable}[]{@{}
  >{\raggedright\arraybackslash}p{(\linewidth - 2\tabcolsep) * \real{0.2667}}
  >{\raggedright\arraybackslash}p{(\linewidth - 2\tabcolsep) * \real{0.7333}}@{}}
\toprule\noalign{}
\begin{minipage}[b]{\linewidth}\raggedright
Identity
\end{minipage} & \begin{minipage}[b]{\linewidth}\raggedright
Biological application
\end{minipage} \\
\midrule\noalign{}
\endhead
\bottomrule\noalign{}
\endlastfoot
\(\log(ab) = \log a + \log b\) & Combining probabilities across linked
loci \\
\(\log(a/b) = \log a - \log b\) & Henderson-Hasselbalch
\(\text{pH} = \text{p}K_a + \log([\text{A}^-]/[\text{HA}])\) \\
\(\log(a^n) = n \log a\) & Hill equation --- slope in log-log plot \\
\(\log_b(b^x) = x\) & Inverting exponential growth:
\(t = \ln(N/N_0)/r\) \\
\end{longtable}
}

\subsubsection*{Worked example --- doubling
time}\label{worked-example-doubling-time}
\addcontentsline{toc}{subsubsection}{Worked example --- doubling time}

A bacterial culture grows at \(r = 0.5\) h\(^{-1}\). When does the
population double?

\[ N(t) = N_0 e^{rt} = 2 N_0 \implies e^{rt} = 2 \implies t = \frac{\ln 2}{r} = \frac{0.693}{0.5} \approx 1.39 \text{ h}  \label{eq:appendices_appendix_math_review_item_2}\]

The \textbf{``ln 2 trick''} recurs throughout biology: half-life
\(t_{1/2} = \ln 2 / k\) (decay); doubling time \(t_d = \ln 2 / r\)
(growth). Memorise \(\ln 2 \approx 0.693\).

\begin{center}\rule{0.5\linewidth}{0.5pt}\end{center}

\subsection*{C.2 Differential and Integral
Calculus}\label{c.2-differential-and-integral-calculus}
\addcontentsline{toc}{subsection}{C.2 Differential and Integral
Calculus}

Biology rarely requires advanced calculus, but the reader must be
comfortable with:

\subsubsection*{Derivatives as rates}\label{derivatives-as-rates}
\addcontentsline{toc}{subsubsection}{Derivatives as rates}

\[ \frac{dN}{dt} = \text{rate of change of } N \text{ over time}  \label{eq:appendices_appendix_math_review_item_3}\]

In population ecology, \(dN/dt > 0\) means the population is growing;
\(dN/dt = 0\) is equilibrium; \(dN/dt < 0\) is decline. In enzyme
kinetics, \(d[P]/dt\) = reaction velocity \(v\).

\subsubsection*{Integration as
accumulation}\label{integration-as-accumulation}
\addcontentsline{toc}{subsubsection}{Integration as accumulation}

\[ \int_0^T \frac{dN}{dt} dt = N(T) - N(0)  \label{eq:appendices_appendix_math_review_item_4}\]

The integral sums up infinitesimal contributions. Examples:

\begin{itemize}
\tightlist
\item
  Total drug dose delivered by an infusion:
  \(\text{Dose} = \int_0^T (\text{infusion rate}) dt\).
\item
  Net photosynthesis over a day: \(\int_0^{24\text{h}} P(t) dt\).
\end{itemize}

\subsubsection*{First-order ODE: exponential growth and
decay}\label{first-order-ode-exponential-growth-and-decay}
\addcontentsline{toc}{subsubsection}{First-order ODE: exponential growth
and decay}

\[ \frac{dN}{dt} = kN \implies N(t) = N_0 e^{kt}  \label{eq:appendices_appendix_math_review_item_5}\]

\begin{itemize}
\tightlist
\item
  \(k > 0\): growth with doubling time \(t_d = \ln 2 / k\).
\item
  \(k < 0\): decay with half-life \(t_{1/2} = \ln 2 / |k|\).
\end{itemize}

\subsubsection*{Logistic growth
(nonlinear)}\label{logistic-growth-nonlinear}
\addcontentsline{toc}{subsubsection}{Logistic growth (nonlinear)}

\[ \frac{dN}{dt} = rN\left(1 - \frac{N}{K}\right)  \label{eq:appendices_appendix_math_review_item_6}\]

The closed-form solution is:

\[ N(t) = \frac{K}{1 + \left(\frac{K-N_0}{N_0}\right) e^{-rt}}  \label{eq:appendices_appendix_math_review_item_7}\]

Inflection (maximum rate) at \(N = K/2\), giving the sigmoid shape.

\begin{center}\rule{0.5\linewidth}{0.5pt}\end{center}

\subsection*{C.3 Basic Probability}\label{c.3-basic-probability}
\addcontentsline{toc}{subsection}{C.3 Basic Probability}

\subsubsection*{Independent events}\label{independent-events}
\addcontentsline{toc}{subsubsection}{Independent events}

\(P(A \cap B) = P(A) \cdot P(B)\) when events are independent. Example:
the probability that both parents transmit the recessive allele at a
locus (\(Aa \times Aa\)) is \((1/2)(1/2) = 1/4\).

\subsubsection*{Conditional probability and Bayes'
theorem}\label{conditional-probability-and-bayes-theorem}
\addcontentsline{toc}{subsubsection}{Conditional probability and Bayes'
theorem}

\[ P(A \mid B) = \frac{P(A \cap B)}{P(B)}  \label{eq:appendices_appendix_math_review_item_8}\]

Bayes' theorem --- central to \cref{sec:unit_0_active_inference} and to
medical testing:

\[ P(\text{disease} \mid +) = \frac{P(+ \mid \text{disease}) \cdot P(\text{disease})}{P(+)}  \label{eq:appendices_appendix_math_review_item_9}\]

A highly sensitive test can still have a low positive predictive value
if the disease is rare --- a counterintuitive result with major clinical
consequences.

\subsubsection*{Binomial distribution}\label{binomial-distribution}
\addcontentsline{toc}{subsubsection}{Binomial distribution}

Number of successes in \(n\) independent trials, each with probability
\(p\):

\[ P(k \text{ successes}) = \binom{n}{k} p^k (1-p)^{n-k}  \label{eq:appendices_appendix_math_review_item_10}\]

Used for Punnett-square ratios, mark-recapture statistics, and
Hardy-Weinberg proportion testing.

\subsubsection*{Poisson distribution}\label{poisson-distribution}
\addcontentsline{toc}{subsubsection}{Poisson distribution}

For rare events at mean rate \texorpdfstring{\ensuremath{\lambda}}{lambda} per interval:

\[ P(k) = \frac{\lambda^k e^{-\lambda}}{k!}  \label{eq:appendices_appendix_math_review_item_11}\]

Used for spontaneous mutation counts per locus, bacterial cell counts in
Petri dishes (when plates are not overgrown), rare-disease incidence.

\begin{center}\rule{0.5\linewidth}{0.5pt}\end{center}

\subsection*{C.4 Dimensional Analysis}\label{c.4-dimensional-analysis}
\addcontentsline{toc}{subsection}{C.4 Dimensional Analysis}

Every physical quantity carries \textbf{units}. Units must balance on
both sides of every equation --- if they don't, the equation is wrong.

\subsubsection*{Example --- diffusion
distance}\label{example-diffusion-distance}
\addcontentsline{toc}{subsubsection}{Example --- diffusion distance}

Fick's rule for mean diffusion distance:

\[ x = \sqrt{2Dt}  \label{eq:appendices_appendix_math_review_item_12}\]

Check units: \(D\) has units of m\texorpdfstring{\textsuperscript{2}}{2} s\texorpdfstring{\textsuperscript{-}}{-}\texorpdfstring{\textsuperscript{1}}{1}; \(t\) has units of s; so
\(\sqrt{2Dt}\) has units of \(\sqrt{\text{m}^2} = \text{m}\). \texorpdfstring{\ensuremath{\checkmark}}{[ok]}

If a student wrote \(x = Dt\), units would be m\texorpdfstring{\textsuperscript{2}}{2} s\texorpdfstring{\textsuperscript{-}}{-}\texorpdfstring{\textsuperscript{1}}{1} \texorpdfstring{\ensuremath{\cdot}}{.} s = m\texorpdfstring{\textsuperscript{2}}{2} ---
dimensionally wrong.

\subsubsection*{Standard SI units used in this
book}\label{standard-si-units-used-in-this-book}
\addcontentsline{toc}{subsubsection}{Standard SI units used in this
book}

{\def\LTcaptype{none} % do not increment counter
\begin{longtable}[]{@{}
  >{\raggedright\arraybackslash}p{(\linewidth - 6\tabcolsep) * \real{0.1667}}
  >{\raggedright\arraybackslash}p{(\linewidth - 6\tabcolsep) * \real{0.1250}}
  >{\raggedright\arraybackslash}p{(\linewidth - 6\tabcolsep) * \real{0.1458}}
  >{\raggedright\arraybackslash}p{(\linewidth - 6\tabcolsep) * \real{0.5625}}@{}}
\toprule\noalign{}
\begin{minipage}[b]{\linewidth}\raggedright
Quantity
\end{minipage} & \begin{minipage}[b]{\linewidth}\raggedright
Symbol
\end{minipage} & \begin{minipage}[b]{\linewidth}\raggedright
SI unit
\end{minipage} & \begin{minipage}[b]{\linewidth}\raggedright
Typical biological magnitude
\end{minipage} \\
\midrule\noalign{}
\endhead
\bottomrule\noalign{}
\endlastfoot
Length & \(l\) & m & cell: 10 \texorpdfstring{\ensuremath{\mu}}{mu}m; neuron: 1 m \\
Mass & \(m\) & kg & ribosome: 3 \texorpdfstring{\ensuremath{\times}}{x} 10\texorpdfstring{\textsuperscript{-}}{-}\texorpdfstring{\textsuperscript{2}}{2}\texorpdfstring{\textsuperscript{1}}{1} kg; cell: 10\texorpdfstring{\textsuperscript{-}}{-}\texorpdfstring{\textsuperscript{1}}{1}\texorpdfstring{\textsuperscript{2}}{2} kg \\
Time & \(t\) & s & AP: 1 ms; cell cycle: 24 h \\
Amount & \(n\) & mol & ATP in a cell: 10\texorpdfstring{\textsuperscript{-}}{-}\texorpdfstring{\textsuperscript{1}}{1}\texorpdfstring{\textsuperscript{5}}{5} mol \\
Concentration & \([X]\) & mol L\texorpdfstring{\textsuperscript{-}}{-}\texorpdfstring{\textsuperscript{1}}{1} = M & cytoplasmic ATP: 5 \texorpdfstring{\ensuremath{\times}}{x} 10\texorpdfstring{\textsuperscript{-}}{-}\texorpdfstring{\textsuperscript{3}}{3} M \\
Energy & \(E\) & J & ATP hydrolysis: 5 \texorpdfstring{\ensuremath{\times}}{x} 10\texorpdfstring{\textsuperscript{-}}{-}\texorpdfstring{\textsuperscript{2}}{2}\texorpdfstring{\textsuperscript{0}}{0} J per molecule, \texorpdfstring{\ensuremath{\approx}}{~}30.5 kJ
mol\texorpdfstring{\textsuperscript{-}}{-}\texorpdfstring{\textsuperscript{1}}{1} \\
Force & \(F\) & N & myosin step force: 5 \texorpdfstring{\ensuremath{\times}}{x} 10\texorpdfstring{\textsuperscript{-}}{-}\texorpdfstring{\textsuperscript{1}}{1}\texorpdfstring{\textsuperscript{2}}{2} N = 5 pN \\
Electric potential & \(V\) & V & resting membrane: \texorpdfstring{\ensuremath{-}}{-}0.07 V = \texorpdfstring{\ensuremath{-}}{-}70 mV \\
\end{longtable}
}

\subsubsection*{Non-SI units common in
biology}\label{non-si-units-common-in-biology}
\addcontentsline{toc}{subsubsection}{Non-SI units common in biology}

{\def\LTcaptype{none} % do not increment counter
\begin{longtable}[]{@{}lll@{}}
\toprule\noalign{}
Unit & SI equivalent & Use \\
\midrule\noalign{}
\endhead
\bottomrule\noalign{}
\endlastfoot
mmHg (Torr) & 133.322 Pa & blood pressure, gas partial pressures \\
cal & 4.184 J & nutrition (``Calorie'' = kcal) \\
pH unit & \(-\log_{10}[\text{H}^+]\) in M & acidity \\
Da (dalton) & 1.661 \texorpdfstring{\ensuremath{\times}}{x} 10\texorpdfstring{\textsuperscript{-}}{-}\texorpdfstring{\textsuperscript{2}}{2}\texorpdfstring{\textsuperscript{7}}{7} kg & molecular mass \\
Kelvin (K) & °C + 273.15 & thermodynamics requires K, not °C \\
\end{longtable}
}

\begin{center}\rule{0.5\linewidth}{0.5pt}\end{center}

\subsection*{C.5 Linear Algebra
Miniature}\label{c.5-linear-algebra-miniature}
\addcontentsline{toc}{subsection}{C.5 Linear Algebra Miniature}

Primarily trace knowledge is required for this book, but the two ideas
below recur.

\subsubsection*{Matrices as transitions}\label{matrices-as-transitions}
\addcontentsline{toc}{subsubsection}{Matrices as transitions}

Allele-frequency change under migration between two populations is a 2 \texorpdfstring{\ensuremath{\times}}{x}
2 matrix multiplication. Markov transition matrices describe stochastic
gene expression (lac operon), state transitions in ion channels, and
Wright-Fisher drift.

\subsubsection*{Eigenvalues and
equilibrium}\label{eigenvalues-and-equilibrium}
\addcontentsline{toc}{subsubsection}{Eigenvalues and equilibrium}

The dominant eigenvalue \(\lambda_1\) of a Leslie matrix equals the
long-term finite rate of increase. Eigenvectors give the stable age
distribution. This is the formal underpinning of
\cref{sec:unit_X_population_ecology}'s demographic analysis.

\begin{center}\rule{0.5\linewidth}{0.5pt}\end{center}

\subsection*{C.6 Common Pitfalls in Biological
Math}\label{c.6-common-pitfalls-in-biological-math}
\addcontentsline{toc}{subsection}{C.6 Common Pitfalls in Biological
Math}

\begin{enumerate}
\def\labelenumi{\arabic{enumi}.}
\tightlist
\item
  \textbf{Confusing concentration and amount.} 1 nM is a concentration;
  1 nmol is an amount. A picomole of hormone in 5 L of blood is 0.2 pM.
\item
  \textbf{Units of rate constants.} First-order: s\texorpdfstring{\textsuperscript{-}}{-}\texorpdfstring{\textsuperscript{1}}{1}. Second-order: M\texorpdfstring{\textsuperscript{-}}{-}\texorpdfstring{\textsuperscript{1}}{1}
  s\texorpdfstring{\textsuperscript{-}}{-}\texorpdfstring{\textsuperscript{1}}{1}. Third-order: M\texorpdfstring{\textsuperscript{-}}{-}\texorpdfstring{\textsuperscript{2}}{2} s\texorpdfstring{\textsuperscript{-}}{-}\texorpdfstring{\textsuperscript{1}}{1}. Dimensional analysis typically exposes an
  order-of-magnitude error.
\item
  \textbf{Exponents and prefixes.} Milli- (10\texorpdfstring{\textsuperscript{-}}{-}\texorpdfstring{\textsuperscript{3}}{3}), micro- (10\texorpdfstring{\textsuperscript{-}}{-}\texorpdfstring{\textsuperscript{6}}{6}), nano-
  (10\texorpdfstring{\textsuperscript{-}}{-}\texorpdfstring{\textsuperscript{9}}{9}), pico- (10\texorpdfstring{\textsuperscript{-}}{-}\texorpdfstring{\textsuperscript{1}}{1}\texorpdfstring{\textsuperscript{2}}{2}), femto- (10\texorpdfstring{\textsuperscript{-}}{-}\texorpdfstring{\textsuperscript{1}}{1}\texorpdfstring{\textsuperscript{5}}{5}), atto- (10\texorpdfstring{\textsuperscript{-}}{-}\texorpdfstring{\textsuperscript{1}}{1}\texorpdfstring{\textsuperscript{8}}{8}). Shifting by
  three orders of magnitude is common in biology; single errors are
  catastrophic.
\item
  \textbf{Logarithmic vs linear differences.} A tenfold change in
  {[}H\texorpdfstring{\textsuperscript{+}}{+}{]} is a one-pH-unit change. ``Twice the acidity'' means \(0.3\)
  pH unit down, not \(2\) pH units.
\item
  \textbf{Error propagation.} Small errors in input variables compound
  through multi-step calculations. The textbook flags this at
  thermodynamics, kinetics, and evolutionary rate estimation in
  particular.
\end{enumerate}

\begin{center}\rule{0.5\linewidth}{0.5pt}\end{center}

\subsection*{Further Reading and Source
Notes}\label{further-reading-and-source-notes-33}
\addcontentsline{toc}{subsection}{Further Reading and Source Notes}

\begin{itemize}
\tightlist
\item
  Keener, J. \& Sneyd, J. (2009). \emph{Mathematical Physiology} (2nd
  ed.). Springer. --- comprehensive biological math.
\item
  Otto, S. P. \& Day, T. (2007). \emph{A Biologist's Guide to
  Mathematical Modeling in Ecology and Evolution}. Princeton. ---
  accessible entry point.
\item
  Strogatz, S. H. (2018). \emph{Nonlinear Dynamics and Chaos} (2nd ed.).
  Westview. --- the bifurcation and limit-cycle literature.
\end{itemize}

\emph{Module: reference only (no code)}

\newpage

\section*{Appendix D --- Units, Physical Constants, and Biological
Ranges}\label{appendix-d-units-physical-constants-and-biological-ranges}
\addcontentsline{toc}{section}{Appendix D --- Units, Physical Constants,
and Biological Ranges}

\label{sec:appendix_units_and_constants}

\begin{quote}
\textbf{Appendix D} \texorpdfstring{\ensuremath{\cdot}}{.} Level 1/3 \texorpdfstring{\ensuremath{\cdot}}{.} 20 min read \texorpdfstring{\ensuremath{\cdot}}{.} Prerequisites: none \texorpdfstring{\ensuremath{\cdot}}{.}
Use as reference
\end{quote}

A compact reference for units, constants, and characteristic magnitudes
encountered throughout the textbook. Use it during worked problems to
sanity-check dimensional analysis, convert between unit systems, and
verify that numerical answers fall in biologically plausible ranges.

\begin{center}\rule{0.5\linewidth}{0.5pt}\end{center}

\subsection*{D.1 SI Base and Derived
Units}\label{d.1-si-base-and-derived-units}
\addcontentsline{toc}{subsection}{D.1 SI Base and Derived Units}

{\def\LTcaptype{none} % do not increment counter
\begin{longtable}[]{@{}llll@{}}
\toprule\noalign{}
Quantity & SI unit & Symbol & In base units \\
\midrule\noalign{}
\endhead
\bottomrule\noalign{}
\endlastfoot
Length & metre & m & --- \\
Mass & kilogram & kg & --- \\
Time & second & s & --- \\
Amount of substance & mole & mol & --- \\
Temperature & kelvin & K & --- \\
Electric current & ampere & A & --- \\
Force & newton & N & kg \texorpdfstring{\ensuremath{\cdot}}{.} m \texorpdfstring{\ensuremath{\cdot}}{.} s\texorpdfstring{\textsuperscript{-}}{-}\texorpdfstring{\textsuperscript{2}}{2} \\
Energy & joule & J & kg \texorpdfstring{\ensuremath{\cdot}}{.} m\texorpdfstring{\textsuperscript{2}}{2} \texorpdfstring{\ensuremath{\cdot}}{.} s\texorpdfstring{\textsuperscript{-}}{-}\texorpdfstring{\textsuperscript{2}}{2} \\
Power & watt & W & J \texorpdfstring{\ensuremath{\cdot}}{.} s\texorpdfstring{\textsuperscript{-}}{-}\texorpdfstring{\textsuperscript{1}}{1} \\
Pressure & pascal & Pa & N \texorpdfstring{\ensuremath{\cdot}}{.} m\texorpdfstring{\textsuperscript{-}}{-}\texorpdfstring{\textsuperscript{2}}{2} \\
Electric charge & coulomb & C & A \texorpdfstring{\ensuremath{\cdot}}{.} s \\
Electric potential & volt & V & J \texorpdfstring{\ensuremath{\cdot}}{.} C\texorpdfstring{\textsuperscript{-}}{-}\texorpdfstring{\textsuperscript{1}}{1} \\
Electric conductance & siemens & S & \texorpdfstring{\ensuremath{\Omega}}{Omega}\texorpdfstring{\textsuperscript{-}}{-}\texorpdfstring{\textsuperscript{1}}{1} \\
Frequency & hertz & Hz & s\texorpdfstring{\textsuperscript{-}}{-}\texorpdfstring{\textsuperscript{1}}{1} \\
\end{longtable}
}

\subsubsection*{Metric prefixes}\label{metric-prefixes}
\addcontentsline{toc}{subsubsection}{Metric prefixes}

{\def\LTcaptype{none} % do not increment counter
\begin{longtable}[]{@{}
  >{\raggedright\arraybackslash}p{(\linewidth - 6\tabcolsep) * \real{0.1875}}
  >{\raggedright\arraybackslash}p{(\linewidth - 6\tabcolsep) * \real{0.1875}}
  >{\raggedright\arraybackslash}p{(\linewidth - 6\tabcolsep) * \real{0.1875}}
  >{\raggedright\arraybackslash}p{(\linewidth - 6\tabcolsep) * \real{0.4375}}@{}}
\toprule\noalign{}
\begin{minipage}[b]{\linewidth}\raggedright
Prefix
\end{minipage} & \begin{minipage}[b]{\linewidth}\raggedright
Symbol
\end{minipage} & \begin{minipage}[b]{\linewidth}\raggedright
Factor
\end{minipage} & \begin{minipage}[b]{\linewidth}\raggedright
Biological use
\end{minipage} \\
\midrule\noalign{}
\endhead
\bottomrule\noalign{}
\endlastfoot
giga- & G & 10\texorpdfstring{\textsuperscript{9}}{9} & Gb base pairs in a genome \\
mega- & M & 10\texorpdfstring{\textsuperscript{6}}{6} & MDa protein complex; Mb DNA fragment \\
kilo- & k & 10\texorpdfstring{\textsuperscript{3}}{3} & kJ mol\texorpdfstring{\textsuperscript{-}}{-}\texorpdfstring{\textsuperscript{1}}{1} (\texorpdfstring{\ensuremath{\Delta}}{Delta}G); kDa protein mass \\
centi- & c & 10\texorpdfstring{\textsuperscript{-}}{-}\texorpdfstring{\textsuperscript{2}}{2} & cM genetic map distance \\
milli- & m & 10\texorpdfstring{\textsuperscript{-}}{-}\texorpdfstring{\textsuperscript{3}}{3} & mM intracellular metabolites; ms AP duration \\
micro- & \texorpdfstring{\ensuremath{\mu}}{mu} & 10\texorpdfstring{\textsuperscript{-}}{-}\texorpdfstring{\textsuperscript{6}}{6} & \texorpdfstring{\ensuremath{\mu}}{mu}m cell size; \texorpdfstring{\ensuremath{\mu}}{mu}M low-abundance metabolite \\
nano- & n & 10\texorpdfstring{\textsuperscript{-}}{-}\texorpdfstring{\textsuperscript{9}}{9} & nm bond length; nM receptor ligand \\
pico- & p & 10\texorpdfstring{\textsuperscript{-}}{-}\texorpdfstring{\textsuperscript{1}}{1}\texorpdfstring{\textsuperscript{2}}{2} & pN molecular motor force; pM low-abundance
hormone \\
femto- & f & 10\texorpdfstring{\textsuperscript{-}}{-}\texorpdfstring{\textsuperscript{1}}{1}\texorpdfstring{\textsuperscript{5}}{5} & fmol amount in a single cell \\
atto- & a & 10\texorpdfstring{\textsuperscript{-}}{-}\texorpdfstring{\textsuperscript{1}}{1}\texorpdfstring{\textsuperscript{8}}{8} & amol; approaches single-molecule detection \\
\end{longtable}
}

\begin{center}\rule{0.5\linewidth}{0.5pt}\end{center}

\subsection*{D.2 Universal Physical
Constants}\label{d.2-universal-physical-constants}
\addcontentsline{toc}{subsection}{D.2 Universal Physical Constants}

{\def\LTcaptype{none} % do not increment counter
\begin{longtable}[]{@{}
  >{\raggedright\arraybackslash}p{(\linewidth - 4\tabcolsep) * \real{0.4211}}
  >{\raggedright\arraybackslash}p{(\linewidth - 4\tabcolsep) * \real{0.3158}}
  >{\raggedright\arraybackslash}p{(\linewidth - 4\tabcolsep) * \real{0.2632}}@{}}
\toprule\noalign{}
\begin{minipage}[b]{\linewidth}\raggedright
Constant
\end{minipage} & \begin{minipage}[b]{\linewidth}\raggedright
Symbol
\end{minipage} & \begin{minipage}[b]{\linewidth}\raggedright
Value
\end{minipage} \\
\midrule\noalign{}
\endhead
\bottomrule\noalign{}
\endlastfoot
Avogadro's number & \(N_A\) & \(6.022 \times 10^{23}\) mol\texorpdfstring{\textsuperscript{-}}{-}\texorpdfstring{\textsuperscript{1}}{1} \\
Boltzmann constant & \(k_B\) & \(1.381 \times 10^{-23}\) J K\texorpdfstring{\textsuperscript{-}}{-}\texorpdfstring{\textsuperscript{1}}{1} =
\(8.617 \times 10^{-5}\) eV K\texorpdfstring{\textsuperscript{-}}{-}\texorpdfstring{\textsuperscript{1}}{1} \\
Gas constant & \(R = N_A k_B\) & 8.314 J K\texorpdfstring{\textsuperscript{-}}{-}\texorpdfstring{\textsuperscript{1}}{1} mol\texorpdfstring{\textsuperscript{-}}{-}\texorpdfstring{\textsuperscript{1}}{1} = 1.987 cal K\texorpdfstring{\textsuperscript{-}}{-}\texorpdfstring{\textsuperscript{1}}{1}
mol\texorpdfstring{\textsuperscript{-}}{-}\texorpdfstring{\textsuperscript{1}}{1} \\
Faraday constant & \(F\) & \(9.649 \times 10^{4}\) C mol\texorpdfstring{\textsuperscript{-}}{-}\texorpdfstring{\textsuperscript{1}}{1} \\
Elementary charge & \(e\) & \(1.602 \times 10^{-19}\) C \\
Speed of light in vacuum & \(c\) & \(2.998 \times 10^{8}\) m s\texorpdfstring{\textsuperscript{-}}{-}\texorpdfstring{\textsuperscript{1}}{1} \\
Planck constant & \(h\) & \(6.626 \times 10^{-34}\) J s \\
Permittivity of free space & \(\varepsilon_0\) &
\(8.854 \times 10^{-12}\) F m\texorpdfstring{\textsuperscript{-}}{-}\texorpdfstring{\textsuperscript{1}}{1} \\
Standard atmospheric pressure & \(P_{\text{atm}}\) & 101,325 Pa = 760
Torr \\
Body temperature & \(T_{\text{body}}\) & 310.15 K = 37.00 °C \\
Thermal energy at body T & \(k_B T\) & \(4.28 \times 10^{-21}\) J \texorpdfstring{\ensuremath{\approx}}{~}
0.027 eV = 2.57 kJ mol\texorpdfstring{\textsuperscript{-}}{-}\texorpdfstring{\textsuperscript{1}}{1} \\
\end{longtable}
}

The rule ``\(k_B T \approx \tfrac{1}{40}\) eV
\(\approx \tfrac{2.5}{1000}\) kJ mol\texorpdfstring{\textsuperscript{-}}{-}\texorpdfstring{\textsuperscript{1}}{1} at room temperature'' is useful:
hydrogen bonds (about 20 kJ mol\texorpdfstring{\textsuperscript{-}}{-}\texorpdfstring{\textsuperscript{1}}{1}) are about 8
\(k_B T\), so they survive thermal fluctuations; van der Waals contacts
(about 1 kJ mol\texorpdfstring{\textsuperscript{-}}{-}\texorpdfstring{\textsuperscript{1}}{1}) are about 0.4 \(k_B T\), so each
one is barely stable but collectively numerous contacts summate.

\begin{center}\rule{0.5\linewidth}{0.5pt}\end{center}

\subsection*{D.3 Energy Conversions}\label{d.3-energy-conversions}
\addcontentsline{toc}{subsection}{D.3 Energy Conversions}

{\def\LTcaptype{none} % do not increment counter
\begin{longtable}[]{@{}lll@{}}
\toprule\noalign{}
From & To & Factor \\
\midrule\noalign{}
\endhead
\bottomrule\noalign{}
\endlastfoot
1 J & cal & 0.239 \\
1 cal & J & 4.184 \\
1 kcal & kJ & 4.184 \\
1 eV & J & \(1.602 \times 10^{-19}\) \\
1 eV & kJ mol\texorpdfstring{\textsuperscript{-}}{-}\texorpdfstring{\textsuperscript{1}}{1} & 96.485 \\
1 kJ mol\texorpdfstring{\textsuperscript{-}}{-}\texorpdfstring{\textsuperscript{1}}{1} & eV & 0.01036 \\
1 J & kg\texorpdfstring{\ensuremath{\cdot}}{.}m\texorpdfstring{\textsuperscript{2}}{2} s\texorpdfstring{\textsuperscript{-}}{-}\texorpdfstring{\textsuperscript{2}}{2} & 1 (definition) \\
1 \(k_B T\) at 37 °C & kJ mol\texorpdfstring{\textsuperscript{-}}{-}\texorpdfstring{\textsuperscript{1}}{1} & 2.577 \\
\end{longtable}
}

\subsubsection*{ATP hydrolysis
benchmarks}\label{atp-hydrolysis-benchmarks}
\addcontentsline{toc}{subsubsection}{ATP hydrolysis benchmarks}

\begin{itemize}
\tightlist
\item
  Standard \texorpdfstring{\ensuremath{\Delta}}{Delta}G°\texorpdfstring{'}{'} for ATP \texorpdfstring{\ensuremath{\to}}{->} ADP + P\(_i\): \textbf{\texorpdfstring{\ensuremath{-}}{-}30.5 kJ mol\texorpdfstring{\textsuperscript{-}}{-}\texorpdfstring{\textsuperscript{1}}{1}}.
\item
  Cellular \texorpdfstring{\ensuremath{\Delta}}{Delta}G (physiological {[}ATP{]}/{[}ADP{]}): \textbf{\texorpdfstring{\ensuremath{-}}{-}57 kJ
  mol\texorpdfstring{\textsuperscript{-}}{-}\texorpdfstring{\textsuperscript{1}}{1}}, roughly 22 \(k_B T\).
\item
  Number of ATP hydrolyses required to lift a 1 kg mass 1 m:
  about 300,000 --- illustrating why muscle needs 10\texorpdfstring{\textsuperscript{8}}{8} ATP
  molecules consumed per second.
\end{itemize}

\begin{center}\rule{0.5\linewidth}{0.5pt}\end{center}

\subsection*{D.4 Pressure Conversions}\label{d.4-pressure-conversions}
\addcontentsline{toc}{subsection}{D.4 Pressure Conversions}

{\def\LTcaptype{none} % do not increment counter
\begin{longtable}[]{@{}lll@{}}
\toprule\noalign{}
From & To & Factor \\
\midrule\noalign{}
\endhead
\bottomrule\noalign{}
\endlastfoot
1 atm & Pa & 101,325 \\
1 atm & mmHg (Torr) & 760 \\
1 atm & bar & 1.01325 \\
1 mmHg & kPa & 0.1333 \\
1 mmHg & Pa & 133.322 \\
\end{longtable}
}

\textbf{Blood pressure}: systolic 120 mmHg \texorpdfstring{\ensuremath{\approx}}{~} 16 kPa. \textbf{Arterial
O\texorpdfstring{\textsubscript{2}}{2}}: \(P_{\text{a}}O_2 \approx 100\) mmHg \texorpdfstring{\ensuremath{\approx}}{~} 13.3 kPa. \textbf{Venous
O\texorpdfstring{\textsubscript{2}}{2}}: \(P_{\text{v}}O_2 \approx 40\) mmHg \texorpdfstring{\ensuremath{\approx}}{~} 5.3 kPa.

\begin{center}\rule{0.5\linewidth}{0.5pt}\end{center}

\subsection*{D.5 Concentration
Conversions}\label{d.5-concentration-conversions}
\addcontentsline{toc}{subsection}{D.5 Concentration Conversions}

1 molar (M) = 1 mol L\texorpdfstring{\textsuperscript{-}}{-}\texorpdfstring{\textsuperscript{1}}{1} = 10\texorpdfstring{\textsuperscript{3}}{3} mmol L\texorpdfstring{\textsuperscript{-}}{-}\texorpdfstring{\textsuperscript{1}}{1} = 10\texorpdfstring{\textsuperscript{6}}{6} \texorpdfstring{\ensuremath{\mu}}{mu}mol L\texorpdfstring{\textsuperscript{-}}{-}\texorpdfstring{\textsuperscript{1}}{1} = 10\texorpdfstring{\textsuperscript{9}}{9} nmol L\texorpdfstring{\textsuperscript{-}}{-}\texorpdfstring{\textsuperscript{1}}{1} =
10\texorpdfstring{\textsuperscript{1}}{1}\texorpdfstring{\textsuperscript{2}}{2} pmol L\texorpdfstring{\textsuperscript{-}}{-}\texorpdfstring{\textsuperscript{1}}{1}.

{\def\LTcaptype{none} % do not increment counter
\begin{longtable}[]{@{}
  >{\raggedright\arraybackslash}p{(\linewidth - 2\tabcolsep) * \real{0.6500}}
  >{\raggedright\arraybackslash}p{(\linewidth - 2\tabcolsep) * \real{0.3500}}@{}}
\toprule\noalign{}
\begin{minipage}[b]{\linewidth}\raggedright
Concentration
\end{minipage} & \begin{minipage}[b]{\linewidth}\raggedright
Example
\end{minipage} \\
\midrule\noalign{}
\endhead
\bottomrule\noalign{}
\endlastfoot
pM (\(10^{-12}\) M) & Circulating hormones (insulin, oxytocin) \\
nM (\(10^{-9}\) M) & Vitamin B\texorpdfstring{\textsubscript{1}}{1}\texorpdfstring{\textsubscript{2}}{2}, transcription factor at a target
site \\
\texorpdfstring{\ensuremath{\mu}}{mu}M (\(10^{-6}\) M) & Glucose in blood (5 \texorpdfstring{\ensuremath{\times}}{x} 10\texorpdfstring{\textsuperscript{-}}{-}\texorpdfstring{\textsuperscript{3}}{3} M = 5 mM, upper end) \\
mM (\(10^{-3}\) M) & Major intracellular ions (Na\texorpdfstring{\textsuperscript{+}}{+}, K\texorpdfstring{\textsuperscript{+}}{+}), ATP
(about 5 mM) \\
M & Cytoplasm water (55 M); seawater NaCl (0.5 M) \\
\end{longtable}
}

Molarity of water:
\(M_{H_2O} = 1000 \text{ g L}^{-1} / 18 \text{ g mol}^{-1} = 55.5\) M.
This is why {[}H\texorpdfstring{\textsubscript{2}}{2}O{]} is treated as constant in biological equilibrium
expressions.

\begin{center}\rule{0.5\linewidth}{0.5pt}\end{center}

\subsection*{D.6 Biological Magnitudes --- Quick
Reference}\label{d.6-biological-magnitudes-quick-reference}
\addcontentsline{toc}{subsection}{D.6 Biological Magnitudes --- Quick
Reference}

\subsubsection*{Time scales}\label{time-scales}
\addcontentsline{toc}{subsubsection}{Time scales}

{\def\LTcaptype{none} % do not increment counter
\begin{longtable}[]{@{}ll@{}}
\toprule\noalign{}
Event & Time \\
\midrule\noalign{}
\endhead
\bottomrule\noalign{}
\endlastfoot
Bond vibration & 10\texorpdfstring{\textsuperscript{-}}{-}\texorpdfstring{\textsuperscript{1}}{1}\texorpdfstring{\textsuperscript{4}}{4} s \\
Enzyme turnover (\(k_{cat}\)) & 10\texorpdfstring{\textsuperscript{-}}{-}\texorpdfstring{\textsuperscript{3}}{3} -- 10\texorpdfstring{\textsuperscript{-}}{-}\texorpdfstring{\textsuperscript{6}}{6} s \\
Action potential & 10\texorpdfstring{\textsuperscript{-}}{-}\texorpdfstring{\textsuperscript{3}}{3} s \\
Synaptic transmission & 10\texorpdfstring{\textsuperscript{-}}{-}\texorpdfstring{\textsuperscript{3}}{3} s \\
Ribosome peptide-bond formation & 0.1 s (5 aa s\texorpdfstring{\textsuperscript{-}}{-}\texorpdfstring{\textsuperscript{1}}{1} per ribosome) \\
Protein folding & 10\texorpdfstring{\textsuperscript{-}}{-}\texorpdfstring{\textsuperscript{3}}{3} -- 1 s \\
Mitosis & 1 h \\
Cell cycle (dividing human cell) & 1 d \\
Hair growth & 1 cm month\texorpdfstring{\textsuperscript{-}}{-}\texorpdfstring{\textsuperscript{1}}{1} \\
Human generation time & about 25 yr \\
Mammalian species lifespan & 10\texorpdfstring{\textsuperscript{6}}{6} -- 10\texorpdfstring{\textsuperscript{7}}{7} yr \\
\end{longtable}
}

\subsubsection*{Length scales}\label{length-scales}
\addcontentsline{toc}{subsubsection}{Length scales}

{\def\LTcaptype{none} % do not increment counter
\begin{longtable}[]{@{}ll@{}}
\toprule\noalign{}
Feature & Size \\
\midrule\noalign{}
\endhead
\bottomrule\noalign{}
\endlastfoot
C--C bond & 0.15 nm \\
Protein (average globular) & 3 nm \\
Ribosome & 25 nm \\
Virus (small RNA) & 30 nm \\
Nuclear pore & 100 nm \\
Bacterium (\emph{E. coli}) & 1 \texorpdfstring{\ensuremath{\mu}}{mu}m \\
Mitochondrion & 2 \texorpdfstring{\ensuremath{\mu}}{mu}m \\
Eukaryotic cell (average) & 10--20 \texorpdfstring{\ensuremath{\mu}}{mu}m \\
Neuron axon (giraffe recurrent laryngeal) & 2 m \\
Animal (blue whale) & 30 m \\
Redwood tree & 100 m \\
\end{longtable}
}

\subsubsection*{Energy scales}\label{energy-scales}
\addcontentsline{toc}{subsubsection}{Energy scales}

{\def\LTcaptype{none} % do not increment counter
\begin{longtable}[]{@{}ll@{}}
\toprule\noalign{}
Interaction / event & Energy \\
\midrule\noalign{}
\endhead
\bottomrule\noalign{}
\endlastfoot
Van der Waals contact & about 1 kJ mol\texorpdfstring{\textsuperscript{-}}{-}\texorpdfstring{\textsuperscript{1}}{1} \\
Hydrogen bond & about 20 kJ mol\texorpdfstring{\textsuperscript{-}}{-}\texorpdfstring{\textsuperscript{1}}{1} \\
C--C covalent bond & about 350 kJ mol\texorpdfstring{\textsuperscript{-}}{-}\texorpdfstring{\textsuperscript{1}}{1} \\
ATP hydrolysis (standard) & \texorpdfstring{\ensuremath{-}}{-}30.5 kJ mol\texorpdfstring{\textsuperscript{-}}{-}\texorpdfstring{\textsuperscript{1}}{1} \\
ATP hydrolysis (cellular) & \texorpdfstring{\ensuremath{-}}{-}57 kJ mol\texorpdfstring{\textsuperscript{-}}{-}\texorpdfstring{\textsuperscript{1}}{1} \\
Glucose complete oxidation & \texorpdfstring{\ensuremath{-}}{-}2870 kJ mol\texorpdfstring{\textsuperscript{-}}{-}\texorpdfstring{\textsuperscript{1}}{1} \\
Photosynthesis per photon (680 nm) & 176 kJ mol\texorpdfstring{\textsuperscript{-}}{-}\texorpdfstring{\textsuperscript{1}}{1} \\
Daily human basal energy & 7000 kJ d\texorpdfstring{\textsuperscript{-}}{-}\texorpdfstring{\textsuperscript{1}}{1} (\texorpdfstring{\ensuremath{\approx}}{~}1700 kcal d\texorpdfstring{\textsuperscript{-}}{-}\texorpdfstring{\textsuperscript{1}}{1}) \\
\end{longtable}
}

\subsubsection*{Organismal physiology reference
scales}\label{organismal-physiology-reference-scales}
\addcontentsline{toc}{subsubsection}{Organismal physiology reference
scales}

{\def\LTcaptype{none} % do not increment counter
\begin{longtable}[]{@{}
  >{\raggedright\arraybackslash}p{(\linewidth - 4\tabcolsep) * \real{0.1538}}
  >{\raggedright\arraybackslash}p{(\linewidth - 4\tabcolsep) * \real{0.3077}}
  >{\raggedright\arraybackslash}p{(\linewidth - 4\tabcolsep) * \real{0.5385}}@{}}
\toprule\noalign{}
\begin{minipage}[b]{\linewidth}\raggedright
System
\end{minipage} & \begin{minipage}[b]{\linewidth}\raggedright
Useful scale
\end{minipage} & \begin{minipage}[b]{\linewidth}\raggedright
Source-backed context
\end{minipage} \\
\midrule\noalign{}
\endhead
\bottomrule\noalign{}
\endlastfoot
Digestive absorption & Small-intestinal villi and microvilli expand
absorptive surface & Nutrient uptake depends on epithelial surface area,
enzymes, bile, transporters, and gut motility
\citep{niddk2024digestivesystem}. \\
Renal excretion & Adult kidneys filter plasma continuously and adjust
urine concentration & Filtration, tubular reabsorption, secretion, and
collecting-duct water handling set volume, electrolyte, pH, and
nitrogen-waste balance \citep{niddk2024kidneys}. \\
Skeletal muscle & Sarcomere force depends on actin-myosin overlap and
cross-bridge cycling & Sliding-filament physiology links molecular
geometry to organismal movement \citep{huxley1954sliding}. \\
Somatosensation & Receptor density sets spatial resolution &
Mechanoreceptor subtype, receptive-field size, and cortical
magnification jointly shape touch discrimination
\citep{abraira2013somatosensory}. \\
Behaviour & Four explanatory levels & Mechanism, development, function,
and phylogeny answer different questions about the same behaviour
\citep{tinbergen1963aims}. \\
Reproduction & Infertility triage starts after 12 months, or 6 months at
age 35+ & Evaluation separates ovulation, sperm production, anatomy,
implantation, endocrine timing, and age-dependent gamete quality
\citep{cdc2024reproductivehealth}. \\
\end{longtable}
}

\begin{center}\rule{0.5\linewidth}{0.5pt}\end{center}

\subsection*{D.7 Fundamental Biological
Constants}\label{d.7-fundamental-biological-constants}
\addcontentsline{toc}{subsection}{D.7 Fundamental Biological Constants}

{\def\LTcaptype{none} % do not increment counter
\begin{longtable}[]{@{}
  >{\raggedright\arraybackslash}p{(\linewidth - 4\tabcolsep) * \real{0.3000}}
  >{\raggedright\arraybackslash}p{(\linewidth - 4\tabcolsep) * \real{0.1667}}
  >{\raggedright\arraybackslash}p{(\linewidth - 4\tabcolsep) * \real{0.5333}}@{}}
\toprule\noalign{}
\begin{minipage}[b]{\linewidth}\raggedright
Parameter
\end{minipage} & \begin{minipage}[b]{\linewidth}\raggedright
Value
\end{minipage} & \begin{minipage}[b]{\linewidth}\raggedright
Source / context
\end{minipage} \\
\midrule\noalign{}
\endhead
\bottomrule\noalign{}
\endlastfoot
Resting cell membrane potential & \texorpdfstring{\ensuremath{\approx}}{~} \texorpdfstring{\ensuremath{-}}{-}70 mV & Neurons, muscle cells \\
Intracellular {[}Na\texorpdfstring{\textsuperscript{+}}{+}{]} / {[}K\texorpdfstring{\textsuperscript{+}}{+}{]} & about 10 / 140 mM &
Maintained by Na\texorpdfstring{\textsuperscript{+}}{+}/K\texorpdfstring{\textsuperscript{+}}{+}-ATPase \\
Extracellular {[}Na\texorpdfstring{\textsuperscript{+}}{+}{]} / {[}K\texorpdfstring{\textsuperscript{+}}{+}{]} & about 145 / 5 mM &
Plasma \\
Blood pH & 7.35--7.45 & Tightly regulated \\
Arterial pCO\texorpdfstring{\textsubscript{2}}{2} & about 40 mmHg & Regulated by respiration \\
Body-water content & 60 \% of body mass & Total body water \\
Haemoglobin O\texorpdfstring{\textsubscript{2}}{2} capacity & 1.34 mL O\texorpdfstring{\textsubscript{2}}{2} / g Hb & Full saturation \\
Blood volume & about 5 L in adult & 7 \% of body mass \\
Body surface area (adult) & about 1.8 m\texorpdfstring{\textsuperscript{2}}{2} & DuBois formula \\
Mitochondrial membrane potential & \texorpdfstring{\ensuremath{\approx}}{~} \texorpdfstring{\ensuremath{-}}{-}180 mV & Drives ATP synthesis \\
Proton gradient \texorpdfstring{\ensuremath{\Delta}}{Delta}pH across inner mito membrane & about 1 unit &
Chemiosmosis \\
Genome size (human) & 3 \texorpdfstring{\ensuremath{\times}}{x} 10\texorpdfstring{\textsuperscript{9}}{9} bp & Haploid \\
Protein-coding genes (human) & about 20,000 & ENCODE
estimate \\
Cells in the adult human body & about 3.7 \texorpdfstring{\ensuremath{\times}}{x} 10\texorpdfstring{\textsuperscript{1}}{1}\texorpdfstring{\textsuperscript{3}}{3} & Bianconi et
al.~2013 \\
Bacteria in human gut & about 10\texorpdfstring{\textsuperscript{1}}{1}\texorpdfstring{\textsuperscript{4}}{4} & about  equal to
human cell count \\
Species described on Earth & about 2 \texorpdfstring{\ensuremath{\times}}{x} 10\texorpdfstring{\textsuperscript{6}}{6} & Estimated 10--30\texorpdfstring{\ensuremath{\times}}{x}
more undescribed \\
\end{longtable}
}

\begin{center}\rule{0.5\linewidth}{0.5pt}\end{center}

\emph{Module: reference only (no code)}

\newpage

\section*{Appendix E --- A Periodic Table for
Biology}\label{appendix-e-a-periodic-table-for-biology}
\addcontentsline{toc}{section}{Appendix E --- A Periodic Table for
Biology}

\label{sec:appendix_periodic_table}

\begin{quote}
\textbf{Appendix E} \texorpdfstring{\ensuremath{\cdot}}{.} Level 1/3 \texorpdfstring{\ensuremath{\cdot}}{.} 15 min read \texorpdfstring{\ensuremath{\cdot}}{.} Prerequisites:
\cref{sec:unit_I_atoms_molecules}
\end{quote}

Only about 25 of the 118 known elements are biologically essential. This
appendix summarises what each essential element does, what the classic
pathologies of deficiency and excess look like, and where to encounter
each in the textbook.

\begin{center}\rule{0.5\linewidth}{0.5pt}\end{center}

\subsection*{E.1 The Six Bulk Elements
(CHNOPS)}\label{e.1-the-six-bulk-elements-chnops}
\addcontentsline{toc}{subsection}{E.1 The Six Bulk Elements (CHNOPS)}

Together these account for about 99 \% of a human's mass.

{\def\LTcaptype{none} % do not increment counter
\begin{longtable}[]{@{}
  >{\raggedright\arraybackslash}p{(\linewidth - 8\tabcolsep) * \real{0.1556}}
  >{\raggedright\arraybackslash}p{(\linewidth - 8\tabcolsep) * \real{0.1333}}
  >{\raggedright\arraybackslash}p{(\linewidth - 8\tabcolsep) * \real{0.2667}}
  >{\raggedright\arraybackslash}p{(\linewidth - 8\tabcolsep) * \real{0.1111}}
  >{\raggedright\arraybackslash}p{(\linewidth - 8\tabcolsep) * \real{0.3333}}@{}}
\toprule\noalign{}
\begin{minipage}[b]{\linewidth}\raggedright
Element
\end{minipage} & \begin{minipage}[b]{\linewidth}\raggedright
Symbol
\end{minipage} & \begin{minipage}[b]{\linewidth}\raggedright
about \% body mass
\end{minipage} & \begin{minipage}[b]{\linewidth}\raggedright
Where
\end{minipage} & \begin{minipage}[b]{\linewidth}\raggedright
Biological role
\end{minipage} \\
\midrule\noalign{}
\endhead
\bottomrule\noalign{}
\endlastfoot
Oxygen & O & 65 \% & water, organic molecules & Electron acceptor in
respiration; H-bond acceptor; part of every major biomolecule \\
Carbon & C & 18 \% & organic biomolecules, CO\texorpdfstring{\textsubscript{2}}{2}/HCO\texorpdfstring{\textsubscript{3}}{3}\texorpdfstring{\textsuperscript{-}}{-} pools & Tetravalent
backbone of proteins, carbohydrates, lipids, and nucleic acids (see
\cref{sec:unit_I_atoms_molecules}) \\
Hydrogen & H & 10 \% & water and most organic biomolecules & Partner in
polar bonds; proton transfer; pH \\
Nitrogen & N & 3 \% & amino acids, nucleic acids, ATP & Codes the
backbone of DNA/RNA; amide in peptide bond; ring atom in bases \\
Phosphorus & P & 1 \% & DNA/RNA backbone, ATP, phospholipids &
Phosphodiester bonds; free-energy carrier (ATP); membrane structure \\
Sulfur & S & 0.25 \% & cysteine, methionine, CoA, Fe-S & Disulfide
bridges in proteins; thioester bonds in metabolism \\
\end{longtable}
}

\begin{center}\rule{0.5\linewidth}{0.5pt}\end{center}

\subsection*{E.2 Essential Macro-Minerals (Electrolytes +
Osteogenics)}\label{e.2-essential-macro-minerals-electrolytes-osteogenics}
\addcontentsline{toc}{subsection}{E.2 Essential Macro-Minerals
(Electrolytes + Osteogenics)}

{\def\LTcaptype{none} % do not increment counter
\begin{longtable}[]{@{}
  >{\raggedright\arraybackslash}p{(\linewidth - 10\tabcolsep) * \real{0.1591}}
  >{\raggedright\arraybackslash}p{(\linewidth - 10\tabcolsep) * \real{0.1364}}
  >{\raggedright\arraybackslash}p{(\linewidth - 10\tabcolsep) * \real{0.2500}}
  >{\raggedright\arraybackslash}p{(\linewidth - 10\tabcolsep) * \real{0.0909}}
  >{\raggedright\arraybackslash}p{(\linewidth - 10\tabcolsep) * \real{0.2273}}
  >{\raggedright\arraybackslash}p{(\linewidth - 10\tabcolsep) * \real{0.1364}}@{}}
\toprule\noalign{}
\begin{minipage}[b]{\linewidth}\raggedright
Element
\end{minipage} & \begin{minipage}[b]{\linewidth}\raggedright
Symbol
\end{minipage} & \begin{minipage}[b]{\linewidth}\raggedright
Common form
\end{minipage} & \begin{minipage}[b]{\linewidth}\raggedright
Role
\end{minipage} & \begin{minipage}[b]{\linewidth}\raggedright
Deficiency
\end{minipage} & \begin{minipage}[b]{\linewidth}\raggedright
Excess
\end{minipage} \\
\midrule\noalign{}
\endhead
\bottomrule\noalign{}
\endlastfoot
Sodium & Na & Na\texorpdfstring{\textsuperscript{+}}{+} & Principal extracellular cation; action potential
upstroke & Hyponatraemia \texorpdfstring{\ensuremath{\to}}{->} seizures & Hypernatraemia \texorpdfstring{\ensuremath{\to}}{->} dehydration \\
Potassium & K & K\texorpdfstring{\textsuperscript{+}}{+} & Principal intracellular cation; resting potential &
Hypokalaemia \texorpdfstring{\ensuremath{\to}}{->} arrhythmia, muscle weakness & Hyperkalaemia \texorpdfstring{\ensuremath{\to}}{->} cardiac
arrest \\
Calcium & Ca & Ca\texorpdfstring{\textsuperscript{2}}{2}\texorpdfstring{\textsuperscript{+}}{+} & Bone matrix; signalling; coagulation; muscle
contraction & Hypocalcaemia \texorpdfstring{\ensuremath{\to}}{->} tetany & Hypercalcaemia \texorpdfstring{\ensuremath{\to}}{->} stones,
constipation \\
Magnesium & Mg & Mg\texorpdfstring{\textsuperscript{2}}{2}\texorpdfstring{\textsuperscript{+}}{+} & ATP co-factor; enzyme cofactor; ribosome
structure & Hypomagnesaemia \texorpdfstring{\ensuremath{\to}}{->} arrhythmia, tetany & Rare; laxative
effect \\
Chloride & Cl & Cl\texorpdfstring{\textsuperscript{-}}{-} & Major anion; gastric HCl; GABA-A/glycine receptors
& Loss with vomiting & Acidosis \\
Phosphate & P & HPO\texorpdfstring{\textsubscript{4}}{4}\texorpdfstring{\textsuperscript{2}}{2}\texorpdfstring{\textsuperscript{-}}{-} & Energy transfer (ATP); DNA backbone; bone
mineralisation & Muscle weakness & Soft tissue calcification (CKD) \\
Sulfate & S & SO\texorpdfstring{\textsubscript{4}}{4}\texorpdfstring{\textsuperscript{2}}{2}\texorpdfstring{\textsuperscript{-}}{-} & Sulfated proteoglycans; sulfation of xenobiotics &
Rare isolated deficiency & Rare \\
\end{longtable}
}

\begin{center}\rule{0.5\linewidth}{0.5pt}\end{center}

\subsection*{E.3 Trace Elements (Enzyme
Cofactors)}\label{e.3-trace-elements-enzyme-cofactors}
\addcontentsline{toc}{subsection}{E.3 Trace Elements (Enzyme Cofactors)}

Each is essential at mg day\texorpdfstring{\textsuperscript{-}}{-}\texorpdfstring{\textsuperscript{1}}{1} to \texorpdfstring{\ensuremath{\mu}}{mu}g day\texorpdfstring{\textsuperscript{-}}{-}\texorpdfstring{\textsuperscript{1}}{1} scale --- too much or too
little causes specific pathology.

{\def\LTcaptype{none} % do not increment counter
\begin{longtable}[]{@{}
  >{\raggedright\arraybackslash}p{(\linewidth - 6\tabcolsep) * \real{0.1321}}
  >{\raggedright\arraybackslash}p{(\linewidth - 6\tabcolsep) * \real{0.0755}}
  >{\raggedright\arraybackslash}p{(\linewidth - 6\tabcolsep) * \real{0.4340}}
  >{\raggedright\arraybackslash}p{(\linewidth - 6\tabcolsep) * \real{0.3585}}@{}}
\toprule\noalign{}
\begin{minipage}[b]{\linewidth}\raggedright
Element
\end{minipage} & \begin{minipage}[b]{\linewidth}\raggedright
Role
\end{minipage} & \begin{minipage}[b]{\linewidth}\raggedright
Pathology of deficiency
\end{minipage} & \begin{minipage}[b]{\linewidth}\raggedright
Pathology of excess
\end{minipage} \\
\midrule\noalign{}
\endhead
\bottomrule\noalign{}
\endlastfoot
Iron (Fe) & Haemoglobin, cytochromes, Fe-S clusters & Anaemia, fatigue &
Hereditary haemochromatosis (iron overload) \\
Copper (Cu) & Cytochrome c oxidase, tyrosinase, lysyl oxidase & Menkes
disease (curly hair, neurodegeneration) & Wilson disease (Cu deposition
in liver, brain) \\
Zinc (Zn) & \textgreater300 enzymes; zinc finger transcription factors &
Impaired growth, taste, immunity & Interferes with Cu absorption \\
Manganese (Mn) & Superoxide dismutase, glutamine synthetase & Rare;
impaired connective tissue & Parkinson-like syndrome (welders) \\
Cobalt (Co) & Vitamin B\texorpdfstring{\textsubscript{1}}{1}\texorpdfstring{\textsubscript{2}}{2} core & Megaloblastic anaemia, neuropathy &
Rare cardiomyopathy (beer cobalt in 1960s) \\
Molybdenum (Mo) & Xanthine oxidase, sulfite oxidase, nitrate reductase &
Rare; sulfite oxidase deficiency & Gout-like syndrome \\
Selenium (Se) & Glutathione peroxidase (GPX); thyroid deiodinase &
Keshan disease (cardiomyopathy) & Selenosis (hair loss, nail
brittleness) \\
Iodine (I) & Thyroid hormones (T3, T4) & Hypothyroidism, goiter,
cretinism & Thyrotoxicosis \\
Chromium (Cr) & Disputed insulin cofactor & Possibly impaired glucose
tolerance & Cr(VI) is carcinogenic \\
Fluorine (F) & Fluoroapatite in enamel & Dental caries & Fluorosis
(mottled enamel) \\
\end{longtable}
}

\begin{center}\rule{0.5\linewidth}{0.5pt}\end{center}

\subsection*{E.4 Ultra-Trace and Context-Dependent
Elements}\label{e.4-ultra-trace-and-context-dependent-elements}
\addcontentsline{toc}{subsection}{E.4 Ultra-Trace and Context-Dependent
Elements}

\begin{itemize}
\tightlist
\item
  \textbf{Nickel (Ni)} --- cofactor in urease (bacterial, plant) and
  hydrogenase; suspected but not proven essential in mammals.
\item
  \textbf{Silicon (Si)} --- structural in diatoms, horsetails, grasses
  (silica phytoliths); supports collagen cross-linking in vertebrates.
\item
  \textbf{Boron (B)} --- essential in plants (cell-wall cross-linking
  via borate--rhamnogalacturonan); possibly beneficial in vertebrates.
\item
  \textbf{Vanadium (V)} --- sea-squirt blood pigment (vanadocytes);
  pharmacology for insulin-mimetic effects.
\item
  \textbf{Tungsten (W)} --- used instead of molybdenum by
  hyperthermophilic archaea.
\end{itemize}

\begin{center}\rule{0.5\linewidth}{0.5pt}\end{center}

\subsection*{E.5 Life Without an Expected
Element}\label{e.5-life-without-an-expected-element}
\addcontentsline{toc}{subsection}{E.5 Life Without an Expected Element}

Biology has bent the rules. Examples of organisms that substitute or
replace commonly ``essential'' elements:

\begin{itemize}
\tightlist
\item
  \textbf{Iron-free bacteria in deep-sea hydrothermal vents} --- use
  manganese peroxidases instead.
\item
  \textbf{Selenocysteine} (the ``21st amino acid'') --- expands the
  amino-acid alphabet using the rare element Se.
\item
  \textbf{Pyrrolysine} (the ``22nd amino acid'') --- found in
  methanogenic archaea; does not require a rare element, but is an
  unusual 22-atom residue.
\item
  \textbf{Arsenic-based life controversy (GFAJ-1, 2010)} --- claim of As
  replacing P in DNA of a \emph{Halomonas} strain; later retracted.
\end{itemize}

The periodic table remains under-explored for biological innovations.
Astrobiological life may use different element sets entirely.

\begin{center}\rule{0.5\linewidth}{0.5pt}\end{center}

\subsection*{E.6 Quick Lookup: Which Chapter Covers Each
Element?}\label{e.6-quick-lookup-which-chapter-covers-each-element}
\addcontentsline{toc}{subsection}{E.6 Quick Lookup: Which Chapter Covers
Each Element?}

{\def\LTcaptype{none} % do not increment counter
\begin{longtable}[]{@{}
  >{\raggedright\arraybackslash}p{(\linewidth - 4\tabcolsep) * \real{0.1750}}
  >{\raggedright\arraybackslash}p{(\linewidth - 4\tabcolsep) * \real{0.3750}}
  >{\raggedright\arraybackslash}p{(\linewidth - 4\tabcolsep) * \real{0.4500}}@{}}
\toprule\noalign{}
\begin{minipage}[b]{\linewidth}\raggedright
Element
\end{minipage} & \begin{minipage}[b]{\linewidth}\raggedright
Primary chapter
\end{minipage} & \begin{minipage}[b]{\linewidth}\raggedright
Supporting chapter
\end{minipage} \\
\midrule\noalign{}
\endhead
\bottomrule\noalign{}
\endlastfoot
C, H, N, O & \cref{sec:unit_I_atoms_molecules} &
\cref{sec:unit_I_macromolecules} \\
P & \cref{sec:unit_IV_dna_replication_and_cell_cycle} &
\cref{sec:unit_III_bioenergetics_and_respiration} \\
S & \cref{sec:unit_I_macromolecules} &
\cref{sec:unit_III_bioenergetics_and_respiration} \\
Na, K & \cref{sec:unit_II_membrane_transport} &
\cref{sec:unit_IX_nervous_system};
\cref{sec:unit_IX_action_potential_synapses} \\
Ca & \cref{sec:unit_II_cell_signaling} &
\cref{sec:unit_IX_circulation_respiration_homeostasis};
\cref{sec:unit_IX_endocrine_and_immune} \\
Mg & \cref{sec:unit_I_macromolecules} &
\cref{sec:unit_III_bioenergetics_and_respiration} \\
Fe & \cref{sec:unit_III_bioenergetics_and_respiration} &
\cref{sec:unit_IX_circulation_respiration_homeostasis} \\
Cu & \cref{sec:unit_III_bioenergetics_and_respiration} & \\
Zn & \cref{sec:unit_I_enzymes_and_kinetics} &
\cref{sec:unit_IV_epigenetics_and_gene_regulation} \\
Se & \cref{sec:unit_III_bioenergetics_and_respiration} & \\
I & \cref{sec:unit_IX_endocrine_and_immune} & \\
Co & --- & \cref{sec:unit_III_metabolic_integration} \\
\end{longtable}
}

\begin{center}\rule{0.5\linewidth}{0.5pt}\end{center}

\emph{Module: reference only (no code)}

\newpage

\section*{Appendix F --- Master Glossary of Biological
Terms}\label{appendix-f-master-glossary-of-biological-terms}
\addcontentsline{toc}{section}{Appendix F --- Master Glossary of
Biological Terms}

\label{sec:glossary}

\emph{Etymology notation: {[}Gk{]} = Greek; \hyperref[l]{L} = Latin;
{[}OFr{]} = Old French; {[}Ar{]} = Arabic.} \emph{Cross-references use
semantic \texttt{\textbackslash{}cref\{...\}} links generated from the
canonical table of contents.}

\emph{Maintenance note: glossary terms must use a Pandoc bracketed span,
\texttt{{[}**Term**{]}\{\#gl:\textless{}slug\textgreater{}\}} (with
\breaktt{markdown+bracketed\_spans}, Pandoc emits a real LaTeX
\texttt{\textbackslash{}label\{gl:...\}} in PDF). Chapters link with
\texttt{{[}**term**{]}(\#gl:\textless{}slug\textgreater{})}. After
editing terms, run
\breaktt{uv\ run\ python\ scripts/link\_glossary.py\ -\/-check} from the
project root and verify \breaktt{tests/test\_build\_invariants.py} plus
markdown validation.}

\begin{center}\rule{0.5\linewidth}{0.5pt}\end{center}

\subsection*{A}\label{a}
\addcontentsline{toc}{subsection}{A}

\protect\phantomsection\label{gl:abiotic}{\textbf{Abiotic}} {[}Gk
\emph{a} = without + \emph{bios} = life + \emph{-ic} = pertaining to{]}
--- Non-living physical and chemical components of an environment
(light, temperature, water, minerals). Contrast with \textbf{biotic}. \texorpdfstring{\ensuremath{\to}}{->}
\cref{sec:unit_X_biomes_and_conservation}

\protect\phantomsection\label{gl:actin}{\textbf{Actin}} {[}Gk
\emph{aktis} = ray or beam{]} --- A 42 kDa globular protein (G-actin)
that polymerises into filaments (F-actin); a major component of the
cytoskeleton, sarcomeres, and lamellipodia. \texorpdfstring{\ensuremath{\to}}{->}
\cref{sec:unit_II_cell_structure}

\protect\phantomsection\label{gl:action-potential}{\textbf{Action
potential}} {[}L \emph{actio} = a doing + \emph{potentia} = power{]} ---
A transient, all-or-none reversal of membrane potential
(about \texorpdfstring{\ensuremath{-}}{-}70 mV to +40 mV) propagated along neuronal axons; gated
by voltage-dependent Na\texorpdfstring{\textsuperscript{+}}{+} and K\texorpdfstring{\textsuperscript{+}}{+} channels. Described mathematically by
the Hodgkin-Huxley equations (1952). \texorpdfstring{\ensuremath{\to}}{->} \cref{sec:unit_IX_nervous_system}

\protect\phantomsection\label{gl:adaptive-immunity}{\textbf{Adaptive
immunity}} {[}L \emph{adaptare} = to fit + \emph{immunis} = exempt{]}
--- Antigen-specific defence mediated by B and T lymphocytes. It is
slower than innate immunity during a first exposure but generates
immunological memory, enabling faster and stronger secondary responses.
\texorpdfstring{\ensuremath{\to}}{->} \cref{sec:unit_IX_endocrine_and_immune}

\protect\phantomsection\label{gl:adenosine-triphosphate}{\textbf{Adenosine
triphosphate (ATP)}} {[}Gk \emph{aden} = gland (as in adenine,
identified in gland tissue) + \emph{tri} = three + \emph{phosphorus} =
light-bearer{]} --- A central phosphoryl-group donor and short-term
energy carrier in cells. Hydrolysis to ADP + P\texorpdfstring{\textsubscript{i}}{i} releases
about 30.5 kJ/mol under biochemical standard conditions, while
the effective cellular free energy is often larger in magnitude because
ATP, ADP, and P\texorpdfstring{\textsubscript{i}}{i} concentrations are far from equilibrium. \texorpdfstring{\ensuremath{\to}}{->}
\cref{sec:unit_III_bioenergetics_and_respiration}

\protect\phantomsection\label{gl:aerobic}{\textbf{Aerobic}} {[}Gk
\emph{aer} = air + \emph{bios} = life{]} --- Requiring molecular oxygen
(O\texorpdfstring{\textsubscript{2}}{2}) as the terminal electron acceptor in cellular respiration.
Opposite: \textbf{anaerobic}. \texorpdfstring{\ensuremath{\to}}{->}
\cref{sec:unit_III_bioenergetics_and_respiration}

\protect\phantomsection\label{gl:allele}{\textbf{Allele}} {[}Gk
\emph{allelon} = of each other; shortened from \emph{allelemorph}{]} ---
One of two or more alternative forms of a gene at the same chromosomal
locus. Defined by a distinct nucleotide sequence. \texorpdfstring{\ensuremath{\to}}{->}
\cref{sec:unit_V_mendelian_genetics}

\protect\phantomsection\label{gl:allostasis}{\textbf{Allostasis}} {[}Gk
\emph{allos} = other + \emph{stasis} = standing{]} --- Predictive
physiological regulation in which the brain adjusts setpoints and
resource allocation in anticipation of demand, rather than merely
correcting error after disturbance. \texorpdfstring{\ensuremath{\to}}{->}
\cref{sec:unit_IX_endocrine_and_immune}

\protect\phantomsection\label{gl:allosteric}{\textbf{Allosteric}} {[}Gk
\emph{allos} = other + \emph{stereos} = solid/space{]} --- Describing a
regulatory site on an enzyme or receptor that is spatially distinct from
the active site; ligand binding there changes the macromolecule's
conformation and activity. \texorpdfstring{\ensuremath{\to}}{->} \cref{sec:unit_I_enzymes_and_kinetics}

\protect\phantomsection\label{gl:anaerobic}{\textbf{Anaerobic}} {[}Gk
\emph{an} = without + \emph{aer} = air + \emph{bios} = life{]} ---
Occurring without oxygen. Anaerobic respiration uses alternative
electron acceptors (nitrate, sulphate); fermentation regenerates NAD\texorpdfstring{\textsuperscript{+}}{+}
without any external acceptor. \texorpdfstring{\ensuremath{\to}}{->}
\cref{sec:unit_III_bioenergetics_and_respiration}

\protect\phantomsection\label{gl:anaphase}{\textbf{Anaphase}} {[}Gk
\emph{ana} = up/apart + \emph{phasis} = appearance{]} --- The stage of
mitosis/meiosis in which sister chromatids (mitosis, meiosis II) or
homologous chromosomes (meiosis I) are pulled to opposite poles by
spindle fibres. \texorpdfstring{\ensuremath{\to}}{->} \cref{sec:unit_IV_dna_replication_and_cell_cycle}

\protect\phantomsection\label{gl:angiosperm}{\textbf{Angiosperm}} {[}Gk
\emph{angeion} = vessel/container + \emph{sperma} = seed{]} --- A
flowering plant whose seeds are enclosed within a fruit derived from the
ovary wall; the dominant plant group (about 300,000 species).
Contrast: \textbf{gymnosperm}. \texorpdfstring{\ensuremath{\to}}{->}
\cref{sec:unit_VIII_plant_structure_and_water}

\protect\phantomsection\label{gl:anticodon}{\textbf{Anticodon}} {[}Gk
\emph{anti} = against or opposite + L \emph{codex} = writing block{]}
--- The three-nucleotide sequence on a tRNA molecule that base-pairs
with the complementary codon on mRNA during translation. \texorpdfstring{\ensuremath{\to}}{->}
\cref{sec:unit_IV_gene_expression}

\protect\phantomsection\label{gl:apoptosis}{\textbf{Apoptosis}} {[}Gk
\emph{apo} = away + \emph{ptōsis} = a falling{]} --- Programmed,
caspase-mediated cell death characterised by chromatin condensation, DNA
fragmentation, and membrane blebbing, without inflammatory response.
Contrast: \textbf{necrosis}. \texorpdfstring{\ensuremath{\to}}{->} \cref{sec:unit_II_cell_signaling}

\protect\phantomsection\label{gl:archaea}{\textbf{Archaea}} {[}Gk
\emph{archaios} = ancient{]} --- One of three domains of life;
prokaryotes with ether-linked isoprenoid membrane lipids, histone-like
proteins, and replication/transcription machinery closer to eukaryotes
than to bacteria. Discovered by Woese \& Fox (1977). \texorpdfstring{\ensuremath{\to}}{->}
\cref{sec:unit_II_cell_theory,sec:unit_VI_phylogenetics}

\protect\phantomsection\label{gl:atp-synthase}{\textbf{ATP synthase}}
{[}see ATP + Gk \emph{syn} = together + \emph{thesis} = a placing{]} ---
Also Complex V; a rotary F\texorpdfstring{\textsubscript{0}}{0}F\texorpdfstring{\textsubscript{1}}{1} molecular motor driven by the
proton-motive force to phosphorylate ADP to ATP. The c-ring rotation of
120° per ATP was directly visualised by Noji et al.~(1997). \texorpdfstring{\ensuremath{\to}}{->}
\cref{sec:unit_III_bioenergetics_and_respiration}

\protect\phantomsection\label{gl:autotroph}{\textbf{Autotroph}} {[}Gk
\emph{autos} = self + \emph{trophe} = nourishment{]} --- An organism
that produces its own organic molecules from inorganic carbon (CO\texorpdfstring{\textsubscript{2}}{2});
either \textbf{photoautotroph} (energy from light) or
\textbf{chemoautotroph} (energy from inorganic chemicals). \texorpdfstring{\ensuremath{\to}}{->}
\cref{sec:unit_II_cell_theory,sec:unit_X_community_ecology}

\protect\phantomsection\label{gl:autonomic-nervous-system}{\textbf{Autonomic
nervous system (ANS)}} {[}Gk \emph{autonomos} = self-governing{]} ---
The visceral motor division of the peripheral nervous system; divided
into \textbf{sympathetic} (thoracolumbar; noradrenaline/adrenaline;
``fight-or-flight'') and \textbf{parasympathetic} (craniosacral;
acetylcholine; ``rest-and-digest'') branches. ANS innervates smooth
muscle, cardiac muscle, and glands. \texorpdfstring{\ensuremath{\to}}{->} \cref{sec:unit_IX_nervous_system}

\protect\phantomsection\label{gl:axon-hillock}{\textbf{Axon hillock}}
{[}Gk \emph{axon} = axis + OE \emph{hyll} = small hill{]} --- The region
of the neuron where the axon originates from the cell body (soma); the
site of action potential initiation due to the high density of
voltage-gated Na\texorpdfstring{\textsuperscript{+}}{+} channels. EPSPs must summate here to reach threshold
(about \texorpdfstring{\ensuremath{-}}{-}55 mV). \texorpdfstring{\ensuremath{\to}}{->} \cref{sec:unit_IX_nervous_system}

\begin{center}\rule{0.5\linewidth}{0.5pt}\end{center}

\subsection*{B}\label{b}
\addcontentsline{toc}{subsection}{B}

\protect\phantomsection\label{gl:base-pair}{\textbf{Base pair (bp)}}
{[}Gk \emph{basis} = step, foundation{]} --- A complementary nucleotide
pair held by hydrogen bonds in double-stranded nucleic acids: A--T (2
H-bonds in DNA) and G--C (3 H-bonds). One bp \texorpdfstring{\ensuremath{\approx}}{~} 0.34 nm along the DNA
helix. \texorpdfstring{\ensuremath{\to}}{->} \cref{sec:unit_IV_dna_replication_and_cell_cycle}

\protect\phantomsection\label{gl:biofilm}{\textbf{Biofilm}} {[}Gk
\emph{bios} = life + L \emph{filum} = thread{]} --- A structured
community of microorganisms attached to a surface and enclosed in a
self-produced extracellular polysaccharide matrix; regulated by quorum
sensing; often up to 1,000-fold more antibiotic tolerant than planktonic
cells because matrix diffusion limits, slow growth, stress responses,
and persister cells overlap. \texorpdfstring{\ensuremath{\to}}{->}
\cref{sec:unit_VII_bacteria_archaea_viruses}

\protect\phantomsection\label{gl:biome}{\textbf{Biome}} {[}Gk
\emph{bios} = life + \emph{-ome} = mass/body{]} --- A large-scale
terrestrial or aquatic ecosystem characterised by a distinctive climate,
vegetation, and fauna. Classified by Whittaker (1975) primarily by mean
annual precipitation (MAP) and mean annual temperature (MAT). \texorpdfstring{\ensuremath{\to}}{->}
\cref{sec:unit_X_biomes_and_conservation}

\protect\phantomsection\label{gl:buffer}{\textbf{Buffer}} {[}OFr
\emph{bufer} = to protect from a blow{]} --- A solution containing a
weak acid and its conjugate base that resists changes in pH by absorbing
added H\texorpdfstring{\textsuperscript{+}}{+} or OH\texorpdfstring{\textsuperscript{-}}{-}. Effective within \texorpdfstring{\ensuremath{\pm}}{+/-}1 pH unit of pK\texorpdfstring{\textsubscript{a}}{a}
(Henderson-Hasselbalch). \texorpdfstring{\ensuremath{\to}}{->} \cref{sec:unit_I_atoms_molecules}

\begin{center}\rule{0.5\linewidth}{0.5pt}\end{center}

\subsection*{C}\label{c}
\addcontentsline{toc}{subsection}{C}

\protect\phantomsection\label{gl:calvin-cycle}{\textbf{Calvin cycle}}
{[}named for Melvin Calvin, who elucidated it using \texorpdfstring{\textsuperscript{1}}{1}\texorpdfstring{\textsuperscript{4}}{4}C radioisotopes,
1950{]} --- The light-independent carbon-fixation reactions of
photosynthesis in the chloroplast stroma; not literally dark, because
they depend on ATP and NADPH produced by the light reactions. Also
called the Calvin-Benson-Bassham cycle. Net: 3 CO\texorpdfstring{\textsubscript{2}}{2} + 9 ATP + 6 NADPH \texorpdfstring{\ensuremath{\to}}{->}
G3P. \texorpdfstring{\ensuremath{\to}}{->} \cref{sec:unit_III_photosynthesis}

\protect\phantomsection\label{gl:caspase}{\textbf{Caspase}} {[}L
\emph{caspase} = cysteine-containing aspartate protease{]} --- A family
of cysteine proteases that execute apoptosis; exists as inactive
zymogens; activated by cleavage at specific aspartate residues in a
proteolytic cascade. \texorpdfstring{\ensuremath{\to}}{->} \cref{sec:unit_II_cell_signaling}

\protect\phantomsection\label{gl:cell-cycle}{\textbf{Cell cycle}} {[}L
\emph{cella} = small room + Gk \emph{kyklos} = circle{]} --- The ordered
sequence of events by which a cell duplicates its contents and divides:
G\texorpdfstring{\textsubscript{1}}{1} \texorpdfstring{\ensuremath{\to}}{->} S (DNA synthesis) \texorpdfstring{\ensuremath{\to}}{->} G\texorpdfstring{\textsubscript{2}}{2} \texorpdfstring{\ensuremath{\to}}{->} M (mitosis) \texorpdfstring{\ensuremath{\to}}{->} cytokinesis. Regulated by
cyclin-CDK complexes and checkpoint kinases. \texorpdfstring{\ensuremath{\to}}{->}
\cref{sec:unit_IV_dna_replication_and_cell_cycle}

\protect\phantomsection\label{gl:camkii}{\textbf{CaMKII
(Calcium/calmodulin-dependent kinase II)}} {[}L \emph{calx} = lime (Ca)
+ \emph{calmodulin} + Gk \emph{kinein} = to move{]} --- A
serine/threonine kinase activated by Ca\texorpdfstring{\textsuperscript{2}}{2}\texorpdfstring{\textsuperscript{+}}{+}/calmodulin;
autophosphorylation at T286 renders it constitutively active
(``molecular memory''). Key effector of NMDA receptor-dependent LTP:
phosphorylates AMPA receptors (increased conductance) and drives AMPA
receptor exocytosis. \texorpdfstring{\ensuremath{\to}}{->} \cref{sec:unit_IX_action_potential_synapses}

\protect\phantomsection\label{gl:cardiac-output}{\textbf{Cardiac output
(CO)}} {[}Gk \emph{kardia} = heart + L \emph{emittere} = to send out{]}
--- The volume of blood ejected by the heart per minute: CO = heart rate
(HR) \texorpdfstring{\ensuremath{\times}}{x} stroke volume (SV). Normal resting CO \texorpdfstring{\ensuremath{\approx}}{~}5.5 L/min; can reach
20--25 L/min during maximal exercise in athletes. Regulated by the
Frank-Starling mechanism and neurohumoral factors. \texorpdfstring{\ensuremath{\to}}{->}
\cref{sec:unit_IX_circulation_respiration_homeostasis}

\protect\phantomsection\label{gl:cdna}{\textbf{cDNA}} {[}complementary
DNA; L \emph{complementum} = that which fills up{]} --- DNA synthesised
from an mRNA template by reverse transcriptase; lacks introns; used in
cloning, expression libraries, and RT-PCR diagnostics. \texorpdfstring{\ensuremath{\to}}{->}
\cref{sec:unit_IV_gene_expression}

\protect\phantomsection\label{gl:centromere}{\textbf{Centromere}} {[}Gk
\emph{kentron} = centre + \emph{meros} = part{]} --- The specialised
chromosomal region where kinetochore proteins assemble and spindle
microtubules attach during mitosis and meiosis. Required for accurate
chromosome segregation. \texorpdfstring{\ensuremath{\to}}{->}
\cref{sec:unit_IV_dna_replication_and_cell_cycle}

\protect\phantomsection\label{gl:chemiosmosis}{\textbf{Chemiosmosis}}
{[}Gk \emph{chemi} = chemistry + \emph{osmos} = push{]} --- The coupling
of electron transport to ATP synthesis via a proton electrochemical
gradient across a membrane (proton-motive force). Proposed by Peter
Mitchell (1961); Nobel Prize 1978. \texorpdfstring{\ensuremath{\to}}{->}
\cref{sec:unit_III_bioenergetics_and_respiration}

\protect\phantomsection\label{gl:chirality}{\textbf{Chirality}} {[}Gk
\emph{cheir} = hand{]} --- The property of a molecule that is
non-superimposable on its mirror image, like a hand. Biology is strongly
homochiral in its core polymers: proteins mostly use L-amino acids,
while nucleic acids and many central metabolites use D-sugars. \texorpdfstring{\ensuremath{\to}}{->}
\cref{sec:unit_I_atoms_molecules}

\protect\phantomsection\label{gl:chlorophyll}{\textbf{Chlorophyll}}
{[}Gk \emph{chloros} = green/yellow-green + \emph{phyllon} = leaf{]} ---
A porphyrin-ring pigment containing Mg\texorpdfstring{\textsuperscript{2}}{2}\texorpdfstring{\textsuperscript{+}}{+} at its centre; the primary
light-harvesting molecule in photosynthesis. Absorbs strongly at 430 nm
(violet/blue) and 680--700 nm (red). \texorpdfstring{\ensuremath{\to}}{->}
\cref{sec:unit_III_photosynthesis}

\protect\phantomsection\label{gl:chromatin}{\textbf{Chromatin}} {[}Gk
\emph{chroma} = colour{]} --- The complex of DNA and proteins (primarily
histones) that compacts the genome within the nucleus. Exists in two
states: \textbf{euchromatin} (open, transcriptionally active) and
\textbf{heterochromatin} (compact, silenced). The fundamental unit is
the nucleosome. \texorpdfstring{\ensuremath{\to}}{->} \cref{sec:unit_IV_epigenetics_and_gene_regulation}

\protect\phantomsection\label{gl:chromosome}{\textbf{Chromosome}} {[}Gk
\emph{chroma} = colour + \emph{soma} = body{]} --- A linear (eukaryotic)
or circular (prokaryotic) DNA molecule associated with proteins
(histones in eukaryotes) that carries the cell's hereditary information.
Humans have 46 chromosomes (23 pairs). \texorpdfstring{\ensuremath{\to}}{->}
\cref{sec:unit_IV_dna_replication_and_cell_cycle}

\protect\phantomsection\label{gl:cpg-island}{\textbf{CpG island}} (C =
cytosine; p = phosphodiester bond; G = guanine) --- A genomic region
(about 500 bp--2 kb) enriched for CG dinucleotides;
observed/expected CpG ratio \textgreater{} 0.6; present at
about 70\% of vertebrate gene promoters. Normally unmethylated;
hypermethylation silences the associated promoter in cancer and
differentiation. \texorpdfstring{\ensuremath{\to}}{->} \cref{sec:unit_IV_epigenetics_and_gene_regulation}

\protect\phantomsection\label{gl:crispr-cas}{\textbf{CRISPR-Cas}} {[}L
\emph{crispus} = clustered + \emph{repetitive} + \emph{palindromic} +
\emph{Cas} = CRISPR-associated protein{]} --- A prokaryotic adaptive
immune system (repeat-spacer arrays + Cas endonucleases) repurposed for
programmable genome editing; Cas9 guided by a synthetic sgRNA cleaves
DNA at a complementary target. \texorpdfstring{\ensuremath{\to}}{->}
\cref{sec:unit_IV_mutations_and_genomics}

\protect\phantomsection\label{gl:codon}{\textbf{Codon}} {[}L
\emph{codex} = block of writing{]} --- A three-nucleotide sequence in
mRNA that specifies one amino acid or a stop signal. The genetic code
has 64 codons encoding 20 amino acids (and 3 stop codons); most amino
acids have multiple synonymous codons (degeneracy). \texorpdfstring{\ensuremath{\to}}{->}
\cref{sec:unit_IV_gene_expression}

\protect\phantomsection\label{gl:cortisol}{\textbf{Cortisol}} {[}L
\emph{cortex} = bark, rind{]} --- The primary glucocorticoid secreted by
the adrenal cortex (zona fasciculata) in response to ACTH; mobilises
glucose (gluconeogenesis, protein catabolism), suppresses immunity
(reduces IL-2, PGE2, and T cell activation), and regulates the stress
response. Circadian peak at 06:00--08:00. Synthetic analogues
(prednisone, dexamethasone) are widely used anti-inflammatory agents. \texorpdfstring{\ensuremath{\to}}{->}
\cref{sec:unit_IX_endocrine_and_immune}

\protect\phantomsection\label{gl:cohesion-tension-theory}{\textbf{Cohesion-tension
theory}} {[}L \emph{cohaerere} = to stick together + Gk \emph{teinein} =
to stretch{]} --- The mechanism by which water ascends xylem vessels in
plants: evaporation at the leaf surface creates tension (negative
pressure) transmitted through the hydrogen-bonded water column via
adhesion to cell walls. Proposed by Dixon \& Joly (1894). \texorpdfstring{\ensuremath{\to}}{->}
\cref{sec:unit_VIII_plant_structure_and_water}

\protect\phantomsection\label{gl:competitive-inhibition}{\textbf{Competitive
inhibition}} {[}L \emph{competere} = to seek together{]} --- Enzyme
inhibition in which the inhibitor binds the active site and competes
with substrate. Raises apparent K\texorpdfstring{\textsubscript{m}}{m} without affecting V\texorpdfstring{\textsubscript{m}}{m}\texorpdfstring{\textsubscript{a}}{a}\texorpdfstring{\textsubscript{x}}{x}; reversed by
increasing \hyperref[s]{S}. \texorpdfstring{\ensuremath{\to}}{->} \cref{sec:unit_I_enzymes_and_kinetics}

\protect\phantomsection\label{gl:cytoplasm}{\textbf{Cytoplasm}} {[}Gk
\emph{kytos} = cell/vessel + \emph{plasma} = something formed{]} --- The
collective cell contents between the plasma membrane and the nucleus (or
nucleoid); includes cytosol, organelles, ribosomes, and the
cytoskeleton. \texorpdfstring{\ensuremath{\to}}{->} \cref{sec:unit_II_cell_structure}

\protect\phantomsection\label{gl:cytoskeleton}{\textbf{Cytoskeleton}}
{[}Gk \emph{kytos} = vessel + \emph{skeletos} = dried up{]} --- The
dynamic network of protein filaments --- actin microfilaments (7 nm),
intermediate filaments (10 nm), and microtubules (25 nm) --- that
maintains cell shape, enables motility, and guides intracellular
transport. \texorpdfstring{\ensuremath{\to}}{->} \cref{sec:unit_II_cell_structure}

\begin{center}\rule{0.5\linewidth}{0.5pt}\end{center}

\subsection*{D}\label{d}
\addcontentsline{toc}{subsection}{D}

\protect\phantomsection\label{gl:diploid}{\textbf{Diploid (2n)}} {[}Gk
\emph{diplous} = double + \emph{eidos} = form{]} --- A cell or organism
with two complete sets of homologous chromosomes; one set from each
parent. Human somatic cells are diploid (2n = 46). Contrast:
\textbf{haploid} (n). \texorpdfstring{\ensuremath{\to}}{->} \cref{sec:unit_V_chromosomal_inheritance}

\protect\phantomsection\label{gl:dna-helicase}{\textbf{DNA helicase}}
{[}Gk \emph{helix} = spiral{]} --- An ATP-dependent enzyme that unwinds
the DNA double helix at the replication fork by breaking hydrogen bonds
between base pairs; the replicative helicase DnaB in bacteria, MCM
complex in eukaryotes. \texorpdfstring{\ensuremath{\to}}{->}
\cref{sec:unit_IV_dna_replication_and_cell_cycle}

\protect\phantomsection\label{gl:dna-ligase}{\textbf{DNA ligase}} {[}L
\emph{ligare} = to bind{]} --- An enzyme that forms phosphodiester bonds
between adjacent nucleotides; joins Okazaki fragments on the lagging
strand; also used in DNA repair and recombinant DNA cloning. \texorpdfstring{\ensuremath{\to}}{->}
\cref{sec:unit_IV_dna_replication_and_cell_cycle}

\protect\phantomsection\label{gl:dominant}{\textbf{Dominant}} {[}L
\emph{dominare} = to rule{]} --- Describing an allele whose phenotypic
effect is expressed in both homozygous and heterozygous states. The
recessive allele is masked in the heterozygote. Mendel (1866). \texorpdfstring{\ensuremath{\to}}{->}
\cref{sec:unit_V_mendelian_genetics}

\protect\phantomsection\label{gl:doubling-time}{\textbf{Doubling time
(t\_d)}} {[}OE \emph{dubbel} = twice + L \emph{tempus} = time{]} --- The
time required for a population or culture to double in number;
\(t_d = \ln 2 / \mu\) where \texorpdfstring{\ensuremath{\mu}}{mu} is the specific growth rate. \texorpdfstring{\ensuremath{\to}}{->}
\cref{sec:unit_VII_bacteria_archaea_viruses}

\begin{center}\rule{0.5\linewidth}{0.5pt}\end{center}

\subsection*{E}\label{e}
\addcontentsline{toc}{subsection}{E}

\protect\phantomsection\label{gl:electronegativity}{\textbf{Electronegativity}}
{[}Gk \emph{elektron} = amber (as source of static charge) + L
\emph{negativus} = denying{]} --- The tendency of an atom to attract
bonding electrons. Pauling scale: F = 3.98; O = 3.44; N = 3.04; C =
2.55; H = 2.20. Determines bond polarity. \texorpdfstring{\ensuremath{\to}}{->}
\cref{sec:unit_I_atoms_molecules}

\protect\phantomsection\label{gl:endocytosis}{\textbf{Endocytosis}}
{[}Gk \emph{endon} = within + \emph{kytos} = vessel{]} --- Uptake of
extracellular material by inward budding of the plasma membrane to form
a vesicle. Subtypes: \textbf{pinocytosis} (fluid/solutes),
\textbf{phagocytosis} (large particles), \textbf{receptor-mediated
endocytosis} (specific ligands via clathrin-coated pits). Contrast:
\textbf{exocytosis}. \texorpdfstring{\ensuremath{\to}}{->} \cref{sec:unit_II_membrane_transport}

\protect\phantomsection\label{gl:endosymbiosis}{\textbf{Endosymbiosis}}
{[}Gk \emph{endon} = within + \emph{sym} = together + \emph{bios} =
life{]} --- A relationship in which one organism lives within another
over evolutionary time, eventually becoming an organelle. The
mitochondrion (about 1.5 Ga ago) and chloroplast
(about 1.2 Ga ago) both originated by endosymbiosis of
bacterial cells (Margulis, 1967). \texorpdfstring{\ensuremath{\to}}{->} \cref{sec:unit_II_cell_theory}

\protect\phantomsection\label{gl:entropy}{\textbf{Entropy (S)}} {[}Gk
\emph{en} = in + \emph{trope} = transformation{]} --- A thermodynamic
state function often described as energy dispersal or the number of
accessible microstates. For spontaneous processes in an isolated system,
total entropy increases (\(\Delta S_{universe} > 0\)). Living organisms
maintain ordered internal structures by exchanging matter and energy
with surroundings. \texorpdfstring{\ensuremath{\to}}{->} \cref{sec:unit_III_bioenergetics_and_respiration}

\protect\phantomsection\label{gl:enzyme}{\textbf{Enzyme}} {[}Gk
\emph{en} = in + \emph{zyme} = leaven/ferment{]} --- A biological
catalyst --- almost always a protein, sometimes RNA (ribozyme) --- that
accelerates a reaction by lowering the activation energy without being
consumed. Characterised by specificity, saturation kinetics
(Michaelis-Menten), and allosteric regulation. \texorpdfstring{\ensuremath{\to}}{->}
\cref{sec:unit_I_enzymes_and_kinetics}

\protect\phantomsection\label{gl:epigenetics}{\textbf{Epigenetics}}
{[}Gk \emph{epi} = above/upon + \emph{genet-} = birth; coined by
Waddington, 1942{]} --- Heritable changes in gene expression that do not
alter the underlying DNA sequence. Molecular carriers: CpG methylation,
histone post-translational modifications (acetylation, methylation,
phosphorylation), and non-coding RNAs (miRNA, lncRNA, piRNA). \texorpdfstring{\ensuremath{\to}}{->}
\cref{sec:unit_IV_epigenetics_and_gene_regulation}

\protect\phantomsection\label{gl:epigenome}{\textbf{Epigenome}} {[}Gk
\emph{epi} = upon + \emph{genos} = origin + \emph{-ome} = total{]} ---
The complete set of epigenetic modifications across a cell's genome;
cell-type specific and environmentally modifiable, unlike the genome
sequence. Mapped by ENCODE and Roadmap Epigenomics projects. \texorpdfstring{\ensuremath{\to}}{->}
\cref{sec:unit_IV_mutations_and_genomics}

\protect\phantomsection\label{gl:euchromatin}{\textbf{Euchromatin}}
{[}Gk \emph{eu} = true/good + \emph{chroma} = colour{]} --- Loosely
packed chromatin associated with H3K4me3 and H3K27ac histone marks;
transcriptionally permissive. Contrast: \textbf{heterochromatin}. \texorpdfstring{\ensuremath{\to}}{->}
\cref{sec:unit_IV_epigenetics_and_gene_regulation}

\protect\phantomsection\label{gl:excitotoxicity}{\textbf{Excitotoxicity}}
{[}L \emph{excitare} = to rouse + Gk \emph{toxikon} = poison{]} ---
Neuronal death caused by excessive glutamate receptor activation
(especially NMDA), leading to massive Ca\texorpdfstring{\textsuperscript{2}}{2}\texorpdfstring{\textsuperscript{+}}{+} influx, mitochondrial
dysfunction, caspase activation, and necrosis/apoptosis. The primary
mechanism of neuronal death in ischaemic stroke. Target of
neuroprotective strategies. \texorpdfstring{\ensuremath{\to}}{->}
\cref{sec:unit_IX_action_potential_synapses}

\protect\phantomsection\label{gl:eukaryote}{\textbf{Eukaryote}} {[}Gk
\emph{eu} = true + \emph{karyon} = nucleus{]} --- An organism whose
cells contain a membrane-enclosed nucleus and membrane-bound organelles.
Includes protists, fungi, plants, and animals. Contrast:
\textbf{prokaryote}. \texorpdfstring{\ensuremath{\to}}{->} \cref{sec:unit_II_cell_theory}

\protect\phantomsection\label{gl:exon}{\textbf{Exon}} {[}Gk \emph{exo} =
outside; coined by Gilbert, 1978{]} --- A coding sequence in a
eukaryotic gene that is retained in the mature mRNA after splicing and
is translated into protein. Contrast: \textbf{intron}. \texorpdfstring{\ensuremath{\to}}{->}
\cref{sec:unit_IV_gene_expression}

\begin{center}\rule{0.5\linewidth}{0.5pt}\end{center}

\subsection*{F}\label{f}
\addcontentsline{toc}{subsection}{F}

\protect\phantomsection\label{gl:fev-fvc}{\textbf{FEV\texorpdfstring{\textsubscript{1}}{1}/FVC}} {[}forced
expiratory volume/forced vital capacity{]} --- The fraction of vital
capacity exhaled in the first second of a forced expiration; obstruction
is best interpreted against the age-dependent lower limit of normal,
with a fixed 0.70 threshold used as a simple COPD screening convention.
Restrictive patterns usually preserve or raise the ratio while reducing
absolute lung volumes. \texorpdfstring{\ensuremath{\to}}{->}
\cref{sec:unit_IX_circulation_respiration_homeostasis}

\protect\phantomsection\label{gl:fatty-acid}{\textbf{Fatty acid}} {[}ME
\emph{fat} + Gk \emph{oxys} = acid/sharp{]} --- A carboxylic acid with a
hydrocarbon chain. Saturated fatty acids pack tightly; unsaturated
\emph{cis} double bonds introduce bends that increase membrane fluidity.
\texorpdfstring{\ensuremath{\to}}{->} \cref{sec:unit_I_macromolecules}

\protect\phantomsection\label{gl:fisheries}{\textbf{Fisheries}} {[}OE
\emph{fisc} = fish{]} --- Human-managed harvest systems for aquatic
populations; modern stock assessment uses age structure, recruitment,
mortality, and maximum sustainable yield models to avoid collapse. \texorpdfstring{\ensuremath{\to}}{->}
\cref{sec:unit_X_population_ecology,sec:unit_X_biomes_and_conservation}

\protect\phantomsection\label{gl:frank-starling-law}{\textbf{Frank-Starling
law}} {[}named for Otto Frank (1895) and Ernest Starling (1914){]} ---
The intrinsic property of cardiac muscle whereby stroke volume increases
with end-diastolic volume (preload): greater stretch of sarcomeres
increases Ca\texorpdfstring{\textsuperscript{2}}{2}\texorpdfstring{\textsuperscript{+}}{+} sensitivity of troponin C, increasing force generation.
Operates up to an optimal sarcomere length (about 2.2 \texorpdfstring{\ensuremath{\mu}}{mu}m);
beyond this, overlap of thick and thin filaments decreases. \texorpdfstring{\ensuremath{\to}}{->}
\cref{sec:unit_IX_circulation_respiration_homeostasis}

\protect\phantomsection\label{gl:facilitated-diffusion}{\textbf{Facilitated
diffusion}} {[}L \emph{facilis} = easy + \emph{diffundere} = to pour
out{]} --- Passive transport of molecules down their concentration
gradient via membrane proteins (channels or carriers); no energy
required. Contrast: \textbf{active transport}. \texorpdfstring{\ensuremath{\to}}{->}
\cref{sec:unit_II_membrane_transport}

\protect\phantomsection\label{gl:fitness}{\textbf{Fitness (w)}} {[}OE
\emph{gefittan} = to be suitable{]} --- In evolutionary biology, the
relative contribution of a genotype to the next generation; often
expressed as \(w = 1 - s\) where \(s\) is the selection coefficient. \texorpdfstring{\ensuremath{\to}}{->}
\cref{sec:unit_V_population_genetics}

\protect\phantomsection\label{gl:free-energy}{\textbf{Free energy (\texorpdfstring{\ensuremath{\Delta}}{Delta}G)}}
{[}Gk \emph{eleutheros} = free{]} --- Gibbs free energy; the energy
available to do work at constant temperature and pressure:
\(\Delta G = \Delta H - T\Delta S\). Spontaneous reactions have
\(\Delta G < 0\) (exergonic). \texorpdfstring{\ensuremath{\to}}{->}
\cref{sec:unit_III_bioenergetics_and_respiration}

\begin{center}\rule{0.5\linewidth}{0.5pt}\end{center}

\subsection*{G}\label{g}
\addcontentsline{toc}{subsection}{G}

\protect\phantomsection\label{gl:glp-1}{\textbf{GLP-1 (glucagon-like
peptide-1)}} {[}L \emph{glucos} = sweet + \emph{lysis} = splitting{]}
--- An incretin hormone secreted by intestinal L cells in response to
nutrient ingestion; potentiates glucose-dependent insulin secretion
(GSIS) from beta cells; suppresses glucagon; delays gastric emptying;
reduces appetite via hypothalamic signalling. GLP-1 receptor agonists
(semaglutide, tirzepatide) are transformative in treating type 2
diabetes and obesity. \texorpdfstring{\ensuremath{\to}}{->} \cref{sec:unit_IX_endocrine_and_immune}

\protect\phantomsection\label{gl:glut4}{\textbf{GLUT4}} {[}glucose
transporter type 4{]} --- The insulin-sensitive glucose transporter
expressed in muscle and adipose tissue; stored in intracellular
vesicles; translocated to the plasma membrane within minutes of insulin
receptor activation (PI3K--Akt--AS160 cascade). Responsible for up to
90\% of post-prandial muscle glucose uptake. Impaired GLUT4 trafficking
underlies insulin resistance in type 2 diabetes. \texorpdfstring{\ensuremath{\to}}{->}
\cref{sec:unit_IX_endocrine_and_immune}

\protect\phantomsection\label{gl:gamete}{\textbf{Gamete}} {[}Gk
\emph{gamete} = wife; \emph{gametes} = husband; from \emph{gamein} = to
marry{]} --- A haploid reproductive cell (sperm or egg) that fuses with
another gamete during fertilisation to restore the diploid number. \texorpdfstring{\ensuremath{\to}}{->}
\cref{sec:unit_V_chromosomal_inheritance}

\protect\phantomsection\label{gl:gene}{\textbf{Gene}} {[}Gk \emph{genos}
= birth/origin; coined by Johannsen, 1909{]} --- A heritable unit of
information encoded in DNA that typically specifies the sequence of a
protein or functional RNA; operationally defined as a sequence occupying
a defined chromosomal locus. \texorpdfstring{\ensuremath{\to}}{->}
\cref{sec:unit_IV_epigenetics_and_gene_regulation}

\protect\phantomsection\label{gl:genetic-code}{\textbf{Genetic code}}
{[}Gk \emph{genos} = birth + L \emph{codex} = writing block{]} --- The
mapping from 64 mRNA codons to 20 amino acids and stop signals. It is
nearly universal, degenerate, and read by tRNA anticodons during
translation. \texorpdfstring{\ensuremath{\to}}{->} \cref{sec:unit_IV_gene_expression}

\protect\phantomsection\label{gl:genome}{\textbf{Genome}} {[}Gk
\emph{genos} = birth + \emph{-ome} = total body of{]} --- The complete
set of genetic information in an organism, including nuclear
chromosomes, mitochondrial DNA, and in plants, chloroplast DNA. Human
haploid genome: about 3.2 Gb, about 20,000
protein-coding genes. \texorpdfstring{\ensuremath{\to}}{->} \cref{sec:unit_IV_mutations_and_genomics}

\protect\phantomsection\label{gl:genome-wide-association-study}{\textbf{Genome-wide
association study (GWAS)}} {[}Eng \emph{genome} + \emph{wide} +
\emph{association}{]} --- A population scan correlating millions of SNP
genotypes with a trait or disease; identifies tag loci (not always
causal genes); requires large samples and correction for multiple
testing. \texorpdfstring{\ensuremath{\to}}{->} \cref{sec:unit_IV_mutations_and_genomics}

\protect\phantomsection\label{gl:glycolysis}{\textbf{Glycolysis}} {[}Gk
\emph{glykys} = sweet + \emph{lysis} = splitting{]} --- The 10-step
cytosolic pathway that converts glucose (6C) to 2 pyruvate (3C),
yielding a net 2 ATP and 2 NADH per glucose under standard textbook
accounting. It does not require oxygen directly and is among the most
ancient and widespread central metabolic pathways, though some organisms
use modified glycolytic routes. \texorpdfstring{\ensuremath{\to}}{->}
\cref{sec:unit_III_bioenergetics_and_respiration}

\protect\phantomsection\label{gl:gram-stain}{\textbf{Gram stain}}
{[}named for Hans Christian Gram, Danish bacteriologist, 1884{]} --- A
differential microbiological stain that distinguishes
thick-peptidoglycan bacteria (Gram-positive: retain crystal violet,
purple) from thin-peptidoglycan + outer membrane bacteria
(Gram-negative: lose crystal violet, pink after safranin counterstain).
\texorpdfstring{\ensuremath{\to}}{->} \cref{sec:unit_II_cell_theory}

\begin{center}\rule{0.5\linewidth}{0.5pt}\end{center}

\subsection*{H}\label{h}
\addcontentsline{toc}{subsection}{H}

\protect\phantomsection\label{gl:hpa-axis}{\textbf{HPA axis
(hypothalamic--pituitary--adrenal axis)}} {[}Gk \emph{hypo} = under +
\emph{thalamos} = inner chamber{]} --- The neuroendocrine stress axis:
CRH (hypothalamus) \texorpdfstring{\ensuremath{\to}}{->} ACTH (anterior pituitary) \texorpdfstring{\ensuremath{\to}}{->} cortisol (adrenal
cortex). Regulated by negative feedback at hypothalamus and pituitary.
Dysregulation implicated in depression, PTSD, and Cushing's syndrome. \texorpdfstring{\ensuremath{\to}}{->}
\cref{sec:unit_IX_endocrine_and_immune}

\protect\phantomsection\label{gl:hpg-axis}{\textbf{HPG axis
(hypothalamic--pituitary--gonadal axis)}} --- GnRH (hypothalamus) \texorpdfstring{\ensuremath{\to}}{->} FSH
+ LH (anterior pituitary) \texorpdfstring{\ensuremath{\to}}{->} sex steroids + gametes (gonads). Pulsatile
GnRH (every 90 min) is required for normal FSH/LH secretion; continuous
GnRH causes receptor downregulation and gonadotropin suppression. \texorpdfstring{\ensuremath{\to}}{->}
\cref{sec:unit_IX_endocrine_and_immune}

\protect\phantomsection\label{gl:hpt-axis}{\textbf{HPT axis
(hypothalamic--pituitary--thyroid axis)}} --- TRH (hypothalamus) \texorpdfstring{\ensuremath{\to}}{->} TSH
(anterior pituitary) \texorpdfstring{\ensuremath{\to}}{->} T4/T3 (thyroid gland). T3 is the active form;
acts on nuclear receptors to increase metabolic rate, thermogenesis, and
cardiac output. Hypothyroidism: TSH \texorpdfstring{\ensuremath{\uparrow}}{up}; hyperthyroidism: TSH \texorpdfstring{\ensuremath{\downarrow}}{down}. \texorpdfstring{\ensuremath{\to}}{->}
\cref{sec:unit_IX_endocrine_and_immune}

\protect\phantomsection\label{gl:haploidy}{\textbf{Haploidy (n)}} {[}Gk
\emph{haplous} = single + \emph{eidos} = form{]} --- Condition of having
one complete set of chromosomes; characteristic of gametes and the
gametophyte generation of plants. \texorpdfstring{\ensuremath{\to}}{->}
\cref{sec:unit_V_chromosomal_inheritance}

\protect\phantomsection\label{gl:hardy-weinberg-equilibrium}{\textbf{Hardy-Weinberg
equilibrium}} {[}named for Godfrey Hardy, mathematician, and Wilhelm
Weinberg, physician, 1908{]} --- The principle that allele and genotype
frequencies remain constant in a population absent of evolutionary
forces (no selection, drift, mutation, migration, non-random mating).
Genotype frequencies: \(p^2 + 2pq + q^2 = 1\). \texorpdfstring{\ensuremath{\to}}{->}
\cref{sec:unit_V_population_genetics}

\protect\phantomsection\label{gl:heterozygous}{\textbf{Heterozygous}}
{[}Gk \emph{heteros} = different + \emph{zygotos} = yoked{]} ---
Possessing two different alleles at a given locus (e.g., \emph{Aa}).
Contrast: \textbf{homozygous} (\emph{AA} or \emph{aa}). \texorpdfstring{\ensuremath{\to}}{->}
\cref{sec:unit_V_mendelian_genetics}

\protect\phantomsection\label{gl:hill-coefficient}{\textbf{Hill
coefficient (n)}} {[}named for Archibald Hill, 1910{]} --- A measure of
cooperativity in ligand binding. \(n = 1\) = no cooperativity; \(n > 1\)
= positive cooperativity (sigmoidal binding curve). Haemoglobin:
\(n \texorpdfstring{\ensuremath{\approx}}{~} 2.8\). \texorpdfstring{\ensuremath{\to}}{->} \cref{sec:unit_I_enzymes_and_kinetics}

\protect\phantomsection\label{gl:homeostasis}{\textbf{Homeostasis}}
{[}Gk \emph{homoios} = similar + \emph{stasis} = standing/state{]} ---
The maintenance of a relatively stable internal environment despite
external fluctuations; achieved by negative feedback loops with a
sensor, comparator (integrator), and effector. Coined by Walter Cannon
(1926, building on Claude Bernard). \texorpdfstring{\ensuremath{\to}}{->} \cref{sec:unit_IX_nervous_system}

\protect\phantomsection\label{gl:heterochromatin}{\textbf{Heterochromatin}}
{[}Gk \emph{heteros} = different + \emph{chroma} = colour{]} --- Densely
compacted chromatin bearing repressive histone marks (H3K9me3, H3K27me3)
and high DNA methylation; transcriptionally silent.
\textbf{Constitutive} heterochromatin (centromeres, telomeres) is
permanent; \textbf{facultative} heterochromatin (e.g., the inactive X
chromosome) is cell-type specific. \texorpdfstring{\ensuremath{\to}}{->}
\cref{sec:unit_IV_epigenetics_and_gene_regulation}

\protect\phantomsection\label{gl:histone}{\textbf{Histone}} {[}Gk
\emph{histos} = loom/web{]} --- One of five abundant nuclear proteins
(H1, H2A, H2B, H3, H4) that form the nucleosome core and linker. Tails
are post-translationally modified by acetylation, methylation,
phosphorylation, and ubiquitination to regulate gene expression (histone
code). \texorpdfstring{\ensuremath{\to}}{->} \cref{sec:unit_IV_epigenetics_and_gene_regulation}

\protect\phantomsection\label{gl:histone-code}{\textbf{Histone code}}
{[}Gk \emph{histos} = loom; L \emph{codex} = writing block{]} --- The
hypothesis (Strahl \& Allis, 2000) that combinations of histone
post-translational modifications specify unique chromatin states;
`written' by HATs/HMTs, `read' by bromodomains/chromodomains, `erased'
by HDACs/KDMs. \texorpdfstring{\ensuremath{\to}}{->} \cref{sec:unit_IV_epigenetics_and_gene_regulation}

\protect\phantomsection\label{gl:holobiont}{\textbf{Holobiont}} {[}Gk
\emph{holos} = whole + \emph{bios} = life{]} --- The host plus its
persistent associated microbiota (and sometimes viruses) treated as an
integrated biological system; useful for analysing nutrition, immunity,
and niche construction. \texorpdfstring{\ensuremath{\to}}{->}
\cref{sec:unit_VII_bacteria_archaea_viruses,sec:unit_VII_microbial_ecology}

\protect\phantomsection\label{gl:hydrogen-bond}{\textbf{Hydrogen bond}}
{[}Gk \emph{hydor} = water + \emph{genes} = born from (as in generated
with water){]} --- A weak electrostatic attraction (2--40 kJ/mol)
between a hydrogen atom covalently bonded to an electronegative atom (N,
O, or F) and another electronegative atom. Critical for DNA
base-pairing, protein secondary structure, and water's unique
properties. \texorpdfstring{\ensuremath{\to}}{->} \cref{sec:unit_I_atoms_molecules}

\protect\phantomsection\label{gl:hydrolysis}{\textbf{Hydrolysis}} {[}Gk
\emph{hydor} = water + \emph{lysis} = loosening{]} --- Cleavage of a
chemical bond by addition of water; central to digestion, ATP turnover,
polymer breakdown, and many enzyme-catalysed reaction mechanisms. \texorpdfstring{\ensuremath{\to}}{->}
\cref{sec:unit_I_macromolecules,sec:unit_I_enzymes_and_kinetics}

\protect\phantomsection\label{gl:hydrophobic-effect}{\textbf{Hydrophobic
effect}} {[}Gk \emph{hydor} = water + \emph{phobos} = fear{]} --- The
entropy-driven clustering of nonpolar surfaces in water as ordered water
cages are released. It drives membrane assembly, protein folding, and
ligand binding. \texorpdfstring{\ensuremath{\to}}{->}
\cref{sec:unit_I_water_and_life,sec:unit_I_macromolecules}

\begin{center}\rule{0.5\linewidth}{0.5pt}\end{center}

\subsection*{I}\label{i}
\addcontentsline{toc}{subsection}{I}

\protect\phantomsection\label{gl:insulin-receptor}{\textbf{Insulin
receptor}} {[}L \emph{insula} = island, as in islets of Langerhans{]}
--- A receptor tyrosine kinase (RTK) composed of two \texorpdfstring{\ensuremath{\alpha}}{alpha} and two \texorpdfstring{\ensuremath{\beta}}{beta}
subunits; insulin binding activates RTK, triggering autophosphorylation
and IRS-1/2 phosphorylation; downstream: PI3K--Akt (GLUT4 translocation,
glycogen synthesis) and MAPK (gene expression, cell growth). Type 2
diabetes involves insulin receptor signalling resistance. \texorpdfstring{\ensuremath{\to}}{->}
\cref{sec:unit_IX_endocrine_and_immune}

\protect\phantomsection\label{gl:independent-assortment}{\textbf{Independent
assortment}} {[}L \emph{independens} = not hanging upon{]} --- Mendel's
second law: genes on non-homologous chromosomes are inherited
independently of each other, producing 9:3:3:1 ratios in dihybrid
crosses. Applies only to genes on different chromosomes or far apart on
the same chromosome. \texorpdfstring{\ensuremath{\to}}{->} \cref{sec:unit_V_mendelian_genetics}

\protect\phantomsection\label{gl:intron}{\textbf{Intron}} {[}L
\emph{intra} = within; coined by Sharp \& Roberts, 1977{]} --- A
non-coding intervening sequence within a eukaryotic gene that is
transcribed but removed from the pre-mRNA by the spliceosome before
translation. Nobel Prize 1993 (Sharp \& Roberts). \texorpdfstring{\ensuremath{\to}}{->}
\cref{sec:unit_IV_gene_expression}

\begin{center}\rule{0.5\linewidth}{0.5pt}\end{center}

\subsection*{K}\label{k}
\addcontentsline{toc}{subsection}{K}

\protect\phantomsection\label{gl:keystone-species}{\textbf{Keystone
species}} {[}L \emph{lapis} = stone; \emph{clavus} = key{]} --- A
species with a disproportionately large effect on its ecosystem relative
to its biomass; its removal causes dramatic shifts in community
structure. Term coined by Robert Paine (1966) based on sea-star
(\emph{Pisaster}) removal experiments. \texorpdfstring{\ensuremath{\to}}{->}
\cref{sec:unit_X_community_ecology}

\protect\phantomsection\label{gl:km}{\textbf{Km (Michaelis constant)}}
{[}after Leonor Michaelis{]} --- The substrate concentration at which an
enzyme operates at half its maximum velocity (V\texorpdfstring{\textsubscript{m}}{m}\texorpdfstring{\textsubscript{a}}{a}\texorpdfstring{\textsubscript{x}}{x}/2); reflects
enzyme-substrate affinity (lower K\texorpdfstring{\textsubscript{m}}{m} = higher affinity). Units: mol/L
(M). \texorpdfstring{\ensuremath{\to}}{->} \cref{sec:unit_I_enzymes_and_kinetics}

\begin{center}\rule{0.5\linewidth}{0.5pt}\end{center}

\subsection*{L}\label{l}
\addcontentsline{toc}{subsection}{L}

\protect\phantomsection\label{gl:ltp}{\textbf{LTP (long-term
potentiation)}} {[}L \emph{longum} = long + \emph{tempus} = time +
\emph{potentia} = power{]} --- A persistent (hours to days) increase in
synaptic efficacy following high-frequency stimulation; requires NMDA
receptor activation \texorpdfstring{\ensuremath{\to}}{->} Ca\texorpdfstring{\textsuperscript{2}}{2}\texorpdfstring{\textsuperscript{+}}{+} influx \texorpdfstring{\ensuremath{\to}}{->} CaMKII autophosphorylation \texorpdfstring{\ensuremath{\to}}{->} AMPA
receptor insertion; the cellular model of learning and memory (Bløndell,
Lømo 1973). \texorpdfstring{\ensuremath{\to}}{->} \cref{sec:unit_IX_action_potential_synapses}

\protect\phantomsection\label{gl:ltd}{\textbf{LTD (long-term
depression)}} {[}L \emph{deprimere} = to press down{]} --- A persistent
decrease in synaptic efficacy; induced by low-frequency stimulation (1
Hz) or by mGluR/endocannabinoid signalling (``endocannabinoid-mediated
LTD''); involves PP1/PP2A phosphatase activation and AMPA receptor
internalisation via dynamin-clathrin endocytosis. \texorpdfstring{\ensuremath{\to}}{->}
\cref{sec:unit_IX_action_potential_synapses}

\protect\phantomsection\label{gl:linkage}{\textbf{Linkage}} {[}ON
\emph{hlenkr} = chain link{]} --- The tendency of genes located on the
same chromosome to be inherited together; the closer the genes, the
stronger the linkage (lower recombination frequency). \texorpdfstring{\ensuremath{\to}}{->}
\cref{sec:unit_V_chromosomal_inheritance}

\protect\phantomsection\label{gl:locus}{\textbf{Locus (plural: loci)}}
{[}L \emph{locus} = place{]} --- The fixed chromosomal position occupied
by a particular gene; defined by its physical map coordinates (base pair
position) or genetic map position (centimorgans from a reference
marker). \texorpdfstring{\ensuremath{\to}}{->} \cref{sec:unit_V_chromosomal_inheritance}

\protect\phantomsection\label{gl:lotka-volterra-equations}{\textbf{Lotka-Volterra
equations}} {[}named for Alfred Lotka (1925) and Vito Volterra (1926){]}
--- A pair of coupled nonlinear ODEs describing predator-prey population
cycles: \(dN/dt = \alpha N - \beta NP\) (prey) and
\(dP/dt = \delta NP - \gamma P\) (predator). Produce neutrally stable
oscillatory dynamics. \texorpdfstring{\ensuremath{\to}}{->} \cref{sec:unit_X_community_ecology}

\begin{center}\rule{0.5\linewidth}{0.5pt}\end{center}

\subsection*{M}\label{m}
\addcontentsline{toc}{subsection}{M}

\protect\phantomsection\label{gl:mrna}{\textbf{mRNA (messenger RNA)}}
{[}L \emph{nuntius} = messenger{]} --- An RNA molecule transcribed from
a gene's template strand; carries the codon sequence from the nucleus to
the ribosome for translation into protein. Processed in eukaryotes by 5'
capping, 3' polyadenylation, and splicing. \texorpdfstring{\ensuremath{\to}}{->}
\cref{sec:unit_IV_gene_expression}

\protect\phantomsection\label{gl:meiosis}{\textbf{Meiosis}} {[}Gk
\emph{meiosis} = lessening{]} --- Two successive cell divisions ---
Meiosis I (reductional; homologue separation) and Meiosis II
(equational; sister chromatid separation) --- that produce four haploid
gametes from one diploid cell. Sources of genetic diversity:
crossing-over and independent assortment. \texorpdfstring{\ensuremath{\to}}{->}
\cref{sec:unit_V_chromosomal_inheritance}

\protect\phantomsection\label{gl:membrane-potential}{\textbf{Membrane
potential (Vm)}} {[}L \emph{membrana} = thin skin + \emph{potentia} =
power{]} --- The electric potential difference across a biological
membrane; in neurons at rest, approximately \texorpdfstring{\ensuremath{-}}{-}70 mV (inside negative).
Governed by the Goldman equation for multiple ion species. \texorpdfstring{\ensuremath{\to}}{->}
\cref{sec:unit_IX_circulation_respiration_homeostasis}

\protect\phantomsection\label{gl:michaelis-menten-kinetics}{\textbf{Michaelis-Menten
kinetics}} {[}after Leonor Michaelis \& Maud Menten, 1913{]} --- The
standard kinetic model for enzyme catalysis:
\(v = V_{max}[S]/(K_m + [S])\). Assumes rapid equilibrium or
steady-state complex formation and one substrate. \texorpdfstring{\ensuremath{\to}}{->}
\cref{sec:unit_I_enzymes_and_kinetics}

\protect\phantomsection\label{gl:microbiome}{\textbf{Microbiome}} {[}Gk
\emph{mikros} + \emph{bios} + \emph{-ome}{]} --- The full complement of
microorganisms (bacteria, archaea, fungi, viruses) and their genes in a
defined environment (e.g., gut, rhizosphere); often profiled by
metagenomic sequencing. \texorpdfstring{\ensuremath{\to}}{->} \cref{sec:unit_VII_bacteria_archaea_viruses}

\protect\phantomsection\label{gl:microbiota}{\textbf{Microbiota}} {[}Gk
\emph{mikros} + \emph{bios} + \emph{-ota} collective{]} --- The
assemblage of microorganisms themselves (contrast \textbf{microbiome},
which often denotes genes and functions as well). \texorpdfstring{\ensuremath{\to}}{->}
\cref{sec:unit_VII_bacteria_archaea_viruses}

\protect\phantomsection\label{gl:microrna}{\textbf{MicroRNA (miRNA)}}
{[}Gk \emph{mikros} = small{]} --- A about 22 nt
single-stranded RNA that directs post-transcriptional gene silencing;
biogenesis: Drosha (nuclear) \texorpdfstring{\ensuremath{\to}}{->} pre-miRNA \texorpdfstring{\ensuremath{\to}}{->} Dicer (cytoplasmic) \texorpdfstring{\ensuremath{\to}}{->}
AGO2/RISC; partial 3'UTR complementarity causes translational repression
and mRNA decay. \textgreater{} 2,000 human miRNAs regulate
about 60\% of protein-coding genes. \texorpdfstring{\ensuremath{\to}}{->}
\cref{sec:unit_IV_epigenetics_and_gene_regulation}

\protect\phantomsection\label{gl:mitosis}{\textbf{Mitosis}} {[}Gk
\emph{mitos} = thread{]} --- Cell division in which one diploid cell
produces two genetically identical diploid daughter cells; consists of
prophase, metaphase, anaphase, and telophase followed by cytokinesis. \texorpdfstring{\ensuremath{\to}}{->}
\cref{sec:unit_IV_dna_replication_and_cell_cycle}

\protect\phantomsection\label{gl:mutation}{\textbf{Mutation}} {[}L
\emph{mutare} = to change{]} --- A heritable change in genetic sequence,
including substitutions, insertions, deletions, inversions,
duplications, and translocations. Mutation creates new sequence
variants; recombination, segregation, gene flow, and selection then
reshape how those variants are combined and distributed in populations.
\texorpdfstring{\ensuremath{\to}}{->} \cref{sec:unit_IV_mutations_and_genomics}

\protect\phantomsection\label{gl:5-methylcytosine}{\textbf{5-Methylcytosine
(5mC)}} --- The ``fifth base'' of DNA; cytosine methylated at C5 by DNMT
enzymes using S-adenosylmethionine as methyl donor. Silences promoters
via MBD protein recruitment; oxidised to 5-hydroxymethylcytosine (5hmC)
by TET enzymes for active demethylation. \texorpdfstring{\ensuremath{\to}}{->}
\cref{sec:unit_IV_epigenetics_and_gene_regulation}

\begin{center}\rule{0.5\linewidth}{0.5pt}\end{center}

\subsection*{N}\label{n}
\addcontentsline{toc}{subsection}{N}

\protect\phantomsection\label{gl:natural-selection}{\textbf{Natural
selection}} {[}L \emph{natura} = nature; \emph{selectio} = choosing{]}
--- The process by which heritable variations that increase survival or
reproduction become more common in a population over generations.
Proposed by Darwin \& Wallace (1858). \texorpdfstring{\ensuremath{\to}}{->}
\cref{sec:unit_V_population_genetics}

\protect\phantomsection\label{gl:niche}{\textbf{Niche}} {[}OFr
\emph{nicher} = to nest{]} --- The totality of an organism's role in its
environment, including physical conditions (temperature, pH, resources)
it can use and interactions with other species. \textbf{Fundamental
niche}: possible without competition. \textbf{Realised niche}: actual
range when competitors are present. \texorpdfstring{\ensuremath{\to}}{->}
\cref{sec:unit_X_community_ecology}

\protect\phantomsection\label{gl:nucleosome}{\textbf{Nucleosome}} {[}L
\emph{nucleus} = nut kernel + Gk \emph{soma} = body{]} --- The
fundamental repeating unit of chromatin: 147 bp of DNA wound
about 1.65 turns around a histone octamer (H2A, H2B, H3, H4 \texorpdfstring{\ensuremath{\times}}{x}
2); diameter 11 nm. First visualised by electron microscopy (Kornberg,
1974; Nobel Prize 2006). \texorpdfstring{\ensuremath{\to}}{->}
\cref{sec:unit_IV_epigenetics_and_gene_regulation}

\protect\phantomsection\label{gl:nucleotide}{\textbf{Nucleotide}} {[}Gk
\emph{nucleus} = kernel + L \emph{nucleus} = nut kernel{]} --- The
monomer of nucleic acids: a nitrogenous base + pentose sugar + one to
three phosphate groups. Adenine, guanine (purines); cytosine, thymine,
uracil (pyrimidines). \texorpdfstring{\ensuremath{\to}}{->}
\cref{sec:unit_IV_dna_replication_and_cell_cycle}

\begin{center}\rule{0.5\linewidth}{0.5pt}\end{center}

\subsection*{O}\label{o}
\addcontentsline{toc}{subsection}{O}

\protect\phantomsection\label{gl:organelle}{\textbf{Organelle}} {[}Gk
\emph{organon} = tool/instrument + \emph{ella} = diminutive suffix{]}
--- A membrane-bound (or in some definitions, functionally discrete)
structure within a eukaryotic cell, performing a specialised function.
Examples: nucleus, mitochondrion, chloroplast, Golgi apparatus,
lysosome. \texorpdfstring{\ensuremath{\to}}{->} \cref{sec:unit_II_cell_structure}

\protect\phantomsection\label{gl:osmosis}{\textbf{Osmosis}} {[}Gk
\emph{osmos} = pushing{]} --- The movement of water across a selectively
permeable membrane from a region of lower solute concentration (higher
water potential) to a region of higher solute concentration (lower water
potential). Van 't Hoff equation: \(\pi = iCRT\). \texorpdfstring{\ensuremath{\to}}{->}
\cref{sec:unit_II_membrane_transport}

\protect\phantomsection\label{gl:oxidative-phosphorylation}{\textbf{Oxidative
phosphorylation}} {[}L \emph{oxidare} = to oxidise + Gk
\emph{phosphoros} = light-bearer{]} --- The ATP-generating process
driven by electron transfer from NADH/FADH\texorpdfstring{\textsubscript{2}}{2} to O\texorpdfstring{\textsubscript{2}}{2} via the electron
transport chain, producing the proton-motive force used by ATP synthase.
Yields about 28 of the about 30--32 ATP per glucose. \texorpdfstring{\ensuremath{\to}}{->}
\cref{sec:unit_III_bioenergetics_and_respiration}

\begin{center}\rule{0.5\linewidth}{0.5pt}\end{center}

\subsection*{P}\label{p}
\addcontentsline{toc}{subsection}{P}

\protect\phantomsection\label{gl:p}{\textbf{P\texorpdfstring{\textsubscript{5}}{5}\texorpdfstring{\textsubscript{0}}{0}}} {[}partial pressure at
which haemoglobin is 50\% saturated{]} --- A measure of haemoglobin-O\texorpdfstring{\textsubscript{2}}{2}
affinity; normal P\texorpdfstring{\textsubscript{5}}{5}\texorpdfstring{\textsubscript{0}}{0} = 26--27 mmHg. Rightward shifts (increased P\texorpdfstring{\textsubscript{5}}{5}\texorpdfstring{\textsubscript{0}}{0}) =
decreased affinity: caused by increased PCO\texorpdfstring{\textsubscript{2}}{2}, H\texorpdfstring{\textsuperscript{+}}{+}, temperature, 2,3-BPG
(Bohr effect). Leftward shifts = increased affinity: caused by decreased
PCO\texorpdfstring{\textsubscript{2}}{2}, alkalosis, hypothermia, foetal Hb. \texorpdfstring{\ensuremath{\to}}{->}
\cref{sec:unit_IX_circulation_respiration_homeostasis}

\protect\phantomsection\label{gl:pan-genome}{\textbf{Pan-genome}} {[}Gk
\emph{pan} = all + \emph{genome}{]} --- The union of the collective
genes found across a species or clade when many strains are sequenced;
comprises a \textbf{core} genome (shared) and \textbf{accessory} genes
(variable). Explains strain-specific virulence and metabolic
capabilities in bacteria. \texorpdfstring{\ensuremath{\to}}{->} \cref{sec:unit_IV_mutations_and_genomics}

\protect\phantomsection\label{gl:poiseuille-s-law}{\textbf{Poiseuille's
law}} {[}named for Jean Léonard Marie Poiseuille, 1840{]} --- Laminar
flow of a viscous fluid in a cylindrical tube:
\(Q = \pi r^4 \Delta P / 8\eta L\), where \(r\) = radius, \(\Delta P\) =
pressure gradient, \texorpdfstring{\ensuremath{\eta}}{eta} = viscosity, \(L\) = length. The \(r^4\) dependence
explains why small changes in vessel radius (e.g., arteriolar
vasodilation) produce large changes in flow and vascular resistance. \texorpdfstring{\ensuremath{\to}}{->}
\cref{sec:unit_IX_circulation_respiration_homeostasis}

\protect\phantomsection\label{gl:peptide-bond}{\textbf{Peptide bond}}
{[}Gk \emph{pepsis} = digestion; peptides were initially thought to be
digestion products{]} --- The covalent bond (--CO--NH--) formed by
condensation between the carboxyl group of one amino acid and the amino
group of another during translation. Has partial double-bond character
(resonance; about 40\% double bond), conferring planarity. \texorpdfstring{\ensuremath{\to}}{->}
\cref{sec:unit_I_macromolecules}

\protect\phantomsection\label{gl:ph}{\textbf{pH}} {[}Danish
\emph{potenz} = power + Gk \emph{hydro} = water; coined by Sørensen,
1909{]} --- The negative common logarithm of hydrogen ion activity:
\(pH = -\log_{10}[H^+]\). Physiological range: blood plasma 7.35--7.45;
lysosomes 4.5--5.0; mitochondrial matrix 7.8--8.0. \texorpdfstring{\ensuremath{\to}}{->}
\cref{sec:unit_I_atoms_molecules}

\protect\phantomsection\label{gl:phenotype}{\textbf{Phenotype}} {[}Gk
\emph{phainein} = to show + \emph{typos} = impression{]} --- The
observable characteristics of an organism, determined by the interaction
of genotype with environment. Contrast: \textbf{genotype}. \texorpdfstring{\ensuremath{\to}}{->}
\cref{sec:unit_V_mendelian_genetics}

\protect\phantomsection\label{gl:phospholipid-bilayer}{\textbf{Phospholipid
bilayer}} {[}Gk \emph{phosphoros} = light-bearer + L \emph{lipidus} =
fat{]} --- A self-assembled membrane structure in which amphipathic
lipids orient hydrophilic heads toward aqueous environments and
hydrophobic tails inward; typical thickness is about 7--8 nm.
It is the dominant membrane architecture in bacteria and eukaryotes,
while archaeal membranes often use ether-linked isoprenoid lipids and
can form monolayers. \texorpdfstring{\ensuremath{\to}}{->} \cref{sec:unit_II_membrane_transport}

\protect\phantomsection\label{gl:photosynthesis}{\textbf{Photosynthesis}}
{[}Gk \emph{phos/photos} = light + \emph{synthesis} = putting
together{]} --- The process by which photosynthetic organisms use light
energy to produce organic compounds from CO\texorpdfstring{\textsubscript{2}}{2} and H\texorpdfstring{\textsubscript{2}}{2}O; in plants: light
reactions (thylakoids; produce ATP + NADPH) + Calvin cycle (stroma; fix
CO\texorpdfstring{\textsubscript{2}}{2} into G3P). \texorpdfstring{\ensuremath{\to}}{->} \cref{sec:unit_III_photosynthesis}

\protect\phantomsection\label{gl:plasmid}{\textbf{Plasmid}} {[}Gk
\emph{plasma} = formed thing{]} --- A small, circular, double-stranded
DNA molecule in bacteria (and some eukaryotes) that replicates
independently of the chromosome; often carries genes for antibiotic
resistance or virulence. \texorpdfstring{\ensuremath{\to}}{->} \cref{sec:unit_II_cell_theory}

\protect\phantomsection\label{gl:polymerase-chain-reaction}{\textbf{Polymerase
chain reaction (PCR)}} {[}Gk \emph{polys} = many + L \emph{merus} =
part; Gk \emph{kyklos} = circle; coined by Kary Mullis, 1983{]} --- An
\emph{in vitro} method for amplifying a specific DNA sequence using
repeated cycles of denaturation, primer annealing, and extension by Taq
polymerase. Transformed molecular biology; Nobel Prize 1993. \texorpdfstring{\ensuremath{\to}}{->}
\cref{sec:unit_IV_mutations_and_genomics}

\protect\phantomsection\label{gl:promoter}{\textbf{Promoter}} {[}L
\emph{promovere} = to move forward{]} --- A DNA sequence upstream of a
gene where RNA polymerase and transcription factors bind to initiate
transcription. In prokaryotes: \texorpdfstring{\ensuremath{-}}{-}10 (Pribnow box) and \texorpdfstring{\ensuremath{-}}{-}35 elements; in
eukaryotes: TATA box, initiator, and Pol II general transcription
factors. \texorpdfstring{\ensuremath{\to}}{->} \cref{sec:unit_IV_gene_expression}

\protect\phantomsection\label{gl:protein}{\textbf{Protein}} {[}Gk
\emph{proteios} = of the first quality; coined by Mulder \& Berzelius,
1838{]} --- A polymer of L-amino acids linked by peptide bonds; encoded
by the gene sequence; functional form depends on primary (1°), secondary
(2°), tertiary (3°), and quaternary (4°) structure. \texorpdfstring{\ensuremath{\to}}{->}
\cref{sec:unit_I_macromolecules}

\begin{center}\rule{0.5\linewidth}{0.5pt}\end{center}

\subsection*{R}\label{r}
\addcontentsline{toc}{subsection}{R}

\protect\phantomsection\label{gl:raas}{\textbf{RAAS
(renin--angiotensin--aldosterone system)}} {[}L \emph{ren} = kidney{]}
--- The hormonal blood pressure and volume regulatory cascade: low renal
perfusion \texorpdfstring{\ensuremath{\to}}{->} juxtaglomerular cells release renin \texorpdfstring{\ensuremath{\to}}{->} cleaves
angiotensinogen to angiotensin I \texorpdfstring{\ensuremath{\to}}{->} ACE (lung) converts to angiotensin II
\texorpdfstring{\ensuremath{\to}}{->} vasoconstriction + aldosterone release (Na\texorpdfstring{\textsuperscript{+}}{+}/H\texorpdfstring{\textsubscript{2}}{2}O retention). ACE
inhibitors (enalapril) and ARBs (losartan) are first-line
antihypertensive drugs. \texorpdfstring{\ensuremath{\to}}{->}
\cref{sec:unit_IX_circulation_respiration_homeostasis}

\protect\phantomsection\label{gl:recombination}{\textbf{Recombination}}
{[}L \emph{re} = again + \emph{combinare} = to combine{]} --- The
process of generating new combinations of alleles through crossing-over
between homologous chromosomes during meiosis I. Frequency (\texorpdfstring{\ensuremath{\theta}}{theta}) is
proportional to genetic map distance (1 cM = 1\% recombination). \texorpdfstring{\ensuremath{\to}}{->}
\cref{sec:unit_V_chromosomal_inheritance}

\protect\phantomsection\label{gl:reduction-potential}{\textbf{Reduction
potential (E°\texorpdfstring{'}{'})}} {[}L \emph{reducere} = to lead back{]} --- The
tendency of a molecule to accept electrons, measured in volts under
standard biochemical conditions (pH 7, 25°C). O\texorpdfstring{\textsubscript{2}}{2}/H\texorpdfstring{\textsubscript{2}}{2}O: +0.816 V;
NAD\texorpdfstring{\textsuperscript{+}}{+}/NADH: \texorpdfstring{\ensuremath{-}}{-}0.320 V. Electrons flow from lower to higher E°\texorpdfstring{'}{'}. \texorpdfstring{\ensuremath{\to}}{->}
\cref{sec:unit_I_atoms_molecules}

\protect\phantomsection\label{gl:ribonucleic-acid}{\textbf{Ribonucleic
acid (RNA)}} {[}L nucleus + suffix \emph{ic}{]} --- A single-stranded
nucleic acid using ribose and uracil (instead of deoxyribose and
thymine); roles include mRNA (messenger), tRNA (transfer), rRNA
(ribosomal), and regulatory (miRNA, siRNA, snRNA). \texorpdfstring{\ensuremath{\to}}{->}
\cref{sec:unit_IV_gene_expression}

\protect\phantomsection\label{gl:ribosome}{\textbf{Ribosome}} {[}named
for ribose + Gk \emph{soma} = body{]} --- The ribonucleoprotein complex
(about 2.5 MDa in eukaryotes; 80S = 40S + 60S subunits) that
catalyses peptide bond formation during translation; catalytic activity
resides in rRNA (a ribozyme). Nobel Prize 2009 (Ramakrishnan, Steitz,
Yonath). \texorpdfstring{\ensuremath{\to}}{->} \cref{sec:unit_IV_gene_expression}

\begin{center}\rule{0.5\linewidth}{0.5pt}\end{center}

\subsection*{S}\label{s}
\addcontentsline{toc}{subsection}{S}

\protect\phantomsection\label{gl:saltatory-conduction}{\textbf{Saltatory
conduction}} {[}L \emph{saltare} = to jump{]} --- The rapid propagation
of action potentials in myelinated axons by jumping from one node of
Ranvier to the next; velocity proportional to internodal distance; 5--70
m/s vs 0.5--2 m/s in unmyelinated C fibres. Enabled by high Na\texorpdfstring{\textsuperscript{+}}{+} channel
density at nodes and insulation by myelin between nodes. \texorpdfstring{\ensuremath{\to}}{->}
\cref{sec:unit_IX_action_potential_synapses}

\protect\phantomsection\label{gl:snare-complex}{\textbf{SNARE complex}}
{[}acronym: \textbf{S}oluble \textbf{N}-ethylmaleimide-sensitive factor
\textbf{A}ttachment protein \textbf{RE}ceptors{]} --- The core vesicle
fusion machinery: v-SNARE (synaptobrevin/VAMP-2, on vesicle) + t-SNAREs
(syntaxin-1A and SNAP-25, on target membrane) zipper together in a
four-helix bundle providing the energy for membrane fusion. Cleaved by
clostridial neurotoxins (botulinum toxin cleaves SNAP-25/VAMP; tetanus
toxin cleaves VAMP in spinal inhibitory interneurons). \texorpdfstring{\ensuremath{\to}}{->}
\cref{sec:unit_IX_action_potential_synapses}

\protect\phantomsection\label{gl:surfactant}{\textbf{Surfactant}} {[}L
\emph{superficies} = surface + \emph{activus} = active{]} --- A complex
mixture of phospholipids (primarily DPPC) and proteins (SP-A, SP-B,
SP-C, SP-D) secreted by alveolar type II pneumocytes; reduces alveolar
surface tension according to Laplace's law (\(P = 2T/r\)), preventing
alveolar collapse (atelectasis) during expiration and reducing the work
of breathing. Deficiency causes neonatal respiratory distress syndrome
(RDS). \texorpdfstring{\ensuremath{\to}}{->} \cref{sec:unit_IX_circulation_respiration_homeostasis}

\protect\phantomsection\label{gl:synaptotagmin}{\textbf{Synaptotagmin}}
{[}Gk \emph{synapsis} = junction + \emph{tagma} = arrangement{]} --- The
Ca\texorpdfstring{\textsuperscript{2}}{2}\texorpdfstring{\textsuperscript{+}}{+} sensor for fast neurotransmitter release; C2A and C2B domains bind
3--5 Ca\texorpdfstring{\textsuperscript{2}}{2}\texorpdfstring{\textsuperscript{+}}{+} ions cooperatively (Hill coefficient n\texorpdfstring{\ensuremath{\approx}}{~}4), explaining Ca\texorpdfstring{\textsuperscript{2}}{2}\texorpdfstring{\textsuperscript{+}}{+}
cooperativity of neurotransmitter release. Synaptotagmin-1 drives
synchronous fast release; synaptotagmin-7 drives asynchronous release
and facilitatory plasticity. \texorpdfstring{\ensuremath{\to}}{->}
\cref{sec:unit_IX_action_potential_synapses}

\protect\phantomsection\label{gl:selection-coefficient}{\textbf{Selection
coefficient (s)}} {[}see \emph{natural selection}{]} --- A quantitative
measure of the fitness disadvantage (or advantage) conferred by an
allele; \(s = 1 - w\) where \(w\) is relative fitness. Dominant lethal:
\(s = 1\); neutral: \(s = 0\). \texorpdfstring{\ensuremath{\to}}{->} \cref{sec:unit_V_population_genetics}

\protect\phantomsection\label{gl:speciation}{\textbf{Speciation}} {[}L
\emph{species} = appearance/kind{]} --- The evolutionary process by
which one ancestral population diverges into two or more reproductively
isolated species. Mechanisms: \textbf{allopatric} (geographic
isolation), \textbf{sympatric} (within a range; polyploidy, disruptive
selection). \texorpdfstring{\ensuremath{\to}}{->} \cref{sec:unit_VI_genetic_drift_and_speciation}

\protect\phantomsection\label{gl:species-area-relationship}{\textbf{Species-area
relationship}} {[}L \emph{species} = kind + \emph{area} = open space{]}
--- The empirical relationship \(S = cA^z\) between species richness
(\(S\)) and area (\(A\)); exponent \(z \texorpdfstring{\ensuremath{\approx}}{~} 0.25\) for continental
habitats, \(z \texorpdfstring{\ensuremath{\approx}}{~} 0.30\) for island habitats. MacArthur \& Wilson (1967).
\texorpdfstring{\ensuremath{\to}}{->} \cref{sec:unit_X_biomes_and_conservation}

\protect\phantomsection\label{gl:surface-area-to-volume-ratio}{\textbf{Surface-area-to-volume
ratio (A/V)}} {[}L \emph{superficies} = surface{]} --- For a sphere:
\(A/V = 3/r\). As cell radius increases, \(A/V\) decreases, limiting
diffusion-based nutrient exchange; this constraint accounts for why
cells are microscopically small. \texorpdfstring{\ensuremath{\to}}{->} \cref{sec:unit_II_cell_theory}

\begin{center}\rule{0.5\linewidth}{0.5pt}\end{center}

\subsection*{T}\label{t}
\addcontentsline{toc}{subsection}{T}

\protect\phantomsection\label{gl:thermodynamics}{\textbf{Thermodynamics}}
{[}Gk \emph{therme} = heat + \emph{dynamis} = power{]} --- The branch of
physical science concerned with heat, work, and energy transformations.
First Law: energy is conserved. Second Law: entropy increases in any
spontaneous process. \texorpdfstring{\ensuremath{\to}}{->} \cref{sec:unit_III_bioenergetics_and_respiration}

\protect\phantomsection\label{gl:transcription}{\textbf{Transcription}}
{[}L \emph{transcribere} = to copy{]} --- The synthesis of an RNA
molecule from a DNA template by RNA polymerase; proceeds 5\texorpdfstring{'}{'} to 3\texorpdfstring{'}{'}
reading the template strand 3\texorpdfstring{'}{'} to 5\texorpdfstring{'}{'}; regulated by transcription
factors, promoters, and enhancers. \texorpdfstring{\ensuremath{\to}}{->} \cref{sec:unit_IV_gene_expression}

\protect\phantomsection\label{gl:translation}{\textbf{Translation}} {[}L
\emph{translatio} = carrying across{]} --- The ribosome-mediated
synthesis of a polypeptide from an mRNA template; proceeds from 5\texorpdfstring{'}{'} to 3\texorpdfstring{'}{'}
(mRNA), adding amino acids from N- to C-terminus at about 15
residues/second (eukaryotes). \texorpdfstring{\ensuremath{\to}}{->} \cref{sec:unit_IV_gene_expression}

\protect\phantomsection\label{gl:transpiration}{\textbf{Transpiration}}
{[}L \emph{trans} = across + \emph{spirare} = to breathe{]} --- The
evaporation of water from leaf mesophyll cells and its diffusion through
stomata; drives the transpiration stream (cohesion-tension mechanism)
and enables mineral transport in xylem. \texorpdfstring{\ensuremath{\to}}{->}
\cref{sec:unit_VIII_plant_structure_and_water}

\protect\phantomsection\label{gl:trophic-level}{\textbf{Trophic level}}
{[}Gk \emph{trophe} = nourishment{]} --- A functional tier in a food web
defined by the number of energy-transfer steps from primary producers;
approximately 10\% of energy passes between levels (Lindeman's
trophic-dynamic theory, 1942). \texorpdfstring{\ensuremath{\to}}{->} \cref{sec:unit_X_ecosystem_ecology}

\begin{center}\rule{0.5\linewidth}{0.5pt}\end{center}

\subsection*{V}\label{v}
\addcontentsline{toc}{subsection}{V}

\protect\phantomsection\label{gl:van-der-waals-forces}{\textbf{Van der
Waals forces}} {[}named for Johannes Diderik van der Waals, 1873{]} ---
Weak, short-range intermolecular attractions (0.1--10 kJ/mol) arising
from transient (London dispersion), permanent-induced (Debye), or
permanent-permanent (Keesom) dipole interactions. Important for packing
of hydrophobic molecules and protein-ligand recognition. \texorpdfstring{\ensuremath{\to}}{->}
\cref{sec:unit_I_atoms_molecules}

\protect\phantomsection\label{gl:virulence}{\textbf{Virulence}} {[}L
\emph{virulentus} = full of poison{]} --- The capacity of a pathogen to
cause disease; quantified as the proportion of infected hosts that
develop disease or die. Virulence factors include toxins, adhesins,
capsules, and immune evasion mechanisms. \texorpdfstring{\ensuremath{\to}}{->}
\cref{sec:unit_VII_bacteria_archaea_viruses}

\begin{center}\rule{0.5\linewidth}{0.5pt}\end{center}

\subsection*{W}\label{w}
\addcontentsline{toc}{subsection}{W}

\protect\phantomsection\label{gl:water-potential}{\textbf{Water
potential (\texorpdfstring{\ensuremath{\Psi}}{Psi})}} {[}named analogy to electrical potential{]} --- A
measure of the free energy of water in a system relative to pure water:
\(\Psi = \Psi_s + \Psi_p\) (solute + pressure components). Water moves
from high \texorpdfstring{\ensuremath{\Psi}}{Psi} to low \texorpdfstring{\ensuremath{\Psi}}{Psi} (down the gradient). Units: megapascals (MPa). \texorpdfstring{\ensuremath{\to}}{->}
\cref{sec:unit_VIII_plant_structure_and_water}

\begin{center}\rule{0.5\linewidth}{0.5pt}\end{center}

\subsection*{X}\label{x}
\addcontentsline{toc}{subsection}{X}

\protect\phantomsection\label{gl:xylem}{\textbf{Xylem}} {[}Gk
\emph{xylon} = wood{]} --- The vascular tissue that conducts water and
dissolved minerals from roots to aerial organs; composed of non-living
tracheids and vessel members (dead at maturity) plus parenchyma and
fibres. Water moves by cohesion-tension mechanism. \texorpdfstring{\ensuremath{\to}}{->}
\cref{sec:unit_VIII_plant_structure_and_water}

\begin{center}\rule{0.5\linewidth}{0.5pt}\end{center}

\subsection*{Z}\label{z}
\addcontentsline{toc}{subsection}{Z}

\protect\phantomsection\label{gl:zygote}{\textbf{Zygote}} {[}Gk
\emph{zygotos} = yoked{]} --- The diploid cell formed by the fusion of
two haploid gametes during fertilisation; the first cell of a new
organism; undergoes mitotic divisions to form the embryo. \texorpdfstring{\ensuremath{\to}}{->}
\cref{sec:unit_V_mendelian_genetics}

\begin{center}\rule{0.5\linewidth}{0.5pt}\end{center}

\emph{This glossary covers about 120 core terms.
Chapter-specific key terms are listed in each chapter's Key Terms
section. For full index, see the combined PDF output.}

\begin{center}\rule{0.5\linewidth}{0.5pt}\end{center}

\subsection*{Additional Terms (Units III --
X)}\label{additional-terms-units-iii-x}
\addcontentsline{toc}{subsection}{Additional Terms (Units III -- X)}

\subsubsection*{C (continued)}\label{c-continued}
\addcontentsline{toc}{subsubsection}{C (continued)}

\protect\phantomsection\label{gl:carrying-capacity}{\textbf{Carrying
capacity (K)}} {[}L \emph{capax} = able to hold{]} --- The maximum
number of individuals an environment can support indefinitely,
determined by resource availability; the equilibrium of the logistic
growth model dN/dt = rN(1 \texorpdfstring{\ensuremath{-}}{-} N/K). \texorpdfstring{\ensuremath{\to}}{->} \cref{sec:unit_X_population_ecology}

\protect\phantomsection\label{gl:chloroplast}{\textbf{Chloroplast}}
{[}Gk \emph{chloros} = green + \emph{plastos} = formed{]} --- A
double-membrane plant organelle derived from a cyanobacterial
endosymbiont; contains the thylakoid membranes (light reactions) and
stroma (Calvin cycle); its own circular genome encodes about 80
proteins. \texorpdfstring{\ensuremath{\to}}{->} \cref{sec:unit_III_photosynthesis}

\protect\phantomsection\label{gl:coevolution}{\textbf{Coevolution}} {[}L
\emph{co-} = together + \emph{evolutio} = unrolling{]} --- The
reciprocal evolutionary change in two or more species exerting strong
selective pressure on each other; classic examples: plant-pollinator
mutualism, host-parasite arms races, flowering plant-butterfly chemical
coevolution (Ehrlich \& Raven, 1964). \texorpdfstring{\ensuremath{\to}}{->}
\cref{sec:unit_X_community_ecology}

\protect\phantomsection\label{gl:competitive-exclusion}{\textbf{Competitive
exclusion}} {[}L \emph{competere} = to strive together +
\emph{excludere} = to shut out{]} --- Gause's competitive exclusion
principle (1934): two species competing for identical limiting resources
cannot coexist indefinitely at equilibrium; one will be excluded. \texorpdfstring{\ensuremath{\to}}{->}
\cref{sec:unit_X_community_ecology}

\protect\phantomsection\label{gl:copepod}{\textbf{Copepod}} {[}Gk
\emph{kope} = oar + \emph{pous/podos} = foot{]} --- A crustacean
subclass of small aquatic organisms forming a major component of marine
and freshwater zooplankton; primary consumers linking phytoplankton to
larger predators in ocean food webs. \texorpdfstring{\ensuremath{\to}}{->}
\cref{sec:unit_X_ecosystem_ecology}

\subsubsection*{D (continued)}\label{d-continued}
\addcontentsline{toc}{subsubsection}{D (continued)}

\protect\phantomsection\label{gl:decomposer}{\textbf{Decomposer}} {[}L
\emph{de} = down + \emph{componere} = put together{]} --- An organism
(primarily bacteria and fungi) that breaks dead organic matter into
inorganic nutrients through enzymatic digestion; drives nutrient cycling
in every ecosystem. \texorpdfstring{\ensuremath{\to}}{->} \cref{sec:unit_X_ecosystem_ecology}

\protect\phantomsection\label{gl:density-dependent-regulation}{\textbf{Density-dependent
regulation}} {[}L \emph{densus} = thick{]} --- A population regulatory
mechanism whose effect (per-capita birth or death rate) changes as a
function of population density; resource competition, disease, and
predation are density-dependent. Contrast: \textbf{density-independent}
factors (e.g., floods, fires). \texorpdfstring{\ensuremath{\to}}{->} \cref{sec:unit_X_population_ecology}

\protect\phantomsection\label{gl:depolarisation}{\textbf{Depolarisation}}
{[}L \emph{de} = away from + Gk \emph{polos} = pivot{]} --- A decrease
in magnitude of membrane potential (movement toward zero or positive
values); in neurons, caused by Na\texorpdfstring{\textsuperscript{+}}{+} influx through voltage-gated channels
during the rising phase of an action potential. Contrast:
\textbf{repolarisation}, \textbf{hyperpolarisation}. \texorpdfstring{\ensuremath{\to}}{->}
\cref{sec:unit_IX_nervous_system}

\subsubsection*{E (continued)}\label{e-continued}
\addcontentsline{toc}{subsubsection}{E (continued)}

\protect\phantomsection\label{gl:effective-population-size}{\textbf{Effective
population size (N\texorpdfstring{\textsubscript{e}}{e})}} {[}Gk \emph{effectivus} = productive{]} --- An
idealised census number equivalent to the actual population in terms of
rate of genetic drift; typically much smaller than the census size due
to unequal sex ratios, variance in reproductive success, and population
size fluctuations. \texorpdfstring{\ensuremath{\sigma}}{sigma}\texorpdfstring{\textsuperscript{2}}{2}(\texorpdfstring{\ensuremath{\Delta}}{Delta}p) = pq/2N\texorpdfstring{\textsubscript{e}}{e}. \texorpdfstring{\ensuremath{\to}}{->}
\cref{sec:unit_VI_genetic_drift_and_speciation}

\protect\phantomsection\label{gl:epiphyte}{\textbf{Epiphyte}} {[}Gk
\emph{epi} = upon + \emph{phyton} = plant{]} --- A plant that grows on
another plant for physical support without parasitising it; absorbs
water and nutrients from rain and air; orchids, bromeliads, and ferns
are prominent epiphytes. \texorpdfstring{\ensuremath{\to}}{->} \cref{sec:unit_X_biomes_and_conservation}

\protect\phantomsection\label{gl:eutrophication}{\textbf{Eutrophication}}
{[}Gk \emph{eutrophos} = well-nourished{]} --- Nutrient enrichment of an
aquatic body (typically N and P from agricultural runoff) triggering
algal blooms; subsequent decomposition of algal biomass depletes
dissolved oxygen, causing hypoxic dead zones. \texorpdfstring{\ensuremath{\to}}{->}
\cref{sec:unit_X_ecosystem_ecology}

\subsubsection*{F (continued)}\label{f-continued}
\addcontentsline{toc}{subsubsection}{F (continued)}

\protect\phantomsection\label{gl:fermentation}{\textbf{Fermentation}}
{[}L \emph{fermentum} = leaven{]} --- Anaerobic metabolic process
regenerating NAD\texorpdfstring{\textsuperscript{+}}{+} by reducing pyruvate to lactate (lactic acid
fermentation) or to ethanol + CO\texorpdfstring{\textsubscript{2}}{2} (alcoholic fermentation); net gain: 2
ATP/glucose. \texorpdfstring{\ensuremath{\to}}{->} \cref{sec:unit_III_bioenergetics_and_respiration}

\protect\phantomsection\label{gl:founder-effect}{\textbf{Founder
effect}} {[}OE \emph{funden} = to establish{]} --- A form of genetic
drift in which allele frequencies in a new population established from a
small group of founders differ from the source population by chance; can
lead to high prevalence of otherwise rare alleles. \texorpdfstring{\ensuremath{\to}}{->}
\cref{sec:unit_VI_genetic_drift_and_speciation}

\protect\phantomsection\label{gl:frond}{\textbf{Frond}} {[}L
\emph{frons/frondis} = leafy branch{]} --- The large leaf of a fern or
other pteridophyte; bears sporangia on its underside (in a sorus);
anatomically equivalent to a compound leaf. \texorpdfstring{\ensuremath{\to}}{->}
\cref{sec:unit_VIII_plant_structure_and_water}

\subsubsection*{G (continued)}\label{g-continued}
\addcontentsline{toc}{subsubsection}{G (continued)}

\protect\phantomsection\label{gl:genetic-drift}{\textbf{Genetic drift}}
{[}ON \emph{drift} = a driving of snow{]} --- Random, chance changes in
allele frequencies in a population from generation to generation due to
sampling error in finite populations; stronger effect in small
populations. \texorpdfstring{\ensuremath{\to}}{->} \cref{sec:unit_VI_genetic_drift_and_speciation}

\protect\phantomsection\label{gl:gpp}{\textbf{GPP (gross primary
productivity)}} {[}Gk \emph{ge} = earth + \emph{protos} = first + L
\emph{producere}{]} --- The total rate of photosynthetic carbon fixation
by producers; GPP \texorpdfstring{\ensuremath{-}}{-} Ra (autotrophic respiration) = NPP (net primary
productivity). \texorpdfstring{\ensuremath{\to}}{->} \cref{sec:unit_X_ecosystem_ecology}

\protect\phantomsection\label{gl:guard-cell}{\textbf{Guard cell}} {[}ME
\emph{garde} = protection + L \emph{cella} = small room{]} --- One of
two kidney-shaped epidermal cells flanking a stoma; changes in turgor
pressure (driven by K\texorpdfstring{\textsuperscript{+}}{+} and water movement) open or close the stomatal
pore to regulate gas exchange and water loss. \texorpdfstring{\ensuremath{\to}}{->}
\cref{sec:unit_VIII_plant_structure_and_water}

\subsubsection*{H (continued)}\label{h-continued}
\addcontentsline{toc}{subsubsection}{H (continued)}

\protect\phantomsection\label{gl:herd-immunity}{\textbf{Herd immunity}}
{[}ME \emph{herd} = flock + L \emph{immunis} = exempt{]} --- Indirect
protection of susceptible individuals when a sufficiently large
proportion of the population is immune (vaccinated or recovered),
reducing pathogen transmission; threshold: p\_c = 1 \texorpdfstring{\ensuremath{-}}{-} 1/R\texorpdfstring{\textsubscript{0}}{0}. \texorpdfstring{\ensuremath{\to}}{->}
\cref{sec:unit_VII_bacteria_archaea_viruses}

\protect\phantomsection\label{gl:herbivory}{\textbf{Herbivory}} {[}L
\emph{herba} = grass + \emph{vorare} = to devour{]} --- Consumption of
plant or algal tissue by an animal or other organism; herbivores occupy
the second trophic level; plant defences (thorns, secondary metabolites,
trichomes) and herbivore counter-adaptations drive coevolution. \texorpdfstring{\ensuremath{\to}}{->}
\cref{sec:unit_X_community_ecology}

\protect\phantomsection\label{gl:hormone}{\textbf{Hormone}} {[}Gk
\emph{hormaein} = to set in motion{]} --- A chemical signal produced in
one part of an organism that travels (via blood in animals, phloem in
plants) to act on distant target cells bearing specific receptors; plant
hormones (phytohormones) include auxin, gibberellin, cytokinin, abscisic
acid, and ethylene. \texorpdfstring{\ensuremath{\to}}{->}
\cref{sec:unit_VIII_plant_reproduction,sec:unit_IX_endocrine_and_immune}

\subsubsection*{I (continued)}\label{i-continued}
\addcontentsline{toc}{subsubsection}{I (continued)}

\protect\phantomsection\label{gl:innate-immunity}{\textbf{Innate
immunity}} {[}L \emph{innatus} = inborn{]} --- The non-specific,
evolutionarily ancient first-line defence against pathogens; includes
physical barriers (skin, mucous membranes), phagocytes (neutrophils,
macrophages), complement, and pattern recognition receptors (TLRs
detecting PAMPs). Acts within minutes to hours. \texorpdfstring{\ensuremath{\to}}{->}
\cref{sec:unit_IX_endocrine_and_immune}

\protect\phantomsection\label{gl:intrinsic-rate-of-increase}{\textbf{Intrinsic
rate of increase (r)}} {[}L \emph{intrinsecus} = inward{]} --- The
per-capita growth rate of a population under ideal (unlimited)
conditions; r = ln(\texorpdfstring{\ensuremath{\lambda}}{lambda}) where \texorpdfstring{\ensuremath{\lambda}}{lambda} is the per-generation finite rate of
increase; positive r = growing population; r = 0 at K. \texorpdfstring{\ensuremath{\to}}{->}
\cref{sec:unit_X_population_ecology}

\protect\phantomsection\label{gl:island-biogeography}{\textbf{Island
biogeography}} {[}OE \emph{ealand} = island + Gk \emph{ge} = earth +
\emph{graphein} = to write{]} --- The study of factors affecting species
richness of islands (or habitat patches); key model: species richness at
equilibrium = balance between immigration rate and extinction rate;
species-area law \(S = cA^z\) (MacArthur \& Wilson, 1967). \texorpdfstring{\ensuremath{\to}}{->}
\cref{sec:unit_X_biomes_and_conservation}

\subsubsection*{K (continued)}\label{k-continued}
\addcontentsline{toc}{subsubsection}{K (continued)}

\protect\phantomsection\label{gl:k-strategist}{\textbf{K-strategist}}
{[}see \emph{carrying capacity} + Gk \emph{strategia} = strategy{]} ---
An organism with a life history strategy maximising competitive success
near carrying capacity: slow growth, long lifespan, large body, few
offspring with high parental investment. Contrast:
\textbf{r-strategist}. \texorpdfstring{\ensuremath{\to}}{->} \cref{sec:unit_X_population_ecology}

\subsubsection*{L (continued)}\label{l-continued}
\addcontentsline{toc}{subsubsection}{L (continued)}

\protect\phantomsection\label{gl:life-table}{\textbf{Life table}} {[}ME
\emph{lyf} + L \emph{tabula} = flat board{]} --- A tabular display of
age-specific survival (lx) and fecundity (mx) values for a cohort; used
to calculate net reproductive rate (R\texorpdfstring{\textsubscript{0}}{0} = \texorpdfstring{\ensuremath{\Sigma}}{Sigma} lx\texorpdfstring{\ensuremath{\cdot}}{.}mx) and intrinsic rate of
increase (r). \texorpdfstring{\ensuremath{\to}}{->} \cref{sec:unit_X_population_ecology}

\protect\phantomsection\label{gl:logistic-growth}{\textbf{Logistic
growth}} {[}Gk \emph{logistikos} = skilled in reckoning{]} ---
Population growth model dN/dt = rN(1 \texorpdfstring{\ensuremath{-}}{-} N/K), producing a sigmoidal
S-shaped curve; growth decelerates as population approaches carrying
capacity K; introduced by Verhulst (1838). \texorpdfstring{\ensuremath{\to}}{->}
\cref{sec:unit_X_population_ecology}

\subsubsection*{M (continued)}\label{m-continued}
\addcontentsline{toc}{subsubsection}{M (continued)}

\protect\phantomsection\label{gl:minimum-viable-population}{\textbf{Minimum
viable population (MVP)}} {[}L \emph{minimus} = smallest + \emph{vita} =
life{]} --- The smallest population size with a specified probability
(e.g., 99\%) of persisting for a given period (e.g., 1,000 years);
accounts for demographic stochasticity, environmental stochasticity, and
genetic factors; used in conservation planning. \texorpdfstring{\ensuremath{\to}}{->}
\cref{sec:unit_X_biomes_and_conservation}

\protect\phantomsection\label{gl:molecular-clock}{\textbf{Molecular
clock}} {[}Gk \emph{molecula} + OE \emph{clokke}{]} --- The approximate
constancy of the rate of neutral nucleotide substitution in lineages;
enables dating of evolutionary divergence events: t = d/(2\texorpdfstring{\ensuremath{\mu}}{mu}) where d =
sequence divergence, \texorpdfstring{\ensuremath{\mu}}{mu} = substitution rate per year. Calibrated using
the fossil record. \texorpdfstring{\ensuremath{\to}}{->} \cref{sec:unit_VI_genetic_drift_and_speciation}

\protect\phantomsection\label{gl:mutualism}{\textbf{Mutualism}} {[}L
\emph{mutuus} = mutual, reciprocal{]} --- A symbiotic relationship
benefiting both participating species (+/+); examples: mycorrhizal fungi
increasing plant phosphorus uptake, nitrogen-fixing \emph{Rhizobium} in
legume root nodules, cleaner fish and reef fish. \texorpdfstring{\ensuremath{\to}}{->}
\cref{sec:unit_X_community_ecology}

\protect\phantomsection\label{gl:mycorrhiza}{\textbf{Mycorrhiza}} {[}Gk
\emph{mykes} = fungus + \emph{rhiza} = root{]} --- A mutualistic
association between soil fungi (Glomeromycota for arbuscular type;
Basidiomycota/Ascomycota for ectomycorrhizal type) and plant roots;
fungi receive photosynthate, plant receives water and phosphorus;
present in about 80\% of land plants. \texorpdfstring{\ensuremath{\to}}{->}
\cref{sec:unit_VIII_plant_structure_and_water}

\subsubsection*{N (continued)}\label{n-continued}
\addcontentsline{toc}{subsubsection}{N (continued)}

\protect\phantomsection\label{gl:neuron}{\textbf{Neuron}} {[}Gk
\emph{neuron} = nerve, sinew{]} --- The basic signalling unit of the
nervous system; consists of a cell body (soma), dendrites (input), and
an axon (output); classified by function (sensory, motor, interneuron)
or morphology; communicates via electrical (action potentials) or
chemical (synaptic) signals. \texorpdfstring{\ensuremath{\to}}{->} \cref{sec:unit_IX_nervous_system}

\protect\phantomsection\label{gl:neutral-theory}{\textbf{Neutral
theory}} {[}L \emph{neutralis} = neither + Gk \emph{theoria} = view{]}
--- Kimura's (1968) theory that the majority of molecular evolutionary
change is selectively neutral (neither advantageous nor deleterious);
neutral mutations drift to fixation at a rate equal to the mutation rate
(u), independent of population size. \texorpdfstring{\ensuremath{\to}}{->}
\cref{sec:unit_VI_genetic_drift_and_speciation}

\protect\phantomsection\label{gl:nitrification}{\textbf{Nitrification}}
{[}L \emph{nitrum} = saltpetre + \emph{facere} = to make{]} --- The
aerobic microbial oxidation of ammonium (NH\texorpdfstring{\textsubscript{4}}{4}\texorpdfstring{\textsuperscript{+}}{+}) to nitrite (NO\texorpdfstring{\textsubscript{2}}{2}\texorpdfstring{\textsuperscript{-}}{-}) by
\emph{Nitrosomonas}, then nitrite to nitrate (NO\texorpdfstring{\textsubscript{3}}{3}\texorpdfstring{\textsuperscript{-}}{-}) by
\emph{Nitrobacter}; a key step in the nitrogen cycle making nitrogen
available to most plants. \texorpdfstring{\ensuremath{\to}}{->} \cref{sec:unit_X_ecosystem_ecology}

\protect\phantomsection\label{gl:nitrogen-fixation}{\textbf{Nitrogen
fixation}} {[}L \emph{nitrum} + \emph{fixatio} = making firm{]} --- The
enzymatic reduction of atmospheric N\texorpdfstring{\textsubscript{2}}{2} to NH\texorpdfstring{\textsubscript{3}}{3} by nitrogenase; performed
by free-living bacteria (\emph{Azotobacter}, cyanobacteria) and
symbionts (\emph{Rhizobium}, \emph{Frankia}); rate about 120 Tg
N/yr globally. \texorpdfstring{\ensuremath{\to}}{->} \cref{sec:unit_VII_bacteria_archaea_viruses}

\subsubsection*{O (continued)}\label{o-continued}
\addcontentsline{toc}{subsubsection}{O (continued)}

\protect\phantomsection\label{gl:operon}{\textbf{Operon}} {[}L
\emph{opus} = work{]} --- A prokaryotic gene cluster transcribed as a
polycistronic mRNA from a single promoter; regulated by
operator/repressor interactions; the \emph{lac} operon (Jacob \& Monod,
1961) was the first characterised gene regulatory system; Nobel Prize
1965. \texorpdfstring{\ensuremath{\to}}{->} \cref{sec:unit_IV_gene_expression}

\subsubsection*{P (continued)}\label{p-continued}
\addcontentsline{toc}{subsubsection}{P (continued)}

\protect\phantomsection\label{gl:parasite}{\textbf{Parasite}} {[}Gk
\emph{parasitos} = one who eats beside another{]} --- An organism that
benefits from living in or on a host organism at the host's expense
(+/\texorpdfstring{\ensuremath{-}}{-}); parasites reduce host fitness without immediately killing the
host (unlike predators); include worms (helminths), protozoa, and many
bacteria and fungi. \texorpdfstring{\ensuremath{\to}}{->} \cref{sec:unit_VII_bacteria_archaea_viruses}

\protect\phantomsection\label{gl:photoperiodism}{\textbf{Photoperiodism}}
{[}Gk \emph{photos} = light + L \emph{periodus} = cycle{]} --- The
response of an organism to the relative duration of day and night
(photoperiod); controls flowering (long-day vs.~short-day
vs.~day-neutral plants) via phytochrome, a red/far-red reversible
photoreceptor. \texorpdfstring{\ensuremath{\to}}{->} \cref{sec:unit_VIII_plant_responses}

\protect\phantomsection\label{gl:phytochrome}{\textbf{Phytochrome}}
{[}Gk \emph{phyton} = plant + \emph{chroma} = colour{]} --- A plant
photoreceptor protein interconvertible between a red-light-absorbing
form (Pr, 660 nm) and a far-red-light-absorbing form (Pfr, 730 nm); Pfr
is the biologically active form; governs seed germination, shade
avoidance, and flowering. \texorpdfstring{\ensuremath{\to}}{->} \cref{sec:unit_VIII_plant_responses}

\protect\phantomsection\label{gl:phytoplankton}{\textbf{Phytoplankton}}
{[}Gk \emph{phyton} = plant + \emph{planktos} = wandering{]} ---
Microscopic photosynthetic organisms (cyanobacteria, diatoms,
dinoflagellates) drifting in the photic zone of aquatic ecosystems;
responsible for about 50\% of Earth's net primary productivity.
\texorpdfstring{\ensuremath{\to}}{->} \cref{sec:unit_X_ecosystem_ecology}

\protect\phantomsection\label{gl:population-genetics}{\textbf{Population
genetics}} {[}L \emph{populus} = people + Gk \emph{genos} = birth +
\emph{statistike} = statecraft{]} --- The study of allele and genotype
frequency distributions and their changes in populations; mathematical
foundation provided by Hardy, Weinberg (1908), Fisher, Wright, and
Haldane (1930s). \texorpdfstring{\ensuremath{\to}}{->} \cref{sec:unit_V_population_genetics}

\protect\phantomsection\label{gl:primary-succession}{\textbf{Primary
succession}} {[}L \emph{primus} = first + \emph{successio} =
following{]} --- Ecological succession on a surface not previously
colonised (e.g., bare rock after a lava flow or retreating glacier);
begins with pioneer species (lichenised fungi, mosses) and proceeds
through intermediate stages to a climax community. \texorpdfstring{\ensuremath{\to}}{->}
\cref{sec:unit_X_community_ecology}

\protect\phantomsection\label{gl:proton-motive-force}{\textbf{Proton-motive
force (\texorpdfstring{\ensuremath{\Delta}}{Delta}p)}} {[}Gk \emph{proton} = first + L \emph{movere} = to move{]}
--- The electrochemical driving force for proton movement across a
membrane, composed of a chemical gradient (\texorpdfstring{\ensuremath{\Delta}}{Delta}pH) and an electrical
gradient (\texorpdfstring{\ensuremath{\Delta}}{Delta}\texorpdfstring{\ensuremath{\Psi}}{Psi}); \texorpdfstring{\ensuremath{\Delta}}{Delta}p (mV) = \texorpdfstring{\ensuremath{\Delta}}{Delta}\texorpdfstring{\ensuremath{\Psi}}{Psi} \texorpdfstring{\ensuremath{-}}{-} 59 \texorpdfstring{\ensuremath{\Delta}}{Delta}pH; drives ATP synthase rotation. \texorpdfstring{\ensuremath{\to}}{->}
\cref{sec:unit_III_bioenergetics_and_respiration}

\subsubsection*{R (continued)}\label{r-continued}
\addcontentsline{toc}{subsubsection}{R (continued)}

\protect\phantomsection\label{gl:reflex-arc}{\textbf{Reflex arc}} {[}L
\emph{reflectere} = to bend back{]} --- The neural pathway of a reflex
response: stimulus \texorpdfstring{\ensuremath{\to}}{->} sensory neuron \texorpdfstring{\ensuremath{\to}}{->} integration centre (interneuron in
spinal cord or brain) \texorpdfstring{\ensuremath{\to}}{->} motor neuron \texorpdfstring{\ensuremath{\to}}{->} effector; characterised by
Sherrington (1906). \texorpdfstring{\ensuremath{\to}}{->} \cref{sec:unit_IX_nervous_system}

\protect\phantomsection\label{gl:reproductive-isolation}{\textbf{Reproductive
isolation}} {[}L \emph{reproducere} = to produce again + \emph{isolatus}
= set apart{]} --- The existence of biological barriers (prezygotic or
postzygotic) preventing interbreeding between populations; the
operational criterion for separate species status in the Biological
Species Concept (Mayr). \texorpdfstring{\ensuremath{\to}}{->}
\cref{sec:unit_VI_genetic_drift_and_speciation}

\protect\phantomsection\label{gl:r-strategist}{\textbf{r-strategist}}
{[}see \emph{r}, rate of increase + Gk \emph{strategia}{]} --- An
organism with a life history strategy maximising reproductive rate (r):
small body, rapid maturation, large clutch sizes, little parental
investment; thrives in unstable, unpredictable environments. Contrast:
\textbf{K-strategist}. \texorpdfstring{\ensuremath{\to}}{->} \cref{sec:unit_X_population_ecology}

\protect\phantomsection\label{gl:resting-potential}{\textbf{Resting
potential}} {[}ME \emph{resten} + L \emph{potentialis}{]} --- The stable
membrane potential of a non-stimulated neuron (about \texorpdfstring{\ensuremath{-}}{-}70 mV
inside relative to outside); maintained by the sodium-potassium pump (3
Na\texorpdfstring{\textsuperscript{+}}{+} out, 2 K\texorpdfstring{\textsuperscript{+}}{+} in per ATP) and selective membrane permeability. \texorpdfstring{\ensuremath{\to}}{->}
\cref{sec:unit_IX_nervous_system}

\subsubsection*{S (continued)}\label{s-continued}
\addcontentsline{toc}{subsubsection}{S (continued)}

\protect\phantomsection\label{gl:saprophyte}{\textbf{Saprophyte}} {[}Gk
\emph{sapros} = rotten + \emph{phyton} = plant{]} --- An organism that
obtains nutrients by absorbing dissolved organic compounds from dead or
decaying matter; now more precisely called saprotrophic or decomposer;
includes most fungi and many bacteria. \texorpdfstring{\ensuremath{\to}}{->}
\cref{sec:unit_X_ecosystem_ecology}

\protect\phantomsection\label{gl:secondary-succession}{\textbf{Secondary
succession}} {[}L \emph{secundus} = following + \emph{successio}{]} ---
Ecological succession in an area where a previous community has been
disrupted (fire, flood, deforestation) but soil and seed bank remain;
faster than primary succession; classic example: forest regeneration
after clear-cutting. \texorpdfstring{\ensuremath{\to}}{->} \cref{sec:unit_X_community_ecology}

\protect\phantomsection\label{gl:self-organization}{\textbf{Self-organization}}
{[}OE \emph{self} + Gk \emph{organon} = tool/instrument{]} --- The
emergence of ordered pattern from local interactions without central
control; examples include biofilms, morphogen gradients, flocking, and
vegetation patching in drylands. \texorpdfstring{\ensuremath{\to}}{->}
\cref{sec:unit_I_water_and_life,sec:unit_X_population_ecology}

\protect\phantomsection\label{gl:shannon-diversity-index}{\textbf{Shannon
diversity index (H')}} {[}after Claude Shannon's information entropy,
1948{]} --- A measure of species diversity: H' = \texorpdfstring{\ensuremath{-}}{-}\texorpdfstring{\ensuremath{\Sigma}}{Sigma}p\texorpdfstring{\textsubscript{i}}{i} ln p\texorpdfstring{\textsubscript{i}}{i} where p\texorpdfstring{\textsubscript{i}}{i} is
the relative abundance of species i; higher H' = greater diversity; H' =
0 for a monoculture; simultaneously measures species richness and
evenness. \texorpdfstring{\ensuremath{\to}}{->} \cref{sec:unit_X_community_ecology}

\protect\phantomsection\label{gl:signal-transduction}{\textbf{Signal
transduction}} {[}L \emph{signum} = mark + \emph{transducere} = to lead
across{]} --- The process by which a receptor converts an extracellular
signal (ligand, light, mechanical force) into an intracellular
biochemical response via cascades of protein-protein interactions and
phosphorylation events. \texorpdfstring{\ensuremath{\to}}{->} \cref{sec:unit_II_cell_signaling}

\protect\phantomsection\label{gl:snp}{\textbf{SNP (single-nucleotide
polymorphism)}} {[}genetics abbreviation{]} --- A common single-base
genomic variant, usually defined by minor allele frequency of at least
1\% in a population; used as a marker in GWAS, ancestry inference, and
risk prediction. \texorpdfstring{\ensuremath{\to}}{->}
\cref{sec:unit_IV_mutations_and_genomics,sec:unit_V_mendelian_genetics}

\protect\phantomsection\label{gl:stomata}{\textbf{Stomata}} (sing.
stoma) {[}Gk \emph{stoma} = mouth{]} --- Microscopic pores in the leaf
epidermis flanked by guard cells; regulate CO\texorpdfstring{\textsubscript{2}}{2} entry for photosynthesis
and water vapour exit (transpiration); respond to light, CO\texorpdfstring{\textsubscript{2}}{2}
concentration, and ABA. \texorpdfstring{\ensuremath{\to}}{->} \cref{sec:unit_VIII_plant_structure_and_water}

\protect\phantomsection\label{gl:succession}{\textbf{Succession}} {[}L
\emph{successio} = a following{]} --- The directional, predictable
change in species composition of a community over time from initial
colonisation to a relatively stable climax community; driven by
facilitation, inhibition, or tolerance mechanisms. \texorpdfstring{\ensuremath{\to}}{->}
\cref{sec:unit_X_community_ecology}

\protect\phantomsection\label{gl:survivorship-curve}{\textbf{Survivorship
curve}} {[}L \emph{super} = beyond + \emph{vivere} = to live{]} --- A
graph of the proportion of a cohort surviving to each age; Type I (most
survive to old age; humans, large mammals); Type II (constant mortality
rate; many birds); Type III (high early mortality; marine invertebrates,
most plants). \texorpdfstring{\ensuremath{\to}}{->} \cref{sec:unit_X_population_ecology}

\protect\phantomsection\label{gl:sympatric-speciation}{\textbf{Sympatric
speciation}} {[}Gk \emph{sym} = together + L \emph{patria} = homeland{]}
--- Speciation occurring within a single geographic area without
physical barriers; common in plants via allopolyploidy; also documented
in cichlids and apple maggot flies. \texorpdfstring{\ensuremath{\to}}{->}
\cref{sec:unit_VI_genetic_drift_and_speciation}

\subsubsection*{T (continued)}\label{t-continued}
\addcontentsline{toc}{subsubsection}{T (continued)}

\protect\phantomsection\label{gl:t-helper-cell}{\textbf{T helper cell
(Th)}} {[}L \emph{thymus} + ME \emph{helpen}{]} --- A CD4+ T lymphocyte
that coordinates adaptive immunity; Th1 cells activate cytotoxic T cells
and macrophages (intracellular pathogens); Th2 cells activate B cells
(antibody production, extracellular pathogens); regulatory T cells
(Treg) suppress autoimmunity. \texorpdfstring{\ensuremath{\to}}{->} \cref{sec:unit_IX_endocrine_and_immune}

\protect\phantomsection\label{gl:telomere}{\textbf{Telomere}} {[}Gk
\emph{telos} = end + \emph{meros} = part{]} --- Repetitive non-coding
DNA sequences (TTAGGG)\texorpdfstring{\textsubscript{n}}{n} capping eukaryotic chromosome ends; protect
against exonuclease degradation, end-to-end chromosome fusions, and loss
of information-carrying DNA; shortened with each S phase
(end-replication problem) unless telomerase is expressed. Nobel Prize
2009 (Blackburn, Greider, Szostak). \texorpdfstring{\ensuremath{\to}}{->}
\cref{sec:unit_IV_dna_replication_and_cell_cycle}

\protect\phantomsection\label{gl:thermoregulation}{\textbf{Thermoregulation}}
{[}Gk \emph{therme} = heat + L \emph{regulare} = to rule{]} ---
Maintenance of body temperature within a physiological range;
\textbf{endotherms} (birds, mammals) generate heat metabolically;
\textbf{ectotherms} (reptiles, fish) rely on behavioural and
environmental sources; homeostatic mechanism with hypothalamic
`thermostat'. \texorpdfstring{\ensuremath{\to}}{->} \cref{sec:unit_IX_nervous_system}

\protect\phantomsection\label{gl:toll-like-receptor}{\textbf{Toll-like
receptor (TLR)}} {[}\emph{Toll} = German interjection of amazement; Toll
pathway discovered in \emph{Drosophila} by Christiane
Nüsslein-Volhard{]} --- A family of pattern-recognition receptors
(TLR1--10 in humans, with additional functional TLRs in other mammals)
expressed on innate immune cells; recognise conserved PAMPs and signal
via MyD88 and TRIF adaptors to activate NF-\texorpdfstring{\ensuremath{\kappa}}{kappa}B and IRF3. Dysregulation
linked to sepsis and autoimmune disease. Nobel Prize 2011 (Hoffmann,
Beutler). \texorpdfstring{\ensuremath{\to}}{->} \cref{sec:unit_IX_endocrine_and_immune}

\protect\phantomsection\label{gl:totipotency}{\textbf{Totipotency}} {[}L
\emph{totus} = whole + \emph{potens} = powerful{]} --- The ability of a
single cell to develop into a complete organism; characterises
fertilised eggs and early-cleavage blastomeres; plant cells frequently
retain totipotency (basis of tissue culture propagation). \texorpdfstring{\ensuremath{\to}}{->}
\cref{sec:unit_VIII_plant_structure_and_water}

\protect\phantomsection\label{gl:transpiration-stream}{\textbf{Transpiration
stream}} {[}L \emph{trans} + \emph{spirare} = breathe{]} --- The
continuous movement of water from soil through roots, stem, and leaves
to the atmosphere; driven by evaporation from mesophyll cell walls;
responsible for transport of dissolved minerals throughout the plant. \texorpdfstring{\ensuremath{\to}}{->}
\cref{sec:unit_VIII_plant_structure_and_water}

\protect\phantomsection\label{gl:trophic-cascade}{\textbf{Trophic
cascade}} {[}Gk \emph{trophe} = nourishment + L \emph{cascata} =
waterfall{]} --- Indirect predator effects on non-prey species,
transmitted down the food chain; top-down cascade: predator suppresses
herbivore \texorpdfstring{\ensuremath{\to}}{->} vegetation recovers; bottom-up cascade: primary productivity
determines higher trophic levels. Documented in Yellowstone (wolves \texorpdfstring{\ensuremath{\to}}{->}
elk \texorpdfstring{\ensuremath{\to}}{->} willows \texorpdfstring{\ensuremath{\to}}{->} beaver \texorpdfstring{\ensuremath{\to}}{->} hydrology). \texorpdfstring{\ensuremath{\to}}{->}
\cref{sec:unit_X_community_ecology}

\protect\phantomsection\label{gl:turgor-pressure}{\textbf{Turgor
pressure}} {[}L \emph{turgere} = to swell{]} --- The pressure exerted by
the contents of a plant cell against its cell wall when the vacuole is
filled with water; drives cell expansion during growth and maintains
tissue rigidity (wilting = loss of turgor). \texorpdfstring{\ensuremath{\to}}{->}
\cref{sec:unit_VIII_plant_structure_and_water}

\subsubsection*{V (continued)}\label{v-continued}
\addcontentsline{toc}{subsubsection}{V (continued)}

\protect\phantomsection\label{gl:vaccine}{\textbf{Vaccine}} {[}L
\emph{vacca} = cow; from Jenner's cowpox inoculation against smallpox,
1796{]} --- A biological preparation conferring immunity against a
specific pathogen; types include live-attenuated (MMR), inactivated
(IPV), subunit (HepB), mRNA (SARS-CoV-2 mRNA-1273, BNT162b2), and viral
vector vaccines. \texorpdfstring{\ensuremath{\to}}{->} \cref{sec:unit_VII_bacteria_archaea_viruses}

\protect\phantomsection\label{gl:viviparity}{\textbf{Viviparity}} {[}L
\emph{vivus} = alive + \emph{parere} = to beget{]} --- Giving birth to
live young rather than laying eggs; characteristic of most mammals and
some fish, reptiles, and invertebrates; requires placentation (in
eutherian mammals) for gas and nutrient exchange. \texorpdfstring{\ensuremath{\to}}{->}
\cref{sec:unit_IX_nervous_system}

\subsubsection*{W (continued)}\label{w-continued}
\addcontentsline{toc}{subsubsection}{W (continued)}

\protect\phantomsection\label{gl:watershed}{\textbf{Watershed}} {[}OE
\emph{wæter} + OE \emph{scēad} = divide{]} --- The land area draining
into a common body of water (stream, lake, or estuary); the ecological
unit for studying nutrient export, hydrology, and land-use impacts on
aquatic ecosystems; Hubbard Brook watershed experiments (Bormann \&
Likens, 1967) pioneered biogeochemical cycle measurement at watershed
scale. \texorpdfstring{\ensuremath{\to}}{->} \cref{sec:unit_X_ecosystem_ecology}

\begin{center}\rule{0.5\linewidth}{0.5pt}\end{center}

\emph{Expanded glossary now covers about 200 terms across all
10 units. Etymologies sourced from Oxford Latin Dictionary,
Liddell-Scott-Jones Greek-English Lexicon, and Online Etymology
Dictionary (etymonline.com).}

\newpage

\section*{Appendix G --- Index of Key
Terms}\label{appendix-g-index-of-key-terms}
\addcontentsline{toc}{section}{Appendix G --- Index of Key Terms}

\label{sec:appendix_index}

\begin{quote}
\textbf{Appendix G} \texorpdfstring{\ensuremath{\cdot}}{.} Level 1/3 \texorpdfstring{\ensuremath{\cdot}}{.} Use as reference
\end{quote}

An index of the 139 glossary terms, each with the chapter(s) where the
term is introduced, bolded, or discussed at length. Terms marked in
\textbf{bold} appear as section headings somewhere in the textbook; the
glossary entry carries the full definition.

Each entry follows the pattern:

\begin{verbatim}
**Term** (`#gl:slug`) with semantic chapter links
\end{verbatim}

The \texttt{\#gl:slug} anchor is the canonical location defined in the
master glossary (\breaktt{manuscript/glossary.md}). Full-text search for
the slug string finds every cross-reference.

\begin{quote}
\textbf{Note on maintenance.} This index is maintained by
\breaktt{scripts/link\_glossary.py} from the glossary file and canonical
ToC labels. To verify after editing \texttt{glossary.md}, run
\breaktt{uv\ run\ python\ scripts/link\_glossary.py\ -\/-check}.
\end{quote}

\begin{center}\rule{0.5\linewidth}{0.5pt}\end{center}

\subsection*{A}\label{a-1}
\addcontentsline{toc}{subsection}{A}

\textbf{Abiotic} (\texttt{\#gl:abiotic}) \texorpdfstring{\ensuremath{\to}}{->}
\cref{sec:unit_X_biomes_and_conservation} \textbf{Actin}
(\texttt{\#gl:actin}) \texorpdfstring{\ensuremath{\to}}{->} \cref{sec:unit_II_cell_structure} \textbf{Action
potential} (\texttt{\#gl:action-potential}) \texorpdfstring{\ensuremath{\to}}{->}
\cref{sec:unit_IX_nervous_system,sec:unit_IX_action_potential_synapses}
\textbf{Adenosine triphosphate} (\texttt{\#gl:adenosine-triphosphate}) \texorpdfstring{\ensuremath{\to}}{->}
\cref{sec:unit_III_bioenergetics_and_respiration} \textbf{Aerobic}
(\texttt{\#gl:aerobic}) \texorpdfstring{\ensuremath{\to}}{->}
\cref{sec:unit_III_bioenergetics_and_respiration} \textbf{Allele}
(\texttt{\#gl:allele}) \texorpdfstring{\ensuremath{\to}}{->}
\cref{sec:unit_V_mendelian_genetics,sec:unit_V_population_genetics}
\textbf{Allosteric} (\texttt{\#gl:allosteric}) \texorpdfstring{\ensuremath{\to}}{->}
\cref{sec:unit_I_enzymes_and_kinetics} \textbf{Anaerobic}
(\texttt{\#gl:anaerobic}) \texorpdfstring{\ensuremath{\to}}{->}
\cref{sec:unit_III_bioenergetics_and_respiration} \textbf{Anaphase}
(\texttt{\#gl:anaphase}) \texorpdfstring{\ensuremath{\to}}{->}
\cref{sec:unit_IV_dna_replication_and_cell_cycle} \textbf{Angiosperm}
(\texttt{\#gl:angiosperm}) \texorpdfstring{\ensuremath{\to}}{->}
\cref{sec:unit_VIII_plant_structure_and_water} \textbf{Anticodon}
(\texttt{\#gl:anticodon}) \texorpdfstring{\ensuremath{\to}}{->} \cref{sec:unit_IV_gene_expression}
\textbf{Apoptosis} (\texttt{\#gl:apoptosis}) \texorpdfstring{\ensuremath{\to}}{->}
\cref{sec:unit_II_cell_signaling} \textbf{Archaea}
(\texttt{\#gl:archaea}) \texorpdfstring{\ensuremath{\to}}{->}
\cref{sec:unit_II_cell_theory,sec:unit_VI_phylogenetics} \textbf{ATP
synthase} (\texttt{\#gl:atp-synthase}) \texorpdfstring{\ensuremath{\to}}{->}
\cref{sec:unit_III_bioenergetics_and_respiration} \textbf{Autotroph}
(\texttt{\#gl:autotroph}) \texorpdfstring{\ensuremath{\to}}{->}
\cref{sec:unit_II_cell_theory,sec:unit_X_community_ecology} \textbf{Axon
hillock} (\texttt{\#gl:axon-hillock}) \texorpdfstring{\ensuremath{\to}}{->}
\cref{sec:unit_IX_nervous_system}

\subsection*{B}\label{b-1}
\addcontentsline{toc}{subsection}{B}

\textbf{Base pair} (\texttt{\#gl:base-pair}) \texorpdfstring{\ensuremath{\to}}{->}
\cref{sec:unit_IV_dna_replication_and_cell_cycle} \textbf{Biofilm}
(\texttt{\#gl:biofilm}) \texorpdfstring{\ensuremath{\to}}{->} \cref{sec:unit_VII_bacteria_archaea_viruses}
\textbf{Biogeochemical cycle} \texorpdfstring{\ensuremath{\to}}{->} \cref{sec:unit_X_ecosystem_ecology}
\textbf{Biome} (\texttt{\#gl:biome}) \texorpdfstring{\ensuremath{\to}}{->}
\cref{sec:unit_X_biomes_and_conservation} \textbf{Bottleneck effect}
(\texttt{\#gl:genetic-drift}) \texorpdfstring{\ensuremath{\to}}{->}
\cref{sec:unit_VI_genetic_drift_and_speciation} \textbf{Buffer}
(\texttt{\#gl:buffer}) \texorpdfstring{\ensuremath{\to}}{->} \cref{sec:unit_I_atoms_molecules}

\subsection*{C}\label{c-1}
\addcontentsline{toc}{subsection}{C}

\textbf{Calvin cycle} (\texttt{\#gl:calvin-cycle}) \texorpdfstring{\ensuremath{\to}}{->}
\cref{sec:unit_III_photosynthesis} \textbf{Carrying capacity}
(\texttt{\#gl:carrying-capacity}) \texorpdfstring{\ensuremath{\to}}{->} \cref{sec:unit_X_population_ecology}
\textbf{Caspase} (\texttt{\#gl:caspase}) \texorpdfstring{\ensuremath{\to}}{->}
\cref{sec:unit_II_cell_signaling} \textbf{Cell cycle}
(\texttt{\#gl:cell-cycle}) \texorpdfstring{\ensuremath{\to}}{->}
\cref{sec:unit_IV_dna_replication_and_cell_cycle} \textbf{Chemiosmosis}
(\texttt{\#gl:chemiosmosis}) \texorpdfstring{\ensuremath{\to}}{->}
\cref{sec:unit_III_bioenergetics_and_respiration} \textbf{Chirality}
(\texttt{\#gl:chirality}) \texorpdfstring{\ensuremath{\to}}{->} \cref{sec:unit_I_atoms_molecules}
\textbf{Chlorophyll} (\texttt{\#gl:chlorophyll}) \texorpdfstring{\ensuremath{\to}}{->}
\cref{sec:unit_III_photosynthesis} \textbf{Chloroplast}
(\texttt{\#gl:chloroplast}) \texorpdfstring{\ensuremath{\to}}{->}
\cref{sec:unit_II_cell_theory,sec:unit_III_photosynthesis}
\textbf{Chromatin} (\texttt{\#gl:chromatin}) \texorpdfstring{\ensuremath{\to}}{->}
\cref{sec:unit_IV_epigenetics_and_gene_regulation} \textbf{Chromosome}
(\texttt{\#gl:chromosome}) \texorpdfstring{\ensuremath{\to}}{->}
\cref{sec:unit_IV_dna_replication_and_cell_cycle,sec:unit_V_chromosomal_inheritance}
\textbf{Clade} \texorpdfstring{\ensuremath{\to}}{->} \cref{sec:unit_VI_phylogenetics} \textbf{Codon}
(\texttt{\#gl:codon}) \texorpdfstring{\ensuremath{\to}}{->} \cref{sec:unit_IV_gene_expression}
\textbf{Competitive exclusion} (\texttt{\#gl:competitive-exclusion}) \texorpdfstring{\ensuremath{\to}}{->}
\cref{sec:unit_X_community_ecology} \textbf{Concentration gradient}
(\texttt{\#gl:facilitated-diffusion}) \texorpdfstring{\ensuremath{\to}}{->}
\cref{sec:unit_II_membrane_transport} \textbf{Cortisol}
(\texttt{\#gl:cortisol}) \texorpdfstring{\ensuremath{\to}}{->} \cref{sec:unit_IX_endocrine_and_immune}
\textbf{CpG island} (\texttt{\#gl:cpg-island}) \texorpdfstring{\ensuremath{\to}}{->}
\cref{sec:unit_IV_epigenetics_and_gene_regulation} \textbf{Cytoskeleton}
(\texttt{\#gl:cytoskeleton}) \texorpdfstring{\ensuremath{\to}}{->} \cref{sec:unit_II_cell_structure}

\subsection*{D}\label{d-1}
\addcontentsline{toc}{subsection}{D}

\textbf{Depolarisation} (\texttt{\#gl:depolarisation}) \texorpdfstring{\ensuremath{\to}}{->}
\cref{sec:unit_IX_nervous_system,sec:unit_IX_action_potential_synapses}
\textbf{Dominant} (\texttt{\#gl:dominant}) \texorpdfstring{\ensuremath{\to}}{->}
\cref{sec:unit_V_mendelian_genetics}

\subsection*{E}\label{e-1}
\addcontentsline{toc}{subsection}{E}

\textbf{Ecosystem} \texorpdfstring{\ensuremath{\to}}{->} \cref{sec:unit_X_ecosystem_ecology}
\textbf{Electronegativity} (\texttt{\#gl:electronegativity}) \texorpdfstring{\ensuremath{\to}}{->}
\cref{sec:unit_I_atoms_molecules} \textbf{Emergence}
(\texttt{\#gl:self-organization}) \texorpdfstring{\ensuremath{\to}}{->} \cref{sec:unit_0_systems_science}
\textbf{Endocytosis} (\texttt{\#gl:endocytosis}) \texorpdfstring{\ensuremath{\to}}{->}
\cref{sec:unit_II_membrane_transport} \textbf{Endoplasmic reticulum}
(\texttt{\#gl:organelle}) \texorpdfstring{\ensuremath{\to}}{->} \cref{sec:unit_II_cell_structure}
\textbf{Endosymbiotic theory} (\texttt{\#gl:endosymbiosis}) \texorpdfstring{\ensuremath{\to}}{->}
\cref{sec:unit_II_cell_theory} \textbf{Enzyme} (\texttt{\#gl:enzyme}) \texorpdfstring{\ensuremath{\to}}{->}
\cref{sec:unit_I_enzymes_and_kinetics} \textbf{Epigenetics}
(\texttt{\#gl:epigenetics}) \texorpdfstring{\ensuremath{\to}}{->}
\cref{sec:unit_IV_epigenetics_and_gene_regulation} \textbf{Epithelial
tissue} \texorpdfstring{\ensuremath{\to}}{->} \cref{sec:unit_II_cell_structure} \textbf{Eukaryote}
(\texttt{\#gl:eukaryote}) \texorpdfstring{\ensuremath{\to}}{->} \cref{sec:unit_II_cell_theory}
\textbf{Evolution} (\texttt{\#gl:natural-selection}) \texorpdfstring{\ensuremath{\to}}{->}
\cref{sec:unit_VI_evolution_and_selection} \textbf{Exocytosis}
(\texttt{\#gl:endocytosis}) \texorpdfstring{\ensuremath{\to}}{->} \cref{sec:unit_II_membrane_transport}

\subsection*{F}\label{f-1}
\addcontentsline{toc}{subsection}{F}

\textbf{Feedback loop} (\texttt{\#gl:homeostasis}) \texorpdfstring{\ensuremath{\to}}{->}
\cref{sec:unit_0_systems_science} \textbf{Fermentation}
(\texttt{\#gl:fermentation}) \texorpdfstring{\ensuremath{\to}}{->}
\cref{sec:unit_III_bioenergetics_and_respiration} \textbf{Fertilisation}
(\texttt{\#gl:zygote}) \texorpdfstring{\ensuremath{\to}}{->}
\cref{sec:unit_V_chromosomal_inheritance,sec:unit_VIII_plant_reproduction}
\textbf{Fitness} (\texttt{\#gl:fitness}) \texorpdfstring{\ensuremath{\to}}{->}
\cref{sec:unit_VI_evolution_and_selection} \textbf{Founder effect}
(\texttt{\#gl:founder-effect}) \texorpdfstring{\ensuremath{\to}}{->}
\cref{sec:unit_VI_genetic_drift_and_speciation}

\subsection*{G}\label{g-1}
\addcontentsline{toc}{subsection}{G}

\textbf{Gamete} (\texttt{\#gl:gamete}) \texorpdfstring{\ensuremath{\to}}{->}
\cref{sec:unit_V_chromosomal_inheritance} \textbf{Gene}
(\texttt{\#gl:gene}) \texorpdfstring{\ensuremath{\to}}{->} \cref{sec:unit_IV_gene_expression} \textbf{Gene
flow} (\texttt{\#gl:genetic-drift}) \texorpdfstring{\ensuremath{\to}}{->}
\cref{sec:unit_VI_genetic_drift_and_speciation} \textbf{Genetic drift}
(\texttt{\#gl:genetic-drift}) \texorpdfstring{\ensuremath{\to}}{->}
\cref{sec:unit_V_population_genetics,sec:unit_VI_genetic_drift_and_speciation}
\textbf{Genome} (\texttt{\#gl:genome}) \texorpdfstring{\ensuremath{\to}}{->}
\cref{sec:unit_IV_mutations_and_genomics} \textbf{Genotype}
(\texttt{\#gl:phenotype}) \texorpdfstring{\ensuremath{\to}}{->} \cref{sec:unit_V_mendelian_genetics}
\textbf{Glucose} (\texttt{\#gl:glycolysis}) \texorpdfstring{\ensuremath{\to}}{->}
\cref{sec:unit_I_macromolecules,sec:unit_III_bioenergetics_and_respiration}
\textbf{GLUT4} (\texttt{\#gl:glut4}) \texorpdfstring{\ensuremath{\to}}{->}
\cref{sec:unit_II_membrane_transport,sec:unit_III_metabolic_integration,sec:unit_IX_endocrine_and_immune}
\textbf{Glycolysis} (\texttt{\#gl:glycolysis}) \texorpdfstring{\ensuremath{\to}}{->}
\cref{sec:unit_III_bioenergetics_and_respiration} \textbf{Golgi
apparatus} (\texttt{\#gl:organelle}) \texorpdfstring{\ensuremath{\to}}{->} \cref{sec:unit_II_cell_structure}

\subsection*{H}\label{h-1}
\addcontentsline{toc}{subsection}{H}

\textbf{Hardy-Weinberg equilibrium}
(\texttt{\#gl:hardy-weinberg-equilibrium}) \texorpdfstring{\ensuremath{\to}}{->}
\cref{sec:unit_V_population_genetics} \textbf{Herd immunity}
(\texttt{\#gl:herd-immunity}) \texorpdfstring{\ensuremath{\to}}{->}
\cref{sec:unit_VII_infectious_disease,sec:unit_IX_endocrine_and_immune}
\textbf{Heterochromatin} (\texttt{\#gl:heterochromatin}) \texorpdfstring{\ensuremath{\to}}{->}
\cref{sec:unit_IV_epigenetics_and_gene_regulation} \textbf{Heterozygous}
(\texttt{\#gl:heterozygous}) \texorpdfstring{\ensuremath{\to}}{->} \cref{sec:unit_V_mendelian_genetics}
\textbf{Histone} (\texttt{\#gl:histone}) \texorpdfstring{\ensuremath{\to}}{->}
\cref{sec:unit_IV_epigenetics_and_gene_regulation} \textbf{Histone code}
(\texttt{\#gl:histone-code}) \texorpdfstring{\ensuremath{\to}}{->}
\cref{sec:unit_IV_epigenetics_and_gene_regulation} \textbf{Homeostasis}
(\texttt{\#gl:homeostasis}) \texorpdfstring{\ensuremath{\to}}{->}
\cref{sec:unit_0_systems_science,sec:unit_IX_circulation_respiration_homeostasis}
\textbf{Hormone} (\texttt{\#gl:hormone}) \texorpdfstring{\ensuremath{\to}}{->}
\cref{sec:unit_IX_endocrine_and_immune} \textbf{Hydrogen bond}
(\texttt{\#gl:hydrogen-bond}) \texorpdfstring{\ensuremath{\to}}{->}
\cref{sec:unit_I_atoms_molecules,sec:unit_I_water_and_life}

\subsection*{I}\label{i-1}
\addcontentsline{toc}{subsection}{I}

\textbf{Insulin receptor} (\texttt{\#gl:insulin-receptor}) \texorpdfstring{\ensuremath{\to}}{->}
\cref{sec:unit_II_cell_signaling,sec:unit_IX_endocrine_and_immune}
\textbf{Island biogeography} (\texttt{\#gl:island-biogeography}) \texorpdfstring{\ensuremath{\to}}{->}
\cref{sec:unit_X_community_ecology}

\subsection*{L}\label{l-1}
\addcontentsline{toc}{subsection}{L}

\textbf{Life table} (\texttt{\#gl:life-table}) \texorpdfstring{\ensuremath{\to}}{->}
\cref{sec:unit_X_population_ecology} \textbf{Logistic growth}
(\texttt{\#gl:logistic-growth}) \texorpdfstring{\ensuremath{\to}}{->} \cref{sec:unit_X_population_ecology}
\textbf{Lysosome} (\texttt{\#gl:organelle}) \texorpdfstring{\ensuremath{\to}}{->}
\cref{sec:unit_II_cell_structure}

\subsection*{M}\label{m-1}
\addcontentsline{toc}{subsection}{M}

\textbf{Meiosis} (\texttt{\#gl:meiosis}) \texorpdfstring{\ensuremath{\to}}{->}
\cref{sec:unit_V_chromosomal_inheritance} \textbf{Metaphase}
(\texttt{\#gl:anaphase}) \texorpdfstring{\ensuremath{\to}}{->}
\cref{sec:unit_IV_dna_replication_and_cell_cycle,sec:unit_V_chromosomal_inheritance}
\textbf{Microtubule} (\texttt{\#gl:cytoskeleton}) \texorpdfstring{\ensuremath{\to}}{->}
\cref{sec:unit_II_cell_structure} \textbf{Mitochondrion}
(\texttt{\#gl:organelle}) \texorpdfstring{\ensuremath{\to}}{->}
\cref{sec:unit_II_cell_structure,sec:unit_III_bioenergetics_and_respiration}
\textbf{Mitosis} (\texttt{\#gl:mitosis}) \texorpdfstring{\ensuremath{\to}}{->}
\cref{sec:unit_IV_dna_replication_and_cell_cycle} \textbf{Molecular
clock} (\texttt{\#gl:molecular-clock}) \texorpdfstring{\ensuremath{\to}}{->}
\cref{sec:unit_V_population_genetics,sec:unit_VI_phylogenetics}
\textbf{Mutation} (\texttt{\#gl:mutation}) \texorpdfstring{\ensuremath{\to}}{->}
\cref{sec:unit_IV_mutations_and_genomics,sec:unit_V_population_genetics}
\textbf{Mutualism} (\texttt{\#gl:mutualism}) \texorpdfstring{\ensuremath{\to}}{->}
\cref{sec:unit_X_community_ecology}

\subsection*{N}\label{n-1}
\addcontentsline{toc}{subsection}{N}

\textbf{Natural selection} (\texttt{\#gl:natural-selection}) \texorpdfstring{\ensuremath{\to}}{->}
\cref{sec:unit_VI_evolution_and_selection} \textbf{Neuron}
(\texttt{\#gl:neuron}) \texorpdfstring{\ensuremath{\to}}{->} \cref{sec:unit_IX_nervous_system}
\textbf{Neutral theory} (\texttt{\#gl:neutral-theory}) \texorpdfstring{\ensuremath{\to}}{->}
\cref{sec:unit_V_population_genetics} \textbf{Niche}
(\texttt{\#gl:niche}) \texorpdfstring{\ensuremath{\to}}{->} \cref{sec:unit_X_community_ecology}
\textbf{Nucleosome} (\texttt{\#gl:nucleosome}) \texorpdfstring{\ensuremath{\to}}{->}
\cref{sec:unit_IV_epigenetics_and_gene_regulation}

\subsection*{O}\label{o-1}
\addcontentsline{toc}{subsection}{O}

\textbf{Operon} (\texttt{\#gl:operon}) \texorpdfstring{\ensuremath{\to}}{->}
\cref{sec:unit_IV_gene_expression} \textbf{Organelle}
(\texttt{\#gl:organelle}) \texorpdfstring{\ensuremath{\to}}{->} \cref{sec:unit_II_cell_structure}

\subsection*{P}\label{p-1}
\addcontentsline{toc}{subsection}{P}

\textbf{Parasite} (\texttt{\#gl:parasite}) \texorpdfstring{\ensuremath{\to}}{->}
\cref{sec:unit_X_community_ecology} \textbf{Peptide bond}
(\texttt{\#gl:peptide-bond}) \texorpdfstring{\ensuremath{\to}}{->} \cref{sec:unit_I_macromolecules}
\textbf{pH} (\texttt{\#gl:ph}) \texorpdfstring{\ensuremath{\to}}{->} \cref{sec:unit_I_atoms_molecules}
\textbf{Phagocytosis} (\texttt{\#gl:endocytosis}) \texorpdfstring{\ensuremath{\to}}{->}
\cref{sec:unit_II_membrane_transport} \textbf{Phenotype}
(\texttt{\#gl:phenotype}) \texorpdfstring{\ensuremath{\to}}{->} \cref{sec:unit_V_mendelian_genetics}
\textbf{Phospholipid bilayer} (\texttt{\#gl:phospholipid-bilayer}) \texorpdfstring{\ensuremath{\to}}{->}
\cref{sec:unit_II_membrane_transport} \textbf{Photosystem}
(\texttt{\#gl:photosynthesis}) \texorpdfstring{\ensuremath{\to}}{->} \cref{sec:unit_III_photosynthesis}
\textbf{Phylogeny} (\texttt{\#gl:molecular-clock}) \texorpdfstring{\ensuremath{\to}}{->}
\cref{sec:unit_VI_phylogenetics} \textbf{Phytoplankton}
(\texttt{\#gl:phytoplankton}) \texorpdfstring{\ensuremath{\to}}{->}
\cref{sec:unit_X_community_ecology,sec:unit_X_ecosystem_ecology}
\textbf{Plasmid} (\texttt{\#gl:plasmid}) \texorpdfstring{\ensuremath{\to}}{->}
\cref{sec:unit_IV_mutations_and_genomics} \textbf{Ploidy}
(\texttt{\#gl:diploid}) \texorpdfstring{\ensuremath{\to}}{->} \cref{sec:unit_V_chromosomal_inheritance}
\textbf{Polymerase} (\texttt{\#gl:transcription}) \texorpdfstring{\ensuremath{\to}}{->}
\cref{sec:unit_IV_dna_replication_and_cell_cycle,sec:unit_IV_gene_expression}
\textbf{Population genetics} (\texttt{\#gl:population-genetics}) \texorpdfstring{\ensuremath{\to}}{->}
\cref{sec:unit_V_population_genetics} \textbf{Primase}
(\texttt{\#gl:dna-helicase}) \texorpdfstring{\ensuremath{\to}}{->}
\cref{sec:unit_IV_dna_replication_and_cell_cycle} \textbf{Prokaryote}
(\texttt{\#gl:eukaryote}) \texorpdfstring{\ensuremath{\to}}{->}
\cref{sec:unit_II_cell_theory,sec:unit_VII_bacteria_archaea_viruses}
\textbf{Promoter} (\texttt{\#gl:promoter}) \texorpdfstring{\ensuremath{\to}}{->}
\cref{sec:unit_IV_gene_expression} \textbf{Protein}
(\texttt{\#gl:protein}) \texorpdfstring{\ensuremath{\to}}{->} \cref{sec:unit_I_macromolecules}

\subsection*{R}\label{r-1}
\addcontentsline{toc}{subsection}{R}

\textbf{Recessive} (\texttt{\#gl:dominant}) \texorpdfstring{\ensuremath{\to}}{->}
\cref{sec:unit_V_mendelian_genetics} \textbf{Replication fork}
(\texttt{\#gl:dna-helicase}) \texorpdfstring{\ensuremath{\to}}{->}
\cref{sec:unit_IV_dna_replication_and_cell_cycle} \textbf{Reproductive
isolation} (\texttt{\#gl:reproductive-isolation}) \texorpdfstring{\ensuremath{\to}}{->}
\cref{sec:unit_VI_genetic_drift_and_speciation} \textbf{Resting
potential} (\texttt{\#gl:resting-potential}) \texorpdfstring{\ensuremath{\to}}{->}
\cref{sec:unit_IX_nervous_system} \textbf{Ribosome}
(\texttt{\#gl:ribosome}) \texorpdfstring{\ensuremath{\to}}{->} \cref{sec:unit_IV_gene_expression}

\subsection*{S}\label{s-1}
\addcontentsline{toc}{subsection}{S}

\textbf{Saltatory conduction} (\texttt{\#gl:saltatory-conduction}) \texorpdfstring{\ensuremath{\to}}{->}
\cref{sec:unit_IX_action_potential_synapses} \textbf{Secondary
succession} (\texttt{\#gl:secondary-succession}) \texorpdfstring{\ensuremath{\to}}{->}
\cref{sec:unit_X_community_ecology} \textbf{Species-area relationship}
(\texttt{\#gl:species-area-relationship}) \texorpdfstring{\ensuremath{\to}}{->}
\cref{sec:unit_X_community_ecology,sec:unit_X_biomes_and_conservation}
\textbf{Speciation} (\texttt{\#gl:speciation}) \texorpdfstring{\ensuremath{\to}}{->}
\cref{sec:unit_VI_genetic_drift_and_speciation} \textbf{Stoma}
(\texttt{\#gl:stomata}) \texorpdfstring{\ensuremath{\to}}{->} \cref{sec:unit_VIII_plant_structure_and_water}
\textbf{Succession} (\texttt{\#gl:succession}) \texorpdfstring{\ensuremath{\to}}{->}
\cref{sec:unit_X_community_ecology} \textbf{Survivorship curve}
(\texttt{\#gl:survivorship-curve}) \texorpdfstring{\ensuremath{\to}}{->}
\cref{sec:unit_X_population_ecology} \textbf{Symbiosis}
(\texttt{\#gl:mutualism}) \texorpdfstring{\ensuremath{\to}}{->} \cref{sec:unit_X_community_ecology}
\textbf{Sympatric speciation} (\texttt{\#gl:sympatric-speciation}) \texorpdfstring{\ensuremath{\to}}{->}
\cref{sec:unit_VI_genetic_drift_and_speciation} \textbf{Synapsis}
(\texttt{\#gl:meiosis}) \texorpdfstring{\ensuremath{\to}}{->} \cref{sec:unit_V_chromosomal_inheritance}

\subsection*{T}\label{t-1}
\addcontentsline{toc}{subsection}{T}

\textbf{Transcription} (\texttt{\#gl:transcription}) \texorpdfstring{\ensuremath{\to}}{->}
\cref{sec:unit_IV_gene_expression} \textbf{Translation}
(\texttt{\#gl:translation}) \texorpdfstring{\ensuremath{\to}}{->} \cref{sec:unit_IV_gene_expression}
\textbf{Trophic cascade} (\texttt{\#gl:trophic-cascade}) \texorpdfstring{\ensuremath{\to}}{->}
\cref{sec:unit_X_community_ecology} \textbf{Trophic level}
(\texttt{\#gl:trophic-level}) \texorpdfstring{\ensuremath{\to}}{->}
\cref{sec:unit_X_community_ecology,sec:unit_X_ecosystem_ecology}
\textbf{Tropism} (\texttt{\#gl:photoperiodism}) \texorpdfstring{\ensuremath{\to}}{->}
\cref{sec:unit_VIII_plant_responses}

\subsection*{V}\label{v-1}
\addcontentsline{toc}{subsection}{V}

\textbf{Vaccine} (\texttt{\#gl:vaccine}) \texorpdfstring{\ensuremath{\to}}{->}
\cref{sec:unit_VII_infectious_disease,sec:unit_IX_endocrine_and_immune}
\textbf{Van der Waals forces} (\texttt{\#gl:van-der-waals-forces}) \texorpdfstring{\ensuremath{\to}}{->}
\cref{sec:unit_I_atoms_molecules} \textbf{Virulence factor}
(\texttt{\#gl:virulence}) \texorpdfstring{\ensuremath{\to}}{->} \cref{sec:unit_VII_infectious_disease}
\textbf{Virus} (\texttt{\#gl:vaccine}) \texorpdfstring{\ensuremath{\to}}{->}
\cref{sec:unit_VII_bacteria_archaea_viruses,sec:unit_VII_infectious_disease}

\subsection*{X}\label{x-1}
\addcontentsline{toc}{subsection}{X}

\textbf{X-linked inheritance} (\texttt{\#gl:linkage}) \texorpdfstring{\ensuremath{\to}}{->}
\cref{sec:unit_V_chromosomal_inheritance}

\begin{center}\rule{0.5\linewidth}{0.5pt}\end{center}

\begin{quote}
\textbf{Using the index.} When a term appears as a bolded,
\texttt{\#gl:}-linked first use in a chapter (see Unit I onward),
clicking the term jumps to the canonical entry in the master glossary.
This appendix provides an author- and instructor-facing overview of
which chapter establishes each concept first, useful for lecture
planning and term-by-term syllabus design.
\end{quote}

\emph{Module: reference only (no code)}



\clearpage

\bibliography{references}
\end{document}
